% Generated solely from DER_FR_annotated_master.tex: DERform=0
\documentclass[11pt,oneside]{article}
\usepackage{fontspec}
\usepackage[english,french]{babel}
\usepackage{amsmath}
\usepackage{mathtools}
\usepackage{amssymb}
\usepackage{amsthm}
\usepackage{mathrsfs}
\usepackage{tikz-cd}
\usepackage{enumitem}
\usepackage{array}
\usepackage{ragged2e}
\usepackage{marginnote}
\usepackage{marginfix}
\usepackage{lmodern}
\usetikzlibrary{babel,arrows.meta,decorations.markings,decorations.pathreplacing}
\usepackage[a4paper,left=27mm,right=35mm,top=28mm,bottom=27mm,headsep=9mm,marginparwidth=25mm,marginparsep=5mm]{geometry}
\usepackage{xcolor}
\usepackage{fancyhdr}
\usepackage{xurl}
\usepackage{adjustbox}
% a diagram wider than the line or taller than the page is scaled down to fit (the widest, a chart of chapter II,
% runs to more than twice the width of the line); a diagram that fits is set as it is
\AddToHook{env/tikzcd/before}{\begin{adjustbox}{max width=\linewidth,max totalheight=\textheight}}
\AddToHook{env/tikzcd/after}{\end{adjustbox}}
\usepackage[hyperfootnotes=false,bookmarksopen=true,bookmarksopenlevel=1]{hyperref}
\hypersetup{pdftitle={Les Dérivateurs (transcription française)},pdfauthor={Alexander Grothendieck},
  pdfsubject={Transcription française de l’édition de M. Künzer, J. Malgoire and G. Maltsiniotis},
  colorlinks=true,linkcolor=black,urlcolor=blue!45!black,pdfdisplaydoctitle=true}
\setlength{\emergencystretch}{2em}
\setlength{\headheight}{14pt}
\makeatletter
\newcommand{\derleqno}{\tagsleft@true}
\renewcommand*{\@pnumwidth}{2.4em}
\renewcommand*{\@tocrmarg}{3.4em}
% \mathllap, \mathrlap, \mathclap as mathtools defines them (one unit uses \mathllap)
\providecommand*{\mathllap}[1]{\mathpalette\der@mathllap{#1}}
\providecommand*{\mathrlap}[1]{\mathpalette\der@mathrlap{#1}}
\providecommand*{\mathclap}[1]{\mathpalette\der@mathclap{#1}}
\def\der@mathllap#1#2{\llap{$\m@th#1{#2}$}}
\def\der@mathrlap#1#2{\rlap{$\m@th#1{#2}$}}
\def\der@mathclap#1#2{\hb@xt@\z@{\hss$\m@th#1{#2}$\hss}}
% the margin marks: p. n where page n of the edition begins, [page n] where page n of the manuscript begins
\newcommand{\derpp}[1]{\marginpar{\raggedright\footnotesize\color{black!55}p.~#1}}
\newcommand{\dermsp}[1]{\marginpar{\raggedright\footnotesize\ttfamily[#1]}}
% the editors' note of an edition file, under the chapter heading
\newenvironment{dernote}{\par\begingroup\small\setlength{\parindent}{0pt}\setlength{\parskip}{0.6ex}}%
  {\par\endgroup\vspace{2mm}\noindent\rule{0.3\textwidth}{0.4pt}\par\vspace{5mm}\@afterheading}
% #1 heading line, #2 title, #3 running head, #4 TOC text
\newcommand{\derchapter}[4]{\clearpage\setcounter{footnote}{0}\phantomsection
  \addcontentsline{toc}{section}{#4}\markboth{#3}{#3}\thispagestyle{plain}%
  \vspace*{10mm}{\centering\ifx\relax#1\relax\else{\large\scshape #1\par}\vspace{3mm}\fi
  {\LARGE\bfseries #2\par}}\vspace{9mm}\@afterheading}
% a later edition file of chapter XII
\newcommand{\derpart}[4]{\clearpage\setcounter{footnote}{0}\phantomsection
  \addcontentsline{toc}{subsection}{#4}\markboth{#3}{#3}%
  \vspace*{4mm}{\centering{\large\scshape #1\par}\vspace{2mm}{\Large\bfseries #2\par}}\vspace{7mm}\@afterheading}
\renewcommand*{\l@section}[2]{\addpenalty{-\@highpenalty}\vskip 0.8em plus 1pt
  {\leftskip 0pt\rightskip\@tocrmarg\parfillskip -\rightskip\parindent 0pt\interlinepenalty\@M
   \leavevmode\normalsize #1\nobreak\leaders\hbox{$\m@th\mkern\@dotsep mu\hbox{.}\mkern\@dotsep mu$}\hfill
   \nobreak\hb@xt@\@pnumwidth{\hss #2}\par}}
\renewcommand*{\l@subsection}{\@dottedtocline{2}{1.5em}{0em}}
\makeatother
% two styles, each setting all its fields: the body (running head, page number at the right of the head) and plain
% (page number in the foot: front matter, chapter openings); a style set with \pagestyle{fancy} would keep the fields
% that \pagestyle{plain} cleared
\fancypagestyle{derbody}{\fancyhf{}\fancyhead[L]{\small\itshape\leftmark}\fancyhead[R]{\small\thepage}%
  \renewcommand{\headrulewidth}{0pt}}
\fancypagestyle{plain}{\fancyhf{}\fancyfoot[C]{\small\thepage}\renewcommand{\headrulewidth}{0pt}}
\pagestyle{derbody}

\usepackage{luacode}
\providecommand{\DERform}{0}
\newcommand{\DERmspin}[1]{\marginnote{\raggedright\scriptsize\ttfamily[#1]}}
\newcommand{\DERinlinepage}[1]{\marginpar{\raggedright\footnotesize\ttfamily[#1]}}
\newcommand{\DERlacuna}{\ifmmode\mathord{[\ldots]}\else[\ldots]\fi}
\newfontfamily\DERappfont{DejaVu Sans}[Scale=MatchLowercase]
\newfontfamily\DERcodefont{DejaVu Sans Mono}[Scale=MatchLowercase]
\DeclareRobustCommand{\DERmark}[1]{\ifmmode\text{\hyperlink{app-#1}{\textsuperscript{\sffamily\tiny[#1]}}}\else\hyperlink{app-#1}{\textsuperscript{\sffamily\tiny[#1]}}\fi\allowbreak}
\newcommand{\DERapparatusitem}[5]{\par\medskip{\DERappfont\noindent\hypertarget{app-#1}{}\textbf{#1}\quad #2\par
  {\footnotesize\noindent\textbf{Print / imprimé:}\ {\DERcodefont #3}\par
  \noindent\textbf{Correction:}\ {\DERcodefont #4}\par
  \noindent\textbf{Source:}\ #5\par}}}
\begin{document}
\pagestyle{plain}
\hypersetup{pageanchor=false}
\begin{titlepage}\centering\vspace*{0.16\textheight}
{\Large Alexander Grothendieck\par}\vspace{2em}
{\Huge\bfseries Les Dérivateurs\par}\vspace{2em}
{\Large Texte diplomatique\par}\vspace{1em}
{\large Français\par}\vfill
{\normalsize Édition de M.~Künzer, J.~Malgoire et G.~Maltsiniotis\par}
\vspace{1em}{\small 12 novembre 2026\par}\end{titlepage}
\hypersetup{pageanchor=true}
\textbf{Crédits.} Auteur : Alexander Grothendieck, \emph{Les Dérivateurs} (1990–1991).
L’édition transcrite est celle de M.~Künzer, J.~Malgoire et G.~Maltsiniotis (Universités de Montpellier~II et Paris~VII).
Le déchiffrement, l’édition et les notes des éditeurs leur appartiennent.
Transcription produite par IA (\emph{AI-transcribed}) : Claude Opus~5.5 (Anthropic).
Commentaire du compagnon et corrections produits par IA (\emph{AI-produced corrections}) : GPT-6 Astra
(OpenAI, ChatGPT Pro mode), vérifiés et intégrés par Claude Opus~5.5 (Anthropic).
Huit blocs de commentaire des parties II-a, III-V et XIIb-a, sept paragraphes d’autres blocs, dix-neuf lignes du registre
des assertions (II-a, III-V, VII-c, VIII-X, XIIb-a) et trois entrées bibliographiques sont produits par Claude Fable~5.2
(Anthropic) à partir de réponses de GPT-6 Astra qu’il a examinées ; chaque ligne concernée indique réponse et examen.
Les corrections des définitions de macros et les modifications des textes vérifiés absentes des tables ont été
comparées aux pages de l’édition par Claude Opus~5.5 (\texttt{transcription/DECISIONS.md}).
Traduction anglaise produite par IA (\emph{AI-translated}) : GPT-6.1 Sol (OpenAI, effort max) ; glossaire de Claude
Fable~5.2 et Claude Opus~5.5 ; consigne, contrôles, corrections, comparaison d’intégralité et livre anglais de
Claude Opus~5.5 (Anthropic). Réponses de recherche brutes produites par GPT-6 Astra ; examens par Claude Fable~5.2.
Le présent assemblage des lecteurs, leur apparat et leur audit sont de GPT-6 Astra (OpenAI, ChatGPT Pro mode).

\medskip\noindent\url{https://webusers.imj-prg.fr/~georges.maltsiniotis/groth/Derivateurs.html}

\bigskip\noindent\textbf{Droits}\par
Les droits sur l'œuvre originale appartiennent à son auteur, à ses représentants ou à ses ayants droit. Là où nous détenons des droits, sur la transcription, la traduction, les corrections et la composition produites par IA, ils sont dédiés au domaine public sous CC0 1.0. Cette édition est publiée à des fins pédagogiques et non lucratives, et nous ne revendiquons aucun droit que nous ne détenons pas. Si vous détenez des droits sur cette œuvre et vous opposez à sa publication ici, ouvrez une issue sur GitHub et nous la retirerons.

\clearpage\textbf{Lecture.} Les marques marginales \texttt{p. n} donnent les pages de l’édition ; les marques
\texttt{[page n]} donnent les pages du manuscrit. Les arguments et l’ordre des ancres sont conservés.
Les mots entre crochets en caractères de machine à écrire appartiennent aux éditeurs de l’édition transcrite.
Les points d’interrogation isolés entre crochets signalent aussi leurs doutes, selon la convention exposée dans leur préface.
Les variantes de notre apparat portent un identifiant FRV ou ENV et un lien vers le relevé final ; elles donnent
la lecture imprimée, la correction et sa source. Pour identifier sans ambiguïté les symboles et les macros propres
à chaque unité, les deux lectures de l’apparat sont citées en code \LaTeX{} exact. Les notes originales ne sont
pas réécrites. Les trois formes proviennent d’un seul maître annoté. Les réparations de composition communes
aux trois formes sont consignées séparément des variantes textuelles.

\clearpage\tableofcontents\clearpage\pagestyle{derbody}
\clearpage\setcounter{footnote}{0}\phantomsection\addcontentsline{toc}{section}{Table des matières de l’édition}\markboth{Table des matières de l’édition}{Table des matières de l’édition}
% BEGIN UNIT dertbmt/dertbmt_p001-005.tex
\begingroup\makeatletter%
\newcommand{\printedpage}[1]{\marginpar{\footnotesize p.~#1}}%
\newcommand{\msp}[1]{\marginpar{\footnotesize [#1]}}%
\newcommand{\tochead}[2]{\par\medskip\noindent #1\hfill #2\par\smallskip}%
\newcommand{\tocl}[4]{\par\noindent\hangindent=\dimexpr #1+3em\relax\hspace*{#1}\makebox[2.5em][r]{#2}\hspace{0.5em}#3\dotfill #4}%
\newcommand{\Ch}{\underline{\mathrm{Ch}}}%
\makeatother\let\printedpage\derpp\let\msp\dermsp%
\let\mspin\DERmspin%


\printedpage{1}
% head: (none)

\begin{center}
{\LARGE\textbf{Les Dérivateurs}}

\bigskip

{\large Alexandre Grothendieck}
\end{center}

\bigskip

\tochead{\textbf{I. Généralités sur les (pré)dérivateurs.}}{[1]-[126]}
\tocl{0em}{1.}{Préliminaires sur les foncteurs images inverses des préfaisceaux sur des catégories d'indices, et leurs adjoints multiples}{[1]}
\tocl{0em}{2.}{Préliminaires heuristiques sur les dérivateurs}{[4]}
\tocl{0em}{}{Morphismes homotopiquement propres, homotopiquement lisses}{[53]}
\tocl{0em}{}{Notions de $\mathbb{D}$-équivalence et de $\mathbb{D}$-contractibilité}{[61]}
\tocl{0em}{}{$\mathbb{D}$-fibrations}{[76]}

\tochead{\textbf{II. Cofinalité (à droite et à gauche) (préliminaire à la cofinalité cohomologique).}}{[1]-[99]}
\tocl{0em}{1.}{Cofinalité cohomologique}{[1]}
\tocl{0em}{2.}{Les catégories localement filtrantes sont les catégories localement totalement 0-connexes}{[11]}
\tocl{0em}{3.}{Retour sur la cofinalité}{[37]}
\tocl{0em}{4.}{Foncteurs biasphériques}{[49]}
\tocl{0em}{5.}{Dualité formelle}{[69]}
\tocl{0em}{6.}{Transporteurs, noyaux, stabiliseurs relatifs à $W$}{[83]}
\tocl{0em}{7.}{Morphismes cofiltrants et localement cofiltrants}{[90]}

\tochead{\textbf{III. Hom externes dans les dérivateurs.}}{[1]-[12]}
\tocl{0em}{1.}{Relation des dérivateurs avec la catégorie Hot}{[1]}

\tochead{\textbf{IV. Diagrammes substantiels.}}{[73]-[83]}
\tocl{0em}{}{Construction de diagrammes substantiels pour un dérivateur}{[73]}

\tochead{\textbf{V. Catégories de chemins et localisation.}}{[1]-[52]}
\tocl{0em}{1.}{La catégorie des chemins}{[1]}
\tocl{0em}{2.}{Calcul d'une catégorie de fractions $X\Sigma^{-1}$}{[5]}
\tocl{0em}{3.}{Passage aux catégories de diagrammes}{[28]}

\tochead{\textbf{VI. HOT.}}{[1]-[204]}
\tocl{0em}{}{Le dérivateur HOT}{[1]}
\tocl{0em}{}{Nouvelle version de la théorie de $\mathrm{HOT}_W$}{[171]}

\par
\printedpage{2}
% head: 2

\tochead{\textbf{VII. Catégories de chemins (2).}}{[1]-[171]}
\tocl{0em}{1.}{Intervalles et leur variantes}{[1]} % [sic] leur variantes -> leurs variantes, accord du pluriel
\tocl{0em}{2.}{Sous-catégories pleines d'un intervalle}{[10]}
\tocl{0em}{3.}{Chemins. Catégorie de chemins d'un type donné $\tau$}{[13]}
\tocl{0em}{4.}{Catégorie de chemins $\Ch(X)$}{[24]}
\tocl{0em}{5.}{Propriétés des morphismes d'intervalles}{[31]}
\tocl{0em}{6.}{Composition dans \underline{Ty}}{[48]}
\tocl{0em}{7.}{Explicitation des catégories de fractions $X\Sigma^{-1}$ en termes des $\Ch(X;a,b)$}{[59]}
\tocl{0em}{8.}{Le lemme d'homotopie pour les morphismes d'intervalles}{[83]}
\tocl{0em}{9.}{Théorème d'homotopie pour les catégories de chemins}{[94]}
\tocl{0em}{10.}{Question de comparaison pour les catégories Cat-fibrées satisfaisant au “théorème d'homotopie” (base filtrante)}{[99]}
\tocl{0em}{11.}{Rapide retour sur les $(\mathcal{M},\Sigma)$, dans l'optique des dérivateurs}{[129]}
\tocl{0em}{12.}{Retour sur les ‘catégories de chemins’ dans Cat, comme catégorie de modèles pour Hot}{[137]}
\tocl{0em}{13.}{Où on prouve enfin que $\Ch(X;-,-)$ est aussi bon (voire meilleur) que $\Ch_\infty(X;-,-)$}{[157]}

\tochead{\textbf{VIII. 1-types d'homotopie relatifs : leur intégration\,\dots}}{[1]-[9]}
\tocl{0em}{1.}{Intégration des 1-types d'homotopie relatifs sur une catégorie}{[1]}
\tocl{0em}{2.}{Questions de changement de base}{[6]}

\tochead{\textbf{IX. Retour sur les catégories $\mathcal{M}W^{-1}$, comparaison avec Quillen.}}{[1]-[25]}
\tocl{0em}{1.}{Fonctorialité de $\mathrm{Ch}^{\Sigma}_{\infty}(a,b)$ en $a$, $b$ dans $\mathcal{M}\Sigma^{-1}$}{[1]}
\tocl{0em}{2.}{Comparaison avec Quillen}{[14]}
\tocl{0em}{3.}{Essai de construction du foncteur $\Omega$}{[24]}

\tochead{\textbf{X. Comparaison de Cat avec $\Delta^\wedge$.}}{[1]-[5]}
\tocl{0em}{1.}{Rappel des notions de $\mathbb{D}$-asphéricité, et lien avec le ‘localisateur fondamental $W_{\mathbb{D}}$’}{[1]}

\tochead{\textbf{XI. Hot-fibrations, foncteurs propres et foncteurs lisses etc.}}{[127]-[136]}
\noindent\hspace*{1.5em}(dans Cat).\par\smallskip
\tocl{0em}{1.}{Hot-fibrations dans Cat}{[127]}

\tochead{\textbf{XII. Caractérisation de $W_\infty$. Foncteurs $W$-propres, $W$-lisses etc. Sommes amalgamées et carrés $W$-cocartésiens dans Cat.}}{[1]-[275]}
\tocl{0em}{I}{Caractérisation de $W_\infty$}{[1]}
\tocl{0em}{II}{Perplexités sur les axiomes et la terminologie}{[18]}
\tocl{0em}{III}{Classes remarquables de foncteurs liés à $W$}{[22]}
\tocl{0em}{IV}{Sur l'axiome W(4) et ses variantes}{[53]}
\tocl{0em}{V}{Sommes amalgamées et carrés $W$-cocartésiens dans Cat}{[59]}
\tocl{2em}{1.}{Généralités sur $U\overset{i}{\hookrightarrow}X\overset{j}{\hookleftarrow}T$ et les sommes amalgamées}{[59]}
\tocl{4.5em}{1.1.}{Données de recollement entre un ouvert et un fermé}{[59]}
\tocl{4.5em}{1.2.}{Relation avec le recollement des topos}{[61]}

\par
\printedpage{3}
% head: 3

\tocl{4.5em}{1.3.}{Flèches dans Cat, de source ou but un objet recollé $(X,U,T)$}{[65]}
\tocl{4.5em}{1.4.}{Construction de sommes amalgamées à partir d'une immersion ouverte où fermée}{[66]} % [sic] où fermée -> ou fermée, conjonction (cf. 2.2)
\tocl{2em}{2.}{Sommes amalgamées et $W$-équivalences}{[75]}
\tocl{4.5em}{2.1.}{Axiome du carré cocartésien et axiome du cube (ou ‘Lemme des cinq’)}{[75]}
\tocl{4.5em}{2.2.}{Application aux carrés cocartésiens déduits d'une immersion ouverte ou fermée}{[80]}
\tocl{4.5em}{2.3.}{Objets totalement $W$-asphériques de Cat, et morphismes totalement $W$-asphériques}{[87]}
\tocl{4.5em}{2.4.}{Application aux immersions ouvertes dwyériennes}{[107]}
\tocl{4.5em}{2.5.}{Complément : critère d'asphéricité totale d'une catégorie $W$-propre sur une autre}{[115]}
\tocl{2em}{3.}{Sommes amalgamées et intégration}{[121]}
\tocl{4.5em}{3.1.}{Réduction canonique au cas d'un diagramme d'immersions fermées}{[121]}
\tocl{4.5em}{3.2.}{Comparaison de l'intégrale avec une somme amalgamée ordinaire. Carrés $W$-homotopiquement cocartésiens}{[128]}
\tocl{4.5em}{3.3.}{Intégration propre versus intégration lisse. Énoncé d'équivalence entre les deux points de vue. Commentaires sur le cas des limites $W$-homotopiques quelconques}{[133]}
\tocl{4.5em}{3.4.}{Critères pour les carrés homotopiquement cocartésiens}{[146]}
\tocl{4.5em}{3.5.}{Applications aux sommes amalgamées dans une catégorie $X^\wedge$ (avec $X\in\mathrm{Ob}\,\mathrm{Cat}$)}{[158]}
\tocl{0em}{VI}{Critère de $W_\infty$-asphéricité pour les $\Omega^n(X,x)$}{[171]}
\tocl{2em}{1.}{Cogitations sur $W_n$}{[171]}
\tocl{2em}{2.}{Critère de $W_n$-asphéricité}{[187]}
\tocl{2em}{3.}{Sur la philosophie des fibres homotopiques}{[191]}
\tocl{2em}{4.}{Axiomatisation des ‘fibres $W$-homotopiques’}{[207]}
\tocl{2em}{5.}{Axiomatique des $W$-types d'homotopie relatifs non localement constants $\mathrm{Ho}^{\mathcal{L}}_{W}(S)$}{[218]}
\tocl{2em}{6.}{Cofibre $W$-homotopique dans Cat}{[227]}
\tocl{2em}{7.}{Retour sur l'axiomatique de la fibre $W$-homotopique}{[240]}
\tocl{2em}{8.}{Appendice : Retour sur les axiomes W(7), W(7 bis) (pages 75, 76)}{[265]}

\tochead{\textbf{XIII. Catégories de modèles (1).}}{[1]-[98]}
\tocl{0em}{1.}{Couples de Quillen $(\Phi,\Psi)$}{[1]}
\tocl{0em}{2.}{Catégories de Quillen faibles, et construction d'icelles (Théorème-scholie conjectural)}{[7]}
\tocl{0em}{3.}{Catégories de Quillen (deuxième mouture)}{[27]}
\tocl{0em}{4.}{Les variances $u^*$, $u_*$, $u_!$ pour $(\mathcal{M}(I),W(I))$ (I variable)}{[44]}
\tocl{0em}{5.}{Catégories de modèles quillénisables}{[50]}
\tocl{0em}{6.}{Catégories de modèles quillénisables, et foncteurs $u_!$, $u_*$}{[57]}
\tocl{0em}{7.}{Filtrations cardinales et parties accessibles d'une grosse catégorie}{[77]}

\par
\printedpage{4}
% head: 4

\tochead{\textbf{XIV. Carrés h-cartésiens et h-cocartésiens.}}{[1]-[8]}
\tocl{0em}{1.}{Axiomes de structures quadrangulées (une tentative avortée)}{[1]}
\tocl{0em}{2.}{Exemples}{[7]}

\tochead{\textbf{XV. Théorèmes de factorisation. (Modèles (2)).}}{[1]-[137]}
\tocl{0em}{1.}{Construction fondamentale : la suite transfinie de factorisation d'une flèche $f$}{[1]}
\tocl{0em}{2.}{Variation des factorisations $f=p_\alpha i_\alpha$ avec la partie $C_0$ génératrice}{[12]}
\tocl{0em}{3.}{Majoration des cardinaux}{[16]}
\tocl{0em}{4.}{Compléments sur les structures de Quillen}{[28]}
\tocl{0em}{5.}{Résultats téchniques préliminaires sur les parties accessibles de $\mathrm{Ob}\,\mathcal{M}$, stables par limites inductives filtrantes}{[43]} % [sic] téchniques -> techniques, accent
\tocl{0em}{6.}{Essai avorté pour un critère d'engendrement des parties L-accessibles de $\mathrm{Fl}(\mathcal{M})$}{[73]}
\tocl{0em}{7.}{Construction de triples de Quillen clos avec $C=\mathrm{Mon}\subset\mathrm{Cof}_W$}{[84]}
\tocl{0em}{8.}{Application aux catégories de modèles $(\widehat{A},W_A)$}{[100]}
\tocl{0em}{9.}{Révision des notations, pour les axiomes d'un localiseur fondamental}{[116]}
\tocl{0em}{10.}{Application au dérivateur $\mathrm{HOT}_W$}{[134]}

\tochead{\textbf{XVI. Localiseurs fondamentaux dans Cat.}}{[1]-[124]}
\tocl{0em}{1.}{Axiomes Loc\,(1) à Loc\,(3) (et variantes): localiseurs fondamentaux}{[1]}
\tocl{0em}{2.}{Axiomes Loc\,(4) et Loc\,(5) (des sommes directes, de connexité)}{[8]}
\tocl{0em}{3.}{Axiome Loc\,(6) du carré cocartésien}{[16$'$]}
\tocl{0em}{4.}{Axiome Loc\,(7) des $W$-fibrations}{[27]}
\tocl{0em}{5.}{Axiome des limites Loc\,(8), et l'axiome Loc\,(9) d'accessibilité}{[71]}
\tocl{0em}{6.}{La condition Loc\,$(5,\mathbb{D}_{W_\omega})$ implique $W\subset\mathrm{W}_\infty$}{[83]}
\tocl{0em}{7.}{Commentaires sur les axiomes}{[97]}

\tochead{\textbf{XVII. Catégories à fibrations et à cofibrations. (Modèles (3)).}}{[1]-[118]}
\tocl{0em}{1.}{Préliminaires}{[1]}
\tocl{0em}{2.}{Description de $C_I\subset\mathrm{Fl}(\mathcal{M}(I^\circ))$, pour $I$ ensemble ordonné}{[25]}
\tocl{0em}{3.}{Catégories de modèles à cofibrations}{[39]}
\tocl{0em}{4.}{Catégories à cofibrations spéciales}{[71]}
\tocl{0em}{5.}{Relation des catégories à cofibrations spéciales avec les catégories de Quillen spéciales}{[84]}
\tocl{0em}{6.}{Application à la construction des dérivateurs (associés aux triples de Quillen spéciaux)}{[98]}

\tochead{\textbf{XVIII. Catégories et ensembles accessibles.}}{[50]-[492]}
\tocl{0em}{4.}{Paratopos et l'opération $\mathcal{M}\boxtimes\mathcal{N}$}{[50]}
\tocl{2em}{4.1.}{Pseudo-topos. Problème du produit}{[50]}
\tocl{2em}{4.2.}{Cafouillages sur $\mathcal{M}\boxtimes\mathcal{N}$}{[57]}
\tocl{2em}{4.3.}{L'opération $\mathcal{M}\boxtimes\mathcal{N}$ et l'axiome d'accessibilité}{[64]}
\tocl{2em}{4.4.}{Étude de $\mathrm{Ind}_\pi(C)$ (pour $C$ petite catégorie stable par $\varinjlim$ de cardinal $\leq\pi$)}{[79]}

\par
\printedpage{5}
% head: 5

\tocl{2em}{4.5.}{Filtration cardinale canonique d'une catégorie $\mathcal{U}$-paratopos}{[106]}
\tocl{2em}{4.6.}{Caractérisation des catégories paratopos (démonstration de 4.3.6)}{[139]}
\tocl{2em}{4.7.}{$\mathrm{Ind}_\pi C$ ($C$ quelconque) et filtrations cardinales généralisées}{[161]}
\tocl{2em}{4.8.}{Catégories $\mathcal{U}$-accessibles}{[187]}
\tocl{2em}{4.9.}{Le théorème de représentation indicielle pour $\mathrm{Ind}_\pi(C)$}{[217]}
\tocl{2em}{4.10.}{Le théorème de représentation indicielle : démonstration}{[243]}
\tocl{2em}{4.11.}{Application à l'existence des limites dans une catégorie de la forme $\mathrm{Ind}_{\pi_0}(C)$}{[267]}
\tocl{2em}{4.12.}{Foncteurs accessibles entre catégories accessibles}{[273]}
\tocl{2em}{4.13.}{Ensembles accessibles d'objets d'une catégorie accessible}{[281]}
\tocl{2em}{4.14.}{Critères d'accessibilité pour grosses catégories}{[370]}
\tocl{2em}{4.15.}{Catégorie accessible sur une autre}{[425]}
\tocl{2em}{4.16.}{Exemples et questions}{[464]}
\tocl{2em}{4.17.}{Retour sur les paratopos, et sur les opérations $\mathcal{M}\boxtimes\mathcal{N}$, $\underline{\mathrm{Hom}}_!(\mathcal{M},\mathcal{N})$}{[492]}
\par\noindent\hspace*{1.5em}\textbf{Appendice}.\hfill[1]-[19]
\tocl{2em}{}{Ultraproduits de corps, et catégories accessibles}{[1]}
\tocl{2em}{}{Application à un exemple d'une sous-catégorie non accessible}{[11]}

\tochead{\textbf{XIX. Modèles (4).}}{[1]-[50]}
\tocl{0em}{1.}{Rappels sur les couples de Quillen (couples de Quillen spéciaux)}{[1]}
\tocl{0em}{2.}{Rappels sur les triples de Quillen clos et les triples de Quillen spéciaux}{[13]}
\tocl{0em}{3.}{Cas des catégories à cofibrations, des catégories à fibrations}{[33]}

\bigskip
\noindent\rule{4cm}{0.4pt}

\par\noindent\textbf{Note des éditeurs de l’édition transcrite.}\par
{\footnotesize\noindent\hangindent=2em \textsc{Note}. Dans cette table de matières, ainsi que dans l'édition électronique provisoire des “Dérivateurs”, on suit le découpage de l'ouvrage effectué a posteriori par Grothendieck. Les numéros des pages sont ceux de l'original, indiquées entre crochets dans le texte numérisé. Le chapitre XI, qui commence à la page 127, était à l'origine la dernière section du chapitre I, qui se termine à la page 126. De même, le chapitre XVIII était la section 4 du chapitre XIX.\par}


\endgroup
% END UNIT dertbmt/dertbmt_p001-005.tex
\clearpage\setcounter{footnote}{0}\phantomsection\addcontentsline{toc}{section}{Chapitre I}\markboth{Chapitre I}{Chapitre I}
% BEGIN UNIT Der-ChI-234/Der-ChI-234_p001-005.tex
\begingroup\makeatletter%
\newcommand{\printedpage}[1]{\marginpar{\footnotesize p.~#1}}%
\newcommand{\msp}[1]{\marginpar{\footnotesize [#1]}}%
\newcommand{\AMSTeX}{$\mathcal{A}\kern-.1667em\lower.5ex\hbox{$\mathcal{M}$}\kern-.125em\mathcal{S}$-\TeX}%
\newcommand{\uHom}{\underline{\mathsf{Hom}}}%
\newcommand{\Ob}{\mathsf{Ob}}%
\newcommand{\Ab}{\mathsf{Ab}}%
\newcommand{\bDelta}{\boldsymbol{\Delta}}%
\makeatother\let\printedpage\derpp\let\msp\dermsp%
\let\mspin\DERmspin%


% ===================== page 1 =====================
\printedpage{--} % no page number printed (title page)
% head: (none)

\begin{center}
{\Large\textsf{LES DÉRIVATEURS}}

\bigskip

\textbf{ALEXANDRE GROTHENDIECK}

\bigskip\bigskip

\textbf{Chapitre I}

\medskip

\textbf{Généralités sur les (pré)dérivateurs}
\end{center}

\par

% ===================== page 2 =====================
\printedpage{--} % no page number printed (editors' note)
% head: (none)

\par\noindent\textbf{Note des éditeurs de l’édition transcrite.}\par
Ce texte a été déchiffré et transcrit en \AMSTeX{} par G. Maltsiniotis. Il a été édité par M. Künzer, J. Malgoire, et G. Maltsiniotis. La transcription est aussi fidèle que possible au manuscrit. Pour les quelques corrections évidentes, ou les rares commentaires des éditeurs, ainsi que pour la numérotation originale des pages du manuscrit, les caractères de machine à écrire \texttt{[typewriter]} entre crochets sont utilisés. Un point d'interrogation entre crochets signifie que l'on n'est pas sûr du mot qui précède.

Cette édition est provisoire. Les remarques, commentaires et corrections sont bienvenus. Envoyer un message à :

\textbf{[e-mail address removed]}

\hfill G. Maltsiniotis

\par

% ===================== page 3 =====================
\printedpage{1}
% head: A. Grothendieck : Dérivateurs, Ch. I, Édition des Universités de Montpellier II et Paris VII
% foot: Version 8 juillet 2001

\msp{page 1}

\section*{Dérivateurs}

\subsection*{\textmd{1.}\ Préliminaires sur les foncteurs images inverses des préfaisceaux sur des catégories d'indices, et leurs adjoints multiples.}

$A$ “catégorie de coefficients”,\\
$I$ “catégorie d'indices”,\\
$A(I)=\uHom(I^{\circ},A)$ préfaisceaux sur $I$ à valeurs dans $A$.

Si $f : I \to J$, on en déduit :
\[
f^{*} : A(J) \to A(I) \quad ,
\]
et on s'intéresse à l'existence des foncteurs adjoints successifs à gauche et à droite. Notation provisoire :
\[
f^{?}\ f_{!}\ f^{*}\ f_{*}\ f^{!} \quad .
\]

Voici les faits principaux :

\textit{a}) Sous réserve de l'existence des $\varprojlim$ (resp. $\varinjlim$) pertinentes, $f_{*}$ (resp. $f_{!}$) existe. Par exemple s'il existe \texttt{[des]} $\varprojlim$ (resp. $\varinjlim$) \textit{finies}, alors $f_{*}$ (resp. $f_{!}$) existe si $I$, $J$ sont des catégories finies.

\textit{b}) Même si $A$ est une catégorie abélienne, avec $\varprojlim$ et $\varinjlim$ (par exemple $\Ab$), les \textit{biadjoints} n'existent pas toujours. Par exemple, dans le cas particulier où $J=e$ est la catégorie ponctuelle, et $I$ l'ensemble ordonné
\[
\begin{tikzcd}
 & a \\
0 \arrow[ur] \arrow[dr] & \\
 & b
\end{tikzcd}
\]
(dans le cas de $f^{!}$), ou l'ensemble ordonné opposé (dans le cas de $f^{?}$). Même si $f$ est une immersion ouverte, $f^{!}$ n'existe pas forcément (par exemple elle n'existe pas si $I$ est comme \texttt{[ci-]}dessus, et
\[
J=\begin{tikzcd}
 & a \arrow[dr] & \\
0 \arrow[ur] \arrow[dr] & & c \\
 & b \arrow[ur] &
\end{tikzcd}
\qquad\qquad \text{),}
\]
et si $f$ est une immersion fermée, $f^{?}$ n'existe pas forcément.

\textit{c}) Si $f$ est une \textit{immersion fermée}, et $A$ est \textit{ponctuée} (\textit{i.e.} admet un objet à la fois initial et final), alors $f^{!}$ existe. Si $f$ est une immersion ouverte, $f^{?}$ existe. (Toujours sous réserve
\printedpage{2}
% head: A. Grothendieck : Dérivateurs, Ch. I, Édition des Universités de Montpellier II et Paris VII
% foot: Version 8 juillet 2001
de l'existence dans $A$ des limites pertinentes). Si $f$ est une immersion, $f_{*}$ et $f_{!}$ \texttt{[sont]} pleinement fidèles \textit{i.e.} $f^{*}$ \texttt{[est]} un foncteur de localisation.

\textit{d}) Supposons seulement l'existence de l'objet final $e$ dans $A$ (plus au besoin certaines $\varprojlim$ finies). Si

\msp{page 2}
\[
i : Y \hookrightarrow S
\]
est une immersion fermée, de complémentaire $U$
\[
Y \overset{i}{\hookrightarrow} S \overset{j}{\hookleftarrow} U
\]
de sorte que $j$ soit une immersion ouverte, alors $A(Y)$ s'identifie par le foncteur $i_{*}$ à la sous-catégorie pleine de $A(S)$ formée des $F \in \Ob\ A(S)$ tels que $F|U \simeq e_{U}$ (préfaisceau final). De même, si l'objet initial $\emptyset$ existe dans $A$, et les $\varinjlim$ pertinentes, alors $A(U)$ s'identifie par $j_{!}$ à la sous-catégorie pleine de $A(S)$ formée des $F$ tels que $F|Y=\emptyset_{Y}$ (préfaisceau initial, prenant partout la valeur $\emptyset$).

Si $A$ est même \textit{ponctuée} ($e=\emptyset$), alors on a, pour tout $F\in \Ob\ A(S)$, un carré à la fois cartésien et cocartésien, fonctoriel en $F$
\[
\begin{tikzcd}
 & F \arrow[dr, "\text{adj}"] & \\
j_{!}j^{*}F \arrow[ur, "\text{adj}"] \arrow[dr] & & i_{*}i^{*}F \\
 & e \arrow[ur] &
\end{tikzcd}
\qquad\qquad \text{.}
\tag*{(1)}
\]
Dans le cas $A$ abélienne, ce diagramme prend la forme d'une suite exacte
\[
0 \longrightarrow j_{!}j^{*}F \longrightarrow F \longrightarrow i_{*}i^{*}F \longrightarrow 0 \quad .
\]

\textit{e}) Mêmes conditions pour $i$ et $j$. On suppose que $A$ a les $\varprojlim$ pertinentes, dont objet final. On a alors un carré cartésien fonctoriel en $F\in \Ob\ A(S)$ \texttt{[il faut supposer la catégorie} $A$ \texttt{ponctuée]}
\[
\begin{tikzcd}
F \arrow[d, "\text{adj}"'] & i_{*}i^{!}F \arrow[l] \arrow[d] \\
j_{*}j^{*}F & e \arrow[l]
\end{tikzcd}
\tag*{(2)}
\]
\textit{i.e.} $i^{!}F$ s'identifie par le foncteur pleinement fidèle $i_{*}$ au \textit{noyau} de $F\to j_{*}j^{*}F$. Dans le cas abélien, cela prend la forme d'une suite exacte
\[
0 \longrightarrow i_{*}i^{!}F \longrightarrow F \longrightarrow j_{*}j^{*}F \quad .
\]

\par

% ===================== page 5 =====================
\printedpage{3}
% head: A. Grothendieck : Dérivateurs, Ch. I, Édition des Universités de Montpellier II et Paris VII
% foot: Version 8 juillet 2001

\msp{page 3}

Je me dispense d'écrire le carré cocartésien dual, quand $A$ a les $\varinjlim$ pertinentes, y compris objet initial. Il semble qu'elle n'ait pas servi en pratique. Les foncteurs adjoints et biadjoints classiquement utiles sont inclus ci-dessous (où les “pas utiles” sont entre crochets)
\[
\begin{aligned}
{}[i_{!}]\ i^{*}\ i_{*}\ i^{!} \\
[j^{?}]\ j_{!}\ j^{*}\ j_{*}
\end{aligned}
\]
\textit{i.e.} les foncteurs $i_{!}$ et $j^{?}$ n'ont pas servi classiquement.

\textit{f}) Cas des catégories $\bDelta_{i}$ (simplexes type, $i\in\mathbb{N}$), et des foncteurs entre ceux-ci (applications semi-simpliciales). Il semblerait que pour une telle $f$, si $A$ est ponctuée, il y ait les \textit{bi-adjoints} à la fois à gauche et à droite. De plus, pour $\bDelta_{i}$ et $\bDelta_{j}$ fixés quand on regarde toutes les $f:\bDelta_{i}\to\bDelta_{j}$ et toutes les $g:\bDelta_{j}\to\bDelta_{i}$, que les suites totales de foncteurs multi-adjoints pour chacun des $f^{*}$, $g^{*}$ correspondants sont toutes les mêmes \texttt{[ce n'est pas vrai]}. Ça a l'air vrai du moins si $j=i+1$ (et je l'ai vérifié pour $i=0$, $i=1$), on doit trouver une “suite d'adjonctions” de longueur $2i+7=(i+1)+(i+2)+2+2$. De façon précise, la suite comporte quatre termes “extrêmes” (les deux premiers, les deux derniers) et les $(i+1)+(i+2)$ termes intermédiaires correspondant aux foncteurs de la forme $f^{*}$ ou $g^{*}$, où $f:\bDelta_{i+1}\to\bDelta_{i}$ est une des $(i+1)$ applications de dégénérescence, et $g:\bDelta_{i}\to\bDelta_{i+1}$ une des $i+2$ applications bord \texttt{[ce n'est pas tout à fait vrai; il faut se limiter aux applications qui sont injectives ou surjectives]}.


\endgroup
% END UNIT Der-ChI-234/Der-ChI-234_p001-005.tex
% BEGIN UNIT Der-ChI-234/Der-ChI-234_p006-010.tex
\begingroup\makeatletter%
\newcommand{\printedpage}[1]{\marginpar{\footnotesize p.~#1}}%
\newcommand{\msp}[1]{\marginpar{\footnotesize [#1]}}%
\newcommand{\Dia}{\mathsf{Dia}}%
\newcommand{\Cat}{\mathcal{C}\mathit{at}}%
\newcommand{\DD}{\mathbb{D}}%
\newcommand{\Hom}{\underline{\mathsf{Hom}}}%
\newcommand{\id}{\mathsf{id}}%
\newcommand{\Ob}{\mathsf{Ob}}%
\newcommand{\Del}{\boldsymbol{\Delta}}%
\makeatother\let\printedpage\derpp\let\msp\dermsp%
\let\mspin\DERmspin%


\printedpage{4}
% head: A. Grothendieck : Dérivateurs, Ch. I, Édition des Universités de Montpellier II et Paris VII
% foot: Version 8 juillet 2001

\msp{page 4}

\section*{2. Préliminaires heuristiques sur les dérivateurs.}

On se donne une sous-catégorie pleine $\Dia$ (comme “diagrammes”) de $\Cat$. Ça peut être tout $\Cat$, ou plus petit. Mais on supposera le plus souvent que $\Dia$ contient au moins les “catégories ordonnées” (correspondant aux ensembles ordonnés, \textit{i.e.} telles que $a$) tout morphisme est un mono et $b$) tout iso est une identité \texttt{[ce n'est pas vrai : la catégorie réduite à une double flèche satisfait à ces conditions, mais n'est pas un ensemble ordonné]}). Par abus de langage, on identifie un ensemble ordonné avec la catégorie correspondante (et une catégorie au topos qu'elle définit?). On considère $\Dia$ comme une 2-\textit{catégorie}.\footnote{N.B. Il faut des conditions de stabilité pour $\Dia$, au moins la stabilité par $I \times J$ et $I \amalg J$, en fait, nettement plus \dots}

Un \textit{dérivateur} $\DD$ (sur $\Dia$, comme “\textit{domaine}” du dérivateur) est un 2-foncteur contravariant de $\Dia$ dans $\Cat$, satisfaisant certaines conditions. Donc
\[ \DD : \Dia^{\circ} \longrightarrow \Cat \quad . \]
Si $I \in \Ob_0 \Dia$ est une catégorie (ou un “type de diagramme”) du domaine de $\DD$, on doit donc avoir une catégorie $\DD(I)$. Les objets de $\DD(I)$ s'appelleront les \textit{objets-coefficients} (ou \textit{coefficients}) de type $\DD$, au dessus de $I$. Pour
\[ f : I \longrightarrow J \]
une 1-flèche de $\Dia$, on doit avoir un foncteur
\[ f^*_{\DD} \text{ ou } f^* : \DD(J) \longrightarrow \DD(I) \quad , \]
dit \textit{“foncteur image inverse”}. Si $F$ est un coefficient de type $\DD$ sur $J$ son “\textit{image inverse}” est donc un coefficient de type $\DD$ sur $I$. On suppose (pour simplifier!) que le 2-foncteur $\DD$ est \textit{strict}, donc qu'on a transitivité stricte
\[ (gf)^* = f^*g^* \ , \qquad \text{si} \qquad I \xrightarrow{f} J \xrightarrow{g} K \qquad \text{dans} \ \ \Dia \ . \]

\msp{page 5}

De plus, si on a deux foncteurs $f, g : I \rightrightarrows J$, et un morphisme entre eux
\[
\begin{tikzcd}[column sep=large]
I \arrow[r, shift left=2, "f"] \arrow[r, shift right=2, "g"'] \arrow[r, phantom, "{\downarrow}u"] & J
\end{tikzcd}
\quad ,
\]
on se donne un morphisme $u^*_{\DD}$ ou $u^*$ entre les foncteurs $g^*$ et $f^*$, d'où
\[
\begin{tikzcd}[column sep=large]
\DD(I) & \DD(J) \arrow[l, shift right=2, "f^*"'] \arrow[l, shift left=2, "g^*"] \arrow[l, phantom, "{\uparrow}u^*"]
\end{tikzcd}
\quad .
\]

\printedpage{5}
% head: A. Grothendieck : Dérivateurs, Ch. I, Édition des Universités de Montpellier II et Paris VII
% foot: Version 8 juillet 2001

On doit encore avoir transitivité pour une composition de morphismes de foncteurs :
\[
\begin{tikzcd}[column sep=large]
I \arrow[r, shift left=4, "f"] \arrow[r, "g"' near start] \arrow[r, shift right=4, "h"'] \arrow[r, phantom, shift left=2, "{\downarrow}u"] \arrow[r, phantom, shift right=2, "{\downarrow}v"] & J
\end{tikzcd}
\]
doit donner
\[ (vu)^* = u^*v^* \]
donc
\[
\begin{tikzcd}[column sep=large]
\DD(I) & \DD(J) \arrow[l, shift right=4, "f^*"'] \arrow[l, "g^*" near end] \arrow[l, shift left=4, "h^*"] \arrow[l, phantom, shift right=2, "{\uparrow}u^*"] \arrow[l, phantom, shift left=2, "{\uparrow}v^*"]
\end{tikzcd}
\quad .
\]
De plus, si $f$ est une identité, $f^*$ aussi
\[ (\id_I)^* = \id_{\DD(I)} \quad , \]
et de même si $u$ est une identité
\[ (\id_f)^* = \id_{f^*} \quad . \]

\msp{page 6}

J'indique tout de suite l'axiome le plus important sur un dérivateur (au-delà de la donnée d'un 2-contra foncteur strict) :

\textbf{Der}$_1$ : Si $f : I \to J$ dans $\Dia$, alors $f^*_{\DD} = f^* : \DD(J) \to \DD(I)$ a un adjoint à droite et un adjoint à gauche.

Ces foncteurs seront notés (du moins provisoirement) $f^{\DD}_*$ ou $f_*$ , et $f^{\DD}_!$ ou $f_!$ , soit
\[ f_! \, , \ f_* : \DD(I) \rightrightarrows \DD(J) \quad , \]
de sorte qu'on trouve une suite de \textit{trois} foncteurs adjoints entre $\DD(I)$ et $\DD(J)$
\[ f_! \, , \ f^*, \ f_* \quad . \]

N.B. Pour la plupart des développements on n'aura pas besoin de l'existence à la fois des $f_!$ et des $f_*$ . Un \textit{dérivateur à gauche} ne suppose que les $f_*$ (adjoints à droite des $f^*$), et l'existence de $f_!$ seulement quand $f$ est une \textit{immersion ouverte}. Dualement, un \textit{dérivateur à droite} ne suppose que les $f_!$ (adjoints à gauche des $f^*$), et $f_*$ seulement si $f$ est une \textit{immersion fermée}. Ce sont sans doute les dérivateurs à gauche qui sont les plus importants.

\msp{page 7}

\textbf{Der}$_2$ : Si $f : I \to J$ est dans $\Dia$, si c'est une immersion \textit{fermée}, le biadjoint à droite $f^!_{\DD}$ ou $f^!$ de $f^*$ existe. Si c'est une immersion ouverte, le biadjoint à gauche $f^?$ existe.\footnote{Il n'y a lieu d'exiger les biadjoints que pour un dérivateur ponctué. \textsf{HOT} n'y satisfait pas.} Dans
\printedpage{6}
% head: A. Grothendieck : Dérivateurs, Ch. I, Édition des Universités de Montpellier II et Paris VII
% foot: Version 8 juillet 2001
les deux cas, on veut que $f^*$ soit un foncteur de localisation, ce qui équivaut à dire que $f_*$ est pleinement fidèle, ou que $f_!$ est pleinement fidèle.\footnote{Il faut de plus caractériser l'image de $f_*$ si $f$ est une immersion fermée (les $F$ tels que $g^*F \overset{\sim}{\to} e_{\DD(U)}$) et celle de $f_!$ si c'est une immersion ouverte (les $F$ tels que $g^*F \simeq \emptyset_{\DD(Y)}$), où $g$ est l'immersion complémentaire.}

N.B. Dans le cas d'un dérivateur à gauche, on suppose seulement l'existence de $f^!$ si $f$ est une immersion fermée. Donc pour résumer, pour un tel dérivateur, on suppose l'existence de
\[
\left\{
\begin{aligned}
& f_* \text{ quel que soit } f : I \to J \text{ dans } \Dia \\
& f_! \text{ si } f \text{ est une immersion ouverte} \\
& f^! \text{ si } f \text{ est une immersion fermée (cas ponctué seulement).}
\end{aligned}
\right.
\]
Dualement, pour un dérivateur à droite, on suppose l'existence de
\[
\left\{
\begin{aligned}
& f_! \text{ quel que soit } f : I \to J \text{ dans } \Dia \\
& f_* \text{ si } f \text{ est une immersion fermée} \\
& f^? \text{ si } f \text{ est une immersion ouverte (cas ponctué seulement).}
\end{aligned}
\right.
\]
Je voudrais formuler des axiomes astucieux, pour le cas d'un couple
\[ Y \overset{i}{\hookrightarrow} I \overset{j}{\hookleftarrow} U \]
d'une immersion ouverte $j$ et d'une immersion fermée $i$ dans $\Dia$, complémentaires l'une de l'autre. D'abord une formulation si on suppose $\DD$ “\textit{ponctué}”, \textit{i.e.} les catégories $\DD(I)$ ponctuées, et les foncteurs $f^*$ compatibles avec les ponctuations. On veut “interpréter” le carré (1) page 2 :

\msp{page 8}

\[
\begin{tikzcd}
& F \arrow[dr, "\mathrm{adj}"] & \\
j_!j^*F \arrow[ur, "\mathrm{adj}"] \arrow[dr] & & i_*i^*F \\
& e \arrow[ur] &
\end{tikzcd}
\quad ,
\tag*{(1)}
\]

\textbf{Der}$'_2$ : Pour tous $I, J \in \Dia$, on a
\[ \DD(I \amalg J) \xrightarrow[(\alpha^*, \beta^*)]{\sim} \DD(I) \times \DD(J) \quad , \]
où $\alpha, \beta$ sont les injections canoniques de $I$ et $J$ dans $I \amalg J$.

\printedpage{7}
% head: A. Grothendieck : Dérivateurs, Ch. I, Édition des Universités de Montpellier II et Paris VII
% foot: Version 8 juillet 2001

carré à la fois cartésien et cocartésien. On cherche donc un carré dans $\DD(I)$, mais mieux que ça. Soit $J = \Del_1 \times \Del_1$ le carré standard
\[
\begin{tikzcd}
& (1,0) \arrow[dl] & \\
(1,1) & & (0,0) \arrow[ul] \arrow[dl] \\
& (0,1) \arrow[ul] &
\end{tikzcd}
\]
et considérons la catégorie $I \times J$. On a
\[
\begin{aligned}
& J \longrightarrow \Hom(I, I \times J) \\
& j \longmapsto \bigl(i \longmapsto (i,j)\bigr)
\end{aligned}
\]
D'autre part, par $\DD$
\[ \Hom(I, I \times J)^{\circ} \longrightarrow \Hom(\DD(I \times J), \DD(I)) \]
d'où un composé
\[ J^{\circ} \longrightarrow \Hom(\DD(I \times J), \DD(I)) \]
ou encore
\[ \DD(I \times J) \longrightarrow \Hom(J^{\circ}, \DD(I)) \quad . \]
Ceci vaut pour tout $J \in \Ob \, \Dia$. On peut donc dire que la donnée d'un “coefficient” $\Phi$ sur une catégorie produit $I \times J$ donne naissance à un diagramme de type $J^{\circ}$ dans la catégorie $\DD(I)$ des coefficients sur $I$. La philosophie
\msp{page 9}
des dérivateurs consiste en grande partie, inversement, chaque fois qu'on a un diagramme dans une catégorie de coefficients $\DD(I)$ (sur un diagramme $I$), disons diagramme de type $K$ : $K \to \DD(I)$, de se demander s'il ne provient pas “canoniquement” d'un objet de $\DD(I \times K^{\circ})$ \textit{i.e.} d'un coefficient sur $I \times K^{\circ}$.

Ainsi pour le carré p. 8. On voudrait qu'il provienne de $\DD(I \times J)$, où $J = \Del_1 \times \Del_1$, et ceci fonctoriellement, donc on désire un foncteur “canonique”
\[ \DD(I) \longrightarrow \DD(I \times \Del_1 \times \Del_1) \]
qui en composant avec
\[ \DD(I \times \Del_1 \times \Del_1) \longrightarrow \Hom((\Del_1 \times \Del_1)^{\circ}, \DD(I)) \]
\printedpage{8}
% head: A. Grothendieck : Dérivateurs, Ch. I, Édition des Universités de Montpellier II et Paris VII
% foot: Version 8 juillet 2001
donne le foncteur décrit par le carré (1). De plus, il faut exprimer dans le contexte présent la condition “bicartésienne” du carré (1), dans le contexte de la page 7. Question analogue pour le carré (2) de la page 2 :
\[
\begin{tikzcd}
& F \arrow[dr] & \\
i_*i^!F \arrow[ur] \arrow[dr] & & j_*j^*F \\
& e \arrow[ur] &
\end{tikzcd}
\quad ,
\tag*{(2)}
\]
où il faut de plus interpréter la condition cartésienne. On va définir en fait un coefficient sur la sous-catégorie ouverte $W$ suivante de $\Del_1 \times \Del_1$
\[
\begin{tikzcd}
(1,0) & \\
& (0,0) \arrow[ul] \arrow[dl] \\
(0,1) &
\end{tikzcd}
\quad ,
\]

\msp{page 10}

plus exactement sur $W \times I$, et si $\alpha$
\[ \alpha : W \times I \longrightarrow \Del_1 \times \Del_1 \times I \]
est l'inclusion canonique, on définira le carré “cartésien” comme l'image par $\alpha_*$ de ce coefficient. Ceci nous ramène donc au problème de “globaliser” d'abord le diagramme partiel
\[
\begin{tikzcd}
F \arrow[dr] & \\
& j_*j^*F \\
e \arrow[ur] &
\end{tikzcd}
\quad .
\]
Et de même pour le carré (1) sauf qu'on est embarrassé si on choisit de le regarder surtout comme cartésien (ce qui suggère à procéder comme tantôt) ou comme \underline{co}cartésien (ce qui % verified transcription correction: the edition prints embarrassé, with double r


\endgroup
% END UNIT Der-ChI-234/Der-ChI-234_p006-010.tex
% BEGIN UNIT Der-ChI-234/Der-ChI-234_p011-015.tex
\begingroup\makeatletter%
\NoAutoSpacing%
\newcommand{\printedpage}[1]{\marginpar{\footnotesize p.~#1}}%
\newcommand{\msp}[1]{\marginpar{\footnotesize [#1]}}%
\newcommand{\DD}{\mathbb{D}}%
\newcommand{\uHom}{\underline{\mathsf{Hom}}}%
\newcommand{\Dia}{\mathsf{Dia}}%
\newcommand{\Ens}{\mathcal{E}ns}%
\newcommand{\bDelta}{\boldsymbol{\Delta}}%
\newcommand{\defeq}{\overset{\text{déf}}{=}}%
\makeatother\let\printedpage\derpp\let\msp\dermsp%
\let\mspin\DERmspin%


\printedpage{9}
% head: A. Grothendieck : Dérivateurs, Ch. I, Édition des Universités de Montpellier II et Paris VII
\noindent suggère de procéder de façon duale). Mais les deux coefficients qu’on trouvera par ces deux procédés, produits soit de
\[
\begin{tikzcd}
F \arrow[dr] & \\
 & i_*i^*F \\
e \arrow[ur] &
\end{tikzcd}
\quad ,
\]
soit de
\[
\begin{tikzcd}
 & F \\
j_!j^*F \arrow[ur] \arrow[dr] & \\
 & e
\end{tikzcd}
\quad ,
\]
devraient être \textit{canoniquement} isomorphes.

\noindent\msp{page 11}Si $\DD$ est un dérivateur, la catégorie
\[
\mathcal{A} = A_{\DD} = \DD(e) % edition reading verified: plain italic A_D is retained; this alphabet variation is not established to be an error
\]
joue un rôle important. On l’appelle la \textit{catégorie de base,} ou \textit{la catégorie fondamentale}, du dérivateur, ces objets sont les \textit{coefficients fondamentaux.} Pour $I$ quelconque dans $\Dia$, on a donc pour $F \in \DD(I)$, coefficient sur $I$ % [sic] ces objets -> ses objets, possessive intended
\[
I^\circ \longrightarrow \mathcal{A} \quad , \qquad i \longmapsto F_i
\]
un diagramme de type $I^\circ$ dans $\mathcal{A}$, fonctoriellement, \textit{i.e.}
\[
\DD(I)\longrightarrow
\uHom(I^\circ,\underset{\substack{\Vert\\\DD(e)}}{\mathcal A})
\]
% the vertical equality sign is printed under the letter \mathcal{A} of \uHom(I^\circ,\mathcal{A}), with \DD(e) beneath it
(\textit{cf.} page 8).

\noindent\textbf{Der}$_3$ : Pour tout $I \in \Dia$, ce foncteur est \textit{conservatif,}\footnote[4]{N.B. Mais ce foncteur n’a aucune envie d’être \textit{fidèle.}} \textit{i.e.} pour que
\[
F \longrightarrow G
\]
\printedpage{10}%
% head: A. Grothendieck : Dérivateurs, Ch. I, Édition des Universités de Montpellier II et Paris VII
dans $\DD(I)$ soit \texttt{[un]} isomorphisme, il faut et il suffit que pour tout $i \in I$,
\[
F_i \longrightarrow G_i
\]
le soit.

\noindent\textbf{Corollaire.} \textit{Le foncteur canonique (pour $I$, $J \in \Dia$)}
\[
\DD(I \times J) \xrightarrow{\ \theta\ } \uHom(J^\circ, \DD(I))
\]
\textit{est conservatif.}

En effet, si $F$ est dans $\DD(I \times J)$, on a pour $(i,j) \in I \times J$
\[
F_{i,j} \simeq \theta(F)(j)_i \quad .
\]

\noindent\msp{page 12}Intuitivement, on regarde un dérivateur $\DD$ comme “étant” la catégorie fondamentale $\mathcal{A}$, avec une “structure” supplémentaire, \texttt{[à]} savoir la donnée des “extensions” de $\mathcal{A}$ (comme catégorie des coefficients ponctuels) aux diagrammes $I \in \Dia$, de façon contravariante en $I$. Un coefficient $F$ dans $\DD(I) \defeq \mathcal{A}_I$ est regardé comme “étant” le diagramme $(F_i)_{i\in I}$ correspondant d’objets fondamentaux, avec une donnée supplémentaire de “globalisation”. Les $F_i$, ou “fibres” de $F$ en les $i \in I$, sont les caractéristiques locaux, du coefficient “global sur $I$” $F$. % [sic] caractéristiques locaux -> caractéristiques locales, agreement (feminine)

Tous les dérivateurs connus sont obtenus de la façon standard suivante. On part d’une “catégorie de modèles”
\[
\mathcal{M} \quad , \quad \Sigma
\]
munie d’un ensemble de flèches, appelées “équivalences”, et satisfaisant aux axiomes pertinents pour le passage à une catégorie des fractions
\[
\mathcal{A} = \mathcal{M}\Sigma^{-1}
\]
qui sera la “catégorie fondamentale”. On définit pour tout $I$ dans $\Dia$
\[
\begin{aligned}
\mathcal{M}_I &= \uHom(I^\circ,\mathcal{M}) \\
\Sigma_I &= \{u \in \mathsf{Fl}(\mathcal{M}_I) \mid \forall i \in I, u_i \in \Sigma\} \\
\mathcal{A}_I &= \mathcal{M}_I\Sigma_I^{-1}
\end{aligned}
\]

\noindent\msp{page 13}Il est clair que $\mathcal{A}_I \defeq \DD_{\mathcal{M},\Sigma}(I)$ dépend de $I$ de façon fonctorielle contravariante, et même de façon 2-fonctorielle.

Bien entendu, pour avoir vraiment un dérivateur, il faut vérifier les axiomes. Le plus fort est l’existence des $f_!$, $f_*$. Elle est prouvée par Thomason lorsque $(\mathcal{M},\Sigma)$ provient d’une “catégorie de modèles” au sens de Quillen. Ce n’est pas clair pour moi qu’alors $(\mathcal{M}_I,\Sigma_I)$ provienne également d’une catégorie de modèles (dont par exemple les fibrations \printedpage{11}%
% head: A. Grothendieck : Dérivateurs, Ch. I, Édition des Universités de Montpellier II et Paris VII
et les cofibrations seraient définies à partir de ces données dans $\mathcal{M}$, comme tantôt pour les équivalences).

Quand on prend pour $(\mathcal{M},\Sigma)$ une catégorie de modèles pour $\mathsf{Hot}$ (par exemple les ensembles semi-simpliciaux $\mathcal{M} = \uHom(\bDelta^\circ,\Ens)$, d’où $\mathcal{M}_I \simeq \uHom(\bDelta^\circ \times I^\circ,\Ens) \simeq \uHom(\bDelta^\circ,\underbrace{\uHom(I^\circ,\Ens)}_{\widehat{I}})$, catégorie des objets semi-simpliciaux dans le topos $\widehat{I}$, on trouve une description d’un dérivateur (non additif), qui jouera un rôle particulièrement important. Je le désigne par % Resolved: exactly two closing parentheses in the formula; the source's outer prose parenthesis remains unclosed.
\[
\mathsf{HOT} \quad ,
\]
sa catégorie des coefficients fondamentaux est $\mathsf{Hot}$, la catégorie homotopique standard.

\noindent\msp{page 14}La notion de “morphisme de Dérivateurs” est claire, c’est un morphisme entre deux 2-foncteurs stricts
\[
\Dia^\circ \longrightarrow \mathcal{C}at
\]
Donc
\[
u : \DD \longrightarrow \DD'
\]
est défini par la donnée, pour tout $I \in \Dia$, de
\[
u_I : \DD(I) \longrightarrow \DD'(I)
\]
avec commutativité (stricte?) pour $f : I \to J$
\[
\begin{tikzcd}
\DD(J) \arrow[r, "f^*_{\DD}"] \arrow[d, "u_J"'] & \DD(I) \arrow[d, "u_I"] \\
\DD'(J) \arrow[r, "f^*_{\DD'}"] & \DD'(I)
\end{tikzcd}
\quad . \tag*{(1)}
\]
On a
\[
\begin{tikzcd}[column sep=large]
\uHom(I,J)^\circ \arrow[r, "\DD_{I,J}"] \arrow[d, "\DD'_{I,J}"'] & \uHom(\DD(J),\DD(I)) \arrow[d, "\phi \mapsto u_I \circ \phi"] \\
\uHom(\DD'(J),\DD'(I)) \arrow[r, "\psi \mapsto \psi \circ u_J"'] & \uHom(\DD(J),\DD'(I))
\end{tikzcd}
\tag*{(2)}
\]
et la commutativité de (1) signifie la commutativité de (2) \textit{pour les objets}. Mais bien sûr on la veut aussi pour les flèches, \textit{i.e.} on veut la commutativité du diagramme (2).\footnote[5]{Donner l’exemple de morphisme de dérivateurs défini par un foncteur de catégories de modèles $(\mathcal{M},\Sigma) \to (\mathcal{M}',\Sigma')$.}

Sans doute on veut un peu plus sur les morphismes de dérivateurs - on veut une commutation aux $f_*$ et (?) aux $f_!$ (cette fois seulement à isomorphisme canonique près) \ldots\footnote[6]{Commutation à $f_*$ (resp. $f_!$) morphisme \textit{exact à gauche} (resp. \textit{exact à droite}). Les morphismes \textit{exacts} sont ceux à la fois exacts à gauche et à droite.}

\printedpage{12}
% head: A. Grothendieck : Dérivateurs, Ch. I, Édition des Universités de Montpellier II et Paris VII
\noindent\msp{page 15}Par exemple, considérons le dérivateur
\[
\mathsf{HOTAB}
\]
dont la catégorie fondamentale est $\mathsf{Hotab}$, la sous-catégorie pleine de la catégorie dérivée $\mathsf{Der}(\mathsf{Ab})$ de $\mathsf{Ab}$, formée des complexes à cohomologie nulle en degré $> 0$. On peut décrire $\mathsf{HOTAB}$ par la catégorie de modèles $\uHom(\bDelta^\circ,\mathsf{Ab})$ des groupes abéliens semi-simpliciaux, avec comme équivalences les quasi-isomorphismes au sens habituel. Le “foncteur abélianisation” canonique
\[
\mathsf{Hot} \longrightarrow \mathsf{Hotab}
\]
se prolonge en un morphisme de dérivateurs
\[
\mathsf{HOT} \longrightarrow \mathsf{HOTAB} \quad .
\]
La question de sa compatibilité aux $f_*$ ne semble pas tautologique. Elle impliquerait, par exemple, que si
\[
\begin{tikzcd}
X \arrow[d, "f"'] & X' \arrow[l, "g"'] \arrow[d, "f'"] \\
Y & Y' \arrow[l, "j"]
\end{tikzcd}
\]
est un diagramme cartésien dans $\mathcal{M} = \uHom(\bDelta^\circ,\Ens)$ avec $f$ une fibration (de Thom dans le cadre des espaces topologiques, avec $f$ par exemple fibration localement triviale) alors le diagramme de $\mathsf{Hotab}$
\[
\begin{tikzcd}
\mathbb{R}\Gamma_!(X) \arrow[d] & \mathbb{R}\Gamma_!(X') \arrow[l] \arrow[d] \\
\mathbb{R}\Gamma_!(Y) & \mathbb{R}\Gamma_!(Y') \arrow[l]
\end{tikzcd}
\]
est $\mathsf{hot}$-cartésien\footnote[7]{C’est faux déjà si $Y$ est l’espace ponctuel : trivialement le $\mathbb{R}\Gamma_!$ ne transforme pas produits en sommes! Ce foncteur est sûrement \texttt{[pas]} “exact à gauche” (comme morphisme de dérivateurs).}. Et de même pour une $\mathsf{hot}$-somme amalgamée (carrée $\mathsf{hot}$-cocartésien dans $\mathsf{Hot}$). % [sic] carrée -> carré, agreement (masculine noun)

\noindent\msp{page 16}Si c’est vrai, et pas connu de plus, ce serait quand même scandaleux!

En procédant comme pour les morphismes de changements de base pour les $Rf_*$, on trouve en tout cas un morphisme canonique, pour $f : I \to J$
\[
u_Jf_*^{\DD} \longrightarrow f_*^{\DD'}u_I
\]
\printedpage{13}%
% head: A. Grothendieck : Dérivateurs, Ch. I, Édition des Universités de Montpellier II et Paris VII
et de même
\[
f_!^{\DD'}u_I \longrightarrow u_Jf_!^{\DD} \quad ,
\]
et la question est de savoir si c’est des isomorphismes. Questions analogues pour les $f^!$, pour $f$ immersion fermée (ou pour $f^?$, pour $f$ immersion ouverte). Comme la notion de morphisme de dérivateurs devrait quand même “coller” à la notion bien connue de morphisme (triangulé) de catégories triangulées\footnote[8]{Non! Le cas des catégories triangulées est “misleading”, car pour celles-ci l’exactitude à gauche et à droite revient au même.}, donc être exact en un sens convenable, cela montre bien que l’on doit exiger la commutation d’un morphisme de dérivateurs aux $f_*$ et aux $f_!$.\footnote[9]{N.B. On conçoit qu’un foncteur entre dérivateurs puisse être exact à gauche ou exact à droite ou les deux.} Je devrais demander à Thomason s’il a regardé, pour un morphisme
\[
\mathcal{M} \longrightarrow \mathcal{M}'
\]
entre catégories de modèles, qui soit disons \textit{exact à gauche}, et “compatible” avec les équivalences, fibrations (et cofibrations (???)).

\noindent\textsc{Notation} : $\mathsf{DER}$ pour la 2-catégorie des \textit{dérivateurs.}

\noindent\msp{page 17}\textbf{Dualité}. Si $\DD$ est un dérivateur, on va définir un dérivateur \textit{dual} $\DD^\circ$ par
\[
\DD^\circ(I) \defeq \DD(I^\circ)^\circ \quad .
\]
Les coefficients sur $I$ de type $\DD^\circ$ sont les coefficient sur la catégorie opposée $I^\circ$ à $I$ de type $\DD$, et en tant que catégorie c’est la catégorie \textit{opposée} de celle des coefficients de type $\DD$ sur $I^\circ$. Pour définir la loi 2-fonctorielle, on note que pour $I$, $J$ dans $\Dia$ on a % [sic] les coefficient -> les coefficients, plural agreement
\[
\begin{tikzcd}[column sep=large]
\uHom(I,J)^\circ \arrow[r, "\DD_{I,J}"] \arrow[d, "\wr"', "\text{dualité}"] & \uHom(\DD(J),\DD(I)) \arrow[d, "\wr"', "\text{dualité}"] \\
\uHom(I^\circ,J^\circ) \arrow[r, "\theta"] & \uHom(\DD(J)^\circ,\DD(I)^\circ)^\circ
\end{tikzcd}
\quad ,
\]
où $\theta$ est l’unique foncteur rendant le diagramme commutatif. Alors $\theta^\circ$est un foncteur
\[
\uHom(I^\circ,J^\circ)^\circ \longrightarrow \uHom(\DD(J)^\circ,\DD(I)^\circ) \quad ,
\]
c’est ça par définition $\DD^\circ_{I,J}$ \texttt{[plutôt} $\DD^\circ_{I^\circ,J^\circ}$\texttt{]}.

Si $\DD$ est défini par $(\mathcal{M},\Sigma)$, alors $\DD^\circ$ par $(\mathcal{M}^\circ,\Sigma^\circ)$, puisque
\[
(\mathcal{M}^\circ)_I \defeq \uHom(I^\circ,\mathcal{M}^\circ) \simeq \uHom(I,\mathcal{M})^\circ \simeq (\mathcal{M}_{I^\circ})^\circ
\]
% foot (pages 9--13): Version 8 juillet 2001


\endgroup
% END UNIT Der-ChI-234/Der-ChI-234_p011-015.tex
% BEGIN UNIT Der-ChI-234/Der-ChI-234_p016-020.tex
\begingroup\makeatletter%
\newcommand{\printedpage}[1]{\marginpar{\footnotesize p.~#1}}%
\newcommand{\msp}[1]{\marginpar{\footnotesize [#1]}}%
\newcommand{\D}{\mathbb{D}}%
\newcommand{\uHom}{\underline{\mathsf{Hom}}}%
\newcommand{\sHom}{\mathsf{Hom}}%
\newcommand{\Dia}{\mathsf{Dia}}%
\newcommand{\DER}{\mathsf{DER}}%
\newcommand{\id}{\mathsf{id}}%
\newcommand{\Cat}{\mathcal{C}at}%
\newcommand{\Ens}{\mathcal{E}ns}%
\newcommand{\sIm}{\mathsf{Im}}%
\makeatother\let\printedpage\derpp\let\msp\dermsp%
\let\mspin\DERmspin%


\printedpage{14}
% head: A. Grothendieck : Dérivateurs, Ch. I, Édition des Universités de Montpellier II et Paris VII
% foot: Version 8 juillet 2001

et $\Sigma^\circ{}_I$ est déduit de $\Sigma_{I^\circ}$, d'où
\[
(\mathcal{M}^\circ)_I(\Sigma^\circ{}_I)^{-1} \simeq (\mathcal{M}_{I^\circ}\Sigma_{I^\circ}^{-1})^\circ = \D(I^\circ)^\circ \quad .
\]

\msp{page 18}

\textbf{Philosophie.} Chaque fois qu'on a un foncteur
\[
\mathcal{A} \xrightarrow{u_0} \mathcal{A}'
\]
entre deux catégories qui peuvent s'interpréter naturellement comme les catégories fondamentales de deux dérivateurs $\D$, $\D'$, il se pose la question d'interpréter $u_0$ comme provenant d'un morphisme de dérivateurs
\[
\D \xrightarrow{u} \D' \quad .
\]
Celui-ci est, en pratique, toujours déterminé canoniquement, à isomorphisme canonique près. (Mais je doute pourtant que le foncteur restriction aux catégories fondamentales
\[
\uHom(\D,\D') \longrightarrow \uHom(\mathcal{A},\mathcal{A}')
\]
soit pleinement fidèle \dots\ )

On va appliquer cette philosophie aux foncteurs $f^*$, $f_!$, $f_*$ (et éventuellement $f^!$, $f^?$), pour $f : I \to J$, donc pour les foncteurs entre $\D(I)$ et $\D(J)$. Cela nous amène donc à regarder $\D(I)$ et $\D(J)$ comme des catégories fondamentales de dérivateurs convenables.

\textbf{Dérivateur induit} par un dérivateur $\D$, sur une catégorie d'indices $I \in \Dia$. C'est un dérivateur $\D_I$ défini par
\[
\D_I(J) = \D(I \times J) \quad .
\]
Le foncteur associé à $\D_I$, pour deux catégories $J$, $J' \in \Dia$

\msp{page 19}
\[
\uHom(J,J')^\circ \xrightarrow{(\D_I)_{J,J'}} \uHom(\D_I(J'),\D_I(J)) = \uHom(\D(I\times J'),\D(I\times J))
\]
est le composé
\[
\uHom(J,J')^\circ \longrightarrow \uHom(I\times J, I\times J')^\circ \xrightarrow{\D_{I\times J,I\times J'}} \uHom(\D(I\times J'),\D(I\times J)) \quad ,
\]
où la première flèche est le foncteur
\[
f \longmapsto \id_I \times f \quad .
\]
En d'autres termes, $\D_I$ est le composé
\[
\begin{aligned}
&\Dia^\circ \longrightarrow \Dia^\circ \xrightarrow{\D} \Cat \\
&J \longmapsto I \times J \qquad .
\end{aligned}
\]
On a
\[
f^*_{\D_I} = (\id \times f)^*_{\D} \quad .
\]
Si les $f_*$ (resp. les $f_!$) existent pour $\D$, ils existent pour $\D_I$. (De même pour toute espèce de multiadjoints.) Comme $J \mapsto I \times J$ transforme immersion fermée (resp. ouverte) en itou, l'axiome \textbf{Der}$_2$ (plus précisément, chacun de ces ingrédients, \textit{cf.} annotations marginales p.~7) est stable par passage de $\D$ à $\D_I$ . De même pour \textbf{Der}$_3$. (Pour chaque nouvel axiome qu'on dégagera, il faudra vérifier la stabilité.)

\printedpage{15}
% head: A. Grothendieck : Dérivateurs, Ch. I, Édition des Universités de Montpellier II et Paris VII
% foot: Version 8 juillet 2001

\msp{page 20}

Soit
\[
f : I \longrightarrow I' \quad .
\]
Considérons entre $\D(I)$, $\D(I')$ les foncteurs
\[
\begin{tikzcd}[column sep=huge]
\D(I) \arrow[r, shift left=3ex, "f_!"] \arrow[r, shift right=3ex, "f_*"'] & \D(I') \arrow[l, "f^*"']
\end{tikzcd}
\]
(sous réserve d'existence des $f_*$ resp. $f_!$). Peut on les interpréter en termes de morphismes de dérivateurs
\[
\begin{tikzcd}[column sep=huge]
\D_I \arrow[r, shift left=3ex, "\overline{f_!}"] \arrow[r, shift right=3ex, "\overline{f_*}"'] & \D_{I'} \arrow[l, "\overline{f^*}"']
\end{tikzcd}
\qquad ?
\]
Donc il faut définir, pour $J$ dans $\Dia$, des foncteurs correspondants entre $\D_I(J)$ et $\D_{I'}(J)$, \textit{i.e.} entre $\D(I \times J)$ et $\D(I' \times J)$. Or on a bien
\[
f \times \id_J : I \times J \longrightarrow I' \times J
\]
et on posera (sous réserve d'existence des deux derniers)
\[
\left\{
\begin{aligned}
&\overline{f^*}(J) = (f \times \id_J)^*_{\D} \\[1ex]
&\overline{f_!}(J) = (f \times \id_J)^{\D}_! \\[1ex]
&\overline{f_*}(J) = (f \times \id_J)^{\D}_* \quad .
\end{aligned}
\right.
\]

\msp{page 21}

Autrement dit, on peut considérer
\[
\Dia \ni I \longmapsto \D_I \quad , \qquad \D \in \DER \quad ,
\]
comme fonctoriel en $I$ de trois façons (une contravariante, les deux autres covariantes), suivant qu'on utilise les $f^*$, les $f_!$ (s'ils existent) ou les $f_*$ (s'ils existent). On désignera par
\[
f^{\D}_! \quad , \qquad f^*_{\D} \quad , \qquad f^{\D}_*
\]
indifféremment les foncteurs entre $\D(I)$ et $\D(I')$ ou (de façon plus fine) entre $\D_I$ et $\D_{I'}$, dont les précédents se déduisent en passant aux catégories fondamentales correspondantes. De même pour $f^!$ ou pour $f^?$ sous réserve d'existence. Par exemple, si on suppose que les $f^!$ existent pour $\D$, quel que soit $f$ immersion fermée, alors on trouve pour un tel $f$ un morphisme de dérivateurs
\[
\D_I \xleftarrow{f^!_{\D}} \D_{I'} \quad .
\]
Il s'impose de regarder si tous ces morphismes de dérivateurs sont exacts, \textit{i.e.} si pour
\[
g : J \longrightarrow J'
\]
dans $\Dia$, ils commutent à $g_!$, $g_*$.

\printedpage{16}
% head: A. Grothendieck : Dérivateurs, Ch. I, Édition des Universités de Montpellier II et Paris VII
% foot: Version 8 juillet 2001

\msp{page 22 vide}

\msp{page 23}

Pour expliciter cette question, on est amenés à regarder le diagramme de catégories
\[
\begin{tikzcd}[column sep=huge, row sep=huge]
I \times J \arrow[r, "G = \id_I \times g"] \arrow[d, "F = f \times \id_J"'] & I \times J' \arrow[d, "F' = f \times \id_{J'}"] \\
I' \times J \arrow[r, "G' = \id_{I'} \times g"] & I' \times J'
\end{tikzcd}
\qquad ,
\]
qui est d'ailleurs un carré cartésien dans $\Cat$. La commutativité de $f^*_{\D}$ aux $g_*$ s'explicite alors par la commutativité dans
\[
\begin{tikzcd}[column sep=huge, row sep=huge]
\D(I \times J) \arrow[r, "G_*", ""{name=T, below, near end}] & \D(I \times J') \\
\D(I' \times J) \arrow[u, "F^*"] \arrow[r, "G'_*"] & \D(I' \times J') \arrow[u, "{F'}^*"', ""{name=R, left}]
\arrow[dashed, from=R, to=T, "\text{iso?}"]
\end{tikzcd}
\qquad .
\]
Cette commutativité sera assurée par des axiomes ultérieures sur les dérivateurs. (cf. plus % [sic] ultérieures -> ultérieurs, accord avec « axiomes »
bas page 59). Ainsi, $f^*_{\D}$ est exact à droite \texttt{[plutôt à gauche]}, et dualement, il doit être aussi exact à gauche \texttt{[plutôt à droite]}. Donc $f^*_D : \D_{I'} \to \D_I$ est \textit{exact.} Considérons $f^{\D}_*$, % [sic] f^*_D -> f^*_{\mathbb{D}}, indice imprimé en D italique au lieu de \mathbb{D}
cette fois c'est la commutativité de
\[
\begin{tikzcd}[column sep=huge, row sep=huge]
\D(I \times J) \arrow[r, "G_*"] \arrow[d, "F_*"'] & \D(I \times J') \arrow[d, "F'_*"] \\
\D(I' \times J) \arrow[r, "G'_*"] & \D(I' \times J')
\end{tikzcd}
\qquad ,
\]
c'est trivial (à isomorphisme canonique près). Ainsi, on
\msp{page 24}
trouve que $f^{\D}_* : \D_I \to \D_{I'}$ est \textit{exact à gauche}, et dualement $f^{\D}_!$ est \textit{exact à droite.} Mais est-ce que $f^{\D}_*$ est aussi exact à \textit{droite}? Cela signifie donc commutativité de
\[
\begin{tikzcd}[column sep=huge, row sep=huge]
\D(I \times J) \arrow[r, "G_!"] \arrow[d, "F_*"'] & \D(I \times J') \arrow[d, "F'_*", ""{name=R2, left}] \\
\D(I' \times J) \arrow[r, "G'_!", ""{name=B, above, near end}] & \D(I' \times J')
\arrow[dashed, from=B, to=R2]
\end{tikzcd}
\qquad .
\]

\printedpage{17}
% head: A. Grothendieck : Dérivateurs, Ch. I, Édition des Universités de Montpellier II et Paris VII
% foot: Version 8 juillet 2001

On arrive bien à définir
\[
G'_!F_* \xrightarrow{?} F'_*G_!
\]
ou ce qui revient au même par adjonction
\[
F_* \xrightarrow{?} {G'}^*F'_*G_!
\]
en utilisant le morphisme de changement de base
\[
{G'}^*F'_* \xrightarrow{\overset{\alpha}{\sim}} F_*G^*
\]
qui est \texttt{[un]} \textit{isomorphisme} (on l'a vu), d'où
\[
\begin{tikzcd}[column sep=large, row sep=normal]
& {G'}^*F'_*G_! \arrow[dd, "\wr"', "\alpha * G_!"] \\
F_* \arrow[ur, "?"] \arrow[dr, "F_* * \beta"'] & \\
& F_*G^*G_!
\end{tikzcd}
\]
où $\beta : \id \to G^*G_!$ est le morphisme d'adjonction, d'où le morphisme voulu. (On peut aussi chercher ${F'}^*(G'_!F_*) \to G_!$, et utiliser la commutation $G_!F^* \overset{\sim}{\to} {F'}^*G'_!$, \textit{i.e.} l'exactitude à \textit{droite} de $f^*_{\D}$.)

\msp{page 25}

Mais ce morphisme est-il un isomorphisme? Je pense que dans le cas additif ça peut être raisonnable, mais dans le cas non additif ça n'a aucune raison d'être vrai. On voit en tout cas que si pour tout $f : I \to I'$ dans $\Dia$, $f^{\D}_* : \D_I \to \D_{I'}$ est (non seulement exact à gauche mais aussi) exact à droite, alors pour tout $g : J \to J'$ dans $\Dia$, $g^{\D}_!$ est (non seulement exact à droite mais même) exacte à gauche. % [sic] exacte -> exact, accord avec « $g^{\D}_!$ est (non seulement exact\dots) »

Prenons le dérivateur des préfaisceaux d'ensembles sur des $I$ variables, sa catégorie fondamentale est $\Ens$. C'est un dérivateur non ponctué. Considérons $I \to I'$ une immersion fermée non surjective, $f_* : \D(I) \to \D(I')$. Si $f^{\D}_*$ était exact à droite $f_*$ devrait transformer objet initial de $\D(I)$ en objet initial de $\D(I')$, sauf erreur (\textit{cf.} plus bas). Or il n'en est rien, puisque pour tout $F \in \sIm\bigl(f^{\D}_* : \D(I) \to \D(I')\bigr)$ la restriction de $F$ à $I' \setminus I$ est le faisceau final $e_U = (i \mapsto e)$, et non le faisceau initial $\emptyset_U = (i \mapsto \emptyset)$.

On va prouver ceci :

\textbf{Proposition.} \textit{Si $\D$ est un dérivateur à gauche (à droite), alors les produits finis (sommes finies) existent dans les $\D(I)$. Tout morphisme $\varphi : \D \to \D'$ de dérivateurs qui est exact à gauche (à droite) commute aux dits produits (sommes), chaque $I$ donnant $\varphi(I) : \D(I) \to \D'(I)$. En particulier, il transforme objet final (initial) en itou.}

\printedpage{18}
% head: A. Grothendieck : Dérivateurs, Ch. I, Édition des Universités de Montpellier II et Paris VII
% foot: Version 8 juillet 2001

\msp{page 26}

Dire que les produits finis indexés par l'ensemble fini $I$ existent dans la catégorie $\mathcal{A} = \D(e)$ (disons) équivaut à dire que le foncteur diagonal
\[
\mathcal{A} \xrightarrow{\Delta} \mathcal{A}^I
\]
admet un adjoint à gauche. Or si
\[
f : I \to e
\]
est la projection de $I$ (considéré comme ensemble ordonné discret) sur l'ensemble ordonné ponctuel $e$, le foncteur $\Delta$ s'identifie à $f^* : \D(e) \to \D(I)$ (grâce à l'axiome \textbf{Der}$_2$ sur le comportement du dérivateur pour une somme discrète de catégories d'indices). Donc l'existence du produit revient à celle de l'adjoint à droite $f_*$ de $f^*$, OK pour dérivateur à gauche. Et un homomorphisme exact à gauche de dérivateurs à gauche, comme il commute aux $f_*$, commute aux produits finis. Itou pour les sommes finies, et les dérivateurs à droite.

L'exemple précédent utilise le fait que le dérivateur envisagé n'est pas ponctué. Mais même pour un dérivateur ponctué $\D$, un morphisme $f^{\D}_* : \D_I \to \D_{I'}$ n'est pas forcément exact à droite. Si on prend le dérivateur associé à une catégorie $\mathcal{M}$ quelconque (avec $\Sigma = \emptyset$),
\[
\D_{\mathcal{M}}(I) = \uHom(I^\circ,\mathcal{M}) \overset{\text{déf}}{=} \mathcal{M}(I)
\]
alors le dérivateur induit
\[
(\D_{\mathcal{M}})_I = \{J \mapsto \D_{\mathcal{M}}(I \times J) = \uHom((I \times J)^\circ,\mathcal{M}) \simeq \uHom(J^\circ,\uHom(I^\circ,\mathcal{M}))\}
\]

\msp{page 27}

s'identifie à
\[
\D_{\uHom(I^\circ,\mathcal{M})} = \D_{\mathcal{M}(I)}
\]
\textit{i.e.}
\[
(\D_{\mathcal{M}})_I = \D_{\mathcal{M}(I)}
\]
N.B. Ici $\mathcal{M}$ et $\D_{\mathcal{M}}$ se déterminent mutuellement et le foncteur
\[
\mathcal{M} \longmapsto \D_{\mathcal{M}}
\]
\[
\Cat \longrightarrow \DER
\]
est 2-fidèle, \textit{i.e.} si $\mathcal{M}$, $\mathcal{M}'$ sont deux catégories
\[
\uHom(\mathcal{M},\mathcal{M}') \xrightarrow{\simeq} \sHom(\D_{\mathcal{M}},\D_{\mathcal{M}'})
\]
est une équivalence de catégories. D'autre part, on vérifie tout de suite que $u : \mathcal{M} \to \mathcal{M}'$ est exact à gauche (ou à droite) si et seulement si $\D_u : \D_{\mathcal{M}} \to \D_{\mathcal{M}'}$ l'est.

Ainsi, pour voir si un morphisme de dérivateurs de la forme
\[
\begin{tikzcd}[column sep=large, row sep=normal]
f^{\D}_* : &[-2.5em] (\D_{\mathcal{M}})_I \arrow[r] \arrow[d, equal] & (\D_{\mathcal{M}})_{I'} \arrow[d, equal] \\
& \D_{\mathcal{M}(I)} & \D_{\mathcal{M}(I')}
\end{tikzcd}
\]
est exact à \textit{droite,} il revient au même de dire


\endgroup
% END UNIT Der-ChI-234/Der-ChI-234_p016-020.tex
% BEGIN UNIT Der-ChI-234/Der-ChI-234_p021-025.tex
\begingroup\makeatletter%
\newcommand{\printedpage}[1]{\marginpar{\footnotesize p.~#1}}%
\newcommand{\msp}[1]{\marginpar{\footnotesize [#1]}}%
\newcommand{\sDia}{\mathsf{Dia}}%
\newcommand{\sOb}{\mathsf{Ob}}%
\newcommand{\sid}{\mathsf{id}}%
\newcommand{\DD}{\mathbb{D}}%
\makeatother\let\printedpage\derpp\let\msp\dermsp%
\let\mspin\DERmspin%


\printedpage{19}
% head: A. Grothendieck : Dérivateurs, Ch. I, Édition des Universités de Montpellier II et Paris VII
% foot: Version 8 juillet 2001
\msp{page 28}
que le foncteur pour les catégories fondamentales associées
\[ \mathcal{M}(I) \xrightarrow{f^{\mathcal{M}}_*} \mathcal{M}(I') \]
est exact à droite. Or ici c'est un foncteur image directe de préfaisceaux ordinaires, et ce foncteur n'a aucune envie d'être exact à \textit{droite.} (Même si $\mathcal{M}$ est \textit{abélienne}, d'ailleurs!) Prenons par exemple pour $f : I \hookrightarrow I'$ une immersion ouverte avec $I' = \boldsymbol{\Delta}_1 \times \boldsymbol{\Delta}_1$, $I = I' - \{c\}$, $c$ objet final de $I'$
\[
I' = \begin{tikzcd}
c & b \arrow[l] \\
a \arrow[u] & 0 \arrow[l] \arrow[u]
\end{tikzcd}
\qquad\qquad
I = \begin{tikzcd}
 & b \\
a & 0 \arrow[l] \arrow[u]
\end{tikzcd}
\]
Le foncteur $f_*$ est celui qui associe à un diagramme
\[
F = \left\{
\begin{tikzcd}
 & F_b \arrow[d] \\
F_a \arrow[r] & F_0
\end{tikzcd}
\right\}
\]
le carré cartésien complété
\[
\overline{F} = \left\{
\begin{tikzcd}
F_a \times_{F_0} F_b \arrow[r] \arrow[d] & F_b \arrow[d, "\beta"] \\
F_a \arrow[r, "\alpha"] & F_0
\end{tikzcd}
\right\}
\]
Or si $F \mapsto \overline{F}$ était exact à droite, $F \mapsto F_a \times_{F_0} F_b$ le serait aussi. Or $F_a \times_{F_0} F_b = K(F) = \mathsf{Ker}(F_a \times F_b \xrightarrow{\alpha pr_1 - \beta pr_2} F_0)$, en tant que foncteur additif en $F$, n'est pas exact à droite, par exemple il ne transforme pas nécessairement épimorphisme en épimorphisme. Prenons par exemple
\[
F' = \left\{
\begin{tikzcd}
 & F_b \arrow[d] \\
F_a \arrow[r] & 0
\end{tikzcd}
\right\}
\]
et l'homomorphisme évident $F \to F'$ qui est l'identité sur
\msp{page 29}
$F_a$ et sur $F_b$. C'est un épimorphisme et l'homomorphisme correspondant $\overline{F} \to \overline{F'}$ est l'inclusion
\[ F_a \times_{F_0} F_b \longrightarrow F_a \times F_b \]
qui n'est un épimorphisme que si $\alpha = \beta = 0$.

\printedpage{20}
% head: A. Grothendieck : Dérivateurs, Ch. I, Édition des Universités de Montpellier II et Paris VII
% foot: Version 8 juillet 2001
\textbf{Perplexité} : On devrait pouvoir définir une notion de dérivateur (additif) “\textit{triangulé}” (qui correspondrait à l'intuition des catégories triangulées; c'est un dérivateur satisfaisant à des axiomes assez forts), de sorte que pour les morphismes entre ceux-ci l'exactitude à gauche et l'exactitude à droite soient \textit{équivalentes.}

\textbf{“Théorème de changement de base”} pour dérivateurs. Soit $\DD$ dérivateur à gauche (existence de $f_*$) et considérons
\[ f : I \longrightarrow J \qquad \text{dans } \sDia \quad , \]
soit $j_0 \in \sOb\ J$, on se propose de “calculer” le foncteur
\[ F \mapsto f_*(F)_{j_0} = \alpha^*_{j_0}(f_*(F)) \quad , \]
où
\[ \alpha_{j_0} : e \longrightarrow J \]
est l'inclusion de la catégorie ponctuelle $e$ dans $J$, de valeur $j_0$.

\msp{page 30}
Considérons le diagramme
\[
\begin{tikzcd}[row sep=large, column sep=large]
I/j_0 \arrow[r, "q"] \arrow[d, "f'"'] & I \arrow[d, "f"] \\
e \arrow[r, "p = \alpha_{j_0}"'] & J
\arrow[from=1-2, to=2-1, dashed, shorten <=14pt, shorten >=14pt, "\gamma"']
\end{tikzcd}
\quad ,
\tag*{(1)}
\]
où $I/j_0$ est la catégorie $I \times_J J/j_0$ des couples
\[ \{(i,\varphi) \mid i \in \sOb\ I,\ \varphi : f(i) \to j_0\} \]
et $q$ est le foncteur $(i,\varphi) \mapsto i$. Ce diagramme n'est pas commutatif, mais il y a un foncteur \texttt{[plutôt morphisme de foncteurs]} canonique
\[ fq \xrightarrow{\gamma} pf' \quad , \]
où $fq$ est le foncteur $(i,\varphi) \mapsto f(i)$ et $pf'$ le foncteur constant $(i,\varphi) \mapsto j_0$ et où $\gamma(i,\varphi) : f(i) \to j_0$ est défini par $\gamma(i,\varphi) = \varphi$.

Plus généralement, considérons un diagramme dans $\sDia$
\[
\begin{tikzcd}[row sep=large, column sep=large]
I' \arrow[r, "q"] \arrow[d, "f'"'] & I \arrow[d, "f"] \\
J' \arrow[r, "p"'] & J
\arrow[from=1-2, to=2-1, dashed, shorten <=14pt, shorten >=14pt, "\gamma"']
\end{tikzcd}
\quad ,
\tag*{(2)}
\]
\printedpage{21}
% head: A. Grothendieck : Dérivateurs, Ch. I, Édition des Universités de Montpellier II et Paris VII
% foot: Version 8 juillet 2001
avec $\gamma : fq \to pf'$. On définit alors un \textit{morphisme de changement de base pour les images directes} relatif à ce diagramme
\[ p^* f_* \xrightarrow{c_\gamma} f'_* q^* \]
\msp{page 31}
de deux façons, qui donnent le même résultat (histoire générale de foncteurs adjoints à expliciter une bonne fois, en termes d'un diagramme de foncteurs
\[
\begin{tikzcd}[row sep=large, column sep=large]
X' & X \arrow[l, "q^*"'] \\
Y' \arrow[u, "{f'}^*"] & Y \arrow[l, "p^*"'] \arrow[u, "f^*"']
\arrow[from=2-1, to=1-2, dashed, shorten <=14pt, shorten >=6pt]
\end{tikzcd}
\]
en supposant que les foncteurs adjoints à gauche $f'_*$, $f_*$ des deux foncteurs verticaux existent). Un cas intéressant est celui où
\[ pf' = fq \quad , \qquad \gamma = \sid_{fq} \]
et où on pose $c_\gamma = c$.

Ici
\[ p^* f_* \xrightarrow{?} f'_* q^* \]
revient à
\[
\begin{tikzcd}[column sep=small]
f_* \arrow[r, "?"] & p_* f'_* q^* \simeq & (pf')_* q^* \\
 & & (fq)_* q^* \simeq (f_* q_*) q^* = f_*(q_* q^*) \quad . \arrow[u, "\gamma_* * q^*"']
\end{tikzcd}
\]
Or on a
\[ \sid \xrightarrow{\alpha} q_* q^* \qquad\qquad \text{(adjonction)} \quad , \]
d'où
\[ f_* \xrightarrow{f_* * \alpha} f_*(q_* q^*) \]
et en composant avec $\gamma * q^*$, on trouve ce qu'on veut. Et on peut au lieu de ceci chercher
\[
\begin{tikzcd}
{f'}^*(p^* f_*) \arrow[r, "?"] \arrow[d, equal] & q^* \\
(pf')^* f_* \arrow[d, "\gamma^* * f_*"] & \\
(fq)^* f_* \simeq q^*(f^* f_*) \arrow[r, "q^* * \gamma'"] & q^*
\end{tikzcd}
\quad ,
\]
où $\gamma' : f^* f_* \to \sid$ est l'adjonction.

\printedpage{22}
% head: A. Grothendieck : Dérivateurs, Ch. I, Édition des Universités de Montpellier II et Paris VII
% foot: Version 8 juillet 2001
\msp{page 32}
\textbf{Der}$_4$ : Le morphisme de changement de base pour images directes $f_*$ associé à un diagramme (1) est toujours un isomorphisme (dérivateurs à gauche). Condition duale pour dérivateurs à droite et le changement de base pour images directes $f_!$, pour le diagramme
\[
\begin{tikzcd}[row sep=large, column sep=large]
j_0 \backslash I \arrow[r] \arrow[d] & I \arrow[d, "f"] \\
e \arrow[r, "p = \alpha_{j_0}"'] & J
\arrow[from=2-1, to=1-2, dashed, shorten <=22pt, shorten >=8pt, "\gamma"]
\end{tikzcd}
\tag*{(1$'$)}
\]
associé à $f : I \to J$ dans $\sDia$, et $j_0 \in \sOb\ J$.

\textbf{Der}$_5$\footnote[10]{\underline{Totalement} canulé, devient raisonnable que pour coefficients constants sur $F$.} : Soit $f : I \to J$ une \textit{hot-équivalence stricte, i.e.} $\forall j_0 \in \sOb\ J$, $I/j_0$ est contractile. Alors $f^*$ est pleinement fidèle, \textit{i.e.} le morphisme d'adjonction
\[ \sid_{\DD(J)} \longrightarrow f_* f^* \]
est un isomorphisme ($\DD$ dérivateur à gauche). Dualement, si $f$ est une équivalence co-stricte, \textit{i.e.} $f^\circ : I^\circ \to J^\circ$ est une équivalence stricte, \textit{i.e.} $\forall j_0 \in \sOb\ J$, $j_0 \backslash I$ est contractile, alors (si $\DD$ est un dérivateur à droite) le foncteur $f^*$ est pleinement fidèle, \textit{i.e.} le morphisme d'adjonction % Resolved: the edition line-break gives co-stricte; retain that spelling.
\[ f_! f^* \longrightarrow \sid_{\DD(J)} \]
est un isomorphisme.

N.B. Donc les foncteurs $f_*$ (resp. $f_!$) sont des foncteurs de localisation.

\textbf{N.B.} Les pages 33 à 40 ont été suprimées, car canulées. % [sic] suprimées -> supprimées, spelling

\msp{page 41}
Je vais reprendre la terminologie et les axiomes relatifs aux dérivateurs.

Un \textit{prédérivateur} est un 2-foncteur
\[ \sDia^\circ \xrightarrow{\DD} \mathcal{C}at \quad , \]
où $\sDia$ est une 2-sous-catégorie pleine de $\mathcal{C}at$.\footnote[11]{$\sDia$ \textit{domaine} du dérivateur. Les objets sont les \textit{catégories d'indices} pour $\DD$. Les $\DD(I)$ sont les catégories de coefficients pour $\DD$, les objets de $\DD(I)$ les “coeficients de type $\DD$ sur $I$”. % [sic] coeficients -> coefficients, spelling
} On suppose tout au moins que $\sDia$ contient les \textit{ensembles ordonnés finis,} que $\sDia$ est stable par $\varprojlim$ finies, par sommes finies. Qu'avec toute catégorie $C$ et toute sous-catégorie strictement pleine ouverte (resp. fermée) elle contienne la sous-catégorie complémentaire.\footnote[12]{Notation $f^*$ : image inverse des coefficients.}

\printedpage{23}
% head: A. Grothendieck : Dérivateurs, Ch. I, Édition des Universités de Montpellier II et Paris VII
% foot: Version 8 juillet 2001
\msp{page 42}
On définit le prédérivateur $\DD_{\mathcal{M},\Sigma}$ associé à un système
\[ (\mathcal{M},\Sigma) \qquad\qquad \mathcal{M} \text{ Catégorie,} \quad \Sigma \subset \mathsf{Fl}(\mathcal{M}) \quad . \] % Resolved: Fl abbreviates fleches, with lowercase l; retain Fl.
Si $\Sigma = \emptyset$, on écrit $\DD_{\mathcal{M}}$.

Les prédérivateurs forment une 2- catégorie $\mathsf{PREDER}$, et
\[ (\mathcal{M},\Sigma) \longmapsto \DD_{\mathcal{M},\Sigma} \]
est un 2-foncteur
\[ \mathsf{PREMOD}
\begin{pmatrix}
\text{“Catégorie des catégories} \\
\text{de modèles” (faibles)}
\end{pmatrix}
\longrightarrow \mathsf{PREDER} \quad . \]
Pour $\DD$ donné, on définit le dérivateur induit ($I \in \sDia$)
\[ \DD_I = \big(J \longmapsto \DD(I \times J)\big) \quad , \]
dépendant de $I$ de façon contravariante, d'où un 2-foncteur
\[ \begin{aligned}
\sDia^\circ &\longrightarrow \mathsf{PREDER} \\
I &\longmapsto \DD_I
\end{aligned} \]
ou encore un bi-2-foncteur
\[ \begin{aligned}
\mathsf{PREDER} \times \sDia^\circ &\longrightarrow \mathsf{PREDER} \\
(\DD, I) &\longmapsto \DD_I \quad .
\end{aligned} \]
Si $\DD$ est défini par $(\mathcal{M},\Sigma)$, alors $\DD_I$ est défini par $(\mathcal{M}(I),\Sigma(I))$ :
\[ (\DD_{\mathcal{M},\Sigma})_I \simeq \DD_{\mathcal{M}(I),\Sigma(I)} \quad , \]
\msp{page 43}
où $\mathcal{M}(I) = \underline{\mathsf{Hom}}(I^\circ,\mathcal{M})$ (préfaisceaux sur $I$ à valeurs dans $\mathcal{M}$), et $\Sigma(I) \subset \mathsf{Fl}\ \mathcal{M}(I)$ est déduit de $\Sigma$ en regardant fibre par fibre (supposant maintenant $\Sigma$ saturé). % Resolved: the repeated arrow-class symbol is Fl, with lowercase l.

Pour tous axiomes qu'on dégagera pour certains prédérivateurs, il faudra vérifier qu'il sont stables par passage aux prédérivateurs induits $\DD_I$, et aussi comment ils se comportent par dualité (passage de $\DD$ à $\DD^\circ$). % [sic] qu'il sont -> qu'ils sont, agreement

Définition du prédérivateur dual $\DD^\circ$
\[ \DD^\circ(I) \overset{\text{déf}}{=} \DD(I^\circ)^\circ \qquad\qquad \textit{etc.} \]
\textbf{Der 1} (Premier axiome de localisation) : \textit{a}) Si $I$, $J \in \sDia$, alors
\[ \DD(I \amalg J) \longrightarrow \DD(I) \times \DD(J) \]
est une équivalence de catégories.

\textit{b}) $\DD(\emptyset)$ est la catégorie finale.

Stable par passage aux prédérivateurs induits, et passage au dual.

\textbf{Corollaire.} \textit{Soit} $(I_\alpha)_{\alpha \in A}$ \texttt{[une]} \textit{famille finie de catégories de} $\sDia$, $I = \coprod\limits_{\alpha \in A} I_\alpha$. \textit{Alors}
\[ \DD(I) \longrightarrow \prod_{\alpha \in A} \DD(I_\alpha) \]
\textit{est une équivalence de catégories.}


\endgroup
% END UNIT Der-ChI-234/Der-ChI-234_p021-025.tex
% BEGIN UNIT Der-ChI-234/Der-ChI-234_p026-030.tex
\begingroup\makeatletter%
\NoAutoSpacing%
\newcommand{\printedpage}[1]{\marginpar{\footnotesize p.~#1}}%
\newcommand{\msp}[1]{\marginpar{\footnotesize [#1]}}%
\newcommand{\D}{\mathbb{D}}%
\newcommand{\Dia}{\mathsf{Dia}}%
\newcommand{\HOT}{\mathsf{HOT}}%
\newcommand{\Ob}{\mathsf{Ob}}%
\makeatother\let\printedpage\derpp\let\msp\dermsp%
\let\mspin\DERmspin%


\printedpage{24}
% head: A. Grothendieck : Dérivateurs, Ch. I, Édition des Universités de Montpellier II et Paris VII

\noindent\textbf{Der 2} (Deuxième axiome de localisation) : Pour tout $I \in \Dia$, le foncteur canonique
\[ \D(I) \longrightarrow \underline{\mathsf{Hom}}(I^{\circ}, \mathcal{A}_{\D}) \quad , \]
$\mathcal{A}_{\D} = \D(e)$ catégorie fondamentale pour le dérivateur $\D$, est conservatif.\\
Stable par induction, et par dualité.

\noindent\msp{page 44}\texttt{[page 44]}

\noindent Un prédérivateur est appelé un \textit{dérivateur à gauche} s'il satisfait \textbf{Der 1} et \textbf{Der 2}, et de plus les axiomes suivants (dont voici le premier).

\noindent\textbf{Der 3} (g) (Existence des images directes) : Pour tout $f : I \to J$ dans $\Dia$, le foncteur $f^{*} : \D(J) \to \D(I)$ admet un \textit{adjoint à droite} $f_{*}$\,.\\
Stable par induction.

On appelle $f_{*}$ le foncteur image directe (cohomologique) pour $f$. L'objet $f_{*}(F) \in \D(J)$ (pour $F \in \D(I)$) s'appelle aussi l'objet de “\textit{cohomologie relative}” de $I$ sur $J$ (par $f$), à coefficients dans $F$. On a des isomorphismes canoniques
\[ (gf)_{*} \simeq g_{*} f_{*} \quad . \]

\noindent\textbf{Corollaire.} \textit{Les catégories $\D(I)$ ont des produits finis, les foncteurs $f^{*}$ et $f_{*}$ y commutent.}

Utilise \textbf{Der 1} pour $I = \emptyset$ et $I = e \amalg e$, et \textbf{Der 3} pour $I \xrightarrow{f} e$.

On a un axiome dual :

\noindent\textbf{Der 3} (d) : existence des adjoints à gauche $f_{!}$ des $f^{*}$, le foncteur $f_{!}$ s'appelle “image directe homologique” par $f$, et on a une notion d'\textit{homologie relative} de $I$ sur $J$ (pour $f$), à coefficients dans $F$.

Par la suite, la mention (g) ou (d) à coté de la désignation \textbf{Der} $n$ d'un axiome, % [sic] coté -> côté, missing circumflex
signifie qu'il s'applique aux dérivateurs à gauche, respectivement à droite. De même pour les énoncés, définitions \textit{etc.}

\noindent\msp{page 45}\texttt{[page 45]}

\noindent\textbf{Der 4}$_{0}$ (g)\footnote[13]{Inutile comme axiome, va devenir un théorème (résulte des autres axiomes).} : Soit $U \overset{j}{\hookrightarrow} X \overset{i}{\hookleftarrow} Y$ dans $\Dia$, immersions ouverte et fermée complémentaires l'une de l'autre. Alors le morphisme d'adjonction dans $\D(Y)$
\[ i^{*} i_{*}(F_{Y}) \longrightarrow F_{Y} \]
est un isomorphisme, donc
\[ i_{*} : \D(Y) \longrightarrow \D(X) \]
est pleinement fidèle, et $i^{*} : \D(X) \to \D(Y)$ est un foncteur de localisation. L'image essentielle de $i_{*}$ est formée des objets $F$ de $\D(X)$ tels que $j^{*}(F)$ soit l'objet final de $\D(U)$.

% foot: Version 8 juillet 2001

\printedpage{25}
% head: A. Grothendieck : Dérivateurs, Ch. I, Édition des Universités de Montpellier II et Paris VII

\noindent\textbf{Der 4}$_{0}$ (d) : Avec les notations précédentes pour $i$, $j$, le morphisme d'adjonction dans $\D(U)$
\[ F_{U} \longrightarrow j^{*} j_{!} F_{U} \]
est un isomorphisme, - $j_{!}$ est pleinement fidèle, $j^{*}$ un foncteur de localisation. L'image essentielle de $j_{!}$ est la sous-catégorie pleine de $\D(X)$ formée des $F$ tels que $i^{*}(F)$ soit objet initial de $\D(Y)$.

Il y a les deux carrés de la page 8 et 9, mais pour interpréter la condition cartésienne ou cocartésienne c'est prématuré ici.

\noindent\msp{page 46}\texttt{[page 46]}

\noindent Axiomes de changement de base.

\noindent\textbf{Der 4} (g) (Fibres d'une image directe) : Soient $f : X \to Y$ dans $\Dia$, $y \in Y$, considérons le diagramme connu
\[
\begin{tikzcd}[row sep=large, column sep=large]
X/y \arrow[r, "q"] \arrow[d, "f'"'] & X \arrow[d, "f"] \\
e \arrow[r, "p = \alpha_{y}"'] & Y
\end{tikzcd}
\]
et le morphisme de changement de base correspondant
\[ f_{*}(F)_{y} \overset{\text{déf}}{=} p^{*} f_{*}(F) \longrightarrow f'_{*}(q^{*}(F))\bigl(\overset{\text{déf}}{=} \Gamma_{\D}(X/y, q^{*}(F))\bigr) \quad . \]

\noindent C'est toujours un isomorphisme.

\noindent\msp{page 47}\texttt{[page 47]}

\noindent\textbf{Der 5} (g) (Axiome de cofinalité pour la cohomologie)\footnote[14]{Il est conséquence d'énoncés plus faibles.} : Soit $\alpha : X' \to X$ dans $\Dia$ “cohomologiquement cofinal” \textit{i.e.} les $x\backslash X'$ sont contractiles ($x \in X$) \textit{i.e.} $\alpha^{\circ} : {X'}^{\circ} \to X^{\circ}$ est une hot-équivalence “localement sur $X^{\circ}$\,”. Alors $\forall F \in \D(X)$, l'homomorphisme
\[ \Gamma_{\D}(X, F) \longrightarrow \Gamma_{\D}(X', \alpha^{*} F) \]
est un isomorphisme (dans $\mathcal{A} = \D(e)$).

\noindent N.B. Si $f : X \to e$, on pose $\Gamma_{\D}(X, F) = f_{*}(F)$ (cohomologie “absolue” de $X$, à coefficients dans $F$).

\noindent Même chose\footnote[15]{Faux!} si $\alpha$ est “homologiquement cofinal”, i.e. les $X'/x$ sont contractiles \texttt{[phrase barrée]}.

% foot: Version 8 juillet 2001

\printedpage{26}
% head: A. Grothendieck : Dérivateurs, Ch. I, Édition des Universités de Montpellier II et Paris VII

\noindent\textbf{Remarque.} On devrait pouvoir renforcer considérablement cet axiome, en exigeant que
\[ \alpha_{!}(\varepsilon_{X'}) \xrightarrow{\sim} \varepsilon_{X} \qquad\qquad \text{dans} \quad \HOT(X) \quad , \]
où $\varepsilon_{X'}$, $\varepsilon_{X}$ \texttt{[sont]} les objets finaux de $\HOT(X')$, $\HOT(X)$, et où $\alpha_{!}(\varepsilon_{X'}) \to \varepsilon_{X}$ est l'homomorphisme d'adjonction, compte tenu que $\varepsilon_{X'} \simeq \alpha^{*}(\varepsilon_{X})$. On dit que $X'$ est \textit{homotopiquement trivial sur $X$}. (Et si la condition duale est satisfaite dans $\HOT(X^{\circ})$, on dit que $X'$ est cohomotopiquement trivial sur $X$.) Cette condition n'est sans doute \underline{pas} autoduale. Elle est vérifiée aussi bien si $\alpha$ est cohomologiquement cofinal, que s'il est homologiquement cofinal (hot-équivalence costricte ou stricte).

\noindent N.B. Pour $f : X \to e$, dire que $f$ est cohomologiquement cofinal, homologiquement cofinal, ou que $X$ est contractile, est équivalent.

\noindent\textbf{Corollaire 1} (g) (Théorème d'asphéricité pour les coefficients constants en cohomologie). \textit{Soit $f : X \to e$ dans} $\Dia$\textit{, $X$ contractile, $M \in \D(e)$. Alors l'homomorphisme d'adjonction}
\[ M \xrightarrow{\sim} f_{*} f^{*}(M) \quad \bigl(\overset{\text{déf}}{=} \Gamma(X, f^{*}(M))\bigr) \]
\textit{est un isomorphisme.}

\noindent\textbf{Corollaire 2} (g)\footnote[16]{\underline{Faux}, sauf pour coefficients constants.} \textit{Soit $f : X \to Y$} \texttt{[une]} \textit{hot-équivalence stricte (i.e. $\forall y \in Y$, $X/y$ contractile). Alors pour $F \in \Ob\,\D(Y)$}
\[ F \overset{\text{adj}}{\xrightarrow{\sim}} f_{*} f^{*} F \]
\texttt{[est un]} \textit{isomorphisme.} \texttt{[Ce corollaire est barré, mais pas sa démonstration.]}

\noindent\msp{page 48}\texttt{[page 48]}

\noindent On applique \textbf{Der 2}, \textbf{Der 4} (g), pour se ramener à regarder, pour $y_{0} \in Y$, l'homomorphisme $\alpha$ du diagramme (commutatif!) suivant :
\[
\begin{tikzcd}[row sep=large, column sep=huge]
\Gamma(\{y_{0}\}, F_{y_{0}}) = F_{y_{0}} \arrow[r, "\alpha"] & \Gamma_{\D}(X/y_{0}, f^{*}(F)) \\
\Gamma(Y/y_{0}, F) \arrow[u, "\varphi\,\wr"] \arrow[ur, "\alpha'"'] &
\end{tikzcd}
\quad ,
\]
où $\alpha'$ est déduit du foncteur induit par $f$
\[ X/y_{0} \longrightarrow Y/y_{0} \]
et où $\varphi$ est déduit du foncteur
\[ \{y_{0}\} \longrightarrow Y/y_{0} \qquad \text{(objet final)}\,. \]

% foot: Version 8 juillet 2001

\printedpage{27}
% head: A. Grothendieck : Dérivateurs, Ch. I, Édition des Universités de Montpellier II et Paris VII

\noindent Donc $\varphi$ est un isomorphisme par \textbf{Der 5} (g)\footnote[17]{\textit{cf.} corollaire 3, page 50, qu'on devrait remonter ici.} (N.B. Si $z \in Z$, d'où $\{z\} \to Z$, ce foncteur est cohomologiquement cofinal si et seulement si $z$ est objet final, homologiquement final \texttt{[plutôt cofinal]} si et seulement si $z$ est objet initial). Donc la conclusion du corollaire équivaut à dire que $\alpha'$ est un isomorphisme. Mais $X/y_{0} \to Y/y_{0}$ est aussi homologiquement cofinal, car si $u : y_{1} \to y_{0}$ est dans $Y/y_{0}$, on a un isomorphisme canonique
\[ (X/y_{0})/(y_{1}, u) \simeq X/y_{1} \]
qui est contractile. Donc par la deuxième partie de \textbf{Der 5} (g),\footnote[18]{(qui est \underline{fausse})} on gagne.

\noindent\msp{page 49}\texttt{[page 49]}

\noindent\textbf{Corollaire} 4 (g) (critère de changement de base pour images directes cohomologiques). \textit{Soit}
\[
\begin{tikzcd}[row sep=huge, column sep=huge]
X' \arrow[r, "q"] \arrow[d, "f'"'] & X \arrow[d, "f"] \arrow[dl, dashed, "\gamma"', shorten <=4ex, shorten >=1ex] \\
Y' \arrow[r, "p"'] & Y
\end{tikzcd}
\]
\textit{dans} $\Dia$ (\textit{cf.} page 30), \textit{d'où l'homomorphisme de changement de base}
\[ p^{*} f_{*}(F) \longrightarrow f'_{*} q^{*}(F) \qquad (F \in \Ob\,\D(X)) \quad . \]
\textit{Pour que ce soit un isomorphisme pour tous $F$, il suffit que pour tout $y' \in \Ob\,Y'$, le foncteur induit par} (p. 9)
\[ X'/y' \longrightarrow X/p(y') \]
\textit{soit une hot-équivalence stricte (ou costricte);\footnote[19]{\underline{Faux} pour “stricte”, ne marche \underline{que} pour “\underline{\underline{co}}stricte”.} cela signifie aussi (dans le cas strict) que pour}
\[ y'_{0} \in \Ob\,Y' \,,\ \ x_{0} \in \Ob\,X \,,\ \ u_{0} : f(x_{0}) \longrightarrow p(y'_{0}) \]
\textit{la catégorie $K(y'_{0}, x_{0}, u_{0})$ des objets de la forme}
\[ (x', v', u) \mid x' \in \Ob\,X',\ \ v' : f'(x') \to y'_{0},\ \ u : q(x') \overset{\dashleftarrow}{\longrightarrow} x_{0} \]
\textit{rendant commutatif le carré dans $Y'$}
\[
\begin{tikzcd}[row sep=huge, column sep=huge]
f(q(x')) \arrow[r, "\gamma(x')"] \arrow[d, shift right=0.5ex, "f(u)"'] & p(f'(x')) \arrow[d, "p(v')"] \\
f(x_{0}) \arrow[r, "u_{0}"'] \arrow[u, dashed, shift right=0.5ex, "\leftarrow\ \text{\scriptsize\shortstack[l]{cas de\\hot-équivalence\\costricte}}"'] & p(y'_{0})
\end{tikzcd}
\tag*{$(*)$}
\]
\textit{est contractile.}

% foot: Version 8 juillet 2001

\printedpage{28}
% head: A. Grothendieck : Dérivateurs, Ch. I, Édition des Universités de Montpellier II et Paris VII

\noindent N.B. Dans le cas de la hot équivalence \underline{co}stricte, c'est la même description de
\[ K = (x_{0}, u_{0})\backslash(X'/y'_{0}) \quad , \]
sauf que la flèche $u$ est en sens inverse $x_{0} \to q(x')$ et itou dans le diagramme commutatif $(*)$.\\
On applique \textbf{Der 2}, \textbf{Der 4} (g) et corollaire 1 \texttt{[plutôt} \textbf{Der 5} (g)\texttt{]}.

\noindent\msp{page 50}\texttt{[page 50]}

\noindent\textbf{Corollaire 5} (g). \textit{Soit $f : X \to Y$ pleinement fidèle dans} $\Dia$\textit{, et $F \in \D(X)$. Alors l'homomorphisme d'adjonction}
\[ f^{*} f_{*}(F) \longrightarrow F \]
\textit{est un isomorphisme.}

En efffet, il suffit de le vérifier fibre par fibre (\textbf{Der 2}). Soit $x \in X$, $y = f(x) \in Y$, % [sic] efffet -> effet, tripled f
alors
\[ f^{*} f_{*}(F)_{x} \simeq f_{*}(F)_{y} \simeq \Gamma_{\D}(X/y, F) \]
or à cause de la pleine fidélité $X/y \simeq X/x$.

\noindent\textbf{Corollaire 3} (remonter page 48).\footnote[20]{Va résulter de \textbf{Der 6} (g) plus fort ci-dessous.} \textit{Soit $Z \in$} $\Dia$\textit{, supposons que $Z$ a un objet final $e_{Z}$. Soit $p : e \to Z$ le foncteur correspondant. Alors l'homomorphisme canonique (pour $F$ dans} $\D(Z)$\,)
\[
\begin{tikzcd}[row sep=huge]
Z \arrow[d, shift left=0.6ex, "f"] \\
e \arrow[u, shift left=0.6ex, "p"]
\end{tikzcd}
\qquad\qquad\qquad\qquad
f_{*}(F) \longrightarrow p^{*}(F)
\]
\textit{est un isomorphisme.}\footnote[21]{Axiome de la catégorie d'indices à objet final (coefficient quelconque).}

Cet homomorphisme est \texttt{[un]} cas particulier de celui-ci, pour
\[
\begin{tikzcd}[row sep=large, column sep=large]
Z' \arrow[rr, "p"] \arrow[dr, "f'"'] & & Z \arrow[dl, "f"] \\
& S &
\end{tikzcd}
\]
(en faisant $Z' = S = e$, $f' = \mathsf{id}_{e}$),
\[ f_{*}(F) \longrightarrow f'_{*}(p^{*}(F)) \quad , \qquad (\, F \in \Ob\,\D(Z)\,) \quad , \]
et c'est un cas particulier\footnote[22]{déjà envisagé dans cor. 2, p. 47} des morphismes de

% foot: Version 8 juillet 2001


\endgroup
% END UNIT Der-ChI-234/Der-ChI-234_p026-030.tex
% BEGIN UNIT Der-ChI-234/Der-ChI-234_p031-035.tex
\begingroup\makeatletter%
\newcommand{\printedpage}[1]{\marginpar{\footnotesize p.~#1}}%
\newcommand{\msp}[1]{\marginpar{\footnotesize [#1]}}%
\newcommand{\Ob}{\mathsf{Ob}}%
\newcommand{\id}{\mathsf{id}}%
\newcommand{\DD}{\mathbb{D}}%
\makeatother\let\printedpage\derpp\let\msp\dermsp%
\let\mspin\DERmspin%


\printedpage{29}
% head: A. Grothendieck : Dérivateurs, Ch. I, Édition des Universités de Montpellier II et Paris VII
% foot: Version 8 juillet 2001

\noindent\msp{page 51}
changement de base pour
\[
\begin{tikzcd}
Z' \arrow[r, "p"] \arrow[d, "f'"'] & Z \arrow[d, "f"] \\
S \arrow[r, "\id"] & S
\end{tikzcd}
\qquad .
\]

\noindent\textsc{Remarque} : On retrouve le fait que $i^*i_*(F_Y) \overset{\sim}{\to} F_Y$ est \texttt{[un]} isomorphisme, pour $i$ une immersion fermée. De plus, avec les notations de \textbf{Der 4} (g), $j^*(i_*(F_Y))$ se calcule par le théorème de changement de base (corollaire 2 \texttt{[plutôt 4]} ci-dessus), qui s'applique dans toute situation commutative
\[
\begin{tikzcd}[row sep=large, column sep=large]
X' \arrow[r, "q"] \arrow[d, "f'"'] \arrow[dr, phantom, "\text{\scriptsize cart}"] & X \arrow[d, "f"] \\
Y' \arrow[r, "p"] & Y
\end{tikzcd}
\]
quand $p$ est une immersion ouverte (et $\gamma = \id$) (alors $X'/y' \to X/p(y')$ est même un isomorphisme!), et on trouve, revenant aux notations de \textbf{Der 4} (g),
\[
\begin{tikzcd}
\emptyset \arrow[r, "q"] \arrow[d, "f'"'] & Y \arrow[d, hook, "i"] \\
U \arrow[r, hook, "j"'] & X
\end{tikzcd}
\]
\textit{i.e.} $j^*(i_*(F_Y)) \simeq f'_*(\varphi)$, où $\varphi = q^*(F_Y)$ est l'unique élément de la catégorie (équivalente à une catégorie ponctuelle, par \textbf{Der 1}) $\DD(\emptyset)$. D'autre part, $f'_*$ transforme objet final en objet final, d'où $j^*i_*(F_Y)$ est objet final.

Inversement, soit $F$ dans $\DD(X)$, tel que $j^*(F) = e_{\DD(U)}$. Considérons
\[
F \xrightarrow{u} i_*i^*F \qquad\quad \text{(adjonction) ;}
\]
je dis que c'est un isomorphisme. En effet, par l'axiome \msp{page 52} de localisation \textbf{Der 2}, il suffit de prouver que $i^*(u)$ et $j^*(u)$ sont \texttt{[des]} isomorphismes. Pour $i^*(u)$ cela résulte du fait que $i^*$ soit foncteur de localisation, et pour $j^*$ du fait que $j^*(F)$ et $j^*(i_*i^*(F))$ sont des objets finaux dans $\DD(U)$.

\printedpage{30}
% head: A. Grothendieck : Dérivateurs, Ch. I, Édition des Universités de Montpellier II et Paris VII
% foot: Version 8 juillet 2001

\noindent\msp{page 53}

\subsection*{Morphismes homotopiquement propres, homotopiquement lisses.}

\noindent Soit $f : X \to Y$ \texttt{[un]} morphisme de catégories, considérons $y_0 \in \Ob\, Y$, et le foncteur canonique
\[
X_{y_0} \xrightarrow{\varphi} X/y_0 \qquad\qquad\qquad \varphi(x) = (x, \id_{f(x)}) \tag*{(1)}
\]
($X_{y_0}$ fibre de $X$ en $y_0$). Soit $(x_0, u_0) \in \Ob\, X/y_0$ donc $x_0 \in \Ob\, X$, $u_0 : f(x_0) \to y_0$. Considérons la catégorie coinduite par $\varphi$ sur $(x_0, u_0)$
\[
(x_0, u_0)\backslash(X_{y_0}) \quad ,
\]
pour tester si $\varphi$ est une hot-équivalence costricte \textit{i.e.} homologiquement \texttt{[plutôt co\-ho\-mo\-lo\-gi\-que\-ment]} cofinal. On va l'interpréter via le foncteur induit par $f$
\[
x_0\backslash X \xrightarrow{\psi} f(x_0)\backslash Y
\]
allant de la catégorie des spécialisations de $x_0$ dans $X$ vers celle des spécialisations de $f(x_0)$ dans $Y$. On trouve alors un isomorphisme canonique de la catégorie coinduite $(x_0, u_0)\backslash X_{y_0}$ par $X_{y_0}$ au dessus de l'objet $(x_0, u_0)$ de $X/y_0$ et de la catégorie fibre $(x_0\backslash X)_{u_0}$ de $x_0\backslash X$ au dessus de $f(x_0)\backslash Y$ en $u_0$. Donc

\noindent\textbf{Proposition.} \textit{Pour que les foncteurs $\varphi : X_{y_0} \to X/y_0$ ($y_0 \in \Ob\, Y$) soient des hot-équivalences costrictes, il faut et il suffit que les fibres des foncteurs $x_0\backslash X \to f(x_0)\backslash Y$ induits par $f$ soient contractiles.}

\noindent\msp{page 54}
On dit alors que $f$ est \textit{homotopiquement propre.} La notion duale s'exprime, soit par la condition que les foncteurs induits par $f$
\[
X_{y_0} \longrightarrow y_0\backslash Y
\]
soient des hot-équivalences strictes \textit{i.e.} qu'ils soient cohomologiquement \texttt{[plutôt ho\-mo\-lo\-gi\-que\-ment]} cofinaux, soit que les fibres des morphismes induits sur les ``catégories localisées'' (catégories des générisations)
\[
X/x_0 \longrightarrow Y/f(x_0)
\]
soient contractiles. (N.B. On utilise ici qu'une catégorie $Z$ est contractile si et seulement si $Z^\circ$ l'est.) On dit alors que $f$ est \textit{homotopiquement lisse.}

\noindent\textbf{Proposition.} \textit{La notion de foncteur propre (resp. lisse) est stable par changement de base, par composition. Les isomorphismes, et plus généralement, les immersions fermées sont propres. L'image d'un morphisme propre est une sous-catégorie fermée.}\footnote[23]{Attention, une équivalence de catégories n'est pas propre! Si $X' \hookrightarrow X$ \texttt{[est]} une inclusion d'une sous-catégorie de $X$, elle n'est propre que si c'est une immersion fermée.}

\noindent N.B. Le foncteur $X \to e$ est toujours propre et lisse, car les catégories $x_0\backslash X$ et $X/x_0$ sont contractiles (ayant objet initial ou objet final). Donc tout foncteur de projection $X \times Y \xrightarrow{pr_2} Y$ est à la fois propre et lisse.

\printedpage{31}
% head: A. Grothendieck : Dérivateurs, Ch. I, Édition des Universités de Montpellier II et Paris VII
% foot: Version 8 juillet 2001

\noindent\msp{page 55}
\textsc{Démonstration}. Stabilité par changement de base :
\[
\begin{tikzcd}
X \arrow[d, "f"'] \arrow[dr, phantom, "\text{\scriptsize cart}"] & X' \arrow[l, "q"'] \arrow[d, "f'"] \\
Y & Y' \arrow[l, "p"]
\end{tikzcd}
\qquad ,
\]
d'où pour $x' \in \Ob\, X'$
\[
\begin{tikzcd}[row sep=large]
X/x \arrow[d] & X'/x' \arrow[l] \arrow[d] \\
Y/y & Y'/y' \arrow[l]
\end{tikzcd}
\qquad ,
\]
où $x$, $y'$, $y$ sont les images de $x'$ dans $X$, $Y'$, $Y$. On vérifie que ce diagramme est \textit{cartésien}, donc les fibres de $X'/x' \longrightarrow Y'/y'$ sont isomorphes à des fibres de $X/x \longrightarrow Y/y$, donc contractiles si celles-ci le sont.

Dans un cas $X \xrightarrow{f} Y$ on vérifie que pour un foncteur bilocalisé de $f$
\[
X/x \xrightarrow{\varphi=f/x} Y/y \quad ,
\]
tout ``bilocalisé''
\[
(X/x)/\xi \longrightarrow (Y/y)/\eta \qquad\qquad (\xi : x' \to x\,, \quad \eta : y' \to y)
\]
de $\varphi$ s'identifie encore à un bilocalisé de $f$
\[
X/x' \longrightarrow Y/y' \quad .
\]
Donc si $f$ est lisse, les bilocalisés $f/x$ ($x \in \Ob\, X$) le sont aussi. Notons d'autre part.

\noindent\textbf{Proposition.}\footnote[24]{\textit{cf.} amélioration dans cor. 2, page suivante.} \textit{Tout morphisme propre (ou dualement lisse) à fibres contractiles est une hot-équivalence.}

\noindent\msp{page 56}
Cela provient du fait que pour un tel foncteur propre, si $F$ est un préfaisceau, ou complexe de préfaisceaux abéliens sur $X$, la formation des $Rf_*(F)$ commute au passage aux fibres, et de même pour $R^1f_*(F)$ si $F$ est un préfaisceau en groupes (pas même abéliens).

\noindent\textbf{Corollaire 1.} \textit{Si $f : X \to Y$ est cohomologiquement propre à fibres contractiles, alors $Y$ est contractile si et seulement si $X$ l'est.}

Grâce à ceci, on trouve que le composé de deux foncteurs propres (lisses) est propre (lisse).

\printedpage{32}
% head: A. Grothendieck : Dérivateurs, Ch. I, Édition des Universités de Montpellier II et Paris VII
% foot: Version 8 juillet 2001

\noindent N.B. Contrairement à ce qui a lieu en topologie pour la notion de propreté, si $gf$ est propre et $g$ propre, $f$ ne l'est pas nécessairement, \textit{i.e.} un $S$-morphisme de $S$-catégories $X$, $Y$ propres sur $S$ n'est pas nécessairement propre. Ainsi une section de $X/S$, pour $X$ propre sur $S$, n'est pas nécessairement propre, \textit{i.e.} n'est pas nécessairement une immersion fermée. C'est déjà faux si $S = e$ (auquel cas $f : X \to e$ est toujours propre).

\noindent\textbf{Corollaire 2.} \textit{Un morphisme propre ou lisse, à fibres contractiles est à la fois une hot-équivalence stricte et costricte.}

En effet, on applique le corollaire 1 à $X/y \to Y/y$ ou $y\backslash X \to y\backslash Y$.

\noindent\textbf{Théorème} (g) \textit{Soit une situation de changement de base dans} \textsf{Dia}
\[
\begin{tikzcd}[row sep=large, column sep=huge]
X \arrow[d, "f"'] \arrow[dr, phantom, "\text{cart}"] & X' \arrow[l, "q"'] \arrow[d, "f'"] \\
Y & Y' \arrow[l, "p"]
\end{tikzcd}
\qquad , \qquad\qquad p^*f_*(F) \xrightarrow{c} f'_*q^*(F)
\]
\textit{($F \in \DD(X)$), $\DD$ dérivateur exact à gauche. Si $f$ est (homotopiquement) propre ou $p$ homotopiquement lisse, alors $c$ est un isomorphisme.}

\noindent\msp{page 57}
La question c'est si pour tout $y' \in Y'$, posant $y = p(y')$ et considérant le foncteur induit par $q$
\[
X/y \longleftarrow X'/y'
\]
le morphisme correspondant
\[
\Gamma_{\DD}(X/y, F) \longrightarrow \Gamma_{\DD}(X'/y', q^*F)
\]
est \texttt{[un]} isomorphisme. Dans le cas $f$ propre, on considère
\[
\begin{tikzcd}[column sep=large]
X_y \arrow[d, "\sim\ \ \text{costricte}"'] & X_{y'} \arrow[l, "\sim"'] \arrow[d, "\sim\ \ \text{costricte}"] \\
X/y & X'/y' \arrow[l]
\end{tikzcd}
\]
d'où la conclusion, puisque par \textbf{Der 5} (g), dans
\[
\begin{tikzcd}[row sep=large, column sep=large]
\Gamma(X_y) \arrow[r] & \Gamma(X_{y'}) \\
\Gamma(X/y) \arrow[u] \arrow[r] & \Gamma(X'/y') \arrow[u]
\end{tikzcd}
\]
les flèches verticales et l'horizontale supérieure sont \texttt{[des]} isomorphismes, donc aussi l'horizontale inférieure.

\printedpage{33}
% head: A. Grothendieck : Dérivateurs, Ch. I, Édition des Universités de Montpellier II et Paris VII
% foot: Version 8 juillet 2001

Dans le cas $p$ lisse, considérons du coté $Y$, $Y'$ le carré de changement de base % [sic] coté -> côté, missing circumflex (cf. "Du côté" below)
\[
\begin{tikzcd}[row sep=large, column sep=large]
Y & Y' \arrow[l] \\
Y/y \arrow[u] & Y'/y' \arrow[l] \arrow[u]
\end{tikzcd}
\]

\noindent\msp{page 58}
d'où le diagramme contigu
\[
\begin{tikzcd}[column sep=small]
X \arrow[dd] & & X' \arrow[ll] \arrow[dd] & \\
& X/y \arrow[ul] & & X'/y' \arrow[ll, crossing over, "\kappa"'] \arrow[ul] \arrow[dd] \\
Y & & Y' \arrow[ll] & \\
& Y/y \arrow[ul] \arrow[from=uu, crossing over] & & Y'/y' \arrow[ll, "\pi"] \arrow[ul]
\end{tikzcd}
\qquad ,
\]
où les carrés latéraux (verticaux) sont cartésiens. On note que $\pi$ est lisse et à fibres contractiles (car $p$ est lisse) donc $\kappa$ aussi. Donc $\kappa$ est une hot-équivalence costricte, donc $\Gamma(X/y) \overset{\sim}{\longrightarrow} \Gamma(X'/y')$, on gagne encore.

\noindent N.B. Sauf erreur, on a une réciproque du théorème, \texttt{[à]} savoir : Si $f : X \to Y$ est tel que le théorème de changement de base pour $f_*(F)$ est vrai pour \textit{tout} dérivateur (il suffit sans doute de \textsf{HOTAB} et \textsf{HOT}) et pour tout $Y' \to Y$, alors $f$ est homotopiquement propre. D'autre part, si $g : Y' \to Y$ est tel que le théorème de changement de base pour $g$ est vrai pout tout $X \to Y$ et tout $F \in \DD(X)$ ($\DD$ dérivateur exact à droite quelconque), alors $g$ est homotopiquement lisse. % [sic] pout tout -> pour tout, misprint

\noindent Du côté homologique les foncteurs lisses $f$ sont exactement ceux pour lesquels $f_!$ commute à tout changement de base. Les foncteurs propres $g$ sont exactement ceux pour lesquels le changement de base par $g$ commute à $f_!$, pour tout $f : X \to Y$ et tout $\DD$ dérivateur exact à gauche.

\noindent\msp{page 59}
Ainsi, les foncteurs propres $f : X \to Y$ se caractérisent de \textit{quatre} façons remarquables
\begin{itemize}[label={}, leftmargin=0pt, itemindent=2\parindent, labelsep=0.3em]
\item[\textit{a})] $f_*$ commute à tout changement de base ($\forall$ dérivateur $\DD$ exact à gauche);
\item[\textit{b})] le changement de base $f^*$ commute à tout $g_!$ ($g : Y' \to Y$, $\DD$ dérivateur exact à droite);
\item[\textit{c})] $\forall y \in \Ob\, Y$, $X_y \to X/y$ est une hot-équivalence costricte \textit{i.e.} est un foncteur homologiquement \texttt{[plutôt co\-ho\-mo\-lo\-gi\-que\-ment]} cofinal;
\item[\textit{d})] $\forall x \in \Ob\, X$ le foncteur induit sur les colocalisés
\[
x\backslash X \longrightarrow f(x)\backslash Y
\]
\end{itemize}


\endgroup
% END UNIT Der-ChI-234/Der-ChI-234_p031-035.tex
% BEGIN UNIT Der-ChI-234/Der-ChI-234_p036-040.tex
\begingroup\makeatletter%
\derleqno%
\NoAutoSpacing%
\newcommand{\printedpage}[1]{\marginpar{\footnotesize p.~#1}}%
\newcommand{\msp}[1]{\marginpar{\footnotesize [#1]}}%
\newcommand{\DD}{\mathbb{D}}%
\newcommand{\cA}{\mathcal{A}}%
\newcommand{\Dia}{\mathsf{Dia}}%
\newcommand{\Ob}{\mathsf{Ob}}%
\newcommand{\Hom}{\mathsf{Hom}}%
\newcommand{\hot}{\mathsf{hot}}%
\newcommand{\Hot}{\mathsf{Hot}}%
\newcommand{\HOT}{\mathsf{HOT}}%
\newcommand{\isoto}{\overset{\sim}{\to}}%
\newcommand{\isolongto}{\overset{\sim}{\longrightarrow}}%
\newcommand{\sHom}{\mathsf{Hom}}%
\newcommand{\uHom}{\underline{\mathsf{Hom}}}%
\makeatother\let\printedpage\derpp\let\msp\dermsp%
\let\mspin\DERmspin%


\printedpage{34}
% head: A. Grothendieck : Dérivateurs, Ch. I, Édition des Universités de Montpellier II et Paris VII
% foot: Version 8 juillet 2001
\noindent est à fibres contractiles.

Autre situation standard où le changement de base marche (mais où les foncteurs qui interviennent n'ont aucune raison d'être propres ou lisses).

Soient
\[
f : I \longrightarrow I' \quad , \qquad g : J \longrightarrow J'
\]
dans $\Dia$, d'où le diagramme (comparer page 23)
\[
\begin{tikzcd}
X' = I \times J \arrow[r, "G'"] \arrow[d, "F'"'] & I \times J' = X \arrow[d, "F"] \\
Y' = I' \times J \arrow[r, "G"] & I' \times J' = Y
\end{tikzcd}
\quad .
\]
Ce diagramme est tel que pour tout dérivateur à gauche

\noindent\msp{page 60}

\noindent on a le changement de base
\[
G^*F_* \simeq F'_*{G'}^* \quad ,
\]
la raison ici étant que les foncteurs induits
\[
X'/y' \longrightarrow X/y
\]
sont des hot-équivalences \underline{\underline{co}}strictes, \textit{i.e.} que les catégories $K$ (page 49) sont contractiles. En effet, si $x_0 = (i_0, j'_0) \in \Ob X$, $y'_0 = (i'_0, j_0) \in \Ob Y'$, $u_0 = (u'_0, v'_0) : F(x_0) \to G(y'_0)$, où $u'_0 : f(i_0) \to i'_0$, $v'_0 : j'_0 \to g(j_0)$ sont donnés, on vérifie que la catégorie $K$ cherchée est isomorphe simplement à $(i_0, u'_0)\backslash(I/i'_0) \times (j_0, v'_0)\backslash(J/j'_0)$ qui est contractile ayant un objet initial \texttt{[plutôt $(i_0, u'_0)\backslash(I/i'_0) \times (j'_0\backslash J)/(j_0, v'_0)$ qui est contractile étant produit d'une catégorie ayant un objet initial et d'une catégorie ayant un objet final]}. Ouf!

Je m'aperçois d'un très gros canular. Je croyais me rappeler que pour $f : X \to Y$ qui est une hot-équivalence stricte, on avait pour tout $F \in \DD(Y)$, $F \isoto f_*f^*F$ (par exemple pour $\DD = \mathsf{Der}(\mathit{Ab})$, ou $\DD = \Hot$), \textit{i.e.} $f^*$ pleinement fidèle, d'où aussi que $\Gamma_{\DD}(Y, F) \to \Gamma_{\DD}(X, f^*F)$ est \texttt{[un]} isomorphisme. Mais ce n'est vrai que si $F$ est un coefficient \textit{constant} sur $Y$, \textit{i.e.} de la forme $q^*(M)$ ($q : Y \to e$, $M \in \DD(e) = \cA_{\DD}$). C'est faux, pour $F$ général, déjà dans le cas de l'immersion $f : \{\varphi\} \to Y$, où $\varphi$ est un objet initial. Alors, $f$ est toujours une hot-équivalence stricte, mais l'homomorphisme $\Gamma(Y, F) \to \Gamma(X, f^*F)$ n'est en général pas un isomorphisme.

\par
\printedpage{35}
% head: A. Grothendieck : Dérivateurs, Ch. I, Édition des Universités de Montpellier II et Paris VII
% foot: Version 8 juillet 2001
\noindent\msp{page 61}

\subsection*{Notions de $\mathbb{D}$-équivalence et de $\mathbb{D}$-contractibilité.}

On va donc reprendre les questions d'acyclicité et de changement de base.

\noindent\textbf{Définition 1.} Soit $\DD$ un prédérivateur, et soit
\[
\begin{tikzcd}
X \arrow[rr, "f"] \arrow[dr, "p"'] & & Y \arrow[dl, "q"] \\
& e &
\end{tikzcd}
\]
dans $\Dia$. On dit que $f$ est une $\DD$-\textit{équivalence} si le foncteur
\[
F \longmapsto f^*F \quad , \qquad \DD(Y) \longrightarrow \DD(X)
\]
induit un foncteur pleinement fidèle sur la sous-catégorie de $\DD(Y)$ formée des $\DD$-coefficients \textit{constants}, \textit{i.e.} si pour tous $M, N \in \Ob \cA = \DD(e)$ (“coefficients absolus”), on a
\[
\sHom(q^*M, q^*N) \isolongto \sHom(p^*M, p^*N) \quad . \tag*{(1)}
\]
Supposons que $p_*$ et $q_*$ existent. On posera, pour un faisceau $F$ sur $X$
\[
H^*_{\DD}(X, F) \overset{\text{déf}}{=} p_*(F) \qquad \text{(et de même } \ H^{\DD}_*(X, F) = p_!(F) \ \text{).}
\]
N.B. Et si $F = p^*M$, $M \in \cA$, on note simplement $H^*_{\DD}(X, M)$, $H^{\DD}_*(X, M)$.

\noindent L'homomorphisme (1) s'interprète par adjonction comme
\[
\begin{tikzcd}[row sep=small]
\sHom(M, q_*q^*N) \arrow[r] \arrow[d, equal] & \sHom(M, p_*p^*N) \arrow[d, equal] \\
\sHom(M, H^*_{\DD}(Y, q^*N)) \arrow[r] & \sHom(M, H^*_{\DD}(X, p^*N))
\end{tikzcd}
\]
déduit de l'homomorphisme
\[
q_*q^*N \longrightarrow p_*p^*N \quad , \qquad \textit{i.e.} \quad H^*_{\DD}(Y, N) \longrightarrow H^*_{\DD}(X, N) \quad . \tag*{(2)}
\]
Donc dire que (1) est un isomorphisme pour tous $M, N$ signifie que (2) est un isomorphisme pour tout $N$. De même, si $p_!$, $q_!$ existent cela équivaut à dire que
\[
p_!p^*N \longrightarrow q_!q^*N \quad , \qquad \textit{i.e.} \quad H^{\DD}_*(X, N) \longrightarrow H^{\DD}_*(Y, N) \tag*{$(2')$}
\]
est un isomorphisme dans $\cA$ pour tout $N \in \cA$.

\par
\printedpage{36}
% head: A. Grothendieck : Dérivateurs, Ch. I, Édition des Universités de Montpellier II et Paris VII
% foot: Version 8 juillet 2001
\noindent\msp{page 62}

\noindent En résumé :

\noindent\textbf{Proposition 1.} \textit{Si $p_*$ et $q_*$ existent, i.e. si $H^*_{\DD}(X, \bullet)$ et $H^*_{\DD}(Y, \bullet)$ existent (il suffit même que ça existe pour des coefficients constants), alors $f$ est une $\DD$-équivalence si et seulement si $f$ induit un isomorphisme pour les $\DD$-cohomologies à coefficients constants de $X$ et $Y$. Si $p_!$ et $q_!$ existent, i.e. si $H^{\DD}_*(X, \bullet)$ et $H^{\DD}_*(Y, \bullet)$ existent, $f$ est une $\DD$-équivalence si et seulement si $f$ induit un isomorphisme pour les $\DD$-homologies à coefficients constants de $X$ et $Y$.}

\noindent\textbf{Remarque 1.} Pour que $f$ soit \texttt{[une]} $\DD$-équivalence, il \textit{suffit} que pour tout $F$ \textit{constant} sur $Y$ le morphisme d'adjonction $F \to f_*f^*F$ soit \texttt{[un]} isomorphisme.

\noindent\textbf{Remarque 2.} Si $f$, $f' : X \to Y$ sont dans une même composante connexe de $\uHom(X, Y)$, alors $f$ est $\DD$-équivalence $\Leftrightarrow$ $f'$ $\DD$-équivalence.

\noindent\textbf{Proposition 2.} \textit{Les isomorphismes sont des $\DD$-équivalences. Si $X \overset{f}{\to} Y \overset{g}{\to} Z$, si deux parmis $f$, $g$, $gf$ sont des $\DD$-équivalences, le troisième aussi.} % [sic] parmis -> parmi, spelling

\noindent\textbf{Définition 2.} On dit que $X$ dans $\Dia$ est $\DD$-\textit{asphérique} (ou $\DD$-\textit{contractile}, ou $\DD$-\textit{trivial}) si $X \overset{p}{\to} e$ est une $\DD$-équivalence, \textit{i.e.} si $p^*$ est pleinement fidèle.

Si $p_*$ existe, cela signifie donc que pour $M \in \Ob(\cA)$
\[
M \overset{\mathrm{adj}}{\isolongto} p_*p^*M \quad , \qquad \textit{i.e.} \quad M \isolongto H^*_{\DD}(X, M) \quad ,
\]
et si $p_!$ existe, cela signifie que pour tout $M$
\[
p_!p^*M \overset{\mathrm{adj}}{\isolongto} M \quad , \qquad \textit{i.e.} \quad H^{\DD}_*(X, M) \isolongto M \quad .
\]
\textbf{Corollaire.} \textit{Si $Y$} (resp. $X$) \textit{est contractile, alors $f : X \to Y$ est une $\DD$-équivalence si et seulement si $X$} (resp. $Y$) \textit{est contractile.}

\noindent\textsc{Exemple}. Prenons $\DD = \HOT$. (On peut prendre $\Dia$ = catégorie des petites catégories). Alors, pour $M \in \cA = \Hot$ et $X$ petite catégorie, on a
\[
H^*_{\HOT}(X, M) \simeq \sHom_{\cA}(\hot(X), M) \quad ,
\]
où
\[
\hot : \mathcal{C}\mathit{at} \longrightarrow \cA = \Hot
\]
est le foncteur canonique qui associe à tout $X$ son “type d'homotopie”. On a
\[
H^{\HOT}_*(X, M) = \hot(X) \times M \quad .
\]

\noindent\msp{page 63}

\noindent On en déduit que $f : X \to Y$ est une $\HOT$-équivalence (au sens de la définition 1) si et seulement si c'est une $\hot$-équivalence (ou “équivalence faible”) \textit{i.e.} si et seulement si
\[
\hot(f) : \hot(X) \isolongto \hot(Y)
\]
\printedpage{37}
% head: A. Grothendieck : Dérivateurs, Ch. I, Édition des Universités de Montpellier II et Paris VII
% foot: Version 8 juillet 2001
est un isomorphisme.

\noindent\textbf{Proposition 2.} \textit{Soit $f : X \to Y$. Conditions équivalentes (si $\DD$} \texttt{[est un]} \textit{dérivateur} à gauche). % Resolved against official supplemental edition p37: a gauche). is upright.

a) \textit{$f$ est une $\DD$-équivalence, et le reste après tout changement de base $Y' \to Y$ qui est un isomorphisme local (par exemple pour $Y/F \to Y$, $F \in \Ob \widehat{Y}$).}

b) \textit{Pour tout $y \in \Ob Y$, le foncteur induit $X/y \to Y/y$ est une $\DD$-équivalence.}

c) \textit{Pour tout $y \in \Ob Y$, $X/y$ est $\DD$-contractile.}\footnote[25]{Autres conditions équivalentes :\par
\indent d) \textit{Pour tout faisceau \underline{constant} $M_Y$ sur $Y$, on a $M_Y \isoto f_*f^*M_Y$.}\par
\indent d$'$) \textit{Comme} d) \textit{avec faisceau \underline{localement} constant.}}

\noindent N.B. On dit que $g : Z \to Y$ est un \textit{isomorphisme local} si pour tout $z \in \Ob Z$, posons $y = g(z)$, le morphisme induit sur les catégories localisées en $z$ et $y$
\[
Z/z \isolongto Y/y
\]
est un isomorphisme. Par exemple toute immersion ouverte, tout foncteur induit $Y/y \to Y$ (foncteur de localisation) sont des isomorphismes locaux.

Il en résulte a) $\Rightarrow$ b) trivialement.

b) $\Leftrightarrow$ c). Résulte du corollaire de la page précédente, si l'on sait que $Y/y$ est $\DD$-contractile\footnote[26]{Cf. cor. 1 de \textbf{Der 5} (g)}. Cela résultera de \textbf{Der 5} ci-dessous.

\noindent\msp{page 64}

c) $\Rightarrow$ a) “Passage du local au global”. Comme l'hypothèse c) est stable par changement de base par immersion ouverte, il suffit de prouver que c) implique que $f$ est une $\DD$-équivalence, \textit{i.e.} que pour tout $F = M_Y$ ($M \in \Ob \cA$), coefficient constant sur $Y$, l'homomorphisme d'adjonction
\[
M_Y \longrightarrow f_*f^*M_Y = f_*M_X
\]
est \texttt{[un]} isomorphisme. Il suffit de prouver que c'est un isomorphisme fibre par fibre (\textbf{Der 2}), donc que pour tout $y \in Y$,
\[
M_y \longrightarrow f_*(M_X)_y
\]
est un isomorphisme. Par \textbf{Der 4} (g), on a
\[
f_*(M_X)_y \isolongto \Gamma(X/y, M) \quad ,
\]
et l'hypothèse que $X/y$ est $\DD$-contractile donne ce qu'on veut.

\noindent\textbf{Proposition 2$'$.} \textit{Dualement, si $\DD$ est un dérivateur à gauche} \texttt{[plutôt à droite]}\textit{, conditions équivalentes sur $f : X \to Y$ :}

a) \textit{$f$ $\DD$-équivalence,} \texttt{[et]} \textit{le reste par tout changement de base $g : Y' \to Y$ qui est un “isomorphisme colocal” (i.e. pour tout $y' \in \Ob Y'$, le foncteur induit sur les catégories “colocalisées” $y'\backslash Y' \isolongto y\backslash Y$ (où $y = g(y')$)} \texttt{[est un isomorphisme]}\textit{).}

b) \textit{Pour tout $y \in Y$, le foncteur induit $y\backslash X \to y\backslash Y$ est une $\DD$-équivalence.}

c) \textit{Pour tout $y \in Y$, $y\backslash Y$ est $\DD$-contractile.}

\noindent Pour le prouver, on doit utiliser que les $y\backslash Y$ sont $\DD$-contractiles (cf. cor. 1 de \textbf{Der 5} (d)).

\par
\printedpage{38}
% head: A. Grothendieck : Dérivateurs, Ch. I, Édition des Universités de Montpellier II et Paris VII
% foot: Version 8 juillet 2001
\noindent\msp{page 65}

\noindent\textbf{Définition 3.} Soit $\DD$ un prédérivateur (peut-être même un dérivateur à gauche ou à droite). On dit que $f$ est une $\DD$-équivalence \textit{stricte} si $f$ satisfait la condition a) de \texttt{[la]} proposition 2. On dit que c'est une $\DD$-équivalence \textit{costricte} si $f$ satisfait celle de \texttt{[la]} proposition $2'$; on dit aussi parfois que $f$ est $\DD$-cofinal. (La condition duale, \textit{i.e.} $\DD$-équivalence stricte, s'appelerait co-$\DD$-cofinal!)\footnote[27]{Meilleure terminologie : $f$ est $\DD$-\textit{asphérique} ($\DD$-équivalence stricte), $\DD$-\textit{coasphérique} ($\DD$-équivalence co-stricte).} % [sic] s'appelerait -> s'appellerait, spelling

\noindent\textbf{Proposition 3.} \textit{Soit $X \overset{f}{\to} Y \overset{g}{\to} Z$. Si $f, g$ sont des $\DD$-équivalences strictes, alors $gf$ aussi. Si $gf$ et $f$ sont des $\DD$-équivalences strictes, alors $g$ aussi. Si $\DD$} \texttt{[est un]} \textit{dérivateur à gauche ou à droite, toute équivalence} \texttt{[de catégories]} \textit{est une $\DD$-équivalence stricte.} (Car elle a un adjoint à droite.)

On utilise le fait que la notion d'isomorphisme local est stable par changement de base.

\noindent\msp{page 66}

\noindent\textbf{Proposition 4.} \textit{Soient $\DD$ un dérivateur à droite, $f : X \to Y$ un foncteur dans} $\Dia$\textit{, tel que pour $y \in Y$, $y\backslash X$ soit $\DD$-contractile (donc $f$ une $\DD$-équivalence costricte, sous réserve que $\DD$ soit un dérivateur à droite \ldots\ ). Alors pour tout $F$ dans $\DD(Y)$ (pas nécessairement à coefficients constant) on a} % [sic] constant -> constants, agreement
\[
H^*_{\DD}(Y, F) \isolongto H^*_{\DD}(X, f^*(F))
\]
\textit{est un isomorphisme. Réciproquement, si la condition est vraie pour tout $F$, alors les $y\backslash X$ sont $\DD$-contractiles.}

\noindent\textsc{Démonstration}. Il revient au même de dire que pour tout $M \in \Ob \cA = \DD(e)$, on a
\[
\sHom(M, H^*_{\DD}(Y, F)) \overset{\mathrm{iso}}{\longrightarrow} \sHom(M, H^*_{\DD}(X, f^*(F))) \ .
\]
Mais par adjonction, cette flèche s'interprète comme
\[
\sHom(M_Y, F) \longrightarrow \sHom(M_X, f^*F) \tag*{$(*)$}
\]
D'autre part, si $f_!$ existe, on a
\[
\sHom(M_X, f^*F) = \sHom(f_!(M_X), F) \ ,
\]
et l'homomorphisme $(*)$ est déduit de l'homomorphisme d'adjonction
\[
f_!(M_X) = f_!(f^*(M_Y)) \longrightarrow M_Y \ , \tag*{$(**)$}
\]
en prenant les $\sHom$ dans $F$. Donc la condition équivaut à dire que cet homomorphisme est \texttt{[un]} isomorphisme, quel que soit le coefficient constant $M_Y$ sur $Y$. Pour voir si c'est un isomorphisme, on regarde fibre par fibre (\textbf{Der 2}), or les fibres s'interprètent comme $H^{\DD}_*(y\backslash X, M_X)$ (par \textbf{Der 4} (d)) et comme $M$.


\endgroup
% END UNIT Der-ChI-234/Der-ChI-234_p036-040.tex
% BEGIN UNIT Der-ChI-234/Der-ChI-234_p041-045.tex
\begingroup\makeatletter%
\newcommand{\printedpage}[1]{\marginpar{\footnotesize p.~#1}}%
\newcommand{\msp}[1]{\marginpar{\footnotesize [#1]}}%
\newcommand{\D}{\mathbb{D}}%
\newcommand{\Dia}{\mathsf{Dia}}%
\newcommand{\Ob}{\mathsf{Ob}}%
\newcommand{\Fl}{\mathsf{Fl}}%
\newcommand{\Hom}{\underline{\mathsf{Hom}}}%
\newcommand{\hot}{\mathsf{hot}}%
\newcommand{\Cat}{\mathcal{C}at}%
\newcommand{\isoto}{\overset{\sim}{\longrightarrow}}%
\makeatother\let\printedpage\derpp\let\msp\dermsp%
\let\mspin\DERmspin%


\printedpage{39}
% head: A. Grothendieck : Dérivateurs, Ch. I, Édition des Universités de Montpellier II et Paris VII
% foot: Version 8 juillet 2001
\msp{page 67}

\noindent Donc la flèche envisagée $(**)$ s'interprète comme
\[
H_*^{\D}(y\backslash X, M_X) \longrightarrow M
\]
et c'est l'homomorphisme d'adjonction $p_!(p^*(M)) \to M$.
\[
\begin{tikzcd}
X \arrow[rr, "f"] \arrow[dr, "p"{swap}] & & Y \arrow[dl, "q"] \\
& e &
\end{tikzcd}
\]

\noindent Donc c'est un isomorphisme pour tout $M$ si et seulement si $y\backslash X$ est $\D$-contractile.

\noindent On va énoncer ce qu'on a prouvé.

\noindent\textbf{Théorème 1.} \textit{Soit}
\[
\begin{tikzcd}
X \arrow[rr, "f"] \arrow[dr, "p"{swap}] & & Y \arrow[dl, "q"] \\
& e &
\end{tikzcd}
\]
\textit{dans} $\Dia$. \textit{On suppose que} $\D$ \textit{est un prédérivateur satisfaisant} \textbf{Der 2}, \textit{que} $p_*, q_*$ \textit{et} $f_!$ \textit{existent, que} $f_!$ \textit{satisfait à} \textbf{Der 4} (d) \textit{pour le calcul des fibres. Conditions équivalentes :}

a) $f$ \textit{est une} $\D$\textit{-équivalence costricte, i.e.} (prop. 2) \textit{pour tout} $y \in \Ob\, Y$, $y\backslash X$ \textit{est} $\D$\textit{-contractile ;}

b) \textit{pour tout} $F \in \Ob\, \D(Y)$, \textit{l'homomorphisme suivant est} \texttt{[un]} \textit{isomorphisme}
\[
H^*_{\D}(Y, F) \isoto H^*_{\D}(X, f^*(F)) \ ;
\]

c) \textit{pour tout} $M \in \Ob(\mathcal{A} = \D(e))$, \textit{l'homomorphisme d'adjonction est} \texttt{[un]} \textit{isomorphisme}
\[
f_!(M_X) = f_!(f^*(M_Y)) \isoto M_Y \ .
\]

Je vais renuméroter les axiomes sur le prédérivateurs. % [sic] le prédérivateurs -> les prédérivateurs, article singulier devant un pluriel

\noindent\textbf{Der 1}, \textbf{Der 2} comme avant (p. 45).

\noindent\textbf{Der 3}\footnote[28]{Axiome des catégories d'indices asphériques, \textit{cf.} variantes \textbf{Der 3} a) b) c) page 97, et surtout pages 111-112.} : Soit $X$ dans $\Dia$, $p : X \to e$. Supposons $X$ contractile.\footnote[29]{N.B. Il suffit de supposer que $X$ a un objet final. Mais ce n'est pas un axiome autodual.} Alors
\[
p^* : M \longmapsto M_X \quad , \qquad p^* : \mathcal{A} = \D(e) \longrightarrow \D(X)
\]
est pleinement fidèle.

\printedpage{40}
% head: A. Grothendieck : Dérivateurs, Ch. I, Édition des Universités de Montpellier II et Paris VII
% foot: Version 8 juillet 2001
\msp{page 68}

\noindent\textbf{Der 4} (g) : existence des $f_*$.

\noindent\textbf{Der 4} (d) : existence des $f_!$.

\noindent\textbf{Der 5} (g) Calcul des fibres de $f_*(F)$.

\noindent\textbf{Der 5} (d) Calcul des fibres de $f_!(F)$.

\noindent L'ancien axiome \textbf{Der 5} (g) résulte de \textbf{Der 2}, \textbf{Der 3}, \textbf{Der 4} (g,d), \textbf{Der 5} (d), grâce au théorème précédent.

Un aspect intéressant du théorème précédent, si on l'envisage comme énonçant l'équivalence de a) et b), c'est que cet énoncé ne fait pas appel à des notions homologiques, mais que sa démonstration y fait appel, sous une forme relativement forte (calcul des fibres de $f_!(F)$). Pour les théorèmes de changement de base pour la $\D$-cohomologie relative, \textit{i.e.} pour les $f_*^{\D}$, on a besoin aussi constamment du calcul des fibres des $f_*(F)$. Donc il semble qu'on ne peut se passer d'aucun des axiomes précédents, y compris \textbf{Der 4} (d) et \textbf{Der 5} (d), même si on ne veut s'intéresser qu'au formalisme des $f^*, f_*$ (aspect cohomologique). Par la suite, on va donc supposer ces axiomes satisfaits.

\noindent\textbf{Théorème 2.} \textit{Soit} $f : X \to Y$ \textit{dans} $\Dia$. \textit{Conditions équivalentes :}

a) \textit{Pour tout} $y_0 \in \Ob\, Y$, \textit{le foncteur canonique}
\[
X_{y_0} \overset{\varphi}{\longrightarrow} X/y_0 \quad (\simeq X \times_Y Y/y_0) \ ,
\]
\textit{où} $X_{y_0}$ \textit{désigne la fibre en} $y_0$, \textit{est une} $\D$\textit{-équivalence co-stricte.}

b) \textit{Pour tout} $x_0 \in \Ob\, X$, \textit{le foncteur induit sur les catégories colocalisées en} $x_0$ \textit{et en} $y_0 = f(x_0)$
\[
x_0\backslash X \longrightarrow y_0\backslash Y
\]
\textit{est à fibres} $\D$\textit{-contractiles.}

\msp{page 69}

c) \textit{La formation des} $f_*F$ \textit{(}$F$ \textit{dans} $\D(X)$\textit{) commute à tout changement de base} $Y' \to Y$, \textit{et ça reste vrai après tout changement de base -- i.e. si on a} $Y'' \to Y' \to Y$, \textit{d'où}
\[
\begin{tikzcd}
X'' \arrow[r] \arrow[d, "f''"{swap}] & X' \arrow[r] \arrow[d, "f'"] & X \arrow[d, "f"] \\
Y'' \arrow[r, "p'"{swap}] & Y' \arrow[r, "p"{swap}] & Y
\end{tikzcd}
\]
\textit{le “théorème de changement de base” est vrai pour} $f'_*$ \textit{et} $p'^{*}$.

c$'$) \textit{La formation des} $f_*(F)$ \textit{(}$F \in \Ob\, \D(X)$\textit{) commute au passage aux fibres et ceci reste vrai après tout changement de base} $p : Y' \to Y$.

c$''$) \textit{Comme} c$'$) \textit{avec} $Y' \to Y$ \textit{un morphisme de localisation} $Y/y \to Y$.

d) \textit{Si} $Y'' \overset{p}{\to} Y' \to Y$, \textit{avec} $p$ \textit{coasphérique, alors} $q : X'' = X \times_Y Y'' \to X' = X \times_Y Y'$ \textit{est coasphérique.}

\printedpage{41}
% head: A. Grothendieck : Dérivateurs, Ch. I, Édition des Universités de Montpellier II et Paris VII
% foot: Version 8 juillet 2001

d$'$) \textit{Comme} d) \textit{avec} $Y'' = e$ (donc $Y'$ est catégorie avec objet final $e$).

d$''$) \textit{Comme} d) \textit{avec} $Y'' = \{1\}$, $Y' = \boldsymbol{\Delta}_1$.

.

\noindent\textsc{Démonstration}. a) $\Leftrightarrow$ b) cf. proposition page 53. Montrons que ces conditions impliquent c). Comme elles sont stables par changement de base (\textit{cf.} démonstration immédiate page 55) donc par \textbf{Der 2} on est ramené à prouver que $H^*_{\D}(X/y, F) \overset{\sim}{\to} H^*_{\D}(Y_y, F)$ $(F \in \Ob\, \D(Y))$, ce qui résulte du théorème 1.\footnote[30]{Donc on utilise \textbf{Der 5} (d).} (Cf. démonstration du théorème page 56, qui marche avec des hypothèses plus générales.)

On a c) $\Leftrightarrow$ c$'$) $\Rightarrow$ c$''$) (l'équivalence c) $\Leftrightarrow$ c$'$) par \textbf{Der 2}).

Prouvons c$''$) $\Rightarrow$ a).

\msp{page 70}

\noindent Considérons le diagramme à carrés cartésiens
\[
\begin{tikzcd}
X_{y_0} \arrow[r, "q"] \arrow[d, "p"{swap}] & X/y_0 \arrow[r] \arrow[d, "f_0"] & X \arrow[d, "f"] \\
\{y_0\} \arrow[r] & Y/y_0 \arrow[r] & Y
\end{tikzcd}
\qquad\qquad .
\]

\noindent L'hypothèse c$''$) prouve que pour tout $F$ dans $\D(X/y_0)$, on a $H^*_{\D}(X/y_0, F) \overset{\sim}{\to} H^*_{\D}(X_{y_0}, F)$. Par le théorème 1, ceci implique que $f$ est un $\D$-équivalence co-stricte. % [sic] un D-équivalence -> une D-équivalence, accord du genre

\noindent\textbf{Corollaire 1.} Énoncé dual du théorème, pour les foncteurs $X_{y_0} \to y_0\backslash X$, $X/x_0 \to Y/y_0$, et le changement de base des $j_!(F)$ \ldots

\noindent\textbf{Définition 4.} Sous les conditions équivalentes du théorème 1, on dit que $f$ est $\D$-\textit{propre.} La condition duale s'appelle $f$ $\D$-\textit{lisse.}

\noindent\textbf{Corollaire 2.} \textit{Soit} $f : X \to Y$ \textit{dans} $\Dia$. \textit{Si} $f$ \textit{est} $\D$\textit{-propre ou} $\D$\textit{-lisse, et à fibres} $\D$\textit{-contractiles, alors} $f$ \textit{est une} $\D$\textit{-équivalence stricte et costricte. Mieux,} $f$ \textit{est une} $\D$\textit{-équivalence} universelle \textit{i.e. est une} $\D$\textit{-équivalence et le reste après tout changement de base.}\footnote[31]{Se généralise : \textbf{Corollaire 2$'$.} \textit{Soit}
\[
\begin{tikzcd}[ampersand replacement=\&]
X \arrow[rr, "f"] \arrow[dr] \& \& Y \arrow[dl] \\
\& S \&
\end{tikzcd}
\]
\textit{avec} $X, Y$, $\D$\textit{-propres sur} $S$. \textit{Si} $f$ \textit{induit une} $\D$\textit{-équivalence sur les fibres c'est une} $\D$\textit{-équivalence. Itou dans le cas lisse.}\par
C'est la proposition 6 ci-dessous.}

\printedpage{42}
% head: A. Grothendieck : Dérivateurs, Ch. I, Édition des Universités de Montpellier II et Paris VII
% foot: Version 8 juillet 2001
\msp{page 71}

\noindent Il suffit de prouver que $f$ est une $\D$-équivalence. Supposons par exemple $f$ $\D$-propre. Il suffit de prouver que si $Y$ est $\D$-contractile, alors $X$ l'est (quitte à l'appliquer aux $Y/y_0$ après changement de base, pour montrer que les $X/y_0$ sont contractiles). Pour ceci, il suffit de prouver que $f$ est une $\D$-équivalence, et on sait que ceci est le cas si pour tout $F$ \textit{constant} sur $Y$, l'homomorphisme $F \overset{\text{adj}}{\longrightarrow} f_*f^*F$ est \texttt{[un]} isomorphisme. En fait, c'est vrai pour \textit{tout} $F$.

\noindent\textbf{Corollaire 3.} \textit{Si} $f : X \to Y$ \textit{est} $\D$\textit{-propre ou} $\D$\textit{-lisse, à fibres} $\D$\textit{-contractiles, alors} $f^* : \D(Y) \to \D(X)$ \textit{est pleinement fidèle, i.e. les homomorphismes d'adjonction}
\[
F \longrightarrow f_*f^*F \quad , \qquad f_!f^*F \longrightarrow F
\]
\textit{dans} $\D(Y)$ \textit{sont des isomorphismes.}

\noindent Il suffit de le voir pour $f$ propre, et cela résulte du calcul des fibres de $f_*f^*F$ \ldots

\noindent\textbf{Proposition 5.} \textit{La notion de foncteur} $\D$\textit{-propre} (resp. $\D$\textit{-lisse}) \textit{est stable par changement de base,} (trivial), \textit{et par composition} (utiliser \texttt{[le]} corollaire 2). \textit{Toute immersion fermée est} $\D$\textit{-propre, toute immersion ouverte} $\D$\textit{-lisse. Plus généralement, tout isomorphisme local est} $\D$\textit{-lisse, tout co-isomorphisme local est} $\D$\textit{-propre.}

\noindent\textbf{Proposition 6.} \textit{Soient} $X, Y$ \textit{sur} $S$, $f : X \to Y$ \textit{un} $S$\textit{-morphisme. Supposons} $X, Y$ \textit{propres sur} $S$, \textit{ou} $X$ \textit{et} $Y$ \textit{lisses sur} $S$, \textit{et que pour tout} $s \in \Ob\, S$, $f_s : X_s \to Y_s$ \textit{soit une} $\D$\textit{-équivalence, alors} $f$ \textit{aussi.}\footnote[32]{\textsc{Démonstration}. On prouvera même que sur la sous-catégorie des coefficients sur $Y$ qui proviennent de $\D(Y)$ \texttt{[plutôt de }$\D(S)$\texttt{]}, $f^* : \D(Y) \to \D(X)$ induit un foncteur pleinement fidèle. Si $f$ \texttt{[est]} propre, je l'interprète comme $q_*q^*F \overset{\sim}{\to} p_*p^*F$ pour tout $F$, et cela résulte du calcul de $p_*$, $q_*$ par passage aux fibres (propre), et de l'hypothèse sur les $f_s$. OK.}

\msp{page 72}

\noindent\textbf{Théorème 3.} \textit{Soit}
\[
\begin{tikzcd}
X \arrow[d, "f"{swap}] & X' \arrow[l, "q"{swap}] \arrow[d, "f'"] \\
Y & Y' \arrow[l, "p"{swap}]
\end{tikzcd}
\]
\textit{cartésien dans} $\Dia$. \textit{Alors le morphisme de changement de base}
\[
p^*f_* \longrightarrow f'_*q^*
\]
\textit{est un isomorphisme si} $f$ \textit{est} $\D$\textit{-propre, ou} $p$ $\D$\textit{-lisse.}

\noindent Le cas $\D$-propre déjà prouvé. Le cas $\D$-lisse : la démonstration p. 56 est OK.

\printedpage{43}
% head: A. Grothendieck : Dérivateurs, Ch. I, Édition des Universités de Montpellier II et Paris VII
% foot: Version 8 juillet 2001

\subsection*{Comparaison de $X$ et $X^{\circ}$.}

Soit $X$ \texttt{[une]} catégorie, on définit une catégorie $\Phi(X)$ et des foncteurs
\[
\begin{tikzcd}[column sep=small]
& & \Phi(X) \arrow[ddl, "p_X"{swap, name=P}] \arrow[ddr, "q_X"{name=Q}] & & \\
|[alias=L]| \text{\footnotesize\begin{tabular}{c}$\Cat$-cofibrant\\(donc propre)\end{tabular}} & & & & |[alias=R]| \text{\footnotesize\begin{tabular}{c}$\Cat$-fibrant\\(donc lisse)\\ \texttt{[plutôt }$\Cat$\texttt{-cofibrant}\\ \texttt{(donc propre)]}\end{tabular}} \\
& X & & X^{\circ} &
\arrow[from=P, to=L]
\arrow[from=Q, to=R]
\end{tikzcd}
\]
qui seront des $\hot$-équivalences, et des $\D$-équivalences si $\D$ est un dérivateur dont le domaine contient $X$ et $\Phi(X)$. On pose
\[
\Ob\, \Phi(X) = \Fl(X) \quad ,
\]
donc mêmes objets que $\mathcal{F}\ell(X) = \Hom(\boldsymbol{\Delta}_1, X)$, mais les morphismes sont différents :

\msp{page 73}
\[
\begin{tikzcd}
Y_1 \arrow[r, "u"] & Y_0 \arrow[d, "v_0"] \\
Y'_1 \arrow[u, "v_1"] \arrow[r, "u'"{swap}] & Y'_0
\end{tikzcd}
\]
un morphisme de $u$ dans $u'$ est un couple $v_0 : Y_0 \to Y'_0$, $v_1 : Y'_1 \to Y_1$ tel que $v_0uv_1 = u'$. Ainsi source et but $s(u)$ et $b(u)$ sont l'une contravariante, l'autre covariant en $u$. On prend pour $q$ le foncteur source, pour $p$ le foncteur but. % [sic] l'autre covariant -> l'autre covariante, accord avec « l'une contravariante »

Si on applique ceci à $X^{\circ}$, on trouve
\[
\begin{tikzcd}
& \Phi(X^{\circ}) \arrow[dl, "p_{X^{\circ}}"{swap}] \arrow[dr, "q_{X^{\circ}}"] & \\
X^{\circ} & & (X^{\circ})^{\circ} = X
\end{tikzcd}
\]
On a $\Ob\, \Phi(X^{\circ}) = \Fl(X^{\circ}) \simeq \Fl(X)$
\[
\begin{tikzcd}
Y_1 \arrow[r, "u"] & Y_0 \arrow[d, "v_1"] & \text{dans}\ \ X \\
Y'_1 \arrow[u, "v_0"] \arrow[r, "u'"{swap}] & Y'_0 & \text{dans}\ \ X
\end{tikzcd}
\]
et mêmes homomorphismes au sens de $\Phi(X)$ et au sens de $\Phi(X^{\circ})$. Donc
\[
\left\{
\begin{aligned}
&\Phi(X^{\circ}) = \Phi(X) \\
&p_{X^{\circ}} = q_X \\
&q_{X^{\circ}} = p_X \quad ,
\end{aligned}
\right.
\]


\endgroup
% END UNIT Der-ChI-234/Der-ChI-234_p041-045.tex
% BEGIN UNIT Der-ChI-234/Der-ChI-234_p046-050.tex
\begingroup\makeatletter%
\NoAutoSpacing%
\newcommand{\printedpage}[1]{\marginpar{\footnotesize p.~#1}}%
\newcommand{\msp}[1]{\marginpar{\footnotesize [#1]}}%
\newcommand{\D}{\mathbb{D}}%
\newcommand{\Ob}{\mathop{\mathsf{Ob}}}%
\newcommand{\Dia}{\textsf{Dia}}%
\newcommand{\bs}{\backslash}%
\makeatother\let\printedpage\derpp\let\msp\dermsp%
\let\mspin\DERmspin%


\printedpage{44}
% head: A. Grothendieck : Dérivateurs, Ch. I, Édition des Universités de Montpellier II et Paris VII
le passage de $X$ à $X^\circ$ revient à échanger $p_X$ et $q_X$. La fibre de $\Phi(X)$ en $x \in \Ob X$ est $\simeq (X/x)^\circ$, celle en $y^\circ \in \Ob X^\circ$ est $\simeq y\bs X$
\[
\left\{
\begin{aligned}
&p_X^{-1}(\{x\}) \simeq (X/x)^\circ\\
&q_X^{-1}(\{y^\circ\}) \simeq y\bs X
\end{aligned}
\right.
\]

N.B. : $p_X$ est $W$-lisse \texttt{[plutôt $W$-propre]} à fibres $W$-asphériques, donc dans $W^u$ \texttt{[$W^u =$ ensemble des morphismes qui sont dans $W$ et le restent après tout changement de base]}, $q_X$ est $W$-propre à fibres $W$-asphériques, donc dans $W^u$.

\msp{page 74}

On vérifie
\[
\left\{
\begin{aligned}
&\Phi(X)/x \simeq \Phi(X/x) &\qquad& x \in \Ob X\\
&\Phi(X)/x^\circ \simeq \Phi(x\bs X) &\qquad& x^\circ \in \Ob X^\circ
\end{aligned}
\right.
\]

Je dis que $p_X$, $q_X$ sont des $\D$-\textit{équivalences strictes} et \textit{costrictes.} Par les formules précédentes, cela signifie que les $\Phi(X/x)$, $\Phi(x\bs X)$ sont $\D$-contractiles. Donc que si $Z$ a un objet final, $\Phi(Z)$ est $\D$-contractile. Mais considérons
\[
\begin{tikzcd}
\Phi(Z) \arrow[dr, "q_Z"] & \\
& Z^\circ
\end{tikzcd}
\quad ;
\]
il suffit de prouver que c'est une $\D$-équivalence stricte (puisque $Z^\circ$, ayant un objet initial, est $\D$-contractil), donc que les $z_0\bs\Phi(Z)$ sont $\D$-contractiles. On a par les formules précédentes % [sic] contractil -> contractile, adjectif ailleurs écrit « contractile »
\[
z_0\bs\Phi(Z) \simeq \Phi(z_0\bs Z) = \Phi(T) \quad ,
\]
où cette fois $T$ a un objet initial $\emptyset$ et un objet final $e$. Considérons alors l'objet
\[
\emptyset \longrightarrow e
\]
de $\Phi(T)$, c'est un objet final, car pour tout $x \to y$ dans $\Phi(T)$, il y a une unique flèche
\[
\begin{tikzcd}
x \arrow[r] & y \arrow[d] \\
\emptyset \arrow[u] \arrow[r] & e
\end{tikzcd}
\]
% foot: Version 8 juillet 2001

\printedpage{45}
% head: A. Grothendieck : Dérivateurs, Ch. I, Édition des Universités de Montpellier II et Paris VII
\msp{page 75}

Pour résumer

\textbf{Proposition 1 :} \textit{$p_X$ et $q_X$ sont des} \textsf{hot}\textit{-équivalences strictes,} i.e. \textsf{Hot}\textit{-asphériques (donc aussi des $\D$-équivalences strictes si $\Phi(X)$ est dans le domaine de $\D$) à fibres les $(X/x)^\circ$ et $y\bs X$ respectivement (donc à fibres contractiles).}

\textbf{Corollaire 1.} \textit{Pour que $X$ soit $\D$-contractile, il faut et il suffit que $X^\circ$ le soit.}

\hspace*{\parindent}Plus généralement :

\textbf{Corollaire 2.} \textit{Soit $f : X \to Y$. Pour que $f$ soit une $\D$-équivalence, il faut et il suffit que $f^\circ$ le soit.}

\hspace*{\parindent}On utilisera que le diagramme
\[
\begin{tikzcd}
& \Phi(X) \arrow[dl, "p_X"'] \arrow[dr, "q_X"] & \\
X & & X^\circ
\end{tikzcd}
\]
est fonctoriel en $X$. On a donc
\[
\begin{tikzcd}
& \Phi(X) \arrow[rr, "\Phi(f)"] \arrow[dl, "p_X"'] & & \Phi(Y) \arrow[dl, "p_Y"] \arrow[dd, "q_Y"] \\
X \arrow[rr, "f" near end] \arrow[from=1-2, to=3-2, "q_X"' near end, crossing over] & & Y & \\
& X^\circ \arrow[rr, "f^\circ"'] & & Y^\circ
\end{tikzcd}
\]
les deux carrés commutatifs. Comme $p_X$ et $p_Y$ sont des $\D$-équivalences, $f$ est une $\D$-équivalence si et seulement si $\Phi(f)$ l'est, et de même pour $f^\circ$. Ces raisonnements supposent que si $X$ est dans le domaine de $\D$, $\Phi(X)$ aussi. C'est le cas si \Dia{} est formé des catégories ordonnées finies.
% foot: Version 8 juillet 2001

\printedpage{46}
% head: A. Grothendieck : Dérivateurs, Ch. I, Édition des Universités de Montpellier II et Paris VII
\msp{page 76}

\subsection*{$\D$-fibrations}

\textbf{Proposition 1.} \textit{Soit $f : X \to Y$ dans} \Dia\textit{, conditions équivalentes :}

\hspace*{\parindent}a) \textit{Pour tout $y \in Y$, l'inclusion $X_y \to X/y$ est une $\D$-équivalence, et pour toute flèche $y_1 \to y_0$ de $Y$, $X/y_1 \to X/y_0$ est une $\D$-équivalence.}

\hspace*{\parindent}b) \textit{Pour tout coefficient} constant \textit{$F$ sur $Y$, $f_*f^*F$ est localement constant,}\footnote[33]{Par là on entend que pour tout $u : y \to y'$ dans $Y$, $G_{y'} \to G_y$ est \texttt{[un]} isomorphisme. Sans axiomes supplémentaires sur $\D$, il n'est pas clair que cela implique que les ``restrictions'' de $G$ aux $Y/y$ ou $y\bs Y$, soient isomorphes à des coefficients constants. Voir cependant l'axiome \textbf{Der 3} a$'$), page 97, qui assure que c'est pareil.} \textit{et ``se calcule fibre par fibre''} i.e. \textit{l'homomorphisme canonique}
\[
f_*f^*(F)_y \overset{\sim}{\longrightarrow} H^*_{\D}(X_y, F|X_y)
\]
\texttt{[est un isomorphisme]}.

Ces conditions sont équivalentes, et impliquent les conditions de (b) pour tout $F$ \textit{localement} constant sur $Y$.

La première condition dans (a) signifie que pour tout $y \in Y$, et tout $M \in \Ob \mathcal{A}$,
\[
H^*_{\D}(X/y, M) \longrightarrow H^*_{\D}(X_y, M)
\]
est un isomorphisme. Or cette flèche s'interprète comme
\[
f_*(M_X)_y \longrightarrow H^*_{\D}(X_y, M)
\]
est un isomorphisme, \textit{i.e.} que $f_*(M_X)$ se calcule fibre par fibre. Comme $M_X \simeq f^*(M_Y)$, c'est la deuxième condition (b). La deuxième condition dans (a), pour $y_1 \to y_0$ signifie que pour $M$ dans $\mathcal{A}$,
\[
H^*_{\D}(X/y_0, M) \longrightarrow H^*_{\D}(X/y_1, M)
\]
est \texttt{[un]} isomorphisme, or cette flèche s'interprète comme
\[
f_*(M_X)_{y_0} \longrightarrow f_*(M_X)_{y_1} \quad ,
\]
que ce soit toujours un isomorphisme, c'est la première condition (b). Que (b) soit valable aussi pour $F$ localement constant sur $Y$ est immédiat, car c'est une condition de nature locale par rapport à $Y$.

\msp{page 77}

La proposition duale fait appel aux foncteurs
\[
y_0\bs X \longrightarrow y_1\bs X \qquad \text{et} \qquad X_y \longrightarrow y\bs X \quad ,
\]
et aux faisceaux $f_!f^*M$, pour $M$ constant, ou localement constant sur $Y$.
% foot: Version 8 juillet 2001

\printedpage{47}
% head: A. Grothendieck : Dérivateurs, Ch. I, Édition des Universités de Montpellier II et Paris VII
\textbf{Définition 1.} Si les conditions équivalentes de la proposition 1 sont satisfaites, on dit que $f$ est une $\D$-\textit{fibration cohomologique} (ou $\D$-\textit{fibration gauche}). Si les conditions duales sont satisfaites, on dit que c'est une $\D$-\textit{fibration homologique} (ou $\D$-\textit{fibration droite}).

N.B. Les $\D$-fibrations cohomologiques sont stables par tout changement de base qui est un isomorphisme local. Les $\D$-fibrations homologiques par tout changement de base qui est un co-isomorphisme local. Mais ce ne sont pas des notions stables par \textit{tout} changement de base\footnote[34]{? donner exemple!}. Mais la notion de $\D$-fibration cohomologique est stable par changement de base $\D$-lisse, et celle de $\D$-fibration homologique par changement de base $\D$-propre.

\hspace*{\parindent}Je vais essayer de construire un exemple où le changement de base canule, en posant $Y = \boldsymbol{\Delta}_1$. Se donner $X$ sur $\boldsymbol{\Delta}_1$ revient à se donner une sous-catégorie ouverte $X_0$, et son complémentaire fermé $X_1$
\[
\begin{tikzcd}
& X & \\
X_1 \arrow[ur, hook] \arrow[d, dash, dashed] & & X_0 \arrow[ul, hook'] \arrow[d, dash, dashed] \\
1 & & 0 \arrow[ll]
\end{tikzcd}
\quad .
\]
Dire que $X \to Y = \boldsymbol{\Delta}_1$ est une $\D$-fibration cohomologique signifie simplement
\msp{page 78}
que les restrictions
\[
X_0 \overset{j}{\longrightarrow} X \quad , \qquad X_1 \overset{i}{\longrightarrow} X
\]
sont des $\D$-équivalences. Prenons un changement de base $Y' \to Y = \boldsymbol{\Delta}_1$, \textit{i.e.} un $Y'$ avec $Y'_0, Y'_1$
\[
\begin{tikzcd}
& Y' & \\
Y'_1 \arrow[ur, hook] \arrow[d, dash, dashed] & & Y'_0 \arrow[ul, hook'] \arrow[d, dash, dashed] \\
1 & & 0 \arrow[ll]
\end{tikzcd}
\quad .
\]
On voit que la condition $Y'$ lisse sur $\boldsymbol{\Delta}_1$ signifie que pour tout $y'_1 \in \Ob Y'_1$, la catégorie $Y'_0/y'_1$ des objets $y_0$ de $Y'_0$ au dessus de $y'_1$ (\textit{i.e.} munis de $y_0 \to y'_1$) est contractile. On voit que si $Y'_1 \neq \emptyset$ alors $Y'_0 \neq \emptyset$ -- mais un exemple avec $Y'_0 = \emptyset$, où $Y'_1 = \emptyset$, n'est pas intéressant
% foot: Version 8 juillet 2001
\printedpage{48}
% head: A. Grothendieck : Dérivateurs, Ch. I, Édition des Universités de Montpellier II et Paris VII
(car une projection $I \times J \to I$ est toujours une $\D$-fibration cohomologique). Prenons par exemple pour $Y'$
\[
\begin{tikzcd}
& a \arrow[dl] \\
b & \\
& a' \arrow[ul] \\
\underbrace{\qquad}_{\displaystyle Y'_1} & \underbrace{\qquad}_{\displaystyle Y'_0}
\end{tikzcd}
\quad ,
\]
somme amalgamée de deux copies de $Y = \boldsymbol{\Delta}_1$, donc $X' = X \times_Y Y'$ peut se voir comme somme amalgamée $X \amalg_{X_1} X$. Mais ça semble \textit{pas} être un contre-exemple. Il faudrait que $X_1 \to X'$ ne soit \textit{pas} une $\D$-équivalence.

\msp{page 79}

\textbf{Proposition 2.} \textit{Soit $f : X \to Y$ dans} \Dia. \textit{Conditions équivalentes :}

\hspace*{\parindent}a) \textit{Pour tout $Y'' \overset{p}{\to} Y' \to Y$, d'où}
\[
\begin{tikzcd}
X'' \arrow[r, "q"] \arrow[d] \arrow[dr, phantom, "\text{cart}"] & X' \arrow[r] \arrow[d] \arrow[dr, phantom, "\text{cart}"] & X \arrow[d, "f"] \\
Y'' \arrow[r, "p"'] & Y' \arrow[r] & Y
\end{tikzcd}
\quad ,
\]
\textit{si $p$ est une $\D$-équivalence,} \texttt{[alors]} \textit{$q$ aussi.}\footnote[35]{Il faut supposer $Y', Y''$ \textsf{Hot}-asphériques, cf. page 96, et supposer vérifié l'axiome Der 3 $a'$) de la page 97.}

\hspace*{\parindent}b) \textit{$f$ est une $\D$-fibration cohomologique ``universelle'',} i.e. \textit{est une fibration cohomologique, et le reste après tout changement de base.}

\hspace*{\parindent}b$'$) \textit{Condition duale de} (b) \textit{: $f$ est une $\D$-fibration homologique universelle.}\footnote[36]{b) $\Rightarrow$ a) \underline{canulée}, \textit{cf.} annotation marginale,\footnotemark[37] \textit{cf.} contre-exemple page 89. Mais ça marche pour les dérivateurs $\D_i$ ($i \in \mathbb{N}$) associés aux ``localisateurs fondamentaux'' $W_i$, \textit{i.e.} à la théorie des $i$-types d'homotopie.}\footnotetext[37]{L'implication b) $\Rightarrow$ a) pas prouvée, seulement sous forme plus faible :\par \hspace*{\parindent}a$'$) si $Y''$, $Y'$ sont $\D$-asphériques et 1-connexes, alors $X'' \to X'$ est une $\D$-équivalence. (Il suffit même de le supposer pour $Y'$, $Y''$ ayant des objets finaux (pour impliquer b) ou ayant des objets initiaux (pour impliquer b$'$).)}

\textit{a}) $\Rightarrow$ \textit{b}) : Supposons \textit{a}), soit $Y' \to Y$ changement de base, d'où
\[
\begin{tikzcd}
X' \arrow[r] \arrow[d, "f'"'] & X \arrow[d, "f"] \\
Y' \arrow[r] & Y
\end{tikzcd}
\quad .
\]
% foot: Version 8 juillet 2001


\endgroup
% END UNIT Der-ChI-234/Der-ChI-234_p046-050.tex
% BEGIN UNIT Der-ChI-234/Der-ChI-234_p051-055.tex
\begingroup\makeatletter%
\newcommand{\printedpage}[1]{\marginpar{\footnotesize p.~#1}}%
\newcommand{\msp}[1]{\marginpar{\footnotesize [#1]}}%
\newcommand{\DD}{\mathbb{D}}%
\newcommand{\Ob}{\mathsf{Ob}}%
\newcommand{\Cat}{\mathcal{C}at}%
\newcommand{\Dia}{\mathsf{Dia}}%
\makeatother\let\printedpage\derpp\let\msp\dermsp%
\let\mspin\DERmspin%


\printedpage{49}
% head: A. Grothendieck : Dérivateurs, Ch. I, Édition des Universités de Montpellier II et Paris VII
% foot: Version 8 juillet 2001

Pour prouver que $f'$ est une $\DD$-fibration cohomologique, on regarde au dessus de $Y'$ les foncteurs de la forme $\{y'\} \to Y'/y'$ ($y' \in \Ob\, Y'$) et $Y'/y'_1 \to Y'/y'_0$ (si $u' : y'_1 \to y'_0$ flèche de $Y'$), et on voit que les flèches dans $\Cat/Y'$ déduites de ceux-ci par le changement de base $X' \xrightarrow{f'} Y'$ restent des $\DD$-équivalences

\msp{page 80}

\[
\begin{tikzcd}[column sep=large, row sep=large]
X'_1 \arrow[r, "q"] \arrow[d, "f'_1"'] \arrow[dr, phantom, "\text{cart}"] & X'_0 \arrow[r] \arrow[d, "f'_0"] \arrow[dr, phantom, "\text{cart}"] & X' \arrow[r] \arrow[d, "f'"] \arrow[dr, phantom, "\text{cart}"] & X \arrow[d, "f"] \\
Y'_1 \arrow[r, "p"'] & Y'_0 \arrow[r] & Y' \arrow[r] & Y & .
\end{tikzcd}
\]
Comme $p$ est une $\DD$-équivalence, cela résulte bien de a).

\indent\textit{b}) $\Rightarrow$ \textit{a}) : Comme la condition \textit{b}) est évidemment stable par changement de base, on est ramené à prouver que si $f$ satisfait \textit{b}), et si on a
\[
\begin{tikzcd}[column sep=large, row sep=large]
X \arrow[d, "f"'] \arrow[dr, phantom, "\text{cart}"] & X' \arrow[l, "q"'] \arrow[d, "f'"] \\
Y & Y' \arrow[l, "p"]
\end{tikzcd}
\tag*{$(*)$}
\]
si $p$ est une $\DD$-équivalence, alors $q$ aussi, \textit{i.e.}
\[
\begin{tikzcd}[row sep=small]
H^*_{\DD}(X,M) \arrow[r, "\sim"] & H^*_{\DD}(X',M) \\
\wr| & \wr| \\
H^*_{\DD}(Y, f_*(M_X)) & H^*_{\DD}(Y', f'_*(M_{X'}))
\end{tikzcd}
\]
pour tout $M$ dans $\mathcal{A} = \DD(e)$. Mais $F = f_*(M_X)$ est localement constant et se calcule “fibre par fibre”, de même pour $F' = f'_*(M_{X'})$.

\msp{page 81}

L'équivalence \textit{a}) $\Leftrightarrow$ \textit{b}$'$) se déduit de \textit{a}) $\Leftrightarrow$ \textit{b}) par dualité (la condition \textit{a}) étant autoduale), ou en reprenant le même raisonnement pour \textit{a}) $\Rightarrow$ \textit{b}$'$), et pour \textit{b}$'$) $\Rightarrow$ \textit{a}) aussi (en remplaçant la cohomologie par l'homologie).

\textbf{Définition 2.} Si les conditions équivalentes de la proposition 2 sont satisfaites, on dit que $f$ est une $\DD$-\textit{fibration}.

N.B. C'est une notion \textit{à la fois} locale et “colocale” sur $Y$ (alors que fibration gauche est locale, fibration droite colocale).

\printedpage{50}
% head: A. Grothendieck : Dérivateurs, Ch. I, Édition des Universités de Montpellier II et Paris VII
% foot: Version 8 juillet 2001

\textbf{Proposition 3.} a) \textit{La notion de $\DD$-fibration est stable par changement de base et par composition.}\footnote[38]{La stabilité par composition \textit{n'est} pas claire.}

\indent b) \textit{Soient}
\[
\begin{tikzcd}[column sep=normal, row sep=large]
X \arrow[rr, "f"] \arrow[dr, "p"'] & \phantom{X} & X' \arrow[dl, "p'"] \\
& Y &
\arrow[from=1-2, to=2-2, phantom, "\text{\footnotesize comm}"]
\end{tikzcd}
\]
\textit{avec $p,p'$ des $\DD$-fibrations (il suffit que $p,p'$ soient des $\DD$-fibrations cohomologiques, ou que ce soient des $\DD$-fibrations homologiques, ou que $p,p'$ soient propres, ou $p,p'$ lisses). Si $f$ induit des $\DD$-équivalences dans les fibres, alors $f$ est une $\DD$-équivalence.}\footnote[39]{Peut se préciser considérablement dans le cas où $p$, $p'$ sont des $\DD$-fibrations, \textit{cf.} prop. 6, page 84.}

\indent\textsc{Démonstration.} \textit{a}) est immédiat sous la condition \textit{a}) (prop. 2) de la notion de $\DD$-fibration. Pour \textit{b}), supposons que $f$ induise des $\DD$-équivalences sur les fibres. On veux prouver que % [sic] On veux -> On veut, verb agreement
\[
H^*_{\DD}(X,M) \longrightarrow H^*_{\DD}(X',M)
\]
est un isomorphisme, pour tout $M$ dans $\mathcal{A}$, mais cette flèche s'identifie à
\[
H^*_{\DD}(Y, p_*(M_X)) \longrightarrow H^*_{\DD}(Y, p'_*(M_{X'}))
\]
déduite de
\[
p_*(M_X) \longrightarrow p'_*(M_{X'}) \qquad\qquad (\text{NB}\ \ M_X \simeq f^*M_{X'}\,)
\]
\msp{page 82}
et il suffit de prouver que cette flèche est un isomorphisme. Pour ceci, on regarde fibre par fibre, et on utilise le fait que “les images directes se calculent fibre par fibre”.

\textbf{Question.} Soit $f : X \to Y$ une $\DD$-fibration, et soit $F \in \DD(Y)$. Alors $f_*f^*F$ se calcule-t-il par passage aux fibres? Même question pour $f_!f^*F$. Pour la première propriété, concernant $f_*$, on montre que cela marche si $F$ est de la forme $p_*(M_Y)$, où $p : Y' \to Y$ est $\DD$-\textit{propre,} et $M \in \Ob\, \mathcal{A}$. D'autre part, la sous-catégorie pleine $C$ de $\DD(Y)$ formée des $F$ pour lesquels ça marche est (sauf erreur) “stable par $\DD$-$\mathsf{hot}$-$\varprojlim$” en le sens suivant : si on a une catégorie $I$ avec objet final $e$ formant une sous-catégorie fermée (\textit{i.e.} $e \to x \Rightarrow x = e$), et si on a $\mathcal{F} \in \DD(Y \times I)$, d'où $I^\circ \to \DD(Y)$, $i \mapsto \mathcal{F}_i$ (diagramme squelettique), si les $\mathcal{F}_i$ pour $i \neq e$ sont dans $C$, et si $\mathcal{F}$ est dans l'image essentielle de $j_*$, où $j : Y \times (I \setminus \{e\}) \hookrightarrow Y \times I$ est l'inclusion, alors $\mathcal{F}_e$ est aussi dans $C$.

\indent On peut cependant considérer que si on veux que $f : X \to Y$ satisfasse à toutes les propriétés cohomologiques et homologiques qui s'attachent à la notion d'une “fibration” (“localement triviale”), on devrait mettre dans les conditions que $f$ est \textit{lisse}, et \textit{propre}, pour que les changements de base par $f$ commutent aux foncteurs $g_*$, $g_!$. % [sic] on veux -> on veut, verb agreement

\printedpage{51}
% head: A. Grothendieck : Dérivateurs, Ch. I, Édition des Universités de Montpellier II et Paris VII
% foot: Version 8 juillet 2001

\msp{page 83}

Signalons :

\textbf{Proposition 4.} \textit{Soit $f : X \to Y$} propre. \textit{Pour que $f$ soit une $\DD$-fibration, il faut et il suffit que pour tout $y_1 \to y_0$ dans $Y$, le foncteur $X/y_1 \to X/y_0$ induit soit une $\DD$-équivalence} (i.e. \textit{que $f$ soit une $\DD$-fibration} cohomologique). \textit{Dualement, si $f$ est $\DD$-lisse, pour que ce soit une $\DD$-fibration, il faut et il suffit que pour tout $y_1 \to y_0$ dans $Y$, $y_0\backslash X \to y_1\backslash X$ soit une $\DD$-équivalence.}

\indent On pourrait appeler $\DD$-fibrations \textit{parfaites} les foncteurs à la fois propres et lisses qui sont des $\DD$-fibrations cohomologiques (ou ce qui revient au même des $\DD$-fibrations homologiques).

\textbf{Corollaire.} \textit{Supposons $f$ $\Cat$-cofibration (donc propre). Pour que $f$ soit une $\DD$-fibration, il faut et il suffit que les foncteurs $X_{y_0} \to X_{y_1}$ (pour $y_0 \to u_1$ dans $Y$) soient des $\DD$-équivalences. Duallement, pour $\Cat$-fibration (donc lisse).} % [sic] u_1 -> y_1, the arrow is y_0 -> y_1 ; [sic] Duallement -> Dualement, spelling

\textbf{Proposition 5.} \textit{Soit $f : X \to Y$ dans} $\Dia$. \textit{Conditions équivalentes :}
\begin{enumerate}[label={}, leftmargin=2em, labelsep=0pt]
\item a) \textit{$f$ est une $\DD$-équivalence, et le reste par tout changement de base ;}
\item b) \textit{$f$ est une $\DD$-fibration à fibres $\DD$-contractiles.}
\end{enumerate}

\textsc{Démonstration.} \textit{a}) $\Rightarrow$ \textit{b}) : Que \textit{a}) implique que $f$ soit une $\DD$-fibration est immédiat sur la définition de celle-ci (condition \textit{a}) de la proposition 2), et que les fibres sont contractiles est tautologique.

\textit{b}) $\Rightarrow$ \textit{a}) : Il suffit de montrer que \textit{b}) implique que $f$ est une $\DD$-équivalence. Mais c'est un cas particulier de la proposition 3 \textit{b}) (en y prenant $X' = Y$).

\msp{page 84}

Quand ces conditions sont satisfaites, on s'attend à ce que le foncteur
\[
f^* : \DD(Y) \longrightarrow \DD(X)
\]
soit pleinement fidèle, \textit{i.e.} que les morphisme d'adjonction % [sic] les morphisme -> les morphismes, plural agreement
\[
F \longrightarrow f_*f^*F \quad , \qquad f_!f^* \longrightarrow F % [sic] f_!f^* -> f_!f^*F, the object F is missing
\]
dans $\DD(Y)$ soient des isomorphismes. (C'est vrai si $f$ est $\DD$-propre ou $\DD$-lisse, cf. cor. 3, p.71.) Cela revient encore à dire que les $f_*f^*F$ se calculent par passage aux fibres, ou encore itou pour les $f_!f^*F$. Cela nous ramène à la question de la page 82.

\textbf{Définition 3.} On dit alors que $f$ est une $\DD$-fibration \textit{triviale}.

\textbf{Proposition 6} (précise prop. 3 page 81). \textit{Considérons un triangle commutatif dans} $\Dia$
\[
\begin{tikzcd}[column sep=small, row sep=large]
X \arrow[rr, "f"] \arrow[dr, "p"'] & & X' \arrow[dl, "p'"] \\
& S & & & .
\end{tikzcd}
\]

\printedpage{52}
% head: A. Grothendieck : Dérivateurs, Ch. I, Édition des Universités de Montpellier II et Paris VII
% foot: Version 8 juillet 2001

\indent a) \textit{Supposons que $p$ et $p'$ soient des $\DD$-fibrations gauches, ou que $p$ et $p'$ soient des $\DD$-fibrations droites. Soit $S_0$ une partie de $S$ qui rencontre toute composante connexe de $S$. Pour que $f$ soit une $\DD$-équivalence, il} suffit \textit{que pour tout $s_0 \in S_0$, $f$ induise une $\DD$-équivalence $X_{s_0} \to Y_{s_0}$.}

\indent b) \textit{Cette condition est aussi nécessaire}\footnote[40]{Pas prouvé!\footnotemark[41] Même si $X' = S$, ce n'est pas clair, même si $\DD = \mathsf{HOTAB}$, $f$ propre et lisse. Mais OK si $S$ est $\DD$-contractile.}\footnotetext[41]{\underline{Faux}, \textit{cf.} contre-exemple, p. 87.} \textit{si $p$ et $p'$ sont des $\DD$-fibrations. Donc dans ce cas, $f$ est une $\DD$-équivalence si et seulement si $f$ induit des $\DD$-équivalences fibre par fibre.}

\msp{page 85}

\textsc{Démonstration.} a) Il suffit, en vertu de la proposition 3, de montrer que l'hypothèse implique que pour \textit{tout} $s \in S$, $f$ induit une $\DD$-équivalence $f_s : X_s \to X'_s$. Cela résulte aussitôt du lemme :

\textbf{Lemme.} \textit{Si $p$ et $p'$ sont toutes deux soit des $\DD$-fibrations gauches, soit des $\DD$-fibrations droites, et si $s_1 \to s_0$ est une flèche dans $S$, alors $f_{s_1}$ est une $\DD$-équivalence si et seulement si $f_{s_0}$ l'est. (Donc l'ensemble $S' \subset \Ob\, S$ des $s \in S$ tels que $f_s$ soit une $\DD$-équivalence est une réunion de composantes connexes de $S$.)}

\indent Supposons par exemple que $p,p'$ soient des $\DD$-fibrations gauches. Considérons, pour $s \in S$, le diagramme
\[
\begin{tikzcd}[row sep=large]
X_s \arrow[r, "f_s"] \arrow[d] & X'_s \arrow[d] \\
X/s \arrow[r, "f/s"'] & X'/s & .
\end{tikzcd}
\]
Les flèches verticales sont des $\DD$-équivalences par l'hypothèse sur $p,p'$, donc $f_s$ est une $\DD$-équivalence si et seulement si $f/s$ l'est. Considérons maintenant
\[
\begin{tikzcd}[row sep=large]
X/s_1 \arrow[r, "f/s_1"] \arrow[d] & X'/s_1 \arrow[d] \\
X/s_0 \arrow[r, "f/s_0"'] & X'/s_0 & .
\end{tikzcd}
\]
Ici encore, l'hypothèse sur $p,p'$ implique que les flèches verticales sont des $\DD$-équivalences, donc $f_{s_1}$ est une $\DD$-équivalence si et seulement si $f_{s_0}$ l'est.

\msp{page 86}

Si $Y' = S$, l'énoncé (non prouvé) dit que si $f : X \to S$ est une $\DD$-fibration et une $\DD$-équivalence, alors les fibres sont $\DD$-contractiles. Ce n'est pas sûr même pour une projection $X = T \times S \to S$. Si $\eta = \mathit{hot}_0(S)$, $\xi = \mathit{hot}_0(T)$, la question se réduit à celle-ci : Si pour tout $M$ dans $\mathcal{A}$, on a : l'homomorphisme
\[
\alpha_M : \xi \times_{\DD} M \longrightarrow M
\]
\printedpage{53}
% head: A. Grothendieck : Dérivateurs, Ch. I, Édition des Universités de Montpellier II et Paris VII
% foot: Version 8 juillet 2001
satisfait
\[
\eta \times_{\DD} \alpha_M : \eta \times_{\DD} (\xi \times_{\DD} M) \xrightarrow{\sim} \eta \times_{\DD} M
\]
est un isomorphisme, alors $\alpha_M$ est lui même un isomorphisme.

\msp{page 87}

Voici un \textit{contre-exemple}\footnote[42]{Voir un contre-exemple encore plus simple pages 93, 94. (Mais il a l'inconvenient de ne pas être à fibres connexes.) % [sic] inconvenient -> inconvénient, missing accent
} que je donne d'abord dans le contexte des espaces topologiques, au lieu de $\Cat$. Je vais construire une fibration localement triviale et non triviale
\[
\begin{tikzcd}[column sep=0.2em, row sep=large]
X \arrow[d, "f"'] & \supset S^1 & \\
S & = S^1 & \qquad ,
\end{tikzcd}
\]
de fibre $S^1$, telle que $f$ soit une $\DD$-équivalence pour tout $\DD = \mathsf{Der}(k\text{-}\mathsf{Mod})$, $k$ étant un anneau où 2 est inversible. Cela signifie que
\[
H_*(X,k) \longrightarrow H_*(S,k)
\]
est un isomorphisme modulo 2-torsion (noyau et conoyau sont annulés par une puissance de 2 - en fait par 2 ici), ou encore $H_*(X,k) \xrightarrow{\sim} H_*(S,k)$ si 2 inversible dans $k$. Mais la fibre $S^1$ n'est pas $\DD$-contractile (sauf si $k = 0$!).\footnote[43]{Exemple beaucoup plus simple encore : on prend $X = S^1$, et $X \to S$ le revêtement à deux feuillets. Si on prend le revêtement à $n$ feuillets, c'est aussi OK si $k$ anneau où $n$ est inversible.}

\indent La fibration est donnée à isomorphisme unique près, en se fixant un point $s \in S$, et en se donnant un automorphisme $u$ de la fibre $X_s$ (“monodromie”). (La classe d'isomorphie du fibré ne dépend que de la classe d'isotopie de $u$). Je prend pour $u$ la symétrie (inversant l'orientation) $z \mapsto z^{-1}$ du cercle standard $\{z \in \mathbb{C} \mid |z| = 1\}$. Cette fibration “se réalise” par des ensembles ordonés $X \to S$, $S^1$ étant représenté (sur la base, comme sur les fibres) par % [sic] Je prend -> Je prends, verb agreement ; [sic] ordonés -> ordonnés, spelling
\[
\begin{tikzcd}[column sep=normal, row sep=normal]
& & a & & \\
S \ = & 0 \arrow[ur] \arrow[dr] & & 0' \arrow[ul] \arrow[dl] & = \ S' \amalg_R S'' \quad , \\
& & b & &
\end{tikzcd}
\]
où $R$ est la catégorie discrète ayant deux éléments $a$, $b$. On définit $X$ comme $X' \amalg_Y X''$, où $X' = X|S'$, $X'' = X|S''$, $Y = X|R$.\footnote[44]{N.B. $S'$, $S''$, $R = S' \cap S''$ fermés dans $S$, $X'$, $X''$, $Y = X' \cap X''$ fermés dans $X$.}


\endgroup
% END UNIT Der-ChI-234/Der-ChI-234_p051-055.tex
% BEGIN UNIT Der-ChI-234/Der-ChI-234_p056-060.tex
\begingroup\makeatletter%
\newcommand{\printedpage}[1]{\marginpar{\footnotesize p.~#1}}%
\newcommand{\msp}[1]{\marginpar{\footnotesize [#1]}}%
\newcommand{\Cat}{\mathcal{C}\mathit{at}}%
\newcommand{\D}{\mathbb{D}}%
\newcommand{\RR}{\mathbb{R}}%
\makeatother\let\printedpage\derpp\let\msp\dermsp%
\let\mspin\DERmspin%


\printedpage{54}
% head: A. Grothendieck : Dérivateurs, Ch. I, Édition des Universités de Montpellier II et Paris VII

\msp{page 88}
On pose
\[ X' = S' \times \bar S \qquad\qquad X'' = S'' \times \bar S \]
($\bar S$ copie conforme de $S$, cf ci-dessous) et il faut, pour recoller $X'$ et $X''$, se donner un isomorphisme entre $X'|R \simeq R \times \bar S = \bar S \amalg \bar S$ et $X''|R \simeq R \times \bar S = \bar S \amalg \bar S$ qui soit un $R$-foncteur, \textit{i.e.} qui provient de deux isomorphismes
\[
\begin{aligned}
\bar S &\xrightarrow{u_a} \bar S \\
\bar S &\xrightarrow{u_b} \bar S \qquad .
\end{aligned}
\]
On prend
\[
\begin{aligned}
u_a &= \mathsf{id}_{\bar S} \\
u_b &= \text{symétrie par rapport à l'axe } \bar 0, \bar 0' \ \textit{i.e.} % Doubt resolved: PDF leaf 18 (printed p.54) reads u_b, with plain b.
\\
&\phantom{=}\quad u_b \text{ laisse fixes } \bar 0, \bar 0' \text{ et intervertit } \bar a \text{ et } \bar b.
\end{aligned}
\]
(C'est plus clair si on désidentifie la base et la fibre, en prenant fibre
\[
\bar S \quad = \quad
\begin{tikzcd}
 & \bar a & \\
\bar 0 \arrow[ur] \arrow[dr] & & \bar 0' \arrow[ul] \arrow[dl] \\
 & \bar b &
\end{tikzcd}
\qquad ).
\]

Pour voir que $X$ est une $\D$-fibration sur $S$ (quelque soit $\D$) prouvons qu'il est \textit{propre}\footnote[45]{N.B. $f : X \to S$ est également \textit{lisse}, car c'est un fibré trivial (un produit) au dessus de deux ouverts $S - \{a\}$, $S - \{b\}$ qui recouvrent $S$. Mieux : $f$ est à la fois $\Cat$-fibrant et $\Cat$-cofibrant (car cette condition n'utilise que les propriétés des $X \times_S \boldsymbol{\Delta}_2 \to \boldsymbol{\Delta}_2$, pour $\boldsymbol{\Delta}_2 \to S$).} sur $S$ - on est ramenés à voir que $X' \to S'$ et $X'' \to S''$ sont propres, ce qui est trivial. Pour voir que c'est une $\D$-fibration, on est ramenés également au cas de $X' \to S'$ et $X'' \to S''$, où c'est aussi trivial. On a utilisé le : % [sic] quelque soit -> quel que soit, spelling

\textbf{Lemme.} \textit{Soient $f : X \to S$, $(S_i)_{i \in I}$ une famille de sous-catégories fermées recouvrant $S$. Pour tout $i$, soit $X_i = X \times_S S_i \xrightarrow{f_i} S_i$. Pour que $f$ soit propre, il faut et il suffit que les $f_i$ le soient. Pour que $f$ soit une $\D$-fibration propre il faut et il suffit que les $f_i$ le soient.}

\msp{page 89}
Voici le petit calcul qui montre que $f : X \to S$ est une $\D$-équivalence. On fait le calcul en cohomologie.\footnote[46]{Mais dans le cas du revêtement à 2 feuillets de $S$, il n'y a pas même de calcul à faire!} On doit prouver que
\[ H^*(S,k) \longrightarrow H^*(X,k) \tag*{$(*)$} \]
\printedpage{55}
% head: A. Grothendieck : Dérivateurs, Ch. I, Édition des Universités de Montpellier II et Paris VII
est un isomorphisme (et de même pour tout coefficient constant $M$ sur $S$). Or
\[ H^*(X,k) \overset{\mathrm{def}}{=} \RR\Gamma_X(k_X) = \RR\Gamma_S(\RR f_*(k_X)) \quad , \]
et $(*)$ est induit par
\[ i : k_S \longrightarrow \RR f_* k_X \quad . \]
Or on a
\[ \RR f_* k_X \simeq \underbrace{k_S}_{= \mathsf{Im}(i)} \quad \oplus \quad \underbrace{k_S[1]}_{k_S \text{ en degré } 1}{}^{\underline{\mathrm{tordu}}} \]
savoir ``tordu'' par l'action de monodromie sur $S = S^1$, laquelle est la symétrie $\lambda \mapsto -\lambda$. On a
\[ \RR\Gamma_S(k_S[1]^{\mathrm{tordu}}) = H^1(\underbrace{\pi_1(S)}_{= \mathbb{Z}}\, , k_S^{\mathrm{tordu}}) = k_{(\mathbb{Z})} = k/(u - \mathsf{id})k = k/2k = 0 \ , \]
d'où le résultat.

Notons que si on regarde
\[ X' = X \times_S X \xrightarrow{f'} S \quad , \]
on a par Künneth relatif
\[ \RR f'_*(k_{X'}) \simeq \RR f_*(k_X) \otimes^{\mathbb{L}} \RR f_*(k_X) \simeq k_S \oplus k_S[1]^t \oplus k_S[1]^t \oplus \underbrace{k_S[1]^t \otimes k_S[1]^t}_{\simeq k_S[2] \text{ non tordu}} \ , \]
d'où
\[ H^*_{\D}(X',k) \simeq \RR\Gamma_{X'}(k_{X'}) \simeq \RR\Gamma_S(\RR f'_*(k_{X'})) \simeq \RR\Gamma_S(k_S) \oplus \RR\Gamma_S(k_S)[2] \]
et
\[ H^*_{\D}(S,k) \longrightarrow H^*_{\D}(X',k) \]
n'est \textit{pas} un isomorphisme.

\msp{page 90}
Ainsi $f : X \to S$ qui est une $\D$ fibration (propre et lisse) satisfait bien les conditions équivalentes $b$), $b'$) de la proposition 2 (page 79), mais pas la condition $a$), puisque le changement de base
\[
\begin{tikzcd}[column sep=large, row sep=large]
X' = X \times_S X \arrow[r, "{q = pr_1}"] \arrow[d, "{f' = pr_2}"'] & X \arrow[d, "f"] \\
S' = X \arrow[r, "{p = f}"'] & S
\end{tikzcd}
\]
est une $\D$-équivalence $S' \to S$, alors que $X' \to X$ ne l'est pas.

Voici un autre contre-exemple, sur le même principe mais un peu plus simple. Soit $G$ un groupe fini (par exemple $G = \mathbb{Z}/2\mathbb{Z}$), faisons le opérer sur un objet $Z$ de $\Cat$, soit $S = [G] = B_G$ la catégorie réduite à un objet $a$, avec $G = \mathsf{End}(a)$. Soit $E_G$ la catégorie (isomorphe à $(B_G)/a$)
\[ \mathsf{Ob}\ E_G = G \quad , \qquad \mathsf{Hom}_{E_G}(g,g') = \{ h \in G \mid g = g'h \} \quad , \]
composition évidente - alors $E_G$ est un groupoïde trivial (1-connexe). Le goupe $G$ y opère à gauche librement par translations, et $B_G = E_G/G$. % [sic] goupe -> groupe, misprint

\printedpage{56}
% head: A. Grothendieck : Dérivateurs, Ch. I, Édition des Universités de Montpellier II et Paris VII

\msp{page 91}
Si $G$ opère à gauche sur $Z$, on peut alors l'utiliser pour tordre $Z$, par
\[ X = E_G \times_G Z = (E_G \times Z)/\text{opérations ``diagonales'' de G} \quad . \]
Donc on a un diagramme cartésien
\[
\begin{tikzcd}
 & X' \arrow[d, equal] & X'' \arrow[d, equal] \\
X \arrow[d, "f"'] & Z \times E_G \arrow[l] \arrow[d, "pr_2"] & Z \arrow[l, dashed, "{\overset{\D}{\sim}}"'] \arrow[d] \\
B_G \arrow[d, equal] & E_G \arrow[l] \arrow[d, equal] & e \arrow[l, dashed, "{\overset{\D}{\sim}}"'] \arrow[d, equal] \\
S & S' & S''
\end{tikzcd}
\]
Comme le changement de base $S' = E_G = (B_G)/a \to B_G = S$ transforme $X$ sur $B_G$ en un produit (lequel est un $\D$-fibré propre et lisse sur $S'$), on voit que $X$ sur $B_G$ est un $\D$-fibré propre et lisse (car c'est là une condition locale sur $B_G$).

On suppose maintenant que $G$ est fini d'ordre $n > 1$, et soit $k$ un anneau où $n$ est invesible, $\D = \mathsf{Der}(k - \mathsf{Mod})$. Alors $E_G \to B_G$ est une $\D$-équivalence, ou ce qui revient au même $e \to B_G$ l'est.\footnote[47]{En fait, $e$, $E_G$, $B_G$ sont $\D$-asphériques.} Pourtant, après changement de base soit par $S' = E_G$, soit par $S'' = e$, $X' \to X$ (ou ce qui revient au même, $X'' \to X$) n'a guère de raison d'en être aussi. On a en cohomologie % [sic] invesible -> inversible, misprint
\[ \RR\Gamma_X(k_X) \simeq \RR\Gamma_{BG}(f_*(k_X)) \simeq \RR\Gamma_G(G, H^*(Z,k)) \]
\msp{page 92}
donc ($n$ étant inversible dans $k$) :
\[ H^*(X,k) = \RR\Gamma_X(k_X) \simeq H^*(Z,k)^G \longrightarrow H^*(Z,k) \quad , \]
ce n'est pas en général un isomorphisme. On peut prendre par exemple $G = \mathbb{Z}/2\mathbb{Z}$, $Z \sim S^1$ comme avant, et $G$ opérant par symétrie de $S^1$ comme précédemment.

Cet exemple est particulièrement convainquant, car c'est un exemple \textit{avec $S$, $S'$ $\D$-asphériques,} et $S' \to S$ un isomorphisme local (en fait, un revêtement principal de $S$). C'est aussi un exemple d'un $\D$-fibré lisse et propre $X \to S$, sur une base \textit{asphérique} $S$, et tel que l'inclusion d'une fibre dans $X$ % [sic] convainquant -> convaincant, spelling of the adjective
\[ X_s \longrightarrow X \]
ne soit \textit{pas} une $\D$-équivalence.

Le seul reproche qu'on puisse faire à cet exemple, en comparaison avec le précédent, c'est que les catégories $X, S$ ne sont pas des ensembles ordonnés. Et $S$ ne peut être un ensemble ordonné fini, car l'espace classifiant d'un groupe fini a de la cohomologie en toutes dimensions (sauf erreur).

\printedpage{57}
% head: A. Grothendieck : Dérivateurs, Ch. I, Édition des Universités de Montpellier II et Paris VII

\msp{page 93}
Mais on peut s'inspirer de cet exemple pour le modifier de façon qu'il se transpose à des ensembles ordonnés. Il suffit de réaliser, pour $G$ donné :

\textit{a}) Une opération \textit{libre} de $G$ sur un objet $E$ de $\Cat$, telle que $H^*(E,k)^G \simeq k$, \textit{i.e.} $B = E/G$ soit $\D$-asphérique.

\textit{b}) Une opération de $G$ sur un objet $Z$, telle que l'inclusion
\[ H^*(Z,k)^G \longrightarrow H^*(Z,k)^G \] % [sic] H^*(Z,k)^G \longrightarrow H^*(Z,k)^G -> H^*(Z,k)^G \longrightarrow H^*(Z,k), an inclusion of the invariants (cf. p. 56)
ne soit pas \texttt{[un]} isomorphisme.

\textit{c}) Choisir pour $E$, $Z$ des ensembles ordonnés, et s'assurer que $B = E/G$ est également représenté par une catégorie ordonnée. Sauf erreur, il en sera de même pour $X = E \times Z/G$. On aura en tout les cas % [sic] en tout les cas -> en tous les cas, agreement
\[
\begin{tikzcd}[column sep=large]
 & Z \times E \arrow[d, equal] & \\
X \arrow[d] & X' \arrow[l] \arrow[d] & Z \arrow[l] \arrow[d] \\
S \arrow[d, equal] & S' \arrow[l] \arrow[d, equal] & S'' = \{s\}\, , \quad s \in E \ . \arrow[l] \\
E/G & E & \text{\footnotesize asph} \\
\text{\footnotesize asph} & \substack{\text{asph} \\ \text{si on peut}} &
\end{tikzcd}
\]
Cela donne encore un exemple avec $X$ $\D$-fibré propre et lisse, sur une base $S$ $\D$-asphérique, tel que l'inclusion d'une fibre $X_s$ dans $X$ ne soit pas une $\D$-équivalence. On peut même se débrouiller pour que $H^*(Z,k)^G \simeq k$ ce qui implique
\msp{page 94}
que $X$ est lui aussi asphérique (mais les fibres de $X \to S$ ne sont pas $\D$-asphériques). Cela donne donc un autre contre-exemple au b) de la prop. 6 (p. 84).

Voici la façon la plus économique, je crois, de réaliser la situation que j'ai dite. Topologiquement, je prend % [sic] je prend -> je prends, conjugation
\[
\begin{aligned}
G &= \{\pm 1\} \\
E &= S^2 &\qquad& \text{sphère euclidienne, sur laquelle } \{\pm 1\} \\
  &      &&      \text{opère par antipodisme } x \mapsto -x \\
B &= S^2/\{\pm 1\} && \text{plan projectif réel} \\
Z &= S^1 && \text{où G opère comme précédemment par la symétrie\,.}
\end{aligned}
\]
C'est donc un exemple très proche de celui du début (p. 87), où on pouvait aussi considérer la base ($B = S^1$) comme quotient d'un $E$ ($= S^1$) par $G = \{\pm 1\}$ opérant librement (par
\printedpage{58}
% head: A. Grothendieck : Dérivateurs, Ch. I, Édition des Universités de Montpellier II et Paris VII
l'antipodisme $z \mapsto -z$ dans $\mathbb{C} \simeq \RR^2$). Mais dans cet exemple, la base n'était pas $\D$-asphérique.

La façon la plus économique de réaliser la sphère $S^2$ par un ensemble ordonné est sans doute par sa décomposition ``globulaire'' par trois paires de cellules de dimension 0, 1, 2, donnant l'ensemble ordonné
\[
I \qquad = \qquad
\begin{tikzcd}[row sep=large, column sep=large]
S_2^- & S_2^+ \\
S_1^- \arrow[u] \arrow[ur] & S_1^+ \arrow[u] \arrow[ul] \\
S_0^- \arrow[u] \arrow[ur] & S_0^+ \arrow[u] \arrow[ul]
\end{tikzcd}
\]
et l'antipodisme échange $S_i^-$ et $S_i^+$, mais le quotient n'est pas ordonné, c'est la catégorie (importante pour la théorie des 2-catégories) :
\msp{page 95}
\[
\begin{tikzcd}[row sep=large]
S_2 \\
S_1 \arrow[u, shift left, "p_1^-"] \arrow[u, shift right, "p_1^+"'] \\
S_0 \arrow[u, shift left, "p_0^-"] \arrow[u, shift right, "p_0^+"']
\end{tikzcd}
\qquad\qquad \text{avec} \qquad\qquad
\begin{cases}
p_1^- p_0^- = p_1^- p_0^+ \\
p_1^+ p_0^- = p_1^+ p_0^+
\end{cases}
\]
Mais prenons la subdivision barycentrique de la décomposition cellulaire précédente de $S$ dont les sommets correspondent aux éléments de $\mathsf{Ob}\ I$, les simplexes aux parties totalement ordonnés de $\mathsf{Ob}\ I$. On trouve l'octaèdre, ayant 6 sommets, 12 arêtes, 8 faces, et $G$ opère librement sur l'ensemble des simplexes, donc le quotient : 3 sommets, 6 arêtes, 4 faces, décomposition cellulaire du plan projectif. (Sa décomposition par trois droites.) On présente $S^1$ par l'ensemble ordonné (comme au début) % [sic] parties totalement ordonnés -> parties totalement ordonnées, agreement
\[
\begin{tikzcd}[row sep=large, column sep=large]
S_1^- & S_1^+ \\
S_0^- \arrow[u] \arrow[ur] & S_0^+ \arrow[u] \arrow[ul]
\end{tikzcd}
\qquad\qquad \textit{i.e.} \qquad\qquad
\begin{tikzcd}
 & S_1^+ & \\
0 = S_0^- \arrow[ur] \arrow[dr] & & S_0^+ = 0' \arrow[ul] \arrow[dl] \\
 & S_1^- &
\end{tikzcd}
\]
et $G$ y opère par symétrie échangeant $S_1^-$ et $S_1^+$. Donc l'ensemble ordonné $X$ aura $4 \times 13 = 52$ sommets et c'est modeste \ldots


\endgroup
% END UNIT Der-ChI-234/Der-ChI-234_p056-060.tex
% BEGIN UNIT Der-ChI-234/Der-ChI-234_p061-065.tex
\begingroup\makeatletter%
\NoAutoSpacing%
\newcommand{\printedpage}[1]{\marginpar{\footnotesize p.~#1}}%
\newcommand{\msp}[1]{\marginpar{\footnotesize [#1]}}%
\newcommand{\Hot}{\textsf{Hot}}%
\newcommand{\Dia}{\textsf{Dia}}%
\newcommand{\Ob}{\textsf{Ob}}%
\newcommand{\Hom}{\textsf{Hom}}%
\newcommand{\DD}{\mathbb{D}}%
\makeatother\let\printedpage\derpp\let\msp\dermsp%
\let\mspin\DERmspin%


\printedpage{59}
% head: A. Grothendieck : Dérivateurs, Ch. I, Édition des Universités de Montpellier II et Paris VII
% foot: Version 8 juillet 2001

\msp{page 96}
Finalement, il semble qu’il ne reste plus rien de la proposition 2 (p.79) introduisant la notion de $\DD$-fibration – pas même l’équivalence des deux conditions $b$) et $b'$) ($\DD$-fibration à gauche universelle, $\DD$-fibration à droite universelle). On songe à leur relation à une condition $a$) modifiée ainsi (pour $f : X \to Y$) :

$a$) Pour tout $Y'' \to Y' \to Y$, avec $Y'$ et $Y''$ \Hot-\underline{asphériques} (et pas seulement $\DD$-asphériques) $X'' = X \times_Y Y'' \to X' = X \times_Y Y'$ est une $\DD$-équivalence. (Quand on suppose que $Y''$ et $Y'$ ont des éléments finaux, c’est la condition que $f$ soit une $\DD$-fibration à gauche universelle; quand on suppose que $Y''$ et $Y'$ ont des éléments initiaux, c’est la condition $\DD$-fibration à droite universelle.)\footnote[48]{Il faut expliciter le raisonnement. Il suffit de tester même sur $Y'' \to Y'$, où $Y''$ est \underline{ponctuelle}, et où $Y'$ a un objet final (resp. initial).}

On a évidemment $a$) $\Rightarrow$ $b$) et $a$) $\Rightarrow$ $b'$) et on aimerait que ce soient des équivalences. Voici un axiome intéressant\footnote[49]{\textit{cf.} \textbf{Der 3} $a'$) page suivante.} sur $\DD$, qui semble vérifié dans tous les cas connus, qui assure l’équivalence souhaitée. Je vais en même temps reformuler \textbf{Der 3} en trois volets (\textit{cf.} p. 67).

\msp{page 97}
\textbf{Der} 3 $a$) Si $Z \in \Dia$ est \Hot-asphérique, alors $p : Z \to e$ est une $\DD$-équivalence, \textit{i.e.} par exemple $p^* : \mathcal{A} = \DD(e) \to \DD(Z)$ est pleinement fidèle.

$b$) Soit $f : Z \to Z'$ une flèche dans \Dia. Si $f$ est une \Hot-équivalence, c’est une $\DD$-équivalence (\textit{i.e.} $f^* : \DD(Z') \to \DD(Z)$ induit une \underline{équivalence} entre les sous-catégories de $Z$, $Z'$ \texttt{[plutôt de} $\DD(Z)$\texttt{,} $\DD(Z')$\texttt{]} formées des objets constants).

$a'$) (Précise considérablement $a$). Si $Z \in \Dia$ est \Hot-asphérique, alors on a $a$), et l’image essentielle de $\DD(e) = \mathcal{A}$ dans $\DD(Z)$ est formée des objets \underline{localement} \underline{constants} de $Z$, \textit{i.e.} des objets tels que le foncteur
\[
Z^{\circ} \longrightarrow \mathcal{A} \qquad z \longmapsto F_z
\]
soit localement constant (\textit{i.e.} transforme toute flèche en isomorphisme) donc isomorphe à un foncteur constant (puisque $Z$ est \Hot-asphérique, donc 1-connexe).

Il faut bien noter que dans $a'$), il ne suffit \underline{pas} de supposer que $Z$ soit $\DD$-asphérique, au lieu de \Hot-asphérique. Dès que la $\DD$-équivalence est strictement plus faible que \Hot, il semble que $a'$) devient toujours faux, en tant que condition sur $\DD$, si on y suppose seulement que $Z$ soit $\DD$-asphérique. Par exemple, même si $\DD = \DD_i$ (théorie des $i$-types d’homotopie relative), la catégorie
\msp{page 98}
des objets localement constants sur $Z$ s’exprime en termes du $(i+1)$-type d’homotopie, mais pas en termes de son $i$-type d’homotopie seul.

Il semble bien que ce soit l’axiome \textbf{Der} 3 ($a'$) qui assure l’existence d’une bonne notion (“bilatère”) de $\DD$-fibration, qui soit stable par changement de base, et qui se réduise soit
\printedpage{60}
% head: A. Grothendieck : Dérivateurs, Ch. I, Édition des Universités de Montpellier II et Paris VII
% foot: Version 8 juillet 2001
à “$\DD$-fibration à gauche universelle” soit à “$\DD$-fibration à droite universelle” et donc qui soit en plus une condition \underline{vérifiable} pratiquement. Dans le cas où $f$ est propre ou lisse, le critère est particulièrement simple, puisque il se borne à la condition que les $f_* f^* a_Y$, ou les $f_! f^* a_Y$ (pour $a \in \mathcal{A}$), soient des objets localements constants sur $Y$ (donc \underline{constants} si $Y$ est \underline{\underline{\Hot}}-asphérique, en vertu de \textbf{Der} 3 ($a'$). % [sic] localements -> localement, plural ending on the adverb; % [sic] "(donc ... ($a'$)." -> closing parenthesis after "($a'$)" missing, unbalanced parenthesis

Je vais également introduire une variante renforcée de \textbf{Der} 3 $b$) :

\textbf{Der} 3 $b'$) Soit $Z' \to Z$ une \Hot-équivalence dans \Dia. Alors $f^* : \DD(Z) \to \DD(Z')$ induit une équivalence entre les catégories des coefficients \underline{localements} \underline{constants} sur $Z$ et sur $Z'$.\footnote[50]{\textbf{Der} 3 $b'_0$) : le foncteur est fidèle\\
\hspace*{6em}$b'_1$) : le foncteur est pleinement fidèle\\
\hspace*{6em}$b'_2$) : le foncteur est une équivalence} % [sic] localements -> localement, plural ending on the adverb

\msp{page 99}
Dans certaines questions il suffit de savoir que le dit foncteur est pleinement fidèle. Cela implique que
\[
H^*_{\DD}(Z,F) \overset{\sim}{\longrightarrow} H^*_{\DD}(Z', f^*F)
\]
quand $F$ est \underline{localement} \underline{constant} sur $Z$\footnote[51]{On pourra dire que $f : Z' \to Z$ est une $\DD$-équivalence 0-\textit{complète à gauche,} et 0-\textit{complète à droite} dans le cas homologique.} (mais il n’est pas clair que cette condition implique toujours la pleine fidélité\footnote[52]{Mais c’est OK si on a des \underline{\Hom} internes dans $\DD(Z)_{\mathrm{loc.const.}}$ satisfaisant
\[
\Hom(F,G) \simeq \Hom(e_Z, \underline{\Hom}(F,G)) \quad ,
\]
où $e$ est l’élément final de $\mathcal{A}$, les \underline{\Hom} étant compatibles aux images inverses.}) et de même pour l’homologie
\[
H_*^{\DD}(Z', f^*F) \longrightarrow H_*^{\DD}(Z,F) \quad .
\]
Moyennant cette condition, on a la :

\textbf{Proposition :} \textit{On suppose} \textbf{Der} 3 \textup{b}$'$) \textit{satisfait. Soit $f : X \to Y$ dans} \Dia. \textit{Pour que $f$ soit une $\DD$-fibration, il faut et il suffit que pour $Y'' \to Y'$ sur $Y$ qui soit une} \Hot\textit{-équivalence, $X'' = X \times_Y Y'' \to X' = X \times_Y Y'$ soit une $\DD$-équivalence.}

À vrai dire, on peut énoncer quelque chose de raisonnable, sans supposer l’axiome \textbf{Der} 3 $b'$). On doit introduire la :

\textbf{Définition :} Soit $f : Z' \to Z$ dans \Dia. On dit que $f$ et une $\DD$-équivalence 1-\textit{complète} si le foncteur $f^* : \DD(Z) \to \DD(Z')$ est pleinement fidèle sur la sous-catégorie pleine de $\DD(Z)$ formée des coefficients localement constants. On dit que $f$ est une $\DD$-équivalence 2-\textit{complète} si elle induit une équivalence entre les catégories des coefficients localement constants. % [sic] f et une -> f est une, verb "est" misprinted as "et"

\msp{page 100}
On dit que $Z$ est \textit{complètement} $\DD$-\textit{asphérique}, si $Z \to e$ est une $\DD$-équivalence 2-complète.

\printedpage{61}
% head: A. Grothendieck : Dérivateurs, Ch. I, Édition des Universités de Montpellier II et Paris VII
% foot: Version 8 juillet 2001

Ainsi, l’axiome \textbf{Der} 3 $a'$) peut se reformuler, en disant que si $Z$ est \Hot-asphérique, alors $Z$ est complètement $\DD$-asphérique. Et de même, \textbf{Der} 3 $b'$) en disant que si $f : X \to Y$ dans \Dia{} est une \Hot-équivalence, alors $f$ est une $\DD$-équivalence 2-complète. Nous allons donner une variante affaiblie, de \textbf{Der} 3 $a'$), qui suffit pour donner une notion raisonnable de $\DD$-fibration :

\textbf{Der} 3 $a''$) Tout $Z$ dans \Dia{} qui a soit un objet final, soit un objet initial, est complètement $\DD$-asphérique.

\textbf{Définition.} On dit que $f : X \to Y$ dans \Dia{} est une $\DD$-fibration, si pour
\[
Y'' \longrightarrow Y' \longrightarrow Y
\]
avec $Y'' \to Y'$ une $\DD$-équivalence 1-complète (par exemple $Y'$ et $Y''$ complètement $\DD$-asphèriques), $X'' = X \times_Y Y'' \longrightarrow X' = X \times_Y Y'$ est une $\DD$-équivalence. % [sic] asphèriques -> asphériques, grave accent instead of acute

Cela équivaut à chacune des conditions $b$), $b'$) de la prop. 2 (page 79), \textit{i.e.} que $f$ soit universellement une $\DD$ fibration à gauche, ou universellement une $\DD$-fibration à droite (grâce à \textbf{Der 3} $b''$ \texttt{[plutôt} $a''$\texttt{]}). % [sic] $\DD$ fibration -> $\DD$-fibration, hyphen missing

\msp{page 101}
\textbf{Définition 1.} $f : X \to Y$ dans \Dia{} est appelé $\DD$-\textit{équivalence cohomologique} (sous-entendu : \textit{complète}) si pour tout $F \in \Ob\,\DD(Y)$, $F$ \underline{localement} \underline{constant}, on a
\[
H^*_{\DD}(Y,F) \overset{\sim}{\longrightarrow} H^*_{\DD}(X, f^*(F)) \quad .
\]
Notion duale $\DD$-\textit{équivalence homologique}
\[
H_*^{\DD}(X, f^*(F)) \overset{\sim}{\longrightarrow} H_*^{\DD}(Y,F) \quad .
\]
Si $Y$ est complètement $\DD$-asphérique, alors les coefficients localement constants sur $Y$ sont constants, et $f$ est une $\DD$-équivalence cohomologique (resp. homologique) si et seulement si c’est une $\DD$-équivalence \textit{i.e.} si $X$ est $\DD$-asphérique.\footnote[53]{L’axiome \textbf{Der} 3 $b'_1$) implique que toute \Hot-équivalence est une $\DD$-équivalence cohomologique.}

\textbf{Proposition 1} (Soritale). \textit{Soient $X \overset{f}{\to} Y \overset{g}{\to} Z$. Si $f$ et $g$ sont des $\DD$-équivalences cohomologiques, alors $gf$ aussi. Si $gf$ et $f$ sont des $\DD$-équivalences cohomologiques, alors $g$ aussi. Si $gf$ est une $\DD$-équivalence cohomologique, et si $g$ est une $\DD$-équivalence \underline{2-complète}} (donc tout $F$ localement constant sur $Y$ provient d’un $G$ localement constant sur $Z$), \textit{alors $f$ est une $\DD$-équivalence cohomologique.}

\textbf{Définition 2.} 1°) $f : X \to Y$ dans \Dia{} est appelé une $\DD$-\textit{fibration cohomologique complète,} je dis bien \underline{complète}, si

$a$) Pour tout $y \in \Ob\, Y$, $X_y \to X/y$ est une $\DD$-équivalence cohomologique (complète).

$b$) Pour tout $y' \to y$ dans $Y$, $X/y' \to X/y$ est une $\DD$-équivalence cohomologique (complète).

\printedpage{62}
% head: A. Grothendieck : Dérivateurs, Ch. I, Édition des Universités de Montpellier II et Paris VII
% foot: Version 8 juillet 2001

\msp{page 102}
$c$) Ces conditions restent satisfaites après tout changement de base $Y' \to Y$.

2°) Notion duale de fibration homologique complète.

\textsc{Exemple.} Si $f$ est \Hot-propre, les flèches des $a$) et $b$) sont des \Hot-équivalences, donc en vertu de \textbf{Der 3} $b'_1$) (qu’il faut admettre), ce sont des $\DD$-équivalences cohomologiques. Cela signifie que $a$) et $b$) sont valables, et $c$) aussi, car la notion de flèche \Hot-propre est stable par changement de base.

\textbf{Proposition 2.} \textit{Soient $f : X \to Y$ $\DD$-fibration cohomologique complète, $F \in \Ob\,\DD(X)$ localement constant sur $X$. Alors $f_*(F)$ est localement constant, sa formation commute à tout changement de base, et en particulier se calcule fibre par fibre.}\footnote[54]{Il y a une réciproque \textit{cf.} prop. 3$'$.}

\textbf{Corollaire 1.} \textit{Soit un diagramme cartésien}
\[
\begin{tikzcd}
X & X' \arrow[l, "q"'] \arrow[d, "f'"] \\
Y & Y' \arrow[l, "p"']
\arrow[from=1-1, to=2-1, "f"']
\end{tikzcd}
\quad .
\]
\textit{On suppose $f$ $\DD$-fibration cohomologique et $p$ une $\DD$-équivalence cohomologique. Alors $q$ est une $\DD$-équivalence cohomologique.}

\msp{page 103}
En effet, on a pour $F$ localement constant sur $X$
\[
H^*_{\DD}(X,F) \simeq H^*_{\DD}(Y, f_*F) \overset{\sim}{\longrightarrow} H^*_{\DD}(Y', p^*f_*F) \simeq H^*_{\DD}(Y', f'_*q^*F) \simeq H^*_{\DD}(X', q^*F)
\]
qed.

\textbf{Proposition 3.} \textit{La notion de $\DD$-fibration cohomologique complète est stable par changement de base} (trivial) \textit{et par composition.}

Pour la stabilité par composition, il suffit de noter le corollaire suivant de la proposition 2.

\textbf{Corollaire 2} (de prop. 2). \textit{Soit $f : X \to Y$. Pour que $f$ soit une $\DD$-fibration cohomologique complète, il faut et il suffit que pour tout $Y'' \overset{p}{\to} Y' \to Y$ sur $Y$, avec $p$ une $\DD$-équivalence cohomologique, la flèche correspondante $X'' = X \times_Y Y'' \longrightarrow X' = X \times_Y Y'$ soit aussi une $\DD$-équivalence cohomologique.}

C’est nécessaire par le corollaire 1, compte tenu que $f' : X' = X \times_Y Y' \to Y'$ est aussi une $\DD$-fibration cohomologique. C’est suffisant, car il suffit de l’appliquer aux cas $\{y\} \to Y/\{y\}$, $Y/y \to Y/y'$, qui sont des $\DD$-équivalences cohomologiques en
\msp{page 104}
vertu de \textbf{Der} 3 $b'_1$) (ou seulement \textbf{Der} 3 $a'$)).

\printedpage{63}
% head: A. Grothendieck : Dérivateurs, Ch. I, Édition des Universités de Montpellier II et Paris VII
% foot: Version 8 juillet 2001

\textbf{Proposition 4.} \textit{On suppose que} \textbf{Der} 3 a$'$) \textit{est satisfait} (\textit{i.e.} tout $Z$ dans \Dia{} qui est \Hot-asphérique, est complètement $\DD$-asphérique). \textit{Soit $f : X \to Y$ dans} \Dia. \textit{Les conditions suivantes sont équivalentes}

1°) \textit{$f$ est une $\DD$-fibration cohomologique \underline{complète}} (conditions $a$), $b$), $c$) de \texttt{[la]} définition 2, page 101).

2°) \textit{Pour tout $F \in$} \Ob\,$\DD(X)$\textit{, localement constant, $f_*(F)$ est localement constant, se calcule fibre par fibre; et ça reste vrai après tout changement de base $Y' \to Y$ (avec $F'$ localement constant sur $X' = X \times_Y Y'$). (Alors la formation de $f_*(F)$ commute à \underline{tout} changement de base).}

3°) \textit{Pour tout $Y'' \overset{p}{\to} Y'$ sur $Y$, si $p$ est une $\DD$-équivalence cohomologique, alors $X'' = X \times_Y Y'' \longrightarrow X' = X \times_Y Y'$ aussi.}

3$'$) \textit{Comme 3, en supposant que $Y''$ est la catégorie ponctuelle, et que $Y'$ a un élément final : il faut que $X'' \to X'$ soit une $\DD$-équivalence cohomologique.}

\msp{page 105}
\textbf{Définition 3.} On dit que $f : X \to Y$ dans \Dia{} est une $\DD$-\textit{fibration complète} si c’est une $\DD$-fibration cohomologique complète et une $\DD$-fibration homologique complète.

C’est stable par changement de base, par composition.

La définition 3 équivaut à ceci : pour tout $p : Y'' \to Y'$ sur $Y$, avec $Y''$ et $Y'$ \Hot-asphériques, le morphisme correspondant $X'' \to X'$ est \underline{à} \underline{la} \underline{fois} une $\DD$-équivalence cohomologique, et une $\DD$-équivalence homologique. (Et il suffit même de le supposer pour $Y'' = e$ et pour $Y'$ avec objet final ou objet initial.) On peut songer à renforcer encore cette condition en exigeant (sans regarder à gauche ni à droite) que $X'' \to X'$ soit une $\DD$-équivalence 1-complète. Mais il n’est pas clair que cette notion renforcée soit stable par composition – à moins de la renforcer encore, en exigeant que si $Y'' \to Y'$ est une $\DD$-équivalence 1-complète, alors $X'' \to X'$ l’est aussi. L’ennui avec cette définition renforcée, c’est que faute d’avoir une description cohomologique ou homologique de la notion de
\msp{page 106}
$\DD$-équivalence 1-complète, je ne me sens pas capable de donner des exemples intéressants (non totalement triviaux) de $\DD$-fibrations en ce sens très fort. Sauf les \Hot-fibrations triviales (\textit{i.e.} à fibres \Hot-asphériques) – mais ce n’est pas drôle!

Mais quand on a des produits $*$ associatifs et unitaires et des \underline{\Hom} internes correspondants dans les $\DD(X)_{\mathrm{loc.\ const.}}$, compatibles avec \texttt{[les]} images inverses, donc on aura sur $X$ si $F$ et $G$ sont localement constants dans $\DD(X)$, et si $\varepsilon_X$ désigne l’unité pour $*$ dans $\DD(X)$,
\[
\Hom(F,G) \simeq \Hom(\varepsilon_X * F, G) \simeq \Hom(\varepsilon_X, \underline{\Hom}(F,G)) \simeq \Hom_{\mathcal{A}}(\varepsilon_e, p_*\, \underline{\Hom}(F,G)) \ ,
\]
car $\varepsilon_X = p^*(\varepsilon_e)$, où $p : X \to e$, donc dans ce cas une flèche $f : X \to Y$ dans \Dia{} est une $\DD$-équivalence 1-complète si (et seulement si) c’est une $\DD$-équivalence cohomologique. Donc dans ce cas, si $f$ est une $\DD$-fibration cohomologique complète, c’est aussi une $\DD$-fibration homologique complète, donc aussi une $\DD$-fibration complète (définition 3). (Et même au sens soit-disant renforcé de tantôt : si $Y'' \to Y'$ sur $Y$ est une $\DD$-équivalence 1-complète ($\Leftrightarrow$ une $\DD$-équivalence cohomologique!) alors $X'' \to X'$ aussi. % [sic] soit-disant -> soi-disant, spelling; % [sic] "(Et même ... aussi." -> closing parenthesis missing, unbalanced parenthesis


\endgroup
% END UNIT Der-ChI-234/Der-ChI-234_p061-065.tex
% BEGIN UNIT Der-ChI-234/Der-ChI-234_p066-070.tex
\begingroup\makeatletter%
\newcommand{\printedpage}[1]{\marginpar{\footnotesize p.~#1}}%
\newcommand{\msp}[1]{\marginpar{\footnotesize [#1]}}%
\newcommand{\HOT}{\textsf{HOT}}%
\newcommand{\Hot}{\textsf{Hot}}%
\newcommand{\Dia}{\textsf{Dia}}%
\newcommand{\Cat}{\mathcal{C}at}%
\newcommand{\DD}{\mathbb{D}}%
\makeatother\let\printedpage\derpp\let\msp\dermsp%
\let\mspin\DERmspin%


\printedpage{64}
% head: A. Grothendieck : Dérivateurs, Ch. I, Édition des Universités de Montpellier II et Paris VII
% foot: Version 8 juillet 2001

\msp{page 107}

Il se pose la question de trouver des $f$ qui soient des $\DD$-fibrations cohomologiques, ou homologiques. Je ne sais si je serai capable d'en trouver, qui ne soient déjà des \HOT-fibrations cohomologiques, ou homologiques (cela revient maintenant au même\footnote[55]{Les \HOT-équivalences cohomologiques $=$ \HOT-équivalences homologiques $=$ \HOT-équivalences sans plus.} ), car pratiquement je ne sais guère dire qu'une flèche $Z' \to Z$ dans $\Cat$ est une $\DD$-équivalence cohomologique complète, que si cette flèche est une \HOT-équivalence (et en admettant l'axiome \textbf{Der} 3 $b'_1$), ou sa variante affaiblie cohomologique).

Je vais donc chercher les \HOT-fibrations, \textit{i.e.} les $f : X \to Y$ tels que pour tout $Y'' \to Y' \to Y$, avec $Y''$ et $Y'$ asphériques (sous-entendu : pour \HOT), $X'' \to X'$ est une \Hot-équivalence. (Il suffit même de supposer $Y''$ réduit à un point, $Y'$ avec objet final $\ldots$ )

Supposons par exemple $f$ propre (\textit{i.e.} \HOT-propre), donc la condition $a$) de la définition 2, p. 101 satisfaite et de même après tout changement de base. Suffit-il de vérifier de plus la condition $b$), \textit{i.e.} que
\[ X/y \longrightarrow X/y' \quad , \qquad \text{pour } y \to y' \text{ dans } Y \ , \]
sont des \Hot-équivalences? (J'entends en restant sur $Y$, sans supposer aussi d'avoir la même chose après changement de base).
\msp{page 108}
Mais cela reste-il vrai après tout changement de base? % [sic] reste-il -> reste-t-il, forme interrogative

La situation semble ici meilleur pour $f$ lisse (au lieu de propre), en interprétant les conditions $a$) et $b$) sous forme homologique, comme signifiant que $f_!(e_X)$ est un objet \underline{localement} \underline{constant} de $\mathsf{HOT}(Y)$, qui se calcule fibre par fibre. Comme $f$ lisse implique que $f_!$ commute à tout changement de base, cette hypothèse sur $f$ (que $f$ soit une \HOT-fibration cohomologique, peut-être pas complète) est stable par changement de base, donc implique que c'est une \HOT-fibration cohomologique complète. Mais cet argument utilise les propriétés générales des dérivateurs pour \HOT, que je n'ai pas encore établies - surtout pas dans le domaine $\Cat$ tout entier! Donc il est sain de me méfier. Je préfère tâter la difficulté en restant d'abord dans le cadre du formalisme \underline{cohomologique} (non commutatif), dont j'ai une certaine maîtrise. Je suppose donc $f$ propre à nouveau, et que les % [sic] meilleur -> meilleure, accord avec « situation »
\[ X_y \longrightarrow X/y \quad , \qquad X/y \longrightarrow X/y' \]
\msp{page 109}
soient des \HOT-équivalences, je veux prouver la même chose après tout changement de base, au moins que dans la situation
\[
\begin{tikzcd}
X \arrow[r, leftarrow, "q"] \arrow[d, "f"'] & X' \arrow[r, leftarrow, "q'"] \arrow[d, "f'"'] & X'' \arrow[d, "f''"] \\
Y \arrow[r, leftarrow, "p"'] & Y' \arrow[r, leftarrow, "p'"'] & Y''
\end{tikzcd}
\]
\printedpage{65}
% head: A. Grothendieck : Dérivateurs, Ch. I, Édition des Universités de Montpellier II et Paris VII
% foot: Version 8 juillet 2001
si $Y'$, $Y''$ sont asphériques, alors $q'$ est une \Hot-équivalence. Si je veux l'établir par critère cohomologique, je dois commencer par prendre sur $X'$ un faisceau d'\textit{ensembles} localement constant $F$, et prouver que $\Gamma(X',F) \overset{\sim}{\to} \Gamma(X'',{q'}^{*}F)$. Comme ce faisceau ne provient pas de $X$, je suis ennuyé pour appliquer l'hypothèse sur $f$, qui concerne à priori des faisceaux d'ensembles constants sur $X$, ou sur les $X/y$. Mais on devrait avoir ceci :

\textbf{Lemme.} \textit{Soit $f : X'' \to X'$. Alors $f$ est une} \HOT-\textit{équivalence} i.e.
\[ H^{*}_{\mathsf{HOT}}(X', a_{X'}) \xrightarrow{\ \sim\ } H^{*}_{\mathsf{HOT}}(X'', a_{X''}) \]
\textit{est} \texttt{[un]} \textit{isomorphisme pour tout $a \in \mathsf{Hot} = \mathsf{HOT}(e)$, si et seulement si $f$ est une} \Hot-\textit{équivalence,}\footnote[56]{Attention ici \texttt{[mot illisible]} la distinction entre \HOT-équivalence et \Hot-équivalence.} \textit{et alors $f$ est même une} \HOT-\textit{équivalence} $2$-\textit{complète} (j'ai besoin seulement de : \HOT-équivalence cohomologique et homologique complète) i.e. \textit{pour tout $\xi \in \DD(X')$ \underline{localement} \underline{constant} sur $X'$, on a}\footnote[57]{Est-ce vrai aussi pour les dérivateurs $\DD_i$?}
\[ H^{*}_{\mathsf{HOT}}(X',\xi) \xrightarrow{\ \sim\ } H^{*}_{\mathsf{HOT}}(X'',{q'}^{*}(\xi)) \quad . \]
\msp{page 110}

La première assertion doit résulter des formules
\[ H^{*}_{\mathsf{HOT}}(Z,a_Z) \simeq \underline{\mathsf{Hom}}_{\mathsf{Hot}}(\mathsf{hot}(Z),a) \tag*{$(*)$} \]
\[ \text{(et aussi} \qquad\qquad H_{*}^{\mathsf{HOT}}(Z,a_Z) \simeq \mathsf{hot}(Z) \times a \qquad\qquad ) \]
donc si
\[ \mathsf{hot}(X'') \xrightarrow{\ \sim\ } \mathsf{hot}(X') \tag*{$(**)$} \]
est \texttt{[un]} isomorphisme, alors
\[ H^{*}_{\mathsf{HOT}}(X', a_{X'}) \longrightarrow H^{*}_{\mathsf{HOT}}(X'', a_{X''}) \]
est \texttt{[un]} isomorphisme, puisque c'est déduit de $(*)$ en prenant $\underline{\mathsf{Hom}}_{\mathsf{Hot}}(\,.\,,a)$. Inversement, si $(*)$ est toujours \texttt{[un]} isomorphisme, exprimant les deux membres \texttt{[mot illisible]} de $\underline{\mathsf{Hom}}_{\mathsf{Hot}}(\,.\,,a)$ et prenant les $\pi_0$, on trouve que
\[ \mathsf{Hom}_{\mathsf{Hot}}(\mathsf{hot}(X'),a) \xrightarrow{\ \sim\ } \mathsf{Hom}_{\mathsf{Hot}}(\mathsf{hot}(X''),a) \]
est un isomorphisme pour tout $a$, donc $(**)$ est un isomorphisme.

Mais le fait que ça implique que $f''$ soit une \HOT-équivalence cohomologique complète, \textit{i.e.} que la relation $(*)$ reste vraie pour un $\xi \in \DD(X')$ \underline{localement} \underline{constant} sur $X'$, me paraît une propriété délicate du dérivateur \HOT, sur laquelle il me faudra revenir.

Donc le lemme est loin d'être établi, faute d'avoir seulement construit le dérivateur \HOT, sans même parler d'en établir les propriétés spéciales délicates! Mais je n'ai aucune raison de douter du lemme, il fait partie de la kyrielle des intuitions de base que je voudrais fonder solidement.

\printedpage{66}
% head: A. Grothendieck : Dérivateurs, Ch. I, Édition des Universités de Montpellier II et Paris VII
% foot: Version 8 juillet 2001

\msp{page 111}

Le lemme implique qu'une $\DD = \mathsf{HOT}$-fibration (au sens de la définition 2 page 81, revue et corrigée à la page 96-97) est en fait déjà une \HOT-fibration cohomologique et homologique complète. Or on a vu que si $f$ est \underline{propre}, $f$ est une $\DD$-fibration si et seulement si les conditions précédentes \underline{sur $Y$} (sans changement de base) sont satisfaites. Et dualement pour $f$ lisse.

\begin{flushright}
11.11.90
\end{flushright}

Je sens le besoin de revoir les notions de $\DD$-fibrations en tous genres (p. 76 ff). D'abord les diverses variantes utiles de l'axiome \textbf{Der} 3 sur les dérivateurs :

\textbf{Der} 3 $a$)\footnote[58]{\textit{cf.} aussi la série \textbf{Der 3} $c \ldots$ , page 124 ff.} Si $Z \in \mathsf{Dia}$ est \Hot-asphérique, il est $\DD$-asphérique \textit{i.e.} $\mathcal{A} = \DD(e) \to \DD(Z)$ est pleinement fidèle. % Doubt resolved: attached PDF leaf 30 (printed p.66), footnote 58 has bold Der 3.

$a'$) Comme $a$), avec en plus la caractérisation de l'image essentielle de $\mathcal{A}$ dans $\DD(Z)$, comme formée des $\DD$-coefficients localements constants. % [sic] localements -> localement, adverbe invariable

\textbf{Der} 3 $b$) Si $f : Z' \to Z$ dans \Dia{} est une \Hot-équivalence, alors c'est une $\DD$-équivalence, \textit{i.e.} le foncteur $f^{*} : \DD(Z) \to \DD(Z')$ induit une équivalence entre les sous-catégories formées des coefficients constants (ou encore un foncteur pleinement fidèle $\DD(Z)_{\text{loc. const.}} \longrightarrow \DD(Z')$).

$b'_1$) Si $f$ comme ci-dessus, alors $f$ est une $\DD$-équivalence 1-complète, \textit{i.e.} le foncteur $f^{*} : \DD(Z) \to \DD(Z')$ induit un foncteur pleinement fidèle
\msp{page 112}
entre les sous-catégories pleines formées des objets localement constants.

$b'_2$) Comme $b'_1$), mais avec $\DD$-équivalence 2-\underline{complète} au lieu de 1-complète, \textit{i.e.} $f^{*}$ induit une équivalence entre les sous-catégories formées des coefficients localement constants.

$b''_{\mathrm{coh}}$) Si $f : Z' \to Z$ dans \Dia{} est une \Hot-équivalence, alors c'est une $\DD$-équivalence cohomologique complète, \textit{i.e.} pour tout $F \in \mathsf{Ob}\ \DD(Z)$ qui est \underline{localement} \underline{constant}, on a
\[ H^{*}_{\DD}(Z,F) \xrightarrow{\ \sim\ } H^{*}_{\DD}(Z',f^{*}F) \quad . \]

$b''_{\mathrm{hom}}$) Si $f : Z' \to Z$ dans \Dia{} est une \Hot-équivalence, c'est une $\DD$-équivalence homologique complète, \textit{i.e.} pour tout $F \in \mathsf{Ob}\ \DD(Z)$ qui est \underline{localement} \underline{constant}, on a\footnote[59]{\textit{cf.} variante $\pm$ spéciale pour \HOT{} pages 124-125.}
\[ H_{*}^{\DD}(Z',f^{*}F) \xrightarrow{\ \sim\ } H_{*}^{\DD}(Z,F) \quad . \]

\printedpage{67}
% head: A. Grothendieck : Dérivateurs, Ch. I, Édition des Universités de Montpellier II et Paris VII
% foot: Version 8 juillet 2001

On a le diagramme d'implications tautologiques pour les variantes de \textbf{Der} 3 :
\[
\begin{tikzcd}
 & b'_2 \arrow[d, Rightarrow] \arrow[drr, Rightarrow] & & & \\
 & b'_1 \arrow[dl, Rightarrow] \arrow[dr, Rightarrow] \arrow[dd, Rightarrow] & & a' \arrow[ddd, Rightarrow] & \\
b''_{\mathrm{coh}} \arrow[dr, Rightarrow] & & b''_{\mathrm{hom}} \arrow[dl, Rightarrow] & & \\
 & b \arrow[drr, Rightarrow] & & & \\
 & & & a & .
\end{tikzcd}
\]
Dans tous les cas à ma connaissance, les dérivateurs satisfont la condition la plus forte \textbf{Der} 3 $b'_2$.
\msp{page 113}

Je rappelle aussi la :

\textbf{Définition 1 :} Soit $f : X \to Y$ dans \Dia. On dit que $f$ est une $\DD$-\textit{équivalence cohomologique} (sous-entendu : “\textit{complète}”) si pour tout $\DD$-coefficient $F$ \underline{localement} \underline{constant} sur $Y$,
\[ H^{*}_{\DD}(Y,F) \xrightarrow{\ \sim\ } H^{*}_{\DD}(X,f^{*}F) \quad ; \]
on dit que $f$ est une $\DD$-\textit{équivalence homologique} (“\textit{complète}”) si pour tout $F$ comme \texttt{[ci-]}dessus,
\[ H_{*}^{\DD}(X,f^{*}F) \xrightarrow{\ \sim\ } H_{*}^{\DD}(Y,F) \quad . \]
On dit que $f$ est une $\DD$-\textit{équivalence} 1-\textit{complète} (resp. 2-\textit{complète}) si le foncteur
\[ f^{*} : \DD(Y)_{\text{loc. const.}} \longrightarrow \DD(X)_{\text{loc. const.}} \]
est pleinement fidèle (resp. une équivalence).
\[
\begin{tikzcd}[column sep=small]
 & \text{$\DD$-équivalence 2-complète} \arrow[d, Rightarrow] & & \\
 & \text{$\DD$-équivalence 1-complète} \arrow[dl, Rightarrow] \arrow[dr, Rightarrow] & & \\
\text{$\DD$-équivalence cohomologique} \arrow[dr, Rightarrow] & & \text{$\DD$-équivalence homologique} \arrow[dl, Rightarrow] & \\
 & \text{$\DD$-équivalence} & & .
\end{tikzcd}
\]
Si dans les sous-catégories $\DD(X)_{\text{loc. const.}}$ il y a une loi $*$ (associative, unitaire, commutative) et des \underline{\textsf{Hom}} internes associés, préservés par les images inverses, alors les $\DD$-équivalences cohomologiques sont déjà 1-complètes et sont donc des $\DD$-équivalences homologiques, donc dans ce cas particulier
\printedpage{68}
% head: A. Grothendieck : Dérivateurs, Ch. I, Édition des Universités de Montpellier II et Paris VII
% foot: Version 8 juillet 2001
\msp{page 114 vide}
\msp{page 115}
\[
\begin{tikzcd}
\text{$\DD$-équivalence 2-complète} \arrow[d, Rightarrow] & & \\
\text{$\DD$-équivalence 1-complète} \arrow[r, Leftrightarrow] & \text{$\DD$-équivalence cohomologique} \arrow[d, Rightarrow] & \\
 & \text{$\DD$-équivalence homologique} \arrow[d, Rightarrow] & \\
 & \text{$\DD$-équivalence} & .
\end{tikzcd}
\]
Dans $\DD = \mathsf{HOT}$, toutes ces notions devraient être équivalentes, et équivalentes à la condition encore plus faible en apparence, que $f : X \to Y$ soit une \Hot-équivalence. Mais je ne connais que \HOT{} comme dérivateur où la $\DD$-équivalence d'une flèche $f : X \to Y$ dans \Dia{} implique déjà la $\DD$-équivalence cohomologique (complète), ou la $\DD$-équivalence homologique (complète), voire même la $\DD$-équivalence 1-complète. (Je reviendrai là-dessus tantôt.)

\textbf{Proposition 1.} \textit{Soit $f : X \to Y$ dans} \Dia. \textit{Conditions équivalentes :}

a) \textit{Pour tout $y \in \mathsf{Ob}\, Y$, l'inclusion $X_y \to X/y$ est une $\DD$-équivalence, et pour toute flèche $y' \to y$ dans $Y$, $X/y' \to X/y$ est une $\DD$-équivalence.}

b) \textit{Pour tout coefficient \underline{constant} $F$ sur $Y$, $f_{*}f^{*}F$ est localement constant, et se calcule “fibre par fibre”.}

(C'est la prop. 1, p. 76).

\textbf{Proposition 1$'$.} Dual.

\msp{page 116}

\textbf{Définition 2.} Sous les conditions équivalentes de la proposition 1, on dit que $f$ est une $\DD$-\textit{fibration cohomologique faible.} Sous les conditions duales, que $f$ est une $\DD$-\textit{fibration homologique faible.}

C'est ce que dans la définition 1, p. 77 j'appelais “$\DD$-fibration cohomologique” ou “$\DD$-fibration à gauche”, et dualement. Maintenant la notion me paraît trop faible pour mériter ce nom, que je réserve pour le cas “complet”.

\textbf{Proposition 2.} \textit{Les $\DD$-fibrations cohomologiques faibles sont stables par changements de base lisses, les $\DD$-fibrations homologiques faibles stables par changements de base propres.}

Je devrais construire un exemple où $f$ est une $\DD$-fibration cohomologique faible, et ne l'est plus après changement de base.


\endgroup
% END UNIT Der-ChI-234/Der-ChI-234_p066-070.tex
% BEGIN UNIT Der-ChI-234/Der-ChI-234_p071-075.tex
\begingroup\makeatletter%
\newcommand{\printedpage}[1]{\marginpar{\footnotesize p.~#1}}%
\newcommand{\msp}[1]{\marginpar{\footnotesize [#1]}}%
\newcommand{\D}{\mathbb{D}}%
\newcommand{\sfup}[1]{\textsf{\upshape #1}}%
\newcommand{\Dia}{\sfup{Dia}}%
\newcommand{\Hot}{\sfup{Hot}}%
\newcommand{\HOT}{\sfup{HOT}}%
\newcommand{\Ob}{\sfup{Ob}}%
\newcommand{\Hom}{\sfup{Hom}}%
\newcommand{\Cat}{$\mathcal{C}$\textit{at}}%
\makeatother\let\printedpage\derpp\let\msp\dermsp%
\let\mspin\DERmspin%


\printedpage{69}
% head: A. Grothendieck : Dérivateurs, Ch. I, Édition des Universités de Montpellier II et Paris VII

\textbf{Proposition 3.} \textit{Soit}
\[
\begin{tikzcd}
X \arrow[rr, "f"] \arrow[dr, "p"'] & & X' \arrow[dl, "p'"] \\
& S &
\end{tikzcd}
\]
\textit{dans \Dia, avec $p$ et $p'$ des $\D$-fibrations cohomologiques faibles} (resp. \textit{des $\D$-fibrations homologiques faibles}), \textit{soit $S_0 \subset S$ tel que $\pi_0(S_0) \to \pi_0(S)$ soit surjective. Pour que $f$ soit une $\D$-équivalence, il suffit que les $f_s : X_s \to X'_s$, pour $s \in S_0$, le soient.}

\textbf{Proposition 4.} \textit{Si $f : X \to Y$ dans \Dia{} est $\D$-propre} (resp. $\D$-\textit{lisse}), \textit{alors $f$ est une $\D$-fibration cohomologique} (resp. \textit{homologique}) \textit{faible si et seulement si pour tout $F$ constant sur $Y$, $f_*f^*F$} (resp. $f_!f^*F$) \textit{est localement constant sur $Y$. Et s'il en est ainsi, pour tout $Y' \to Y$ dans \Dia, $f' : X' = X \times_Y Y' \to Y'$ est encore une $\D$-fibration cohomologique} (resp. \textit{homologique}) \textit{faible.}

Immédiat sur le critère \textit{b}) de la proposition 1.

\msp{page 117}\textbf{Proposition 5.} \textit{Soit $f : X \to Y$ dans \Dia. Considérons les conditions suivantes :}

a) \textit{Pour tout $Y'' \to Y'$ sur $Y$ dans \Dia, si $Y''$ et $Y'$ sont \Hot-asphériques, alors}
\[
X'' = X \times_Y Y'' \longrightarrow X' = X \times_Y Y'
\]
\textit{est une $\D$-équivalence.}
\[
\begin{tikzcd}
X & X' \arrow[l] & X'' \arrow[l] \\
Y \arrow[from=u] & Y' \arrow[l] \arrow[from=u] & Y'' \arrow[l, "p"] \arrow[from=u]
\end{tikzcd}
\]

a$'$) \textit{Comme} a), \textit{mais en supposant $p : Y'' \to Y'$ une \Hot-équivalence (au lieu de supposer $Y'$, $Y''$ \Hot-asphériques).}

a$_1$) \textit{Comme} a), \textit{mais en supposant (au lieu de l'asphéricité) que $Y'$ a un objet final, et $Y''$ est la catégorie unité.}

a$_2$) \textit{Comme} a), \textit{mais en supposant que $Y'$ a un objet cofinal, et $Y''$ est la catégorie unité.}

b$_1$) \textit{$f$ est une $\D$-fibration cohomologique faible, et le reste après tout changement de base.}

b$_2$) \textit{$f$ est une $\D$-fibration homologique faible, et le reste après tout changement de base.}

\printedpage{70}
% head: A. Grothendieck : Dérivateurs, Ch. I, Édition des Universités de Montpellier II et Paris VII
\textit{On a alors le diagramme d'implications tautologiques}
\[
\begin{tikzcd}[column sep=small]
& & & a' \arrow[d, Rightarrow] & & & & \\
& & & a \arrow[dlll, Rightarrow] \arrow[drrr, Rightarrow] & & & & \\
a_1 & \Longleftrightarrow & b_1 & & b_2 & \Longleftrightarrow & a_2 & .
\end{tikzcd}
\]
\textit{Lorsque $\D$ satisfait à l'axiome} \textbf{Der} 3 a$'$ (p. 111), \textit{alors toutes ces conditions, sauf} a$'$), \textit{sont équivalentes. Lorsque $\D$ satisfait à l'axiome} \textbf{Der} 3 b$'_1$ \textit{(ou seulement} \textbf{Der} 3 b$''_{\mathrm{coh}}$ ou b$''_{\mathrm{hom}}$\textit{),} \textit{alors toutes les conditions (y inclus} a$'$\textit{) sont équivalentes.}

Pour simplifier, on va par la suite supposer que l'axiome \textbf{Der} 3 $a'$) est satisfait.

\textbf{Définition 3.} Sous les conditions équivalentes $a$, $a_1$, $b_1$, $a_2$, $b_2$ de la proposition 4, on dit que $f$ est une $\D$-\textit{fibration}. % [sic] proposition 4 -> proposition 5, the conditions a, a_1, b_1, a_2, b_2 are those of proposition 5

\msp{page 118}\textbf{Proposition 6.} \textit{Supposons que $f : X \to Y$ dans \Dia{} soit $\D$-propre} (resp. $\D$-\textit{lisse}). \textit{Conditions équivalentes :}

a) \textit{$f$ est une $\D$-fibration.}

b) \textit{$f$ est une $\D$-fibration cohomologique} (resp. \textit{homologique}) \textit{faible.}

c) \textit{Pour tout $y' \to y$ dans $Y$, $X/y' \to X/y$} (resp. $y\backslash X \to y'\backslash X$) \textit{est une $\D$-équivalence.}

d) \textit{Pour tout $F$ constant sur $Y$, $f_*f^*F$} (resp. $f_!f^*F$) \textit{est localement constant sur $Y$.}

d$'$) \textit{Pour tout $F$ localement constant sur $Y$, $f_*f^*F$} (resp. $f_!f^*F$) \textit{est localement constant sur $Y$.}

Cela résulte formellement de Prop. 4 et Prop. 5.

\textbf{Corollaire.} \textit{Supposons $f : X \to Y$} \Cat\textit{-fibrant} (resp. \Cat\textit{-cofibrant}). \textit{Pour que $f$ soit une $\D$-fibration, il faut et il suffit que pour tout $u : y \to y'$ dans $Y$, $u^* : X_{y'} \to X_y$} (resp. $u_* : X_y \to X_{y'}$) \textit{soit une $\D$-équivalence.}

Il est immédiat que la notion de $\D$-fibration est stable par changement de base, mais il est douteux qu'elle le soit par composition.\footnote[60]{N.B. Si les $\D$-équivalences sont déjà des \Hot-équivalences, alors la notion coïncide avec celle de \Hot-fibration, et elle est stable par composition. Également, si les $\D$-équivalences sont au sens des $\D$-équivalences cohomologiques, ou homologiques, alors la notion coïncide avec celle des $\D$-fibrations cohomologiques (resp. homologiques) ci-dessous, et est également stable par composition.} (Ici encore, j'ai omis de faire un contre-exemple, avec des dérivateurs $\D = \sfup{Der}(k\text{-}\sfup{Mod})$ disons.) Autre lacune de la notion : il n'est pas clair, si $f$ est une $\D$-fibration et $F$ un $\D$-coefficient sur $Y$, si
\[
f_*f^*F \qquad \text{et} \qquad f_!f^*F
\]
\printedpage{71}
% head: A. Grothendieck : Dérivateurs, Ch. I, Édition des Universités de Montpellier II et Paris VII
se calculent “fibre par fibre”, comme on s'attendrait dans le cas d'un “fibré”.

Par contre, voici un résultat sympa :

\msp{page 119}\textbf{Proposition 7.} \textit{Soit $f : X \to Y$ dans \Dia. Conditions équivalentes :}

a) \textit{$f$ est une $\D$-équivalence “universelle”,} i.e. \textit{pour tout $Y' \to Y$ dans \Dia, $f' : X' = X \times_Y Y' \to Y'$ est une $\D$-équivalence.}

b) \textit{$f$ est une $\D$-fibration à fibres $\D$-asphériques.}

\textbf{Définition 4.} Sous les conditions équivalentes de la proposition 7, on dit que $f$ est une $\D$-fibration \textit{triviale}.

Mais attention, $f$ peut-être une $\D$-fibration et une $\D$-équivalence, sans que la $\D$-fibration soit triviale \textit{i.e.} sans que les fibres soient $\D$-asphériques, cf. p. 87 et suivantes. % [sic] peut-être -> peut être, verb "pouvoir" followed by infinitive
Et également (devrait remonter) même si $\D$ satisfait \textbf{Der} 3 $b_2$ (la plus forte des variantes de \textbf{Der} 3), donc quand \underline{toutes} les conditions de la proposition 5, y inclus $a'$), sont équivalentes, celles-ci (\textit{i.e.} que $f$ soit une $\D$-fibration) n'impliquent pourtant pas, en général, que dans la condition $a'$) on puisse se borner à $p : Y'' \to Y'$ une $\D$-équivalence (au lieu d'une \Hot-équivalence). En d'autres termes, si on a un diagramme cartésien dans \Dia
\[
\begin{tikzcd}
X \arrow[d, "f"'] & X' \arrow[l, "q"'] \arrow[d, "f'"] \\
Y & Y' \arrow[l, "p"] & ,
\end{tikzcd}
\]
avec $f$ une $\D$-fibration, et $p$ une $\D$-équivalence, il ne s'ensuit pas forcément que $q$ soit une $\D$-équivalence. (Mais ce sera pourtant le cas, “bien sûr”, si $\D = \HOT$.)

\msp{page 120}Par contre, on a :

\textbf{Proposition 8.} \textit{Soit $f : X \to Y$ dans \Dia. Conditions équivalentes :}

a) \textit{$f$ est une $\D$-fibration.}

b) \textit{Pour tout $p : Y'' \to Y'$ sur $Y$, si $p$ est une $\D$-équivalence} 1-\textit{complète} (déf. 1, p. 113), \textit{alors $q : X'' \to X'$ est une $\D$-équivalence.}
\[
\begin{tikzcd}
X \arrow[d, "f"'] & X' \arrow[l] \arrow[d, "f'"'] & X'' \arrow[l, "q"'] \arrow[d, "f''"] \\
Y & Y' \arrow[l] & Y'' \arrow[l, "p"]
\end{tikzcd}
\]
\textbf{Corollaire.} \textit{Supposons que dans \Dia, toute $\D$-équivalence soit} 1-\textit{complète.} (N.B. Je ne connais que le dérivateur \HOT{} pour satisfaire à cette condition, qui paraît quasiment caractéristique de \HOT.) \textit{Alors $f : X \to Y$ est une $\D$-fibration si et seulement si pour tout $p : Y'' \to Y'$ sur $Y$, si $p$ est une $\D$-équivalence, alors $q : X'' \to X'$ aussi. Dans ce cas exceptionnel parmi les dérivateurs, la notion de $\D$-fibration est stable par composition} (ce qui est tautologique sur le critère que je viens de donner.)

\printedpage{72}
% head: A. Grothendieck : Dérivateurs, Ch. I, Édition des Universités de Montpellier II et Paris VII

On va renforcer la notion de $\D$-fibration, pour la rendre plus souple, et notamment stable par composition.

\textbf{Proposition 9.} \textit{Soit $f : X \to Y$ dans \Dia. Conditions équivalentes :}

a) \textit{Pour tout $y \in Y$, $X_y \to X/y$ est une $\D$-\underline{équivalence} \underline{cohomologique}} (déf. 1, p. 113), \textit{et de même pour $X/y \to X/y'$ pour toute flèche $y \to y'$ dans $Y$. Et ces conditions sur $f$ restent valables après tout changement de base} (N.B. il suffit que la première de ces conditions le reste).

\msp{page 121}b) \textit{pour tout $F \in \Ob\ \D(X)$ \underline{localement} \underline{constant}, $f_*F$ est localement constant sur $Y$ et se calcule “fibre par fibre” -- et itou après tout changement de base.} % Doubt resolved: attached PDF leaf 36 (printed p.72) has an en dash, encoded --.

c) \textit{Pour tout $p : Y'' \to Y'$ sur $Y$, si $p$ est une $\D$-équivalence cohomologique, alors il en est de même de $q : X'' \to X'$.}

c$'$) \textit{Comme} c, \textit{mais en supposant comme hypothèse sur $p$ que $Y''$ soit la catégorie finale, et $Y'$ ait un objet final.}

\textsc{Démonstration}. Dans l'ordre ($\pm$ tautologique)
\[
\mathrm{a} \Longrightarrow \mathrm{b} \Longrightarrow \mathrm{c} \Longrightarrow \mathrm{c}' \Longrightarrow \mathrm{a} \quad .
\]

\textbf{Définition 5.} Sous les conditions équivalentes de la proposition 9, on dit que $f$ est une $\D$-\textit{fibration cohomologique}. Sous les conditions duales, on dit que $f$ est une $\D$-\textit{fibration homologique}. On dit que $f$ est une $\D$-\textit{fibration forte}, si c'est à la fois une $\D$-fibration cohomologique et une $\D$-fibration homologique.

Ces notions renforcent la notion de $\D$-fibration. Elles ont l'avantage de plus d'être \underline{stables} \underline{par} \underline{composition}, non seulement par changement de base.

\msp{page 122}La notion de $\D$-fibration forte a été choisie pour coiffer les deux notions duales de $\D$-fibration cohomologique et homologique (à défaut d'une meilleur notion) tout en restant stable par composition. On pourrait exiger que pour tout % [sic] meilleur -> meilleure, agreement with "notion"
\[
Y'' \longrightarrow Y'
\]
sur $Y$, avec $Y'$, $Y''$ \Hot-asphériques, $q : X'' \to X'$ soit une $\D$-équivalence 1-complète (voire même 2-complète), appeler ça les $\D$-\textit{fibrations} 1-\textit{complètes} (resp. 2-\textit{complètes}). Ça renforce de façon assez élégante la notion de $\D$-fibration forte, mais a le grand tort de ne pas être stable par composition (que je sache!), sauf dans le cas particulier où on a des \underline{\Hom} internes sympas (cf. bas de la page 113), auquel cas la notion de $\D$-équivalence implique déjà celle de $\D$-équivalence 1-complète,
\printedpage{73}
% head: A. Grothendieck : Dérivateurs, Ch. I, Édition des Universités de Montpellier II et Paris VII
\msp{page 123}donc aussi celle de $\D$-équivalence homologique. Donc dans ce cas
\[
\begin{tikzcd}
\text{$\D$-fibration cohomologique} \arrow[r, Leftrightarrow] \arrow[d, Rightarrow] & \text{$\D$-fibration 1-complète} \\
\text{$\D$-fibration homologique} \quad , &
\end{tikzcd}
\]
et il n'y a pas à chercher plus loin que les $\D$-fibrations cohomologiques.

On pourrait songer, il est vrai, à définir (dans le cas général) une notion plus forte, en exigeant que pour tout $p : Y'' \to Y'$ sur $Y$ qui soit une $\D$-équivalence 1-complète, $q : X'' \to X'$ soit aussi une $\D$-équivalence 1-complète. Cette fois-ci, la notion est stable par composition, mais elle me paraît néanmoins inutilisable (en dehors du cas qu'on vient de dire, où elle coïncide avec la notion de $\D$-fibration cohomologique, ou encore de $\D$-fibration forte), car n'étant pas en mesure de décrire les $\D$-équivalences 1-complètes, je suis aussi incapable de donner des exemples non triviaux de $\D$-fibrations “strictement 1-complètes”.

Pour vérifier qu'une $f : X \to Y$ dans \Dia{} est une $\D$-fibration cohomologique, il faut
\msp{page 124}savoir vérifier que des flèches $q : Z' \to Z$ dans \Dia{} sont des $\D$-équivalences cohomologiques, donc que $H^*_{\D}(Z,F) \overset{\sim}{\rightarrow} H^*_{\D}(Z',q^*(F))$ est \texttt{[un]} isomorphisme pour tout $F \in \Ob\ \D(Z)$ qui est localement constant sur $Z$. Or je ne sais vérifier un énoncé de cette nature (en dehors du cas où je sais d'avance que $Z$ est complètement $\D$-asphérique, \textit{i.e.} que les coefficients localement constants sur $Z$ sont en fait constant) \underline{que} si $q$ est déjà une \Hot-équivalence. % [sic] constant) -> constants), agreement with "coefficients"
C'est dire que je ne sais exhiber des $\D$-fibrations cohomologiques, que quand ce sont déjà des \HOT-fibrations cohomologiques. Je veux, dans la suite, donner des critères utilisables pour des \HOT-fibrations. Mais on va axiomatiser ces développements, en gardant un $\D$ général, sur lequel je vais faire une hypothèse si draconienne que pratiquement, je ne connais que le dérivateur \HOT{} qui y satisfait. Il s'agit de $c_1$ ci dessous.

\textbf{Der} 3 $c_1$) Toute $\D$-équivalence est une $\D$-équivalence cohomologique.

$c_2$) Toute $\D$-équivalence est une $\D$-équivalence homologique.

\msp{page 125}$c'$) Toute $\D$-équivalence est 1-complète.

$c''$) Toute $\D$-équivalence est 2-complète.

N.B. On a
\[
\begin{tikzcd}
& c'' \arrow[d, Rightarrow] & \\
& c' \arrow[dl, Rightarrow] \arrow[dr, Rightarrow] & \\
c_1 & & c_2
\end{tikzcd}
\]


\endgroup
% END UNIT Der-ChI-234/Der-ChI-234_p071-075.tex
% BEGIN UNIT Der-ChI-234/Der-ChI-234_p076-076.tex
\begingroup\makeatletter%
\newcommand{\printedpage}[1]{\marginpar{\footnotesize p.~#1}}%
\newcommand{\msp}[1]{\marginpar{\footnotesize [#1]}}%
\newcommand{\HOT}{\textsf{HOT}}%
\newcommand{\Hot}{\textsf{Hot}}%
\newcommand{\Dia}{\textsf{Dia}}%
\newcommand{\Cat}{\mathcal{C}at}%
\newcommand{\D}{\mathbb{D}}%
\makeatother\let\printedpage\derpp\let\msp\dermsp%
\let\mspin\DERmspin%


\printedpage{74}
% head: A. Grothendieck : Dérivateurs, Ch. I, Édition des Universités de Montpellier II et Paris VII
% foot: Version 8 juillet 2001
tautologiquement. Je ne connais que \HOT{} pour satisfaire l'ensemble des deux plus faibles conditions $c_1$, $c_2$. Et en fait, sous réserve d'écrire la théorie du dérivateur \HOT, il doit être vrai qu'elle satisfait même la forme la plus forte $c''$. (Et les \HOT-équivalences doivent être simplement les \Hot-équivalences.) (Cf. le “lemme” p.~109.)

La proposition suivante est une tautologie.

\noindent\textbf{Proposition 10.} \textit{Supposons que $\D$ satisfasse la condition \textup{\textbf{Der} 3 $\mathrm{c}_1$ (p.~124),} par exemple $\D = \text{\HOT}$. Alors la notion de $\D$-fibration coïncide avec celle de $\D$-fibration cohomologique. Elle est donc stable par composition. Pour que $f : X \to Y$ soit une $\D$-fibration, il faut et il suffit que pour $p : Y'' \to Y'$ sur $Y$, si $p$ est une $\D$-équivalence, il en est de même de $q : X'' \to X'$.}

\noindent\textbf{Corollaire 1.} \textit{Si $f : X \to Y$ dans \textup{\Dia} est $\D$-propre, alors $f$ est une $\D$-fibration cohomologique si et seulement si pour toute flèche $y' \to y$ dans $Y$, $X/y' \to X/y$ est une $\D$-équivalence \textup{(i.e.} si et seulement si $f$ est une $\D$-fibration cohomologique faible\textup{).}}

Appliquer le critère $c$) de \texttt{[la]} prop.~6.

\msp{page 126}

\noindent\textbf{Corollaire 2.} \textit{Supposons que $\D$ satisfasse \textup{\textbf{Der} 3 $(\mathrm{c}_1, \mathrm{c}_2)$ (par exemple $\D = \text{\HOT}$).} Alors pour une flèche $f : X \to Y$ dans \textup{\Dia}, il revient au même que $f$ soit une $\D$-fibration, une $\D$-fibration cohomologique, ou une $\D$-fibration homologique. Si $f$ est propre \textup{(resp.} lisse\textup{),} il revient au même de dire que $f$ est une $\D$-fibration cohomologique \textup{(resp.} homologique\textup{),} faible.}

On voit donc que dans $\Cat$, on doit avoir une seule notion raisonnable de \HOT-fibration, qui coïncidera avec celle de \HOT-fibration cohomologique ou homologique. Mais pour que cette notion soit dès à présent disponible, il faudrait que le dérivateur \HOT{} soit déjà construit sur le domaine $\text{\Dia} = \Cat$ tout entier. Or ce n'est pas chose faite encore (et il n'a donc pas pu être question de prouver les propriétés de \HOT{} dont il vient d'être question, disons la plus forte \textbf{Der} 3 $c''$).

Néanmoins, je voudrais développer de façon directe (“élémentaire”) la notion de \Hot-fibration, sans me réferer à la théorie du dérivateur \HOT, et aux développements théoriques précédents. % [sic] réferer -> référer, missing accent on the second e


\endgroup
% END UNIT Der-ChI-234/Der-ChI-234_p076-076.tex
\clearpage\setcounter{footnote}{0}\phantomsection\addcontentsline{toc}{section}{Chapitre II}\markboth{Chapitre II}{Chapitre II}
% BEGIN UNIT Der-ChII-234/Der-ChII-234_p001-005.tex
\begingroup\makeatletter%
\newcommand{\printedpage}[1]{\marginpar{\footnotesize p.~#1}}%
\newcommand{\msp}[1]{\marginpar{\footnotesize [#1]}}%
\newcommand{\Vcat}[2]{\underset{#2}{\underset{\vee}{#1}}}%
\makeatother\let\printedpage\derpp\let\msp\dermsp%
\let\mspin\DERmspin%


% ===================== PDF page 1 =====================
\printedpage{--}
% no page number printed (title page)
% head: (none)

\begin{center}
\vspace*{4cm}
{\Large\textsf{LES DÉRIVATEURS}}

\bigskip\bigskip

\textbf{ALEXANDRE GROTHENDIECK}

\bigskip\bigskip\bigskip

\textbf{Chapitre II}

\textbf{Cofinalité (à droite et à gauche)}

\textbf{(préliminaire à la cofinalité cohomologique)}
\end{center}

\par

% ===================== PDF page 2 =====================
\printedpage{--}
% no page number printed (editorial note)
% head: (none)

\vspace*{6cm}

\noindent \par\noindent\textbf{Note des éditeurs de l’édition transcrite.}\par
Ce texte a été déchiffré et transcrit en \LaTeXe{} par M. Künzer. Il a été édité par M. Künzer, J. Malgoire, et G. Maltsiniotis. La transcription est aussi fidèle que possible au manuscrit. Pour les quelques corrections évidentes, ou rares commentaires des éditeurs, ainsi que pour la numérotation originale des pages du manuscrit, les caractères de machine à écrire \texttt{[typewriter]} entre crochets sont utilisés. Un point d'interrogation entre crochets signifie que l'on n'est pas sûr du mot qui précède.

\noindent Cette édition est provisoire. Les remarques, commentaires et corrections sont bienvenus. Envoyer un message à :

\noindent\textbf{[e-mail address removed]}

\hfill G. Maltsiniotis\hspace*{2cm}

\par

% ===================== PDF page 3 =====================
\printedpage{1}
% head: A. Grothendieck : Dérivateurs, Ch. II, Édition des Universités de Montpellier II et Paris VII
% foot: Version 8 juillet 2001

\msp{page 1}

\section*{1\quad Cofinalité cohomologique}

Soit
\[ \alpha : C' \longrightarrow C \]
un foncteur de (petites) catégories, $\mathcal{M}$ une catégorie ayant les $\varprojlim$ pertinentes, $F$ un préfaisceau (donc contravariant en $i$) sur $C$ à valeurs dans $\mathcal{M}$. On a un homomorphisme canonique dans $\mathcal{M}$
\[ \varprojlim_{C} F \longrightarrow \varprojlim_{C'} \underbrace{\alpha^*(F)}_{F\circ\alpha} \tag*{(1)} \]
et on veut des conditions sur $\alpha$ pour que ce soit toujours un isomorphisme (\footnote{On dit alors que $\alpha$ est \textit{cofinal,} cf. SGA 4, I 8.1, ou 3.3, p. 45.}). On est ramené aussitôt au cas où $\mathcal{M} = (\mathrm{Ens})$. Disons qu'une catégorie $I$ est \textit{faiblement filtrante à gauche} si $I \neq \emptyset$ et si $\forall\, i', i'' \in \operatorname{Ob} I$, $\exists\, i \in \operatorname{Ob} I$ et des flèches
\[
\begin{tikzcd}[row sep=small, column sep=small]
 & i' \\
i \arrow[ur] \arrow[dr] & \\
 & i''.
\end{tikzcd}
\]
Cela signifie aussi que $\operatorname{Ob} I$ en tant qu'ensemble ordonné est filtrant à gauche.

\textbf{Proposition.} \textit{Pour que} (1) \textit{soit toujours un isomorphisme} ($\alpha$ \textit{`anodin pour} $\varprojlim$ \textit{'), il suffit que pour tout} $x_0 \in \operatorname{Ob} C$\textit{, la catégorie} $x_0\backslash C'$ \textit{des couples} $(x', u)$\textit{, avec} $x' \in \operatorname{Ob} C'$\textit{,} $u : \alpha(x') \longrightarrow x$\textit{, soit faiblement filtrante à gauche} (\footnote{Cf. amélioration des remarques p. 5.}, \footnote{Cf. 3.2, p. 43-44, pour conditions nécessaires et suffisantes de cofinalité.}). % [sic] u : \alpha(x') \longrightarrow x -> x_0 in place of x, the category is x_0\backslash C'

Ceci n'implique pas la conclusion plus forte
\[ F \underbrace{\xrightarrow{\;\sim\;}}_{\text{adj.}} \alpha_*\alpha^* F \]
que le morphisme d'adjonction canonique soit un isomorphisme. Pour qu'il soit ainsi, il suffit que $\forall\, x_0 \in C$, le foncteur $\alpha/x_0 : C'/x_0 \longrightarrow C/x_0$ induit par $\alpha$ soit `anodin pour les $\varprojlim$ ', ce qui signifie aussi que pour toute flèche $u : x_1 \longrightarrow x_0$ dans $C$, la catégorie $\Vcat{C'}{u}$ formée des triples $(x', \xi, \eta)$ avec $x' \in \operatorname{Ob} C'$ et $\xi : x_1 \longrightarrow \alpha(x')$, $\eta : \alpha(x') \longrightarrow x_0$ rendant commutatif
% printed: C' inside a V-shaped wedge, with u printed below the wedge
\msp{page 2}
le diagramme
\printedpage{2}
% head: A. Grothendieck : Dérivateurs, Ch. II, Édition des Universités de Montpellier II et Paris VII
% foot: Version 8 juillet 2001
\[
\begin{tikzcd}
x_1 \arrow[rr, "u"] \arrow[dr, "\xi"'] & & x_0 \\
 & \alpha(x'), \arrow[ur, "\eta"'] &
\end{tikzcd}
\]
est faiblement filtrante à gauche.

Plus généralement, considérons
\[
\begin{tikzcd}
X' = C' \arrow[rr, "\alpha"] \arrow[dr, "f'"'] & & C = X \arrow[dl, "f"] \\
 & S &
\end{tikzcd}
\]
et posons nous la question quand l'homomorphisme fonctoriel, par rapport à un préfaisceau $F$ variable sur $X$,
\[ f_*(F) \longrightarrow f'_*(\alpha^*(F)) \]
est toujours un isomorphisme. (Le cas précédent est celui où $S = X$, $f = \mathrm{id}_X$.) On trouve qu'il suffit que $\forall\, s_0 \in \operatorname{Ob} S$, le foncteur
\[ \alpha/s_0 : X'/s_0 \longrightarrow X/s_0 \]
soit `anodin pour les $\varprojlim$ ', et pour ceci, il suffit que pour \texttt{[tout]} objet $(x_0, u_0)$ de $X/s_0$ (où
\[ f(x_0) \xrightarrow{u_0} s_0) \]
la catégorie $(x_0, u_0)\backslash(X'/s_0)$ soit faiblement filtrante à gauche.

\msp{page 3}

Cette catégorie s'interprète comme celle des triples $(x', \xi, u)$ avec
\[
\left\{
\begin{aligned}
& x' \in \operatorname{Ob} C' \\
& u : g(x') = f(\alpha(x')) \longrightarrow s_0 \\
& \xi : x_0 \longrightarrow \alpha(x')
\end{aligned}
\right.
\]
\texttt{[plutôt} $f'(x')$ \texttt{au lieu de} $g(x')$\texttt{]} de façon à rendre commutatif le diagramme
\[
\begin{tikzcd}
f(x_0) \arrow[rr, "u_0"] \arrow[dr, "f(\xi)"'] & & s_0 \\
 & f(\alpha(x')). \arrow[ur, "u"'] &
\end{tikzcd}
\]
Je voudrais dégager des variantes de ces énoncés dans le cadre cohomologique, les foncteurs $\varprojlim$ étant remplacés par des $R\Gamma$ (de préfaisceaux abéliens p.ex. - mais on pense aussi aux `cas non commutatifs', définis par des dérivateurs notamment.)

\par

% ===================== PDF page 5 =====================
\printedpage{3}
% head: A. Grothendieck : Dérivateurs, Ch. II, Édition des Universités de Montpellier II et Paris VII
% foot: Version 8 juillet 2001

Si $I$ est une catégorie \textit{ordonnée finie}, la condition \texttt{[faiblement]} filtrante à gauche équivaut à l'existence d'un objet initial, ce qui implique que $I$ est contractile. Dans les énoncés précedents, il se pourrait que la bonne philosophie soit de remplacer partout `faiblement filtrante à gauche' par `contractile'. Si p.ex. $C$ est la catégorie finale, alors remplacer la condition que $C'$ soit filtrante à gauche, par l'acyclicité. C'est en tous cas en accord avec l'esprit de l'axiome $\mathrm{Der}_5$ (cas où $J$ est la catégorie finale). En effet, dans le cas où $C = e$, la condition que $\alpha$ soit anodin pour les $\varprojlim$ équivaut visiblement à dire que le foncteur $\alpha^* : \mathcal{M} \longrightarrow \underline{\mathrm{Hom}}(C'^{o}, \mathcal{M})$, associant à tout $F \in \mathcal{M}$ le préfaisceau constant de valeur $F$ % [sic] précedents -> précédents, missing accent on the second e
\msp{page 4}
sur $C'$, donne un morphisme d'adjonction
\[ \varprojlim_{e} F = F \longrightarrow \alpha_*(\alpha^*(F)) = \varprojlim_{C'} \alpha^* F \]
qui est un isomorphisme. Et l'axiome $\mathrm{Der}_5$ assure qu'il doit en être ainsi si $C'$ est contractile.

La condition filtrante dans la proposition est d'ailleurs une condition \textit{suffisante} d'anodinité, mais nullement nécessaire. Si p.ex. $C$ est la catégorie finale, alors la condition d'anodinité pour $C' \longrightarrow C$ est simplement que $C'$ soit connexe non vide, i.e. $\pi_0(C') = \{0\}$. (Condition qui visiblement néglige les $\pi_i$ supérieurs.)

Soit \fbox{$\alpha : X' \longrightarrow X$} un foncteur de petites catégories. La condition suggérée par la proposition, que les $x_0\backslash X'$ soient contractiles, signifie que $\alpha$ est une hot-équivalence co-stricte, i.e. que $\alpha^{o} : X'^{o} \longrightarrow X^{o}$ est une hot-équivalence `localement sur $X^{o}$', donc aussi globalement. Si $\alpha$
\msp{page 5}
lui-même était une hot-équivalence stricte, au lieu de $\alpha^{o}$, alors on aurait par $\mathrm{Der}_5$, que pour $F \in \mathbf{D}(X)$, $F \xrightarrow{\;\sim\;} \alpha_*\alpha^* F$ ($\mathbf{D}$ dérivateur à gauche), ce qui est bien plus fort que $\Gamma_{\mathbf{D}}(X, F) \xrightarrow[\sim]{\alpha^*} \Gamma_{\mathbf{D}}(X', \alpha^* F)$. La condition duale implique (si $\mathbf{D}$ \texttt{[est un]} dérivateur à droite) $\alpha_!\alpha^* F \xrightarrow{\;\sim\;} F$. Mais est-il raisonnable de s'attendre qu'elle implique aussi la condition `du côté cohomologie, au lieu de l'homologie',
\[ \Gamma_{\mathbf{D}}(X, F) \xrightarrow{\;\sim\;} \Gamma_{\mathbf{D}}(X', \alpha^*(F)) \ \ ? \]

Le premier cas, qui devrait emporter conviction, est celui du dérivateur DerAb, i.e. des catégories dérivées des faisceaux abéliens :

\textbf{Question.} \textit{Soit} $\alpha : X' \longrightarrow X$ \textit{une} hot-\textit{équivalence co-stricte,} $F$ \textit{un préfaisceau abélien sur} $X$\textit{. A-t-on alors} $H^i(X, F) \xrightarrow[\sim]{\alpha^*} H^i(X', \alpha^*(F))$ \textit{pour tout} $i \in \mathbf{Z}$ \textit{?}

Si $X$ admet un objet initial $\emptyset$, alors $\emptyset\backslash X' \simeq X'$, donc $X'$ contractile, donc c'est OK pour $F$ constant.

\textbf{Remarque.} Je m'aperçois que dans la proposition 1, au lieu de la condition que les $x\backslash C'$ soient faiblement filtrantes à gauche, il suffirait qu'elles soient non vides et connexes, i.e.


\endgroup
% END UNIT Der-ChII-234/Der-ChII-234_p001-005.tex
% BEGIN UNIT Der-ChII-234/Der-ChII-234_p006-010.tex
\begingroup\makeatletter%
\newcommand{\printedpage}[1]{\marginpar{\footnotesize p.~#1}}%
\newcommand{\msp}[1]{\marginpar{\footnotesize [#1]}}%
\newcommand{\Homint}{\operatorname{Hom\,int}}%
\makeatother\let\printedpage\derpp\let\msp\dermsp%
\let\mspin\DERmspin%


\printedpage{4}
% head: A. Grothendieck : Dérivateurs, Ch. II, Édition des Universités de Montpellier II et Paris VII
% foot: Version 8 juillet 2001
0-connexes (\footnote[4]{Cela signifie aussi que $\alpha$ est $W_0$-coasphérique. Pour un traitement ab ovo, cf. prop. 3.2, page 43.}). Le fait que les $x\backslash C'$ non vide implique que $\varprojlim\limits_{C} \longrightarrow \varprojlim\limits_{C'}$ soit \textit{injectif.} Pour la sujectivité, si $\xi' \in \varprojlim\limits_{C'}$ (cas ensembliste), on définit $\xi \in \varprojlim\limits_{C}$ par les $\xi(x) \in F(x)$ à l'aide de % [sic] sujectivité -> surjectivité, lettre r omise
\msp{page 6}
\[ x \xrightarrow{u} \alpha(x') \]
en prenant $u^*(\xi(x'))$ (où $\xi(x') \in (\alpha^*(F))(x') = F(\alpha(x')))$, le point essentiel est de montrer que les $u^*(\xi(x'))$ ne dépendent que de $x$, pas du choix de $(x',u) \in x\backslash C'$. Or pour $(x',u)$ variable dans $x\backslash C'$, l'application
\[
\begin{aligned}
\varphi : \operatorname{Ob}(x\backslash C') &\quad\longrightarrow\quad F(x)\\
(x',u) &\quad\longmapsto\quad u^*(\xi(x'))
\end{aligned}
\]
est telle que s'il existe une flèche
\[ (x',u) \xrightarrow{u'} (\bar{x}',\bar{u}) \qquad \text{où} \qquad u' : x' \longrightarrow \bar{x}' \]
(donc le diagramme
\[
\begin{tikzcd}
 & \alpha(x') \arrow[dd, "\alpha(u')"] \\
x \arrow[ur, "u"] \arrow[dr, "\bar{u}"'] & \\
 & \alpha(\bar{x}')
\end{tikzcd}
\]
est commutatif), alors $\varphi(x',u) = \varphi(\bar{x}',\bar{u})$. Donc $\varphi$ est constant sur les composantes connexes de $x\backslash C'$, donc constant si $x\backslash C'$ est 0-connexe, d'où $\xi(x)$. On prouve alors aussitôt que pour $x$ variable,
\[ \xi \overset{\text{déf}}{=} (\xi(x))_{x\in\operatorname{Ob} C} \simeq \varprojlim F, \qquad \text{\texttt{[plutôt}}\ \ \in \varprojlim F, \text{\texttt{]}} \]
et que $\alpha^*(\xi) = \xi'$, cqfd.

On peut dire que pour $x_0 \in \operatorname{Ob} C_0$ fixé, si on regarde le préfaisceau \textit{constant} sur $x_0\backslash C'$ de valeur $F(x_0)$, soit $\Phi_{x_0}$, l'application
\[ (x',u) \longmapsto \xi^*(u') : x_0\backslash C' \longrightarrow F(x_0) \]
peut être regardée comme une section de ce préfaisceau sur $x\backslash C'$, défini par $\xi' \in \varprojlim\limits_{C'} \alpha^*F$, donc on a
\[ \varprojlim_{C'} \alpha^*F \xrightarrow{c_{x_0}} \Gamma(x_0\backslash C', \Phi_{x_0}) \qquad (\xleftarrow{\ \sim\ } F(x_0) \text{ si } x_0\backslash C \text{ 0-connexe}). \]
Si on a \fbox{$x_0 \xrightarrow{v} x_1$} dans $C'$, on en déduit
\[ x_1\backslash C' \xrightarrow{V} x_0\backslash C' \]
(les $x\backslash C'$ sont les fibres d'une \textit{catégorie fibrée} $C\backslash C'$ sur $C$ \dots) et on a $F(x_1) \longrightarrow F(x_0)$, d'où \fbox{$\Phi_{x_1} \longrightarrow V^*\Phi_{x_0}$} pour les faisceaux constants correspondants sur $x_0\backslash C'$ resp. $x_1\backslash C'$
\printedpage{5}
% head: A. Grothendieck : Dérivateurs, Ch. II, Édition des Universités de Montpellier II et Paris VII
% foot: Version 8 juillet 2001
\msp{page 7}
et on a commutativité dans
\[
\begin{tikzcd}[column sep=large]
\varprojlim\limits_{C'} \alpha^*F \arrow[r, "c_{x_1}"] \arrow[d, "c_{x_0}"'] & \Gamma(x_0\backslash C', F(x_0)) \simeq \operatorname{Hom}(\pi_0(x_0\backslash C'), F(x_0)) \arrow[d, "V^*"] & \\
\operatorname{Hom}(\pi_0(x_1\backslash C'), F(x_1)) \simeq \Gamma(x_1\backslash C', F(x_1)) \arrow[r, "\substack{v^*\\ (F(x_1) \longrightarrow F(x_0))}"'] & \Gamma(x_1\backslash C', F(x_0)) \simeq \operatorname{Hom}(\pi_0(x_1\backslash C'), F(x_0)) & \\
\pi_0(x_1\backslash C') \arrow[r, "\pi_0(V)"] \arrow[d, dotted] & \pi_0(x_0\backslash C') \arrow[d, dotted] & \text{commutatif.} \\
F(x_1) \arrow[r, "F(v)"] & F(x_0) &
\end{tikzcd}
\]
On considère le préfaisceau d'ensembles
\[ x \longmapsto \pi_0(x\backslash C') \qquad \text{sur } C, \]
appelons-le $\pi_0(C'/C)$, on aura un morphisme canonique
\[ \boxed{\varprojlim_{C'} \alpha^*(F) \xrightarrow{\ \sim\ } \operatorname{Hom}_{\substack{\text{préfaisceaux}\\ \text{d'ensembles sur $C$}}} (\pi_0(C'/C), F),} \]
sans doute toujours \texttt{[un]} isomorphisme.

Si on travaille dans la catégorie dérivée des préfaisceaux des $k$-modules ($k$ un anneau), la formule précédente suggère conjecturalement
\[ \boxed{\mathbf{R}\Gamma(C', \alpha^*F) \simeq \mathbf{R}\underbrace{\operatorname{Hom}}_{\text{Hom internes}}\ (\underbrace{\mathbf{L}\alpha_!(k_{C'})}_{\substack{\text{homologie}\\ \text{relative de}\\ \text{$C'$ sur $C$}}}, F).} \]

\msp{page 8}
L'augmentation canonique
\[ \mathbf{L}\alpha_!(k_{C'}) \longrightarrow k_C \]
redonne, en prenant les $\mathbf{R}\mathrm{Hom}$ dans $F$,
\[ \mathbf{R}\Gamma(C,F) \longrightarrow \mathbf{R}\operatorname{Hom}(\mathbf{L}\alpha_!(k_{C'}), F) \xleftarrow{\ \sim\ } \mathbf{R}\Gamma(C', \alpha^*F) \]
l'homomorphisme canonique de `contravariance' de la cohomologie.

Si par contre on travaille avec le dérivateur non additif HOT, donc $F$ un `type d'homotopie relatif' sur $C$, on doit avoir
\[ \boxed{\mathbf{R}\Gamma_{\mathrm{HOT}}(C', \alpha^*F) \xrightarrow{\ \sim\ } \mathbf{R}\Gamma(C\,,\ \underbrace{\operatorname{Hom}}_{\substack{\text{Hom interne}\\ \text{de la catégorie des}\\ \text{types d'homotopie}\\ \text{relatifs sur $C$}}}\ (\ \alpha_!\ (\ \underbrace{e_{C'}}_{\substack{\text{objet final}\\ \text{de HOT($C'$),}\\ \text{type d'homotopie}\\ \text{relatif final}\\ \text{sur $C'$}}}\ ),\ F)).} \]

\printedpage{6}
% head: A. Grothendieck : Dérivateurs, Ch. II, Édition des Universités de Montpellier II et Paris VII
% foot: Version 8 juillet 2001
En fait, ces formules sont quasi-tautologiques, si on admet l'existence des foncteurs $\alpha_!$ comme adjoints \textit{à gauche} des $\alpha^*$. Dans le cas abélien, disons,
\msp{page 9}
l'isomorphie des $H^0$ des deux membres est la formule d'adjonction
\[ \operatorname{Hom}_{C'}(k_{C'}, \alpha^*F) \simeq \operatorname{Hom}_C(\alpha_!(k_{C'}), F), \]
d'où l'isomorphie pour les $H^i$ des deux membres, par décalage. De même dans le cas de HOT, l'isomorphie canonique des $\pi_0$ des deux membres est la formule d'adjonction. Pour avoir l'isomorphisme plus fin dans $\mathbf{D}(e)$, la catégorie fondamentale du dérivateur envisagé, en termes de Hom internes, on doit admettre l'existence d'objets unité (les $k_X$ dans un cas, les $e_X$ \texttt{[plutôt} $\varepsilon_X$\texttt{]} dans l'autre), compatibles avec images inverses, de telle façon \texttt{[que]} $\Homint(\varepsilon_X, \underbrace{F}_{\in\mathbf{D}_X}) \simeq F$. On aura
\[ \Gamma(C', \alpha^*F) \overset{\text{déf}}{=} f'_*(\alpha^*(F)) \simeq f_*(\alpha_*\alpha^*(F)), \]
et la question est donc de donner une interprétation de $\alpha_*\alpha^*F$ comme
\[ \alpha_*\alpha^*F \simeq \Homint(\alpha_!(\varepsilon_{C'}), F) \qquad \text{dans } \mathbf{D}(C). \]
Pour ceci, on va écrire
\[ \alpha^*(F) \simeq \Homint(\varepsilon_{C'}, F), \]
et la formule doit résulter d'une `formule d'adjonction interne'
\msp{page 10}
\[ ? \qquad\qquad \boxed{\alpha_*(\Homint_{C'}(G', \alpha^*F)) \simeq \Homint_C(\alpha_!G', F).} \]
(\textit{Formule d'adjonction `interne',} pour $\alpha_!$ et $\alpha^*$, valable en principe pour \textit{tout} $\alpha : C' \longrightarrow C$ dans Diag, tout $G'$ dans $\mathbf{D}(C')$ et tout $F$ dans $\mathbf{D}(C)$.)

Comment prouver cette formule ? En prenant l'opération produit $\otimes$ associé à l'opération Hom\,int, on est ramené (en prenant les foncteurs contravariants associés) à une formule
\[ G \otimes \alpha_!G' \xrightarrow{\ \sim\ } \alpha_!(\alpha^*G \otimes G'). \]
On trouve par adjonction un homomorphisme canonique, il faut voir s'il est toujours \texttt{[un]} isomorphisme.

\printedpage{7}
% head: A. Grothendieck : Dérivateurs, Ch. II, Édition des Universités de Montpellier II et Paris VII
% foot: Version 8 juillet 2001
\msp{page 11}

\section*{2\quad Les catégories localement filtrantes sont les catégories localement totalement $0$-connexes}

Soit $\mathcal{T}$ un topos. On dit que $\mathcal{T}$ est \textit{irréductible} si l'ensemble ordonné $\operatorname{Ouv}(\mathcal{T})$ des sous-objets de $e_{\mathcal{T}}$ a la propriété suivante : a) $\operatorname{Ouv}(\mathcal{T})$ n'est pas réduit à un élément, i.e. $\emptyset_{\mathcal{T}} \longrightarrow e_{\mathcal{T}}$ n'est pas un isomorphisme ($\mathcal{T}$ n'est pas identique à son sous-topos vide). b) Si $U'$, $U''$ sont deux objets `non vides', i.e. non initiaux, de $\operatorname{Ouv}(\mathcal{T})$, leur intersection est `non vide'. Il revient au même pour b) de dire que si $\mathcal{T}'$, $\mathcal{T}''$ sont deux sous-topos fermés distinct de $\mathcal{T}$, alors leur supremum $\mathcal{T}' \cup \mathcal{T}''$ est distinct de $\mathcal{T}$, ou encore que tout ouvert de $\mathcal{T}$ est connexe. On dit que $\mathcal{T}$ est \textit{localement irréductible} si a) $\mathcal{T}$ est localement connexe (i.e. chaque faisceau sur $\mathcal{T}$ est somme de faisceaux 0-connexes), et si b) pour tout faisceau 0-connexe $U$ sur $\mathcal{T}$, le topos induit $\mathcal{T}/U$ est irréductible (i.e. l'ensemble ordonné $\operatorname{Ouv}(U)$ des sous-objets de $U$ a la propriété b) ci-dessus ( a) signifie simplement que $U \not\simeq \emptyset_{\mathcal{T}}$, i.e. que $U$ n'est pas objet initial dans $\mathcal{T}^\wedge$ ; ou plus simplement) : tout sous-faisceau d'un faisceau connexe est connexe (\footnote[5]{\textbf{NB} Irréductible n'implique pas localement irréductible, pas plus que connexe n'implique pas localement connexe etc.}) (\footnote[6]{$\mathcal{T}$ est localement irréductible si et seulement s'il a un ensemble des générateurs $(U_\alpha)$, avec les $U_\alpha$ irréductibles (donc 0-connexes, donc $\mathcal{T}$ est localement connexe).}).

Quand $\mathcal{T}$ est de la forme $\operatorname{Top}(T)$, où $T$ est un espace topologique, alors $\mathcal{T}$ irréductible signifie que $T$ est irréductible (car alors $\operatorname{Ouv}(\mathcal{T}) \simeq \operatorname{Ouv}(T)$). Cela implique déjà trivialement que $T$ est localement connexe (tout ouvert est somme disjoint d'ouverts connexes), donc $\operatorname{Top}(\mathcal{T})$ est localement connexe. De plus, cela implique aussi que $\operatorname{Top}(T)$ est localement irréductible (\footnote[7]{expliciter la démonstration.}). (\textbf{NB} Les objets 0-connexes % [sic] \operatorname{Top}(\mathcal{T}) -> \operatorname{Top}(T), Top s'applique à l'espace T et non au topos \mathcal{T}
\msp{page 12}
de $\operatorname{Top}(T)^\wedge$, i.e. les faisceaux d'ensembles 0-connexes sur $T$, sont ceux dont la `fibre générique' a exactement un point \dots) Aussi $\mathcal{T} = \operatorname{Top}(T)$ est localement irréductible si et seulement si $T$ \texttt{[est]} localement connexe, et les composantes connexes de $T$ sont irréductibles, i.e. si et seulement si $T = \coprod_i T_i$, les $T_i$ irréductibles.

Si $\mathcal{T}$ est de la forme $\operatorname{Top}(X)$, $X$ une catégorie, on sait que $\mathcal{T}$ est localement connexe (et même localement $\infty$-connexe). On voit d'autre part que $\mathcal{T}$ est irréductible si et seulement si $\operatorname{Esp}(X)$ (l'espace topologique associé, dont l'ensemble sous-jacent est $\operatorname{Ob} X$) est irréductible (\footnote[8]{car les ouverts de $\operatorname{Top}(X)$, ceux de $\operatorname{Esp}(X)$ et ceux de $X$ se correspondent.}) ; et on dit alors que $X$ est irréductible (\footnote[9]{Cela signifie aussi que pour $x$, $y$ dans $X$, $\exists\, z$ avec \begin{tikzcd}[ampersand replacement=\&, row sep=0.2em, column sep=1.2em, baseline=(\tikzcdmatrixname-2-1.base)] \& x \\ z \arrow[ur] \arrow[dr] \& \\ \& y \end{tikzcd}, i.e. $\operatorname{Ord}(X)$ est filtrant à gauche.}). Comme les composantes connexes de $X$ et celles de $\operatorname{Esp}(X)$ ou $\operatorname{Top}(X) = \mathcal{T}$ se correspondent, on voit donc que les conditions suivantes sont équivalentes : a) le topos $\operatorname{Top}(X)$ est localement irréductible, b)
\printedpage{8}
% head: A. Grothendieck : Dérivateurs, Ch. II, Édition des Universités de Montpellier II et Paris VII
% foot: Version 8 juillet 2001
$\operatorname{Esp}(X)$ est localement irréductible, c) les composantes connexes de $X$ sont irréductibles, d) les $X/x$ \texttt{[sont]} irréductibles. On dit alors que $X$ est localement irréductible (\footnote[10]{Faux, on a (a $\Longleftrightarrow$ d) $\Longrightarrow$ (b $\Longleftrightarrow$ c) $X$ irréductible n'implique \textit{pas} que les $X/x$ le soient, exemple $X = (0 \overset{\alpha}{\underset{\beta}{\rightrightarrows}} 1)$ $(\alpha \neq \beta)$, on a $X/1 =$ \begin{tikzcd}[ampersand replacement=\&, row sep=0.2em, column sep=1.2em, baseline=(\tikzcdmatrixname-2-2.base)] 0' \arrow[dr] \& \\ \& 1 \\ 0'' \arrow[ur] \& \end{tikzcd}, qui n'est \textit{pas} irréductible, ni localement irréductible puisqu'il est connexe. Pourtant il est étale sur $X$.}).

Cela implique donc que pour tout $X'$ étale sur $X$, avec $X'$ 0-connexe, $X'$ est irréductible.

\textit{Dire conditions équivalentes.}
\begin{enumerate}
\item[a)] $\operatorname{Top}(X)$ est localement irréductible.
\item[b)] Les $X/x$ \texttt{[sont]} irréductibles.
\item[c)] Tout objet 0-connexe de $X^\wedge$ est irréductible.
\item[d)] Les composantes connexes de $X$ sont faiblement cofiltrantes (i.e. faiblement filtrantes à gauche).
\item[e)] Condition PS $1_{\mathrm{g}}$ de SGA 4, I 2.
\end{enumerate}
Si $X$ est \texttt{[une]} catégorie, $X$ irréductible signifie aussi que $X \neq \emptyset$ et $\forall\, x, y \in \operatorname{Ob} X$, \texttt{[il]} existe \texttt{[un]} $s$ avec \begin{tikzcd}[row sep=0.2em, column sep=1.2em, baseline=(\tikzcdmatrixname-2-1.base)] & x \\ s \arrow[ur] \arrow[dr] & \\ & y \end{tikzcd}. (Cela signifie aussi que $X$ est \textit{faiblement filtrante} (\footnote[11]{Il est mieux l'appeler \textit{pseudo-filtrante,} et réservons faiblement filtrante à : localement faiblement filtrante + connexe.}) à gauche, ou faiblement cofiltrante.) $X$ localement irreductible signifie que pour $x$, $y$ \textit{appartenant à une même composante connexe,} \texttt{[il]} existe \texttt{[un]} $s$ avec \begin{tikzcd}[row sep=0.2em, column sep=1.2em, baseline=(\tikzcdmatrixname-2-1.base)] & x \\ s \arrow[ur] \arrow[dr] & \\ & y \end{tikzcd}, % [sic] irreductible -> irréductible, accent omis
\msp{page 13}
i.e. $X/[x] \cap X/[y] \neq \emptyset$ (où $[x]$ est considéré comme élément de l'ensemble ordonné $\operatorname{Ob} X$). Il revient encore au même de dire que les catégories induites $X/x$ sont irréductibles (\footnote[12]{faux, cf. annotation marginale plus haut.}).

\textbf{Exemple.} Une catégorie \textit{ayant un objet initiale,} ou seulement un `plus petit objet', est irréductible. Inversement, si l'ensemble \textit{ordonné} $\operatorname{Ord}(X)$ associé à $X$ (quotient de $\operatorname{Ob}(X)$ par la relation d'équivalence associé à son préordre) est fini, p.ex. $X$ fini, ou seulement localement fini, p.ex. $X$ localement fini (i.e. les catégories $X/x$ finies), alors si $X$ est irréductible, $\operatorname{Ob} X$ a un plus petit élément, i.e. $X$ a un plus petit objet. C'est vrai plus généralement encore, si $\operatorname{Ord}(X)$ est artinien, i.e. toute suite décroissante dans $\operatorname{Ord}(X)$ est stationnaire. % [sic] initiale -> initial, accord avec « objet »

Mais la condition intéressante est la condition topossique de \textit{locale} irréductibilité, qui équivaut : $I$ pseudo-filtrante à gauche, i.e. les composantes connexes de $I$ filtrantes, i.e. les conditions PS 1, PS 2

PS 1\enspace Tout diagramme \begin{tikzcd}[row sep=0.2em, column sep=1.2em, baseline=(\tikzcdmatrixname-2-2.base)] x \arrow[dr] & \\ & s \\ y \arrow[ur] & \end{tikzcd} se complète en \begin{tikzcd}[row sep=0.2em, column sep=1.2em, baseline=(\tikzcdmatrixname-2-1.base)] & x \arrow[dr] & \\ z \arrow[ur] \arrow[dr] & & s \\ & y \arrow[ur] & \end{tikzcd} \textit{commutatif.}

PS 2\enspace Pour toute double flèche $x \rightrightarrows y$, \texttt{[il]} existe un égalisateur $z \longrightarrow x \rightrightarrows y$.


\endgroup
% END UNIT Der-ChII-234/Der-ChII-234_p006-010.tex
% BEGIN UNIT Der-ChII-234/Der-ChII-234_p011-015.tex
\begingroup\makeatletter%
\newcommand{\printedpage}[1]{\marginpar{\footnotesize p.~#1}}%
\newcommand{\msp}[1]{\marginpar{\footnotesize [#1]}}%
\newcommand{\sX}{\mathcal{X}} % resolved against PDF: CMSY calligraphic X
\makeatother\let\printedpage\derpp\let\msp\dermsp%
\let\mspin\DERmspin%


\printedpage{9}
% head: A. Grothendieck : Dérivateurs, Ch. II, Édition des Universités de Montpellier II et Paris VII
% foot: Version 8 juillet 2001

\msp{page 14}

En résumé, on a la

\noindent\textbf{Proposition 2.1.} \textit{Soit $X$ un objet de} Cat. \textit{Conditions équivalentes :}
\begin{enumerate}[label={},leftmargin=3em]
\item[\textit{a)}] \textit{Le topos associé} $\operatorname{Top}(X)$ \textit{est localement irréductible, i.e. (étant localement connexe) pour tout objet $F$ de $X^{\wedge}$ et tout sous-objet $F'$ de $F$, si $F$ est connexe, $F'$ l'est.}
\item[\textit{a$'$)}] \textit{Pour tout $X'$ étale sur $X$, si $X'$ est connexe, tout ouvert de $X'$ est connexe.}
\item[\textit{a$''$)}] \textit{Tout ouvert d'une catégorie localisée $X/a$ ($a \in \operatorname{Ob} X$) est connexe (i.e. $X/a$ est une catégorie irréductible, i.e. le topos} $\operatorname{Top}(X/a)$ \textit{l'est).}
\item[\textit{b)}] \textit{Pour tout objet $a$ de $X$,} $\operatorname{Esp}(X/a)$ \textit{est un espace topologique irréductible.}
\item[\textit{c)}] \textit{Les catégories $X/a$ sont faiblement filtrantes à gauche (ou faiblement cofiltrantes), i.e. pour deux objets $\xi$, $\eta$ d'une telle catégorie,} \texttt{[il]} \textit{existe un objet $\zeta$ et deux flèches}
\begin{tikzcd}[ampersand replacement=\&, row sep=tiny, column sep=small, baseline=(\tikzcdmatrixname.center)]
 \& \xi \\
\zeta \arrow[ur] \arrow[dr] \& \\
 \& \eta
\end{tikzcd}
\textit{(i.e. l'ensemble ordonné Ord($X/a$) associé à $X/a$ est filtrant à gauche)}
\item[\textit{c$'$)}] \textit{Pour tout diagramme en traits pleins dans $X$,}
\begin{tikzcd}[ampersand replacement=\&, row sep=tiny, column sep=small, baseline=(\tikzcdmatrixname.center)]
 \& x \arrow[dr] \& \\
z \arrow[ur, dotted] \arrow[dr, dotted] \& \& s \\
 \& y \arrow[ur] \&
\end{tikzcd}
\textit{, on peut le compléter en un carré} commutatif \textit{comme indiqué. (Condition} PS 1 \textit{des catégories localement filtrantes.)}
\end{enumerate}

\noindent Ces conditions impliquent que toute composante connexe $X_{\alpha}$ de $X$ est irréductible, ou ce qui revient au même, que $\operatorname{Esp}(X_{\alpha})$ est irréductible,

\msp{page 15}

\noindent mais ne sont pas impliquées par celle-ci.

Contrexemple, la catégorie 0-connexe $0 \underset{\beta}{\overset{\alpha}{\rightrightarrows}} 1$, avec $\alpha \neq \beta$. On a $X/1 \simeq$
\begin{tikzcd}[ampersand replacement=\&, row sep=tiny, column sep=small, baseline=(\tikzcdmatrixname.center)]
0 \arrow[dr] \& \\
 \& 1 \\
0' \arrow[ur] \&
\end{tikzcd}, qui n'est \textit{pas} irréductible, puisque $\Phi \backslash \{1\} \simeq \varepsilon$ est un ouvert disconnexe \texttt{[}$\Phi = X/1$\texttt{,} $\varepsilon = \bullet\ \bullet$\texttt{]}.

J'aimerais une condition `topossique' sur $X$, équivalente à $X$ `pseudo-filtrante à gauche', i.e. les composantes connexes filtrantes, i.e. PS 1\textsuperscript{opp} (condition c$'$) ci-dessus) + PS 2\textsuperscript{opp} (condition de la double flèche), et qui de plus implique que tout objet 0-connexe de $X^{\wedge}$ (i.e. tout objet $X'$ de Cat étale sur $X$ et 0-connexe) soit $W_{\infty}$-asphérique. Et je songe à la notion de `totale asphéricité' de Pursuing Stacks.

\noindent\textbf{Définition 2.2.} (\footnote[13]{Cf. tableau récapitulatif pages 97, 98 et surtout 99.})\ \ Soit $X$ une catégorie.
\begin{enumerate}[label={},leftmargin=3em]
\item[a)] On dit que $X$ est \textit{faiblement filtrante} (\footnote[14]{dire plutôt \textit{pseudo-filtrante}, réserver `faiblement filtrante' à `localement faiblement filtrante' (ou \textit{irréductible}) + connexe, condition PS $1'_{\mathrm{g}}$.}) si $X \neq \emptyset$, et si pour $a, b \in \operatorname{Ob} X$, $\exists\, c$ et des flèches $c$
\begin{tikzcd}[ampersand replacement=\&, row sep=tiny, column sep=small, baseline=(\tikzcdmatrixname.center)]
 \& a \\
{} \arrow[ur, phantom] \& \\
 \& b
\end{tikzcd}
\hspace{-2.6em}% layout aid only
\begin{tikzcd}[ampersand replacement=\&, row sep=tiny, column sep=small, baseline=(\tikzcdmatrixname.center)]
 \& {} \\
{} \arrow[ur] \arrow[dr] \& \\
 \& {}
\end{tikzcd}
, i.e. $X/[a] \cap X/[b] \neq \emptyset$. (\textbf{NB} $X$ est faiblement filtrante (\footnote[15]{pseudo-filtrante.}) à gauche si et seulement si $\operatorname{Ord}(X)$ l'est, i.e. est filtrante à gauche. \textbf{ReNB} $X$ faiblement filtrante $\Longrightarrow$ $X$ 0-connexe.)

\item[b)] \printedpage{10}
% head: A. Grothendieck : Dérivateurs, Ch. II, Édition des Universités de Montpellier II et Paris VII
% foot: Version 8 juillet 2001
On dit que $X$ est \textit{localement faiblement filtrante à gauche} ou plus simplement préfiltrante à gauche si les catégories $X/x$ ($x \in \operatorname{Ob} X$) sont faiblement filtrantes à gauche, i.e. si $X$ satisfait la condition : $\forall$ diagrammes
\begin{tikzcd}[ampersand replacement=\&, row sep=tiny, column sep=small, baseline=(\tikzcdmatrixname.center)]
a \arrow[dr] \& \\
 \& x \\
b \arrow[ur] \&
\end{tikzcd}, $\exists\, c$ dans $X$ et des flèches de $c$ dans $a$, $b$ rendant \textit{commutatif} le carré
\[
\begin{tikzcd}[ampersand replacement=\&, row sep=tiny, column sep=small]
 \& a \arrow[dr] \& \\
c \arrow[ur, dotted] \arrow[dr, dotted] \& \& x \\
 \& b \arrow[ur] \&
\end{tikzcd}
\]
(\footnote[16]{Condition PS $1_{\mathrm{g}}$.}). (\textbf{NB} Cette condition implique que $\operatorname{Ord}(X)$ est localement faiblement filtrant, mais n'équivaut pas à cette dernière. Exemple, $X = 0 \underset{\beta}{\overset{\alpha}{\rightrightarrows}} 1$, $X/1 =$
\begin{tikzcd}[ampersand replacement=\&, row sep=tiny, column sep=small, baseline=(\tikzcdmatrixname.center)]
0' \arrow[dr] \& \\
 \& 1 \\
0'' \arrow[ur] \&
\end{tikzcd}
n'est pas faiblement filtrante à gauche.)

\msp{page 16}

\textbf{NB} Si $X$ localement faiblement filtrante à gauche, tout $X'$ étale sur $X$ itou.

\textbf{NB} (\footnote[17]{\textbf{Faux.}}) $X$ localement faiblement filtrante à gauche $\Longleftrightarrow$ les composantes connexes sont faiblement filtrantes à gauche, ou encore si $X$ est 0-connexe, $X$ localement faiblement filtrante à gauche $\Longleftrightarrow$ $X$ est faiblement filtrante (\footnote[18]{pseudo-filtrante.}) à gauche. On a en effet $\Longleftarrow$ tautologiquement (\footnote[19]{\textbf{faux,} ainsi $0 \underset{\beta}{\overset{\alpha}{\rightrightarrows}} 1$ est 0-connexe et faiblement filtrante à gauche, mais non pas localement faiblement filtrante à gauche car
\begin{tikzcd}[ampersand replacement=\&, row sep=tiny, column sep=small, baseline=(\tikzcdmatrixname.center)]
0 \arrow[dr, "\alpha"] \& \\
 \& 1 \\
0 \arrow[ur, "\beta"'] \&
\end{tikzcd}
ne se complète pas en un carré \textit{commutatif.}}), et $\Longrightarrow$ en montrant que si PS 1 \texttt{[est]} satisfait, alors $PS(a,b)$ défini par $X/[a] \cap X/[b] \neq \emptyset$ est une relation d'équivalence, impliquée par la relation $C(a,b)$ : $\{\operatorname{Hom}(a,b) \neq \emptyset$ \textit{ou} $\operatorname{Hom}(b,a) \neq \emptyset\}$, donc la relation d'équivalence (homotopie) engendrée par $C(a,b)$ est plus fine que $PS(a,b)$.

\item[c)] On dit que $X$ est \textit{filtrante à gauche}, ou \textit{cofiltrante,} si elle est faiblement filtrante (\footnote[20]{pseudo-filtrante.}) à gauche (donc \textit{non vide}), et satisfait (en plus de la condition PS $1'_{\mathrm{g}}$ de a)) la condition PS $2_{\mathrm{g}}$ suivante :

Pour toute double flèche $a \underset{v}{\overset{u}{\rightrightarrows}} b$, \texttt{[il]} existe \texttt{[une]} flèche $c \xrightarrow{w} a$ telle que $uw = vw$ (\footnote[21]{i.e. trois conditions a) $X \neq \emptyset$ b) PS $1'$ habituelle c) PS 2 (double flèche).}).

Il revient au même de dire que $X$ satisfait les trois conditions a$'$) $X$ est 0-connexe (plus forte que $X \neq \emptyset$), b$'$) PS 1, i.e. $X$ localement faiblement filtrante à gauche et c$'$) PS 2 comme avant.

\item[d)] On dit que $X$ est \textit{localement filtrante à gauche} si les composantes connexes sont filtrantes à gauche, i.e. si $X$ satisfait les conditions PS 1 et PS 2 (`pseudo-cofiltrante' dans la terminologie de SGA 4). \textbf{NB} Si $X$ est localement filtrante à gauche, il en est de même de tout $X'$ étale sur $X$, car c'est évident pour PS 1 qui concerne les $X/x$, et pour PS 2 cela résulte du fait que $X' \longrightarrow X$ est fidèle. Donc si $X$ est localement cofiltrante, les $X/a$ sont filtrantes, mais la réciproque est fausse.
\end{enumerate}

\printedpage{11}
% head: A. Grothendieck : Dérivateurs, Ch. II, Édition des Universités de Montpellier II et Paris VII
% foot: Version 8 juillet 2001

\noindent On a les notions duales : \textit{faiblement filtrante à droite} (ou simplement faiblement filtrante), et \textit{filtrante à droite} ou simplement filtrante. (Notions duales de a), c) ci-dessus.) Pour la notion duale de b), savoir que les $x\backslash X$ soient faiblement filtrantes, on ne peut plus l'appeler `localement faiblement filtrante', mais ce serait \textit{co}localement faiblement filtrante.

\noindent\textbf{Exemple.} $X = \Bigl(0 \xrightarrow{\gamma} 1 \underset{\beta'}{\overset{\alpha'}{\rightrightarrows}} 2 \underset{\beta}{\overset{\alpha}{\rightrightarrows}} 3\Bigr)$ avec $(\alpha\alpha' = \beta\beta') \neq (\alpha\beta' = \beta\alpha')$.

\msp{page 17}

\noindent Seule la notion a) (faiblement filtrante (\footnote[22]{pseudo-filtrant.}) à gauche) ne dépend que de $\operatorname{Ord} X$ - mais bien sûr si $X$ est \textit{pré}ordonnée, toutes trois ne dépendent que de $\operatorname{Ord} X$. Notons que dire que ça ne dépend que de $\operatorname{Ord} X$, ou que de $\mathrm{Esp}\,\mathrm{sep}X$, revient au même, puisque les deux se déterminent mutuellement. La condition a) signifie que $\operatorname{Esp} X$, ou aussi $\operatorname{Top}(X)$, est \textit{irréductible.} On dira que $X$ est \textit{irréductible} (au lieu de pseudo-filtrante). Les trois autres notions ne peuvent s'exprimer en termes de $\operatorname{Esp} X$ ou de $\operatorname{Ord} X$ seul, mais elles s'expriment en termes du $\mathrm{Topos}\, X$. La condition b) ($X$ localement faiblement filtrante) signifie que $\operatorname{Top}(X)$ est \textit{localement irréductible,} et on dira aussi que $X$ est \textit{localement irréductible} (au lieu de localement faiblement filtrante). Ça signifie que pour tout $X'$ étale sur $X$ et 0-connexe, $X'$ est faiblement filtrante à gauche, i.e. est irréductible.

Reste à exprimer la condition filtrante elle-même en termes de $\operatorname{Top} X$. Je rappelle à ce propos :

\noindent\textbf{Proposition 2.3.} \textit{Les conditions suivantes sur un topos $\sX$ sont équivalentes (quand $S$ (comme `site') désigne une partie génératrice de} $\mathrm{Fais}(\sX) = S^{\wedge}$, formée d'objets 0-connexes (\footnote[23]{donc $\sX$ est localement 0-connexe.}).\textit{)}
\begin{enumerate}[label={},leftmargin=3em]
\item[\textit{a)}] \textit{Si $U$, $V$ sont deux objets $0$-connexes de $\operatorname{Ob} \sX$, $U \times V$ est $0$-connexe.}
\item[\textit{b)}] \textit{Si $U$ est un objet $0$-connexe de $\sX^{\wedge}$, alors $U \longrightarrow e$ est `$W_0$-asphérique', i.e. pour tout $V$, la projection $V \times U \longrightarrow V$ induit $\pi_0(U \times V) \xrightarrow{\sim} \pi_0(U)$.}
\item[\textit{c)}] \textit{Si $U$, $V$} \texttt{[sont]} \textit{dans $\sX^{\wedge}$, alors}
\[
\pi_0(U \times V) \xrightarrow{\sim} \pi_0(U) \times \pi_0(V).
\]
\item[\textit{d)}] \textit{Comme a), avec $U$, $V$ dans $S$. (Mais on suppose les objets de $S$ $0$-connexes.)}
\end{enumerate}
\noindent(\footnote[24]{Il faudrait expliciter la démonstation avec celle de 2.9 ci-dessous (p.~25).}). % [sic] démonstation -> démonstration, letter r missing

\msp{page 18}

\noindent De plus, ces conditions impliquent que $\sX$ est connexe, donc 0-connexe si $\sX \neq \emptyset$. Cela implique de plus que $\sX$ est $W_{\infty}$-asphérique, cf. Pursuing Stacks §35 (Corollary p.~52) (\footnote[25]{\textbf{Faux,} loc. cit. ne dit rien de tel.}).

C'est essentiellement Proposition 1 de Pursuing Stacks §65, cf. loc. cit. p.~170-171 (\footnote[26]{Mais je m'y borne aux topos de la forme $\operatorname{Top}(A)$, avec $A$ dans Cat.}). Cf. aussi Pursuing Stacks, §35, pages 50-55.

\printedpage{12}
% head: A. Grothendieck : Dérivateurs, Ch. II, Édition des Universités de Montpellier II et Paris VII
% foot: Version 8 juillet 2001

\noindent\textbf{Définition 2.4.} On dit que $\sX$ est \textit{totalement $0$-connexe} s'il est localement connexe, non vide et satisfait les conditions équivalentes de 2.3. (NB c'est un cas particulier de la notion `totalement $W$-asphérique' de loc. cit. On a ainsi la notion de totalement $n$-connexe.)

Notons que pour un objet $X$ de Cat, dire que $\operatorname{Top}(X)$ est totalement 0-connexe revient à dire que $X \neq \emptyset$, et que pour $x, y \in \operatorname{Ob} X$, $x \times y$ dans $X^{\wedge}$ est 0-connexe, i.e. $X/x \times y$ est 0-connexe (\footnote[27]{Cela signifie aussi que $\mathrm{diag}_X : X \longrightarrow X \times X$ est co-cofinale, i.e. $X^o \longrightarrow X^o \times X^o$ est cofinale, et implique qu'il en est de même de $X \longrightarrow X^E$ pour tout ensemble \textit{fini} $E$.}). C'est une condition très forte. On dira aussi que $X$ est \textit{totalement $0$-connexe.} Notons ici

\noindent\textbf{Proposition 2.5.} \textit{Soit $\sX$ un topos. Si $\sX$ est totalement $0$-connexe, alors $\sX$ est irréductible. (En fait la conclusion vaut, tautologiquement, en supposant seulement que le produit de deux objets $0$-connexes de $\sX^{\wedge}$ est non vide, et il suffit même de le supposer pour des sous-objets de l'objet final $e_{\sX}$.) En particulier, si $\sX = \mathit{Top}(X)$, où $X$ est dans} Cat\textit{, alors $X$ est faiblement filtrante à gauche.}

Mais bien sûr, cela n'implique \textit{pas} que $\sX$ soit localement irréductible.

\msp{page 19}

\noindent Par exemple la catégorie $\boldsymbol{\Delta}$, ou toute petite catégorie stable par produits binaires, est totalement 0-connexe, mais $\boldsymbol{\Delta}$ n'est pas localement irréductible (les catégories induites $\boldsymbol{\Delta}/\boldsymbol{\Delta}^n$, pour $n > 0$, ne sont visiblement \textit{pas} irréductibles). (Mais rappelons que si $\sX = \operatorname{Top}(T)$, $T$ un espace topologique, alors $T$ irréductible implique que $\sX$ est aussi \textit{localement} irréductible. Donc $T$ totalement 0-connexe implique $\operatorname{Top}(T)$ globalement \textit{et} localement irréductible, i.e. pour tout $T'$ étale sur $T$, et 0-connexe, $T'$ est irréductible.)

\noindent\textbf{Théorème 2.6.} \textit{Soit $\sX$ un topos localement connexe, $S \subseteq \sX^{\wedge}$ une petite partie génératrice formée d'objets $0$-connexes. Conditions équivalentes (}sauf \textit{c, c$'$, qui semblent} (\footnote[28]{\textbf{sont.}}) \textit{plus faibles, cf. démonstration).)}
\begin{enumerate}[label={},leftmargin=3em]
\item[\textit{a)}] \textit{Tout topos induit $\sX/U$ sur $U$ $0$-connexe dans $\sX^{\wedge}$ est totalement $0$-connexe.}
\item[\textit{a$'$)}] \textit{Tout produit fibré $U \times_W V$ d'objets $0$-connexes dans $\sX^{\wedge}$ est $0$-connexe.}
\item[\textit{b)}] \textit{$\sX$ est localement irréductible, et ses composantes connexes sont totalement $0$-connexes.}
\item[\textit{c)}] \textit{Comme a), avec $U$ dans $S$, i.e. les $\sX/U$ ($U \in \operatorname{Ob} S$) sont totalement $0$-connexes.}
\item[\textit{c$'$)}] \textit{Comme a$'$), avec $U$, $V$, $W$ dans $S$.}
\item[\textit{d)}] \textit{$S$ est localement filtrante à gauche, i.e. $S$ satisfait les deux conditions}

PS$^*$ 1\ \ \textit{Tout diagramme}
\begin{tikzcd}[ampersand replacement=\&, row sep=tiny, column sep=small, baseline=(\tikzcdmatrixname.center)]
a \arrow[dr] \& \\
 \& x \\
b \arrow[ur] \&
\end{tikzcd}
\textit{dans $S$ s'insère dans}
\begin{tikzcd}[ampersand replacement=\&, row sep=tiny, column sep=small, baseline=(\tikzcdmatrixname.center)]
 \& a \arrow[dr] \& \\
c \arrow[ur, dotted] \arrow[dr, dotted] \& \& x \\
 \& b \arrow[ur] \&
\end{tikzcd}
\textit{commutatif.}

PS$^*$ 2\ \ \textit{Tout diagramme $a \rightrightarrows b$ s'insère dans un diagramme commutatif $c \longrightarrow a \rightrightarrows b$.}
\end{enumerate}
\noindent(\footnote[29]{NB On a équivalence de \tikz[baseline=(n.base)]\node[draw, rounded corners, inner sep=1.5pt](n){a), a$'$), b), d),}; qui impliquent c), c$'$), lesquels sont équivalents. Mais il semble (\footnotemark[30]) que c), c$'$) soient strictement plus faible. Cf. cependant les variantes e), e$'$) dans 2.12 (p.~28, 29).}). % [sic] plus faible -> plus faibles, adjective does not agree with the plural subject c), c')
\footnotetext[30]{Il s'avère. C'est bien vrai, c) et c$'$) n'impliquent \textit{pas} les autres conditions, cf. annotation marginale
% footnote 30 continues at the foot of printed page 13:
p.~28 (cas de $X = B_G$).}

\printedpage{13}
% head: A. Grothendieck : Dérivateurs, Ch. II, Édition des Universités de Montpellier II et Paris VII
% foot: Version 8 juillet 2001

\msp{page 20}

\noindent Ces conditions impliquent que les objets 0-connexes de $\sX^{\wedge}$ (en particulier les éléments de $S$) sont $\infty$-connexes.

\noindent\textbf{Corollaire 2.7} (de l'équivalence b $\Longleftrightarrow$ d). \textit{Soit $X$ un objet de} Cat. \textit{Pour que $X$ soit filtrant à gauche il faut et il suffit que $\mathit{Top}(X)$ soit totalement $0$-connexe et localement (et donc aussi globalement) irréductible (ou encore que $X$ soit totalement $0$-connexe et que pour tout $x \in \operatorname{Ob} X$, la catégorie localisée $X/x$ soit irréductible (elle est même totalement $0$-connexe), ou encore que pour deux objets $a$, $b$ dans $X$, $a \times b \in X^{\wedge}$ soit $0$-connexe, et pour}
\begin{tikzcd}[ampersand replacement=\&, row sep=tiny, column sep=small, baseline=(\tikzcdmatrixname.center)]
a \arrow[dr] \& \\
 \& x \\
b \arrow[ur] \&
\end{tikzcd}
\textit{dans $X$, $a \times_x b$ (dans $X^{\wedge}$) soit $0$-connexe)} (\footnote[31]{Exemple de $X$ totalement 0-connexe, mais non cofiltrant, i.e. non localement irréductible ? La catégorie $\boldsymbol{\Delta}$ des simplexes type est totalement 0-connexe, mais n'est pas cofiltrante car $\boldsymbol{\Delta}^0 = e \underset{\partial_1}{\overset{\partial_0}{\rightrightarrows}} \boldsymbol{\Delta}^1$ ne peut être égalisé ! Donc \texttt{[la catégorie opposée]} $\boldsymbol{\Delta}^o = I$ est telle que $I \longrightarrow I^E$ soit cofinal pour tout ensemble fini $E$, mais $I$ non filtrante.}).

\noindent\textsc{Démonstration} de 2.6. On a le diagramme d'implications $\pm$ tautologiques
\[
\begin{tikzcd}[ampersand replacement=\&, column sep=large, row sep=large]
 \& \mathrm{a} \arrow[r, Leftrightarrow] \arrow[dl, Rightarrow, "{\substack{\text{par}\\ \text{2.5}}}"'] \arrow[d, Rightarrow] \& \mathrm{a}' \arrow[d, Rightarrow] \\
\mathrm{b} \arrow[d, Rightarrow] \& \mathrm{c} \arrow[r, Leftrightarrow, "{(\Leftarrow \text{ par 2.3.})}"] \& \mathrm{c}' \arrow[dll, Rightarrow] \\
\text{PS 1.} \& \&
\end{tikzcd}
\]
Montrons quand même les deux implications c$'$) $\Longrightarrow$ PS 1 et b $\Longrightarrow$ PS 1.

\msp{page 21}

\noindent C'est contenu dans le

\noindent\textbf{Corollaire 2.8.} \textit{Soient $\sX$, $S$ comme dans 2.6. Conditions équivalentes :}
\begin{enumerate}[label={},leftmargin=3em]
\item[\textit{a)}] \textit{$\sX$ est localement irréductible, i.e. les topos induits $\sX/U$, où $U$ est $0$-connexe, sont irréductibles.}
\item[\textit{b)}] \textit{Tout produit fibré $U \times_W V$ d'objets $0$-connexes $U$, $V$ est `non vide'.}
\item[\textit{b$'$)}] \textit{Tout produit fibré $U \times_W V$, avec $U$, $V$ `non vide', $W$ $0$-connexe, est non vide.}
\item[\textit{c)}] \textit{Les $\sX/x$ avec $x \in S$ sont irréductibles.}
\item[\textit{d)}] \textit{Comme b), avec $U$, $V$, $W$ dans $S$.}
\item[\textit{e)}] \textit{$S$ est localement faiblement filtrante à gauche, i.e. satisfait PS$^*$ 1.}
\end{enumerate}
\noindent Ici les implications $\pm$ évidentes sont


\endgroup
% END UNIT Der-ChII-234/Der-ChII-234_p011-015.tex
% BEGIN UNIT Der-ChII-234/Der-ChII-234_p016-020.tex
\begingroup\makeatletter%
\newcommand{\printedpage}[1]{\marginpar{\footnotesize p.~#1}}%
\newcommand{\msp}[1]{\marginpar{\footnotesize [#1]}}%
\newcommand{\X}{\mathcal{X}}%
\makeatother\let\printedpage\derpp\let\msp\dermsp%
\let\mspin\DERmspin%


\printedpage{14}
% head: A. Grothendieck : Dérivateurs, Ch. II, Édition des Universités de Montpellier II et Paris VII
% foot: Version 8 juillet 2001

\[
\begin{tikzcd}
\mathrm{b}' \arrow[r, Rightarrow] & \mathrm{b} \arrow[rr, Rightarrow] \arrow[d, Rightarrow] & & \mathrm{a} \arrow[d, Rightarrow] \\
 & \mathrm{d} \arrow[r, Leftrightarrow] & \mathrm{e} \arrow[r, Leftrightarrow] & \mathrm{c}.
\end{tikzcd}
\]

Prouvons les implications de la ligne du bas. Considérons d). Si $\begin{tikzcd}[cramped, row sep=0.1em, column sep=0.8em, baseline=(current bounding box.center)] a \arrow[dr] & \\ & x \\ b \arrow[ur] & \end{tikzcd}$ \texttt{[est]} dans $S$, dire que le produit fibré dans $\X$ est `non vide' signifie, puisque $S$ est un système de générateurs et en particulier \textit{non vide,} qu'il existe $c$ dans $S$ et $c\longrightarrow a\times_x b$, ou ce qui revient au même, $\begin{tikzcd}[cramped, row sep=0.1em, column sep=0.8em, baseline=(current bounding box.center)] & a \\ c \arrow[ur, dotted] \arrow[dr, dotted] & \\ & b \end{tikzcd}$, rendant commutatif le diagramme $\begin{tikzcd}[cramped, row sep=0.1em, column sep=0.8em, baseline=(current bounding box.center)] & a \arrow[dr] & \\ c \arrow[ur, dotted] \arrow[dr, dotted] & & x \\ & b \arrow[ur] & \end{tikzcd}$. Donc d) équivaut à PS 1, i.e. à e) ($\X$ localement faiblement filtrant à gauche). Considérons la condition c) : les $\X/a=(S/a)^\sim$ sont irréductibles. Cela signifie que pour deux sous-objets \textit{non vides} $U,V\subseteq a$ dans $\X^\wedge$, on a $U\cap V\neq\emptyset$, ou encore que pour deux cribles

\msp{page 22}$\underline{U}$, $\underline{V}$ de $S/x$, si $\underline{U}$, $\underline{V}$ non vides (ce qui signifie que les sous-objets de $x$ dans $\X^\wedge$ qu'ils définissent sont `non vides', ou que les objets de $S$ sont `non vides' dans $\X^\wedge$), alors $\underline{U}\cap\underline{V}$ est non vide. On peut supposer que $\underline{U}$, $\underline{V}$ est engendré chacun par une seule flèche $a\xrightarrow{u}x$, $b\xrightarrow{v}x$. Le crible intersection est formé des objets $c\longrightarrow x$ de $S/x$ tels que $c\longrightarrow x$ se factorise par un $\begin{tikzcd}[cramped, row sep=0.1em, column sep=0.8em, baseline=(current bounding box.center)] & a \\ c \arrow[ur, dotted] \arrow[dr, dotted] & \\ & b \end{tikzcd}$, donc dire que ce crible n'est pas vide, signifie que PS 1 est satisfait pour les deux flèches données $u$, $v$. Donc c) équivaut à la condition PS 1\texttt{[, i.e. à]} e.

Ainsi on a établi
\[
\mathrm{b}'\Longrightarrow\mathrm{b}\Longrightarrow\mathrm{a}\Longrightarrow(\mathrm{c}\Longleftrightarrow\mathrm{d}\Longleftrightarrow\mathrm{e}),
\]
reste à prouver que c $\Longrightarrow$ b$'$. Soit donc $\begin{tikzcd}[cramped, row sep=0.1em, column sep=0.8em, baseline=(current bounding box.center)] U \arrow[dr] & & V \arrow[dl] \\ & W & \end{tikzcd}$ un diagramme d'objets `non vides', avec $W$ 0-connexe, il faut prouver que le produit fibré est non vide. Considérant $a\longrightarrow U$ et $b\longrightarrow V$ avec $a,b\in S$, on est ramené au cas où $U=a$, $V=b$ dans $S$. Prouvons que si $a\times_W b$ est `vide', alors $W$ est disconnexe. Pour ceci, considérons la sous-catégorie $C$ de $S/W$ formée des $x\longrightarrow W$ tels que $a\times_W x$ soit `non vide'. Désignant par $w(a)$, $w(x)$ l'image de $a$, $x$ dans $W$, cela signifie que $w(a)\cap w(x)\neq\emptyset$.

\msp{page 23}Notons que tout quotient dans $\X^\wedge$ d'un objet irréductible est irréductible, donc $w(a)$, $w(x)$ sont des objets irréductibles. Cela montre que pour $x$ dans $C$ et $x'$ dans $S/W$, on a $x'$ dans $C$, i.e. $w(a)\cap w(x')\neq\emptyset$ si et seulement si $w(x)\cap w(x')\neq\emptyset$. (Lemme direct : dans un objet $W$ d'un topos $\X$, si on a \texttt{[des]} sous-objets $W_1$, $W_2$, $W_3$ irréductibles, si l'un des deux \texttt{[plutôt trois]} rencontre les deux autres, ceux-ci se rencontrent.) Cela montre que $C$ est un crible, puisque $x'\longrightarrow x$ implique $w(x')\subseteq w(x)$, donc $w(x')$ rencontre $w(x)$, donc aussi $w(a)$. Mais alors le co-crible $C'=(S/W)\backslash C$ est également un crible, car si $c$
\printedpage{15}
% head: A. Grothendieck : Dérivateurs, Ch. II, Édition des Universités de Montpellier II et Paris VII
% foot: Version 8 juillet 2001
est dans $C'$, i.e. \texttt{[si]} $w(c)$ ne rencontre pas $w(x)$, a fortiori pas tout $x'$ au-dessus de $x$, pour lequel on aura $w(x')\subseteq w(x)$. Ainsi, $S/W$ est décomposé en somme de deux ouverts $C$, $C'$, non vides puisque l'un contient $a$, l'autre $b$. Par suite, $W$ est décomposé en somme de deux sous-objets non vides, donc $W$ est non 0-connexe.

On a prouvé 2.8, donc aussi (revenant à 2.6) les implications vers PS 1 dans le diagramme page 20. Il reste à prouver
\[
\left\{
\begin{aligned}
&\mathrm{b}\Longrightarrow\mathrm{PS}\ 2,\ \mathrm{c}'\Longrightarrow\mathrm{PS}\ 2\\
&\mathrm{PS}\ 1+\mathrm{PS}\ 2\ (\text{i.e. d})\Longrightarrow\mathrm{a}'
\end{aligned}
\right.
\]
(\footnote[32]{NB finalement on ne prouve pas c$'\Longrightarrow$ PS 2.}).

\msp{page 24}Nous nous plaçons donc dans l'hypothèse où PS 1 est satisfait (puisque b), c$'$) et d) impliquent PS 1), donc où on est sous les conditions équivalentes de 2.8 - $\X$ est localement irréductible. Considérons alors la condition PS 2, et prouvons qu'elle est impliquée tant par b), que par c$'$), et qu'elle implique a$'$). Elle signifie que pour une double flèche $a\overset{u}{\underset{v}{\rightrightarrows}}b$ dans $S$, le conoyau \texttt{[plutôt noyau]} dans $S^\wedge$, ou, ce qui revient au même, dans $\X^\wedge=S^\sim\subseteq S^\wedge$, est `non vide'. On peut l'écrire sous la forme d'un produit fibré
\[
\begin{tikzcd}
a \arrow[r, "{\scriptstyle (u,v)}"] & b\times b \\
\operatorname{Ker}(u,v) \arrow[u] \arrow[r] & b. \arrow[u, "\mathrm{diag}"']
\end{tikzcd}
\]
Ici par hypothèse $a$ et $b$ sont des objets 0-connexes de $\X^\wedge$ (ceux de $S$ l'étant), c) est vérifié (($\X$ localement irréductible et) ses composantes connexes \texttt{[sont]} totalement 0-connexes), alors $b\times b$ est 0-connexe comme produit de deux objets 0-connexes sur une \textit{même composante connexe,} laquelle est totalement 0-connexe. Comme $\X\times\X$ est localement irréductible, il s'ensuit que le produit fibré est non vide (2.8), donc b) $\Longrightarrow$ d). Prouvons d) $\Longrightarrow$ a), i.e. si

\msp{page 25}$S$ est localement cofiltrante, alors les topos induits $\X/W$, où $W$ est 0-connexe, sont totalement 0-connexes. Ce topos induit est décrit par le site $S/W$, i.e. $S/W$ est un système de générateurs. Or si $S$ est localement filtrante, $S/W$ (étale sur $S$) l'est aussi. Donc on peut supposer $W=e$, et à prouver que si $S$ est localement cofiltrante et $\X$ 0-connexe, alors $\X$ est totalement 0-connexe. Notons que $\X$ 0-connexe équivaut à $S$ 0-connexe - il faudrait quand même l'expliciter :

\textbf{Lemme 2.9.} \textit{Soit $\X$ un topos, $S\subseteq\X^\wedge$ une sous-catégorie pleine génératrice formée d'objets} 0-connexes. \textit{Alors $\X$ est $0$-connexe si et seulement si $S$ l'est - plus généralement, l'ensemble ordonné des sous-objets à la fois ouverts et fermés de $S$ dans} Cat \textit{(ou encore des
\printedpage{16}
% head: A. Grothendieck : Dérivateurs, Ch. II, Édition des Universités de Montpellier II et Paris VII
% foot: Version 8 juillet 2001
sous-objets ouverts $=$ cribles de $S$ qui sont aussi fermés, i.e. des cocribles) est isomorphe à celui des sommandes directs de $\X$.}

Je l'ai utilisé tacitement dans la démonstration de 2.8, pour le topos induit $\X/W$. Je n'explicite pas la démonstration.

Ainsi (revenant à la démonstration de d) $\Longrightarrow$ c) dans 2.6.) $S$ étant 0-connexe et localement cofiltrante,

\msp{page 26}\texttt{[elle]} est en fait cofiltrante. Donc pour deux objets $a$, $b$ de $S$, $\exists\ c$ avec $c\longrightarrow a$, $c\longrightarrow b$ (condition PS 1$'$, plus forte que PS 1). Pour prouver que $\X$ est totalement 0-connexe, on est réduit à prouver par 2.3 que les $a\times b$ dans $\X^\wedge$ sont 0-connexes. Par 2.9, il revient au même de prouver qu'ils sont 0-connexes dans $S^\wedge$. Donc on est ramené à prouver le

\textbf{Lemme 2.10.} \textit{Soit $S$ une catégorie cofiltrante. Si $a$, $b$ sont deux objets de $S$, $a\times b$ est un objet} $0$\textit{-connexe de $S^\wedge$.}

On voit de suite qu'il suffit de prouver (et en fait il revient au même, compte tenu que $a\times b$ est localement irréductible) que si on a deux flèches
\[
\begin{tikzcd}
x \arrow[r] & a\times b \\
 & y, \arrow[u]
\end{tikzcd}
\]
alors $x\times_{a\times b}y\neq\emptyset$, i.e. $\exists\ z$ dans $S$ rendant commutatif
\[
\begin{tikzcd}
x \arrow[r] & a\times b \\
z \arrow[u] \arrow[r] & y. \arrow[u]
\end{tikzcd}
\]
Or par PS 1$'$ il existe bien un $z$ s'envoyant dans $x$ et dans $y$, mais les deux flèches composées
\[
z\rightrightarrows a\times b,
\]
i.e. chacune des double flèches
\[
z\rightrightarrows a \tag*{$(*)$}
\]
\[
z\rightrightarrows b \tag*{$(*')$}
\]
peuvent n'être pas égales. Mais par PS 2,

\msp{page 27}on peut trouver
\[
z'\longrightarrow z
\]
\printedpage{17}
% head: A. Grothendieck : Dérivateurs, Ch. II, Édition des Universités de Montpellier II et Paris VII
% foot: Version 8 juillet 2001
égalisant la double flèche $(*)$, donc on peut supposer que celle-ci est déjà triviale. Appliquant encore une fois PS 2, par un $z''\longrightarrow z'$, on égalise de même la double flèche $(*')$, et on gagne. On a donc prouvé, finalement,
\[
\begin{tikzcd}
\mathrm{a} \arrow[r, Leftrightarrow] \arrow[d, Rightarrow] & \mathrm{a}' \arrow[r, Rightarrow] & \mathrm{b} \arrow[r, Leftrightarrow] & \mathrm{d}\ (=\mathrm{PS}\ 1+\mathrm{PS}\ 2) \\
\mathrm{c} \arrow[r, Leftrightarrow] & \mathrm{c}',
\end{tikzcd}
\]
mais pas l'implication inverse. En d'autres termes, si les topos induits $\X/a$, pour $a$ dans $S$, sont totalement 0-connexes, je n'ai pas prouvé qu'il en est de même des $\X/W$, $W$ un objet 0-connexe \textit{quelconque} de $\X^\wedge$. Ainsi, on en est au point suivant

\textbf{Perplexité 2.11.} \textit{Soit $X$ dans} Cat. \textit{Si $X$ est localement filtrante à gauche, alors les produits fibrés $a\times_x b$ d'objets dans $X$ sont des objets} $0$\textit{-connexes de $X^\wedge$, et cette condition à son tour implique PS 1, i.e. que $X$ est localement irréductible. Mais en présence de cette hypothèse PS 1, ou $X$ localement irréductible, la condition sur les produits fibrés $a\times_x b$ équivaut à la suivante : pour tout diagramme} \texttt{[en]} \textit{traits pleins}

\msp{page 28}
\[
\begin{tikzcd}
c \arrow[r] & a\times_x b \\
e \arrow[u, dotted] \arrow[r, dotted] & d, \arrow[u]
\end{tikzcd}
\qquad\qquad
\textit{avec $x$, $a$, $b$, $c$, $d$ dans $X$,}
\]
\textit{on peut compléter en un diagramme} commutatif, \textit{avec $e$ dans $X$. En d'autre termes, tout diagramme commutatif à traits pleins ci-dessous se complète en y rajoutant un `objet initial' $e$} % [sic] d'autre termes -> d'autres termes, missing plural s
\[
\begin{tikzcd}[column sep=small, row sep=small]
 & & a \arrow[ddr] & \\
 & c \arrow[ur] \arrow[dddr] & & \\
e \arrow[ur, dotted] \arrow[dr, dotted] & & & x \\
 & d \arrow[uuur] \arrow[dr] & & \\
 & & b \arrow[uur] &
\end{tikzcd}
\]
(\footnote[33]{Il suffit d'exiger pour $e$ la commutativité des deux carrés $(e\ c\ b\ d)$ et $(e\ d\ a\ c)$.}).

Cette condition est donc impliquée par $X$ localement cofiltrante, et elle est satisfaite aussi par ce qui précède plus généralement si et seulement si (\footnote[34]{si et seulement si ?}) les $X/x$ sont cofiltrantes (\footnote[35]{Elle signifie que $X$ satisfait PS 1, et la condition

\textbf{PS$'$ 2} Tout diagramme commutatif
\[
a\rightrightarrows b\longrightarrow x
\]
% footnote continued on printed page 18
s'insère dans un diagramme commutatif
\[
y\longrightarrow a\rightrightarrows b\longrightarrow x.
\]
}). Mais la réciproque est-elle vraie - cette condition de totale 0-connexité des $X/x$ implique-t-elle que $X$ soit totalement 0-connexe ? Mais je n'ai pas construit d'exemple où les $X/x$
\printedpage{18}
% head: A. Grothendieck : Dérivateurs, Ch. II, Édition des Universités de Montpellier II et Paris VII
% foot: Version 8 juillet 2001
sont cofiltrantes sans que $X$ soit localement cofiltrante (ou ce qui revient au même, de $X$ 0-connexe avec les $X/x$ cofiltrantes sans que $X$ en soit cofiltrante)(\footnote[36]{Mais il y a un \textit{contre-exemple} évident, si $X$ est un groupoïde 0-connexe, les $X/x$ sont $\approx e$, donc cofiltrantes, mais $X$ n'est cofiltrant (i.e. \texttt{[mot(s) illisible(s)]}) que si $X$ est 1-connexe. P.ex. si $C$ est un groupoïde $\neq 1$, alors $X=B_G$ est tel que les $X/x$ sont cofiltrantes, mais $X$ non cofiltrant. % [sic] si $C$ est un groupoïde -> si $G$ est un groupoïde, the groupoid is then written $B_G$
}) ; ni d'exemple où les $X/x$ sont totalement 0-connexes, sans que les $X/x$ soient cofiltrantes. Mais notons le

\textbf{Corollaire 2.12 :} \textit{Les conditions équivalentes a), a$'$), b), d) (à l'exclusion de c), c$'$), plus faibles apparemment) du théorème 2.6}

\msp{page 29}\textit{équivalent encore à la suivante}
\begin{itemize}[leftmargin=2.5em]
\item[\textit{e)}] \textit{$\forall\ x$ dans $S$, le topos induit $\X/x$ et le topos induit $\X/x\times x$ est irréductible} (\footnote[37]{NB $\X/x$ inutile.}),
\end{itemize}
\textit{où encore à la suivante} % [sic] où encore -> ou encore, conjunction "ou" without accent
\begin{itemize}[leftmargin=2.5em]
\item[\textit{e$'$)}] \textit{$\forall\ x$ dans $S$, $S^\wedge/x$ et $S^\wedge/x{\times}x$ sont irréductibles, i.e. $S/x$ et $S/x{\times}x$ sont irréductibles} (\footnote[38]{$S/x$ inutile, car $x\xrightarrow{\mathrm{diag}}x\times x$ est une immersion ouverte non vide, donc si $x\times x$ est irréductible, $x$ l'est aussi. Mais il revient au même de dire que les $S/x$ sont irréductibles (condition PS$^*$ 1, cf. 2.8, p. 21) \textit{et} que les $x\xrightarrow{\mathrm{diag}}x\times x$ sont dominants (ce que implique en effet que $x\times x$ est irréductible). Ceci redonne immédiatement la forme PS$^*$ 2 de la condition cofiltrante. % [sic] ce que implique -> ce qui implique, relative pronoun subject
}).
\end{itemize}
C'est la forme `topologique' la plus faible en apparence que je peux trouver pour la condition $S$ localement cofiltrante, i.e. $\X$ (ou Top$(S)$, au choix) localement totalement 0-connexe.

\textbf{Définition 2.13.} On dit qu'un topos $\X$ est \textit{localement totalement} $0$\textit{-connexe} s'il satisfait aux conditions équivalentes a, a$'$, b, d du théorème 2.6, i.e. si $\X$ est localement connexe et si de plus tout topos induit 0-connexe est totalement 0-connexe - le produit fibré de deux objets 0-connexes sur un troisième est 0-connexe. Il revient au même de dire que toute (ou que une) sous-catégorie pleine génératrice $S$ de $\X^\wedge$ formée d'objets 0-connexes est localement cofiltrante (en supposant toujours $\X$ localement connexe, ou ce qui revient au même, l'existence d'un $S$ formé d'objets 0-connexes).

On n'a pas terminé encore la démonstration de 2.6 - il reste à prouver que si $\X$ (supposé localement totalement 0-connexe) est 0-connexe, il est $\infty$-connexe - ou encore que les composantes connexes de $\X$, supposé localement totalement 0-connexe, sont $\infty$-connexes, i.e. $W_\infty$-asphériques.

\msp{page 30}\textbf{Corollaire 2.14.} \textit{Notons quand même que si $\X$ est de la forme Top$(T)$, $T$ un espace topologique, alors $\X$ est localement totalement} $0$\textit{-connexe si et seulement si il est localement irréductible, i.e. si et seulement si $T$ est localement irréductible.}


\endgroup
% END UNIT Der-ChII-234/Der-ChII-234_p016-020.tex
% BEGIN UNIT Der-ChII-234/Der-ChII-234_p021-025.tex
\begingroup\makeatletter%
\newcommand{\printedpage}[1]{\marginpar{\footnotesize p.~#1}}%
\newcommand{\msp}[1]{\marginpar{\footnotesize [#1]}}%
\newcommand{\Hom}{\operatorname{Hom}}%
\newcommand{\Ker}{\operatorname{Ker}}%
\newcommand{\Ob}{\operatorname{Ob}}%
\makeatother\let\printedpage\derpp\let\msp\dermsp%
\let\mspin\DERmspin%


\printedpage{19}
% head: A. Grothendieck : Dérivateurs, Ch. II, Édition des Universités de Montpellier II et Paris VII
% foot: Version 8 juillet 2001

On est ramené à prouver que cette dernière condition implique la condition d) de 2.6, et on peut prendre pour $S$ l'ensemble ordonné des ouverts connexes non vides de $S$, cela devient alors évident.

Ceci dit, il faut se rappeler que la théorie cohomologique des espaces topologiques irréductibles n'est nullement triviale - comme nous le montre FAC de \textsc{Serre}\,! D'autre part, c'est \textsc{Serre} qui m'avait fait observer que pour des coefficients \textit{constants,} la cohomologie d'un espace irréductible est triviale.

Dans le cas actuel, prouvons déjà, si $\mathcal{X}$ est 0-connexe et localement totalement 0-connexe, que tout revêtement de l'objet final de $\mathcal{X}$ est trivial. Il suffit de voir que si $U$ est 0-connexe, $U \longrightarrow e$ est un isomorphisme. Mais considérons $U \hookrightarrow U \times U$, on sait que c'est un sommande direct, et comme $U \times U$ est connexe, c'est un isomorphisme. Donc $U \longrightarrow e$ \texttt{[est un]} monomorphisme, donc sommande direct, et comme $U \neq \emptyset$, on a $U \longrightarrow e$ isomorphisme.

Reste à prouver que pour tout faisceau \msp{page 31} abélien constant $F$ sur $\mathcal{X}$, on a $H^i(X,F)=0$ pour $i>0$. On procède par récurrence sur $i$, c'est déjà fait pour $i=1$. Supposons que ce soit prouvé pour les $j<i$ ($i$ fixé) pour tout topos \textit{totalement} 0\textit{-connexe.} Soit $\xi \in H^i(\mathcal{X},F)$, soit $U \longrightarrow e$ \texttt{[un]} épimorphisme tel que $\xi|U=0$, et considérons la suite spectrale % [sic] H^i(X,F) -> H^i(\mathcal{X},F), plain X printed where the topos \mathcal{X} is meant elsewhere
\[
E_2^{pq} = H^p(C(U^*,\mathcal{H}^q(F))) \Longrightarrow H^*(\mathcal{X},F).
\]
On a par hypothèse totalement 0-connexe
\[
\pi_0(U^q) \simeq \pi_0(U)^q, \tag*{$(*)$}
\]
et les composantes connexes de $U^q$ sont totalement 0-asphériques,
\begin{center}
\begin{tikzpicture}[scale=0.8]
\draw[->] (0,0) -- (0,5);
\node[left] at (0,5) {$q$};
\draw[->] (0,0) -- (5,0);
\node at (4.3,-0.6) {$p$};
\node[left] at (0,4) {$i$};
\foreach \x in {0,1,2,3,4} {\fill (\x,4) circle (2pt); \fill (\x,0) circle (2pt);}
\end{tikzpicture}
\end{center}
donc par hypothèse de récurrence
\[
E_1^{pq} = 0 \qquad \text{si } 0<q<i,
\]
donc on trouve une suite exacte
\[
0 \longrightarrow H^i(C(U^*,F)) \longrightarrow H^i(\mathcal{X},F) \longrightarrow H^0(\ \underbrace{H^i(U^*,F)}_{\substack{\text{complexe}\\ \text{co-semisimplicial}\\ \text{pour } i\\ \text{variable}}}\ ) = \Ker(H^i(U,F) \rightrightarrows H^i(U^2,F)).
\]

\printedpage{20}
% head: A. Grothendieck : Dérivateurs, Ch. II, Édition des Universités de Montpellier II et Paris VII
% foot: Version 8 juillet 2001

Par hypothèse sur $U$, l'image de $\xi$ dans le troisième terme est nulle, donc $\xi$ appartient au premier terme. Mais pour $(*)$, le complexe $C(U^*,F)$ est le complexe \textit{trivial} des puissances de $I=\pi_0(U)$, à coefficients dans $M$ (si $F=M_{\mathcal{X}}$). Sa cohomologie \msp{page 32} en degrés $>0$ est nulle, donc $\xi=0$, q.e.d.

\noindent\textbf{Proposition 2.15.} (\footnote[39]{C'est déjà dans SGA 4, I 8.}) \textit{Soit $X$ un objet de} Cat\textit{. Conditions équivalentes :}
\begin{enumerate}[label={}, leftmargin=2em, itemsep=0pt]
\item \textit{a) $X$ est filtrant, i.e. $X$ est $0$-connexe et $X^o$ satisfait les conditions équivalentes du théorème 2.6 ($X^{o,\wedge}$ est localement totalement $0$-connexe).}
\item \textit{b) L'objet final de $X^\wedge$ est ind-représentable.}
\end{enumerate}
\noindent\textsc{Démonstration}. On sait qu'on a
\[
e_{X^\wedge} \simeq \varinjlim_{x \in X} x, \tag*{(2.15.1)}
\]
limite inductive prise dans $X^\wedge$. En effet, cela signifie aussi que pour tout objet $a$ dans $X$, on a
\[
\varinjlim_{x \in X} \ \underbrace{x(a)}_{\Hom_X(a,x)} \ \xrightarrow{\ \sim\ } e \qquad \text{(ensemble ponctuel)}
\]
Mais considérons l'objet $X' = a\backslash X \longrightarrow X$ coétale sur $X$, qui est donc Cat-cofibré sur $X$, on sait que l'on a
\[
\pi_0(X') \simeq \varinjlim_{x \in X} \pi_0(X'_x),
\]
où $X'_x = \Hom(a,x)$, donc on trouve
\[
\varinjlim_{x \in X} \Hom(a,x) = \pi_0(a\backslash X) = e
\]
car $a\backslash X$ est 0-connexe.

La formule (2.15.1) montre que si \msp{page 33} $X$ est filtrant, $e_{X^\wedge}$ est limite inductive filtrante de foncteurs représentables. On sait d'ailleurs qu'on peut réindexer par une catégorie \textit{ordonnée} filtrante (cf. SGA 4, I 8.1.6 (c'est dû à P. \textsc{Deligne}\,!)).

On a prouvé a $\Longrightarrow$ b, prouvons b $\Longrightarrow$ a. Supposons donc
\[
e_{X^\wedge} = \varinjlim_{i \in I} x_i,
\]
où $I$ est un (petit) ensemble ordonné filtrant, et $I \longrightarrow X$ un foncteur. Prouvons PS 1$'$ (forme droite), i.e. pour $a, b$ dans $X$, \texttt{[il]} existe $\vcenter{\hbox{\begin{tikzcd}[row sep=0.1em, column sep=1.2em]
a \arrow[dr] & \\
& x \\
b \arrow[ur] &
\end{tikzcd}}}$. En effet, comme $\varinjlim\limits_{I} \Hom(a,x_i) = e$, il s'ensuit que $\Hom(a,x_i) \neq \emptyset$ pour $i$ grand, de même $\Hom(b,x_i) \neq \emptyset$, donc il suffit de prendre $\vcenter{\hbox{\begin{tikzcd}[row sep=0.1em, column sep=1.2em]
a \arrow[dr] & \\
& x_i \\
b \arrow[ur] &
\end{tikzcd}}}$. Prouvons PS 2$_{\mathrm{d}}$, i.e. que pour \texttt{[une]} double flèche $a \overset{u}{\underset{v}{\rightrightarrows}} b$, \texttt{[il]} existe $b \longrightarrow x$ qui égalise.

\printedpage{21}
% head: A. Grothendieck : Dérivateurs, Ch. II, Édition des Universités de Montpellier II et Paris VII
% foot: Version 8 juillet 2001

\msp{page 34}
Soit $\xi_{i_0}$ dans $\Hom(b,x_{i_0})$, ses deux images dans $\Hom(a,x_{i_0})$ ne sont pas nécessairement égales, mais le deviennent dans la limite inductive, réduite à $e$, donc deviennent égales en composant avec $x_{i_0} \longrightarrow x_i$ pour $i$ assez grand, donc le composé $b \longrightarrow x_{i_0} \longrightarrow x_i$
\[
a \rightrightarrows b \longrightarrow x_{i_0} \longrightarrow x_i
\]
égalise la double flèche.

\noindent\textbf{Corollaire 2.16.} \textit{Soit $W$ un localiseur fondamental satisfaisant l'axiome des limites} Loc (8)\textit{. Si $X$ dans} Cat \textit{est une catégorie filtrante ou cofiltrante, alors $X$ est $W$-asphérique.}

Plus généralement et agréablement :

\noindent\textbf{Corollaire 2.17.} \textit{Soit $W$ comme dans 2.16. Alors toute catégorie filtrante est} totalement $W$-asphérique (\footnote[40]{cf. XII, §V pour cette notion, que j'ai beaucoup peiné à dégager et à comprendre\,!}), \textit{i.e. pour toute flèche $f : X \longrightarrow Y$ dans} Cat\textit{, l'objet $f_!(e_X)$ de $e_Y$ (défini aussi comme $y \longmapsto \pi_0(y\backslash X)$, et intervenant dans la factorisation canonique de $f$ en $ip$, avec $p$ étale et les cofibres} \texttt{[?]} \textit{de $p$ $0$-connexes)} \msp{page 35} \textit{est $W$-asphérique. En fait, $Y/f_!(e_{X^\wedge})$ est même une catégorie filtrante.}

En effet, le foncteur $f_!$ commute aux limites inductives, donc $f_! e_{X^\wedge}$ est limite inductive filtrante des $f_!(x_i) = f(x_i)$. Donc $Y/f_!(e_{X^\wedge})$ est limite inductive filtrante des $Y/f(x_i)$. Donc c'est même une catégorie filtrante, a fortiori elle est $W$-asphérique.

Donc finalement, on trouve des exemples non triviaux de catégories totalement $W$-asphériques, i.e. des exemples donc des catégories sans objet final. % [sic] donc -> dont (?), the conjunction "donc" makes no sense here

\noindent\textbf{Corollaire 2.18.} (\footnote[41]{déjà dans SGA 4, I 8.}) \textit{Soit $X$ dans} Cat\textit{, $F$ dans $X^\wedge$. Pour que $F$ soit ind-représentable il faut et il suffit que $X/F$ soit filtrante.}

En effet, par (2.15), $X/F$ est filtrante si et seulement si l'objet final $F$ de $X/F$ provient d'un ind-objet de $X/F$, mais cela signifie aussi que $F$ est la limite inductive dans $X^\wedge$ d'un ind-objet de $X$, i.e. que $F$ est ind-représentable.

\msp{page 36}
\noindent\textbf{Question.} \textit{Si $X$ dans} Cat \textit{est totalement $W_\infty$-asphérique, i.e. les $f_!(e_{X^\wedge})$ sont $W_\infty$-asphériques pour toute flèche $f : X \longrightarrow Y$ dans} Cat\textit{, ou encore pour toute telle flèche, avec les $y\backslash Y$ $0$-connexes ($y \in \Ob Y$), $Y$ est $W_\infty$-asphérique - cela implique-t-il que $X$ soit filtrante ?} Ça ne m'étonnerait pas - et même en supposant seulement les $f_!(e_{X^\wedge})$ 1-connexes (\footnote[42]{Ils le sont toujours (cf. XII, p. 101), c'est 2-\textit{connexe} qu'il faut lire ici\,!}) (ils sont 0-connexes, dès que $X$ l'est). Mais je ne veux pas m'attarder à présent sur cette question.

La réponse à la question précédente est négative, comme on voit par la :

\printedpage{22}
% head: A. Grothendieck : Dérivateurs, Ch. II, Édition des Universités de Montpellier II et Paris VII
% foot: Version 8 juillet 2001

\noindent\textbf{Proposition 2.19.} (\footnote[43]{\textbf{Faux.}}) \textit{Soit $X$ dans} Cat \textit{telle que $\mathit{diag}_X : X \longrightarrow X \times X$ soit cofinal} (\footnote[44]{cf. définition 3.3, p. 45.})\textit{, i.e. telle que $X^o$ soit totalement $0$-connexe. Alors $X$ est totalement $W_\infty$-asphérique (donc aussi totalement $W$-asphérique pour tout localiseur fondamental $W$ tel que $W \supseteq W_\infty$).}

Soit en effet un foncteur cofinal $f : X \longrightarrow Y$, il faut prouver que $Y$ est $W_\infty$-asphérique. Or on sait que totalement 0-connexe implique $W_\infty$-asphérique (\footnote[45]{C'est \textbf{faux.}}) (cf. proposition 2.3, p. 17-18), donc il suffit de voir que $Y^o$ est totalement 0-connexe, i.e. $Y \xrightarrow{d_Y} Y \times Y$ cofinal. Or on a le diagramme commutatif
\[
\begin{tikzcd}[row sep=large, column sep=large]
X \arrow[r, "d_X"] \arrow[d, "f"'] & X \times X \arrow[d, "f \times f"] \\
Y \arrow[r, "d_Y"] & Y \times Y,
\end{tikzcd}
\]
où par hypothèse $f$, donc aussi $f \times f$, et $d_X$ sont cofinaux. Donc $\varphi = (f \times f) \circ p_X$ \texttt{[plutôt} $\varphi = (f \times f) \circ d_X$\texttt{]} est cofinal, et comme $\varphi = d_Y f$ et $f$ est aussi cofinal, on conclut que $d_Y$ l'est, q.e.d.

\noindent\textbf{Corollaire.} $\boldsymbol{\Delta}$ \textit{étant totalement $0$-connexe, $\boldsymbol{\Delta}^o$ est totalement $W_\infty$-asphérique, pourtant elle n'est pas filtrante.}

\noindent\textbf{Question.} \textit{$X$ totalement $W_\infty$-asphérique $\Longrightarrow$ $X^o$ totalement $0$-connexe ? Ou du moins implique PS 1$'$ ?}

\msp{page 36bis}
\noindent\textbf{Contrexemple 2.20} à la proposition fausse 2.19.

Soit $X$ un objet de Cat, rajoutons-lui la catégorie finale $e$ comme fermé, i.e. considérons un objet $X'$ de Cat dans lequel $X$ soit un sous-objet ouvert, $e$ le fermé complémentaire :
\[
X \overset{i}{\hookrightarrow} X' \overset{j}{\hookleftarrow} e.
\]
On sait (cf. XII) qu'une telle opération est déterminée, à isomorphisme près, par la donnée d'un objet
\[
H \in \Ob X^\wedge
\]
en posant
\[
H(x) = \Hom_{X'}(x,e),
\]
la catégorie $X/e$ s'identifiant (à isomorphisme près) à $X/H$. Elle est donc 0-connexe si et seulement si l'objet $H$ de $X^\wedge$ est 0-connexe. Cela signifie aussi que le morphisme $i$ est cocofinal, i.e. que $i^o : X^o \longrightarrow X'^o$ est cofinal. Supposons maintenant $X$ $W_\infty$-asphérique et totalement 0-connexe, i.e. l'actuel Loc (6), cf. XVI. La proposition 2.19 affirme que, sous ces conditions, $X'$ (qui est lui aussi totalement 0-connexe) est même $W$-asphérique. Il s'ensuivrait donc que $i \in W_\infty$. Mais par XII, th. 4, cor. 4 (p. 149), il s'ensuivrait que
\printedpage{23}
% head: A. Grothendieck : Dérivateurs, Ch. II, Édition des Universités de Montpellier II et Paris VII
% foot: Version 8 juillet 2001
\msp{page 36ter}
$X/e$ est $W_\infty$-asphérique, i.e. $H$ $W_\infty$-asphérique. Ainsi, si $X$ est totalement 0-asphérique, tout objet 0-connexe de $X^\wedge$ serait $\infty$-connexe. Mais c'est manifestement faux, p.ex. si on prend pour $X$ une catégorie-test pour $W_\infty$, p.ex. $X = \boldsymbol{\Delta}$.

Cela me fait soupçonner que la conjecture
\[
X \text{ totalement } W_\infty\text{-asphérique} \Longrightarrow X \text{ filtrante}
\]
pourrait être valable.


\endgroup
% END UNIT Der-ChII-234/Der-ChII-234_p021-025.tex
% BEGIN UNIT Der-ChII-234/Der-ChII-234_p026-030.tex
\begingroup\makeatletter%
\newcommand{\printedpage}[1]{\marginpar{\footnotesize p.~#1}}%
\newcommand{\msp}[1]{\marginpar{\footnotesize [#1]}}%
\newcommand{\mspin}[1]{\texttt{[#1]}}%
\newcommand{\HOT}{\mathrm{HOT}}%
\newcommand{\Cat}{\mathrm{Cat}}%
\newcommand{\Fl}{\mathrm{Fl}}%
\newcommand{\Ob}{\mathrm{Ob}}%
\newcommand{\Ens}{\mathrm{Ens}}%
\newcommand{\Hom}{\underline{\mathrm{Hom}}}%
\newcommand{\defeq}{\overset{\text{déf}}{=}}%
\newcommand{\pl}{\texttt{[plutôt}}%
\newcommand{\uHom}{\underline{\operatorname{Hom}}}%
\makeatother\let\printedpage\derpp\let\msp\dermsp%
\let\mspin\DERmspin%


\printedpage{24}
% head: A. Grothendieck : Dérivateurs, Ch. II, Édition des Universités de Montpellier II et Paris VII
% foot: Version 8 juillet 2001
\msp{page 37}

\section*{3\quad Retour sur la cofinalité}

Il est instructif de relier les questions de cofinalité (pour une flèche $u : I' \longrightarrow I$ dans Cat) au dérivateur $\HOT_{W_0}$, où $W_0 \subseteq \Fl\,\Cat$ est le plus grossier des localiseurs dans Cat satisfaisant l'axiome de connexité Loc 5, i.e.
\[
W_0 = \{ f \in \Fl\,\Cat \mid \pi_0(f) \text{ bijectif}\}. \tag*{(3.1)}
\]
On trouve alors, pour $I$ dans Cat, une équivalence canonique, 2-fonctorielle en $I$,
\[
\HOT_{W_0}(I) \xrightarrow{\;\sim\;} I^{\wedge} \defeq \uHom(I,(\Ens)), \quad \text{\pl}\ \uHom(I^o,\Ens)\,\text{\texttt{]}} \tag*{(3.2)}
\]
en d'autres termes, on a une équivalence
\[
\HOT_{W_0} \simeq \mathbf{D}_{(\Ens)}. \tag*{(3.3)}
\]
Quand $\HOT_{W_0}$ est défini par localisation à partir de la catégorie $\Cat(I) \defeq \uHom(I^o,\Cat)$, la flèche (3.2) est celle déduite, par localisation, de la flèche
\[
\begin{array}{rcl}
\pi_0^I : \uHom(I^o,\Cat) & \longrightarrow & \uHom(I^o,(\Ens)) \\[1ex]
X & \longmapsto & \pi_0 \circ X,
\end{array} \tag*{(3.4)}
\]
où
\[
\pi_0 : \Cat \longrightarrow (\Ens) \tag*{(3.5)}
\]
est le foncteur ‘composantes connexes’. Si on décrit $\HOT_{W_0}(I)$ comme localisée de $\Cat/I$, (3.2) est déduit de la flèche
\[
\left\{
\begin{array}{rcl}
\Cat/I & \longrightarrow & I^{\wedge} = \uHom(I^o,(\Ens)) \\[1ex]
X & \longmapsto & \underline{\pi_0}(X/I) = (i \longmapsto \pi_0(i\backslash X)).
\end{array}
\right. \tag*{(3.6)}
\]

\msp{page 38}
Ainsi, les foncteurs familiers, pour une flèche
\[
u : I' \longrightarrow I
\]
dans Cat
\[
\left\{
\; I'^{\wedge} \quad
\begin{array}{c}
\xrightarrow{\;u_!\;} \\[1.5ex]
\xleftarrow{\;u^*\;} \\[1.5ex]
\xrightarrow{\;u_*\;}
\end{array}
\quad I^{\wedge},
\right. \tag*{(3.7)}
\]
\printedpage{25}
% head: A. Grothendieck : Dérivateurs, Ch. II, Édition des Universités de Montpellier II et Paris VII
% foot: Version 8 juillet 2001
s'interprètent comme les foncteurs de même nom pour le dérivateur $\HOT_{W_0}$, ou plus directement, pour le dérivateur (qu'on pourrait faussement considérer comme canularesque) $\mathbf{D}_{(\Ens)}$.

Cela donne un sens à la cohorte des notions cohomologiques associées à un dérivateur, ou plus particulièrement, à un localiseur fondamental $W$ sur Cat (donnant naissance au dérivateur $\HOT_W$). On les désignera par le nom de $W_0$-fourbi (au lieu de $\HOT_{W_0}$-fourbi, ou $\mathbf{D}_{(\Ens)}$-fourbi), notamment pour une flèche $u : I' \longrightarrow I$ de Cat telle la notion
\[
\left\{
\begin{array}{lll}
\text{$W_0$-équivalence} & \Bigl(\iff & \text{$\pi_0(u)$ bijectif.}\Bigr) \\
\text{$W_0$-asphérique} & \Bigl(\iff & \text{les $I'/i$ pour $i \in \Ob I$ sont} \\
 & & \text{$0$-connexes.}\Bigr) \\
\text{$W_0$-coasphérique} & \Bigl(\iff & \text{les $i\backslash I'$ pour $i \in \Ob I$ sont} \\
 & & \text{$0$-connexes.}\Bigr) \\
\text{$W_0$-propre (\footnotemark[46])} & \Bigl(\iff & \text{les inclusions $I'_i \longrightarrow I'/i$ ($i \in \Ob I$)} \\
 & & \text{sont $W_0$-coasphériques, i.e. les fibres} \\
 & & \text{des colocalisées $i'\backslash I' \longrightarrow i\backslash I$} \\
 & & \text{($i' = u(i)$) \pl\ ($i = u(i')$)]} \\
 & & \text{sont $0$-connexes.}\Bigr) \\
\text{\mspin{page 39}} & & \\
\text{$W_0$-lisse (\footnotemark[47])} & \Bigl(\iff & \text{les inclusions $I'_i \longrightarrow i\backslash I'$} \\
 & & \text{sont $W_0$-asphériques, i.e. les fibres} \\
 & & \text{des flèches localisées $I'/i' \longrightarrow I/i$} \\
 & & \text{($i = u(i')$) sont $0$-connexes.}\Bigr) \\
\text{$W_0$-parfait} & \Bigl(\iff & \text{$W_0$-propre et $W_0$-lisse.}\Bigr) \\
\text{$W_0$-équivalence universelle (\footnotemark[48])} & \Bigl(\iff & \text{$\pi_0$ bijectif et le reste par tout} \\
 & & \text{changement de base. Ce sont aussi} \\
 & & \text{les $W_0$-fibrations (ou $W_0$-fibrations} \\
 & & \text{faibles à fibres $W_0$-asphériques,} \\
 & & \text{i.e. $0$-connexes.}\Bigr) \\
\text{$W_0$-fibration} & \Bigl(\iff & \text{$u^* : \Cat/I \longrightarrow \Cat/I'$ transforme} \\
 & & \text{$W_0$-équivalences en $W_0$-équivalences.}\Bigr) \\
\text{$W_0$-fibration faible (\footnotemark[49])} & \Bigl(\iff & \text{le changement de base} \\
 & & \text{$u^* : \Cat/I \longrightarrow \Cat/I'$ transforme} \\
 & & \text{homotopismes $J_1 \longrightarrow J_2$ en} \\
 & & \text{$W_0$-équivalences - et il suffit de} \\
 & & \text{l'exiger pour l'inclusion d'un objet} \\
 & & \text{final (ou au choix initial) et pour} \\
 & & \text{l'inclusion d'un point dans $\boldsymbol{\Delta}^1$.}\Bigr)
\end{array}
\right. \tag*{(3.8)}
\]
\footnotetext[46]{Exemple : les Cat-cofibrations.}
\footnotetext[47]{Exemple : les Cat-fibrations.}
\footnotetext[48]{Exemple : les flèches $W_0$-lisses ou $W_0$-propres, à fibres 0-connexes.}
\footnotetext[49]{Exemple : les flèches $W_0$-lisses telles que $\pi_0(I'/I) = (i \longmapsto \pi_0(i\backslash I')) = u_!(e_{I'^{\wedge}})$ soit localement constant sur $I$, i.e. $\underline{\pi_0}(I'/I)$ soit un \textit{revêtement} étale de $I$, ou si $u^{\mathrm{opp}}$ satisfait cette condition.}

\printedpage{26}
% head: A. Grothendieck : Dérivateurs, Ch. II, Édition des Universités de Montpellier II et Paris VII
% foot: Version 8 juillet 2001
\texttt{[Nous avons une]} implication \textit{stricte} (et comment)
\[
\text{$W_0$-équivalence universelle} \Longrightarrow \text{$W_0$-fibration}
\]
\texttt{[et une]} implication \textit{stricte} (et comment\,!) (\footnote[50]{Mais si les fibres sont 0-connexes, les deux notions de $W_0$-fibrations sont équivalentes et signifient qu'on a une $W_0$-équivalence universelle \dots})
\[
\text{$W_0$-fibration} \Longrightarrow \text{$W_0$-fibration faible.}
\]

Signalons, comme exemple d'utilisation de ces notions, que pour un diagramme cartésien dans Cat
\[
\begin{tikzcd}
I' \arrow[d, "u"'] & J' \arrow[l, "p'"'] \arrow[d, "v"] \\
I & J, \arrow[l, "p"']
\end{tikzcd} \tag*{(3.9)}
\]
la flèche de changement de base pour $u_*$, $v_*$ pour les préfaisceaux est \texttt{[un]} isomorphisme si $u$ est $W_0$-propre, ou $p$ $W_0$-lisse, et la flèche de changement de base pour $u_!$, $v_!$ est \texttt{[un]} isomorphisme si $u$ est $W_0$-lisse ou $p$ $W_0$ propre (\footnote[51]{Sauf erreur, ca donne \textit{deux caractérisations} des morphismes $W_0$-propres, ou des morphismes $W_0$-lisses.}). Je signale à ce propos que ces énoncés, qui a priori se formulent pour les préfaisceaux à valeurs dans (Ens), restent valables % [sic] ca -> ça (in footnote 51), missing cedilla
\msp{page 40}
pour des préfaisceaux à valeurs dans une catégorie $\mathcal{M}$ quelconque (sous réserve d'existence des operations d'images directes cohomologiques ou homologiques envisagées). Cela se voit p.ex. en notant que % [sic] operations -> opérations, missing accent
\[
\text{$W_0$-équivalence} \;\Longrightarrow\; \text{$\mathbf{D}_{\mathcal{M}}$-équivalence} \tag*{(3.10)}
\]
pour toute $\mathcal{U}$-catégorie $\mathcal{M}$, d'où résulte que toute (\footnote[52]{Pas tout à fait exacte pour \textit{toute} ‘$W_0$-notion’, il faut excepter les deux notions de $W_0$-fibration forte, et faible.}) ‘$W_0$-notion’ implique la $\mathbf{D}_{\mathcal{M}}$-notion correspondante, et permet donc d'appliquer à $\mathbf{D}_{\mathcal{M}}$ les résultats le concernant. (Il y a bien sûr un argument direct, moins savant, consistant à se ramener au cas où $\mathcal{M}$ est petite, et à utiliser alors le plongement pleinement fidèle $\mathcal{M} \hookrightarrow \mathcal{M}^{\wedge}$ resp. $\mathcal{M} \hookrightarrow ((\mathcal{M}^o)^{\wedge})^o$.)

Je rappelle pour mémoire une (troisième\,?) caractérisation des flèches $u : I' \longrightarrow I$ $W_0$-propres (resp. $W_0$-lisses) : ce sont celles pour lesquelles le foncteur changement de base $u^* : \Cat/I \longrightarrow \Cat/I'$ transforme flèches $W_0$-coasphériques (resp. $W_0$-asphériques) en itou. Cela gagne en intérêt, du fait que la notion de foncteur $W_0$-coasphérique (resp. $W_0$-asphérique) s'interprète comme celle de foncteur cofinal (resp. co-cofinal), cf. plus bas.

\printedpage{27}
% head: A. Grothendieck : Dérivateurs, Ch. II, Édition des Universités de Montpellier II et Paris VII
% foot: Version 8 juillet 2001
\msp{page 41}
La notion de $W_0$-fibration apparaît comme extrêmement forte - beaucoup plus forte à certains égards que celle (disons) de $W_\infty$-fibration ou de $W_\omega$-fibration. Elle équivaut à ceci (comme on voit aussitôt) : pour toute composante connexe $I'_\alpha$ de $I'$, désignant par $I_\alpha$ la composante connexe correspondante de $I$, et pour tout $J_\alpha$ 0-connexe et au dessus de $I_\alpha$, le produit fibré $I'_\alpha \times_{I_\alpha} J_\alpha$ est encore 0-connexe. (Ça ressemble aux histoires de totale 0-connexité d'un topos\,!) Mais cela signifie aussi que $I'_\alpha \longrightarrow I_\alpha$ est une $W_0$-équivalence universelle. Cela implique que ses fibres sont 0-connexes, et la réciproque est vraie si $u : I' \longrightarrow I$ est $W_0$-propre ou $W_0$-lisse. Ainsi, la condition essentielle pour une $W_0$-fibration est celle que les fibres des flèches induites $I'_\alpha \longrightarrow I_{u(\alpha)}$ \pl\ $I'_\alpha \longrightarrow I_\alpha$\texttt{]} soient 0-connexes. En fait, on a ceci

\noindent\textbf{Proposition 3.1.} \textit{Soit $u : I' \longrightarrow I$ une flèche dans} Cat.
\begin{itemize}
\item[\textit{a)}] \textit{Pour que ce soit une $W_0$-équivalence universelle, il faut et il suffit que pour toute $\boldsymbol{\Delta}^1 \longrightarrow I$, i.e. toute flèche $\alpha$ de $I$, la catégorie $I' \times_I \boldsymbol{\Delta}^1 \defeq I'|\alpha$ soit $0$-connexe, ou encore que les fibres de $I$ soient $0$-connexes, et que pour toute flèche} non identique \textit{$\alpha$ de $I$, $I'|\alpha$ soit $0$-connexe.}
\item[\textit{b)}] \textit{Pour que $u : I' \longrightarrow I$ soit une $W_0$-fibration, il faut et il suffit que pour toute composante connexe $I'_{\alpha'}$ de}
\msp{page 42}
\textit{$I'$, posant $\alpha = u(\alpha')$, la flèche induite $I'_{\alpha'} \longrightarrow I_\alpha$ soit une $W_0$-équivalence universelle (cf. critère a) pour cette notion).}
\item[\textit{c)}] \textit{Pour que ce soit une $W_0$-fibration faible, il} faut \textit{que $\forall\, i \in I$, l'inclusion $I'_i \hookrightarrow i\backslash I'$ soit une $W_0$-équivalence, i.e. $\pi_0(I'_i) \xrightarrow{\;\sim\;} \pi_0(i\backslash I')$ soit bijectif,} et \textit{que $i \longmapsto \pi_0(i\backslash I')$ soit un préfaisceau localement constant sur $I$ (i.e. $\underline{\pi_0}(I'/I)$ soit un} revêtement \textit{étale de $I$). Cette condition est aussi suffisante, si $u$ est $W_0$-lisse. De plus (pour mémoire) $u$ est une $W_0$-fibration faible si et seulement si $u^{\mathrm{opp}}$ l'est.}
\end{itemize}

La plus grande différence entre la notion de $W_0$-fibration, et celle de $W_0$-fibration faible, me paraît que les morphismes $W_0$-parfaits (i.e. $W_0$-lisses et $W_0$-propres), et a fortiori les $W$-parfaits pour un localiseur fondamental $W \subseteq W_0$, sont forcément des $W$-fibrations faibles, mais (\textit{même si $W = W_\omega$,} le plus fin de tous les localiseurs fondamentaux) ne sont pas des $W$-fibrations. L'exemple le plus frappant à cet égard est sans doute celui des revêtements étales, qui sont $W_\omega$-parfaits (donc
\msp{page 43}
des $W_\omega$-fibrations faibles et a fortiori des $W_0$-fibrations faibles), mais pas des $W_0$-fibrations\,; car dans le cas où $I'$, $I$ sont 0-connexes, il faudrait que les fibres soient 0-connexes (cf. b) ci-dessus), i.e. que $u$ soit un isomorphisme\,!

Une propriété sympathique des morphismes $W_0$-lisses (resp. $W_0$-propres), c'est que ce sont des applications \textit{ouvertes} (resp. fermées) pour les espaces topologiques sous-jacents, ce qui signifie en effet que les fibres locales, i.e. des flèches localisées en haut $I'/i' \longrightarrow I/i$ (resp. colocales, i.e. les flèches colocalisées en bas $i'\backslash I' \longrightarrow i\backslash I$), sont surjectives, i.e. à fibres non-vides (ce qui est impliqué par la condition que celles-ci soient 0-connexes).

\printedpage{28}
% head: A. Grothendieck : Dérivateurs, Ch. II, Édition des Universités de Montpellier II et Paris VII
% foot: Version 8 juillet 2001
Dans tout ça je n'ai encore guère parlé des flèches $W_0$-asphériques resp. co-asphériques, qui pourtant sont ici une motivation principale pour introduire $W_0$\,!

\noindent\textbf{Proposition 3.2.} \textit{Conditions équivalentes sur la flèche $u : I' \longrightarrow I$ dans} Cat :
\begin{itemize}
\item[a)] \textit{Pour toute $\mathcal{U}$-catégorie $\mathcal{M}$ et tout foncteur $F : I \longrightarrow \mathcal{M}$, la flèche canonique (sous réserve que les petites limites inductives soient représentables dans $\mathcal{M}$),}
\[
\varinjlim_{I'} u^*(F) \longrightarrow \varinjlim_{I} F \tag*{(3.2.1)}
\]
\msp{page 44}
\textit{est} \texttt{[un]} \textit{isomorphisme.}
\item[a$'$)] \textit{Pour toute $\mathcal{U}$-catégorie $\mathcal{M}$ où les petites limites projectives sont représentables et tout $F \in \mathcal{M}(I) = \uHom(I^o,\mathcal{M})$, i.e. tout préfaisceau sur $I$ à valeurs dans $\mathcal{M}$, la flèche canonique}
\[
\varprojlim_{I} F \longrightarrow \varprojlim_{I'} u^*(F) \tag*{(3.2.2)}
\]
\textit{est} \texttt{[un]} \textit{isomorphisme.}
\item[a$''$)] (\footnote[53]{C'est une condition sur les morphismes de topos, savoir que si $F$ \texttt{[est]} dans $\mathcal{X}^{\wedge}$, on a
\[
H^0(\mathcal{X},F) \xrightarrow{\;\sim\;} H^0(\mathcal{X}',u^*F).
\]
Il y a \texttt{[une]} notion similaire pour les $H^i$ ($i \leq n$), à dégager. Notion de morphismes $W_0$-coasphériques de topos\,! C'est dans la nature d'une condition de 0-connexité relative, duale de la notion $W_0$-asphérique, qui signifie simplement que pour $U$ 0-connexe dans $U^{\wedge}$ \pl\ $\mathcal{X}^{\wedge}$\texttt{]}, $f^*(U)$ \pl\ $u^*(U)$\texttt{]} est 0-connexe. De même, on a sans doute des notions (‘duales’ des précédentes) de $W_n$-coasphéricité (bien comprise en principe si $\mathcal{X}$ est localement $n$-asphérique) (La $W_n$-coasphéricité, étant celle évoquée plus haut.)}) \textit{Comme a$'$, avec $\mathcal{M} = (\Ens)$, i.e. $F \in \Ob I^{\wedge}$.}~(\footnote[54]{\textbf{NB} On peut introduire
\begin{itemize}
\item[a$'''$)] Comme a), mais avec $\mathcal{M} = (\Ens)$.
\item[a$^{IV}$)] Comme a), avec le foncteur canonique $I' \longrightarrow I'^{\wedge}$.
\end{itemize}}).
\item[b)] \textit{La flèche $u$ est $W_0$-coasphérique, i.e. pour tout $i \in I$, $i\backslash I'$ est $0$-connexe.}
\end{itemize}
L' équivalence de a, a$'$, a$''$ est $\pm$ tautologique en termes des définitions (cf. SGA 4, I, 8\dots). La forme a$''$ s'écrit, posant $\mathbf{D} = \HOT_{W_0}$, que $\forall\, F \in \mathbf{D}(I)$, on veut que
\[
\underbrace{H^{\bullet}_{\mathbf{D}}(I,F)}_{\defeq p_{I,*}(F)} \longrightarrow \underbrace{H^{\bullet}_{\mathbf{D}}(I',u^*F)}_{\defeq p_{I',*}(F)} \qquad \text{\pl}\quad \underbrace{H^{\bullet}_{\mathbf{D}}(I',u^*F)}_{\defeq p_{I',*}(u^*F)}\ \text{\texttt{]}} \tag*{(3.2.3)}
\]
($p_I : I \longrightarrow e$ et $p_{I'} : I' \longrightarrow e$ les foncteurs canoniques) soit \texttt{[un]} isomorphisme. On sait que cela équivaut à $u$ $\mathbf{D}$-coasphérique. La démonstration est un jeu amusant entre $u_!$, $u^*$, $u_*$ à la fois. Il faut exprimer que le Hom d'un objet $M$ de $\mathbf{D}(e)$ ($=$ (Ens) en l'occurence) dans la flèche (3.2.3) est bijective. Ici, il % [sic] occurence -> occurrence, French spelling


\endgroup
% END UNIT Der-ChII-234/Der-ChII-234_p026-030.tex
% BEGIN UNIT Der-ChII-234/Der-ChII-234_p031-035.tex
\begingroup\makeatletter%
\newcommand{\printedpage}[1]{\marginpar{\footnotesize p.~#1}}%
\newcommand{\msp}[1]{\marginpar{\footnotesize [#1]}}%
\newcommand{\Hom}{\operatorname{Hom}}%
\newcommand{\uHom}{\underline{\operatorname{Hom}}}%
\newcommand{\Ob}{\operatorname{Ob}}%
\newcommand{\bs}{\backslash}%
\newcommand{\cM}{\mathcal{M}}%
\newcommand{\dbl}[2]{\mathrel{\overset{#1}{\underset{#2}{\rightrightarrows}}}}%
\newcommand{\defeq}{\overset{\text{déf}}{=}}%
\makeatother\let\printedpage\derpp\let\msp\dermsp%
\let\mspin\DERmspin%


\printedpage{29}
% head: A. Grothendieck : Dérivateurs, Ch. II, Édition des Universités de Montpellier II et Paris VII
% foot: Version 8 juillet 2001

\msp{page 45}
est plus naturel de prendre $M = e$ (ensemble ponctuel) et d'interpréter la flèche (3.2.3) comme
\[
\Hom(e_I, F) \longrightarrow \Hom(e_{I'}, u^*F)
\tag*{(3.2.4)}
\]
et par adjonction de $u_!$, $u^*$ le deuxième membre s'interprète comme $\Hom(u_!(e_{I'}), F)$, et la flèche est déduite de
\[
u_!e_{I'} = u_!u^*(e_I) \xrightarrow{\text{adj.}} e_I
\tag*{(3.2.5)}
\]
en prenant le Hom dans $F$. Donc dire que (3.2.4) est \texttt{[un]} isomorphisme pour tout $F$ signifie que (3.2.5) est \texttt{[un]} isomorphisme. Passant aux valeurs des deux membres en $i$, on trouve $\pi_0(i\bs I') \xrightarrow{\sim} 1$, i.e. que $i\bs I'$ est 0-connexe, q.e.d.

\noindent\textbf{Définition 3.3.}
\begin{itemize}[label={}]
\item[a)] Quand les conditions équivalentes de 3.2 sont satisfaites, donc quand les $i\bs I'$ ($i \in \Ob I$) sont 0-connexes, on dit que $u : I' \longrightarrow I$ est un foncteur \textit{cofinal.} On dit que $u$ est \textit{co-cofinal} si $u^o : I'^o \longrightarrow I^o$ est cofinal, i.e. si les localisées $I'/i$ de $I'$ sur $I$ ($i \in \Ob I$) sont 0-connexes.
\item[b)] Si $I'$ est une \textit{sous-catégorie} de $I$, on dit que $I'$ \textit{est cofinale dans} $I$ si le foncteur d'inclusion est cofinal, \textit{co-cofinale} si ce foncteur est co-cofinal.
\end{itemize}

\noindent\textbf{Corollaire 3.4.} (\footnote[55]{\textbf{NB} Pour un énoncé complémentaire de cofinalité, cf. XVIII prop. 4.9.6.3 (p. 224-225, 228-242). Il faudrait un énoncé qui coiffe les deux.})
\begin{itemize}[label={}]
\item[\textit{a)}] \textit{Pour que $u$ soit cofinal, il} faut \textit{que les catégories colocalisées $i\bs I'$ soient non-vides, i.e. qu'on ait}

\msp{page 46}
\begin{itemize}[label={}]
\item[\textit{F 0)}] $\pi_0(f)$ \texttt{[plutôt $\pi_0(u)$]} \textit{est bijectif.}
\item[\textit{F 1)}] \textit{Pour tout} $i \in \Ob I$, \texttt{[il]} \textit{existe} \texttt{[un]} $i' \in I'$ \textit{et une flèche} $i \longrightarrow u(i')$ \textit{dans} $I$.
\end{itemize}
(\footnote[56]{\textbf{NB} F 1 implique déjà \texttt{[la]} surjectivité de $\pi_0$, donc sa bijectivité si $I'$ est 0-connexe - dans ce cas F 0 résulte de F 1.}).
\item[\textit{b)}] \textit{Si $I'$ est faiblement filtrante, pour que $u$ soit cofinale, il} suffit \textit{qu'elle satisfasse F 1, ainsi que la condition :}
\begin{itemize}[label={}]
\item[\textit{F 2)}] \textit{Pour toute double flèche} $i \dbl{\beta}{\gamma} u(i')$ \textit{(avec} $i \in \Ob I$, $i' \in \Ob I'$\textit{),} \texttt{[il]} \textit{existe une flèche} $i' \xrightarrow{\alpha'} j'$ \textit{dans} $I'$\textit{, de telle façon que} $u(\alpha')$ \textit{égalise} $(\beta, \gamma)$\textit{, i.e.} $u(\alpha')\beta = u(\alpha')\gamma$.
\end{itemize}
\textit{et ces conditions F 1, F 2 impliquent que les $i\bs I'$ sont faiblement filtrantes, et $I$ aussi} (\footnote[57]{Si $I'$ est localement faiblement filtrante (= préfiltrante), alors pour la cofinalité de $u$, il suffit qu'on ait F 0, F 1, F 2, et alors les $i\bs I'$ sont faiblement filtrantes, et $I$ est préfiltrante tout comme $I'$.}).

\printedpage{30}
% head: A. Grothendieck : Dérivateurs, Ch. II, Édition des Universités de Montpellier II et Paris VII
% foot: Version 8 juillet 2001
\item[\textit{c)}] \textit{Si $I'$ est filtrante, alors $u$ est cofinal} si et seulement s'\textit{il satisfait les deux conditions F1, F2, ou encore si et seulement si les $i\bs I'$ sont} filtrantes \textit{(et non seulement} 0-\textit{connexes). Et cela implique que $I'$ est également filtrante} (\footnote[58]{Si $I'$ est localement filtrante, $u$ est cofinal si et seulement si elle satisfait F 0, F 1, F 2, i.e. si et seulement si les $i\bs I'$ sont filtrantes, et alors $I$ est localement filtrante.}).
\item[\textit{d)}] \textit{Si $I$ est préfiltrante (= localement faiblement filtrante, i.e. PS 1), et $u$} pleinement fidèle, \textit{pour que $u$ soit cofinal, il} faut \textit{et il} suffit \textit{que $u$ satisfasse la condition F 1 plus haut. Et alors $I'$ est également préfiltrante, et les $i\bs I'$ sont faiblement filtrantes. Si de plus $I$ est filtrante, les $i\bs I'$ et $I'$ le sont également. (Comparer SGA 4, I 8.1.3.)}
\end{itemize}

\msp{page 47}
\textsc{Démonstration.} a) F 1 est évident, car 0-connexe $\Longrightarrow$ non vide. F 0 car $W_0$-coasphérique $\Longrightarrow$ $W_0$-équivalence, i.e. $\pi_0(u)$ bijectif.

Pour \textit{prouver} b), il suffit de prouver que F 1, F 2 (si $I'$ est faiblement filtrante) impliquent que les $i\bs I'$ sont faiblement filtrantes (a fortiori 0-connexes) et de plus que $I$ l'est également. On sait déjà qu'ils sont non vides par F 1, reste à voir que si $(i', \alpha)$, $(j', \beta)$ sont dans $i\bs I'$, alors ils sont `majorés' par un troisième objet de $i\bs I'$. Comme $I'$ est faiblement filtrante, on trouve
\begin{tikzcd}[row sep=0pt, column sep=small, baseline=(k.base)]
i' \arrow[dr, "r"] & \\
& |[alias=k]| k' \\
j' \arrow[ur, "s"'] &
\end{tikzcd}, d'où $u(r)$, $u(s)$ dans $I$ (cf. diagramme)
\[
\begin{tikzcd}[row sep=small, column sep=small]
& i' \arrow[rr, "r"] \arrow[dd, dash] & & k' \arrow[r, dotted, "t"] \arrow[dd, dash] & l' \arrow[dd, dash] \\
& & j' \arrow[ur, "s"] \arrow[dd, dash] & & \\
& u(i') \arrow[rr, "u(r)"] & & u(k') \arrow[r, dotted, "u(t)"] & u(l') \\
i \arrow[ur, "\alpha"] \arrow[rr, "\beta"'] & & u(j'), \arrow[ur, "u(s)"'] & &
\end{tikzcd}
\]
appliquant F 2 à la double flèche $u(r)\alpha, u(s)\beta : i \rightrightarrows u(k')$, on trouve $k' \xrightarrow{t} l'$ (en pointillés sur \texttt{[le]} diagramme) telle que $(u(t)u(r))\alpha = (u(t)u(s))\beta$, soit $\gamma$, ce qui montre que l'objet $(l', \gamma : i \longrightarrow u(l'))$ majore $(i', \alpha)$ et $(j', \beta)$.

La démonstration que si $I'$ est faiblement filtrante et $u$ satisfait F 1, alors $I$ est faiblement filtrante, est immédiate.

\textit{Prouvons} c). Il reste à prouver que si $I'$ est \textit{filtrante} (non seulement faiblement filtrante), alors si $u$ \texttt{[est]} cofinal, alors F 2 est vérifiée, et de plus les $i\bs I'$ et $I$ sont filtrantes (pas seulement faiblement filtrantes). Mais notons tout de suite que $i\bs I' \longrightarrow I'$ est coétale, donc $I'$ filtrante, donc localement filtrante implique que $i\bs I'$ est aussi localement filtrante\,; comme par hypothèse les $i\bs I'$ sont 0-connexes, ils sont donc filtrantes. La condition F 2 sort facilement à présent :

\printedpage{31}
% head: A. Grothendieck : Dérivateurs, Ch. II, Édition des Universités de Montpellier II et Paris VII
% foot: Version 8 juillet 2001
\msp{page 48}
\[
\begin{tikzcd}[row sep=small, column sep=small]
& i' \arrow[rr, "r"] \arrow[dd, dash] & & j' \arrow[r, dotted, "t"] \arrow[dd, dash] & k' \arrow[dd, dash] \\
& & i' \arrow[ur, "s"] \arrow[dd, dash] & & \\
& u(i') \arrow[rr, "u(r)"] & & u(j') \arrow[r, dotted, "u(t)"] & u(k') \\
i \arrow[ur, "\alpha"] \arrow[rr, "\beta"'] & & u(i'). \arrow[ur, "u(s)"'] & &
\end{tikzcd}
\]
Comme $i\bs I'$ est filtrante, donc faiblement filtrante (PS 1), on peut trouver $i' \dbl{r}{s} j'$ telles que
\[
u(r)\alpha = u(s)\beta,
\]
et comme $I'$ est filtrante, on peut trouver $t : j' \longrightarrow k'$ qui égalise la double flèche $(r, s)$ :
\[
tr = ts(\defeq \tau),
\]
alors
\[
\underbrace{u(t)u(r)}_{u(tr)}\alpha = \underbrace{u(t)u(s)}_{u(ts)}\beta \qquad\quad \text{i.e.} \qquad\quad u(\tau)\alpha = u(\tau)\beta,
\]
i.e. on a vérifié F 2.

Reste à prouver que $I$ est filtrante (si $I'$ l'est et \texttt{[si]} $u : I' \longrightarrow I$ est cofinal). On sait déjà que $I$ est faiblement filtrante par b), reste à vérifier la condition de la double flèche
\[
i \dbl{\alpha}{\beta} j
\]
(PS 2). Soit $i'$ dans $I'$ et $j \longrightarrow u(i')$, et considérons
\[
i \dbl{\gamma\alpha}{\gamma\beta} u(i'),
\]
par hypothèse F 2 $\exists\, r : i' \longrightarrow j'$ dans $I'$ telle que
\msp{page 49}
$u(r)(\gamma\alpha) = u(r)(\gamma\beta)$, donc $u(r)\gamma$ égalise $(\alpha, \beta)$, q.e.d.

La vérification de d) par AQT \texttt{[abréviation dont le sens échappe aux éditeurs]}.

\printedpage{32}
% head: A. Grothendieck : Dérivateurs, Ch. II, Édition des Universités de Montpellier II et Paris VII
% foot: Version 8 juillet 2001
\section*{4\quad Foncteurs biasphériques}

Soit encore
\[
u : I' \longrightarrow I
\]
une flèche de Cat, soit $\cM$ une catégorie, d'où
\[
u^* : \cM(I) \longrightarrow \cM(I');
\tag*{(4.1)}
\]
nous cherchons des conditions sur $u$, $\cM$ pour que ce foncteur soit pleinement fidèle. Si le foncteur adjoint à gauche $u_!$ (resp. adjoint à droite $u_*$) existe, on sait qu'il revient au même de dire que le morphisme d'adjonction
\[
u_!u^*F \longrightarrow F
\tag*{(4.2)}
\]
resp.
\[
u_*u^*F \longleftarrow F
\tag*{(4.2 bis)}
\]
soit \texttt{[un]} isomorphisme pour tout $F$ dans $\cM(I)$. Disons alors que $u$ est $\cM$-connexe, ou $\cM$-biasphérique, cette terminologie sous-entendant que c'est là une propriété qui se réfère au dérivateur $\mathbf{D}_{\cM} = (S \longmapsto \cM(S))$ sur Cat défini par $\cM$. Notons que
\[
\begin{aligned}
&\text{Si $u$ est $\cM$-connexe, alors il est aussi $\cM(S)$-connexe, pour tout $S$ dans Cat,}\\
&\text{puisque $\cM(S)(I) \simeq \cM(I)(S) \defeq \uHom(S^o, \cM(I))$,}
\end{aligned}
\tag*{(4.3)}
\]

\msp{page 50}
et
\[
\text{Si $\cM' \hookrightarrow \cM$ est pleinement fidèle, alors $\cM$-connexe implique $\cM'$-connexe,}
\tag*{(4.4)}
\]
enfin, inversement,
\[
\begin{aligned}
&\text{pour que $u$ soit $\cM$-connexe, il faut et il suffit que pour toute \textit{petite}}\\
&\text{sous-catégorie pleine $\cM'$ de $\cM$, $u$ soit $\cM'$-connexe.}
\end{aligned}
\tag*{(4.5)}
\]
Comme toute petite $\cM'$ se plonge dans $\cM'^\wedge = \operatorname{Ens}(\cM')$, il résulte de (4.3), (4.4), (4.5) que
\[
\text{Si $u$ est Ens-connexe, il est $\cM$-connexe pour toute catégorie $\cM$,}
\tag*{(4.6)}
\]
car c'est le cas si $\cM$ est petite par (4.3), donc pour tout $\cM$ par (4.5).

Disons qu'un foncteur $u$ dans Cat est $\cM$-\textit{cofinal} ou $\cM$-coasphérique si l'homomorphisme fonctoriel
\[
\varprojlim_{I} F \longrightarrow \varprojlim_{I'} u^*(F)
\tag*{(4.7)}
\]
(flèche dans $\uHom(\cM(I), \cM)$ si les $\varprojlim$ de types $I$, $I'$ existent dans $\cM$, sinon flèche dans $\uHom(\cM(I), \cM^\wedge)$) est un isomorphisme. En se ramenant au cas de petites catégories $\cM$, chacune plongée dans $\cM^\wedge$, on voit par le critère (4.2 bis) que
\[
\text{$u$ est $\cM$-connexe si pour tout $i \in \Ob I$,}
\tag*{(4.8)}
\]

\printedpage{33}
% head: A. Grothendieck : Dérivateurs, Ch. II, Édition des Universités de Montpellier II et Paris VII
% foot: Version 8 juillet 2001
\msp{page 51}
le foncteur induit $I'/i \longrightarrow I/i$ est $\cM$-cofinal. La réciproque n'est pas claire, car $u$ est $\cM$-connexe si et seulement si pour tout $i \in \Ob I$ et tout $\cM$-préfaisceau $F_i$ sur $I/i$ \textit{induit par un préfaisceau $F$ sur $I$,} l'homomorphisme $\varprojlim\limits_{I/i} F_i \longrightarrow \varprojlim\limits_{I'/i} (u/i)^*(F/i)$ est \texttt{[un]} isomorphisme. On ne peut donc affirmer d'emblée que cela implique itou pour \textit{tout} $F_i$ sur $I/i$, que si on sait que $(u/i)^* : \cM(I) \longrightarrow \cM(I/i)$ est essentiellement surjectif. Il en est ainsi pour tout $i$ quand $I$ est \textit{ordonné,} car alors $I/i \longrightarrow I$ est pleinement fidèle, donc $\cM(I/i) \longrightarrow \cM(I)$ \texttt{[est]} pleinement fidèle, donc $\cM(I) \longrightarrow \cM(I/i)$ \texttt{[est]} un foncteur de localisation, en particulier essentiellement surjectif.

\msp{page 52}
Mais quand $\cM = (\text{Ens})$, cette perplexité disparaît. On trouve la

\noindent\textbf{Proposition 4.1.} \textit{Soit $u : I' \longrightarrow I$ une flèche dans} Cat. \textit{Conditions équivalentes :}
\begin{itemize}[label={}]
\item[\textit{a)}] \textit{Pour toute catégorie $\cM$, $u^*_{\cM} : \cM(I) \longrightarrow \cM(I')$ est pleinement fidèle (i.e. $u$ est $\cM$-connexe).}
\item[\textit{a$'$)}] \textit{Pour toute catégorie $\cM$ où les limites inductives de type $I$, $I'$ existent, la flèche d'adjonction pour $F$ dans $\cM(I)$} % Transcription corrected: official printed33 reads a prime.
\[
u_!u^*F \xrightarrow{\sim} F
\]
\textit{est} \texttt{[un]} \textit{isomorphisme.}
\item[\textit{a$''$)}] \textit{Dual de a$'$), avec} \texttt{[la]} \textit{flèche d'adjonction}
\[
u_*u^* \xleftarrow{\sim} F. % [sic] u_*u^* -> u_*u^*F, the object F is missing on the left side (compare a') and (4.2 bis))
\]
\item[\textit{b) b$'$) b$''$)}] \textit{Comme a) a$'$) a$''$), avec $\cM =$} (Ens).
\item[\textit{c)}] \textit{Comme b$'$), avec $F$ représentable, i.e. de la forme $j \longmapsto$} $\Hom_I(j, i)$.
\item[\textit{d)}] \textit{Pour tout $I$ dans $I$, $I'/i \longrightarrow I/i$ est cofinal.} % [sic] Pour tout I dans I -> Pour tout i dans I, the variable of the statement is i (compare d'))
\item[\textit{d$'$)}] \textit{Pour tout $i$ dans $I$, $i\bs I' \longrightarrow i\bs I$ est cofinal.}
\item[\textit{d$''$)}] \textit{Pour toute flèche $i_0 \xrightarrow{\alpha} i_1$ dans $I$, la catégorie} $i_0\bs \underset{\alpha}{I'}/i_1 \defeq I' \times_I \underbrace{i_0\bs \underset{\alpha}{I}/i_1}_{\substack{\text{catégorie biïnduite}\\ (i_0, \alpha)\bs(I/i_1) \simeq (i_0\bs I)/(i_1, \alpha)}}$ \qquad \textit{est}

0-\textit{connexe.}
\end{itemize}

\msp{page 53}
\textsc{Démonstration.} On sait déjà, par les préliminaires, l'équivalence de a, a$'$, a$''$, b, b$'$, b$''$, soit A cette condition - nous dirons que $u$ est $W_0$-\textit{connexe} ou $W_0$-\textit{biasphérique,} si ces conditions équivalentes sont satisfaites. On a de plus l'autodualité
\[
\text{$u$ est $W_0$-connexe si et seulement si $u^o$ l'est,}
\tag*{(4.8)} % [sic] (4.8) -> another number, the label (4.8) is already used for a display on p. 32
\]
car on a évidemment
\[
\text{$u$ est $\cM$-connexe si et seulement si $u^o$ est $\cM^o$-connexe.}
\tag*{(4.9)}
\]


\endgroup
% END UNIT Der-ChII-234/Der-ChII-234_p031-035.tex
% BEGIN UNIT Der-ChII-234/Der-ChII-234_p036-040.tex
\begingroup\makeatletter%
\newcommand{\printedpage}[1]{\marginpar{\footnotesize p.~#1}}%
\newcommand{\msp}[1]{\marginpar{\footnotesize [#1]}}%
\newcommand{\Hom}{\operatorname{Hom}}%
\newcommand{\Ob}{\operatorname{Ob}}%
\newcommand{\bD}{\mathbf{D}}%
\newcommand{\bDelta}{\boldsymbol{\Delta}}%
\makeatother\let\printedpage\derpp\let\msp\dermsp%
\let\mspin\DERmspin%


\printedpage{34}
% head: A. Grothendieck : Dérivateurs, Ch. II, Édition des Universités de Montpellier II et Paris VII
Ceci montre l'équivalence de d, d$'$, qui se voit aussi par l'equivalence de chacune de ces deux conditions avec d$''$, par le critère de cofinalité du paragraphe précedent. Soit D cette condition sur $u$. On a par (4.8) D $\Longrightarrow$ A, donc % [sic] l'equivalence -> l'équivalence, accent missing; précedent -> précédent, accent missing
\[
\mathrm{D} \Longrightarrow \mathrm{A} \Longrightarrow \mathrm{c},
\]
il reste à prouver que c $\Longrightarrow$ D, plus précisément c $\Longrightarrow$ d$''$ et c$'$ $\Longrightarrow$ d$''$ \texttt{[c$'$?]}. Explicitons le morphisme d'adjonction (pour $i$ dans $I$)
\[
u_!\underbrace{u^*i} \longrightarrow i,
\]
i.e. pour $j$ dans $I$
\[
\underbrace{(u_!u^*i)(j)}_{\pi_0(j\backslash(I'/u^*(i)))} \longrightarrow i(j) = \Hom(j,i).
\]

\msp{page 54}Or $I'/u^*(i) \simeq I'/i$, donc on a la flèche
\[
(u_!u^*(i))(j) \simeq \pi_0(j\backslash(I'/i)) \longrightarrow \Hom_I(j,i). \tag*{(4.10)}
\]
Or on voit aussitôt que pour $\alpha : j \longrightarrow i$ dans $\Hom(j,i)$, la fibre de (4.10) en $u$ s'identifie à $\pi_0(j\backslash\underset{\alpha}{I'}/i\ )$, de façon plus précise, $j\backslash\underset{\alpha}{I'}/i$\ \ est la sous-catégorie ouverte et fermée de $J = j\backslash(I'/i)$ correspondant à la partie du $\pi_0(J)$ qui est au dessus de $\alpha$ dans (4.10). Cela achève la démonstration de 4.1.

NB Une autre façon de prouver que c $\Longrightarrow$ A c'est de noter que les deux membres dans le morphisme d'adjonction $u_!u^*F \longrightarrow F$ dans $I^{\wedge}$ sont des foncteurs $I^{\wedge} \longrightarrow I^{\wedge}$ qui commutent aux $\varinjlim$ quelconques. Donc pour vérifier que c'est \texttt{[un]} isomorphisme, il suffit de le vérifier pour $F$ dans $I$, puisque $I$ engendre $I^{\wedge}$ par $\varinjlim$ . Par contre, un argument de ce type ne semble pas prouver que pour vérifier que l'on a un \textit{isomorphisme} $F \longrightarrow u_*u^*F$, il suffise de le faire pour $F$ représentable, i.e. de vérifier b$''$ pour $F$ représentable.

Il resterait à voir si c$'$ $\Longrightarrow$ d$''$ \texttt{[c$'$?]}.

\msp{page 55}Explicitons le morphisme d'adjonction
\[
i \longrightarrow u_*u^*i
\]
(pour $i \in \Ob I$) par ses valeurs sur $j$
\[
\underbrace{i(j)}_{\Hom(j,i)} \longrightarrow \Gamma(I'/j, u^*i). \tag*{$(\ast)$}
\]
Soit donc
\[
u_j : I'/j \longrightarrow I,
\]
% foot: Version 8 juillet 2001
\printedpage{35}
% head: A. Grothendieck : Dérivateurs, Ch. II, Édition des Universités de Montpellier II et Paris VII
on a
\[
\begin{aligned}
\Gamma(I'/j, \underbrace{u^*(i)|I'/j}_{u_j^*(i)}) &\simeq \Hom(e_{I'/j}, u^*(i)) \\ % ?? uncertain: extent of the underbrace -- over "u^*(i)|I'/j" or including the closing parenthesis
&\overset{\mathrm{adj}}{\simeq} \Hom(u_{j,!}e_{I'/j}, i),
\end{aligned}
\]
et la flèche $(\ast)$, fonctorielle en $i$, est déduite d'une flèche
\[
u_{j,!}(e_{I'/j}) \longrightarrow j. \tag*{$(\ast\ast)$}
\]
On voit tout de suite, comme ci-dessus, que $(\ast\ast)$ est \texttt{[un]} isomorphisme si et seulement si pour tout $\alpha : k \longrightarrow j$, la catégorie $k\backslash\underset{\alpha}{I'}/j$ est 0-connexe, et si on l'exige pour tout $j$, $k$, $\alpha$, c'est la condition d$''$. Mais il n'est pas clair que pour qu'on ait $(\ast\ast)$ isomorphisme, il suffise que le Hom des deux membres à valeurs dans tout $i$ dans $I$ (et non tout $F$ dans $I^{\wedge}$) soit \texttt{[un]} isomorphisme. La difficulté provient du fait que $u_{j,!}(e_{I'/j})$ n'est pas forcément représentable, et qu'un $F$ quelconque dans $I^{\wedge}$ \msp{page 56}n'est pas forcément $\varprojlim$ de foncteurs représentables. Donc il est douteux que c$'$ \texttt{[c$'$?]} suffise à entraîner que $u$ est $W_0$-connexe.

\textbf{Corollaire 4.2.} \textit{Pour que $u$ soit $W_0$-connexe, il} suffit \textit{qu'il soit une $W_0$-équivalence universelle, i.e. (cf. 3.1) que les fibres soient \textup{0}-connexes, et que pour toute flèche non identique dans $I$, définissant $\bDelta^1 \longrightarrow I$, $I' \times_I \bDelta^1$ soit \textup{0}-connexe.}

4.3. Cette condition de $W_0$-équivalence universelle est cependant \textit{strictement} plus forte que la $W_0$-connexité, comme on le voit sur l'exemple où $u$ est une \textit{équivalence de catégories} non surjective sur les objets (p.ex. l'inclusion d'un sommet dans un groupoïde 1-connexe ayant plus d'un objet). Plus généralement, tout foncteur de localisation est $W_0$-connexe, alors que ses fibres ne sont pas forcément connexes, et a fortiori ce n'est pas une $W_0$-équivalence universelle.

Cependant, si tout endomorphisme $p$ d'un $i$ dans $I$, tel que $p^2 = p$ (projecteur) est l'identité, et si \msp{page 57}$I' \longrightarrow I$ est `transportable', p.ex. si dans $I$ tout isomorphisme est une identité, alors $u$ $W_0$-connexe implique que les fibres de $u$ sont 0-connexes, comme on voit en exprimant d$''$ pour $\alpha$ flèche identique. Les conditions envisagées sur $I$ sont satisfaites si $I$ est ordonnée. Mais même si $I'$ et $I$ sont toutes deux ordonnées, $u$ peut être $W_0$-connexe, sans être une $W_0$-équivalence universelle. Voici un exemple avec $I = \bDelta^2$ :
\[
\begin{array}{r@{\;}l@{\quad}c}
I' = & \left\{\begin{tikzcd}[ampersand replacement=\&]
 \& 1' \arrow[r] \arrow[d] \& 2 \\
0 \arrow[r] \& 1 \&
\end{tikzcd}\right\} & \\[1ex]
 & & \Big\downarrow \\[1ex]
I = & \left\{\begin{tikzcd}[ampersand replacement=\&]
0 \arrow[r] \& 1 \arrow[r] \& 2
\end{tikzcd}\right\}. &
\end{array}
\]
% foot: Version 8 juillet 2001

\printedpage{36}
% head: A. Grothendieck : Dérivateurs, Ch. II, Édition des Universités de Montpellier II et Paris VII
Les catégories bilocalisées de $I$ sont les trois sommets $\{0\}$, $\{1\}$, $\{2\}$, et les trois segments $[0,1]$, $[1,2]$ et $[0,2] = I$, les images inverses sont respectivement
\[
\bullet\ ,\quad
\begin{tikzcd}[row sep=normal] {} \arrow[d] \\ {} \end{tikzcd}\ ,\quad
\bullet\ ,\quad
\begin{tikzcd}[row sep=normal] & {} \arrow[d] \\ {} \arrow[r] & {} \end{tikzcd}\ ,\quad
\begin{tikzcd}[row sep=normal] {} \arrow[r] \arrow[d] & {} \\ {} & \end{tikzcd}\ ,\quad
I'\quad ,
\]
elles sont toutes 0-connexes, i.e. $W_0$-asphériques, et même contractiles donc $W_\omega$-asphériques (donc $u$ est même $W_\omega$-connexe, i.e. $W_\omega$-biasphérique). Mais ce n'est pas une $W_0$-équivalence universelle, car l'image inverse de $I'$ par $\bDelta^1 \xrightarrow{\partial} \bDelta^2$, $\partial(0) = 0$, $\partial(1) = 2$, \msp{page 58}est la catégorie disconnexe $\varepsilon = \{\bullet\ \ \bullet\}$.

J'avais espéré un moment que les foncteurs $W_0$-connexes $I' \xrightarrow{u} I$ seraient ceux qui se factorisent en $I' \xrightarrow{u_0} I'' \xrightarrow{v} I$, avec $u_0$ une équivalence, et $v$ une $W_0$-équivalence universelle. Mais l'exemple précédent le met en défaut, avec $u$ même $W_\omega$-connexe. En effet, $I''$ serait déduit d'un ensemble préordonné, et $v$ se factoriserait par l'ensemble ordonné correspondant, savoir $I'$, donc serait de la forme
\[
I'' \longrightarrow I' \xrightarrow{u} I,
\]
où $I'' \longrightarrow I'$ est un passage au quotient, et l'image inverse de $\bDelta^1 \xrightarrow{\partial} I$ serait encore disconnexe (un $\pi_0$ de cardinal 2). Donc c'est vraiment sans espoir.

Cependant

\textbf{Corollaire 4.3.} \textit{Supposons que $u$ soit $W_0$-lisse, ou $W_0$-propre, ou une $W_0$-fibration faible. Alors $u$ est $W_0$-connexe si et seulement si c'est une $W_0$-équivalence universelle, i.e. si et seulement si a) ses fibres sont \textup{0}-connexes et b) $I' \times_I \bDelta^1$} \texttt{[est]} \textit{\textup{0}-connexe pour tout $\bDelta^1 \longrightarrow I$ (correspondant à une flèche non identique).}

\msp{page 59}Dans les trois cas on sait que $u$ est une $W_0$-équivalence universelle si et seulement si ses fibres sont $W_0$-connexes. Or la $W_0$-lissité implique que $\forall i \in \Ob I$, l'inclusion
\[
I'_i \longrightarrow i\backslash I'
\]
est une $W_0$-équivalence, d'autre part $i\backslash I' \longrightarrow i\backslash I$ étant $W_0$-asphérique, donc une $W_0$-équivalence par hypothèse, $i\backslash I'$ est 0-connexe, donc $I'_i$ aussi. Le cas $u$ $W_0$-propre est dual. Dans le cas \texttt{[où]} $u$ \texttt{[est]} une $W_0$-fibration faible, on sait encore que $I'_i \longrightarrow i\backslash I'$ est une $W_0$-équivalence, et on conclut comme dans le cas lisse (ou, au choix, via $I'_i \longrightarrow I'/i$).

Mais l'idée me vient que les foncteurs $W_0$-connexes pourraient être exactement les foncteurs de localisation ? Il revient au même de dire qu'un foncteur $u : I' \longrightarrow I$ \textit{conservatif} (i.e. $u(\alpha)$ isomorphisme $\Longrightarrow$ $\alpha$ isomorphisme) qui est $W_0$-connexe, est une équivalence. Mais c'est faux, comme on voit en prenant $I = e$, $I'$ un groupoïde 0-connexe et non 1-connexe : $I' \longrightarrow e$ est $W_0$-connexe (et même une $W_0$-équivalence universelle), mais ce n'est pas une équivalence (et pas un foncteur de localisation).
% foot: Version 8 juillet 2001

\printedpage{37}
% head: A. Grothendieck : Dérivateurs, Ch. II, Édition des Universités de Montpellier II et Paris VII
\msp{page 60}Je voudrais généraliser la notion de flèche $W_0$-biasphérique au cas d'un dérivateur $\bD$ quelconque.

\textbf{Proposition 4.4.} \textit{Soit $\bD$ un dérivateur sur} Diag, \textit{et $u : I' \longrightarrow I$ une flèche dans} Diag. \textit{Supposons que $u_*$ et $u_!$ existent, et se calculent de façon standard, pour tout flèche $u$ dans} Diag. \textit{(En fait, il suffit que ce soit le cas pour $u$ et ses localisées $u/i : I'/i \longrightarrow I/i$, ou $u$ et ses colocalisées $i\backslash u : i\backslash I' \longrightarrow i\backslash I$.) Considérons les conditions :} % [sic] pour tout flèche -> pour toute flèche, gender agreement
\begin{itemize}[leftmargin=4em]
\item[\textit{a)}] $u^* : \bD(I) \longrightarrow \bD(I')$ \texttt{[est]} \textit{pleinement fidèle.}
\item[\textit{a$'$)}] \textit{$\forall$ $F$ dans $\bD(I)$, la flèche}
\[
u_!u^*F \longrightarrow F
\]
\textit{est un isomorphisme.}
\item[\textit{a$''$)}] \textit{$\forall$ $F$ dans $\bD(I)$, la flèche}
\[
u_*u^*F \longleftarrow F
\]
\textit{est un isomorphisme.}
\item[\textit{b)}] \textit{$\forall$ $i$ dans $I$, la flèche $i\backslash I' \longrightarrow i\backslash I$ est $\bD$-asphérique.}
\item[\textit{b$'$)}] \textit{$\forall$ $i$ dans $I$, la flèche $I'/i \longrightarrow I/i$ est $\bD$-coasphérique.}
\item[\textit{b$''$)}] \textit{$\forall$ $\alpha : i_0 \longrightarrow i_1$ dans $I$, la catégorie $i_0\backslash\underset{\alpha}{I'}/i_1$ est $\bD$-asphérique.}
\end{itemize}
\textit{On a alors les implications}
\[
(b \Longleftrightarrow b' \Longleftrightarrow b'') \Longrightarrow (a \Longleftrightarrow a' \Longleftrightarrow a'').
\]
\textit{On a $a \Longrightarrow b$ (i.e. toutes les conditions licites envisagées sont équivalentes)} \msp{page 61}\textit{dans chacun des cas suivants.}
\begin{itemize}[leftmargin=4em]
\item[$1^o$] \textit{$I$ est préordonné.}
\item[$2^o$] \textit{$\forall$ $i$ dans $I$, $\alpha_i^* : \bD(I) \longrightarrow \bD(I/i)$ est essentiellement surjectif.}
\item[$3^o$] \textit{$\forall$ $i$ dans $I$, considérant $\alpha_i : I/i \longrightarrow I$, le foncteur $\alpha_{i,!} : \bD(I/i) \longrightarrow \bD(I)$ est conservatif.}
\item[$4^o$] $\mathcal{D} = \mathrm{HOT}_{W_0} \approx \bD_{\mathrm{Ens}}$.
\end{itemize}
\textsc{Démonstration}. L'équivalence de a, a$'$, a$''$ est claire, soit $A$. De même l'equivalence de b, b$'$, b$''$, soit $B$. % [sic] l'equivalence -> l'équivalence, accent missing

Prouvons
\[
B \Longrightarrow A
\]
en supposant p.ex. l'existence de $u_*$. On admet tacitement, avec l'existence de $u_*$ : $\bD(I') \longrightarrow \bD(I)$, le mode de calcul des fibres de $u_*(-)$. Donc $A$, via a$''$, équivaut à
\[
H^{\bullet}_{\bD}(I'/i, u^*F) \xleftarrow{\mathit{iso}} H^{\bullet}_{\bD}(I/i, F) \simeq F(i)
\]
% foot: Version 8 juillet 2001
\printedpage{38}
% head: A. Grothendieck : Dérivateurs, Ch. II, Édition des Universités de Montpellier II et Paris VII
pour tout $F$ dans $\bD(I)$, tout $i$ dans $I$.

Soit
\[
\alpha_i : I/i \longrightarrow F, \qquad F_i = \alpha_i^*(F), % [sic] I/i \longrightarrow F -> I/i \longrightarrow I, alpha_i is defined as I/i -> I in 3^o
\]
on doit donc exprimer que pour tout $F_i$ provenant d'un $F$, on a
\[
H^{\bullet}_{\bD}(I/i, F_i) \xrightarrow{\mathrm{iso}} H^{\bullet}_{\bD}(I'/i, (u/i)^*(F_i)). \tag*{$(\ast)$}
\]
Or sous réserve d'existence de $(u/i)_!$, dire que $(\ast)$ est \texttt{[un]} isomorphisme pour \textit{tout} $F_i$ de $\bD(I/i)$, signifie que
\[
u/i : I'/i \longrightarrow I/i
\]
\msp{page 62}est $\bD$-coasphérique, d'où b$'$ $\Longrightarrow$ a$''$. Il est clair que l'implication inverse est valable si $\alpha_i^*$ est essentiellement surjectif, c'est le cas $2^o$. Or $2^o$ est satisfait si $I$ est préordonné, d'où le cas $1^o$. Le cas $4^o$ a déjà été traité (4.1). Il faut encore traiter le cas $3^o$, qui contient d'ailleurs le cas $2^o$, donc le cas $1^o$ (et aussi le cas $4^o$, cf. plus bas).

Il faut exprimer l'isomorphie de la flèche $(\ast)$ pour $F_i$ de la forme $\alpha_i^*(F)$. Soit $\xi$ \texttt{[un]} objet de $\mathcal{A} = \mathcal{D}(e)$, il faut exprimer l'isomorphie de $\Hom(\xi, -)$ dans les deux membres de $(\ast)$, i.e. de
\[
\begin{aligned}
\Hom(I/i, F_i) \quad &\longrightarrow \quad \Hom(\xi'_{I'/i}, (u/i)^*(F_i)) \\
&\overset{\text{adj.}}{\simeq} \quad \Hom((u/i)_!(\xi'_{I'/i}), F_i),
\end{aligned}
\tag*{$(\ast\ast)$}
\]
\texttt{[plutôt} $\xi_{I'/i}$\texttt{]} flèche qui se déduit de
\[
(u/i)_!(u/i)^*(\xi_{I'/i}) \xrightarrow{\text{adj.}} \xi_{I'/i} \tag*{$(\ast\ast\ast)$}
\]
en prenant le Hom dans $F_i$. La condition $u/i$ coasphérique (i.e. b$'$) signifie que $(\ast\ast\ast)$ est un isomorphisme pour tout $\xi$ dans $\mathcal{A}$. \textit{Suffit}-il pour ceci que son Hom dans un $F_i$ soit bijectif, quand on restreint $F_i$ dans $\bD(I/i)$ à être dans l'image de $\bD(I)$ par $\alpha_i^*$ ? Or cette flèche $(\ast\ast)$, pour $F_i = \alpha_i^*F$, s'explicite par la formule d'adjonction relative \msp{page 63}au foncteur $\alpha_i$, comme
\[
\Hom(\alpha_{i,!}(\xi_{I/i}), F) \longrightarrow \Hom(\alpha_{i,!}(\ldots), F).
\]
Dire que c'est un isomorphisme pour tout $F$, signifie donc \textit{exactement} que la transformée de la flèche $(\ast\ast\ast)$ par $\alpha_!$ est un isomorphisme. Pour en conclure que $(\ast\ast\ast)$ elle-même est \texttt{[un]} isomorphisme, la condition naturelle à imposer est donc que $\alpha_{i,!} : \bD(I/i) \longrightarrow \bD(I)$ soit conservatif. C'est le cas $3^o$. Cela achève la démonstration de 4.4.

4.5. Montrons que la condition $3^o$ sur $I$ n'est pas toujours satisfaite, dans le cas de $\bD = (\mathrm{HOT}_{W_0})^{\mathrm{opp}} \approx (\bD_{\mathrm{Ens}})^{\mathrm{opp}}$, bien que dans ce cas, il soit pourtant vrai que pour toute $I'$ sur
% foot: Version 8 juillet 2001


\endgroup
% END UNIT Der-ChII-234/Der-ChII-234_p036-040.tex
% BEGIN UNIT Der-ChII-234/Der-ChII-234_p041-045.tex
\begingroup\makeatletter%
\newcommand{\printedpage}[1]{\marginpar{\footnotesize p.~#1}}%
\newcommand{\msp}[1]{\marginpar{\footnotesize [#1]}}%
\newcommand{\D}{\mathbf{D}}%
\makeatother\let\printedpage\derpp\let\msp\dermsp%
\let\mspin\DERmspin%


\printedpage{39}
% head: A. Grothendieck : Dérivateurs, Ch. II, Édition des Universités de Montpellier II et Paris VII

$I$ satisfaisant la condition a$''$ (autoduale), la condition b$'$ (elle aussi autoduale) est aussi satisfaite (par cas $4^o$ ci-dessus). Je dis donc qu'il existe $I$ et $i$ dans $I$, tels que
\[
\beta_{i,*} : (i\backslash I)^\wedge \longrightarrow I^\wedge
\]
ne soit pas conservatif. (\textbf{NB} Par contre $\alpha_{i,!} : (I/i)^\wedge \longrightarrow I^\wedge$ est toujours conservatif, ce qui redonne le cas $4^o$ comme cas particulier de $3^o$). Comme $\beta_i : i\backslash I \longrightarrow I$ est $W_0$-propre (et même $W_\omega$-propre !), $\beta_{i,*}$ `se calcule fibre par fibre', donc on a
\msp{page 64}
\[
\beta_{i,*}(G)(j) \simeq \prod_{\alpha\in\mathrm{Hom}(i,j)} G(\alpha),
\]
la fibre en question en $j$ étant discrète. Prenons
\[
I = 0 \underset{\beta}{\overset{\alpha}{\rightrightarrows}} 1, \qquad\quad \text{avec } \alpha \neq \beta,
\]
alors $0\backslash I$ est $\Psi = {}$\begin{tikzcd}[ampersand replacement=\&,row sep=0.1em, column sep=1.2em, baseline=(\tikzcdmatrixname-2-1.base)]
 \& 1 \\
0 \ar[ur, "\alpha'"] \ar[dr, "\beta'"'] \& \\
 \& 1'
\end{tikzcd}, et $\Psi^\wedge$ est la catégorie des diagrammes d'ensembles $G = {}$\begin{tikzcd}[ampersand replacement=\&,row sep=0.1em, column sep=1.2em, baseline=(\tikzcdmatrixname-2-1.base)]
 \& G_1 \ar[dl, "\lambda"'] \\
G_0 \& \\
 \& G_{1'} \ar[ul, "\mu"]
\end{tikzcd}. Le foncteur $\beta_*$ déduit de $\beta : i\backslash I \longrightarrow I$ transforme ce diagramme en
\[
G_1 \times G_{1'} \underset{\lambda\,\mathrm{pr}_2}{\overset{\lambda\,\mathrm{pr}_1}{\rightrightarrows}} G_0.
\]
\texttt{[plutôt} $\mu\,\mathrm{pr}_2$\texttt{]}. Si on se borne aux diagrammes $G$ pour lesquels $G_{1'} = \emptyset$, $\beta_*$ sur ceux-ci donne
\[
\emptyset \rightrightarrows G_0,
\]
donc la restriction de $\beta_*$ à cette sous-catégorie (équivalente à $(\boldsymbol{\Delta}^1)^\wedge$) n'est pas fidèle, ni conservative, ne tenant compte que de la composante $G_0$, et non pas de $G_1$.

\msp{page 63bis}

\textbf{Definition 4.6.} \textit{On dit que $u : I' \longrightarrow I$ dans} Diag \textit{est} $\D$-biasphérique \textit{si elle satisfait les conditions équivalentes b, b$'$, b$''$ de 4.4,} $\D$-connexe \textit{si elle satisfait la condition a, ou encore (sous réserve d'existence de $u_!$ resp. de $u_*$) les conditions équivalentes a$'$, a$''$.} % [sic] Definition -> Définition, accent missing in print

Ainsi, $\D$-biasphérique implique $\D$-connexe, et la réciproque est vraie quand $I$ satisfait (par rapport à $\D$) la condition $3^o$ de 4.4 (donc aussi, si elle satisfait l'une des conditions plus faibles $2^o$, $1^o$), ou la condition duale de celle-ci (vu que les notions $\D$-biasphérique, $\D$-connexe sont autoduales). À cet égard, notons le

\textbf{Corollaire 4.7.} \textit{Supposons que $\mathcal{A} = \D(e)$ satisfasse la condition suivante : le foncteur}
\[
\left\{
\begin{array}{lcccl}
\mathcal{A} & \times & \mathcal{A} & \longrightarrow & \mathcal{A} \\[1ex]
(\xi & , & \eta) & \longrightarrow & \xi \amalg \eta
\end{array}
\right.
\]
\textit{est conservatif (resp. fidèle). Alors pour toute flèche étale $u : J \longrightarrow I$ dans} Diag \textit{(et en particulier pour $I/i \longrightarrow I$), le foncteur $u_! : \D(J) \longrightarrow \D(I)$ est conservatif (resp. fidèle)} (\footnote{\textbf{NB} Cette condition est satisfaite si $\forall\, \xi, \eta \in \mathcal{A}$ on a $\mathrm{Hom}(\xi,\eta) \neq \emptyset$, (car on se ramène à ceci : si $u_1$,
% footnote continues at the foot of p. 40
$u_2$ sont deux flèches dans Ens entre ensembles \textit{non vides,} et si $u_1 \times u_2$ est \texttt{[un]} isomorphisme, $u_1$, $u_2$ sont \texttt{[des]} isomorphismes). Cela s'applique en particulier si $\mathcal{A}$ admet un objet neutre, p.ex. si $\mathcal{A}$ est additive.}). \textit{Si cette condition, ou la condition duale,}
\printedpage{40}
% head: A. Grothendieck : Dérivateurs, Ch. II, Édition des Universités de Montpellier II et Paris VII
\msp{page 64bis}
\textit{impliquant les produits $\xi\times\eta$ ($\xi$, $\eta$ dans $\mathcal{A}$), est satisfaite, alors pour toute flèche $u : I' \longrightarrow I$ dans} Diag, \textit{$u$ est $\D$-connexe si et seulement si $u$ est $\D$-biasphérique.}

\textsc{Démonstration.} L'hypothèse de conservativité, ou de fidélité se transporte aux sommes indexées par un ensemble quelconque (sous réserve d'existence), comme on voit en écrivant $\coprod\limits_{\lambda} \xi_\lambda = \xi_{\lambda_0} \amalg \coprod\limits_{\lambda\neq\lambda_0} \xi_\lambda$. L'assertion faite sur $u_!$ résulte alors du calcul de $u_!$ fibre par fibre (qui est valable vu que $u$ est étale donc lisse), en tenant compte que les fibres sont discrètes.

Je ne connais pas d'exemple d'un couple $(\D, u)$ d'un dérivateur $\D$ et d'une flèche $u$ dans Diag, pour lesquels $u$ soit $\D$-connexe sans être $\D$-asphérique. On a donné des conditions sur $\D$, qui assurent équivalence des deux notions. Si $\D$ est de la forme $\D_{\mathcal{M}}$, et si $\mathcal{M}$ satisfait la condition dite dans 4.7 ou sa duale, alors par 4.7 c'est OK. Donc un contrexemple avec un $\D_{\mathcal{M}}$ devrait en tenir compte. Je ne vais pas m'obstiner pour tirer au clair ce qui en est.

\msp{page 65}

\textbf{Corollaire 4.8.} \textit{Une flèche $u$ de} Diag \textit{qui est une $\D$-équivalence universelle est $\D$-biasphérique, et \textup{a fortiori} $\D$-connexe. Inversement, si $u$ est soit $\D$-lisse, soit $\D$-propre, soit une $\D$-fibration faible, et si $u$ est $\D$-biasphérique, c'est une $\D$-équivalence universelle.}

La démonstration est celle du cas particulier 4.3. On a donc un diagramme d'implications :
\[
\begin{tikzcd}[column sep=small, row sep=large]
\text{Cat-équivalence} \ar[dr, Rightarrow] & \D\text{-équivalence universelle } (W_{\mathrm{univ}}) \ar[d, Rightarrow, "{\text{\footnotesize\begin{tabular}{@{}l@{}}implication stricte ; équivalence\\ dans le cas où $u$ est soit $\mathbf{D}$-propre,\\ soit $\mathbf{D}$-lisse, soit une $\mathbf{D}$-fibration faible.\end{tabular}}}"] & \\
 & \D\text{-biasphérique} \ar[dl, Rightarrow] \ar[d, Rightarrow, "{\text{\footnotesize\begin{tabular}{@{}l@{}}équivalence\\ dans tous les\\ cas connus\end{tabular}}}"] \ar[dr, Rightarrow] & \\
\D\text{-asphérique} \ar[dr, Rightarrow] & \D\text{-connexe} \ar[d, Rightarrow] & \D\text{-coasphérique} \ar[dl, Rightarrow] \\
 & \D\text{-équivalences } (W) &
\end{tikzcd}
\]
(\footnote{\textbf{NB} Ces sept classes de morphismes dans Cat sont stables par composition sauf la notion de $\D$-connexité, toutes ces notions ne dépendent que de l'ensemble $W = W_{\D}$ des $\D$-équivalences.}). Rappelons que pour $u : I' \longrightarrow I$ donné
\begin{enumerate}[leftmargin=*, labelindent=1.5em]
\item[a)] $u$ est $\D$-asphérique si et seulement si $\forall\, F$ \textit{constant} dans $\D(I)$
\[
F \xrightarrow[\sim]{\text{adj.}} u_*u^*F \qquad\quad \text{est \texttt{[un]} isomorphisme.} \tag*{$(\ast)$}
\]
\printedpage{41}
% head: A. Grothendieck : Dérivateurs, Ch. II, Édition des Universités de Montpellier II et Paris VII
\item[b)] $u$ est $\D$-coasphérique si et seulement si $\forall\, F$ \textit{constant} dans $\D(I)$
\[
F \xleftarrow[\sim]{\text{adj.}} u_!u^*F \qquad\quad \text{est \texttt{[un]} isomorphisme.} \tag*{$(\ast\ast)$}
\]
\item[c)] $u$ est $\D$-connexe si et seulement si pour \textit{tout} $F$ dans $\D(I)$ (pas forcément constant, ni localement constant), $(\ast)$ est \texttt{[un]} isomorphisme, ou ce qui revient au même, pour tout $F$ $(\ast\ast)$ est \texttt{[un]} isomorphisme.
\end{enumerate}

\msp{page 66}

Rappelons que la classe $W_{\mathrm{asph}}$ \texttt{[des]} flèches $W$-asphériques et celle, $\mathrm{Liss}_W$, des flèches $W$-lisses, se déterminent mutuellement, ainsi
\[
\left\{
\begin{array}{lcl}
(f : X \longrightarrow Y) \in W_{\mathrm{asph}} & \Longleftrightarrow & \forall\, (Y' \xrightarrow{\ u\ } Y) \in \mathrm{Liss}_W, \\
 & & X' = X \times_Y Y' \longrightarrow X' \text{ est } \in W. % [sic] X' = X \times_Y Y' \longrightarrow X' -> X' = X \times_Y Y' \longrightarrow Y', the target must be Y' (compare (4.12))
\\
(f : Y' \longrightarrow Y) \in \mathrm{Liss}_W & \Longleftrightarrow & \text{Le foncteur } X \longmapsto X' = X \times_Y Y' \\
 & & \text{de } \mathrm{Cat}/Y \text{ dans } \mathrm{Cat}/Y' \text{ transforme} \\
 & & \text{flèches dans } W_{\mathrm{asph}} \text{ dans \texttt{[plutôt en]}} \\
 & & \text{flèches dans } W_{\mathrm{asph}}.
\end{array}
\right.
\tag*{(4.10)}
\]

Pour toute partie $A$ de Fl(Cat) stable par isomorphismes, on peut associer une partie $\mathbf{L}(A)$ de Fl(Cat), par la relation
\[
\begin{array}{rcl}
(u : Y' \longrightarrow Y) \in \mathbf{L}(A) & \Longleftrightarrow & u^* : \mathrm{Cat}/Y \longrightarrow \mathrm{Cat}/Y' \\
 & & \text{transforme flèches} \\
 & & \in A \text{ en flèches } \in A.
\end{array}
\tag*{(4.11)}
\]

On voit que $\mathbf{L}(A)$ contient les isomorphismes, est stable par composition, et est stable aussi par changement de base (comme c'est le cas de $\mathrm{Liss}_W$). Inversement, à une partie $L$ de Fl(Cat), stable par isomorphismes, on associe une partie $\mathbf{A}_W(L)$ par
\[
\begin{array}{rcl}
(f : X \longrightarrow Y) \in \mathbf{A}_W(L) & \Longleftrightarrow & \forall\, (Y' \longrightarrow Y) \text{ dans } L, \\
 & & X' = X \times_Y Y' \longrightarrow Y' \\
 & & \text{est dans } W.
\end{array}
\tag*{(4.12)}
\]

Si $L$ est stable par changement de base, alors $\mathbf{A}_W(L)$ est stable par composition. Il contient les isomorphismes (car $W$ les contient), est contenu dans $W$ si $L$ contient les identités. Mais
\msp{page 67}
aucune raison qu'il soit stable par changement de base quelconque (et en fait, $W_{\mathrm{asph}}$ ne l'est pas).

(\footnote{Canulé \texttt{[jusqu'à la fin du} \S4\texttt{]}, cf. \S6.}). On a, pour un $A \subseteq W \subseteq \mathrm{Fl(Cat)}$ donné, stable par isomorphismes,
\[
A \subseteq \mathbf{A}_W(\mathbf{L}(A)) \overset{\text{déf}}{=} A^\natural \subseteq W, \qquad \text{\footnotesize\begin{tabular}{@{}l@{}}partie stable par composition,\\ contient \texttt{[les]} isomorphismes,\end{tabular}}
\tag*{(4.13)}
\]
\printedpage{42}
% head: A. Grothendieck : Dérivateurs, Ch. II, Édition des Universités de Montpellier II et Paris VII
et pour toute partie $L$ de Fl(Cat) qui est stable par changement de base, et par composition
\[
L \subseteq \mathbf{L}(\mathbf{A}_W(L)) \overset{\text{déf}}{=} L^\flat, \qquad \text{\footnotesize\begin{tabular}{@{}l@{}}partie de Fl\,Cat stable\\ par composition et par changement\\ de base, contenant les isomorphismes.\end{tabular}}
\tag*{(4.14)}
\]
On a
\[
\left\{
\begin{array}{lcl}
A^{\natural\natural} & = & A^\natural, \quad \text{car } \mathbf{A}_W(L)^\natural = \mathbf{A}_W(L) \\[1ex]
L^{\flat\flat} & = & L^\flat, \quad \text{car } \mathbf{L}(A)^\flat = \mathbf{L}(A).
\end{array}
\right.
\]
Exemples :
\begin{enumerate}[leftmargin=*, labelindent=1em]
\item[$1^o$)] $L = \text{Ét}$, $\mathbf{A}_W(L) = W_{\mathrm{asph}}$, $\mathbf{L}(\mathbf{A}_W(L)) = L^\flat = \mathrm{Liss}_W$.
\item[$2^o$)] $L = \text{Coét}$, $\mathbf{A}_W(L) = W_{\mathrm{coas}}$, $\mathbf{L}(\mathbf{A}_W(L)) = L^\flat = \mathrm{Prop}_W$.
\item[$3^o$)] $A = W$, $\mathbf{L}(A) = \mathrm{Fib}_W$, $A^\natural = W$.
\item[$4^o$)] $L = \mathrm{Fl\,Cat}$, $\mathbf{A}_W(L) = W_{\mathrm{univ}}$, $L^\flat = L$.
\end{enumerate}
\msp{page 68 vide}

\par
\printedpage{43}
% head: A. Grothendieck : Dérivateurs, Ch. II, Édition des Universités de Montpellier II et Paris VII
\msp{page 69}

\section*{5\quad Dualité formelle}

Soient $\mathcal{P}$, $\mathcal{Q}$ deux ensembles ordonnés. J'essaye de dégager une notion de `dualité' (formelle) reliant $\mathcal{P}$ et $\mathcal{Q}$. Ce sera une relation entre $\mathcal{P}$ et $\mathcal{Q}$
\[
\mathcal{R} \subseteq \mathcal{P} \times \mathcal{Q}
\tag*{(5.1)}
\]
ayant les propriétés suivantes
\begin{enumerate}[leftmargin=*, labelindent=0.5em]
\item[(D)] Pour toute partie $A$ de $\mathcal{P}$, l'ensemble
\end{enumerate}
\[
A^\cup \overset{\text{déf}}{=} \{ q \in \mathcal{Q} \mid (p,q) \in \mathcal{R} \text{ pour tout } p \in A,\ \text{ i.e. } A \times \{q\} \subseteq \mathbf{R} \} % [sic] \mathbf{R} -> \mathcal{R}, the relation is written \mathcal{R} everywhere else
\]
est une partie ouverte de $\mathcal{Q}$ ayant un plus grand élément $\beta(A) : A^\cup = \mathcal{Q}/\beta(A)$. Dualement, pour toute partie $B$ de $\mathcal{Q}$, l'ensemble
\[
B^\cap \overset{\text{déf}}{=} \{ p \in \mathcal{P} \mid (p,q) \in \mathcal{R} \text{ pour tout } q \in B,\ \text{ i.e. } \{p\} \times B \subseteq \mathcal{R} \}
\]
est une partie ouverte de $\mathcal{P}$ ayant un plus grand élément $\alpha(B)$, i.e. $B^\cap = \mathcal{P}/\alpha(B)$.

Le plus souvent, $\mathcal{P}$ et $\mathcal{Q}$ seront stables par sup quelconques (donc aussi par inf quelconques). Dans ce cas, (D) équivaut, du côté des $A^\cup$ (pour $A \in \mathfrak{P}(\mathcal{P})$) à ceci :
\begin{enumerate}[leftmargin=*, labelindent=0em]
\item[(D$_0$)] Pour tout $p \in \mathcal{P}$, $\mathcal{R}(p) = \{ q \in Q \mid (p,q) \in \mathcal{R} \}$ est une partie ouverte \textit{filtrante} de $\mathcal{Q}$, ou ce qui revient au même, une partie non vide ($\Longleftrightarrow$ contenant le plus petit élément $\emptyset_{\mathcal{Q}}$ de $\mathcal{Q}$) stable par sup de deux éléments. De plus, $\mathcal{R}(p)$ est stable par sup filtrants, ou ce qui revient au même (étant lui-même filtrant), il contient un plus grand élément. Il revient % [sic] Q -> \mathcal{Q}, plain italic Q printed in "q \in Q" instead of the script \mathcal{Q}
\msp{page 70}
au même de dire que $\mathcal{R}(p)$ est un \textit{ouvert} de $\mathcal{Q}$ stable par sup quelconques (ce qui implique qu'il contient $\emptyset_{\mathcal{Q}}$). Dualement, pour tout $q \in \mathcal{Q}$, ${}^{s}\mathcal{R}(q) = \{ p \in \mathcal{P} \mid (p,q) \in \mathcal{R} \}$ est une partie \textit{ouverte} de $\mathcal{P}$, stable par sup quelconques, i.e. une partie ouverte ayant un plus grand élément. % doubt resolved against PDF leaf 15: raised s, transpose relation
\end{enumerate}

En d'autres termes, quand on postule la condition (D) pour $A$ resp. $B$ réduit à un élément, alors elle reste vérifiée pour toute partie. Cela provient des formules
\[
\left\{
\begin{array}{lcl}
(\bigcup_i A_i)^\cup & = & \bigcap_i A_i^\cup \\[1ex]
(\bigcup_j B_j)^\cap & = & \bigcap_j B_j^\cap,
\end{array}
\right.
\tag*{(5.2)}
\]
et du fait que l'intersection de parties ouvertes de la forme $\mathcal{P}/\alpha_i$ dans $\mathcal{P}$ est de la forme $\mathcal{P}/\alpha$, où $\alpha = \inf_i \alpha_i$, et de même dans $\mathcal{Q}$. Ainsi on trouve
\[
\left\{
\begin{array}{lcl}
\beta(\bigcup_i A_i) & = & \inf_i \beta(A_i) \\[1ex]
\alpha(\bigcup_j B_j) & = & \inf_j \alpha(B_j).
\end{array}
\right.
\tag*{(5.3)}
\]


\endgroup
% END UNIT Der-ChII-234/Der-ChII-234_p041-045.tex
% BEGIN UNIT Der-ChII-234/Der-ChII-234_p046-050.tex
\begingroup\makeatletter%
\newcommand{\printedpage}[1]{\marginpar{\footnotesize p.~#1}}%
\newcommand{\msp}[1]{\marginpar{\footnotesize [#1]}}%
\newcommand{\cP}{\mathcal{P}}%
\newcommand{\cQ}{\mathcal{Q}}%
\newcommand{\cR}{\mathcal{R}}%
\newcommand{\fP}{\mathfrak{P}}%
\newcommand{\sR}{{}^{s}\mathcal{R}}%
\newcommand{\circled}[1]{\tikz[baseline=(c.base)]\node[draw,rounded corners=2pt,inner sep=1.5pt](c){#1};}%
\makeatother\let\printedpage\derpp\let\msp\dermsp%
\let\mspin\DERmspin%


\printedpage{44}
% head: A. Grothendieck : Dérivateurs, Ch. II, Édition des Universités de Montpellier II et Paris VII

On posera, pour une partie $A$ de $\cP$ resp. $B$ de $\cQ$
\[
A^{\natural} = (A^{\cup})^{\cap}, \qquad B^{\flat} = (B^{\cap})^{\cup}
\tag*{(5.4)}
\]
de sorte qu'on aura
\[
A \subseteq A^{\natural}, \qquad B \subseteq B^{\flat}
\tag*{(5.5)}
\]
On dira que $A$ est $\cR$-\textit{clos} si $A = A^{\natural}$ (ou encore si $A$ est de la forme $B'^{\cap}$ pour $B' \subseteq \cQ$, cf. 5.6), que
\msp{71}
$B$ est $\cR$-\textit{clos} si $B = B^{\flat}$. On aura pour tout $A \subseteq \cP$, $B \subseteq \cQ$
\[
\left\{
\begin{array}{rcll}
(A^{\cup})^{\flat} & = & A^{\cup}, & \text{i.e. } A^{\cup} \text{ est } \cR\text{-clos,}\\[1ex]
(B^{\cap})^{\natural} & = & B^{\cap}, & \text{i.e. } B^{\cap} \text{ est } \cR\text{-clos.}
\end{array}
\right.
\tag*{(5.6)}
\]
En particulier
\[
(A^{\natural})^{\natural} = A^{\natural}, \qquad (B^{\flat})^{\flat} = B^{\flat},
\tag*{(5.7)}
\]
i.e. la $\cR$-clôture $A^{\natural}$ de $A \subseteq \cP$ est $\cR$-clos, et de même la $\cR$-cloture $B^{\flat}$ de $B \subseteq \cQ$. % [sic] cloture -> clôture, accent circonflexe manquant (cf. « clôture » dans la même ligne)

En fait, je m'aperçois que je suis en train de mélanger deux espèces de faits différents. Les relations (5.2), et (5.4) à (5.7), ne dépendent pas de la donné de relations d'ordre sur $\cP$, $\cQ$, ni d'axiomes reliant celles-ci à la relation donnée $\cR$, mais seulement de cette relation, laquelle n'a pas à être soumise à aucune condition. Notons que % [sic] la donné -> la donnée, accord manquant
\[
\emptyset^{\cup} = \cQ, \qquad \emptyset^{\cap} = \cP,
\tag*{(5.8)}
\]
tandis que
\[
\emptyset^{\natural} = \cQ^{\cap}, \qquad \emptyset^{\flat} = \cP^{\cup}
\tag*{(5.9)}
\]
sont respectivement la plus petite partie $\cR$-close de $\cP$ et de $\cQ$.

L'application
\[
A \longmapsto A^{\natural}
\tag*{(5.10)}
\]
est une application \textit{croissante} (et \textit{idempotente}) de $\fP(\cP)$ dans lui-même, ayant donc la propriété
\msp{72}
\[
(\bigcup_i A_i)^{\natural} = (\bigcup_i A_i^{\natural})^{\natural}.
\]

\printedpage{45}
% head: A. Grothendieck : Dérivateurs, Ch. II, Édition des Universités de Montpellier II et Paris VII

Pour tout $A \subseteq \cP$, $A^{\natural}$ est la plus petite partie $\cR$-close de $\cP$ contenant $A$, et de même pour tout $B \subseteq \cQ$, $B^{\flat}$ est la plus petite partie $\cR$-close de $Q$ contentant $B$. On a de plus (c'était même à dire avant de dire que $A \longmapsto A^{\natural}$ est croissante) % [sic] de $Q$ -> de $\mathcal{Q}$, lettre droite au lieu de la lettre calligraphique ; [sic] contentant -> contenant, faute de frappe
\[
\left\{
\begin{array}{rcl}
A \subseteq A' & \Longrightarrow & A'^{\cup} \subseteq A^{\cup}\\[1ex]
B \subseteq B' & \Longrightarrow & B'^{\cap} \subseteq B^{\cap},
\end{array}
\right.
\tag*{(5.11)}
\]
i.e. les applications
\[
\begin{tikzcd}
\fP(\cP) \arrow[r, shift left, "\beta"] & \fP(\cQ) \arrow[l, shift left, "\alpha"]
\end{tikzcd}
\qquad
\begin{array}{l}
A \longmapsto A^{\cup}\\
B \longmapsto B^{\cap}
\end{array}
\tag*{(5.12)}
\]
sont décroissantes (ce qui implique que les deux composés $\alpha\beta = (A \longmapsto A^{\natural})$ et $\beta\alpha = (B \longmapsto B^{\flat})$ sont croissantes). Leurs restrictions aux sous-ensembles $\fP_{\cR}(\cP)$ et $\fP_{\sR}(\cQ)$ des parties $\cR$-closes de $\cP$ resp. de $\cQ$, sont, par définition de celles-ci comme
\[
\begin{array}{rcl}
\fP_{\cR}(\cP) & = & \{A \subseteq \cP \mid A^{\natural} \ (\overset{\text{déf}}{=} \alpha\beta(A)) = A\}\\
\fP_{\sR}(\cQ) & = & \{B \subseteq \cQ \mid B^{\flat} \ (\overset{\text{déf}}{=} \beta\alpha(B)) = B\},
\end{array}
\]
inverse l'une de l'autre. On trouve ainsi

\textbf{Proposition 5.1.} \textit{Soient $\cP$, $\cQ$ deux ensembles, $\cR \subseteq \cP \times \cQ$ une relation. Alors les applications $\alpha$, $\beta$ de (5.12) induisent des bijections inverses l'}
\msp{73}
\textit{une de l'autre, entre l'ensemble $\mathit{Cl}_{\cR}(\cP)$ des parties $\cR$-closes de $\cP$, et l'ensemble $\mathit{Cl}_{\sR}(\cQ)$ des parties ${}^{s}\cR$-closes (nous dirons $\cR$-closes par abus de langage) de $\cQ$. Ces bijections sont décroissantes, donc des antiisomorphismes pour les relations d'ordre naturelles (d'inclusion).}

\textbf{Corollaire 5.2.} \textit{Soit $(A_i)_i$ une famille de parties $\cR$-closes de $\cP$. Alors $\bigcap A_i$ est une partie $\cR$-close de $\cP$, et}
\[
(\bigcap_i A_i)^{\cup} = \underbrace{(\bigcup_i A_i^{\cup})^{\flat}}_{\substack{\text{sup des}\\ \text{parties } \cR\text{-closes}\\ A_i^{\cup} \text{ de } \cQ}} \quad .
\tag*{(5.2.1)}
\]

Mais la relation $\cR$ se `prolonge' en une relation $\bar{\cR}$ entre
\[
\bar{\cP} = \fP(\cP), \qquad \bar{\cQ} = \fP(\cQ)
\tag*{(5.13)}
\]
par la prescription
\[
\underbrace{\bar{\cR}}_{\subseteq \bar{\cP} \times \bar{\cQ}} = \{(A,B) \in \bar{\cP} \times \bar{\cQ} \mid B \subseteq A^{\cup},\ \text{i.e. } A \subseteq B^{\cap},\ \text{i.e. } A \times B \subseteq \cR\}.
\tag*{(5.14)}
\]
Les ensembles $\bar{\cP}$, $\bar{\cQ}$ sont stables par sup quelconques (réunion de parties de $\cP$ resp. $\cQ$), et la relation $\bar{\cR}$ sur ces ensembles ordonnés satisfait à la condition (D$_0$)
\printedpage{46}
% head: A. Grothendieck : Dérivateurs, Ch. II, Édition des Universités de Montpellier II et Paris VII
\msp{74}
ci-dessus (p. 69), de façon évidente : pour $A$ fixé, les $B$ tels que $(A,B) \in \bar{\cR}$ sont ceux qui sont contenus dans $A^{\cup}$, donc l'ensemble $\bar{\cR}(A)$ n'est autre que $\bar{\cQ}/A^{\cup}$, et dualement ${}^{s}\bar{\cR}(B)$ n'est autre que $\bar{\cP}/B^{\cap}$. Ainsi
\[
\beta(A) = A^{\cup}, \qquad \alpha(B) = B^{\cap}.
\tag*{(5.14)} % [sic] (5.14) -> numéro distinct, numéro déjà employé pour la formule (5.14) de la page précédente
\]
Si j'ai songé à introduire l'axiome (D$_0$), pour un couple d'ensembles ordonnés quelconques, c'est à cause des cas où ceux-ci ne sont \textit{pas} de la forme $\fP(\cP)$, $\fP(\cQ)$, et (alors même qu'ils le seraient) où la relation de base n'est pas déduite comme ci-dessus d'une relation quelconque entre ensembles $\cP$, $\cQ$.

Je reviens donc à l'axiomatique du début, avec $\cP$, $\cQ$ ordonnés, et $\cR$ satisfaisant (D). Dans tous les cas auxquels je peux songer, $\cP$ et $\cQ$ sont stables par sup quelconques, donc (D) équivaut à (D$_0$). Notons que l'application
\[
\cP \longrightarrow \fP(\cP), \qquad p \longmapsto \cP/p
\]
est injective, et de même pour $\cQ \longrightarrow \fP(\cQ)$. Ainsi, la connaisance de $A \longmapsto A^{\cup}$ équivaut à celle de % [sic] connaisance -> connaissance, faute d'orthographe
\[
\begin{array}{rcl}
A & \longmapsto & \beta(A)\\[0.5ex]
\beta : \fP(\cP) & \longrightarrow & \cQ,
\end{array}
\]
et de même celle de $B \longmapsto B^{\cap}$ équivaut à
\msp{75}
celle de
\[
\alpha : \fP(\cQ) \longrightarrow \cP.
\]
En fait, il suffit de se donner
\[
\begin{array}{rcccl}
\beta_0 : & \cP & \longrightarrow & \cQ & \qquad \text{applications}\\
\alpha_0 : & \cQ & \longrightarrow & \cP & \qquad \textit{décroissantes}
\end{array}
\tag*{(5.15)}
\]
définies par
\[
\left\{
\begin{array}{rcl}
\beta_0(p) & = & \beta(\{p\})\\[1ex]
\alpha_0(q) & = & \alpha(\{q\}),
\end{array}
\right.
\]
puisqu'on en déduit $\beta$, $\alpha$ par
\[
\left\{
\begin{array}{rcl}
\beta(A) & = & \inf\nolimits_{p \in A} \beta_0(p)\\[1ex]
\alpha(B) & = & \inf\nolimits_{q \in B} \alpha_0(q).
\end{array}
\right.
\tag*{(5.16)}
\]
La relation $\cR$ se déduit de la donnée de $\beta_0$, ou au choix de $\alpha_0$, par
\[
(p,q) \in \cR \iff q \leq \beta_0(p) \iff p \leq \alpha_0(q).
\tag*{(5.17)}
\]

\printedpage{47}
% head: A. Grothendieck : Dérivateurs, Ch. II, Édition des Universités de Montpellier II et Paris VII

Cela implique donc que $\alpha_0$, $\beta_0$ se déterminent mutuellement, par les formules
\[
\left\{
\begin{array}{rcl}
\alpha_0(q) & = & \sup \underbrace{\{p \in \cP \mid q \leq \beta_0(p)\}}_{\{q\}^{\cap}}\\[3ex]
\beta_0(p) & = & \sup \underbrace{\{q \in \cQ \mid p \leq \alpha_0(q)\}}_{\{p\}^{\cup}}.
\end{array}
\right.
\tag*{(5.18)}
\]
Reste à se demander quelles sont les applications $\alpha_0 : \cP \longrightarrow \cQ$, ou $\beta_0 : \cQ \longrightarrow \cP$, qui correspondent à une \textit{`relation dualisante'} % [sic] \alpha_0 : \cP \to \cQ, \beta_0 : \cQ \to \cP -> \beta_0 : \cP \to \cQ, \alpha_0 : \cQ \to \cP, les rôles de \alpha_0 et \beta_0 sont échangés par rapport à (5.15)
\msp{76}
$\cR$ sur $\cP \times \cQ$. Voici les conditions sur $\beta_0 : \cP \longrightarrow \cQ$ :
\begin{itemize}[leftmargin=4em, labelsep=0.8em]
\item[\circled{a}] $\beta_0$ est \textit{décroissante},
\end{itemize}
ce qui implique que la relation
\[
\underbrace{\cR_{\beta_0}}_{\subseteq \cP \times \cQ} = \{(p,q) \in \cP \times \cQ \mid q \leq \beta_0(p)\}
\]
satisfait à la condition (D$_0$) du côté gauche, et il faut encore la vérifier du côté droit, i.e. regarder, pour $q$ fixé, l'ensemble
\[
\{q\}^{\cap} = \{p \in \cP \mid q \leq \beta_0(p)\}.
\]
Comme $\beta_0$ est décroissant, on voit que $p \in \{q\}^{\cap}$, et $p' \leq p \Longrightarrow p' \in \{q\}^{\cap}$, i.e. $\{q\}^{\cap}$ est \textit{ouvert.} On doit vérifier encore (supposant l'existence des sup quelconques dans $\cP$) que cet ensemble est stable par sup quelconques - il en résultera qu'il est de la forme $\cP/\alpha(q)$. Cela équivaut à la condition
\begin{itemize}[leftmargin=4em, labelsep=0.8em]
\item[\circled{b}] $\beta_0$ transforme sup quelconques dans $\cP$ en inf dans $\cQ$.
\end{itemize}
Si on regarde $\beta_0$ comme un foncteur contravariant $\cP^{o} \longrightarrow \cQ$, cela signifie qu'il commute aux limites projectives, ou encore qu'il a un adjoint à gauche $\cQ \longrightarrow \cP^{o}$, lequel commute aux limites inductives, et correspond
\msp{77}
à l'application décroissante $\alpha_0 : \cQ \longrightarrow \cP$, transformant sup en inf.

Ici il est commode de considérer $\natural$ et $\flat$ comme des opérations sur $\cP$ resp. $\cQ$, par
\[
p^{\natural} = \alpha_0\beta_0(p), \qquad q^{\flat} = \beta_0\alpha_0(q).
\tag*{(5.19)}
\]
La relation entre ces endomorphismes, et l'opération de même nom sur $\fP(\cP)$, $\fP(\cQ)$, est donnée par
\[
\left\{
\begin{array}{rcl}
p^{\natural} & = & \text{Plus grand élément de } \{p\}^{\natural}\\[1ex]
A^{\natural} & = & \sup\nolimits_{p \in A} p^{\natural},
\end{array}
\right.
\tag*{(5.20)}
\]
\printedpage{48}
% head: A. Grothendieck : Dérivateurs, Ch. II, Édition des Universités de Montpellier II et Paris VII
et symétriquement du côté de $\cQ$.

On dit que $p$ est $\cR$-\textit{clos,} ou $\beta_0$-clos, ou $\alpha_0$-clos, si $p = p^{\natural}$, et de même $q \in \cQ$ est dit $\cR$-\textit{clos} etc. si $q = q^{\flat}$. Les applications
\[
\left\{
\begin{array}{ccccccc}
p & \longmapsto & p^{\natural} & : & \cP & \longrightarrow & \cP\\[1ex]
q & \longmapsto & q^{\flat} & : & \cQ & \longrightarrow & \cQ
\end{array}
\right.
\tag*{(5.21)}
\]
sont \textit{croissantes} et \textit{idempotentes.} On peut caractériser $p^{\natural}$ comme le plus petit élément clos de $\cP$ majorant $p$, et de même pour $q^{\flat}$. On trouve alors, par (5.19), que $\alpha_0$ et $\beta_0$ induisent des bijections inverses l'une de l'autre,
\msp{78}
entre l'ensemble des éléments clos de $\cP$, et l'ensemble des éléments clos de $\cQ$
\[
\begin{tikzcd}
\cP^{\natural} \arrow[r, shift left, "\bar{\beta}_0"] \arrow[r, phantom, "\scriptstyle\sim"] & \cQ^{\flat}. \arrow[l, shift left, "\bar{\alpha}_0"]
\end{tikzcd}
\tag*{(5.22)}
\]
Mais en pratique, les applications $\alpha_0$, $\beta_0$ ne sont pas obtenues directement, mais on dispose de la relation $\cR$, satisfaisant la condition (D$_0$), dont on déduit $\alpha_0$, $\beta_0$. Dans tous les cas auxquels je pense ici, ou bien $\cP$, $\cQ$ sont de la forme $\fP(P_0)$, $\fP(Q_0)$, et $\cR$ se déduit d'une relation $R_0 \subseteq P_0 \times Q_0$ comme ci-dessus (de loin le plus fréquent, il me semble), ou bien $\cP$ et $\cQ$ apparaissent comme sous-ensembles convenables d'ensembles $\fP(P_0)$, $\fP(Q_0)$, munis de l'ordre induit, et la relation $\cR$ sur $\cP \times \cQ$ s'obtient d'une façon plus complexe.

\noindent\textbf{Exemples.}
\begin{enumerate}[label={\arabic*)}]
\item Relations d'\textit{orthogonalité} entre parties de deux modules sur un anneau, muni d'un accouplement de dualité. Dans le cas des \texttt{[espaces]} vectoriels de dimension finie et d'une vraie dualité, les parties closes sont les sous-espaces vectoriels, dans le cas infini, ce sont les sous-espaces vectoriels fermés pour la `topologie faible' définie par l'accouplement.
\item Relations d'\textit{orthogonalité} entre parties de deux groupes topologiques commutatifs localement compacts duales l'une de l'autre, au sens de Pontrjagin.
\msp{79}

Les parties closes sont les sous-groupes fermés.
\item Pour deux espaces vectoriels `en dualité' sur un corps valué complet, tel $\mathbf{R}$ où $\mathbf{C}$, la relation % [sic] où -> ou, accent grave fautif
\[
|\langle x, x' \rangle| \leq 1.
\]
Dans le cas réel ou complexe, les parties closes sont les parties convexes cerclées [\,?] faiblement fermées.
\item Pour les espaces vectoriels réels, la relation de semi-polarité
\[
\langle x, x' \rangle \geq -1.
\]
Les parties closes sont les parties convexes faiblement fermées contenant l'origine.
\end{enumerate}


\endgroup
% END UNIT Der-ChII-234/Der-ChII-234_p046-050.tex
% BEGIN UNIT Der-ChII-234/Der-ChII-234_p051-055.tex
\begingroup\makeatletter%
\newcommand{\printedpage}[1]{\marginpar{\footnotesize p.~#1}}%
\newcommand{\msp}[1]{\marginpar{\footnotesize [#1]}}%
\newcommand{\Fl}{\mathrm{Fl}}%
\newcommand{\Pf}{\mathfrak{P}}%
\setcounter{footnote}{61}%
\makeatother\let\printedpage\derpp\let\msp\dermsp%
\let\mspin\DERmspin%


\printedpage{49}
% head: A. Grothendieck : Dérivateurs, Ch. II, Édition des Universités de Montpellier II et Paris VII
% foot: Version 8 juillet 2001

\begin{itemize}[leftmargin=2.5em]
\item[5)] Soit $K$ un corps (pour simplifier), $L$ une extension galoisienne de groupe $G$. Relation entre $L$ et $G$
\[
\gamma\cdot x=x,\qquad \text{i.e. } x\in L^{\gamma}.
\]
Les parties closes de $L$ sont les sous-extensions, les parties closes de $G$ les sous-groupes de $G$.
\end{itemize}

Il y a des analogues divers dans des cas inséparables (à coups de $p$-algèbres de Lie par exemple), et des énoncés de descente en tous genres.

\msp{page 80}Tous ces exemples sont de 1\textsuperscript{ère} espèce, les ensembles ordonnés qui interviennent sont de la forme $\Pf(P_0)$, $\Pf(Q_0)$, et la relation entre eux est déduite d'une relation $R_0$ entre $P_0$ et $Q_0$. Voici maintenant des exemples plus complexes.

\begin{itemize}[leftmargin=2.5em]
\item[6)] Soit $\mathcal{M}$ une catégorie (‘grosse’, le plus souvent \dots). On prend
\[
\mathcal{P}=\mathcal{Q}=\Pf(\Fl(\mathcal{M})).
\]
La relation $\mathcal{R}\subseteq\mathcal{P}\times\mathcal{Q}$ est la relation de Quillen, que j'écris, pour $\Phi,\Psi\in\Pf$
\[
\Psi\subseteq\Phi_{*},
\]
ou encore de façon équivalente
\[
\Phi\subseteq\Psi^{*},
\]
($\Psi$ a la LLP \texttt{[left lifting property]} pour $\Psi$, $\Phi$ la RLP \texttt{[right lifting property]} pour $\Psi$). Donc les applications $\alpha_0$, $\beta_0$ sont données par % [sic] Psi a la LLP pour Psi -> Psi a la LLP pour Phi, the relation is between Psi and Phi (cf. "Phi la RLP pour Psi")
\[
\beta_0(\Phi)=\Phi_{*},\qquad \alpha_0(\Psi)=\Psi^{*}.
\]
L'idée ici ne viendrait à personne d'introduire $\Pf(\mathcal{P})$, $\Pf(\mathcal{Q})$ - c'est déjà assez gros comme ça ! La notation pour $\Phi^{\natural}$ est ici $\tilde{\Phi}$, pour $\Psi^{\flat}$ c'est $\bar{\Psi}$. Les parties closes sont celles que j'appelle Q-closes à gauche et à droite.

La relation $\mathcal{R}$, que j'écris
\[
\Phi\longleftrightarrow\Psi,
\]
\msp{page 81}n'est pas déduite d'une relation entre flèches de $\mathcal{M}$, i.e. d'une partie de $\Fl(\mathcal{M})\times\Fl(\mathcal{M})$, mais est de nature plus complexe. Pour la décrire, il faut introduire l'ensemble $\mathcal{D}$ des diagrammes carrés commutatifs (traits pleins) de la forme
\[
\begin{tikzcd}
A \arrow[r] \arrow[d, "i"] & X \arrow[d, "p"] \\
B \arrow[r] \arrow[ur, dotted, "f"] & Y
\end{tikzcd}
\tag*{$(*)$}
\]
\printedpage{50}
% head: A. Grothendieck : Dérivateurs, Ch. II, Édition des Universités de Montpellier II et Paris VII
% foot: Version 8 juillet 2001
(qu'on peut interpréter comme $\Fl(\underline{\Fl}(\mathcal{M}))$, ensemble des flèches dans $\underline{\Fl}(\mathcal{M})$). On a une partie $\mathcal{D}_0$ de $\mathcal{D}$, formée des diagrammes pour lesquels \texttt{[il]} existe \texttt{[un]} $f$ avec deux triangles commutatifs (c'est donc l'image de $\underbrace{\Fl_3(\mathcal{M})}_{\mathrm{Hom}(\boldsymbol{\Delta}^3,\mathcal{M})}$ dans $\mathrm{Hom}(\boldsymbol{\Delta}^1\times\boldsymbol{\Delta}^1,\mathcal{M})$, transposée de $\boldsymbol{\Delta}^1\times\boldsymbol{\Delta}^1\longrightarrow\boldsymbol{\Delta}^3$ donné par $(0,0)\longmapsto 0$, $(0,1)\longmapsto 1$, $(1,0)\longmapsto 2$, $(1,1)\longmapsto 3$), et une application
\[
\mathcal{D}_0\subseteq\mathcal{D}\xrightarrow{\ \varphi\ }\Fl(\mathcal{M})\times\Fl(\mathcal{M}),
\]
associant à $(*)$ le couple $(i,j)$ \texttt{[plutôt} $(i,p)$\texttt{]} (source et but de $(*)$, considéré comme flèche dans $\underline{\Fl}(\mathcal{M})$ de $i$ dans $p$). La relation $\mathcal{R}$ est alors déduite de $\varphi$, $\mathcal{D}_0$ par la condition
\msp{page 82}
\[
\Phi\longleftrightarrow\Psi\qquad\Longleftrightarrow\qquad\varphi^{-1}(\Phi\times\Psi)\subseteq\mathcal{D}_0.
\]
Mais à vrai dire, je m'aperçois que c'est quand-même déduit d'une relation entre éléments de $\Fl(\mathcal{M})$, et cette relation entre $i : A\longrightarrow B$ et $p : X\longrightarrow Y$ peut s'écrire ainsi :
\begin{itemize}[leftmargin=3em]
\item[$(**)$] pour toutes flèches $A\longrightarrow X$, $B\longrightarrow Y$ rendant commutatif le diagramme $(*)$, la flèche $\varphi$ pointillée dans $(*)$ existe. % [sic] la flèche $\varphi$ -> la flèche $f$, the dotted arrow in $(*)$ is labelled $f$
\end{itemize}
\end{itemize}

Cela me fait penser que si on a une relation $\mathcal{R}$ entre $\Pf(P)$ et $\Pf(Q)$, satisfaisant $(\mathrm{D}_0)$, donc décrite par une application décroissante
\[
\Pf(P)\longrightarrow\Pf(Q)
\]
transformant sup en inf, la donné d'une telle application équivaut à celle d'une application quelconque % [sic] la donné -> la donnée, feminine noun
\[
P\longrightarrow\Pf(Q),
\]
ou ce qui revient au même, d'une relation $\mathcal{R}$ sur $P\times Q$. Donc cette structure doit être regardée comme ‘de première espèce’. La situation change à partir du moment où pour décrire la relation $\mathcal{R}$ avec les propriétés voulues, on se trouve contraint de remplacer $\Pf(P)$, $\Pf(Q)$ par des sous-ensembles convenables - quand on ne regarde que certaines parties de $P$, $Q$. C'est ce qu'on
\msp{page 83}va faire au § suivant. Le cas que je vais y traiter est le seul de son espèce que j'ai rencontré - un cas de ‘dualité formelle’ qui s'introduit de façon naturelle, et qui ne soit ‘de 1\textsuperscript{ère} espèce’.

\printedpage{51}
% head: A. Grothendieck : Dérivateurs, Ch. II, Édition des Universités de Montpellier II et Paris VII
% foot: Version 8 juillet 2001
\section*{6\quad Transporteurs, noyaux, stabiliseurs relatifs à $W$}

Soit $\mathcal{M}$ une catégorie, munie d'une partie
\[
W\subseteq\Fl(\mathcal{M}).
\tag*{(6.1)}
\]
Je pense au cas d'un localiseur, mais n'introduirai d'hypothèse sur $W$ qu'au besoin. Pour simplifier, je suppose $\mathcal{M}$ stable par produits fibrés.

Je pose
\[
\left\{
\begin{array}{rcl}
\mbox{(\footnotemark)}\mathcal{P} & = & \Pf(W)=\{A\in\Pf(\Fl(\mathcal{M}))\mid A\subseteq W\}\\[1ex]
\mathcal{Q} & = & \{L\in\Pf(\Fl(\mathcal{M}))\mid\
\parbox[c]{17em}{$L$ \textit{stable} par changement de\\
base (par qui je sous-entends aussi,\\
stable par isomorphie)\\
et par composition. $\}$.}
\end{array}
\right.
\tag*{(6.2)}
\]
\footnotetext{cf. rectification p. 86.}
J'introduis une relation
\[
\mathcal{R}\subseteq\mathcal{P}\times\mathcal{Q}
\]
de la façon suivante. Pour $A\in\mathcal{P}$, $L\in\mathcal{Q}$, on pose
\[
(A,L)\in\mathcal{R}\iff
\begin{array}[t]{l}
\text{Pour tout diagramme cartésien}\\
\begin{tikzcd}[ampersand replacement=\&]
X \arrow[d, "f"] \& X' \arrow[l] \arrow[d, "f'"] \\
Y \& Y' \arrow[l, "p"]
\end{tikzcd}\\
\text{avec $f\in A$, $p\in L$, on a $f'\in W$}\\
\text{(i.e. $A$ est stable par changement de base par $p\in L$).}
\end{array}
\tag*{(6.3)}
\]

\msp{page 84}Quand j'ordonne $\mathcal{P},\mathcal{Q}\subseteq\Pf(\Fl(\mathcal{M}))$ par la relation d'ordre induite par l'inclusion de parties de $\Fl(\mathcal{M})$, on trouve des sous-ensembles ordonnés de $\Pf(\Fl(\mathcal{M}))$, stables par réunions quelconques, et aussi par intersections quelconques, avec la seul réserve que dans le cas de $\mathcal{P}$, la famille d'indices doit être $\neq\emptyset$ (puisqu'alors on trouve comme intersection $\Fl(\mathcal{M})$ tout entier, qui n'est pas en général contenu dans $W$). % [sic] la seul réserve -> la seule réserve, feminine agreement

On va vérifier la condition $(\mathrm{D}_0)$ de la page 69, pour $\mathcal{R}$ et les relations d'ordre. Il y a deux vérifications à faire.

1$^{o}$) \texttt{[où est ‘2$^{o}$’ ?]} Donnons nous $A\in\mathcal{P}$, i.e. $A\subseteq W$, et considérons l'ensemble $\mathcal{R}(A)$ des $L\in\mathcal{Q}$ tels que l'on ait $(A,L)\in\mathcal{R}$. Il est clair que c'est un ouvert (et de même, que ${}^{s}\mathcal{R}(L)$ est un ouvert de $\mathcal{Q}$ pour tout $L$ dans $\mathcal{Q}$) - i.e. $L\in\mathcal{R}(A)$, $L'\subseteq L$ implique % ?? uncertain: left superscript before R(L) (also on p. 53) -- italic s or upright s % [sic] ouvert de $\mathcal{Q}$ -> ouvert de $\mathcal{P}$, cf. "l'ouvert ${}^{s}\mathcal{R}(L)$ de $\mathcal{P}$" p. 53
\printedpage{52}
% head: A. Grothendieck : Dérivateurs, Ch. II, Édition des Universités de Montpellier II et Paris VII
% foot: Version 8 juillet 2001
$L'\in\mathcal{R}(A)$. Reste à trouver un plus grand élément de $\mathcal{R}(A)$. Notons que si $L\in\mathcal{R}(A)$, alors pour tout $p\in L$, $p : Y'\longrightarrow Y$, le changement de base
\[
p^{*} : \mathrm{Cat}/Y\longrightarrow\mathrm{Cat}/Y' : X\longmapsto X\times_{Y}Y'
\]
doit transformer flèches $f\in A$ en flèches $g\in W$, comme il résulte du fait que $L$ est stable par changement de base.
\[
\begin{tikzcd}
X_1 \arrow[d, "f"'] & X'_1 \arrow[l, "q_1"'] \arrow[d, "g"] \\
X_2 \arrow[d] & X'_2 \arrow[l, "q_2"'] \arrow[d] \\
Y & Y'. \arrow[l, "p"']
\end{tikzcd}
\]
Donc $L\subseteq\mathbf{L}_W(A)$, où on pose
\msp{page 85}
\[
\begin{aligned}
\mathbf{L}_W(A)=\{(p : Y'\longrightarrow Y)\in\Fl(\mathcal{M})\quad\mid\quad & p^{*} : \mathrm{Cat}/Y\longrightarrow\mathrm{Cat}/Y'\ \text{transforme}\\
& \text{flèches de $A$ en flèches \textit{de} $W$.}\}
\end{aligned}
\]
(\footnote{NB le dire plutôt en termes de diagrammes cartésiens, vu qu'on n'a pas fait d'hypothèse de stabilité de $W$ par isomorphismes.}).
\[
\begin{tikzcd}
X_1 \arrow[d, "f"'] \arrow[dr, phantom, "\text{cart.}"] & X'_1 \arrow[l] \arrow[d, "f'"] \\
X_2 \arrow[d] \arrow[dr, phantom, "\text{cart.}"] & X'_2 \arrow[l] \arrow[d] \\
Y & Y' \arrow[l, "p"']
\end{tikzcd}
\]
\[
f\in A\implies f'\in W.
\]
Il est clair d'autre part que cette partie de $\Fl(\mathcal{M})$ est stable par changement de base, donc $\mathbf{L}_W(A)\in\mathcal{Q}$, et que $(A,\mathbf{L}_W(A))\in\mathcal{R}$. Donc $\mathbf{L}_W(A)$ est un plus grand élément de $\mathcal{R}(A)$, donc on trouve
\[
\mathcal{R}=\{(A,L)\in\mathcal{P}\times\mathcal{Q}\mid L\subseteq\mathbf{L}_W(A)\}.
\tag*{(6.5)}
\]
\printedpage{53}
% head: A. Grothendieck : Dérivateurs, Ch. II, Édition des Universités de Montpellier II et Paris VII
% foot: Version 8 juillet 2001
\[
\begin{tikzcd}
X_1 \arrow[d, "f"'] \arrow[dr, phantom, "\text{cart.}"] & X'_1 \arrow[l, "q_1"'] \arrow[d, "f'"] \\
X_2 \arrow[d] \arrow[dr, phantom, "\text{cart.}"] & X'_2 \arrow[l, "q_2"'] \arrow[d] \\
Y_1 \arrow[d] \arrow[dr, phantom, "\text{cart.}"] & Y'_1 \arrow[l, "p_1"'] \arrow[d] \\
Y & Y'. \arrow[l, "p"']
\end{tikzcd}
\]
Inversement, partons de $L\in\mathcal{Q}$, et montrons que l'ouvert ${}^{s}\mathcal{R}(L)$ de $\mathcal{P}$ a un plus grand élément. Pour ceci, posons
\[
\begin{aligned}
\mathbf{A}_W(L)=\{(f : X\longrightarrow Y)\in W\subseteq\Fl(\mathcal{M})\quad\mid\quad\forall\qquad
&\begin{tikzcd}[ampersand replacement=\&]
X \arrow[d, "f"'] \arrow[dr, phantom, "\textit{cart.}"] \& X' \arrow[l, "q"'] \arrow[d, "f'"] \\
Y \& Y', \arrow[l, "p"]
\end{tikzcd}\\
&\text{avec $p\in L$, on a $f'\in W$}\}.
\end{aligned}
\tag*{(6.6)}
\]
(NB Quand $L$ contient les identités, comme c'est le cas notamment s'il est de la forme $\mathbf{L}_W(A)$ - alors il contient les isomorphismes, du moins si $A$ les contient - alors $\mathbf{A}_W(L)$ est l'ensemble de $f\in\Fl(\mathcal{M})$ satisfaisant la condition énoncée après $|$, sans avoir à spécifier $f\in W$.) On a $\mathbf{A}_W(L)\in\mathcal{P}$, et on vérifie que $(\mathbf{A}_W(L),L)\in\mathcal{R}$ (en utilisant la stabilité de $L$ par composition), enfin il est clair que $(A,L)\in\mathcal{R}$ implique $A\subseteq\mathbf{A}_W(L)$ (en utilisant $A\subseteq W$), donc $\mathbf{A}_W(L)$ est un plus grand élément de ${}^{s}\mathcal{R}(L)$.

\msp{page 86}Notons que si $W$ \textit{est stable par composition} (ce que nous allons supposer désormais), il en est de même de $\mathbf{A}_W(L)$ (ce qui utilise encore le fait que $L$ soit stable par changement de base). De plus, si $W$ \textit{contient les isomorphismes} (ce que nous supposerons également), $\mathbf{A}_W(L)$ les contient aussi. Quant à $\mathbf{L}_W(A)$, il contient les isomorphismes \textit{pourvu que} $A$ soit stable par isomorphismes ; quand $A$ est stable par composition et contient les isomorphismes, cette condition sera automatiquement satisfaite. Nous tenons à cette condition, c'est pourquoi du côté de $A$, dans la définition de $\mathcal{P}$, je vais introduire cette condition supplémentaire que $A$ soit \textit{stable par isomorphismes.} Donc je pose, en corrigeant (6.2),
\[
\left\{
\begin{array}{rcl}
\mathcal{P} & = & \{A\in\Pf(\Fl(\mathcal{M}))\quad\mid\quad\mathrm{Is}(\mathcal{M})\subseteq A\subseteq W,\quad A\text{ stable par composition}\}\\[1ex]
\mathcal{Q} & = & \{L\in\Pf(\Fl(\mathcal{M}))\quad\mid\quad\mathrm{Is}(\mathcal{M})\subseteq L,\quad L\text{ stable par changement de base}\}.
\end{array}
\right.
\tag*{(6.7)}
\]


\endgroup
% END UNIT Der-ChII-234/Der-ChII-234_p051-055.tex
% BEGIN UNIT Der-ChII-234/Der-ChII-234_p056-060.tex
\begingroup\makeatletter%
\newcommand{\printedpage}[1]{\marginpar{\footnotesize p.~#1}}%
\newcommand{\msp}[1]{\marginpar{\footnotesize [#1]}}%
\newcommand{\Fl}{\mathrm{Fl}}%
\newcommand{\Is}{\mathrm{Is}}%
\newcommand{\Fib}{\mathrm{Fib}}%
\makeatother\let\printedpage\derpp\let\msp\dermsp%
\let\mspin\DERmspin%


\printedpage{54}
% head: A. Grothendieck : Dérivateurs, Ch. II, Édition des Universités de Montpellier II et Paris VII
On est bien sûr intéressé surtout aux éléments clos de $\mathcal{P}$, $\mathcal{Q}$, i.e. les parties $A$, $L$ de $\Fl(\mathcal{M})$, contenant les isomorphismes, stables par composition (et $L$ stable aussi par changement de base, $A \subseteq W$) telles que l'on ait
\[
A = \mathbf{A}_W(L), \qquad L = \mathbf{L}(A) \tag*{(6.8)}
\]
\texttt{[plutôt }$\mathbf{L}_W(A)$\texttt{]}.

\msp{page 87}

Je garderai les notations $\natural$, $\flat$ du paragraphe précédent. Quels sont les deux couples extrêmes ?

Partant de $A_0 = \Is(\mathcal{M})$ (plus petit des éléments de $\mathcal{P}$), on trouve $\mathbf{L}_W(A_0) = \Fl(\mathcal{M})$ (plus grand élément de $\mathcal{Q}$), et $A_0^{\natural} = \mathbf{A}_W(\Fl(\mathcal{M})) = W_{\mathrm{univ}}$, ensemble des flèches qui sont dans $W$ et le restent par tout changement de base :
\[
A_0^{\natural} = W_{\mathrm{univ}}, \qquad \beta_0(A_0) \overset{\text{déf}}{=} \mathbf{L}_W(A_0) = \Fl(\mathcal{M}), \qquad \text{\scriptsize où $A_0 = \Is(\mathcal{M})$ plus petit élément de $\mathcal{P}$.} \tag*{(6.9)}
\]
Inversement, partons de $L_0 = \Is(\mathcal{M})$, le plus petit élément de $\mathcal{Q}$. On a $\mathbf{A}_W(L_0) = W$ ($W$ étant stable par composition et contenant les isomorphismes), et $L_0^{\flat} = \mathbf{L}_W(W) = \Fib_W$, les $W$-fibrations :
\[
L_0^{\flat} = \Fib_W, \qquad \alpha_0(L_0) = W, \qquad \text{\scriptsize où $L_0 = \Is(\mathcal{M})$, plus petit élément de $\mathcal{Q}$.} \tag*{(6.10)}
\]
Donc les parties $\mathcal{R}$-closes à gauche de $\Fl(\mathcal{M})$ se promènent entre $W_{\mathrm{univ}}$ et $W$, les parties $\mathcal{R}$-closes à droite entre $\Fib_W$ et $\Fl(\mathcal{M})$.

\textbf{Proposition 6.1.} \textit{Soient $A \in \mathcal{P}$, $L \in \mathcal{Q}$ tels que $A = \alpha_0(L)$, $L = \beta_0(A)$ (i.e. $A = \mathbf{A}_W(L)$, $L = \mathbf{L}_W(A)$). Pour que $L$ soit stable par composition, il faut et il suffit que $A$ soit stable par changement de base par les flèches $f \in L$. (Plus généralement, si $L \in \mathcal{Q}$ est un élément quelconque de $\mathcal{Q}$, et si on suppose de} \msp{page 88} \textit{plus que $L$ est stable par composition, alors $A = \mathbf{A}_W(L)$ est stable par changement de base par $L$.)}

Terminologie : pour $A \in \mathcal{P}$, i.e. $A$ une partie de $\Fl(\mathcal{M})$ contenue dans $W$ contenant les isomorphismes et stable par composition, $\mathbf{L}_W(A)$ s'appelle le \textit{transporteur de $A$ dans $W$}. (Il y a une notion duale de \textit{cotransporteur.}) C'est un élement de $\mathcal{Q}$, donc une partie de $\Fl(\mathcal{M})$ contenant les isomorphismes et stable par changement de base. D'autre part, pour $L$ dans $\mathcal{Q}$, les $f \in \mathbf{A}_W(L)$ s'appellent flèches \textit{$W$-$L$-asphériques,} ou les \textit{$W$-$L$-équivalences,} ou les \textit{$W$-équivalences $L$-universelles,} et $\mathbf{A}_W(L)$ s'appelle le \textit{$L$-noyau de $W$}. On pourrait le noter plus simplement $W_L$ (quand cette notation ne conflicte pas avec $W_S \subseteq \Fl(\mathcal{M}/S)$ pour $S \in \mathrm{Ob}\,\mathcal{M}$). C'est une partie de $W$ qui (comme $W$ lui-même) contient les isomorphismes et est stable par composition. % [sic] élement -> élément, missing accent

Enfin, il y a lieu aussi d'introduire le \textit{stabiliseur} $\mathrm{Stab}(A)$, ou transporteur de $A$ dans $A$, qui ne dépend pas de la donné de $W$, mais seulement de la partie $A$ de $\Fl(\mathcal{M})$, contenant les isomorphismes et stable par composition. On pourrait aussi le noter $\Fib_A$, et appeler ses éléments des \textit{$A$-fibrations} ($A$ jouant à présent le rôle de $W$). On a $\mathrm{Stab}(A) \subseteq \mathbf{L}_W(A)$, et il semblerait qu'en % [sic] donné -> donnée, feminine noun
% foot: Version 8 juillet 2001
\printedpage{55}
% head: A. Grothendieck : Dérivateurs, Ch. II, Édition des Universités de Montpellier II et Paris VII
\msp{page 89}
pratique, ce sont ces stabiliseurs, plus que les transporteurs $\mathbf{L}_W(A)$, qui sont importants (quand les deux ne coïncident pas). C'est ainsi que dans Cat, si $W$ est un localiseur fondamental, les morphismes $W$-lisses (resp. $W$-propres) $\mathrm{Liss}_W$ ($\mathrm{Prop}_W$), qui sont les stabiliseurs des ensembles $W_{\mathrm{as}}$ ($W_{\mathrm{coas}}$) des ensembles de morphismes $W$-asphériques (resp. $W$-coasphériques), jouant un rôle essentiel, alors que je ne sais pas caractériser de façon sensible les flèches des transporteurs $\mathbf{L}_W(W_{\mathrm{as}})$, $\mathbf{L}_W(W_{\mathrm{coas}})$, lesquels contiennent, en plus des morphismes $W$-lisses (resp. $W$-propres) les $W$-fibrations $\Fib_W$. Je ne sais pas comment caractériser une classe de morphismes qui apparaît aussi hétéroclite! De même, il y a lieu d'introduire et de caractériser le stabiliseur de $W_{\mathrm{bias}}$ (ensemble des flèches $W$-biasphériques), comme les morphismes `$W$-nets' $\mathrm{Net}_W$, par une description similaire à celle des $\mathrm{Liss}_W$, $\mathrm{Prop}_W$, utilisant les bilocalisés $\backslash\underset{u}{X}/$, $\backslash\underset{v}{Y}/$ , au lieu des localisés $X/x$, $Y/y$ ou des colocalisés $x\backslash X$, $y\backslash Y$. Il y a lieu enfin de déterminer le stabiliseur de l'ensemble Equ des Cat-équivalences (ce qui ne dépend pas de la donné d'un $W$ - mais on a $\mathrm{Equ} \subseteq \mathcal{P}(W)$ pour tout $W$ fondamental). % [sic] donné -> donnée, feminine noun
% foot: Version 8 juillet 2001

\printedpage{56}
% head: A. Grothendieck : Dérivateurs, Ch. II, Édition des Universités de Montpellier II et Paris VII
\msp{page 90}

\section*{7\quad Morphismes cofiltrants et localement cofiltrants}

Mais je ne veux pas m'attarder à présent sur ces développements, qui ne me paraissent pas urgents, et voudrais revenir aux questions liées à la notion de catégories filtrantes ou cofiltrantes.

\textbf{Definition 7.1.} Soit $f : X \longrightarrow Y$ une flèche dans Cat. On dit que $f$ est \textit{cofiltrante} (resp. \textit{localement cofiltrante}, resp. \textit{faiblement cofiltrante} (\footnote[64]{J'ai modifié la terminologie : maintenant faiblement cofiltrant $=$ 0-connexe $+$ localement faiblement cofiltrant, \textit{stable par passage} à $X'$ sur $X$, $X'$ étale et $\infty$-connexe. Revoir la terminologie dans le cas faiblement cofiltrant, qui est une notion traîtresse, faute d'être stable par passage \texttt{[à]} un $X'$ étale sur $X$.}), resp. \textit{localement faiblement cofiltrante} ou encore \textit{précofiltrante}) si pour tout $y$ dans $Y$, la catégorie colocalisé $y\backslash X$ est cofiltrante (resp. \dots). (Dualement, on dit que $f$ est \textit{filtrante}, resp. \dots, si les catégories $X/y$ sont filtrantes, resp. \dots, i.e. si $f^o : X^o \longrightarrow Y^o$ est cofiltrante resp. \dots.) % [sic] Definition -> Définition, missing accent; colocalisé -> colocalisée, agreement with catégorie

\textbf{Remarque 7.2.} Notions stables par changement de base coétale. Contient les isomorphismes et même les Cat-équivalences. Les morphismes \textit{étales} sont localement filtrants.

\textbf{Proposition 7.3.} \textit{Soit $f : X \longrightarrow Y$ un morphisme cofiltrant (resp. localement cofiltrant, resp. faiblement cofiltrant, resp. précofiltrant), et supposons que $Y$ soit cofiltrant (resp. localement cofiltrant, resp. faiblement cofiltrant, resp. précofiltrant). Alors $X$ lui-même est cofiltrant (resp. \dots).}

\textsc{Démonstration}. Il s'agit, dans les quatre cas envisagés, de vérifier pour $X$ la condition $\underbrace{\mathrm{PS}_g1' \;/\; \mathrm{PS}_g1}_{\text{deux derniers cas}} \;/\; \mathrm{PS}_g1' + \mathrm{PS}_g2 \;/\; \mathrm{PS}_g1 + \mathrm{PS}_g2$. Mais \msp{page 91} pour chacune des trois conditions $\mathrm{PS}_g1$, $\mathrm{PS}_g1'$, $\mathrm{PS}_g2$ séparément, on voit immédiatement que si elle est vérifiée pour $Y$ et pour les $y\backslash X$, elle l'est pour $X$, d'où la proposition.

\textbf{Corollaire 7.3.} \textit{Les quatre notions envisagées dans (7.1) sont stables par composition.}

\textsc{Démonstration} immédiate, en utilisant la remarque (7.2).
\[
\begin{tikzcd}
X \arrow[d] & z\backslash X \arrow[l] \arrow[d] \\
Y \arrow[d] & z\backslash Y \arrow[l] \arrow[d] \\
Z & z\backslash Z. \arrow[l]
\end{tikzcd}
\]
% foot: Version 8 juillet 2001

\printedpage{57}
% head: A. Grothendieck : Dérivateurs, Ch. II, Édition des Universités de Montpellier II et Paris VII
\textbf{Remarque 7.4.}
\begin{enumerate}
\item[1)] Notons que la notion de morphisme cofiltrant ou faiblement cofiltrant n'est \textit{pas} stable par changement de base quelconque, ni même par une immersion ouverte. Exemple : $e \longrightarrow \boldsymbol{\Delta}^1$, appliquant $e$ en 1, est cofiltrant,
\[
\begin{tikzcd}
 & e \arrow[d, mapsto] \\
0 \arrow[r] & 1
\end{tikzcd}
\]
mais sa fibre en 0 est vide, donc n'est ni cofiltrante, ni faiblement cofiltrante. Je n'ai pas vu d'exemple montrant que la notion de flèche \textit{localement} cofiltrante resp. précofiltrante n'est pas stable par changement de base, et ne voudrais pas m'y attarder, n'ayant pas besoin de connaître la réponse.
\item[2)] Pour que $X$ soit cofiltrante (resp. \dots), il faut et il suffit que $X \longrightarrow e$ le soit. Ainsi la notion absolue est ramenée à une notion relative.
\item[3)] Si $f : X \longrightarrow Y$ est cofiltrante, elle est \msp{page 92} $W$-coasphérique (donc $\in W$) pour tout $W$ satisfaisant l'axiome Loc (8) des limites inductives. En particulier, elle est $W_0$-asphérique, i.e. un foncteur cofiltrant (et même seulement faiblement cofiltrant) est \textit{cofinal.}
\end{enumerate}
La raison principale pour m'intéresser aux flèches cofiltrantes est contenue dans la

\textbf{Proposition 7.4.} \textit{Soit $f : X \longrightarrow Y$ dans} Cat \textit{localement cofiltrant (resp. localement faiblement cofiltrant). Alors :}
\begin{enumerate}
\item[\textit{a)}] \textit{Pour tout $X'$ étale sur $X$, le composé $f' : X' \longrightarrow X \longrightarrow Y$ est localement cofiltrant (resp. \dots), et pour toute factorisation de $f'$ en $X' \xrightarrow{g'} Y' \longrightarrow Y$ avec $Y' \longrightarrow Y$ étale, $g'$ est localement cofiltrant (resp. \dots).}
\item[\textit{b)}] \textit{$\forall\ x$ dans $X$, le morphisme sur les localisés}
\[
f/x : X/x \longrightarrow Y/x \qquad (y = f(x)) % [sic] Y/x -> Y/y, since y = f(x)
\]
\textit{est cofiltrant (resp. faiblement cofiltrant).}
\item[\textit{c)}] \textit{Si $x$, $y$ comme dans b), et $U$ est un ouvert de $X/x$, désignant par $V$ l'ouvert de $Y/y$ engendré par l'image de $f/x$ (i.e. l'ouvert formé des $z$ dans $Y/y$ tels que $z\backslash(X/x)$ soit non vide), la flèche induite $q : U \longrightarrow V$} \texttt{[est]} \textit{cofiltrante (resp. faiblement cofiltrante), a fortiori cofinale (et même $W$-coasphérique, pour tout $W$ satisfaisant l'axiome des limites Loc (8))} (\footnote[65]{Il suffit ici qu'on soit dans l'hypothèse respée ($f$ \textit{faiblement} localement cofiltrant) sauf pour le NB ici \texttt{[`le NB ici' est `(et même ...Loc (8))']}. % [sic] respée -> intended word unclear (perhaps « respectée » or « resp. »), not a French word
}). % doubt resolved: PDF prints q here; later g is an edition naming inconsistency (recorded in corrections.csv).

\msp{page 93}
\textit{Donc pout tout} cofaisceau \textit{$G$ sur $V$, à valeurs dans une catégorie $\mathcal{M}$ avec $\varinjlim$ , on a} \texttt{[un]} \textit{isomorphisme} % [sic] pout -> pour, misprint
\[
\varinjlim_U g^*(G) \xrightarrow{\ \sim\ } \varinjlim_V G,
\]
% foot: Version 8 juillet 2001
\printedpage{58}
% head: A. Grothendieck : Dérivateurs, Ch. II, Édition des Universités de Montpellier II et Paris VII
\textit{et pour tout préfaisceau $F$ sur $V$, à valeurs dans une catégorie $\mathcal{M}$ avec $\varprojlim$ , on a} \texttt{[un]} \textit{isomorphisme}
\[
\varprojlim_U g^*(F) \xleftarrow{\ \sim\ } \varprojlim_V F.
\]
\end{enumerate}

\textsc{Démonstration}. Je vais d'abord expliciter le

\textbf{Lemme 7.5.} \textit{Soit $h : X' \longrightarrow X$ un morphisme étale dans} Cat. \textit{Alors $h$ est localement cofiltrant. De façon plus précise, pour tout $x$ dans $X$, on a}
\[
\pi_0(x\backslash X') \xrightarrow{\ \sim\ } \mathrm{Ob}\,X'_x,
\]
\textit{toute composante connexe de $x\backslash X'$ contient un élément et un seul de la catégorie discrète $X'_x \subseteq x\backslash X'$. Si $x'$ est dans $X'_x$ (i.e. $x'$ objet de $X'$ tel que $h(x') = x$), c'est un objet} initial \textit{de la composante connexe qui le contient. (Donc cette composante est cofiltrante.)}

Ceci prouve (et au delà) la première assertion de 7.4, compte tenu de 7.3 (stabilité par composition), ou parce qu'une \msp{page 94} catégorie étale (en l'occurrence $y\backslash X'$) au dessus d'une catégorie localement cofiltrante (resp. localement faiblement cofiltrante) (en l'occurrence $y\backslash X$) est itou.
\[
\begin{tikzcd}[row sep=large]
y\backslash X' \arrow[r] \arrow[d, "\text{ét.}"'] \arrow[dr, phantom, "\text{cart.}"] & X' \arrow[d, "\text{ét.}"] \\
y\backslash X \arrow[r] \arrow[d] \arrow[dr, phantom, "\text{cart.}"] & X \arrow[d] \\
y\backslash Y \arrow[r] & Y.
\end{tikzcd}
\]
Pour la deuxième assertion dans a), on peut maintenant supposer $X' = X$. Il faut prouver que pour tout $y'$ dans $Y'$, $y'\backslash X \overset{\text{déf}}{=} X \times_{Y'} y'\backslash Y'$ est localement cofiltrante (resp. localement faiblement cofiltrante). Soit $y$ dans $Y$ l'image de $y'$, et considérons le diagramme de trois carrés cartésiens ci-contre.
\[
\begin{tikzcd}[row sep=large]
y'\backslash X \arrow[r] \arrow[d, "\text{ét.}"] \arrow[dr, phantom, "\textit{cart.}"] & y\backslash X \arrow[r] \arrow[d, "\text{ét.}"] \arrow[dr, phantom, "\textit{cart.}"] & X \arrow[d, "\text{ét.}"] \\
y'\backslash Y' \arrow[r] & y\backslash Y' \arrow[r] \arrow[d] \arrow[dr, phantom, "\textit{cart.}"] & Y' \arrow[d] \\
 & y\backslash Y \arrow[r] & Y
\end{tikzcd}
\]
Le Lemme 7.5 implique que $y'\backslash Y' \longrightarrow y\backslash Y'$ est l'inclusion dans $y\backslash Y'$ d'une composante connexe de $y\backslash Y'$, c'est donc une immersion ouverte. Donc $y'\backslash X \longrightarrow y\backslash X$ est une immersion
% foot: Version 8 juillet 2001


\endgroup
% END UNIT Der-ChII-234/Der-ChII-234_p056-060.tex
% BEGIN UNIT Der-ChII-234/Der-ChII-234_p061-065.tex
\begingroup\makeatletter%
\newcommand{\printedpage}[1]{\marginpar{\footnotesize p.~#1}}%
\newcommand{\msp}[1]{\marginpar{\footnotesize [#1]}}%
\newcommand{\margmark}[1]{\marginpar{\footnotesize #1}}%
\newcommand{\PS}{\mathrm{PS}}%
\newcommand{\Ob}{\mathrm{Ob}}%
\newcommand{\tl}[1]{\begin{tabular}{@{}l@{}}#1\end{tabular}}%
\newcommand{\tls}[1]{\text{\footnotesize\begin{tabular}{@{}l@{}}#1\end{tabular}}}%
\newcommand{\circled}[1]{\tikz[baseline=(c.base)]\node[draw,circle,inner sep=1pt](c){#1};}%
\makeatother\let\printedpage\derpp\let\msp\dermsp%
\let\mspin\DERmspin%


%======================================================================
\printedpage{59}
% head: A. Grothendieck : Dérivateurs, Ch. II, Édition des Universités de Montpellier II et Paris VII

ouverte. Comme $y\backslash X$ est localement cofiltrant (resp. localement faiblement cofiltrant) par hypothèse, il en est donc de même de $y\backslash X$. Cela prouve d).

Appliquant a) à la situation $X/x \longrightarrow Y/y$ dans b), on voit que les colocalisés $z\backslash(X/x)$ de
\margmark{!} % margin: "!" with a triple vertical rule alongside the lines "X/x sur Y/y ... donc est 0-connexe."
$X/x$ sur $Y/y$ ($z \in \Ob\,(Y/y)$) sont localement cofiltrants (resp. faiblement). Mais comme $X/x$ a un objet final au dessus de l'objet final de $Y/y$, il s'ensuit que $z\backslash(X/x)$ a un objet final, donc est $0$-\textit{connexe.}

\msp{page 95}
Donc ces catégories, étant localement cofiltrantes (resp. faiblement), sont en fait cofiltrantes (resp. faiblement cofiltrantes).

Prouvons c). Notons que pour $z$ dans $V$, $z\backslash U$ s'identifie à l'image inverse, par $z\backslash(X/x) \longrightarrow X/x$, de l'ouvert $U$ de $X/x$, c'est donc un \textit{ouvert} de $z\backslash(X/x)$, non vide par construction de $V$. Mais un ouvert non vide d'une catégorie cofiltrante (resp. faiblement cofiltrante) est encore cofiltrant (resp. faiblement cofiltrant). Cela signifie que $g : U \longrightarrow V$ est cofiltrant (resp. faiblement cofiltrant), a fortiori cofinal (et $W$-coasphérique, pour tout $W$ satisfaisant Loc (8)). Cela achève la démonstration.

\noindent\textbf{Remarque 7.5.}
\begin{itemize}
\item[1)] Le même argument que dans c) montre que si $f : X \longrightarrow Y$ est cofiltrant (resp. faiblement cofiltrant), alors pour tout ouvert $U$ dans $X$, si on désigne par $V$ l'ouvert de $Y$ engendré par $f(U)$, i.e. l'ouvert formé des $y$ dans $Y$ tels que $y\backslash X \neq \emptyset$, alors la flèche induite $g : U \longrightarrow V$ est cofiltrante (resp. faiblement cofiltrante), et a fortiori $g$ est cofinal (et plus généralement, si $f$ est cofiltrant, $g$ est $W$-coasphérique pour tout $W$ satisfaisant la condition Loc (8)).

\msp{page 96}
\item[2)] Un cas particulier intéressant de 7.4 est celui où on suppose que les $y\backslash X$ ($y \in \Ob\,Y$) sont sommes directes de catégories ayant chacune un \textit{objet initial.} Alors on voit que cette condition est stable par les opérations de 7.4 a) [\,?], et par conséquent, les colocalisés $z\backslash(X/x)$ des morphismes localisés en haut $X/x \longrightarrow Y/y$ sont des catégories à objet initial. Ceci implique alors qu'il en est de même des colocalisés $z\backslash U$ ($z \in \Ob\,V$) dans $g : U \longrightarrow V$. Il s'ensuit que $g$ est $W$-coasphérique pour \textit{tout} localiseur fondamental $W$ (sans avoir à se soucier de l'axiome des limites).

La condition envisagée sur $f : X \longrightarrow Y$ est satisfaite notamment dans les deux cas suivants (qui me paraissent de loin les plus importants) : a) $f$ est étale et b) $X$ est discrète.
\end{itemize}
% foot: Version 8 juillet 2001

%======================================================================
\par
\printedpage{60}
% head: A. Grothendieck : Dérivateurs, Ch. II, Édition des Universités de Montpellier II et Paris VII

\msp{page 97}

\begin{center}
{\small
\begin{tikzcd}[ampersand replacement=\&, column sep=small, row sep=large]
\left\{\begin{tabular}{@{}c@{\ }l@{}}
 & Top $X$\\
 & localement totalement 0-connexe\\
 & et 0-connexe\\
$\Longleftrightarrow$ & totalement 0-connexe\\
 & et localement irréductible\\
$\Longleftrightarrow$ & $X$ \textit{cofiltrante,}\\
 & i.e. $\PS^*1' + \PS^*2$\\
$\Longleftrightarrow$ & $X$ 0-connexe\\
 & $+ \PS^*1 + \PS^*2$\\
$\Longleftrightarrow$ & $\left\{\begin{tabular}{@{}l@{}}
$\forall\, x \in \Ob\,X$, $X/x$ et\\
$X/x \times x$ sont irréductible,\\ % [sic] irréductible -> irréductibles, agreement with "sont"
$X$ 0-connexe
\end{tabular}\right.$\\
$\Longleftrightarrow$ & $\left\{\begin{tabular}{@{}l@{}}
$X$ 0-connexe,\\
les $X/a \times_x b$ non vides\\
(et même 0-connexes),\\
les $X/a \times_{x\times x} b$ 0-connexes
\end{tabular}\right.$
\end{tabular}\right.
\;
\left[\begin{tabular}{@{}l@{}}
stable par\\
passage\\
à $X' \longrightarrow X$\\
étale sur $X$, si\\
$X'$ 0-connexe,\\
\textit{mais} si les\\
$X/x$ sont cofiltrantes\\
et $X$ 0-connexe,\\
cela n'implique\\
\textit{pas} $X$ cofiltrante\\
Ex. $B_G$
\end{tabular}\right]
\arrow[rrrd, Rightarrow, "{\tl{stricte\\(ex. $\bullet\ \bullet$)}}"]
\arrow[rrdd, Rightarrow, "{\tl{stricte\\(ex. $\boldsymbol{\Delta}$)}}"]
\arrow[ddd, Rightarrow, "{\tl{stricte (ex.\\$B_G$)}}"']
\& \& \& |[draw,rounded corners=6pt]| \varepsilon = \bullet\ \bullet \\
\& \& \&
\begin{tabular}{@{}l@{}}
\begin{tabular}{@{}c@{\ }l@{}}
 & Top $X$ localement\\
 & totalement 0-connexe\\
$\Longleftrightarrow$ & les composantes connexes de $X$\\
 & sont cofiltrantes\\
$\Longleftrightarrow$ & $\PS^*1 + \PS^*2$\\
 & (X \textit{localement cofiltrante}) % [sic] X (upright roman) -> $X$ (italic), inconsistent with all other occurrences
\end{tabular}\\[2ex]
$\left[\begin{tabular}{@{}l@{}}
Stable par passage à $X' \longrightarrow X$\\
étale, mais si les $X'/x$\\
y satisfont, pas forcément $X'$!
\end{tabular}\right]$
\end{tabular}
\arrow[d, Rightarrow] \\
\& \&
\begin{tabular}{@{}l@{}}
\begin{tabular}{@{}c@{\ }l@{}}
 & Top $X$ totalement 0-connexe\\
 & (on dit $X$ \textit{totalement} 0-\textit{connexe})\\
$\Longleftrightarrow$ & $X^o \longrightarrow X^o \times X^o$ est cofinal\\
 & \textit{et} $X \neq \emptyset$\\
$\Longleftrightarrow$ & $\forall\, x, y \in \Ob\,X$, $x \times y$ 0-connexe de $X^\wedge$\\
 & i.e. $X/x \times y$ 0-connexe\\
 & \textit{et} $X \neq \emptyset$
\end{tabular}\\[2ex]
$\left[\begin{tabular}{@{}l@{}}
NB n'est pas stable par passage aux $X/x$ (\footnotemark[66]), et si\\
les $X/x$ sont totalement 0-connexes et $X \neq \emptyset$,\\
il ne s'ensuit pas que\\
$X$ le soit (\footnotemark[67]).
\end{tabular}\right]$
\end{tabular}
\arrow[dddl, Rightarrow, "{\tl{stricte (ex. $B_G$, ou\\$0 \rightrightarrows 1$)}}"]
\&
\tl{Top $X$\\localement irréductible} \\
\begin{tabular}{@{}c@{\ }l@{}}
 & Top $X$ localement\\
 & irréductible\\
 & $+$ 0-connexe\\
$\Longleftrightarrow$ & $X$ \textit{pseudofiltrante}\\
 & $+$ 0-\textit{connexe,}\\
 & i.e. $\PS^*1$\\
 & $+$ 0-connexe
\end{tabular}
\arrow[ddr, Rightarrow, "{\tl{stricte\\(ex. $0 \rightrightarrows 1$)}}"']
\&
\left[\begin{tabular}{@{}l@{}}
$\PS^*1$ ne dépend\\
que des catégories\\
induites $X/x$,\\
donc stable par passage\\
aux $X' \longrightarrow X$ étales,\\
donc la condition\\
$\PS^*1 +{}$0-connexe stable\\
par passage aux\\
revêtements\\
$X' \longrightarrow X$ 0-\textit{connexes.}
\end{tabular}\right]
\& \& \\
|[draw,rounded corners=6pt]| B_G \& \& \& \\
\&
\begin{tabular}{@{}c@{\ }l@{}}
 & Top $X$ irréductible\\
 & ($X$ \textit{irréductible})\\
$\Longleftrightarrow$ & $\PS^*1'$
\end{tabular}
\& \& \\
|[draw,rounded corners=6pt]| 0 \rightrightarrows 1
\&
\left[\begin{tabular}{@{}l@{}}
ne dépend que \texttt{[de]} $\mathrm{Ord}\,X$ ou de $\mathrm{Esp}(X)$,\\
signifie que $\mathrm{Ord}\,X$ cofiltrant,\\
ou $\mathrm{Esp}(X)$ irréductible.\\
N'est pas stable par passage à\\
$X' \longrightarrow X$ étale 0-connexe,\\
p.ex. par passage à un $X/x$,\\
p.ex. $X = (0 \rightrightarrows 1)$ est irréductible,\\
mais $X/1 = \left(\begin{tabular}{@{}l@{}}$\cdot$\\ $\;\searrow$\\ $\qquad\cdot$\\ $\;\nearrow$\\ $\cdot$\end{tabular}\right)$ ne l'est\\
pas.
\end{tabular}\right]
\& \& \\
\end{tikzcd}
}
\end{center}
% footnote marks (66), (67) are set inside the diagram; texts below
\footnotetext[66]{\textbf{Ex.} \circled{$\boldsymbol{\Delta}$}.}
\footnotetext[67]{\textbf{Ex.} $B_G$.}
% foot: Version 8 juillet 2001

%======================================================================
\par
\printedpage{61}
% head: A. Grothendieck : Dérivateurs, Ch. II, Édition des Universités de Montpellier II et Paris VII

\msp{page 98}

\begin{center}
{\small
\begin{tikzcd}[ampersand replacement=\&, column sep=small, row sep=large]
\& \tl{Top $X$ localement totalement 0-connexe\\$+$ 0-connexe\\($X$ \textit{cofiltrante})}
\arrow[ddl, Rightarrow] \arrow[ddr, Rightarrow] \arrow[rrdd, Rightarrow]
\& \& \\
\& \& \& |[draw,rounded corners=6pt]| \varepsilon = \bullet\ \bullet \\
\tl{Top $X$ localement irréductible\\$+$ 0-connexe\\($X$ \textit{pseudofiltrante} et 0-\textit{connexe})\\$\PS^*1 +{}$ 0-connexe}
\arrow[ddr, Rightarrow]
\& \&
\tl{Top $X$ totalement 0-connexe\\($X$ \textit{totalement} 0-\textit{connexe})\\$X^o \longrightarrow X^o \times X^o$ cofinal}
\arrow[ddl, Rightarrow]
\&
\tl{Top $X$ localement totalement\\0-connexe\\($X$ \textit{localement cofiltrante})\\$\PS^*1 + \PS^*2$}
\arrow[dd, Rightarrow] \\
|[draw,rounded corners=6pt]| B_G \& \& |[draw,circle,inner sep=1pt]| \boldsymbol{\Delta} \& \\
\& \tl{Top $X$ irréductible\\($X$ \textit{irréductible})\\$\PS^*1'$}
\& \&
\tl{Top $X$ localement irréductible\\($X$ \textit{pseudo-cofiltrante} [\,?])\\$\PS^*1$} \\
\& |[draw,rounded corners=6pt]| 0 \rightrightarrows 1 \& \&
|[draw,rounded corners=12pt]| \begin{tabular}{@{}c@{}}$B_G \amalg e$\\avec $G \neq 1$\end{tabular} \\
\end{tikzcd}
}
\end{center}

\begin{center}
{\small
\begin{tikzcd}[ampersand replacement=\&, column sep=small, row sep=large]
\text{(\footnotemark[68]).}
\& \tl{$X$ \textit{cofiltrante}\\$\PS^*1' + \PS^*2$\\$= \PS^*1 + \PS^*2 +{}$ 0-connexe}
\arrow[ddl, Rightarrow] \arrow[ddr, Rightarrow] \arrow[rrdd, Rightarrow]
\& \& \\
\& \& \& |[draw,rounded corners=6pt]| \bullet\ \bullet \\
\tl{$X$ \textit{pseudo-cofiltrante} et 0-\textit{connexe}\\($\PS^*1 +{}$ 0-connexe)}
\arrow[ddr, Rightarrow] \arrow[ddrrr, Rightarrow]
\& \&
\tl{$X$ \textit{totalement} 0-\textit{connexe}\\($X^o \longrightarrow X^o \times X^o$ cofinal)}
\arrow[ddl, Rightarrow]
\&
\tl{$X$ \textit{localement cofiltrante}\\$\PS^*1 + \PS^*2$}
\arrow[dd, Rightarrow] \\
|[draw,rounded corners=6pt]| B_G \& \& |[draw,circle,inner sep=1pt]| \boldsymbol{\Delta} \& \\
\& \tl{$X$ \textit{irréductible}\\$\PS^*1'$}
\& \&
\tl{$X$ \textit{pseudo-cofiltrante}\\$\PS^*1$} \\
\& |[draw,rounded corners=6pt]| \bullet \rightrightarrows \bullet \& \&
|[draw,rounded corners=6pt]| B_G \amalg e \\
\end{tikzcd}
}
\end{center}
\footnotetext[68]{\textbf{NB} Les cinq catégories typiques $B_G$ ($G \neq 1$), $\boldsymbol{\Delta}$, $0 \rightrightarrows 1$, $\varepsilon = \bullet\ \bullet$, $B_G \amalg e$ ($G \neq 1$) sont telles qu'elles appartiennent à la classe envisagée $C$ d'objets de Cat, mais à aucune des cinq autres classes sauf celles qui sont reliées à $C$ par $C \Longrightarrow C'$ sur le diagramme d'implications (ou par une implication composée).}
% foot: Version 8 juillet 2001

%======================================================================
\par
\printedpage{62}
% head: A. Grothendieck : Dérivateurs, Ch. II, Édition des Universités de Montpellier II et Paris VII

\noindent\begin{tabular}{@{}p{2.8cm}@{}p{11.5cm}@{}}
$\bullet \rightrightarrows \bullet$ & n'est ni pseudo-cofiltrante, ni totalement 0-connexe.\\[1ex]
$B_G$ & n'est ni localement cofiltrante, ni totalement 0-connexe.\\[1ex]
$\boldsymbol{\Delta}$ & n'est pas pseudo-filtrante.\\[1ex]
$B_G \amalg e$ & n'est pas localement cofiltante, car $B_G$ ne l'est pas, ni appartient à une des 4 classes d'objets connexes.\\[1ex] % [sic] cofiltante -> cofiltrante, letter r missing
Itou pour & $\varepsilon = \bullet\ \bullet$
\end{tabular}

\msp{page 99}

\begin{center}
{\small
\begin{tikzcd}[ampersand replacement=\&, column sep=small, row sep=large]
\& \tl{$X$ \textit{cofiltrante}\\($\PS^*1' + \PS^*2$\\$\Longleftrightarrow \PS^*1 + \PS^*2$\\$+$ 0-connexe)}
\arrow[ddl, Rightarrow, "{\tls{$\PS^*1'$\\($\Longrightarrow$ 0-connexe)}}"]
\arrow[dd, Rightarrow, "{\tls{$\PS^*1$\\(les $x\backslash X$ satisfont\\$\PS^*1'$)}}"]
\arrow[ddr, Rightarrow]
\& \\
% layout: the two small labels stand between the arrows; attribution to the arrows as in this layout
\& \& |[draw,rounded corners=6pt]| \varepsilon = \bullet\ \bullet \\
\tl{$X$ totalement 0-connexe\\($X^o \longrightarrow X^o \times X^o$ cofinal)}
\arrow[ddd, Rightarrow]
\&
\tl{$X$ \textit{pseudo-cofiltrante}\\0-connexe\\($\PS^*1 +{}$ 0-connexe)\\(manque $\PS^*2$)}
\arrow[dddl, Rightarrow] \arrow[dddr, Rightarrow]
\&
\tl{$X$ \textit{localement cofiltrante}\\(composantes connexes filtrantes)\\($\PS^*1 + \PS^*2$)}
\arrow[ddd, Rightarrow] \\
|[draw,circle,inner sep=1pt]| \boldsymbol{\Delta} \& |[draw,rounded corners=6pt]| B_G \& \\
\& G \neq \{1\} \& \\
\tl{$X$ irréductible\\($\PS^*1'$)}
\& \&
\tl{$X$ \textit{pseudo-cofiltrante}\\ou encore \textit{localement irréductible}\\(composantes connexes\\pseudo-filtrantes)\\($\PS^*1$)} \\
|[draw,rounded corners=6pt]| \bullet \rightrightarrows \bullet \& \& |[draw,rounded corners=6pt]| B_G \amalg e \\
\& \& G \neq \{1\} \\
\end{tikzcd}
}
\end{center}

et dualement on a

\begin{center}
{\small
\begin{tikzcd}[ampersand replacement=\&, column sep=small, row sep=huge]
\& \text{$X$ filtrante}
\arrow[dl, Rightarrow] \arrow[d, Rightarrow] \arrow[dr, Rightarrow]
\& \\
\text{$X^o$ totalement 0-connexe}
\arrow[d, Rightarrow]
\&
\tl{$X$ pseudo-filtrante\\0-connexe}
\arrow[dl, Rightarrow] \arrow[dr, Rightarrow]
\&
\text{$X$ localement filtrante}
\arrow[d, Rightarrow] \\
\text{$X$ coirréductible} \& \& \text{$X$ pseudo-cofiltrante} \\
\end{tikzcd}
}
\end{center}
% foot: Version 8 juillet 2001

%======================================================================
\par
\printedpage{63}
% head: A. Grothendieck : Dérivateurs, Ch. II, Édition des Universités de Montpellier II et Paris VII

donc on a une terminologie directe sauf pour `$X^o$ totalement 0-connexe' - je n'ose dire `$X$ totalement co-connexe'!!
% foot: Version 8 juillet 2001


\endgroup
% END UNIT Der-ChII-234/Der-ChII-234_p061-065.tex
\clearpage\setcounter{footnote}{0}\phantomsection\addcontentsline{toc}{section}{Chapitre III}\markboth{Chapitre III}{Chapitre III}
% BEGIN UNIT Der-ChIII-567/Der-ChIII-567_p001-005.tex
\begingroup\makeatletter%
\newcommand{\printedpage}[1]{\marginpar{\footnotesize p.~#1}}%
\newcommand{\msp}[1]{\marginpar{\footnotesize [#1]}}%
\newcommand{\Hot}{\mathrm{Hot}}%
\newcommand{\D}{\mathbf{D}}%
\newcommand{\uHom}{\underline{\mathrm{Hom}}}%
\newcommand{\hm}{\mathrm{hom}}%
\makeatother\let\printedpage\derpp\let\msp\dermsp%
\let\mspin\DERmspin%


% ===================== page 1 (title page) =====================
\printedpage{--}
% no page number printed on this page
% head: (none)

\begin{center}
{\LARGE\sffamily LES DÉRIVATEURS}

\bigskip

\textbf{ALEXANDRE GROTHENDIECK}

\bigskip\bigskip

\textbf{Chapitre III}

\medskip

\textbf{Hom externes dans les dérivateurs}
\end{center}

\par

% ===================== page 2 (editorial note) =====================
\printedpage{--}
% no page number printed on this page
% head: (none)

\par\noindent\textbf{Note des éditeurs de l’édition transcrite.}\par
Ce texte a été déchiffré et transcrit en \LaTeXe{} par M. Künzer. Il a été édité par M. Künzer, J. Malgoire, et G. Maltsiniotis. La transcription est aussi fidèle que possible au manuscrit. Pour les quelques corrections évidentes, ou rares commentaires des éditeurs, ainsi que pour la numérotation originale des pages du manuscrit, les caractères de machine à écrire \texttt{[typewriter]} entre crochets sont utilisés. Un point d'interrogation entre crochets signifie que l'on n'est pas sûr du mot qui précède.

Cette édition est provisoire. Les remarques, commentaires et corrections sont bienvenus. Envoyer un message à :

\textbf{[e-mail address removed]}
\begin{flushright}
G. Maltsiniotis\qquad\qquad
\end{flushright}

\par

% ===================== page 3 (printed 1) =====================
\printedpage{1}
% head: A. Grothendieck : Dérivateurs, Ch. III, Édition des Universités de Montpellier II et Paris VII
% foot: Version 2 février 2002

\msp{page 1}

\section*{1\quad Relation des dérivateurs avec la catégorie $\Hot$}

Soit $\D$ un dérivateur, $\mathcal{A}$ sa catégorie fondamentale,
\[
\mathcal{A} = \D(e).
\]
Je me propose de définir, si faire se peut, des foncteurs
\[
\left\{
\begin{array}{rccccccl}
\Hot & \times & \mathcal{A} & \xrightarrow{\pi_{\D}} & \mathcal{A} & (\xi, M) & \longmapsto & \xi \times M \ \text{ ou } \ \xi \otimes M \quad \text{(suivant les cas)}\\
\Hot^o & \times & \mathcal{A} & \xrightarrow{\varpi_{\D}} & \mathcal{A} & (\xi, M) & \longmapsto & \hm(\xi, M).
\end{array}
\right.
\tag*{(1)}
\]
Ces foncteurs doivent s'échanger par dualité : si on prend les foncteurs $\pi^o$, $\varpi^o$ sur les catégories opposées, on trouve
\[
\left\{
\begin{array}{rcccc}
\Hot & \times & \mathcal{A}^o & \xrightarrow{\varpi^o_{\D}} & \mathcal{A}^o\\
\Hot^o & \times & \mathcal{A}^o & \xrightarrow{\pi^o_{\D}} & \mathcal{A}^o,
\end{array}
\right.
\tag*{(1$'$)}
\]
et on doit avoir
\[
\left\{
\begin{array}{rcl}
\varpi^o_{\D} & = & \pi_{\D^o}\\
\pi^o_{\D} & = & \varpi_{\D^o}.
\end{array}
\right.
\]
À vrai dire, considérant (1) comme des foncteurs
\[
\begin{array}{rcl}
\Hot & \xrightarrow{\pi_{\D}} & \uHom(\mathcal{A}, \mathcal{A})\\
\Hot^o & \xrightarrow{\varpi_{\D}} & \uHom(\mathcal{A}, \mathcal{A}),
\end{array}
\]
on s'attend à ce que ces foncteurs s'étendent en
\[
\left\{
\begin{array}{rcl}
\Hot & \xrightarrow{\pi_{\D}} & \uHom(\D, \D)\\
\Hot^o & \xrightarrow{\varpi_{\D}} & \uHom(\mathcal{A}, \mathcal{A})
\end{array}
\right.
\tag*{(2)}
\]
\texttt{[plutôt deux fois} $\uHom(\D, \D)$\texttt{]}. Cela signifie, essentiellement, que pour $\xi \in \mathrm{Ob}\,\Hot$ donné, les opérations $\xi \times_{\D} M$ et $\uHom_{\D}(\xi, M)$ \texttt{[plutôt} $\hm_{\D}(\xi, M)$\texttt{]} sont définies non seulement dans $\mathcal{A} = \D(e)$, mais dans tous les $\D(I)$, et ceci de façon à commuter aux foncteurs $f^*$ :

\msp{page 2}
Si
\[
f : I' \longrightarrow I,
\]
\printedpage{2}
% head: A. Grothendieck : Dérivateurs, Ch. III, Édition des Universités de Montpellier II et Paris VII
% foot: Version 2 février 2002
alors pour $F \in \mathrm{Ob}\,\D(I)$, on doit avoir des isomorphismes canoniques
\[
\left\{
\begin{array}{rcl}
f^*(\xi \times_{\D} M) & \simeq & \xi \times_{\D} f^*(M)\\
f^*\hm_{\D}(\xi, M) & \simeq & \hm_{\D}(\xi, f^*(M)),
\end{array}
\right.
\tag*{(3)}
\]
avec les compatibilités habituelles pour un composé $fg$, relatif à $I'' \xrightarrow{g} I' \xrightarrow{f} I$.

Mais je me place de préférence dans $\mathcal{A}$, ce qui devrait impliquer les constructions similaires dans tous les $\D(I)$, quitte à appliquer ce qui a été dit au dérivateur $\D_I$.

L'intuition algébrique fondamentale est celle-ci, en désignant par
\[
X \longmapsto \mathrm{hot}(X) : (\mathrm{Cat}) \longrightarrow (\Hot)
\]
le foncteur de localisation de (Cat) dans (Hot), obtenu en inversant les flèches dans (Cat) qui sont des `équivalences faibles'. Supposons $X$ dans Diag, soit
\[
p : X \longrightarrow e
\]
le foncteur à valeurs dans la catégorie ponctuelle. On veut que
\[
\left\{
\begin{array}{rcccl}
\mathrm{hot}(X) \times_{\D} M & \simeq & p_!(p^*(M)) & = & H^{\D}_{\bullet}(X, M)\\
\hm_{\D}(\mathrm{hot}(X), M) & \simeq & p_*(p^*(M)) & = & H^{\bullet}_{\D}(X, M)
\end{array}
\right.
\tag*{(4)}
\]
(isomorphismes canoniques).

\msp{page 3}
Ceci donnerait un principe de définition pour les opérations cherchées. Une difficulté vient du fait qu'en général, Diag n'est pas égal à (Cat), et que les types d'homotopie provenant de Diag ne donnent pas (même à isomorphisme près) \textit{tous} les types d'homotopie. Intuitivement, on ne trouve que des types d'homotopie `assez petits'. Mais on peut localiser Diag par rapport aux équivalences faibles (il vaut mieux dire `quasi-isomorphismes'), on trouve alors une catégorie $(\Hot)_0$, et un foncteur canonique
\[
\Hot_0 \longrightarrow \Hot,
\]
si on a de la chance, ce foncteur sera pleinement fidèle, et on essayera de décrire les foncteurs (1) ou (2) en y remplaçant Hot par la sous-catégorie pleine des `petits' types d'homotopie. On peut aussi oublier (provisoirement!) Hot, et travailler avec $\Hot_0$. C'est ce qu'on va faire pour le moment.

Pour pouvoir définir les foncteurs
\[
\left\{
\begin{array}{lccccccl}
\Hot_0 & \times & \mathcal{A} & \longrightarrow & \mathcal{A} & (\xi, F) & \longmapsto & \xi \times_{\D} F\\
\Hot_0^o & \times & \mathcal{A} & \longrightarrow & \mathcal{A} & (\xi, F) & \longmapsto & \hm_{\D}(\xi, F),
\end{array}
\right.
\tag*{(5)}
\]
je note que pour $F$ \texttt{[plutôt} $M$\texttt{]} fixé, $X$ variable dans Diag, les
\printedpage{3}
% head: A. Grothendieck : Dérivateurs, Ch. III, Édition des Universités de Montpellier II et Paris VII
% foot: Version 2 février 2002
\msp{page 4}
derniers membres dans (4)
\[
H^{\D}_{\bullet}(X, M), \qquad H^{\bullet}_{\D}(X, M)
\tag*{(6)}
\]
sont bien, l'un covariant, l'autre contravariant en $X$. Donc il y a , à isomorphisme unique près, une seule paire de bifoncteurs (5) satisfaisant (bifonctoriellement) les isomorphismes (4), si et seulement si les expressions (6), considérées comme fonctorielles en $X$, transforment quasi-isomorphismes en isomorphismes. Cela équivaut donc à la validité du % [sic] "il y a ," -> "il y a,", espace imprimée avant la virgule

\textbf{Théorème} (?) (\footnote{Nouvel axiome ?}). \textit{Soit $f : X' \longrightarrow X$ dans Diag. Si $f$ est un quasi-isomorphisme, c'est une $\D$-équivalence.}

Quand $X$ est la catégorie ponctuelle, ce n'est là autre que l'axiome Der 3 (p. 67). J'ignore si c'est une conséquence des axiomes Der 1 à Der 5 (d). Si ce n'est pas le cas, il faudrait l'ajouter aux axiomes, je crois. Je n'ai pas même prouvé qu'une équivalence de catégories dans Diag (au sens des catégories!) est une $\D$-équivalence!

Notons la \textit{formule d'adjonction}
\[
\mathrm{Hom}(\xi \times_{\D} M, N) \xrightarrow{\sim} \mathrm{Hom}(M, \hm_D(\xi, N)).
\tag*{(7)}
\] % [sic] \hom_D (D italique maigre) -> \hom_{\mathbf{D}}, le dérivateur est noté en gras partout ailleurs
Si $\xi$ provient de $X \in \mathrm{Diag}$, alors cela s'écrit
\[
\mathrm{Hom}(p_!p^*M, N) \xrightarrow{\sim} \mathrm{Hom}(M, p_*p^*N),
\tag*{(8)}
\]
\msp{page 5}
or par adjonction les deux membres sont canoniquement isomorphes à
\[
\mathrm{Hom}_{\D(X)}(p^*M, p^*N).
\]

\textbf{Exemple.} Prenant pour $\xi$ le type d'homotopie du cercle standard (représenté p.ex. par l'ensemble ordonné
\[
\begin{tikzcd}[row sep=small, column sep=small]
 & b & \\
0 \arrow[ur] \arrow[dr] & & 0' \arrow[ul] \arrow[dl] \ , \\
 & a &
\end{tikzcd}
\]
où on convient de prendre 0 comme origine, et d'orienter de 0 vers $a$), les deux foncteurs en question sont (sauf erreur) les foncteurs suspension $\Sigma$ et `cosuspension' ou foncteur `espace des lacets' $\Omega$. Sous réserve de vérifier que $\Omega F$ coïncide bien avec ce que j'ai décrit ailleurs, à propos des carrés substantiels cartésiens dans $\mathcal{A}$.

Dans $\Hot_0$ on a une opération produit, en fait ça doit être le produit catégorique dans $\Hot_0$. Mais sur les modèles il est donné par
\[
\mathrm{hot}_0(X) \times \mathrm{hot}_0(Y) \ \simeq \ \mathrm{hot}_0(X \times Y)
\tag*{(8)}
\] % [sic] (8) -> (9), le numéro (8) est déjà employé plus haut sur cette page


\endgroup
% END UNIT Der-ChIII-567/Der-ChIII-567_p001-005.tex
% BEGIN UNIT Der-ChIII-567/Der-ChIII-567_p006-010.tex
\begingroup\makeatletter%
\newcommand{\printedpage}[1]{\marginpar{\footnotesize p.~#1}}%
\newcommand{\msp}[1]{\marginpar{\footnotesize [#1]}}%
\newcommand{\D}{\mathbf{D}}%
\newcommand{\Hot}{\mathrm{Hot}}%
\newcommand{\Hom}{\mathrm{Hom}}%
\newcommand{\hm}{\mathrm{hom}}%
\newcommand{\homext}{\mathrm{hom}\,\mathrm{ext}}%
\newcommand{\HOT}{\mathrm{HOT}}%
\newcommand{\Ob}{\mathrm{Ob}\,}%
\newcommand{\Ltens}{\overset{\mathbf{L}}{\otimes}}%
\newcommand{\RHom}{\mathbf{R}\mathrm{Hom}}%
\newcommand{\ab}{\mathrm{ab}}%
\newcommand{\defeq}{\overset{\text{déf}}{=}}%
\makeatother\let\printedpage\derpp\let\msp\dermsp%
\let\mspin\DERmspin%


\printedpage{4}
% head: A. Grothendieck : Dérivateurs, Ch. III, Édition des Universités de Montpellier II et Paris VII
% foot: Version 2 février 2002

\texttt{[deux fois} (8)\texttt{]}. Ceci posé, on a les isomorphismes canoniques, pour $\xi,\eta\in\Ob\Hot_0$,
\[
\left\{
\begin{aligned}
\xi\times_{\D}(\eta\times_{\D}M) &\;\simeq\; (\xi\times\eta)\times_{\D}M\\
\hm_{\D}(\xi,\hm_{\D}(\eta,M)) &\;\simeq\; \hm_{\D}(\xi\times\eta,M).
\end{aligned}
\right.
\tag*{(9)}
\]

Si $\xi$, $\eta$ sont représentés par $X$, $Y$, la première formule p.ex.
\msp{page 6}
s’établit en calculant le deuxième membre comme
\[
\varphi_!(M_{X\times Y})\simeq p_!(\;\underbrace{q'_!(M_{X\times Y})}_{\substack{\simeq\ q_!(M_Y)_X\\ \text{car $q$ est $\mathbf{D}$-lisse}}}\;)\simeq\xi\times_{\D}(\eta\times_{\D}M),
\]
O.K.
\[
\begin{tikzcd}
X \ar[d,"p"'] & X\times_e Y \ar[l,"q'"'] \ar[d,"p'"] \ar[dl,dotted,"\varphi"'] \\
e & Y \ar[l,"q"]
\end{tikzcd}
\]

J’ai aussi envie de prouver :

\textbf{Proposition 1.} \textit{Soit $\xi\in\Ob\Hot_0$. Alors $F\longmapsto\xi\times_{\D}F$, en tant que morphisme de dérivateurs $\D\longrightarrow\D$, est} exact à droite, \textit{et $F\longmapsto\hm_{\D}(\xi,F)$ est exact à gauche.}

On peut supposer que $\xi$ est défini par $X$ dans Diag. Soit $f:I\longrightarrow J$, soit $F$ dans $\D(I)$, il faut prouver
\[
f_*(\hm_{\D}(\xi,F))\;\simeq\;\hm_{\D}(\xi,f_*(F)).
\]
Considérons donc
\[
\begin{tikzcd}
X\times I \ar[r,"f'"] \ar[d,"{p=(x,i)\longmapsto i}"'] & X\times J \ar[d,"{q=(x,j)\longmapsto j}"] \\
I \ar[r,"f"] & J.
\end{tikzcd}
\]

On a
\[
\hm_{\D}(\xi,F)=p_*(p^*F),
\]
donc
\[
f_*(\hm_{\D}(\xi,F))=\underbrace{f_*p_*}_{(fp)_*}(p^*(F))=(qf')_*p^*(F)\simeq q_*(f'_*p^*(F)).
\]
Mais
\[
f'_*p^*(F)\;\simeq\;q^*f_*(F),
\]
formule de changement de base des $f_*$, pour le changement de base $q$, qui est lisse, car déduit de $X\longrightarrow e$, qui est lisse, par changement de base $J\longrightarrow e$. On trouve
\printedpage{5}
% head: A. Grothendieck : Dérivateurs, Ch. III, Édition des Universités de Montpellier II et Paris VII
% foot: Version 2 février 2002
\msp{page 7}
donc
\[
f_*(\hm_{\D}(\xi,F))\simeq q_*q^*(f_*(F))\defeq\hm_{\D}(\xi,f_*(F)),
\]
q.e.d. L’énoncé pour $F\longmapsto\xi\times_{\D}F$ est dual de celui qu’on vient de prouver.

Reprenons, pour $\xi\in(\mathrm{hot})_0$ \texttt{[}$=\Hot_0$\texttt{]}, la formule d’adjonction
\[
\Hom(\xi\times_{\D}M,N)\simeq\Hom(M,\hm_{\D}(\xi,N))\defeq\Hom_{\xi}(M,N).
\tag*{$(*)$}
\]

Si $\xi$ est défini par $X$ dans Diag, les deux membres de $(*)$ s’identifient à
\[
\Hom_{\D(X)}(p^*M,p^*N),
\]
où $p:X\longrightarrow e$ est le foncteur structural.

Fixons $M$, $N$, et considérons l’expression $(*)$ comme contrafoncteur
\[
(\Hot_0)^o\;\longrightarrow\;\mathrm{Ens}.
\]

Il semblerait que dans tous les cas connus, ce foncteur soit ‘représentable’ par un objet, sinon de $\Hot_0$ lui-même, du moins un objet dans Hot. Si on ne peut le définir comme objet $H$ de $\Hot_0$, le fait qu’on ait un isomorphisme canonique
\[
\Hom_{\xi}(M,N)\;\simeq\Hom_{\Hot}(\xi,H),
\tag*{(10)}
\]
fonctoriel en $\xi\in\Ob\Hot_0$, ne \texttt{[le]} définit pas forcément
\msp{page 8}
à isomorphisme unique près. En fait, quand
\[
\D=\D_{\mathcal{M},\Sigma}
\]
est défini par une catégorie de modèles $(\mathcal{M},\Sigma)$, alors il y a une définition naturelle d’une opération $\homext$ (Hom ‘externe’) à valeurs dans Hot, savoir un foncteur
\[
\left\{
\begin{aligned}
&\mathcal{A}^o\times\mathcal{A} &&\longrightarrow &&\mathrm{Hot}\\
&(M,N) &&\longmapsto &&\homext(M,N),
\end{aligned}
\right.
\tag*{(11)}
\]
où $\mathcal{A}=\mathcal{M}\Sigma^{-1}$ est la catégorie fondamentale de $\D=\D_{\mathcal{M},\Sigma}$. Il est vrai que je n’ai pas vérifié, si p. ex. $(\mathcal{M},\Sigma)$ est sous-jacente à une catégorie de modèles quillénienne, de sorte que (par \textsc{Thomason}) $\D$ est bien un dérivateur, que cette notion de $\homext$ satisfait bien à une formule de la forme (10). Ici je vais le poser [?], au besoin, ces données supplémentaires : donnée de (11), et d’un isomorphisme fonctoriel
\[
\underbrace{\Hom_{\xi}(M,N)}_{\begin{subarray}{l}=\ \Hom(\xi\times_{\D}M,N)\\ =\ \Hom(M,\hm_{\D}(\xi,N))\end{subarray}}
\quad\simeq\;\Hom_{\Hot}(\xi,\homext(M,N)).
\tag*{(12)}
\]

\printedpage{6}
% head: A. Grothendieck : Dérivateurs, Ch. III, Édition des Universités de Montpellier II et Paris VII
% foot: Version 2 février 2002
\msp{page 9}
Si maintenant $\eta$ est un deuxième élément de $\Hot_0$, on a envie d’écrire
\[
\left\{
\begin{aligned}
\homext(\eta\times_{\D}M,N) &\;\simeq\;\homext(M,\hm_{\D}(\eta,N))\\
&\;\simeq\;\hm_{\HOT}(\eta,\homext(M,N)),
\end{aligned}
\right.
\tag*{(13)}
\]
où dans le deuxième membre, $\hm_{\HOT}$ est l’opération (1) dans le cas où $\D=\mathrm{HOT}$, et qui est défini aussi en termes du produit cartésien $\xi\times\eta$ dans Hot, par
\[
\Hom_{\Hot}(\xi,\hm_{\HOT}(\eta,\zeta))\;\simeq\;\Hom_{\Hot}(\xi\times\eta,\zeta),
\]
en d’autres termes, pour $\eta$ fixé, le foncteur
\[
\zeta\;\longmapsto\;\hm_{\HOT}(\eta,\zeta)
\]
est défini comme le foncteur adjoint à droite du foncteur
\[
\xi\;\longmapsto\;\xi\times\eta.
\]

Pour vérifier heuristiquement la formule (13) (vérification valable ‘en forme’ si les trois objets envisagés sont en fait définis comme objets de $\Hot_0\ \ldots$), on prend $\Hom_{\Hot}(\xi,-)$ à valeurs dans les trois membres, pour $\xi$ donné dans $\Hot_0$. Pour le premier, on trouve
\[
\begin{aligned}
\Hom_{\xi}(\eta\times_{\D}M,N) &\;=\;\Hom(\xi\times_{\D}(\eta\times_{\D}M),N)\\
&\;=\;\Hom((\xi\times\eta)\times_{\D}M,N).
\end{aligned}
\]
Pour le deuxième
\[
\begin{aligned}
\Hom_{\xi}(M,\hm_{\D}(\eta,N)) &\;=\;\Hom_{\mathcal{A}}(\xi\times_{\D}M,\hm_{\D}(\eta,N))\\
&\;\simeq\;\Hom_{\mathcal{A}}(\eta\times_{\D}(\xi\times_{\D}M),N)\\
&\;\simeq\;\Hom_{\mathcal{A}}((\eta\times\xi)\times_{\D}M,N),
\end{aligned}
\]
\msp{page 10}
et enfin pour le troisième
\[
\begin{aligned}
\Hom_{\Hot}(\xi\times\eta,\homext(M,N)) &\;\simeq\;\Hom_{\xi\times\eta}(M,N)\\
&\;\simeq\;\Hom((\xi\times\eta)\times_{\D}M,N),
\end{aligned}
\]
donc toujours la même chose ! Pas de doute, ça doit marcher.

S’il faut se donner le $\homext$ (11) comme donnée supplémentaire, avec de plus un isomorphisme trifonctoriel (12), il faut aussi se donner les isomorphismes trifonctoriels (13), \textit{‘compatibles’} avec l’isomorphisme (12) au sens du calcul précédent.

\printedpage{7}
% head: A. Grothendieck : Dérivateurs, Ch. III, Édition des Universités de Montpellier II et Paris VII
% foot: Version 2 février 2002

Par exemple, prenant $\eta=$ type d’homotopie du cercle $S^1$, les formules (13) se spécialisent en
\[
\begin{aligned}
\homext(\Sigma M,N) &\;\simeq\;\homext(M,\Omega N)\\
&\;\simeq\;\Omega\homext(M,N)\\
&\;\defeq\;\hm_{\HOT}(S^1,\homext(M,N)).
\end{aligned}
\tag*{(14)}
\]

Mais cette formule dans Hot n’est sans doute pas tout à fait la bonne, car le bon formalisme de $\Omega$, $\Sigma$ se fait dans les catégories \textit{ponctuées,} et en particulier dans la catégorie homotopique ponctué $\Hot_{\bullet}$. % [sic] ponctué -> ponctuée, agreement with « catégorie »

Je vais expliciter un exemple, en prenant le dérivateur
\[
\mathrm{Dér}(A)
\]
‘dérivé’ d’une catégorie abélienne $A$. Donc $\D(e)=$ catégorie dérivée de $A$. Pour éviter des
\msp{page 11}
difficultés accessoires ici, je vais supposer que $A$ est la catégorie des $k$-modules, $k$ un anneau quelconque.

À tout $\xi\in\Ob\Hot$ est associé un ‘abélianisé’
\[
\xi_{\ab}\defeq H_{\bullet}(\xi)\in\underbrace{D(\mathrm{Ab})}_{\text{ou }\mathbf{D}(\mathbf{Z}\text{-Mod})}\;,\qquad\quad\text{catégorie dérivée des groupes abéliens,}
\]
objet qui n’a de cohomologie qu’en degré $\le 0$. Ceci dit, si $M$ est dans $D(k\text{-Mod})$, on aura
\[
\left\{
\begin{aligned}
\xi\times_{\D}M &\;\simeq\;\xi_{\ab}\Ltens_{\mathbf{Z}}M\\
\hm_{\D}(\xi,M) &\;\simeq\;\RHom_{\mathbf{Z}}(\xi_{\ab},M),
\end{aligned}
\right.
\tag*{(15)}
\]
formule qui est en cheville avec la formule de Künneth pour l’homologie
\[
(\xi\times\eta)_{\ab}\simeq\xi_{\ab}\Ltens_{\mathbf{Z}}\eta_{\ab},
\tag*{(16)}
\]
ce qui redonne un aspect familier aux formules (9), qui se réduisent à
\[
\left\{
\begin{aligned}
\xi_{\ab}\Ltens_{\mathbf{Z}}(\eta_{\ab}\Ltens_{\mathbf{Z}}M) &\;\simeq\;(\xi_{\ab}\Ltens_{\mathbf{Z}}\eta_{\ab})\Ltens_{\mathbf{Z}}M\\
\RHom_{\mathbf{Z}}(\xi_{\ab},\RHom_{\mathbf{Z}}(\eta_{\ab},M)) &\;\simeq\;\RHom_{\mathbf{Z}}(\xi_{\ab}\Ltens_{\mathbf{Z}}\eta_{\ab},M).\qquad\text{(\footnotemark[2])}
\end{aligned}
\right.
\]
\footnotetext[2]{Expliciter aussi la formule d’adjonction (7) (p. 4).}

D’autre part, existe-t-il bien un $\homext(M,N)$, si $M$, $N$ dans $D(k\text{-Mod})$ ? On pense à $\RHom_k(M,N)$, mais pour en faire un type d’homotopie, on doit ‘tronquer en degrés $>0$’. On pourrait donc prendre
\[
\homext(M,N)\;\simeq\;\tau_{\le 0}(\RHom_k(M,N)).
\]

\printedpage{8}
% head: A. Grothendieck : Dérivateurs, Ch. III, Édition des Universités de Montpellier II et Paris VII
% foot: Version 2 février 2002
\msp{page 12}
La formule fondamentale (10) est-elle alors valable ? Le deuxième membre est en effet
\[
\Hom_{\Hot}(\xi,H)\;\simeq\;\Hom_{\mathrm{Hotab}}(\xi_{\ab},H)\qquad\quad\text{(\footnotemark[3])}
\]
\footnotetext[3]{? ? est-ce vrai ? ? ?}
car $H$ est un objet abélien de Hot (cf. le tapis sur l’abélianisation \ldots), à partir de là c’est gagné :
\[
\begin{aligned}
\Hom_{\underbrace{\scriptstyle\mathrm{Hotab}}_{\begin{subarray}{l}\approx\ D(\mathbf{Z}\text{-Mod})^{\le 0},\\ \text{sous-catégorie \textit{pleine}}\\ \text{de }D(\mathbf{Z}\text{-Mod})\end{subarray}}}(\xi_{\ab},\tau_{\le 0}\RHom_k(M,N)) &\;\simeq\;\Hom_{D(\mathbf{Z}\text{-Mod})}(\xi_{\ab},\RHom_k(M,N))\\
&\;\simeq\;\Hom_{D(k\text{-Mod})}(\xi_{\ab}\times M,N),
\end{aligned}
\]
O.K.

D’autre part, le bifoncteur
\[
(M,N)\;\longmapsto\;\homext(M,N)\;=\;\tau_{\le 0}\RHom_k(M,N)
\]
se factorise canoniquement par Hotab et a fortiori par $\Hot_{\bullet}$. Le ‘bon’ $\homext(M,N)$ serait ici plutôt celui dans Hotab, voire même dans la catégorie $D^{-}(\mathbf{Z}\text{-Mod})$, déduite de celle là ‘en rendant le foncteur translation inversible’. Il est évident d’ailleurs qu’ici le $\xi_{\ab}\times_{\D}M$, où $\xi=S^1$, n’est pas ‘la bonne’ suspension, qui doit être le décalage des degrés – $\xi_{\ab}$ ne devrait avoir de la cohomologie qu’en degré $-1$, et non aussi en degré $0$.


\endgroup
% END UNIT Der-ChIII-567/Der-ChIII-567_p006-010.tex
\clearpage\setcounter{footnote}{0}\phantomsection\addcontentsline{toc}{section}{Chapitre IV}\markboth{Chapitre IV}{Chapitre IV}
% BEGIN UNIT Der-ChIV-567/Der-ChIV-567_p001-005.tex
\begingroup\makeatletter%
\newcommand{\printedpage}[1]{\marginpar{\footnotesize p.~#1}}%
\newcommand{\msp}[1]{\marginpar{\footnotesize [#1]}}%
\newcommand{\D}{\mathbf{D}}%
\newcommand{\cA}{\mathcal{A}}%
\newcommand{\bDelta}{\boldsymbol{\Delta}}%
\newcommand{\uHom}{\underline{\mathrm{Hom}}}%
\makeatother\let\printedpage\derpp\let\msp\dermsp%
\let\mspin\DERmspin%


\printedpage{}% no page number printed on this page
% head: (none)

\begin{center}
{\LARGE\textsf{LES DÉRIVATEURS}}

\bigskip

\textbf{ALEXANDRE GROTHENDIECK}

\bigskip\bigskip

\textbf{Chapitre IV}

\medskip

\textbf{Diagrammes substantiels}
\end{center}

\par

\printedpage{}% no page number printed on this page
% head: (none)

\noindent \par\noindent\textbf{Note des éditeurs de l’édition transcrite.}\par
Ce texte a été déchiffré et transcrit en \LaTeXe\ par M. Künzer. Il a été édité par M. Künzer, J. Malgoire, et G. Maltsiniotis. La transcription est aussi fidèle que possible au manuscrit. Pour les quelques corrections évidentes, ou rares commentaires des éditeurs, ainsi que pour la numérotation originale des pages du manuscrit, les caractères de machine à écrire \texttt{[typewriter]} entre crochets sont utilisés. Un point d'interrogation entre crochets signifie que l'on n'est pas sûr du mot qui précède.

\noindent Cette édition est provisoire. Les remarques, commentaires et corrections sont bienvenus. Envoyer un message à :

\noindent\textbf{[e-mail address removed]}

\hfill G. Maltsiniotis\hspace*{2cm}

\par

\printedpage{1}
% head: A. Grothendieck : Dérivateurs, Ch. IV, Édition des Universités de Montpellier II et Paris VII
% foot: Version 2 février 2002

\noindent\texttt{[Il n'y a pas de pages 1-72]}

\msp{page 73}

\subsection*{Construction de diagrammes substantiels pour un dérivateur.}

\noindent Soit $\D$ un prédérivateur, $I$ dans Diag, considérons

\noindent $\varphi_I : \mathcal{D}(I) \longrightarrow \uHom(I^o,\cA)$,\ \ où $\cA=\D(e)$ est la catégorie fondamentale des coefficients. % [sic] \mathcal{D}(I) -> \mathbf{D}(I), elsewhere the prederivator is the bold D

Les objets de $\uHom(I^o,\cA)$ s'interprètent comme des \textit{diagrammes de type $I^o$} dans $\cA$. Ce langage est surtout commode si $I$ est ordonnée. Notons, pour $F\in\mathrm{Ob}\,\D(I)$, par $(F_i)_{i\in\mathrm{Ob}\,I}$ le $I^o$-diagramme associé dans $\cA$. Par l'axiome Dér 2, le foncteur
\[
\varphi : F \longmapsto (F_i)_{i\in\mathrm{Ob}\,I}
\]
est conservatif. Mais il n'a aucune envie d'être fidèle, encore moins pleinement fidèle. Quand on connaît le diagramme $(F_i)_i$, on est loin de connaître $F$, pas même à isomorphisme non unique près, je présume. Il n'est nullement clair que si $F$, $F'$ dans $\D(I)$ sont tels que les diagrammes associés soient isomorphes, ils soient eux-mêmes isomorphes.

Pour insister sur ce point, les objets de $\D(I)$ s'appelleront les `\textit{diagrammes substantiels} de type $I^o$ dans $\cA$' (ou $\D$-substantiel, s'il y a risque de confusion). Pour étayer l'intuition et simplifier les explications, un objet $F$ de $\D(I)$ ($\D$-coefficients sur $I$) sera représenté graphiquement par le diagramme `faible' ou `squelettique' correspondant -- on le regarde
\msp{page 74}
comme une sorte de (pâle) `squelette' du diagramme substantiel dont il provient. P. ex. si $I=\bDelta_0 : 0\longrightarrow 1$ \texttt{[plutôt $\bDelta_1$]}, on note un diagramme substantiel de type $(\bDelta_0)^o$ \texttt{[plutôt $(\bDelta_1)^o$]} par le même graphisme, que le diagramme ordinaire :
\[
F_0 \longleftarrow F_1.
\]
Si
\[
I = \begin{tikzcd}
a \arrow[r] & c \\
0 \arrow[u] \arrow[r] & b, \arrow[u]
\end{tikzcd}
\]
les diagrammes substantiels de type $I$ se notent
\[
\begin{tikzcd}
F_a \arrow[d] & F_c \arrow[l] \arrow[d] \\
F_0 & F_b. \arrow[l]
\end{tikzcd}
\]

\par

\printedpage{2}
% head: A. Grothendieck : Dérivateurs, Ch. IV, Édition des Universités de Montpellier II et Paris VII
% foot: Version 2 février 2002

Il faudra bien sûr préciser chaque fois si ce graphisme désigne un diagramme squelettique, ou substantiel. Souvent on part d'un squelette, et on essaye de `l'étoffer' avec un peu de chair, pour en faire un diagramme substantiel -- si possible, de façon canonique. C'est là un point de philosophie dérivatique essentiel pour la pratique. On procédera souvent de proche en proche, en `substantifiant' des sous-diagrammes correspondants à une suite strictement croissante $I_0\subset I_1\subset\cdots\subset I_n=I$ de $I$, au mieux de la situation particulière.

\msp{page 75}

En voici un exemple typique, en supposant $\D$ un dérivateur \textit{ponctué} (les $\D(I)$ sont ponctuées, et les $f^*$ transforment objets nuls en objets nuls). Il s'agit de la définition du foncteur `cosuspension' ou `espace de lacets' dans $\cA$,
\[
\Omega : \cA \longrightarrow \cA.
\]
Il est entendu que pour $M$ dans $\cA$, $\Omega(M)$ doit s'insérer dans un carré commutatif
\[
\begin{tikzcd}
e \arrow[d] & \Omega M \arrow[l] \arrow[d] \\
M & e, \arrow[l]
\end{tikzcd}
\tag*{(C)}
\]
où $e$ désigne les éléments unités. Ce carré cependant devrait être substantiel, plus précisément, être un carré substantiel `$\D$-\textit{cartésien}'. Étant entendu qu'un objet de $\D(I)$ (carré substantiel dans $\cA$) est dit $\D$-cartésien, s'il est dans l'image essentielle de $\D(I_1)$ par $i_{1*}$, où
\[
i_1 : I_1 \hookrightarrow I,
\]
$I_1$ étant la sous-catégorie ouverte de $I$ définie par
\[
I_1 = \begin{tikzcd}
a \\
0 \arrow[u] \arrow[r] & b.
\end{tikzcd}
\]

\msp{page 76}

La restriction du carré substantiel cherché (C) à la sous-catégorie $I_1$ est donc un diagramme substantiel de squelette
\[
\begin{tikzcd}
e \arrow[d] \\
M & e, \arrow[l]
\end{tikzcd}
\tag*{(Sq)}
\]
\printedpage{3}
% head: A. Grothendieck : Dérivateurs, Ch. IV, Édition des Universités de Montpellier II et Paris VII
% foot: Version 2 février 2002
dont il se déduit (à isomorphisme canonique près) en appliquant $i_*$ à ce dernier, de façon à le compléter en $\begin{tikzcd}
e \arrow[d] & \Omega M \arrow[l, dotted] \arrow[d, dotted] \\
M & e \arrow[l]
\end{tikzcd}$. Il reste donc à décrire une substantiation canonique du squelette (Sq), en termes de $M$. Cela nous conduit à regarder la sous-catégorie ouverte
\[
I_0 = \{0\} :\quad
\begin{tikzcd}
{} \arrow[r, dotted] & {} \\
0 \arrow[u, dotted] \arrow[r, dotted] & {} \arrow[u, dotted]
\end{tikzcd}
\]
de $I_1$ (donc aussi de $I$) réduite à l'objet initial, et l'inclusion
\[
I_0 \overset{i_0}{\hookrightarrow} I_1,
\]
et de substantifier le squelette (Sq) par $(i_0)_!(M)$.

On peut dire a posteriori que, si on connaissait le carré substantiel, il n'y a, à isomorphisme unique près, qu'une seule façon de l'obtenir comme un $i_{1*}i_{0!}(?)$, avec ? dans $\D(I_0)\simeq\cA$, car $i_{0!}$ et $i_{1*}$ étant pleinement fidèles, il en est de même de leur composé.

\msp{page 77}

En fait, on voit que le foncteur
\[
i_{1*}i_{0!} :\ \cA\simeq\D(I_0) \longrightarrow \D(I)
\]
induit une équivalence de $\cA$ avec la sous-catégorie pleine de $\D(I)$ formée des carrés $\D$-cartésiens dont les fibres en $a$ et $b$ sont nulles. (Le foncteur quasi-inverse évident est $i^*$, où $i : U_0 \hookrightarrow I$ \texttt{[plutôt $I_0 \hookrightarrow I$]}.) Cela signifie donc que dans ce cas, la donnée du squelette de nature très particulière ($F_a\simeq F_b\simeq 0_{\cA}$), plus la contrainte supplémentaire $\D$-cartésienne, sur le carré cherché, détermine celui-ci à isomorphisme unique près en termes de son squelette : il y a une substantiation non seulement `canonique' de celui-ci, mais même unique à isomorphisme unique près, à cause de la contrainte $\D$-cartésienne.

Je voudrais commencer par regarder le cas plus simple encore où $I=\bDelta_1=(0\longrightarrow 1)$, donc les squelettes envisagés sont simplement les flèches
\[
F_1 \longrightarrow F_0 \qquad \text{dans } \cA.
\]
La philosophie nous oblige, pour les flèches les plus fondamentales dans la théorie des dérivateurs $\D$ (flèches dans $\cA$, ou par extension dans une $\D(I)\simeq\cA_{\D_I}$, catégorie fondamentale du dérivateur induit $\D_I$) de se demander s'il n'y a pas


\endgroup
% END UNIT Der-ChIV-567/Der-ChIV-567_p001-005.tex
% BEGIN UNIT Der-ChIV-567/Der-ChIV-567_p006-009.tex
\begingroup\makeatletter%
\newcommand{\printedpage}[1]{\marginpar{\footnotesize p.~#1}}%
\newcommand{\msp}[1]{\marginpar{\footnotesize [#1]}}%
\newcommand{\Hom}{\operatorname{Hom}}%
\newcommand{\uHom}{\underline{\operatorname{Hom}}}%
\newcommand{\Ob}{\operatorname{Ob}}%
\newcommand{\id}{\mathrm{id}}%
\newcommand{\bDelta}{\boldsymbol{\Delta}}%
\makeatother\let\printedpage\derpp\let\msp\dermsp%
\let\mspin\DERmspin%


\printedpage{4}
% head: A. Grothendieck : Dérivateurs, Ch. IV, Édition des Universités de Montpellier II et Paris VII
% foot: Version 2 février 2002
\msp{page 78}
une façon ‘canonique’, ‘naturelle’, ‘voulue par Dieu’, de le ‘remplumer’ en un diagramme substantiel. Les flèches les plus évidentes qui se présentent sont relatives à un foncteur
\[ f : I \longrightarrow J \]
dans Diag, et aux deux morphismes d'adjonction pour le couple $(f^*, f_*)$, savoir
\[
\begin{aligned}
G &\longrightarrow f_*f^*G && \text{dans } D(J)\\
f^*f_*F &\longrightarrow F && \text{dans } D(I). \quad (\text{\footnotemark})
\end{aligned}
\]
\footnotetext{Bien sûr, problème dual pour $(f_!, f^*)$.}

Je vais essayer de remplumer ces deux flèches, et vais généraliser en envisageant la situation ci-contre \texttt{[ci-dessous]}
\[
\begin{tikzcd}
I \arrow[d, "f"'] & F \\
J & G, \arrow[u, "u"']
\end{tikzcd}
\]
où
\[ \Hom^f(G,F) \overset{\text{déf}}{=} \Hom(f^*G,F) \overset{\text{adj.}}{\simeq} \Hom(G,f_*F). \]
Donc le même $u$ s'interprète par \textit{deux} flèches squelettiques
\[
\left\{
\begin{aligned}
f^*(G) &\longrightarrow F && \text{dans } \mathbf{D}(I)\\
G &\longrightarrow f_*(F) && \text{dans } \mathbf{D}(J).
\end{aligned}
\right.
\tag*{(1)}
\]
Les remplumer, c'est les interpréter (si faire se peut) en termes de deux objets (si possible canoniques) dans $\mathbf{D}(I \times \bDelta_1)$ et $\mathbf{D}(J \times \bDelta_1)$ respectivement.

Je vais me guider sur le cas particulier d'un dérivateur trivial
\[ \mathbf{D} = \mathbf{D}_{\mathcal{M}}, \]
où $\mathcal{M}$ est une catégorie de modèles ‘nue’ $(\Sigma = \emptyset)$.

\msp{page 79}
Pour $f$ fixé, on va construire une catégorie auxiliaire $K(f)$, de telle façon que pour toute catégorie $\mathcal{M}$, les systèmes $(G,F,u)$ d'un $\mathcal{M}$-préfaisceau $G$ sur $J$, d'un autre $F$ sur $I$, et d'un $f$-homomorphisme $G \longrightarrow F$, s'identifient (fonctoriellement) aux $\mathcal{M}$-préfaisceaux sur
\printedpage{5}
% head: A. Grothendieck : Dérivateurs, Ch. IV, Édition des Universités de Montpellier II et Paris VII
% foot: Version 2 février 2002
$K(f)$. Ceci dit, on va essayer d'interpréter les flèches (1) squelettiques, mais peu importe ici -- tous les $\mathbf{D}$-squelettes sont maintenant remplumés) en termes de la géométrie de $K(f)$. % [sic] remplumés) -> remplumés, closing parenthesis printed without any opening parenthesis

\[
\left\{
\begin{aligned}
&\Ob K(f) = \Ob K(I) \amalg \Ob K(J) \qquad \text{\texttt{[plutôt} $\Ob I \amalg \Ob J$\texttt{ ]}}\\
&\text{$I$, $J$ s'identifiant à des sous-catégories pleines de $K(f)$,}\\
&\text{il n'y a pas des flèches d'un $j \in J$ dans un $i \in I$.}\\
&\text{(Donc $I$ est ouverte dans $K(f)$, $J$ son complémentaire fermé.)}\\
&\text{Pour $i \in \Ob I$, $j \in \Ob J$, on pose $\Hom_{K(f)}(i,j) = \Hom_J(f(i),j)$.}
\end{aligned}
\right.
\]

La composition des flèches dans $K(f)$ est évidente. On va définir une rétraction canonique
\[ \bar{f} : K(f) \longrightarrow J \]
de $K(f)$ sur sa sous-catégorie fermée $J$. Donc $\bar{f}$ est connu sur $\Ob J$, et sur $\Ob I$ ce sera $f$ :
\[
\left\{
\begin{aligned}
\bar{f}(j) &= j && \text{si } j \in \Ob J\\
\bar{f}(i) &= f(i) && \text{si } i \in \Ob I.
\end{aligned}
\right.
\]
Il reste à définir, pour $x, y \in \Ob K(f)$,
\msp{page 80}
\[ \Hom_{K(f)}(x,y) \longrightarrow \Hom_J(\bar{f}(x), \bar{f}(y)). \]
Si $x, y \in \Ob J$, c'est l'identité, si $x, y \in I$, c'est $f$. Si $x \in I$, $y \in J$, on doit définir une flèche
\[ \Hom(f(x),y) \longrightarrow \Hom(f(x),y), \]
et on prend encore l'identité. On a
\[
\begin{tikzcd}
J \arrow[r, hook, "j"] \arrow[dr, equal] & K(f) \arrow[d, "\bar{f}"] & I \arrow[l, hook', "i"'] \\
& J &
\end{tikzcd}
\]
\[
\left\{
\begin{aligned}
\bar{f} \circ j &= \id_J\\
j \circ \bar{f} &= \left\{
\begin{aligned}
&j && \text{sur } J\\
&jf && \text{sur } I,
\end{aligned}
\right.
\end{aligned}
\right.
\qquad \text{i.e.} \quad
\boxed{
\begin{aligned}
j\bar{f}j &= j && \text{sur } J\\
j\bar{f}i &= jf && \text{sur } I
\end{aligned}
}\ ,
\]
et on a \texttt{[}$K = K(f)$\texttt{]}
\[
\begin{aligned}
\alpha :\ \id_K &\longrightarrow j\bar{f} && \text{i.e. homomorphisme fonctoriel}\\
x &\xrightarrow{\alpha(x)} j(\bar{f}(x)), && x \in \Ob K,
\end{aligned}
\]
\printedpage{6}
% head: A. Grothendieck : Dérivateurs, Ch. IV, Édition des Universités de Montpellier II et Paris VII
% foot: Version 2 février 2002
défini ainsi :
\[
\left\{
\begin{aligned}
&\text{Si } x \in J,\ \alpha(x) = \id_x\ :\ x \xrightarrow{\sim} j(\bar{f}(x)) = j(x) = x\\
&\text{Si } x \in I,\ \alpha(x) : \underbrace{x \longrightarrow f(x)}_{\text{dans } K(f)},\ \text{ou } \underbrace{f(x) \longrightarrow f(x)}_{\text{dans } J},\ \alpha(x) = \id_{f(x)}.
\end{aligned}
\right.
\]
On doit vérifier la fonctorialité pour $x$ variable dans $K$ \dots

Soit $\Phi$ un $\mathcal{M}$-préfaisceau sur $K$. Il définit
\[
\begin{aligned}
F &= i^*(\Phi) && \in \Ob \mathbf{D}_{\mathcal{M}}(I)\\
G &= j^*(\Phi) && \in \Ob \mathbf{D}_{\mathcal{M}}(J).
\end{aligned}
\]
D'autre part $f^*(G)$ s'interprète comme $(j\bar{f})^*(\Phi)|I$, car
\[
\begin{aligned}
f^*(G)(i) &= G(\ \underbrace{f(i)}_{= \bar{f}(i) \,\in\, J}\ )\\
&= j^*(\Phi)(\bar{f}(i))\\
&= \Phi(j\bar{f}(i))\\
&= (j\bar{f})^*(\Phi)(i).
\end{aligned}
\]

\msp{page 81}
Donc
\[ f^*(G) \simeq (j\bar{f})^*(\Phi)|I. \]
Or la flèche
\[ \alpha : \id_K \longrightarrow j\bar{f} \]
donne
\[ \underbrace{\id_K(\Phi)}_{= \Phi} \xleftarrow{\Phi(\alpha)} (j\bar{f})^*(\Phi), \qquad \text{fonctoriel en } \Phi, \]
d'où, en prenant les restrictions à $I$,
\[ i^*(\Phi) \overset{\text{déf}}{=} F \xleftarrow{\ u\ } f^*G. \]

En fait, sans tellement se tortiller, on devrait pouvoir montrer que pour toute catégorie $\mathcal{M}$, la catégorie $\uHom(K^o, \mathcal{M})$ des $\mathcal{M}$-préfaisceaux sur $K$ est isomorphe à celle des triples $(F,G,u)$, avec $F \in \Ob \uHom(I^o, \mathcal{M})$, $G \in \uHom(J^o, \mathcal{M})$, et $u : f^*(G) \longrightarrow F$. Ainsi, $K$ apparaît comme solution d'un problème universel
\[ \uHom(K^o, \mathcal{M}) \longrightarrow \{F,\ G,\ u \text{ comme dessus}\}(\mathcal{M}). \]

Mais pour y préciser la flèche $\longrightarrow$ de façon conceptuelle simple, il est bon de se donner $K$ avec les structures supplémentaires \begin{tikzcd}[cramped, column sep=normal] I \arrow[r, hook, shift left, "j"] & K \arrow[l, shift left, "\bar{f}"] & I \arrow[l, hook', "i"'] \end{tikzcd}, $\bar{f}j = \id_J$, $\id_K \xrightarrow{\alpha} j\bar{f}$. Et ce système % [sic] I (first object, source of j and target of f-bar) -> J, since j : J -> K(f) and f-bar : K(f) -> J
\printedpage{7}
% head: A. Grothendieck : Dérivateurs, Ch. IV, Édition des Universités de Montpellier II et Paris VII
% foot: Version 2 février 2002
\msp{page 82}
$\{K,\ i,\ j,\ \bar{f},\ \alpha\}$ est (sauf erreur) défini à isomorphisme canonique près par la condition de représenter le foncteur $\mathcal{M} \longmapsto \{F,\ G,\ u\}^o(\mathcal{M})$ \dots

En particulier, le système défini par $G \in \Ob \mathbf{D}_{\mathcal{M}}(J)$
\[ (f^*(G),\ G,\ \id : f^*(G) \longrightarrow f^*(G)) \]
est défini sauf erreur par $\bar{f}^*(G)$ sur $K(f)$. On se propose [?] de trouver une interprétation similaire pour
\[ (F,\ f_*(F),\ f^*f_*F \longrightarrow F) \]
$(F \in \Ob \mathbf{D}(I))$. Le candidat naturel serait ici $\Phi = i_*(F)$. Admettons ce point (i.e. surtout la vérification du fait que $j^*(i_*(F)) \simeq f_*(F)$). Tout ceci montre que si $\mathbf{D}$ est à nouveau un dérivateur quelconque, on peut, par les formules précédentes, trouver un objet canonique $\Phi$ de $\mathbf{D}(K)$, en termes soit d'un $\mathbf{D}$-coefficient $F$ sur $I$ (en prenant $\Phi = i_*(F)$), soit d'un $\mathbf{D}$-coefficient $G$ sur $J$ (en prenant $\Phi = \bar{f}^*(G)$). D'autre part, on peut espérer définir canoniquement
\[
\begin{aligned}
\mathbf{D}(K) &\longrightarrow \mathbf{D}(I \times \bDelta_1)\\
\mathbf{D}(K) &\longrightarrow \mathbf{D}(J \times \bDelta_1)
\end{aligned}
\]
dont les composés avec les foncteurs ‘squelettes’
\[
\begin{aligned}
\mathbf{D}(I \times \bDelta_1) &\longrightarrow \uHom(\bDelta_1^o, \mathbf{D}(I))\\
\mathbf{D}(J \times \bDelta_1) &\longrightarrow \uHom(\bDelta_1^o, \mathbf{D}(J))
\end{aligned}
\]
\msp{page 83}
redonnent (pour $\Phi \in \Ob \mathbf{D}(K)$) les flèches squelettiques
\[
\left\{
\begin{aligned}
f^*(G) &\xrightarrow{\ u\ } F\\
G &\xrightarrow{\ v\ } f_*F
\end{aligned}
\right.
\]
associées à $(F,G,u)$. De façon précise, $f^*(G) \longrightarrow F$ se définit en termes de $\mathbf{D}$, et $\Phi \in \Ob \mathbf{D}(K)$, comme ci-dessus, et $v$ s'en déduit par adjonction.

L'idée naturelle ici, serait de ramener le problème de l'étoffement des flèches d'adjonction sqelettiques, au cas des immersions ouvertes ou fermées. Notons que les sous-catégories ouvertes de $K$ correspondent justement aux foncteurs $K \longrightarrow \bDelta_1$, en prenant les images inverses de l'ouvert générique. Il y a donc un unique % [sic] sqelettiques -> squelettiques, missing u (spelled squelettiques elsewhere in the text)
\[ \varphi :\ K(f) \longrightarrow \bDelta_1 \]
avec $I = \varphi^{-1}(\{0\})$. Donc \texttt{[fin du chapitre]}


\endgroup
% END UNIT Der-ChIV-567/Der-ChIV-567_p006-009.tex
\clearpage\setcounter{footnote}{0}\phantomsection\addcontentsline{toc}{section}{Chapitre V}\markboth{Chapitre V}{Chapitre V}
% BEGIN UNIT Der-ChV-234/Der-ChV-234_p001-005.tex
\begingroup\makeatletter%
\NoAutoSpacing%
\newcommand{\printedpage}[1]{\marginpar{\footnotesize p.~#1}}%
\newcommand{\msp}[1]{\marginpar{\footnotesize [#1]}}%
\newcommand{\Ob}{\operatorname{Ob}}%
\newcommand{\Fl}{\operatorname{Fl}}%
\newcommand{\uCh}{\underline{\mathrm{Ch}}}%
\newcommand{\uHom}{\underline{\mathrm{Hom}}}%
\newcommand{\bPhi}{\boldsymbol{\Phi}}%
\makeatother\let\printedpage\derpp\let\msp\dermsp%
\let\mspin\DERmspin%


\printedpage{}
% head: (none); no page number printed on this page

\begin{center}
\vspace*{6cm}

{\LARGE\textsf{LES DÉRIVATEURS}}

\vspace{1.5cm}

{\large\textbf{ALEXANDRE GROTHENDIECK}}

\vspace{2.5cm}

{\large\textbf{Chapitre V}}

\vspace{0.4cm}

{\large\textbf{Catégories de chemins et localisation}}
\end{center}

\par

\printedpage{}
% head: (none); no page number printed on this page

\vspace*{6cm}

\par\noindent\textbf{Note des éditeurs de l’édition transcrite.}\par
Ce texte a été déchiffré et transcrit en \LaTeXe{} par M. Künzer. Il a été édité par M. Künzer, J. Malgoire, et G. Maltsiniotis. La transcription est aussi fidèle que possible au manuscrit. Pour les quelques corrections évidentes, ou rares commentaires des éditeurs, ainsi que pour la numérotation originale des pages du manuscrit, les caractères de machine à écrire \texttt{[typewriter]} entre crochets sont utilisés. Un point d'interrogation entre crochets signifie que l'on n'est pas sûr du mot qui précède.

Cette édition est provisoire. Les remarques, commentaires et corrections sont bienvenus. Envoyer un message à :

\textbf{[e-mail address removed]}\\
\hspace*{\fill}G. Maltsiniotis\hspace*{2cm}

\par

\printedpage{1}
% head: \textbf{A. Grothendieck :} \textit{Dérivateurs,} Ch. V, Édition des Universités de Montpellier II et Paris VII
% foot: Version 9 septembre 2001

\msp{page 1}

\section*{1\quad La catégorie des chemins}

Soit $\Phi$ la catégorie suivante :
\[
\left\{
\begin{array}{rcll}
\Ob\Phi & = & \mathbf{Z} & \\
\Fl\Phi & = & \Ob\Phi\sqcup\mathbf{Z} & \qquad (\Ob\Phi\leftrightarrow\text{identités})
\end{array}
\right.
\]
On note $\{\alpha\}$, ou simplement $\alpha$, l'objet qui correspond à $\alpha\in\mathbf{Z}$. On note $u_\alpha$ la flèche qui correspond à l'entier $\alpha\in\mathbf{Z}$.
\[
\begin{array}{rcl@{\qquad}l@{\quad}l@{\quad}l}
s(u_\alpha) & = & \alpha & \text{si } \alpha \text{ pair}, & \alpha+1 & \text{si } \alpha \text{ impair}\\[1ex]
b(u_\alpha) & = & \alpha+1 & \text{si } \alpha \text{ pair}, & \alpha-1 & \text{si } \alpha \text{ impair}
\end{array}
\]
\[
\texttt{[plutôt}\qquad b(u_\alpha) \;=\; \alpha+1 \quad \texttt{si}\ \alpha\ \texttt{pair,}\quad \alpha \quad \texttt{si}\ \alpha\ \texttt{impair]}
\]

Il n'y pas à définir de compositions, car si deux flèches sont composables, elles sont l'identité. En fait, $\Phi$ est une catégorie \textit{ordonnée}, $\{\alpha\}<\{\beta\}$ si et seulement si $\alpha$ pair, $\beta=\alpha+1$ \texttt{[plutôt} $\beta=\alpha+1$ \texttt{ou} $\beta=\alpha-1$\texttt{]}. % [sic] Il n'y pas -> Il n'y a pas, verb "a" missing

On note, pour $n\in\mathbf{N}$,
\[
\Phi_n\subseteq\Phi,\quad \Ob\Phi_n=\{\alpha\in\Ob\Phi\mid -n\le\alpha\le+n\},
\]
la sous-catégorie pleine de $\Phi$ dont les objets sont les $\{\alpha\}$ avec $-n\le\alpha\le+n$. Ainsi
\[
\Phi_1=\{-1\longleftarrow 0\longrightarrow 1\},
\]
la duale est
\[
\Psi_1=\Phi_1^{o}=\{-1\longrightarrow 0\longleftarrow 1\}.
\]
On peut considérer $\Phi$ comme somme amalgamée d'une suite ordonnée infinie (de deux côtés) de copies de $\Phi_1$, ou aussi d'une suite similaire de copies de $\Psi_1$.
\[
\begin{tikzcd}[column sep=small, execute at end picture={
\draw[decorate,decoration={brace,amplitude=4pt}] ([yshift=10pt]\tikzcdmatrixname-1-2.center) -- ([yshift=10pt]\tikzcdmatrixname-1-4.center) node[midway,above=4pt]{$\scriptstyle\Phi_1$};
\draw[decorate,decoration={brace,amplitude=4pt}] ([yshift=10pt]\tikzcdmatrixname-1-4.center) -- ([yshift=10pt]\tikzcdmatrixname-1-6.center) node[midway,above=4pt]{$\scriptstyle\Phi_1$};
\draw[decorate,decoration={brace,amplitude=4pt}] ([yshift=10pt]\tikzcdmatrixname-1-6.center) -- ([yshift=10pt]\tikzcdmatrixname-1-8.center) node[midway,above=4pt]{$\scriptstyle\Phi_1$};
\draw[decorate,decoration={brace,amplitude=4pt}] ([yshift=10pt]\tikzcdmatrixname-1-8.center) -- ([yshift=10pt]\tikzcdmatrixname-1-10.center) node[midway,above=4pt]{$\scriptstyle\Phi_1$};
\draw[decorate,decoration={brace,amplitude=4pt}] ([yshift=10pt]\tikzcdmatrixname-1-10.center) -- ([yshift=10pt]\tikzcdmatrixname-1-12.center) node[midway,above=4pt]{$\scriptstyle\Phi_1$};
\draw[decorate,decoration={brace,mirror,amplitude=4pt}] ([yshift=-10pt]\tikzcdmatrixname-1-3.center) -- ([yshift=-10pt]\tikzcdmatrixname-1-5.center) node[midway,below=4pt]{$\scriptstyle\Psi_1$};
\draw[decorate,decoration={brace,mirror,amplitude=4pt}] ([yshift=-10pt]\tikzcdmatrixname-1-5.center) -- ([yshift=-10pt]\tikzcdmatrixname-1-7.center) node[midway,below=4pt]{$\scriptstyle\Psi_1$};
\draw[decorate,decoration={brace,mirror,amplitude=4pt}] ([yshift=-10pt]\tikzcdmatrixname-1-7.center) -- ([yshift=-10pt]\tikzcdmatrixname-1-9.center) node[midway,below=4pt]{$\scriptstyle\Psi_1$};
\draw[decorate,decoration={brace,mirror,amplitude=4pt}] ([yshift=-10pt]\tikzcdmatrixname-1-9.center) -- ([yshift=-10pt]\tikzcdmatrixname-1-11.center) node[midway,below=4pt]{$\scriptstyle\Psi_1$};
}]
\cdots & {} & -4 \arrow[l] \arrow[r] & -3 & -2 \arrow[l] \arrow[r] & -1 & 0 \arrow[l] \arrow[r] & 1 & 2 \arrow[l] \arrow[r] & 3 & 4 \arrow[l] \arrow[r] & {} & \cdots
\end{tikzcd}
\]

Les $\Phi_n$, pour $n$ impair, sont de même somme de copies de $\Phi_1$ mises bout à bout, et les $\Psi_n$, pour $n$ pair $>0$, somme de copies de $\Psi_1$.

\par

\printedpage{2}
% head: \textbf{A. Grothendieck :} \textit{Dérivateurs,} Ch. V, Édition des Universités de Montpellier II et Paris VII
% foot: Version 9 septembre 2001

Les $\Phi_n$ forment un système inductif avec foncteurs de transition pleinement fidèles, dont $\Phi$ est la limite inductive
\[
\Phi=\varinjlim_n\Phi_n.
\]
Mais pour la suite on va faire des $\Phi_n$ un système \textit{projectif}, en définissant pour tout $n\le m$ une \textit{rétraction} $\varphi_{n,m}$ de $\Phi_m$ sur $\Phi_n$, par

\msp{page 2}
\[
\left\{
\begin{array}{l}
\varphi_{n,m}:\Phi_m\longrightarrow\Phi_n\subset\Phi_m\\[1ex]
\varphi_{n,m}(\{\alpha_i\})=
\left\{
\begin{array}{rcll@{\quad}l@{\ \ }r}
\alpha_i & = & i & \in\Ob\Phi_n, & \text{i.e.} & -n\le i\le n\\[1ex]
\alpha_i & = & n & & \text{si} & n\le i\le m\\[1ex]
\alpha_i & = & -n & & \text{si} & -m\le i\le -n
\end{array}
\right.
\end{array}
\right.
\]

\texttt{[plutôt}
\[
\begin{array}{rcl}
\varphi_{n,m}:\Phi_m & \longrightarrow & \Phi_n\subset\Phi_m\\[1ex]
\varphi_{n,m}(\{i\}) & = &
\left\{
\begin{array}{l@{\qquad}l@{\quad}l@{\ \ }r}
\{i\} & \texttt{si}\ i\in\Ob\Phi_n, & \texttt{i.e} & -n\le i\le n\\[1ex]
\{n\} & & \texttt{si} & n\le i\le m\\[1ex]
\{-n\} & & \texttt{si} & -m\le i\le -n\ \texttt{]}
\end{array}
\right.
\end{array}
\]
On note le système projectif d'ensembles ordonnés ainsi obtenu par
\[
\bPhi=(\Phi_n)_{n\in\mathbf{N}}\in\mathrm{Pro}(\mathrm{Cat}).
\]
Soit $X$ une catégorie, alors pour $n\in\mathbf{N}$
\[
\uCh_n(X)\overset{\text{déf}}{=}\uHom(\Phi_n,X)
\]
est la catégorie des diagrammes de $X$ de la forme
\[
A_{-n}\cdots\longleftarrow A_{-2}\xrightarrow{u_{-2}}A_{-1}\xleftarrow{u_{-1}}A_0\xrightarrow{u_0}A_1\xleftarrow{u_1}A_2\xrightarrow{u_2}A_3\longleftarrow\cdots A_n
\tag*{$(*)$}
\]
avec $2n$ flèches bout à bout $u_i$ $(-n\le i\le n)$ \texttt{[plutôt (}$-n\le i\le n-1$\texttt{)]}, deux flèches consécutives allant toujours en sens opposé,
\[
\text{soit}\ \ \cdot\longrightarrow\cdot\longleftarrow\cdot,\ \ \text{soit}\ \ \cdot\longleftarrow\cdot\longrightarrow\cdot\,.
\]
Le foncteur canonique déduit de $\varphi_{n,m}:\Phi_m\longrightarrow\Phi_n$
\[
\varphi_X^{m,n}:\uCh_n(X)\longrightarrow\uCh_m(Y),\qquad n\le m,
\]
\printedpage{3}
% head: \textbf{A. Grothendieck :} \textit{Dérivateurs,} Ch. V, Édition des Universités de Montpellier II et Paris VII
% foot: Version 9 septembre 2001
\texttt{[plutôt} $\uCh_m(X)$\texttt{]} est le `foncteur de dégénérescence', consistant à compléter un diagramme $(*)$, par $m-n$ termes $A_{-n}$ avec des identités à gauche, et par $m-n$ termes \texttt{[}$A_n$\texttt{]} avec des identités à droite, donc pour $n$ pair on trouve
\[
\underbrace{A_{-n}\cdots\xrightarrow{\text{id}}A_{-n}}_{m-n\text{ termes}}\xleftarrow{\text{id}}A_{-n}\longrightarrow A_{-(n-1)}\cdots\cdots A_{n-1}\longleftarrow A_n\xrightarrow{\text{id}}\underbrace{A_n\xleftarrow{\text{id}}A_n\longrightarrow\cdots A_n}_{m-n\text{ termes}}.
\]
C'est un foncteur pleinement fidèle (ce qui correspond au fait que $\varphi_{n,m}$ est surjectif sur les objets et sur les flèches. % [sic] flèches. -> flèches)., opening parenthesis never closed

\msp{page 3}
On pose
\[
\underbrace{\uCh_\infty(X)}_{\text{ou simplement }\uCh(X)}\qquad=\uHom(\bPhi,X)\overset{\text{déf}}{=}\varinjlim_n\uCh_n(X)
\]
(\footnote{L'ensemble des objets de cette catégorie se notera $\mathrm{Ch}_\infty(X)$ ou $\mathrm{Ch}(X)$, de même on définit les $\mathrm{Ch}_n(X)$.}) (\footnote{dans cette formule: $\bPhi$ procatégorie}). On a une inclusion pleinement fidèle
\[
\uCh_\infty(X)\lhook\joinrel\longrightarrow\uHom(\Phi,X),
\]
où la deuxième catégorie est celle des diagrammes infinis (dans les deux directions)
\[
\cdots\longleftarrow A_{-2}\xrightarrow{u_{-2}}A_{-1}\xleftarrow{u_{-1}}A_0\xrightarrow{u_0}A_1\xleftarrow{u_1}A_2\xrightarrow{u_2}\cdots
\]
et où $\uCh_\infty(X)$ s'identifie à la sous-catégorie pleine formées des objets pour lesquels les flèches $u_i$ sont des identités, pour $|i|$ assez grand. % [sic] formées -> formée, agreement with "la sous-catégorie"

On appellera $\uCh_\infty(X)$ la \textit{catégorie des chemins} dans $X$, ses objets seront les chemins dans $X$. Ceux qui proviennent de $\uCh_n(X)$ seront appelés \textit{chemins standards de longueur} $2n$. (Le terme `standard' s'applique ici à nos conventions de notations. Ainsi nous ne considérons que des chemins de longueur paire, et pour une longueur déterminée, nous choisissons l'une parmi les deux espèces raisonnables, par exemple pour les chemins de longueur 2 nous avons pris seulement
\[
A_{-1}\longleftarrow A_0\longrightarrow A_1
\]
et non pas
\[
A_{-1}\longrightarrow A_0\longleftarrow A_1.)
\]

\msp{page 4}
Je me rends compte que mes notations $u_i$ ne sont pas commodes pour la dualité, je vais les changer, en $u_i$ pour les flèches entre objets $\{i\}$ avec $i\ge0$, et $v_i$ la flèche symétrique de $u_i$ par rapport à l'origine. (Le groupe des automorphismes des $\Phi_n$ est réduit au groupe à deux éléments, le seul automorphisme non trivial étant la symétrie $\{i\}\longrightarrow\{-i\}$. Donc nos chemins seront notés désormais % [sic] {-i}. -> {-i}). , opening parenthesis "(Le groupe" not closed on this page
\[
A_{-2}\xrightarrow{v_1}A_{-1}\xleftarrow{v_0}A_0\xrightarrow{u_0}A_1\xleftarrow{u_1}A_2\longrightarrow\cdots.
\tag*{$(*)$}
\]


\endgroup
% END UNIT Der-ChV-234/Der-ChV-234_p001-005.tex
% BEGIN UNIT Der-ChV-234/Der-ChV-234_p006-010.tex
\begingroup\makeatletter%
\newcommand{\printedpage}[1]{\marginpar{\footnotesize p.~#1}}%
\newcommand{\msp}[1]{\marginpar{\footnotesize [#1]}}%
\newcommand{\uCh}{\underline{\mathrm{Ch}}}%
\newcommand{\uHom}{\underline{\mathrm{Hom}}}%
\newcommand{\uFl}{\underline{\mathrm{Fl}}}%
\newcommand{\Hom}{\mathrm{Hom}}%
\makeatother\let\printedpage\derpp\let\msp\dermsp%
\let\mspin\DERmspin%


\printedpage{4}
% head: A. Grothendieck : Dérivateurs, Ch. V, Édition des Universités de Montpellier II et Paris VII
% foot: Version 9 septembre 2001

(Ce n'est pas non plus enthousiasmant \dots)

Soit
\[
\varepsilon=\{-1,+1\}
\]
l'ensemble ordonné \textit{discret,} réduit aux deux éléments $-1$, $+1$, et définissons des foncteurs d'inclusion (pleinement fidèles)
\[
\begin{aligned}
\alpha_n:\varepsilon\longrightarrow\Phi_n\quad & \alpha_n(-1)=-n\\
& \alpha_n(1)=n,
\end{aligned}
\]
ces foncteurs sont compatible avec les $\varphi_{n,m}$, d'où % [sic] compatible -> compatibles, accord avec « ces foncteurs »
\[
\alpha_\infty:\varepsilon\longrightarrow\boldsymbol{\Phi}\in\mathrm{ProCat},
\]
d'où pour tout $X$ un foncteur
\[
\kappa:\uHom(\boldsymbol{\Phi},X)\longrightarrow\uHom(\varepsilon,X)\simeq X\times X
\]
i.e.

\msp{page 5}
\[
\begin{aligned}
\uCh(X) &\longrightarrow X\times X\\
c &\longrightarrow (\sigma(x),\beta(x)).
\end{aligned}
\]

\texttt{[plutôt} $(\sigma(c),\beta(c))$\texttt{]} Pour tout chemin $c$ dans $X$ i.e. un diagramme $(*)$ (satisfaisant à la condition de dégénérescence à l'infini), on a une source ou origine du chemin, et un but ou une \textit{extrémité}
\[
\left\{
\begin{aligned}
\sigma(x) &= A_i\quad\text{pour } i\leq -\mathrm{long}(x)\\
\beta(x) &= A_i\quad\text{pour } i\geq \mathrm{long}(x)
\end{aligned}
\right.
\]
\texttt{[remplacer} $x$ \texttt{par} $c$\texttt{; pour long}$(c)$ \texttt{défini convenablement]}. La fibre du foncteur précédent en $x$, $y$ sera notée
\[
\uCh(X;x,y)\qquad\text{ou simplement }\uCh(x,y),
\]
\textit{catégorie des chemins de $x$} à \textit{$y$.}

\printedpage{5}
% head: A. Grothendieck : Dérivateurs, Ch. V, Édition des Universités de Montpellier II et Paris VII
% foot: Version 9 septembre 2001

\section*{2\quad Calcul d'une catégorie de fractions $X\Sigma^{-1}$}

Avant de continuer le formalisme des chemins (composition etc.), je vais tester l'utilité de la notion de catégorie des chemins, pour le `calcul' des $X\Sigma^{-1}$, où
\[
\Sigma\subset\mathrm{Fl}(X),\qquad\Sigma\text{ contient les identités}
\]
est une partie autrement \textit{quelconque} de $\mathrm{Fl}(X)$. On sait que $\mathrm{Ob}\,X\Sigma^{-1}=\mathrm{Ob}\,X$, et tout revient à définir les
\[
\Hom_{X\Sigma^{-1}}(x,y),\qquad x,y\in\mathrm{Ob}\,X,
\]
\msp{page 6}
et pour $x,y,z\in\mathrm{Ob}\,X$ la composition
\[
\Hom_\Sigma(x,y)\times\Hom_\Sigma(y,z)\xrightarrow{(u,v)\mapsto vu}\Hom_\Sigma(x,z),
\]
(où on note $\Hom_\Sigma$ au lieu de $\Hom_{X\Sigma^{-1}}$, pour alléger les notations).

Grâce à $\Sigma$, on définit une sous-catégorie pleine
\[
\uCh_\Sigma(X)\hookrightarrow\uCh(X),
\]
(\footnote[3]{L'ensemble de ses objets est noté $\mathrm{Ch}(X)$. Ses éléments sont les $\Sigma$-chemins de $x$ à $y$ \texttt{[de} $x$ \texttt{à} $y$\texttt{?]}.}) par la condition sur $c\in\mathrm{Ch}(X)$ : Pour $i$ impair, \texttt{[plutôt pair]} i.e. $i-1\xleftarrow{u_i}i$ ($u_i$ est dite `rétrograde'), $c(u_i)\in\Sigma$, i.e. les `tronçons' rétrogrades dans le chemin $c$ doivent être dans $\Sigma$. On définit de même les sous-catégories pleines
\[
\uCh_\Sigma(X;x,y)=\{c\in\uCh(X;x,y)\mid c\in\uCh_\Sigma(X)\},
\]
ou encore la fibre en $(x,y)$ du foncteur source-but pour $\uCh_\Sigma(X)$
\[
\uCh_\Sigma(X)\longrightarrow X\times X.
\]

On définit une application canonique
\[
\mathrm{Ch}_\Sigma(X,x,y)\longrightarrow\Hom_\Sigma(x,y),
\]
\msp{page 7}
en associant `a tout chemin de $x$ dans $y$ le $\Sigma$-homomorphisme de $x$ dans $y$ obtenu en \textit{inversant} (dans $X\Sigma^{-1}$) les tronçons rétrogrades (ce qui est possible, puisqu'ils sont dans $\Sigma$), de sorte qu'on trouve dans $X\Sigma^{-1}$ un `chemin direct' % [sic] `a -> à, accent imprimé comme guillemet ouvrant
\[
x=A_{-n}\longrightarrow\cdots\longrightarrow A_{-1}\longrightarrow A_0\longrightarrow A_1\longrightarrow A_2\cdots A_n=A_{n+1}=\cdots=y.
\]
et on compose tous ces morphismes.

En fait, on trouve aussitôt que l'on définit ainsi un \textit{foncteur}
\[
\uCh_\Sigma(X)\longrightarrow\uFl(X\Sigma^{-1}),
\]
\texttt{[}$\uFl$ \texttt{catégorie des flèches]} du fait que si on a un diagramme

\printedpage{6}
% head: A. Grothendieck : Dérivateurs, Ch. V, Édition des Universités de Montpellier II et Paris VII
% foot: Version 9 septembre 2001
\[
\begin{tikzcd}
A_{i-1} \arrow[d, "f_{i-1}"] & A_i \arrow[l, "u_i"'] \arrow[d, "f_i"] \\
A'_{i-1} & A'_i \arrow[l, "u'_i"']
\end{tikzcd}
\qquad\qquad (i\text{ pair})
\]
avec $u_i$, $u'_i$ dans $\Sigma$, ce diagramme, en tant que diagramme dans $X\Sigma^{-1}$, est commutatif avec flèches horizontales des \textit{isomorphismes,} et reste donc commutatif en inversant lesdites flèches. D'où un foncteur induit
\[
\uCh_\Sigma(X;x,y)\longrightarrow\uFl(\underbrace{X\Sigma^{-1}}_{X'};x,y).
\]

Mais pour toute catégorie $X'$, et des objets $x$, $y$ dans $X'$, la catégorie $\uFl(X',x,y)$ est
\msp{page 8}
une catégorie \textit{discrète,} ayant comme ensemble d'objets $\Hom_{X'}(x,y)$. Ainsi le foncteur précédant définit en fait une application canonique % [sic] précédant -> précédent, adjectif
\[
\boxed{\pi_0\uCh_\Sigma(X;x,y)\longrightarrow\Hom_\Sigma(x,y).}
\]
Je dis que cette application est toujours bijective. Pour le voir, nous allons \textit{constuire} `from scratch' la catégorie $X\Sigma^{-1}$, en prenant comme objets ceux de $X$, et en \textit{définissant} % [sic] constuire -> construire, lettre omise
\[
\boxed{\Hom_\Sigma(x,y)=\pi_0\uCh_\Sigma(X;x,y).}
\]
Il faut donc définir la composition
\[
\Hom_\Sigma(x,y)\times\Hom_\Sigma(y,z)\longrightarrow\Hom_\Sigma(x,z).
\]
Pour ceci, on va définir des foncteurs `composition de chemins'
\[
\begin{aligned}
\uCh_m(X;x,y)\times\uCh_n(X;y,z) &\longrightarrow \uCh_{m+n}(X;x,z)\\
(c,c') &\longmapsto c'\circ c,
\end{aligned}
\]
mais en nous bornant pour simplifier au cas où $m$, $n$ ont \textit{la même parité,} de telle façon qu'en `mettant bout à bout' $c$ et $c'$ (dans cet ordre), on trouve bien un chemin de type $n+m$.

\msp{page 9}
Cette opération est `transposée' du couple de foncteurs
\[
\begin{aligned}
\Phi_n &\xrightarrow{\alpha} \Phi_{n+m}\quad i\longmapsto (i-m)\\
\Phi_m &\xrightarrow{\beta} \Phi_{n+m}\quad i\longmapsto (i+n),
\end{aligned}
\]
\printedpage{7}
% head: A. Grothendieck : Dérivateurs, Ch. V, Édition des Universités de Montpellier II et Paris VII
% foot: Version 9 septembre 2001
où on aura
\[
\alpha(\{n\})=\beta(\{-m\})=n-m,
\]
et on aura
\[
\Phi_{n+m}\simeq\Phi_n\sqcup_e\Phi_m,
\]
i.e. $\Phi_{n+m}$ est la `somme amalgamée' dans (Cat), obtenue en `recollant' $\Phi_n$ et $\Phi_m$, par identification de l'`élément but' de $\Phi_n$ avec l'`élément source' de $\Phi_m$. Le foncteur composition de chemins induit \texttt{[dans ce qui précède, à partir du début de la page 9, il faut inverser} $m$ \texttt{et} $n$\texttt{]}
\[
\uCh_{m,\Sigma}(x,y)\times\uCh_{n,\Sigma}(y,z)\longrightarrow\uCh_{m+n,\Sigma}(x,z),
\]
d'où en passant aux $\pi_0$
\[
\Hom_{\Sigma,m}(x,y)\times\Hom_{\Sigma,n}(y,z)\longrightarrow\Hom_{\Sigma,m+n}(x,z),
\]
où on dénote, pour $a,b\in\mathrm{Ob}\,X$, par
\[
\Hom_{\Sigma,n}(a,b)\overset{\text{déf}}{=}\pi_0(\uCh_{m,\Sigma}(a,b)).
\]

\texttt{[plutôt} $\uCh_{n,\Sigma}$\texttt{]}

\msp{page 10}
Pour en déduire une application sur les $\varinjlim$
\[
\underbrace{\Hom_\Sigma(x,y)}_{=\ \varinjlim\ \Hom_{\Sigma,m}(x,y)}\quad\times\quad\underbrace{\Hom_\Sigma(y,z)}_{=\ \varinjlim\ \Hom_{\Sigma,n}(y,z)}\quad\longrightarrow\quad\underbrace{\Hom_\Sigma(x,z)}_{=\ \varinjlim\ \Hom_{\Sigma,n+m}(x,z)}
\]
il faut prouver, pour $n\leq n'$, $m\leq m'$, la commutativité dans
\[
\begin{tikzcd}
\Hom_{\Sigma,m}(x,y)\times\Hom_{\Sigma,n}(y,z) \arrow[r] \arrow[d] & \Hom_{\Sigma,m+n}(x,y) \arrow[d] \\
\Hom_{\Sigma,m'}(x,y)\times\Hom_{\Sigma,n'}(y,z) \arrow[r] & \Hom_{\Sigma,m'+n'}(x,y)
\end{tikzcd}
\]
\texttt{[plutôt} $\Hom_{\Sigma,m+n}(x,z)$\texttt{,} $\Hom_{\Sigma,m'+n'}(x,z)$\texttt{]} En d'autres termes, que pour $c:x\longrightarrow y$ et $c':y\longrightarrow z$ deux $\Sigma$-chemins composables de longueur $m$ et $n$ respectivement, le composé $c'c:x\longrightarrow z$ ne change pas, `à homotopie près' dans $\uCh_\Sigma(x,z)$, quand on remplace $c$, $c'$ par des chemins $\bar{c}$, $\bar{c}'$ de longueur $m'$, $n'$ déduits de $c$ et $c'$ par `dégénérescence aux bouts'.

Plus précisement, % [sic] précisement -> précisément, accent omis
\[
\bar{c}'\bar{c}\sim\overline{c'c},
\]
le chemin composé est homotope à celui déduit de $c'c$ par `dégénérescence aux bouts'.

Je viens de regarder pour $n=m=1$, $n'=m'=2$, et de constater qu'il n'y a pas de flèche naturelle entre $\bar{c}'\bar{c}$ et $\overline{c'c}$ - l'homotopie cherchée doit
\printedpage{8}
% head: A. Grothendieck : Dérivateurs, Ch. V, Édition des Universités de Montpellier II et Paris VII
% foot: Version 9 septembre 2001
\msp{page 11}
se définir à l'aide d'un `chemin' dans $\uCh_{\Sigma,m',n'}(X;x,z)$. Dans le cas général, on va en fait trouver une homotopie
\[
h=h_{(n,m),(n',m')}
\]
`rendant commutatif' le diagramme de foncteurs
\[
\begin{tikzcd}[column sep=large]
\mathcal{A}=\uCh_{\Sigma,m}(X;x,y)\times\uCh_{\Sigma,n}(X;y,z) \arrow[r, "*_{m,n}"] \arrow[d, "i_{m,m'}\times j_{n,n'}"] & \uCh_{\Sigma,m+n}(X;x,z) \arrow[d, "k_{m+n,m'+n'}", ""'{name=K}] \\
\uCh_{\Sigma,m'}(X;x,y)\times\uCh_{\Sigma,n'}(X;y,z) \arrow[r, "*_{m',n'}", ""{name=S}] & \uCh_{\Sigma,m'+n'}(X;x,z)=\mathcal{B},
\arrow[from=K, to=S, bend right, "h"']
\end{tikzcd}
\]
où les $i_{\alpha,\beta}$, $j_{\alpha,\beta}$, $k_{\alpha,\beta}$ sont les foncteurs de transition
\[
\uCh_{\Sigma,\alpha}(X;a,b)\longrightarrow\uCh_{\Sigma,\beta}(X;a,b)
\]
$(\alpha\leq\beta)$ pour les $(x,y)$, $(y,z)$ et $(x,z)$ (en guise de $(a,b)$) respectivement. L'homotopie $h$ sera un chemin dans la catégorie $\uHom(\mathcal{A},\mathcal{B})$, allant de
\[
k_{m+n,m'+n'}\circ{*}_{m+n}\longrightarrow{*}_{m'+n'}\circ(i_{m,m'}\times j_{n,n'}).
\]
L'existence d'une telle homotopie signifie ni plus ni moins que le diagramme carré ci-dessus est commutatif, en tant que diagramme dans la catégorie (Cat\,hot) des `catégories à homotopie près,' dont les objets sont
\msp{page 12}
les (petites) catégories, et où on définit
\[
\Hom_{\mathrm{Cat\,hot}}(\mathcal{A},\mathcal{B})\overset{\text{déf}}{=}\underbrace{\pi_0(\uHom(\mathcal{A},\mathcal{B}))}_{\begin{array}{@{}l@{}}\text{classes d'homotopie}\\ \text{de foncteurs de }\mathcal{A}\\ \text{dans }\mathcal{B}\end{array}}\,.
\]
C'est cette conclusion qui sera pour nous utile, et non la description explicite du chemin $h$. Cette existence admise (cf. ci-dessous la preuve), on voit donc que si on regarde, pour $a,b\in\mathrm{Ob}\,X$, le système des catégories $\uCh_{\Sigma,r}(X;a,b)$ $(r\in\mathbf{N})$ comme un \textit{ind-objet de} Cat\,hot\textit{,} qu'on désigne par
\[
\uCh^{\mathrm{ind}}_{\Sigma}\,(X;a,b)\in\mathrm{Ob}\,\mathrm{Ind}(\mathrm{Cat\,hot}),
\]
on aura un morphisme canonique dans Ind(Cat\,hot)
\[
\uCh^{\mathrm{ind}}_{\Sigma}\,(X;x,y)\times\uCh^{\mathrm{ind}}_{\Sigma}\,(X;y,z)\xrightarrow{\ *\ }\uCh^{\mathrm{ind}}_{\Sigma}\,(X;x,z),
\]
appelé \textit{`ind-morphisme de composition de chemins.'} On aurait aimé le déduire d'un honnête \textit{foncteur} `composition de chemins'
\[
\begin{aligned}
\uCh_\Sigma(X;x,y)\times\uCh_\Sigma(X;y,z) &\xrightarrow{\ ?\ } \uCh_\Sigma(X;x,z),\\
\text{qu'on noterait }(c,c') &\longmapsto c'\circ c\text{ ou }c'\cdot c,
\end{aligned}
\]
mais je doute qu'un tel foncteur existe,


\endgroup
% END UNIT Der-ChV-234/Der-ChV-234_p006-010.tex
% BEGIN UNIT Der-ChV-234/Der-ChV-234_p011-015.tex
\begingroup\makeatletter%
\newcommand{\printedpage}[1]{\marginpar{\footnotesize p.~#1}}%
\newcommand{\msp}[1]{\marginpar{\footnotesize [#1]}}%
\newcommand{\Ch}{\underline{\mathrm{Ch}}}%
\newcommand{\Hom}{\underline{\mathrm{Hom}}}%
\newcommand{\Fl}{\mathrm{Fl}}%
\newcommand{\card}{\mathrm{card}}%
\newcommand{\uHom}{\underline{\mathrm{Hom}}}%
\makeatother\let\printedpage\derpp\let\msp\dermsp%
\let\mspin\DERmspin%


\printedpage{9}
% head: A. Grothendieck : Dérivateurs, Ch. V, Édition des Universités de Montpellier II et Paris VII

\msp{page 13}
qui par passage aux ind-objets redonne le morphisme précédent, ou, ce qui revient au même, tel que pour tout $m,n\in\mathbf{N}$ (de même parité), il existe un $N\geq m+n$ assez grand tel que l'on ait un diagramme
\[
\begin{tikzcd}
\Ch_{\Sigma,m}(X;x,y)\times\Ch_{\Sigma,n}(X;y,z) \arrow[r,"*_{m,n}"] \arrow[dd,hook] \arrow[dr] & \Ch_{\Sigma,m+n}(X;x,z) \arrow[d,hook] \\
 & \Ch_{\Sigma,N}(X;x,z) \arrow[d,hook] \\
\Ch_{\Sigma}(X;x,y)\times\Ch_{\Sigma}(X;y,z) \arrow[r,"?"] & \Ch_{\Sigma}(X;x,z)
\end{tikzcd}
\]
et où le triangle supérieur est \textit{commutatif dans} Cat\,hot. (Il faut supposer que pour $N$ assez grand, le composé
\begin{tikzcd}[cramped, sep=small]
{} \arrow[d,hook] & \\
{} \arrow[r,"?"] & {}
\end{tikzcd}
se factorise par $\Ch_{\Sigma,N}(X;x,z)$ \dots)
Pour voir que les foncteurs composition $*_{m,n}$ définissent un homomorphisme des ind-objets dans (Cat\,hot), il suffit de le prouver dans les cas particuliers $(m',n')=(m+1,n+1)$. On va faire la figure en supposant $m$, $n$ \textit{impairs.}

Je m'aperçois ici que je me suis canulé stupidement, la composition de deux chemins de longueur \textit{impaire $m$, $n$ n'est}

\msp{page 14}
\textit{pas un chemin} de longueur $m+n$, i.e. un diagramme de type $\Phi_n$, mais un diagramme de type \textit{dual} $\Phi_n^o=\Psi_n$. Par exemple, pour $m=n=1$
\[
\underbrace{A_0\longleftarrow A_1\longrightarrow A_2}_{c}=\underbrace{B_0\longleftarrow B_1\longrightarrow B_2}_{c'}
\]
est de type $\Phi_2^o$! Donc pour ne pas trop compliquer la vie, on va dans les développements précédents se borner à $m$, $n$ \textit{pairs,} donc des chemins de longueur au moins 4. C'est un peu ridicule! On est ramené à prouver la commutativité à homotopie près dans le cas où on laisse $c'$ fixé et où on envoie $\Ch_{\Sigma,m}(X;x,y)$ dans $\Ch_{\Sigma,m+2}(X;x,y)$, ou inversement $c$ fixe et $\Ch_{\Sigma,n}(X;y,z)$ dans $\Ch_{\Sigma,n+2}$ \texttt{[lire:}$\Ch_{\Sigma,n+2}(X;y,z)$\texttt{]}. Prenons le premier cas, et notons un chemin $c$ sans indices négatifs (ce n'est pas une notation bien heureuse!)
\[
A_0\longrightarrow A_1\longleftarrow A_2\cdots\longleftarrow A_{2m}
\]
(NB à cause de la condition $n$ pair, tous nos chemins commencent désormais par une flèche de $A_0$ dans $A_1$, jamais en sens inverse).

Je vois que l'explicitation de l'homotopie, sans offrir de difficulté, est très fastidieuse à expliquer. Je vais essayer de donner une démonstration conceptuelle, sans `calcul' fastidieux.

% foot: Version 9 septembre 2001

\printedpage{10}
% head: A. Grothendieck : Dérivateurs, Ch. V, Édition des Universités de Montpellier II et Paris VII

\msp{page 15}
Les deux foncteurs composés
\[
f=k\circ *_{m,n}\ ,\quad g=*_{m',n'}\circ(i\times j)
\]
dans le carré page 11, dont il faut prouver qu'ils sont homotopes dans $\uHom(\mathcal{A},\mathcal{B})$, peuvent être vus comme obtenus de la façon suivante, en termes de deux foncteurs
\[
\varphi,\psi:\Phi_{m'+n'}\rightrightarrows\Phi_{m+n}, \tag*{(1)}
\]
d'où deux foncteurs `transposés'
\[
\begin{tikzcd}
\uHom(\Phi_{m,n},X) \arrow[r,shift left,"F"] \arrow[r,shift right,"G"'] & \uHom(\Phi_{m',n'},X)
\end{tikzcd} \tag*{(2)}
\]
\texttt{[plutôt} $\Phi_{m+n}$\texttt{,} $\Phi_{m'+n'}$\texttt{]} lesquels foncteurs `induisent' $f$ et $g$. De façon précise, désignant par $a$, $b$ les objets source et but de $\Phi_{m+n}$, par $c$ l'objet qui correspond au but du chemin $c$ de longueur $m$ \texttt{[trop de} $c$\texttt{'s]}, à la source du chemin $c'$ de longueur $n$ qu'il s'agit de composer, enfin par $a_0$, $b_0$ les éléments source et but de $\Phi_{m'+n'}$, les foncteurs envisagés (1) appliquent $a_0$ dans $a$, $b_0$ dans $b$, donc définissent un diagramme commutatif
\[
\begin{tikzcd}
\Phi_{m'+n'} \arrow[r,shift left,"\varphi"] \arrow[r,shift right,"\psi"'] & \Phi_{m+n} \\
\underset{\text{\footnotesize catégorie discrète}}{\varepsilon=\{a_0,b_0\}.} \arrow[u,hook,"\alpha'"] \arrow[ur,hook,"\alpha"] &
\end{tikzcd} \tag*{(3)}
\]
Ceci dit, la catégorie but $\mathcal{B}$ pour $f$, $g$, à savoir $\Ch_{\Sigma,m'+n'}(X;x,z)$, s'interprète comme la sous-catégorie pleine de $\uHom(\Phi_{m+n},X)$ formée

\msp{page 16}
des foncteurs $c':\Phi_{m'+n'}\longrightarrow X$ dont le composé avec $\alpha'$ est le foncteur
\[
\beta:\varepsilon\longrightarrow X\qquad\qquad
\begin{aligned}
a_0&\longrightarrow x\\
b_0&\longrightarrow z,
\end{aligned} \tag*{$(*)$}
\]
et de même la catégorie source $\mathcal{A}$ s'interprète comme la sous-catégorie pleine de $\uHom(\Phi_{m+n},X)$ formée des $c:\Phi_{m+n}\longrightarrow X$ dont la restriction à $\{a,b,c\}$ est le foncteur
\[
a\longrightarrow x,\ b\longrightarrow y,\ c\longrightarrow z.
\]
\texttt{[plutôt} $b\longrightarrow z$\texttt{,} $c\longrightarrow y$\texttt{]} Elle est contenue dans la sous-catégorie pleine $\mathcal{A}'$ des foncteurs dont la restriction à $\{a,b\}$ soit $a\longrightarrow x,\ b\longrightarrow z$, i.e. dont le composé avec $\alpha$ soit le foncteur $\beta:\varepsilon\longrightarrow X$ déjà envisagé ci-dessus. Ainsi les foncteurs $f,g:\mathcal{A}\rightrightarrows\mathcal{B}$ se factorisent par $\mathcal{A}'$
% foot: Version 9 septembre 2001
\printedpage{11}
% head: A. Grothendieck : Dérivateurs, Ch. V, Édition des Universités de Montpellier II et Paris VII
\[
\begin{tikzcd}
\mathcal{A} \arrow[r,hook] \arrow[dr,shift left,"f"] \arrow[dr,shift right,"g"'] & \mathcal{A}' \arrow[d,shift right,"g'"'] \arrow[d,shift left,"f'"] \\
 & \mathcal{B},
\end{tikzcd}
\]
foncteurs déduits de $\varphi$, $\psi$ ci-dessus, donnant lieu aux triangles commutatifs (3), et du foncteur $\beta:\varepsilon\longrightarrow X$, de la façon dite. Pour prouver que $f$ et $g$ sont homotopes, il suffit de le prouver pour $f'$, $g'$. Et pour ceci, il suffit de montrer que $\varphi$ et $\psi$ sont homotopes en tant qu'objets de la catégorie $\uHom(\Phi_{m+n},\Phi_{m'+n'};\alpha)$ formée des foncteurs dont la restriction à $\varepsilon$ est $\alpha$ \texttt{[plutôt de la catégorie} $\uHom(\Phi_{m'+n'},\Phi_{m+n};\alpha')$ \texttt{formée des foncteurs dont la restriction à} $\varepsilon$ \texttt{est} $\alpha'$ \texttt{]}.

\msp{page 17}
D'ailleurs, les deux assertions (sur $f'$, $g'$, quels que soient $X$, $x$, $y$, $z$, \texttt{[plutôt} $x$\texttt{,} $z$\texttt{]} et sur $\varphi$, $\psi$, $\alpha$ \texttt{[plutôt} $\alpha'$\texttt{]}) sont en fait équivalentes comme on voit en faisant $X=\Phi_{m+n}$, $x=a$, $y=c$, $z=b$ \texttt{[plutôt} $x=a$\texttt{,} $z=b$\texttt{]}. Ceci sera contenu dans le lemme plus général. (\footnote[4]{Dans les considérations qui suivent, on considère en fait tacitement que les foncteurs envisagés sont \textit{croissants} sur les sommets pour l'ordre de leur numérotation. Donc il faudrait écrire $\uHom_{cr}$ au lieu de $\uHom$, dans les énoncés qui suivent, y compris dans le lemme de connexité ci-contre.})

\textbf{Lemme.} \textit{Soient $k$ et $l$ deux entiers. Alors la sous-catégorie pleine de $\uHom(\Phi_k,\Phi_l)$ formée des foncteurs $\Phi_k\longrightarrow\Phi_l$ appliquant source en source, but en but, est connexe.} (NB Elle est non vide si et seulement si $k\geq l$.)

On va décrire entièrement ces foncteurs (= applications croissantes) $\Phi_k\xrightarrow{\varphi}\Phi_l$, en identifiant $\Fl(\Phi_k)$ \texttt{[plutôt l'ensemble des flèches non identiques de} $\Phi_k$\texttt{]} à l'intervalle d'entiers $[1,2k]$, et en associant à $\varphi$ la partie
\[
E_\varphi\subset\Fl(\Phi_k)=[1,2k]\qquad\qquad E_\varphi=\{u\in\Fl(\Phi_k)\mid\varphi(u)\text{ une identité}\}
\]
formée des $u\in\Fl(\Phi_k)$ qui sont `contractées' par $\varphi$ en une identité. Si $E$ est un ensemble d'entiers, on y introduit la notion de composantes connexes (par la relation d'équivalence dans $E$ engendrée par la relation de succession $y=x+1$). On voit aussitôt que les composantes connnexes de $E_\varphi$ qui ne contiennent ni la première ni la dernière flèche de $\Phi_k$ sont paires, et que % [sic] connnexes -> connexes, triple n
\[
\underbrace{\card\Fl(\Phi_k)}_{2k}=\card\,(E_\varphi)+\underbrace{\card\Fl(\Phi_l)}_{2l},
\]

\msp{page 18}
i.e.
\[
\card\,(E_\varphi)=2(k-l).
\]

% foot: Version 9 septembre 2001

\printedpage{12}
% head: A. Grothendieck : Dérivateurs, Ch. V, Édition des Universités de Montpellier II et Paris VII

Il est immédiat qu'un foncteur $\varphi$ est surjectif sur les objets et sur les flèches, donc identifie $\Phi_l$ à une catégorie quotient de $\Phi_k$ - et la relation d'équivalence dans $\Phi_k$ définie par $\varphi$ est celle obtenue en identifiant tous les objets qui sont source ou but de flèches d'une même composante connexe $I$ de $E_\varphi$
\[
\cdots\longleftarrow\ \underbrace{\longrightarrow\ \longleftarrow\ \longrightarrow\ \longleftarrow}_{I}\ \longrightarrow\cdots
\]
et les flèches dans cette composante à l'identité correspondante. Les conditions énoncées sur $E$ (indépendamment d'un $\varphi$ qui lui donnerait naissance assurent que la catégorie quotient est isomorphe (avec \textit{préservation} de source et de but) soit à $\Phi_l$, soit à la duale $\Psi_l=\Phi_l^o$, suivant la direction de la première flèche de $\bar{\Phi}_k$ (quotient de $\Phi_k$). Cette `direction initiale' est la même que celle pour $\Phi_k$ si la première flèche de $\Phi_k$ n'appartient pas à $E$, ou bien si elle y appartient et sa composante connexe dans $E$ est de cardinal pair. Dans le cas inverse ($u\in E$ et sa composante connnexe dans $E$ est de cardinal impair), la direction initiale de $\bar{\Phi}_k$ est opposée à celle de $\Phi_k$. % [sic] "(indépendamment d'un $\varphi$ qui lui donnerait naissance" -> parenthèse fermante manquante, the opening parenthesis is never closed; [sic] connnexe -> connexe, triple n

\msp{page 19}
On voit donc que $E$ doit satisfaire encore à une condition supplémentaire, pour que le $\bar{\Phi}_k$ correspondant soit isomorphe à $\Phi_l$, i.e. pour que $E$ provienne d'un foncteur $\varphi$. Je résume toutes les conditions

\begin{itemize}
\item[a)] $\card E=2(k-l)$.
\item[b)] Les composantes connexes de $E$ qui ne contiennent pas de flèche initiale ou finale sont de cardinal pair.
\item[c)] Si $\Phi_k$ et $\Phi_l$ ont même direction initiale ($k$ et $l$ de même parité), alors l'éventuelle composante de $E$ contenant la flèche initiale doit être elle aussi de cardinal pair (ce qui, à cause de a) \texttt{[plutôt a) et b)]}, équivaut à la même condition pour l'éventuelle composante connexe contenant la flèche finale de $\Phi_k$). Si $\Phi_k$ et $\Phi_l$ ont des directions initiales différentes, il faut que la flèche initiale de $\Phi_k$ appartienne à une composante connexe de cardinal \textit{impair} de $E$ (ce qui équivaut à itou pour la flèche finale).
\end{itemize}

D'autre part, le foncteur $\varphi$ est uniquement déterminé par la connaissance de $E$, car il y a un \textit{unique} isomorphisme $\bar{\Phi}_k\xrightarrow{\sim}\Phi_l$ compatible avec sources et buts, à cause du

\textbf{Lemme} (immédiat) : \textit{Un automorphisme de $\Phi_l$ qui est l'identité sur l'objet source (ou sur l'objet but) est l'identité partout.}

Les foncteurs $\varphi$, $\psi$ de tantôt correspondaient au choix, l'un de $E_\varphi=\{1,2,k-1,k\}$ \texttt{[plutôt} $E_\varphi=\{1,2,2k-1,2k\}$\texttt{]}, l'autre de $E_\psi=\{1,2,i,i+1\}$, 2 composantes connexes

\msp{page 20}
dans un cas, trois \texttt{[plutôt deux]} dans l'autre.) % [sic] "l'autre.)" -> parenthèse fermante sans parenthèse ouvrante, no matching opening parenthesis

% foot: Version 9 septembre 2001

\printedpage{13}
% head: A. Grothendieck : Dérivateurs, Ch. V, Édition des Universités de Montpellier II et Paris VII

Dans le cas général, il y a lieu d'expliciter la relation d'ordre de l'ensemble $\uHom^{!}(\Phi_k,\Phi_l)$ (! signifie : compatible avec source et but), en termes des ensembles représentatifs $E_\varphi$, $E_\psi$. Je ne vais pas expliciter cette relation d'ordre, mais seulement une propriété de celle-ci qui nous sera suffisante. Soit $\varphi_E$ le foncteur défini par une partie $E\subset[1,2k]$. Tout d'abord :

\textbf{Lemme 1} (immédiat) : \textit{Soient $E,E'\subset[1,2k]$ admissibles, et soit $j\in[1,2k]$ tel que $E\cap[1,j]=E'\cap[1,j]$. Alors $\varphi_E$ et $\varphi_{E'}$ coïncident sur le but de $u_j$ ($j$ une flèche de $\Phi_k$) et sur tous les objets de $\Phi_k$ antérieurs à celle-ci (au sens de la numérotation des objets, et non de la relation d'ordre de la catégorie $\Phi_k$ bien sûr.) Dualement, si $E\cap[j,2k]=E'\cap[j,2k]$, alors $\varphi_E$ et $\varphi_{E'}$ coïncident sur la source de $y$, et sur tous les objets de $\Phi_k$ postérieurs à celui-ci.} % [sic] la source de $y$ -> la source de $u_j$, apparently intended by duality with "le but de $u_j$"

Considérons maintenant une partie connexe $I$ de $E$ (pas forcément une composante connexe), $I=[r,s]\subset[1,2k]$ ($r$ le plus petit élément de $I$, $s$ son plus grand élément) et supposons que

\begin{itemize}
\item[a)] $\card I$ est pair (i.e. $s-r$ est pair) \texttt{[plutôt} $s-r$ \texttt{impair]}
\item[b)] $r\neq1$, et $r-1\notin E$ (i.e. $r-1\in[1,2k]\backslash E$)
\end{itemize}

et soit
\[
E'=(E\backslash\{s\})\cup\{r-1\}=E\backslash I\cup\overbrace{[r-1,s-1]}^{I'}
\]
déduit de $I$ en `translatant' sur la gauche d'un cran.

\msp{page 21}
Il est immédiat que si $E$ est admissible, $E'$ aussi (et inversement). Ceci dit, on a le

\textbf{Lemme 2.} \textit{$E$ et $E'$ étant comme ci-dessus, il existe une homomorphisme de foncteurs entre $\varphi_E$ et $\varphi_{E'}$} (je ne dis pas dans quel sens va la flèche!). % [sic] une homomorphisme -> un homomorphisme, wrong gender

\textbf{Corollaire.} Énoncé dual, quand $E'$ est déduit de $E$ en translatant à droite d'un cran un intervalle connexe de longueur paire $I=[r,s]$ de $E$, $s\neq2k$ et $s+1\notin E$.

En effet, dans ce cas $E$ est déduit de $E'$ par une opération du type explicité dans le lemme 2.

La connexité de la catégorie $\uHom^{!}(\Phi_k,\Phi_l)$ est maintenant claire. Si $k$ et $l$ ont même parité, en poussant les composantes connexes de $E$ vers la gauche, on écrit $E$ sous la forme `normale'
\[
E_0=[1,2(k-l)].
\]
Si $k$ et $l$ sont de parité inverse, on ne peut pousser toute la composante connexe de $E$ qui contient l'élément dernier $2k$, car elle est de cardinal impair. Il faut laisser $2k$ dans $E$ et pousser le reste, et on ramène $E$ à la forme normale
\[
E_0=[1,2(k-l)-1]\cup\{2k\}.
\]

\hfill cqfd!

% foot: Version 9 septembre 2001


\endgroup
% END UNIT Der-ChV-234/Der-ChV-234_p011-015.tex
% BEGIN UNIT Der-ChV-234/Der-ChV-234_p016-020.tex
\begingroup\makeatletter%
\newcommand{\printedpage}[1]{\marginpar{\footnotesize p.~#1}}%
\newcommand{\msp}[1]{\marginpar{\footnotesize [#1]}}%
\newcommand{\uCh}{\underline{\mathrm{Ch}}}%
\newcommand{\uHom}{\underline{\mathrm{Hom}}}%
\newcommand{\uFl}{\underline{\mathrm{Fl}}}%
\newcommand{\bDelta}{\boldsymbol{\Delta}}%
\newcommand{\Hom}{\mathrm{Hom}}%
\makeatother\let\printedpage\derpp\let\msp\dermsp%
\let\mspin\DERmspin%


\printedpage{14}
% head: A. Grothendieck : Dérivateurs, Ch. V, Édition des Universités de Montpellier II et Paris VII
Ce qui précède justifie notre raisonnement, qui donnait par composition des chemins des applications de composition bien définies
\[
\underbrace{\Hom_\Sigma(x,y)}_{=\pi_0\uCh_\Sigma(x,y)} \times \underbrace{\Hom_\Sigma(y,z)}_{=\pi_0\uCh_\Sigma(y,z)} \longrightarrow \underbrace{\Hom_\Sigma(x,z)}_{=\pi_0\underline{\textstyle\mathrm{Ch}}_\Sigma(x,z)}.
\]
% (the third "Ch" under the brace is printed in a larger size than the two others)

\msp{page 22}Je dis que cette composition est associative. Ceci résulte du fait que si on a trois chemins (d'indices \textit{pairs}!)
\[
x \xrightarrow{c} y \xrightarrow{c'} z \xrightarrow{c''} t,
\]
on a (sans avoir à effectuer des homotopies \dots) associativité stricte
\[
c''(c'c) = (c''c')c.
\]

Je dis qu'il y a des identités. On définit
\[
1_x \in \Hom_\Sigma(x,x)
\]
par le chemin d'indice 0 (qui est pair!), i.e. par le diagramme réduit à $x$, sans flèche (ou si on veut interpréter le diagramme comme un foncteur $e \longrightarrow X$, il y a la flèche identique sous-entendue). Nos raisonnements n'excluaient pas le cas $n = 0$ ou $m = 0$, et il est tautologique que pour un chemin
\[
x \xrightarrow{c} y
\]
on a
\[
c \cdot 1_x = 1_y \cdot c = c.
\]
Donc on a bien défini une catégorie, que je vais appeler $X_\Sigma$. Ja vais définir un foncteur canonique % [sic] Ja -> Je, misprint of the pronoun
\[
\theta : X \longrightarrow X_\Sigma
\]
qui sera l'identité sur les objets, il faut le définir sur les flèches
\[
\Hom_X(x,y) \longrightarrow \Hom_{X_\Sigma}(x,y) = \pi_0\mathrm{Ch}_\Sigma(x,y)
\]
\texttt{[plutôt} $\pi_0\uCh_\Sigma(x,y)$\texttt{]}. Malheureusement, une flèche
\[
x \xrightarrow{u} y
\]
\msp{page 23}ne peut s'interpréter directement comme un chemin à notre sens, il faut payer nos fuites! Mais on va à une telle flèche associer le chemin
\[
x \longrightarrow y \xleftarrow{\mathrm{id}} y \xrightarrow{\mathrm{id}} y \xleftarrow{\mathrm{id}} y,
\]
% foot: Version 9 septembre 2001
\printedpage{15}
% head: A. Grothendieck : Dérivateurs, Ch. V, Édition des Universités de Montpellier II et Paris VII
élément de $\uCh_{\Sigma,2}(x,y)$, ce qui définit l'application cherchée. Il faut vérifier la compatibilité avec les identités et sous la composition. À la flèche $1_x : x \longrightarrow x$ est associé l'élément de $\Phi_2(x,x)$
\[
x \xrightarrow{\mathrm{id}} x \xleftarrow{\mathrm{id}} x \xrightarrow{\mathrm{id}} x \xleftarrow{\mathrm{id}} x,
\]
il faut prouver que dans $\pi_0(\uCh_{\Sigma,2}(x,x))$ il devient égal à l'élément déduit de $1_x \in \Phi_{\Sigma,0}(x,x)$, par le foncteur de transition. Mais cet élément n'est autre justement que
\[
x \xrightarrow{\mathrm{id}} x \xleftarrow{\mathrm{id}} x \xrightarrow{\mathrm{id}} x \xleftarrow{\mathrm{id}} x,
\]
c'est le même (sans avoir même à passer au $\pi_0$)!

Pour la compatibilité avec la composition
\[
x \xrightarrow{u} y \xrightarrow{v} z,
\]
on a d'une part
\[
x \xrightarrow{vu} z \xleftarrow{\mathrm{id}} z \xrightarrow{\mathrm{id}} z \xleftarrow{\mathrm{id}} z,
\]
d'autre part
\[
x \xrightarrow{u} y \xleftarrow{\mathrm{id}} y \xrightarrow{\mathrm{id}} y \xleftarrow{\mathrm{id}} y \xrightarrow{v} z \xleftarrow{\mathrm{id}} z \xrightarrow{\mathrm{id}} z \xleftarrow{\mathrm{id}} z.
\]
On complète le premier chemin par dégénérescence aux deux bouts, pour le comparer au deuxième en tant qu'élément de $\uCh_4(x,z)$,
\[
x \xrightarrow{\mathrm{id}} x \xleftarrow{\mathrm{id}} x \xrightarrow{vu} z \xleftarrow{\mathrm{id}} z \xrightarrow{\mathrm{id}} z \xleftarrow{\mathrm{id}} z \xrightarrow{\mathrm{id}} z \xleftarrow{\mathrm{id}} z.
\]

\msp{page 24}Il faut prouver que les deux sont homotopes. On a donc deux foncteurs
\[
f,g : \underbrace{\overbrace{\uFl(X;x,y)}^{\text{catégorie discrète}} \times \overbrace{\uFl(X;y,z)}^{\text{catégorie discrète}}}_{\mathcal{A}} \underset{g}{\overset{f}{\rightrightarrows}} \underbrace{\uCh_{4,\Sigma}(X;x,z)}_{\mathcal{B}}, \tag*{$(\ast)$}
\]
on veut prouver qu'ils sont homotopes. Pour ceci, on les regarde (comme tantôt) comme `transposés' des deux foncteurs
\[
\begin{tikzcd}[column sep=large]
\bDelta_1 \sqcup \bDelta_1 \simeq \bDelta_2 & \Phi_4 \arrow[l, shift right, "\varphi"'] \arrow[l, shift left, "\psi"] \\
& \varepsilon, \arrow[ul, "\alpha"] \arrow[u, "\alpha'"']
\end{tikzcd}
\]
vu que les catégories $\mathcal{A}$, $\mathcal{B}$ des deux membres de (*) s'interprètent comme des sous-catégories pleines de $\uHom(\bDelta_2, X)$ et $\uHom(\Phi_4, X)$ respectivement; à savoir les sous-catégories pleines formées des foncteurs dont le composé avec $\alpha$ resp. $\alpha'$ est donné par $x$ et par $z$, et de plus (pour $\mathcal{A}$) dont la valeur en l'objet médian de $\bDelta_2$ soit $y$. On peut encore remplacer $\mathcal{A}$ par la sous-catégorie pleine plus grande $\mathcal{A}'$ de $\uHom(\bDelta_2, X)$, obtenue en ignorant cette dernière condition, et il suffit de prouver que les deux foncteurs `plus gros'
% foot: Version 9 septembre 2001
\printedpage{16}
% head: A. Grothendieck : Dérivateurs, Ch. V, Édition des Universités de Montpellier II et Paris VII
\[
\mathcal{A}' \underset{g'}{\overset{f'}{\rightrightarrows}} \mathcal{B}'
\]
sont homotopes. Et pour ceci, il suffit de
\msp{page 25}voir encore que $\varphi$ et $\psi$ sont homotopes dans la catégories des foncteurs $\Phi_4 \longrightarrow \bDelta_2$ qui sont compatibles avec la source, et le but. On voit donc que notre lemme d'homotopie de tantôt n'était pas assez général. Il faudrait définir les \textit{catégories intervalles généralisés,} de longueur $n$ quelconque ($n =$ nombre de flèches bout à bout), une catégorie intervalle généralisé \textit{standard} de longueur $n$ a pour ensemble d'objets l'intervalle d'entiers % [sic] la catégories -> la catégorie, singular article with plural noun
\[
[0,n],
\]
et pour tout couple d'entiers consécutifs $i-1$, $i$ parmi les objets ($1 \leq i \leq n$), il y a exactement une flèche qui relie $i$, $i+1$ \texttt{[plutôt} $i-1$\texttt{,} $i$\texttt{]}. On affecte l'intervalle $[i,i+1]$ du signe $+$ si $i \longrightarrow i+1$ (flèche en sens direct), du signe $i$ \texttt{[plutôt} $-$\texttt{]} si $i \longleftarrow i+1$ (flèche en sens rétrograde) - donc les catégories standards sont définies, pour $n$ fixé, par un système de signes $+$ ou $-$, dont on affecte les $n$ `tronçons' consécutifs $A_i := [i-1,i]$ ($1 \leq i \leq n$) de l'intervalle d'entiers $[0,n]$. Il y en a donc exactement $2^n$.

On s'attend à ce que l'énoncé suivant soit vrai :

\textbf{Lemme.} \textit{Soient $I$, $I'$ deux intervalles généralisés, munis d'objets source et objet but (de sorte que $I$, $I'$ soient} canoniquement \textit{isomorphes aux intervalles genéralisés standards correspondants). Soit} $\uHom^!(I,I')$ % [sic] genéralisés -> généralisés, missing accent on the first e
\msp{page 26}\textit{la catégorie des foncteurs de $I$ dans $I'$, compatibles avec source et but. Si cette catégorie est $\neq \emptyset$, elle est} 0\textit{-connexe} (et, je présume, même contractile?).

Je vais admettre provisoirement qu'il en est ainsi, du moins dans le cas particulier de $\Phi_4$ (intervalle généralisé de longueur 8) et $\bDelta_2$. Cela donne donc un foncteur
\[
\theta : X \longrightarrow X_\Sigma.
\]
On a gagné, car la propriété universelle de ce foncteur est pratiquement évidente. Mais j'ai oublié de vérifier que $\theta$ transforme éléments de $\Sigma$ en flèches inversibles. Soit donc
\[
x \xrightarrow{u} y
\]
avec $u \in \Sigma$, on va définir un chemin en sens inverse
\[
y \xrightarrow{c} x
\]
qui définira dans $X_\Sigma$ une flèche inverse de $u$. L'idée naturelle est de prendre le `chemin'
\[
y \xleftarrow{u} x
\]
% foot: Version 9 septembre 2001
\printedpage{17}
% head: A. Grothendieck : Dérivateurs, Ch. V, Édition des Universités de Montpellier II et Paris VII
de $y$ dans $x$, mais ce n'est pas un chemin au sens de nos définitions, on va le compléter en
\[
c : y \xrightarrow{\mathrm{id}} y \xleftarrow{u} x \xrightarrow{\mathrm{id}} x \xleftarrow{\mathrm{id}} x.
\]
Je dis que $c$ fait l'affaire. On doit prouver deux relations d'homotopie, de $c \cdot \bar{u}$ et de $\bar{u} \cdot c$ avec les chemins identiques en $x$ \texttt{[resp. en} $y$\texttt{]} de longueur 8 (indice 4, i.e. dans $\uCh_{\Sigma,4}(X;x,x)$ resp. $\uCh_{\Sigma,4}(X;y,y)$),
\msp{page 27}où $\bar{u}$ est le chemin associé à $u$
\[
x \xrightarrow{u} y \xleftarrow{\mathrm{id}} y \xrightarrow{\mathrm{id}} y \xleftarrow{\mathrm{id}} y.
\]

On a donc à regarder
\[
x \xrightarrow{u} y \xleftarrow{\mathrm{id}} y \xrightarrow{\mathrm{id}} y \xleftarrow{\mathrm{id}} y \xrightarrow{\mathrm{id}} y \xleftarrow{u} x \xrightarrow{\mathrm{id}} x \xleftarrow{\mathrm{id}} x
\]
et voir qu'il est homotope au chemin trivial
\[
x \xrightarrow{\mathrm{id}} x \xleftarrow{\mathrm{id}} \cdots x \xrightarrow{\mathrm{id}} x \xleftarrow{\mathrm{id}} x
\]
(des identités partout), et itou pour
\[
y \xrightarrow{\mathrm{id}} y \xleftarrow{u} x \xrightarrow{\mathrm{id}} x \xleftarrow{\mathrm{id}} x \xrightarrow{u} y \xleftarrow{\mathrm{id}} y \xrightarrow{\mathrm{id}} y \xleftarrow{\mathrm{id}} y
\]
et le chemin trivial en $y$. Cette fois, on a encore deux foncteurs
\[
\begin{tikzcd}[column sep=large]
\bDelta_1 & \Phi_4 \arrow[l, shift right, "\varphi"'] \arrow[l, shift left, "\psi"] \\
& \varepsilon, \arrow[ul, "\alpha'"] \arrow[u, "\alpha"']
\end{tikzcd}
\]
compatibles avec objets sources et objets buts, et on doit prouver qu'ils sont homotopes dans $\uHom^!(\Phi_4, \bDelta_1)$. Cela doit donc résulter du même lemme général. On a gagné - par une interprétation simple de $X\Sigma^{-1}$, et en particulier des $\Hom_{X\Sigma^{-1}}(x,y)$ et de leur composition, en termes de la géométrie des espaces de chemins :
\[
\left\{
\begin{aligned}
&\Hom_{X\Sigma^{-1}} \simeq \pi_0(\uCh_\Sigma(X;x,y)) \\ % [sic] Hom_{XΣ^{-1}} -> Hom_{XΣ^{-1}}(x,y), arguments (x,y) missing
&\text{Composition des homomorphismes dans } X\Sigma^{-1} \text{ correspond à} \\
&\text{la composition des } \Sigma\text{-chemins dans } X
\end{aligned}
\right.
\]
% foot: Version 9 septembre 2001

\printedpage{18}
% head: A. Grothendieck : Dérivateurs, Ch. V, Édition des Universités de Montpellier II et Paris VII
\noindent\msp{page 28}

\section*{3\quad Passage aux catégories de diagrammes}

Nous aurions dû noter la catégorie $X$ $\mathcal{M}$, car par la suite elle jouera le rôle d'une `catégorie de modèles'. On désignera par contre par $X$ les `catégories d'indices' sur lesquelles on veut construire les `catégories de diagrammes' de type $X$
\[
\mathcal{M}(X) = \uHom(X,\mathcal{M}).
\]
On se donne donc
\[
\underbrace{\Sigma}_{\text{quasi-isomorphismes}} \subseteq \mathrm{Fl}(\mathcal{M}), \qquad \Sigma \supset \text{ensemble des flèches identiques},
\]
et pour toute catégorie $X$ on définit
\[
\Sigma_X \subset \mathrm{Fl}\,\mathcal{M}(X)
\]
ainsi: si $F, G \in \mathrm{Ob}\,\mathcal{M}(X)$ ($F, G : X \rightrightarrows \mathcal{M}$), une flèche $u : F \longrightarrow G$ (homomorphisme de foncteurs) est dans $\Sigma_X$ si et seulement si $\forall x \in X$, $u(x) : F(x) \longrightarrow G(x)$ est dans $\Sigma$ (`critère de quasi-isomorphie fibre par fibre', pour l'homomophisme entre les copréfaisceaux $F$, $G$ envisagés). % [sic] homomophisme -> homomorphisme, letter r missing

Soient
\[
a, b \in \mathrm{Ob}\,\mathcal{M}
\]
et considérons les foncteurs constants correspondants sur $X$
\[
a_X, b_X \in \mathrm{Ob}\,\mathcal{M}(X).
\]

\msp{page 29}On se propose d'appliquer la théorie du \S2 au `calcul' de
\[
\Hom_{\Sigma_X}(a_X, b_X),
\]
on a donc une bijection canonique
\[
\boxed{\Hom_{\Sigma_X}(a_X, b_X) \simeq \pi_0\uCh_{\Sigma_X}(\mathcal{M}(X); a_x, b_x).} % [sic] a_x, b_x -> a_X, b_X, lowercase index against a_X, b_X elsewhere
\]
On se propose d'expliciter la catégorie de chemins du second membre. On a, abstraction faite des données de $\Sigma$, $a$, $b$ dans $\mathcal{M}$, l'expression pour la catégorie de \textit{tous} les chemins
% foot: Version 9 septembre 2001


\endgroup
% END UNIT Der-ChV-234/Der-ChV-234_p016-020.tex
% BEGIN UNIT Der-ChV-234/Der-ChV-234_p021-025.tex
\begingroup\makeatletter%
\NoAutoSpacing%
\newcommand{\printedpage}[1]{\marginpar{\footnotesize p.~#1}}%
\newcommand{\msp}[1]{\marginpar{\footnotesize [#1]}}%
\newcommand{\cM}{\mathcal{M}}%
\newcommand{\Ch}{\underline{\mathrm{Ch}}}%
\newcommand{\Chh}{\underline{\underline{\mathrm{Ch}}}}%
\newcommand{\Hom}{\underline{\mathrm{Hom}}}%
\newcommand{\KK}{\underline{\underline{K}}}%
\newcommand{\dfeq}{\overset{\text{déf}}{=}}%
\newcommand{\Cathot}{\mathrm{Cat\,hot}}%
\newcommand{\uHom}{\underline{\mathrm{Hom}}}%
\makeatother\let\printedpage\derpp\let\msp\dermsp%
\let\mspin\DERmspin%


\printedpage{19}
% head: A. Grothendieck : Dérivateurs, Ch. V, Édition des Universités de Montpellier II et Paris VII

dans $\cM(X)$
\[
\begin{aligned}
\Ch(\cM(X)) \quad &\simeq \quad \uHom(\ \underbrace{\boldsymbol{\Phi}}_{\substack{\text{pro-catégorie}\\ (\Phi_n)_{n\in\mathbf{N}}}}\ ,\ \underbrace{\cM(X)}_{\uHom(X,M)}\ ) \\ % reading resolved: italic M is printed; category-valued Hom underlines restored.
&\dfeq \quad \varinjlim_n \uHom(\Phi_n,\cM(X)) \\
&\simeq \quad \varinjlim_n \uHom(\Phi_n\times X,\cM) \\
&\simeq \quad \varinjlim_n \uHom(X,\uHom(\Phi_n,\cM)) \\
&\dfeq \quad \varinjlim_n \uHom(X,\Ch_n(\cM)) \\
&\simeq \quad \uHom^{!}(X,\Ch(\cM)),
\end{aligned}
\]
où le signe ! signifie qu'on prend la sous-catégorie pleine des foncteurs
\[
X\longrightarrow \Ch(\cM),
\]
\msp{page 30}
qui se factorisent par une des sous-catégories pleines $\Ch_n(\cM)$. On peut dire aussi que si on considère le système inductif des catégories $\Ch_n(\cM)$ comme un ind-objet de Cat, qu'on va noter
\[
\Chh(\cM) \qquad\qquad\qquad\qquad \text{ind-objet de Cat,}
\]
on a
\[
\Ch(\cM(X)) \xrightarrow{\ \sim\ } \uHom(X,\Chh(\cM)).
\]
Si on en vient à $\Sigma$, $\Sigma_X$, on trouve aussitôt qu'un chemin dans $\cM(X)$ est un $\Sigma_X$-chemin, si et seulement si le foncteur correspondant $X\longrightarrow\Ch(\cM)$ se factorise par $\Ch_{\Sigma}(\cM)$. Ainsi on trouve
\[
\Ch_{\Sigma_X}(\cM(X))\simeq\uHom^{!}(X,\Ch_{\Sigma}(\cM)),
\]
où le signe ! a le même sens que plus haut. Si on introduit le ind-objet $\Chh_{\Sigma}(\cM)$ défini par les $\Ch_{\Sigma,n}(\cM)$, on trouve donc aussi
\[
\Ch_{\Sigma_X}(\cM(X))\xrightarrow{\ \sim\ }\uHom(X,\Chh_{\Sigma}(\cM)).
\]

\msp{page 31}
Soit $c$ un chemin dans $\cM(X)$, et
\[
X\longrightarrow\Ch(\cM)
\]
le foncteur qu'il définit, composons-le avec les foncteurs source-but
\[
\Ch(\cM)\longrightarrow\cM\times\cM,
\]
on trouve un foncteur
\[
X\longrightarrow\cM\times\cM,
\]
i.e. un objet $(F_0,F_1)$ de $\cM(X)\times\cM(X)$. Je dis que $F_0$ et $F_1$ sont respectivement la source et le but du chemin $c$. De façon un peu plus précise, on a commutativité (stricte) dans le diagramme de foncteurs
% foot: Version 9 septembre 2001
\printedpage{20}
% head: A. Grothendieck : Dérivateurs, Ch. V, Édition des Universités de Montpellier II et Paris VII
\[
\begin{tikzcd}[column sep=large, row sep=large]
\Ch_{\Sigma_X}(\cM(X)) \arrow[r, "\sim"] \arrow[d, "{\sigma_{\cM(X)},\,\beta_{\cM(X)}}"] & \uHom^{!}(X,\Ch_{\Sigma}(\cM)) \arrow[d, "{\text{déduit de }\sigma_{\cM},\,\beta_{\cM}}"] \\
\cM(X)\times\cM(X) \arrow[r, "\sim"] & \uHom(X,\cM\times\cM),
\end{tikzcd}
\]
où les foncteurs $\sigma$ sont les foncteurs source, les foncteurs $\beta$ les foncteurs but.

Si on se donne arbitrairement les objets $F_0$, $F_1$ dans $\cM(X)$, donc un foncteur
\[
X\longrightarrow\cM\times\cM,
\]
$\Ch_{\Sigma_X}(F_0,F_1)$ s'identifie donc à la catégorie des foncteurs ‘admissibles’ (i.e. se factorisant par un $\Ch_{\Sigma,n}(\cM)$) $X\xrightarrow{\ C\ }\Ch_{\Sigma}(\cM)$ qui \textit{relèvent}
\msp{page 32 vide}
\msp{page 33}
le foncteur $(\sigma_{\cM},\beta_{\cM})$ \texttt{[plutôt} $(F_0,F_1)$\texttt{]} précédent, i.e. qui donnent lieu à un triangle commutatif (strict)
\[
\begin{tikzcd}[row sep=large]
 & \Ch_{\Sigma}(\cM) \arrow[d, "{(\sigma_{\cM},\,\beta_{\cM})}"] \\
X \arrow[ur, "C"] \arrow[r, "{(F_0,F_1)}"] & \cM\times\cM.
\end{tikzcd}
\]
Mais si les foncteurs $F_0$, $F_1$ sont constants, de valeur $a$, $b$ comme tantôt, i.e. $(F_0,F_1)$ constant de valeur $(a,b)$, alors la catégorie des foncteurs $C$ en question est tautologiquement isomorphe à celle des foncteurs de $X$ dans la fibre de $(\sigma_{\cM},\beta_{\cM})$ en $(a,b)$, d'où l'isomorphisme de catégories fondamental
\[
\boxed{\Ch_{\Sigma_X}(X;a_X,b_X)\simeq\uHom^{!}(X,\Ch_{\Sigma}(\cM;a,b)).}
\]

On va expliciter le contenu essentiel de cette formule en une

\textbf{Scholie.} Soient $\cM$ une catégorie avec un ensemble $\Sigma$ de flèches de localisation, $a,b\in\mathrm{Ob}\,\cM$. Alors on peut trouver un objet $\KK=K(a,b)$, $K=(K_n)$, de Ind\,Cat (indexé par $\mathbf{N}$ mais peu importe ce fait), et un isomorphisme canonique pour toute catégorie $X$ (isomorphisme fonctoriel en $X$)
\[
\begin{aligned}
\mathrm{Hom}_{\Sigma_X}(a_X,b_X) \quad &\simeq \quad \pi_0(\uHom(X,\KK)) \\
&\dfeq \quad \pi_0(\ \varinjlim \uHom(X,K_n)) \\
&\simeq \quad \varinjlim \pi_0(\uHom(X,K_n)).
\end{aligned}
\]
% foot: Version 9 septembre 2001

\printedpage{21}
% head: A. Grothendieck : Dérivateurs, Ch. V, Édition des Universités de Montpellier II et Paris VII

\msp{page 34}
Je vais épiloguer sur ce résultat, sans me soucier pour l'instant de la façon particulière dont la ind-catégorie $K$ \texttt{[plutôt} $\KK$\texttt{]} a été obtenue en termes de $\cM$, $\Sigma$, $a$, $b$.

Comme les $\pi_0\uHom(X,K_n)=\mathrm{Hom}_{\Cathot}(X,K_n)$ ne dépendent que de $X$ et de $K_n$ en tant qu'objets de Cat\,hot, je vais interpréter $\KK$ comme un ind-objet de Cat\,hot, et écrire la formule fondamentale sous la forme
\[
\mathrm{Hom}_{\Sigma_X}(a_X,b_X)\simeq\mathrm{Hom}_{\mathrm{Ind\,Cat\,hot}}(X,\KK).
\]

Ainsi, l'énoncé du scholie implique deux choses :

\begin{enumerate}
\item[a)] Le contrafoncteur $X\longmapsto\mathrm{Hom}_{\Sigma_X}(a_X,b_X)$ sur Cat se factorise canoniquement par Cat\,hot, via un foncteur canonique
\[
\Cathot^{o}\ \xrightarrow{\ H=H^{\cM,\Sigma}_{a,b}\ }\mathrm{Ens}\ ;
\]
\item[b)] Ce dernier contrafoncteur sur Cat\,hot est \textit{ind-représentable} par un ind-objet $\KK$ de Cat\,hot.
\end{enumerate}

C'est cet ind-objet de Cat\,hot qui est défini à isomorphisme unique près par $\cM$, $\Sigma$, $a$, $b$, indépendamment de tout artifice de calcul pour expliciter les $\mathrm{Hom}_{\Sigma_X}(a_X,b_X)$.

Notons aussi la conclusion suivante, nullement évidente à priori :

\msp{page 35}
Notons d'abord (ceci étant implicite dans le scholie) que toute flèche dans Cat
\[
X\xrightarrow{\ f\ }Y
\]
définit
\[
\cM(Y)\xrightarrow{\cM(f)}\cM(X) \qquad\qquad\qquad \text{appliquant } \Sigma_Y \text{ dans } \Sigma_X,
\]
d'où
\[
\cM(Y)\Sigma_Y^{-1}\dfeq\mathbf{D}_{\cM,\Sigma}(Y)\longrightarrow\cM_X\Sigma_X^{-1}=\mathbf{D}_{\cM,\Sigma}(X),
\]
\texttt{[plutôt} $\cM(X)$\texttt{]} d'où, pour tout $a,b\in\mathrm{Ob}\,\cM$, vu que $\cM(f)$ transforme $a_Y$, $b_Y$ en $a_X$, $b_X$ :
\[
\boxed{\mathrm{Hom}_{\Sigma_Y}(a_Y,b_Y)\longrightarrow\mathrm{Hom}_{\Sigma_X}(a_X,b_X).}
\]

Ceci posé, on a :

\textbf{Corollaire :} \textit{Si $f$ est un homotopisme, alors quelque soient $a,b\in\cM$, la flèche précédente est un isomorphisme.} % [sic] quelque soient -> quels que soient, grammatical form of the locution

La question naturelle, c'est de savoir si la même conclusion est encore valable, en supposant seulement que $f$ soit un ‘\textit{quasi-isomorphisme}’ (ou ‘équivalence faible’), ici que $f$ définisse un isomorphisme dans la catégorie de fractions Hot de Cat. (À ne pas confondre avec Cat\,hot,
% foot: Version 9 septembre 2001
\printedpage{22}
% head: A. Grothendieck : Dérivateurs, Ch. V, Édition des Universités de Montpellier II et Paris VII
\msp{page 36}
nettement plus fine). S'il en était ainsi, le contrafoncteur
\[
X\longmapsto\mathrm{Hom}_{\Sigma_X}(a_X,b_X)
\]
se factoriserait non seulement par Cat\,hot, mais même par la catégorie de fractions plus grossière Hot. Et dès lors se poserait la question si ce contrafoncteur hypothétique
\[
\mathrm{Hot}^{o}\longrightarrow\mathrm{Ens}
\]
est encore ind-représentable; et de préférence (\footnote[5]{ce sera automatique si on a ind-représentabilité comme foncteur sur Hot, cf. ci-dessous.}) ind-représentable par le ind-objet image de $\KK$ dans Ind\,Hot par
\[
\mathrm{Ind}(\Cathot)\longrightarrow\mathrm{Ind}(\mathrm{Hot}),
\]
déduit du foncteur de localisation $\Cathot\longrightarrow\mathrm{Hot}$. (Car cet ind-objet est un candidat tout trouvé.)

\msp{page 37}
Je vais faire quelques rappels catégoriques, généralisant la situation du foncteur canonique de localisation
\[
\Cathot\xrightarrow{\ \varphi\ }\mathrm{Hot},
\]
dont on sait qu'il admet un adjoint à droite \texttt{[si un lecteur a une démonstration de ce fait, prière de contacter G. Maltsiniotis : [e-mail address removed]]}
\[
\Cathot\xleftarrow{\ \psi\ }\mathrm{Hot}, \qquad\qquad\qquad \mathrm{Hom}_{\mathrm{Hot}}(\varphi X,Y)\simeq\mathrm{Hom}_{\Cathot}(X,\psi Y),
\]
où $\psi$ est le foncteur ‘enveloppe de Kan’ ou ‘représentant de Kan’. On va donc considérer la situation plus générale
\[
\begin{tikzcd}[column sep=large]
C \arrow[r, shift right=1, "\varphi"'] & C' \arrow[l, shift right=1, "\psi"']
\end{tikzcd}
\qquad\qquad \mathrm{Hom}_{C'}(\varphi X,X')\simeq\mathrm{Hom}_{C}(X,\psi X')
\]
i.e.
\[
\varphi,\psi \text{ couple de foncteurs adjoints.}
\]
On suppose de plus
\[
\left\{
\begin{aligned}
&\varphi \text{ un foncteur localisation, i.e. } \psi \text{ pleinement fidèle,}\\
&\text{i.e. } \varphi\psi\xrightarrow[\sim]{\text{adj}}\mathrm{id}_{C'} \quad .
\end{aligned}
\right.
\]
Soit $W=W_{\varphi}\subset\mathrm{Fl}\,C$, défini par $\varphi$ : $W_{\varphi}=\{u\in\mathrm{Fl}\,C\mid\varphi(u)\text{ isomorphisme}\}$. Voici la proposition que j'avais oubliée, et que j'ai eu du mal à reconstituer, hélas!

\textbf{Proposition 1.} \textit{Soit $K\in \mathit{Ob}C$. Les conditions suivantes sur $K$ sont équivalentes :}
% foot: Version 9 septembre 2001

\printedpage{23}
% head: A. Grothendieck : Dérivateurs, Ch. V, Édition des Universités de Montpellier II et Paris VII

\begin{enumerate}
\item[(i)] \textit{Le foncteur}
\[
X\longmapsto K(X)\dfeq\mathrm{Hom}_C(X,K) \qquad\quad C^{o}\longrightarrow\mathrm{Ens}
\]
\textit{transforme flèches de $W$ (‘quasi-isomorphismes’) en bijections.}
\item[(i$'$)] \textit{Même condition pour les flèches d'adjonction}
\[
X\longrightarrow\psi\varphi X \qquad\quad (\forall X\in\mathrm{Ob}\ \mathrm{C})
\]
\item[(i$''$)] \textit{Même condition pour la flèche d'adjonction} (\footnote[6]{N.B. Il suffit même que l'élément $\mathrm{id}_K$ de $F_K(K)$ provienne de $F_K(\psi\varphi K)$ \texttt{[où} $F_K=K(-)=\mathrm{Hom}(-,K)$\texttt{]}, i.e.\par (i$'''$) Il existe une rétraction de $\psi\varphi X$ sur $X$ \texttt{[plutôt} $\psi\varphi K$ \texttt{sur} $K$\texttt{]} (NB Celle-ci sera même un isomorphisme).})
\[
K\longrightarrow\psi\varphi K.
\]
\item[(ii)] \textit{$K$ appartient à l'image essentielle de $\psi$}
\item[(ii$'$)] \textit{$K\longrightarrow\psi\varphi K$ (adjonction) est un isomorphisme.}
\item[(ii$''$)] \textit{L'homomorphisme fonctoriel en $X$}
\[
\mathrm{Hom}_C(X,K)\longrightarrow\mathrm{Hom}_{C'}(\varphi X,\varphi K)
\]
\textit{est un isomorphisme.}
\end{enumerate}

\textsc{Démonstration.} On a d'abord comme quasi-tautologies (ii$'$)$\Longrightarrow$ (ii)$\Longrightarrow$ (ii$''$)$\Longrightarrow$ (ii$'$), i.e. l'équivalence des trois dernières conditions. En effet, (ii$'$)$\Longrightarrow$ (ii) tautologique, (ii)$\Longrightarrow$ (ii$''$) car si $K=\psi(K')$ on a % reading resolved: single and double mathematical primes as printed.
\msp{page 38}
par adjonction
\[
\mathrm{Hom}_C(X,K)=\mathrm{Hom}_C(X,\psi(K'))\overset{\text{adj}}{\simeq}\mathrm{Hom}_{C'}(\varphi X,K'),
\]
or $\varphi(K)=\varphi\psi(K')\xrightarrow{\ \sim\ }K'$ (adjonction, et $\psi$ pleinement fidèle).
\[
\begin{tikzcd}[column sep=small, row sep=large]
\mathrm{Hom}_C(X,K) \arrow[rr, "\sim"] \arrow[dr, "\varphi_{X,K}"] & & \mathrm{Hom}_C(\varphi X,K') & \text{\footnotesize qui se complète en} \\ % [sic] Hom_C(\varphi X,K') -> Hom_{C'}(\varphi X,K'), \varphi X and K' are objects of C' (cf. the display above)
& \mathrm{Hom}_{C'}(\varphi X,\ \underbrace{\varphi K}_{=\varphi\psi(K')}\ ) \arrow[ur, "\sim", "\text{adj}"'] & &
\end{tikzcd}
\]
ce qui prouve que $\varphi_{X,K}$ est un isomorphisme.

(ii$''$) $\Longrightarrow$ (ii$'$). Car la flèche
% foot: Version 9 septembre 2001


\endgroup
% END UNIT Der-ChV-234/Der-ChV-234_p021-025.tex
% BEGIN UNIT Der-ChV-234/Der-ChV-234_p026-030.tex
\begingroup\makeatletter%
\newcommand{\printedpage}[1]{\marginpar{\footnotesize p.~#1}}%
\newcommand{\msp}[1]{\marginpar{\footnotesize [#1]}}%
\newcommand{\Hom}{\operatorname{Hom}}%
\newcommand{\Ind}{\operatorname{Ind}}%
\newcommand{\Ens}{\mathrm{Ens}}%
\newcommand{\Cat}{\mathrm{Cat}}%
\newcommand{\Cathot}{\mathrm{Cat\,hot}}%
\newcommand{\Hot}{\mathrm{Hot}}%
\newcommand{\Ob}{\mathrm{Ob}}%
\newcommand{\id}{\mathrm{id}}%
\newcommand{\Ch}{\mathrm{Ch}}%
\newcommand{\Chh}{\underline{\underline{\mathrm{Ch}}}}%
\makeatother\let\printedpage\derpp\let\msp\dermsp%
\let\mspin\DERmspin%


\printedpage{24}
% head: A. Grothendieck : Dérivateurs, Ch. V, Édition des Universités de Montpellier II et Paris VII
% foot: Version 9 septembre 2001

\[
\begin{tikzcd}
\Hom_C(X,K) \arrow[r,"\alpha_X"] \arrow[dr,dashed] & \Hom_{C'}(\varphi X,\varphi K) \arrow[d,dash,"\wr"',"\mathrm{adj}"] \\
 & \Hom_C(X,\psi\varphi K)
\end{tikzcd}
\]
est déduite de la flèche d'adjonction
\[
K \longrightarrow \psi\varphi K
\]
en prenant $\Hom_C(X,-)$, donc $\alpha_X$ est un isomorphisme pour tout $X$ si et seulement si cette flèche d'adjonction est un isomorphisme. Donc
\[
\text{(ii)} \Longleftrightarrow \text{(ii$'$)} \Longleftrightarrow \text{(ii$''$)}.
\]
D'autre part,
\[
\text{(ii$''$)} \Longrightarrow \text{(i)} \Longrightarrow \text{(i$'$)} \Longrightarrow \text{(i$''$)}
\]
sont tautologiques, en tenant compte pour (i) $\Longrightarrow$ (i) \texttt{[probablement (i) $\Longrightarrow$ (i$'$)]} que les flèches d'adjonction
\[
X \xrightarrow{u} \psi\varphi X \qquad\qquad\qquad \text{dans } C
\]
sont bien dans $W$, car on a un triangle commutatif
\[
\begin{tikzcd}[column sep=large]
\varphi(X) \arrow[r,"\varphi(u)"] \arrow[dr,"\id_{\varphi(X)}"] & \varphi(\psi\varphi X) = (\varphi\psi)(\varphi X) \arrow[d,"{\begin{array}{l}\wr\ \mathrm{adj}\ \varphi\psi \longrightarrow 1,\ \text{qui est un}\\ \text{isomorphisme car } \psi \text{ est pleinement fidèle}\end{array}}"] \\
 & \varphi(X),
\end{tikzcd}
\]
donc $\varphi(u)$ est bien un isomorphisme.

Reste à prouver (i$''$) $\Longrightarrow$ (ii$'$), i.e. on suppose

\noindent\msp{page 39}%
que le foncteur $X \longrightarrow \Hom(X,K)$ transforme le morphisme d'adjonction particulier
\[
K \xrightarrow{i} \psi\varphi K
\]
en
\[
\begin{aligned}
u \quad &\longrightarrow \quad u\circ i\\
\Hom(\psi\varphi(K),K) \quad &\longrightarrow \quad \Hom(K,K),
\end{aligned}
\]
qui est une bijection. Donc $\id_K$ est dans l'image, soit $p$ un élément du premier membre qui lui donne naissance

\printedpage{25}
% head: A. Grothendieck : Dérivateurs, Ch. V, Édition des Universités de Montpellier II et Paris VII
% foot: Version 9 septembre 2001

\[
\begin{tikzcd}
K \arrow[r,hook,shift left,"i"] & \psi\,\underbrace{\varphi(K)}_{K'} \arrow[l,shift left,"p"]
\end{tikzcd}
\qquad\qquad {\scriptstyle \text{donc } p\circ i = \id_K}
\]
Considérons
\[
ip = q : \psi(K') \longrightarrow \psi(K') \qquad\qquad\qquad (\text{NB } q^2 = q),
\]
comme $\psi$ est pleinement fidèle, il existe un unique
\[
q_0 : K' \longrightarrow K' \qquad\qquad\qquad (K' \overset{\text{déf}}{=} \varphi(K))
\]
tel que l'on ait
\[
q = \psi(q_0).
\]
La relation $pi = \id_K$ donne, en appliquant $\varphi$,
\[
\varphi(p)\varphi(i) = \id,
\]
donc comme $\varphi(i)$ est un isomorphisme (c'est
\[
\begin{tikzcd}[column sep=large]
\varphi(K) \arrow[r,"\sim"] \arrow[dr,"\sim"{sloped,pos=0.45},"\id"{pos=0.6}] & \varphi(\psi\varphi)(K) = (\varphi\psi)(\varphi K) \arrow[d,"\wr"',"\mathrm{adj}"] \\
 & \varphi(K),
\end{tikzcd}
\]
rendant commutatif le triangle ci-contre), $\varphi(p)$ l'est. Donc $\varphi(q) = \varphi(i)\varphi(p)$ l'est, puisque $\varphi(i)$ et $\varphi(i)$ \texttt{[plutôt $\varphi(p)$]} le sont. Mais
\[
\varphi(q) = \varphi\psi(q_0) \quad \underbrace{\xrightarrow{\ \sim\ }}_{\mathrm{adj}\ \varphi\psi\ \xrightarrow{\sim}\ \id} \quad q_0,
\]
donc $q_0$ est un isomorphisme, donc aussi $q = \psi q_0$. Comme $q^2 = q$, cela implique $q = \id$. Donc $i$ et $p$ sont des isomorphismes inverses l'un de l'autre, ce qui établit (ii$''$) \texttt{[plutôt (ii$'$)]} et achève la démonstration.

\textbf{Corollaire :} \textit{Considérons des foncteurs $C^o \xrightarrow{H} \Ens$, et $C'^o \xrightarrow{H'} \Ens$, tels que $H \simeq H'\circ\varphi$. Pour que $H$ soit représentable, il faut et il suffit que $H'$ le soit, et dans ce cas, l'objet $K$ de $C$}

\noindent\msp{page 40}%
\textit{qui représente $H$ satisfait aux conditions équivalentes de la proposition. Si $K$, $K'$ les représentent, on a des isomorphismes canoniques $K' \simeq \varphi K$, $K \simeq \psi K'$.}

Si $H'$ est représentable par $K' \in \Ob C'$, on aura donc
\[
H'\circ\varphi(X) \simeq \Hom_{C'}(\varphi X,K') \overset{\mathrm{adj}}{\simeq} \Hom_C(X,\psi(K'))
\]
\printedpage{26}
% head: A. Grothendieck : Dérivateurs, Ch. V, Édition des Universités de Montpellier II et Paris VII
% foot: Version 9 septembre 2001
donc $H$ est représentable par $K = \psi(K')$, satisfaisant la condition (ii) de la proposition. Inversement, si $H$ est représentable par un objet $K$, alors $K$ satisfait à la condition (i), donc par (ii$''$) on a
\[
\Hom_C(X,K) \xrightarrow{\sim} \Hom_{C'}(\varphi X,\varphi K),
\]
donc le foncteur $H' : C^o \longrightarrow \Ens$ qui factorise $H$ s'explicite comme $X' \longmapsto \Hom_{C'}(X',\varphi K)$, donc est représentable par $\varphi K$. % [sic] C^o -> C'^o, H' is a functor on C'^o (cf. the Corollaire)

Passons aux ind-objets. Les foncteurs $\varphi$, $\psi$ passent aux ind-objets, d'où des foncteurs
\[
\begin{tikzcd}[column sep=large]
\Ind C \arrow[r,shift left,"\Psi = \Ind\psi"] \arrow[r,shift right,"\Phi = \Ind\varphi"'] & \Ind C'.
\end{tikzcd} % [sic] both arrows printed pointing to Ind C' -> the arrow Psi = Ind psi apparently intended from Ind C' to Ind C, psi goes from C' to C
\]
On aura encore la formule d'adjonction pour $X = (X_i)_{i\in I}$ dans $\Ind C$, $X' = (X'_{i'})_{i'\in I'}$ dans $\Ind C'$, car
\[
\begin{aligned}
\Hom(X,\Psi X') &= \varprojlim_i \varinjlim_{i'} \Hom(X_i,\psi(X'_{i'}))\\
&= \varprojlim_i \varinjlim_{i'} \Hom(\varphi X_i,X'_{i'})\\
&= \Hom(\Phi X,X').
\end{aligned}
\]

\noindent\msp{page 41}%
On prouve de même que $\Psi$ est pleinement fidèle, car si $X' = (X'_i)_{i\in I}$, $Y' = (Y'_j)_{j\in J}$ sont des objets de $\Ind C'$, on aura
\[
\begin{tikzcd}[column sep=small]
\Hom(X',Y') \arrow[d,"\Psi_{X',Y'}"] & \simeq \varprojlim\limits_i \varinjlim\limits_j & \Hom_{C'}(X'_i,Y'_j) \arrow[d,"\psi_{X'_i,Y'_j}"] \\
\Hom(\Psi X',\Psi Y') & \simeq \varprojlim\limits_i \varinjlim\limits_j & \Hom_C(\psi(X'_i),\psi(Y'_j)),
\end{tikzcd}
\]
donc $\Psi_{X',Y'}$ est un isomorphisme si les $\psi_{X'_i,Y'_j}$ le sont.

Donc la proposition 1 peut s'appliquer à la situation $\Phi$, $\Psi$ entre ind-objets. Je vais l'expliciter sous la forme la plus commode :

\textbf{Proposition 2.} \textit{Soit $K = (K_i)_{i\in I}$ un ind-objet de $C$, et soit $H$ le foncteur}
\[
\left\{
\begin{aligned}
H : X \quad &\longmapsto \quad \underbrace{\Hom_C(X,K)}_{K(X)} \overset{\text{déf}}{=} \varinjlim_i \Hom_C(X,K_i)\\
C^o \quad &\longrightarrow \quad \Ens
\end{aligned}
\right.
\]
\textit{qu'il ind-représente. Conditions équivalentes}

\begin{enumerate}[label={}]
\item[(i)] \textit{Le foncteur $H$ se factorise (à isomorphisme près) par $C'$ via $\varphi$, i.e. transforme les $u \in W_\varphi$ en isomorphismes.}
\printedpage{27}
% head: A. Grothendieck : Dérivateurs, Ch. V, Édition des Universités de Montpellier II et Paris VII
% foot: Version 9 septembre 2001
\item[(i$'$)] \textit{Même condition pour les flèches d'adjonction $X \longrightarrow \psi\varphi(X)$ (lesquelles sont dans $W$).}
\item[(i$''$)] \textit{Même condition pour la flèche d'adjonction $K \longrightarrow \psi\varphi(K)$.}
\item[(ii)] \textit{$K$ appartient à l'image essentielle de} $\Ind C'$ \textit{par $\Psi$.}
\item[(ii$'$)] $K \longrightarrow \Psi\Phi(K)$ \textit{est un} isomorphisme \textit{dans} $\Ind C$.) % [sic] ".)" -> ".", closing parenthesis without opening parenthesis
\item[(ii$''$)] \textit{L'homomorphisme fonctoriel}
\[
K(X) = \Hom_C(X,K) \longrightarrow \Hom_{C'}(\varphi(X),\varphi(K))
\]
\texttt{[plutôt} $\Hom_{\mathrm{Ind}(C)}(X,K) \longrightarrow \Hom_{\mathrm{Ind}(C')}(\Phi(X),\Phi(K))$\texttt{, en identifiant l'objet} $X$ \texttt{de} $C$ \texttt{à son image dans} $\mathrm{Ind}(C)$\texttt{]} \textit{est un isomorphisme.}
\end{enumerate}

\textsc{Démonstration.} L'équivalence de (ii), (ii$'$), (ii$''$) est la proposition 1 appliquée à $\Phi$ et $\Psi$, à cela près qu'il faut noter que (ii$''$) reste évidemment valable si on y remplace $X$ par un ind-objet $(X_j)_{j\in J}$ quelconque. On a

\noindent\msp{page 42}%
\[
\text{(ii$''$)} \Longrightarrow \text{(i)} \Longrightarrow \text{(i$'$)}
\]
tautologiquement, il reste à prouver que (i$'$) implique (i$''$) et (i$''$) (ii$'$). On ne peut pas appliquer directement la proposition 1, car l'hypothèse (i$'$) ici ne porte que sur les objets de $\Ind C$ qui proviennent de $C$, et il n'est pas évident à priori qu'elle passe aux objets quelconques. Mais cela l'est dans le cas d'un homomorphisme de ind-objets $(X_i)$, $(Y_i)$ sur le même ensemble d'indices, donné par un homomorphisme ordinaire (terme à terme) de systèmes inductifs sur $I$, $X_i \longrightarrow Y_i$: si les $X_i \longrightarrow Y_i$ sont dans $W$, l'hypothèse (i) implique que $H(Y) \longrightarrow H(X)$ est un isomorphisme, et de même pour (i$'$) pour $\underbrace{X}_{=(X_i)_i} \;\longrightarrow\; \underbrace{\Psi\Phi(X)}_{=(\psi\varphi(X_i))_i}$\,, où l'homomorphisme d'adjonction est donné terme à terme par les homomorphismes d'adjonction $X_i \longrightarrow \psi\varphi(X_i)$. Donc l'hypothèse (i$'$) ici implique l'hypothèse (i$'$) de la proposition 1 pour $\Phi$ et $\Psi$, donc aussi (ii$'$) (par la proposition 1), et on gagne. % [sic] "et (i'') (ii')" -> "et (i'') => (ii')", the implication sign or verb is missing

Revenons à la situation où $C = \Cathot$, $C' = \Hot$, $\varphi$ le foncteur \texttt{[de]} localisation. Si $K \in \Ob\Cathot$, alors $\psi\varphi(K)$ s'appelle \textit{l'enveloppe de Thomason-Kan} de $K$ - si $K$ est une catégorie, son enveloppe $\bar{K}$ et la flèche $K \longrightarrow \bar{K}$ ne sont définis que dans la catégorie Cat\,hot, i.e. $\bar{K}$ n'est défini qu'à `homotopisme près'. % reading resolved: no underline is printed under psi phi(K).

\noindent\msp{page 43}%
Mais on peut toujours \textit{choisir} un objet $\bar{K}$ de Cat, et une flèche
\[
i : K \longrightarrow \bar{K} \qquad\qquad\qquad \textit{dans } \Cat,
\]
\printedpage{28}
% head: A. Grothendieck : Dérivateurs, Ch. V, Édition des Universités de Montpellier II et Paris VII
% foot: Version 9 septembre 2001
pas seulement dans Cat\,hot, de telle façon de plus que $i$ soit une `\textit{cofibration} de Thomason' (et en particulier une immersion ouverte?). Il en résulte que si on a un système inductif
\[
K = (K_n)_{n\in\mathbf{N}}
\]
indexé par $\mathbf{N}$, et si on choisit pour chaque $n$ une enveloppe de Thomason-Kan
\[
K_n \overset{\alpha_n}{\hookrightarrow} \bar{K}_n
\]
de cette façon, on peut trouver des foncteurs de transition
\[
\bar{K}_n \xrightarrow{\bar{i}_n} \bar{K}_{n+1}
\]
tels que les carrés
\[
\begin{tikzcd}
K_n \arrow[r,hook,"\alpha_n"] \arrow[d,"i_n"'] & \bar{K}_n \arrow[d,"\bar{i}_n"] \\
K_{n+1} \arrow[r,hook,"\alpha_{n+1}"] & \bar{K}_{n+1}
\end{tikzcd}
\]
soient commutatifs, de telle sorte que l'on ait donc un homomorphisme de \textit{systèmes inductifs de} Cat
\[
K \longrightarrow \bar{K}
\]
et a fortiori un homomorphisme de ind-objets de Cat. Ceci posé, les conditions (i), (ii$'$) de la proposition 2 montrent que le foncteur ind-représenté par $K$ \textit{dans} $C = \Cathot$ transforme quasi-isomorphismes

\noindent\msp{page 44}%
ordinaires (équivalences faibles) en bijections, \textit{si et seulement si} l'homomorphisme précédent est un isomorphisme de ind-objets de Ind\,Cat\,hot.

Ce sera donc cela la condition naturelle à imposer aux ind-catégories de chemins
\[
\Chh_{\Sigma}(\mathcal{M};a,b) \qquad\qquad\qquad (a,b \in \mathrm{ob}\,\mathcal{M})
\]
provenant d'une `catégorie de modèles' $\mathcal{M}$, $\Sigma$.

Bien sûr, disposant d'un système inductif explicite de catégories (pas seulement d'objets de Cat\,hot)
\[
K_n = \underline{\Ch}_{\Sigma,n}(\mathcal{M};a,b),
\]
il est tentant de passer à la limite inductive
\[
K_\infty = \varinjlim_n \underline{\Ch}_{\Sigma,n}(\mathcal{M};a,b) \overset{\text{déf}}{=} \underline{\Ch}_{\Sigma}(\mathcal{M};a,b),
\]
et ceci d'autant plus que les foncteurs de transition sont des monomorphismes et pleinement fidèles, de sorte que les $K_n$ s'identifient à une suite de sous-catégories pleines de


\endgroup
% END UNIT Der-ChV-234/Der-ChV-234_p026-030.tex
% BEGIN UNIT Der-ChV-234/Der-ChV-234_p031-034.tex
\begingroup\makeatletter%
\newcommand{\printedpage}[1]{\marginpar{\footnotesize p.~#1}}%
\newcommand{\msp}[1]{\marginpar{\footnotesize [#1]}}%
\newcommand{\Cathot}{\mathrm{Cat\,hot}}%
\newcommand{\uHom}{\underline{\mathrm{Hom}}}%
\newcommand{\uuK}{\underline{\underline{K}}}%
\newcommand{\uCh}{\underline{\mathrm{Ch}}}%
\newcommand{\Hom}{\mathrm{Hom}}%
\newcommand{\Hot}{\mathrm{Hot}}%
\newcommand{\Ens}{\mathrm{Ens}}%
\newcommand{\Ch}{\mathrm{Ch}}%
\newcommand{\defeq}{\overset{\text{déf}}{=}}%
\newcommand{\Chh}{\underline{\underline{\mathrm{Ch}}}}%
\makeatother\let\printedpage\derpp\let\msp\dermsp%
\let\mspin\DERmspin%


\printedpage{29}
% head: A. Grothendieck : Dérivateurs, Ch. V, Édition des Universités de Montpellier II et Paris VII
% foot: Version 9 septembre 2001
$K_\infty$, dont la réunion est $K_\infty$. Mais d'autre part je ne sais donner aucune interprétation convaincante de $K_\infty$, comme de celle du ind-objet de Cat\,hot defini par $\uuK = (K_n)_{n\in\mathbf{N}}$, qui ind-représente le foncteur remarquable % [sic] defini -> défini, missing accent
\[
H = H_{\mathcal{M};x,y} : \Cathot^{o} \longrightarrow \Ens
\]
factorisant le foncteur
\msp{page 45}
\[
X \longmapsto H_{\mathcal{M};x,y}(X) \defeq \Hom_{\mathcal{M}(X)\Sigma_X^{-1}}(x_X, y_X),
\]
par la formule donc
\[
H(X) = \varinjlim_{n} \underbrace{\pi_0\uHom(X, K_n)}_{\Hom_{\Cathot}(X,K_n)} .
\]
Je ne sais donner d'interprétation de ce type ni de la catégorie $K_\infty$ elle-même, ni de son image dans Cat\,hot, ni de son image dans Hot. Il est vrai que si je restreins le foncteur $H$ à la sous-catégorie pleine de (Cat) formée des catégories \textit{finies} (Ob\,$X$ et Fl\,$X$ finis), alors on aura
\[
\uHom^{!}(X, K_\infty) = \uHom(X, K_\infty),
\]
et l'isomorphisme précédant s'interprète comme % [sic] précédant -> précédent, adjective intended
\[
\Hom_{\mathcal{M}(X)\Sigma_X^{-1}}(a_X, b_Y) \defeq H(X) \simeq \pi_0\uHom(X, K_\infty) \defeq \Hom_{\Cathot}(X, K_\infty).
\]
\texttt{[plutôt} $(a_X, b_X)$\texttt{]} Cela prouve du moins que la connaissance de $K_\infty$ est utile. Mais cette formule, pour $X$ restreint aux catégories \textit{finies,} ne définira pas $K_\infty$ comme objet de Cat\,hot, à isomorphisme unique près dans cette catégorie - \textit{sauf} dans le cas où $K_\infty$ est elle-même homotope à une catégorie
\msp{page 46}
finie. (Je ne me rends pas bien compte si ça arrive très souvent \ldots\ Mais on ne peut fonder une théorie générale des dérivateurs et des catégories de modèles `dérivables' sur une hypothèse de ce genre.)

Mais supposons que le foncteur $H$ envisagé se factorise même par Hot, i.e. qu'il transforme quasi-isomorphismes en bijections. On a vu que cela équivaut à la condition qu'il puisse se représenter par un système inductif sur $\mathbf{N}$ formé de catégories de Thomason-Kan. Il doit être facile de plus, si on y tient, de s'arranger pour que les foncteurs de transition soient des monomorphismes et pleinement fidèles, pour avoir une $\varinjlim$ sympathique. Ce sera bien sûr une catégorie de Thomason-Kan. Cette catégorie est-elle déterminée par le foncteur $H : \Cathot^{o} \longrightarrow \Ens$ à isomorphisme unique près dans Cat\,hot? Bien sûr, il n'est pas question qu'elle représente le foncteur $H$, sauf \texttt{[si]} celui-ci est bel et bien représentable, i.e. si le ind-objet $\uuK$ de Cat\,hot est essentiellement constant (ce qui ne doit pratiquement jamais être le cas). Notons déjà
\printedpage{30}
% head: A. Grothendieck : Dérivateurs, Ch. V, Édition des Universités de Montpellier II et Paris VII
% foot: Version 9 septembre 2001
\msp{page 47}
que si $u$ est un homomorphimse de ind-objets \textit{dans} Cat, ou, ce qui revient maintenant au même, un homomorphisme dans Cat % [sic] homomorphimse -> homomorphisme, transposed letters
\[
K_\infty \xrightarrow{u_\infty} K'_\infty
\]
tel que $\forall n$, $u_\infty(K_n) \subseteq K'_{n'}$ pour $n' = \tau(n)$ convenable, et si on suppose que l'homomorphisme de ind-objets induit un isomorphisme pour les foncteurs ind-représentables \textit{sur} Cat\,hot (par exemple, si $\uuK$ et $\uuK'$ ind-représentent le même foncteur $H$, et l'homomorphisme de ind-objets est celui qui correspond à $\mathrm{id}_H$), alors $u_\infty$ est forcément un isomorphisme \textit{dans} Hot, donc aussi (sauf erreur) dans Cat\,hot puisque $K_\infty$, $K'_\infty$ sont des catégories de Thomason-Kan. Mais ce qui n'est pas évident, c'est l'existence d'un homomorphisme $u_\infty$ comme \texttt{[ci-]}dessus, et s'il existe, qu'il soit unique (ou du moins canonique?) comme flèche de Cat\,hot.

Pour l'existence d'un $u_\infty$, on est ramené à prouver que si on a un diagramme à flèches pleines
\[
\begin{tikzcd}
K_n \arrow[r, "i_n", "\text{cofib.}"'] \arrow[d, "u_n"] & K_{n+1} \arrow[dr, dashed] & \\
K'_{\tau(n)} \arrow[rr, "\text{transition}"'] & & K'_N ,
\end{tikzcd}
\]
et si une flèche $u_{n+1}$ est donnée `en pointillés' \textit{dans} Cat\,hot, de façon à rendre le carré commutatif,
\msp{page 48}
alors on peut la réaliser par une flèche \textit{dans} Cat qui rende le diagramme commutatif pour de bon (dans Cat). On est ramené au même problème dans une situation
\[
\begin{tikzcd}
X \arrow[r, hook, shift right, "i"'] & Y , \arrow[l, dashed, shift right, "p"']
\end{tikzcd}
\]
où $X$ est une catégorie de Thomason, $i$ une cofibration de Thomason, $p$ une rétraction de $Y$ dans $X$ \textit{dans} Cat\,hot, on veut la réaliser par une vraie rétraction (\footnote[7]{\texttt{[un ?\ \ dans la marge]}}).

En admettant que ça soit le cas, donc que tout homomorphisme des ind-objets $\uuK$, $\uuK'$ \textit{dans} Cat\,hot se représente par un homomorphisme entre les ind-objets \textit{dans} Cat, donc aussi par un foncteur $K_\infty \longrightarrow K'_\infty$, reste la question si l'homomorphisme correspondant dans Cat\,hot est unique. Donc si on a deux foncteurs $f$, $g$ qui induisent le même homomorphisme sur les foncteurs ind-représentés sur Cat\,hot, si $f$ et $g$ sont homotopes. Donc

\textbf{Question}. Deux foncteurs $f, g : K_\infty \longrightarrow K'_\infty$ tels que pour tout $n \in \mathbf{N}$, $f_n$, $g_n$ qui s'en déduisent par restriction à $K_n$ sont homotopes (et même homotopes en tant que foncteurs de $K_n$ dans un $K'_N$ convenable), sont-ils homotopes? (\footnote[8]{$K$, $K'$ catégories de Thomason, les $K_n$, $K'_n$ aussi.})

\printedpage{31}
% head: A. Grothendieck : Dérivateurs, Ch. V, Édition des Universités de Montpellier II et Paris VII
% foot: Version 9 septembre 2001
\msp{page 49}
Je présume que non - qu'il peut arriver que la longueur minimum requise $l(n)$ pour un chemin dans $\uHom(K_n, K'_{\tau(n)})$ reliant $f_n$ et $g_n$, tende vers l'infini si $n \longrightarrow \infty$. L'hypothèse que les $K'_n$ sont des catégories de Thomason y change-t-il quelque chose (deux objets homotopes dans $\uHom(X, K')$, si $K'$ est de Thomason, sont-ils reliés par une flèche?)

Bien sûr, on peut prétendre que les catégories $\uCh_{\Sigma,n}(\mathcal{M}; a, b)$ ont été définies `canoniquement', et leur limite inductive $\uCh_{\Sigma}(\mathcal{M}; a, b)$ aussi. Mais la définition était assez arbitraire et abracadabrante - je suis convaincu que dix mathématiciens, se proposant de développer cette même idée des catégories de chemins, dans le même but que moi, arriveraient à dix définitions différentes! Pourtant, il n'est pas exclu que l'on finisse par trouver un système transitif d'isomorphismes dans Cat\,hot, entre les catégories de chemins `infinis' (i.e. de longueur non précisée) $K_\infty = K_\infty(\mathcal{M}; a, b)$ que nous aurons construites, dans le but
\msp{page 50}
principal (disons) de représenter le foncteur $H$ \textit{sur les catégories finies.} Et que les isomorphismes obtenus seront, bien entendu, des homomorphismes dans Ind(Cat\,hot). J'ai l'impression de ne pas avoir encore axiomatisé à fond la notion même de chemins, qui après tout a un sens géométrique assez évident et important par lui-même. Donc par ce qu'il faut entendre par `\textit{une} théorie des chemins dans une catégorie'. Une telle théorie $\mathcal{T}$ associerait donc à $\mathcal{M}$, $\Sigma$, $a$, $b$ un système inductif de catégories $\Ch^{\mathcal{T}}_{n,\Sigma}(\mathcal{M}; a, b)$. Et les contraintes géométriques naturelles à imposer à de telles théories, en vue aussi de leur utilisation pour expliciter
\[
\Hom_{\mathcal{M}(X)\Sigma_X^{-1}}(a_X, b_X) \xrightarrow{\ \sim\ } \pi_0\uHom(X, \Chh^{\mathcal{T}}(\mathcal{M}; a, b)) % reading resolved: two underlines, the ind-path object.
\]
donnerait un système transitif d'isomorphismes non seulement entre les ind-objets de Cat\,hot, mais peut-être aussi entre les $\Ch^{\mathcal{T}}_{\infty,\Sigma}(\mathcal{M}; a, b)$, pour des théories differentes. Si on dispose d'un tel énoncé, la catégorie $\Ch_{\infty,\Sigma}(\mathcal{M}; a, b)$, ou plus exactement ses images dans Cat\,hot et Hot, % [sic] differentes -> différentes, missing accent
\msp{page 51}
prendraient un sens convaincant.

Il y a d'autre part un autre point de vue, qui peut-être pourrait `canonifier' les catégories $\uCh_\infty$ en tant qu'objet de Cat\,hot et surtout de Hot (on passerait à Cat\,hot en appliquant le foncteur $\psi$, enveloppe de Thomason-Kan \ldots). En effet, on est intéressé à construire, en termes de $\mathcal{M}$, $\Sigma$, un `dérivateur' (ou un prédérivateur) $\mathbf{D}_{\mathcal{M},\Sigma}$ surtout sur le `domaine' des catégories finies. L'étendre à un domaine plus large est peut-être important dans un `deuxième souffle', mais peut être considéré comme du luxe dans le premier mouvement de fondation d'une théorie des dérivateurs. Ceci dit, on peut se poser des questions du type suivant.

Soit
\[
\Hot_0 \subset \Hot
\]
\printedpage{32}
% head: A. Grothendieck : Dérivateurs, Ch. V, Édition des Universités de Montpellier II et Paris VII
% foot: Version 9 septembre 2001
la sous-catégorie pleine de Hot, définie par les catégories finies (\footnote[9]{\textit{NB} Hot$_0$ est équivalente à une petite catégorie.}) - appelons les types d'homotopie `de type fini' (ou de \textit{présentation} finie?). On a un foncteur évident
\msp{page 52}
\[
\begin{aligned}
\Hot \quad &\longrightarrow \quad \Hot_0^{\wedge} = \uHom(\Hot_0^{o}, (\Ens)) \, ,\\
X \quad &\longmapsto \quad \tilde{X} = (\xi \longmapsto \Hom_{\Hot}(\xi, X)).
\end{aligned}
\]
Ce foncteur est \textit{conservatif} \texttt{[cela semble faux]}, puisque un homomorphisme $X \longrightarrow X'$ dans Hot est un isomorphisme s'il induit des isomorphismes pour les foncteurs $\tilde{X} \longrightarrow \tilde{X}'$, et même seulement pour leurs valeurs pour les sphères $S^n$ ($n \geq 0$). (Sauf erreur. Y a-t-il des canulars, parce que le critère d'équivalence faible par les $\pi_i$ ne s'applique pas directement, en l'absence de points bases. Ça ne doit pas être très sérieux \ldots\ \texttt{[cela semble en fait sérieux]}). Mais est-il peut-être même \textit{pleinement fidèle} (\footnote[10]{\texttt{[un ?\ \ dans la marge]}})? De sorte que Hot se réaliserait comme sous-catégorie pleine de $\Hot_0^{\wedge}$, contenant les foncteurs représentables?


\endgroup
% END UNIT Der-ChV-234/Der-ChV-234_p031-034.tex
\clearpage\setcounter{footnote}{0}\phantomsection\addcontentsline{toc}{section}{Chapitre VI}\markboth{Chapitre VI}{Chapitre VI}
% BEGIN UNIT Der-ChVI-345/Der-ChVI-345_p001-005.tex
\begingroup\makeatletter%
\NoAutoSpacing%
\newcommand{\printedpage}[1]{\marginpar{\footnotesize p.~#1}}%
\newcommand{\msp}[1]{\marginpar{\footnotesize [#1]}}%
\newcommand{\deq}{\overset{\text{déf}}{=}}%
\makeatother\let\printedpage\derpp\let\msp\dermsp%
\let\mspin\DERmspin%


% ===================== page 1 of the PDF (title page) =====================
\begin{center}
\printedpage{[sans numéro]}% no page number printed on this page
% head: (none)

\vspace*{6cm}

{\Large\textsf{LES DÉRIVATEURS}}

\bigskip

\textbf{ALEXANDRE GROTHENDIECK}

\vspace{2cm}

\textbf{Chapitre VI}

\medskip

\textbf{HOT}
\end{center}

\par

% ===================== page 2 of the PDF (editorial note) =====================
\vspace*{6cm}

\noindent\printedpage{[sans numéro]}% no page number printed on this page
% head: (none)
\par\noindent\textbf{Note des éditeurs de l’édition transcrite.}\par
Ce texte a été déchiffré et transcrit en \LaTeXe{} par M. Künzer. Il a été édité par M. Künzer, J. Malgoire, et G. Maltsiniotis. La transcription est aussi fidèle que possible au manuscrit. Pour les quelques corrections évidentes, ou rares commentaires des éditeurs, ainsi que pour la numérotation originale des pages du manuscrit, les caractères de machine à écrire \texttt{[typewriter]} entre crochets sont utilisés. Un point d'interrogation entre crochets signifie que l'on n'est pas sûr du mot qui précède.

\noindent Cette édition est provisoire. Les remarques, commentaires et corrections sont bienvenus. Envoyer un message à :

\noindent\textbf{[e-mail address removed]}

\hfill G. Maltsiniotis\hspace*{2cm}

\par

% ===================== printed page 1 =====================
\noindent\printedpage{1}
% head: A. Grothendieck : Dérivateurs, Ch. VI, Édition des Universités de Montpellier II et Paris VII
% foot: Version 9 novembre 2001
\msp{page 1}

\section*{\normalsize Le dérivateur \textmd{HOT}}

\begin{flushright}
14.11.90
\end{flushright}

\begin{flushright}
(Les premières feuilles, d'il y a deux semaines, ont été larguées, dépassées par VII)
\end{flushright}

\noindent Je définis Hot à partir de la catégorie de modèles
\[
\mathcal{M} = \mathrm{Cat},\ W \subseteq \mathrm{Fl}(\mathrm{Cat}),
\]
$W$ étant en principe les Hot-équivalences. Mais je vais dans la mesure du possible axiomatiser plus, en supposant donné un ensemble $W$ plus général, appelé `localisateur fondamental' (basic localiser), p.~ex. $W = W_{\mathbf{D}}$, où $\mathbf{D}$ est un dérivateur déjà donné (mais en supposant que son domaine soit (Cat) tout entier). Nous allons noter au fur et à mesure de quelles propriétés de $W$ nous aurons besoin. Aux fins de référence ultérieure, je vais rappeler ici celles que j'ai dégagées dans Pursuing Stacks.

\noindent\textbf{L1}\ \ Saturation de $W$

\begin{itemize}[leftmargin=2.5em]
\item[a)] \textit{Les identités} \texttt{[sont]} \textit{dans $W$.}
\item[b)] \textit{Si deux parmi $gf$, $f$, $g$} \texttt{[sont]} \textit{dans $W$, le troisième aussi, $X \xrightarrow{f} Y \xrightarrow{g} Z$.}
\item[c)] \textit{Si $gf$ et $fg$ sont dans $W$, alors $f$ et $g$ aussi, $X \overset{f}{\underset{g}{\rightleftarrows}} Y$ .}
\end{itemize}

\noindent\textbf{L2}\ \ $\boldsymbol{\Delta}^1$ \textit{est universellement $W$-asphérique, i.e. $\boldsymbol{\Delta}^1 \times X \longrightarrow X$ est dans $W$ pour tout $X$ dans} Cat. Cela implique

\noindent\textbf{L2 bis}\ \ $\boldsymbol{\Delta}^1$ \textit{est $W$-asphérique.} (C'était L3$'$ dans Pursuing Stacks.)

\noindent\textbf{L3}\ \ \textit{Si $X$ a un objet final, $X$ est $W$-asphérique.} (\footnote[1]{Il y a un L3$'$ dual de L3.})

\noindent Je note ici \texttt{[des]} variantes renforcées.

\noindent\textbf{L3 bis}\ \ \textit{Si $X$ est} Hot-\textit{asphérique, il est $W$-asphérique.}

\noindent\textbf{L3 ter}\ \ \textit{Si $X$ est} Hot-\textit{asphérique, il est universellement $W$-asphérique, i.e. $X \times Y \xrightarrow{\mathrm{pr}_2} Y$ est dans $W$ pour tout $Y$ dans} Cat. (\footnote[2]{Suffit pour entraîner L2, L3.})

\par

% ===================== printed page 2 =====================
\noindent\printedpage{2}
% head: A. Grothendieck : Dérivateurs, Ch. VI, Édition des Universités de Montpellier II et Paris VII
% foot: Version 9 novembre 2001
\textbf{L3 quater}\ \ \textit{Toute} Hot-\textit{équivalence est dans $W$.}

\noindent L3 quater $\Longrightarrow$ L3 ter $\Longrightarrow$ L3 bis $\Longrightarrow$ L3.

\noindent L3 ter $\Longrightarrow$ L2 $\Longrightarrow$ L2 bis.

\noindent\msp{page 2}

\begin{itemize}[leftmargin=2.5em]
\item[\textbf{L4}] \textit{Si $f : X \longrightarrow Y$ est $W$-asphérique (i.e. les $X/y \longrightarrow Y/y$ dans $W$, i.e. (si on a L3) si les $X/y$ sont $W$-asphériques), alors $f$ aussi.} (Axiome de localisation.) (\footnote[3]{il y a un L4$'$ dual de L4.}) \texttt{[`aussi'} \texttt{= dans} $W$\texttt{]}

\item[\textbf{L5}] (Axiome de fibration). \textit{Si}
\[
\begin{tikzcd}[column sep=small]
X \arrow[rr, "f"] \arrow[dr] & & Y \arrow[dl] \\
& S &
\end{tikzcd}
\]
\textit{avec $X$ et $Y$ fibrées sur $S$ et $f$ cartésien, si pour tout $s \in S$ $f_s : X_s \longrightarrow Y_s$ est dans $W$, $f$ aussi.}
\end{itemize}

\noindent Dans Pursuing Stacks j'ai supposé toujours \fbox{L1, L3, L4} (impliquant L2) - c'étaient ces conditions qui définissaient la notion de `localisateur' (fondamental) dans Cat. (\footnote[4]{Si $W$ \texttt{[est]} un localisateur, alors les catégories contractiles sont $W$-asphériques, les homotopismes sont dans $W$, les applications homotopes sont égales dans (Cat)$W^{-1} \deq \mathrm{Hot}_W$.}) D'autre part, \fbox{L1, L2 bis, L5} impliquent L1,2,3,4,5.

\noindent\textit{On va supposer ici aussi qu'on ait au moins} L1, L3, L4 (i.e. $W$ un localisateur). En appliquant la construction de I, p.~72-\textbf{75,} on voit que $f : X \longrightarrow Y$ est dans $W$ si et seulement si $f^o : X^o \longrightarrow Y^o$ l'est. Donc L1, L3, L4 impliquent L3$'$, L4$'$ (axiomes duaux de L3, L4).

\noindent Il y avait des axiomes complémentaires

\begin{itemize}[leftmargin=2.5em]
\item[\textbf{La}] $f \in W \Longrightarrow \pi_0(f)$ \textit{bijectif.}

\item[\textbf{Lb}] \textit{Le foncteur canonique de localisation}
\[
\mathrm{Cat} \longrightarrow (\mathrm{Cat})W^{-1} \deq \mathrm{Hot}_W
\]
\textit{commute aux sommes finies.}

\item[\textbf{Lc}] \textit{Itou pour produits finis.}
\end{itemize}

\par

% ===================== printed page 3 =====================
\noindent\printedpage{3}
% head: A. Grothendieck : Dérivateurs, Ch. VI, Édition des Universités de Montpellier II et Paris VII
% foot: Version 9 novembre 2001
Il faudrait quand même dégager des conditions concrètes sur $W$ pour la validité de Lb, Lc. Les conditions L3 quater et La s'écrivent respectivement
\[
W_\infty \subseteq W,\ \ W \subseteq W_0.
\]

\noindent\msp{page 3}

\noindent Leur conjonction est l'axiome des plus raisonnables
\[
W_\infty \subseteq W \subseteq W_0.
\]
Bien sûr, l'axiome L5 a maintenant l'air un peu étriqué, i.e. faiblard. On va voir très vite de quoi on aura besoin!

\noindent On veut donc définir un dérivateur
\[
\mathrm{HOT}_W
\]
à partir de la catégorie de modèles
\[
(\mathrm{Cat}, W \subseteq \mathrm{Fl}(\mathrm{Cat})).
\]
Donc pour définir
\[
\mathrm{HOT}_W(S) \qquad\qquad\qquad (S \in \mathrm{Ob}\,\mathrm{Cat}),
\]
la récette générale est de prendre la catégorie % [sic] récette -> recette, misspelling of "recette"
\[
\mathcal{M}(S) = \mathrm{Cat}(S) = \underline{\mathrm{Hom}}(S^o, (\mathrm{Cat})),
\]
d'y introduire l'ensemble de flèches
\[
\left\{
\begin{aligned}
& W_S \subseteq \mathrm{Fl}(\mathcal{M}(S)) \\
& u \in W_S \Longleftrightarrow u_s \in W \ \ \forall\, s \in \mathrm{Ob}\,S.
\end{aligned}
\right.
\]
On doit poser
\[
\mathrm{HOT}_W(S) \deq \underline{\mathrm{Hom}}(S^o, \mathrm{Cat})W_S^{-1}.
\]
Mais la catégorie $\underline{\mathrm{Hom}}(S^o, \mathrm{Cat})$ s'interprète, à isomorphisme canonique près, comme celle
\[
\underline{\mathrm{Fibsc}}\,S
\]
des catégories $X$ Cat-fibrées sur $S$, munies d'un `scindage', i.e. d'un choix de foncteurs changement de base $u^* : X_t \longrightarrow X_s$ (pour $t \xrightarrow{u} s$ dans $S$) \texttt{[plutôt} $u^* : X_s \longrightarrow X_t$\texttt{]}, avec de plus \textit{transitivité stricte}
\[
(vu)^* = u^* v^* \qquad\qquad\qquad \text{si}\ \ t \xrightarrow{u} s \xrightarrow{v} r\ \ \text{dans}\ S
\]

\noindent\msp{page 4}


\endgroup
% END UNIT Der-ChVI-345/Der-ChVI-345_p001-005.tex
% BEGIN UNIT Der-ChVI-345/Der-ChVI-345_p006-010.tex
\begingroup\makeatletter%
\newcommand{\printedpage}[1]{\marginpar{\footnotesize p.~#1}}%
\newcommand{\msp}[1]{\marginpar{\footnotesize [#1]}}%
\newcommand{\ul}[1]{\underline{\mathrm{#1}}}%
\newcommand{\xline}[1]{\mathrel{\overset{#1}{\rule[0.6ex]{2em}{0.4pt}}}}%
\makeatother\let\printedpage\derpp\let\msp\dermsp%
\let\mspin\DERmspin%


\printedpage{4}
% head: A. Grothendieck : Dérivateurs, Ch. VI, Édition des Universités de Montpellier II et Paris VII
et non seulement transitivité à isomorphisme canonique près. Ainsi on peut écrire aussi
\[
\mathrm{HOT}_W(S) \xrightarrow{\ \sim\ } (\ul{Fibsc}\,S)W_S^{-1}, \tag*{$(*)$}
\]
où cette fois $W_S$ est interprété comme
\[
W_S \subseteq \mathrm{Fl}(\ul{Fibsc}\,S),
\]
savoir
\[
W_S = \{u \in \mathrm{Fl}(\ul{Fibsc}\,S) \mid u_s \in W \text{ pour tout } s \in \mathrm{Ob}\,S\}
\]
(foncteurs strictement cartésiens, i.e. commutant aux scindages, qui induisent des $W$-équivalences entre les catégories fibres).

La flèche dans $(*)$ est bel et bien un isomorphisme de catégories, pas seulement une équivalence. Mais on va définir d'autres catégories, plus commodes à certains égards, qui seront \textit{équivalentes} (on l'espère), mais sûrement pas isomorphes, à $\mathrm{HOT}_W(S)$. (On s'occupera plus tard de la variance de $\mathrm{HOT}_W(S)$ en $S$.)

Les catégories $\ul{Hom}(S^o,\mathrm{Cat})$, $\ul{Fibsc}\,S$, ainsi que les variantes que nous allons introduire maintenant, sont en fait des 2-catégories. Mais pour décrire $\mathrm{HOT}_W(S)$, seule les structures de 1-catégories sont utilisées. On a l'impression que c'est bien % [sic] seule les structures -> seules les structures, accord du pluriel
\msp{page 5}
grossier, qu'on perd quelque chose ce faisant. Mais ce qui est perdu va se retrouver, et surabondamment, dans la structure des \underline{Hom} internes dans les $\mathrm{HOT}(S)$, si on se `rappelle' que ces \underline{Hom} internes, donnant lieu à des `types d'homotopie relatifs' $\ul{Hom}_{\mathrm{HOT}_W}(S;\xi,\eta)$ (si $\xi$, $\eta$ en sont), peuvent s'interpréter en termes de $\infty$-champs en groupoïdes sur $S$, avec toute la richesse de structure qui leur appartient. Mais ceci est une autre histoire\dots

En fait, je vois un diagramme de 13 catégories de `categories sur $S$\,', à examiner à présent % [sic] categories -> catégories, accent manquant
\[
\begin{tikzcd}[column sep=small]
 & \ul{Fibsc}_0S \arrow[dl, hook] \arrow[d, "\mathrm{fid}"] & & \ul{Cofibsc}_0S \arrow[d] \arrow[dr, hook] & \\
\ul{Fibsc}\,S \arrow[d, "\mathrm{fid}"] & \ul{Fib}_0S \arrow[dl, hook] & & \ul{Cofib}_0S \arrow[dr, hook] & \ul{Cofibsc}\,S \arrow[d, "\mathrm{fid}"] \\
\ul{Fib}\,S \arrow[dr, hook] & & \ul{Parf}\,S \arrow[dl, hook] \arrow[d, hook] \arrow[dr, hook] & & \ul{Cofib}\,S \arrow[dl, hook] \\
 & \ul{Liss}\,S \arrow[dr, hook] & \ul{Fibhot}\,S \arrow[d, hook] & \ul{Prop}\,S \arrow[dl, hook] & \\
 & & \mathrm{Cat}/S & &
\end{tikzcd}
\]
% foot: Version 9 novembre 2001

\printedpage{5}
% head: A. Grothendieck : Dérivateurs, Ch. VI, Édition des Universités de Montpellier II et Paris VII
où \underline{Fib}, \underline{Cofib} désignent les catégories Cat-fibrées, Cat-cofibrées sur $S$, \underline{Liss} et \underline{Prop} les catégories
\msp{page 6}
lisses resp. propres sur $S$, \underline{Parf} celles qui sont à la fois lisses et propres (fibrations parfaites), enfin \underline{Fibhot} les catégories Hot-fibrées sur $S$. Toutes les flèches sont des foncteurs strictement pleinement fidèles, à l'exception des deux (plutôt quatre) flèches verticales supérieures, qui sont fidèles sans plus. Je viens d'en rajouter encore 4, de catégories : $\ul{Fibsc}_0$, $\ul{Fib}_0$, $\ul{Cofibsc}_0$, $\ul{Cofib}_0$ sont les sous-catégories pleines des catégories correspondantes sans indice, formées des catégories fibrées resp. cofibrées $X$ sur $S$, pour lesquelles les foncteurs changement de base resp. cochangement de base
\[
u^* : X_t \longrightarrow X_s, \qquad\quad \text{resp.} \qquad\quad u_* : X_s \longrightarrow X_t,
\]
pour $s \xrightarrow{u} t$ dans $S$, sont des Hot-équivalences. En termes de foncteurs
\[
S^o \xrightarrow{F} \mathrm{Cat} \qquad\quad \text{resp.} \qquad\quad S \xrightarrow{F'} \mathrm{Cat},
\]
cela signifie que ce foncteur transforme flèches quelconques dans $S$ en éléments de $W$ dans Cat, i.e. qu'on a un diagramme commutatif
\[
\begin{tikzcd}[row sep=large]
S^o \arrow[r, "F"] \arrow[d, "\mathrm{can}"'] & \mathrm{Cat} \arrow[d, "\mathrm{can}"] \\
\Pi(S^o) \simeq (\Pi S)^o = S^o(\mathrm{Fl}\,S^o)^{-1} \arrow[r, "\bar F"'] & \mathrm{Hot}_W = (\mathrm{Cat})W^{-1}
\end{tikzcd}
\]
(et de même pour $F'$), i.e. qu'on a un `système
\msp{page 7}
local' $\bar F$ sur $S$, à valeurs dans $\mathrm{Hot}_W$, défini à l'aide de $F$. Je rappelle qu'on trouve également un tel système local dans le cas d'une catégorie Hot-fibrée sur $S$ (et même un système local à valeurs dans Hot lui-même, au lieu de $\mathrm{Hot}_W$ - on utilise ici tacitement que $W_\infty \subseteq W$, i.e. \fbox{L3 quater} (\footnote[5]{On va le supposer vrai dans toute la suite.})). Et il en est de même a fortiori pour les objets de $\ul{Parf}\,S$. Pour $\ul{Liss}\,S$ resp. $\ul{Prop}\,S$, on trouve simplement un foncteur
\[
S^o \xrightarrow{\bar F} \mathrm{Hot} \qquad\quad \text{resp.} \qquad\quad S \xrightarrow{\bar F'} \mathrm{Hot},
\]
d'où, en composant avec $\mathrm{Hot} \longrightarrow \mathrm{Hot}_W$, des foncteurs à valeurs dans $\mathrm{Hot}_W$.

Dans chacune de ces 14 \texttt{[13?]} catégories, on introduit un ensemble de flèches, qui serviront à localiser, et on le désignera toujours par la même notation $W_S$. Dans toutes ces catégories sauf seulement $\mathrm{Cat}/S$, $W_S$ désigne la $W$-équivalence `fibre par fibre'. Dans $\mathrm{Cat}/S$ d'autre part, ce sera la $W$-équivalence sans plus. Toutes les flèches du diagramme sont compatibles avec les ensembles localisants de flèches.
% foot: Version 9 novembre 2001

\printedpage{6}
% head: A. Grothendieck : Dérivateurs, Ch. VI, Édition des Universités de Montpellier II et Paris VII
\msp{page 8}
\[
\begin{tikzcd}[column sep=small, row sep=large]
 & \ul{Fibsc}_0S \arrow[dl, hook, "\alpha'"'] \arrow[d] & & \ul{Cofibsc}_0S \arrow[d] \arrow[dr, hook, "\beta'"] & \\
\ul{Fibsc}S \arrow[d] & \ul{Fib}_0S \arrow[dl, hook, "\alpha"'] \arrow[ddr, dotted, bend right=20] & & \ul{Cofib}_0S \arrow[dr, hook, "\beta"] \arrow[ddl, dotted, bend left=20] & \ul{Cofibsc}\,S \arrow[d] \\
\ul{Fib}S \arrow[dr, hook] & & \ul{Parf}\,S \arrow[dl, hook] \arrow[d, hook] \arrow[dr, hook] & & \ul{Cofib}\,S \arrow[dl, hook] \\
 & \ul{Liss}\,S \arrow[dr, hook] & \ul{Fibhot}\,S \arrow[d, hook] \arrow[uuul, dotted, bend left=10, "\Phi_0"] \arrow[uuur, dotted, bend right=10, "\Psi_0"'] & \ul{Prop}\,S \arrow[dl, hook] & \\
 & & \mathrm{Cat}/S \arrow[uuull, dotted, bend left=60, "\Phi"] \arrow[uuurr, dotted, bend right=60, "\Psi"'] & &
\end{tikzcd}
\]
% ?? uncertain: starting point of the dotted arrows labelled Phi_0 and Psi_0 (the dotted lines run alongside Parf S down to Fibhot S) -- Fibhot S or Parf S

\texttt{[Dans ce diagramme, il y a aussi} $\ul{Liss}_0$ \texttt{et} $\ul{Prop}_0$ \texttt{mais sans connexion aux autres objets.]}

Je suis quasiment sûr du

\textbf{Theorème.} (\footnote[6]{C'est totalement canulé, cf. commentaires page 32 ff., solution des perplexités pages 41 ff. Voir surtout le Théorème-Scholie récapitulatif p. 70-81.}) \textit{Tous les foncteurs précédents induisent des équivalences entre les catégories localisées.} % [sic] Theorème -> Théorème, accent manquant

Comme il y a 16 foncteurs, ça a l'air d'un exercise de patience. Le principe général, c'est qu'on arrive à définir des foncteurs en sens inverse, compatibles avec les $W_S$, de telle façon que les composés soient reliées aux foncteurs identiques par des flèches qui sont dans $W_S$, de sorte que par passage aux localisés ces foncteurs induisent des (équivalences) quasi-inverses des foncteurs déduits de ceux du diagramme. Tous ces foncteurs en sens opposé résultent essentiellement par composition, de trois foncteurs bien familiers. Deux de ceux-ci sont indiqués en pointillés sur le diagramme, comme $\Phi$, $\Psi$. % [sic] exercise -> exercice, orthographe française ; % [sic] reliées -> reliés, accord avec « les composés »

\msp{page 9}
Le troisième est un foncteur canonique
\[
C : \mathrm{Cat}/S \longrightarrow \ul{Parf}\,S,
\]
qu'on va expliciter également.
% foot: Version 9 novembre 2001

\printedpage{7}
% head: A. Grothendieck : Dérivateurs, Ch. VI, Édition des Universités de Montpellier II et Paris VII
Le foncteur $\Phi$ (resp. $\Psi$) appliqué à $\ul{Fib}_0S$ (resp. $\ul{Cofib}_0S$), induit un foncteur
\[
\ul{Fib}_0S \longrightarrow \ul{Fibsc}_0S \qquad\quad \text{resp.} \qquad\quad \ul{Cofib}_0S \longrightarrow \ul{Cofibsc}_0S \tag*{$(*)$}
\]
(ce qu'il faudra vérifier). Cela fournit toutes les flèches `inverses' dont on a besoin, à l'exception seulement de $\alpha$, $\beta$, $\alpha'$, $\beta'$. Mais si on a prouvé, par la méthode indiquée, que toutes les flèches du diagramme induisent des équivalences sur les catégories localisées, pour prouver qu'il en est encore de même pour $\alpha$, $\beta$, $\alpha'$, $\beta'$, il suffit de le prouver pour les deux foncteurs
\[
\ul{Fibsc}_0S \longrightarrow \ul{Fib}\,S, \qquad\qquad \ul{Cofibsc}_0S \longrightarrow \ul{Cofib}\,S,
\]
donc il reste à définir les foncteurs en sens inverse
\[
\ul{Fib}\,S \longrightarrow \ul{Fibsc}_0S, \qquad\qquad \ul{Cofib}\,S \longrightarrow \ul{Cofibsc}_0S,
\]
et comme les sources sont des sous-catégories pleines de $\mathrm{Cat}/S$, il suffira de définir
\[
\Phi_0 : \mathrm{Cat}/S \longrightarrow \ul{Fibsc}_0S, \qquad\qquad \Psi_0 : \mathrm{Cat}/S \longrightarrow \ul{Cofibsc}_0S.
\]
Pour ceci, on posera
\[
\Phi_0 = \Phi r C, \qquad\qquad \Psi_0 = \Psi r C, \tag*{$(**)$}
\]
\msp{page 10}
où $C : \mathrm{Cat}/S \longrightarrow \ul{Parf}\,S$, $\Phi$ et $\Psi$ sont les 3 foncteurs `rétrogrades' inversés plus haut, et $r : \ul{Parf}(S) \hookrightarrow \mathrm{Cat}/S$ est l'inclusion. Il faut donc vérifier que $\Phi_0$, $\Psi_0$ définis par $(**)$, qui a priori ont comme catégories but $\ul{Fibsc}\,S$ et $\ul{Cofibsc}\,S$, prennent leur valeur dans les sous-catégories avec indice 0.

Je vais donner un énoncé d'équivalence unique, qui couvrira (presque) tous les 16 cas à la fois :

\textbf{Corollaire.} \textit{Considérons le diagramme}
\[
\begin{tikzcd}[column sep=small, row sep=small]
 & \ul{Parf}\,S \arrow[dd, "r"'] & \\
\ul{Fibsc}_0S \arrow[dr, "p"'] & & \ul{Cofibsc}_0S \arrow[dl, "q"] \\
 & \mathrm{Cat}/S &
\end{tikzcd}
\]
\textit{et soit $\mathcal{F}$ une sous-catégorie pleine de $\mathrm{Cat}/S$, telle que l'un des foncteurs $p$, $q$, $r$ se factorise par $\mathcal{F}$. Soit $W_{\mathcal{F}}$ induit sur $\mathcal{F}$ par $W_S$ dans $\mathrm{Cat}/S$. Alors $\mathcal{F} \hookrightarrow \mathrm{Cat}/S$ induit une équivalence sur les catégories localisées.}

\msp{page 11}
\textsc{Démonstration}. Je dis qu'il suffit de prouver, pour les trois foncteurs $p$, $q$, $r$, et pour les foncteurs en sens inverse
% foot: Version 9 novembre 2001
\printedpage{8}
% head: A. Grothendieck : Dérivateurs, Ch. VI, Édition des Universités de Montpellier II et Paris VII
\[
\begin{tikzcd}[column sep=small, row sep=small]
 & \ul{Parf}\,S & \\
\ul{Fibsc}_0S & & \ul{Cofibsc}_0S \\
 & \mathrm{Cat}/S \arrow[uu, "C"] \arrow[ul, "\Phi_0"] \arrow[ur, "\Psi_0"'] &
\end{tikzcd}
\]
que les six composés
\[
p\Phi_0, \Phi_0p, q\Psi_0, \Psi_0q, rC, Cr
\]
sont reliés aux foncteurs identiques (dans la catégorie pertinente) par une flèche qui est dans le localisateur correspondant. Pour chacun de ces trois couples de foncteurs $(p,\Phi_0)$, $(q,\Psi_0)$, $(r,C)$, cela prouvera déjà qu'en passant aux localisés, ils induisent des équivalences quasi-inverses l'une de l'autre. Mais on a mieux à l'aide du lemme formel suivant.

\textbf{Lemme.} \textit{Soient $(C,W)$, $(C',W')$ deux catégories avec localisateurs $W \subseteq \mathrm{Fl}\,C$, $W' \subseteq \mathrm{Fl}\,C'$. Soit}
\[
p : C \longrightarrow C'
\]
\textit{un foncteur compatible avec les localisateurs, et supposons qu'il existe un foncteur $\varphi : C' \longrightarrow C$ également compatible, et tel que l'on ait des flèches (à direction non précisée)}
\[
\varphi p \xline{u} id_C, \qquad\qquad p\varphi \xline{v} id_{C'},
\]
\textit{qui soient dans $W$ resp. $W'$ `argument par argument'.}

\msp{page 12}
\textit{Alors $p$ et $\varphi$ induisent des équivalences quasi-inverses l'une de l'autre en passant aux catégories localisées (elles n'ont pas de merite). De plus, pour toute sous-catégorie pleine $C'_0$ de $C'$, telle que $p$ se factorise (en $p_0 : C \longrightarrow C'_0$) par $C'_0$, si on désigne par $W'_0$ le localisateur induit dans $C'_0$ par $C'$, alors $p_0 : C \longrightarrow C'_0$ et $\varphi_0 = \varphi|C'_0 : C' \longrightarrow C$} \texttt{[plutôt} $C'_0 \longrightarrow C$\texttt{]} \textit{induisent également des équivalences quasi-inverses l'une de l'autre sur les localisés.} % [sic] merite -> mérite, accent manquant

En effet, $\varphi_0p_0 = (\varphi i)p_0 = \varphi(ip_0) = \varphi p \xline{u} \mathrm{id}_C$ (ici $i : C'_0 \hookrightarrow C'$ désigne l'inclusion), et $F \overset{\text{déf}}{=} p_0\varphi_0 = p_0\varphi i$ satisfait à
\[
iF = \underbrace{ip_0}_{p}\varphi i = (p\varphi)i.
\]
Pour définir $F \xline{v_0} \mathrm{id}_{C'_0}$, qui soit dans $W'_0$ argument par argument,
\msp{page 13}
il suffit de définir
\[
\underbrace{iF}_{=\,(p\varphi)i} \xline{\bar v_0} i,
\]
qui soit dans $W'$ argument par argument, puisque $W'_0$ induit par $W'$ et que $F \longmapsto iF$
\[
\ul{Hom}(C'_0,C'_0) \longrightarrow \ul{Hom}(C'_0,C')
\]
% foot: Version 9 novembre 2001


\endgroup
% END UNIT Der-ChVI-345/Der-ChVI-345_p006-010.tex
% BEGIN UNIT Der-ChVI-345/Der-ChVI-345_p011-015.tex
\begingroup\makeatletter%
\newcommand{\printedpage}[1]{\marginpar{\footnotesize p.~#1}}%
\newcommand{\msp}[1]{\marginpar{\footnotesize [#1]}}%
\newcommand{\cat}[1]{\underline{\mathrm{#1}}}%
\newcommand{\circled}[1]{\tikz[baseline=(c.base)]\node[draw,circle,inner sep=1pt](c){#1};}%
\makeatother\let\printedpage\derpp\let\msp\dermsp%
\let\mspin\DERmspin%


\printedpage{9}
% head: A. Grothendieck : Dérivateurs, Ch. VI, Édition des Universités de Montpellier II et Paris VII
est pleinement fidèle ($i$ l'étant). Or il suffit de prendre
\[
\bar{v}_0 = v * i : (p\varphi)\circ i \mathrel{\relbar\joinrel\relbar} \mathrm{id}_{C'}\circ i.
\] % ?? uncertain: the arrow between (p\varphi)\circ i and id_{C'}\circ i is printed as a plain line with no visible head -- plain line / arrow

Pour résumer, il faut montrer que le composé des couples de foncteurs quasi-inverses marche dans les cinq cas suivants
\[
\begin{tikzcd}[column sep=large, row sep=huge]
\cat{Fibsc} \arrow[rrd, shift left] & \cat{Fibsc}_0 \arrow[l, hook'] \arrow[rd, shift left] & \cat{Parf} \arrow[d, shift left] & \cat{Cofibsc}_0 \arrow[r, hook] \arrow[ld, shift left] & \cat{Cofibsc} \arrow[lld, shift left] \\
& & \mathrm{Cat}/S \arrow[llu, shift left, "\Phi"] \arrow[lu, shift left, "\Phi_0"] \arrow[u, shift left, "C"] \arrow[ru, shift left, "\Psi_0"] \arrow[rru, shift left, "\Psi"] & &
\end{tikzcd}
\]

D'ailleurs par dualité, on est ramené aux trois cas de $\cat{Fibsc}$, $\cat{Fibsc}_0$, $\cat{Parf}$. (Peut-être le cas $\cat{Fibsc}_0$, avec $\Phi_0$ défini en termes de $\Phi$ et de $C$, se ramènera-t-il aux deux cas qui l'entourent, de $\cat{Fibsc}$ et $\cat{Parf}$.)

Je vois que je dois rectifier le tir, tout d'abord parce que j'ai oublié que les foncteurs
\msp{page 14}
$\Phi$, $\Psi$ que j'avais en vue ne sont \textit{pas} compatibles avec les localiseurs \texttt{[= localisateurs]} - une flèche $f : X \longrightarrow Y$ de $\mathrm{Cat}/S$ telle que la flèche correspondante dans Cat soit dans $W$, ne va pas induire une flèche dans $\cat{Fibsc}\,S$ ou $\cat{Cofibsc}\,S$ qui soit dans $W$ sur chaque fibre. Ce n'est pas grave, car on peut remplacer les foncteurs par $\Phi' = \alpha'\Phi_0$, $\Psi' = \beta'\Psi_0$, qui sauf erreur doivent marcher. Mais en vue du Lemme précédent, il y a une perplexité, car j'ai l'air d'avoir prouvé que si $\mathcal{F}$ est une sous-catégorie pleine de $\mathrm{Cat}/S$ dont je suppose seulement qu'elle contient l'image essentielle d'une des trois catégories $\cat{Fib}_0S$, $\cat{Cofib}_0S$, ou $\cat{Parf}\,S$, alors l'inclusion dans $\mathrm{Cat}/S$ induit une équivalence des catégories localisées, en prenant dans $\mathcal{F}$ comme localisateur celui induit par $\mathrm{Cat}/S$, i.e.
\begin{tikzcd}[column sep=small, row sep=small, baseline=(\tikzcdmatrixname-2-2.base)]
f : X \arrow[rr] \arrow[dr] & & Y \arrow[dl] \\
& S &
\end{tikzcd}
est dans $W_{\mathcal{F}}$ si et seulement si en tant que flèche de Cat, $f$ est dans $W$. Or ceci coïncide bien (via un théorème non trivial sur lequel je devrai revenir dans le cadre d'un localisateur fondamental général $W$) avec l'ensemble de flèches de localisation `fibre par fibre', dans le cas de $\mathcal{F} = \cat{Parf}\,S$, $\cat{Fibhot}\,S$, $\cat{Fib}_0S$, $\cat{Cofib}_0S$.

\msp{page 15}
Par contre, ce n'est nullement le cas en prenant $\mathcal{F} = \cat{Fib}\,S, \cat{Liss}\,S$ ou $\cat{Cofib}\,S$. Ayant choisi dans ces catégories comme localisateurs la notion de $W$-équivalences `fibre par fibre', j'ai donc l'air de dire que je trouve pour ces quatre catégories les \textit{mêmes} localisées (à équivalence près) que si je prenais le localisateur beaucoup plus grossier hérité par $\mathrm{Cat}/S$. Cela semble difficile à croire, et il faudra que je fasse très attention en démontrant les énoncés-clef.

Pour voir que les flèches en traits pleins du grand diagramme p.~6 sont toutes compatibles avec les localis(at)eurs \texttt{[sic]}, donc aussi les flèches composées au moyen de celles-ci, j'ai à utiliser déjà ceci (pour le cas des flèches $\cat{Liss} \longrightarrow \mathrm{Cat}/S$, $\cat{Prop} \longrightarrow \mathrm{Cat}/S$):
% foot: Version 9 novembre 2001

\printedpage{10}
% head: A. Grothendieck : Dérivateurs, Ch. VI, Édition des Universités de Montpellier II et Paris VII
\circled{1}\ \ Si $X, Y$ lisses ou propres sur $S$, et $f : X \longrightarrow Y$ un $S$-foncteur, si $f_s : X_s \longrightarrow Y_s$ est dans $W$ pour tout $s \in S$, alors $f \in W$. (\footnote[7]{OK si on suppose L3 quater et L4.}).

Cela renforce considérablement la condition L5 sur $W$, où on supposait que $f$ est un foncteur \textit{cartésien} entre catégories \textit{fibrées} sur $S$. (À dégager, pour les différentes propriétés dont on a besoin pour $W$, à quelles propriétés standard on peut les ramener. Tout sera OK si $W = W_\infty$.)

\msp{page 16}
I) \textbf{Définition de} $\Phi$, $\Psi$.

On va calquer la construction canonique de $\cat{Ch}(X) \longrightarrow X \times X$ (ou $\cat{Ch}_\infty(X) \longrightarrow X \times X$) factorisant $X \xrightarrow{\mathrm{diag}} X \times X$, et la factorisation canonique correspondante d'un $f : X' \longrightarrow X$ en $X' \xrightarrow{i} \bar{X}' \xrightarrow{\bar{f}} Y$ \texttt{[plutôt} $\bar{X}' \xrightarrow{\bar{f}} X$\texttt{]}. On remplacera simplement $\cat{Ch}(X)$ par la sous-catégorie $\cat{Ch}_+(X)$, laquelle n'est autre que la catégorie des flèches de $X$
\[
\cat{Fl}(X) = \cat{Hom}(\boldsymbol{\Delta}_1, X),
\]
muni des foncteurs
\[
X \xrightarrow{i_X} \cat{Fl}(X) \xrightarrow{p_X=(\sigma_X,\beta_X)} X \times X,
\]
où $i_X$ est la flèche $x \longmapsto \mathrm{id}_x$, et $\sigma_X$, $\beta_X$ sont les foncteurs source et but. La fibre de $p_X$ en $(a,b) \in \mathrm{Ob}(X \times X)$ est la catégorie discrète définie par l'ensemble $\mathrm{Hom}(a,b)$, ça dépend grandement de $(a,b)$ et pas seulement de sa composante connexe dans $X \times X$ - donc $p_X$ n'a rien d'une Hot-fibration. Il n'est en géneral ni propre ni lisse ni rien, par contre les deux flèches composées avec $\mathrm{pr}_1$ et $\mathrm{pr}_2$ % [sic] géneral -> général, missing accent
\[
\begin{tikzcd}[column sep=small]
& \cat{Fl}(X) \arrow[ld, "\sigma_X"'] \arrow[rd, "\beta_X"] & \\
X & & X
\end{tikzcd}
\]
sont sympathiques, on a en effet

\msp{page 17}
\textbf{Lemme.} \textit{On a}
\[
\sigma_X^{-1}(\{x\}) \simeq x\backslash X, \qquad\qquad \beta_X^{-1}(\{x\}) \simeq X/x.
\]
\textit{Le foncteur} $\sigma_X$ \textit{est} Cat-\textit{fibrant à fibres contractiles (donc asphériques, donc} $W$\textit{-asphériques), le foncteur} $\beta_X$ \textit{cofibrant à fibres contractiles (donc asphériques, donc} $W$\textit{-asphériques). Ce sont donc aussi des} Hot-\textit{équivalences, a fortiori des} $W$\textit{-équivalences} (\footnote[8]{En fait, ce sont même des homotopismes, cf. remarque marginale plus bas (plus générale - on retrouve le cas de $\sigma_X$ en faisant $X' = X$.)}).

\textbf{Corollaire 1.} \textit{Le foncteur} $i_X : X \longrightarrow \cat{Fl}(X)$ \textit{est une} Hot-\textit{équivalence} (puisque son composé avec la Hot-équivalence $\sigma_X$ est l'identité).
% foot: Version 9 novembre 2001

\printedpage{11}
% head: A. Grothendieck : Dérivateurs, Ch. VI, Édition des Universités de Montpellier II et Paris VII
Voici la construction générale de la factorisation de $f : X' \longrightarrow X$, correspondant au choix précédent d'une `catégorie de chemins' $\mathcal{C}(X) = \cat{Fl}(X)$
\[
\begin{tikzcd}
\bar{X}' \arrow[d, "\pi_f"] \arrow[rd] \arrow[rrd, "\bar{f}"] & & \\
X' \arrow[u, dotted, bend left=60, "i_f"] \arrow[rd, "f"'] \arrow[r, phantom, "\text{\footnotesize Cart.}"] & \cat{Fl}(X) \arrow[r, "\beta_X"'] \arrow[d, "\sigma_X"'] & X \\
& X &
\end{tikzcd}
\]
où $i_f$ est déduit de la section $i$ \texttt{[plutôt} $i_X$\texttt{]} de $\cat{Fl}(X)$ sur $X$ (\footnote[9]{\textbf{NB} $\pi_f i_f = \mathrm{id}$, et $i_f\pi_f \longrightarrow \mathrm{id}$, donc $i_f$, $\pi_f$ sont des homotopismes inverses l'un de l'autre dans Cat\,hot. Cela précise \texttt{[le]} cor.~3, page 18.}). Ainsi, $\bar{X}'$ est la catégorie des couples
\[
\{(x',u) \mid x' \in \mathrm{Ob}\,X',\ \underbrace{u}_{\in\mathrm{Fl}\,(X)} : f(x') \longrightarrow y\},
\]
le foncteur $\bar{f}$ est défini par
\[
\bar{f}(x',u) = \mathrm{but}(u) = y
\]
et $i_f$ par
\[
i_f(x') = (x', \mathrm{id}_{f(x')} : f(x') \xrightarrow{\mathrm{id}} f(x')),
\]
donc on a bien
\[
f(x') = \bar{f} i_f(x').
\]

\msp{page 18}
\textbf{Corollaire 2.} \textit{Le foncteur}
\[
\bar{f} : \bar{X}' \longrightarrow X
\]
\textit{a comme fibres les catégories $X'/x$, il est cofibrant, les foncteurs de changement de base sont les foncteurs $X'/x \longrightarrow X'/y$ associés aux flèches $x \longrightarrow y$ dans $X$.}

En fait, on voit que cette catégorie est non seulement cofibrée, mais canoniquement \textit{scindée.}

(\footnote[10]{\textbf{Corollaire 3.} \textit{La flèche $i_f : X' \longrightarrow \bar{X}'$ est un $X$-morphisme (si $X'$ sur $X$ par $f$, $\bar{X}'$ par $\bar{f}$), i.e. $f = \bar{f}\circ i_f$, et c'est une} Hot-\textit{équivalence} (car $\pi_f : \bar{X}' \longrightarrow X'$ est Cat-fibrant, à fibres contractiles, donc une Hot-équivalence).}).

Appliquant ceci à des $f : X \longrightarrow S$, on trouve une construction
\[
\mathrm{Cat}/S \xrightarrow{\Psi} \cat{Cofibsc}\,S,
\]
c'est bel et bien fonctoriel en $X$ variant dans $\mathrm{Cat}/S$. Si
\[
\psi : X \longrightarrow Y
\]
% foot: Version 9 novembre 2001
\printedpage{12}
% head: A. Grothendieck : Dérivateurs, Ch. VI, Édition des Universités de Montpellier II et Paris VII
dans $\mathrm{Cat}/S$, alors $\psi$ induit des
\[
\psi/s : X/s \longrightarrow Y/s\ ,
\]
pour $s \in S$, et ceci de façon compatible avec les
\msp{page 19}
foncteurs `cochangement de base' associés aux $u : s \longrightarrow t$ dans $S$
\[
\begin{tikzcd}[column sep=large, row sep=large]
X/s \arrow[r, "u_*^X"] \arrow[d, "\psi/s"'] & X/t \arrow[d, "\psi/t"] \\
Y/s \arrow[r, "u_*^Y"] & Y/t \quad .
\end{tikzcd}
\]
Cela définit donc un morphisme de catégories scindées sur $S$
\[
\Psi(X) \longrightarrow \Psi(Y),
\]
c'est lui!

On a un morphisme fonctoriel en $X$
\[
X \xrightarrow{i_X} \Psi(X)
\]
dans $\mathrm{Cat}/X$, et c'est une Hot-équivalence, donc $\in W$. Sur les fibres, le morphisme $i_X$ induit
\[
X_s \xrightarrow{i_{X,s}} X/s,
\]
l'inclusion canonique.

\msp{page 20}
\textbf{Corollaire 1.} \textit{Si $X \xrightarrow{f} Y$ dans} $\mathrm{Cat}/S$ \textit{est dans $W$, alors $\Psi(X) \xrightarrow{\Psi(f)} \Psi(Y)$ aussi.} Mais on se gardera de croire que c'est dans $W_S$ au sens de $\cat{Cofib}\,S$ ou $\cat{Cofibsc}\,S$.

Car sur les fibres de $\Psi(X)$, $\Psi(Y)$ on trouve les flèches
\[
f/s : X/s \longrightarrow Y/s
\]
induits par $f$, qui n'ont aucune raison d'être dans $W$. Qu'elles le soient signifie que $f : X \longrightarrow Y$ est dans $W$ `localement au dessus de S' (\footnote[11]{cf. \circled{1} page 15.}), ce qui (pour les $W$ qui nous intéressent) est beaucoup plus fort que d'être dans $W$. (Si $Y = S$, cela signifie en fait que $p : X \longrightarrow S$ est $W$-asphérique, beaucoup plus fort que $p \in W$.)

Donc pour prouver que
\[
\cat{Cofib}\,S \longrightarrow \mathrm{Cat}/S \qquad\quad \text{ou} \qquad\quad \cat{Cofibsc}\,S \longrightarrow \mathrm{Cat}/S
\]
% foot: Version 9 novembre 2001
\printedpage{13}
% head: A. Grothendieck : Dérivateurs, Ch. VI, Édition des Universités de Montpellier II et Paris VII
induisent des équivalences pour les catégories localisées, le foncteur $\Psi$ directement ne sera pas utilisable. Cependant :

\textbf{Corollaire 2.} \textit{Supposons $X$, $Y$ dans} $\cat{Prop}\,S$ \textit{ou dans} $\cat{Fibhot}\,S$\textit{, et considérons}
\[
\Psi(f) : \Psi(X) \longrightarrow \Psi(Y),
\]
\textit{pour que ce soit une $W$-équivalence dans} $\cat{Cofibsc}\,S$\textit{,}
\msp{page 21}
\textit{i.e. une $W$-équivalence fibre par fibre, il faut et il suffit que $f$ soit une $W$-équivalence fibre par fibre} (et cela doit impliquer (\footnote[12]{cf. \circled{1} page 15.}) que $f$ est une $W$-équivalence).

En effet, dans le diagramme commutatif
\[
\begin{tikzcd}[column sep=large, row sep=huge]
X_s \arrow[r, "f_s"] \arrow[d] & Y_s \arrow[d] \\
X/s \arrow[r, "f/s"] & Y/s
\end{tikzcd}
\]
les flèches verticales sont dans $W$, car $X$, $Y$ \texttt{[sont]} propres sur $S$ (ou $X$, $Y$ dans $\cat{Fibhot}\,S$), donc $f_s$ est dans $W$ si et seulement si $f/s$ l'est.

\texttt{[Le foncteur]} $\Phi$ se définit de façon duale à $\Psi$, ce qui revient à échanger les rôles de $\sigma_X$, $\beta_X$ dans le diagramme page 17.

(\footnote[13]{\textbf{Corollaire 3.} \textit{Si $X \in \mathit{Ob}\,$}$\cat{Fibhot}\,S$\textit{, alors $\Psi(X)$ est dans} $\cat{Fibsc}_0S$.\par Cela signifie en effet que les $X/s \longrightarrow X/t$ sont dans $W$, or l'hypothèse sur $X$ implique que c'est dans $W_\infty$. OK.})

\textbf{II) Définition de} $C : \mathrm{Cat}/S \longrightarrow \cat{Parf}\,S$

Cette fois, c'est la factorisation des flèches $p : X \longrightarrow S$ dans Cat, en
\[
X \xrightarrow{j_{X/S}} \bar{X} \xrightarrow{\bar{p}} S,
\]
où $i_X$ \texttt{[plutôt} $j_{X/S}$\texttt{]} est une Hot-équivalence (donc dans $W$), et $\bar{p}$ parfait. Cela se construit formellement comme tantôt, en utilisant $\cat{Ch}_\infty(S)$ au lieu de $\cat{Fl}(S)$, et la factorisation
\[
S \xrightarrow{j_S} \cat{Ch}_\infty(S) \xrightarrow{p_S=(\sigma_S,\beta_S)} S \times S
\]
de $\mathrm{diag}_S : S \longrightarrow S \times S$. La situation est plus jolie, car $p_S$ est déjà `parfait', ce qui implique formellement que le $\bar{p}$ qui s'en déduit, pour factoriser $p : X \longrightarrow S$, est lui-même parfait.
% foot: Version 9 novembre 2001


\endgroup
% END UNIT Der-ChVI-345/Der-ChVI-345_p011-015.tex
% BEGIN UNIT Der-ChVI-345/Der-ChVI-345_p016-020.tex
\begingroup\makeatletter%
\NoAutoSpacing%
\newcommand{\printedpage}[1]{\marginpar{\footnotesize p.~#1}}%
\newcommand{\msp}[1]{\marginpar{\footnotesize [#1]}}%
\newcommand{\und}[1]{\underline{\mathrm{#1}}}%
\newcommand{\Cat}{\mathrm{Cat}}%
\newcommand{\circled}[1]{\tikz[baseline=(c.base)]\node[draw,circle,inner sep=1pt](c){#1};}%
\makeatother\let\printedpage\derpp\let\msp\dermsp%
\let\mspin\DERmspin%


\printedpage{14}
% head: A. Grothendieck : Dérivateurs, Ch. VI, Édition des Universités de Montpellier II et Paris VII
% foot: Version 9 novembre 2001

\msp{page 22}
Ici encore, la construction est fonctorielle, on trouve le foncteur
\[ C : \Cat/S \longrightarrow \und{Parf}\,S \]
qu'il faut composer au foncteur en sens inverse
\[ r : \Cat/S \longleftarrow \und{Parf}\,S. \]
Tout d'abord, vérifions que $C$, tout comme $r$, est compatible avec les localiseurs. Soit $f : X \longrightarrow Y$ dans $\Cat/S$, supposons $f$ est dans $W$, prouvons $C(f)$ est dans $W_S$. On a commutativité dans
\[
\begin{tikzcd}[row sep=large, column sep=large]
X \arrow[r, "f"] \arrow[d, "j_X"'] & Y \arrow[d, "j_Y"] \\
CX \arrow[r, "Cf"] & CY
\end{tikzcd}
\]
et $j_X$, $j_Y$ sont dans $W$, donc $f$ est dans $W$ si et seulement si $Cf$ l'est. Mais dans le cas $W = W_\infty$ on sait que cela équivaut à ce que $Cf$ soit dans $W_\infty$ fibre par fibre. Admettant \circled{1} p. 15 pour $W$, il reste à s'assurer d'une réciproque dans \underline{Parf}, i.e.

\circled{2} Soit
\begin{tikzcd}[baseline=(\tikzcdmatrixname-2-2.base), row sep=small, column sep=small]
X \arrow[rr, "f"] \arrow[dr] & & Y \arrow[dl] \\
& S &
\end{tikzcd}
dans $\und{Parf}\,S$ (ou seulement dans $\und{Fibhot}\,S$). Si $f$ est dans $W$, alors les $f_s : X_s \longrightarrow Y_s$ aussi.

\msp{page 23}
Essayons de le prouver, en calquant la démonstration dans le cas $C = C_\infty$ (XI, th.~2). On peut supposer $S$ connexe (\footnote[14]{ne sert à rien}), en utilisant La (i.e. $W \subseteq W_0$) et en admettant le

\textbf{Lemme.} \textit{Soit $f : X \longrightarrow Y$ dans} Cat. \textit{Alors $f \in W$ si et seulement si $\pi_0(f)$ est bijectif, et pour toute composante connexe $X_i$ de $X$, si $Y_i$ est la composante connexe correspondante de $Y$, $f_i : X_i \longrightarrow Y_i$ induit par $f$ est dans $W$.}

Utilisons alors l'existence de
\[ S' \longrightarrow S \]
Hot-fibration avec $S'$ Hot-asphérique (\footnote[15]{\textbf{NB} Il suffit d'avoir une Hot-fibration $S' \longrightarrow S_0 \hookrightarrow S$, où $S_0$ est la composante connexe de $S$ contenant l'élément $s \in \mathrm{Ob}\,S$ envisagé.}). On peut supposer même que $S' \longrightarrow S$ est parfait. Supposons qu'on sache qu'un changement de base par de telles fibrations parfaites transforme $W$-équivalences en itou. (Ça doit être vrai pour les $W_i$, mais on a vu que c'était faux pour l'équivalence $k$-homologique, cf. I p.~87 ff.) Alors $X' \xrightarrow{f'} Y'$ sur $S'$ est
\printedpage{15}
% head: A. Grothendieck : Dérivateurs, Ch. VI, Édition des Universités de Montpellier II et Paris VII
% foot: Version 9 novembre 2001
une $W$-équivalence. Or pour $s' \in S'$, comme $\{s'\} \longrightarrow S'$ est une Hot-équivalence ($S'$ Hot-asphérique), et $X'$, $Y'$ étant Hot-fibrés sur $S$, il s'ensuit que $X'_{s'} \longrightarrow X'$, $Y'_{s'} \longrightarrow Y'$ sont des Hot-équivalences, donc aussi $X'_{s'} \xrightarrow{f'_{s'}} Y'_{s'}$, comme on voit sur le carré commutatif
\[
\begin{tikzcd}[row sep=large, column sep=large]
X'_{s'} \arrow[r, "f'_{s'}"] \arrow[d, "\wr"'] & Y'_{s'} \arrow[d, "\wr"'] \\
X' \arrow[r, "{f'}"] & Y'.
\end{tikzcd}
\]

\msp{page 24}
Comme $S' \longrightarrow S$ est surjectif, cela montre que les $f_s : X_s \longrightarrow Y_s$ sont dans $W$, qed.

Ainsi, (2) sera conséquence de \circled{2.} Soit $p : Y' \longrightarrow Y$ une Hot-fibration (ou même un foncteur parfait, si ça peut aider \dots). Soit $f : X \longrightarrow Y$ une $S$-flèche (dans $\Cat/S$), alors $f' : X' \longrightarrow Y'$ par changement de base. Si $f \in W$, $f' \in W'$
\[
\begin{tikzcd}[row sep=large]
X' \arrow[r, "q"] \arrow[d, "{f'}"'] \arrow[dr, phantom, "\scriptstyle\text{cart.}"] & X \arrow[d, "f"] \\
Y' \arrow[r, "p"] & Y
\end{tikzcd}
\qquad\qquad
f \in W \text{ et } p \text{ parfait} \Longrightarrow f' \in W
\]

En tous cas, il nous faut absolument \circled{2} pour pouvoir conclure que $C$ est compatible avec les localiseurs.

Considérons $rC : \Cat/S \longrightarrow \Cat/S$, on a une flèche
\[
\left\{
\begin{array}{l}
\mathrm{id}_{\Cat/S} \xrightarrow{j} rC \\[1.5ex]
j_X : X \longrightarrow CX \quad \text{la flèche canonique } i_{X/S}\ \text{\texttt{[plutôt} } j_{X/S}\text{\texttt{]}} \text{ (p.~21),}
\end{array}
\right.
\]
et $j_X$ est toujours dans $W$.

Considérons $Cr : \und{Parf}\,S \longrightarrow \und{Parf}\,S$, i.e. appliquons la construction $C$ à un $X$ parfait sur $S$. Alors on a encore
\[
\left\{
\begin{array}{l}
\mathrm{id}_{\und{Parf}\,S} \xrightarrow{j'} Cr \\[1ex]
j'_X = j_{r(X)} : X = r(X) \xrightarrow{j_X} CX.
\end{array}
\right.
\]
Ce morphisme est dans $W$, donc aussi une $W$-équivalence fibre par fibre, puisque source et but sont parfaits, donc des Hot-fibrés sur $S$.

\msp{page 25}
Cela prouve déjà que les foncteurs $r$ et $C$ induisent des équivalences quasi-inverses l'une de l'autre sur les catégories localisées. Par le lemme p.~11, cela montre que pour toute sous-catégorie pleine $\mathcal{F}$ de $\Cat/S$ contenant $\und{Parf}\,S$, l'inclusion de $\mathcal{F}$ dans $\Cat/S$ induit une
\printedpage{16}
% head: A. Grothendieck : Dérivateurs, Ch. VI, Édition des Universités de Montpellier II et Paris VII
% foot: Version 9 novembre 2001
équivalence sur les localisés, en prenant comme localiseur dans $\mathcal{F}$ le localiseur `grossier' induit par $\Cat/S$. Cela s'applique notamment aux sous-catégories pleines
\[ \und{Liss}\,S, \quad \und{Fibhot}\,S, \quad \und{Prop}\,S, \]
mais en prenant dans $\und{Liss}\,S$ et $\und{Prop}\,S$ non pas le localiseur fin `fibre par fibre', comme je l'avais annoncé à tort, mais le localiseur grossier induit par $\Cat/S$. (Pour $\und{Fibhot}\,S$, il n'y a pas de risque d'ambiguïté). Il faudrait que je me convainque que si je prends la $W$-équivalence fine (dans le cas $W = W_\infty$, disons), je ne trouve \textit{pas} les `bons' localisés, i.e. le foncteur des localisés, vers le localisé de $\Cat/S$, n'est pas une équivalence. Ceci en fait suppose que dans $\und{Fib}\,S$ et $\und{Cofib}\,S$ également, c'est l'équivalence \textit{grossière}
\msp{page 26}
qu'il convient de prendre. (Dans $\und{Fib}_0S$ et $\und{Cofib}_0S$, ou $\und{Fibsc}_0S$ et $\und{Cofibsc}_0S$, les deux équivalences sont identiques, puisque ces quatre catégories sont contenues dans $\und{Fibhot}\,S$.) La suite nous montrera s'il en est bien ainsi. Par contre, dans $\und{Fibsc}\,S$ et $\und{Cofibsc}\,S$, il n'y a pas de doute, suivant les principes généraux de construction d'un dérivateur à partir d'un $(\mathcal{M}, W)$, qu'il faut prendre l'équivalence fine, fibre par fibre. À ce sujet, il faut faire attention que les foncteurs verticaux du diagramme p.~8
\[ \und{Fibsc}\,S \longrightarrow \und{Fib}\,S, \quad \und{Cofibsc}\,S \longrightarrow \und{Cofib}\,S, \]
quoique essentiellement surjectifs et fidèles, ne sont pas pleinement fidèles. Ils le seraient, et mêmes des équivalences de catégories, si dans $\und{Fib}\,S$ resp. $\und{Cofib}\,S$ je me bornais aux $S$-foncteurs \textit{cartésiens} resp. \textit{cocartésiens}. Donc on peut interpréter, à équivalence près, $\und{Fibsc}\,S$ et $\und{Cofibsc}\,S$ comme $\und{Fib}\,S$ et $\und{Cofib}\,S$, où on aurait restreint les morphismes aux seuls $S$-foncteurs \textit{cartésiens} resp. \textit{cocartésiens.} Il semblerait donc qu'en restreignant ainsi les flèches, mais en restreignant de façon plus draconienne encore le localiseur, on trouve des catégories localisées équivalentes. % [sic] et mêmes des équivalences -> et même des équivalences, adverb "même" printed with plural s

\msp{page 27}
C'est assez étrange - on va voir si c'est bien vrai!

\subsection*{III) Construction de $\Phi_0, \Psi_0 :$ \begin{tikzcd}[ampersand replacement=\&, baseline=(\tikzcdmatrixname-2-1.base), row sep=small, column sep=normal]
\& \und{Fibsc}_0S \\
\Cat/S \arrow[ur] \arrow[dr] \& \\
\& \und{Cofibsc}_0S.
\end{tikzcd}}

On doit avoir
\[ \Phi_0(X) = \Phi r C(X) \qquad\qquad\qquad \Psi_0(X) = \Psi r C(X), \]
où
\[ r : \und{Parf}\,S \longrightarrow \Cat/S \]
est l'inclusion canonique. En tant que foncteurs à valeurs dans \underline{Fibsc} resp. \underline{Cofibsc}, cela définit déjà $\Phi_0$, $\Psi_0$ par
\[ \Phi_0 = \Phi r C \qquad\qquad\qquad \Psi_0 = \Psi r C, \]
\printedpage{17}
% head: A. Grothendieck : Dérivateurs, Ch. VI, Édition des Universités de Montpellier II et Paris VII
% foot: Version 9 novembre 2001
mais il faut voir que $\Phi_0$, $\Psi_0$ prennent leurs valeurs dans les sous-catégories pleines à indice 0. Cela équivaut donc au Lemme suivant, appliqué à $Y = rC(X)$:

\textbf{Lemme.} \textit{Soit $Y$ parfait sur $S$ (ou seulement un \textup{Hot}-fibré sur $S$). Alors $\Phi(Y)$ est dans $\und{Fibsc}_0S$, $\Psi(Y)$ dans $\und{Cofibsc}_0S$.}

Il faut prouver (pour $\Psi(Y)$) que si $s \longrightarrow t$ est une flèche dans $S$, alors
\[ \Psi(Y)_s \longrightarrow \Psi(Y)_t \]
est une $W$-équivalence. Or cette flèche s'identifie à
\[ Y/s \longrightarrow Y/t \]
et on sait bien que c'est une Hot-équi\msp{page 28}valence, par définition des Hot-fibrés, donc c'est une $W$-équivalence.

On a donc les foncteurs
\[
\begin{tikzcd}[column sep=large]
\und{Fibsc}_0S \arrow[r, shift left, "p"] & \Cat/S, \arrow[l, shift left, "\Phi_0"]
\end{tikzcd}
\]
il faut voir qu'ils satisfont aux conditions générales du lemme p.~11 (et ce sera pareil pour $q$, $\Psi_0$). On sait déjà que $p$ est compatible avec les localiseurs. Prouvons le pour $\Psi_0$. Soit
\[ f : X \longrightarrow Y \]
dans $\mathrm{Cat/S}$ une $W$-équivalence, alors $rC(f)$ l'est aussi. Il suffit de prouver ceci : % [sic] Cat/S (upright S) -> Cat/$S$, the S printed upright unlike elsewhere

Soit
\[ q : X' \longrightarrow Y' \]
un morphisme dans $\und{Fibhot}\,S$ (par exemple un morphisme de $\und{Parf}\,S$), qui soit $W$-équivalence, donc \texttt{[qui induit]} des $W$-équivalences sur les fibres. Alors $\Psi(q) : \Psi(X') \longrightarrow \Psi(Y')$ est une $W$-équivalence fibre par fibre, i.e.
\[ X'/s \longrightarrow Y'/s \]
une $W$-équivalence pour tout $s \in S$ (a été vu dans le corollaire 2, page 20).

\msp{page 29}
Considérons
\[ p\Phi_0 : \Cat/S \longrightarrow \Cat/S, \]
ça transforme donc $X$ en $\Phi C(X)$, regardé comme objet de $\Cat/S$ sans plus. On a un morphisme fonctoriel
\[ X \xrightarrow{j_X} C(X) \qquad\qquad\qquad j : \mathrm{id}_{\Cat/S} \longrightarrow C, \]
\printedpage{18}
% head: A. Grothendieck : Dérivateurs, Ch. VI, Édition des Universités de Montpellier II et Paris VII
% foot: Version 9 novembre 2001
d'où un morphisme fonctoriel
\[ \Phi(X) \xrightarrow{\Phi(j_X)} \Phi C(X) \qquad\qquad\qquad \Phi * j : \Phi \longrightarrow \Phi_0 \]
et également un morphisme fonctoriel
\[ X \xrightarrow{i_X} \Phi(X), \]
d'où en composant
\[
\begin{tikzcd}
X \arrow[r, "i_X"] \arrow[rr, bend right, "\alpha_X"'] & \Phi(X) \arrow[r, "\Phi(j_X)"] & \Phi C(X),
\end{tikzcd}
\]
i.e.
\[ \alpha : \mathrm{id}_{\Cat/S} \longrightarrow p\Phi_0. \]
Pour $X$ fixé dans $\Cat/S$, on sait que $j_X$ est dans $W$, donc aussi $\Phi(j_X)$ (quoique non fibre par fibre), cf. cor.~1 page 20. On sait aussi que $i_X$ est dans $W$, donc aussi le composé $\alpha_X$.

Considérons
\[ \Phi_0 p : \und{Fibsc}_0\,S \longrightarrow \und{Fibsc}_0\,S\ \ . \]
On part de $X$ dans $\und{Fibsc}_0\,S$, on oublie le scindage et le fait même qu'il soit Cat-fibrant
\msp{page 30}
(on a appliqué $p$), on lui applique $\Phi_0$, i.e. on regarde $\Phi CX$ comme objet de $\und{Fibsc}\,S$ (et même de $\und{Fibsc}_0\,S$). On veut définir encore
\[ X \xrightarrow{\beta_X} \Phi_0 p(X) = \Phi C(X) \qquad\quad \text{\textbf{dans} } \und{Fibsc}_0\text{\textbf{,}} \] % ?? uncertain: weight of "Fibsc_0" after bold "dans" -- regular or bold
fonctoriel en $X$. Or on a comme tantôt
\[ X \xrightarrow{j_X} C(X), \]
d'où en appliquant $\Phi$
\[ \Phi(X) \xrightarrow{\Phi(j_X)} \Phi C(X), \]
mais regardé cette fois comme flèche dans $\und{Fibsc}_0\,S$ (N.B. on a bien $\Phi(X)$ dans $\und{Fibsc}_0\,S$, car $X$, provenant de $\und{Fibsc}_0\,S$, est dans $\und{Fibhot}\,S$, et on applique la remarque `Cor.~3' page 21.) D'autre part, on a une flèche dans $\Cat/S$
\[ X \xrightarrow{i_X} \Phi(X), \]
est-elle une flèche dans $\und{Fibsc}_0(S)$, i.e. est-elle cartésienne? On voit tout de suite que non, que l'on a seulement, pour une flèche $u : s \longrightarrow t$ dans $S$, un homomorphisme de commutation (pas un isomorphisme) dans


\endgroup
% END UNIT Der-ChVI-345/Der-ChVI-345_p016-020.tex
% BEGIN UNIT Der-ChVI-345/Der-ChVI-345_p021-025.tex
\begingroup\makeatletter%
\newcommand{\printedpage}[1]{\marginpar{\footnotesize p.~#1}}%
\newcommand{\msp}[1]{\marginpar{\footnotesize [#1]}}%
\newcommand{\uu}[1]{\underline{\mathrm{#1}}}%
\newcommand{\circled}[1]{\tikz[baseline=(c.base)]\node[draw,circle,inner sep=1pt](c){#1};}%
\newcommand{\lline}{\mathrel{\relbar\joinrel\relbar\joinrel\relbar}}%
\makeatother\let\printedpage\derpp\let\msp\dermsp%
\let\mspin\DERmspin%


\printedpage{19}
% head: A. Grothendieck : Dérivateurs, Ch. VI, Édition des Universités de Montpellier II et Paris VII
% foot: Version 9 novembre 2001

\[
\begin{tikzcd}[row sep=large, column sep=large]
X_s \arrow[r, "u_*^X"] \arrow[d, "i_s"] & X_t \arrow[d, "i_t"] \\
X/s \arrow[r, "(\mathrm{id}_X)/u", "= u_*^{\Phi(X)}"'] \arrow[ur, bend left=30, shorten <=28pt, shorten >=12pt, "\alpha"] & X/t
\end{tikzcd}
\]

\msp{page 31}
Donc il me faut remballer! On peut considérer $i_X$ au mieux comme une flèche de $\uu{Fib}_0\, S$, et en prenant les composés on trouve une flèche dans $\uu{Fib}_0\, S$ :
\[
\begin{tikzcd}
X \arrow[r, "i_X"] \arrow[rr, bend right=25, "\beta_X(=\alpha_X)"'] & \Phi(X) \arrow[r, "\Phi(j_X)"] & \Phi C(X),
\end{tikzcd}
\]
i.e. (regardant maintenant $\Phi_0$ comme foncteur à valeurs dans $\uu{Fib}_0\, S$) un homomorphisme de foncteurs
\[
\beta : \mathrm{id}_{\uu{Fib}_0} S \longrightarrow \Phi_0 p_0, % [sic] \mathrm{id}_{\uu{Fib}_0} S -> \mathrm{id}_{\uu{Fib}_0 S}, the S is printed on the baseline, not in the subscript
\]
où $p_0 : \uu{Fib}_0 S \hookrightarrow \mathrm{Cat}/S$ est l'inclusion canonique. Ce $\beta$ est dans $W$ argument par argument, et de même pour $\alpha$ page 29, interprété cette fois comme
\[
\mathrm{id}_{\mathrm{Cat}/S} \longrightarrow p_0\Phi_0.
\]
Ce qu'on a prouvé finalement est ceci:

\textbf{Théorème-remballage.} \textit{Soit $\mathcal{F}$ une sous-catégorie pleine de $\mathrm{Cat}/S$ qui contient une des trois sous-catégories pleines $\uu{Fib}_0 S$, $\uu{Parf}\, S$, $\uu{Cofib}_0 S$, et munissant $\mathcal{F}$ du localiseur $W$ `grossier' (induit par celui de $\mathrm{Cat}/S$). Alors l'inclusion}
\[
\mathcal{F} \longrightarrow \mathrm{Cat}/S
\]
\textit{induit une équivalence sur les localiseurs.}

\msp{page 32}
Par contre, les catégories localisées des quatre catégories de fibrés scindés
\[
\uu{Fibsc}\, S, \quad \uu{Cofibsc}\, S, \quad \uu{Fibsc}_0 S, \quad \uu{Cofibsc}_0 S
\]
échappent à nos arguments. \textit{Et pour cause!} Il est totalement ouvert, sur $S$ connexe disons (p. ex. $S = \boldsymbol{\Delta}^1$), que pour un type d'homotopie relatif $\xi$ sur $S$, les types induits $\xi_s$ ($s \in S$) forment un système local sur $S$, à valeurs dans $\mathrm{Hot}_W$, comme ce serait le cas si le théorème annoncé au début était vrai. Au contraire, \textit{toute} flèche de $\mathrm{Hot}_W$ doit provenir d'un objet de $\mathrm{Hot}_W(\boldsymbol{\Delta}^1)$. Les localisés des 9 autres catégories du diagramme p. 8, et de toutes celles dont il est question dans le théorème qu'on vient de prouver, doivent correspondre `visiblement' aux sous-catégories pleines de $\mathrm{HOT}_W(S)$ formées des HOT-coefficients localement constants sur $S$.

Aussi, je me rends compte que c'était absurde de m'attendre à ce que $\uu{Fibsc}\, S$ et $\uu{Cofibsc}\, S$ donnent des catégories localisées canoniquement équivalentes. Car les éléments $\xi$ de la première
\printedpage{20}
% head: A. Grothendieck : Dérivateurs, Ch. VI, Édition des Universités de Montpellier II et Paris VII
% foot: Version 9 novembre 2001
\msp{page 33}
des deux catégories donnent naissance chacun à un foncteur contravariant sur $S$
\[
\bar{\xi} : S^o \longrightarrow \mathrm{Hot}_W,
\]
ceux $\eta$ de la seconde (qui n'est autre que $(\uu{Fibsc}(S^o))^o$) donnent naissance chacun à un foncteur \textit{covariant}
\[
\bar{\eta} : S \longrightarrow \mathrm{Hot}_W,
\]
ce qui n'est pas pareil! Ce n'est `pareil' que si ce sont des coefficients locaux, ce qui signifie exactement que $\xi$ est dans $\uu{Fibsc}_0 S$ resp. $\eta$ dans $\uu{Cofibsc}_0 S$. On s'attend donc bien \circled{$1^o$} que les flèches
\[
\uu{Fibsc}_0 S \longrightarrow \uu{Fib}_0 S, \quad \uu{Cofibsc}_0 S \longrightarrow \uu{Cofib}_0 S
\]
induisent des équivalences sur les localisés (chose que j'ai été incapable de prouver tantôt) (\setcounter{footnote}{15}\footnote{Les quasi-inverses étant bel et bien donnés par $\Phi_0$, $\Psi_0$ - car [?] par quoi d'autre donc ???}). Par contre, il est pire que déraisonnable de s'attendre à ce qu'il en soit de même pour les foncteurs
\[
\uu{Fibsc}_0 S \longrightarrow \uu{Fibsc}\, S, \quad \uu{Cofibsc}_0 S \longrightarrow \uu{Cofibsc}\, S,
\]
ça devrait donner par localisation les foncteurs pleinement fidèles d'inclusion, des catégories d'objets localement constants dans $\mathrm{HOT}_W(S)$ resp. dans $(\mathrm{HOT}_W(S^o))^o$. Et de même pour
\msp{page 34}
les foncteurs induits sur les localisés par
\[
(*) \qquad \uu{Fibsc}\, S \longrightarrow \uu{Fib}\, S, \quad \uu{Cofibsc}\, S \longrightarrow \uu{Cofib}\, S.
\]
Comme les localisés des catégories but de ces deux foncteurs ont tout l'air de s'identifier aux sous-catégories précédentes d'objets localement constants, il semblerait donc que : \circled{$2^o$} Les foncteurs précédents définissent des rétractions de la catégorie $\mathrm{HOT}_W(S)$ resp. $\mathrm{HOT}_W(S^o)^o$ sur la sous-catégorie des objets localement constants sur $S$, les foncteurs localisés des foncteurs d'inclusion
\[
\uu{Fibsc}_0 S \longrightarrow \uu{Fibsc}\, S, \quad \uu{Cofibsc}_0 S \longrightarrow \uu{Cofibsc}\, S
\]
étant pleinement fidèles et ayant comme images essentielles les sous-catégories en question.

Cette rétraction m'apparait comme une chose tout à fait mystérieuse. C'est lié au fait que pour le dérivateur très spécial $\mathbf{D} = \mathrm{HOT}$, quand on a un $\xi \in \mathbf{D}(S)$, on en déduit (à $\mathrm{HOT}/S$-équivalence près) un type de diagramme $X$ au dessus de $S$,
\msp{page 35}
qu'on peut regarder alors (via la construction $C$) comme définissant un type d'homotopie relatif \textit{localement constant} sur $S$.

\printedpage{21}
% head: A. Grothendieck : Dérivateurs, Ch. VI, Édition des Universités de Montpellier II et Paris VII
% foot: Version 9 novembre 2001

Donc pour élucider la situation, i.e. les relations entre les 13 catégories localisées provenant du diagramme de la page 8, il reste à regarder si les assertions hypothétiques que je viens de faire sont bien valables, savoir

\begin{enumerate}[leftmargin=*]
\item[\circled{$1^o$}] Le foncteur déduit de
\[
\uu{Fibsc}_0 S \longrightarrow \uu{Fib}_0 S
\]
par passage aux localisés est une équivalence de catégories (et on peut dans cette question remplacer $\uu{Fib}_0 S$ par n'importe quelle autre sous-catégorie pleine de $\mathrm{Cat}/S$ la contenant - cela revient au même), et

\item[\circled{$2^o$}] Le foncteur déduit de
\[
\uu{Fibsc}_0 S \longrightarrow \uu{Fibsc}\, S
\]
par localisation est pleinement fidèle. (Que son image soit exactement formée des objets localement constants du premier localisé, interprété comme $\mathrm{HOT}_W(S)$, est tautologique.)
\end{enumerate}

Il n'y a plus lieu de se soucier du côté
\msp{page 36}
cofibrés, puisque tout ce qu'on pourra en dire revient à ce que je viens d'expliciter, pour $X^o$ au lieu de $X$.

Dans la direction de \circled{$1^o$}, nos arguments précédents à coups de $\Phi_0$, $p$ nous donnent déjà un renseignement utile. Mais il sera plus commode à présent de laisser tomber $\uu{Fib}_0 S$ ou plutôt, de le remplacer par $\uu{Fib}\, S$ (ce qui est licite, on l'a dit dans \circled{$1^o$}), de sorte qu'on a maintenant un diagramme commutatif (traits pleins)
\[
\begin{tikzcd}
\uu{Fibsc}\, S \arrow[dr, "j"'] & & \uu{Fibsc}_0 S \arrow[ll, hook', "i"'] \arrow[dl, "k"'] \\
& \uu{Fib}\, S, \arrow[ur, dotted, bend right=40, "\Phi_0=\psi"'] &
\end{tikzcd}
\]
où $i$ est une inclusion (strictement) pleinement fidèle, $j$ (donc aussi $k$) fidèle. Enfin, on a un foncteur canonique $\psi$ ($= \Phi_0$) (en pointillés). Les localiseurs sur la ligne du haut sont les localiseur fins (pour $\uu{Fibsc}_0 S$ c'est pareil que la $W$-équivalence grossière), sur $\uu{Fib}\, S$ % [sic] les localiseur fins -> les localiseurs fins, missing plural s
par contre on prend la $W$-équivalence grossière. (On sait que le localisé est le même, à équivalence près, que si on se restreint à la sous-catégorie pleine $\uu{Fib}_0 S$, ce qui est assez remarquable. % [sic] "(On sait ... remarquable." -> "(On sait ... remarquable.)", the opening parenthesis is never closed

\msp{page 37}
Tous les quatre foncteurs $i, j, k, \Phi_0$ sont compatibles aux localiseurs, donc passent aux localisés
\printedpage{22}
% head: A. Grothendieck : Dérivateurs, Ch. VI, Édition des Universités de Montpellier II et Paris VII
% foot: Version 9 novembre 2001
\[
\begin{tikzcd}
\mathrm{HOT}_W(S) \arrow[dr, "\bar{j}"'] & & \mathrm{HOT}^{?}_W(S) \arrow[ll, "\bar{\imath}"'] \arrow[dl, "\bar{k}"'] \\
& \mathrm{HOT}^{\mathrm{lc}}_W(S), \arrow[ur, bend right=40, "\bar{\psi}"'] &
\end{tikzcd}
\]
où l'exposant ? dans la notation signifie que je ne sais interpréter cette catégorie, sauf si je prouve que $\bar{k}$ est une équivalence, et l'exposant lc dans la catégorie du bas dois suggerer l'intuition qu'il s'agit en tous cas d'une sorte de `types d'homotopie relatifs' (dans un sens modifié) \textit{localement constants} sur $S$. Le problème est de comprendre ce diagramme. La conjecture naturelle est % [sic] dois suggerer -> doit suggérer, verb agreement and missing accent

\begin{enumerate}[leftmargin=*]
\item[$1^o$)] $\bar{k}$ est une équivalence. (Ce qui implique que $\bar{\psi}$ est l'équivalence quasi-inverse, que $\bar{k}$, $\bar{\psi}$ permettent essentiellement d'identifier $\mathrm{HOT}^{?}_W$ et $\mathrm{HOT}^{\mathrm{lc}}_W$, de sorte que $\bar{\imath}$ et $\bar{\jmath}$ deviennent des foncteurs entre $\mathrm{HOT}_W$ et $\mathrm{HOT}^{\mathrm{lc}}_W$.)

\item[$2^o$)] $\bar{\imath}$ est pleinement fidèle (NB l'image essentielle est bien comprise de toutes façons.)
\end{enumerate}

\msp{page 38}
En direction de $1^o$), je sais pour l'instant que l'on a
\[
\bar{k}\bar{\psi} \simeq \mathrm{id}_{\mathrm{HOT}^{\mathrm{lc}}_W(S)},
\]
donc $\bar{k}$ est essentiellement surjectif (sur les objets), et induit des \textit{surjections} sur les $\mathrm{Hom}(-,-)$, tandis que $\bar{\psi}$ est injectif sur les classes d'isomorphismes d'objets, et sur les $\mathrm{Hom}(-,-)$ (i.e. \textit{fidèle}).

Si les relations entre $\mathrm{HOT}^{?}_W$ et $\mathrm{HOT}^{\mathrm{lc}}_W$ s'avèrent imbitables (p. ex. $\bar{\psi}$ pas même pleinement fidèle, i.e. $\bar{k}$ pas même fidèle), alors la catégorie $\mathrm{HOT}^{?}_W$ mérite d'être larguée dans la corbeille, et on ne gardera que
\[
\begin{tikzcd}[row sep=large]
\mathrm{HOT}_W(S) \arrow[d, shift right=1.5, "\bar{j}"'] \\
\mathrm{HOT}^{\mathrm{lc}}_W(S), \arrow[u, shift right=1.5, "\bar{\psi}'"']
\end{tikzcd}
\]
où on définit $\bar{\psi}'$ comme
\[
\bar{\psi}' = \bar{\imath}\bar{\psi} = \overline{i\psi}.
\]
Et la question (plus modeste) devient à présent :

\begin{enumerate}[leftmargin=*]
\item[$3^o$)] $\bar{\psi}'$ est-il pleinement fidèle? (L'image \textit{essentielle} est la même que celle de $\uu{Fib}_0$ ($\approx$ $\uu{Fib}$), elle est formée sauf erreur des objets localement constants - à vérifier au besoin \dots)
\end{enumerate}

\printedpage{23}
% head: A. Grothendieck : Dérivateurs, Ch. VI, Édition des Universités de Montpellier II et Paris VII
% foot: Version 9 novembre 2001

\msp{page 39}
Il est difficile d'imaginer que $\bar{k}$, ou ce qui revient au même $\bar{\psi}$, puisse être une équivalence, sans qu'il existe au moins une `$W_S$-homotopie'
\[
\mathrm{id}_{\uu{Fibsc}_0 S} \ \leftrightarrow\ \leftrightarrow^{h}\!\leftrightarrow\ \leftrightarrow\ \psi k
\]
dans la catégorie des endofoncteurs de $\uu{Fibsc}_0 S$. À vrai dire, je me rends compte maintenant que la description de $\bar{\psi}$ à partir de $\psi : X \longmapsto \Phi_0 X = \Phi C X$ est inutilement compliquée, car on avait fait le détour par $C(X)$ uniquement pour être sûr de tomber dans $\uu{Fibsc}_0 S$, pas seulement dans $\uu{Fibsc}\, S$. Mais si on laisse tomber $\uu{Fibsc}_0 S$, on s'en fout, et le foncteur $\Phi : X \longmapsto \Phi(X)$ est beaucoup plus naturel. Et si on ne laisse pas tomber, il n'y a pas d'inconvénient de remplacer la source $\uu{Fib}\, S$ par $\uu{Fib}_0 S$ (qui a même localisé, à équivalence près), et sur $\uu{Fib}_0 S$ à nouveau il n'y a pas lieu de s'encombrer de $CX$, pour le plaisir de lui appliquer $\Phi$, car $\Phi(X)$ lui-même est déjà dans $\uu{Fibsc}_0 S$.

Donc, la question est celle-ci : Si $X$ varie
\msp{page 40}
% layout: the arrows nearrow, uparrow, nwarrow are printed under each "Fibsc_0 S", pointing up to it; the arrow labelled c from X to Phi(X) is a long line with its arrowhead printed in the middle
\[
\begin{array}{l}
\text{dans } \underset{\nearrow\uparrow\nwarrow}{\uu{Fibsc}_0 S}, \quad \text{y-a-t'il une homotopie dans } \underset{\nearrow\uparrow\nwarrow}{\uu{Fibsc}_0 S}, \quad \text{`naturelle', i.e. fonctorielle en } X, \\[2ex]
\qquad\qquad X \xrightarrow{\qquad c \qquad} \Phi(X) = \ \text{Catégorie fibrée scindée de fibres les } s\backslash X
\end{array}
\]
Mais un chemin qui soit formé de flèches
\[
X = X_0 \lline X_1 \lline X_2 \cdots \lline X_n = \Phi(X)
\]
qui soient toutes cartésiennes (et de plus, bien sûr, dans $W$!). Cela signifie donc qu'on aurait des catégories fibrées scindées sur $S$, $X_1$, \dots, $X_{n-1}$, dépendant \textit{fonctoriellement} de $X$, qui relient
\[
X = X_0 \qquad \text{à} \qquad \Phi(X) = X_n
\]
i.e.
\[
X_s \qquad \text{à} \qquad s\backslash X \qquad \text{???}
\]
Je ne vois rien à l'horizon.


\endgroup
% END UNIT Der-ChVI-345/Der-ChVI-345_p021-025.tex
% BEGIN UNIT Der-ChVI-345/Der-ChVI-345_p026-030.tex
\begingroup\makeatletter%
\newcommand{\printedpage}[1]{\marginpar{\footnotesize p.~#1}}%
\newcommand{\msp}[1]{\marginpar{\footnotesize [#1]}}%
\makeatother\let\printedpage\derpp\let\msp\dermsp%
\let\mspin\DERmspin%


\printedpage{24}
% head: A. Grothendieck : Dérivateurs, Ch. VI, Édition des Universités de Montpellier II et Paris VII

\msp{page 41}

\begin{flushright}
15.11.90
\end{flushright}

Je crois qu j'ai enfin compris, cette nuit, la situation exacte. Il y a lieu, dans Cat$/S$ et dans les autres catégories du diagramme p. 8 (s'envoyant toutes dans Cat$/S$) d'envisager quatre sortes de ‘$W$-équivalences’. % [sic] qu j'ai -> que j'ai, letter missing

\begin{itemize}
\item[a)] La \textit{$W$-équivalence ‘grossière’} : \begin{tikzcd}[baseline=(s.base), row sep=small, column sep=small]
f : X \arrow[rr] \arrow[dr] & & Y \arrow[dl] \\
& |[alias=s]| S &
\end{tikzcd} dans Cat$/S$ est une $W$-équivalence grossière si en tant que flèche de Cat, c'est une $W$-équivalence. (\footnote[17]{stable par changement de base à la fois lisse et propre, plus généralement si $S' \longrightarrow S$ est une Hot-fibration.}).

\item[b)] La \textit{$W$-équivalence ‘par fibres’} (sur $S$) : $\forall\ s \in S$, $f_s : X_s \longrightarrow Y_s$ est une $W$-équivalence. (\footnote[18]{stable par tout chagement de base $S' \longrightarrow S$}). % [sic] chagement -> changement, letter missing

\item[c)] La \textit{$W$-équivalence ‘locale sur $S$’} : $\forall\ s \in S$, $f/s : X/s \longrightarrow Y/s$ est une $W$-équivalence. Signifie que c'est une $W$-équivalence, et le reste par tout changement de base $S' \longrightarrow S$ qui est un isomorphisme local. (\footnote[19]{devrait être stable par changement de base \textit{lisse} $S' \longrightarrow S$}).

\item[c$'$)] La \textit{$W$-équivalence ‘colocale sur $S$’} : $\forall\ s \in S$, $s\backslash f : s\backslash X \longrightarrow s\backslash Y$ est une $W$-équivalence. Signifie que c'est une $W$-équivalence, et le reste par tout changement de base $S' \longrightarrow S$ qui est un co-isomorphisme local. (\footnote[20]{devrait être stable par changement de base \textit{propre} sur $S$}).
\end{itemize}

Je vais résumer dans un énoncé récapitulatif les relations entre ces quatre notions de $W$-équivalence, en pensant au cas $W = W_\infty$. Une fois la chose faite, je me préoccuperai de dégager les conditions minimales sur un $W$ général qui assurent que c'est vrai pour $W$.

\textbf{Proposition 1.} \textit{Soit} \begin{tikzcd}[baseline=(s.base), row sep=small, column sep=small]
f : X \arrow[rr] \arrow[dr] & & Y \arrow[dl] \\
& |[alias=s]| S &
\end{tikzcd} \textit{dans} Cat$/S$.

\begin{itemize}
\item[1)] \textit{La $W$-équivalence locale sur $S$, et la $W$-équivalence colocale sur $S$, impliquent la $W$-équivalence grossière :}
\[
\begin{tikzcd}
\text{c} \arrow[dr, Rightarrow] & & \text{c}' \arrow[dl, Rightarrow] \\
& \text{a} &
\end{tikzcd}
\]

\printedpage{25}
% head: A. Grothendieck : Dérivateurs, Ch. VI, Édition des Universités de Montpellier II et Paris VII

\item[2)] \textit{Si $X$, $Y$ sont propres sur $S$, ou $X$, $Y$ lisses sur $S$, ou $X$, $Y$ des} Hot-\textit{fibrés sur $S$, alors la $W$-équivalence par fibres implique que c'est une $W$-équivalence} universelle \textit{sur $S$ ($W$-équivalence, et le reste par} tout \textit{changement de base $S' \longrightarrow S$), donc dans ces trois cas, on a}
\[
\begin{tikzcd}
& \text{b} \arrow[dl, Rightarrow] \arrow[dr, Rightarrow] & \\
\text{c} \arrow[dr, Rightarrow] & & \text{c}' \arrow[dl, Rightarrow] \\
& \text{a} &
\end{tikzcd}
\]
(\footnote[21]{Il suffit $X$, $Y$ $W$-propres sur $S$, ou $W$-lisses sur $S$, ou $W$-fibrés sur $S$.}).

Mais on peut préciser chacun de ces trois cas :

\item[(2 fib)] \textit{Si $X$, $Y$ sont} Hot-\textit{fibrés sur $S$, alors la $W$-équivalence grossière de $f$ implique que c'est une $W$-équivalence par fibres, donc les quatre notions de $W$-équivalence sont équivalentes pour $f$.}
\[
\text{a} \Longleftrightarrow \text{c} \Longleftrightarrow \text{c}' \Longleftrightarrow \text{b}
\]
(\footnote[22]{a $\Longrightarrow$ b seulement moyennant L5, voir ci-dessous p. 43.})

\item[(2 lisse)] \msp{page 42}\textit{Si $X$, $Y$ sont lisse sur $S$, alors la $W$-équivalence par fibres sur $S$ équivaut à la $W$-équivalence colocale sur $S$ et implique la $W$-équivalence} universelle \textit{sur $S$, donc dans ce cas} % [sic] lisse -> lisses, plural agreement with X, Y
\[
\begin{tikzcd}
\text{b} \arrow[r, Leftrightarrow] \arrow[d, Rightarrow] & \text{c}' \\
\text{c} \arrow[d, Rightarrow] & \\
\text{a} &
\end{tikzcd}
\]

\printedpage{26}
% head: A. Grothendieck : Dérivateurs, Ch. VI, Édition des Universités de Montpellier II et Paris VII
(\footnote[23]{il suffit même $X$, $Y$ $W$-lisses sur $S$.}).

\item[(2 prop)] \textit{Si $X$, $Y$ sont propres sur $S$, alors la $W$-équivalence par fibres sur $S$ équivaut à la $W$-équivalence locale sur $S$ :}
\[
\begin{tikzcd}
\text{b} \arrow[r, Leftrightarrow] \arrow[d, Rightarrow] & \text{c} \\
\text{c}' \arrow[d, Rightarrow] & \\
\text{a} &
\end{tikzcd}
\]
(\footnote[24]{Il suffit même $X$, $Y$ $W$-propres sur $S$.}).
\end{itemize}

\textsc{Démonstration} 1) n'est autre que L4 (\footnote[25]{Non, il y a confusion, L4 est plus faible, cf. p. 44 l'axiome L4 bis.}) et son dual L4$'$ pour un localiseur fondamental $W$, et on a vu que si on suppose L1, L3, alors L4 équivaut à L4$'$, et dans la notion de localiseur fondamental on supposera toujours que c'est vérifié. Également, nous supposons que L3 est vrai même sous la forme plus forte L3 quater (p. 1) : $W_\infty \subseteq W$, i.e. toute Hot-équivalence est une $W$-équivalence.

(2 lisse) Si $X$, $Y$ \texttt{[sont]} lisses sur $S$, alors les $X_s \longrightarrow s\backslash X$, $Y_s \longrightarrow s\backslash Y$ sont des $W_\infty$-équivalences, donc des $W$-équivalences (il suffirait même que $X$, $Y$ soient $W$-lisses$/S$), donc le diagramme commutatif
\[
\begin{tikzcd}
X_s \arrow[r, "f_s"] \arrow[d, "i_s^X"'] & Y_s \arrow[d, "i_s^Y"] \\
s\backslash X \arrow[r, "s\backslash f"] & s\backslash Y
\end{tikzcd}
\]
nous montre que $f_s$ est une $W$-équivalence si et seulement si $s\backslash f$ l'est, d'où b $\Longleftrightarrow$ c$'$. Pour voir que la $W$-équivalence par fibres est même une $W$-équivalence \textit{universelle} sur $S$, on est ramené à montrer que c'est une $W$-équivalence grossière (puisque la lissité \texttt{[est]} stable par changement de base sur $S$), et on gagne puisque c$'$ $\Longrightarrow$ a par 1).

Le cas (2 propre) est dual.

(2 fib) Si $X$, $Y$ sont des Hot-fibrés sur $S$, alors
\printedpage{27}
% head: A. Grothendieck : Dérivateurs, Ch. VI, Édition des Universités de Montpellier II et Paris VII
\msp{page 43}
les inclusions
\[
X_s \longrightarrow s\backslash X, \qquad X_s \longrightarrow X/s,
\]
et de même pour $Y$, sont des $W_\infty$-équivalences, donc des $W$-équivalences (il suffirait même que $X$, $Y$ soient des $W$-fibrations sur $S$ (\footnote[26]{comparer déf. 3, p. 117, dans (I).}), i.e. par définition que les foncteurs précédents soient des $W$-équivalences, et que la même condition soit satisfaite après tout changement de base sur $S$), l'argument dans le cas lisse resp. propre montre déjà que
\[
\begin{tikzcd}
\text{c} & \text{b} \arrow[l, Rightarrow] \arrow[r, Leftrightarrow] \arrow[d, Rightarrow] & \text{c}' \\
& \text{a,} &
\end{tikzcd}
\]
et que la $W$-équivalence par fibres (ou la $W$-équivalence locale sur $S$, resp. colocale sur $S$) implique déjà la $W$-équivalence \textit{universelle} sur $S$. (\footnote[27]{valable en supposant seulement $X$, $Y$ des $W$-fibrés sur $S$.})

La question a $\Longrightarrow$ b est par contre nettement plus délicate, et on a vu que ce n'est nullement le cas pour tout localisateur fondamental - sans doute pas pour celui associé au dérivateur Der($\mathbf{Z}$-Mod), en tous cas pas à celui associé à Der($k$-Mod), quand 2 est inversible dans $k$. (Il suffirait même que $\exists\ n \in \mathbf{N}$, $n \geq 2$, tel que $n$ inversible dans $k$ - le contrexemple typique est $X = S^1 \xrightarrow{\ n\ } S^1 = Y = S$, le revêtement étale d'ordre $n$ du cercle.)

L'implication a $\Longrightarrow$ b est donc une condition très forte sur $W$, je vais la noter L5 (l'ancienne L5, cf. page 2, ne méritant pas de survivre, car c'est une côte mal taillée).

\textbf{L5} \textit{Soit} \begin{tikzcd}[baseline=(s.base), row sep=small, column sep=small]
X \arrow[rr, "f"] \arrow[dr] & & Y \arrow[dl] \\
& |[alias=s]| S &
\end{tikzcd} \textit{dans} Cat\textit{, avec $X$, $Y$ des} Hot-\textit{fibrés sur $S$. Si $f$ est une $W$-équivalence (‘grossière’ par rapport à $S$), alors c'est une $W$-équivalence par fibres} (\footnote[28]{Dans XII (p. 12) on voit que si $W$ satisfait à L1, L3 quater, L4 bis (cf. page suivante), L5, alors \fbox{$W = W_\infty$}. Donc c'est parfaitement illusoire, si on suppose L5 de continuer à traîner un $W$.}).

On a vu qu'à un $\varepsilon$ près (‘lemme’ p. 23), L5 est une conséquence de

\textbf{L6} \textit{Soit $f : X \longrightarrow Y$ une $W$-équivalence, $Y' \longrightarrow Y$ un morphisme parfait (i.e. lisse et propre), alors $f' : X' = X \times_Y Y' \longrightarrow Y'$ est une $W$-équivalence.}

\msp{page 44}
Notons la variante bien naturelle

\textbf{L6 bis} \textit{Comme L6, mais en supposant seulement $Y' \longrightarrow Y$ une} Hot-\textit{fibration.}

\printedpage{28}
% head: A. Grothendieck : Dérivateurs, Ch. VI, Édition des Universités de Montpellier II et Paris VII

Il y a des chances que pratiquement, quand L6 (ou seulement L5) est verifié, L6 bis l'est aussi. % [sic] verifié -> vérifié, missing accent

On va introduire un axiome

\textbf{L7} \ a) \textit{Soit $(f_i : X_i \longrightarrow Y_i)_{i \in I}$ une famille de flèches dans} Cat\textit{, soit $f : \coprod X_i \longrightarrow \coprod Y_i$ la flèche correspondante sur les catégories somme. Pour que $f \in W$, il faut et il suffit que $f_i \in W$ pour tout $i \in I$.}

\hspace*{2em}b) \textit{Soit $f : X \longrightarrow Y$ dans $W$. Si $Y$ est} $0$\textit{-connexe, $X$ aussi.}

Cet axiome équivaut à la conjonction de La ($f \in W \Longrightarrow \pi_0 f$ bijectif) et de la validité du ‘Lemme’ page 23, que je remets ici pour mémoire :

\textbf{Lemme.} \textit{L'axiome L7 précédent sur le localiseur fondamental $W$ équivaut à la condition suivante : $f : X \longrightarrow Y$ est dans $W$ si et seulement si $\pi_0(f)$ est bijectif, et $f$ induit des $W$-équivalences $f_i : X_i \longrightarrow Y_i$ sur les composantes connexes.}

Il me faut revenir encore sur le 1) de la proposition, il n'est pas vrai que ‘c'est’ la condition L4 sur $W$ (que j'admets, ai-je dit), c'est L4 dans le cas particulier où $Y = S$. Réflexion faite, je le vois comme une variante renforcée de L4, qui ne semble pas devoir découler de L4 :

\textbf{L4 bis} \textit{Soit} \begin{tikzcd}[baseline=(s.base), row sep=small, column sep=small]
X \arrow[rr, "f"] \arrow[dr] & & Y \arrow[dl] \\
& |[alias=s]| S &
\end{tikzcd} \textit{dans} Cat\textit{, supposons que $f$ soit une $W$-équivalence locale sur $S$ (i.e. $f/s : X/s \longrightarrow Y/s$ une $W$-équivalence $\forall\ s \in S$), alors $f$ est une $W$-équivalence} (\footnote[29]{équivaut à l'axiome dual, en termes des $s\backslash f$}).

\msp{page 45}
Par la suite, par un ‘localiseur fondamental’, nous entendrons un ensemble
\[
W \subseteq \mathrm{Fl}\,\mathrm{Cat}
\]
satisfaisant (tout au moins) les axioms % [sic] axioms -> axiomes, French spelling
\[
\text{L1, L3 quater, L4 bis, L7,}
\]
i.e.

\begin{itemize}
\item[a)] \textit{Saturation} de $W$ (conditions a) b) c) p. 1)

\item[b)] $W_\infty \subseteq W$ $[\subseteq W_0]$ (\footnote[30]{Je ne suis plus sûr s'il est vraiment utile d'imposer $W_\infty \subseteq W$. Prendre au lieu de ça l'axiome de l'objet final, soit b$'$). Si on veut a) b$'$) c) + L5, alors on aura $W \subseteq W_\infty$. \textbf{NB} a) b) c) + L5 implique $W = W_\infty$, cf. annotation page 43.}) (\footnote[31]{cf. d)}).

\item[c)] Si $f : X \longrightarrow Y$ sur $S$ est localement sur $S$ une $W$-équivalence, alors $f \in W$ (L4 bis)
\end{itemize}


\endgroup
% END UNIT Der-ChVI-345/Der-ChVI-345_p026-030.tex
% BEGIN UNIT Der-ChVI-345/Der-ChVI-345_p031-035.tex
\begingroup\makeatletter%
\newcommand{\printedpage}[1]{\marginpar{\footnotesize p.~#1}}%
\newcommand{\msp}[1]{\marginpar{\footnotesize [#1]}}%
\newcommand{\uc}[1]{\underline{\mathrm{#1}}}%
\newcommand{\HOT}{\mathrm{HOT}}%
\newcommand{\Cat}{\mathrm{Cat}}%
\newcommand{\defeq}{\overset{\text{déf}}{=}}%
\makeatother\let\printedpage\derpp\let\msp\dermsp%
\let\mspin\DERmspin%


\printedpage{29}
% head: A. Grothendieck : Dérivateurs, Ch. VI, Édition des Universités de Montpellier II et Paris VII
% foot: Version 9 novembre 2001

\begin{enumerate}[leftmargin=*]
\item[d)] Pour que $f : X \longrightarrow Y$ soit dans $W$, il faut et il suffit que $\pi_0(f)$ soit bijectif, et que les $f_i : X_i \longrightarrow Y_i$ induits sur les composantes connexes soient dans $W$.
\end{enumerate}

Ces axiomes sont vérifiés par tous les $W$ qui je connais, et qu' on aurait envie de regarder comme localiseurs. Par contre, L5, L6, L6 bis (p. 43, 44) sont des axiomes de nature beaucoup plus spéciale. J'ai envie de dire qu'ils distinguent les localiseurs de nature proprement \textit{`homotopique'} de ceux qui sont de nature \textit{`homologique'} (ou `cohomologique'). Il faudrait que je m'assure qu'ils sont satisfaits pour les $W_i$ ($i$-équivalence d'homotopie) (\footnote[32]{C'est \textit{faux}, cf. annotation ci-dessus.}). % [sic] qui je connais -> que je connais, relative pronoun object requires "que"

Le fait que L4 bis, i.e. c) ci-dessus, implique l'axiome dual, résulte du fait que L4 bis implique ceci

\textbf{Proposition 2.} \textit{Soit $f : X \longrightarrow Y$ dans} $\Cat$\textit{, alors $f \in W$ si et seulement si $f^o : X^o \longrightarrow Y^o$ est dans $W$.}

Se voit par l'argument de I, page 75 (cor. 2) - n'utilise que L4, pas même L1, L3.

\msp{page 46}

Je vais essayer de résumer ce que je crois avoir compris :

\begin{enumerate}[leftmargin=*]
\item[1)] La $W$-équivalence `par fibres' n'a d'intérêt (pour construire des catégories localisées) que dans le cas où on sait qu'elle implique déjà, soit la $W$-équivalence locale, soit la $W$-équivalence colocale (soit même la $W$-équivalence universelle) sur $S$, donc quand on suppose $X$, $Y$ pour le moins tous deux soit $W$-lisses, soit $W$-propres, soit $W$-fibrés sur $S$.

\item[2)] On pourrait appeler la $W$-équivalence locale sur $S$ la \textit{`$W$-équivalence fine à gauche'} sur $S$, et la $W$-équivalence colocale sur $S$ la \textit{`$W$-équivalence fine à droite'} sur $S$. Elles coïncident dans $\uc{Fibhot}\,S$, a fortiori dans $\uc{Parf}\,S$ et dans les quatre catégories (cf. diagramme p. 8) $\uc{Fib}_0S$, $\uc{Fibsc}_0S$, $\uc{Cofib}_0S$, $\uc{Cofibsc}_0S$ (\footnote[33]{On pourrait les appeler les six catégories \textit{internes} du diagramme.}). Par contre, elles ne coïncident pas sur les catégories \textit{bordantes} du diagramme, notamment sur $\uc{Liss}\,S$, $\uc{Prop}\,S$, $\uc{Fib}\,S$, $\uc{Cofib}\,S$, $\uc{Fibsc}\,S$, $\uc{Cofibsc}\,S$. À cet égard, c'est la \textit{$W$-équivalence fine} \textbf{droite} (en termes des $s\backslash X$) qui est intéressante sur le côté (hélas) gauche du diagramme, i.e. pour $\uc{Liss}\,S$, $\uc{Fib}\,S$, $\uc{Fibsc}\,S$, alors que c'est la \textit{$W$-équivalence fine} \textbf{gauche} (en termes des $X/s$) qui est pertinente du côté droit, pour $\uc{Prop}\,S$, $\uc{Cofib}\,S$, $\uc{Cofibsc}\,S$. On n'aura jamais à regarder simultanément les deux équivalences fines sur une même catégorie du diagramme, à la seule exception de $\Cat/S$ tout entier.
\end{enumerate}

\msp{page 47}

\begin{enumerate}[leftmargin=*]
\item[3)] La $W$-équivalence grossière, elle, est `bonne partout' - mais ça dépend pour qui. Pour les 6 catégories qui sont au dessus de $\uc{Fibhot}\,S$, i.e.
\[
\uc{Fibhot}\,S,\quad \uc{Parf}\,S,\quad \uc{Fib}_0S,\quad \uc{Fibsc}_0S,\quad \uc{Cofib}_0S,\quad \uc{Cofibsc}_0S,
\]
\printedpage{30}
% head: A. Grothendieck : Dérivateurs, Ch. VI, Édition des Universités de Montpellier II et Paris VII
% foot: Version 9 novembre 2001
elle coïncide avec toutes les autres équivalences (moyennnant l'axiome L5), pour les trois catégories bordantes supérieures gauches % [sic] moyennnant -> moyennant, triple n
\[
\uc{Liss}\,S,\quad \uc{Fib}\,S,\quad \uc{Fibsc}\,S
\]
du diagramme, elle coïncide avec la $W$-équivalence fine droite, et dualement sur
\[
\uc{Prop}\,S,\quad \uc{Cofib}\,S,\quad \uc{Cofibsc}\,S
\]
elle coïncide avec la $W$-équivalence fine gauche. Quant à $\Cat/S$ tout entier, c'est la seule des catégories du diagramme où il y ait \textit{trois} notions intéressantes différentes de $W$-équivalence, à savoir : la grossière, et les deux fines (à gauche, à droite). Sur les autres six catégories bordantes du diagramme, il y en a en principe deux : la grossière, et une spécifiée des deux fines (défendu de regarder l'autre - à ses risques et périls \ldots) : la droite du côté gauche, la gauche du côté droit. (NB Je devrais refaire le diagramme, décidément, en mettant les fibrations à droite, pas à gauche.) Quant aux six catégories internes au diagramme, dans elles il n'y a plus qu'une seule notion de $W$-équivalence (puisque sur elles les quatre notions coïncident).
\end{enumerate}

\msp{page 48 (vide)}

\msp{page 49} (\textbf{NB} il n'y a pas de page 48)

Je résume maintenant ce que je crois avoir compris sur les catégories localisées correspondantes. Il y en a, à équivalence près, \textit{trois} différentes seulement en tout et pour tout, quand on prend les localiseurs `admissibles' dans les catégories, comme je viens de le dire (c'étant ça en somme l'intérêt objectif des prescriptions!) De façon précise :

\begin{enumerate}[leftmargin=*]
\item[$1^o$)] La \textit{$W$-équivalence grossière} donne `la même' catégorie localisée pour les 13 catégories du diagramme, i.e. toutes les flèches de transition du diagramme induisent des équivalences pour cette $W$-équivalence grossière. Le localisé sera noté \fbox{$\HOT^{\mathrm{lc}}_W(S)$}, où l'exposant lc signifie: ($W$-types d'homotopie relatif sur $S$) \textit{`localement constants.'} (Il ne me reste qu'un petit doute pour $\uc{Fibsc}\,S$ et $\uc{Cofibsc(S)}$ (\footnote[34]{Ce doute n'est pas fondé, tout marche parfaitement!}), si l'équivalence grossière y est bel et bien raisonnable. On verra en faisant la démonstration.)

\item[$2^o$)] La \textit{$W$-équivalence fine à droite} donne le même résultat pour toutes les catégories du diagramme où elle est considérée comme admise, \textit{et où elle ne coïncide pas avec la $W$-équivalence grossière} déjà considérée, donc sur les quatre catégories bordantes gauches
\[
\Cat/S,\quad \uc{Liss}\,S,\quad \uc{Fib}\,S,\quad \uc{Fibsc}\,S
\]
(\footnote[35]{\textbf{NB} Sur les trois dernières cette équivalence est la $W$-équivalence par fibres.}). Le localisé est noté \fbox{$\HOT_W(S)$}.
\end{enumerate}

\msp{page 50}

\printedpage{31}
% head: A. Grothendieck : Dérivateurs, Ch. VI, Édition des Universités de Montpellier II et Paris VII
% foot: Version 9 novembre 2001

\begin{enumerate}[leftmargin=*]
\item[$2'^o$)] Dualement, la \textit{$W$-équivalence fine gauche} donne le même résultat sur toutes les catégories du diagramme où elle est admise, et où elle ne coïncide pas avec la grossière, i.e. sur les 4 catégories
\[
\Cat/S,\quad \uc{Prop}\,S,\quad \uc{Cofib}\,S,\quad \uc{Cofibsc}\,S.
\]
Il y a lieu de noter la catégorie des fractions commune par le sigle
\[
\HOT_W(S^o)^o \qquad \text{(ou aussi $\HOT^o_W(S)$)},
\]
puisqu'elle est isomorphe à celle notée par ce sigle.
\end{enumerate}

Si on veut des notations sans ambiguïté, désignant des catégories bien définies, et non pas seulement `à équivalence près', je choisirai
\[
\HOT_W(S) \defeq \uc{Fibsc}(S)(W^{\mathrm{d}}_S)^{-1}
\]
d'où
\[
\HOT_W(S^o) \qquad \underbrace{\simeq}_{\text{\footnotesize isomorphisme de catégories}} \qquad \uc{Cofibsc}(S)(W^{\mathrm{g}}_S)^{-1},
\]
où $W^{\mathrm{d}}_S$, $W^{\mathrm{g}}_S$ désignent les ensembles de flèches correspondants à la $W$-équivalence fine droite resp. gauche. Pour $\HOT^{\mathrm{lc}}_W(S)$, je préfère une définition aussi symétrique que possible en $S$, $S^o$, et de plus une catégorie aussi petite que possible, ce qui me conduit à choisir
\[
\HOT^{\mathrm{lc}}_W(S) \defeq (\uc{Parf}\,S)W^{-1}_S,
\]
où on a noté par $W_S$ les $W$-équivalences dans $\uc{Parf}\,S$. (Inutile cette fois de préciser lesquelles - toutes les quatre \texttt{[notions]} coïncident dans $\uc{Parf}\,S$ !)

\msp{page 51}

Notons l'isomorphisme
\[
\uc{Parf}(S^o) \simeq \uc{Parf}\,S
\]
(provenant du fait que $f : X \longrightarrow Y$ dans $\Cat$ est une fibration parfaite si et seulement si $f^o$ l'est). Cet isomorphisme $X \longmapsto X^o$ est compatible avec les localiseurs (car $f \in W \Longleftrightarrow f^o \in W$), d'où
\[
\HOT^{\mathrm{lc}}_W(S^o) \simeq \HOT^{\mathrm{lc}}_W(S)
\]
(isomorphisme de catégories). On aura deux foncteurs canoniques entre les trois catégories de $W$-types d'homotopie relatifs sur $S$
\[
\begin{tikzcd}[column sep=small]
(\uc{Fibsc}_0S)W_S^{-1} \arrow[rr, "\approx"] \arrow[dr, dotted, "i'"'] & & \HOT^{\mathrm{lc}}_W(S) \arrow[dl, "i"] \arrow[dr, "j"'] & & (\uc{Cofibsc}_0S)W_S^{-1} \arrow[ll, "\approx"'] \arrow[dl, dotted, "j'"] \\
& \HOT_W(S) & & \HOT^o_W(S) \defeq \HOT_W(S^o). &
\end{tikzcd}
\]
\printedpage{32}
% head: A. Grothendieck : Dérivateurs, Ch. VI, Édition des Universités de Montpellier II et Paris VII
% foot: Version 9 novembre 2001
\texttt{[plutôt} $\defeq \HOT_W(S^o)^o$\texttt{]} Ces flèches $i$, $j$ sont suggerées par les flèches en pointillés $i'$, $j'$, provenant des inclusions % [sic] suggerées -> suggérées, missing accent
\[
\uc{Fibsc}_0S \lhook\joinrel\longrightarrow \uc{Fibsc}\,S,\quad \uc{Cofibsc}_0S \lhook\joinrel\longrightarrow \uc{Cofibsc}\,S,
\]
qui définiraient $i$, $j$ à l'aide du choix de foncteurs quasi-inverses pour les flèches horizontales. Mais $i$, $j$ ne seraient définis qu'à isomorphisme (unique) près. Il sera plus joli de choisir les quasi-inverses de façon canonique, à l'aide des flèches
\[
\uc{Fibsc}_0S \xleftarrow{\Phi_0} \uc{Parf}\,S \xrightarrow{\Psi_0} \uc{Cofibsc}_0S
\]
(où $\Phi_0 = \Phi|\uc{Parf}\,S$, $\Psi_0 = \Psi|\uc{Parf}\,S$) du diagramme p. 8. La question principale
\msp{page 52}
qui reste non élucidée, c'est si ces foncteurs $i$, $j$ sont bien \textit{pleinement fidèles,} comme je m'y attends certes.

Je veux maintenant prouver les affirmations précédentes sur les localisés.

\textcircled{1} Le localisé pour la $W$-équivalence grossière. Il suffira, à un $\varepsilon$ près, d'établir les conditions du lemme p. 11, dans les trois cas suivants.
\begin{tikzcd}[baseline=(\tikzcdmatrixname-1-1.base)]
\uc{Parf}\,S \arrow[r, "r"] & \Cat/S \arrow[l, dotted, bend left, "C"]
\end{tikzcd}
(ça donnera le cas de $\uc{Fibhot}\,S$, $\uc{Liss}\,S$ et $\uc{Prop}\,S$ sur $\Cat/S$) et
\begin{tikzcd}[baseline=(\tikzcdmatrixname-1-1.base)]
\uc{Fibsc}_0S \arrow[r, "p"] & \Cat/S \arrow[l, dotted, bend left]
\end{tikzcd},
\begin{tikzcd}[baseline=(\tikzcdmatrixname-1-1.base)]
\uc{Cofibsc}_0S \arrow[r, "q"] & \Cat/S \arrow[l, dotted, bend left]
\end{tikzcd},
qui donneront respectivement les cas de
\[
\uc{Fib}_0S,\ \uc{Fib}\,S,\ \uc{Liss}\,S \qquad \text{et} \qquad \uc{Cofib}_0S,\ \uc{Cofib}\,S,\ \uc{Prop}\,S \qquad \text{sur } S.
\]
Mais l'ennui pour $p$ et $q$, c'est que les foncteurs en sens inverse que j'avais construit (et noté encore $\Phi_0$, $\Psi_0$ (mais c'étaient essentiellement $\Phi\circ r\circ C$ et $\Psi\circ r\circ c$, à ne pas confondre avec les foncteurs $\Phi_0$, $\Psi_0$ en bas de la page précédente, et marqués aussi dans le diagramme p. 8), notons les plutôt $\Phi'$, $\Psi'$, marchent bien du côté $\Cat/S$, i.e. $p\Phi'$, $q\Psi'$ reçoivent le foncteur $\mathrm{id}_{\Cat/S}$ par une $W$-flèche, mais pas du côté des catégories but - on n'arrive pas à trouver une flèche naturelle entre $\Phi'p$ resp. $\Psi'q$, et les foncteurs identiques de $\uc{Fibsc}_0S$ resp. $\uc{Cofibsc}_0S$. Je ne vois
\msp{page 53}
pas, en effet, qu'il y en ait - la difficulté, par cette voie, parait irrémédiable. Par contre, si on regarde $\Phi'$, $\Psi'$ comme foncteurs à valeurs dans $\uc{Fib}_0S$, $\uc{Cofib}_0S$ respectivement (donc dans les catégories où on n'exige pas que les flèches soient catésiennes resp. cocartésiennes), on trouve encore des $W$-morphismes % [sic] catésiennes -> cartésiennes, missing r
\[
\mathrm{id} \longrightarrow \Phi'p, \qquad \mathrm{id} \longrightarrow \psi'q, % [sic] \psi'q -> \Psi'q, lowercase psi where \Psi' is used everywhere else
\]
ce qui liquide au moins déjà le cas de toutes les catégories des diagrammes p. 8 qui sont des
\printedpage{33}
% head: A. Grothendieck : Dérivateurs, Ch. VI, Édition des Universités de Montpellier II et Paris VII
% foot: Version 9 novembre 2001
sous-catégories pleines de $\Cat/S$ (\footnote[36]{Mais c'est idiot de prendre des foncteurs si compliqués, parce qu'on veut tout ramener à $\Cat/S$. Il est de plus simple de comparer $\uc{Fibhot}\,S$ directement à $\uc{Fibsc}_0S$ et à $\uc{Cofibsc}_0S$. $\Phi_0$ et $\Psi_0$ (qui sont en fait définis sur $\uc{Fibhot}$, pas seulement sur $\uc{Parf}$), et on trouve comme ci-dessous $\mathrm{id} \xrightarrow{W} \pi_0\Phi_0$ et $\Phi_0\pi_0 \xrightarrow{W} \mathrm{id}$ (itou pour $\Psi_0$ et $\uc{Cofibsc}_0S$), où $\pi_0 : \uc{Fibsc}_0S \longrightarrow \uc{Fibhot}\,S$.}). Il ne reste à traiter que le cas des quatre catégories de structures `scindées'
\[
\uc{Fibsc}\,S,\quad \uc{Fibsc}_0S \qquad \text{et} \qquad \uc{Cofibsc}\,S,\quad \uc{Cofibsc}_0S.
\]
Bien sûr il suffit de le faire d'un côté, disons du côté des fibrés, par dualité ($S \longmapsto S^o$) cela établira le résultat symétrique. On va donc regarder d'un peu plus près le carré de foncteurs
\[
\begin{tikzcd}[column sep=huge, row sep=large]
\uc{Fibsc}\,S \arrow[d, "p"] & \uc{Fibsc}_0S \arrow[l, hook, "\alpha'"'] \arrow[d, "p_0"'] \\
\uc{Fib}\,S \arrow[u, dotted, bend left, "\varphi"] & \uc{Fib}_0S \arrow[l, hook, "\alpha"'] \arrow[u, dotted, bend right, "\varphi_0"']
\end{tikzcd}
\]
(\footnote[37]{NB Il s'agit de prouver seulement que $p$, $p_0$ induisent des équivalences pour les localisés (pour la $W$-équivalence grossière). On sait déjà qu'il en est ainsi pour $\alpha$, et ça en résultera pour $\alpha'$.}, \footnote[38]{inutile de regarder ce carré, cf. annotations marginales ci-dessus (haut de la page)}). Le foncteur
\[
\Phi : \Cat/S \longrightarrow \uc{Fibsc}\,S
\]
induit des foncteurs en sens inverse de $p$, $p_0$ sur les sous-catégories pleines $\uc{Fib}\,S$, $\uc{Fib}_0S$ de $\Cat/S$, foncteurs marqués en pointillés comme $\varphi$, $\varphi_0$.

\msp{page 54}

Il faut seulement vérifier que si $X/S$ est dans $\uc{Fib}_0S$, alors $\varphi(X) = \Phi(X)$ est dans la sous-catégorie pleine $\uc{Fibsc}_0S$ de $\uc{Fibsc}\,S$. Mais on a vu que c'est le cas, pour peu seulement que $X/S$ soit dans $\uc{Fibhot}\,S$, car les morphismes de changement de base dans $\Phi(X)$, associés aux flèches $u : s \longrightarrow t$ dans $S$
\[
u^*_{\Phi_X} : t\backslash X \longrightarrow s\backslash X
\]
sont des $W_\infty$-équivalences si $X$ est Hot-fibré sur $S$. D'autre part, pour $X$ quelconque dans $\Cat/S$, on a une flèche, fonctorielle en $X$,
\[
X \longrightarrow \Phi(X),
\]
i.e. $\mathrm{id}_{\Cat/S} \longrightarrow \pi\Phi$ où $\pi : \uc{Fibsc}\,S \longrightarrow \Cat/S$ est le foncteur canonique, qui est une $W_\infty$-équivalence (même un homotopisme), donc une $W$-équivalence. A fortiori les foncteurs $p\varphi$, $p_0\varphi_0$ donnent lieu ($\uc{Fib}\,S$ et $\uc{Fib}_0S$ étant des sous-catégories pleines de $\Cat/S$) à des flèches
\[
\mathrm{id} \longrightarrow p\varphi, \qquad \mathrm{id} \longrightarrow p_0\varphi_0.
\]


\endgroup
% END UNIT Der-ChVI-345/Der-ChVI-345_p031-035.tex
% BEGIN UNIT Der-ChVI-345/Der-ChVI-345_p036-040.tex
\begingroup\makeatletter%
\newcommand{\printedpage}[1]{\marginpar{\footnotesize p.~#1}}%
\newcommand{\msp}[1]{\marginpar{\footnotesize [#1]}}%
\newcommand{\ul}[1]{\underline{\mathrm{#1}}}%
\newcommand{\cercle}[1]{\tikz[baseline=(c.base)]\node[draw,circle,inner sep=1pt](c){#1};}%
\makeatother\let\printedpage\derpp\let\msp\dermsp%
\let\mspin\DERmspin%


\printedpage{34}
% head: A. Grothendieck : Dérivateurs, Ch. VI, Édition des Universités de Montpellier II et Paris VII

Il reste à trouver une flèche (et $W$-équivalence grossière)
\[\varphi p \longrightarrow \mathrm{id},\]
\msp{page 55}
laquelle induira automatiquement
\[\varphi_0 p_0 \longrightarrow \mathrm{id}\]
puisque $\underline{\mathrm{Fibsc}}_0 S$ est \texttt{[une]} sous-catégorie pleine de $\ul{Fibsc}\,S$. En fait, on va définir $\varphi p \longrightarrow \mathrm{id}$, donc
\[\boxed{\gamma_X : \Phi(X) \longrightarrow X}\]
quand $X$ est fibrée scindée sur $X$ - une flèche \textit{cartésienne,} donc (si on veut) compatible avec les scindages de $\Phi(X)$, et de $X$. En fait, on aura une telle flèche, fonctorielle en $X$, pour $X$ variable dans $\ul{Fib}\,S$ (donc sans se borner à des $S$-morphismes cartésiens entre catégories fibrées sur $S$), mais la flèche $\longrightarrow$ étant bien sûr \textit{cartésienne.} (Ça a un sens, sans choisir de scindage sur $X$, et sans qu'il en existe forcément.) La définition de cette flèche $\gamma$ fait appel de façon essentielle à l'hypothèse que $X$ soit Cat-fibrant sur $S$. Bien sûr, composé à la flèche d'inclusion évidente % [sic] fibrée scindée sur X -> fibrée scindée sur S, the base category is S
\[X \xrightarrow{i_X} \Phi(X)\]
on va avoir
\[\gamma_X i_X = \mathrm{id}_X\]
(\footnote[39]{NB Comme on sait que $i_X \in W$, il en résultera bien $\gamma_X \in W$ comme on le veut.}). On définit $\gamma$ fibre par fibre,
\[\gamma_{X,s} : \Phi(X)_s \overset{\text{déf}}{=} s\backslash X \longrightarrow X_s,\]
donc par une \textit{rétraction} de $s\backslash X$ sur la sous-catégorie pleine $X_s$. Un objet de $s\backslash X$ est un couple $(x,u)$, $x \in \mathrm{Ob}\, X$, $u : s \longrightarrow p_X(x)$ ($p_X : X \longrightarrow S$ le morphisme structural de $X$). On pose
\[\gamma(x,u) = u_X^*(x) \in \mathrm{Ob}\, X_s,\]
\msp{page 56}
ce qui suppose donc qu'on a choisi des foncteurs image inverse $u_X^* : X_t \longrightarrow X_s$, pour la catégorie Cat-fibrée $X/S$. (Mais à isomorphisme unique près, $\gamma$ ne dépend pas de ces choix.) Je me dispense de la vérification fastidieuse que cette loi est bien ‘fonctorielle en $s$’, i.e. définit bien une flèche cartésienne
\[\gamma_X : \Phi(X) \longrightarrow X,\]
et de vérifier que celle-ci est fonctorielle en $X$.

\printedpage{35}
% head: A. Grothendieck : Dérivateurs, Ch. VI, Édition des Universités de Montpellier II et Paris VII

Peut-être est-il plus élégant de montrer que $\pi : \ul{Fibsc}\,S \longrightarrow \mathrm{Cat}/S$ et $\Phi : \mathrm{Cat}/S \longrightarrow \ul{Fibsc}\,S$ sont adjoints l'un de l'autre :
\[\boxed{\mathrm{Homcart}_S(\Phi(\underbrace{\;X\;}_{\in\mathrm{Cat}/S}),\underbrace{\;Y\;}_{\in\ul{Fibsc}}) \simeq \mathrm{Hom}_S(X,\pi Y),}\]
d'où résulteront des flèches d'adjonction
\[\Phi\pi \longrightarrow \mathrm{id}_{\ul{Fibsc}}, \qquad\qquad \pi\Phi \longleftarrow \mathrm{id}_{\ul{Fib}},\]
\texttt{[plutôt} $\pi\Phi \longleftarrow \mathrm{id}_{\mathrm{Cat}/S}$\texttt{]} et on aura gagné. Mais pour prouver cette formule d'adjonction, j'ai l'impression qu'il faut utiliser $\varphi p \longrightarrow \mathrm{id}$ \ldots\ Il y a des vérifications fastidieuses \ldots

Finalement, je vois maintenant qu'on a cinq cas à traiter, et trois si on joue sur les symétries :

\msp{page 57}

\[
\begin{tikzcd}[row sep=large]
\ul{Parf}\,S \arrow[d, "r"'] \\
\mathrm{Cat}/S \arrow[u, dotted, bend right, "C"']
\end{tikzcd}
\qquad\qquad\qquad
\left\{\begin{array}{lcl}
rC & \longleftarrow & \mathrm{id} \\
Cr & \longleftarrow & \mathrm{id}
\end{array}\right.
\]

\[
\begin{tikzcd}
\ul{Fibsc}\,S \arrow[dr, "\lambda"] & & \ul{Cofibsc}\,S \arrow[dl, "\mu"'] \\
& \mathrm{Cat}/S \arrow[ul, dotted, bend left, "\Phi"] \arrow[ur, dotted, bend right, "\Psi"'] &
\end{tikzcd}
\qquad
\begin{array}{l}
\left\{\begin{array}{lcl}
\lambda\Phi & \longleftarrow & \mathrm{id} \\
\Phi\lambda & \longrightarrow & \mathrm{id}
\end{array}\right. \\[2ex]
\left\{\begin{array}{lcl}
\mu\Psi & \longleftarrow & \mathrm{id} \\
\Psi\mu & \longrightarrow & \mathrm{id}
\end{array}\right.
\end{array}
\]

\[
\begin{tikzcd}
\underline{\mathrm{Fibsc}}_0 S \arrow[dr, "\lambda_0"] & & \underline{\mathrm{Cofibsc}}_0 S \arrow[dl, "\mu_0"'] \\
& \ul{Fibhot}\,S \arrow[ul, dotted, bend left, "\Phi_0"] \arrow[ur, dotted, bend right, "\Psi_0"'] &
\end{tikzcd}
\qquad
\begin{array}{l}
\left\{\begin{array}{lcl}
\lambda_0\Phi_0 & \longleftarrow & \mathrm{id} \\
\Phi_0\lambda_0 & \longrightarrow & \mathrm{id}
\end{array}\right. \\[2ex]
\left\{\begin{array}{lcl}
\mu_0\Psi_0 & \longleftarrow & \mathrm{id} \\
\Psi_0\mu_0 & \longrightarrow & \mathrm{id}
\end{array}\right.
\end{array}
\]

(\footnote[40]{\textbf{NB} Tous les foncteurs marqués sont compatibles avec la $W$-équivalence grossière, et des dix homomorphismes fonctoriels sont des $W$-équivalences grossières.})

NB Les quatre derniers homomorphismes de foncteurs, relatifs au troisième diagramme (avec les indices 0) sont simplement induits par ceux relatifs au diagramme qui le précède. Donc les vérifications-clef sont relatives aux deux premiers diagrammes, et par symétrie, on n'a finalement qu'à regarder les \textit{deux} couples de foncteurs $(r,C)$, $(\lambda,\Phi)$, et pour chacun de ces couples, exhiber les deux homomorphismes fonctoriels, des deux composés avec les foncteurs identiques.

\printedpage{36}
% head: A. Grothendieck : Dérivateurs, Ch. VI, Édition des Universités de Montpellier II et Paris VII

\msp{page 58}

\cercle{2} Les localisés pour la \textit{$W$-équivalence fine droite.} Il n'y a plus lieu, on l'a dit, de regarder les catégories internes du diagramme p. 8, où ladite équivalence coïncide avec la $W$-équivalence grossière, pour laquelle on connait déjà les localisés. Il reste à regarder les catégories
\[\mathrm{Cat}/S,\ \ \ul{Liss}\,S,\ \ \ul{Fib}\,S,\ \ \ul{Fibsc}\,S.\]
Le lemme de la page 11 nous dit qu'il suffit de montrer que dans la situation
\[\ul{Fibsc}\,S \xrightarrow{\lambda} \mathrm{Cat}/S,\]
il y a un foncteur en sens inverse, avec les propriétés specifiées. Il en résultera que toute sous-catégorie pleine $\mathcal{F}$ de $\mathrm{Cat}/S$ qui contient $\ul{Fib}\,S$, quand on la munit du localiseur de la $W$-équivalence fine droite, donne un localisé équivalent, par le foncteur induit par l'inclusion, au localisé correspondant de $\mathrm{Cat}/S$, i.e. essentiellement à $\mathrm{HOT}_W(S)$. % [sic] specifiées -> spécifiées, missing accent

Comme foncteur en sens inverse, il s'impose de prendre encore $\Phi$. Il y a à vérifier trois choses :

a) $\Phi$ \texttt{[est]} compatible avec les localiseurs, i.e. si
\begin{tikzcd}[baseline=(\tikzcdmatrixname-2-2.base), row sep=small, column sep=small]
X \arrow[rr, "f"] \arrow[dr] & & Y \arrow[dl] \\
& S &
\end{tikzcd}
sur $S$ est une $W$-équivalence fine droite, i.e. les $s\backslash X \longrightarrow s\backslash Y$ sont des $W$-équivalences, alors $\Phi(f)$ aussi. Mais ladite équivalence coïncide
\msp{page 59}
dans $\ul{Fibsc}\,S$ (comme aussi dans $\ul{Fib}\,S$ et $\ul{Liss}\,S$) avec la $W$-équivalence par fibres, et d'autre part $\Phi(f)_s$ n'est autre que la flèche précédente $s\backslash X \longrightarrow s\backslash Y$. Donc il est évident que
\[f \in \underbrace{W_S^{\mathrm{d}}}_{W\text{-équivalence droite}} \iff \Phi(f) \in \underbrace{W_S^{\mathrm{f}}}_{W\text{-équivalence par fibres}} .\]

b) Que $X \xrightarrow{i_X} \lambda\Phi(X)$ est dans $W_S^{\mathrm{d}}$, pour tout $X$ dans $\mathrm{Cat}/S$. Cela signifie donc que $i_X$ induit
\[s\backslash X \longrightarrow s\backslash\Phi(X)\]
qui est une $W$-équivalence. Ce serait O.K. s'il était clair que $X \longmapsto \Phi(X)$ commute à tout changement de base propre sur $S$ (donc en particulier à $s\backslash S \longrightarrow S$), puisqu'on sait que $X \longrightarrow \Phi(X)$ est une $W$-équivalence. C'est sans doute trop demander, mais du moins que ça commute au changement de base qui est un co-isomorphisme local (telles les $s\backslash S \longrightarrow S$).

Pour un tel changement de base $S' \xrightarrow{p} S$, vérifier que la fibre de $\Phi(X')$ en $s'$ est isomorphe à celle de $\Phi(X)$ en $s = p(s')$, je m'en contenterai. Or cette fibre est $s'\backslash X'$, déduit de $X'$ sur $S'$ par le changement de base $s'\backslash S' \longrightarrow S'$, donc déduit de $X$ par $s'\backslash S' \longrightarrow S$. Or $s'\backslash S' \longrightarrow s\backslash S$ est un isomorphisme par hypothèse
\printedpage{37}
% head: A. Grothendieck : Dérivateurs, Ch. VI, Édition des Universités de Montpellier II et Paris VII
\msp{page 60}
sur $p$ (co-isomorphisme local), donc le changement de base $s'\backslash S' \longrightarrow S$ s'identifie à $s\backslash S \longrightarrow S$, et l'image inverse est bien isomorphe à $s\backslash X$, OK.

c) Que de même
\[\Phi\lambda \longrightarrow \mathrm{id}\]
est non seulement dans $W_S$ (grossier), mais que c'est une $W$-équivalence par fibres. Donc on part de $X$ fibrée sur $S$. On regarde
\[\Phi(X) \longrightarrow X,\]
pour voir que c'est une $W$-équivalence par fibres, il suffit de voir qu'il en est ainsi de $X \longrightarrow \Phi(X)$,
\[X_s \longrightarrow s\backslash X\]
est une $W$-équivalence pour tout $s \in S$. Mais ceci est déjà vrai dès que $X$ est lisse sur $S$.

Comme le cas de la $W$-équivalence fine gauche se ramène au cas actuel par $S \longmapsto S^{o}$, on a donc prouvé tout ce qu'était attendu, à la seule exception de la pleine fidèlité des foncteurs $i$, $j$ de la page 51. % [sic] ce qu'était -> ce qui était, grammar; fidèlité -> fidélité, accent

Mais dans le résumé prospectif de la situation pour les localiseurs, p. 49-51, j'ai oublié d'autres foncteurs que $i$, $j$, de nature si possible encore plus triviale. Ce sont
\msp{page 61}
ceux qu' on peut décrire à partir du foncteur identique, sur une des catégories du diagramme où aient cours deux, voire trois, notions de $W$-équivalence différentes. Ce sont, je le rappelle

\begin{enumerate}[label={}]
\item[$1^{o}$)] Les quatre catégories
\[(*) \qquad\qquad \ul{Fibsc}\,S \longrightarrow \ul{Fib}\,S \longrightarrow \ul{Liss}\,S \longrightarrow \mathrm{Cat}/S, \qquad\qquad W_S^{\mathrm{d}} \subseteq W\]
avec la $W$-équivalence fine droite d'une part, grossière de l'autre, et dualement
\[(*') \qquad\qquad \ul{Cofibsc}\,S \longrightarrow \ul{Cofib}\,S \longrightarrow \ul{Prop}\,S \longrightarrow \mathrm{Cat}/S, \qquad\qquad W_S^{\mathrm{g}} \subseteq W\]
pour la $W$-équivalence fine gauche d'une part, la grossière de l'autre, et de plus
\item[$2^{o}$)] $\mathrm{Cat}/S$, qui était inclue dans les deux suites précédentes, et qui est la seule catégorie du diagramme donnant lieu à trois notions de $W$-équivalence différentes % [sic] inclue -> incluse, past participle of inclure
\end{enumerate}

\[
\begin{tikzcd}[row sep=small, column sep=large]
W_S^{\mathrm{d}} \arrow[dr, hook] & \\
& W_S \\
W_S^{\mathrm{g}} \arrow[ur, hook] &
\end{tikzcd}
\qquad\qquad .
\]

\printedpage{38}
% head: A. Grothendieck : Dérivateurs, Ch. VI, Édition des Universités de Montpellier II et Paris VII

Dans chacun des quatre cas ($C$, munie de deux localiseurs $\Sigma \subseteq \Sigma'$) de la suite $(*)$, le foncteur
\[C\Sigma^{-1} \longrightarrow C\Sigma'^{-1}\]
est essentiellement ‘le même’, à cause des théorèmes de comparaison. Il suffit de regarder ce foncteur dans un seul das, p. ex. un des deux extrêmes $\ul{Fibsc}\,S$ ou $\mathrm{Cat}/S$. Même remarque pour la suite $(*')$. Donc finalement il reste deux foncteurs remarquables % [sic] das -> cas, misprint
\msp{page 62}
à examiner, qui n'ont pas l'air du tout de provenir des foncteurs du diagramme p. 8 ou de quelque composé (ça ne donne que des équivalences canoniques ou bien les foncteurs $i$, $j$ de page 51, quand on va des catégories internes au diagramme vers les catégories bordantes). Le plus joli est sans doute de les regarder à partir de $\mathrm{Cat}/S$, puisque c'est là seulement qu'on les a tous les deux à la fois
\[
\begin{tikzcd}[row sep=small, column sep=tiny]
(\mathrm{Cat}/S)(W_S^{\mathrm{d}})^{-1} \approx \mathrm{HOT}_W(S) \arrow[dr, "\rho"] & \\
& (\mathrm{Cat}/S)W_S^{-1} \approx \mathrm{HOT}_W^{\mathrm{lc}}(S). \\
(\mathrm{Cat}/S)(W_S^{\mathrm{g}})^{-1} \approx \mathrm{HOT}_W^{o}(S) \arrow[ur, "\sigma"'] &
\end{tikzcd}
\]
Par contre, les foncteurs (qu'on a envie d'appeler ‘foncteurs d'inclusion’) $i$, $j$ de la page 51, en sens inverse de $\rho$, $\sigma$, sont ceux qui proviennent, par localisation $W_S^{\mathrm{d}}$ ou $W_S^{\mathrm{g}}$ selon le cas, d'un des foncteurs du diagramme page 8 (ou d'un composé de tels foncteurs), allant d'une des six catégories internes (où toutes les $W$-équivalences envisagées coïncident) vers l'une des sept catégories bordantes, sauf quand cette dernière est $\mathrm{Cat}/S$, est déterminé déjà avec quelle $W$-équivalence, $W_S^{\mathrm{d}}$ ou $W_S^{\mathrm{g}}$, on va travailler \texttt{[structure de cette phrase?]}. La détermi-\msp{page 63}% hyphen kept as printed: the word is divided at the manuscript page marker
nation peut-être la plus jolie est par
\[(*) \qquad\qquad \ul{Parf}\,S \longrightarrow \mathrm{Cat}/S\]
qui donne le foncteur $i$ quand on travaille avec $W_S^{\mathrm{d}}$, le foncteur $j$ quand on travaille avec $W_S^{\mathrm{g}}$. On voit alors tout de suite que si on compose avec le foncteur pertinent, soit $\rho$ soit $\sigma$, défini par l'inclusion $W_S^{\mathrm{d}} \subseteq W_S$ resp. $W_S^{\mathrm{g}} \subseteq W_S$ (comme parties de $\mathrm{Fl}(\mathrm{Cat}/S)$), on trouve le foncteur
\[(**) \qquad\qquad (\ul{Parf}\,S)W_S^{-1} \xrightarrow{\approx} (\mathrm{Cat}/S)W_S^{-1}\]
qui est l'\textit{équivalence canonique} entre ces deux catégories. Ainsi le foncteur correspondant à $i$, $j$ page 51, exprimé en termes de $\mathrm{Cat}/S$ comme foncteur en sens inverse de $\rho$, $\sigma$ (p. 62), est-il le composé de l'équivalence quasi-inverse de $(**)$, laquelle s'explicite par le foncteur $C$, et du foncteur sur le localisé $W_S^{\mathrm{d}}$ ou $W_S^{\mathrm{g}}$ déduit de $(*)$, i.e.


\endgroup
% END UNIT Der-ChVI-345/Der-ChVI-345_p036-040.tex
% BEGIN UNIT Der-ChVI-345/Der-ChVI-345_p041-045.tex
\begingroup\makeatletter%
\newcommand{\printedpage}[1]{\marginpar{\footnotesize p.~#1}}%
\newcommand{\msp}[1]{\marginpar{\footnotesize [#1]}}%
\newcommand{\HOT}{\mathrm{HOT}}%
\newcommand{\Cat}{\mathrm{Cat}}%
\newcommand{\Ob}{\mathrm{Ob}}%
\newcommand{\Fl}{\mathrm{Fl}}%
\newcommand{\id}{\mathrm{id}}%
\newcommand{\cA}{\mathcal{A}}%
\newcommand{\ul}[1]{\underline{\mathrm{#1}}}%
\newcommand{\circled}[1]{\tikz[baseline=(c.base)]\node[draw,circle,inner sep=1pt](c){#1};}%
\makeatother\let\printedpage\derpp\let\msp\dermsp%
\let\mspin\DERmspin%


\printedpage{39}
% head: A. Grothendieck : Dérivateurs, Ch. VI, Édition des Universités de Montpellier II et Paris VII

\[
\begin{tikzcd}[column sep=large, row sep=large]
 & \HOT_W(S)\approx(\Cat/S)(W_S^{\mathrm{d}})^{-1} \\
\HOT^{\mathrm{lc}}_W(S)\approx(\Cat/S)W_S^{-1} \arrow[ur, "i"] \arrow[r, "C"] \arrow[dr, "j"'] & (\ul{Parf}\,S)W_S^{-1} \arrow[u, "\text{\footnotesize car $W_S=W_S^{\mathrm{d}}$ sur $\ul{Parf}$}"'] \arrow[d, "\text{\footnotesize car $W_S=W_S^{\mathrm{g}}$ sur $\ul{Parf}$}"] \\
 & \HOT^{o}_W(S)\approx(\Cat/S)(W_S^{\mathrm{g}})^{-1}.
\end{tikzcd}
\]

\msp{page 64}
On voit par là qu'on a des isomorphismes canoniques
\[
\rho i\simeq\id,\qquad\sigma j\simeq\id
\]
dans la catégorie des endofoncteurs de $(\Cat/S)W_S^{-1}$ ($\approx\HOT^{\mathrm{lc}}_W(S)$), isomorphisme qui s'explicite ici par le fait que $\rho i$, $\sigma j$ sont déduits du même foncteur $C$ par passage aux localisés $W_S^{-1}$ (donc en fait $\rho i=\sigma j$ !), et qu'on a une flèche fonctorielle en $X$
\[
X\xrightarrow{i_X}C(X),
\]
qui est dans $W$ pour tout $X$. Ainsi, $\rho$ et $\sigma$ sont à isomorphisme près des \textit{rétractions} des catégories $\HOT_W(S)$, $\HOT^{o}_W(S)$ sur la catégorie $\HOT^{\mathrm{lc}}_W(S)$. De plus, ces rétractions sont ce qu'on appelle des foncteurs de localisation (comme on voit sur le diagramme de la page 62). (On dit que $C\xrightarrow{F}C'$ est un \textit{foncteur de localisation} si, désignant par $\Sigma\subseteq\Fl\,C$ l'ensemble des flèches de $C$ transformées en des isomorphismes par $F$, le foncteur
\[
C\Sigma^{-1}\xrightarrow{\bar{F}}C'
\]
qui factorise $F$ est une équivalence de catégories.)

\msp{page 65}
On sait d'ailleurs que si un foncteur $f$ a un adjoint $g$ (à gauche ou à droite), alors $f$ est un foncteur \texttt{[de]} localisation si et seulement si l'adjoint $g$ est pleinement fidèle. Pour prouver que $i$, $j$ sont pleinement fidèles, il serait donc tentant de prouver qu'ils sont adjoints des foncteurs $\rho$, $\sigma$.

Pour nous y reconnaître, je vais poser
\[
\left\{
\begin{aligned}
&\cA=\Cat/S,\ \ \Sigma_0=W_S^{\mathrm{d}},\ \ \Sigma=W_S\qquad\quad(\subseteq\Fl(\cA))\\
&C:\cA\longrightarrow\cA\ \text{ le foncteur familier, qui a la vertu}\\
&\qquad\qquad\qquad\quad C(\Sigma)\subseteq\Sigma_0\quad.
\end{aligned}
\right.
\]

On a donc la situation
\printedpage{40}
% head: A. Grothendieck : Dérivateurs, Ch. VI, Édition des Universités de Montpellier II et Paris VII
\[
(*)\qquad\qquad(\cA,\Sigma_0)\ \underset{C}{\overset{\rho_0\,=\,\id_{\cA}}{\rightleftarrows}}\ (\cA,\Sigma),
\]
induisant par passage aux localisés
\[
\cA\Sigma_0^{-1}\quad\underset{\bar{C}}{\overset{\rho\,=\,\bar{\rho}_0}{\rightleftarrows}}\ \cA\Sigma^{-1},
\]
on veut vérifier que $\rho$ et $C$ \texttt{[plutôt $\bar{C}$]} sont adjoints l'un de l'autre, $C$ \texttt{[plutôt $\bar{C}$]} \textit{sûrement à droite,} ce qui signifie que l'on doit avoir des homomorphismes fonctoriels
\[
\left\{
\begin{aligned}
&\id_{\cA\Sigma^{-1}}&&\rightleftarrows&&\rho\,\bar{C}\quad(=\overline{\rho_0C})\\
&\id_{\cA\Sigma_0^{-1}}&&\rightleftarrows&&\bar{C}\rho\quad(=\overline{C\rho_0})
\end{aligned}
\right.
\]
(\footnote[41]{mais si $C$ est à \textit{droite,} les flèches ici doivent aller en sens inverse et c'est bien à quoi je dois me résoudre p.~67.}) \texttt{[les flèches en bas vont en direction corrigée, mais les flèches en haut ne sont pas effacées]} satisfaisant deux relations de compatibilité sur lesquelles on reviendra. On peut espérer trouver ces homomorphismes fonctoriels à partir de
\[
\left\{
\begin{aligned}
&\id_{\cA}&&\xrightarrow{\ i\ }&&\rho_0C&&=&&C\\
&\id_{\cA}&&\xleftarrow{\ p\ }&&C\rho_0&&=&&C
\end{aligned}
\right.
\]

\msp{page 66}
et dans notre cas on dispose bel et bien de
\[
i:\id_{\cA}\longrightarrow C\qquad\text{i.e. de }\ i_X:X\longrightarrow C(X)
\]
fonctoriel en $X\in\Ob\,\cA$ (et de plus $i_X$ est dans $\Sigma$). Il n'y a pas cependant de flèche naturelle en sens inverse $p_X$
\[
C(X)\xrightarrow{\ ?\ }X.
\]
J'avais songé à la projection $\pi_X:C(X)\longrightarrow X$, déduite de la construction $C(X)\simeq X\times_S(\ul{Ch}_\infty(S),\sigma)$, mais ce n'est malheureusement pas un $S$-morphisme. Comme on cherche une flèche dans $\cA\Sigma_0^{-1}$, on pourrait se contenter d'avoir une flèche fonctorielle en sens inverse, du moment qu'elle soit dans $\Sigma_0$, donc puisse s'inverser quand on travaille dans $\cA\Sigma_0^{-1}$. Le candidat evident est $i_X:X\longrightarrow C(X)$ dont on vient de parler. L'ennui, c'est que $i_X$ est seulement dans $\Sigma$, pas dans $\Sigma_0$ (\footnote[42]{S'il n'en était pas ainsi, le type d'homotopie relatif sur $S$ défini par $X/S$ serait localement constant (puisqu'il en est ainsi pour celui défini par $C(X)$), ce qui n'est pas toujours le cas.}). Pourtant, quand $X$ est lui-même de la forme $C(Y)$ (plus généralement, s'il est un Hot-fibré sur $S$), alors le fait que $X\longrightarrow C(X)$ soit dans $\Sigma$ implique déjà qu'il soit dans $\Sigma_0$. Notons donc ceci % [sic] evident -> évident, accent missing
\[
\left\{
\begin{array}{lrl}
i:\id_{\cA}\longrightarrow C\ \text{tel que} & 1^o) & i_X:X\longrightarrow C(X)\ \textit{dans}\ \Sigma\\[1ex]
 & & \text{pour tout }X\in\Ob\,\cA\\[1ex]
 & 2^o) & i_X\ \text{dans }\Sigma_0\text{ si }X\text{ de la}\\[1ex]
 & & \text{forme }C(Y)
\end{array}
\right.
\]

\printedpage{41}
% head: A. Grothendieck : Dérivateurs, Ch. VI, Édition des Universités de Montpellier II et Paris VII

\msp{page 67}
L'idée me vient, après des essais merdouillants \texttt{[sic]}, que $\bar{C}$ pourrait être l'adjoint à droite, et non à gauche \texttt{[c'est la version corrigée (droite et gauche échangées) cf.~p.~65]}, i.e. que les flèches cherchées pourraient aller en sens inverse
\[
\left\{
\begin{aligned}
&\id_{\cA\Sigma^{-1}}&&\xleftarrow{\ p\ }&&\rho\,\bar{C}\quad(=\overline{\rho_0C})\\
&\id_{\cA\Sigma_0^{-1}}&&\xrightarrow{\ i\ }&&\bar{C}\rho\quad(=\overline{C\rho_0})
\end{aligned}
\right.
\]
(j'inverse du coup les notations pour $p$ et $i$), donc on est amené si possible à chercher des flèches
\[
\begin{aligned}
&\id_{\cA}&&\xleftarrow{\ p\ }&&\rho_0C&&=&&C\\
&\id_{\cA}&&\xrightarrow{\ i\ }&&C\rho_0&&=&&C,
\end{aligned}
\]
donc des flèches fonctorielles en $X$
\[
X\xrightarrow{i_X}C(X)\xrightarrow{p_X}X.
\]
On tient bel et bien une $i_X$ fonctorielle - toujours la même! Pour $p_X$, on n'en a pas plus qu'avant, mais se rappelant qu'il suffit de le trouver dans $\cA\Sigma^{-1}$ (et non plus $\cA\Sigma_0^{-1}$!), il suffit d'inverser la flèche $i_X$, qui est en effet un \textit{isomorphisme} dans $\cA\Sigma^{-1}$ (mais non dans $\cA\Sigma_0^{-1}$). On posera donc
\[
p_X=i_X^{-1}:C(X)\longrightarrow X\qquad\quad\text{dans }\cA\Sigma^{-1}.
\]
Il faut prouver deux compatibilités seulement, pour les foncteurs $\bar{C}$ et $\rho$ entre $\cA\Sigma_0^{-1}$ et $\cA\Sigma^{-1}$ (et j'écrirai $C$ au lieu de $\bar{C}$) : les deux composés
\[
\left\{
\begin{array}{l}
\begin{tikzcd}[ampersand replacement=\&]
\rho \arrow[r, "\rho*i"] \arrow[rr, bend left=20, dotted, "\id"] \& \rho\,C\rho \arrow[r, "p*\rho"] \& \rho
\end{tikzcd}
\qquad\quad\text{\footnotesize (relation pour foncteurs à valeurs dans $\cA\Sigma^{-1}$)}\\[2ex]
\begin{tikzcd}[ampersand replacement=\&]
C \arrow[r, "i*C"] \arrow[rr, bend right=20, dotted, "\id"'] \& C\rho\,C \arrow[r, "C*p"] \& C
\end{tikzcd}
\qquad\quad\text{\footnotesize (relation pour foncteurs à valeurs dans $\cA\Sigma_0^{-1}$)}
\end{array}
\right.
\]
doivent être l'identité $\id_\rho$ resp. $\id_C$.

\msp{page 68}
La première relation se vérifie tautologiquement, elle revient à dire que dans $\cA\Sigma^{-1}$, on a un composé identité des
\[
X\xrightarrow{i_X}C(X)\xrightarrow{i_X^{-1}}X\qquad\quad!
\]
La seconde signifie que pour tout $X$, on a un composé identité dans $\cA\Sigma_0^{-1}$
\[
C(X)\xrightarrow{i_{C(X)}}C(C(X))\xrightarrow{C(p_X)}C(X).
\]
Or $p_X:C(X)\longrightarrow X$ avait été défini dans $A\Sigma^{-1}$ comme inverse $i_X^{-1}$ de $i_X$, d'autre part le foncteur $C$ sur $\cA$ définit par hypothèse $\cA\Sigma^{-1}\longrightarrow\cA\Sigma_0^{-1}$, plus précisément $C(\Sigma)\subseteq\Sigma_0$ % [sic] $A\Sigma^{-1}$ (italic A) -> $\mathcal{A}\Sigma^{-1}$, elsewhere script A
\printedpage{42}
% head: A. Grothendieck : Dérivateurs, Ch. VI, Édition des Universités de Montpellier II et Paris VII
ainsi $C(i_X)$ est inversible dans $\cA\Sigma_0^{-1}$, et $C(p_X)=C(i_X)^{-1}$. Donc la compatibilité revient à la relation
\[
C(i_X)^{-1}i_{C(X)}=\id\qquad\quad\textit{dans}\ \ \cA\Sigma_0^{-1},
\]
i.e.
\[
\boxed{\,i_{C(X)}=C(i_X)\,}:C(X)\longrightarrow C^2(X).
\]
En d'autre termes, il faut prouver le : % [sic] d'autre termes -> d'autres termes, plural

\textbf{Lemme.} \textit{Soit $X$ dans $\Cat/S$, considérons les deux flèches}
\[
i_{C(X)},\ C(i_X):C(X)\rightrightarrows C(C(X)).
\]
\textit{Ces flèches sont égales, sinon au sens absolu, du moins dans $(\Cat/S)(W_S^{\mathrm{d}})^{-1}$.}

Ici $f_X:X\longrightarrow S$ est la projection standard et
\[
\begin{aligned}
C(X)\ &=\ \quad\{(x,c)\ \mid\ x\in\Ob\,X,\ c:f_X(x)\ \underbrace{\longrightarrow}_{\text{chemin dans }S}\ t\}\\
C^2(X)\ &=\ \{(x,c,c')\ \mid\ x\in\Ob\,X,\ \underbrace{f_X(x)\xrightarrow{c}t\xrightarrow{c'}t'}_{c,\ c'\text{ chemins dans }S}\},
\end{aligned}
\]
les projections dans $S$ sont données par
\[
\begin{aligned}
(x,c)&&\longmapsto&&t&&=&&\text{but }c\\
(x,c,c')&&\longmapsto&&t'&&=&&\text{but }c'
\end{aligned}
\]

\msp{page 69}
On trouve
\[
i_{C(X)}(x,c)=(x,c,1_t):
\begin{tikzcd}
x \arrow[d] & & \\
f_X(x) \arrow[r, "c"] & t \arrow[r, "1_t", "\text{chemin vide}"'] & t
\end{tikzcd}
\]
\[
C(i_X)(x,c)=(x,1_{f_X(x)},c):
\begin{tikzcd}
x \arrow[d] & & \\
f_X(x) \arrow[r, "1_{f_X(x)}", "\text{chemin vide}"'] & f_X(x) \arrow[r, "c"] & t
\end{tikzcd}
\]
Mais on a le choix entre deux théories de chemins, pour construire $C(X)$ : soit avec les $\ul{Ch}_\infty(S)$, soit avec les $\ul{Ch}(S)$. ‘Ça donne les mêmes résultats’, mais suivant les situations l'un est plus commode que l'autre. Je vais travailler avec les $\ul{Ch}(S)$, auquel cas les $1_S$ sont des unités à gauche et à droite, et on a une loi de composition pour les chemins. Cette loi nous permet de définir
\[
\begin{aligned}
C^2(X)&&\xrightarrow{\ \gamma\ }&&C(X)\\
(x,c,c')&&\longmapsto&&(x,c'\circ c),
\end{aligned}
\]
et les composés de $i_{C(X)}$ et de $C_{i(X)}$ avec cette flèche sont $\id_{C(X)}$ dans les deux cas % [sic] $C_{i(X)}$ -> $C(i_X)$, the arrow considered is C(i_X)
\printedpage{43}
% head: A. Grothendieck : Dérivateurs, Ch. VI, Édition des Universités de Montpellier II et Paris VII
\[
\begin{tikzcd}
C(X) \arrow[r, shift left] \arrow[r, shift right] \arrow[rr, bend right=15, "\id"'] & C^2(X) \arrow[r, "\gamma"] & C(X).
\end{tikzcd}
\]

Donc pour montrer que $i_{C(X)}=C(i_X)$ dans $(\Cat/S)(W_S^{\mathrm{d}})^{-1}$, il suffit de prouver que $\gamma$ est dans $W_S^{\mathrm{d}}$. Mais on sait que le foncteur C transforme $W_S$ en $W_S^{\mathrm{d}}$, donc $C(i_X)$ est dans $W_S^{\mathrm{d}}$, et comme $\gamma\circ C(i_X)=\id$, $\gamma$ est dans $W_S^{\mathrm{d}}$, OK! % [sic] foncteur C (upright roman C) -> foncteur $C$, math italic elsewhere

\msp{page 70}
J'ai envie de grouper ce qui a été prouvé dans un énoncé récapitulatif.

\textbf{Théorème-Scholie.} On considère le diagramme de catégories de la page 8 (\footnote[43]{On suppose sur $W$ les conditions a) b) c) de p.~45, de plus L5 (p.~43). Cf. commentaires p.~81.}).

\circled{I} Quand on munit toutes les 13 catégories du diagramme du localiseur $W_S$ ($W$-équivalence grossière, i.e. $W$-équivalence dans Cat, sans plus), \textit{tous} les foncteurs du diagramme induisent des équivalences sur les catégories localisées (\footnote[44]{De plus, si $\mathcal{F}$ est une sous-catégorie pleine de $\Cat/S$ qui contient une des deux sous-catégories $\ul{Fib}_0S$, $\ul{Cofib}_0S$ (ce qui inclut le cas des 8 catégories des diagrammes, autres que $\Cat/S$ elle-même), l'inclusion $(\mathcal{F},W_S|\mathcal{F})$ dans $(\Cat/S,W_S)$ induit une équivalence pour les localisés.}). Désignons par
\[
\HOT^{\mathrm{lc}}_W(S)
\]
‘la’ catégorie localisée commune.

\circled{II} Munissons les quatre catégories bordantes gauches
\[
(*)\qquad\qquad
\begin{tikzcd}
\Cat/S & \ul{Liss}\,S \arrow[l, hook'] & \ul{Fib}\,S \arrow[l, hook'] & \ul{Fibsc}\,S \arrow[l, hook']
\end{tikzcd}
\]
du localiseur $W_S^{\mathrm{d}}$ de la $W$-équivalence fine \textit{droite} sur $S$ ( $\begin{tikzcd}[row sep=small, column sep=small]
X \arrow[rr, "f"] \arrow[dr] & & Y \arrow[dl]\\
 & S &
\end{tikzcd}$ est dans $W_S^{\mathrm{d}}$ si et seulement si les $s\backslash X\longrightarrow s\backslash Y$ le sont \texttt{[plutôt, sss ils sont dans $W$]}, pour tout $s\in\Ob\,S$). Les foncteurs du diagramme, i.e. ceux de $(*)$, entre ces catégories, donc aussi leurs composés, induisent des équivalences entre les localisées (\footnote[45]{Plus généralement, ça marche pour toute sous-catégorie pleine de $\Cat/S$ qui contient $\ul{Fib}\,S$ (et dualement, avec $W_S^{\mathrm{g}}$ au lieu de $W_S^{\mathrm{d}}$, si $\mathcal{F}$ contient $\ul{Cofib}\,S$).}). Désignons par
\[
\HOT_W(S)
\]
‘la’ catégorie localisée commune.

Enoncé dual pour
\[
(*')\qquad\qquad
\begin{tikzcd}
\Cat/S & \ul{Prop}\,S \arrow[l, hook'] & \ul{Cofib}\,S \arrow[l, hook'] & \ul{Cofibsc}\,S, \arrow[l, hook']
\end{tikzcd}
\]

% foot (each page): Version 9 novembre 2001


\endgroup
% END UNIT Der-ChVI-345/Der-ChVI-345_p041-045.tex
% BEGIN UNIT Der-ChVI-345/Der-ChVI-345_p046-050.tex
\begingroup\makeatletter%
\newcommand{\printedpage}[1]{\marginpar{\footnotesize p.~#1}}%
\newcommand{\msp}[1]{\marginpar{\footnotesize [#1]}}%
\newcommand{\HOT}{\mathrm{HOT}}%
\newcommand{\Cat}{\mathrm{Cat}}%
\newcommand{\Hot}{\mathrm{Hot}}%
\newcommand{\hot}{\mathrm{hot}}%
\newcommand{\Ob}{\operatorname{Ob}}%
\newcommand{\Hom}{\operatorname{Hom}}%
\newcommand{\ul}[1]{\underline{\mathrm{#1}}}%
\newcommand{\cercle}[1]{\tikz[baseline=(c.base)]\node[draw,circle,inner sep=1pt](c){#1};}%
\newcommand{\ovale}[1]{\tikz[baseline=(c.base)]\node[draw,rounded corners=0.8ex,inner sep=2pt](c){#1};}%
\makeatother\let\printedpage\derpp\let\msp\dermsp%
\let\mspin\DERmspin%

\setcounter{footnote}{45}

\printedpage{44}
% head: A. Grothendieck : Dérivateurs, Ch. VI, Édition des Universités de Montpellier II et Paris VII
% foot: Version 9 novembre 2001
et les localiseurs $W_S^{\mathrm{g}}$ dans ces catégories ($f\in W_S^{\mathrm{g}}$ si et seulement si les $f/s : X/s\longrightarrow Y/s$ sont dans $W$). On trouve une `même catégorie localisée', notée
\[ \HOT^o_W(S), \]
et on a tautologiquement
\msp{page 71}
\[ \HOT^o_W(S)\approx\HOT_W(S^o). \]
\texttt{[plutôt} $(\HOT_W(S^o))^o$ \texttt{cf. p. 50]}

\ovale{III}\quad Sur les six catégories internes du diagramme, les trois notions de $W$-équivalence, en termes de $W_S$, $W_S^{\mathrm{d}}$, $W_S^{\mathrm{g}}$ coïncident, et coïncident avec la $W$-équivalence par fibres $W_S^{\mathrm{f}}$. Sur les trois catégories de $(*)$ autres que $\Cat/S$, les localiseurs $W_S^{\mathrm{f}}$, $W_S^{\mathrm{d}}$ coïncident; sur les trois catégories de $(*')$ autres que $\Cat/S$, les localiseurs $W_S^{\mathrm{f}}$, $W_S^{\mathrm{g}}$ coïncident. Mais en général (\footnote{sans doute dans tous les cas où $S$ n'est pas équivalente à une catégorie dicrète.% [sic] dicrète -> discrète, faute d'orthographe
}), sur ces six catégories, ces localiseurs $W_S^{\mathrm{f}}$ ne sont pas égaux à $W_S$, plus précisément donnent des localisées non équivalentes, elles sont donc munies chacune de \textit{deux} localiseurs essentiellement distincts, $W_S^{\mathrm{f}}$ et $W_S$. La catégorie $\Cat/S$ est la seule catégorie du diagramme sur laquelle les localiseurs $W_S$, $W_S^{\mathrm{d}}$ et $W_S^{\mathrm{g}}$ sont tous les trois essentiellment distincts. % [sic] essentiellment -> essentiellement, lettre omise
C'est la seule qui permette d'obtenir simultanément les trois catégories localisées, comme
\[
\begin{tikzcd}[column sep=small]
 & \HOT^{\mathrm{lc}}_W(S)\simeq(\Cat/S)W_S^{-1} & \\
\HOT_W(S)\simeq(\Cat/S)(W_S^{\mathrm{d}})^{-1} \arrow[ur,"\rho"] & & \HOT^o_W(S)\simeq(\Cat/S)(W_S^{\mathrm{g}})^{-1} \arrow[ul,"\sigma"']
\end{tikzcd}
\]

Ici, $\rho$ et $\sigma$ sont des foncteurs de localisation (comme il est évident sur la définition).

Tout foncteur ou foncteur composé du diagramme page 8, allant d'une catégorie interne au diagramme vers une catégorie bordante, induit un foncteur pour les localisées pour $W_S^{\mathrm{d}}$ (quand la catégorie but fait partie des quatre gauches,
\msp{page 72}
rappelées dans $(*)$, p. 70), soit pour $W_S^{\mathrm{g}}$ (quand ces catégories font partie des quatre du côté droit, rappelées dans $(*')$), lesquels localiseurs coïncident aussi avec $W_S^{\mathrm{f}}$ sauf dans le seul cas de la catégorie $\Cat/S$. Tous les foncteurs ainsi obtenus se réduisent, à équivalence près, à deux, suivant que la catégorie localisée but est $\HOT_W(S)$ (cas de la localisation $W_S^{\mathrm{d}}$), ou qu'elle est $\HOT^o_W(S)$ (cas de la localisation $W_S^{\mathrm{g}}$). On trouve ainsi des foncteurs (en sens inverse de $\rho$ et de $\sigma$)
\printedpage{45}
% head: A. Grothendieck : Dérivateurs, Ch. VI, Édition des Universités de Montpellier II et Paris VII
% foot: Version 9 novembre 2001
\[
\begin{tikzcd}
 & \HOT^{\mathrm{lc}}_W(S) \arrow[dl,"i"'] \arrow[dr,"j"] & \\
\HOT_W(S) & & \HOT^o_W(S).
\end{tikzcd}
\]

Les foncteurs $\rho$, $\sigma$ s'interchangent par dualité $S\longmapsto S^o$, et de même pour $i$, $j$.

\ovale{IV}\quad Les foncteurs $\rho$, $i$ sont adjoints l'un de l'autre dans cet ordre, i.e. on a pour deux objets $\xi\in\Ob\HOT_W(S)$, $\eta\in\Ob\HOT^{\mathrm{lc}}_W(S)$
\[ \Hom(\rho(\xi),\eta)\simeq\Hom(\xi,i(\eta)) \]
(isomorphisme de bifoncteurs). Et dualement les foncteurs $\sigma$, $j$ sont adjoint l'un de l'autre % [sic] adjoint -> adjoints, accord du pluriel
dans cet ordre : pour $\xi\in\Ob\HOT^o_W(S)$, $\eta\in\Ob\HOT^{\mathrm{lc}}_W(S)$
\[ \Hom(\xi,j(\eta))\simeq\Hom(\sigma(\xi),\eta). \]

\msp{page 73}
Il s'ensuit notamment que \textit{les foncteurs $i$, $j$ sont pleinement fidèles.}

\cercle{V}\quad On a des foncteurs évidents
\[
\begin{aligned}
&(*) &\qquad& \HOT_W(S) &\quad& \longrightarrow &\quad& \underline{\mathrm{Hom}}(S^o,\Hot_W)\\
&(*') && \HOT^o_W(S) && \longrightarrow && \underline{\mathrm{Hom}}(S,\Hot_W).
\end{aligned}
\]
Le premier s'interprète sur les quatre catégories $(*)$ (bordantes gauches) de page 70, qui servent indifférement à décrire $\HOT_W(S)$, par % [sic] indifférement -> indifféremment, faute d'orthographe
\[ s\longmapsto\hot_W(X_s) \]
dans le cas des trois catégories distinctes de $\Cat/S$, et par
\[ s\longmapsto\hot_W(s\backslash X) \]
dans le cas de $\Cat/S$. La dépendence fonctorielle par rapport à $s$ est évidente dans le cas % [sic] dépendence -> dépendance, faute d'orthographe
de $\Cat/S$ (donnée par $s\longmapsto\hot_W(s\backslash X)$), et dans le cas de $\ul{Fibsc}\,S$, $\ul{Fib}\,S$, où on dispose, pour $u : s\longrightarrow t$ dans $S$, du changement de base
\[ u^* : X_t\longrightarrow X_s. \]
Dans le cas de $\ul{Liss}\,S$, la dépendence fonctorielle se déduit de celle de $\hot_W(s\backslash X)$, grâce % [sic] dépendence -> dépendance, faute d'orthographe
au fait que l'inclusion
\[ X_s\longrightarrow s\backslash X \]
est alors une $W$-équivalence.

On a une interprétation duale de l'application
\printedpage{46}
% head: A. Grothendieck : Dérivateurs, Ch. VI, Édition des Universités de Montpellier II et Paris VII
% foot: Version 9 novembre 2001
\msp{page 74}
$(*)'$, \texttt{[mot illisible]}
\[ s\longmapsto\hot_W(X_s)\qquad\text{ou}\qquad s\longmapsto\hot_W(s\backslash X), \]
avec la loi fonctorielle en $s$ déduite du foncteur $t\backslash X\longrightarrow s\backslash X$ \texttt{[plutôt} $X/s\longrightarrow X/t$\texttt{]} (pour $u : s\longrightarrow t$) dans le cas du $\Cat/S$, du foncteur cochangement de base dans les cas de $\ul{Cofib}\,S$ et $\ul{Cofibsc}\,S$, et de la variante pour $\Cat/S$ et de la $W$-équivalence
\[ X_s\longrightarrow X/s \]
dans le cas de $\ul{Prop}\,S$.

On trouve alors ceci : pour qu'un objet $\xi$ de $\HOT_W(S)$ (resp. $\HOT^o_W(S)$) soit dans l'image essentielle de $\HOT^{\mathrm{lc}}_W(S)$, il faut et il suffit que le foncteur associé $S^o\longrightarrow\Hot_W$ (resp. $S\longrightarrow\Hot_W$) soit un système local (\footnote{il faudra encore écrire la vérification}), i.e. transforme toute flèche de $S$ en un isomorphisme, i.e. se factorise par $(\Pi S)^o\longrightarrow\Hot_W$ resp. $\Pi S\longrightarrow\Hot_W$ (\footnote{où $\Pi S$ est le groupoïde fondamental de $S$}). Donc pour un $X$ dans $\Cat/S$, cela signifie soit que les foncteurs
\[ t\backslash X\longrightarrow s\backslash X \]
associés à des $u : s\longrightarrow t$ dans $S$ soient tous des $W$-équivalences (cas de $\HOT_W(S)$), soit que les foncteurs
\[ X/s\longrightarrow X/t \]
le soient (cas de $\HOT^o_W(S)$). (NB. ces conditions ne sont nullement équivalentes. Mais si $X$ \texttt{[est]} lisse \texttt{[sur]} $S$,
\msp{page 75}
la première implique la seconde, et l'inverse est vrai si $X$ est propre sur $S$, et dans ce cas, la condition équivaut à ce que $X$ soit dans $\ul{Fibhot}$. Enfin, si $X$ est déjà dans $\ul{Fibhot}\,S$, alors les deux conditions sont automatiquement satisfaites (\footnote{On aurait dû dans le diagramme introduire aussi $\ul{Liss}_0S$, $\ul{Prop}_0S$ !}).
\[
\begin{tikzcd}[column sep=small, row sep=small]
\ul{Fibsc}_0S \arrow[dd,hook] \arrow[drr] & & & & & & & & \ul{Cofibsc}_0S \arrow[dll] \arrow[dd,hook] \\
 & & \ul{Fib}_0S \arrow[dd,hook] \arrow[dr,hook] & & \ul{Parf}\,S \arrow[dl,hook] \arrow[dd,hook] \arrow[dr,hook] & & \ul{Cofib}_0S \arrow[dl,hook] \arrow[dd,hook] & & \\
\ul{Fibsc}\,S \arrow[drr] & & & \ul{Liss}_0S \arrow[dr,hook] \arrow[dd,hook] & & \ul{Prop}_0S \arrow[dl,hook] \arrow[dd,hook] & & & \ul{Cofibsc}\,S \arrow[dll] \\
 & & \ul{Fib}\,S \arrow[dr,hook] & & \ul{Fibhot}\,S \arrow[dd,hook] \arrow[uuullll,dotted,bend right,"\Phi_0"'] \arrow[uuurrrr,dotted,bend left,"\Psi_0"'] & & \ul{Cofib}\,S \arrow[dl,hook] & & \\
 & & & \ul{Liss}\,S \arrow[dr,hook] & & \ul{Prop}\,S \arrow[dl,hook] & & & \\
 & & & & \Cat/S \arrow[uuullll,dotted,bend left,"\Phi"] \arrow[uuurrrr,dotted,bend right,"\Psi"'] & & & &
\end{tikzcd}
\]
\printedpage{47}
% head: A. Grothendieck : Dérivateurs, Ch. VI, Édition des Universités de Montpellier II et Paris VII
% foot: Version 9 novembre 2001
(\footnote{Le diagramme de \textit{quinze} catégories remarquables.}). De plus, si deux éléments $\xi\in\Ob\HOT_W(S)$, $\xi'\in\Ob\HOT^o_W(S)$ proviennent d'un même $\eta\in\Ob\HOT^{\mathrm{lc}}_W(S)$, alors les deux foncteurs correspondants
\[ (\Pi S)^o\xrightarrow{\ \varphi\ }\Hot_W\xleftarrow{\ \varphi'\ }\Pi S \]
sont \textit{associés,} i.e. chacun est isomorphe à `l'inverse' de l'autre.
\[
\left\{
\begin{array}{lcl}
\varphi^o(s) & = & \varphi(s)\\
\varphi^o(u) & = & \varphi(u)^{-1}.
\end{array}
\right.
\]

\ovale{VI}\quad NB Dans la nouvelle version du diagramme, où j'ai rajouté des catégories $\ul{Liss}_0S$ et $\ul{Prop}_0S$ pour rendre le tracé plus transparent, les catégories appelés `bordantes' deviennent % [sic] appelés -> appelées, accord du participe
le `rez-de-chausseée', et les catégories `internes', le `premier étage' ou `étage'. Tous % [sic] rez-de-chausseée -> rez-de-chaussée, lettre en trop
les 7 carrés du diagramme sont catésiens. Toutes les flèches sont des inclusions pleinement % [sic] catésiens -> cartésiens, lettre omise
fidèles, sauf les quatre flèches horizontales extérieures des catégories de Cat-fibrés ou de Cat-cofibrés \textit{scindés}, avec morphismes exclusivement cartésiens ou cocartésiens et même compatibles avec les scindages, vers les catégories correspondantes de structures non scindées.

\msp{page 76}
\[
\begin{tikzcd}[column sep=small, row sep=small]
 & & & & \ul{Fibhot}\,S \arrow[dd,hook,"\delta"] \arrow[dddllll,dotted,bend right,"\Phi_0"'] \arrow[dddrrrr,dotted,bend left,"\Psi_0"] & & & & \\
 & & & \ul{Liss}_0S \arrow[ur,hook] \arrow[dd,hook] & & \ul{Prop}_0S \arrow[ul,hook] \arrow[dd,hook] & & & \\
 & \ul{Fib}_0S \arrow[urr,hook] \arrow[dd,hook] \arrow[ddrrr,dotted,bend left] & & & \Cat/S \arrow[dd,"C"] \arrow[dddllll,dotted,bend right,"\Phi"'] \arrow[dddrrrr,dotted,bend left,"\Psi"] & & & \ul{Cofib}_0S \arrow[ull,hook] \arrow[dd,hook] \arrow[ddlll,dotted,bend right] & \\ % Reading verified against PDF leaf 9: sources Fib_0 S and Cofib_0 S, target middle Parf S.
\ul{Fibsc}_0S \arrow[ur] \arrow[dd,hook] & & & \ul{Liss}\,S \arrow[ur,hook] & & \ul{Prop}\,S \arrow[ul,hook] & & & \ul{Cofibsc}_0S \arrow[ul] \arrow[dd,hook] \\
 & \ul{Fib}\,S \arrow[urr,hook] & & & \ul{Parf}\,S \arrow[uuul] \arrow[uuur] \arrow[dd,dotted,"\mathrm{id}"] & & & \ul{Cofib}\,S \arrow[ull,hook] & \\
\ul{Fibsc}\,S \arrow[ur] & & & & & & & & \ul{Cofibsc}\,S \arrow[ul] \\
 & & & & \ul{Parf}\,S \arrow[uuul,dotted,hook] \arrow[uuur,dotted,hook] & & & &
\end{tikzcd}
\]

J'ai tracé en pointillés verts \texttt{[couleur dans l'original seulement]} les cinq foncteurs $\Phi$, $\Psi$, $\Phi_0$, $\Psi_0$, $C$ qui par passage aux localisés donnent des quasi-inverses pour toutes les équivalences de catégories énoncées dans \cercle{I} \cercle{II} du présent Scholie. De façon précise, le foncteur
\[ \Phi : \Cat/S\longrightarrow\ul{Fibsc}\,S \]
associe à toute $X$ dans $\Cat/S$ la catégorie scindée sur $S$
\[ s\longmapsto s\backslash X \]
\printedpage{48}
% head: A. Grothendieck : Dérivateurs, Ch. VI, Édition des Universités de Montpellier II et Paris VII
% foot: Version 9 novembre 2001
sur $S$. Ce foncteur, et ses restrictions à $\ul{Liss}\,S$, $\ul{Fib}\,S$ et, plus généralement, à toute sous-catégorie pleine $\mathcal{F}$ de $\Cat/S$ contenant $\ul{Fib}\,S$, est compatible avec les localiseurs $W_S$ d'une part (sur $\mathcal{F}$ et $\ul{Fibsc}\,S$), $W_S^{\mathrm{d}}$ d'autre part, et passe donc aux localisés correspondants, pour fournir le quasi-inverse cherché de
\[ (\ul{Fibsc}\,S)\Sigma^{-1}\longrightarrow\mathcal{F}\Sigma^{-1}\qquad(\Sigma=W_S\text{ ou }W_S^{\mathrm{d}}) \]

On peut aussi, pour toute sous-catégorie pleine $\mathcal{F}'$ intermédiaire
\[ \ul{Fib}\,S\subseteq\mathcal{F}'\subseteq\mathcal{F}\subseteq\Cat/S, \]
regarder $\Phi$ comme un foncteur $\mathcal{F}\longrightarrow\mathcal{F}'$ (plus précisément, prendre le composé
\[ \mathcal{F}\xrightarrow{\Phi|\mathcal{F}}\ul{Fibsc}\,S\longrightarrow\ul{Fib}\,S\hookrightarrow\ \mathcal{F}'\ ), \]
soit $\Phi_{\mathcal{F}',\mathcal{F}}$, et de ce qui vient d'être dit résulte que les deux foncteurs
\[
\begin{tikzcd}
\mathcal{F}' \arrow[r,hook] & \mathcal{F} \arrow[l,bend left,"\Phi_{\mathcal{F}',\mathcal{F}}"]
\end{tikzcd}
\]
\msp{page 77}
établissent des équivalences quasi-inverses l'une de l'autre
\[
\begin{tikzcd}
\mathcal{F}'\Sigma^{-1} \arrow[r,"\approx"] & \mathcal{F}\Sigma^{-1} \arrow[l,bend left,"\approx"',"\Phi_{\mathcal{F}',\mathcal{F}}"]
\end{tikzcd}
\qquad\qquad(\Sigma=W_S\text{ ou }W_S^{\mathrm{d}})
\]

Énoncé dual pour le foncteur
\[ \Psi : \Cat/S\longrightarrow\ul{Cofibsc}\,S, \]
associant à $X$ dans $\Cat/S$ la catégorie cofibrée scindée sur $S$
\[ s\longmapsto X/s. \]
Il donne des couples de foncteurs quasi-inverses
\[
\begin{tikzcd}
(\ul{Cofibsc}\,S)\Sigma^{-1} \arrow[r] & \mathcal{F}\Sigma^{-1} \arrow[l,bend left,"\Psi_{\mathcal{F}}"]
\end{tikzcd}
\qquad\qquad(\Sigma=W_S\text{ ou }W_S^{\mathrm{g}})
\]
pour tout sous-catégorie pleine $\mathcal{F}$ de $\Cat/S$ telle que % [sic] tout -> toute, accord avec sous-catégorie
\[ \ul{Cofib}\,S\subseteq\mathcal{F}\subseteq\Cat/S, \]
et des couples de foncteurs quasi-inverses


\endgroup
% END UNIT Der-ChVI-345/Der-ChVI-345_p046-050.tex
% BEGIN UNIT Der-ChVI-345/Der-ChVI-345_p051-055.tex
\begingroup\makeatletter%
\newcommand{\printedpage}[1]{\marginpar{\footnotesize p.~#1}}%
\newcommand{\msp}[1]{\marginpar{\footnotesize [#1]}}%
\newcommand{\usl}[1]{\underline{\mathrm{#1}}}%
\newcommand{\encadre}[1]{\tikz[baseline=(n.base)]\node[draw,rounded corners=3pt,inner sep=1.5pt](n){#1};}%
\makeatother\let\printedpage\derpp\let\msp\dermsp%
\let\mspin\DERmspin%


\printedpage{49}
% head: A. Grothendieck : Dérivateurs, Ch. VI, Édition des Universités de Montpellier II et Paris VII
% foot: Version 9 novembre 2001

\[
\begin{tikzcd}
\mathcal{F}'\Sigma^{-1} \arrow[r, "\approx"] & \mathcal{F}\Sigma^{-1} \arrow[l, bend left=40, "\approx"', "\Psi_{\mathcal{F}',\mathcal{F}}"]
\end{tikzcd}
\qquad\qquad (\Sigma = W_S \text{ ou } W_S^{\mathrm{g}})
\]

pour deux sous-catégories pleines $\mathcal{F}'$, $\mathcal{F}$
\[
\usl{Cofib}\,S \subseteq \mathcal{F}' \subseteq \mathcal{F} \subseteq \mathrm{Cat}/S.
\]
Les foncteurs $\Phi$, $\Psi$ envoient la sous-catégorie pleine $\usl{Fibhot}\,S$ de ‘l’étage’, dans les sous-catégories $\usl{Fibsc}_0S$, $\usl{Cofibsc}_0S$, respectivement, donc induisent à l’étage des foncteurs $\Phi_0$, $\Psi_0$
\[
\usl{Fibsc}_0S \xleftarrow{\Phi_0} \usl{Fibhot}\,S \xrightarrow{\Psi_0} \usl{Cofibsc}_0S,
\]
jouant à l’étage exactement le même rôle que $\Phi$, $\Psi$ au rez-de-chaussée. La seule différence, c’est que à l’étage on a $W_S = W_S^{\mathrm{d}} = W_S^{\mathrm{g}}$ $(= W_S^{\mathrm{f}})$,
\msp{page 78}
donc il n’y a qu’une suite de catégories localisées à examiner. On trouve ainsi
\[
\begin{tikzcd}
(\usl{Fibsc}_0S)W_S^{-1} \arrow[r, "\approx"] & \mathcal{F}W_S^{-1} \arrow[l, bend left=40, "\approx"', "\Phi_{0,\mathcal{F}}"]
\end{tikzcd}
\]
\[
\begin{tikzcd}
\mathcal{F}'W_S^{-1} \arrow[r, "\approx"] & \mathcal{F}W_S^{-1} \arrow[l, bend left=40, "\approx"', "\Phi_{0,\mathcal{F}',\mathcal{F}}"]
\end{tikzcd}
\]
pour $\mathcal{F}$ (ou $\mathcal{F}$, $\mathcal{F}'$) sous-catégorie(s) pleine(s) de $\usl{Fibhot}\,S$ telle(s) que
\[
\usl{Fib}_0S \subseteq \mathcal{F} \subseteq \usl{Fibhot}\,S \qquad\text{resp.}\qquad \usl{Fib}_0S \subseteq \mathcal{F}' \subseteq \mathcal{F} \subseteq \usl{Fibhot}\,S,
\]
et de même du côté droit pour les sous-catégories pleines $\mathcal{F}$ ou $\mathcal{F}$, $\mathcal{F}'$ de $\usl{Fibhot}\,S$ satisfaisant
\[
\usl{Cofib}_0S \subseteq \mathcal{F} \subseteq \usl{Fibhot}\,S \qquad\text{resp.}\qquad \usl{Cofib}_0S \subseteq \mathcal{F}' \subseteq \mathcal{F} \subseteq \usl{Fibhot}\,S,
\]
NB. La catégorie $\usl{Parf}\,S$ contient à la fois $\usl{Fib}_0S$ et $\usl{Cofib}_0S$, et fait donc partie des catégories admissibles tout comme $\mathcal{F}$, ou [?] $\mathcal{F}'$.

Ainsi, les quatre foncteurs $\Phi$, $\Psi$, $\Phi_0$, $\Psi_0$ donnent tous les quasi-inverses qu’on peut souhaiter au rez-de-chaussée d’une part ($\Phi$ et $\Psi$), à l’étage de l’autre ($\Phi_0$, $\Psi_0$). Ils ramènent donc toutes les localisées pertinentes aux quatre suivantes
\[
\begin{gathered}
(\mathrm{Cat}/S)W_S^{-1},\ \ (\mathrm{Cat}/S)(W_S^{\mathrm{d}})^{-1},\ \ (\mathrm{Cat}/S)(W_S^{\mathrm{g}})^{-1} \qquad \text{(au rez-de-chaussée)}\\
(\usl{Fibhot}\,S)W_S^{-1} \qquad \text{(à l’étage).}
\end{gathered}
\]

\encadre{VII} Il reste à comparer la localisée à l’étage avec les localisées au rez-de-chaussée. Les flèches verticales descendantes de l’étage au rez-de-chaussée, sont compatibles avec $W_S^{\mathrm{d}}$ du côté gauche du diagramme, avec $W_S^{\mathrm{g}}$ du côté droit, donc elles induisent essentiellement \textit{deux} foncteurs différents
\printedpage{50}
% head: A. Grothendieck : Dérivateurs, Ch. VI, Édition des Universités de Montpellier II et Paris VII
% foot: Version 9 novembre 2001
\msp{page 79}
$i$, $j$ sur $(\usl{Fibhot}\,S)W_S^{-1}$, dans $(\mathrm{Cat}/S)(W_S^{\mathrm{d}})^{-1}$ et $(\mathrm{Cat}/S)(W_S^{\mathrm{g}})^{-1}$ :
\[
\begin{tikzcd}[column sep=small, row sep=large]
& (\usl{Fibhot}\,S)W_S^{-1} \arrow[dl, "i"'] \arrow[dr, "j"] \arrow[dd, dotted, "\approx"', "\delta"] & \\
\mathrm{HOT}_W(S) \overset{\text{déf}}{\approx} (\mathrm{Cat}/S)(W_S^{\mathrm{d}})^{-1} \arrow[dr, dotted, "\rho"'] & & (\mathrm{Cat}/S)(W_S^{\mathrm{g}})^{-1} \overset{\text{déf}}{\approx} \mathrm{HOT}^{o}_{W}(S) \arrow[dl, dotted, "\sigma"] \\
& \begin{array}{l} (\mathrm{Cat}/S)W_S^{-1} \\ \overset{\text{déf}}{\approx} \mathrm{HOT}^{\mathrm{lc}}_{W}(S) \end{array} &
\end{tikzcd}
\qquad ,
\]
de façon à faire commuter le carré complété par les pointillés, représentant les foncteurs de localisation $\rho$, $\sigma$. On a bien sûr
\[
\begin{aligned}
(*) \qquad \rho i = \sigma j : (\usl{Fibhot}\,S)W_S^{-1} \longrightarrow (\mathrm{Cat}/S)W_S^{-1}\\
\text{induit par l’inclusion}\\
\delta : \usl{Fibhot}\,S \hookrightarrow \mathrm{Cat}/S
\end{aligned}
\]
Nous avons énoncé dans \encadre{I} que cette flèche est une équivalence, de sorte que $\rho$ et $\sigma$ peuvent être interprétées, à isomorphisme de foncteurs près, comme des rétractions de $\mathrm{HOT}_W(S)$ et de $\mathrm{HOT}^{o}_{W}(S)$ sur une sous-catégorie commune, s’interprétant soit comme $(\usl{Fibhot}\,S)W_S^{-1}$, soit comme $(\mathrm{Cat}/S)W_S^{-1}$. Il reste à décrire le foncteur quasi-inverse de $(*)$, ce qui donnera en même temps les quasi-inverses de tous les foncteurs verticaux (de l’étage vers le rez-de-chaussée) dans le diagramme p. 76, plus exactement pour leurs localisés, pour $W_S$.

\msp{page 80}
C’est ce qui est obtenu à l’aide du foncteur ‘montant’ (du rez-de-chaussée vers l’étage) $C$, qu’on peut regarder comme un foncteur $\mathrm{Cat}/S \longrightarrow \usl{Fibhot}\,S$, mais qui en fait prend ses valeurs dans $\usl{Parf}\,S$. (C’est là un détail sans grande importance, me semble-t-il.)
\[
\mathrm{Cat}/S \xrightarrow{C} \usl{Parf}\,S \hookrightarrow \usl{Fibhot}\,S.
\]
Ce foncteur est compatible avec les localiseurs $W_S$ des deux côtés, et établit des équivalences (quasi-inverses des équivalences
\[
\begin{gathered}
(\usl{Parf}\,S)W_S^{-1} \xrightarrow{\sim} (\mathrm{Cat}/S)W_S^{-1} \qquad\text{et}\qquad (\usl{Fibhot}\,S)W_S^{-1} \xrightarrow{\sim} (\mathrm{Cat}/S)W_S^{-1}),\\
(\mathrm{Cat}/S)W_S^{-1} \xrightarrow{\sim} (\usl{Parf}\,S)W_S^{-1} \left( \xrightarrow{\sim} (\usl{Fibhot}\,S)W_S^{-1} \right).
\end{gathered}
\]

\printedpage{51}
% head: A. Grothendieck : Dérivateurs, Ch. VI, Édition des Universités de Montpellier II et Paris VII
% foot: Version 9 novembre 2001
Donc, quand les trois catégories HOT pour $S$ sont interprétées en termes de localisées de $\mathrm{Cat}/S$, c’est le foncteur $C$ (plus exactement le composé $\mathrm{Cat}/S \xrightarrow{C} \usl{Parf}\,S \hookrightarrow \mathrm{Cat}/S$) qui fournit à la fois le foncteur d’inclusion (!) $i$, et le foncteur d’inclusion $j$ dans $\mathrm{HOT}_W(S)$ et $\mathrm{HOT}^{o}_{W}(S)$ respectivement. Le foncteur en question
\[
C' : \mathrm{Cat}/S \longrightarrow \mathrm{Cat}/S
\]
satisfait bien
\[
C'(W_S) \subseteq W_S^{\mathrm{d}} \cap W_S^{\mathrm{g}}
\]
et définit les deux foncteurs
\[
\begin{tikzcd}[column sep=large]
& (\mathrm{Cat}/S)(W_S^{\mathrm{d}})^{-1} \approx \mathrm{HOT}_W(S) \\
\begin{array}{l} (\mathrm{Cat}/S)W_S^{-1} \\ \approx \mathrm{HOT}^{\mathrm{lc}}_{W}(S) \end{array} \arrow[ur, "i"] \arrow[dr, "j"'] & \\
& (\mathrm{Cat}/S)(W_S^{\mathrm{g}})^{-1} \approx \mathrm{HOT}^{o}_{W}(S)
\end{tikzcd}
\]
Cela nous permet d’expliciter les formules d’adjonction de \encadre{IV}, quand on y interprète $\xi$, $\eta$
\msp{page 81}
comme provenant de $X$, $Y$ dans $\mathrm{Cat}/S$ :
\[
\boxed{\mathrm{Hom}_{W_S^{\mathrm{d}}}(X,CY) \simeq \mathrm{Hom}_{W_S^{\mathrm{g}}}(X,CY) \simeq \mathrm{Hom}_{W_S}(X,Y),}
\]
où $\mathrm{Hom}_{W_S^{\mathrm{d}}}$, $\mathrm{Hom}_{W_S^{\mathrm{g}}}$ et $\mathrm{Hom}_{W_S}$ désignent respectivement les Hom dans les catégories localisées pour $W_S^{\mathrm{d}}$, $W_S^{\mathrm{g}}$ et $W_S$, i.e. dans $\mathrm{HOT}_W(S)$, $\mathrm{HOT}^{o}_{W}(S)$ et $\mathrm{HOT}^{\mathrm{lc}}_{W}(S)$ respectivement (\footnote[51]{Cette formule suggère que si $X \in \mathrm{Ob}\,\mathrm{Cat}/S$, $Y \in \mathrm{Ob}\,\usl{Fibhot}\,S$, alors
\[
\mathrm{Hom}_{W_S^{\mathrm{d}}}(X,Y) \simeq \mathrm{Hom}_{W_S^{\mathrm{g}}}(X,Y) \simeq \mathrm{Hom}_{W_S}(X,Y),
\]
mais ça résulte de la formule d’adjoinction et de $Y \longrightarrow C(Y)$ dans $W_S^{\mathrm{d}} \cap W_S^{\mathrm{g}}$ pour $Y$ dans $\usl{Fibhot}\,S$.}). % [sic] adjoinction -> adjonction, spelling (elsewhere "adjonction")

\textbf{Commentaire.} Toutes les constructions et les énoncés de \encadre{I} à \encadre{VI}, à l’exception du fait que les foncteurs descendants induisent des équivalences pour les localisés \textit{grossiers}, et des formules d’adjonction de \encadre{IV}, sont de nature quasiment soritale, et n’exigent d’autres conditions sur le localiseur fondamental $W \subseteq \mathrm{Fl}\,\mathrm{Cat}$ que L1, L3 quater, L4 bis, récapitulées comme a) b) c) de la page 45. (On n’a même pas à utiliser d).) Il est vrai qu’en l’absence de la condition L5 (page 43), on ignore si à l’étage les deux localiseurs pertinents, $W_S$ et $W_S^{\mathrm{f}}$ (égal à $W_S^{\mathrm{d}}$ du côté gauche, à $W_S^{\mathrm{g}}$ du côté droit) sont égaux. Donc à priori on trouve à l’étage \textit{deux} catégories localisées essentiellement distinctes, et non pas une seule, savoir $(\usl{Fibhot}\,S)\Sigma^{-1}$ avec $\Sigma = W_S$ ou $\Sigma = W_S^{\mathrm{f}}$ $(= W_S^{\mathrm{d}} = W_S^{\mathrm{g}})$.

\printedpage{52}
% head: A. Grothendieck : Dérivateurs, Ch. VI, Édition des Universités de Montpellier II et Paris VII
% foot: Version 9 novembre 2001
Par contre, le fait que les foncteurs de descente induisent des équivalences entre les localisés grossiers est plus profond, il exige (me semble-t-il) qu’on connaisse un foncteur ‘montant’ $C$, disons de $\mathrm{Cat}/S$ vers $\usl{Fibhot}\,S$, compatible avec $W_S$, et muni de
\[
X \xrightarrow{i_X} C(X),
\]
qui soit dans $W_S$ pour $X$ quelconque. Quand on sait que Cat est une catégorie de Quillen, même sans avoir les propriétés des foncteurs $\usl{Ch}_\infty$ ou $\usl{Ch}$ (\footnote[52]{ni même être sûr de l’existence d’une factorisation \textit{fonctorielle} de $X \xrightarrow{f} S$ en $X \xrightarrow{i_X} C(X) \xrightarrow{\bar{f}} S$ avec $i_X \in W$ et $\bar{f}$ parfait.}), on doit sans doute
\msp{page 82}
pouvoir se débrouiller. Mais la formule d’adjonction me semble rester inattendue et profonde, même en admettant les résultats de \textsc{Quillen} et \textsc{Thomason} (\footnote[53]{Voir cependant le théorème scholie 2, p. 108-113. Ce n’est finalement pas si profond qu’il me semblait, mais quasiment sorital!}). Elle parait liée ici à une propriété assez particulière du foncteur $C$, savoir le lemme p. 68
\[
\begin{gathered}
i_{C(X)} \underset{\big\uparrow}{=} C(i_X) : C(X) \longrightarrow C(C(X)),\\
\scriptstyle \text{dans } \underbrace{(\mathrm{Cat}/S)(W_S^{\mathrm{d}})^{-1}}_{\mathrm{HOT}_W(S)} \text{ et dans } \underbrace{(\mathrm{Cat}/S)(W_S^{\mathrm{g}})^{-1}}_{\mathrm{HOT}^{o}_{W}(S)}
\end{gathered}
\]
ce qui n’a pas l’air purement formel. Il serait intéressant d’examiner dans quelle mesure cette formule garde un sens dans une ‘catégorie de modèles’ $(\mathcal{M},W)$ plus ou moins quelconque, ou provenant d’un triple de Quillen $(W,\mathrm{Cof},\mathrm{Fib})$ dans $\mathcal{M}$. Il est clair que la factorisation $X \xrightarrow{f} S$ en $X \xrightarrow{i_X} C(X) \xrightarrow{\bar{f}} S$ $(i_X \in W,\ \bar{f} \in \mathrm{Parf})$, jouant le rôle des ‘fibrations’ dans le contexte Quillen, a un sens intrinsèque indépendant du choix d’une construction particulière, chose qu’on peut préciser de façon parfaite. C’est la description des localiseurs $W_S^{\mathrm{d}}$, $W_S^{\mathrm{g}}$ dans $\mathcal{M}/S$ (pour $S \in \mathrm{Ob}\,\mathcal{M}$) qui n’est pas comprise de façon générale, en dehors du cas $\mathcal{M} = \mathrm{Cat}$. En d’autre termes, je ne comprends pas comment construire des catégories $\mathrm{HOT}_{\mathcal{M},W}(S)$, $\mathrm{HOT}^{o}_{\mathcal{M},W}(S)$, même si $(\mathcal{M},W)$ provient d’un triple de Quillen, je vois seulement la description d’une catégorie $\mathrm{HOT}^{\mathrm{lc}}_{\mathcal{M},W}(S) = (\mathcal{M}/S)(W_S^{\text{grossière}})^{-1}$, et un foncteur de localisation $\mathcal{M}/S \longrightarrow \mathrm{HOT}^{\mathrm{lc}}_{\mathcal{M},W}(S)$. Il faut arriver à décrire dans $\mathcal{M}/S$ des $W$-équivalences (\textit{‘relatives à $S$’}) plus fines que la grossière donnée par $W$ sur la source $X$ sans plus. Dans Cat j’y arrive en utilisant soit la famille des isomorphismes locaux (pour décrire $W_S^{\mathrm{d}}$), soit celle des isomorphismes colocaux (pour décrire $W_S^{\mathrm{g}}$) - en tous les cas, il y a une notion de ‘localisation’ (soit sur $S$, soit sur $S^{o}$ \dots), avec % [sic] En d’autre termes -> En d’autres termes, fixed expression requires plural
\msp{page 83}
toute la richesse des intuitions liées à une telle notion. Peut-être, sans le savoir ai-je toujours traité Cat non comme une catégorie ‘nue’, mais comme un site, ou plutôt (puisqu’il
\printedpage{53}
% head: A. Grothendieck : Dérivateurs, Ch. VI, Édition des Universités de Montpellier II et Paris VII
% foot: Version 9 novembre 2001
y a l’involution $X \longmapsto X^{o}$ qui ne peut être éludée), comme un ‘bi-site’ - on y a \textit{deux} manières de localiser, soit par les $X/x$ (topologie de Zariski!) soit par les $x\backslash X$ (topologie co-Zariskienne). Il y a aussi la topologie lisse (toute proche sans doute de l’intuition de la cohomologie étale), et la topologie propre - et on s’attend que la lisse donne essentiellement les mêmes résultats que la Zariskienne (vue que la topologie étale lui est égale - il n’y a pas lieu d’en faire un plat \dots); et la topologie propre les mêmes résultats (p. ex. la même cohomologie?) que la topologie co-Zariskienne. % [sic] vue que -> vu que, conjunctional locution

Il semble bien que dans toutes les catégories de modèles utilisées jusqu’à présent pour décrire Hot comme une localisée, on dispose de notions de localisation : (Top), $\boldsymbol{\Delta}^\wedge$ (ou $A^\wedge$, avec $A$ catégorie-test quelconque), (Cat); et dans Cat on dispose même de \textit{deux} telles notions. Je pense par contre aux $\infty$-groupoïdes, c’est peut-être la seule catégorie de modèles à tel point proche déjà de Hot elle-même, qu’on n’y discerne pas de notion de localisation, et que l’idée même de vouloir s’y localiser (si ce n’est en ‘localisant’ Hot lui-même) paraît quasiment
\msp{page 84}
saugrenue.

Cette rapide digression ne convainc à présent que la compatibilité dite pour la construction de $C(X)$ et $i_X : X \longrightarrow C(X)$, ne peut en aucun cas être purement formelle, une pure histoire de catégories de Quillen disons. Ça donnera au mieux la compatibilité dans $(\mathcal{M}/S)W_S^{-1}$, ce qui est ridicule! Mais si on prend comme localiseur $\Sigma$ dans $\mathcal{M}/S$ les isomorphismes sans plus, la relation de compatibilité va devenir fausse à tous les coups! Comme ensemble de flèches $\Sigma$ un peu plus raisonnable, qui soit encore très petit, on peut songer p. ex. à la $W$-équivalence \textit{universelle} sur $S$ (nettement plus fine encore que $W_S^{\mathrm{d}}$, $W_S^{\mathrm{g}}$ dans le cas particulier de Cat). Aurait-on encore compatibilité pour cet ensemble de flèches? Je suis aussi un peu perplexe par le fait (pas bien compris) que les formules d’adjonction pour les topologies Zariskienne et co-Zariskienne soient en sens opposé. (Ou me serai-je trompé en dualisant à la va-vite?) Mais à vrai dire, cette compatibilité à elle seule, même quand on utilise une structure de site sur $\mathcal{M}$, ne suffit pas à écrire une formule d’adjonction. On a deux catégories localisées, % [sic] ne convainc -> me convainc, the verb needs an object pronoun
\[
\mathrm{HOT}^{\mathrm{lc}}_{\mathcal{M}} = (\mathcal{M}/S)W_S^{-1}, \qquad \mathrm{HOT}_W(S) = (\mathcal{M}/S)(W_S^{\mathrm{f}})^{-1},
\]
où $W_S^{\mathrm{f}}$ désigne une $W$-équivalence ‘fine’, correspondant à la notion de localisation, plus exactement aux changements de base ‘admissibles’ qu’on a choisis, et on a un foncteur de localisation
\msp{page 85}
\[
\mathrm{Hot}_W(S) \longrightarrow \mathrm{Hot}^{\mathrm{lc}}_{W}(S),
\]
mais pas de foncteur en sens inverse. Mais si, mais si, en travaillant avec la catégorie similaire $\usl{Fibhot}\,S$, ça doit garder un sens tout au moins quand on a une bonne notion de ‘fibration’. Je vois un petit travail de dégrossissage conceptuel à faire, avec peut-être une formule d’adjonction à la clef, mais je ne vais pas m’y lancer à présent.


\endgroup
% END UNIT Der-ChVI-345/Der-ChVI-345_p051-055.tex
% BEGIN UNIT Der-ChVI-345/Der-ChVI-345_p056-060.tex
\begingroup\makeatletter%
\newcommand{\printedpage}[1]{\marginpar{\footnotesize p.~#1}}%
\newcommand{\msp}[1]{\marginpar{\footnotesize [#1]}}%
\newcommand{\cM}{\mathcal{M}}%
\newcommand{\cF}{\mathcal{F}}%
\newcommand{\Fib}{\underline{\mathrm{Fib}}\,}%
\newcommand{\Fibhot}{\underline{\mathrm{Fibhot}}\,}%
\newcommand{\Ho}{\mathrm{Ho}}%
\newcommand{\Fl}{\mathrm{Fl}}%
\newcommand{\Ob}{\mathrm{Ob}}%
\makeatother\let\printedpage\derpp\let\msp\dermsp%
\let\mspin\DERmspin%


\printedpage{54}
% head: A. Grothendieck : Dérivateurs, Ch. VI, Édition des Universités de Montpellier II et Paris VII

\begin{flushright}
16.11.90
\end{flushright}

Je vais quand même essayer de faire un travail ‘de dégrossissage’, car le diagramme des
15 catégories a l'air un peut touffu, % [sic] peut -> peu, adverb "un peu" intended
finalement! Je me donne
\[
(\cM,W) \qquad W\subseteq \Fl\,\cM,
\]
un $S\in\Ob\,\cM$, d'où $\cM/S$, avec $W_S\subseteq\Fl(\cM/S)$ induit par $W$. On va poser
\[
\Ho^{\mathrm{lc}}(S)=(\cM/S)W_S^{-1},
\]
ça jouera le rôle des types d'homotopie relativement constants. Je suppose d'autre part
que l'on a distingué une classe $F\subseteq\Fl(\cM)$ de flèches, jouant le rôle de ‘fibrations’ (ce
seront les morphismes Hot-fibrants dans Cat, donnant lieu à la sous-catégorie $\Fibhot S$
de $\mathrm{Cat}/S$). Je prévois qu'il faudra
\msp{page 86}supposer, outre que $\cM$ est stable par $\varprojlim$ finie, que $F$ est stable par changement de base
quelconque, et stable par composition, qu'il contient les isomorphismes. Et surtout qu'on
a un théorème de factorisation
\[
\begin{tikzcd}
X \arrow[rr, "f"] \arrow[dr, "i_X"'] & & S \\
& \bar{X} \arrow[ur, "\bar{f}"'] &
\end{tikzcd}
\qquad\qquad\qquad
i_X\in W,\ \bar{f}\in F
\]
pout % [sic] pout -> pour, misprint
toute flèche $f$ dans $\cM$. Chez \textsc{Quillen} on exige même que $i_X$ ait la propriété de
relèvement à droite relativement à $F$, mais ça ne devrait pas être essentiel, puisque je n'ai
pas eu à utiliser ce fait dans Cat. Quand je travaille dans la catégorie $\cM/S$, où $f$ est
sous-entendu pour un objet $X$ de ladite, je noterai (par abus de langage) $C(X)$ pour $\bar{X}$,
bien qu'il n'est pas dit qu'on ait un choix canonique de $C(X)$, encore moins que $C(X)$
dépende fonctoriellement de $X$.

Désignons par
\[
\Fib S\subseteq\cM/S \tag*{$(*)$}
\]
la sous-catégorie de $\cM/S$ formée des $X$ dont le morphisme structural $f_X:X\longrightarrow S$ est
dans $F$. Désignons (par abus de notation) par $W_S$ le localiseur induit sur $\Fib S$ par $W_S$,
alors l'inclusion $(*)$ définit
\[
(\Fib S)W_S^{-1}\xrightarrow{\ \delta\ }(\cM/S)W_S^{-1}\overset{\text{déf}}{=}\Ho^{\mathrm{lc}}_W(S).
\]

\msp{page 87}Une première question, c'est de voir que c'est une équivalence. On voit que c'est essentiellement
surjectif grâce à l'existence des $C(X)$, $i_X$. C'est peu! On va essayer de définir
une flèche en sens inverse
\[
\bar{C}:(\cM/S)W_S^{-1}\longrightarrow(\Fib S)W_S^{-1},
\]
% foot: Version 9 novembre 2001
\printedpage{55}
% head: A. Grothendieck : Dérivateurs, Ch. VI, Édition des Universités de Montpellier II et Paris VII
i.e. un \textit{foncteur}
\[
\gamma:\cM/S\longrightarrow(\Fib S)W_S^{-1}
\]
qui transforme $W_S\subseteq\Fl(\cM/S)$ en isomorphismes. On choisira, pour $X$ dans $\cM/S$, un
$C(X)$ dans $\Fib S$ et une
\[
i_X:X\longrightarrow C(X)\ , \qquad i_X\in W_S,
\]
et on posera
\[
\gamma(X)=C(X).
\]
La difficulté, c'est de montrer que ça ne dépend pas des choix faits, à isomorphisme
canonique près, et que c'est fonctoriel en $X$ - et le deuxième implique le premier.

Pour voir que ça marche, visiblement il faut supposer quelque chose de plus. Je vois deux
hypothèses qui, l'une et l'autre, pourraient faire marcher : l'une, c'est qu'on se donne
déjà $C$ comme foncteur et $i_X$ fonctoriel en $X$ (c'est ce qui était le cas dans Cat tantôt).
L'autre, c'est que
\msp{page 88}les $i_X$ ont, par rapport à $F$, la propriété de relèvement à là Quillen. % [sic] à là -> à la, accent on article "la"

Je commence par l'hypothèse que $C$ est déjà un foncteur. Le diagramme
\[
\begin{tikzcd}[row sep=large]
X \arrow[r, "i_X"] \arrow[d, "u"'] & CX \arrow[d, "C(u)"] \\
X' \arrow[r, "i_{X'}"] & CX',
\end{tikzcd}
\]
où $i_X,\ i_{X'}\in W$, montre que $C$ transforme $W_S$ en $W_S$, donc en des isomorphismes de
$\Ho^{\mathrm{lc}}_W(S)$. Donc on tient le foncteur cherché $\gamma$. On a en fait des foncteurs
\[
\begin{tikzcd}
\Fib S \arrow[r, hook, "\overset{\text{incl.}}{D}"] & \cM/S \arrow[l, bend left=50, "C"]
\end{tikzcd}
\]
compatibles avec les localiseurs $W_S$, induisant $\delta$ et $\gamma$, et on a
\[
\mathrm{id}_{\cM/S}\longrightarrow DC, \qquad \text{dans $W_S$ argument par argument,}
\]
car on a
\[
X\xrightarrow{\ i_X\ }DC(X)=C(X)
\]
avec $i_X\in W_S$. De même, on a
\[
\mathrm{id}_{\Fib S}\longrightarrow CD \qquad \text{dans $W_S$ argument par argument,}
\]
% foot: Version 9 novembre 2001
\printedpage{56}
% head: A. Grothendieck : Dérivateurs, Ch. VI, Édition des Universités de Montpellier II et Paris VII
avec le même
\[
X\xrightarrow{\ i_X\ }CX \qquad \text{(cette fois $X$ dans $\Fib S$).}
\]

\msp{page 89}Donc on a sans larmes les deux foncteurs $\delta$, $\gamma$ comme quasi-inverses l'un de l'autre. Et le
lemme page 11 montre même que pour \textit{toute} sous-catégorie pleine $\cF$ de $\cM/S$, telle que
\[
\Fib S\subseteq\cF\subseteq\cM/S,
\]
les inclusions induisent des équivalences sur les localisés $W_S$
\[
(\Fib S)W_S^{-1}\xrightarrow{\ \sim\ }\cF W_S^{-1}\xrightarrow{\ \sim\ }\underbrace{(\cM/S)W_S^{-1}}_{\overset{\text{déf}}{=}\Ho^{\mathrm{lc}}_W(S)}.
\]
Je vais traiter maintenant le cas de type Quillen, pour me convaincre que l'existence du
foncteur $C$, et le choix particulier qu'on peut en faire, est irrelevant. (Ce dernier point
est d'ailleurs évident a posteriori, quand on sait que $\delta$ est une équivalence, car le quasi-inverse
$\gamma$ sera alors unique à isomorphisme unique près, et le choix d'un foncteur $C$, et
d'un homomorphisme fonctoriel $i:\mathrm{id}\longrightarrow C$, n'est qu'une façon parmi d'autres d'expliciter
le quasi-inverse.)

\msp{page 90}On se place donc dans l'hypothèse quillénienne sur les $i_X$. On choisit $C(X)$, $i_X$ pour
chaque $X$ dans $\cM/S$, et on pose comme tantôt
\[
\gamma(X)=C(X) \qquad \text{comme objet de $(\cM/S)W_S^{-1}$.}
\]
Montrons que c'est fonctoriel en $X$. Soit
\[
u:X\longrightarrow X'
\]
dans $\cM/S$, considérons le diagramme anodin (commutatif)
\[
\begin{tikzcd}[row sep=large]
X \arrow[rr, "u"] \arrow[d, "i_X"'] \arrow[drr, loosely dotted, "v=i_{X'}u"] & & X' \arrow[d, "i_{X'}"] \\
C(X) \arrow[rr, dotted, "\bar{u}"] \arrow[dr] & & C(X') \arrow[dl, "\in F"] \\
& S, &
\end{tikzcd}
\]
complétons le par la flèche en gros pointillés, $v=i_{X'}u$ (rendant le triangle supérieur
commutatif). (\footnote[54]{Lemme : soit
% footnote continues on printed p. 57
\[
\begin{tikzcd}[ampersand replacement=\&, row sep=large]
X \arrow[r, "u"] \arrow[d, "i"'] \& X' \arrow[d] \\
C \arrow[r, dotted] \& C'
\end{tikzcd}
\]
dans $\cM/S$ avec $i\in\mathrm{Cof}_0$, $(C'\longrightarrow S)\in F=\mathrm{Fib}$, alors la flèche en pointillés existe dans $\cM/S$, de façon
à rendre le carré commutatif. \texttt{[cf. Lemme p. 91]}}). L'hypothèse sur $i_X$ implique que la flèche $\bar{u}$ en pointillés fins existe,
de façon à avoir encore deux triangles commutatifs. Mais ça signifie exactement qu'il y a
un $S$-\textit{morphisme} $\bar{u}$, i.e. une flèche dans $\cM/S$, rendant commutatif le carré
% foot: Version 9 novembre 2001
\printedpage{57}
% head: A. Grothendieck : Dérivateurs, Ch. VI, Édition des Universités de Montpellier II et Paris VII
\[
\begin{tikzcd}[row sep=large]
X \arrow[r, "u"] \arrow[d, "i_X"'] & X' \arrow[d, "i_{X'}"] \\
C(X) \arrow[r, "\bar{u}"] & C(X'),
\end{tikzcd}
\]
\msp{page 91}et il s'impose de prendre
\[
\gamma(u)=\ \text{image de $\bar{u}$ dans $\Fl((\Fib S)W_S^{-1})$.}
\]
Il faut prouver que cette flèche ne dépend pas du choix de $\bar{u}$. On doit dégager un Lemme,
où on désigne par $\mathrm{Cof}_0$ l'ensemble des flèches de $\cM$ qu'ont % [sic] qu'ont -> qui ont, relative pronoun
la propriété de relèvement par
rapport aux éléments de $\mathrm{Fib}=F$. ($\mathrm{Cof}_0$ joue le rôle des cofibrations triviales.)

\textbf{Lemme.} \textit{Soit un diagramme en traits pleins dans $\cM/S$}
\[
\begin{tikzcd}[row sep=large]
X \arrow[r, "u"] \arrow[d, "i"'] & X' \arrow[d, "i'"] \\
C \arrow[r, dotted, "v"] & C'
\end{tikzcd}
\]
\textit{avec $i\in$ Cof$_0$, et $C'\in$ Ob$\,\Fib S$ i.e. $(C'\longrightarrow S)\in$ Fib. Alors $v$ existe, de façon à rendre
le carré commutatif. De plus, si $i\in W$ et $C\in\Ob\,\Fib S$, la classe de $v$ comme flèche de
$(\Fib S)W_S^{-1}$ en est indépendante du choix fait.}

Il reste à prouver cette indépendance,
\msp{page 92}ce qui résultera aussitôt du corollaire :

\textbf{Corollaire.} \textit{Soit un diagramme dans $\cM/S$}
\[
\begin{tikzcd}
X \arrow[r, "i"] & C \arrow[r, shift left, "v_1"] \arrow[r, shift right, "v_2"'] & C',
\end{tikzcd}
\]
\textit{avec $C$ et $C'$ dans $\Fib S$ (i.e. $C\longrightarrow S$ et $C'\longrightarrow S$ dans Fib), et $i\in$ Cof$_0$ et $i\in W$. Si
$v_1i=v_2i$, alors $v_1$ et $v_2$ définissent la même flèche de $(\Fib S)W_S^{-1}$.}
% foot: Version 9 novembre 2001

\printedpage{58}
% head: A. Grothendieck : Dérivateurs, Ch. VI, Édition des Universités de Montpellier II et Paris VII
\textsc{Démonstration}. Soit
\[
v=(v_1,v_2):C\longrightarrow C'\times C',
\]
\texttt{[plutôt} $C'\times_S C'$\texttt{]} qui rend commutatif par hypothèse le diagramme
\[
\begin{tikzcd}[row sep=large]
X \arrow[r, "w"] \arrow[d, "i"'] & C' \arrow[d, "\mathrm{diag}_{C'}"] \\
C \arrow[r, "v"] & C'\times_S C',
\end{tikzcd}
\]
où $w=v_1i=v_2i$. Factorisons $\mathrm{diag}_{C'}$ en
\[
C'\xrightarrow{\ i'\ }\bar{C}'\xrightarrow{\ \delta'\ }C'\times_S C', \qquad \text{avec $\delta'\in\mathrm{Fib}$,}
\]
et $i'\in W$. On trouve un carré commutatif
\[
\begin{tikzcd}[row sep=large]
X \arrow[r, "w'=i'w"] \arrow[d, "i"'] & \bar{C}' \arrow[d, "\delta'"] & \\
C \arrow[r, "v"] \arrow[ur, dotted, "\bar{v}"] & C'\times_S C' \arrow[r, shift left, "\mathrm{pr}_1"] \arrow[r, shift right, "\mathrm{pr}_2"'] & C'
\end{tikzcd}
\]
avec $i\in\mathrm{Cof}_0$, $\delta'\in\mathrm{Fib}$, d'où l'existence de la flèche en pointillés, rendant les deux triangles
commutatifs. On veut prouver $\mathrm{pr}_1\circ v=\mathrm{pr}_2\circ v$ dans $(\Fib S)W_S^{-1}$, comme $v=\delta'\bar{v}$ il suffit
de prouver que l'on a $\mathrm{pr}_1\delta'=\mathrm{pr}_2\delta'$ dans $(\Fib S)W_S^{-1}$.

\msp{page 93}Il faut seulement vérifier d'abord (pour justifier l'argument) que cette assertion a un sens,
i.e. que $\bar{C}'$ est lui aussi dans $\Fib S$. Mais les hypothèses de stabilité de $F=\mathrm{Fib}$ par
changement de base et composition assurent que $1^o$) $C'\times_S C'$ est dans $\Fib S$ (qui est
stable par produit), et par suite, $2^o$) que $\bar{C}'$ l'est aussi, puisque $\delta'\in\mathrm{Fib}$. Comme $C'$ est
aussi dans $\Fib S$ et $C'\xrightarrow{\ i'\ }\bar{C}'$ est dans $W_S$, pour vérifier $\mathrm{pr}_1\delta'=\mathrm{pr}_2\delta'$ dans le localisé de
$\Fib S$, il suffit de voir que $(\mathrm{pr}_1\delta')i'=(\mathrm{pr}_2\delta')i'$, or comme $\delta'i'=\mathrm{diag}_{C'}$, on a l'égalité même
dans $\Fib S$ (sans localiser).

On a donc défini $\gamma(X)$, $\gamma(u)$ pour objets et flèches de $\cM/S$. La transitivité pour $\gamma(u)$ est
claire, et aussi $\gamma(\mathrm{id})=\mathrm{id}$. Donc on tient le foncteur
\[
\gamma:\cM/S\longrightarrow(\Fib S)W_S^{-1}.
\]
Soit maintenant $\cF$ une sous-catégorie pleine de $\cM/S$, avec
\[
\Fib S\subseteq\cF\subseteq\cM/S,
\]
considérons $\gamma$ comme foncteur à valeurs dans $\cF W_S^{-1}$. (J'ai oublié de vérifier que $\gamma$ transforme
$W_S$ dans $W_S$, mais c'est tautologique.) Ainsi on a des foncteurs
% foot: Version 9 novembre 2001


\endgroup
% END UNIT Der-ChVI-345/Der-ChVI-345_p056-060.tex
% BEGIN UNIT Der-ChVI-345/Der-ChVI-345_p061-065.tex
\begingroup\makeatletter%
\NoAutoSpacing%
\newcommand{\printedpage}[1]{\marginpar{\footnotesize p.~#1}}%
\newcommand{\msp}[1]{\marginpar{\footnotesize [#1]}}%
\newcommand{\Fibu}{\underline{\mathrm{Fib}}\,}%
\newcommand{\Ho}{\mathrm{Ho}}%
\newcommand{\WL}{W_S^{\mathcal{L}}}%
\newcommand{\defeq}{\overset{\text{déf}}{=}}%
\makeatother\let\printedpage\derpp\let\msp\dermsp%
\let\mspin\DERmspin%


\printedpage{59}
% head: A. Grothendieck : Dérivateurs, Ch. VI, Édition des Universités de Montpellier II et Paris VII
\[
\begin{tikzcd}
\mathcal{F}W_S^{-1} \arrow[r, shift left, "\delta"] & (\mathcal{M}/S)W_S^{-1}, \arrow[l, shift left, "\gamma"]
\end{tikzcd}
\]
il faut voir qu'ils sont quasi-inverse l'une de \msp{page 94} l'autre.

Du côté de $\mathcal{M}/S$ : on doit comparer $X$ et $\delta\gamma(X) = C(X)$, c'est fait par
\[
i_X : X \longrightarrow C(X)
\]
qui est dans $W_S$, d'où l'isomorphisme dans le localisé
\[
X \simeq \delta\gamma(X) \qquad \text{dans } (\mathcal{M}/S)W_S^{-1}.
\]
Il faut vérifier que c'est fonctoriel en $X$.

Du côté de $\mathcal{F}$, comparer $X$ et $\gamma\delta(X)$, i.e. $X$ et $C(X)$. Même tabac, on choisit $i_X : X \simeq C(X) = \gamma\delta(X)$. À vérifier la fonctorialité.

Il s'agit maintenant d'introduire un localisé plus fin de $\mathcal{M}/S$. Pour ceci, on suppose donné une autre classe de morphismes, je l'appelle $\mathcal{L}$ (comme `localisation')
\[
\mathcal{L} \subseteq \mathrm{Fl}(\mathcal{M}),
\]
p. ex. dans le cas de $\mathcal{M} = \mathrm{Cat}$, $\mathcal{L}$ sera formée des isomorphismes locaux (\footnote[55]{Ou aussi, des morphismes \textit{lisses}. Ça doit donner le même résultat.}) (s'il s'agit de définir $W_S^{\mathrm{g}}$), ou des isomorphismes colocaux (\footnote[56]{Ou aussi des morphismes propres.}) (s'il s'agit de définir $W_S^{\mathrm{d}}$). On supposera que $\mathcal{L}$ contient les isomorphismes, et $\mathcal{L}$ stable par composition (\footnote[57]{Il n'en est peut-être pas besoin.}), mais non par changement de base. On définit alors
\[
\WL \subseteq \mathrm{Fl}(\mathcal{M}/S)
\]
par
\[
f \in \WL \overset{\text{déf}}{\Longleftrightarrow} (f \times_S S') \in W,\ \ \forall (S' \longrightarrow S) \in \mathcal{L}.
\]
\msp{page 95}

Comme $\mathcal{L}$ contient $\mathrm{id}_S$, on aura
\[
\WL \subseteq W_S,
\]
d'où un foncteur canonique
\[
\underbrace{(\mathcal{M}/S)(\WL)^{-1}}_{\defeq \Ho_W^{\mathcal{L}}(S)} \xrightarrow{\ \rho\ } \underbrace{(\mathcal{M}/S)(W_S)^{-1}}_{\defeq \Ho_W^{\mathrm{lc}}\mathcal{L}(S)} % [sic] \Ho_W^{\mathrm{lc}}\mathcal{L}(S) -> \Ho_W^{\mathrm{lc}}(S), cf. the definition \Ho_W^{\mathrm{lc}}(S) \defeq (\mathcal{M}/S)W_S^{-1} on p. 61
\]
% foot: Version 9 novembre 2001
\printedpage{60}
% head: A. Grothendieck : Dérivateurs, Ch. VI, Édition des Universités de Montpellier II et Paris VII
qui est un foncteur de localisation. Je voudrais définir un foncteur en sens inverse. Pour ceci, il est naturel de supposer (en analogie avec la situation de Cat)

\textbf{Hypothèse 1.} \textit{Sur $\Fibu S$, $W_S = \WL$, i.e. si $f : X \longrightarrow Y$ est dans $W$ et $X$, $Y$ dans $\Fibu S$, alors pour tout changement de base $S' \longrightarrow S$ qui est dans $\mathcal{L}$, $f' : X_{S'} \longrightarrow Y_{S'}$ est encore dans $W$.}

Dans le contexte de Quillen, cette hypothèse a une nette tendance à être vérifiée quelque soit $\mathcal{L}$, i.e. pout \textit{tout} changement de base $S' \longrightarrow S$ sans restriction. Il suffit en fait que la structure de Quillen soit `propre', plus précisément (et moins restrictivement), qu'un \texttt{[morphisme]} $\in W$ reste tel après changement de base par un morphisme $\in$ Fib. Dégageons ceci comme un Lemme. % [sic] pout -> pour, misprint

\textbf{Lemme.} \textit{L'hypothèse est vérifiée si on suppose} \msp{page 96} \textit{que $F =$ Fib et $W$ sont reliés par la condition supplémentaire (en plus de la factorisation $i_X$) : pour tout carré cartésien}
\[
\begin{tikzcd}
X \arrow[d, "f"'] & X' \arrow[l, "q"'] \arrow[d, "f'"] \\
S & S', \arrow[l, "p"']
\end{tikzcd}
\]
\textit{si $f$ est dans Fib et $p$ dans $W$, alors $q \in W$.}

\textsc{Démonstration}. Soit
\[
\begin{tikzcd}
X \arrow[rr, "f"] \arrow[dr, "p"'] & & Y \arrow[dl, "q"] \\
& S &
\end{tikzcd}
\]
avec $p$, $q \in \mathrm{Fib}$, $f \in W$, montrons que pour \textit{tout} changement de base $g : S' \longrightarrow S$, $f_{S'} : X_{S'} \longrightarrow Y_{S'}$ est encore dans $W$. On factorise $g$ en
\[
\begin{tikzcd}
S' \arrow[r, "i"] \arrow[rr, bend right=15, "g"'] & \bar{S}' \arrow[r, "\bar{g}"] & S,
\end{tikzcd}
\qquad\qquad i \in W,\ g \in \mathrm{Fib},
\]
\texttt{[plutôt} $\bar{g} \in$ \texttt{Fib]} et on est ramené à prouver que l'assertion est vraie dans les deux cas particuliers où $g \in \mathrm{Fib}$, ou $g \in W$. Le cas $g \in \mathrm{Fib}$ résulte aussitôt de l'hypothèse du Lemme (appliquée au carré cartésien
\[
\begin{tikzcd}
X' \arrow[r, "f'"] \arrow[d, "q"'] & Y' \arrow[d, "p"] \\
X \arrow[r, "f"] & Y,
\end{tikzcd}
\]
% foot: Version 9 novembre 2001
\printedpage{61}
% head: A. Grothendieck : Dérivateurs, Ch. VI, Édition des Universités de Montpellier II et Paris VII
où $p \in \mathrm{Fib}$.)

Dans le cas où $S' \longrightarrow S$ est dans $W$, il s'ensuit par l'hypothèse du Lemme appliquée à
\[
\begin{tikzcd}
Y \arrow[d, "\in\mathrm{Fib}"] & Y' \arrow[l, "p"'] \arrow[d] \\
S & S' \arrow[l]
\end{tikzcd}
\]
que \msp{page 97} $p : Y' \longrightarrow Y$ est $\in W$, et de même pour $q : X' \longrightarrow X$, et $f \in W$ implique encore $f' \in W$.

Revenons à l'hypothèse p. 95. On aura donc un foncteur
\[
(\Fibu S, W_S = \WL) \longrightarrow (\mathcal{M}/S, \WL),
\]
d'où pour les localisés un foncteur horizontal
\[
\begin{tikzcd}
(\Fibu S)W_S^{-1} \arrow[r, "i_0"] \arrow[d, "\approx"] & (\mathcal{M}/S)(\WL)^{-1} \defeq \Ho_W^{\mathcal{L}}(S) \arrow[dl] \\
\Ho_W^{\mathrm{lc}}(S) \defeq (\mathcal{M}/S)W_S^{-1} &
\end{tikzcd}
\]
qui factorise le foncteur vertical induit par l'inclusion $(\Fibu S, W_S) \longrightarrow (\mathcal{M}/S, W_S)$ - foncteur dont on a vu (dans les deux cas envisagés) que c'est une équivalence. En inversant ce dernier et le composant avec $i_0$, on trouve donc un foncteur
\[
\underbrace{\Ho_W^{\mathrm{lc}}(S)}_{\defeq (\mathcal{M}/S)W_S^{-1}} \quad \xrightarrow{\ i\ } \quad \underbrace{\Ho_W^{\mathcal{L}}(S)}_{\defeq (\mathcal{M}/S)(\WL)^{-1}} \qquad ,
\]
tel que
\[
\rho i \simeq \mathrm{id},
\]
donc une \textit{rétraction} de $\Ho_W^{\mathcal{L}}(S)$ sur $\Ho_W^{\mathrm{lc}}(S)$. On s'attend que $i$ soit pleinement fidèle, et plus précisément encore, \msp{page 98} que $\rho$ et $i$ soient adjoints, $i$ étant l'adjoint à droite de $\rho$ :
\[
\begin{tikzcd}
\Ho_W^{\mathrm{lc}}(S) \arrow[r, shift right, "i"'] & \Ho_W^{\mathcal{L}}(S), \arrow[l, shift right, "\rho"']
\end{tikzcd}
\]
% foot: Version 9 novembre 2001
\printedpage{62}
% head: A. Grothendieck : Dérivateurs, Ch. VI, Édition des Universités de Montpellier II et Paris VII
i.e. on cherche
\[
\begin{aligned}
i\rho &\quad \underbrace{\xleftarrow{\ i\ }}_{\textit{pas}\ \text{iso}} \quad \mathrm{id} \quad \text{dans } \Ho_W^{\mathcal{L}} \quad (\text{localisation par } \WL)\\
\rho i &\quad \underbrace{\xrightarrow{\ p\ }}_{\text{iso}} \quad \mathrm{id} \quad \text{dans } \Ho_W^{\mathrm{lc}} \quad (\text{localisation par } W_S),
\end{aligned}
\]
\texttt{[trop de} $i$\texttt{'s]} donc pour $X$ dans $\mathcal{M}/S$, on cherche
\[
\begin{aligned}
i_X :&\quad X \quad \longrightarrow \quad C(X) \quad \text{tout au moins dans } (\mathcal{M}/S)(\WL)^{-1}\\
p_X :&\quad X \quad \longleftarrow \quad C(X) \quad \text{tout au moins dans } (\mathcal{M}/S)(W_S)^{-1}.
\end{aligned}
\]
On choisira pour $i_X$ la flèche `canonique' de $X$ dans $C(X)$, définie dans $\mathcal{M}/S$ lui-même, sans avoir à localiser. Mais dans le `cas Quillen' (quand $C$ n'est pas donné comme foncteur $\mathcal{M}/S \longrightarrow \Fibu S$), il faut voir que cet $i_X$ est fonctoriel en $X$, quand on le regarde dans le localisé $(\mathcal{M}/S)W_S^{-1}$. Mais c'est contenu dans la construction de $C$ comme \textit{foncteur}
\[
\mathcal{M}/S \longrightarrow (\Fibu S)W_S^{-1}.
\]
On notera que $i_X$ est dans $W_S$ sans plus, pas dans $\WL$, donc ne définit pas en général un isomorphisme de $\Ho_W^{\mathcal{L}}$.

On définit $p_X$ à l'aide de $i_X$, comme l'inverse de $i_X$ \textit{dans} $(\mathcal{M}/S)W_S^{-1}$, qui existe car \msp{page 99} $i_X$ est dans $W_S$, donc \texttt{[est un]} isomorphisme dans $(\mathcal{M}/S)W_S^{-1} = \Ho_W^{\mathrm{lc}}$, et $p$ est donc \texttt{[un]} isomorphisme. On retrouve donc $\rho i \simeq \mathrm{id}$, donc que $\rho$ est une \textit{rétraction} de $\Ho_W^{\mathcal{L}}$ sur $\Ho_W^{\mathrm{lc}}$.

Pour prouver l'adjonction, en procédant comme dans le cas de Cat, on est ramené à prouver que pour tout $X$,
\[
\left\{
\begin{aligned}
&C(i_X) = i_{C(X)} : C(X) \rightrightarrows C(C(X))\\
&\textit{dans } \Ho_W^{\mathcal{L}} \ (\text{localisation par } \WL).
\end{aligned}
\right.
\]
J'y étais arrivé dans Cat, en utilisant un morphisme
\[
C(C(X)) \xrightarrow{\ c\ } C(X)
\]
(essentiellement, un morphisme de composition des chemins) tel que $c \in \WL$, et que $c \circ C(i_X) = c \circ i_{C(X)}$. Ça avait l'air très délicat, car c'était à grand peine que j'étais arrivé à construire un foncteur \underline{Ch} (`chemins') raisonnable, où on puisse bel et bien composer les chemins! Je renonce pour le moment à axiomatiser les propriétés du foncteur \underline{Ch} dont j'ai besoin pour faire marcher la machine (je devrais y revenir tantôt). Je vais commencer plutôt par le `cas Quillen', avec des $i_X \in \mathrm{Cof}_0$.
% foot: Version 9 novembre 2001

\printedpage{63}
% head: A. Grothendieck : Dérivateurs, Ch. VI, Édition des Universités de Montpellier II et Paris VII
\msp{page 100}

On reprend la description du `foncteur' $C$ (à valeurs dans $(\Fibu S)W_S^{-1}$, et alors sans guillemets), on trouve que $C(i_X)$, où
\[
X \xrightarrow{u = i_X} X' = C(X)
\]
est défini simplement (dans $(\Fibu S)W_S^{-1}$) par `la' flèche $C(u)$ rendant commutatif le diagramme (dans $\Fibu S$)
\[
\begin{tikzcd}
X \arrow[r, "u\ (=i_X)"] \arrow[d, "i_X"'] & X'\ (= C(X)) \arrow[d, "i_{X'}\ (=i_{C(X)})"] \\
C(X) \arrow[r, "C(u)\ (=C(i_X))"] & C(X')\ (= C(C(X))),
\end{tikzcd}
\]
de sorte qu'on a dans $\mathcal{M}/S$ le diagramme
\[
\begin{tikzcd}
X \arrow[r, "\overbrace{i_X}^{\in\mathrm{Cof}_0\cap W}"] & \overbrace{C(X)}^{\in\mathrm{Ob}\,\Fibu S} \arrow[r, shift left, "i_{C(X)}"] \arrow[r, shift right, "C(i_X)"'] & \overbrace{C(C(X))}^{\in\mathrm{Ob}\,\Fibu S},
\end{tikzcd}
\]
commutatif, i.e. $i_{C(X)}i_X = C(i_X)i_X$. Le corollaire page 92 nous dit que les flèches $i_{C(X)}$ et $C(i_X)$ sont égales dans $(\Fibu S)W_S^{-1}$. Mais par hypothèse (p. 95), dans $\Fibu S$ on a $W_S = \WL$, donc les flèches sont égales dans $(\Fibu S)(\WL)^{-1}$, et a fortiori dans $(\mathcal{M}/S)(\WL)^{-1}$, et on gagne.

Voyons ce qu'on peut faire dans le cas où on dispose d'un \textit{foncteur} $C$. On suppose bien sûr que ce foncteur dans $\mathcal{M}/S$ \msp{page 101} n'est pas lié au $S$ particulier dans $\mathcal{M}$ qu'on a choisi comme base, que sa construction provient à la \textsc{Cartan-Serre} d'un foncteur
\[
X \longmapsto \mathrm{Ch}(X) : \mathcal{M} \longrightarrow \mathcal{M},
\]
(`objet des chemins') et de morphismes fonctoriels
\[
X \xrightarrow{\alpha_X} \mathrm{Ch}(X) \xrightarrow{p_X} X \times X, \qquad\qquad \boxed{p_X\alpha_X = \mathrm{diag}_X}
\]
(\footnote[58]{\textbf{NB} On pose $\mathrm{pr}_1 \circ p_X = \sigma_X$, $\mathrm{pr}_2 \circ p_X = \beta_X$ (flèches `source' et `but' pour l'objet $X$).}). Je présume qu'il faudra supposer surtout ceci

\textbf{Hypothèse 2.} \textit{$\alpha_X \in W$, $p_X \in$} Fib.

Pour un $X$ sur $S$, on construira $C_S(X)$ ou $C(X)$ ainsi
% foot: Version 9 novembre 2001


\endgroup
% END UNIT Der-ChVI-345/Der-ChVI-345_p061-065.tex
% BEGIN UNIT Der-ChVI-345/Der-ChVI-345_p066-070.tex
\begingroup\makeatletter%
\newcommand{\printedpage}[1]{\marginpar{\footnotesize p.~#1}}%
\newcommand{\msp}[1]{\marginpar{\footnotesize [#1]}}%
\newcommand{\Fib}{\underline{\mathrm{Fib}}}%
\newcommand{\Ch}{\mathrm{Ch}}%
\newcommand{\pr}{\mathrm{pr}}%
\newcommand{\circled}[1]{\tikz[baseline=(c.base)]\node[draw,circle,inner sep=1pt](c){#1};}%
\makeatother\let\printedpage\derpp\let\msp\dermsp%
\let\mspin\DERmspin%


\printedpage{64}
% head: A. Grothendieck : Dérivateurs, Ch. VI, Édition des Universités de Montpellier II et Paris VII

\[
\begin{tikzcd}[column sep=large]
 & \Ch(S) \arrow[d, "p_S"] \arrow[dl, dotted, "\beta_S"'] \arrow[dd, dotted, bend right=40, "\sigma_S"'] \arrow[dr, phantom, "\text{cart.}"]
 & C_S(X)=C(X) \arrow[l, "\pi"'] \arrow[d, "p'"] \arrow[dll, bend right=40, "\bar f_X=f_{C(X)}"'] \\
S & S\times S \arrow[l, "\pr_2^S"'] \arrow[d, "\pr_1^S"] \arrow[dr, phantom, "\text{cart.}"]
 & X\times S \arrow[l, "\pi_1=f_X\times\mathrm{id}_S"] \arrow[d, "\pr_1"] \\
 & S & X(=S'). \arrow[l, "f_X"'] \arrow[uu, dotted, bend right=40, "i_X"']
\end{tikzcd}
\]

Il est l'image inverse de $(\Ch(S),\pr_1^Sp_S=\sigma_S)$ par $X\xrightarrow{f_X}S$, mais on a explicité la factorisation de $\sigma_S$ sur le diagramme pour réinterpréter la flèche structurale $\bar f_X=f_{C(X)}$, définie comme
\[
f_{C(X)}=\beta_S\pi(=\pr_2^S\underbrace{p_S\pi}_{=\pi_1p'})
\]
par la formule
\[
f_{C(X)}=\underbrace{(\pr_2^S\pi_1)}_{\text{projection } \pr_2 \text{ de } X\times S \text{ dans } S}p'.
\]

\msp{page 102}
Or
\[
\begin{aligned}
&p'\in\mathrm{Fib}, && \text{car déduit de } p_S\in\mathrm{Fib} \text{ par changement de base } \pi_1,\\
&(\pr_2 : X\times S\longrightarrow S)\in\mathrm{Fib} && \text{pourvu que } X\longrightarrow\underbrace{e}_{\text{objet final de } \mathcal{M}}\quad \text{soit dans Fib (même raison),}
\end{aligned}
\]
donc si $X\in\Fib\,e$, on a $f_{C(X)}\in\mathrm{Fib}$. Comme c'est une propriété essentielle qu'on veut avoir pour \textit{tout} $X$ dans $\mathcal{M}/S$, on va faire l'hypothèse supplémentaire

\textbf{Hypothèse 3.} \textit{$\forall\,X\in\mathit{Ob}\,\mathcal{M}$, $X\longrightarrow e$ est dans Fib.}

Quant à
\[
i_X : X\longrightarrow C(X),
\]
il est déduit de la section $\alpha_S$ de $(\Ch(S),\sigma_S)$ sur $S$ par changement de base $X\longrightarrow S$. Notons maintenant que
\[
\begin{tikzcd}
\Ch(S) \arrow[r, shift left, "\sigma_S"] \arrow[r, shift right, "\beta_S"'] & S
\end{tikzcd}
\qquad\qquad\qquad \text{sont dans } W,
\]
car leurs composés avec $S\xrightarrow{\alpha_S}\Ch(S)$ sont l'identité, et $\alpha_S\in W$. D'autre part, $\sigma_S=\pr_1^S\circ p$ est dans Fib, comme composé de deux éléments dans Fib (NB $S\longrightarrow e$ dans Fib par Hypothèse 3). Ainsi $\sigma_S$, $\beta_S$ sont dans $\mathrm{Fib}\cap W$. On aura besoin de % [sic] \pr_1^S\circ p -> \pr_1^S\circ p_S, la flèche est notée p_S dans le diagramme et plus haut (\pr_1^Sp_S=\sigma_S)

\textbf{Hypothèse 4.} \textit{L'ensemble de flèches $\mathit{Fib}\cap W$ est stable par tout changement de base.}

% foot: Version 9 novembre 2001

\printedpage{65}
% head: A. Grothendieck : Dérivateurs, Ch. VI, Édition des Universités de Montpellier II et Paris VII

\msp{page 103}
Cela implique donc que
\[
C(X)\xrightarrow{\pr_1\pi'}X
\]
\texttt{[plutôt} $\pr_1p'$\texttt{]} est $\in W$ (étant déduit de $\sigma_S : \Ch(S)\longrightarrow S$ par changement de base), donc toute section $i_X : X\longrightarrow C(X)$ de $C(X)$ sur $X$ est dans $i_X$ \texttt{[plutôt dans} $W$\texttt{]}. On a donc bien une factorisation, fonctorielle en $X$
\[
\begin{tikzcd}
 & C(X) \arrow[dr, "\bar f_X=f_{C(X)}"] & \\
X \arrow[ur, "i_X"] \arrow[rr, "f_X"] & & S
\end{tikzcd}
\]
avec
\[
i_X\in W,\qquad\qquad f_{C(X)}\in\mathrm{Fib}.
\]
C'est tout ce qu'on avait utilisé pour établir
\[
(\Fib\,S)W_S^{-1}\xrightarrow{\sim}(\mathcal{M}/S)W_S^{-1}.
\]
Mais il s'agit maintenant, sous l'hypothèse où on introduit un $\mathcal{L}$, de prouver que pour tout $X$ dans $\mathcal{M}/S$, les deux flèches
\[
\begin{tikzcd}
C(X) \arrow[r, shift left, "C(i_X)"] \arrow[r, shift right, "i_{C(X)}"'] & C(C(X))
\end{tikzcd}
\]
sont égales dans $(\mathcal{M}/S)(W_S^{\mathcal{L}})^{-1}$. Comme par hypothèse on a $W_S^{\mathcal{L}}=W_S$ dans $\Fib\,S$ (Hypothèse 1), et que $C(X),\ C(C(X))\in\mathrm{Ob}\,\Fib\,S$,
\msp{page 104}
il suffit de prouver que les deux flèches sont égales \textit{dans} $(\Fib\,S)W_S^{-1}$. Cet énoncé est indépendant de la donnée de $\mathcal{L}$, et on a envie de le prouver.

Peut-être y a-t-il moyen d'expliciter un langage géométrique, pour s'y reconnaître. Identifions les objets $X$ de $\mathcal{M}$ aux foncteurs $\mathcal{M}^o\longrightarrow(\mathrm{Ens})$ qu'ils définissent,
\[
X\longmapsto(T\longmapsto X(T)=\mathrm{Hom}_{\mathcal{M}}(T,X))\,.
\]
Ainsi, pout tout $T$ % [sic] pout -> pour, faute de frappe
\[
\Ch(X)(T)
\]
est un ensemble, appelé \textit{l'ensemble des chemins dans $X(T)$.} Il s'envoie (par $p_X$) fonctoriellement (en $T$) dans $(X\times X)(T)=X(T)\times X(T)$, i.e. chaque chemin $c$ dans $X(T)$ a une origine $\sigma(c)$, et un but $\beta(c)$. De plus, par $\alpha_X$, $X(T)$ s'envoie (fonctoriellement en $T$) dans l'ensemble des chemins de $X(T)$ : pour chaque élément $x$ dans $X(T)$, il y a donc un chemin (noté $1_x$ ou $\alpha_x$) bien défini de source et but $x$, qu'on appelle le chemin trivial, ou vide. Ceci dit,

% foot: Version 9 novembre 2001

\printedpage{66}
% head: A. Grothendieck : Dérivateurs, Ch. VI, Édition des Universités de Montpellier II et Paris VII

\msp{page 105}
pour $X$ dans $\mathcal{M}/S$, $C(X)$ représente le contra-foncteur suivant dans $\mathcal{M}$ : pour tout $T$ dans $\mathcal{M}$,
\[
C(X)(T)=\Bigl\{(x,c)\Bigm| x\in X(T),\ c : f_X(x)\xrightarrow{c}s\ \text{ un chemin dans } S(T), \text{ de source } f_X(x)\Bigr\}\,,
\]
où $f_X : X\longrightarrow S$ est le morphisme structural.

Les ‘calculs’ faits, il y a deux (?) jours montrent qu'on a de même une interprétation
\[
\begin{aligned}
C(C(X))(T)=\Bigl\{(x,c,c')\Bigm| x\in X(T),\ f_X(x)\xrightarrow{c}s\xrightarrow{c'}t,\\
c\text{ un chemin de source } f_X(x),\\
c'\text{ un chemin de source égale à } \beta(c)\Bigr\}
\end{aligned}
\]
et les applications $i_{C(X)}$ et $C(i_X)$ s'identifient respectivement à
\[
\begin{aligned}
(x,c)\quad&\longmapsto\quad(x,1_{f(x)},c)\quad f(x)\xrightarrow{1_x}f(x)\xrightarrow{c}s\\
(x,c)\quad&\longmapsto\quad(x,c,1_{\beta(c)})\quad f(x)\longrightarrow s\xrightarrow{1_s}s
\end{aligned}
\]
\texttt{[plutôt} $f(x)\xrightarrow{1_{f(x)}}f(x)\xrightarrow{c}s$\texttt{]}. Pour les comparer, on a une envie irrésistible d'écrire
\[
c\circ 1_x=1_s\circ c=c,
\]
mais a-t-on le droit? Il faudrait pouvoir composer les chemins, et c'est très, très difficile!

Mais on peut toujours axiomatiser. On suppose donc que pour $X$ dans $\mathcal{M}$, $\Ch(X)$
\msp{page 106}
n'est pas seulement une structure de (soi-disants) ‘chemins’ dans $X$ amorphe, seulement la donnée de la source et du but des chemins et des chemins triviaux (ces structures sur $\Ch(X)$ dépendant fonctoriellement de $X$), mais qu'on a une composition des chemins quand ceux-ci sont ‘composables’, i.e. une flèche d'un produit fibré
\[
(c,c')\longmapsto c'\circ c : (\Ch(X),\beta_X)\times_X(\Ch(X),\sigma_X)\xrightarrow{\pi_X}\Ch(X),
\]
i.e. on peut composer des chemins dans $X(T)$
\[
x\xrightarrow{c}y\xrightarrow{c'}z
\]
si et seulement si $\text{but}\,c=\text{source}\,c'$, et de plus on exige bien sûr que $\text{source}(c'\circ c)=\text{source}\,c$, $\text{but}(c'\circ c)=\text{but}\,c'$, ce qui s'exprime par un diagramme commutatif
\[
\begin{tikzcd}
(\Ch(X),\beta_X)\times_X(\Ch(X),\sigma_X) \arrow[rr] \arrow[dr] & & \Ch(X) \arrow[dl, "p_X"] \\
 & X\times X. &
\end{tikzcd}
\]

% foot: Version 9 novembre 2001

\printedpage{67}
% head: A. Grothendieck : Dérivateurs, Ch. VI, Édition des Universités de Montpellier II et Paris VII

La pensée naturelle, c'est qu'avec cette composition et ces identités, les $X(T)$, $\Ch(T)$ deviennent les objets et les flèches d'une \textit{catégorie} - i.e. que $(X,\Ch(X),\alpha_X,p_X,\pi_X)$ définissent une structure de ‘catégorie dans $\mathcal{M}$’. Mais pour ce qu'on veut faire maintenant, on n'a sans doute pas besoin de tout, i.e.
\msp{page 107}
de l'axiome pertinent d'associativité de la composition et celui sur les unités, mais seulement que les $1_x$ soient des unités bilatères pour la composition.

En supossant donné ce supplément de structure sur le foncteur Ch, et revenant à la considération du foncteurs associé % [sic] supossant -> supposant, faute de frappe ; foncteurs -> foncteur, accord avec « du » et « associé »
\[
C : \mathcal{M}/S\longrightarrow\Fib\,S,
\]
on trouve un morphisme canonique
\[
C(C(X))\xrightarrow{\gamma_X}C(X)
\]
par
\[
(x,c,c')\longmapsto(x,c'\circ c),
\]
et l'axiome des unités bilatères nous assure aussitôt que
\[
(\ast)\qquad\qquad\gamma_X\circ i_{C(X)}=\gamma_X\circ C(i_X)=\mathrm{id}_{C(X)}.
\]
Comme $i_{C(X)}\in W$, cela prouve déjà que $\gamma_X\in W$, donc $\gamma_X\in W_S$. Donc $\gamma_X$ devient un isomorphisme dans $(\Fib\,S)W_S^{-1}$, et la relation $(\ast)$ implique bien $i_{C(X)}=C(i_X)$ dans le localisé, q.e.d.

Je vais à nouveau résumer ce que j'ai dégagé.

\msp{page 108}
\textbf{Théorème-Scholie 2.} Soit $\mathcal{M}$ une catégorie, munie de deux ensembles de flèches
\[
W,\ \mathrm{Fib}\subseteq\mathrm{Fl}(\mathcal{M}).
\]
On fait les hypothèses suivantes.

\begin{enumerate}[label={}]
\item[$1^o$] $W$ contient les isomorphismes, si $f$, $g$ sont composables, si deux parmi $f$, $g$, $gf$ sont dans $W$, le troisième aussi.
\item[$2^o$] $\mathcal{M}$ est stable par $\varprojlim$ finies, Fib contient les isomorphismes, est stable par composition et par changement de base (\footnote[59]{Ces hypothèses de stabilité sur Fib ne serviront que dans II, et encore seulement dans le cas A), repris in toto sous la forme A$'$ plus bas (p.~112).}).
\end{enumerate}

Pour $S$ dans $\mathcal{M}$, \fbox{$\Fib\,S$} désigne la sous-catégorie pleine de $\mathcal{M}/S$ formée des $X\longrightarrow S$ qui sont dans Fib. On distingue deux versions différentes d'un axiome de factorisation d'une $f : X\longrightarrow S$ quelconque dans $\mathcal{M}$ en $X\xrightarrow{i}\bar X\xrightarrow{\bar f}S$, avec $i\in W$, $\bar f\in\mathrm{Fib}$.

% foot: Version 9 novembre 2001

\printedpage{68}
% head: A. Grothendieck : Dérivateurs, Ch. VI, Édition des Universités de Montpellier II et Paris VII

\begin{enumerate}[label={}]
\item[A)] On se donne, pour tout $S$ dans $\mathcal{M}$, un foncteur
\[
\mathcal{M}/S\longrightarrow\mathcal{M}/S\qquad\qquad X\longmapsto C(X)\ \ (=\bar X)
\]
et un morphisme de foncteurs en $\mathrm{X}$
\[
X\xrightarrow{i_X}C(X)
\]
avec
\[
i_X\in W,\qquad\qquad C(X)\in\mathrm{Ob}\,\Fib\,S.
\]
\item[B)] Désignant par
\[
\mathrm{Cof}_0\subseteq\mathrm{Fl}\,\mathcal{M}
\]
l'ensemble des flèches de $\mathcal{M}$ qui possèdent la propriété de relèvement à droite \texttt{[plutôt à gauche]} par rapport à Fib, on suppose ‘à la \textsc{Quillen}’, que toute flèche $f : X\longrightarrow S$ se factorise en
\[
X\xrightarrow{i_X}C(X)\xrightarrow{\bar f}S,
\]
avec
\[
i\in\mathrm{Cofib}_0\cap W,\qquad\qquad\bar f\in\mathrm{Fib}.
\]
\texttt{[où} $\mathrm{Cof}_0=\mathrm{Cofib}_0$\texttt{]}.
\end{enumerate}

\msp{page 109}
Donnons nous $S$ dans $\mathcal{M}$, considérons
\[
W_S\subseteq\mathrm{Fl}(\mathcal{M}/S)
\]
l'ensemble de flèches induit par $W$ (
\begin{tikzcd}[column sep=small, row sep=small, baseline={([yshift=.5ex]current bounding box.south)}]
f : X \arrow[rr] \arrow[dr] & & Y \arrow[dl] \\
 & S &
\end{tikzcd}
est dans $W_S$ si et seulement si il l'est en oubliant $S$), et désignons par $W_S$ (par abus de notation) l'ensemble de flèches induit par $W_S$ sur $\Fib\,S$, ou plus généralement sur toute sous-catégorie pleine $\mathcal{F}$ de $\mathcal{M}/S$, avec
\[
\Fib\,S\subseteq\mathcal{F}\subseteq\mathcal{M}/S.
\]

\circled{I} Le foncteur
\[
(\ast)\qquad\qquad(\Fib\,S)W_S^{-1}\xrightarrow{\sim}(\mathcal{M}/S)W_S^{-1}
\]
induit par l'inclusion est une équivalence de catégories, plus généralement pour toute $\mathcal{F}$ comme dessus
\[
(\Fib\,S)W_S^{-1}\longrightarrow\mathcal{F}W_S^{-1}\longrightarrow(\mathcal{M}/S)W_S^{-1}
\]
sont des équivalences de catégories. Un foncteur quasi-inverse est obtenu (dans les deux cas A) et B)) à l'aide du foncteur
\[
\gamma : \mathcal{M}/S\longrightarrow(\Fib\,S)W_S^{-1}
\]

% foot: Version 9 novembre 2001


\endgroup
% END UNIT Der-ChVI-345/Der-ChVI-345_p066-070.tex
% BEGIN UNIT Der-ChVI-345/Der-ChVI-345_p071-075.tex
\begingroup\makeatletter%
\newcommand{\printedpage}[1]{\marginpar{\footnotesize p.~#1}}%
\newcommand{\msp}[1]{\marginpar{\footnotesize [#1]}}%
\newcommand{\circled}[1]{\tikz[baseline=(c.base)]\node[draw,circle,inner sep=1pt](c){#1};}%
\newcommand{\Fib}{\underline{\mathrm{Fib}}\,}%
\newcommand{\Fibsc}{\underline{\mathrm{Fibsc}}\,}%
\newcommand{\Cofib}{\underline{\mathrm{Cofib}}\,}%
\newcommand{\Cofibsc}{\underline{\mathrm{Cofibsc}}\,}%
\newcommand{\Fibhot}{\underline{\mathrm{Fibhot}}\,}%
\newcommand{\Parf}{\underline{\mathrm{Parf}}\,}%
\newcommand{\Liss}{\underline{\mathrm{Liss}}}%
\newcommand{\Prop}{\underline{\mathrm{Prop}}}%
\newcommand{\Ho}{\mathrm{Ho}}%
\newcommand{\Hom}{\mathrm{Hom}}%
\newcommand{\HOT}{\mathrm{HOT}}%
\newcommand{\hot}{\mathrm{hot}}%
\newcommand{\Fl}{\mathrm{Fl}}%
\newcommand{\Cat}{\mathrm{Cat}}%
\newcommand{\M}{\mathcal{M}}%
\makeatother\let\printedpage\derpp\let\msp\dermsp%
\let\mspin\DERmspin%

\setcounter{footnote}{59}

\printedpage{69}
% head: A. Grothendieck : Dérivateurs, Ch. VI, Édition des Universités de Montpellier II et Paris VII
qui est donné sur les objets par
\[\gamma(X) = C(X),\]
et qui sur les flèches se définit de façon évidente dans le cas A) (fonctorialité de $C(X)$ en $X$), et dans le cas de B) se construit ad hoc, comme la seule loi fonctorielle pour laquelle les homomorphismes
\[i_X : X \longrightarrow C(X)\]
\msp{page 110}
(interprétés comme flèches dans $(\M/S)W_S^{-1}$) soient eux-mêmes fonctoriels en $X$.

On désigne par
\[\Ho^{\mathrm{lc}}_W(S) \qquad\qquad \text{(lc comme ‘localement constant’)}\]
la catégorie des fractions $(\M/S)W_S^{-1}$, équivalente à $(\Fib S)W_S^{-1}$.

\circled{II} Donnons nous un localiseur $W_S^0$ dans $\M/S$
\[W_S^0 \subseteq W_S \subseteq \Fl(\M/S),\]
avec l'hypothèse formelle similaire à celle de $W$ (cf. 1$^o$). On suppose que

\begin{itemize}[leftmargin=4em]
\item[3$^o$)] $W_S^0 \cap \Fl(\Fib S) = W_S|\Fib S$.
\end{itemize}

(Cf. pages 94-96 pour une façon générale d'obtenir des tels localiseurs. Il faut bien voir que pour définir des ‘types d'homotopie’ relatifs sur $S$ le localiseur $W_S$ est beaucoup trop grossier, moralement il ne donne qu'une catégorie des types d'homotopie relatifs localement constants, et c'est justement ça qui est montré par l'equivalence de \circled{I} , où $\Fib S$ est visualisé comme formé des ‘fibrations localement triviales’ sur $S$.) On a alors un deuxième localisé % [sic] l'equivalence -> l'équivalence, accent manquant
\[\Ho_{W_S^0}(S) \overset{\text{déf}}{=} (\M/S)(W_S^0)^{-1},\]
et deux foncteurs
\[
\begin{tikzcd}[column sep=large]
\Ho_{W_S^0}(S) \arrow[r, shift left, "\rho"] & \Ho^{\mathrm{lc}}_W(S), \arrow[l, shift left, "i"]
\end{tikzcd}
\tag*{$(*)$}
\]
où $\rho$ est le foncteur de localisation (strict) évident, dû à l'inclusion $W_S^0 \subseteq W_S$, et où $i$ est
\msp{page 111}
essentiellement le foncteur déduit de l'inclusion
\[(\Fib S, \underbrace{W_S|\Fib S}_{=W_S^0|\Fib S}) \longrightarrow (\M/S, W_S^0)\]
par localisation \textit{au moyen de $W_S^0$}, donc le foncteur
\[(\Fib S)(W_S^0)^{-1} \longrightarrow (\M/S)(W_S^0)^{-1}, \tag*{$(**)$}\]
% foot: Version 9 novembre 2001
\printedpage{70}
% head: A. Grothendieck : Dérivateurs, Ch. VI, Édition des Universités de Montpellier II et Paris VII
compte tenu de l'équivalence $(*)$ de \circled{I} , i.e. $i$ est le composé du quasi-inverse de $(*)$, avec le foncteur $(**)$ précédent. Explicitement, il est défini par
\[i(X) = \text{classe de C(X) dans } (\M/S)(W_S^0)^{-1},\]
en tenant compte du fait que $X \longmapsto C(X)$ transforme $W_S$ en $W_S^o$, grâce à la condition 3$^o$) sur $W_S^0$. % [sic] W_S^o -> W_S^0, lettre o au lieu du chiffre 0 utilisé partout ailleurs

Ceci posé, \textit{sous l'hypothèse} B (‘Quillen’), on a : dans le couple de foncteurs $(\rho, i)$ de $(*)$, page 110, \textit{$i$ est adjoint à droite de $\rho$,} i.e. on a pour $\xi \in \mathrm{Ob}\,\Ho_{W_S^0}(S)$, $\eta \in \mathrm{Ob}\,\Ho^{\mathrm{lc}}_W(S)$
\[\Hom_{W_S^0}(\xi, i(\eta)) \simeq \Hom_{W_S}(\rho(\xi), \eta),\]
ce qui s'explicite, pour $X$, $Y$ dans $\M/S$, par
\[\boxed{\Hom_{W_S^0}(X, C(Y)) \simeq \Hom_{W_S}(X, Y)}\]
(\footnote{\textbf{NB} Le premier membre est indépendant du choix de $W_S^0$! (Pourvu seulement que l'hypothèse $W_S^0 \cap \Fl\,\Fib S = W_S \cap \Fl\,\Fib S$, en plus de $W_S^0 \subseteq W_S$, soit satisfaite.)}). Il s'ensuit que \textit{le foncteur $i$ est pleinement fidèle.}

La même formule d'adjonction, et la
\msp{page 112}
pleine fidélité de $i$, sont valables dans le cas A, quand on renforce l'hypothèse A (relative à un $S$ particulier) en une hypothèse plus restrictive et plus compréhensive, que voici.

\begin{itemize}[leftmargin=3em]
\item[A$'$)] Le foncteur $C : \M/S \longrightarrow \M/S$ et le morphisme fonctoriel $i_X : X \longrightarrow C(X)$ sont construits ‘à la \textsc{Cartan-Serre}’ à partir d'un foncteur défini sur $\M$ entier
\[\M \xrightarrow{\mathrm{Ch}} \M,\]
du type ‘chemins’, i.e. muni en plus des données suivantes
\begin{itemize}[leftmargin=3em]
\item[a)] \[X \xrightarrow{\alpha_X} \mathrm{Ch}(X) \xrightarrow{p_X = (\sigma_X, \beta_X)} X \times X\]
(structure de source et but des chemins, et de chemins triviaux), et
\item[b)] Une composition des chemins, sans exiger ici l'associativité de la composition mais seulement que les chemins triviaux soient unités bilatères pour la composition. De plus, on suppose que pout tout $X$ % [sic] pout -> pour, faute de frappe
\[\begin{cases} \alpha_X \in W,\ p_X \in \mathrm{Fib} \\ (X \longrightarrow e) \in \mathrm{Fib}. \end{cases}\]
\end{itemize}
\end{itemize}
% foot: Version 9 novembre 2001

\printedpage{71}
% head: A. Grothendieck : Dérivateurs, Ch. VI, Édition des Universités de Montpellier II et Paris VII
\msp{page 113}
\textbf{Corollaire.} Soient $X$ dans $\M/S$, $Y$ dans $\Fib S$, alors
\[\Hom_{W_S^0}(X, Y) \xrightarrow{\sim} \Hom_{W_S}(X, Y).\]
\textsc{Démonstration}. Par la formule d'adjonction (p. 111), le deuxième membre est $\Hom_{W_S^0}(X, CY)$. Mais comme $Y$ et $C(Y)$ sont dans $\Fib S$, où $W_S$ et $W_S^0$ coïncident, $Y \xrightarrow{i_Y} C(Y)$ est non seulement dans $W_S$, mais dans $W_S^0$, donc $\Hom_{W_S^0}(X, Y) \xrightarrow{\sim} \Hom_{W_S^0}(X, CY)$, d'où le résultat.

Revenons à Cat, Cat$/S$. La digression précédente sur $(\M, W, \mathrm{Fib})$ a permis de mieux situer ce qui concerne le foncteur $C$, et les relations entre les catégories localisées de
\[\Cat/S \ \ (\text{pour } W_S,\ W_S^{\mathrm{d}},\ W_S^{\mathrm{g}}) \qquad\quad \text{et} \qquad\quad \Parf S \ \ (\text{pour } W_S^{\mathrm{f}} = W_S = W_S^{\mathrm{d}} = W_S^{\mathrm{g}})\]
ainsi que pour les catégories intermédiaires
\[\Parf S \subseteq \mathcal{F} \subseteq \Cat/S\]
pleines de Cat$/S$ (sur le diagramme p. 76, il y en a cinq autres, savoir $\Fibhot S$, $\Liss\,S$, $\Liss_0 S$, $\Prop\,S$, $\Prop_0\,S$). Pour ce qui concerne les relations entre les localisés de ces sept catégories, qui se trouvent au centre du diagramme, les énoncés du théorème-scholie 1 (p. 70-81) ne donnent rien de plus que le contenu du théorème-scholie 2 (p. 108-113). Pour les établir, on a donc le choix entre les deux hypothèses A) ou B) dudit théorème. Soit utiliser un foncteur $C_S$ explicite, à l'aide d'un
\msp{page 114}
foncteur $X \longmapsto \underline{\mathrm{Ch}}(X)$ (‘chemins dans $X$’) sur Cat, donc les structures de source, but et composition pour les chemins, et les chemins triviaux (agissant comme unités bilatères pour la composition). La construction d'un tel foncteur $C$ n'a par l'air trivial et semble ne pas être connue - j'ai peiné pour y arriver! (Avec, en plus, l'associativité de la composition \ldots). Ou bien utiliser, à la \textsc{Quillen}, l'existence d'une classe de morphismes dits Q-fibrants, avec des bonnes propriétés de factorisation. Ce n'est pas non plus évident dans Cat, et c'est établi par \textsc{Thomason}, par une voie un peu compliquée. Peut-être serai-je amené à revenir dessus - mais seulement en cas de besoin! % [sic] n'a par l'air -> n'a pas l'air, faute de frappe

Par contre, les assertions concernant les catégories $\Fib S$, $\Fibsc S$, ou $\Cofib S$, $\Cofibsc S$, et leurs variantes à indices 0 (donc les huit catégories-ailes du diagramme), et leurs relations à Cat$/S$ d'une part, à $\Fibhot S$ de l'autre (pour les catégories à indice 0), ne peuvent être contenues dans un tapis général sur des catégories de modèles $(\M, W)$, me semble-t-il, vu la nature très spéciale des quatre catégories désignées ci-dessus, de leur rôle très particulier dans le formalisme catégorique. Or tout au moins les catégories ‘scindées’ $\Fibsc S$, $\Cofibsc S$, et par suite aussi leurs variantes à indice 0, jouent un rôle essentiel pour la théorie des HOT$(S)$, vu que c'est elles qui
\msp{page 115}
interviennent dualement quand on applique la recette générale de construction d'un dérivateur à partir d'une catégorie de modèles $(\M, W)$. Donc je ne vois pas de simplification
% foot: Version 9 novembre 2001
\printedpage{72}
% head: A. Grothendieck : Dérivateurs, Ch. VI, Édition des Universités de Montpellier II et Paris VII
substantielle à apporter au théorème-scholie 1. Tout au plus pourrait-on élaguer le diagramme de catégories envisagées, en laissant tomber $\Liss\,S$, $\Liss_0 S$, $\Prop\,S$, $\Prop_0 S$, qui n'y jouent pas un rôle vraiment utile. Pour ce qui concerne les résultats énoncés, ce qui les concerne est couvert par ce qui est dit pour les sous-catégories pleines $\mathcal{F}$ de Cat$/S$ qui contiennent (disons) $\Parf S$ (ou au choix celles qui contiennent $\Fib S$ ou $\Cofib S$).

La description des catégories du type HOT$(S)$, pour $S$ fixé, me paraît à présent parfaitement bien comprise. Il est temps de faire varier $S$, et de comprendre le caractère contravariant en $S$. On se donne donc
\[S' \xrightarrow{f} S,\]
et on se propose de décrire les foncteurs image inverse pour les trois types de catégories HOT$(S)$. Mais nous allons nous borner à présent aux $\HOT_W(S)$ et $\HOT^{\mathrm{lc}}_W(S)$, puisque tout ce qui concerne $\HOT^o_W(S) = \HOT_W(S^o)$ \texttt{[plutôt} $\HOT_W(S^o)^o$\texttt{, cf. p. 50]} s'en déduira par dualité $S \longmapsto S^o$. Donc il s'agit de décrire
\msp{page 116}
\[f^* : \HOT_W(S) \longrightarrow \HOT_W(S') \tag*{$(*)$}\]
et le foncteur induit
\[f^* : \HOT^{\mathrm{lc}}_W(S) \longrightarrow \HOT^{\mathrm{lc}}_W(S').\]
Quand on décrit les catégories HOT$(S)$, conformément à la recette générale, à l'aide de foncteurs
\[F : S^o \longrightarrow \Cat\]
(préfaisceaux en catégories sur $S$), ou, ce qui revient au même, de catégories scindées, la description de la flèche $(*)$ est évidente : elle revient à la loi
\[\left|\begin{array}{rcl}
F & \longmapsto & F \circ f^o \\[1ex]
\underline{\mathrm{Hom}}(S^o, \Cat) & \longrightarrow & \underline{\mathrm{Hom}}(S'^o, \Cat),
\end{array}\right.\]
en passant aux localisés pour la $W$-équivalence argument par argument. Dans le langage des catégories fibrées, il est bien connu que cette opération triviale sur les préfaisceaux en catégories correspond à l'opération \textit{image inverse} de catégories sur $S$ - laquelle opération passe aux scindages. Cela nous donne donc
\[f^*(\hot_{W,S}(X)) = \hot_{W,S'}(\underbrace{X \times_S S'}_{X'}),\]
pour toute $X$ fibrée scindée sur $S$. Or on sait que le scindage de $X$ ne joue aucun rôle, finalement, dans la description de la classe
\[\hot_{W,S}(X) \in \HOT_W(S)\]
% foot: Version 9 novembre 2001
\printedpage{73}
% head: A. Grothendieck : Dérivateurs, Ch. VI, Édition des Universités de Montpellier II et Paris VII
\msp{page 117}
associée à $X$. On peut oublier le scindage, et même (pour étudier les $\Hom_{\mathrm{Hot}_W(S)}(X, Y)$) travailler avec des $S$-foncteurs \textit{quelconque} entre catégories fibrées sur $S$ - pas la peine même de les supposer cartésiens, voir même compatibles avec des scindages. Cela me paraît une chose très importante pour le formalisme des $\HOT_W(S)$, et très spéciale au cas de la catégorie de modèles Cat. C'est un peu tacite dans le théorème-scholie - mais mes doutes initiaux concernant la catégorie $\underline{\mathrm{Fibsc}}$ elle-même montrent bien que je sentais qu'il y avait là quelque chose de pas évident, ni même (presque) \texttt{[?]} de vraisemblable. Et pourtant l'existence du foncteur $\Phi$ montre qu'il en est bien ainsi.

On peut donc interpréter $f^*$ en termes de $\Fib S$, $\Fib S'$, par le foncteur changement de base
\[f^* : \Fib S \longrightarrow \Fib S', \qquad\quad X \longmapsto X \times_S S'/S',\]
lequel foncteur préserve évidemment la $W$-équivalence par fibres $W_S^{\mathrm{f}}$, $W_{S'}^{\mathrm{f}}$, qu'est celle pertinente ici
\[f^*(W_S^{\mathrm{f}}) \subseteq W_{S'}^{\mathrm{f}},\]
de sorte que $f^*$ induit bien
\[(\Fib S)(W_S^{\mathrm{f}})^{-1} \longrightarrow (\Fib S')(W_S^{\mathrm{f}})^{-1}\]
\texttt{[plutôt} $(\Fib S')(W_{S'}^{\mathrm{f}})^{-1}$\texttt{]},
\msp{page 118}
i.e. le foncteur $(*)$
\[f^* : \HOT_W(S) \longrightarrow \HOT_W(S').\]
D'autre part, $\HOT_W(S)$ peut être décrit aussi comme localisé de Cat$/S$ tout entier, par rapport au localiseur $W_S^{\mathrm{d}}$, lequel sur $\Fib S$, et même sur $\Liss\,S$, coïncide avec $W_S^{\mathrm{f}}$. Mais on se bute ici à la difficulté que le foncteur image inverse
\[f^* : \Cat/S \longrightarrow \Cat/S', \qquad\quad X \longmapsto X \times_S S'/S',\]
n'est \textit{pas} compatible en général (si $S$ n'est isomorphe à la catégorie ponctuelle, pour $S'$ quelconque) avec les localiseurs $W_S^{\mathrm{d}}$, $W_{S'}^{\mathrm{d}}$. On le voit p.ex. en prenant $S' = \{s\}$, et en se rappelant que $W_S^{\mathrm{d}}$ n'implique \textit{pas}, sur Cat$/S$ tout entier, la $W$-équivalence par fibres.

Ces ennuis n'ont pas lieu sûr $\Liss\,S$, où $W_S^{\mathrm{d}}$ et $W_S^{\mathrm{f}}$ coïncident. On peut donc décrire $f^*$ sur les HOT aussi à l'aide de % [sic] sûr -> sur, accent fautif
\[f^* : \Liss\,S \longrightarrow \Liss\,S', \qquad\quad X \longmapsto X \times_S S'/S',\]
en passant aux localisés
\[f^* : (\Liss\,S)(W_S^{\mathrm{f}})^{-1} \longrightarrow (\Liss\,S')(W_{S'}^{\mathrm{f}})^{-1}.\]
En d'autres termes :

\textbf{Proposition 1.} \textit{Si $X$ dans} Cat$/S$ \textit{est lisse sur $S$, alors on a un isomorphisme canonique}
\[f^*(\hot_{W,S}(X)) \simeq \hot_{W,S'}(X'), \qquad\quad \text{\textit{où} } X' = X \times_S S'/S'.\]
% foot: Version 9 novembre 2001


\endgroup
% END UNIT Der-ChVI-345/Der-ChVI-345_p071-075.tex
% BEGIN UNIT Der-ChVI-345/Der-ChVI-345_p076-080.tex
\begingroup\makeatletter%
\NoAutoSpacing%
\newcommand{\printedpage}[1]{\marginpar{\footnotesize p.~#1}}%
\newcommand{\msp}[1]{\marginpar{\footnotesize [#1]}}%
\newcommand{\und}[1]{\underline{\mathrm{#1}}}%
\newcommand{\HOT}{\mathrm{HOT}}%
\newcommand{\hot}{\mathrm{hot}}%
\newcommand{\Cat}{\mathrm{Cat}}%
\makeatother\let\printedpage\derpp\let\msp\dermsp%
\let\mspin\DERmspin%


\printedpage{74}
% head: A. Grothendieck : Dérivateurs, Ch. VI, Édition des Universités de Montpellier II et Paris VII
% foot: Version 9 novembre 2001

\msp{page 119}
Une remarque similaire s'applique à $\und{Fibhot}\,S$, où $W_S^{\mathrm{d}} = W_S^{\mathrm{f}}$ également. Le foncteur changement de base induit
\[
\und{Fibhot}\,S \longrightarrow \und{Fibhot}\,S', \qquad X \longmapsto X \times_S S'/S',
\]
compatible avec les localiseurs relatifs pertinents (tous ceux envisagés jusqu'à présent coïncident sur ces catégories, avec les hypothèses faites sur $W$), d'où
\[
f^* : (\und{Fibhot}\,S)W_S^{-1} \longrightarrow (\und{Fibhot}\,S')W_{S'}^{-1}.
\]
Les localisés cette fois ne sont pas les $\HOT_W$, mais les $\HOT_W^{\mathrm{lc}}$, donc on trouve une description sympathique de
\[
f^* : \HOT_W^{\mathrm{lc}}(S) \longrightarrow \HOT_W^{\mathrm{lc}}(S'),
\]
plus avantageuse encore qu'en termes de $\und{Liss}_0S$, $\und{Liss}_0S'$, car s'appliquant à des catégories un peu plus vastes (NB $\und{Liss}_0 \subseteq \und{Fibhot}$). Notons donc le

\textbf{Corollaire 1.} \textit{La formule de la proposition 1 (p. 118) reste valable si $X$ est supposé} Hot-\textit{fibrant sur $S$ (au lieu de lisse sur $S$).} (On trouve alors des types d'homotopie relatifs \textit{localement constants} sur $S$, $S'$.)

On a donc une description de $f^*$ qui s'applique à toutes les catégories du côté gauche
\msp{page 120}
du diagramme \texttt{[page]} 76, à la seule exception de $\Cat/S$ elle-même. Si $X$ est dans $\Cat/S$, sans être ni lisse ni Hot-fibrée sur $S$, comment trouver l'image inverse dans $\HOT_W(S')$? On a le choix entre $\Phi$ et $C$ pour améliorer la situation! On sait que $X$ définit la même classe dans $\HOT_W(S)$ que $\Phi(X)$, $C(X)$, qui sont l'une lisse (et même fibrée scindée), l'autre Hot-fibrée (et même parfaite) sur $S$. Donc

\textbf{Corollaire 2.} \textit{Si $X$ dans} $\Cat/S$ \textit{est quelconque, alors}
\[
f^*(\hot_{W,S}(X)) \simeq \hot_{W,S'}(\Phi(X) \times_S S') \simeq \hot_{W,S'}(C(X) \times_S S').
\]

\bigskip

On peut s'attendre, d'autre part, que pour certains changements de base particuliers, la formule de la proposition 1 soit valable sans restriction sur $X$ dans $\Cat/S$. Il doit revenir $\pm$ au même que le foncteur image inverse
\[
X \longmapsto X \times_S S'
\]
soit compatible avec les localiseurs $W_S^{\mathrm{d}}$, $W_{S'}^{\mathrm{d}}$. La condition de $W$-équivalence `colocale sur $S$\,' doit être vue comme une condition de nature `homologique', pour le dérivateur HOT. Elle s'associe, pour un dérivateur quelconque $\mathbf{D}$, à la condition (pour $f : X \longrightarrow Y$ dans $\Cat/S$) que les $s\backslash X \longrightarrow s\backslash Y$ sont des $\mathbf{D}$-équivalences,
\printedpage{75}
% head: A. Grothendieck : Dérivateurs, Ch. VI, Édition des Universités de Montpellier II et Paris VII
% foot: Version 9 novembre 2001
\msp{page 121}
ce qui s'écrit ici élégamment en termes d'homologie relative, par la condition que $\forall\ a \in \mathcal{A} = \mathbf{D}(e)$, la flèche induit par $f$ % [sic] induit -> induite, accord avec « la flèche »
\[
\underbrace{H_\bullet^{\mathbf{D}}(X/S, a_X)}_{\overset{\text{déf}}{=}(p_X)_!(a_X)} \longrightarrow \underbrace{H_\bullet^{\mathbf{D}}(Y/S, a_Y)}_{\overset{\text{déf}}{=}(p_Y)_!(a_Y)}
\]
soit \texttt{[un]} isomorphisme. (Utiliser le `calcul' des faisceaux d'homologie relative $(p_X)_!(a_X)$, $(p_Y)_!(a_Y)$, par la caractérisation des fibres dans Der 4 bis). Or, les changements de base $\mathbf{D}$-\textit{propres} $S' \longrightarrow S$ commutent aux $p_!$, de sorte qu'on trouve le

\textbf{Lemme.} \textit{La notion d'équivalence} $\mathbf{D}$\textit{-colocale sur $S$ dans} $\Cat/S$ \textit{est stable par tout changement de base $S' \longrightarrow S$} $\mathbf{D}$\textit{-propre. (Quel que soit le dérivateur} $\mathbf{D}$\textit{.)}

On se rappelera que les $s\backslash S \longrightarrow S$ sont justement des changements des base propres! On % [sic] rappelera -> rappellera, orthographe; des base -> de base, accord
peut dire que l'équivalence $\mathbf{D}$-colocale peut se décrire en exigeant que $f : X \longrightarrow Y$ soit transformé en $W$-équivalences par certains types de changement de base $S' \longrightarrow S$, en prenant au choix l'un des trois types

\begin{itemize}
\item[a)] Les changements de base $s\backslash S \longrightarrow S$ (ce qui redonne la définition que nous avions prise).
\item[b)] Les changements de base qui sont des co-isomorphismes locaux.
\item[c)] Les changements de base propres.
\item[d)] Les changements de base $\mathbf{D}$-propres.
\end{itemize}

On ne peut appliquer ces résultats qu'au dérivateur
\msp{page 122}
HOT, faute de l'avoir encore construit, et verifié, les axiomes pertinents. On peut poser % [sic] verifié -> vérifié, accent
simplement la

\textbf{Conjecture} (sûre!) : \textit{Soit $p : S' \longrightarrow S$ une flèche propre (sous-entendu :} HOT\textit{-propre) dans} Cat. \textit{Alors le foncteur changement de base}
\[
\Cat/S \longrightarrow \Cat/S', \qquad X \longmapsto X \times_S S'/S'
\]
\textit{est compatible avec les localiseurs $W_S^{\mathrm{d}}$, $W_{S'}^{\mathrm{d}}$. (Pour un foncteur $S' \longrightarrow S$ lisse, le changement de base serait compatible avec $W_S^{\mathrm{g}}$, $W_{S'}^{\mathrm{g}}$.)} (\footnote[61]{NB Cette `conjecture' est prouvé sans lemme à la page 168-169, sans avoir besoin de rien qui ne soit % [sic] prouvé -> prouvée, accord avec « conjecture »
déjà connu ici. Cf. théorème 4, cas b) p. 162.}).

Mais pout tout changement de base $S' \longrightarrow S$ (propre ou non) qui a la propriété précédente, % [sic] pout -> pour, faute de frappe
donc qui induit un foncteur
\[
\underbrace{(\Cat/S)(W_S^{\mathrm{d}})^{-1}}_{\approx\,\HOT_W(S)} \longrightarrow \underbrace{(\Cat/S')(W_{S'}^{\mathrm{d}})^{-1}}_{\approx\,\HOT_W(S')}\ ,
\]
\printedpage{76}
% head: A. Grothendieck : Dérivateurs, Ch. VI, Édition des Universités de Montpellier II et Paris VII
% foot: Version 9 novembre 2001
il est immédiat, à l'aide du théorème de comparaison (p.ex. avec $\und{Liss}\,S$, $\und{Liss}\,S'$) sur $S$ et sur $S'$, que ce foncteur n'est autre (à isomorphisme canonique près) que le foncteur $f^*$ image inverse des types d'homotopie relatifs. Donc énoncons le % [sic] énoncons -> énonçons, cédille

\textbf{Corollaire 3.} \textit{La formule de la proposition 1 reste valable sans hypothèse aucune sur $X$ dans} $\Cat/S$\textit{, pourvu que la flèche de changement de base $f : S' \longrightarrow S$ soit telle que le foncteur $X \longmapsto X \times_S S'$ de} $\Cat/S$ \textit{dans} $\Cat/S'$
\msp{page 123}
\textit{transforme $W$-équivalences colocales sur $S$ en $W$-équivalences colocales sur $S'$.} (Ce qui devrait être le cas si (et seulement si?) $f$ est $W$-propre, donc aussi si $f$ est (`absolument') propre (\footnote[62]{cf. I, pages 55-59.}).)

On n'a travaillé jusqu'ici qu'avec les localiseurs $W_S^{\mathrm{d}}$, $W_{S'}^{\mathrm{d}}$, à l'exclusion de $W_S$, $W_{S'}$. (Sauf sur les catégories contenues dans $\und{Fibhot}\,S$, où les deux coïncident.) Quand on ne fait pas d'hypothèse draconienne sur $p : S' \longrightarrow S$, du type Hot-fibration, on n'a aucune chance pour que $X \longmapsto X \times_S S'$ transforme $W$-équivalences grossières dans itou, donc aucune chance d'attraper l'image inverse du type d'homotopie relatif `grossier' de $X$, en prenant celui de $X' = X \times_S S'$. Cela est lié au fait que pour un $S' \longrightarrow S$ général, les foncteurs de localisation
\[
\HOT_W(S) \xrightarrow{\rho_S} \HOT_W^{\mathrm{lc}}(S), \qquad \HOT_W(S') \xrightarrow{\rho_{S'}} \HOT_W^{\mathrm{lc}}(S')
\]
n'ont aucune envie de commuter au foncteur image inverse $f^*$ sur les $\HOT_W$. On peut hasarder la

\textbf{Proposition 2.} \textit{Conditions équivalentes sur le foncteur $S' \longrightarrow S$ :}

\begin{itemize}
\item[\textit{a)}] \textit{Le diagramme de foncteurs}
\[
\begin{tikzcd}
\HOT_W(S) \arrow[r, "\rho_S"] \arrow[d, "f^*"] & \HOT_W^{\mathrm{lc}}(S) \arrow[d, "f^*_{\mathrm{lc}}"] \\
\HOT_W(S') \arrow[r, "\rho_{S'}"] & \HOT_W^{\mathrm{lc}}(S')
\end{tikzcd}
\]
\textit{commute (à isomorphisme de foncteurs composés près).}

\msp{page 124}

\item[\textit{b)}] \textit{Le foncteur changment de base} % [sic] changment -> changement, faute de frappe
\[
\und{Liss}\,S \longrightarrow \und{Liss}\,S', \qquad X \longmapsto X \times_S S'/S'
\]
\textit{transforme $W_S$ en $W_{S'}$.}
\printedpage{77}
% head: A. Grothendieck : Dérivateurs, Ch. VI, Édition des Universités de Montpellier II et Paris VII
% foot: Version 9 novembre 2001
\item[\textit{b$'$)}] \textit{Énoncé analogue pour} $\und{Fib}\,S \longrightarrow \und{Fib}\,S'$\textit{.}
\item[\textit{b$''$)}] \textit{Énoncé analogue pour} $\und{Fibsc}\,S \longrightarrow \und{Fibsc}\,S'$\textit{.}
\end{itemize}

(\footnote[63]{Pas prouvée, seulement que b, b$'$, b$''$ impliquent a). Pour l'implication inverse, il faudrait savoir que toute flèche de $\mathcal{F}'$ qui devient \texttt{[un]} isomorphisme dans $\mathcal{F}'W_{S'}^{-1}$, est dans $W_{S'}$ (saturation forte de $W_{S'}$ dans $\mathcal{F}'$), chose que je ne sais pas. Mais je crois que ce doit être vrai.}, \footnote[64]{\textbf{NB} Si on suppose $p$ propre, on doit pouvoir prendre aussi $\mathcal{F} = \Cat/S$.}).

\textsc{Démonstration.} Si $\mathcal{F}$ est une des trois catégories $\und{Liss}\,S$, $\und{Fib}\,S$ ou $\und{Fibsc}\,S$, le foncteur $\rho_S$ peut se décrire comme le foncteur de localisation canonique
\[
\mathcal{F}(W_S^{\mathrm{d}})^{-1} \longrightarrow \mathcal{F}W_S^{-1}
\]
(ça marcherait pour n'importe quelle sous-catégorie pleine $\mathcal{F}$ de $\Cat/S$ qui contient $\und{Fib}\,S$). Or si b), b$'$) ou b$''$) est satisfait, on trouve un diagramme commutatif de foncteurs induits par le changement de base
\[
\begin{tikzcd}
\mathcal{F}(W_S^{\mathrm{d}})^{-1} \arrow[r, "\rho_S"] \arrow[d, "f^*_{W_S^{\mathrm{d}}}"'] & \mathcal{F}(W_S)^{-1} \arrow[d, "f^*_{W_S}"] \\
\mathcal{F}'(W_{S'}^{\mathrm{d}})^{-1} \arrow[r, "\rho_{S'}"] & \mathcal{F}'(W_{S'})^{-1},
\end{tikzcd}
\]
sous réserve, il est vrai, dans le cas $\mathcal{F} = \Cat/S$, que le foncteur image inverse soit aussi compatible avec $W_S^{\mathrm{d}}$, $W_{S'}^{\mathrm{d}}$ (ce qui est automatique si $\mathcal{F} = \und{Fib}\,S$ ou $\und{Fibsc}\,S$). On sait que $f^*_{W_S^{\mathrm{d}}}$ décrit le foncteur $f^*$ sur les types d'homotopie relatifs, et pour prouver a) comme conséquence
\msp{page 125}
de b) (renforcé), b$'$ ou b$''$, il reste à s'assurer que le foncteur $f^*_{W_S}$ décrit de même le foncteur induit par le précédent sur les types d'homotopie relatifs localement constants. Or on a les inclusions canoniques $\mathcal{F}_0 \subseteq \mathcal{F}$, $\mathcal{F}_0' \subseteq \mathcal{F}'$, et sur $\mathcal{F}_0$, $\mathcal{F}_0'$ \texttt{[on a]} $W_S = W_S^{\mathrm{d}} = W_S^{\mathrm{f}}$, et les catégories sont OK pour décrire $\HOT_W^{\mathrm{lc}}(S)$ et $\HOT_W^{\mathrm{lc}}(S')$, et le foncteur $f^*$ induit entre eux par images inverses (\footnote[65]{NB $(\Cat/S)_0 = \und{Fibhot}\,S$, pour les autres cas, c'est clair.}) (\textbf{NB} Le foncteur changement de base $X \longmapsto X \times_S S'$ est compatible avec les localiseurs sur ces catégories, car c'est trivial sous la forme $W_S^{\mathrm{f}}$.) On a alors, par images inverses $X \longmapsto X \times_S S'$, un diagramme commutatif
\[
\begin{tikzcd}
\mathcal{F}_0(W_S)^{-1} \arrow[r, "\approx"] \arrow[d, "(f^*_{W_S})_0"'] & \mathcal{F}(W_S)^{-1} \arrow[d, "f^*_{W_S}"] \\
\mathcal{F}_0'(W_{S'})^{-1} \arrow[r, "\approx"] & \mathcal{F}'(W_{S'})^{-1},
\end{tikzcd}
\]
\printedpage{78}
% head: A. Grothendieck : Dérivateurs, Ch. VI, Édition des Universités de Montpellier II et Paris VII
% foot: Version 9 novembre 2001
où les flèches horizontales sont des \textit{équivalences} de catégories, par le théorème-scholie 1 \textcircled{I}~. Cela prouve que $f^*_{W_S}$ lui aussi est bon pour calculer les images inverses de types d'homotopie relatifs.

On a prouvé que b, b$'$, b$''$ impliquent chacun a), prouvons que a) les implique. Or a) dit exactement que les images inverses $f^*$ des éléments de $\HOT_W^{\mathrm{lc}}(S)$
\msp{page 126}
sont, à isomorphisme près, obtenus par le foncteur image inverse sur $\mathcal{F}$ tout entier, en interprétant $\HOT_W^{\mathrm{lc}}(S)$ comme $\mathcal{F}W_S^{-1}$. Mais cela implique bien que $\mathcal{F} \longrightarrow \mathcal{F}'$ doit être compatible avec les localiseurs. Mais je charrie - on trouve seulement que si $u : X \longrightarrow Y$ dans $\mathcal{F}$ est dans $W_S$, alors la flèche $u' : X' \longrightarrow Y'$ donne un isomorphisme dans $\mathcal{F}'(W_{S'})^{-1}$. Je ne sais pas, d'autre part, si cela implique que $u'$ est dans $W_{S'}$.

\textbf{Corollaire 1.} \textit{Supposons que $p : S' \longrightarrow S$ soit une} Hot\textit{-fibration. Alors le foncteur image inverse $X \longmapsto X \times_S S'$ de} $\Cat/S$ \textit{dans $\mathit{Cat}/S'$ applique $W_S$ dans $W_{S'}$, donc le foncteur $p^* : \HOT_W(S) \longrightarrow \HOT_W(S')$ commute aux foncteurs de localisation $\rho_S$, $\rho_{S'}$ dans ces catégories.} (\textbf{NB} Il faut faire l'hypothèse adéquate sur $W$, savoir L6 bis page 44. C'est OK pour $W_\infty$.)

Le fait que $X \longmapsto X \times_S S'$ soit compatible avec $W_S$, $W_{S'}$ est une tautologie, compte tenu de l'hypothèse L6 bis. Il ne semble pas qu'il soit aussi compatible à $W_S^{\mathrm{d}}$, $W_{S'}^{\mathrm{d}}$ (je devrais faire un contrexemple, avec bien sûr $S'$ non propre sur $S$). Donc on est sous les conditions de b), qui implique a).

\msp{page 127}

Je ne me sens pas capable d'exhiber un foncteur autre qu'une Hot-fibration (ou tout au moins une $W$-fibration), pour lequel le foncteur image inverse sur $\HOT_W$ soit compatible aux foncteurs de localisation $\rho_S$, $\rho_{S'}$. Bien sûr, si $X$ est quelconque dans $\Cat/S$, rien n'empêche de regarder son localisé pour $W_S$, i.e. $\rho_S(\hot_{W,S}(X))$, et de demander comment trouver son image inverse. Si $p : S' \longrightarrow S$ est une Hot-fibration, alors il suffit de prendre $\hot_{W,S'}(X')$, $X' = S \times_S S'$. Dans tout autre cas (j'entends, sans hypothèse sur $X/S$, plus % [sic] X' = S \times_S S' -> X' = X \times_S S', cf. la définition de X' plus haut
précisément sans supposer que $X$ est lui-même Hot-fibré sur $S$), cette formule est canulée sans espoir. On doit, pour calculer $f^*\rho_S(\hot_{W,S}(X))$, `calculer' $\rho_S(\hot_{W,S}(X))$ en `montant à l'étage', comme $\hot_{W,S}(C(X))$. On trouve alors l'image inverse $\hot_{W,S'}(C(X) \times_S S')$ (par \texttt{[le]} corollaire 1, p. 119).

Maintenant que la loi contravariante de $\HOT_W(S)$ en $S \in \mathrm{Ob}\,\Cat$ semble bien comprise, il s'agit de voir qu'on a bien un dérivateur HOT. Je ne m'attends pas à ce que les $\HOT_W^{\mathrm{lc}}(S)$
\msp{page 128}
définissent eux aussi un dérivateur, mais il faudrait m'en assurer expressément. (C'est pour l'existence de $f_!$, $f_*$ que je doit chercher.) % [sic] je doit -> je dois, conjugaison

À dire vrai, il faudrait, pour bien faire, regarder aussi la structure de 2-\textit{foncteur} de
\[
S \longmapsto \HOT_W(S) \qquad \Cat^o \longrightarrow \Cat,
\]


\endgroup
% END UNIT Der-ChVI-345/Der-ChVI-345_p076-080.tex
% BEGIN UNIT Der-ChVI-345/Der-ChVI-345_p081-085.tex
\begingroup\makeatletter%
\NoAutoSpacing%
\newcommand{\printedpage}[1]{\marginpar{\footnotesize p.~#1}}%
\newcommand{\msp}[1]{\marginpar{\footnotesize [#1]}}%
\newcommand{\HOT}{\mathrm{HOT}}%
\newcommand{\Hot}{\mathrm{Hot}}%
\newcommand{\Cat}{\mathrm{Cat}}%
\newcommand{\Ob}{\mathrm{Ob}}%
\newcommand{\hot}{\mathrm{hot}}%
\newcommand{\uHom}{\underline{\mathrm{Hom}}}%
\newcommand{\uFib}{\underline{\mathrm{Fib}}}%
\newcommand{\uFibsc}{\underline{\mathrm{Fibsc}}}%
\newcommand{\uCh}{\underline{\mathrm{Ch}}}%
\newcommand{\circled}[1]{\tikz[baseline=(c.base)]\node[draw,rounded corners=4pt,inner sep=1.5pt](c){#1};}%
\makeatother\let\printedpage\derpp\let\msp\dermsp%
\let\mspin\DERmspin%


\printedpage{79}
% head: A. Grothendieck : Dérivateurs, Ch. VI, Édition des Universités de Montpellier II et Paris VII
% foot: Version 9 novembre 2001

donc pour $S$, $S'$ données dans Cat, définir
\[ \uHom(S',S)^o \longrightarrow \uHom(\HOT_W(S),\HOT_W(S')). \]
Cela revient surtout, pour
\[ f,g : S' \rightrightarrows S \]
deux flèches de $S'$ dans $S$, et un homomorphisme de foncteurs
\[ u : f \longrightarrow g, \]
à décrire un homomorphisme de foncteurs
\[ u^* : g^* \longrightarrow f^* \]
entre $\HOT_W(S)$ et $\HOT_W(S')$. C'est d'ailleurs chose faite, dans le cadre général des $(\mathcal{M},W)$, quand on décrit $\HOT_W(S)$ et $\HOT_W(S')$ par localisation de $\mathrm{Hom}(S^o,\Cat)$ et $\mathrm{Hom}(S'^o,\Cat)$. Ces catégories s'interprètent respectivement comme $\uFibsc\,S$ et $\uFibsc\,S'$, donc comme des catégories de catégories sur $S$ resp. $S'$, on voudrait donc interpréter $u^*$ en termes de telles catégories, et tant qu'à faire, pour des catégories $X$ sur $S$ aussi générales que possible. Nous nous
\msp{page 129}
attendrons alors à devoir faire des hypothèses, soit sur $X$, soit sur $f$ et sur $g$, qui assurent que les images de $\hot_{W,S}(X)$ par $f^*$ et $g^*$ ne sont autres que $\hot_{W,S}(f^*X)$, $\hot_{W,S}(g^*X)$ respectivement, où on pose
\[ f^*(X) = X \times_S (S',f), \qquad g^*(X) = X \times_S (S',g). \]
Donc on devra sans doute se placer dans l'une des hypothèses suivantes:
\begin{enumerate}
\item[a)] $X$ est lisse sur $S$.
\item[b)] $X$ est Hot-fibrée sur $S$.
\item[c)] $f$ et $g$ sont propres (et sous réserve de la validité de la conjecture page 122).
\end{enumerate}
On s'attend alors à ce que $u^*$ soit associé à un homomorphisme
\[ u^* : g^*(X) \longrightarrow f^*(X) \qquad \textit{dans } \Cat/S' \]
défini en termes de $u$. Y a-t-il moyen de le définir bel et bien, sans supposer $X$ au moins fibrée sur $S$? Prenons $S'=e$, donc $f$, $g$ correspondant à des objets $s,t\in\Ob\,S$, et $u$ à une fléche % [sic] fléche -> flèche, wrong accent
\[ u : s \longrightarrow t, \]
les images inverses de $X/S$ par $f$, $g$ sont les fibres $X_s$, $X_t$, et on cherche
\[ u^* : X_t \longrightarrow X_s. \]
On n'a aucune chance d'en trouver canoniquement, si on ne suppose $X$ fibrée sur $S$! Donc ce ne sont pas a), b) ni c) qui peuvent nous servir, mais unique%
\printedpage{80}%
% head: A. Grothendieck : Dérivateurs, Ch. VI, Édition des Universités de Montpellier II et Paris VII
% foot: Version 9 novembre 2001
\msp{page 130}%
ment l'hypothèse $X\in\Ob\,\uFib\,S$. Donc on ne peut espérer trouver une description de $u^*$ entre $\HOT_W(S)$ et $\HOT_W(S')$ en termes d'opérations sur les catégories sur $S$, $S'$, comme suggéré à l'instant, qu'en se bornant à des `modèles' $X$ qui sont dans $\uFib\,S$, tout au moins. Et je ne doute pas en effet que ça marche sur $\uFib\,S$ pas moins bien que directement sur $\uFibsc\,S$, en termes de $\uHom(S^o,\Cat)$.

Pour les axiomes essentiels sur les dérivateurs, je renvoie à I p. 67-68 et p. 111-112. On a vu que Der 1, \circled{Der 2} (\footnote[66]{Canulé pour \underline{Der 2}, je ne sais si c'est vrai même en supposant $W$ \textit{fortement} saturé, i.e. tel que $\hot_W(f)$ isomorphisme $\Longrightarrow f\in W$.}) sont satisfaits pour tout prédérivateur défini par un couple $(\mathcal{M},W\subseteq\mathrm{Fl}\,\mathcal{M})$, donc c'est OK ici pour $\HOT_W$. Il s'ensuit que c'est OK aussi pour $\HOT_W^{\mathrm{lc}}$, vu la caractérisation des $\HOT_W^{\mathrm{lc}}(S)$ en tant que sous-catégories de $\HOT_W(S)$. Pour Der 1, ça marche-t-il même pour des sommes quelconque de catégories, i.e. a-t-on % [sic] quelconque -> quelconques, agreement with sommes; Canulé (footnote 66) is verified authorial wording; see corrections.csv
\[ \HOT_W(\coprod_i S_i) \xrightarrow{\sim} \prod_i \HOT_W(S_i)\,? \]
On a en tous cas
\[ \uFib(\overbrace{\coprod\limits_i S_i}^{S}) \xrightarrow{\sim} \prod_i \uFib(S_i), \]
et le localiseur $W_S^{\mathrm{f}}$ dans $\uFib(S)$ s'identifie
\msp{page 131}
bel et bien au produit des localiseurs $W_{S_i}^{\mathrm{f}}$ dans les $\uFib\,S_i$. Mais dans le cas d'une famille peut-être infinie de catégories $\mathcal{F}_i$, munies de localiseurs $W_i$, si on pose
\[ \mathcal{F} = \prod\mathcal{F}_i, \qquad W = \prod W_i \subseteq \mathrm{Fl}(\mathcal{F}), \]
le foncteur canonique
\[ \mathcal{F}W^{-1} \longrightarrow \prod_i(\mathcal{F}_iW_i^{-1}) \]
est-il toujours une équivalence? Même dans \texttt{[`Même dans' barré]} le cas $I$ fini il faut (il me semble) un argument assez soigneux, soit par la propriété universelle par récurrence sur card\,$I$ (on est ramené au cas de deux foncteurs), soit à coups de catégories de chemins, en notant que
\[ \uCh_\infty^W(\mathcal{F};(s_i),(t_i)) \xrightarrow{\sim} \prod_i \uCh_\infty^{W_i}(\mathcal{F}_i;s_i,t_i), \tag*{$(*)$} \]
et en passant aux $\pi_0$, qui donne
\[ \mathrm{Hom}_W((s_i),(t_i)) \xrightarrow{\sim} \prod_i \mathrm{Hom}_{W_i}(s_i,t_i). \]
Mais la formule $(*)$, visiblement, ne marche que pour $I$ fini. Dans le cas général, on trouve la sous-catégorie $\prod_i^!$ de $\prod_i$ formée des systèmes $(c_i)_{i\in I}$ de chemins qui sont de longueur bornée.

Donc la question de la validité de Der 1 sans hypothèse de finitude sur l'ensemble d'indices $I$
\printedpage{81}
% head: A. Grothendieck : Dérivateurs, Ch. VI, Édition des Universités de Montpellier II et Paris VII
% foot: Version 9 novembre 2001
\msp{page 132}
reste en suspens pour l'instant. (Je soupçonne que c'est vrai, mais il faudrait, pour le prouver, tenir une description des flèches dans $\HOT_W(S)$ par des chemins de longueur bornée, disons de longueur deux ou trois. C'est une question intéressante, mais que je réserve pour plus tard, si le besoin m'y contraint $\ldots\,$)

Plus importante est la question de la validité de Der 3 et de ses variantes (I, p. 111-112). Le minimum qu'on demande d'un dérivateur $\mathbf{D}$, c'est que si $S'\longrightarrow S$ dans Diag est une Hot-équivalence, ce soit une $\mathbf{D}$-équivalence, i.e. (dans le cas actuel) que le foncteur
\[ f^* : \HOT_W(S) \longrightarrow \HOT_W(S') \]
soit pleinement fidèle sur la sous-catégorie des objets \textit{constants} dans $\HOT_W(S)$. Cette condition devient plus tangible quand on dispose déjà soit des foncteurs cohomologie $H^\bullet_{\mathbf{D}}(S,-)$, $H^\bullet_{\mathbf{D}}(S',-)$, soit de l'homologie $H_\bullet^{\mathbf{D}}(S,-)$, $H_\bullet^{\mathbf{D}}(S',-)$. En termes de cette dernière, elle signifie que
\[ p_!p^*(a) \xrightarrow{\sim} p'_!p'^*(a) \qquad \text{\texttt{[est un]} isomorphisme,} \]
\msp{page 133}
pour tout $a$ dans $\mathcal{A}=\mathbf{D}(e)$. (Où $p:S\longrightarrow e$, $p':S'\longrightarrow e$ sont les morphismes structuraux sur la catégorie ponctuelle.) On a vu la formule heuristique générale
\[ p_!\underbrace{p^*(a)}_{a_S} = \hot(S)\times a, \]
où $\hot\,S$ est la classe de $S\in\Ob\,\Cat$ dans $\Hot=(\Cat)W_\infty^{-1}$, et où le produit $\times$ est une opération externe de Hot sur $\mathcal{A}=\mathbf{D}(e)$
\[ \Hot\times\mathcal{A} \longrightarrow \mathcal{A}\ . \]
Dans le cas actuel, $\mathcal{A}=\Hot$, et l'opération
\[ \Hot\times\Hot \longrightarrow \Hot \]
ne doit être autre que le foncteur produit cartésien, lequel est déduit du foncteur similaire dans Cat. Quand on dispose de ce formalisme, le résultat cherché pour un $f:S'\longrightarrow S$ se réduit simplement à dire que si $f$ est dans $W$, alors il en est de même de $f\times\mathrm{id}_X$ pour tout $X$ dans Cat (en prenant $X$ tel que $\hot(X)\simeq a$). Or là c'est un fait bien connu, résultant des axiomes mis sur $W$.

J'ai donc l'impression que je n'ai aucune chance de prouver Der 3a$'$ (p. 111), ni même Der 3a) (cas où $\mathrm{S}=\mathrm{e}$), sans avoir une maîtrise du formalisme des $p_!$. Et j'ai
\printedpage{82}
% head: A. Grothendieck : Dérivateurs, Ch. VI, Édition des Universités de Montpellier II et Paris VII
% foot: Version 9 novembre 2001
\msp{page 134}
bon espoir que cetter maîtrise va s'obtenir à bon compte, alors que la question des $p_*$ semble beaucoup plus délicate. % [sic] cetter -> cette, misspelling

Je vais donc me tourner d'abord vers l'axiome Der 4$'$ d'existence de
\[ f_! : \HOT_W(S') \longrightarrow \HOT_W(S), \]
comme adjoint à gauche de $f^*$. Je vais essayer de construire ad-hoc un foncteur $f_!$, puis les deux flèches d'adjonction
\[ \begin{aligned} \mathrm{id} &\quad\longrightarrow\quad f^*f_! \\ \mathrm{id} &\quad\longleftarrow\quad f_!f^*, \end{aligned} \]
et enfin, de vérifier les deux compatibilités pertinentes.

L'idée géométrique `évidente', c'est pour $X'\in\Cat/S'$, représentant $\xi'\in\Ob\,\HOT_W(S')$, de définir $f_!(\xi')$ par
\[ f_!(\hot_{S',W}(X')) = \hot_{S,W}(X'), \]
où dans le second membre, $X'$ est regardé comme objet de $\Cat/S$, grâce à $f:S'\longrightarrow S$ et au foncteur
\[ f_! : \Cat/S' \longrightarrow \Cat/S \]
qu'il induit. Le seul problème, c'est de voir si ce foncteur est compatible avec les
\msp{page 135}
localiseurs $W_{S'}^{\mathrm{d}}$, $W_S^{\mathrm{d}}$, et sinon, s'il l'est tout au moins sur la sous-catégorie $\uFib\,S'$ ou $\uFibsc\,S'$, qui suffit pour calculer $\HOT_W(S')$ par localisation. Il faut donc voir si un $X'_1\longrightarrow X'_2$ dans $\Cat/S'$, qui est dans $W_{S'}^{\mathrm{d}}$, i.e. qui est dans $W$ et le reste par tout changement de base de la forme $s'\backslash S'\longrightarrow S'$, plus généralement pour tout co-isomorphisme local, est aussi dans $W_S^{\mathrm{d}}$, i.e. si les changements de base de la forme $s\backslash S\longrightarrow S$ le transforment en $W$-équivalences (grossières). Or les co-isomorphismes locaux sont stables par \textit{tout changement de base} (cf. I, p. 55 dans le cas dual des isomorphismes locaux), donc on gagne!
\[
\begin{tikzcd}
\overline{S'} = s\backslash S' \arrow[r] \arrow[d] & S' \arrow[d] \\
\overline{S} = s\backslash S \arrow[r] & S
\end{tikzcd}
\quad .
\]
Je vais quand même le vérifier encore. Si
\[
\begin{tikzcd}
\overline{S'} \arrow[r] \arrow[d] & S' \arrow[d] \\
\overline{S} \arrow[r] & S
\end{tikzcd}
\]
\printedpage{83}
% head: A. Grothendieck : Dérivateurs, Ch. VI, Édition des Universités de Montpellier II et Paris VII
% foot: Version 9 novembre 2001
est un carré cartésien dans Cat, on a dit dans loc. cit., essentiellement, que pour $\overline{s'}$ dans $\overline{S'}$, ayant des images $s'$, $\overline{s}$, $s$ dans $S'$, $\overline{S}$, $S$ respectivement, le carré des colocalisés
\[
\begin{tikzcd}
\overline{s'}\backslash\overline{S'} \arrow[r] \arrow[d] & s'\backslash S' \arrow[d] \\
\overline{s}\backslash\overline{S} \arrow[r] & s\backslash S
\end{tikzcd}
\]
est encore
\msp{page 136}
cartésien. Je viens encore le de revérifier à l'instant, pour les objets tout au moins. Dire que $\overline{S}\longrightarrow S$ est un isomorphisme local, signifie que les $\overline{s}\backslash\overline{S}\longrightarrow s\backslash S$ sont \texttt{[des]} isomorphismes. Mais le carré cartésien précédent montre qu'alors les $\overline{s'}\backslash\overline{S'}\longrightarrow s'\backslash S'$ sont aussi \texttt{[des]} isomorphismes, on gagne! % [sic] le de revérifier -> de le revérifier, word order

\textbf{Théorème 1.} \textit{Pour toute flèche $f:S'\longrightarrow S$ dans} Cat\textit{, le foncteur induit}
\[ f_! : \Cat/S' \longrightarrow \Cat/S \]
\textit{est compatible avec les localiseurs $(W_{S'}^{\mathrm{d}},W_S^{\mathrm{d}})$, donc dualement avec les localiseurs $(W_{S'}^{\mathrm{g}},W_S^{\mathrm{g}})$, et trivialement avec les localiseurs $(W_{S'},W_S)$. Il induit donc un diagramme commutatif de foncteurs sur les localisés}
\[
\begin{tikzcd}
\HOT_W(S') \arrow[r, "\rho_{S'}"] \arrow[d, "f_!"] & \HOT_W^{\mathrm{lc}}(S') \arrow[d, "f_!^{\mathrm{lc}}"] & \HOT_W^o(S') \arrow[l, "\sigma_{S'}"'] \arrow[d, "f_!^o"'] \\
\HOT_W(S) \arrow[r, "\rho_S"] & \HOT_W^{\mathrm{lc}}(S) & \HOT_W^o(S). \arrow[l, "\sigma_S"']
\end{tikzcd}
\]

\textbf{Définition 1.} \textit{On a défini ainsi les foncteurs verticaux $f_!$, $f_!^{\mathrm{lc}}$, $f_!^o=(f^o)_!$, appelés foncteurs image directe homologique, ou foncteurs d'intégration pour $f$ (et les types d'homotopie relatifs).}

\msp{page 137}
On a obtenu en même temps, gratuitement, la compatibilité de ces foncteurs aux foncteurs \texttt{[de]} localisation $\rho_{S'}$, $\sigma_{S'}$ pour $\HOT(S')$, et $\rho_S$, $\sigma_S$ pour $\HOT(S)$. Par contre, on ne s'attendra pas qu'ils commutent avec les foncteurs pleinement fidèles d'inclusion des $\HOT_W^{\mathrm{lc}}$ dans $\HOT_W$. En effet, si on ne suppose pas $f$ dans \underline{Fibhot} (\footnote[67]{Ce cas doit être regardé avec soin ultérieurement.}), il est tout à fait immoral de s'attendre à ce que $f_!$ dans $\HOT_W(S')$ transforme types localement constants en itou. La situation est donc exactement inverse de celle pour le foncteur $f^*$, qui était compatible avec les foncteurs d'inclusion $i_S$, $i_{S'}$ et $j_S$, $j_{S'}$ - c'était à tel point trivial qu'il n'a pas été question même de le dire, voire en faire une proposition! La situation est


\endgroup
% END UNIT Der-ChVI-345/Der-ChVI-345_p081-085.tex
% BEGIN UNIT Der-ChVI-345/Der-ChVI-345_p086-090.tex
\begingroup\makeatletter%
\newcommand{\printedpage}[1]{\marginpar{\footnotesize p.~#1}}%
\newcommand{\msp}[1]{\marginpar{\footnotesize [#1]}}%
\newcommand{\HOT}{\mathrm{HOT}}%
\newcommand{\hot}{\mathrm{hot}}%
\newcommand{\Cat}{\mathrm{Cat}}%
\newcommand{\id}{\mathrm{id}}%
\newcommand{\pr}{\mathrm{pr}}%
\newcommand{\Ob}{\mathrm{Ob}}%
\newcommand{\Wd}{W^{\mathrm{d}}}%
\makeatother\let\printedpage\derpp\let\msp\dermsp%
\let\mspin\DERmspin%


\printedpage{84}
% head: A. Grothendieck : Dérivateurs, Ch. VI, Édition des Universités de Montpellier II et Paris VII
même encore meilleure pour le foncteur $f_!$ entre les HOT, puisque la description de ce foncteur peut se faire directement sur $\Cat/S'$ tout entier, sans avoir à se borner à des sous-catégories plus ou moins \msp{page 138}astucieuses, où le foncteur induit serait compatible avec les localiseurs ($\underline{\mathrm{Liss}}\,S$ et $\underline{\mathrm{Fibhot}}\,S$ dans le cas de $f^*$).

On a de façon immédiate la transitivité de $f_!$, $f_!^{\mathrm{lc}}$, $f_!^{o}$, i.e.
\[
(gf)_! = g_!f_!, \qquad (gf)_!^{\mathrm{lc}} = g_!^{\mathrm{lc}}f_!^{\mathrm{lc}}, \qquad (gf)_!^{o} = g_!^{o}f_!^{o}.
\]

Mais on ne va pas faire bouger $S$, $S'$, $f$ pour le moment, car il faut maintenant, avant toute autre chose, essayer d'établir l'adjonction de $f_!$ et de $f^*$ (dans cet ordre).

Pour le moment, occupons-nous de $\HOT_W$ seulement - le cas de $\HOT_W^{o}$ s'en déduira en appliquant ce qu'on saura à $f^{o}$. Et on s'occupera de $\HOT_W^{\mathrm{lc}}$ plus tard.

1) \textit{L'homomorphisme}
\[
\alpha : \id_{\HOT_W(S')} \longrightarrow f^*f_!.
\]

Si $X'$ est dans $\Cat/S'$, $f_!(\hot_{S',W}(X'))$ est $\hot_{S,W}(X')$, mais on ne peut trouver son image inverse $f^*$ sur $S'$ en prenant l'image inverse catégorique. Car même si on fait des hypothèses draconiennes sur $X'$ en tant qu'objet de $\Cat/S'$ (p.ex. $X'$ fibrée sur $S'$), celles-ci seront réduites à néant quand on aura composé avec $f$, sur lequel nous ne voulons ici faire aucune hypothèse. Il \msp{page 139}faut donc d'abord appliquer le foncteur $\Phi_S$ à $X'$ sur $S$, puis prendre son image inverse sur $S'$ :
\[
f^*f_!(\hot_{S',W}(X')) = \hot_{S',W}(\Phi_S(X'/S)\times_S S').
\]
Il faut donc trouver une flèche
\[
\alpha_{X'} : X' \longrightarrow \Phi_S(X'/S)\times_S S'
\]
dans le localisé $(\Cat/S')(\Wd_{S'})^{-1}$, voire même une flèche dans $\Cat/S'$ lui-même. Mais par les propriétés d'adjonction des foncteurs $f_!$ et $f^*$ au \textit{niveau catégorique,} une telle flèche revient à une flèche
\[
f_!(X') \longrightarrow \Phi_S(X'/S)
\]
dans $\Cat/S$, i.e. à une flèche qui rend commutatif le carré
\[
\begin{tikzcd}
X' \arrow[r] \arrow[d] & \Phi_S(X'/S') \arrow[d] \\ % [sic] \Phi_S(X'/S') -> \Phi_S(X'/S), the arrow f_!(X') -> \Phi_S(X'/S) just above has X'/S
S' \arrow[r] & S.
\end{tikzcd}
\]
% foot: Version 9 novembre 2001

\printedpage{85}
% head: A. Grothendieck : Dérivateurs, Ch. VI, Édition des Universités de Montpellier II et Paris VII
Or on a une telle flèche canonique, savoir $i_{X'/S}$. Elle est fonctorielle en $X'$ sur $S$ (sans avoir même à se soucier si $X'$ sur $S$ provient d'un $X'$ sur $S'$, i.e. si la flèche structurale $X'\longrightarrow S$ se factorise par $S'\longrightarrow S$).

\leavevmode\msp{page 140}2) \textit{L'homomorphisme}
\[
\beta : \id_{\HOT_W(S)} \longleftarrow f_!f^*.
\]

On peut décrire $\HOT_W(S)$ par $(\underline{\mathrm{Fib}}\,S)(\Wd_S)^{-1}$ ou $(\underline{\mathrm{Liss}}\,S)(\Wd_S)^{-1}$ par exemple, de sorte que $f^*$ a une description simple par
\[
f^*(\hot_{S,W}(X)) = \hot_{S',W}(X\times_S S')
\]
(pour $X$ lisse sur $S$). On a ensuite
\[
f_!f^*(\hot_{S,W}(X)) = \hot_{S,W}(X\times_S S'),
\]
mais attention, $X\times_S S'$ n'est pas a priori dans $\underline{\mathrm{Liss}}\,S$, seulement dans $\Cat/S$ sans plus. Mais cela n'empêche qu'il donne ce qu'il faut, par localisation $(\Wd_S)^{-1}$ dans $\Cat/S$. Pour définir $\beta_X$, il suffit donc de trouver une flèche
\[
\beta_X : X \longleftarrow X\times_S S'
\]
dans $\Cat/S$, et on voit bien $\pr_1^X$ (\footnote[68]{quand $\xi\in\HOT_W(S)$ est défini par $X$ dans $\Cat/S$ non nécessairement lisse, on doit trouver $\Phi_S(X)\times_S S'\longrightarrow X$ dans $(\Cat/S)(\Wd_S)^{-1}$, on prend $(i_X)^{-1}\pr_1^{\Phi_S(X)}$ (cf. page 143).}). On a donc
\[
\left\{
\begin{array}{lcl}
\alpha : \id & \longrightarrow & f^*f_! \\
\beta : \id & \longleftarrow & f_!f^*
\end{array}
\right.
\]

3) \textit{La compatibilité}
\[
\begin{tikzcd}
f_! \arrow[r, "f_!*\alpha"] \arrow[rr, bend right, "\id"'] & f_!f^*f_! \arrow[r, "\beta*f_!"] & f_!.
\end{tikzcd}
\]

On doit donc prendre $X'$ sur $S'$, et expliciter les flèches à partir de l'objet
\[
\xi' = \hot_{S',W}(X').
\]
% foot: Version 9 novembre 2001

\printedpage{86}
% head: A. Grothendieck : Dérivateurs, Ch. VI, Édition des Universités de Montpellier II et Paris VII
\leavevmode\msp{page 141}On a
\[
f_!(\xi') = \hot_{S,W}(X'/S),
\]
et pour calculer $f^*f_!(\xi')$, il faut comme tantôt prendre $\Phi_S(X')$ et on trouve
\[
f^*f_!(\xi') = \hot_{S',W}(\Phi_S(X')\times_S S').
\]
Enfin,
\[
f_!f^*f_!(\xi') = \hot_{S,W}(\Phi_S(X')\times_S S').
\]
Il faut expliciter les flèches
\[
\begin{tikzcd}[row sep=small]
\underbrace{f_!(\xi')} \arrow[r, "{f_!(\alpha_{\xi'})}"] & \underbrace{f_!f^*f_!(\xi')} \arrow[r, "{\beta_{f_!(\xi')}}"] & \underbrace{f_!(\xi')}. \\
X'/S \arrow[r, "{\alpha_{X'}}", "{\text{déduit de } i_{X'}}"'] & \Phi_S(X')\times_S S' \arrow[r, dotted, "{\beta_{X'}}"] & X'/S \\
 & \text{sur } S &
\end{tikzcd}
\tag*{$(*)$}
\]

$\alpha_{X'}$ est déduit de $i_{X'} : X'\longrightarrow\Phi_S(X')$ comme l'unique $S'$-morphisme rendant commutatif le diagramme
\[
\begin{tikzcd}
X' \arrow[r, "{\alpha_{X'}}"] \arrow[d] \arrow[rr, bend left, "{i_{X'}}"] & \Phi_S(X')\times_S S' \arrow[r, "\pr_1"] \arrow[dl, dotted] & \Phi_S(X') \arrow[d] \\ % Reading verified against PDF leaf 48: the curved i_X' arrow is solid.
S' \arrow[rr] & & S.
\end{tikzcd}
\tag*{$(**)$}
\]

Pour décrire $\beta_{f_!(\xi')}$, où $f_!(\xi') = \hot_{S,W}(X')$, il faut donner une description générale de $\beta_\xi$ qui soit valable pour \textit{tout} représentant $X$ de $\xi$, pas seulement dans le cas $X$ lisse sur $S$. Pour calculer $f^*(\xi)$, il faut donc prendre le $\hot_{S',W}$ de $\Phi(X)\times_S S'$, puis pour $f_!f^*(\xi)$, il \msp{page 142}faut prendre le $\hot_{S,W}$ du même $\Phi_S(X)\times_S S'$. Il faut définir, si on peut, une flèche
\[
\Phi_S(X)\times_S S' \longrightarrow X \qquad\qquad \text{dans } \Cat/S,
\]
ou du moins dans $(\Cat/S)(\Wd_S)^{-1}$. Mais on a
\[
i_X : X \longrightarrow \Phi_S(X)
\]
qui est dans $\Wd_S$, et qui par produit cartésien définit
\[
i_X\times_S\id_{S'} : X\times_S S' \longrightarrow \Phi_S(X)\times_S S'
\]
(en d'autres termes on a pris l'image inverse de $i_X$ par $S'\longrightarrow S$). Cette flèche est-elle encore dans $\Wd_S$? (\footnote[69]{Non, elle ne l'est pas, car si $S'=\{s\}$, cela signifierait que $X'_s\longrightarrow s\backslash X$ est dans $W$, ce qui est faux en général.}). Si oui, on a le diagramme
% foot: Version 9 novembre 2001
\printedpage{87}
% head: A. Grothendieck : Dérivateurs, Ch. VI, Édition des Universités de Montpellier II et Paris VII
\[
\begin{tikzcd}
\Phi_S(X)\times_S S' \arrow[r, dotted, "\beta_X ?"] & X \\
X\times_S S' \arrow[u, "{i_X\times_S\id_{S'}}", "{=\pi_X}"'] \arrow[ur, "\pr_1^X"] &
\end{tikzcd}
\]
et on définit la flèche cherchée $\beta_X$ \textit{dans} $(\Cat/S)(\Wd_S)^{-1}$ par la commutativité de ce triangle, donc comme
\[
\beta_X = \pr_1^X\circ\underbrace{\pi_X^{-1}}_{\text{inverse dans }(\Cat/S)(\Wd_S)^{-1}} \qquad\qquad .
\]

\leavevmode\msp{page 143}Mais ça ne marche pas - cf. annotation! On va faire alors l'inverse
\[
\begin{tikzcd}
\Phi_S(X)\times_S S' \arrow[r, "\pr_1^{\Phi_S(X)}"] \arrow[dr, dotted, "\beta_X ?"'] & \Phi_S(X) \\
 & X \arrow[u, "{i_X \qquad \text{dans } \Wd_S}"']
\end{tikzcd}
\]

(\footnote[70]{Description de $\beta_\xi$ pour $\xi=\hot_{S,W}(X)$, avec $X$ quelconque dans $\Cat/S$, à expliciter dans définition en forme.}), on définit $\beta_X$ comme $(i_X)^{-1}\pr_1^{\Phi_S(X)}$, où $(i_X)^{-1}$ est défini dans $(\Cat/S)(\Wd_S)^{-1}$.

Pour prouver la commutativité de $(*)$ (p. 141), comme $i_{X'} : X'\longrightarrow\Phi_S(X')$ est dans $\Wd_S$, il suffit de prouver, au lieu de $\beta_{X'}\alpha_{X'} = \id_{X'}$, la relation déduite en composant avec $i_{X'} : X'\longrightarrow\Phi_S(X')$
\[
\underbrace{i_{X'}\beta_{X'}}_{\pr_1^{\Phi_S(X')}}\alpha_{X'} = i_{X'},
\]
i.e.
\[
\pr_1^{\Phi_S(X')}\circ\alpha_{X'} = i_{X'},
\]
ce qui n'est autre que la relation de définition de $\alpha_{X'}$ (cf. $(**)$ p. 141), on gagne!

4) \textit{Compatibilité}
\[
\begin{tikzcd}
f^* \arrow[r, "\alpha*f^*"] \arrow[rr, bend right, "\id"'] & f^*f_!f^* \arrow[r, "f^**\beta"] & f^*.
\end{tikzcd}
\]
% foot: Version 9 novembre 2001

\printedpage{88}
% head: A. Grothendieck : Dérivateurs, Ch. VI, Édition des Universités de Montpellier II et Paris VII
\leavevmode\msp{page 144}Partons donc d'un $\xi\in\Ob\HOT_W(S)$, qu'on peut supposer défini par $X\in\Ob\underline{\mathrm{Fib}}\,S$. Alors
\[
\begin{aligned}
f^*(\xi) &\;=\; \hot_{S',W}(X\times_S S'),\\
f_!f^*(\xi) &\;=\; \hot_{S,W}(X\times_S S'),\\
f^*f_!f^*(\xi) &\;=\; \hot_{S',W}(\Phi_S(X\times_S S')\times_S S').
\end{aligned}
\]
Il faut expliciter maintenant
\[
\underbrace{f^*(\xi)}_{\hot_{S',W}(X')} \xrightarrow{\alpha_{f^*(\xi)}} \underbrace{f^*f_!f^*(\xi)}_{\hot_{S',W}(\Phi_S(X')\times_S S')} \xrightarrow{f^*(\beta_\xi)} \underbrace{f^*(\xi)}_{\hot_{S',W}(X')} \quad ,
\]
et vérifier que le composé est l'identité. On a posé
\[
X' = X\times_S S',
\]
et le diagramme est donc aussi
\[
\begin{tikzcd}
X' \arrow[r, "{\alpha_{X'}}"] \arrow[rr, bend left, "{\id_{X'}}"] & \Phi_S(X')\times_S S' \arrow[r, "f^*(\beta_\xi)"] & X' % Reading verified against PDF leaf 50: the curved identity arrow is solid.
\end{tikzcd}
\]
dans $(\Cat/S')(\Wd_{S'})^{-1}$. Pour vérifier la commutativité, on peut composer avec
\[
X\times_S S' = X' \xrightarrow{i_X\times_S\id_{S'}} \Phi_S(X')\times_S S' \qquad\qquad \text{(où } i_X : X\longrightarrow\Phi_S(X)\text{)}
\]
qui est dans $\Wd_{S'}$ (\footnote[71]{\textbf{NB} Comme $X$ est lisse$/S$, et $\Phi_S(X)$ aussi, $X\longrightarrow\Phi_S(X)$ reste une $W$-équivalence colocale par tout changement de base.}), or on vérifie, sur la construction générale de $\beta_\xi=\beta_X$ (page précédente), que l'on a commutativité dans
\[
\begin{tikzcd}
\Phi_S(X')\times_S S' \arrow[r, "f^*(\beta_\xi)"] \arrow[dr, "\Phi_S(\pr_1^X)\times_S S'"'] & X' \arrow[d, "i_X\times_S S'"] \\
 & \Phi_S(X)\times_S S'.
\end{tikzcd}
\]

\leavevmode\msp{page 145}On est ramené à la commutativité de
\[
\begin{tikzcd}
X\times_S S' = X' \arrow[r, "{\alpha_{X'}}"] \arrow[dr, "i_X\times_S S'"', "?"] & \Phi_S(X')\times_S S' \arrow[d, "\Phi_S(\pr_1^X)\times_S S'"] \\
 & \Phi_S(X)\times_S S'
\end{tikzcd}
\tag*{$(*)$}
\]
% foot: Version 9 novembre 2001


\endgroup
% END UNIT Der-ChVI-345/Der-ChVI-345_p086-090.tex
% BEGIN UNIT Der-ChVI-345/Der-ChVI-345_p091-095.tex
\begingroup\makeatletter%
\newcommand{\printedpage}[1]{\marginpar{\footnotesize p.~#1}}%
\newcommand{\msp}[1]{\marginpar{\footnotesize [#1]}}%
\newcommand{\Cat}{\mathrm{Cat}}%
\newcommand{\HOT}{\mathrm{HOT}}%
\newcommand{\hot}{\mathrm{hot}}%
\newcommand{\pr}{\mathrm{pr}}%
\newcommand{\id}{\mathrm{id}}%
\newcommand{\Ob}{\mathrm{Ob}\,}%
\newcommand{\Hom}{\mathrm{Hom}}%
\newcommand{\uHom}{\underline{\mathrm{Hom}}}%
\newcommand{\Liss}{\underline{\mathrm{Liss}}\,}%
\newcommand{\cA}{\mathcal{A}}%
\newcommand{\dd}{\mathrm{d}}%
\newcommand{\lc}{\mathrm{lc}}%
\newcommand{\adj}[2]{\mathrel{\underset{#2}{\overset{#1}{\rightleftarrows}}}}%
\makeatother\let\printedpage\derpp\let\msp\dermsp%
\let\mspin\DERmspin%


\printedpage{89}
% head: A. Grothendieck : Dérivateurs, Ch. VI, Édition des Universités de Montpellier II et Paris VII

dans $(\Cat/S')(W^{\dd}_{S'})^{-1}$. Essayons de prouver que c'est commutatif dans $\Cat/S'$, pas seulement dans le localisé par $W^{\dd}_{S'}$. On peut le regarder comme diagramme dans $\Cat/S$ et prouver la commutativité à ce titre. Cette relation de commutativité $vu = w$ se décompose alors en deux, suivant les deux facteurs $\Phi_S(X)$ et $S'$ de l'objet but $\Phi_S(X) \times_S S'$. Suivant le facteur $S'$, la relation signifie simplement que $\alpha_{X'}$ est un $S'$-morphisme, c'est OK. Pour le facteur $\Phi_S(X)$, on prolonge le diagramme $(*)$ à droite par le carré
\[
\begin{tikzcd}
X' \arrow[r, dotted, "\alpha_{X'}"] \arrow[dr, dotted] & \Phi_S(X') \times_S S' \arrow[r, "\pr_1^{\Phi_S(X')}"] \arrow[d, "\Phi_S(\pr_1 X)\times_S S'"] & \Phi_S(X') \arrow[d, "\Phi_S(\pr_1^X)"] \\
 & \Phi_S(X) \times_S S' \arrow[r, "\pr_1^{\Phi_S(X)}"] & \Phi_S(X),
\end{tikzcd}
\] % [sic] \Phi_S(\pr_1 X)\times_S S' -> \Phi_S(\pr_1^X)\times_S S', the X is printed on the line instead of as superscript of \pr_1
\msp{page 146}
et compte tenu que le composé de la première ligne horizontale est $i_{X'}$, par construction de $\alpha_{X'}$, on est ramené à la commutativité de
\[
\begin{tikzcd}
X \times_S S' = X' \arrow[r, "i_{X'}"] \arrow[d, "i_X\times_S S'"'] & \Phi_S(X') \arrow[d, "\Phi_S(\pr_1^X)"] \\
\Phi_S(X) \times_S S' \arrow[r, "\pr_1^{\Phi_S(X)}"] & \Phi_S(X).
\end{tikzcd}
\tag*{$(*)$}
\]

Or le morphisme $i_X : X \longrightarrow \Phi_S(X)$ dans $\Cat/S$ donne lieu au carré commutatif (comme toute flèche $i : X \longrightarrow Y$)
\[
\begin{tikzcd}
X \times_S S' \arrow[r, "\pr_1^X"] \arrow[d, "i_X\times_S S'"'] & X \arrow[d, "i_X"] \\
\Phi_S(X) \times_S S' \arrow[r, "\pr_1^{\Phi_S(X)}"] & \Phi_S(X),
\end{tikzcd}
\]
de sorte que le composé inférieur de $(*)$ n'est autre que celui de $(**)$, donc égal au composé supérieur de ce dernier, et la commutativité de $(*)$ revient à celle de
\[
\begin{tikzcd}
X \times_S S' = X' \arrow[r, "i_{X'}"] \arrow[d, "\pr_1^X"'] & \Phi_S(X') \arrow[d, "\Phi_S(\pr_1^X)"] \\
X \arrow[r, "i_X"] & \Phi_S(X),
\end{tikzcd}
\]
qui est évidente par fonctorialité de $Z \xrightarrow{i_Z} \Phi_S(Z)$, appliquée à la flèche $\pr_1^X : X' \longrightarrow X$ de $\Cat/S$.

% foot: Version 9 novembre 2001

\printedpage{90}
% head: A. Grothendieck : Dérivateurs, Ch. VI, Édition des Universités de Montpellier II et Paris VII

\msp{page 147}
On a ainsi prouvé que les foncteurs
\[
\HOT_W(S') \adj{f_!}{f^*} \HOT_W(S)
\tag*{$(*)$}
\]
sont adjoints l'un de l'autre, avec des morphismes d'adjonction
\[
\left\{
\begin{array}{lcl}
\id & \xrightarrow{\ \alpha\ } & f^*f_! \\[1ex]
\id & \xleftarrow{\ \beta\ } & f_!f^*
\end{array}
\right.
\]
précisés plus haut. Je présume que la même démonstration, en travaillant avec les localiseurs grossiers $W_S$, $W_{S'}$, au lieu de $W^{\dd}_S$, $W^{\dd}_{S'}$, donnerait que les foncteurs correspondants
\[
\HOT^{\lc}_W(S') \adj{f_!^{\lc}}{f^*_{\lc}} \HOT^{\lc}_W(S)
\tag*{$(**)$}
\]
sont adjoints l'un de l'autre, avec essentiellement ‘les mêmes’ morphismes d'adjonction $\alpha$, $\beta$. Je n'ai pas envie de reprendre les raisonnements précédents et préfère le \textit{déduire} de l'adjonction dans la situation $(*)$. Cela me donnera l'occasion de mieux comprendre les traits formels essentiels de la situation. Je poserai pour simplifier
\msp{page 148}
\[
\begin{array}{rcl@{\qquad\qquad}rcl}
\HOT_W(S) & = & \cA, & \HOT_W(S') & = & \cA', \\
\HOT^{\lc}_W(S) & = & \cA_0, & \HOT^{\lc}_W(S') & = & \cA'_0,
\end{array}
\]
et j'utilise les couples de foncteurs adjoints
\[
\cA \adj{\rho}{i} \cA_0, \qquad\qquad \cA' \adj{\rho'}{i'} \cA'_0,
\]
où $\rho$, $\rho'$ sont des foncteurs de localisation, i.e. $i$, $i'$ des foncteurs pleinement fidèles. Je dispose d'un couple de foncteurs adjoints
\[
\cA' \adj{f_!}{f^*} \cA,
\]
et je suppose de plus

\begin{enumerate}
\item[a)] $f_!$ est compatible avec les localiseurs de $\cA'$, $\cA$ définis par $\rho'$, $\rho$ respectivement, donc définit un foncteur (bien déterminé à isomorphisme unique près)
\[
f_!^0 : \cA'_0 \longrightarrow \cA_0,
\]
donnant donc lieu à un diagramme commutatif (à isomorphisme fonctoriel près)

% foot: Version 9 novembre 2001

\printedpage{91}
% head: A. Grothendieck : Dérivateurs, Ch. VI, Édition des Universités de Montpellier II et Paris VII
\[
\begin{tikzcd}
\cA' \arrow[r, "\rho'"] \arrow[d, "f_!"'] & \cA'_0 \arrow[d, "f_!^0"] \\
\cA \arrow[r] & \cA_0.
\end{tikzcd}
\tag*{(I)}
\]
\texttt{[plutôt} $\cA \xrightarrow{\rho} \cA_0$\texttt{]}

\item[b)] ‘Dualement’, $f^* : \cA \longrightarrow \cA'$ est compatible avec les sous-catégories pleines images essentielles des foncteurs pleinement fidèles $i$, $i'$, et induit
\msp{page 149}
alors un foncteur, défini à isomorphisme unique près :
\[
f_0^* : \cA_0 \longrightarrow \cA'_0,
\]
donnant lieu à un carré de foncteurs, commutatif à isomorphisme fonctoriel canonique près
\[
\begin{tikzcd}
\cA_0 \arrow[r, "i"] \arrow[d, "f_0^*"'] & \cA \arrow[d, "f^*"] \\
\cA'_0 \arrow[r, "i'"] & \cA'.
\end{tikzcd}
\tag*{(II)}
\]
\end{enumerate}

\textit{Sous ces conditions, le couple de foncteurs}
\[
\cA_0 \adj{f_!^0}{f_0^*} \cA'_0, \qquad \text{\texttt{[plutôt} } \cA'_0 \adj{f_!^0}{f_0^*} \cA_0, \text{\texttt{]}}
\]
\textit{est un couple de foncteurs adjoints.} (NB quand je parle de couples de foncteurs adjoints, celui écrit au-dessus de l'autre est l'adjoint à gauche.)

\textsc{Démonstration}. Soient
\[
\eta_0 \in \Ob \cA_0, \qquad\qquad \xi'_0 \in \Ob \cA'_0,
\]
il faut définir une bijection bifonctorielle
\[
\Hom_{\cA_0}(f_!^0(\xi'_0), \eta_0) \simeq \Hom_{\cA'_0}(\xi'_0, f_0^*(\eta_0))
\]
Pour cela, on va écrire $\xi'_0$ sous la forme
\[
\xi'_0 = \rho'(\xi')\,, \qquad\qquad \xi' \in \Ob \cA',
\]
il suffira d'établir la bijection bifonctorielle en remplaçant $\xi'_0$ par $\rho'(\xi')$, i.e.
\[
\Hom_{\cA_0}(f_!^0(\rho'(\xi')), \eta_0) \simeq \Hom_{\cA'_0}(\rho'(\xi'), f_0^*(\eta_0))
\]

% foot: Version 9 novembre 2001

\printedpage{92}
% head: A. Grothendieck : Dérivateurs, Ch. VI, Édition des Universités de Montpellier II et Paris VII

\msp{page 150}
(En effet, la formule d'adjonction s'interprète comme un isomorphisme de bifoncteurs, ou, ce qui revient au même, de foncteurs $\cA'_0 \mathrel{\underset{\psi}{\overset{\varphi}{\rightrightarrows}}} \cA_0^{\wedge}$, or
\[
\begin{array}{ccc}
\uHom(\cA'_0, \cA_0^{\wedge}) & \longrightarrow & \uHom(\cA', \cA_0^{\wedge}) \\[1ex]
\varphi & \longmapsto & \varphi \circ \rho'
\end{array}
\]
est pleinement fidèle, puisque $\rho'$ est un foncteur de localisation, donc il revient au même de se donner un isomorphisme $\varphi \simeq \psi$, ou $\varphi \circ \rho' \simeq \psi \circ \rho'$.) On doit donc transformer l'expression
\[
\begin{array}{cccc}
 & \Hom_{\cA_0}(\underbrace{f_!^0(\rho'(\xi'))}_{\rho f_!(\xi')}, \eta_0) & \underbrace{\simeq}_{\text{\footnotesize commutativité de (I)}} & \Hom_{\cA_0}(\rho f_!(\xi'), \eta_0) \\[3ex]
\underbrace{\simeq}_{\text{\footnotesize adjonction $\rho$, $i$}} & \Hom_{\cA}(f_!(\xi'), i(\eta_0)) & \underbrace{\simeq}_{\text{\footnotesize adjonction $f_!$, $f^*$}} & \Hom_{\cA'}(\xi', \underbrace{f^*i(\eta_0)}_{i'(f_0^*(\eta_0))}) \\[3ex]
\underbrace{\simeq}_{\text{\footnotesize commutativité de (II)}} & \Hom_{\cA'}(\xi', i'(f_0^*(\eta_0))) & \underbrace{\simeq}_{\text{\footnotesize adjonction $\rho'$, $i'$}} & \Hom_{\cA'_0}(\underbrace{\rho'(\xi')}_{\xi'_0}, f_0^*(\eta_0)),
\end{array}
\]
ce qui est la bijection annoncée.

Je résume l'essentiel des résultats obtenus sur $f_!$, $f_!^{\lc}$ :

\textbf{Théorème 2.} \textit{Soit $f : S' \longrightarrow S$ une flèche dans} $\Cat$\textit{, considérons les foncteurs}
\[
\begin{aligned}
\HOT_W(S') &\adj{f_!}{f^*} \HOT_W(S) \\
\HOT^{\lc}_W(S') &\adj{f_!^{\lc}}{f^*_{\lc}} \HOT^{\lc}_W(S),
\end{aligned}
\]
\msp{page 151}
\textit{où les foncteurs $f_!$, $f_!^{\lc}$ sont décrits dans le théorème} 1 \textit{(p. 136). Les morphismes d'adjonction}
\[
\left\{
\begin{array}{lcl}
\id & \longrightarrow & f^*f_! \\[1ex]
\id & \longleftarrow & f_!f^*
\end{array}
\right.
\]
\textit{sont décrits p. 138 à 140.}

On voudrait expliciter la formule sympathique, obtenue en passant
\[
\boxed{\;f_!f^*(\xi) = \hot_{S,W}(S'/S) \times_S \xi\;}\ ,
\]
pour ceci il faut expliciter le deuxième membre. Notons la

\textbf{Proposition 4.} \textit{Il existe un unique foncteur}
\[
(\xi, \eta) \longmapsto \xi \times_S \eta : \HOT_W(S) \times \HOT_W(S) \longrightarrow \HOT_W(S)
\]

% foot: Version 9 novembre 2001

\printedpage{93}
% head: A. Grothendieck : Dérivateurs, Ch. VI, Édition des Universités de Montpellier II et Paris VII

\textit{avec la propriété que pour $X$, $Y$ variables dans} $\Liss S$\textit{, on ait un isomorphisme bifonctoriel en $X$, $Y$}
\[
\hot_{S,W}(X) \times_S \hot_{S,W}(Y) \simeq \hot_{S,W}(X \times_S Y)
\tag*{$(*)$}
\]
(\footnote[72]{\textbf{NB} La multiplication $\times_S$ dans $\HOT_W(S)$ est associative, commutative, unitaire \textit{etc., etc.}, puisque c'est le cas dans $\Liss S$.}).

C'est formel, en utilisant seulement que sur $\Liss S$, le bifoncteur
\[
(X, Y) \longrightarrow X \times_S Y
\]
est compatible avec les localiseurs $W^{\dd}_S \times W^{\dd}_S$ et $W^{\dd}_S$. Pour le voir, on est ramené à voir que si $X \xrightarrow{u} X'$ dans $\Liss S$ est dans $W^{\dd}_S$, alors il en est de même de $X \times_S Y \xrightarrow{\ u\times_S \id_Y\ } X' \times_S Y$, pour tout
\msp{page 152}
$Y$ dans $\Liss S$. En fait, on n'utilise que la lissité de $X$, $X'$, pas celle de $Y$, et c'est immédiat, car si $X \longrightarrow X'$ est dans $W^{\dd}_S$, $X \times_S Y \longrightarrow X' \times_S Y$ est dans $W^{\dd}_Y$, donc aussi dans $W^{\dd}_S$ (cf. théorème 1 appliqué à $Y \longrightarrow S$). Cette remarque suggère en même temps :

\textbf{Corollaire 1. (Faux)} \textit{La formule bifonctorielle} $(*)$ \textit{reste valable pour $X$ dans} $\Liss S$\textit{, $Y$ quelconque dans} $\Cat/S$\textit{, et aussi dans le cas inverse.} (\footnote[73]{Cet énoncé est sûrement \textit{faux!} Contrexemple p. 153, cf. corollaire 1 rectifié page 154-156.})

Dans le premier cas, considérons
\[
Y \xrightarrow{i_Y} \Phi(Y)
\]
dans $W^{\dd}_S$, donc un isomorphisme dans $(\Cat/S)(W^{\dd}_S)^{-1}$, on aura donc
\[
\begin{array}{lcl}
\hot_{S,W}(X) \times_S\ \hot_{S,W}(Y) & \xrightarrow{\ \sim\ } & \hot_{S,W}(X) \times_S\ \hot_{S,W}(\Phi(Y)) \\[1ex]
 & \underbrace{\xrightarrow{\ \sim\ }}_{\text{\footnotesize car $X$ \textit{et} $\Phi(Y)$ lisses sur $S$}} & \hot_{S,W}(X \times_S \Phi(Y)),
\end{array}
\]
or je dis que
\[
X \times_S Y \xrightarrow{\ X\times_S i_Y\ } X \times_S \Phi(Y)
\]
est encore dans $W^{\dd}_S$. Il suffit de voir que c'est dans $W^{\dd}_X$ (cf. théorème 1), et pour ceci, que le changement de base lisse $X \longrightarrow S$ induit $\Cat/S \longrightarrow \Cat/X$ compatible avec les localiseurs $W^{\dd}_S$, $W^{\dd}_X$. Mais c'est \textit{faux} (prendre $X = \{s\}$).

\msp{page 153}
Montrons que les corollaire est faux. Je vais prendre $Y = \{s\}$, catégorie ponctuelle en $s$, alors $\Phi(Y)$ est la catégorie fibrée à fibres discrètes définie par le foncteur représentable $F$ défini par $s$, donc c'est $S/F$ et, en l'occurence $S/s$: % [sic] les corollaire -> le corollaire, agreement; occurence -> occurrence, spelling
\[
\Phi(\{s\}) \simeq S/s \qquad\quad \text{comme catégorie sur } S.
\]

% foot: Version 9 novembre 2001


\endgroup
% END UNIT Der-ChVI-345/Der-ChVI-345_p091-095.tex
% BEGIN UNIT Der-ChVI-345/Der-ChVI-345_p096-100.tex
\begingroup\makeatletter%
\newcommand{\printedpage}[1]{\marginpar{\footnotesize p.~#1}}%
\newcommand{\msp}[1]{\marginpar{\footnotesize [#1]}}%
\newcommand{\hot}{\mathrm{hot}}%
\newcommand{\HOT}{\mathrm{HOT}}%
\newcommand{\Cat}{\mathrm{Cat}}%
\newcommand{\Liss}{\underline{\mathrm{Liss}}}%
\newcommand{\multhot}{\mathrm{multhot}}%
\newcommand{\mult}{\mathrm{mult}}%
\makeatother\let\printedpage\derpp\let\msp\dermsp%
\let\mspin\DERmspin%


\printedpage{94}
% head: A. Grothendieck : Dérivateurs, Ch. VI, Édition des Universités de Montpellier II et Paris VII
Si p.ex. $s$ est élément final $e$ de $S$, on trouve
\[\Phi(\underset{S}{e}) \simeq S,\]
donc
\[\hot_{S,W}(X)\times_S \hot_{S,W}(e)=\hot_{S,W}(X),\]
d'autre part
\[\hot_{S,W}(X\times_S e)=\hot_{S,W}(X_e).\]

\msp{154}
Si les deux étaient isomorphes dans $\HOT_W(S)$ par la flèche canonique
\[X_e\longrightarrow X,\]
en appliquant le foncteur $p_!$ (où $p : S\longrightarrow e$), il s'ensuiverait que $X_e\longrightarrow X$ est dans $W$. Si par exemple $S=\boldsymbol{\Delta}^1$, cela impliquerait que pour tout $X$ lisse sur $S$, $X_1\longrightarrow X$ est dans $W$, alors que normalement la lissité implique que $X_0\longrightarrow X$ est dans $W$. Donc ça reviendrait à dire que $\Liss\,\boldsymbol{\Delta}^1=\Liss_0\boldsymbol{\Delta}^1$, d'où aussitôt $\Liss\,S=\Liss_0S$, pour tout $S$ dans Cat, ça se saurait! Le contrexemple le plus trivial est de prendre $X=\{0\}\hookrightarrow S=\boldsymbol{\Delta}^1$, l'inclusion de l'ouvert ‘générique’, immersion ouverte, donc lisse, la fibre $X_1$ est vide! % [sic] s'ensuiverait -> s'ensuivrait, nonstandard verb form

% marginal mark: handwritten sign in the left margin beside the following statement
\textbf{Corollaire 1 rectifié.} \textit{Si $X$ est dans} $\Cat/S$\textit{, pas nécessairement lisse sur $S$, alors le foncteur}
\[\mult_{X/S} : Y\longmapsto Y\times_S X : \Cat/S\longrightarrow \Cat/S\]
\textit{induit un foncteur}
\[\Liss\,S\longrightarrow \Cat/S\]
\textit{compatible avec les localiseurs $W_S^{\mathrm{d}}$, $W_S^{\mathrm{d}}$, donc induit un foncteur}
\[(\Liss\,S)(W_S^{\mathrm{d}})^{-1}\longrightarrow(\Cat/S)(W_S^{\mathrm{d}})^{-1},\]

\msp{155}
\textit{ou encore (compte tenu de l'équivalence induite par l'inclusion}
\[(\underline{\mathrm{Liss}}\,S)(W_S^{\mathrm{d}})^{-1}\xrightarrow{\ \sim\ }(\Cat/S)(W_S^{\mathrm{d}})^{-1}\quad\textit{),}\] % doubt resolved in DER_VI-c_corrections.csv; input note: underline of Liss -- double underline or single underline
\textit{un foncteur, encore noté} $\multhot_{X/S}$,
\[\multhot_{X/S} : \HOT_W(S)\longrightarrow\HOT_W(S).\]

Ce foncteur satisfait donc à la relation
\[\multhot_{X/S}(\hot_{S,W}(Y))\qquad\underbrace{\ \simeq\ }_{\text{isomorphisme canonique}}\qquad\hot_{S,W}(X\times_S Y)\]
% foot: Version 9 novembre 2001
\printedpage{95}
% head: A. Grothendieck : Dérivateurs, Ch. VI, Édition des Universités de Montpellier II et Paris VII
% marginal mark: handwritten sign in the left margin beside the following two lines
pour $Y$ \textit{lisse sur $S$,} et pour $Y$ quelconque dans $\Cat/S$,
\[\multhot_{X/S}(\hot_{S,W}(Y))\qquad\underbrace{\ \simeq\ }_{\text{isomorphisme canonique}}\qquad\hot_{S,W}(X\times_S \Phi_S(Y)).\]

Quand $X$ est lui-même lisse sur $S$, ce foncteur n'est autre que la multiplication par $\hot_{S,W}(X)$ :
\[\multhot_{X/S}=\mult_{\hot_{S,W}(X)}\qquad\text{si $X$ \textit{lisse sur $S$.}}\]

% marginal mark: handwritten sign in the left margin beside the following paragraph
Mais sa valeur pour $\hot_{S,W}(Y)$, quand $Y$ n'est \textit{pas} lisse sur $S$, ne doit pas être confondue avec $\hot_{S,W}(X\times_S Y)$.

Ainsi, si $X$ et $Y$ sont dans $\Cat/S$, considérons
\[X\times_S Y\qquad\text{et}\qquad\boxed{\hot_{S,W}(X\times_S Y).}\] % the frame is printed with rounded corners

\begin{enumerate}
\item[a)] Si ni $X$ ni $Y$ ne sont supposés lisses sur $S$, cet élément ‘n'a aucune interprétation naturelle en termes d'opérations sur les catégories $\HOT_W(S)$’ (\footnote[74]{Sauf dans le cas où l'un des deux soit dans $\underline{\mathrm{Fibhot}}\,S$, cf. annotation page suivante.}).

\msp{156}

\item[b)] Si $X$ et $Y$ sont lisses sur $S$, cet élément s'interprète comme
\[\hot_{S,W}(X)\times_S\hot_{S,W}(Y)\]
(\footnote[75]{Mais aussi si l'\textit{un} des deux éléments est dans $\underline{\mathrm{Fibhot}}\,S$ (et moyennant l'axiome L6 bis sur $W$).}).

\item[c)] Si l'un des deux $X$, $Y$ est lisse sur $S$, sans rien supposer sur l'autre, on a une interprétation en termes de $\mult_X$ (si $Y$ lisse), ou $\mult_Y$ (si $X$ lisse), notamment
\[\begin{aligned}
\hot_{S,W}(X\times_S Y)\ &\simeq\ \multhot_X(\hot_{S,W}(Y))\quad\text{si $Y$ lisse sur $S$}\\
&\simeq\ \multhot_Y(\hot_{S,W}(X))\quad\text{si $X$ lisse sur $S$.}
\end{aligned}\]
\end{enumerate}

\textbf{Corollaire 2.} \textit{Soit $f : S'\longrightarrow S$ une flèche dans} Cat\textit{, d'où $f_!$, $f^*$ entre} $\HOT_W(S')$ \textit{et} $\HOT_W(S)$\textit{. On a la relation}
\[f_!f^*(\xi)\simeq\multhot_{S'}(\xi),\]
\textit{donc si $S'$ est lisse sur $S$,}
\[f_!f^*(\xi)\simeq\hot_{S,W}(S')\times_S\xi.\]

Nous sommes maintenant en mesure de vérifier les axiomes Der 4$'$b, Der 5$'$ relatifs aux images directes (I, p.~67, 68), et de plus Der 3b (I, p.~111). L'axiome Der 4$'$ n'est autre que l'existence des adjoints à gauche $f_!$ des $f^*$. L'axiome Der 5$'$ permet le ‘calcul’ des fibres d'un objet $f_!(\xi)$. Il dit qu'une certaine
% foot: Version 9 novembre 2001
\printedpage{96}
% head: A. Grothendieck : Dérivateurs, Ch. VI, Édition des Universités de Montpellier II et Paris VII
\msp{157}
flèche du changement de base, relatif au diagramme commutatif dans Cat
\[\begin{tikzcd}[row sep=large]
S' \arrow[d,"f"'] & s\backslash S' \arrow[l,"q"'] \arrow[d,"f'"] \\
S & e\simeq\{s\} \arrow[l,"p"']
\end{tikzcd}\qquad(s\in\mathrm{Ob}\,S),\]
savoir
\[p^*(f_!(\xi'))\longrightarrow f'_!(q^*(\xi'))\]
est un isomorphisme (pour $\xi'\in\mathbf{D}(S')$). Soit donc
\[\xi'\in\mathrm{Ob}\,\HOT_W(S'),\qquad\xi'=\hot_{S',W}(X')\]
avec $X'$ dans $\Cat/S'$. Comme
\[f_!(\xi')=\hot_{S,W}(X'),\]
et que $X'$ n'a aucune raison d'être lisse sur $S$ (si on ne suppose pas $S'$ lui-même lisse sur $S$, auquel cas il pourrait être utile de prendre $X'$ lisse sur $S'$, donc aussi sur $S$), pour calculer $p^*(f_!(\xi'))$ il n'est pas question de prendre la fibre $X'_s$, mais il faut prendre $s\backslash X'=\Phi_S(X')_s$ :
\[p^*(f_!(\xi'))\simeq\hot_W(s\backslash X').\]
\[\begin{tikzcd}[row sep=large]
X' \arrow[d] & s\backslash X' \arrow[l] \arrow[d] \\
S' \arrow[d] & s\backslash S' \arrow[l,"q"'] \arrow[d] \\
S & s\backslash S \arrow[l]
\end{tikzcd}\]

Considérons le diagramme de changement de base par $s\backslash S\longrightarrow S$ ci-contre \texttt{[i.e. ci-dessus]}, à carrés cartésiens. On voit d'abord que $q^*(\xi')$ n'est autre que la classe dans $\HOT_W(s\backslash S')$ de l'image inverse de $X'$ par le changement de base $s\backslash S'\longrightarrow S'$. Pour l'assurer, il suffirait de prendre $X'$ lisse sur $S'$ (ce qui
\msp{158}
est licite) - alors la description de l'image inverse est OK pour tout changement de base. Mais on a vu que dans le cas d'un co-isomorphisme local comme $s\backslash S'\longrightarrow S'$, il n'y a pas même besoin de cet artifice - la formule est OK sans restriction de lissité. Ainsi
\[q^*(\xi')=\hot_{s\backslash S',W}(s\backslash X'),\]
% foot: Version 9 novembre 2001
\printedpage{97}
% head: A. Grothendieck : Dérivateurs, Ch. VI, Édition des Universités de Montpellier II et Paris VII
d'où
\[f'_!(q^*(\xi'))\simeq\hot_{S,W}(s\backslash X')=\hot_W(s\backslash X').\]
Donc on a bien trouvé
\[p^*f_!\simeq f'_!q^*.\]

On peut maintenant prouver aussi Der 3b. On a vu qu'une fois qu'on dispose des adjoints à gauche $f_!$ des $f^*$, cet axiome s'énonce sous \textit{forme homologique} ainsi : Si
\[f : S'\longrightarrow S\]
est dans $W$, il faut prouver que pour tout
\[\xi\in\HOT_W(e)\overset{\text{déf}}{=}\mathrm{Hot}_W,\]
la flèche correspondante
\[(*)\qquad p'_!p'^*(\xi)\longrightarrow p_!p^*(\xi)\]
est \texttt{[un]} isomorphisme dans $\mathrm{Hot}_W$, où $p : S\longrightarrow e$, $p' : S'\longrightarrow e$ sont les flèches structurales. Or comme tout objet de Cat est lisse sur $e$, le corollaire 2 précédent permet d'identifier cette flèche $(*)$ à
\[\eta'\times_S\xi\longrightarrow\eta\times_S\xi\]
\msp{159}
déduite de
\[\hot_W(f) :\quad\underbrace{\ \eta'\ }_{\overset{\text{déf}}{=}\,\hot_W(S')}\quad\longrightarrow\quad\underbrace{\ \eta\ }_{\overset{\text{déf}}{=}\,\hot_W(S)}\quad.\]
Comme celle-ci est \texttt{[un]} isomorphisme par hypothèse, on gagne.

Donc on a prouvé

\textbf{Théorème 3.} \textit{Le prédérivateur} HOT \textit{satisfait aux axiomes} Der 1, Der 2 (\footnote[76]{\textit{pas prouvé}, le ‘résultat’ auquel je pensais est canulé, cf. complément proposition 5, page 164.}) \textit{(pour mémoire), et} Der 3b (\footnote[77]{Et même la variante la plus forte Der 3b$'_2$ de Der 3, cf. démonstration page suivante. \textbf{NB} Pour Der 3b$'_2$, on a utilisé l'axiome L6 sur $W$ - et c'est essentiel. On voit que la validité de Der 3b$'_2$ implique L6, et même la forme plus forte L6 bis, et même si on suppose seulement Der 3b$'_1$, i.e. $f_0^*$ pleinement fidèle.}), Der 4$'$, Der 5$'$\textit{.} % [sic] canulé (footnote 76) -> annulé (?), not a French word

On a envie de prouver aussi les variantes renforcées de Der 3, la forme la plus forte étant celle-ci (Der 3b$'_2$) : si $f\in W$, alors $f^*$ induit une \textit{équivalence} de catégories
\[f_0^* : \HOT^{\mathrm{lc}}_W(S)\longrightarrow\HOT^{\mathrm{lc}}_W(S').\]

La pleine fidélité du foncteur $f_0^*$ doit résulter de la variante naturelle Der 3b$''_{\mathrm{coh}}$ de Der 3b, une fois qu'on aura développé le formalisme des Hom internes dans les $\HOT_W(S)$. Mais pour donner un sens à Der 3b$''_{\mathrm{coh}}$, on a besoin des foncteurs \textit{cohomologie,} i.e. de $p_*$ et $p'_*$ (au lieu de $p_!$, $p'_!$) pour HOT, donc c'est encore prématuré pour l'instant. La construction des $f_*$ et celle des Hom internes se fera d'ailleurs sûrement $\pm$ conjointement.
% foot: Version 9 novembre 2001

\printedpage{98}
% head: A. Grothendieck : Dérivateurs, Ch. VI, Édition des Universités de Montpellier II et Paris VII
\msp{160}
Mais je m'aperçois à l'instant qu'on trouve quasiment ‘à l'œil’ le fait que $f_0^*$ soit bien une \textit{équivalence} de catégories, et que $f_!^0$ soit l'équivalence quasi-inverse, en utilisant l'adjonction des foncteurs $f_!^0$, $f_0^*$, et se rappelant que pour prouver que ce sont des équivalences quasi-inverses l'une de l'autre, il revient au même de montrer que les morphismes d'adjonction
\[\begin{aligned}
\mathrm{id}&\quad\xrightarrow{\ \alpha_0\ }\quad f_0^*f_!^0\\
\mathrm{id}&\quad\xleftarrow{\ \beta_0\ }\quad f_!^0f_0^*
\end{aligned}\]
sont \texttt{[des]} isomorphismes. Je vais essayer de décrire ces flèches d'adjonction dans le cas général (sans supposer $f\in W$), puis de montrer qu'elles sont dans $W$ si $f$ l'est.

\begin{enumerate}
\item[\texttt{[1)}] \texttt{Soit]} $\xi=\hot^{\mathrm{lc}}_{S,W}(X)$ dans $\HOT_W(S)$. On peut supposer $X\in\mathrm{Ob}\,\underline{\mathrm{Parf}}\,S$. Comme \texttt{[‘Comme’ barré]} $X$ lisse sur $S$, on a
\[f_0^*(\xi)=\hot^{\mathrm{lc}}_{S',W}(X\times_S S'),\qquad f_!^0f_0^*(\xi)=\hot^{\mathrm{lc}}_{S,W}(X\times_S S').\]
Donc
\[\beta_{0,\xi} : f_!^0f_0^*(\xi)\longrightarrow\xi\]
est une flèche
\[X\times_S S'\longrightarrow X\qquad\text{dans }(\Cat/S)(W_S)^{-1}\]
et je présume que c'est $\mathrm{pr}_1$ (\footnote[78]{\textbf{NB} On voit que l'axiome L6 est $\pm$ \textit{nécessaire} pour que le résultat sur $f_0^*$, $f_!^0$ soit valable.}). Est-ce dans $W$ si $f : S'\longrightarrow S$ l'est? Oui (car $X$ est parfait

\msp{161}

sur $S$), sous réserve que $W$ satisfait l'axiome L6 (cf. page 43).

\item[2)] Soit
\[\xi'=\hot^{\mathrm{lc}}_{S',W}(X')\qquad\text{dans }\HOT_W(S'),\]
on a
\[f_!^0(\xi')=\hot^{\mathrm{lc}}_{S,W}(X'),\]
mais pour expliciter $f_0^*$ de cet élément, comme $X'$ n'a aucune raison d'être un Hot-fibré sur $S$ (auquel cas il suffirait de prendre $f^*(X')\overset{\text{déf}}{=}X\times_S S'$), il faut d'abord prendre $C_S(X')$, puis appliquer $f^*$ :
\[f_0^*f_!^0(X')=\hot^{\mathrm{lc}}_{S',W}(C_S(X')\times_S S').\]

Il faut donc définir
\[X'\xrightarrow{\beta_{0,X'}}C_S(X')\times_S S'\qquad\text{dans }(\Cat/S')W_{S'}^{-1}.\]

Or on trouve une flèche canonique dans $\Cat/S'$, déduite du $S$-morphisme $i_{X'}$ : $X'\longrightarrow C_S(X')$, i.e. du carré commutatif en traits pleins
\end{enumerate}
% foot: Version 9 novembre 2001


\endgroup
% END UNIT Der-ChVI-345/Der-ChVI-345_p096-100.tex
% BEGIN UNIT Der-ChVI-345/Der-ChVI-345_p101-105.tex
\begingroup\makeatletter%
\newcommand{\printedpage}[1]{\marginpar{\footnotesize p.~#1}}%
\newcommand{\msp}[1]{\marginpar{\footnotesize [#1]}}%
\newcommand{\HOT}{\mathrm{HOT}}%
\newcommand{\hot}{\mathrm{hot}}%
\newcommand{\Hot}{\mathrm{Hot}}%
\newcommand{\Cat}{\mathrm{Cat}}%
\newcommand{\Hom}{\mathrm{Hom}}%
\newcommand{\pr}{\mathrm{pr}}%
\newcommand{\D}{\mathbf{D}}%
\newcommand{\bs}{\backslash}%
\makeatother\let\printedpage\derpp\let\msp\dermsp%
\let\mspin\DERmspin%


\noindent\printedpage{99}%
% head: A. Grothendieck : Dérivateurs, Ch. VI, Édition des Universités de Montpellier II et Paris VII
\[
\begin{tikzcd}
X' \arrow[rr, "i_{X'}"] \arrow[ddd, "p_{X'}"'] \arrow[dr, dotted, "\beta_{0,X'}"'] & & C_S(X') \arrow[ddd] \\
& C_S(X') \times_S S' \arrow[ur, dotted, "\pr_1"'] \arrow[ddl, dotted] & \\
& \text{\footnotesize commutatif} & \\
S' \arrow[rr, "f"] & & S,
\end{tikzcd}
\]

\begin{quote}
comme l'unique flèche pointillée $\beta_{0,X'}$ qui soit un $S'$-morphisme (triangle gauche commutatif) et qui rende commutatif le triangle supérieur.
\end{quote}

Si $f \in W$, alors $\pr_1$, déduit de $f$ par changement de base $C_S(X') \longrightarrow S$, qui est parfait, est dans $W$ par L6, on sait qu'il en est de même de $i_{X'}$, donc aussi de $\beta_{0,X'}$, on gagne.

\msp{page 162}Pour terminer les propriétés générales des foncteurs $f_!$ pour $\HOT_W$, je voudrais donner les énoncés raisonnables de changement de base. Ce serait un

\noindent\textbf{Théorème 4.} \textit{Soit}
\[
\begin{tikzcd}
X \arrow[d, "f"'] & X' \arrow[l, "q"'] \arrow[d, "f'"] \\
S & S' \arrow[l, "p"']
\end{tikzcd}
\]
\textit{un diagramme cartésien dans} $\Cat$\textit{, considérons la flèche de changement de base (dans} $\underline{\Hom}(\HOT_W(X), \HOT_W(S'))$\textit{) :}
\[
p^* f_! \longleftarrow f'_! q^*.
\]
\textit{C'est un isomorphisme dans chacun des deux cas suivants.}

\begin{enumerate}
\item[\textit{a)}] \textit{$f$ est lisse.}
\item[\textit{b)}] \textit{$p$ est propre.}
\end{enumerate}

\textsc{Démonstration.} \textit{Cas a).} On teste avec $\xi$ dans $\HOT_W(X)$, on peut supposer $\xi$ de la forme
\[
\xi = \hot_{X,W}(Z) \qquad \text{avec } Z \in \mathrm{Ob}\, \underline{\mathrm{Liss}}\, X.
\]
Notons que $Z$ est alors aussi lisse sur $S$, puisque $f$ est lisse.
\[
\begin{aligned}
f_!(\xi) &= \hot_{S,W}(Z), &\qquad p^* f_!(\xi) &= \hot_{S',W}(\quad \underbrace{Z \times_S S'}_{\text{$Z'$, cf. diagramme}} \quad) \\
q^*(\xi) &= \hot_{X',W}(Z'), & f'_! q^*(\xi) &= \hot_{S',W}(Z'),
\end{aligned}
\]
OK!

\noindent\printedpage{100}%
% head: A. Grothendieck : Dérivateurs, Ch. VI, Édition des Universités de Montpellier II et Paris VII
\[
\begin{tikzcd}
Z \arrow[d] & Z' \arrow[l] \arrow[d] \\
X \arrow[d, "f"'] & X' \arrow[l, "q"] \arrow[d, "f'"] \\
S & S' \arrow[l, "p"]
\end{tikzcd}
\]
\textit{Cas b).} Il faut admettre que \msp{page 163}le foncteur changement de base par $p$ :
\[
Z \longmapsto Z \times_S S' : \Cat/S \longrightarrow \Cat/S'
\]
est compatible avec les localiseurs $W^{\mathrm{d}}_{S}$, $W^{\mathrm{d}}_{S'}$ (cf. plus bas) - c'est ça, formellemet, la propriété pertinente pour $p$ (\footnote[79]{Ça doit être \textit{équivalent} à $p$ propre. Comparer avec le corollaire 1, page 166.}). Il en résultera que % [sic] formellemet -> formellement, letter omitted
\[
p^* : \HOT_W(S) \longrightarrow \HOT_W(S')
\]
se calcule, pour un élément dans $\HOT_W(S)$ donné par un $Z$ \textit{quelconque} dans $\Cat/S$ (pas nécessairement lisse sur $S$), par
\[
p^*(\hot_{S,W}(Z)) = \hot_{S',W}(Z \times_S S').
\]
Partons alors, comme dans a), d'un élément $\xi$ de $\HOT_W(X)$ donné par un $Z$ de $\Cat/X$ (inutile cette fois de supposer $Z$ lisse), le calcul immédiat précédent établit encore la formule de commutation.

Mais il ne semble pas que la propriété de compatibilité de $X \longmapsto X \times_S S'$, avec l'équivalence colocale sur $S$ et sur $S'$, résulte de façon $\pm$ évidente de la condition que $p : S' \longrightarrow S$ soit propre (\footnote[80]{Mais si, cf. démonstration page 168. Les développements des quatre pages qui suivent sont inutiles par cette démonstration.}). J'ai envie de le déduire de fourbis généraux sur le formalisme des dérivateurs, à savoir du théorème de changement de base des $f_!^{\D}$ pour un morphisme $\D$-propre.

\msp{page 164}Sous forme duale d'un théorème de changement de base pour $f_*^{\D}$, par changement de base $\D$-lisse, c'est fait dans I, p. 56 à 58, et on doit pouvoir l'appliquer ici. Notons que dans la démonstration du théorème loc. cit., on n'utilise que L4 \texttt{[plutôt} Der 4\texttt{]}, i.e. l'existence des $f_*^{\D}$, à l'exclusion des $f_!^{\D}$ (L4$'$ \texttt{[plutôt} Der 4$'$\texttt{]}), plus L5 \texttt{[plutôt} Der 5\texttt{]} bien sûr, donc la forme duale ne fait appel à l'existence des $f_*^{\D}$, seulement des $f_!^{\D}$. On a donc une bonne chance de pouvoir l'appliquer à $\HOT_W$ (où je n'ai pas encore établi l'existence des $f_*^{\mathrm{HOT}}$). % [sic] ne fait appel -> ne fait pas appel, "pas" apparently omitted (sense and "seulement des")

\noindent\printedpage{101}%
% head: A. Grothendieck : Dérivateurs, Ch. VI, Édition des Universités de Montpellier II et Paris VII
Notons d'abord

\noindent\textbf{Proposition 5.} \textit{Soit} $f : S' \longrightarrow S$ \textit{une flèche dans} $\Cat$. \textit{Alors $f$ est une} $\HOT_W$\textit{-équivalence} si et seulement si $\hot_W(f)$ \textit{est} \texttt{[un]} \textit{isomorphisme. Si $W$ est fortement saturé, cela signifie donc que $f \in W$ (et en tous cas $f \in W \Longrightarrow f$} \texttt{[est]} \textit{une} $\HOT_W$\textit{-équivalence)} (\footnote[81]{devrait remonter vers le théorème 3, p. 159.}).

C'est contenu dans la démonstration de Der 3b tantôt. Je rappelle que $f$ est une HOT-équivalence, par définition, si et seulement si pour $a, b \in \Hot_W (= \HOT_W(e))$,
\[
\Hom_{\HOT_W(S)}(a_S, b_S) \longrightarrow \Hom_{\HOT_W(S')}(a_{S'}, b_{S'}) \tag*{$(*)$}
\]
est \texttt{[un]} isomorphisme, où $a_S$ etc. désigne $p_S^*(a)$ etc. ($p_S : S \longrightarrow e$). Les formules d'adjonction de $p_{S!}$, $p_S^*$ et de $p_{S'!}$, $p_{S'}^*$ permettent d'interpréter la flèche $(*)$ comme
\[
\Hom_{\HOT_W(S)}(p_{S!}(a_S), b) \longrightarrow \Hom_{\HOT_W(S')}(p_{S'!}(a_{S'}), b),
\]
\msp{page 165}déduit par ‘transposition’ $\Hom(-, b)$ de la flèche
\[
p_{S'!}(a_{S'}) \longrightarrow p_{S!}(a_S)
\]
induite par $f$ - donc il faut prouver que celle-ci est un isomorphisme (c'est la forme que nous avions utilisé tantôt). Or cette flèche s'identifie à
\[
\eta' \times a \longrightarrow \eta \times a
\]
induite par
\[
\hot_W(f) : \underbrace{\eta'}_{= \hot_W(S')} \quad \longrightarrow \quad \underbrace{\eta}_{= \hot_W(S)} \quad ,
\]
et dire que c'est \texttt{[un]} isomorphisme quel que soit $a$ dans $\Hot_W$ revient à dire que $\eta' \longrightarrow \eta$ l'est (puisque la multiplication dans $\Hot_W$ a une élément unité \dots). % [sic] une élément -> un élément, gender agreement

Maintenant je me rappelle que pour avoir l'axiome Der 2 des dérivateurs, je devrais également supposer $W$ \textit{fortement} saturé (i.e. égal au localiseur défini par le foncteur canonique $\Cat \longrightarrow (\Cat)W^{-1} \overset{\text{déf}}{=} \Hot_W$). (C'est ce que je supposais d'emblée, en partant d'un $(\mathcal{M}, W)$ général pour définir un prédérivateur, pour avoir l'axiome Der 2.)

En appliquant le dual du théorème de I, p. 56, on trouve donc le théorème de changement de base pour $f_!$, pour changement de base $f : S' \longrightarrow S$ \textit{$W$-propre.}

% the following paragraph is marked in the left margin by a double vertical rule
\msp{page 166}\marginpar{$\Vert$}Mais la démonstration utilise sûrement Der 2 (critère d'isomorphisme par fibres, dans $\HOT_W(S)$), que je n'ai pas prouvé encore. Elle est incomplète pour le moment. Mais je peux noter dès maintenant le

\noindent\printedpage{102}%
% head: A. Grothendieck : Dérivateurs, Ch. VI, Édition des Universités de Montpellier II et Paris VII
\textbf{Corollaire 1} (du théorème 4, p. 162). \textit{Si la flèche} $S' \xrightarrow{p} S$ \textit{dans} $\Cat$ \textit{donne lieu (pour tout} $f : X \longrightarrow S$ \textit{dans} $\Cat/S$\textit{) à la formule de changement de base pour}
\[
f_! : \HOT_W(X) \longrightarrow \HOT_W(S),
\]
\textit{alors le foncteur} $p^* : X \longmapsto X \times_S S'$ \textit{de} $\Cat/S$ \textit{dans} $\Cat/S'$ \textit{est compatible avec les localiseurs} $W^{\mathrm{d}}_{S}$\textit{,} $W^{\mathrm{d}}_{S'}$\textit{.} (L'implication inverse a été vue déjà, cf. pages 161-163.)

Cela résultera du

\noindent\textbf{Lemme.} \textit{Soit} $u : X' \longrightarrow X$ \textit{une flèche dans} $\Cat/S$\textit{. Pour que} $u \in W^{\mathrm{d}}_{S}$\textit{, il faut et il suffit (si $W$} \texttt{[est]} \textit{fortement saturé) que pour tout} $\xi \in \Hot_W$\textit{, la flèche correspondante}
\[
u_* : p'_!(\xi_{X'}) \longrightarrow p_!(\xi_X)
\]
\textit{soit} \texttt{[un]} \textit{isomorphisme dans} $\HOT_W(S)$\textit{.}

C'est suffisant, sous réserve de saturation forte de $W$, car si $u_*$ est \texttt{[un]} isomorphisme, passant aux fibres, on trouve par Der 5$'$
\[
H_\bullet^{\HOT_W}(s\bs X', \xi) \longrightarrow H_\bullet^{\HOT_W}(s\bs X, \xi),
\]
si c'est \texttt{[un]} isomorphisme, alors $s\bs X' \longrightarrow s\bs X$ est par définition une HOT-équivalence, ce qui implique que c'est dans $W$ (proposition 5).

\msp{page 167}C'est nécessaire, car si $\xi = \hot_W(T)$, donc $\xi_S = \hot_{S,W}(T_S)$ (où $T_S = T \times S/S$), et on a donc
\[
p'_!(\xi_{X'}) = \hot_{S,W}(X' \times T), \qquad p_!(\xi_X) = \hot_{S,W}(X \times T),
\]
et il faut prouver que si $u : X' \longrightarrow X$ est dans $W^{\mathrm{d}}_{S}$, alors $\underbrace{X' \times T}_{X' \times_S T_S} \longrightarrow \underbrace{X \times T}_{X \times_S T_S}$ aussi, ce qui signifie simplement que
\[
s\bs(X' \times T) \longrightarrow s\bs(X \times T) \tag*{$(*)$}
\]
est dans $W$ pour tout $s \in S$. Mais le changement de base $s\bs S \longrightarrow S$ commute aux produits fibrés, donc transforme $X \times T = X \times_S T_S$ en $(s\bs X) \times_{s\bs S} T_{(s\bs S)} = s\bs X \times T$, et de même pour $X'$. Donc la flèche s'identifie à
\[
(s\bs X') \times T \longrightarrow (s\bs X) \times T,
\]
déduite de
\[
s\bs X' \longrightarrow s\bs X
\]
par le foncteur $? \times T$. Comme ce foncteur transforme $W$ en $W$, on gagne.

\noindent\printedpage{103}%
% head: A. Grothendieck : Dérivateurs, Ch. VI, Édition des Universités de Montpellier II et Paris VII
\msp{page 168}Finalement, je m'aperçois que j'ai eu un blocage stupide, pour la question que le changement de base propre soit compatible à la $W$-équivalence colocale. Soit
\[
\begin{tikzcd}[column sep=small]
X \arrow[rr, "f"] \arrow[dr] & & Y \arrow[dl] & \\
& S & & S', \arrow[ll, "g"']
\end{tikzcd}
\]
avec $f \in W^{\mathrm{d}}_{S}$ et $g$ propre, il faut voir que $f' : X' \longrightarrow Y'$ déduit de $f$ par changement de base $S' \longrightarrow S$ est encore dans $W^{\mathrm{d}}_{S'}$. Cela signifie que tout changement de base de la forme $s'\bs S' \longrightarrow S'$ le transforme en $\in W$. Mais soit $s = g(s')$, on a le carré commutatif
\[
\begin{tikzcd}
S & S' \arrow[l] \\
\underbrace{s\bs S}_{= Z} \arrow[u] & \underbrace{s'\bs S'}_{= Z'} \arrow[l] \arrow[u] \ ,
\end{tikzcd}
\]
et il suffit de voir que pour toute flèche $\varphi : \mathcal{X} \longrightarrow \mathcal{Y}$ dans $\Cat/Z$ qui est dans $W$, son image inverse par changement de base $Z' \longrightarrow Z$ est itou. (En appliquant ceci au cas $\mathcal{X} = X \times_S (s\bs S)$, $\mathcal{Y} = Y \times_S (s\bs S)$, $\varphi = s\bs f$.) Or ceci résulte du fait que, par \msp{page 169}hypothèse de propreté, $Z' \longrightarrow Z$ est propre et à fibres Hot-asphériques, donc c'est une Hot-équivalence universelle. Il suffirait même que $g$ soit $W$-propre, pour impliquer que $Z' \longrightarrow Z$ est une $W$-équivalence universelle, i.e. qui reste telle après tout changement de base. Or on a le

\noindent\textbf{Lemme.} Si $g : Z' \longrightarrow Z$ est une $W$-équivalence universelle, alors le foncteur
\[
g^* : X \longmapsto X \times_Z Z' : \Cat/Z \longrightarrow \Cat/Z'
\]
est compatible avec les localiseurs $W_Z$, $W_{Z'}$, i.e. si $f : X \longrightarrow Y$ dans $\Cat/Z$ est dans $W$, de même $f' : X' \longrightarrow Y'$ dans $\Cat/Z'$.
\[
\begin{tikzcd}
X \arrow[dr, "f"'] \arrow[ddr] & & X' \arrow[ll, "g_X"'] \arrow[dr, "f'"] \arrow[ddr] & \\
& Y \arrow[d] & & Y' \arrow[ll, "g_Y"', crossing over] \arrow[d] \\
& Z & & Z' \arrow[ll, "g"']
\end{tikzcd}
\]
Cela est évident sur le carré supérieur du diagramme ci-dessus, où $g_X$, $g_Y$ sont dans $W$ par l'hypthèse sur $g$, donc $f$ est dans $W$ si et seulement si $f'$ l'est, cqfd. % [sic] hypthèse -> hypothèse, letter omitted

\msp{page 170 manque}\texttt{[page 170 manque]}


\endgroup
% END UNIT Der-ChVI-345/Der-ChVI-345_p101-105.tex
% BEGIN UNIT Der-ChVI-345/Der-ChVI-345_p106-110.tex
\begingroup\makeatletter%
\newcommand{\printedpage}[1]{\marginpar{\footnotesize p.~#1}}%
\newcommand{\msp}[1]{\marginpar{\footnotesize [#1]}}%
\newcommand{\ul}[1]{\underline{\mathrm{#1}}}%
\newcommand{\circled}[1]{\tikz[baseline=(c.base)]\node[draw,rounded corners,inner sep=2pt](c){#1};}%
\makeatother\let\printedpage\derpp\let\msp\dermsp%
\let\mspin\DERmspin%

\setcounter{footnote}{81}

\printedpage{104}
% head: A. Grothendieck : Dérivateurs, Ch. VI, Édition des Universités de Montpellier II et Paris VII
% foot: Version 9 novembre 2001

\msp{page 171}

\subsection*{Nouvelle version de la théorie de $\mathrm{HOT}_W$.}

Je me suis aperçu que l'hypothèse $W \supseteq W_\infty$ (L3 quater, p. 1), que j'avais fini par mettre sur $W$, est artificielle, et que l'hypothèse L5 (p. 43), voire L6, L6 bis (p. 43-44), est inutile pour une grande partie de la théorie de $\mathrm{HOT}_W$. J'ai été ainsi amené à revoir la théorie (i.e. ce que j'en ai fait jusqu'à présent) sous les hypothèses suivantes sur $W$ (cf. XII) :

\begin{itemize}[leftmargin=4em]
\item[W(1)] (\textit{Saturation} de $W$) : comme avant (L1 page 1)
\item[W(2)] (localisation) : c'est l'ancien L4 bis (\footnote{i.e. le L5 de la page 2. L'axiome est repris (sous une forme voisine et équivalente) au bas de la page 43.}).
\item[W(3)] (objet final) : Si $X$ a un objet final, $X$ est $W$-asphérique (ancien L3).
\end{itemize}

Ces trois axiomes sont suffisants pout tout ce que j'avais fait dans Pursuing Stacks, et dans une grande partie de ce qui a été fait ici sur HOT. Parfois cependant il faut l'ancien L 5, que je réénonce ici (comparer XII, théorème \texttt{[numéro du théorème manque]}) : % [sic] pout -> pour, misprint

\begin{itemize}[leftmargin=4em]
\item[W(4b)] Soient $X$, $Y$ $W$-fibrés sur $S$, $f : X \longrightarrow Y$ un $S$-morphisme. Si $f \in W$, alors $\forall s \in \mathrm{Ob}\, S$, $f_s : X_s \longrightarrow Y_s$ est dans $W$.
\end{itemize}

Pour la notion de $W$-fibré (\textit{faible}), j'ai été long à trouver enfin \textit{la} bonne notion, bonne pour un $W$ ne satisfaisant que $\mathrm{W}(1,2,3)$ (et même seulement $\mathrm{W}(1, 2 \text{ bis}, 3)$, où W(2 bis) est l'ancien L4). Je renvoie pour cela à (XII, théorème 3, p. 51), où la notion est bien cernée à coups de conditions équivalentes. En fait, pour les applications principales à $\mathrm{HOT}_W$, vues jusqu'à présent, il suffit (\footnote{C'est finalement une erreur, semble-t-il à présent. Je crois que l'axiome W(4b$'$) ne sert finalement à strictement rien (si ce n'est en fait à entretenir des idées fausses.)}) d'une variante affaiblie de W(4b) :

\begin{itemize}[leftmargin=4em]
\item[W(4b$'$)] Comme W(4b), mais en exigeant de plus que $X$, $Y$ soient dans $\ul{Parf}_\omega\, S$.
\end{itemize}

\msp{page 172}

Pour la notion de morphismes $W$-propres, $W$-lisses, $W$-parfaits dans Cat, je renvoie à XII, partie III (p. 22-52), et les notations $\ul{Liss}_W S$, $\ul{Prop}_W S$, $\ul{Parf}_W S$ (sous-catégories de $\mathrm{Cat}/S$), et $\mathrm{Parf}_W \subseteq \mathrm{Fl}(\mathrm{Cat})$. Dans l'énoncé de W(4b$'$) on a posé
\[
\ul{Parf}_\omega S = \ul{Parf}_{W_\omega} S, \quad
\begin{array}[t]{l}
\text{où } W_\omega \subseteq \mathrm{Fl}(\mathrm{Cat}) \text{ est la plus petite partie de } \mathrm{Fl}(\mathrm{Cat})\\
\text{satisfaisant } \mathrm{W}(1, 2\mathrm{bis}, 3).
\end{array}
\]

Je note de même
\[
W'_\omega \subseteq \mathrm{Fl}(\mathrm{Cat}), \qquad \text{plus petite partie de } \mathrm{Fl}(\mathrm{Cat}) \text{ satifaisant } \mathrm{W}(1, 2, 3, 4\mathrm{b}). % [sic] satifaisant -> satisfaisant, misprint
\]

\printedpage{105}
% head: A. Grothendieck : Dérivateurs, Ch. VI, Édition des Universités de Montpellier II et Paris VII
% foot: Version 9 novembre 2001

L'hypothèse $W \supseteq W_\infty$ que nous avions fini par imposer sur $W$, et qui est apparu artificielle au cours de la réflexion du XII, a avantage à être remplacée par % [sic] apparu -> apparue, agreement with « l'hypothèse »
\[
W \supseteq W'_\omega \qquad (\text{ou par } W \supseteq \overline{W'_\omega}, \text{ où } \overline{W'_\omega} \text{ est le saturé fort de } W'_\omega)
\]
(\footnote{$\overline{W'_\omega}$ est le plus petit sous-ensemble \textit{fortement saturé} de $\mathrm{Fl}(\mathrm{Cat})$, satisfaisant W(4b), et il satisfait même la condition W$'$(4a), variante la plus forte des axiomes du type W(4).}). Mais nous ne l'imposons pas comme condition impérative, au même titre que $\mathrm{W}(1,2,3)$, pas plus que W(4b). Le rôle que j'ai fait jouer à $W_\infty$ dans les semaines passées, semble revenir plus naturellement à $W'_\omega$ ou $\overline{W'_\omega}$. Il en est ainsi notamment aussi, il me semble, dans l'énoncé des conditions du type Der 3 sur un dérivateur (I, p. 111-112), qu'il faudrait revoir dans cette lumière. Mais tout dépend, si (comme j'en ai l'impression) on peut faire vraiment de l'algèbre homotopique substantielle, avec notamment les suites exactes longues de fibration et de suspension, pour $\mathrm{HOT}_W$, avec des $W$ qui sont \textit{strictement} contenus dans $W_\infty$. Avant de même être assuré, il est trop tôt encore pour me prononcer sur la `bonne' formulation de Der 3,

\msp{page 173}

qui est l'axiome mettant en relation le dérivateur $\mathbf{D}$ avec $W_\infty$, ou avec quelque autre localiseur fondamental.

Je vais faire à présent un résumé de la théorie de $\mathrm{HOT}_W$, revue dans l'optique que je viens de dire : on suppose pour $W$ \textit{les conditions} $\mathrm{W}(1,2,3)$ \textit{sans plus,} sauf mention explicite.

\circled{1} \textbf{Localiseurs principaux} sur $\mathrm{Cat}/S$ et sur les principales catégories de $S$-catégories (cf. diagramme p. 76).

Ils seront tous induits par les trois localiseurs familiers
\[
\begin{tikzcd}[column sep=small]
W_S^{\mathrm{d}} \arrow[dr, hook] & & W_S^{\mathrm{g}} \arrow[dl, hook'] \\
& W_S &
\end{tikzcd}
\]
sur $\mathrm{Cat}/S$. Il y a lieu aussi de regarder $W_S^{\mathrm{f}}$, mais jamais sur $\mathrm{Cat}/S$ tout entier, mais sur d'autres catégories du diagramme, et toujours dans des cas où $W_S^{\mathrm{f}}$ coïncide soit avec $W_S^{\mathrm{d}}$, soit avec $W_S^{\mathrm{g}}$. Mais le point de vue $W_S^{\mathrm{f}}$ est souvent plus directement commode, et a l'avantage de plus d'être symétrique.

\circled{a} Dans $\ul{Liss}_W S$, on a

\printedpage{106}
% head: A. Grothendieck : Dérivateurs, Ch. VI, Édition des Universités de Montpellier II et Paris VII
% foot: Version 9 novembre 2001

\[
\begin{tikzcd}[row sep=large, column sep=0pt]
W_S^{\mathrm{f}} \arrow[d, Rightarrow] & {} = W_S^{\mathrm{d}} = W_S^{\mathrm{univ}} \\
W_S^{\mathrm{g}} \arrow[d, Rightarrow] & \\
W_S &
\end{tikzcd}
\]

Dualement,

\circled{b} dans $\ul{Prop}_W S$, on a
\[
\begin{tikzcd}[row sep=large, column sep=0pt]
W_S^{\mathrm{f}} \arrow[d, Rightarrow] & {} = W_S^{\mathrm{g}} = W_S^{\mathrm{univ}} \\
W_S^{\mathrm{d}} \arrow[d, Rightarrow] & \\
W_S, &
\end{tikzcd}
\]

\msp{page 174}

enfin,

\circled{c} dans $\ul{Fib}_W S$, on a
\[
\begin{tikzcd}[row sep=large]
W_S^{\mathrm{f}} \arrow[r, Leftrightarrow] \arrow[d, Rightarrow] & W_S^{\mathrm{d}} \arrow[r, Leftrightarrow] & W_S^{\mathrm{g}} \arrow[r, Leftrightarrow] & W_S^{\mathrm{univ}} \\
W_S. & & &
\end{tikzcd}
\]

Cette implication $W_S^{\mathrm{f}} \Longrightarrow W_S$ est stricte si on ne fait pas sur $W$ l'axiome W(4b), dont l'utilité ici est exactement d'assurer
\[
\left|\;
\begin{array}{l}
W_S^{\mathrm{f}} \Longleftrightarrow W_S \ \text{ dans } \ul{Fib}_W S \ \textit{si}\ \mathrm{W}(4\mathrm{b})\ \texttt{[est]}\ \text{satisfait}\\[1.5ex]
\text{et alors les cinq localiseurs envisagés ci-dessus}\\[1.5ex]
\text{dans } \ul{Fib}_W S \text{ sont tous égaux.}
\end{array}
\right.
\]

Les relations entre [?] localiseurs dans $\ul{Liss}_W S$, $\ul{Prop}_W S$, $\ul{Fib}_W S$ se transportent automatiquement aux sous-catégories de celles-ci, plus généralement à toutes celles qui s'envoient
\printedpage{107}
% head: A. Grothendieck : Dérivateurs, Ch. VI, Édition des Universités de Montpellier II et Paris VII
% foot: Version 9 novembre 2001
par un foncteur d'inclusion dans une de celles-ci. Dans le diagramme p. 76, toutes les catégories, à la seule exception de $\mathrm{Cat}/S$ tout entier, sont dans ce cas.

Pour les résultats principaux développés jusqu'à présent sur $\mathrm{HOT}_W$, quand l'axiome W(4b) est utilisé, la forme faible W(4b$'$) suffit. On a
\[
\left|\;
\begin{array}{l}
W_S^{\mathrm{f}} \Longleftrightarrow W_S \ \text{ dans } \ul{Parf}_\omega S \ \textit{si}\ \mathrm{W}(4\mathrm{b}')\ \texttt{[est]}\ \text{satisfait}\\[1.5ex]
\text{et alors les cinq localiseurs } W_S^{\mathrm{f}},\, W_S^{\mathrm{d}},\, W_S^{\mathrm{g}},\, W_S^{\mathrm{univ}},\, W_S\\[1.5ex]
\text{sont égaux dans } \ul{Parf}_\omega S.
\end{array}
\right.
\]

\circled{2} \textbf{Principales catégories de catégories sur $S$.}

Ce sont, à des modifications d'éclairage près, celles du diagramme p. 76. Les catégories $\ul{Fibsc}\,S$, $\ul{Fib}\,S$, $\ul{Cofibsc}\,S$, $\ul{Cofib}\,S$ et leurs variantes à indice 0, sont inchangées - leur rôle est d'être aussi petites

\msp{page 175}

que possible pour remplir leur fonction (de description de $\mathrm{HOT}_W(S)$, $\mathrm{HOT}^{o}_W(S)$, $\mathrm{HOT}^{\mathrm{lc}}_W(S)$, ou $\mathrm{HOT}^{\mathrm{lc}!}_W(S)$ ci-dessous). La catégorie $\ul{Parf}\,S$ du diagramme remplissait une fonction du même type, c'est pourquoi nous la remplaçons à présent (en surenchérissant) par $\ul{Parf}_\omega S$. Les autres catégories du diagramme sont, en plus de $\mathrm{Cat}/S$, $\ul{Liss}\,S$ et $\ul{Prop}\,S$ et leurs variantes à indice 0, ainsi que $\ul{Fibhot}\,S$. (Toutes ces catégories, ainsi que $\ul{Parf}\,S$, étaient donc construites par référence à $W_\infty$, non pas à $W$.) Leur rôle est d'être aussi grandes que possible pour donner lieu à certaines propiétés et relations. Or il se trouve qu'ici la référence à $W_\infty$ était artificielle, et il y a lieu de mettre plutôt les catégories $\ul{Liss}_W S$, $\ul{Prop}_W S$, $\ul{Fib}_W S$, et en plus les variantes $\ul{Liss}_{0,W} S$, $\ul{Prop}_{0,W} S$. Ainsi le diagramme devient % [sic] propiétés -> propriétés, misprint
\[
\begin{tikzcd}[column sep=tiny, row sep=small]
 & & & & \ul{Fib}_W S \arrow[dd, hook, "\delta"] \arrow[dddllll, dotted, bend right=35, "\Phi_0"'] \arrow[dddrrrr, dotted, bend left=35, "\Psi_0"] & & & & \\
 & & \ul{Liss}_{0,W} S \arrow[urr, hook] \arrow[dd, hook] & & & & \ul{Prop}_{0,W} S \arrow[ull, hook'] \arrow[dd, hook] & & \\
 & \ul{Fib}_{0,W} S \arrow[ur, hook] \arrow[dd, hook] & & & \mathrm{Cat}/S \arrow[dd, "C"] \arrow[dddllll, dotted, bend right=30, "\Phi"] \arrow[dddrrrr, dotted, bend left=30, "\Psi"'] & & & \ul{Cofib}_{0,W} S \arrow[ul, hook'] \arrow[dd, hook] & \\
\ul{Fibsc}_{0,W} S \arrow[ur] \arrow[dd, hook] & & \ul{Liss}_W S \arrow[urr, hook] & & & & \ul{Prop}_W S \arrow[ull, hook'] & & \ul{Cofibsc}_{0,W} S \arrow[ul] \arrow[dd, hook] \\
 & \ul{Fib}\,S \arrow[ur, hook] & & & \ul{Parf}_\omega S \arrow[uuull, hook] \arrow[uuurr, hook'] & & & \ul{Cofib}\,S \arrow[ul, hook'] & \\
\ul{Fibsc}\,S \arrow[ur] & & & & & & & & \ul{Cofibsc}\,S \arrow[ul]
\end{tikzcd}
\]
(\footnote{NB la nouvelle notation serait plutôt $\ul{Fibf}_W S$, car il s'agit de $W$-fibrations qu'à présent j'appelle `faibles' (et définies seulement dans Cat, à l'opposé des $W$-fibrations, définissables dans toute catégorie de modèles).}) (tous les carrés \texttt{[sont]} cartésiens, les foncteurs notés $\hookrightarrow$ sont pleinement fidèles),

\printedpage{108}
% head: A. Grothendieck : Dérivateurs, Ch. VI, Édition des Universités de Montpellier II et Paris VII
% foot: Version 9 novembre 2001

\msp{page 176}

où j'ai encore marqué en vert \texttt{[dans l'original seulement]} les cinq \textit{`foncteurs d'inversion'}
\[
\Phi,\ \Psi,\ \Phi_0,\ \Psi_0,\ C,
\]
qui fournissent des quasi-inverses des foncteurs sur les catégories localisées déduits par les \textit{foncteurs d'inclusion.}

\circled{3} \textbf{Les quatre types de catégories} $\mathrm{HOT}_W$\textbf{.}

Sur chacune des catégories de l'etage, il y a \textit{deux} localiseurs à regarder, % [sic] l'etage -> l'étage, missing accent (printed « l'étage » below)
\[
\begin{tikzcd}[row sep=huge, column sep=0pt]
W_S^{\mathrm{f}} \arrow[d, Rightarrow] & {} = W_S^{\mathrm{g}} = W_S^{\mathrm{d}} = W_S^{\mathrm{univ}} \\
W_S, &
\end{tikzcd}
\]
donnant lieu à deux localisés distincts, notés
\[
\mathrm{HOT}^{\mathrm{lc}}_W(S) \qquad \text{et} \qquad \mathrm{HOT}^{\mathrm{lc}!}_W(S),
\]
il se trouve que (grâce à $\Phi_0$ et $\Psi_0$) ces localisés `sont les mêmes' pour les 8 catégories de l'étage, à l'exception de $\ul{Parf}_\omega S$; plus exactement, les foncteurs d'inclusion (à l'exception peut-être de ceux ayant pour source $\ul{Parf}_\omega S$) induisent des équivalences sur les deux types de localisés, les équivalences quasi-inverses étant obtenues par $\Phi_0$ (aile gauche) ou $\Psi_0$ (aile droite). Quand on localise \textit{pour} $W_S$ (à l'exclusion de $W_S^{\mathrm{f}}$), on trouve aussi des quasi-inverses par $C|\ul{Liss}_{0,W} S$ resp. $C|\ul{Prop}_{0,W} S$, pour les deux inclusions
\[
\begin{tikzcd}
\ul{Liss}_{0,W} S & \ul{Parf}_\omega S \arrow[l, hook'] \arrow[r, hook] & \ul{Prop}_{0,W} S
\end{tikzcd}
\]
(\footnote{cf. plus bas, ce qu'on peut dire pour les localisés pour $W_S^{\mathrm{f}}$.}). Par raison de symétrie, le plus joli semble de définir
\[
\left\{
\begin{array}{lcl}
\mathrm{HOT}^{\mathrm{lc}}_W(S) & = & (\ul{Fib}_W S)(W_S^{\mathrm{f}})^{-1}\\
\mathrm{HOT}^{\mathrm{lc}!}_W(S) & = & (\ul{Fib}_W S)(W_S)^{-1},
\end{array}
\right.
\]
et on a un foncteur de localisation
\[
\mathrm{HOT}^{\mathrm{lc}}_W(S) \xrightarrow{\ \lambda\ } \mathrm{HOT}^{\mathrm{lc}!}_W(S)
\]
induit par l'inclusion
\[
W_S^{\mathrm{f}} \subseteq W_S.
\]


\endgroup
% END UNIT Der-ChVI-345/Der-ChVI-345_p106-110.tex
% BEGIN UNIT Der-ChVI-345/Der-ChVI-345_p111-115.tex
\begingroup\makeatletter%
\NoAutoSpacing%
\newcommand{\printedpage}[1]{\marginpar{\footnotesize p.~#1}}%
\newcommand{\msp}[1]{\marginpar{\footnotesize [#1]}}%
\newcommand{\Fib}{\underline{\mathrm{Fib}}}%
\newcommand{\Parf}{\underline{\mathrm{Parf}}}%
\newcommand{\Prop}{\underline{\mathrm{Prop}}}%
\newcommand{\Liss}{\underline{\mathrm{Liss}}}%
\newcommand{\Fibsc}{\underline{\mathrm{Fibsc}}}%
\newcommand{\Cofibsc}{\underline{\mathrm{Cofibsc}}}%
\newcommand{\HOT}{\mathrm{HOT}}%
\newcommand{\Cat}{\mathrm{Cat}}%
\newcommand{\Hom}{\mathrm{Hom}}%
\newcommand{\Ob}{\mathrm{Ob}}%
\newcommand{\lcs}{\mathrm{lc}\,!}%
\newcommand{\cercle}[1]{\tikz[baseline=(c.base)]\node[draw,circle,inner sep=1pt](c){#1};}%
\newcommand{\troisbarres}{\leavevmode\llap{$|||$\quad}}%
\makeatother\let\printedpage\derpp\let\msp\dermsp%
\let\mspin\DERmspin%


\printedpage{109}
% head: A. Grothendieck : Dérivateurs, Ch. VI, Édition des Universités de Montpellier II et Paris VII

\noindent\msp{page 177}
Pour toute sous-catégorie pleine
\[
\mathcal{F} \subseteq \Fib_W S
\]
qui contient l’une des trois sous-catégories
\[
\Fib_{0,W}S,\ \Parf_\omega S,\ \Prop_{0,W}S,
\]
l’inclusion de $\mathcal{F}$ dans $\Fib_W S$ définit des équivalences
\[
\mathcal{F}(W_S)^{-1} \approx \HOT^{\lcs}_W(S),
\]
et si $\mathcal{F}$ contient l’une des deux sous-catégories $\Fib_{0,W}S$, $\Prop_{0,W}S$ (à l’exclusion de $\Parf_\omega S$), on trouve aussi
\[
\mathcal{F}(W^{\mathrm{f}}_S)^{-1} \approx \HOT^{\mathrm{lc}}_W(S).
\]
Le rez-de-chaussée contribue deux autres variantes de HOT, en localisant soit par $W^{\mathrm{d}}_S$ (aile gauche), soit par $W^{\mathrm{g}}_S$ (aile droite). Les foncteurs d’inclusion définissent des équivalences pour les localisées soit pour $W^{\mathrm{d}}_S$ (aile gauche), soit pour $W^{\mathrm{g}}_S$ (aile droite). Le plus joli (symétrie) semble de définir ces deux localisés par $\Cat/S$
\[
\HOT_W(S) = (\Cat/S)(W^{\mathrm{d}}_S)^{-1}, \qquad \HOT^{o}_W(S) = (\Cat/S)(W^{\mathrm{g}}_S)^{-1}.
\]
Les quasi-inverses pour les équivalences de catégories localisées, induites par les foncteurs d’inclusion, sont donnés soit par $\Phi$ (aile gauche), soit par $\Psi$ (aile droite).

Mais $\Phi$ et $\Psi$ fournissent des quasi-inverses également quand on localise pour $W_S$, la $W$-équivalence grossière au lieu de la fine - pareil qu’à l’étage. On trouve donc au rez-de-chaussée un (unique) troisième localisé, qu’on peut décrire plus joliment par
\[
(\Cat/S)(W_S)^{-1},
\]
mais il est équivalent à $(\Fib_W S)W_S^{-1}$ par le
\msp{page 178}
foncteur induit par l’inclusion. (Un quasi-inverse étant fourni par $C$.) C’est donc $\HOT^{\lcs}_W(S)$. On trouve ainsi, par le diagramme commutatif de catégories avec localiseurs
\[
\begin{tikzcd}[column sep=small]
 & (\Fib_W S, W^{\mathrm{f}}_S) \arrow[dl] \arrow[dd] \arrow[dr] & \\
(\Cat/S, W^{\mathrm{d}}_S) \arrow[ddr] & & (\Cat/S, W^{\mathrm{g}}_S) \arrow[ddl] \\
 & (\Fib_W S, W_S) \arrow[d] & \\
 & (\Cat/S, W_S) &
\end{tikzcd}
\]
% foot: Version 9 novembre 2001
\printedpage{110}
% head: A. Grothendieck : Dérivateurs, Ch. VI, Édition des Universités de Montpellier II et Paris VII
(où la flèche verticale courte induit une équivalence sur les localisés), le diagramme commutatif entre les catégories $\HOT_W$ liées à $S$ :
\[
\begin{tikzcd}[column sep=small]
 & \HOT^{\mathrm{lc}}_W(S) \arrow[dl, "i"'] \arrow[dd, "\lambda"'] \arrow[dr, "j"] & \\
\HOT_W(S) \arrow[dr, "\rho"'] & & \HOT^{o}_W(S)(\simeq \HOT_W(S^{o})) \arrow[dl, "\sigma"] \\
 & \HOT^{\lcs}_W(S), &
\end{tikzcd}
\tag*{$(\ast)$}
\]
où $\rho$, $\sigma$, $\lambda$ sont des foncteurs de localisation, et où $i$, $j$ sont ressentis comme des foncteurs d’inclusion. Ce diagramme $(\ast)$ semble épuiser les foncteurs naturels qui relient entre elles les quatre catégories envisagées.

\noindent\msp{page 179}
Les propriétés des foncteurs $\Phi$, $\Psi$, $C$ qui ont été utilisées, pour fournir tous les quasi-inverses voulus de foncteurs (entre catégories localisées) induits par les foncteurs d’inclusion du diagramme, sont les suivantes.

\cercle{a} On a un homomorphisme fonctoriel dans $\Cat/S$
\[
X \xrightarrow{i_X} \Phi(X), \qquad \underline{i_X \in W^{\mathrm{d}}_S}.
\]
Quand $X$ est dans $\Fib_W S$, $\Phi(X)$ est dans $\Fibsc_0 S$ \texttt{[plutôt }$\Fibsc_{0,W}S$\texttt{, cf. diag. p. 175]}, donc $i_X \in W^{\mathrm{f}}_S$ ($= W^{\mathrm{d}}_S$ sur $\Fib_W S$). Ainsi $\Phi$ induit le foncteur $\Phi_0$ du diagramme.

Dualement, on a
\[
X \xrightarrow{i_X} \Psi(X), \qquad \underline{i_X \in W^{\mathrm{g}}_S},
\]
et quand $X$ est dans $\Fib_W S$, $\Psi(X)$ est dans $\Cofibsc_0 S$ \texttt{[plutôt }$\Cofibsc_{0,W}S$\texttt{]}, d’où le foncteur $\Psi_0$.

\cercle{b} On a un homomorphisme fonctoriel
\[
\left|
\begin{aligned}
& X \xrightarrow{i_X} C(X), \qquad \underline{i_X \in W_S} \\
& \text{et } C(X) \in \Ob\,\Parf_\omega S.
\end{aligned}
\right.
\]

\troisbarres Mais on fera attention qu’en l’absence de l’axiome W(4b) sur $W$, on ne peut \textit{pas} affirmer $i_X \in W^{\mathrm{f}}_S$, même si $X$ est dans $\Fib_W S$ (voire, dans $\Parf_\omega S$), seul cas (de toutes façons) où on est en droit de s’y attendre. Je note à ce propos le diagramme \textit{cartésien,} pour $f : X \longrightarrow Y$ dans $\Cat/S$
% foot: Version 9 novembre 2001
\printedpage{111}
% head: A. Grothendieck : Dérivateurs, Ch. VI, Édition des Universités de Montpellier II et Paris VII
\[
\begin{tikzcd}
C(Y)_s \arrow[d, "\pi^{s}_{Y}"'] & C(X)_s \arrow[l, "C(f)_s"'] \arrow[d, "\pi^{s}_{X}"] \\
Y & X, \arrow[l, "f"']
\end{tikzcd}
\qquad \text{\footnotesize flèches verticales dans $\Parf_\omega \subseteq \Fib_W$}
\]
qui montre par exemple que
\msp{page 180}
l’on ne peut pas s’attendre que $f \in W_S \implies C(f) \in W^{\mathrm{f}}_S$ sans disposer de l’axiome W(4a$'$) \texttt{[qu’est-ce que c’est?]} (ou d’un axiome de force comparable, tel W(4b)).

Ce sont ces faits négatifs (en l’absence de W(4b)) qui sont cause que je ne peux rien dire de particulier sur les foncteurs induits par \textit{localisation fine} $(W^{\mathrm{f}}_S)^{-1}$ des deux foncteurs d’inclusion
\[
\Liss_0 S \longleftarrow \Parf_\omega S \longrightarrow \Prop_0 S,
\]
ou, ce qui revient au même, de celui induit par le composé commun de ceux-là avec les inclusions dans $\Fib_W S$
\[
\begin{tikzcd}[column sep=tiny, row sep=small]
 & \Parf_\omega S \arrow[dl] \arrow[dd] \arrow[dr] & \\
\Liss_0 S \arrow[dr] & & \Prop_0 S \arrow[dl] \\
 & \Fib_W S, &
\end{tikzcd}
\]
en d’autres termes, je ne peux rien dire sur les propriétés de fidélité du foncteur
\[
(\Parf_\omega S)(W^{\mathrm{f}}_S)^{-1} \longrightarrow (\Fib_W S)(W^{\mathrm{f}}_S)^{-1},
\]
et encore moins que ce soit une équivalence de catégories. Donc, en l’absence de W(4b), pour les sous-catégories pleines $\mathcal{F}$ de $\Fib_W S$ qui contiennent $\Parf_\omega S$, sans contenir ni $\Liss_0 S$ ni $\Prop_0 S$, il se pourrait a priori que les localisées
\[
\mathcal{F}(W^{\mathrm{f}}_S)^{-1}
\]
soient toutes différentes. Je pense notamment à $\Parf_W S$, aussi intéressante en elle-même, après tout, que $\Parf_\omega S$. Donc en plus des quatre localisées du diagramme $(\ast)$, p. 178, il y a toute une série infinie ‘tératologique’ de localisées
\msp{page 181}
qui refusent de se laisser ‘boucler’ \texttt{[?]} par $\Phi_0$ ou $\Psi_0$, et pour lesquelles même $C$ n’est d’aucun secours, faute de l’axime W(4b). % [sic] axime -> axiome, missing letter

\subsection*{\cercle{4} Cas où l’axiome \textmd{W(4b)} est satisfait.}
% foot: Version 9 novembre 2001

\printedpage{112}
% head: A. Grothendieck : Dérivateurs, Ch. VI, Édition des Universités de Montpellier II et Paris VII
Alors les perplexités que je viens de dire disparaissent. On a aussitôt, en vertu de W(4b) :
\[
|||\qquad
\begin{array}{@{}r@{}}
\text{Si } X \in \Ob\,\Fib_W S, \text{ alors } i_X : X \longrightarrow C(X) \text{ est dans } W^{\mathrm{f}}_S \\
(\textit{en supposant l’axiome } \text{W(4b)})
\end{array}
\]
et $C$ fournit des quasi-inverses pour les foncteurs ‘d’inclusion’ entre localisés fins (pour $W^{\mathrm{f}}_S$) envisagés à la page précédente. Il en fournit \texttt{[‘Il en fournit’ barré]} un également pour le foncteur $\lambda$ envisagé plus haut ($(\ast)$ page 178)
\[
\underbrace{(\Fib_W S)(W^{\mathrm{f}}_S)^{-1}}_{\overset{\text{déf}}{=}\HOT^{\mathrm{lc}}_W(S)} \longrightarrow \underbrace{(\Cat/S)(W_S)^{-1}}_{\simeq \HOT^{\lcs}_W(S)},
\]
lequel est donc une équivalence. Une autre façon, plus évidente, de dire la même chose, c’est que sur chacune des catégories de l’étage les deux localiseurs $W^{\mathrm{f}}_S$ et $W_S$ sont désormais identiques (puisqu’elles sont toutes contenues dans $\Fib_W S$, où elles \texttt{[plutôt ils]} le sont justement par l’axiome W(4b)), donc le foncteur $\lambda$ entre ces deux localisés (définition originale de $\lambda$, dans le cas $\mathcal{F} = \Fib_W S$, cf. page 176) est même le \textit{foncteur identique,} les deux catégories localisées étant \textit{égales :}
\[
\HOT^{\mathrm{lc}}_W(S) = \HOT^{\lcs}_W(S) \qquad\qquad !
\]
Ainsi le diagramme $(\ast)$ p. 178 prend la forme
\[
\begin{tikzcd}[column sep=small]
 & \HOT^{\mathrm{lc}}_W(S) \arrow[dl, shift left=1.2, "i"] \arrow[dr, shift right=1.2, "j"'] & \\
\HOT_W(S) \arrow[ur, shift left=1.2, "\rho"] & & \HOT^{o}_W(S). \arrow[ul, shift right=1.2, "\sigma"']
\end{tikzcd}
\tag*{$(\ast)$}
\]

\noindent\msp{page 182}
Quand on interprète non seulement les deux catégories latérales $\HOT_W(S)$, $\HOT^{o}_W(S)$, mais aussi celle du milieu, comme des localisées de $\Cat/S$, pour les localiseurs $W^{\mathrm{d}}_S$, $W^{\mathrm{g}}_S$ et $W_S$ respectivement, les foncteurs $\rho$, $\sigma$ sont ceux déduits des deux inclusions
\[
\begin{tikzcd}[column sep=small, row sep=small]
W^{\mathrm{d}}_S \arrow[dr, hook] & & W^{\mathrm{g}}_S \arrow[dl, hook'] \\
 & W_S &
\end{tikzcd}
\qquad \text{\footnotesize comme parties de Fl(Cat/$S$),}
\]
tandis que les foncteurs $i$, $j$ sont déduits, par passage aux localisés, d’un seul et même foncteur
\[
\Cat/S \xrightarrow{C} \Cat/S,
\]
ce foncteur ayant la vertu d’appliquer $W_S$ dans $W^{\mathrm{d}}_S \cap W^{\mathrm{g}}_S$,
\[
f \in W_S \overset{\text{moyennant W(4b)}}{\Longrightarrow} C(f) \in W^{\mathrm{d}}_S \cap W^{\mathrm{g}}_S \qquad (\text{et même } \in W^{\mathrm{univ}}_S).
\]
% foot: Version 9 novembre 2001

\printedpage{113}
% head: A. Grothendieck : Dérivateurs, Ch. VI, Édition des Universités de Montpellier II et Paris VII
(Et pour cause, puisque $C$ se factorise par $\Fib_W S$ (et même par $\Parf_\omega S$), où $W_S$ est identique à $W^{\mathrm{d}}_S$, $W^{\mathrm{g}}_S$, $W^{\mathrm{univ}}_S$. Que $f \in W_S$ implique $C(f) \in W$ est indépendant de l’axiome W(4b), et se voit sur le diagramme
\[
\begin{tikzcd}
X \arrow[r, "f"] \arrow[d, "i_X"'] & Y \arrow[d, "i_Y"'] \\
C(X) \arrow[r, "C(f)"] & C(Y),
\end{tikzcd}
\]
où les flèches $i_X$, $i_Y$ sont dans $W_S$.)

Ceci dit, il est vrai que les couples
\[
(\rho, i),\ (\sigma, j)
\]
sont des couples de foncteurs adjoints, donc ($\rho$ et $\sigma$ étant des foncteurs de localisation) que $i$ et $j$ sont \textit{pleinement fidèles.} La formule d’adjonction peut se lire comme
\[
\boxed{\Hom_{W^{\mathrm{g}}_S}(X, CY) \simeq \Hom_{W^{\mathrm{d}}_S}(X, CY) \simeq \Hom_{W_S}(X, Y)}
\tag*{$(\ast)$}
\]
\msp{page 183}
pour $X$, $Y$ dans $\Cat/S$ sans plus. Un $\Hom_\Sigma$ indique le Hom dans $(\Cat/S)\Sigma^{-1}$.

\noindent\textbf{Corollaire.} \textit{Si $X$ dans} $\Cat/S$\textit{, $Y$ dans} $\Fib_W S$\textit{, alors}
\[
\begin{tikzcd}[column sep=small]
\Hom_{W^{\mathrm{g}}_S}(X, Y) \arrow[dr, "\sim"] & & \Hom_{W^{\mathrm{d}}_S}(X, Y) \arrow[dl, "\sim"'] \\
 & \Hom_{W_S}(X, Y) &
\end{tikzcd}
\]
\textit{sont bijectifs.}

En effet, $Y$ dans $\Fib_W(S)$ implique que $Y \xrightarrow{i_Y} C(Y)$ est dans $W^{\mathrm{univ}}_S \subseteq W^{\mathrm{d}}_S \cap W^{\mathrm{g}}_S$, donc
\[
\Hom_{W^{\mathrm{g}}_S}(X, Y) \xrightarrow{\sim} \Hom_{W^{\mathrm{g}}_S}(X, CY) \underbrace{\xrightarrow{\sim}}_{\text{adjonction}} \Hom_{W_S}(X, Y),
\]
et itou pour $\Hom_{W^{\mathrm{d}}_S}(X, Y)$.

\subsection*{\cercle{5} Perplexités sur les variantes affaiblies de l’axiome \textmd{W(4b)}.}

Que reste-t-il de toutes ces belles choses (plus que trois catégories HOT, au lieu de $4+\infty$!) quand on ne suppose plus que W(4b$'$) au lieu de W(4b), i.e. l’égalité de $W_S$ et $W^{\mathrm{f}}_S$ dans
% foot: Version 9 novembre 2001


\endgroup
% END UNIT Der-ChVI-345/Der-ChVI-345_p111-115.tex
% BEGIN UNIT Der-ChVI-345/Der-ChVI-345_p116-120.tex
\begingroup\makeatletter%
\newcommand{\printedpage}[1]{\marginpar{\footnotesize p.~#1}}%
\newcommand{\msp}[1]{\marginpar{\footnotesize [#1]}}%
\newcommand{\Parf}{\underline{\mathrm{Parf}}}%
\newcommand{\Fib}{\underline{\mathrm{Fib}}}%
\newcommand{\Liss}{\underline{\mathrm{Liss}}}%
\newcommand{\Prop}{\underline{\mathrm{Prop}}}%
\newcommand{\Bifib}{\underline{\mathrm{Bifib}}}%
\newcommand{\Cofibsc}{\underline{\mathrm{Cofibsc}}}%
\newcommand{\Fibsc}{\underline{\mathrm{Fibsc}}}%
\newcommand{\HOT}{\mathrm{HOT}}%
\newcommand{\margbar}{\marginpar{\(\|\)}}%
\makeatother\let\printedpage\derpp\let\msp\dermsp%
\let\mspin\DERmspin%


\printedpage{114}
% head: A. Grothendieck : Dérivateurs, Ch. VI, Édition des Universités de Montpellier II et Paris VII
% foot: Version 9 novembre 2001

$\Parf_\omega S$? Finalement, il me semble (contrairement à ce que j'avais cru voir) qu'il ne reste presque rien (\footnote{\textbf{NB} Si $W$ satisfait l'axiome des limites, alors W(4b$'$) implique W(4b), cf. théorème récapitulatif dans XVI 4.4.13 (p. 58-62), notamment $2^{o}$). La condition W(4b) est la condition B de cet énoncé, et W(4b$'$) la condition D. (Elle est impliquée par W(4a) $=$ C, ce que je soupçonnais nullement en écrivant XII \dots)}). On ne sait alors toujours pas, pour $X$ dans $\Fib_W S$, si
\[ i_X : X \longrightarrow C(X) \]
n'est dans $W^{\mathrm{f}}_S$, et non seulement dans $W_S$; on ne le sait que pour $X \in \Parf_\omega S$, mais dans ce cas on n'en a pas besoin, car l'utilité du foncteur $C(X)$ et de l'inclusion $i_X : X \longrightarrow C(X)$, c'est ‘d'améliorer’ $X$ et de le rendre parfait, voire même $W$-parfait. Mais s'il l'est déjà, c'est plus la peine. Il y aura, comme
\msp{page 184}
avant, 4 localisées essentiellement distinctes $\HOT_W$, $\HOT^{o}_W$, $\HOT^{\mathrm{lc}}_W$, $\HOT^{\mathrm{lc}\,!}_W$, plus une infinité d'autres $\mathcal{F}(W^{\mathrm{f}}_S)^{-1}$ pour $\mathcal{F}$ entre $\Parf_\omega S$ et $\Fib_W S$ (cf. page 180), avec le seul bénéfice qu'on peut maintenant exclure le cas $\mathcal{F} = \Parf_\omega S$, car dans ce cas $\mathcal{F}(W^{\mathrm{f}}_S)^{-1} = \mathcal{F}(W_S)^{-1} = \HOT^{\mathrm{lc}\,!}_W(S)$ fait partie des quatre canoniques. La chose décisive, c'est qu'on n'a toujours aucun moyen de prouver que
\[ \lambda : \HOT^{\mathrm{lc}}_W(S) \longrightarrow \HOT^{\mathrm{lc}\,!}_W(S) \]
est une équivalence, et je doute que ce soit le cas. Il faudrait regarder avec $W = W_{\mathbf{D}}$ la $\mathbf{D}$-équivalence, où $\mathbf{D}$ est un dérivateur Der($k$-Mod), $k$ un anneau donné (p.ex. $k = \mathbf{Z}$, ou $k$ un corps $\mathbf{Z}/p\mathbf{Z}$ ou $\mathbf{Q}$ \dots). Je ne vais pas le faire maintenant cependant. (Il est vrai que l'axiome W(4b$'$) n'est pas non plus satisfait dans ce cas \dots)

Si dans l'énoncé de l'axiome, on remplaçait $\Parf_\omega S$ par $\Parf_W S$, rien d'essentiel ne serait changé - tout au plus je remplacerais $\Parf_\omega S$ par $\Parf_W S$ dans le brillant diagramme p. 175. Finalement, je suis à présent convaincu qu'il faut tout ou rien : ou bien \textit{quatre} catégories $\HOT(S)$ pour $W$, en se contentant de W$(1,2,3)$, ou bien rien que trois, i.e. pas de $\HOT^{\mathrm{lc}\,!}_W(S)$ - et pour ça, il ne semble pas qu'on puisse se contenter de moins de W(4b).

\msp{page 185}
Il suffirait, il est vrai, de supposer que $W_S = W^{\mathrm{f}}_S$ sur l'une quelconque des sept catégories de l'étage autres que $\Parf_\omega S$. Cela impliquerait alors que dans les six autres, on aura $\overline{W_S} = \overline{W^{\mathrm{f}}_S}$ (mêmes saturés forts dans $\mathrm{Fl}(\mathcal{F})$). Un des deux les plus faibles des axiomes en question serait : % doubt resolved in DER_VI-c_corrections.csv; input note: line under "étage" -- read as the overline of $\overline{W^{\mathrm{f}}_S}$ on the next line, not an underline of "étage"

\begin{quote}
Si $f : X \longrightarrow Y$ est un foncteur \textit{cartésien} de catégories Cat-fibrées et $W$-fibrées (\footnote{‘$W$-fibrées’ signifiant ici que pour $u \in \mathrm{Fl}\,S$, les foncteurs $u^*$ de changement de base $X_{s'} \longrightarrow X_s$, $Y_{s'} \longrightarrow Y_s$ sont dans $W$.}) sur une autre $S$, si $f \in W$, alors $f_s \in W$ pour tout $s \in S$,
\end{quote}

(et l'autre est la version duale, pour les foncteurs cocartésiens de catégories cofibrées). Sauf erreur, cet axiome doit déjà impliquer que $\overline{W}$ satisfait la version W(4a) la plus forte
\printedpage{115}
% head: A. Grothendieck : Dérivateurs, Ch. VI, Édition des Universités de Montpellier II et Paris VII
% foot: Version 9 novembre 2001
de toutes (a fortiori W(4b)) des axiomes W(4). Cet axiome me paraît donc à présent beaucoup plus raisonnable que W(4b$'$), et ceci d'autant plus qu'il se réfère directemment à la description de $\HOT^{\mathrm{lc}}_W(S)$ en termes de foncteurs (ou de ‘pseudofoncteurs’, si on ne veut supposer les catégories scindées et les foncteurs compatibles aux scindages) $S^{o} \longrightarrow \mathrm{Cat}$. % [sic] directemment -> directement, spelling
D'ailleurs, bien qu'il n'en ait pas l'air, cet axiome \textit{est} autodual, puisque la version duale, pour $S$ fixée, est la version directe pour $S^{o}$! Cela rend cet axiome encore plus appétissant. Donc pour formuler W(4a$'$), j'ai envie à présent de laisser tomber $\Parf_\omega$, et de faire pareil. Si on garde deux versions W(4a$'$) (l'ancienne) et W(4a$''$) (la nouvelle), l'une et l'autre ont la vertu d'impliquer W(4b). Et W(4b$'$) (nouvelle version) doit impliquer W(4a) (version ultraforte), sinon pour $W$, du moins pour $\overline{W}$.

Mes idées étaient tenacement [?] embrouillées
\msp{page 186}
jusqu'à aujourd'hui même, car j'avais l'idée que, grâce au mirifique foncteur $C$, la catégorie $\Parf_\omega S$, comme les sept autres catégories de l'étage, pourrait servir à définir $\HOT^{\mathrm{lc}}_W(S)$. Or il n'en est rien, semble-t-il, à moins de supposeer l'axiome W(4b), ou la variante W(4b$'$) (nouvelle version). L'ancienne version de W(4b$'$) s'avère parfaitement inutilisable. % [sic] supposeer -> supposer, spelling

À ce propos, je n'ai toujours pas répondu clairement pourtant, pas un ‘oui’ ou un ‘non’, à la question :

\begin{quote}
\textbf{Question :} \textit{les inclusions}
\[
\begin{tikzcd}
 & \Liss_{0,W}S \\
\Parf_W S \arrow[ur] \arrow[dr] & \\
 & \Prop_{0,W}S
\end{tikzcd}
\]
\textit{induisent-elles, même en l'absence de W(4b) (ou de ses variantes affaiblies) des équivalences pour les localisées fines (par $W^{\mathrm{f}}_S$)? Même question pour} $\Parf_\omega S$ \textit{au lieu de} $\Parf_W S$ (\footnote{OK en présence de l'axiome W(6bis) sur les $\varinjlim$ , cf. p. 192 ff.}). \textit{Même question pour}
\[ \Bifib\, S \subseteq \Parf_\omega S, \]
\textit{la catégorie des catégories sur $S$ qui sont à la fois} Cat-\textit{fibrées et} Cat-\textit{cofibrées.}
\end{quote}

Je ne me rends pas compte, même en admettant (cette fois) l'axiome W(4a), si $\Bifib\, S$ est assez grosse pour pouvoir décrire tous le types d'homotopie relatifs localement constants sur $S$, i.e. si % [sic] tous le types -> tous les types, missing letter
\[ (\Bifib\, S)(W^{\mathrm{f}}_S)^{-1} \longrightarrow (\Fib_W S)(W^{\mathrm{f}}_S)^{-1} \]
est au moins essentiellement surjectif, sinon une équivalence de catégories. Il manque un foncteur
\[ \mathrm{Cat}/S \longrightarrow \Bifib\, S \qquad \text{ou du moins} \qquad \underbrace{\Fib_W S}_{\text{ou $\Fib_\omega S$\,?}} \longrightarrow \Bifib\, S, \]
\printedpage{116}
% head: A. Grothendieck : Dérivateurs, Ch. VI, Édition des Universités de Montpellier II et Paris VII
% foot: Version 9 novembre 2001
qui puisse jouer le rôle de $\Phi$ et de $\Psi$ \textit{simultané\msp{page 187}ment.} Je n'ai pas même vérifié que $\Phi$ ne transforme \textit{pas} objets de $\Prop_W S$ ou de $\Prop_{0,W}(S)$ en objets $W$-propres (donc $W$-parfaits, étant de toutes façons $W$-lisse et même $W_\omega$-lisse) sur $S$ (\footnote{Voir réponse à cette intéressante question page 192.}). Si c'était pourtant vrai sur $\Prop_W S$, ça montrerait que l'image essentielle de $\Prop_W S$ dans $\HOT_W(S)$ ($= (\mathrm{Cat}/S)(W^{\mathrm{d}}_S)^{-1}$) est contenu dans $\HOT^{\mathrm{lc}}_W(S)$. % [sic] contenu -> contenue, agreement with "l'image"
Prenant $S = \boldsymbol{\Delta}^1$, cela signifierait que si dans $X$ on a une sous-catégorie fermée $X_1$ telle que $X_1 \overset{i}{\hookrightarrow} X$ soit $W_\omega$-asphérique, alors $1\backslash X \longrightarrow 0\backslash X$ est dans $W$ - or $0\backslash X \simeq X$, $1\backslash X \simeq X_1$, la flèche est $X_1 \longrightarrow X$, qui est bien dans $W$ étant $W$-asphérique! C'est étrange et quasiment immoral (\footnote{On observe plus bas que c'est ainsi chaque fois que $S$ a un objet final.}). Mais à force de gribouillis je viens de me convaincre que même si $X$ est cofibrée sur $S$ (donc aussi $W$-propre), $\Phi(X)$ n'est pas $W$-propre (ce qui impliquerait $W$-parfait), ni
\msp{page 188}
même un $W$-fibré, ce qui signifierait que
\[ s'\backslash X \longrightarrow s\backslash X \qquad \text{associé à } s \longrightarrow s' \text{ dans } S \]
serait toujours dans $W$. Posant $Z' = s'\backslash S$, $Z = s\backslash S$, cela signifierait que si on a un \textit{co-isomorphisme local}
\[ i : \underbrace{Z'}_{\simeq z_0\backslash Z} \longrightarrow Z, \]
où $Z'$, $Z$ ont des éléments initiaux, et si $X$ est cofibrée sur $Z$, $X'$ son image inverse sur $Z'$, alors $X' \longrightarrow X$ est dans $W$. En fait, si $Z$ a aussi un élément \textit{final} $e$ (comme c'était le cas tantôt pour $\boldsymbol{\Delta}^1$), alors $Z' \simeq z_0\backslash Z$ en a aussi en $e'$, qui s'envoie sur $e$, et on aura le diagramme commutatif
\[
\begin{tikzcd}
X' \arrow[r, "j"] & X \\
X'_{e'} \arrow[u, "W"] \arrow[r, "\sim"] & X_e, \arrow[u, "W"']
\end{tikzcd}
\]
d'où $j \in W$. Donc pour faire un contrexemple,
\msp{page 189}
il faut prendre $Z$ sans élément final.

Mais prenons $Z$ \texttt{[un]} ensemble ordonné, avec un objet initial $0$, mais sans objet final, et prenons pour $z_0$ un élément \textit{maximal} de $Z$, de sorte que $Z' = z_0\backslash Z$ soit isomorphe à $\{z_0\}$. Il faut donner dans Cat un exemple où l'inclusion de $X_{z_0}$ dans $X$ n'est pas dans $W$. On va prendre
\printedpage{117}
% head: A. Grothendieck : Dérivateurs, Ch. VI, Édition des Universités de Montpellier II et Paris VII
% foot: Version 9 novembre 2001
\[
\begin{tikzcd}
 & & a \\
Z = & 0 \arrow[ur] \arrow[dr] & \\
 & & b = z_0,
\end{tikzcd}
\]
donc la donnée de $X$ cofibrée coscindée sur $Z$ équivaut à un diagramme dans Cat
\[
\begin{tikzcd}
 & X_a \\
X_0 \arrow[ur, "\alpha"] \arrow[dr, "\beta"] & \\
 & X_b.
\end{tikzcd}
\]
Je pense que sous des conditions assez fortes (p.ex. si $\alpha$, $\beta$ sont des immersions ouvertes, ou des immersions fermées - à vérifier!), la catégorie $X$ est isomorphe dans Hot à la somme amalgamée dans Cat
\[ X_a{\textstyle\coprod_{X_0}}X_b, \]
et il n'y \texttt{[a]} aucune raison (sauf si on suppose $\alpha \in W$, et encore \dots) que $X_b \longrightarrow X_a \coprod_{X_0} X_b$ soit dans $W$. Prenons $X_0 = X_b = e$, $\beta = \mathrm{id}_e$, donc $\alpha$ est donné par la donnée d'un élément $t$ de $\mathrm{Ob}(X_a = T)$. La somme amalgamée est $T$, qui n'est pas isomorphe à $e$ dans $\mathrm{Hot}_W$,
\msp{page 190}
si $T$ n'a pas été prise $W$-asphérique. Mais je vais essayer de montrer directement que l'inclusion
\[ T = X_a \longrightarrow X \]
est $\in W$. Pour ceci je regarde la sous-catégorie ouverte $Z''$ de $Z$, complémentaire de la sous-catégorie pleine fermée $Z'$, et je note que $X'' = X|Z''$ étant propre sur $Z''$ et $a$ étant objet final de $Z''$, on a
\[ X_a \longrightarrow X'' \qquad \text{dans } W, \]
il reste à voir qu'il en est de même de $X'' \longrightarrow X$. Or $X''$ est déduite de $X$ en ajoutant un unique élément $\xi$ (celui de $X_b$), lequel ne s'envoie dans aucun élément de $X''$, et qui ne reçoit aucun élément de $X_a$ - par contre pour l'unique élément $\xi_0$ de $X_0$ il y a une unique flèche de celui-ci dans $\xi$. On a donc une rétraction de $X$ sur $X''$, envoyant $\xi$ sur $\xi_0$ et la flèche précédente sur $\mathrm{id}_{\xi_0}$. Est-ce un homotopisme? Il suffit de \texttt{[montrer]} qu'en tant qu'endomorphisme $\varphi$ de $X$, c'est $\sim \mathrm{id}_X$. Or on a bien $\mathrm{id}_X \longleftarrow \varphi$, dont la valeur sur $X''$ est l'identité, et la valeur sur $\xi$ est l'unique $\varphi(\xi) = \xi_0 \longrightarrow \xi$. C'est un peu bref, mais ça doit être vrai. (De toutes façons, $X$ est la réunion de deux ouverts $X''$, et de l'ouvert $X|\{0,b\} = X'''$, avec $X''' \simeq \boldsymbol{\Delta}^1$, et $X'' \cap X''' \simeq e$, donc au moins pour la cohomologie à coefficients localement constants de $X$ je suis à l'aise par Mayer-Vietoris, et sais bien que c'est celle du morceau $X''$, donc celle de $X_0$ \texttt{[ou $X_a$, indice illisible]} nullement celle de $e$!)

\printedpage{118}
% head: A. Grothendieck : Dérivateurs, Ch. VI, Édition des Universités de Montpellier II et Paris VII
% foot: Version 9 novembre 2001
\msp{page 191}
\margbar On a donc vu au moins ceci : Si $X$ est cofibrée donc $W$-propre sur $S$, il ne s'ensuit nullement que $\Phi(X)$ soit encore $W$-propre sur $S$ (\footnote{voir résultat plus précis page suivante!}), ni que ce soit un $W$-fibré. Mais que se passe-t-il si $X \in \Prop_{0,W}S$, par exemple s'il provient de $\Cofibsc_0 S$? Alors $X$ est dans $\Fib_W S$, donc on sait déjà que $\Phi(X)$ est dans $\Fibsc_{0,W}S$, donc ce n'est pas par là qu'on fera un contrexemple. Mais considérons encore $X$ cofibrée sur
\[
\begin{tikzcd}
 & & a \\
S = & 0 \arrow[ur] \arrow[dr] & \\
 & & b,
\end{tikzcd}
\]
donnée par le diagramme
\[
\begin{tikzcd}
 & X_a \\
X_0 \arrow[ur, "u"] \arrow[dr, "v"] & \\
 & X_b
\end{tikzcd}
\]
dans Cat, en supposant cette fois $u, v \in W$ ($X \in \Cofibsc_{0,W}S$). Alors $\Phi(X)$ est la catégorie fibrée scindée sur $S$ définie par le diagramme contravariant sur $S$
\[
\begin{tikzcd}
 & X_a \arrow[dl, hook, "i"'] \\
X_0 & \\
 & X_b, \arrow[ul, hook, "j"']
\end{tikzcd}
\]
où $i$, $j$ sont les inclusions (qui sont des immersions fermées). Dire que $\mathcal{X} = \Phi(X)$ est propre sur $S$ signifie aussi que ses restrictions au dessus des sous-catégories ouvertes $\{0,a\}$ et $\{0,b\}$ de $S$ le sont (car toute flèche $\boldsymbol{\Delta}^1 \longrightarrow S$ se factorise par l'une ou l'autre). Cela
\msp{page 192}
signifie donc aussi que l'inclusion des fibres fermées
\[ X_a \longrightarrow X, \qquad X_b \longrightarrow X \]
sont coasphériques. Mais $X_a \longrightarrow X$ ne peut être coasphérique que si $X_b = \emptyset$, car pour $x \in X_b \subseteq X$, $x\backslash X_a$ (ainsi que $X_a/x$) est vide, donc \textit{pas} $W$-asphérique (si on ne suppose $W = \mathrm{Fl}(\mathrm{Cat})$, cas que nous excluons). Donc $\Phi(X)$ n'est propre sur $S$ \textit{que} si $X_a$ et $X_b = \emptyset$ (résultat valable pour \textit{toute} $X$ sur $S$, qu'elle soit cofibrée sur $S$ ou $W$-propre, ou non). Si on veut $u, v \in W$, cela implique donc que $X_0 = \emptyset$, donc on a prouvé :


\endgroup
% END UNIT Der-ChVI-345/Der-ChVI-345_p116-120.tex
% BEGIN UNIT Der-ChVI-345/Der-ChVI-345_p121-125.tex
\begingroup\makeatletter%
\NoAutoSpacing%
\newcommand{\printedpage}[1]{\marginpar{\footnotesize p.~#1}}%
\newcommand{\msp}[1]{\marginpar{\footnotesize [#1]}}%
\newcommand{\und}[1]{\underline{\mathrm{#1}}}%
\newcommand{\circnum}[1]{\tikz[baseline=(c.base)]\node[draw,circle,inner sep=1pt](c){#1};}%
\makeatother\let\printedpage\derpp\let\msp\dermsp%
\let\mspin\DERmspin%


\printedpage{119}
% head: A. Grothendieck : Dérivateurs, Ch. VI, Édition des Universités de Montpellier II et Paris VII
% foot: Version 9 novembre 2001

% the following paragraph is marked by a triple vertical rule in the left margin
\noindent\llap{$|||$\quad}Soit $X \in \mathrm{Cat}/S$, $S = \begin{tikzcd}[row sep=tiny, column sep=small] & a \\ 0 \arrow[ur] \arrow[dr] & \\ & b \end{tikzcd}$. Si $\Phi(X)$ est propre sur $S$, alors $X = \emptyset$ (si on suppose $W \neq \mathrm{Fl}(\mathrm{Cat})$).

Mais chemin faisant est apparue une idée bizarre. Considérons le diagramme infini, pour $X$ dans $\mathrm{Cat}/S$ :
\[
X \longrightarrow \Phi(X) \longrightarrow \Psi\Phi(X) \longrightarrow \Phi\Psi\Phi(X) \longrightarrow \cdots
\]
Si $X$ est quelconque, ces flèches sont alternativement dans $W_S^{\mathrm{d}}$ et $W_S^{\mathrm{g}}$, donc leurs composées sont dans $W_S$ sans plus. Mais si $X$ est dans $\und{Fib}_W S$, comme cette catégorie est \textit{stable} par $\Phi$, $\Psi$, et que sur $\und{Fib}_W S$ $W_S^{\mathrm{d}} = W_S^{\mathrm{g}} = W_S^{\mathrm{f}} = W_S^{\mathrm{univ}}$, on voit que dans ce cas les flèches sont toutes dans $W_S^{\mathrm{f}}$. On a envie maintenant de poser
\[
\Phi_\infty(X) = \varinjlim \text{ de tout ça}
\]
\msp{page 193}
et on définirait symétriquement
\[
\Psi_\infty(X) = \varinjlim (X \longrightarrow \Psi(X) \longrightarrow \Phi\Psi(X) \longrightarrow \cdots).
\]

Pour que ce soit utilisable, il faut poser l'axiome \underline{W(6 bis)} (comparer XII, pages 14, 15). La limite inductive d'une suite de morphismes $f_i : X_i \longrightarrow Y_i$, pour un morphisme de systèmes inductifs dans Cat indexés par $\mathbf{N}$, est dans $W$ si les $f_i$ le sont.

Cet axiome est voisin de W(6) de loc. cit., mais n'est impliqué par lui que moyennant W(4b). Il est plus fort dans le cas de systèmes inductifs indexés par $\mathbf{N}$ (donc aussi, plus généralement, de systèmes inductifs sur des ensembles ordonnés filtrants ayant un ensemble cofinal dénombrable), mais ne dit rien sur les systèmes inductifs sur des ensembles ordonnés d'indices plus généraux.

Cet axiome implique aussi qu'une limite inductive filtrante de morphismes $f_i$ est $W$-lisse, resp. $W$-propre, si les $f_i$ le sont. On trouve alors que $\Phi_\infty(X)$ est à la fois $W$-lisse (comme limite inductive des $\Phi(X)$, $\Phi\Psi\Phi(X)$, \ldots, $\Phi(\Psi\Phi)^n(X)$, lesquelles sont $W$-lisses sur $S$) et $W$-propre (comme limite inductive des $\Psi\Phi(X)$, \ldots, $(\Psi\Phi)^n(X)$ \ldots). D'autre part, que $X \longrightarrow \Phi_\infty(X)$ est dans $W_S^{\mathrm{f}}$ - car $W_S^{\mathrm{f}}$ est stable par limites inductives filtrantes dénombrables dans $\mathrm{Cat}/S$. On a donc trouvé un foncteur
\msp{page 194}
\[
\Phi_\infty : \und{Fib}_W S \longrightarrow \und{Parf}_W S \subseteq \und{Fib}_W S
\]
et un morphisme fonctoriel
\[
i_X : X \longrightarrow \Phi_\infty(X), \qquad \boxed{i_X \in W_S^{\mathrm{f}}} \qquad \text{si } X \in \und{Fib}_W S.
\]
% the box around i_X \in W_S^f is printed with rounded corners

Ce foncteur permet alors, en l'absence de W(4b), de faire ce que le foncteur $C$ refusait absolument. Il est vrai que $C$ a l'avantage d'être défini sur $\mathrm{Cat}/S$ tout entier, et d'y être
\printedpage{120}
% head: A. Grothendieck : Dérivateurs, Ch. VI, Édition des Universités de Montpellier II et Paris VII
% foot: Version 9 novembre 2001
raisonnable et même fort utile, quand on a W(4b). Mais $\Phi_\infty$ est, lui aussi, défini sur $\mathrm{Cat}/S$ tout entier, ainsi que $i_X$
\[
\begin{aligned}
&\Phi_\infty : \mathrm{Cat}/S \longrightarrow \und{Parf}_W S\\
&i : \mathrm{id} \longrightarrow \Phi_\infty \qquad \text{i.e.}\ \begin{cases} i_X : X \longrightarrow \Phi_\infty(X)\\ i_X \in W_S, \end{cases}
\end{aligned}
\]
avec le seul handicap (inévitable) que si $X$ n'est pas dans $\und{Fib}_W S$, on ne peut dire mieux que $i_X \in W_S$, à l'exclusion de $i_X \in W_S^{\mathrm{f}}$. Mais pour $C(X)$ c'était forcément pareil. Donc on a l'impression que $\Phi_\infty$, au choix, $\Psi_\infty$, bonnet blanc, blanc bonnet, joue en présence de W(6 bis) le même rôle que C en présence de W(4b). Tout au plus (en présence des deux axiomes à la fois) le foncteur $C$ a-t-il l'avantage esthétique de prendre ses valeurs dans $\und{Parf}_\omega S$, au lieu de $\und{Parf}_W S$ seulement (\footnote[93]{\textit{En fait} $\Phi_\infty = C$, si $C$ est défini via des chemins infinis d'un seul côté simplement. Si on définit $C$ via les chemins \textit{finis} de tous types $c$, alors il prend en effet ses valeurs dans $\und{Parf}_\omega S$.}). Mais il s'impose de % [sic] "que C en présence" (C in roman type) -> $C$, italic everywhere else
\msp{page 195}
se rendre compte à présent que $\Phi_\infty$ a un avantage exactement du même ordre sur $C$, en introduisant
\[
W''_\omega = \text{plus petite partie } \Sigma \text{ de } \mathrm{Fl}(\mathrm{Cat}) \text{ satisfaisant } \mathrm{W}(1, 2\text{ bis}, 3, 6\text{ bis})
\]
et noter que
\[
\Phi_\infty(X) \in \mathrm{Ob}\,\und{Parf}_{W''_\omega} S, \qquad i_x \in W''_\omega % [sic] i_x -> i_X, the morphism is i_X everywhere else
\]
- mais là je m'aperçois qu'on a
\[
W_\omega \subseteq W''_\omega,
\]
donc $C$ est quand même un peu meilleur à cet égard \ldots

Quand on dispose de $\Phi_\infty$ avec les propriétés qu'on vient d'expliciter, on trouve que l'inclusion
\[
\und{Parf}_W S \longrightarrow \und{Fib}_W(S)
\]
induit une équivalence non seulement pour les localisés pour $W_S$, mais aussi pour les localisés fins
\[
(\und{Parf}_W S)(W_S^{\mathrm{f}})^{-1} \longrightarrow (\und{Fib}_W(S))(W_S^{\mathrm{f}})^{-1} \overset{\text{déf}}{=} \mathrm{HOT}^{\mathrm{lc}}_W(S)
\]
(\footnote[94]{On peut y remplacer même $\und{Parf}_W S$ par $\und{Parf}_{W''_\omega} S$, \textit{ou} par toute autre sous-catégorie pleine de $\und{Fib}_W S$ contenant $\und{Parf}_{W''_\omega} S$.}). Mais on ne peut pour autant réduire le nombre des localisés utiles à trois, en l'absence de W(4b). Tout au plus, ayant un foncteur $\Phi_\infty$ qui est plus simple que $C$, en l'absence de W(4b), du fait qu'on sait de toutes façons que
\[
i_X : X \longrightarrow \Phi_\infty(X)
\]
est dans $W_S^{\mathrm{f}}$ si $X$ est dans $\und{Fib}_W S$, peut on
\printedpage{121}
% head: A. Grothendieck : Dérivateurs, Ch. VI, Édition des Universités de Montpellier II et Paris VII
% foot: Version 9 novembre 2001
\msp{page 196}
espérer s'en tirer avec la variante affaiblie de W(4b) obtenue en se bornant dans l'énoncé de W(4b) (p. 171) au cas où $X$, $Y$ sont dans $\und{Parf}_{W''_\omega} S$. Cette fois, ça suffit en effet à impliquer que le foncteur
\[
\lambda : (\und{Fib}_W S)(W_S^{\mathrm{f}})^{-1} \longrightarrow (\und{Fib}_W S)(W_S)^{-1}
\]
est une équivalence de catégories, puisque on peut le ‘calculer’ aussi bien en utilisant $\und{Parf}_{W''_\omega} S$ au lieu de $\und{Fib}_W S$, or l'hypothèse W(4b($\mathrm{W}''_\omega$)) implique que sur $\und{Parf}_{W''_\omega} S$ les deux localiseurs $W_S^{\mathrm{f}}$, $W_S$ sont \textit{égaux.}

Il faut une conclusion pratique utilisable de toutes ces perplexités et observations. Essayons

\textbf{Proposition.} \textit{Supposons que l'on soit sous l'une des conditions suivantes.}

\begin{enumerate}[label={}]
\item[$1^o$)] \textit{L'axiome} W(4b) \textit{est satisfait dans le cas de catégories $X$, $Y$ qui sont dans} $\und{Fibsc}_{0,W} S$ (\footnote[95]{mais c'est \textit{équivalent} à W(4b)! Cf. XVI 4.9. loc. cit.}) \textit{(cette condition sur $W$ est en fait autoduale).}

\item[$2^o$)] \textit{L'axiome} W(4b) \textit{est satisfait dans le cas de catégories $X$, $Y$ qui sont dans} $\und{Parf}_{W''_\omega} S$\textit{, où $W''_\omega$ est le plus petit ensemble $\Sigma \subseteq$ Fl}(Cat) \textit{satisfaisant les conditions} $\mathrm{W}(1, 2\text{ bis}, 3, 6\text{ bis})$. \textit{(On peut aussi se contenter de la condition, plus forte mais plus compréhensible, que} W(4b) \textit{est OK si $X$, $Y$ sont dans} $\und{Parf}_W S$ (\footnote[96]{Mais cela implique aussi W(4b), cf. loc. cit.}).\textit{) De plus, on suppose que $W$ satisfait} W(6 bis).
\end{enumerate}

\textit{Sous ces conditions, le foncteur correspondant [?] $\lambda$}
\[
\mathrm{HOT}^{\mathrm{lc}}_W(S) \longrightarrow \mathrm{HOT}^{\mathrm{lc!}}_W(S)
\]
\textit{est une équivalence de catégories (cf. diagramme p. 178), donc $(\rho, i)$, $(\sigma, j)$ peuvent être regardés comme des couples de foncteurs entre deux catégories (cf. $(*)$ p. 181).} Ce sont des couples de foncteurs adjoints, \textit{donc $i$, $j$ sont pleinement fidèles, et on a la formule d'adjonction $(*)$, page 182.}

\msp{page 197}
Pour prouver l'adjonction de $(\rho, i)$ (ou dualement de $(\sigma, j)$), on exprime encore les deux catégories en présence $\mathrm{HOT}_W(S)$, $\mathrm{HOT}^{\mathrm{lc}}_W(S)$ comme $(\mathrm{Cat}/S)(W_S^{\mathrm{d}})^{-1}$, $(\mathrm{Cat}/S)(W_S)^{-1}$, le foncteur $\rho$ étant la localisation canonique allant du premier vers le second, et le foncteur $i$ étant donné par un foncteur d'inversion
\[
\gamma : (\mathrm{Cat}/S, W_S) \longrightarrow (\mathrm{Cat}/S, \overline{W_S^{\mathrm{d}}})
\]
qui est $C$ dans le cas $1^o$, $\Phi_\infty$ dans le cas $2^o$. Les morphismes d'adjonction
\[
\rho i \longrightarrow \mathrm{id}, \qquad i\rho \longleftarrow \mathrm{id}
\]
doivent se décrire alors comme dans le cas où on suppose tout bonnement W(4b), et les compatibilités se voient de même : on est ramené à prouver que deux flèches
\printedpage{122}
% head: A. Grothendieck : Dérivateurs, Ch. VI, Édition des Universités de Montpellier II et Paris VII
% foot: Version 9 novembre 2001
\[
\begin{tikzcd}
\gamma(X) \arrow[r, shift left, "i_{\gamma(X)}"] \arrow[r, shift right, "\gamma(i_X)"'] & \gamma(\gamma(X))
\end{tikzcd}
\]
sont égales dans $(\mathrm{Cat}/S)(W_S^{\mathrm{d}})^{-1}$. La vérification dans le cas $1^o$), où $\gamma = C$, ne doit pas être changée. Dans le cas $2^o$) où $\gamma = \Phi_\infty$, je ne l'ai pas fait. Je ne le ferai que sous la pression des besoins.

\subsection*{\circnum{6} Loi contravariante de $\mathrm{HOT}_W(S)$, $\mathrm{HOT}^{\mathrm{lc}}_W(S)$.}

Il y a pour
\[
f : S' \longrightarrow S
\]
un homomorphisme du diagramme
\[
\begin{tikzcd}[column sep=small]
& \mathrm{HOT}^{\mathrm{lc}}_W(S') \arrow[dl] \arrow[dr] & & \text{\texttt{[Plutôt}} & & \mathrm{HOT}^{\mathrm{lc}}_W(S) \arrow[dl] \arrow[dr] & \\
\mathrm{HOT}_W(S') & & \mathrm{HOT}^{o}_W(S') & & \mathrm{HOT}_W(S) & & \mathrm{HOT}^{o}_W(S)\ \text{\texttt{]}}
\end{tikzcd}
\]
\msp{page 198}
dans le diagramme similaire pour $S$ \texttt{[plutôt} $S'$\texttt{]}, par des foncteurs $f^*$ qui sont définis respectivement sur
\[
\begin{aligned}
\und{Liss}_W S &\quad \text{pour}\quad \mathrm{HOT}_W(S) & \quad\longrightarrow\quad & \mathrm{HOT}_W(S')\\
\und{Prop}_W S &\quad \text{pour}\quad \mathrm{HOT}^{o}_W(S) & \quad\longrightarrow\quad & \mathrm{HOT}^{o}_W(S')\\
\und{Fib}_W S &\quad \text{pour}\quad \mathrm{HOT}^{\mathrm{lc}}_W(S) & \quad\longrightarrow\quad & \mathrm{HOT}^{\mathrm{lc}}_W(S'),
\end{aligned}
\]
par l'intermédiare des foncteurs % [sic] intermédiare -> intermédiaire, missing i
\[
f^* : X \longmapsto X \times_S S' : \left\{
\begin{aligned}
\und{Liss}_W S &\quad\longrightarrow\quad \und{Liss}_W S'\\
\und{Prop}_W S &\quad\longrightarrow\quad \und{Prop}_W S'\\
\und{Fib}_W S &\quad\longrightarrow\quad \und{Fib}_W S',
\end{aligned}
\right.
\]
lequel foncteur est compatible avec les localiseurs pretinents, qui sont $W_S^{\mathrm{f}}$ dans tous les cas. Par contre, en l'absence d'hypothèses draconiennes sur $W$ ou sur $S$, on ne peut dire que $f^*$ soit compatible avec $W_S$, $W_{S'}$ sur une quelconque des catégories du diagramme p. 175, permettant de définir $\mathrm{HOT}^{\mathrm{lc!}}_W(S)$, de façon à nous donner un foncteur % [sic] pretinents -> pertinents, transposed letters
\[
\mathrm{HOT}^{\mathrm{lc!}}_W(S) \longrightarrow \mathrm{HOT}^{\mathrm{lc!}}_W(S').
\]
En l'absence d'axiomes sur $W$, le seul cas qui me soit connu est celui où $f$ est \textit{une $W$-équivalence universelle,} i.e. une $W$-fibration à fibres $W$-asphériques, auquel cas il est évident que pour \textit{toute} flèche $u : X \longrightarrow Y$ dans $\mathrm{Cat}/S$, $u \in W_S$ si et seulement si $(u' : X' \longrightarrow Y') \in W_{S'}$. On trouve alors un homomorphisme du diagramme complet $(*)$ p. 176 \texttt{[plutôt p. 178]} pour $S$, dans le diagramme correspondant pour $S'$.

Si on suppose W(4b), alors $W_S = W_S^{\mathrm{f}}$ sur $\und{Fib}_W S$, donc sur cette catégorie-là le foncteur $f^*$ est compatible (sans mérite!) avec les localiseurs $W_S$, $W_{S'}$, mais ce n'est pas nouveau vraiment
\printedpage{123}
% head: A. Grothendieck : Dérivateurs, Ch. VI, Édition des Universités de Montpellier II et Paris VII
% foot: Version 9 novembre 2001
\msp{page 199}
par rapport à ce qui a été dit plus haut, pour $\und{Fib}_W S$, $\und{Fib}_W S'$, et pour les $\mathrm{HOT}^{\mathrm{lc}}_W$.

Mais si on suppose W(4a), et $f \in \und{Fib}_W$, alors $f^*$ sur $\mathrm{Cat}/S$ tout entier est compatible avec les localiseurs $W_S$, $W_{S'}$, et on trouve un calcul de $f^*$ sur $\mathrm{HOT}^{\mathrm{lc!}}_W(S) \approx \mathrm{HOT}^{\mathrm{lc}}_W(S)$ par
\[
f^*(\mathrm{hot}^{\mathrm{lc}}_{S,W}(X)) = \mathrm{hot}^{\mathrm{lc}}_{S',W}(X \times_S S').
\]

Donc c'est là un cas, exceptionellement, où les foncteurs $f^*$ pour les $\mathrm{HOT}_W$ sont compatibles non seulement aux foncteurs d'inclusion $i$, $j$ (de la page 176 \texttt{[plutôt p. 178]}), mais aussi aux foncteurs de localisation $\rho$, $\sigma$. % [sic] exceptionellement -> exceptionnellement, French spelling has double n

\textbf{Oubli :} Si $p : S' \longrightarrow S$ est $W$-propre, alors $p^* : \mathrm{Cat}/S \longrightarrow \mathrm{Cat}/S'$ est compatible avec $W_S^{\mathrm{d}}$, $W_{S'}^{\mathrm{d}}$, et $p^* : \mathrm{HOT}_W(S) \longrightarrow \mathrm{HOT}_W(S')$ peut se ‘calculer’ par la formule
\[
p^*(\mathrm{hot}_{S,W}(X)) = \mathrm{hot}_{S',W}(X \times_S S')
\]
pour \textit{tout} $X$ dans $\mathrm{Cat}/S$.

\subsection*{\circnum{7} Produit dans $\mathrm{HOT}_W(S)$.}

Le foncteur
\[
(X, Y) \longmapsto X \times_S Y : \und{Liss}_W S \times \und{Liss}_W S \longrightarrow \und{Liss}_W S
\]
est compatible avec les localiseurs $W_S^{\mathrm{f}} \times W_S^{\mathrm{f}}$ et $W_S^{\mathrm{f}}$, donc passe aux localisés en
\[
\left\{
\begin{aligned}
\mathrm{HOT}_W(S) \times \mathrm{HOT}_W(S) &\quad\longrightarrow\quad \mathrm{HOT}_W(S)\\
(\xi, \eta) &\quad\longmapsto\quad \xi \times_S \eta
\end{aligned}
\right.
\]
(\footnote[97]{\textbf{NB} Si $X$, $Y$ sont \textit{quelconques} dans $\mathrm{Cat}/S$, on a $\mathrm{hot}_{S,W}(X) \times_S \mathrm{hot}_{S,W}(Y) = \mathrm{hot}_{S,W}(\Phi_S(X) \times_S \Phi_S(Y))$. Si $X$ ou $Y$ est lisse sur $S$, on peut y remplacer $\Phi(X)$ resp. $\Phi(Y)$ par $X$ resp. $Y$.}). Ce produit est associatif, unitaire, commute avec toutes les compatibilités habituelles, comme il résulte du fait que c'est comme ça dans $\und{Liss}_S$. L'élément unité est la classe de $S \xrightarrow{\mathrm{id}} S$. La sous-catégorie $\mathrm{HOT}^{\mathrm{lc}}_W(S)$ est stable par le produit, qui induit
\msp{page 200}
\[
\mathrm{HOT}^{\mathrm{lc}}_W(S) \times \mathrm{HOT}^{\mathrm{lc}}_W(S) \longrightarrow \mathrm{HOT}^{\mathrm{lc}}_W(S),
\]
lui-même induit par le foncteur produit $(X, Y) \longrightarrow X \times_S Y$ dans $\und{Fib}_W S$.

Si $X$ est quelconque dans $\mathrm{Cat}/S$, le foncteur
\[
\mathrm{mult}_{X/S} : Y \longmapsto X \times_S Y : \und{Liss}_W S \longrightarrow \mathrm{Cat}/S
\]
est compatible avec les localiseurs $W_S^{\mathrm{d}}$, car ce foncteur peut s'écrire comme le composé du foncteur correspondent \textit{à valeurs dans $\und{Liss}_W X$,} lequel est compatible avec les localiseurs % [sic] correspondent -> correspondant, adjective intended


\endgroup
% END UNIT Der-ChVI-345/Der-ChVI-345_p121-125.tex
% BEGIN UNIT Der-ChVI-345/Der-ChVI-345_p126-129.tex
\begingroup\makeatletter%
\newcommand{\printedpage}[1]{\marginpar{\footnotesize p.~#1}}%
\newcommand{\msp}[1]{\marginpar{\footnotesize [#1]}}%
\newcommand{\cercle}[1]{\tikz[baseline=(c.base)]\node[draw,circle,inner sep=1pt](c){#1};}%
\makeatother\let\printedpage\derpp\let\msp\dermsp%
\let\mspin\DERmspin%


\setcounter{footnote}{97}

\printedpage{124}
% head: A. Grothendieck : Dérivateurs, Ch. VI, Édition des Universités de Montpellier II et Paris VII
% foot: Version 9 novembre 2001

$W_S^{\mathrm{d}}$, $W_S^{\mathrm{d}}$, et le foncteur ‘oubli de $X$’ de $\underline{\mathrm{Liss}}_W X$ (même de $\mathrm{Cat}/X$) dans $\mathrm{Cat}/S$, lequel est compatible avec $W_X^{\mathrm{d}}$ et $W_S^{\mathrm{d}}$. On trouve ainsi un foncteur % [sic] W_S^d, W_S^d -> apparently W_S^d, W_X^d (or W_X^d, W_S^d), the same symbol is printed twice
\[
\left\{\begin{aligned}
&\mathrm{mult}_{X/S} : \mathrm{HOT}_W(S) \longrightarrow \mathrm{HOT}_W(S)\\
&\mathrm{mult}_{X/S}(\mathrm{hot}_{S,W}(Y)) \overset{\text{déf}}{=} \mathrm{hot}_{S,W}(X\times_S Y) \qquad \text{si $Y$ \textit{lisse} sur $S$}
\end{aligned}\right.
\]
(\footnote{Cette formule est également valable si $Y$ est dans $\underline{\mathrm{Fib}}_W\, S$, cf. $(*)$, p. 202.}), puisque $\mathrm{HOT}_W(S)$ peut s'obtenir comme localisé tant de $\underline{\mathrm{Liss}}_W S$ que de $\mathrm{Cat}/S$, pour $W_S^{\mathrm{d}}$. Quand $X$ est lui-même $W$-lisse sur $S$, ce foncteur n'est autre que
\[
\left\{\begin{aligned}
&\mathrm{mult}_\xi : \eta \longmapsto \xi\times_S \eta,\\
&\text{où }\ \xi = \mathrm{hot}_{S,W}(X).
\end{aligned}\right.
\]

Mais il n'en est pas de même si $X$ n'est pas lisse, i.e. il faut distinguer dans ce cas, même pour $Y$ $W$-lisse sur $S$,
\[
\mathrm{hot}_{S,W}(X\times_S Y) \qquad\text{et}\qquad \mathrm{hot}_{S,W}(X)\times_S \mathrm{hot}_{S,W}(Y).
\]

Le deuxième terme peut s'expliciter comme $\mathrm{hot}_{S,W}(\Phi(X)\times_S Y)$ qui reçoit $\mathrm{hot}_{S,W}(X\times_S Y)$, i.e.
\msp{page 201}
on a un homomorphisme canonique
\[
\underbrace{\mathrm{hot}_{S,W}(X\times_S Y)}_{=\,\mathrm{mult}_{X/S}(Y)} \longrightarrow \underbrace{\mathrm{hot}_{S,W}(X)\times_S \mathrm{hot}_{S,W}(Y)}_{=\,\mathrm{hot}_{S,W}(\Phi(X)\times_S Y)},
\]
induit par
\[
X\times_S Y \longrightarrow \Phi(X)\times_S Y,
\]
mais cette dernière flèche n'a aucune raison d'être dans $W_S^{\mathrm{d}}$, et n'induit donc pas nécessairement un isomorphisme dans $\mathrm{HOT}_W(S)$. (Contrexemple page 154.)

Il faut regarder à part le cas \textit{où $X$ est dans} $\underline{\mathrm{Fib}}_W S$, j'ai envie de dire que dans ce cas on a
\[
\mathrm{hot}_{S,W}(X\times_S Y) \xrightarrow{\ \sim\ } \underbrace{\mathrm{hot}_{S,W}(X)\times_S \mathrm{hot}_{S,W}(Y)}_{=\,\mathrm{hot}_{S,W}(\Phi(X)\times_S\, \Phi(Y))} \tag*{$(*)$} % doubt resolved in DER_VI-c_corrections.csv; input note: under the brace "hot" is printed larger than "(Phi(X)x_S Phi(Y))", which stands at subscript level -- argument of hot_{S,W} or part of its subscript
\]
(\footnote{Cette formule n'est établie qu'en supposant ou bien $Y$ $W$-lisse (en plus de $X$ $W$-fibré sur $S$), ou bien $X$ $W$-propre, p.ex. $X$ $W$-parfait.}) pour \textit{tout} $Y$ dans $\mathrm{Cat}/S$ (pas nécessairement $W$-lisse), a fortiori que dans ce cas
\[
\left\{\begin{aligned}
&\mathrm{mult}_{X/S} = \mathrm{mult}_\xi \qquad \text{dans } \mathrm{HOT}_W(S)\\
&\text{où } \xi = \mathrm{hot}_{S/W}(X/S) \text{ (si $X$ dans $\underline{\mathrm{Fib}}_W S$).} % [sic] hot_{S/W} -> hot_{S,W}, the notation used everywhere else
\end{aligned}\right. \tag*{$(**)$}
\]

\printedpage{125}
% head: A. Grothendieck : Dérivateurs, Ch. VI, Édition des Universités de Montpellier II et Paris VII
% foot: Version 9 novembre 2001

Dans le cas où $Y$ est lisse sur $S$ (ce qui suffit pour établir $(**)$), il suffit de se rappeler que $X\longrightarrow \Phi(X)$ est dans $W_S^{\mathrm{f}}$ (résulte du fait que $X$ est $W$-fibré sur $S$ (\footnote{il suffit même $W$-prélisse sur $S$}), donc $X_s \longrightarrow s\backslash X$ est dans $W$, $\forall\ s\in S$), donc $\underbrace{X\times_S Y}_{X_Y} \longrightarrow \underbrace{\Phi(X)\times_S Y}_{\Phi(X)_Y}$ est aussi dans $W_Y^{\mathrm{f}}$, donc (comme $X_Y$, $\Phi(X)_Y$ sont $W$-fibrés sur $Y$) dans $W_Y^{\mathrm{d}}$, donc aussi dans $W_S^{\mathrm{d}}$. En fait, même sans supposer $Y$ lisse, ceci montre que l'on a
\msp{page 202}
\[
\mathrm{hot}_{S,W}(X\times_S Y) \xrightarrow{\ \sim\ } \mathrm{mult}_{Y/S}(\mathrm{hot}_{S,W}(X)) \qquad \text{si $X$ est dans $\underline{\mathrm{Fib}}_W S$.} \tag*{$(*)$}
\]

Pour avoir $(*)$ pour $Y$ quelconque, on peut par ce qui précède supposer $X$ lisse (quitte à le remplacer par $\Phi(X)$), donc $X$ dans $\underline{\mathrm{Liss}}_{0,W}S$. Il faut voir que
\[
X\times_S Y \longrightarrow X\times_S \Phi(Y)
\]
est encore dans $W_S^{\mathrm{d}}$. Or si $X$ est \textit{$W$-propre} sur $S$, le changement de base $X\longrightarrow S$ est compatible avec $W_S^{\mathrm{d}}$, $W_X^{\mathrm{d}}$ (sur $\mathrm{Cat}/S$ tout entier), et on gagne. (Ça doit être faux sans hypothèse de $W$-propreté sur $X$, \textit{ou} de $W$-lissité sur $Y$ \dots)

Quand on a
\[
f : S' \longrightarrow S,
\]
alors
\[
f^* : \mathrm{HOT}_W(S) \longrightarrow \mathrm{HOT}_W(S')
\]
est compatible avec les structures de produit.

\subsection*{\protect\cercle{8} Les foncteurs $f_!$ : théorème d'adjonction, propriétés du dérivateur $\mathrm{HOT}_W$.}

Soit
\[
f : S' \longrightarrow S
\]
dans Cat, alors le foncteur induit
\[
f_! : \mathrm{Cat}/S' \longrightarrow \mathrm{Cat}/S
\]
est compatible avec les localiseurs $(W_{S'}^{\mathrm{d}}, W_S^{\mathrm{d}})$, ou $(W_{S'}^{\mathrm{g}}, W_S^{\mathrm{g}})$, ou $(W_{S'}, W_S)$. On trouve donc un foncteur induit pour les localisées, et \texttt{[un]} diagramme commutatif
\msp{page 203}
\[
\begin{tikzcd}[column sep=small]
\mathrm{HOT}_W(S') \arrow[rrd, "\rho'"] \arrow[dd, "f_!"'] & & & & \mathrm{HOT}^{o}_W(S') \arrow[lld, "\sigma'"'] \arrow[dd, "f^{o}_!"] \\
& & \mathrm{HOT}^{\mathrm{lc}\,!}_W(S') \arrow[dd, "f^{\mathrm{lc}\,!}_!"] & & \\
\mathrm{HOT}_W(S) \arrow[rrd, "\rho"] & & & & \mathrm{HOT}^{o}_W(S). \arrow[lld, "\sigma"'] \\
& & \mathrm{HOT}^{\mathrm{lc}\,!}_W(S) & &
\end{tikzcd}
\]

\printedpage{126}
% head: A. Grothendieck : Dérivateurs, Ch. VI, Édition des Universités de Montpellier II et Paris VII
% foot: Version 9 novembre 2001

On n'a pas de foncteur $f_!^{\mathrm{lc}}$ au niveau des $\mathrm{HOT}_W^{\mathrm{lc}}$, sauf si on fait une hypothèse, telle \underline{W(4b)}, qui assure que $\mathrm{HOT}_W^{\mathrm{lc}\,!} \approx \mathrm{HOT}_W^{\mathrm{lc}}$.

Le théorème principal, que le foncteur $f_!$ est adjoint à gauche de $f^*$, est valable \textit{sous la seule hypothèse} W$(1,2,3)$ (\footnote{i.e. avec ces nouvelles notations Loc$(1,2,3\ \text{bis})$ (cf. XVI).}). Donc dualement pour $f_!^{o}$, $f_{o}^*$. Mais comme on n'a pas de foncteur $f_{\mathrm{lc}\,!}^*$, la question d'adjonction ne se pose pas pour $f$ général, et sans hypothèse W(4b) sur $W$. Elle ne se pose, apparemment, que pour $f$ un $W$-fibré trivial (cf. fin de \cercle{6}, p. 198), et dans ce cas l'adjonction ne doit pas faire de difficultés. Mais tant que je n'ai pas de sentiment bien net au sujet de la signification des $\mathrm{HOT}_W^{\mathrm{lc}\,!}$, ce n'est pas là un résultat bien bandant \texttt{[sic]}.

L'adjonction de $f_!$, $f^*$ pour $\mathrm{HOT}_W$ permet déjà de montrer que le prédérivateur $\mathrm{HOT}_W$
\msp{page 204}
satisfait les conditions Der 1, Der $4'$, Der $5'$ des dérivateurs. De plus, on a ce qu'on peut appeler la propriété Der 3b pour $\mathrm{HOT}_W$ en tant que $W$-dérivateur, i.e.

\begin{quote}
\textit{Si $f\in W$, alors $f$ est une} $\mathrm{HOT}_W$\textit{-équivalence} (\footnote{On le prouve en interprétant la notion de $\mathrm{HOT}_W$-équivalence par la $\mathrm{HOT}_W$\textit{-homologie.} Cela revient finalement à ceci : Si $f : S' \longrightarrow S$ est dans $W$, alors $\forall\, T$, $S'\times T \longrightarrow S\times T$ aussi!}) (moyennant les seuls axiomes W$(1,2,3)$ sur $W$). Plus précisément, $f$ est une $\mathrm{HOT}_W$-équivalence si et seulement si $\mathrm{hot}_W(f)$ est \texttt{[un]} isomorphisme, i.e. $f\in \overline{W}$, le saturé fort de $W$.
\end{quote}

Si $W$ \textit{satisfait} W(4b)\textit{,} alors on a également un couple de foncteurs adjoints
\[
\begin{tikzcd}
\mathrm{HOT}_W^{\mathrm{lc}}(S') \arrow[r, shift left, "f_!^{\mathrm{lc}}"] & \mathrm{HOT}_W^{\mathrm{lc}}(S) \arrow[l, shift left, "f^*_{\mathrm{lc}}"]
\end{tikzcd}
\]
pout toute flèche $S'\longrightarrow S$ dans Cat. On voit aussi dans ce cas que % [sic] pout -> pour, misprint

\begin{quote}
% printed with a vertical rule at the left of the block
\begin{tabular}{|p{0.9\linewidth}}
si $f : S' \longrightarrow S$ est dans $W$, alors $f^*_{\mathrm{lc}}$ et $f_!^{\mathrm{lc}}$ entre $\mathrm{HOT}_W^{\mathrm{lc}}(S)$ et $\mathrm{HOT}_W^{\mathrm{lc}}(S')$ sont des équivalences quasi-inverses l'une de l'autre (moyennant W(4b)),
\end{tabular}
\end{quote}

ce qui est l'axiome Der 3b$_2$ ‘version $W$-dérivateurs’. On ne peut s'attendre à mieux, du côté Der 3!

\textbf{Oubli :} Soit $f : S' \longrightarrow S$, on a alors $\boxed{f_!f^*} : \mathrm{HOT}_W(S) \longrightarrow \mathrm{HOT}_W(S)$ est égal à $\mathrm{mult}_{S'/S}$ (cf. p. 200). Il est donc égal à $\mathrm{mult}_\xi$, où $\xi = \mathrm{hot}_{S,W}(S')$, si on suppose ou bien $S'$ $W$-lisse sur $S$, ou bien $S'$ $W$-fibré sur $S$ (cf. $(**)$, p. 201).

\subsection*{\protect\cercle{9} Théorème de changement de base pour $f_!$ dans $\mathrm{HOT}_W$.}

Si on a
\printedpage{127}
% head: A. Grothendieck : Dérivateurs, Ch. VI, Édition des Universités de Montpellier II et Paris VII
% foot: Version 9 novembre 2001
\[
\begin{tikzcd}
X \arrow[d, "f"'] & X' \arrow[l, "q"'] \arrow[d, "f'"] \\
S & S', \arrow[l, "g"']
\end{tikzcd}
\]
avec $f$ $W$-lisse \textit{ou} $g$ $W$-propre \texttt{[la suite manque]}


\endgroup
% END UNIT Der-ChVI-345/Der-ChVI-345_p126-129.tex
\clearpage\setcounter{footnote}{0}\phantomsection\addcontentsline{toc}{section}{Chapitre VII}\markboth{Chapitre VII}{Chapitre VII}
% BEGIN UNIT Der-ChVII-456/Der-ChVII-456_p001-005.tex
\begingroup\makeatletter%
\newcommand{\printedpage}[1]{\marginpar{\footnotesize p.~#1}}%
\newcommand{\msp}[1]{\marginpar{\footnotesize [#1]}}%
\newcommand{\Fltr}{\operatorname{Fl}\operatorname{tr}}%
\newcommand{\Ob}{\operatorname{Ob}}%
\newcommand{\bDelta}{\boldsymbol{\Delta}}%
\makeatother\let\printedpage\derpp\let\msp\dermsp%
\let\mspin\DERmspin%


% ===================== page 1 =====================
\printedpage{}
% head: (none; title page, no page number printed)

\begin{center}
{\Large\textsf{LES DÉRIVATEURS}}

\bigskip

\textbf{ALEXANDRE GROTHENDIECK}

\bigskip\bigskip

\textbf{Chapitre VII}

\medskip

\textbf{Catégories de chemins (2)}
\end{center}

\par

% ===================== page 2 =====================
\printedpage{}
% head: (none; no page number printed)

\par\noindent\textbf{Note des éditeurs de l’édition transcrite.}\par
Ce texte a été déchiffré et transcrit en \LaTeXe\ par M. Künzer. Il a été édité par M. Künzer, J. Malgoire, et G. Maltsiniotis. La transcription est aussi fidèle que possible au manuscrit. Pour les quelques corrections évidentes, ou rares commentaires des éditeurs, ainsi que pour la numérotation originale des pages du manuscrit, les caractères de machine à écrire \texttt{[typewriter]} entre crochets sont utilisés. Un point d'interrogation entre crochets signifie que l'on n'est pas sûr du mot qui précède.

Cette édition est provisoire. Les remarques, commentaires et corrections sont bienvenus. Envoyer un message à :

\textbf{[e-mail address removed]}

\hfill G. Maltsiniotis\hspace*{3em}

\par

% ===================== page 3 =====================
\printedpage{1}
% head: A. Grothendieck : Dérivateurs, Ch. VII, Édition des Universités de Montpellier II et Paris VII
% foot: Version 22 décembre 2001

\msp{page 1}
Comme le formalisme ‘simplifié’ des chemins que j'ai voulu dégager à la va-vite s'est révélé d'une lourdeur insupportable, de nature à le rendre quasiment inutilisable, je vais essayer de dégager un formalisme souple, collant à l'intuition, au prix seulement d'un effort conceptuel supplémentaire.

\section*{1\quad Intervalles et leur variantes} % [sic] leur variantes -> leurs variantes, plural agreement

On appelle \textit{catégorie intervalle standard} une catégorie définie par un ensemble ordonné, dont l'ensemble des élements est un intervalle d'entiers $[0,n]$ ($n\in\mathbf{N}$, on n'exclut pas le cas $n=0$, i.e. de la catégorie ponctuelle $\{0\}$) : % [sic] élements -> éléments, missing accent
\[
\Ob I=[0,n]\subseteq\mathbf{N},
\]
et satisfaisant l'hypothèse suivante :

\begin{enumerate}
\item[a)] pour toute paire de deux sommets consécutifs $i-1$, $i$ de $I$ ($1\leq i\leq n$), il y a une flèche qui les relie, donc soit $i\longrightarrow i+1$, soit $i\longleftarrow i+1$ \texttt{[plutôt} $i-1\longrightarrow i$\texttt{,} $i-1\longleftarrow i$\texttt{]} (cette flèche est notée $\underline{u}_i$,
\[
(i-1)\xrightarrow{\underline{u}_i}i\qquad\quad\text{ou}\qquad\quad(i-1)\xleftarrow{\underline{u}_i}i\ \ ).
\]
\item[b)] La catégorie $I$ est engendrée par l'ensemble de ses flèches $\underline{u}_i$ (appelées \textit{flèches de transition}).
\end{enumerate}

La flèche de transition est dite \textit{directe} si elle va de $i-1$ à $i$, \textit{rétrograde} dans le sens inverse. Si on affecte l'intervalle $[i-1,i]$ du signe $+$ ou $-$ suivant que $\underline{u}_i$ est directe ou rétrograde, on voit donc que les
\msp{page 2}
intervalles standard de longueur donnée $n$ correspondent biunivoquement aux systèmes de signes $+$ ou $-$ sur les $n$ intervalles de transition $[i-1,i]$ ($1\leq i\leq n$), il y en a donc $2^n$. Le sommet $0$ est l'\textit{origine} ou la \textit{source} de l'intervalle, le sommet $n$ l'\textit{extrémité} ou le \textit{but} de l'intervalle.

Un intervalle standard est \textit{orienté,} du fait qu'on à la relation d'ordre \textit{total} naturelle sur l'ensemble de ses sommets, savoir celle des entiers dans $[0,n]$. Quand on parlera de la relation d'ordre sur $I$, qu'un sommet $i$ en \textit{‘précède’} ou \textit{‘suit’} un autre $j$ ($i\leq j$ ou $i\geq j$ respectivement), c'est de cette relation d'ordre qu'on voudra parler. Celle associée à la structure de catégorie ne sera jamais regardée comme une relation d'ordre. Il y a également une relation d'ordre total naturelle sur l'ensemble des flèches de transition $\underline{u}_i$ ($1\leq i\leq n$), donnée par la relation d'ordre des entiers $i$, $1\leq i\leq n$, qui les indexent. On dira qu'une flèche $i$ \texttt{[plutôt} $\underline{u}_i$\texttt{]} de transition \textit{précède} ou \textit{suit} une autre \texttt{[}$\underline{u}_j$\texttt{]}, ou qu'elle est \textit{successeur immédiat} ou \textit{prédécesseur immédiat} d'une autre (si $i=j+1$ ou $i=j-1$). % [sic] qu'on à -> qu'on a, verb; [sic] un sommet i en ‘précède’ -> un sommet i ‘précède’, superfluous "en"

% Reader layout: continuous text; edition coordinate retained below.
% ===================== page 4 =====================
\printedpage{2}
% head: A. Grothendieck : Dérivateurs, Ch. VII, Édition des Universités de Montpellier II et Paris VII
% foot: Version 22 décembre 2001

Si on a un ensemble d'entiers
\[
E\subseteq\mathbf{Z},
\]
on considérera dans $E$ la relation d'équivalence engendrée par la relation entre deux éléments $(i,j)$ de $E$ : $j$ est successeur immédiat de $i$ \textit{dans} $\mathbf{Z}$ (i.e. $j=i+1$). Deux éléments distinct $i$, $j$ de $E$ sont congrus pour cette relation si et seulement si l'intervalle d'entiers qu'ils déterminent est $\subseteq E$, donc % [sic] distinct -> distincts, plural agreement
\[
\text{si } i<j,\ i \text{ et } j \text{ sont \textit{‘connectés’} dans } E \text{ si et seulement si } [i,j]\subseteq E.
\]

\msp{page 3}
Quand on parlera des ‘composantes connexes’ de $E$, il s'agira toujours des classes suivant cette relation d'équivalence.

Nous identifierons l'ensemble des flèches de transition $\Fltr(I)$ de $I$ à l'intervalle d'entiers $[1,n]$, donc à une partie de $\mathbf{Z}$. Si on a une partie $E\subseteq\Fltr(I)$, on y a donc la notion de flèches de transition ‘connectées’ dans $E$, et la notion de connexité pour $E$. Si $i<j$, $\underline{u}_i$ et $\underline{u}_j$ sont connectées si et seulement si toutes les flèches de transition intermédiaires entre $\underline{u}_i$ et $\underline{u}_j$ sont dans $E$. Et $E$ est connexe si et seulement s'il est soit vide, soit formée d'une suite de flèches de transition dont chacune est successeur immédiat de celle qui la précède, % [sic] soit formée -> soit formé, agreement with "il"
\[
E=\{\underline{u}_i,\underline{u}_{i+1}\ldots,\underline{u}_{i+l-1}\}, % [sic] \underline{u}_{i+1}\ldots -> \underline{u}_{i+1},\ldots, missing comma
\]
où $l$ est le cardinal de $E$. Une telle $E$ est appelée une \textit{chaîne} de flèches de transition, et $l$ est sa longueur. Une chaîne est dite \textit{composable} si le composé $\underline{u}_{i+l-1}\circ\cdots\circ\underline{u}_{i+1}\circ\underline{u}_i$ \texttt{[ou} $\underline{u}_i\circ\underline{u}_{i+1}\circ\ldots\circ\underline{u}_{i+l-1}$\texttt{]} a un sens, i.e. si elles sont toutes de même sens, soit directe, soit rétrograde.

On désigne par $\Fltr^+(I)$, $\Fltr^-(I)$ l'ensemble des flèches de transition directes, ou rétrogrades, respectivement. Leurs composantes connexes ne sont autres que les chaînes composables maximales. Ainsi, les chaînes composables maximales forment une partition de l'ensemble des flèches de transition, en intervalles d'entiers mutuellement disjoints. Les chaînes forment elles-mêmes un ensemble totalement ordonné
\[
C_1,\ C_2,\ \ldots\ C_r\qquad\quad(1\leq r\leq n).
\]
L'entier $r$ s'appelle la \textit{‘longueur utile’} de l'intervalle
\msp{page 4}
$I$, c'est donc le cardinal de l'ensemble de ses chaînes composables maximales.

Deux chaînes composables maximales \texttt{[consécutives]} sont toujours de signe opposé, si l'une est directe l'autre est rétrograde, et inversement. La première chaîne composable maximale $C_1$ (disons simplement \textit{‘chaîne structurale’}) a le même signe que $\underline{u}_1$; \texttt{[on]} dit que la \textit{direction initiale} de $I$ est directe, resp. rétrograde, s'il en est ainsi de $\underline{u}_1$, ou ce qui revient au même, de $C_1$. De même on définit la \textit{direction finale} de $I$ comme celle de $\underline{u}_n$, ou encore celle de $C_r$. Direction initiale et finale sont égales si la longueur utile $r$ est
\printedpage{3}
% head: A. Grothendieck : Dérivateurs, Ch. VII, Édition des Universités de Montpellier II et Paris VII
% foot: Version 22 décembre 2001
impaire, inégales si elle est paire. (\textbf{NB} Ce n'est que pour $n=0$, i.e. quand $I=\{0\}$ est la catégorie ponctuelle, que direction initiale et direction finale ne sont pas définies.)

Les chaînes composables correspondent aux sous-catégories pleines de $I$ qui sont isomorphes à une catégorie $\bDelta_i$ (où $i$ est la longueur de la chaîne). Notons que $\bDelta_i$ n'a pas d'autre automorphisme que l'identité, donc les chaînes maximales de longueur donnée $i\geq 1$ correspondent exactement aux foncteurs
\[
\bDelta_i\lhook\joinrel\longrightarrow\ I
\]
qui sont injectifs sur les sommets. Un tel foncteur est pleinement fidèle, et un monomorphisme. Ce foncteur définit une application strictement croissante ou strictement décroissante sur
\msp{page 5}
les sommets, suivant que la chaîne est directe ou rétrograde.

Notons que tout simplexe type $\bDelta_i$, $i\geq 0$, peut être vu comme un intervalle, de longueur $i$, de longueur utile $1$ (si $i\geq 1$, $0$ si $i=0$). Son origine est $0$, son extrémité est $i$. Mais appelons une catégorie quelconque une \textit{catégorie intervalle} si elle est isomorphe à un intervalle standard. Notons ici

\textbf{Proposition 1.} \textit{Soit $I$ un intervalle standard de longueur $n\geq 1$.}

\begin{enumerate}
\item[\textit{a)}] \textit{Soit $i\in I$. Pour que $i\in\{0,n\}$, i.e. pour que $i$ soit un des deux} ‘bouts’ \textit{(ou} ‘extrémités’\textit{) (source ou but) de $I$, il faut et il suffit que l'une des deux catégories $I/i$ ou $i\backslash I$ soit la catégorie ponctuelle et que l'autre soit isomorphe à un} $\bDelta_j$ (\footnote{où $j$ est la longueur de l'unique chaîne structurale qui ‘touche’ $i$ (i.e. telle que $i$ soit dans la sous-catégorie pleine définie par ladite chaîne).}). \textit{(Caractérisation intrinsèque des bouts d'un intervalle.) Pour que $i$ soit la source, il faut et il suffit que $i\backslash I$} \texttt{[phrase incomplète]}
\item[\textit{b)}] \textit{Tout foncteur $I\longrightarrow I$ qui fixe les deux extrémités de $I$ est l'identité.}
\end{enumerate}

\textbf{Corollaire.} \textit{Le groupe des automorphismes de $I$ a au plus deux éléments. S'il en a deux, l'automorphisme non trivial échange les deux extrémités. (On dit alors que l'intervalle est} symétrique.\textit{)}

\msp{page 6}
NB Si $i$ n'est pas un bout, et si néanmoins l'une des deux catégories $i\backslash I$, $I/i$ est ponctuelle, disons $I/i$, alors au voisinage immédiat de $i$, $I$ a la forme
\[
\begin{tikzcd}[row sep=0pt]
\bullet & \bullet \arrow[l, "\underline{u}_i"'] \arrow[r, "\underline{u}_{i+1}"] & \bullet \\
i-1 & i & i+1
\end{tikzcd}
\]
($\underline{u}_i$ rétrograde, $\underline{u}_{i+1}$ directe - si c'était $i\backslash I$ la ponctuelle, alors $\underline{u}_i$ serait directe, $\underline{u}_{i+1}$ rétrograde - en tous les cas ces deux flèches sont de sens opposé, i.e. non composables). On trouve donc que $i\backslash I$ est une catégorie de la forme


\endgroup
% END UNIT Der-ChVII-456/Der-ChVII-456_p001-005.tex
% BEGIN UNIT Der-ChVII-456/Der-ChVII-456_p006-010.tex
\begingroup\makeatletter%
\newcommand{\printedpage}[1]{\marginpar{\footnotesize p.~#1}}%
\newcommand{\msp}[1]{\marginpar{\footnotesize [#1]}}%
\newcommand{\bD}{\boldsymbol{\Delta}}%
\newcommand{\iddots}{\mathinner{\mkern1mu\raise1pt\hbox{.}\mkern2mu\raise4pt\hbox{.}\mkern2mu\raise7pt\hbox{.}\mkern1mu}}%
\makeatother\let\printedpage\derpp\let\msp\dermsp%
\let\mspin\DERmspin%


\printedpage{4}
% head: A. Grothendieck : Dérivateurs, Ch. VII, Édition des Universités de Montpellier II et Paris VII
% foot: Version 22 décembre 2001

\[
\begin{tikzcd}[column sep=small, row sep=small]
 & & & i \arrow[dl] \arrow[dr] & & & \\
 & & i-1 & & i+1 & & \\
 & \iddots \arrow[dl] & & & & \ddots \arrow[dr] & \\
 i-r & & & & & & i+r, % [sic] i+r -> i+s, s est la longueur de la chaîne contenant u_{i+1}
\end{tikzcd}
\]

où $r$ est la longueur de la chaîne structurale contenant $\underline{u}_i$, $s$ celle de la chaîne structurale contenant $\underline{u}_{i+1}$. Cette catégorie est une somme amalgamée $\bD_r \amalg \bD_s$ (recollées suivant leur origine), mais n'est pas isomorphe à une $\bD_j$, p.~ex. faute d'avoir un objet final.

Cette observation prouve a). Pour prouver b), il suffit de prouver que l'ordre total sur $I$ est déterminé de façon intrinsèque, par le choix de \textit{l'origine} $s_0$ de $I$ (parmi ces % [sic] ces -> ses, possessif
deux extrémités). Supposons que l'on ait déjà déterminé intrinsèquement les sommets $s_0$, $s_1$, $\ldots$, $s_i$ ($i \geq 0$). Comment déterminer $s_{i+1}$? (Quand $i < n$.) Soit $I'$ la sous-catégorie pleine de $I$
\msp{page 7}
formée des $j \in I$ qui ne sont pas parmi les $s_j$ \texttt{[plutôt} $s_k$\texttt{]} déjà construits, sauf le dernier $s_i$ (donc $s_i \in I'$, $\operatorname{card} I' \geq 2$). Alors $I'$ est isomorphe à un intervalle standard, avec $s_i$ correspondant à l'origine. Une des catégories $i\backslash I$, $I/i$ est ponctuelle, l'autre isomorphe à une catégorie $\bD_j$, et de plus, $i$ est objet initial ou final de $\bD_j$. Si c'est l'objet initial $0$, on prend pour $s_{i+1}$ le successeur immédiat de $i$ dans $\bD_j$. Si c'est l'objet final $j$, on prend pour $s_{i+1}$ le prédécesseur immédiat $j-1$.

Si on était parti du but $n$ de $I$, au lieu de partir de la source, on aurait trouvé la relation d'ordre total opposée. Affectant l'élément but du numéro d'ordre $0$, on trouve une renumération de $I$, $i$ devenant $n-i$, et les flèches de transition qui étaient directes pour la relation d'ordre initiale, deviennent rétrogrades dans la relation d'ordre opposée, et inversement. On voit ainsi :

\textbf{Proposition 2.} \textit{Si $I$ est un intervalle standard, il existe un unique couple $(I',\theta)$, où $I'$ est un intervalle standard, et $\theta$ un isomorphisme $I \xrightarrow{\sim} I'$ qui transforme le but de $I$ dans la source de $I'$ (et la source de $I$ en le but de $I'$).}

Cet intervalle $I'$ est appelé l'intervalle \textit{opposé} de $I$. Il faut le voir, intrinsèquement, comme la `même' categorie $I$, % [sic] categorie -> catégorie, accent manquant
mais avec une `orientation' opposée.

\msp{page 8}

Une façon plus claire de voir ceci est la suivante. Soit $J$ une catégorie intervalle, i.e. isomorphe à un intervalle standard. Alors il y a exactement (si $I \not\simeq e$) \textit{deux} couples $(I,\theta)$ d'un intervalle standard $I$ et d'un isomorphisme $\theta : I \xrightarrow{\sim} J$. Ces deux isomorphismes induisent sur $J$ deux ordres totaux, opposés l'un de l'autre. Chaque couple est déterminé de façon unique par ledit ordre total, qu'on appelera % [sic] appelera -> appellera, orthographe
une \textit{orientation} de l'intervalle $J$. Ainsi $J$ a deux orientations, et chacune détermine un \textit{isomorphisme orienté} (i.e. respectant
\printedpage{5}
% head: A. Grothendieck : Dérivateurs, Ch. VII, Édition des Universités de Montpellier II et Paris VII
% foot: Version 22 décembre 2001
les orientations) de $J$ avec un intervalle standard (\footnote[2]{\textbf{NB} On parlera de l'intervalle \textit{orienté opposé} $I^{\mathrm{opp}}$ d'un intervalle \textit{orienté} $I$, qu'il faut se garder de confondre avec la catégorie opposée $I^o$. Ici $I^{\mathrm{opp}}$ a même catégorie sous-jacente que $I$, mais avec une orientation opposée.}). Ces deux intervalles sont les mêmes si et seulement s'il existe un automorphisme de $J$ (nécessairement unique) qui échange ses deux orientations, ou ce qui revient au même, qui échange entre elles les deux extrémités. (Ainsi les catégories opposées l'une de l'autre
\[
\begin{tikzcd}[column sep=small, row sep=small]
 & \bullet \arrow[dl] \arrow[dr] & & & & \bullet & \\
 \bullet & & \bullet & & \bullet \arrow[ur] & & \bullet \arrow[ul]
\end{tikzcd}
\]
sont chacunes % [sic] chacunes -> chacune, orthographe
des intervalles symétriques.) Choisir une orientation revient aussi `a % [sic] `a -> à, accent imprimé comme guillemet
choisir l'origine parmi les deux extrémités.

Notons que la \textit{catégorie opposée} à un intervalle est un intervalle, et la catégorie opposée à un intervalle standard est l'intervalle standard qui s'en déduit en changeant tous les signes sur les intervalles de transition $[i-1,i]$ ($1 \leq i \leq n$). On se gardera de confondre la propriété $I \simeq I^o$, avec la propriété de symétrie.

\msp{page 9}

Ainsi les deux catégories ($\Psi_1$ et $\Phi_1$) ci-dessus sont symétriques, mais opposées l'une de l'autre sans être isomorphes. D'autre part, les $\bD_i$ ne sont pas symétriques, mais chacune isomorphe à la catégorie opposée.

Si $I$ est isomorphe à la catégorie opposée $I^o$, et si $\theta$ est un tel isomorphisme (nécessairement unique), alors il est immédiat que $\theta$ renverse l'orientation (sinon ce serait l'identité sur les sommets, et les flèches de transition seraient à la fois directes et indirectes), donc $\theta(0) = n$, $\theta(1) = n-1$, et on veut que $\underline{u}_1$ et $\underline{u}_n$ sont de même signe. Mais si $I$ est symétrique, $\underline{u}_1$ et $\underline{u}_n$ sont de signe opposé. Donc $I$ ne peut être à la fois symétrique et isomorphe à l'intervalle opposé \texttt{[plutôt à la catégorie opposée; exception :} $I = \{0\}$\texttt{]}.

Pour toute catégorie $\bD_i$, celle-ci étant un intervalle, il y a (si $i > 0$) \textit{deux orientations.} En tant qu'intervalles orientés, il y a $\bD_i^+$ (avec l'orientation ordinaire, donné par l'ordre ordinaire sur $\operatorname{Ob} \bD_i = [0,i]$), et $\bD_i^-$ avec l'orientation opposée.

\printedpage{6}
% head: A. Grothendieck : Dérivateurs, Ch. VII, Édition des Universités de Montpellier II et Paris VII
% foot: Version 22 décembre 2001

\msp{page 10}

\section*{2\quad Sous-catégories pleines d'un intervalle}

Soit $I$ une catégorie intervalle orientée, de longueur $n \geq 1$. Nous identifierons $\operatorname{Ob} I$ à l'ensemble des entiers $[0,n]$. Pour une chaîne $\underline{u}_i$, $\underline{u}_{i+1}$, $\ldots$, $\underline{u}_{i+l-1}$ de flèches de transition dans $I$, on appelle \textit{domaine} de la chaîne, l'ensemble des sources et de buts de ces éléments. C'est donc un intervalle d'entiers $[i-1, i+l-1]$, de longueur $l$ égale à celle de la chaîne, dont l'origine est déterminé % [sic] déterminé -> déterminée, accord avec l'origine
par la première flèche $\underline{u}_i$ de la chaîne (c'est sa source, si $\underline{u}_i$ est directe, et son but si $\underline{u}_i$ est rétrograde), et dont l'extrémité est déterminée par la dernière $\underline{u}_{i+l-1}$. Les intervalles d'entiers (vus comme des sommets de $I$) correspondant aux chaînes structurales sont appelés les \textit{intervalles structuraux} ou \textit{intervalles critiques} de $I$. Contrairement aux chaînes structurales, ces intervalles ne forment pas une partition de $I$, puisque deux intervalles structuraux peuvent avoir une extrémité commune. Si les chaînes structurales sont
\[
C_1,\ C_2,\ \ldots\ ,\ C_r \qquad (r = \text{ longueur utile de } I)
\]
de longueurs
\[
l_1,\ l_2,\ \ldots\ ,\ l_r,
\]
alors les intervalles structuraux sont
\[
\underbrace{[0,l_1]}_{I_1},\ \underbrace{[l_1,l_1+l_2]}_{I_2},\ \ldots\ ,\ \underbrace{[l_1+\cdots+l_{r-1},l_1+\cdots+l_{r-1}+l_r]}_{I_r},
\]
une suite d'intervalles d'entiers dont deux consécutifs ont une extrémité commune, et tels que si $i < j$, chaque élément de $I$ % [sic] I -> I_i, comparaison des éléments de I_i et de I_j
soit $\leq$ chaque élément de $I_j$. On appelle \textit{sommet critique} de $I$ les extrémités des intervalles structuraux.

\textbf{Proposition 3.} \textit{Soit $s \in I$, soit $I' = I/s$ (sous-catégorie pleine de $I$ des $s'$ tels qu'il existe $s' \longrightarrow s$), $I'' = I/s$ \textup{\texttt{[plutôt} $I'' = s\backslash I$\texttt{]}}, $I'_0 = (I/s)\backslash\{s\}$, $I''_0 = (s\backslash I)\backslash\{s\}$} (\footnote[3]{i.e. qu'on ait : il n'y a pas de flèches entre un élément de $I'_0$ et un élément de $I''_0$.}). Pour que $s$ soit critique, il faut et il suffit que $I\backslash\{s\}$ soit catégorie somme de $I'_0$ et $I''_0$.

\msp{page 11}

\textbf{Proposition 4.} \textit{Soient $i,j$ dans $I$, avec $i < j$. Conditions équivalentes :}
\begin{enumerate}
\item[\textit{a)}] \textit{Il existe une flèche entre $i$ et $j$ (soit $i \longrightarrow j$, soit $j \longrightarrow i$).}
\item[\textit{b)}] \textit{$i$ et $j$ appartiennent à un même intervalle structural (nécessairement unique).}
\item[\textit{c)}] \textit{La sous-catégorie pleine de $I$ définie par l'ensemble d'objets $[i,j]$ est isomorphe à une $\bD_\alpha$ (on aura $\alpha = j-i$).}
\end{enumerate}

\printedpage{7}
% head: A. Grothendieck : Dérivateurs, Ch. VII, Édition des Universités de Montpellier II et Paris VII
% foot: Version 22 décembre 2001
On dira alors que $i$ et $j$ sont \textit{directement connectés} dans $I$. (Et on utilisera la même terminologie dans une catégorie quelconque : $i$ et $j$ sont directement connectés si et seulement si $\operatorname{Hom}_I(i,j)$ \textit{ou} $\operatorname{Hom}_I(j,i)$ est non vide.)

\textbf{Proposition 5.} \textit{Soit $J \subseteq I$ une partie de $I$, considérons la sous-catégorie pleine de $I$ qu'elle définit. Alors :}
\begin{enumerate}
\item[\textit{a)}] \textit{Les composantes connexes $J_\alpha$ de $J$ sont des intervalles et l'ordre induit sur un $J_\alpha$ par celui de $I$ est une orientation (dite \textup{orientation induite}) de $J_\alpha$.}
\item[\textit{b)}] \textit{Supposons $\operatorname{card} J > 1$. Soit $i$ le plus petit élément de $J$, $j$ le plus grand élément, donc $i < j$. Pour que $J$ lui-même soit un intervalle, il faut et il suffit que tout sommet critique dans $[i,j]$ soit élément de $J$.}
\end{enumerate}

(\footnote[4]{On appelle \textit{sous-intervalle} de $I$ une sous-categorie % [sic] sous-categorie -> sous-catégorie, accent manquant
pleine connexe, i.e. qui est elle-même un intervalle. Notons d'ailleurs que \textit{toute} sous-catégorie $I'$ de $I$ qui est un intervalle, est nécessairement pleine, cf. plus bas \texttt{[paragraphe]} des foncteurs entre intervalles.}).

\textbf{NB} Si $s \in [i,j]$ est un sommet critique de $I$ qui n'est pas dans $J$, alors par la proposition 3 il disconnecterait $J$, i.e. $J$ serait somme directe de deux sous-catégories $J'$ et $J''$ non vides ($i \in J'$, $j \in J''$), donc $J$ ne serait pas un intervalle.

\msp{page 12}

\textbf{Corollaire.} \textit{L'ensemble des sommets critiques de $I$ est aussi l'ensemble des sommets du plus petit sous-intervalle de $I$ contenant les extrémités. Plus généralement, si $i < j$ sont deux sommets de $I$, le plus petit sous-intervalle de $I$ contenant $i$, $j$ est formé de $i$, $j$ et des sommets critiques $s$ tels que $i < s < j$.}

\textbf{Définition.} Un intervalle est dit \textit{réduit} (ou \textit{irrédondant}) si tous ses sommets sont critiques, ou ce qui revient au même, si ses chaînes structurales sont de longueur 1.

Un intervalle orienté \textit{réduit} est connu, à isomorphisme près, quand on connaît
\begin{enumerate}
\item[a)] sa longueur $\mathrm{n}$,
\item[b)] sa direction initiale $+$ (directe) ou $-$ (rétrograde).
\end{enumerate}

\textbf{NB} Cette direction est la même que la direction finale si $n$ est impair, elle est opposée si $n$ est pair.

Ainsi, suivant les deux cas de direction initiale, on trouve
\[
\begin{aligned}
0 \longrightarrow 1 \longleftarrow 2 \longrightarrow \cdots \qquad\qquad \cdots\ n & \qquad (\text{Avec } & n-1 \longrightarrow n & \quad \text{si } n \text{ impair,}\\
& & n-1 \longleftarrow n & \quad \text{si } n \text{ pair.)}\\
0 \longleftarrow 1 \longrightarrow 2 \longleftarrow \cdots \qquad\qquad \cdots\ n & \qquad (\text{Avec } & n-1 \longleftarrow n & \quad \text{si } n \text{ impair,}\\
& & n-1 \longrightarrow n & \quad \text{si } n \text{ pair.)}
\end{aligned}
\]

\printedpage{8}
% head: A. Grothendieck : Dérivateurs, Ch. VII, Édition des Universités de Montpellier II et Paris VII
% foot: Version 22 décembre 2001

\textbf{Proposition 6.} \textit{Un sous-intervalle $J$, ayant $i$ et $j$ ($i < j$) comme plus petit et plus grand sommet, est réduit si et seulement si tout $s \in J$ distinct de $i$ et de $j$ est critique.}

\textbf{Corollaire.} \textit{La sous-catégorie pleine de $I$ formée des sommets critiques est un sous-intervalle. On l'appelle l'intervalle réduit associé à $I$, noté $I_{\textup{réd}}$.}


\endgroup
% END UNIT Der-ChVII-456/Der-ChVII-456_p006-010.tex
% BEGIN UNIT Der-ChVII-456/Der-ChVII-456_p011-015.tex
\begingroup\makeatletter%
\newcommand{\printedpage}[1]{\marginpar{\footnotesize p.~#1}}%
\newcommand{\msp}[1]{\marginpar{\footnotesize [#1]}}%
\newcommand{\Ch}{\underline{\mathrm{Ch}}}%
\newcommand{\uHom}{\underline{\mathrm{Hom}}}%
\newcommand{\uFl}{\underline{\mathrm{Fl}}}%
\newcommand{\bDelta}{\boldsymbol{\Delta}}%
\newcommand{\opp}{\mathrm{opp}}%
\makeatother\let\printedpage\derpp\let\msp\dermsp%
\let\mspin\DERmspin%


\printedpage{9}
% head: A. Grothendieck : Dérivateurs, Ch. VII, Édition des Universités de Montpellier II et Paris VII
% foot: Version 22 décembre 2001

\msp{page 13}

\section*{3\quad Chemins. Catégorie de chemins d'un type donné $\tau$.}

Soient $I$ un intervalle orienté, et $X$ une catégorie. On appelle \textit{chemin de type $I$ dans $X$} tout foncteur $I \longrightarrow X$, et \textit{catégorie de chemins de type $I$ dans $X$} la catégorie
\[
\Ch_I(X) = \uHom(I,X).
\]
Comme $I$ est \textit{canoniquement} isomorphe à une catégorie intervalle standard $I_0$, on peut identifier les chemins de type $I$, au de type $I_0$. Le type $\tau$ d'un chemin le longueur $n$ $(= \operatorname{card} I)$ dans $X$ peut donc être décrit par une suite de $n$ signes $+$ ou $-$ % [sic] au de type -> à ceux de type, word(s) missing; [sic] le longueur -> de longueur, wrong preposition
\[
\tau_1,\ \tau_2,\ \dots\ ,\ \tau_n \qquad (\tau_i \in \{-1,+1\} \text{ pour } 1 \leq i \leq n).
\]
À cause du fait que $I$ est engendré pas ses flèches de transition, un foncteur $c : I \longrightarrow X$ est déterminé quand on connaît % [sic] pas ses flèches -> par ses flèches, misprint

\begin{enumerate}
\item[a)] Les $c(i)$ $(i \in I)$, i.e. une suite de $n+1$ objets de $X$ :
\[
a_0,\ a_1,\ \dots\ ,\ a_n\ (\text{\footnotemark[5]});
\]
\footnotetext[5]{$a_0$ est dit la \textit{source} du chemin, $a_n$ son but.}

\item[b)] Les $c(\underline{u}_i)$ (où les $\underline{u}_i$ sont les flèches de transition dans $I$), i.e. pour tout $i \in [1,n]$ une flèche $u_i$ dans $X$ entre $a_{i-1}$ et $a_i$, allant de $a_{i-1}$ dans $a_i$ si $\underline{u}_i$ est directe $(\tau_i = +1)$, et de $a_i$ dans $a_{i-1}$ si $\underline{u}_i$ est rétrograde $(\tau_i = -1)$,
\end{enumerate}

et à toute telle donnée a), b) correspond un foncteur. Ainsi on pourra identifier les chemins de type $I$ \msp{page 14} dans $X$ aux diagrammes du type précisé dans $X$. Et les morphismes entre chemins correspondent de même aux morphismes entre diagrammes $a$ et $a'$, une flèche $a \xrightarrow{\alpha} a'$ étant donc une suite de $n+1$ flèches
\[
\alpha_i : a_i \longrightarrow a'_i, \qquad i \in [0,n],
\]
satisfaisant la condition de commutativité pour les $n$ carrés
\[
\begin{tikzcd}
a_{i-1} \arrow[r, "u_i"] \arrow[d, "\alpha_{i-1}"'] & a_i \arrow[d, "\alpha_i"] \\
a'_{i-1} \arrow[r, "u'_i"] & a'_i
\end{tikzcd}
\qquad\qquad \text{(cas $\underline{u}_i$ directe)}
\]

\printedpage{10}
% head: A. Grothendieck : Dérivateurs, Ch. VII, Édition des Universités de Montpellier II et Paris VII
% foot: Version 22 décembre 2001
ou
\[
\begin{tikzcd}
a_{i-1} \arrow[d, "\alpha_{i-1}"'] & a_i \arrow[l, "u_i"'] \arrow[d, "\alpha_i"] \\
a'_{i-1} & a'_i. \arrow[l, "u'_i"']
\end{tikzcd}
\qquad\qquad \text{(cas $\underline{u}_i$ rétrograde)}
\]
Si on a une flèche entre deux sommets quelconque $i$, $j$ dans $I$, avec $i < j$ et $j$ non successeur immédiat de $i$, alors $i$, $j$ appartiennent à un même intervalle critique, et dans $I$ on aura % [sic] quelconque -> quelconques, agreement
\[
i \xrightarrow{\underline{u}_{i+1}} i+1 \xrightarrow{\underline{u}_{i+2}}\ \cdots\ j-1 \xrightarrow{\underline{u}_j} j \qquad \text{une chaîne de longueur } j-i \geq 2 % doubt resolved: PDF p. 12 (edition p. 10) clearly underlines u; the existing reading is retained.
\]
ou
\[
i \longleftarrow i+1 \longleftarrow\ \ \cdots\ j-1 \longleftarrow j,
\]
et la flèche correspondante entre $a_i$ et $a_j$ s'obtient en composant les flèches $u_i$ correspondantes.

\msp{page 15}

\textbf{Exemples.}

\begin{enumerate}
\item[1.] $I = \bDelta_0 = e$, $\Ch_I(X) \simeq X$. Les chemins de longueur 0 s'identifient aux objets $a$ de $X$. L'unique élément $\mathrm{id}_e$ de $\mathrm{Fl}(I)$ donne $\mathrm{id}_a$ dans $X$.

\item[2.] $I = \bDelta_1^+$, $\Ch_I(X) \simeq \uFl(X)$. Les chemins de type $\bDelta_1^+$ sont ceux de longueur 1 et \textit{directs,} i.e. donnés par une flèche
\[
a_0 \longrightarrow a_1
\]
dont la source est aussi la source $a_0$ du chemin. Par contre, les chemins de type $\bDelta_1^-$ sont donnés par des diagrammes
\[
a_0 \longleftarrow a_1,
\]
i.e. par des flèches dont le but est considéré comme source du chemin. Bien sûr, les catégories $\Ch_{\bDelta_1^+}(X)$ et $\Ch_{\bDelta_1^-}(X)$ sont canoniquement isomorphes, mais par un isomorphisme qui \textit{échange source et but du chemin.} Nous nous garderons donc bien de confondre les deux chemins précédents.

\msp{page 16}

\item[3.] $\bDelta_2^+$, $\bDelta_2^-$, correspondant aux diagrammes
\[
a_0 \longrightarrow a_1 \longrightarrow a_2
\]
resp.
\[
a_0 \longleftarrow a_1 \longleftarrow a_2,
\]
on peut faire une remarque similaire que pour $\bDelta_1^+$, $\bDelta_1^-$.

\printedpage{11}
% head: A. Grothendieck : Dérivateurs, Ch. VII, Édition des Universités de Montpellier II et Paris VII
% foot: Version 22 décembre 2001

\item[4.] $\Phi_1$, $\Psi_1$, ce sont respectivement les intervalles
\[
\Phi_1 :\ 0 \xleftarrow{\underline{u}_0} 1 \xrightarrow{\underline{u}_1} 2 \qquad\qquad \text{ou} \qquad\qquad
\begin{tikzcd}[row sep=small, column sep=small]
 & 0 \\
1 \arrow[ur] \arrow[dr] & \\
 & 2
\end{tikzcd}
\qquad , % doubt resolved: PDF p. 13 (edition p. 11) clearly underlines u in both labels; the existing reading is retained.
\]
et
\[
\Psi_1 :\ 0 \xrightarrow{\underline{u}_0} 1 \xleftarrow{\underline{u}_1} 2 \qquad\qquad \text{ou} \qquad\qquad
\begin{tikzcd}[row sep=small, column sep=small]
 & 0 \arrow[dl] \\
1 & \\
 & 2 \arrow[ul]
\end{tikzcd}
\qquad .
\]
Ils correspondent respectivement aux diagrammes
\[
a_0 \xleftarrow{u_0} a_1 \xrightarrow{u_1} a_2
\]
et
\[
a_0 \longrightarrow a_1 \longleftarrow a_2.
\]
Les catégories $\Phi_1$, $\Psi_1$ sont opposées l'une de l'autre, et ne sont pas isomorphes. Aussi n'y a-t-il en général aucun isomorphisme entre $\Ch_{\Phi_1}(X)$ et $\Ch_{\Psi_1}(X)$.
\end{enumerate}

Soit $I$ un intervalle orienté, et considérons l'\textit{intervalle opposé $I^{\opp}$}, ayant même catégorie sous-jacente mais l'orientation opposé. Comme la définition de la catégorie des chemins du type $I$ ne fait pas intervenir l'orientation de $I$, on a donc un isomorphisme canonique % [sic] l'orientation opposé -> l'orientation opposée, agreement
\[
\Ch_I(X) \xrightarrow{\sim} \Ch_{I^{\opp}}(X).
\]

\msp{page 17}

Néanmoins on se gardera de confondre un chemin de type $I$ avec le chemin opposé de type $I^{\opp}$, la source de l'un étant le but de l'autre et inversement. C'est dire qu'un chemin doit être vu comme \textit{plus} qu'un simple foncteur
\[
I \longrightarrow X,
\]
il faut de plus se donner l'\textit{orientation de $I$} comme faisant partie de la structure de chemin. Donc la définition formelle raisonnable, qui colle à l'intuition du chemin comme allant d'une \textit{source} bien définie vers un \textit{but} également bien défini, c'est qu'un chemin (défini en termes d'une catégorie intervalle $I$) est la donné d'un couple $(\omega,c)$, où $\omega$ est une orientation de $I$, et $c$ un foncteur $c : I \longrightarrow X$. Quand on regarde la catégorie des chemins d'un `type' déterminé, ce \textit{type} est $(I,\omega)$, et non $I$. C'est quand $\omega$ est fixé une bonne fois que cette catégorie s'identifie impunément à $\uHom(I,X)$ sans plus. % [sic] la donné -> la donnée, missing final e

Il faut voir cette catégorie de chemins comme munie de la structure essentielle
\[
\begin{tikzcd}
\Ch_{(I,\omega)}(X) \arrow[d, "{(\sigma,\beta)}"] \\
X \times X
\end{tikzcd}
\qquad
\begin{array}{l}
c \longmapsto (\sigma(c),\beta(c)) \\
\text{ou $\sigma(c)$ est la \textit{source,}} \\
\qquad \text{$\beta(c)$ le \textit{but} de $c$.}
\end{array} % [sic] ou sigma(c) -> où sigma(c), missing accent
\]

\printedpage{12}
% head: A. Grothendieck : Dérivateurs, Ch. VII, Édition des Universités de Montpellier II et Paris VII
% foot: Version 22 décembre 2001

\msp{page 18}

Ce foncteur s'identifie au foncteur de composition
\[
c \longmapsto c \circ \alpha
\]
où $\alpha$ est le foncteur d'inclusion
\[
\left\{
\begin{array}{lll}
\varepsilon \xrightarrow{\alpha} I & & \\[1ex]
\varepsilon = \{0,1\}, & \multicolumn{2}{l}{\text{catégorie discrète à deux éléments,}} \\[1ex]
\alpha \text{ donné par} & \alpha(0) = & \text{source de } I \\[1ex]
 & \alpha(1) = & \text{but de } I.
\end{array}
\right.
\]
Si $\omega'$ est l'orientation opposée de $I$, on a un diagramme commutatif
\[
\begin{tikzcd}
\Ch_{(I,\omega)}(X) \arrow[r, "\varphi", "\sim"'] \arrow[d, "{(\sigma,\beta)}"] & \Ch_{(I,\omega')}(X) \arrow[d, "{(\sigma',\beta')}"] \\
X \times X \arrow[r, "s"] & X \times X,
\end{tikzcd}
\]
où $s$ est le foncteur de symétrie
\[
s(x,y) = (y,x)
\]
(\footnote[6]{Le foncteur $\varphi$ est celui qui associe à tout chemin $c$, de type $\tau$, le \textit{chemin opposé $c^{\opp}$} ou $c^{-1}$, de type $\tau^{\opp}$. Si $\tau = \tau^{\opp}$, c'est donc un chemin \textit{de même type,} et $\varphi$ est un \textit{automorphisme} de $\Ch_\tau(X)$. Cet automorphisme (fonctoriel en $X$) est déduit d'un automorphisme $\theta$ de $I$, cf. ci-dessous.}). Donc l'isomorphisme $\varphi$ ne respecte pas les flèches structurales `source-but' $(\sigma,\beta)$.

Il y a cependant un embarras quand il existe un isomorphisme d'intervalles \textit{orientés}
\[
\theta : (I,\omega) \xrightarrow{\sim} (I,\omega'),
\]
i.e. quand $I$ est symétrique. Ce n'est pas le cas pour les $\bDelta_i$ $(i \geq 1)$ (ces catégories n'ont pas d'autres automorphismes que l'identité), mais c'est le cas pour les catégories $\Phi_1$, $\Psi_1$ de l'exemple $4^o$, plus généralement \msp{page 19} pour les catégories réduites de longueur \textit{paire.} Comme chaque fois qu'on a un isomorphisme $\theta$ d'intervalles \textit{orientés,} on trouve un isomorphisme entre les catégories de chemins de type $(I,\omega)$ et $(I,\omega)^{\opp} = (I,\omega')$, \textit{compatible} cette fois avec les foncteurs $(\sigma,\beta)$ :
\[
\begin{tikzcd}
\Ch_{(I,\omega)}(X) \arrow[rr, "\Psi_\theta", "\sim"'] \arrow[dr, "{(\sigma,\beta)}"'] & & \Ch_{(I,\omega')}(X) \arrow[dl, "{(\sigma',\beta')}"] \\
 & X \times X. &
\end{tikzcd}
\]

\printedpage{13}
% head: A. Grothendieck : Dérivateurs, Ch. VII, Édition des Universités de Montpellier II et Paris VII
% foot: Version 22 décembre 2001

\msp{page 20}

Il est bon de regarder un \textit{type} de chemins comme une donné non catégorique, mais combinatoire, une simple suite $\tau$ de $n$ signes $\tau_i$ $(1 \leq i \leq n)$ $+$ ou $-$ ($n =$ longueur des chemins envisagés). Cette donné combinatoire définit un intervalle orienté standard, $I_\tau$. Quand on a un intervalle orienté quelconque $(I,\omega)$, il définit un type $\tau$, et on l'identifie canoniquement au type $I_\tau$. Il vaut mieux écrire $\Ch_\tau(X)$, plutôt que $\Ch_{(I,\omega)}(X)$, pour bien indiquer que pour $X$ donné, c'est là une catégorie totalement `épinglée', sans automorphismes liés à la nature particulière de $\tau$ (à des automorphismes de $\tau$!). % [sic] une donné -> une donnée, missing final e; [sic] Cette donné -> Cette donnée, missing final e

On peut préciser ce point ainsi. Soit $\tau'$ un autre type d'intervalles, et proposons-nous de déterminer des `lois' qui permettent à tout chemin de type $\tau$ dans une catégorie arbitraire $X$, d'associer un chemin de type $\tau'$ de cette même catégorie, et ceci de façon `fonctorielle', i.e. compatible avec les applications horizontales
\[
\begin{tikzcd}
\mathrm{Ch}_\tau(X) \arrow[r, "\mathrm{Ch}_\tau(f)"] \arrow[d, "L_X"'] & \mathrm{Ch}_\tau(X') \arrow[d, "L_{X'}"] \\
\mathrm{Ch}_{\tau'}(X) \arrow[r, "\mathrm{Ch}_{\tau'}(f)"] & \mathrm{Ch}_{\tau'}(X')
\end{tikzcd}
\]
associées à un foncteur $f : X \longrightarrow X'$.

\msp{page 21}

Par exemple, prenant le type $\bDelta_2^+$ (i.e. $+1,+1$), correspondant aux chemins
\[
a_0 \xrightarrow{u_1} a_1 \xrightarrow{u_2} a_2,
\]
on peut associer à un tel chemin le chemin de type \texttt{[}$\bDelta_1^+$\texttt{]}
\[
a_0 \xrightarrow{u_2u_1} a_2,
\]
ou aussi le chemin de type $\bDelta_2^+$
\[
a_0 \xrightarrow{u_2u_1} a_2 \xrightarrow{\mathrm{id}} a_2,
\]
etc. Comme le foncteur
\[
\left\{
\begin{array}{cccl}
X & \longmapsto & \mathrm{Ch}_\tau(X) & \qquad (= \mathrm{Ob}\, \Ch_\tau(X)) \\[1ex]
\mathrm{Cat} & \longrightarrow & (\mathrm{Ens}) &
\end{array}
\right.
\]
est \textit{représenté} par l'objet $I_\tau$ de $X$ \texttt{[plutôt de} Cat\texttt{]}, le sorite des foncteurs représentables nous dit que de telles lois de formation de chemins de type $\tau$ à partir des chemins de type $\tau'$ \texttt{[plutôt} $\tau'$ \texttt{à partir des chemins de type} $\tau$\texttt{]} correspondent biunivoquement aux flèches dans Cat
\[
I_{\tau'} \xrightarrow{\varphi} I_\tau\ ,
\]


\endgroup
% END UNIT Der-ChVII-456/Der-ChVII-456_p011-015.tex
% BEGIN UNIT Der-ChVII-456/Der-ChVII-456_p016-020.tex
\begingroup\makeatletter%
\newcommand{\printedpage}[1]{\marginpar{\footnotesize p.~#1}}%
\newcommand{\msp}[1]{\marginpar{\footnotesize [#1]}}%
\newcommand{\Hom}{\mathrm{Hom}}%
\newcommand{\uHom}{\underline{\mathrm{Hom}}}%
\newcommand{\Ch}{\mathrm{Ch}}%
\newcommand{\uCh}{\underline{\mathrm{Ch}}}%
\newcommand{\Ty}{\mathrm{Ty}}%
\newcommand{\uTy}{\underline{\mathrm{Ty}}}%
\newcommand{\Cat}{\mathrm{Cat}}%
\newcommand{\bDelta}{\boldsymbol{\Delta}}%
\newcommand{\circled}[1]{\text{\textcircled{$\scriptstyle #1$}}}%
\makeatother\let\printedpage\derpp\let\msp\dermsp%
\let\mspin\DERmspin%


\printedpage{14}
% head: A. Grothendieck : Dérivateurs, Ch. VII, Édition des Universités de Montpellier II et Paris VII
% foot: Version 22 décembre 2001

en prenant simplement la ‘loi’ correspondante à $\varphi$
\[
c \longmapsto c\circ\varphi \qquad \text{pour } c : I_\tau \longrightarrow X.
\]
De façon plus précise, l'application naturelle
\[
\underbrace{\Hom_{\Cat}(I_{\tau'},I_\tau)}_{=\mathrm{Ob}\,\uHom(I_{\tau'},I_\tau)} \longrightarrow \Hom(\mathrm{Ch}_\tau,\mathrm{Ch}_{\tau'}),
\]
\msp{page 22}
où le deuxième Hom est considéré comme un Hom d'objets dans $\Cat^{\vee}=\uHom(\Cat,\mathrm{Ens})$. Ainsi, dans les deux exemples plus haut, où $\tau=\bDelta_2^+$ et $\tau'=\bDelta_1^+$ ou $\bDelta_2^+$, les foncteurs en sens inverse sont
\[
\begin{tikzcd}
\bDelta_2^+\ \Big\{ & 0 \arrow[r] & 1 \arrow[r] & 2 \\
\bDelta_1^+\ \Big\{ & 0 \arrow[u, mapsto] \arrow[rr] & & 1 \arrow[u, mapsto]
\end{tikzcd}
\]
et
\[
\begin{tikzcd}
\bDelta_2^+\ \Big\{ & 0 \arrow[r] & 1 \arrow[r] & 2 \\
\bDelta_2^+\ \Big\{ & 0 \arrow[u, mapsto] \arrow[r] & 1 \arrow[ur, mapsto] \arrow[r] & 2\,. \arrow[u, mapsto]
\end{tikzcd}
\]
C'est là un principe général pour ramener toutes les ‘opérations’ $\pm$ canoniques sur des chemins dans des catégories arbitraires, et des énoncés les concernant, à des questions où $X$ disparaît et où on travaille dans les catégories intervalles elles-mêmes. Un des avantages techniques, c'est que ce sont des catégories fort petites, donnant lieu de plus à une intuition topologique sûre (qui risque de se volatiliser par moments en travaillant dans des catégories de chemins), et aussi que dans de telles catégories, donc aussi dans les catégories de foncteurs entre elles, il y a toujours au plus
\msp{page 23}
\textit{une seule flèche} allant d'un objet donné vers un autre. En d'autres termes, dans de telles catégories tous les diagrammes sont commutatifs.

Bien sûr, on voit aussi que si on a une loi de formation $L$ de chemins de type $\tau'$ à partir de chemins de type $\tau$ (avec la condition de fonctorialité par rapport à la catégorie $X$ où on prend les chemins), cette loi est nécessairement \textit{fonctorielle par rapport aux chemins :} si on a un homomorphisme
\[
c \xrightarrow{v} c' \qquad \text{dans } \uCh_\tau(X),
\]
\printedpage{15}
% head: A. Grothendieck : Dérivateurs, Ch. VII, Édition des Universités de Montpellier II et Paris VII
% foot: Version 22 décembre 2001
alors la loi $L$ ‘se prolonge’ à de tels $v$ pour donner
\[
L(c) \xrightarrow{L(v)} L(c'),
\]
avec les relations habituelles
\[
L(\mathrm{id}_C)=\mathrm{id}_{L(c)}, \qquad L(v'v)=L(v')L(v). % [sic] \mathrm{id}_C -> \mathrm{id}_c, capital C printed while the chemin is c (cf. L(c))
\]
Cela résulte de la description de $L$ comme une opération $L_\varphi$ en termes d'un foncteur bien déterminé
\[
I_{\tau'} \xrightarrow{\varphi} I_\tau,
\]
lequel à son tour donne pour toute $X$ un \textit{foncteur}
\[
\left\{
\begin{array}{rcl}
\uCh_\tau(X) & \longrightarrow & \uCh_{\tau'}(X)\\[1ex]
c & \longmapsto & c\circ\varphi
\end{array}
\right.
\]
plus précis que simplement un \textit{application} $\mathrm{Ch}_\tau(X)\longrightarrow\mathrm{Ch}_{\tau'}(X)$ entre les ensembles d'objets correspondants. % [sic] un application -> une application, gender agreement

\msp{page 24}
On a le choix maintenant entre

\begin{enumerate}[label=\alph*)]
\item Une étude plus ou moins systématique des foncteurs entre intervalles. Je préfère attendre pour cela d'avoir besoin de ceci et de cela, et de ne faire que le strict minimum pour les constructions que j'ai en vue.
\item Une description de la compositions des chemins. % [sic] la compositions -> la composition, number agreement
\item Un formalisme des catégories de chemins \textit{de type non précisé.}
\end{enumerate}

Comme mon objectif est surtout de décrire de telles catégories $\uCh(X)$ (sans indice $\tau$!), avec des bonnes propriétés, je commencerai par c), regarderai ensuite b) comme un test crucial pour l'utilisabilité des notions obtenues, et enfin m'occuperai de a) quand je saurai de quoi j'aurai en besoin.

\printedpage{16}
% head: A. Grothendieck : Dérivateurs, Ch. VII, Édition des Universités de Montpellier II et Paris VII
% foot: Version 22 décembre 2001

\section*{4\quad Catégorie de chemins $\uCh(X)$}

Je veux ‘inclure’ toutes les catégories $\uCh_\tau(X)$ de chemins de tous les types qu'on veut, dans une même catégorie $\uCh(X)$. Rappelons
\[
\uCh_\tau(X)=\uHom(I_\tau,X),
\]
où $I_\tau$ est l'intervalle (orienté) standard de type $\tau$.

L'ensemble de tous les types d'intervalles
\msp{page 25}
peut être identifié au monoïde libre Ty à deux générateurs $\Delta^+$ et $\Delta^-$, dont les éléments sont les suites finies quelconques (eventuellement vides) d'objets alternativement égales à $\Delta^+$ et à $\Delta^-$. Si le premier élément est $\Delta^+$, on a un type de chemin à direction initiale directe, rétrograde dans le cas contraire. La décomposition de $I_\tau$, ou de $\mathrm{Fl}\,(I_\tau)$, en intervalles critiques correspond à l'écriture réduite d'un mot comme % [sic] eventuellement -> éventuellement, missing accent
\[
(\Delta^+)^{l_1}(\Delta^-)^{l_2}\ \cdots\ (\Delta^?)^{l_r} \qquad (\Delta^?=\Delta^+ \text{ si } r \text{ impair, } \Delta^?=\Delta^- \text{ si } r \text{ pair}),
\]
ou comme
\[
(\Delta^-)^{l_1}(\Delta^+)^{l_2}\ \cdots\ (\Delta^?)^{l_r} \qquad (\Delta^?=\Delta^- \text{ si } r \text{ impair, } \Delta^?=\Delta^+ \text{ si } r \text{ pair}).
\]
Les $l_i$ sont les longueurs des intervalles critiques successifs (ou des chaines structurales successives), leur nombre est la longueur réduite de l'intervalle.

On va considérer Ty comme l'ensemble des objets d'une catégorie $\uTy$ :
\[
\left\{
\begin{array}{l}
\mathrm{Ob}\,\uTy=\Ty=\ \text{monoïde libre à deux générateurs } \Delta^+,\ \Delta^- \text{ précédents.}\\[1ex]
\text{Si } \tau,\tau'\in\Ty,\ \Hom(\tau,\tau')=\Hom_{\Cat}(I_\tau,I_{\tau'}).
\end{array}
\right.
\]

(\footnote[7]{Définition un peu canulée, il faut se borner aux homomorphismes de $I_\tau$ dans $I_{\tau'}$ compatibles avec source et but - cf. p.~29.}). On notera qu'on a non seulement un ensemble $\Hom_{\Cat}(I_\tau,I_{\tau'})$, mais même une catégorie $\uHom(I_\tau,I_{\tau'})$, ce qui revient ici à dire que $\Hom(I_\tau,I_{\tau'})$ (i.e. $\Hom(\tau,\tau')$) peut être vu comme un ensemble \textit{ordonné.} Mais pour l'instant nous allons ignorer cette relation d'ordre.

\msp{page 26}
La composition des morphismes est la composition des foncteurs entre catégories $I_\tau$. Ainsi on a un foncteur \textit{pleinement fidèle} (\footnote[8]{faux, cf. annotation marginale page précédente - on ne prend pas \textit{tous} les foncteurs dans Cat entre $I_\tau$ et $I_{\tau'}$.}) canonique
\[
\tau \longmapsto I_\tau : \uTy \xrightarrow{\rho} (\Cat) \qquad \text{(\footnotemark[9]),}
\]
\footnotetext[9]{foncteur \textit{réalisation} d'un type (combinatoire) d'intervalles.}% footnote text printed at the foot of p. 17
\printedpage{17}
% head: A. Grothendieck : Dérivateurs, Ch. VII, Édition des Universités de Montpellier II et Paris VII
% foot: Version 22 décembre 2001
par quoi la catégorie $I_\tau$ \texttt{[plutôt} $\uTy$\texttt{]} pourrait s'identifier à une sous-catégorie pleine de (Cat) (celle des $I_\tau$). Mais je préfère ne pas identifier $\tau\in\Ty$ (donnée combinatoire) avec $I_\tau$ (donnée catégorique).

Soit $X$ une (petite?) catégorie. On a un contrafoncteur
\[
\left\{
\begin{array}{rcl}
S & \longmapsto & \uHom(S,X)\\[1ex]
(\Cat)^o & \longrightarrow & (\Cat)
\end{array}
\right.
\]
défini par $X$, d'où en composant avec $\rho$ un contrafoncteur
\[
\uTy^{o} \xrightarrow{\mathrm{CH}_X} \Cat.
\]
C'est le contrafoncteur
\[
\left\{
\begin{array}{l}
\tau \longmapsto \uHom(I_\tau,X) \overset{\text{déf}}{=} \uCh_\tau(X),\\[1ex]
\text{i.e. } \mathrm{CH}_X(\tau)=\uCh_\tau(X).
\end{array}
\right.
\]
Mais un contrafoncteur sur une petite catégorie $\uTy$ donne lieu à une \textit{catégorie fibrée} (et même scindée) sur celle-ci.

\msp{page 27}
Je vais désigner cette catégorie fibrée par $\uCh(X)$. (Provisoirement - si elle s'avère trop grosse, je la restreindrai \dots) Ainsi on a
\[
\begin{tikzcd}[row sep=large]
\uCh(X) \arrow[d, "\pi_X"] \\
\uTy\ .
\end{tikzcd}
\]
La fibre de $\uCh(X)$ en $\tau\in\mathrm{Ob}\,\uTy$ est la catégorie $\uCh_\tau(X)$ des chemins de type $\tau$ dans $X$. Cela nous dit donc quelles sont les flèches dans $\uCh(X)$ au dessus des flèches identiques dans $\uTy$. Mais que sont les flèches plus générales? On se donne une flèche dans $\uTy$
\[
\varphi : \tau' \longrightarrow \tau, \qquad \text{i.e. } I_{\tau'} \xrightarrow{\varphi} I_\tau,
\]
quelles sont les flèches de $\uCh(X)$ au dessus de $\varphi$? Par le sorite général des catégories fibrées, si
\[
\begin{array}{lclcl}
c\in\uCh(X)_\tau & = & \uCh_\tau(X) & \overset{\text{déf}}{=} & \uHom(I_\tau,X)\\[1ex]
c'\in\uCh(X)_{\tau'} & = & \uCh_{\tau'}(X) & \overset{\text{déf}}{=} & \uHom(I_{\tau'},X)
\end{array}
\]
on a
\[
\boxed{\Hom_\varphi(c',c)\simeq\Hom(c',\underbrace{\varphi^*(c)}_{=c\circ\varphi}).}
\]
Donc : les homomorphismes du
\printedpage{18}
% head: A. Grothendieck : Dérivateurs, Ch. VII, Édition des Universités de Montpellier II et Paris VII
% foot: Version 22 décembre 2001
\msp{page 28}
chemin $c'$ de type $\tau'$ dans $X$, dans le chemin $c$ de type $\tau$, ‘au dessus’ du morphisme $\varphi$ de $\tau'$ dans $\tau$, i.e. du foncteur
\[
\varphi : I_{\tau'} \longrightarrow I_\tau,
\]
est un homomorphisme $v$ de $c'$ dans $c\circ\varphi$, i.e.
\[
\begin{tikzcd}[row sep=large]
I_{\tau'} \arrow[r, "\varphi"] \arrow[dr, "c'"', ""{name=A}] & I_\tau \arrow[d, "c", ""'{name=B}] \\
& X.
\arrow[from=A, to=B, bend left=20, shorten <=6pt, shorten >=2pt, "\circled{v}"]
\end{tikzcd}
\]
La composition des flèches est évidente :
\[
\begin{tikzcd}[row sep=large]
I_{\tau''} \arrow[r, "\varphi'"] \arrow[drr, "c''"', ""{name=C}] & I_{\tau'} \arrow[r, "\varphi"] \arrow[dr, "c'"', ""{name=A}] & I_\tau \arrow[d, "c", ""'{name=B}] \\
& & X,
\arrow[from=A, to=B, bend left=20, shorten <=6pt, shorten >=2pt, "v"]
\arrow[from=C, to=A, shorten <=4pt, shorten >=4pt, "v'"]
\end{tikzcd}
\]
on a
\[
\begin{array}{rccl}
v : & c' & \longrightarrow & c\circ\varphi, \quad \text{d'où } c'\circ\varphi' \xrightarrow{v*\varphi'} \underbrace{(c\circ\varphi)\circ\varphi'}_{=c\circ(\varphi\varphi')}\\[3ex]
v' : & c'' & \longrightarrow & c'\circ\varphi'
\end{array}
\]
et
\[
v''=(v*\varphi')\circ v', \qquad \text{donc} \qquad
\begin{tikzcd}
c'' \arrow[r, "v'"] \arrow[rr, bend right=12, "\circled{v''}"'] & c'\circ\varphi' \arrow[r, "v*\varphi'"] & c\circ(\varphi\varphi').
\end{tikzcd}
\]
Mais a-t-on des foncteurs source et but
\[
\uCh(X) \overset{\sigma}{\underset{\beta}{\rightrightarrows}} X \qquad
\begin{array}{ccl}
c & \longmapsto & \text{source}(c)\\[1ex]
c & \longmapsto & \text{but}(c)
\end{array}
\qquad ?
\]
Il faudrait donc qu'un homomorphisme $v$ comme dessus
\[
v : c' \longrightarrow c\circ\varphi
\]
définisse
\[
\begin{array}{lrcl}
\sigma(c) : & \text{source}(c') & \longrightarrow & \text{source}(c)\\[1ex]
\beta(c) : & \text{but}(c') & \longrightarrow & \text{but}(c).
\end{array}
\]


\endgroup
% END UNIT Der-ChVII-456/Der-ChVII-456_p016-020.tex
% BEGIN UNIT Der-ChVII-456/Der-ChVII-456_p021-025.tex
\begingroup\makeatletter%
\newcommand{\printedpage}[1]{\marginpar{\footnotesize p.~#1}}%
\newcommand{\msp}[1]{\marginpar{\footnotesize [#1]}}%
\newcommand{\ul}[1]{\underline{\mathrm{#1}}}%
\newcommand{\bDelta}{\boldsymbol{\Delta}}%
\newcommand{\red}{\text{\upshape réd}}%
\newcommand{\circled}[1]{\tikz[baseline=(c.base)]\node[draw,circle,inner sep=0.4pt](c){\scriptsize #1};}%
\makeatother\let\printedpage\derpp\let\msp\dermsp%
\let\mspin\DERmspin%


\printedpage{19}
% head: A. Grothendieck : Dérivateurs, Ch. VII, Édition des Universités de Montpellier II et Paris VII
% foot: Version 22 décembre 2001

\msp{page 29}
Mais comme
\[
\left\{
\begin{array}{l}
\sigma(c)=c(0_\tau),\ \beta(c)=c(1_\tau)\\[1ex]
\text{où } 0_\tau,\ 1_\tau \text{ sont les objets source et but de l'intervalle orienté standard } I_\tau\ ,
\end{array}
\right.
\]
et de même pour $\tau'$, $v$ n'induira bien
\[
\underbrace{\sigma(c')}_{=c'(0_{\tau'})} \longrightarrow c(\varphi(0_{\tau'})) \overset{?}{=} \underbrace{\sigma(c)}_{c(0_\tau)},
\]
et itou pour $\beta(c')$ et $\beta(c)$, que si on impose
\[
\boxed{
\left\{
\begin{array}{lcl}
\varphi(0_{\tau'}) & = & 0_\tau\\
\varphi(1_{\tau'}) & = & 1_\tau,
\end{array}
\right.
}
\]
i.e. que le foncteur envisagé $\varphi$ soit \textit{compatible avec la source et le but.} C'est ce visiblement qu'il fallait imposer dans la définition de % [sic] C'est ce visiblement -> C'est visiblement ce, word order
\[
\mathrm{Hom}_{\ul{Ty}}(\tau',\tau)=\underline{\mathrm{Hom}}^{\circled{int}}(I_{\tau'},I_\tau)\quad\text{-}
\]
il faut prendre les foncteurs $\varphi$ respectant les structures d'intervalle, i.e. rendant commutatif le diagramme
\[
\begin{tikzcd}[column sep=small]
I_{\tau'} \arrow[rr] & & I_\tau \\
& \varepsilon=\{0,1\} \arrow[ul, hook', "i_{\tau'}"] \arrow[ur, hook, "i_\tau"'] &
\end{tikzcd}
\]
définissant l'orientation de $I_{\tau'}$ et de $I_\tau$ (\footnote[10]{En fait, il faut sans doute restreindre encore plus les morphismes en exigeant qu'il soient \textit{croissants} sur les objets - cf. page 31.}). Cette rectification faite, j'ai bon espoir que la catégorie $\ul{Ch}(X)$ qu'on vient de construire est `la bonne'. On pourrait désigner par $\ul{CH}(X)$ la catégorie plus grande, sur la catégorie $\ul{Typ}$ ayant mêmes objets que $\ul{Ty}$, mais plus de morphismes : à savoir % [sic] (in footnote 10) qu'il soient -> qu'ils soient, missing plural s
\msp{page 30}
\[
\begin{array}{lcl}
\mathrm{Hom}_{\ul{Typ}}(\tau,\tau') & = & \mathrm{Hom}_{\mathrm{Cat}}(I_\tau,I_{\tau'})\\[0.5ex]
\mathrm{Hom}_{\ul{Ty}}(\tau,\tau') & = & \mathrm{Hom}^{\mathrm{int}}_{\mathrm{Cat}}(I_\tau,I_{\tau'})\\[0.5ex]
 & & \text{(homomorphismes respectant la structure d'intervalle).}
\end{array}
\]
On aura
\[
\ul{Ch}(X)\simeq\ul{CH}(X)\times_{\ul{Typ}}\ul{Ty},
\]
\printedpage{20}
% head: A. Grothendieck : Dérivateurs, Ch. VII, Édition des Universités de Montpellier II et Paris VII
% foot: Version 22 décembre 2001
i.e. $\ul{Ch}(X)$ est la \textit{restriction} de la catégorie fibrée $\ul{CH}(X)$ sur $\ul{Typ}$, à la sous-catégorie (non-pleine) $\ul{Ty}$ de $\ul{Typ}$.

\textbf{NB} Les catégories fibrées $\ul{CH}(X)$, $\ul{Ch}(X)$ sur $\ul{Typ}$, $\ul{Ty}$ dépendent fonctoriellement de $X$, variant dans (Cat), comme il se doit. Comme de juste, on s'y intéresse surtout quand $X$ \textit{n'est pas} dans (Cat), i.e. quand $X$ est une `grosse' catégorie (telle que (Cat) elle-même). Je me dispense d'expliciter ici les questions pertinentes d'Univers.

Deux tests essentiels, pour s'assurer si la notion est bonne
\begin{enumerate}
\item[a)] La théorie des $\mathcal{M}\Sigma^{-1}$.
\item[b)] La composition des chemins doit bien marcher.
\end{enumerate}

\msp{page 31}
On a le diagramme de foncteurs canoniques
\[
\begin{tikzcd}
\ul{Ty} \arrow[r, hook] \arrow[d, hook] & \varepsilon\backslash\mathrm{Cat} \arrow[d] \\
\ul{Typ} \arrow[r, hook] & \mathrm{Cat},
\end{tikzcd}
\]
où les flèches horizontales sont pleinement fidèles.

\printedpage{21}
% head: A. Grothendieck : Dérivateurs, Ch. VII, Édition des Universités de Montpellier II et Paris VII
% foot: Version 22 décembre 2001

\section*{5\quad Propriétés des morphismes d'intervalles}

Je me suis plus ou moins convaincu que pour se prêter à des calculs et à des raisonnements simples, il faut encore plus restreindre les morphismes dans la définition de la catégorie $\ul{Ty}$. Je vais exiger que les morphismes $\tau\longrightarrow\tau'$ soient des foncteurs
\[
w : I_\tau\longrightarrow I_{\tau'}
\]
qui sont compatibles avec la structure d'\textit{intervalles orientés} de $I_\tau$, $I_{\tau'}$, i.e. qu'ils transforment source en source, but en but, mais de plus qu'ils soient \textit{croissants} sur les objets (pour la relation d'ordre de progression des chemins). C'est une restriction qui géométriquement me paraît assez naturelle. (Elle est similaire à la restriction analogue, dans la définition pré-catégorique des morphismes entre $\bDelta_i$ \ldots\,) C'est cette hypothèse qui permet de faire une théorie de structure assez simple des morphismes d'intervalles. Mais pour l'instant je vais distinguer `morphismes' et `morphismes croissants'.

\msp{page 32}
Soit
\[
f : I_\tau\longrightarrow I_{\tau'}
\]
un morphisme d'intervalles orientés. Considérons les sous-catégories pleines, isomorphes à des $\Delta_i^+$ ou $\Delta_i^-$ comme intervalles orientés, définis par les intervalles critiques de $I_\tau$. Ils forment une suite $\Sigma_\tau$ de sous-catégories de $I_\tau$
\[
\Sigma_\tau=(\bDelta^{\varepsilon_1}_{l_1}(\tau),\bDelta^{\varepsilon_2}_{l_2}(\tau),\ldots,\bDelta^{\varepsilon_r}_{l_r}(\tau))
\]
dont $I_\tau$ est la somme amalgamée (on reviendra sur ce point plus \texttt{[bas]}), $r$ étant la longueur utile de $\tau$. On les appelle les \textit{sous-intervalles critiques} de $I_\tau$. Ainsi $f$ se décompose en une suite de morphismes
\[
f_i : \bDelta^{\varepsilon_i}_{l_i}\longrightarrow I_{\tau'},\qquad i\in[1,r],
\]
tels que
\begin{enumerate}
\item[a)] pour $i\in[1,r-1]$, le but de $f_i$ soit égal à la source de $f_{i+1}$.

De plus,
\item[b)] $f_1$ transforme source en source, et $f_r$ transforme but en but.
\end{enumerate}

Inversement, quand on se donne des $f_i$ satisfaisant ces deux conditions, ils définissent un morphisme d'intervalles $f$ unique. (Et si la condition b) n'est pas exigée, on trouve une description des morphismes de catégories $I_\tau\longrightarrow I_{\tau'}$.)

\msp{page 33}
On va d'abord regarder un foncteur quelconque
\[
f : \bDelta^{\varepsilon}_i\longrightarrow I_{\tau'}\qquad(\varepsilon\in\{\pm1\})\quad .
\]

\printedpage{22}
% head: A. Grothendieck : Dérivateurs, Ch. VII, Édition des Universités de Montpellier II et Paris VII
% foot: Version 22 décembre 2001
Supposons pour fixer les idées que $\varepsilon=+1$, i.e. que le premier membre soit un intervalle critique direct \ldots{} On trouve que $f$ se factorise par un sous-intervalle critique de $I_{\tau'}$, soit $\bDelta^{\varepsilon'}_j$. Il est donc défini par une application croissante de $\bDelta^{\varepsilon}_i$ dans $\bDelta^{\varepsilon'}_j$ (si $\varepsilon$, $\varepsilon'$ \texttt{[sont]} de même signe), une application décroissante dans le sens inverse (\footnote[11]{\textbf{Lemme.} \textit{Si $I_{\tau'}$ est un intervalle, toute partie $J$ de $I_{\tau'}$ telle que pour deux éléments de $J$, il existe toujours une flèche de l'un dans l'autre, est contenue dans un intervalle critique (unique si} $\operatorname{card} J\geq 2$\textit{).}}). Mais si on exige que $f$ soit croissant, alors il faut que \textit{ou bien} $f$ soit constant (auquel cas, si sa valeur n'est pas une extrémité de $I_{\tau'}$, il y a exactement \textit{deux} intervalles critiques de $I_{\tau'}$ par lesquels se factorise $f$) \textit{ou bien} $f$ non constant et $\varepsilon=\varepsilon'$ (et en tous cas, si $f$ non constant, alors l'intervalle critique par lequel se factorise $f$ est unique). Pour résumer :

\textbf{Proposition 7.} \textit{Soient $I_\tau$, $I_{\tau'}$ deux intervalles orientés,}
\[
f : I_\tau\longrightarrow I_{\tau'}
\]
\textit{un foncteur. Alors pour tout intervalle critique $J\subseteq I_\tau$ tel que $f|J$ soit non constant, il existe un unique intervalle critique $J'\subseteq I_{\tau'}$ tel que $f|J$ se factorise par $J'$, i.e. tel que $f(J)\subseteq J'$. Si $f$ est croissant, alors $J$ et $J'$ sont de même signe,}
\msp{page 34}
\textit{i.e. tous deux directs, ou tous deux rétrogrades.}

Ainsi, si on regarde la suite $s_0,s_1,\ldots,s_n$ des sommets critiques de $\tau$, alors
\[
t_0=f(s_0),\ t_1=f(s_1),\ \ldots\ ,\ t_n=f(s_n)
\]
forment une suite de sommets de $I_{\tau'}$, telle que deux consécutifs soient reliés par une flèche de $I_{\tau'}$. Donc la sous-catégorie pleine de $I_{\tau'}$ définie par ces $t_i$ est connexe, c'est donc un sous-intervalle. Si $f$ est croissant, ce sous-intervalle est $[t_0,t_n]$. Si $f$ est un morphisme d'intervalles, il contient tous les sommets critiques de $I_{\tau'}$. Dans tous les cas, prenant $t_i$, $t_j$ $(i<j)$ et appliquant ce qui précède à $f|J$, où $J$ est le sous-intervalle de $I$ d'origine $s_i$, d'extrémité $s_j$, on voit que tout sommet critique de $I_{\tau'}$ compris entre $t_i$ et $t_j$ est parmi les $t_k$ $(i\leq k\leq j)$. Pour résumer :

\textbf{Corollaire.} \textit{Soient $s_i\leq s_j$ deux sommets critiques de $I_{\tau'}$. Alors tout sommet critique de $I_{\tau'}$ compris entre $f(s_i)=t_i$ et $f(s_j)=t_j$, est l'image d'un sommet critique $s_k$ compris entre $s_i$ et $s_j$ (i.e. $i\leq k\leq j$). En particulier, si $f$ est un morphisme d'intervalles,} % [sic] sommets critiques de I_{\tau'} (first occurrence) -> de I_\tau, the s_i are the critical vertices of \tau
\[
f(I_{\tau\,\red})\supseteq I_{\tau'\,\red}
\]
\textit{(où $I_{\tau\,\red}$ désigne l'ensemble des sommets critiques de $I_\tau$ ou aussi le sous-intervalle réduit qu'il définit).}

\msp{page 35}
On fera attention que même si $f$ est croissant, on \textit{n'a pas} en général $f(I_{\tau\,\red})\subseteq I_{\tau'\,\red}$, i.e. $f$ ne transforme pas forcément sommets critiques en sommets critiques, Exemple minimal :
\printedpage{23}
% head: A. Grothendieck : Dérivateurs, Ch. VII, Édition des Universités de Montpellier II et Paris VII
% foot: Version 22 décembre 2001

\noindent
\begin{minipage}[c]{0.6\textwidth}
\[
\begin{tikzcd}
\Phi_3^+\ : & s_0 \arrow[r] \arrow[d, mapsto] & s_1 \arrow[d, mapsto] & s_2 \arrow[l] \arrow[r] \arrow[dl, mapsto] & s_3 \arrow[dl, mapsto] \\
\bDelta_2^+\ : & {} \arrow[r] & {} \arrow[r] & {} &
\end{tikzcd}
\]
\end{minipage}\hfill
\begin{minipage}[c]{0.35\textwidth}
$\Phi_3^+$ est réduite, les deux sommets critiques transformés en le sommet non critique 1.
\end{minipage}

Désignons par $\mathrm{Cr}(\tau)$ l'ensembles des intervalles critiques de $I_\tau$, en correspondance biunivoque (respectant l'ordre total) avec l'ensemble des flèches de transition de $I_{\tau\,\red}$, % [sic] l'ensembles -> l'ensemble, spurious plural s
\[
\mathrm{Cr}(\tau)\simeq\mathrm{Fl}\,\mathrm{tr}\,(I_{\tau\,\red}).
\]
Soit, pour un foncteur
\[
f : I_\tau\longrightarrow I_{\tau'},
\]
$\mathrm{Cr}(\tau,f)$ l'ensemble des intervalles critiques de $I_\tau$ sur lesquels $f$ ne soit pas constant. On a alors une application
\[
\mathrm{Cr}(\tau,f)\longrightarrow\mathrm{Cr}(\tau') \tag*{$(*)$} % label (*) printed at the left margin
\]
induite par $f$ via la proposition 7, associant à tout $J\in\mathrm{Cr}(\tau,f)$ l'unique $J'\in\mathrm{Cr}(\tau')$ tel que $f(J)\subseteq J'$. On a de plus

\textbf{Corollaire 2.}
\begin{enumerate}
\item[\textit{a)}] \textit{L'image de $(*)$ est un segment de l'ensemble totalement ordonné} $\mathrm{Cr}(\tau')$\textit{, i.e. si deux intervalles critiques $J'$, $J''$ dans $I'_{\tau'}$ sont dans l'image, il en est de même de tous} % [sic] I'_{\tau'} -> I_{\tau'}, spurious prime
\msp{page 36}
\textit{ceux qui sont compris entre eux.}
\item[\textit{b)}] \textit{Cette image contient aussi l'intervalle critique de $I_{\tau'}$ (il y en a au plus un) qui contient $f(0_\tau)$ (resp. $f(1_\tau)$) et dont l'intersection avec le segment $[f(0_\tau),f(1_\tau)]$ est de cardinal $\geq 2$, ainsi que tous les intervalles critiques contenus dans ledit segment.}
\item[\textit{c)}] \textit{Si $f$ est un morphisme d'intervalles orientés, alors l'application induite}
\[
\mathrm{Cr}(f):\mathrm{Cr}(\tau,f)\longrightarrow\mathrm{Cr}(\tau')
\]
\textit{est surjective. De façon plus précise, soit $J'\subseteq$} $\mathrm{Cr}(\tau')$\textit{, et soit $J=f^{-1}(J')$. Alors} $\mathrm{Cr}(f)^{-1}(\{J'\})$ \textit{est formé des intervalles critiques dans} $\mathrm{Cr}(\tau,f)$ contenus dans $J$. \textit{On a de plus,} si $f$ est croissant, % [sic] J' \subseteq Cr(\tau') -> J' \in Cr(\tau'), J' is an element of Cr(\tau')
\[
\mathrm{long}(J')=\sum_{K\in\mathrm{Cr}(f)^{-1}(\{J'\})}\pi_{\mathrm{n}}(K,f),
\]
\textit{où $\pi_{\mathrm{n}}(K,f)$ est la `partie' (numérique) de $f|K$, i.e. la longueur de l'intervalle $[f(0_K),f(1_K)]$. (}\textbf{NB} \textit{Si $K$ est un intervalle critique de $I$ tel que $K\notin Cr(\tau,f)$, i.e. $f|K$ est constant, on posera $\pi_{\mathrm{n}}(K,f)=0$. Donc la formule précédente se lit aussi}
\[
\mathrm{long}(J')=\sum\limits_{\underbrace{K\in\mathrm{Cr}(I_\tau),\ K\subseteq J}_{\text{(i.e. }K\,\in\,\mathrm{Cr}(f^{-1}(J'))\text{)}}}\pi_{\mathrm{n}}(K,f).\ \textit{)}
\]
\end{enumerate}


\endgroup
% END UNIT Der-ChVII-456/Der-ChVII-456_p021-025.tex
% BEGIN UNIT Der-ChVII-456/Der-ChVII-456_p026-030.tex
\begingroup\makeatletter%
\newcommand{\printedpage}[1]{\marginpar{\footnotesize p.~#1}}%
\newcommand{\msp}[1]{\marginpar{\footnotesize [#1]}}%
\newcommand{\Tr}{\mathrm{Tr}}%
\newcommand{\Cr}{\mathrm{Cr}}%
\newcommand{\Hom}{\mathrm{Hom}}%
\newcommand{\Ty}{\underline{\mathrm{Ty}}}%
\newcommand{\Ob}{\mathrm{Ob}}%
\newcommand{\lng}{\operatorname{long}}%
\newcommand{\card}{\operatorname{card}}%
\newcommand{\bD}{\boldsymbol{\Delta}}%
\newcommand{\uTy}{\underline{\mathrm{Ty}}}%
\makeatother\let\printedpage\derpp\let\msp\dermsp%
\let\mspin\DERmspin%


\printedpage{24}
% head: A. Grothendieck : Dérivateurs, Ch. VII, Édition des Universités de Montpellier II et Paris VII

\textbf{Corollaire 3.} \textit{Si $f$ est un morphisme} croissant \textit{d'intervalles orientés, alors pour tout $J' \in \Cr(\tau')$, on a}
\[
(0 \leq)\ \card \Cr(f)^{-1}(\{J\}) \longrightarrow \lng(J'),
\] % [sic] \longrightarrow -> \leq, the next sentence speaks of equality; and {J} -> {J'}, as in Cr(f)^{-1}({J'}) below
\msp{page 37}
\textit{l'égalité n'étant obtenue que si les parties de $f$ sur les intervalles critiques de $I_\tau$ qui sont au dessus de $J'$ sont toutes égales à $0$ ou $1$ (donc si et seulement si sur les $J \in \Cr(f)^{-1}(\{J'\})$, elles sont égales à $1$).}

\textbf{Corollaire 4.} \textit{Supposons de plus $I'$} \texttt{[}$= I_{\tau'}$\texttt{]} \textit{réduit. Alors}
\[
\mathit{Cr}(f) : \mathit{Cr}(\tau, f) \longrightarrow \mathit{Cr}(\tau')\ (\simeq \mathrm{Fl}\,\mathrm{tr}(\tau'))
\]
\textit{est} bijective.

Donc dans ce cas, $\tau'$ réduit, on obtient une application en sens inverse
\[
\begin{tikzcd}[column sep=small]
\Tr(\tau') \simeq \Cr(\tau') \arrow[rr, "\Cr^*(f)"] \arrow[dr, "\Cr(f)^{-1}"'] & & \Cr(\tau) \\
 & \Cr(\tau, f), \arrow[ur, hook, "\text{inclusion}"'] &
\end{tikzcd}
\]
application qui transforme l'intervalle critique $J'$ $(= \{i - 1, i\})$ en l'\textit{unique} intervalle critique $J$ de $I_\tau$ tel que
\[
f(J) = J'.
\]

\msp{page 38}
Mais $J \simeq \bD_i^{\varepsilon}$, et on a donc un morphisme \textit{surjectif}
\[
f_0 : \bD_i^{\varepsilon} \longrightarrow J' \simeq \bD_1^{\varepsilon}.
\]
On voit tout de suite qu'il existe alors une \textit{unique} flèche de transition $u$ de $\bD_i^{\varepsilon}$, telle que $f(u) \neq \mathrm{id}$. Donc on trouve le corollaire suivant, plus précis que le corollaire 4 :

\textbf{Corollaire 5.} \textit{Soit $f : I_\tau \longrightarrow I_{\tau'}$ un morphisme croissant d'intervalles orientés, avec $I_{\tau'}$ réduit, et soit $\Tr(\tau, f)$ l'ensemble des flèches de transition $\underline{u}_i$ de $I_\tau$ telle que $f(\underline{u}_i) \neq \mathrm{id}$, i.e. telles que $\pi(u_i, f) = 1$. Alors $f$ induit une application} % [sic] telle -> telles, agreement with "flèches" (cf. "telles que" just after)
\[
\Tr(\tau, f) \longrightarrow \Tr(\tau')
\]
\textit{qui est bijective, et (bien sûr) croissante, donc un isomorphisme d'ensembles ordonnés.}

Ainsi l'application $\Tr(\tau') \longrightarrow \Cr(\tau, f)$ de tantôt se précise par une application bijective
\[
\Tr(f)^* : \Tr(\tau') \longrightarrow \Tr(\tau, f),
\]
reliée à $\Cr(f)^*$ plus haut par le triangle commutatif
\[
\begin{tikzcd}
 & \Tr(\tau, f) \arrow[dd, "\wr\ \alpha"] \\
\Tr(\tau') \arrow[ur, "\Tr(f)^*"] \arrow[dr, "\Cr(f)^*"'] & \\
 & \Cr(\tau, f),
\end{tikzcd}
\]
% foot: Version 22 décembre 2001
\printedpage{25}
% head: A. Grothendieck : Dérivateurs, Ch. VII, Édition des Universités de Montpellier II et Paris VII
\msp{page 39}
l'application $\alpha$ associant à toute $u_i \in \Tr(\tau, f)$ l'unique intervalle critique qui ‘contienne’ $u_i$ (i.e. tel que $u_i$ soit une flèche dudit intervalle, ou encore, fasse partie de la ‘chaîne critique’ associée \dots).

On a mieux encore :

\textbf{Proposition 8.} \textit{Soient $\tau, \tau' \in \mathrm{Ty}$, $\tau'$} réduit. \textit{Pour tout morphisme croissant d'intervalles orientés}
\[
f : I_\tau \longrightarrow I_{\tau'}
\]
\textit{(on va dire simplement que $f$ est une flèche dans $\uTy$, et écrire $f \in \Hom_{\uTy}(I_\tau, I'_\tau)$), associons lui la partie $\Tr(\tau, f)$ des flèches de transition de $I_\tau$ non dégénérées pour $f$. Ces flèches forment un ensemble totalement ordonné affecté de signes $+$ ou $-$ (caractère direct ou rétrograde des flèches), et $f$ induit une bijection d'ensembles totalement ordonnés} % [sic] I'_\tau -> I_{\tau'}, the target of f is I_{\tau'}
\[
\underbrace{E_f}_{\subseteq \Tr(\tau)} = \Tr(\tau, f) \simeq \Tr(\tau'),
\]
\textit{et même un isomorphisme d'ensembles totalement ordonnés ‘affectés de signes’. Ainsi, à toute $f$ on associe une partie $E_f$ de $\Tr(\tau)$, et telle que son ‘type’ soit égal à $\tau'$. Désignons par}
\[
\mathfrak{P}_{\tau'}(\Tr(\tau)) \subseteq \mathfrak{P}(\Tr(\tau))
\]
\textit{l'ensemble des parties $E$ de $\Tr(\tau)$ dont le type soit $\tau'$.}

\msp{page 40}
\textit{Alors l'application}
\[
f \longmapsto \Tr(\tau, f) : \Hom_{\uTy}(\tau, \tau') \longrightarrow \mathfrak{P}_{\tau'}(\Tr(\tau))
\]
\textit{est} bijective \textit{: se donner un morphisme $f : \tau \longrightarrow \tau'$ dans $\uTy$,} revient au même \textit{que de se donner une suite strictement croissante de flèches de transition}
\[
\underline{u}_{i_1},\ \underline{u}_{i_2},\ \dots\ ,\underline{u}_{i_{n'}} \qquad (n' = \lng \tau')
\]
\textit{dans $I_\tau$ $(i_1 < i_2 < \cdots < i_{n'})$, telle que la suite des signatures correspondante}
\[
(\varepsilon_{i_1}, \varepsilon_{i_2}, \dots, \varepsilon_{i_{n'}}) = \tau'
\]
\textit{soit égale à $\tau'$.}

\textbf{Corollaire 1.} \textit{Pour tout $\tau$, on a des bijections canoniques}
\[
\left\{
\begin{array}{lcl}
\Tr^+(\tau) & \simeq & \Hom_{\uTy}(\tau, \Delta^+) \\
\Tr^-(\tau) & \simeq & \Hom_{\uTy}(\tau, \Delta^-),
\end{array}
\right.
\]
fonctorielles en $\tau$ \textit{(variant dans} $\uTy$\textit{).}
% foot: Version 22 décembre 2001

\printedpage{26}
% head: A. Grothendieck : Dérivateurs, Ch. VII, Édition des Universités de Montpellier II et Paris VII

Pour donner un sens au corollaire, il faut définir
\[
\tau \longmapsto \Tr^+(\tau), \qquad \tau \longmapsto \Tr^-(\tau)
\]
comme \textit{foncteurs} contravariant en $\tau$, donc pour un morphisme
\[
f : I_{\tau_2} \longrightarrow I_{\tau_1}
\]
dans $\uTy$, définir des applications
\msp{page 41}
\[
\begin{array}{rcl}
\Tr^+(f)^* : \Tr^+(\tau_1) & \longrightarrow & \Tr^+(\tau_2) \\
\Tr^-(f)^* : \Tr^-(\tau_1) & \longrightarrow & \Tr^-(\tau_2),
\end{array}
\]
en d'autres termes, une application
\[
\Tr(f)^* : \Tr(\tau_1) \quad \longrightarrow \quad \Tr(\tau_2)
\]
‘compatible avec les signatures’. Dans le corollaire 5 plus haut, on a défini une telle application seulement pour $\tau'$ réduit \texttt{[}$\tau_1 = \tau'$\texttt{,} $\tau_2 = \tau$\texttt{]} (\footnote[12]{Ce qui suit ne fait pas intervenir les signatures - c'est un sorite sur les applications croissantes entre ensembles finis totalement ordonnés.}). Pour une flèche $u$ de $I_\tau$, appelons \textit{partie ensembliste} de $f$ sur $u$, où de $u$ pour $f$, la partie ensembliste de la flèche $f(u)$ dans $I_{\tau'}$, i.e. le segment $[f(i-1), f(i)] \subseteq I_{\tau'}$ (ici $u = \underline{u}_i$). Soit % [sic] où -> ou, alternative "de f sur u, ou de u pour f"
\[
\underline{u}_{i_1},\ \dots\ ,\underline{u}_{i_N} \qquad \text{ou } (v_1, v_2, \dots, v_n)
\] % [sic] v_n -> v_N, the sequence has N elements (cf. v_N below)
la suite croissante des éléments de $\Tr(\tau, f)$, on voit alors par récurrence sur $i$ que pour $i \in [1, N-1]$, on a
\[
\begin{array}{rcl}
f(\mathrm{but}(v_i)) & = & f(\mathrm{source}(v_{i+1}))\ \ (= t_i), \\
f(\mathrm{source}(v_1)) & = & 0_{\tau'}, \\
f(\mathrm{but}(v_N)) & = & 1_{\tau'}.
\end{array}
\]
Considérons la suite
\[
\{\overbrace{0_{\tau'}}^{= t_0}, t_1, t_2, \dots, t_{N-1}, \overbrace{1_{\tau'}}^{= t_N}\},
\]
elle est définie par le sous-ensemble correspondant de $I_{\tau'}$, qu'on appelera l'ensemble des \textit{sommets critiques de $I_{\tau'}$ pour $f$.} C'est une suite contenant $0_{\tau'}$, $1_{\tau'}$, et telle que deux éléments consécutifs sont reliés dans $I_{\tau'}$ par une flèche, ce qui équivaut à dire qu'elle contient les sommets critiques (intrinsèques) de $I_{\tau'}$. D'autre part, pour tout intervalle de transition $[i-1, i]$ de $I_{\tau'}$ $(i \in [1, n'])$, % [sic] appelera -> appellera, French spelling
% foot: Version 22 décembre 2001
\printedpage{27}
% head: A. Grothendieck : Dérivateurs, Ch. VII, Édition des Universités de Montpellier II et Paris VII
\msp{page 42}
il existe un et un seul intervalle de transition $[t_{j-1}, t_j]$ pour cette suite $(j \in [1, N])$, telle qu'on ait $[i-1, i] = [t_{j-1}, t_j]$. C'est dire qu'il existe une et une seule $v_j \in \Tr(\tau, f)$ telle que la partie de $f$ sur $v_j$ (savoir justement $[j-1, j]$) contienne (celle de) $\underline{u}_i$. Pour résumer :

\textbf{Corollaire 2.} \textit{Soit $f : \tau \longrightarrow \tau'$ dans $\uTy$. Alors pour toute $u \in \Tr(\tau') = \Tr(I_{\tau'})$, il existe une et une seule $v \in \Tr(\tau)$ telle que la partie de $f$ sur $v$ contienne $u$. De plus, les signatures de $u$ et de $v$ sont égales.}

C'est ainsi qu'on obtient donc une application
\[
\Tr(f)^* : \Tr(\tau') \quad \longrightarrow \quad \Tr(\tau) \qquad (\text{se factorisant par } \Tr(\tau, f))
\]
induisant
\[
\begin{array}{rcl}
\Tr^+(f)^* : \Tr^+(\tau') & \longrightarrow & \Tr^+(\tau) \\
\Tr^-(f)^* : \Tr^-(\tau') & \longrightarrow & \Tr^-(\tau).
\end{array}
\]
La transitivité pour un composé
\[
\tau \xrightarrow{f} \tau' \xrightarrow{g} \tau''
\]
ou
\[
I_\tau \xrightarrow{f} I_{\tau'} \xrightarrow{g} I_{\tau''}
\]
est immédiate. Si $u$, $u'$, $u''$ sont des flèches de transition dans $I_\tau$, $I_{\tau'}$, $I_{\tau''}$ respectivement, telles que
\[
\begin{array}{rcll}
\pi(g, u') & \supseteq & \pi(u''), & \qquad \text{i.e. } u' = \Tr(g)^*(u'') \\
\pi(f, u) & \supseteq & \pi(u'), & \qquad \text{i.e. } u = \Tr(f)^*(u'),
\end{array}
\]
alors on a
\[
\pi(gf, u) \ \supseteq\ \pi(u''), \qquad \text{i.e. } u = \Tr(gf)^*(u'').
\]

\msp{page 43}
Si les flèches en question ont comme partie intrinsèque $[i-1, i]$, $[i'-1, i']$, $[i''-1, i'']$ respectivement, l'hypothèse est
\[
\begin{array}{rclrcl}
g(i'-1) & \leq & i''-1, & g(i') & \geq & i'' \\
f(i-1) & \leq & i'-1, & f(i) & \geq & i'
\end{array}
\]
et comme $g$ est croissant, on a déduit
\[
\begin{array}{rcll}
gf(i-1) & \leq & g(i'-1) & (\leq i''-1) \\
gf(i) & \geq & g(i') & (\geq i''),
\end{array}
\]
\begin{flushright}
c.q.f.d.
\end{flushright}
% foot: Version 22 décembre 2001

\printedpage{28}
% head: A. Grothendieck : Dérivateurs, Ch. VII, Édition des Universités de Montpellier II et Paris VII

Ainsi, le corollaire 1 plus haut obtient un sens précis et il est démontré du même coup, vu que les correspondances du corollaire 1 sont données par
\[
\underbrace{f}_{\in \Hom_{\uTy}(\tau, \Delta^{\varepsilon})} \qquad \longmapsto \Tr(f)^*(\underline{u}),
\]
où $\underline{u}$ est l'unique flèche de transition de $\Delta^{\varepsilon}$.

On va utiliser ces sorites pour donner, pour tout $\tau \in \uTy$, une description combinatoire simple des catégories $\uTy/\tau$ et $\tau\backslash\uTy$, ou du moins de leurs ensembles d'objets.

\textbf{1$^o$ Description de} $\Ob\,(\tau\backslash\uTy)$\textbf{.}

À $f \in \Ob(\tau\backslash\uTy)$, i.e.
\[
f : \tau \longrightarrow \tau' \qquad \text{i.e. } f : I_\tau \longrightarrow I_{\tau'}
\]
dans $\uTy$, est associé $E = E_f = \Tr(\tau, f) \subseteq \Tr(\tau)$, l'ensemble des flèches de transition critiques (i.e. non dégénérées) pour $f$. À chaque $u \in E$, associons l'entier $\lng(\pi(f, u))$,
\msp{page 44}
la \textit{partie} de $f$ sur $u$. Alors $(\tau', f)$ est déterminé de façon unique par le couple
\[
(E, \pi), \qquad E \in \mathfrak{P}(\Tr(\tau)), \qquad \underbrace{\pi}_{\text{pondération}} : E \ \longrightarrow\ \underbrace{\mathbf{N}^*}_{\substack{= \text{ensemble des} \\ \text{entiers} \geq 1}} \quad .
\]
Ou si on préfère, par l'application $\pi = \pi_f$
\[
\Tr(\tau) \xrightarrow{\pi} \mathbf{N}, \qquad u \longmapsto \pi(u) = \pi(u, f) = \lng(\pi_{\tau'}(f(u)))
\]
(qui redonne $E$ comme
\[
E = \pi^{-1}(\mathbf{N}^*) \qquad ).
\]
Cette application n'est soumise à aucune restriction,
\[
\Ob\,(\tau\backslash\uTy) \ \xrightarrow[\sim]{\text{bij.}}\ \Hom_{\mathrm{Ens}}(\Tr(\tau), \mathbf{N}).
\]
On récupère $\tau'$ à partir de $\pi$ ainsi. On prend $E = \pi^{-1}(\mathbf{N}^*)$ (flèches de transition non dégénérées pour $\pi$), $E$ est totalement ordonné par l'ordre induit par celui de $\Tr(\tau)$, il est de plus signaturé \texttt{[sic]} par la signaturation \texttt{[sic]} de $\Tr(\tau)$, soit
\[
(\varepsilon_1, n_1),\ \dots\ ,\ (\varepsilon_r, n_r)
\]
la suite des signatures et pondérations, de la suite des éléments de $E$ $(\varepsilon_i \in \{\pm 1\}$, $n_i \in \mathbf{N}^*$ pour $i \in [1, r])$. Alors,
\[
\tau' = \Delta_{\varepsilon_1}^{n_1} \Delta_{\varepsilon_2}^{n_2} \cdots \Delta_{\varepsilon_r}^{n_r}, \qquad r = \card E,
\]
dans le monoïde libre $\mathrm{Ty}$, où on a désigné par $\Delta_+$, $\Delta_-$ (au lieu de $\Delta^+$, $\Delta^-$) les deux générateurs. L'application
\[
I_\tau \ \longrightarrow\ I_{\tau'}
\]
% foot: Version 22 décembre 2001


\endgroup
% END UNIT Der-ChVII-456/Der-ChVII-456_p026-030.tex
% BEGIN UNIT Der-ChVII-456/Der-ChVII-456_p031-035.tex
\begingroup\makeatletter%
\newcommand{\printedpage}[1]{\marginpar{\footnotesize p.~#1}}%
\newcommand{\msp}[1]{\marginpar{\footnotesize [#1]}}%
\newcommand{\Ty}{\mathrm{Ty}}%
\newcommand{\uTy}{\underline{\mathrm{Ty}}}%
\newcommand{\uTyp}{\underline{\mathrm{Typ}}}%
\newcommand{\uCh}{\underline{\mathrm{Ch}}}%
\newcommand{\ovalbox}[1]{\tikz[baseline=(X.base)]\node[draw,rounded corners=6pt,inner sep=2pt](X){#1};}%
\makeatother\let\printedpage\derpp\let\msp\dermsp%
\let\mspin\DERmspin%


\printedpage{29}
% head: A. Grothendieck : Dérivateurs, Ch. VII, Édition des Universités de Montpellier II et Paris VII
\msp{page 45}est décrite de façon unique en disant qu'elle est \textit{constante} sur chacune des `composantes connexes' (au sens de la numérotation des sommets de $I_\tau$) de l'ensemble $E' = E'(f) \subseteq I_\tau = [0,n]$ des sommets qui sont sources ou buts de flèches de transition dégénérées pour $f$, i.e. dans $\mathrm{Tr}(\tau)\backslash E$. Et sur le $\alpha$-ème élément de $E$ ($\alpha \in [1, r = \operatorname{card} E]$), d'indice $i_\alpha$ dans $\mathrm{Tr}(\tau)$, $f$ induit le foncteur
\[
\boldsymbol{\Delta}_{\varepsilon_\alpha} \simeq \underbrace{J_{i_\alpha}}_{\substack{\text{intervalle}\\ \text{de transition}\\ \text{d'indice } i_\alpha\\ \text{dans } I_\tau}} \longrightarrow I_{\tau'}
\]
donné par \textit{l'unique} \underline{Ty}-\textit{flèche} dans $I_\tau$ de $J_{i_\alpha}$ dans le $\alpha$-ème facteur, de type $\Delta^{n_\alpha}_{\varepsilon_\alpha}$, de $\tau'$, donné par $(*)$
\[
\Delta^{n_\alpha}_{\varepsilon_\alpha} \simeq \text{type de } \boldsymbol{\Delta}^{\varepsilon_\alpha}_{n_\alpha}\ \ \ldots
\]

\subsection*{2$^o$ Détermination de $\mathrm{Ob}\,(\uTy/\tau')$.}

Si $(\tau, f) \in \mathrm{Ob}\,(\uTy/\tau')$, donc
\[
f : I_\tau \longrightarrow I_{\tau'},
\]
on lui associe tout d'abord la partie
\[
\mathrm{Cr}(\tau', f) = E'_f = E' \subseteq I_{\tau'}
\]
formée des sommets (dits `critiques pour $f$', cf. page 41) de $I_{\tau'}$ qui sont source ou but d'une flèche $f(u)$, où $u \in \mathrm{Tr}(\tau, f)$, i.e. $f(u)$ n'est pas dégénérée (\footnote[13]{dégager un lemme $f$ croissante inutile.}). $E'$ contient $f(I_{\tau\,\text{réd}})$ (i.e. les images par $f$ des sommets critiques de $I_\tau$),
\msp{page 46}mais peut être plus grand. (P.ex. si $\tau = \tau'$, $f = \mathrm{id}$, $E' = I_{\tau'}$, qui est donc strictement plus grand que $I_{\tau'\,\text{réd}}$ si $\tau'$ n'est pas réduit.) Pour tout $i \in E'$, considérons $f^{-1}(\{i\}) \subseteq I_\tau$. Comme $f$ \texttt{[est]} croissant, c'est un intervalle de l'ensemble totalement ordonné $I_\tau$, donc c'est un sous-intervalle de la catégorie intervalle $I_\tau$. Comme tel, il a une signature $\tau_i \in \Ty$. On a donc une application
\[
\eta_f \text{ ou } \eta : E' \longrightarrow \Ty.
\]
Je dis que $(\tau, f) \in \mathrm{Ob}\,(\uTy/\tau')$ est déterminé de façon unique par ce couple
\[
(E', \eta) \begin{cases} E' \in \mathfrak{P}(I_{\tau'}) = \mathfrak{P}([0,n']) \\[1ex] \eta : E' \longrightarrow \Ty\ . \end{cases}
\]
De plus, celui-ci est soumis à la seule restriction que
\[
E' \supseteq \mathrm{Cr}(\tau') = \text{ensemble des sommets critiques de } I_{\tau'}
\]
\printedpage{30}
% head: A. Grothendieck : Dérivateurs, Ch. VII, Édition des Universités de Montpellier II et Paris VII
(où ce qui revient au même, que $E'$ contient $0_{\tau'}$, $1_{\tau'}$ et que deux éléments consécutifs de $E'$ sont reliés par une flèche de $I_{\tau'}$) : % [sic] où -> ou, « ou ce qui revient au même »
\[
\left\{
\begin{array}{lcl}
\mathrm{Ob}\,(\Ty/\tau') & \simeq & \text{ensemble des couples } (E',\eta), \text{ avec} \\[1ex]
 & & E' \subseteq [0, \underbrace{\ n'\ }_{=\,\mathrm{long}\,\tau'}], \quad E' \supseteq \mathrm{Cr}(\tau'), \\[1ex]
 & & \eta : E' \longrightarrow \Ty\ .
\end{array}
\right.
\]
Soit
\[
0 = t_0 < t_1 < \ \cdots\ < t_r = n'
\]
la suite des éléments de $E'$, soit pour \ovalbox{$i \in [1,N]$}
\[
\left\|
\begin{array}{l}
l_i = t_i - t_{i-1} \\[1ex]
\varepsilon_i \quad \text{signature de } [t_{i-1}, t_i] \text{ dans } I_{\tau'}\ ,
\end{array}
\right.
\]
\msp{page 47}et pour \ovalbox{$i \in [0,N]$}\,, soit
\[
\|\ \ \tau_i = \eta(t_i).
\]
Alors le type $\tau$ est donné par
\[
\tau = \tau_0 \Delta_{\varepsilon_1} \tau_1 \Delta_{\varepsilon_2} \cdots \Delta_{\varepsilon_N} \tau_N, \tag*{$(*)$}
\]
\texttt{[de]} longueur $N + \sum_{i=0}^N \underbrace{\mathrm{long}(\tau_i)}_{\lambda_i \geq 0}$.

\subsection*{3$^o$ Description de $\mathrm{Fl}(\uTy)$.}

Si
\[
f : \tau \longrightarrow \tau' \qquad \texttt{[est]} \ \text{dans } \uTy\ ,
\]
on lui associe comme tantôt un entier \underline{$N \geq 0$} ($N = \operatorname{card}\mathrm{Cr}(\tau', f){-}1$, $N = 0$ si et seulement si $\tau' = 1$, type de l'intervalle ponctuel), des applications
\[
\left\{
\begin{array}{lcl@{\qquad\qquad}ccl}
{[1,N]} & \longrightarrow & [+1,-1] \times \mathbf{N}^*, & i & \longmapsto & (\varepsilon_i, l_i) \\[1ex]
{[0,N]} & \longrightarrow & \Ty & i & \longmapsto & \tau_i
\end{array}
\right.
\]
(\footnote[14]{\textbf{NB} Si $N = 0$, $[1,N] = \emptyset$, et il n'y a pas de $\varepsilon_i$, $l_i$; l'application $\underbrace{[0,N]}_{=\,\{0\}} \longrightarrow \Ty$ se réduit à la donnée de $\tau_0 \in \Ty$.}). La flèche $f$ dans $\uTy$ est determinée par ces données. Le type $\tau$ est donné par la formule $(*)$ ci-dessus, le type $\tau'$ par % [sic] determinée -> déterminée, accent manquant
\[
\tau' = \Delta^{l_1}_{\varepsilon_1} \Delta^{l_2}_{\varepsilon_2}\ \cdots\ \Delta^{l_N}_{\varepsilon_N}, \qquad \text{longueur } n' = \textstyle\sum l_i, \tag*{$(*')$}
\]
\printedpage{31}
% head: A. Grothendieck : Dérivateurs, Ch. VII, Édition des Universités de Montpellier II et Paris VII
et l'application croissante de $I_\tau = [0, N + \sum \mathrm{long}(\tau_i)]$ dans $I_{\tau'} = [0, \sum l_i]$ est donnée ainsi :
\[
\begin{array}{l@{\qquad}cl}
\text{constante sur } [0, 0+\lambda_0] & \longmapsto & 0 \\[1ex]
\text{constante sur } [\lambda_0+1, \lambda_0+1+\lambda_1] & \longmapsto & l_1 \\[1ex]
\text{constante sur } [\lambda_0+\lambda_1+2, \lambda_0+\lambda_1+2+\lambda_2] & \longmapsto & l_2 \\[1ex]
\ldots & & \\[1ex]
\text{constante sur } [\lambda_0+\lambda_1+\cdots+\lambda_i+(i+1), \lambda_0+\cdots+\lambda_i+\lambda_{i+1}+(i+1)] & \longmapsto & l_i.
\end{array}
\]

\printedpage{32}
% head: A. Grothendieck : Dérivateurs, Ch. VII, Édition des Universités de Montpellier II et Paris VII
\msp{page 48}

\section*{6\quad Composition dans $\uTy$}

On a une loi de composition (partout définie) dans $\mathrm{Ob}\,(\uTy) = \Ty$, c'est celle du monoïde libre à deux générateurs $\Delta_+$, $\Delta_-$. Cette loi de composition s'interprète, au niveau catégorique (et non plus seulement combinatoire) comme une somme amalgamée
\[
I_{\tau\tau'} = I_\tau \amalg_{(e, \underline{1}_\tau, \underline{0}_{\tau'})} I_{\tau'},
\]
où $e$ est la catégorie unité, et où $\underline{1}_\tau$, $\underline{0}_{\tau'}$ représentent maintenant les morphismes $e \longrightarrow I_\tau$, $e \longrightarrow I_{\tau'}$ , ayant comme image l'élément but et l'élément source respectivement de l'intervalle orienté envisagé. On a donc un carré cocartésien dans (Cat)
\[
\begin{tikzcd}[column sep=small, row sep=small]
 & e \arrow[dl, hook', "\underline{1}_\tau"'] \arrow[dr, hook, "\underline{0}_{\tau'}"] & \\
I_\tau \arrow[dr, hook, "i_{\tau,\tau'}"'] & \text{\scriptsize cocart.} & I_{\tau'} \arrow[dl, hook', "j_{\tau,\tau'}"] \\
 & I_{\tau\tau'}, &
\end{tikzcd}
\]
où toutes les flèches sont des foncteurs mono et pleinement fidèles. Que ce diagramme (évidemment commutatif) soit bien cocartésien dans (Cat) revient, par définition, à dire que pour toute $X$ dans Cat, le diagramme transposé
\[
\begin{tikzcd}[column sep=small, row sep=small]
 & X & \\
\uCh_\tau(X) \arrow[ur, "\beta_\tau" near end, "\text{\scriptsize (but)}" near start] & & \uCh_{\tau'}(X) \arrow[ul, "\sigma_{\tau'}"' near end, "\text{\scriptsize (source)}"' near start] \\
 & \uCh_{\tau\tau'}(X), \arrow[ul] \arrow[ur] &
\end{tikzcd}
\]

\texttt{[est cartésien]}

\msp{page 49}ce qui revient à la tautologie suivante qu'un chemin $C$ dans $X$ de type $\tau\tau'$ `n'est autre chose' qu'un couple $(c,c')$ d'un chemin $c$ de type $\tau$ et d'un chemin $c'$ de type $\tau'$ `mis bout à bout' (dans cet ordre), i.e. tels que
\[
\mathrm{but}(c) = \mathrm{source}(c')
\]
(et on a alors
\[
\begin{array}{rcll}
\mathrm{source}(C) & = & \mathrm{source}(c) & \\[1ex]
\mathrm{but}(C) & = & \mathrm{but}(c') & \qquad\quad ).
\end{array}
\]

\printedpage{33}
% head: A. Grothendieck : Dérivateurs, Ch. VII, Édition des Universités de Montpellier II et Paris VII
Il se pose ici le sempiternel problème si on va noter ce chemin $c{\cdot}c'$, ce qui semble à première vue s'imposer, ou plutôt $c' \circ c$, regardant le chemin $c$ de $x = \mathrm{source}(c)$ vers $y = \mathrm{but}(c)$ comme une sorte de morphisme de $x$ dans $y$ (cf. plus bas le formalisme des catégories de fractions), $c'$ comme un morphisme de $y = \mathrm{source}(c')$ vers $z = \mathrm{but}(c')$, et en interprétant le composé $C$ comme le composé de ces deux morphismes $c$, $c'$ (dans cet ordre du point de vue des opérations à effectuer), qu'on note traditionellement $c' \circ c$ ou $c' \cdot c$, et non $c \circ c'$ ou $c \cdot c'$, hélas! C'est cette notation que nous allons suivre plus bas. % [sic] traditionellement -> traditionnellement, orthographe

Mais pour l'instant on ne va pas se soucier des chemins dans des catégories quelconques $X$, mais revenir au formalisme des $\tau$ et des $I_\tau$ eux-mêmes. L'interprétation catégorique de $\tau\tau'$ comme une somme amalgamée nous assure que cette opération de composition est fonctorielle :

\msp{page 50}
\[
\tau_1 \xrightarrow{\ f\ } \tau_2, \quad \tau'_1 \xrightarrow{\ f'\ } \tau'_2 \qquad\quad \text{donne } \tau_1\tau'_1 \xrightarrow{f*f'} \tau_2\tau'_2
\]
(on n'écrira pas $f \cdot f'$ pour la flèche notée $f * f'$, à cause de la confusion possible avec la composition des flèches dans $\uTy$). On a transitivité de façon évidente : si de plus $\tau_2 \xrightarrow{\ g\ } \tau_3$, $\tau'_2 \xrightarrow{\ g'\ } \tau'_3$, \texttt{[on a]}
\[
(gf) * (g'f') = (g * g') \circ (f * f').
\]
Donc on a un bifoncteur de composition
\[
\uTy \times \uTy \longrightarrow \uTy \qquad\quad (c,c') \longmapsto c \cdot c' \text{ ou } c * c'.
\]
(En fait, on a aussi un bifoncteur
\[
\uTyp \times \uTyp \longrightarrow \uTyp,
\]
mais on préfère oublier maintenant la catégorie \underline{Typ}, qui donne un formalisme nettement moins simple.)

Cette composition est \textit{strictement} associative et \textit{strictement unitaire} (l'unité étant bien sûr l'élément unité 1 du monoïde libre Ty, correspondant à la catégorie ponctuelle $\{0\}$), et pas seulement `à isomorphisme près'. Et pour cause, puisque dans $\uTy$ les isomorphismes sont des identités. \textit{Donc} $\uTy$ \textit{est un monoïde dans la catégorie} (Cat) des petites catégories!

\msp{page 49$'$ (il y a un décalage dans la numérotation)}Cela implique que non seulement $\mathrm{Ob}\,\uTy = \Ty$, mais aussi $\mathrm{Fl}(\uTy)$, et plus généralement tous les ensembles de la forme $\mathrm{Hom}_{\mathrm{Cat}}(X, \uTy)$ sont des monoïdes, et ceci fonctoriellement en $X$, par exemple les composantes $\mathrm{Hom}_{\mathrm{Cat}}(\boldsymbol{\Delta}_i, \uTy)$ de l'ensemble semi-simplicial `nerf' de $\uTy$,
\[
\mathrm{Nerf}_i(\uTy) = \mathrm{Hom}_{\mathrm{Cat}}(\boldsymbol{\Delta}_i, \uTy),
\]
sont des monoïdes, et les applications semi-simpliciales entre elles sont des homomorphismes de monoïdes : le nerf de $\uTy$ est un \textit{monoïde semi-simplicial.}


\endgroup
% END UNIT Der-ChVII-456/Der-ChVII-456_p031-035.tex
% BEGIN UNIT Der-ChVII-456/Der-ChVII-456_p036-040.tex
\begingroup\makeatletter%
\newcommand{\printedpage}[1]{\marginpar{\footnotesize p.~#1}}%
\newcommand{\msp}[1]{\marginpar{\footnotesize [#1]}}%
\newcommand{\uTy}{\underline{\mathrm{Ty}}}%
\makeatother\let\printedpage\derpp\let\msp\dermsp%
\let\mspin\DERmspin%


\printedpage{34}
% head: A. Grothendieck : Dérivateurs, Ch. VII, Édition des Universités de Montpellier II et Paris VII
On va décrire ces monoïdes $\mathrm{Nerf}_i(\uTy)$, et constater que ce sont tous des monoïdes \textit{libres,} tout comme $\mathrm{Nerf}_0(\uTy) = \mathrm{Ob}\,\uTy = \mathrm{Ty}$. Les applications semi-simpliciales entre ces monoïdes sont donc connus quand on les connaît sur les \textit{générateurs} de ces monoïdes libres (lesquels générateurs sont aussi les éléments indécomposables, donc définis intrinsèquement en termes de la structure multiplicative). % [sic] connus -> connues, agreement with « applications »

Pour $\mathrm{Nerf}_0(\uTy) = \mathrm{Ty}$, c'est une chose vue, les générateurs sont
\[
\Delta_+,\Delta_-,
\]
correspondant respectivement aux intervalles orientés
\[
\begin{array}{rccc}
\boldsymbol{\Delta}_+ : & 0 & \longrightarrow & 1\\[4pt]
\boldsymbol{\Delta}_- : & 0 & \longleftarrow & 1.
\end{array}
\]

\msp{page 50$'$}
Pour $\mathrm{Nerf}_1(\uTy) = \mathrm{Fl}(\uTy)$, i.e. pour les flèches, le résultat $3^o$) du paragraphe précédent (détermination de $\mathrm{Fl}(\uTy)$) s'interprète aussi en disant que $\mathrm{Fl}(\uTy)$ est le monoïde libre engendré par un système de générateurs infini
\[
\mathfrak{f}_{\varepsilon,n}\ , \qquad\quad (\varepsilon,n)\in\{+1,-1\}\times\mathbf{N},
\]
indexé par l'ensemble produit $\{+1,-1\}\times\mathbf{N}$. Ici $\mathfrak{f}_{\varepsilon,n}$ désigne l'unique flèche dans $\uTy$
\[
\mathfrak{f}_{\varepsilon,l} : \Delta_\varepsilon \longrightarrow \Delta^l_\varepsilon
\]
(où bien sûr l'exposant $l$ désigne une exponentation pour la loi de composition dans $\uTy$
\[
\Delta^l_\varepsilon = \underbrace{\Delta_\varepsilon\circ\Delta_\varepsilon\circ\ \cdots\ \circ\Delta_\varepsilon}_{l\text{ copies}},
\]
et non la puissance cartésienne dans Cat, où dans $\uTy$ (où elle n'existe que si $l\in\{0,1\}$)). % [sic] où dans -> ou dans, conjunction intended
Dans la description de $\mathrm{Fl}(\uTy)$ dans le paragraphe précédent, on a écrit le `terme général' du monoïde libre en groupant ensemble les groupes de foncteurs successifs de la forme $f_{+,0}$, $f_{-,0}$, i.e. corréspondant aux deux flèches de dégénérescence indécomposables (pour $*$) % [sic] corréspondant -> correspondant, accent
\[
\left\{
\begin{array}{lccc}
f_{+,0} : & \Delta_+ & \longrightarrow & e\\[4pt]
f_{-,0} : & \Delta_- & \longrightarrow & e
\end{array}
\right.
\]
(\footnote[15]{\textbf{NB} Le `mot' vide de $\mathrm{Fl}(\uTy)$ correspond à la flèche $\mathrm{id}_1$, où $1$ est l'élément unité de Ty, réalisant l'`intervalle' ponctuel.}). Le sous-monoïde libre de $\mathrm{Fl}(\uTy)$ engendré par ces deux éléments est le monoïde des flèches de dégénérescence (pure) quelconques
% foot: Version 22 décembre 2001
\printedpage{35}
% head: A. Grothendieck : Dérivateurs, Ch. VII, Édition des Universités de Montpellier II et Paris VII
\msp{page 51}
dans $\uTy$ :
\[
f_\tau : I_\tau \longrightarrow e \qquad\quad (\tau\in\mathrm{Ty}),
\]
et on a
\[
f_\tau * f_{\tau'} = f_{\tau\tau'},
\]
i.e. l'application
\[
\tau \longmapsto f_\tau :\quad \underbrace{\mathrm{Ob}\,(\uTy)}_{=\,\mathrm{Ty}\,=\,\mathrm{Nerf}_0(\uTy)} \quad\longrightarrow\quad \underbrace{\mathrm{Fl}(\uTy)}_{\mathrm{Nerf}_1(\uTy)}
\]
est un homomorphisme de monoïdes, transformant $\Delta_\varepsilon$ ($\varepsilon\in\{+1,-1\}$) dans $\mathfrak{f}_{\varepsilon,0}$. Ce n'est pas l'application de dégénérenscence de $\mathrm{Nerf}_0$ dans $\mathrm{Nerf}_1$, i.e. l'application $\tau\longmapsto\mathrm{id}_\tau$. Celle-ci est donnée par % [sic] dégénérenscence -> dégénérescence, spelling
\[
\begin{array}{ccc}
\mathrm{Ob}\,(\uTy) & \longrightarrow & \mathrm{Fl}(\uTy)\\[4pt]
\tau & \longmapsto & \mathrm{id}_\tau
\end{array}
\ :\
\left|\
\begin{array}{rcccl}
 & \Delta_\varepsilon & \longmapsto & \mathfrak{f}_{\varepsilon,1} & \qquad(\varepsilon\in\{+1,-1\})\\[4pt]
\text{i.e.} & \Delta_+ & \longmapsto & \mathfrak{f}_{+,1} & \\[4pt]
 & \Delta_- & \longmapsto & \mathfrak{f}_{-,1}\ . &
\end{array}
\right.
\]
Quant aux homomorphismes source et but
\[
\begin{tikzcd}[column sep=large]
\mathrm{Fl}(\uTy) \arrow[r, shift left, "\sigma\text{ (source)}"] \arrow[r, shift right, "\beta\text{ (but)}"'] & \mathrm{Ob}\,(\uTy),
\end{tikzcd}
\]
ils sont évidents sur les générateurs
\[
\left\{
\begin{array}{lcl}
\sigma(\mathfrak{f}_{\varepsilon,l}) & = & \Delta_\varepsilon\\[4pt]
\beta(\mathfrak{f}_{\varepsilon,l}) & = & \Delta^l_\varepsilon.
\end{array}
\right.
\]

Je vais décrire maintenant $\mathrm{Nerf}_2(\uTy)$, i.e. les systèmes
\[
\tau \xrightarrow{\ f\ } \tau' \xrightarrow{\ g\ } \tau''.
\]

\msp{page 52}
Se donner $f\in\mathrm{Fl}(\uTy)$ revient à se donner une suite ordonnée de générateurs $f_{\varepsilon,l}$, donc une suite
\[
\begin{aligned}
&\mathfrak{f}_{\varepsilon_1,l_1},\ \mathfrak{f}_{\varepsilon_2,l_2},\ \ \dots\ \ ,\mathfrak{f}_{\varepsilon_n,l_n}\\
&((\varepsilon_i,l_i))_{1\leq i\leq n}\in[\{\pm1\}\times\mathbf{N}]^n.
\end{aligned}
\]
On a alors
\[
\begin{aligned}
x &= \Delta_{\varepsilon_1}\Delta_{\varepsilon_2}\cdots\Delta_{\varepsilon_n} &\qquad\qquad& \mathrm{long}\,x = n\\
y &= \Delta^{l_1}_{\varepsilon_1}\Delta^{l_2}_{\varepsilon_2}\cdots\Delta^{l_n}_{\varepsilon_n} && \mathrm{long}\,y = \textstyle\sum_{i\in[1,n]} l_i.
\end{aligned}
\]
% foot: Version 22 décembre 2001

\printedpage{36}
% head: A. Grothendieck : Dérivateurs, Ch. VII, Édition des Universités de Montpellier II et Paris VII
Se donner $q : y\longrightarrow z$, i.e. un objet de $y\backslash\uTy$, revient à se donner une application de l'ensemble des $m=\sum l_i$ flèches de transition de $y$ dans $\mathbf{N}$. En fait, on a % [sic] q -> g, the map is called g below (« la donnée de g »)
\[
\mathrm{Tr}(y) = \mathrm{Tr}(\Delta^{l_1}_{\varepsilon_1})\amalg\mathrm{Tr}(\Delta^{l_2}_{\varepsilon_2})\amalg\ \cdots\ \amalg\mathrm{Tr}(\Delta^{l_n}_{\varepsilon_n})
\]
(`ensemble somme', i.e. réunion disjointe), donc
\[
\mathrm{Hom}_{\mathrm{Ens}}(\mathrm{Tr}(y)) \simeq \prod_{i\in[1,n]}\mathrm{Hom}_{\mathrm{Ens}}(\underbrace{\mathrm{Tr}(\Delta^{l_i}_{\varepsilon_i})}_{[1,l_i]},\mathbf{N}), % [sic] Hom_Ens(Tr(y)) -> Hom_Ens(Tr(y),N), second argument missing
\]
donc la donnée de $g$ équivaut à la donnée d'un système de $n$ applications
\[
\left\{
\begin{array}{l}
\pi_i : [1,l_i]\longrightarrow\mathbf{N} \qquad\quad (\text{pondération de }\Delta^{l_i}_{\varepsilon_i})\\[6pt]
\text{pour } i\in[1,n].
\end{array}
\right.
\]
Mais pour $i\in[1,n]$, considérons l'élément de $\mathrm{Nerf}_2$
\[
\Delta_{\varepsilon_1}\xrightarrow{\mathfrak{f}_{\varepsilon_i,l_i}}\Delta^{l_i}_{\varepsilon_i}\xrightarrow{g_{\varepsilon,\pi_i}}\underbrace{\Delta^{\pi_i(1)}_{\varepsilon_i}\Delta^{\pi_i(2)}_{\varepsilon_i}\cdots\Delta^{\pi_i(l_i)}_{\varepsilon_i}}_{\Delta_{\varepsilon_i}^{\Sigma_j\pi_i(j)\overset{\text{déf}}{=}\mathrm{long}(\pi_i)}}, % [sic] \Delta_{\varepsilon_1} -> \Delta_{\varepsilon_i}, and g_{\varepsilon,\pi_i} -> g_{\varepsilon_i,\pi_i}, index i intended
\tag*{(1)}
\]

\msp{page 53}
où $g_{\varepsilon_i,\pi_i}$ est la flèche définie en termes de $\pi_i$ au paragraphe précédent. Mais ce n'est autre aussi que
\[
g_{\varepsilon_i,\pi_i} = \mathfrak{f}_{\varepsilon_i,\pi_i(1)}*\mathfrak{f}_{\varepsilon_i,\pi_i(2)}*\cdots*\mathfrak{f}_{\varepsilon_i,\pi_i(l_i)}.
\tag*{(2)}
\]
L'élément de $\mathrm{Nerf}_2$ envisagé ne dépend que de
\[
(\varepsilon_i,l_i,\pi_i), \qquad\quad \text{où } (\varepsilon_i,l_i)\in\{\pm1\}\times\mathbf{N},\ \pi_i : [1,l_i]\longrightarrow\mathbf{N}.
\]
(\textbf{NB} Si $l_i=0$, $[1,l_i]=[1,0]=\emptyset$ par convention, et $\pi_i$ est l'unique application $\emptyset\longrightarrow\mathbf{N}$, donc la donnée $(\varepsilon_i,l_i,\pi_i)$ revient alors à $(\varepsilon_i,l_i)$.)

Si
\[
(\varepsilon,l,\pi)\ \text{ comme dessus,} \qquad\quad (\varepsilon,l)\in\{\pm1\}\times\mathbf{N},\ \pi : [1,l]\longrightarrow\mathbf{N},
\]
désignons par
\[
\mathfrak{f}_{\varepsilon,l,\pi}
\]
l'élément de $\mathrm{Nerf}_2$ donné par
\[
\Delta_\varepsilon\xrightarrow{\mathfrak{f}_{\varepsilon,l}}\Delta^l_\varepsilon\ \xrightarrow{\quad\mathfrak{f}_{\varepsilon,\pi(1)}\ \ *\ \cdots\ *\ \ \mathfrak{f}_{\varepsilon,\pi(l)}\quad}\Delta^{\pi(1)}_\varepsilon\Delta^{\pi(2)}_\varepsilon\cdots\Delta^{\pi(l)}_\varepsilon.
\]
On peut formuler la construction précédente en disant que tout élément de $\mathrm{Nerf}_2$ peut se décomposer de façon unique en un produit
\[
(f,g) = \mathfrak{f}_{\varepsilon_1,l_1,\pi_1}*\mathfrak{f}_{\varepsilon_2,l_2,\pi_2}*\cdots*\mathfrak{f}_{\varepsilon_n,l_n\pi_n}, % [sic] l_n\pi_n -> l_n,\pi_n, comma missing in index
\]
ce qui prouve que $\mathrm{Nerf}_2(\uTy)$ est le monoïde libre sur les générateurs précédents $\mathfrak{f}_{\varepsilon,l,\pi}$. (\textbf{NB} à la suite vide correspond
\[
1\xrightarrow{\mathrm{id}_1}1\xrightarrow{\mathrm{id}_1}1 \qquad\qquad .)
\]
% foot: Version 22 décembre 2001

\printedpage{37}
% head: A. Grothendieck : Dérivateurs, Ch. VII, Édition des Universités de Montpellier II et Paris VII
\msp{page 54}
Si l'objet de $\mathrm{Nerf}_2$ est
\[
x\xrightarrow{\ f\ }y\xrightarrow{\ g\ }z,
\]
la construction précédente nous dit, en termes de la donnée combinatoire, i.e. de la suite de données
\[
(\varepsilon_1,l_1,\pi_1),\ \dots\ ,\ (\varepsilon_n,l_n,\pi_n),
\]
quels sont $x$, $y$, $z$ : chacun d'eux est le produit dans $\mathrm{Nerf}_0=\mathrm{Ty}$ des $n$ facteurs $x_i$, $y_i$ ou $z_i$ respectivement, pour $i\in[1,n]$, donnés dans la définition des $\mathfrak{f}_{\varepsilon_i,l_i,\pi_i}$, i.e. de
\[
\Delta_{\varepsilon_i},\ \Delta^{l_i}_{\varepsilon_i},\ \text{ et }\Delta^{\pi_i(1)}_{\varepsilon_i}\Delta^{\pi_i(2)}_{\varepsilon_i}\cdots\Delta^{\pi_i(l_i)}_{\varepsilon_i}=\Delta^{\mathrm{long}(\pi_i)}_{\varepsilon_i}
\]
(où $\mathrm{long}(\pi_i)=\sum_{j\in[1,l_i]}\pi_i(j)$). Donc
\[
\left\{
\begin{array}{lcl}
x & = & \prod^{\mathrm{Ty}}_{i\in[1,n]}\Delta_{\varepsilon_i}\\[4pt]
y & = & \prod^{\mathrm{Ty}}_{i\in[1,n]}\Delta^{l_i}_{\varepsilon_i}\\[4pt]
z & = & \prod^{\mathrm{Ty}}_{i\in[1,n]}\Delta^{\mathrm{long}(\pi_i)}_{\varepsilon_i}
\end{array}
\right.
\]
(\footnote[16]{attention, les produits doivent être pris en respectant l'ordre des facteurs donné par la notation indicielle en $i$.}), et on trouve de même $f$, $g$ comme des produits (au sens $*$) de $n$ facteurs $f_i$ resp. $g_i$ ($i\in[1,n]$) explicités dans les formules (1), (2) (p. 52, 53) :
\[
\begin{array}{lclccl}
f_i & = & \mathfrak{f}_{\varepsilon_i,l_i} : & \Delta_{\varepsilon_i} & \longrightarrow & \Delta^{l_i}_{\varepsilon_i}\\[4pt]
g_i & = & g_{\varepsilon_i,l_i} : & \Delta^{l_i}_{\varepsilon_i} & \longrightarrow & \Delta^{\mathrm{long}(\pi_i)}_{\varepsilon_i}, % [sic] g_{\varepsilon_i,l_i} -> g_{\varepsilon_i,\pi_i}, notation of (1), (2)
\end{array}
\]
$g_i$ lui-même étant représenté comme un $*$-produit de $l_i$ facteurs indécomposables, de sorte que $f$ et $g$ apparaissent sous forme
\msp{page 55}
de $*$-produits de $n$, resp. de $\sum_{i\in[1,n]}l_i$ facteurs de la forme $\mathfrak{f}_{\varepsilon,\nu}$. Pour terminer de tenir en mains l'objet semi-simplicial tronqué $(\mathrm{Nerf}_i(\uTy))_{0\leq i\leq 2}$, il reste à décrire les deux applications de dégénérescence $\mathrm{Nerf}_1\longrightarrow\mathrm{Nerf}_2$ (cas où $f$ ou $g$ est une flèche identique) et surtout l'application de \textit{composition} des flèches
\[
\mathrm{Nerf}_2(\uTy)\xrightarrow{\ \gamma\ }\mathrm{Nerf}_1(\uTy)=\mathrm{Fl}(\uTy).
\]
Il suffit de donner $\gamma$ sur les générateurs, et on a tautologiquement, par définition des $\mathfrak{f}_{\varepsilon,l,\pi}$
\[
\gamma(\mathfrak{f}_{\varepsilon,l,\pi})=\mathfrak{f}_{\varepsilon,\mathrm{long}(\pi)}.
\]
(\textbf{NB} Il n'existe au plus qu'un seul morphisme de $\Delta_\varepsilon$ (de $\Delta_+$ ou de $\Delta_-$) dans un intervalle $I_\tau$. Il existe si et seulement si $\tau$ est de la forme $\Delta^l_\varepsilon$, $l\in\mathbf{N}$.)
% foot: Version 22 décembre 2001

\printedpage{38}
% head: A. Grothendieck : Dérivateurs, Ch. VII, Édition des Universités de Montpellier II et Paris VII
Le cas $g=\mathrm{id}$ correspond au cas où
\[
\pi_i(j)=1 \qquad\quad \forall j\in[1,l_i],
\]
les $\mathfrak{f}_{\varepsilon_i,l_i,\pi_i}$ sont alors détérminés en fonction des $\mathfrak{f}_{\varepsilon_i,l_i}$, et l'application de dégénérescence est % [sic] détérminés -> déterminés, accent
\[
\mathfrak{f}_{\varepsilon,l}\longmapsto\mathfrak{f}_{\varepsilon,l,1_l} \qquad\quad (\text{où } 1_l \text{ est la fonction constante de valeur 1 sur } [1,l]).
\]
Le $f=\mathrm{id}$ est celui où les $l_i$ sont égaux à 1. Alors chaque $\pi_i : [1,1]\longrightarrow\mathbf{N}$ peut être vu comme un élément de $\mathbf{N}$, la donnée de $(\varepsilon_i,l_i,\pi_i)$ revient à celle des $(\varepsilon_i,n_i)\in\{\pm1\}\times\mathbf{N}$, i.e. la donnée % [sic] Le f = id -> Le cas f = id, word « cas » missing
\msp{page 56}
d'une suite de $n$ générateurs parmi les $\mathfrak{f}_{\varepsilon,l}$. L'application de dégénérescence $\mathrm{Nerf}_1\longrightarrow\mathrm{Nerf}_2$ est ici donnée sur les générateurs par
\[
\mathfrak{f}_{\varepsilon,l}\longmapsto\mathfrak{f}_{\varepsilon,1,l}.
\]
Heureusement, on n'a plus à prouver l'associativité de la composition.

\textbf{NB} La simplicité de ces descriptions renforce mon sentiment que la catégorie $\uTy$ définie à présent, avec les \textit{morphismes d'intervalles orientés croissants} $I_\tau\longrightarrow I_{\tau'}$ comme flèches de $\uTy$, est bien `la bonne'. On va déjà voir si elle donne une description sympa des catégories de fractions $X\Sigma^{-1}$, sans tortillements comme précédemment.

Mais on a promis de déterminer $\mathrm{Nerf}_r(\uTy)$, pour tout $r\in\mathbf{N}$. Le générateur général de $\mathrm{Nerf}_r$ est de la forme (pour $\varepsilon\in\{\pm1\}$ fixé)
\[
\underbrace{\Delta_\varepsilon}_{=\,\Delta^{l_0}_\varepsilon\text{avec } l_0=1}\xrightarrow{\ f_1\ }\Delta^{l_1}_\varepsilon\xrightarrow{\ f_2\ }\Delta^{l_2}_\varepsilon\longrightarrow\ \cdots\ \xrightarrow{\ f_r\ }\Delta^{l_r}_\varepsilon,
\]
où les $f_i : \Delta^{l_{i-1}}_\varepsilon\longrightarrow\Delta^{l_i}_\varepsilon$, $i\in[1,r]$, sont donnés par des applications
\[
\pi_{i-1} : [1,l_{i-1}]\longrightarrow\mathbf{N} \qquad\quad \text{telles que }\ \underbrace{\mathrm{long}(\pi_{i-1})}_{\overset{\text{déf}}{=}\ \Sigma_{j\in[1,l_i]}\pi_i(j)}\ = l_i. % [sic] \Sigma_{j\in[1,l_i]}\pi_i(j) -> \Sigma_{j\in[1,l_{i-1}]}\pi_{i-1}(j), indices shifted with respect to long(\pi_{i-1})
\]
(NB $\pi_0 : [1,1]=\{1\}\longrightarrow\mathbf{N}$ est identifié à $\mathrm{long}(\pi_0)=\pi_0(1)=l_1$.)

\msp{page 57}
Pour ne pas faire dépendre la définition du domaine des $\pi_i$ ($i\in[0,r-1]$) des entiers $l_i$, on peut aussi considérer les $\pi_i$ comme des applications
\[
\mathbf{N}^*\xrightarrow{\ \pi\ }\mathbf{N} \qquad\quad (\mathbf{N}^*=\{i\in\mathbf{N}\mid i\geq 1\})
\]
à supports finis (i.e. tels que $\pi(i)=0$ pour $i$ grand), et exiger de façon précise
\[
\begin{aligned}
&\mathrm{supp}\,\pi_0\ \subseteq\ \{1\}\\
&\mathrm{supp}\,\pi_1\ \subseteq\ [1,\mathrm{long}(\pi_0)]\\
&\qquad\ \cdots\\
&\boxed{\mathrm{supp}\,\pi_i\ \subseteq\ [1,\mathrm{long}(\pi_{i-1})]}, \qquad\quad \text{pour }\ \boxed{i\in[0,r-1]}\quad .
\end{aligned}
\tag*{$(**)$}
\]
% foot: Version 22 décembre 2001


\endgroup
% END UNIT Der-ChVII-456/Der-ChVII-456_p036-040.tex
% BEGIN UNIT Der-ChVII-456/Der-ChVII-456_p041-045.tex
\begingroup\makeatletter%
\newcommand{\printedpage}[1]{\marginpar{\footnotesize p.~#1}}%
\newcommand{\msp}[1]{\marginpar{\footnotesize [#1]}}%
\newcommand{\Ch}{\underline{\mathrm{Ch}}}%
\newcommand{\Ty}{\underline{\mathrm{Ty}}}%
\newcommand{\Nerf}{\mathrm{Nerf}}%
\makeatother\let\printedpage\derpp\let\msp\dermsp%
\let\mspin\DERmspin%


\printedpage{39}
% head: A. Grothendieck : Dérivateurs, Ch. VII, Édition des Universités de Montpellier II et Paris VII
D'autre part, posant $l_i = \mathrm{long}(\pi_{i-1})$ (pour $i \in [1,r]$),
\[
f_i : \Delta_\varepsilon^{l_{i-1}} \longrightarrow \Delta_\varepsilon^{l_i\ (=\mathrm{long}(\pi_{i-1}))}
\]
est défini comme
\[
f_i = \mathfrak{f}_{\varepsilon,\pi_{i-1}(1)} * \mathfrak{f}_{\varepsilon,\pi_{i-1}(2)} * \cdots * \mathfrak{f}_{\varepsilon,\pi_{i-1}(l_{i-1})},
\]
de sorte que l'élément $(*)$ de $\Nerf_r(\Ty)$ est déterminé de façon unique par la donnée des $\varepsilon \in \{\pm 1\}$ et de
\[
\pi_\bullet = (\pi_i)_{i\in[0,r-1]} \in (\mathrm{Hom}_{\mathrm{Ens}}(\mathbf{N}^*,\mathbf{N}))^r
\]
satisfaisant les conditions $(**)$. Soit $F_{\varepsilon,\pi_\bullet}$ cet élément de $\Nerf_r(\Ty)$, et soit $N_r$ l'ensemble des $\pi_\bullet$ qu'on vient d'expliciter. Alors
\[
(F_{\varepsilon,\pi_\bullet})_{(\varepsilon,\pi_\bullet)\in\{\pm1\}\times N_r}
\]
est un système de générateurs libre du monoïde $\Nerf_r(\Ty)$.

Si $l_i = 0$ (i.e. \fbox{$\pi_{i-1} = 0$}), on doit avoir $\mathrm{supp}(\pi_{i+1}) = \emptyset$, i.e. $\pi_{i+1} = 0$, d'où $l_{i+1} = 0$, et de même pour les $l_j$ et $\pi_j$ suivants. Géométriquement l'hypothèse s'exprime par la condition que le but
\msp{page 58}
de $f_i$ est l'élément 1 de Ty (correspondant à l'intervalle ponctuel), et comme $\mathrm{Hom}(1,\tau) = \emptyset$ pour $\tau \neq 1$, cela signifie que les flèches qui suivent $f_i$ doivent être toutes égales à la flèche $\mathrm{id}_1 : 1 \longrightarrow 1$.

\textbf{Remarque.} Dans la description combinatoire de $\Ty$ et de $\Nerf_\bullet(\Ty)$, l'ensemble des `signatures' $\{\pm1\}$ joue un rôle assez bidon. Il peut être remplacé par n'importe quel ensemble $\Sigma$ (dit `ensemble des signatures'), et on construit alors une catégorie monoïdale (ou un monoïde catégorique) ayant comme ensemble d'objets le monoïde libre engendré par $\Sigma$, interprétés comme des intervalles d'entiers $[0,n]$, munis de `signatures' dans $\Sigma$ pour chaque intervalle de transition, ou ce qui revient au même, des `graphes' orientés $\Sigma$-pondérés de type combinatoire un intervalle $\boldsymbol{\Delta}_i$, ce qui revient à la donnée d'une suite (ordonnée) de $n$ éléments $\varepsilon_1,\ \varepsilon_2,\ \ldots\ ,\ \varepsilon_n \in \Sigma$, i.e. d'un `mot' $\varepsilon_1\varepsilon_2\ldots\varepsilon_n$, qu'on préfère écrire $\Delta_{\varepsilon_1}\Delta_{\varepsilon_2}\cdots\Delta_{\varepsilon_n}$. On peut définir de façon évidente les homomorphismes \textit{croissants} de tels graphes pondérés orientés (sous-entendu : avec l'élément source et l'élément but dans la structure, à respecter par les morphismes), sans avoir besoin d'une interprétation catégorique. La fonctorialité de l'opération de composition des intervalles pondérés, pour cette notion de morphisme, est évidente. Le
\msp{page 59}
choix $\Sigma = \{\pm1\}$ est à ma connaissance le seul utile, à cause de son lien avec la notion de chemins et le calcul de $\pi_0$ d'un objet de $X$ \texttt{[plutôt d'un objet} $X$ \texttt{de} Cat\texttt{]}, etc.
% foot: Version 22 décembre 2001

\printedpage{40}
% head: A. Grothendieck : Dérivateurs, Ch. VII, Édition des Universités de Montpellier II et Paris VII
\section*{7\quad Explicitation des catégories de fractions $X\Sigma^{-1}$ en termes des $\Ch(X;a,b)$}

Si $X$ est une catégorie, on désigne par $\Ch(X)$ la catégorie des chemins dans $X$ (cf. §4), et pour $a,b \in \mathrm{Ob}\,X$, par
\[
\Ch(X;a,b) \qquad \text{ou simplement } \Ch(a,b)
\]
celle des chemins de $a$ vers $b$, i.e. de source $a$, but $b$. L'intuition géométrique (qu'on explicitera), c'est qu'un tel chemin $c$ définit un chemin `direct' de $a$ vers $b$ (i.e. un chemin de type $\tau'$ `direct', i.e. toutes les flèches de transition sont directes, i.e. $\tau'$ de la forme $\Delta_+^n$), de même longueur que le chemin initial $c$, obtenu en inversant toutes les flèches rétrogrades dans le chemin $c$. Cette opération est effectivement possible dans $X$ lui-même si et seulement si lesdites flèches sont inversibles. On obtient alors une suite composable de flèches, dont le composé total (qu'on peut appeler `l'effet utile' du chemin)
\msp{page 60}
est une flèche de $a$ vers $b$. Dans le cas général, les chemins de $a$ vers $b$ peuvent être vus comme une sorte de `flèches généralisées' allant de $a$ vers $b$, dont la composition serait la composition évidente, et les identités les `chemins vides' relatifs aux $a \in \mathrm{Ob}\,X$. On trouve d'ailleurs bien ainsi une catégorie, ayant mêmes objets que $X$, mais où les
\[
\mathrm{Hom}_{\mathrm{Ch}}(a,b) = \mathrm{Ch}(a,b) = \mathrm{Ob}\,\Ch(a,b)
\]
sont prohibitivement gros. Cela s'exprime notamment par le fait que les seules flèches inversibles dans $X_{\mathrm{Ch}}$ sont les identités ! On n'a pas de foncteur naturel entre $X_{\mathrm{Ch}}$ et $X$ qui soit l'identité sur les objets. Car à un chemin $a \xrightarrow{c} b$ on ne peut associer de $X$-morphismes $a \longrightarrow b$ (sauf si $c$ \texttt{[est]} de type $\Delta_+$), et à un $X$-morphisme $a \xrightarrow{u} b$ on peut associer un chemin $c(u)$ de type $\Delta_+$, mais si $a \xrightarrow{u} b \xrightarrow{v} c$, $c(vu)$ est encore de type $\Delta_+$ alors que $c(v)c(u)$ est de type $\Delta_+^2$, ils ne sont pas égaux. Par contre, il y a une flèche naturelle entre ces deux éléments de $\Ch(a,c)$, flèche au dessus de l'unique homomorphisme
\[
\Delta_+ \xrightarrow{\mathfrak{f}_{+,2}} \Delta_+^2.
\]
\msp{page 61}
On a alors
\[
\mathfrak{f}_{+,2}^*(\underbrace{a \xrightarrow{u} b \xrightarrow{v} c}_{c(v)c(u)}) = (\underbrace{a \xrightarrow{vu} c}_{c(vu)}),
\]
d'où une flèche canonique
\[
c(vu) \longrightarrow c(v)c(u) \qquad \text{dans } \Ch(a,c).
\]
Cela montre que $c(vu)$ et $c(v)c(u)$ sont tout au moins dans une même composante connexe de $\Ch(a,c)$, ou (comme on dit encore) qu'ils sont `homotopes'. Cela suggère ici de modifier la définition de $\mathrm{Hom}_{\mathrm{Ch}}(a,b)$ en
\[
\pi(a,b) = \pi_0(\Ch(a,b)).
\]
% foot: Version 22 décembre 2001

\printedpage{41}
% head: A. Grothendieck : Dérivateurs, Ch. VII, Édition des Universités de Montpellier II et Paris VII
Le foncteur de composition des chemins
\[
\begin{array}{ccccc}
\Ch(a,b) & \times & \Ch(b,c) & \longrightarrow & \Ch(a,c)\\
(c & , & c') & \longmapsto & c' \circ c
\end{array}
\]
\texttt{[il y a trop de} $c$\texttt{]} (\footnote[17]{Ce foncteur est déduit des foncteurs
\[
\Ch_\tau(a,b) \times \Ch_{\tau'}(b,c) \longrightarrow \Ch_{\tau\tau'}(a,c)
\]
évidents, cf. page 77.}) induit une application sur les $\pi_0$, ce qui donne une loi de composition
\[
\pi(a,b) \times \pi(b,c) \longrightarrow \pi(a,c)
\]
sur les flèches. Il est immédiat qu'elle est associative (la composition des chemins l'étant), et unitaire. Elle admet comme flèche identique en $x$ soit le `lacet vide en $x$', lequel est (directement) homotope dans $\Ch(x,x)$ au chemin `dégénéré' $x \xrightarrow{\mathrm{id}} x$
\msp{page 62}
de type $\Delta_+$, ou aussi au chemin dégénéré de n'importe quel type $\tau$, disons
\[
x \xleftarrow{\mathrm{id}} x \xrightarrow{\mathrm{id}} x \xrightarrow{\mathrm{id}} x \xleftarrow{\mathrm{id}} x \xleftarrow{\mathrm{id}} x,
\]
un tel chemin n'étant autre que $p_\tau^*(x)$, où
\[
p_\tau : \tau \longrightarrow 1
\]
est le morphisme de dégénérescence (l'unique morphisme dans $\Ty$ vers 1, qui est objet final), d'où
\[
p_\tau^* : \Ch_1(x,x) = \{x\} \longrightarrow \Ch_\tau(x,x) \subseteq \Ch(x,x).
\]
La catégorie $\Pi(X)$ ainsi obtenu doit s'appeler le \textit{groupoïde fondamental} de $X$. En effet, % [sic] obtenu -> obtenue, agreement with « La catégorie »

\begin{enumerate}[label=\alph*)]
\item toutes les flèches de $\Pi(X)$ sont inversibles, i.e. $\Pi(X)$ est bien un groupoïde, et
\item le foncteur canonique
\[
X \longrightarrow \Pi(X)
\]
qui est l'identité sur les objets, et qui est défini sur les flèches par $u \longmapsto c(u)$ comme ci-dessus, est \textit{universel} (dans Cat) pour les foncteurs (i.e. les Cat-morphismes) de $X$ vers des \textit{groupoïdes.} (Plus généralement encore, il est universel pour les foncteurs $X \xrightarrow{F} Y$ tels que $F$ transforme toute flèche de $X$ en \texttt{[une]} flèche inversible.)
\end{enumerate}
% foot: Version 22 décembre 2001

\printedpage{42}
% head: A. Grothendieck : Dérivateurs, Ch. VII, Édition des Universités de Montpellier II et Paris VII
\msp{page 63}
Pour prouver que toute flèche dans $\Pi(X)$ est inversible, comme tout chemin non vide est composé de chemins `élémentaires' de type $\Delta_+$ ou $\Delta_-$, il suffit de prouver que ces chemins sont inversibles. En fait, ils se correspondent biunivoquement,
\[
\Ch_{\Delta_+}(a,b) \simeq \Ch_{\Delta_-}(b,a),
\]
en associant au chemin de $a$ vers $b$
\[
c(u) : a \xrightarrow{u} b \qquad \text{de type } \Delta_+
\]
le chemin
\[
c(u)^o : b \xleftarrow{u} a \qquad \text{de type } \Delta_-
\]
de $b$ vers $a$. Je dis que
\[
\begin{array}{rcl}
c(u)^o c(u) & \sim & 1_a\\
c(u) c(u)^o & \sim & 1_b
\end{array}
\]
\[
\left.
\begin{tikzcd}
 & a \arrow[r, "u"] & b & a \arrow[l, "u"'] \\
 {} & a \arrow[u, "\mathrm{id}"'] \arrow[r, "\mathrm{id}"] & a \arrow[u, "u"'] & a \arrow[u, "\mathrm{id}"'] \arrow[l, "\mathrm{id}"'] \\
 {} \arrow[u] & & a &
\end{tikzcd}
\right\}\
\begin{array}{l}
\text{on a un homomorphisme dans } \Ch_{\Delta_+\Delta_-}(a,a)\\
\text{de } 1_a \text{ dans } c(u)^o c(u).
\end{array}
\]
\[
\left.
\begin{tikzcd}
 & b & a \arrow[l, "u"'] \arrow[r, "u"] & b \\
 {} & a \arrow[u, "u"'] & a \arrow[u, "\mathrm{id}"'] \arrow[l, "\mathrm{id}"'] \arrow[r, "\mathrm{id}"] & a \arrow[u, "u"'] \\
 {} \arrow[u] & & a &
\end{tikzcd}
\right\}\
\begin{array}{l}
\text{itou dans } \Ch_{\Delta_-\Delta_+}(b,b),\\
\text{de } 1_b \longrightarrow c(u)c(u)^0 % [sic] c(u)^0 -> c(u)^o, digit 0 printed instead of the letter o used elsewhere
\end{array}
\]
\texttt{[plutôt} $1_b \longleftarrow c(u)c(u)^0$\texttt{, viz.} % [sic] c(u)^0 -> c(u)^o, digit 0 printed instead of the letter o used elsewhere
\[
\begin{tikzcd}
b \arrow[d, "\mathrm{id}"] & a \arrow[l, "u"'] \arrow[r, "u"] \arrow[d, "u"] & b \arrow[d, "\mathrm{id}"] \\
b & b \arrow[l, "\mathrm{id}"'] \arrow[r, "\mathrm{id}"] & b \ \texttt{].}
\end{tikzcd}
\]
% foot: Version 22 décembre 2001

\printedpage{43}
% head: A. Grothendieck : Dérivateurs, Ch. VII, Édition des Universités de Montpellier II et Paris VII
\msp{page 64}
Notons que, plus généralement, pour tout chemin
\[
c : a \longrightarrow b, \qquad \text{de type } \tau,
\]
on définit un \textit{chemin opposé}
\[
c^o : b \longrightarrow a \qquad \text{de type } \tau^o,
\]
où
\[
\tau \longmapsto \tau^o : \mathrm{Ty} \longrightarrow \mathrm{Ty}
\]
est l'antiïnvolution du monoide libre Ty qui échange les deux générateurs $\Delta_-$, $\Delta_+$, et qui est donc donné sur l'élément général par
\[
(\Delta_{\varepsilon_1}\Delta_{\varepsilon_2}\cdots\Delta_{\varepsilon_n})^o = \Delta_{-\varepsilon_n}\Delta_{-\varepsilon_{n-1}}\cdots\Delta_{-\varepsilon_1}.
\]
Par récurrence sur la longueur de $c$ (\footnote[18]{inutile, résulte directement de la formule $(*)$ \texttt{[de la]} page suivante, en décomposant $c$ en composés de chemins `élémentaires' de type $\Delta_+$ ou $\Delta_-$.}), en utilisant le résultat précédent, que $c(u)$, $c(u)^o$ sont inverses l'un de l'autre dans $\Pi(X)$, on trouve
\[
\left.
\begin{array}{rcl}
c^o c & = & \mathrm{id}_a\\
c c^o & = & \mathrm{id}_b
\end{array}
\right\}\text{dans } \Pi(X).
\]

\textbf{Exemple.} Si
\[
\begin{array}{rcl}
\tau & = & \Delta_+\Delta_-\Delta_+^2\Delta_-\\
\tau^o & = & \Delta_+\Delta_-^2\Delta_+\Delta_-,
\end{array}
\]
et si on a un chemin
\[
c : a_0 \xrightarrow{u_1} a_1 \xleftarrow{u_2} a_2 \xrightarrow{u_3} a_3 \xrightarrow{u_4} a_4 \xleftarrow{u_5} a_5
\]
de type $\tau$, on en déduit le chemin
\[
c^o : a_5 \xrightarrow{u_5} a_4 \xleftarrow{u_4} a_3 \xleftarrow{u_3} a_2 \xrightarrow{u_2} a_1 \xleftarrow{u_1} a_0
\]
de type $\tau^o$. L'opération $c \longmapsto c^o$ satisfait en
\msp{page 65}
effet \texttt{[pour]}
\[
\begin{array}{ccccc}
a & \xrightarrow{\gamma} & b & \xrightarrow{\gamma'} & c\\
c & \xrightarrow{\gamma'^o} & b & \xrightarrow{\gamma^o} & a
\end{array}
\]
la relation
\[
(\gamma'\gamma)^o = \gamma^o\gamma'^o, \tag*{$(*)$}
\]
% foot: Version 22 décembre 2001


\endgroup
% END UNIT Der-ChVII-456/Der-ChVII-456_p041-045.tex
% BEGIN UNIT Der-ChVII-456/Der-ChVII-456_p046-050.tex
\begingroup\makeatletter%
\newcommand{\printedpage}[1]{\marginpar{\footnotesize p.~#1}}%
\newcommand{\msp}[1]{\marginpar{\footnotesize [#1]}}%
\newcommand{\Ch}{\underline{\mathrm{Ch}}}%
\newcommand{\Ty}{\mathrm{Ty}}%
\newcommand{\uTy}{\underline{\mathrm{Ty}}}%
\newcommand{\Fl}{\mathrm{Fl}}%
\newcommand{\Hom}{\mathrm{Hom}}%
\newcommand{\ch}{\mathrm{ch}}%
\newcommand{\hot}{\mathrm{hot}}%
\newcommand{\id}{\mathrm{id}}%
\makeatother\let\printedpage\derpp\let\msp\dermsp%
\let\mspin\DERmspin%


\printedpage{44}
% head: A. Grothendieck : Dérivateurs, Ch. VII, Édition des Universités de Montpellier II et Paris VII
% foot: Version 22 décembre 2001

ce qui est en accord avec l'intuition que l'opération $\gamma \longmapsto \gamma^o$ joue le rôle d'un passage à l'inverse (`à homotopie près'), et on a la formule
\[
(\gamma'\gamma)^{-1} = \gamma^{-1}\gamma'^{-1}
\]
pour deux flèches inversibles $\gamma$, $\gamma'$ composables (dans l'ordre $\gamma'$, $\gamma$).

Pour la propriété universelle de
\[
X \longrightarrow \Pi(X)
\]
(qui permettra d'écrire
\[
\Pi(X) \simeq X\Sigma^{-1}, \qquad \text{avec } \Sigma = \Fl(X) \ \ ),
\]
autant l'établir tout de suite dans le cadre général où on se donne une partie
\[
\Sigma \subseteq \Fl(X),
\]
et on définit, pour $\tau \in \Ty$,
\[
\Ch_{\Sigma,\tau}(X) \subseteq \Ch_{\tau}(X)
\]
comme la sous-catégorie pleine de $\Ch_{\tau}(X)$ formée des chemins dans $X$ dont \textit{toutes les flèches rétrogrades}

\msp{page 66}

sont dans $X$ \texttt{[plutôt dans $\Sigma$]}. On notera que pour qu'un chemin $c$ et son opposé $c^o$ soient tous les deux des $\Sigma$-\textit{chemins,} il faut et il suffit que \textit{toutes} les flèches du chemin $c$ soient dans $\Sigma$ (on dira alors que $c$ est un $\Sigma$-\textit{chemin stricte}), puisque les flèches rétrogrades de $c^o$ sont les flèches directes de $c$.

Pour que les sous-catégories pleines $\Ch_{\tau,\Sigma}(X)$ soient stables par les foncteurs images inverses
\[
f^* : \Ch_{\tau'}(X) \longrightarrow \Ch_{\tau}(X)
\]
associés aux flèches
\[
f : \tau \longrightarrow \tau'
\]
dans $\uTy$, ou simplement pour qu'y ait stabilité pour les deux foncteurs $f^*$ associés à un $f : \Delta_- \longrightarrow 1$ et $\Delta_- \longrightarrow \Delta_-^2$, il faut et il suffit que $\Sigma$ soit stable par composition et contienne toutes les flèches identiques % [sic] qu'y ait -> qu'il y ait, missing word
\[
\left\{
\begin{aligned}
& x \in \mathrm{Ob}\,X \ \Longrightarrow \ \id_x \in \Sigma \\
& u, v \in \Sigma, \ \ u,\, v \text{ composables} \Longrightarrow uv \in \Sigma.
\end{aligned}
\right.
\]
En effet, les flèches dans un chemin $f^*(c') \in \Ch_{\tau}(X)$ sont soit des flèches identiques (qui correspondent aux flèches de transition de $I_\tau$ dégénérées pour $f$), soit des flèches parmi celles de $\tau$, ou des composés d'une suite de flèches consécutives composables de $\tau'$ (qui correspondent aux flèches de transition dans $I_{\tau'}$ incluses dans la partie d'une flèche de transition de $I_\tau$), avec conservation des signatures.

\printedpage{45}
% head: A. Grothendieck : Dérivateurs, Ch. VII, Édition des Universités de Montpellier II et Paris VII
% foot: Version 22 décembre 2001

\msp{page 67}

On suppose ces conditions satisfaites, de sorte que les $\Ch_{\tau,\Sigma}(X)$ forment les fibres d'une sous-catégorie fibrée pleine de $\Ch(X)$ sur $\uTy$, qu'on dénotera par
\[
\Ch_{\Sigma}(X) \hookrightarrow \Ch(X).
\]
Il est évident que cette sous-catégorie est stable par l'opération de composition des chemins (mais non pas, en général, on l'a dit, par passage de $c$ à $c^o$ (il faudrait pour cela que $\Sigma$ contienne \textit{toutes} les flèches de $X$, puisque pour toute $u : a \longrightarrow b$, qui est donc dans $\Ch_{\Delta_+,\Sigma}(X)$, $u^o$ donné par $b \xleftarrow{\ u\ } a$ dans $\Ch_{\Delta_-}(X)$ devrait être un $\Sigma$-chemin, d'où $u \in \Sigma$)).

On va définir une catégorie canonique
\[
\Pi_\Sigma(X),
\]
ayant même objets que $X$, et un foncteur canonique
\[
X \longrightarrow \Pi_\Sigma(X)
\]
qui est l'identité sur les objets, et qui transforme les $u \in \Sigma$ en flèches inversibles. Pour cela on calque la description du groupoïde fondamental (mais $\Pi_\Sigma(X)$ ne sera pas en général un groupoïde), en posant
\[
\Hom_\Sigma(x,y) = \pi_0(\Ch_\Sigma(x,y)),
\]

\msp{page 68}

et en définissant la composition des `$\Sigma$-flèches' en termes de la composition des chemins. L'argument qui prouvait, dans le cas de $\Pi(X)$, que \textit{toute} flèche de $X$, vue comme chemin $c$ de type $\Delta_+$, était inversible dans $\Pi(X)$, ayant comme inverse $c^o$ de type $\Delta_-$, est encore valable, pourvu que $c^o$ soit encore bien un $\Sigma$-chemin, ce qui exige que $u \in \Sigma$.

On va l'interpréter en termes d'un foncteur canonique
\[
\kappa : X \longrightarrow \Pi_\Sigma(X) \qquad \text{(\footnotemark[19])}
\tag*{$(*)$}
\]
\footnotetext[19]{Sur les flèches, $\kappa$ est défini, pour $u : a \longrightarrow b$, par $\kappa(u) =$ classe d'homotopie de $\ch(u)$ dans $\Ch_\Sigma(a,b)$, où $\ch(u)$ est le chemin de type $\Delta_+$ de $a$ vers $b$ défini par $u$.}
qui se décrit encore comme dans le cas du groupoïde fondamental, en associant à toute flèche $u$ de $X$ la classe d'homotopie du chemin $c(u)$ qu'elle définit. On voit donc que ce foncteur canonique transforme les $u \in \Sigma$ en flèches inversibles.

Quand les flèches de $\Sigma$ sont déjà inversibles dans $X$, je dis que le foncteur $(*)$ est un isomorphisme, i.e. que pour tout $a, b \in \mathrm{Ob}\,X$,
\[
\alpha : \Hom_X(a,b) \longrightarrow \pi_0\Ch_\Sigma(a,b), \qquad u \longmapsto \underbrace{\ch_\Sigma(u)}_{= \ch(u)}
\]
est bijectif. Pour le voir, on va définir
\[
\beta : \pi_0(\Ch_\Sigma(a,b)) \longrightarrow \Hom_X(a,b),
\]
qui sera inverse de $\beta$. On définit % [sic] inverse de $\beta$ -> inverse de $\alpha$, $\beta$ is meant to be the inverse of $\alpha$

\printedpage{46}
% head: A. Grothendieck : Dérivateurs, Ch. VII, Édition des Universités de Montpellier II et Paris VII
% foot: Version 22 décembre 2001

\msp{page 69}

$\beta(c)$, pour une $\Sigma$-flèche de type $\tau$, le longueur $n$, en inversant toutes les flèches rétrogrades (ce qui est possibles, puisqu'elles sont dans $\Sigma$), de façon à obtenir un chemin $c^\natural$ de type $\Delta_+^n$, et en prenant le composé (dans $X$) des flèches de ce chemin $c^\natural$. Il est évident alors que $\beta\alpha = \id$, il faut prouver que l'on a aussi $\alpha\beta = \id$, i.e. que dans $\Pi_\Sigma(X)$, les flèches $u = \int c^\natural$ et $c$ sont égales, i.e. $u$ et $c$ sont homotopes dans $\Ch_\Sigma(a,b)$. Considérons le chemin $c^\natural$, de type $\Delta_+^n$ : % [sic] le longueur -> la longueur, gender; possibles -> possible, number
\[
a = a_0 \xrightarrow{u_1} a_1 \xrightarrow{u_2} a_2 \longrightarrow \ \cdots \ \longrightarrow a_{n-1} \xrightarrow{u_n} a_n = b,
\]
et soit $u_i$ ($i \in [1,n]$) une flèche de ce chemin, correspondant à une flèche de transition $\underline{u}_i$ \textit{rétrograde} de $\tau$, donc par construction
\[
u_i = v_i^{-1} \qquad \text{dans } X\text{, avec } v_i \in \Sigma
\]
la flèche d'indice $i$ dans le chemin initial $c$ de type $\tau$. Si $\underline{u}_i$ par contre est directe, on aura
\[
u_i = v_i.
\]
(NB On désigne par $v_i$ la $i$-flèche dans le chemin $c$.) Or on a
\[
c = c_n * c_{n-1} * \cdots * c_1 \qquad \text{dans } \Pi_\Sigma(X)
\tag*{$(*)$}
\]
et le facteur $c_i$ est soit la classe d'homotopie de $\ch(v_i) = \ch(u_i)$ (cas où la flèche de transition $\underline{u}_i$ de $\tau$ est directe), soit celle du chemin \textit{opposé} au chemin $\ch(v_i) : a_i \xrightarrow{v_i} a_{i-1}$ de type $\Delta_+$, donc le chemin $\ch(v_i)^o : a_{i-1} \xleftarrow{v_i} a_i$ de type $\Delta_-$. Je dis que
\[
\boxed{\ \ch(v_i)^o = \ch(u_i) \qquad \text{\textit{dans} } \Pi(X),\ }
\tag*{(1)}
\]
ce qui prouve que dans $\Pi_\Sigma(X)$, $c$ est égal au composé des $\ch(u_i)$ ($i \in [1,n]$), donc à $c^\natural$. Il faut prouver de plus que
\[
c^\natural = \ch(u) \qquad (u = u_nu_{n-1}\cdots u_1) \qquad \text{dans } \Pi(X).
\tag*{(2)}
\]

\msp{page 70}

Donc il y a deux vérifications à faire (où je suppose à nouveau $\Sigma$ quelconque) :

\begin{enumerate}[label={}]
\item[$1^o$)] Soit $u \in \Sigma$, $u : a \longrightarrow b$, supposons $\Sigma$ inversible dans $X$. Alors
\[
\ch(u)^o = \ch(u^{-1}) \qquad \text{dans } \Pi_\Sigma(X),
\]
i.e. les deux chemins de $b$ vers
\[
\begin{array}{cccl}
b & \xleftarrow{\ u\ } & a & \text{de type } \Delta_- \\
b & \xrightarrow{u^{-1}} & a & \text{de type } \Delta_+
\end{array}
\]
\printedpage{47}
% head: A. Grothendieck : Dérivateurs, Ch. VII, Édition des Universités de Montpellier II et Paris VII
% foot: Version 22 décembre 2001
sont homotopes dans $\Ch_\Sigma(b,a)$. Or on sait déjà, comme $\kappa : X \longrightarrow \Pi_\Sigma(X)$ est un foncteur, que
\[
\hot\,\ch(u^{-1}) = (\hot\,\ch(u))^{-1} \qquad \text{dans } \Pi_\Sigma(X)
\]
où $\hot(\bullet)$ désigne la classe d'homotopie d'une chaine (ici classe dans $\Ch_\Sigma(b,a)$ et dans $\Ch_\Sigma(a,b)$ respectivement). Donc l'énoncé à prouver c'est que $\hot(\ch(u)^o)$ est égal à l'\textit{inverse} de $\hot(\ch(u))$ dans $\Pi_\Sigma(X)$. C'est donc ce qu'on a prouvé pour \textit{tout} $u \in \Fl(X)$, dans le cas de $\Pi(X)$. L'argument doit marcher, du moment qu'il a un sens, i.e. que $\ch(u)^o$ est bel et bien un $\Sigma$-chemin, i.e. $u \in \Sigma$.

\item[$2^o$)] La formule (2) signifie que pour tout chemin $c$ de type $\Delta_+^n$ (donc tautologiquement un \msp{page 71} $\Sigma$-chemin), allant de $a$ vers $b$, sa classe dans $\Hom_\Sigma(a,b)$ n'est autre que celle de $\ch(u)$, où $u$ est le composé $u_nu_{n-1}\cdots u_1$. Mais on a
\[
c = c_1 * c_2 * \cdots * c_n,
\]
d'où dans $\Pi_\Sigma(X)$
\[
\hot(c) = \hot(c_n)\hot(c_{n-1})\cdots\hot(c_1),
\]
or $\hot(c_i) = \hot\,\ch(u_i) = \kappa(u_i)$, donc
\[
\begin{aligned}
\hot(c) \ &= \ \kappa(u_n)\kappa(u_{n-1})\cdots\kappa(u_1) \\
&= \ \kappa(\underbrace{u_nu_{n-1}\cdots u_1}_{u}),
\end{aligned}
\]
puisque $\kappa$ est un foncteur, donc
\[
\hot(c) = \hot\,\ch(u),
\]

\hfill c.q.f.d.
\end{enumerate}

Notons maintenant que si $(X', \Sigma')$ est un deuxième couple d'une catégorie $X'$ et d'un ensemble \textit{`admissible'} de flèches de $X'$, et si
\[
f : X \longrightarrow X'
\]
est un foncteur tel que
\[
f(\Sigma) \subseteq \Sigma',
\]
alors il existe un \textit{unique} foncteur
\[
\bar{f} = \Pi_\Sigma(f) : \Pi_\Sigma(X) \longrightarrow \Pi_{\Sigma'}(X')
\]
qui `prolonge' $f$, i.e. tel que le carré
\printedpage{48}
% head: A. Grothendieck : Dérivateurs, Ch. VII, Édition des Universités de Montpellier II et Paris VII
% foot: Version 22 décembre 2001
\msp{page 72}
\[
\begin{tikzcd}
X \arrow[r, "\kappa_X"] \arrow[d, "{\scriptstyle f}"] & \Pi_\Sigma(X) \arrow[d, "{\scriptstyle \bar{f} = \Pi_\Sigma(f)}"] \\
X' \arrow[r, "\kappa_{X'}"] & \Pi_{\Sigma'}(X')
\end{tikzcd}
\]
soit commutatif. L'unicité résulte du fait que sur les objets $\bar{f}$ doit être donné par $f$, et que sur les flèche la commutativité du carré implique que $\bar{f}$ est connu sur les flèches de la forme $\kappa_X(u)$, ainsi donc que sur les inverses de telles flèches, quand ces inverses sont inversibles, donc en particulier sur les $\kappa_X(u)^{-1}$ pour $u \in \Sigma$. Or toute flèche de $\Pi_\Sigma(X)$, provenant d'un chemin, est composé d'une suite de flèches de cette forme. Pour l'existence de $\bar{f}$, on note que $f$ induit des foncteurs % [sic] les flèche -> les flèches, plural
\[
\Ch_\tau(X; a, b) \longrightarrow \Ch_\tau(X'; f(a), f(b))
\]
qui commutent aux opérations $\varphi^*$ sur les chemins associées aux foncteurs
\[
\varphi : \tau \longrightarrow \tau',
\]
d'où un homomorphisme de catégories fibrées sur $\uTy$
\[
f_{a,b} : \Ch(X; a, b) \longrightarrow \Ch(X'; f(a), f(b)).
\]

\msp{page 73}

D'où une application sur les $\pi_0$, i.e.
\[
\bar{f}_{a,b} : \Hom_\Sigma(a,b) \longrightarrow \Hom_\Sigma(f(a), f(b)),
\]
donnée par
\[
\bar{f}_{a,b}(\hot(c)) = \hot(f_{a,b}(c)).
\]
On a évidemment, pour deux chemins $\gamma : a \longrightarrow b$ et $\gamma' : b \longrightarrow c$,
\[
f_{a,c}(\gamma\gamma') = f_{b,c}(\gamma')f_{a,b}(\gamma),
\]
ce qui prouve que $\bar{f}$ est compatible avec la composition. Il est trivial qu'il est compatible avec les identités (car $f_{a,a}$ transforme le chemin vide en $a$ en le chemin vide en $f(a)$), donc on a un foncteur $\bar{f}$, et il rend bel et bien le carré commutatif.

On peut maintenant prouver la propriété universelle de $\Pi_\Sigma(X)$ et de
\[
X \xrightarrow{\ \kappa\ } \Pi_\Sigma(X),
\]
savoir que tout foncteur $f : X \longrightarrow X'$ tel que $u \in \Sigma \Longrightarrow f(u)$ inversible, se factorise \textit{de façon unique} par $\Pi_\Sigma(X)$. L'unicité se voit comme tantôt (\footnote[20]{C'est inutile, cf. ci-dessous.}). Pour l'existence, on note que $f$ applique $\Sigma$ dans l'ensemble $\Sigma'$


\endgroup
% END UNIT Der-ChVII-456/Der-ChVII-456_p046-050.tex
% BEGIN UNIT Der-ChVII-456/Der-ChVII-456_p051-055.tex
\begingroup\makeatletter%
\newcommand{\printedpage}[1]{\marginpar{\footnotesize p.~#1}}%
\newcommand{\msp}[1]{\marginpar{\footnotesize [#1]}}%
\newcommand{\Ch}{\underline{\mathrm{Ch}}}%
\newcommand{\CH}{\underline{\mathrm{CH}}}%
\newcommand{\Ty}{\underline{\mathrm{Ty}}}%
\newcommand{\Hom}{\underline{\mathrm{Hom}}}%
\newcommand{\red}{\mathrm{r\acute{e}d}}%
\newcommand{\uHom}{\underline{\mathrm{Hom}}}%
\makeatother\let\printedpage\derpp\let\msp\dermsp%
\let\mspin\DERmspin%

\setcounter{footnote}{20}

\printedpage{49}
% head: A. Grothendieck : Dérivateurs, Ch. VII, Édition des Universités de Montpellier II et Paris VII
\msp{page 74}
des flèches inversibles de $X'$ (lequel ensemble satisfait bien aux deux conditions d'admissibilité). Donc par ce qui précède on a un diagramme commutatif
\[
\begin{tikzcd}
X \ar[r,"\kappa"] \ar[d,"f"'] & \Pi_{\Sigma}(X) \ar[d,"\bar f = \Pi_{\Sigma}(f)"] \ar[dl,dotted,"\tilde f"'] \\
X' \ar[r,"\kappa'","\sim"'] & \Pi_{\Sigma'}(X'),
\end{tikzcd}
\]
où $\bar f$ est uniquement déterminé. Comme $\kappa'$ est un isomorphisme (comme on a vu tantôt, puisque $\Sigma'$ \texttt{[est]} formé des flèches inversibles), l'existence et l'unicité de $\bar f$ équivaut à la conclusion annoncée,

\hfill c.q.f.d.

\textbf{Encore un sorite} sur les classes d'homotopie d'éléments d'une catégorie $X$. De façon générale, on voit que $X$ et $\Pi_{\Sigma}(X) = X\Sigma^{-1}$ ont mêmes composantes connexes, car si
\msp{page 75}
$a$ et $b$ dans $X$ sont reliés par une flèche de $X\Sigma^{-1}$, i.e. par un $\Sigma$-chemin, alors ils sont dans une même composante connexe de $X$. Mais l'inverse est également vrai, comme on voit p.~ex. en prenant $\Sigma = \mathrm{Fl}(X)$, et en disant que dans le \textit{groupoïde} $\Pi(X)$, où deux sommets sont dans une même composante \texttt{[connexe]} si et seulement si $\exists\ x \longrightarrow y$ (relation qui est non seulement transitive et reflexive, mais aussi symétrique dans un groupoïde). Un chemin
\[
a \xrightarrow{\quad c \quad} b
\]
sera appelé aussi une \textit{homotopie} de $a$ vers $b$. On dit que $a$ et $b$ sont \textit{homotopes} s'il existe une homotopie $c$ de $a$ vers $b$, auquel cas il existe aussi une homotopie `inverse' $c^{o}$ de $b$ vers $a$. Cela signifie aussi que $a$ et $b$ sont dans la même composante connexe de $X$, i.e. définissent le même élément de $\pi_0(X)$.

L'homotopie sera dite \textit{simple} si $c$ est de type $\Delta_{+}$ ou $\Delta_{-}$, i.e. si elle est donnée par une flèche $u : a \longrightarrow b$ ou $v : b \longrightarrow a$. Si on cherche une homotopie entre $a$ et $b$, on essaiera de la trouver la plus `courte' possible. Si on a une homotopie $c$ de longueur $l$ de $a$ vers $b$, alors en remplaçant $\tau$ par $\tau_{\red}$, donc $I_{\tau}$ par $I_{\tau\,\red}$, le sous-intervalle réduit
\msp{page 76}
correspondant, on trouve une homotopie de $a$ à $b$ de type $\tau_{\red}$, dont la longueur est égale à la \textit{`longueur utile'} du type $\tau$ de $c$ (ce qui justifie la terminologie). Notons que si
\[
\tau = \Delta_{\varepsilon}^{l_1}\Delta_{-\varepsilon}^{l_2}\Delta_{\varepsilon}^{l_3}\cdots\Delta_{(-1)^{r}\varepsilon}^{l_r},
\]
alors
\[
\tau_{\red} = \Delta_{\varepsilon}\Delta_{-\varepsilon}\Delta_{\varepsilon}\cdots\Delta_{(-1)^{r}\varepsilon},
\]
% foot: Version 22 décembre 2001
\printedpage{50}
% head: A. Grothendieck : Dérivateurs, Ch. VII, Édition des Universités de Montpellier II et Paris VII
\texttt{[plutôt deux fois} $(-1)^{r-1}$\texttt{]} où $r$ est la longueur utile de $\tau$ (la longueur pas utile étant $n = \sum\limits_{i\in[1,r]} l_i$).

Quand on part d'une homotopie entre \textit{foncteurs}
\[
X \longrightarrow Y,
\]
il est entendu qu'il s'agit d'homotopies dans la catégorie $\uHom(X,Y)$. (Ou, suivant le contexte, dans une sous-catégorie précisée de $\uHom(X,Y)$.)

\msp{page 77}
À la page 61, j'ai invoqué l'existence d'un \textit{foncteur} de \textit{`composition des chemins'}
\[
\Ch_{\Sigma}(X;a,b)\times\Ch_{\Sigma}(X;b,c) \ \longrightarrow\ \Ch_{\Sigma}(X;a,c),
\]
qui sur la fibre du premier membre au dessus de $(\tau,\tau') \in \operatorname{Ob}(\Ty\times\Ty)$, serait donné par
\[
\Ch_{\Sigma,\tau}(X;a,b)\times\Ch_{\Sigma,\tau'}(X;b,c) \ \longrightarrow\ \Ch_{\Sigma,\tau\tau'}(X;a,c), \tag*{$(*)$}
\]
la composition des chemins de types précisés $\tau$, $\tau'$. Qu'on ait là un \textit{foncteur,} pas seulement une application sur les ensembles d'objets, est assez évident, et peut s'exprimer par exemple à l'aide du diagramme d'amalga\texttt{[ma]}tion
\[
\begin{tikzcd}[column sep=small]
 & \varepsilon = \{0,1\} \ar[dl,hook'] \ar[dr,hook] & \\
I_{\tau} \ar[dr,hook] & & I_{\tau'} \ar[dl,hook'] \\
 & I_{\tau\tau'} &
\end{tikzcd}
\]
en le transformant par $\uHom(\cdot,X)$ (au lieu de le transformer par $\mathrm{Hom}_{\mathrm{Cat}}(\cdot,X)$ seulement), du moins dans le cas particulier où $\Sigma = \mathrm{Fl}(X)$ (qui était celui de la page 61). Mais le passage aux flèches est assez évident - pour deux composés $c'_1\circ c_1$ et $c'_2\circ c_2$, de
\[
\begin{tikzcd}[column sep=small]
a \ar[rr,"c_1"] \ar[d,equal] & {} \ar[d,"w"'] & b \ar[rr,"c'_1"] \ar[d,equal] & {} \ar[d,"w'"] & c \ar[d,equal] \\
a \ar[rr,"c_2"] & {} & b \ar[rr,"c'_2"] & {} & c
\end{tikzcd}
\]
\texttt{[on]} voit que les homomorphismes $w''$ du premier dans le second, dans $\Ch_{\Sigma}(X;a,c)$, qui sont l'identité sur le sommet d'indice $n+1$ des deux chemins (où $n = \operatorname{long}\tau$), correspondant à $1_{\tau}$ et à $0_{\tau'}$, `sont' exactement les couples
% foot: Version 22 décembre 2001
\printedpage{51}
% head: A. Grothendieck : Dérivateurs, Ch. VII, Édition des Universités de Montpellier II et Paris VII
\msp{page 78}
formés d'un homomorphisme $w : c_1 \longrightarrow c_2$ et $w' : c'_1 \longrightarrow c'_2$ dans $\Ch_{\Sigma,\tau}(a,b)$ et $\Ch_{\Sigma,\tau'}(b,c)$ respectivement, et si $w''$, $w$, $w'$ se correspondent ainsi, on posera
\[
w'' = w * w'.
\]
Il y a une transitivité évidente à vérifier pour voir que ça définit un \textit{foncteur} $(*)$. On peut exprimer la situation à l'aide du diagramme
\[
\begin{tikzcd}[column sep=small]
 & \Ty\times\Ty \ar[rrr,"\pi\ (\text{composition})"] \ar[dl,"p_1"'] \ar[dr,"p_2"] & & & \Ty \\
\Ty & & \Ty, & &
\end{tikzcd}
\]
où $p_1$, $p_2$ sont les deux projections. On considère sur $\Ty\times\Ty$ les trois catégories fibrés % [sic] fibrés -> fibrées, agreement with "catégories"
\[
p_1^*(\CH_{\Sigma,\tau}(a,b)),\ \ p_2^*(\CH_{\Sigma,\tau'}(b,c)),\ \ \pi^*(\Ch_{\Sigma,\tau\tau'}(a,c))
\]
\texttt{[plutôt} $\pi^*(\CH_{\Sigma,\tau\tau'}(a,c)$\texttt{]}, dont les fibres en $(\tau,\tau')$ sont respectivement % [sic] \pi^*(\CH_{\Sigma,\tau\tau'}(a,c) -> \pi^*(\CH_{\Sigma,\tau\tau'}(a,c)), one closing parenthesis missing
\[
\Ch_{\Sigma,\tau}(a,b),\ \ \Ch_{\Sigma,\tau'}(b,c),\ \ \Ch_{\Sigma,\tau\tau'}(a,c).
\]
(On écrit $\CH$ au lieu de $\Ch$, pour indiquer que l'on considère les catégories envisagées comme des catégories (fibrées) \textit{au dessus de} $(\Ty)$ (\footnote{Mais il s'agit pour $p_1^*$, $p_2^*$ et $\pi^*$ de produits cartésiens ordinaires dans Cat, il n'y a pas à avoir peur et à écrire des $\CH$ \ldots}).) Au dessus de $\Ty\times\Ty$, on a un homomorphisme (cartésien) de catégories fibrées
\[
p_1^*(\CH_{\tau}(a,b))\times_{\Ty\times\Ty} p_2^*(\CH_{\tau}(b,c)) \ \longrightarrow\ \pi^*(\CH_{\tau}(a,c)) \tag*{$(**)$}
\]
\texttt{[plutôt} $p_1^*(\CH_{\tau}(a,b))\times_{\Ty\times\Ty} p_2^*(\CH_{\tau'}(b,c)) \ \longrightarrow\ \pi^*(\CH_{\tau\tau'}(a,c))$\texttt{]},
\msp{page 79}
car c'est OK sur les fibres ($(*)$ page 77), et si
\[
(\varphi,\varphi') : (\tau_1,\tau'_1) \longrightarrow (\tau,\tau') \qquad (\varphi : \tau_1 \longrightarrow \tau,\ \varphi' : \tau'_1 \longrightarrow \tau')
\]
est une flèche dans $\Ty\times\Ty$, on en déduit un diagramme commutatif
\[
\begin{tikzcd}
\Ch_{\Sigma,\tau}(a,b)\times\Ch_{\Sigma,\tau'}(b,c) \ar[r,"\text{comp.}"] \ar[d,"\varphi^*_{a,b}\times(\varphi'_{b,c})^*"] & \Ch_{\Sigma,\tau}(a,c) \ar[d,"(\varphi*\varphi')^*_{a,c}"] \\
\Ch_{\Sigma,\tau_1}(a,b)\times\Ch_{\Sigma,\tau'_1}(b,c) \ar[r,"\text{comp.}"] & \Ch_{\Sigma,\tau_1\tau'_1}(a,c)
\end{tikzcd}
\]
% foot: Version 22 décembre 2001
\printedpage{52}
% head: A. Grothendieck : Dérivateurs, Ch. VII, Édition des Universités de Montpellier II et Paris VII
\texttt{[plutôt} $\Ch_{\Sigma,\tau\tau'}(a,c)$\texttt{]}, i.e. la composition des chemins d'extrémités fixés est compatible avec les images inverses de chemins par flèches entre les types de chemins.

Mais l'homomorphisme cartésien $(**)$ de catégories fibrées sur $\Ty\times\Ty$ n'est pas tout à fait la même chose qu'un anodin foncteur
\[
\Ch_{\Sigma}(a,b)\times\Ch_{\Sigma}(b,c) \ \longrightarrow\ \Ch_{\Sigma}(a,c). \tag*{$(*)$}
\]
La formule $(*)$ définit en tout cas une application sur les ensembles d'objets, puisque les ensembles d'objets de chacune des trois catégories envisagées se décompose en somme d'ensembles $\Ch_{\Sigma,\tau}(\bullet,\bullet)$, où $\tau$ parcourt $\mathrm{Ty} = \operatorname{Ob}\Ty$. Mais cela se prolonge-t-il de façon naturelle aux morphismes?

\msp{page 80}
Mais voilà que je me rappelle que ce fameux produit cartésien $p_1^*(\text{fourbi})\times p_2^*(\text{machin})$ n'est autre que le produit cartésien tout bête
\[
\Ch_{\Sigma}(a,b)\times\Ch_{\Sigma}(b,c),
\]
envoyé sur $\Ty\times\Ty$ par le produit cartésien des deux foncteurs de projection
\[
\begin{aligned}
p_{a,b} : \Ch_{\Sigma}(a,b) &\ \longrightarrow\ \ \Ty\\
p_{b,c} : \Ch_{\Sigma}(b,c) &\ \longrightarrow\ \ \Ty,
\end{aligned}
\]
d'où
\[
\begin{tikzcd}
\Ch_{\Sigma}(a,b)\times\Ch_{\Sigma}(b,c) \ar[d,"p_{a,b}\times p_{b,c}"] \\
\Ty\times\Ty
\end{tikzcd}
\]
\msp{page 81}
qui est bel et bien fibrant (comme produit cartésien de deux foncteurs fibrants) et a les fibres qu'il faut! On trouve donc bien, par le sorite général sur les foncteurs fibrants (ici sur $\Ty\times\Ty$) un foncteur de catégories fibrées sur $\Ty\times\Ty$
\[
\Ch_{\Sigma}(a,b)\times\Ch_{\Sigma}(b,c) \ \longrightarrow\ \qquad \underbrace{\pi^*(\CH_{\Sigma}(a,c))}_{=\,(\Ty\times\Ty)\times_{\Ty}\Ch_{\Sigma}(a,c)\ \ \text{(\footnotemark)}} \qquad .
\]
\footnotetext{produit fibré par $\pi : \Ty\times\Ty \longrightarrow \Ty$.}
Comme on a d'autre part un foncteur du dernier membre (produit fibré) vers le facteur $\Ch_{\Sigma}(a,c)$, on gagne. (On l'aurait retrouvé par le raisonnement direct que je m'apprêtais de faire, et que j'ai biffé - mais je suis plus heureux comme ça!)

% foot: Version 22 décembre 2001
\printedpage{53}
% head: A. Grothendieck : Dérivateurs, Ch. VII, Édition des Universités de Montpellier II et Paris VII
Pour bien faire, il faudrait prouver l'\textit{associativité} des foncteurs de composition, pas seulement au niveau des objets, par des arguments comme ceci, mais ça revient simplement à une relation du type
\[
(w''*w')*w = w''*(w'*w)
\]
pour un système de trois flèches dans
\[
\Ch_{\Sigma}(a,b),\ \Ch_{\Sigma}(b,c)\ \ \text{et}\ \ \Ch_{\Sigma}(c,d)
\]
respectivement. Ça ne peut être que trivialement
\msp{page 82}
vrai, mais fastidieux à expliciter.

Les sueurs froides sont terminées!

En résumé, le formalisme des $\Ch_{\Sigma}(X)$ et des $\Ch_{\Sigma}(X;a,b)$ développé jusqu'à présent donne satisfaction complète sur les points suivants :

\begin{enumerate}
\item[a)] Bon formalisme de composition des chemins (qui était dégueulasse avant).

\item[b)] Traitement élégant de la construction des catégories de fractions $\mathcal{M}\Sigma^{-1}$. Mais je m'attends que ce sera moins joli qu'avant, quand il s'agira de remplacer un $(\mathcal{M},\Sigma)$ par des $(\ \underbrace{\mathcal{M}(X)}_{=\,\uHom(X,\mathcal{M})}\ ,\Sigma_X)$ ($X$ une catégorie d'indices) \texttt{[plutôt} $\uHom(X^{o},\mathcal{M})$\texttt{]}. C'est pour y pallier qu'il me faudra développer des énoncés de comparaison, entre la catégorie $\Ch_{\Sigma}$ construite ici, et d'autres dont le type dans Hot se prête aux calculs de façon plus directe, comme dans l'état premier de la question. De plus, on sera sûrement obligé d'en venir à des Ind-catégories, pour pouvoir suivre la variance de $\mathcal{M}(X)\Sigma_X^{-1}$ pour $X$ variable.

\item[c)] Bon formalisme des `homotopies' et des compositions des homotopies (ce qui est une autre façon de parler des chemins et de leur composition).
\end{enumerate}
% foot: Version 22 décembre 2001


\endgroup
% END UNIT Der-ChVII-456/Der-ChVII-456_p051-055.tex
% BEGIN UNIT Der-ChVII-456/Der-ChVII-456_p056-060.tex
\begingroup\makeatletter%
\newcommand{\printedpage}[1]{\marginpar{\footnotesize p.~#1}}%
\newcommand{\msp}[1]{\marginpar{\footnotesize [#1]}}%
\newcommand{\uHom}{\underline{\mathrm{Hom}}}%
\newcommand{\Ty}{\mathrm{Ty}}%
\newcommand{\uTy}{\underline{\mathrm{Ty}}}%
\newcommand{\ul}[1]{\underline{#1}}%
\newcommand{\red}{\text{\upshape réd}}%
\newcommand{\triv}{\mathrm{triv}}%
\makeatother\let\printedpage\derpp\let\msp\dermsp%
\let\mspin\DERmspin%


\printedpage{54}
% head: A. Grothendieck : Dérivateurs, Ch. VII, Édition des Universités de Montpellier II et Paris VII
% foot: Version 22 décembre 2001

\msp{page 83}

\section*{8\quad Le lemme d'homotopie pour les morphismes d'intervalles}

C'est le suivant :

\textbf{Lemme 1.} \textit{Soient $\tau,\tau'\in\mathrm{Ty}$, et $f,g:\tau\rightrightarrows\tau'$ deux morphismes dans $\uTy$. Alors $f$, $g$ sont homotopes dans la catégorie $\uHom\,\underline{\mathrm{int}}(I_\tau,I_{\tau'})$ des morphismes croissants d'intervalles orientés de $I_\tau$ dans $I_{\tau'}$.}

Cette catégorie pourrait licitement se désigner par $\uHom_{\uTy}(\tau,\tau')$, ou $\uHom(\tau,\tau')$. C'est une sous-catégorie de la catégorie $\uHom(I_\tau,I_{\tau'})$ de \textit{tous} les foncteurs de $I_\tau$ dans $I_{\tau'}$, et une sous-catégorie pleine de la catégorie $\uHom^\varepsilon(I_\tau,I_{\tau'})$ des foncteurs $f$ qui sont compatibles avec la donnée des sources et des buts, i.e. qui rendent commutatif le diagramme
\[
\begin{tikzcd}[column sep=small]
 & \varepsilon=\{0,1\} \arrow[dl, hook', "{(\sigma_\tau,\beta_\tau)}"'] \arrow[dr, hook, "{(\sigma_{\tau'},\beta_{\tau'})}"] & \\
I_\tau \arrow[rr, "f"] & & I_{\tau'}.
\end{tikzcd}
\]

Si $f$, $g$ sont dans cette catégorie (pas forcément morphismes croissants d'ailleurs), une flèche de $f$ dans $g$ est donc la donnée, pour tout $i\in I_\tau$, d'une flèche $f(i)\xrightarrow{w_i}g(i)$ dans $I_{\tau'}$, avec les deux conditions :

\begin{enumerate}[label={}, leftmargin=*]
\item[a)] Si $i\in\{0_\tau,1_\tau\}$, $w_i=\mathrm{id}$, et
\item[b)] le carré commutatif habituel pour une flèche $i\longrightarrow j$ dans $I_\tau$. (Il suffit d'avoir cette compatibilité pour les flèches de transition $\ul{u}_i:i-1\longrightarrow i$ dans $I_\tau$ ($i\in[1,\operatorname{long}\tau]$).)
\end{enumerate}

Une autre façon de décrire cette catégorie $\uHom^\varepsilon$, c'est de dire que c'est la catégorie
\msp{page 84}
\[
\underline{\mathrm{Ch}}_\tau(I_{\tau'};0_{\tau'},1_{\tau'})\qquad(\simeq\uHom^\varepsilon(I_\tau,I_{\tau'})),
\]
mais $\uHom_{\uTy}(\tau,\tau')$ est la sous-catégorie pleine de cette catégorie formée des \textit{chemins croissants} dans $I_{\tau'}$. Le lemme d'homotopie signifie aussi que cette catégorie $\uHom_{\uTy}(\tau,\tau')$ est toujours \textit{connexe.} Elle n'est d'ailleurs pas forcément 0-connexe, car elle peut être vide. Notons ici, pour plus tard, la

\textbf{Proposition 1 :} \textit{Soient $\tau,\tau'\in\mathrm{Ty}$, soient $r$, $r'$ leurs longueurs utiles. Pour que $\mathrm{Hom}_{\uTy}(\tau,\tau')$ soit non vide, il faut et il suffit que l'une ou l'autre des deux conditions suivantes soit satisfaite :}

\begin{enumerate}[label={}, leftmargin=*]
\item[\textit{a)}] $r\geq r'+1$

\printedpage{55}
% head: A. Grothendieck : Dérivateurs, Ch. VII, Édition des Universités de Montpellier II et Paris VII
% foot: Version 22 décembre 2001

\item[\textit{b)}] \textit{$r=r'$, et (si $r\neq0$) $\tau$ et $\tau'$ ont même direction initiale (de sorte que $\tau_{\red}=\tau'_{\red}$).}
\end{enumerate}

(Démonstration $\pm$ immédiate.)

Notons que si $\tau$, $\tau'$ n'ont pas même direction initiale, alors tout morphisme $\tau\longrightarrow\tau'$ est forcément dégénéré sur le premier intervalle critique de $I_\tau$. On aura donc $\tau=\Delta_\varepsilon^{l_0}\cdot\tau_0$ (où $l_0$ est la longueur du premier intervalle critique, et où $\tau_0$ a \texttt{[la]} direction initiale opposée de celle de $\tau$), et un isomorphisme
\[
\uHom(\tau,\tau')\simeq\uHom(\tau_0,\tau'),
\]
où cette fois $\tau_0$, $\tau'$ ont même direction initiale. On supposera donc désormais que \textit{la direction initiale de $\tau$, $\tau'$ est la même.} La condition $\uHom(\tau,\tau')\neq\emptyset$ équivaut maintenant à $\underline{r\geq r'}$, ce que nous pouvons supposer. Je vais alors décrire un élément canonique
\msp{page 85}
$F_0$ de $\uHom(\tau,\tau')$, ainsi. Soient
\[
I_1,I_2,\dots,I_r
\]
les intervalles critiques de $I_\tau$, et
\[
I'_1,I'_2,\dots,I'_{r'}\qquad(r'\leq r)
\]
ceux de $I_{\tau'}$. Le foncteur $F_0$ proviendra de morphismes d'intervalles orientés (comme
\[
\left\{
\begin{array}{ccc}
I_1 & \xrightarrow{f_1} & I'_1\\[4pt]
I_2 & \xrightarrow{f_2} & I'_2\\[4pt]
 & \dots & \\[4pt]
I_{r'} & \xrightarrow{f_{r'}} & I'_{r'}
\end{array}
\right.\qquad)
\]
et sera donc forcément trivial sur les intervalles critiques $I_\alpha$ pour $r'<\alpha\leq r$ (il y en a si et seulement si $r>r'$). Comme $I_\tau$, $I_{\tau'}$ ont même direction initiale, il s'ensuit que pour $\alpha\in[1,r']$, $I_\alpha$ et $I'_\alpha$ ont même signature (à savoir $\varepsilon_\alpha=(-1)^{\alpha+1}\varepsilon_1$, où $\varepsilon_1$ est la signature commune de $I_1$ et $I'_1$). Donc
\[
I_\alpha=\Delta^{l_\alpha}_{\varepsilon_\alpha},\qquad I'_\alpha=\Delta^{l'_\alpha}_{\varepsilon_\alpha},
\]
et on définit $f_\alpha$ comme
\[
f_\alpha=\mathfrak{f}_{\varepsilon_\alpha,l'_\alpha}*\triv_\alpha,
\]
où
\[
\mathfrak{f}_{\varepsilon_\alpha,l'_\alpha}:\Delta_{\varepsilon_\alpha}\longrightarrow\Delta^{l'_\alpha}_{\varepsilon_\alpha}
\]
et où
\[
\triv_\alpha:\Delta^{l_{\alpha-1}}_{\varepsilon_\alpha}\longrightarrow1. % [sic] l_{\alpha-1} -> l_\alpha-1, the exponent is printed as l with subscript \alpha-1, while I_\alpha = \Delta^{l_\alpha} = \Delta * \Delta^{l_\alpha-1}
\]
En d'autres termes, pour $\varepsilon=+1$ (p. ex.) on doit définir un morphisme des intervalles type $\underbrace{\boldsymbol{\Delta}^+_{l_\alpha}}_{=[0,l_\alpha]}\longrightarrow\underbrace{\boldsymbol{\Delta}^+_{l'_\alpha}}_{=[0,l'_\alpha]}$, et on pose (avec les numérotations internes à $\boldsymbol{\Delta}_{l_\alpha}$, $\boldsymbol{\Delta}_{l'_\alpha}$)
\[
f_\alpha(0)=0,\qquad f_\alpha(i)=l'_\alpha\ \text{ pour }i\in[1,l_\alpha].
\]

\printedpage{56}
% head: A. Grothendieck : Dérivateurs, Ch. VII, Édition des Universités de Montpellier II et Paris VII
% foot: Version 22 décembre 2001

\msp{page 86}
L'idée donc, c'est d'être dans les dispositions d'un voyageur qui, tenu à un plan de voyage par étapes décrit par $\tau$, voudrait arriver aussi vite que possible à $1_{\tau'}$, en partant de $0_{\tau'}$. Il va dès la première étape $\Delta_\varepsilon\subseteq I_1$, se porter jusqu'au bout de $I'_1$, regrettant ne pas pouvoir aller plus loin (à cause du plan de voyage). Il doit rester inactif jusqu'à ce qu'il soit arrivé au moment marqué par l'extrémité finale de $I_1$ (si celle-ci est au delà de celle de sa première flèche de transition, i.e. si $l_1>1$), puis il fait son deuxième grand pas jusqu'à l'extrémité terminale de $I'_2$ etc. Il peut espérer, à chaque moment de son voyage, être en avance (ou du moins pas en retard) sur les autres voyageurs, qui auraient suivi de tactiques de progression différentes. Si cet espoir était justifié, cela signifierait que $F_0$ serait un plus grand élément de $\uHom(\tau,\tau')$ (\textbf{NB} toutes ces catégories $\uHom$ entre intervalles sont ‘ordonnées’, i.e. il y a au plus une flèche entre deux objets), i.e. un objet final. Cela impliquerait que $\uHom(\tau,\tau')$ serait même contractile - chose peut-être vraie, mais sûrement pas de façon si simpliste je crains. Je me suis d'ailleurs canulé : $F_0$ est sans doute bien un ‘plus grand’ élément, mais pour l'ordre défini par la numérotation des sommets de $I_{\tau'}$ (donc $F\leq G$ si et seulement si $F(i)\leq G(i)$ $\forall i\in I_\tau$), et nullement pour l'ordre défini par la structure de catégorie de $\uHom(\tau,\tau')$.

\msp{page 87}
L'idée pour prouver la connexité de $\uHom(\tau,\tau')$, c'est de ramener tout $F\in\mathrm{Hom}(\tau,\tau')$, par une suite d'homotopies simples, à la forme normale. Rappelons que $F$ est donné par une suite croissante de flèches de transition
\[
\ul{u}_{i_1},\ \ul{u}_{i_2},\ \dots\ ,\ \ul{u}_{i_s},\qquad i_1<i_2<\cdots<i_s,
\]
et une suite d'entiers $\geq1$
\[
\lambda_1,\ \lambda_2,\ \dots\ ,\ \lambda_s
\]
avec la condition que
\[
\Delta^{\lambda_1}_{\varepsilon_{i_1}}\Delta^{\lambda_2}_{\varepsilon_{i_2}}\cdots\Delta^{\lambda_s}_{\varepsilon_{i_s}}=\tau'. \tag*{$(*)$}
\]
Les $\ul{u}_{i_\alpha}$ sont les flèches de transition de $I_\tau$ \textit{critiques} pour $F$, i.e. non dégénérées pour $F$, et $\lambda_\alpha$ est la \textit{partie} de $\ul{u}_{i_\alpha}$ pour $F$. Le foncteur normalisé (‘pressé’) $F_0$ est distingué par les conditions suivantes a) et b) \texttt{[et c)]}. Tout d'abord :

\begin{enumerate}[label={}, leftmargin=*]
\item[a)] La suite des signes $\varepsilon_{i_1},\varepsilon_{i_2},\dots,\varepsilon_{i_s}$ est ‘alternée’, i.e. $\varepsilon_{i_\alpha}\neq\varepsilon_{i_{\alpha+1}}$ ($\alpha\in[1,s-1]$). (Cela traduit le fait que le voyageur pressé ne fait pas en plusieurs pas ce qu'il pourrait faire en un seul.)

Comme le premier de ces signes est forcément $\varepsilon'_1=\varepsilon_1$, cela signifie que la suite des signes est la même que celle des $\varepsilon_1$, $\varepsilon_2$, \dots, $\varepsilon_{r'}$ (à la longueur près). Mais mieux, que l'expression $(*)$ de $\tau'$ est son unique expression réduite, donc que $s=r'$ (longueur utile de $\tau'$), et que les $\lambda_i$ sont les longueurs $l'_i$ des intervalles critiques de $I_{\tau'}$.

\printedpage{57}
% head: A. Grothendieck : Dérivateurs, Ch. VII, Édition des Universités de Montpellier II et Paris VII
% foot: Version 22 décembre 2001

\msp{page 88}
Soit donc $\alpha\in[1,s=r']$, et considérons la flèche de transition $F$-critique $\ul{u}_{i_\alpha}$, soit $j_\alpha$ l'indice marquant l'intervalle critique de $I_\tau$ contenant la partie de $\ul{u}_{i_\alpha}$ dans $I_\tau$,
\[
\ul{u}_{i_\alpha}\in\mathrm{Fl}(I_{j_\alpha}),
\]
alors $F$ induit des morphismes (croissants) d'intervalles orientés
\[
F_\alpha=F|I_{j_\alpha}:I_{j_\alpha}\longrightarrow I'_\alpha
\]
et est trivial sur les autres intervalles critiques (ceux pas de la forme $I_{j_\alpha}$). La donnée des $F_\alpha$ détermine donc $F$ sur les $I_{j_\alpha}$, donc sur \textit{tous} les $I_j$, donc partout (par l'opération $*$ de composition par recollement). De plus la condition a) assure que chaque $F_\alpha$ est de nature assez particulière : une seule flèche de transition critique $\ul{u}_{i_\alpha}$. De plus, pour qu'on ait $F=F_0$, il faudra :

\item[b)] $\forall\alpha\in[1,r']$, $\ul{u}_{i_\alpha}$ est la \textit{première} flèche de transition de l'intervalle critique $I_{j_\alpha}$ qui la contient.

Et enfin :

\item[c)] La suite des intervalles critiques de $I_\tau$
\[
I_{j_1},\ I_{j_2},\ \dots\ ,\ I_{j_{r'}}\qquad(\text{de signes alternés }\varepsilon_1,\ -\varepsilon_1,\ \varepsilon_1,\ \ldots)
\]
associés aux flèches $F$-critiques, n'est autre que celle
\[
I_1,\ I_2,\ \dots\ ,\ I_{r'}
\]
des $r'$ premiers intervalles critiques de $I_\tau$, i.e.
\[
(j_1,j_2,\dots,j_{r'})=(1,2,\dots,r').
\]
\end{enumerate}

\msp{page 89}
Ceci vu, l'idée naturelle pour ‘déformer’ $F$ en $F_0$, c'est de trouver une suite d'homotopies qui transforment $F$ de façon à lui faire acquérir successivement les caractéristiques a), b), c) du voyageur pressé. Au besoin, pour aller plus vite au but, on utilisera l'hypothèse de récurrence sur $n=\operatorname{long}\tau$, que le théorème de connexité est déjà prouvé quand $I$ est de longueur $\leq n-1$. (NB Ce théorème est trivial pour $n\leq1$, car alors les catégories $\uHom(\tau,\tau')$ sont toutes ou bien vides, ou bien ponctuelles.)

\msp{page 90}
1) On veut assurer a). Pour ceci, on utilise l'hypothèse de récurrence, pour remplacer un groupe de \texttt{[mot illisible]} flèches critiques consécutives de même signe, par une seule flèche critique. Cette méthode ne foire que si les flèches critiques sont toutes de même signe $\varepsilon$ et incluent $\ul{u}_1$ et $\ul{u}_n$. Mais l'hypothèse de récurrence appliquée à $I=[1,n-1]\subseteq I_\tau$ nous ramène au cas où il n'y a que deux flèches critiques $\ul{u}_1$ et $\ul{u}_n$. Si le premier intervalle critique n'est pas réduit à $\Delta_\varepsilon$, i.e. soit de longueur 2, alors l'hypothèse de récurrence appliquée à $I=[1,n]\subseteq I_\tau$ nous ramène au cas où les deux flèches critiques sont $\ul{u}_1$ et $\ul{u}_2$, et on est
\printedpage{58}
% head: A. Grothendieck : Dérivateurs, Ch. VII, Édition des Universités de Montpellier II et Paris VII
% foot: Version 22 décembre 2001
\msp{page 91}
alors ramené au cas où $\boxed{I=\Delta^2_\varepsilon,\ \tau'=\Delta^2_\varepsilon.}$ Sinon, au cas où $\boxed{\tau=\Delta_\varepsilon\Delta^{l_2}_{-\varepsilon}\Delta_\varepsilon,\ \tau'=\Delta^2_\varepsilon,}$ et où les deux flèches critiques sont encore la première et la dernière, et $\tau'=\Delta^2_\varepsilon$. Ce sont là les deux cas irréductibles, qu'il faudra traiter à part.

Ces cas admis, on a donc moyen pour définir $F$ de façon à satisfaire a). Ensuite on veut :

2) Assurer que la suite des intervalles critiques $I_{j_1}$, $I_{j_2}$, \dots, $I_{j_{r'}}$ qui contient les flèches critiques $\ul{u}_{i_1}$, \dots, $\ul{u}_{i_{r'}}$, soit la suite $I_1$, \dots, $I_{r'}$. C'est encore immédiat par l'hypothèse de récurrence, le seul cas irréductible étant celui où $r'=1$ (une seule flèche critique), et où il s'agit de faire descendre cette flèche critique jusqu'à $\ul{u}_1$. Si $I_1$ est de longueur $>1$, alors on fait descendre en $\ul{u}_2$ par hypothèse de récurrence, et le cas irréductible est celui où $\boxed{\tau=\Delta^2_\varepsilon,\ \tau'=\Delta_\varepsilon,}$ une flèche critique $\ul{u}_2$, la faire descendre en $\ul{u}_1$. Si par contre $I_1$ \texttt{[est]} de longueur 1, le cas irréductible est $\boxed{\tau=\Delta_\varepsilon\Delta^{l_2}_{-\varepsilon}\Delta_\varepsilon,\ \tau'=\Delta_\varepsilon,}$ une flèche critique $\ul{u}_{l_2+2}$, la faire descendre en $\ul{u}_1$.

Sous réserve de traiter les deux nouveaux cas irréductibles, on peut déformer $F$ en un foncteur qui satisfait à a) et à c). Reste à normaliser chacun des $F_\alpha:I_\alpha\longrightarrow I'_\alpha$ ($\alpha\in[1,r']$). Donc on est ramené à $\Delta^n_\varepsilon\longrightarrow\Delta_\varepsilon$ : une seule flèche critique $\ul{u}_\alpha$, $\alpha\in[1,n]$, à faire descendre en $\ul{u}_1$. Si $n\geq2$, par hypothèse de récurrence elle descend en $\ul{u}_2$, cas supposé traité en $2^o$). Donc
\msp{page 92}
on gagne, sous réserve de traiter les 4 cas irréductibles.

A) Cas $\tau=\Delta^2_\varepsilon$, $\tau'=\Delta^2_\varepsilon$ ou $\Delta_\varepsilon$.

La catégorie $\uHom(\Delta^2_\varepsilon,\Delta_\varepsilon)$ a deux éléments, les morphismes $\Delta^2_\varepsilon\longrightarrow\Delta_\varepsilon$ correspondent aux flèches de transition de $\Delta^2_\varepsilon$ ($=\boldsymbol{\Delta}^+_2$ ou $\boldsymbol{\Delta}^-_2$). Dans le cas $\varepsilon=+1$, ce sont les deux applications de $\{0,1,2\}$ dans $\{0,1\}$, qui appliquent 0 en 0, 2 en 1, et 1 soit en 1 ($F_0$), soit en 1 \texttt{[plutôt en 0]} ($F_1$).
\[
\begin{tikzcd}
0 \arrow[r] & 1 \arrow[r] & 2
\end{tikzcd}
\]
\[
\begin{tikzcd}
0 \arrow[r] & 1
\end{tikzcd}
\]

On a évidemment une flèche $F_1\longrightarrow F_0$ (le voyageur pressé arrive plus vite \ldots). Similaire pour $\varepsilon=-1$,
\[
\uHom(\Delta^2_\varepsilon,\Delta_\varepsilon)\simeq\boldsymbol{\Delta}^1,\qquad\text{catégorie 0-connexe.}
\]

La catégorie $\uHom(\Delta^2_\varepsilon,\Delta^2_\varepsilon)$ a trois éléments, s'identifiant à l'ensemble des parties non vide de l'ensemble $\mathrm{Tr}(\tau)$ à deux éléments des flèches de transition $\ul{u}_1$, $\ul{u}_2$ de $\tau$ : il y a $F_{\ul{u}_1}=F_0$, $F_{\ul{u}_2}$, $F_{\ul{u}_1,\ul{u}_2}$. Si $\varepsilon=+1$, il s'agit des applications croissantes (sans plus) de $\{0,1,2\}$ dans $\{0,1,2\}$, appliquant 0 en 0, 2 en 2, et 1 en ce qu'elle voudra. % [sic] parties non vide -> parties non vides, agreement


\endgroup
% END UNIT Der-ChVII-456/Der-ChVII-456_p056-060.tex
% BEGIN UNIT Der-ChVII-456/Der-ChVII-456_p061-065.tex
\begingroup\makeatletter%
\newcommand{\printedpage}[1]{\marginpar{\footnotesize p.~#1}}%
\newcommand{\msp}[1]{\marginpar{\footnotesize [#1]}}%
\newcommand{\uHom}{\underline{\mathrm{Hom}}}%
\newcommand{\uCh}{\underline{\mathrm{Ch}}}%
\newcommand{\uTy}{\underline{\mathrm{Ty}}}%
\makeatother\let\printedpage\derpp\let\msp\dermsp%
\let\mspin\DERmspin%


\printedpage{59}
% head: A. Grothendieck : Dérivateurs, Ch. VII, Édition des Universités de Montpellier II et Paris VII
% foot: Version 22 décembre 2001

\begin{tabular}{lll}
Donc & $1\longmapsto 2$ & correspond à \quad $F_0$\\[1ex]
 & $1\longmapsto 1$ & correspond à \quad $F_1$\\[1ex]
 & $1\longmapsto 0$ & correspond à \quad $F_2$
\end{tabular}

et on a bien sur $F_2\longrightarrow F_1\longrightarrow F_0$, % [sic] bien sur -> bien sûr, accent missing
\[
\uHom(\Delta^2_\varepsilon,\Delta^2_\varepsilon)\simeq\boldsymbol{\Delta}^2,\qquad\quad\text{catégorie 0-connexe.}
\]

Dans ces deux cas, deux éléments de $\mathrm{Hom}(\tau,\tau')$ sont reliés par une homotopie \textit{simple.}

\msp{page 93}

B) Cas $\tau=\Delta_\varepsilon\Delta^l_{-\varepsilon}\Delta_\varepsilon$, $\tau'=\Delta^2_\varepsilon$ ou $\Delta_\varepsilon$.

Notons que l'on a, pour $\tau,\tau'\in\uTy$ quelconque, un isomorphisme de catégories
\[
\uHom(\tau,\tau')\longrightarrow\uHom(\tau'^o,\tau^o),
\]
exprimant simplement que
\[
\tau\longmapsto\tau^o
\]
est un \texttt{[2]}-foncteur de $\uTy$ dans $\uTy$, et même un automorphisme de cette \texttt{[2]}-catégorie. Donc dire que $\uHom(\tau,\tau')$ est connexe revient à dire que $\uHom(\tau^o,\tau'^o)$ l'est. Donc on peut supposer ci-dessus que $\varepsilon=+1$ (et pareil dans A!).
\[
\begin{aligned}
\tau :\quad & \longrightarrow\ \underbrace{\longleftarrow\ \longleftarrow\ \cdots\ \longleftarrow}_{l}\ \longrightarrow\\
\tau' :\quad & \longrightarrow\ \longrightarrow\quad\text{ou}\quad\longrightarrow\ .
\end{aligned}
\]

Mais tous les morphismes $\tau\longrightarrow\tau'$ sont triviaux sur l'intervalle critique $I_2$ (rétrograde), ce qui implique que
\[
\uHom(\tau,\tau')\xleftarrow{\ \sim\ }\uHom(\Delta^2_\varepsilon,\tau'),
\]
via le foncteur de contraction canonique
\[
I_\tau\longrightarrow I_{\Delta^2_\varepsilon}=\boldsymbol{\Delta}^2.
\]

Or dans A on a vu que $\uHom(\Delta^2_\varepsilon,\tau')$ est connexe (pour $\tau'$ de la forme $\Delta_\varepsilon$ ou $\Delta^2_\varepsilon$).

On a gagné!

\par
\printedpage{60}
% head: A. Grothendieck : Dérivateurs, Ch. VII, Édition des Universités de Montpellier II et Paris VII
% foot: Version 22 décembre 2001

\msp{page 94}

\section*{9\quad Théorème d'homotopie pour les catégories de chemins}

Soient $X$ une catégorie, $\Sigma\subseteq\mathrm{Fl}(X)$ \texttt{[un]} ensemble admissible de flèches dans $X$, $a,b\in\mathrm{Ob}\,X$. Considérons pour $\tau\in\mathrm{Ty}$ variable, les catégories
\[
\uCh_{\Sigma,\tau}(X),\qquad\quad\uCh_{\Sigma,\tau}(X;a,b),
\]
qui forment les fibres des catégories Cat-fibrées
\[
\begin{tikzcd}[column sep=small]
\uCh_\Sigma(X) \arrow[dr] & & \uCh_\Sigma(X;a,b) \arrow[ll, "\text{incl.}"'] \arrow[dl]\\
 & \uTy &
\end{tikzcd}
\]
sur la catégorie $\uTy$. Ainsi, pour $\tau,\tau'\in\mathrm{Ty}$, et pour toute
\[
\varphi:\tau\longrightarrow\tau',
\]
on a un carré commutatif d'homomorphismes
\[
\begin{tikzcd}
\uCh_{\Sigma,\tau'}(X) \arrow[d, "\varphi^*_X"'] & \uCh_{\Sigma,\tau'}(X;a,b) \arrow[l, "\text{incl.}"'] \arrow[d, "\varphi^*_{a,b}"]\\
\uCh_{\Sigma,\tau}(X) & \uCh_{\Sigma,\tau}(X;a,b), \arrow[l, "\text{incl.}"']
\end{tikzcd}
\]
avec transitivité pour $\varphi$ variable (composition $\psi\circ\varphi$ de flèches dans $\uTy$).

En fait, pour $\varphi$ variable dans $\uHom(\tau,\tau')$, les foncteurs images inverses $\varphi^*_X$, $\varphi^*_{a,b}$ dépendent \textit{fonctoriellement} de $\varphi$; i.e. on a des \textit{foncteurs}
\[
\begin{tikzcd}[row sep=small]
 & \uHom(\uCh_{\Sigma,\tau'}(X),\uCh_{\Sigma,\tau}(X))\\
\uHom(\tau,\tau') \arrow[ur] \arrow[dr] & \\
 & \uHom(\uCh_{\Sigma,\tau'}(X;a,b),\uCh_{\Sigma,\tau}(X;a,b))
\end{tikzcd}
\tag*{$(*)$}
\]

\par
\printedpage{61}
% head: A. Grothendieck : Dérivateurs, Ch. VII, Édition des Universités de Montpellier II et Paris VII
% foot: Version 22 décembre 2001

\msp{page 95}

Cela provient de l'inclusion
\[
\uHom(\tau,\tau')\hookrightarrow\uHom^{\varepsilon}(I_\tau,I_{\tau'})
\]
et de la description des catégories $\uCh_\tau$, $\uCh_{\tau'}$ comme catégories de foncteurs, telle
\[
\uCh_{\Sigma,\tau}(X)\simeq\uHom_{\Sigma_\tau,\Sigma}(I_\tau,X),
\]
où $\Sigma_\tau$ désigne l'ensemble des flèches dans $I_\tau$ qui sont soit inversibles, soit \textit{rétrogrades.} (Cet ensemble de flèches est admissible, car stable par composition.) On désigne par $\uHom_{\Sigma_\tau,\Sigma}$ la sous-catégorie pleine de $\uHom$ formée des foncteurs qui appliquent $\Sigma_\tau$ dans $\Sigma$. Or de façon générale, si
\[
(I,\Sigma_I),\ \ (I',\Sigma_{I'}),\ \ (X,\Sigma)
\]
sont trois catégories, munies d'ensembles admissibles de flèches, considérant les foncteurs $\varphi:I\longrightarrow I'$ avec $\varphi(\Sigma_I)\subseteq\Sigma_{I'}$, alors le foncteur familier
\[
\uHom(I,I')\longrightarrow\uHom(\uHom(I',X),\uHom(I,X))
\]
(transposé du foncteur de composition
\[
\begin{matrix}
\uHom(I,I') & \times & \uHom(I',X) & \longrightarrow & \uHom(I,X) & \\
(\varphi & , & F') & \longmapsto & F'\circ\varphi & ),
\end{matrix}
\]
se généralise en un foncteur
\[
\uHom_{\Sigma_I,\Sigma_{I'}}(I,I')\longrightarrow\uHom(\uHom_{\Sigma_{I'},\Sigma}(I',X),\uHom_{\Sigma_I,\Sigma}(I,X)),
\tag*{$(*)$}
\]
transposé du foncteur de composition
\[
\uHom_{\Sigma_I,\Sigma_{I'}}(I,I')\times\uHom_{\Sigma_{I'},\Sigma}(I',X)\longrightarrow\uHom_{\Sigma_I,\Sigma}(I,X)
\]
induit par $(*)$.

\msp{page 96}

Comme toute
\[
\varphi:\tau\longrightarrow\tau'
\]
induit
\[
\varphi:I_\tau\longrightarrow I_{\tau'},
\]
transformant flèches rétrogrades en flèches identiques ou rétrogrades, i.e.
\[
\varphi(\Sigma_\tau)\subseteq\Sigma_{\tau'},
\]
\printedpage{62}
% head: A. Grothendieck : Dérivateurs, Ch. VII, Édition des Universités de Montpellier II et Paris VII
% foot: Version 22 décembre 2001
on trouve le premier des deux foncteurs $(*)$ annoncés page 94, qui est associé par transposition au foncteur $\gamma$
\[
\begin{tikzcd}[column sep=large]
\uHom(\tau,\tau')\times\ \underbrace{\uCh_{\Sigma,\tau'}(X)}_{\simeq\,\uHom_{\Sigma_{\tau'},\Sigma}(I_{\tau'},X)} \arrow[r, "\gamma"] \arrow[d, "\text{can.}"'] & \uCh_{\Sigma,\tau}(X) \arrow[d, equal]\\
\uHom_{\Sigma_\tau,\Sigma_{\tau'}}(I_\tau,I_{\tau'})\times\uHom_{\Sigma_{\tau'},\Sigma}(I_{\tau'},X) \arrow[r, dotted, "\gamma_0"] & \uHom_{\Sigma_\tau,\Sigma}(I_\tau,X)
\end{tikzcd}
\tag*{$(*)$}
\]
déduit du foncteur composition $\gamma_0$ dans ce diagramme : $\gamma=\gamma_0\circ\mathrm{can}.$ D'autre part, ce foncteur de composition $\gamma$ induit visiblement un foncteur sur les sous-catégories de chemins à source et but fixés
\[
\uHom(\tau,\tau')\times\uCh_{\Sigma,\tau'}(X;a,b)\longrightarrow\uCh_{\Sigma,\tau}(X;a,b),
\]
d'où la deuxième des deux flèches annoncées dans $(*)$ p. 94.

Cela montre donc que si
\[
\varphi_1,\varphi_2:\tau\rightrightarrows\tau'
\]
sont deux morphismes de $\tau$ dans $\tau'$ homotopes dans $\uHom(\tau,\tau')$, alors les foncteurs induits entre catégories de $\Sigma$-chemins ou entre catégories de $\Sigma$-chemins à extrémités précisés : % [sic] précisés -> précisées, agreement with extrémités

\msp{page 97}
\[
\begin{aligned}
\varphi^*_{1\,X},\varphi^*_{2\,X} : && \uCh_{\Sigma,\tau'}(X) &\quad\rightrightarrows\quad\uCh_{\Sigma,\tau}(X),\\
(\varphi_1)^*_{a,b},(\varphi_2)^*_{a,b} : && \uCh_{\Sigma,\tau'}(X;a,b) &\quad\rightrightarrows\quad\uCh_{\Sigma,\tau}(X;a,b)
\end{aligned}
\]
sont également homotopes.

Or le lemme d'homotopie nous assure que l'hypothèse $\varphi_1\sim\varphi_2$ est \textit{toujours} satisfaite. Donc :

\textbf{Théorème d'homotopie sur les opérations universelles sur les chemins.} \textit{Soient $X$, $\Sigma$, $a$, $b$ comme dessus, et $\tau,\tau'\in\mathrm{Ty}$ telles que $\mathrm{Hom}(\tau,\tau')\neq\emptyset$. Alors les foncteurs $\varphi^*_X$, $\varphi^*_{X;a,b}$ induits par un $\varphi\in\mathrm{Hom}(\tau,\tau')$}
\[
\left\{
\begin{aligned}
\varphi^*_X : && \uCh_{\Sigma,\tau'}(X) &\quad\longrightarrow\quad\uCh_{\Sigma,\tau}(X)\\
\varphi^*_{X;a,b} : && \uCh_{\Sigma,\tau'}(X;a,b) &\quad\longrightarrow\quad\uCh_{\Sigma,\tau}(X;a,b)
\end{aligned}
\right.
\]
\textit{ne dépendent pas de $\varphi$,} à homotopie près.

On peut préciser cet énoncé, de telle façon que l'énoncé pour les $\varphi^*_{X;a,b}$ apparaisse, comme il devrait, comme corollaire pour celui sur $\varphi^*_X$. Pour ceci, on va regarder les catégories de chemins dans $X$ comme catégories \textit{au dessus de $X$,} par les foncteurs `source et but' :
\[
\begin{tikzcd}[column sep=small]
\uCh_{\Sigma,\tau}(X) \arrow[dr, "(\sigma_\tau{,}\beta_\tau)"'] & & \uCh_{\Sigma,\tau'}(X) \arrow[dl, "(\sigma_{\tau'}{,}\beta_{\tau'})"]\\
 & X\times X &
\end{tikzcd}
\]

\par
\printedpage{63}
% head: A. Grothendieck : Dérivateurs, Ch. VII, Édition des Universités de Montpellier II et Paris VII
% foot: Version 22 décembre 2001

\msp{page 98}

Alors le foncteur $(*)$ (p. 94) se factorise par la sous-catégorie $\uHom_{\mathrm{Cat}/X\times X}(\cdot,\cdot)$ de $\uHom_{\mathrm{Cat}}(\cdot,\cdot)$ formée des foncteurs qui commutent aux foncteurs de projection dans $X\times X$, et les homomorphismes entre tels foncteurs qui s'effectuent par des flèches `verticales' de la catégorie but (i.e. par des flèches au dessus de flèches identiques de la catégorie base $X\times X$) :
\[
\uHom(\tau,\tau')\longrightarrow\uHom_{\mathrm{Cat}/X\times X}(\uCh_{\Sigma,\tau'}(X),\uCh_{\Sigma,\tau}(X)).
\]

Donc dans l'énoncé du théorème d'homotopie, on peut préciser que l'homotopie se fait dans cette sous-catégorie, i.e. s'effectue par une suite d'homotopies simples qui sont des morphismes de foncteurs entre catégories au dessus de $X\times X$. Cela implique alors l'énoncé d'homotopie sur les $\varphi^*_{X;a,b}$.


\endgroup
% END UNIT Der-ChVII-456/Der-ChVII-456_p061-065.tex
% BEGIN UNIT Der-ChVII-456/Der-ChVII-456_p066-070.tex
\begingroup\makeatletter%
\NoAutoSpacing%
\newcommand{\printedpage}[1]{\marginpar{\footnotesize p.~#1}}%
\newcommand{\msp}[1]{\marginpar{\footnotesize [#1]}}%
\newcommand{\Cat}{\mathrm{Cat}}%
\newcommand{\Cathot}{\mathrm{Cat\,hot}}%
\newcommand{\Hot}{\mathrm{Hot}}%
\newcommand{\Hom}{\mathrm{Hom}}%
\newcommand{\Ob}{\mathrm{Ob}\,}%
\newcommand{\uTy}{\underline{\mathrm{Ty}}}%
\newcommand{\uCh}{\underline{\mathrm{Ch}}}%
\newcommand{\uHom}{\underline{\mathrm{Hom}}}%
\newcommand{\sHom}{\underline{\mathcal{H}\mathit{om}}}%
\newcommand{\To}{T_{\mathrm{o}}}%
\newcommand{\circled}[1]{\tikz[baseline=(c.base)]\node[draw,circle,inner sep=1pt](c){#1};}%
\makeatother\let\printedpage\derpp\let\msp\dermsp%
\let\mspin\DERmspin%


\leavevmode\printedpage{64}
% head: A. Grothendieck : Dérivateurs, Ch. VII, Édition des Universités de Montpellier II et Paris VII
\msp{page 99}

\section*{10\quad Question de comparaison pour les catégories \textmd{Cat}-fibrées satisfaisant au ‘théorème d'homotopie' (base filtrante)}

Nous allons généraliser maintenant au cas d'une catégorie (Cat)-fibrée sur une autre
\[
\begin{tikzcd}
C \arrow[d, "\pi"'] \\
T,
\end{tikzcd}
\]
en pensant surtout au cas $T = \uTy$, $C = \uCh_{\Sigma}(X;a,b)$. L'hypothèse essentielle est la suivante :

\noindent\begin{tabular}{@{}p{0.3\textwidth}p{0.64\textwidth}@{}}
(Hypothèse d'homotopie) & Si $t,t' \in \Ob T$ sont tels que $\Hom(t,t') \neq \emptyset$, alors les foncteurs $\varphi^{*} : C_{t'} \longrightarrow C_{t}$ associés aux $\varphi \in \Hom(t,t')$ sont homotopes.
\end{tabular}

Ils définissent donc une seule et même flèche dans Cat\,hot. Je désigne par $T_{\mathrm{ord}}$ ou $\To$ l'ensemble $\Ob(T)$, avec comme relation d'ordre celle déduite de la structure de catégorie :
\[
t \leq t' \qquad\Longleftrightarrow\qquad \Hom_{T}(t,t') \neq \emptyset,
\]
de sorte qu'on a un foncteur
\[
T \longrightarrow \To
\]
qui est \textit{universel} pour les foncteurs de $T$ dans des catégories ‘ordonnées' (i.e. où toute flèche est monomorphique \texttt{[non, cf. }$\bullet\rightrightarrows\bullet$\texttt{; plutôt, où toujours }$\#\Hom(a,b) \leq 1$ \texttt{(admettant des isomorphismes non identiques)]}) ou pour les foncteurs $F$ sur $T$ tels que $F$ soit constant sur tout $\Hom(t,t')$. L'hypothèse d'homotopie peut s'exprimer en disant que le foncteur composé
\[
\begin{tikzcd}
T^{o} \arrow[r] & \Cat \arrow[r] & \Cathot \\
t \arrow[r, mapsto] & C_{t} &
\end{tikzcd}
\]
\msp{page 100}
se factorise (évidemment de façon unique) en un foncteur
\[
H_{C} : \To^{o} \longrightarrow \Cathot,
\]
i.e. qu'on a un diagramme commutatif
\printedpage{65}
% head: A. Grothendieck : Dérivateurs, Ch. VII, Édition des Universités de Montpellier II et Paris VII
\[
\begin{tikzcd}
T^{o} \arrow[r, "t\,\longmapsto\, C_{t}"] \arrow[d] & \Cat \arrow[d] \\
\To^{o} \arrow[r, "H_{C}"] & \Cathot.
\end{tikzcd}
\]

Nous allons supposer de plus :

\noindent\begin{tabular}{@{}p{0.36\textwidth}p{0.58\textwidth}@{}}
(Hypothèse de filtration) & L'ensemble ordonné $\To$ associé à $T$ est \textit{filtrant} décroissant (filtrant à gauche), i.e. si $t,t' \in \Ob T$, il existe $t'' \in \Ob T$ et des flèches $t'' \longrightarrow t$, $t'' \longrightarrow t'$.
\end{tabular}

Cette hypothèse est bien satisfaite dans $\uTy$, car pour $\tau,\tau' \in \Ob \uTy = \mathrm{Ty}$, il suffit de prendre $\tau\tau'$, qui s'envoie dans $\tau$ par $\mathrm{id}_{\tau} * p_{\tau'}$, dans $\tau'$ par $p_{\tau} * \mathrm{id}_{\tau'}$, où $p_{\tau} : \tau \longrightarrow 1$ est le morphisme de $\tau$ dans l'objet final.

Ainsi, $H_{C}$ définit un ind-objet de Cat\,hot, que nous allons désigner par $\underline{\underline{C}}$. Si on suppose de plus :

\noindent\begin{tabular}{@{}p{0.36\textwidth}p{0.58\textwidth}@{}}
(Hypothèse de dénombrabilité, toujours satisfaite si $\mathrm{card}\,\To \leq \aleph_{0}$) & $\To$ admet un sous-ensemble \textit{dénombrable} cofinal à gauche.
\end{tabular}

alors cet ind-objet de Cat\,hot peut être défini, de beaucoup de façons, en termes
\msp{page 101}
d'un système inductif dans Cat (pas seulement dans Cat\,hot)
\[
(C_{n})_{n\in\mathbf{N}},
\]
d'où on récupère l'ind-objet $\underline{\underline{C}}$ dans Cat\,hot en passant au système inductif correspondant dans Cat\,hot, puis à l'ind-objet associé.

Pour construire le système inductif $C_{n}$, il suffit de prendre une chaîne infinie
\[
t_{0} \longleftarrow t_{1} \longleftarrow t_{2} \longleftarrow \;\cdots\; \longleftarrow t_{n} \longleftarrow t_{n+1} \longleftarrow \;\cdots
\]
de morphismes dans $T$, telle que l'ensemble des $t_{i}$ soit cofinal à gauche dans $\To$. On va supposer de plus $C$ \textit{strictement fibrée} sur $T$ (i.e. provenant d'un véritable foncteur $T^{o} \longrightarrow \Cat$, i.e. on peut choisir les foncteurs images inverses
\[
\varphi^{*} : C_{t'} \longrightarrow C_{t}
\]
de façon transitive stricte, pas seulement transitive à isomorphisme près) - cas auquel on peut toujours se ramener d'ailleurs par l'astuce de \textsc{Giraud}. Cette hypothèse est d'ailleurs vérifiée dans le cas $\uCh_{\Sigma}(X;a,b)$ sur $\uTy$. Ainsi, les
\[
(C_{t_{i}})_{i\in\mathbf{N}}
\]
forment un système inductif de catégories indexées par $\mathbf{N}$, le morphisme de transition
\[
C_{n} \longrightarrow C_{n'} \qquad \text{i.e. } C_{t_{n}} \longrightarrow C_{t_{n'}}
\]
\printedpage{66}
% head: A. Grothendieck : Dérivateurs, Ch. VII, Édition des Universités de Montpellier II et Paris VII
\msp{page 102}
étant le foncteur image inverse $\varphi^{*}_{n,n'}$, où
\[
\varphi_{n,n'} : t_{n'} \longrightarrow t_{n}
\]
est le composé des $t_{n'} \longrightarrow t_{n'-1} \longrightarrow \;\cdots\; \longrightarrow t_{n}$. (Si on n'avait pas transitivité stricte, on définirait $\varphi^{*}_{n,n'}$ comme le composé des foncteurs image inverse par des flèches de transition successives $t_{n'} \longrightarrow t_{n'-1} \longrightarrow \;\cdots\; \longrightarrow t_{n+1} \longrightarrow t_{n}$, ce ne serait $\varphi^{*}_{n,n'}$ qu'à isomorphisme canonique près.)

Soit
\[
C_{\infty} = \varinjlim_{n} C_{n} = \varinjlim_{n} C_{t_{n}}
\]
la catégorie $\varinjlim$ . (Elle sera surtout sympathique si on a pu choisir le système $(t_{n})_{n\in\mathbf{N}}$ de telle façon que les foncteurs de transition
\[
C_{t_{n}} \longrightarrow C_{t_{n-1}}
\]
soient pleinement fidèles, et peut-être même \texttt{[des]} monomorphismes, de telle façon que les $C_{n}$ s'identifient à des sous-catégories pleines de leur limite inductive $C_{\infty}$.) Je voudrais comprendre les relations précises entre les trois genres d'objets suivants :

\begin{itemize}[leftmargin=2.5em]
\item[\circled{a}] $C$ elle-même (comme objet dans Cat, pas de $\Cat/T$), et surtout les objets de Cat\,hot et de Hot qui lui correspondent.

\item[\circled{b}] \msp{page 103}$C_{\infty}$, et les objets de Cat\,hot et de Hot qui lui correspondent. Dans quelle mesure ces deniers % [sic] deniers -> derniers, missing letter
sont-ils indépendant % [sic] indépendant -> indépendants, agreement with the plural subject
du choix du système projectif cofinal à gauche $(t_{n})_{n\in\mathbf{N}}$ dans $T$? Quel est le rôle de $C_{\infty}$? Est-elle canoniquement isomorphe à $C$ dans Cat\,hot, ou dans Hot? (Auquel cas la question de son ‘rôle' me semblerait résolue!)

\item[\circled{c}] Le ind-objet % [sic] Le ind-objet -> L'ind-objet, missing elision
$H_{C}$ dans Cat\,hot, et l'ind-objet de Hot qu'il définit.

Je note que l'ind-objet $H_{C}$, et aussi les catégories $C_{\infty}$ construites comme dessus, ne changent pas quand on remplace $T$ par une sous-catégorie $T'$ cofinale à gauche, ou aussi par une catégorie de la forme $T/t_{0}$ ($t_{0} \in \Ob T$), qui donne lieu à l'ensemble ordonnée % [sic] ordonnée -> ordonné, agreement with "ensemble"
associé
\[
(T/t)_{\mathrm{o}} = \To/t = \{t \in \To \mid t \leq t_{0}\}
\]
\texttt{[plutôt} $(T/t_{0})_{\mathrm{o}} = \To/t_{0} = \{t \in \To \mid t \leq t_{0}\}$\texttt{]}. Si la relation entre les objets dans a), et ceux dans b) et dans c), est très étroite, on en concluerait % [sic] concluerait -> conclurait, spelling of the conditional
que le ‘type' de $C$ dans Cat\,hot, ou dans Hot suivant le cas, ne change pas en remplaçant $C$ par
\printedpage{67}
% head: A. Grothendieck : Dérivateurs, Ch. VII, Édition des Universités de Montpellier II et Paris VII
\msp{page 104}
$C \times_{T} T'$ ou par $C/t_{0} \simeq C \times_{T} T/t_{0}$. Bien mieux, elle ne changerait pas en remplaçant $T$ par $(\mathbf{N})^{o}$, la catégorie (ordonnée) opposée à $\mathbf{N}$ (ayant donc un unique objet final), et la catégorie fibrée $C$ sur $T$ par la catégorie $C_{\underline{t}}$ induite par $C$ sur $(\mathbf{N})^{o}$, via le foncteur
\[
\underline{t} : (\mathbf{N})^{o} \longrightarrow T
\]
défini par le système projectif $\underline{t} = (t_{n})_{n\in\mathbf{N}}$ dans $T$. Cela signifierait donc que pour comprendre le type homotopique (au sens strict, i.e. dans Cat\,hot, ou large, i.e. dans Hot) de $C$, on peut le ‘nettoyer' de façon draconienne, la base $T$ étant vue comme une sorte de gangue un peu vaseuse, ou du moins pléthorique.

\item[\circled{d}] De toutes façons, si $T'$ est une autre catégorie filtrante à gauche, et
\[
\alpha : T' \longrightarrow T
\]
un foncteur cofinal à gauche, de sorte que $C$ sur $T$, et $C' = C \times_{T} T'$ sur $T'$, définissent essentiellement le ‘même' objet de Ind\,Cat\,hot, et que la catégorie $C'_{\infty}$ pour $C'$ sur $T'$ fait partie des catégories $C_{\infty}$ pour $C$ sur $T$ - on peut se demander
\msp{page 105}
si le foncteur canonique
\[
C' \;\longrightarrow\; C
\]
est également toujours un isomorphisme dans Cat\,hot, ou sinon, du moins dans Hot.
\end{itemize}

Un intérêt particulier s'attache à $\pi_{0}(C)$, et de plus à sa variance pour $C$ sur $T$ variable de la façon suivante. Si $S$ est dans Cat, on regarde les catégories
\[
\uHom_{\Cat}(S,C_{t}) \qquad (t \in T),
\]
on trouve un foncteur
\[
t \longmapsto \uHom_{\Cat}(S,C_{t}), \qquad T^{o} \longrightarrow (\Cat),
\]
d'où une catégorie fibrée sur $T$, que je vais désigner par
\[
\sHom_{T}(S_{T},C)
\]
(Sauf erreur, elle s'interprète comme la catégorie fibrée sur $T$ correspondante au champ sur $\mathrm{Top}(T) = T^{\wedge}$ du morphisme de $T$-champs de $S_{T}$ dans $C$, où $S_{T} = S \times T$ (considérée comme catégorie fibrée sur $T$ à fibre essentiellement constante $S$) dans $C$. Ici j'identifie catégorie fibrée sur $T$, et champs sur le topos $T^{\wedge}$.) Il est immédiat que la propriété d'homotopie (p. 99) est stable par passage de $C$ à $\sHom_{T}(S_{T},C)$. On s'intéresse à la variation de
\[
\pi_{0}(\sHom_{T}(S_{T},C))
\]
\printedpage{68}
% head: A. Grothendieck : Dérivateurs, Ch. VII, Édition des Universités de Montpellier II et Paris VII
\msp{page 106}
comme contrafoncteur en $S$,
\[
H : S \longmapsto \pi_{0}\sHom_{T}(S_{T},C) \qquad (\Cat)^{o} \longrightarrow (\mathrm{Ens}),
\]
et à une éventuelle expression de ce contrafoncteur $H$ en termes de l'ind-objet $H_{C}$ de Cat\,hot, par la formule hypothétique
\[
H(S) \simeq \Hom_{\mathrm{Ind}(\Cathot)}(\gamma(S),H_{C}),
\]
où
\[
\gamma : \Cat \longrightarrow \Cathot
\]
est le foncteur canonique. Si on avait un tel isomorphisme canonique fonctoriel en $S$, cela donnerait une description (à isomorphisme unique près) de l'ind-objet $H_{C}$ en termes du foncteur précédent $H$ sur Cat.

Il me vient un doute ici, en vue de l'exemple $T = \uTy$, $C = \uCh_{\Sigma}(X;a,b)$. Les trois types d'objets a) b) c) ont le même $\pi_{0}$, et il est sûr que ‘c'est le bon'. D'autre part, pour une catégorie paramètre variable $S$, il est sûr aussi que le bon foncteur en $S$ est
\[
S \longmapsto \pi_{0}(\uHom_{T}(S_{T},C)) \simeq \Hom_{\mathrm{Ind}\,\Cathot}(S,H_{C}),
\]
isomorphe aussi, \textit{si $S$ est finie,} à
\[
\Hom_{\Cathot}(S,C_{\infty}),
\]
ce qui montre que, du moins pour des arguments $S \in (\Cat)$ qui sont des catégories \textit{finies}, les objets
\msp{page 107}
$C_{\infty}$ et $H_{C}$ sont ‘bons', il nous permettent % [sic] il nous permettent -> ils nous permettent, agreement with the plural subject
de décrire un foncteur important qui domine la situation envisagée. Moyennant une hypothèse très forte, mais satisfaite dans les cas importants, on peut dans ces deux formules remplacer Cat\,hot par Hot.

Par contre, sauf pour le cas où $S$ est réduit à un point et où $\Hom_{\Cathot}$ ou $\Hom_{\Hot}$ de $S$ dans $C$ n'est autre que $\pi_{0}(C)$ ($\simeq \pi_{0}(C_{\infty}) \simeq \pi_{0}(H_{C})$), donc effectivement ‘bon', je ne sais pas si
\[
\underbrace{\Hom_{\Cathot}(S,C)}_{=\,\pi_{0}\uHom(S,C)} \qquad \text{ou} \qquad \Hom_{\Hot}(S,C)
\]
est lié de façon raisonnable aux ensembles précédents (déterminés par $H_{C}$, resp. par $C_{\infty}$), p. ex. s'ils lui sont isomorphes de façon naturelle, et s'ils ont une signification géométrico-topologique raisonnable. Je commence à suspecter que \textit{non} - donc qu'on aurait ici un cas où une relation de type très général concernant un $\pi_{0}$, ici donc
\[
\pi_{0}(C) \simeq \pi_{0}(C_{\infty}) \simeq \pi_{0}(H_{C}),
\]
ne se prolonge pas en un isomorphisme sur les $\pi_{i}$ supérieurs, et en une propriété de Hot-équivalence.


\endgroup
% END UNIT Der-ChVII-456/Der-ChVII-456_p066-070.tex
% BEGIN UNIT Der-ChVII-456/Der-ChVII-456_p071-075.tex
\begingroup\makeatletter%
\newcommand{\printedpage}[1]{\marginpar{\footnotesize p.~#1}}%
\newcommand{\msp}[1]{\marginpar{\footnotesize [#1]}}%
\newcommand{\Hotz}{\mathrm{Hot}^{\wedge}_{0}}%
\newcommand{\Ch}{\underline{\mathrm{Ch}}}%
\newcommand{\Ty}{\underline{\mathrm{Ty}}}%
\newcommand{\CNo}{C_{\mathbf{N}^{o}}}%
\newcommand{\lab}[2][l]{\text{\footnotesize\begin{tabular}{@{}#1@{}}#2\end{tabular}}}%
\makeatother\let\printedpage\derpp\let\msp\dermsp%
\let\mspin\DERmspin%


\printedpage{69}
% head: A. Grothendieck : Dérivateurs, Ch. VII, Édition des Universités de Montpellier II et Paris VII

\msp{page 108}
Je ne trouve d'ailleurs aucune flèche directe entre $C$ et l'un des objets $C_{\infty}$, $H_{C}$; tout ce que je vois dans ce sens, c'est
\[
\begin{tikzcd}[column sep=large, row sep=huge]
\CNo \arrow[r, "{\substack{\text{foncteur de}\\ \text{localisation}}}"] \arrow[d] & C_{\infty} \\
C, &
\end{tikzcd}
\]
où la flèche verticale est la flèche de projection de $\CNo = C \times_{T} \mathbf{N}^{o}$ (avec $\mathbf{N}^{o} \xrightarrow{\ t\ } T$). Ceci n'établit une relation topologique étroite entre $C$ (en tant qu'objet de Cat\,hot ou de Hot simplement) et $C_{\infty}$, que si les deux flèches précédentes sont tout au moins des Hot-équivalences. Et le cas-type (problématique) où la question de la relation entre $C$ et $C_{\infty}$ se pose, est donc celui où $T = \mathbf{N}^{o}$, i.e. celui d'un système inductif quelconque de (petites) catégories, indexé par $\mathbf{N}$, et interprété comme défini par une catégorie fibrée $C = \CNo$ sur $\mathbf{N}^{o}$. Dans ce cas particulier, on a du moins une flèche canonique
\[
C \longrightarrow C_{\infty}
\]
sur laquelle on peut se poser la question si c'est
\msp{page 109}
une Hot-équivalence. Sinon, cela prouvera que dans le cas général, le Hot-type de $C$ (et a fortiori sa classe d'isomorphie dans Cat\,hot) n'a pas de sens géométrique intéressant. C'est à dire que $C$ est intéressante uniquement \textit{en tant que catégorie fibrée sur $T$}, permettant notamment de construire l'ind-objet $H_{T}$ de Cat\,hot ou de Hot, mais pas en tant que `catégorie nue'.

Mais je viens de me convaincre que $\CNo \longrightarrow C_{\infty}$ est bien une Hot-équivalence. Donc la flèche $\CNo \longrightarrow C$ définit une flèche \textit{dans} Hot
\[
C_{\infty} \longrightarrow C
\]
(en composant avec l'inverse dans Hot de la précédente). Est-ce un isomorphisme? Donc pour cette question, on est ramené à la question pour $\CNo \longrightarrow C$. Mais si on prend $C = T$, cela pose déjà la même question pour $\mathbf{N}^{o} \longrightarrow T$, donc la question si $T$ est asphérique, sous la seule hypothèse que $T_{\mathrm{o}}$ soit filtrante à gauche et dénombrable à l'infini (à gauche). Mais bien sûr c'est faux (p. ex. si $T$ n'a qu'un seul objet, il n'est pas pour autant asphérique !).

Cela me convainc donc qu'à défaut
% foot: Version 22 décembre 2001
\printedpage{70}
% head: A. Grothendieck : Dérivateurs, Ch. VII, Édition des Universités de Montpellier II et Paris VII
\msp{page 110}
d'hypothèse supplémentaire draconienne sur $T$, la catégorie $C$ `nue' n'offre pas d'intérêt particulier, à part son $\pi_{0}$ (qui est donné aussi par $C_{\infty}$, et par $H_{C}$). Il n'est pas dit, certes, que \underline{Ty} ne satisfasse pas à ces `conditions draconiennes'. Mais comme $C$ `nue' ne s'insère pas dans un formalisme sympathique (p.ex. on n'a pas d'interprétation géométrique de $\mathrm{Hom}_{\mathrm{Cat\,hot}}(S,C)$, ou de $\mathrm{Hom}_{\mathrm{Hot}}(S,C)$), je crois que par la suite je peux en faire abstraction. Exit a) !

\begin{flushright}
{\footnotesize 5.11.90}
\end{flushright}

J'ai quand même passé la journée entière hier à essayer de comprendre la relation de $C$ et de son type dans Hot, à celui de $C_{\infty}$. En effet, comme on a finalement une flèche dans Hot
\[
C_{\infty} \longrightarrow C,
\]
déduite du diagramme dans (Hot)
\[
\begin{tikzcd}[column sep=huge, row sep=huge]
\CNo \arrow[r, "\sim"] \arrow[d] & C_{\infty} \\
C, &
\end{tikzcd}
\]
de prouver que $C_{\infty} \longrightarrow C$ est un isomorphisme dans Hot
\msp{page 111}
semble la façon la plus simple, et peut-être unique, de prouver que l'image de $C_{\infty}$ dans Hot ne dépend pas, à isomorphisme unique près, des choix faits pour construire $C_{\infty}$. Sans compter que dans le cas motivant, $C = \Ch_{\Sigma}(X)$ ou $\Ch_{\Sigma}(X;a,b)$, $T = \Ty$, c'est la catégorie des chemins $C$ qui donne lieu à une loi de composition des chemins impeccable, fonctorielle, associative et tout, alors qu'il n'en est rien pour $C_{\infty}$. On ne trouve pas une loi de composition
\[
C_{\infty} \times C_{\infty} \longrightarrow C_{\infty}
\]
dans Cat, ni même dans Cat\,hot ou dans Hot - ou plutôt, si on arrive pourtant à construire une flèche $C_{\infty} \times C_{\infty} \longrightarrow C_{\infty}$ dans Cat, c'est encore laborieusement à coups de choix $\pm$ arbitraires, cran par cran, et sans savoir si ce qu'on obtient a une propriété d'unicité, ne serait ce que dans Hot; on ne trouve une interprétation géométrique tangible de cette composition que dans $\Hotz$ (où on a une véritable loi de groupoïde), ou dans $H_{T}$ en tant qu'objet de Ind\,Cat\,hot
% foot: Version 22 décembre 2001
\printedpage{71}
% head: A. Grothendieck : Dérivateurs, Ch. VII, Édition des Universités de Montpellier II et Paris VII
\msp{page 112}
ou de Ind\,Hot. (Donc la limite inductive dans $\Hotz$ n'est autre que l'élément de $\Hotz$ défini par $C_{\infty}$.)

Mes efforts pour prouver que
\[
C_{\infty} \longrightarrow C
\]
était un isomorphisme dans Hot, ou ne serait-ce qu'un 1-isomorphisme (i.e. un isomorphisme non seulement pour les $\pi_{0}$, ce qui est toujours vrai, mais aussi pour les $\pi_{1}$ de toutes les composantes), ont échoués. J'ai trouvé qu'il suffisait que le foncteur % [sic] ont échoués -> ont échoué, accord fautif du participe
\[
T \longrightarrow T_{\mathrm{o}}
\]
soit ou bien lisse et coasphérique, ou bien propre et asphérique - donc tout est OK si $T = T_{\mathrm{o}}$, i.e. si $T$ est ordonnée. Mais dans le cas $T = \Ty$, le foncteur $T \longrightarrow T_{\mathrm{o}}$ n'est ni propre, ni lisse (pas même 0-propre ou 0-lisse).

Finalement j'en suis arrivé à la conviction que dans le cas qui nous occupe ($T = \Ty$, $C = \Ch_{\Sigma}(X;a,b)$), la flèche $C_{\infty} \longrightarrow C$ n'est en général \textit{pas} un isomorphisme dans Hot. En effet, dans le cas $a = b$, si cela \texttt{[était vrai]},
\msp{page 113}
on aurait une description du Hot-groupe $\Omega(X,a) = \Ch(X,a)$ (\footnote[23]{On fait ici $\Sigma = \mathrm{Fl}(X)$.}) en termes d'une catégorie $C$ munie d'une loi de composition \textit{associative} (et unitaire) \textit{dans} Cat, redonnant la loi de groupe de $\Omega$ en passant de (Cat) à Hot. Prenant le nerf de $C$, on trouverait un monoïde semi-simplicial qui, dans Hot, donnerait la loi de groupe de $\Omega(X,a)$. Or je me rappelle (\footnote[24]{C'est un erreur - ça doit être pour d'\textit{autres} objets groupe de Hot que les groupes de lacets. J'ai vu dans Quillen que $\Omega(X,x)$ doit se décrire par un groupe semi-simplicial (quand $X$ est un ensemble semi-simplicial, $x$ un sommet). Donc le problème si $C_{\infty} \simeq C$ dans Hot reste entier.}) qu'on m'a expliqué qu'en général, la loi de groupe de $\Omega(X,a)$ ne pourrait pas être obtenue par une loi \textit{associative} sur un ensemble semi-simplicial. Ainsi, l'insuccès de mes efforts semble tenir, non à ma maladresse, mais à la nature des choses. % [sic] C'est un erreur (note 24) -> C'est une erreur, genre

Je vais résumer les faits positifs auxquels je suis arrivé dans le diagramme suivant, qui a lieu dans la situation générale $C \longrightarrow T$ envisagée au début du présent paragraphe :

\msp{page 114}
% foot: Version 22 décembre 2001
\printedpage{72}
% head: A. Grothendieck : Dérivateurs, Ch. VII, Édition des Universités de Montpellier II et Paris VII
\[
\begin{tikzcd}[column sep=small, row sep=huge]
 & & C_{t'} \arrow[d, "{u^{*}\quad\lab{(pour $u : t \longrightarrow t'$\\ dans $T$)}}"] \arrow[ldd, dotted] \arrow[rdd, dotted] & \\
 & & C_{t} \arrow[ld, swap, "{\lab{flèche définie dans\\ Ind\,Cat\,hot\\ (\texttt{[ou dans]}\\ IndHot,\\ Ind\,($\Hotz$))}}"] \arrow[rd, "{\lab[r]{flèche dans Cat\\ (donc dans\\ Cat\,hot,\\ Hot,\\ $\Hotz$)}}"] & \\
\varinjlim C_{t} = & |[draw, circle]| H_{T} \arrow[rd, swap, "{\lab{flèche\\ définie dans\\ Ind\,Cat\,hot\\ ou dans Ind\,Hot\\ ou dans Ind\,$\Hotz$}}"] & \lab{diagramme\\ commutatif\\ dans Ind\,Hot\\ ou Ind\,($\Hotz$)} & C \\
\lab{dans Ind\,Cat\,hot\\ (ou dans Ind\,Hot\\ ou dans Ind\,($\Hotz$))} & & |[draw, circle]| C_{\infty} \arrow[ru, swap, "{\lab[r]{flèche\\ définie\\ dans Hot,\\ donc dans $\Hotz$,\\ en général pas un\\ isomorphisme, même\\ dans $\Hotz$}}"] & \\
 & & = \varinjlim H_{T} \text{ dans } \Hotz &
\end{tikzcd}
\]

\texttt{[Les annotations dans ce diagramme concernent les flèches en traits pleins.]}

Les objets importants sont $H_{T}$ et $C_{\infty}$, alors que $C$ (sur $T$) et les $C_{t}$ prennent un rôle technique intermédiaire, pour construire $H_{T}$ et $C_{\infty}$. On peut considérer $C_{\infty}$ comme une catégorie (objet de (Cat)), mais dont la construction dépend de nombreux arbitraires, à l'inverse de $H_{T}$ dans Ind\,Cat\,hot. C'est l'objet de $\Hotz$ déduit de $C_{\infty}$ qui a un sens canonique, et peut être décrit comme la limite inductive de $H_{T}$ comme objet de Ind\,$\Hotz$, i.e. comme représentant le même foncteur sur $\mathrm{Hot}_{0}$ que $H_{T}$.

\msp{page 115}
Dans le cas géométrique $T = \Ty$, et $C$ de la forme $\Ch_{\Sigma}(X;a,b)$, donnant lieu à $C_{\infty}(a,b)$, on \textit{n'a pas} de définition naturelle d'un bifoncteur de composition des chemins
\[
C_{\infty}(a,b) \times C_{\infty}(b,c) \ \longrightarrow\ C_{\infty}(a,c).
\]
On peut définir, certes, des accouplements
\[
\Ch_{\Sigma,\tau}(a,b) \times \Ch_{\Sigma,\tau'}(b,c) \ \longrightarrow\ \Ch_{\Sigma,\tau\tau'}(a,c),
\]
ou, avec les notations relatives à la donnée de
\[
\begin{aligned}
&\underline{t} : \mathbf{N}^{o} \longrightarrow \Ty,\\
&\Ch_{n}(a,b) \times \Ch_{m}(b,c) \longrightarrow \Ch_{n+m}(a,c),
\end{aligned}
\]
% foot: Version 22 décembre 2001
\printedpage{73}
% head: A. Grothendieck : Dérivateurs, Ch. VII, Édition des Universités de Montpellier II et Paris VII
mais ces accouplements \textit{ne sont pas} compatibles avec les morphismes de transition des trois systèmes inductifs de catégories $(\Ch_{n})_{n\in\mathbf{N}}$, ils ne le sont \textit{qu'à homotopie près.} Cela suffit cependant pour avoir des accouplements parfaits
\[
H_{T}(a,b) \times H_{T}(b,c) \ \longrightarrow\ H_{T}(a,c),
\]
avec associativité, et (si $\Sigma = \mathrm{Fl}(X)$ tout entier) existence de \textit{vrais} inverses, i.e. de flèches (dans Ind\,Cat\,hot) de `passage au chemin opposé'

\msp{page 116}
\[
H_{T}(a,b) \ \longrightarrow\ H_{T}(b,a), \qquad c \longmapsto c^{o}, \tag*{$(*)$}
\]
donnant lieu au diagramme commutatif exprimant que c'est bien un inverse pour la composition. Une autre façon de le dire, c'est que pour tout objet $Y$ de Cat\,hot (ou de Ind\,Cat\,hot), les ensembles
\[
\mathrm{Hom}_{\mathrm{Cat\,hot}}(Y, H_{T}(a,b)) \overset{\text{déf}}{=} H_{T}^{Y}(a,b)
\]
forment, pour $a,b \in X$ variables, les flèches d'un \textit{groupoïde} (\footnote[25]{Si $\Sigma = \mathrm{Fl}(X)$. Si $\Sigma$ est plus petit, au lieu d'un groupoïde, on trouve une \textit{catégorie} sans plus.}) $\mathrm{Ch}_{\Sigma}^{Y}(X)$ ayant pour ensemble de sommets $\mathrm{Ob}\,X$, et $H_{T}^{Y}(a,b)$ étant formé des flèches de $a$ vers $b$. L'application
\[
H_{T}^{Y}(a,b) \longrightarrow H_{T}^{Y}(b,a), \qquad c \longmapsto c^{o},
\]
déduite du foncteur $(*)$ plus haut n'est autre que l'application de passage à l'inverse. Cette catégorie n'est d'ailleurs autre que celle déduit de % [sic] celle déduit -> celle déduite, accord
\[
X(Y)\Sigma_{Y}^{-1}, \qquad \text{où } \left\{
\begin{aligned}
&X(Y) = \underline{\mathrm{Hom}}_{\mathrm{Cat}}(Y,X)\\
&\Sigma_{Y} \subseteq \mathrm{Fl}(X(Y)),\\
&\text{défini en termes de } \Sigma \text{ `argument par argument'},
\end{aligned}
\right.
\]
à l'aide du foncteur
\[
X \longrightarrow X(Y), \qquad a \longmapsto a_{Y},
\]
\msp{page 117}
associant à tout $a \in X$ le foncteur \textit{constant} sur $Y$ de valeur $a$, en prenant
\[
H_{T}(a,b) \simeq \mathrm{Hom}_{X(Y)\Sigma_{Y}^{-1}}(a_{Y}, b_{Y}) \ .
\]

Dans le même cas géométrique, je me suis convaincu d'ailleurs que $C_{\infty}$ était, en tant qu'objet de Hot (pas seulement de $\Hotz$) \textit{unique à} Hot-\textit{isomorphisme près,} mais (chose étrange) peut-être pas à Hot-isomorphisme \textit{unique} près. De façon générale, revenant aux
% foot: Version 22 décembre 2001


\endgroup
% END UNIT Der-ChVII-456/Der-ChVII-456_p071-075.tex
% BEGIN UNIT Der-ChVII-456/Der-ChVII-456_p076-080.tex
\begingroup\makeatletter%
\NoAutoSpacing%
\newcommand{\printedpage}[1]{\marginpar{\footnotesize p.~#1}}%
\newcommand{\msp}[1]{\marginpar{\footnotesize [#1]}}%
\newcommand{\Hot}{\mathrm{Hot}}%
\newcommand{\Hotz}{\mathrm{Hot}_0^{\wedge}}%
\newcommand{\Hc}{\underline{\mathrm{Hom\,cof}}}%
\newcommand{\Nop}{\mathbf{N}^o}%
\makeatother\let\printedpage\derpp\let\msp\dermsp%
\let\mspin\DERmspin%


\printedpage{74}
% head: A. Grothendieck : Dérivateurs, Ch. VII, Édition des Universités de Montpellier II et Paris VII
conditions du début du paragraphe, $C_\infty$ doit s'écrire $C_\infty(\underline{t})$, elle dépend du choix du foncteur cofinal
\[
\underline{t} : \Nop \longrightarrow T.
\]
Quand $\underline{t}$ varie dans la catégorie
\[
\Hc(\Nop, T),
\]
les $C_\infty(\underline{t})$ varient de façon contravariante en $\underline{t}$ (donc définissent une catégorie fibrée sur $\Hc(\Nop, T)$). De plus, le foncteur
\[
C_\infty(\underline{t'}) \longrightarrow C_\infty(\underline{t})
\]
associé à un morphisme de systèmes projectifs
\[
\underline{t} \longrightarrow \underline{t'}
\]
est un \textit{isomorphisme} dans Hot. Donc le
\msp{page 118}
foncteur
\[
\underline{t} \longmapsto C_\infty(\underline{t}), \qquad \Hc(\Nop, T) \longrightarrow (\Hot)
\]
se factorise par le groupoïde fondamental
\[
\Pi(\Hc(\Nop, T)) \longrightarrow (\Hot).
\]
Quand ce groupoïde, i.e. quand la catégorie $\Hc$ elle-même est connexe, on peut donc conclure que tous les $C_\infty(\underline{t})$ sont isomorphes entre eux dans Hot. Pour conclure même à un système \textit{transitif} d'isomorphismes entre les $C_\infty(\underline{t})$, i.e. que $C_\infty(\underline{t})$ est ‘unique à isomorphisme \textit{unique} près’, il faudrait savoir de plus que le groupoïde fondamental est équivalent à la catégorie ponctuelle (groupe fondamental trivial), ou, comme on dit, que $\Hc$ soit \textit{simplement connexe,} ou 1-connexe.

La connexité de $\Hc$ est assurée quand la catégorie $T$ a (comme c'est le cas pour $\underline{\mathrm{Ty}}$) une loi de composition $(t,t') \longmapsto t * t'$, avec des morphismes fonctoriels en $t, t'$
\[
\begin{tikzcd}[row sep=small]
 & t \\
t * t' \arrow[ru] \arrow[rd] & \\
 & t'
\end{tikzcd}
\]
(cette dernière condition étant automatique
% foot: Version 22 décembre 2001
\printedpage{75}
% head: A. Grothendieck : Dérivateurs, Ch. VII, Édition des Universités de Montpellier II et Paris VII
\msp{page 119}
si la loi de composition admet l'objet final de $T$, supposé existant, comme unité bilatère). Dans ce cas, $\underline{t} = (t_n)_n$ et $\underline{t'} = (t'_n)_n$ sont reliés entre eux par $\underline{t} * \underline{t'} = (t_n * t'_n)_n$, puisqu'on a un morphisme de systèmes projectifs
\[
\begin{tikzcd}[column sep=small]
 & \underline{t} * \underline{t'} \arrow[ld] \arrow[rd] & \\
\underline{t} & & \underline{t'}.
\end{tikzcd}
\]
Je ne sais d'ailleurs pas si la flèche dans Hot
\[
C_\infty(\underline{t}) \longrightarrow C_\infty(\underline{t'})
\]
associée à $\underline{t'} \longrightarrow \underline{t}$ dans $\Hc$ ne dépend que de $\underline{t}$, $\underline{t'}$, et pas du choix de $\underline{t'} \longrightarrow \underline{t}$. Ce qui est clair, c'est que son image dans $\Hotz$ est unique, c'est l'identité quand on identifie $C_\infty(\underline{t})$ et $C_\infty(\underline{t'})$ au seul et même élément de $\Hotz$ défini par $\varinjlim H_T$.

Et pourtant, je m'aperçois à l'instant que, sous l'hypothèse de l'existence de la loi de composition et de son unité bilatère, et admettant de plus que la loi en question est associative à isomorphisme près, \textit{on a transitivité} dans les isomorphismes (canoniques!) entre les $C_\infty(\underline{t})$ qu'on vient d'indiquer !

\msp{page 120}
(Sans avoir pour autant à calculer le groupe fondamental de $\Hc(\Nop, T)$!) Si on a trois éléments $\underline{t}$, $\underline{t'}$, $\underline{t''}$ de $\Hc$, il suffit de regarder le diagramme déduit du diagramme commutatif suivant
\[
\begin{tikzcd}[column sep=small]
 & & \underline{t} * \underline{t'} * \underline{t''} \arrow[ld] \arrow[rd] & & \\
 & \underline{t} * \underline{t'} \arrow[ld] \arrow[rd] & & \underline{t'} * \underline{t''} \arrow[ld] \arrow[rd] & \\
\underline{t} & & \underline{t'} & & \underline{t''}
\end{tikzcd}
\]
en lui appliquant le contrafoncteur $\underline{s} \longrightarrow C_\infty(\underline{s})$ de $\Hc$ dans Hot.

Donc dans le cas des $C(a,b) = \underline{\mathrm{Ch}}_\Sigma(X; a, b)$, je n'ai plus aucun scrupule à considérer la catégorie $C_\infty(a,b)$ comme un objet \textit{canonique} de Hot, relié à l'ind-objet canonique $H_T(a,b)$ dans Hot par un homomorphisme canonique
\[
C_\infty(a,b) \longleftarrow H_T(a,b), \tag*{$(*)$}
\]
qui dans $\Hotz$ donne un isomorphisme
\[
C_\infty(a,b) \simeq \varinjlim H_T(a,b), \tag*{$(**)$}
\]
i.e. la flèche $(*)$ induit, pour $\xi \in \mathrm{Ob}\,\Hot$, une application ensembliste
\[
\mathrm{Hom}_{\Hot}(\xi, C_\infty(a,b)) \longleftarrow \underbrace{\mathrm{Hom}_{\mathrm{Ind}\,\Hot}(\xi, H_T(a,b))}_{\underset{n}{\varinjlim}\, \mathrm{Hom}_{\Hot}(\xi, C_n(a,b))},
\]
qui est \textit{bijective} si $\xi$ est dans Hot.
% foot: Version 22 décembre 2001

\printedpage{76}
% head: A. Grothendieck : Dérivateurs, Ch. VII, Édition des Universités de Montpellier II et Paris VII
\msp{page 121}
Bien entendu, les flèches
\[
C_\infty(\underline{t'}) \longrightarrow C_\infty(\underline{t}) \tag*{$(*)$}
\]
dans Hot, associés aux flèches $\underline{t} \longrightarrow \underline{t'}$ dans $\Hc(\Nop, T)$, sont en fait des flèches dans % [sic] associés -> associées, accord avec « les flèches »
(Cat) elle-même - et même des foncteurs pleinement fidèles, dans le cas ‘géométrique’ et pour les deux flèches canoniques
\[
\begin{tikzcd}[column sep=small]
 & \underline{t} * \underline{t'} \arrow[ld] \arrow[rd] & \\
\underline{t} & & \underline{t'}.
\end{tikzcd}
\]
Elles sont isomorphiques \textit{dans} Hot, parce qu'elles définissent un homomorphisme sur les ind-objets associés de Cat\,hot, donc aussi pour les ind-objets associés dans Hot et dans $\Hotz$, d'où résulte que
\[
\mathrm{Hom}_{\Hot}(Y, C_\infty(\underline{t'})) \longrightarrow \mathrm{Hom}_{\Hot}(Y, C_\infty(\underline{t}))
\]
associé à $(*)$ est une bijection si $Y$ \textit{est finie} (\footnote[26]{Ce qui suffit à assurer que $(*)$ est \texttt{[un]} isomorphisme dans Hot : $\Hot \longrightarrow \Hotz$ est conservatif!}), cette flèche s'identifie alors à
\[
\underset{n}{\varinjlim}\, \mathrm{Hom}_{\Hot}(Y, C_{t'_n}) \longrightarrow \underset{n}{\varinjlim}\, \mathrm{Hom}_{\Hot}(Y, C_{t_n}).
\]
(Car $Z \longmapsto \mathrm{Hom}_{\Hot}(Y, Z)$ commute aux limites inductives filtrantes, si $Y$ est une catégorie \textit{finie.} Cela provient du fait qu'une $\varinjlim$ filtrante de catégories de Thomason-Kan est itou \dots) Mais j'ignore par contre si $(*)$ est \texttt{[un]} isomorphisme dans Cat\,hot. Seulement que pour une catégorie \textit{finie} $Y$, elle induit
\[
\mathrm{Hom}_{\mathrm{Cat\,hot}}(Y, C_\infty(\underline{t'})) \xrightarrow{\ \sim\ } \mathrm{Hom}_{\mathrm{Cat\,hot}}(Y, C_\infty(\underline{t})),
\]
les deux membres s'identifiant alors à
\msp{page 122}
\[
H_T(Y)_{\mathrm{Cat\,hot}} = \mathrm{Hom}_{\mathrm{Ind\,Cat\,hot}}(Y, H_T).
\]
Pour résumer ces lumières et perplexités sur les $C_\infty(\underline{t})$ :

\begin{enumerate}
\item[a)] À tout $\underline{t} = (t_n)_{n \in \Nop} \in \Hc(\Nop, T)$ est associé un système inductif, indexé par $\mathbf{N}$, de catégories $C_n(\underline{t})$ $(= C_{t_n})$, et leur limite inductive
\[
C_\infty(\underline{t}).
\]
% foot: Version 22 décembre 2001

\printedpage{77}
% head: A. Grothendieck : Dérivateurs, Ch. VII, Édition des Universités de Montpellier II et Paris VII
\item[b)] Ce système inductif, et aussi sa limite est \texttt{[un]} foncteur contravariant en $\underline{t}$. Si $\underline{t}, \underline{t'} \in \mathrm{Ob}\,\Hc$, alors on peut comparer $C_\infty(\underline{t})$, $C_\infty(\underline{t'})$ et les systèmes inductifs correspondants, grâce à
\[
\begin{tikzcd}[column sep=small]
 & \underline{t} * \underline{t'} \arrow[ld] \arrow[rd] & \\
\underline{t} & & \underline{t'}.
\end{tikzcd}
\]

\item[c)] L'objet
\[
(C_n(\underline{t}))_{n \in \mathbf{N}} \in \mathrm{Ob\,Ind\,Cat\,hot}
\]
défini par le système inductif $(C_n(\underline{t}))_{n \in \mathbf{N}}$ de Cat est (à isomorphisme unique près) indépendant du choix de $\underline{t}$, soit $H_C$ ou $H$. Les homomorphismes
\[
C_\infty(\underline{t'}) \xrightarrow{u^*} C_\infty(\underline{t}) \qquad \text{associé à } u : \underline{t} \longrightarrow \underline{t'} \tag*{$(*)$}
\]
sont compatibles avec cette donnée, i.e. on a commutativité dans
\[
\begin{tikzcd}
\underset{\mathrm{Ind\,Cat\,hot}}{\varinjlim} C_*(\underline{t'}) \arrow[rr] & & \underset{\mathrm{Ind\,Cat\,hot}}{\varinjlim} C_*(\underline{t}) \\
 & H \arrow[lu, "\sim"] \arrow[ru, "\sim"'] &
\end{tikzcd}
\]

\msp{page 123}
\item[d)] Il s'ensuit que les homomorphismes $(*)$ sont des \textit{isomorphismes} dans Hot, donc tous les $C_\infty(\underline{t})$ sont isomorphes \textit{dans} Hot. L'astuce des $t * t'$ (dans le cas où on a une loi $t * t'$ associative et admettant l'objet final comme objet unité) donne même des isomorphismes \textit{canoniques}
\[
C_\infty(\underline{t'}) \longrightarrow C_\infty(\underline{t}) \qquad \text{dans Hot,}
\]
et c'est là un système transitif d'isomorphismes. Donc on a dans Hot un objet canonique
\[
C_\infty \in \mathrm{Ob}\,\Hot,
\]
dont la relation à $H \in \mathrm{Ob\,Ind\,Hot}$ est déduite de c), c'est une flèche canonique
\[
H \longrightarrow C_\infty \qquad \text{dans Ind\,Hot (\footnotemark[27]).}
\]
\footnotetext[27]{\textbf{NB} $H$ n'est pratiquement jamais essentiellement constant, donc la flèche ci-contre \texttt{[n'est]} pas un isomorphisme dans Ind\,Hot.}

Cette flèche induit un isomorphisme pour les objets correspondants de $\Hotz$; mais en regardant $H$ non comme ind-objet de $\Hotz$, mais prenant le foncteur $H_0$ induit sur $\Hot_0$ par le foncteur ind-représentable $H$ sur Hot. Ce foncteur induit est ‘représentable’ par l'objet $C_\infty$ de la catégorie plus grande
\[
H_0 \simeq C_\infty \qquad \text{dans } \Hotz.
\]

(Si le foncteur canonique $\Hot \longrightarrow \Hotz$ est pleinement fidèle, alors cette formule décrit $C_\infty$ comme objet de Hot, à isomorphisme unique près, en termes de $H_0$, donc aussi en termes de $H \in \mathrm{Ob\,Ind\,Cat\,hot}$.)
% foot: Version 22 décembre 2001

\printedpage{78}
% head: A. Grothendieck : Dérivateurs, Ch. VII, Édition des Universités de Montpellier II et Paris VII
\msp{page 124}
\item[e)] J'ignore par contre si les flèches
\[
C_\infty(\underline{t'}) \longrightarrow C_\infty(\underline{t})
\]
sont des isomorphismes dans Cat\,hot, bien qu'elles soient définies en termes d'homomorphismes de systèmes inductifs de catégories, qui sont des isomorphismes dans Ind\,Cat\,hot. Il en est ainsi, même pour les flèches $C_\infty(\underline{t}), C_\infty(\underline{t'}) \rightrightarrows C_\infty(\underline{t} * \underline{t'})$ utilisées pour comparer $C_\infty(\underline{t})$ et $C_\infty(\underline{t'})$ canoniquement. Ainsi, je ne vois pas comment définir des flèches en sens inverse dans Cat, qui seraient candidats à être des isomorphismes inverses dans Cat\,hot.

\item[f)] Dans le cas géométrique ($T = \underline{\mathrm{Ty}}$, situation décrite en termes de $X$, $\Sigma$, et $a, b \in \mathrm{Ob}\,X$, d'où des $C_\infty(\underline{t}; a, b)$, $C_\infty(\overline{a, b})$, $H(a,b)$ etc.), on \textit{n'a pas} d'opération de composition de chemins au niveau des \textit{catégories} $C_\infty(\underline{t}; a, b)$ et $C_\infty(\underline{t}; b, c)$, mais seulement au niveau des objets $H(a,b)$, $H(b,c)$ de Ind\,Cat\,hot ou de Ind\,Hot. Suis-je seulement sûr de savoir définir \textit{dans} Hot
\[
C_\infty(a,b) \times C_\infty(b,c) \longrightarrow C_\infty(a,c) \ \ ?
\]

Sûrement, il n'y a pas de problèmes pour les images dans $\Hotz$, puisque là ce sont les mêmes que celles de $H(a,b)$, $H(b,c)$, $H(a,c)$ respectivement, et on a des accouplements parfaits, associatifs, avec unité, et avec inverses là ou il doit y % [sic] là ou -> là où, accent omis sur le relatif
\msp{page 125}
en avoir. Mais pas pour les $C_\infty(a,b)$ eux-mêmes?
\end{enumerate}

C'est un peu vexant qu'avec tout ce travail il semblerait qu'on n'ait pas même été foutu de définir l'analogue dans Cat du classique espace des lacets $\Omega(X, x)$ comme un groupe \textit{dans} Hot. (Que fait \textsc{Thomason}?) Mais voilà une idée pensable [?]. Pour définir \textit{quand-même}
\[
C_\infty(a,b) \times C_\infty(b,c) \longrightarrow C_\infty(a,c)
\]
dans Hot, on va décrire les deux foncteurs $C_\infty(a,b)$ et $C_\infty(b,c)$ par des systèmes projectifs différents $\underline{t}$, $\underline{t'}$. On peut l'expliciter en partant d'un $\tau_0$ de longueur utile $\geq 2$, et en prenant dans les deux cas
\[
t_n = t'_n = \tau_0^n,
\]
mais des morphismes de transition différents
\[
\begin{aligned}
(t_{n+p} \longrightarrow t_n) &= \pi_p * \mathrm{id}_{\tau_0^n} \\
(t'_{n+p} \longrightarrow t'_n) &= \mathrm{id}_{\tau_0^n} * \pi_p,
\end{aligned}
\]
où $\pi_p : \tau_0^p \longrightarrow 1$ est la dégénérescence totale. On définira $C_\infty(a,c)$ en termes de $\underline{t''} = \underline{t} * \underline{t'}$, donc
\[
\left\{
\begin{aligned}
t''_n &= \tau_0^{2n} \\
(t''_{n+p} \longrightarrow t''_n) &= \pi_p * \mathrm{id}_{\tau_0^{2n}} * \pi_p.
\end{aligned}
\right.
\]
% foot: Version 22 décembre 2001


\endgroup
% END UNIT Der-ChVII-456/Der-ChVII-456_p076-080.tex
% BEGIN UNIT Der-ChVII-456/Der-ChVII-456_p081-085.tex
\begingroup\makeatletter%
\newcommand{\printedpage}[1]{\marginpar{\footnotesize p.~#1}}%
\newcommand{\msp}[1]{\marginpar{\footnotesize [#1]}}%
\newcommand{\Ch}{\underline{\mathrm{Ch}}}%
\newcommand{\uHom}{\underline{\mathrm{Hom}}}%
\newcommand{\comp}{\operatorname{comp}}%
\newcommand{\Hot}{\mathrm{Hot}}%
\makeatother\let\printedpage\derpp\let\msp\dermsp%
\let\mspin\DERmspin%


\printedpage{79}
% head: A. Grothendieck : Dérivateurs, Ch. VII, Édition des Universités de Montpellier II et Paris VII
% foot: Version 22 décembre 2001

Cette fois, les foncteurs de composition
\[ C_n(\underline{t};a,b)\times C_n(\underline{t'};b,c) \longrightarrow C_n(\underline{t''};a,c), \]
sont compatibles avec les foncteurs transition pour $n$ variable, et on trouve
\msp{page 126}
\[ C_\infty(\underline{t};a,b)\times C_\infty(\underline{t'};b,c) \longrightarrow C_\infty(\underline{t''};a,c), \]
donc une flèche dans Hot
\[ \comp_{\tau_0} : C_\infty\times C_\infty \longrightarrow C_\infty. \]
Il faut voir qu'elle ne dépend pas du choix de $\tau_0$, mais si on a $\tau_0$ et $\tau'_0$, on les compare par $\tau_0\times\tau'_0 \begin{array}{l}\longrightarrow \tau_0\\ \longrightarrow \tau'_0\end{array}$, il suffit donc de voir que s'il existe une flèche $\tau_0\longrightarrow\tau'_0$, alors $\comp_{\tau_0} = \comp_{\tau'_0}$. Mais une telle flèche induit des flèches $\underline{t}_{\tau_0}\longrightarrow\underline{t}_{\tau'_0}$, $\underline{t'}_{\tau_0}\longrightarrow\underline{t'}_{\tau'_0}$, $\underline{t''}_{\tau_0}\longrightarrow\underline{t''}_{\tau'_0}$, d'où un diagramme dans Cat
\[
\begin{tikzcd}
C_\infty(\underline{t}_{\tau'_0};a,b)\times C_\infty(\underline{t'}_{\tau'_0};b,c) \arrow[r, "{\comp_{\tau'_0}}"] \arrow[d] & C_\infty(\underline{t''}_{\tau'_0};a,c) \arrow[d] \\
C_\infty(\underline{t}_{\tau_0};a,b)\times C_\infty(\underline{t'}_{\tau_0};b,c) \arrow[r, "{\comp_{\tau_0}}"] & C_\infty(\underline{t''}_{\tau_0};a,c),
\end{tikzcd}
\]
dont on vérifie tout de suite qu'il est commutatif. Or ce diagramme peut se réécrire dans Hot comme
\[
\begin{tikzcd}
C_\infty(a,b)\times C_\infty(b,c) \arrow[r, "{\comp_{\tau'_0}}"] \arrow[d, "{\mathrm{id}}"] & C_\infty(a,c) \arrow[d, "{\mathrm{id}}"] \\
C_\infty(a,b)\times C_\infty(b,c) \arrow[r, "{\comp_{\tau_0}}"] & C_\infty(a,c),
\end{tikzcd}
\]
\msp{page 127}
d'où $\comp_{\tau_0} = \comp_{\tau'_0}$.

Je recule devant la vraie difficulté - celle de prouver l'associativité de cette loi de composition. Elle devient associative dans $\Hot^{\wedge}_{0}$ - mais pour que deux flèches dans Hot soient égales (ici donc
\[ C_\infty(a,b)\times C_\infty(b,c)\times C_\infty(c,d) \rightrightarrows C_\infty(a,d)\quad ), \]
suffit-il qu'il en soit ainsi de leurs images dans $\Hot^{\wedge}_{0}$ - en d'autres termes, le foncteur
\[ \Hot\longrightarrow\Hot^{\wedge}_{0}, \]
est-il (en plus d'être conservatif) au moins fidèle (sinon pleinement fidèle)?

\par
\printedpage{80}
% head: A. Grothendieck : Dérivateurs, Ch. VII, Édition des Universités de Montpellier II et Paris VII
% foot: Version 22 décembre 2001

Il y a aussi que cette loi de composition est compatible avec celle des $H(a,b)$ dans Ind\,Hot, pour les flèches canoniques
\[ H(a,b)\longrightarrow C_\infty(a,b) \qquad \text{dans Ind\,Hot}, \]
si on saurait que ce sont des épimorphismes, et le restent par produits cartésiens, on aurait gagné.

Bien sûr, on saura que c'est associatif, une fois qu'on saura que la composititon définie ici est compatible avec la composition des chemins dans le cadre semi-simplicial. Mais comment fait-on dans ce cadre, pour composer les chemins de longueur non bornée? N'y a-t-il pas là exactement la même difficulté? À moins de s'en % [sic] composititon -> composition, misspelling % Reading resolved against the supplied PDF; see corrections.csv.
\msp{page 128}
tirer en se ramenant tout de suite, pour la théorie des opérations $\Omega$, aux complexes de Kan, et aux chemins de type $\Delta_+$. Mais pour les composer \textit{fonctoriellement,} il me semble y avoir toujours le même problème - on trouvera des chemins de type $\Delta_+\Delta_+ = \boldsymbol{\Delta}_2^+$, et non pas $\Delta_+$, mais on peut alors composer les homotopies directes (revient à s'induire sur $\Delta_+\simeq\{0,2\}\subseteq\boldsymbol{\Delta}_2^+$). Or j'ai fait le formalisme des catégories de chemins sans me soucier de conditions du type Kan. Je présume que dans le cadre semi-simplicial, les $\uHom(\boldsymbol{\Delta}_1,X)$ ne donnent \textit{pas} le bon objet des chemins, pour $X$ pas de Kan. (Car sinon, il semblerait que pour $X=X_0$ à composantes \textit{finies,} le calcul des $\pi_i$ devrait être fini, ce qui n'est pourtant pas le cas, m'a-t-on toujours assuré \dots)

Je présume pourtant, quand on prend $\Sigma=\mathrm{Fl}(X)$, que les objets $\Ch_\infty(X;a,b)$ et $\Ch_\infty(\bar{X};a,b)$ sont isomorphes dans Hot, quand on immerge $X$ dans une enveloppe de Thomason-Kan $\bar{X}$.

\par
\printedpage{81}
% head: A. Grothendieck : Dérivateurs, Ch. VII, Édition des Universités de Montpellier II et Paris VII
% foot: Version 22 décembre 2001

\msp{page 129}

\section*{11\quad Rapide retour sur les $(\mathcal{M},\Sigma)$, dans l'optique des dérivateurs}

Quand on s'intéresse au formalisme des $\Ch_\Sigma(X;a,b)$ en vue d'un \textit{dérivateur} $\mathbf{D}_{X,\Sigma}$ qu'on voudrait définir en termes de $X$, $\Sigma$, il semblerait que le prédérivateur qu'on obtient bel et bien ne mérite le nom de dérivateur que lorsqu'on suppose que pour toute flèche
\[ Y\longrightarrow Y' \]
dans Cat qui est une Hot-équivalence, les flèches correspondantes
\[ \mathrm{Hom}_{\mathrm{Ind\,Cat\,hot}}(Y',H(a,b)) \longrightarrow \mathrm{Hom}_{\mathrm{Ind\,Cat\,hot}}(Y,H(a,b)), \]
ou encore
\[ \varinjlim_n \mathrm{Hom}_{\mathrm{Cat\,hot}}(Y',\Ch_{n,\Sigma}(a,b)) \longrightarrow \varinjlim_n \mathrm{Hom}_{\mathrm{Cat\,hot}}(Y,\Ch_{n,\Sigma}(a,b)) \]
sont des bijections. J'ai regardé dans V ce que ça signifie pour l'objet $H(a,b) = (\Ch_{n,\Sigma}(a,b))_n$ de Ind\,Cat\,hot. C'est cette condition seulement qui assure, du point de vue des dérivateurs, une interprétation convaincante de $\Ch_{\infty,\Sigma}(a,b)$,
\[ \Ch_{\infty,\Sigma}(a,b) \;=\; \varinjlim_n \Ch_{n,\Sigma}(a,b) \]
en tant qu'objet de Hot, ou plus exactement de $\Hot^{\wedge}_{0}$, car on a alors pour $Y$ une catégorie \textit{finie}
\[ \underbrace{\mathrm{Hom}_{\mathrm{Cat\,hot}}(Y,H(a,b))}_{\simeq\,\mathrm{Hom}_{X(Y)\Sigma_Y^{-1}}(a_X,b_X)} \;\simeq\; \mathrm{Hom}_{\Hot}(Y,\Ch_{\infty,\Sigma}(a,b)), \]
\msp{page 130}
car je veux quand-même que le type dans Hot soit le bon, qu'il correspond bien à l'homotopie des espaces de chemins. Donc on voudrait qu'un morphisme
\[ X\xrightarrow{\;f\;}X' \]
dans (Cat), qui est une Hot-équivalence, définisse
\[ \Ch_\infty(X;a,b) \longrightarrow \Ch_\infty(X';f(a),f(b)), \]
qui soient eux aussi des Hot-équivalences. D'autre part, c'est évidemment compatible avec la brillante composition des chemins, que j'ai réussi à définir entre ces objets de Hot. Donc (sous réserve d'avoir vérifié ce que je viens de dire) la vérification de l'associativité se ramenerait au cas où $X$ est de Thomason-Kan. Mais dans ce cas, il y a des chances que, tout comme dans le cas semi-simplicial, les $\Ch_\infty(X;a,b)$ peuvent se définir (comme objets de (Hot)) comme étant simplement les chemins de type $\Delta_+$ de $a$ à $b$, qu'on compose alors de la façon triviale. Mais non, ces catégories sont discrètes! % [sic] ramenerait -> ramènerait, missing accent

Je renonce pour l'instant à poursuivre plus la question d'associativité. Demander à \textsc{Thomason} pour les propriétés du foncteur $\Hot\longrightarrow\Hot^{\wedge}_{0}$.

\par
\printedpage{82}
% head: A. Grothendieck : Dérivateurs, Ch. VII, Édition des Universités de Montpellier II et Paris VII
% foot: Version 22 décembre 2001

\msp{page 131}
Cela signifie, à très peu de choses près, qu'on a une interprétation de $\Ch_{\infty,\Sigma}(a,b)$, en tant qu'objet de Hot, comme un `Hom externe' $\mathrm{Hom\,ext}(a,b)$ à valeurs dans Hot, défini sur la catégorie $\mathcal{M}=X$ à l'aide du seul ensemble admissible de flèches $\Sigma\subseteq\mathrm{Fl}(\mathcal{M})$, opération satisfaisant la condition
\[ \boxed{\mathrm{Hom}_{\Hot}(Y,\mathrm{Hom\,ext}(a,b)) \;\simeq\; \mathrm{Hom}_{\mathcal{M}(Y)\Sigma_Y^{-1}}(a_X,b_X)} \tag*{$(*)$} \]
(qui peut s'interpréter aussi, soit comme
\[ \mathrm{Hom}_{\mathcal{M}}(p_!(a_Y),b) \]
si $p_! : \underbrace{\mathcal{M}(Y)\Sigma_Y^{-1}}_{\mathbf{D}(Y)} \longrightarrow \mathcal{M}$ existe, soit comme
\[ \mathrm{Hom}_{\mathcal{M}}(a,p_*(b_Y)) \]
si $p_* : \underbrace{\mathcal{M}(Y)\Sigma_Y^{-1}}_{\mathbf{D}(Y)} \longrightarrow \underbrace{\mathcal{M}}_{=\,\mathbf{D}(*)}$ existe); avec le seul petit ennui que cette formule n'est valable que pour $Y$ \textit{finie.} Si on veut exprimer le second membre pour $Y$ variable quelconque dans Cat, il ne suffit pas d'un objet de Hot ($\mathrm{Hom\,ext}(a,b)$) dans le premier membre, mais il faut l'ind-objet $H(a,b) = (\Ch_n(a,b))_{n\in\mathbf{N}}$ dans Hot. Si donc on tient à travailler dans un `domaine' Diag du dérivateur $\mathbf{D}=\mathbf{D}_{\mathcal{M},\Sigma}$ aussi grand que possible, en se bornant à des catégories finies,
\msp{page 132}
et qu'on veut qu'une formule $(*)$ reste valable pour toute $Y$ dans le domaine, il faut `en toute rigueur' définir le $\mathrm{Hom\,ext}(a,b)$ non comme objet de Hot, mais de Ind\,Hot. Il se décrirait donc alors comme
\[ \mathrm{Hom\,ext}(a,b) = (\Ch_{n,\Sigma}(a,b))_n \qquad \text{(ind-objet dans Hot)}\ , \]
et l'objet $\Ch_{\infty,\Sigma}$ dans Hot en est une sorte de limite inductive.

Je me suis demandé si la condition préalable, savoir que les foncteurs $H(a,b)$ sur Cat transforment Hot-équivalences $Y'\longrightarrow Y$ en bijections, était bel et bien une (forte) restriction. Je me suis convaincu que c'est bien le cas, déjà dans le cas particulier où $\Sigma=\mathrm{Fl}(X)$. Il me semble que lorsque les composantes connexes de $\Omega(X,a)$ ne sont pas contractiles, ladite condition n'est pratiquement jamais satisfaite. J'ai pensé p.~ex. au cas où $X$ a le type d'homotopie (dans Hot) d'une 2-sphère, de sorte que son $\Omega(X,a)=\Ch_\infty(X,a,a)$ doit avoir un $\pi_1\simeq\mathbf{Z}$. Prenons $Y$ ayant le type dans Hot d'un cercle, disons
\[
\begin{tikzcd}
 & & a & & \\
Y\colon & 0 \arrow[ur] \arrow[dr] & & 0' \arrow[ul] \arrow[dl] & , \\
 & & b & &
\end{tikzcd}
\]

\par
\printedpage{83}
% head: A. Grothendieck : Dérivateurs, Ch. VII, Édition des Universités de Montpellier II et Paris VII
% foot: Version 22 décembre 2001

\msp{page 133}
donc les foncteurs
\[ f : Y\longrightarrow \Ch_\infty(X;\bullet,\bullet) \overset{\text{déf}}{=} \Omega(X,\bullet) \]
correspondent à \textit{certains} `lacets' de $\Omega(X,\bullet)$ (mais lacets à origine non précisée), ils ont un type d'homotopie au sens topologique (`nombre de tours' autour de la sphère $X$) donné par un entier $n\in\mathbf{Z}$ (le `degré' du lacet $\lambda$). L'intuition évidente, c'est que l'ensemble des degrés $\delta$ qu'on peut obtenir ainsi, avec une catégorie \textit{finie} donnée $Y$ qui représente le cercle $S^1\in(\Hot)$, est fini - $|\delta|$ majoré sûrement par le nombre de ces flèches (donc au moins par 4 dans le cas envisagé, chaque flèche de $Y$ faisant faire \textit{au plus} un tour autour du pseudo-cercle $\Omega(S^2,\bullet)$) (\footnote[28]{probablement on ne trouve tout au plus que les degrés $\delta$ entre $-2$ et $+2$.}). D'autre part, quand on représente $S^1$ par des catégories $Y_n$ arbitrairement grandes (disons déduites des intervalles $\Phi_n=(\Delta_+\Delta_-)^n$ en identifiant la source et le but), on doit trouver des degrés arbitrairement élevés. Or toute $Y_n$ s'envoie dans $Y(=Y_2)$ par une flèche
\[ Y_n = Y' \longrightarrow Y = Y_2 \]
\msp{page 134}
qui est de degré 1, donc qui est une Hot-équivalence. Mais l'application correspondante
\[ \pi_0\uHom(Y,\Omega(X,\bullet)) \longrightarrow \pi_0\uHom(Y',\Omega(X,\bullet)) \]
n'est pas bijective. Ce genre d'argument me semble devoir marcher chaque fois que l'ensemble des groupes d'homotopie $\pi_j(X,x)$, pour $j\geq 2$ (donc l'ensemble des groupes d'homotopie $\pi_i$ de $\Omega(X,x)$, pour $i\geq 1$) est \textit{infini,} en prenant pour $Y$, $Y'$ différents représentants $Y_n$ de la sphère $S^i$, disons les $Y_n$ déduites de $\Phi_n$ $(n\geq 2)$ en prenant $\Phi_n^i = \Phi_n\times\Phi_n\times\cdots\times\Phi_n$ (qui est une sorte de $i$-cube combinatoire), et en identifiant son `bord' à un seul point. Je présume que le nombre d'élements dans $\pi_i(\Omega(X,x))$ qu'on peut obtenir à partir de flèches $Y_n\longrightarrow\Omega(X,x)$ est limité par la taille de $Y_n$, mais on peut sûrement les obtenir chacun, en prenant $n$ assez grand. Si donc $\xi\in\pi_i(\Omega)$ échappe à $Y_{n_0}$ ($n_0=0$, disons), on prend un $n'$ assez grand pour que $\xi$ n'échappe pas à $Y_{n'}$, et on choisit % [sic] d'élements -> d'éléments, missing accent
\[ Y_{n'} \longrightarrow Y_n\ , \]
\texttt{[}$n=n_0$\texttt{?]}

\msp{page 135}
qui soit une Hot-équivalence (p.~ex. à l'aide d'une rétraction de $\Phi_{n'}$ sur $\Phi_n$, compatible avec sources et buts).

Ceci attire fortement l'attention sur le fait que si $Y$, $C$ sont deux catégories, la catégorie $\uHom(Y,C)$ \textit{n'est pas} un représentant de $\uHom_{\Hot}(\mathrm{hot}(Y),\mathrm{hot}(C))$, on a seulement une flèche canonique
\[ \mathrm{hot}(\uHom(Y,C)) \longrightarrow \uHom_{\Hot}(\mathrm{hot}(Y),\mathrm{hot}(C)), \tag*{$(*)$} \]


\endgroup
% END UNIT Der-ChVII-456/Der-ChVII-456_p081-085.tex
% BEGIN UNIT Der-ChVII-456/Der-ChVII-456_p086-090.tex
\begingroup\makeatletter%
\newcommand{\printedpage}[1]{\marginpar{\footnotesize p.~#1}}%
\newcommand{\msp}[1]{\marginpar{\footnotesize [#1]}}%
\newcommand{\uHom}{\underline{\mathrm{Hom}}}%
\newcommand{\uCh}{\underline{\mathrm{Ch}}}%
\newcommand{\hot}{\mathrm{hot}}%
\makeatother\let\printedpage\derpp\let\msp\dermsp%
\let\mspin\DERmspin%


\printedpage{84}
% head: A. Grothendieck : Dérivateurs, Ch. VII, Édition des Universités de Montpellier II et Paris VII
% foot: Version 22 décembre 2001

i.e.
\[
\underbrace{\hot(\uHom(Y,C))\times\hot(Y)}_{\hot(\uHom(Y,C)\times Y)}\longrightarrow\hot(C),
\]
laquelle flèche est simplement déduite de
\[
\uHom(Y,C)\times Y\longrightarrow C.
\]
Mais la flèche $(*)$ n'a aucune envie d'être \texttt{[un]} isomorphisme. C'est parce qu'elle ne l'est déjà pas dans le cas particulier $Y=\boldsymbol{\Delta}^1$, essentiellement, qui nous oblige, pour faire une théorie des chemins dans les catégories qui soit topologiquement raisonnable, de ne pas nous borner à $\boldsymbol{\Delta}^1$ comme type de chemins, mais à prendre des intervalles composés $\tau\in\mathrm{Ty}$.

\msp{page 136}

Par contre, $(*)$ est \texttt{[un]} isomorphisme quand $C$ est une catégorie de Thomason-Kan. Il suffit de voir que pour tout $Z$ dans (Cat), le Hom de $\hot(Z)$ dans les deux membres est isomorphe par la flèche déduite de $(*)$. Or si $C$ est de Thomason-Kan, de même $\uHom(Y,C)$ (immédiat sur la définition de Thomason-Kan en termes de prolongements de
\begin{tikzcd}[baseline=(\tikzcdmatrixname-1-1.base), row sep=small]
A_0 \arrow[r, hook, "i"] \arrow[dr] & A \arrow[d, dotted] \\
& C
\end{tikzcd}
$\ldots$(\footnote[29]{un peu vif! Il faut savoir que les applications cofibrantes $i$ de \textsc{Thomason} le restent après multiplication cartésienne par $Y$ \ldots})), donc
\[
\mathrm{Hom}_{\mathrm{Hot}}(\hot(Z),\hot(\uHom(Y,C)))\xleftarrow{\ \sim\ }\pi_0(\underbrace{\uHom_{\mathrm{Cat}}(Z,\uHom(Y,C))}_{\uHom_{\mathrm{Cat}}(Z\times Y,C)})\simeq\pi_0\uHom(Z\times Y,C), % Reading resolved against the supplied PDF; see corrections.csv.
\]
et
\[
\begin{aligned}
&\mathrm{Hom}_{\mathrm{Hot}}(\hot(Z),\uHom_{\mathrm{Hot}}(\hot(Y),\hot(C)))\simeq\mathrm{Hom}_{\mathrm{Hot}}(\underbrace{\hot(Z)\times\hot(Y)}_{\hot(Z\times Y)},\hot(C))\\
&\qquad\underbrace{\xleftarrow{\ \sim\ }}_{\text{car $C$ est de Thomason-Kan}}\qquad\quad\pi_0(\uHom_{\mathrm{Cat}}(Z\times Y,C)),
\end{aligned}
\]
O.K.

On voit donc que les catégories $\uCh_{\infty,\Sigma}(X;a,b)$ n'ont aucune tendance à être des catégories de Thomason-Kan. Pourtant, son type dans Hot doit être `le bon', sans avoir d'abord à remplacer $X$ par une enveloppe de Thomason-Kan $\bar{X}$.

\printedpage{85}
% head: A. Grothendieck : Dérivateurs, Ch. VII, Édition des Universités de Montpellier II et Paris VII
% foot: Version 22 décembre 2001

\msp{page 137}

\section*{12\quad Retour sur les `catégories de chemins' dans \textmd{Cat}, comme catégorie de modèles pour \textmd{Hot}}

Je vois exactement deux applications intéressantes du formalisme des chemins dans les catégories. L'une, qui m'a surtout forcé la main, est pour pouvoir construire des Hom\,ext, Hom \textit{externes} à valeurs dans Hot, pour tout prédérivateur $\mathbf{D}$ défini par un couple $(\mathcal{M},\Sigma)$ (et je ne connais pas de prédérivateur qui ne soit défini ainsi).

L'autre application est pour avoir dans (Cat) une \textit{construction fonctorielle,} en termes d'un objet $X$ quelconque ou d'une flèche $X\longrightarrow Y$, de l'analogue du formalisme de \textsc{Cartan-Serre} sur les espaces de chemins. Dans ce cadre, il est entendu qu'on prend toujours $\Sigma_X=\mathrm{Fl}(X)$, ensemble de \textit{toutes} les flèches de $X$. Pour construire fonctoriellement en $X$ dans (Cat)
\[
\begin{tikzcd}[row sep=huge]
\uCh_\infty(X) \arrow[d] \\
X\times X,
\end{tikzcd}
\]
il faut avoir fixé un système projectif
\[
\underline{t}=(t_n)_{n\in\mathbf{N}}
\]
\msp{page 138}
dans \underline{Ty} qui soit cofinal dans \underline{Ty}, ce qui signifie simplement que parmi les $t_n$ il y en a de longueur utile arbitrairement grande (et comme
\[
n\longmapsto\mathrm{long.ut.}(t_n)
\]
est croissante, cela implique que long.ut.$(t_n)$ tend vers $+\infty$ pour $n\longrightarrow+\infty$). Les faits principaux à vérifier ensuite (si faire se peut) sont les propriétés essentielles des flèches dans le diagramme
\printedpage{86}
% head: A. Grothendieck : Dérivateurs, Ch. VII, Édition des Universités de Montpellier II et Paris VII
% foot: Version 22 décembre 2001
\[
\begin{tikzcd}[column sep=large, row sep=large]
& X \arrow[d, "\tau"] & \\
& \uCh_\infty(X) \arrow[d, "p"] \arrow[ddl, "\sigma"'] \arrow[ddr, "\beta"] & \\
& X\times X \arrow[dl, "\mathrm{pr}_1"] \arrow[dr, "\mathrm{pr}_2"'] & \\
X & & X\ ,
\end{tikzcd}
\qquad\qquad
\left|\begin{array}{l}
\sigma\tau=\beta\tau=\mathrm{id}_X\\
p\tau=\mathrm{diag}_X
\end{array}\right.
\]
et plus généralement, pour une flèche
\[
X'\xrightarrow{\ f\ }X
\]
dans Cat, des flèches dans le diagramme correspondant
\[
\begin{tikzcd}
& \overline{X'}=C(f) \arrow[dl] \arrow[dr, "g"] \arrow[ddrr, bend left=30, "\bar{f}"] & & \\
X' \arrow[ur, dotted, bend left=40, "\alpha"] \arrow[dr, "f"'] & |[draw, rounded corners]| \text{cart.} & \uCh_\infty(X) \arrow[dl, "\sigma"] \arrow[dr, "\beta"] & \\
& X & & X,
\end{tikzcd}
\qquad
\bar{f}\alpha=(\beta g)\alpha=f,
\]
et plus particulièrement de $\beta g$. La flèche
\msp{page 139}
en pointillés $\alpha$ est la section de $\overline{X'}$ sur $X'$ déduite de la section $\tau$ de $\uCh_\infty(X)$ sur $X$ qui réalise la décomposition de $f$ en
\[
X'\xrightarrow{\ \alpha\ }\overline{X'}\xrightarrow{\ \bar{f}\ }X,
\]
où $\bar{f}$ devrait être une flèche \textit{`fibrante'} et une Hot-équivalence, $\alpha$ une flèche \textit{`cofibrante'} au sens de \textsc{Thomason}. (Mais s'il en est ainsi, est à voir!) Voici pour le moment les propriétés essentielles que je voudrais vérifier.

\begin{enumerate}[label={}, leftmargin=*]
\item[1$^o$)] $p$ est lisse et propre, et est une HOT-fibration.

\item[2$^o$)] Donc il en est de même de $\sigma=\mathrm{pr}_1p$ et $\beta=\mathrm{pr}_2p$. De plus, ce sont des HOT-fibrations \textit{triviales,} i.e. à fibres asphériques.

\printedpage{87}
% head: A. Grothendieck : Dérivateurs, Ch. VII, Édition des Universités de Montpellier II et Paris VII
% foot: Version 22 décembre 2001

\item[3$^o$)] Plus généralement (pour ce qui concerne $\beta$), $\overline{X'}\xrightarrow{\ \bar{f}\ }X$ devrait encore être propre et lisse et une HOT-fibration, mais bien sûr en général pas triviale. (Elle devrait être triviale si et seulement si $f$ est une Hot-équivalence.) (\footnote[30]{résulte formellement de 1$^o$ et 2$^o$, cf. page 170.}).
\end{enumerate}

Un mot sur la `dualité'. Les énoncés sur $\sigma$ (resp. $\beta$) se traduisent en énoncés sur $\beta$ (resp. $\sigma$) par passage de $\underline{t}$ à $\underline{t}^o=(t^o_n)$, où $\tau\longmapsto\tau^o$ est le foncteur dans \underline{Ty} `passage à l'intervalle opposé'. On a en effet % Reading resolved against the supplied PDF; see corrections.csv.

\msp{page 140}

un isomorphisme fonctoriel en $X$ (passage au chemin opposé)
\[
\boxed{\begin{array}{ccc}
c & \longmapsto & c^o\\[4pt]
\uCh^{\underline{t}}_\infty(X) & \longrightarrow & \uCh^{\underline{t}^o}_\infty(X)
\end{array}}
\qquad\qquad
\text{(\footnotemark[31])}
\]
\footnotetext[31]{cf. page 143 pour le passage de $\underline{t}$ à $\underline{t}^-$.}
qui `échange source et but'. Bien sûr, on a intérêt à prendre $\underline{t}$ tel qu'on ait $\underline{t}=\underline{t}^o$ (c'est ce que j'avais fait, par un sain instinct, dès le début dans V), de sorte qu'on a un foncteur involutif canonique
\[
\begin{array}{ccc}
c & \longmapsto & c^o\\[4pt]
\uCh_\infty(X) & \longrightarrow & \uCh_\infty(X)
\end{array}
\]
(je laisse tomber à nouveau l'exposant $\underline{t}$), qui `échange source et but'. Donc on trouve, à isomorphisme canonique près, la `même' factorisation de $f:X'\longrightarrow X$ en $\bar{f}\circ\alpha$, si on fait comme dans page 38 pour définir $\overline{X'}$ et $\alpha$, ou si dans ce diagramme on intervertit les rôles de $\sigma$ et de $\beta$. % [sic] page 38 -> page 138, the diagram defining \overline{X'} and \alpha stands on manuscript page 138 (apparently)

Enfin :

\begin{enumerate}[label={}, leftmargin=*]
\item[4$^o$)] Si
\[
c:a\longrightarrow b
\]
est un chemin quelconque dans $X$, on voudrait en conclure un isommorphisme dans Hot (de `transport par $c$') % [sic] isommorphisme -> isomorphisme, doubled m
\[
\overline{X'}_a\xrightarrow[\sim]{\bar{c}}\overline{X'}_b,
\]
avec bien sûr transitivité pour la composition,

\msp{page 141}

ce qui donnera aussi
\[
\overline{c^{\mathrm{opp}}}=\overline{c}^{-1}
\]
(en admettant que les chemins triviaux $a\xrightarrow{\mathrm{id}}a$ donnent l'identité $\overline{X'}_a\xrightarrow[\sim]{\mathrm{id}}\overline{X'}_a$ !) (\footnotemark[32]\textsuperscript{,}\footnotemark[33]). Bien sûr, $\overline{c}$ ne doit dépendre que de la \textit{classe d'homotopie} de $\overline{c}$ (dans % [sic] de $\overline{c}$ (dans -> de $c$ (dans, the homotopy class is apparently that of the path c
\footnotetext[32]{Surtout pas se borner au cas où $c$ serait un des $t_n$!}
\footnotetext[33]{décrire quand-même l'effet sur les $\pi_0$!}
\printedpage{88}
% head: A. Grothendieck : Dérivateurs, Ch. VII, Édition des Universités de Montpellier II et Paris VII
% foot: Version 22 décembre 2001
$\uCh_\tau(X;a,b)$, ou même dans $C=\uCh(X;a,b)$, catégorie de tous les chemins dans $X$, de $a$ vers $b$, de type non précisé). En d'autres termes, on doit avoir un foncteur
\[
\begin{array}{rcl}
\Pi(X) & \longrightarrow & \mathrm{Hot}\\[4pt]
a & \longmapsto & \hot(\overline{X'}_a)
\end{array}
\]
du groupoïde fondamental de $X$ dans Hot. À vrai dire, ceci n'a rien à voir avec la situation particulière étudiée ici - ça doit marcher chaque fois qu'on a une HOT-fibration, ici $\overline{X'}\longrightarrow X$, mais aussi $\uCh_\infty(X)\longrightarrow X\times X$ (\footnote[34]{Ici aussi, il faudrait décrire l'effet du transport $\overline{c}$ sur les $\pi_0$ des fibres.}). Et pour une $\mathbf{D}$-fibration $X'\longrightarrow X$ ($\mathbf{D}$ un dérivateur quelconque), et $a\in\mathbf{D}(e)$, on doit avoir de même deux foncteurs
\[
\Pi(X)\rightrightarrows\mathbf{D}(e)
\]
correspondant à
\[
\left\{\begin{array}{lcl}
x & \longmapsto & H^\bullet_{\mathbf{D}}(X'_x,\bullet)\\[4pt]
x & \longmapsto & H_\bullet^{\mathbf{D}}(X'_x,\bullet).
\end{array}\right.
\]

\msp{page 142}

Le cas qui nous occupe ici serait celui où $\mathbf{D}=\mathrm{HOT}$, où on prend $H^{\mathbf{D}}_\bullet$, et $a=$ objet final de $\mathbf{D}(e)=\mathrm{Hot}$.)

\item[5$^o$)] Si 1$^o$, 2$^o$, 3$^o$ sont vérifiés, on est moralement sûr que la construction des $\uCh_\infty(X)$ est bel et bien l'équivalent exact du formalisme de \textsc{Cartan-Serre}. Pour en avoir le c\oe ur totalement net, i.e. raccrocher explicitement cette théorie aux constructions classiques, il faudrait prouver un \textit{énoncé de comparaison} avec la théorie semi-simpliciale, qui dirait essentiellement ceci : Si on désigne, pour $X$ dans Cat, par $\tilde{X}$ l'objet associé dans $\boldsymbol{\Delta}^\wedge$ (le `nerf' de $X$), alors la flèche
\[
\uCh_\infty(X)\tilde{}\longrightarrow\underbrace{(X\times X)\tilde{}}_{\simeq\,\tilde{X}\times\tilde{X}}
\]
dans $\boldsymbol{\Delta}^\wedge/(X\times X)\tilde{}$, est $W$-équivalente sur $(X\times X)\tilde{}$ à la flèche
\[
\uCh(\tilde{X})\longrightarrow\tilde{X}\times\tilde{X}
\]
de la théorie semi-simpliciale des chemins (à supposer qu'il y ait une telle flèche classique définie pour \textit{tout} objet $X$ de $\boldsymbol{\Delta}^\wedge$). Ou qu'on ait tout au moins des isomorphismes canoniques dans Hot
\[
\hot_{\mathrm{Cat}}(\uCh_\infty(X;a,b))\ \simeq\ \hot_{\boldsymbol{\Delta}^\wedge}(\uCh(\tilde{X};\tilde{a},\tilde{b})).
\]
\end{enumerate}


\endgroup
% END UNIT Der-ChVII-456/Der-ChVII-456_p086-090.tex
% BEGIN UNIT Der-ChVII-456/Der-ChVII-456_p091-095.tex
\begingroup\makeatletter%
\newcommand{\printedpage}[1]{\marginpar{\footnotesize p.~#1}}%
\newcommand{\msp}[1]{\marginpar{\footnotesize [#1]}}%
\newcommand{\Ch}{\underline{\mathrm{Ch}}}%
\newcommand{\Ty}{\underline{\mathrm{Ty}}}%
\newcommand{\dashto}[1]{\mathrel{\overset{#1}{\rule[0.6ex]{2em}{0.4pt}}}}%
\makeatother\let\printedpage\derpp\let\msp\dermsp%
\let\mspin\DERmspin%


\printedpage{89}
% head: A. Grothendieck : Dérivateurs, Ch. VII, Édition des Universités de Montpellier II et Paris VII
% foot: Version 22 décembre 2001

\msp{page 143}
Dans $1^o$ et $3^o$, on veut prouver que certaines flèches ($p$ et $\bar{f}$) sont à la fois \textit{propre} et \textit{lisse} - ce qui est vraiment idéal ! Je dis qu'il suffit à chaque fois de prouver l'un des deux, que l'autre s'en déduit (comme `cas particulier') par dualité. Pour ceci, notons
\[
\tau \longmapsto \tau^{-}
\]
l'automorphisme (\textit{pas anti}-automorphisme) de $\Ty$ qui consiste à changer de signe toutes les flèches de transition - c'est donc, au niveau des objets, l'unique automorphisme du monoïde libre Ty, qui échange $\Delta_+$ et $\Delta_-$. Au système projectif $\underline{t}$ dans $\Ty$ correspond le système projectif $\underline{t}^-$. Ceci dit, on a un isomorphisme de catégories
\[
\boxed{\Ch^{\underline{t}}_{\infty}(X^o) \ \simeq\ (\Ch^{\underline{t}^-}_{\infty}(X))^o} \tag*{$(^{35})$} % Reading resolved against the supplied PDF; see corrections.csv.
\]
\footnotetext[35]{cf. page 140 pour le passage de $\underline{t}$ à $\underline{t}^o$.}
et le diagramme commutatif
\[
\begin{tikzcd}
\Ch^{\underline{t}}_{\infty}(X^o) \arrow[r, "\sim"] \arrow[d, "p^{\underline{t}}_{X^o}"] & \Ch^{\underline{t}^-}_{\infty}(X)^o \arrow[d, "(p^{\underline{t}^-}_{X})^o"] \\
X^o \times X^o \arrow[r, equal] & (X \times X)^o.
\end{tikzcd}
\]
Ainsi, dire que $p^{\underline{t}}_{X}$ \texttt{[plutôt} $p^{\underline{t}^-}_{X}$\texttt{]} est lisse, i.e. que $(p^{\underline{t}}_{X})^o$ \texttt{[plutôt} $(p^{\underline{t}^-}_{X})^o$\texttt{]} est propre, équivaut à dire que $p^{\underline{t}}_{X^o}$ est propre.

\msp{page 144}
Il nous reste donc, à ce sujet, que deux faits à vérifier : la lissité de $p : \Ch_{\infty}(X) \longrightarrow X \times X$, et celle de $\bar{f} : \overline{X'} \longrightarrow X$.

a) Lissité de $p : \Ch_{\infty}(X) \longrightarrow X \times X$. On se donne une flèche
\[
w : (a,b) \longrightarrow (a',b'), \qquad \text{i.e.} \quad \left\{ \begin{array}{lccc} u : & a & \longrightarrow & a' \\[1ex] v : & b & \longrightarrow & b', \end{array} \right.
\]
sur la base, un objet $c'$ au dessus du but, i.e.
\[
c' \in \Ch_{\infty}(X; a', b'),
\]
et on regarde la catégorie des flèches dans $\Ch_{\infty}(X)$, de but $c'$, qui relèvent $w$,
\[
\begin{tikzcd}
\underset{\text{\scriptsize variable}}{c} \arrow[r, "\varphi"] \arrow[d, dash] & \underset{\text{\scriptsize fixe}}{c'} \arrow[d, dash] \\
(a,b) \arrow[r, "w"] & (a',b'),
\end{tikzcd}
\]
\printedpage{90}
% head: A. Grothendieck : Dérivateurs, Ch. VII, Édition des Universités de Montpellier II et Paris VII
% foot: Version 22 décembre 2001
i.e. la catégorie des couples $(c,\varphi)$, où
\[
\left\{ \begin{array}{l}
c \in \Ch_{\infty}(X; a, b) \\[1ex]
\varphi : c \longrightarrow c' \ \ (\in \Ch_{\infty}(X; a', b')) \\[1ex]
\text{une flèche induisant } \begin{array}[t]{l} u \ \text{ sur les sources} \\ v \ \text{ sur les buts,} \end{array}
\end{array} \right.
\]
\[
\begin{tikzcd}
a \arrow[rrrr, dotted, dash, "c"] \arrow[d] & {} \arrow[d, dotted] & {} \arrow[d, dotted] & {} \arrow[d, dotted] & b \arrow[d] \\
a' \arrow[rrrr, dash, "c'"] & {} & {} & {} & b'.
\end{tikzcd}
\]

\msp{page 145}
\begin{flushright}
{\footnotesize 9.11.90}
\end{flushright}

Pour la question précédente, relative aux flèches dans le diagramme
\[
\begin{tikzcd}
& X \arrow[d, "i"] & \\
& \Ch_{\infty}(X) \arrow[d, "p"] \arrow[ddl, "\sigma"'] \arrow[ddr, "\beta"] & \\
& X \times X \arrow[dl, "\mathrm{pr}_1"] \arrow[dr, "\mathrm{pr}_2"'] & \\
X & & X,
\end{tikzcd}
\]
il faut préciser absolument la déscription de $\Ch_{\infty}(X)$ \textit{dans} Cat, donc il faut préciser le système projectif $\underline{t} = (t_n)_{n \in \mathbf{N}}$. Nous allons prendre en effet % [sic] déscription -> description, misspelling
\[
t_n = (\underbrace{\Delta_+ \Delta_-}_{\tau_0})^{2n},
\]
avec les morphismes de transition décrits par ailleurs
\[
\underbrace{t_{n+p}}_{=\tau_0^{2n+2p}} \ \longrightarrow \ \underbrace{t_n}_{\tau_0^{n}}, % [sic] \tau_0^{n} -> \tau_0^{2n}, since t_n = (\Delta_+\Delta_-)^{2n} = \tau_0^{2n}
\]
\printedpage{91}
% head: A. Grothendieck : Dérivateurs, Ch. VII, Édition des Universités de Montpellier II et Paris VII
% foot: Version 22 décembre 2001
étant donnés par $\delta_p * \mathrm{id}_{\tau_o^{2n}} * \delta_p$, où $\delta_p : \tau_0^p \longrightarrow 1$ est la dégénérescence de $\tau_0^p$. Ainsi on peut considérer $\Ch_{\infty}(X)$ comme étant la catégorie des chemins biïnfinis dans $X$ % [sic] \tau_o^{2n} -> \tau_0^{2n}, letter o printed in the index instead of 0
\[
\cdots \longrightarrow a_{-5} \longleftarrow \overbrace{a_{-4} \longrightarrow a_{-3} \xleftarrow{u_{-2}} \overbrace{a_{-2} \xrightarrow{u_{-1}} a_{-1} \xleftarrow{u_0} \overbrace{a_0}^{\Ch_0} \xrightarrow{u_1} a_1 \xleftarrow{u_2} a_2}^{\Ch_1} \xrightarrow{u_3} a_3 \longleftarrow a_4}^{\Ch_2} \longrightarrow a_5 \longleftarrow \cdots,
\]
avec les $u_i$ des identités pour $|i|$ assez grand.

La sous-catégorie (pleine)
\[
\Ch^{W}_{t_n}(X) \ = \ \Ch^{W}_{(\Delta_+\Delta_-)^{2n}}
\]
s'identifie à celle pour laquelle on a
\begin{quote}
$u_i$ une identité pour $i \notin [-2i+1, 2i]$, i.e. les seuls $u_i$ `significatifs' sont les $4i$ flèches $u_{-2i+1}$, $u_{-2i+2}$, $\ldots$, $u_{2i}$. % [sic] i (in [-2i+1,2i], 4i, u_{-2i+1}, u_{-2i+2}, u_{2i}) -> n, the index of t_n is n
\end{quote}

Mais il sera plus commode de regarder $\Ch_{\infty}(X)$ comme la limite inductive d'une
\msp{page 146}
famille à deux indices de sous-catégories pleines
\[
\Ch_{n,m}(X), \qquad n, m \in \mathbf{Z},
\]
en définissant cette catégorie comme \textit{vide} si $n > m$, et pour $n \leq m$ comme formée des objets
\[
\left\{ \begin{array}{l}
(a_i, u_i)_{i \in \mathbf{Z}} \in \Ch^{W}_{n,m}(X) \iff u_i \text{ une identité pour } i \notin [n, m-1] \\[1ex]
\overset{\mathrm{id}}{\cdots}\, a_{n-1} \dashto{\mathrm{id}} \underbrace{a_n \dashto{u_n} a_{n+1} \cdots \dashto{} a_{m-1} \dashto{u_{m-1}} a_m}_{\text{\scriptsize `partie utile' du chemin}} \dashto{\mathrm{id}} a_{m+1} \overset{\mathrm{id}}{\cdots}\,.
\end{array} \right.
\]
Donc pour $n = m$, on trouve la sous-catégorie pleine isomorphe à $X$, formée des chemins infinis `totalement dégénérés', i.e. tels que \textit{toutes} les $u_i$ soient des identités. (Cette sous-catégorie de $\Ch_{\infty}$ ne dépend donc pas du choix de $n = m$.)
\[
\xrightarrow{\mathrm{id}} \underbrace{a_{-1}}_{=\, a} \xleftarrow{\mathrm{id}} \underbrace{a_0}_{=\, a} \xrightarrow{\mathrm{id}} \underbrace{a_1}_{=\, a} \xleftarrow{\mathrm{id}}
\]
Tous les sommets d'un tel chemin sont égaux à un même objet de $X$, et les flèches de transition entre eux sont identiques.

La structure de $\Ch^{W}_{n,m}(X) \ \simeq \ \Ch^{W}_{\tau_{n,m}}(X)$ ne dépend que de la longueur $m-n$ de $\tau_{n,m}$ et de la \textit{parité} de $n$. Si $n$ est pair, la direction initiale de $\tau_{n,m}$ est positive, et
\[
\tau_{n,m} = \underbrace{\Delta_+ \Delta_- \cdots}_{m-n \text{ facteurs en alternance}} = \left\{ \begin{array}{ll}
(\Delta_+ \Delta_-)^{\frac{m-n}{2}} & \text{si } m \text{ pair} \\[1ex]
(\Delta_+ \Delta_-)^{\frac{m-n}{2}-1} \Delta_+ & \text{si } m \text{ impair}
\end{array} \right.
\]
\texttt{[plutôt} $\frac{m-n-1}{2}$ \texttt{au lieu de} $\frac{m-n}{2} - 1$\texttt{]}, tandis que pour $n$ impair c'est juste l'inverse
\printedpage{92}
% head: A. Grothendieck : Dérivateurs, Ch. VII, Édition des Universités de Montpellier II et Paris VII
% foot: Version 22 décembre 2001
\msp{page 147}
\[
\tau_{n,m} = \underbrace{\Delta_- \Delta_+ \cdots}_{m-n \text{ facteurs alternés}} = \left\{ \begin{array}{ll}
(\Delta_- \Delta_+)^{\frac{m-n}{2}} & \text{si } m \text{ impair} \\[1ex]
(\Delta_- \Delta_+)^{\frac{m-n}{2}-1} \Delta_- & \text{si } m \text{ pair}
\end{array} \right.
\]
\texttt{[plutôt} $\frac{m-n-1}{2}$ \texttt{au lieu de} $\frac{m-n}{2} - 1$\texttt{]}. La direction initiale est conforme à la parité de $n$, i.e. égale à $(-1)^n$, tandis que la direction terminale est opposée à la parité de $m$, i.e. égale à $(-1)^m$ \texttt{[plutôt} $(-1)^{m+1}$\texttt{]}.

Je note que l'on a $\tau_0 = (\tau_0)^o$, d'où
\[
\underline{t} = \underline{t}^o, \qquad \text{mais } \underline{t} \neq \underline{t}^-,
\]
et pourtant on a des isomorphismes fonctoriels en $X$
\[
\Ch_{\underline{t}}(X) \ \simeq \ \Ch_{\underline{t}^-}(X)
\]
respectant source et but, à cause d'un isomorphisme
\[
\underline{t} \ \simeq \ \underline{t}^- \qquad \text{dans Pro}(\Ty).
\]
On peut dire p. ex. que le \textit{foncteur translation} $T^p$ d'amplitude $p$ (pour $p \in \mathbf{Z}$) est défini sur $\Ch^{W}_{\underline{t}}(X)$, par
\[
\left\{ \begin{array}{l}
T^p(((a_i),(u_i))_{i \in \mathbf{Z}}) = ((b_j),(v_j))_{j \in \mathbf{Z}} \\[1ex]
\text{avec} \left\{ \begin{array}{lcl} b_j & = & a_{j-p} \\[1ex] v_j & = & u_{j-p} \end{array} \right.
\end{array} \right.
\]
et que
\[
\begin{aligned}
T^p &: \left\{ \begin{array}{lcl} \Ch_{t}(X) & \xrightarrow{\sim} & \Ch_{-t}(X) \\[1ex] \Ch_{-t}(X) & \xrightarrow{\sim} & \Ch_{t}(X) \end{array} \right\} \ \text{si } p \text{ impair} \\
T^p &: \left\{ \begin{array}{lcl} \Ch_{t}(X) & \xrightarrow{\sim} & \Ch_{t}(X) \\[1ex] \Ch_{-t}(X) & \xrightarrow{\sim} & \Ch_{-t}(X) \end{array} \right\} \ \text{si } p \text{ pair.}
\end{aligned} % Reading resolved against the supplied PDF; see corrections.csv.
\]

\msp{page 148}
En d'autres termes, on peut formellement regarder la catégorie somme
\[
\Ch_{\pm}(X) = \Ch_{+}(X) \amalg \Ch_{-}(X)
\]
et l'opération translation $T$ comme un automorphisme qui échange les deux morceaux, et $T^p$ n'est autre que son itéré $p$-ème, qui échange les morceaux $+$ et $-$ si $p$ impair, qui les stabilise si $p$ pair.

\printedpage{93}
% head: A. Grothendieck : Dérivateurs, Ch. VII, Édition des Universités de Montpellier II et Paris VII
% foot: Version 22 décembre 2001

Il y a aussi l'automorphisme involutif de passage au chemin opposé, qui échange source et but, et que je définis (conformément à l'isomorphisme canonique identique
\[
\underline{t} \xrightarrow{\sim} \underline{t}^o \quad )
\]
comme la symétrie $S$
\[
\left\{ \begin{array}{l}
S((a_i),(u_i)) = ((b_j),(u_j)) \ \ \text{avec} \\[1ex]
\begin{array}{lcl} b_j & = & a_{-j} \\[1ex] u_j & = & u_{-j+1}. \end{array}
\end{array} \right.
\]
Si je le regarde comme opérant sur $\Ch^{W}_{\pm}(X)$, il est lié au foncteur $T$ par
\[
\left\{ \begin{array}{rcl}
S^2 & = & \mathrm{id} \\[1ex]
STS^{-1} & = & T^{-1},
\end{array} \right.
\]
de sorte que c'est le groupe semidirect de $\mathbf{Z}/2\mathbf{Z}$ par $\mathbf{Z}$ (groupe affine de dimension 1 de l'anneau $\mathbf{Z}$) qui opère sur $\Ch_{\pm}$.

\msp{page 149}
On est embarassé de trouver un automorphisme $T^2$ de la situation $\Ch_{\infty}(X) \longrightarrow X \times X$, sans qu'il soit évident que dans Hot, ce foncteur induise l'identité! Je ne vois aucun moyen de relier un chemin $c$ et $T^2 c$ par une flèche dans $\Ch_{\infty}$, ni même pas un chemin `canonique'. Mais je me dis que cet automorphisme est décrit par un automorphisme, non pas du système projectif % [sic] embarassé -> embarrassé, misspelling
\[
(\tau_{n,m})_{n,m \in \mathbf{Z}^o \times \mathbf{Z},\ n \leq m}
\]
(où $\mathbf{Z}^o$ désigne l'ensemble ordonné \textit{opposé} de $\mathbf{Z}$), mais du \textit{pro-objet} de $\Ty$ qu'il définit (qui est bien sûr `le même' pour tous les systèmes projectifs cofinaux à gauche dans $\Ty$). Il en résulte déjà que $T^2$ induit l'identité sur l'ind-objet de Cat\,hot, $\underline{\underline{\mathrm{Ch}}}_{\infty}(X)$, et a fortiori sur l'ind-objet correspondant de Hot, ou de $\mathrm{Hot}^{\wedge}_{0}$, ou enfin pour l'objet de $\mathrm{Hot}^{\wedge}_{0}$ lui-même. De même pour les $\Ch_{\infty}(X; a, b)$. Il s'ensuit, revenant à $\Ch_{\infty}(X)$ lui-même et aux $\Ch_{\infty}(X; a, b)$, que l'automorphisme induit par $T^2$ est l'identité dans $\mathrm{Hot}^{\wedge}_{0}$. J'espère que $\mathrm{Hot} \longrightarrow \mathrm{Hot}^{\wedge}_{0}$ est
\msp{page 150}
tout au moins fidèle, de sorte que ça prouvera que $T^2$ induit l'identité dans Hot également.

Je reviens aux questions $1^o$ à $4^o$ d'il y a quelques jours.

Pour $1^o$), il faut prouver que
\[
p : \Ch_{\infty}(X) \longrightarrow X \times X
\]
est lisse et propre et une Hot-fibration. Par dualité, $p$ propre résultera de $p$ lisse. Mais $\Ch_{\infty}(X)$ est $\varinjlim$ des $\Ch_{n,m}(X)$, avec $n \longrightarrow -\infty$, $m \longrightarrow +\infty$, et où on peut se borner à $n$ et $m$ \textit{pairs.} Or on voit que pour $n$, $m$ pairs, i.e. $\tau_{n,m}$ à flèche initiale directe, à flèche finale


\endgroup
% END UNIT Der-ChVII-456/Der-ChVII-456_p091-095.tex
% BEGIN UNIT Der-ChVII-456/Der-ChVII-456_p096-100.tex
\begingroup\makeatletter%
\newcommand{\printedpage}[1]{\marginpar{\footnotesize p.~#1}}%
\newcommand{\msp}[1]{\marginpar{\footnotesize [#1]}}%
\newcommand{\Ch}{\underline{\mathrm{Ch}}}%
\newcommand{\CCh}{\underline{\underline{\mathrm{Ch}}}}%
\newcommand{\Xb}{\overline{X'}}%
\newcommand{\Xbb}{\overline{\overline{X'}}}%
\newcommand{\Hotz}{\mathrm{Hot}^{\wedge}_{0}}%
\newcommand{\Ty}{\underline{\mathrm{Ty}}}%
\makeatother\let\printedpage\derpp\let\msp\dermsp%
\let\mspin\DERmspin%


\printedpage{94}
% head: A. Grothendieck : Dérivateurs, Ch. VII, Édition des Universités de Montpellier II et Paris VII
rétrograde, $\Ch_{n,m}(X) \simeq \Ch_{\tau_{n,m}}(X)$ est Cat-\textit{fibrée} sur $X \times X$, donc lisse (i.e. HOT-lisse), et cela passe aux $\varinjlim$ (\footnote[36]{\textbf{NB} $\Ch_\infty(X)$ n'est ni fibrée ni cofibrée sur $X \times X$ - car les foncteurs d'inclusion entre les $\Ch_{n,m}$ ne sont pas cartésiens resp. cocartésiens.}).

Ainsi $\Ch_\infty(X) \xrightarrow{p} X \times X$ est à la fois propre et lisse (\footnote[37]{12.11. Mais ça implique déjà que c'est une fibration !}). Dire que c'est une HOT-fibration, revient p. ex. à dire que pour
\[
(a,b) \xrightarrow{(u,v)} (a',b')
\]
dans $X \times X$, la flèche correspondante
\[
\Ch_\infty(X)/(a,b) \longrightarrow \Ch_\infty(X)/(a',b')
\]
est une Hot-équivalence. Or, comme $p$ est propre, les inclusions des fibres en $(a,b)$ et en $(a',b')$ dans l'un et l'autre membre sont des Hot-équivalences :
\msp{page 151}
\[
\begin{tikzcd}[row sep=large]
\Ch_\infty(X)/(a,b) \arrow[r, "\rho"] & \Ch_\infty(X)/(a',b') \\
\Ch_\infty(X;a,b) \arrow[u, hook, "\wr"] \arrow[r, dotted, "\varphi"] & \Ch_\infty(X;a',b') \arrow[u, hook, "\wr"]
\end{tikzcd}
\tag*{$(*)$}
\]
Ainsi on trouve une flèche en pointillés dans Hot, dont il faut prouver que c'est une Hot-équivalence. Mais en considérant
\[
u : a \longrightarrow a' \qquad v : b \longrightarrow b'
\]
comme des \textit{chemins} dans $X$, et
\[
u^o : a' \longleftarrow a
\]
comme un chemin rétrograde de $a'$ à $a$, la théorie de variance générale (dans IX) nous donne bel et bien un isomorphisme dans (Hot) :
\[
\Ch_\infty(X;a,b) \longrightarrow \Ch_\infty(X;a',b'),
\]
on a envie de dire que c'est $\varphi$, i.e. qu'en l'insérant à la place de $\varphi$, le carré ci-dessus devient commutatif (ce qui prouverait ce qu'on veut). Je pense que si on travaille avec les ind-objets $\CCh_\infty(X)/(a,b)$, $\CCh_\infty(X;a,b)$ etc., pour lesquels on a des accouplements qui s'explicitent agréablement aux crans finis, et donnent naissance aux $\varphi$, cela devrait être facile (quoiqu'un peu fastidieux). Cela ne prouverait pas pour autant la commutativité de $(*)$ dans Hot, mais
% foot: Version 22 décembre 2001
\printedpage{95}
% head: A. Grothendieck : Dérivateurs, Ch. VII, Édition des Universités de Montpellier II et Paris VII
\msp{page 152}
seulement dans $\Hotz$, ce qui n'est peut-être pas suffisant. Mais cela suffit pour prouver que la flèche $\rho$ dans $(*)$ devient \texttt{[un]} isomorphisme dans $\Hotz$, donc elle est \texttt{[un]} isomorphisme.

2$^o$) Il faut montrer que les HOT-fibrations
\[
\Ch_\infty(X) \overset{\sigma}{\underset{\beta}{\rightrightarrows}} X
\]
sont triviales, ou ce qui revient au même, qu'elles sont à fibres asphériques. Il suffit de le voir pour $\sigma$, donc que
\[
\Ch_\infty(X;a,?)
\]
est asphérique. Mais cette catégorie est la limite inductive des $\Ch_n(X;a,?)$, i.e. des
\[
\underline{\mathrm{Hom}}^{\text{source fixée}}(\tau^n, X).
\]
Mais comme l'intervalle $I_{\tau^n}$ se rétracte par déformation sur son origine, il est immédiat que cette catégorie \underline{Hom} est même contractile.

3$^o$) Considérons
\[
\begin{tikzcd}[column sep=small, row sep=large]
 & \Xb = C(f) \arrow[dl] \arrow[dr, "g"] \arrow[ddrr, bend left, "\bar f"] & & \\
X' \arrow[ur, dotted, bend left, "\alpha"] \arrow[dr, "f"'] & \underline{\text{cart.}} & \Ch_\infty(X) \arrow[dl, "\sigma"'] \arrow[dr, "\beta"] & \\
 & X & & X,
\end{tikzcd}
\]
\msp{page 153}
je voudrais prouver que $\bar f$ est une \textbf{D}-fibration propre et lisse (\footnote[38]{Ça résulte formellement de ce qui a été déjà vu, cf. pages 170-171.}).

$\Xb$ est la catégorie des couples
\[
(x',c) \ \in \ \mathrm{Ob}\,X' \times \mathrm{Ch}_\infty(X)
\]
tels que
\[
\sigma(c) = f(x')
\]
et
\[
\bar f(x',c) = \beta(c), \qquad \text{but de } c.
\]
La fibre de $\Xb$, $\bar f$ en $x_0 \in X$ est donc la catégorie des couples $(x',c) \in \mathrm{Ob}\,X' \times \mathrm{Ch}_\infty(X)$, tels que
\[
\sigma(c) = f(x'), \qquad \beta(c) = x_0.
\]
% foot: Version 22 décembre 2001

\printedpage{96}
% head: A. Grothendieck : Dérivateurs, Ch. VII, Édition des Universités de Montpellier II et Paris VII
On peut regarder encore $\Xb$ comme limite inductive des sous-catégories $\Xb_{n,m}$, correspondante aux couples $(x',c)$ avec $c \in \mathrm{Ch}_{n,m}(X)$. On trouve que pour $m$ pair (direction terminale de $\tau_{n,m}$ rétrograde) les $\Xb_{n,m}$ sont Cat-fibrées sur $X$, pour $m$ impair (direction terminale de $\tau_{n,m}$ directe) elles sont Cat-cofibrées. Donc on conclut encore comme tantôt, $\Xb$ étant (au choix) limite inductive filtrante de sous-catégories lisses sur $X$, ou propres sur $X$. % [sic] correspondante -> correspondantes, agreement with « sous-catégories »

\msp{page 154}
Il faut prouver enfin que $\Xb \xrightarrow{\bar f} X$ est une \textbf{D}-fibration (\footnote[39]{Cela résulte déjà du fait que c'est propre et lisse!}), ce qui revient p. ex. à dire que pour une flèche
\[
a \xrightarrow{v} b
\]
dans $X$, la flèche induite
\[
\Xb/a \xrightarrow{\bar v} \Xb/b
\]
est une Hot-équivalence. On sait déjà comme tantôt que l'on a un diagramme à flèches verticales des Hot-équivalences
\[
\begin{tikzcd}[row sep=large]
\Xb/a \arrow[r] & \Xb/b \\
\Xb_a \arrow[u, hook, "\wr"] \arrow[r, dotted] & \Xb_b, \arrow[u, hook, "\wr"]
\end{tikzcd}
\]
et tout revient à montrer qu'une certaine flèche entre $\Xb_a$ et $\Xb_b$ dans Hot est \texttt{[un]} isomorphisme. J'ai envie de définir directement une telle flèche dans Hot, voire dans Cat. Or pour les $m$ impairs, les $\Xb_{n,m}$ sont Cat-cofibrées sur $X$. Pour $\Xb_{-m,m}$ ($m$ impair) on trouve donc
\[
(\Xb_{-m,m})_a \xrightarrow{v_*} (\Xb_{-m,m})_b,
\]
\msp{page 155}
l'ennui est qu'elles ne sont pas compatibles avec les morphismes de transition pour $m$ variable. Elles le sont cependant \textit{à homotopie près} sûrement, donc on trouvera au moins une flèche d'ind-objets dans Cat\,hot
\[
\Xbb_a \xrightarrow{v_*} \Xbb_b,
\]
ce qui par passage à $\Hotz$ donne quand-même une flèche
\[
\Xb_a \xrightarrow{v_*} \Xb_b.
\]
En travaillant avec les $m$ impairs, on trouve de même des foncteurs en sens inverses % [sic] m impairs -> m pairs, the Cat-fibrées (giving v^*) are stated above for m pair
\[
(\Xb_{-m,m})_a \xleftarrow{v^*} (\Xb_{-m,m})_b,
\]
% foot: Version 22 décembre 2001
\printedpage{97}
% head: A. Grothendieck : Dérivateurs, Ch. VII, Édition des Universités de Montpellier II et Paris VII
qui ne se recollent pas pour $m$ variable, mais qui se recollent à homotopie près, pour donner une flèche dans Ind\,Cat\,hot
\[
\Xbb_a \xleftarrow{v^*} \Xbb_b,
\]
d'où aussi
\[
\Xb_a \xleftarrow{v^*} \Xb_b \qquad \text{dans } \Hotz.
\]
D'autre part, on doit pouvoir montrer soritalement,
\msp{page 156}
cran par cran, que $v_*v^*$ et $v^*v_*$ dans Ind\,Cat\,hot sont égales à l'identité, donc $v_*$ et $v^*$ sont des isomorphismes inverses l'un de l'autre. On doit vérifier d'autre part soritalement que
\[
\begin{tikzcd}[row sep=large]
\Xbb/a \arrow[r, "\bar v"] & \Xbb/b \\
\Xbb_a \arrow[u] \arrow[r, "v_*"] & \Xbb_b, \arrow[u]
\end{tikzcd}
\]
dans Ind\,Cat\,hot est commutatif. Passant au diagramme correspondant dans $\Hotz$, en prenant les limites inductives, on trouve un carré commutatif dans $\Hotz$
\[
\begin{tikzcd}[row sep=large]
\Xb/a \arrow[r, "\bar v"] & \Xb/b \\
\Xb_a \arrow[u, "\wr"'] \arrow[r, "v_*", "\sim"'] & \Xb_b, \arrow[u, "\wr"']
\end{tikzcd}
\]
d'où résulte que $\bar v$ est \texttt{[un]} isomorphisme dans $\Hotz$, donc aussi dans Hot.

Pour vérifier que la HOT-fibration $\bar f$ est triviale si et seulement si $f$ est une HOT-équivalence, il faudrait avoir développé dans Cat le formalisme de la suite exacte des lacets - on est tout près visiblement !
% foot: Version 22 décembre 2001

\printedpage{98}
% head: A. Grothendieck : Dérivateurs, Ch. VII, Édition des Universités de Montpellier II et Paris VII
\msp{page 157}

\begin{flushright}
\footnotesize 13.11.90
\end{flushright}

\section*{13\quad Où on prouve enfin que $\Ch(X;-,-)$ est aussi bon (voire meilleur) que $\Ch_\infty(X;-,-)$.}

Je crois finalement que les catégories $\Ch(X)$, $\Ch(X;a,b)$ (sans indice $\infty$!) sont `bonnes', que les flèches canoniques dans Hot
\[
\left\{
\begin{aligned}
\Ch_\infty(X) \quad &\xrightarrow{\sim} \quad \Ch(X)\\
\Ch_\infty(X;a,b) \quad &\xrightarrow{\sim} \quad \Ch(X;a,b)
\end{aligned}
\right.
\]
sont des isomorphismes. (La première est même \texttt{[un]} isomorphisme dans $\mathrm{HOT}(X \times X)$, ce qui implique la deuxième relation, en passant aux fibres \dots) Je crois que je suis à présent (après XI) en mesure de le prouver. Cela résultera des deux résultats-clef suivants :

\begin{enumerate}
\item[1)] Le foncteur canonique
\[
X \longrightarrow \Ch(X)
\]
(associant à tout $x \in \mathrm{Ob}\,X$ le chemin vide de source et but $X$) est une Hot-équivalence, et : % [sic] but $X$ -> but $x$, the object is $x \in \mathrm{Ob}\,X$
\item[2)] Le foncteur source-but
\[
\Ch(X) \longrightarrow X \times X
\]
est \textit{lisse} et \textit{propre} (donc une `Hot-fibration parfaite').
\end{enumerate}

Cela doit donner ce qu'on cherche. Comme la $X \times X$-flèche $\Ch_\infty(X) \longrightarrow \Ch(X)$ (et les flèches déduites $\Ch_\infty(X;a,b) \longrightarrow \Ch(X;a,b)$) a été définie par l'intermédiaire de la catégorie
\[
\mathcal{C}_{\mathbf{N}^o} = \mathbf{N}^o \times_{\Ty} \mathcal{C}
\qquad
\begin{aligned}
&\text{(où $\mathcal{C} = \Ch(X)$ \texttt{[est]} considérée comme catégorie fibrée sur $\Ty$,}\\
&\text{et $\mathbf{N}^o \longrightarrow \Ty$ est le système projectif des $\underbrace{(\Delta_+\Delta_-)^{2n}}_{t_n} \dots$)}
\end{aligned}
\]
\msp{page 158}
à l'aide des deux flèches canoniques
\[
\begin{tikzcd}[row sep=huge, column sep=large]
\mathcal{C}_{\mathbf{N}^o} \arrow[r] \arrow[d] & \Ch_\infty(X) \\
\Ch(X), &
\end{tikzcd}
\]
% foot: Version 22 décembre 2001


\endgroup
% END UNIT Der-ChVII-456/Der-ChVII-456_p096-100.tex
% BEGIN UNIT Der-ChVII-456/Der-ChVII-456_p101-105.tex
\begingroup\makeatletter%
\newcommand{\printedpage}[1]{\marginpar{\footnotesize p.~#1}}%
\newcommand{\msp}[1]{\marginpar{\footnotesize [#1]}}%
\newcommand{\Ch}{\underline{\mathrm{Ch}}}%
\newcommand{\Ty}{\underline{\mathrm{Ty}}}%
\newcommand{\cC}{\mathcal{C}}%
\newcommand{\cX}{\mathcal{X}}%
\newcommand{\NN}{\mathbf{N}^{o}}%
\newcommand{\Ob}{\operatorname{Ob}}%
\makeatother\let\printedpage\derpp\let\msp\dermsp%
\let\mspin\DERmspin%


\printedpage{99}
% head: A. Grothendieck : Dérivateurs, Ch. VII, Édition des Universités de Montpellier II et Paris VII
on sera obligé de prouver aussi les résultats qui correspondent à 1) 2) pour $\cC_{\NN}$, i.e.

\begin{itemize}[label={}]
\item[1$'$)]
\[
\begin{tikzcd}
X \ar[r] \ar[d,"\wr"'] & \cC_{\NN} \\
\cC_{t_0} = \cC_0 &
\end{tikzcd}
\]
(l'inclusion canonique de $\cC_0$, fibre de $\cC_{\NN}$ (fibrée sur $\mathbf{N}^o$) en $0 \in \mathbf{N}^o$) est une Hot-équivalence, et :

\item[2$'$)] $\cC_{\NN} \longrightarrow X \times X$ est propre et lisse.
\end{itemize}

Moyennant ces quatre résultats techniques, que je vais prouver tantôt (et $1'$ et $2'$ ne font aucune surprise pour moi), on gagne. Considérons en effet le diagramme
\[
\begin{tikzcd}
& X \ar[dl,"\alpha_1"'] \ar[d,"\alpha_0"] \ar[dr,"\alpha_2"] & \\
\Ch(X) \ar[dr,"p_1"'] & \cC_{\NN} \ar[l,"\pi_1"'] \ar[r,"\pi_2"] \ar[d,"p_0"] & \Ch_{\infty}(X) \ar[dl,"p_2"] \\
& X \times X. &
\end{tikzcd}
\]

Par 1) et $1'$) (concernant $\alpha_1$ et $\alpha_0$) et le résultat
\msp{page 159}
similaire déjà connu pour $\alpha_2$, ces trois foncteurs sont des Hot-équivalences. Donc il en est de même de $\pi_1$, $\pi_2$. Comme $p_1$, $p_0$, $p_2$ sont des Hot-fibrations (étant propres et lisses), il s'ensuit que $\pi_1$ et $\pi_2$ induisent des Hot-équivalences sur les fibres
\[
\Ch(X;a,b) \quad \xleftarrow{\ \pi_1(a,b)\ } \quad \cC_{\NN}(a,b) \quad \xrightarrow{\ \pi_2(a,b)\ } \quad \Ch_{\infty}(X;a,b).
\]
Je l'avais vérifié (sans mal) déjà pour $\pi_2$ et les $\pi_2(a,b)$, ce qui avait permis, en inversant $\pi_2$ dans Hot, de définir \textit{dans} Hot
\[
\left\{
\begin{aligned}
\Ch_{\infty}(X) &\quad\longrightarrow\quad \Ch(X) \\
\Ch_{\infty}(X;a,b) &\quad\longrightarrow\quad \Ch(X;a,b),
\end{aligned}
\right.
\]
et dire que ce sont des isomorphismes signifie justement que $\pi_1$ et les $\pi_1(a,b)$ sont des isomorphismes.

Je ne crois pas que $\cC_{\NN}$ ait en lui-même d'autre interêt, que de permettre de comparer % [sic] interêt -> intérêt, missing acute accent
les $X \times X$-catégories $\Ch_{\infty}(X)$ et $\Ch(X)$, à l'aide du diagramme précédent. Par contre, $\Ch_{\infty}(X)$ et $\Ch(X)$ m'aparaissent comme les deux ingrédients essentiels du formalisme des % [sic] m'aparaissent -> m'apparaissent, spelling
chemins dans Cat.
% foot: Version 22 décembre 2001

\printedpage{100}
% head: A. Grothendieck : Dérivateurs, Ch. VII, Édition des Universités de Montpellier II et Paris VII
\msp{page 160}
Je vais prouver 1, 2, $1'$, $2'$.

1) Considérons le diagramme
\[
\begin{tikzcd}
X = (X_{\Ty})_e \ar[rr,hook,"i"] \ar[dr,dotted] & & X_{\Ty} = X \times \Ty \ar[rr,"j"] \ar[dr] & & \Ch(X) \ar[dl] \\
& e \ar[rr,dotted] & & \Ty\ , &
\end{tikzcd}
\]
où $X_{\Ty}$ et $\Ch(X)$ sont considérées comme catégories Cat-fibrées sur $\Ty$, et le foncteur $j$ est déduit des foncteurs canoniques
\[
X \longrightarrow \Ch(X)_{\tau} \overset{\text{déf}}{=} \Ch_{\tau}(X),
\]
associant à $x \in \Ob X$ le chemin $p_{\tau}^{*}(x)$, où $p_{\tau} : \tau \longrightarrow e$ est la projection de $\tau$ vers l'élément unité (et élément final) $e = 1_{\Ty}$ de $\Ty$, et $\Ch_{e}(X)$ est identifiée à $X$. Il faut vérifier le

\textbf{Lemme.} \textit{Pour tout $\tau \in \Ty$, le foncteur précédent est une} Hot-\textit{équivalence (et même un homotopisme).}

Mais on sait que pour deux catégories $\cX$, $\cC$ lisses sur une autre, $T$, tout $T$-morphisme de $\cX$ dans $\cC$ qui induit des Hot-équivalences sur les fibres, est une Hot-équivalence. Ainsi $j$ est
\msp{page 161}
est une Hot-équivalence. D'autre part, $\Ty$, ayant un objet final, est asphérique, donc % [sic] est ... est -> est, word repeated across the manuscript page break
$X \times \Ty \xrightarrow{\mathrm{pr}_1} X$ est une Hot-équivalence, d'où résulte aussitôt que $i$ est une Hot-équivalence. Donc de même $ji$, or $ji = \alpha_1$, q.e.d.

2) Prouvons p. ex. la lissité de $\Ch(X) \longrightarrow X \times X$, la propriété s'en déduira en appliquant ce résultat à $X^{\mathrm{opp}}$. Soit donc
\[
(a',b') \ \longrightarrow\ (a,b), \qquad \text{i.e. } a' \xrightarrow{u} a,\ \ b' \xrightarrow{v} b,
\]
une flèche dans $X \times X$, et $c$ un objet de $\Ch(X)$ au dessus de $(a,b)$ :
\[
\begin{tikzcd}[column sep=large]
a' \ar[r,dotted,"c'","\scriptstyle\rightarrow"'] \ar[d,"u"'] & b' \ar[d,"v"] \\
a \ar[r,"c","\scriptstyle\rightarrow"'] & b.
\end{tikzcd}
\]

On regarde la catégorie $\cC(c;u,v)$ des couples $(c',w)$ qui relèvent $(c,(u,v))$, où donc $c'$ est un chemin (de type non précisé) de $a'$ vers $b'$, et
\[
w : c' \longrightarrow c
\]
une flèche dans $\Ch(X)$ qui sur les objets source induise $u$, et sur les objets buts induise $v$ (ainsi $c'$ est un chemin de $a'$ vers $b'$). Il faut
% foot: Version 22 décembre 2001
\printedpage{101}
% head: A. Grothendieck : Dérivateurs, Ch. VII, Édition des Universités de Montpellier II et Paris VII
\msp{page 162}
prouver que cette catégorie $\cC(c;u,v)$ est asphérique. Pour cela, on regarde le chemin composé
\[
\tilde{c} = \mathrm{ch}(u)\ c\ \mathrm{ch}(v)^{o},
\]
où $\mathrm{ch}(u)$ est le chemin (direct) de $a'$ vers $a$ défini par $u : a' \longrightarrow a$, et $\mathrm{ch}(v)^{o}$ désigne donc le chemin (rétrograde) de $b$ à $b'$ défini par $v : b' \longrightarrow b$. Si $c$ est de type $\tau$, $\tilde{c}$ est donc de type $\Delta_{+}\tau\Delta_{-}$. On va comparer $\cC(c;u,v)$ et la catégorie $\Ch(X)/\tilde{c}$, à l'aide de deux foncteurs
\[
\begin{tikzcd}
\cC(c;u,v) \ar[r,shift left,"\varphi"] & \Ch(X;a',b')/\tilde{c}\ , \ar[l,shift left,"\psi"]
\end{tikzcd}
\]
dont on prouvera qu'ils définissent des flèches inverses l'une de l'autre dans Hot. Comme $\Ch(X)/\tilde{c}$ est contractile (ayant un objet final), donc asphérique, il en est de même de $\cC(c;u,v)$, et on gagnera.

On définit $\varphi$ par $\boxed{\varphi(c',w') = (\tilde{c}',\tilde{w}'),}$ où ($c'$ étant du type $\tau'$) $\tilde{c}'$ est le chemin de type $\Delta_{+}\tau'\Delta_{-}$ déduit de $c'$ comme image inverse par la dégénérescence $\Delta_{+}\tau'\Delta_{-} \longrightarrow \tau'$, i.e. on prolonge $c'$ d'un `cran $\Delta_{+}$' $a' \xrightarrow{\mathrm{id}} a'$ à gauche et d'un `cran $\Delta_{-}$' $b' \xleftarrow{\mathrm{id}} b'$ à droite, cf. figure.
\[
\overbrace{\underbrace{\begin{tikzcd}[ampersand replacement=\&]
a' \ar[r,"\mathrm{id}"] \ar[d,"\mathrm{id}"'] \ar[dr,phantom,"\ast"] \& a' \ar[rr,dotted,no head,"c'"] \ar[d,"u"'] \& {} \ar[d,dotted,"w'"] \& b' \ar[d,"v"'] \ar[dr,phantom,"\ast"] \& b' \ar[l,"\mathrm{id}"'] \ar[d,"\mathrm{id}"] \\
a' \ar[r,"u"] \& a' \ar[rr,no head,"c'"'] \& {} \& b' \& b' \ar[l,"v"'] % [sic] a' ... b' with label c' in the bottom row -> a ... b with label c, the bottom row of the figure is the path c~ = ch(u) c ch(v)^o
\end{tikzcd}}_{\tilde{c}}}^{\tilde{c}'}
\]

\msp{page 163}
On peut dire que
\[
\tilde{c}' = a'_{\Delta_{+}} \circ c' \circ b'_{\Delta_{-}},
\]
où de façon générale, si $x \in X$ et $\tau \in \mathrm{Ty}$, $x_{\tau}$ désigne le `chemin constant' de valeur $x$, type $\tau$. On a d'autre part des morphismes évidents
\[
a'_{\Delta_{+}} \longrightarrow \mathrm{ch}(u), \qquad b'_{\Delta_{-}} \longrightarrow \mathrm{ch}(v)^{\mathrm{opp}},
\]
comme montré [?] par les flèches verticales $(\mathrm{id}_{a'},u)$ et $(v,\mathrm{id}_{b'})$ dans les deux carrés marqués $\ast$ du diagramme, flèches qui sont l'une dans $\Ch_{\Delta_{+}}(X)$, l'autre dans $\Ch_{\Delta_{-}}(X)$. Ceci dit, on pose
\[
\tilde{w}' : \tilde{c}' \longrightarrow \tilde{c}, \qquad \tilde{w} = \lambda(u) \ast w' \ast \mu(v), % [sic] \tilde{w} -> \tilde{w}', the morphism just named is \tilde{w}'
\]
où
\[
\left\{
\begin{aligned}
\lambda(u) :\ \ a'_{\Delta_{+}} &\quad\longrightarrow\quad \mathrm{ch}(u) \\
\mu(v) :\ \ b'_{\Delta_{-}} &\quad\longrightarrow\quad \mathrm{ch}(v)^{\mathrm{opp}}
\end{aligned}
\right.
\]
% foot: Version 22 décembre 2001
\printedpage{102}
% head: A. Grothendieck : Dérivateurs, Ch. VII, Édition des Universités de Montpellier II et Paris VII
(cf. figure). Le couple
\[
(\tilde{c}',\tilde{w}') \in \Ob \Ch(X)/\tilde{c}
\]
ainsi défini dépend fonctoriellement de $(c',w')$ variable dans $\cC(c;u,v)$ : Si $(c'_1,w'_1)$ est un deuxième objet de cette catégorie, une flèche du premier dans le second est une flèche
\[
c' \xrightarrow{\alpha} c'_1
\]
dans $\Ch(X;a',b')$, telle qu'on ait commutativité dans
\[
\begin{tikzcd}
c' \ar[rr,"\alpha"] \ar[dr,"w'"'] & & c'_1 \ar[dl,"w'_1"] \\
& c, &
\end{tikzcd}
\qquad\qquad w' = w'_1\alpha.
\]

Une telle flèche $\alpha$ définit
\[
\tilde{\alpha} = a'_{\Delta_{+}} \ast \alpha \ast b'_{\Delta_{-}} : \tilde{c}' \longrightarrow \tilde{c}'_1
\]
dans $\Ch(X;a',b')$, et on vérifie
\msp{page 164}
aussitôt que l'on a commutativité dans
\[
\begin{tikzcd}
\tilde{c}' \ar[rr,"\tilde{\alpha}"] \ar[dr,"\tilde{w}'"'] & & \tilde{c}'_1 \ar[dl,"\tilde{w}'_1"] \\
& \tilde{c}, &
\end{tikzcd}
\qquad\qquad \text{i.e. } \tilde{w}' = \tilde{w}'_1\tilde{\alpha},
\]
i.e. que $\tilde{\alpha}$ est bien une flèche dans $\Ch(X;a',b')/\tilde{c}$. De même, que $\alpha \longmapsto \tilde{\alpha}$ définit une dépendance fonctorielle, i.e. $(\beta\alpha)\tilde{} = \tilde{\beta}\tilde{\alpha}$, est évident sur la définition de $\tilde{\alpha}$. Donc, on a bien défini % Reading resolved against the supplied PDF; see corrections.csv.
\[
\cC(c;u,v) \ \xrightarrow{\ \varphi\ }\ \Ch(X;a',b')/\tilde{c},
\]
définissons maintenant
\[
\cC(c;u,v) \ \xleftarrow{\ \psi\ }\ \Ch(X;a',b')/\tilde{c}.
\]
Soit $(c'',w'')$ un objet de $\Ch(X;a',b')/\tilde{c}$,
\[
\overbrace{\underbrace{\begin{tikzcd}[ampersand replacement=\&]
a'_0 = a' \ar[r] \ar[d,"\mathrm{id}"'] \& a'_1 \ar[rr,no head,"\rightarrow"] \ar[d,"w''_1"'] \& {} \ar[d,dotted,"w''"] \& a'_n \ar[d,"w''_n"'] \& b' \ar[l] \ar[d,equal] \\
a' \ar[r,"u"] \& a \ar[rr,no head] \& {} \& b \& b', \ar[l,"v"']
\end{tikzcd}}_{c}}^{c''}
\]
% foot: Version 22 décembre 2001

\printedpage{103}
% head: A. Grothendieck : Dérivateurs, Ch. VII, Édition des Universités de Montpellier II et Paris VII
donc
\[
w'' : c'' \longrightarrow \tilde{c}, \qquad\quad \text{flèche dans } \Ch(X;a',b'),
\]
notons qu'on a une flèche de rétraction évidente
\[
\tilde{c} \xrightarrow{\rho_c} c \qquad\quad \text{associée à } \Delta_{+}\tau\Delta_{-} \longrightarrow \tau \text{ (rétraction),}
\]
induisant $u$ et $v$ sur la source et le but :
\[
\begin{tikzcd}
a' \ar[r,"u"] \ar[d,"u"] & a \ar[rrr,no head,"c"] \ar[d,"\mathrm{id}"] & {} \ar[d,"\mathrm{id}"] & {} \ar[d,"\mathrm{id}"] & b \ar[d,"\mathrm{id}"] & b' \ar[l,"v"'] \ar[d,"v"] & \Big\}\ \tilde{c} & \in \Ch_{\Delta_{+}\tau\Delta_{-}}(a',b') \\
a \ar[r,"\mathrm{id}"] & a \ar[rrr,no head,"c"] & {} & {} & b & b' \ar[l,"v"'] & \Big\}\ \delta^{*}_{\tau}(c) & \in \Ch_{\Delta_{+}\tau\Delta_{-}}(a,b) \\
& a \ar[rrr,no head,"c"] & & & b & & \Big\}\ c & \in \Ch_{\tau}(a,b).
\end{tikzcd}
\]

\msp{page 165}
Donc on définit un objet
\[
(c',w') = \psi(c'',w'')\ \in \Ob \cC(c;u,v)
\]
par
\[
c' = c'', \qquad\quad w' = \rho_c w'',
\]
et c'est fonctoriel en $(c'',w'')$.

Le composé $\varphi\psi$ étant un endomorphisme de la catégorie contractile $\Ch(X;a',b')/\tilde{c}$, est homotope à l'identité. Il faut vérifier qu'il en est de même de $\psi\varphi$, endomorphisme de $\cC(c;u,v)$. Ce foncteur associe à $(c',w')$, $(\tilde{c}',\overline{w'})$, où $\tilde{c}'$ est défini comme plus haut, $\tilde{c}' = a'_{\Delta_{+}} \circ c' \circ b'_{\Delta_{-}}$, et où $\overline{w'}$ est le composé
\[
\underbrace{\tilde{c}'}_{= \delta^{\prime *}_{\tau'}(c')}\ \ \underbrace{\xrightarrow{\ \rho_{c'}\ }}_{\substack{\text{flèche}\\ \text{canonique}}}\ \ c' \xrightarrow{w'} c.
\]

On a donc
\[
(\psi\varphi)(c',w') = (\tilde{c}',\overline{w'}) \xrightarrow{\rho_{c'}} (c',w'),
\]
c'est bien un morphisme dans $\cC(c;u,v)$, i.e. on a commutativité dans
\[
\begin{tikzcd}
\tilde{c}' \ar[rr,"\rho_{c'}"] \ar[dr,"\overline{w'}"'] & & c' \ar[dl,"w'"] \\
& c &
\end{tikzcd}
\qquad\qquad \text{(flèche de } \Ch(X;a',b')\text{).}
\]
% foot: Version 22 décembre 2001


\endgroup
% END UNIT Der-ChVII-456/Der-ChVII-456_p101-105.tex
% BEGIN UNIT Der-ChVII-456/Der-ChVII-456_p106-109.tex
\begingroup\makeatletter%
\newcommand{\printedpage}[1]{\marginpar{\footnotesize p.~#1}}%
\newcommand{\msp}[1]{\marginpar{\footnotesize [#1]}}%
\newcommand{\CN}{\mathcal{C}_{\mathbf{N}^o}}%
\tikzset{zigzag/.style={-, postaction={decorate},%
  decoration={markings, mark=at position 0.12 with {\arrow{>}},%
  mark=at position 0.88 with {\arrow{<}}}}}%
\makeatother\let\printedpage\derpp\let\msp\dermsp%
\let\mspin\DERmspin%


\printedpage{104}
% head: A. Grothendieck : Dérivateurs, Ch. VII, Édition des Universités de Montpellier II et Paris VII
\msp{page 166}

Ainsi, le foncteur $\psi\varphi$ est bien homotope au foncteur identique. On a gagné !

Passons à $1'$, $2'$) pour $\CN$. L'assertion $1'$) se démontre exactement de la même manière que 1) pour \underline{Ch}$(X)$, en regardant $\CN$ comme fibrée sur $\mathbf{N}^o$ (qui a un objet final 0), et en utilisant le diagramme
\[
\begin{tikzcd}
X = (X_{\mathbf{N}^o})_e \arrow[rr, "i'"] & & X_{\mathbf{N}^o} = X \times \mathbf{N}^o \arrow[rr, "j'"] \arrow[dr] & & \mathcal{C}_{\mathbf{N}^o} \arrow[dl] \\
 & & & \mathbf{N}^o \ , &
\end{tikzcd}
\]
où $j'$ est (cartésien et) une Hot-équivalence sur chaque fibre d'après le lemme, donc une Hot-équivalence globalement, et $i'$ est une Hot-équivalence car $\mathbf{N}^o$ est asphérique (ayant un objet final). Donc $j'i'$ est asphérique, q.e.d.

Il reste à prouver que
\[ \CN \longrightarrow X \times X \]
est propre et lisse. Ici encore,
\msp{page 167}
c'est exactement le même argument que tantôt, je présume. Il y a cependant une complication technique. Nous avions pris, pour définir
\[ \underline{t} : \mathbf{N}^o \longrightarrow \underline{\mathrm{Ty}}, \qquad \underline{t} = (t_n)_n, \]
le choix
\[ t_n = (\Delta_+\Delta_-)^{2n}, \]
ou si on préfère
\[ t_n = (\Delta_+\Delta_-)^n, \]
en tous cas les $t_n$ sont à direction initiale $+$, terminale $-$, et ce sont des types réduits. Si on prend $\Delta_+t_n\Delta_-$, on trouve un type non réduit, qui n'est pas parmi les $t_n$. Il ne servirait à rien de changer la définition de $t_n$ en $t_n = (\Delta_-\Delta_+)^{\text{quelque chose}}$, car cette fois $\Delta_+t_n\Delta_-$ sera bien réduit, mais n'aura pas les bons signes initial et terminal. Donc pour définir $\tilde{c}$ en termes de $c$, comme objet de $\CN$, il faudra, si $c$ est de type $\tau = t_n$, prendre $\tilde{c}$ de type $\Delta_+\Delta_-\tau\Delta_+\Delta_- = t_{n+1}$ (en prenant $t_{n+1} = (\Delta_+\Delta_-)^{2n}$), et de même pour définir $\tilde{c}'$ en termes de $c'$. Les définitions pertinentes du passage de $(c',w')$ à $(\tilde{c}',\tilde{w}')$ sont décrites dans la figure : % [sic] t_{n+1} = (\Delta_+\Delta_-)^{2n} -> t_{n+1} = (\Delta_+\Delta_-)^{2n+2}, exponent inconsistent with t_n = (\Delta_+\Delta_-)^{2n}

\msp{page 168}
% foot: Version 22 décembre 2001

\printedpage{105}
% head: A. Grothendieck : Dérivateurs, Ch. VII, Édition des Universités de Montpellier II et Paris VII
\[
\tilde{w}'\Big\downarrow
\left\{
\overbrace{\underbrace{
\begin{tikzcd}[ampersand replacement=\&]
a' \arrow[r, "\mathrm{id}"] \arrow[d, "\mathrm{id}"'] \& a' \arrow[d, "u"'] \& a' \arrow[l, "\mathrm{id}"'] \arrow[d, "u"'] \arrow[rrrr, dotted, no head, "c'", "\longrightarrow"'] \& {} \arrow[d, "w'"] \& {} \arrow[d, "w'"] \& {} \arrow[d, "w'"] \& b' \arrow[r, "\mathrm{id}"] \arrow[d, "v"'] \& b' \arrow[d, "v"] \& b' \arrow[l, "\mathrm{id}"'] \arrow[d, "\mathrm{id}"] \\
a' \arrow[r, "u"] \& a \& a \arrow[l, "\mathrm{id}"'] \arrow[rrrr, no head, "c", "\longrightarrow"'] \& {} \& {} \& {} \& b \arrow[r, "\mathrm{id}"] \& b \& b'. \arrow[l, "v"']
\end{tikzcd}
}_{\tilde{c}}}^{\tilde{c}'}
\right\}
\tilde{w}'
\]

Ainsi
\[
\left\{
\begin{array}{lcl ll}
\tilde{c} & = & \xi_u \circ c \circ \eta_v & \text{où} &
\left\{
\begin{array}{lccccc}
\xi_u : & a' & \xrightarrow{u} & a & \xleftarrow{\mathrm{id}} & a \\[1ex]
\eta_v : & b & \xrightarrow{\mathrm{id}} & b & \xleftarrow{v} & b'
\end{array}
\right\}
\text{ type } \Delta_+\Delta_- \\[4ex]
\tilde{c}' & = & a'_{\Delta_+\Delta_-} \circ c' \circ b'_{\Delta_+\Delta_-} & & \\[4ex]
\tilde{w}' & = & \lambda(u) * w * \mu(v) & \text{avec} & % [sic] w -> w', the homomorphism of the pair (c',w')
\left\{
\begin{array}{lccc}
\lambda(u) : & a'_{\Delta_+\Delta_-} & \longrightarrow & \xi_u \\[1ex]
\lambda(v) : & b'_{\Delta_+\Delta_-} & \longrightarrow & \eta_v % [sic] \lambda(v) -> \mu(v), the formula above has \mu(v)
\end{array}
\right\}
\text{ comme dans \texttt{[la]} figure.}
\end{array}
\right.
\]
On trouve ainsi, mutatis mutandis,
\[ \CN(c;u,v) \xrightarrow{\varphi} \CN(a',b')/\tilde{c}, \qquad \varphi(c',w') = (\tilde{c}',\tilde{w}'), \]
et on définit le foncteur
\[ \CN(c;u,v) \xleftarrow{\psi} \CN(a',b')/\tilde{c} \]
ainsi : à $(c'',w'')$ dans $\CN(a',b')/\tilde{c}$, donc un chemin $c''$ de type $t_n = (\Delta_+\Delta_-)^{2n}$ de $a'$ à $b'$ et un homomorphisme $w'' : c'' \longrightarrow \tilde{c}$ à extrémités fixes, on associe $(c',w')$ dans $\CN(c;u,v)$, avec
\[ c' = c'', \qquad w' = \rho_c w'', \]
où
\[ \rho_c : \tilde{c} \longrightarrow c \]
est défini de façon similaire que tantôt, par
\msp{page 169}
\[
\begin{tikzcd}
a' \arrow[r, "u"] \arrow[d, "u"'] & a \arrow[d, "\mathrm{id}"] & a \arrow[l, "\mathrm{id}"'] \arrow[d, "\mathrm{id}"] \arrow[rr, zigzag] & {} \arrow[d, "\mathrm{id}"'] & b \arrow[r, "\mathrm{id}"] \arrow[d, "\mathrm{id}"] & b \arrow[d, "\mathrm{id}"] & b' \arrow[l, "v"'] \arrow[d, "v"] & \tilde{c},\ \text{ type } t_{n+1}. \\
a \arrow[r, "\mathrm{id}"] & a & a \arrow[l, "\mathrm{id}"'] \arrow[rr, zigzag, "c"] & {} & b \arrow[r, "\mathrm{id}"] & b & b' \arrow[l, "\mathrm{id}"'] & \varphi^*_{n,n+1}(c).
\end{tikzcd}
\]

On doit seulement vérifier que l'endofoncteur $\psi\varphi$ de $\CN(c;u,v)$ est homotope à l'identité. Mais ce foncteur transforme $(c',w')$ en $(\tilde{c}' = a'_{\Delta_+\Delta_-}c'b'_{\Delta_+\Delta_-}, \overline{w'})$, où $\overline{w'}$ est encore le composé
% foot: Version 22 décembre 2001
\printedpage{106}
% head: A. Grothendieck : Dérivateurs, Ch. VII, Édition des Universités de Montpellier II et Paris VII
\[
\begin{tikzcd}
\tilde{c}' \arrow[r, "\text{can.}"] \arrow[rr, bend right=20, "\overline{w'}"'] & c' \arrow[r, "w'"] & c.
\end{tikzcd}
\]

On a donc
\[ \psi\varphi \longrightarrow \mathrm{id}, \]
$\psi\varphi$ est homotope au foncteur identique.

Cela prouve la lissité de $\CN \longrightarrow X \times X$. Pour la propreté, il faut prouver la lissité de le cas de $X^o$, et du système projectif $\underline{t}^{\mathrm{opp}} \simeq (t_n^{\mathrm{opp}})_n$. Mais on a $\underline{t}^{\mathrm{opp}} = \underline{t}$, et il suffit d'appliquer ce qu'on a déjà prouvé au cas de $X^o$, pour voir que $\CN \longrightarrow X \times X$ est propre. % [sic] de le cas -> dans le cas, grammatical

Il ne reste plus que le lemme à prouver. Ça doit être essentiellement
\msp{page 170}
trivial, et je renonce à rédiger une démonstration.

Je vais maintenant montrer `formellement' que la factorisation `canonique' d'une flèche de Cat $f : X' \longrightarrow X$ en
\[ X' \xrightarrow{i_f} \overline{X'} \xrightarrow{\bar{f}} X, \qquad i_f \text{ Hot-équivalence,} \]
avec (on l'espérait) $\bar{f}$ propre et lisse, à partir d'une factorisation similaire pour la flèche diagonale $X \xrightarrow{\ \ \mathrm{diag}_X\ \ } X \times X$ en
\[
\begin{tikzcd}
X \arrow[d, "i_X\ \ \text{\footnotesize Hot-équivalence}"] \\
\mathcal{C}_X\ \ \text{`chemins'} \arrow[d, "p_X"] \\
X \times X
\end{tikzcd}
\]
est chose formelle. (Cf. page 139 pour les notations et \texttt{[la]} problématique.) On regarde le diagramme
\[
\begin{tikzcd}
 & \mathcal{C}_X \arrow[d, "p_X"] \arrow[dl, "\beta_X"'] \arrow[dr, phantom, "\text{cart.}"] & C(f) \arrow[l, "\mathrm{pr}'_2"'] \arrow[d, "p_{X,f}"] \arrow[dll, bend right=50, "\bar{f}"'] & [-1.5em] \simeq X' \times_X (C_X, \sigma = \mathrm{pr}_1 \circ p_X) \\ % Reading resolved against the supplied PDF; see corrections.csv.
X & X \times X \arrow[l, "\mathrm{pr}_2^X"] \arrow[d, "\mathrm{pr}_1^X"'] \arrow[dr, phantom, "\text{cart.}"] & X' \times X \arrow[l, "f \times \mathrm{id}_X = f'"] \arrow[d, "\mathrm{pr}_1"] & \\
 & X & X' \arrow[l, "f"'] \arrow[uu, bend right=60, dotted, "i_f"'] &
\end{tikzcd}
\]
% foot: Version 22 décembre 2001
\printedpage{107}
% head: A. Grothendieck : Dérivateurs, Ch. VII, Édition des Universités de Montpellier II et Paris VII
(\footnote[40]{NB $C(f)$ est défini comme l'image inverse sur $X'$ du fibré $C_X$ sur $X$ (par $\sigma_X = \mathrm{pr}_1^X \circ p_X$) par $f : X' \longrightarrow X$.}). Comme $p_X$ est propre et lisse, i.e. `parfait', son image inverse sur $X' \times X$, savoir $p_{X,f}$, l'est aussi; d'autre part $\mathrm{pr}_2^X \circ f'$ est la
\msp{page 171}
projection $\mathrm{pr}_2$ de $X' \times X$ sur $X$, donc c'est aussi un morphisme parfait (puisque $X' \longrightarrow e$ est parfait, donc il en est de même du composé
\[ (\mathrm{pr}_2^X \underbrace{f') \circ p_{X,f}} = \underbrace{\mathrm{pr}_2^X(p_X}\, \mathrm{pr}'_2) = \beta_X \mathrm{pr}'_2 \overset{\text{déf}}{=} \bar{f} \ \ ). \]

Je rappelle que la section $i_f$ de $C(f)$ sur $X'$ est définie à partir de la section $i_X$ de $(\mathcal{C}_X,\sigma_X)$ sur $X$. C'est une Hot-équivalence, car $\sigma_X = \mathrm{pr}_1^X p_X$ est parfait à fibres asphériques, donc universellement une Hot-équivalence, donc $\mathrm{pr}_1 p_{X,f} : C(f) \longrightarrow X'$ est aussi une Hot-équivalence, et toute section en est donc une aussi.
% foot: Version 22 décembre 2001


\endgroup
% END UNIT Der-ChVII-456/Der-ChVII-456_p106-109.tex
\clearpage\setcounter{footnote}{0}\phantomsection\addcontentsline{toc}{section}{Chapitre VIII}\markboth{Chapitre VIII}{Chapitre VIII}
% BEGIN UNIT Der-ChVIII-567/Der-ChVIII-567_p001-005.tex
\begingroup\makeatletter%
\newcommand{\printedpage}[1]{\marginpar{\footnotesize p.~#1}}%
\newcommand{\msp}[1]{\marginpar{\footnotesize [#1]}}%
\newcommand{\Grp}{\underline{\mathrm{Grp}}}%
\makeatother\let\printedpage\derpp\let\msp\dermsp%
\let\mspin\DERmspin%


% ===================== page 1 =====================
\printedpage{}
% head: (aucun ; page de titre non numérotée)

\begin{center}
{\Large\textsf{LES DÉRIVATEURS}}

\bigskip

\textbf{ALEXANDRE GROTHENDIECK}

\bigskip\bigskip

\textbf{Chapitre VIII}

\medskip

\textbf{1-types d'homotopie relatifs : leur intégration \dots}
\end{center}

\par

% ===================== page 2 =====================
\printedpage{}
% head: (aucun ; page non numérotée)

\par\noindent\textbf{Note des éditeurs de l’édition transcrite.}\par
Ce texte a été déchiffré et transcrit en \LaTeXe{} par M. Künzer. Il a été édité par M. Künzer, J. Malgoire, et G. Maltsiniotis. La transcription est aussi fidèle que possible au manuscrit. Pour les quelques corrections évidentes, ou rares commentaires des éditeurs, ainsi que pour la numérotation originale des pages du manuscrit, les caractères de machine à écrire \texttt{[typewriter]} entre crochets sont utilisés. Un point d'interrogation entre crochets signifie que l'on n'est pas sûr du mot qui précède.

Cette édition est provisoire. Les remarques, commentaires et corrections sont bienvenus. Envoyer un message à :

\textbf{[e-mail address removed]}

\hfill G. Maltsiniotis\hspace*{2cm}

\par

% ===================== page 3 =====================
\printedpage{1}
% head: A. Grothendieck : Dérivateurs, Ch. VIII, Édition des Universités de Montpellier II et Paris VII
% foot: Version 2 février 2002

\msp{page 1}

\section*{1\quad Intégration des $1$-types d'homotopie relatifs sur une catégorie.}

Les types d'homotopie relatifs sur une (petite) catégorie $S$ sont définis par les \textit{champs} (en catégories) sur $S^{\wedge}$, qui correspondent essentiellement aux contrafoncteurs $S^o \longrightarrow (\mathrm{Cat})$, ou encore aux catégories $X$ sur $S$ qui sont Cat-fibrées sur $S$. (Dans $(\mathrm{Cat}/S)_{\mathrm{fib}}$, on inverse les flèches qui sont des Hot-équivalences sur les fibres.) D'autre part, les 1-types d'homotopie relatifs sur $S$ doivent se décrire par des \textit{champs en groupoïdes} sur $S$. Si on a un type d'homotopie relatif, défini par
\[ X \longrightarrow S, \qquad\qquad \text{catégorie Cat-fibrée sur } S, \]
ou par
\[
\begin{aligned}
S^o \quad &\longrightarrow \quad \mathrm{Cat}\\
s \quad &\longmapsto \quad X_s\ ,
\end{aligned}
\]
le 1-type d'homotopie associé est défini par
\[ s \longmapsto \Pi(X_s), \qquad\qquad S^o \longrightarrow (\mathrm{Grp}), \qquad\qquad \text{catégorie des groupoïdes,} \]
où $\Pi(X_s)$ est le groupoïde fondamental de $X_s$. On note la catégorie fibrée correspondante sur $S$ par \fbox{$\Pi(X/S)$}. Il est vrai que si
\[ X \ \longrightarrow\ X' \]
est un homomorphisme de catégories fibrées sur $S$ qui induise une Hot-équivalence sur les fibres, elle induit une équivalence fibrée
\[ \Pi(X/S) \ \longrightarrow\ \Pi(X'/S), \]
mais non un isomorphisme. Il y aurait donc du travail à faire pour définir un foncteur
\[
\underbrace{\mathrm{Hot}(S)}_{\text{\footnotesize\begin{tabular}{c}catégorie des\\ types d'homotopie\\ relatifs sur $S$\end{tabular}}}
\quad\xrightarrow{\ \xi\,\longmapsto\,\Pi(\xi/S)\ }\quad
\underbrace{\Grp(S)}_{\text{\footnotesize\begin{tabular}{c}catégorie des\\ groupoïdes\\ fibrés sur $S$\\ (ou catégories\\ des champs en\\ groupoïdes\\ sur $S^{\wedge}$)\end{tabular}}}
\quad .
\]

Mais ce n'est pas pour l'instant mon propos. Plutôt, de donner une formule de transitivité en homologie non commutative, analogue à $R(gf)_* \simeq Rg_*Rf_*$ en cohomologie.

\par

% ===================== page 4 =====================
\printedpage{2}
% head: A. Grothendieck : Dérivateurs, Ch. VIII, Édition des Universités de Montpellier II et Paris VII
% foot: Version 2 février 2002

\msp{page 2}

Si
\[ f : S \longrightarrow T \]
est un foncteur, on va définir un foncteur
\[ f_! :\ \Grp(S) \ \longrightarrow\ \Grp(T) \]
(image directe de 1-types d'homotopie relatifs). Mais notons d'abord que si
\[ p :\ X \ \longrightarrow\ S \]
est une $S$-catégorie \textit{quelconque} (pas nécessairement Cat-fibrée), on lui associe néanmoins un type d'homotopie relatif, via la catégorie fibrée $\tilde{X}^{(S)}$, ou $\tilde{X}$, ayant comme fibres les $s\backslash X$ (en notant que $s' \longrightarrow s$ dans $S$ définit bien $s\backslash X \longrightarrow s'\backslash X$) (\footnote{Ayant $\tilde{X}$, on définit $\Pi(X/S) \overset{\text{déf}}{=} \Pi(\tilde{X}/S)$.}). On a un $S$-foncteur canonique
\[ X \xrightarrow{\ i\ } \tilde{X}, \qquad\qquad x \longmapsto (x, p(x), \underbrace{p(x)}_{=\,\text{‘}\,s\,\text{’}} \xrightarrow{\ \mathrm{id}\ } p(x)). \]
C'est une $S$-Hot-équivalence si $X$ était déjà fibrée. En fait, sur chaque fibre c'est un homotopisme car on définit alors $\tilde{X}_{s_0} \xrightarrow{\ j\ } X_{s_0}$ de façon évidente, de façon que
\[
X_{s_0} \underset{j}{\overset{i}{\rightleftarrows}} \tilde{X}_{s_0}, \qquad ji \simeq \mathrm{id}, \qquad ij \underbrace{\xrightarrow{\text{homotopie simple}}}_{\text{\footnotesize\begin{tabular}{l}NB cette flèche\\ n'est pas \texttt{[un]}\\ isomorphisme\\ en général.\end{tabular}}} \mathrm{id}.
\]
Revenons alors à $f : S \longrightarrow T$, et considérons un groupoïde fibré sur $S$,
\[ G \ \longrightarrow\ S, \]
ce n'est pas une catégorie fibrée sur $T$, mais on peut lui associer
\[ t \longmapsto t\backslash G, \qquad\qquad T^o \longrightarrow (\mathrm{Cat}), \]
d'où une catégorie fibrée sur $T$, dont on prend le $\Pi$ relatif. Mais il est plus simple de prendre directement
\[ t \ \longmapsto\ \Pi(t\backslash G), \]
qui définit une catégorie fibrée en groupoïdes sur $T$.

\par

% ===================== page 5 =====================
\printedpage{3}
% head: A. Grothendieck : Dérivateurs, Ch. VIII, Édition des Universités de Montpellier II et Paris VII
% foot: Version 2 février 2002

\msp{page 3}

Ainsi on trouve bien un foncteur (en fait, même un 2-foncteur de 2-catégories)
\[ f_! :\ \Grp(S) \ \longrightarrow\ \Grp(T). \]
Si
\[ u :\ G \ \longrightarrow\ G' \]
est un homomorphisme de groupoïdes fibrés sur $S$, qui induise des Cat-équivalences sur les fibres, alors
\[ f_!(u) :\ f_!(G) \ \longrightarrow\ f_!(G') \]
est aussi une Cat-équivalence sur les fibres. (Cela est formel, à partir du fait que $f_!$ est un 2-foncteur, compte tenu que si on a une équivalence sur les fibres, alors il y a un quasi-inverse -- c'est une histoire générale de catégories fibrées \dots)

Soit maintenant
\[ p :\ X \ \longrightarrow\ S, \]
et prenons le cas $G = \Pi(X/S)$. Je dis qu'on a alors une $T$-équivalence
\[ \boxed{\ f_!(\Pi(X/S)) \ \simeq\ \Pi(X/T),\ } \tag*{$(*)$} \]
en d'autres termes, pour connaître $\Pi(X/T)$ (à $T$-équivalence près), il suffit de connaître $\Pi(X/S)$, et de lui appliquer le foncteur $f_!$ d'intégration des 1-types d'homotopie relatifs, pour $f : S \longrightarrow T$.

Notons que l'on a trivialement
\[ \Pi(X/S) \ =\ p_!(e_X), \]
où $e_X$ désigne le $X$-groupoïde fibré trivial (toutes les fibres ponctuelles). Donc la formule précédente s'écrit comme une formule de transitivité
\[ f_!(p_!(\Gamma)) \ \simeq\ (fp)_!(\Gamma), \]
\msp{page 4}
où $\Gamma$ est le groupoïde fibré trivial sur $X$. Cela suggère de généraliser la formule $(*)$ en
\[ \boxed{\ (fp)_!(\Gamma) \ \simeq\ f_!(p_!(\Gamma))\ } \]
pour \textit{tout} groupoïde fibré sur $X$, dans une situation
\[ X \xrightarrow{\ p\ } S \xrightarrow{\ f\ } T. \]

Ce sera formel, si on établit une propriété de 2-adjonction pour le foncteur $f_!$ :


\endgroup
% END UNIT Der-ChVIII-567/Der-ChVIII-567_p001-005.tex
% BEGIN UNIT Der-ChVIII-567/Der-ChVIII-567_p006-010.tex
\begingroup\makeatletter%
\newcommand{\printedpage}[1]{\marginpar{\footnotesize p.~#1}}%
\newcommand{\msp}[1]{\marginpar{\footnotesize [#1]}}%
\newcommand{\uHom}{\underline{\mathrm{Hom}}}%
\newcommand{\uGrp}{\underline{\mathrm{Grp}}}%
\newcommand{\uuGrp}{\underline{\underline{\mathrm{Grp}}}}%
\newcommand{\fl}{\longrightarrow}%
\makeatother\let\printedpage\derpp\let\msp\dermsp%
\let\mspin\DERmspin%


\printedpage{4}
% head: A. Grothendieck : Dérivateurs, Ch. VIII, Édition des Universités de Montpellier II et Paris VII
% foot: Version 2 février 2002

\textbf{Théorème.} \textit{Soit $f : S \fl T$ une flèche dans} (Cat)\textit{, $G$ un champ en groupoïdes sur $S$, $H$ un champ en groupoïdes sur $T$. On a alors une équivalence de catégories canonique}
\[
\boxed{\ \uHom_{\uGrp(S)}(G, f^*H) \ \simeq\ \uHom_{\uGrp(T)}(f_!G, H).\ }
\]

Les flèches d'adjonction
\[
\left\{
\begin{array}{lclcrcl}
f_!f^*H & \fl & H & \qquad\qquad & H & \in & \mathrm{Ob}\,\uGrp(T)\\[4pt]
f^*f_!G & \longleftarrow & G & \qquad\qquad & G & \in & \mathrm{Ob}\,\uGrp(S)
\end{array}
\right.
\]
devraient se définir sans problème. Pour la deuxième, il faut définir pour $s \in S$
\[
G_s \ \fl\ (f^*(f_!(G)))_s = (f_!(G))_{f(s)} \overset{\text{déf}}{=} \Pi(s\backslash G),
\]
or on a
\[
G_s \fl s\backslash G,
\]
d'où
\[
\underbrace{\Pi(G_s)}_{\begin{subarray}{l} =\ G_s\\ \text{car } G_s\\ \text{est un}\\ \text{groupoïde}\end{subarray}} \ \fl \Pi(s\backslash G),
\]
O.K. C'est fonctoriel en $s$ \dots

\msp{page 5}
Pour la flèche $f_!f^*H \fl H$ sur $T$, il faut trouver pour tout $t \in \mathrm{Ob}\,T$, et fonctoriellement en $t$,
\[
\underbrace{(f_!f^*H)_t}_{\overset{\text{déf}}{=}\Pi(t\backslash f^*(H))} \quad \fl \quad \underbrace{H_t}_{=\,\Pi(H_t)} \ ,
\]
il suffirait donc de définir
\[
t\backslash f^*(H) \ \fl\ H_t,
\]
mais on a le diagramme cartésien
\[
\begin{tikzcd}[column sep=large, row sep=large]
f^*H = H \times_S T \arrow[r] \arrow[d] & H \arrow[d] \\
S \arrow[r, "f"] & T \ni t
\end{tikzcd}
\]
\texttt{[plutôt} $H \times_T S$\texttt{]} induisant
\[
f^*H \ \fl\ H,
\]
\printedpage{5}
% head: A. Grothendieck : Dérivateurs, Ch. VIII, Édition des Universités de Montpellier II et Paris VII
% foot: Version 2 février 2002
d'où
\[
t\backslash f^*H \ \fl\ t\backslash H \qquad\quad \text{(au lieu de s'envoyer dans $H_t$),}
\]
mais d'autre part, comme $H$ est fibrée sur $T$, donc lisse sur $T$, l'inclusion
\[
H_t \fl t\backslash H
\]
est une Hot-équivalence. Ou disons plutôt : $f^*H \fl H$ est un morphisme de catégories sur $T$, induisant
\[
\underbrace{\Pi(f^*H/T)}_{\overset{\text{déf}}{=}\ f_!(f^*H)} \ \fl\ \Pi(H/T),
\]

\msp{page 6}
or le $T$-foncteur canonique
\[
H \fl \Pi(H/T) \qquad\quad (= \text{catégorie sur } T \text{ définie par } t \longmapsto \Pi(t\backslash H))
\]
est une équivalence, parce que $H$ est un fibré en groupoïdes sur $T$. Donc on peut trouver un $T$-quasi-inverse
\[
\Pi(H/T) \fl H
\]
et en composant on trouve
\[
f_!(f^*H) \ \fl\ \Pi(H/T) \ \fl\ H
\]
la flèche cherchée.

Ayant les 2-flèches d'adjonction, je m'arrête là dans les vérifications !

\printedpage{6}
% head: A. Grothendieck : Dérivateurs, Ch. VIII, Édition des Universités de Montpellier II et Paris VII
% foot: Version 2 février 2002

\section*{2\quad Questions de changement de base}

\[
\begin{tikzcd}[column sep=large, row sep=large]
S & S' \arrow[l, "q"'] \arrow[d, "f'"] \\
T \arrow[from=1-1, "f"'] & T' \arrow[l, "p"]
\arrow[from=2-2, to=1-1, bend right=50, shorten <=18pt, shorten >=18pt, "\gamma"']
\end{tikzcd}
\qquad\qquad\qquad
p^*f_!(G) \longleftarrow f'_!q^*(G)
\tag*{$(*)$}
\]

La flèche est définie, pourvu qu'on ait commutativité ou une flèche de commutativité $\gamma : pf' \fl fq$ (pas nécessairement \texttt{[un]} isomorphisme). Par adjonction, il revient au même de donner
\[
\begin{tikzcd}[row sep=large]
p_!f'_!q^*G \arrow[r] \arrow[d, dash, "\wr"', "\text{transitivité}"] & f_!(G) \\
(pf')_!q^*G \arrow[d, "\gamma"] & \\
(fq)_!q^*G \arrow[d, dash, "\text{transitivité}\ \wr"'] & \\
f_!(q_!q^*)G. \arrow[uuur, dashed, "\text{car } q_!q^*G \longrightarrow G"'] &
\end{tikzcd}
\]

On se demande quand la flèche est \texttt{[un]} isomorphisme.

\msp{page 7}
Le théorème suivant doit être vrai, et résulte de mes fourbis généraux.

% vertical bar printed in the left margin beside the following theorem
\textbf{Théorème.} \textit{Si le carré $(*)$ est} cartésien \textit{dans} (Cat)\textit{, $\gamma$ un isomorphisme, et si on suppose que $f$ soit} 1\textit{-lisse, ou $p$} 1\textit{-propre, alors le $T'$-foncteur de changement de base est une équivalence de $T'$-catégories.}

Les ‘fourbis’ généraux seraient en termes d'un dérivateur $\mathbf{D}$ ad hoc, celui des 1-types d'homotopie relatifs. J'y reviendrai par la suite.

On se pose la question quand le 2-foncteur associé à $f : S \fl T$
\[
f^* :\ \uuGrp(T) \ \fl\ \uuGrp(S)
\]
\printedpage{7}
% head: A. Grothendieck : Dérivateurs, Ch. VIII, Édition des Universités de Montpellier II et Paris VII
% foot: Version 2 février 2002
est 0-fidèle, 1-fidèle, 2-fidèle (\footnote[2]{La $i$-fidélité de $f^*$ ($i \in \{0,1,2\}$) équivaut à celle des foncteurs (obtenus en prenant $G = e_T$)
\[
\underline{\Gamma}(H/T) \ \fl\ \underline{\Gamma}(f^*H/S).
\]}). Si $G$, $H$ sont des fibrés en groupoïdes sur $T$, on considère
\[
\uHom_{\uuGrp(T)}(G,H) \ \fl\ \uHom_{\uuGrp(S)}(f^*G, f^*H),
\]
et suivant que ce foncteur est toujours fidèle, ou pleinement fidèle, ou une équivalence de catégories, on dit que $f^*$ est 0-, 1-, ou 2-fidèle. Dans le cas 2-fidèle, cela implique que si $K$ est un fibré en groupoïdes sur $S$, s'il est équivalent à un $f^*G$, alors $G$ est déterminé à $T$-équivalence près, elle même déterminée à isomorphisme unique près.

\msp{page 8}
Soit $i \in \mathbf{N}$. On dit que $f$ est $i$-asphérique si $\forall t \in T$, $Z = S/t$ est $i$-asphérique, i.e.
\[
\pi_0(Z),\ \dots,\ \pi_i(Z) \qquad\quad \text{triviaux.}
\tag*{$(*)$}
\]

On définit aussi $(-1)$-asphérique : les $S/t$ sont non-vides. De même pour la notion de $i$-coasphérique.

En fait, pour $i \in [-1,+\infty[ \ \subseteq \mathbf{Z}$, on définit un dérivateur $\mathbf{D}_i$ des $i$-types d'homotopie relatifs (cf. plus bas). D'où les notions d'objet $\mathbf{D}_i$-asphérique ($=$ $i$-asphérique), de $\mathbf{D}_i$-équivalence, de foncteur $\mathbf{D}_i$-asphérique et $\mathbf{D}_i$-coasphérique, de foncteur $\mathbf{D}_i$-lisse et $\mathbf{D}_i$-propre, de $\mathbf{D}_i$-fibration et de $\mathbf{D}_i$-fibration \textit{triviale} (i.e. à fibres $\mathbf{D}_i$-asphériques). On doit avoir le

\textbf{Théorème 3.} \textit{Soit $f : S \fl T$ dans} (Cat)\textit{, $i \in \{0,1,2\}$. Supposons l'une des conditions suivantes satisfaite :}

\begin{itemize}[leftmargin=3em]
\item[\textit{a)}] \textit{$f$ est $(i-1)$-propre, et $(i-1)$-asphérique (s'exprime par la condition que certaines catégories, telles les $S/t$, soient $(i-1)$-asphériques).}
\item[\textit{b)}] \textit{$f$ est $(i-1)$-lisse, et $(i-1)$-coasphérique.}
\end{itemize}

\textit{Alors :}

\begin{itemize}[leftmargin=3em]
\item[\textit{1\textsuperscript{o})}] \textit{Le foncteur $f^*$ image inverse}
\[
f^* : \uGrp(T) \ \fl\ \uGrp(S)
\]
\textit{est $i$-fidèle.}
\item[\textit{2\textsuperscript{o})}] \textit{Dans le cas $i = 2$ ($f$ pleinement fidèle), on}
\printedpage{8}
% head: A. Grothendieck : Dérivateurs, Ch. VIII, Édition des Universités de Montpellier II et Paris VII
% foot: Version 2 février 2002
\msp{page 9}
\textit{a de plus ceci : si $G \in \mathrm{Ob}\,\uGrp(S)$ est un fibré en groupoïdes sur $S$, alors $G$ est équivalent à un fibré en groupoïdes de la forme $f^*(H)$ si et seulement si pour tout $t \in T$, la restriction $G_t$ de $G$ à $S_t$ est équivalente à un fibré en groupoïdes constant sur $S_t$. (}\textbf{NB} \textit{La condition a) ou b) implique déjà que $S_t$ est} 1\textit{-connexe, i.e. que son groupoïde fondamental est équivalent à $e$.)}
\item[\textit{3\textsuperscript{o})}] \textit{Soit $Z$ un objet} 1\textit{-connexe dans} (Cat)\textit{, $G$ un fibré en groupoïdes sur $Z$. Pour que $G$ soit $Z$-équivalent à un groupoïde constant sur $Z$, il faut et il suffit que pour toute flèche $u$ dans $Z$, le morphisme de groupoïdes $u^*$ induit sur les fibres de $G$ soit une $Z$-équivalence} (\footnote[3]{On dit alors que $G$ est \textit{essentiellement localement constant} sur $Z$.}).
\end{itemize}

\textbf{Corollaire.} \textit{La condition nécessaire et suffisante de 2\textsuperscript{o}) s'écrit aussi ainsi : pour toute flèche $u$ de $S$ ‘verticale’ sur $T$, i.e. telle que $f(u)$ soit une flèche identique, le foncteur $u^*$ induit par $u$ sur les fibres de $G$ soit une équivalence de groupoïdes.}


\endgroup
% END UNIT Der-ChVIII-567/Der-ChVIII-567_p006-010.tex
\clearpage\setcounter{footnote}{0}\phantomsection\addcontentsline{toc}{section}{Chapitre IX}\markboth{Chapitre IX}{Chapitre IX}
% BEGIN UNIT Der-ChIX-678/Der-ChIX-678_p001-005.tex
\begingroup\makeatletter%
\newcommand{\printedpage}[1]{\marginpar{\footnotesize p.~#1}}%
\newcommand{\msp}[1]{\marginpar{\footnotesize [#1]}}%
\newcommand{\Ch}{\mathrm{Ch}}%
\newcommand{\Hot}{\mathrm{Hot}}%
\newcommand{\Hom}{\mathrm{Hom}}%
\makeatother\let\printedpage\derpp\let\msp\dermsp%
\let\mspin\DERmspin%


\printedpage{[non numérotée]}
% head: (aucun titre courant imprimé)

\begin{center}
{\LARGE\textsf{LES DÉRIVATEURS}}

\bigskip

\textbf{ALEXANDRE GROTHENDIECK}

\bigskip\bigskip

\textbf{Chapitre IX}

\textbf{Retour sur les catégories de fractions $\boldsymbol{\mathcal{M}W^{-1}}$,\\
comparaison avec Quillen}
\end{center}

\par

\printedpage{[non numérotée]}
% head: (aucun titre courant imprimé)

\par\noindent\textbf{Note des éditeurs de l’édition transcrite.}\par
Ce texte a été déchiffré et transcrit en \LaTeXe{} par M. Künzer. Il a été édité par M. Künzer, J. Malgoire, et G. Maltsiniotis. La transcription est aussi fidèle que possible au manuscrit. Pour les quelques corrections évidentes, ou rares commentaires des éditeurs, ainsi que pour la numérotation originale des pages du manuscrit, les caractères de machine à écrire \texttt{[typewriter]} entre crochets sont utilisés. Un point d'interrogation entre crochets signifie que l'on n'est pas sûr du mot qui précède.

Cette édition est provisoire. Les remarques, commentaires et corrections sont bienvenus. Envoyer un message à :

\textbf{[e-mail address removed]}

\hfill G. Maltsiniotis\hspace*{2cm}

\par

\printedpage{1}
% head: A. Grothendieck : Dérivateurs, Ch. IX, Édition des Universités de Montpellier II et Paris VII
% foot: Version 12 mars 2002

\msp{page 1}

\section*{1\quad Fonctorialité de $\Ch^{\Sigma}_{\infty}(a,b)$ en $a$, $b$ dans $\mathcal{M}\Sigma^{-1}$}

On se donne une catégorie $\mathcal{M}$ avec \texttt{[un]} ensemble admissible de flèches $\Sigma \subseteq \mathrm{Fl}(\mathcal{M})$, d'où (en prenant p.~ex. $\underline{t} = (t_n)$, avec $t_n = (\underbrace{\Delta_+\Delta_-}_{\tau_0})^{2n}$, et $\underbrace{t_{n+1}}_{=\,\tau_0^{2n+2}} \longrightarrow \underbrace{t_n}_{=\,\tau_0^{2n}}$ défini comme $p * \mathrm{id}_{\tau_0^{2n}} * p$, où $p : \tau_0 \longrightarrow 1$ la dégénérescence.) On a alors une théorie des $\underline{\Ch}_{\Sigma,\infty}(\mathcal{M},a,b)$, que je note simplement $\Ch_\infty(a,b)$ en tant qu'objet de Hot. Des accouplements
\[
\Ch_\infty(a,b) \times \Ch_\infty(b,c) \longrightarrow \Ch_\infty(a,c) \qquad \text{dans Hot}
\]
avec (nous allons l'admettre) \textit{associativité} dans Hot. Cela définit donc des accouplements sur les $\pi_0(\Ch_\infty(a,b)) \overset{\text{déf}}{=} \pi_0(a,b)$, qui n'est autre que la composition des flèches dans $\mathcal{M}\Sigma^{-1}$. D'autre part, on a pour $a,b \in \mathrm{Ob}\,\mathcal{M}$
\[
\Hom_{\mathcal{M}}(a,b) \longrightarrow \pi_0(\Ch_\infty(a,b)) = \pi_0(a,b),
\]
compatible avec les compositions. Mais on peut l'ignorer pour l'instant, et regarder $\mathrm{Ob}\,\mathcal{M}$, ou simplement $\mathcal{M}_0$, avec la seule structure
\[
\left\{
\begin{array}{ccccccc}
\mathcal{M}_0 \times \mathcal{M}_0 & \longrightarrow & \Hot & \qquad & (a,b) & \longmapsto & \Ch_\infty(a,b) \\[1ex]
\mathcal{M}_0 \times \mathcal{M}_0 \times \mathcal{M}_0 & \longrightarrow & \mathrm{Fl}\,\Hot & \qquad & (a,b,c) & \longmapsto & (\Ch_\infty(a,b) \times \Ch_\infty(b,c) \underbrace{\longrightarrow}_{\substack{\text{flèche}\\ \text{dans Hot}}} \Ch_\infty(a,c))
\end{array}
\right.
\]
avec condition \texttt{[d']}associativité, et aussi une condition d'unités :
\[
\Ch_\infty(a,a) \text{ est muni d'un } 1_a \in \pi_0\Ch_\infty(a,a) = \pi_0(a,a),
\]
jouant le rôle d'une `identité bilatère', dans le sens suivant. Notons que les accouplements
\[
\begin{array}{ccccc}
\Ch_\infty(a,b) & \times & \Ch_\infty(b,c) & \longrightarrow & \Ch_\infty(a,c) \\[1ex]
P & \times & Q & \longrightarrow & R
\end{array}
\]
\msp{page 2}
définissent
\[
\begin{array}{ccc}
P & \longrightarrow & \underline{\Hom}_{\Hot}(Q,R) \\[1ex]
Q & \longrightarrow & \underline{\Hom}_{\Hot}(P,R),
\end{array}
\]
d'où en passant aux $\pi_0$
\[
\begin{array}{ccc}
\pi_0(P) & \longrightarrow & \Hom_{\Hot}(Q,R) \\[1ex]
\pi_0(Q) & \longrightarrow & \Hom_{\Hot}(P,R),
\end{array}
\]
\printedpage{2}
% head: A. Grothendieck : Dérivateurs, Ch. IX, Édition des Universités de Montpellier II et Paris VII
% foot: Version 12 mars 2002
en d'autres termes,
\[
\begin{array}{ccc}
\pi_0(a,b) & \longrightarrow & \Hom_{\Hot}(\Ch_\infty(b,c),\Ch_\infty(a,c)) \\[1ex]
\pi_0(b,c) & \longrightarrow & \Hom_{\Hot}(\Ch_\infty(a,b),\Ch_\infty(a,c)).
\end{array}
\]

On en déduit en particulier pour $a,b \in \mathcal{M}_0$
\[
\left\{
\begin{array}{cccc}
1) & \pi_0(a,a) & \longrightarrow & \Hom_{\Hot}(\Ch_\infty(a,b),\Ch_\infty(a,b)) \\[1ex]
2) & \pi_0(a,a) & \longrightarrow & \Hom_{\Hot}(\Ch_\infty(b,a),\Ch_\infty(b,a)),
\end{array}
\right.
\tag*{$(*)$}
\]
et un objet $\alpha_a$ de $\pi_0(a,a)$ est unité bilatère si ses deux images sont les homomorphismes identiques dans $\Ch_\infty(a,b)$ resp. $\Ch_\infty(b,a)$. Passant aux $\pi_0$ de ceux-ci, cela implique que $\alpha_a$ est unité bilatère pour la composition des $\pi_0(x,y)$, relativement à $b$. Si on veut que ce soit le cas pour \textit{tout $b$,} cela signifie (prenant $b = a$) que $\alpha_a$ est un objet unité bilatère pour la structure de composition associative (donc de monoïde) dans $\pi_0(a,a)$, donc que $\pi_0(a,a)$ est un monoïde \textit{unitaire,} ayant $\alpha_a$ pour objet unité. L'élément $\alpha_a$ est donc déterminé de façon unique. Et on veut que $\alpha_a$ existe
\msp{page 3}
pour tout $a \in \mathcal{M}_0$, donc $\mathcal{M}_0$, muni des $\pi_0(a,b)$ et de leur composition, devient une \textit{catégorie,} soit $\mathcal{M}_\Pi$; et les flèches identiques de celle-ci doivent jouer le rôle d'unité bilatère également pour les `Hom externes' $\Ch_\infty(a,b)$, à valeurs dans Hot.

On peut dire que pour $a = a_0$ fixé, les $\Ch_\infty(a_0,b)$ ($b$ variable) forment les valeurs d'un \textit{foncteur} covariant
\[
\mathcal{M}_\Pi \longrightarrow \Hot \qquad b \longmapsto \Ch_\infty(a_0,b),
\]
en associant à
\[
b \xrightarrow{\ u\ } b', \qquad u \in \pi_0(b,b'),
\]
la flèche
\[
\Ch_\infty(a,b) \longrightarrow \Ch_\infty(b') % [sic] \Ch_\infty(b') -> \Ch_\infty(a_0,b'), argument manquant (cf. la note des éditeurs qui suit)
\]
\texttt{[plutôt} $\Ch_\infty(a_0,b) \longrightarrow \Ch_\infty(a_0,b')$\texttt{]} correspondante $*_1$ [?]. Que ce soit bien là une loi fonctorielle doit résulter formellement de l'associativité des accouplements sur les $\Ch_\infty$. De même, on a une loi fonctorielle, pour $b$ fixé,
\[
\mathcal{M}^o_\Pi \longrightarrow \Hot, \qquad a \longmapsto \Ch_\infty(a,b),
\]
grâce à l'associativité. En fait, l'associativité impliquera aussi qu'on a un \textit{bifoncteur}
\[
\begin{array}{ccc}
\mathcal{M}^o_\Pi \times \mathcal{M}_\Pi & \longrightarrow & \Hot \\[1ex]
(a\,,\,b) & \longmapsto & \Ch_\infty(a,b),
\end{array}
\]
\printedpage{3}
% head: A. Grothendieck : Dérivateurs, Ch. IX, Édition des Universités de Montpellier II et Paris VII
% foot: Version 12 mars 2002
\msp{page 4}
car si
\[
a' \xrightarrow{\ u\ } a, \qquad b \xrightarrow{\ v\ } b',
\]
on a le diagramme
\[
\begin{tikzcd}[column sep=large, row sep=large]
\Ch_\infty(a,b) \arrow[r, "\substack{\text{comp. avec}\\ v}"] \arrow[d, "\text{comp. avec } u"'] & \Ch_\infty(a,b') \arrow[d, "\text{comp. avec } u"] \\
\Ch_\infty(a',b) \arrow[r, "\substack{\text{comp. avec}\\ v}"] & \Ch_\infty(a',b'),
\end{tikzcd}
\]
qui est commutatif grâce à l'associativité pour
\[
\Ch_\infty(a',a),\ \Ch_\infty(a,b),\ \Ch_\infty(b,b').
\]

\textbf{Corollaire} de cette dépendance bifonctorielle : \textit{Si $u : a' \longrightarrow a$, $v : b \longrightarrow b'$ sont dans $\mathcal{M}_\Pi$, alors la flèche correspondante}
\[
\Ch_\infty(a,b) \longrightarrow \Ch_\infty(a',b') \qquad \textit{dans } \Hot
\]
\textit{est un} isomorphisme.

Les mêmes arguments devraient marcher si au lieu des $\Hom\,\mathrm{ext}$ à valeurs dans Hot, on les prenait dans Ind\,Hot, ou dans Cat\,hot, ou dans Ind\,Cat\,hot -- je pense bien sûr aux $\mathrm{CH}_\infty(a,b)$ (dans la théorie de la localisation) dans Ind\,Cat\,hot, qui peuvent être regardés comme provenant d'un bifoncteur (où $\mathcal{M}_\Pi = \mathcal{M}\Sigma^{-1}$)
\[
\mathcal{M}_\Pi \times \mathcal{M}_\Pi \longrightarrow \mathrm{Ind}\,\mathrm{Cat}\,\mathrm{hot}.
\]

\msp{page 5}
Voici l'axiomatisation définitive (?). On remplace Hot par une catégorie

\begin{itemize}[label={$\mathcal{H}$,}]
\item catégorie avec produit $*$ associatif, commutatif (?) (\footnote{On ne doit pas avoir besoin de la commutativité.}), unitaire, et toutes les compatibilités habituelles au besoin,
\end{itemize}

l'objet unité pour $*$ sera noté 1, et on notera
\[
\left\{
\begin{array}{ccc}
\pi_0(\xi) & \overset{\text{déf}}{=} & \Hom_{\mathcal{H}}(1,\xi) \\[1ex]
\pi_0 : \mathcal{H} & \longrightarrow & \mathrm{Ens}.
\end{array}
\right.
\]

Si on a
\[
P * Q \xrightarrow{\ u\ } R,
\]


\endgroup
% END UNIT Der-ChIX-678/Der-ChIX-678_p001-005.tex
% BEGIN UNIT Der-ChIX-678/Der-ChIX-678_p006-010.tex
\begingroup\makeatletter%
\NoAutoSpacing%
\newcommand{\printedpage}[1]{\marginpar{\footnotesize p.~#1}}%
\newcommand{\msp}[1]{\marginpar{\footnotesize [#1]}}%
\newcommand{\Hot}{\mathrm{Hot}}%
\newcommand{\HotO}{\mathrm{Hot}_0^{\wedge}}%
\newcommand{\MOP}{\mathcal{M}_{0\,\Pi}}%
\makeatother\let\printedpage\derpp\let\msp\dermsp%
\let\mspin\DERmspin%


\printedpage{4}
% head: A. Grothendieck : Dérivateurs, Ch. IX, Édition des Universités de Montpellier II et Paris VII
% foot: Version 12 mars 2002

on en déduit
\[
\left\{
\begin{array}{lcl}
\pi_0(P) & \xrightarrow{c_u^{\mathrm{l}}} & \mathrm{Hom}_{\mathcal{H}}(Q,R)\\
\pi_0(Q) & \xrightarrow{c_u^{\mathrm{r}}} & \mathrm{Hom}_{\mathcal{H}}(P,R)
\end{array}
\right.
\]
ainsi :
\[
\begin{array}{rlcl}
\text{Si } \alpha\in\pi_0(P),\ \text{i.e. } 1\xrightarrow{\ \alpha\ }P,\ \text{alors} & \multicolumn{3}{l}{\alpha*\mathrm{id}_Q : Q\longrightarrow P*Q}\\
\text{et} & c^{\mathrm{l}}_u(\alpha) & = & u\circ(\alpha*\mathrm{id}_Q)\\
\text{et dualement} & c^{\mathrm{r}}_u(\beta) & = & u\circ(\mathrm{id}_P*\beta).
\end{array}
\]
Si on a trois objets $P_0$, $P_1$, $P_2$ dans $\mathcal{H}$, des accouplements
\[
\left\{
\begin{array}{rcl}
P_0*P_1 & \xrightarrow{u_1} & Q_0\\
P_1*P_2 & \xrightarrow{u_2} & Q_1\\
P_0*Q_1 & \xrightarrow{u_3} & R\\
Q_0*P_2 & \xrightarrow{u_4} & R
\end{array}
\right.
\]
\[
\begin{tikzcd}
a \arrow[r, "P_0"'] \arrow[rr, bend left=20, "Q_0"] \arrow[rrr, bend left=35, "R"] & b \arrow[r, "P_1"'] \arrow[rr, bend right=20, "Q_1"'] & c \arrow[r, "P_2"'] & d
\end{tikzcd}
\]

\msp{page 6}
avec une relation d'associativité, disant que
\[
\begin{tikzcd}
P_0*P_1*P_2 \arrow[r] \arrow[d] & Q_0*P_2 \arrow[d]\\
P_0*Q_1 \arrow[r] & R
\end{tikzcd}
\]
est commutatif, alors on en déduit deux relations de transitivité et une de commutation, pour les opérations associées aux éléments de $\pi_0(P_0)$, $\pi_0(P_1)$, $\pi_0(P_2)$, qu'on n'explicitera pas, sauf l'une :
\[
\begin{array}{rrrcl}
\multicolumn{5}{l}{\text{Soit } \beta\in P_1,\ \gamma\in P_2,}\\
\text{d'où par } \beta : & c_{u_1}(\beta) : & P_0 & \longrightarrow & Q_0,\\
\text{et par } \gamma : & c_{u_4}(\gamma) : & Q_0 & \longrightarrow & R,\\
\text{et par } \beta *_{u_2}\gamma\in\pi_0(Q_1) : & c_{u_3}(\beta\,\texttt{[}{*_{u_2}}\texttt{]}\gamma) : & P_0 & \longrightarrow & R,
\end{array}
\]

\printedpage{5}
% head: A. Grothendieck : Dérivateurs, Ch. IX, Édition des Universités de Montpellier II et Paris VII
% foot: Version 12 mars 2002
\texttt{[plutôt} $\beta\in\pi_0(P_1)$\texttt{,} $\gamma\in\pi_0(P_2)$\texttt{]}, et bien on a
\[
\boxed{\;c(\beta *_{u_2}\gamma) \;=\; c(\gamma)c(\beta).\;}
\]
Il faut encore définir $\beta *_{u_2}\gamma$, via l'accouplement $P_1*P_2\xrightarrow{u_2}Q_1$, on a en effet
\[
\begin{array}{ccc}
\pi_0(P_1)\times\pi_0(P_2) & \longrightarrow & \pi_0(P_1*P_2)\\
(\beta\,,\,\gamma) & \longmapsto & \beta*\gamma,
\end{array}
\]
et on compose avec
\[
\pi_0(u_2) : P_1*P_2 \longrightarrow Q_1
\]
\texttt{[plutôt} $\pi_0(u_2) : \pi_0(P_1*P_2) \longrightarrow \pi_0(Q_1)$\texttt{]}.

\msp{page 7}
On se donne un ensemble $\mathcal{M}_0$, et
\[
\left\{
\begin{array}{rclccrcl}
\mathcal{M}_0\times\mathcal{M}_0 & \xrightarrow{\mathrm{Ch}} & \mathcal{H} & \text{noté} & & (a,b) & \longmapsto & \mathrm{Ch}(a,b)\ \text{ou}\ \mathrm{Hom\,ext}_{\mathcal{H}}(a,b)\\
\mathcal{M}_0\times\mathcal{M}_0\times\mathcal{M}_0 & \xrightarrow{\text{comp.}} & \mathrm{Fl}\,\mathcal{H} & & & (a,b,c) & \longmapsto & \bigl(\mathrm{Ch}(a,b)\times\mathrm{Ch}(b,c)\longrightarrow\mathrm{Ch}(a,c)\bigr)
\end{array}
\right.
\]
avec deux axiomes

\begin{enumerate}
\item[a)] Associativité pour les compositions, pour $(a,b,c,d)\in\mathcal{M}_0^4$
\[
\begin{tikzcd}
\mathrm{Ch}(a,b)\times\mathrm{Ch}(b,c)\times\mathrm{Ch}(c,d) \arrow[r] \arrow[d] & \mathrm{Ch}(a,c)\times\mathrm{Ch}(c,d) \arrow[d]\\
\mathrm{Ch}(a,b)\times\mathrm{Ch}(b,d) \arrow[r] & \mathrm{Ch}(a,d).
\end{tikzcd}
\]
\item[b)] Existence d'objets unité bilatères
\[
1_a\in\pi_0(\mathrm{Ch}(a,a)) \qquad \text{pour tout } a\in\mathcal{M}_0
\]
\end{enumerate}

Alors les $\pi_0(\mathrm{Ch}(a,b))$ et leur composition définissent sur $\mathcal{M}_0$ comme ensemble d'objets une structure de \textit{catégorie} $\MOP$, ayant les $1_a$ comme identités, et on trouve que
\[
(a,b)\longmapsto\mathrm{Ch}(a,b)
\]
peut s'interpréter à partir d'un \textit{bifoncteur} canonique
\[
\mathcal{M}^o_{0\,\Pi}\times\MOP \longrightarrow \mathcal{H},
\]
dont le composé avec
\[
\mathcal{H}\xrightarrow{\pi_0}\mathrm{Ens}
\]
n'est autre que le bifoncteur $(a,b)\longrightarrow\mathrm{Hom}_{\MOP}(a,b)$,
\printedpage{6}
% head: A. Grothendieck : Dérivateurs, Ch. IX, Édition des Universités de Montpellier II et Paris VII
% foot: Version 12 mars 2002
\msp{page 8}
\[
\mathrm{Hom}_{\MOP}(a,b) \;\simeq\; \pi_0(\mathrm{Ch}(a,b)).
\]
Cette relation est la \textit{définition} de $\underline{\mathrm{Hom}}_{\MOP}$, mais ici la relation est entendue comme une % Reading resolved against PDF leaf 18: a single underline; see corrections.csv.
bijection \textit{fonctorielle} en $a$, $b$.

On trouve notamment que si
\[
u : a'\longrightarrow a, \qquad v : b\longrightarrow b'
\]
sont des flèches \textit{inversibles} dans $\MOP$, alors
\[
\mathrm{Ch}(a,b)\longrightarrow\mathrm{Ch}(a',b')
\]
correspondant dans $\mathcal{H}$ est \texttt{[un]} \textit{isomorphisme.}

On appliquera tout ceci surtout aux deux situations déduites d'une situation $(\mathcal{M},\Sigma)$, avec
\[
\begin{array}{rlcllcl}
1^o) & \mathcal{H} & = & \mathrm{Hot}, & \mathrm{Ch}(a,b) & = & \mathrm{Ch}^{\Sigma}_{\infty}(\mathcal{M};a,b)\\
2^o) & \mathcal{H} & = & \mathrm{Ind\,Cat\,hot}, & \mathrm{Ch}(a,b) & = & \mathrm{CH}^{\Sigma}_{\infty}(\mathcal{M};a,b)\ \text{(\footnotemark[2])}.
\end{array}
\]
\footnotetext[2]{$*$ dans Hot et Ind\,Cat\,hot est le produit cartésien sans plus.}
Dans les deux cas, on a donc un foncteur sur $(\mathcal{M}\Sigma^{-1})^o\times(\mathcal{M}\Sigma^{-1})$, à valeurs soit dans Hot, soit dans Ind\,Cat\,hot. (\textbf{NB} Les deux structures de catégorie sur $\mathcal{M}_0$ déduites des % [sic] des -> de, « déduites des ces deux foncteurs »
ces deux foncteurs sont canoniquement isomorphes.)
\[
\begin{tikzcd}[column sep=small]
 & \mathrm{Hot} & (a,b)\longmapsto\mathrm{Ch}^{\Sigma}_{\infty}(a,b)\\
(\mathcal{M}\Sigma^{-1})^o\times(\mathcal{M}\Sigma^{-1}) \arrow[ur] \arrow[dr] & & \\
 & \mathrm{Ind\,Cat\,hot} & (a,b)\longmapsto\mathrm{CH}^{\Sigma}_{\infty}(a,b)
\end{tikzcd}
\]

\msp{page 9}
Comme les foncteurs canoniques
\[
\begin{tikzcd}
\mathrm{Ind\,Cat\,hot} \arrow[r, "\mathrm{Ind}\,\varphi"] & \mathrm{Ind\,Hot} \arrow[r, "\mathrm{Ind}\,\psi"] & \mathrm{Ind\,Hot}_0^{\wedge} \arrow[r, "\varinjlim"] & \mathrm{Hot}_0^{\wedge},\\
\mathrm{Cat\,hot} \arrow[u, hook] \arrow[r, "\varphi"] & \mathrm{Hot} \arrow[u, hook] \arrow[r, "\psi"] & \mathrm{Hot}_0^{\wedge} \arrow[u, hook] \arrow[ur, "\mathrm{id}"'] &
\end{tikzcd}
\]
commutant aux produits finis, toute théorie de `Hom externes' à valeurs dans une de ces catégories, en définit une à valeurs dans chacune de celles dans laquelle elle s'envoie.
\printedpage{7}
% head: A. Grothendieck : Dérivateurs, Ch. IX, Édition des Universités de Montpellier II et Paris VII
% foot: Version 12 mars 2002
D'ailleurs, ces foncteurs ne sont pas tous pleinement fidèles (c'est en tous cas faux pour $\varphi$ et $\mathrm{Ind}\,\varphi$, et j'ignore ce qui en est de $\psi$ et $\mathrm{Ind}\,\psi$ -- les flèches verticales par contre sont pleinement fidèles (\footnote[3]{et également, $\mathrm{Ind\,Hot}_0^{\wedge}\xrightarrow{\varinjlim}\mathrm{Hot}_0^{\wedge}$ n'est \textit{pas} pleinement fidèle.}).) Mais ils induisent des bijections pour les $\pi_0$. Donc la catégorie $\MOP$ déduite d'une théorie de Hom externes à valeurs dans une de ces catégoires, et celle % [sic] catégoires -> catégories, faute d'orthographe
déduite dans une autre où elle s'envoie, est la même à isomorphisme canonique près. Pour ce qui est de cette définition de $\MOP$, rien n'est changé donc quand on regarde les $\mathrm{CH}_{\infty}$ comme des objets de Ind\,Hot, ou seulement de $\mathrm{Ind\,Hot}_0^{\wedge}$, ou même seulement de $\HotO$, ou quand on regarde les $\mathrm{Ch}_{\infty}$ comme des objets
\msp{page 10}
dans $\HotO$. Mais bien sûr, la donnée des accouplements de composition perd chaque fois en précision (sauf peut-être en passant de Hot à $\HotO$, si $\psi$ est pleinement fidèle \dots\,). Quand on prend les images des $\mathrm{CH}^{\Sigma}_{\infty}(a,b)$ et $\mathrm{Ch}^{\Sigma}_{\infty}(a,b)$ dans $\HotO$ (la dernière catégorie du diagramme), on trouve des objets canoniquement isomorphes (fonctoriellement en $a$, $b$ dans $\MOP\simeq\mathcal{M}\Sigma^{-1}$). Quand on les regarde comme étant respectivement dans Ind\,Hot et dans $\mathrm{Hot}\hookrightarrow\mathrm{Ind\,Hot}$, on trouve un homomorphisme fonctoriel
\[
\mathrm{CH}^{\Sigma}_{\infty}(a,b)\longrightarrow\mathrm{Ch}^{\Sigma}_{\infty}(a,b) \qquad\quad \text{dans Ind\,Hot ou dans } \mathrm{Ind\,Hot}_0^{\wedge} \text{ au choix.} \tag*{$(*)$}
\]
La $\varinjlim$ de $\mathrm{CH}^{\Sigma}_{\infty}(a,b)$ dans Hot n'existe pas en général (sauf si l'ind-objet est essentiellement constant), mais elle existe dans $\HotO$, et l'isomorphisme canonique
\[
\varinjlim_{\HotO}\mathrm{CH}^{\Sigma}_{\infty}(a,b)\xrightarrow{\ \sim\ }\mathrm{Ch}^{\Sigma}_{\infty}(a,b) \qquad\quad \text{dans } \HotO \tag*{$(**)$}
\]
est la flèche déduite de $(*)$ dans $\mathrm{Ind\,Hot}_0^{\wedge}$.

\msp{page 11}
Dans les deux cas CH et Ch (qu'on va désigner par le symbole commun $C_{\infty}(a,b)$), la fonctorialité de cet objet $C_{\infty}(a,b)$ (même au sens le plus faible, comme objet de $\HotO$) par rapport à $a$ et $b$ (même au sens le plus faible, comme foncteur $\mathcal{M}^o\times\mathcal{M}\longrightarrow\HotO$ et pas même $(\mathcal{M}\Sigma^{-1})^o\times(\mathcal{M}\Sigma^{-1})\longrightarrow\HotO$) n'est pas évidente du tout sur la définition directe en termes du système inductif explicite dans Cat
\[
C_{\infty}(a,b) \;=\; (\mathrm{Ch}_{\tau_0^{2n}}(a,b))_{n\in\mathbf{N}},
\]
car une flèche $a'\longrightarrow a$ ou $b\longrightarrow b'$ dans $\mathcal{M}$, et même qu'elle serait dans $\Sigma$, ne définit aucunement une flèche dans Cat
\[
\begin{array}{rcl}
C_{\infty}(a,b) & \longrightarrow & C_{\infty}(a',b)\\
\text{ou}\ \ C_{\infty}(a,b) & \longrightarrow & C_{\infty}(a,b').
\end{array}
\]
Cette difficulté rappelle celle qu'avait à surmonter \textsc{Quillen}, dans Homotopical Algebra, pour affronter la différence entre `left homotopies' et `right homotopies'. Je me rappelle comment
\printedpage{8}
% head: A. Grothendieck : Dérivateurs, Ch. IX, Édition des Universités de Montpellier II et Paris VII
% foot: Version 12 mars 2002
\msp{page 12}
on peut néanmoins définir de telles flèches, comme résultant d'une flèche (beaucoup plus riche) de `composition' de chemins entre $a'$ et $a$, et $a$ et $b$, ou entre $a$ et $b$, et $b$ et $b'$, soit en général donc l'accouplement
\[
C_{\infty}(a,b)\times C_{\infty}(b,c)\longrightarrow C_{\infty}(a,c),
\]
dont la définition directe n'est pas moins problématique, et ne peut en tous cas \textit{pas} se donner comme une flèche dans Cat lui-même, ou dans Ind\,Cat. Interprétant p.~ex. $\mathrm{Ch}_{\infty}(a,b)$ en termes de chemins \textit{infinis} (des deux cotés) de $a$ à $b$, en un sens évident (mais % [sic] cotés -> côtés, accent omis
`essentiellement finis' des deux cotés), on peut définir aussi % [sic] cotés -> côtés, accent omis
\[
\mathrm{Ch}^{-}_{\infty}(a,b), \qquad \mathrm{Ch}^{+}_{\infty}(a,b)
\]
comme les sous-catégories pleines formées respectivement de chemins qui sont totalement dégénérés en degrés $\geq 0$ (i.e. $a_0\longleftarrow a_1\longrightarrow a_2\longleftarrow a_3\cdots$ sont tous des identités), ce qui revient aux chemins infinis seulement à gauche, i.e. s'arrêtant à $a_0$, et symétriquement pour $\mathrm{Ch}^{+}_{\infty}(a,b)$. Même tabac pour les $\mathrm{CH}^{+}$, $\mathrm{CH}^{-}$.

\msp{page 13}
Ceci dit, on a ceci (revenant à la notion $C_{\infty}$ ambivalente) :

\begin{enumerate}
\item[a)] les inclusions pleinement fidèles
\[
C^{-}_{\infty}(a,b)\hookrightarrow C_{\infty}(a,b)\hookleftarrow C^{+}_{\infty}(a,b)
\]
sont des isomorphismes dans Ind\,Cat\,hot resp. dans Hot, et
\item[b)] On a bel et bien une composition des chemins
\[
C^{-}_{\infty}(a,b)\times C^{+}_{\infty}(b,c)\longrightarrow C_{\infty}(a,c).
\]
\end{enumerate}

L'\textit{associativité} dans le cas du CH n'offre d'ailleurs pas de difficulté, à cause de la possibilité de se ramener aux
\[
C_m(a,b)\times C_n(b,c)\times C_p(c,d)\rightrightarrows C_N(a,d)
\]
pour $N$ grand, et prouver l'egalité comme flèches dans Cat\,hot -- ce qui résultera du lemme % [sic] egalité -> égalité, accent omis
d'homotopie sur les intervalles dans Cat. Dans le cas du $\mathrm{Ch}_{\infty}$ par contre je n'ai pas trouvé de démonstration. (Faute seulement d'avoir essayé sérieusement de m'y mettre? Je n'ai pas trop envie, vu qu'il se pourrait que $\mathrm{Hot}\longrightarrow\HotO$ soit pleinement fidèle, auquel cas le travail devient inutile \dots\,)


\endgroup
% END UNIT Der-ChIX-678/Der-ChIX-678_p006-010.tex
% BEGIN UNIT Der-ChIX-678/Der-ChIX-678_p011-015.tex
\begingroup\makeatletter%
\NoAutoSpacing%
\newcommand{\printedpage}[1]{\marginpar{\footnotesize p.~#1}}%
\newcommand{\msp}[1]{\marginpar{\footnotesize [#1]}}%
\newcommand{\M}{\mathcal{M}}%
\newcommand{\Hot}{\mathrm{Hot}}%
\newcommand{\Ch}{\mathrm{Ch}}%
\newcommand{\uCh}{\underline{\mathrm{Ch}}}%
\newcommand{\Homext}{\mathrm{Hom\,ext}}%
\newcommand{\Lotimes}{\overset{\mathbf{L}}{\otimes}}%
\makeatother\let\printedpage\derpp\let\msp\dermsp%
\let\mspin\DERmspin%


\printedpage{9}
% head: A. Grothendieck : Dérivateurs, Ch. IX, Édition des Universités de Montpellier II et Paris VII
% foot: Version 12 mars 2002

\msp{page 14}

\section*{2\quad Comparaison avec Quillen}

Il se pose un certain nombre de questions sur ce foncteur $(\M\Sigma^{-1})^o \times (\M\Sigma^{-1}) \longrightarrow \Hot$, et sa relation aux constructions classiques, notamment à la théorie de \textsc{Quillen}. Je vais maintenant noter
\[
\begin{aligned}
\Ch^{\Sigma}_{\infty}(a,b) \;&=\; \Homext_{\Hot}(a,b) && \text{ou simplement}\\
&\phantom{\;=\;} \Homext(a,b).
\end{aligned}
\]

1$^o$). Soient $f,g : a \rightrightarrows b$ dans $\M$, alors (si $\M$ est une `model category') \textsc{Quilllen} définit un ensemble % [sic] Quilllen -> Quillen, three l's printed
\[
\pi_1(a,b;f,g) \qquad \text{ou} \qquad \pi_1(f,g)
\]
des classes d'homotopie de $f$ à $g$, de telle façon que $\mathrm{Hom}_{\M}(a,b)$ devient l'ensemble des objets d'un groupoïde, dont les flèches sont les éléments de $\pi_1(f,g)$, pour une notion de composition naturelle. Je suspecte que l'on a
\[
\pi_1(f,g) \;=\; \mathrm{Hom}_{\underbrace{\Pi(\uCh^{\Sigma}_{\infty}(a,b))}_{\substack{\text{groupoïde}\\ \text{fodamental}\\ \text{de } \uCh^{\Sigma}_{\infty}(a,b)}}}(f,g), % [sic] fodamental -> fondamental, missing n
\]
quand $f$ et $g$ sont considérés comme des `chemins infinis' (en insérant chaque $f : a \longrightarrow b$ dans
\[
\cdots \underbrace{a_{-2}}_{=a} \xleftarrow{\ \mathrm{id}\ } \underbrace{a_{-1}}_{=a} \xrightarrow{\ f\ } \underbrace{a_{0}}_{=b} \xleftarrow{\ \mathrm{id}\ } \underbrace{a_{1}}_{=b} \longrightarrow \cdots\ \ )
\]

\msp{page 15}

Donc la définition des $\pi_1(f,g)$ (si c'est bien ça) se prolonge de façon naturelle aux chemins infinis (essentiellement finis) arbitraires.

2$^o$) Supposons que $\M$ admette un objet à la fois initial et final $0_{\M}$ ($\M$ `ponctuée'). Je dis qu'alors toutes les catégories $\uCh^{\Sigma}_{\infty}(a,b)$ sont ponctuées de façon naturelle : à savoir, par le chemin infini de $a$ à $b$ défini (comme plus haut) à partir de $f = 0_{a,b} : a \longrightarrow b$, l'application nulle de $a$ à $b$. On a envie alors de regarder les $\uCh^{\Sigma}_{\infty}(a,b)$ comme définissant des éléments dans le $\Hot_{\bullet}$ (`Hot ponctué'). C'est plus fin que de dire qu'on se donne une classe neutre
\[
e \;\longrightarrow\; \Ch(a,b) \qquad \text{dans Hot,}
\]
car le foncteur de $\Hot_{\bullet}$ dans Hot (ou dans $e\backslash\Hot$) n'est pas fidèle. Le problème, c'est de définir les flèches
\[
\Ch(a,b)\times\Ch(b,c) \;\longrightarrow\; \Ch(a,c)
\]
dans $\Hot_{\bullet}$, et pas seulement dans Hot, et avec associativité (et les unités) dans $\Hot_{\bullet}$ également. Si on applique la définition avec $\Ch^{-}(a,b)$ et $\Ch^{+}(b,c)$, le composé des éléments nuls $a \xrightarrow{\ 0\ } b$ de $\Ch^{-}(a,b)$ et $b \xleftarrow{\ \mathrm{id}\ } b \xrightarrow{\ 0\ } c$ de $\Ch^{+}(b,c)$
\printedpage{10}
% head: A. Grothendieck : Dérivateurs, Ch. IX, Édition des Universités de Montpellier II et Paris VII
% foot: Version 12 mars 2002
\msp{page 16}
est l'élément de $\Ch_{\infty}(a,c)$ que voici
\[
\begin{tikzcd}[row sep=huge]
 & a_{-2} & a_{-1} \arrow[l] \arrow[r] & a_{0} \arrow[d, equal] & a_{1} \arrow[l] \arrow[r] & a_{2} & \cdots \\
{}_{\mathrm{id}}\ \cdots & a & a \arrow[l, "\mathrm{id}"'] \arrow[r, "0"'] & b & b \arrow[l, "\mathrm{id}"'] \arrow[r, "0"'] & c & \cdots_{\mathrm{id}} \arrow[l, "\mathrm{id}"']
\end{tikzcd}
\]
ce n'est \textit{pas} l'élément nul de $\Ch_{\infty}(a,c)$ ! Mais on a des flèches canoniques
\[
e^{-}_{a,b}\circ e^{+}_{b,c} \;\;\rightleftarrows\;\; e_{a,c} \qquad \text{dans } \Ch_{\infty}(a,c)
\]
qui sont l'identité sur tous les sommets sauf celui d'indice 0. Il y a un sorite à développer, comme quoi une paire
\[
(f : K_{\bullet} \longrightarrow L_{\bullet}, c)
\]
d'un foncteur entre catégories ponctuées, et d'un chemin $c : f(e_{K_{\bullet}}) \longrightarrow e_{L_{\bullet}}$, définit dans Hot une flèche canonique
\[
\mathrm{hot}_{\bullet}(K_{\bullet}) \;\longrightarrow\; \mathrm{hot}_{\bullet}(L_{\bullet})
\]
avec transitivité si on a
\[
(f' : L_{\bullet} \longrightarrow M_{\bullet}, c' : f'(e_{L_{\bullet}}) \longrightarrow e_{M_{\bullet}}),
\]

\msp{page 17}

d'où
\[
f'' = f'\circ f : K_{\bullet} \longrightarrow M_{\bullet}, \qquad \text{et}\ \ c'' : f''(e_{K_{\bullet}}) \longrightarrow e_{M_{\bullet}}), % [sic] closing parenthesis without opening parenthesis -> (f'' = f'\circ f : ..., c'' : ...), unbalanced parenthesis
\]
définit ainsi
\[
\begin{tikzcd}
f'(f(e_{K_{\bullet}})) \arrow[r, "f'(c)"] \arrow[rr, bend right=12, "{c''\ \text{par définition}}"'] & f'(e_{L_{\bullet}}) \arrow[r, "c'"] & e_{M_{\bullet}}
\end{tikzcd}
\]

\noindent\llap{?\quad}Cela nous donne donc bel et bien des accouplements dans $\Hot_{\bullet}$ (et aussi dans $\mathrm{Cat\,hot}_{\bullet}$?). Il est difficile de croire qu'ils pourraient ne pas être associatifs! D'ailleurs $\Hot_{\bullet}$, avec le produit cartésien ordinaire (\footnote[4]{Non, il faut prendre le produit $\xi_{\bullet} \divideontimes \eta_{\bullet}$ ($\xi_{\bullet}\times\eta_{\bullet}$, avec $\xi_{\bullet}\vee\eta_{\bullet}$ contracté en $e$).}), satisfait à l'axiomatique de tantôt pour $\mathcal{H}$, et $\Hot_{\bullet} \longrightarrow \Hot$ est compatible avec ce produit (\footnote[5]{Non, mais on s'en tire quand même \ldots}), et avec $\pi_0$. % Reading resolved against PDF leaf 22: \divideontimes, meaning smash; see corrections.csv.

Si tout ceci marche, on aurait donc une théorie des $\Homext_{\bullet}$ à valeurs dans $\Hot_{\bullet}$; et sans doute aussi pour les $\mathrm{CH}_{\infty}$ comme objets de $\mathrm{Ind\,Hot}_{\bullet}$, et non seulement de $\mathrm{Ind\,Hot}$. Mais
\printedpage{11}
% head: A. Grothendieck : Dérivateurs, Ch. IX, Édition des Universités de Montpellier II et Paris VII
% foot: Version 12 mars 2002
de toutes façons on ne peut en général laisser tomber Cat hot, qui reste importante -- donc il faudrait avoir des objets $\mathrm{CH}_{\infty}(a,b)$ dans $\mathrm{Ind\,Cat\,hot}_{\bullet}$, pour avoir une côte bien taillée.

Les $\pi_0(C(a,b))$ seront à présent des ensembles ponctués, et ce sera par l'image de l'application nulle $a \xrightarrow{\ 0\ } b$, i.e. le composé
\[
a \longrightarrow 0 \longrightarrow b
\]

\msp{page 18}

qu'il faut composer avec l'elément déduit du chemin tronqué % [sic] l'elément -> l'élément, missing accent
\[
a \longrightarrow \underbrace{0}_{=a_0} \longleftarrow b \qquad \ldots
\]

Il faudrait montrer que 0 est aussi un objet neutre de $\M\Sigma^{-1}$, i.e. que
\[
\pi_0(a,0) = \pi_0(0,a) = \{0\} \qquad \text{pour tout } a \in \mathrm{Ob}\,\M
\]
et interpréter l'élément nul dans $\pi_0(a,b)$ comme l'application nulle dans $\M\Sigma^{-1}$. De façon générale, il faut des conditions qui assurent (sans que $\M$ soit forcément ponctuée) que
\[
\M \xrightarrow{\ \gamma\ } \M\Sigma^{-1}
\]
commute aux sommes finies, aux produits finis (donc $\gamma(\varphi_{\M}) = \varphi_{\M\Sigma^{-1}}$, resp. $\gamma(e_{\M}) = e_{\M\Sigma^{-1}}$) (\footnote[6]{à faire.}).

Supposons enfin que $\M$ soit \textit{additive,} et même $k$-linéaire ($k$ un anneau fixé), i.e. on a de plus un homomorphisme d'anneaux
\[
k \;\longrightarrow\; \mathrm{End}(\mathrm{id}_{\M}).
\]

On doit supposer quelque compatibilité de $\Sigma$ à la structure additive, p. ex. que la somme directe de deux flèches dans $\Sigma$ est dans $\Sigma$, et sans doute quelque chose de plus fort, du genre `lemme des cinq' pour $\Sigma$. (Mais on ne suppose pas forcément $\M$ abélienne, donc on aurait dû mal à formuler un lemme des cinq?) % [sic] dû mal -> du mal, accent on du

\msp{page 19}

Mais peut-être ici faut-il un axiome plus subtil, p.ex. (si $\underbrace{\Sigma}_{\text{suspension}}$ et $\Omega$ existent (\footnote[7]{cf. plus bas ce que ça signifie.})) que l'ensemble des morphismes d'adjonction $\Sigma\Omega \longrightarrow \mathrm{id}$, $\mathrm{id} \longrightarrow \Omega\Sigma$ soit \texttt{[formé d']}isomorphismes, ce qui implique que $\Omega$ resp. $\Sigma$ est pleinement fidèle. En tous cas, l'attente naturelle ici, c'est que $\Ch^{\Sigma}_{\infty}(a,b)$ soit non seulement dans $\Hot_{\bullet}$, mais même dans la catégorie $D^{\leq 0}(k\text{-Mod})$, sous-catégorie pleine de la catégorie dérivée $D^{-}(k\text{-Mod})$ formée des $K$ tels que $H^{i}(K) = 0$ si $i > 0$. On peut aussi la noter $D_{\geq 0}(k\text{-Mod})$, l'interprétant en
\printedpage{12}
% head: A. Grothendieck : Dérivateurs, Ch. IX, Édition des Universités de Montpellier II et Paris VII
% foot: Version 12 mars 2002
termes de complexes de `chaînes' plutôt que de cochaînes. Ici les accouplements doivent correspondre à l'opération $* = \Lotimes_k$ dans $D_{\geq 0}(k\text{-Mod})$, admettant $k$ comme objet unité, et
\[
\pi_0(K_{\bullet}) \;\simeq\; H_0(K_{\bullet}).
\]
Donc on doit avoir
\[
\Homext(a,b) \Lotimes_k \Homext(b,c) \;\longrightarrow\; \Homext(a,c),
\]
où cette fois Hom ext désigne l'objet de $D_{\geq 0}(k\text{-Mod})$, au lieu d'un objet de Hot seulement, ou de $\Hot_{\bullet}$. On utilise les foncteurs
\[
D_{\geq 0}(k\text{-Mod}) \longrightarrow \Hot_{\bullet} \longrightarrow \Hot
\]
\msp{page 20}
pour faire le lien avec la théorie ponctuée dans $\Hot_{\bullet}$, ou \texttt{[la]} théorie générale dans Hot. Si on désigne par
\[
\gamma :\ D_{\geq 0}(k\text{-Mod}) \;\longrightarrow\; \Hot_{\bullet}
\]
le premier de ces foncteurs, il ne commute pas aux produits $*$ ($\Lotimes_k$ et $\times$ respectivement), mais on a, pour $P_{\bullet}$, $Q_{\bullet}$ variable dans $D_{\geq 0}(k\text{-Mod})$, une flèche fonctorielle canonique % [sic] variable -> variables, agreement with P, Q
\[
\gamma(P_{\bullet} \Lotimes_k Q_{\bullet}) \xleftarrow{\ \mathit{can.}\ } \gamma(P_{\bullet})\times\gamma(Q_{\bullet})
\]
`$k$-bilinéaire', ce qui montre que de donner des accouplements
\[
P_{\bullet} * Q_{\bullet} \xrightarrow{\ c_k\ } R_{\bullet},
\]
d'où
\[
\gamma(P_{\bullet} * Q_{\bullet}) \xrightarrow{\ \gamma(c_k)\ } \gamma(R_{\bullet}), \tag*{$(*)$}
\]
est l'un et l'autre plus fin que de se donner seulement
\[
\gamma(P_{\bullet}) * \gamma(Q_{\bullet}) \xrightarrow{\ c_{\bullet}\ } \gamma(R_{\bullet}), \tag*{$(*')$}
\]
vu qu'une donnée $(*)$ impliquera une donne $(*')$ par % [sic] donne -> donnée, missing letters
\[
\begin{tikzcd}
 & \gamma(P_{\bullet} * Q_{\bullet}) \arrow[dr, "\gamma(c)"] & \\
\gamma(P_{\bullet}) * \gamma(Q_{\bullet}) \arrow[ur, "\text{can.}"] \arrow[rr, "c_{\bullet}"] & & \gamma(R_{\bullet}).
\end{tikzcd}
\]
Ceci vu, on s'attend bien à avoir une théorie \textit{forte} des accouplements
\printedpage{13}
% head: A. Grothendieck : Dérivateurs, Ch. IX, Édition des Universités de Montpellier II et Paris VII
% foot: Version 12 mars 2002
\msp{page 21}
\[
\Homext_k(a,b) \Lotimes_k \Homext_k(b,c) \;\longrightarrow\; \Homext_k(a,c),
\]
d'où on déduira les accouplements plus faibles dans $\Hot_{\bullet}$ (au lieu de $D_{\geq 0}(k\text{-Mod})$)
\[
\Homext_{\Hot_{\bullet}}(a,b) \times \Homext_{\Hot_{\bullet}}(b,c) \;\longrightarrow\; \Homext_{\Hot_{\bullet}}(a,c).
\]

3$^o$) Revenant au cas $\M$ ponctuée sans plus, \textsc{Quillen} définit, pour $a$, $b$ dans $\M$,
\[
\pi_1(a,b) \;=\; \pi_1(a,b\,;0_a,0_b)
\]
comme bifoncteur sur $(\M\Sigma^{-1})^o\times(\M\Sigma^{-1})$ (avec les notations de 1$^o$)). Si ce que j'avance sans preuve dans 1$^o$) est vrai, il doit s'ensuivre qu'on a un isomorphisme canonique
\[
\pi_1(a,b) \;\simeq\; \pi_1(\Ch^{\Sigma}_{\infty}(a,b)),
\]
où $\Ch^{\Sigma}_{\infty}(a,b)$ est vu comme objet de $\Hot_{\bullet}$, donc son $\pi_1$ est défini. On a en tous cas, sans aucune référence à des cofibrations et des fibrations, un bifoncteur à valeurs dans Groupes, composé dans
\[
(\M\Sigma^{-1})^o\times(\M\Sigma^{-1}) \xrightarrow{\ \ (a,b)\longmapsto \Ch^{\Sigma}_{\infty}(a,b)\ \ } \Hot_{\bullet} \xrightarrow{\ \pi_1\ } (\text{Groupes}).
\]
Si pour $a$ fixé, $b$ variable,
\[
b \;\longmapsto\; \pi_1(a,b)
\]
est représentable, on note l'objet qui le représente par $\Sigma(a)$ (\textit{suspension} -- ne pas confondre ce $\Sigma$ avec $\Sigma \subseteq \mathrm{Fl}(\M)$, qu'il vaut mieux maintenant noter $W$, comme `weak equivalence'). Donc

\msp{page 22}

\[
\mathrm{Hom}(\Sigma a,b) \;\simeq\; \pi_1(a,b),
\]
et $\Sigma a$ dépend fonctoriellement de $a$, sous réserve d'existence. S'il existe toujours, on a
\[
\Sigma :\ \M W^{-1} \;\longrightarrow\; \M W^{-1}.
\]
On définit de même, sous réserve d'existence,
\[
\Omega :\ \M W^{-1} \;\longrightarrow\; \M W^{-1}
\]
par
\[
\mathrm{Hom}(a,\Omega b) \;\simeq\; \pi_1(a,b),
\]
et si tous les deux existent, on a la formule d'adjonction
\[
\mathrm{Hom}(\Sigma a,b) \;\simeq\; \mathrm{Hom}(a,\Omega b) \qquad (\simeq \pi_1\Homext_{\Hot_{\bullet}}(a,b)).
\]


\endgroup
% END UNIT Der-ChIX-678/Der-ChIX-678_p011-015.tex
% BEGIN UNIT Der-ChIX-678/Der-ChIX-678_p016-018.tex
\begingroup\makeatletter%
\newcommand{\printedpage}[1]{\marginpar{\footnotesize p.~#1}}%
\newcommand{\msp}[1]{\marginpar{\footnotesize [#1]}}%
\newcommand{\cercle}[1]{\tikz[baseline=(c.base)]\node[draw,circle,inner sep=1pt](c){#1};}%
\newcommand{\Homext}{\mathrm{Hom}\,\mathrm{ext}}%
\newcommand{\Chu}{\underline{\mathrm{Ch}}}%
\newcommand{\Hot}{\mathrm{Hot}}%
\makeatother\let\printedpage\derpp\let\msp\dermsp%
\let\mspin\DERmspin%


\printedpage{14}
% head: A. Grothendieck : Dérivateurs, Ch. IX, Édition des Universités de Montpellier II et Paris VII
% foot: Version 12 mars 2002

L'existence de
\[
\begin{aligned}
\Omega :\ & \mathcal{M}W^{-1} \ \longrightarrow\ \mathcal{M}W^{-1}, \quad \text{ou de}\\
\Sigma :\ & \mathcal{M}W^{-1} \ \longrightarrow\ \mathcal{M}W^{-1},
\end{aligned}
\]
voire même des deux à la fois, apparaît comme une hypothèse \cercle{a} extrêmement forte sur $(\mathcal{M}, W)$. Cette condition \textit{plus} celle \cercle{b} que les $\Chu_{\infty}(a,b)$ soient homotopes à des catégories de Thomason-Kan, ou (ce qui revient au même, à peu de choses près?) que $Y \longmapsto \mathrm{Hom}_{\mathcal{M}(Y)W_Y^{-1}}(a_Y, b_Y)$ transforme Hot-équivalences de Cat en bijections (ne serait-ce que pour $Y$ fixé \dots), semble le test crucial pour savoir si `$(\mathcal{M}, \Sigma)$ \texttt{[plutôt} $(\mathcal{M}, W)$\texttt{]} se prête à une théorie homotopique `dans $\mathcal{M}$' '. Je doute que

\msp{page 23}

ces conditions puissent être satisfaites pour $\mathcal{M}$ finie, sauf dans le cas trivial où $\mathcal{M} = \emptyset$, où $\mathcal{M}\Sigma^{-1}$ \texttt{[plutôt} $\mathcal{M}W^{-1}$\texttt{]} est équivalente à la catégorie ponctuelle.

Un troisième test \cercle{c} qui me paraît important (mais peut-être ça résulte du reste), c'est les formules dans $\Hot_\bullet$
\[
\left\{
\begin{aligned}
\Homext_\bullet(\Sigma a, b) \ &\simeq\ \Homext_\bullet(a, \Omega b)\\
&\simeq\ \Omega\Homext_\bullet(a, b),
\end{aligned}
\right.
\tag*{$(*)$}
\]
où le $\Homext_\bullet$ désigne le Hom externe à valeurs dans $\Hot_\bullet$. Il faut le voir séparément (sous condition d'existence soit de $\Omega$, soit de $\Sigma$) comme deux formules `duales' l'une de l'autre,
\[
\left\{
\begin{aligned}
\Homext_\bullet(a, \Omega b) \ &\simeq\ \Omega\Homext_\bullet(a, b)\\
\Homext_\bullet(\Sigma a, b) \ &\simeq\ \Omega\Homext_\bullet(a, b),
\end{aligned}
\right.
\tag*{$(**)$}
\]
isomorphismes de \textit{groupes} dans $\Hot_\bullet$. (\textbf{NB} $\Omega b$ est un objet groupe de $\mathcal{M}W^{-1}$, car il représente un contrafoncteur à valeurs dans Gr, et de même $\Sigma a$ est un cogroupe de $\mathcal{M}W^{-1}$. Le foncteur $y \longmapsto \Homext(a, y)$ doit commuter aux produits finis, et donc transformer objets groupes en itou. Et de même $x \longmapsto \Homext(x, b)$ doit transformer sommes finies en produits finis, donc cogroupes en groupes (\footnote[8]{Le foncteur $\Omega$ doit provenir d'un endomorphisme du dérivateur $\mathbf{D} = \mathbf{D}_{\mathcal{M},W}$, et celui-ci doit être exact à gauche, et dualement pour $\Sigma$.}).)

Je me suis convaincu que les isomorphismes $(**)$ n'ont de chance d'exister que si on a la condition préalable \cercle{b}.

\par

\printedpage{15}
% head: A. Grothendieck : Dérivateurs, Ch. IX, Édition des Universités de Montpellier II et Paris VII
% foot: Version 12 mars 2002

\msp{page 24}

\section*{3\quad Essai de construction du foncteur $\Omega$}

Il n'y a sans doute guère de chance de prouver quelque chose sur le foncteur
\[
\Omega :\ \mathcal{M}W^{-1} \longrightarrow \mathcal{M}W^{-1},
\]
sans tenir en mains une \textit{construction} plus ou moins explicite de celui-ci. Je vais essayer, en m'inspirant de \textsc{Quillen}, de dégager des conditions générales sur $(\mathcal{M}, W)$ (plus générales que celles de \textsc{Quillen}) qui assurent l'existence de $\Omega$, et permettent de le construire. L'idée géométrique remonte à \textsc{Serre-Cartan}. Il s'agit, pour un objet $S$ de $\mathcal{M}$, de construire un objet $\mathrm{Ch}(S)$ \textit{dans $\mathcal{M}$,} jouant le rôle d'un `espace de chemins', et qui doit s'inclure dans un diagramme de factorisation de $\mathrm{diag}_S : S \longrightarrow S \times S$ (on suppose que dans $\mathcal{M}$ les $\varprojlim$ finies, donc les produits finis, existent)
\[
\begin{tikzcd}
S \arrow[d, "i_S"] \arrow[dd, bend right=40, "\mathrm{diag}_S"'] & \\
\mathrm{Ch}(S) \arrow[d, "p_S"] & \\
S \times S & e. \arrow[l]
\end{tikzcd}
\]
La propriété essentielle, c'est que $p$ soit une \textit{`fibration'} (dans un sens convenable, à dégager).

\msp{page 25}

À titre accessoire il peut être utile (peut-être) que $i$ soit une `cofibration' (\footnote[9]{et peut-être que $S \longrightarrow e$ soit également une fibration ?}). On définit $\Omega(S)$ comme la fibre de $\mathrm{Ch}(S)$ sur $S{\times}S$ (par rapport à l'unique morphisme $e = e_{\mathcal{M}} \longrightarrow S\times S$, où $e_{\mathcal{M}}$ est l'objet nul de $\mathcal{M}$).
\[
\begin{tikzcd}[column sep=large]
\mathrm{Ch}(S) \arrow[d] & \Omega S \arrow[l] \arrow[d]\\
S \times S & e. \arrow[l]
\end{tikzcd}
\]

\par

\printedpage{16}
% head: A. Grothendieck : Dérivateurs, Ch. IX, Édition des Universités de Montpellier II et Paris VII
% foot: Version 12 mars 2002

C'est un objet de $\mathcal{M}$, qu'on regardera surtout comme un objet de $\mathcal{M}W^{-1}$. Il faut s'assurer qu'on a une bijection fonctorielle en $T$ dans $\mathcal{M}W^{-1}$
\[
\underbrace{\pi_1\Chu^W_\infty(T, S)}_{\overset{\text{déf}}{=}\,\pi_1(T,S)} \ \simeq\ \underbrace{\mathrm{Hom}_{\mathcal{M}W^{-1}}(T, \Omega S)}_{\overset{\text{déf}}{=}\,\pi_0\Chu^W_\infty(T,\Omega S)}.
\]
Pour cela, il s'impose de comparer
\[
\Chu^W_\infty(T, \Omega S) \qquad\quad \text{et} \qquad\quad \underbrace{\Omega_{\Hot_\bullet}(\Chu^W_\infty(T, S))}_{\overset{\substack{\text{`par définition'}\\ \text{pour $\Omega$ de $\mathrm{Hot}_\bullet$}}}{=}\ \Chu_\infty(\mathrm{Ch}^W_\infty(T,S);0,0)} \qquad ,
\]
puisqu'il s'agit de construire une bijection entre leurs $\pi_0$.


\endgroup
% END UNIT Der-ChIX-678/Der-ChIX-678_p016-018.tex
\clearpage\setcounter{footnote}{0}\phantomsection\addcontentsline{toc}{section}{Chapitre X}\markboth{Chapitre X}{Chapitre X}
% BEGIN UNIT Der-ChX-789/Der-ChX-789_p001-005.tex
\begingroup\makeatletter%
\newcommand{\printedpage}[1]{\marginpar{\footnotesize p.~#1}}%
\newcommand{\msp}[1]{\marginpar{\footnotesize [#1]}}%
\newcommand{\D}{\mathbf{D}}%
\newcommand{\Cat}{\mathrm{Cat}}%
\newcommand{\Fl}{\operatorname{Fl}}%
\newcommand{\Ob}{\operatorname{Ob}}%
\newcommand{\bDelta}{\boldsymbol{\Delta}}%
\makeatother\let\printedpage\derpp\let\msp\dermsp%
\let\mspin\DERmspin%


% ===================== page 1 (p001) =====================
\printedpage{}
% no page number printed (title page)
% head: (none)

\begin{center}
{\LARGE \textsf{LES DÉRIVATEURS}}

\bigskip

{\large \textbf{ALEXANDRE GROTHENDIECK}}

\vspace{3em}

{\large \textbf{Chapitre X}}

\medskip

{\large \textbf{Comparaison de} Cat \textbf{avec} $\bDelta^{\wedge}$}
\end{center}

\par

% ===================== page 2 (p002) =====================
\printedpage{}
% no page number printed (editorial note)
% head: (none)

\par\noindent\textbf{Note des éditeurs de l’édition transcrite.}\par
Ce texte a été déchiffré et transcrit en \LaTeXe{} par M. Künzer. Il a été édité par M. Künzer, J. Malgoire, et G. Maltsiniotis. La transcription est aussi fidèle que possible au manuscrit. Pour les quelques corrections évidentes, ou rares commentaires des éditeurs, ainsi que pour la numérotation originale des pages du manuscrit, les caractères de machine à écrire \texttt{[typewriter]} entre crochets sont utilisés. Un point d’interrogation entre crochets signifie que l’on n’est pas sûr du mot qui précède.

Cette édition est provisoire. Les remarques, commentaires et corrections sont bienvenus. Envoyer un message à :

\textbf{[e-mail address removed]}

\hfill G. Maltsiniotis\hspace*{3em}

\par

% ===================== page 3 (p003), printed 1 =====================
\printedpage{1}
% head: A. Grothendieck : Dérivateurs, Ch. X, Édition des Universités de Montpellier II et Paris VII
% foot: Version 9 avril 2002

\msp{page 1}

\section*{1\quad Rappel des notions de D-asphéricité, et lien avec le ‘localisateur fondamental $W_{\mathbf{D}}$’}

Pour un dérivateur $\D$, il s’agit des notions suivantes (cf. I) :

\begin{itemize}
\item[(i)] \textit{Objet} $\D$-\textit{asphérique} de Cat (i.e. $Z \longrightarrow e$ est une $\D$-équivalence). NB Je préfère dire ‘$\D$-asphérique’ au lieu de ‘$\D$-contractile’, qui prête parfois à confusion.

\item[\textcircled{ii}] $\D$-équivalence pour une flèche $f : X \longrightarrow Y$ dans Cat.

\item[(iii)] \textit{Foncteur} $f : X \longrightarrow Y$ $\D$-\textit{asphérique} (i.e. les $X/y$ sont $\D$-asphériques) et $\D$-\textit{coasphérique} (les $y\backslash X$ sont $\D$-asphériques). C’est ce que j’avais appelé ‘$\D$-équivalence stricte’ (ou ‘costricte’), ou foncteur homologiquement $\D$-cofinal (ou cohomologiquement $\D$-cofinal), mais la terminologie ‘asphérique’ me semble la plus simple et la meilleure. (Pêchée dans PS, Pursuing Stacks.)

\item[(iv)] \textit{Foncteur} $f : X \longrightarrow Y$ $\D$-\textit{lisse,} $\D$-\textit{propre} (i.e. tel que les fibres des foncteurs induits $X/x \longrightarrow Y/f(x)$, respectivement $x\backslash X \longrightarrow f(x)\backslash Y$, soient $\D$-asphériques).

\item[(v)] $\D$-\textit{fibration} $f : X \longrightarrow Y$ (notion autoduale), i.e. telle que pour $Y'' \xrightarrow{u} Y'$ au dessus de $Y$, si $Y''$, $Y'$ \texttt{[sont]} Hot-asphériques, alors $u_X : X'' \longrightarrow X'$ dans
\[
\begin{tikzcd}
X \arrow[d] & X' \arrow[l] \arrow[d] & X'' \arrow[l] \arrow[d] \\
Y & Y' \arrow[l] & Y'' \arrow[l]
\end{tikzcd}
\]
est une $\D$-équivalence (\footnote{On dira aussi que $f$ est $\D$-\textit{fibrant.} (\textbf{NB} Il n’y a pas de notion ‘$\D$-cofibrant’ à ma connaissance.)}). La $\D$-\textit{fibration} est dite \textit{triviale} (ou \textit{asphérique}) si et seulement si ses fibres sont asphériques, ou encore si et seulement si $f$ est asphérique, ou si et seulement si $f$ est coaspérique. Cela signifie aussi que $f$ est une $\D$-équivalence \textit{universelle,} % [sic] coaspérique -> coasphérique, letter h missing
\msp{page 2}
i.e. que pour tout $Y' \longrightarrow Y$, $X' = X \times_Y Y' \longrightarrow Y'$ est une $\D$-équivalence.

\item[(vi)] On a aussi des variantes affaiblies de la notion de $\D$-fibrations, savoir les ‘$\D$-\textit{fibrations à gauche’,} ou les $f$ ‘\textit{cohomologiquement} $\D$-\textit{fibrants’,} et de même les ‘$\D$-\textit{fibrations à droite’,} i.e. les $f$ ‘\textit{homologiquement} $\D$-\textit{fibrants’} (\footnote{Nouvelle terminologie : fibrations cohomologiques (resp. homologiques) \textit{faibles,} cf. I, pages 111-126 pour une révision corrigée des diverses notions de ‘$\D$-fibrations’.}). La notion ‘gauche’
\printedpage{2}
% head: A. Grothendieck : Dérivateurs, Ch. X, Édition des Universités de Montpellier II et Paris VII
% foot: Version 9 avril 2002
ou cohomologique pour $f : X \longrightarrow Y$ signifie que les foncteurs $X_y \hookrightarrow X/y$ ($y \in \Ob Y$) et $X/y \longrightarrow X/y'$ (pour $(u : y \longrightarrow y') \in \Fl Y$) sont des $\D$-équivalences. Et dualement avec les $X_y \hookrightarrow y\backslash X$, $y'\backslash X \longrightarrow y\backslash X$ pour la notion ‘droite’ ou ‘homologique’. Les $\D$-fibrations ne sont autres que les $\D$-fibrations gauches (cohomologiques) universelles ou les $\D$-fibrations droites (homologiques) universelles.
\end{itemize}

Finalement, toute cette kyrielle de notions se déduit de la seule notion de $\D$-équivalence (ii). Si je désigne par
\[
W_{\D} \subseteq \Fl(\Cat)
\]
l’ensemble de ces $\D$-équivalences, on trouve ce que dans Pursuing Stacks j’ai appelé un ‘localisateur’ (localiseur) dans Cat (\footnote{Mais attention, je suppose ici que le domaine Diag de $\D$ est Cat tout entier!}). Dans Pursuing Stacks c’était une telle partie
\[
W \subseteq \Fl(\Cat),
\]
et non la donnée d’un dérivateur, qui
\msp{page 3}
était pris comme la donnée de départ fondamental pour faire de l’algèbre homotopique en tous genres. J’avais d’ailleurs développé une kyrielle de notions du même genre (i) à (vi), en partant de cette donnée $W$. Il se pourrait que les énoncés pour les notions que j’ai développé dans I, en termes d’un dérivateur, soient valables plus généralement en termes d’un localisateur $W$, dans la mesure du moins où ils gardent un sens. (Donc pour l’instant je mets à part les énoncés proprement cohomologiques ou homologiques, qui font directement appel à $\D$.) Encore qu’il se pourrait que moyennant des conditions à expliciter sur $W$, qui impliquent que $(\Cat, W)$ définissent bel et bien un dérivateur $\D_{\Cat,W} = \D$, on arrive encore à interpréter $W$ comme un $W_{\D}$, de sorte que les énoncés cohomologiques et homologiques gardent un sens.

J’ai envie d’abord de regarder dans Pursuing Stacks les axiomes que j’ai dégagé pour les localisateurs, et vérifier s’ils sont tous vérifiés pour
\msp{page 4}
$W_{\D}$. Ces conditions L 1 à L 5 sont énoncées dans Pursuing Stacks, §64 (page 164).

\textbf{L 1}\enspace Saturation de $W_D$. O.K.

Cela résulte formellement du fait qu’on a un foncteur
\[
\Cat \xrightarrow{\varphi} C \qquad (= \Cat^o)
\]
tel que pour $u \in \Fl(\Cat)$, $u \in W_{\D} \iff \varphi(u)$ isomorphisme. (Prendre pour $\varphi(X)$ la sous-catégorie de $\D(X)$ formée des $\D$-coefficients \textit{constants} sur $X$).

\textbf{L 2}\enspace “$\bDelta^1$ est $W$-asphérique sur $e$”, i.e. pour tout $X$ dans Cat, $X \times \bDelta^1 \longrightarrow X$ est dans $W$.

Cela signifie aussi que $\bDelta^1 \longrightarrow e$ est une $\D$-fibration triviale, ce qui est bien le cas.

\par

% ===================== page 5 (p005), printed 3 =====================
\printedpage{3}
% head: A. Grothendieck : Dérivateurs, Ch. X, Édition des Universités de Montpellier II et Paris VII
% foot: Version 9 avril 2002

\textbf{L 3}\enspace Si $X \in \Ob\Cat$ a un objet final, alors $X \longrightarrow e$ est dans $W$.

En fait, nous exigeons pour les dérivateurs que ceci soit vrai en supposant seulement $X$ asphérique (au sens absolu), et on en conclut même que $X \longrightarrow e$ est $W$-asphérique, i.e. $X \times Y \xrightarrow{\mathrm{pr}_2} Y$ est dans $W$ pour tout $Y$ dans Cat (\footnote{En fait L 2 $\Longrightarrow$ L 3.}).

\textbf{L 4}\enspace Si $f : X \longrightarrow Y$ est tel que $X/y \longrightarrow Y/y$ sont dans $W$, alors $f$ est dans $W$.

C’est OK.

\msp{page 5}
\textbf{L 5}\enspace Si $X$, $Y$ sont Cat-fibrées sur $S$, et $X \longrightarrow Y$ est un $S$-morphisme \textit{cartésien,} induisant sur les fibres des $X_s \longrightarrow Y_s$ qui sont dans $W$, alors $f$ est dans $W$.

Pour $W_{\D}$, c’est la proposition 6, page 71 de I.

Ces axiomes L 1 à L 5 (équivalent à L 1, L 3$'$ ($\bDelta^1 \longrightarrow e$ est dans $W$) et L 5) sont dans la nature d’axiomes de stabilité, assurant que $W$ est assez grand. Il y a une autre condition qui limite $W$ à ne pas être trop grand :

\textbf{L a)}\enspace Si $f \in W$, alors $\pi_0(f)$ est bijectif.

Ceci n’est pas satisfait pour tous les $W_{\D}$ – p. ex. si $\D$ est le dérivateur \textit{nul,} alors $W_{\D} = \Fl\Cat$! Mais pour les dérivateurs utiles en pratique, cela semble toujours satisfait.

\textbf{L b), c)}\enspace Le foncteur
\[
\Cat \longrightarrow \Cat\, W^{-1}
\]
commute aux sommes finies et aux produits finis (\footnote{Ça marche si $(\Cat, W)$ provient d’un triple de Quillen (avec des Q-fibrations, Q-cofibrations, Q-équivalences faibles).}).

À examiner ce que ça signifie en termes concrets sur $W$ – je n’ai pas vérifié que c’est O.K. pour les $W_{\D}$.

\textbf{L d)}\enspace Si $f : X \longrightarrow Y$ dans $W$, alors $f^o : X^o \longrightarrow Y^o$ aussi.

C’est le corollaire 2, page 75 de I.


\endgroup
% END UNIT Der-ChX-789/Der-ChX-789_p001-005.tex
\clearpage\setcounter{footnote}{0}\phantomsection\addcontentsline{toc}{section}{Chapitre XI}\markboth{Chapitre XI}{Chapitre XI}
% BEGIN UNIT Der-ChXI-890/Der-ChXI-890_p001-005.tex
\begingroup\makeatletter%
\newcommand{\printedpage}[1]{\marginpar{\footnotesize p.~#1}}%
\newcommand{\msp}[1]{\marginpar{\footnotesize [#1]}}%
\newcommand{\Hot}{\mathrm{Hot}}%
\newcommand{\Cat}{\mathrm{Cat}}%
\newcommand{\Ob}{\operatorname{Ob}}%
\makeatother\let\printedpage\derpp\let\msp\dermsp%
\let\mspin\DERmspin%


% ===================== page 1 =====================
\printedpage{--}
% head: (none) -- title page, no page number printed

\begin{center}
\vspace*{5cm}
{\LARGE\textsf{LES DÉRIVATEURS}}

\vspace{1.5em}
{\large\textbf{ALEXANDRE GROTHENDIECK}}

\vspace{3em}
{\large\textbf{Chapitre XI}}

\vspace{0.5em}
{\large $\Hot$\textbf{-fibrations, foncteurs propres et foncteurs lisses etc.}\\
\textbf{(dans} $\Cat$\textbf{)}}
\end{center}

\par

% ===================== page 2 =====================
\printedpage{--}
% head: (none) -- no page number printed

\vspace*{6cm}
\noindent \par\noindent\textbf{Note des éditeurs de l’édition transcrite.}\par
Ce texte a été déchiffré et transcrit en \LaTeXe\ par M. Künzer. Il a été édité par M. Künzer, J. Malgoire, et G. Maltsiniotis. La transcription est aussi fidèle que possible au manuscrit. Pour les quelques corrections évidentes, ou rares commentaires des éditeurs, ainsi que pour la numérotation originale des pages du manuscrit, les caractères de machine à écrire \texttt{[typewriter]} entre crochets sont utilisés. Un point d'interrogation entre crochets signifie que l'on n'est pas sûr du mot qui précède.

\noindent Cette édition est provisoire. Les remarques, commentaires et corrections sont bienvenus. Envoyer un message à :

\noindent\textbf{[e-mail address removed]}

\hfill G. Maltsiniotis\hspace*{2.5em}

\par

% ===================== page 3 =====================
\printedpage{1}
% head: A. Grothendieck : Dérivateurs, Ch. XI, Édition des Universités de Montpellier II et Paris VII
% foot: Version 28 mai 2002

\noindent\texttt{[Il n'y a pas de pages 1-126]}

\msp{page 127}

\section*{1\quad $\Hot$-fibrations dans $\Cat$}

\noindent\textbf{Théorème 1.} \textit{Soit $f : X \longrightarrow Y$ dans} $\Cat$\textit{. Conditions équivalentes :}

\begin{itemize}[leftmargin=3em]
\item[\textit{a)}] \textit{Pour tout $Y'' \xrightarrow{p} Y'$ sur $Y$, d'où $X'' = X \times_Y Y'' \xrightarrow{q} X' = X \times_Y Y'$, si $p$ est une} $\Hot$\textit{-équivalence, $q$ l'est aussi.}

\item[\textit{a$'$)}] \textit{Comme a), avec $Y'' = e$, $Y'$ ayant un objet final.}

\item[\textit{a$''$)}] \textit{Comme a), avec $Y'' = e$, $Y'$ ayant un objet initial} (\footnote{garder a$''$) et b$''$), conditions duales de a$'$), b$'$) \texttt{[a$''$) et b$''$) étaient barrés]}.})\textit{.}

\item[\textit{b$'$)}] \textit{Pour tout $y \in \Ob Y$, $X_y \longrightarrow X/y$ est une} $\Hot$\textit{-équivalence, pour tout $y \longrightarrow y'$ dans $Y$, $X/y \longrightarrow X/y'$ est une} $\Hot$\textit{-équivalence; enfin ces conditions sont encore vérifiées pour $f' : X' \longrightarrow Y'$, pour tout changement de base $Y' \longrightarrow Y$ dans} $\Cat$\textit{. (En fait, il suffit de le demander pour la première des deux conditions, concernant $X_y \longrightarrow X/y$.)}

\item[\textit{b$''$)}] \textit{Dual de b$'$), avec $X_y \longrightarrow y\backslash X$ et $y'\backslash X \longrightarrow y\backslash X$.}

\item[\textit{c)}] \textit{Pour tout faisceau d'ensembles localement constant $F$ sur $X$, $f_*(F)$ est localement constant et se calcule `fibre par fibre', de même pour $R^1 f_*(F)$ si $F$ est un faisceau en groupes, et pour les $R^i f_*(F)$ si $F$ est un faisceau abélien. De plus, on veut que ces conditions restent satisfaites après tout changement de base.} (\footnote{\textit{c$'$) (version homologique). 1$^o$) en termes d'images directes de groupoïdes localement constants, et 2$^o$) en termes des $R^i f_!(L)$} \texttt{[plutôt $L_i f_!(L)$]}\textit{, pour $L$ faisceau abélien localement constant.}})
\end{itemize}

\noindent\textsc{Démonstration}. On a les implications $\pm$ tautologiques
\[
\mathrm{a} \implies \mathrm{a}' \implies \mathrm{b}' \implies \mathrm{c},
\]
et il reste à prouver que c $\implies$ a. Comme a est autoduale, étant équivalente à a$'$, b$'$, elle est aussi équivalente aux conditions duales de celles-ci, a$''$ et b$''$ (cf. cor. 2, page 75). On est ramené à prouver que si c) est satisfait, et si on a
\[
\begin{tikzcd}[row sep=large, column sep=large]
X \arrow[d, "f"'] & X' \arrow[l, "q"'] \arrow[d, "f'"] \\
Y & Y' \arrow[l, "p"]
\end{tikzcd}
\]
\printedpage{2}
% head: A. Grothendieck : Dérivateurs, Ch. XI, Édition des Universités de Montpellier II et Paris VII
% foot: Version 28 mai 2002
cartésien dans $\Cat$, si $p$ est une $\Hot$-équivalence, $q$ aussi. On utilise le critère d'Artin-Mazur.

% ===================== page 4 (printed 2) =====================

\noindent 1$^o$) Si $F$ est un faisceau d'ensembles localement constant sur $X$, il faut prouver que $H^0(X,F) \xrightarrow{\sim} H^0(X', q^*F)$.

\msp{page 128}

\noindent Mais
\[
H^0(X,F) \simeq H^0(Y,\ \underbrace{f_*(F)}_{\substack{\text{localement constant}\\ \text{par hypothèse c)}}}\ )
\qquad
\underbrace{\xrightarrow{\qquad\sim\qquad}}_{\substack{\text{car } p\\ \text{Hot-équivalence}}}
\qquad
H^0(Y', p^* f_*(F)),
\]
mais
\[
p^* f_*(F) \xrightarrow{\ \sim\ } f'_*(q^*(F)),
\]
car $f_*(F)$ et $f'_*(q^*F)$ se calculent `fibre par fibre' par hypothèse. Donc
\[
H^0(X,F) \simeq H^0(Y', f'_*(q^*(F))) \simeq H^0(X', q^*(F)),
\]

\hfill q.e.d.

\noindent 2$^o$) Si $F$ est un faiscau en groupes localement constant sur $X$, il faut prouver % [sic] faiscau -> faisceau, misspelling
\[
H^1(X,F) \xrightarrow{\ \sim\ } H^1(X', q^*F).
\]
Le fait que c'est injectif résulte déjà de 1$^o$), qui implique que le foncteur
\[
\begin{array}{rcl}
F\text{-torseurs sur } X & \longrightarrow & q^*(F)\text{-torseurs sur } X' \\
P & \longmapsto & f^*P % [sic] f^*P -> q^*P, the functor is the inverse image by q from X to X'
\end{array}
\]
est pleinement fidèle, par l'argument bien connu utilisant les faisceaux $\mathcal{H}$om. Il faut voir que ce foncteur est essentiellement surjectif.

\msp{page 129}

\noindent On a un homomorphisme de suites exactes d'ensembles ponctués
\[
\begin{tikzcd}[column sep=large, row sep=large,
  every matrix/.append style={name=m},
  execute at end picture={
    \coordinate (lb) at ([xshift=0.4em]m-2-3.east);
    \coordinate (lt) at ([xshift=1.8em]m-2-3.east);
    \draw (lt |- m-1-3.north) -- (lb |- m-1-3.north) -- (lb |- m-2-3.south) -- (lt |- m-2-3.south);
    \coordinate (rb) at ([xshift=0.6em]m-2-4.east);
    \coordinate (rt) at ([xshift=-0.6em]m-2-4.east);
    \draw (rt |- m-1-3.north) -- (rb |- m-1-3.north) -- (rb |- m-2-3.south) -- (rt |- m-2-3.south);
  }]
H^1(Y, f_*F) \arrow[r, "u"] \arrow[d, "\alpha"] & H^1(X,F) \arrow[r, "v"] \arrow[d, "\beta"] & \overbrace{H^0(Y, R^1 f_*(F))}^{\ni\, \eta} \arrow[r, "w"] \arrow[d, "\gamma"] & H^2(Y, f_*F) \arrow[d, "\delta"] \\
H^1(Y', f'_*F) \arrow[r, "u'"] & \underbrace{H^1(X', F')}_{\ni\, \xi'} \arrow[r, "v'"] & \underbrace{H^0(Y', R^1 f'_*(F'))}_{\ni\, \eta'} \arrow[r, "w'"] & H^2(Y', f'_*F') % [sic] H^1(Y', f'_*F) -> H^1(Y', f'_*F'), the other entries of the lower row are in F'
\end{tikzcd}
\]
(\footnote{où on a posé $F' = q^*(F)$.}), où les faisceaux $f_*F$, $R^1 f_*(F)$ sont localement constants par hypothèse et se calculent fibre par fibre, les faisceaux $f'_*(F')$ et $R^1 f'_*(F')$ aussi, qui implique que ces derniers sont
\printedpage{3}
% head: A. Grothendieck : Dérivateurs, Ch. XI, Édition des Universités de Montpellier II et Paris VII
% foot: Version 28 mai 2002
canoniquement isomorphes (par les flèches de changement de base) aux images inverses de leurs homologues $f_*(F)$, $R^1 f_*(F)$. Comme $p$ est une $\Hot$-équivalence, cela implique que les flèches $\alpha$, $\gamma$ sont \texttt{[des]} isomorphismes. Donc si $\xi' \in H^1(X', F')$, son image $\eta'$ dans $H^0(Y', R^1 f'_*(F'))$ provient d'un $\eta \in H^0(Y, R^1 f_*(F))$ bien déterminé. Il faut voir que cet $\eta$ provient d'un $\xi \in H^1(X,F)$. Ceci (si je me rappelle bien) est possible si et seulement si une certaine obstruction $w(\eta)$ dans $H^2(Y, f_*(F))$ est dans le sous-ensemble neutre du $H^2$. \textbf{NB} Le champ sur $Y^\wedge$ des $F$-torseurs sur les $X^\wedge/U$ qui remontent $\eta$ est une $\mathcal{F} = f_*F$-gerbe $\mathcal{G}_\eta$, et il faut voir qu'elle est neutre. Mais il y a mieux à faire :

% ===================== page 5 (printed 3) =====================

\msp{page 130}

\noindent On note que cette gerbe $\mathcal{G}$ est localement constante (puisqu'elle correspond à un lien $f_*(G)$ localement constant) et que $\mathcal{G}'$ est son image inverse. Mais comme $p$ est une $\Hot$-équivalence, il s'ensuit que le foncteur
\[
\underline{\Gamma}(Y, \mathcal{G}) \longrightarrow \underline{\Gamma}(Y', \mathcal{G}')
\]
est une équivalence de catégories, donc $\xi'$ provient (à isomorphisme près) d'un unique $\xi$ de $\underline{\Gamma}(Y, \mathcal{G})$, q.e.d.

\noindent 3$^o$) Soit $F$ faisceau abélien localement constant sur $X$, il faut voir \texttt{[que]}
\[
H^i(X,F) \xrightarrow{\ \sim\ } H^i(X', F') \qquad (F' = q^*F), \tag*{$(*)$}
\]
pour tout $i \in \mathbf{N}$. Mais les deux termes sont l'aboutissement des suites spectrales de Leray, de termes initiaux
\[
\begin{aligned}
E_2^{pq} &\simeq H^p(Y, R^q f_*(F)) \\
E^{\prime pq}_2 &\simeq H^p(Y', R^q f'_*(F')),
\end{aligned}
\]
et $(*)$ est associé à un homomorphisme de suites spectrales, qui pour les termes initiaux est
\[
H^p(Y, R^q f_*(F)) \xrightarrow{\sim} H^p(Y', R^q f'_*(F)), \tag*{$(**)$} % [sic] R^q f'_*(F) -> R^q f'_*(F'), cf. the E'_2 term just above
\]
associé à
\[
p^*(R^q f_*(F)) \longrightarrow R^q f'_*(F'). \tag*{$(*{*}*)$}
\]
Mais par hypothèse les $R^q f_*(F)$, $R^q f'_*(F')$ sont localement constants et se calculent fibre par fibre,

\msp{page 131}

\noindent d'où il résulte que les homomorphismes de changement de base $(*{*}*)$ sont \texttt{[des]} isomorphismes, et comme il s'agit de faisceaux localement constants et que $p$ est une $\Hot$-équivalence, il s'ensuit que $(**)$ est \texttt{[un]} isomorphisme, donc aussi $(*)$ sur les aboutissements, q.e.d.

\noindent\textbf{Définition 1.} Sous les conditions équivalentes du théorème 1, on dit que $f$ est une $\Hot$-\textit{fibration.}


\endgroup
% END UNIT Der-ChXI-890/Der-ChXI-890_p001-005.tex
% BEGIN UNIT Der-ChXI-890/Der-ChXI-890_p006-008.tex
\begingroup\makeatletter%
\derleqno%
\newcommand{\printedpage}[1]{\marginpar{\footnotesize p.~#1}}%
\newcommand{\msp}[1]{\marginpar{\footnotesize [#1]}}%
\newcommand{\Hot}{\mathrm{Hot}}%
\newcommand{\Ch}{\underline{\mathrm{Ch}}}%
\newcommand{\Ob}{\mathrm{Ob}\,}%
\newcommand{\Cat}{\mathrm{Cat}}%
\makeatother\let\printedpage\derpp\let\msp\dermsp%
\let\mspin\DERmspin%


\printedpage{4}
% head: A. Grothendieck : Dérivateurs, Ch. XI, Édition des Universités de Montpellier II et Paris VII
% foot: Version 28 mai 2002

\textbf{Proposition 1.} \textit{Cette notion est stable par changement de base (trivial), par composition (formel sur a)). Pour que $f$ soit une \textup{Hot}-fibration, il faut et il suffit que $f$ \textup{\texttt{[plutôt}} $f^{o}$\textup{\texttt{]}} le soit.}

\textbf{Proposition 2.} \textit{Soit $f : X \longrightarrow S$ une \textup{Hot}-fibration. Alors $s \longmapsto \mathrm{hot}(X_s)$ est un système local sur $S$, à valeurs dans \textup{Hot}. Donc les $X_s$ pour $s$ dans une même composante connexe de $S$ sont \textup{Hot}-isomorphes entre eux (mais pas canoniquement, si cette composante n'est pas \textup{1}-connexe).}

% (a vertical change bar is printed in the left margin beside Théorème 2)
\textbf{Théorème 2.} {\itshape Soit
\begin{tikzcd}[ampersand replacement=\&, baseline=(\tikzcdmatrixname-1-1.base), column sep=small, row sep=small]
X' \arrow[rr, "f"] \arrow[dr, "p'"'] \& \& X \arrow[dl, "p"] \\
\& S \&
\end{tikzcd}
un diagramme commutatif dans \textup{Cat}, avec $p$, $p'$ des \textup{Hot}-fibrations, et soit $S_0 \subseteq S$ une partie de $S$ telle que $S_0 \longrightarrow \pi_0(S)$ soit surjectif. Pour que $f$ soit une \textup{Hot}-équivalence, il faut et il suffit que les $f_s : X_s \longrightarrow Y_s$ $(s \in S_0)$ le soient.} % [sic] f_s : X_s -> Y_s -> f_s : X'_s -> X_s, the diagram has X' and X (cf. f_s : X'_s -> X_s p. 6)

a) La suffisance se prouve comme tantôt, avec le critère de Artin-Mazur.

\msp{page 132}

Notons que l'hypothèse sur $S_0$ assure que si $F$, $F'$ sont deux faisceaux d'ensembles localement constants sur $S$, et $u : F \longrightarrow F'$ un homomorphisme tel que $u_s : F_s \xrightarrow{\ \sim\ } F'_s$ pour $s \in S$, alors $u$ est \texttt{[un]} isomorphisme. % [sic] s \in S -> s \in S_0, otherwise the hypothesis on S_0 plays no role

1$^{o}$) Soit $F$ un faisceau d'ensembles localement constant sur $X$, $F'$ son image inverse sur $X'$, il faut prouver
\[
H^0(X,F) \ \xrightarrow{\ \sim\ }\ H^0(X',F').
\]
Or
\[
H^0(X,F) \simeq H^0(Y, \underbrace{f_*(F)}_{\substack{\text{localement}\\ \text{constant sur } S}}) \ \longrightarrow\ H^0(Y, f'_*(F')) \simeq H^0(X',F'),
\] % [sic] H^0(Y, f_*(F)), H^0(Y, f'_*(F')) -> H^0(S, p_*(F)), H^0(S, p'_*(F')), the brace says "localement constant sur S" and (*) uses p_*, p'_*
il faut prouver que
\[
p_*(F) \ \xrightarrow{\ \sim\ } p'_*(F') \tag*{$(*)$}
\]
est \texttt{[un]} isomorphisme. Or par hypothèse sur $p$, $p'$, les deux faisceaux sont localement constants et se calculent fibre par fibre, d'où résulte que $(*)$ est un isomorphisme pour les fibres en $s \in S_0$, donc un isomorphisme partout. OK.

2$^{o}$) Il en résulte déjà, si $F$ est un faisceau en \textit{groupes,} que le foncteur
\[
(F\text{-torseurs sur } X) \ \longrightarrow\ (F\text{-torseurs sur } X')
\]
est pleinement fidèle, donc $H^1(X,F) \longrightarrow H^1(X',F')$ \texttt{[est]} injectif; il reste à prouver que ce foncteur est essentiellement surjectif. À cause de la pleine fidélité, c'est \textit{local sur $S$} (\textbf{NB} On voit directement que l'hypothèse fait sur les $s \in S_0$ implique la même hypothèse pour tout $s \in S$, en utilisant le fait que $p$, $p'$ sont des fibrations) (\footnote{Les hypothèses faites sont bel et bien locales sur $S$ (du moins si $S_0 = S$).}). Donc on est ramené au cas où $S = S/s_0$ a un objet final $s_0$.

\printedpage{5}
% head: A. Grothendieck : Dérivateurs, Ch. XI, Édition des Universités de Montpellier II et Paris VII
% foot: Version 28 mai 2002

\msp{page 133}

On a alors
\[
\begin{aligned}
H^1(X,F) &\ \simeq\ H^0(S, R^1p_*(F)) \\
H^1(X',F') &\ \simeq\ H^0(S, R^1p'_*(F')),
\end{aligned}
\]
et on voit encore que
\[
R^1p_*(F) \ \longrightarrow\ R^1p'_*(F')
\]
est \texttt{[un]} isomorphisme (car il se calcule fibre par fibre, et on utilise l'hypothèse sur les $f_s$), donc induit un isomorphisme
\[
H^1(X,F) \ \longrightarrow\ H^1(X',F').
\]

3$^{o}$) Si $F$ est même abélien, on doit prouver
\[
H^i(X,F) \ \xrightarrow{\ \sim\ }\ H^i(X',F')
\]
pour tout $i$, et cela résulte encore de la suite spectrale de Leray pour $(p,F)$ et pour $(p',F')$.

b) Inversement, supposons que $f$ soit une Hot-équivalence, et prouvons que les $f_s : X'_s \longrightarrow X_s$ sont des Hot-équivalences. Je vais utiliser le

% (a vertical change bar is printed in the left margin beside the Lemme)
\textbf{Lemme.} \textit{Soit $S \in \Ob\Cat$, $S$ \textup{0}-connexe. Alors il existe un $\overline{S} \in \Ob\Cat$, $\overline{S}$ \textup{Hot}-asphérique, et $\overline{S} \longrightarrow S$ une \textup{Hot}-fibration.}

Voyons comment ceci implique le théorème 2 (partie réciproque). On peut supposer $S$ 0-connexe, et on fait le changement de base par $\overline{S} \longrightarrow S$

\msp{page 134}

comme dans le lemme :
\[
\begin{tikzcd}[column sep=small, row sep=small]
X' \arrow[rr, "f"] \arrow[dr, "p'"'] & & X \arrow[dl, "p"] & & & & \\
& S & & & \overline{X'} \arrow[rr, "\overline{f}"] \arrow[dr, "\overline{p'}"'] & & \overline{X} \arrow[dl, "\overline{p}"] \\
& & & & & \overline{S} \arrow[ullll, "\varphi"] &
\end{tikzcd}
\quad .
\]

Alors $\overline{p'}$ et $\overline{p}$ sont des Hot-fibrations, et comme $\varphi$ est une Hot-fibration, le changement de base par $\varphi$ transforme la Hot-équivalence $f$ en une Hot-équivalence $\overline{f}$. Considérons alors $\overline{s} \in \Ob\overline{S}$, et le diagramme
\[
\begin{tikzcd}
\overline{X'} \arrow[r, "\overline{f}"] & \overline{X} \\
\overline{X'}_{\overline{s}} \arrow[u, hook, "i'"] \arrow[r, "\overline{f}_{\overline{s}}"] & \overline{X}_{\overline{s}} \arrow[u, hook, "i"']
\end{tikzcd}
\qquad .
\]

\printedpage{6}
% head: A. Grothendieck : Dérivateurs, Ch. XI, Édition des Universités de Montpellier II et Paris VII
% foot: Version 28 mai 2002

Comme $\{\overline{s}\} \longrightarrow \overline{S}$ est une Hot-équivalence, la flèche $i$ qui s'en déduit par le changement de base $\overline{p} : \overline{X} \longrightarrow \overline{S}$ l'est aussi, puisque $\overline{p}$ est une Hot-fibration, et de même pour $i'$. Ainsi dans le carré commutatif précédent, $\overline{f}$, $i$ et $i'$ sont des Hot-équivalences, donc aussi $\overline{f}_{\overline{s}}$. Mais si $s = \varphi(\overline{s})$, $\overline{f}_{\overline{s}}$ s'identifie à $f_s : X'_s \longrightarrow X_s$. D'ailleurs, il est immédiat que $\varphi$, étant une Hot-fibration

\msp{page 135}

à source non-vide, et à but connexe, est surjective. (\textbf{NB} toutes les fibres d'une même composante connexe de $S$ sont Hot-équivalentes, donc si l'une est non-vide, toutes le sont.) Cela achève la démonstration, modulo le lemme.

Il reste à prouver celui-ci. Comme $S = \emptyset$ \texttt{[plutôt} $S \neq \emptyset$\texttt{]}, je peux choisir une ‘origine’
\[
s_0 \in S,
\]
et regarder la catégorie
\[
\Ch_{\infty}(S; s_0, -)
\]
des chemins infinis dans $S$, d'origine $s_0$, fibre de $\Ch_{\infty}(S)$ par le foncteur-source
\[
\sigma : \Ch_{\infty}(S) \ \longrightarrow\ S.
\]
On sait que les fibres de $\sigma$ sont Hot-asphériques (VII), et la catégorie construite ici n'est autre que (avec les notations de VII, p.~152) que la catégorie $\overline{S'}$ qui entre dans la factorisation canonique de
\[
f : \{s_0\} = S' \ \longrightarrow\ S
\]
en
\[
\{s_0\} = S' \ \underbrace{\longrightarrow}_{\text{Hot-équivalence}}\ \overline{S'} \ \xrightarrow{\ \overline{f}\ }\ S,
\]
où ici $\overline{f}$ s'interprète simplement comme le foncteur but (i.e.\ le foncteur induit sur

\msp{page 136}

$\Ch_{\infty}(S; s_0, -)$ par le foncteur but
\[
\beta : \Ch_{\infty}(S) \ \longrightarrow\ S \ \ ).
\]
On a prouvé dans loc.\ cit.\ que $\overline{f}$ est une Hot-fibration (‘propre’ et ‘lisse’), mais dans un contexte qui semblait faire appel tacitement à la théorie du dérivateur HOT. Donc pour mettre les résultats ‘topologiques’ de loc.\ cit.\ sur une base solide, il faut développer les notions pertinentes de foncteur propre, foncteur lisse, sans référence à une théorie HOT. C'est ce que j'avais d'ailleurs fait déjà en 1983, du temps de Pursuing Stacks, mais sans rédaction en forme à ce sujet. Il me faut donc à présent revenir sur ces notions, de façon un peu systématique. D'ailleurs, j'aurais dû mal à construire un seul foncteur Hot-fibrant (sauf les exemples triviaux, tels les isomorphismes, ou les produits $Y \times Z \xrightarrow{\mathrm{pr}_1} Y$), si je ne disposais d'une théorie des foncteurs lisses et des foncteurs propres (\footnote{J'exagère un peu quand même -- après tout il doit être assez évident que les foncteurs de Thomason-Kan sont des Hot-fibrations (mais il y en a encore un peu trop peu !).}). % [sic] dû mal -> du mal, "avoir du mal à" (partitive article, no accent)


\endgroup
% END UNIT Der-ChXI-890/Der-ChXI-890_p006-008.tex
\clearpage\setcounter{footnote}{0}\phantomsection\addcontentsline{toc}{section}{Chapitre XII, 1–4}\markboth{Chapitre XII, 1–4}{Chapitre XII, 1–4}
% BEGIN UNIT Der-ChXIIa-012/Der-ChXIIa-012_p001-005.tex
\begingroup\makeatletter%
\newcommand{\printedpage}[1]{\marginpar{\footnotesize p.~#1}}%
\newcommand{\msp}[1]{\marginpar{\footnotesize [#1]}}%
\newcommand{\Cat}{\mathrm{Cat}}%
\newcommand{\Dl}{\boldsymbol{\Delta}}%
\newcommand{\fl}{\longrightarrow}%
\newcommand{\Ob}{\mathrm{Ob}\,}%
\newcommand{\Hot}{\mathrm{Hot}}%
\newcommand{\hot}{\mathrm{hot}}%
\newcommand{\uHom}{\underline{\mathrm{Hom}}}%
\newcommand{\id}{\mathrm{id}}%
\newcommand{\pr}{\mathrm{pr}}%
\makeatother\let\printedpage\derpp\let\msp\dermsp%
\let\mspin\DERmspin%


\printedpage{--}
% head: (none; no running head and no page number printed on this page)

\begin{center}
{\Large \textsf{LES DÉRIVATEURS}}

\bigskip

\textbf{ALEXANDRE GROTHENDIECK}

\bigskip\bigskip

\textbf{Chapitre XII}

\medskip

\textbf{Caractérisation de $W_\infty$. Foncteurs $W$-propres,\\
$W$-lisses etc. Sommes amalgamées et carrés\\
$W$-cocartésiens dans $\Cat$.}
\end{center}

\par

\printedpage{--}
% head: (none; no running head and no page number printed on this page)

\par\noindent\textbf{Note des éditeurs de l’édition transcrite.}\par
Ce texte a été déchiffré et transcrit en \LaTeXe{} par M. Künzer. Il a été édité par M. Künzer, J. Malgoire, et G. Maltsiniotis. La transcription est aussi fidèle que possible au manuscrit. Pour les quelques corrections évidentes, ou rares commentaires des éditeurs, ainsi que pour la numérotation originale des pages du manuscrit, les caractères de machine à écrire \texttt{[typewriter]} entre crochets sont utilisés. Un point d'interrogation entre crochets signifie que l'on n'est pas sûr du mot qui précède.

Cette édition est provisoire. Les remarques, commentaires et corrections sont bienvenus. Envoyer un message à :

\textbf{[e-mail address removed]}

\hfill G. Maltsiniotis\hspace*{2em}

\par

\printedpage{1}
% head: A. Grothendieck : Dérivateurs, Ch. XII, Édition des Universités de Montpellier II et Paris VII
% foot: Version 7 décembre 2002

\msp{page 1}

\hfill 19.11.90

\section*{I\quad Caractérisation de $W_\infty$.}

Je voudrais caractériser $W_\infty$ en tant que partie de $\mathrm{Fl}(\Cat)$, par ses propriétés comme localiseur, et sans faire appel à la connaissance de notions utilisant déjà $W_\infty$, comme la propreté etc. Je voudrais tout au moins des propriétés sur un localiseur $W$, familières pour $W_\infty$, qui soient si fortes qu'elles impliquent $W \subseteq W_\infty$. Bien sûr, dans les démonstrations je me permets d'utiliser tout ce que je sais sur $W_\infty$.

Bien sûr, il faudra au moins une forme de saturation, peut-être très faible, je présume le suivant suffira (\footnote{(non, j'ai dû le renforcer en 2 [?]).}).

\begin{itemize}[leftmargin=*]
\item[\textbf{W(1)}] (Axiome de saturation.) a) Les identités sont dans $W$, b) si $f$, $g$ \texttt{[sont]} composables et deux parmi $f$, $g$, $gf$ sont dans $W$, le troisième aussi, c) si $fg \in W$, $gf \in W$, alors $f, g \in W$.
\end{itemize}

Une propriété particulièrement cruciale, qui déjà distingue $W_\infty$ de l'ensemble des homotopismes (notion trop stricte à beaucoup d'égards), est l'axiome que j'ai appelé L4 bis :

\begin{itemize}[leftmargin=*]
\item[\textbf{W(2)}] (Axiome de localisation.) Soit $f : X \fl Y$ une flèche dans $\Cat/S$, telle que $\forall\ s \in \Ob S$, $f_s : X/s \fl Y/s$ soit dans $W$ (on dira que $f$ est dans $W$ `localement sur $S$\,'). Alors $f \in W$.
\end{itemize}

\texttt{[Les axiomes W(1) etc. sont appelés aussi} \textcircled{1} \texttt{etc. dans l'original.]}

Dans le cas $Y = S$, cela donne

\textbf{Corollaire :} \textit{Si $f : X \fl Y$ est $W$-asphérique, i.e. $y\backslash X \fl y\backslash Y$} \texttt{[plutôt $X/y \fl Y/y$]} \textit{dans $W$ pour tout $y \in \Ob Y$, alors $f \in W$} (\footnote{On peut considérer la validité du corollaire comme une variante affaiblie W(2 bis) de W(2) ci-dessus. Dans la proposition 1, on n'utilise que W(2 bis), forme faible de W(2).}).

\textbf{Proposition 1.} \textit{Si $W$ satisfait} W(1) \textit{et} W(2)\textit{, les propriétés suivantes sur $W$ sont équivalentes :}

\begin{itemize}[leftmargin=*]
\item[\textit{a)}] \textit{Pour tout $X$ dans $\Cat$, $X \times \Dl^1 \xrightarrow{\pr_1^X} X$ est dans $W$ (i.e. $\Dl^1 \fl e$ est `universellement dans $W$\,').}

\item[\textit{b)}] \textit{Si $X$ dans $\Cat$ a un objet final, alors $X$ est $W$-asphérique, i.e. $X \fl e$ est dans $W$.}

\item[\textit{c)}] \textit{Tout homotopisme est dans $W$.}
\end{itemize}

\par

\printedpage{2}
% head: A. Grothendieck : Dérivateurs, Ch. XII, Édition des Universités de Montpellier II et Paris VII
% foot: Version 7 décembre 2002

\textsc{Démonstration.} b) $\Longrightarrow$ a) (\footnote{n'utilise W(1) que sous forme faible, sans la condition de saturation c). NB Seul b) $\Longrightarrow$ c) utilise W(2 bis), les implications a) $\Longrightarrow$ c) $\Longrightarrow$ b) sont valables sous la seule hypothèse W(1).}), car pour voir que $\pr_1 \in W$, en appliquant le corollaire ci-dessus, il faut prouver que $\forall\ x \in X$
\[ (X \times \Dl^1)/x \fl X/x \tag*{$(*)$} \]
est dans $W$. Or
\[ (X \times \Dl^1)/x \simeq X/x \times \Dl^1, \]
\[
\begin{tikzcd}
\Dl^1 \arrow[d] & X \times \Dl^1 \arrow[l] \arrow[d] & X/x \times \Dl^1 \arrow[l] \arrow[d] \\
e & X \arrow[l] & X/x \arrow[l]
\end{tikzcd}
\]
\texttt{[la flèche verticale à droite est ajoutée]} et les deux catégories $X/x$, $X/x \times \Dl^1$ sont $W$-asphériques car elles ont un objet final, donc $(*)$ est dans $W$ par saturation faible W(1).

a) $\Longrightarrow$ c). Je rappelle que
\[ f : X \fl Y \]
\msp{page 2}
est un homotopisme s'il existe
\[ g : Y \fl X \]
telle que
\[ gf \sim \id_X, \quad fg \sim \id_Y, \tag*{$(*)$} \]
où la relation $(\sim)$ désigne l'homotopie dans $\uHom(X,X)$ et $\uHom(Y,Y)$ respectivement (définie par chemins allant d'un élément à l'autre). Prouvons le

\textbf{Lemme.} \textit{Supposons que $W$ satisfait} W(1) \textit{et le} a) \textit{de la proposition, i.e. $X \times \Dl^1 \fl X$ dans $W$ pour tout $X$ dans} $\Cat$. \textit{Alors si $X$, $Y$} \texttt{[sont]} \textit{dans} $\Cat$ \textit{et $f, g$ dans} $\uHom(X,Y)$ \textit{sont homotopes, on a}
\[ \hot_W(f) = \hot_W(g). \]
\textit{De plus, $f \in W$ si et seulement si $g \in W$.}

Il suffit de le prouver quand il existe un homomorphisme $f \fl g$. Cela peut s'interpréter par l'existence d'une flèche
\[ X \times \Dl^1 \xrightarrow{F} Y \]
dont les restrictions à $X \times \{0\}$ et $X \times \{1\}$ respectivement soient $f$, $g$. Pour prouver que $\hot_W(f) = \hot_W(g)$, il suffit de prouver que
\[ i_0, i_1 : X \fl X \times \Dl^1 \qquad x \longmapsto (x,0),\ x \longmapsto (x,1) \]
\printedpage{3}
% head: A. Grothendieck : Dérivateurs, Ch. XII, Édition des Universités de Montpellier II et Paris VII
% foot: Version 7 décembre 2002
sont égales dans $\Hot_W$, puisque
\[ f = F \circ i_0, \quad g = F \circ i_1. \]
Or cela résulte du fait que ce sont des sections de $X \times \Dl^1$ sur $X$, et que dans $\Hot_W$ la projection $X \times \Dl^1 \xrightarrow{p} X$ est \texttt{[un]} isomorphisme (donc $i_0 = i_1 = p^{-1}$ dans $\Hot_W$). De plus, le diagramme
\[
\begin{tikzcd}
X \arrow[r, "i_0"] \arrow[dr, "f"'] & X \times \Dl^1 \arrow[d, "F"] \\
& Y,
\end{tikzcd}
\]
où $i_0 \in W$, montre que $f \in W$ si et seulement si $F \in W$, et de même $g \in W$ si et seulement si $F \in W$, donc $f \in W \Longleftrightarrow g \in W$.

Revenons à $(*)$, on voit que dans $\Hot_W$ $g \circ f = \id_X$, $f \circ g = \id_Y$, donc $f$ et $g$ sont des $\Hot_W$-isomorphismes inverses l'un de l'autre. D'autre part, $gf \in W$, $fg \in W$, puisque $\id_X, \id_Y \in W$. Donc on gagne, si on a formulé W(1) de façon assez forte. On a prouvé donc a) $\Longrightarrow$ c).

c) $\Longrightarrow$ b) est trivial, puisque si $X$ a un objet final, $X \fl e$ est un homotopisme.

\msp{page 3}

On doit bien sûr par la suite supposer que $W$ satisfait aux conditions équivalentes de la proposition. On prendra p. ex. la forme

\begin{itemize}[leftmargin=*]
\item[\textbf{W(3)}] (Axiome de l'élément final.) Si $X$ a un objet final, alors $X$ est $W$-asphérique.
\end{itemize}

\textbf{Corollaire 1.} \textit{Pour que $f : X \fl Y$ soit $W$-asphérique, il faut et il suffit que les $X/y$ soient $W$-asphériques} ($y \in \Ob Y$).

Car pour que
\begin{tikzcd}[baseline=(\tikzcdmatrixname-1-1.base), column sep=small, row sep=small]
X/y \arrow[rr] \arrow[dr] & & Y/y \arrow[dl] \\
& e &
\end{tikzcd}
soit dans $W$, comme $Y/y$ est $W$-asphérique (objet final), il faut et il suffit que $X/y$ le soit aussi.

\textbf{Corollaire 2.} \textit{Si $X$ dans} $\Cat$ \textit{a un objet initial, il est $W$-asphérique.} (En effet, $X \fl e$ est un homotopisme, et on applique prop. 1.) (\footnote{C'est donc l'énoncé dual de l'axiome W(3) de l'objet final.})

Je mettrai pour mémoire

\textbf{Corollaire 3.} \textit{Tout homotopisme dans} $\Cat$ \textit{est dans $W$. Si $f, g : X \rightrightarrows Y$ dans} $\Cat$ \textit{sont homotopes, alors elles sont égales dans} $\Hot_W$\textit{, et $f$ est dans $W$ si et seulement si $g$ l'est.}

\textbf{Proposition 2.} Soit $f : X \fl Y$ dans $\Cat$. Si $f \in W$, $f^\circ \in W^\circ$ \texttt{[plutôt $f^\circ \in W$].}


\endgroup
% END UNIT Der-ChXIIa-012/Der-ChXIIa-012_p001-005.tex
% BEGIN UNIT Der-ChXIIa-012/Der-ChXIIa-012_p006-010.tex
\begingroup\makeatletter%
\newcommand{\printedpage}[1]{\marginpar{\footnotesize p.~#1}}%
\newcommand{\msp}[1]{\marginpar{\footnotesize [#1]}}%
\newcommand{\Ch}{\underline{\mathrm{Ch}}}%
\makeatother\let\printedpage\derpp\let\msp\dermsp%
\let\mspin\DERmspin%


\printedpage{4}
% head: A. Grothendieck : Dérivateurs, Ch. XII, Édition des Universités de Montpellier II et Paris VII
On utilise la démonstration de I, p. 75, à l'aide de la catégorie nommée $\Phi(X)$ ($\mathrm{Ob}\,\Phi(X) = \mathrm{Fl}(X)$, mais avec une notion de flèche dans $\Phi(X)$ un peu particulière). On trouve que $f \in W \iff \Phi(f) \in W$, et de même $f^\circ \in W \iff \Phi(f) \in W$.

NB. Dans la démonstration de la proposition 2, comme dans celle de la proposition 1 (pour l'implication b) $\Longrightarrow$ a)), on n'a utilisé la propriété de localisation W(2) que sous la forme faible du corollaire, i.e. sous la forme $f$ $W$-asphérique $\Longrightarrow f \in W$.

\textbf{Corollaire} (de la prop. 2) : \textit{Soit $W \subseteq \mathrm{Fl}(\mathrm{Cat})$ satisfaisant l'axiome} W(1) (saturation). \textit{Alors la conjonction des axiomes} W(2), W(3) \textit{équivaut à celle des axiomes duaux} W($2'$), W($3'$) (\footnote[5]{Itou si on remplace dans cet énoncé W(2) par W(2 bis).}).

Il suffit de prouver que 1,2,3 $\Longrightarrow$ $2'$,$3'$.

\msp{page 4}

Pour $2'$ \texttt{[plutôt} $3'$\texttt{]}, c'est le corollaire 2 ci-dessus. Pour $3'$ \texttt{[plutôt} $2'$\texttt{]}, cela résulte aussitôt de la proposition 2. (NB Cet énoncé resterait valable, si on aurait pris W(2) sous la forme faible seulement, cf. le NB page précédente.)

Je dois encore mettre un axiome de la forme L5, L6 (cf. VI, p. 43-44). Mais l'énoncé de ces axiomes utilise la notion de Hot-fibrés, ou de foncteurs parfaits, i.e. fait appel à la connaissance de $W_\infty$, ce que j'aimerais éviter (si je peux). Je dois donc introduire la notion de $W$-fibration, sur le modèle de ce que j'ai fait dans le contexte des dérivateurs, dans I (cf. p. 111-126). On introduit la notion de \textbf{D}-fibration à la page 117 (prop. 5), par \texttt{[des]} conditions équivalentes, dont certaines ont un sens en termes de la donnée de $W$ sans plus. Soit donc
\[ f : X \longrightarrow Y \]
dans Cat et considérons les conditions

\begin{enumerate}
\item[a)] Si $Y'' \longrightarrow Y'$ dans $\mathrm{Cat}/Y$ est dans $W$, alors $X'' \longrightarrow X'$ aussi (le changement de base par $X \longrightarrow Y$ transforme $W$ dans $W$).
\item[a$'$)] Comme a), mais avec $Y''$, $Y'$ $W$-asphériques.
\item[a$_1$)] Comme a), dans le cas particulier où $Y'$ a un objet final et $Y''$ est la catégorie ponctuelle.
\item[a$_2$)] Dual de a$_1$), avec objet initial de $Y'$.
\end{enumerate}

On a donc
\[
\begin{tikzcd}[row sep=large, column sep=1.4em]
 & & \mathrm{a} \arrow[d, Rightarrow] & \text{$W$-fibration ultraforte} & \\
 & & \mathrm{a}' \arrow[dl, Rightarrow] \arrow[dr, Rightarrow] & \text{$W$-fibration forte} & \\
\text{$W$-fibration gauche} & \mathrm{a}_1 & & \mathrm{a}_2 & \text{$W$-fibration droite}
\end{tikzcd}
\]
% foot: Version 7 décembre 2002
\printedpage{5}
% head: A. Grothendieck : Dérivateurs, Ch. XII, Édition des Universités de Montpellier II et Paris VII
et dans le cas $W = W_\infty$ toutes ces conditions sont équivalentes. J'ai envie \texttt{[de]} prendre la plus forte de ces conditions, mais je dois me donner des moyens pour trouver des $W$-fibrations en un sens aussi fort. Par exemple j'aurais envie de savoir que si $X$
\msp{page 5}
est dans Cat, alors $\Ch_\infty(X) \longrightarrow X \times X$ ou $\Ch(X) \longrightarrow X \times X$ est une $W$-fibration. Dans les deux cas, et (disons) pour $W_\infty$, je le prouve en prouvant que $p$ est à la fois lisse, et propre. Il faudrait donc revoir les notions de lissité et de propreté dans le contexte de la donnée d'un $W$.

Si je me rappelle bien, la propriété a$'$) \texttt{[plutôt a}$_1$\texttt{)]} était quelque chose d'assez aisément vérifiable, il suffisait de supposer que les
\[ X_y \longrightarrow X/y \qquad \text{et} \qquad X/y \longrightarrow X/y' \]
(pour $y \in \mathrm{Ob}\,Y$ resp. $u : y \longrightarrow y'$ dans $\mathrm{Fl}\,Y$) sont dans $W$, et que ça reste vrai par tout changement de base - ce qui revient à la condition a$_1$) sur $f$. On pourrait donc formuler une condition sur $W$ ainsi :

\begin{enumerate}[leftmargin=*]
\item[\textbf{W(4)}] (Axiome des $W$-fibrations (\footnote[6]{Voir variante affaiblie suffisante, dans le NB p. 10. Voir encore [?] de variantes pages 13, 14 ci-dessous.}).) Soit $f : X \longrightarrow Y$ dans Cat satisfaisant la condition de $W$-fibration a$_1$) ci-dessus (p. 4), alors elle satisfait même à a), i.e. le changement de base par $f$ transforme flèches dans $W$ en flèches dans $W$ (\footnote[7]{i.e. une $W$-fibration à gauche est une fibration ultraforte.}).
\end{enumerate}

Grâce à la proposition 2, il s'ensuit que la condition duale a$_2$) sur $f$ implique la condition de $W$-fibration forte, donc l'axiome W(4), en apparence asymétrique, signifie aussi que les quatre conditions de fibration sur $f$ sont équivalentes, il est donc en fait symétrique.

\textbf{NB} Je ne vois pas, sans axiomes supplémentaires sur $W$, comment passer de a$_1$ à a$'$ seulement. La difficulté, c'est de montrer, moyennant a$_1$ seulement, que si $Y$ est $W$-asphérique, l'inclusion $X_y \longrightarrow X$ (où $y \in \mathrm{Ob}\,Y$) est dans $W$. Dans le contexte des dérivateurs \textbf{D}, cela se voyait par voie \textbf{D}-cohomologique.

Cet axiome très fort suffit à impliquer que si $X$, $Y$ sont $W$-fibrés sur $S$, et si $f : X \longrightarrow Y$ est un $S$-morphisme qui est dans $W$, alors $f_s : X_s \longrightarrow Y_s$ est dans $W$ ($\forall\ s \in S$). Montrons déjà au moins l'inverse.

\msp{page 6}

\textbf{Proposition 3.} \textit{Soient $X$, $Y$ dans} Cat$/S$\textit{, $f : X \longrightarrow Y$ une $S$-flèche. Supposons ($X$, $Y$ $W$-fibrés) que $\forall\ s \in S$, $f_s : X_s \longrightarrow Y_s$ soit dans $W$, alors $f \in W$} (\footnote[8]{\textit{Sans} axiome W(4). C'est la première fois qu'on a besoin de W(2) sous sa forme forte, à l'exclusion de W(2 bis).}).

Il s'impose d'utiliser l'axiome de localisation, donc de regarder les $f/s$ induits par $f$
% foot: Version 7 décembre 2002
\printedpage{6}
% head: A. Grothendieck : Dérivateurs, Ch. XII, Édition des Universités de Montpellier II et Paris VII
\[
\begin{tikzcd}[row sep=large]
X/s \arrow[r, "f/s"'] & Y/s \\
X_s \arrow[u, "i_s^X"] \arrow[r, "f_s"] & Y_s \arrow[u, "i_s^Y"'] \quad .
\end{tikzcd}
\]

La condition a$_1$ suffit à impliquer que $i_s^X$ et $i_s^Y$ sont dans $W$, donc $f/s$ est dans $W$ si et seulement si $f_s$ l'est. On n'a pas eu à utiliser l'axiome W(4), et il suffisait de supposer que $X$, $Y$ soient fibrés ‘$W$-cohomologiques’ (ou $W$-fibrés à gauche) sur $S$, ou dualement, des $W$-fibrés à droite.

De plus :

\textbf{Corollaire} (\footnote[9]{Sans axiome W(4).}). \textit{Soit $S_0 \subseteq \mathrm{Ob}\,S$ une partie de $S$ telle que $S_0 \longrightarrow \pi_0(S)$ surjectif. Alors dans la proposition 3 il suffit de supposer $f_s \in W$ pour $s \in S_0$.}

Il suffit de voir que si $u : s \longrightarrow s'$ est une flèche dans $S$, alors $f_s \in W \iff f_{s'} \in W$. On utilise le fait que $f_s \in W \iff f/s \in W$, et de même pour $f_{s'}$. Donc à voir que $f/s \in W \iff f/s' \in W$. Or on a
\[
\begin{tikzcd}[row sep=large]
X/s \arrow[r, "f/s"] \arrow[d, "X/u"'] & Y/s \arrow[d, "Y/u"] \\
X/s' \arrow[r, "f/s'"] & Y/s' \quad ,
\end{tikzcd}
\]
où $X/u$ et $Y/u$ sont dans $W$, d'où le résultat.

Mais revenons à la réciproque : si $f \in W$, alors a-t-on $f_s \in W_S$? \texttt{[plutôt} $f_s \in W$\texttt{?]} Il y a un argument dans ce sens dans VI, p. 23, que je vais reprendre. L'idée est celle-ci. On considère la composante connexe $Y_0$ de $Y$ \texttt{[plutôt} $S_0$ \texttt{de} $S$\texttt{]} contenant $s$, et on construit un
\msp{page 7}
$W$-fibré (au \textit{sens fort} a))
\[ S' \longrightarrow S_0 \subseteq S \]
avec $S'$ \textit{$W$-asphérique.} Comme $X \longrightarrow Y$ sur $S$ est dans $W$, il en sera donc de même de $X' \longrightarrow Y'$ sur $S'$. Comme les fibres de $X'$, $Y'$ en $s'$ s'identifient à celles de $X$, $Y$ en l'image de $s'$ dans $S$, on est ramené au cas où $S$ est $W$-asphérique. Mais alors on a
\[
\begin{tikzcd}[row sep=large]
X \arrow[r, "f"] & Y \\
X_s \arrow[u, "i_s^X"] \arrow[r, "f_s"] & Y_s \arrow[u, "i_s^Y"'] \quad ,
\end{tikzcd}
\]
% foot: Version 7 décembre 2002
\printedpage{7}
% head: A. Grothendieck : Dérivateurs, Ch. XII, Édition des Universités de Montpellier II et Paris VII
où $i_s^X$, $i_s^Y$ sont dans $W$, car déduites de $\{s\} \longrightarrow S$ (qui est dans $W$, $S$ étant $W$-asphérique) par les changements de base $W$-fibrants $X \longrightarrow S$, $Y \longrightarrow S$. Donc $f \in W \implies f_s \in W$.

NB On n'a utilisé la propriété de fibration pour $X$, $Y$ qu'au sens a$'$) (\footnote[10]{Et même, dans a$'$) on pourrait exiger que $Y'' = e$, $Y'$ $W_\omega$-asphérique, où $W_\omega$ \texttt{[est]} défini page 14, $W_\omega \subseteq W \cap W_\infty \ldots$}).

Il faut donc prouver, si on peut, le

\textbf{Lemme} (?) (\footnote[11]{Est prouvé moyennant l'axiome W(4) (+ des petites vérifications ultérieures).}). \textit{Soit $S$ un objet} $0$\textit{-connexe de} Cat\textit{. Alors il existe un objet $W$-asphérique $S'$ de} Cat\textit{, et une $W$-fibration (au sens} fort \textit{de a), page 4) $S' \longrightarrow S$.}

L'idée naturelle, c'est de prendre, pour $s_0 \in \mathrm{Ob}\,S$ fixé,
\[ S' = \Ch(S; s_0, -) \qquad \text{ou} \qquad S' = \Ch_\infty(S; s_0, -), \]
et la flèche
\[ \beta : c \longmapsto \text{but } c : S' \longrightarrow S. \]
On sait que $S'$ est $W_\infty$-aphérique et que $\beta$ est $W_\infty$-fibrante (et même parfaite). Il faut voir que $S'$ est aussi $W$-asphérique, et $\beta$ $W$-fibrante. Il faut donc revoir la démonstration dans le cas de $W_\infty$, et prouver qu'elle se transforme à $W$, % [sic] aphérique -> asphérique, lettre s omise
\msp{page 8}
moyennant les axiomes qu'il faut sur $W$. J'ai envie de regarder d'abord le choix $S' = \Ch(S; s_0, -)$, sans passage à une limite inductive qui m'obligerait à mettre un axiome de limites sur $W$. Les propriétés homotopiques délicates et essentielles de $\Ch(S)$ sont établies dans VII, p. 157-169, savoir que dans le diagramme
\[ S \xrightarrow{i_S} \Ch(S) \xrightarrow{p_S} S \times S, \]
$i_S \in W_\infty$ et $p_S$ est une $W_\infty$-fibration (en fait, \textit{propre} et \textit{lisse} pour $W_\infty$) (\footnote[12]{Dans l'énoncé des ‘propriétés essentielles’ de $\Ch(S)$, on peut travailler avec $W_\omega$, la plus petite partie de Fl(Cat) satisfaisant les axiomes W(1), W(2 bis), W(3) (cf. page 14), donc $W_\omega \subseteq W \cap W_\infty$.}). Le fait que $i_S \in W$ \texttt{[plutôt} $i_S \in W_\infty$\texttt{]} implique que les foncteurs source et but
\[ \begin{tikzcd} \Ch(S) \arrow[r, shift left, "\sigma"] \arrow[r, shift right, "\beta"'] & S \end{tikzcd} \qquad \sigma = \mathrm{pr}_1 \circ p_S,\ \beta = \mathrm{pr}_2 \circ p_S \]
sont dans $W_\infty$ (puisque $\sigma i_S = \beta i_S = \mathrm{id}_S$). Comme ce sont des Hot-fibrations, il en résulte que les fibres $\Ch(S; s, -)$ resp. $\Ch(S; -, s)$ sont Hot-asphériques. Mais même en admettant que $\sigma$, $\beta$ sont dans $W$ et sont des $W$-fibrations, je ne peux en conclure tout de suite que les fibres sont $W$-asphériques - il y aurait cercle vicieux! Il me faut donc une démonstration directe.

Considérons le foncteur
\[
\begin{tikzcd}[row sep=huge]
\Ch(S; s, -) \arrow[d, "p"'] \\
\underline{Ty} \quad ,
\end{tikzcd}
\]
% foot: Version 7 décembre 2002
\printedpage{8}
% head: A. Grothendieck : Dérivateurs, Ch. XII, Édition des Universités de Montpellier II et Paris VII
c'est un foncteur Cat-fibrant, donc $W$-lisse (je devrai développer la notion de $W$-lissité (\footnote[13]{$f : X \longrightarrow Y$ est $W$-lisse si les $X_y \longrightarrow y\backslash X$ sont $W$-asphériques.}) avec soin, j'admettrai pour le moment les propriétés formelles familières, et les justifierai par la suite). À cause de la $W$-lissité, pour prouver $p \in W$, il suffit de
\msp{page 9}
prouver que les fibres de $p$ sont $W$-asphériques (\footnote[14]{On utilise ici W(2 bis), pas $W(2)$.}). Or les fibres sont les $Z_\tau = \Ch_\tau(S; s, -)$, qui sont des catégories \textit{contractiles} (i.e. $Z_\tau \longrightarrow e$ est un homotopisme), donc $W$-asphériques par W(3) (cf. cor. 3, page 3).

D'autre part, si on admet que $\Ch(S) \xrightarrow{p_S} S{\times}S$ est $W$-lisse (et $W$-propre), donc $\Ch(S) \xrightarrow{\sigma_S} S$ aussi, étant à fibres $W$-asphériques comme on vient de voir, il s'ensuit que $\sigma_S \in W$, d'où $S \xrightarrow{i_S} \Ch(S)$ aussi dans $W$ (puisque $\sigma_S \circ i_S = \mathrm{id}_S$). Donc il reste à prouver seulement que
\[ \Ch(S) \xrightarrow{p_S} S \times S \tag*{$(*)$} \]
est $W$-propre et $W$-lisse. Il en résultera que
\[ \Ch(S; s_0, -) \xrightarrow{\beta} S \]
est lui aussi $W$-propre et $W$-lisse, car déduit de $p_S$ par le changement de base
\[ S \longrightarrow S \times S \qquad s \longmapsto (s_0, s). \]

Il faut ensuite voir que $W$-lisse et $W$-propre implique la propriété de $W$-fibration \textit{sous sa forme la plus forte} a) (alors que je prévois qu'a priori on ne l'aura que sous les formes faibles a$_1$, a$_2$ - pas même sous la forme a$'$). C'est ici donc que l'axiome W(4) des $W$-fibrations sera essentiel (\footnote[15]{sous la forme affaiblie W(4a$''$) (p. 13).})!

Voici l'argument pour la lissité de $(*)$, dans loc. cit. p. 161 ff. Il faut prouver que pour $c \in \mathrm{Ob}\,\Ch(S)$, et une flèche $u \in \mathrm{Fl}(S \times S)$ \texttt{[plutôt} $(u,v) \in$ \texttt{Fl(}$S \times S$\texttt{)]}, une certaine catégorie de ‘relèvements de $u$ de source $c$ ’ \texttt{[plutôt ‘relèvements de} $(u,v)$ \texttt{de but} $c$ \texttt{’]}, $\mathcal{C} = \mathcal{C}(c,u,v)$, est $W$-asphérique (\footnote[16]{En fait, on prouve que $\mathcal{C}(c,u,v)$ est \textit{contractile}, donc elle est $W_\omega$-asphérique.}). Pour ceci, on la compare à une catégorie $\mathcal{C}'$ à objet final, donc contractile et a fortiori $W$-asphérique, à l'aide de deux foncteurs
\[ \begin{tikzcd} \mathcal{C} \arrow[r, shift left, "\varphi"] & \mathcal{C}', \arrow[l, shift left, "\psi"] \end{tikzcd} \]
et des homotopies $\varphi\psi \sim \mathrm{id}_{\mathcal{C}'}$ (sans mérite, $\mathcal{C}'$ étant contractile)
\msp{page 10}
et $\psi\varphi \sim \mathrm{Id}_{\mathcal{C}}$ (on trouve même une flèche $\psi\varphi \longrightarrow \mathrm{id}_{\mathcal{C}}$). Ainsi $\varphi$ est un homotopisme, donc une $W$-équivalence (cor. 3, p. 3), donc $\mathcal{C}$ est $W$-asphérique, $\mathcal{C}'$ l'étant. La propreté de $p_S$ équivaut à la lissité de $p_{(S^\circ)}$, elle est donc prouvée du même coup.
% foot: Version 7 décembre 2002


\endgroup
% END UNIT Der-ChXIIa-012/Der-ChXIIa-012_p006-010.tex
% BEGIN UNIT Der-ChXIIa-012/Der-ChXIIa-012_p011-015.tex
\begingroup\makeatletter%
\NoAutoSpacing%
\newcommand{\printedpage}[1]{\marginpar{\footnotesize p.~#1}}%
\newcommand{\msp}[1]{\marginpar{\footnotesize [#1]}}%
\newcommand{\Ch}{\underline{\mathrm{Ch}}}%
\makeatother\let\printedpage\derpp\let\msp\dermsp%
\let\mspin\DERmspin%


\printedpage{9}
% head: A. Grothendieck : Dérivateurs, Ch. XII, Édition des Universités de Montpellier II et Paris VII

Ainsi, le lemme de la page 7 est prouvé, moyennant l'axiome W(4) qui assure qu'une flèche qui est à la fois $W$-lisse et $W$-propre, satisfait à la condition très forte a) des $W$-fibrations.

\textbf{NB} En fait, il suffirait, au lieu de l'axiome W(4) sous sa forme forte, de supposer qu'une $f : X \longrightarrow Y$ qui est à la fois `homotopiquement lisse' et `homotopiquement propre', i.e. les fibres des $X/x \longrightarrow Y/y$ et des $x\backslash X \longrightarrow y\backslash Y$ sont contractiles (donc $W$-asphériques pour tout $W$ satisfaisant W(1,2,3)), soit $W$-fibration forte. Énonçons quand même ce que nous avons prouvé.

\textbf{Proposition 4.} \textit{Si $W$ satisfait} W(1,2,3) (\footnote[17]{Pour W(2), il le suffit sous la forme W(2 bis).}) (nous le supposons toujours par la suite) \textit{plus} W(4) (sous une forme un peu affaiblie suffisante, cf. remarque précédente), \textit{on a ceci: Soit $f : X \longrightarrow Y$ une flèche dans} Cat$/S$\textit{, avec $X$, $Y$ des $W$-fibrés sur $S$. Si $f \in W$, alors $f_s : X_s \longrightarrow Y_s$ est dans $W$ pour tout $s \in S$.}

Je me demande si j'en sais assez pour pouvoir prouver
\[ W \subseteq W_\infty . \]
Il faut quand même éliminer des cas canularesques, comme $W = \mathrm{Fl}(\mathrm{Cat})$, donc mettre l'axiome

\begin{itemize}[leftmargin=4em, labelwidth=3em, labelsep=0.6em]
\item[W(5)] (Axiome de connexité (\footnote[18]{Voir à ce sujet le corollaire page 22.}).) Si $f : X \longrightarrow Y$ est dans $W$, $\pi_0(f) : \pi_0(X) \longrightarrow \pi_0(Y)$ \texttt{[est]} bijectif.
\end{itemize}

\msp{page 11}Je factorise alors $f$ en (\footnote[19]{À nouveau, dans ces arguments on a intérêt à travailler avec $W_\omega$.})
\[ X \xrightarrow{i} \bar{X} \xrightarrow{\bar{f}} Y, \]
avec $\bar{f}$ une $W_\infty$-fibration et une $W$-fibration (\footnote[20]{en fait, $\bar{f}$ est à la fois $W_\infty$-parfait (i.e. $W_\infty$-lisse et propre) et $W$-parfait (en fait, $W_\omega$-parfait).}), $i \in W \cap W_\infty$ - cela marche finallement, % [sic] finallement -> finalement, spelling
par la construction de VII, p. 170-171 (de \textsc{Cartan-Serre}), à partir des propiétés de % [sic] propiétés -> propriétés, missing letter
$Y \longrightarrow \Ch(Y) \longrightarrow Y \times Y$ sans plus. Comme $\bar{f} i = f \in W$ et $i \in W$, on en déduit $\bar{f} \in W$. Comme $i \in W_\infty$, on en conclut $f \in W_\infty \iff \bar{f} \in W_\infty$. Donc pour prouver que $f$ est dans $W_\infty$, on est ramené au cas où $f$ est une $W$-fibration et une $W_\infty$-fibration. Mais alors la proposition 4 nous montre que les fibres de $f$ sont $W$-asphériques (\footnote[21]{on n'a utilisé ici que cette propriété sur les fibres, comme consequence de W(4). (On utilise W(4) sous sa forme la plus faible W(4b$''$), cf. pages 13, 14).}), et il suffit de voir qu'elles sont $W_\infty$-asphériques (proposition 3 (\footnote[22]{on n'utilise la proposition 3 que dans le cas $Y = S$, pour lequel W(2 bis) suffit.}) appliquée à $W_\infty$). Ainsi on est ramené au cas où $Y = e$, i.e. à prouver que si $X$ dans Cat est $W$-asphérique, il est $W_\infty$-asphérique.
% [sic] (footnote 21) consequence -> conséquence, missing accent

Soit $x_0 \in \mathrm{Ob}\, X$, considérons comme tantôt
\[ E = \Ch(X, x_0, -) \xrightarrow{p} X, \]
% foot: Version 7 décembre 2002
\printedpage{10}
% head: A. Grothendieck : Dérivateurs, Ch. XII, Édition des Universités de Montpellier II et Paris VII
on sait que $E$ est $W \cap W_\infty$-asphérique, et que c'est une fibration pour $W$ et pour $W_\infty$. Comme $X$ est $W$-asphérique (comme $E$), $p \in W$, donc à nouveau par la proposition 4 (\footnote[23]{en fait, la condition W(4b$''$) sur $W$.}), les fibres de $p$ sont $W$-asphériques, en particulier (par W(5)) 0-connexes. On a donc prouvé le

\msp{page 12}\textbf{Lemme.} (\footnote[24]{Utilise W(1, 2 bis, 3, 4b$''$) sur $W$.}) \textit{Soit $(X, x_0)$ dans} Cat$_\bullet$\textit{. Si $X$ est $W$-asphérique, alors $\Omega(X, x_0) \overset{\text{déf}}{=} \Ch(X; x_0, x_0)$ l'est aussi} (moyennant les axiomes W(1,2,3,4) sur $W$)\textit{. Donc (par itération) les $\Omega^n(X, x_0)$ sont $W$-asphériques} ($\forall\, n \in \mathbf{N}$).

\textbf{Corollaire.} \textit{Si de plus $W$ satisfait l'axiome} W(5) \textit{de connexité, alors $X$ est $W_\infty$-asphérique.}

En effet, `on sait' que cela signifie que tous les $\Omega^n(X, x_0)$ sont 0-connexes. Or l'axiome de connexité W(5) assure que $W$-asphérique implique 0-connexe.

Mais le `on sait' est un tout petit peu léger (\footnote[25]{C'est prouvé fort simplement dans le corollaire du `theorème 1', page 190.}). Pour le justifier, il faudrait, ou bien % [sic] (footnote 25) theorème -> théorème, missing accent
que je fasse le lien avec la théorie 1/2-simpliciale classique, et prouver que les $\Omega^n$ que je construis sont `les bons', i.e. que leur nerf redonne bien (à isomorphisme près dans Hot) les $\Omega^n$ classiques (??). En particulier, que le $\pi_0(\Omega^n(X, x_0))$ est bien le $\pi_n(X, x_0)$ `de tout le monde'. Ou bien que je prouve directement que la 0-connexité des $\Omega^n(X, x_0)$ équivaut à la $W_\infty$-asphéricité de $X$, définie cohomologiquement à la \textsc{Artin-Mazur}. C'est donc un travail qu'il me faudra faire absolument. Mais ce n'est plus qu'un simple travail de routine, et je peux énoncer dès maintenant le

\textbf{Théorème 1.} (\footnote[26]{Il suffit même de W(4) sous une forme affaiblie, avec des foncteurs `fibrants' dans un sens `absolu'. Cf. le NB p. 10 - et surtout, cf. pages 13, 14.},\footnote[27]{NB Au lieu de W(5), il suffit de supposer $W \neq \mathrm{Fl}(\mathrm{Cat})$, cf. th. 2, p. 20.},\footnote[28]{Les hypothèses utiles sur $W$ sont finalement W(1), W(2 bis), W(3), W(4b$''$), et $W \neq \mathrm{Fl}(\mathrm{Cat})$.}) \textit{Soit $W \subseteq$} Fl(Cat) \textit{satisfaisant les conditions} W(1,2,3,4,5) (axiome de saturation, de localisation, de l'objet final, de fibration, de connexité)\textit{. Alors}
\[ W \subseteq W_\infty . \]

\textbf{Corollaire.} \textit{On peut caractériser $W_\infty$ comme la plus grande partie de} Fl(Cat) \textit{satisfaisant aux}
\msp{page 13}\textit{conditions} W(1) \textit{à} W(5) \textit{précédentes.}

Bien sûr, on aimerait un axiome 6 plus intéressant, qui joint aux précédents assure que $W = W_\infty$ (\footnote[29]{Cf. axiome W(6) page suivante.}). Ou est-ce que ces propriétés assurent déjà que $W = W_\infty$? Pour le voir, il faudrait peut-être que j'en sache plus sur la structure des objets $W_\infty$-asphériques? Si c'était faux, je songe à un axiome de `descente homotopique', qui impliquerait qu'un objet $W_\infty$-asphérique est aussi $W$-asphérique.

% foot: Version 7 décembre 2002
\printedpage{11}
% head: A. Grothendieck : Dérivateurs, Ch. XII, Édition des Universités de Montpellier II et Paris VII

Je note dès maintenant que si des $W_i$ ($i \in I$) satisfont aux axiomes W(1) à W(5), il en est de même de $\bigcap_i W_i$ (\footnote[30]{C'est vrai pour tout ensemble d'axiomes dans la nature d'axiomes de stabilité.}). Donc il y a un plus petit ensemble de flèches $W_\omega$ satisfaisant à ces conditions. S'il n'est égal à $W_\infty$, i.e. si les axiomes W(1) à W(5) ne caractérisent déjà $W_\infty$, ce serait un localiseur vraiment très intéressant.

Je me rends compte que je n'ai utilisé qu'une forme assez faible de l'axiome W(4) - essentiellement, la validité de la proposition 4. Je vais résumer différentes variantes de l'axiome de fibration.

\textbf{W(4)}

\begin{itemize}[leftmargin=4em, labelwidth=2.5em, labelsep=0.6em, align=right]
\item[a$_1$)] Si $f : X \longrightarrow Y$ dans Cat est une $W$-fibration gauche, c'est une $W$-fibration ultraforte.

\item[a$_2$)] Si $f : X \longrightarrow Y$ est une $W$-fibration droite, c'est une $W$-fibration ultraforte.

\item[a$'$)] Si $f : X \longrightarrow Y$ est $W$-parfait, i.e. à la fois $W$-propre et $W$-lisse (\footnote[31]{ce qui implique que c'est une $W$-fibration gauche \textit{et} droite.}), alors $f$ est une $W$-fibration ultraforte.

\item[a$''$)] Si $f : X \longrightarrow Y$ est à la fois homotopiquement propre et homotopiquement lisse (cf. NB page 10), $f$ est une $W$-fibration ultraforte. (Il vaut mieux dire : Si $f$ est $W_\omega$-parfait, c'est une $W$-fibration ultraforte.)

\item[b)] Soit $f : X \longrightarrow Y$ une flèche dans Cat$/S$, avec $X$ et $Y$ des $W$-fibrés sur $S$. Si $f \in W$, alors $\forall\, s \in S$, $f_s : X_s \longrightarrow Y_s$ est \texttt{[dans]} $W$.

\item[b$'$)] Soit $f : X \longrightarrow Y$ une flèche de Cat$/S$ qui soit $W$-parfaite et $W_\infty$-parfaite, alors $f \in W$ implique que les fibres de $f$ sont $W$-asphériques. (Il vaut mieux dire : Si $f$ est $W_\omega$-parfait, il est une $W$-fibration ultraforte.) \texttt{[Plutôt : Soit $f : X \longrightarrow Y$ une flèche dans \textrm{Cat}$/S$ telle que $X \longrightarrow S$ et $Y \longrightarrow S$ soient $W_\omega$-parfaites. Si $f \in W$, alors $\forall\, s \in S$, $f_s : X_s \longrightarrow Y_s$ est dans $W$.]}

\item[b$''$)] \msp{page 14}Soit $W_\omega \subseteq W \cap W_\infty$ la plus petite partie de Fl(Cat) satisfaisant les conditions W(1,2,3) (\footnote[32]{Il vaut mieux prendre W(2 bis) au lieu de W(2), pour trouver $W_\omega$ la plus petite possible!}) (axiomes de saturation, de localisation, de l'objet final). Soit $f : X \longrightarrow Y$ une flèche dans Cat qui soit $W_\omega$-parfaite. Si $f \in W$, alors les fibres de $f$ sont $W_\omega$-asphériques \texttt{[plutôt $W$-asphériques]}.
\end{itemize}

Entre ces propriétés de $W$, on a les implications tautologiques (sauf a$''$) $\Longrightarrow$ b), qui est essentiellement la proposition 4)

% foot: Version 7 décembre 2002
\printedpage{12}
% head: A. Grothendieck : Dérivateurs, Ch. XII, Édition des Universités de Montpellier II et Paris VII

\[
\begin{tikzcd}[row sep=huge, column sep=0pt]
\mathrm{a} \arrow[d, Rightarrow] & {}= (\mathrm{a}_1 \Longleftrightarrow \mathrm{a}_2) & \hspace{2.5cm} \\
\mathrm{a}' \arrow[d, Rightarrow] & & \smash{\left.\rule[-1.8cm]{0pt}{3.6cm}\right\}\ \text{\begin{tabular}{l}En fait,\\ toutes ces\\ conditions sont\\ équivalentes\end{tabular}}} \\
\mathrm{a}'' \arrow[d, Rightarrow] & & \\
\mathrm{b} \arrow[d, Rightarrow] & & \\
\mathrm{b}' \arrow[d, Rightarrow] & & \\
\mathrm{b}'' & &
\end{tikzcd}
\]

(\footnote[33]{Je reviens sur les diverses variantes de W(4) plus bas, p. 53-56.}).

Eh bien, dans le lemme page 12, il suffisait de supposer que $W$ satisfait à la version la plus faible b$''$) de l'axiome des fibrations (sept versions!). En effet, les fibrations $\bar{X} \longrightarrow Y$ et $E \longrightarrow X$ utilisées dans la démonstration du lemme sont $W_\omega$-parfaites, et il suffisait de savoir que si elles sont dans $W$, alors les fibres sont $W$-asphériques.

Mais je ne pense pas que même la version la plus forte W(4a) de l'axiome des fibrations, jointe aux trois axiomes précédents et à W(5) (connexité), implique $W = W_\infty$. Il faudrait tout au moins y joindre l'axiome

\begin{itemize}[leftmargin=4em, labelwidth=3em, labelsep=0.6em]
\item[\textbf{W(6)}] (Axiome des limites filtrantes.) Soit $X = \varinjlim\limits_{i} X_i$ une limite inductive filtrante dans Cat. Si les $X_i$ sont $W$-asphériques, $X$ aussi.
\end{itemize}

Ceci devrait être équivalent à la propriété, en apparence plus forte, en termes de flèches de $W$ :

\textbf{Proposition 5.} \textit{Supposons qu'on ait} W(6) (en plus de W(1) à W(4) (\footnote[34]{W(4) sous la forme la plus faible W(4b$''$) .})),
\msp{page 15}\textit{et soit}
\[ (f_i)_{i\in I} : (X_i)_{i\in I} \longrightarrow (Y_i)_{i\in I} \]
% foot: Version 7 décembre 2002
\printedpage{13}
% head: A. Grothendieck : Dérivateurs, Ch. XII, Édition des Universités de Montpellier II et Paris VII
\textit{un homomorphisme de systèmes inductifs dans} Cat\textit{, indexé par un ensemble ordonné filtrant, soient $X = \varinjlim X_i$, $Y = \varinjlim Y_i$, $f = \varinjlim f_i$, $f : X \longrightarrow Y$. Si $f_i$ est dans $W$ pour tout $i$, alors $f \in W$.}

On factorise chaque $f_i$ en
\[ X_i \xrightarrow{\alpha_i} \bar{X}_i \xrightarrow{\bar{f}_i} Y_i, \qquad \alpha_i \in W_\omega,\ \bar{f}_i \in \mathrm{Parf}\, W_\omega \]
(\footnote[35]{$W_\omega$ introduit dans la condition W(4b$''$), page 14.}) suivant le système breveté. La construction étant fonctorielle, les $\bar{X}_i$ forment eux aussi un système inductif. De plus, comme le foncteur $S \longrightarrow \Ch(S)$ qui donne naissance à la factorisation canonique des morphismes, commute aux $\varinjlim$ filtrantes (immédiat), il s'ensuit que, si on pose
\[ \bar{X} = \varinjlim_{i} \bar{X}_i \]
et qu'on regarde les flèches
\[ X \xrightarrow{\alpha} \bar{X} \xrightarrow{\bar{f}} Y, \qquad \alpha = \varinjlim \alpha_i,\ \bar{f} = \varinjlim \bar{f}_i, \]
avec $\bar{f}\alpha = \bar{f}$, cette factorisation de $f$ est bien sa factorisation canonique. Donc $\alpha \in W$, % [sic] $\bar{f}\alpha = \bar{f}$ -> $\bar{f}\alpha = f$, the composite should be f
$\bar{f} \in \mathrm{Parf}\, W_\omega$. Ceci vu, en vertu de prop. 3, $f \in W$, i.e. $\bar{f} \in W$, est vrai si les fibres de $\bar{f}$ sont $W$-asphériques. Mais les fibres de $\bar{f}$ s'interprètent comme limites inductives filtrantes de fibres des $\bar{f}_i$, lesquelles sont $W$-asphériques par la condition W(4b$''$) des fibrations. Donc l'axiome W(6) donne ce qu'il faut.

\textbf{Conjecture} (hardie? sotte?). \textit{Si $W \subseteq$} Fl(Cat) \textit{satisfait les axiomes} W(1) \textit{à} W(6) \textit{(avec} W(4) \textit{sous sa forme la plus faible} W(4b$''$) (\footnote[36]{et peut-être W(2) sous la forme plus faible W(2 bis).})\textit{), alors $W = W_\infty$.}

\msp{page 16}Mais il est bien possible qu'il faille renforcer un peu W(1) (saturation), en exigeant que tout rétracte d'un $f \in W$ soit dans $W$.

\textbf{Remarque.} Considérons, pour $n \in \mathbf{N}$, l'ensemble $W_n \subseteq$ Fl(Cat) des `$n$-équivalences d'homotopie' dans Cat :
\[ f \in W_n \iff \pi_0(f) \text{ bijectif, et } \pi_i(f, x) : \pi_i(X, x) \longrightarrow \pi_i(Y, f(x)) \text{ bijectif pour } 0 < i \leq n. \]

Je croyais que ça satisfaisait à tout ce qu'on peut demander à un localiseur, sans avoir pris la peine de l'écrire avec soin. Ma ça ne peut satisfaire à W(1) jusqu'à W(4b$''$), puisqu'on % [sic] Ma -> Mais, word truncated
n'a pas $W_n \subseteq W_\infty$!

W(1) est trivial sur la définition (et on a sans larmes même la \textit{saturation forte} de $W_i$, % [sic] $W_i$ -> $W_n$ (also in the following lines), the remark concerns W_n
i.e. si $f$ dans Cat est \texttt{[un]} isomorphisme dans $\mathrm{Hot}_{W_i} = (\mathrm{Cat})W_i^{-1}$, alors $f \in W_i$). Je crois me rappeler avoir vérifié W(2), par la même méthode que pour $W_\infty$, en utilisant un critère cohomologique à la \textsc{Artin-Mazur}, pour exprimer $f \in W_i$, et en utilisant les suites spectrales de Leray pour $X \longrightarrow S$ et $Y \longrightarrow S$. La condition W(3) est trivialement satisfaite puisque $W_n \supseteq W_\infty$ et qu'elle l'est pour $W_\infty$. Donc ça canule là où je m'y
% foot: Version 7 décembre 2002


\endgroup
% END UNIT Der-ChXIIa-012/Der-ChXIIa-012_p011-015.tex
% BEGIN UNIT Der-ChXIIa-012/Der-ChXIIa-012_p016-020.tex
\begingroup\makeatletter%
\NoAutoSpacing%
\newcommand{\printedpage}[1]{\marginpar{\footnotesize p.~#1}}%
\newcommand{\msp}[1]{\marginpar{\footnotesize [#1]}}%
\newcommand{\Fl}{\mathrm{Fl}}%
\newcommand{\Cat}{\mathrm{Cat}}%
\newcommand{\card}{\operatorname{card}}%
\makeatother\let\printedpage\derpp\let\msp\dermsp%
\let\mspin\DERmspin%


\printedpage{14}
% head: A. Grothendieck : Dérivateurs, Ch. XII, Édition des Universités de Montpellier II et Paris VII
attendais le moins : l'axiome des fibrations. Prenons la forme faible : Soit $f : X \longrightarrow Y$ une $W_\omega$-fibration et $f \in W_i$, les fibres $X_s$ sont elles $W_i$-asphériques (\footnote[37]{$\omega > \infty$, $W_\omega \subset W_\infty \subset W_n$})? On peut supposer $X$ et $Y$ 0-connexes, soit $x \in X_s$, la suite exacte d'homotopie donne
\[
\pi_{i+1}(X) \xrightarrow[\text{si } i \leq n-1]{\text{iso}} \pi_{i+1}(Y) \longrightarrow \pi_i(X_s) \longrightarrow \pi_i(X) \xrightarrow[\text{si } i \leq n]{\text{iso}} \pi_i(Y),
\]
donc
\[
\pi_i(X_s) = 1 \qquad \text{si } 0 \leq i \leq n-1,
\]
mais pour $i = n$ on trouve\msp{page 17}
\[
\pi_n(X_s) = \mathrm{Coker}\left(\pi_{n+1}(X) \longrightarrow \pi_{n+1}(Y)\right),
\]
il n'y a aucune raison que ce soit nul, comme ça devrait si $X_s \longrightarrow e$ était dans $W_n$.

Il est vrai que parfois on renforce justement la condition pour $W_n$, disons qu'on définit $W'_n \subseteq W_n$ par
\[
f \in W'_n \iff \left\{
\begin{array}{ll}
\text{a)} & \pi_0(f) \text{ bijectif,} \\[1ex]
\text{b)} & \pi_i(f,x) : \pi_i(X,x) \longrightarrow \pi_i(Y,f(x)) \text{ bijectif pour } 0 < i \leq n \text{ et } x \in X, \\[1ex]
\text{c)} & \pi_{n+1}(f,x) : \pi_{n+1}(X,x) \longrightarrow \pi_{n+1}(Y,f(x)) \ \textit{surjectif,}
\end{array}
\right.
\]
c'est à dire aussi
\[
f \in W'_n \iff \left\{
\begin{array}{l}
\text{les `fibres homotopiques' } Z \text{ de } X \longrightarrow Y \text{ (i.e. les fibres de } \bar{X} \xrightarrow{\bar{f}} Y\text{)} \\
\text{sont } n\text{-connexes, i.e. 0-connexe et } \pi_i(Z_y,z) = 0 \text{ pour } 1 \leq i \leq n.
\end{array}
\right.
\]
Alors $X$ dans Cat sera $W'_n$-asphérique si et seulement s'il est $W_n$-asphérique, i.e. $n$-connexe. Donc par la forme 2 de la définition de $W'_n$, l'axiome W(4b$''$) est satisfait. Il faut croire que ça pèche ailleurs! Peut-être déjà la saturation.

\hspace*{3em}Si $gf$ et $g$ sont dans $W'_i$, il n'est pas forcément de même de $f$,

car si $\bar{g}\bar{f}$ (et donc $\bar{g}$) est surjectif, il n'en est par forcément de même pour $\bar{f}$! Donc $W'_n$ ne mérite pas le nom de `localiseur'. % [sic] par forcément -> pas forcément, misprint

Par contre, $W_n$ me paraît excellent, si on ne lui demande pas l'impossible - de satisfaire à un `axiome de fibration' qui semble bien (à peu de choses près) caractériser le seul $W_\infty$ parmi les `bons' localiseurs. En tous cas l'axiome de connexité W(5) est satisfait, et l'axiome des limites filtrantes W(6) aussi. Par ailleurs W(1) (et même la saturation forte), W(2), W(3) sont satisfaits par tout $W_{\mathbf{D}}$ associé à un dérivateur - ce qui confirme à nouveau leur importance. Les axiomes W(4) (l'axiome `draconien' des fibrations, $\pm$ spécial à $W_\infty$), W(5) (connéxité) et % [sic] connéxité -> connexité, misspelling (written connexité elsewhere)
\printedpage{15}
% head: A. Grothendieck : Dérivateurs, Ch. XII, Édition des Universités de Montpellier II et Paris VII
\msp{page 18}W(6) (limites) m'apparaissent comme de nature plus spéciale, je ne songerais à imposer aucun d'eux pour l'usage terminologique de `localiseur fondamental'.

\par
\printedpage{16}
% head: A. Grothendieck : Dérivateurs, Ch. XII, Édition des Universités de Montpellier II et Paris VII
\section*{II\quad Perplexités sur les axiomes et la terminologie.}

Dans I et la plupart du temps, en travaillant avec les dérivateurs, j'ai supposé \fbox{$\ W_{\mathbf{D}} \supseteq W_\infty\ $}, hypothèse qui m'a parue essentielle (à tort (\footnote[38]{à \textit{tort} il s'avère maintenant.}) ou à raison). On voit donc que si $W_{\mathbf{D}}$ satisfait à l'axiome de fibrations W(4), même sous sa forme la plus faible W(4b$''$), il est \textit{égal} à $W_\infty$. Cela montre bien que cet axiome doit être vu comme pratiquement une \textit{caractérisation} de $W_\infty$ (parmi les localiseurs fondamentaux qui contiennent $W_\infty$ -- et ces localisateurs me paraissent les plus importants, et peut-être les seuls vraiment utiles).

La théorie des \textbf{D}-fibrations dans I, p. 111-126, a été faite dans l'esprit qu'on vient de dire, en supposant $W_{\mathbf{D}} \supseteq W_\infty$. La notion de \textbf{D}-fibrations à laquelle je m'étais finalement arreté était par suite différente des notions envisagées ici dans un esprit différent (page 4). Elle revient à tester la condition a) (trop forte presque toujours, sauf pour $W = W_\infty$), caractérisant les $W$-fibrations ultrafortes, par des $Y'' \longrightarrow Y'$, où $Y''$, $Y'$ sont tous deux $W_\infty$-asphériques (ce qui est plus exigeant que $W$-asphérique). C'est seulement quand $Y''$ et $Y'$ sont $W_\infty$-asphériques (au lieu de $W$-asphériques), qu'on demande que $X'' \longrightarrow X'$ soit dans $W$. On trouve donc\msp{page 19} une notion moins exigeante que celle de $W$-fibration forte. Si on dénote ici cette condition par b), on trouve (sous réserve que $W_\infty \subseteq W$), avec les notations de la page 4 (\footnote[39]{Cf. la bonne notion de $W$-fibration dans un contexte où on ne suppose pas $W_\infty \subseteq W$, dans th. 3 et déf. 3, p. 51, 52.}) % [sic] arreté -> arrêté, missing circumflex
% In the print, a curved arrow starts at "b" in the diagram below and points up to "b)" in the line "par b), on trouve ..." above; it is drawn here to the empty node at the top left.
\[
\begin{tikzcd}[column sep=small, row sep=large]
{} & & & & \\
& & \mathrm{a} \arrow[d, Rightarrow] & & W\text{-fibration ultraforte} \\
& & \mathrm{a}' \arrow[d, Rightarrow] & & W\text{-fibration forte} \\
& & \mathrm{b} \arrow[dl, Rightarrow] \arrow[dr, Rightarrow] \arrow[uuull, bend left=30] & & W\text{-fibration} \\
W\text{-fibration gauche} & \mathrm{a}_1 & & \mathrm{a}_2 & W\text{-fibration droite.}
\end{tikzcd}
\]
a$_1$), a$_2$) étaient aussi appelées \textbf{D}-fibrations cohomologique resp. \textbf{D}-fibrations homologiques (\footnote[40]{Non, c'est une notion encore plus forte, à coups de coefficients localement constants, cf. I, p. 121.}). Quand \textbf{D} satisfait Der 3a$'$ (I, p. 111) sur les \textbf{D}-coefficients localement constants sur les $X$ $W_\infty$-asphériques, alors les notions b), a$_1$), a$_2$) sont équivalentes. (Je présume que c'est le cas pour tous les \textbf{D} et tous les $W$ vraiment intéressants.) Par contre, j'ai l'impression que la notion de $W$-fibration forte, voire ultraforte, n'est guère utilisable que dans le cas $W = W_\infty$, où elle coïncide avec les précédentes. C'est dire qu'on peut probablement oublier la terminologie, et se rappeler par contre que dans le cas $W = W_\infty$, les % [sic] cohomologique -> cohomologiques, agreement with "D-fibrations"
\printedpage{17}
% head: A. Grothendieck : Dérivateurs, Ch. XII, Édition des Universités de Montpellier II et Paris VII
$W$-fibrations ont des propriétés d'une force extrême, s'exprimant par les deux propriétés dégagées dans W(4a) et W(4b).

\msp{page 20}Dans la théorie de $\mathrm{HOT}_W$, influencé par la réflexion générale sur les dérivateurs et l'axiome Der 3b, j'ai travaillé sous l'hypothèse $W \supseteq W_\infty$, mais en utilisant un axiome de fibration L5 $=$ W(4b) qui, il s'avère, implique finalement $W = W_\infty$. J'ai l'impression que je n'ai jusqu'à présent (propriétés de $f_!$ pour HOT) pas utilisé $W \supseteq W_\infty$, donc ça peut être utile, comme `ça ne coûte pas plus cher' (et oblige d'être plus attentif à ce qu'on utilise), de laisser tomber $W \supseteq W_\infty$ (impliquant $W = W_\infty$), et de mettre par contre sur $W$ des axiomes aussi forts (notamment de fibration) que les propriétés de $W_\infty$, `sans se priver en rien'. Je supposerai donc W(1), W(2), W(3), W(4a) (la forme la plus forte de l'axiome de fibration, cf. pages 13, 14) - c'est tout ce qui a été utile jusqu'à présent. Si on y joint W(5) (axiome de connexité) ça implique $W \subseteq W_\infty$. Mais pour l'instant je ne crois pas en avoir un besoin.

Mais méditant sur la question, je crois avoir trouvé le

\textbf{Théorème 2} (!) (\footnote[41]{Cf. pour un résultat plus fort XV 9.12, p. 133 bis. \textbf{NB} Le résultat est plus général, mais pas vraiment plus fort, car il ne donne pas $W = \Fl(\Cat)$.}) : \textit{Soit} $W \subseteq \Fl(\Cat)$ \textit{satisfaisant} W(1,2,3,4b$''$)\textit{. Si} $W$ \textit{ne satisfait pas à} W(5)\textit{, alors} $W = \Fl(\Cat)$.

\textsc{Démonstration}. Soit $f : X \longrightarrow Y$ dans Cat, tel que $\pi_0(f)$ ne soit pas bijectif. Soit $I$ la catégorie discrète ayant pour ensemble d'objets l'ensemble $\pi_0(Y)$, on a donc un diagramme commutatif
\[
\begin{tikzcd}[column sep=small]
X \arrow[rr, "f"] \arrow[dr, "p"'] & & Y \arrow[dl, "q"] \\
& I, &
\end{tikzcd}
\]
d'autre part un foncteur d'une catégorie dans une catégorie discrète est toujours $W_\omega$-parfait, donc $p$, $q$ sont des $W$-fibrations. Comme\msp{page 21} $f \in W$, il s'ensuit $f_i \in W$ pour tout $i \in I$ (\footnote[42]{On utilise W(4) sous forme plus forte que W(4b$''$)!}).

Si $\pi_0(f)$ n'était pas surjective, prenant $i$ qui n'est pas dans l'image, on aurait $X_i = \emptyset \longrightarrow Y_i \neq \emptyset$, qui serait dans $W$. Soit donc
\[
f : X = \emptyset \longrightarrow Y
\]
dans $W$. Alors $\forall$ $Z$ dans Cat, $f \times \mathrm{id}_Z : \underbrace{X \times Z}_{= \emptyset} \longrightarrow Y \times Z$ est dans $W$ (car $Z$ est
\printedpage{18}
% head: A. Grothendieck : Dérivateurs, Ch. XII, Édition des Universités de Montpellier II et Paris VII
$W$-fibrée sur $e$), donc $\emptyset \longrightarrow Y \times Z$ est dans $W$. Donc aussi la projection
$\begin{tikzcd}[column sep=small, row sep=small, baseline=(\tikzcdmatrixname-1-1.base)]
Y \times Z \arrow[rr] & & Y. \\
& \emptyset \arrow[ul] \arrow[ur] &
\end{tikzcd}$
Comme cette projection est une $W$-fibration, il s'ensuit par W(4b$''$) que ses fibres sont $W$-asphériques. Comme $Y \neq \emptyset$, il y a une telle fibre, elle est $\simeq Z$. Donc $Z$ est $W$-asphérique. Donc on trouve que \textit{toute} catégorie $Z$ dans Cat est $W$-asphérique. Donc toute flèche $Z \longrightarrow Z'$ dans Cat est dans $W$, q.e.d.

Supposons maintenant $\pi_0(f)$ surjective, mais non injective. Il y aurait une fibre de $\pi_0(X) \longrightarrow \pi_0(Y) = I$ qui serait de cardinal $> 1$. En prenant le $f_i$ correspondant, on trouve $X_i \longrightarrow Y_i$ avec $Y_i$ 0-connexe et $\card \pi_0(X_i) \geq 2$. Donc on est ramené au cas
\[
f : X \longrightarrow Y \qquad\qquad \left\{
\begin{array}{l}
f \in W \\[1ex]
Y \text{ 0-connexe} \\[1ex]
\card \pi_0(X) \geq 2.
\end{array}
\right.
\]
Considérons la factorisation canonique
\[
X \xrightarrow{i_f} \bar{X} \xrightarrow{\bar{f}} Y,
\]
avec $i_f \in W$, donc $\bar{f} \in W$. Donc par W(4b$''$) les fibres de $\bar{f}$ sont $W$-asphériques. Mais on voit tout de suite, en décomposant $\bar{X}$ suivant les composantes connexes $X_i$ de $X$, que ces fibres ont un $\pi_0$ de cardinalité $\geq \card \pi_0(X) \geq 2$. Donc on est ramené au cas où $Y = e$, i.e.
\[
X \ W\text{-asphérique}, \qquad \card \pi_0(X) \geq 2.
\]
Considérons $X \times X$, qui est aussi $W$-asphérique, et l'application diagonale\msp{page 22}
\[
d : X \longrightarrow X \times X = Y,
\]
qui est donc dans $W$. Mais $\pi_0(d)$ s'identifie à l'application diagonale $\pi_0(X) \longrightarrow \pi_0(X) \times \pi_0(X)$, qui n'est pas surjective puisque $\card \pi_0(X) \geq 2$. Donc pour la première partie, $W = \Fl(\Cat)$.

Donc sous forme plus positive :

\textbf{Corollaire.} \textit{Soit} $W \subseteq \Fl(\Cat)$ \textit{qui satisfait à} W(1,2,3,4b$''$)\textit{, et telle que} $W \neq \Fl(\Cat)$. \textit{Alors $W$ satisfait l'axiome de connexité} W(5)\textit{, et par suite (th. 1, page 12)} $W \subseteq W_\infty$.


\endgroup
% END UNIT Der-ChXIIa-012/Der-ChXIIa-012_p016-020.tex
% BEGIN UNIT Der-ChXIIa-012/Der-ChXIIa-012_p021-025.tex
\begingroup\makeatletter%
\derleqno%
\newcommand{\printedpage}[1]{\marginpar{\footnotesize p.~#1}}%
\newcommand{\msp}[1]{\marginpar{\footnotesize [#1]}}%
\newcommand{\Ob}{\mathrm{Ob}\,}%
\makeatother\let\printedpage\derpp\let\msp\dermsp%
\let\mspin\DERmspin%


\printedpage{19}
% head: A. Grothendieck : Dérivateurs, Ch. XII, Édition des Universités de Montpellier II et Paris VII
% foot: Version 7 décembre 2002

\section*{III\quad Classes remarquables de foncteurs liés à $W$.}

Finalement, j'appelerai \textit{‘localiseur fondamental’} une partie $W$ de Fl(Cat) satisfaisant les trois conditions suivantes : % [sic] appelerai -> appellerai, spelling

\begin{enumerate}
\item[(i)] Saturation W(1) (cf. rappel p. 1).
\item[(ii)] Axiome de l'élément final W(3) (page 3).
\item[(iii)] La variante affaiblie W(2 bis) de l'axiome de localisation W(2) donnée dans le corollaire page 1 (anciennement c'était L4$'$ dans Pursuing Stacks).
\end{enumerate}

Quand j'aurai besoin de plus sur $W$, je le dirai explicitement. Ces trois axiomes ont suffi pour pratiquement tout ce que j'ai fait avec $W$ dans Pursuing Stacks, la seule exception étant un usage (si je me rappelle) unique, quelque part, d'une variante affaiblie de W(2) (pour les $S$-foncteurs \textit{cartésiens} entre $S$-catégories \textit{fibrées} - si sur toute fibre il est dans $W$, il est lui-même dans $W$).

On suppose donc par la suite donné un localiseur fondamental
\[
W \subseteq \mathrm{Fl}(\mathrm{Cat}).
\]

\textbf{Proposition 1.} \textit{Soit $f : X \longrightarrow Y$ dans} \textup{Cat}. \textit{Conditions équivalentes :}

\begin{enumerate}
\item[\textit{a)}] \textit{Pour tout isomorphisme local $Y' \longrightarrow Y$, $f' : X' \longrightarrow Y'$ est $\in W$.}
\item[\textit{b)}] \textit{Même condition, dans le cas particulier $Y' = Y/y \longrightarrow Y$, donc $\forall\ y \in Y$, $X/y \simeq X \times_Y (Y/y)$ est $W$-asphérique.}
\end{enumerate}

\msp{page 23}\textbf{Définition 1.} Sous ces conditions équivalentes a), b) de la proposition 1, on dit que $f$ est \textit{$W$-asphérique.} On dit \texttt{[que]} $f$ est \textit{$W$-coasphérique} si $f^{\circ}$ est $W$-asphérique, i.e. si les $y\backslash X$ sont $W$-asphériques (pour $y \in \Ob Y$).

\textsc{Démonstration.} Il suffit de prouver b) $\Longrightarrow$ a). Par définition d'un isomorphisme local, on voit que la condition b) est stable par changement de base $Y' \longrightarrow Y$, avec $Y'$ un isomorphisme local (puisque les $Y'/y' \xrightarrow{\sim} Y/y$, donc $X'/y' \xrightarrow{\sim} X/y$). Il suffit donc de voir que si $f$ satisfait à b), alors $f \in W$. Mais c'est l'axiome W(2 bis).

\begin{center}
\begin{tikzcd}[row sep=large]
X \arrow[d] & & & \\
Y & & Y' \arrow[ll] & \\
& Y/y \arrow[ul] & & Y'/y' \arrow[ll] \arrow[ul]
\end{tikzcd}
\end{center}

\printedpage{20}
% head: A. Grothendieck : Dérivateurs, Ch. XII, Édition des Universités de Montpellier II et Paris VII
% foot: Version 7 décembre 2002

\textbf{Remarque.} Si $f$ a un adjoint à droite $g$, alors pour tout $y \in Y$, $X/y \simeq X/g(y)$, donc $X/y$ a un objet final et est donc $W$-asphérique. Donc si $f$ a un adjoint à droite (à gauche), alors $f$ est $W$-asphérique ($W$-coasphérique).

Mais même si $f$ n'a pas d'adjoint à droite, il y a quand même un adjoint généralisé
\[
g : Y \longrightarrow X^{\wedge}
\]
et on a
\[
X/y \simeq X/g(y) \qquad \text{(pour l'inclusion pleinement fidèle $X \hookrightarrow X^{\wedge}$),}
\]
donc $\mathrm{Top}(X/y)$ s'identifie au topos induit $\mathrm{Top}(X)/g(y)$. La condition $f$ $W$-asphérique signifie donc que les faisceaux $g(y)$ dans $\mathrm{Top}(X)$ sont
\msp{page 24}
$W$-asphériques. Pour de nombreuses autres formulations équivalentes, je renvoie à Pursuing Stacks.

\textbf{Proposition 2.} \textit{Soit $X \xrightarrow{f} Y \xrightarrow{g} Z$ dans} \textup{Cat}. \textit{Si $f$ et $g$ sont $W$-asphériques ($W$-coasphériques), de même $gf$. Si $gf$ et $f$ sont $W$-asphériques ($W$-coasphériques), $g$ aussi.}

\textbf{Rappel.} Un morphisme $p : S' \longrightarrow S$ dans Cat est un \textit{isomorphisme local} si $\forall\ s' \in \Ob S'$, le morphisme induit
\[
p_{s'} : S'/s' \xrightarrow{\sim} S/s \qquad (s = p(s'))
\]
est un isomorphisme : $p$ induit un isomorphisme sur les ‘localisées’ de $S'$ et $S$ en $s'$ et $s = p(s')$ respectivement. Il est immédiat que cette notion est stable par composition, que si $gf$ et $g$ sont des isomorphismes locaux, alors $f$ l'est aussi, que tout isomorphisme est un isomorphisme local. De plus, les isomorphismes locaux sont \textit{stables par changement de base,} ce qui résulte du fait (trivial et important) que si on a un carré cartésien dans Cat

\begin{center}
\begin{tikzcd}[row sep=large]
X & X' \arrow[l, "q"'] \arrow[d, "f'"] \\
S \arrow[from=u, "f"'] & S', \arrow[l, "p"']
\end{tikzcd}
\end{center}

et si $x' \in \Ob X'$ a pour image $x$, $s'$, $s$ dans $X$, $S'$, $S$, alors le diagramme

\begin{center}
\begin{tikzcd}[row sep=large]
X/x \arrow[d, "f/x"'] & X'/x' \arrow[l, "q/x'"'] \arrow[d, "f'/x'"] \\
S/s & S'/s' \arrow[l, "p/s'"']
\end{tikzcd}
\end{center}

est également cartésien. (Donc si $p/s'$ est \texttt{[un]} isomorphisme, $q/x'$ aussi.)

\printedpage{21}
% head: A. Grothendieck : Dérivateurs, Ch. XII, Édition des Universités de Montpellier II et Paris VII
% foot: Version 7 décembre 2002

\msp{page 25}Revenons à la proposition 2. La stabilité de la notion de morphisme asphérique par composition résulte aussitôt de la forme a) de la condition. De même pour la condition que si $gf$ et $f$ sont $W$-asphériques, $g$ aussi - il faut montrer que si $Z' \longrightarrow Z$ est \texttt{[un]} isomorphisme local, alors $g'$ est dans $W$, or l'hypothèse implique que $g'f'$ et $f'$ le sont.

\begin{center}
\begin{tikzcd}[row sep=large]
X \arrow[d, "f"'] & X' \arrow[l, "r"'] \arrow[d, "f'"] \\
Y \arrow[d, "g"'] & Y' \arrow[l, "q"'] \arrow[d, "g'"] \\
Z & Z' \arrow[l, "p"']
\end{tikzcd}
\end{center}

\textbf{Corollaire 1.} \textit{L'image inverse d'un $f : X \longrightarrow Y$ qui est $W$-asphérique (resp. $W$-coasphérique) par un isomorphisme local (resp. par un coïsomorphisme local) est $W$-asphérique ($W$-coasphérique).}

\begin{center}
\begin{tikzcd}[row sep=huge]
X \arrow[d] & X' \arrow[l] \arrow[d] & X'' \arrow[l] \arrow[d] \\
Y & Y' \arrow[l, "\text{iso local}"] & Y'' \arrow[l, "\text{iso local}"]
\end{tikzcd}
\end{center}

Résulte formellement de la définition par a) et du fait qu'un composé de deux isomorphismes locaux (resp. coïsomorphismes locaux) est itou.

\textbf{NB} Bien sûr, la notion de foncteur $W$-asphérique n'est pas stable par changement de base quelconque. Par exemple, l'inclusion $X = \{y_0\} \longrightarrow Y$ est $W$-asphérique si et seulement si les catégories discrètes $\mathrm{Hom}(y_0, y)$ sont $W$-asphériques, i.e. (sous réserve que l'on ait W(5), axiome de connexité) si et seulement si les $\mathrm{Hom}(y_0, y)$ sont de cardinal 1, i.e. $y_0$ \textit{objet initial} de $Y$. Mais si $Y \neq \{y_0\}$, prenant $y$ dans $Y$, $y \neq y_0$, et posant $Y' = \{y\} \longrightarrow Y$, on a $X' = X \times_Y Y' = \emptyset$ et $\emptyset \longrightarrow e$ n'est pas asphérique!

\msp{page 26}Par contre, on s'attend que les morphismes $W$-asphériques soient stables par changements de base \textit{$W$-lisses} (qui généralisent les isomorphismes locaux), et les morphismes $W$-coasphériques par changements de base $W$-propres. Cf. plus bas.

\textbf{Corollaire 2.} \textit{Soit $f : X \longrightarrow Y$ dans} \textup{Cat} \textit{une} \textup{équivalence} \textit{de catégories. Alors $f$ est $W$-asphérique et $W$-coasphérique.}

En effet, $f$ a un adjoint à gauche et un adjoint à droite, et on utilise la remarque page 23.

\textbf{Proposition 3.} \textit{Soit $f : X \longrightarrow Y$ dans} \textup{Cat}. \textit{Conditions équivalentes :}

\printedpage{22}
% head: A. Grothendieck : Dérivateurs, Ch. XII, Édition des Universités de Montpellier II et Paris VII
% foot: Version 7 décembre 2002

\begin{enumerate}
\item[\textit{a)}] \textit{Pour tout $y \in \Ob Y$, le foncteur d'inclusion}
\[
X_y \longrightarrow X/y \tag*{$(*)$}
\]
\textit{est $W$-}co\textit{asphérique.}
\item[\textit{b)}] \textit{Pour tout $x_0 \in X$, le morphisme induit pour les colocalisées}
\[
x_0\backslash X \longrightarrow y_0\backslash Y \qquad \textit{(où $y_0 = f(x_0)$)}
\]
\textit{est à fibres $W$-asphériques.}
\item[\textit{c)}] \textit{Pour tout morphisme $\boldsymbol{\Delta}^1 = Y' \longrightarrow Y$, considérons $X' = X \times_Y Y'$ sur $\boldsymbol{\Delta}^1$, l'inclusion $X'_1 \longrightarrow X'$ de la fibre fermée de $X'$ (fibre sur le point fermé de $Y'$) est $W$-coasphérique.}
\item[\textit{d)}] \textit{Pour toute flèche $u_0 : y_0 \longrightarrow y_1$ dans $Y$, et tout $x_0$ dans $X_{y_0}$, la catégorie $X(x_0;u_0)$ formée des flèches $u' : x_0 \longrightarrow x$ dans $X$ de source $x_0$, et qui relèvent $u$} \texttt{[plutôt} $u_0$\texttt{]}\textit{, est $W$-asphérique.}
\end{enumerate}

\msp{page 27}\textsc{Démonstration.} Pour expliciter la condition a), il faut expliciter la catégorie $\xi_0\backslash(X_y)$, pour un objet quelconque $\xi_0 = (x_0, u_0)$ de $X/y$, $x_0 \in \Ob X$,
\[
u_0 : y_0 \longrightarrow y_1 = y \quad \text{dans $Y$, avec $y_0 = f(x_0)$.}
\]

Cette catégorie $\xi_0\backslash X_{y_1}$ s'explicite alors comme la catégorie des $x_1$ dans $X_{y_1}$, \textit{munis} d'une flèche de $\xi_0$ dans l'image $\xi_1 = (x_1, \mathrm{id}_{f(x_1)} = \mathrm{id}_{y_1})$ de $x_1$ dans $X/y_1$ dans $\xi_0$ \texttt{[‘dans} $\xi_0$\texttt{’ à effacer]}, i.e. d'une flèche $u$ de $x_0$ dans $x_1$, telle que $f(u)$ donne lieu au triangle commutatif dans $Y$, i.e. $f(u) = u_0$. % AUDIT resolved: regular CMTT12 editorial bracket; see corrections.csv, PDF leaf 32

\begin{center}
\begin{tikzcd}
x_1 \arrow[dd, dash] & & \\
& & x_0 \arrow[ull, "u"'] \arrow[dd, dash] \\
y_1 & y_1 \arrow[l, "\mathrm{id}"'] & \\
& & y_0 \arrow[ul, "u_0"'] \arrow[ull, "f(u)"]
\end{tikzcd}
\end{center}

Donc c'est la catégorie notée $X(x_0;u_0)$ dans d). On voit tout de suite qu'elle est aussi isomorphe à la fibre de $x_0\backslash X \longrightarrow y_0\backslash Y$ en l'élément $u_0$ de $y_0\backslash Y$, dont il est question dans b). D'où l'équivalence de a) b) c). D'autre part, la donnée de $u_0$ dans $Y$ équivaut à celle d'un foncteur $\boldsymbol{\Delta}^1 \longrightarrow Y$, comme dans c), et pour exprimer alors que $\underbrace{X'_1}_{\simeq X_{y_1}} \longrightarrow X'$ est coasphérique, il suffit d'exprimer que les $x'\backslash X'_1$ ($x' \in \Ob X'$) sont $W$-asphériques. Or c'est trivial si $x' \in \Ob X'_1$ (objet initial), donc il suffit de le vérifier pour $x'$ dans $X'\backslash X'_1 = X'_0 \simeq X_{y_0}$. Alors $x'$ s'identifie à un élément $x_0$ de $X_{y_0}$, et on vérifie tout de suite que $x_0\backslash X'_1$ est isomorphe à la catégorie $X(x_0;u_0)$ de d). Cela achève la démonstration.

\printedpage{23}
% head: A. Grothendieck : Dérivateurs, Ch. XII, Édition des Universités de Montpellier II et Paris VII
% foot: Version 7 décembre 2002

\msp{page 28}\textbf{Définition 2.} Sous les conditions de la proposition 3, on dit que $f$ est \textit{$W$-propre.} On dit que $f$ est \textit{$W$-lisse} si $f^{\circ}$ est $W$-propre, i.e. si les morphismes d'inclusion $(*)$ \texttt{[plutôt} $X_y \longrightarrow y\backslash X$\texttt{]} sont $W$-asphériques etc. (conditions duales de celles de la proposition 3).

\textbf{Remarque.} La condition a) de la proposition 3 est satisfaite en particulier si l'inclusion $(*)$ a un adjoint à gauche (cf. remarque page 23). Cela signifie que les catégories $X(x_0;u)$ ont un objet initial. Cela signifie aussi que la catégorie $X$ sur $Y$ est Cat-cofibrée sur $Y$ au sens faible, i.e. qu'il existe, pour $u_0 : y_0 \longrightarrow y_1$ flèche dans $Y$, et $x_0$ dans $X_{y_0}$, un objet $u_{0*}(x_0)$ dans $X_{y_1}$ qui représente le foncteur covariant $X_{y_1} \longrightarrow (\mathrm{Ens})$, $x_1 \longmapsto \mathrm{Hom}_u(x_0, x_1)$. Ainsi, si $X$ est Cat-cofibrée au sens faible sur $Y$, et a fortiori si elle est Cat-cofibrée, alors $f : X \longrightarrow Y$ est $W$-propre. Dualement, si $f$ est faiblement Cat-fibrée, $f$ est \texttt{[}$W$\texttt{-]}lisse.

\textbf{Proposition 4.} (\footnote[43]{Dire aussi : Si $X \xrightarrow{f} Y \xrightarrow{g} Z$, si $gf$ et $f$ $W$-lisses, alors $g$ est lisse en les points de $\mathrm{Im}\, f$, idem pour propre.}) \textit{La notion de morphisme $W$-propre est stable par composition, par changement de base. Les coisomorphismes locaux sont propres (trivial sur b)), a fortiori les isomorphismes sont propres, les immersions fermées sont propres.}

\textit{Énoncé dual pour les morphismes lisses : stabilité par changement de base, par composition, les isomorphismes locaux sont lisses, notamment les immersions ouvertes.}

\msp{page 29}La stabilité par changement de base est immédiate sur la forme c). La stabilité par composition \texttt{[est]} moins triviale :

\begin{center}
\begin{tikzcd}[row sep=large]
{\scriptstyle x} &[-1.5em] X \arrow[d, "f"] & x\backslash X \arrow[l] \arrow[d] & \bar{X} \arrow[l] \arrow[d] \\
{\scriptstyle y} & Y \arrow[d, "g"] & y\backslash Y \arrow[l] \arrow[d] & \bar{Y} \arrow[l] \arrow[d] \\
{\scriptstyle z} & Z & z\backslash Z \arrow[l] & \{u\} \arrow[l]
\end{tikzcd}
\end{center}

Dans le diagramme à carrés cartésiens ci-contre \texttt{[plutôt dont les deux carrés de droite sont cartésiens]} (déduit de $x \in \Ob X$, d'où $y = f(x)$, $z = g(y) = gf(x)$, et de $u \in \Ob z\backslash Z$), il faut prouver que $\bar{X}$ est $W$-asphérique. On sait que $\bar{Y}$ l'est ($g$ propre) et que $\bar{X} \longrightarrow \bar{Y}$ est à fibres $W$-asphériques. Il suffit donc d'utiliser la proposition 5 ci-dessous, et la remarque immédiate % AUDIT resolved: regular CMTT12 editorial bracket; see corrections.csv, PDF leaf 33

\textbf{Lemme.} \textit{Si $f$ est propre, il en est de même de $x\backslash f : x\backslash X \longrightarrow y\backslash Y$, pour tout $x$ dans $X$, $y = f(x)$.}


\endgroup
% END UNIT Der-ChXIIa-012/Der-ChXIIa-012_p021-025.tex
% BEGIN UNIT Der-ChXIIa-012/Der-ChXIIa-012_p026-030.tex
\begingroup\makeatletter%
\newcommand{\printedpage}[1]{\marginpar{\footnotesize p.~#1}}%
\newcommand{\msp}[1]{\marginpar{\footnotesize [#1]}}%
\newcommand{\Cat}{\mathrm{Cat}}%
\newcommand{\Hot}{\mathrm{Hot}}%
\newcommand{\hot}{\mathrm{hot}}%
\newcommand{\Ob}{\mathrm{Ob}}%
\makeatother\let\printedpage\derpp\let\msp\dermsp%
\let\mspin\DERmspin%


\printedpage{24}
% head: A. Grothendieck : Dérivateurs, Ch. XII, Édition des Universités de Montpellier II et Paris VII
% foot: Version 7 décembre 2002
En effet, si $u'$ est dans $x\backslash X$ et $v'$ dans $y\backslash Y$, et si $x'$, $y'$ sont leurs images dans $X$, $Y$ respectivement, alors
\[
u'\backslash(x\backslash X)\simeq x'\backslash X,\qquad v'\backslash(y\backslash Y)\simeq y'\backslash Y,
\]
et la flèche entre les deux s'identifie à la flèche de localisation
\[
x'\backslash X\longrightarrow y'\backslash Y.
\]

\textbf{Proposition 5.} \textit{Soit $f : X\longrightarrow Y$ dans $\Cat$, supposons $f$ $W$-propre ou $W$-lisse, et à fibres $W$-asphériques, alors $f\in W$ (et même $f$ est} \textup{universellement} \textit{dans $W$, i.e. il est dans $W$ et le reste par tout changement de base).}

\leavevmode\msp{page 30}Comme les hypothèses sont stables par changement de base, il suffit de prouver $f\in W$. Supposons p.~ex. $f$ propre, prouvons que $f$ est $W$-\textit{asphérique}, i.e. que les $X/y$ ($y\in\Ob\,Y$) sont $W$-asphériques. Mais $X_y\longrightarrow X/y$, étant $W$-coasphérique par hypothèse de propreté, est dans $W$ (par W(2 bis)), et comme $X_y$ est $W$-asphérique par hypothèse, $X/y$ aussi. Mais par W(2 bis), $f$ $W$-asphérique $\Longrightarrow f\in W$. Le cas $f$ lisse est dual.

\textbf{Remarque.} Ainsi, si $f$ est propre, s'il est à fibres $W$-asphériques il est `universellement $W$-asphérique', i.e. est et reste $W$-asphérique par tout changement de base. Inversement, si $f$ est universellement dans $W$ (pas forcément propre), alors en particulier ses fibres sont $W$-asphériques. Mais ça n'implique bien sûr pas que $f$ soit propre, ou lisse (\footnote[44]{Mais `universellement dans $W$' implique déjà `universellement $W$-asphérique', puisque la notion de $W$-asphéricité est définie par \texttt{[la]} proposition 1 a).}).

\textbf{Corollaire.} \textit{Si $W$ satisfait} \textup{W(2) (pas seulement W(2 bis))}\textit{, alors pour un $S$-morphisme $X\longrightarrow Y$ de catégories qui sont toutes deux $W$-propres sur $S$, ou toutes deux $W$-lisses, si $f_s : X_s\longrightarrow Y_s$ est dans $W$ pour tout $s\in S$, alors $f\in W$.}

\leavevmode\msp{page 31}Supposons p.~ex. $f$ propre, considérons, pour $s\in\Ob\,S$
\[
\begin{tikzcd}[row sep=large, column sep=large]
X_s \arrow[r, "f_s"] \arrow[d, "i_s^X"'] & Y_s \arrow[d, "i_s^Y"] \\
X/s \arrow[r, "f/s"] & Y/s,
\end{tikzcd}
\]
comme $X$, $Y$ \texttt{[sont]} propres, les flèches verticales sont dans $W$ (il suffirait de beaucoup moins que la propreté \dots), donc $f_s\in W\Longrightarrow f/s\in W$. Ainsi, `$f$ est dans $W$ localement sur $S$', ce qui implique $f\in W$ par W(2).

\printedpage{25}
% head: A. Grothendieck : Dérivateurs, Ch. XII, Édition des Universités de Montpellier II et Paris VII
% foot: Version 7 décembre 2002
\textbf{Proposition 6.} \textit{Soit}
\[
\begin{tikzcd}[row sep=large, column sep=large]
X & X' \arrow[l, "q"'] \arrow[d, "f'"] \\
S \arrow[from=u, "f"'] & S', \arrow[l, "p"]
\end{tikzcd}
\]
\textit{cartésien dans} $\Cat$. \textit{Si $f$} \texttt{[est]} \textit{$W$-asphérique, $p$ $W$-lisse, alors $f'$ est $W$-asphérique.}

Il faut prouver \texttt{[que]} $X'\times_{S'}S'/s'$ est $W$-asphérique ($s'\in\Ob\,S'$). Soit $s=p(s')$, on a
\[
\begin{tikzcd}[column sep=small]
S & & S' \arrow[ll, "p"'] & \\
& S/s \arrow[ul] & & S'/s', \arrow[ul] \arrow[ll, "p_0"]
\end{tikzcd}
\]
avec $p_0$ $W$-lisse (lemme p.~28 \texttt{[plutôt 29]}) et à fibres $W$-asphériques, donc une $W$-équivalence universelle (prop.~5). Comme l'image inverse de $X$ sur $S/s$ est $W$-asphérique (hypothèse pour $f$), il \texttt{[s'ensuit]}

\leavevmode\msp{page 32}que dans le carré cartésien ci-contre \texttt{[ci-dessous]}
\[
(*)\qquad\qquad
\begin{tikzcd}[row sep=large, column sep=large]
X/s \arrow[d] & X'/s' \arrow[l, "W"'] \arrow[d] \\
S/s & S'/s' \arrow[l, "W"']
\end{tikzcd}
\]
les flèches horizontales sont dans $W$, et la première flèche verticale l'est par hypothèse sur $f$, donc la seconde aussi, donc $X'/s'$ \texttt{[est]} $W$-asphérique.

\textbf{Corollaire} (généralise \texttt{[la]} prop.~6, obtenue en faisant $Y=S$). \textit{Soit $p : S'\longrightarrow S$} \texttt{[un]} \textit{morphisme $W$-lisse dans} $\Cat$. \textit{Alors le foncteur changement de base}
\[
\Cat/S\longrightarrow\Cat/S'
\]
\textit{est compatible avec les localiseurs $W^{\mathrm{g}}_S$, $W^{\mathrm{g}}_{S'}$, i.e. transforme $f : X\longrightarrow Y$ qui est dans $W$} localement \textit{sur $S$, dans $f' : X'\longrightarrow Y'$ dans $W$ localement sur $S'$.}

Cette fois, on regarde le diagramme similaire à $(*)$
\[
\begin{tikzcd}[column sep=small]
X/s \arrow[dd] \arrow[dr, "f/s"] & & X'/s' \arrow[ll, "W"'] \arrow[dd] \arrow[dr, "f'/s'"] & \\
& Y/s \arrow[dl] & & Y'/s' \arrow[ll, "W"', crossing over] \arrow[dl] \\
S/s & & S'/s', \arrow[ll, "W"'] &
\end{tikzcd}
\]

\printedpage{26}
% head: A. Grothendieck : Dérivateurs, Ch. XII, Édition des Universités de Montpellier II et Paris VII
% foot: Version 7 décembre 2002
où les flèches horizontales sont dans $W$, donc $f/s\in W$ (hypothèse) implique $f'/s'\in W$.

\leavevmode\msp{page 33}Dualement, si $p$ est $W$-propre, le foncteur image inverse est compatible avec les localiseurs $W^{\mathrm{d}}_S$, $W^{\mathrm{d}}_{S'}$.

Si $f : X\longrightarrow S$ est $W$-propre, on va définir un foncteur canonique
\[
\begin{aligned}
\hot_*^{X/S,W} : S &\quad\longrightarrow\quad \Hot_W\overset{\text{déf}}{=}(\Cat)W^{-1}\\
\text{donné sur les objets par}\qquad s &\quad\longmapsto\quad \hot_W(X_s).
\end{aligned}
\]
Il faut définir, pour une flèche $u$ dans $S$, $u : s\longrightarrow s'$, une flèche
\[
\hot_W(X_s)\longrightarrow\hot_W(X_{s'})
\]
dans $\Hot_W$ - flèche que nous pouvons noter $\hot_*^{X/S,W}(u)$. On la définit par le diagramme
\[
\begin{tikzcd}[row sep=large, column sep=large]
X_s \arrow[r, dotted] \arrow[d, "i_s^X"] & X_{s'} \arrow[d, "i_{s'}^X"] \\
X/s \arrow[r, "X/u"] & X/s',
\end{tikzcd}
\]
où les flèches verticales sont $W$-coasphériques, donc dans $W$. Ce sont donc des isomorphismes dans $\Hot_W$, d'où l'unique flèche en pointillés dans $\Hot_W$, rendant le carré commutatif. Sur cette définition il est immédiat que
\[
u\longmapsto\hot_W(u)
\]
est compatible avec composition, et unités.

Ce foncteur est donc défini en supposant

\leavevmode\msp{page 34}seulement que les
\[
X_y\longrightarrow X/y
\]
sont dans $W$, ce qui est nettement moins fort que propre. Si on suppose que cette hypothèse est satisfaite `universellement' (i.e. reste satisfaite après tout changement de base $S'\longrightarrow S$) (\footnote[45]{On dira que $f$ est $W$-\textit{prépropre.}}), alors on a donc, pour tout $S'$ dans $\Cat/S$, un foncteur
\[
\hot_*^{X'/S',W} : S'\longrightarrow\Hot_W,
\]
et on vérifie tout de suite que ce n'est autre que le composé dans

\printedpage{27}
% head: A. Grothendieck : Dérivateurs, Ch. XII, Édition des Universités de Montpellier II et Paris VII
% foot: Version 7 décembre 2002
\[
(*)\qquad\qquad
\begin{tikzcd}[column sep=small]
S' \arrow[rr, "p"] \arrow[dr, "\hot_W^{X'/S'}"'] & & S \arrow[dl, "\hot_W^{X/S}"] \\
& \Hot_W &
\end{tikzcd}
\]
\texttt{[plutôt} $\hot_*^{X'/S',W}$ \texttt{resp.} $\hot_*^{X/S,W}$\texttt{]}. On a des observations duales si $X$ est lisse sur $S$, plus généralement si les inclusions
\[
X_s\longrightarrow s\backslash X
\]
sont dans $W$. On trouve alors un contrafoncteur
\[
\left\{
\begin{aligned}
S^\circ &\quad\longrightarrow\quad \Hot_W\\
s &\quad\longmapsto\quad \hot_W(X_s)\\
u &\quad\longmapsto\quad \hot^*_{X/S,W}(u)\ \text{ou}\ \hot^*_{X/S}(u)
\end{aligned}
\right.
\]

\leavevmode\msp{page 35}et une compatibilité similaire à $(*)$, si l'hypothèse est satisfaite \textit{universellement} sur $S$, p.~ex. $f$ $W$-lisse.

\textbf{Proposition 7.} \textit{Supposons $f : X\longrightarrow S$ à la fois $W$-propre et $W$-lisse, plus généralement que les morphismes}
\[
X_s\longrightarrow X/s,\qquad X_s\longrightarrow s\backslash X\qquad(s\in S)
\]
\textit{soient dans $W$, et que cette hypothèse reste verifiée après tout changement de base. Alors les foncteurs correspondants} % [sic] verifiée -> vérifiée, accent missing
\[
\begin{aligned}
\lambda &\;=\;\hot_*^{X/S,W} : &&\quad S \quad\longrightarrow\quad \Hot_W\\
\lambda^\circ &\;=\;\hot^*_{X/S,W} : &&\quad S^\circ \quad\longrightarrow\quad \Hot_W
\end{aligned}
\]
\textit{transforment toute flèche $u$ de $S$ en} \texttt{[un]} \textit{isomorphisme de} $\Hot_W$\textit{, et on a}
\[
\lambda^\circ(u)=\lambda(u)^{-1}\qquad\textit{dans}\quad\Hot_W,
\]
\textit{i.e. $\lambda(u)$ et $\lambda^c irc(u)$ sont inverses l'un de l'autre.} % [sic] λ^c irc(u) -> λ°(u), command \circ misprinted as superscript c followed by "irc"

\textsc{Démonstration}. La donnée de $u : s_0\longrightarrow s_1$ équivaut à celle de
\[
S'=\boldsymbol{\Delta}^1\longrightarrow S,
\]
d'où
\[
X'=X\times_S S'\qquad\text{sur } S'=\boldsymbol{\Delta}^1,
\]
avec les deux fibres
\[
X'_0\simeq X_{s_0},\qquad X'_1\simeq X_{s_1}.
\]
L'hypothèse assurant l'existence de $\lambda$ s'exprime par la condition que
\[
X'_1\longrightarrow X'/1\simeq X'
\]
\printedpage{28}
% head: A. Grothendieck : Dérivateurs, Ch. XII, Édition des Universités de Montpellier II et Paris VII
% foot: Version 7 décembre 2002
\msp{page 36}est dans $W$, de sorte que le diagramme
\[
X'_0\xrightarrow{i_0}X'\xleftarrow[W]{i_1}X'_1
\]
définit une flèche $\lambda(u) : X'_0\longrightarrow X'_1$ dans $\Hot_W$, savoir $\lambda(u)=i_1^{-1}i_0$. L'hypothèse qui assure l'existence de $\lambda^\circ$ s'exprime par la condition symétrique $i_0\in W$, qui permet de définir la flèche $\lambda^\circ(u) : X'_1\longrightarrow X'_0$ par $\lambda^\circ(u)=i_0^{-1}i_1$. Quand les deux conditions sont satisfaites simultanément, on voit donc que $\lambda(u)$, $\lambda^\circ(u)$ sont inverses l'un de l'autre.

\textbf{Remarques.} Pour définir $\lambda$ `universellement', il n'était pas indispensable de considérer les $X/s$ et les inclusions $X_s\longrightarrow X/s$, au lieu de cela il suffisait de regarder les changements de base $S'=\boldsymbol{\Delta}^1\xrightarrow{u}S$, et de faire l'hypothèse que pour tout tel changement de base, l'inclusion $X'_1\xrightarrow{i_1}X'=X_{(\boldsymbol{\Delta}^1,u)}$ est dans $W$. Cette hypothèse est automatiquement stable par changement de base arbitraire. (La $W$-propreté de $f : X\longrightarrow S$ signifie que $X'_1\longrightarrow X'$ est non seulement dans $W$, mais même $W$-coasphérique.) Dualement, pour la définition `universelle' de $\lambda^\circ$, il suffirait que pour tout changement de base $\boldsymbol{\Delta}^1\longrightarrow S$ comme dessus, l'inclusion $X'_0\xrightarrow{i_0}X'$ est dans $W$.

\leavevmode\msp{page 37}La notion la plus faible raisonnable de $W$-fibré, me semble être celle où on suppose les deux choses à la fois. Les hypothèses qu'on a faites jusqu'à présent, à l'exclusion de la forme forte de $W$-propreté ou de $W$-lissité, reviennent toujours à supposer que le \textit{changement de base par $f : X\longrightarrow S$} transforme certaines flèches dans $\Cat/S$ en \texttt{[élément de]} $W$ : Pour définir $\lambda$, ces flèches sont
\[
\text{soit } \{s\}\longrightarrow S/s,\qquad \text{soit } \{1\}\longrightarrow\boldsymbol{\Delta}^1(\longrightarrow S),
\]
pour définir $\lambda^\circ$, ce sont
\[
\text{soit } \{s\}\longrightarrow s\backslash S,\qquad \text{soit } \{0\}\longrightarrow\boldsymbol{\Delta}^1(\longrightarrow S),
\]
le `soit' suivant la définition qu'on a adoptée. Dans ces exemples, il s'agit toujours de l'inclusion de l'objet final d'une catégorie $S'$, dans ladite, ou inversement de l'objet initial. Cela montre d'ailleurs (si on exige l'universalité), que la première hypothèse (assurant $\lambda$) est \textit{exactement} cela (cas de l'inclusion d'un objet final $s'$ - puisque pour un tel objet, on a $S'/s'\simeq S'$), tandis que la deuxième hypothèse exige moins, seulement pour le cas de $\boldsymbol{\Delta}^1$. Et dualement pour $\lambda'$ \texttt{[plutôt $\lambda^\circ$]}. Mais il convient de noter que pratiquement, la façon la plus commode, et de loin,

\leavevmode\msp{page 38}pour vérifier de telles propriétés, c'est soit par la $W$-propreté de $f$, soit par la $W$-lissité - ce qui ramène à vérifier que certaines catégories bien précisées sont $W$-asphériques, sans avoir à se soucier de changements de base et de vérifier que les flèches construites en termes de changements de base arbitraires sont dans $W$. Et l'emploi de la $W$-propreté et de la $W$-lissité est d'autant plus sympathique, que ces notions sont stables par composition,


\endgroup
% END UNIT Der-ChXIIa-012/Der-ChXIIa-012_p026-030.tex
% BEGIN UNIT Der-ChXIIa-012/Der-ChXIIa-012_p031-035.tex
\begingroup\makeatletter%
\derleqno%
\newcommand{\printedpage}[1]{\marginpar{\footnotesize p.~#1}}%
\newcommand{\msp}[1]{\marginpar{\footnotesize [#1]}}%
\newcommand{\fl}{\longrightarrow}%
\newcommand{\uHom}{\underline{\mathrm{Hom}}}%
\newcommand{\uFib}{\underline{\mathrm{Fib}}}%
\newcommand{\Ob}{\operatorname{Ob}}%
\newcommand{\HOT}{\mathrm{HOT}}%
\newcommand{\Cat}{\mathrm{Cat}}%
\makeatother\let\printedpage\derpp\let\msp\dermsp%
\let\mspin\DERmspin%


\printedpage{29}
% head: A. Grothendieck : Dérivateurs, Ch. XII, Édition des Universités de Montpellier II et Paris VII
alors que toutes les notions dérivées, ainsi que les notions de fibrations à la clef, ont une fâcheuse tendance à ne \textit{pas} être stable par composition.

On pourrait appeler pseudo-$W$-propre resp. quasi-$W$-propre (un peu mieux que pseudo) ce qu'on vient de dégager - j'entends des ricanements! Et si on veut les deux à la fois - $\lambda$ et $\lambda^{\circ}$ associés comme on a dit? Je vois le même principe de description avec des $S''\fl S'$ sur $S$, d'exiger que $X''\fl X'$ soit transformé en des \texttt{[éléments de]} $W$. Avec les choix suivants qui ont l'air sympathique

\begin{itemize}
\item[a$_1$)] \msp{page 39}$Y'$ a un objet final, $Y''=e$. La définition revient donc à la condition que les
\[ X_s\fl X/s, \qquad X/s'\fl X/s \]
(pour $s\in\Ob S$ où $u : s'\fl s$ dans $\mathrm{S}$) soient dans $W$, et que ça reste comme ça après tout changement de base $S'\fl S$.

\item[a$_2$)] $Y'$ a un objet initial, $Y''=e$, c'est la condition duale à coups de
\[ X_s\fl s\backslash X, \qquad s\backslash X\fl s'\backslash X. \]

\item[a$_0$)] $Y'=\boldsymbol{\Delta}^1$, $Y''=e$ (\footnote[46]{NB a$_1$), a$_2$) sont notés pareil p. 4, mais a$_0$) a été oublié.}).

\item[b)] $Y'$ est contractile, $Y''$ aussi (ou au choix, $Y'$ contractile, $Y''=e$)
\end{itemize}

donc la condition revient à ceci : pour toute $Y'$ sur $Y$ \textit{contractile} (notion indépendante de $W$, c'est ça son principal avantage) les inclusions des fibres de $X'=X\times_Y Y'$ sur $Y'$ dans $X'$ sont \texttt{[dans]} $W$.

Je sens que même pour des $W$ générales, b) reste du domaine du vérifiable, tandis qu'on mollit si on passe à l'extrême en généralité dans ce sens :

\begin{itemize}
\item[b$'$)] Comme b), mais en supposant seulement $Y'$ $W$-asphérique (\footnote[47]{C'est essentiellement le a$'$) de p. 4}).
\end{itemize}

Enfin (la condition a) de la page 4)

\begin{itemize}
\item[c)] Si $Y''\fl Y'$ est \texttt{[dans]} $W$ sans plus, $X''\fl X'$ aussi.
\end{itemize}
% foot: Version 7 décembre 2002

\printedpage{30}
% head: A. Grothendieck : Dérivateurs, Ch. XII, Édition des Universités de Montpellier II et Paris VII
\msp{page 40}La seule de ces notions qui ait la moindre envie d'être stable par composition, c'est la plus forte c). Par contre, la notion de flèche $W$-parfaite (cf. déf. 3) est stable spontanément par composition :

\textbf{Définition 3.} On dit que $f : X\fl S$ est $W$-\textit{parfaite} si elle est à la fois $W$-propre et $W$-lisse.

Dans le cas $W=W_\infty$, toutes les variantes de la page précédente, à l'exclusion seulement de a$_0$), sont équivalentes. Ce serait intéressant de dégager des conditions minimales sur $W$ pour assurer qu'il en soit ainsi. (Et il se pose la question si dans le cas de $W_\infty$, a$_0$) suffit pour impliquer le reste - je ne m'y attends guère \dots)

Dans le cas général, on a le diagramme d'implications (et encore je me suis limité)
\[
\begin{tikzcd}[column sep=small]
 & & \mathrm{c} \arrow[d, Rightarrow] & & \text{$W$-fibration ultraforte} \\
 & & \mathrm{b}' \arrow[dd, Rightarrow] & & \text{$W$-fibration $W$-forte} \\
 & \text{$W$-parfait} \arrow[dr, Rightarrow] & & & \\
 & & \mathrm{b} \arrow[dl, Rightarrow] \arrow[dr, Rightarrow] & & \text{$W$-fibration hot-forte (\footnotemark[48])} \\
\text{$W$-fibration gauche} & \mathrm{a}_1 \arrow[dr, Rightarrow] & & \mathrm{a}_2 \arrow[dl, Rightarrow] & \text{$W$-fibration droite} \\
 & & \mathrm{a}_0 & & \text{$W$-pseudofibration}
\end{tikzcd}
\]
\footnotetext[48]{Il semblerait que ce soit ça, finalement, la ‘bonne notion’ de $W$-fibration, \textit{équivalente} à a$_1$ et a$_2$, lesquelles seraient équivalentes entre elles. Cf. prop. 9, p. 48.}

\texttt{[‘$W$-parfait $\Longrightarrow$ b’ encerclé]}.

\msp{page 41}\textbf{Proposition 8.} \textit{Soit $f : X\fl S$ dans} Cat. \textit{Les conditions suivantes sont équivalentes.}

\begin{itemize}
\item[\textit{a)}] \textit{$f$ est une $W$-fibration gauche, i.e. la condition} a$_1$ \textit{de la page 39 est vérifiée.}

\item[\textit{b)}] \textit{$f$ est $W$-prépropre, et $\forall\; s_0\fl s_1$ dans $S$, $X/s_0\fl X/s_1$ est dans $W$.}

\item[\textit{c)}] \textit{$f$ est $W$-prépropre, et une $W$-pseudofibration, i.e. la condition a$_0$ de la page 39 est vérifiée.} % AUDIT resolved: CMTI12 text-italic a with CMR8 subscript; see corrections.csv, PDF leaf 40, ms41
\end{itemize}

(\footnote[49]{\textbf{Corollaire.} \textit{Quand $f$ est} Cat-\textit{cofibrant, alors ces conditions équivalent aussi à ceci : $\forall\; u : s_0\fl s_1$ dans $S$, $u_* : X_{s_0}\fl X_{s_1}$ est dans $W$.}}).
% foot: Version 7 décembre 2002

\printedpage{31}
% head: A. Grothendieck : Dérivateurs, Ch. XII, Édition des Universités de Montpellier II et Paris VII
\textbf{NB} Ces conditions sont satisfaites si $f$ est en même temps $W$-prélisse, cf. prop. 7.

\textsc{Démonstration}. On a a) $\Longrightarrow$ b) trivialement (faire $Y''=Y/s$, $Y''=Y/s'$ \texttt{[lien avec l'assertion? Cf. p. 4.]}).

b) $\Longleftrightarrow$ c) par le raisonnement de la proposition 7, car en termes de $X'=X_u=X\times_S(\boldsymbol{\Delta}^1,u)$, la condition b) pour $u : s_0\fl s_1$ donné s'interprète dans le diagramme commutatif
\[
\begin{tikzcd}[column sep=small, row sep=small]
X'_0 \arrow[dr, "i'_0\in W"] \arrow[dd, "\wr"] & & X'_1 \arrow[dr, "i'_1\in W"] \arrow[dd, "\wr"] & \\
 & \underset{(\simeq X'_0)}{X'/0} \arrow[rr, "\varphi'", crossing over] \arrow[dd, "p_0"] & & X'/1\ (\simeq X') \arrow[dd, "p_1"] \\
X_{s_0} \arrow[dr, "i_0\in W"'] & & X_{s_1} \arrow[dr, "i_1\in W"] & \\
 & X/s_0 \arrow[rr, "\varphi"] & & X/s_1
\end{tikzcd}
\]

\texttt{[il y a quelques $W$ redondants dans le diagramme original]} où $i_0$, $i_1$, $i'_0$, $i'_1$ sont dans $W$ à cause de l'hypothèse de $W$-prépropreté (stable par changement de base).

Donc $p_0$, $p_1$ sont \texttt{[dans]} $W$, donc $\varphi\in W$ si et seulement si $\varphi'\in W$, or $\varphi'\in W$ est la condition a$_0$.

\msp{page 42}Prouvons enfin que b), c) impliquent a). Comme c) (donc aussi b)) est stable par changement de base, il suffit de montrer ceci : Si $S$ a un objet final $e$, alors pour tout $s\in\Ob S$, l'inclusion $X_s\xrightarrow{\ i\ }X$ est dans $W$. Or on a le diagramme commutatif
\[
\begin{tikzcd}
X_s \arrow[r, "i_s\in W"] \arrow[dr, "i"'] & X/s \arrow[d, "u\,\in\,W\text{ par c) }\texttt{[plutôt par b)]}"] \\
 & X/e=X,
\end{tikzcd}
\]
qui montre bien $i=ui_s\in W$.

Dualement, si $f$ est $W$-prélisse, alors la condition a$_0$ ($f$ une $W$-pseudofibration) est équivalente à a$_2$ ($f$ une $W$-fibration droite), et signifie aussi que les
\[ s_1\backslash X\fl s_0\backslash X \]
sont dans $W$. Mais la notion de ‘$W$-fibration’ sur laquelle on tombe (si on ne se contente du minable a$_0$!) est différente dans les deux cas duaux l'un de l'autre. On aimerait dans les deux cas trouver des $W$-fibrations ‘hot-fortes’ (condition b), tout au moins en mettant le paquet (s'il le faut) - i.e. pour $f$ $W$-parfait. Ça revient donc à ceci.

\textbf{Question :} \textit{Soient $X\xrightarrow{\ f\ }S$, avec $f$ $W$-parfait, $S$ contractile, $s\in\Ob S$. L'inclusion $X_s\fl X$, est-elle dans $W$?}
% foot: Version 7 décembre 2002

\printedpage{32}
% head: A. Grothendieck : Dérivateurs, Ch. XII, Édition des Universités de Montpellier II et Paris VII
\msp{page 43}On va supposer seulement que $f$ est $W$-lisse et une $W$-fibration gauche \texttt{[plutôt que $f$ est simplement une $W$-fibration à droite, i.e que $f$ est $W$-prélisse et une $W$-pseudofibration]} (cf. prop. 8). On va se ramener au cas où $X$ est même fibrant sur $S$, en introduisant $\Phi_S(X)$ \texttt{[catégorie fibrée définie par le foncteur $S^{\circ}\fl\Cat$, $s\longmapsto s\backslash X$]} et % [sic] i.e que -> i.e. que, missing period of the abbreviation
\[ X\xrightarrow{\ i_X\ }\Phi_S(X). \]

L'ennui, c'est que j'ai fait la théorie de $\HOT_W$ jusqu'à maintenant avec l'hypothèse (qui s'avère idiote) que $W\supseteq W_\infty$ - je serai donc obligé de revoir rapidement sous quelles hypothèses plus raisonnables sur $W$ les résultats familiers sont valables. Par exemple j'aimerais que $i_X\in W$ (\footnote[50]{en fait, $i_X\in(W_\omega)^{\mathrm{d}}_S\subseteq W^{\mathrm{d}}_S\subseteq W_S$.}), de sorte qu'on ait
\[
\begin{tikzcd}[column sep=large, row sep=large]
X_s \arrow[r, "i_s\in W", "\text{car $X$ $W$-lisse}"'] \arrow[d, "j^X_s"'] & \Phi_S(X)_s=s\backslash X \arrow[d, "j_s^{\Phi_S(X)}"] \\
X \arrow[r, "i_X\in W"] & \Phi_S(X),
\end{tikzcd}
\]
on voit donc que pour prouver que l'inclusion $j^X_s$ est \texttt{[dans]} $W$, il suffit de prouver la même chose pour $\Phi_S(X)$, qui est fibrant (pas seulement $W$-lisse). De plus, l'hypothèse pseudo-fibration sur $X$ passe sûrement à $\Phi_S(X)$ (je réserve les deux vérifications à plus tard (\footnote[51]{O.K.})). On est donc ramené au cas de
\[ X\in\Ob\uFib_0S, \qquad S\text{ contractile,} \]
à prouver que l'inclusion
\[ X_s\fl X \]
\msp{page 44}est dans $W$. Je vais prouver ceci :

\textbf{Lemme.} \textit{Soient $S$ dans} Cat, \textit{$X\in\Ob\uFib_0S$, et considérons, pour un $S'$ dans} Cat,
\[
\begin{aligned}
f &\quad\longmapsto\quad f^*(X/S)=X\times_S(S',f)\\
\uHom(S',S)^{\circ} &\quad\longrightarrow\quad \uFib_0S'.
\end{aligned}
\]
\textit{(Ce foncteur est défini grâce au fait que $X$ est fibrant sur $S$, et peut s'interpréter donc, à $S$-équivalence près, en termes d'un foncteur $S^{\circ}\fl$} Cat \textit{\dots\ (\footnote[52]{\textbf{NB} $\Phi_S(X)$ était même déjà scindée, définie par le foncteur $s\longmapsto s\backslash X$.})). Les flèches de $\uHom(S',S)$ sont transformées en flèches $\in W^{\mathrm{d}}_{S'}$.}
% foot: Version 7 décembre 2002

\printedpage{33}
% head: A. Grothendieck : Dérivateurs, Ch. XII, Édition des Universités de Montpellier II et Paris VII
En effet, en termes de foncteurs
\[ S'^{\circ}\fl\Cat, \]
si $X$ est définie par
\[ S^{\circ}\xrightarrow{\ \mathscr{X}\ }\Cat, \]
alors $f^*(X)$ est définie par $\mathscr{X}\circ f$, et la flèche
\[ \mathscr{X}\circ f\fl\mathscr{X}\circ g \]
définie par $u : f\fl g$ n'est autre que $\mathscr{X}*u$. Sa valeur en $s'\in\Ob S'$ est obtenue en prenant
\[ u(s') : f(s')\fl g(s'), \]
d'où par $\mathscr{X}$
\[ \mathscr{X}(u(s')) : \underbrace{\mathscr{X}(g(s'))}_{(\mathscr{X}\circ g)(s')}\fl\underbrace{\mathscr{X}(f(s'))}_{(\mathscr{X}\circ f)(s')} \]
(\footnote[53]{à dégager un énoncé à part.}). Mais l'hypothèse $X\in\uFib_0S$ signifie que les $\mathscr{X}(v)$, $v\in\mathrm{Fl}(S)$, sont tous dans $W$. Donc
\msp{page 45}$\mathscr{X}$ est dans $W$ fibre par fibre, donc dans $W^{\mathrm{d}}_{S'}$.

\textbf{Corollaire 1.} \textit{Le foncteur du lemme, composé avec}
\[ \uFib_0S'\fl\HOT^{\mathrm{lc}}_W(S')\overset{\text{déf}}{=}(\uFib_0S')(W^{\mathrm{d}}_{S'})^{-1} \]
\textit{se factorise par le groupoïde fondamental de} $\uHom$ \textit{:}
\[
\begin{tikzcd}[row sep=huge]
\uHom(S',S)^{\circ} \arrow[r] \arrow[d] & \uFib_0S' \arrow[d] \\
\textstyle\prod(\uHom(S',S))^{\circ} \arrow[r] & \HOT^{\mathrm{lc}}_W(S').
\end{tikzcd}
\]
\textit{Si $f$, $g$ sont dans une même composante connexe de $\uHom(S',S)$, alors $f^*X$ et $g^*X$ sont isomorphes (non canoniquement, en général) dans} $\HOT^{\mathrm{lc}}_W(S')$.

Il est plus astucieux de ne pas fixer $X$ et de considérer
\[ \uHom(S',S)\fl\uHom(\HOT^{\mathrm{lc}}_W(S),\HOT^{\mathrm{lc}}_W(S')) \tag*{$(*)$} \]
correspondant à
\[
\begin{aligned}
\uHom(S',S) &\quad\times\quad \HOT^{\mathrm{lc}}_W(S) &&\quad\longrightarrow\quad &&\HOT^{\mathrm{lc}}_W(S')\\
(f &\quad,\quad\ \xi) &&\quad\longmapsto\quad && f^*(\xi).
\end{aligned}
\]
Le lemme dira que les flèches de $S'\fl S$ sont transformées par $(*)$ en isomorphismes, d'où
\[ \prod\uHom(S',S)\fl\uHom(\HOT^{\mathrm{lc}}_W(S),\HOT^{\mathrm{lc}}_W(S')), \]
% foot: Version 7 décembre 2002


\endgroup
% END UNIT Der-ChXIIa-012/Der-ChXIIa-012_p031-035.tex
% BEGIN UNIT Der-ChXIIa-012/Der-ChXIIa-012_p036-040.tex
\begingroup\makeatletter%
\newcommand{\printedpage}[1]{\marginpar{\footnotesize p.~#1}}%
\newcommand{\msp}[1]{\marginpar{\footnotesize [#1]}}%
\newcommand{\HOTlc}{\mathrm{HOT}^{\mathrm{lc}}_W}%
\newcommand{\Hot}{\mathrm{Hot}}%
\newcommand{\Hom}{\mathrm{Hom}}%
\newcommand{\hot}{\mathrm{hot}}%
\newcommand{\id}{\mathrm{id}}%
\newcommand{\Wd}{W^{\mathrm{d}}_S}%
\newcommand{\Wf}{W^{\mathrm{f}}_S}%
\makeatother\let\printedpage\derpp\let\msp\dermsp%
\let\mspin\DERmspin%


\printedpage{34}
% head: A. Grothendieck : Dérivateurs, Ch. XII, Édition des Universités de Montpellier II et Paris VII
% foot: Version 7 décembre 2002

\msp{page 46}
donc si $f,g : S' \rightrightarrows S$ sont deux foncteurs homotopes, alors
\[ f^* \simeq g^* : \HOTlc(S) \rightrightarrows \HOTlc(S'). \]

Il s'ensuit aussitôt ceci

\textbf{Corollaire 2.} \textit{Si $f$ est un homotopisme, $f^* : \HOTlc(S) \longrightarrow \HOTlc(S')$ est une équivalence de catégories.}

En effet, prenant $g : S \longrightarrow S'$ telle que
\[ fg \sim \id_S, \qquad gf \sim \id_{S'}, \]
on trouve que
\[ \begin{aligned}
g^*f^* \simeq (\id_S)^* &= \id_{\HOTlc(S)} \\
f^*g^* \simeq (\id_{S'})^* &= \id_{\HOTlc(S')},
\end{aligned} \]
donc on trouve que (quitte à ajuster ces isomorphismes fonctoriels) $f^*$ et $g^*$ sont des équivalences quasi-inverses l'une de l'autre.

J'ai l'impression (\footnote[54]{O.K.}) que tout ça ne doit pas utiliser plus que les axiomes standard sur $W$ comme localiseur fondamental, tout au plus W(2) au lieu de la forme faible W(2 bis).

Revenons à $X$ lisse sur $S$ contractile. Le corollaire 2 implique que
\[ \begin{aligned}
\Hot_W &\longrightarrow \HOTlc(S) \\
\xi &\longmapsto p^*_S(\xi),
\end{aligned} \]
où $p_S : S \longrightarrow e$, est une équivalence. Donc dans la catégorie $\HOTlc(S)$, $X$ est isomorphe \msp{page 47} à un Cat-fibré de la forme $Z \times S \xrightarrow{\mathrm{pr}_2} S$, avec $Z$ dans Cat. Or pour $X_1 = Z \times S \longrightarrow S$, il est immédiat que $X_{1\,s} \longrightarrow X_1$ est dans $W$, car un produit cartésien de flèches \texttt{[dans]} $W$ (savoir $\id_Z$ et $s \longrightarrow S$) est dans $W$. (En fait, l'inclusion est même un homotopisme!) Notons maintenant que si $X_0$, $X_1$ sont deux objets dans $\underline{\mathrm{Liss}}_0S$, tels qu'il existe une $S$-flèche $X_0 \longrightarrow X_1$ qui soit dans $\Wd$, alors $X_{0\,s} \longrightarrow X_0$ est dans $W$ si et seulement si $X_{1\,s} \longrightarrow X_1$ l'est. En effet, on a
\[ \begin{tikzcd}
X_{0\,s} \arrow[r, "\in W"] \arrow[d, "i_s^{X_0}"'] & X_{1\,s} \arrow[d, "i_s^{X_1}"] & \\
X_0 \arrow[r, "\in W", ""{name=A, below}] & X_1 & \\
 & & \text{\footnotesize\begin{tabular}{@{}l@{}}ici on utilise sans doute\\ W(2), non seulement W(2 bis),\end{tabular}}
\arrow[from=A, to=3-3, bend right]
\end{tikzcd} \]
\printedpage{35}
% head: A. Grothendieck : Dérivateurs, Ch. XII, Édition des Universités de Montpellier II et Paris VII
% foot: Version 7 décembre 2002
qui prouve le résultat.

Il est vrai que le fait que $X = X_0$ et $X_1 = Z \times S$ soient isomorphes dans $\HOTlc(S)$ n'implique pas, par lui-même, l'existence d'une \msp{page 48} chaîne d'isomorphismes provenant de flèches dans $\Wd$. Mais en fait, si dans l'énoncé du corollaire 2 ci-dessus on compare $X$ dans $\underline{\mathrm{Fib}}_0S$ à $g^*f^*(X)$, grâce à une homotopie $\id_S \sim fg$, on trouve bel et bien une suite de foncteurs $\varphi_i : S \longrightarrow S$, avec $\varphi_0 = \id_S$, $\varphi_N = fg$, et $\varphi_i$ et $\varphi_{i+1}$ reliés par une flèche dans $\underline{\mathrm{Hom}}(S,S)$, donc on a une suite d'objets $X_i = \varphi_i^*(X)$ de $\underline{\mathrm{Fib}}_0S$ (avec $X_0 = X$, $X_N = g^*f^*X$), et $X_i$ et $X_{i+1}$ reliés par une flèche \textit{dans} $\underline{\mathrm{Fib}}_0S$ (pas seulement dans $\HOTlc(S)$), laquelle flèche est dans $\Wd = \Wf$. Donc on gagne, sous réserve de vérifications de routine diverses (\footnote[55]{O.K.}). Donc modulo celles-ci, on aura

\textbf{Proposition 9} (\footnote[56]{C'est ok en supposant que $W$ est \textit{fortement} saturé, i.e. $\hot_W(f)$ isomorphisme $\Longrightarrow f \in W$.},\footnote[57]{Donc $W$-fibration à gauche $\Longleftrightarrow$ $W$-fibration à droite $\Longleftrightarrow$ $W$-fibration.}). \textit{Supposons que $X$ dans $\mathrm{Cat}/S$ soit une $W$-fibration à gauche (cf. prop. 8). Alors c'est même une $W$-fibration (au sens de b), page 40, où c'est appelé `$W$-fibration \textup{hot}-forte'), en d'autres termes, si $S'' \longrightarrow S'$ est sur $S$ avec $S''$, $S'$ contractiles, alors $X'' \longrightarrow X'$ est dans $W$ (\footnote[58]{on trouve seulement que $\hot_W(X'') \longrightarrow \hot_W(X')$ \texttt{[est un]} isomorphisme.}). On a sans doute mieux : si $S'' \longrightarrow S'$ est un homotopisme, alors $X'' \longrightarrow X'$ est dans $W$.}

\msp{page 49}
Prouvons (si nous pouvons) ce complément sympathique. Je précise que j'ai eu à faire l'hypothèse de $W$-prélissité (où de $W$-prépropreté) pour être sûr que l'image inverse $X \times_S S'$, par $f' : S' \longrightarrow S$, coïncide bien avec $f^*$ sur les $\HOTlc$ - autrement je n'ai aucune chance. % [sic] où -> ou, the conjunction "or" is meant
`On sait' (tout est à revoir!) que c'est O.K. dans les trois cas $X$ $W$-lisse, $X$ $W$-propre, $X$ une $W$-fibration (mais en quel sens???). La vérification à faire (quand on travaille avec les $\Wd$) est celle-ci :
\[ \begin{aligned}
&X \xrightarrow{i_X} \Phi_S(X) = Y \text{ est dans } \Wd \text{ (sans mérite), et} \\
&\text{le reste par tout changement de base.}
\end{aligned} \tag*{$(*)$} \]

Mais pour des $X, Y \in \mathrm{Ob}\,\mathrm{Cat}/S$ qui sont $W$-prélisse, la $\Wd$-équivalence coïncide avec $\Wf$, laquelle est préservée par changement de base, ainsi que l'hypothèse sur $X$, $Y$. Donc on est tranquille.

On est ramené à prouver que si $X$ est $W$-prélisse sur $S$ et une $W$-pseudofibration, et si $S' \longrightarrow S$ est un homotopisme, alors $X' \longrightarrow X$ est dans $W$. Par l'argument de tantôt utilisant $X \longrightarrow \Phi_S(X)$, on se ramène au cas où $X$ est fibrée scindée.

\printedpage{36}
% head: A. Grothendieck : Dérivateurs, Ch. XII, Édition des Universités de Montpellier II et Paris VII
% foot: Version 7 décembre 2002

\msp{page 50}
\[ \begin{tikzcd}
X' \arrow[rr] \arrow[d] & & X \arrow[d] \\
S' \arrow[rr, "f"] \arrow[dr, "p'"'] & & S \arrow[dl, "p"] \\
 & e &
\end{tikzcd} \]

La flèche $X' = f^*(X) \longrightarrow X$ s'interprète dans $\Hot_W$ comme
\[ p'_!(f^*(\xi)) \longrightarrow p_!(\xi) \ , \qquad \text{où } \xi = \hot_{S,W}(X). \tag*{$(*)$} \]

Je dis que cette flèche est \texttt{[un]} isomorphisme, i.e. que pour tout $\eta \in \Hot_W$, la flèche transposée
\[ \Hom_W(p_!(\xi),\eta) \longrightarrow \Hom_W(p'_!(f^*(\xi)),\eta) \tag*{$(**)$} \]
est \texttt{[une]} bijection. Admettant la formule d'adjonction, la flèche $(**)$ s'interprète comme
\[ \Hom_W(\xi,\underbrace{p^*(\eta)}_{\eta_S}) \longrightarrow \Hom_W(f^*(\xi),\underbrace{p'^*(\eta)}_{f^*(\eta_S)}), \]
i.e.
\[ \Hom_W(\xi,\eta_S) \longrightarrow \Hom_W(f^*(\xi),f^*(\eta_S)), \]
que ce soit \texttt{[un]} isomorphisme résulte donc de $f^*$ pleinement fidèle sur les $\HOTlc$. Or on a vu que $f^*$ est une équivalence (cor. 2, page 46).

\msp{page 51}
Je vais reprendre les propositions 9 et 8, un peu mal foutues.

\textbf{Théorème 3.} \textit{On suppose que $W$ satisfait à \textup{W(1,2,3)}. Soit $f : X \longrightarrow S$. Considérons les conditions suivantes.}

\begin{itemize}[leftmargin=4em]
\item[\textit{a)}] \textit{$\forall\ S'' \longrightarrow S'$ sur $S$, avec $S''$, $S'$} contractiles, \textit{$X'' \longrightarrow X'$ est dans $W$.}

\item[\textit{a}$'$\textit{)}] \textit{$\forall\ S'$ sur $S$, avec $S'$ contractile, et $s' \in S'$, l'inclusion} % AUDIT resolved: the label is a prime, not d prime; see corrections.csv, PDF leaf 46
\[ X_{s'} \hookrightarrow X' = X \times_S S' \]
\textit{est dans $W$.}

\item[\textit{a}$_1$\textit{)}] \textit{Comme a$'$), mais en supposant que $S'$ a un objet final (au lieu de $S'$ contractile).} % AUDIT resolved: the reference is a prime; see corrections.csv, PDF leaf 46

\item[\textit{a}$_2$\textit{)}] \textit{Comme a$'$), mais en supposant que $S'$ a un objet initial.} % AUDIT resolved: the reference is a prime; see corrections.csv, PDF leaf 46

\printedpage{37}
% head: A. Grothendieck : Dérivateurs, Ch. XII, Édition des Universités de Montpellier II et Paris VII
% foot: Version 7 décembre 2002

\item[\textit{b}$_1$\textit{)}] \textit{$f$ est $W$-}prépropre \textit{(i.e. a$'$) est satisfait quand $s'$ est un objet final de $S'$, sans autre condition sur $S'$ dans $\mathrm{Cat}/S$), et de plus pour toute flèche $s_0 \longrightarrow s_1$ dans $S$, la flèche induite $X/s_0 \longrightarrow X/s_1$ est dans $W$ (i.e. a$'$ \textup{\texttt{[plutôt a]}}) est satisfait pour $Y/s_0 \longrightarrow Y/s_1$ \textup{\texttt{[plutôt} $S/s_0 \longrightarrow S/s_1$\texttt{]}}).}

\item[\textit{b}$_2$\textit{)}] \textit{$f$ est $W$-}prélisse \textit{(i.e. a$'$) est satisfait quand $s'$ est un objet initial de $S'$, sans autre condition sur $S'$), et de plus les flèches $s_0\backslash X \longrightarrow s_1\backslash X$ \textup{\texttt{[plutôt} $s_1\backslash X \longrightarrow s_0\backslash X$\texttt{]}} (pour $s_0 \longrightarrow s_1$ dans $S$) sont dans $W$, (i.e. a$'$ \textup{\texttt{[plutôt a]}}) satisfait pour $s_1\backslash Y \longrightarrow s_0\backslash Y$ \textup{\texttt{[plutôt} $s_1\backslash S \longrightarrow s_0\backslash S$\texttt{]}}).}

\item[\textit{c}$_1$\textit{)}] \textit{$f$ est $W$-prépropre et une $W$-pseudofibration, i.e. la condition a$'$) est satisfaite dans les deux cas suivants: 1$^\circ$) $s'$ est un objet final de $S'$, 2$^\circ$) $S' = \boldsymbol{\Delta}^1$.}

\item[\textit{c}$_2$\textit{)}] \textit{$f$ est $W$-prélisse et une $W$-pseudofibration, i.e. a$'$) dans les deux cas: 1$^\circ$) $s'$ est objet initial de $S'$, 2$^\circ$) $S' = \boldsymbol{\Delta}^1$.}
\end{itemize}

\textit{Toutes ces 8 conditions sont équivalentes. Elles impliquent la suivante :}

\begin{itemize}[leftmargin=4em]
\item[\textit{d)}] \textit{Soit $S'' \xrightarrow{u} S'$ un homotopisme dans $\mathrm{Cat}/S$, et considérons $u_X : X'' \longrightarrow X'$. Alors $\hot_W(u_X)$ est \textup{\texttt{[un]}} isomorphisme. (Donc si $W$ est fortement saturé, $u_X \in W$. Donc sous cette condition sur $W$, d) équivaut aux conditions précédentes.)}
\end{itemize}

\msp{page 52}
\textbf{Définition 4.} On dit que $f : X \longrightarrow Y$ dans Cat est une \textit{$W$-fibration}, si elle satisfait les conditions équivalentes a) à c$_2$) de la proposition 9.

\textbf{Corollaire 1.} \textit{Si $f$ est $W$-propre ou $W$-lisse, pour que $f$ soit une $W$-fibration, il faut et il suffit que $f$ soit une $W$-pseudofibration, i.e. que pour $S' = \boldsymbol{\Delta}^1 \longrightarrow S$, les inclusions des deux fibres}
\[ X'_0 \xrightarrow{i_0} X' \xleftarrow{i_1} X'_1 \]
\textit{de $X' = X \times_S S'$ soient dans $W$.}

\textbf{Corollaire 2.} \textit{Supposons $f$ \textup{Cat}-fibrant. Alors $f$ est une $W$-fibration si et seulement si pout toute flèche $u : s \longrightarrow s'$ dans $S$, $u^* : X_{s'} \longrightarrow X_s$ est dans $W$. Énoncé dual si $f$ est \textup{Cat}-cofibrant.} % [sic] pout -> pour, misprint

\textbf{Corollaire 3.} \textit{Supposons $f$ $W$-prépropre et $W$-prélisse (par exemple $f$ $W$-parfait). Alors $f$ est une $W$-fibration.}

\textbf{Proposition 10.} \textit{Soit $f : X \longrightarrow Y$ une flèche dans $\underline{\mathrm{Fib}}_WS$ (\footnote[59]{$X$ et $Y$ tous deux prépropres ou tous deux prélisses suffit.}) (sous-catégorie pleine de $\mathrm{Cat}/S$ formée des $W$-fibrés sur $S$). Si $\forall\ s \in S$, $f_s : X_s \longrightarrow Y_s$ est dans $W$, alors $f \in W$.}

Considérons
\printedpage{38}
% head: A. Grothendieck : Dérivateurs, Ch. XII, Édition des Universités de Montpellier II et Paris VII
% foot: Version 7 décembre 2002
\[ \begin{tikzcd}
X \arrow[rr, "f"] \arrow[dr, "i_X"'] & & Y \arrow[dr, "i_Y"] & \\
 & \Phi_S(X) \arrow[rr, "\Phi_S(f)"'] & & \Phi_S(Y),
\end{tikzcd} \]
où $i_X, i_Y \in W$, donc $\Phi_S(f) \in W \Rightarrow f \in W$. Or le fait que $X$ \texttt{[est]} $W$-fibrée sur $S$ implique que $i_X \in \Wf$ (il suffit même que $X$ soit prépropre \texttt{[plutôt prélisse]} sur $S$), itou pour $Y$, $i_Y$. Donc l'hypothèse $f \in \Wf$ implique $\Phi_S(f) \in \Wf$, donc $\Phi_S(f) \in W$ par \texttt{[la]} proposition 5 (page 29).


\endgroup
% END UNIT Der-ChXIIa-012/Der-ChXIIa-012_p036-040.tex
% BEGIN UNIT Der-ChXIIa-012/Der-ChXIIa-012_p041-044.tex
\begingroup\makeatletter%
\NoAutoSpacing%
\newcommand{\printedpage}[1]{\marginpar{\footnotesize p.~#1}}%
\newcommand{\msp}[1]{\marginpar{\footnotesize [#1]}}%
\newcommand{\dashedbox}[1]{\tikz[baseline=(n.base)]\node[draw,dashed,inner sep=2pt](n){#1};}%
\newcommand{\Fib}{\mathrm{Fib}}%
\newcommand{\Parf}{\mathrm{Parf}}%
\newcommand{\HOT}{\mathrm{HOT}}%
\newcommand{\Hot}{\mathrm{Hot}}%
\newcommand{\hot}{\mathrm{hot}}%
\newcommand{\Hom}{\mathrm{Hom}}%
\makeatother\let\printedpage\derpp\let\msp\dermsp%
\let\mspin\DERmspin%


\printedpage{39}
% head: A. Grothendieck : Dérivateurs, Ch. XII, Édition des Universités de Montpellier II et Paris VII

\msp{page 53}

\section*{IV\quad Sur l'axiome W(4) et ses variantes.}

(\footnote[60]{Cf. aussi partie VI (HOT$_W$), p. 186-196}). J'ai passé des jours à essayer de prouver que toutes ces variantes sont équivalentes. Je vais donner une liste de celles, parmi ces variantes, qui avec le recul me paraissent encore avoir une utilité. (Avec des notations différentes de celles des pages 13-14.)

\begin{itemize}[leftmargin=5em,labelsep=0.5em]
\item[\fbox{W(4a)}] Toute $W$-fibration (cf. th. 3, p. 51, déf. 3, p. 52) est ultraforte, ce qui revient à ceci: pour
\[
\begin{tikzcd}
X & X' \arrow[l,"q"'] \arrow[d,"f'"] \\
Y \arrow[from=u,"f"'] & Y' \arrow[l,"p"']
\end{tikzcd}
\]
cartésien dans Cat, si $f \in \underline{\Fib}_W$, $p \in W$, alors $q \in W$.

\item[W(4a$'$)] Comme a), en supposant $f \in \Parf_\omega$, i.e. $f$ $W_\omega$-parfait (donc une $W_\omega$-fibration), où $W_\omega \subseteq W$ est la plus petits partie $\Sigma$ de Fl(Cat) satisfaisant W(1, 2 bis, 3). % [sic] petits -> petite, accord avec « partie »

\item[\fbox{W(4b)}] Soit un diagramme commutatif
$\begin{tikzcd}[ampersand replacement=\&, column sep=small, row sep=small, baseline=(X.base)]
|[alias=X]| X \arrow[rr,"f"] \arrow[dr,"p"'] \& \& Y \arrow[dl,"q"] \\
\& S \&
\end{tikzcd}$
dans Cat, avec $p, q \in \Fib_W$, $f \in W$.
Alors $\forall s \in S$, $f_s : X_s \longrightarrow Y_s$ \texttt{[est dans} $W$\texttt{]}.

\item[W(4b$'$)] Comme b), en supposant \texttt{[de plus]} $X$ et $Y$ fibrés scindés sur $S$, $f$ compatible avec \texttt{[les]} scindages (donc foncteur cartésien).

\item[W(4b$''$)] Comme b), en supposant $p$, $q$ $W$\textit{-parfaits.}

\item[\fbox{W(4c)}] Soit $f : X \longrightarrow Y$, avec $f \in \Fib_W \cap W$, alors les fibres de $f$ sont $W$-asphériques.

\item[\dashedbox{W(4c$'$)}] Comme c), en supposant $f \in \Parf_\omega \cap W$. (C'est l'ancien W(4b$''$), dépouillé du $W_\infty$-parasite.)
\end{itemize}

\msp{page 54}

J'ai cerclé les notions \texttt{[encadrées ici]} qui ont été vraiment utiles. Pratiquement, on peut parier que c'est la notion la plus forte, W(4a), qui est la seule importante. Voici les implications tautologiques (à l'exception de W(4a$'$) $\Longrightarrow$ W(4b), qui est démontré dans la démonstration de la proposition 4 (pages 6-10) (\footnote[61]{Cette démonstration utilise une notion plus forte de $W$-fibration que celle du théorème 3, p. 51.})).

\printedpage{40}
% head: A. Grothendieck : Dérivateurs, Ch. XII, Édition des Universités de Montpellier II et Paris VII

\[
\begin{tikzcd}[column sep=small]
& \mathrm{a} \arrow[d,Rightarrow] & \\
& \mathrm{a}' \arrow[dd,Rightarrow,"{\text{\rule[0.6ex]{2cm}{0.4pt}\ \begin{tabular}[t]{@{}l@{}}implication douteuse,\\ il faut affaiblir\\ un peu b), en\\ renforçant la\\ notion de\\ $W$-fibration\\ utilisée\\ pour énoncer b)\end{tabular}}}"] & \\
& & \\
& \mathrm{b} \arrow[dl,Rightarrow] \arrow[dr,Rightarrow] \arrow[dd,Rightarrow] & \\
\mathrm{b}' & & \mathrm{b}'' \\
& \mathrm{c} \arrow[d,Rightarrow] & \\
& \mathrm{c}' &
\end{tikzcd}
\]

J'ai surtout sué pour montrer que c$'$) $\Longrightarrow$ c), ou c) $\Longrightarrow$ b), sans y arriver. L'effort m'a permis de me débarrasser de diverses idées fausses, et de les remplacer par des justes, cf. plus bas. Il y a quand même un résultat positif qui surnage :

\textbf{Proposition 1.} \textit{La condition b) (et sans doute même b$'$) et b$''$)} (\footnote[62]{et en supposant dans le cas de b$''$, que $W$ satisfait la condition des limites inductives W(6 bis), cf. prop. 5 du `corollaire' de la, page 14.})\textit{, mais je ne vais pas le prouver ici \dots) implique sans doute a) pour $\bar{W}$, le saturé fort de $W$. En tous cas, si $W$ est fortement saturé, les conditions a), a$'$), b) (et sûrement aussi b$'$) et b$''$)) sont équivalentes.}

\textsc{Démonstration} de b) $\Longrightarrow$ a), quand $W$ est fortement saturé. En fait, moyennant b) (et sans même supposer $W$ fortement saturé), nous allons prouver que si on a
\[
\begin{tikzcd}
X & X' \arrow[l,"q"'] \arrow[d,"f'"] \\
Y \arrow[from=u,"f"'] & Y' \arrow[l,"p"']
\end{tikzcd}
\]
cartésien dans Cat, avec
\msp{page 55}
$f$ une $W$-fibration, et $p \in W$, alors $\underline{q \in \bar{W}}$, i.e. $\hot_W(q)$ \texttt{[est un]} isomorphisme. Pour ceci, nous interprétons la flèche $q$ dans $\Hot_W$ de la façon suivante. Soit
\[
\xi = \hot_{Y,W}(X/Y), \qquad \xi' = p^*(\xi),
\]
\printedpage{41}
% head: A. Grothendieck : Dérivateurs, Ch. XII, Édition des Universités de Montpellier II et Paris VII
qui dit que $q$ s'identifie à la flèche canonique
\[
H_\bullet^{\HOT_W}(Y',\xi') \longrightarrow H_\bullet^{\HOT_W}(Y,\xi) \qquad \text{dans } \Hot_W = \HOT_W(e), \tag*{$(*)$}
\]
où on note de plus que $\xi \in \HOT_W(Y)$ est un coefficient \textit{localement constant} sur $Y$.

Établissons d'abord cette interprétation de $q$. Notons que comme $X \in \underline{\Fib}_W Y$, l'image inverse $\xi' = p^*(\xi)$ n'est autre que $\hot_{Y',W}(X' = X \times_Y Y')$. D'autre part, si $g : Y \longrightarrow e$, $g' : Y' \longrightarrow e$ sont les flèches structurales, les deux membres de $(*)$ s'interprètent resp. comme $g'_!(\xi')$ et $g_!(\xi)$ (les $g'_!$ et $g_!$ au sens du prédérivateur $\HOT_W$). Or le calcul de $g'_!$, $g_!$ nous donne que les deux membres s'identifient resp. à
\[
\hot_W(X'), \qquad \hot_W(X),
\]
et nous admettons que la flèche
\[
\hot_W(X') \longrightarrow \hot_W(X)
\]
n'est autre que $\hot_W(q)$. D'ailleurs, comme $X$ est $W$-fibrée sur $S$, par définition pratiquement, $\xi$ est dans $\HOT^{\mathrm{lc}}_W(S)$. Donc il nous suffira de prouver le

\msp{page 56}

\textbf{Corollaire.} \textit{Supposons que $W$ satisfait à W(4b), et soit $p : S' \longrightarrow S$ dans $W$, $\xi \in \HOT^{\mathrm{lc}}_W(S)$, $\xi' = p^*(\xi)$. Alors}
\[
H_\bullet(S',\xi') \xrightarrow{\ \sim\ } H_\bullet(S,\xi)
\]
\textit{est un isomorphisme dans} $\Hot_W$.

\textsc{Démonstration}. Il suffit de prouver que pour tout $\eta \in \mathrm{Ob}\,\Hot_W$ on a une bijection pour les $\Hom(-,\eta)$. Par la formule d'adjonction sur $S$, $S'$, cette application entre les Hom dans $\Hot_W$ s'identifie \texttt{[à]}
\[
\Hom_{\HOT_W(S)}(\xi,\eta_S) \longrightarrow \Hom_{\HOT_W(S')}(\xi',\eta_{S'}),
\]
où $\eta_S = p^*(\eta)$, $\eta_{S'} = p'^*(\eta)$ \texttt{[plutôt} $\eta_S = g^*(\eta)$\texttt{,} $\eta_{S'} = g'^*(\eta)$\texttt{, où} $g : S \longrightarrow e$\texttt{,} $g' : S' \longrightarrow e$ \texttt{sont les flèches structurales]}, donc $\eta_{S'} = p^*(\eta_S)$, et la flèche $(*)$ est la flèche sur les Hom induite par le foncteur
\[
p^*_{\mathrm{lc}} : \HOT^{\mathrm{lc}}_W(S) \longrightarrow \HOT^{\mathrm{lc}}_W(S').
\]
Or moyennant l'axiome W(4b) (que nous n'avons pas utilisé jusqu'à présent), ledit foncteur est une \textit{équivalence de catégories} si $p \in W$ (cf. VI, p. 204).

\hfill q.e.d.

De peur de les oublier, je vais maintenant faire une liste des petites questions que j'ai élucidées chemin faisant en gribouillant.

\begin{itemize}[leftmargin=5em,labelsep=0.5em]
\item[1°)] Si $X$, $Y$ sont $W$-lisses sur $S$, et $f : X \longrightarrow Y$ un $S$-morphisme tel que $f_s : X_s \longrightarrow Y_s$ soit $W$-lisse pour tout $s \in \mathrm{Ob}\,S$, il ne s'ensuit \textit{pas} que $f$ soit $W$-lisse. Mais c'est pourtant vrai quand $X$ et $Y$ sont fibrées scindées sur $S$ (ce
\printedpage{42}
% head: A. Grothendieck : Dérivateurs, Ch. XII, Édition des Universités de Montpellier II et Paris VII
\msp{page 57}
qui implique déjà leur $W$-lissité sur $S$), et $f$ compatible avec les scindages et a fortiori cartésien. Énoncé dual pour les catégories cofibrées scindées et les foncteurs strictement cartésiens entre elles qui sont $W$-propres fibre par fibre : ils sont $W$-propres.

\item[2°)] J'étais intéressé, dans un cas de catégories fibrées scindées $X$, $Y$ sur $S$ et d'un foncteur strictement cartésien $f : X \longrightarrow Y$, supposé $W$-parfait fibre par fibre, de conclure que $f$ soit aussi $W$-parfait. (C'était un cas où de plus on sait que $X$, $Y$ sont aussi des $W$-fibrés sur $S$, i.e. les fonteurs changement de base $X_{s'} \longrightarrow X_s$, $Y_{s'} \longrightarrow Y_s$ sont dans $W$. Mais ça n'a pas beaucoup servi \dots) La question qui restait, en 1°), était celle de la $W$\textit{-propreté} de $f$. Mais c'est sans espoir (même pour $Y = S$), en vertu de la brillante % [sic] fonteurs -> foncteurs, lettre omise

\textbf{Proposition 2.} \textit{Supposons que} $S = 0$
$\begin{tikzcd}[ampersand replacement=\&, row sep=tiny, column sep=small, baseline=(Z.base)]
\& a \\
|[alias=Z]| \phantom{0} \arrow[ur] \arrow[dr] \& \\
\& b
\end{tikzcd}$
\textit{, et soit $Z$ une catégorie sur $S$. Alors $\Phi(Z)$ est propre sur $S$ si et seulement si $Z = \emptyset$.}

Prenons alors dans 2°) ci-dessus $X = \Phi(Z)$, $Y = S$, $f : \Phi(Y) \longrightarrow S$ \texttt{[plutôt} $f : \Phi(Z) \longrightarrow S$\texttt{]} le morphisme structural (donc les $f_s$ sont $W$-parfaits), alors $f$ n'est propre que si $X = \emptyset$. Pourtant, si on prend $Z$ $W$-fibrée sur $S$, $\Phi(Z)$ l'est aussi, et on a un contre-exemple à 2°).
\end{itemize}

Mais on s'est quand même convaincu du résultat suivant :

\textbf{Proposition 3.} \textit{Soient $X$, $Y$ catégories fibrées scindées sur $S$, $f : X \longrightarrow Y$ $S$-foncteur strictement cartésien, tel que}
\msp{page 58}
\textit{pour tout $s \in S$, $f_s : X_s \longrightarrow Y_s$ soit} $W$-fibrant. \textit{Alors,} si $W$ satisfait W(4a), $f$ \textit{est $W$-fibrant.}

(J'espérais avoir que $f$ est $W$-fibrant `à l'\oe{}il', pour prouver que W(4c) implique W(4b). Mais j'étais obligé de supposer W(4a) \dots)


\endgroup
% END UNIT Der-ChXIIa-012/Der-ChXIIa-012_p041-044.tex
\clearpage\setcounter{footnote}{0}\phantomsection\addcontentsline{toc}{section}{Chapitre XII, 5}\markboth{Chapitre XII, 5}{Chapitre XII, 5}
% BEGIN UNIT Der-ChXIIb-012/Der-ChXIIb-012_p001-005.tex
\begingroup\makeatletter%
\newcommand{\printedpage}[1]{\marginpar{\footnotesize p.~#1}}%
\newcommand{\msp}[1]{\marginpar{\footnotesize [#1]}}%
\newcommand{\uHom}{\underline{\mathrm{Hom}}}%
\makeatother\let\printedpage\derpp\let\msp\dermsp%
\let\mspin\DERmspin%


\printedpage{43}
% head: A. Grothendieck : Dérivateurs, Ch. XII, Édition des Universités de Montpellier II et Paris VII ; foot: Version 7 décembre 2002

\msp{page 59}

\section*{V\quad Sommes amalgamées et carrés $W$-cocartésiens dans $\mathrm{Cat}$}

\subsection*{1\quad Généralités sur $U \overset{i}{\hookrightarrow} X \overset{j}{\hookleftarrow} T$, et les sommes amalgamées.}

\subsubsection*{1.1\quad Données de recollement entre un ouvert et un fermé.}

Soit $X$ dans Cat, $U$ un ouvert dans $X$, $T$ le complémentaire fermé. (\textbf{NB} Les ouverts correspondent bien aux ouverts de $\mathrm{Ob}\,X$, topologisé par sa relation de préordre naturelle $x \leq y \iff \mathrm{Hom}_X(x,y) \neq \emptyset$, de même pour les fermés. Les topologies associées aux relations de préordre sont celles où toute intersection d'ouverts est encore ouverte, i.e. toute réunion de fermés est encore fermée. Donc les ouverts d'une telle topologie sont les fermés d'une topologie duale, correspondant à la relation de préordre opposée. La relation de préordre associée à la topologie est
\[
x \leq y \qquad\qquad \iff \qquad\qquad x \in \overline{\{y\}},
\]
la topologie associée au préordre sur $I$ est celle pour laquelle les ouverts sont les réunions d'ensembles de la forme $I/x$ (lesquels forment une base d'ouverts), ou encore les \textit{`cribles'} ($x \in U$, $y \leq x \Longrightarrow y \in U$), les fermés sont les réunions d'ensembles $x\backslash I$, c'est-à-dire les cocribles. Pour toute partie $A$ de $I$, il y a un plus petit ouvert contenant $A$, qui est aussi son plus petit voisinage, c'est la réunion des $I/x$ pour $x \in A$. Et dualement il y a un plus grand fermé contenu dans $A$, dont le complémentaire est le plus petit ouvert contenant $B = I\backslash A$, donc $\bigcup_{x\in B} I/x$. L'adhérence de $A$ est la réunion des adhérences $x\backslash I$ de ses points.)

\msp{page 60}

La situation $(X,U,T)$ se reconstitue, à isomorphisme unique près, par la connaissance de

\begin{enumerate}[label={}]
\item[$1^\circ$)] Les deux catégories $U$, $T$.
\item[$2^\circ$)] Le bifoncteur $\displaystyle\left\{\begin{aligned}
&H : U^\circ \times T \longrightarrow \mathrm{Ens}\\[1ex]
&H(u,t) = \mathrm{Hom}_X(u,t).
\end{aligned}\right.$
\end{enumerate}

De plus, ces données $1^\circ$) $2^\circ$) en sont soumises à aucune restriction. On a, plus exactement % [sic] en sont -> ne sont, negation required by the sense

\textbf{Proposition 1.} \textit{Un isomorphisme} \texttt{[plutôt une équivalence de catégories]} \textit{entre la catégorie $\mathcal{T}$ dont les objets sont les triples $(X,U,T)$, avec $X \in \mathrm{Ob}\,\mathrm{Cat}$, $U \in \mathrm{Ouv}(X)$, $T \in \mathrm{Ferm}(X)$, $U$ et $T$ complémentaires l'un de l'autre, les morphismes de $(X,U,T)$ dans $(X',U',T')$ étant les foncteurs $f : X \longrightarrow X'$ qui induisent $U \longrightarrow U'$, $T \longrightarrow T'$ (d'où
\printedpage{44}
% head: A. Grothendieck : Dérivateurs, Ch. XII, Édition des Universités de Montpellier II et Paris VII ; foot: Version 7 décembre 2002
$U = f^{-1}(U')$, $T = f^{-1}(T')$), avec la catégorie $\mathcal{T}'$ des triples $(U,T,H)$, de deux catégories $U$, $T$ et d'un bifoncteur}
\[
H : U^\circ \times T \longrightarrow \mathrm{Ens}
\]
\hspace*{2em}\textit{(ou encore un préfaiscau $U^\circ \longrightarrow T^\vee$ sur $U$, donc un objet de} $\uHom(U^\circ,T^\vee)$\textit{)),} % [sic] préfaiscau -> préfaisceau, misspelling

\textit{un morphisme de $(U,T,H)$ dans $(U',T',H')$ étant un triple}
\[
(f_U, f_T, u) \qquad\qquad
\left\{\begin{array}{l}
\left.\begin{array}{llll}
f_U : & U & \longrightarrow & U' \\[1ex]
f_T : & T & \longrightarrow & T'
\end{array}\right\}\text{des foncteurs}\\[2ex]
u : H \longrightarrow H' \circ (f_U^\circ \times f_T),\\
\text{i.e. } H(u,t) \longrightarrow H'(f_U(u), f_T(t))\\
\text{un homomorphisme fonctoriel}\\
\text{(flèche dans } (U \times T^\circ)^\wedge \simeq \uHom(U^\circ,T^\vee)\text{)),}
\end{array}\right.
\]
\textit{composition des flèches évidente.}

Si $U \hookrightarrow X \hookleftarrow T$ est un triple, et
\[
f_{U'} : U' \longrightarrow U,\ f_{T'} : T' \longrightarrow T
\]
des flèches dans Cat, il y a une façon privilégiée de `recoller' $U'$ et $T'$, de façon à ce que l'objet $(X',U',T')$ s'envoie dans $(X,U,T)$ en induisant

\msp{page 61}

$f_{U'}$, $f_{T'}$:
\[
\begin{tikzcd}[column sep=large, row sep=large]
U' \arrow[r, hook, "i'"] \arrow[d, "f_{U'}"'] & X' \arrow[d, "f"] & T' \arrow[l, hook', "j'"'] \arrow[d, "f_{T'}"] \\
U \arrow[r, hook, "i"] & X & T, \arrow[l, hook', "j"']
\end{tikzcd}
\]
elle s'obtient à l'aide des bifoncteurs
\[
\left\{\begin{aligned}
&H'(u',t') \overset{\text{déf}}{=} H(f_{U'}(u'), f_{T'}(t'))\\
&\text{i.e. } H' = H \circ (f_{U'}^\circ \times f_{T'}).
\end{aligned}\right.
\]
Ce triple, i.e. cet objet de $\mathcal{T}$, représente le foncteur suivant sur $\mathcal{T}$ :
\[
(X'',U'',T'') \longmapsto \mathrm{Hom}_{\mathrm{Cat}}(U'',U) \times \mathrm{Hom}_{\mathrm{Cat}}(T'',T),
\]
i.e. on a pour tout objet $(X'',U'',T'')$ de $\mathcal{T}$ une bijection canonique, fonctorielle en cet objet
\[
\mathrm{Hom}_{\mathcal{T}}((X'',U'',T''),(X',U',T')) \simeq \mathrm{Hom}_{\mathrm{Cat}}(U'',U) \times \mathrm{Hom}_{\mathrm{Cat}}(T'',T). \tag*{$(*)$}
\]

\printedpage{45}
% head: A. Grothendieck : Dérivateurs, Ch. XII, Édition des Universités de Montpellier II et Paris VII ; foot: Version 7 décembre 2002

Nul doute qu'on a même
\[
\uHom_{\mathcal{T}}((X'',U'',T''),(X',U',T')) \simeq \uHom_{\mathrm{Cat}}(U'',U) \times \uHom_{\mathrm{Cat}}(T'',T), \tag*{$(**)$}
\]
tout aussi fonctoriel, dont $(*)$ se déduit en passant aux ensembles d'objets. On dit que le triple $(X',U',T')$ est le triple \textit{induit} par le triple $(X,U,T)$, via les foncteurs $f_{U'}$, $f_{T'}$.

\textbf{NB} Si $f_{U'}$, $f_{T'}$ sont pleinement fidèles, alors il en est de même de $f : X' \longrightarrow X$. Nous utiliserons cette opération de recollement privilégié surtout quand $U' = U$, $f_{U'} = \mathrm{id}_U$, ou inversement $T' = T$, $f_{T'} = \mathrm{id}_T$.

\subsubsection*{1.2\quad Relation avec le recollement des topos.}

La situation d'un triple
\[
(\mathcal{X},\mathcal{U},\mathcal{T})
\]
d'un topos $\mathcal{X}$, d'un ouvert $\mathcal{U}$ de $\mathcal{X}$, et d'un

\msp{page 62}

sous-topos fermé complémentaire, est toute similaire (\textbf{NB} Si $\mathcal{X} = \mathrm{Top}(X)$, avec $X$ dans Cat, alors $\mathrm{Ouv}(\mathcal{X}) \simeq \mathrm{Ouv}(X)$, puisque les ouverts correspondent aux sous-faisceaux du faisceau final sur $\mathcal{X}$, i.e. du préfaisceau final sur $X$. Cette bijection est compatible avec le passage aux complémentaires.) La théorie des topos dit que les couples $(\mathcal{X},\mathcal{U},\mathcal{T})$ \texttt{[plutôt triples]} se reconstituent, à équivalence définie à isomorphisme unique près, à l'aide d'un triple
\[
\left\{\begin{aligned}
&\mathcal{U},\ \mathcal{T} \text{ deux topos}\\[1ex]
&\varphi : \mathrm{Fais}\,(\mathcal{U}) \longrightarrow \mathcal{T} \text{ un foncteur exact à gauche (`accessible'),}
\end{aligned}\right.
\]
ce foncteur $\varphi$ étant obtenu, quand on part de
\[
\mathcal{U} \xrightarrow{\;i\;} \mathcal{X} \xleftarrow{\;j\;} \mathcal{T},
\]
comme
\[
\varphi = j^*i_*.
\]
De plus, le triple $(\mathcal{U},\mathcal{T},\varphi)$ peut être choisi arbitrairement, avec la seule restriction (pratiquement anodine, car toujours satisfaite pour tous les foncteurs naturels qu'on rencontre) que $\varphi$ soit (en plus d'être exact à gauche) `accessible'.

Dans le cas actuel $\mathcal{X} = \mathrm{Top}(X)$, $\mathcal{U} = \mathrm{Top}(U)$, $\mathcal{T} = \mathrm{Top}(T)$, le foncteur $\varphi = j^*i_*$ est non seulement exact à gauche, mais même il commute aux $\varprojlim$ quelconques, puisqu'il en est ainsi (toujours) de $i_*$, et aussi (parce que le $j$ topossique est induit par le $j : T \longrightarrow X$ catégorique) de $j^*$. En fait, on
\printedpage{46}
% head: A. Grothendieck : Dérivateurs, Ch. XII, Édition des Universités de Montpellier II et Paris VII ; foot: Version 7 décembre 2002
\msp{page 63}
trouve ceci :

\textbf{Proposition 2 :} \textit{Considérons une situation de recollement de deux topos $\mathcal{U}$, $\mathcal{T}$ par $\varphi : \mathrm{Fais}\,\mathcal{U} \longrightarrow \mathrm{Fais}\,\mathcal{T}$, foncteur exact à gauche. Supposons que $\mathcal{U} = \mathrm{Top}(U)$, $\mathcal{T} = \mathrm{Top}(T)$, avec $U$ et $T$ dans $\mathrm{Cat}$. Pour que le topos recollé $\mathcal{X}$ soit équivalent à un topos $\mathrm{Top}(X)$, avec $X$ dans $\mathrm{Cat}$, il faut et il suffit que le foncteur de recollement $\varphi$ commute aux $\varprojlim$ quelconques, ou ce qui revient au même, qu'il ait un adjoint à gauche :}
\[
\psi : \underbrace{\mathrm{Fais}\,\mathcal{T}}_{= T^\wedge} \longrightarrow \underbrace{\mathrm{Fais}\,\mathcal{U}}_{= U^\wedge}.
\]

On s'attend dès lors que la donnée de recollement $\varphi$, ou ce qui revient au même le foncteur $\psi$ ayant un adjoint à droite, i.e. commutant aux $\varinjlim$ quelconques, revient à celle d'une donnée de recollement $H$
\[
H : U^\circ \times T \longrightarrow \mathrm{Ens}
\]
pour $U$ et $T$. Mais on sait bien que les foncteurs $\psi : T^\wedge \longrightarrow \mathcal{C}$ dans une catégorie $\mathcal{C}$ avec $\varinjlim$ quelconques, qui commutent auxdites $\varinjlim$, forment une catégorie équivalente (par le foncteur $\psi \longmapsto \psi|T(\subseteq T^\wedge)$) à celle des foncteurs
\[
\psi_0 : T \longrightarrow \mathcal{C} \quad\ (= \underbrace{U^\wedge}_{\uHom(U^\circ,\mathrm{Ens})} \text{ en l'occurrence}),
\]
\msp{page 64}
et la donnée d'un tel foncteur $\psi_0$ revient bien à celle d'un $H$. Donc

\textbf{Corollaire.} \textit{La catégorie des données de recollement topossique $\varphi : U^\wedge \longrightarrow T^\wedge$ entre} $\mathrm{Top}(U)$, $\mathrm{Top}(T)$\textit{, telles que le topos recollé soit un `topos} paradisiaque\textit{', i.e. équivalent à un topos de la forme $X^\wedge$ (avec $X$ dans} Cat\textit{), est équivalente à celle des données de recollement $H : U^\circ \times T \longrightarrow$} Ens \textit{entre $U$ et $T$. La donnée topossique $\varphi$ associée à $H$ est le foncteur adjoint à droite de $T^\wedge \longrightarrow U^\wedge$ déduit de $T \xrightarrow{H} U^\wedge$ par prolongement commutant aux $\varinjlim$, ou encore elle est donnée par}
\[
\varphi(\underbrace{F}_{\in U^\wedge})(t) \simeq \mathrm{Hom}_{U^\wedge}(H(t),F),
\]
\textit{où dans le second membre $H$ est interprété comme un foncteur $T \longrightarrow U^\wedge$. Le topos recollé par $\varphi$ est canoniquement équivalent à} $\mathrm{Top}(X)$\textit{, où $X$ est l'objet de} Cat \textit{obtenu par recollement à l'aide de $H$.}

Du point de vue de la théorie des faisceaux, cela donne l'interprétation de $H$, comme exprimant le foncteur
\[
\varphi = j^*i_* : U^\wedge \longrightarrow T^\wedge
\]
par la formule précédente.

\printedpage{47}
% head: A. Grothendieck : Dérivateurs, Ch. XII, Édition des Universités de Montpellier II et Paris VII ; foot: Version 7 décembre 2002

\msp{page 65}

\subsubsection*{1.3\quad Flèches dans $\mathrm{Cat}$, de source ou but un objet recollé $(X,U,T)$.}

Quand $(X,U,T)$ est donné par une donnée de recollement
\[
U,\ T,\ H : U^\circ \times T \longrightarrow \mathrm{Ens},
\]
les flèches
\[
f : X' \longrightarrow X
\]
dans Cat se décrivent immédiatement ainsi : la donnée de $f$ définit d'abord
\[
U' = f^{-1}(U), \qquad\qquad T' = f^{-1}(T),
\]
couple d'un ouvert de $X'$ et du fermé complémentaire, de sorte que cette donnée équivaut à celle de

\begin{enumerate}[label={}]
\item[$1^\circ$)] un couple $(U',T')$ dans $X'$
\item[$2^\circ$)] un morphisme (dans $\mathcal{T}$) de $(X',U',T')$ dans $(X,U,T)$, ce qui équivaut par 1.1 à un triple
\[
\left\{\begin{array}{rccl}
f_{U'} : & U' & \longrightarrow & U\\[1ex]
f_{T'} : & T' & \longrightarrow & T\\[1ex]
u : & H_{U',T'} & \longrightarrow & H \circ (f_{U'}^\circ \times f_{T'}),
\end{array}\right.
\]
où on a posé
\[
\left\{\begin{aligned}
&H_{U',T'} : U^{\prime\circ} \times T' \longrightarrow \mathrm{Ens}\\[1ex]
&H_{U',T'}(u',t') = \mathrm{Hom}_{X'}(u',t'),
\end{aligned}\right.
\]
i.e. $H_{U',T'}$ est la donnée de recollement qui conrrespond au triple $(X',U',T')$ dans $\mathcal{T}$. % [sic] conrrespond -> correspond, misspelling
\end{enumerate}

\msp{page 66}

Nous serons plus intéressés par la situation inverse, i.e. par la question de déterminer les foncteurs
\[
\underbrace{X}_{= X_{U,T,H}} \quad \longrightarrow X'
\]
de $X$, obtenu par $(U,T,H)$, dans un objet quelconque de Cat. On trouve tautologiquement

\textbf{Proposition 3.} $\mathrm{Hom}_{\mathrm{Cat}}(X,X')$ \textit{est en correspondance bijective canonique, fonctorielle en $X'$ dans} Cat\textit{, avec l'ensemble des triples}


\endgroup
% END UNIT Der-ChXIIb-012/Der-ChXIIb-012_p001-005.tex
% BEGIN UNIT Der-ChXIIb-012/Der-ChXIIb-012_p006-010.tex
\begingroup\makeatletter%
\newcommand{\printedpage}[1]{\marginpar{\footnotesize p.~#1}}%
\newcommand{\msp}[1]{\marginpar{\footnotesize [#1]}}%
\newcommand{\Hom}{\mathrm{Hom}}%
\newcommand{\uHom}{\underline{\mathrm{Hom}}}%
\newcommand{\Cat}{\mathrm{Cat}}%
\newcommand{\Ens}{\mathrm{Ens}}%
\newcommand{\id}{\mathrm{id}}%
\newcommand{\Ob}{\mathrm{Ob}\,}%
\makeatother\let\printedpage\derpp\let\msp\dermsp%
\let\mspin\DERmspin%


\printedpage{48}
% head: A. Grothendieck : Dérivateurs, Ch. XII, Édition des Universités de Montpellier II et Paris VII
% foot: Version 7 décembre 2002

\[
\left\{
\begin{aligned}
&\left.\begin{aligned} f_U &: U \longrightarrow X' \\ f_T &: T \longrightarrow X' \end{aligned}\right\}\text{flèches dans Cat}\\
&u : H \longrightarrow H_{X'}(f_U^{\circ}\times f_T),\ \text{flèche dans } \uHom(U^{\circ}\times T,\Ens),
\end{aligned}
\right.
\tag*{\texttt{[}$(*)$\texttt{]}}
\]
\textit{où pour tout $X'$ dans} Cat\textit{, on a posé}
\[
H_{X'} = ((x,y)\longmapsto \Hom_{X'}(x,y)) : X'^{\circ}\times X' \longrightarrow \Ens.
\]

Sûrement on aura même un isomorphisme de catégories
\[
\uHom(X,X') \simeq \text{catégorie des triples } (f_U,f_T,u) \text{ comme dans } (*).
\]

\subsection*{1.4\quad Construction de sommes amalgamées à partir d'une immersion ouverte où fermée.} % [sic] où fermée -> ou fermée, conjonction « ou » et non adverbe « où »

Pour une construction d'une sorte de somme amalgamée, en termes d'un triple $(X,U,T)$ $+\ \ldots$, \msp{page 67} qui soit symétrique en $U$, $T$, j'ai envie de regarder d'abord la situation ‘mixte’, d'un diagramme en traits pleins
\[
\begin{tikzcd}
U \arrow[r, hook, "i"] \arrow[d, "f_U"'] & X \arrow[d, dotted, "f"] & T \arrow[l, hook', "j"'] \arrow[d, "f_T"] \\
U' \arrow[r, hook, dotted, "i'"] & X' & T', \arrow[l, hook', dotted, "j'"']
\end{tikzcd}
\tag*{$(*)$}
\]
qu'on aimerait compléter ‘canoniquement’, si faire se peut, en un morphisme de triples (i.e. en une flèche dans $\mathcal{T}$). L'idée naturelle, c'est que le $X'$ qui s'y insère serait la limite inductive, dans Cat, du diagramme initial. Si $T'=T$, $f_T=\id_T$, cela signifie donc que $X'$ s'identifie à la somme amalgamée $X\amalg_U U'$, et inversement si $U'=U$, $f_U=\id_U$.

Désignons par $X'$ la somme amalgamée \texttt{[plutôt la limite inductive]}, et explicitons, pour $X''$ quelconque dans Cat, l'ensemble des flèches $g$ de $X'$ dans $X''$, comme foncteur en $X''$. C'est l'ensemble des triples
\[
(g_X, g_{U'}, g_{T'}) \qquad
\left\{\begin{aligned} g_X &: X \longrightarrow X'' \\ g_{U'} &: U' \longrightarrow X'' \\ g_{T'} &: T' \longrightarrow X''\end{aligned}\right.
\qquad \text{avec} \quad
\begin{aligned} g_X\circ i &= g_{U'} f_U \\ g_X \circ j &= g_{T'} f_T. \end{aligned}
\tag*{$(*)$}
\]
\printedpage{49}
% head: A. Grothendieck : Dérivateurs, Ch. XII, Édition des Universités de Montpellier II et Paris VII
% foot: Version 7 décembre 2002
\[
\begin{tikzcd}
U \arrow[r, hook, "i"] \arrow[d, "f_U"'] & X \arrow[dd, "g_X"'] & T \arrow[l, hook', "j"'] \arrow[d, "f_T"] \\
U' \arrow[dr, "g_{U'}"'] & & T' \arrow[dl, "g_{T'}"] \\
& X'' &
\end{tikzcd}
\]

Notons que pour un tel triple, $g_X$ est partiellement connu déjà en termes de $g_{U'}$, $g_{T'}$, par ses restrictions $g_X\circ i$ et $g_X\circ j$ à $U$ et $T$ respectivement, par les formules
\[
g_U = g_{U'}f_U, \qquad g_T = g_{T'}f_T.
\]

\msp{page 68} Donc la donnée d'un triple $(*)$ équivaut à celle d'un triple
\[
(g_{U'}, g_{T'}, g_X) \qquad
\left\{\begin{aligned}
&g_X \text{ une flèche } X\longrightarrow X'' \text{ correspondant à}\\
&g_U = g_{U'}\circ f_U \text{ et } g_T = g_{T'}\circ f_T \text{ déjà donnés (via } g_{U'}\text{, } g_{T'}\text{),}
\end{aligned}\right.
\]
donc en vertu de 1.3. cela correspond à
\[
(g_{U'}, g_{T'}, u) \qquad
\left\{\begin{aligned}
g_{U'} &: U' \longrightarrow X''\\
g_{T'} &: T' \longrightarrow X''\\
u &: H \longrightarrow \underbrace{H_{X''}\left((g_{U'}\circ f_U)^{\circ}\times(g_{T'}\circ f_T)\right)}_{(u,t)\longmapsto \Hom_{X''}(g_{U'}f_U(u),g_{T'}f_T(t))}.
\end{aligned}\right.
\tag*{$(*)$}
\]
On va interpréter la donnée de $u$, en termes d'un foncteur
\[
F^{*} = (H'\longmapsto H'\circ F) : \uHom(U'^{\circ}\times T',\Ens) \longrightarrow \uHom(U^{\circ}\times T,\Ens)
\tag*{$(**)$}
\]
où
\[
F : U^{\circ}\times T \longrightarrow U'^{\circ}\times T' \qquad \text{est donné par } F = f_U^{\circ}\times f_T.
\]
On a un objet $H''=H''_{g_{U'},g_{T'}}$ de la catégorie source de $(**)$ déterminé en termes du couple $(g_{U'},g_{T'})$, par
\[
H''(u',t') = \Hom_{X''}(g_{U'}(u'),g_{T'}(t')),
\]
et le bifoncteur but de la flèche $u$ dans $(*)$ s'interprète comme
\[
F^{*}(H'') = H''\circ F : (u,t)\longmapsto \underbrace{H''(f_U(u),f_T(t))}_{=\Hom(g_{U'}f_U(u),g_{T'}f_T(t))} \quad .
\]
Ainsi, $u$ s'interprète comme une flèche
\[
u : H \longrightarrow F^{*}(H'').
\]
Or le foncteur $F^{*}$ a un adjoint à gauche $F_{!}$, i.e. on a
\[
\Hom(H,F^{*}(H'')) \simeq \Hom(F_{!}(H),H'').
\tag*{$(*{*}*)$}
\]

\printedpage{50}
% head: A. Grothendieck : Dérivateurs, Ch. XII, Édition des Universités de Montpellier II et Paris VII
% foot: Version 7 décembre 2002
\msp{page 69} avec
\[
F_{!} : \uHom(U^{\circ}\times T,\Ens) \longrightarrow \uHom(U'^{\circ}\times T',\Ens).
\]
Posons
\[
H' = F_{!}(H) : U'^{\circ}\times T' \longrightarrow \Ens,
\]
c'est un bifoncteur bien déterminée par la seule donnée du diagramme initial $(*)$ p. 67, ou ce qui revient au même, en termes de % [sic] déterminée -> déterminé, accord avec « bifoncteur » (masculin)
\[
(U,T,H,f_U : U\longrightarrow U', f_T : T\longrightarrow T').
\]

Et on voit que la donnée d'une flèche $X'\longrightarrow X''$, $X''$ objet variable de Cat, équivaut à celle de
\[
g_{U'} : U'\longrightarrow X'',\ g_{T'} : T'\longrightarrow X'',\ u : H'\longrightarrow H'' = H_{X''}(g_{U'}^{\circ}\times g_{T'}),
\]
c'est à dire, en vertu de 1.3., à celle d'une flèche
\[
X'_0 \longrightarrow X'',
\]
où $X'_0$ désigne l'objet de Cat déduit de $U'$, $T'$ par recollement au moyen de $H'$. Cela montre donc que la limite inductive cherchée n'est autre que $X'_0$. On a ainsi trouvé le

\textbf{Théorème 1.} \textit{La limite inductive $X'$ dans} Cat \textit{du diagramme} $(*)$ \textit{page 67, défini en termes d'un triple $(X,U,T)$ dans $\mathcal{T}$ et de flèches $f_U : U\longrightarrow U'$, $f_T : T\longrightarrow T'$ dans} Cat\textit{, existe (sans mérite)} \msp{page 70} \textit{et les flèches canoniques}
\[
U' \overset{i'}{\hookrightarrow} X' \overset{j'}{\hookleftarrow} T'
\]
\textit{sont des immersions ouvertes et fermées complémentaires l'une de l'autre. La donnée de recollement qui définit $X'$ via $U'$, $T'$ est donnée par}
\[
H' = (\underbrace{f_U^{\circ}\times f_T}_{F})_{!}(H) : U'^{\circ}\times T' \longrightarrow \Ens,
\]
\textit{où $H$ est la donnée de recollement qui définit $X$, i.e. $H : U^{\circ}\times T\longrightarrow$} Ens \textit{est le bifoncteur $(u,t)\longmapsto \Hom_X(u,t)$.}

Ainsi on a, pour $u'$ dans $U'$ et $t'$ dans $T'$
\[
\Hom_{X'}(u',t') = F_{!}(H)(u',t'),
\]
et il reste à expliciter $F_{!}(H)$. Ici $F=f_U^{\circ}\times f_T$ est un foncteur covariant sur $U^{\circ}\times T$, on peut l'interpréter comme un préfaisceau sur $U\times T^{\circ}$, et écrire dès lors
\[
\underbrace{F_{!}(H)(u',t')}_{\Hom_{X'}(u',t')} = \varinjlim_{(u',t')\backslash U\times T^{\circ}} \underbrace{\Hom_X(u,t)}_{H(u,t)} \ .
\tag*{$(*)$}
\]

\printedpage{51}
% head: A. Grothendieck : Dérivateurs, Ch. XII, Édition des Universités de Montpellier II et Paris VII
% foot: Version 7 décembre 2002
La catégorie $(u',t')\backslash U\times T^{\circ}$ sur laquelle on prend la $\varinjlim$ s'explicite comme un produit
\[
(u',t')\backslash U\times T^{\circ} \simeq (u'\backslash U)\times(T/t')^{\circ}.
\]
Les applications
\[
\Hom_X(u,t)\longrightarrow \Hom_{X'}(u',t') \qquad (=F_{!}(H)(u',t'))
\]
\msp{page 71} déduites de cette formule ne sont autres que les suivantes : $u$, $t$ sont données comme objets de $u'\backslash U$, $T/t'$, i.e. avec des flèches
\[
u'\xrightarrow{\alpha} f_U(u), \qquad f_T(t)\xrightarrow{\beta} t',
\]
d'où
\[
\Hom_X(u,t) \xrightarrow{\text{via } f=f_X:X\to X'} \Hom_{X'}(\underbrace{f(u)}_{f_U(u)},\underbrace{f(t)}_{f_T(t)}) \xrightarrow{\text{via } \alpha,\ \beta} \Hom_{X'}(u',t').
\]

Notons pour mémoire que pour tout objet $Z$ de Cat, et tout préfaisceau $G$ sur $Z$, on a la formule
\[
\varinjlim_{Z} G \simeq \pi_0(Z/G),
\]
où $Z/G$ s'interprète en termes de l'immersion canonique $Z\longrightarrow Z^{\wedge}$ $(G\in \Ob Z^{\wedge})$. C'est donc une catégorie fibrée sur $Z$ à fibres discrètes, et on a en fait un isomorphisme entre la sous-catégorie pleine de $\Cat/Z$ formée de telles catégories fibrées à fibres discrètes, avec $Z^{\wedge}$. Dans ce dictionnaire parfait entre préfaisceaux d'ensembles et catégories fibrées à fibres discrètes, le foncteur $\varinjlim$ sur $Z^{\wedge}$ correspond simplement au foncteur $\pi_0$ sur la catégorie des objets au dessus de $Z$ dans Cat, lui-même induit par le foncteur (indépendant de $Z$)
\[
\Cat \xrightarrow{\pi_0} \Ens.
\]

\msp{page 72} On peut donc réécrire la formule fondamentale $(*)$ p. 70, en termes du foncteur $H$ définissant le recollement $X$ de $U$ et de $T$, par
\[
\boxed{\underbrace{\Hom_{X'}(u',t')}_{H'(u',t')} \simeq \pi_0(\underbrace{(u',t')\backslash U\times T^{\circ}/H}_{((u'\backslash U)\times(T/t')^{\circ})/H}).}
\tag*{$(*)$}
\]

Par la suite, nous n'aurons sans doute qu'à utiliser le cas où on a soit $T'=T$, $f_T=\id_T$, soit l'inverse. Nous travaillerons surtout dans le premier cas, en déduisant ce dont nous aurons besoin pour le cas dual, par principe de dualité (automorphisme $X\longmapsto X^{\circ}$ de la
\printedpage{52}
% head: A. Grothendieck : Dérivateurs, Ch. XII, Édition des Universités de Montpellier II et Paris VII
% foot: Version 7 décembre 2002
catégorie Cat). Dans ce cas, dans la formule fondamentale $(*)$ l'élément $t'$ de $T'$ s'écrit comme un élément $t$ de $T$, et la formule se retranscrit comme
\[
\boxed{\Hom_{X'}(u',t) \simeq \pi_0(\underbrace{u'\backslash U/t}_{\text{catégorie des triples } (u,\alpha,\beta'),\ \left\{\begin{aligned} &u\in \Ob U\\ &\alpha : u\longrightarrow t\\ &\beta' : u'\longrightarrow f_U(u)\end{aligned}\right.}).}
\tag*{$(**)$}
\]

C'est la formule que j'avais trouvée directement dans ce cas, sans me soucier encore de faire varier $T$ en même temps que $U$. Je voudrais faire le lien entre les deux formules $(*)$ et $(**)$, par voie topologique, donc établir une bijection canonique, fonctorielle etc.
\[
\pi_0(u'\backslash U/t) \xrightarrow{\sim} \pi_0(((u'\backslash U)\times(T/t)^{\circ})/H).
\]

\msp{page 73} On a sûrement un foncteur
\[
u'\backslash U/t \longrightarrow ((u'\backslash U)\times(T/t)^{\circ})/H
\]
qui associe à un triple
\[
(u,\alpha,\beta') : u\in \Ob U,\ \alpha : u\longrightarrow t,\ \beta' : u'\longrightarrow f_U(u)
\]
le triple $(\bar u,\bar t,\xi)$ $(\xi\in H(\bar u,\bar t))$ dans le second membre, avec $\bar u$ dans $u'\backslash U$, $\bar t$ dans $T/t$, défini ainsi :
\[
\bar u = (u,\beta'),\ \bar t = (t,\id_t), \underbrace{\xi}_{=\alpha} \in H(\bar u,\bar t) = \Hom(u,t).
\]
En somme
\[
u'\backslash U/t = \text{fibre de } ((u'\backslash U)\times(T/t)^{\circ})/H \text{ en l'objet initial } \bar t \text{ de } (T/t)^{\circ}
\]
par la projection naturelle, flèche composée
\[
\begin{tikzcd}
((u'\backslash U)\times(T/t)^{\circ})/H \arrow[d] \\
(u'\backslash U)\times(T/t)^{\circ} \arrow[d, "\pi"] \\
(T/t)^{\circ}.
\end{tikzcd}
\]

Il faut voir que l'inclusion de la fibre donne une bijection sur les $\pi_0$, ou ce qui revient au même, que l'inclusion de la fibre $u'\backslash U$ \textit{de} $\pi$ en l'objet initial $(t,\id_t)$ de $(T/t)^{\circ}$ donne une bijection sur les $\varinjlim$ du foncteur $H$. Donc l'inclusion a intérêt à être asphérique, ou du moins co-cofinalle. % [sic] co-cofinalle -> co-cofinale, orthographe


\endgroup
% END UNIT Der-ChXIIb-012/Der-ChXIIb-012_p006-010.tex
% BEGIN UNIT Der-ChXIIb-012/Der-ChXIIb-012_p011-015.tex
\begingroup\makeatletter%
\newcommand{\printedpage}[1]{\marginpar{\footnotesize p.~#1}}%
\newcommand{\msp}[1]{\marginpar{\footnotesize [#1]}}%
\newcommand{\Ob}{\mathrm{Ob}\,}%
\newcommand{\Hom}{\mathrm{Hom}}%
\makeatother\let\printedpage\derpp\let\msp\dermsp%
\let\mspin\DERmspin%


% ===================== page 11 =====================
\printedpage{53}
% head: A. Grothendieck : Dérivateurs, Ch. XII, Édition des Universités de Montpellier II et Paris VII

\msp{page 74}
Mais l'inclusion d'un objet initial est toujours asphérique, et le reste par produit avec un objet quelconque (i.e. $u'\backslash U$) de Cat (et même par tout changement de base lisse \dots), d'où la conclusion.

Dans la suite, nous nous intéressons surtout à la situation de sommes amalgamées (où on change de notation un peu, à présent)
\[
\begin{tikzcd}
U \arrow[r, hook, "i"] \arrow[d, "p"'] & X \arrow[d, "q"'] & T \arrow[l, hook', dotted, "j"'] \arrow[d, dotted] \\
U' \arrow[r, hook, "i'"] & X' & T' \simeq T. \arrow[l, hook', dotted]
\end{tikzcd}
\tag*{$(*)$}
\]

Le fermé $T'$ complémentaire de $U'$ dans $X'$ est alors canoniquement isomorphe au fermé $T$ complémentaire de $U$ dans $X$, l'isomorphisme étant induit par $q$. Souvent nous identifierons $T'$ et $T$, ce qui est licite en regardant $X'$ comme déduit du recollement de $U'$ et de $t$ \texttt{[plutôt }$T$\texttt{]}, par le bifoncteur $U'^{\circ}\times T\longrightarrow \mathrm{Ens}$ explicité dans la formule $(**)$ p.~72. Nous nous intéressons surtout à dégager des conditions qui assurent les deux implications suivantes

\begin{itemize}
\item[a)] $i\in W\Longrightarrow i'\in W$ (et même $\Longleftrightarrow$).
\item[b)] $p\in W\Longrightarrow q\in W$ (et même $\Longleftrightarrow$).
\end{itemize}

\msp{page 75}
\textbf{Proposition 4.} \textit{Soit $t\in\Ob T$. Pour que $t\in\Ob\bar{U}$, il faut et il suffit que $q(t)\in\Ob\overline{U'}$. En d'autres termes, le `bord' $\dot{T}=T\cap\bar{U}$ de $T$ dans $X$ est `le même' que son bord dans $X'$.}

\textsc{Démonstration}. Dire que $t\in\Ob\bar{U}$ signifie aussi
\[
U/t\neq\emptyset
\]
et implique, par continuité, que $q(t)\in\overline{q(U)}\subseteq\overline{U'}$. Inversement, si $q(t)\in\overline{U'}$, i.e. s'il existe $u'\in U'$ tel que
\[
\Hom(u',q(t))(=\pi_0(u'\backslash U/t))\neq\emptyset,
\]
il en résulte que $u'\backslash U/t\neq\emptyset$, a fortiori $U/t$ (qui reçoit $u'\backslash U/t)\neq\emptyset$, q.e.d.

\textbf{Proposition 5.} \textit{Soit}
\[
\begin{tikzcd}
Y \arrow[r, "i"] \arrow[d, "p"'] & X \arrow[d, "q"] \\
Y' \arrow[r, "i'"] & X'
\end{tikzcd}
\]
% foot: Version 7 décembre 2002
% ===================== page 12 =====================
\printedpage{54}
% head: A. Grothendieck : Dérivateurs, Ch. XII, Édition des Universités de Montpellier II et Paris VII
\textit{un carré cocartésien dans \textup{Cat}. Si $i$ est une immersion ouverte (resp. fermée), alors $i'$ aussi, et le carré est aussi cartésien. Idem pour $p$ et $q$. Si $i$ et $p$ sont des immersions ouvertes (\footnote{\textbf{NB} Nous utilisons \textit{`immersion ouverte'}, \textit{`immersion fermée'} au \textit{sens strict}\,: c'est un \textit{monomorphisme} pleinement fidèle dont l'image est une sous-catégorie \textit{pleine,} soit ouverte, soit fermée.}) (resp. \dots), alors la situation est isomorphe à celle d'un objet $X'$ de \textup{Cat}, de deux ouverts (resp. fermés) $X$, $Y'$ de $X$ ayant pour réunion $X'$, pour intersection} % [sic] de $X$ -> de $X'$, X and Y' are opens of X' (their union is X')
\textit{$Y$. Inversement, dans une telle situationm le carré correspondant est cocartésien, i.e. $X$} % [sic] situationm -> situation, (comma), typing error
\texttt{[plutôt }$X'$\texttt{]} \textit{s'identifie à la somme amalgamée}
\[
X'\simeq X\amalg_{Y}Y'.
\]

Donc aussi on a la proposition 15, page 158, et ses corollaires.
% foot: Version 7 décembre 2002

% ===================== page 13 =====================
\printedpage{55}
% head: A. Grothendieck : Dérivateurs, Ch. XII, Édition des Universités de Montpellier II et Paris VII

\section*{2\quad Sommes amalgamées et $W$-équivalences.}

\subsection*{2.1\quad Axiome du carré cocartésien et axiome du cube (ou `Lemme des cinq').}

Nous supposons à présent donné un localiseur fondamental
\[
W\subseteq\mathrm{Fl}(\mathrm{Cat}),
\]
donc $W$ satisfait à W(1, 2 bis, 3). Nous ne supposons pas pour le moment pour les autres axiomes standard W(4) (axiome de fibration), W(5) (axiome de connexité), W(6) (axiome des limites inductives). Mais nous allons introduire à présent deux axiomes de stabilité supplémentaire, de nature différente, concernant les sommes amalgamées.

\begin{description}[leftmargin=5em, style=nextline]
\item[W(7 bis)] (Axiome du carré cocartésien.) Soient $X$ un objet de Cat, $U$, $U'$ deux ouverts de $X$ tels que $U\cup U'=X$, considérons le diagramme cocartésien correspondant (cf. prop. 5 en marge de la page \texttt{[ci-dessus]}) :

\msp{page 76}
\[
\begin{tikzcd}[column sep=small]
& U\cap U' \arrow[dl, hook', "i'_0"'] \arrow[dr, hook, "i_0"] & \\
U \arrow[dr, hook, "i"'] & & U' \arrow[dl, hook', "i'"] \\
& X &
\end{tikzcd}
\tag*{$(*)$}
\]

Alors
\[
i_0\in W\qquad\Longleftrightarrow\qquad i\in W
\]
(et symétriquement
\[
i'_0\in W\qquad\Longleftrightarrow\qquad i'\in W).
\]
\end{description}

\textbf{NB} Par dualité, cet axiome est équivalent à l'axiome similaire, avec des \textit{fermés} $T$, $T'$ de $X$ de réunion $X$.

\textbf{Corollaire.} \textit{Supposons \textup{W(7 bis)}. Considérons les conditions suivantes sur le carré précédent.}

\begin{itemize}
\item[\textit{a)}] \textit{$U\cap U'$ $W$-asphérique.}
\item[\textit{b)}] \textit{$U$, $U'$ $W$-asphériques.}
\item[\textit{c)}] \textit{$X$ $W$-asphérique.}
\end{itemize}

\textit{Alors (a et b) implique c, et (b et c) implique a.}

\textbf{NB} Moralement, il faudrait aussi que (c et a) implique b, i.e. si $U\cap U'$ et $X$ sont $W$-asphériques, il en est de même de $U$ et de $U'$. Mais il ne semble pas que l'axiome du carré cocartésien permette de le prouver. Par contre, l'axiome qui suit impliquera le corollaire précédent sous sa forme complète.
% foot: Version 7 décembre 2002

% ===================== page 14 =====================
\printedpage{56}
% head: A. Grothendieck : Dérivateurs, Ch. XII, Édition des Universités de Montpellier II et Paris VII

\begin{description}[leftmargin=5em, style=sameline]
\item[W(7)] (`Axiome du cube', ou `Lemme des cinq'.) Soit un diagramme commutatif cubique
\[
\begin{tikzcd}[row sep=small, column sep=small]
U_1\cap U'_1 \arrow[rr] \arrow[dr] \arrow[dd, "f_{U_1\cap U'_1}"'] & & U_1 \arrow[dr] \arrow[dd, "f_{U_1}"{swap,pos=.25}] & \\
& U'_1 \arrow[rr, crossing over] & & X_1 \arrow[dd, "f"] \\
U_2\cap U'_2 \arrow[rr] \arrow[dr] & & U_2 \arrow[dr] & \\
& U'_2 \arrow[rr] \arrow[from=uu, crossing over, "f_{U'_1}"{pos=.25}] & & X_2
\end{tikzcd}
\]
représentant un morphisme (flèches verticales) d'un

\msp{page 77}
carré cocartésien\quad $\vcenter{\hbox{\begin{tikzcd}
U_1\cap U'_1 \arrow[r] \arrow[d] & U_1 \arrow[d] \\
U'_1 \arrow[r] & X_1
\end{tikzcd}}}$\quad du type envisagé dans W(7 bis) ($X_1$ réunion de deux ouverts $U_1$, $U'_1$) dans un autre (à indices 2). Considérons les conditions suivantes.

\begin{itemize}
\item[a)] $f_{U_1\cap U'_1}\in W$.
\item[b)] $f_{U_1}$ et $f_{U'_1}$ sont \texttt{[dans]} $W$.
\item[c)] $f\in W$.
\end{itemize}

Alors deux parmi ces conditions impliquent la troisième. En d'autres termes, quand l'une de ces conditions est satisfaite, les deux autres sont équivalentes entre elles.
\end{description}

\textbf{NB} Cet axiome est équivalent à l'axiome dual, en termes d'objets $X_1$, $X_2$ représentés chacun comme réunion de deux fermés.

\textbf{Corollaire.} \textit{Supposons \textup{W(7)} satisfait. Alors le corollaire à \textup{W(7 bis)} (axiome qu'on n'a pas à utiliser ici), sous sa forme complétée (cf. le NB qui le suit) est valable : deux parmi les conditions a), b), c) impliquent la troisième.}

Il suffit de regarder l'unique morphisme du diagramme $(*)$, p.~76, dans le diagramme trivial\quad $\vcenter{\hbox{\begin{tikzcd}[column sep=small]
& e \arrow[dl] \arrow[dr] & \\
e \arrow[dr] & & e, \arrow[dl] \\
& e &
\end{tikzcd}}}$\quad correspondant au cas où $X$ est la catégorie finale $e$, et $U$, $U'$ égaux à $X$. Alors le corollaire n'est autre que W(7) appliqué à ce cas particulier.

Cet argument fait sentir que W(7) est sans doute plus fort que W(7 bis). Incité par ce sentiment, je m'aperçois avec plaisir qu'il est fondé :
% foot: Version 7 décembre 2002

% ===================== page 15 =====================
\printedpage{57}
% head: A. Grothendieck : Dérivateurs, Ch. XII, Édition des Universités de Montpellier II et Paris VII

\msp{page 78}
\textbf{Proposition 6.} \textit{L'axiome \textup{W(7)} (lemme des cinq) implique \textup{W(7 bis)} (axiome du carré cocartésien).}

On regarde le diagramme cubique suivant
\[
\begin{tikzcd}[row sep=small, column sep=small]
U\cap U' \arrow[rr, "i_0"] \arrow[dr, "i'_0"] \arrow[dd, "f_{U\cap U'}=i_0"'] & & U' \arrow[dr, "i'"] \arrow[dd, "f_{U'}=\mathrm{id}"{swap,pos=.25}] & \\
& U \arrow[rr, crossing over, "i"] & & X \arrow[dd, "f=\mathrm{id}_X"] \\
U' \arrow[rr, "\mathrm{id}"] \arrow[dr, "i'"] & & U' \arrow[dr, "i'"] & \\
& X \arrow[rr, "\mathrm{id}"] \arrow[from=uu, crossing over, "f_U=i"{pos=.25}] & & X,
\end{tikzcd}
\]
du carré cocartésien donné dans le carré cocartésien déduit de $X$ et de ses deux ouverts $X$, $U'$, d'intersection $U'$. Ici les deux flèches verticales $f$ et $f_{U'}$ sont des isomorphismes, donc \texttt{[dans]} $W$. Avec les notations de W(7), on a donc c) ($f\in W$), donc a) et b) sont équivalentes, i.e. $i_0\in W$ et $i\in W$ sont équivalentes.

\hfill q.e.d.

\textbf{Théorème 2.} \textit{Le localiseur fondamental $W_\infty$ satisfait à l'axiome \textup{W(7)}.}

C'est sûrement une conséquence du critère cohomologique d'\textsc{Artin-Mazur}, compte tenu de la suite exacte de Mayer-Vietoris pour l'un et l'autre des deux carrés cocartésiens, et les variantes non commutatives de Mayer-Vietoris. Je n'ai pas écrit la vérification en forme, et devrais sans doute réparer cette omission.

\msp{page 79}
\textbf{Remarque.} La version duale (au sens de \textsc{Quillen}) de l'axiome W(7 bis) du carré cocartésien, i.e. l'axiome correspondant pour les carrés cartésiens, n'est pas valable pour $W_\infty$. L'énoncé pertinent serait le suivant : soit
\[
\begin{tikzcd}
X & X' \arrow[l, "q"'] \arrow[d, "f'"] \\
S & S' \arrow[l, "p"']
\arrow[from=1-1, to=2-1, "f"']
\end{tikzcd}
\]
un carré cartésien, avec $f$ une \textit{$W$-fibration.} Alors $f\in W\Longleftrightarrow f'\in W$. Si $W$ satisfait à W(4c) (ce qui est le cas pour $W_\infty$), alors $f\in W$ implique bien que ses fibres sont $W$-asphériques, donc que c'est une $W$-équivalence universelle, i.e. $f'\in W\ \forall\, p:S'\longrightarrow S$. Mais pour l'implication inverse $f'\in W\Longrightarrow f\in W$, il faut supposer visiblement que
\[
\pi_0(p):\pi_0(S')\longrightarrow\pi_0(S)
\]
% foot: Version 7 décembre 2002


\endgroup
% END UNIT Der-ChXIIb-012/Der-ChXIIb-012_p011-015.tex
% BEGIN UNIT Der-ChXIIb-012/Der-ChXIIb-012_p016-020.tex
\begingroup\makeatletter%
\newcommand{\printedpage}[1]{\marginpar{\footnotesize p.~#1}}%
\newcommand{\msp}[1]{\marginpar{\footnotesize [#1]}}%
\newcommand{\circled}[1]{\tikz[baseline=(c.base)]\node[draw,circle,inner sep=1pt](c){#1};}%
\newcommand{\Ob}{\operatorname{Ob}}%
\makeatother\let\printedpage\derpp\let\msp\dermsp%
\let\mspin\DERmspin%


\printedpage{58}
% head: A. Grothendieck : Dérivateurs, Ch. XII, Édition des Universités de Montpellier II et Paris VII
% foot: Version 7 décembre 2002

est \textit{surjectif.} (P. ex. si $S'=\emptyset$, on a $X'=\emptyset$, donc $f'\in W$, sans que ça implique $f\in W$!) Le fait que dans le cas cocartésien de W(7 bis), $i_0\in W\Longrightarrow i\in W$ doit être considéré comme ‘allant de soi’ - autrement le carré ne mérite pas, du point de vue $W$-homotopique, le nom de ‘cocartésien’. Par contre, l’implication inverse $i\in W\Longrightarrow i_0\in W$ m’apparaît comme un trait très particulier du dérivatuer HOT$_W$ (quand il le possède), et notamment % [sic] dérivatuer -> dérivateur, misprint
de HOT $=$ HOT$_{W_\infty}$. Le début
\msp{page 80}
de la remarque signifie simplement que le dérivateur dual HOT$^\circ$ ne possède pas cette même propriété. Je dois renvoyer à ce sujet aux deux axiomes pour les dérivateurs correspondants à W(7 bis), W(7), ce que je ferai sans doute sous peu, de façon à donner un sens plus précis aux présents commentaires.

\subsection*{2.2 Application aux carrés cocartésiens déduits d’une immersion ouverte ou fermée.}

Les axiomes W(7), W(7 bis) (du cube, et du carré cocartésien) concernent des carrés cocartésiens de nature très particulière. Nous allons voir à présent comment en déduire des énoncés du même type pour des carrés cocartésiens beaucoup plus généraux, y compris pour des carrés qui n’ont rien de cocartésien (dans Cat) - et que nous appelerons % [sic] appelerons -> appellerons, misspelling
‘homotopiquement cocartésiens’ (relativement à $W$), ou simplement $W$-cocartésiens. (Il y aura également des carrés $W$-cartésiens un jour \dots) Dans le cas où $W=W_\infty$, je les appelerai aussi Hot-cocartésiens. % [sic] appelerai -> appellerai, misspelling

Je commencerai par le cas particulier d’un carré cocartésien associé à une immersion ouverte (diagramme ($*$), p. 74) :
\[
(*)\qquad\qquad
\begin{tikzcd}
U \arrow[r, hook, "i"] \arrow[d, "p"'] & X \arrow[d, "q"'] & T \arrow[l, hook', dotted, "j"'] \arrow[d, dotted] \\
U' \arrow[r, hook, "i'"] & X' & T'\simeq T. \arrow[l, hook', dotted, "j'"']
\end{tikzcd}
\]
\msp{page 81}

Le grand traquenard ici qu’il faut arriver à exorciser, c’est qu’il peut arriver que $i\in W$, sans qu’on ait $i'\in W$, et aussi $p\in W$ sans qu’on ait $q\in W$. Je m’en étais convaincu en 1983 en écrivant Pursuing Stacks (durant les semaines de ‘scratchwork’), et n’ai pas envie de rechercher un contre-exemple ici, s’il ne se présente de lui-même (ce qu’il fera sans doute).

Je vais donner un cas important où ces ennuis disparaissent :

\textbf{Proposition 6 :} \textit{Considérons les adhérences $\bar U$ et $\overline{U'}$ de $U$ et $U'$ dans $X$, $X'$ respectivement, et supposons que les inclusions}
\[
U\xrightarrow{i_0}\bar U,\qquad U'\xrightarrow{i'_0}\overline{U'}
\]
\printedpage{59}
% head: A. Grothendieck : Dérivateurs, Ch. XII, Édition des Universités de Montpellier II et Paris VII
% foot: Version 7 décembre 2002
\textit{soient dans $W$} (\footnote[64]{cf. th. 3 p. 107 pour des conditions générales assurant que cette hypothèse est satisfaite.}). \textit{Alors}
\[
\left\{
\begin{array}{llll}
\circled{\textit{a}}\ i\in W & \iff & i'\in W & \textit{si on a }\mathrm{W}(7\text{ bis})\\
\circled{\textit{b}}\ p\in W & \iff & q\in W & \textit{si on a }\mathrm{W}(7).
\end{array}
\right.
\]

\textsc{Démonstration}. Considérons d’abord le cas \circled{a} de $i$, $i'$. Soit
\[
k:\bar U\hookrightarrow X
\]
l’inclusion. Le triangle
\begin{tikzcd}
U \arrow[r, "i_0"] \arrow[rr, bend right=8, "i"'] & \bar U \arrow[r, "k"] & X
\end{tikzcd}
, où $i_0\in W$, montre que
\[
i\in W\iff k\in W.
\]

D’autre part, $X$ est la réunion des fermés $\bar U$, $T$, donc on a un carré cocartésien dans Cat
\[
(*)\qquad\qquad
\begin{tikzcd}
\bar U\cap T=\dot T \arrow[r, "k_0"] \arrow[d, "j_0"'] & T \arrow[d, "j"] \\
\bar U \arrow[r, "k"] & X.
\end{tikzcd}
\]

L’axiome W(7 bis) (sous sa forme duale (i.e. appliqué au carré formé par les catégories opposées \dots)) implique
\msp{page 82}
que
\[
k\in W\iff k_0\in W,
\]
d’où
\[
i\in W\iff (k_0:\dot T\longrightarrow T)\in W.
\]
De même on trouve, en appliqunt ceci à $(X',U')$ au lieu de $(X,U)$ % [sic] appliqunt -> appliquant, misprint
\[
i'\in W\iff (k'_0:\dot T'\longrightarrow T')\in W.
\]
Or on a $T\xrightarrow{\sim}T'$ (isomorphisme), et par la proposition 4 cet isomorphisme transforme $\dot T$ en $\dot T'$. Donc
\[
k_0\in W\iff k'_0\in W\qquad\quad\text{(trivialement)},
\]
d’où a).

Prouvons b). Le carré commutatif induit par ($*$), p. 80,
\printedpage{60}
% head: A. Grothendieck : Dérivateurs, Ch. XII, Édition des Universités de Montpellier II et Paris VII
% foot: Version 7 décembre 2002
\[
\begin{tikzcd}
U \arrow[r, hook, "i_0"] \arrow[d, "p"'] & \bar U \arrow[d, "\bar p"] \\
U' \arrow[r, hook, "i'_0"] & \overline{U'},
\end{tikzcd}
\]
où $i_0,i'_0\in W$ par hypothèse, montre que
\[
p\in W\iff\bar p\in W,
\]
donc tout revient à prouver que
\[
\bar p\in W\overset{?}{\iff}q\in W.
\]
Mais $q$ induit un homomorphisme du carré ($*$) déduit de $(X,\bar U,T)$ (avec $X=\bar U\cup T$), dans le carré similaire pour $X'$
\[
\begin{tikzcd}
\overline{U'}\cap T'=\dot T' \arrow[r, "k'_0"] \arrow[d, "j'_0"'] & T \arrow[d, "j'"] \\ % [sic] T -> T', the square for X' should have the corner T' (although T'\simeq T)
\overline{U'} \arrow[r, "k'"] & X',
\end{tikzcd}
\]
les flèches verticales étant notées $q_{\dot T}$, $q_T$, $\bar p$, $q$.
\msp{page 83}
On sait que $q_{\dot T}$, $q_T$ sont des isomorphismes, donc dans $W$, donc le ‘lemme des cinq’ W(7) nous assure que $\bar p\in W\iff q\in W$, q.e.d.

\textbf{Commentaire.} Cette proposition nous amène à chercher des conditions sur un ouvert $U$ d’un $X$, qui assurent que l’inclusion $i_0:U\longrightarrow\bar U$ soit \texttt{[dans]} $W$, et \textit{que cette propriété soit conservée par tout cochangement de base $U\longrightarrow U'$.} De plus, nous recherchons des conditions qui soient de nature locale sur $\bar U$, plus exactement sur $\bar U\cap T=\dot U=\dot T$. Nous pouvons motiver ce souci (a posteriori) en disant que nous voudrions que la propriété $i_0\in W$ soit conservée, non seulement par cochangement de base quelconque $U\longrightarrow U'$, mais aussi, de plus, par passage de $X$ à toute sous-catégorie pleine $X_1$ contenant $U$, donnant donc naissance à une situation
\[
U\overset{i_Y}{\hookrightarrow}X_1,\qquad U\xrightarrow{(i_Y)_0}\bar U^{X_1}=\bar U\cap X_1
\]
\texttt{[}$Y=X_1$\texttt{?]}, similaire à la situation de départ. Plus généralement encore, nous voudrions qu’elle soit conservée en passant du triple $(X,U,T)$ à un triple “induit” $(X_1,U,T_1)$, par une flèche quelconque $g:T_1\longrightarrow T$ de but $T$ dans Cat, donnant naissance au couple de flèches
\[
U\xrightarrow{\mathrm{id}}U,\ g:T_1\longrightarrow T
\]
(cf. pages 60, 61). Le cas précédent est celui où
\printedpage{61}
% head: A. Grothendieck : Dérivateurs, Ch. XII, Édition des Universités de Montpellier II et Paris VII
% foot: Version 7 décembre 2002
\msp{page 84}
$T_1$ est une sous-catégorie pleine de $T$, $g$ la flèche d’inclusion, Notons à ce sujet la proposition soritale % [sic] inclusion, Notons -> inclusion. Notons, comma before a capital letter

\textbf{Proposition 7.} \textit{Conditions équivalentes sur l’ouvert $U$ d’un objet $X$ de} Cat (moyennant W(7 bis)) :

\begin{enumerate}[label={}, leftmargin=*]
\item[\textit{a)}] \textit{L’inclusion $i_0:U\longrightarrow\bar U$ est $W$-asphérique, i.e. pour tout $t\in\dot U=\bar U\backslash U$, la catégorie $U/t$ est $W$-asphérique.}
\item[\textit{b)}] \textit{L’inclusion $i_0:U\longrightarrow\bar U$ est dans $W$, et cette condition reste satisfaite, quand on remplace $(X,U)$ par $(X_1,U)$, où $(X_1,U)$ est déduit de $(X,U)$ au moyen d’une flèche $T_1\longrightarrow T(=X\backslash U)$ dans} Cat\textit{, comme rappelé ci-dessus.}
\item[\textit{b$'$)}] \textit{Comme b), dans le cas où $T_1$ est la catégorie ponctuelle $e$, et où de plus son image est dans $\dot U$.}
\end{enumerate}

\textsc{Démonstration}. Il est évident que si dans b), b$'$) on remplace la condition “$i_0\in W$” par la condition plus forte ‘$i_0$ est $W$-asphérique’, les quatre \texttt{[plutôt trois]} conditions a) à b$'$) sont équivalentes. Il reste à prouver que la condition que les $U/t$ soient $W$-asphériques ($t\in\bar U$ - mais c’est trivial pour $t\in U$, donc il suffit de l’exprimer pour $t\in\bar U\backslash U=\dot U$) est déjà impliquée par la plus faible b$'$ et on devine qu’on a ceci :

\msp{page 85}

\textbf{Corollaire.} (Moyennant W(7 bis).) \textit{Soit $X$ dans} Cat\textit{, $U$ un ouvert de $X$ tel que $X\backslash U=\overline{\{t\}}$ soit la catégorie finale. Alors l’inclusion $U\xrightarrow{i}X$ est dans $W$ si et seulement si elle est $W$-asphérique, i.e. si et seulement si $U/t$ est $W$-asphérique (où $t$ est l’unique élément de $T$).}

Nous démontrons ultérieurement ce corollaire (\footnote[65]{cf. cor. 4 de la proposition 15, page 149.}). La proposition 7 ne nous servira pas d’ici là - elle n’aura qu’un rôle purement esthétique à jouer, pour justifier notre propos délibéré de chercher des conditions sur $i:U\longrightarrow X$ (pour assurer que $U\longrightarrow\bar U$ est dans $W$ et le reste, quand on remplace $(X,U)$ par d’autres couples $(X_1,U_1)$ qui s’en déduisent par certaines opérations \dots) qui
\msp{page 86}
soient de \textit{nature locale} sur $\bar U$ (en fait, en les points de $\dot U$), et qui en tout état de cause entrainent la \textit{$W$-asphéricité} de $U\longrightarrow\bar U$.

J’avais espéré que cette condition d’asphéricité (suggérée par la notion ‘d’immersion de Dwyer’) serait stable par cochangement de base quelconque $U\longrightarrow U'$, ce qui aurait assuré en même temps qu’après une telle opération, l’immersion ouverte $i':U'\longrightarrow X'$ obtenue satisfait encore à la condition que $i'_0:U'\longrightarrow\overline{U'}$ est dans $W$. Nous verrons que cette attente n’est pas fondée - la condition de $W$-asphéricité n’est \textit{pas} stable par cochangement de base, et elle ne suffit \textit{pas} même à assurer qu’après un cochangement de base, $i'_0:U'\longrightarrow\overline{U'}$ soit
\printedpage{62}
% head: A. Grothendieck : Dérivateurs, Ch. XII, Édition des Universités de Montpellier II et Paris VII
% foot: Version 7 décembre 2002
dans $W$. Il faudra une notion plus forte, qui est suggérée par le cas où $T=\{t\}$ est la catégorie finale et contenue dans $\bar U$ (cas auquel nous avons essayé de nous ramener dans la proposition 7). Dans ce cas, on l’a dit (corollaire précédent), la condition que $i:U\longrightarrow\bar U=X$ soit dans $W$, équivaut à la $W$-asphéricité de $i$, i.e. à la $W$-asphéricité de $U/t$, et de même pour $i':U'\longrightarrow\overline{U'}=X'$ : elle est dans $W$ si et seulement si $U'/t'$ est $W$-asphérique, où
\msp{page 87}
$t'$ est l’unique élément de $T'(\simeq T)=X'\backslash U'$. Donc il faut des conditions sur $U/t$ qui assurent que non seulement $U/t$, mais aussi toutes les catégories similaires $U'/t$ \texttt{[plutôt} $U'/t'$\texttt{]}, soient $W$-asphériques. Dans le cas général où $T$ est quelconque, c’est cette condition-là que nous allons imposer aux $U/t$ ($t\in\Ob T$), et il s’avérera que cette fois-ci, elles sont stables \texttt{[plutôt elle est stable]} par tout cochangement de base $U\xrightarrow{p}U'$, i.e. que les $U'/t'$ ($t'\in\Ob T'$) satisferont alors à la même condition. C’est cette condition que nous allons à présent reformuler.

\subsection*{2.3 Objets totalement $W$-asphériques de \textnormal{Cat}, et morphismes totalement $W$-asphériques.}

Reprenons la situation
\[
U\overset{i}{\hookrightarrow}X\overset{j}{\hookleftarrow}T,\qquad\qquad\text{avec }T=\{t\}\text{, catégorie finale, }t\in\bar U\text{, i.e. }X=\bar U,
\]
et celle d’un cochangement de base $U\longrightarrow U'$, donnant naissance à un carré cocartésien dans Cat
\[
(*)\qquad\qquad
\begin{tikzcd}
U \arrow[r, hook, "i"] \arrow[d, "p"'] & X \arrow[d, "q"'] & T=\{t\} \arrow[l, hook', dotted, "j"'] \arrow[d, dotted, "\mathrm{iso}"] \\
U' \arrow[r, hook, "i'"] & X' & T'=\{t'\}. \arrow[l, hook', dotted, "j'"']
\end{tikzcd}
\]
La donnée de recollement $H$ entre $U$ et $T$, $H:\underbrace{U^\circ\times T}_{\simeq U^\circ}\longrightarrow\mathrm{Ens}$, peut s’interpréter comme un préfaisceau sur $U$, $H\in\Ob U^\wedge:U^\circ\longrightarrow\mathrm{Ens}$,
\msp{page 88}
et on a
\[
U/t\simeq U/H,
\]
où $U/H$ est interprété en termes du plongement canonique
\[
U\hookrightarrow U^\wedge.
\]
De même, le recollement entre $U'$ et $T'$ sera donné par un préfaisceau
\[
H'\in\Ob U^{\prime\wedge}:U^{\prime\circ}\longrightarrow\mathrm{Ens},
\]


\endgroup
% END UNIT Der-ChXIIb-012/Der-ChXIIb-012_p016-020.tex
% BEGIN UNIT Der-ChXIIb-012/Der-ChXIIb-012_p021-025.tex
\begingroup\makeatletter%
\newcommand{\printedpage}[1]{\marginpar{\footnotesize p.~#1}}%
\newcommand{\msp}[1]{\marginpar{\footnotesize [#1]}}%
\newcommand{\cercle}[1]{\tikz[baseline=(c.base)]\node[draw,circle,inner sep=1pt](c){#1};}%
\newcommand{\defeq}{\overset{\text{déf}}{=}}%
\newcommand{\upi}{\underline{\pi}}%
\newcommand{\Cat}{\mathrm{Cat}}%
\newcommand{\Ob}{\mathrm{Ob}\,}%
\makeatother\let\printedpage\derpp\let\msp\dermsp%
\let\mspin\DERmspin%


\printedpage{63}
% head: A. Grothendieck : Dérivateurs, Ch. XII, Édition des Universités de Montpellier II et Paris VII
et on aura
\[ U'/t' \simeq U/H'. \]
La relation entre $H'$ et $H$ a été explicitée au §1, comme
\[ H' = p_!H, \qquad \text{où } p_! : U^{\wedge} \longrightarrow U'^{\wedge} \]
est l'adjoint à gauche de $p^* : U'^{\wedge} \longrightarrow U^{\wedge}$, donné par
\[ p^*F' = F' \circ p. \]
On a
\[ p_!(H)(u') \simeq \pi_0(u'\backslash U/H). \]

\textbf{Définition 1.} \cercle{a} Soit $U$ un objet de Cat, $H \in \Ob U^{\wedge}$ un préfaisceau sur $U$. On dit que $H$ est \textit{totalement $W$-asphérique} si pour toute flèche $p : U \longrightarrow U'$ de Cat, $p_!(H)$ dans $U'^{\wedge}$ est un objet $W$-asphérique, i.e. $U'/p_!H$ est $W$-asphérique. \cercle{b} On dit qu'un objet $V$ de \msp{page 89} Cat est \textit{totalement $W$-asphérique,} si le foncteur final $e_V$ sur $V$ ($v \longmapsto \{\emptyset\}$) est totalement $W$-asphérique, i.e. si pour toute flèche $f : V \longrightarrow V'$ dans Cat, le préfaisceau
\[ \pi_0(f) \defeq f_!(e_V) \defeq (v' \longmapsto \pi_0(v'\backslash V)) \]
sur $V'$ est $W$-asphérique, i.e. si la catégorie
\[ \upi_0(f) \ \text{ ou } \ \upi_0(V/V') \defeq V'/\pi_0(f) \]
est $W$-asphérique. \texttt{[Cette notion est différente de la notion de totale asphéricité de Pursuing Stacks.]}

Notons que la définition des foncteurs $p_!$ (ou $f_!$) comme adjoints, et leur `calcul' en termes de $\pi_0$ de certaines catégories, ne dépend pas du localiseur $W$. Donnons quelques précisions à ce sujet. On a un morphisme d'adjonction dans $V^{\wedge}$
\[ f^*f_!(F) \longleftarrow F, \]
d'où
\[ V/F \longrightarrow V/f^*f_!(F) \simeq V \times_{V'} (V'/f_!(F)), \tag*{$(*)$} \]
le dernier isomorphisme provenant du fait que l'on a commutativité, à isomorphisme canonique près, dans le diagramme de foncteurs
\[
\begin{tikzcd}[column sep=14em, row sep=large]
V'^{\wedge} \arrow[r, "f^*=(F'\longmapsto F'\circ f)"] \arrow[d, "F'\longmapsto V'/F'"'] & V^{\wedge} \arrow[d, "F\longmapsto F/V"] \\
\Cat/V' \arrow[r, "f^*_{\Cat}=(W'\longmapsto W'\times_{V'}V)"] & \Cat/V,
\end{tikzcd}
\]
% foot: Version 7 décembre 2002
\printedpage{64}
% head: A. Grothendieck : Dérivateurs, Ch. XII, Édition des Universités de Montpellier II et Paris VII
\msp{page 90}
en d'autres termes, l'opérateur $f^*$ sur les préfaisceaux s'interprète, en termes géométriques, comme l'opération $W' \longmapsto W' \times_{V'} V$ de changement de base par $V \longrightarrow V'$, sur les catégories au dessus de $V'$ qui sont fibrées à fibres discrètes (propriété qui se conserve par changement de base).

Ceci dit, la donnée d'une flèche $(*)$ équivaut, par la propriété universelle du changement de base, à celle d'une flèche (canonique en l'occurrence)
\[ f/F : V/F \longrightarrow V'/f_!(F) \]
rendant commutatif le carré
\[
\begin{tikzcd}[column sep=large, row sep=huge]
V/F \arrow[r, "\text{oubli de }F", "\substack{\text{ou foncteur}\\ \text{source}}"'] \arrow[d, "f/F"'] & V \arrow[d, "f"] \\
V'/(F' = f_!(F)) \arrow[r, "\text{oubli de }F'"] & V'.
\end{tikzcd}
\tag*{$(*)$}
\]
Dans le cas qui nous intéresse surtout ici, $F = e_V$, $V/F \longrightarrow V$ est un isomorphisme, et on obtient une factorisation canonique de foncteurs \msp{page 91}
\[
\begin{tikzcd}[row sep=large]
 & V \arrow[dl, "\tilde{f}=f/e_V"'] \arrow[dd, "f"] \\
\begin{array}{r} \upi_0(V/V') \\ = V'/f_!(e_V) \end{array} \arrow[dr, "\text{can.}"'] & \\
 & V'.
\end{tikzcd}
\tag*{$(**)$}
\]
Explicitons le $V'$-morphisme $\tilde{f}$ sur les fibres,
\[ \tilde{f}_{v'} : V_{v'} \longrightarrow \upi_0(V/V')_{v'} \defeq \pi_0(v'\backslash V), \tag*{$(***)$} \]
où dans le dernier membre, on a identifié l'ensemble $\pi_0(v'\backslash V)$ à la catégorie discrète qu'il définit. Pour expliciter cette flèche $(***)$, notons qu'on a une inclusion canonique
\[ V_{v'} \longrightarrow v'\backslash V \]
et pour toute catégorie $W$, un morphisme canonique (d'adjonction) $W \longrightarrow \pi_0(W)$, d'où un composé
\[ V_{v'} \longrightarrow v'\backslash V \longrightarrow \pi_0(v'\backslash V), \]
qui n'est autre que $\tilde{f}_{v'}$ dans $(***)$.

Il doit être facile de montrer que la factorisation $(**)$ de $f$ est caractérisée, à isomorphisme unique près, en tant que
% foot: Version 7 décembre 2002
\printedpage{65}
% head: A. Grothendieck : Dérivateurs, Ch. XII, Édition des Universités de Montpellier II et Paris VII
\msp{page 92}
factorisation dans $\Cat/V'$
\[
\begin{tikzcd}[column sep=small]
 & \tilde{V} \arrow[dr, "\pi"] & \\
V \arrow[ur, "\tilde{f}"] \arrow[rr, "f"'] & & V',
\end{tikzcd}
\]
par les propriétés suivantes :
\begin{enumerate}[label={}, leftmargin=*]
\item[a)] $\pi$ est fibrée à fibres discrètes, i.e. un isomorphisme local (i.e. $\tilde{V}$ est isomorphe dans $\Cat/V'$ à une catégorie $V/F'$, avec $F'$ dans $V'^{\wedge}$). % [sic] V/F' -> V'/F', une catégorie au-dessus de V' définie par F' dans V'^∧ est de la forme V'/F'
\item[b)] $\tilde{f}_!(e_V) = e_{\tilde{V}}$, i.e. pout tout $\tilde{v} \in \Ob \tilde{V}$, la catégorie $\tilde{v}\backslash V$ est 0-connexe : les colocalisés de $V$ au dessus de $\tilde{V}$ sont 0-connexes. % [sic] pout -> pour, faute de frappe
\end{enumerate}

Notons ici, en passant, que pour \textit{tout} morphisme $\pi : W' \longrightarrow W$ dans Cat qui soit fibré à fibres discrètes, i.e. tel que $W'$ soit isomorphe dans $\Cat/W$ à une catégorie de la forme $W/F$ ($F \in \Ob W^{\wedge}$), on a
\[ \pi_!(e_{W'}) = F, \qquad \text{i.e. } W' \underbrace{\ \simeq\ }_{\text{dans } \Cat/W} W/\pi_!(e_{W'}). \tag*{$(*)$} \]
Il faut vérifier que le foncteur
\[ w \longmapsto \pi_0(\underbrace{w\backslash(W/F)}_{\defeq w\backslash W/F}) \]
est canoniquement isomorphe à $F$. Mais l'inclusion
\[ \underbrace{W'_w}_{= F(w) \text{ (catégorie discrète)}} \qquad \longrightarrow w\backslash W' \]
\msp{page 93}
définit bien comme tantôt
\[ F(w) = W'_w \longrightarrow \pi_0(w\backslash W'), \]
il faut prouver que c'est bijectif. Or, $W'$ étant $W_\infty$-lisse sur $W$, on sait que l'inclusion $W'_w \hookrightarrow w\backslash W'$ est dans $W_\infty$, a fortiori induit un isomorphisme pour les $\pi_0$.

Ceci posé, si on a une factorisation de $f : V \longrightarrow V'$ en $f = \pi\tilde{f}$ ayant les propriétés a) b) de la page précédente, on voit déjà que l'on a
\[ f_!(e_V) \simeq \pi_!(\, \underbrace{\tilde{f}_!(e_V)}_{\simeq e_{\tilde{V}} \text{ par b)}} \,) \simeq \pi_!(e_{\tilde{V}}), \]
ce qui déjà détermine $\tilde{V}$ à isomorphisme canonique près (grâce à a) comme
\[ \tilde{V} \simeq V'/f_!(e_V) \defeq \upi_0(V'/V), \] % [sic] \underline{\pi}_0(V'/V) -> \underline{\pi}_0(V/V'), notation employée partout ailleurs (Définition 1, (**))
% foot: Version 7 décembre 2002
\printedpage{66}
% head: A. Grothendieck : Dérivateurs, Ch. XII, Édition des Universités de Montpellier II et Paris VII
et il ne doit pas être bien dur de se convaincre que $\tilde{f}$ n'est autre que la flèche canonique de la page 91.

\textbf{Proposition 8.} \textit{Pour qu'un objet $H$ de $U^{\wedge}$ soit totalement $W$-asphérique, il suffit que l'objet $U/H$ de} Cat \textit{soit totalement $W$-asphérique.}

\msp{page 94}
Supposons en effet cette dernière condition satisfaite et prouvons que $H$ est totalement $W$-asphérique, i.e. que pour tout $p : U \longrightarrow U'$, $U'/p_!(H)$ soit $W$-asphérique.
\[
\begin{tikzcd}[row sep=large]
V \defeq U/H \arrow[r, "\alpha"] \arrow[dr, "f"'] & U \arrow[d, "p"] \\
 & U'
\end{tikzcd}
\]
Soit $f$ le foncteur composé du diagramme ci-contre \texttt{[ci-dessus]}, on a $H \simeq \alpha_!(e_V)$ (par $(*)$ p. 92), d'où
\[ p_!(H) \simeq p_!\alpha_!(e_V) \simeq (p\alpha)_!(e_V) \simeq f_!(e_V), \]
qui est $W$-asphérique par hypothèse sur $V = U/H$.

La réciproque de \texttt{[la]} prop. 8 est-elle valable, i.e. si $H$ est totalement $W$-asphérique, en est-il de même de $V = U/H$? La difficulté provient du fait que cela signifie que pour \textit{toute} flèche $f : V \longrightarrow V'$ dans Cat de source $V$, on ait $f_!(e_V)$ $W$-asphérique, alors que l'hypothèse équivaut à le postuler seulement pour les flèches qui se factorisent par $U$. Notons que si on ne fait pas cette hypothèse sur $f$, mais qu'on suppose au lieu de cela que $f$ s'insère dans un carré commutatif \msp{page 95}
\[
\begin{tikzcd}[row sep=large]
V \arrow[r, "\alpha"] \arrow[d, "f"'] & U \arrow[d, "p"] \\
V' \arrow[r, "\alpha'"] & U',
\end{tikzcd}
\]
avec $\alpha'$ lui aussi un isomorphisme local, alors on peut conclure que $\upi_0(f) = \upi_0(V/V')$ est bien asphérique. En effet, on a un isomorphisme canonique dans $\Cat/U'$,
\[ \upi_0(\alpha'f) \simeq \upi_0(f), \qquad \text{i.e. } \ \upi_0(V/V') \simeq \upi_0(V/U'), \]
du fait p. ex. que la factorisation canonique de $f$
% foot: Version 7 décembre 2002
\printedpage{67}
% head: A. Grothendieck : Dérivateurs, Ch. XII, Édition des Universités de Montpellier II et Paris VII
\[
\begin{tikzcd}[row sep=normal]
 & V \arrow[dl, "\tilde{f}"'] \arrow[dd, "f"] & \\
\upi_0(V/V') \arrow[dr, "\pi"'] & & \\
 & V' \arrow[dr, dotted, "\alpha'"'] & \\
 & & U'
\end{tikzcd}
\]
fournit une factorisation du même type pour $\alpha'f$, en
\[ \alpha'f = (\alpha'\pi) \circ \tilde{f}. \]
Ceci dit, on peut appliquer l'hypothèse à $p$ et on trouve bien la $W$-asphéricité de $\upi_0(f) \simeq \upi_0(\alpha'f) = \upi_0(p\alpha)$.

\msp{page 96}
L'idée naturelle, pour une réciproque à \texttt{[la]} prop. 8, serait de prendre pour $U'$ la somme amalgamée de $U$ et $V'$ sur $V$, ce qui marchera pourvu qu'on sache que $\alpha'$ est bel et bien un isomorphisme local tout comme $\alpha$. C'est bien le cas p. ex. si $V \longrightarrow U$ est une immersion ouverte, i.e. $H$ un sous-préfaisceau du préfaisceau final. Mais je crains que ce ne soit pas vrai en général - ça a l'air un peu trop beau pour être vrai! C'est en effet déjà faux dans le cas
\[
\begin{tikzcd}[row sep=huge]
V = E_G \arrow[r, "\alpha\ \text{can}"] \arrow[d, "f"'] & B_G = U \\
V' = B_{G'}, &
\end{tikzcd}
\]
dans le cas p. ex. où $G = \{\pm 1\}$, groupe à deux éléments, et où l'élément $\alpha$ de $G'$, image par $f$ de l'une des deux flèches non identiques de $E_G$ est $\neq 1$ \texttt{[plutôt de carré} $\neq 1$\texttt{]}. Si $H'$ est le sous-groupe invariant engendré par $\alpha$ \texttt{[plutôt par} $\alpha^2$\texttt{]} dans $G'$, la somme amalgamée s'identifie à $B_{G'/H'}$, or $B_{G'} \longrightarrow B_{G'/H'}$ n'est \textit{pas} un isomorphisme local (car $(B_{G'})/t_0 \longrightarrow (B_{G'/H'})/t'_0$ n'est pas bijectif sur les objets, s'identifiant à $G' \overset{q}{\longrightarrow} G'/H' \ \ldots$).

Mais d'autre part la condition de totale $W$-asphéricité sur un objet $H$ de $U^{\wedge}$ paraît déjà à tel point exigeante, que je me vois mal con\msp{page 97}struire un exemple où elle est satisfaite, sans que $U/H$ elle-même ne soit totalement $W$-asphérique.

L'existence d'objets totalement $W$-asphériques dans Cat n'est pas tellement évidente a priori. Qu'il y en ait un bon paquet est assuré par la
% foot: Version 7 décembre 2002


\endgroup
% END UNIT Der-ChXIIb-012/Der-ChXIIb-012_p021-025.tex
% BEGIN UNIT Der-ChXIIb-012/Der-ChXIIb-012_p026-030.tex
\begingroup\makeatletter%
\newcommand{\printedpage}[1]{\marginpar{\footnotesize p.~#1}}%
\newcommand{\msp}[1]{\marginpar{\footnotesize [#1]}}%
\newcommand{\upi}{\underline{\pi}}%
\newcommand{\uCh}{\underline{\mathrm{Ch}}}%
\newcommand{\Ob}{\operatorname{Ob}}%
\newcommand{\Hom}{\operatorname{Hom}}%
\newcommand{\bDelta}{\boldsymbol{\Delta}}%
\makeatother\let\printedpage\derpp\let\msp\dermsp%
\let\mspin\DERmspin%

\setcounter{footnote}{65}

\printedpage{68}
% head: A. Grothendieck : Dérivateurs, Ch. XII, Édition des Universités de Montpellier II et Paris VII

\noindent\textbf{Proposition 9.}

\begin{enumerate}[label=\textit{\alph*)}]
\item \textit{Tout objet $V$ de} Cat \textit{ayant un objet final est totalement $W$-asphérique - donc tout préfaisceau d'ensembles $H$} représentable \textit{sur une catégorie $U$ (la représentabilité signifiant justement que $U/H$ a un objet final) est totalement $W$-asphérique.}

\item (\footnote{!! Sans doute \textit{faux} tel quel, cf. rectification page suivante.}) \textit{(Supposons} W(7 bis) \textit{satisfait.) Soient $V$ dans} Cat, \textit{$V'$, $V''$ deux ouverts de réunion $V$. Si $V'$, $V''$, $V'\cap V''$ sont totalement $W$-asphériques, il en est de même de $V$.}
\end{enumerate}

\noindent\textsc{Démonstration} de a). Soit $V$ avec objet final et $f : V \longrightarrow V'$, d'où la factorisation
\[
\begin{tikzcd}[column sep=small, row sep=small]
 & \tilde{V} = \upi_0(V/V') \arrow[dr, "\pi"'] & \\
V \arrow[ur, "\tilde{f}"] \arrow[rr, "f"'] & & V'
\end{tikzcd}
\]
de $f$. On démontre sans mal (\footnote{Voici une démonstration simple : On a pour toute $V$ dans Cat $e_V = \varinjlim_{v \text{ dans } V} v$ (limite inductive du foncteur d'inclusion de $V$ dans $V^\wedge$), car $\forall\, x$ dans $V$
\[
\varinjlim_{v \text{ dans } V} \Hom(x,v) \simeq \pi_0(x\backslash V) = e.
\]
Comme $f_!$ commute aux $\varinjlim$ , on a donc $f_!(e_V) = \varinjlim_{v \text{ dans } V} f(v)$. Si $v_0$ est un objet final de $V$, cela montre que $f_!(e_V) = f(v_0)$, q.e.d.}) (histoire de cofinalité) ceci :

\noindent\textbf{Corollaire.} \textit{L'image par $\tilde{f}$ d'un objet final de $V$ est un objet final de $\tilde{V} = \upi_0(V/V')$.}

Comme un objet de Cat avec objet final est $W$-asphérique par W(3), cela établit a).

Prouvons b). J'ai envie de dire que si on a $f : V \longrightarrow W$ dans Cat, alors $\upi_0(V/W)$ est réunion de deux ouverts isomorphes à $\upi_0(V'/W)$, $\upi_0(V''/W)$, dont l'intersection est isomorphe à $\upi_0(V'\cap V''/W)$. Cela signifierait ni plus ni
\msp{page 98}
moins que ceci : $f_!$ transforme les monomorphismes $V' \longrightarrow V$, $V'' \longrightarrow V$ (interprétés comme flèches dans $V^\wedge$, de but $e_V$) en monomorphismes, et le produit fibré $V'\cap V'' = V' \times_V V''$ en le produit fibré correspondant dans $W^\wedge$. Mais le foncteur $f_!$ commute aux $\varinjlim$ quelconques (étant un adjoint à gauche), et je ne vois aucune raison a priori pour qu'il commute (disons) aux produits fibrés (ce qui impliquerait aussi qu'il transforme monomorphismes en itou). Si je me rappelle bien, c'est ultrafaux en général. Donc je doute fort à présent que l'énoncé b) soit vrai, sans hypothèse supplémentaire. Pour assurer la monomorphie, on songe à imposer l'existence de \textit{rétractions} de $V$ sur $V'$, $V''$ et de $V'$ ou $V''$ sur $V'\cap V''$.

Je viens de vérifier la chose, et donne le
% foot: Version 7 décembre 2002
\printedpage{69}
% head: A. Grothendieck : Dérivateurs, Ch. XII, Édition des Universités de Montpellier II et Paris VII
\textbf{Corollaire 1} (rectification de \texttt{[la]} prop. 9, partie b) (\footnote{NB Il y a un énoncé légèrement plus général, en partant d'une catégorie $U$ et d'un diagramme
$\begin{tikzcd}[ampersand replacement=\&, row sep=small]
V'\cap V'' \arrow[r] \arrow[d] \& V' \arrow[d] \\
V'' \arrow[r] \& V
\end{tikzcd}$
\textit{dans} $U^\wedge$ : si $V'$, $V''$, $V'\cap V''$ sont totalement $W$-asphériques, il en est de même de $V$.}) : \textit{L'énoncé est valable si on suppose qu'il existe une} rétraction \textit{$r'_0$ de $V''$ sur $V'\cap V''$, et de même une rétraction $r''_0$ de $V'$ sur $V'\cap V''$.} (\footnote{Corollaire rectifié faux également - les rétractions ne sont pas des flèches de $V^\wedge$, et ne passent dans $p_!$.}) % [resolved: see DER_XIIb-a_corrections.csv]
\[
\begin{tikzcd}
V'\cap V'' \arrow[r, "i''_0"'] \arrow[d, "i'_0"] & V' \arrow[d, "i'"] \arrow[l, dotted, bend right=50, "r''_0"'] \\
V'' \arrow[r, "i''"] \arrow[u, dotted, bend left=50, "r'_0"] & V \arrow[u, dotted, bend right=50, "r'"'] \arrow[l, dotted, bend left=50, "r''"]
\end{tikzcd}
\]

(\textbf{NB} Nous verrons plus tard qu'il suffit de l'existence d'\textit{une} de ces rétractions pour assurer la conclusion.)

\textsc{Démonstration}. Comme $V$ est somme amalgamée de $V'$ et $V''$, il existe une rétraction $r'$ de $V$ sur $V'$ (donnée par $\mathrm{id}_{V'}$ sur $V'$, par $i''_0r'_0$ sur $V''$) qui induit $r'_0$ sur $V''$, et de même il existe
\msp{page 99}
une rétraction $r''$ de $V$ sur $V''$, induisant $r''_0$ sur $V'$. L'existence de ces rétractions assure déjà que pour tout foncteur $f : V \longrightarrow W$, le diagramme transformé de
\[
\begin{tikzcd}[row sep=large]
V'\cap V'' \arrow[r] \arrow[d] & V' \arrow[d] \\
V'' \arrow[r] & V
\end{tikzcd}
\]
(intérprété comme diagramme dans $V^\wedge$) est formé de monomorphismes (\footnote{\textbf{faux,} déjà dans le cas où $V = \Phi$, $V'$, $V''$ ses deux ouverts canoniques isomorphes à $\bDelta^1$, et $V'{\cap}V'' = e$.}) % [sic] intérprété -> interprété, misspelt accent
\[
\begin{tikzcd}[row sep=large]
f_!(V'\cap V'') \arrow[r, hook] \arrow[d, hook] & f_!(V') \arrow[d, hook] \\
f_!(V'') \arrow[r, hook] & f_!(V),
\end{tikzcd}
\]
et de plus que $f_!(V') \cap f_!(V'')$ (qui contient $f_!(V'\cap V'')$) lui est en fait \textit{égal.} Il en résulte que l'on a un diagramme à la fois cartésien et cocartésien d'immersions ouvertes dans Cat
% foot: Version 7 décembre 2002
\printedpage{70}
% head: A. Grothendieck : Dérivateurs, Ch. XII, Édition des Universités de Montpellier II et Paris VII
\[
\begin{tikzcd}[row sep=large]
W/f_!(V'\cap V'') \arrow[r, hook] \arrow[d, hook] & W/f_!(V') \arrow[d, hook] \\
W/f_!(V'') \arrow[r, hook] & W/\Bigl(f_!(V) \overset{\text{déf}}{=} f_!(e_V)\Bigr),
\end{tikzcd}
\]
lequel s'identifie aussi à
\[
\begin{tikzcd}[row sep=large]
\upi_0(V'\cap V''/W) \arrow[r, hook] \arrow[d, hook] & \upi_0(V'/W) \arrow[d, hook] \\
\upi_0(V''/W) \arrow[r, hook] & \underline{\pi_0}(V/W).
\end{tikzcd}
\]

\msp{page 100}
L'hypothèse sur $V'$, $V''$, $V'\cap V''$ implique que les trois sommets du carré autres que $\upi_0(V/W)$ sont $W$-asphériques, donc par W(7 bis) il en est de même de ce dernier, q.e.d.

\textbf{Corollaire 2} (\footnote{Plus généralement, si $I$ est une catégorie \textit{ordonnée finie,} ayant un objet initial, alors}). \textit{Moyennant} W(7 bis)\textit{, la catégorie} $\Psi = \begin{tikzcd}[ampersand replacement=\&, row sep=small, column sep=small, baseline=(B.base)]
 \& a \\
|[alias=B]| 0 \arrow[ur] \arrow[dr] \& \\
 \& b
\end{tikzcd}$ \textit{(ou $0 \longleftarrow 1 \longrightarrow 2$) est totalement $W$-asphérique} (\footnote{\textbf{Faux,} $\Psi$ n'est \textit{pas} totalement $W$-asphérique.}).

On applique le corollaire 1 à $\Psi$ et à ses deux sous-catégories ouvertes, $\Psi/a$, $\Psi/b$ (isomorphes à $\bDelta^1$), et à leur intersection $\{0\} \simeq \bDelta^0 = e$. Comme ces trois catégories ont un objet final (contrairement à $\Psi$ elle-même, qui n'a qu'un objet initial \dots), elles sont totalement $W$-asphériques. Le corollaire s'applique, l'existence des rétractions $r'_0$, $r''_0$ étant ici évidente.

C'est surtout la partie b) (rectifiée) de la page 9 qui assure l'existence de catégories totalement $W$-asphériques autres que celles ayant un objet final (\footnote{Je conjecture à présent que ce sont les seules! (\footnotemark)}\footnotetext{Faux, il y a les catégories filtrantes (si $W$ satisfait l'axiome des limites).}). Il est vrai qu'en sens inverse il faudrait quand même s'assurer qu'un objet $X$ $W$-asphérique de Cat n'est pas \textit{toujours} totalement $W$-asphérique! Le contre-exemple n'est pas tellement évident finalement -
\msp{page 101}
j'ai passé des heures avant d'en trouver un et d'y voir un peu clair! Il suffit en somme d'exhiber un exemple d'une flèche dans Cat
\[
f : X \longrightarrow Y
\]
telle que
% foot: Version 7 décembre 2002
\printedpage{71}
% head: A. Grothendieck : Dérivateurs, Ch. XII, Édition des Universités de Montpellier II et Paris VII
\begin{enumerate}[label=\alph*)]
\item $\underline{\pi_0}(f) \overset{\text{déf}}{=} \upi_0(X/Y) \xrightarrow{\ \sim\ } Y$, i.e. les $\pi_0(y\backslash X)$ \texttt{[plutôt les} $y\backslash X$\texttt{]} sont 0-connexes (pour $y \in \Ob Y$),
\end{enumerate}

et telle que de plus

\begin{enumerate}[label=\alph*), start=2]
\item $X$ est $W$-asphérique, et $Y$ non.
\end{enumerate}

% a vertical bar is printed in the left margin beside the next paragraph (lines "La difficulté ..." to "... le foncteur")
La difficulté est la suivante, quand on cherche des contre-exemples `idiots' : c'est que l'hypothèse a) implique déjà, quand $X$ est $W_\infty$-asphérique (ou seulement 1-connexe) que $Y$ est 1-connexe : le foncteur
\[
Y' \longmapsto X' = X \times_Y Y'
\]
de la catégorie des revêtements étales $Y'$ de $Y$ vers celle des revêtements étales $X'$ de $X$ est pleinement fidèle (vérification immédiate), d'où le résultat annoncé. Donc un contre-exemple doit faire appel à l'existence d'un $\pi_i(Y,y)$ non nul, avec $i \geq 2$. Cela implique
\msp{page 102}
qu'il n'y a pas de contre-exemple avec des catégories de la forme $Y = Z/F$ (déduites d'un foncteur $q : X \longrightarrow Z$ en prenant $F = q_!(e_X) = \pi_0(X/Z)$), où $Z$ est une catégorie telle que $\bDelta^1$ (?), $\Phi$, $\Psi$, ou $\begin{tikzcd}[ampersand replacement=\&, sep=small, baseline=(M.base)]
 \& {} \& \\
{} \arrow[ur] \arrow[dr] \& |[alias=M]| \phantom{x} \& {} \arrow[ul] \arrow[dl] \\
 \& {} \&
\end{tikzcd}$ (ayant le type d'homotopie du cercle) etc., car les composantes connexes de telles catégories $Y$ sont des $K(\pi,1)$ - les $\pi_i$ supérieures ($i \geq 2$) sont nuls.

Cette difficulté vue, le contre-exemple devient évident. Je pars d'un objet 1-connexe \textit{quelconque} $Y$ de $X$ \texttt{[plutôt de} Cat\texttt{]}, et je prends $y_0 \in \Ob Y$,
\[
X = \uCh(Y; y_0, -) \xrightarrow{\ f\ } Y.
\]
On sait que $X$ est $W_\infty$-asphérique, et que $f$ est $W_\infty$-parfait, donc les inclusions
\[
X_y \longrightarrow y\backslash X
\]
sont des $W_\infty$-équivalences. L'hypothèse a) que les $y\backslash X$ soient 0-connexes équivaut donc à la même hypothèse sur les fibres $X_y$, ou
\msp{page 103}
encore sur une seule fibre ($f$ étant parfait, donc une $W_\infty$-fibration, et $Y$ 0-connexe), soit
\[
X_{y_0} \simeq \Omega(Y, y_0).
\]
Sa 0-connexité équivaut à l'hypothèse de 1-connexité de $Y$ (on n'a utilisé encore que celle de 0-connexité). Donc pour un contre-exemple, il suffit de prendre $Y$ 1-connexe et non $\infty$-connexe. Le cas le plus simple est celui où $Y$ a le type d'homotopie de la sphère $S^2$, réalisé p. ex. par la catégorie ordonnée ci-contre \texttt{[ci-dessous]} à 6 sommets.
% foot: Version 7 décembre 2002

\printedpage{72}
% head: A. Grothendieck : Dérivateurs, Ch. XII, Édition des Universités de Montpellier II et Paris VII
\[
\begin{tikzcd}[row sep=small]
e_2^+ \bullet & \bullet\, e_2^- \\
e_1^+ \bullet \arrow[ur] & \bullet\, e_1^- \arrow[ul] \\
e_0^+ \bullet \arrow[ur] & \bullet\, e_0^- \arrow[ul]
\end{tikzcd}
\]
\texttt{[de plus,} $e_0^+ \leq e_1^+ \leq e_2^+$\texttt{,} $e_0^- \leq e_1^- \leq e_2^-$\texttt{]} (\footnote{C'est en somme cet exemple (à homotopie topologique près) que je réalise (sans d'abord m'en apercevoir dans l'Appendice p. 10-12, mais avec des catégories finies).}). Je résume le résultat obtenu.

\textbf{Proposition 10.} \textit{Soit $Y$ une petite catégorie} 1\textit{-connexe ($y_0 \in \Ob Y$). Alors} $\uCh(Y; y_0, -)$ \textit{est $W_\infty$-asphérique, mais n'est pas totalement $W_\infty$-asphérique si $Y$ elle-même n'est $\infty$-connexe, i.e. $W_\infty$-asphérique.} (\footnote{Mais un exemple plus percutant est $\Psi = \begin{tikzcd}[ampersand replacement=\&, row sep=small, column sep=small, baseline=(A.base)]
 \& b \\
|[alias=A]| a \arrow[ur] \arrow[dr] \& \\
 \& c
\end{tikzcd}$ , qui est $W_\infty$-asphérique, mais pas totalement!}).

On peut se demander si $Y$ $W_\infty$-asphérique implique que $\uCh(Y; y_0, -)$ est totalement $W_\infty$-asphérique. J'en doute fort!

\msp{page 104}
Revenons maintenant à la situation du début de la section, diagramme $(*)$, p.~87,
\[
\begin{tikzcd}
U \arrow[r, hook, "i"] \arrow[d, "p"'] & X \arrow[d, "q"'] & T = \{t\} \arrow[l, hook', dotted, "j"'] \arrow[d, dotted, "\mathrm{iso}"] \\
U' \arrow[r, hook, "i'"] & X' & T' = \{t'\}. \arrow[l, hook', dotted, "j'"']
\end{tikzcd}
\tag*{$(*)$}
\]
Nous pouvons donner maintenant un exemple où $i$ est dans $W$ (et même $W$-asphérique) sans que $i'$ le soit :
\[
i \in W, \qquad i' \notin W. \tag*{$(**)$}
\]
On peut même supposer que $t$ soit objet final de $X$ (ce qui signifie que l'objet $H$ de $U^\wedge$ est isomorphe à $e_U$, l'objet final), et $t'$ objet final de $X'$ (i.e. $p_!(e_U) = e_{U'}$, i.e. les $u'\backslash U$ sont 0-asphériques). Dans ce cas, comme $X$, $X'$ sont $W$-asphériques, les relations $(**)$ équivalent à
\[
U \text{ est } W\text{-asphérique, et } U' \text{ non,} \tag*{$(***)$}
\]
ce qui est justement la situation que nous venons de construire (p. 101-103).

On aimerait aussi trouver un exemple où dans $(*)$ on ait $p \in W$, $q \notin W$.
% foot: Version 7 décembre 2002


\endgroup
% END UNIT Der-ChXIIb-012/Der-ChXIIb-012_p026-030.tex
% BEGIN UNIT Der-ChXIIb-012/Der-ChXIIb-012_p031-035.tex
\begingroup\makeatletter%
\newcommand{\printedpage}[1]{\marginpar{\footnotesize p.~#1}}%
\newcommand{\msp}[1]{\marginpar{\footnotesize [#1]}}%
\newcommand{\hra}{\lhook\joinrel\longrightarrow}%
\newcommand{\Cat}{\mathrm{Cat}}%
\makeatother\let\printedpage\derpp\let\msp\dermsp%
\let\mspin\DERmspin%

\setcounter{footnote}{76}

\printedpage{73}
% head: A. Grothendieck : Dérivateurs, Ch. XII, Édition des Universités de Montpellier II et Paris VII
% foot: Version 7 décembre 2002

\msp{page 105}

Je prendrai encore $(X,U,\{t\})$ tel que $t$ soit objet final de $X$, mais si on veut $q\notin W$, on ne peut supposer $t'$ élément final de $X'$. Au lieu de cela, je vais supposer

\begin{enumerate}[label=\alph*)]
\item $U$, $U'$ (et aussi $X$, ayant objet final $t$) $W$-asphériques, donc $i\in W$ (et même $i$ $W$-asphérique), $p\in W$, mais
\item $X'$ non $W$-asphérique, donc $i',q\notin W$.
\end{enumerate}

Pour ceci, notons que
\[ U'/t'\simeq U'/p_!(e_U)\overset{\text{déf}}{=}\underline{\pi}_0(U/U'). \]
Le corollaire de la proposition 7 (cf. p. 85), non encore démontré, nous dit que $i'\in W$ (i.e. ici, $X'$ $W$-asphérique) si et seulement si $i'$ est $W$-asphérique, i.e. $U'/t'$ est $W$-asphérique. Nous allons l'admettre ici. Tout revient donc à trouver une flèche
\[ p:U\longrightarrow U' \]
dans Cat avec les propriétés suivantes.

\begin{enumerate}[label=\alph*)]
\item $U$ et $U'$ sont $W$-asphériques,
\item $\underline{\pi}_0(U/U')$ n'est pas $W$-asphérique.
\end{enumerate}

Nous allons reprendre l'exemple précédent
\[ f:X\longrightarrow Y \]
(p. 101-103), en prenant
\[ U=X,\qquad \underline{\pi}_0(U/U')=Y, \]

\msp{page 106}

donc tout revient, à partir de l'exemple donné $X\xrightarrow{f}Y$, à construire une flèche
\[ \pi:Y\longrightarrow U' \tag*{$(*)$} \]
dans Cat qui satisfasse les conditions

\begin{enumerate}[label=\alph*)]
\item $\pi$ est un isomorphisme local,
\item $U'$ est $W$-asphérique.
\end{enumerate}

On prendra alors
\[ p=\pi f. \]
Mais par la façon dont $X$ a été construit à partir d'un $Y$ 1-connexe et non $\infty$-connexe, par ailleurs quelconque, on est ramené finalement à trouver $(*)$ satisfaisant a), b), et de plus

\begin{enumerate}[label=\alph*), start=3]
\item \printedpage{74}
% head: A. Grothendieck : Dérivateurs, Ch. XII, Édition des Universités de Montpellier II et Paris VII
% foot: Version 7 décembre 2002
$Y$ est 1-connexe, mais non $\infty$-connexe.
\end{enumerate}

(NB On suppose maintenant $W=W_\infty$.) Il revient au même de trouver

% printed with a vertical rule at the left of this block
\begin{quote}
$U'$ objet \textit{asphérique} de Cat, $F\in\operatorname{Ob}U'^\wedge$ tels que $U'/F$ soit 1-connexe et non $\infty$-connexe.
\end{quote}

On songe par exemple à prendre simplement une sous-catégorie ouverte $U'_0$ de $U'$, qui soit 1-connexe et non $\infty$-connexe. Partons d'une telle catégorie $U'_0$ - par exemple une réalisation homotopique dans Cat, de la sphère $S^2$. On l'immerge comme ouvert dans un $U'$ asphérique, en lui ajoutant un objet final. On a gagné.

\msp{page 107}

Faisons un énoncé en forme.

\textbf{Proposition 11.} \textit{On peut trouver une situation} $(*)$ (p. 124 \texttt{[plutôt 104]}) \textit{avec}

\begin{enumerate}[label=\textit{\alph*)}]
\item \textit{$t$ objet final de $X$.}
\item \textit{$U$, $U'$ (et $X$) $W_\infty$-asphériques, donc $i,p\in W_\infty$ (et même $i$ $W_\infty$-asphérique).}
\item \textit{$X'$ non $W_\infty$-asphérique, donc $i',q\notin W_\infty$.}
\end{enumerate}

Il suffit pour cela de partir d'un objet 1-connexe $Y$ de Cat, de lui adjoindre un objet final $u'_0$, ce qui donne $U'$, de prendre $U=\underline{\mathrm{Ch}}(Y;y_0,-)$, $p:U\longrightarrow U'$ le composé $U\xrightarrow{\text{but}}Y\underset{\text{can.}}{\overset{\text{incl.}}{\hra}}U'$, et $X$ déduit de $U$ en lui ajoutant un objet final $t$. On a alors $U'/t'\simeq Y$, d'où
\[ X'\ W_\infty\text{-asphérique (i.e. }i'\in W_\infty)\Longleftrightarrow Y\ W_\infty\text{-asphérique.} \]
On a le contre-exemple cherché en prenant $Y$ (1-connexe mais) non $\infty$-connexe.

\subsection*{2.4\quad Application aux immersions ouvertes dwyériennes.}

Donnons quand même l'enoncé \textit{positif} auquel nous en avions. % [sic] l'enoncé -> l'énoncé, missing accent

\textbf{Théorème 3.} \textit{Soient $(X,U,T)$ un triple $\in\mathcal{T}$ ($U$ ouvert de $X$ dans} Cat\textit{, $T=X\backslash U$), et pour toute flèche $p:U\longrightarrow U'$ flèche de} Cat\textit{, le carré cocartésien dans} Cat

\[
\begin{tikzcd}
U \arrow[r, hook, "i"] \arrow[d, "p"'] & X \arrow[d, "q"'] & T \arrow[l, dotted, hook'] \arrow[d, dotted, "\text{iso}"] \\
U' \arrow[r, hook, "i'"] & X' & T'. \arrow[l, dotted, hook']
\end{tikzcd}
\]

\printedpage{75}
% head: A. Grothendieck : Dérivateurs, Ch. XII, Édition des Universités de Montpellier II et Paris VII
% foot: Version 7 décembre 2002

\msp{page 108}

\textit{Supposons que pour tout $t\in\dot{T}=\dot{U}$, le préfaisceau d'ensembles}
\[ i^*(t)=(u\longmapsto\operatorname{Hom}(u,t))\in U^\wedge \]
\textit{soit totalement $W$-asphérique (déf. 1, p. 88)} (\footnote{\textbf{NB} On dit sous ces conditions que $i_0:U\longrightarrow\bar{U}$ est \textit{strictement $W$-asphérique} (cf. déf. 2 page suivante).}) \textit{- p. ex. que la catégorie $U/t$ soit totalement $W$-asphérique (prop. 8, page 93). (NB le cas où les $U/t$ ont un objet final, cf. prop. 9 a), i.e. celui où l'inclusion $U\longrightarrow\bar{U}$ admet un adjoint à droite, est celui envisagé par Dwyer, celui des immersions ouvertes de Dwyer} (\footnote{Je ne connais pas d'autre cas où th. 3 s'applique - faute de connaître de catégorie totalement $W$-asphérique qui n'ait un objet final, ni d'objet $F$ d'une catégorie $U^\wedge$ qui soit totalement $W$-asphérique sans être représentable, i.e. $U/F$ a un objet final. À présent, je sais que les catégories \textit{filtrantes} sont totalement $W$-asphériques si $W$ satisfait l'axiome des limites inductives.})\textit{.) Alors pour toute $p:U\longrightarrow U'$, l'immersion ouverte}
\[ i'_0:U'\longrightarrow\overline{U'} \]
\textit{est $W$-asphérique (et même, tout comme $i_0$, strictement $W$-asphérique), et a fortiori elle est} \texttt{[dans]} \textit{$W$. Par suite (prop. 6} \texttt{[p. 81]}\textit{) on a ceci.}

\[
\begin{aligned}
&\text{\textcircled{a}}\qquad && i\in W && \Longleftrightarrow\quad && i'\in W\quad && \textit{moyennant } \mathrm{W}(7\text{ bis}).\\
&\text{\textcircled{b}}\qquad && p\in W && \Longleftrightarrow\quad && q\in W\quad && \textit{moyennant } \mathrm{W}(7).\ (\footnotemark)
\end{aligned}
\]
\footnotetext{Démonstration page 112.}

\textbf{Corollaire 1.} \textit{Supposons} W(7)\textit{. Si $i:U\hra X$ est une immersion ouverte telle que $i_0:U\hra\bar{U}$ soit totalement} \texttt{[il suffit strictement]} \textit{$W$-asphérique, alors $i:U\longrightarrow X$ est une $W$-cofibration (`ultraforte'), par quoi j'entends ceci : pour toute flèche $U'\xrightarrow{f}U''$ dans $U\backslash$}Cat\textit{,}

\msp{page 109}

\textit{telle que $f\in W$, la flèche $g:X'\longrightarrow X''$ qui s'en déduit par cochangement de base à partir de l'objet $X$ de $U\backslash$}Cat\textit{, est aussi dans $W$.}

\textsc{Démonstration du corollaire} : on applique le théorème 3 à l'immersion ouverte $U'\overset{i'}{\hra}X'$ déduite de $U\overset{i}{\hra}X$ par le cochangement de base $U\xrightarrow{f}U'$, en notant de plus que l'hypothèse du théorème 3 sur $i$ est préservée en passant de $i$ à $i'$.

Il faut quand même introduire une terminologie adéquate :

\textbf{Définition 2 :}

\begin{enumerate}[label=\alph*)]
\item Une flèche $f:X\longrightarrow Y$ est dite \textit{strictement $W$-asphérique} si pour toute $y\in\operatorname{Ob}Y$, $f^*(y)=(x\longmapsto\operatorname{Hom}_Y(f(x),y))$ est un objet totalement $W$-asphérique de $X^\wedge$ (cf. déf. 1, page 88). Elle est dite \textit{totalement $W$-asphérique} si pour tout $y\in\operatorname{Ob}Y$, $X/y$ ($\simeq X/f^*(y)$) est un objet totalement $W$-asphérique de Cat.

\item \printedpage{76}
% head: A. Grothendieck : Dérivateurs, Ch. XII, Édition des Universités de Montpellier II et Paris VII
% foot: Version 7 décembre 2002
Une immersion ouverte $i:U\longrightarrow X$ est dite \textit{$W$-dwyérienne} (resp. \textit{totalement} \texttt{[$W$-]} \textit{dwyérienne}) si l'immersion induite $i_0:U\longrightarrow\bar{U}$ est strictement (resp. totalement) $W$-asphérique. Une immersion fermée $j:T\longrightarrow X$ est dite \textit{$W$-dwyérienne} (resp. \textit{totalement $W$-dwyérienne}) si et seulement si $j^\circ:T^\circ\longrightarrow X^\circ$ est une immersion ouverte $W$-dwyérienne (resp totalement $W$-dwyérienne).

\item Une flèche $i:U\longrightarrow X$ de Cat (pas nécessairement une immersion)

\msp{page 110}

est appelée une \textit{$W$-cofibration,} si pour toute flèche $f:U'\longrightarrow U''$ dans $U\backslash$Cat, la flèche $g:X'\longrightarrow X''$ qui s'en déduit par le cochangement de base par $X$ (objet de $U\backslash$Cat) est dans $W$ quand $f$ l'est.
\end{enumerate}

\textbf{NB} Cette notion de $W$-cofibration est stable par composition et par cochangement de base, et les isomorphismes de Cat sont des $W$-cofibrations. Si $f$ est une $W$-cofibration, $f^\circ$ l'est et inversement.

(Le corollaire s'énonce en disant qu'une immersion ouverte ou fermée qui est $W$-dwyérienne est une $W$-cofibration.)

\begin{enumerate}[label=\alph*), start=4]
\item Une flèche $i:U\longrightarrow X$ dans Cat est appelée une \textit{$W$-cofibration triviale}, ou une \textit{$W$-équivalence couniverselle,} si pour toute $U'$ dans $U\backslash$Cat, i.e. toute flèche $p:U\longrightarrow U'$ dans Cat, la flèche déduite de $i$,
\[ i':U'\longrightarrow X'=U'\amalg_U X'\quad\texttt{[plutôt }U'\amalg_U X\texttt{]}\ \ , \]
est dans $W$. (NB Il est immédiat que cela implique bien que $i$ est une $W$-cofibration.)
\end{enumerate}

(Cette notion de $W$-cofibration triviale est elle aussi stable par composition et par cochangement de base, et les isomorphismes dans Cat en sont.)

\msp{page 111}

\textbf{Corollaire 2.} (Moyennant W(7).) \textit{Soit $i:U\longrightarrow X$ une immersion ouverte $W$-dwyérienne. Pour que $i$ soit une $W$-cofibration triviale, il faut et il suffit que $i\in W$, ou ce qui revient au même} (cf. page 82), \textit{que $\dot{T}\hra T$ soit dans $W$.}

C'est évidemment nécessaire, et la suffisance résulte du théorème 3 a).

\textbf{Questions ouvertes.}

\begin{enumerate}[label=\alph*)]
\item J'ignore si une immersion ouverte $i$ qui est une $W$-cofibration, voire une $W$-cofibration triviale, est forcément $W$-dwyérienne (\footnote{et peut-être même une immersion de Dwyer?}). Mais par le corollaire de la proposition 7 (p. 85), je sais que c'est le cas si on suppose $T$ isomorphe à la catégorie finale $e$. Il s'ensuit toute au moins ceci : Pour que non seulement $i:U\longrightarrow X$, mais aussi toutes les immersions ouvertes $i_1:U\longrightarrow X_1$ qui s'en déduisent par le procédé de prop. 7 b) au moyen de flèches quelconques $T_1\longrightarrow T$ de but $T$ dans Cat, soient des $W$-cofibrations, il faut et il suffit que $i$ soit dwyérienne. % [sic] toute au moins -> tout au moins, agreement error

\item \printedpage{77}
% head: A. Grothendieck : Dérivateurs, Ch. XII, Édition des Universités de Montpellier II et Paris VII
% foot: Version 7 décembre 2002
Y-a-t-il une bonne notion d' immersion dwyérienne, applicable à toutes les immersions

\msp{page 112}

et incluant le cas des immersions ouvertes et des immersions fermées?
\end{enumerate}

\textsc{Démonstration} du théorème 3 (p. 107). Il suffit de prouver que si $i$ est $W$-dwyérienne (i.e. $i_0$ strictement $W$-asphérique). alors $i'$ l'est également - car il en résultera que $i'$ est $W$-asphérique et a fortiori $\in W$. Il faut donc vérifier que pour tout $t'\in\dot{U}'=\dot{T}$, l'élément $i^*(t')=(u'\longmapsto\operatorname{Hom}_{X'}(u',t'))$ de $U'^\wedge$ est totalement $W$-asphérique. Or $t'$ provient d'un unique élément $t$ de $\dot{T}$, et on sait qu'on a % [sic] asphérique). alors -> asphérique), alors, period before lower-case word

\[ i'^*(t)=p_!(i^*(t))\quad\texttt{[plutôt }i'^*(t')=p_!(i^*(t))\texttt{]},\qquad\text{où }\ p_!:U^\wedge\longrightarrow U'^\wedge. \]

Par hypothèse, $F=i^*(t)$ est un élément totalement $W$-asphérique de $U^\wedge$, ce qui signifie que pour toute flèche $p:U\longrightarrow U'$ dans Cat de souce $U$, $p_!(F)$ est un élément $W$-asphérique de $U'$. En fait, on a le lemme immédiat (par transitivité $(p'p)_!\simeq p'_!p_!$). % [sic] souce -> source, missing letter

\textbf{Lemme.} \textit{Soit $U$ dans} Cat\textit{, et $F$ dans $U^\wedge$. Si $F$ est totalement $W$-asphérique (resp. si $U/F$ est totalement $W$-asphérique), alors pour toute flèche $p:U\longrightarrow U'$ dans} Cat \textit{de source $U$, $F'=p_!(F)$ est un objet totalement $W$-asphérique de $U'^\wedge$ (resp.}

\msp{page 113}

\textit{$U'/F'$ est un objet totalement $W$-asphérique de} Cat\textit{.)}

\textbf{Corollaire 3} (du théorème 3). \textit{Si $i:U\longrightarrow X$ est une immersion ouverte ou fermée $W$-dwyérienne (resp. totalement $W$-dwyérienne), alors pour toute flèche $p:U\longrightarrow U'$ dans} Cat\textit{, $i':U'\longrightarrow X'=U'\amalg_U X'$} \texttt{[plutôt $U'\amalg_U X$]} \textit{l'est aussi.}

(NB Le cas non respé est essentiellement le théorème 3. Il ne dépend d'aucun axiome supplémentaire sur $W$, mais n'a d'intérêt qu'en vue des équivalences \textcircled{a} et \textcircled{b} (p. 108) qui en résultent, et qui sont valables moyennant W(7 bis) resp. W(7) seulement.)

\textbf{Commentaire de motivation :} les immersions ouvertes ou fermées $W$-dwyériennes, ou totalement $W$-dwyériennes, sont stables par cochangement de base (d'après le corollaire précédent), mais n'ont pas du tout l'air de l'être par composition. À cet égard, les `immersions de Dwyer' (où on suppose l'existence de l'adjoint à droite de $i_0:U\longrightarrow\bar{U}$, dans le cas d'une immersion ouverte) se comportent mieux, étant stables par composition. (\textbf{NB} Il ne semble pas

\msp{page 114}

non plus clair que les flèches strictement $W$-asphériques, ou les totalement $W$-asphériques, soient stables par composition - sorry!) Ma motivation pour m'échiner sur la notion $W$-dwyérienne est double.

\begin{enumerate}[label=\alph*)]
\item J'aimerais avoir des conditions à la fois compréhensibles, et aussi générales que possible, sur une immersion (ouverte p. ex), pour que ce soit une $W$-cofibration, propriété visiblement importante (toute comme la notion duale de $W$-fibration). % [sic] toute comme -> tout comme, agreement error
\end{enumerate}


\endgroup
% END UNIT Der-ChXIIb-012/Der-ChXIIb-012_p031-035.tex
% BEGIN UNIT Der-ChXIIb-012/Der-ChXIIb-012_p036-040.tex
\begingroup\makeatletter%
\NoAutoSpacing%
\newcommand{\printedpage}[1]{\marginpar{\footnotesize p.~#1}}%
\newcommand{\msp}[1]{\marginpar{\footnotesize [#1]}}%
\makeatother\let\printedpage\derpp\let\msp\dermsp%
\let\mspin\DERmspin%


\printedpage{78}
% head: A. Grothendieck : Dérivateurs, Ch. XII, Édition des Universités de Montpellier II et Paris VII
\begin{enumerate}
\item[b)] J'ai envie, pour des applications ultérieures, de disposer d'une classe d'immersions ouvertes ayant les propriétés suivantes :
\begin{enumerate}
\item[1°)] Ce sont des $W$-cofibrations, et
\item[2°)] Si $U \overset{i}{\hookrightarrow} X$ est dans la classe, et si $\varepsilon \xrightarrow{j} I$ est l'inclusion ‘source-but’ de la catégorie discrète $\{0,1\}$ dans l'intervalle
\[
\Phi_n = 0 \longrightarrow 1 \longleftarrow 2 \longrightarrow \dots 2n-1 \longleftarrow 2n
\]
$(j(0)=0$, $j(1)=2n)$, $n\geq 2$, l'inclusion $\alpha$ de $(U\times I)\cup(X\times\varepsilon)$ dans $X\times I$ est dans la classe.
\end{enumerate}
Or cette dernière propriété de stabilité n'est \textit{pas} vérifiée par les immersions de Dwyer. (Prendre $i : \{0\} \longrightarrow \boldsymbol{\Delta}^1 = \{0,1\}$, l'inclusion $\alpha$ n'est pas immersion de Dwyer.)
\end{enumerate}

\msp{page 115}

\subsection*{2.5\quad Complément : critère d'asphéricité totale d'une catégorie $W$-propre sur une autre.}

Soit
\[
f : X \longrightarrow I
\]
une catégorie propre sur une catégorie d'indices $I$. Si $I$ a un objet final $i_0$, alors $X/i_0$ s'identifie à $X$, et par suite l'inclusion
\[
X_{i_0} \longrightarrow X
\]
de la fibre de $X$ en $i_0$ est $W$-asphérique \texttt{[plutôt }$W$\texttt{-coasphérique]}, et a fortiori est dans $W$. Donc $X$ est asphérique si et seulement si $X_{i_0}$ l'est. Nous voulons étendre ce résultat en un critère de \textit{totale} $W$-asphéricité, de façon précise prouver la

\textbf{Proposition 11 :} \textit{Soit $f : X \longrightarrow I$ $W$-propre, et soit $i_0$ un objet final de $I$. Si $X_{i_0}$ est totalement $W$-asphérique} (déf. 1, page 88), \textit{il en est de même de $X$.} (Sous réserve que $W$ satisfasse l'axiome de connexité W(5).)

La démonstration sera donnée plus bas, après quelques préliminaires ayant un intérêt indépendant. Tout d'abord, rappelons que si $X$ est $W$-propre sur $I$, alors il donne naissance à un foncteur
\[
\begin{aligned}
I \quad &\longrightarrow \quad \mathrm{Hot}_W\\
i \quad &\longmapsto \quad \mathrm{hot}_W(X_i)(\simeq \mathrm{hot}_W(X/i));
\end{aligned}
\]
en le composant avec le foncteur
\[
\pi_0 : \mathrm{Hot}_W \longrightarrow \mathrm{Ens}
\]
(sous réserve que l'axiome de connexité W(5) (p. 10) soit satisfait)
% foot: Version 7 décembre 2002
\printedpage{79}
% head: A. Grothendieck : Dérivateurs, Ch. XII, Édition des Universités de Montpellier II et Paris VII
\msp{page 116}
on trouve un foncteur
\[
i \longmapsto \underbrace{\pi_0(X_i)}_{\simeq\, \pi_0(X/i)} : I \longrightarrow \mathrm{Ens}.
\]
Ceci posé, on a le résultat sorital

\textbf{Proposition 12 :} \textit{Soit $f : X \longrightarrow I$ $W$-propre, et supposons que $W$ satisfasse l'axiome de connexité} W(5) (sinon, la $W$-propreté ne nous sera d'aucune utilité \dots) (\footnote[81]{Il vaut mieux dire : Soit $f$ $W_0$-propre, où $W_0$ est le localiseur de la 0-connexité \dots}). \textit{On a alors une bijection canonique}
\[
\pi_0(X) \xleftarrow{\ \sim\ } \varinjlim_{I} \pi_0(X_i), \tag*{$(*)$}
\]
\textit{où la flèche est induite par les applications canoniques $\pi_0(X_i) \longrightarrow \pi_0(X)$.}

\textsc{Démonstration.} Passant à $f^\circ : X^\circ \longrightarrow I^\circ$, on est ramené à prouver le même résultat pour $X$ \textit{$W$-lisse} sur $I$, avec la seule différence que le foncteur
\[
i \longmapsto \pi_0(X_i)
\]
sur $I$ sera contravariant, i.e. un préfaisceau, ce qui offre l'avantage de l'intuition attachée au formalisme des préfaisceaux. Soit
\[
g : I \longrightarrow e, \qquad \text{d'où } gf = p : X \longrightarrow e,
\]
il est immédiat, par le calcul standard
\[
p_!(F) = \pi_0(X/F) \qquad (F \in \mathrm{Ob}\, X^\wedge)
\]
(où on a identifié $e^\wedge$ à (Ens)), qu'on a
\[
\pi_0(X) = p_!(e_X) = g_!(f_!(e_X)).
\]
Or le préfaisceau
\[
i \longmapsto \pi_0(X_i) \qquad (\simeq \pi_0(i\backslash X), \text{ car } X\ W\text{-lisse sur } I)
\]
\msp{page 117}
n'est autre que $f_!(e_X)$. D'autre part, pour un préfaisceau quelconque $G$ sur $I$, on a
\[
g_!(G) \simeq \pi_0(I/G) \simeq \varinjlim_{I} G(i),
\]
et appliquant ceci à $f_!(e_X)$, on trouve la formule $(*)$ de la proposition 12.

Revenons à la situation avec $f$ \textit{$W$-propre,} et considérons un foncteur
\[
p : X \longrightarrow Y,
\]
nous sommes intéressés à ‘calculer’
\[
p_!(e_X) = (y \longmapsto \pi_0(y\backslash X)) \in \mathrm{Ob}\, Y^\wedge.
\]
Considérons pour ceci le diagramme suivant
% foot: Version 7 décembre 2002
\printedpage{80}
% head: A. Grothendieck : Dérivateurs, Ch. XII, Édition des Universités de Montpellier II et Paris VII
\[
\begin{tikzcd}
y\backslash X_i \arrow[r] \arrow[d] & y\backslash X \arrow[r] \arrow[d] & y\backslash Y \arrow[d] \\
X_i \arrow[r] \arrow[d] & X \arrow[r, "p"] \arrow[d, "f"] & Y \\
\underset{\simeq\, e}{\{i\}} \arrow[r] & I\rlap{\quad ,} &
\end{tikzcd}
\]
où les trois carrés sont cartésiens. Comme $y\backslash Y \longrightarrow Y$ est $W$-propre, il en est de même de $y\backslash X \longrightarrow X$ par changement de base, donc du composé $y\backslash X \longrightarrow X \longrightarrow I$ par transitivité. Le diagramme montre d'autre part que la fibre de $y\backslash X$ en $i \in \mathrm{Ob}\, I$ s'identifie à $y\backslash X_i$.

\msp{page 118}
La proposition 12 nous donne donc la formule
\[
\pi_0(y\backslash X) = \varinjlim_{I} \pi_0(y\backslash X_i),
\]
laquelle formule ne peut s'empêcher d'être fonctorielle en $y$, tout comme celle de prop. 12 l'est dans $\underline{\mathrm{Prop}}_W I$. Or pour $y$ variable dans $Y$, le préfaisceau sur $Y$
\[
y \longmapsto \pi_0(y\backslash X_i)
\]
n'est autre que $p_{i\,!}(e_{X_i})$, où
\[
p_i : X_i \longrightarrow Y
\]
est le composé $X_i \longrightarrow X \xrightarrow{p} Y$ (induit par $p$). On a ainsi prouvé

\textbf{Corollaire 1} (de prop. 12). \textit{Soit $f : X \longrightarrow I$ $W$-propre} (\footnote[82]{dire $W_0$-propre \dots}), \textit{et $p : X \longrightarrow Y$ un foncteur quelconque. Soit, pour tout $i \in \mathrm{Ob}\, I$, $p_i : X_i \longrightarrow Y$ le foncteur induit sur la fibre $X_i$ de $X$ en $i$. On a alors une loi} fonctorielle \textit{naturelle}
\[
i \longmapsto p_{i\,!}(e_{X_i}) : I \longrightarrow Y^\wedge,
\]
\textit{et on a la formule}
\[
p_!(e_X) \xleftarrow{\ \sim\ } \varinjlim_{I} p_{i\,!}(e_{X_i}), \tag*{$(*)$}
\]
\textit{où la flèche est induite par les morphismes canoniques}
\[
p_{i\,!}(e_{X_i}) \longrightarrow p_!(e_X)
\]
% foot: Version 7 décembre 2002
\printedpage{81}
% head: A. Grothendieck : Dérivateurs, Ch. XII, Édition des Universités de Montpellier II et Paris VII
\msp{page 119}
\textit{induites par les inclusions $X_i \longrightarrow X$.}

La loi fonctorielle de $p_{i\,!}(e_{X_i})$ en $i$ a été explicitée en cours de démonstration.

Notons que lorsque $X$ sur $I$ est cofibrée et scindée, i.e. provient d'un foncteur
\[
i \longmapsto X_i, \quad I \longrightarrow \mathrm{Cat},
\]
alors la loi fonctorielle de $p_{i\,!}(e_{X_i})$ en $i$ s'explicite de façon particulièrement agréable : à une flèche
\[
\alpha : i \longrightarrow j, \qquad \text{d'où } \alpha_* : X_i \longrightarrow X_j,
\]
est associé
\[
p_!(\alpha) : p_{i\,!}(e_{X_i}) \longrightarrow p_{j\,!}(e_{X_j}),
\]
en notant qu'on a un diagramme
\[
\begin{tikzcd}
X_i \arrow[rr, "\alpha_*"] \arrow[dr, "p_i"', ""{name=A}] & & X_j \arrow[dl, "p_j", ""'{name=B}] \\
& Y &
\arrow[from=A, to=B, bend left, "\lambda_\alpha"] % [resolved: see DER_XIIb-a_corrections.csv]
\end{tikzcd}
\]
avec une flèche de commutativité
\[
\lambda_\alpha : p_i \longrightarrow p_j\alpha_*
\]
(pas \texttt{[un]} isomorphisme en général), d'où
\[
p_{i\,!}(e_{X_i}) \xrightarrow{(\lambda_\alpha)_!} p_{j\,!}\underbrace{\alpha_{*\,!}(e_{X_i})}_{\longrightarrow\, e_{X_j}} \longrightarrow p_{j\,!}(e_{X_j}),
\]
plus de doute c'est lui! (\footnote[83]{Cette flèche \texttt{[la flèche }$(\lambda_\alpha)_!$\texttt{]} devrait être explicitée avec soin!}).

\msp{page 120}
\textbf{Corollaire 2 :} \textit{Sous les conditions du corollaire 1, supposons que $I$ ait un objet final $i_0$. Alors}
\[
(p_{i_0})_!(e_{X_{i_0}}) \longrightarrow p_!(e_X)
\]
\textit{est un isomorphisme dans $Y^\wedge$.}

D'où aussitôt la proposition 11 : si $X_{i_0}$ est totalement $W$-asphérique, alors $X$ l'est, car pour tout $p : X \longrightarrow Y$, $p_!(e_X) \simeq p_{i_0\,!}(X_{i_0})$, et ce dernier élément de $Y^\wedge$ est $W$-asphérique % [sic] p_{i_0 !}(X_{i_0}) -> p_{i_0 !}(e_{X_{i_0}}), elsewhere the argument of p_{i!} is e_{X_i}
par hypothèse sur $X_{i_0}$.

Notons le cas particulier :

\textbf{Corollaire 3.} \textit{Soit $Z$ un objet totalement $W$-asphérique dans} Cat. \textit{Alors son produit $Z\times I$ par toute $I$ dans} Cat \textit{ayant un objet final est aussi totalement $W$-asphérique.}
% foot: Version 7 décembre 2002
\printedpage{82}
% head: A. Grothendieck : Dérivateurs, Ch. XII, Édition des Universités de Montpellier II et Paris VII

\msp{page 121}

\section*{3\quad Sommes amalgamées et intégration\quad \begin{tikzcd}[ampersand replacement=\&, row sep=small, column sep=small, baseline=(\tikzcdmatrixname-2-1.base)]
\& X_1 \\
\int X_0 \arrow[ur] \arrow[dr] \& \\
\& X_2
\end{tikzcd}}

\subsection*{3.1\quad Réduction canonique au cas d'un diagramme d'immersions fermées}

Considérons un diagramme
\[
\begin{tikzcd}[row sep=small, column sep=small]
& X_1 \\
X_0 \arrow[ur, "\alpha_1"] \arrow[dr, "\alpha_2"'] & \\
& X_2
\end{tikzcd}
\tag*{$(*)$}
\]
dans Cat. Nous pouvons le regarder comme un contrafoncteur
\[
\Phi^\circ \longrightarrow \mathrm{Cat},
\]
où
\[
\Phi = \left( \vcenter{\hbox{\begin{tikzcd}[row sep=small, column sep=small]
& 1 \arrow[dl] \\
0 & \\
& 2 \arrow[ul]
\end{tikzcd}}} \right),
\]
donc comme définissant une catégorie fibrée (scindée) $\mathscr{X}$ sur $\Phi$. Il est parfois commode de représenter cette catégorie $\mathscr{X}$ par le ‘signe d'integration’ % [sic] integration -> intégration, accent missing in the print
\[
\begin{tikzcd}[row sep=small, column sep=small]
& X_1 \\
\int X_0 \arrow[ur] \arrow[dr] & \\
& X_2,
\end{tikzcd}
\]
et la regarder comme une sorte de ‘somme amalgamée homotopique’ de $X_1$ et $X_2$ sous $X_0$, ou comme une ‘limite inductive homotopique’ du diagramme de départ $(*)$.

\msp{page 122}
Il faut entendre par là que (pour $W \subseteq \mathrm{Fl}(\mathrm{Cat})$ donné) l'image $\mathrm{hot}_W(\mathscr{X})$ de $\mathscr{X}$ dans $\mathrm{Hot}_W$ a un ‘sens homotopique’ remarquable en termes de $(*)$, de façon précise, le diagramme $(*)$ définit un objet $\xi$ de $\mathrm{HOT}_W(\Phi)$ (qui est une catégorie de fractions de la catégorie de tous les diagrammes du type $(*)$ \dots), et si
\[
\pi : \Phi \longrightarrow e
\]
désigne la projection absolue de $\Phi$, on a
\[
\pi_!(\xi) = \mathrm{hot}_W(\mathscr{X}),
\]
% foot: Version 7 décembre 2002


\endgroup
% END UNIT Der-ChXIIb-012/Der-ChXIIb-012_p036-040.tex
% BEGIN UNIT Der-ChXIIb-012/Der-ChXIIb-012_p041-045.tex
\begingroup\makeatletter%
\newcommand{\printedpage}[1]{\marginpar{\footnotesize p.~#1}}%
\newcommand{\msp}[1]{\marginpar{\footnotesize [#1]}}%
\newcommand{\Hot}{\mathrm{Hot}}%
\newcommand{\HOT}{\mathrm{HOT}}%
\newcommand{\Cat}{\mathrm{Cat}}%
\newcommand{\cX}{\mathcal{X}}%
\makeatother\let\printedpage\derpp\let\msp\dermsp%
\let\mspin\DERmspin%


\printedpage{83}
% head: A. Grothendieck : Dérivateurs, Ch. XII, Édition des Universités de Montpellier II et Paris VII
d'après le calcul connu du foncteur $\pi_!$, adjoint à gauche du foncteur $\pi^*$
\[
\HOT_W(e) \overset{\text{déf}}{=} \Hot_W \longrightarrow \HOT_W(\Phi)
\]
(\footnote[84]{Cela suppose qu'on a W(2), non seulement W(2 bis).}). Or pour tout objet $I$ dans Cat, de projection absolue $p_I : I \longrightarrow e$, le foncteur
\[
p_{I\,!} : \HOT_W(I) \longrightarrow \HOT_W(e) = \Hot_W
\]
doit être regardé comme l'équivalent `homotopique' du foncteur limite inductive $\varinjlim\limits_{I}$ des diagrammes de type $I^{\circ}$ dans la catégorie $\Hot_W$ - lesquelles limites inductives \textit{n'existent pas} dans $\Hot_W$ (sauf
\msp{page 123}
pour des catégories très particulières - les seules que je connaisse sont les catégories discrètes (qu'il faut peut-être même supposer finies en plus)).

Notons qu'on a un diagramme de foncteurs canoniques
\[
\begin{tikzcd}[column sep=tiny, row sep=small, ampersand replacement=\&]
 \& X_0 \arrow[dl, "\alpha_1"'] \arrow[dr, "\alpha_2"] \& \\
X_1 \arrow[dr, "i_1"'] \& {\shortstack{\scriptsize non\\ \scriptsize commuta-\\ \scriptsize tif!}} \& X_2 \arrow[dl, "i_2"] \\
 \& \cX, \&
\end{tikzcd}
\tag*{$(*)$}
\]
lequel diagramme n'a rien de commutatif dans Cat. \textit{Il est cependant commutatif dans} $\Hot_W$ (quel que soit le localiseur fondamental $W$), et même dans Cat\,hot, car si on désigne par
\[
i_0 : X_0 \longrightarrow \cX
\]
l'inclusion canonique, on a des homomorphismes canoniques
\[
i_1\alpha_1 \longrightarrow i_0 \longleftarrow i_2\alpha_2,
\]
ce qui montre que $i_1\alpha_1$ et $i_2\alpha_2$ sont homotopes, donc définissent la même flèche dans $\Hot_W$.

\msp{page 124}
Considérons dans $\Phi$ les deux fermés isomorphes à $\boldsymbol{\Delta}^1$
\[
\begin{aligned}
F_1 &= 1\backslash\Phi &&= \{0,1\}\\
F_2 &= 2\backslash\Phi &&= \{0,2\},
\end{aligned}
\]
on a
\[
F_1\cap F_2 = 0\backslash\Phi = \{0\}, \qquad F_1\cup F_2 = \Phi.
\]

\printedpage{84}
% head: A. Grothendieck : Dérivateurs, Ch. XII, Édition des Universités de Montpellier II et Paris VII
Soient
\[
\begin{aligned}
1\backslash\cX &= \cX\times_{\Phi} 1\backslash\Phi &&= \cX\times_{\Phi} F_1\\
2\backslash\cX &= \cX\times_{\Phi} 2\backslash\Phi &&= \cX\times_{\Phi} F_2
\end{aligned}
\]
les images inverses de $F_1$, $F_2$ dans $\cX$, on a donc
\[
1\backslash\cX\cap 2\backslash\cX = X_0, \qquad 1\backslash\cX\cup 2\backslash\cX = \cX,
\]
donc le diagramme
\[
\begin{tikzcd}[column sep=small, row sep=small, ampersand replacement=\&]
 \& X_0 \arrow[dl, hook', "\alpha'_1"'] \arrow[dr, hook, "\alpha'_2"] \& \\
1\backslash\cX \arrow[dr, hook, "i'_1"'] \& \& 2\backslash\cX \arrow[dl, hook', "i'_2"] \\
 \& \cX, \&
\end{tikzcd}
\tag*{$(**)$}
\]
déduit du diagramme similaire
\[
\begin{tikzcd}[column sep=small, row sep=small, ampersand replacement=\&]
 \& \{0\} \arrow[dl, hook'] \arrow[dr, hook] \& \\
F_1 \arrow[dr, hook] \& \& F_2 \arrow[dl, hook'] \\
 \& \Phi, \&
\end{tikzcd}
\]
est (comme ce dernier) un carré d'immersions fermées à la fois cartésien et cocartésien. D'autre part, on a un homomorphisme
\msp{page 125}
canonique du diagramme $(*)$ dans $(**)$, qui est l'identité sur les sommets $X_0$ et $\cX$, et qui sur $X_1$, $X_2$ est donné par les inclusions canoniques
\[
\underbrace{X_1}_{\text{\scriptsize $\cX_1$ fibre de $\cX$ en 1}} \longrightarrow 1\backslash\cX, \qquad
\underbrace{X_2}_{\text{\scriptsize $\cX_2$ fibre de $\cX$ en 2}} \longrightarrow 2\backslash\cX.
\]
Ces inclusions, on le sait, sont $W$-asphériques (elles ont même un adjoint à droite), a fortiori elles sont dans $W$. Donc le diagramme non commutatif $(*)$ est isomorphe, en tant que diagramme dans $\Hot_W$, au diagramme commutatif cartésien et cocartésien $(**)$ formé d'immersions fermées. Comme $\cX$ est bel et bien somme amalgamée de ses deux fermées % [sic] fermées -> fermés, cf. « les deux fermés » p. 83
$1\backslash\cX$, $2\backslash\cX$ sous leur intersection $X_0$, cela étaye l'intuition qui voit $\cX$ comme une `somme amalgamée homotopique'. Donnons tout de suite l'application suivante :

\textbf{Proposition 13.} \textit{Soit un diagramme} $\begin{tikzcd}[column sep=small, row sep=tiny, ampersand replacement=\&, baseline=(\tikzcdmatrixname-2-1.base)]
 \& X_1 \\
X_0 \arrow[ur, "\alpha_1"] \arrow[dr, "\alpha_2"'] \& \\
 \& X_2
\end{tikzcd}$ \textit{dans} Cat\textit{, d'où le diagramme} $(*)$\textit{, p. 123,}

\printedpage{85}
% head: A. Grothendieck : Dérivateurs, Ch. XII, Édition des Universités de Montpellier II et Paris VII
\[
\begin{tikzcd}[column sep=small, row sep=small, ampersand replacement=\&]
 \& X_0 \arrow[dl, "\alpha_1"'] \arrow[dr, "\alpha_2"] \& \\
X_1 \arrow[dr, "i_1"'] \& \& X_2 \arrow[dl, "i_2"] \\
 \& \cX, \&
\end{tikzcd}
\]
\msp{page 126}
\textit{commutatif dans} $\Hot_W$ \textit{(mais non dans} Cat\textit{, sauf si $X_0 = \emptyset$). Supposons que $W$ satisfait} W(7 bis) (axiome du carré cocartésien, cf. p. 75). \textit{Alors on a}
\[
\left\{
\begin{aligned}
\alpha_1\in W &\iff i_2\in W\\
\alpha_2\in W &\iff i_1\in W. \text{ (\footnotemark[85])}
\end{aligned}
\right.
\]
\footnotetext[85]{\textbf{NB} Pour l'implication $\Longrightarrow$ on n'a besoin que de la forme directe de W(7 bis), i.e. (avec les notations de XVI) de Loc(6$'$a).}
\textsc{Démonstration}. Utilisant l'homomorphisme précédent de $(*)$ dans $(**)$, on voit que chacune des flèches $\alpha_1$, $\alpha_2$, $i_1$, $i_2$ est dans $W$, si et seulement si la flèche correspondante $\alpha'_1$, $\alpha'_2$, $i'_1$ ou $i'_2$ l'est. Donc les équivalences ci-dessus équivalent à
\[
\begin{aligned}
\alpha'_1\in W &\iff i'_2\in W\\
\alpha'_2\in W &\iff i'_1\in W,
\end{aligned}
\]
lesquelles sont vraies en vertu de W(7 bis) (formulation duale) appliqué à $\cX$ et aux deux fermés $1\backslash\cX$, $2\backslash\cX$ de $\cX$.

\textbf{Corollaire :} \textit{Considérons} deux \textit{diagrammes de type} $\Phi^{\circ}$\textit{, et un homomorphisme} $f = (f_0, f_1, f_2)$ \textit{de l'un dans l'autre (cf. diagramme en traits pleins ci-contre} \texttt{[ci-dessous]}\textit{):}
\[
\begin{tikzcd}[column sep=small, row sep=small, ampersand replacement=\&]
 \& X_1 \arrow[dd, "f_1"'] \arrow[drr, dotted, "i_1"] \& \& \\
X_0 \arrow[ur, "\alpha_1"] \arrow[drr, "\alpha_2"] \arrow[dd, "f_0"'] \& \& \& \cX \arrow[dd, dotted, "f_{\cX}"] \\
 \& X'_1 \arrow[drr, dotted, "i'_1"] \& X_2 \arrow[ur, dotted, "i_2"] \arrow[dd, "f_2"'] \& \\
X'_0 \arrow[ur, "\alpha'_1"] \arrow[drr, "\alpha'_2"] \& \& \& \cX', \\
 \& \& X'_2 \arrow[ur, dotted, "i'_2"] \&
\end{tikzcd}
\]
\textit{d'où un homomorphisme}
\msp{page 127}
\[
f_{\cX} = \int f : \quad \cX = \int
\begin{tikzcd}[column sep=small, row sep=tiny, ampersand replacement=\&, baseline=(\tikzcdmatrixname-2-1.base)]
 \& X_1 \\
X_0 \arrow[ur] \arrow[dr] \& \\
 \& X_2
\end{tikzcd}
\qquad\longrightarrow\qquad
\cX' = \int
\begin{tikzcd}[column sep=small, row sep=tiny, ampersand replacement=\&, baseline=(\tikzcdmatrixname-2-1.base)]
 \& X'_1 \\
X'_0 \arrow[ur] \arrow[dr] \& \\
 \& X'_2,
\end{tikzcd}
\]

\printedpage{86}
% head: A. Grothendieck : Dérivateurs, Ch. XII, Édition des Universités de Montpellier II et Paris VII
\textit{s'insérant dans un diagramme cubique (cf. diagramme précédent complété par des pointillés). Considérons les conditions suivantes.}

\begin{itemize}
\item[\textit{a)}] $f_0\in W$\textit{.}
\item[\textit{b)}] $f_1$ \textit{et} $f_2$ \textit{sont} \texttt{[dans]} $W$\textit{.}
\item[\textit{c)}] $f_{\cX} = \int(f_0, f_1, f_2)\in W$\textit{.}
\end{itemize}

\textit{Supposons l'axiome} W(7) (`lemme des 5', cf. page 76) \textit{satisfait. Alors deux parmi les conditions a), b), c) impliquent la troisième}
\[
\left\{
\begin{aligned}
a \mathbin{\&} b &\quad\Longrightarrow\quad c\\
b \mathbin{\&} c &\quad\Longrightarrow\quad a\\
c \mathbin{\&} a &\quad\Longrightarrow\quad b.
\end{aligned}
\right.
\]

En effet, on a un homomorphisme du diagramme cubique ci-dessus dans le diagramme similaire, où $X_1$, $X_2$, $X'_1$, $X'_2$ sont remplacés respectivement pas les fermés % [sic] pas -> par, misprint
$1\backslash\cX$, $2\backslash\cX$ de $\cX$, et $1\backslash\cX'$, $2\backslash\cX'$ de $\cX'$, les deux carrés supérieur et inférieur étant des inclusions d'ouverts \texttt{[plutôt de fermés]} de $\cX$ resp. de $\cX'$. Ainsi, les conditions a), b), c) sont équivalentes respectivement aux conditions similaires pour le nouveau diagramme cubique
\msp{page 128}
(a) et c) n'ayant d'ailleurs pas changé, et b) donnant
\begin{itemize}
\item[b$'$)] $f'_1 : \cX/1 \longrightarrow \cX'/1$\ \ et\ \ $f'_2 : \cX/2 \longrightarrow \cX'/2$\ \ sont dans $W$.
\end{itemize}
\texttt{[plutôt} $f'_1 : 1\backslash\cX \longrightarrow 1\backslash\cX'$ \texttt{et} $f'_2 : 2\backslash\cX \longrightarrow 2\backslash\cX'$\texttt{]}). Alors la conclusion du corollaire est équivalente à W(7) (formulation duale) dans le cas du nouveau diagramme cubique, q.e.d.

\subsection*{3.2\quad Comparaison de \ $\displaystyle\int
\begin{tikzcd}[column sep=small, row sep=tiny, ampersand replacement=\&, baseline=(\tikzcdmatrixname-2-1.base)]
 \& X_1 \\
X_0 \arrow[ur] \arrow[dr] \& \\
 \& X_2
\end{tikzcd}$ \ avec une somme amalgamée ordinaire.\\
Carrés $W$-homotopiquement cocartésiens.}

La catégorie intégrée $\cX = \displaystyle\int
\begin{tikzcd}[column sep=small, row sep=tiny, ampersand replacement=\&, baseline=(\tikzcdmatrixname-2-1.base)]
 \& X_1 \\
X_0 \arrow[ur] \arrow[dr] \& \\
 \& X_2
\end{tikzcd}$ est solution du problème universel suivant, variante 2-catégorique de celui qui définit la limite inductive des diagrammes. Il s'agit, pour $Y$ objet variable dans Cat, de considérer les systèmes de trois foncteurs $f_0$, $f_1$, $f_2$ et des deux homomorphismes fonctoriels $\lambda_1$, $\lambda_2$ :
\printedpage{87}
% head: A. Grothendieck : Dérivateurs, Ch. XII, Édition des Universités de Montpellier II et Paris VII
\[
\begin{tikzcd}[column sep=large, row sep=large, ampersand replacement=\&]
 \& X_0 \arrow[dl, "\alpha_1"'] \arrow[dr, "\alpha_2"] \arrow[dd, "f_0", ""{name=M, pos=0.62}] \& \\
X_1 \arrow[dr, "f_1"', ""{name=A, pos=0.45}] \& \& X_2 \arrow[dl, "f_2", ""{name=B, pos=0.45}] \\
 \& Y, \&
\arrow[from=A, to=M, bend left, shorten=3pt, "\lambda_1"]
\arrow[from=B, to=M, bend right, shorten=3pt, "\lambda_2"']
\end{tikzcd}
\tag*{$(*)$}
\]
\[
\begin{aligned}
\lambda_1 :\quad & f_1\alpha_1 \quad\longrightarrow\quad f_0\\
\lambda_2 :\quad & f_2\alpha_2 \quad\longrightarrow\quad f_0.
\end{aligned}
\]

\msp{page 129}
Se donner une flèche de Cat
\[
\cX \longrightarrow Y
\]
`revient au même' que de se donner un système $(f_0, f_1, f_2, \lambda_1, \lambda_2)$ comme dessus, et ceci `fonctoriellement en $Y$'. C'est là le cas particulier de la détermination générale des foncteurs de la catégorie fibrée `intégrée' d'un contrafoncteur
\[
I^{\circ} \longrightarrow \Cat, \qquad
\begin{aligned}
i &\quad\longmapsto\quad X_i\\
u &\quad\longmapsto\quad u^*
\end{aligned}
\]
vers une catégorie $Y$, en termes d'un système de foncteurs
\[
X_i \xrightarrow{f_i} Y,
\]
et d'homomorphismes fonctoriels
\[
f_i\circ u^* \xrightarrow{\lambda_u} f_j
\]
pour $u : i \longrightarrow j$, d'où $u^* : X_j \longrightarrow X_i$, avec des conditions de transitivité pour des flèches composées, dont on n'a pas à tenir compte dans le cas actuel, vu qu'il n'y a pas de flèches composées à regarder dans $I = \Phi$ (sauf des composées avec une flèche identique, qu'on peut se dispenser de regarder ici \dots).

\msp{page 130}
Si en particulier on a un diagramme \textit{commutatif} dans Cat


\endgroup
% END UNIT Der-ChXIIb-012/Der-ChXIIb-012_p041-045.tex
% BEGIN UNIT Der-ChXIIb-012/Der-ChXIIb-012_p046-050.tex
\begingroup\makeatletter%
\newcommand{\printedpage}[1]{\marginpar{\footnotesize p.~#1}}%
\newcommand{\msp}[1]{\marginpar{\footnotesize [#1]}}%
\newcommand{\sX}{\mathscr{X}}%
\newcommand{\tX}{\tilde{\mathscr{X}}}%
\makeatother\let\printedpage\derpp\let\msp\dermsp%
\let\mspin\DERmspin%


\printedpage{88}
% head: A. Grothendieck : Dérivateurs, Ch. XII, Édition des Universités de Montpellier II et Paris VII
\[
Q \quad = \qquad
\begin{tikzcd}
 & X_0 \arrow[dl, "\alpha_1"'] \arrow[dr, "\alpha_2"] & \\
X_1 \arrow[dr, "f_1"'] & & X_2 \arrow[dl, "f_2"] \\
 & Y, &
\end{tikzcd}
\]
on peut le compléter par
\[
f_0 \overset{\text{déf}}{=} f_1\alpha_1 = f_2\alpha_2
\]
en un diagramme
\[
\begin{tikzcd}[row sep=large, column sep=large]
 & X_0 \arrow[dl, "\alpha_1"'] \arrow[dr, "\alpha_2"] \arrow[dd, "f_0", ""'{name=F0, pos=0.65}] & \\
X_1 \arrow[dr, "f_1"', ""{name=F1, pos=0.4}] & & X_2 \arrow[dl, "f_2", ""'{name=F2, pos=0.4}] \\
 & Y, &
\arrow[from=F1, to=F0, bend left, "\mathrm{id}", shorten <=2pt, shorten >=2pt]
\arrow[from=F2, to=F0, bend right, "\mathrm{id}"', shorten <=2pt, shorten >=2pt]
\end{tikzcd}
\]
ce qui montre qu'on a une donnée du type $(*)$ ci-dessus, définissant donc une flèche
\[
\varphi_Q : \quad \sX = \int
\begin{tikzcd}[row sep=small, column sep=small]
 & X_1 \\
X_0 \arrow[ur] \arrow[dr] & \\
 & X_2
\end{tikzcd}
\qquad \longrightarrow \qquad Y
\]
qui donne lieu à un diagramme cubique

\msp{page 131}
\[
\begin{tikzcd}[row sep=small, column sep=small]
 & X_1 \arrow[dd, "\mathrm{id}"'] \arrow[rrd, "i_1", ""'{name=I1}] & & \\
X_0 \arrow[ru, "\alpha_1"] \arrow[rrr, dotted, "i_0", ""'{name=I0, pos=0.8}] \arrow[rrd, "\alpha_2"] \arrow[dd, "\mathrm{id}"'] & & & \sX \arrow[dd, "\varphi"] \\
 & X_1 \arrow[rrd, "f_1"] & X_2 \arrow[ru, "i_2"', ""{name=I2}] \arrow[dd, "\mathrm{id}"'] & \\
X_0 \arrow[ru, "\alpha_1"] \arrow[rrr, dotted, "f_0"] \arrow[rrd, "\alpha_2"] & & & Y, \\
 & & X_2 \arrow[ru, "f_2"'] &
\arrow[from=I1, to=I0, "\lambda_1"', shorten <=1pt, shorten >=1pt]
\arrow[from=I2, to=I0, bend left, "\lambda_2", shorten <=1pt, shorten >=1pt] % Resolved from PDF leaf 11: a solid curve, unlike the genuinely dotted i_0 line.
\end{tikzcd}
\]
% foot: Version 7 décembre 2002

\printedpage{89}
% head: A. Grothendieck : Dérivateurs, Ch. XII, Édition des Universités de Montpellier II et Paris VII
commutatif sauf le carré supérieur, i.e. on a
\[
\varphi i_1 = f_1, \qquad \varphi i_2 = f_2, \qquad \text{et de plus} \quad \varphi i_0 = f_0,
\]
$\varphi$ transformant de plus les homomorphismes fonctoriels
\[
\lambda_1 : i_1\alpha_1 \longrightarrow i_0, \qquad \lambda_2 : i_2\alpha_2 \longrightarrow i_0
\]
en les homomorphismes fonctoriels identiques pertinents.

\setcounter{footnote}{85}
\textbf{Définition 3.} On dit que le diagramme commutatif
\[
\begin{tikzcd}[row sep=small, column sep=small]
 & X_0 \arrow[dl, "\alpha_1"'] \arrow[dr, "\alpha_2"] & \\
X_1 \arrow[dr, "f_1"] & & X_2 \arrow[dl, "f_2"'] \\
 & Y &
\end{tikzcd}
\tag*{$(*)$}
\]
est \textit{$W$-homotopiquement cocartésien} (\footnote{ou simplement $W$-cocartésien, ou qu'il définit $Y$ comme somme amalgamée $W$-homotopique de etc.}) si l'homomorphisme canonique
\[
\sX = \int
\begin{tikzcd}[row sep=small, column sep=small]
 & X_1 \\
X_0 \arrow[ur] \arrow[dr] & \\
 & X_2
\end{tikzcd}
\qquad \xrightarrow{\ \varphi\ } \qquad Y
\]
qu'il définit est dans $W$. (\footnote{Introduire $\varphi$ dans une proposition en forme, et donner ses propriétés essentielles : flèches strictement cartésiennes sont transformées en flèches \textit{identiques} $+$ lemmes 1 et 2 des p. 152, 154.})

\msp{page 132}
On a alors le :

\textbf{Corollaire 1.} \textit{Supposons} W(7 bis)\textit{. Si le carré} $(*)$ \textit{est $W$-homotopiquement cocartésien, alors}
\[
\left\{
\begin{aligned}
\alpha_1 \in W &\quad\iff\quad f_2 \in W \\
\alpha_2 \in W &\quad\iff\quad f_1 \in W.
\end{aligned}
\right.
\]

Cela résulte aussitôt du diagramme cubique de la page précédente, où toutes les flèches verticales sont dans $W$, et de l'énoncé similaire concernant la face supérieure du cube, donné dans la proposition 13, p. 125.

\textbf{Corollaire 2.} \textit{Supposons} W(7)\textit{. Considérons un diagramme cubique}
% foot: Version 7 décembre 2002

\printedpage{90}
% head: A. Grothendieck : Dérivateurs, Ch. XII, Édition des Universités de Montpellier II et Paris VII
\[
\begin{tikzcd}[row sep=small, column sep=small]
 & X_1 \arrow[dd, "p_1"'] \arrow[rrd, "f_1"] & & \\
X_0 \arrow[ru, "\alpha_1"] \arrow[rrd, "\alpha_2"] \arrow[dd, "p_0"'] & & & Y \arrow[dd, "p"] \\
 & X'_1 \arrow[rrd, "f'_1"] & X_2 \arrow[ru, "f_2"] \arrow[dd, "p_2"'] & \\
X'_0 \arrow[ru, "\alpha'_1"] \arrow[rrd, "\alpha'_2"] & & & Y', \\
 & & X'_2 \arrow[ru, "f'_2"] &
\end{tikzcd}
\]
\textit{où les deux faces inférieure et supérieure du cube sont des carrés $W$-homotopiquement cocartésiens. Considérons les conditions}

\msp{page 133}
\begin{enumerate}[label={}]
\item[\textit{a)}] $p_0 \in W$\textit{.}
\item[\textit{b)}] $p_1$ \textit{et} $p_2 \in W$\textit{.}
\item[\textit{c)}] $p \in W$\textit{.}
\end{enumerate}

\textit{Alors deux de ces conditions impliquent la troisième.}

Résulte aussitôt de la définition 3, et du résultat similaire (cor. page 126) pour le diagramme cubique similaire, avec $Y$, $Y'$ remplacées par $\sX = \int
\begin{tikzcd}[row sep=small, column sep=small]
 & X_1 \\
X_0 \arrow[ur] \arrow[dr] & \\
 & X_2
\end{tikzcd}$
et $\sX' = \int
\begin{tikzcd}[row sep=small, column sep=small]
 & X'_1 \\
X'_0 \arrow[ur] \arrow[dr] & \\
 & X'_2
\end{tikzcd}$\;.

\section*{3.3\quad Intégration propre versus intégration lisse. Énoncé d'équivalence entre les deux points de vue. Commentaires sur le cas des limites $W$-homotopiques quelconques.}

La définition 3, et les corollaires précédents, n'ont d'intérêt que dans la mesure où on dispose de critères maniables, permettant d'affirmer qu'un carré commutatif est bel et bien homotopiquement cocartésien. C'est ce qu'on va essayer de faire tantôt. Il y a aussi une perplexité, concernant le rôle privilégié joué, dans la théorie développée ici, par les immersions fermées, à l'exclusion des immersions ouvertes. Nous avons en somme ramené, du
% foot: Version 7 décembre 2002
\printedpage{91}
% head: A. Grothendieck : Dérivateurs, Ch. XII, Édition des Universités de Montpellier II et Paris VII
\msp{page 134}
point de vue homotopique, les diagrammes quelconques $\begin{tikzcd}[row sep=small, column sep=small]
 & X_1 \\
X_0 \arrow[ur] \arrow[dr] & \\
 & X_2
\end{tikzcd}$ dans Cat, au cas d'un diagramme d'immersions fermées, où la somme amalgamée prend des allures familières. Mais pourquoi pas aussi bien le ramener au cas d'immersions ouvertes? On considérerait alors le diagramme de départ comme définissant une catégorie \textit{cofibrée} sur
\[
\Phi^{\circ} = \Psi = \
\begin{tikzcd}[row sep=small, column sep=small]
 & 1 \\
0 \arrow[ur] \arrow[dr] & \\
 & 2,
\end{tikzcd}
\]
soit $\tX$. On a dans $\Psi$ les deux ouverts
\[
\underbrace{\Psi/1}_{=\{0,1\}} = U_1, \qquad \underbrace{\Psi/2}_{=\{0,2\}} = U_2, \qquad U_1 \cap U_2 = \{0\},
\]
d'où dans $\tX$ deux ouverts images inverses
\[
\tX/1, \qquad \tX/2, \qquad \text{avec } \tX/1 \cap \tX/2 = X_0,
\]
d'où un diagramme commutatif d'inclusions ouvertes
\[
\begin{tikzcd}
 & X_0 \arrow[dl, "\alpha'_1"'] \arrow[dr, "\alpha'_2"] & \\
\tX/1 \arrow[dr, "i'_1"'] & & \tX/2 \arrow[dl, "i'_2"] \\
 & \tX, &
\end{tikzcd}
\tag*{$(*)$}
\]

\msp{page 135}
lequel reçoit le diagramme non commutatif
\[
\begin{tikzcd}[row sep=large, column sep=large]
 & X_0 \arrow[dl, "\alpha_1"'] \arrow[dr, "\alpha_2"] \arrow[dd, dotted, "i_0", ""'{name=J0, pos=0.6}] & \\
X_1 \arrow[dr, "i_1"', ""{name=J1, pos=0.4}] & & X_2 \arrow[dl, "i_2", ""'{name=J2, pos=0.4}] \\
 & \tX, &
\arrow[from=J0, to=J1, dotted, bend right, "\tilde{\lambda}_1"', shorten <=2pt, shorten >=2pt]
\arrow[from=J0, to=J2, dotted, bend left, "\tilde{\lambda}_2", shorten <=2pt, shorten >=2pt]
\end{tikzcd}
\]
% foot: Version 7 décembre 2002
\printedpage{92}
% head: A. Grothendieck : Dérivateurs, Ch. XII, Édition des Universités de Montpellier II et Paris VII
les flèches
\[
X_1 \longrightarrow \tX/1, \qquad X_2 \longrightarrow \tX/2
\]
étant dans $W$ (et même $W$-coasphériques - elles ont même un adjoint à gauche).

Donc on trouve, comme précédemment
\[
\left\{
\begin{aligned}
\alpha_1 \in W &\quad\iff\quad i_2 \in W \\
\alpha_2 \in W &\quad\iff\quad i_1 \in W,
\end{aligned}
\right.
\]
et de même on peut formuler un ‘lemme des cinq’ comme le corollaire page 125 \texttt{[plutôt 126]}, mais avec $\tX$, $\tX'$ au lieu de $\sX$, $\sX'$.

De plus, $\tX$ est la solution d'un problème universel dans Cat, en termes de données
\[
\begin{tikzcd}[row sep=large, column sep=large]
 & X_0 \arrow[dl, "\alpha_1"'] \arrow[dr, "\alpha_2"] \arrow[dd, "f_0", ""'{name=G0, pos=0.6}] & \\
X_1 \arrow[dr, "f_1"', ""{name=G1, pos=0.4}] & & X_2 \arrow[dl, "f_2", ""'{name=G2, pos=0.4}] \\
 & Y, &
\arrow[from=G0, to=G1, bend right, "\tilde{\lambda}_1"', shorten <=2pt, shorten >=2pt]
\arrow[from=G0, to=G2, bend left, "\tilde{\lambda}_2", shorten <=2pt, shorten >=2pt]
\end{tikzcd}
\]

\msp{page 136}
la seule différence avec celui formulé à la page 128 étant la direction opposée des flèches $\lambda_1$, $\lambda_2$. Pour tout carré commutatif
\[
\begin{tikzcd}
 & X_0 \arrow[dl, "\alpha_1"'] \arrow[dr, "\alpha_2"] & \\
X_1 \arrow[dr, "f_1"'] & & X_2 \arrow[dl, "f_2"] \\
 & Y, &
\end{tikzcd}
\tag*{$(*)$}
\]
on trouve encore une flèche canonique
\[
\tilde{\varphi} : \tX \longrightarrow Y.
\]

Il serait tout aussi raisonnable de dire que le carré $(*)$ est $W$-homotopiquement cocartésien, lorsque $\tilde{\varphi}$ est \texttt{[dans]} $W$, que d'adopter la définition 3, page 131 - on trouverait exactement les mêmes corollaires. Les deux définitions seraient elles équivalentes?? % [sic] seraient elles -> seraient-elles, hyphen of the inverted subject missing
% foot: Version 7 décembre 2002


\endgroup
% END UNIT Der-ChXIIb-012/Der-ChXIIb-012_p046-050.tex
% BEGIN UNIT Der-ChXIIb-012/Der-ChXIIb-012_p051-055.tex
\begingroup\makeatletter%
\newcommand{\printedpage}[1]{\marginpar{\footnotesize p.~#1}}%
\newcommand{\msp}[1]{\marginpar{\footnotesize [#1]}}%
\newcommand{\Cat}{\mathrm{Cat}}%
\newcommand{\Hom}{\mathrm{Hom}}%
\newcommand{\HOT}{\mathrm{HOT}}%
\newcommand{\Hot}{\mathrm{Hot}}%
\newcommand{\Fl}{\mathrm{Fl}}%
\newcommand{\sX}{\mathscr{X}}%
\newcommand{\tsX}{\tilde{\mathscr{X}}}%
\newcommand{\circled}[1]{\tikz[baseline=(char.base)]{\node[shape=circle,draw,inner sep=1pt] (char) {#1};}}%
\makeatother\let\printedpage\derpp\let\msp\dermsp%
\let\mspin\DERmspin%


\printedpage{93}
% head: A. Grothendieck : Dérivateurs, Ch. XII, Édition des Universités de Montpellier II et Paris VII
% foot: Version 7 décembre 2002

Je note que $\tsX$ étant cofibrée sur $\Psi=\Phi^\circ$, $\tsX^\circ$ est fibrée sur $\Psi^\circ=\Phi$. De façon précise, on aura
\[
\tsX^\circ = \int
\begin{tikzcd}[row sep=small, column sep=small]
 & (X_1)^\circ \\
(X_0)^\circ \arrow[ur, "\alpha_1^\circ"] \arrow[dr, "\alpha_2^\circ"'] & \\
 & (X_2)^\circ,
\end{tikzcd}
\]
\msp{page 137}
et il s'agit donc de comparer les deux intégrales
\[
\sX = \int
\begin{tikzcd}[row sep=small, column sep=small]
 & X_1 \\
X_0 \arrow[ur, "\alpha_1"] \arrow[dr, "\alpha_2"'] & \\
 & X_2
\end{tikzcd}
, \qquad
\tsX^\circ = \int
\begin{tikzcd}[row sep=small, column sep=small]
 & (X_1)^\circ \\
(X_0)^\circ \arrow[ur, "\alpha_1^\circ"] \arrow[dr, "\alpha_2^\circ"'] & \\
 & (X_2)^\circ
\end{tikzcd}
,
\]
notamment comme objets de $\HOT_W(\Phi)$. Pour cela, je reprends la construction du foncteur
\[
X \longmapsto \Phi(X) : \Cat \longrightarrow \Cat
\]
de I, p. 72-75, mais je préfère ici noter ce foncteur $\Sigma(X)$ (car $\Phi$ désigne déjà une petite catégorie $\begin{tikzcd}[row sep=tiny, column sep=small]
 & 1 \arrow[dl] \\
0 & \\
 & 2 \arrow[ul]
\end{tikzcd}$). On utilisera les homomorphismes fonctoriels en $X$ :
\[
\begin{tikzcd}[column sep=small]
 & \Sigma(X) \arrow[dl, "p_X"'] \arrow[dr, "q_X"] & \\
X & & X^\circ,
\end{tikzcd}
\]
qui sont $W$-asphériques, donc $\in W$. De plus, $\Sigma(X^\circ)=\Sigma(X)$, $p_{X^\circ}=q_X$, $q_{X^\circ}=p_X$. Ainsi, si
\[
\begin{aligned}
X_* \quad &= \qquad
\begin{tikzcd}[ampersand replacement=\&, row sep=small, column sep=small]
 \& X_1 \\
X_0 \arrow[ur, "\alpha_1"] \arrow[dr, "\alpha_2"'] \& \\
 \& X_2
\end{tikzcd}
\qquad \in \quad \Hom(\Phi^\circ,\Cat)\\
X_*^\circ \quad &= \qquad
\begin{tikzcd}[ampersand replacement=\&, row sep=small, column sep=small]
 \& (X_1)^\circ \\
(X_0)^\circ \arrow[ur, "\alpha_1^\circ"] \arrow[dr, "\alpha_2^\circ"'] \& \\
 \& (X_2)^\circ
\end{tikzcd}
\qquad \in \quad \Hom(\Phi^\circ,\Cat),
\end{aligned}
\]
\msp{page 138}
on trouve
\[
\Sigma(X_*) =
\left(
\begin{tikzcd}[row sep=small, column sep=small]
 & \Sigma(X_1) \\
\Sigma(X_0) \arrow[ur, "\Sigma(\alpha_1)"] \arrow[dr, "\Sigma(\alpha_2)"'] & \\
 & \Sigma(X_2)
\end{tikzcd}
\right)
\]
et les flèches
\printedpage{94}
% head: A. Grothendieck : Dérivateurs, Ch. XII, Édition des Universités de Montpellier II et Paris VII
% foot: Version 7 décembre 2002
\[
\begin{tikzcd}[column sep=small]
 & \Sigma(X_*) \arrow[dl, "p_{X_*}"'] \arrow[dr, "q_{X_*}"] & \\
X_* & & X_*^\circ,
\end{tikzcd}
\]
dans $\underline{\Hom}(\Phi^\circ,\Cat)$. Ces flèches sont dans $W$ terme à terme, donc les homomorphismes dans $\Cat$
\[
\begin{tikzcd}[column sep=small]
 & \int\Sigma(X_*) \arrow[dl, "\int p_{X_*}"'] \arrow[dr, "\int q_{X_*}"] & \\
\sX=\int X_* & & \int X_*^\circ=\tsX^\circ
\end{tikzcd}
\]
sont dans $W$.

Finalement, l'innocent carré commutatif
\[
Q \quad = \qquad
\begin{tikzcd}[row sep=small, column sep=small]
 & X_0 \arrow[dl] \arrow[dr] & \\
X_1 \arrow[dr] & & X_2 \arrow[dl] \\
 & Y &
\end{tikzcd}
\tag*{$(*)$}
\]
(et le diagramme correspondant
\[
\widehat{Q} \quad = \qquad
\begin{tikzcd}[row sep=small, column sep=small]
 & X_0 \arrow[dl] \arrow[dr] & \\
X_1 \arrow[dr] & & X_2 \arrow[dl] \\
 & \sX=\int X_* \arrow[d, "\varphi_Q"] & \\
 & Y &
\end{tikzcd}
\qquad\qquad )
\tag*{$(**)$}
\]
donne naissance, en le transformant
\msp{page 139}
par $X\longmapsto X^\circ$ et par $X\longmapsto\Sigma(X)$, et en tenant compte des homomorphismes fonctoriels en $X$ $\begin{tikzcd}[row sep=small, column sep=small]
 & \Sigma(X) \arrow[dl] \arrow[dr] & \\
X & & X^\circ
\end{tikzcd}$, à trois diagrammes du type $(**)$, i.e. modelés sur la catégorie opposée de
\[
\Gamma \quad = \qquad
\begin{tikzcd}[row sep=small, column sep=small]
 & 0 & \\
1 \arrow[ur, "\alpha_1"] & & 2 \arrow[ul, "\alpha_2"'] \\
 & 3 \arrow[ul, "i_1"] \arrow[ur, "i_2"'] & \\
 & 4 \arrow[u, "\varphi"'] &
\end{tikzcd}
\qquad
\text{\parbox[c]{0.35\textwidth}{\footnotesize (i.e. $\boldsymbol{\Delta}^1\times\boldsymbol{\Delta}^1$ (\footnotemark[88]), auquelle on a rajouté un objet initial)}} % [sic] auquelle -> à laquelle, grammatical form (catégorie is feminine)
\]
\footnotetext[88]{faux, car le carré n'est pas commutatif}
\printedpage{95}
% head: A. Grothendieck : Dérivateurs, Ch. XII, Édition des Universités de Montpellier II et Paris VII
% foot: Version 7 décembre 2002
(\footnote[89]{On a $\alpha_1i_1\neq\alpha_2i_2$, mais la relation $\alpha_1i_1\varphi=\alpha_2i_2\varphi$.}), savoir
\[
\widehat{Q}, \qquad \widehat{\Sigma(Q)}, \qquad \widehat{Q^\circ},
\]
et deux homomorphismes de diagrammes (\texttt{[dans]} $\underline{\Hom}(\Gamma^\circ,\Cat)$)
\[
\begin{tikzcd}[column sep=small]
 & \widehat{\Sigma(Q)} \arrow[dl] \arrow[dr] & \\
\widehat{Q} & & \widehat{Q^\circ}
\end{tikzcd}
\qquad
\text{(déduit de }
\begin{tikzcd}[column sep=small]
 & \Sigma(Q) \arrow[dl] \arrow[dr] & \\
Q & & Q^\circ
\end{tikzcd}
\text{),}
\tag*{$(***)$}
\]
lesquels sont dans $W$ terme à terme (i.e. sont dans $W_\Gamma\subseteq\Fl(\underline{\Hom}(\Gamma^\circ,\Cat))$). Explicitement, en les sommets $i=0,1,2$, resp. $4$ de $\Gamma$, $(***)$ n'est autre que
\[
D_i =
\left(
\begin{tikzcd}[column sep=small]
 & \Sigma(X_i) \arrow[dl, "p_{X_i}"'] \arrow[dr, "q_{X_i}"] & \\
X_i & & X_i^\circ
\end{tikzcd}
\right)
\qquad \text{resp.} \qquad
\begin{tikzcd}[column sep=small]
 & \Sigma(Y) \arrow[dl, "p_Y"'] \arrow[dr, "q_Y"] & \\
Y & & Y^\circ,
\end{tikzcd}
\]
en le sommet 3 c'est
\[
\begin{tikzcd}[column sep=small]
 & \int\Sigma(X_*) \arrow[dl, "\int p_{X_*}"'] \arrow[dr, "\int q_{X_*}"] & \\
\sX=\int X_* & & \int X_*^\circ=\tsX^\circ,
\end{tikzcd}
\]
obtenu par intégration des flèches
\[
\begin{tikzcd}[column sep=small]
 & \Sigma(X_*) \arrow[dl, "p_{X_*}"'] \arrow[dr, "q_{X_*}"] & \\
X_* & & X_*^\circ.
\end{tikzcd}
\]
\msp{page 140}
Regardant la restriction des deux morphismes de diagrammes $(***)$ aux parties des trois diagrammes de type $\Gamma^\circ$ au dessus de la queue $\{3,4\}$ de $\Gamma$ on trouve le diagramme en traits pleins que voici
\[
\begin{tikzcd}[column sep=3em, row sep=normal]
 & \int\Sigma(X_*) \arrow[dl, "W"'] \arrow[dr, "W"] \arrow[dd, "\varphi_{\Sigma(Q)}"'] & & \Sigma\tsX \arrow[dl, dotted, "W"'] \arrow[dr, dotted, "W"] \arrow[dd, dotted, "\Sigma\tilde{\varphi}_Q"] & \\
\sX \arrow[dd, "\varphi_Q"'] & & \tsX^\circ \arrow[dd, "\varphi_{Q^\circ}=(\tilde{\varphi}_Q)^\circ"'] & & \tsX \arrow[dd, dotted, "\tilde{\varphi}_Q"] \\
 & \Sigma(Y) \arrow[dl, "W"'] \arrow[dr, "W"] & & \Sigma(Y) \arrow[dl, dotted, "W"'] \arrow[dr, dotted, "W"] & \\
Y & & Y^\circ & & Y
\end{tikzcd}
\]
\printedpage{96}
% head: A. Grothendieck : Dérivateurs, Ch. XII, Édition des Universités de Montpellier II et Paris VII
% foot: Version 7 décembre 2002
qu j'ai complété par un deuxième volet en pointillés. Dans ce diagramme de quatre carrés accolés, les côtés horizontaux des carrés représentent des flèches de $\Cat$ qui sont dans $W$. On trouve ainsi : % [sic] qu j'ai -> que j'ai, missing letter
\[
\varphi_Q\in W \iff \varphi_{\Sigma Q}\in W \iff (\varphi_{Q^\circ}=\tilde{\varphi}_Q^\circ)\in W \iff \Sigma\tilde{\varphi}_Q\in W \iff \tilde{\varphi}_Q\in W.
\]
On trouve aussi que les cinq flèches verticales sont isomorphes, en tant que flèches de $\Hot_W$. Pour résumer :

\textbf{Proposition 14.} \textit{Considérons un diagramme (carré commutatif)}
\[
Q \quad = \qquad
\begin{tikzcd}[row sep=small, column sep=small]
 & X_0 \arrow[dl] \arrow[dr] & \\
X_1 \arrow[dr] & & X_2 \arrow[dl] \\
 & Y &
\end{tikzcd}
\]
\textit{dans} $\Cat$\textit{, et le diagramme `opposé' déduit par} $X\longmapsto X^\circ$
\[
Q^\circ \quad = \qquad
\begin{tikzcd}[row sep=small, column sep=small]
 & X_0^\circ \arrow[dl] \arrow[dr] & \\
X_1^\circ \arrow[dr] & & X_2^\circ \arrow[dl] \\
 & Y^\circ, &
\end{tikzcd}
\]
\msp{page 141}
\textit{et les flèches correspondantes}
\[
\left\{
\begin{array}{rlclcl}
\varphi_Q : & \sX_Q & = &
\left(\displaystyle\int
\begin{tikzcd}[ampersand replacement=\&, row sep=small, column sep=small]
 \& X_1 \\
X_0 \arrow[ur] \arrow[dr] \& \\
 \& X_2
\end{tikzcd}
\right)
& \longrightarrow & Y \\[2ex]
\varphi_{Q^\circ} : & \sX_{Q^\circ} & = &
\left(\displaystyle\int
\begin{tikzcd}[ampersand replacement=\&, row sep=small, column sep=small]
 \& X_1^\circ \\
X_0^\circ \arrow[ur] \arrow[dr] \& \\
 \& X_2^\circ
\end{tikzcd}
\right)
& \longrightarrow & Y^\circ.
\end{array}
\right.
\]
\textit{Enfin, considérons aussi les flèches $\tilde{\varphi}_Q$, $\tilde{\varphi}_{Q^\circ}$ canoniques, relative à l'intégrale propre $\int_{\mathrm{pr}}$} % [sic] relative -> relatives, agreement with « les flèches »
\printedpage{97}
% head: A. Grothendieck : Dérivateurs, Ch. XII, Édition des Universités de Montpellier II et Paris VII
% foot: Version 7 décembre 2002
\textit{(au lieu de l'intégrale lisse $\int=\int_{\mathrm{lis}}$)}
\[
\left\{
\begin{array}{rlclcl}
\tilde{\varphi}_Q : & \tsX_Q & = &
\left(\displaystyle\int_{\mathrm{pr}}
\begin{tikzcd}[ampersand replacement=\&, row sep=small, column sep=small]
 \& X_1 \\
X_0 \arrow[ur] \arrow[dr] \& \\
 \& X_2
\end{tikzcd}
\right)
& \longrightarrow & Y \\[2ex]
\tilde{\varphi}_{Q^\circ} : & \tsX_{Q^\circ} & = &
\left(\displaystyle\int_{\mathrm{pr}}
\begin{tikzcd}[ampersand replacement=\&, row sep=small, column sep=small]
 \& X_1^\circ \\
X_0^\circ \arrow[ur] \arrow[dr] \& \\
 \& X_2^\circ
\end{tikzcd}
\right)
& \longrightarrow & Y^\circ.
\end{array}
\right.
\]
\textit{On a alors ce qui suit :}

\begin{enumerate}[label={}, leftmargin=2.5em]\itshape
\item[\circled{\textit{a}}] $\varphi_{Q^\circ}\simeq(\tilde{\varphi}_Q)^\circ$, $\varphi_Q\simeq(\tilde{\varphi}_{Q^\circ})^\circ$ (d'où résulte déjà que $\left\{\begin{array}{lcl}\varphi_Q\in W & \Longleftrightarrow & \tilde{\varphi}_{Q^\circ}\in W\\ \varphi_{Q^\circ}\in W & \Longleftrightarrow & \tilde{\varphi}_Q\in W\end{array}\right.$ (\footnote[90]{\upshape et que les flèches $\varphi_Q$ et $\tilde{\varphi}_{Q^\circ}$ d'une part, $\varphi_{Q^\circ}$ et $\tilde{\varphi}_Q$ d'autre part, sont isomorphes comme flèches dans $\Hot_W$.})).

\item[\circled{\textit{b}}] Les relations
\[
\varphi_Q\in W, \qquad \varphi_{Q^\circ}\in W, \qquad \tilde{\varphi}_Q\in W, \qquad \tilde{\varphi}_{Q^\circ}\in W
\]
sont équivalentes, en d'autres termes : pour que le carré soit $W$-homotopiquement cocartésien (version intégrale lisse), il faut et il suffit que le carré opposé $Q^\circ$ le soit, et ces conditions équivalent aussi aux conditions correspondantes en termes des intégrales propres (au lieu des intégrales lisses).

\item[\circled{\textit{c}}] Les quatre flèches $\varphi_Q$, $\varphi_{Q^\circ}$, $\tilde{\varphi}_Q$, $\tilde{\varphi}_{Q^\circ}$ de $\Cat$ en tant que flèches dans $\Hot_W$, sont isomorphes dans $\underline{\Fl}(\Hot_W)$. Plus précisément, en tant que flèches de $\Cat$,
\msp{page 142}
elles définissent des objets isomorphes de $\HOT_W(\boldsymbol{\Delta}^1)$.

\item[\circled{\textit{d}}] Plus précisément encore, les quatre diagrammes dans $\Cat$ de type $\Gamma^\circ$, où
\[
\Gamma = \;
\begin{tikzcd}[row sep=tiny, column sep=small]
 & 1 \arrow[dl, "\alpha_1"'] & & \\
0 & & 3 \arrow[ul, "i_1"'] \arrow[dl, "i_2"] & 4, \arrow[l, "\varphi"'] \\
 & 2 \arrow[ul, "\alpha_2"] & &
\end{tikzcd}
\]
$\alpha_1i_1\varphi=\alpha_2i_2\varphi$, dans lesquels s'insèrent les quatre flèches qu'on vient d'envisager, savoir
\[
\widehat{Q} = 
\begin{tikzcd}[row sep=tiny, column sep=small]
 & X_1 \arrow[dr, "i_1"] & & \\
X_0 \arrow[ur, "\alpha_1"] \arrow[dr, "\alpha_2"'] \arrow[rr, dotted] & & \sX_Q \arrow[r, "\varphi_Q"] & Y, \\
 & X_2 \arrow[ur, "i_2"'] & &
\end{tikzcd}
\]
\end{enumerate}


\endgroup
% END UNIT Der-ChXIIb-012/Der-ChXIIb-012_p051-055.tex
% BEGIN UNIT Der-ChXIIb-012/Der-ChXIIb-012_p056-060.tex
\begingroup\makeatletter%
\newcommand{\printedpage}[1]{\marginpar{\footnotesize p.~#1}}%
\newcommand{\msp}[1]{\marginpar{\footnotesize [#1]}}%
\newcommand{\Cat}{\mathrm{Cat}}%
\newcommand{\uHom}{\underline{\mathrm{Hom}}}%
\newcommand{\HOT}{\mathrm{HOT}}%
\newcommand{\Hot}{\mathrm{Hot}}%
\newcommand{\tld}{^{\tilde{}}}%
\makeatother\let\printedpage\derpp\let\msp\dermsp%
\let\mspin\DERmspin%

\setcounter{footnote}{90}

\printedpage{98}
% head: A. Grothendieck : Dérivateurs, Ch. XII, Édition des Universités de Montpellier II et Paris VII
% foot: Version 7 décembre 2002
\[
\begin{tikzcd}[row sep=small, column sep=small]
 & X_1^\circ \arrow[dr,"i^{\prime}_{1}"] & & \\
\widehat{Q^\circ} = X_0^\circ \arrow[ur,"\alpha_1^\circ"] \arrow[dr,"\alpha_2^\circ"'] \arrow[rr,dotted] & & \mathcal{X}_{Q^\circ} \arrow[r] & Y^\circ, \\
 & X_2^\circ \arrow[ur,"i^{\prime}_{2}"'] & &
\end{tikzcd}
\]
\[
\begin{tikzcd}[row sep=small, column sep=small]
 & X_1 \arrow[dr] & & \\
\widehat{Q}\tld = X_0 \arrow[ur] \arrow[dr] \arrow[rr,dotted] & & \tilde{\mathcal{X}}_{Q} \arrow[r] & Y, \\
 & X_2 \arrow[ur] & &
\end{tikzcd}
\]
\[
\begin{tikzcd}[row sep=small, column sep=small]
 & X_1^\circ \arrow[dr] & & \\
\widehat{Q^\circ}\tld = X_0^\circ \arrow[ur] \arrow[dr] \arrow[rr,dotted] & & \tilde{\mathcal{X}}_{Q^\circ} \arrow[r] & Y^\circ, \\
 & X_2^\circ \arrow[ur] & &
\end{tikzcd}
\]
(\footnote{\textbf{NB} On devrait encore compléter ces diagrammes (non commutatifs) par les flèches en pointillés \dots}) \textit{sont isomorphes dans} $\Hot_W(\Gamma)$\textit{. De façon plus précise encore, on a les isomorphismes dans} $\uHom(\Gamma^\circ,\Cat)$
\[
(\widehat{Q}\tld)^\circ \simeq \widehat{Q^\circ},
\]
\textit{d'où en l'appliquant à $Q^\circ$ au lieu de $Q$,}
\[
(\widehat{Q^\circ}\tld)^\circ \simeq \widehat{Q},
\]
\textit{ce qui implique (via les isomorphismes canoniques valables dans} $\HOT_W(S)$ \textit{pour tout objet $S$ dans} $\Cat$ \textit{et tout} $X_* \in \mathrm{Ob}\,\uHom(S,\Cat)$\textit{,} $X_* \simeq X_*^\circ$ (\footnote{via le foncteur $\Sigma$ utilisé plus haut \dots})\textit{), qu'on a les isomorphismes canoniques}
\[
\left.\begin{aligned}
\widehat{Q} &\simeq \widehat{Q^\circ}\tld \\
\widehat{Q^\circ} &\simeq \widehat{Q}\tld
\end{aligned}\right\}
\qquad \text{\textit{dans} } \HOT_W(\Gamma),
\]
\textit{mais aussi directement}
\[
\widehat{Q} \simeq \widehat{Q^\circ}, \qquad \widehat{Q}\tld \simeq \widehat{Q^\circ}\tld \qquad \text{\textit{dans} } \HOT_W(\Gamma)
\]
\msp{page 143}\textit{(et il reste à vérifier qu'on a ainsi un système} transitif \textit{d'isomorphismes canoniques dans} $\HOT_W(\Gamma)$\textit{, entre les quatre objets envisagés, ce qui doit être essentiellement sorital \dots).}

\subsection*{Remarques.}

\begin{enumerate}
\item[1)] Ce qu'on vient de dire dans le cas des diagrammes dans Cat de type $\Phi^\circ = \Psi$, a l'air de garder un sens pour les diagrammes de type quelconque
\[
X_* : S^\circ \longrightarrow \Cat, \qquad s \longmapsto X_s
\]
\printedpage{99}
% head: A. Grothendieck : Dérivateurs, Ch. XII, Édition des Universités de Montpellier II et Paris VII
% foot: Version 7 décembre 2002
($S$ un objet de Cat). On peut l'intégrer de façon lisse, $\int X_* = \int_{\mathrm{lis}} X_*$, obtenant une catégorie fibrée (scindée) sur $S$, ou de façon propre, $\int_{\mathrm{pr}} X_*$, obtenant une catégorie cofibrée scindée sur $S^\circ$. On peut également considérer $X_*^\circ$, composé de $X_*$ avec l'automorphisme $Y \longmapsto Y^\circ$ dans Cat, et l'intégrer de façon soit propre, soit lisse. Dans tous ces cas, on trouve un diagramme de type $\widehat{S}^\circ$, où $\widehat{S}$ est déduit de $S$ en ajoutant un élément qui s'envoie par une flèche ‘canonique’ (donnée) dans chacun des autres (\textit{sans} relations de commutation!) Un
\msp{page 144}\textit{système inductif} de morphismes dans Cat
\[
f \qquad X_s \xrightarrow{f_s} Y
\]
donne alors un morphisme
\[
\int_{\mathrm{lis}} X_* \xrightarrow{\varphi_Q} Y,
\]
où $Q$ désigne le diagramme, de type $\widehat{S}^\circ$, formé des $X_s$, de $Y$, des $u^* : X_t \longrightarrow X_s$ (pour $u : s \longrightarrow t$ dans $S$) et des $f_s$. De même on aura
\[
\int_{\mathrm{prop}} X_* \xrightarrow{\tilde{\varphi}_Q} Y,
\]
et si
\[
\Gamma = (\widehat{S} + \text{objet initial}),
\]
on trouve des diagrammes $\widehat{Q}$, $\widehat{Q}\tld$ de type $\Gamma^\circ$ dans Cat. Sauf erreur, les relations entre $\widehat{Q}$, $\widehat{Q}\tld$, $\widehat{Q^\circ}$, $\widehat{Q^\circ}\tld$ sont les mêmes que celles explicitées dans la proposition 14 ci-dessus. Je ne vais pas m'arrêter là-dessus avec le soin requis, sauf en cas de besoin. On dira qu'un diagramme $Q$ définit $Y$ comme $\varinjlim\limits_{I}$ $W$-homotopique \texttt{[où} $I = S^\circ$\texttt{]} si $\varphi_Q \in W$. Il revient au même de dire que $\varphi_{Q^\circ} \in W$, ou que $\tilde{\varphi}_Q \in W$ ou que $\tilde{\varphi}_{Q^\circ} \in W$ \dots
\msp{page 145}

\item[2)] Le cas ‘universel’, pour
\begin{tikzcd}[baseline=(\tikzcdmatrixname-2-1.base), row sep=small, column sep=small]
 & X_1 \\
X_0 \arrow[ur] \arrow[dr] & \\
 & X_2
\end{tikzcd}
donné, d'un carré commutatif $Q = {}$%
\begin{tikzcd}[baseline=(\tikzcdmatrixname-2-1.base), row sep=tiny, column sep=small]
 & X_1 \arrow[dr] & \\
X_0 \arrow[ur] \arrow[dr] & & Y \\
 & X_2 \arrow[ur] &
\end{tikzcd}
dans Cat, donnant naissance à
\[
\begin{tikzcd}[row sep=small, column sep=small]
 & & & X_1 & & & \\
\varphi_Q : &[1em] \left(\vrule height 6ex depth 6ex width 0pt\right. &[-1.5em] \int X_0 \arrow[ur] \arrow[dr] & & \left.\vrule height 6ex depth 6ex width 0pt\right) &[2em] \longrightarrow &[2em] Y, \\
 & & & X_2 & & &
\end{tikzcd}
\]
est celui de $Y = X_1 \amalg_{X_0} X_2$, donnant naissace à % [sic] naissace -> naissance, letter n omitted
\printedpage{100}
% head: A. Grothendieck : Dérivateurs, Ch. XII, Édition des Universités de Montpellier II et Paris VII
% foot: Version 7 décembre 2002
\[
\begin{tikzcd}[row sep=small, column sep=small]
 & X_1 \arrow[dr] & \\
Q_0 = Q(X_*) = X_0 \arrow[ur] \arrow[dr] & & X_1 \amalg_{X_0} X_2. \\
 & X_2 \arrow[ur] &
\end{tikzcd}
\]

Le cas général s'y ramène, car on aura un diagramme commutatif
\[
\begin{tikzcd}[row sep=small, column sep=small]
 & & X_1 & & & \\
\left(\vrule height 6ex depth 6ex width 0pt\right. &[-1.5em] \int X_0 \arrow[ur] \arrow[dr] & & \left.\vrule height 6ex depth 6ex width 0pt\right) \arrow[rr,"{\varphi_{Q_0}=\varphi_{Q(X_*)}}"] \arrow[ddrr,"\varphi_Q"] &[4em] \hphantom{MMMM} &[4em] X_1 \amalg_{X_0} X_2 \arrow[dd,"\alpha_Q"] \\
 & & X_2 & & & \\
 & & & & & Y.
\end{tikzcd}
\]
\end{enumerate}

Mais il doit y avoir des cas intéressants où $\varphi_{Q_0}$ n'est \textit{pas} dans $W$, i.e. le carré Cat-cocartésien
\begin{tikzcd}[baseline=(\tikzcdmatrixname-2-1.base), row sep=tiny, column sep=small]
 & X_1 \arrow[dr] & \\
X_0 \arrow[ur] \arrow[dr] & & X_1 \amalg_{X_0} X_2 \\
 & X_2 \arrow[ur] &
\end{tikzcd}
n'est pas $W$-homotopiquement cocartésien, et où pourtant on peut trouver un carré commutatif $Q$
\msp{page 146}complétant $X_*$, de telle façon que ce carré soit $W$-homotopiquement commutatif \texttt{[plutôt cocartésien]}, i.e. que $\varphi_Q = \alpha_Q \varphi_{Q_0}$ soit dans $W$. C'est d'autant plus plausible que dans Cat, les sommes amalgamées ont tendance à ne pas avoir des bonnes propriétés homotopiques, même si l'une des flèches $\alpha_1$, $\alpha_2$ est \texttt{[un]} monomorphisme (voire une immersion ouverte), contrairement à ce qui se passe dans $\boldsymbol{\Delta}^\wedge$, et plus généralement dans les catégories de la forme $Z^\wedge$ ($Z$ dans Cat).

\subsection*{3.4\quad Critères pour les carrés homotopiquement cocartésiens.}

Les critères que j'ai dégagé sont tous contenus dans l'énoncé suivant :

\textbf{Théorème 4} (\footnote{\textbf{NB} Dans ce théorème, et les corollaires 1, 2, 3, on n'a besoin que de la \textit{forme faible} de W(7 bis), savoir $i_0 \in W \Longrightarrow i \in W$ (cf. p. 79) (Loc(6$'$a) avec les notations de XVI).},\footnote{\textbf{NB} Au total, le théorème (partie directe) est valable sans hypothèse W(7 bis) ou W(7 bis faible) sur $W$, et il en est de même de ses corollaires 1 à 3 (à l'exclusion des corollaires 4, 5, qui eux utilisent W(7 bis) dans toute sa force).}). \textit{Considérons un diagramme carré commutatif dans} Cat
\[
Q \quad = \quad
\begin{tikzcd}
X_0 \arrow[r,"\alpha_1"] \arrow[d,"\alpha_2"'] \arrow[dr,dotted,"f_0"] & X_1 \arrow[d,"f_1"] \\
X_2 \arrow[r,"f_2"'] & Y,
\end{tikzcd}
\tag*{$(*)$}
\]
\printedpage{101}
% head: A. Grothendieck : Dérivateurs, Ch. XII, Édition des Universités de Montpellier II et Paris VII
% foot: Version 7 décembre 2002
\textit{et considérons le morphisme canonique}
\[
\begin{tikzcd}[row sep=small, column sep=small]
 & & X_1 & & & \\
\varphi_Q : \quad \mathcal{X} = \Bigl(\int &[-1.2em] X_0 \arrow[ur,"\alpha_1"] \arrow[dr,"\alpha_2"'] & & \Bigr) &[2em] \longrightarrow &[2em] Y, \\
 & & X_2 & & &
\end{tikzcd}
\]
\msp{page 147}\textit{l'intégrale étant l'intégrale} lisse. \textit{(Il y a un énoncé dual concernant l'intégrale propre de}\newline
\begin{tikzcd}[baseline=(\tikzcdmatrixname-2-1.base), row sep=small, column sep=small]
 & X_1 \\
X_0 \arrow[ur] \arrow[dr] & \\
 & X_2
\end{tikzcd}
\textit{.)}

\begin{enumerate}
\item[\textit{a)}] \textit{Pour que $\varphi_Q$ soit $W$-lisse, il faut et il suffit que $f_1$, $f_2$ et $f_0 = f_1\alpha_1 = f_2\alpha_2$ le soient.} (Moyennant W(7 bis) (\footnote{Pour la partie directe ($f_0$, $f_1$, $f_2$ lisses $\Longrightarrow \varphi_Q$ lisse), on n'a besoin que de W(7 bis) (\footnotemark), forme faible. Ce n'est que pour la réciproque (qui ne m'a encore jamais servie) qu'on a besoin de W(7 bis), dans toute sa force.}\footnotetext{On n'a pas même besoin de W(7 bis) faible, cf. annotation page 156. On utilise une hypothèse restrictive sur $W$ dans la réciproque seulement ($\varphi_Q$ lisse $\Longrightarrow f_0$, $f_1$, $f_2$ lisses).})).

\item[\textit{b)}] \textit{La fibre de $\varphi_Q$ en un élément $y$ de $Y$ est isomorphe à}
\[
\begin{tikzcd}[row sep=small, column sep=small]
 & & X_{1\,y} & \\
\int &[-1.2em] X_{0\,y} \arrow[ur] \arrow[dr] & &[2em] (\simeq \mathcal{X}_y). \\
 & & X_{2\,y} &
\end{tikzcd}
\tag*{$(**)$}
\]
\end{enumerate}

\textbf{Corollaire 1.} (\footnote{Valable sans hypothèse sur $W$ (autre que Loc(1,2,3 bis)).}) \textit{Supposons $f_1$, $f_2$ et $f_1\alpha_1 = f_2\alpha_2$ $W$-lisses, et les catégories intégrées $(**)$ $W$-asphériques pour tout $y \in \mathrm{Ob}\,Y$. (NB Il s'agit encore de l'intégrale lisse, mais peu importe au fond, car on a vu (cf. partie supérieure du diagramme p. 140, qui implique que $\mathcal{X}$} \texttt{[est]} \textit{$W$-asphérique si et seulement si $\tilde{X}$ l'est) que l'une est $W$-asphérique si et seulement si l'autre l'est.) Alors $\varphi_Q \in W$, i.e. le carré $(*)$ est $W$-homotopiquement cartésien} \texttt{[plutôt cocartésien]}\textit{. Par dualité, même conclusion si on suppose $f_1$, $f_2$, $f_1\alpha_1 = f_2\alpha_2$} propres \textit{au lieu de lisses. (Car $Q$ $W$-homotopiquement trivial} \texttt{[plutôt cocartésien]}\textit{, si et seulement si $Q^\circ$ l'est, par} \texttt{[la]} \textit{prop. 14.)}

\msp{page 148}En effet, un morphisme dans Cat qui est $W$-lisse et à fibres $W$-asphériques est dans $W$ (et c'est même une $W$-équivalence universelle \dots).

\textbf{Corollaire 2}. \textit{On suppose que $W$ satisfait} W(7 bis) (axiome du carré cocartésien) (\footnote{inutile.}).
\printedpage{102}
% head: A. Grothendieck : Dérivateurs, Ch. XII, Édition des Universités de Montpellier II et Paris VII
% foot: Version 7 décembre 2002
\textit{Supposons le carré $Q$ cocartésien dans} Cat (\footnote{NB On n'utilise pas vraiment que $Q$ soit cocartésien, plutôt que $Q$ est \textit{cartésien,} et $f_2$ (donc aussi $\alpha_1$) une immersion ouverte (resp. fermée), et que pour $y \in Y\backslash X_2$, $X_{1\,y}$ soit $W$-asphérique. (Dans le cas où $Q$ est cocartésien, $X_{1\,y}$ est $\simeq e$.)})\textit{, et que $\alpha_1$ soit une immersion ouverte (resp. fermée) et $f_1$ lisse (resp. propre). Alors $Q$ est $W$-homotopiquement cocartésien.} (\footnote{On a besoin ici de W(7 bis), forme directe, i.e. de Loc(6$'$a) (notation de XVI).})

Il suffit d'établir le résultat non respé, l'autre en résultera par dualité (puisque $Q$ est $W$-cocartésien si et seulement si $Q^\circ$ l'est). Par hypothèse $f_1$ est $W$-lisse, d'autre part $\alpha_1$ et $f_2$ sont des immersions ouvertes, donc $W$-lisses, donc $f_1\alpha_1$ est $W$-lisse aussi. (D'ailleurs $\alpha_2$ aussi, car le carré est aussi cartésien, donc $\alpha_2$ déduit de $f_1$ (qui est $W$-lisse par hypothèse) par changement de base $X_2 \longrightarrow Y$.) Donc $\varphi_Q$ est \texttt{[}$W$\texttt{-]}lisse, et il suffit de voir que ses fibres sont $W$-asphériques. Le calcul $(**)$ donne, pour $y$ dans $Y\backslash X_2$,
\[
\begin{tikzcd}[row sep=small, column sep=small]
 & & e & \\
\mathcal{X}_y \simeq \quad \int &[-1.2em] \emptyset \arrow[ur] \arrow[dr] & &[1em] \simeq e, \qquad \text{O.K.} \\
 & & \emptyset &
\end{tikzcd}
\]
Et pour $y$ dans $X_2$
\[
\begin{tikzcd}[row sep=small, column sep=small]
 & & X_{1\,y} & \\
\mathcal{X}_y \simeq \int &[-1.2em] X_{0\,y} \arrow[ur,"\sim"] \arrow[dr] & &[2em] . \\
 & & e &
\end{tikzcd}
\]
\msp{page 149}Comme $X_{0\,y} \longrightarrow X_{1\,y}$ est un isomorphisme, donc est dans $W$, il s'ensuit que $e \longrightarrow \mathcal{X}_y$ l'est aussi (prop. 13, p. 125), moyennant W(7 bis), donc $\mathcal{X}_y$ est $W$-asphérique, O.K. (\footnote{On n'a besoin que de la forme directe de W(7 bis), i.e. (notation de XVI) de Loc(6$'$a).}).

\textbf{Corollaire 3.} (Moyennant W(7 bis) (\footnote{inutile.}).) \textit{Soit $X$ dans} Cat\textit{, $U_1$, $U_2$ deux ouverts de $X$ tels que $X = U_1 \cup U_2$, $U_0 = U_1 \cap U_2$ l'ouvert intersection. Alors le carré d'immersions ouvertes}
\[
\begin{tikzcd}[row sep=large, column sep=large]
U_0 \arrow[r,hook] \arrow[d,hook] & U_1 \arrow[d,hook] \\
U_2 \arrow[r,hook] & X
\end{tikzcd}
\]
\textit{est $W$-cocartésien. (Énoncé dual pour deux fermés $T_1$, $T_2$ tels que $T_1 \cup T_2 = X$, avec $T_0 = T_1 \cap T_2$.)}

On applique le corollaire 2. On sait que le carré est cocartésien dans Cat (\footnote{On n'a pas même besoin de le savoir, cf. annotation marginale du corollaire 2.}), de plus $U_1 \lhook\joinrel\longrightarrow X$ est bien $W$-lisse, étant une immersion ouverte.


\endgroup
% END UNIT Der-ChXIIb-012/Der-ChXIIb-012_p056-060.tex
% BEGIN UNIT Der-ChXIIb-012/Der-ChXIIb-012_p061-065.tex
\begingroup\makeatletter%
\newcommand{\printedpage}[1]{\marginpar{\footnotesize p.~#1}}%
\newcommand{\msp}[1]{\marginpar{\footnotesize [#1]}}%
\newcommand{\hto}[1][]{\overset{#1}{\lhook\joinrel\longrightarrow}}%
\newcommand{\circled}[1]{\tikz[baseline=(c.base)]\node[draw,circle,inner sep=1pt](c){#1};}%
\newcommand{\intdiag}[3]{\begin{tikzcd}[row sep=small,column sep=small,ampersand replacement=\&,baseline=(\tikzcdmatrixname-2-1.base)]%
 \& #2 \\%
 \displaystyle\int #1 \arrow[ur] \arrow[dr] \& \\%
 \& #3%
\end{tikzcd}}%
\makeatother\let\printedpage\derpp\let\msp\dermsp%
\let\mspin\DERmspin%


\printedpage{103}
% head: A. Grothendieck : Dérivateurs, Ch. XII, Édition des Universités de Montpellier II et Paris VII

\noindent\textbf{Corollaire 4.} (Moyennant W(7 bis) (\footnote[104]{Ici on a besoin de W(7 bis) forme forte, i.e. $i_0\in W \iff i\in W$ (cf. annotation th. 4).}).) \textit{Soit $U \hto[i] X$ une immersion ouverte telle que $T=X\backslash U$ soit isomorphe à la catégorie ponctuelle $e$. Pour que $i\in W$, il faut et il suffit que $i$ soit $W$-asphérique, i.e. que $U/t$ (où $t$ est l'unique objet de $T$) soit $W$-asphérique.}

Il suffit bien sûr de prouver que
\[
i\in W \Longrightarrow U/t \ W\text{-asphérique.}
\]

\msp{page 150}
Considérons le carré cartésien
\[
\begin{tikzcd}[row sep=large,column sep=large]
U/t \arrow[r, hook, "i/t"] \arrow[d, "p"'] & X/t \arrow[d, "q"] \\
U \arrow[r, hook, "i"] & X,
\end{tikzcd}
\]
nous allons prouver qu'il est $W$-cocartésien. Moyennant W(7 bis), il s'ensuivra que $i\in W$ implique $i/t\in W$, et comme $X/t$ est $W$-asphérique (ayant un objet final), cela signifie que $U/t$ est $W$-asphérique, ce qu'on veut prouver.

Il suffit, en vertu du corollaire 2, de prouver que le carré est cocartésien dans Cat, puisque $i/t$ est une immersion ouverte ($i$ l'étant), et $X/t\longrightarrow X$ est un isomorphisme local, donc $W$-lisse. On voit que $(X/t)\backslash(U/t)\simeq e$ (catégorie ponctuelle), d'ailleurs l'unique élément de $(X/t)\backslash(U/t)$ est l'élément final $t_0 : t\xrightarrow{\mathrm{id}_t} t$ de $X/t$. Ainsi $H_{t_0}\simeq e_{U/t}$ (préfaisceau final), et $p_{!}(H_{t_0})\simeq H'_t \overset{\text{déf}}{=} (u\longmapsto \mathrm{Hom}_X(u,t))$, ce qui caractérise la somme amalgamée dans Cat
\dots

Le corollaire 4 n'est autre que le corollaire non démontré de la proposition 7 (p. 85).

Voici une variante du corollaire 4.

\noindent\textbf{Corollaire 5.} (Moyennant W(7 bis) (\footnote[105]{forme forte.}).) \textit{Soit $i : U \hto X$ une immersion ouverte, soit $t_0\in\mathrm{Ob}\,T$ un point \textup{fermé} de $T=X\backslash U$, tel que $\{t_0\}$}
\msp{page 151}
\textit{soit isomorphe à la catégorie ponctuelle, i.e. $\mathrm{Hom}_X(t_0,t_0)=\{\mathrm{id}\}$. Supposons que pour tout $t\in T\backslash\{t_0\}$, $U/t$ soit $W$-asphérique. Pour que l'on ait $i\in W$, il faut et il suffit que $i$ soit asphérique, i.e. que $U/t_0$ soit également $W$-asphérique.}

Il faut prouver que $i\in W$ implique que $U/t_0$ est $W$-asphérique. Considérons l'inclusion
\[
i_0 : U\longrightarrow U' \overset{\text{déf}}{=} X\backslash\{t_0\},
\]
l'hypothèse sur les $U/t$, pour $t\in T\backslash\{t_0\}$, signifie que $i_0$ est $W$-asphérique. Donc $i_0$ est dans $W$, et comme le composé $i : U\xrightarrow{i_0} U'\xrightarrow{i'} X$ est dans $W$, cela implique que
% foot: Version 7 décembre 2002
\printedpage{104}
% head: A. Grothendieck : Dérivateurs, Ch. XII, Édition des Universités de Montpellier II et Paris VII
$i' : U'\longrightarrow X$ est dans $W$. Nous sommes maintenant dans la situation du corollaire 4, qui implique que $i'$ est $W$-asphérique. Donc $i=i'i_0$ est $W$-asphérique, comme composé de deux morphismes $W$-asphériques, q.e.d.

\noindent\textsc{Démonstration de th. 4.}

\noindent\circled{a} Critère de lissité de $\varphi_Q : \mathcal{X} = \intdiag{X_0}{X_1}{X_2} \longrightarrow Y$.

Comme $\{1\}$, $\{2\}$ sont des ouverts de $\Phi = \begin{tikzcd}[row sep=small,column sep=small,baseline=(\tikzcdmatrixname-2-1.base)]
 & 1 \arrow[dl] \\
 0 & \\
 & 2 \arrow[ul]
\end{tikzcd}$, $i_1 : X_1\longrightarrow\mathcal{X}$ et $i_2 : X_2\longrightarrow\mathcal{X}$ sont des immersions ouvertes. Comme $\varphi_Q\circ i_1 = f_1$,
\msp{page 152}
$\varphi_Q\circ i_2=f_2$, on voit déjà que
\[
f_1 \text{ et } f_2 \text{ lisses} \iff \varphi_Q \text{ lisse sur l'ouvert réunion de } i_1(X_1) \text{ et } i_2(X_2)\text{, i.e. sur } \mathcal{X}\backslash i_0(X_0).
\]
Il reste à exprimer la lissité de $\varphi_Q$ en les points $x_0$ de $i_0(X_0)$. Soient $x_1$, $x_2$, $y$ les images de $x_0$ dans $X_1$, $X_2$, $Y$ respectivement. Il faut exprimer que les fibres de
\[
(\varphi_Q)/x_0 : \mathcal{X}/x_0\longrightarrow Y/y \tag*{$(*)$}
\]
sont asphériques. Nous allons pour cela expliciter $\mathcal{X}/x_0$ et la flèche $(*)$ en termes du carré commutatif
\[
(**)\qquad\qquad
\begin{tikzcd}[row sep=large,column sep=large]
X_0/x_0 \arrow[r, "\alpha_1/x_0"] \arrow[d, "\alpha_2/x_0"'] & X_1/x_1 \arrow[d, "f_1/x_1"] \\
X_2/x_2 \arrow[r, "f_2/x_2"] & Y/y,
\end{tikzcd}
\]
induit par $Q$. Je dis

\noindent\textbf{Lemme 1.} \textit{Soit un carré commutatif}
\[
Q \quad = \quad
\begin{tikzcd}[row sep=large,column sep=large]
X_0 \arrow[r, "\alpha_1"] \arrow[d, "\alpha_2"'] & X_1 \arrow[d, "f_1"] \\
X_2 \arrow[r, "f_2"] & Y
\end{tikzcd}
\]
\textit{dans} Cat, \textit{$x_0\in\mathrm{Ob}\,X_0$, $x_1$, $x_2$, $y$ ses images dans $X_1$, $X_2$, $Y$ respectivement, $Q/x_0$ le carré $(**)$ ci-dessus, induit par $Q$ sur les localisées $X_0/x_0$ etc., $\mathcal{X} = \intdiag{X_0}{X_1}{X_2}$. On a alors un}
% foot: Version 7 décembre 2002
\printedpage{105}
% head: A. Grothendieck : Dérivateurs, Ch. XII, Édition des Universités de Montpellier II et Paris VII
\msp{page 153}
\textit{isomorphisme canonique}
\[
\mathcal{X}/x_0 \quad\simeq\quad \intdiag{X_0/x_0}{X_1/x_1}{X_2/x_2} ,
\]
\textit{et la flèche $(\varphi_Q)/x_0$, localisée de $\varphi_Q : \mathcal{X}\longrightarrow Y$ en $x_0$, s'identifie à la flèche $\varphi_{(Q/x_0)}$ :}
\[
\varphi_{(Q/x_0)} : \quad \intdiag{X_0/x_0}{X_1/x_1}{X_2/x_2} \qquad\longrightarrow\qquad Y/y.
\]

\noindent\textsc{Démonstration} du Lemme. On sait que
\[
\mathcal{X}/x_0\longrightarrow\mathcal{X}
\]
est Cat-fibrée (scindée) (à fibres discrètes), et comme
\[
\mathcal{X}\longrightarrow\Phi
\]
est aussi Cat-fibrée (scindée), il s'ensuit qu'il en est de même de la flèche composée
\[
\mathcal{X}/x_0\longrightarrow\Phi.
\]
Donc pour `calculer' $\mathcal{X}/x_0$, en tant que catégorie sur $\Phi$, il faut expliciter ses trois fibres, et les deux foncteurs de changement de base qui les relient. On trouve aussitôt que ces fibres sont canoniquement isomorphes à
\[
X_0/x_0, \qquad X_1/x_1, \qquad X_2/x_2,
\]
et on admettra (!) que les foncteurs changement de base
\msp{page 154}
\[
X_0/x_0\longrightarrow X_1/x_1, \qquad X_0/x_0\longrightarrow X_2/x_2
\]
(définis à isomorphisme unique près) ne sont autres que les localisés $\alpha_1/x_0$ et $\alpha_2/x_0$ de $\alpha_1 : X_0\longrightarrow X_1$ et $\alpha_2 : X_0\longrightarrow X_2$. Cela montre donc qu'on a un isomorphisme
\[
\mathcal{X}/x_0 \quad\simeq\quad \intdiag{X_0/x_0}{X_1/x_1}{X_2/x_2} ,
\]
% foot: Version 7 décembre 2002
\printedpage{106}
% head: A. Grothendieck : Dérivateurs, Ch. XII, Édition des Universités de Montpellier II et Paris VII
et l'on admettra aussi que la flèche $(\varphi_Q)/x_0 : \mathcal{X}/x_0\longrightarrow Y/y$ s'identifie à $\varphi_{(Q/x_0)}$.

Pour exprimer la $W$-lissité de $\varphi_Q$ en $x_0\in\mathrm{Ob}\,X_0\subseteq\mathrm{Ob}\,\mathcal{X}$, on est donc ramené à exprimer la $W$-asphéricité des fibres de
\[
\varphi_{(Q/x_0)} : \quad \intdiag{X_0/x_0}{X_1/x_1}{X_2/x_2} \qquad\longrightarrow\qquad Y/y.
\]
Pour expliciter ces fibres, on a besoin d'un autre lemme auxiliaire

\noindent\textbf{Lemme 2} (\footnote[106]{NB C'est le b) du th. 4.}). \textit{Soit un carré commutatif dans} Cat
\[
Q \quad = \quad
\left(
\begin{tikzcd}[row sep=large,column sep=large]
X_0 \arrow[r, "\alpha_1"] \arrow[d, "\alpha_2"'] & X_1 \arrow[d, "f_1"] \\
X_2 \arrow[r, "f_2"] & Y
\end{tikzcd}
\right),
\]
\textit{d'où}
\[
\varphi_Q : \quad \left( \intdiag{X_0}{X_1}{X_2} \right) \qquad\longrightarrow\qquad Y,
\]

\msp{page 155}
\textit{et soit $y\in\mathrm{Ob}\,Y$. Alors la fibre de $\varphi_Q$ en $y$}
\[
\mathcal{X}_y = \varphi_Q^{-1}(\ \underbrace{\{y\}}_{\text{catégorie ponctuelle}}\ )
\]
\textit{est donnée par (comme catégorie au dessus de $\Phi$)}
\[
\mathcal{X}_y \quad\simeq\quad \intdiag{X_{0\,y}}{X_{1\,y}}{X_{2\,y}} .
\]

\noindent\textsc{Démonstration.} Il est immédiat qu'on a
\[
\mathcal{X}_y\cap\underbrace{\mathcal{X}_s}_{X_s}\simeq(X_s)_y \qquad\quad \text{pour } s\in\mathrm{Ob}\,\Phi=\{0,1,2\},
\]
% foot: Version 7 décembre 2002
\printedpage{107}
% head: A. Grothendieck : Dérivateurs, Ch. XII, Édition des Universités de Montpellier II et Paris VII
et (le hasard fait bien les choses) les foncteurs changement de base
\[
X_0\longrightarrow X_1, \qquad X_0\longrightarrow X_2
\]
envoient la sous-catégorie (non pleine) $X_{0\,y}$ de $X_0$ dans les sous-catégories (non pleines) $X_{1\,y}$ et $X_{2\,y}$ de $X_1$ et $X_2$ respectivement.
\[
\begin{tikzcd}[row sep=large,column sep=large]
X_{s\,y} \arrow[r] \arrow[d] \arrow[dr, phantom, "\scriptstyle\text{cart.}"] & \mathcal{X}_y \arrow[r] \arrow[d] \arrow[dr, phantom, "\scriptstyle\text{cart.}"] & \{y\} \arrow[d] \\
X_s \arrow[r] \arrow[d] \arrow[dr, phantom, "\scriptstyle\text{cart.}"] & \mathcal{X} \arrow[r] \arrow[d] & Y \\
\{s\} \arrow[r] & \Phi &
\end{tikzcd}
\]
On a donc une sous-catégorie fibrée de $\mathcal{X}$, sous réserve de vérifier que les flèches strictement cartésiennes
\[
x_0\longrightarrow\alpha_1(x_0), \qquad x_0\longrightarrow\alpha_2(x_0)
\]
dans $\mathcal{X}$, pour $x_0$ dans $X_{0\,y}$, sont dans $\mathcal{X}_y$, i.e. sont au-dessus de la flèche identique de $y$. Or c'est là la propriété essentielle
\msp{page 156}
du foncteur $\varphi_Q : \mathcal{X}\longrightarrow Y$, qu'on a oublié d'expliciter dans 3.2. Ça a vraiment l'air de marcher (comme sur des roulettes!)

Appliquons le Lemme 2 au diagramme carré $Q/x_0$ et à $\varphi_{(Q/x_0)}$ (page 154), on voit donc que les fibres à examiner en les $u\in\mathrm{Ob}\,Y/y$ ($u : y'\longrightarrow y$) peuvent s'écrire comme les intégrales
\[
\mathcal{X}(x_0,u) \quad\overset{\text{déf}}{=}\quad \intdiag{(X_0/x_0)_u}{(X_1/x_1)_u}{(X_2/x_2)_u} .
\]
Utilisons maintenant le fait que $f_1$, $f_2$ sont $W$-lisses (exprimant que $\varphi_Q$ est $W$-lisse dans $\mathcal{X}\backslash X_0$), donc que $(X_1/x_1)_u$ et $(X_2/x_2)_u$ sont $W$-asphériques. Dans le carré $W$-cocartésien
\[
\begin{tikzcd}[row sep=large,column sep=small]
 & (X_1/x_1)_u \arrow[dr, "{i_1(x_0,u)}"] & \\
(X_0/x_0)_u \arrow[ur, "{\alpha_1(x_0,u)}"] \arrow[dr, "{\alpha_2(x_0,u)}"'] & & (\mathcal{X}/x_0)_u\simeq\mathcal{X}(x_0,u) \\
 & (X_2/x_2)_u \arrow[ur, "{i_2(x_0,u)}"'] &
\end{tikzcd}
\]
% foot: Version 7 décembre 2002


\endgroup
% END UNIT Der-ChXIIb-012/Der-ChXIIb-012_p061-065.tex
% BEGIN UNIT Der-ChXIIb-012/Der-ChXIIb-012_p066-070.tex
\begingroup\makeatletter%
\newcommand{\printedpage}[1]{\marginpar{\footnotesize p.~#1}}%
\newcommand{\msp}[1]{\marginpar{\footnotesize [#1]}}%
\newcommand{\Ob}{\operatorname{Ob}}%
\newcommand{\Fl}{\operatorname{Fl}}%
\newcommand{\Cat}{\mathrm{Cat}}%
\newcommand{\Fibdisc}{\underline{\mathrm{Fib\ disc}}}%
\newcommand{\cercle}[1]{\tikz[baseline=(c.base)]\node[draw,circle,inner sep=1pt](c){#1};}%
\makeatother\let\printedpage\derpp\let\msp\dermsp%
\let\mspin\DERmspin%


\printedpage{108}
% head: A. Grothendieck : Dérivateurs, Ch. XII, Édition des Universités de Montpellier II et Paris VII
(\footnote[107]{attention, ce n'est pas un carré commutatif. Dire plutôt que $(\mathcal{X}/x_0)_u \simeq \int \begin{tikzcd}[ampersand replacement=\&, row sep=0.4em, column sep=0.4em, baseline=(\tikzcdmatrixname-2-1.base)] \& {} \\ {} \ar[ur]\ar[dr] \& \\ \& {} \end{tikzcd}$ est lisse sur $\Phi$ à fibres $W$-asphériques (si $f_0$, $f_1$, $f_2$ sont lisses), donc $W$-asphérique puisque $\Phi$ l'est.}), il faut exprimer que $(\mathcal{X}/x_0)_u$ est $W$-asphérique, sachant que $(X_1/x_1)_u$ et $(X_2/x_2)_u$ le sont. Donc cela équivaut à % Resolved from PDF leaf 31: the original diagram has exactly these two unlabelled arrows and no vertex labels.
\[
i_1(x_0,u) \in W \qquad (\text{ou } i_2(x_0,u) \in W),
\]
donc aussi (moyennant W(7 bis)) à
\msp{page 157}
\[
\alpha_1(x_0,u) \in W, \qquad \text{i.e. } (X_0/x_0)_u \ W\text{-asphérique.}
\]
Dire que c'est vrai pour tout $u$ dans $Y/y$, signifie que $f_0 : X_0 \longrightarrow Y$ est lisse en $u$ \texttt{[plutôt en $x_0$]}.

Donc on a prouvé a), sous la forme plus précise :

\textbf{Corollaire 6.} \textit{Soit $x \in \Ob \mathcal{X}$ (où $\mathcal{X} = \int \begin{tikzcd}[ampersand replacement=\&, row sep=tiny, column sep=small, baseline=(\tikzcdmatrixname-2-1.base)] \& X_1 \\ X_0 \ar[ur]\ar[dr] \& \\ \& X_2 \end{tikzcd}$.) On veut exprimer la $W$-lissité de}
\[
\varphi_Q : \mathcal{X} \longrightarrow Y
\]
\textit{en $x$.}

\begin{enumerate}[label={}, leftmargin=2.5em]
\item[\textit{a)}] \textit{Si $x \in \Ob X_i$ ($i \in \{1,2\}$), $\varphi_Q$ est $W$-lisse en $x$ si et seulement si $f_i : X_i \longrightarrow Y$ l'est.}
\item[\textit{b)}] \textit{Supposons} \textup{W(7 bis)} (\footnote[108]{Inutile pour la partie directe, en supposant \texttt{[phrase incomplète]}.})\textit{. Soit $x_0 \in \Ob X_0$, et supposons $f_1$, $f_2$ lisses en $x_1$, $x_2$. Alors $\varphi_Q$ est $W$-lisse en $x_0$ si et seulement si $f_0 : X_0 \longrightarrow Y$ l'est.}
\end{enumerate}

\textit{Donc} (moyennant W(7 bis)) \textit{dans les trois cas $x \in \Ob X_i$ ($i \in \{0,1,2\}$) on trouve que $\varphi_Q$ est $W$-lisse en $x$ si et seulement si $f_i : X_i \longrightarrow Y$ l'est} \texttt{[pas vraiment pour $i=0$]}.

\cercle{b} Quant à b) du théorème 4, c'est le lemme 2 ci-dessus. Donc tout est prouvé!

\msp{page 158}

\subsection*{3.5\quad Applications aux sommes amalgamées dans une catégorie $X^{\wedge}$ (avec $X \in \Ob \Cat$).}

D'abord une caractérisation des sommes amalgamées impliquant une immersion ouverte :

\textbf{Proposition 15.} \textit{Soit}
% foot: Version 7 décembre 2002
\printedpage{109}
% head: A. Grothendieck : Dérivateurs, Ch. XII, Édition des Universités de Montpellier II et Paris VII
\[
Q : \qquad
\begin{tikzcd}
U \ar[r, "i"] \ar[d, "p"'] & X \ar[d, "q"] \\
U' \ar[r, "i'"] & X'
\end{tikzcd}
\]
\textit{un carré commutatif dans} $\Cat$\textit{. Pour que ce carré soit cocartésien dans} $\Cat$\textit{, avec $i$ (donc $i'$) une immersion ouverte, il faut et il suffit que les conditions suivantes soient satisfaites :}

\begin{enumerate}[label={}, leftmargin=2.5em]
\item[\textit{a)}] \textit{$Q$ est cartésien.}
\item[\textit{b)}] \textit{$i'$ est une immersion ouverte (donc $i$ aussi).}
\item[\textit{c)}] \textit{Si $T' = X' \backslash U'$, et $T = X \times_{X'} T' = X \backslash U$, $T \longrightarrow T'$ est un isomorphisme.}
\item[\textit{d)}] \textit{Pour tout $t \in T$, si $t'$ est son image dans $T'$, le foncteur induit par $p$}
\[
p/t : \underbrace{U/t}_{= V_t} \longrightarrow \underbrace{U'/t'}_{= V'_{t'}}
\]
\textit{satisfait à $(p/t)_!(e_{U/t}) \simeq e_{U'/t'}$, ou ce qui revient au même, les `fibres colocales' de $p/t$, $v' \backslash V_t$, avec $v'$ dans $V'_{t'}$, sont} $0$\textit{-connexes.}
\end{enumerate}

\msp{page 159}
Démonstration immédiate - ça devrait plutôt remonter au §1.

\textbf{Corollaire 1.} \textit{Soit $Q$ un carré} cartésien (cf. page 158), \textit{avec $i'$ une} immersion ouverte \textit{(donc $i$ aussi), et $q$} $W$-lisse \textit{(donc $p$ aussi). On suppose que $W$ satisfait l'axiome} W(5) \textit{de connexité} (\footnote[109]{Donc il suffit $q$ $W_0$-lisse, où $W_0 = \{f \in \Fl \Cat \mid \pi_0(f) \text{ bijectif}\}$.})\textit{. Soit $T' = X' \backslash U'$, $T = X \backslash U$, supposons que $T' \longrightarrow T$ soit un isomorphisme. Alors $Q$ est cocartésien.}

Il suffit de prouver la condition d) de prop. 15. Considérons le carré
\[
\begin{tikzcd}
U/t \ar[r, hook, "i/t"] \ar[d, "p/t"'] & X/t \ar[d, "q/t"] \\
U'/t' \ar[r, hook, "i'/t'"] & X'/t',
\end{tikzcd}
\tag*{$(*)$}
\]
je dis qu'il est \textit{cartésien.} Cela résulte de la seule hypothèse que $Q$ soit cartésien, sans hypothèse de lissité sur $q$ ni que $i'$ soit une immersion ouverte, et se voit p. ex. en regardant le diagramme cubique :
% foot: Version 7 décembre 2002
\printedpage{110}
% head: A. Grothendieck : Dérivateurs, Ch. XII, Édition des Universités de Montpellier II et Paris VII
\[
\begin{tikzcd}[row sep=small, column sep=small]
& U/t \ar[rr] \ar[dl] \ar[dd] & & X/t \ar[dl] \ar[dd] \\
U'/t' \ar[rr, crossing over] \ar[dd] & & X'/t' & \\
& U \ar[rr] \ar[dl] & & X \ar[dl] \\
U' \ar[rr] & & X' \ar[from=uu, crossing over] &
\end{tikzcd}
\]
(\footnote[110]{NB Ici $t$ est un élément quelconque de $\Ob X$, $t'$ son image dans $X'$.}). Dans ce diagramme, la face antérieure et postérieure du cube est cartésienne, par définition de $U/t$ et de $U'/t'$
\msp{page 160}
respectivement, et la base l'est par hypothèse. Il en résulte que la face supérieure du cube est également un carré cartésien.

Ceci vu, commme $q$ est lisse, il en est de même de $q/t$ dans le carré cartésien $(*)$, donc aussi de $p/t$ par changement de base. D'autre part, la lissité de $q$ implique que $q/t$ est à fibres $W$-asphériques, donc par changement de base il en est de même de $p/t$ (qui est donc une $W$-équivalence universelle, i.e. une $W$-fibration triviale). De plus, à cause de la lissité de $p/t$, pour $v'$ dans $U'/t'$, l'inclusion % [sic] commme -> comme, repeated letter
\[
(U/t)_{v'} \longrightarrow v' \backslash (U/t)
\]
est dans $W$, donc (par W(5)) définit une bijection sur les $\pi_0$. Comme $(U/t)_{v'}$ est $W$-asphérique par ce qui précède, il s'ensuit (par W(5)) que son $\pi_0$ est de cardinal 1, donc il en est de même de $v' \backslash (U/t)$, i.e. celui-ci est 0-connexe. Or c'est là la condition d) de la proposition 15, qu'il fallait établir.

\msp{page 161}
\textbf{Corollaire 2.} \textit{Soit $Q$ un carré cartésien} (p. 158), \textit{avec $i'$ (donc aussi $i$) une immersion ouverte, et $q$ (donc aussi $p$) un} isomorphisme local. \textit{Soit $T = X \backslash U$, $T' = X' \backslash U'$ (donc $T \simeq X \times_{X'} T'$), et supposons $T' \longrightarrow T$ isomorphisme. Alors $Q$ est cocartésien.}

C'est un cas particulier du corollaire 1, mais sans référence explicite à un localiseur $W$. On peut pendre $W = W_{\infty}$, ou $W = W_0$, pour se ramener au corollaire 1. Mais il est plus raisonnable de dire que dans le carré cartésien $(*)$, l'hypothèse que $q$ soit un isomorphisme local implique que $q/t$ est un isomorphisme, donc $p/t$ l'est aussi. Il en résulte bien que les $v' \backslash (U/t)$ sont 0-connexes, cette catégorie étant isomorphe à $v' \backslash (U'/t')$, donc a un objet initial. % [sic] pendre -> prendre, missing letter

Soit maintenant $X$ un objet de Cat, et considérons un diagramme cocartésien dans $X^{\wedge}$ :
\[
\begin{tikzcd}
G \ar[r, "i"] \ar[d, "p"'] & F \ar[d, "q"] \\
G' \ar[r, "i'"] & F',
\end{tikzcd}
\tag*{$(*)$}
\]
% foot: Version 7 décembre 2002
\printedpage{111}
% head: A. Grothendieck : Dérivateurs, Ch. XII, Édition des Universités de Montpellier II et Paris VII
\msp{page 162}
où $i$ est un monomorphisme. On sait que le caractère cocartésien du carré se vérifie `argument par argument' par le caractère cocartésien de
\[
\begin{tikzcd}
G(x) \ar[r, "i(x)"] \ar[d, "p(x)"'] & F(x) \ar[d, "q(x)"] \\
G'(x) \ar[r, "i'(x)"] & F'(x),
\end{tikzcd}
\tag*{$(*)$}
\]
dans Ens, $\forall\, x$ dans $X$. Or un tel carré cocartésien dans Ens, avec $i(x)$ \texttt{[un]} monomorphisme, est aussi cartésien. Cela implique que le carré $Q$ dans $X^{\wedge}$ est également cartésien.

Considérons, d'autre part, le foncteur canonique
\[
\varphi_X : F \longmapsto X/F : X^{\wedge} \longrightarrow \mathrm{Cat}.
\]
Il vaut mieux, d'ailleurs, le regarder comme un foncteur
\[
\varphi_X : X^{\wedge} \longrightarrow \mathrm{Cat}/X,
\]
et je rappelle, pour mémoire, que ce foncteur est pleinement fidèle, et que son image essentielle est formée des catégories $X'$ sur $X$ telles que
\[
X' \longrightarrow X
\]
\msp{page 163}
soit un isomorphisme local, ou ce qui revient au même, que $X'$ soit \textit{fibrée} sur $X$ à fibres discrètes :
\[
\varphi_X : X^{\wedge} \xrightarrow{\ \sim\ } \underbrace{\Fibdisc\, X}_{\substack{\text{catégories fibrées}\\ \text{à fibres discrètes}}} \ .
\]
Notons que toute flèche $X' \longrightarrow X''$ dans $\Fibdisc\, X$ est un isomorphisme local (tautologique sur $\begin{tikzcd}[ampersand replacement=\&, row sep=small, column sep=tiny, baseline=(\tikzcdmatrixname-1-1.base)] X'/x' \ar[rr] \ar[dr, "\text{iso.}"'] \& \& X''/x'' \ar[dl, "\text{iso.}"] \\ \& X/x \& \end{tikzcd}$), ce qui montre que si on a un diagramme
\[
\begin{tikzcd}[row sep=small, column sep=small]
X_1 \ar[dr] & & X_2 \ar[dl] \\
& X_0 &
\end{tikzcd}
\]
dans $\Fibdisc\, X$, le produit fibré $X_3 = X_1 \times_{X_0} X_2$ dans Cat est tel que $X_3 \longrightarrow X_2$ est un isomorphisme local, donc aussi $X_3 \longrightarrow X$. C'est-à-dire $\Fibdisc(X)$ est stable par produits fibrés. Ceci implique que les produits fibrés existent dans $\Fibdisc\, X$, et que ce sont les produits fibrés dans Cat, ou dans Cat$/X$.
% foot: Version 7 décembre 2002

\printedpage{112}
% head: A. Grothendieck : Dérivateurs, Ch. XII, Édition des Universités de Montpellier II et Paris VII
\msp{page 164}
De même pour les produits finis dans $\Fibdisc\, X$, ce sont les produits finis dans Cat$/X$. Bref

\textbf{Lemme 1.} \textit{La sous-catégorie pleine} $\Fibdisc\, X$ \textit{de} $\Cat/X$ \textit{est stable par} $\varprojlim$ \textit{finies, donc les} $\varprojlim$ \textit{existent dans} $\Fibdisc\, X$ \textit{et le foncteur d'inclusion}
\[
\Fibdisc\, X \longrightarrow \Cat/X
\]
\textit{y commute, et en particulier transforme carrés cartésiens de} $\Fibdisc\, X$ \textit{en carrés cartésiens de} $\Cat/X$\textit{, ou ce qui revient au même, de} $\Cat$\textit{. De plus, en vertu de l'équivalence}
\[
\varphi_X : X^{\wedge} \xrightarrow{\ \sim\ } \Fibdisc\, X \hookrightarrow \Cat/X,
\]
\textit{le foncteur $\varphi_X$ de $X^{\wedge}$ dans} $\Cat/X$ \textit{est exacte à gauche, en particulier il transforme carrés cartésiens de $X^{\wedge}$ en carrés cartésiens de} $\Cat/X$\textit{, ou encore de} $\Cat$\textit{.} % [sic] exacte -> exact, agreement with "le foncteur"

\msp{page 165}
Notons aussi qu'étant exact à gauche, le foncteur $\varphi_X$ de $X^{\wedge}$ dans Cat$/X$ transforme monomorphismes en monomorphismes. D'autre part :

\textbf{Lemme 2.} \textit{Les monomorphismes qui sont des isomorphismes locaux sont exactement les immersions ouvertes. Donc dans} $\Fibdisc\, X$\textit{, les monomorphismes sont les immersions ouvertes, et le foncteur $\varphi_X : F \longmapsto X/F$ transforme monomorphismes en immersions ouvertes.}

Revenons alors au diagramme cocartésien $(*)$ (p. 161), et au diagramme correspondant dans Cat (et même dans $\Fibdisc\, X$)
\[
\begin{tikzcd}
X/G \ar[r, hook, "X/i"] \ar[d, "X/p"'] & X/F \ar[d, "X/q"] & T \ar[l, hook', dotted] \ar[d, dotted] \\
X/G' \ar[r, hook, "X/i'"] & X/F' & T'. \ar[l, hook', dotted]
\end{tikzcd}
\]
Comme $(*)$ est cartésien, le carré précédent est cartésien, par ce qui précède. Je veux prouver qu'il est également cocartésien
\msp{page 166}
(dans Cat$/X$, ou dans Cat, cela revient au même), et pour ceci j'applique le corollaire 2 de la proposition 15 (p. 161), compte tenu que $X/q$ est un isomorphisme local, et $X/i'$ (étant un isomorphisme local \texttt{[et un]} monomorphisme) une immersion ouverte. Il suffit de prouver que, $T'$ et $T$ étant les sous-catégories fermées complémentaires des ouverts $X/G'$ de $X/F'$ et de $X/G$ de $X/F$, que $T \longrightarrow T'$ est \texttt{[un]} isomorphisme.

{\small \texttt{[Maintenant, $F$ et $G$, et $F'$ et $G'$, mais pas $p$ et $q$, échangent leur rôles.]}} % [sic] leur rôles -> leurs rôles, agreement

La vérification est immédiate : $T$ est la catégorie des flèches $x \longrightarrow G$ dans $X^{\wedge}$ (de source $x$ dans $X$) qui ne se factorisent pas par $G$ \texttt{[plutôt par $F$]}, $T'$ celle des flèches $x \longrightarrow G'$
% foot: Version 7 décembre 2002


\endgroup
% END UNIT Der-ChXIIb-012/Der-ChXIIb-012_p066-070.tex
% BEGIN UNIT Der-ChXIIb-012/Der-ChXIIb-012_p071-072.tex
\begingroup\makeatletter%
\newcommand{\printedpage}[1]{\marginpar{\footnotesize p.~#1}}%
\newcommand{\msp}[1]{\marginpar{\footnotesize [#1]}}%
\makeatother\let\printedpage\derpp\let\msp\dermsp%
\let\mspin\DERmspin%


\printedpage{113}
% head: A. Grothendieck : Dérivateurs, Ch. XII, Édition des Universités de Montpellier II et Paris VII
qui ne se factorisent pas par $F'$, le foncteur $T \longrightarrow T'$ étant le foncteur composition de $x \longrightarrow G$ avec $G \xrightarrow{q} G'$
\[
\begin{tikzcd}[row sep=large]
F \arrow[r, hook] \arrow[d] & G \arrow[d] \\
F' \arrow[r, hook] & G'
\end{tikzcd}
\]
(NB Le composé ne se factorise pas par $F'$, car autrement, à cause du caractère cartésien de $Q$, $x \longrightarrow G$ se factoriserait par $F$.) Mais pour $x$ donné, la donnée de $x \longrightarrow G$ ne se factorisant pas par $F$ est celle d’un élément \msp{page 167} de $G(x)\backslash F(x)$, et à cause du caractère cocartésien de $Q$,
\[
(G(x)\backslash F(x)) \longrightarrow (G'(x)\backslash F'(x))
\]
est bijectif, i.e. ‘il revient au même’ de se donner $x \longrightarrow G$ ne se factorisant pas par $F$, ou $x \longrightarrow G'$ ne se factorisant pas par $F'$.

Ainsi on a prouvé le

\textbf{Théorème 5.} \textit{Soit $X$ dans} Cat\textit{, et}
\[
Q : \qquad
\begin{tikzcd}[row sep=large, baseline=(current bounding box.center)]
F \arrow[r, hook, "i"] \arrow[d, "p"'] & G \arrow[d, "q"] \\
F' \arrow[r, hook, "i'"] & G'
\end{tikzcd}
\]
\textit{un carré cocartésien dans $X^\wedge$, avec $i$} \texttt{[un]} \textit{monomorphisme. Alors $i'$ est} \texttt{[un]} \textit{monomorphisme, le carré $Q$ est aussi cartésien, et le carré correspondant dans} Cat$/X$ \textit{ou dans} Cat
\[
\begin{tikzcd}[row sep=large]
X/F \arrow[r, hook, "X/i"] \arrow[d, "X/p"'] & X/G \arrow[d, "X/q"] \\
X/F' \arrow[r, hook, "X/i'"] & X/G'
\end{tikzcd}
\]
\msp{page 168} \textit{est à la fois cartésien et cocartésien. De plus, si $W$ est un localiseur fondamental dans} Cat\textit{, satisfaisant} W(7 bis) (en plus de W(1,2,3))\textit{, alors le carré $X/Q$ dans} Cat \textit{est $W$-cocartésien.}

\printedpage{114}
% head: A. Grothendieck : Dérivateurs, Ch. XII, Édition des Universités de Montpellier II et Paris VII
Ce dernier énoncé résulte de ce qui précédait et du corollaire 2 du théorème 4, compte tenu que $X/q$ est un isomorphisme local, donc est $W$-lisse.

\textbf{Corollaire :} (\footnote[111]{\textbf{NB} Si on veut seulement $i \in W_X \Longrightarrow i' \in W_X$, $p \in W_X \Longrightarrow q \in W_X$, il suffit que l’on ait W(7 bis) \textit{faible}, i.e. $i_0 \in W \Longrightarrow i \in W$ dans l’énoncé de W(7 bis), p. 76 (Loc(6$'$a) avec les notations de XVI).}) \textit{On a} (moyennant W(7 bis)) \textit{les équivalences}
\[
\begin{aligned}
i \in W_X \quad &\Longleftrightarrow \quad i' \in W_X \\
p \in W_X \quad &\Longleftrightarrow \quad q \in W_X,
\end{aligned}
\]
\textit{où}
\[
W_X \subseteq \mathit{Fl}(X^\wedge)
\]
\textit{est défini par (cf. Pursuing Stacks \ldots)}
\[
u \in W_X \qquad\qquad \Longleftrightarrow \qquad\qquad X/u : X/F \longrightarrow X/G \quad \textit{est dans } W.
\]

Ainsi, dans une catégorie $X^\wedge$, les sommes amalgamées impliquant un monomorphisme se comportent de façon parfaite du point de vue homotopique, plus exactement pour les localiseurs $W_X$, pourvu \msp{page 169} seulement que $W$ satisfasse à W(7 bis). Notamment, les monomorphismes dans $X^\wedge$ sont des ‘$W_X$-cofibrations’, en un sens similaire de celui de la définition 2 c) (pages 99, 100 \texttt{[plutôt 109, 110]}).

\texttt{[page 170 n’existe pas.]}


\endgroup
% END UNIT Der-ChXIIb-012/Der-ChXIIb-012_p071-072.tex
\clearpage\setcounter{footnote}{0}\phantomsection\addcontentsline{toc}{section}{Chapitre XII, 6}\markboth{Chapitre XII, 6}{Chapitre XII, 6}
% BEGIN UNIT Der-ChXIIc-30-03-03/Der-ChXIIc-30-03-03_p001-005.tex
\begingroup\makeatletter%
\newcommand{\printedpage}[1]{\marginpar{\footnotesize p.~#1}}%
\newcommand{\msp}[1]{\marginpar{\footnotesize [#1]}}%
\newcommand{\Ch}{\underline{\mathrm{Ch}}}%
\newcommand{\Hotp}[1]{\mathrm{Hot}_{\bullet\,#1}}%
\newcommand{\hotp}[1]{\mathrm{hot}_{\bullet\,#1}}%
\newcommand{\Ob}{\mathrm{Ob}\,}%
\newcommand{\HomEns}{\mathrm{Hom}_{\mathrm{Ens}}}%
\newcommand{\utpi}{\underset{\sim}{\pi}}%
\makeatother\let\printedpage\derpp\let\msp\dermsp%
\let\mspin\DERmspin%


\printedpage{115}
% head: A. Grothendieck : Dérivateurs, Ch. XII, Édition des Universités de Montpellier II et Paris VII
% foot: Version 30 mars 2003

\msp{page 171}

\section*{VI\quad Critère de $W_\infty$-asphéricité pour les $\Omega^n(X,x)$.}

\subsection*{1\quad Cogitations sur $W_n$.}

On va définir, pour $n \in \mathbf{N}$,
\[
W_n \subseteq \mathrm{Fl}(\mathrm{Cat})
\]
par les conditions suivantes (où $f : X \longrightarrow Y$ est une flèche dans Cat).

\noindent\begin{tabular}{@{}l c l >{\RaggedRight\arraybackslash}p{0.64\textwidth}@{}}
$f \in W_0$ & $\Longleftrightarrow$ & a) & $\pi_0(f)$ bijectif.\\
 & $\Longleftrightarrow$ & b) & Pour tout préfaisceau d'ensembles $F$ \textit{constant} sur $Y$, $H^0(Y,F) \longrightarrow H^0(X,f^*F)$ est bijectif.\\
 & $\Longleftrightarrow$ & c) & Le foncteur $F \longmapsto f^*F$ des préfaisceaux d'ensembles \textit{constants} sur $Y$ vers les faisceaux d'ensembles constants sur $X$ est pleinement fidèle.
\end{tabular}

\texttt{[Les termes de faisceau ou de préfaisceau sont utilisés ici indifféremment.]}

\textbf{NB} Si $F_0$ est un ensemble, $F_{0\,Y}$ le faisceau d'ensembles constant défini par $F$ sur $Y$, on a
\[
\begin{aligned}
H^0(Y,F_{0\,Y}) &\simeq \HomEns(\pi_0(Y),F_0) && \text{(fonctoriel en $Y$)}\\
\mathrm{Hom}_{Y^\wedge}(F_{0\,Y},G_{0\,Y}) &\simeq \HomEns(\pi_0(Y),\HomEns(F_0,G_0)) && \text{(fonctoriel en $Y$),}
\end{aligned}
\]
ce qui implique que a) $\Longrightarrow$ c) $\underset{\substack{\text{en prenant}\\ \mathrm{Hom}_{Y^\wedge}(e_Y,F)}}{\overset{\text{tautologiquement}}{\Longrightarrow}}$ b) $\Longrightarrow$ a).

\noindent\begin{tabular}{@{}l c l >{\RaggedRight\arraybackslash}p{0.64\textwidth}@{}}
$f \in W_1$ & $\Longleftrightarrow$ & a) & $\pi_0(f)$ bijectif, et pour tout $x \in \Ob X$, $\pi_1(f,x) : \pi_1(X,x) \longrightarrow \pi_1(Y,f(x))$ est bijectif. (NB Il suffit de le vérifier pour \textit{un} $x$ dans chaque composante connexe de $X$.)\\
 & $\Longleftrightarrow$ & a$'$) & $\Pi(X) \xrightarrow{\Pi f} \Pi(Y)$ est une équivalence de catégories.\\
 & $\Longleftrightarrow$ & b) & $H^0(Y,F) \longrightarrow H^0(X,f^*F)$ bijectif pour tout faisceau d'ensembles \textit{localement constant} sur $Y$, et $H^1(Y,F) \longrightarrow H^1(X,f^*F)$ bijectif si de plus $F$ est un faisceau en groupes.\\
 & $\Longleftrightarrow$ & b$'$) & Comme b), mais pour le $H^1$ on se borne à $F$ \textit{constant}, et on exige $H^1(f)$ \textit{surjectif} seulement.\\
 & $\Longleftrightarrow$ & c) & Le foncteur $F \longmapsto f^*F$ de la catégorie des faisceaux d'ensembles localement constants sur $Y$ vers celle des faisceaux localement constants sur $X$ est une équivalence de catégories.
\end{tabular}

\par

\printedpage{116}
% head: A. Grothendieck : Dérivateurs, Ch. XII, Édition des Universités de Montpellier II et Paris VII
% foot: Version 30 mars 2003

\msp{page 172}

\textbf{NB} Le groupe $\pi_1(X,x)$ est défini via la théorie du groupoïde fondamental. On sait qu'il est en bijection canonique avec $\pi_0(\Omega(X,x))$, où $\Omega(X,x) = \Ch(X;x,x)$ (ce qui permet d'avoir une loi de composition monoïdale sur $\Omega(X,x)$, donnant la loi de groupe sur le $\pi_0$), ou au choix $\Omega(X,x) = \Ch_\infty(X;x,x)$ (plus commode dans certaines situations, mais ne donnant pas directement la loi de composition dans $\pi_0(\Omega(X,x))$).

Je vais poser, pour $n \in \mathbf{N}$,
\[
\pi_n(X,x) = \pi_0(\Omega^n(X,x)),
\]
$\Omega^n$ désignant l'itéré $n$-ième du foncteur $(X,x) \longmapsto \Omega(X,x)$ de la catégorie $\mathrm{Cat}_\bullet$ des objets ponctués de Cat dans elle-même. Je préfère ici travailler avec la théorie \underline{Ch} au lieu de $\Ch_\infty$. Tôt ou tard il faudra donner un énoncé de comparaison \dots

\textbf{Théorème 1} (\footnote[112]{Sans doute \textit{faux}. On doit avoir a) $\Longrightarrow$ b), mais b) $\Longrightarrow$ a) \texttt{[est]} sans doute faux déjà pour $n = 2$ (sauf le cas $X$, $Y$ 1-connexes). Cf. énoncé rectifié \texttt{[p. 178]}.}). \textit{Soit $f : X \longrightarrow Y$ dans} Cat\textit{, et soit $n \in \mathbf{N}$, $n \geq 2$, $S$ une partie de} $\Ob X$ \textit{telle que le composé $S \longrightarrow \Ob X \longrightarrow \pi_0(X)$ soit} \texttt{[un]} \textit{épimorphisme. Les conditions suivantes sont équivalentes.}

\begin{itemize}
\item[\textit{a)}] \textit{$\pi_i(f,x)$ isomorphe pour tout $x \in \Ob X$ et tout $i \in [0,n]$.}

\item[\textit{a$'$)}] \textit{Comme a), mais avec $x \in S$.}

\item[\textit{b)}] \textit{On a}
\[
\begin{array}{rcll}
H^0(Y,F) & \xrightarrow{\ \sim\ } & H^0(X,f^*F) & \textit{pour $F$ faisceau d'ensembles}\\
 & & & \textit{localement constant sur $Y$.}\\
H^1(Y,F) & \xrightarrow{\ \sim\ } & H^1(X,f^*F) & \textit{si de plus $F$ faisceau en groupes.}\\
H^i(Y,F) & \xrightarrow{\ \sim\ } & H^i(X,f^*F) & \textit{si de plus $F$ faisceau abélien, $i \in [2,n]$.}
\end{array}
\]
\end{itemize}

\msp{page 173}

\textbf{NB} L'équivalence de a) et a$'$) devrait être triviale, mais pour cela il faut avoir défini les $\pi_n(X,x)$, pour $x$ variable dans $\Ob X$, comme système local sur $X$, chose que je n'ai pas faite encore. Je vais essayer de le faire dans un courant d'air. Tout d'abord, les $\Omega(X,x)$ sont les fibres d'une $W_\omega$-fibration parfaite (\footnote[113]{\textbf{NB} $W_\omega$ est le plus petit localiseur de Cat (satisfaisant W(1, 2 bis, 3)).}), s'insérant dans le carré cartésien
\[
\begin{tikzcd}
\Ch(X) \arrow[d, "{p_X=(\sigma_X,\beta_X)}"'] & \Omega(X,-) \arrow[l] \arrow[d] \\
X \times X & X. \arrow[l, "\mathrm{diag}_X"'] \arrow[ul, dotted, "i_X"']
\end{tikzcd}
\]

\par

\printedpage{117}
% head: A. Grothendieck : Dérivateurs, Ch. XII, Édition des Universités de Montpellier II et Paris VII
% foot: Version 30 mars 2003

On sait alors que l'on a un foncteur
\[
\begin{aligned}
\Pi X &\longrightarrow \mathrm{Hot}_{W_\omega}\\
x &\longmapsto \mathrm{hot}_{W_\omega}(\Omega(X,x)).
\end{aligned}
\tag*{$(*)$}
\]
Or on a, pour $i \geq 1$,
\[
\pi_i(X,x) = \pi_{i-1}(\ \underbrace{\Omega(X,x)}_{\text{objet ponctué}}\ ).
\]
Cela donne le système local cherché des $\pi_i$, si seulement on avait un foncteur un peu plus précis que $(*)$, savoir
\[
\Pi X \longrightarrow \Hotp{W_\omega} \qquad \text{(catégorie $W_\omega$-homotopique \textit{ponctuée})}
\]
(\footnote[114]{Mais il faut aussi vérifier que les $\pi_n(X,x)$ ($n \geq 1$) sont des foncteurs sur $\Hotp{W_\infty}$, i.e. transforment $W_\infty$-équivalences en itou. Cela résulte de W(4a) pour $W_\infty$, qui implique que $(X,x) \mapsto \Omega(X,x)$ se factorise par $\Hotp{W_\infty}$ (en tant que foncteur à valeurs dans $\Hotp{W_\infty}$).}). Or la flèche $i_X : X \longrightarrow \Ch X$ relevant $p_X$ donne une section canonique de $\Omega(X,-)$ sur $X$ -
\msp{page 174}
celle justement qui associe à tout $x \in X$ le point marqué $1_x$ de $\Omega(X;x,x)$ (`lacet vide en $x$'). La question est donc celle-ci :

\begin{quote}
On a une $W$-fibration $\Omega \longrightarrow X$, avec une section donné $\sigma$. Cela définit-il un foncteur canonique % [sic] donné -> donnée, accord avec « section »
\[
X \longrightarrow \Hotp{W},
\]
transformant toute flèche en \texttt{[un]} isomorphisme?
\end{quote}

L'idée naturelle est de prendre
\[
x \longmapsto \hotp{W}(\Omega/x)
\]
en ponctuant $\Omega/x$ par le point marqué $\sigma(x)$ de $\Omega_x \hookrightarrow \Omega/x$. L'ennui, c'est que les morphismes de transition $\Omega/x \longrightarrow \Omega/x'$ (pour $x \xrightarrow{u} x'$ dans $X$) ne sont pas compatibles avec les points marqués. Mais on a une flèche (savoir $\sigma(u)$) allant
\msp{page 175}
de l'image $\sigma(x)/x'$ du point marqué de $\Omega/x$ vers le point marqué $\sigma(x')/x'$ de $\Omega/x'$. Ceci doit suffire pour définir la flèche
\[
(\Omega/x,\sigma(x)/x) \longrightarrow (\Omega/x',\sigma(x')/x')
\]
comme flèche dans $\Hotp{W}$, avec transitivité pour une composée $x \xrightarrow{u} x' \xrightarrow{v} x''$ dans $X$. Je vais pour l'instant me borner
\printedpage{118}
% head: A. Grothendieck : Dérivateurs, Ch. XII, Édition des Universités de Montpellier II et Paris VII
% foot: Version 30 mars 2003
\msp{page 176}
à une indication, quitte à faire des vérifications circonstanciées par la suite. Au besoin, j'utiliserai dès maintenant que pour $i \in \mathbf{N}$, $i \geq 1$,
\[
\pi_i(f,x) : \pi_i(X,x) \longrightarrow \pi_i(Y,f(x))
\]
peuvent s'interpréter comme provenant d'un homomorphisme de systèmes locaux
\[
\utpi_i(X,-) \longrightarrow f^*(\utpi_i(Y,-))
\]
sur $X$. La condition a) du théorème 1 s'interprète en disant que cet homomorphisme est un \textit{isomorphisme} pour $i \in [1,n]$, et la condition a$'$) se borne alors à rappeler qu'il suffit de tester qu'il en est ainsi pour \textit{un} point de chaque composante connexe de $X$.

L'énoncé du théorème 1 garde un sens pour $n = 1$, et dans ce cas a été `rappelé' précédemment (et est $\pm$ tautologique). Cela va nous permettre de procéder par récurrence sur $n$, en supposant que le théorème est prouvé pour $n-1$. Cela signifie que l'on peut supposer que les conditions a) \textit{et} b) sont satisfaites par $f$ pour $i \in [1,n-1]$, et se borner à prouver que, sous ces conditions, la condition a) (ou a$'$), pour $i = n$, équivaut à la condition b) pour $i = n$.

\msp{page 177}

Donc il faut prouver simplement ceci :

\textbf{Lemme} (\footnote[115]{Sans doute faux.}). \textit{Supposons que $\pi_i(f,x)$} \texttt{[est un]} \textit{isomorphisme pour $i \in [0,n-1]$ et b) satisfait pour $i \in [0,n-1]$, où $n \in \mathbf{N}$, $n \geq 2$, est fixé. Alors les conditions suivantes sont équivalentes :}

\begin{itemize}
\item[\textit{a)}] $\pi_n(f,x) : \pi_n(X,x) \longrightarrow \pi_n(Y,f(x))$ \texttt{[est un]} \textit{isomorphisme pour tout $x \in \Ob X$.}

\item[\textit{b)}] $H^n(Y,F) \xrightarrow{\sim} H^n(X,f^*F)$ \texttt{[est un]} \textit{isomorphisme pour tout faisceau abélien localement constant sur $Y$.}
\end{itemize}

Notons que $f$ induit une flèche dans $\mathrm{Cat}_\bullet$,
\[
\Omega(f,x) : \Omega(X,x) \longrightarrow \Omega(Y,f(x)),
\]
d'où
\[
\begin{tikzcd}[row sep=small]
\pi_i(\Omega(X,x)) \arrow[r] \arrow[d, dash, "\wr"'] & \pi_i(\Omega(Y,f(x))) \arrow[d, dash, "\wr"'] \\
\pi_{i+1}(X,x) \arrow[r] & \pi_{i+1}(Y,f(x)),
\end{tikzcd}
\]
s'identifiant à $\pi_{i+1}(f,x)$ :
\[
\pi_i(\Omega(f,x)) \simeq \pi_{i+1}(f,x).
\]


\printedpage{119}
% head: A. Grothendieck : Dérivateurs, Ch. XII, Édition des Universités de Montpellier II et Paris VII
% foot: Version 30 mars 2003

L'hypothèse a) du lemme implique que $\pi_i(\Omega(X,x)) \xrightarrow{\sim} \pi_i(\Omega(Y,y))$ pour $i \in [0,n-1]$. Il s'agit ici des $\pi_i$ basés en les points unités. On devra sans doute prouver que, du fait qu'il s'agit d'un morphisme de monoïdes dans Cat (satisfaisant certaines conditions?) ça reste vrai pour les $\pi_i$ en \textit{tout} point $\xi$ de $\Omega(X,n)$. Par l'hypothèse de récurrence, on saura donc que $\Omega(X,x) \longrightarrow \Omega(Y,y)$ % [sic] \Omega(X,n) -> \Omega(X,x), le point de base est x
\msp{page 178}
induit un isomorphisme pour la cohomologie à coefficients localement constants, jusqu'en dimension $n-1$ inclus.

Finalement, je me suis $\pm$ convaincu que l'énoncé du théorème 1 est faux tel quel. Voici un énoncé que j'éspère vrai, et que je vais essayer de prouver dans un esprit `utilitaire', en utilisant tout ce que je sais en homotopie et cohomologie (et pas seulement ce que j'ai développé `from scratch' dans le modélisateur Cat). % [sic] j'éspère -> j'espère, orthographe

\textbf{Théorème 1 rectifié.} \textit{Considérons, pour un entier $n \geq 2$ donné, les conditions suivantes sur une flèche $f : X \longrightarrow Y$ dans} Cat \textit{(ou une application continue d'espaces topologiques, ou un morphisme dans $\boldsymbol{\Delta}^\wedge$ etc. etc.) :}

\begin{itemize}
\item[$a_n$)] \textit{Si $Z$ est une `fibre homotopique' de $f$ (NB il y en a essentiellement une par composante connexe de $Y$), alors $Z$ est $n$-connexe, i.e.}
\msp{page 179}
\textit{(p. ex) $\pi_i(Z) = 0$ pour $i \in [0,n]$.}

\item[$b_n$)] \textit{$\pi_0(f)$ est bijectif, et $\forall\ x \in X$, $\pi_i(f,x) : \pi_i(X,x) \longrightarrow \pi_i(Y,f(x))$ est} \texttt{[un]} \textit{isomorphisme pour $i \in [1,n]$. (}\textbf{NB} \textit{Il suffit de le dire pour} \textup{un} \textit{point $x$ de chaque composante connexe de $X$.)}

\item[$c_n$)] $H^i(Y,F) \longrightarrow H^i(X,f^*F)$ \texttt{[est un]} \textit{isomorphisme pour $F$ faisceau localement constant sur $Y$ (avec $F$ abélien si $i \geq 2$, $F$ faisceau en groupes si $i = 1$).}
\end{itemize}

\textit{On a alors}
\[
a_n \ \Longrightarrow\ b_n \ \Longrightarrow \text{(\footnotemark[116])}\ c_n \ \Longrightarrow \text{(\footnotemark[117])}\ a_{n-1}
\]
\footnotetext[116]{si $X$, $Y$ 1-connexes, c'est vraiment prouvé, sinon, il y a des vérifications à faire (qui ont l'air soritales \dots).}%
\footnotetext[117]{\textbf{NB} Même si $X$, $Y$ 1-connexes, il n'est pas clair que $\mathrm{c}_n \Longrightarrow \mathrm{b}_n$, car il faut montrer que $\pi_n(Z) \longrightarrow \pi_n(X)$ est nul, et on trouve seulement que le composé avec $\pi_n(X) \longrightarrow H_n(X)$ est nul. Je présume qu'il y a des contre-exemples.}%
\textit{(et je présume que chacune de ces implications est stricte).}

\textsc{Démonstration}. $\mathrm{a}_n \Longrightarrow \mathrm{b}_n$ est conséquence immédiate de la suite exacte d'homotopie pour une Hot-fibration $\bar{X} \xrightarrow{\bar{f}} Y$ qui factorise $X \xrightarrow{f} Y$ en $f = \bar{f} \circ i$, $i \in W_\infty$. En fait, $\mathrm{a}_n$ \textit{équivaut} à $\mathrm{b}_n$ + la surjectivité des $\pi_{n+1}(X,x) \longrightarrow \pi_{n+1}(Y,f(x))$, laquelle bien sûr ne résulte \textit{pas} de $\mathrm{b}_n$, i.e. de la bijectivité des $\pi_i(f,x)$ pour $i \leq n$ (prendre des $K(\pi,n+1)$). Donc l'implication est bien stricte.


\endgroup
% END UNIT Der-ChXIIc-30-03-03/Der-ChXIIc-30-03-03_p001-005.tex
% BEGIN UNIT Der-ChXIIc-30-03-03/Der-ChXIIc-30-03-03_p006-010.tex
\begingroup\makeatletter%
\NoAutoSpacing%
\newcommand{\printedpage}[1]{\marginpar{\footnotesize p.~#1}}%
\newcommand{\msp}[1]{\marginpar{\footnotesize [#1]}}%
\makeatother\let\printedpage\derpp\let\msp\dermsp%
\let\mspin\DERmspin%


\printedpage{120}
% head: A. Grothendieck : Dérivateurs, Ch. XII, Édition des Universités de Montpellier II et Paris VII

\msp{page 180}
$\mathrm{b}_n \Longrightarrow \mathrm{c}_n$. On peut supposer que $f$ est une Hot-fibration (et même Hot-parfait si ça nous chante \dots). On sait déjà que pour
\[
H^i(Y,F) \longrightarrow H^i(X,f^*F) \tag*{$(*)$}
\]
tout est O.K. si $i \in \{0,1\}$. Donc on peut supposer $F$ abélien. Si ça nous avance, on peut supposer, par récurrence, (que l’implication $\mathrm{b}_{n-1} \Longrightarrow \mathrm{c}_{n-1}$ est connue, donc) que $(*)$ est \texttt{[un]} isomorphisme pour $i<n$. Reste à le voir pour $i=n$. On est ramené bien sûr, comme $\pi_0(f)$ est bijectif, au cas où $X$, $Y$ sont 0-connexes, donc il y a essentiellement une seule fibre $Z$. La suite exacte d’homotopie nous apprend que
\[
\pi_i(Z,z)=0 \qquad \text{pour } i \le n-1 \ (Z \text{ est } (n-1)\text{-connexe})
\]
et de plus
\[
\pi_n(Z,z) \longrightarrow \pi_n(X,z) \qquad \text{est nul}
\]
(car $\pi_n(X,z) \longrightarrow \pi_n(Y,y)$ \texttt{[est un]} isomorphisme donc \texttt{[un]} monomorphisme), donc
\[
\pi_{n+1}(Y,y) \longrightarrow \pi_n(Z,z) \qquad \text{\texttt{[est un]} épimorphisme.}
\]
Comme $Z$ est $(n-1)$-connexe (donc aussi 1-connexe), on a
\[
H^i(Z,G)=0\ , \qquad \text{si } 0<i<n\ , \tag*{$(*)$}
\]
pour $G$ localement constant (donc constant) sur $Z$.

\msp{page 181}
Écrivons la suite spectrale de Leray pour $f : X \longrightarrow Y$ et $f^*F = F_X$ :
\[
H^*(X,f^*F) \Longleftarrow E_2^{pq} = H^p(Y,H^q(Z,f^*F)),
\]
où ‘$H^q(Z)$’ doit être vu comme un système local sur $Y$. Donc par $(*)$ on a
\[
E_2^{pq}=0 \qquad \text{si } 0<q<n, \qquad
\begin{tikzpicture}[baseline=0pt, x=0.6cm, y=0.6cm]
\draw[->] (0,0) -- (0,4);
\draw[->] (0,0) -- (6,0);
\node[left] at (0,3) {$n$};
\foreach \i in {0,...,4} { \fill (\i,3) circle (2.5pt); \fill (\i,0) circle (2.5pt); }
\node at (1,1.5) {$0$};
\node at (3,1.5) {$0$};
\end{tikzpicture}
\qquad ,
\]
donc on trouve la suite exacte
\[
0 \longrightarrow H^n(Y,F) \longrightarrow H^n(X,F_X) \xrightarrow{\ \alpha\ } H^0(Y,\ \underbrace{H^n(Z,\overbrace{F_Z}^{M_Z})}_{\begin{array}{l}\scriptstyle = \ \mathrm{Hom}(\pi_n(Z,z),M)\\ \scriptstyle \text{par \texttt{[le]} th. de Hurewicz}\end{array}}\ ) \longrightarrow H^{n+1}(Y,F).
\]

\printedpage{121}
% head: A. Grothendieck : Dérivateurs, Ch. XII, Édition des Universités de Montpellier II et Paris VII

Ici je fixe un point $x$ de $X$, $Z$ désigne la fibre de $f$ en $y=f(x)$, et $M$ celle de $F$ en $y$, $M=F_y$. Le groupe $\pi_1(Y)$ y opère, d’ailleurs $\pi \overset{\text{déf}}{=} \pi_1(X,x) \simeq \pi_1(Y,y)$ aussi. Le $H^0(Y,H^n(Z,F_Z))$ peut s’interpréter, sauf erreur, comme
\[
\mathrm{Hom}_\pi(\pi_n(Z,x),M),
\]
en tenant compte de l’opération de $\pi=\pi_1(X,x)$ sur $\pi_n(Z,x)$ également. (À vérifier.) Il faut interpréter la flèche
\[
H^n(X,F_X) \longrightarrow \mathrm{Hom}_\pi(\pi_n(Z,x),M),
\]
via l’application composée \texttt{[où $z=x$, vu comme point de $Z$]}
\[
\pi_n(Z,z) \longrightarrow \pi_n(X,x) \longrightarrow H_n(X,\mathbf{Z}). \tag*{$(*)$}
\]
Si $F$ est \textit{constant} sur $Y$ (p. ex. si $Y$ \texttt{[est]} 1-connexe), donc de valeur $M$, on a les flèches canoniques
\[
H^n(X,F_X) = H^n(X,M) \longrightarrow \mathrm{Hom}(H_n(X,\mathbf{Z}),M) \longrightarrow \mathrm{Hom}(\pi_n(Z,z),M)\text{(\footnotemark[118])} \tag*{$(*')$}
\]
\footnotetext[118]{grâce à la composition avec $(*)$.}
\texttt{[lire : }$\mathrm{Hom}_\pi$\texttt{]},

\msp{page 182}
et sûrement (en dehors de toute hypothèse de type $\mathrm{b}_n$ ou autre) l’application canonique directe ‘de localisation sur $Y$’, suivi de \dots,
\[
H^n(X,F_X) \longrightarrow H^0(Y,H^n(Z,F_Z)) \longrightarrow \mathrm{Hom}_\pi(\pi_n(Z),M) \tag*{$(**)$}
\]
n’est autre que la composée des applications canoniques ci-dessus. Dans le cas actuel où on suppose $\mathrm{b}_n$, la deuxième flèche dans $(**)$ est \texttt{[un]} isomorphisme, donc la composée ‘n’est autre’ que la flèche de localisation elle-même. D’autre part $\pi_n(Z,z) \longrightarrow \pi_n(X,x)$ dans $(*)$ est nul, donc aussi $\pi_n(Z,z) \longrightarrow H_n(X,\mathbf{Z})$, donc aussi la flèche composée $(*')$ qui s’en déduit. Revenons alors à la suite exacte de Leray (page précédente), on trouve que $H^n(Y,F) \xrightarrow{\ \sim\ } H^n(X,F_X)$.

Comment récupérer le cas où $F$ n’est \textit{pas} constant? Dans le cas $n=2$, cela marchait, en passant aux revêtements universels $\tilde{X}$, $\tilde{Y}$ de $X$, $Y$ (par rapport aux points $x$, $y$ choisis), et en écrivant les suites spectrales
\[
\begin{tikzcd}[row sep=small, column sep=small]
H^*(Y,F) \arrow[d] & \Longleftarrow & E_2^{pq} \arrow[d] & = & H^p(\pi,H^q(\tilde{Y},F_{\tilde{Y}})) \\
H^*(X,F_X) & \Longleftarrow & E^{\prime pq}_2 & = & H^p(\pi,H^q(\tilde{X},F_{\tilde{X}})).
\end{tikzcd}
\]
L’hypothèse $\mathrm{b}_n$ est conservée en passant de $f : X \longrightarrow Y$ à $\tilde{f} = \tilde{f}_x : \tilde{X} \longrightarrow \tilde{Y}$, donc on sait que
\printedpage{122}
% head: A. Grothendieck : Dérivateurs, Ch. XII, Édition des Universités de Montpellier II et Paris VII
\msp{page 183}
\[
E_2^{pq} \xrightarrow{\ \sim\ } E^{\prime pq}_2 \qquad \text{\texttt{[est un]} isomorphisme si } q \le n. \tag*{$(*)$}
\]
Dans le cas $n=2$, utilisant que
\[
E_2^{p1} = E^{\prime p1}_2 = 0 \qquad \text{car } H^1(\tilde{X}) = H^1(\tilde{Y}) = 0, \qquad
\begin{tikzpicture}[baseline=0pt, x=0.5cm, y=0.5cm]
\draw[->] (0,0) -- (0,2.8);
\draw[->] (0,0) -- (5,0);
\node[above left] at (0,2.6) {$q$};
\node[left] at (0,2) {$2$};
\foreach \i in {0,...,3} { \fill (\i,2) circle (2.5pt); \fill (\i,0) circle (2.5pt); }
\node at (0,1) {$0$};
\node at (1,1) {$0$};
\node at (2,1) {$0$};
\end{tikzpicture}
\qquad ,
\]
on trouve deux suites exactes
\[
\begin{tikzcd}
0 \arrow[r] & H^2(\pi,H^0(\tilde{Y})) \arrow[r] \arrow[d, "\text{iso}"] & H^2(Y) \arrow[r] \arrow[d] & H^0(\pi,H^2(\tilde{Y})) \arrow[r] \arrow[d, "\text{iso}"] & H^3(\pi,H^0(\tilde{Y})) \arrow[d, "\text{iso}"] \\
0 \arrow[r] & H^2(\pi,H^0(\tilde{X})) \arrow[r] & H^2(X) \arrow[r] & H^0(\pi,H^2(\tilde{X})) \arrow[r] & H^3(\pi,H^0(\tilde{X})),
\end{tikzcd}
\]
et un homomorphisme de l’une dans l’autre, qui permet de conclure par le lemme des 5 que $H^2(Y,F) \longrightarrow H^2(X,F_X)$ est \texttt{[un]} isomorphisme. Mais pour $n$ quelconque, la relation $(*)$ sur les termes $E_2$ (pour $q \le n$) ne me semble pas impliquer la relation similaire pour les $E_n$ supérieures, et de là pour les termes $E_\infty$, ce qui permettrait de conclure.

Je reviens donc à la suite exacte p. 181, avec $F$ faisceau \textit{tordu.} J’ai envie de trouver une application canonique
\[
H^n(X,G) \longrightarrow \overbrace{\mathrm{Hom}_\pi(\pi_n(X,x),G_x)}^{\text{\scriptsize Homs \textit{additifs} bien sûr, i.e. de $\mathbf{Z}[\pi]$-modules}} \qquad ,
\]
où $G$ est un faisceau abélien localement constant quelconque sur $X$, et $\pi = \pi_1(X,x)$. Si j’interprète
\msp{page 184}
$\pi_n(X,x)$ par des classes d’homotopie d’applications ponctuées
\[
(S^n,x_0) \xrightarrow{\ \varphi\ } (X,x)
\]
j’ai bien, pour une telle application $\varphi$,
\[
\varphi^* : H^n(X,G) \longrightarrow H^n(S^n,\ \overbrace{\varphi^*G}^{\simeq\,(G_x)_{S^n}}\ ) \simeq G_x
\]
(tenant compte que $S^n$ est 1-connexe, et que sa $n$-cohomologie à valeurs dans un faisceau constant $M$ est $M$.) Sûrement l’application
\[
H^n(X,G) \longrightarrow G_x
\]
\printedpage{123}
% head: A. Grothendieck : Dérivateurs, Ch. XII, Édition des Universités de Montpellier II et Paris VII
ainsi définie ne dépend que de la classe d’homotopie ponctuée de $\varphi$, d’où
\[
\pi_n(X,x) \times H^n(X,G) \longrightarrow G_x,
\]
et il faut de plus prouver a) que c’est compatible avec les opérations de $\pi = \pi_1(X,x)$, et b) que c’est additif en $\pi_n(X,x)$, pour avoir l’application cherchée
\[
H^n(X,G) \longrightarrow \mathrm{Hom}_\pi(\pi_n(X,x),G_x).
\]
Si j’admets ceci, il n’y a qu’un pas pour admettre aussi que la flèche qui nous intéresse vraiment (dans la suite exacte p. 181)
\msp{page 185}
\[
H^n(X,F_X) \longrightarrow \mathrm{Hom}_\pi(\pi_n(Z,z),\underbrace{M}_{=\,F_y})
\]
n’est autre que la composée de la flèche précédente (page précédente) pour $G = F_X$, avec
\[
\mathrm{Hom}_\pi(\pi_n(X,x),M) \longrightarrow \mathrm{Hom}_\pi(\pi_n(Z,x),M)
\]
déduite de
\[
\pi_n(Z,x) \longrightarrow \pi_n(X,x).
\]
Or cette flèche est nulle, par l’hypothèse $\mathrm{b}_n$. Donc la flèche $\alpha$ dans la suite exacte page 181 est nulle, donc $H^n(Y,F) \longrightarrow H^n(X,F_X)$ est \texttt{[un]} isomorphisme, q.e.d. (mais il y a pas mal de travail de vérification à faire \dots).

$\mathrm{c}_n \Longrightarrow \mathrm{a}_{n-1}$. On peut procéder par récurrence sur $n$. (NB L’enoncé a un sens même pour $n=1$, où il est prouvé, et même $\mathrm{c}_1 \Longrightarrow \mathrm{b}_1$, et $\mathrm{b}_1 \Longrightarrow \mathrm{a}_0$, plus généralement $\mathrm{b}_n \Longrightarrow \mathrm{a}_{n-1}$, immédiat par \texttt{[la]} suite exacte d’homotopie.) Donc on sait déjà, en plus de l’hypothèse $\mathrm{c}_n$, qu’on a $\mathrm{a}_{n-2}$, i.e. $Z$ est $(n-2)$-connexe. Il faut prouver qu’il est même $(n-1)$-connexe, i.e. $\pi_{n-1}(Z)=0$. Le cas $n=2$ demande % [sic] L’enoncé -> L’énoncé, missing accent on é
\msp{page 186}
peut-être une attention particulière, car on ne sait pas même encore que $Z$ est 1-connexe. Mais on sait que $\mathrm{c}_2$ implique $\mathrm{c}_1$, qui équivaut \texttt{[?]} à $\mathrm{b}_1$, donc on a
\[
\pi_1(X,x) \longrightarrow \pi_1(Y,y) \qquad \text{isomorphisme,}
\]
donc
\[
\pi_1(Z,x) \longrightarrow \pi_1(X,x) \qquad \text{\texttt{[est]} \textit{nul,}}
\]
donc
\[
\pi_2(Y,y) \longrightarrow \pi_1(Z,x) \qquad \text{\texttt{[est un]} épimorphisme,}
\]
ce qui implique tout au moins que $\pi_1(Z,x)$ est \textit{abélien.} Il faut le voir comme un $\pi$-module, et prouver sa nullité comme tel (moyennant $\mathrm{c}_2$).

\printedpage{124}
% head: A. Grothendieck : Dérivateurs, Ch. XII, Édition des Universités de Montpellier II et Paris VII

Passons au cas général. On a pour $F$ localement constant abélien sur $Y$
\[
H^*(X,F_X) \Longleftarrow \underbrace{E_2^{pq} = H^p(Y,H^q(Z,F_Z))}_{\begin{array}{l}\scriptstyle = 0 \text{ si } 0<q<n-1\\ \scriptstyle \text{car } Z \text{ est } (n-2)\text{-connexe}\end{array}}, \qquad
\begin{tikzpicture}[baseline=0pt, x=0.6cm, y=0.6cm]
\draw[->] (0,0) -- (0,5);
\draw[->] (0,0) -- (6,0);
\node[left] at (0,4) {$n-1$};
\foreach \i in {0,...,4} { \fill (\i,4) circle (2.5pt); \fill (\i,0) circle (2.5pt); }
\node at (1,2) {$0$};
\node at (3,2) {$0$};
\end{tikzpicture}
\qquad ,
\]
D’où
\[
0 \longrightarrow H^{n-1}(Y,F) \longrightarrow H^{n-1}(X,F_X) \longrightarrow \underbrace{H^0(Y,H^{n-1}(Z,F_Z))}_{\begin{array}{l}\scriptstyle \mathrm{Hom}_\pi(\pi_{n-1}(Z,x),F_x)\\ \scriptstyle \text{par Hurewicz}\end{array}} \longrightarrow H^n(Y,F) \longrightarrow H^n(X,F_X).
\]
L’hypothèse $\mathrm{c}_n$ implique qu’on a
\msp{page 187}
\[
\mathrm{Hom}_\pi(\pi_{n-1}(Z,x),\underbrace{M}_{=\,F_x}) = 0
\]
et ceci \textit{quel que soit le $\pi$-module $M$,} puisque $\pi = \pi_1(X,x) \xrightarrow{\ \sim\ } \pi_1(Y,y)$. Cela implique (en prenant p. ex. $M = \pi_{n-1}(Z,x)$ !) que $\pi_{n-1}(Z,x)=0$. (Mais en faisant attention, dans le cas $n=2$, qu’on sait déjà que $\pi_1(Z,x)$ est abélien).

Donc le théorème rectifie page 178 a bien l’air vrai. Il est un peu triste, puisque dans chacune des implications % [sic] rectifie -> rectifié, past participle intended
\[
\mathrm{a}_n \ \Longrightarrow\ \mathrm{b}_n \ \Longrightarrow\ \mathrm{c}_n \ \Longrightarrow\ \mathrm{a}_{n-1}
\]
on perd un chouia, et en fin de compte on sort avec la magnifique implication
\[
\mathrm{a}_n \ \Longrightarrow\ \mathrm{a}_{n-1}.
\]
De toutes façons, il est $\pm$ clair que le “bon” localiseur $W_n$, celui qui doit définir les “$n$-types d’homotopie” comme objets de $\mathrm{Hot}_{W_n}$, est défini par la condition $\mathrm{a}_n$.

\section*{2\quad Critère de $W_n$-asphéricité.}

Je vais essayer de faire mieux quand $Y$ est réduit à un point. Cette fois, je


\endgroup
% END UNIT Der-ChXIIc-30-03-03/Der-ChXIIc-30-03-03_p006-010.tex
% BEGIN UNIT Der-ChXIIc-30-03-03/Der-ChXIIc-30-03-03_p011-015.tex
\begingroup\makeatletter%
\newcommand{\printedpage}[1]{\marginpar{\footnotesize p.~#1}}%
\newcommand{\msp}[1]{\marginpar{\footnotesize [#1]}}%
\newcommand{\Ch}{\underline{\mathrm{Ch}}}%
\newcommand{\Fib}{\underline{\mathrm{Fib}}}%
\newcommand{\Fl}{\underline{\mathrm{Fl}}}%
\newcommand{\Hom}{\underline{\mathrm{Hom}}}%
\newcommand{\Cat}{\mathrm{Cat}}%
\makeatother\let\printedpage\derpp\let\msp\dermsp%
\let\mspin\DERmspin%


\printedpage{125}
% head: A. Grothendieck : Dérivateurs, Ch. XII, Édition des Universités de Montpellier II et Paris VII
% foot: Version 30 mars 2003

\msp{page 188}
veux être strict, et donner (si je peux) des démonstrations complètes dans le cadre où je me place à présent, celui du modéliseur Cat, et en n'utilisant que ce que je sais y démontrer directement à présent, sans me référer à la théorie semi-simpliciale ou topologique.

\textbf{Théorème 2} (?). \textit{Soit $X$ dans} Cat, \textit{$n \in \mathbf{N}$ un entier $\geq 1$. Alors les conditions suivantes sont équivalentes :}

\begin{itemize}
\item[\textit{$a_n$)}] \textit{$X$ est $0$-connexe, et si $x$ est dans $X$, $\pi_i(X,x) = 0$ pour $i \in [1,n]$ (condition qui ne dépend pas du choix de $x$), i.e. les $\Omega^i(X,x)$ pour $i \in [1,n]$ sont $0$-connexes.}
\item[\textit{$b_n$)}] \textit{Pour tout groupe $G$, $G \longrightarrow H^0(X,G_X)$ est un isomorphisme, et $H^1(X,G_X) = 0$. De plus, si $G$ est abélien, on a aussi $H^i(X,G_X) = 0$, pour $i \in [1,n]$.}
\end{itemize}

\textsc{Démonstration}. Pour $n = 1$ cette équivalence est $\pm$ tautologique et a été rappelée (sous une forme légèrement différente) au début (p.~71). On procède par récurrence. On peut donc supposer que $n \geq 2$, que $X$ satisfait $\mathrm{a}_{n-1}$, $\mathrm{b}_{n-1}$, et prouver qu'alors $\mathrm{a}_n$, $\mathrm{b}_n$ sont équivalentes, i.e. $\pi_n(X,x) = 0$ équivaut à $H^n(X,M) = 0$ pour tout groupe abélien $M$.

\msp{page 189}
On utilise la $W_\infty$-fibration (parfaite)
\[
\begin{tikzcd}[row sep=large]
E = \Ch(X;x,-) \arrow[d] & E\ W_\infty\text{-asphérique} \\
X &
\end{tikzcd}
\]
de fibre $Z$ en $x$ égale à $\Omega(X,x)$.

Supposons $\mathrm{a}_n$) pour $X$, donc $a_{n-1}$) pour $\Omega(X,x)$. Par hypothèse de récurrence, on a donc
\[
H^i(Z,M) = 0 \qquad \text{pour } i \in [1,n-1],
\]
$M$ groupe abélien. Donc dans la suite spectrale de Leray,
\[
H^*(E,M) \Longleftarrow E_2^{pq} = H^p(X,H^q(Z,M))
\]
on a
\[
E_2^{pq} = 0 \qquad \text{si } q \in [1,n-1], \qquad
\begin{tikzpicture}[baseline=0]
\draw[->] (0,0) -- (2.2,0);
\draw[->] (0,0) -- (0,2.1);
\foreach \x in {0,0.35,0.7,1.05,1.4} {
  \fill (\x,0) circle (1.5pt);
  \fill (\x,1.4) circle (1.5pt);
}
\node[left] at (-0.05,1.4) {$n$};
\node at (0.35,0.7) {$0$};
\node at (1.05,0.7) {$0$};
\end{tikzpicture}
\quad .
\]

\printedpage{126}
% head: A. Grothendieck : Dérivateurs, Ch. XII, Édition des Universités de Montpellier II et Paris VII
% foot: Version 30 mars 2003

On a donc la suite exacte
\[
0 \longrightarrow H^n(X,M) \longrightarrow H^n(E,M)
\]
et comme $E$ est $W_\infty$-asphérique, on a $H^n(E,M) = 0$, donc $H^n(X,M) = 0$.

Inversement, supposons $\mathrm{b}_n$), prouvons $\mathrm{a}_n$). On sait déjà qu'on a $\mathrm{a}_{n-1}$), donc $\mathrm{a}_{n-2}$ pour $\Omega(X,x)$, ce qui implique
\[
H^i(Z,M) = 0 \qquad \text{pour } i \in [1,n-2],
\]
donc
\[
E_2^{pq} = 0 \qquad \text{si } q \in [1,n-2].
\]
On a donc la suite exacte

\msp{page 190}
\[
0 \longrightarrow \underbrace{H^{n-1}(X,M)}_{=\,0} \longrightarrow \underbrace{H^{n-1}(E,M)}_{=\,0} \longrightarrow H^0(X,H^{n-1}(Z,M)) \longrightarrow \underbrace{H^n(X,M)}_{=\,0} \longrightarrow H^n(E,M),
\]
donc
\[
H^0(X,\ \underbrace{H^{n-1}(Z,M)}_{\begin{array}{l}\text{système local}\\ \textit{constant}\text{ sur } X,\\ \text{car } X \text{ est 1-connexe}\end{array}}\ ) = 0,
\]
donc
\[
H^{n-1}(Z,M) = 0,
\]
ce qui, par hypothèse de récurrence, implique $\pi_{n-1}(Z,z) = 0$, donc $\pi_n(X,x) = 0$, q.e.d.

\textbf{Corollaire.} \textit{Pour que $X$ soit $W_\infty$-asphérique, i.e. pour qu'on ait $G \xrightarrow{\ \sim\ } H^0(X,G)$, $H^1(X,G) = 0$} \texttt{[pour]} \textit{tout groupe} \texttt{[}$G$\texttt{]}\textit{, et $H^i(X,G) = 0$ si $i \geq 2$ pour $G$ abélien, il faut et il suffit que} $1^\circ$) \textit{$X$ soit} $0$\textit{-connexe, et ($x$ étant un point de $X$)} $2^\circ$) \textit{les $\Omega^i(X,x)$} \texttt{[soient]} $0$\textit{-connexes pour $i \geq 1$.}

\msp{page 191}

\section*{3\quad Sur la philosophie des fibres homotopiques.}

Le critère d'asphéricité (pour $W_\infty$) que je viens d'obtenir, grâce aux foncteur $\Ch(X)$ (ou au choix, $\Ch_\infty(X)$) et ses dérivés $\Omega(X,x)$ et les itérés, enlève tout doute (s'il pouvait en rester encore) que les `espaces de chemins' que j'ai construit (assez péniblement!) sont bel et bien `les bons'. Je voudrais maintenant me rendre compte de façon la plus précise possible de leur `canonicité', dans quelle mesure ces constructions sont indépendant d'un choix (p.ex. celui entre la théorie $\Ch$, et $\Ch_\infty$). % [sic] aux foncteur -> au foncteur, number agreement

Tout revient, finalement, à des questions de factorisation d'une flèche quelconque
\[
f : X \longrightarrow Y
\]
\printedpage{127}
% head: A. Grothendieck : Dérivateurs, Ch. XII, Édition des Universités de Montpellier II et Paris VII
% foot: Version 30 mars 2003
de Cat en
\[
X \xrightarrow{\ i_f\ } \bar{X} \xrightarrow{\ \bar{f}\ } Y,
\]
avec
\[
i_f \in W, \qquad \bar{f} \in W\text{-fib}
\]
(le fait que $\bar{f}$ puisse même être choisie $W_\omega$-parfait, ne me paraissant pas essentiel ici).

C'est bien de cela qu'il s'agit pour construire

\msp{page 192}
un objet `chemins dans $X$'
\[
\Ch(X),
\]
en factorisant l'application diagonale
\[
X \xrightarrow{\ \mathrm{diag}_X\ } X \times X,
\]
et aussi quand on construit $\Omega(X,x)$ comme fibre en $x$ d'une $W$-fibration
\[
E(X,x) \longrightarrow (X,x)
\]
qui factorise l'inclusion
\[
\{x\} \longrightarrow X
\]
(donc $E(X,x)$ étant $W$-asphérique).

On peut être plus ou moins exigeant, et regarder, soit pour un `but' $Y$ (ou plutôt un $S$) fixé dans Cat, une factorisation fonctorielle
\[
C : \Cat/S \longrightarrow \Fib_W S \hookrightarrow \Cat/S
\]
du morphisme structural $X \longrightarrow S$ d'un objet, en
\[
X \xrightarrow{\ i_X\ } C(X) \longrightarrow S,
\]
i.e. on se donne le foncteur $C$, \textit{plus}

\msp{page 193}
le morphisme fonctoriel
\[
i_C : \mathrm{id}_{\Cat/S} \longrightarrow C
\]
avec
\[
(i_X : X \longrightarrow C(X)) \in W \qquad \forall\ X \text{ dans } \Cat/S,
\]
soit être encore plus exigeant et faire varier en même temps la source et le but de $f : X \longrightarrow Y (= S)$, et demander un foncteur
\[
\begin{aligned}
C : \Fl(\Cat) \quad &\longrightarrow \quad \Cat \\
f \quad &\longmapsto \quad C(f),
\end{aligned}
\]
\printedpage{128}
% head: A. Grothendieck : Dérivateurs, Ch. XII, Édition des Universités de Montpellier II et Paris VII
% foot: Version 30 mars 2003
avec des morphismes fonctoriels en $f$
\[
\underbrace{s(f)}_{=\,X} \xrightarrow{\ i_f\ } \underbrace{C(f)}_{=\,\bar{X}} \xrightarrow{\ p_f\ } \underbrace{b(f)}_{=\,Y}
\]
($i$, $p$ des flèches dans $\Hom(\Fl(\Cat),\Cat)$), avec
\[
i_f \in W, \qquad p_f \in W\text{-fib}
\]
pour toute flèche $f$ dans Cat.

\msp{page 194}
Dans les deux points de vue, dans quel sens a-t-on \textit{unicité} d'une telle construction?

Prenons d'abord le premier point de vue, factorisation des flèches à but fixé $S$, donc \texttt{[c'est une]} question dans $\Cat/S$. J'ai oublié une propriété importante, mais elle résulte de ce qui a été dit : le foncteur
\[
\begin{tikzcd}[column sep=small]
X \longmapsto C(X) : & \Cat/S \arrow[rr] \arrow[dr, dotted] & & \Cat/S \\
& & \Fib_W S \arrow[ur, dotted, hook] &
\end{tikzcd}
\]
\textit{transforme} \texttt{[les]} \textit{flèches} \texttt{[de]} \textit{$W$ en itou}, comme on voit sur le diagramme commutatif
\[
\begin{tikzcd}[row sep=large, column sep=large]
X \arrow[r, "i_X"] \arrow[d, "f"'] & C(X) \arrow[d, "C(f)"] \\
Y \arrow[r, "i_Y"] & C(Y),
\end{tikzcd}
\]
où $i_X$, $i_Y$ sont dans $W$, donc
\[
f \in W \qquad \Longleftrightarrow \qquad C(f) \in W.
\]
Soit
\[
(C',i')
\]
un deuxième foncteur, $\Cat/S \longrightarrow \Cat/S$,

\msp{page 195}
muni d'un morphisme
\[
i' : \mathrm{id}_{\Cat/S} \longrightarrow C'
\]
qui est dans $W$ argument par argument, et tel que $C'$ se factorise par $\Fib_W S$. On peut comparer $C$ et $C'$ moyennant $C \circ C'$, car
\[
\begin{tikzcd}[column sep=small]
& C \circ C' & \\
C = C \circ \mathrm{id} \arrow[ur, "C\ast i'"] & & C' = \mathrm{id} \circ C', \arrow[ul, "i\ast C'"']
\end{tikzcd}
\]
\printedpage{129}
% head: A. Grothendieck : Dérivateurs, Ch. XII, Édition des Universités de Montpellier II et Paris VII
% foot: Version 30 mars 2003
i.e. on a, pour tout $X$ dans $\Cat/S$,
\[
\begin{tikzcd}[column sep=small]
& C(C'(X)) & \\
C(X) \arrow[ur, "C(i'_X)"] & & C'(X), \arrow[ul, "i_{C'(X)}"']
\end{tikzcd}
\]
le deux flèches étant dans $W$ (celle de gauche parce que $C(W) \subseteq W$). Ce sont des flèches dans $\Fib_W S$, et si on suppose que $W$ satisfait W(4b), elles induisent des $W$-équivalences en chaque fibre. Cela montre donc déjà % [sic] le deux flèches -> les deux flèches, typo

\msp{page 196}
qu'à isomorphisme (peut-être non-canonique) près, le foncteur composé défini par $C$,
\[
\begin{tikzcd}
X \arrow[r, mapsto] & C(X) \arrow[r, mapsto] & (s \longmapsto \mathrm{hot}_W(C(X)_s)) \\
\Cat/S \arrow[r] \arrow[rr, bend right, "\varphi_C"'] & \Fib_W S \arrow[r] & \Hom(S^\circ,\mathrm{Hot}_W),
\end{tikzcd}
\]
ne dépend pas du choix particulier de $C$, i.e. deux tels foncteurs `fibres homotopiques' (correspondant à deux choix $C$, $C'$) sont toujours isomorphes. Mais au lieu d'utiliser $C \circ C'$ pour définir un tel isomorphisme, on aurai pu utiliser $C' \circ C$, ou quelque autre procédé. On aimerait comprendre s'il y a unicité véritable, une caractérisation intrinsèque de ces isomorphismes $\varphi_C \xrightarrow{\ \sim\ } \varphi_{C'}$, indépendamment de telle construction ou de telle autre. % [sic] aurai pu -> aurait pu, typo

Je présume que $C \circ C'$ est lui aussi muni de
\[
\mathrm{id}_{\Cat/S} \longrightarrow C \circ C',
\]

\msp{page 197}
et même a priori de deux façons possible
\[
\begin{tikzcd}[column sep=large, row sep=large]
& C(C'(X)) & \\
C(X) \arrow[ur, "C(i'_X)"] & & C'(X) \arrow[ul, "i_{C'(X)}"'] \\
& X, \arrow[ul, "i_X"] \arrow[ur, "i'_X"'] \arrow[uu, dotted] &
\end{tikzcd}
\]
soit par $C(i'_X) \circ i_X$, soit par $i_{C(X)} \circ i'_X$, heureusement elles sont égales, car pour toute flèche $X \xrightarrow{\ i'\ } X'$ dans $\Cat/S$, le carré % [sic] i_{C(X)} -> i_{C'(X)}, the diagram above has i_{C'(X)}


\endgroup
% END UNIT Der-ChXIIc-30-03-03/Der-ChXIIc-30-03-03_p011-015.tex
% BEGIN UNIT Der-ChXIIc-30-03-03/Der-ChXIIc-30-03-03_p016-020.tex
\begingroup\makeatletter%
\newcommand{\printedpage}[1]{\marginpar{\footnotesize p.~#1}}%
\newcommand{\msp}[1]{\marginpar{\footnotesize [#1]}}%
\newcommand{\Hom}{\underline{\mathrm{Hom}}}%
\newcommand{\Fib}{\underline{\mathrm{Fib}}}%
\newcommand{\Cat}{\mathrm{Cat}}%
\newcommand{\HOT}{\mathrm{HOT}}%
\newcommand{\Hot}{\mathrm{Hot}}%
\newcommand{\hot}{\mathrm{hot}}%
\newcommand{\id}{\mathrm{id}}%
\makeatother\let\printedpage\derpp\let\msp\dermsp%
\let\mspin\DERmspin%


\printedpage{130}
% head: A. Grothendieck : Dérivateurs, Ch. XII, Édition des Universités de Montpellier II et Paris VII

\[
\begin{tikzcd}[column sep=small, row sep=small]
 & C(X') & \\
C(X) \arrow[ur, "C(i')"] & & X' \arrow[ul, "i_{X'}"'] \\
 & X \arrow[ul, "i_X"] \arrow[ur, "i'"'] &
\end{tikzcd}
\]

est commutatif (par fonctorialité de $i_X : X \longrightarrow C(X)$), et on applique ceci au cas de $X' = C'(X)$, $i' = i'_X$.

Cette remarque nous encourage à faire un saut dans l'abstraction, et d'introduire la catégorie\msp{page 198}
\[
\mathcal{C} = \mathcal{C}_S = \begin{cases}
\text{catégorie des couples } (C,i), \text{ avec} \\
C : \Cat/S \longrightarrow \Cat/S \text{ se factorisant par } \Fib_W S, \\
i : \id_{\Cat/S} \longrightarrow C \text{ tel que } i_X \in W \ \forall\ X \text{ dans } \Cat/S,
\end{cases}
\tag*{$(*)$}
\]
une flèche
\[
(C,i) \longrightarrow (C',i')
\]
étant une flèche
\[
C \xrightarrow{u} C' \qquad \text{(dans } \Hom(\Cat/S, \Cat/S))
\]
rendant commutatif le carré \texttt{[lire : le diagramme]}
\[
\begin{tikzcd}[column sep=small]
 & \id_{\Cat/S} \arrow[dl, "i"'] \arrow[dr, "i'"] & \\
C \arrow[rr, "u"] & & C'.
\end{tikzcd}
\]

En d'autres termes, $\mathcal{C}$ est défini comme une sous-catégiorie pleine de % [sic] sous-catégiorie -> sous-catégorie, misspelled word
\[
\id_{\Cat/S}\backslash\Hom(\Cat/S, \Cat/S),
\]
à savoir celle dont les objets $(C,i)$ satisfont les deux conditions $(*)$ plus haut (impliquant l'une et l'autre le localiseur $W$).

\msp{page 199}Ces conditions impliquent ceci

\begin{enumerate}
\item[a)] $\forall\ X$ dans Cat/$S$, $(u_X : C(X) \longrightarrow C'(X)) \in W$, puisque
\[
\begin{tikzcd}[column sep=small]
 & X \arrow[dl, "i_X\in W"'] \arrow[dr, "i'_X\in W"] & \\
C(X) \arrow[rr, "u_X"] & & C'(X),
\end{tikzcd}
\]
\printedpage{131}
% head: A. Grothendieck : Dérivateurs, Ch. XII, Édition des Universités de Montpellier II et Paris VII
et

\item[b)] si $W$ satisfait W(4b), de sorte que $u_X \in W^{\mathrm{f}}_S$, alors ce définit un \textit{isomorphisme} entre $\varphi_C(X)$ et $\varphi_{C'}(X)$ % [sic] ce définit -> ceci définit, missing word
\[
\begin{tikzcd}[row sep=small]
s \arrow[r, mapsto] & \hot_W(C(X)_s) \arrow[d, "\wr\, u_{X,s}"] \\
s \arrow[r, mapsto] & \hot_W(C'(X)_s).
\end{tikzcd}
\]
\end{enumerate}

Mieux encore, dans ce cas, le foncteur composé
\[
\mathcal{C} \times \Cat/S \longrightarrow \Fib_W S \longrightarrow \HOT^{\mathrm{lc}}_W(S)
\]
donne naissance à
\[
\mathcal{C} \longrightarrow \Hom(\Cat/S, \HOT^{\mathrm{lc}}_W(S))
\]
qui transforme toute flèche de $\mathcal{C}$ en \texttt{[un]} \textit{isomorphisme,} donc se factorise par le groupoïde fondamental de $\mathcal{C}$,
\[
\Pi\mathcal{C} \longrightarrow \Hom(\Cat/S, \HOT^{\mathrm{lc}}_W(S)).
\]

\msp{page 200}Quant au foncteur envisagé tantôt
\[
\begin{aligned}
\mathcal{C} &\quad\longrightarrow\quad \underbrace{\Hom\,(\Cat/S, \Hom(S^\circ, \Hot_W))}_{\simeq\Hom((\Cat/S)\times S^\circ, \Hot_W)} \\
C &\quad\longmapsto\quad ((X,s) \longmapsto \hot_W(C(X)_s))\,,
\end{aligned}
\]
il se déduit du précédent en composant avec
\[
\begin{tikzcd}[column sep=small]
\HOT_W(S) \arrow[rr] \arrow[dr, dotted] & & \Hom(S^\circ, \Hot_W), \\
 & \Hom(\Pi S^\circ, \Hot_W) \arrow[ur, dotted] &
\end{tikzcd}
\]
plus exactement, avec le foncteur qui s'en déduit
\[
\Hom(\Cat/S, \HOT_W(S)) \longrightarrow \Hom(\Cat/S, \Hom(S^\circ, \Hot_W)).
\]

On voit poindre ici un principe plausible de canonicité d'isomorphismes, en ce qui concerne les foncteurs
\[
C_W : \Cat/S \longrightarrow \HOT^{\mathrm{lc}}_W(S)
\]
définis par les $C \in \mathrm{Ob}\,\mathcal{C}$, par la relation
\[
\Pi\mathcal{C} \approx \ \text{catégorie ponctuelle, i.e. } \mathcal{C} \text{ est 1-connexe.}
\]

De plus, les arguments utilisés jusqu'ici sont de nature à tel point formelle, que les aspects spécifiques du Modéliseur Cat
\printedpage{132}
% head: A. Grothendieck : Dérivateurs, Ch. XII, Édition des Universités de Montpellier II et Paris VII
\msp{page 201}semblent n'y jouer pratiquement aucun rôle, ou un rôle très effacé.

Le carré de la page 197 montre surtout que la catégorie $\mathcal{C}$ est filtrante - $(C,i)$, $(C,i')$ s'envoient l'un et l'autre dans $(C'',i'')$, avec $C'' = C \circ C'$, et $i'' = (C * i') \circ i = (i * C') \circ i'$. Il y a même une loi de composition dans $\mathcal{C}$, nullement commutative ni unitaire c'est sûr, mais sûrement associative, et (si cette composition est notée $*$) avec des morphismes fonctoriels % [sic] (C,i') -> (C',i'), the second couple of the following formula is (C',i')
\[
\xi \xrightarrow{\alpha_{\xi,\xi'}} \xi * \xi' \xleftarrow{\beta_{\xi,\xi'}} \xi'
\]
qui assurent le caractère filtrant croissant de $\mathcal{C}$. Ces flèches semblent remplacer en quelque sorte l'unité bilatère absente, et on s'attend à ce qu'elles aient les mêmes excellentes propriétés que si on avait une telle unité, concernant les composés doubles $(\xi * \xi') * \xi'' \simeq \xi * (\xi' * \xi'')$,

\msp{page 202}
\[
(\xi * \xi') * \xi'' \xrightarrow{\sim} \xi * (\xi' * \xi'').
\]

Il y a deux façons d'envoyer un produit binaire
\[
\xi * \xi', \qquad \xi' * \xi'', \qquad \xi * \xi''
\]
dans le produit ternaire, et bien elles doivent dans chacun des trois cas, être égales. Et de même il y a, ceci admis, encore deux façons d'envoyer $\xi$ dans le produit ternaire, en factorisant soit par $\xi * \xi'$, soit par $\xi * \xi''$, et de même pour les deux autres foncteurs $\xi'$, $\xi''$; et là encore, ces deux façons doivent être égales, dans chacun de ces trois cas - donc $3 + 3 = 6$ égalités en tout à écrire et à vérifier - on a vu pire!

La philosophie calculatoire devrait être celle-ci : Si on a un produit $n$-uple
\[
\xi_1 * \xi_2 * \cdots * \xi_n
\]
(où je me dispense de mettre des parenthèses \dots), pour tout produit $m$-uple partiel $(m \in [1, n-1])$\msp{page 203}
\[
\xi_{i_1} * \xi_{i_2} * \cdots * \xi_{i_m} \qquad\qquad (i_1 < i_2 < \cdots < i_m),
\]
en utilisant à tour de bras les isomorphismes d'associativité et les $\alpha_{\xi,\xi'}$, $\beta_{\xi,\xi'}$, il n'y a pourtant qu'une seule façon de l'envoyer dans le produit $n$-uple. J'ai bien l'impression que ça doit être vrai. (Est-ce que ça implique déjà que $\mathcal{C}$ est 1-connexe?)

En fait, il y a bien une unité bilatère, mais elle n'est pas dans $\mathcal{C}$, mais dans la catégorie plus grande
\[
\id_{\Cat/S}\backslash\Hom(\Cat/S, \Cat/S),
\]
qui elle, comme toute catégorie
\[
\id_{\mathcal{X}}\backslash\Hom(\mathcal{X}, \mathcal{X}),
\]
\printedpage{133}
% head: A. Grothendieck : Dérivateurs, Ch. XII, Édition des Universités de Montpellier II et Paris VII
est une catégorie monoïdale, i.e. avec produit associatif \textit{et} unitaire. Les flèches $\alpha_{\xi,\xi'}$, $\beta_{\xi,\xi'}$ sont les flèches qui résultent de la présence de cette unité (vu que $\mathcal{C}$ est une sous-catégorie pleine), et alors il doit devenir clair qu'on a les compatibilités dites.

\msp{page 204}Mais même une catégorie monoïdale, est elle 1-connexe? Par exemple $B_G$, pour $G$ un groupe disons, est bien une catégorie monoïdale et même ‘groupale’, une $G$-catégorie (tout ce qu'il y a de strict), et pourtant elle n'est \textit{pas} 1-connexe!

Bien sûr, on peut court-circuiter la question de la 1-connexité de $\mathcal{C}$, en notant qu'après tout on a déjà a priori un objet privilégié dans
\[
\Hom(\Cat/S, \HOT^{\mathrm{lc}}_W(S)),
\]
i.e. un foncteur canonique
\[
\Cat/S \longrightarrow \HOT^{\mathrm{lc}}_W(S),
\]
savoir le foncteur composé
\[
\Cat/S \underbrace{\longrightarrow}_{\substack{\text{localisation}\\ \text{par } W_S}} (\Cat/S)W_S^{-1} \longrightarrow \underbrace{\HOT^{\mathrm{lc}}_W(S)}_{\overset{\text{déf}}{=} (\Fib_W S)(W^{\mathrm{f}}_S)^{-1}} \quad ,
\tag*{$(*)$}
\]
où la deuxième flèche est définie comme le quasi-inverse de
\[
(\Fib_W S)(W^{\mathrm{f}}_S)^{-1} \longrightarrow (\Cat/S)W_S^{-1}
\tag*{$(**)$}
\]

\msp{page 205}déduit de l'inclusion
\[
(\Fib_W S, W^{\mathrm{f}}_S) \hookrightarrow (\Cat/S, W_S).
\]

Le fait que $(*)$ soit une équivalence, se prouve justement par le fait de l'existence d'un foncteur $C : \Cat/S \longrightarrow \Cat/S$ et d'un $i : \id_{\Cat/S} \longrightarrow C$, ayant les propriétés qu'on sait, i.e. du fait que $C \neq \emptyset$. De plus, la donnée d'un objet $C$ de $\mathcal{C}$ donne un moyen commode d'expliciter le quasi-inverse dudit foncteur. Mais ce quasi-inverse a une existence canonique et indépendante du choix d'un $C$, i.e. d'un mode de calcul. Le fait que le foncteur composé $\varphi$ dans $(*)$ puisse se ‘relever’ (à isomorphisme près) en un foncteur % [sic] C \neq \emptyset -> \mathcal{C} \neq \emptyset, the non-emptiness concerns the category \mathcal{C}
\[
C : \Cat/S \longrightarrow \Fib_W S
\]
entre $\varphi(X) = \hot_{W,S}(X)$ et $\hot_{W,S}(C(X))$ et que l'isomorphisme puisse s'expliciter, lui aussi, par une flèche $X$-fonctorielle
\[
i_X : X \longrightarrow C(X), \qquad i_X \in W,
\]
dans Cat/$S$ lui-même (et non seulement dans le localisé $(\Cat/S)W_S^{-1}$),
\printedpage{134}
% head: A. Grothendieck : Dérivateurs, Ch. XII, Édition des Universités de Montpellier II et Paris VII
\msp{page 206}est une vertu particulière du \texttt{[mot illisible]}, qui ne doit pas pour autant faire oublier le caractère intrinsèque des flèches $(*)$, indépendamment de tout mode de calcul.

{\reversemarginpar\marginnote{$\|$}}% double vertical bar printed in the left margin beside the first lines of this paragraph
La clef de la situation, et des questions de canonicité, me semble être dans le fait que $(**)$ est une équivalence de catégories - et il faut avouer que, tout familier et quasiment banal qu'il soit finalement devenu pour moi, c'est là un fait tout à fait remarquable et nullement évident a priori. (Mais sans doute peut-on l'avoir aussi par les fourbis généraux à la \textsc{Quillen-Thomason}, sans se battre avec un formalisme de chemins \dots) Les ‘isomorphismes canoniques’ qu'on peut établir entre une ‘théorie $C$ ’ et une ‘théorie $C'$ ’ (via deux constructions différentes, p. ex., de catégories de chemins, ou via l'approche de \textsc{Quillen-Thomason}, sans chemins) doivent toujours revenir à ceux

\msp{page 207}qu'on peu définir via la comparaison de la ‘théorie $C$ ’ d'une part, la ‘théorie $C'$ ’ de l'autre, avec ce que nous donne directement la flèche composée $(*)$ (p. 204). Finalement, il me semble que \textit{l'existence} d'une théorie des ‘fibres homotopiques’ (ou ‘$W$-homotopiques’) d'une flèche $X \longrightarrow S$ dans Cat, résulte de l'équivalence $(**)$ (p. 204), et cette théorie pourrait sans doute dans une certaine mesure se développer à partir de là, sans choix d'un foncteur $C$ ni d'une structure de catégorie de modèles au sens de \textsc{Quillen}. % [sic] qu'on peu définir -> qu'on peut définir, verb form

\section*{4\quad Axiomatisation des ‘fibres $W$-homotopiques’.}

La reflexion de hier me fait revenir sur l'axiomatique dans des ‘catégories de modèles’ % [sic] reflexion -> réflexion, missing accent
\[
(\mathcal{M}, W)
\]
(\footnote[119]{$W$ satisfait l'axiome de saturation faible : les isomorphismes \texttt{[sont]} dans $W$, et si deux parmi $f$, $g$, $gf$ \texttt{[sont]} dans $W$, le troisième aussi.}) plus générales que Cat, dans l'esprit de VI (cf. théorème-scholie 2, p. 108-113). Pour le moment, on va supposer seulement $\mathcal{M}$ stable par produits fibrés (on n'a pas besoin d'objet final). On introduit l'ensemble
\[
\mathrm{Fib}_W \subseteq \mathrm{Fl}(\mathcal{M})
\]

\msp{page 208}des \textit{$W$-fibrations} (sous-entendu ‘ultrafortes’!) par la condition

\begin{quote}
$f : X \longrightarrow Y$ est dans $\mathrm{Fib}_W$ si et seulement si le foncteur changement de base par $X$,
\[
\begin{aligned}
Y' &\quad\longmapsto\quad X' = X \times_Y Y' \\
\mathcal{M}/Y &\quad\longrightarrow\quad \mathcal{M}/X
\end{aligned}
\]
\end{quote}


\endgroup
% END UNIT Der-ChXIIc-30-03-03/Der-ChXIIc-30-03-03_p016-020.tex
% BEGIN UNIT Der-ChXIIc-30-03-03/Der-ChXIIc-30-03-03_p021-025.tex
\begingroup\makeatletter%
\derleqno%
\NoAutoSpacing%
\newcommand{\printedpage}[1]{\marginpar{\footnotesize p.~#1}}%
\newcommand{\msp}[1]{\marginpar{\footnotesize [#1]}}%
\newcommand{\cM}{\mathcal{M}}%
\newcommand{\cF}{\mathcal{F}}%
\newcommand{\Ho}{\mathrm{Ho}}%
\newcommand{\ho}{\mathrm{ho}}%
\newcommand{\Fl}{\mathrm{Fl}}%
\newcommand{\Fib}{\underline{\mathrm{Fib}}}%
\newcommand{\Wu}{W^{\mathrm{u}}}%
\newcommand{\lc}{\mathrm{lc}}%
\newcommand{\circled}[1]{\tikz[baseline=(c.base)]\node[draw,circle,inner sep=1pt](c){#1};}%
\makeatother\let\printedpage\derpp\let\msp\dermsp%
\let\mspin\DERmspin%


\printedpage{135}
% head: A. Grothendieck : Dérivateurs, Ch. XII, Édition des Universités de Montpellier II et Paris VII

\begin{quote}
transforme $W_X$ dans $W_Y$, où $W_X \subseteq \Fl(\cM/X)$ \texttt{[est]} défini par

$u \in W_X \iff \sigma(U) \in W$ \ ($\sigma : \cM/X \longrightarrow \cM$ le foncteur ‘source’, ou ‘oubli de $X$’). % [sic] \sigma(U) -> \sigma(u), the arrow is called u
\end{quote}

(\footnote[120]{\textbf{NB} Donner ici la proposition: $\mathrm{Fib}_W$ \textit{stable par composition et par changement de base, et contient les isomorphismes, stable par facteurs directs.}}). Posons, pour tout $S$ dans $\cM$,
\[
\Ho^!_W(S) = (\cM/S)W_S^{-1},
\]
qui joue le rôle d'une catégorie de $W$-types d'homotopie relatifs \textit{localement constants} sur $S$. On aimerait décrire une loi contravariante en $S$ pour $\Ho^!_W(S)$, i.e. pour
\[
f : S' \longrightarrow S
\]
flèche quelconque dans $\cM$, définir
\[
f^*_W : \Ho^!_W(S) \xrightarrow{?} \Ho^!_W(S'). \tag*{$(*)$}
\]
La difficulté est, bien sûr, que le foncteur changement de base
\[
f^* : \cM/S \longrightarrow \cM/S'
\]
n'est \textit{pas} en général compatible avec les \msp{page 209} localiseurs $W_S$, $W_{S'}$. La notion de $W$-fibration sert donc en premier lieu à dégager la condition la plus évidente sur une flèche $f : S' \longrightarrow S$ dans $\cM$, pour que le foncteur changement de base soit compatible avec les localiseurs $W_S$, $W_{S'}$ (c'est \textit{ça} la définition), donc induit $(*)$ sur les localisées.

Le \textit{‘problème des fibres homotopiques’} est justement de définir $(*)$ pour \textit{toute} flèche dans $\cM$, pas seulement pour les $W$-fibrations, de tellle façon que les $\Ho^!_W(S)$ forment, pour $S$ % [sic] tellle -> telle, letter doubled
variable, les fibres d'une catégorie fibrée remarquable sur $\cM$. Quand $\cM$ admet un objet final, de sorte que l'on peut regarder les flèches
\[
p : e \longrightarrow S \qquad (p \text{ comme ‘point’})
\]
dans $\cM$ comme des ‘points’ de $S$, la ‘fibre $W$-homotopique’ d'un objet quelconque $X/S$ de $\cM/S$ en ledit point sera par définition l'élément de
\[
\Ho_W = \Ho_W(\cM) \simeq \Ho^!_W(e)
\]
défini par
\[
p^*(\ho^!_{W/S}(X/S)),
\]
\printedpage{136}
% head: A. Grothendieck : Dérivateurs, Ch. XII, Édition des Universités de Montpellier II et Paris VII
\msp{page 210}
où
\[
\ho^!_{W/S} : \cM/S \longrightarrow \Ho^!_W(S) \overset{\text{déf}}{=} (\cM/S)(W_S)^{-1}
\]
est le foncteur de localisation canonique.

Pour tout $S$ dans $\cM$, soit
\[
\Fib_W S \subseteq \cM/S
\]
la sous-catégorie pleine formée des objets $X/S$ de $\cM/S$ tels que $X \longrightarrow S$ soit $W$-fibrant. On va définir une variante
\[
\Ho^{\lc}_W(S)
\]
de $\Ho^!_W(S)$, comme
\[
\Ho^{\lc}_W(S) \overset{\text{déf}}{=} (\Fib_W S)(\Wu_S)^{-1},
\]
où

\begin{quote}
$\Wu_S \subseteq \Fl(\Fib_W S)$ est le localiseur formé des flèches $u : X \longrightarrow Y$ qui sont ‘universellement sur $S$ dans $W$’, i.e. telles que pour toute flèche $S' \longrightarrow S$ de $\cM$ de but $S$, la flèche déduite par changement de base,
\[
u' : X' = X \times_S S' \longrightarrow Y' = Y \times_S S'
\]
est dans $W_{S'}$.
\end{quote}

On a ‘fait ce qu'il fallait’ pour que $\Ho^{\lc}_W(S)$ soit contravariant en $S$ pour des flèches \textit{quelconques} $f : S' \longrightarrow S$ de $\cM$, elles définissent
\[
f^{*\,\lc}_W : \Ho^{\lc}_W(S) \longrightarrow \Ho^{\lc}_W(S'), \tag*{$(**)$}
\]
\msp{page 211}
puisque les $W$-fibrations sont stables par changement de base, et que par suite $f^* : \cM/S \longrightarrow \cM/S'$ induit
\[
f^*_{\Fib_W} : \Fib_W S \longrightarrow \Fib_W S'
\]
qui, par définition de $\Wu_S$ et $\Wu_{S'}$, est compatible avec les localiseurs
\[
f^*_{\Fib_W} : (\Fib_W S, \Wu_S) \longrightarrow (\Fib_W S', \Wu_{S'}),
\]
et passe donc aux catégories des fractions en une flèche $(*)$.

On peut se demander pourquoi ne pas être plus généreux, et définir $\Ho^{\lc}_W(S)$ comme $(\cM/S)(\Wu_S)^{-1}$, en définissant $\Wu_S$ dans $\cM/S$ de la même façon que tantôt dans $\Fib_W S$ - ce qui ne coûte pas plus cher. Une première réponse est d'expérience le localiseur $\Wu_S$ dans $\cM/S$ tout entier s'avère pratiquement inutilisable, d'une part parce que l'appartenance à $\Wu_S$ (pour une flèche entre objets de $\cM/S$ qui ne sont pas très particuliers, p.ex. $W$-fibrés sur $S$) est invérifiable (sauf pour des isomorphismes et des choses comme ça \dots), et que d'autre part la catégorie de fractions est prohibitivement
\printedpage{137}
% head: A. Grothendieck : Dérivateurs, Ch. XII, Édition des Universités de Montpellier II et Paris VII
\msp{page 212}
grosse et ne peut servir (semble-t-il) à rien. Une autre réponse, c'est qu'on veut que la catégorie obtenue réponde bien à l'intuition de types d'homotopie relatifs \textit{localement constants} sur $S$ - et on veut donc que pour un changement de base quelconque $f : S' \longrightarrow S$ par un

\begin{quote}
$S'$ dans $\cM/S$, d'où un foncteur composé
\[
\begin{tikzcd}[column sep=small]
\cF_S(\Wu_S)^{-1} \arrow[r, "f^{*\,\lc}_W"] \arrow[rr, dotted, bend right=15, "\varphi_{S'}"'] & \cF_{S'}(\Wu_{S'})^{-1} \arrow[r] & \Ho_W
\end{tikzcd}
\]

(\footnote[121]{$\cF$ désigne ici une sous-catégorie fibrée de $\cM/\cM = \underline{\Fl}(\cM)$ sur $\cM$, substitut éventuel de celle formée de $\Fib_W S$ ou de $\cM/S$ tout entier.}) déduit des foncteurs de catégories avec localiseurs
\[
(\cF_S, \Wu_S) \xrightarrow{f^*} (\cF_{S'}, \Wu_{S'}) \xrightarrow{\text{oubli de } S'} (\cM, W)
\]
\end{quote}

(foncteur composé $\varphi_{S'}$ qui peut être vu comme le foncteur \textit{‘fibre homotopique en l'objet $S'$ sur $S$’}, cet objet étant visualisé comme un ‘point’ de $S$ ‘à valeurs dans $S'$ ’ suivant le yoga fonctoriel général \dots), le foncteur
\[
\begin{aligned}
S' &\longmapsto \varphi_{S'}(\xi) \qquad \text{pour } \xi \text{ dans } \cF_S(\Wu_S)^{-1}\\
\cM/S &\longrightarrow \Ho_W
\end{aligned}
\]
transforme flèches dans $W_S$ en isomorphismes : la ‘fibre $W$-homotopique’ en $S'/S$ d'un objet de $\cF_S(\Wu_S)^{-1}$ (vu comme $W$-type d'homotopie \textit{localement constant} sur $S$) doit être ‘la même’, si $S'$ remplacé par $S''$ avec $S'' \longrightarrow S'$ donné et dans $W_S$. C'est pour cela qu'il est raisonnable de prendre $\cF_S$ ‘assez petit’ (savoir $\Fib_W S$, qui est \msp{page 213} exactement ajusté pour cela), plutôt que de prendre $\cM/S$ tout entier.

On peut donc dire que la catégorie $\Fib_W S$, et le localiseur $\Wu_S$ ont été construits ad-hoc, le plus simplement possible, pour donner lieu à un localisé $\Ho^{\lc}_W(S)$ qui réponde à l'intuition de $W$-types d'homotopie relatifs localement constants des deux façons suivantes :

\begin{enumerate}[label={}]
\item[1°)] Il y a spontanément une loi contravariante en $S$, pour des flèches quelconques $S' \longrightarrow S$ dans Cat (on a défini $\Wu_S$ juste exprès en conséquence).
\item[2°)] que la notion de ‘fibre $W$-homotopique en $S'/S$’ d'un objet $\xi$ de $\Ho^{\lc}_W$ satisfasse à la condition précédente, i.e. une flèche $S'' \longrightarrow S'$ dans $\cM/S$ qui est dans $W_S$ définisse un \textit{isomorphisme} (dans $\Ho_W$) des fibres homotopiques
\[
\underbrace{\xi_{S''}}_{\varphi_{S''}(\xi)} \longrightarrow \underbrace{\xi_{S'}}_{\varphi_{S'}(\xi)}
\]
(et on a choisi $\cF_S = \Fib_W(S)$ juste exprès pour ça).
\end{enumerate}

\printedpage{138}
% head: A. Grothendieck : Dérivateurs, Ch. XII, Édition des Universités de Montpellier II et Paris VII

On a donc une excellente notion de fibres $W$-homotopiques pour des $X$ $W$-fibrés sur $S$ - cette fois excellente pour des changements de base \textit{quelconque} - et il n'y a aucun mérite à ça puisqu'on a tout défini exprès pour ça! % [sic] quelconque -> quelconques, agreement with "des changements"

\msp{page 214}
Pour le dire autrement, on a défini \textit{deux} notions de $W$-types d'homotopie relatifs sur $S$, l'une de façon à donner un tel type d'homotopie $\ho^!_{W,S}(X)$ pour \textit{tout} objet de $\cM/S$, mais il n'y a a priori une bonne notion de changement de base que pour les changements de base $S' \longrightarrow S$ qui sont $W$-fibrés, l'autre, qui correspond plus directement à l'intuition d'un $W$-type d'homotopie localement constant sur $S$, n'associe un tel type $\ho^{\lc}_{W,S}(X)$ qu'aux $X$ de $\cM/S$ qui sont $W$-fibrés sur $S$, en revanche cette notion admet une bonne notion d'image inverse et de ‘fibre $W$-homotopique’ pour des changements de base quelconques.

Une ‘bonne’ théorie des fibres $W$-homotopiques doit réunir en une seule les deux notions précédentes, de sorte que les avantages des deux s'ajoutent, et les désavantages s'éliminent! Elle s'obtient à l'aide d'un axiome draconien sur la situation (alors que jusqu'à maintenant on n'a rien utilisé sur $W \subseteq \Fl(\cM)$, si ce n'est tout au plus la saturation faible, en plus de l'existence des produits fibrés) :

\msp{page 215}
\circled{a} (\textit{Axiome de la fibre homotopique.}) Pour tout objet $S$ dans $\cM$, le foncteur canonique
\[
\Ho^{\lc}_W(S) \longrightarrow \Ho^!_W(S)
\]
déduit de l'inclusion
\[
(\Fib_W(S), \Wu_S) \longrightarrow (\cM/S, W_S)
\]
est une \textit{équivalence de catégories.}

Cet axiome implique que dans $\Fib_W S$, $\Wu_S$ et le localiseur induit par $W_S$ (que je désigne encore par $W_S$ par abus de notation) ont même saturé fort. En fait, il est naturel de compléter l'axiome précédent par le suivant.

\circled{b} (Axiome adjoint de l'axiome de la fibre homotopique.) Sur $\Fib_W S$, les localiseurs $\Wu_S$ et $W_S$ coïncident, i.e. on a ceci :

\begin{quote}
Soient $X$, $Y$ dans $\Fib_W S$, $X \xrightarrow{u} Y$ un $S$-morphisme tel que $u \in W_S$, alors $u \in \Wu_S$ (i.e. le morphisme $u' : X' \longrightarrow Y'$ déduit par tout changement de base $S' \longrightarrow S$ est encore dans $W$).
\end{quote}

\msp{page 216}
Je présume que a) sera rarement satisfait sans que b) le soit (et inversement d'ailleurs), et que la démonstration de a) et b) se fera pratiquement en une seule haleine.

Si $\cM = \mathrm{Cat}$, $W$ un localiseur fondamental, alors la notion adaptée ici de $W$-fibration (lire : ‘ultraforte’, dans la terminologie spécifique développée dans Cat) est à tel point forte, qu'elle implique déjà pratiquement b) (moyennant W(4a$'$), il est vrai, cf. p. 53), et
\printedpage{139}
% head: A. Grothendieck : Dérivateurs, Ch. XII, Édition des Universités de Montpellier II et Paris VII
a) n'est guère vérifiable que si on a W(4a) (\footnote[122]{Non, il me semble que W(4b) doit suffire \dots}) (loc. cit.). Donc ici c'est ‘l'axiome des fibrations’ sur $W$, sous sa forme forte W(4a), qui est la clef de voûte d'une bonne théorie des fibres homotopiques.

Une telle théorie (revenant au cas général) revient à dire, pour l'essentiel, qu'il y a ‘assez’ de $W$-fibrations (‘ultrafortes’) pour pouvoir décrire \textit{tout} objet d'un $\cM/S$, à $W_S$-équivalence près, comme un objet $W$-fibrant sur $S$. Dans les cas à ma connaissance, on a de plus ceci : on \msp{page 217} a pour tout $S$ dans $\cM$ une sous-catégorie $\cF_{0\,S}$ de $\Fib_W S$ (je pense par exemple à $\underline{\mathrm{Parf}}_W S$, voire $\underline{\mathrm{Parf}}_{W_\omega} S$, dans le cas où $\cM = \mathrm{Cat}$), telle que pour toute sous-catégorie pleine $\cF$ de $\cM/S$ telle que
\[
\cF_{0\,S} \subseteq \cF \subseteq \cM/S
\]
(a fortiori, pour une $\cF$ pleine telle que $\Fib_W S \subseteq \cF \subseteq \cM/S$), les inclusions induisent des équivalences de catégories
\[
\cF_{0\,S}(W_S)^{-1} \xrightarrow{\sim} \cF(W_S)^{-1} \xrightarrow{\sim} \underbrace{(\cM/S)(W_S)^{-1}}_{\overset{\text{déf}}{=} \mathrm{Hot}^!_W(S)}. % [sic] Hot -> Ho, elsewhere the notation is Ho^!_W(S)
\]
Le choix de $\cF_{0\,S}$ ne semble pas pouvoir se fixer de façon ‘la meilleure possible’, une fois pour toutes dans le contexte axiomatique adopté ici. C'est un élément de structure supplémentaire dans $\cM$ (en plus du localiseur $W$), au même titre que la donnée des classes de $Q$-fibrations et $Q$-cofibrations, à la \textsc{Quillen}. J'ignore si les axiomes a), b) ci-dessus (axiome de la fibre homotopique) impliquent déjà \msp{page 218} la variante de \circled{a} :

\circled{c} Pour toute sous-catégorie pleine $\cF$ de $\cM/S$ telle que
\[
\Fib_W S \subseteq \cF \subseteq \cM/S
\]
les inclusions induisent des équivalences
\[
\underbrace{(\Fib_W S)(W_S)^{-1}}_{\begin{subarray}{l} = \ \Ho^{\lc}_W(S)\\ \text{moyennant a)}\end{subarray}} \xrightarrow{\sim} \cF(W_S)^{-1} \xrightarrow{\sim} \underbrace{(\cM/S)(W_S)^{-1}}_{=\Ho^!_W(S)}.
\]

Mais je présume que pratiquement, quand on pourra vérifier \circled{a}, on pourra vérifier \circled{c} également (\footnote[123]{Cf. page 240 ff. Il faudrait ici développer un peu plus le formalisme de $\Ho^{\lc}_W(S)$ et les fibres homotopiques, sous l'hypothèse \circled{a}\,\circled{b} \ plus haut.}).


\endgroup
% END UNIT Der-ChXIIc-30-03-03/Der-ChXIIc-30-03-03_p021-025.tex
% BEGIN UNIT Der-ChXIIc-30-03-03/Der-ChXIIc-30-03-03_p026-030.tex
\begingroup\makeatletter%
\newcommand{\printedpage}[1]{\marginpar{\footnotesize p.~#1}}%
\newcommand{\msp}[1]{\marginpar{\footnotesize [#1]}}%
\newcommand{\circled}[1]{\tikz[baseline=(c.base)]\node[draw,circle,inner sep=1pt](c){#1};}%
\newcommand{\deff}{\overset{\text{déf}}{=}}%
\makeatother\let\printedpage\derpp\let\msp\dermsp%
\let\mspin\DERmspin%


\printedpage{140}
% head: A. Grothendieck : Dérivateurs, Ch. XII, Édition des Universités de Montpellier II et Paris VII

\section*{5 Axiomatique des $W$-types d'homotopie relatifs non localement constants $\mathrm{Ho}^{\mathcal{L}}_W(S)$.}

On n'a rien dit jusqu'à ce point d'éventuels $W$-types d'homotopie relatifs (sur un objet $S$ de $\mathcal{M}$) qui ne seraient \textit{pas} ``localement constants''. Il en était question pourtant dans loc. cit. (VI, théorème-scholie 2, partie II du théorème pages 109-113.) Je profite de l'occasion pour revenir là dessus. Il faut introduire, dans chaque $\mathcal{M}/S$, un deuxième localiseur
\[
W^0_S \subseteq \mathrm{Fl}(\mathcal{M}/S), \qquad W^0_S \subseteq W_S
\]
(plus fin que $W_S$), ce qui permettra de définir
\msp{page 219}
\[
\mathrm{Ho}_W(S) \deff (\mathcal{M}/S)(W^0_S)^{-1}.
\]
Qu'il n'y ait pas de choix \textit{canonique} qui s'impose ici, déduit de la seule structure $(\mathcal{M},W)$, est claire sur l'exemple de $(\mathrm{Cat},W)$, auquel cas on trouve \textit{deux} choix symétriques, transformés l'un en l'autre par l'involution
\[
X \longmapsto X^\circ : (\mathrm{Cat},W) \xrightarrow{\sim} (\mathrm{Cat},W),
\]
à savoir $W^{\mathrm{g}}_S$ et $W^{\mathrm{d}}_S$. Le choix de l'un ou de l'autre, disons (pour fixer les idées) de $W^{\mathrm{g}}_S$, est associé à une notion d'\textit{isomorphisme local} (ou de `coisomorphisme local' dans le cas de $W^{\mathrm{d}}_S$) que je voudrais axiomatiser. Il semble bien que pour les principales catégories de modèles pour Hot (et pour ses variantes), on dispose de façon assez naturelle d'une telle notion d'isomorphisme local (sans l'embarras, le plus souvent, d'un automorphisme involutif fournissant une notion symétrique \dots).

On supposera donc, comme dans loc. cit. p. 94-96, donné une classe
\[
\mathcal{L} \subseteq \mathrm{Fl}(\mathcal{M})
\]
de morphismes, jounant le rôle de `flèches de % [sic] jounant -> jouant, faute de frappe
\msp{page 220}
localisation' ou `d'isomorphismes locaux', satisfaisant les conditions suivantes.

\begin{enumerate}
\item[$(*)$] $\mathcal{L}$ contient les isomorphismes, est stable par composition, et par changement de base.
\end{enumerate}

On pose, pour $S$ dans $\mathcal{M}$,
\[
\mathcal{L}_S \subseteq \mathcal{M}/S \qquad \text{sous-catégorie pleine des } S'/S \text{ telle que } S' \longrightarrow S \text{ \texttt{[soit]} dans } \mathcal{L}.
\]
On définit alors pour tout $S$ dans $\mathcal{M}$
\[
\begin{aligned}
&W^{\mathcal{L}}_S \subseteq \mathrm{Fl}(\mathcal{M}/S)\\
&(X \xrightarrow{u} Y) \in W^{\mathcal{L}}_S \iff \forall\, S' \in \mathcal{L}_S,\ (u_{S'} : X \times_S S' \longrightarrow Y \times_S S') \text{ est dans } W.
\end{aligned}
\]
% foot: Version 30 mars 2003

\printedpage{141}
% head: A. Grothendieck : Dérivateurs, Ch. XII, Édition des Universités de Montpellier II et Paris VII

C'est donc une variante du localiseur $W^{\mathrm{u}}_S$, où maintenant on ne prend comme changements de base que les flèches `$\mathcal{L}$-localisantes', i.e. $\in \mathcal{L}$. On a donc
\[
W^{\mathrm{u}}_S \subseteq W^{\mathcal{L}}_S \subseteq W_S
\]
en tant que parties de $\mathrm{Fl}(\mathcal{M}/S)$, et de même dans toute sous-catégorie pleine $\mathcal{F}$ de $\mathcal{M}/S$. Dans le cas de $\underline{\mathrm{Fib}}_W S$, moyennant l'axiome \circled{b} ci-dessus ($W^{\mathrm{u}}_S = W_S$ dans $\underline{\mathrm{Fib}}_W S$), les trois localiseurs coïncident.

On pose
\[
\mathrm{Ho}^{\mathcal{L}}_W(S) \text{ (ou } \mathrm{Ho}_W(S)\text{)} \ \deff (\mathcal{M}/S)(W^{\mathcal{L}}_S)^{-1},
\]
et on a donné dans loc. cit. (partie II du théorème-scholie) un résultat intéressant sur cet invariant, pour $S$ fixé. Mon propos
\msp{page 221}
présent est de voir ce qui se passe pour $S$ variable. Visiblement, une notion de $W$-type d'homotopie relatif sur un objet quelconque $S$ de $\mathcal{M}$ n'est satisfaisante, que si elle permet des images inverses par des flèches \textit{quelconques}
\[
f : S' \longrightarrow S
\]
dans $\mathcal{M}$, qui doivent définir absolument
\[
f^{*\mathcal{L}}_W \text{ ou } f^*_W : \mathrm{Ho}^{\mathcal{L}}_W(S) \longrightarrow \mathrm{Ho}^{\mathcal{L}}_W(S').
\]
On est donc confronté à la même perplexité que tantôt - et ce n'est \textit{pas,} cette fois, l'introduction de $\underline{\mathrm{Fib}}_W(S)$ qui va sauver la situation, car
\[
(\underline{\mathrm{Fib}}_W S)(\underbrace{W^{\mathcal{L}}_S}_{\substack{= W_S \\ = W^{\mathrm{u}}_S}})^{-1} \longrightarrow (\mathcal{M}/S)(W^{\mathcal{L}}_S)^{-1}
\]
n'a aucune envie, cette fois, d'être une équivalence de catégories, mais bien plutôt une inclusion (pleinement fidèle) des $W$-types d'homotopie relatifs constants sur $S$, dans les $W$-types d'homotopie quelconques. Il faut visiblement un axiome qui remplace l'axiome `de la fibre homotopique'.

\msp{page 222}

Dans le cas de Cat, la classe $\mathcal{L}$ des isomorphismes locaux (disons) permet de définir, comme ci-dessus, le localiseur $W^{\mathcal{L}}_S = W^{\mathrm{g}}_S$ dans les $\mathrm{Cat}/S$, comme formé des $u : X \longrightarrow Y$ sur $S$ qui sont dans $W$ et le restent par tout changement de base $S' \longrightarrow S$ qui est dans $\mathcal{L}$. Mais il y a lieu de \textit{saturer} $\mathcal{L}$, en regardant la classe $\mathcal{L}^\wedge$ des flèches $f : S \longrightarrow S'$ de Cat telles que le changement de base
\[
f^* : \mathrm{Cat}/S \longrightarrow \mathrm{Cat}/S'
\]
envoie $W^{\mathrm{g}}_S$ dans $W^{\mathrm{g}}_{S'}$. Ce $\mathcal{L}^\wedge$ est tautologiquement stable par composition, et contient les isomorphismes (mieux, contient $\mathcal{L}$ = ensemble des isomorphismes locaux). Il n'est pas
% foot: Version 30 mars 2003
\printedpage{142}
% head: A. Grothendieck : Dérivateurs, Ch. XII, Édition des Universités de Montpellier II et Paris VII
clair, par la définition donnée ici, qu'il soit stable par changement de base - on va renforcer la définition en conséquence, en exigeant, pour que $f : S' \longrightarrow S$ soit dans $\mathcal{L}^\wedge$, que pour tout $T$ sur $S$, considérons $f_T : T' = S' \times_S T \longrightarrow T$, $f^*_T : \mathrm{Cat}/T \longrightarrow \mathrm{Cat}/T'$ soit compatible avec les localiseurs $W^{\mathrm{g}}_T$, $W^{\mathrm{g}}_{T'}$. Alors $\mathcal{L}^\wedge$ est (tout comme $\mathcal{L}$) stable par composition et par changement de base, et de plus il contient $\mathcal{L}$. De plus, pour tout $S$, on a $W^{\mathcal{L}}_S = W^{\mathcal{L}^\wedge}_S$. Ce petit sorite pourrait se faire bien sûr pour tout modéliseur $(\mathcal{M},W)$. Mais la chose intéressante ici, c'est que $\mathcal{L}^\wedge$ contient l'ensemble $\mathrm{Liss}_W$ des flèches $W$-lisses. (J'ignore s'il est égal.) On aurait donc pu
\msp{page 223}
aussi bien partir de l'ensemble de flèches $\mathrm{Liss}_W$, pour définir les localiseurs $W^{\mathrm{g}}_S$ comme
\[
W^{\mathrm{g}}_S = W^{\mathrm{Liss}_W}_S \qquad \text{au lieu de } W^{\mathrm{g}}_S = W^{\mathrm{Isoloc}}_S,
\]
ce qui est équivalent.

D'autre part, en travaillant avec les localisés $\mathrm{HOT}^\circ_W(S) = (\mathrm{Cat}/S)(W^{\mathrm{g}}_S)^{-1}$, le rôle de $\underline{\mathrm{Fib}}_W S$ pour la théorie des $\mathrm{HOT}^{\mathrm{lc}}_W(S)$ est joué ici par $\underline{\mathrm{Prop}}_W(S)$. On voit donc que dans la théorie des $W$-types d'homotopie relatifs non localement constants, mode $\mathrm{HOT}^\circ_W$, les ensembles de flèches $\mathrm{Liss}_W$ et $\mathrm{Prop}_W$ jouent un rôle en quelque sorte complémentaire : les flèches $f : S' \longrightarrow S$ de $\mathrm{Liss}_W$ sont (pratiquement) celles qui permettent directement de définir
\[
f^{*\circ}_W : \mathrm{HOT}^\circ_W(S) \longrightarrow \mathrm{HOT}^\circ_W(S')
\]
comme induit par
\[
f^* : (\mathrm{Cat}/S, W^{\mathrm{g}}_S) \longrightarrow (\mathrm{Cat}/S', W^{\mathrm{g}}_{S'}),
\]
tandis que les flèches de $\mathrm{Liss}_W$ permettent de définir pour tout $S$ dans Cat la sous-catégorie pleine
\[
\underline{\mathrm{Prop}}_W S \subseteq \mathrm{Cat}/S,
\]
donnant lieu à une équivalence de catégories
\[
\left\{
\begin{aligned}
&(\underline{\mathrm{Prop}}_W S)(W^{\mathrm{u}}_S)^{-1} \xrightarrow{\sim} \underbrace{(\mathrm{Cat}/S)(W^{\mathrm{g}}_S)^{-1}}_{= \mathrm{HOT}^0_W(S)} % ?? uncertain: exposant de HOT sous l'accolade -- 0 (zéro) ou \circ
\\
&\text{avec}\\
&W^{\mathrm{u}}_S = W^{\mathrm{g}}_S \text{ sur } \underline{\mathrm{Prop}}_W S
\end{aligned}
\right.
\]
(ce qui visiblement est le substitut annoncé de ``l'axiome de la fibre $W$-homotopique''). Cette
\msp{page 224}
équivalence, qui permet de décrire $\mathrm{Hot}^\circ_W(S)$ en termes de $\underline{\mathrm{Prop}}_W S \subseteq \mathrm{Cat}/S$ et d'un localiseur $W^{\mathrm{u}}_S$ \textit{stable par} \textbf{tout} \textit{changement de base,} permet d'autre part de définir la loi contravariante attendue pour $\mathrm{HOT}^\circ_W(S)$.
% foot: Version 30 mars 2003

\printedpage{143}
% head: A. Grothendieck : Dérivateurs, Ch. XII, Édition des Universités de Montpellier II et Paris VII

D'autre part, quand on échange le rôle de $\mathrm{Liss}_W$ et $\mathrm{Prop}_W$, on trouve la théorie complémentaire du $\mathrm{HOT}_W(S)$,
\[
\left\{
\begin{aligned}
&\mathrm{HOT}_W(S) \deff (\mathrm{Cat}/S)(W^{\mathrm{d}}_S)^{-1} \xleftarrow{\sim} (\underline{\mathrm{Liss}}_W S)(W^{\mathrm{u}}_S)^{-1},\\
&\text{sur } \underline{\mathrm{Liss}}_W S,\ W^{\mathrm{d}}_S = W^{\mathrm{u}}_S,
\end{aligned}
\right.
\]
la description de $\mathrm{HOT}_W(S)$ en termes de $\underline{\mathrm{Liss}}_W S$ et de $W^{\mathrm{u}}_S$ permet d'en définir le caractère contravariant en $S$, pour \textit{toute} flèche $f : S' \longrightarrow S$ dans Cat. Tandis que si on veut pouvoir passer aux catégories de fractions sur $\mathrm{Cat}/S$ et $\mathrm{Cat}/S'$ tout entier (via $W^{\mathrm{d}}_S$, $W^{\mathrm{d}}_{S'}$), il faut se limiter aux changements de base $W$-propres.

Du coup, me vient la perplexité si cette complémentaire, voire cette dualité, entre deux ensembles de flèches, les `lisses' et les `propres', ne serait pas un fait général qui se retrouvait tout au moins dans les catégories de modèles proche de la description de HOT ou d'un $\mathrm{Hot}_W$.

\msp{page 225}

Il se pourrait bien que dans $\boldsymbol{\Delta}^\wedge$ on ait une dualité parfaite entre ces deux espèces de notions (il faudra que j'examine la chose), peut-être même dans toute catégorie $\mathcal{A}^\wedge$ (mais dans $\boldsymbol{\Delta}^\wedge$ il y a un automorphisme involutif $K \longmapsto K^\circ$, correspondant au $X \longmapsto X^\circ$ dans Cat \dots). Et après tout, même dans le contexte des espaces topologiques, les notions de propreté d'une part, d'asphéricité locale d'autre part, pour une flèche $X \longrightarrow Y$, sont bien définies, et pourraient bien s'ajuster comme le font les notions correspondantes dans Cat, et (qui sait) donner lieu à deux versions complémentaires, ou `duales' l'une de l'autre, d'une notion de type d'homotopie relatif.

N'en ayant pas le cœur net, je suis un peu indécis comment présenter une description axiomatique de la théorie des $\mathrm{Ho}^{\mathcal{L}}_W(S)$, pour une catégorie de modèles générale $(\mathcal{M},W)$. Dois-je me donner d'emblée
\[
\mathrm{Liss}, \mathrm{Prop} \subseteq \mathrm{Fl}(\mathcal{M}),
\]
stables par composition, par changement de base, contenant les isomorphismes, et tels que (peut-être)
\[
\mathrm{Liss} \cap \mathrm{Prop} \subseteq \mathrm{Fib}_W,
\]
\msp{page 226}
mais surtout ils permettent de définir, pour $S$ de $\mathcal{M}$
\[
\left\{
\begin{array}{lcccl}
W^{\mathrm{g}}_S & \deff & W^{\mathrm{Liss}}_S & \subseteq & \mathcal{F}(\mathcal{M}/S)\\
W^{\mathrm{d}}_S & \deff & W^{\mathrm{Prop}}_S & \subseteq & \mathcal{F}(\mathcal{M}/S)
\end{array}
\right.
\]
avec les propriétés suivantes :

\begin{enumerate}
\item[{\circled{a}}] Si $f : S' \longrightarrow S$ est dans Prop (resp. Liss), alors $f^* : \mathcal{M}/S \longrightarrow \mathcal{M}/S'$ transforme $W^{\mathrm{d}}_S$ dans $W^{\mathrm{d}}_{S'}$ (resp. $W^{\mathrm{g}}_S$ dans $W^{\mathrm{g}}_{S'}$),
\end{enumerate}
% foot: Version 30 mars 2003
\printedpage{144}
% head: A. Grothendieck : Dérivateurs, Ch. XII, Édition des Universités de Montpellier II et Paris VII
\noindent ce qui est une pure tautologie, et ne concerne pas encore les relations de Prop avec Liss, mais par contre :

\begin{enumerate}
\item[{\circled{b}}]
\[
\left\{
\begin{aligned}
&(\underline{\mathrm{Liss}}S)(W^{\mathrm{u}}_S)^{-1} \xrightarrow{\sim} (\mathcal{M}/S)(W^{\mathrm{d}}_S)^{-1} \text{ \texttt{[est]} une équivalence}\\
&\text{et } W^{\mathrm{u}}_S = W^{\mathrm{d}}_S \text{ sur } \underline{\mathrm{Liss}}S,
\end{aligned}
\right.
\]
et dualement
\[
\left\{
\begin{aligned}
&(\underline{\mathrm{Prop}}S)(W^{\mathrm{u}}_S)^{-1} \xrightarrow{\sim} (\mathcal{M}/S)(W^{\mathrm{g}}_S)^{-1} \text{ \texttt{[est]} une équivalence}\\
&\text{et } W^{\mathrm{u}}_S = W^{\mathrm{g}}_S \text{ sur } \underline{\mathrm{Prop}}S.
\end{aligned}
\right.
\]
\end{enumerate}

De plus, il y a lieu d'énoncer \circled{c} deux propriétés d'adjonction (duales l'une de l'autre), pour lesquelles je renvoie à VI, théorème-scholie 2, p. 108-113. Il se pose ici la question subsidiaire de simplifier les conditions de validité de la formule d'adjonction, un peu hétéroclites dans loc. cit.

Il se pose d'autre part la question de faire la théorie des lois covariantes $f_!$ et $f_*$ pour $\mathrm{Ho}^{\mathrm{lc}}_W(S)$ et pour $\mathrm{Ho}^{\mathrm{Prop}}_W(S)$, $\mathrm{Ho}^{\mathrm{Liss}}_W(S)$.

\msp{page 227}

J'ai bon espoir que le foncteur $f_!$ peut être décrit, et son adjonction avec $f^*$ établie, à aussi bon compte dans ce contexte général que dans Cat. Mais je n'ai pas envie de m'y arrêter à présent.

\section*{6 Cofibre $W$-homotopique dans $\mathrm{Cat}$.}

Par contre, j'ai envie de regarder ici sous quelles conditions sur le localiseur fondamental $W$, il existe dans $(\mathrm{Cat},W)$ une théorie duale de celle de la fibre $W$-homotopique.

La question-clef est, bien sûr, celle d'avoir une bonne prise sur les $W$-cofibrations. Or je ne connais jusqu'à présent \textit{que} les immersions ouvertes de Dwyer, les immersions fermées de Dwyer, et ce qui s'en déduit par composition etc. Le travail de \textsc{Thomason} peut motiver l'espoir que la connaissance de ces immersions-là, qui semblent ici jouer un rôle similaire (dual) à celui des morphismes $W$-parfaits, suffira pour faire une bonne théorie des `cofibres $W$-homotopiques'.

% trois traits verticaux imprimés dans la marge gauche de ce paragraphe
{\reversemarginpar\marginnote{\textbar\textbar\textbar}}Je suppose dans toute cette section que $W$ satisfait W(7 bis) (axiome des carrés cocartésiens), pour avoir des bonnes propriétés des immersions ouvertes de Dwyer (\footnote[124]{\begin{enumerate}
\item[a)] Ce sont des $W$-cofibrations (résulte de b) et c)).
\item[b)] Si
\end{enumerate}}), et plus généralement, pour les $W$-cofibrations. % ?? uncertain: coupure en fin de ligne "co-cartésiens" -- cocartésiens ou co-cartésiens
% foot: Version 30 mars 2003


\endgroup
% END UNIT Der-ChXIIc-30-03-03/Der-ChXIIc-30-03-03_p026-030.tex
% BEGIN UNIT Der-ChXIIc-30-03-03/Der-ChXIIc-30-03-03_p031-035.tex
\begingroup\makeatletter%
\newcommand{\printedpage}[1]{\marginpar{\footnotesize p.~#1}}%
\newcommand{\msp}[1]{\marginpar{\footnotesize [#1]}}%
\makeatother\let\printedpage\derpp\let\msp\dermsp%
\let\mspin\DERmspin%


\printedpage{145}
% head: A. Grothendieck : Dérivateurs, Ch. XII, Édition des Universités de Montpellier II et Paris VII
% foot: Version 30 mars 2003

\msp{page 228}
Il se pose donc la question de la factorisation duale d'une flèche $f : X \longrightarrow Y$ de Cat en
\[
\begin{tikzcd}[column sep=huge]
X \arrow[r, hook, "W\text{-cofibration}", "i_f"'] & \tilde{Y} \arrow[r, "\in W", "p_f"'] & Y.
\end{tikzcd}
\]

Les résultats de \textsc{Thomason} assurent que ça existe, avec $i$ une immersion ouverte de Dwyer, et $p$ une $W$-fibration - du moins dans le cas $W = W_\infty$. Je vais essayer de le voir directement, en calquant dualement la construction de \textsc{Cartan-Serre}.

Il s'agit donc d'abord de factoriser
\[
X \amalg X \xrightarrow{\ \mathrm{codiag}_X\ } X,
\]
et pour ceci je choisis un intervalle fixe
\[
I \ = \ \Phi_n \ = \ (0 \longrightarrow 1 \longleftarrow 2 \longrightarrow \cdots \longrightarrow 2n-1 \longleftarrow 2n)
\]
avec
\[
n \geq 2
\]
de sorte que l'inclusion
\[
\varepsilon = \text{catégorie discrète } \{0,1\} \xrightarrow{\partial} I, \qquad \partial(0) = 0, \quad \partial(1) = 2n,
\]
soit une immersion ouverte de Dwyer. Cette propriété est stable par changement de base Cat-fibrant, en particulier en prenant un produit
\[
\underbrace{X \times \varepsilon}_{\simeq X \coprod X} \xrightarrow{i_X} X \times I,
\]
% continuation of a footnote begun on the preceding page of the edition (no footnote mark on this page)
\begingroup\renewcommand{\thefootnote}{}\footnotetext{%
\[
\begin{tikzcd}[ampersand replacement=\&]
U \arrow[r, hook, "i"] \arrow[d, "p"'] \& X \arrow[d, "q"] \\
U' \arrow[r, hook, "i'"] \& X'
\end{tikzcd}
\]
avec $i$ immersion ouverte de Dwyer, alors
\[
\begin{aligned}
i \in W \quad &\Longleftrightarrow \quad i' \in W \\
p \in W \quad &\Longleftrightarrow \quad q \in W,
\end{aligned}
\]
j'ai donc besoin de W(7 bis) \textit{et} même de W(7) pour la deuxième implication.
\begin{enumerate}
\item[c)] Les immersions ouvertes de Dwyer \texttt{[sont]} stables par cochangement de base et par composition, contiennent les isomorphismes.
\end{enumerate}}\endgroup
\printedpage{146}
% head: A. Grothendieck : Dérivateurs, Ch. XII, Édition des Universités de Montpellier II et Paris VII
% foot: Version 30 mars 2003
\msp{page 229}
d'autre part
\[
X \times I \xrightarrow{\mathrm{pr}_1} X
\]
est $W_\omega$-parfait, d'où la factoriastion cherchée % [sic] factoriastion -> factorisation, misprint
\[
X \amalg X \xrightarrow{i_X} X \times I \ \xrightarrow{\ p_X = \mathrm{pr}_1^X\ } X
\]
- il n'y a pas plus simple! Comment l'utiliser pour construire la factorisation d'une flèche $f : X \longrightarrow Y$ quelconque, en dualisant les constructions relatives aux espaces de chemins? On définit $Z(f)$ par le carré cocartésien dans Cat
\[
\begin{tikzcd}[column sep=5em, row sep=4em]
X \arrow[r, dotted, "\mathrm{inj}_1^X"'] & X \times I \arrow[r, "f'"] & Z(f) \arrow[d, dotted, bend left=45, "\,p_f"] \\
& X \arrow[u, hook, "\mathrm{inj}_0^X"'] \arrow[r, "f"] & Y, \arrow[u, hook, "W"']
\end{tikzcd}
\]
on a
\[
Z(f) \xrightarrow{\alpha} Y \qquad \text{induit par } \mathrm{pr}_1 : X \times I \longrightarrow X,
\]
et la factorisation cherchée de $f$ est
\[
X \ \xrightarrow{\ i_f = f' \circ \mathrm{inj}_1^X\ } Z(f) \xrightarrow{p_f} Y.
\]

Je dis que
\[
\left\{
\begin{aligned}
& i_f \text{ est immersion ouverte de Dwyer, donc une } W\text{-cofibration.} \\
& p_f \in W.
\end{aligned}
\right.
\]

\msp{page 230}
On voit sur le diagramme plus complet
\[
\begin{tikzcd}[column sep=6.5em, row sep=4.5em]
& X \times I \arrow[r, "f'"] \arrow[dd, dotted, bend left=45, "\mathrm{pr}_1"{pos=.60}] \arrow[dr, phantom, "\substack{\text{cocart.}\\ \text{et cart.}}"] & Z(f) \arrow[dd, dotted, bend left=60, "p_f"] \\
X \arrow[ur, dotted, "\mathrm{inj}_1^X"] \arrow[r, hook, "i_1^X"] & X \amalg X \arrow[u, hook, "i_X"'] \arrow[r, "f \coprod \mathrm{id}_X"] \arrow[dr, phantom, "\substack{\text{cocart.}\\ \text{et cart.}}"] & Y \amalg X \arrow[u, hook, "j"'] \\
& X \arrow[u, hook, "i_0^X"'] \arrow[r, "f"] & Y, \arrow[u, hook]
\end{tikzcd}
\]
\printedpage{147}
% head: A. Grothendieck : Dérivateurs, Ch. XII, Édition des Universités de Montpellier II et Paris VII
% foot: Version 30 mars 2003
\[
i_f = f' \circ \mathrm{inj}_1^X = f' i_X i_1^X = j \ \ \underbrace{(f \amalg \mathrm{id}_X) i_1^X}_{\substack{\text{inclusion canonique}\\ \text{de } X \text{ dans}\\ Y \coprod X}} \ \ ,
\]
c'est \texttt{[une]} immersion ouverte de Dwyer, comme composé de $Y \longrightarrow Y \amalg X$ et $j : Y \amalg X \longrightarrow Z(f)$ qui le sont ($j$ obtenu par cochangement de base à partir de $i_X$, qui est immersion ouverte de Dwyer).

D'autre part, $X \xrightarrow{\mathrm{inj}_0^X} X \times I$ est \texttt{[une]} immersion ouverte de Dwyer et dans $W$, donc coüniversellement dans $W$, donc $Y \longrightarrow Z(f)$ est elle aussi dans $W$, donc aussi la rétraction $p_f$.

D'autre part, la proposition duale de
\[
\underline{\mathrm{Fib}}_W S \cap W_S = W^{\mathrm{u}}_{S\,\mathrm{univ}}
\]
est valable pour les $W$-cofibrations :

\msp{page 231}
\textbf{Proposition 16} (\footnote[125]{N'a rien à voir avec Cat - une chose générale quand il y a assez de $W$-cofibrations.}). \textit{Sur la sous-catégorie pleine $\underline{\mathrm{Cofib}}_W(S)$ de $S\backslash\mathrm{Cat}$, le localiseur ${}_SW$ est égal au $({}_SW)^{\mathrm{univ}}$ ($W$-équivalences couniverselles), i.e. on a ceci : Soient}
\[
S \longrightarrow X, \qquad S \longrightarrow Y
\]
\textit{des $W$-cofibrations, et soit $f : X \longrightarrow Y$ un morphisme de $S\backslash\mathrm{Cat}$ tel que $f \in W$. Alors pour tout cochangement de base $S \longrightarrow S'$, la flèche correspondante}
\[
f' : X' = X \amalg_S S' \longrightarrow Y' = Y \amalg_S S'
\]
\textit{est dans $W$.}

\textsc{Démonstration}. On factorise $S \longrightarrow S'$ en
\[
S \xrightarrow{i} S'_0 \xrightarrow{j} S',
\]
avec $i$ immersion ouverte de Dwyer (donc $W$-cofibration), et $j \in W$. On a donc le diagramme à carrés cocartésiens
\[
\begin{tikzcd}
S \arrow[r] \arrow[d, hook, "i"'] & X \arrow[r, "f"] \arrow[d, hook] & Y \arrow[d, hook] \\
S'_0 \arrow[r] \arrow[d, "\in W"'] & X'_0 \arrow[r, "f'_0", "W"'] \arrow[d, "W"] & Y'_0 \arrow[d, "W"] \\
S' \arrow[r] & X' \arrow[r, "f'"] & Y'.
\end{tikzcd}
\]
Comme $i$ est une $W$-cofibration et $f \in W$, on a $f'_0 \in W$. D'autre part, comme $S \longrightarrow X$, donc $S'_0 \longrightarrow X'_0$ est
\printedpage{148}
% head: A. Grothendieck : Dérivateurs, Ch. XII, Édition des Universités de Montpellier II et Paris VII
% foot: Version 30 mars 2003
\msp{page 232}
une $W$-cofibration, $X'_0 \longrightarrow X'$ est dans $W$, et pour la même raison, $Y'_0 \longrightarrow Y$ l'est aussi. Donc $f'$ l'est aussi. % [sic] Y'_0 \longrightarrow Y -> Y'_0 \longrightarrow Y', the diagram has Y'_0 \to Y'

\textbf{NB} Ici, c'est un argument formel, valable (sous forme duale) dans tout $(\mathcal{M}, W)$ où on a un énoncé de factorisation pour une flèche quelconque $f : X \longrightarrow Y$ en $X \xrightarrow{i_f} \bar{X} \xrightarrow{\alpha} Y$, avec $i_f \in W$ et $\alpha$ une $W$-fibration. Alors, si $X$, $Y$ sont $W$-fibrés sur $S$, toute $S$-flèche de $X$ dans $Y$ qui est dans $W$, est dans $W^{\mathrm{u}}_S$. Et dualement, si toute flèche $X \longrightarrow Y$ se factorise en $X \longrightarrow \tilde{Y} \longrightarrow Y$, avec $X \longrightarrow \tilde{Y}$ une $W$-cofibration et $\tilde{Y} \longrightarrow Y$ dans $W$.

Il me vient ici une (petite) inquiétude : ai-je au moins prouvé que si
\[
S \longrightarrow X
\]
est une $W$-cofibration, alors pour toute flèche $S' \longrightarrow S$ dans Cat, le carré cocartésien dans Cat % [sic] S' \longrightarrow S -> S \longrightarrow S', the square below has S \to S'
\[
\begin{tikzcd}
S \arrow[r] \arrow[d] & X \arrow[d] \\
S' \arrow[r] & X'
\end{tikzcd}
\]
est aussi $W$-cocartésien (\footnote[126]{dans le sens ancien, spécial à Cat (déf. 3, p. 131). Pour le comparer avec la théorie générale, cf. prop. 20, p. 262.})? Autrement,
\msp{page 233}
à quoi ce carré-là serait-il bon?! Pour le prouver, factorisons encore $S \longrightarrow S'$ en
\[
\underset{\substack{\text{immersion}\\ \text{ouverte de}\\ \text{Dwyer}}}{S \hookrightarrow} S'_0 \underset{\in W}{\longrightarrow} S',
\]
d'où le diagramme de 2 carrés cocartésiens
\[
\begin{tikzcd}
S \arrow[r, "\mathrm{Cof}_W"] \arrow[d, hook, "D_W"'] & X \arrow[d, hook, "D_W"] \\
S'_0 \arrow[r, "\mathrm{Cof}_W"] \arrow[d, "\in W"'] & X'_0 \arrow[d, "\in W"] \\
S' \arrow[r, "\mathrm{Cof}_W"] & X',
\end{tikzcd}
\]
où $X'_0 \longrightarrow X'$ est $\in W$ d'après l'hypothèse sur $X$ sous $S$. Il faut donc vérifier :
\printedpage{149}
% head: A. Grothendieck : Dérivateurs, Ch. XII, Édition des Universités de Montpellier II et Paris VII
% foot: Version 30 mars 2003
\begin{enumerate}
\item[a)] Un carré cocartésien construit sur une flèche qui est \texttt{[une]} immersion de Dwyer, est $W$-cocartésien.
\item[b)] Un carré commutatif dans lequel deux côtés opposés sont dans $W$, est $W$-cocartésien.
\item[c)] Le composé de deux carrés $W$-cocartésiens est $W$-cocartésien.
\end{enumerate}

Admettons provisoirement b) et c), et prouvons a), pour un carré cocartésien
\[
\begin{tikzcd}
U \arrow[r, hook, "i"] \arrow[d, "p"'] & X \arrow[d, "q"] \\
U' \arrow[r, hook, "i'"] & X',
\end{tikzcd}
\]
où $i$ est immersion ouverte de Dwyer. On procède encore comme tantôt, en factorisant $p$ en $U \overset{j}{\hookrightarrow} U'_0 \xrightarrow{p'} U'$, où $j$ est \texttt{[une]} immersion ouverte de Dwyer, $p' \in W$. On a donc un diagramme de deux carrés cocartésiens accolés
\msp{page 234}
\[
\begin{tikzcd}
U \arrow[r, hook, "i \in D_W"] \arrow[d, hook, "D_W"'] & X \arrow[d, hook] \\
U'_0 \arrow[r, hook, "D_W"] \arrow[d, "p' \in W"'] & X'_0 \arrow[d, "q' \in W"] \\
U' \arrow[r, hook] & X',
\end{tikzcd}
\]
où le premier carré est formé d'immersions ouverts. On sait alors qu'il est $W$-cocartésien. Dans le deuxième, on a $p' \in W$, donc $q' \in W$ puisque $i$ est une immersion de Dwyer, donc une $W$-cofibration. Donc il est également $W$-cocartésien. On conclut alors grâce à c) ci-dessus. % [sic] immersions ouverts -> immersions ouvertes, agreement

\noindent Prouvons b), pour
\begin{tikzcd}[baseline=(\tikzcdmatrixname-1-1.base)]
X_0 \arrow[r, "i_1"] \arrow[d, "i_2 \in W"'] & X_1, \arrow[d, "f_1 \in W"] \\
X_2 \arrow[r, "f_2"] & Y
\end{tikzcd}
en l'insérant dans
\begin{tikzcd}[baseline=(\tikzcdmatrixname-1-1.base)]
X_0 \arrow[r, "i_1"] \arrow[d, "i_2 \in W"'] & X_1, \arrow[d, "j_1"] \arrow[ddr, "f_1"] & \\
X_2 \arrow[r, "j_2"] \arrow[drr, "f_2"'] & \mathcal{X} \arrow[dr, "\varphi"] & \\
& & Y
\end{tikzcd}
où $\mathcal{X} = \displaystyle\int X_0$
\begin{tikzcd}[baseline=(\tikzcdmatrixname-2-1.base), row sep=small]
& X_1 \\
{} \arrow[ur] \arrow[dr] & \\
& X_2
\end{tikzcd}
.

Comme $i_2 \in W$, on sait qu'il en résulte $j_1 \in W$, et comme $f_1 \in W$, il en résulte que $\varphi \in W$, q.e.d.


\endgroup
% END UNIT Der-ChXIIc-30-03-03/Der-ChXIIc-30-03-03_p031-035.tex
% BEGIN UNIT Der-ChXIIc-30-03-03/Der-ChXIIc-30-03-03_p036-040.tex
\begingroup\makeatletter%
\newcommand{\printedpage}[1]{\marginpar{\footnotesize p.~#1}}%
\newcommand{\msp}[1]{\marginpar{\footnotesize [#1]}}%
\newcommand{\cX}{\mathscr{X}}%
\newcommand{\Ob}{\operatorname{Ob}}%
\makeatother\let\printedpage\derpp\let\msp\dermsp%
\let\mspin\DERmspin%


\printedpage{150}
% head: A. Grothendieck : Dérivateurs, Ch. XII, Édition des Universités de Montpellier II et Paris VII
% foot: Version 30 mars 2003

Reste à prouver c). Je vais le déduire d'une

\textbf{Proposition 17.} \textit{Soit un carré commutatif}
\[
Q \quad = \quad
\begin{tikzcd}[ampersand replacement=\&]
X_0 \arrow[r, "i_1"] \arrow[d, "i_2"'] \& X_1 \arrow[d, "f_1"] \\
X_2 \arrow[r, "f_2"] \& Y
\end{tikzcd}
\]
\textit{dans} Cat\textit{, et $\cX_Q$ la catégorie fibrée sur}
\[
\tilde{\Phi} \quad (\simeq \boldsymbol{\Delta}^1 \times \boldsymbol{\Delta}^1) \quad = \quad
\begin{tikzcd}[ampersand replacement=\&]
0 \& 1 \arrow[l] \\
2 \arrow[u] \& 3 \arrow[l] \arrow[u]
\end{tikzcd}
\]
\textit{qui l'intègre. Soit $\cX$ l'image inverse du fermé}
\[
\begin{tikzcd}[ampersand replacement=\&]
0 \& 1 \arrow[l] \\
2 \arrow[u] \&
\end{tikzcd}
\]
\msp{page 235}
\textit{dans $\cX_Q$, qui intègre le diagramme}
\[
\begin{tikzcd}[ampersand replacement=\&, row sep=small]
 \& X_1 \\
X_0 \arrow[ru] \arrow[rd] \& \\
 \& X_2
\end{tikzcd}
\]
\textit{. Pour que $Q$ soit $W$-cocartésien, il faut et il suffit que l'inclusion}
\[
i : \cX \longrightarrow \cX_Q
\]
\textit{soit dans $W$.}

Comme $\cX_Q$ est lisse sur $\tilde{\Phi}$, et 3 est objet initial de $\tilde{\Phi}$, il s'ensuit que l'inclusion de la fibre de $\cX_Q$ en le sommet 3 de $\tilde{\Phi}$ est dans $W$,
\[
Y \xrightarrow{i_Y} \cX_Q \qquad\qquad \in W.
\]

Considérons alors le diagramme
\[
\begin{tikzcd}[ampersand replacement=\&, column sep=small]
\cX \arrow[rr, "\varphi_Q"] \arrow[rd, "i"', ""{name=direct,pos=.55,above right}] \& \& Y \arrow[ld, "i_Y\in W", ""{name=via,pos=.30,below left}] \\
 \& \cX_Q. \& \arrow[from=via,to=direct,dotted,shorten <=1pt,shorten >=1pt,"\alpha"']
\end{tikzcd}
\]

\printedpage{151}
% head: A. Grothendieck : Dérivateurs, Ch. XII, Édition des Universités de Montpellier II et Paris VII
% foot: Version 30 mars 2003

Ce diagramme n'est pas commutatif, mais il y a une flèche de commutativité
\[
\alpha : i_Y\varphi_Q \longrightarrow i.
\]

On sait que ceci implique que
\[
i \in W \qquad\qquad \Longleftrightarrow \qquad\qquad (i_Y\varphi_Q) \in W
\]
(cf. lemme p. 2), et comme $i_Y \in W$,
\[
(i_Y\varphi_Q) \in W \qquad\qquad \Longleftrightarrow \qquad\qquad \varphi_Q \in W,
\]
d'où
\[
i \in W \qquad\qquad \Longleftrightarrow \qquad\qquad \varphi_Q \in W,
\]

\hfill q.e.d.

Nous pouvons maintenant prouver

\msp{page 236}

\textbf{Proposition 18} (\footnote[127]{devrait être vrai dans toute catégorie de modèles ayant assez de cofibrations.}). (Moyennant W(7 bis).) \textit{Considérons un diagramme commutatif}
\[
\begin{tikzcd}[ampersand replacement=\&]
X_0 \arrow[r] \arrow[d] \arrow[rd, phantom, "1"] \& X_1 \arrow[r] \arrow[d] \arrow[rd, phantom, "2"] \& X_2 \arrow[d] \\
X'_0 \arrow[r] \& X'_1 \arrow[r] \& X'_2
\end{tikzcd}
\tag*{$(\ast)$}
\]
\textit{d'où un carré composé}
\[
\begin{tikzcd}[ampersand replacement=\&]
X_0 \arrow[r] \arrow[d] \arrow[rd, phantom, "3"] \& X_2 \arrow[d] \\
X'_0 \arrow[r] \& X'_2.
\end{tikzcd}
\]
\textit{Supposons que le carré} 1 \textit{soit $W$-cocartésien. Alors le carré} 2 \textit{est $W$-cocartésien si et seulement si le carré} 3 \textit{l'est.}

\textsc{Démonstration.}

\texttt{[Maintenant, }$X_i$\texttt{ et }$X'_i$\texttt{ échangent leurs rôles.]}

\printedpage{152}
% head: A. Grothendieck : Dérivateurs, Ch. XII, Édition des Universités de Montpellier II et Paris VII
% foot: Version 30 mars 2003

Soit
\[
I \quad = \quad
\left(
\begin{tikzcd}[ampersand replacement=\&]
0 \arrow[d] \& 1 \arrow[l] \arrow[d] \& 2 \arrow[l] \arrow[d] \\
0' \& 1' \arrow[l] \& 2' \arrow[l]
\end{tikzcd}
\right)
\]
et soit $\cX$ l'intégrale du diagramme $(\ast)$, catégorie fibrée sur $\cX$. On désigne par $\cX_{010'}$ la sous-catégorie pleine de $\cX$ image inverse de la sous-catégorie pleine % [sic] fibrée sur $\cX$ -> fibrée sur $I$, a category fibred over itself; the base is the category $I$ just defined
\[
\underbrace{I_{010'}}_{=\{0,1,0'\}} \ \subseteq I,
\]
et de même pour toute autre partie $A$ de $\Ob I = \{0,1,2,0',1',2'\}$, donnant une sous-catégorie pleine $\cX_A$ de $\cX$. On pose
\[
U \quad = \quad
\begin{tikzcd}[ampersand replacement=\&]
0 \arrow[d] \& 1 \arrow[l] \& 2 \arrow[l] \\
0' \& \&
\end{tikzcd}
, \qquad
V \quad = \quad
\begin{tikzcd}[ampersand replacement=\&]
0 \arrow[d] \& 1 \arrow[l] \arrow[d] \& 2 \arrow[l] \\
0' \& 1' \arrow[l] \&
\end{tikzcd}
,
\]
d'où
\[
\cX_U = \cX_{0120'}, \qquad\qquad \cX_V = \cX_{0120'1'}.
\]
Les flèches entre les $\cX_A$ sont les inclusions

\msp{page 237}

déduites d'inclusions $A \hookrightarrow A'$ dans $\mathfrak{P}(\Ob I)$. Ainsi, le caractère $W$-cocartésien des carrés 1, 2, 3 s'exprime respectivement par la condition que les inclusions
\[
\begin{aligned}
\cX_{010'} \quad &\xrightarrow{\varphi_1} \quad \cX_{010'1'} \\
\cX_{121'} \quad &\xrightarrow{\varphi_2} \quad \cX_{121'2'} \\
\cX_{020'} \quad &\xrightarrow{\varphi_1} \quad \cX_{020'2'} % [sic] $\varphi_1$ -> $\varphi_3$, this inclusion is called $\varphi_3$ in the diagram that follows
\end{aligned}
\]
soient dans $W$. Considérons alors le diagramme de flèches d'inclusion
\printedpage{153}
% head: A. Grothendieck : Dérivateurs, Ch. XII, Édition des Universités de Montpellier II et Paris VII
% foot: Version 30 mars 2003
\[
\begin{tikzcd}[ampersand replacement=\&]
\cX_{020'} \arrow[r, "\varphi_3"] \arrow[d, "\alpha"'] \& \cX_{020'2'} \arrow[d, "\beta"] \\
\cX_U \arrow[r, "\psi_3"] \arrow[d, "\psi_1"'] \& \cX \arrow[d, "\mathrm{id}"] \\
\cX_V \arrow[r, "\psi_2"] \& \cX \\
\cX_{121'} \arrow[r, "\varphi_2"] \arrow[u, "\alpha'"] \& \cX_{121'2'}. \arrow[u, "\beta'"']
\end{tikzcd}
\]

Je dis que dans ce diagramme on a
\[
\alpha, \beta, \alpha', \beta' \in W
\tag*{$(\ast)$}
\]
et
\[
\varphi_1 \in W \ \ \text{(i.e. carré 1 $W$-cocartésien)} \qquad\qquad \Longleftrightarrow \qquad\qquad \psi_1 \in W.
\tag*{$(\ast\ast)$}
\]

Les relations $(\ast)$ impliquent déja que % [sic] déja -> déjà, missing grave accent
\[
\begin{aligned}
\varphi_2 \in W \quad &\Longleftrightarrow \quad \psi_2 \in W \\
\varphi_3 \in W \quad &\Longleftrightarrow \quad \psi_3 \in W,
\end{aligned}
\]
et sous l'hypothèse $W$-cocartésienne $(\ast\ast)$, impliquant $\psi_1 \in W$,

\msp{page 238}

on trouve
\[
\psi_2 \in W \qquad\qquad \Longleftrightarrow \qquad\qquad \psi_3 \in W,
\]
d'où
\[
\varphi_2 \in W \qquad\qquad \Longleftrightarrow \qquad\qquad \varphi_3 \in W,
\]
ce qu'il fallait prouver.

Il reste à établir $(\ast)$ et $(\ast\ast)$. Les relations $(\ast)$ proviennent de la remarque que si $i : A \longrightarrow A'$ est une inclusion de sous-catégories pleines de $I$, telle que $i$ soit $W$-\textit{asphérique} (i.e. $A/j$ $W$-asphérique pour $j \in \Ob A'\backslash\Ob A$), alors $\cX_A \longrightarrow \cX_{A'}$ l'est aussi (changement de base lisse), donc est \texttt{[dans]} $W$. On doit vérifier la $W$-asphéricité pour les quatre inclusions
\[
\begin{array}{rcll}
\{0,2,0'\} & \longrightarrow & U = \{0,1,2,0'\} & \qquad\qquad \boldsymbol{\Delta}^0 \\[1ex]
\{0,2,0',2'\} & \longrightarrow & I & \qquad\qquad \boldsymbol{\Delta}^0/\boldsymbol{\Delta}^1 \\[1ex]
\{1,2,1'\} & \longrightarrow & V = \{0,1,2,0',1'\} & \qquad\qquad \boldsymbol{\Delta}^1,\ \Phi \\[1ex]
\{1,2,1',2'\} & \longrightarrow & I & \qquad\qquad \boldsymbol{\Delta}^1,\ \boldsymbol{\Delta}^1 \times \boldsymbol{\Delta}^1.
\end{array}
\]

\printedpage{154}
% head: A. Grothendieck : Dérivateurs, Ch. XII, Édition des Universités de Montpellier II et Paris VII
% foot: Version 30 mars 2003

(NB J'ai marqué en marge les catégories qu'on obtient comme $A/j$, à isomorphisme près, dans les quatre cas envisagés - elles ont toutes un objet initial, donc sont $W$-asphériques.)

Reste à prouver $(\ast\ast)$. Soit
\[
U' = \{0,1,0',1'\}, \qquad\qquad \text{donc } V = U \cup U',\ \ U \cap U' = \{0,1,0'\},
\]
où $U$, $U'$ sont des \textit{fermés} de $V$ qui recouvrent $V$. Donc on aura de même
\[
\cX_V = \cX_U \cup \cX_{U'}, \qquad\qquad \cX_U \cap \cX_{U'} = \cX_{010'},
\]

\msp{page 239}

\[
\begin{tikzcd}[ampersand replacement=\&]
\cX_{U\cap U'} \arrow[r] \arrow[d, "\varphi_1"] \& \cX_U \arrow[d, "\psi_1"] \\
\cX_{U'} \arrow[r] \& \cX_V.
\end{tikzcd}
\]

Donc par W(7 bis) (cas dual des immersions fermées), on a
\[
\varphi_1 \in W \qquad\qquad \Longleftrightarrow \qquad\qquad \psi_1 \in W,
\]

\hfill q.e.d.

Notons pour mémoire le corollaire qui nous avait motivé.

\textbf{Corollaire.} \textit{Supposons} W(7 bis) (comme dans toute cette section). \textit{Soit un carré cocartésien dans} Cat
\[
\begin{tikzcd}[ampersand replacement=\&]
X_0 \arrow[r, "i_1"] \arrow[d, "i_2"'] \& X_1 \arrow[d] \\
X_2 \arrow[r] \& X_3,
\end{tikzcd}
\]
\textit{avec $i_1$ (ou $i_2$) une $W$-cofibration (p.ex. une immersion de Dwyer). Alors ce carré est $W$-cocartésien.}

Cette digression (p. 232 ff.) n'est pas nécessaire d'ailleurs pour formuler dès maintenant la conclusion de nos cogitations sur les cofibres $W$-homotopiques : \textit{Il y a une telle théorie dans} Cat, pourvu seulement que $W$ satisfasse à W(7 bis). En effet, tout résulte du fait que toute flèche $f : X \longrightarrow Y$ dans Cat se factorise en
\[
X \xrightarrow{i_f} \tilde{Y} \xrightarrow{p_f} Y,
\]
avec $i_f$ une $W$-cofibration (et même une imm-


\endgroup
% END UNIT Der-ChXIIc-30-03-03/Der-ChXIIc-30-03-03_p036-040.tex
% BEGIN UNIT Der-ChXIIc-30-03-03/Der-ChXIIc-30-03-03_p041-045.tex
\begingroup\makeatletter%
\derleqno%
\newcommand{\printedpage}[1]{\marginpar{\footnotesize p.~#1}}%
\newcommand{\msp}[1]{\marginpar{\footnotesize [#1]}}%
\newcommand{\cercle}[1]{\tikz[baseline=(c.base)]\node[draw,circle,inner sep=1pt](c){#1};}%
\newcommand{\Fibu}{\underline{\mathrm{Fib}}}%
\newcommand{\Homu}{\underline{\mathrm{Hom}}}%
\newcommand{\Ho}{\mathrm{Ho}}%
\newcommand{\ho}{\mathrm{ho}}%
\newcommand{\lc}{\mathrm{lc}}%
\makeatother\let\printedpage\derpp\let\msp\dermsp%
\let\mspin\DERmspin%


\printedpage{155}
% head: A. Grothendieck : Dérivateurs, Ch. XII, Édition des Universités de Montpellier II et Paris VII
% foot: Version 30 mars 2003

\msp{page 240}
ersion ouverte de Dwyer), et $p_f \in W$. % beginning of a word continued from the preceding page (outside this range), as printed

\section*{7\quad Retour sur l'axiomatique de la fibre $W$-homotopique.}

Je reviens sur la réflexion du n$^\circ$ 4. Finalement, pour
\[
(\mathcal{M}, W \subseteq \mathrm{Fl}(\mathcal{M}))
\]
donnés avec $W$ faiblement saturé, il apparaît qu'une excellente théorie des fibres $W$-homotopiques, satisfaisant aux conditions \cercle{a}\cercle{b}\cercle{c} du n$^\circ$ 4 (pages 215-218), existe dès que l'axiome d'apparence anodine suivant est satisfait (en plus de l'existence de l'objet final et des produits fibrés (\footnote[128]{et encore! il me semble qu'il suffit de supposer que les $X \longrightarrow Y$ qui sont $W$-fibrants sont `quarrables'.})).

\textbf{Factorisation gauche :} (\footnote[129]{Sous ces conditions, $\mathcal{M}$ muni de $(W, F_W)$ est une catégorie à cofibrations de K.S. \textsc{Brown-Anderson}, et `tout est O.K.', mais il faudrait vérifier a) b) c) ci-dessus (p. 215, 218), je ne suis pas sûr de l'avoir fait.},\footnote[130]{\textbf{NB} L'axiome de factorisation gauche n'est suffisant pour l'existence d'une bonne théorie des fibres homotopiques (semble-t-il) que dans le cas où tout $X$ dans $\mathcal{M}$ est $W$-fibrant, i.e. est dans $\Fibu_W e$.}). Toute flèche $f : X \longrightarrow Y$ de $\mathcal{M}$ se factorise en
\[
X \xrightarrow{i_f} \bar{X} \xrightarrow{p_f} Y,
\]
avec
\[
i_f \in W, \qquad p_f \in \mathrm{Fib}_W.
\]

L'axiome dual donne lieu à une bonne théorie des cofibres $W$-homotopiques.

\textbf{Factorisation droite :} Toute flèche $f : X \longrightarrow Y$ dans $\mathcal{M}$ se factorise en
\[
X \xrightarrow{j_f} \tilde{Y} \xrightarrow{q_f} Y,
\]
avec
\[
q_f \in W, \qquad j_f \in \mathrm{Cofib}_W.
\]

\msp{page 241}
Je voudrais à présent développer quelque peu le formalisme des fibres $W$-homotopiques, sous l'hypothèse (disons) que $W$ satisfait l'axiome de factorisation gauche.

Considérons un diagramme cartésien
\[
\begin{tikzcd}[ampersand replacement=\&]
X_1 \arrow[d, "p_1"'] \& Y \arrow[l, "q_2"'] \arrow[d, "q_1"] \\
S \& X_2 \arrow[l, "p_2"]
\end{tikzcd}
\]
\printedpage{156}
% head: A. Grothendieck : Dérivateurs, Ch. XII, Édition des Universités de Montpellier II et Paris VII
% foot: Version 30 mars 2003
dans $\mathcal{M}$. Si $p_1$ est dans $\mathrm{Fib}_W$, alors $Y$ sur $X_2$ peut s'interpréter comme la fibre $W$-homotopique de $X_1/S$ en $X_2/S$. Plus précisément, il en est ainsi de la classe $\ho^{\lc}_{W,X_2}(Y)$,
\[
\ho^{\lc}_{W,X_2}(Y) \simeq p_2^*(\ho^{\lc}_{W,S}(X_1/S)). \tag*{$(*)$}
\]
Mais en tant qu'objet sur $X_1$, $Y$ peut s'interpréter également comme fibre $W$-homotopique de $X_2/S$ en $X_1/S$ :
\[
\ho^{\lc}_{W,X_1}(Y/X_1) \simeq p_1^*(\ho^{\lc}_{W,S}(X_2/S)). \tag*{$(**)$}
\]
À dire vrai, il faudrait le prouver. Pour ceci, considérons une factorisation
\[
X_2 \xrightarrow{i_2} \bar{X}_2 \xrightarrow{\bar{p}_2} S \qquad i_2 \in W\ p_2 \in \mathrm{Fib}_W, % [sic] p_2 \in \mathrm{Fib}_W -> \bar{p}_2 \in \mathrm{Fib}_W, the factorisation just written is through \bar{p}_2
\]
de sorte qu'on trouve un diagramme de carrés cartésiens
\[
\begin{tikzcd}[ampersand replacement=\&]
X_1 \arrow[d] \& \bar{Y} \arrow[l, "\bar{q}_2"'] \arrow[d] \& Y \arrow[l, "j_2"', "W"] \arrow[d] \\
S \& \bar{X}_2 \arrow[l, "\bar{q}_2"'] \& X_2. \arrow[l, "i_2"', "W"] % [sic] \bar{q}_2 (bottom left arrow) -> \bar{p}_2, the arrow \bar{X}_2 \to S is \bar{p}_2 of the factorisation above
\end{tikzcd}
\]

\msp{page 242}
Par définition de $p_1^*$, on a
\[
p_1^*(\ho^{\lc}_{W,S}(X_2/S)) \simeq \ho^{\lc}_{W,X_1}(\bar{Y}/X_1),
\]
d'autre part, comme $X_1 \longrightarrow S$ est dans $\mathrm{Fib}_W$, et $i_2 \in W$, on a $j_2 \in W$, donc
\[
\ho^{\lc}_{W,X_1}(Y/X_1) \xrightarrow{\sim} \ho^{\lc}_{W,X_1}(\bar{Y}/X_1),
\]
d'où la formule $(**)$.

Ainsi, l'objet $Y$ a deux interprétations homotopiques différentes, suivant qu'on le regarde comme étant sur $X_1$, ou sur $X_2$. On aimerait alors une interprétation qui englobe les deux à la fois. On va considérer le diagramme
\[
\begin{tikzcd}[ampersand replacement=\&]
X_1 \arrow[d] \& \\
S \& X_2 \arrow[l]
\end{tikzcd}
\]
de type $\Psi^\circ$, ou
\[
\Psi \quad = \quad
\begin{tikzcd}[ampersand replacement=\&]
1 \& \\
0 \arrow[u] \arrow[r] \& 2
\end{tikzcd}
\qquad \text{ou} \qquad
\begin{tikzcd}[ampersand replacement=\&, row sep=small, column sep=small]
 \& 1 \\
0 \arrow[ur] \arrow[dr] \& \\
 \& 2
\end{tikzcd}
\quad ,
\]
% the last "2" of the first diagram is printed overlapping the page number 156
\printedpage{157}
% head: A. Grothendieck : Dérivateurs, Ch. XII, Édition des Universités de Montpellier II et Paris VII
% foot: Version 30 mars 2003
dont $Y$ se déduit par produit fibré, en complétant ce diagramme en un diagramme de type $\tilde{\Psi}^\circ$, où
\[
\tilde{\Psi} =
\begin{tikzcd}[ampersand replacement=\&, row sep=small, column sep=small]
 \& 1 \arrow[dr] \& \\
0 \arrow[ur] \arrow[dr] \& \& 3 \\
 \& 2 \arrow[ur] \&
\end{tikzcd}
\qquad (\simeq \boldsymbol{\Delta}^1 \times \boldsymbol{\Delta}^1).
\]

\msp{page 243}
Pour toute catégorie $I$, soit
\[
\Ho_{\mathcal{M},W}(I) \text{ \ ou \ } \Ho_W(I) \overset{\text{déf}}{=} \mathcal{M}(I)W(I)^{-1},
\]
où
\[
\left\{
\begin{array}{rcl}
\mathcal{M}(I) & = & \Homu(I^\circ, \mathcal{M}) \\[1ex]
W(I) & \subseteq & \mathrm{Fl}(\mathcal{M}(I)) \text{ est formé des } u : F \longrightarrow G \text{ telles que} \\[1ex]
 & & u_i : F(i) \longrightarrow G(i) \text{ dans } W \text{ pour tout } i \text{ dans } \mathrm{Ob}\, I.
\end{array}
\right.
\]

Ainsi, le diagramme $\vcenter{\hbox{\begin{tikzcd}[ampersand replacement=\&, row sep=small, column sep=small]
 \& X_1 \arrow[dl] \\
S \& \\
 \& X_2 \arrow[ul]
\end{tikzcd}}}$ définit un objet $\xi$ de $\mathcal{M}(\Psi)$, donc de $\Ho_W(\Psi)$, et le diagramme $\vcenter{\hbox{\begin{tikzcd}[ampersand replacement=\&, row sep=small, column sep=small]
 \& X_1 \arrow[dl] \& \\
S \& \& Y \arrow[ul] \arrow[dl] \\
 \& X_2 \arrow[ul] \&
\end{tikzcd}}}$ peut de même s'interpréter comme un objet $\tilde{\xi}$ de $\Ho_W(\tilde{\Psi})$, dont la restriction à $\Psi$ est $\xi$.

Ceci posé, j'aimerais définir un foncteur
\[
i_* : \Ho_W(\Psi) \longrightarrow \Ho_W(\tilde{\Psi})
\]
(où
\[
i : \Psi \longrightarrow \tilde{\Psi}
\]
est l'inclusion, qui est une immersion ouverte), lequel doit correspondre, dans le cas d'un objet $\vcenter{\hbox{\begin{tikzcd}[ampersand replacement=\&, row sep=small, column sep=small]
 \& X_1 \arrow[dl] \\
S \& \\
 \& X_2 \arrow[ul]
\end{tikzcd}}}$ de $\Ho_W(\Psi)$ avec $p_1 \in \mathrm{Fib}_W$ \texttt{[}$:$ $X_1 \longrightarrow S$\texttt{]}, à l'opération de produit fibré (\footnote[131]{\textbf{NB} Cette construction, qui se fait ici assez lourdement, devient immédiate, et aussi ses principales propriétés, si on admet que la catégorie à fibrations $(\mathcal{M}, W, \mathrm{Fib}_W)$ définit un \textit{dérivateur} sur le domaine des ensembles ordonnés finis, ce qui se prouve en effet assez simplement (avec pas plus de mal que je ne me donne ici).}).

\printedpage{158}
% head: A. Grothendieck : Dérivateurs, Ch. XII, Édition des Universités de Montpellier II et Paris VII
% foot: Version 30 mars 2003

\msp{page 244}
La construction de ce foncteur $i_*$ devrait être essentiellement sorital, avec ce dont nous disposons. Nous allons définir d'abord, pour $S$ \textit{fixé,} un foncteur canonique
\[
\Ho^{\lc}_W(S) \times \Ho^{\lc}_W(S) \longrightarrow \Ho_{W,S}(\Psi), \tag*{$(*)$}
\]
où le deuxième membre est défini, par analogie avec les $\Ho_W(I)$, par
\[
\Ho_{W,S}(I) \overset{\text{déf}}{=} ((\mathcal{M}/S)(I))(W_S(I))^{-1}
\]
- i.e. c'est la même définition, sauf que $(\mathcal{M}, W)$ est remplacé par $(\mathcal{M}/S, W_S)$. Quant à $\Phi = \vcenter{\hbox{\begin{tikzcd}[ampersand replacement=\&, row sep=small, column sep=small]
1 \arrow[dr] \& \\
 \& 3 \\
2 \arrow[ur] \&
\end{tikzcd}}}$ , on le considère comme une sous-catégorie fermée de $\tilde{\Psi}$, qui s'en déduit en lui rajoutant un objet initial 0 (tout comme $\tilde{\Psi}$ se déduit de $\Psi$ en lui rajoutant un objet final 3). D'autre part, on a un foncteur tautologique
\[
\Ho_{W,S}(\Phi) \longrightarrow \Ho_W(\tilde{\Psi}), \tag*{$(*')$}
\]
qui sur les objet consiste à compléter un diagramme $\vcenter{\hbox{\begin{tikzcd}[ampersand replacement=\&, row sep=small, column sep=small] % [sic] les objet -> les objets, plural agreement
X_1 \& \\
 \& Y \arrow[ul] \arrow[dl] \\
X_2 \&
\end{tikzcd}}}$ de $\mathcal{M}/S$ en le diagramme carré $\vcenter{\hbox{\begin{tikzcd}[ampersand replacement=\&, row sep=small, column sep=small]
 \& X_1 \arrow[dl] \& \\
S \& \& Y \arrow[ul] \arrow[dl] \\
 \& X_2 \arrow[ul] \&
\end{tikzcd}}}$ de $\mathcal{M}$.

\msp{page 245}
Le composé de $(*)$ et de $(*')$ donnera donc un foncteur
\[
\Ho^{\lc}_W(S) \times \Ho^{\lc}_W(S) \longrightarrow \Ho_W(\tilde{\Psi}), \tag*{$(**)$}
\]
très proche déjà du foncteur $i_*$ à construire. Pour construire $(*)$, on interprète (à équivalence de catégories près) le premier membre comme
\[
(\Fibu_W S \times \Fibu_W S)(W_S \times W_S)^{-1},
\]
donc il faut construire un foncteur
\[
\Fibu_W S \times \Fibu_W S \longrightarrow \Ho_{W,S}(\Phi),
\]
qui transforme flèches dans $W_S \times W_S$ en isomorphismes. On prendra le foncteur
\[
(X_1, X_2) \ \longmapsto \
\left(
\begin{tikzcd}[ampersand replacement=\&]
X_1 \& X_1 \times_S X_2 \arrow[l] \arrow[d] \\
 \& X_2
\end{tikzcd}
\right),
\]
\printedpage{159}
% head: A. Grothendieck : Dérivateurs, Ch. XII, Édition des Universités de Montpellier II et Paris VII
% foot: Version 30 mars 2003
et il faut seulement vérifier que si
\[
X_1 \xrightarrow{\alpha_1} X_1', \qquad X_2 \xrightarrow{\alpha_2} X_2'
\]
sont deux flèches dans $\Fibu_W S$, d'où
\[
X_1 \times_S X_2 \xrightarrow{\alpha = \alpha_1 \times_S \alpha_2} X_1' \times_S X_2',
\]
alors
\[
\alpha_1 \in W,\ \alpha_2 \in W \qquad \Longrightarrow \qquad \alpha \in W,
\]
ce qui est immédiat sur la factorisation
\[
X_1 \times_S X_2 \quad \underbrace{\longrightarrow}_{\substack{\in W \\ \text{car } X_2 \in \Fibu_W S}} \quad X_1' \times_S X_2 \quad \underbrace{\longrightarrow}_{\substack{\in W \\ \text{car } X_1' \in \Fibu_W S}} \quad X_1' \times_S X_2'.
\]

\msp{page 246}
On a prouvé la

\textbf{Proposition 18.} \textit{Soit $S$ dans $\mathcal{M}$. Il existe, à isomorphisme unique près, un foncteur et un seul}
\[
\Ho^{\lc}_W(S) \times \Ho^{\lc}_W(S) \xrightarrow{\pi} \Ho_{W,S}(\Phi)
\]
\textit{donnant lieu à un carré commutatif à isomorphisme près}
\[
\begin{tikzcd}[ampersand replacement=\&, column sep=large]
 \& {\scriptstyle (X_1,X_2)\ \longmapsto\ \begin{array}{@{}c@{}c@{}c@{}} \scriptstyle X_1 & & \\ & \nwarrow & \\ & & \scriptstyle X_1 \times_S X_2 \\ & \swarrow & \\ \scriptstyle X_2 & & \end{array}} \& \\[-3ex]
\Fibu_W S \times \Fibu_W S \arrow[rr] \arrow[d] \& \& (\mathcal{M}/S)(\Phi) \arrow[d] \\
\Ho^{\lc}_W(S) \times \Ho^{\lc}_W(S) \arrow[rr] \& \& \Ho_{W,S}(\Phi) \\[-4ex]
(= (\mathcal{M}/S)(W_S)^{-1} \times (\mathcal{M}/S)(W_S)^{-1}) \& \& (= \mathcal{M}/S(\Phi)(W_S(\Phi))^{-1}).
\end{tikzcd}
\]
\textit{Pour $X_1$, $X_2$ dans $\mathcal{M}/S$, le diagramme}
\[
\pi(X_1, X_2) \overset{\text{déf}}{=} \pi(\ho^{\lc}_{W,S}(X_1), \ho^{\lc}_{W,S}(X_2))
\]
\textit{se `calcule' ainsi, à isomorphisme canonique près : on choisit}
\[
X_1 \xrightarrow{\alpha_1} \bar{X}_1, \qquad X_2 \xrightarrow{\alpha_2} \bar{X}_2 \qquad \textit{flèches dans } \mathcal{M}/S
\]
\textit{avec}
\[
\alpha_1, \alpha_2 \in W_S, \qquad \bar{X}_1, \bar{X}_2 \in \mathrm{Ob}\, \Fibu_W S,
\]
\textit{ce qui définit}
\[
\pi(\alpha_1, \alpha_2) : \pi(X_1, X_2) \xrightarrow{\sim} \pi(\bar{X}_1, \bar{X}_2) = \ho_{S,W,\Phi}
\left(
\begin{tikzcd}[ampersand replacement=\&]
\bar{X}_1 \& \bar{X}_1 \times_S \bar{X}_2 \arrow[l] \arrow[d] \\
 \& \bar{X}_2
\end{tikzcd}
\right).
\]


\endgroup
% END UNIT Der-ChXIIc-30-03-03/Der-ChXIIc-30-03-03_p041-045.tex
% BEGIN UNIT Der-ChXIIc-30-03-03/Der-ChXIIc-30-03-03_p046-050.tex
\begingroup\makeatletter%
\newcommand{\printedpage}[1]{\marginpar{\footnotesize p.~#1}}%
\newcommand{\msp}[1]{\marginpar{\footnotesize [#1]}}%
\newcommand{\Ho}{\mathrm{Ho}}%
\newcommand{\lc}{\mathrm{lc}}%
\newcommand{\uFib}{\underline{\mathrm{Fib}}}%
\makeatother\let\printedpage\derpp\let\msp\dermsp%
\let\mspin\DERmspin%

\setcounter{footnote}{131}

\printedpage{160}
% head: A. Grothendieck : Dérivateurs, Ch. XII, Édition des Universités de Montpellier II et Paris VII
% foot: Version 30 mars 2003

\msp{page 247}

\textbf{Corollaire.} {\itshape Description du foncteur
\[
(\xi_1,\xi_2)\longmapsto\xi_1\times_S\xi_2 : \Ho^{\lc}_W(S)\times\Ho^{\lc}_W(S)\longrightarrow\Ho^{\lc}_W(S).
\]
Ce foncteur est bel et bien un produit cartésien dans} $\Ho^{\lc}_W(S)$ (\footnote{mais c'est un cas particulier des cas de $\Ho_W$ (si $\mathcal{M}$ a des produits binaires), en remplaçant $\mathcal{M}$ par $\mathcal{M}/S$.}), {\itshape où les produits finis existent} (\footnote{Dire aussi que les foncteurs $f^* : \Ho^{\lc}_W(S)\longrightarrow\Ho^{\lc}_W(S')$ (pour $f : S'\longrightarrow S$) commutent aux produits finis.}).

Pour prouver ce dernier point, posant
\[
\mathcal{H}=\Ho^{\lc}_W(S),
\]
il faut montrer que les foncteurs
\[
\begin{tikzcd}[ampersand replacement=\&]
\mathcal{H} \arrow[r, shift left, "\delta"] \& \mathcal{H}\times\mathcal{H} \arrow[l, shift left, "\pi"]
\end{tikzcd}
\]
sont adjoints, avec
\[
\begin{array}{rclcrcl}
\delta & = & \mathrm{diag}_{\mathcal{H}} & : & \xi & \longmapsto & (\xi,\xi)\\
\pi & = & \text{`produit'} & : & (\xi_1,\xi_2) & \longmapsto & \xi_1\times_S\xi_2.
\end{array}
\]
On va définir les homomorphismes d'adjonction
\[
\begin{array}{rcllrcl}
\mathrm{id}_{\mathcal{H}} & \xrightarrow{\alpha} & \pi\delta & \text{i.e.} & \xi & \xrightarrow{\alpha_\xi} & \xi\times_S\xi\\
\mathrm{id}_{\mathcal{H}\times\mathcal{H}} & \xleftarrow{\beta} & \delta\pi & \text{i.e.} & (\xi_1,\xi_2) & \xleftarrow{\beta_{\xi_1,\xi_2}} & (\xi_1\times_S\xi_2,\xi_1\times_S\xi_2).
\end{array}
\]
Pour ceci, on décrit $\mathcal{H}$ comme
\[
\mathcal{H}=(\uFib_W S)(W_S)^{-1}
\]
et on note que $\delta$, $\pi$ proviennent de foncteurs adjoints au niveau de $\mathcal{H}_0=\uFib_W S$
\[
\begin{tikzcd}[ampersand replacement=\&]
\mathcal{H}_0 \arrow[r, shift left, "\delta_0"] \& \mathcal{H}_0\times\mathcal{H}_0 \arrow[l, shift left, "\pi_0"]
\end{tikzcd},
\]
lesquels donnent lieu à $\alpha_0$, $\beta_0$
\[
\begin{array}{rcl}
\mathrm{id}_{\mathcal{H}_0} & \xrightarrow{\alpha_0} & \pi_0\delta_0\\
\mathrm{id}_{\mathcal{H}_0\times\mathcal{H}_0} & \xleftarrow{\beta_0} & \delta_0\pi_0.
\end{array}
\]

\printedpage{161}
% head: A. Grothendieck : Dérivateurs, Ch. XII, Édition des Universités de Montpellier II et Paris VII
% foot: Version 30 mars 2003

\msp{page 248}

On définit $\alpha$, $\beta$ par passage aux catégories de fractions pour $\alpha_0$, $\beta_0$. Les 2 compatibilités, étant vérifiées pour $\alpha_0$, $\beta_0$, le sont aussi pour $\alpha$, $\beta$. Q.e.d.

Je dégage ici le

\textbf{Lemme.} {\itshape Soit \begin{tikzcd}[ampersand replacement=\&]
\mathcal{X} \arrow[r, shift left, "f"] \& \mathcal{Y} \arrow[l, shift left, "g"]
\end{tikzcd} un couple de foncteurs adjoint, et soient $\Sigma$, $\Sigma'$ des localiseurs dans $\mathcal{X}$, $\mathcal{Y}$ tels que $f(\Sigma)\subseteq\Sigma'$, $g(\Sigma')\subseteq\Sigma$, de sorte qu'on a un couple de foncteurs
\[
\begin{tikzcd}[ampersand replacement=\&]
\mathcal{X}\Sigma^{-1} \arrow[r, shift left, "\bar{f}"] \& \mathcal{Y}\Sigma'^{-1} \arrow[l, shift left, "\bar{g}"]
\end{tikzcd}.
\]
Alors $\bar{f}$, $\bar{g}$ sont adjoints, les flèches d'adjonction étant déduites de celles pour $f$, $g$.}

\textbf{Corollaire.} {\itshape Soit $(\mathcal{M},W)$ une catégorie de modèles, supposons que dans $\mathcal{M}$ les produits binaires (resp. les produits finis) existent, et que pour $u_1 : X_1\longrightarrow X'_1$ et $u_2 : X_2\longrightarrow X'_2$, si $u_1,u_2\in W$, alors $u_1\times u_2 : X_1\times X_2\longrightarrow X'_1\times X'_2$ est dans $W$. Alors dans $\mathcal{M}W^{-1}$ les produits binaires (resp. produits finis) existent, et le foncteur localisation $\mathcal{M}\longrightarrow\mathcal{M}W^{-1}$ y commute.}

On applique le lemme à
\[
\begin{tikzcd}[ampersand replacement=\&]
\mathcal{M} \arrow[r, shift left, "\delta"] \& \mathcal{M}\times\mathcal{M} \arrow[l, shift left, "\pi"]
\end{tikzcd}
\qquad \text{et } W,\ W\times W,
\]
et pour l'objet final,
\[
\begin{tikzcd}[ampersand replacement=\&]
\mathcal{M} \arrow[r, shift left, "e"] \& e \arrow[l, shift left, "\bullet"]
\end{tikzcd}
\qquad\qquad \text{et } W,\ W_0=\{\mathrm{id}_e\}.
\]

\msp{page 249}

Je vais maintenant construire le foncteur
\[
i_* : \Ho_W(\Psi)\longrightarrow\Ho_W(\tilde{\Psi}).
\]
Il me faut d'abord prouver le

\textbf{Lemme 1.} {\itshape Soit $\mathcal{F}$ la sous-catégorie pleine de $\mathcal{M}(\Psi)$, catégorie des diagrammes \begin{tikzcd}[ampersand replacement=\&, row sep=small, column sep=small]
 \& X_1 \arrow[dl, "p_1"'] \\
S \& \\
 \& X_2 \arrow[ul, "p_2"]
\end{tikzcd} dans $S$ \texttt{[plutôt dans }$\mathcal{M}$\texttt{, sur }$S$\texttt{]}, formée de ceux pour lesquels $p_1$, $p_2$ sont dans Fib$_W$. Soit $W_{\mathcal{F}}$ le localiseur induit sur $\mathcal{F}$ par $W(\Psi)$ ($W$-équivalences terme à terme). Alors
\[
\mathcal{F}W_{\mathcal{F}}^{-1}\xrightarrow{\sim}\mathcal{M}(\Psi)W(\Psi)^{-1}\ (\overset{\text{déf}}{=}\Ho_W(\Psi)).
\]}

\printedpage{162}
% head: A. Grothendieck : Dérivateurs, Ch. XII, Édition des Universités de Montpellier II et Paris VII
% foot: Version 30 mars 2003

Nous allons prouver ce lemme plus bas. En l'admettant, pour construire $i_*$, il suffit de construire
\[
\mathcal{F}\longrightarrow\Ho_W(\tilde{\Psi})
\]
qui applique $W_{\mathcal{F}}$ en des isomorphismes. Je prendrai, bien sûr,
\[
\begin{gathered}
\begin{tikzcd}[ampersand replacement=\&, row sep=small, column sep=small]
 \& X_1 \arrow[dl] \& \& \& \& X_1 \arrow[dl] \& \\
S \& \& \longmapsto \& \& S \& \& X_1\times_S X_2 \arrow[ul] \arrow[dl] \\
 \& X_2 \arrow[ul] \& \& \& \& X_2 \arrow[ul] \&
\end{tikzcd}\\
\mathcal{F}\quad\longrightarrow\quad\Ho_W(\Psi).
\end{gathered}
\] % [sic] \Ho_W(\Psi) -> \Ho_W(\tilde{\Psi}), the image is a diagram of type \tilde{\Psi} and the functor to be constructed is \mathcal{F}\longrightarrow\Ho_W(\tilde{\Psi})
Vérifions qu'il transforme flèches de $W_{\mathcal{F}}$ en isomorphismes.

\msp{page 250}

\[
\begin{tikzcd}[ampersand replacement=\&, row sep=small, column sep=small]
 \& X_1\times_S X_2 \arrow[dl] \arrow[dd] \& \& X'_1\times_{S'}X'_2 \arrow[ll, "f_3"'] \arrow[dl] \arrow[dd] \\
X_1 \arrow[dd] \& \& X'_1 \arrow[ll, crossing over, "f_1"'] \& \\
 \& X_2 \arrow[dl] \& \& X'_2 \arrow[ll, "f_2"'] \arrow[dl] \\
S \& \& S' \arrow[ll, "f_0"'] \arrow[from=uu, crossing over] \&
\end{tikzcd}
\]

Il faut prouver que si $f_0,f_1,f_2\in W$, alors $f_3\in W$. Mais considérons l'image inverse du carré cartésien $Q$ par $f_0 : S'\longrightarrow S$, d'où un carré cartésien $Q_{S'}$ et une factorisation
\[
\begin{tikzcd}[ampersand replacement=\&]
Q \& Q_{S'} \arrow[l, "\tilde{f}"'] \& Q' \arrow[l, "g"'] \arrow[ll, bend left, "f"]
\end{tikzcd}
\]
de $Q'\longrightarrow Q$. Comme $f_0\in W$, et que $X_1$, $X_2$, $X_1\times_S X_2$ sont dans $\uFib_W S$, les flèches horizontales réalisant $Q_{S'}\longrightarrow Q$ sont toutes dans $W$. Ceci, plus le fait que $f_1,f_2\in W$, impliquent que dans le morphisme $g$, les flèches $g_1$, $g_2$ sont dans $W$. Il en résulte que la flèche
\[
g_3 : X'_1\times_{S'}X'_2\longrightarrow(X_1\times_S X_2)\times_S S'
\]
l'est aussi (théorie du produit dans $\uFib_W S'$ et $\Ho^{\lc}_W(S)$), donc $f_3=\tilde{f}_3g_3\in W$, q.e.d.

\msp{page 251}

Notons qu'on a un isomorphisme canonique
\[
\alpha : i^*i_*\xrightarrow{\sim}\mathrm{id}_{\Ho_W(\Phi)}. \tag*{$(*)$}
\]
J'ai envie de prouver que les foncteurs
\[
\begin{tikzcd}[ampersand replacement=\&]
\Ho_W(\Phi) \arrow[r, shift right, "i_*"'] \& \Ho_W(\tilde{\Phi}) \arrow[l, shift right, "i^*"']
\end{tikzcd}
\]
\printedpage{163}
% head: A. Grothendieck : Dérivateurs, Ch. XII, Édition des Universités de Montpellier II et Paris VII
% foot: Version 30 mars 2003
sont adjoints, et que $(*)$ est une des flèches d'adjonction. Il faut définir l'autre flèche d'adjonction
\[
\beta : i_*i^*\longleftarrow\mathrm{id}_{\Ho_W(\tilde{\Phi})}. \tag*{$(**)$}
\]
Mais pour ceci, je dois utiliser le

\textbf{Lemme 2.} Soit $\mathcal{G}$ la sous-catégorie pleine de $\mathcal{M}(\tilde{\Phi})$, formée des \begin{tikzcd}[ampersand replacement=\&]
X_1 \arrow[d, "p_1"'] \& X_3 \arrow[l, "q_2"'] \arrow[d, "q_1"] \\
S \& X_2 \arrow[l, "p_2"]
\end{tikzcd} tels que $p_1,p_2\in\mathrm{Fib}_W$ et soit $W_{\mathcal{G}}$ le localiseur induit par $W(\tilde{\Phi})$ ($W$-équivalences terme à terme). Alors l'inclusion $\mathcal{G}\longrightarrow\mathcal{M}(\tilde{\Phi})$ induit une équivalence
\[
\mathcal{G}W_{\mathcal{G}}^{-1}\xrightarrow{\sim}\mathcal{M}(\tilde{\Phi})W(\tilde{\Phi})^{-1}\overset{\text{déf}}{=}\Ho_W(\tilde{\Phi}).
\]

Je vais pour le moment admettre ce

\msp{page 252}

lemme, tout comme le lemme 1 (p. 249). Je vais considérer $i^*$, $i_*$ comme des flèches entre les categories $\mathcal{F}W_{\mathcal{F}}^{-1}$, $\mathcal{G}W_{\mathcal{G}}^{-1}$. Elles sont induites par les flèches % [sic] categories -> catégories, missing accent
\[
\begin{tikzcd}[ampersand replacement=\&]
\mathcal{F} \arrow[r, shift right, "i_{0\,*}"'] \& \mathcal{G} \arrow[l, shift right, "i_0^*"']
\end{tikzcd}
\]
entre $\mathcal{F}$ et $\mathcal{G}$ eux mêmes, % [sic] eux mêmes -> eux-mêmes, missing hyphen
\[
\begin{tikzcd}[ampersand replacement=\&]
 \& X_1 \arrow[dd] \& X_3 \arrow[l] \arrow[dd] \& \& X_1 \arrow[dd] \& \\
i_0^* :\quad \mathcal{G}\quad\longrightarrow\quad\mathcal{F} : \& \& \& \longmapsto \& \& \\
 \& S \& X_2 \arrow[l] \& \& S \& X_2 \arrow[l]
\end{tikzcd}
\]
\[
\begin{tikzcd}[ampersand replacement=\&]
 \& X_1 \arrow[dd] \& \& \& X_1 \arrow[dd] \& X_1\times_S X_2 \arrow[l] \arrow[dd] \\
i_{0\,*} :\quad \mathcal{F}\quad\longrightarrow\quad\mathcal{G} : \& \& \& \longmapsto \& \& \\
 \& S \& X_2 \arrow[l] \& \& S \& X_2. \arrow[l]
\end{tikzcd}
\]
Ces deux foncteurs sont adjoints l'un de l'autre, et compatibles aux localiseurs $W_{\mathcal{F}}$, $W_{\mathcal{G}}$. Donc les foncteurs déduits de $i_0^*$, $i_{0\,*}$ par passage aux catégories de fractions sont également adjoints (cf. lemme p. 248).

Ainsi, modulo la démonstration des lemmes 1 et 2, on a prouvé ceci :

\printedpage{164}
% head: A. Grothendieck : Dérivateurs, Ch. XII, Édition des Universités de Montpellier II et Paris VII
% foot: Version 30 mars 2003

\msp{page 253}

\textbf{Théorème.} {\itshape Soit $(\mathcal{M},W)$ une catégorie de modèles ($W\subseteq\mathit{Fl}\,\mathcal{M}$) satisfaisant l'axiome de factorisation gauche (p. 240), donnant donc lieu à une bonne théorie des fibres $W$-homotopiques. Soient
\[
\begin{tikzcd}[ampersand replacement=\&, row sep=small, column sep=small]
 \& 1 \& \& \& \& 1 \arrow[dr] \& \\
\Psi=0 \arrow[ur] \arrow[dr] \& \& {} \arrow[r, "i"] \& {} \& \tilde{\Psi}=0 \arrow[ur] \arrow[dr] \& \& 3 \\
 \& 2 \& \& \& \& 2 \arrow[ur] \&
\end{tikzcd}
\]
(immersions ouvertes dans} $\mathrm{Cat}${\itshape ), et considérons le foncteur restriction
\[
i^* : \Ho_W(\tilde{\Psi})\longrightarrow\Ho_W(\Psi).
\]
Ce foncteur admet un adjoint à droite
\[
i_* : \Ho_W(\Psi)\longrightarrow\Ho_W(\tilde{\Psi}).
\]
Ce foncteur $i_*$ se `calcule' ainsi : sa valeur en \begin{tikzcd}[ampersand replacement=\&]
X_1 \arrow[d, "p_1"'] \& \\
S \& X_2 \arrow[l, "p_2"]
\end{tikzcd} est canoniquement isomorphe (dans} $\Ho_W(\tilde{\Psi})${\itshape ) à \begin{tikzcd}[ampersand replacement=\&]
X_1 \arrow[d] \& X_1\times_S X_2 \arrow[l] \arrow[d] \\
S \& X_2 \arrow[l]
\end{tikzcd}, si on suppose que $p_1$ ou $p_2$ est une $W$-fibration. Dans le cas général il se calcule en factorisant $p_1$ en
\[
X_1\xrightarrow{i_1}\bar{X}_1\xrightarrow{p'_1}S,\qquad i_1\in W,\ p'_1\in\mathit{Fib}_W,
\]
et en prenant \begin{tikzcd}[ampersand replacement=\&]
\bar{X}_1 \arrow[d] \& \bar{X}_1\times_S X_2 \arrow[l] \arrow[d] \\
S \& X_2 \arrow[l]
\end{tikzcd}, ou en

\msp{page 254}

procédant de façon symétrique, ou en factorisant simultanement $p_1$ et $p_2$, et prenant % [sic] simultanement -> simultanément, missing accent
\[
\begin{tikzcd}[ampersand replacement=\&]
\bar{X}_1 \arrow[d] \& \bar{X}_1\times_S\bar{X}_2 \arrow[l] \arrow[d] \\
S \& \bar{X}_2. \arrow[l]
\end{tikzcd}
\]}


\endgroup
% END UNIT Der-ChXIIc-30-03-03/Der-ChXIIc-30-03-03_p046-050.tex
% BEGIN UNIT Der-ChXIIc-30-03-03/Der-ChXIIc-30-03-03_p051-055.tex
\begingroup\makeatletter%
\NoAutoSpacing%
\tikzcdset{ampersand replacement=\&}%
\newcommand{\printedpage}[1]{\marginpar{\footnotesize p.~#1}}%
\newcommand{\msp}[1]{\marginpar{\footnotesize [#1]}}%
\newcommand{\Ho}{\mathrm{Ho}}%
\newcommand{\Fib}{\mathrm{Fib}}%
\makeatother\let\printedpage\derpp\let\msp\dermsp%
\let\mspin\DERmspin%


\printedpage{165}
% head: A. Grothendieck : Dérivateurs, Ch. XII, Édition des Universités de Montpellier II et Paris VII
% foot: Version 30 mars 2003

\textit{La loi fonctorielle s'obtient à l'aide du lemme 1 p. 249, qui permet de se borner aux arguments $\begin{tikzcd}[ampersand replacement=\&] X_1 \arrow[d, "p_1"'] \& \\ S \& X_2 \arrow[l, "p_2"] \end{tikzcd}$ tels que $p_1,p_2\in \mathit{Fib}_W$, de sorte que l'on trouve un foncteur}
\[
\begin{tikzcd}[ampersand replacement=\&]
X_1 \arrow[d] \& \\
S \& X_2 \arrow[l]
\end{tikzcd}
\qquad\longmapsto\qquad
\begin{tikzcd}[ampersand replacement=\&]
X_1 \arrow[d] \& X_1\times_S X_2 \arrow[l] \arrow[d] \\
S \& X_2, \arrow[l]
\end{tikzcd}
\]
\textit{compatible aux localiseurs, et passant aux catégories de fractions} $\Ho_W(\Psi)$ \textit{et} $\Ho_W(\tilde{\Psi})$ \textit{respectivement.}

\textbf{NB} En procédant de façon similaire, plus simple (sans besoin d'un recours au lemme 1), on trouve l'existence d'un \textit{adjoint à gauche} $i_!$ de $i^* : \Ho_W(\Psi)\longrightarrow \Ho_W(\tilde{\Psi})$, sous réserve de l'existence d'un objet initial $\varphi$ dans $\mathcal{M}$. Cet foncteur $i_!$ est donné par % [sic] Cet foncteur -> Ce foncteur, "foncteur" begins with a consonant

\msp{page 255}
\[
i_!\left(
\begin{tikzcd}[ampersand replacement=\&]
X_1 \arrow[d] \& \\
S \& X_2 \arrow[l]
\end{tikzcd}
\right)
\qquad = \qquad
\left(
\begin{tikzcd}[ampersand replacement=\&]
X_1 \arrow[d] \& \varphi \arrow[l] \arrow[d] \\
S \& X_2 \arrow[l]
\end{tikzcd}
\right),
\]
la loi fonctorielle étant évidente.

Ce n'est pas encore la fin du sorite! Il faut encore interpréter le foncteur composé
\[
\Ho_W(\Psi)\xrightarrow{\ i_*\ }\Ho_W(\tilde{\Psi})\xrightarrow{\ \varepsilon_3^*\ }\Ho_W(e)=\Ho_W,
\]
où
\[
\varepsilon_3 : e\xrightarrow{\ \sim\ }\{3\}
\qquad\longrightarrow\qquad
\begin{tikzcd}[ampersand replacement=\&, column sep=small, row sep=small]
\& 1 \arrow[dl] \& \\
3 \& \& 0 \arrow[ul] \arrow[dl] \\
\& 2 \arrow[ul] \&
\end{tikzcd}
\]
est l'inclusion du sommet 3 dans $\tilde{\Psi}$. Si $\Ho_W(\Psi)$ est interprété, conformément au lemme 1, comme $\mathcal{F}W_{\mathcal{F}}^{-1}$, ce foncteur composé n'est autre que
\[
\begin{tikzcd}[ampersand replacement=\&]
X_1 \arrow[d] \& \\
S \& X_2 \arrow[l]
\end{tikzcd}
\qquad\longmapsto\qquad X_1\times_S X_2
\]
\printedpage{166}
% head: A. Grothendieck : Dérivateurs, Ch. XII, Édition des Universités de Montpellier II et Paris VII
% foot: Version 30 mars 2003
(pour $\left(\begin{tikzcd}[ampersand replacement=\&] X_1 \arrow[d, "p_1"'] \& \\ S \& X_2 \arrow[l, "p_2"] \end{tikzcd}\right)\in \mathrm{Ob}\,\mathcal{F}$, i.e. $p_1,p_2\in\Fib_W$.)

On veut interpréter ce foncteur comme un $\varprojlim$ $W$-homotopique, i.e. vu comme un `produit fibré $W$-homotopique'. De façon précise, on a ceci :

\msp{page 256}
\textbf{Corollaire 1.} \textit{Considérons $p_\Psi : \Psi\longrightarrow e$, et le foncteur}
\[
p_\Psi^* : \Ho_W(e)=\Ho_W\longrightarrow\Ho_W(\Psi),
\]
\textit{déduit par passage aux catégories de fractions du foncteur}
\[
\begin{array}{ccc}
\mathcal{M} & \longrightarrow & \mathcal{M}(\Psi) \\
X & \longrightarrow & \left(
\begin{tikzcd}[ampersand replacement=\&]
X \arrow[d] \& \\
X \& X \arrow[l]
\end{tikzcd}
\right).
\end{array}
\]
\textit{Ce foncteur admet le foncteur $\varepsilon_3^*i_*$ précédent comme adjoint à droite, i.e.}
\[
p_{\Psi\,*}=\varepsilon_3^*i_* : \Ho_W(\Psi)\longrightarrow\Ho_W.
\]

\textsc{Démonstration}. On peut regarder $p_\psi$ comme le composé % [sic] p_\psi -> p_\Psi, the subscript is the capital \Psi everywhere else
\[
\Psi\xrightarrow{\ i\ }\tilde{\Psi}\xrightarrow{\ p_{\tilde{\Psi}}\ }e,
\]
et on a donc
\[
p_\Psi^*\simeq i^*(p_{\tilde{\Psi}})^*.
\]
Donc l'existence et la détermination d'un adjoint à droite pour $p_\Psi^*$ résultera de l'existence et de la détermination d'adjoints à droite $i_*$, $(p_{\tilde{\Psi}})_*$ de $i^*$, $(p_{\tilde{\Psi}})^*$, via
\[
(p_\Psi)_*\simeq(p_{\tilde{\Psi}})_*i_*.
\]

\msp{page 257}
Donc le corollaire 1 résulte de la formule
\[
p_{\tilde{\Psi}\,*}=\varepsilon_3^*,
\]
qui elle-même est un cas particulier du

\printedpage{167}
% head: A. Grothendieck : Dérivateurs, Ch. XII, Édition des Universités de Montpellier II et Paris VII
% foot: Version 30 mars 2003

\textbf{Lemme.} \textit{Soit $I$ dans} Cat\textit{, avec objet final $e$, considérons l'inclusion}
\[
\varepsilon : e\longrightarrow I,\qquad \textit{et } p : I\longrightarrow e.
\]
\textit{Soit $(\mathcal{M},W)$ une catégorie de modèles, et considérons}
\[
\Ho_W(I)\xleftarrow{\ p^*\ }\Ho_W.
\]
\textit{Ce foncteur admet un adjoint à gauche} \texttt{[plutôt à droite]} \textit{$p_*$, lequel est isomorphe à $\varepsilon^*$.}

\textsc{Démonstration}. Considérons le couple de foncteurs adjoints
\[
\begin{tikzcd}[ampersand replacement=\&, column sep=large]
\mathcal{M}(I) \arrow[r, shift right, "p_*=\varprojlim\limits_I"'] \& \mathcal{M} \arrow[l, shift right, "p^*"']
\end{tikzcd}\ ;
\]
comme $I$ admet l'objet final $e$, le foncteur $p_*$ est isomorphe à $\varepsilon^*$. Ces foncteurs $p^*$, $\varepsilon^*$ sont compatibles aux localiseurs $W(I)$, $W$, donc on conclut par le lemme page 248.

Il ne reste plus qu'à prouver les lemmes 1 et 2. Je le ferai demain.

\msp{page 258}
Aujourd'hui, j'ai envie encore de donner les sorites des carrés $W$-cartésiens dans une

\textbf{Proposition 19 :} \textit{Soit $\begin{tikzcd}[ampersand replacement=\&] X_1 \arrow[d, "p_1"'] \& X_3 \arrow[l, "f_1"'] \arrow[d, "f_2"] \\ S \& X_2 \arrow[l, "p_2"] \end{tikzcd}$ un carré commutatif dans $\mathcal{M}$. Considérons des factorisations de $p_1$, $p_2$ en}
\[
X_1\xrightarrow{\ i_1\ }\bar{X}_1\xrightarrow{\ \bar{p}_1\ }S,\qquad X_2\xrightarrow{\ i_2\ }\bar{X}_2\xrightarrow{\ \bar{p}_2\ }S,
\]
\textit{$i_1,i_2\in W$, $\bar{p}_1,\bar{p}_2\in F_W$} \texttt{[$=$ Fib$_W$]}\textit{, et considérons les flèches induites par $X_3\longrightarrow X_1\times_S X_2$ (déduite de $Q$) via $\begin{tikzcd}[ampersand replacement=\&, row sep=small] \& \bar{X}_1\times_S X_2 \\ X_1\times_S X_2 \arrow[ur] \arrow[r] \arrow[dr] \& \bar{X}_1\times_S\bar{X}_2 \\ \& X_1\times_S\bar{X}_2 \end{tikzcd}$, de sorte qu'on trouve un diagramme commutatif}
\[
\begin{tikzcd}[ampersand replacement=\&, row sep=small]
\& \bar{X}_1\times_S X_2 \arrow[d, "j_2"'] \\
X_3 \arrow[ur, "\pi_1"] \arrow[r, "\pi_0"] \arrow[dr, "\pi_2"] \& \bar{X}_1\times_S\bar{X}_2 \\
\& X_1\times_S\bar{X}_2, \arrow[u, "j_1"]
\end{tikzcd}
\]
\textit{où $j_1=i_1\times_S\mathrm{id}_{\bar{X}_2}$ et $j_2=\mathrm{id}_{\bar{X}_1}\times i_2$ sont dans $W$ (car $i_1$ et $i_2$ le sont, et $\bar{X}_1$ et $\bar{X}_2$ sont dans $\underline{\mathrm{Fib}}_W S$). On a donc}
\[
(\pi_1\in W)\qquad\iff\qquad(\pi_2\in W)\qquad\qquad(\iff\ \pi_0\in W),
\]
\textit{ce qui prouve que la condition $\pi_1\in W$, qui semble dépendre du choix de la factorisation par $i_1$, n'en dépend pas, et idem pour $\pi_2\in W$, et pour $\pi_0\in W$ (qui semble dépendre à la fois de $i_1$, et de $i_2$).}

\printedpage{168}
% head: A. Grothendieck : Dérivateurs, Ch. XII, Édition des Universités de Montpellier II et Paris VII
% foot: Version 30 mars 2003

\msp{page 259}
\textbf{Définition.} \textit{Quand ces conditions sont satisfaites, on dit que le carré $Q$ est $W$-cartésien.}

\textbf{Corollaire 1.} \textit{Si $Q$ est $W$-cartésien, alors l'image de $Q$ dans} $\Ho_W(\tilde{\Psi})$ \textit{est dans l'image essentielle de}
\[
i_* : \Ho_W(\Psi)\longrightarrow\Ho_W(\tilde{\Psi}),
\]
\textit{i.e. la flèche canonique dans} $\Ho_W(\tilde{\Psi})$
\[
Q\longrightarrow i_*i^*Q
\]
\textit{est} \texttt{[un]} \textit{isomorphisme. La réciproque est vraie si $W$ est fortement saturé.}

\textbf{Corollaire 2.} \textit{Considérons un diagramme commutatif dans $\mathcal{M}$}
\[
\begin{tikzcd}[ampersand replacement=\&]
X_0' \arrow[d] \& X_1' \arrow[l] \arrow[d] \arrow[dl, phantom, "\textit{1}"] \& X_2' \arrow[l] \arrow[d] \arrow[dl, phantom, "\textit{2}"] \\
X_0 \& X_1 \arrow[l] \& X_2, \arrow[l]
\end{tikzcd}
\tag*{(D)}
\]
\textit{d'où deux carrés commutatifs 1, 2, et un carré composé}
\[
\begin{tikzcd}[ampersand replacement=\&]
X_0' \arrow[d] \& X_2' \arrow[l] \arrow[d] \arrow[dl, phantom, "\textit{3}"] \\
X_0 \& X_2. \arrow[l]
\end{tikzcd}
\]
\textit{Supposons le carré 1 $W$-cartésien. Alors le carré 2 est $W$-cartésien si et seulement si le carré 3 l'est.}

\textsc{Démonstration}. Considérons une factorisation de $\begin{tikzcd}[ampersand replacement=\&, row sep=small] X_0' \arrow[d] \\ X_0 \end{tikzcd}$ en
\[
\begin{tikzcd}[ampersand replacement=\&, column sep=small, row sep=small]
\bar{X}_0' \arrow[dr] \& X_0' \arrow[l] \arrow[d] \\
\& X_0
\end{tikzcd}
\]
avec $(X_0'\longrightarrow\bar{X}_0')\in W$, et $(\bar{X}_0'\longrightarrow X_0)\in\Fib_W$. On peut alors remplacer dans (D) $X_0'$ par $\bar{X}_0'$, obtenant un diagramme $(\bar{D})$, et on voit aussitôt
\printedpage{169}
% head: A. Grothendieck : Dérivateurs, Ch. XII, Édition des Universités de Montpellier II et Paris VII
% foot: Version 30 mars 2003
\msp{page 260}
que chacun des carrés 1(D), 2(D), 3(D) est $W$-cartésien si et seulement si le carré correspondant $1(\bar{D})$, $2(\bar{D})$ ou $3(\bar{D})$ l'est. Cela montre qu'on peut supposer $X_0\longrightarrow S$ dans $\Fib_W S$. Considérons alors le diagramme commutatif déduit de (D)
\[
\begin{tikzcd}[ampersand replacement=\&, column sep=small]
X_0' \arrow[d] \& \tilde{X}_1'\overset{\text{déf}}{=}X_0'\times_{X_0}X_1 \arrow[l] \arrow[d] \& X_1' \arrow[l, "\alpha"'] \& X_2' \arrow[l] \arrow[d] \\
X_0 \& X_1 \arrow[l] \& \& X_2, \arrow[ll]
\end{tikzcd}
\]
où $\alpha\in W$ en vertu de l'hypothèse $W$-cartésienne sur le carré 1. On voit encore que le carré 2 est $W$-cartésien si et seulement si le carré correspondant avec $X_1'$ remplacé par $\tilde{X}_1'$ l'est, quant aux carrés 1, 3, ils ne changent pas par cette modification. Cela montre qu'on peut supposer le carré 1 \textit{cartésien.} Considérons alors le diagramme
\[
\begin{tikzcd}[ampersand replacement=\&]
X_0' \arrow[d, "\mathrm{Fib}_W"'] \& X_1' \arrow[l] \arrow[d, "\mathrm{Fib}_W"] \arrow[dl, phantom, "\text{cart.}"] \& \tilde{X}_2'\overset{\text{déf}}{=}X_1'\times_{X_1}X_2 \arrow[l] \arrow[d, "\mathrm{Fib}_W"] \arrow[dl, phantom, "\text{cart.}"] \& X_2' \arrow[l, "\beta"'] \\
X_0 \& X_1 \arrow[l] \& X_2. \arrow[l]
\end{tikzcd}
\]
Alors la proposition 19 et la définition montrent que le carré 2 est $W$-cartésien si et seulenment si $\beta\in W$, et de même pour le carré 3. Ce qui achève la démonstration. % [sic] seulenment -> seulement, misspelling

\textbf{Corollaire 2.} \textit{Soit $\begin{tikzcd}[ampersand replacement=\&] X_1 \arrow[d, "p_1"'] \& X_3 \arrow[l, "f_2"'] \arrow[d, "f_1"] \\ S \& X_2 \arrow[l, "p_2"] \end{tikzcd}$ un carré $W$-cartésien. Alors}
\[
\begin{aligned}
p_1\in W\quad&\Longrightarrow\quad f_1\in W\\
p_2\in W\quad&\Longrightarrow\quad f_2\in W,
\end{aligned}
\]
\textit{plus précisément, si $p_i\in W$, alors $Q$ est $W$-cartésien si et seulement si $f_i\in W$.}

\msp{page 261}
Supposons $p_1\in W$. Alors dans la factorisation de $p_1$ en $X_1\longrightarrow\bar{X}_1\longrightarrow S$, on peut prendre $\bar{X}_1=S$. Donc le carré est $W$-cartésien si et seulement si
\[
(X_3\longrightarrow\bar{X}_1\times_S X_2=X_2)\in W,
\]
i.e. $f_1\in W$, q.e.d.

Bien sûr, on ne peut en général affirmer les implications inverses : $f_i\in W$ n'implique nullement $p_i\in W$.


\endgroup
% END UNIT Der-ChXIIc-30-03-03/Der-ChXIIc-30-03-03_p051-055.tex
% BEGIN UNIT Der-ChXIIc-30-03-03/Der-ChXIIc-30-03-03_p056-060.tex
\begingroup\makeatletter%
\NoAutoSpacing%
\newcommand{\printedpage}[1]{\marginpar{\footnotesize p.~#1}}%
\newcommand{\msp}[1]{\marginpar{\footnotesize [#1]}}%
\newcommand{\Ho}{\mathrm{Ho}}%
\newcommand{\Cof}{\mathrm{Cof}}%
\makeatother\let\printedpage\derpp\let\msp\dermsp%
\let\mspin\DERmspin%


\printedpage{170}
% head: A. Grothendieck : Dérivateurs, Ch. XII, Édition des Universités de Montpellier II et Paris VII

Tous ces résultats, et la notion de carré $W$-cartésien, se dualisent (moyennant l'axiome idoine de factorisation à droite, cf. page 240) en une théorie des sommes amalgamées $W$-homotopiques, et des carrés $W$-cocartésiens. Cette fois on travaille avec la catégorie
\[
\Phi \;=\; \left(
\begin{tikzcd}[ampersand replacement=\&, row sep=small, column sep=small]
 \& 1 \arrow[dl] \\
0 \& \\
 \& 2 \arrow[ul]
\end{tikzcd}
\right)
\qquad\quad \overset{j}{\hookrightarrow} \qquad\quad
\tilde{\Phi} \;=\; \left(
\begin{tikzcd}[ampersand replacement=\&, row sep=small, column sep=small]
 \& 1 \arrow[dl] \& \\
0 \& \& 3 \arrow[ul] \arrow[dl] \\
 \& 2 \arrow[ul] \&
\end{tikzcd}
\right),
\]
$j$ étant une immersion fermée. On a un foncteur restriction
\[
j^* : \Ho_W(\tilde{\Phi}) \longrightarrow \Ho_W(\Phi),
\]
dont on construit un adjoint à gauche
\[
j_! : \Ho_W(\Phi) \longrightarrow \Ho_W(\tilde{\Phi}).
\]
Et on trouve une notion naturelle de
\msp{page 262}
carré $W$-cocartésien, qui paraît irréprochable à tout points de vue, et d'une simplicité parfaite. % [sic] tout points -> tous points, agreement of "tout"

Cela amène à revoir la notion en question dans Cat, où elle avait l'air très astucieuse, profonde et un peu obscure, avec des intégrales et tout ça. C'est le moment de se rendre compte que `ce n'est que ça'. Donc de prouver la

\textbf{Proposition 20.} \textit{Supposons} $\mathcal{M} = \mathrm{Cat}$\textit{, $W$ un localiseur fondamental satisfaisant W(7 bis). Considérons un carré commutatif dans} Cat
\[
Q = \left(
\begin{tikzcd}[ampersand replacement=\&]
S \arrow[r] \arrow[d] \& X_1 \arrow[d] \\
X_2 \arrow[r] \& X_3
\end{tikzcd}
\right).
\]
\textit{Pour qu'il soit $W$-cartésien} \texttt{[plutôt} $W$\texttt{-cocartésien]} \textit{au sens de la théorie développée ici (pour des catégories de modèles générales), il faut et il suffit que}
\[
\varphi_Q : \quad \mathcal{X}_Q = \left(
\begin{tikzcd}[ampersand replacement=\&, row sep=small, column sep=small]
 \& X_1 \arrow[dl] \\
\int S \& \\
 \& X_2 \arrow[ul]
\end{tikzcd}
\right) \qquad\quad \longrightarrow \qquad\quad X_3
\]

\texttt{[plutôt}
\[
\varphi_Q : \quad \mathcal{X}_Q = \left(
\begin{tikzcd}[ampersand replacement=\&, row sep=small, column sep=small]
 \& X_1 \\
\int S \arrow[ur] \arrow[dr] \& \\
 \& X_2
\end{tikzcd}
\right) \qquad\quad \longrightarrow \qquad\quad X_3\text{\texttt{]}}
\]
\printedpage{171}
% head: A. Grothendieck : Dérivateurs, Ch. XII, Édition des Universités de Montpellier II et Paris VII
\textit{soit dans $W$.} (NB L'intérêt de ce critère est qu'il est `intrinsèque' à $Q$, ne faisant pas appel à une factorisation de $\beta_1$ ou de $\beta_2$ \texttt{[où} $\beta_i : S \longrightarrow X_i$\texttt{]} en $S \xrightarrow{\;\in \Cof_W\;} \tilde{X}_i \xrightarrow{\in W} X_i$. Cf. application intéressante dans section V 3.4., p. 146 ff., th. 4 et corollaires.)

\msp{page 263}
\textsc{Démonstration.} Considérons une factorisation de $S \longrightarrow X_1$ en
\[
S \xrightarrow{i_1} \tilde{X}_1 \xrightarrow{p_1} X, % [sic] X -> X_1, the factorisation is of S -> X_1
\]
avec $i_1$ immersion ouverte de Dwyer, $p_1 \in W$. Donc on a un carré commutatif
\[
\tilde{Q} = \left(
\begin{tikzcd}[ampersand replacement=\&]
S \arrow[r, hook, "\mathrm{Dw}"] \arrow[d] \& \tilde{X}_1 \arrow[d] \\
X_2 \arrow[r] \& X_3
\end{tikzcd}
\right)
\]
et un homomorphisme $\tilde{Q} \xrightarrow{\alpha} Q$ qui est l'identité sur les trois sommets distincts de $\tilde{X}_1$, et qui est $p_1$ en $\tilde{X}_1$. On en déduit un diagramme commutatif
\[
\begin{tikzcd}[ampersand replacement=\&, row sep=small, column sep=small]
\mathcal{X}_Q \arrow[rr, "\mathcal{X}_\alpha = \lambda"] \arrow[dr, "\varphi_Q"'] \& \& \mathcal{X}_{\tilde{Q}} \arrow[dl, "\varphi_{\tilde{Q}}"] \\
 \& X_3, \&
\end{tikzcd}
\]
et on a $\lambda \in W$ (car c'est une $\Phi$-flèche de catégories fibrées sur $\Phi$, qui induit une $W$-équivalence sur chacune de trois fibres.) Donc $\varphi_Q \in W \iff \varphi_{\tilde{Q}} \in W$, d'autre part il est clair que $Q$ est $W$-cocartésien si et seulement si $\tilde{Q}$ l'est. Donc on peut supposer que $S \hookrightarrow X_1$ est \texttt{[une]} immersion ouverte (de Dwyer) (\footnote[134]{On utilise ici W(3) (forme forte), i.e. Loc(3 bis) avec les notations de XVI. Mais on peut s'en passer, cf. plus bas.}). Appliquant ceci
\msp{page 264}
en changeant les rôles de $X_1$ et $X_2$, on voit qu'on peut supposer de plus que $S \longrightarrow X_2$ est également \texttt{[une]} immersion ouverte de Dwyer. Considérons alors
\[
\begin{tikzcd}[ampersand replacement=\&]
S \arrow[r, hook] \arrow[d, hook] \arrow[dr, phantom, "Q'"] \& X_1 \arrow[d, hook] \& \\
X_2 \arrow[r, hook] \& X_1 \amalg_S X_2 \arrow[r, "\pi"] \& X_3 \\
 \& \mathcal{X}_Q \arrow[u, "\varphi_{Q'}"] \arrow[ur, "\varphi_Q"'] \&
\end{tikzcd}
\]
\printedpage{172}
% head: A. Grothendieck : Dérivateurs, Ch. XII, Édition des Universités de Montpellier II et Paris VII
On a $\varphi_Q = \pi\varphi_{Q'}$, d'autre part on sait que $\varphi_{Q'} \in W$ (cor. 3, p. 149). Donc
\[
\pi \in W \qquad\quad \iff \qquad\quad \varphi_Q \in W,
\]
où $\pi \in W$ signifie que $Q$ est $W$-cocartésien, q.e.d.

\msp{page 265}

\section*{8\quad Appendice : Retour sur les axiomes W(7), W(7 bis) (pages 75, 76).}

(\footnote[135]{Reflexion provisoire, supplanté par XVI.}). On a vu (prop. 6, p. 78) que W(7 bis) implique W(7). Montrons une réciproque partielle. % [sic] Reflexion -> Réflexion, missing accent; supplanté -> supplantée, agreement (in footnote 135)

\textbf{Proposition.} \textit{Supposons} W(7 bis) (axiome du carré cocartésien) \textit{satisfait, et plaçons nous sous les conditions de l'énoncé de} W(7), \textit{et plus généralement, dans le cas d'un diagramme cubique}
\[
\begin{tikzcd}[ampersand replacement=\&, row sep=small, column sep=small]
 \& X_2 \arrow[rr] \arrow[dd, "f_2" near end] \& \& X_3 \arrow[dd, "f_3"] \\
X_0 \arrow[ur] \arrow[rr, crossing over] \arrow[dd, "f_0"'] \& \& X_1 \arrow[ur] \& \\
 \& Y_2 \arrow[rr] \& \& Y_3 \\
Y_0 \arrow[ur] \arrow[rr] \& \& Y_1 \arrow[ur] \arrow[from=uu, crossing over, "f_1"' near end] \&
\end{tikzcd}
\]
\textit{où les deux carrés $X$ et $Y$ (face supérieure et inférieure du cube) sont $W$-cocartésiens. Considérons les conditions suivantes sur ce diagramme}
\begin{enumerate}[label={}]
\item[\textit{a)}] \textit{$f_0 \in W$.}
\item[\textit{b)}] \textit{$f_1$ et $f_2 \in W$.}
\item[\textit{c)}] \textit{$f_3 \in W$.}
\end{enumerate}
\textit{Alors on a les implications}
\[
a \text{ \textit{\&} } b \Longrightarrow c, \qquad\quad b \text{ \textit{\&} } c \Longrightarrow a.
\]
\textit{De plus, si $W$ est stable par facteurs directs (p. ex. $W$ fortement saturé), alors on a aussi}
\[
a \text{ \textit{\&} } c \Longrightarrow b,
\]
\textit{donc la conjonction} \texttt{[de]} \textit{deux quelconques parmi les conditions a) b) c) implique la troisième.}

\textsc{Démonstration.} Utilisant la factorisation
\printedpage{173}
% head: A. Grothendieck : Dérivateurs, Ch. XII, Édition des Universités de Montpellier II et Paris VII
\msp{page 266}
fonctorielle d'une flèche quelconque (ici les $f_i$) en $pi$, avec $i$ immersion ouverte de Dwyer et $p \in W$ (p. 229), on est remené au cas où les $f_i$ sont des immersions ouvertes de Dwyer. Mais pour toute flèche $g : S \longrightarrow T$ qui est dans $\Cof_W$, soient $Z(g)$ son mapping-cone % [sic] remené -> ramené, spelling
\[
\begin{tikzcd}[ampersand replacement=\&, row sep=small]
S \arrow[rr, "g"] \arrow[dr] \& \& T \arrow[dr, "\alpha(g)"] \& \\
 \& e \arrow[rr, "\bar{g}"] \& \& Z(g),
\end{tikzcd}
\]
et
\[
\alpha(g) : T = \mathrm{but}(g) \longrightarrow Z(g)
\]
la flèche canonique (\footnote[136]{C'est un peu vif, car il n'est pas clair que $X_1' \amalg_{X_0} X_2' \longrightarrow X_1 \amalg_{X_0} X_2$ soit \texttt{[dans]} $W$, sans utiliser déjà un des trois ingrédients de W(7) (savoir $f_0, f_1, f_2 \in W \implies f_3 \in W$). Il semble donc que notre argument ne marche que moyennant une hypothèse supplémentaire de cet ordre \dots Mais il y a le \textit{lemme de Brown,} cf. XVI \S3, notamment p. 20.}). Celle-ci est fonctorielle en $g$, de sorte que le diagramme cubique donné se prolonge en un double cube
\[
\begin{tikzcd}[ampersand replacement=\&, row sep=small, column sep=small]
 \& X_2 \arrow[rr] \arrow[dd, "f_2" near end] \& \& X_3 \arrow[dd, "f_3"] \\
X_0 \arrow[ur] \arrow[rr, crossing over] \arrow[dd, "f_0"'] \& \& X_1 \arrow[ur] \& \\
 \& Y_2 \arrow[rr] \arrow[dd, "\alpha_2" near end] \& \& Y_3 \arrow[dd, "\alpha_3"] \\
Y_0 \arrow[ur] \arrow[rr, crossing over] \arrow[dd, "\alpha_0"'] \& \& Y_1 \arrow[ur] \arrow[from=uu, crossing over, "f_1"' near end] \& \\
 \& Z_2 \arrow[rr] \& \& Z_3 \\
Z_0 \arrow[ur] \arrow[rr] \& \& Z_1 \arrow[ur] \arrow[from=uu, crossing over, "\alpha_1"' near end] \&
\end{tikzcd}
\]
avec
\[
Z_i = Z(f_i),\ \alpha_i = \alpha(f_i) \qquad\quad \text{pour } i \in \{0,1,2,3\}.
\]
Notons que W(7 bis) implique que si $g \in \Cof_W$,
\msp{page 267}
on a
\[
g \in W \qquad\quad \iff \qquad\quad \bar{g} \in W, \qquad\quad \text{i.e. } Z(g)\ W\text{-asphérique},
\]
en particulier
\[
f_i \in W \qquad\quad \iff \qquad\quad Z_i\ W\text{-asphérique}.
\]
Ainsi, l'énoncé à prouver sur le cube initial se ramène à un énoncé sur le carré
\printedpage{174}
% head: A. Grothendieck : Dérivateurs, Ch. XII, Édition des Universités de Montpellier II et Paris VII
\[
\begin{tikzcd}[ampersand replacement=\&, row sep=small, column sep=small]
 \& Z_1 \arrow[dr] \& \\
Z_0 \arrow[ur] \arrow[dr] \& \& Z_3, \\
 \& Z_2 \arrow[ur] \&
\end{tikzcd}
\]
savoir que les trois conditions
\begin{enumerate}[label={}]
\item[a$'$)] $Z_0$ $W$-asphérique.
\item[b$'$)] $Z_1$ et $Z_2$ $W$-asphériques.
\item[c$'$)] $Z_3$ $W$-asphérique.
\end{enumerate}
donnant lieu aux implications
\[
\mathrm{a}' \text{ \& } \mathrm{b}' \Longrightarrow \mathrm{c}', \qquad\quad \mathrm{b}' \text{ \& } \mathrm{c}' \Longrightarrow \mathrm{a}',
\]
et que si $W$ est stable par facteurs directs, alors
\[
\mathrm{a}' \text{ \& } \mathrm{c}' \Longrightarrow \mathrm{b}'.
\]

C'est à peu de choses près le corollaire de la page 76, complété à la page 77 (corollaire), à cela près qu'on y suppose les flèches des immersions ouvertes, et le carré cartésien (donc aussi cocartésien) dans Cat, i.e. $Z_3 = Z_1 \cup Z_2$, $Z_0 = Z_1 \cap Z_2$. Mais
\msp{page 268}
ce qui importe pour la validité de cet énoncé, c'est (on s'en doute) que $Z_*$ soit $W$-cocartésien. Donc il faut d'abord prouver ceci.

\textbf{Lemme.} \textit{Si le diagramme cubique p. 265 est tel que les faces $X$, $Y$ soient $W$-cocartésiennes, et les $f_i$ ($i \in \{0,1,2,3\}$) des $W$-cofibrations, alors le carré $Z$ des cofibres des $f_i$ est lui aussi $W$-cocartésien.}

Ce lemme résulte de la théorie du `déploiement des diagrammes cubiques', relatif à une catégorie de modèles $(\mathcal{M}, W)$ admettant une théorie des cofibres $W$-homotopiques, qu'il me faudra bien écrire! (Mais l'existence pour $(\mathrm{Cat}, W)$ d'une telle théorie des cofibres $W$-homotopiques, et par suite de celle des déploiements des diagrammes dans $\mathcal{M}$ de type $(\boldsymbol{\Delta}^1)^I$, $I$ ensemble d'indices fini, ne dépend que de W(7 bis) - on n'a pas à faire appel à W(7).)

\msp{page 269}
Ce `deploiement' est un diagramme infini dont toutes les suites infinies, dans les trois directions des axes de coordonnées, sont des `suites de suspension'. Je reproduis la partie utile du déploiement. % [sic] deploiement -> déploiement, missing accent


\endgroup
% END UNIT Der-ChXIIc-30-03-03/Der-ChXIIc-30-03-03_p056-060.tex
% BEGIN UNIT Der-ChXIIc-30-03-03/Der-ChXIIc-30-03-03_p061-064.tex
\begingroup\makeatletter%
\newcommand{\printedpage}[1]{\marginpar{\footnotesize p.~#1}}%
\newcommand{\msp}[1]{\marginpar{\footnotesize [#1]}}%
\newcommand{\Asph}{\mathrm{Asph}}%
\newcommand{\Cof}{\mathrm{Cof}}%
\makeatother\let\printedpage\derpp\let\msp\dermsp%
\let\mspin\DERmspin%


\printedpage{175}
% head: A. Grothendieck : Dérivateurs, Ch. XII, Édition des Universités de Montpellier II et Paris VII
% foot: Version 30 mars 2003

% dangling arrows (from X_2, X_3, Y_2, Y_3, Z_2, Z_3) as printed
\[
\begin{tikzcd}[ampersand replacement=\&, column sep=small, row sep=small]
 \& \& {} \& \& {} \& \& U \arrow[dd, "\mu"] \\
 \& X_2 \arrow[rr] \arrow[dd] \arrow[ur, shorten >=1ex] \& \& X_3 \arrow[rr] \arrow[dd] \arrow[ur, shorten >=1ex] \& \& T \arrow[ur] \arrow[dd] \& \\
X_0 \arrow[rr, crossing over] \arrow[ur] \arrow[dd] \& \& X_1 \arrow[rr, crossing over] \arrow[ur] \arrow[dd, crossing over] \& \& S \arrow[ur, "\lambda"] \arrow[dd, crossing over] \& \& U' \arrow[dd] \\
 \& Y_2 \arrow[rr] \arrow[dd] \arrow[ur, shorten >=1.5ex] \& \& Y_3 \arrow[rr] \arrow[dd] \arrow[ur, shorten >=1.5ex] \& \& T' \arrow[ur] \arrow[dd] \& \\
Y_0 \arrow[rr, crossing over] \arrow[ur] \arrow[dd] \& \& Y_1 \arrow[rr, crossing over] \arrow[ur] \arrow[dd, crossing over] \& \& S' \arrow[ur, "\lambda'"] \arrow[dd, crossing over] \& \& U'' \\
 \& Z_2 \arrow[rr] \arrow[ur, shorten >=1.5ex] \& \& Z_3 \arrow[rr] \arrow[ur, shorten >=1.5ex] \& \& T'' \arrow[ur] \& \\
Z_0 \arrow[rr] \arrow[ur] \& \& Z_1 \arrow[rr] \arrow[ur] \& \& S'' \arrow[ur, "\lambda''"] \& \&
\end{tikzcd}
\]

La condition que $X$ soit $W$-cocartésien s'exprime par $\lambda \in W$, ou encore par $U \in \Asph_W$ (car $U$ s'identifie à la cofibre homotopique $Z(\lambda)$ de $\lambda$ \dots\,), de même $Y$ $W$-cocartésien s'exprime par $U' \in \Asph_W$, et $Z$ $W$-cocartésien par $U'' \in \Asph_W$. Or $U'' \sim Z(\mu)$, où $\mu : U \longrightarrow U'$ est dans $W$ si $U$, $U'$ sont $W$-asphériques, donc $U'' = Z(\mu)$ est alors $W$-asphérique. Cela prouve le lemme.

Cela montre que, moyennant W(7 bis) (en fait, on a utilisé seulement la partie ‘universelle’ de W(7 bis)), savoir que dans
\begin{tikzcd}[ampersand replacement=\&, column sep=small, row sep=small, baseline=(\tikzcdmatrixname-2-2.base)]
T_0 \arrow[rr] \arrow[dr, "i"'] \& \& T_2 \arrow[dr, "i'"] \& \\
 \& T_1 \arrow[rr] \& \& T_3
\end{tikzcd}
diagramme cocartésien d'immersions ouvertes, $i \in W \Longrightarrow i' \in W$ (au lieu de $\Longleftrightarrow$),

\msp{page 270}

la validité de W(7) se ramène au cas où le carré $Y$ est le carré trivial
\begin{tikzcd}[ampersand replacement=\&, column sep=small, row sep=small, baseline=(\tikzcdmatrixname-2-1.base)]
 \& e \arrow[dr] \& \\
e \arrow[ur] \arrow[dr] \& \& e, \\
 \& e \arrow[ur] \&
\end{tikzcd}
$e$ l'objet final de Cat.

On peut maintenant conclure : On a (moyennant la partie directe de W(7 bis))
\[
\mathrm{a}' \ \& \ \mathrm{b}' \Longrightarrow \mathrm{c}',
\]
car si $Z_0$, $Z_1$ sont $\in \Asph_W$, alors $i_1 \in W$, donc $j_1 \in W$, et si de plus $Z_2 \in \Asph_W$, on a donc $Z_3 \in \Asph_W$.
\[
\begin{tikzcd}[ampersand replacement=\&, column sep=small, row sep=small]
 \& Z_1 \arrow[dr, "j_2"] \& \\
Z_0 \arrow[ur, "i_1"] \arrow[dr, "i_2"'] \& \& Z_3 \\
 \& Z_2 \arrow[ur, "j_1"'] \&
\end{tikzcd}
\]

Si on a aussi la partie inverse (‘inhabituelle’) de W(7 bis), alors
\[
\mathrm{b}' \ \& \ \mathrm{c}' \Longrightarrow \mathrm{a}',
\]
car si $Z_1$ et $Z_3$ sont $W$-asphériques, alors $j_2 \in W$, donc $i_2 \in W$, et si de plus $Z_2$ est $W$-asphérique, on conclut $Z_0$ $W$-asphérique.

\printedpage{176}
% head: A. Grothendieck : Dérivateurs, Ch. XII, Édition des Universités de Montpellier II et Paris VII
% foot: Version 30 mars 2003

Enfin, supposons que $W$ soit \textit{stable par facteurs directs,} et W(7 bis) satisfait sous forme directe. Je dis qu'alors
\[
\mathrm{c}' \ \& \ \mathrm{a}' \Longrightarrow \mathrm{b}',
\]
i.e.
\[
Z_0 \text{ et } Z_3 \ \in \Asph_W \qquad\quad \Longrightarrow \qquad\quad Z_1 \text{ et } Z_2 \ \in \Asph_W.
\]
Factorisant $i_1$ et $i_2$ en
\[
\begin{aligned}
Z_0 &\xrightarrow{i'_1} Z'_1 \xrightarrow{\alpha_1} Z_1, \\
Z_0 &\xrightarrow{i'_2} Z'_2 \xrightarrow{\alpha_2} Z_2,
\end{aligned}
\]
avec $i'_1$, $i'_2$ des immersions ouvertes de Dwyer, et $\alpha_1, \alpha_2 \in W$, on est ramené au cas où

\msp{page 271}

$i_1$, $i_2$ sont des \textit{immersions ouvertes de Dwyer,} et où le carré est cocartésien (ce qui implique qu'il est $W$-cocartésien). Faisant alors le cochangement de base par $Z_0 \xrightarrow{p_0} Z'_0 = e$, on trouve un homomorphisme de carrés cocartésiens
\[
\begin{tikzcd}[ampersand replacement=\&, column sep=small, row sep=small]
 \& Z_1 \arrow[rr] \arrow[dd, "p_1" near end] \& \& Z_3 \arrow[dd, "p_3"] \\
Z_0 \arrow[ur] \arrow[rr, crossing over] \arrow[dd, "p_0"'] \& \& Z_2 \arrow[ur] \& \\
 \& Z'_1 \arrow[rr] \& \& Z'_3, \\
e = Z'_0 \arrow[ur] \arrow[rr] \& \& Z'_2 \arrow[ur] \arrow[from=uu, crossing over, "p_2"' near end] \&
\end{tikzcd}
\]
où toutes les faces du cube sont $W$-cocartésiennes, et comme $p_0 \in W$ ($Z_0$ étant $W$-asphérique), les autres flèches verticales $p_1$, $p_2$, $p_3$ sont \texttt{[dans]} $W$, donc le carré $Z'_*$ satisfait aussi $Z'_0$ et $Z'_3$ \texttt{[}$W$\texttt{-]}asphériques (et de plus $Z'_0 = e$), et il faut prouver $Z'_1$, $Z'_2$ $W$-asphériques (d'où itou pour $Z_1$, $Z_2$). Donc, on est ramené au cas particulier où $Z_0 = e$.

Mais dans ce cas il existe une rétraction (triviale!) de $Z_1$ sur $Z_0$, donc aussi de $Z_3$ sur $Z_2$. Ainsi, $Z_2$ est un facteur

\msp{page 272}

direct d'un objet $W$-asphérique, ce qui implique, par l'hypothèse sur $W$, que $Z_2$ est $W$-asphérique. Par symétrie, $Z_1$ l'est aussi. On a gagné!

\textbf{NB} En fait, ces arguments sont généraux, et donnent ceci :

\setcounter{footnote}{136}
\textbf{Théorème.} \textit{Soit $(\mathcal{M}, W)$ une catégorie des modèles, admettant une théorie de la cofibre $W$-homotopique \textup{(}\footnote{Mais on ne sait vérifier l'existence d'une telle théorie dans $(\mathrm{Cat}, W)$ que si on admet la forme faible de W(7), i.e. $f_0, f_1, f_2 \in W \implies f_3 \in W$, ou quelque cas particulier de celle-ci. Mais le \textit{lemme de Brown} sauve la situation, cf. XV \S3.}\textup{)} (i.e. on a une factorisation de toute flèche $f : X \longrightarrow Y$ en $pi$, avec $i \in \Cof_W$, $p \in W$.) On a donc une notion de carré $W$-cocartésien dans $\mathcal{M}$ \textup{(}\footnote{à vérifier que la factorisation suffit, même si tous les objets ne sont pas fibrants sur $e$. Mais on peut alors remplacer $\mathcal{M}$ par $\mathcal{M}_f$, et développer ce formalisme des fibres homotopiques en conséquence \dots}\textup{)}. Soit
\printedpage{177}
% head: A. Grothendieck : Dérivateurs, Ch. XII, Édition des Universités de Montpellier II et Paris VII
% foot: Version 30 mars 2003
alors $f : X \longrightarrow Y$ un homomorphisme d'un carré $W$-cocartésien $X$ dans un autre $Y$, $f = (f_0, f_1, f_2, f_3)$.}
\[
\begin{tikzcd}[ampersand replacement=\&, column sep=small, row sep=small]
 \& X_2 \arrow[rr] \arrow[dd, "f_2" near end] \& \& X_3 \arrow[dd, "f_3"] \\
X_0 \arrow[ur] \arrow[rr, crossing over] \arrow[dd, "f_0"'] \& \& X_1 \arrow[ur] \& \\
 \& Y_2 \arrow[rr] \& \& Y_3 \\
Y_0 \arrow[ur] \arrow[rr] \& \& Y_1 \arrow[ur] \arrow[from=uu, crossing over, "f_1"' near end] \&
\end{tikzcd}
\]

\textit{Considérons les conditions sur $f$}

\begin{enumerate}
\item[\textit{a)}] \textit{$f_0 \in W$.}
\item[\textit{b)}] \textit{$f_1$ et $f_2 \in W$.}
\item[\textit{c)}] \textit{$f_3 \in W$.}
\end{enumerate}

\textit{Alors}

\begin{enumerate}
\item[\textit{1°)}] \textit{a \& b $\Longrightarrow$ c.}

\msp{page 273}

\item[\textit{2°)}] \textit{Supposons que $\mathcal{M}$ ait un objet final $e$. Alors l'implication}
\[
\textit{b \& c} \Longrightarrow \textit{a}
\]
\textit{pour}~tout \textit{homomorphisme $f : X \longrightarrow Y$ de carrés $W$-cocartésiens équivaut à l'axiome suivant sur $W$ :}

Axiome des carrés $W$-cocartésiens conservatifs (\footnote{cf. argument page 78 qui montre que l'implication b \& c $\Longrightarrow$ a entraine l'axiome conservatif.}) : \textit{Si}
\[
\begin{tikzcd}[ampersand replacement=\&, column sep=small, row sep=small]
 \& Z_1 \arrow[dr, "j_2"] \& \\
Z_0 \arrow[ur, "i_1"] \arrow[dr, "i_2"'] \& \& Z_3 \\
 \& Z_2 \arrow[ur, "j_1"'] \&
\end{tikzcd}
\]
\textit{est $W$-cocartésien, alors $j_1 \in W \Longrightarrow i_1 \in W$ (donc en fait $i_1 \in W \Longleftrightarrow j_1 \in W$ et $i_2 \in W \Longleftrightarrow j_2 \in W$).}

\textit{Cet axiome est équivalent à la conjonction des deux axiomes suivants.}

\textbf{Cons 1.} \textit{Si on a un carré cocartésien}
\printedpage{178}
% head: A. Grothendieck : Dérivateurs, Ch. XII, Édition des Universités de Montpellier II et Paris VII
% foot: Version 30 mars 2003
\[
\begin{tikzcd}[ampersand replacement=\&, column sep=small, row sep=small]
S \arrow[rr] \arrow[dr] \& \& T \arrow[dr] \& \\
 \& e \arrow[rr] \& \& U,
\end{tikzcd}
\]
\textit{avec $f \in \Cof_W$, alors $U \in \Asph_W \Longrightarrow f \in W$ (donc $U \in \Asph_W \Longleftrightarrow f \in W$).}

\textbf{Cons 2.} \textit{Si on a un carré cartésien,}
\[
\begin{tikzcd}[ampersand replacement=\&, column sep=small, row sep=small]
 \& Z_1 \arrow[dr] \& \\
Z_0 \arrow[ur, "i_1"] \arrow[dr, "i_2"'] \& \& Z_3, \\
 \& Z_2 \arrow[ur] \&
\end{tikzcd}
\]
\textit{avec $i_1, i_2 \in \Cof_W$, alors : Si $Z_1$, $Z_2$, $Z_3$ sont $W$-asphériques, $Z_0$ l'est aussi.}

\msp{page 274}

\item[\textit{3°)}] \textit{Supposons que tout facteur direct d'un objet $W$-asphérique soit $W$-asphérique, ce qui est le cas notamment si $W$ est lui-même stable par facteurs directs, a fortiori si $W$ est fortement saturé. Supposons de plus que $\mathcal{M}$ ait} \texttt{[un]} \textit{objet final. Alors on a (sous les conditions générales du début)}
\[
\textit{c \& a} \Longrightarrow \textit{b,}
\]
\textit{i.e. si $f_0$, $f_3$ sont} \texttt{[dans]} \textit{$W$, de même $f_1$, $f_2$.}
\end{enumerate}

Cette dernière partie peut se déployer ainsi :

\textbf{Corollaire 1.} \textit{Sous l'hypothèse supplémentaire} \textup{3\,°)} \textit{sur $W$, la validité de c \& a $\Longrightarrow$ b (pour toute flèche $f$ entre carrés $W$-cocartésiens) équivaut à chacun des deux autres axiomes suivants sur $W$ (qui sont donc équivalents) :}

\begin{enumerate}
\item[\textit{1)}] \textit{Si on a un carré $W$-cocartésien}
\[
\begin{tikzcd}[ampersand replacement=\&, column sep=small, row sep=small]
 \& Z_1 \arrow[dr] \& \\
Z_0 \arrow[ur] \arrow[dr] \& \& Z_3, \\
 \& Z_2 \arrow[ur] \&
\end{tikzcd}
\]
\textit{alors $Z_0$ et $Z_3$ $W$-asphériques $\Longrightarrow$ $Z_1$ et $Z_2$ $W$-asphériques.}
\item[\textit{1$'$)}] \textit{Comme 1), avec $Z_0 = e$ (dans le cas où on suppose que $\mathcal{M}$ admet} \texttt{[un]} \textit{objet final $e$).}
\end{enumerate}

\msp{page 275}

\textbf{Corollaire 2.} \textit{Supposons que $W$ satisfasse l'axiome des carrés $W$-cocartésiens conservatifs, que tout facteur direct d'un objet $W$-asphérique soit $W$-asphérique (p. ex. $W$ stable par facteurs directs), enfin que $\mathcal{M}$ admette un objet final. Alors pout tout homomorphisme $f : X \longrightarrow Y$ de carrés $W$-cocartésiens, la conjonction de deux parmi les conditions a) b) c) de p. 272 entraîne la troisième.} % [sic] pout -> pour, typographical error


\endgroup
% END UNIT Der-ChXIIc-30-03-03/Der-ChXIIc-30-03-03_p061-064.tex
\clearpage\setcounter{footnote}{0}\phantomsection\addcontentsline{toc}{section}{Chapitre XII, appendice}\markboth{Chapitre XII, appendice}{Chapitre XII, appendice}
% BEGIN UNIT Der-ChXIIapp-19-06-03/Der-ChXIIapp-19-06-03_p001-005.tex
\begingroup\makeatletter%
\newcommand{\printedpage}[1]{\marginpar{\footnotesize p.~#1}}%
\newcommand{\msp}[1]{\marginpar{\footnotesize [#1]}}%
\newcommand{\Hom}{\mathrm{Hom}}%
\newcommand{\Fl}{\mathrm{Fl}}%
\newcommand{\Ob}{\mathrm{Ob}}%
\newcommand{\Cat}{\mathrm{Cat}}%
\newcommand{\Ens}{\mathrm{Ens}}%
\newcommand{\pr}{\mathrm{pr}}%
\newcommand{\Ker}{\mathrm{Ker}}%
\newcommand{\Coker}{\mathrm{Coker}}%
\newcommand{\mon}{\mathrm{mon}}%
\newcommand{\deq}{\mathrel{\overset{\text{déf}}{=}}}%
\makeatother\let\printedpage\derpp\let\msp\dermsp%
\let\mspin\DERmspin%


\printedpage{}
% no page number printed (title page)
% head: (none)

\begin{center}
{\LARGE\textsf{LES DÉRIVATEURS}}

\bigskip

\textbf{ALEXANDRE GROTHENDIECK}

\bigskip\bigskip

{\large\textbf{Appendice au Chapitre XII}}
\end{center}

\par

\printedpage{1}
% head: A. Grothendieck : Dérivateurs, Ch. XII, Édition des Universités de Montpellier II et Paris VII
% foot: Version 20 juin 2003

\msp{page 1}
\textbf{Proposition} (?) (\footnote{Cf. pages 11 ff. pour les contrexemples-clef.}). \textit{Soit $X$ dans} $\Cat$\textit{, avec $b$, $c$ dans $X$ tels que}

\begin{itemize}
\item[\textit{a)}] \textit{$b\times c$ est dans $X$.}
\item[\textit{b)}] \textit{$\forall$ $x$ dans $X$, $b(x)$ ($\deq$} $\Hom_X(x,b)$\textit{) et $c(x)$ sont $\neq\emptyset$, i.e. $X/b \longrightarrow X$ et $X/c \longrightarrow X$ sont surjectifs sur les objets.}
\end{itemize}

\textit{Alors $X$ est $\mathbf{W}_\infty$-asphérique.} (\footnote{Faux, cf. plus bas. $X$ est 1-connexe, mais pas nécessairement 2-connexe.})

Le fait que $b$ soit un plus grand élément de $\Ob X$ assure déjà la 0-connexité.

Pour la 1-connexité, prouvons que dans le groupoïde fondamental $\Pi(X)$, $a \deq b\times c$ devient un élément final.
\[
\begin{tikzcd}
b\times c \arrow[r] \arrow[d] & b \\
c &
\end{tikzcd}
\]

Pour ceci, prouvons d'abord que $\forall$ $x$ dans $X$, tous les éléments de $a(x) = b(x)\times c(x)$ sont égaux dans $\Fl\Pi(X)$. En effet, comme $b\times c \longrightarrow b$ est \texttt{[un]} isomorphisme, donc \texttt{[un]} monomorphisme dans $\Pi(X)$, cela montre déjà que si $\beta,\beta'\in b(x)$, $\gamma,\gamma'\in c(x)$, d'où $(\beta,\gamma)$, $(\beta,\gamma')$, $(\beta',\gamma')\in b(x)\times c(x)$, on a $(\beta,\gamma) = (\beta,\gamma')$ dans $\Pi(X)$ (puisque $\pr_1(\beta,\gamma) = \pr_1(\beta,\gamma')$, et de même (en utilisant que $\pr_2$ est \texttt{[un]} isomorphisme, donc \texttt{[un]} monomorphisme) $(\beta,\gamma') = (\beta',\gamma')$). Donc tous les $(\beta,\gamma)\in b(x)\times c(x) = a(x)$ sont égaux dans $\Fl\Pi(X)$. Pour achever de prouver que $a$ est objet final de $\Pi(X)$, il suffit de voir que toute flèche \textit{dans} $\Pi(X)$ d'un $x\in\Ob X$ dans $a$, provient d'une flèche $x\longrightarrow a$ dans $X$. Mais une flèche dans $\Pi(X)$ est obtenue par un chemin dans $X$
\[
\begin{tikzcd}
x = x_0 \arrow[r, "u_1"] \arrow[dr, "\alpha_0"] & x_1 \arrow[d, "\alpha_1"'] & x_2 \arrow[l, "u_2"'] \arrow[r] \arrow[dl, "\alpha_2"'] & {} & \cdots & {} & x_{2N} = a \arrow[l, "u_{2N}"'] \\
& a. & & & & &
\end{tikzcd}
\]

Prenons une flèche $\alpha_0 : x\longrightarrow a_0$ dans $X$, je dis que la % [sic] a_0 -> a, the arrows α_i go to a (no a_0 is defined)
\msp{page 2}
classe d'homotopie du chemin $c$ (i.e. son image dans $\Fl(\Pi(X))$ est égale à celle de $u_0$. % [sic] u_0 -> \alpha_0 (?), no u_0 is defined; also the parenthesis opened at "(i.e." is not closed
Pour le voir, choisissons pour tout $i\in[0,2N]$ une flèche $\alpha_i : x_i\longrightarrow a$. (Ça existe grâce à l'hypothèse b)). Je dis que pour chaque $i\in[0,2N]$, le triangle
\printedpage{2}
% head: A. Grothendieck : Dérivateurs, Ch. XII, Édition des Universités de Montpellier II et Paris VII
% foot: Version 20 juin 2003
\[
\begin{tikzcd}[column sep=small]
x_0 \arrow[rr, "c_i"] \arrow[dr, "\alpha_0"] & & x_i \arrow[dl, "\alpha_i"'] \\
& a &
\end{tikzcd}
\]
dans $\Pi(X)$ est commutatif, où $c_i$ désigne la classe d'homotopie du chemin de $x_0$ à $x_i$ définie par les $u_i$ ($i\in[1,i]$) \texttt{[plutôt} ($i\in[1,2N]$)\texttt{]}. Nous le faisons par récurrence sur $i$, l'assertion est triviale si $i = 0$ ($c_0$ représente l'application identitique). Le diagramme % [sic] identitique -> identique, non-existent word
\[
\begin{tikzcd}[column sep=small]
x_0 \arrow[rr, "c_i"] \arrow[dr, "\alpha_0"] \arrow[rrrr, bend left=32, "c_{i+1}"] & & x_i \arrow[dl, "\alpha_i"'] \arrow[rr, "u_{i+1}"] & & x_{i+1} \arrow[dlll, "\alpha_{i+1}"'] \\ % verified: solid curved arrow; bend increased for label clearance
& a & & &
\end{tikzcd}
\]
montre qu'il suffit de prouver que pour $i\in[0,2N-1]$ le triangle
\[
\begin{tikzcd}[column sep=small]
x_i \arrow[rr, "u_{i+1}"] \arrow[dr, "\alpha_i"'] & & x_{i+1} \arrow[dl, "\alpha_{i+1}"] \\
& a &
\end{tikzcd}
\]
est commutatif dans $\Pi(X)$. (\textbf{NB} $u_{i+1}$ est en sens directe ou rétrograde, suivant que % [sic] directe -> direct, agreement with "sens"
$i$ est paire ou impaire.) Mais cela résulte du fait que l'image de $\Hom_X(x_i,a)$ (resp. $\Hom_X(x_{i+1},a)$) % [sic] paire ou impaire -> pair ou impair, agreement with "i"
\msp{page 3}
dans $\Fl\Pi(X)$ est réduite à un seul élément.

Ainsi on trouve que $X$ \textit{est} 1-\textit{connexe.} Il y a donc de l'espoir! Il reste à prouver
\[
H^i(X,M) = 0, \qquad i\geq 2, \quad M \text{ coefficient constant.}
\]

J'ai regardé le cas lorsqu'on a ceci : $b = c$, $\Ob X = \{b, b\times b\}$, $\Hom_X(b\times b)$ \texttt{[plutôt} $\Hom(b\times b, b)$\texttt{]} $= \{$ réunion disjointe des ensembles $u\circ\pr_1$ et $u\circ\pr_2$, avec $u\in\Hom_X(b,b) = G$ un monoïde quelconque $\}$. Alors $X^\circ$ est fibrée, via $I\longmapsto b^I$, sur la sous-catégorie pleine $S$ de $\Ens$ formée de $\{0\}$ et $\{0,1\}$, et donc $X$ est cofibrée (en particulier propre) sur $S^\circ \deq T$. (NB $S^\circ = T$ a un objet initial, mais non un objet final. Sa cohomologie à valeurs dans un préfaisceau quelconque n'est pas toujours triviale, cf. ci-dessous.)

On trouve \texttt{[`On trouve' barré]}
\[
H^n(X,M) \Longleftarrow H^p(T, \underbrace{R^q f_*(M_X)}_{= \mathcal{H}^q}), \qquad \text{où } f : X \longrightarrow S^\circ
\]
\texttt{[}$n = p+q$\texttt{]}. J'admet ici (provisoirement) ceci : % [sic] J'admet -> J'admets, verb ending
\[
\begin{aligned}
H^p(T,F) &= 0 \ , &&\text{pour } p\geq 2, \text{ tout } F\in T^\wedge \ ,\\
H^0(T,F) &= \Ker\delta \ , &&\\
H^2(T,F) &= \Coker\delta \ , && % [sic] H^2(T,F) -> H^1(T,F), the first line gives H^p(T,F) = 0 for p >= 2
\end{aligned}
\]
\printedpage{3}
% head: A. Grothendieck : Dérivateurs, Ch. XII, Édition des Universités de Montpellier II et Paris VII
% foot: Version 20 juin 2003
où $\delta : F_0 \longrightarrow F_1$ est défini comme la différence $\delta = p_1^* - p_2^*$, où $p_1$, $p_2$ sont les deux projections de $\{0,1\}\simeq\{0\}^2$ dans $T = S^\circ$. On a de plus
\[
(\mathcal{H}^q)_0 = H^q(G,M), \qquad (\mathcal{H}^q)_1 = H^q(G\times G, M),
\]
et
\[
p_1^*,\ p_2^*\ :\ H^q(G,M) \rightrightarrows H^q(G{\times}G, M) \qquad \text{sont les transposés de } \pr_1,\ \pr_2 : G\times G\longrightarrow G.
\]

\msp{page 4}
On vérifie aisément que
\[
H^0(T,\mathcal{H}^q) = 0 \qquad \text{si } q\geq 1,
\]
et il reste
\[
H^n(X,M) \simeq H^1(S,\mathcal{H}^{n-1}) = \Coker(\delta : H^{n-1}(G) \longrightarrow H^{n-1}(G\times G)). % [sic] H^1(S, ...) -> H^1(T, ...), the cohomology is taken on T = S°
\]
Cela donne, pour $n = 2$,
\[
\begin{aligned}
H^2(X,M) &\simeq \Coker\Big(\delta : \Hom_\mon(G,M) \longrightarrow \underbrace{\mathit{Hom}_\mon(G\times G, M)}_{\Hom_\mon(G,M)\times\Hom_\mon(G,M)} \ \Big)\\
&\simeq \Coker\Big(u\longmapsto (u,-u) : H \longrightarrow H\times H\Big),
\end{aligned}
\]
où $H$ est le groupe abélien $\Hom_\mon(G,M)$.

D'où
\[
H^2(X,M) \ \simeq\ H = \Hom_\mon(G,M).
\]
Donc
\begin{quote}
$X$ est 2-connexe si et seulement si le groupe \textit{abélien} enveloppant du monoïde $G$ est nul.
\end{quote}

Ainsi, si $G$ est un groupe abélien $\neq 0$, $X$ \textit{n'est pas} 2-connexe.

\textbf{Corollaire.} $\Psi = {}$\begin{tikzcd}[ampersand replacement=\&, row sep=tiny, column sep=small, baseline=(A.base)]
\& b_0 \\
|[alias=A]| a_0 \arrow[ur] \arrow[dr] \& \\
\& c_0
\end{tikzcd}
\textit{n'est} pas \textit{totalement} $\mathbf{W}_\infty$\textit{-asphérique; i.e. si on a un diagramme} \begin{tikzcd}[ampersand replacement=\&, row sep=tiny, column sep=small, baseline=(B.base)]
\& b \\
|[alias=B]| a \arrow[ur] \arrow[dr] \& \\
\& c
\end{tikzcd}
\textit{dans une catégorie} $X\in\Ob\Cat$\textit{, tel que sa somme amalgamée dans $X^\wedge$ soit égale à $e_X$, il ne s'ensuit pas que $X$ soit $\mathbf{W}_\infty$-asphérique.} (Mais peut-être voudra bien être 1-connexe?)

\printedpage{4}
% head: A. Grothendieck : Dérivateurs, Ch. XII, Édition des Universités de Montpellier II et Paris VII
% foot: Version 20 juin 2003

\msp{page 5}
\textbf{Question.} \textit{Y-a-t-il d'autres objets totalement} $\mathbf{W}_\infty$\textit{-asphériques de} $\Cat$\textit{, que les catégories ayant un objet final?} Je suspecte à présent que non -- donc que l'on ne peut améliorer substantiellement sur la notion déjà connue d'immersion ouverte de Dwyer.

Dans le diagramme
\[
\begin{tikzcd}
U \arrow[r, hook, "i"] \arrow[d, "p"'] \arrow[dr, phantom, "\substack{\text{cocart.}\\ \text{et cart.}}"] & X \arrow[d, "q"'] & T \arrow[l, hook', dotted, "j"'] \arrow[d, dotted, "\text{iso}"] \\
U' \arrow[r, hook, "i'"] & X' & T', \arrow[l, hook', dotted, "j'"']
\end{tikzcd}
\tag*{$(*)$}
\]
prenons
\[
\begin{aligned}
U = \Psi, \quad & T = \{t\} = e \text{ catégorie finale,}\\
& t \text{ objet final de } X,\\
& \text{i.e. } H_t = (u\longmapsto \Hom_X(u,t)) = e_U \in \Ob U^\wedge \ \mbox{(\footnotemark)}.
\end{aligned}
\]
\footnotetext{i.e. $X$ déduit de $U = \Psi$ en lui rajoutant un objet final $t$ -- donc $X\simeq\boldsymbol{\Delta}^1\times\boldsymbol{\Delta}^1$,
\[
\begin{tikzcd}[ampersand replacement=\&,
execute at end picture={
\draw[rounded corners] ([xshift=-10pt,yshift=10pt]A.center) -- ([xshift=-10pt,yshift=-10pt]C.center) -- ([xshift=10pt,yshift=-10pt]D.center) -- ([xshift=10pt,yshift=3pt]D.center) -- ([xshift=3pt,yshift=3pt]C.center) -- ([xshift=3pt,yshift=10pt]A.center) -- cycle;
\node (P) at ([xshift=-35pt,yshift=-22pt]C.center) {$\Psi$};
\draw[-stealth] (P) -- ([xshift=-11pt,yshift=-4pt]C.center);
}]
|[alias=A]| {} \arrow[r] \& {} \\
|[alias=C]| {} \arrow[u] \arrow[r] \& |[alias=D]| {} \arrow[u]
\end{tikzcd}
\ .
\]}
Prenons pour $U'$ la catégorie dont il est question dans la `proposition' hypothétique (et fausse) p. 1 avec $b = c$, plus précisément une catégorie avec deux objets $b$ et $b\times b$, un groupe (voire un monoïde) $G$ d'automorphismes \texttt{[(voire d'endomorphismes)]} de $b$, et les autres flèches celles qui s'en déduisent, en exigeant que $b\times b$ soit bel et bien le carré cartésien de $b$. On trouve une catégorie cofibrée (donc propre) sur la catégorie $T$ opposée de celle des ensembles finis de la forme $\{0\}\cup\{0,1\}$ (deux objets) avec comme fibres $B_G$ et $B_{G\times G}$ respectivement ($B_\pi =$ catégorie avec un seul objet, et monoïde d'endomorphismes $\pi$). Le foncteur
\[
U = \Psi \longrightarrow U'
\]
est celui donné par le diagramme

\msp{page 6}
\[
\begin{tikzcd}[row sep=small]
& b \\
b\times b \arrow[ur, "\pr_1"] \arrow[dr, "\pr_2"'] & \\
& b
\end{tikzcd}
\]


\endgroup
% END UNIT Der-ChXIIapp-19-06-03/Der-ChXIIapp-19-06-03_p001-005.tex
% BEGIN UNIT Der-ChXIIapp-19-06-03/Der-ChXIIapp-19-06-03_p006-010.tex
\begingroup\makeatletter%
\tikzcdset{every diagram/.append style={baseline={([yshift=-.5ex]current bounding box.center)}}}%
\newcommand{\printedpage}[1]{\marginpar{\footnotesize p.~#1}}%
\newcommand{\msp}[1]{\marginpar{\footnotesize [#1]}}%
\newcommand{\Hom}{\mathrm{Hom}}%
\newcommand{\id}{\mathrm{id}}%
\newcommand{\Winf}{\mathbf{W}_{\infty}}%
\makeatother\let\printedpage\derpp\let\msp\dermsp%
\let\mspin\DERmspin%


\printedpage{5}
% head: A. Grothendieck : Dérivateurs, Ch. XII, Édition des Universités de Montpellier II et Paris VII
% foot: Version 20 juin 2003

dans $U'$, et les hypothèses a), b) de la page 1 assurent que la somme amalgamée $\varinjlim$ de ce diagramme dans $U'^{\wedge}$ est $e_{U'}$, i.e. que
\[ p_!(e_U) = e_{U'}. \]
Il s'ensuit que dans la somme amalgamée, $X'$ envisagée dans $(*)$, l'unique élément $t'$ de $T'$ est objet final de $X'$. On a vu que $U'$ est 1-connexe, mais non 2-connexe si le groupe abélien enveloppant de $G$ (i.e. $G/$(groupe des commutateurs) dans le cas où $G$ est un groupe) n'est pas réduit à 1. L'exemple le plus économique s'obtient en prenant $G = \{\pm 1\}$ le groupe à deux éléments. C'est donc un cas où dans $(*)$ on a
\[ i \in W \ \ \text{(et même $i$ asphérique),} \qquad \text{mais } i' \notin W, \]
et qui a l'avantage sur l'exemple plus sophistiqué de la page 103 (prop. 10) de ne mettre en jeu que des catégories finies -- et même fort petites !

Procédons maintenant comme dans la proposition 11, en désignant par $U''$ la catégorie déduite de $U'$ en lui ajoutant un objet
\msp{page 7}
final, et soit
\[ p' : U' \longrightarrow U'' \]
l'inclusion, de sorte que
\[
\begin{aligned}
p'_!(e_{U'}) \quad &= \quad (p'p)_!(e_U) = U' \quad && \text{en tant qu'objet dans } U''^{\wedge}\\
& && \text{(sous-préfaisceau du préfaisceau final } e_{U''}\text{).}
\end{aligned}
\]
Considérons
\[
\begin{tikzcd}
U \arrow[r, hook, "i"] \arrow[d, "{p''=p'p}"'] \arrow[dr, phantom, "\begin{smallmatrix}\text{cocart.}\\ \text{et cart.}\end{smallmatrix}"] & X \arrow[d, "q''"'] & T \arrow[l, hook', dotted, "j"'] \arrow[d, dotted, "\text{iso}"] \\
U'' \arrow[r, hook, "i''"] & X'' & T''. \arrow[l, hook', dotted, "j''"']
\end{tikzcd}
\]
Cette fois $U$, $X$ \textit{et} $U''$ sont $\Winf$-asphériques (et même contractiles toutes trois), donc $i$, $p''$ sont dans $W$, pourtant $X''$ n'est pas $\Winf$-asphérique (donc $i'', q'' \notin \Winf$), car pour qu'elle le soit, i.e. pour que $i'' \in \Winf$, il faudrait (prop. 7, cor. page 85 (\footnote[4]{pas encore démontré !})) que $i''$ soit $\Winf$-asphérique, i.e. que $U''/t''$ soit $\Winf$-asphérique. Mais par construction $U''/t'' \simeq U'$, qui n'est \textit{pas} $\Winf$-asphérique.

Voilà un exemple encore plus simple d'un foncteur $f : \Psi \longrightarrow X$ tel que $f_!(e_{\Psi})$ ne soit pas $\Winf$-asphérique, i.e. d'un diagramme
\[
\begin{tikzcd}
& b \\
a \arrow[ur, "\beta_0"] \arrow[dr, "\gamma_0"'] & \\
& c
\end{tikzcd}
\]
\printedpage{6}
% head: A. Grothendieck : Dérivateurs, Ch. XII, Édition des Universités de Montpellier II et Paris VII
% foot: Version 20 juin 2003
\msp{page 8}
dans une catégorie $V$, telle que la limite inductive $L = b \amalg_a c$ \textit{dans} $V^{\wedge}$ soit un objet non asphérique de $V^{\wedge}$, i.e. $V/L$ soit non asphérique. Je sais qu'il faut, pour un tel contre-exemple, que $a \longrightarrow b$ et $a \longrightarrow c$ ne soient pas des monomorphismes. L'idée, c'est de prendre un exemple universel où ces flèches ne soient pas \texttt{[des]} monomorphismes, en prenant
\msp{page 9}
\[
V = \left(
\begin{tikzcd}
& & b \\
0 \arrow[r, shift left, "\lambda"] \arrow[r, shift right, "\mu"'] & a \arrow[ur, "\beta"] \arrow[dr, "\gamma"'] & \\
& & c
\end{tikzcd}
\right)
\]
avec
\[ \beta\lambda = \beta\mu \qquad \text{et} \qquad \gamma\lambda = \gamma\mu \]
les deux seules flèches non identiques, en plus des quatre flèches marquées $\beta$, $\gamma$, $\lambda$, $\mu$.

On trouve que
\[ b \amalg_a c \xrightarrow{\ \sim\ } e_X \qquad \text{(somme amalgamée dans } V^{\wedge}\text{).} \] % [sic] e_X -> e_V, the category considered here is V

Pourtant $V$ n'est pas asphérique, ni même 1-connexe : il est homotope à la sous-catégorie ouverte $V_0 = (0 \overset{\lambda}{\underset{\mu}{\rightrightarrows}} a)$, car on a une rétraction $V \xrightarrow{\ r\ } V_0$ (envoyant $b$ et $c$ sur $a$), et une flèche $ir \longrightarrow \id_V$ (où $i : X_0 \hookrightarrow X$ est l'inclusion). Donc $V$ a le même type d'homotopie \textit{que} $V_0$, qui est un cercle. Donc on a intérêt à prendre ce $V$, où $\tilde{V}$ déduit de $V$ en rajoutant un élément final $e$, % [sic] i : X_0 \hookrightarrow X -> i : V_0 \hookrightarrow V, the inclusion of V_0 in V is meant
\[
\tilde{V} = \left(
\begin{tikzcd}
& & b \arrow[dr] & \\
0 \arrow[r, shift left] \arrow[r, shift right] & a \arrow[ur] \arrow[dr] & & e \\
& & c \arrow[ur] &
\end{tikzcd}
\right),
\]
comme $U'$, pour les contre-exemples pages 6 et 7.

\msp{page 10}
\subsection*{2.6. Contre-exemples pour la totale asphéricité.}

Je me suis finalement convaincu que
\[
\Psi = \left\{
\begin{tikzcd}
& b_0 \\
a_0 \arrow[ur] \arrow[dr] & \\
& c_0
\end{tikzcd}
\right\}
\]
\printedpage{7}
% head: A. Grothendieck : Dérivateurs, Ch. XII, Édition des Universités de Montpellier II et Paris VII
% foot: Version 20 juin 2003
n'est \textit{pas} totalement asphérique (pour le localiseur $\Winf$ -- dans le présent numéro, on supposera $W = \Winf$), ou ce qui revient au même, qu'on peut avoir un diagramme de type $\Psi$
\[
\begin{tikzcd}
& b \\
a \arrow[ur, "\beta"] \arrow[dr, "\gamma"'] & \\
& c
\end{tikzcd}
\tag*{$(*)$}
\]
dans une catégorie $Z \in \mathrm{Ob}\,\mathrm{Cat}$, dont la limite inductive $b \amalg_a c$ dans $Z^{\wedge}$ soit le préfaisceau final $e_Z$, i.e. tel que l'on ait
\[ b(z) \amalg_{a(z)} c(z) \xrightarrow{\ \sim\ } \text{ensemble ponctuel} \]
pour tout objet $z$ de $Z$ (où $a(z) = \Hom_Z(a,z)$ \texttt{[plutôt} $\Hom_Z(z,a)$\texttt{]}) sans que pour autant $Z$ soit asphérique, ni même 1-connexe (\footnote[5]{faux -- il est 1-connexe forcément (\footnotemark[6]), mais non 2-connexe.}).\footnotetext[6]{cf. commentaire XII, p. 101.} L'exemple le plus simple que j'ai trouvé, est l'exemple ‘universel’ d'un diagramme $(*)$ où ni l'une ni l'autre des deux flèches $a \longrightarrow b$, $a \longrightarrow c$ n'est \texttt{[un]} monomorphisme (car si elle l'est, alors on peut prouver que $Z$ est bel et bien asphérique, comme nous verrons par la suite). C'est la catégorie
\[
Z = \left\{
\begin{tikzcd}
& & b_0 \\
0_0 \arrow[r, shift left, "\lambda_0"] \arrow[r, shift right, "\mu_0"'] & a_0 \arrow[ur, "\beta_0"] \arrow[dr, "\gamma_0"'] & \\
& & c_0
\end{tikzcd}
\right\}
\]
\msp{page 11}
déduite de $\Psi$ en lui rajoutant un objet $0_0$, relié à $a_0$ par deux flèches distinctes $\lambda_0$, $\mu_0$, telles que
\[ \beta_0\lambda_0 = \beta_0\mu_0 \qquad \text{et} \qquad \gamma_0\lambda_0 = \gamma_0\mu_0, \]
ces deux flèches composées étant, avec les quatre flèches $\beta_0$, $\gamma_0$, $\lambda_0$, $\mu_0$ marquées sur le diagramme, les seules (six) flèches non identiques de $Z$. $\Psi$ immergée dans $Z$ comme un fermé, de complémentaire la catégorie ponctuelle réduite à $0_0$. (NB L'immersion fermée $\Psi \longrightarrow Z$ n'est pas immersion de Dwyer, i.e. (comme $Z$ est le plus petit ouvert de $Z$ contenant $\Psi$) n'est pas coasphérique, car la catégorie $0\backslash\Psi$, cofibrée sur $\Psi$ avec fibres discrètes, de cardinal 2 en $a_0$ et 1 en $b_0$ et $c_0$, est la catégorie
\[
\begin{tikzcd}
& \tilde{b}_0 \\
a'_0 \arrow[ur] \arrow[ddr] & \\
a''_0 \arrow[uur, crossing over] \arrow[dr] & \\
& \tilde{c}_0
\end{tikzcd}
\]
\printedpage{8}
% head: A. Grothendieck : Dérivateurs, Ch. XII, Édition des Universités de Montpellier II et Paris VII
% foot: Version 20 juin 2003
i.e.
\[
\begin{tikzcd}
& {} & \\
{} \arrow[ur] \arrow[dr] & & {} \arrow[ul] \arrow[dl] \\
& {} &
\end{tikzcd}
\quad ,
\]
qui a le type d'homotopie du cercle.) Ni l'immersion ouverte $\{0_0\} \longrightarrow Z$, car $\{0_0\}/a_0$ est la catégorie discrète de cardinalité 2.

L'inclusion $i : Z_0 \longrightarrow Z$ de
\[ Z_0 \ = \ (\ 0_0 \overset{\lambda_0}{\underset{\mu_0}{\rightrightarrows}} a_0\ ) \]
dans $Z$ est un homotopisme (\footnote[7]{faux, il n'y a pas de rétraction du tout! En fait, $Z$ est la \textit{suspension} de la $\rightrightarrows$, qui a le type d'homotopie (topologique) du cercle, et il a donc le type d'homotopie de la 2-sphère.}) (car il existe une rétraction évidente $r$ de $Z$ sur $Z_0$ (\footnotemark[7]), et une flèche $ir \longrightarrow \id_Z$), donc $Z$ est homotope à $Z_0$, qui a le type d'homotopie du cercle, donc n'est pas 1-connexe.

\msp{page 12}
D'autre part, la vérification de
\[ b_0 \amalg_{a_0} c_0 \ = \ e_Z \qquad \text{dans } Z^{\wedge} \]
est immédiate.

On va utiliser cette flèche
\[ \Psi \longrightarrow Z \]
pour construire des exemples intéressants (et déchirants) de sommes amalgamées
\[
\begin{tikzcd}
U \arrow[r, hook, "i"] \arrow[d, "p"'] & X \arrow[d, "q"'] \\
U' \arrow[r, hook, "i'"] & X',
\end{tikzcd}
\]
avec $i$ immersion ouverte, telle qu'on ait même
\[ i :\ U \longrightarrow \overline{U} \qquad \text{est asphérique, donc dans } W, \]
alors que pourtant on aura
\[ i' :\ U' \longrightarrow \overline{U'} \qquad \notin W, \]
et ce qui mieux (ou pis) est,
\[ p \in W \qquad \text{mais} \qquad q \notin W \]
\printedpage{9}
% head: A. Grothendieck : Dérivateurs, Ch. XII, Édition des Universités de Montpellier II et Paris VII
% foot: Version 20 juin 2003
(et éventuellement, en même temps, $i \in W$ mais $i' \notin W$). De plus (contrairement aux constructions des propositions 10, 11 (p. 101-107)), les catégories envisagées pourront toutes être finies. Je suis particulièrement intéressé par le cas où $X$ est un ‘carré’
\[ X = \Phi_n \times \Phi_m \qquad (n,m \geq 2), \]
et où $U$ est son ‘bord’
\[ U = \dot{X} \overset{\text{déf}}{=} (\dot{\Phi}_n \times \Phi_m) \cup (\Phi_n \times \dot{\Phi}_m). \]

\msp{page 13}
\[
\begin{tikzcd}
0 \arrow[r] \arrow[d] & a \arrow[d] & a' \arrow[l] \arrow[r] \arrow[d] & \cdots \\
b \arrow[r] & t_1 & \cdots \arrow[l] & \\
b' \arrow[u] \arrow[r] \arrow[d] & {} \arrow[u] & \cdots \arrow[l] \arrow[u] & \\
\cdots & & &
\end{tikzcd}
\]

\textbf{Proposition.} \textit{Soient $n,m \geq 2$, et considérons l'inclusion $U = \partial X \overset{i}{\hookrightarrow} X$ du bord du ‘carré’ $X = \Phi_n \times \Phi_m$ dans $X$. C'est une immersion ouverte $W$-asphérique, de façon précise, les $U/t$ ($t \in \dot{U} = \bar{U}\backslash U$) sont isomorphes à $e$ (catégorie ponctuelle) ou à $\Phi =$ \begin{tikzcd}[ampersand replacement=\&, row sep=small, column sep=small] \& {} \arrow[dl] \\ {} \& \\ \& {} \arrow[ul] \end{tikzcd} (catégorie avec élément final), sauf pour les quatre éléments $t_i$ des quatre ‘coins’ de $\dot{U}$, où c'est la catégorie}
\[
\tilde{\Psi} = \left(
\begin{tikzcd}
& & a' \arrow[dl] \\
& a & \\
0 \arrow[ur] \arrow[dr] & & \\
& b & \\
& & b' \arrow[ul]
\end{tikzcd}
\right)
\]
\textit{n'ayant pas d'élément final. Donc $U \xrightarrow{\ i\ } X$ n'est pas immersion ouverte de Dwyer (mais $U \longrightarrow X\backslash\{t_1,t_2,t_3,t_4\}$ l'est \ldots). Ceci posé, je dis que $i : U \longrightarrow X$ n'est pas une $W$-cofibration.}


\endgroup
% END UNIT Der-ChXIIapp-19-06-03/Der-ChXIIapp-19-06-03_p006-010.tex
% BEGIN UNIT Der-ChXIIapp-19-06-03/Der-ChXIIapp-19-06-03_p011-014.tex
\begingroup\makeatletter%
\newcommand{\printedpage}[1]{\marginpar{\footnotesize p.~#1}}%
\newcommand{\msp}[1]{\marginpar{\footnotesize [#1]}}%
\newcommand{\dotto}{\mathrel{\tikz[baseline=-0.5ex]\draw[dotted,-latex] (0,0)--(1.6em,0);}}%
\makeatother\let\printedpage\derpp\let\msp\dermsp%
\let\mspin\DERmspin%


\printedpage{10}
% head: A. Grothendieck : Dérivateurs, Ch. XII, Édition des Universités de Montpellier II et Paris VII
% foot: Version 20 juin 2003
\textsc{Démonstration}. Soit
\[ F \hookrightarrow U \]
l'inclusion dans $U = \partial X$ de la catégorie-coin \texttt{[}$V =$\texttt{]} $\vcenter{\hbox{\begin{tikzcd}[row sep=small, column sep=small]
0 \arrow[r] \arrow[d] & a \\
b &
\end{tikzcd}}}$ (cf. figure), isomorphe à $\Psi$. On voit aussitôt que c'est une immersion fermée de Dwyer (le plus petit ouvert de $U$ contenant $\Psi$ est justement la catégorie $\tilde{\Psi} = U/t_1$ de tantôt), donc une $W$-cofibration. Si $U \longrightarrow X$ était une $W$-cofibration, il en serait de même du composé $V \longrightarrow X$. On va voir que ce n'est pas le cas.

Soit $\overline{V} = \vcenter{\hbox{\begin{tikzcd}[row sep=small, column sep=small]
0 \arrow[r] \arrow[d] & a \arrow[d] \\
b \arrow[r] & t_1
\end{tikzcd}}}$ l'adhérence de $V$ dans $X$,
\msp{page 14}et considérons l'inclusion
\[ \overline{V} \longrightarrow X, \]
je dis que \texttt{[c]}'est une immersion fermée de Dwyer, i.e. que pour tout $s \in W\backslash\overline{V}$, où $W$ est le plus petit ouvert de $X$ contenant $\overline{V}$, i.e.
\[
W = \quad
\begin{tikzcd}[row sep=small, column sep=small]
0 \arrow[r] \arrow[d] & a \arrow[d] & a' \arrow[l] \arrow[d] \\
b \arrow[r] & {} & a'' \arrow[l] \\
b' \arrow[u] \arrow[r] & b'' \arrow[u] & c, \arrow[l] \arrow[u]
\end{tikzcd}
\]
la catégorie $s\backslash\overline{V}$ a un objet initial. (Les catégories $s\backslash\overline{V}$ qu'on obtient, à partir des cinq choix possibles pour $s$, sont $\boldsymbol{\Delta}^1$ (si $s = a'$ ou $b'$) ou $e$ (pour $s = a''$, $b''$ ou $c$)). Donc $\overline{V} \longrightarrow X$ est une $W$-cofibration.

Ceci posé, on sait qu'il existe, pour [?] l'inclusion (immersion ouverte)
\[ V \hookrightarrow \overline{V} \qquad \text{i.e. } \Psi \hookrightarrow \boldsymbol{\Delta}^1 \times \boldsymbol{\Delta}^1, \]
un foncteur $V \longrightarrow V'$, donnant lieu à un carré cocartésien
\[
\begin{tikzcd}
V \arrow[r, hook, "i"] \arrow[d, "p"'] & \overline{V} \arrow[d, "\overline{p}"] \\
V' \arrow[r, hook, "i'"] & \overline{V'}
\end{tikzcd}
\tag*{$(*)$}
\]
avec
\[ i, p \in W \ \text{(car $V$, $V'$, $\overline{V}$ $W$-asphérique), mais } i', \overline{p} \notin W \ \text{(i.e. $\overline{V'}$ \textit{pas} $W$-asphérique).} \]

Je rappelle qu'on prend d'abord
\printedpage{11}
% head: A. Grothendieck : Dérivateurs, Ch. XII, Édition des Universités de Montpellier II et Paris VII
% foot: Version 20 juin 2003
\[
V'_0 = \left(
\begin{tikzcd}[row sep=small]
 & & a \\
0' \arrow[r, shift left, "\lambda"] \arrow[r, shift right, "\mu"'] & 0 \arrow[ur, "\alpha"] \arrow[dr, "\beta"'] & \\
 & & b
\end{tikzcd}
\right)
\qquad \text{avec} \quad
\left\{
\begin{aligned}
\alpha\lambda &= \alpha\mu \\
\beta\lambda &= \beta\mu
\end{aligned}
\right.
\]
et l'immersion fermée ‘identique’ de $V$ dans $V'$ \texttt{[plutôt} $V'_0$\texttt{]}
\[ p_0 : V \longrightarrow V'_0 \]
\msp{page 15}(qui est telle que
\[ p_{0!}(e_V) = e_{V'_0} \ \ ), \]
puis on prend pour $V'$ la catégorie obtenue en ajoutant un objet final à $V'_0$
\[
V' = \left(
\begin{tikzcd}[row sep=small, column sep=small]
 & & a \arrow[dr, "\rho"] & \\
0' \arrow[r, shift left, "\lambda"] \arrow[r, shift right, "\mu"'] & 0 \arrow[ur, "\alpha"] \arrow[dr, "\beta"'] & & e \\
 & & b \arrow[ur, "\sigma"'] &
\end{tikzcd}
\right)
\qquad \text{avec} \quad
\left\{
\begin{aligned}
\alpha\lambda &= \alpha\mu \\
\beta\lambda &= \beta\mu \\
\rho\alpha &= \sigma\beta
\end{aligned}
\right.
\]
et pour $p$ le foncteur composé
\[
\begin{tikzcd}
V \arrow[r, "p_0"] \arrow[rr, bend right, "p"'] & V'_0 \arrow[r, "p_1"] & V',
\end{tikzcd}
\]
où $p_1$ est l'inclusion (ouverte). On trouve alors le carré $(*)$, avec $\overline{V'}$ non $W$-asphérique, ou ce qui revient au même, $i' \notin W$. Car comme $\overline{V'}\backslash V' \simeq \overline{V}\backslash V$ est la catégorie ponctuelle, $i' \in W$ équivaut à $V'/t$ $W$-asphérique (\footnote[8]{moyennant (W 7 bis).}) (cf. cor. 4, p. 149 dans XII). Or par construction $V'/t \simeq V'_0$, qui n'est pas $W$-asphérique (dans le cas, disons, où $W \subseteq W_1$, car $\pi_1(V'_0) \simeq \mathbf{Z}$).

\msp{page 16}Considérons alors le diagramme de carrés cocartésiens
\[
\begin{tikzcd}
V \arrow[r, hook, "i"] \arrow[d, "p"'] & \overline{V} \arrow[r, hook] \arrow[d, "\overline{p}"] & X \arrow[d, "q"] \\
V' \arrow[r, hook, "i'"] & \overline{V'} \arrow[r, hook] & X'
\end{tikzcd}
\qquad .
\]

\printedpage{12}
% head: A. Grothendieck : Dérivateurs, Ch. XII, Édition des Universités de Montpellier II et Paris VII
% foot: Version 20 juin 2003
Comme $\overline{V} \hookrightarrow X$ est \texttt{[une]} immersion fermée de Dwyer, le carré de droite est $W$-cocartésien, donc $(\overline{p} \in W) \iff (q \in W)$, donc (comme $\overline{p} \notin W$) on a
\[ p \in W, \qquad \text{mais} \qquad q \notin W, \]
ce qui montre bien que $V \longrightarrow X$ n'est pas une $W$-cofibration. Factorisant $V \longrightarrow X$ par $U$ (au lieu de $\overline{V}$), on trouve
\[
\begin{tikzcd}
V \arrow[r] \arrow[d, "p\in W"'] & U \arrow[r] \arrow[d, "q_U\in W"] & X \arrow[d, "q\notin W"] \\
V' \arrow[r] & U' \arrow[r] & X'
\end{tikzcd}
\qquad ,
\]
où $q_U \in W$ puisque $p \in W$ et $V \longrightarrow U$ est \texttt{[une]} immersion de Dwyer. Cela rend explicite que $U \longrightarrow X$ n'est pas une $W$-cofibration, puisque
\[ q_U \in W, \qquad \text{mais} \qquad q \notin W, \]
\hfill q.e.d.

\msp{page 17}\textbf{Recollement dans} Cat \textbf{d'un ouvert et d'un fermé complémentaire.}

\texttt{[Supposons donné]}
\[ X_0 \underbrace{\hookrightarrow}_{\text{ouvert}} X \ \underbrace{\hookleftarrow}_{\substack{\text{fermé}\\ \text{complé-}\\ \text{mentaire}}} X_1 \]
dans Cat, \texttt{[et soit]}
\[ \tilde{X}_{01} \subseteq \mathrm{Fl}(X) \]
\texttt{[la]} sous-catégorie pleine des flèches de source dans $X_0$, but dans $X_1$.

\textbf{Question.} \textit{Ce diagramme}
\[
\begin{tikzcd}[column sep=small]
 & & \tilde{X}_{01} \arrow[dll, "p"' near start, "\text{fibrant, fibres les $x_0\backslash X_1$}"' near end] \arrow[drr, "q" near start, "\text{cofibrant, fibres les $X_0/x_1$}" near end] & & \\
X_0 \arrow[ddrr, "i"'] & & & & X_1 \arrow[ddll, "j"] \\
 & {} \arrow[rr, bend left] & & {} & \\
 & & X, & &
\end{tikzcd}
\tag*{$(*)$}
\]
\printedpage{13}
% head: A. Grothendieck : Dérivateurs, Ch. XII, Édition des Universités de Montpellier II et Paris VII
% foot: Version 20 juin 2003
\textit{est-il $W$-cocartésien?}

OK si $i$ $W$-asphérique ou j $W$-coasphérique, i.e. $X \longrightarrow \boldsymbol{\Delta}^1$ lisse ou propre -- c'est un cas où on a soit $X_0 \xrightarrow{W} X$, soit $X_1 \xrightarrow{W} X$. Mais même en supposant $i$ ou $j$ dans $W$, et de plus $X$ asphérique (donc $X_0$ et $X$ asphériques, ou $X_1$ et $X$ asphériques), et $W = \mathbf{W}_\infty$, je n'ai pu prouver que le diagramme est cocartésien, i.e. $q \in W$ resp. $p \in W$ (\footnote[9]{Je suppose que $W$ satisfait Loc (6).}). (Pas non plus si $X_0 \simeq e$, $X$, $X_1$ $W$-asphérique $\Longrightarrow$ $e\backslash X$ $W$-asphérique, i.e. $X_1 \longrightarrow X$ $W$-coasphérique $\ldots$\,) Mais j'ai prouvé du moins, dans tous les cas, que % [sic] j (upright) -> $j$, letter denoting a functor printed in roman type
\[
\begin{tikzcd}[row sep=small]
 & X_0 & & \\
\mathcal{X} \overset{\text{déf}}{=} \displaystyle\int \tilde{X}_{01} \arrow[ur] \arrow[dr] & & {} \arrow[r, "\varphi"] & X \\
 & X_1 & &
\end{tikzcd}
\tag*{$(**)$} % verified: CMSY10 glyph, \mathcal{X}; corrections.csv, PDF leaf 14
\]
est dans $\mathbf{W}_1$.

Chose \textit{fausse} (sauf si $i$ ou $j$ dans $W = \mathbf{W}_\infty$) si on remplace $\tilde{X}_{01}$ par $\tilde{X}'_{01}$, formée des triples $(x_0, x_1, c)$, avec $x_0 \in \mathrm{Ob}\,X_0$, $x_1 \in \mathrm{Ob}\,X_1$, $c : x_0 \dotto x_1$ un \textit{chemin} de $x_0$ à $x_1$ dans $X$! On a
\msp{page 18}donc une situation ‘topologique’ tout à fait spéciale à la catégorie de modèles Cat, qui ne semble pas avoir d'analogue dans des catégories telles que Esp ou $\boldsymbol{\Delta}^\wedge$, où pourtant les notions d'ouvert et de fermé complémentaire (avec grain de sel dans $\boldsymbol{\Delta}^\wedge$) ont un sens.

L'argument catégorique quasi-trivial qui établit que $(**)$ est dans $\mathbf{W}_1$, i.e. induit une équivalence de catégories entre $\mathrm{Rev}(X)$ et
\[ \mathrm{Rev}(\mathcal{X}) \ \simeq \ \mathrm{Rev}(X_0) \overset{2}{\times}_{\mathrm{Rev}(\tilde{X}_{01})} \mathrm{Rev}(X_1), \]
me semble devoir s'étendre aux dimensions supérieures, à coups de $n$-gerbes sur $X$ -- mais les fondements sur la théorie des $n$-gerbes manquent. (Et je n'ai pas eu le courage d'expliciter l'argument même pour $n = 2$ (où pourtant \textsc{Giraud} fournit ce qu'il faut) (\footnote[10]{Mais j'ai fait quand même une esquisse de vérification pour $n = 2$, qui me convainc qu'on a bien $\varphi \in \mathbf{W}_2$. Et je n'ai pas guère de doute qu'on a même $\varphi \in \mathbf{W}_\infty$.}).) Aussi je conjecture que $\varphi$ est tout au moins dans $\mathbf{W}_\infty$, i.e. que le diagramme $(*)$ est $W$-cocartésien si $W \supseteq \mathbf{W}_\infty$ (et satisfait Loc (6 a$'$)). % [sic] je n'ai pas guère (footnote 10) -> je n'ai guère, double negation pas + guère


\endgroup
% END UNIT Der-ChXIIapp-19-06-03/Der-ChXIIapp-19-06-03_p011-014.tex
\clearpage\setcounter{footnote}{0}\phantomsection\addcontentsline{toc}{section}{Chapitre XIII}\markboth{Chapitre XIII}{Chapitre XIII}
% BEGIN UNIT Der-ChXIII-25-01-04/Der-ChXIII-25-01-04_p001-005.tex
\begingroup\makeatletter%
\newcommand{\printedpage}[1]{\marginpar{\footnotesize p.~#1}}%
\newcommand{\msp}[1]{\marginpar{\footnotesize [#1]}}%
\newcommand{\Fl}{\mathrm{Fl}}%
\newcommand{\M}{\mathcal{M}}%
\makeatother\let\printedpage\derpp\let\msp\dermsp%
\let\mspin\DERmspin%


\printedpage{--} % no page number printed on this page
% head: (none)

\begin{center}
\vspace*{5cm}
{\LARGE\sffamily LES DÉRIVATEURS}

\vspace{1.5cm}
{\large\textbf{ALEXANDRE GROTHENDIECK}}

\vspace{2.5cm}
{\large\textbf{Chapitre XIII}}

\vspace{0.5cm}
{\large\textbf{Catégories de modèles (1)}}
\end{center}

\par

\printedpage{--} % no page number printed on this page
% head: (none)

\vspace*{6cm}
\noindent \par\noindent\textbf{Note des éditeurs de l’édition transcrite.}\par
Ce texte a été déchiffré et transcrit en \LaTeXe{} par M. Künzer. Il a été édité par M. Künzer, J. Malgoire, et G. Maltsiniotis. La transcription est aussi fidèle que possible au manuscrit. Pour les quelques corrections évidentes, ou rares commentaires des éditeurs, ainsi que pour la numérotation originale des pages du manuscrit, les caractères de machine à écrire \texttt{[typewriter]} entre crochets sont utilisés. Un point d'interrogation entre crochets signifie que l'on n'est pas sûr du mot qui précède. Les index, la liste des axiomes, le glossaire, et les références bibliographiques à la fin du chapitre ont été compilés par les éditeurs. Les numéros des pages dans les index sont ceux de l'original.

\noindent Cette édition est provisoire. Les remarques, commentaires et corrections sont bienvenus. Envoyer un message à :

\noindent\textbf{[e-mail address removed]}

\hfill G. Maltsiniotis\hspace*{2cm}

\par

\printedpage{1}
% head: A. Grothendieck : Dérivateurs, Ch. XIII, Édition des Universités de Montpellier II et Paris VII
% foot: Version 25 janvier 2004

\msp{page 1}

\section*{1\quad Couples de Quillen $(\Phi, \Psi)$}

Je veux d'abord donner une variante affaiblie de la notion de triple de Quillen (``model category'' de Quillen). Ce ne sera pas une notion tout-à-fait autoduale -- ainsi il y aura \textit{deux} variantes affaiblies, duales l'une de l'autre.

Mais avant encore, un sorite sur les LLP (left lifting property) et RLP (right lifting property) (\footnote[1]{donner définition en forme, avec renvoi à \textsc{Quillen}.}). Soit $\M$ une catégorie quelconque (le plus souvent, $\M$ sera grande). On considère des parties
\[
\Phi,\ \Psi\ \subset\ \Fl(\M)
\]
et on définit des applications
\[
\boxed{\begin{array}{ccc}
\Phi & \longmapsto & \Phi_* \\[1ex]
\Psi & \longmapsto & \Psi^*
\end{array}}
\qquad
\Fl(\M) \longrightarrow \Fl(\M)
\]
\texttt{[plutôt} $\mathfrak{P}(\Fl(\M) \longrightarrow \mathfrak{P}(\Fl(\M))$\texttt{,}\newline\texttt{où pour tout ensemble} $E$\texttt{,} $\mathfrak{P}(E)$ \texttt{désigne l'ensemble des parties de} $E$\,\texttt{]} par % [sic] \mathfrak{P}(\Fl(\M) -> \mathfrak{P}(\Fl(\M)), a closing parenthesis is missing in the first term
\[
\begin{aligned}
\Phi_* \ &= \ \{\text{ensemble des } p \in \Fl(\M) \text{ qui ont la RLP par rapport à } \Phi\} \\
\Psi^* \ &= \ \{\text{ensemble des } i \in \Fl(\M) \text{ qui ont la LLP par rapport à } \Psi\},
\end{aligned}
\]
et ensuite
\[
\boxed{\begin{array}{ccc}
\Phi & \longmapsto & \widetilde{\Phi} \\[1ex]
\Psi & \longmapsto & \overline{\Psi}
\end{array}}
\qquad
\Fl(\M) \longrightarrow \Fl(\M)
\]
\texttt{[plutôt} $\mathfrak{P}(\Fl(\M) \longrightarrow \mathfrak{P}(\Fl(\M))$\texttt{]} par % [sic] \mathfrak{P}(\Fl(\M) -> \mathfrak{P}(\Fl(\M)), a closing parenthesis is missing in the first term
\[
\begin{aligned}
\widetilde{\Phi} \ &= \ (\Phi_*)^* \\
\overline{\Psi} \ &= \ (\Psi^*)_* \ .
\end{aligned}
\]

\msp{page 2}

Je résume les propriétés les plus immédiates de ces opérations dans une

\noindent\textbf{Proposition 1.}

\begin{enumerate}[label={}]
\item[a)] $\Phi \longmapsto \Phi_*$, $\Psi \longmapsto \Psi^*$ \textit{sont des applications décroissantes de} $\Fl(\M)$ \texttt{[plutôt} $\mathfrak{P}(\Fl(\M))$\texttt{]} \textit{dans lui-même, pour la relation d'ordre} $\subset$\textit{, et} $\Phi \longmapsto \widetilde{\Phi}$\textit{,} $\Psi \longmapsto \overline{\Psi}$ \textit{des applications croissantes.}

\printedpage{2}
% head: A. Grothendieck : Dérivateurs, Ch. XIII, Édition des Universités de Montpellier II et Paris VII
% foot: Version 25 janvier 2004

\item[b)] \textit{On a}
\[
\Phi\ \subset\ \widetilde{\Phi},
\]
\textit{et pour qu'on ait $\Phi = \widetilde{\Phi}$, il faut et il suffit que $\Phi$ soit de la forme $\Psi^*$. On dit alors que $\Phi$ est} \textup{clos à gauche} (\footnote[2]{Donner définition en forme.})\textit{. L'ensemble des parties de} $\Fl(\M)$ \textit{closes à gauche est stable par intersections quelconques} (\footnote[3]{de façon précise, $\bigcap_i(\Psi_i^*) = (\bigcup_i \Psi_i)^*$}), \textit{et $\widetilde{\Phi}$ est la plus petite partie close de} $\Fl(\M)$ \textit{qui contienne $\Phi$ et la plus grande partie $\Phi'$ de} $\Fl(\M)$ \textit{telle que $\Phi'_* = \Phi_*$ (ou telle que $\Phi'_* \supset \Phi_*$).}

\textit{Énoncés symétriques pour $\Psi$ : on a}
\[
\Psi\ \subset\ \overline{\Psi},
\]
\textit{on a $\Psi = \overline{\Psi}$ si et seulement si $\Psi$} \texttt{[est]} \textit{de la forme $\Phi_*$, on dit alors que $\Psi$ est} \textup{clos à droite.} \textit{L'ensemble des parties de} $\Fl(\M)$ \textit{closes à droite est stable par intersections quelconques} (\footnote[4]{de façon précise, on a $\bigcap_i(\Phi_{i\,*}) = (\bigcup_i \Phi_i)_*$}), \textit{et $\overline{\Psi}$ est la plus petite partie close à droite contenant $\Psi$, et aussi la plus grande partie $\Psi'$ de} $\Fl(\M)$ \textit{telle que $\Psi'^* = \Psi^*$ (ou telle que $\Psi'^* \supset \Psi^*$).}

\item[] \hspace*{-2.5em}\msp{page 3}

\item[c)] \textit{Les applications}
\[
\Phi \ \longmapsto\ \Phi_*, \qquad\quad \Psi\ \longmapsto\ \Psi^*
\]
\textit{établissent des bijections inverses l'une de l'autre entre l'ensemble des parties $\Phi$ closes à gauche de} $\Fl(\M)$\textit{, et l'ensemble des parties $\Psi$ closes à droite. C'est un anti-isomorphisme d'ensembles ordonnés. On a alors}
\[
(\Phi \cap \Phi')_* \ = \quad
\begin{array}{l}
\textit{partie close à droite de } \Fl(\M) \ \textit{``engendrée par } \Phi_*,\ \Phi'_*\textit{'',} \\
\text{i.e. \texttt{[la]} \textit{plus petite partie close à droite contenant} } \Phi_*, \\
\Phi'_*, \text{ i.e. } \overline{\Phi_* \cup \Phi'_*}
\end{array}
\]
\texttt{[pour} $\Phi$\texttt{,} $\Phi'$ \texttt{parties closes à gauche]}\textit{, plus généralement}
\[
(\bigcap_i \Phi_i)_* \ = \ \overline{\bigcup_i (\Phi_{i\,*})}
\]
\textit{pour toute famille $(\Phi_i)_{i \in I}$ de parties closes à gauche de} $\Fl(\M)$\textit{. Formules duales} \texttt{[pour des parties closes à droite]}
\[
\left\{
\begin{aligned}
(\Psi \cap \Psi')^* \ &= \ \widetilde{\Psi^* \cup \Psi'^*} \\
(\bigcap_j \Psi_j)^* \ &= \ \widetilde{\bigcup_j (\Psi_j^*)}\ .
\end{aligned}
\right.
\]

\item[d)] \textit{Si} $\Phi, \Psi \subset \Fl(\M)$\textit{, alors on a}
\[
\Phi \subset \Psi^* \iff \Psi \subset \Phi_*, \qquad\quad \textit{on écrit alors } \Phi \longleftrightarrow \Psi,
\]
\textit{cela signifie que $\Phi$ a la LLP pour $\Psi$, ou que $\Psi$ a la RLP pour $\Phi$.}
\end{enumerate}

\printedpage{3}
% head: A. Grothendieck : Dérivateurs, Ch. XIII, Édition des Universités de Montpellier II et Paris VII
% foot: Version 25 janvier 2004

\msp{page 4}

\noindent\textbf{Proposition 2.} \textit{Soit $\Phi$ une partie close à gauche de} $\Fl(\M)$\textit{. Alors $\Phi$ satisfait les conditions} a) b) c) \textit{ci-dessous :}

\begin{enumerate}[label={}]
\item[a)] \textit{$\Phi$ contient les isomorphismes, et est stable par composition, et par cochangement de base (}i.e. \textit{si on a}
$\vcenter{\hbox{\begin{tikzcd}
A \arrow[r, "i"] \arrow[d] & B \arrow[d] \\
A' \arrow[r, "i'"] & B'
\end{tikzcd}}}$
\textit{carré cocartésien dans $\M$, avec $i \in \Phi$, on a $i' \in \Phi$).}

\item[b)] \textit{$\Phi$ stable par facteurs directs (en tant que partie de l'ensemble des objets de} $\underline{\Fl}(\M)$\textit{).}

\item[c)] \textit{$\Phi$ ``stable par limites inductives ordinales'', dans le sens suivant : Soit $I$ un ensemble bien ordonné, $(A_i)_{i \in I}$ un système inductif dans $\M$, tel que les $u_{ji} : A_i \longrightarrow A_j$ soient dans $\Phi$, et $A_{i_0} \xleftarrow{\ \sim\ } \varinjlim\limits_{i<i_0} A_i$ si $i_0 \in I$ est un élément limite. Supposons que $A = \varinjlim\limits_{I} A_i$ existe. Alors $\forall\ i_0 \in I$, la flèche canonique $A_{i_0} \longrightarrow A$ est dans $\Phi$.}
\end{enumerate}

Il y a un énoncé dual pour une partie $\Psi$ close à droite; elle contient les isomorphismes, est stable par composition et par changement de base, par facteurs directs, et enfin stable par limites projectives ordinales.

\msp{page 5}

On aura une réciproque partielle dans les catégories accessibles. C'est pour l'énoncé de cette réciproque que j'ai mis à part la condition \textit{b}) de stabilité par facteurs directs, au lieu de la bloquer avec \textit{a}).

\noindent\textbf{Définition 1.}

\begin{enumerate}[label={}]
\item[\textit{a})] (Pour mémoire.) RLP, LLP, la relation $\Phi \longleftrightarrow \Psi$ entre parties de $\Fl(\M)$ (relation non symétrique!).

\item[\textit{b})] (Pour mémoire.) Partie de $\Fl(\M)$ close à gauche, close à droite.

\item[\textit{c})] Soit $\Phi \longleftrightarrow \Psi$. On dit que le couple $(\Phi, \Psi)$ a la \textit{propriété de factorisation de Quillen}, si \textit{toute} flèche $f$ de $\M$ se factorise en
\[
f = pi, \qquad\quad i \in \Phi,\ p \in \Psi.
\]
On dit que $(\Phi, \Psi)$ est un \textit{couple de Quillen} si c'est un couple associé par \texttt{[les]} propriétés LLP et RLP ($\Phi \longleftrightarrow \Psi$), et \texttt{[si]} elles \texttt{[les parties} $\Phi$ \texttt{et} $\Psi$\texttt{]} sont closes (\textit{i.e.} $\Phi = \Psi^*$, $\Psi = \Phi_*$), et si elles ont la propriété de factorisation de Quillen.
\end{enumerate}

J'énonce dès maintenant le résultat technique essentiel, assurant l'existence de beaucoup de couples de Quillen :


\endgroup
% END UNIT Der-ChXIII-25-01-04/Der-ChXIII-25-01-04_p001-005.tex
% BEGIN UNIT Der-ChXIII-25-01-04/Der-ChXIII-25-01-04_p006-010.tex
\begingroup\makeatletter%
\newcommand{\printedpage}[1]{\marginpar{\footnotesize p.~#1}}%
\newcommand{\msp}[1]{\marginpar{\footnotesize [#1]}}%
\newcommand{\cM}{\mathcal{M}}%
\newcommand{\Fl}{\mathrm{Fl}}%
\newcommand{\Cof}{\mathrm{Cof}}%
\newcommand{\Fib}{\mathrm{Fib}}%
\newcommand{\Ob}{\mathrm{Ob}}%
\makeatother\let\printedpage\derpp\let\msp\dermsp%
\let\mspin\DERmspin%

\setcounter{footnote}{4}

\printedpage{4}
% head: A. Grothendieck : Dérivateurs, Ch. XIII, Édition des Universités de Montpellier II et Paris VII
% foot: Version 25 janvier 2004
\msp{page 6}\textbf{Théorème 1.} \textit{Supposons que $\cM$ soit} accessible, \textit{et stable par (petites) limites inductives. (Il doit suffire stabilité par sommes amalgamées, et par $\varinjlim$ filtrantes, voire ordinales \dots).} \textit{Soit $\Phi_0$ une} petite \textit{partie de} $\Fl(\cM)$\textit{, $\Psi=(\Phi_0)_*$, et $\Phi$ une partie de} $\Fl(\cM)$ \textit{telle que}
\[
\Phi_0\ \subset\ \Phi\ \subset\Psi^*\ (=\widetilde{\Phi}_0).
\]
\textit{Alors :}

\begin{itemize}[leftmargin=2.5em]
\item[a)] \textit{Pour que $\Phi$ soit close à gauche,} i.e. $\Phi=\widetilde{\Phi}_0$\textit{, il faut et il suffit qu'elle satisfasse les conditions} a) b) c) \textit{de la proposition 2. Quand il en est ainsi, le couple $(\Phi,\Psi)$ est un couple de Quillen.}

\item[b)] \textit{Supposons seulement que $\Phi$ satisfasse aux conditions} a), c) \textit{de la proposition 2 (à l'exclusion de la condition} b) \textit{des facteurs directs) – donc on ne peut affirmer que $\Phi$ soit close à gauche. Néanmoins, $(\Phi,\underbrace{\Psi}_{=\,\Phi_*})$ satisfait à la propriété de factorisation de Quillen.}
\end{itemize}

\textbf{Remarque.} Dans la démonstration, l'existence de $\Phi_0\subset\Phi$, qui joue le rôle d'un petit ensemble de “générateurs” (à gauche) pour $\Phi$, joue un rôle essentiel. Il y aurait lieu de développer la notion de \textit{parties accessibles}
\msp{page 7}
de $\Fl(\cM)$, qui impliquerait l'existence d'un petit ensemble $\Phi_0\subset\Phi$ tel que $\Phi\subset\widetilde{\Phi}_0$, et qui aurait la propriété que
\[
\begin{aligned}
&\Phi\ \text{accessible}\quad\Longrightarrow\quad\Phi_*\ \text{accessible}\\
&\Psi\ \text{accessible}\quad\Longrightarrow\quad\Psi^*\ \text{accessible.}
\end{aligned}
\]
Tout couple $(\Phi,\Psi)$ de parties closes associées ($\Phi=\Psi^*$, $\Psi=\Phi_*$) et accessible, serait donc un couple de Quillen.

\section*{2\quad Catégories de Quillen faibles, et construction d'icelles (Théorème-scholie conjectural)}

Dans toute la suite, on se donne une \textit{catégorie de modèles}
\[
\left\{
\begin{array}{ll}
\mbox{$(\cM,W)$, où $W\subset\Fl(\cM)$} & \\[3ex]
\mbox{\textit{$W$ faiblement saturée} (\footnotemark) :} &
\begin{array}[c]{@{}l@{\quad}p{6.5cm}@{}}
a) & $W$ contient les isomorphismes.\\[1ex]
b) & Si deux parmi $f$, $g$, $gf$ \texttt{[sont]} dans $W$, le troisième aussi.
\end{array}
\end{array}
\right.
\]
\footnotetext{C'est l'axiome M 5 de \textsc{Quillen}.}

\printedpage{5}
% head: A. Grothendieck : Dérivateurs, Ch. XIII, Édition des Universités de Montpellier II et Paris VII
% foot: Version 25 janvier 2004
On veut décrire des données supplémentaires un peu plus faibles que celles d'une catégorie de modèles de Quillen. On suppose données des parties
\[
C,\ F,\ TC,\ TF\ \subset\ \Fl(\cM),
\]
avec les propriétés suivantes. Comme axiome préliminaire, je pose

\begin{itemize}[leftmargin=3em,labelsep=0.6em]
\item[\textbf{M 0}] Stabilité de $\cM$ par $\varprojlim$ finies et $\varprojlim$ \texttt{[plutôt $\varinjlim$]} finies.

\item[\textbf{M 1}]\msp{page 8}Propriété de relèvement-prolongement
\[
C\longleftrightarrow TF,\qquad TC\longleftrightarrow F.
\]

\item[\textbf{M 2}] Propriété de factorisation pour $(C,TF)$, $(TC,F)$.

\item[\textbf{M 3}] Stabilité de $F$ par composition et par changement de base, de $C$ par composition et cochangement de base, $F$ et $C$ contiennent les isomorphismes.

\item[\textbf{M 4}] Stabilité de $TF$ par composition et changement de base, de $TC$ par composition et cochangement de base, et $TF$ et $TC$ contiennent les isomorphismes.

\item[\textbf{M}$'$] $TC\subset C\cap W$, $TF\subset F\cap W$.

\item[\textbf{M}$''$] $C\subset\ \underbrace{\Cof_W}_{\text{$W$-cofibrations}}\ ,\ F\subset\ \underbrace{\Fib_W}_{\text{$W$-fibrations}}\ .$
\end{itemize}

\texttt{[Un morphisme $p:X\longrightarrow Y$ de $\cM$ est une $W$-fibration si pour tout diagramme de carrés cartésiens de la forme :}
\[
\begin{tikzcd}
X \ar[d,"p"'] & X' \ar[l] \ar[d] & X'' \ar[l,"t"'] \ar[d] \\
Y & Y' \ar[l] & Y'' \ar[l,"s"] \ ,
\end{tikzcd}
\]
\texttt{si $s\in W$, alors $t\in W$. La notion de $W$-cofibration est la notion duale.]}

Dans le cas de \textsc{Quillen} (Q-\textit{catégories de modèles}), M$'$ est renforcé en
\[
TC=C\cap W,\qquad TF=F\cap W,
\]
et il n'y a pas d'axiome M$''$. L'axiome M$''$ équivaut à la propriété de “propriety”, correspondant aux Q-catégories de modèles \textit{propres.}

\textbf{Proposition 3.} \textit{Supposons les conditions} M 0 \textit{(pour mémoire),} M 1, M 2, M$'$, M$''$ \textit{satisfaites. Alors en remplaçant $C$, $TC$ par les saturés pour les propriétés de stabilité de contenir les isomorphismes, d'être stable par composition et par cochangement de base, et $F$, $TF$ par celles de contenir les isomorphismes et d'être stable par composition et par changement de base, on trouve $C'$, $TC'$, $F'$, $TF'$ satisfaisant}
\printedpage{6}
% head: A. Grothendieck : Dérivateurs, Ch. XIII, Édition des Universités de Montpellier II et Paris VII
% foot: Version 25 janvier 2004
\msp{page 9}\textit{à tous les axiomes} M 1 à M 4, M$'$, $M''$. % [sic] $M''$ (italic M) -> M$''$ (roman M), the axiom is printed with roman M everywhere else

\textsc{Démonstration}. La proposition 2 nous assure que M 1 continue à être satisfait, et pour M 2 c'est trivial. M 3 et M 4 seront valables par construction. D'autre part, M$''$ reste valable, car $\Cof_W$ a les stabilités requises pour $C'$, et $\Fib_W$ celles requises pour $F'$. Reste à prouver M$'$, \textit{i.e.}
\[
TC'\ \subset\ C'\cap W,\qquad TF'\subset F'\cap W,
\]
et pour ceci il suffit de voir que $C'\cap W$ resp. $F'\cap W$ a les stabilités requises pour $TC'$ resp. $TF'$. Cela résulte du

\textbf{Lemme.} \textit{Soit $C$ une partie de} $\Cof_W$\textit{. Si $C$ contient les isomorphismes, il en est de même de $C\cap W$. Si de plus $C$ est stable par composition resp. par cochangement de base, il en est de même} \texttt{[de]} \textit{$C\cap W$. Énoncé dual pour une partie $F$ de} $\Fib_W$\textit{.}

Cela provient du fait que
\[
\boxed{\ \Cof_W\cap W=W^{\mathrm{univ}}\ (\mbox{\footnotemark})\ }
\stackrel{\text{déf}}{=}
\begin{array}{l}
\text{ensemble des flèches $A\longrightarrow B$ telles que pour tout}\\
A\longrightarrow A',\ A'\longrightarrow B'=A'\vee_A B\ \ \text{est dans $W$.}
\end{array}
\]
\footnotetext{flèches \textit{coüniversellement} dans $W$, ou $W$-\textit{équivalences coüniverselles}.}
est lui-même stable par composition et par cochangement de base, et que l'on a bien sûr, si $C\subset\Cof_W$,
\[
C\cap W\ =\ C\cap(\Cof_W\cap W)\ =\ C\cap W^{\mathrm{univ}}.
\]

\msp{page 10}Donc il reste à prouver le

\textbf{Corollaire 1.} \textit{Supposons qu'il existe des parties $C$, $F$, $TC$, $TF$ $\subset$} $\Fl(\cM)$ \textit{satisfaisant} M 2, M$'$ (\footnote{On n'a pas besoin de tout M$'$, seulement de $TC\subset W$, $TF\subset W$.}), M$''$. \textit{Alors}
\[
\begin{aligned}
\Cof_W\cap W\ &=\ W^{\mathrm{univ}}\\
\Fib_W\cap W\ &=\ W_{\mathrm{univ}}\ \ (\mbox{\footnotemark}).
\end{aligned}
\]
\footnotetext{Flèches $f:X\longrightarrow Y$ \textit{universellement} dans $W$, \textit{i.e.} telles que pour tout $Y'\longrightarrow Y$, $f':X'=X\times_Y Y'\longrightarrow Y'$ soit dans $W$.}

Pour voir ceci, montrons p. ex. le deuxième (le premier s'en déduira par dualité).
\[
\begin{tikzcd}
X \ar[d,"f"'] & X' \ar[l] \ar[d,"f'"] \\
Y & Y' \ar[l]
\end{tikzcd}
\qquad\qquad\qquad
f\in\Fib_W,\ f\in W\Longrightarrow f'\in W.
\]

\printedpage{7}
% head: A. Grothendieck : Dérivateurs, Ch. XIII, Édition des Universités de Montpellier II et Paris VII
% foot: Version 25 janvier 2004
On a mieux :

\textbf{Proposition 4.} \textit{Supposons que toute flèche de $\cM$ se factorise en $pi$, avec $i\in W$ et $p\in$} $\Fib_W$\textit{. Alors, si $X$, $Y$ sont dans} $\underline{\Fib}_W S\subset\cM/S$ \texttt{[sous-catégorie pleine de $\cM/S$ formée des $r:Z\longrightarrow S$, où $r\in\Fib_W$]}\textit{, alors toute $f:X\longrightarrow Y$ dans $\cM/S$ qui est dans} $W_S$ \texttt{[formé des flèches au-dessus de $S$ qui sont dans $W$]} \textit{est dans $W^{\mathrm{u}}_S$,} i.e. \textit{elle reste dans $W$ après tout changement de base $S'\longrightarrow S$.}
\[
\begin{tikzcd}[column sep=small]
X \ar[rr,"f"] \ar[dr,"p"'] & & Y \ar[dl,"q"] \\
& S &
\end{tikzcd}
\qquad\qquad\qquad
p,q\in\Fib_W
\]

\textit{Énoncé dual pour les flèches de} $S\backslash\underline{\Cof}_W\subset S\backslash\cM$ \texttt{[sous-catégorie pleine de $S\backslash\cM$ formée des $i:S\longrightarrow Z$, où $i\in\Cof_W$]}\textit{, une telle flèche qui est dans $W$ est $S$-coüniversellement dans $W$, si toute flèche de $\cM$ se factorise en $qj$, avec $j\in$} $\Cof_W$ \textit{et $q\in W$.}

\msp{page 11}\textsc{Démonstration}.
\[
\begin{tikzcd}[column sep=small]
X \ar[dr,"f\in W"] \ar[dd] & & \overline{X'} \ar[ll] \ar[dr,"\overline{f'}\in W"] \ar[dd] & & X' \ar[ll,"i_X"'] \ar[dr,"f'"] \ar[dd] & \\
& Y \ar[dl] & & \overline{Y'} \ar[ll,crossing over] \ar[dl] & & Y' \ar[ll,"i_Y"',crossing over] \ar[dl] \\
S & & \overline{S'} \ar[ll,"F"] & & S' \ar[ll,"i"',"TC"] &
\end{tikzcd}
\]
On factorise $S'\longrightarrow S$ en $S'\overset{i}{\longrightarrow}\overline{S'}\overset{p}{\longrightarrow}S$, avec $i\in W$, $p\in\Fib_W$, par M 2. Donc $f\in W\Longrightarrow\overline{f'}\in W$. Comme $X,Y\in\Ob\,\underline{\Fib}_W S$, on a $\overline{X'},\overline{Y'}\in\Ob\,\underline{\Fib}_W\overline{S'}$, et comme $i:S'\longrightarrow\overline{S'}$ est dans $W$, donc aussi $i_X$, $i_Y$, on a $f'\in W$ par M 5 (saturation faible de $W$).

\textbf{Commentaire sur Proposition 3.} On peut être encore plus exigeant, en demandant dans M 3, M 4 également la stabilité par facteurs directs, et (quand seuls M 1, M2, M$'$, M$''$ % [sic] M2 -> M 2, printed without the space used everywhere else
sont satisfaits) en saturant $C$, $TC$, $F$, $TF$ également pour cette stabilité supplémentaire. La proposition 2 assure encore que M 1 continue à être vérifié (et M 2, M 3, M 4 trivialement), on voit qu'il en est ainsi encore de M$'$, M$''$ à condition qu'on renforce l'axiome de saturation faible M 5 sur $W$ par
\[
\textit{stabilité de $W$ par facteurs directs,}
\]
ce qui implique en effet que $\Cof_W$, $\Fib_W$ sont stables par facteurs directs, ainsi que $W^{\mathrm{univ}}$ et $W_{\mathrm{univ}}$.

\msp{page 12}Il ne semble pas que les axiomes M 0 à M 4, plus M$'$ et M$''$ suffisent à développer les suites exactes de suspension et de cosuspension. Il faut un petit chouia de plus :

\printedpage{8}
% head: A. Grothendieck : Dérivateurs, Ch. XIII, Édition des Universités de Montpellier II et Paris VII
% foot: Version 25 janvier 2004
\textbf{Définition 2.} J'appelle \textit{pré-Q-catégorie de modèles à gauche,} la donnée de $\cM$ et de $W$, $C$, $F$, $TC$, $TF$ $\subset$ $\Fl(\cM)$, satisfaisant les axiomes M 0 à M 5, M$'$, M$''$, plus la forme renforcée de la première inclusion dans M$'$
\[
\text{M}'_0\ \text{(g)}\qquad\qquad\qquad TC=C\cap W\qquad\quad(\text{non seulement } TC\subset C\cap W).
\]
De même, on définit la \textit{pré-Q-catégorie de modèles à droite,} par la condition
\[
\text{M}'_0\ \text{(d)}\qquad\qquad\qquad\qquad\qquad\qquad TF=F\cap W.
\]

NB Quand les deux conditions sont satisfaites à la fois, alors la donnée de $TC$, $TF$ résulte déjà de celle de $W$, $C$, $F$, et on a simplement une Q-catégorie de modèles “propre” (au sens de \textsc{Quillen}).

\textbf{Scholie.} Pour moi, l'introduction de $C$, $TC$, $F$, $TF$ est purement auxiliaire – ce qui m'importe, ce sont les propriétés de $W$. Ainsi, l'existence de $C$, $TC$, $F$, $TF$ satisfaisant aux axiomes faibles M 0 à M 4, M$'$, M$''$, du fait qu'elle donne les factorisations
\msp{page 13}
dont il est question dans \texttt{[la]} proposition 3, entraîne la propriété remarquable de $W$ énoncée dans cette proposition, et dans le corollaire qui la précède. De même, ces propriétés de factorisation (\footnote{sous une forme légèrement renforcée il est vrai – il faut une factorisation \textit{fonctorielle} pour $f$ \dots}) assurent l'existence d'une bonne théorie des fibres et cofibres $W$-homotopiques sans besoin de référence dans cette théorie à $C$, $TC$, $F$, $TF$ (dont l'introduction donne tout au plus des petits compléments \dots). Mais je ne vois pas comment, à partir des seules hypothèses de factorisation, prouver la propriété cruciale d'exactitude pour les carrés $W$-cartésiens, et $W$-cocartésiens, nécessaire pour pouvoir développer la suite \textit{exacte} de cosuspension et de suspension. C'est pour avoir cette exactitude que l'introduction auxiliaire, et transitoire, de l'attirail $(C,TC,F,TF)$ est utile, l'axiome supplémentaire à gauche $TC=C\cap W$ assurant (sauf erreur) l'exactitude des carrés $W$-cartésiens, et l'axiome à droite $TF=F\cap W$ celle des carrés $W$-cocartésiens. (Chose à vérifier avec plus de soin que je ne l'ai fait à présent.) Une fois ceci obtenu, et éventuellement itou pour les $(\cM(I),W(I))$ pour toute catégorie-diagramme $I$, on peut
\msp{page 14}
totalement oublier ces ensembles auxiliaires – j'ai envie en tous cas de les oublier.

Mon espoir, c'est que pour une $W$ donnée, même si on n'arrive pas à trouver une structure de triple de Quillen $(W,C,F)$ associée (p. ex. parce qu'il n'y en aurait pas), on peut trouver à la fois des quintuplets $(W,C,F,TC,TF)$ qui soient des pré-Quillen à gauche, et des quintuplets (d'\textit{autres} quintuplets!) qui soient des pré-Quillen à droite.

Voici le genre d'énoncé en forme, pour le moment conjectural voire totalement heuristique, auquel je songe :

\textbf{Théorème-scholie conjectural :} \textit{On suppose que $\cM$ est stable par $\varinjlim$ “quelconques”, et par $\varprojlim$ finies, et que $W$ est stable par $\varinjlim$ filtrantes (en plus de la saturation faible).}


\endgroup
% END UNIT Der-ChXIII-25-01-04/Der-ChXIII-25-01-04_p006-010.tex
% BEGIN UNIT Der-ChXIII-25-01-04/Der-ChXIII-25-01-04_p011-015.tex
\begingroup\makeatletter%
\newcommand{\printedpage}[1]{\marginpar{\footnotesize p.~#1}}%
\newcommand{\msp}[1]{\marginpar{\footnotesize [#1]}}%
\newcommand{\circled}[1]{\tikz[baseline=(c.base)]\node[draw,circle,inner sep=1pt](c){#1};}%
\newcommand{\Cof}{\mathrm{Cof}}%
\newcommand{\Fib}{\mathrm{Fib}}%
\newcommand{\Fl}{\mathrm{Fl}}%
\newcommand{\univ}{\mathrm{univ}}%
\makeatother\let\printedpage\derpp\let\msp\dermsp%
\let\mspin\DERmspin%


\printedpage{9}
% head: A. Grothendieck : Dérivateurs, Ch. XIII, Édition des Universités de Montpellier II et Paris VII

\begin{itemize}[leftmargin=2.5em]
\item[\circled{1}] \textit{Alors} $\Cof_W$ \textit{et} $W^{\univ}$ \textit{ont les propriétés de stabilité} a) \textit{et} c) \textit{de la proposition} 2 (p.~4), \textit{et aussi la propriété de stabilité} b) \textit{par facteurs directs, pourvu que} $W$ \textit{lui-même l'a.}

\item[\circled{2}] \textit{Supposons} $\mathcal{M}$ accessible, \textit{de sorte qu'on dispose du théorème} 1. \textit{Conditions équivalentes sur} $W$ \textit{:}

\msp{page 15}

\begin{itemize}[leftmargin=2.5em]
\item[a)] \textit{Il existe une partie petite} $C_0$ \textit{de} $\Cof_W$\textit{, et} $TC_0$ \textit{de} $W^{\univ}$\textit{, telles que}
\[ TF \overset{\text{déf}}{=} C_{0*} \subset W, \qquad (TC_0)_* \overset{\text{déf}}{=} F \subset \Fib_W. \]

\item[a$'$)] \textit{Il existe un quadruplet} $(C,TC,F,TF)$ \textit{de parties de} $\Fl(\mathcal{M})$\textit{, satisfaisant les conditions} M~1 \textit{à} M~4, M$'$, M$''$ \textit{de la page} 6 (\footnote[10]{avec $C$, $TC$ clos à gauche, $F$, $TF$ clos à droite.}).

\item[a$''$)] \textit{Comme} a$'$)\textit{, en supposant} M$'$ \textit{sous forme renforcée avec} M$'_0$ (g), i.e. \textit{on a une structure de pré-Q-catégorie de modèles à gauche (}$TC = C \cap W$\textit{).}

\item[b)] \textit{On a}
\[ (\Cof_W)_* \ \subset\ W, \qquad (W^{\univ})_* \ \subset\ \Fib_W \]
\textit{(donc} $(\Cof_W)_* \subset \Fib_W \cap W$ \textit{puisque} $(\Cof_W)_* \subset (W^{\univ})_*$ (\footnote[11]{$W^{\univ} \subset \Cof_W$.})\textit{).}

\item[b$'$)] \textit{On a}
\[ (\Cof_W)_* \ \subset\ W_{\univ}, \qquad (W^{\univ})_* \ \subset\ \Fib_W. \]
\end{itemize}

\item[\circled{3}] \textit{Pour qu'il existe un quadruplet qui, avec} $W$\textit{, définisse une pré-Q-catégorie de modèles à droite (}$TF = F \cap W$\textit{)} (\footnote[12]{On suppose toujours $\mathcal{M}$ accessible dans 3° et 4°, avec en plus $F$, $TF$ (et $C$, $TC$) closes à droite (resp. à gauche).})\textit{, il faut et il suffit qu'on ait} b$'$) \textit{ci-dessus, ainsi que}
\[ (W^{\univ})_* \cap W_{\univ} \ \ \text{\itshape est clos à droite.} \]

\msp{page 16}

\item[\circled{4}] \textit{Les hypothèses préliminaires sur} $(\mathcal{M}, W)$\textit{, y inclus l'accessibilité, ainsi que les conditions équivalentes dans}\circled{2}\textit{, sont stables par passage de} $(\mathcal{M}, W)$ \textit{à} $(\mathcal{M}(I), W(I))$\textit{. Mais qu'en est-il de la condition dans}\circled{3} \textit{? Est-elle conséquence des conditions dans} \circled{2} \textit{(}i.e. \textit{leurs est-elle équivalente) ?} % [sic] leurs est-elle -> lui est-elle, grammar
\end{itemize}

Je vais esquisser la démonstration du théorème-scholie, pour cerner le travail technique qui reste à faire.

1°) Est de la vérification de routine.

2°) Voici les implication $\pm$ tautologiques entre les conditions envisagées : % [sic] les implication -> les implications, agreement
\[ a'' \Longrightarrow a' \overset{?}{\Longrightarrow} a \Longrightarrow b \Longleftarrow b', \tag*{$(*)$} \]
% foot: Version 25 janvier 2004
\printedpage{10}
% head: A. Grothendieck : Dérivateurs, Ch. XIII, Édition des Universités de Montpellier II et Paris VII
mais je vois que $a' \Longrightarrow a$ n'est nullement tautologique, et peut-être même faux après tout~! (\footnote[13]{C'est pourquoi je l'ai modifié, en renforçant $a'$) par \texttt{[la condition de la note 10]}.}). Il faudrait supposer que dans \texttt{[le]} quadruplet $(C,F,TC,TF)$, $F$ et $TF$ sont clos à droite, et plus précisément, que $F = (TC)_*$, $TF = C_*$, ce qui sera le cas p.~ex. si on a une structure de Q-catégorie de modèles “close” (“closed model category”), au sens de \textsc{Quillen}. Les autres trois implications dans $(*)$ sont bel et bien tautologiques.

$a \Longrightarrow a'$. On peut dans $a$) supposer
\[ TC_0 \ \subset\ C_0 \]
(quitte à remplacer $C_0$ par $C_0 \cup TC_0$), puis, grâce à
\msp{page 17}
1°) ($W^{\univ}$ et $\Cof_W$ sont clos à gauche), posant
\[ \overline{TC_0} = TC, \qquad \overline{C_0} = C, \qquad \text{\texttt{[plutôt}}\ \ \widetilde{TC_0} = TC, \qquad \widetilde{C_0} = C, \text{\texttt{]}} \]
on aura $TC \subset W^{\univ}$, $C \subset \Cof_W$. Passant de $TC_0$ à $TC$, de $C_0$ à $C$, on a perdu seulement que ces parties de $\Fl(\mathcal{M})$ soient petites. Mais par le théorème 1, posant
\[ F \ =\ (TC)_*, \qquad TF \ =\ C_*, \]
les couples de parties closes (à gauche -- à droite)
\[ (C,\, TF), \qquad (TC,\, F) \]
satisfont à la condition de factorisation de Quillen. On aura alors les propriétés M~1 à M~4, plus les inclusions $TC \ \subset\ C \cap W$ (car $TC \ \subset\ W^{\univ} \ \subset\ W$, et $TC \ \subset\ C$ par construction), et $C \subset \Cof_W$ de M$'$ et M$''$. Il reste à vérifier $TF \subset F \cap W$, $F \subset \Fib_W$, mais $F$ $(= (TC)_*)$ $\ \subset \Fib_W$ par hypothèse, et $TF$ $(= C_*)$ $\ \subset W$ aussi~; d'autre part $TF \subset F$ résulte de $TC \subset C$.

$a' \Longrightarrow a''$. On veut donc, en plus des conditions M~1 à M~4 et M$'$, M$''$, la condition
\[ TC \ =\ C \cap W\ (= C \cap W^{\univ}). \]
L'idée, c'est de remplacer le couple
\msp{page 18}
\[ C,\ TC \qquad (C \subset \Cof_W,\ TC \subset C \cap W), \]
donnant naissance à
\[ F = (TC)_*, \qquad TF = C_*, \]
par
\[ C,\ C \cap W\ (\supset TC), \]
% foot: Version 25 janvier 2004
\printedpage{11}
% head: A. Grothendieck : Dérivateurs, Ch. XIII, Édition des Universités de Montpellier II et Paris VII
ce qui ne change pas le couple $(C, TF)$, mais bien le couple $(TC, F)$, remplacé par
\[ \underbrace{(TC)'}_{\supset TC} \ =\ \underbrace{C \cap W}_{= C\cap W^{\univ}}, \qquad F' \ =\ \underbrace{(TC')_*}_{\subset F} \]
(on a remplacé $F$ par un $F'$ plus petit) (\footnote[14]{NB $TC'$ sera encore clos à gauche comme intersection de deux parties closes à gauche.}). Le seul problème, c'est si $(TC', F')$ satisfait encore à la condition de factorisation de Quillen. Ce sont des parties closes (à gauche pour $TC'$, à droite pour $F'$) “duales” l'une de l'autre,
\[ TC' = F'^*, \qquad F' = TC'_* \ , \]
mais comme on ne sait pas si $TC'$ est de la forme $\widetilde{TC'_0}$, avec $TC'_0$ petite, on n'est pas sûr que la conclusion du théorème 1 s'applique. Mais si on suppose dans $a'$) que $C$ est une partie \textit{accessible} de $\Fl(\mathcal{M})$ (\textit{i.e.}, étant déjà close à gauche, que $C = \widetilde{C_0}$, $C_0$ petit),
\msp{page 19}
alors je présume qu'on pourra prouver qu'il en est de même de $C \cap W$ ($W$ elle-même étant accessible, grâce à la condition de stabilité par petites $\varinjlim$ filtrantes), donc le théorème 1 s'applique bel et bien à $TC' = C \cap W$. Mais si on garde $a'$) telle quelle, alors il faudrait sans doute prouver que $a'$) \textit{implique} $a'$) sous la forme plus forte, avec $C$ \textit{accessible.} Mais comme on a $a' \Longrightarrow a$, et qu'on vient d'établir que $a \Longrightarrow a'$ sous cette forme plus forte, on gagne (modulo le fait conjectural admis sur les parties admissibles \texttt{[plutôt accessibles]} de $\Fl(\mathcal{M})$).

Il reste donc à prouver
\[ b \ \Longrightarrow\ a, \qquad b \ \Longrightarrow\ b', \]
et il suffit de prouver la première implication, car on sait que $a$) $\Longrightarrow$ $a'$), et que $a'$) implique que
\[ \Fib_W \cap W \ =\ W_{\univ} \]
(cor.~1, page 10), donc $(\Cof_W)_* \subset W$ et $(W^{\univ})_* \subset \Fib_W$ (\textit{i.e.} $b$)) impliquent, comme $(\Cof_W)_* \subset (W^{\univ})_*$ (car $W^{\univ} \subset \Cof_W$), $(\Cof_W)_* \subset W \cap \Fib_W = W_{\univ}$, ce qui est la partie à établir de $b'$) (l'autre $(W^{\univ})_* \subset \Fib_W$ étant tautologiquement contenue dans $b$).

\msp{page 20}

Notons maintenant que $b$) n'est autre que la variante affaiblie de $a$), quand on omet de spécifier que $C_0$, $TC_0$ doivent être petits. Donc le seul problème, pour $b \Longrightarrow a$, c'est de voir ceci :

\textbf{Lemme conjectural} (\footnote[15]{Ce lemme a l'air un peu faux, mais ça doit pouvoir s'arranger \dots}).

\begin{itemize}[leftmargin=2.5em]
\item[\circled{1}] \textit{Si on a une partie close à gauche} $\Phi$ \textit{de} $\Fl(\mathcal{M})$\textit{, telle que} $\Phi_* \subset \Fib_W$ \textit{(resp.} $\Phi_* \subset W$\textit{), alors il existe une partie} petite $\Phi_0 \subset \Phi$\textit{, telle que} $\Phi_{0*} \subset \Fib_W$ \textit{(resp.} $\Phi_0 \subset W$\textit{).}

\item[\circled{2}] \textit{Si} $C$\textit{,} $C'$ \textit{sont deux parties closes à gauche de} $\Fl(\mathcal{M})$\textit{, si} $C$ \textit{et} $C'$ \textit{sont accessibles, il en est de même de} $C \cap C'$\textit{.}
\end{itemize}
% foot: Version 25 janvier 2004

\printedpage{12}
% head: A. Grothendieck : Dérivateurs, Ch. XIII, Édition des Universités de Montpellier II et Paris VII
Je crois que c'est bel et bien vrai sous les hypothèses faites sur $\mathcal{M}$, $W$ -- qui doivent impliquer que dans un sens convenable, $W$ est une partie accessible de l'ensemble $\Fl(\mathcal{M})$ des flèches de la catégorie accessible $\mathcal{M}$. Il doit en résulter qu'une partie de $\mathcal{M}$ construite explicitement en termes de $W$, telle $\Fib_W$, est elle-même accessible, et le lemme doit sortir de là.

Ainsi, 2° est prouvé, modulo ce lemme d'accessibilité, qu'il me faudra prouver $\pm$ en même temps que le théorème 1 (que je n'ai fait ici qu'énoncer).

\msp{page 21}

Prouvons 3°. Supposons d'abord qu'il existe un quadruplet “clos” $(C, TC, F, TF)$, s'insérant dans un “pré-quintuplet de Quillen clos” qui soit une Q-catégorie de modèles à droite, \textit{i.e.}
\[ TF \ =\ F \cap W \ =\ F \cap W_{\univ}. \]
Prouvons alors \texttt{[que]}
\[ \underbrace{(W^{\univ})_* \cap W_{\univ}}_{A_W} \qquad \text{est clos à droite.} \tag*{$(*)$} \]
En effet, on a
\[ TC \ \subset\ W^{\univ}, \]
d'où
\[ (W^{\univ})_* \ \subset\ (TC)_* \ =\ F, \]
donc
\[ A = (W^{\univ})_* \cap W_{\univ} \ \subset\ F \cap W_{\univ} \ \overset{\text{(par hyp.)}}{=}\ TF, \]
or $TF$ est clos à droite et contenu dans $W_{\univ}$, d'où
\[ \overline{A} \ \subset\ TF \ \subset\ W_{\univ}, \]
et comme $(W^{\univ})_*$ est clos à droite, on a aussi
\[ \overline{A} \ \subset\ (W^{\univ})_*, \]
d'où $\overline{A} \subset W_{\univ} \cap (W^{\univ})_* = A$, donc $A$ est clos à droite, q.e.d.

Prouvons la réciproque, \textit{i.e.} supposons satisfaites les conditions de 2°, \textit{plus} la condition $(*)$ ci-dessus. On doit construire le quadruplet.

\msp{page 22}

Notons d'abord, pour nous y reconnaître, que les relations $b'$) de 2°), et les relations transposées (compte tenu que $\Cof_W$ et $W^{\univ}$ sont clos à gauche, par 1°), savoir
\[ \left\{ \begin{array}{lllllll}
(\Cof_W)_* & \subset & W_{\univ} & , & (W^{\univ})_* & \subset & \Fib_W \\
(W_{\univ})^* & \subset & \Cof_W & , & (\Fib_W)^* & \subset & W^{\univ},
\end{array} \right. \]
donnant (\footnote[16]{compte tenu de $W^{\univ} \subset \Cof_W$, $W_{\univ} \subset \Fib_W$ et des inclusions transposées.}) les deux diagrammes carrés d'inclusions
% foot: Version 25 janvier 2004
\printedpage{13}
% head: A. Grothendieck : Dérivateurs, Ch. XIII, Édition des Universités de Montpellier II et Paris VII
\[
\begin{tikzcd}[column sep=small, row sep=small]
 & W^{\univ} \arrow[dr, hook] & \\
(\Fib_W)^* \arrow[ur, hook] \arrow[dr, hook] & & \Cof_W \\
 & (W_{\univ})^* \arrow[ur, hook] &
\end{tikzcd}
\qquad\qquad (\text{\footnotemark[17]})
\tag*{(I)}
\]
\footnotetext[17]{\textbf{NB} On a ces carrés d'inclusions, chaque fois qu'on a un système de Quillen $(W, C, F, TC, TF)$ \textit{clos.}}
et
\[
\begin{tikzcd}[column sep=small, row sep=small]
 & (W^{\univ})_* \arrow[dl, hook'] & \\
\Fib_W & & (\Cof_W)_*\ , \arrow[ul, hook'] \arrow[dl, hook'] \\
 & W_{\univ} \arrow[ul, hook', "\alpha"'] &
\end{tikzcd}
\tag*{(II)}
\]
dont le premier est transposé (par $\Phi \longmapsto \Phi^*$) du second, et où les 8 ensembles de flèches envisagés, à l'exception de $\Fib_W$ et $W_{\univ}$ (reliés par $\alpha$ dans le carré II) sont clos (les quatre ensembles de I clos à gauche, et les deux ensembles de II non intéressés par $\alpha$, clos à droite).

\msp{page 23}

Les “pré-quintuplets de Quillen clos”, associés à $W$, correspondent au choix de parties closes à gauche $TC$ et $C$, s'insérant dans le diagramme (I) comme montré dans I$'$ ci-dessous, et donnent lieu par transposition à $F$ et $TF$ comme dans II$'$ ci-dessous,
\[
\begin{tikzcd}[column sep=small, row sep=small]
 & & W^{\univ} \arrow[dr, hook] & \\
 & TC \arrow[ur, hook] \arrow[dr, hook] & & \Cof_W \\
(\Fib_W)^* \arrow[ur, hook] \arrow[dr, hook] & & C \arrow[ur, hook] & \\
 & (W_{\univ})^* \arrow[ur, hook] & &
\end{tikzcd}
\tag*{(I$'$)}
\]
\[
\begin{tikzcd}[column sep=small, row sep=small]
 & & (W^{\univ})_* \arrow[dl, hook'] & \\
 & F \arrow[dl, hook'] & & (\Cof_W)_* \arrow[ul, hook'] \arrow[dl, hook'] \\
\Fib_W & & TF \arrow[ul, hook'] \arrow[dl, hook'] & \\
 & W_{\univ}\ , \arrow[ul, hook'] & &
\end{tikzcd}
\tag*{(II$'$)}
\]
le “hic” étant donc de choisir $TC$ et $C$ assez grandes pour qu'on ait
\[ \underbrace{(TC)_*}_{= F} \ \subset\ \Fib_W, \qquad \underbrace{C_*}_{= TF} \ \subset\ W_{\univ} \qquad \text{--} \]
% foot: Version 25 janvier 2004


\endgroup
% END UNIT Der-ChXIII-25-01-04/Der-ChXIII-25-01-04_p011-015.tex
% BEGIN UNIT Der-ChXIII-25-01-04/Der-ChXIII-25-01-04_p016-020.tex
\begingroup\makeatletter%
\newcommand{\printedpage}[1]{\marginpar{\footnotesize p.~#1}}%
\newcommand{\msp}[1]{\marginpar{\footnotesize [#1]}}%
\newcommand{\univ}{\mathrm{univ}}%
\newcommand{\Cof}{\mathrm{Cof}}%
\newcommand{\Fib}{\mathrm{Fib}}%
\newcommand{\cercle}[1]{\mbox{\tikz[baseline=(c.base)]\node[draw,circle,inner sep=1pt](c){#1};}}%
\newcommand{\clarrow}{\mathrel{\leftarrow\mkern-4mu{+}\mkern-4mu\rightarrow}}%
\makeatother\let\printedpage\derpp\let\msp\dermsp%
\let\mspin\DERmspin%


\printedpage{14}
% head: A. Grothendieck : Dérivateurs, Ch. XIII, Édition des Universités de Montpellier II et Paris VII
% foot: Version 25 janvier 2004

\msp{page 24}
\noindent et les inclusions stipulées dans 2$^\circ$, $b'$) assurent qu'on peut bel et bien trouver $TC$ et $C$, quitte à les prendre les plus grandes possibles, $TC = W^{\univ}$, $C = \Cof_W$ -- auquel cas, d'ailleurs, on aura
\[
TC \;=\; C \cap W^{\univ},
\]
\textit{i.e.} on a une pré-Q-catégorie de modèles à gauche. Mais il s'agit d'assurer maintenant, par contre, que
\[
TF \;=\; F \cap W_{\univ},
\]
ce qui implique, on vient de le voir, (en notant $(W^{\univ})_* \subset F$) que
\[
\overline{(W^{\univ})_* \cap W_{\univ}} \;\subset\; \overline{F \cap W_{\univ}} = \overline{TF} = TF \subset W_{\univ}. \tag*{$(*)$}
\]
Inversement, si $(*)$ est satisfait, prenons
\[
F = (W^{\univ})_*
\]
le plus petit possible, \textit{i.e.} $C = W^{\univ}$ \texttt{[plutôt} $TC = W^{\mathtt{univ}}$\texttt{]} le plus grand possible. On doit donc prendre
\[
TF \;=\; F \cap W_{\univ} \;=\; (W^{\univ})_* \cap W_{\univ},
\]
et cet ensemble est bien clos à droite par hypothèse.

\textbf{Corollaire.} \textit{Supposons la condition} b$'$) \textit{du théorème-scholie satisfaite. Les conditions suivantes} a), b) \textit{sur $W$ sont équivalentes (et impliquent la condition de} 3$^\circ$ \textit{du théorème-scholie).}

\msp{page 25}
\begin{itemize}[leftmargin=3em]
\item[\cercle{a}] $(W^{\univ})_* \cap \overline{W_{\univ}} = (\Cof_W)_*$.

\item[\cercle{b}] $\widetilde{W^{\univ} \cup (W_{\univ})}^{*} = \Cof_W$, i.e. $\Cof_W$ \textit{est la partie close à gauche engendrée par} $W^{\univ}$ \textit{et} $(W_{\univ})^*$\textit{,} % doubt resolved: printed tilde ends before the final star; preserve the edition reading, see corrections.csv (PDF leaf 16).
\end{itemize}

\noindent\textit{et ces conditions impliquent}

\begin{itemize}[leftmargin=3em]
\item[\cercle{c}] \textit{qu'il existe une structure de} triple clos de Quillen $(W, C, F)$ \textit{incluant} $W$ \textit{(donc avec} $TC = W \cap C$, $TF = W \cap F$\textit{) (lequel sera clos au sens de Quillen si et seulement si} $W$ \texttt{[est]} \textit{stable par facteurs directs -- car $C$ et $F$ sont déjà stables spontanément).}
\end{itemize}

\noindent Il suffit en effet de prendre alors
\[
TC \;=\; W^{\univ}, \qquad\quad C \;=\; \Cof_W .
\]
Mais je doute que l'existence d'un triple clos de Quillen $(W, C, F)$ suffise à impliquer les conditions \textit{a}), \textit{b}) ci-dessus, ni même la condition plus faible
\[
(W^{\univ})_* \cap W_{\univ} \;=\; (\Cof_W)_*
\]
qui signifie \textit{exactement} que $C = \Cof_W$, $TC = W^{\univ}$, et leurs transposés $TF = (\Cof_W)_*$, $F = (W^{\univ})_*$ forment une catégorie de Quillen close.

\par
\printedpage{15}
% head: A. Grothendieck : Dérivateurs, Ch. XIII, Édition des Universités de Montpellier II et Paris VII
% foot: Version 25 janvier 2004

\msp{page 26}
\noindent Conditions sur $(\mathcal{M}, W)$ avec
\begin{itemize}[leftmargin=3em]
\item[\textit{a})] $\mathcal{M}$ stable par $\varprojlim$ finies et $\varinjlim$ (petites),
\item[\textit{b})] $W$ stable par $\varinjlim$ filtrantes,
\item[\textit{c})] $\mathcal{M}$ accessible.
\end{itemize}

\[
\mathrm{I}\ \left\{
\begin{array}{@{}l@{\quad}p{11cm}@{}}
\cercle{a} & existence d'un système de Quillen faible $(W, C, TC, F, TF)$ (clos si on y tient) \\[1ex]
\Updownarrow & \\[1ex]
\cercle{b} & existence d'un système de Quillen faible exact à gauche ($TC = C \cap W$). \\[1ex]
\Updownarrow & \\[1ex]
\cercle{c} & $\underbrace{(\Cof_W)_* \subset W,\ (W^{\univ})_* \subset \Fib_W}_{(*)}$ ($\Longleftrightarrow$ $(\Cof_W)_* \subset W_{\univ}$, $(W^{\univ})_* \subset \Fib_W$).
\end{array}\right.
\]

\noindent$\Uparrow$
\[
\mathrm{II}\ \left\{
\begin{array}{@{}l@{\quad}p{11cm}@{}}
\cercle{a} & existence d'un système de Quillen faible exact à droite ($TF = F \cap W$) \\[1ex]
\Updownarrow & \\[1ex]
\cercle{d} & On a $(*)$, et de plus $(W^{\univ})_* \cap W_{\univ}$ clos à droite. \\[1ex]
\Updownarrow & \\[1ex]
\cercle{b} & existence d'un système de Quillen clos (closed model category), $TC = C \cap W$, $TF = F \cap W$, $F = (TC)_*$, $TF = C_*$, $TC = F^*$, $C = (TF)^*$. \\[1ex]
\Updownarrow & \\[1ex]
\cercle{c} & Posant $TC = W^{\univ}$, $F = (TC)_*$, $TF = F \cap W_{\univ}$, $C = (TF)_*$, on trouve avec $W$ un système de Quillen clos. % [sic] $C = (TF)_*$ -> $C = (TF)^*$, cf. the line b) above where $C = (TF)^*$
\end{array}\right.
\]

\noindent$\Uparrow$
\[
\mathrm{III}\ \left\{
\begin{array}{@{}p{11.5cm}@{}}
$C = \Cof_W$ et $TC = W^{\univ}$ font partie (avec $W$) d'un système de Quillen clos (donc $F = (W^{\univ})_*$, $TF = (\Cof_W)_*$, $W = TF \circ TC$). \\[1ex]
$\Updownarrow$ \\[1ex]
\texttt{[On a]} $(*)$, et de plus $(W^{\univ})_* \cap W_{\univ} = (\Cof_W)_*$.
\end{array}\right.
\]

\noindent$\Uparrow$
\[
\mathrm{IV}\ \left\{
\begin{array}{@{}p{11.5cm}@{}}
On a $(*)$ et $\Cof_W = \widetilde{(W_{\univ})^* \cup W^{\univ}}$\ \ . \\[1ex]
$\Updownarrow$ \\[1ex]
On a $(*)$ et $(W^{\univ})_* \cap \overline{W_{\univ}} = (\Cof_W)_*$.
\end{array}\right.
\]

\par
\printedpage{16}
% head: A. Grothendieck : Dérivateurs, Ch. XIII, Édition des Universités de Montpellier II et Paris VII
% foot: Version 25 janvier 2004

\noindent Cas $\mathcal{M} = \mathrm{Cat}$, $W = \mathbf{W}_\infty$. On sait que II est satisfait, par \textsc{Thomason}. Donc on a
\[
(\Cof_W)_* \subset W_{\univ}, \qquad\quad (W^{\univ})_* \subset \Fib_W, \qquad\quad (W^{\univ})_* \cap W_{\univ} \ \text{ clos à droite.}
\]
J'ai pu prouver directement les deux premières inclusions (\textit{i.e.} I), mais pas que $(W^{\univ})_* \cap W_{\univ}$ soit clos à droite. (Je devrais au moins essayer!) (\footnote[18]{Il suffirait de savoir que $\overline{W_{\univ}} \cap \Fib_W = W_{\univ}$, ce qui impliquerait que III et IV sont équivalents (\footnotemark[19]).}\footnotetext[19]{Mais c'est faux, p. ex. si $X$ \texttt{[est une]} catégorie discrète non vide, $X \longrightarrow e$ est dans $\overline{W_{\univ}} \cap \Fib_W$.}). On n'a pas IV, \textit{i.e.} on a
\[
(W^{\univ})_* \cap \overline{W_{\univ}} \;\neq\; (\Cof_W)_*
\]
car
\[
(W^{\univ})_* \cap \overline{W_{\univ}} \not\subset W_{\univ}
\]
(contient les $X \longrightarrow e$ avec $X$ discrète), donc on a
\[
(W^{\univ})_* \cap \overline{W_{\univ}} \;\;\underset{\neq}{\supset}\;\; \underbrace{(W^{\univ})_* \cap W_{\univ}}_{\text{n'est \textit{pas} l'orthogonal de } W^{\univ} \cup (W_{\univ})^*} \quad . % doubt resolved: printed strict-superset sign confirmed; see corrections.csv (PDF leaf 18).
\]

\msp{page 27}

\section*{3\quad Catégorie de Quillen (deuxième mouture)}

On se donne toujours un couple
\[
(\mathcal{M}, W), \qquad\quad W \subset \mathrm{Fl}(\mathcal{M}), \qquad\quad W \text{ faiblement saturé,}
\]
$\mathcal{M}$ satisfait M 0 (stabilité par $\varprojlim$ finies et $\varinjlim$ finies).

\noindent\textbf{Définition 1 :} Une \textit{structure de Quillen faible} associée à la catégorie de modèles $(\mathcal{M}, W)$ est la donnée de quatre parties $(C, TC, F, TF)$ de $\mathrm{Fl}(\mathcal{M})$, satisfaisant les conditions suivantes.

\begin{itemize}[leftmargin=3em,labelsep=0.5em,align=left]
\item[\textbf{M 1}] $C \longleftrightarrow TF$, $TC \longleftrightarrow F$ (\textit{i.e.} $TF \subset C_*$, $F \subset (TC)_*$).
\item[\textbf{M 2}] Propriété de factorisation de Quillen pour $(C, TF)$ et $(TC, F)$.
\item[\textbf{M 3}] Stabilité de $C$ par composition et par cochangement de base, de $F$ par composition et changement de base -- $C$ et $F$ contiennent les isomorphismes.
\item[\textbf{M 4}] Stabilité similaire pour $TC$ et $TF$.
\item[$\mathbf{M}'$] $TC \subset C \cap W$, $TF \subset F \cap W$.
\item[$\mathbf{M}''$] $C \subset \Cof_W$, $F \subset \Fib_W$.
\end{itemize}

\noindent On dit que la structure de Quillen est \textit{exacte} si elle satisfait à l'axiome supplémentaire

\par
\printedpage{17}
% head: A. Grothendieck : Dérivateurs, Ch. XIII, Édition des Universités de Montpellier II et Paris VII
% foot: Version 25 janvier 2004

\noindent\textbf{M 2 bis}\qquad $W = TF \circ TC$.

\noindent On dit qu'elle est \textit{exacte à gauche} si elle satisfait l'axiome supplémentaire

\noindent\textbf{M$'_0$ (g)}\qquad $TC = C \cap W$ (au lieu de $TC \subset C \cap W$ seulement),

\noindent\textit{exacte à droite} si elle satisfait

\noindent\textbf{M$'_0$ (d)}\qquad $TF = F \cap W$ (au lieu de $TF \subset F \cap W$ seulement) (\footnote[20]{\textbf{NB} Si $C$, $F$, $W$ \texttt{[sont]} stables par facteurs directs (donc aussi $C \cap W$ et $F \cap W$), alors l'exactitude à gauche $TC = C \cap W$ \textit{équivaut} à l'exactitude à droite $TF = F \cap W$, et implique que $C$, $TC$, $F$, $TF$ sont \textit{clos} (à gauche pour $C$, $TC$, à droite pour $F$, $TF$).}).

\msp{page 28}
\noindent\textbf{Remarque.}
\begin{itemize}[leftmargin=3em]
\item[a)] Une structure de Quillen exacte à gauche ou exacte à droite est exacte.
\item[b)] Si on a une donnée $(C, TC, F, TF)$ satisfaisant M 1, M 2, M$'$, $M''$, quand on sature $C$, $TC$ à gauche et $F$, $TF$ à droite, pour les stabilités précisées dans M 2, M 3, on trouve une structure de Quillen. (Tous les axiomes sont évidents, sauf M$'$, pour lequel je renvoie à p. 9, 10, notamment prop. 4.) Si on avait M 2 bis, cela sera encore vérifié (a fortiori) après saturation, \textit{i.e.} on a une structure de Quillen exacte.
\end{itemize}

\noindent\textbf{Définition 1 bis.} La structure de Quillen faible est dite \textit{structure de Quillen} (tout court) si elle est exacte à gauche et à droite, ou (comme on dit) “biexacte”, alors $(C, F)$ determine déjà $(C, TC, F, TF)$. (Ça correspond exactement aux “model categories” de \textsc{Quillen} \texttt{[propres]}.) On dit que la structure de Quillen faible est \textit{close} si on a % [sic] determine -> détermine, missing accent
\[
C \clarrow TF\ , \qquad\quad TC \clarrow F\ ,
\]
\textit{i.e.} $TF = C_*$, $C = (TF)^*$ et $F = (TC)_*$, $TC = F^*$ (\footnote[21]{Cela signifie aussi que $C$, $TF$, $TC$, $F$ sont stables par facteurs directs. On dit que la structure de Quillen faible est \textit{close à gauche} (resp. \textit{à droite}) si $C$, $TC$ (resp. $F$, $TF$) sont clos à gauche (resp. à droite).}). Une structure de Quillen faible à la fois biexacte (\textit{i.e.} pas faible du tout!) et close est appelée structure de Quillen \textit{parfaite}. (Ça correspond exactement à la “closed model category” \texttt{[propre]} de \textsc{Quillen}.)

\msp{page 29}
\noindent\textbf{Remarque.} Pour qu'une structure de Quillen faible soit parfaite, il suffit qu'elle soit biexacte et qu'on ait
\[
C = (TF)^*, \qquad\quad F = (TC)_*
\]
(ce qui implique déjà que $TC = F^*$, $TF = C_*$, \textit{i.e.} que $TC$ est clos à gauche et $TF$ clos à droite). C'est prouvé dans Quillen. Ça implique aussi que $W$ est fortement saturé -- c'est le pays de cocagne !

\par
\printedpage{18}
% head: A. Grothendieck : Dérivateurs, Ch. XIII, Édition des Universités de Montpellier II et Paris VII
% foot: Version 25 janvier 2004

\noindent\textbf{Scholie :} Le couple $TC \longleftrightarrow F$, satisfaisant à la condition de factorisation de Quillen, et satisfaisant $TC \subset W$, $F \subset \Fib_W$, permet déjà de développer une théorie satisfaisante des fibres $W$-homotopiques dans $\mathcal{M}$ -- mais sans propriété d'exactitude pour les carrés $W$-cartésiens. Dualement, la donnée de $C \longleftrightarrow TF$, avec factorisation, et $TF \subset W$, $C \subset \mathrm{Cofib}_W$, donne une bonne théorie des cofibres $W$-homotopiques, mais sans propriété d'exactitude pour les carrés $W$-cocartésiens. La structure de Quillen faible réunit ces deux données duales l'une de l'autre, avec en plus le lien supplémentaire entre ces deux donnés, par $TC \subset C$, $TF \subset F$ (stipulé dans $M'$). Cette donnée assure déjà que % [sic] donnés -> données, feminine noun as in "ces deux données duales"
\[
\underbrace{\mathcal{M}_{\mathrm{cf}}}_{\substack{\text{sous-catégorie pleine}\\ \text{de $\mathcal{M}$ formée des}\\ \text{objets fibrants et}\\ \text{cofibrants à la fois}}} W_{\mathrm{cf}}^{-1} \;\xrightarrow{\ \sim\ }\; \mathcal{M}W^{-1} = \mathrm{Ho}_W
\]
\msp{page 30}
est une équivalence, mais non pas que le foncteur localisation
\[
\mathcal{M}_{\mathrm{cf}} \longrightarrow \mathcal{M}_{\mathrm{cf}}W_{\mathrm{cf}}^{-1}
\]
est un simple passage au quotient sur les flèches, \textit{i.e.} que si $X$, $Y$ sont dans $\mathcal{M}_{\mathrm{cf}}$, alors
\[
\mathrm{Hom}_{\mathcal{M}}(X, Y) \;\longrightarrow\; \underbrace{\mathrm{Hom}_W(X, Y)}_{\overset{\text{déf}}{=}\,\mathrm{Hom}_{\mathrm{Ho}_W}(X, Y)} \tag*{$(*)$}
\]
soit surjectif. Ceci est assuré par l'axiome d'exactitude M 2 bis,
\[
W \;=\; TF \circ TC.
\]
Avec ce seul axiome, on n'a pas de caractérisation sympathique de la relation d'équivalence dans $\mathrm{Hom}_{\mathcal{M}}(X, Y)$ définie par la flèche $(*)$ (pour $X$, $Y$ dans $\mathcal{M}_{\mathrm{cf}}$).

Si on a exactitude à gauche,
\[
TC = C \cap W \qquad\quad (\Longleftrightarrow\ TC = C \cap W^{\univ}),
\]
alors cette relation d'équivalence est celle engendrée par la relation de \textit{Q-homotopie à droite} (\footnote[22]{Il semblerait que la relation de Q-homotopie à droite est une relation d'équivalence, dès que $Y$ \texttt{[est]} dans $\mathcal{M}_{\mathrm{f}}$ et sans supposer la structure de Quillen faible exacte à droite ou à gauche.}, \footnote[23]{\textbf{NB} D'après \textsc{Baues}, il semblerait que sur $\mathcal{M}_{\mathrm{cf}}$, cette relation d'homotopie à droite équivaut à la relation d'homotopie à gauche, et que celle-ci ne dépend pas du cylindre choisi \ldots}) (impliquée par la Q-homotopie à gauche), et la catégorie quotient ainsi définie
\[
\pi_{\mathrm{d}}\mathcal{M}_{\mathrm{cf}} \;=\; \pi\mathcal{M}_{\mathrm{cf}} \;\simeq\; \mathrm{Ho}_W,
\]
s'envoie dans $\pi_{\mathrm{d}}\mathcal{M}_{\mathrm{c}}$,
\[
\pi\mathcal{M}_{\mathrm{cf}} \longrightarrow \pi_{\mathrm{d}}\mathcal{M}_{\mathrm{c}},
\]
et ce foncteur peut être interprété (sauf erreur)


\endgroup
% END UNIT Der-ChXIII-25-01-04/Der-ChXIII-25-01-04_p016-020.tex
% BEGIN UNIT Der-ChXIII-25-01-04/Der-ChXIII-25-01-04_p021-025.tex
\begingroup\makeatletter%
\derleqno%
\newcommand{\printedpage}[1]{\marginpar{\footnotesize p.~#1}}%
\newcommand{\msp}[1]{\marginpar{\footnotesize [#1]}}%
\newcommand{\cM}{\mathcal{M}}%
\newcommand{\circled}[1]{\tikz[baseline=(c.base)]\node[draw,circle,inner sep=1pt](c){#1};}%
\makeatother\let\printedpage\derpp\let\msp\dermsp%
\let\mspin\DERmspin%


\setcounter{footnote}{23}

\printedpage{19}
% head: A. Grothendieck : Dérivateurs, Ch. XIII, Édition des Universités de Montpellier II et Paris VII
% foot: Version 25 janvier 2004
\msp{page 31}
comme un \textit{adjoint à droite} (pleinement fidèle) du foncteur de localisation
\[
\pi_{\mathrm{d}}\cM_{\mathrm{c}} \longrightarrow (\pi_{\mathrm{d}}\cM_{\mathrm{c}})W_{\mathrm{c}}^{-1} \;\simeq\; \mathrm{Ho}_W,
\]
et $\mathrm{Ho}_W$ se déduit de $\pi_{\mathrm{d}}\cM_{\mathrm{c}}$ par un \textit{calcul de fractions à gauche.} Dualement, si on a exactitude à droite, \textit{i.e.} $TF = F\cap W$ ($\Longleftrightarrow$ $TF = F\cap W_{\mathrm{univ}}$), alors la relation d'équivalence pertinente sur $\mathrm{Hom}_{\cM}(X,Y)$, pour $X$, $Y$ dans $\cM_{\mathrm{cf}}$, est la Q-homotopie à gauche (impliquée par la Q-homotopie à droite), et cette fois la flèche naturelle de $\pi_{\mathrm{g}}\cM_{\mathrm{cf}} = \pi\cM_{\mathrm{cf}} \simeq \mathrm{Ho}_W$ dans $\pi_{\mathrm{d}}\cM_{\mathrm{f}}$ \texttt{[plutôt }$\pi_{\mathtt{g}}\cM_{\mathtt{f}}$\texttt{]},
\[
\underbrace{\pi\cM_{\mathrm{cf}}}_{\simeq\,\mathrm{Ho}_W} \longrightarrow \pi_{\mathrm{d}}\cM_{\mathrm{f}} \qquad \texttt{[plutôt }\pi_{\mathtt{g}}\cM_{\mathtt{f}}\texttt{]}
\]
est un \textit{adjoint à gauche} (pleinement fidèle) du foncteur de localisation
\[
\pi_{\mathrm{d}}\cM_f \longrightarrow \cM_{\mathrm{f}}W_{\mathrm{f}}^{-1} \;\simeq\; \mathrm{Ho}_W \qquad \texttt{[plutôt }\pi_{\mathtt{g}}\cM_f \longrightarrow \cM_{\mathtt{f}}W_{\mathtt{f}}^{-1} \;\simeq\; \mathrm{Ho}_W\texttt{]}, % [sic] \cM_f (italic f, twice) -> \cM_{\mathrm{f}}, elsewhere the index f is upright
\]
et la localisée s'explicite par un calcul de fractions à droite.

Dans le cas biexacte (\textit{i.e.} celui des véritables structures de Quillen), les deux relations de Q-homotopie dans $\cM_{\mathrm{cf}}$ (plus généralement dans $\mathrm{Hom}(X,Y)$, si $X$ dans $\cM_{\mathrm{c}}$ et $Y$ dans $\cM_{\mathrm{f}}$)
\msp{page 32}
coïncident, et deviennent même particulièrement simples. (Car si $f, g : X \rightrightarrows Y$, avec $X$ cofibrant et $Y$ fibrant, alors $f\sim g$ peut se tester sur \textit{n'importe quel} “objet cylindre” pour $X$, ou sur \textit{n'importe quel} “objet chemins” pour $Y$.)

Dans le cas où la structure de Quillen faible est exacte à gauche, \textit{i.e.} où la relation d'équivalence pertinente dans $\mathrm{Fl}\,\cM_{\mathrm{cf}}$ est la Q-homotopie à droite, alors l'axiome d'exactitude est vérifié pour les carrés $W$-cartésiens, mais je doute qu'il en soit de même pour les carrés $W$-cocartésiens (\footnote{Il semble pourtant que c'est le cas, d'après \textsc{Baues}, cf. annotation marginale correspondante page 30.}). (Si on ne fait une hypothèse supplémentaire, p. ex. l'exactitude à droite de la structure de Quillen.) Si on n'a que l'exactitude (mais ni à gauche ni à droite), il semble qu'on ne peut en conclure l'exactitude ni des carrés $W$-cartésiens, ni des carrés $W$-cocartésiens, donc pas celle des suites de cosuspension (loop sequences), ni des suites de suspension. On pourrait donc penser que la condition d'exactitude, à elle seule, ne sert pas à grand
\msp{page 33}
chose pour le formalisme homotopique dans les catégories de fractions $\mathrm{Ho}_W$ etc. Mais on a pourtant ceci : soit $\Psi = {}$\begin{tikzcd}[row sep=tiny, column sep=small, baseline=(\tikzcdmatrixname-2-1.base)]
 & 1 \\
0 \arrow[ru] \arrow[rd] & \\
 & 2
\end{tikzcd}, et supposons que $(\cM,W)$ \textit{et} $(\cM(\Psi),W(\Psi))$ \texttt{[où pour toute petite catégorie }$I$\texttt{, }$\cM(I)$\texttt{ désigne la catégorie des préfaisceaux sur }$I$\texttt{ à valeurs dans }$\cM$\texttt{, et }$W(I)$\texttt{ est la partie de }\texttt{Fl}$(\cM(I))$\texttt{ formée des morphismes}
\printedpage{20}
% head: A. Grothendieck : Dérivateurs, Ch. XIII, Édition des Universités de Montpellier II et Paris VII
% foot: Version 25 janvier 2004
\texttt{de préfaisceaux qui sont dans }$W$\texttt{ "argument par argument"]} admettent une structure de Quillen faible exacte. Alors dans $\cM$ les carrés $W$-cartésiens sont $W$-exactes, donc on a dans $\cM$ des suites exactes longues de cosuspension (exact loop sequences) (\footnote{il semble que je me suis canulé ici, et je doute que le résultat évoqué soit exact.}). (L'hypothèse sur $(\cM(\Psi),W(\Psi))$ sert uniquement à assurer que toute flèche de $\mathrm{Ho}_{W(\Psi)} = \cM(\Psi)W(\Psi)^{-1}$ est isomorphe à une flèche provenant d'une flèche de $\cM(\Psi)$ lui-même. Mais c'est là une propriété apparemment très délicate, et je ne connais d'autre façon de l'obtenir que via l'existence d'une structure de Quillen faible exacte.) Dualement, si $(\cM,W)$ et $(\cM(\Phi),W(\Phi))$ admettent des structures de Quillen faibles exactes, alors les carrés $W$-cocartésiens dans $\cM$ sont $W$-exacts, et on trouve (si $\cM$ est ponctué) des suites exactes longues de suspension. % [sic] W-exactes -> W-exacts, agreement with "carrés" (cf. "W-exacts" below)

\msp{page 34}
Ceci montre que si, pour une catégorie de modèles $(\cM,W)$ donnée, on s'intéresse à “faire de l'algèbre homotopique” pour tous les couples $(\cM(I),W(I))$ à la fois, avec $I$ dans Cat, il suffit (si on veut avoir une bonne théorie des fibres et cofibres $W$-homotopiques dans ces catégories, y inclus la propriété de $W$-exactitude pour les carrés $W$-cartésiens et pour les carrés $W$-cocartésiens) de disposer pour chaque $(\cM(I),W(I))$ d'une structure de Quillen faible exacte sans plus.

Ceci nous montre donc l'intérêt de développer des critères, pour une $(\cM,W)$, assurant qu'il en est bien ainsi. À ce sujet, je note que je ne sais comment vérifier l'exactitude d'une structure de Quillen \textit{que si} je dispose de l'exactitude à gauche, ou à droite. C'est dire que je dois dégager en fait des conditions sur $(\cM,W)$ qui assurent (disons) que les $(\cM(I),W(I))$ admettent des structures de Quillen faibles exactes à gauche.

\msp{page 35}
Les structures de Quillen faibles pour les $(\cM(I),W(I))$ seront déduites d'une structure de Quillen faible pour $(\cM,W)$, de la façon suivante : elles seront toutes \textit{“closes à droite”,} \textit{i.e.}
\[
F(I) \;=\; (TC(I))_*, \qquad TF(I) \;=\; (C(I))_*, \tag*{(1)}
\]
en particulier (pour $I=e$)
\[
F \;=\; (TC)_*, \qquad TF \;=\; C_*,
\]
ce qui implique déjà (puisque $C\subset \mathrm{Cof}_W$, $TC\subset W^{\mathrm{univ}}$)
\[
(\mathrm{Cof}_W)_* \;\subset\; W, \qquad (W^{\mathrm{univ}})_* \;\subset\; \mathrm{Fib}_W \tag*{(2)}
\]
(d'où $(\mathrm{Cof}_W)_*\subset W_{\mathrm{univ}}$). En vertu de (1), pour définir les structures de Quillen faibles pour les $(\cM(I),W(I))$, il suffit de définir $C(I)$, $TC(I)$, et même, comme on veut l'exactitude à gauche, \textit{i.e.}
\[
TC(I) \;=\; C(I)\cap W(I) \tag*{(3)}
\]
\printedpage{21}
% head: A. Grothendieck : Dérivateurs, Ch. XIII, Édition des Universités de Montpellier II et Paris VII
% foot: Version 25 janvier 2004
(en particulier
\[
TC \;=\; C\cap W \quad ), \tag*{(4)}
\]
il reste à définir $C$. On définira $C(I)$ “argument par argument”, \textit{i.e.} pour $f\in\mathrm{Fl}(\cM(I))$
\[
f\in C(I) \qquad\Longleftrightarrow\qquad \alpha_s^*(f)\in C \ \text{ pour tout } s\in\mathrm{Ob}\,I, \tag*{(5)}
\]
où
\[
\alpha_s : e\longrightarrow I \qquad \texttt{[est]}\ \text{défini par } \alpha_s(e)=s.
\]
Il en résulte d'ailleurs aussitôt, par (3),
\[
f\in TC(I) \qquad\Longleftrightarrow\qquad \alpha_s^*(f)\in TC \ \text{ pour tout } s\in\mathrm{Ob}\,I. \tag*{(6)}
\]

\msp{page 36}
On suppose donc que $(\cM,W)$ admet une structure de Quillen faible exacte à gauche et close à droite $(C,TC,F,TF)$, définie en termes de
\[
C\;\subset\;\mathrm{Fl}(\cM),
\]
par
\[
TC \;=\; C\cap W, \qquad F \;=\; (TC)_*, \qquad TF = C_*.
\]
Pour que cela donne bel et bien une structure de Quillen faible (forcément exacte à gauche et close à droite), il faut simplement vérifier les conditions suivantes.

\begin{itemize}[leftmargin=3em]
\item[\circled{i}] Factorisation pour $(C, TF=C_*)$ et $(TC=C\cap W, F=(TC)_*)$ – on reviendra sur ce point.
\item[\circled{ii}] Stabilité de $C$ par composition et cochangement de base, et $C$ contient les isomorphismes (d'où résultera la stabilité similaire de $TC = C\cap W$, compte tenu que $C\subset\mathrm{Cof}_W$) (NB Les stabilités idoines pour $F$, $TF$ sont triviales par définition $F=(TC)_*$, $TF=C_*$.)
\item[\circled{iii}] $C\subset\mathrm{Cof}_W$, $\underbrace{C_*}_{=\,TF}\subset W$, $(C\cap W)_*\subset\mathrm{Fib}_W$.
\end{itemize}

Pour vérifier, de même, que $(C(I),TC(I),F(I),TF(I))$ soit une structure de Quillen faible pour $W(I)$, il faut vérifier les conditions similaires à \circled{i}\,\circled{ii}\,\circled{iii} . Or (\textit{ii}) pour $C(I)$ résulte aussitôt de la condition (\textit{ii}) pour $C$. Il reste à examiner (\textit{iii}) et (\textit{i}).

\printedpage{22}
% head: A. Grothendieck : Dérivateurs, Ch. XIII, Édition des Universités de Montpellier II et Paris VII
% foot: Version 25 janvier 2004
\msp{page 37}
Les inclusions de (\textit{iii}).

L'inclusion
\[
C\;\subset\;\mathrm{Cof}_W
\]
implique que toute $f\in C(I)$ est dans $\mathrm{Cof}_W$ argument par argument. Cela implique que $f$ est dans $\mathrm{Cof}_{W(I)}$, comme on voit aussitôt sur la définition – je l'explicite en le

\textbf{Lemme 1 :} \textit{Soit $f : X_\bullet\longrightarrow Y_\bullet$ une flèche de $\cM(I)$. Pour que $f$ soit dans} $\mathrm{Cof}_{W(I)}$ (resp. \textit{dans} $\mathrm{Fib}_{W(I)}$), \textit{il} suffit \textit{que pour tout $s$ dans} $\mathrm{Ob}\,I$, \textit{$\alpha_s^*(f)$ soit dans} $\mathrm{Cof}_W$ (resp. \textit{dans} $\mathrm{Fib}_W$)\textit{. Itou (pour mémoire) pour une condition pour que $f\in W(I)$ – là il faut et il suffit que $\alpha_s^*(f)\in W\ \forall\, s\in\mathrm{Ob}\,I$.}

Nous allons appliquer ce même lemme pour établir les inclusions
\[
C(I)_* \;\subset\; W(I), \qquad (C(I)\cap W(I))_* \;\subset\; \mathrm{Fib}_{W(I)}.
\]
Donc il suffit de voir que
\[
\begin{aligned}
\alpha_s^*(C(I)_*) \;&\subset\; C_* &&\qquad (\subset\; W)\\
\alpha_s^*((\underbrace{C(I)\cap W(I)}_{(C\cap W)(I)})_*) \;&\subset\; (C\cap W)_* &&\qquad (\subset\; \mathrm{Fib}_W),
\end{aligned}
\]
et plus généralement
\[
\alpha_s^*(\Phi(I)_*) \;\subset\; \Phi_* \tag*{(7)}
\]
pour toute ( ?) partie $\Phi$ de $\mathrm{Fl}(\cM)$, définissant donc $\Phi(I)\subset\mathrm{Fl}(\cM(I))$ terme à terme.

\msp{page 38}
En explicitant la chose, on trouve la condition
\[
\alpha_{s\,!}(\Phi)\subset\Phi(I) \qquad \text{pour tout } s\in\mathrm{Ob}\,I,
\]
pour que ça marche, \textit{i.e.} on veut que
\[
i\in\Phi \qquad\Longrightarrow\qquad \alpha_{s\,!}(i)\in\Phi(I), \quad \textit{i.e.}\ \ \alpha_t^*(\alpha_{s\,!}(i))\in\Phi\ \ \forall\, t\in\mathrm{Ob}\,I.
\]
Or le foncteur $\alpha_t^*\alpha_{s\,!}$ s'explicite comme
\[
A \longmapsto \underbrace{\mathrm{Hom}_I(t,s)\times A}_{\substack{\text{somme de } \mathrm{Hom}_I(t,s)\\ \text{copies de } A \text{ dans } \cM}},
\]
donc il faut supposer $\Phi$ \textit{stable par sommes directes} (du type envisagé $\mathrm{Hom}_I(t,s)\times A$). C'est bien le cas pour $\Phi = C$ et $\Phi = C\cap W$ et pour des sommes directes \textit{finies,} grâce à (\textit{ii}). On trouve donc le

\printedpage{23}
% head: A. Grothendieck : Dérivateurs, Ch. XIII, Édition des Universités de Montpellier II et Paris VII
% foot: Version 25 janvier 2004
\textbf{Lemme 2.}

\begin{itemize}[leftmargin=3em]
\item[a)] \textit{Soit $I$ dans} Cat \textit{telle que les sommes directes de type} $\mathrm{Hom}_I(t,s)$ \textit{existent dans $\cM$. (P. ex. on suppose que $I$ est} finie, \textit{plus généralement que les} $\mathrm{Hom}_I(t,s)$ \textit{sont des ensembles finis ; ou que les “petites” sommes directes existent dans $\cM$.) Soit $\Phi$ une partie de} $\mathrm{Fl}(\cM)$\textit{, stable par ce type de sommes directes, et soit} $\Phi(I)\subset\mathrm{Fl}(\cM(I))$ \textit{la partie correspondante. Alors pour tout $s\in\mathrm{Ob}\,I$, on a}

\msp{page 39}
\[
\alpha_s^*(\Phi(I)_*) \;\subset\; \Phi_*.
\]
\item[b)] \textit{Supposons que les} $\mathrm{Hom}_I(t,s)$ \textit{soient finis (p. ex. $I$ finie). Alors l'hypothèse de} a) \textit{s'applique pour $\Phi=C$ et pour $\Phi=C{\cap}W$, donc la condition} (iii) \textit{p.} 36 \textit{est satisfaite pour la structure $(W(I),C(I),TC(I),F(I),TF(I))$ dans $\cM(I)$. Même conclusion sans hypothèse de finitude sur les} $\mathrm{Hom}_I(t,s)$\textit{, quand on suppose $\cM$ stable par “petites” sommes directes, et $C$ et $W$ stables également par de telles sommes directes.}
\end{itemize}

\textbf{NB} Vu le scholie plus haut, il est parfaitement raisonnable de se borner au besoin aux catégories $I$ \textit{finies} – vu que la catégorie de ces $I$ est stable par les produits $I\longmapsto I\times\Psi$, $I\longmapsto I\times\Phi$.

Il reste, pour être heureux, à vérifier la condition de factorisation \circled{i} pour $(C,TF)$, $(TC,F)$, mais aussi et surtout pour les $(C(I),TF(I),TC(I),F(I))$. C'est dire que de partir même d'une structure de Quillen close (pays de cocagne)
\msp{page 40}
ne nous avance guère à lui tout seul, faute de savoir que l'axiome de factorisation se conserve, en passant de $(\cM,W,\ldots)$ à $(\cM(I),W(I),\ldots)$.

C'est ici que le théorème de factorisation du no.~1 doit s'avérer providentiel. Mais pour pouvoir l'appliquer aux couples $(C(I),C(I)_*)$, $(TC(I),TC(I)_*)$, il faut (en plus de l'accessibilité de $\cM$, ce qui est une condition fort anodine) des conditions de stabilité très fortes pour $C(I)$, $TC(I)$ (\footnote{Plus les \textit{petits} systèmes de générateurs pour $C$, $TC$.}). À vrai dire, il reste la stabilité de $C$ et de $C\cap W$ par limites inductives ordinales (ce qui impliquera la même stabilité pour les $C(I)$, $(C{\cap}W)(I) = C(I){\cap}W(I)$).

À dire vrai, je vois peu de chances de trouver une structure de Quillen faible $(C,TC,F,TF)$ associé à $W$, et satisfaisant toutes les conditions dites, y compris celle de stabilité de $C$ et de $C\cap W$ par $\varinjlim$ ordinales, sans supposer que $W$ elle-même ne soit déjà stable par (petites) limites inductives filtrantes. % [sic] associé -> associée, agreement with "structure"

\msp{page 41}
On peut songer à utiliser une factorisation \textit{fonctorielle} des flèches dans $\cM$ en
\[
\begin{aligned}
f = pi \quad \text{avec}\quad & i\in TC, && p\in F\\
\text{ou}\quad & i\in C, && p\in TF,
\end{aligned}
\]


\endgroup
% END UNIT Der-ChXIII-25-01-04/Der-ChXIII-25-01-04_p021-025.tex
% BEGIN UNIT Der-ChXIII-25-01-04/Der-ChXIII-25-01-04_p026-030.tex
\begingroup\makeatletter%
\newcommand{\printedpage}[1]{\marginpar{\footnotesize p.~#1}}%
\newcommand{\msp}[1]{\marginpar{\footnotesize [#1]}}%
\newcommand{\limind}{\mathop{\underrightarrow{\mathrm{lim}}}}%
\newcommand{\limproj}{\mathop{\underleftarrow{\mathrm{lim}}}}%
\newcommand{\M}{\mathcal{M}}%
\newcommand{\Fl}{\mathrm{Fl}}%
\newcommand{\Ob}{\mathrm{Ob}}%
\newcommand{\Cof}{\mathrm{Cof}}%
\newcommand{\Fib}{\mathrm{Fib}}%
\newcommand{\univ}{\mathrm{univ}}%
\newcommand{\HO}{\mathrm{HO}}%
\makeatother\let\printedpage\derpp\let\msp\dermsp%
\let\mspin\DERmspin%


\printedpage{24}
% head: A. Grothendieck : Dérivateurs, Ch. XIII, Édition des Universités de Montpellier II et Paris VII
pour déduire une factorisation similaire pour les flèches $f$ de $\M(I)$, en
\[
\begin{aligned}
f = pi \quad \text{avec}\quad & i \in TC(I), && p \in F(I)\\
\text{ou}\quad & i \in C(I), && p \in TF(I),
\end{aligned}
\]
où $C(I)$, $TC(I)$, $F(I)$, $TF(I)$ sont ici définis via $C$, $TC$, $F$, $TF$ “argument par argument”. Mais l'ennui, c'est que en procédant de cette façon naïve, on n'aura plus nécessairement $C(I) \longleftrightarrow TF(I)$, $TC(I) \longleftrightarrow F(I)$. Il faut considérablement rapetisser $TF(I)$ et $F(I)$, pour avoir une propriété RLP par rapport à $C(I)$ resp. $TC(I)$.

Donc je ne vois guère d'autre moyen de s'en tirer, dans ce contexte général, qu'en supposant la stabilité de $W$ par $\limind$ filtrantes (ce qui n'est après tout pas très méchant).

\noindent\msp{page 42}

Bon gré, mal gré, on revient à la situation envisagée dans le paragraphe 2. On va réénoncer l'essentiel de ce qui a été pressenti alors :

\textbf{Théorème 2} (encore un peu conjectural). {\itshape Soit $(\M,W)$ une catégorie de modèles, avec $W$ (comme toujours) faiblement saturée, mais de plus on suppose $\M$ stable par petites $\limind$ et par $\limproj$ finies, et $W$ stable par $\limind$ filtrantes, enfin $\M$ accessible. Considérons les conditions suivantes :

\begin{enumerate}
\item[\textup{a)}] Il existe une structure de Quillen faible pour $W$, exacte à gauche, close à droite, $(C,TC,F,TF)$.
\item[\textup{a$'$)}] Comme \textup{a)}, mais avec de plus $C$, $TC$ clos à gauche (\textup{i.e.} on a une structure de Quillen faible close), et $C$, $TC$ de la forme $\widetilde{C_0}$, $\widetilde{TC_0}$, où $C_0$ et $TC_0$ sont petits.
\item[\textup{a$''$)}] Il existe $C_0 \subset \Cof_W$, $TC_0 \subset W^{\univ}$, parties petites, telles que $C_{0*} \subset W$, $(TC_0)_* \subset \Fib_W$.
\item[\textup{b)}] On a
\[
(\Cof_W)_* \;\subset\; W_{\univ}, \qquad\quad (W^{\univ})_* \;\subset\; \Fib_W
\]
(d'où résulte, $\Cof_W$ et $W^{\univ}$ étant clos à gauche, que
\[
(W_{\univ})^* \subset \Cof_W, \qquad\quad (\Fib_W)^* \subset W^{\univ} \qquad\quad \text{\itshape ).}
\]
\end{enumerate}

\noindent\msp{page 43}

Ces conditions impliquent que pour toute $I$ dans \textup{Cat}, il existe une structure de Quillen faible exacte à gauche et close pour $W(I)$ dans $\M(I)$. Une telle structure se définit, en termes d'une structure similaire $C$, $TC = C \cap W$, $F = (TC)_*$, $TF = C_*$ pour $W$, par}
\[
C(I) \;=\; \Bigl\{ f \in \Fl(\M(I)) \Bigm| \alpha_s^*(f) \in C \ \ \forall\, s \in \Ob I \Bigr\},
\]

\printedpage{25}
% head: A. Grothendieck : Dérivateurs, Ch. XIII, Édition des Universités de Montpellier II et Paris VII
\noindent\textit{puis}
\[
\begin{aligned}
TC(I) \;&=\; C(I)\cap W(I) \;=\; \Bigl\{ f \in \Fl(\M(I)) \Bigm| \alpha_s^*(f) \in TC \ \ \forall\, s \in \Ob I \Bigr\}\\
F(I) \;&=\; (TC(I))_*\\
TF(I) \;&=\; C(I)_*\,.
\end{aligned}
\]

\textbf{Corollaire-Scholie.} \textit{Ces conditions impliquent que pour tout $I$ dans} Cat, \textit{$(\M(I),W(I))$ donne lieu à une bonne théorie des fibres et cofibres $W(I)$-homotopiques, avec carrés $W(I)$-cartésiens $W(I)$-exacts, et carrés $W(I)$-cocartésiens $W(I)$-exacts – donc avec des suites} exactes \textit{longues de cosuspension et de suspension, dans le cas où $\M$, donc aussi $\M(I)$, est une catégorie avec objet neutre.}

\noindent\msp{page 44}

Je m'occuperai de la démonstration du théorème précédent, ainsi que du théorème de factorisation du paragraphe 1, quand j'aurai mes notes sur les catégories et foncteurs accessibles et filtrations cardinales. En attendant, j'ai envie déjà d'expliciter un peu plus le formalisme auquel on parvient moyennant les hypothèses faites ci-dessus. (Je suis à peu près sûr tout au moins que la condition $a''$) suffit à impliquer toutes les autres. Et c'est bien sous la forme $a''$) que ces conditions vont se vérifier en pratique, je présume, notamment dans le cas où $\M = $ Cat.)

\section*{4\quad Les variances $u^*$, $u_*$, $u_{!}$ pour $(\M(I),W(I))$ (I variable)}

Je suppose les conditions du paragraphe précédent vérifiées pour $(\M,W)$ (cf. théorème p. 42), et choisis une structure de Quillen faible close et exacte à gauche $(C,TC,F,TF)$ pour $W$, d'où pour tout $I$ dans Cat

\noindent\msp{page 45}

\noindent une structure de Quillen faible close et exacte à gauche $(C(I),TC(I),F(I),TF(I))$ pour $W(I)$, dans $\M(I)$, décrite page 43. Soit maintenant
\[
u \;:\; I' \longrightarrow I
\]
une flèche dans Cat, d'où une flèche
\[
u^* \;:\; (\M(I),W(I)) \;\longrightarrow\; (\M(I'),W(I')).
\]
On sait, vu que $\M$ est stable par $\limind$ quelconques et aussi (étant de plus accessible) par $\limproj$ quelconques, que $u^*$ admet des foncteurs adjoints à gauche et à droite,
\[
u_{!},\ u_* \;:\; \M(I') \;\longrightarrow\; \M(I),
\]
\printedpage{26}
% head: A. Grothendieck : Dérivateurs, Ch. XIII, Édition des Universités de Montpellier II et Paris VII
et je me propose ici de résumer les propriétés les plus immédiates de $u^*$ et de $u_{!}$, $u_*$ par rapport aux structures de Quillen faibles envisagées.

\textbf{Proposition 1 :} \textit{Avec les notations précédentes, on a les relations :}
\[
u^*(C(I)) \subset C(I'), \qquad u^*(TC(I)) \subset TC(I'). \tag*{(1)}
\]
\[
u_*(F(I')) \subset F(I), \qquad u_*(TF(I')) \subset TF(I). \tag*{(2)}
\]
\[
\begin{aligned}
u^*(F(I)) \subset F(I') &\iff u_{!}(TC(I')) \subset TC(I)\\
u^*(TF(I)) \subset TF(I') &\iff u_{!}(C(I')) \subset C(I).
\end{aligned} \tag*{(3)}
\]

D'autre part, je ne sais rien dire de positif pour les inclusions
\[
\begin{aligned}
u_*(C(I')) \;&\subset\; C(I), \;\; u_*(TC(I')) \;\subset\; TC(I)\\
u_{!}(F(I')) \;&\subset\; F(I), \;\; u_{!}(TF(I')) \;\subset\; TF(I).
\end{aligned}
\]
Si ce n'est que pour un $u$ tant soit peu général ($u$ pas un isomorphisme, disons), $u_*$ n'a aucune envie

\noindent\msp{page 46}

\noindent d'être compatible avec les structures $C$, $TC$ (propriété d'exactitude homotopique à droite, moralement, vu que $u_*$ est exact à gauche et en général pas du tout à droite), ni $u_{!}$ avec les structures $F$, $TF$ (commentaire dual à donner, $u_{!}$ est exacte à \textit{droite}, et n'a pas envie d'être “homotopiquement exact à gauche”). % [sic] exacte -> exact, accord avec « foncteur » / « exact à gauche »

\textsc{Démonstration} de (1), (2), (3). Les relations (1) sont évidentes par définition. Les relations (2), (3) le sont à peine moins, vu la propriété soritale générale suivante, dont la place serait au paragraphe 1 :

\textbf{Proposition 2.} \textit{Soit}
\[
\begin{tikzcd}
C' \arrow[r, shift left, "u_{!}"] & C \arrow[l, shift left, "u^*"]
\end{tikzcd}
\]
\textit{un couple de foncteurs adjoints $(u_{!},u^*)$, et soit $\Phi' \subset \Fl(C')$, $\Psi \subset \Fl(C)$. Alors}
\[
(u_{!}(\Phi') \longleftrightarrow \Psi) \qquad\quad \iff \qquad\quad (\Phi' \longleftrightarrow u^*(\Psi)).
\]

\textbf{Corollaire.} (C'est le même énoncé, avec seulement des notations changées.) \textit{Supposons que $u^* : C \longrightarrow C'$ admette un adjoint à} droite \textit{$u_*$ (pas forcément un adjoint à gauche)}
\[
\begin{tikzcd}
C' \arrow[r, shift right, "u_*"'] & C, \arrow[l, shift right, "u^*"']
\end{tikzcd}
\]

\printedpage{27}
% head: A. Grothendieck : Dérivateurs, Ch. XIII, Édition des Universités de Montpellier II et Paris VII
\noindent\textit{soit $\Phi \subset \Fl(C)$, $\Psi' \subset \Fl(C')$, alors}
\[
(u^*(\Phi) \longleftrightarrow \Psi') \qquad\quad \iff \qquad\quad (\Phi \longleftrightarrow u_*(\Psi')).
\]

C'est la proposition 2 pour \begin{tikzcd} C'_0 \arrow[r, shift left, "u_{0\,!}"] & C_0 \arrow[l, shift left, "u_0^*"] \end{tikzcd}, avec $C'_0 = C$, $C_0 = C'$, $(u_{0\,!} : C \longrightarrow C') = u^*$, et $(u_0^* : C' \longrightarrow C) = u_*$.

\noindent\msp{page 47}

Pour les inclusions (2), on utilise le corollaire de la proposition 2 pour les couples $\Phi = TC(I)$, $\Psi' = F(I')$, resp. $\Phi = C(I)$, $\Psi' = F(I')$ \texttt{[plutôt} $\Psi' = TF(I')$\texttt{]}, et on applique (1) qui donne $u^*(\Phi) \longleftrightarrow \Psi'$, d'où $\Phi \longleftrightarrow u_*(\Psi')$, ce qui n'est autre que (2). Pour les équivalences (3), on utilise la proposition 2, en faisant $\Psi = F(I)$, $\Phi' = TC(I')$ resp. $\Psi = TF(I)$, $\Phi' = C(I')$.

\textbf{Corollaire 1} (de la proposition 1) : \textit{Supposons $u : I' \longrightarrow I$ telle que pour tout $i \in \Ob I$, la catégorie $i\backslash I'$ soit discrète} (\footnote[27]{Plus généralement, il suffit que toute composante connexe $Z_\alpha$ de $i\backslash I'$ ait un objet initial $e_\alpha$, de sorte que $\limind\limits_{i\backslash I'} X'(i') \simeq \coprod_{\alpha\in\pi_0(i\backslash I')} X'(e_\alpha)$. C'est le cas notamment si $I' \longrightarrow I$ est un isomorphisme local, p. ex. une immersion ouverte.}) \textit{(p. ex. $I'$ discrète). Alors on a}
\[
\left\{
\begin{aligned}
u_{!}(C(I')) \;&\subset\; C(I), \;\; u_{!}(TC(I')) \;\subset\; TC(I)\\
u^*(F(I)) \;&\subset\; F(I'), \;\; u^*(TF(I)) \;\subset\; TF(I').
\end{aligned}
\right.
\]

En effet, les relations de la deuxième ligne équivalent (à cause de (3)) à celles de la première. Il faut prouver que si $f$ dans $\Fl\M(I')$ est dans $C(I')$ resp. $TC(I')$, $u_{!}(f)$ est dans $C(I)$ resp. $TC(I)$, ce qui signifie aussi que $\alpha_s^*(u_{!}(f))$ est dans $C$ resp. $TC$, pour tout $s \in \Ob I$. Or le foncteur
\[
\alpha_s^* u_{!} \;:\; \M(I') \;\longrightarrow\; \M
\]
s'explicite par la formule bien connue
\[
\alpha_s^*(u_{!}(X')) \;\overset{\text{déf}}{=}\; u_{!}(X')(s) \;\simeq\; \limind\limits_{i'\in s\backslash I'} X'(i').
\]

Comme

\noindent\msp{page 48}

\noindent $s\backslash I'$ est discrète, on a donc l'isomorphisme fonctoriel en $X'$ dans $\M(I')$
\[
\alpha_s^*(u_{!}(X')) \;\simeq \coprod_{(s',\, x \,:\, s \,\longrightarrow\, u(s')) \text{ dans } \Ob\,(s\backslash I')} X'(s'),
\]
donc si $f \,:\, X' \longrightarrow Y'$, $\alpha_s^*(u_{!}(f))$ s'explicite comme une somme directe dans $\underline{\mathrm{Fl}}(\M)$ de flèches de la forme
\[
f_{s'} \;:\; X'(s') \;\longrightarrow\; Y'(s'),
\]
\printedpage{28}
% head: A. Grothendieck : Dérivateurs, Ch. XIII, Édition des Universités de Montpellier II et Paris VII
lesquelles sont dans $C$ resp. $TC$ puisque $f$ est dans $C(I')$ resp. $TC(I')$. Comme $C$, $TC$ sont stables par sommes directes quelconques ($C$, $TC$ étant clos à gauche), on gagne.

\textbf{Corollaire 2.} \textit{Soit $I$ dans} Cat, \textit{et $f \in \Fl(\M(I))$. Pour que $f$ soit dans $F(I)$} (resp. \textit{dans $TF(I)$\,), il} faut \textit{que $\forall\, s \in \Ob I$, $f_s : X(s) \longrightarrow Y(s)$ le soit.} (Mais cette condition n'a pas du tout l'air d'être suffisante !) \textit{En particulier, soit $X \in \Ob\M(I)$, alors $X \in \Ob\M(I)_f \Longrightarrow X(s) \in \Ob\M_f \ \ \forall\, s \in \Ob I$.}

On applique le corollaire 1 à $\alpha_s : e \longrightarrow I$.

\textbf{Scholie.} La proposition 1 (relation (1)) montre que le foncteur $u^*$ transforme carrés $W$-cocartésiens en itou, suites exactes longues de suspension en itou – c'est “l'exactitude à droite homotopique” de $u^*$. On s'attend à ce qu'il en soit de même pour les carrés cartésiens et les suites exactes

\noindent\msp{page 49}

\noindent longues de cosuspension, mais pour le prouver il faudra un travail, qu'on fera dans la suite.

Les relations (2) montrent que le foncteur $u_*$ transforme carrés $W$-cartésiens en itou, donc suites exactes longues de cosuspension en itou. Je ne m'attends \textit{pas} à ce qu'il en soit de même en général pour les carrés $W$-cocartésiens – \textit{i.e.} que $u_*$ soit “homotopiquement exact à droite”, et non seulement à gauche. Mais il faudrait au moins que je fasse des exemples (avec $I'$, $I$ des catégories finies ordonnées).

Quant à $u_{!}$, la proposition ne dit rien de bien palpable à son sujet. Néanmoins on s'attend à ce que $u_{!}$ transforme carrés $W$-cocartésiens en itou – c'est la moindre des choses ! Ou peut-être je fais erreur, en confondant les foncteurs $u_*$, $u_{!}$ envisagés ici, avec les foncteurs similaires
\[
\begin{tikzcd}
\HO_W(I') \arrow[r, shift left, "u_{!}^W"] \arrow[r, shift right, "u_*^W"'] & \HO_W(I),
\end{tikzcd}
\]
adjoints à gauche resp. à droite de
\[
\begin{tikzcd}
\HO_W(I') & \HO_W(I), \arrow[l, "u_W^*"']
\end{tikzcd}
\]
induit par
\[
u^* \;:\; (\M(I),W(I)) \;\longrightarrow\; (\M(I'),W(I')).
\]
Mais en général $u_*$, $u_{!}$ ne sont pas compatibles

\noindent\msp{page 50}

\noindent avec les localiseurs $W(I')$, $W(I)$, donc ils ne doivent pas être confondus avec les foncteurs $u_{!}^W$, $u_*^W$. Donc j'ai été peut-être hâtif, dans ce que j'ai dit au sujet de $u_*$ et de $u_{!}$. Je serai obligé d'étudier la situation de très près, quand j'essayerai de construire les foncteurs $u_*^W$ et $u_{!}^W$. Pour le moment, je n'ai encore aucune assurance que les conditions envisagées pour $(\M,W)$ suffisent à assurer l'existence de $u_*^W$, $u_{!}^W$ – pas même pour $I$, $I'$ des ensembles ordonnés finis !

\textbf{Corollaire 1.} \textit{Soit $X \in \M(I)_{\mathrm{f}}$, alors pour toute $\alpha : s \longrightarrow t$ dans $I$ qui est} \texttt{[un]} monomorphisme, \textit{$X(\alpha) : X(t) \longrightarrow X(s)$ est dans $F$.}

% foot: Version 25 janvier 2004


\endgroup
% END UNIT Der-ChXIII-25-01-04/Der-ChXIII-25-01-04_p026-030.tex
% BEGIN UNIT Der-ChXIII-25-01-04/Der-ChXIII-25-01-04_p031-035.tex
\begingroup\makeatletter%
\newcommand{\printedpage}[1]{\marginpar{\footnotesize p.~#1}}%
\newcommand{\msp}[1]{\marginpar{\footnotesize [#1]}}%
\newcommand{\circled}[1]{\tikz[baseline=(c.base)]\node[draw,circle,inner sep=1pt](c){#1};}%
\makeatother\let\printedpage\derpp\let\msp\dermsp%
\let\mspin\DERmspin%


\printedpage{29}
% head: A. Grothendieck : Dérivateurs, Ch. XIII, Édition des Universités de Montpellier II et Paris VII
(On applique \texttt{[le]} corollaire 1 (p. 47) à $\tilde{\alpha} : \boldsymbol{\Delta}_1 \longrightarrow I$.)

\textbf{Proposition 3.} (\footnote[28]{Comparer avec le lemme p. 63 et sa démonstration (plus exactement, le corollaire du lemme).}) \textit{Soit $X$ dans $\mathcal{M}(\boldsymbol{\Delta}_1)$, $X = (X_0 \xleftarrow{\alpha} X_1)$. Alors $X$ est dans $\mathcal{M}(\boldsymbol{\Delta}_1)_{\mathrm{f}}$ si et seulement si $X_0$, $X_1$} \texttt{[sont]} \textit{dans $\mathcal{M}_{\mathrm{f}}$ et $\alpha \in F$.}

\textbf{Corollaire 2.} \textit{Soit $I' \overset{u}{\hookrightarrow} I$} \texttt{[une]} immersion ouverte. \textit{Alors pour $X \in \operatorname{Ob}\mathcal{M}(I)_{\mathrm{f}}$, $\varprojlim\limits_{I} X(i) \longrightarrow \varprojlim\limits_{I'} X(i)$ est dans $F$.}

\section*{5\quad Catégories de modèles quillénisables}

Je reviens au cas général d'une catégorie de modèles $(\mathcal{M},W)$.

% marginal vertical bar printed beside the following Définition 1
\textbf{Définition 1.} On dit que $(\mathcal{M},W)$ est \textit{quillénisable} s'il existe une structure de Quillen faible \textit{exacte} $(C,TC,F,TF)$ pour $\mathcal{M}$.

C'est là une condition autoduale, mais que je ne sais vérifier en pratique que moyennant l'exactitude soit à gauche, soit à droite -- donc par une condition plus forte qui détruit la symétrie. C'est pour la sauvegarder dans la mesure du possible que j'introduis la définition précédente, ainsi que la

\textbf{Définition 2.} Soit, en plus de $(\mathcal{M},W)$, Diag $\subset$ Cat une sous-catégorie pleine de Cat. On dit que
\msp{page 51}$W$ est \textit{quillénisable sur} Diag, si pour tout $I$ dans Diag, $(\mathcal{M}(I),W(I))$ est quillénisable.

En pratique, j'utiliserai cette définition dans le cas où Diag $=$ Cat, ou alors, qu'il satisfasse au moins les propriétés de stabilité utiles pour le domaine d'un dérivateur : stable par $\varprojlim$ finies, par sommes finies -- eventuellement par $\varinjlim$ finies, etc (\footnote[29]{Je supposerai que Diag contient tout au moins les catégories provenant d'ensembles ordonnés finis.}). % [sic] eventuellement -> éventuellement, missing accent

\textbf{Exemple 1 :} Celui des paragraphes 3, 4, où on a dégagé des conditions sur $(\mathcal{M},W)$, avec $\mathcal{M}$ une ``grosse'' catégorie et $\mathcal{M}$ accessible, pour que $W$ soit quillénisable sur Cat tout entier. C'est le cas qui me paraît à présent le plus intéressant.

\textbf{Exemple 2 :} Soient $\mathcal{A}$ une catégorie abélienne, et $\mathcal{M} = K^{\bullet}(A)$ la catégorie des complexes de cochaînes à valeurs dans $\mathcal{A}$, $W = W_{\mathcal{A}}$ l'ensemble des quasi-isomorphismes, donc % [sic] K^\bullet(A) -> K^\bullet(\mathcal{A}), plain italic A where calligraphic A is used everywhere else
\[ \mathcal{M}W^{-1} \;=\; K^{\bullet}(\mathcal{A})W_{\mathcal{A}}^{-1} \;\simeq\; D^{\bullet}(\mathcal{A}), \qquad \text{catégorie dérivée de $\mathcal{A}$.} \]
On définit aussi $\mathcal{M}^{+}$, $\mathcal{M}^{-}$, $\mathcal{M}^{\mathrm{b}}$, avec les localiseurs induits $W_{\mathcal{A}}^{+}$, $W_{\mathcal{A}}^{-}$, $W_{\mathcal{A}}^{\mathrm{b}}$, et les catégories de fractions $D^{+}(\mathcal{A})$, $D^{-}(\mathcal{A})$, $D^{\mathrm{b}}(\mathcal{A})$. Ou
\msp{page 52}encore, pour $I$ dans Cat,
\[ \mathcal{M}(I) \;=\; K^{\bullet}(\mathcal{A}(I)) \]
% foot: Version 25 janvier 2004
\printedpage{30}
% head: A. Grothendieck : Dérivateurs, Ch. XIII, Édition des Universités de Montpellier II et Paris VII
\[ W(I) \;=\; \underbrace{W_{\mathcal{A}(I)}}_{\text{quasi-isomorphismes}} \;\subset\; \mathrm{Fl}(K^{\bullet}(\mathcal{A}(I))), \]
et de même pour $\mathcal{M}^{\alpha}(I)$, $W^{\alpha}(I)$ $(\alpha \in \{+,-,\mathrm{b}\})$, du moins \textit{si $I$ est finie.} Il y a aussi les sous-catégories $K^{\geq 0}(\mathcal{A})$ et $K^{\leq 0}(\mathcal{A})$ des complexes de cochaînes resp. de chaînes au sens strict (avec $K^{i} = 0$ si $i \geq 0$, resp. si $i \leq 0$ \texttt{[plutôt si $i < 0$, resp. si $i > 0$]}). On aura pour $I$ quelconque
\[ \left\{ \begin{aligned} K^{\geq 0}(\mathcal{A})(I) &\;\simeq\; K^{\geq 0}(\mathcal{A}(I)) \\ W_{\mathcal{A}}^{\geq 0}(I) &\;=\; W_{\mathcal{A}(I)}^{\geq 0}, \end{aligned} \right. \]
et de même pour $K^{\leq 0}$, $W^{\leq 0}$.

D'après \textsc{Quillen} page 1.2, $(K^{+}(\mathcal{A}),W_{\mathcal{A}}^{+})$ est quillénisable si $\mathcal{A}$ a assez d'injectifs, $(K^{-}(\mathcal{A}),W_{\mathcal{A}}^{-})$ l'est si $\mathcal{A}$ a assez de projectifs. Il ne dit rien pour $(K^{\bullet}(\mathcal{A}),W_{\mathcal{A}})$ sans plus -- pourtant on fait de l'algèbre cohomologique avec $D^{\bullet}(\mathcal{A})$ sans plus, et même sans faire d'hypothèse d'existence d'injectifs ou de projectifs. J'ai un peu oublié ce que \textsc{Verdier} a fait à ce sujet -- je devrais quand même regarder sa thèse (!), ou bien l'exposé de \textsc{Hartshorne}.

\msp{page 53}Toujours est-il qu'on a là des exemples de type tout différent de ceux envisagés dans 1°. Ainsi, il n'y a pas à supposer que $\mathcal{M}$ soit stable par petites $\varinjlim$ -- $\mathcal{M}$ peut même être une petite catégorie, elle l'est si et seulement si $\mathcal{A}$ elle même l'est. Si $\mathcal{A}$ est stable par (petites) $\varinjlim$ filtrantes, et est accessible (p. ex. a une petite famille génératrice, et les $\varinjlim$ filtrantes sont exactes), alors $K^{\bullet}(\mathcal{A})$ etc. (\footnote[30]{!!}) ont les mêmes propriétés, et $(K^{+}(\mathcal{A}),W_{\mathcal{A}}^{+})$ est quillénisable sur Cat tout entier (\footnote[31]{attention, $K^{+}(\mathcal{A})$ n'est \textit{pas} stable par $\varinjlim$ filtrantes infinies -- donc on ne peut appliquer le tapis des paragraphes 2, 3 directement à $K^{+}$.}). (C'est le cas où on sait que $\mathcal{A}$ a assez d'injectifs, et l'algèbre cohomologique se fait surtout à coups de résolutions injectives \dots) Dans ce cas on sait que les % [sic] elle même -> elle-même, missing hyphen
\[ u_{*} \;:\; D^{+}(I',\mathcal{A}) \longrightarrow D^{+}(I,\mathcal{A}) \]
existent comme adjoints à droite des
\[ u^{*} \;:\; D^{+}(I,\mathcal{A}) \longrightarrow D^{+}(I',\mathcal{A}), \]
mais je ne suis plus trop sûr si on sait définir $u_{*}$ aussi sur $D(I',\mathcal{A})$ tout entier, et je doute que $u_{!}$ existe sans hypothèse sur $u$ (genre : dimension cohomologique finie). Sans doute on doit toujours avoir
\[ u_{!} \;:\; D^{-}(I',\mathcal{A}) \longrightarrow D^{-}(I,\mathcal{A}), \]
\msp{page 54}adjoint à gauche de
\[ u^{*} \;:\; D^{-}(I,\mathcal{A}) \longrightarrow D^{-}(I',\mathcal{A}). \]
% foot: Version 25 janvier 2004

\printedpage{31}
% head: A. Grothendieck : Dérivateurs, Ch. XIII, Édition des Universités de Montpellier II et Paris VII
Mais je ne sais plus trop s'il faut s'attendre à pouvoir définir $u_{!}$ sur $D(I',\mathcal{A})$ tout entier, en dehors d'hypothèses genre dimension cohomologique finie. La triste vérité, c'est que j'ai $\pm$ oublié la théorie des foncteurs dérivés ! (Il faudrait que je demande encore à \textsc{Jean Malgoire} de m'envoyer des textes canoniques.)

\textbf{Scholie.} Si $(\mathcal{M},W)$ est quillénisable sur Diag $\subset$ Cat (et même sans qu'il soit quillénisable), il définit un prédérivateur $\mathbf{D}_{\mathcal{M},W}$ sur Diag. On posera parfois
\[ \mathrm{HO}_{\mathcal{M},W}(I) \;\text{ ou }\; \mathrm{HO}_{W}(I) \;=\; \mathbf{D}_{\mathcal{M},W}(I). \]
L'hypothèse quillénisable implique, sur chacune de ces catégories homotopiques, la structure (relativement faible) ``triangulée'', provenant des suites \textit{exactes} de suspension et de cosuspension (\footnote[32]{Dans le cas où $\mathcal{M}$ est ponctuée par \texttt{[un]} objet neutre. Sinon, il y a la structure associée aux carrés cartésiens et cocartésiens. Il serait assez intéressant d'axiomatiser une telle espèce de structure.}). (Le caractère quillénisable assure l'exactitude de ces suites \dots) Les foncteurs $u^{*} : \mathrm{HO}_{W}(I) \longrightarrow \mathrm{HO}_{W}(I')$
\msp{page 55}sont ``homotopiquement exacts'' dans le sens des ces structures ``triangulées'' (non commutatives), \textit{i.e.} transforment suites exactes de suspension resp. de cosuspension en itou. % [sic] des ces -> de ces, grammar

Quels sont les axiomes des dérivateurs qui sont vérifiés sous la seule hypothèse quillénisable ? \textbf{Der} 1 \texttt{[voir chapitre I, p. 43]} l'est sans hypothèse aucune sur la catégorie de modèles $(\mathcal{M},W)$ :
\[ \mathrm{HO}_{W}(\coprod_{\alpha} I_{\alpha}) \;\simeq\; \prod_{\alpha} \mathrm{HO}_{W}(I_{\alpha}) \]
pour une famille \textit{finie} d'objets $I_{\alpha}$ dans Diag. L'axiome de localisation \textbf{Der} 2 \texttt{[voir chapitre I, p. 43]} l'est sous l'hypothèse que $W$ soit \textit{fortement saturé,} et quillénisable. On utilise le fait que toute flèche de $\mathrm{HO}_{W}(I)$ provient, à isomorphisme près, d'une flèche dans $\mathcal{M}(I)$. (C'était là un des points qui restait en suspens dans VI, dans la théorie de $\mathrm{HOT}_{W}$, $W \subset \mathrm{Fl}(\mathrm{Cat})$ un localiseur fondamental \dots)

L'axiome \textbf{Der} 3, nouvelle mouture \texttt{[voir chapitre I, pp. 111-112]}, implique un localiseur fondamental $\mathbf{W} \subset \mathrm{Fl}(\mathrm{Cat})$ (pas forcément $\mathbf{W}_{\infty}$) (\footnote[33]{Il est naturel ici de supposer que $\mathbf{W}$ satisfait à tous les ``bons axiomes'' sur les localiseurs fondamentaux W (1, 2, 3, 4, 5, 6, 7) \texttt{[voir chapitre XII]}.}). La forme la plus forte serait celle-ci : si
\[ u \;:\; I' \longrightarrow I \]
\msp{page 56}est dans $\mathbf{W}$, on veut que le foncteur $u^{*}$ induise une \textit{équivalence}
\[ u_{\mathrm{lc}}^{*} \;:\; \mathrm{HO}_{W}^{\mathrm{lc}}(I) \xrightarrow{\;\sim\;} \mathrm{HO}_{W}^{\mathrm{lc}}(I'). \]
Mais je doute qu'on soit armé pour vérifier un axiome aussi fort, sans avoir au moins une théorie des opérations $u_{!}$ ou $u_{*}$. On pourrait à tout hasard poser cet axiome en prenant $\mathbf{W} = \mathbf{W}_{\alpha} \overset{\text{déf}}{=}$ le plus petit localiseur fondamental satisfaisant W (1 à 7) (le plus grand de
% foot: Version 25 janvier 2004
\printedpage{32}
% head: A. Grothendieck : Dérivateurs, Ch. XIII, Édition des Universités de Montpellier II et Paris VII
ces localiseurs étant $\mathbf{W}_{\infty}$). Pour l'instant, la façon de vérifier cet axiome (en termes de la donnée $(\mathcal{M},W)$) est pour moi un mystère $\pm$ total. (Mais je sais le vérifier pour $\mathrm{HOT}_{W}$, quand $W \subset \mathrm{Fl}(\mathrm{Cat})$ satisfait W (1, 2, 3, 4, 5, \dots).)

Je vais m'en occuper, une fois que j'aurai une maîtrise raisonnable sur le formalisme de $u_{*}$, $u_{!}$ relatif aux $u : I' \longrightarrow I$ dans Diag. L'hypothèse quillénisable seule semble trop grossière pour assurer l'existence de $u_{!}$ ou de $u_{\star}$ dans des cas intéressants. % [sic] u_\star -> u_*, star symbol where u_* is used elsewhere

\noindent\msp{page 57}

\section*{6\quad Catégorie de modèles quillénisables, et foncteurs $u_{!}$, $u_{*}$}

\textbf{Définition 1 :}

\begin{itemize}[leftmargin=2.5em]
\item[\circled{a}] On dit que $(\mathcal{M},W)$ est \textit{quillénisable uniformément à gauche} sur Diag, si on peut trouver une structure de Quillen faible $(C,TC,F,TF)$ pour $\mathcal{M}$ qui soit \textit{close,} et qui de plus ait la propriété suivante : pour $I$ dans Cat, soient $C(I)$, $TC(I) \subset \mathrm{Fl}(\mathcal{M}(I))$ définis par
\[ \begin{aligned} &(u : X \longrightarrow Y) \in C(I) \text{ (resp. } TC(I)) \text{ si et seulement si } \ \forall\, s \in \operatorname{Ob} I, \\ &u_{s} : X_{s} \longrightarrow Y_{s} \text{ est dans } C(I) \text{ (resp. } TC(I)). \end{aligned} \]
Soit de plus
\[ F(I) \;=\; TC(I)_{*}, \qquad TF(I) \;=\; C(I)_{*}. \]
La condition envisagée, c'est que pour tout $I$ \textit{dans} Diag, $(C(I),TC(I),F(I),TF(I))$ est une structure de Quillen faible \textit{close exacte} pour $W(I)$ (\footnote[34]{$(C,TC,F,TF)$ est appelé alors une \textit{quillénisation uniforme à gauche} de $(\mathcal{M},W)$ sur Diag.}).

\item[\circled{b}] On dit que $(\mathcal{M},W)$ est \textit{quillénisable sur Diag uniformément et exactement} \texttt{[à]} \textit{gauche}, si on a de plus $TC = C \cap W$, \textit{i.e.} si $(C,TC,F,TF)$ est exacte à gauche, donc défini par $C$ seul, via
\[ TC = C \cap W, \qquad F = (TC)_{*}, \qquad TF = C_{*} \qquad \text{(\footnotemark[35])} \]
\footnotetext[35]{On dit que $(C,\ldots)$, ou simplement $C$, définit une \textit{quillénisation uniforme et exacte à gauche} de $(\mathcal{M},W)$ sur Diag.}
\end{itemize}

On a vu au paragraphe 3 que si $\mathcal{M}$ est stable par petites $\varinjlim$ et accessible, et $W$ stable par petites $\varinjlim$ filtrantes, enfin si
% foot: Version 25 janvier 2004
\printedpage{33}
% head: A. Grothendieck : Dérivateurs, Ch. XIII, Édition des Universités de Montpellier II et Paris VII
\msp{page 58}$W$ admet une structure de Quillen faible $(C,TC,F,TF)$ qui soit close et exacte à gauche -- alors les conditions de la définition 1 sont satisfaites.

D'autre part, notons que la proposition 1 page 45 et ses corollaires, y inclus la proposition 3 (page 50) et ses corollaires, sont valables dès lors que les quillénisations des $(\mathcal{M}(I),W(I))$ envisagés sont définies ``uniformément à gauche'' par une quillénisation uniforme à gauche (pas nécessairement exacte à gauche) de $(\mathcal{M},W)$ sur Diag.

La quillénisation uniforme gauche nous servira à établir l'existence des images directes homotopiques $u_{*}$, pour certains $u$. La quillénisation uniforme droite, dualement, aura comme conséquence principale l'existence des $u_{!}$ homotopiques pour certains $u : I' \longrightarrow I$ dans Diag.

Je vais résumer ce que j'ai pu établir jusqu'à présent dans ce sens, dans le

\msp{page 59}\textbf{Théorème-scholie :} \textit{Soit $(C,TC,F,TF)$ une quillénisation uniforme gauche de $W$ sur} Diag $\subset$ Cat.

\begin{itemize}[leftmargin=3em]
\item[1)] \textit{Soit $u : I' \longrightarrow I$ une flèche dans} Diag. \textit{On suppose} Diag \textit{accordé à $\mathcal{M}$ de façon que pour $u$, flèche dans} Diag, \textit{les $\varprojlim$ dans $\mathcal{M}$ (de type $I'/i$) requises pour assurer l'existence de}
\[ u_{*} \;:\; \mathcal{M}(I') \longrightarrow \mathcal{M}(I) \]
\textit{existent dans $\mathcal{M}$. (C'est O.K. si $I$, $I'$} \texttt{[sont]} \textit{finies, $\mathcal{M}$ étant supposée stable par $\varinjlim$ finies de toutes façons.) On sait} (prop. 1, page 45) \textit{que $u_{*}$ induit}
\[ u_{*}^{\mathrm{f}} \;:\; \mathcal{M}(I')_{\mathrm{f}} \longrightarrow \mathcal{M}(I)_{\mathrm{f}}\,. \tag*{$(*)$} \]
\textit{On dit que $u$ est} admissible (à gauche) pour la quillénisation uniforme gauche donnée de $W$ sur Diag, \textit{si $u_{*}^{\mathrm{f}}$ est compatible avec les localiseurs $W(I')_{\mathrm{f}}$, $W(I)_{\mathrm{f}}$,} i.e. \textit{si $f' : X' \longrightarrow Y'$ dans $\mathcal{M}(I')$, $X'$, $Y'$ fibrants, $f' \in W(I')$ implique $u_{*}(f') \in W(I)$. On dit que $Z$ dans} Diag \textit{est} admissible (pour la quillénisation uniforme gauche donnée de $W$ sur Diag \dots) \textit{si le foncteur $Z \longrightarrow e$ est admissible. Ceci posé :}
\begin{itemize}[leftmargin=2.5em]
\item[a)] \textit{Pour que $u$ soit admissible, il suffit que les $I'/s$ pour $s \in \operatorname{Ob} I$, le soient. NB il y a une réciproque, quand $I$ est ordonné, donc les $I'/s \longrightarrow I'$ des immersions ouvertes.}
\item[b)] \msp{page 60}\textit{Tout ensemble ordonné fini est admissible. Plus généralement (sauf erreur) tout ensemble ordonné $I$ tel que l'espace topologique associé à $I^{\circ}$ soit noethérien (}i.e. \textit{toute suite croissante d'ouverts de $I^{\circ}$ est stationnaire, ou encore toute suite croissante de fermées de $I$ est stationnaire), et si de plus toute partie fermée irréductible de $I^{\circ}$ admet un point générique,} i.e. \textit{$I^{\circ}$ est (non seulement un espace noethérien, mais aussi) un espace topologique} sobre. \textit{(Pour $I$, cela signifie que toute partie ouverte filtrante croissante de $I$ admet un plus grand élément.)} % [sic] fermées -> fermés, agreement with masculine (ouverts / fermés)
\end{itemize}
\end{itemize}
% foot: Version 25 janvier 2004


\endgroup
% END UNIT Der-ChXIII-25-01-04/Der-ChXIII-25-01-04_p031-035.tex
% BEGIN UNIT Der-ChXIII-25-01-04/Der-ChXIII-25-01-04_p036-040.tex
\begingroup\makeatletter%
\newcommand{\printedpage}[1]{\marginpar{\footnotesize p.~#1}}%
\newcommand{\msp}[1]{\marginpar{\footnotesize [#1]}}%
\newcommand{\cM}{\mathcal{M}}%
\newcommand{\HO}{\mathrm{HO}}%
\newcommand{\Hom}{\mathrm{Hom}}%
\newcommand{\Ob}{\mathrm{Ob}}%
\newcommand{\rf}{\mathrm{f}}%
\newcommand{\rc}{\mathrm{c}}%
\makeatother\let\printedpage\derpp\let\msp\dermsp%
\let\mspin\DERmspin%


\printedpage{34}
% head: A. Grothendieck : Dérivateurs, Ch. XIII, Édition des Universités de Montpellier II et Paris VII

\begin{itemize}[leftmargin=3em]
\item[c)] \textit{En particulier, si $I'$, $I$ sont des ensembles ordonnés finis, alors tout $u : I' \longrightarrow I$ est admissible. Il en est de même, plus généralement, si $I$, $I'$ sont des ensembles ordonnés et $I'^{\circ}$ est un espace noethérien sobre.}
\end{itemize}

\begin{itemize}[leftmargin=2em]
\item[2)] \textit{Supposons $u : I' \longrightarrow I$ admissible. Considérons}
\[
u_*^W \ : \ \underbrace{\cM(I')_{\rf}W(I')_{\rf}^{-1}}_{\simeq\, \HO_W(I')} \ \longrightarrow \ \underbrace{\cM(I)_{\rf}W(I)_{\rf}^{-1}}_{\HO_W(I)},
\]
\textit{induit par $u_*$. Ce foncteur est adjoint à droite de $u_W^* : \HO_W(I) \longrightarrow \HO_W(I')$.}

\msp{page 61}
\textit{De plus, l'homomorphisme canonique (pour $s \in \Ob I$, $X' \in \Ob \HO_W(I')$)}
\[
u_*^W(X')(s) \ \simeq \ p_{s\,*}(X'|I'/s) \ \overset{\text{déf}}{=} \ \underset{I'/s}{\underleftarrow{\mathrm{holim}}}{}_W(X'|I'/s)
\]
\textit{(où $p_s : I'/s \longrightarrow e$) est un isomorphisme, si on suppose que les $I'/s$ sont admissibles (on dira alors que $u$ est} strictement \textit{admissible).}
\end{itemize}

\textbf{Corollaire.} \textit{Sur le domaine formé des ensembles ordonnés finis, le prédérivateur} $\HO_W$ \textit{satisfait (en plus des axiomes} \textbf{Der}~1, \textbf{Der}~2\textit{) aux axiomes} \textbf{Der}~4 \textit{et} \textbf{Der}~5 \textit{(existence des images directes $u_*$, adjoints à droite des foncteurs $u^*$, et calcul standard des fibres d'une image directe).}

\textsc{Démonstration} du théorème-scholie.

$1^\circ$) \textit{a}) est immédiat. Pour \textit{b}), on procède par récurrence sur le cardinal $n$ de l'ensemble ordonné fini $I$, les cas $n \in \{0,1\}$ étant triviaux. Supposons donc $n \geq 2$, et le théorème prouvé pour les $n' < n$. Je distingue deux cas : si $I$ admet un plus grand élément $a$, alors, désignant par $\alpha : e \longrightarrow I$ le foncteur de valeur $a$, par $p : I \longrightarrow e$ la projection, on sait que l'on a
\[
p_* \ \simeq \ \alpha^*, \qquad \textit{i.e.} \ \varprojlim_{I} X_i \simeq X_a,
\]
d'où le fait que sur $\cM(I)$ tout entier
\msp{page 62}
(pas seulement sur $\cM(I)_{\rf}$), $p_*$ est compatible aux localiseurs $W(I)$, $W$.

Si $I$ n'a pas de plus grand élément, alors l'ensemble $M$ des ses éléments maximaux est de cardinal $\geq 2$, on peut donc l'écrire $M = M_1 \amalg M_2$, et on aura % [sic] des ses -> de ses, typographical slip
\[
I \ = \ I_1 \cup I_2, \qquad \text{où } I_1 = I/M_1,\ I_2 = I/M_2
\]
sont deux ouverts de $I$, différents de $I$, donc de cardinaux $n_1$, $n_2 \ < \ n$. Considérons aussi $I_0 = I_1 \cap I_2$, c'est un ouvert de cardinal $< n$. Le diagramme
\[
\begin{tikzcd}[row sep=small]
 & I_1 \arrow[dr, dotted] & \\
I_0 \arrow[ur] \arrow[dr] & & I = I_1 \cup I_2 \\
 & I_2 \arrow[ur, dotted] &
\end{tikzcd}
\]
% foot: Version 25 janvier 2004
\printedpage{35}
% head: A. Grothendieck : Dérivateurs, Ch. XIII, Édition des Universités de Montpellier II et Paris VII
correspond à une flèche
\[
\begin{tikzcd}[row sep=small, column sep=small]
 & & & 1 & \\
u \ : \ I \arrow[r] & \Psi = 0 \arrow[urr, start anchor=east, end anchor=south west] \arrow[drr, start anchor=east, end anchor=north west] & & & , \\
 & & & 2 &
\end{tikzcd}
\]
avec
\[
I_1 \ = \ I/1, \qquad I_2 \ = \ I/2, \qquad I_0 \ = \ I/0.
\]
L'hypothèse de récurrence et \textit{a}) impliquent que $u$ est admissible. Comme un composé de foncteurs admissibles est admissible (tautologiquement, via $(vu)_* \simeq v_*u_*$), on est réduit finalement au
\msp{page 63}

\textbf{Lemme.} \textit{La catégorie} \begin{tikzcd}[row sep=small, column sep=small, baseline=(\tikzcdmatrixname-2-1.base)]
 & 1 \\
\Psi = 0 \arrow[ur] \arrow[dr] & \\
 & 2
\end{tikzcd} \quad \textit{est admissible.}

Ceci va résulter du

\textbf{Corollaire} (\footnote[36]{\textbf{NB} Il y a une variante évidente, caractérisant les $X$ tels que $X \longrightarrow e$ soit dans $TF(\Psi)$ comme ceux pour lesquels $X_0$, $X_1$, $X_2$ sont fibrants \textit{triviaux,} et $\alpha_1$, $\alpha_2 \in TF$.}). \textit{Soit} \begin{tikzcd}[row sep=small, column sep=small, baseline=(\tikzcdmatrixname-2-1.base)]
 & X_1 \arrow[dl, "\alpha_1"'] \\
X_\bullet = X_0 & \\
 & X_2 \arrow[ul, "\alpha_2"]
\end{tikzcd} \quad \textit{dans $\cM(\Psi)$. Pour que $X_\bullet$ soit fibrant, il faut et il suffit que l'on ait}

\begin{itemize}[leftmargin=3em]
\item[a)] \textit{$X_0$, $X_1$, $X_2$ sont fibrants, et}
\item[b)] \textit{$\alpha_1$, $\alpha_2 \in F$.}
\end{itemize}

La nécessité de \textit{a}) est cas particulier du corollaire 2, page 48, celle de \textit{b}) est cas particulier du corollaire 1 page 50. Pour la suffisance (dont on n'aura pas besoin pour le théorème), on sait montrer que dans le diagramme commutatif à traits pleins
\[
% Audit resolved: PDF leaf 7 confirms interrupted vertical sides A_0 -> X_0 and A_2 -> X_2, obscured by the displaced B labels; the printed geometry is retained. See corrections.csv.
\begin{tikzcd}
A_1 \arrow[r] \arrow[dr, "TC"] \arrow[dd] & A_0 \arrow[dr, "TC"] \arrow[d, dash] & A_2 \arrow[l] \arrow[dr, "TC"] \arrow[d, dash] & \\
 & B_1 \arrow[r] \arrow[dl, dotted, "h_1"'] \arrow[d] & B_0 \arrow[dl, dotted, "h_0"'] \arrow[d] & B_2 \arrow[l] \arrow[dl, dotted, "h_2"'] \\
X_1 \arrow[r, "\alpha_1"'] & X_0 & X_2 \ , \arrow[l, "\alpha_2"] &
\end{tikzcd}
\]
avec trois flèches marquées $TC$, on peut trouver $h_1$, $h_0$, $h_2$ laissant le diagramme commutatif (trois triangles et deux carrés commutatifs à respecter). Pour ceci, on commence par choisir $h_0$ de façon à rendre le triangle vertical médian commutatif -- c'est possible car $X_0$ est fibrant, $A_0 \longrightarrow B_0$ dans $TC$.

% foot: Version 25 janvier 2004
\printedpage{36}
% head: A. Grothendieck : Dérivateurs, Ch. XIII, Édition des Universités de Montpellier II et Paris VII
\msp{page 64}

Ensuite, on voit séparément pour $h_1$, $h_2$, qu'on peut les trouver en satisfaisant chacune aux deux contraintes de commutativité (1 triangle vertical et un carré penché). Pour $h_1$ p.~ex. l'existence se voit sur le diagramme carré auxiliaire (traits pleins)
\[
\begin{tikzcd}
A_1 \arrow[r, "TC"] \arrow[dd] & B_1 \arrow[dd] \arrow[ddl, dotted, "h_1"'] \arrow[dr, dotted, "\alpha_1^B"] & \\
 & & B_0 \arrow[dl, dotted, "h_0"] \\
X_1 \arrow[r, "F", "\alpha_1^X"'] & X_0 \ \ . &
\end{tikzcd}
\]

\textbf{NB} C'est la démonstration de la proposition 3, page 50.

Le corollaire (partie “il faut”) étant vu, le lemme en résulte, car on sait que si on a un homomorphisme $f$ de carrés cartésiens $X$, $Y$
\[
\begin{tikzcd}[row sep=small, column sep=small]
 & X_1 \arrow[dl, "p_1^X"'] \arrow[dd, "f_1"'] & & X_3 \arrow[ll] \arrow[dl] \arrow[dd, "f_3"] \\
X_0 \arrow[dd, "f_0"'] & & X_2 \arrow[ll, crossing over, "p_2^X"'] & \\
 & Y_1 \arrow[dl, "p_1^Y"'] & & Y_3 \arrow[ll] \arrow[dl] \\
Y_0 & & Y_2 \ \ , \arrow[ll, "p_2^Y"'] \arrow[from=uu, crossing over, "f_2"] &
\end{tikzcd}
\]
avec
\[
p_1^X,\ p_2^X,\ p_1^Y,\ p_2^Y \ \ \text{fibrants \texttt{[appartenant à} $F$\texttt{]} (il suffit même $p_1^X$ et $p_1^Y$ fibrant),}
\]
et
\[
f_0,\ f_1,\ f_2 \ \in \ W,
\]
alors $f_3 \in W$. (Ça a été vu dans la théorie des produits fibrés $W$-homotopiques \texttt{[voir chapitre XII, p. 253]}, et la vérification
\msp{page 65}
directe est immédiate.) De l'hypothèse initiale pertinente que \begin{tikzcd}[row sep=small, column sep=small, baseline=(\tikzcdmatrixname-2-1.base)]
 & X_1 \arrow[dl] \\
X_0 & \\
 & X_2 \arrow[ul]
\end{tikzcd} et \begin{tikzcd}[row sep=small, column sep=small, baseline=(\tikzcdmatrixname-2-1.base)]
 & Y_1 \arrow[dl] \\
Y_0 & \\
 & Y_2 \arrow[ul]
\end{tikzcd}
sont fibrants, on n'a besoin d'utiliser qu'un petit bout, p.~ex. que $p_1^X$ et $p_1^Y$ sont fibrants, qui assure que les deux carrés $X$, $Y$ sont $W$-homotopiquement cartésiens, ce qui suffit à impliquer alors que $f_0$, $f_1$, $f_2 \in W \Longrightarrow f_3 \in W$.

Je vais essayer de généraliser ce raisonnement, pour traiter le cas d'un ensemble ordonné $I$ qui peut être infini, mais tel que $I^\circ$ soit noethérien et sobre, donc toute suite croissante de parties fermées de $I$, et toute suite décroissante de parties ouvertes, est stationnaire, et de plus toute partie ouverte filtrante croissante a un plus grand élément.

% foot: Version 25 janvier 2004
\printedpage{37}
% head: A. Grothendieck : Dérivateurs, Ch. XIII, Édition des Universités de Montpellier II et Paris VII

% margin: a vertical bar is printed in the left margin beside the lines "prouver ceci : (*) Si I est tel que ... alors I aussi."
On veut prouver que $I$ est admissible, et par “récurrence noethérienne” on est ramené à prouver ceci : $(*)$ Si $I$ est tel que pour tout ouvert $I' \subset I$, avec $I' \neq I$, $I'$ soit admissible, alors $I$ aussi.

$(*)$ admis, prouvons que $I = I_0$ est admissible.

\msp{page 66}
S'il ne l'était, par $(*)$ il existerait $I_1 \subset I_0$, $I_1 \neq I_0$, tel que $I_1$ non admissible. De même, il y aurait $I_2 \subset I_1$, $I_2 \neq I_1$, non admissible etc. -- on aurait construit une suite strictement décroissante d'ouverts $I_n$ de $I$, contrairement au caractère noethérien de $I^\circ$.

Prouvons $(*)$, et distinguons encore deux cas.

\textit{a}) $I^\circ$ est irréductible, donc admet un point générique, \textit{i.e.} un plus petit élément, \textit{i.e.} $I$ admet un plus grand élément. Ce cas a déjà été traité tantôt -- toute catégorie avec objet final est admissible.

\textit{b}) $I^\circ$ est réductible, donc réunion de deux fermés $I_1^\circ$, $I_2^\circ$ distincts de $I^\circ$, donc $I$ est réunion de deux ouverts $I_1$, $I_2$ distincts de $I$. Alors par l'hypothèse dans $(*)$, $I_1$, $I_2$ et $I_0 = I_1 \cap I_2$ sont admissibles. On termine alors comme tantôt, étant ramené à l'admissibilité de \begin{tikzcd}[row sep=small, column sep=small, baseline=(\tikzcdmatrixname-2-1.base)]
 & 1 & \\
\Psi = 0 \arrow[ur] \arrow[dr] & & ! \\
 & 2 &
\end{tikzcd}

\msp{page 67}
Prouvons la partie \textit{c}) du théorème-scholie, partie $1^\circ$). Cela résulte de \textit{a}), une fois vu que pour $s \in I$, $I'/s$ satisfait aux mêmes hypothèses que $I'$. Or $I$ étant ordonné, $I/s$ est une sous-catégorie ouverte de $I$, donc $I'/s$ une sous-catégorie ouverte de $I'$, dont $(I'/s)^\circ$ un sous-ensemble ordonné fermé de $I'^{\circ}$. Ce dernier étant noethérien et sobre, il en est de même de tout fermé, d'où la conclusion.

Prouvons la partie $2^\circ$) du théorème-scholie. Considérons
\[
\begin{tikzcd}[row sep=large]
\cM(I')_{\rc} \arrow[r, hook, dotted] & \cM(I') & \cM(I')_{\rf} \arrow[l, hook', dotted] \arrow[d, "u_*^{\rf}"'] \\
\cM(I)_{\rc} \arrow[r, hook, dotted] \arrow[u, "u_{\rc}^*"] & \cM(I) \arrow[u, "u^*"] & \cM(I)_{\rf} \ , \arrow[l, hook', dotted]
\end{tikzcd}
\]
où $u^*$ et a fortiori $u_{\rc}^*$ sont compatibles avec $W(I)$, $W(I')$, et par hypothèse $u_*$ est compatible avec $W(I')_{\rf}$, $W(I)_{\rf}$. En passant aux catégories de fractions, on trouve

% foot: Version 25 janvier 2004
\printedpage{38}
% head: A. Grothendieck : Dérivateurs, Ch. XIII, Édition des Universités de Montpellier II et Paris VII
\msp{page 68}
\[
\begin{tikzcd}[row sep=large]
\cM(I')_{\rc}W(I')_{\rc}^{-1} \arrow[r, dotted, "\sim"] & \HO_W(I') & \cM(I')_{\rf} \arrow[l, dotted, "\sim"'] \arrow[d, "u_*^W"'] \\
\cM(I)_{\rc}W(I)_{\rc}^{-1} \arrow[r, dotted, "\sim"] \arrow[u, "u_W^*"] & \HO_W(I) & \cM(I)_{\rf} \arrow[l, dotted, "\sim"']
\end{tikzcd}
\]
\texttt{[plutôt} $\cM(I')_{\rf}W(I')_{\rf}^{-1}$ \texttt{(resp.} $\cM(I)_{\rf}W(I)_{\rf}^{-1}$\texttt{) au lieu de} $\cM(I')_{\rf}$ \texttt{(resp.} $\cM(I)_{\rf}$\texttt{)]} et on veut montrer que $u_W^*$ et $u_*^W$ sont adjoints, de façon précise trouver un isomorphisme bifonctoriel, pour $F$ \textit{dans} $\cM(I)_{\rc}$ et $G'$ \textit{dans} $\cM(I')_{\rf}$
\[
\Hom_W(u^*F, G') \xrightarrow{\ \sim\ } \Hom_W(F, u_*G') \tag*{$(*)$}
\]
(où $\Hom_W$ désigne le Hom dans $\HO_W(I')$ resp. $\HO_W(I)$). Comme $F$ est dans $\cM(I)_{\rc}$, $u^*F$ est dans $\cM(I')_{\rc}$, et comme $G'$ est dans $\cM(I')_{\rf}$, $u_*G'$ est dans $\cM(I)_{\rf}$. Il s'ensuit que les deux membres de $(*)$ sont des quotients des deux membres dans la bijection canonique
\[
\Hom_{\cM(I')}(u^*F, G') \xrightarrow[\ \sim\ ]{\alpha} \Hom_{\cM(I)}(F, u_*G'). \tag*{$(**)$}
\]
Il faut simplement montrer que cette dernière transforme la relation d'équivalence pertinente $R'$ du premier membre,
\msp{page 69}
en celle $R$ du second.

Je dis que l'on a de toute façon
\[
R \ \leq \ \alpha(R')
\]
\textit{sans} hypothèse d'admissibilité sur $u$, \textit{i.e.} que la bijection invese de $(**)$ donne par passage aux quotients un homomorphisme % [sic] invese -> inverse, missing letter
\[
\Hom_W(u^*F, G') \longleftarrow \Hom_W(F, u_*G'), \tag*{$(***)$}
\]
forcément surjectif, et de plus fonctoriel en $F$ \texttt{[,} $G'$ \texttt{barré]}, variables dans $\HO_W(I)$ et $\HO_W(I')$ respectivement. En effet, on définit la flèche $(*$ bis$)$ \texttt{[plutôt} $(***)$\texttt{]} comme composé
\[
\begin{tikzcd}[column sep=large]
\Hom_W(F, u_*G') \arrow[r, "\underbrace{\hspace{2.6cm}}_{\text{\footnotesize\begin{tabular}{@{}c@{}}car $u^*$ est\\ compatible avec\\ $W(I)$, $W(I')$ sur\\ $\cM(I)$ tout entier,\\ et on interprète\\ les deux membres\\ comme des Hom\\ dans $\cM(I)W(I)^{-1}$\\ et $\cM(I')W(I')^{-1}$.\end{tabular}}}"'] & \Hom_W(u^*F, u^*u_*G') \arrow[dr, "\text{\footnotesize\begin{tabular}{@{}l@{}}composition\\ avec l'homomorphisme\\ d'adjonction $u^*u_*G' \longrightarrow G'$\end{tabular}}"] & \\
 & & \Hom_W(u^*F, G').
\end{tikzcd}
\]
% foot: Version 25 janvier 2004


\endgroup
% END UNIT Der-ChXIII-25-01-04/Der-ChXIII-25-01-04_p036-040.tex
% BEGIN UNIT Der-ChXIII-25-01-04/Der-ChXIII-25-01-04_p041-045.tex
\begingroup\makeatletter%
\newcommand{\printedpage}[1]{\marginpar{\footnotesize p.~#1}}%
\newcommand{\msp}[1]{\marginpar{\footnotesize [#1]}}%
\newcommand{\M}{\mathcal{M}}%
\newcommand{\Hom}{\mathrm{Hom}}%
\newcommand{\HO}{\mathrm{HO}}%
\newcommand{\Ob}{\mathrm{Ob}\,}%
\newcommand{\Fl}{\mathrm{Fl}\,}%
\newcommand{\Filt}{\mathrm{Filt}}%
\newcommand{\card}{\mathrm{card}\,}%
\newcommand{\rf}{\mathrm{f}}%
\makeatother\let\printedpage\derpp\let\msp\dermsp%
\let\mspin\DERmspin%


\printedpage{39}
% head: A. Grothendieck : Dérivateurs, Ch. XIII, Édition des Universités de Montpellier II et Paris VII
% foot: Version 25 janvier 2004

Cela implique donc que le foncteur sur $\HO_W(I)$ $F \longmapsto \Hom_{\HO_{W}(I')}(u^*F,G')$ est, sinon représentable, du moins quotient terme à terme du foncteur représenté par $u_*G'$ (\footnote[37]{Mais il n'est pas clair, sans hypothèse supplémentaire sur la situation, que ce foncteur représentable puisse être choisi en dépendance fonctorielle de $G'$, envisagé dans $\HO_W(I')$.}).

\msp{page 70}

Mais supposons maintenant que $u$ soit admissible. Alors je vais définir la flèche en sens initial $(*)$ (p. 68) de la façon suivante : Je choisis une flèche
\[
i' \ : \ u^*F \longrightarrow F' \qquad \text{avec } i' \in TC \subset C \cap W,\ F' \in \Ob \M(I')_{\rf}
\]
(donc en fait $F' \in \Ob \M(I')_{\mathrm{cf}}$, car $u^*F$ est cofibrant), et par adjonction la flèche $i'$ correspond à
\[
i \ : \ F \longrightarrow u_*F'.
\]
Je considère le diagramme (commutatif!)
\[
\begin{tikzcd}[column sep=large, row sep=large]
\Hom(F',G') \arrow[r, "\alpha"] \arrow[d, "{\scriptstyle \lambda:f'\,\longmapsto\, f'\circ i'}"] & \Hom(u_*F',u_*G') \arrow[d, "{\scriptstyle \mu:f\,\longmapsto\, f\circ i}"] \\
\Hom(u^*F,G') \arrow[r, "\text{adj.}"] & \Hom(F,u_*G')\ .
\end{tikzcd}
\]

On note que l'on a encore $F'$ cofibrant (car $u^*F$ l'est et $i' \in TC \subset C$), donc $\Hom_W(F',G')$ est un quotient de $\Hom(F',G')$, et $\lambda$ induit un isomorphisme
\[
\Hom_W(F',G') \xrightarrow{\ \sim\ } \Hom_W(u^*F',G') % [sic] u^*F' -> u^*F, the target of lambda is Hom(u^*F,G')
\]
(car $i' \in TC \subset W$). Comme $F'$, $G' \in \Ob \M(I')_{\rf}$ et que $u_*^{\rf}$ est compatible avec $W(I')_{\rf}$, $W(I)_{\rf}$, il s'ensuit que la flèche horizontale $\alpha$
\msp{page 71}
est compatible avec les relations d'équivalence définies par
\[
\Hom(F',G') \longrightarrow \Hom_{W(I')}(F',G') \quad \text{resp.} \ \ \Hom(u_*F',u_*G') \longrightarrow \Hom_{W(I)}(u_*F',u_*G')\ .
\]
Et il en est de même de la flèche $\mu : f \longmapsto f \circ i$, trivialement. Donc la flèche d'adjonction $\Hom(u^*F,G') \longrightarrow \Hom(F,u_*G')$ est compatible avec les relations $R$, $R'$, ce qui achève de prouver que $(**)$, page 68, induit une bijection $(*)$, page 68.

Ceci montre donc

\begin{enumerate}
\item[1°)] que pour un objet fixé $G'$ dans $\HO_W(I')$ (qu'on peut toujours supposer, à isomorphisme près, provenir de $\M(I')_{\rf}$), le foncteur
\[
F \ \longmapsto\ \Hom_{\HO_{W}(I')}(u^*_W F,G')
\]
est représentable -- \textit{i.e.} le foncteur $u^*_W$ admet un adjoint à gauche, et
\item[2°)] que cet adjoint à gauche se calcule comme $u_*(G')$ (quand l'argument $G'$ est pris fibrant) -- vu que l'homomorphisme $({*}{*}{*})$, page 69, est \textit{fonctoriel en $F$}. (A posteriori il s'ensuit qu'il est aussi fonctoriel en $G'$ \dots)
\end{enumerate}

\printedpage{40}
% head: A. Grothendieck : Dérivateurs, Ch. XIII, Édition des Universités de Montpellier II et Paris VII
% foot: Version 25 janvier 2004

\msp{page 72}

Le dernier énoncé dans 2° du théorème-scholie résulte du calcul de $u_*^W$, compte tenu que $I'/s \longrightarrow I'$ est un isomorphisme local, donc le foncteur image inverse envoie
\[
\M(I')_{\rf} \longrightarrow \M(I'/s)_{\rf},
\]
donc si $G'$ est dans $\M(I')_{\rf}$,
\[
u_*(G')(s) \ \simeq\ p_{s\,*}(G'|I'/s)
\]
est un calcul de $p^W_{s\,*}(G'|I'/s)$.

Cela achève la démonstration du théorème-scholie sur les images directes.

\textbf{Commentaires.}

1) Ce théorème me paraît loin d'être satisfaisant, c'est un pis-aller. J'avais espéré, avec les hypothèses du paragraphe 3, qu'on aurait l'existence des $u_*^W$ pour \textit{toute} flèche $u : I' \longrightarrow I$ dans Cat. Je doute à présent que ce soit le cas, mais n'ai pas encore réussi de faire un contre-exemple. Les cas test les plus importants pour la question d'admissibilité d'une catégorie $I$ me paraissent les suivants :

\begin{enumerate}
\item[\textit{a})] $I$ catégorie discrète infinie (disons dénombrable). La question revient ici à celle-ci : soient $f_i : X_i \longrightarrow Y_i$ dans $W$, avec $X_i$, $Y_i$ fibrants, \textit{i.e.} dans $F$ \texttt{[\textit{i.e.}} $X_i \longrightarrow e$\texttt{,} $Y_i \longrightarrow e$ \texttt{dans} $F$\texttt{]}. Alors $f = \prod f_i : \prod X_i \longrightarrow \prod Y_i$ est-il dans $W$ ?

\msp{page 73}

\item[\textit{b})] I catégorie % [sic] I (upright roman) -> $I$, italic everywhere else for the category
\[
0 \longrightarrow 1 \longrightarrow 2 \longrightarrow 3 \longrightarrow \cdots,
\]
\textit{i.e.} la catégorie ordonnée associée à l'ensemble ordonné $\mathbf{N}$, pour la relation d'ordre ordinaire $m \leq n$. Donc $\M(I)$ est formée des systèmes \textit{projectifs} dans $\M$, indexés par $\mathbf{N}$. On prouve facilement (en procédant comme dans le corollaire du lemme p. 63) que $\M(I)_{\rf}$ est formée des systèmes projectifs
\[
X_0 \longleftarrow X_1 \longleftarrow X_2 \longleftarrow \cdots \longleftarrow X_n \xleftarrow{\ p_n\ } X_{n+1} \longleftarrow \cdots
\]
tels que les $X_i$ soient dans $\M_{\rf}$, et les $p_n$ dans $F$. La question est donc : si $X_\bullet$, $Y_\bullet$ sont fibrants et $f_\bullet : X_\bullet \longrightarrow Y_\bullet$ est telle que $f_i \in W$ pour $i \in \mathbf{N}$, alors a-t-on $\varprojlim f_i : X_\infty \longrightarrow Y_\infty$ dans $W$ ?

Si la reponse est affirmative pour $\mathbf{N}$ (pour un $W$ et une quillénisation donnée), elle % [sic] reponse -> réponse, accent missing
l'est aussi pour toute catégorie discrète dénombrable, comme on le voit aussitôt. Mais je ne vois pas comment déduire même \textit{a}) des axiomes, disons pour un triplet de Quillen clos (closed model category) tel que $\M$ soit de plus stable par $\varinjlim$ filtrantes, que celles-ci soient exactes, et qu'il y ait
\printedpage{41}
% head: A. Grothendieck : Dérivateurs, Ch. XIII, Édition des Universités de Montpellier II et Paris VII
% foot: Version 25 janvier 2004
\msp{page 74}
un petit ensemble de générateurs dans $\M$ (donc $\M$ est accessible), que $W$ soit stable par $\varinjlim$ filtrantes. Je devrais demander à \textsc{Thomason} s'il connait un contre-exemple. % [sic] connait -> connaît, circumflex missing

\item[\textit{c})] Les deux exemples précédents étaient des catégories ordonnées. Comme exemple typique de catégorie non-ordonnée, où je ne vois pas ce qui se passe, je signale $I = B_G$, $G$ un groupe (donc un seul objet, avec un groupe d'automorphismes $G$). Aussi le cas de $\boldsymbol{\Delta}[1] = \{\boldsymbol{\Delta}_0, \boldsymbol{\Delta}_1\}$ (sous-catégorie pleine de $\boldsymbol{\Delta}$) n'est pas clair pour moi.
\end{enumerate}

2) On aimerait pouvoir traiter de façon duale l'existence des $u_!^W$, à défaut de mieux. Dans l'hypothèse de départ du paragraphe 4 sur $(\M,W)$, cela amène donc à changer de tactique. Cette fois, on devra supposer l'existence d'une quillénisation close \textit{exacte à droite} pour $W$ (ce qui est plus exigeant, semble-t-il, qu'une telle structure exacte à gauche, dans le cas envisagé), donc
\[
TF \ = \ F \cap C,
\]
\msp{page 75}
et de plus encore $C$, $TC$ tels qu'on ait
\[
C \ = \ \widetilde{C_0}, \qquad TC \ = \ \widetilde{TC_0}
\]
avec $C_0$, $TC_0$ petits (\footnote[38]{On dit alors que $C$, $TC$ (en tant que parties closes à droite \texttt{[plutôt à gauche]}) sont “accessibles”.}). On posera maintenant
\[
\begin{aligned}
F(I) \ &= \ \{f \in \Fl \M(I) \ \mid \ f_s \in F \quad\ \ \forall\, s \in \Ob I\}\\
TF(I) \ &= \ \{f \in \Fl \M(I) \ \mid \ f_s \in TF \ \ \forall\, s \in \Ob I\},
\end{aligned}
\]
puis
\[
\begin{aligned}
C(I) \ &= \ TF(I)^*\\
TC(I) \ &= \ F(I)^*\ .
\end{aligned}
\]
L'espoir (que je devrai vérifier ultérieurement), c'est qu'on obtient ainsi une quillénisation uniforme à droite de $W$, sur Cat tout entier. Le seul problème, c'est la validité des axiomes de factorisation de $f \in \Fl \M(I)$ en $f = pi$, avec $i \in TC$, $f \in F$, et aussi avec $i \in C$, $f \in TF$. % [sic] f \in F, f \in TF -> p \in F, p \in TF, the factorisation is f = pi
J'ai bon espoir que c'est vrai -- cela se réduit essentiellement à voir que $(C(I),TF(I))$ et $(TC(I),F(I))$, soit $(\Phi(I),\Psi(I))$, satisfont en plus de $\Phi(I) = \Psi(I)^*$, à $\Psi(I) = \Phi(I)_*$ (\textit{i.e.} $\Psi(I)$ est clos à droite), et que $\Phi(I)$ est accessible (tout comme $\Phi = \Phi(e)$). Si ça marche, on aura donc l'existence
\printedpage{42}
% head: A. Grothendieck : Dérivateurs, Ch. XIII, Édition des Universités de Montpellier II et Paris VII
% foot: Version 25 janvier 2004
\msp{page 76}
de $u_!$ pour $u$ un foncteur entre catégories ordonnées finies, et sous des conditions légèrement plus générales. (P. ex. $u : I' \longrightarrow I$ avec $I$ ordonnée, $I'$ ordonnée et à espace sous-jacent noethérien et sobre.) Peut-être que cette fois, l'hypothèse que $W$ est stable par $\varinjlim$ filtrantes va quand même servir! Ainsi, le cas des catégories discrètes, et des catégories filtrantes à gauche comme $\mathbf{N}^\circ$, ne va plus poser de problèmes à cause de cela. Il y a donc de l'espoir! Tout ensemble ordonné qui est $\varinjlim$ filtrante de sous-ensembles ordonnés \textit{fermés} noethériens et sobres, est alors $W$-admissible à droite. Mais qu'en est-il de catégories de type $\Gamma = B_G$ p. ex. ? ? ?

\msp{page 77}

\section*{7\quad Filtrations cardinales et parties accessibles d'une grosse catégorie}

Je me réfère à la théorie des foncteurs et catégories accessibles de SGA 4, I, §9 (p. 138-178). Je me donne une grosse catégorie (toujours une $\mathfrak{U}$-catégorie) $E$ avec une \textit{filtration cardinale} (croissante)
\[
\Bigl(\Filt^{\pi}(E)\Bigr)_{\pi \geq \pi_0}, \qquad \pi_0 \text{ un cardinal infini fixé,}
\]
par les sous-catégories strictement pleines $\Filt^{\pi}(E)$. On a \texttt{[SGA 4, I, 9.12]}

\begin{enumerate}
\item[\textit{a})] $\Filt^{\pi}(E)$ essentiellement petite, \textit{i.e.} équivalente à une petite catégorie (\textit{i.e.} à une catégorie appartenant à $\mathfrak{U}$).
\item[\textit{b})] $E$ satisfait $L_{\pi_0}$ (\textit{i.e.} $E$ est stable par $\varinjlim$ de type $I$, quand $I$ est un ensemble ordonné grand devant $\pi_0$), et $\Filt^{\pi}(E)$ est stable par les limites inductives filtrantes indexées par $I$ grand devant $\pi_0$, et tel que $\card I \leq \pi$.
\item[\textit{c})] Pour tout $\pi \geq \pi_0$, et pour tout $X \in \Ob E$, on a
\[
X \ \simeq\ \varinjlim_{i \in I} X_i,
\]
où les $X_i$ \texttt{[sont]} dans $\Filt^{\pi}(E)$, et où $I$ est un ensemble ordonné grand devant $\pi$. De plus, si $X \in \Ob \Filt^{\pi'}(E)$, avec $\pi' > \pi$, on peut prendre $I$ tel que $\card I \leq \pi'^{\pi}$ (donc si $\pi' = 2^{c'}$ avec $c' \geq \pi$, d'où $(\pi')^{\pi} = (2^{c'})^{\pi} = 2^{c'\pi} = 2^{c'} = \pi'$, on aura $\card I \leq \pi'$). (\textbf{NB} Si on admet l'hypothèse du continu généralisée, on aura toujours $\pi'^{\pi} = \pi'$, donc $\card I \leq \pi'$.) (\footnote[39]{NB \textit{b}) et \textit{c}) impliquent que $E = \bigcup_{\pi \geq \pi_0} \Filt^{\pi}(E)$.})
\end{enumerate}

\printedpage{43}
% head: A. Grothendieck : Dérivateurs, Ch. XIII, Édition des Universités de Montpellier II et Paris VII
% foot: Version 25 janvier 2004

\msp{page 78 barrée}

\msp{page 79}

Considérons la représentation de \textit{c}),
\[
X \ = \ \varinjlim_{I} X_i \qquad\qquad
\left\{
\begin{aligned}
& X_i \in \Filt^{\pi},\\
& I \text{ grand devant } \pi.
\end{aligned}
\right.
\]
Soit
\[
\tilde{I} \ = \ \{I' \subset I \mid \card I' \leq \pi,\ \ I' \text{ grand devant } \pi_0\},
\]
muni de la relation d'ordre d'inclusion. Pour
\[
I' \ \in\ \tilde{I}, \qquad \text{soit}\ \ X_{I'} = \varinjlim_{I'} X_i\ ,
\]
cette limite existe par $L_{\pi_0}$ (hypothèse \textit{b}), elle est dans $\Filt^{\pi}$ par \textit{b}), donc on a un système inductif
\[
(X_{I'})_{I' \in \tilde{I}}, \qquad X_{I'} \in \Filt^{\pi},
\]
avec
\[
\card \tilde{I} \ \leq\ (\card I)^{\pi},
\]
et si $X \in \Filt^{\pi'}$ $(\pi' > \pi)$, prenant
\[
\card I \ \leq\ \pi'^{\pi},
\]
on aura $\card \tilde{I} \leq (\pi'^{\pi})^{\pi} = \pi'^{\pi^2} = \pi'^{\pi}$, donc
\[
\card \tilde{I} \ \leq\ \pi'^{\pi}.
\]
Je dis que $\tilde{I}$ est lui aussi \textit{grand devant $\pi$,} tout comme $I$. En effet,
\msp{page 80}
soit
\[
(I_\alpha)_{\alpha \in A}, \qquad I_\alpha \in \tilde{I}, \qquad \card A \ \leq\ \pi,
\]
prouvons qu'il existe $I' \in \tilde{I}$ qui majore les $I_\alpha$, \textit{i.e.} qui contient $I'_0 = \bigcup_{\alpha \in A} I_\alpha$. On aura $\card I'_0 \leq \pi^2 = \pi$ (car $\card A \leq \pi$ et $\card I_\alpha \leq \pi$ pour tout $\alpha$), et il suffit de voir ceci :

\texttt{[En fait, l'assertion est évidente : comme} $\card I'_0 \leq \pi$\texttt{, et comme} $I$ \texttt{est grand devant} $\pi$\texttt{,} $I'_0$ \texttt{admet un majorant} $i'_0$ \texttt{dans} $I$\texttt{, et on peut prendre} $I' = I'_0 \cup \{i'_0\}$\texttt{.]}

\textbf{Lemme 1 :} \textit{Soient $\pi_0$, $\pi$ deux cardinaux, avec $\pi \geq \pi_0$, $I$ un ensemble ordonné grand devant $\pi$, et $I'_0$ une partie de $I$ de cardinal $\leq \pi$. Alors il existe une partie $I'$ de $I$ satisfaisant les conditions :} (\footnote[40]{Pourvu qu'on suppose $\pi > \pi_0$, et de plus $\pi^{\pi_0} = \pi$ (p. ex. $\pi = 2^c$ avec $c \geq \pi_0$).})

\begin{enumerate}
\item[a)] $I' \supset I'_0$.
\item[b)] $\card I' \leq \pi$.
\end{enumerate}


\endgroup
% END UNIT Der-ChXIII-25-01-04/Der-ChXIII-25-01-04_p041-045.tex
% BEGIN UNIT Der-ChXIII-25-01-04/Der-ChXIII-25-01-04_p046-050.tex
\begingroup\makeatletter%
\newcommand{\printedpage}[1]{\marginpar{\footnotesize p.~#1}}%
\newcommand{\msp}[1]{\marginpar{\footnotesize [#1]}}%
\newcommand{\card}{\operatorname{card}}%
\newcommand{\Filt}{\mathrm{Filt}}%
\newcommand{\Ob}{\operatorname{Ob}}%
\newcommand{\circled}[1]{\tikz[baseline=(c.base)]\node[draw,circle,inner sep=1pt](c){#1};}%
\makeatother\let\printedpage\derpp\let\msp\dermsp%
\let\mspin\DERmspin%


\printedpage{44}
% head: A. Grothendieck : Dérivateurs, Ch. XIII, Édition des Universités de Montpellier II et Paris VII

\begin{itemize}
\item[c)] \textit{$I'$ grand devant $\pi_0$.}
\end{itemize}

\texttt{[Le lemme tel qu'il est énoncé est évident pour les mêmes raisons que ci-dessus (même sans l'hypothèse de la note de bas de page). La démonstration qui suit prouve en fait un résultat plus fin :}

\texttt{\textsc{Lemme} 1$'$ : \textit{Soient $\pi_0$, $\pi$ deux cardinaux, avec $\pi_0<\pi$ et $\pi^{\pi_0}=\pi$, $I$ un ensemble ordonné grand devant $\pi_0$, et $I'_0$ une partie de $I$ de cardinal $\le\pi$. Alors il existe une partie $I'$ de $I$ satisfaisant les conditions :}}

\begin{itemize}
\item[\texttt{a)}] $I'\supset I'_0$\texttt{.}
\item[\texttt{b)}] $\card I'\le\pi$\texttt{.}
\item[\texttt{c)}] \texttt{\textit{$I'$ grand devant $\pi_0$. ]}}
\end{itemize}

On va essayer de construire $I'$ par récurrence transfinie, en définissant les $I'_\alpha\subset I$ ainsi :
\[
\begin{aligned}
I'_0 &= I'_0\;, \\
I'_\alpha &= \textstyle\bigcup_{\beta<\alpha} I'_\beta \quad \text{si $\alpha$ est un ordinal limite,} \\
I'_{\alpha+1} &= \left\{\parbox{0.6\textwidth}{déduit de $I'_\alpha$ ainsi : pour toute partie $J$ de $I'_\alpha$ telle que $\card J\le\pi_0$, qui n'est pas majorée dans $I'_\alpha$, on choisit un majorant $i_J$ dans $I$ (\texttt{[qui]} existe, car $I$ \texttt{[est]} grand devant $\pi$, donc devant $\pi_0$) et $I'_{\alpha+1}$ s'obtient en adjoignant ces $i_J$ à $I'_\alpha$.}\right.
\end{aligned}
\]

\msp{page 81}
Soit alors $\alpha_0$ le premier ordinal qui soit de cardinal $>\pi_0$, et considérons
\[
I' \;=\; I'_{\alpha_0} \;=\; \bigcup\limits_{\substack{\alpha<\alpha_0\\ \textit{i.e.}\ \card\alpha\le\pi_0}} I'_\alpha .
\]
Alors $I'$ contient $I'_0$, il est grand devant $\pi_0$, mais quel est son cardinal ? Prouvons par récurrence transfinie que
\[
\card I'_\alpha \;\le\pi \qquad\quad \text{si } \alpha<\alpha_0,\ \ \textit{i.e.}\ \ \card\alpha\le\pi_0 .
\]
C'est O.K. pour $\alpha=0$ par hypothèse sur $I'_0$. Si c'est vrai pour $I'_\alpha$, est-ce vrai pour $I'_{\alpha+1}$ ? On adjoint à $I'_\alpha$ un ensemble d'éléments $i_J$ qui est de cardinal $\le\card\mathfrak{P}_{\le\pi_0}(I'_\alpha)$, donc de cardinal $\le\card(I'_\alpha)^{\pi_0}\le\pi^{\pi_0}$, et si on suppose
\[
\pi^{\pi_0} \;=\; \pi,
\]
on gagne. Supposons enfin que $\alpha$ soit un ordinal limite (de cardinal $\le\pi_0$). Donc on a
\[
I'_\alpha \;=\; \bigcup_{\beta<\alpha} I'_\beta ,
\]
donc
\printedpage{45}
% head: A. Grothendieck : Dérivateurs, Ch. XIII, Édition des Universités de Montpellier II et Paris VII
\msp{page 82}
\[
\card I'_\alpha \;\le\; \underbrace{(\card\alpha)\pi}_{\le\pi\pi=\pi^2=\pi} \qquad\quad \text{puisque les } \card I'_\beta\le\pi,
\]
on gagne encore.

On trouve enfin, par le même argument,
\[
\card I'_{\alpha_0} \;\le\; \underbrace{(\card\alpha_0)\pi\le\pi^2}_{\substack{\text{si }\card\alpha_0\le\pi,\\ \textit{i.e.}\text{ si }\pi_0<\pi}} = \pi,
\]
donc on gagne encore.

Ainsi, on a trouvé un deuxième système inductif
\[
(X_{\tilde\imath})_{\tilde\imath\in\tilde I}
\]
du même type que le système initial, ayant comme limite inductive $X$ (qu'il ait même limite inductive est immédiat). Mais qu'a-t-on gagné ? Ceci : dans $\tilde I$ toute partie de cardinal $\le\pi$ et grande devant $\pi_0$ est non seulement \textit{majorée} dans $\tilde I$, mais même y admet une borne supérieure. Cela résulte du % margin: double vertical rule beside the four lines "inductive est immédiat" ... "supérieure. Cela résulte du"

\msp{page 83}
\textbf{Lemme 2 :} \textit{Soit $I$ grand devant $\pi>\pi_0$, avec $\pi^{\pi_0}=\pi$} \texttt{[il suffit $\pi\ge\pi_0$]} \textit{(cf. lemme 1). Soit $(I'_\alpha)_{\alpha\in A}$ un ensemble de parties de $I$ qui sont dans $\tilde I$,} i.e.

\begin{itemize}
\item[a)] \textit{de cardinal $\le\pi$, et}
\item[b)] \textit{grandes devant $\pi_0$,}
\end{itemize}

\textit{et supposons que} $\card A\le\pi$ \textit{et que $A$, ordonné par inclusion, soit grand devant $\pi_0$. Alors $I'=\bigcup_{\alpha\in A} I'_\alpha$ est encore dans $\tilde I$,} i.e. \textit{$I'$ est de cardinal $\le\pi$ (car de cardinal $\le\pi^2=\pi$), et de plus, grande devant $\pi_0$.} (Vérification immédiate.)

Ainsi, on a prouvé ceci :

% margin: single vertical rule beside Proposition 1, from here to the end of the page
\textbf{Proposition 1.} \textit{Soit $E$} \texttt{[une]} \textit{grosse catégorie, avec filtration cardinale} $(\Filt^\pi)_{\pi\ge\pi_0}$. \textit{Soit $\pi$ un cardinal avec}
\[
\pi>\pi_0, \qquad\quad \pi^{\pi_0}=\pi \qquad\quad \textit{(p. ex. $\pi=2^c$, avec $c\ge\pi_0$)} \quad \texttt{[}\pi\ge\pi_0\ \texttt{suffit].}
\]
\textit{Alors pour tout $X$ dans $E$, on peut trouver un système inductif $(X_i)_{i\in I}$, de limite $X$, avec}

\begin{itemize}
\item[a)] \[ \left\{ \begin{aligned} & X_i\in\Filt^\pi \ \ \textit{pour tout $i\in I$,} \\ & \ I \ \textit{grand devant $\pi$,} \end{aligned} \right. \]
\end{itemize}

\textit{satisfaisant aux conditions supplémentaires suivantes :}

\begin{itemize}
\item[b)] \textit{Toute partie $J$ de $I$ de cardinal $\le\pi$, et grande devant $\pi_0$, admet une borne supérieure $i_J$ dans $I$.}
\item[c)] \textit{On a, sous ces conditions}
\end{itemize}

\printedpage{46}
% head: A. Grothendieck : Dérivateurs, Ch. XIII, Édition des Universités de Montpellier II et Paris VII
\msp{page 84}
\[
X_{i_J} \xleftarrow{\ \sim\ } \varinjlim_{i\in J} X_i .
\]
\textit{De plus, si $X\in\Filt^{\pi'}$ avec $\pi'>\pi$, on peut prendre $I$ de cardinal $\le\pi'^{\pi}$ (donc de cardinal $\le\pi'$ si $\pi'^{\pi}=\pi'$, p. ex. si $\pi'=2^{c'}$, avec $c'\ge\pi$).}

\textbf{Définition.} On appelle un système inductif $(X_i)_{i\in I}$, satisfaisant les conditions \textit{a}), \textit{b}), \textit{c}) ci-dessus, un \textit{système inductif bien adapté} (à la filtration cardinale donnée, pour les cardinaux $\pi_0$, $\pi$ donnés), et la représentation correspondante
\[
X \xleftarrow{\ \sim\ } \varinjlim X_i
\]
une \textit{représentation bien adaptée} de $X$ en termes d'objets de $\Filt^\pi$.

\textbf{NB} Quand $E$ est stable par petites $\varinjlim$ filtrantes, la situation se simplifie considérablement. On peut alors prendre
\[
\tilde I \;=\; \{I'\subset I \mid \card I'\le\pi,\ I' \text{ filtrant}\}
\]
(sans s'encombrer de la condition $I'$ grand devant $\pi_0$) et définir encore le système inductif $(X_{I'})_{I'\in\tilde I}$ par
\[
X_{I'} \;=\; \varinjlim_{i\in I'} X_i \;.
\]
On a encore
\msp{page 85}
\[
\card\tilde I \;\le\; (\card I)^\pi \;\le\; \pi'^{\pi} \qquad\quad \text{si }\ \card I\le\pi'^{\pi}
\]
(notamment, si $X\in\Filt^{\pi'}$, et qu'on prenne $I$ comme indiqué dans la condition $c$ des filtrations cardinales). De plus, cette fois, le fait que $\tilde I$ soit grand devant $\pi$ est trivial, on a même mieux, à part que $I$ soit filtrante :

\begin{itemize}
\item[$(*)$] \textit{Toute} partie $(I'_\alpha)_{\alpha\in A}$ de $\tilde I$ filtrante et de cardinal $\le\pi$ admet une borne supérieure $\tilde\imath_A$,
\end{itemize}

savoir la réunion (qui est bien filtrante). De plus, le système inductif $(X_{\tilde\imath})_{\tilde\imath\in\tilde I}$ a la propriété plus forte que
\[
X_{\tilde\imath_A} \xleftarrow{\ \sim\ } \varprojlim_{\tilde\imath\in A} X_{\tilde\imath} \;. % [sic] \varprojlim (arrow under lim pointing left) -> \varinjlim, inductive limit of the inductive system X_ĩ over the filtrant part A
\]
Ainsi on a la

\textbf{Proposition 1 bis.} \textit{Supposons $E$ stable par petites $\varinjlim$ filtrantes, et} $\Filt^\pi$ \textit{stable par celles-ci, où $\pi\ge\pi_0$ est donné. Alors pour tout $X$ dans $E$, on peut trouver un système inductif $(X_i)_{i\in I}$ de limite inductive $X$, avec}

\begin{itemize}
\item[a)] $X_i\in\Filt^\pi$ \textit{pour tout $i\in I$, $I$ ensemble ordonné grand devant $\pi$.}
\item[b)] \textit{Toute partie $J$ de $I$ filtrante de cardinal $\le\pi$ admet une borne supérieure $i_J$ dans $I$.}
\item[c)] \textit{Sous ces conditions}
\end{itemize}

\printedpage{47}
% head: A. Grothendieck : Dérivateurs, Ch. XIII, Édition des Universités de Montpellier II et Paris VII
\msp{page 86}
\[
X_{i_J} \;=\; \varinjlim_{i\in J} X_i \;.
\]
De tels systèmes inductifs seront dits \textit{“parfaitement adaptés”} (à la filtration cardinale donnée, pour le cardinal $\pi$ donné – $\pi_0$ n'y figure plus, si ce n'est comme début de la filtration cardinale sans plus).

Comme pratiquement, on aura à manier surtout des $E$ stables par petites $\varinjlim$ filtrantes, et des filtrations cardinales stables en conséquence (renforcement de la condition \textit{b}) des filtrations cardinales), c'est le cas, et la notion de \textit{représentation parfaitement adaptée} de $X$ comme $\varinjlim\limits_{I} X_i$ ($X_i\in\Filt^\pi$), que je mettrai en avant. Le cas “général”, où on suppose seulement \texttt{[que]} $E$ satisfait $L_{\pi_0}$, est plus alambiqué et se traite moins joliment (\footnote[41]{On parlera de filtration cardinale \textit{stricte,} si les $\underline{\Filt}^\pi$ sont stables par $\varinjlim$ filtrantes de cardinal $\le\pi$ (et que ces limites filtrantes petites existent dans $E$).}).

\textbf{Définition 2.} Soit $E$ munie d'une filtration cardinale stricte $(\underline{\Filt}^\pi)_{\pi\ge\pi_0}$. Soit $M$ une partie de $\Ob E$, $\underline{M}$ la sous-catégorie pleine de $E$ qu'elle définit, et soit $\pi$ un cardinal $\ge\pi_0$. On dit que $M$ est \textit{$\pi$-accessible,} où que $\underline{M}$ est \textit{$\pi$-accessible}, % [sic] où que -> ou que, disjunction "ou" not the relative "où"
\msp{page 87}
si les conditions suivantes sont satisfaites. Pour tout système inductif
\[
(X_i)_{i\in I}
\]
parfaitement adapté à la filtration cardinale pour $\pi$ (donc $I$ ensemble ordonné filtrant tel que toute partie filtrante $J$ de cardinal $\le\pi$ ait une borne supérieure $i_J$, et $X_{i_J}\xleftarrow{\ \sim\ }\varinjlim\limits_{j\in J} X_j$) ; soit $I_M$ l'ensemble des $i\in I$ tels que $X_i\in M$. Ceci posé, on veut ceci : pour que $X$ ($=\varinjlim X_i$) soit dans $M$, il faut et il suffit que $I_M$ soit cofinal dans $I$ (de sorte que $X$ apparaît comme $\varinjlim$ filtrante, avec ensemble d'indices $I_M$ grand devant $\pi$, d'éléments de $M\cap\underline{\Filt}^\pi\overset{\text{déf}}{=}M^\pi$).

\msp{page 88}
J'ai envie de reprendre la définition de la $\pi$-accessibilité de $M$, pour une filtration cardinale donnée, à coups de conditions équivalentes. Pour tout système inductif $(X_i)_{i\in I}$ dans $M$ \texttt{[plutôt dans $E$]}, je garde la notation
\[
I_M \;=\; \{i\in I \mid X_i\in M\}.
\]
Je pose, pour tout cardinal $\pi\ge\pi_0$,
\[
M^\pi \;=\; M\cap\Filt^\pi .
\]

\begin{itemize}
\item[\circled{a$_\pi$}] $M$ est l'ensemble des objets de la forme
\[
\varinjlim_{i\in I} X_i \;,
\]
où les $X_i$ sont dans $M^\pi$, et $I$ est grand devant $\pi$.
\end{itemize}

\printedpage{48}
% head: A. Grothendieck : Dérivateurs, Ch. XIII, Édition des Universités de Montpellier II et Paris VII

\begin{itemize}
\item[\circled{b$_\pi$}] Soit $(X_i)_{i\in I}$ un système inductif dans $M$ \texttt{[plutôt dans $E$]} parfaitement adapté à la filtration cardinale pour $\pi$, \textit{i.e.} $I$ filtrante et toute partie filtrante $J$ de $I$ de cardinal $\le\pi$ admet une borne supérieure $i_J$, et $X_{i_J}\xleftarrow{\ \sim\ }\varinjlim\limits_{J} X_i$, et enfin $X_i\in\Filt^\pi$ pour tout $i\in I$. Soit $X=\varinjlim X_i$. Pour que l'on ait $X\in M$, il faut et il suffit que $I_M$ soit cofinal dans $I$.
\item[\circled{c$_\pi$}] $\underline{M}^\pi=\underline{M}\cap\underline{\Filt}^\pi$ est une sous-catégorie de $\underline{M}$ génératrice par épimorphismes stricts, $\underline{M}$ est stable par $\varinjlim$ grandes devant $\pi$, enfin pour tout $X\in M$, $\underline{M}^\pi/X$ est une catégorie filtrante.
\end{itemize}

\msp{page 89}
Je présume que ces conditions sont équivalentes.

$b_\pi\Longrightarrow a_\pi$. On suppose \textit{b}) vérifié, on veut prouver \textit{a}). D'abord que tout $X\in M$ est de la forme dite $\varinjlim\limits_{I} X_i$, avec les $X_i\in M^\pi$ et $I$ grand devant $\pi$. Par \texttt{[la]} proposition 1, on sait qu'on a une présentation parfaite adaptée pour $\pi$, $X=\varinjlim\limits_{i\in I} X_i$. Par \textit{b}) “il faut”, la partie $I_M$ de $I$ est cofinale dans $I$, donc comme $I$ est grand devant $\pi$, $I_M$ aussi, et on aura (posant $J=I_M$)
\[
X \;=\; \varinjlim_{i\in J} X_i, \qquad\quad X_i\in M^\pi,\ \ J \text{ grand devant } \pi,
\]
O.K. Inversement, prouvons que si $X$ se met sous la forme dite dans \textit{a}), il est dans $M$. Cela va résulter du

\textbf{Lemme 1 :} \textit{Si la condition} b) (\textit{ou} a)) \textit{est satisfaite, $M$ est stable par limites inductives grandes devant $\pi$.}

Il faut prouver que si
\msp{page 90}
\[
X \;=\; \varinjlim_{i\in I} X_i \;,
\]
avec $I$ grand devant $\pi$, les $X_i$ dans $M$, alors $X\in M$. La condition \textit{b}) “il suffit” l'assure quand on suppose de plus les $X_i$ dans $\Filt^\pi$ (donc dans $M^\pi$), de sorte que l'on aura $I_M=I$, donc $I_M$ cofinal, donc $X\in M$ (\footnote[42]{Pas tout à fait – on y suppose plus, sur le système inductif. On ne peut se dispenser, même dans le cas $X_i\in M^\pi$ $\forall\, i\in I$, faire le raisonnement des pages suivantes.}). Si on ne suppose pas les $X_i$ dans $\Filt^\pi$, il faut travailler un peu plus.

Passer tout de suite à la partie de la démonstration ci-dessous, page 96, qui prouve \textit{b}) $\Longrightarrow$ \textit{a}). Le reste de la démonstration va prouver le

\textbf{Lemme 2.} \textit{Si tout élément de Filt$^\pi$ est $\pi$-accessible, alors} b$_\pi$) \textit{(et même} a$_\pi$\textit{) implique que $M$ est stable par $\varinjlim$ grandes devant $\pi$.}


\endgroup
% END UNIT Der-ChXIII-25-01-04/Der-ChXIII-25-01-04_p046-050.tex
% BEGIN UNIT Der-ChXIII-25-01-04/Der-ChXIII-25-01-04_p051-055.tex
\begingroup\makeatletter%
\newcommand{\printedpage}[1]{\marginpar{\footnotesize p.~#1}}%
\newcommand{\msp}[1]{\marginpar{\footnotesize [#1]}}%
\newcommand{\barlongleftrightarrow}{\mathrel{\text{\ooalign{$\longleftrightarrow$\cr\hidewidth$\mid$\hidewidth\cr}}}}%
\newcommand{\idx}[2]{\noindent #1\dotfill #2\par}%
\makeatother\let\printedpage\derpp\let\msp\dermsp%
\let\mspin\DERmspin%


\printedpage{49}
% head: A. Grothendieck : Dérivateurs, Ch. XIII, Édition des Universités de Montpellier II et Paris VII
% foot: Version 25 janvier 2004

\msp{page 91}Soit
\[ \mathcal{C} \;=\; \underline{M}^\pi/\mathcal{X} \]
la catégorie suivante. Ses objets sont les triples
\[ (Z,i,\alpha) \quad \text{tels que} \quad Z\in M^\pi,\ i\in I,\ \alpha : Z\longrightarrow X_i \ , \]
les flèches
\[ (Z,i,\alpha) \;\longrightarrow\; (Z',i',\alpha') \]
sont les couples $(u,\lambda)$ avec
\[ u : Z\longrightarrow Z', \qquad \lambda : i\longrightarrow i', \qquad (\textit{i.e.}\ \text{on a } i\le i' \text{ et } \lambda \text{ est l'unique flèche de } i \text{ dans } i'), \]
tels que le diagramme
\[
\begin{tikzcd}[row sep=large, column sep=large]
Z \arrow[r,"\alpha"] \arrow[d,"u"'] & X_i \arrow[d,"\text{transition } \lambda_{i',i}"] \\
Z' \arrow[r,"\alpha'"] & X_{i'}
\end{tikzcd}
\]
commute. On peut donc définir $\mathcal{C}$ comme le produit fibré dans le carré cartésien
\[
\begin{tikzcd}[row sep=large, column sep=large]
 & \mathcal{C} \arrow[r,"\text{cofibré}"] \arrow[d] \arrow[dl,dotted,"\mathcal{Z}"'] & I \arrow[d,"\mathcal{X} : X \longmapsto X_i"] \\ % [sic] X \longmapsto X_i -> i \longmapsto X_i, the functor \mathcal{X} goes from I, so the variable should be i
\underline{M}^\pi & \underline{M}^\pi/E \arrow[l,dotted,"\text{source}"'] \arrow[r,"\text{but}","\text{cofibré}"'] & E.
\end{tikzcd}
\]
La fibre $\mathcal{C}_i$ de $\mathcal{C}$ en $i\in I$ est $\mathrm{Filt}^\pi/X_i$ \texttt{[plutôt} $M^\pi/X_i$\texttt{]}. On a, comme \textit{b}) $\Longrightarrow$ première moitié de \textit{a}), que $\underline{M}^\pi$ est un
\msp{page 92}système de générateurs stricts de $\underline{M}$ \texttt{[cela n'est pas clair sans autres hypothèses sur} $\underline{M}$\texttt{]}, et $X_i\in\operatorname{Ob}\underline{M}$ par hypothèse,
\[ X_i \;=\; \varinjlim_{Z\in\operatorname{Ob}\underline{M}^\pi/X_i} Z \;=\; \varinjlim_{\mathcal{C}_i}\mathcal{Z}|\mathcal{C}_i, \]
donc le système inductif donné $(X_i)_{i\in I}$ sur $I$ à valeurs dans $X$, se déduit du système inductif (non ordonné en général) sur $\mathcal{C}$, à valeurs dans $\underline{M}^\pi$, en ``intégrant sur les fibres'' de $\mathcal{C}\longrightarrow I$. D'ailleurs, comme $\underline{M}^\pi/E\longrightarrow E$ est cofibrée, donc propre, il en est de même de $\mathcal{C}\longrightarrow I$, donc les $\varinjlim$ de $\mathcal{Z}$ sur les fibres $\mathcal{C}_i$ ($=$ $\underline{M}^\pi/X_i$) forment bien un foncteur covariant $I\longrightarrow E$, dont la limite inductive est celle de $\mathcal{Z}$ (\footnote[43]{\textbf{NB} Si on admet la validité de $a$), alors il en résulte déjà que $\varinjlim\limits_{i\in I} X_i \simeq \varinjlim\limits_{\alpha\in C} Z_\alpha$ est dans $M$. (Stabilité de $M$ par $\varinjlim$ grandes devant $\pi$).}).

\printedpage{50}
% head: A. Grothendieck : Dérivateurs, Ch. XIII, Édition des Universités de Montpellier II et Paris VII
% foot: Version 25 janvier 2004

\msp{page 93}Soit, pour tout $i\in I$, $\tilde{\mathcal{C}}_i$ l'ensemble ordonné
\[ \tilde{\mathcal{C}}_i = \{\text{sous-catégories } \mathcal{C}'_i \text{ de } \mathcal{C}_i = \underline{M}^\pi/X_i \mid \operatorname{card}\mathcal{Z}(\mathcal{C}'_i)\le\pi,\ \ \mathcal{C}'_i \text{ filtrante}\}. \]
J'admets pour l'instant le

\textbf{Lemme} (\footnote[44]{Utilise l'hypothèse tacite que les éléments de $\mathrm{Filt}^\pi$ sont $\pi$-accessibles.}). \textit{Pour tout $X\in M$, la catégorie $\underline{M}^\pi/X$ est filtrante. (Sous l'hypothèse} b), \textit{ou que tout objet de $M$ est de la forme dite dans} a).\textit{)}

Il en résulte, comme $X_i\in M$, que $\mathcal{C}_i$ est filtrante, donc la réunion des $\mathcal{C}'_i\in\tilde{\mathcal{C}}_i$ est $\mathcal{C}_i$ ; l'ensemble ordonné $\tilde{\mathcal{C}}_i$ est filtrant, et toute partie $(\mathcal{C}'_{i_\alpha})_{\alpha\in A}$ filtrante et de cardinal $\le\pi$ de $\tilde{\mathcal{C}}_i$ y admet une borne supérieure $\mathcal{C}'_{i_A}$, savoir la réunion des $\mathcal{C}'_{i_\alpha}$. $\tilde{\mathcal{C}}_i$ est donc un ensemble ordonné $\pi$-adapté.

On peut considérer les ensembles ordonnés $\tilde{\mathcal{C}}_i$, pour $i\in I$ variable, comme les fibres d'une application croissante d'ensembles ordonnés
\[ p \;:\; \tilde{\mathcal{C}} \;\longrightarrow\; \mathcal{C}, \]
en écrivant
\[ \left\{ \begin{array}{l}
(\mathcal{C}'_i\subset\mathcal{C}_i) \;\le\; (\mathcal{C}'_j\subset\mathcal{C}_j) \\
\text{si et seulement si } i\le j, \\
\text{et le foncteur canonique de transition } \mathcal{C}_i\longrightarrow\mathcal{C}_j \text{ applique } \mathcal{C}'_i \text{ dans } \mathcal{C}'_j.
\end{array}\right. \]
On voit donc que $\tilde{\mathcal{C}}$ est lui aussi filtrant et grand devant $\pi$.

\msp{page 94}On définit un système inductif sur $\tilde{\mathcal{C}}$, à valeurs dans $\underline{\mathrm{Filt}}^\pi$, en posant
\[ Z_{\mathcal{C}'_i} \;=\; \varinjlim_{\alpha\,\in\,\mathcal{C}'_i\ (\subset\,\mathcal{C}_i\,=\,\underline{M}^\pi/X_i)} Z_\alpha \]
(où un $\alpha\in\mathcal{C}_i$ est formé, je le rappelle, de $Z_\alpha\xrightarrow{u_\alpha}X_i$, avec $Z_\alpha\in M^\pi\subset\underline{\mathrm{Filt}}^\pi$). Ces $Z_{\mathcal{C}'_i}$ ne sont plus forcément dans $M^\pi$, mais seulement dans $\mathrm{Filt}^\pi$. (Alors à quoi ça a servi de définir les $\mathcal{C}_i$ eux-mêmes comme $\underline{M}^\pi/X_i$, au lieu de $\underline{\mathrm{Filt}}^\pi/X_i$ ?)

Le système inductif
\[ (Z_{\mathcal{C}'_i})_{\mathcal{C}'_i\in\tilde{\mathcal{C}}_i} \]
est parfaitement bien adapté à la filtration cardinale pour le cardinal $\pi$, et il est immédiat que sa $\varinjlim$ est isomorphe à $\varinjlim\limits_{\alpha\in\operatorname{Ob}\mathcal{C}_i=\underline{M}^\pi/X_i} Z_\alpha = X_i\in M$ . Donc par \textit{b}), l'ensemble $\tilde{\mathcal{C}}_M$ des $\mathcal{C}'_i\in\tilde{\mathcal{C}}_i$ tels que $Z_{\mathcal{C}'_i}\in M$ (donc $Z_{\mathcal{C}'_i}\in M^\pi$) est cofinal dans $\tilde{\mathcal{C}}_i$. Il s'ensuit donc que l'ensemble $\tilde{\mathcal{C}}_M$ est lui aussi cofinal dans $\tilde{\mathcal{C}}$.

\printedpage{51}
% head: A. Grothendieck : Dérivateurs, Ch. XIII, Édition des Universités de Montpellier II et Paris VII
% foot: Version 25 janvier 2004

\msp{page 95}On vérifie aussi immédiatement que
\[ \varinjlim_{\alpha\in\tilde{\mathcal{C}}} Z_\alpha \;=\; \varinjlim X_i \;=\; X, \]
donc, $\tilde{\mathcal{C}}_M$ étant cofinal,
\[ \varinjlim_{\alpha\in\tilde{\mathcal{C}}_M} Z_\alpha \;=\; X. \]
Au total, on est arrivé à représenter $X$ sous la forme
\[ X \;=\; \varinjlim_{\alpha\in A} X_\alpha, \]
où cette fois les $X_\alpha$ sont \textit{dans} $M^\pi$ (pas seulement dans $M$), et $A$ est grand devant $\pi$. Mais on ne sait pas que toute partie filtrante $B$ de cardinal $\le\pi$ de $A$ admet une borne supérieure $\alpha_B$ dans $A$, et encore moins que $X_{\alpha_B} = \varinjlim\limits_{\alpha\in B} X_\alpha$ (\footnote[45]{Si on suppose $(a_\pi)$, cet argument prouve du moins que $M$ est stable par $\varinjlim$ grandes devant $\pi$.}). Donc on ne peut encore appliquer directement \textit{b}). S'il ne s'agissait que de prouver la deuxième moitié de \textit{a}), où on suppose déjà les $X_i\in M^\pi$ et $I$ grand devant $\pi$, on n'a rien gagné du tout par nos labeurs !

\msp{page 96}Mais revenons à ce cas
\[ X \;=\; \varinjlim X_i \;\left|\; \begin{array}{l} X_i\in M^\pi \\[1.5ex] I \text{ grand devant } \pi, \end{array}\right. \]
et introduisons à nouveau
\[ \tilde{I} \;=\; \{I'\subset I \mid I' \text{ filtrant, de cardinal } \le\pi\}, \]
donc cette fois $\tilde{I}$ est bien adapté à $\pi$ (filtrant, et toute partie filtrante de cardinal $\le\pi$ a une borne supérieure), et si on pose
\[ X_{\tilde{\imath}} \;=\; \varinjlim_{i\in I'_{\tilde{\imath}}} X_i, \]
on voit que le système inductif
\[ (X_{\tilde{\imath}})_{\tilde{\imath}\in\tilde{I}} \]
est parfaitement adapté à $\pi$ et la filtration cardinale. De plus, je dis que l'ensemble $\tilde{I}_M$ est \textit{cofinal} dans $\tilde{I}$. En effet, $\tilde{I}_M$ contient le sous-ensemble $\tilde{I}_0$ formé des $I'\in\tilde{I}$ qui ont un plus grand élément -- et $\tilde{I}_0$ est bien cofinal ! Maintenant, on peut appliquer
\printedpage{52}
% head: A. Grothendieck : Dérivateurs, Ch. XIII, Édition des Universités de Montpellier II et Paris VII
% foot: Version 25 janvier 2004
\msp{page 97}le ``il suffit'' dans \textit{b}), qui conclue que $X = \varinjlim X_i\in M$, q.e.d. % [sic] conclue -> conclut, verb form (indicative of conclure)

On a prouvé $b\Longrightarrow a$, sans utiliser l'hypothèse que tout objet de $\mathrm{Filt}^\pi$ est $\pi$-accessible. Mais on a utilisé cette hypothèse pour prouver que $b$ (ou $a$) implique que $M$ est stable par $\varinjlim$ grandes devant $\pi$ (lemme 2, p. 90).

$a\overset{?}{\Longrightarrow} b$. Le ``il suffit'' dans \textit{b}) est conséquence immédiate de \textit{a}). Il reste à prouver que si, dans le système inductif $\pi$-adapté $(X_i)_{i\in I}$ de limite $X$, on a $X\in M$, alors $I_M$ est cofinal dans $I$. Mais par hypothèse \textit{a}), on aura $X = \varinjlim\limits_J Y_j$, avec $J$ grand devant $\pi$, les $Y_j$ dans $M^\pi$. Admettons que $\mathrm{Filt}^\pi\subset E_\pi$, \textit{i.e.} tout élément de $\mathrm{Filt}^\pi$, en particulier les $X_i$, sont $\pi$-accessibles.

Je pense que je n'arriverai pas à prouver $a_\pi\Longrightarrow b_\pi$, mais seulement ceci : Si $\pi'$ est un cardinal $>\pi$ (disons le successeur de $\pi$ dans la suite ordinale des cardinaux), alors $a_\pi\Longrightarrow b_\pi$.

\msp{page 98}Ceci va résulter du

\textbf{Lemme 3.} \textit{Soient $X\simeq\varinjlim\limits_I X_i$, $X\simeq\varinjlim\limits_J Y_j$ deux représentations de $X$ sous forme des $\varinjlim$ filtrantes, avec $I$ grand devant $\pi'$, $J$ grand devant $\pi$, avec $\pi$, $\pi'$ deux cardinaux tels que $\pi'>\pi$. On suppose $(X_i)_{i\in I}$ $\pi'$-adapté,} i.e. \textit{pour toute partie filtrante $I'$ de $I$ de cardinal $\le\pi'$, $I'$ a une borne supérieure $i_{I'}$, et $X_{i_{I'}} = \varinjlim\limits_{i\in I'} X_i$. Alors pour tout $i\in I$, on peut trouver une partie $I'$ de $I$, contenant $i$,} grande devant $\pi$, \textit{de cardinal $\le\pi'$, et une partie $J'$ de $J$, enfin un isomorphisme $I'\xrightarrow[\sim]{\lambda}J'$ et une application strictement croissante $\mu : J'\longrightarrow I'$ et des homomorphismes} \texttt{[fin du chapitre]}

\par
\printedpage{53}
% head: A. Grothendieck : Dérivateurs, Ch. XIII, Édition des Universités de Montpellier II et Paris VII
% foot: Version 25 janvier 2004

\section*{Index des notations}

\idx{LLP, RLP}{1}
\idx{$\Phi_*$, $\Psi^*$}{1}
\idx{$\widetilde{\Phi}$, $\overline{\Psi}$}{1}
\idx{$\Phi \longleftrightarrow \Psi$}{3}
\idx{$\mathrm{Cof}_W$, $\mathrm{Fib}_W$}{8}
\idx{$W^{\mathrm{univ}}$}{9}
\idx{$W_{\mathrm{univ}}$}{10}
\idx{$\underline{\mathrm{Fib}}_W S\subset\mathcal{M}/S$}{10}
\idx{$W_S$}{10}
\idx{$W_S^{\mathrm{u}}$}{10}
\idx{$S\backslash\underline{\mathrm{Cof}}_W$}{10}
\idx{$\Phi \barlongleftrightarrow \Psi$}{28}
\idx{$\mathcal{M}_{\mathrm{cf}}$, $W_{\mathrm{cf}}$}{29}
\idx{$\mathrm{Ho}_W$}{29}
\idx{$\mathrm{Hom}_W$}{30}
\idx{$\pi_{\mathrm{d}}$, $\pi$}{30}
\idx{$\mathcal{M}_{\mathrm{c}}$}{30}
\idx{$W_{\mathrm{c}}$}{31}
\idx{$\pi_{\mathrm{g}}$}{31}
\idx{$\mathcal{M}_{\mathrm{f}}$}{31}
\idx{$f\sim g$}{32}
\idx{$\Psi,\ \Phi,\ \mathcal{M}(\Psi),\ W(\Psi),\ \mathcal{M}(\Phi),\ W(\Phi)$}{33}
\idx{$\mathcal{M}(I)$, $W(I)$}{34}
\idx{$F(I),\ TF(I),\ C(I),\ TC(I)$}{35}
\idx{$\alpha_s$}{35}
\idx{$u^*$, $u_*$, $u_!$}{45}
\idx{$u^*_W$, $u^W_*$, $u^W_!$}{49}
\idx{$\mathrm{HO}_W(I)$}{49}
\idx{Diag}{50}
\idx{$\mathrm{HO}_{\mathcal{M},W}(I)$, $\mathbf{D}_{\mathcal{M},W}(I)$}{54}
\idx{$\mathrm{HO}^{\mathrm{lc}}_W(I)$, $u^*_{\mathrm{lc}}$}{56}
\idx{$u^{\mathrm{f}}_*$}{59}
\idx{$u^*_{\mathrm{c}}$}{67}
\idx{$\mathfrak{U}$}{77}
\idx{$\mathrm{Filt}^\pi(E)$}{77}


\endgroup
% END UNIT Der-ChXIII-25-01-04/Der-ChXIII-25-01-04_p051-055.tex
% BEGIN UNIT Der-ChXIII-25-01-04/Der-ChXIII-25-01-04_p056-060.tex
\begingroup\makeatletter%
\newcommand{\printedpage}[1]{\marginpar{\footnotesize p.~#1}}%
\newcommand{\msp}[1]{\marginpar{\footnotesize [#1]}}%
\newcommand{\entry}[2]{\noindent #1\dotfill\ #2\par}%
\newcommand{\Fl}{\mathrm{Fl}}%
\makeatother\let\printedpage\derpp\let\msp\dermsp%
\let\mspin\DERmspin%


\printedpage{54}
% head: A. Grothendieck : Dérivateurs, Ch. XIII, Édition des Universités de Montpellier II et Paris VII
% foot: Version 25 janvier 2004

\entry{$L_\pi$}{77}
\entry{$\tilde{I}$}{79}
\entry{$\underline{M}$}{86}
\entry{$I_M$}{87}
\entry{$M^\pi$}{87}

\par

\printedpage{55}
% head: A. Grothendieck : Dérivateurs, Ch. XIII, Édition des Universités de Montpellier II et Paris VII
% foot: Version 25 janvier 2004

\section*{Index terminologique}

\entry{accessible (partie close à gauche)}{75}
\entry{admissible (à gauche)}{59}
\entry{catégorie de modèles}{7}
\entry{catégorie de Quillen close}{25}
\entry{clos à gauche, à droite}{2}
\entry{couple de Quillen}{5}
\entry{faiblement saturée}{7}
\entry{filtration cardinale}{77}
\entry{filtration cardinale stricte}{86}
\entry{flèche coüniversellement dans $W$}{9}
\entry{flèche universellement dans $W$}{10}
\entry{grand devant $\pi$}{77}
\entry{objet chemin}{32}
\entry{objet cylindre}{32}
\entry{$\pi$-accessible (partie d'une catégorie munie d'une filtration cardinale stricte)}{86}
\entry{pré-Q-catégorie de modèles à gauche, à droite}{12}
\entry{pré-Quillen à gauche, à droite}{14}
\entry{pré-quintuplet de Quillen clos}{21}
\entry{propre}{8}
\entry{propriété de factorisation de Quillen}{5}
\entry{Q-catégorie de modèles}{8}
\entry{Q-catégorie de modèles à droite}{21}
\entry{Q-catégorie de modèles “close”}{16}
\entry{Q-homotopie à gauche, à droite}{30}
\entry{quillénisable}{50}
\entry{quillénisable sur Diag}{51}
\entry{quillénisable sur Diag uniformément et exactement à gauche}{57}
\entry{quillénisable uniformément à gauche sur Diag}{57}
\entry{quillénisation uniforme à droite sur Diag}{58}
\entry{quillénisation uniforme à gauche sur Diag}{57}
\entry{quillénisation uniforme et exacte à gauche sur Diag}{57}
\entry{représentation bien adaptée}{84}
\entry{représentation parfaitement adaptée}{86}
\entry{strictement admissible}{61}
\entry{structure de Quillen}{28}

\par

\printedpage{56}
% head: A. Grothendieck : Dérivateurs, Ch. XIII, Édition des Universités de Montpellier II et Paris VII
% foot: Version 25 janvier 2004

\entry{structure de Quillen exacte}{27}
\entry{structure de Quillen exacte à gauche, à droite}{27}
\entry{structure de Quillen faible}{27}
\entry{structure de Quillen faible biexacte}{28}
\entry{structure de Quillen faible close, close à gauche, à droite}{28}
\entry{structure de Quillen parfaite}{28}
\entry{système de Quillen clos}{26}
\entry{système de Quillen faible}{26}
\entry{système de Quillen faible exact à gauche, à droite}{26}
\entry{système inductif bien adapté}{84}
\entry{système inductif parfaitement adapté}{86}
\entry{triple clos de Quillen}{25}
\entry{triple de Quillen}{14}
\entry{$W$-cofibration}{8}
\entry{$W$-équivalence coüniverselle}{9}
\entry{$W$-fibration}{8}

\par

\printedpage{57}
% head: A. Grothendieck : Dérivateurs, Ch. XIII, Édition des Universités de Montpellier II et Paris VII
% foot: Version 25 janvier 2004

\section*{Liste des axiomes}

\noindent $\mathcal{M}$ catégorie

\noindent $W \ \subset\ \Fl(\mathcal{M})$ , \qquad $C$, $F$, $TC$, $TF \ \subset\ \Fl(\mathcal{M})$

\begin{description}[leftmargin=4.5em,labelwidth=4em]
\item[M 0] La catégorie $\mathcal{M}$ admet des limites projectives finies et des limites inductives finies.

\item[M 1] Les flèches appartenant à $C$ ont la propriété de relèvement à gauche relativement à celles appartenant à $TF$, et les flèches appartenant à $F$ ont la propriété de relèvement à droite relativement à celles appartenant à $TC$.

\item[M 2] Toute flèche de $\mathcal{M}$ se décompose en une flèche de $C$ suivie d'une flèche de $TF$, et en une flèche de $TC$ suivie d'une flèche de $F$.

\item[M 3] La partie $F$ de $\Fl(\mathcal{M})$ est stable par composition et changement de base, $C$ est stable par composition et cochangement de base, $F$ et $C$ contiennent les isomorphismes.

\item[M 4] La partie $TF$ de $\Fl(\mathcal{M})$ est stable par composition et changement de base, $TC$ est stable par composition et cochangement de base, $TF$ et $TC$ contiennent les isomorphismes.

\item[M 5] Si dans un triangle commutatif de $\mathcal{M}$ deux des trois flèches sont dans $W$, il en est de même de la troisième ; $W$ contient les isomorphismes.

\item[M$'$] $TC \ \subset\ C \cap W$, $TF \ \subset\ F \cap W$.

\item[M$''$] Les flèches appartenant à $C$ sont des $W$-cofibrations, et les flèches appartenant à $F$ sont des $W$-fibrations.

\item[M$'_{\mathbf{0}}$(g)] $TC = C \cap W$.

\item[M$'_{\mathbf{0}}$(d)] $TF = F \cap W$.

\item[M 2\,bis] $W = TF \circ TC$.
\end{description}

\par

\printedpage{58}
% head: A. Grothendieck : Dérivateurs, Ch. XIII, Édition des Universités de Montpellier II et Paris VII
% foot: Version 25 janvier 2004

\setcounter{footnote}{0}
\renewcommand{\thefootnote}{\fnsymbol{footnote}}
\section*{Glossaire}

\subsection*{Section 1}

\begin{itemize}
\item $\Phi \ \subset\ \Fl(\mathcal{M})$

$\Phi$ est \textit{clos à gauche} : il existe \ $\Psi \ \subset\ \Fl(\mathcal{M})$ tel que $\Phi$ soit formé des flèches ayant la propriété de relèvement à gauche relativement aux flèches de $\Psi$.

\item $\Psi \ \subset\ \Fl(\mathcal{M})$

$\Psi$ est \textit{clos à droite} : il existe \ $\Phi \ \subset\ \Fl(\mathcal{M})$ tel que $\Psi$ soit formé des flèches ayant la propriété de relèvement à droite relativement aux flèches de $\Phi$.

\item $\Phi$, $\Psi \ \subset\ \Fl(\mathcal{M})$

$(\Phi, \Psi)$ a la \textit{propriété de factorisation de Quillen} :
\begin{itemize}
\item[a)] les flèches de $\Phi$ ont la propriété de relèvement à gauche relativement aux flèches de $\Psi$ ;
\item[b)] toute flèche $f$ de $\mathcal{M}$ se factorise en
\[
f = pi \quad , \qquad i \in \Phi \ , \quad p \in \Psi \ .
\]
\end{itemize}

\item $(\Phi, \Psi)$ \textit{couple de Quillen} :
\begin{itemize}
\item[a)] $(\Phi, \Psi)$ a la propriété de factorisation de Quillen ;
\item[b)] $\Phi$ est clos à gauche et $\Psi$ clos à droite.
\end{itemize}
\end{itemize}

\subsection*{Section 2}

\begin{itemize}
\item $W \ \subset\ \Fl(\mathcal{M})$

\textit{$W$ faiblement saturé} :
\begin{itemize}
\item[a)] $W$ contient les isomorphismes ;
\item[b)] si deux parmi les trois flèches d'un triangle commutatif sont dans $W$, la troisième aussi.
\end{itemize}

\item $W \ \subset\ \Fl(\mathcal{M})$

Un morphisme $p : X \longrightarrow Y$ de $\mathcal{M}$ est une \textit{$W$-fibration} ($p \in \mathrm{Fib}_W$) si pour tout diagramme de carrés cartésiens
\end{itemize}


\endgroup
% END UNIT Der-ChXIII-25-01-04/Der-ChXIII-25-01-04_p056-060.tex
% BEGIN UNIT Der-ChXIII-25-01-04/Der-ChXIII-25-01-04_p061-065.tex
\begingroup\makeatletter%
\newcommand{\printedpage}[1]{\marginpar{\footnotesize p.~#1}}%
\newcommand{\msp}[1]{\marginpar{\footnotesize [#1]}}%
\newcommand{\Fl}{\mathrm{Fl}}%
\newcommand{\Filt}{\mathrm{Filt}}%
\newcommand{\Mo}{\mathcal{M}}%
\newcommand{\Mzg}{$\mathbf{M}'_{\mathbf{0}}\mathbf{(g)}$}%
\newcommand{\Mzd}{$\mathbf{M}'_{\mathbf{0}}\mathbf{(d)}$}%
\renewcommand{\thefootnote}{\fnsymbol{footnote}}%
\makeatother\let\printedpage\derpp\let\msp\dermsp%
\let\mspin\DERmspin%
\edef\derfnsaved{\the\value{footnote}}\setcounter{footnote}{0}%


\printedpage{59}
% head: A. Grothendieck : Dérivateurs, Ch. XIII, Édition des Universités de Montpellier II et Paris VII

\[
\begin{tikzcd}
X \arrow[d, "p"'] & X' \arrow[l] \arrow[d] & X'' \arrow[l, "t"'] \arrow[d] \\
Y & Y' \arrow[l] & Y'' \arrow[l, "s"]
\end{tikzcd}
\quad ,
\]

si $s \in W$, alors $t \in W$.

\begin{itemize}
\item[\textbullet] $W \ \subset\ \Fl(\Mo)$

Un morphisme $i : A \longrightarrow B$ de $\Mo$ est une $W$-\textit{cofibration} ($i \in \mathrm{Cof}_W$) si pour tout diagramme de carrés cocartésiens
\[
\begin{tikzcd}
A'' \arrow[d] & A' \arrow[l, "s"'] \arrow[d] & A \arrow[l] \arrow[d, "i"] \\
B'' & B' \arrow[l, "t"] & B \arrow[l]
\end{tikzcd}
\quad ,
\]
si $s \in W$, alors $t \in W$.

\item[\textbullet] $W \ \subset\ \Fl(\Mo)$

Un morphisme $p : X \longrightarrow Y$ de $\Mo$ est \textit{universellement dans} $W$, ou une $W$-\textit{équivalence universelle} ($p \in W_{\mathrm{univ}}$) si pour tout carré cartésien
\[
\begin{tikzcd}
X \arrow[d, "p"'] & X' \arrow[l] \arrow[d, "p'"] \\
Y & Y' \arrow[l]
\end{tikzcd}
\quad ,
\]
la flèche $p'$ appartient à $W$.

\item[\textbullet] $W \ \subset\ \Fl(\Mo)$

Un morphisme $i : A \longrightarrow B$ de $\Mo$ est \textit{coüniversellement dans} $W$, ou une $W$-\textit{équivalence coüniverselle} ($i \in W^{\mathrm{univ}}$) si pour tout carré cocartésien
\[
\begin{tikzcd}
A' \arrow[d, "i'"'] & A \arrow[l] \arrow[d, "i"] \\
B' & B \arrow[l]
\end{tikzcd}
\quad ,
\]
la flèche $i'$ appartient à $W$.
\end{itemize}
% foot: Version 25 janvier 2004

\printedpage{60}
% head: A. Grothendieck : Dérivateurs, Ch. XIII, Édition des Universités de Montpellier II et Paris VII

\begin{itemize}
\item[\textbullet] \textit{Pré-Q-catégorie de modèles à gauche} (resp. \textit{à droite}) :

$\Mo$ et $W$, $C$, $F$, $TC$, $TF$ $\subset$ $\Fl(\Mo)$ satisfaisant aux axiomes \textbf{M~0} à \textbf{M~5}, $\mathbf{M}'$, $\mathbf{M}''$ et \Mzg{} (resp. \Mzd).

\item[\textbullet] \textit{Q-catégorie de modèles}, \textit{triple de Quillen} :

catégorie de modèles au sens de Quillen.

\item[\textbullet] \textit{Q-catégorie de modèles “close”}, \textit{triple clos de Quillen}, \textit{catégorie de Quillen close}, \textit{système de Quillen clos} :

catégorie de modèles fermée au sens de Quillen.

\item[\textbullet] \textit{Pré-Quillen à gauche} (resp. \textit{à droite}) :

pré-Q-catégorie de modèles à gauche (resp. à droite).

\item[\textbullet] \textit{Pré-quintuplet de Quillen clos} :

$W$, $C$, $F$, $TC$, $TF$ $\subset$ $\Fl(\Mo)$ satisfaisant aux axiomes \textbf{M~0} à \textbf{M~5}, $\mathbf{M}'$, $\mathbf{M}''$, $C$, $TC$ étant clos à gauche et $F$, $TF$ clos à droite.

\item[\textbullet] \textit{Système de Quillen faible} :

$W$, $C$, $F$, $TC$, $TF$ $\subset$ $\Fl(\Mo)$ satisfaisant aux axiomes \textbf{M~0} à \textbf{M~5}, $\mathbf{M}'$, $\mathbf{M}''$.

\item[\textbullet] \textit{Système de Quillen faible exact à gauche} (resp. \textit{à droite}) :

$W$, $C$, $F$, $TC$, $TF$ $\subset$ $\Fl(\Mo)$ satisfaisant aux axiomes \textbf{M~0} à \textbf{M~5}, $\mathbf{M}'$, $\mathbf{M}''$ et \Mzg{} (resp. \Mzd).
\end{itemize}

\section*{Section 3}

\begin{itemize}
\item[\textbullet] $(\Mo,W)$ \textit{catégorie de modèles} :

$\Mo$ catégorie satisfaisant à l'axiome \textbf{M~0} ;

$W$ $\subset$ $\Fl(\Mo)$ faiblement saturé.

\item[\textbullet] \textit{Structure de Quillen faible associée à la catégorie de modèles} $(\Mo,W)$ :

quadruplet $(C,F,TC,TF)$ de parties de $\Fl(\Mo)$ satisfaisant aux axiomes \textbf{M~1} à \textbf{M~5}, $\mathbf{M}'$, $\mathbf{M}''$.

\item[\textbullet] \textit{Structure de Quillen faible exacte} :

structure de Quillen faible satisfaisant en outre à l'axiome \textbf{M~2\,bis}.

\item[\textbullet] \textit{Structure de Quillen faible exacte à gauche} (resp. \textit{à droite}) :

structure de Quillen faible satisfaisant en outre à l'axiome \Mzg{} (resp. \Mzd).
\end{itemize}
% foot: Version 25 janvier 2004

\printedpage{61}
% head: A. Grothendieck : Dérivateurs, Ch. XIII, Édition des Universités de Montpellier II et Paris VII

\begin{itemize}
\item[\textbullet] \textit{Structure de Quillen faible close à gauche} (resp. \textit{à droite}) :

structure de Quillen faible $(C,F,TC,TF)$ telle que $C$, $TC$ (resp. $F$, $TF$) soient clos à gauche (resp. à droite).

\item[\textbullet] \textit{Structure de Quillen faible close} :

structure de Quillen faible à la fois close à gauche et à droite.

\item[\textbullet] \textit{Structure de Quillen} :

structure de Quillen faible “biexacte” :

structure de Quillen faible à la fois exacte à gauche et à droite

$=$ catégorie de modèles de Quillen propre.

\item[\textbullet] \textit{Structure de Quillen parfaite} :

structure de Quillen faible “biexacte” et close

$=$ catégorie de modèles fermée de Quillen propre.
\end{itemize}

\section*{Section 5}

\begin{itemize}
\item[\textbullet] Catégorie de modèles $(\Mo,W)$ \textit{quillénisable} :

il existe une structure de Quillen faible \textit{exacte} $(C,F,TC,TF)$ associée à $(\Mo,W)$.

\item[\textbullet] $(\Mo,W)$ \textit{quillénisable sur} Diag $\subset$ Cat :

pour tout $I$ dans Diag, la catégorie des modèles $(\Mo(I),\, W(I))$ est quillénisable.
\end{itemize}

\section*{Section 6}

\begin{itemize}
\item[\textbullet] Catégorie de modèles $(\Mo,W)$ \textit{quillénisable uniformément à gauche sur} Diag :

il existe une structure de Quillen faible \textit{close} $(C,F,TC,TF)$ pour $(\Mo,\, W)$ telle que pour tout $I$ dans Diag, si l'on définit $C(I)$ et $TC(I)$ “argument par argument” (\footnote{Un morphisme de foncteurs $u : X \longrightarrow Y$ de $I^{\circ}$ vers $\Mo$ étant dans $C(I)$ (resp. $TC(I)$) si et seulement si pour tout objet $i$ de $I$, $u_i$ est dans $C$ (resp. $TC$).}), et $F(I)$ et $TF(I)$ par la propriété de relèvement à droite relativement à $TC(I)$ et $C(I)$ respectivement (\footnote{Un morphisme de $\Mo(I)$ étant dans $F(I)$ (resp. $TF(I)$) si et seulement si il a la propriété de relèvement à droite relativement aux flèches appartenant à $TC(I)$ (resp. $C(I)$).}), le quadruplet $(C(I),TC(I),F(I),TF(I))$ soit une structure de Quillen faible \textit{close exacte} pour $(\Mo(I),\, W(I))$. On dit alors que $(C,F,TC,TF)$ est une \textit{quillénisation uniforme à gauche} de $(\Mo,\, W)$ sur Diag).
\end{itemize}
% foot: Version 25 janvier 2004

\printedpage{62}
% head: A. Grothendieck : Dérivateurs, Ch. XIII, Édition des Universités de Montpellier II et Paris VII

\begin{itemize}
\item[\textbullet] Catégorie de modèles $(\Mo,W)$ \textit{quillénisable sur} Diag \textit{uniformément et exactement à gauche} :

$(\Mo,W)$ admet une quillénisation uniforme à gauche sur Diag $(C,F,TC,TF)$ telle que $TC = C \cap W$ (autrement dit, telle que $(C,F,TC,TF)$ soit une strucure de Quillen faible \textit{exacte à gauche}). On dit alors que $(C,F,TC,TF)$ est une \textit{quillénisation uniforme et exacte à gauche} de $(\Mo,\, W)$ sur Diag. % [sic] strucure -> structure, letter t omitted
\end{itemize}

\section*{Section 7}

\begin{itemize}
\item[\textbullet] $I$ ensemble ordonné filtrant, \quad $\pi$ cardinal

$I$ \textit{grand devant} $\pi$ :

toute partie de $I$ de cardinal inférieur ou égal à $\pi$ admet un majorant dans $I$.

\item[\textbullet] $E$ catégorie

Une \textit{filtration cardinale} de $E$ est une filtration croissante
\[
\Big(\Filt^{\pi}(E)\Big)_{\pi\geq\pi_0} \qquad , \qquad\qquad \pi_0 \text{ un cardinal infini fixé,}
\]
de $E$ par des sous-catégories strictement pleines $\Filt^{\pi}(E)$, indexée par les cardinaux $\pi \geq \pi_0$ appartenant à un univers donné, et satisfaisant aux conditions suivantes :
\begin{itemize}
\item[\textit{a})] pour tout $\pi \geq \pi_0$, $\Filt^{\pi}(E)$ est équivalente à une petite catégorie ;
\item[\textit{b})] la catégorie $E$ admet des limites inductives filtrantes indexées par des ensembles ordonnés $I$ grands devant $\pi_0$, et pour tout $\pi \geq \pi_0$, $\Filt^{\pi}(E)$ est stable par les dites limites si $\mathrm{card}\, I \leq \pi$ ;
\item[\textit{c})] Pour tout $\pi \geq \pi_0$, et pour tout $X \in \mathrm{Ob}\, E$, il existe un isomorphisme
\[
X \ \simeq\ \varinjlim_{i\in I} X_i \ ,
\]
où $(X_i)_{i\in I}$ est un système inductif filtrant dans $E$, indexé par un ensemble ordonné $I$ grand devant $\pi$, et où les $X_i$ sont des objets de $\Filt^{\pi}(E)$. De plus, si $X \in \mathrm{Ob}\,\Filt^{\pi'}(E)$, $\pi' \geq \pi$, on peut prendre $I$ tel que $\mathrm{card}\, I \leq \pi'^{\pi}$.
\end{itemize}

\item[\textbullet] $E$ catégorie, \quad $(\Filt^{\pi}(E))_{\pi\geq\pi_0}$ filtration cardinale de $E$, \quad $\pi \geq \pi_0$ cardinal

Un \textit{système inductif bien adapté} (à la filtration cardinale donnée, pour les cardinaux $\pi_0$, $\pi$ donnés) est un système inductif filtrant $(X_i)_{i\in I}$ satisfaisant aux conditions suivantes :
\begin{itemize}
\item[\textit{a})] $I$ est grand devant $\pi$, et pour tout $i \in I$, $X_i \in \Filt^{\pi}(E)$ ;
\item[\textit{b})] Toute partie $J$ de $I$ de cardinal $\leq \pi$, et grande devant $\pi_0$, admet une borne supérieure $i_J$, et on a
\[
X_{i_J} \ \xleftarrow{\ \sim\ }\ \varinjlim_{i\in J} X_i \ .
\]
% foot: Version 25 janvier 2004
\end{itemize}

\printedpage{63}
% head: A. Grothendieck : Dérivateurs, Ch. XIII, Édition des Universités de Montpellier II et Paris VII

Si $X$ est un objet de $E$, une \textit{représentation bien adaptée} de $X$, en termes d'objets de $\Filt^{\pi}(E)$, est la donnée d'un système inductif bien adapté $(X_i)_{i\in I}$, à la filtration cardinale donnée, pour les cardinaux $\pi_0$, $\pi$, et d'un isomorphisme
\[
X \ \xleftarrow{\ \sim\ }\ \varinjlim_{i\in I} X_i \ .
\]

\item[\textbullet] $E$ catégorie admettant des petites limites inductives filtrantes

Une \textit{filtration cardinale stricte} de $E$ est une filtration cardinale $(\Filt^{\pi}(E))_{\pi\geq\pi_0}$ telle que pour tout cardinal $\pi \geq \pi_0$, $\Filt^{\pi}(E)$ soit stable par limites inductives filtrantes.

\item[\textbullet] $E$ catégorie admettant des petites limites inductives filtrantes

$(\Filt^{\pi}(E))_{\pi\geq\pi_0}$ filtration cardinale stricte de $E$, \quad $\pi \geq \pi_0$ cardinal

Un \textit{système inductif parfaitement adapté} (à la filtration cardinale donnée, pour le cardinal $\pi$) est un système inductif filtrant $(X_i)_{i\in I}$ satisfaisant aux conditions suivantes :
\begin{itemize}
\item[\textit{a})] $I$ est grand devant $\pi$, et pour tout $i \in I$, $X_i \in \Filt^{\pi}(E)$ ;
\item[\textit{b})] Toute partie $J$ de $I$ filtrante de cardinal $\leq \pi$, admet une borne supérieure $i_J$, et on a
\[
X_{i_J} \ \xleftarrow{\ \sim\ }\ \varinjlim_{i\in J} X_i \ .
\]
\end{itemize}
Si $X$ est un objet de $E$, une \textit{représentation parfaitement adaptée} de $X$, en termes d'objets de $\Filt^{\pi}(E)$, est la donnée d'un système inductif bien adapté $(X_i)_{i\in I}$, à la filtration cardinale donnée, pour le cardinal $\pi$, et d'un isomorphisme
\[
X \ \xleftarrow{\ \sim\ }\ \varinjlim_{i\in I} X_i \ .
\]
\end{itemize}
% foot: Version 25 janvier 2004


\setcounter{footnote}{\derfnsaved}%
\endgroup
% END UNIT Der-ChXIII-25-01-04/Der-ChXIII-25-01-04_p061-065.tex
% BEGIN UNIT Der-ChXIII-25-01-04/Der-ChXIII-25-01-04_p066-066.tex
\begingroup\makeatletter%
\newcommand{\printedpage}[1]{\marginpar{\footnotesize p.~#1}}%
\newcommand{\msp}[1]{\marginpar{\footnotesize [#1]}}%
\makeatother\let\printedpage\derpp\let\msp\dermsp%
\let\mspin\DERmspin%


\printedpage{64}
% head: A. Grothendieck : Dérivateurs, Ch. XIII, Édition des Universités de Montpellier II et Paris VII
% foot: Version 25 janvier 2004

\section*{Références}

\noindent H. J. Baues, ``Algebraic Homotopy'', Cambridge studies in advanced mathematics, vol. 15 (1989).

\noindent R. Hartshorne, ``Residues and Duality'', Lecture Notes in Mathematics, vol. 20, Springer-Verlag (1966).

\noindent D. G. Quillen, ``Homotopical Algebra'', Lecture Notes in Mathematics, vol. 43, Springer-Verlag (1967).

\noindent R. W. Thomason, \textit{Cat as a Closed Model Category}, Cahiers de Topologie et Géométrie Différentielle, Vol. XXI-3, pp. 305-324 (1980).

\noindent J.-L. Verdier, ``Des catégories dérivées des catégories abéliennes'', Astérisque, vol. 239 (1996).

\noindent SGA 4, I, ``Théorie des Topos et Cohomologie Étale des Schémas'', Séminaire de Géométrie Algébrique du Bois-Marie, dirigé par M. Artin, A. Grothendieck, J.-L. Verdier, Lecture Notes in Mathematics, vol. 269, Springer-Verlag (1972).


\endgroup
% END UNIT Der-ChXIII-25-01-04/Der-ChXIII-25-01-04_p066-066.tex
\clearpage\setcounter{footnote}{0}\phantomsection\addcontentsline{toc}{section}{Chapitre XIV}\markboth{Chapitre XIV}{Chapitre XIV}
% BEGIN UNIT Der-ChXIV-07-10-10/Der-ChXIV-07-10-10_p001-005.tex
\begingroup\makeatletter%
\newcommand{\printedpage}[1]{\marginpar{\footnotesize p.~#1}}%
\newcommand{\msp}[1]{\marginpar{\footnotesize [#1]}}%
\makeatletter%
\renewcommand\@makefnmark{\hbox{\textsuperscript{\normalfont(\@thefnmark)}}}%
\renewcommand\@makefntext[1]{\parindent 1em\noindent\hb@xt@1.8em{\hss\textsuperscript{\normalfont\@thefnmark}}#1}%
\makeatother%
\makeatother\let\printedpage\derpp\let\msp\dermsp%
\let\mspin\DERmspin%


% ===================== PDF page 1 =====================
\printedpage{}% no page number printed on this page
% head: (none)

\begin{center}
\vspace*{6cm}

{\LARGE \textsf{LES DÉRIVATEURS}}

\vspace{1cm}

{\large \textbf{ALEXANDRE GROTHENDIECK}}

\vspace{2cm}

{\large \textbf{Chapitre XIV}}

\vspace{0.3cm}

{\large \textbf{Carrés h-cartésiens et h-cocartésiens}}
\end{center}

\par

% ===================== PDF page 2 =====================
\printedpage{}% no page number printed on this page
% head: (none)

\vspace*{6cm}

\noindent \par\noindent\textbf{Note des éditeurs de l’édition transcrite.}\par
Ce texte a été déchiffré et transcrit en \LaTeXe{} par M. Künzer. Il a été édité par M. Künzer, J. Malgoire, et G. Maltsiniotis. La transcription est aussi fidèle que possible au manuscrit. Pour les quelques corrections évidentes, ou rares commentaires des éditeurs, ainsi que pour la numérotation originale des pages du manuscrit, les caractères de machine à écrire \texttt{[typewriter]} entre crochets sont utilisés. Un point d'interrogation entre crochets signifie que l'on n'est pas sûr du mot qui précède.

\noindent Ce chapitre est inachevé. Dans la première section, Grothendieck introduit une notion de “catégorie quadrangulée” à gauche (resp. à droite). Le texte s'interrompt brutalement au début de la deuxième section, consacrée aux exemples, quand Grothendieck réalise, sans doute, que l'exemple principal qu'il a en vue (la catégorie homotopique d'une catégorie de modèles munie des carrés homotopiquement cartésiens (resp. cocartésiens)) ne satisfait \textit{pas} à l'axiomatique des catégories quadrangulées. Cela est brièvement expliqué dans une postface des éditeurs.

\noindent Cette édition est provisoire. Les remarques, commentaires et corrections sont bienvenus. Envoyer un message à :

\noindent \textbf{[e-mail address removed]}

\hfill G. Maltsiniotis\hspace*{2cm}

\par

% ===================== PDF page 3 =====================
\printedpage{1}
% head: A. Grothendieck : Dérivateurs, Ch. XIV, Édition des Universités de Montpellier II et Paris VII
% foot: Version 7 octobre 2010

\msp{page 1}

\section*{1\quad Axiomes de structures quadrangulées (une tentative avortée)}

Soit $(\mathcal{M},W)$ une catégorie de modèles admettant une théorie de la fibre homotopique. Alors on a dans $\mathrm{Ho}_W=\mathcal{M}W^{-1}$ une notion de \textit{carré $W$-cartésien} (ou $W$-homotopiquement cartésien), par quoi on entend un diagramme carré
\[
\begin{tikzcd}
X_1 \arrow[d] & X_3 \arrow[l] \arrow[d] \\
X_0 & X_2 \arrow[l]
\end{tikzcd}
\]
dans $\mathrm{Ho}_W$ provenant, à isomorphisme près, d'un carré $W$-cartésien dans $\mathcal{M}$ – donc isomorphe, dans $\underline{\mathrm{Hom}}(\boldsymbol{\Delta}_1\times\boldsymbol{\Delta}_1,\mathrm{Ho}_W)$, a un carré \textit{cartésien} \texttt{[dans $\mathcal{M}$]} % [sic] a un carré -> à un carré, preposition printed without its grave accent
\[
\begin{tikzcd}
Y_1 \arrow[d, "p_1"'] & Y_3\ (\simeq Y_1\times_{Y_0}Y_2) \arrow[l] \arrow[d] & \\
Y_0 & Y_2 \arrow[l, "p_2"] & ,
\end{tikzcd}
\]
avec $p_1,\ p_2\in\mathrm{Fib}_W$. (En fait, il suffit de $p_1$ ou $p_2$ dans $\mathrm{Fib}_W$ pour assurer le caractère $W$-cartésien du carré cartésien envisagé.) Je voudrais axiomatiser cette situation, histoire surtout de faire la liste des propriétés serviables essentielles de cette notion, et de la notion duale de carré $W$-cocartésien (ou $W$-homotopiquement cocartésien), dans le cas où $(\mathcal{M},W)$ admet une théorie de la cofibre homotopique. L'espèce de structure, savoir celle des “catégories quadrangulées” (à gauche, à droite, bilatères) est une variante, dans un cas \msp{page 2} non additif, de la notion de catégorie triangulée dégagée par \textsc{Verdier}. Comme celle-ci, cette notion est en fait trop faible – ce n'est \textit{pas} “la bonne” – mais elle a un intérêt pratique provisoire. La “bonne” notion (s'il en est une à part celle de dérivateur) devrait travailler avec une notion plus fine de carré et de carré $W$-cartésien, en travaillant systématiquement dans $\mathrm{HO}_W(\boldsymbol{\Delta}_1\times\boldsymbol{\Delta}_1)=\mathcal{M}(\boldsymbol{\Delta}_1\times\boldsymbol{\Delta}_1)W(\boldsymbol{\Delta}_1\times\boldsymbol{\Delta}_1)^{-1}$ (où vivent les “vrais” carrés $W$-cartésiens !), et catégories similaires, telles que $\mathrm{HO}_W(\Psi)$, $\mathrm{HO}_W(\Phi)$. Mais les inclure dans une axiomatique ne faisant pas mention d'une catégorie de modèles $(\mathcal{M},W)$, revient manifestement à mettre le doigt dans l'engrenage de la notion de dérivateur, et ce n'est pas mon propos pour l'instant.

\noindent\textbf{Définition 1.} J'appelle \textit{“catégorie à carrés h-cartésiens”,} ou \textit{“catégorie quadrangulée à gauche”,} une catégorie $\mathcal{H}$, munie d'un sous-ensemble Cart de l'ensemble $\mathrm{Hom}(\boldsymbol{\Delta}_1\times\boldsymbol{\Delta}_1,\mathcal{H})$
\printedpage{2}
% head: A. Grothendieck : Dérivateurs, Ch. XIV, Édition des Universités de Montpellier II et Paris VII
% foot: Version 7 octobre 2010
des diagrammes carrés \texttt{[par carré on entend carré commutatif]}
\[
\begin{tikzcd}
X_1 \arrow[d] & X_3 \arrow[l] \arrow[d] \\
X_0 & X_2 \arrow[l]
\end{tikzcd}
\]
dans $\mathcal{H}$, carrés appelés \textit{“h-cartésiens”} (h comme “homotopiquement”), cette \msp{page 3} donnée étant soumise aux axiomes qui suivent, Cart (1) à Cart (\ldots).

\noindent Comme je ne suis pas sûr de connaître d'avance toutes les propriétés essentielles, le nombre de ces axiomes reste pour le moment ouvert.

\noindent\textbf{Cart (1) (Symétrie).} Si un carré
\[
\begin{tikzcd}
X_1 \arrow[d] & X_3 \arrow[l] \arrow[d] \\
X_0 & X_2 \arrow[l]
\end{tikzcd}
\]
est h-cartésien, de même le carré symétrique (obtenu en échangeant le rôle de $X_1$ et de $X_2$).

\noindent\textbf{Cart (2) (Carrés triviaux).} Soit
\[
\begin{tikzcd}
X_1 \arrow[d, "p_1"'] & X_3 \arrow[l, "q_2"'] \arrow[d, "q_1"] \\
X_0 & X_2 \arrow[l, "p_2"]
\end{tikzcd}
\]
un carré dans $\mathcal{H}$, avec $p_1$ un isomorphisme. Alors le carré est h-cartésien si et seulement si $q_1$ est un isomorphisme.

\noindent\textbf{Cart (3) (Carrés composés).} Soit un diagramme commutatif
\[
\begin{tikzcd}
X'_0 \arrow[d] \arrow[dr, phantom, "\mathrm{I}"] & X'_1 \arrow[l] \arrow[d] \arrow[dr, phantom, "\mathrm{II}"] & X'_2 \arrow[l] \arrow[d] & \\
X_0 & X_1 \arrow[l] & X_2 \arrow[l] & ,
\end{tikzcd}
\]
d'où les carrés I, II et III :
\[
\begin{tikzcd}
X'_0 \arrow[d] & X'_2 \arrow[l] \arrow[d] \\
X_0 & X_2 \arrow[l]
\end{tikzcd}
\]
\printedpage{3}
% head: A. Grothendieck : Dérivateurs, Ch. XIV, Édition des Universités de Montpellier II et Paris VII
% foot: Version 7 octobre 2010
(qu'on appelle \textit{carré composé} (“horizontal”) des carrés I et II). Supposons I h-cartésien. Alors II est h-cartésien si et seulement si III l'est \texttt{[voir postface des éditeurs]}.

\noindent\textbf{Corollaire.} \textit{Soit un diagramme commutatif cubique, qu'on peut considérer comme homomorphisme d'un carré} (\textit{face supérieure du cube}) \textit{dans un autre} (\textit{face inférieure}) :
\[
\begin{tikzcd}
 & {} \arrow[dl] \arrow[dd] & & {} \arrow[ll] \arrow[dl] \arrow[dd] \\
{} \arrow[dd] & & {} \arrow[ll, crossing over] & \\
 & {} \arrow[dl] & & {} \arrow[ll] \arrow[dl] \\
{} & & {} \arrow[ll] \arrow[from=uu, crossing over] & .
\end{tikzcd}
\]
\textit{Si toutes les faces du cube}\footnote{En fait, il suffit de la face inférieure, et des deux faces parallèles latérales (p. ex. celle de devant et celle de derrière).} \textit{sauf la face supérieure sont} \msp{page 4} \textit{des carrés h-cartésiens (savoir la face inférieure et les quatre faces latérales), il en est de même de la face supérieure.}\footnote{\textbf{Corollaire} (de Cart (3) et Cart (2)) : \textit{tout carré isomorphe à un carré h-cartésien est h-cartésien.}}

\noindent Ces axiomes sont dans la nature d'axiomes de stabilité, nous donnant des critères affirmant que certains carrés sont h-cartésiens, si telles conditions sont remplies (p. ex. que certains autres carrés sont déjà h-cartésiens). J'ai l'impression d'avoir épuisé ici les propriétés pertinentes de ce type (\textit{i.e.} que les autres devraient pouvoir s'en déduire). À présent, je m'occupe de propriétés d'existence et d'unicité de carrés h-cartésien. % [sic] carrés h-cartésien -> carrés h-cartésiens, missing plural agreement

\noindent\textbf{Cart (4) (Axiome du produit fibré).} Tout diagramme
\[
\begin{tikzcd}
X_1 \arrow[d] & \\
X_0 & X_2 \arrow[l]
\end{tikzcd}
\]
dans $\mathcal{H}$ peut se compléter en un carré h-cartésien
\[
\begin{tikzcd}
X_1 \arrow[d] & X_3 \arrow[l] \arrow[d] & \\
X_0 & X_2 \arrow[l] & .
\end{tikzcd}
\]
L'axiome suivant (\textit{cf.} son corollaire) assurera entre autres que le carré h-cartésien en question est \textit{unique à isomorphisme près} – mais on fera attention qu'il n'est en général pas unique “à isomorphisme unique près”.


\endgroup
% END UNIT Der-ChXIV-07-10-10/Der-ChXIV-07-10-10_p001-005.tex
% BEGIN UNIT Der-ChXIV-07-10-10/Der-ChXIV-07-10-10_p006-010.tex
\begingroup\makeatletter%
\newcommand{\printedpage}[1]{\marginpar{\footnotesize p.~#1}}%
\newcommand{\msp}[1]{\marginpar{\footnotesize [#1]}}%
\newcommand{\uHom}{\underline{\mathrm{Hom}}}%
\newcommand{\bDelta}{\boldsymbol{\Delta}}%
\newcommand{\ho}{\mathrm{ho}}%
\makeatother\let\printedpage\derpp\let\msp\dermsp%
\let\mspin\DERmspin%


\printedpage{4}
% head: A. Grothendieck : Dérivateurs, Ch. XIV, Édition des Universités de Montpellier II et Paris VII

\textbf{Cart (5) (Axiome d'exactitude)} Considérons deux carrés h-cartésiens
\[
X = \left(\begin{tikzcd}
X_1 \ar[d] & X_3 \ar[l] \ar[d] \\
X_0 & X_2 \ar[l]
\end{tikzcd}\right)
\qquad \text{et} \qquad
Y = \left(\begin{tikzcd}
Y_1 \ar[d] & Y_3 \ar[l] \ar[d] \\
Y_0 & Y_2 \ar[l]
\end{tikzcd}\right)
\]
et un homomorphisme $p = (p_0, p_1, p_2)$ entre les diagrammes de type
\[
\Phi \ = \
\begin{tikzcd}
1 \ar[d] & \\
0 & 2 \ar[l]
\end{tikzcd}
\]
associés (\textit{cf.} diagramme en traits pleins ci-dessous) :

\msp{page 5}
\[
\begin{tikzcd}[row sep=small, column sep=small]
& X_1 \ar[dl] \ar[dd, "p_1" near start] & & X_3 \ar[ll] \ar[dl] \ar[dd, dotted, "p_3"] \\
X_0 \ar[dd, "p_0"'] & & X_2 \ar[ll, crossing over] & \\
& Y_1 \ar[dl] & & Y_3 \ar[ll] \ar[dl] \\
Y_0 & & Y_2 \ar[ll] \ar[from=uu, crossing over, "p_2"' near end] & .
\end{tikzcd}
\]

Alors :
\begin{itemize}[leftmargin=2.5em]
\item[a)] Il existe un homomorphisme $\bar{p} : X \longrightarrow Y$ induisant $p$, \textit{i.e.} il existe une flèche $p_3$ de façon à compléter le diagramme cubique commutatif (en respectant donc la contrainte de commutativité pour les deux carrés adjacents à l'arête $p_3$).
\item[b)] Si de plus $p_0$, $p_1$, $p_2$ sont des isomorphismes, il en est de même de $p_3$ \texttt{[voir postface des éditeurs]}.
\end{itemize}

Prenant le cas où $p_0$, $p_1$, $p_2$ sont des identités, on trouve :

\textbf{Corollaire 1.} \textit{Sous les conditions de} Cart (4)\textit{, si on a} deux \textit{carrés h-cartésiens $X$, $X'$ qui complètent}
\[
\begin{tikzcd}[row sep=small]
& X_1 \ar[dl] & \\
X_0 & & , \\
& X_2 \ar[ul] &
\end{tikzcd}
\]
\textit{il existe un homomorphisme de $X$ dans $X'$ qui induise l'identité sur les sommets d'indices} 0, 1, 2\textit{, et un tel homomorphisme est toujours un isomorphisme.}

Mais cet isomorphisme n'est pas forcément unique, en d'autres termes, un automorphisme d'un carré h-cartésien peut induire l'identité sur les sommets $X_0$, $X_1$, $X_2$, sans être l'identité.
% foot: Version 7 octobre 2010

\printedpage{5}
% head: A. Grothendieck : Dérivateurs, Ch. XIV, Édition des Universités de Montpellier II et Paris VII

\msp{page 6}
\textbf{Corollaire 2.} \textit{Considérons un carré commutatif}
\[
Y = \left(\begin{tikzcd}
Y_1 \ar[d] & Y_3 \ar[l] \ar[d] \\
Y_0 & Y_2 \ar[l]
\end{tikzcd}\right) \quad ,
\]
\textit{et un objet $T$ de $\mathcal{H}$, d'où (posant} $\mathrm{Hom}(T,X) = X(T)$\textit{) un carré commutatif dans} (Ens)
\[
\begin{tikzcd}
Y_1(T) \ar[d] & Y_3(T) \ar[l] \ar[d] \\
Y_0(T) & Y_2(T) \ar[l]
\end{tikzcd} \quad ,
\]
\textit{et par suite}
\[
Y_3(T) \longrightarrow Y_1(T) \times_{Y_0(T)} Y_2(T) \quad . \tag*{$(*)$}
\]
\textit{Si $Y$ est h-cartésien, cette application} $(*)$ \textit{est} surjective \textit{quel que soit $T$.}

On applique Cart (5) à $Y$ et au carré constant $X$ de valeur $T$
\[
X = \quad
\begin{tikzcd}
T \ar[d, "\mathrm{id}"'] & T \ar[l, "\mathrm{id}"'] \ar[d, "\mathrm{id}"] \\
T & T \ar[l, "\mathrm{id}"']
\end{tikzcd} \quad ,
\]
en notant que la donnée de $p = (p_0, p_1, p_2)$ de
\[
\begin{tikzcd}
T \ar[d] & \\
T & T \ar[l]
\end{tikzcd}
\qquad\qquad \text{dans} \qquad\qquad
\begin{tikzcd}
Y_1 \ar[d] & \\
Y_0 & Y_2 \ar[l]
\end{tikzcd}
\]
revient à la donnée d'un élément du deuxième membre de $(*)$,
\[
\begin{tikzcd}
Y_1 \ar[d] & T \ar[l, "p_1"'] \ar[dl, "p_0"] \ar[d, "p_2"] \\
Y_0 & Y_2 \ar[l]
\end{tikzcd} \quad ,
\]
et que les éléments de $Y_3(T)$ qui s'envoient (par $(*)$) dans cet élément ne sont autres que les $p_3$ dont l'axiome Cart (5) affirme l'existence.

\textbf{N.B.} On notera que l'application surjective $(*)$ n'est en général \textit{pas} bijective, \textit{i.e.} le carré h-cartésien n'est en général pas cartésien.
% foot: Version 7 octobre 2010

\printedpage{6}
% head: A. Grothendieck : Dérivateurs, Ch. XIV, Édition des Universités de Montpellier II et Paris VII

\msp{page 7}
\textbf{N.B.} Il y a une notion duale de \textit{catégorie à carrés h-cocartésiens,} ou \textit{catégorie quadrangulée à droite.} Le plus souvent, les deux structures sont présentes simultanément. Il y a lieu de dégager des propriétés qui relient les deux structures duales sur une même catégorie. Ainsi, sauf erreur, dans le cas $\mathcal{H}$ additive, les deux espèces de carrés, les h-cartésiens et les h-cocartésiens, devraient \textit{coïncider.} (Il me faudra revenir là-dessus par la suite.) Dans le cas $\mathcal{H}$ ``ponctuée'', il y a l'adjonction des foncteurs $\Sigma$, $\Omega$ liés l'un à la structure quadrangulée droite, l'autre à la structure quadrangulée gauche. Et il y a, quand il y a des \underline{Hom} internes ou externes, des propriétés éventuelles de $\uHom(X,Y)$ par rapport aux quadrangles h-cocartésiens en $X$, h-cartésiens en $Y$ -- mais c'est à peine un lien entre les deux structures \ldots

\section*{2\quad Exemples}

Un premier exemple trivial, est celui où $\mathcal{H}$ est stable par produits fibrés, et où les carrés h-cartésiens sont définis comme les carrés cartésiens. Les axiomes Cart (1) à Cart (5) sont alors des propriétés bien connues, $\pm$ tautologiques de la notion de carré cartésien. Et
\msp{page 8}
en un sens, le cas général doit se ramener à ce cas particulier banal.

Le seul type d'autre cas que je connaisse, est celui où $\mathcal{H} = \mathcal{M}W^{-1}$, où $W$ est un localiseur admettant une théorie de la fibre $W$-homotopique (\textit{cf.} chapitre XII). Dans ce cas, on l'a dit, on prend comme carrés h-cartésiens dans $\mathcal{H}$ ceux qui sont isomorphes, dans $\uHom(\bDelta_1 \times \bDelta_1, \mathcal{H})$, à un carré $W$-cartésien dans $\mathcal{M}$. Encore faut-il vérifier les axiomes Cart (1) à Cart (5). Je pense que je n'ai guère de chance d'y arriver que moyennant une

\textbf{Proposition 1.} \textit{Moyennant des hypothèses sur $(\mathcal{M}, W)$ qui vont être dégagées ci-dessous, vérifiées en tous cas si $(\mathcal{M}, W)$ est quillénisable, sur la catégorie des ensembles ordonnés finis} (\textit{cf.} XIII, p. 51)\textit{, on a ceci : soit}
\[
X = \left(\begin{tikzcd}
X_1 \ar[d] & X_3 \ar[l] \ar[d] \\
X_0 & X_2 \ar[l]
\end{tikzcd}\right)
\]
\textit{un carré dans $\mathcal{M}$, et soit} $\ho_W(X)$ \textit{son image dans} $\uHom(\bDelta_1 \times \bDelta_1, \mathcal{H})$\textit{. Pour que} $\ho_W(X)$ \textit{soit h-cartésien, il} faut (\textit{et il suffit, bien sûr}) \textit{que $X$ soit $W$-cartésien} \texttt{[voir postface des éditeurs]}\textit{.}
% foot: Version 7 octobre 2010

\printedpage{7}
% head: A. Grothendieck : Dérivateurs, Ch. XIV, Édition des Universités de Montpellier II et Paris VII

\section*{Postface des éditeurs}

La proposition ci-dessus est fausse. On rappelle que si $\mathcal{M}$ est une catégorie de modèles fermée de Quillen de classe d'équivalences faibles $W$, un carré commutatif
\[
X = \left(\begin{tikzcd}
X_1 \ar[d] & X_3 \ar[l] \ar[d] \\
X_0 & X_2 \ar[l]
\end{tikzcd}\right)
\]
de $\mathcal{M}$ est $W$-cartésien si et seulement si il existe un carré \textit{cartésien}
\[
Y = \left(\begin{tikzcd}
Y_1 \ar[d] & Y_3 \ar[l] \ar[d] \\
Y_0 & Y_2 \ar[l]
\end{tikzcd}\right)
\]
avec $Y_1 \longrightarrow Y_0$ une fibration, et $Y_0$ et $Y_2$ fibrants (cette dernière condition étant superflue si la catégorie de modèles $\mathcal{M}$ est propre à droite), et un morphisme de carrés $s : X \longrightarrow Y$, $s = (s_0, s_1, s_2, s_3)$, \textit{i.e.} un cube commutatif
\[
\begin{tikzcd}[row sep=small, column sep=small]
& X_1 \ar[dl] \ar[dd, "s_1" near start] & & X_3 \ar[ll] \ar[dl] \ar[dd, "s_3"] \\
X_0 \ar[dd, "s_0"'] & & X_2 \ar[ll, crossing over] & \\
& Y_1 \ar[dl] & & Y_3 \ar[ll] \ar[dl] \\
Y_0 & & Y_2 \ar[ll] \ar[from=uu, crossing over, "s_2"' near end] & ,
\end{tikzcd}
\]
tel que $s$ soit une équivalence faible argument par argument, \textit{i.e.} $s_0, s_1, s_2, s_3 \in W$. Il est équivalent de demander que pour \textit{tout} diagramme commutatif à flèches verticales des équivalences faibles
\[
\begin{tikzcd}
X_1 \ar[r] \ar[d] & X_0 \ar[d] & X_2 \ar[l] \ar[d] \\
Y_1 \ar[r] & Y_0 & Y_2 \ar[l]
\end{tikzcd}
\]
tel que $Y_1 \longrightarrow Y_0$ soit une fibration, et $Y_0$, $Y_2$ des objets fibrants, la flèche canonique $X_3 \longrightarrow Y_1 \times_{Y_0} Y_2$ soit une équivalence faible.

Supposons maintenant que la catégorie de modèles $\mathcal{M}$ soit ponctuée, autrement dit qu'elle admette un objet $*$ à la fois initial et final, et soit $X$ un objet fibrant de $\mathcal{M}$. Choisissons une décomposition de la flèche $* \longrightarrow X$, en une équivalence faible suivie d'une fibration
\[
* \longrightarrow C \longrightarrow X \quad ,
\]
et formons le carré cartésien
% foot: Version 7 octobre 2010
\printedpage{8}
% head: A. Grothendieck : Dérivateurs, Ch. XIV, Édition des Universités de Montpellier II et Paris VII
\[
\begin{tikzcd}
C \ar[d] & \Omega X \ar[l] \ar[d] \\
X & * \ar[l]
\end{tikzcd} \quad ,
\]
où par définition $\Omega X = C \times_X *$ est l'\textit{espace des lacets} de $X$. Ce carré est $W$-cartésien. Par ailleurs, il est isomorphe dans $\uHom(\bDelta_1 \times \bDelta_1,\ \mathcal{M}W^{-1})$ au carré trivial
\[
\begin{tikzcd}
* \ar[d] & \Omega X \ar[l] \ar[d] \\
X & * \ar[l]
\end{tikzcd} \quad .
\]

Néanmoins, ce dernier n'est \textit{pas} en général $W$-cartésien. En effet, s'il l'était, l'unique flèche $s_3 : \Omega X \longrightarrow \Omega X$ de $\mathcal{M}$ rendant commutatif le cube
\[
\begin{tikzcd}[row sep=small, column sep=small]
& * \ar[dl] \ar[dd] & & \Omega X \ar[ll] \ar[dl] \ar[dd, dotted, "s_3"] \\
X \ar[dd, "\parallel"'] & & * \ar[ll, crossing over] & \\ % Audit resolved: vertical identity mark on X; see corrections.csv, PDF leaf 10.
& C \ar[dl] & & \Omega X \ar[ll] \ar[dl] \\
X & & * \ar[ll] \ar[from=uu, crossing over] &
\end{tikzcd}
\]
serait une équivalence faible. Or, comme la flèche $s_3$ se factorise par $* \times_X * \simeq *$, cela impliquerait que $\Omega X \to *$ est une équivalence faible. Cette condition est très restrictive. Si, par exemple, $\mathcal{M}$ est la catégorie de modèles des espaces topologiques pointés, et $X$ un CW-complexe pointé, elle équivaut à la contractibilité de la composante connexe de $X$ contenant le point base. De même, si $\mathcal{M}$ est la catégorie de modèles des complexes d'une catégorie abélienne de Grothendieck et $X$ un complexe, elle signifie que $X$ est acyclique.

Des considérations analogues montrent que les carrés $h$-cartésiens de $\mathcal{M}W^{-1}$, définis dans la section 2 comme étant les carrés commutatifs de $\mathcal{M}W^{-1}$ qui sont isomorphes dans $\uHom(\bDelta_1 \times \bDelta_1,\ \mathcal{M}W^{-1})$ à un carré $W$-cartésien, ne satisfont pas à toutes les conditions de Cart (3) et Cart (5) de la section 1. En effet, supposons toujours que $\mathcal{M}$ soit ponctuée, et gardons les notations ci-dessus. Considérons le diagramme commutatif dans $\mathcal{M}W^{-1}$
\[
\begin{tikzcd}
* \ar[d] & \Omega X \ar[l] \ar[d] & \Omega X \ar[l, "*_{\Omega X}"'] \ar[d] \\
X & * \ar[l] & * \ar[l]
\end{tikzcd} \quad ,
\]
où $*_{\Omega X}$ désigne l'unique endomorphisme de $\Omega X$ se factorisant par $*$. Il résulte de ce qui précède que le carré de gauche et le carré composé sont $h$-cartésiens. En revanche, le carré de droite ne l'est pas en général. En effet, si ce carré était isomorphe dans $\uHom(\bDelta_1 \times \bDelta_1,\ \mathcal{M}W^{-1})$ à un carré $W$-cartésien, la flèche $*_{\Omega X}$ serait un isomorphisme
% foot: Version 7 octobre 2010


\endgroup
% END UNIT Der-ChXIV-07-10-10/Der-ChXIV-07-10-10_p006-010.tex
% BEGIN UNIT Der-ChXIV-07-10-10/Der-ChXIV-07-10-10_p011-011.tex
\begingroup\makeatletter%
\newcommand{\printedpage}[1]{\marginpar{\footnotesize p.~#1}}%
\newcommand{\msp}[1]{\marginpar{\footnotesize [#1]}}%
\newcommand{\cM}{\mathcal{M}}%
\makeatother\let\printedpage\derpp\let\msp\dermsp%
\let\mspin\DERmspin%


\printedpage{9}
% head: A. Grothendieck : Dérivateurs, Ch. XIV, Édition des Universités de Montpellier II et Paris VII
% foot: Version 7 octobre 2010
de $\cM W^{-1}$, donc une équivalence faible de $\cM$, ce qui impliquerait les mêmes conditions très restrictives que ci-dessus. Ainsi l'une des implications de l'axiome Cart 3 n'est pas satisfaite.

De même le cube commutatif dans $\cM W^{-1}$
\[
\begin{tikzcd}[row sep=small, column sep=small]
 & * \arrow[dl] \arrow[dd] & & \Omega X \arrow[ll] \arrow[dl] \arrow[dd, "*_{\Omega X}"] \\
X \arrow[dd, "\shortparallel"'] & & * \arrow[ll, crossing over] & \\
 & * \arrow[dl] & & \Omega X \arrow[ll] \arrow[dl] \\
X & & * \arrow[ll] \arrow[from=uu, crossing over] & \mathllap{,}
\end{tikzcd}
\]
dont les faces horizontales sont $h$-cartésiennes, montre que la condition $(b)$ de l'axiome Cart (5) n'est pas satisfaite en général.


\endgroup
% END UNIT Der-ChXIV-07-10-10/Der-ChXIV-07-10-10_p011-011.tex
\clearpage\setcounter{footnote}{0}\phantomsection\addcontentsline{toc}{section}{Chapitre XV}\markboth{Chapitre XV}{Chapitre XV}
% BEGIN UNIT Der-ChXV-05-11-10/Der-ChXV-05-11-10_p001-005.tex
\begingroup\makeatletter%
\derleqno%
\newcommand{\printedpage}[1]{\marginpar{\footnotesize p.~#1}}%
\newcommand{\msp}[1]{\marginpar{\footnotesize [#1]}}%
\newcommand{\M}{\mathcal{M}}%
\newcommand{\Fl}{\mathrm{Fl}}%
\newcommand{\uFl}{\underline{\mathrm{Fl}}}%
\newcommand{\uNerf}{\underline{\mathrm{Nerf}}}%
\makeatother\let\printedpage\derpp\let\msp\dermsp%
\let\mspin\DERmspin%


% ===================== PDF page 1 =====================
\printedpage{---}
% no page number printed on this page
% head: (none)

\begin{center}
{\Large\textsf{LES DÉRIVATEURS}}

\bigskip

\textbf{ALEXANDRE GROTHENDIECK}

\vspace{2cm}

\textbf{Chapitre XV}

\medskip

\textbf{Théorèmes de factorisation}

\textbf{(Modèles (2))}
\end{center}

\par

% ===================== PDF page 2 =====================
\printedpage{---}
% no page number printed on this page
% head: (none)

\par\noindent\textbf{Note des éditeurs de l’édition transcrite.}\par
Ce texte a été déchiffré et transcrit en \LaTeXe{} par M. Künzer. Il a été édité par M. Künzer, J. Malgoire, et G. Maltsiniotis. La transcription est aussi fidèle que possible au manuscrit. Seules quelques corrections mineures évidentes ont été effectuées. Pour les rares ajouts ou commentaires des éditeurs, ainsi que pour la numérotation originale des pages du manuscrit, les caractères de machine à écrire \texttt{[typewriter]} entre crochets sont utilisés.

Cette édition est provisoire. Les remarques, commentaires et corrections sont bienvenus. Envoyer un message à :

\textbf{[e-mail address removed]}

\hfill G. Maltsiniotis\hspace*{2cm}

\par

% ===================== PDF page 3 =====================
\printedpage{1}
% head: A. Grothendieck : Dérivateurs, Ch. XV, Édition des Universités de Montpellier II et Paris VII
% foot: Version 5 novembre 2010

\msp{page 1}

\section*{1\quad Construction fondamentale : la suite transfinie de factorisation d'une flèche $\boldsymbol{f}$.}

Soit $\M$ une $\mathfrak{U}$-catégorie, stable par “petites” (\textit{i.e.} $\mathfrak{U}$-petites) $\varinjlim$ filtrantes. Soit
\[
C_0 \;\subset\; \Fl(\M)
\tag*{(1.1)}
\]
un \textit{petit} ensemble de flèches, on se propose de construire pour toute flèche $f$ de $\M$, une factorisation
\[
f \;=\; p\,i \,, \qquad i \;\in\; \widetilde{C_0} = ((C_0)_*)^* \;, \;\; p \;\in\; (C_0)_* \;.
\tag*{(1.2)}
\]
Pour ceci, on va construire une suite ordinale de foncteurs
\[
\left\{
\begin{array}{rcl}
f & \longmapsto & X_\alpha(f) \;=\; X_\alpha(C_0,f) \\
X_\alpha : \uFl(\M) & \longrightarrow & \M
\end{array}
\right.
\tag*{(1.3)}
\]
et des homomorphismes fonctoriels
\[
i_\alpha(f) : \mathrm{source}\, f \;\longrightarrow\; X_\alpha(f) \;,
\tag*{(1.4)}
\]
\[
p_\alpha(f) : X_\alpha(f) \;\longrightarrow\; \mathrm{but}\, f \;,
\tag*{(1.5)}
\]
de sorte qu'on aura le diagramme (où $X = \mathrm{source}\, f$, $Y = \mathrm{but}\, f$)
\[
\begin{tikzcd}[column sep=small]
X \arrow[rr, "f"] \arrow[dr, "i_\alpha(f)"'] & & Y \\
& X_\alpha(f) \arrow[ur, "p_\alpha(f)"'] &
\end{tikzcd}
\qquad ,
\tag*{(1.6)}
\]
ce diagramme étant commutatif
\[
f \;=\; p_\alpha(f) i_\alpha(f) \;.
\tag*{(1.7)}
\]
On peut dire aussi qu'on a une section du foncteur canonique de “composition”, transposé de $\boldsymbol{\Delta}_1 \longrightarrow \boldsymbol{\Delta}_2$, donné par $0 \longmapsto 0$, $1 \longmapsto 2$ :
\[
\left\{
\begin{array}{l}
\begin{tikzcd}[ampersand replacement=\&]
\uNerf_2(\M) \arrow[r, shift left, "\gamma"] \& \uNerf_1(\M) \arrow[l, dotted, shift left, "\Phi_\alpha"]
\end{tikzcd}
\;=\; \uFl(\M) \;, \\[1ex]
\\[1ex]
\\
\end{array}
\right.
\tag*{(1.8)}
\]
$\Phi$ comme “factorisation”,
\printedpage{2}
% head: A. Grothendieck : Dérivateurs, Ch. XV, Édition des Universités de Montpellier II et Paris VII
% foot: Version 5 novembre 2010
\msp{page 2}
où
\[
\Phi_\alpha(f) \;=\; \{\; X \xrightarrow{i_\alpha(f)} X_\alpha(f) \xrightarrow{p_\alpha(f)} Y \;\}
\tag*{(1.9)}
\]
avec $p_\alpha(f)i_\alpha(f) = f$. La fonctorialité de 1.4 et 1.5, ou celle de $\Phi_\alpha$, \textit{i.e.} de 1.9, s'exprime par le fait que tout carré commutatif
\[
\begin{tikzcd}[row sep=large]
X \arrow[r, "f"] \arrow[d, "u_X"'] & Y \arrow[d, "u_Y"] \\
X' \arrow[r, "{f'}"] & Y'
\end{tikzcd}
\tag*{(1.10)}
\]
dans $\M$, interprété comme flèche $u = (u_X, u_Y)$ de $f$ dans $f'$ dans $\uFl(\M)$, donne naissance au double carré commutatif
\[
\begin{tikzcd}[row sep=large]
X \arrow[r, "i_\alpha(f)"] \arrow[d, "u_X"] & X_\alpha(f) \arrow[r, "p_\alpha(f)"] \arrow[d, "X_\alpha(u)"] & Y \arrow[d, "u_Y"] \\
X' \arrow[r, "{i_\alpha(f')}"] & X_\alpha(f') \arrow[r, "{p_\alpha(f')}"] & Y' \;.
\end{tikzcd}
\tag*{(1.11)}
\]
\textbf{N.B.} Quand on voudra expliciter la dépendance de ces constructions par rapport à $C_0$, on écrira
\[
X_\alpha(C_0,f) \;, \qquad i_\alpha(C_0,f) \;, \qquad p_\alpha(C_0,f) \;, \qquad \Phi_\alpha(C_0,f) \;,
\]
ou
\[
X_\alpha(C_0,-) \;\;\text{ou}\;\; X_\alpha^{C_0} \;, \qquad i_\alpha(C_0,-) \;\;\text{ou}\;\; i_\alpha^{C_0} \qquad \textit{etc.} \;,
\]
au lieu de $X_\alpha(f)$, $i_\alpha(f)$, \dots

Voici les propriétés essentielles du système inductif ordinal des $X_\alpha(f)$ ou $\Phi_\alpha(f)$.

\msp{page 3}

\textbf{1.} $X_0(f) = X$, $i_0(f) = \mathrm{id}_X$, $p_0(f) = f$ (\textit{i.e.} la factorisation $\Phi_0(f)$ de $f$ est la factorisation triviale $X \xrightarrow{\mathrm{id}_X} X \xrightarrow{f} Y$).

\textbf{2.}
\[
\begin{aligned}
X_{\alpha+1}(f) &\;=\; X_1(p_\alpha(f)) \\
i_{\alpha+1}(f) &\;=\; i_1(p_\alpha(f)) \circ i_\alpha(f) \\
p_{\alpha+1}(f) &\;=\; p_1(p_\alpha(f))
\end{aligned}
\]
\[
\begin{tikzcd}
\underbrace{X}_{= \,\mathrm{source}\, f} \arrow[r, "i_\alpha(f)"] \arrow[rrr, bend left=20, "f"] & X_\alpha(f) \arrow[rr, "p_\alpha(f)"] \arrow[dr, bend right=40, "i_1(p_\alpha(f))"'] & & Y \\
& & \underbrace{X_1(p_\alpha(f))}_{= X_{\alpha+1}(f)} \arrow[ur, bend right=40] &
\end{tikzcd}
\]

\par

% ===================== PDF page 5 =====================
\printedpage{3}
% head: A. Grothendieck : Dérivateurs, Ch. XV, Édition des Universités de Montpellier II et Paris VII
% foot: Version 5 novembre 2010

ou mieux
\[
\begin{tikzcd}[row sep=large]
X \arrow[rrr, "f"] \arrow[d, "\mathrm{id}"] & & & Y \arrow[d, "\mathrm{id}"] \\
X \arrow[r, "i_\alpha(f)"] \arrow[d, "\mathrm{id}"] & X_\alpha(f) \arrow[rr, "p_\alpha(f)"] \arrow[d, "\mathrm{id}"] & & Y \arrow[d, "\mathrm{id}"] \\
X \arrow[r] & X_\alpha(f) \arrow[r, "i_1(p_\alpha(f))"] & \underbrace{X_1(p_\alpha(f))}_{= X_{\alpha+1}(f)} \arrow[r, "p_1(p_\alpha(f))"] & Y \;.
\end{tikzcd}
\tag*{(1.12)}
\]
Le morphisme de transition $X_\alpha(f) \longrightarrow X_{\alpha+1}(f)$ est
\[
X_\alpha(f) \;\; \xrightarrow{\;\;i_1(p_\alpha(f))\;\;} \;\; X_1(p_\alpha(f)) \;=\; X_{\alpha+1}(f) \;.
\]

\textbf{3.} Si $\alpha$ est un ordinal limite, on a
\[
\left\{
\begin{array}{rcl}
X_\alpha(f) & = & \varinjlim\limits_{\alpha'<\alpha} X_{\alpha'}(f). \\[3ex]
i_\alpha(f) & = & \parbox[t]{8cm}{morphisme canonique de $X = X_0(f)$ dans la limite inductive.} \\[5ex]
p_\alpha(f) & = & \text{limite projective des } p_{\alpha'}(f) : X_{\alpha'}(f) \longrightarrow Y.
\end{array}
\right.
\]
(On peut considérer la limite inductive définissant $X_\alpha(f)$ comme une limite dans $X\backslash\M/Y\dots$)

Les morphismes de transition
\[
X_{\alpha'}(f) \;\longrightarrow\; X_\alpha(f) \qquad (\; \alpha' \;<\; \alpha \;)
\]
sont les morphismes canoniques relatifs au système inductif des $X_{\alpha'}(f)$.

\msp{page 4}

On peut donc dire que la suite ordinale des factorisations fonctorielles $\Phi_\alpha(f)$ de $f$ variable, se déduit de la factorisation fonctorielle $\Phi_1(f)$ de $f$ en
\[
\begin{tikzcd}
X \arrow[r, "i_1(f)"] \arrow[rr, bend left=20, "f"] & X_1(f) \arrow[r, "p_1(f)"] & Y
\end{tikzcd}
\]
par “itération transfinie à droite”, \textit{i.e.} en regardant les opérations $\Phi_\alpha$ comme des foncteurs
% footnote mark (1) printed beside the following display; \footnotemark/\footnotetext used because \footnote does not work inside math
\[
\left\{
\begin{array}{rcl}
p_\alpha :\; \uFl(\M) & \longrightarrow & \uFl(\M) \\
f & \longmapsto & p_\alpha(f) \;,
\end{array}
\right.
\qquad\qquad \text{\footnotemark}
\]
\footnotetext{\textbf{N.B.} $p_\alpha$ respecte les structures de catégories \textit{sur} $\M$, via les foncteurs but : $\uFl(\M) \longrightarrow \M$, \textit{i.e.} $\uFl(\M)$ est considérée comme objet de $\mathrm{CAT}/\M$.}
munis d'homomorphismes fonctoriels (au-dessus de $\M$)
\[
\mathrm{id}_{\uFl(\M)} \;\longrightarrow\; p_\alpha \;.
\]


\endgroup
% END UNIT Der-ChXV-05-11-10/Der-ChXV-05-11-10_p001-005.tex
% BEGIN UNIT Der-ChXV-05-11-10/Der-ChXV-05-11-10_p006-010.tex
\begingroup\makeatletter%
\newcommand{\printedpage}[1]{\marginpar{\footnotesize p.~#1}}%
\newcommand{\msp}[1]{\marginpar{\footnotesize [#1]}}%
\newcommand{\cM}{\mathcal{M}}%
\newcommand{\Flu}{\underline{\mathrm{Fl}}}%
\newcommand{\Fl}{\mathrm{Fl}}%
\newcommand{\fU}{\mathfrak{U}}%
\newcommand{\fV}{\mathfrak{V}}%
\makeatother\let\printedpage\derpp\let\msp\dermsp%
\let\mspin\DERmspin%


\printedpage{4}
% head: A. Grothendieck : Dérivateurs, Ch. XV, Édition des Universités de Montpellier II et Paris VII
% foot: Version 5 novembre 2010

Ainsi, si on considère
\[
\left\{
\begin{array}{rl}
\Flu(\cM) & \text{comme objet de CAT}/\cM \text{ par le foncteur } f \longmapsto \text{but}\, f\\[1ex]
\text{où} & \text{CAT est la catégorie des catégories \textit{éléments de} } \fV\text{\textit{,} univers}\\
 & \text{tel que } \fU \in \fV\text{, d'où Cat } \overset{\text{(déf)}}{=} \mathrm{Cat}^{\fU},\ \cM,\ \Flu(\cM) \in \text{CAT} ,
\end{array}
\right.
\]
on a un endofoncteur $p_1$ de $\Flu(\cM)$ sur $\cM$, et un homomorphisme fonctoriel
\[
i_1 \ : \ \mathrm{id}_{\Flu(\cM)/\cM} \longrightarrow p_1 \ , \qquad \textit{i.e.} \quad
\begin{tikzcd}[column sep=small]
X \arrow[rr, "f"] \arrow[dr, "i_1(f)"'] & & Y \\
 & \underbrace{X_1(f)}_{=\,\text{source}\, p_1(f)} \arrow[ur, "p_1(f)"'] &
\end{tikzcd} \ ,
\tag*{(1.13)}
\]
et la description de la suite ordinale des endofoncteurs $p_\alpha$ de $\Flu\,\cM/\cM$ résulte de cette seule donnée, compte tenu que les petites limites filtrantes existent dans les fibres de $\Flu\,\cM \xrightarrow{\text{but}} \cM$.

\msp{page 5}

\[
\begin{tikzcd}[column sep=normal]
X \arrow[r, "i_1(f)"] \arrow[dr, "f"'] & X_1(f) \arrow[r, "i_1(p_1(f))"] \arrow[d, "p_1(f)"] & X_2(f) \arrow[r] \arrow[dl, "p_2(f)"] & \cdots \arrow[r] & X_\alpha(f) \arrow[r, "i_1(p_\alpha(f))"] \arrow[dlll, "p_\alpha(f)"] & X_{\alpha+1}(f) \arrow[r] \arrow[dllll, "p_{\alpha+1}(f) = p_1(p_\alpha(f))"] & \cdots \\
 & Y & & & & &
\end{tikzcd}
\tag*{(1.14)}
\]

Donc, pour achever de définir la suite ordinale des $X_\alpha(f)$ (munie des $i_\alpha(f)$, $p_\alpha(f)$), ou des $p_\alpha(f)$ (munie des $i_\alpha(f)$), il reste à décrire $X_1(f)$, $i_1(f)$, $p_1(f)$.

La propriété essentielle de ce triple dont nous aurons besoin est la suivante :

\begin{itemize}[leftmargin=5em, labelwidth=4em, labelsep=1em, align=left]
\item[(1.15)] \textbf{Fact}$(C_0)$ (propriété de factorisation par rapport à $C_0$).

Pour tout diagramme commutatif
\[
D : \qquad
\begin{tikzcd}
A \arrow[r] \arrow[d, "i"'] & X \arrow[d, "f"] \\
B \arrow[r] & Y
\end{tikzcd}
\qquad \text{avec } i \in C_0 \ ,
\]
d'où un diagramme traits pleins
\[
\begin{tikzcd}
A \arrow[r] \arrow[d, "i"'] & X \arrow[d, "i_D"] \arrow[dr, dotted, "i_1(f)"] \arrow[dd, dotted, bend right=60, "f"'] & \\
B \arrow[r] & X_D \overset{\text{déf}}{=} X \amalg_A B \arrow[r, dotted, "u_D"] \arrow[d, "p_D"] & X_1(f) \arrow[dl, dotted, "p_1(f)"] \\
 & Y \ , &
\end{tikzcd}
\]
\printedpage{5}
% head: A. Grothendieck : Dérivateurs, Ch. XV, Édition des Universités de Montpellier II et Paris VII
% foot: Version 5 novembre 2010
il existe un homomorphisme de la factorisation $f = p_D i_D$ de $f$ dans la factorisation $f = p_1(f)i_1(f)$.
\end{itemize}

Pour formuler cette condition, j'ai utilisé l'existence des sommes amalgamées $X_D$, \textit{i.e.} la condition sur $C_0$ :

\msp{page 6}

\begin{itemize}[leftmargin=5em, labelwidth=4em, labelsep=1em, align=left]
\item[(1.16)] Les flèches $i \in C_0$ sont coquarrables, \textit{i.e.} pour un diagramme
\[
\begin{tikzcd}
A \arrow[r, "u"] \arrow[d, "i"'] & X \\
B &
\end{tikzcd}
\]
avec $i \in C_0$, la somme amalgamée existe.
\end{itemize}

Mais on peut formuler 1.15 aussi ainsi, indépendamment de l'hypothèse (1.16) :

\begin{itemize}[leftmargin=5em, labelwidth=4em, labelsep=1em, align=left]
\item[(1.17)] Pour tout carré commutatif
\[
\begin{tikzcd}
A \arrow[r, "u"] \arrow[d, "i"'] & X \arrow[d, "f"] \\
B \arrow[r, "v"] & Y
\end{tikzcd}
\qquad ,
\]
avec $i \in C_0$, il existe un carré commutatif traits pleins
\[
\begin{tikzcd}
A \arrow[r, "u"] \arrow[d, "i"'] & X \arrow[d, "i_1(f)"] \\
B \arrow[r, "v_1"] \arrow[dr, dotted, "v"] & X_1(f) \arrow[d, dotted, "p_1(f)"] \\
 & Y
\end{tikzcd}
\]
tel que $v = p_1(f)v_1$.
\end{itemize}

Ceci a pour conséquence le

\textbf{Corollaire 1} (de 1.17). \textit{Soit} $f \in \Fl(\cM)$\textit{, et} $\alpha$ \textit{un ordinal} $(\in \fU)$. \textit{Pour tout carré commutatif}
\[
\begin{tikzcd}
A \arrow[r, "u_\alpha"] \arrow[d, "i"'] & X_\alpha(f) \arrow[d, "p_\alpha(f)"] \\
B \arrow[r, "v"] & Y
\end{tikzcd}
\]
\printedpage{6}
% head: A. Grothendieck : Dérivateurs, Ch. XV, Édition des Universités de Montpellier II et Paris VII
% foot: Version 5 novembre 2010
\textit{il existe un carré commutatif}
\[
\begin{tikzcd}
A \arrow[r, "u_\alpha"] \arrow[d, "i"'] & X_\alpha(f) \arrow[d, "\text{trans.}"] \arrow[dd, dotted, bend left=60, "p_\alpha(f)"] \\
B \arrow[r, "v_1"] \arrow[dr, dotted, "v"] & X_{\alpha+1}(f) \arrow[d, dotted, "p_{\alpha+1}(f)"] \\
 & Y
\end{tikzcd}
\]
\textit{tel que}
\[
v \ = \ p_{\alpha+1}(f)\ v_1 \ .
\]

\msp{page 7}

\textbf{Corollaire 2.} \textit{Soit $\pi$ un cardinal} $(\in \fU)$ \textit{tel que pour toute flèche $i : A \longrightarrow B$ dans $C_0$, la source $A$ soit $\pi$-accessible. Soit $\alpha$ un ordinal grand devant $\pi$} (i.e. \textit{tel que l'ensemble ordonné $I_\alpha$ des ordinaux $\alpha' < \alpha$ soit grand devant $\pi$}) -- \textit{p. ex. $\alpha$ le premier ordinal dont le cardinal soit $> \pi$. Alors pour toute flèche $f \in \Fl(\cM)$, on a $p_\alpha(f) \in (C_0)_*$.}

\textsc{Démonstration}. Cela signifie que pour tout carré commutatif
\[
\begin{tikzcd}
A \arrow[r, "u"] \arrow[d, "i"'] & X_\alpha(f) \arrow[d, "p_\alpha(f)"] \\
B \arrow[r, "v"] \arrow[ur, dotted, "w"] & Y \ ,
\end{tikzcd}
\tag*{$(*)$}
\]
avec $i \in C_0$, la flèche pointillée $w$ (avec $wi = u$, $p_\alpha(f)w = v$) existe. Or
\[
X_\alpha(f) = \varinjlim_{\alpha' \in I_\alpha} X_{\alpha'}(f) \ ,
\]
donc $A$ étant $\pi$-accessible et $I_\alpha$ grand devant $\pi$, $u$ provient d'une flèche $u_{\alpha'} : A \longrightarrow X_{\alpha'}(f)$, et (quitte à remplacer $\alpha'$ par $\alpha'' > \alpha'$ et $u_{\alpha'}$ par le composé $u_{\alpha''} : A \longrightarrow X_{\alpha'}(f) \longrightarrow X_{\alpha''}(f)$) on peut supposer $p_{\alpha'}(f)u_{\alpha'} = vi$, \textit{i.e.} que $(*)$ provient d'un carré commutatif
\[
\begin{tikzcd}
A \arrow[r, "u_{\alpha'}"] \arrow[d, "i"'] \arrow[rr, dotted, bend left=30, "u"] & X_{\alpha'}(f) \arrow[r, dotted, "\text{trans.}"'] \arrow[d, "p_{\alpha'}(f)"'] & X_\alpha(f) \arrow[dl, dotted, "p_\alpha(f)"] \\
B \arrow[r, "v"] & Y \ , &
\end{tikzcd}
\]

\printedpage{7}
% head: A. Grothendieck : Dérivateurs, Ch. XV, Édition des Universités de Montpellier II et Paris VII
% foot: Version 5 novembre 2010
et appliquant le corollaire 1, on trouve
\[
\begin{tikzcd}
A \arrow[r, "u_{\alpha'}"] \arrow[d, "i"'] & X_{\alpha'}(f) \arrow[d, "\text{trans.}"] & \\
B \arrow[r, "v_1"] \arrow[dr, dotted, "v"] & X_{\alpha'+1}(f) \arrow[r, dotted, "\text{trans.}"] \arrow[d, "p_{\alpha+1}(f)"] & X_\alpha(f) \arrow[dl, dotted, "p_\alpha(f)"] \\ % [sic] p_{\alpha+1}(f) -> p_{\alpha'+1}(f), the arrow starts at X_{\alpha'+1}(f); ?? uncertain: subscript of the label, tightly set -- \alpha+1 or \alpha'+1
 & Y &
\end{tikzcd}
\]
avec $v = p_{\alpha+1}(f)v_1$, et on prendra % [sic] p_{\alpha+1}(f)v_1 -> p_{\alpha'+1}(f)v_1, corollaire 1 is applied to \alpha' (arrow from X_{\alpha'+1}(f))
\[
w \ = \ \mathrm{trans}_{\alpha,\alpha'+1} \circ v_1 \ .
\]

\msp{page 8}

Mais on veut une factorisation de $f$, flèche quelconque dans $\cM$, en
\[
f \ = \ p\ i \ , \qquad \text{avec} \ \ p \in (C_0)_* \ \ \textit{et} \ \ i \in \widetilde{C_0} \ ,
\]
donc il faut encore énoncer une propriété, assurant que les $i_\alpha(f)$ sont dans $\widetilde{C_0}$. On aura même quelque chose de plus précis. Introduisons

\begin{itemize}[leftmargin=10em, labelwidth=9em, labelsep=1em, align=left]
\item[(1.18)\hfill $C'_0 \ = $] plus petit sous-ensemble de $\Fl(\cM)$ contenant $C_0$, stable par cochangement de base (chaque fois qu'il existe).
\item[\hfill $C''_0 \ = $] ensemble des flèches qui sont soit dans $C'_0$, soit des sommes amalgamées d'une famille \textit{non vide} de telles flèches $A \xrightarrow{i_\alpha} B_\alpha$ de même source $A$, $i_\alpha \in C'_0$ ; la flèche $A \longrightarrow B = \displaystyle\coprod_{\alpha}{}_{(A)} B_\alpha$ est dans $C''_0$.
\end{itemize}

(La condition 1.16 implique que les éléments de $C'_0$ sont coquarrables, et cela suffit, avec la condition de stabilité de $\cM$ par $\varinjlim$ filtrantes, à assurer que les sommes amalgamées infinies envisagées existent dans $\cM$.)

\textbf{N.B.} $C''_0$ est \textit{égal} au plus petit ensemble de flèches contenant $C_0$ et stable par cochangement de base et par petites sommes amalgamées. \texttt{[Si les petites sommes existent dans $\cM$, alors $C''_0$ est aussi l'ensemble des flèches obtenues par cochangement de base à partir de sommes de flèches de $C_0$.]}

\setcounter{footnote}{1}% the edition numbers these footnotes 2 and 3
\begin{itemize}[leftmargin=10em, labelwidth=9em, labelsep=1em, align=left]
\item[(1.19)\hfill $C'''_0 \ = $] plus petit sous-ensemble de $\Fl(\cM)$ contenant $C''_0$, et stable de plus par limites inductives ordinales strictes, \textit{i.e.} pour un système inductif ordinal $(A_\alpha)_{\alpha<\alpha_0}$, tel que pour $\alpha$ ordinal limite on ait $A_\alpha \xleftarrow{\ \sim\ } \varinjlim\limits_{\alpha' \in I_\alpha} A_{\alpha'}$, si les morphismes de transition $i_{\beta,\alpha} : A_\alpha \longrightarrow A_\beta$ sont dans $C'''_0$, de même les morphismes canoniques
\[
A_\alpha \xrightarrow{\ i_\alpha\ } \varinjlim_{\beta \in I_{\alpha_0}} A_\beta \ . \quad \footnotemark
\]
\end{itemize}
\footnotetext{$C_0 \subset C'_0 \subset C''_0 \subset C'''_0$}% footnote 2: mark printed in the display, text printed at the foot of p. 8

\printedpage{8}
% head: A. Grothendieck : Dérivateurs, Ch. XV, Édition des Universités de Montpellier II et Paris VII
% foot: Version 5 novembre 2010
\msp{page 9}

La proposition 2, p. 4 de XIII assure que l'on a
\[
C'''_0 \ \subset \ \widetilde{C_0} \ .
\]
Ceci posé, voici la deuxième propriété essentielle de la factorisation $f = p_1(f)i_1(f)$ (la première étant 1.17) :

\begin{itemize}[leftmargin=5em, labelwidth=4em, labelsep=1em, align=left]
\item[(1.20)] \qquad\qquad Pour toute flèche $f \in \Fl(\cM)$, on a $i_1(f) \in C''_0$.
\end{itemize}

Il en résulte aussitôt, par récurrence ordinale :

\begin{itemize}[leftmargin=5em, labelwidth=4em, labelsep=1em, align=left]
\item[(1.21)] \qquad Pour toute flèche $f \in \Fl(\cM)$ et tout ordinal $\alpha \in \fU$, on a $i_\alpha(f) \in C'''_0$, a fortiori $i_\alpha(f) \in \widetilde{C_0}$.
\end{itemize}

Donc si $\alpha$ est grand devant $\pi$ (\textit{cf.} cor. 2, page 7), la factorisation $f = p_\alpha(f)i_\alpha(f)$ satisfait
\[
\left\{
\begin{array}{l}
f = p_\alpha(f)i_\alpha(f)\text{, avec } i_\alpha(f) \in C'''_0 \subset \widetilde{C_0} = \Phi \text{ \texttt{[à ne pas}}\\
\text{\texttt{confondre avec }}\Phi_\alpha\text{\texttt{]}, } p_\alpha(f) \in C_{0*} = \Psi \text{, lorsque } \alpha \text{ est}\\
\text{grand devant } \pi\text{, les sources des } i \in C_0 \ \pi\text{-accessibles.}
\end{array}
\right.
\tag*{(1.22)}
\]

Il reste à construire le foncteur $f \longmapsto X_1(f)$, avec $i_1(f)$, $p_1(f)$, satisfaisant les conditions (1.17) et (1.20), sous l'hypothèse (1.16) que les $i \in C_0$ sont coquarrables. On considère alors, pour $f \in \Fl(\cM)$ fixé, l'ensemble $\Delta(f)$ (ou $\Delta(C_0, f)$) de \textit{tous} les diagrammes $D$ de la forme (1.15), avec $i \in C_0$, et on pose
\[
X_1(f) \ = \ \coprod_{D \in \Delta(f)}{}_{(X)} X_D
\tag*{(1.23)}
\]
(somme amalgamée sous $X$, \textit{i.e.} somme indexée
\msp{page 10}
par $\Delta(f)$ dans $X\backslash\cM$), le morphisme $X \xrightarrow{i_1(f)} X_1(f)$ étant le morphisme costructural, et le morphisme $X_1(f) \xrightarrow{p_1(f)} Y$ étant celui ayant pour composantes les $p_D$. En somme, on fait “au plus économique” (et au plus naturel) pour assurer la validité de (1.17), sous la forme équivalente (1.15). D'autre part, par définition de $C'_0$, $C''_0$, les $i_D$ sont dans $C'_0$, et $i_1(f)$ est dans $C''_0$, ce qui est (1.20)\footnote{Il faudrait quand même expliciter la dépendance fonctorielle de $X_1(f)$ par rapport à $f$\,!!!}.

Pour résumer :

\textbf{Théorème 1.} \textit{Soient $\cM$ une $\fU$-catégorie stable par $\fU$-limites inductives filtrantes, $C_0 \subset \Fl(\cM)$ un petit ensemble de flèches de $\cM$. On suppose les flèches $i \in C_0$ coquarrables. Cela assure la définition d'un endofoncteur $f \longmapsto p_1(f)$ de $\Flu(\cM)$, et d'un homomorphisme fonctoriel $i_1 : \mathrm{id}_{\Flu(\cM)} \longrightarrow p_1$,} i.e. \textit{$i_1(f) : f \longrightarrow p_1(f)$, induisant l'identité sur les objets buts de $f$ et de $p_1(f)$,}
\[
\begin{tikzcd}[column sep=small]
X \arrow[rr, "i_1(f)"] \arrow[dr, "f"'] & & X_1(f) \arrow[dl, "p_1(f)"] \\
 & Y \ , &
\end{tikzcd}
\]


\endgroup
% END UNIT Der-ChXV-05-11-10/Der-ChXV-05-11-10_p006-010.tex
% BEGIN UNIT Der-ChXV-05-11-10/Der-ChXV-05-11-10_p011-015.tex
\begingroup\makeatletter%
\newcommand{\printedpage}[1]{\marginpar{\footnotesize p.~#1}}%
\newcommand{\msp}[1]{\marginpar{\footnotesize [#1]}}%
\newcommand{\uFl}{\underline{\mathrm{Fl}}}%
\newcommand{\Fl}{\mathrm{Fl}}%
\newcommand{\Filt}{\mathrm{Filt}}%
\newcommand{\cM}{\mathcal{M}}%
\newcommand{\card}{\operatorname{card}}%
\newcommand{\circled}[1]{\tikz[baseline=(c.base)]\node[draw,circle,inner sep=0.5pt](c){#1};}%
\newcommand{\pfn}[2]{\textsuperscript{(}\footnote[#1]{#2}\textsuperscript{)}}%
\makeatother\let\printedpage\derpp\let\msp\dermsp%
\let\mspin\DERmspin%


\printedpage{9}
% head: A. Grothendieck : Dérivateurs, Ch. XV, Édition des Universités de Montpellier II et Paris VII
% foot: Version 5 novembre 2010

{\itshape d'où par itération un système inductif ordinal d'endofoncteurs $f\longmapsto p_\alpha(f)$ de $\uFl(\cM)$, et d'homomorphismes $i_\alpha : \mathrm{id}_{\uFl(\cM)} \longrightarrow p_\alpha$, \textup{i.e.}}
\[
\begin{tikzcd}[column sep=small]
X \arrow[rr, "i_\alpha(f)"] \arrow[dr, "f"'] & & X_\alpha(f) \arrow[dl, "p_\alpha(f)"] \\
& Y\,, &
\end{tikzcd}
\]
{\itshape formant un système inductif ordinal de factorisations de $f$ fonctorielles en $f$.}

\msp{page 11}

{\itshape Ceci posé, on a :}
\begin{enumerate}[label={}, leftmargin=*]
\item[a)] {\itshape Pour tout ordinal $\alpha$, on a $i_\alpha(f)\in\Phi$.}
\item[b)] {\itshape Supposons que les sources des $i\in C_0$ soient accessibles et soit alors $\pi$ un cardinal dans $\mathfrak{U}$ tel qu'elles soient $\pi$-accessibles. Alors pour tout nombre ordinal $\alpha$ grand devant $\pi$, on a $p_\alpha(f)\in C_{0*}$, donc}
\[
f \;=\; p_\alpha(f)\, i_\alpha(f) \qquad\quad \textit{avec} \qquad\quad i_\alpha(f)\in\Phi\,, \qquad\quad p_\alpha(f)\in\Psi\,.
\]
\end{enumerate}

On a un résultat légèrement plus précis, que voici :

\textbf{Corollaire 1.} \textit{Soit $\Phi_0$ une partie de} $\Fl(\cM)$ \textit{telle que}
\[
C_0''' \;\subset\; \Phi_0 \;\subset\; \Phi \qquad\quad (\textit{d'où}\ \ C_0\subset\Phi_0\subset\Phi)\,,
\tag*{(1.24)}
\]
\textit{où $C_0'''\subset\Phi$ est définie dans} 1.19, \textit{p.} 8. \textit{Par exemple, si $\Phi_0$ est une partie de} $\Fl(\cM)$ \textit{satisfaisant}
\[
C_0 \;\subset\; \Phi_0 \;\subset\; \Phi\,,
\]
\textit{et stable par cochangement de base, par petites sommes amalgamées} (\textit{cf.} (1.18), p.~8), \textit{par $\varinjlim$ ordinales strictes} (\textit{cf.} (1.19), p.~8), \textit{elle satisfait} (1.24). (\textbf{N.B.} \textit{Si $\cM$ admet des sommes amalgamées simples, de tels $\Phi_0$ existent trivialement, il suffit de prendre $\Phi_0=\Phi$.})

\msp{page 12}

\textit{Ceci posé, pour toute flèche $f$ dans $\cM$, et tout ordinal $\alpha$, on a}
\[
i_\alpha(f) \;\in\; \Phi_0\,.
\]

\textbf{Corollaire 2.} \textit{Supposons que les sources des flèches $i\in C_0$ soient accessibles, et soit $\Phi_0$ comme dans le corollaire} 1 (p.~ex. $\Phi_0=C_0'''$). \textit{Alors le couple $(\Phi_0,\Psi)$ est un précouple de Quillen,} i.e.
\begin{enumerate}[label={}, leftmargin=*]
\item[a)] $\Phi_0\longleftrightarrow\Psi$, i.e. $\Psi\subset\Phi_{0*}$ (\textit{ou ce qui revient au même, $\Phi_0\subset\Psi^*$}), \textit{et}
\item[b)] \textit{on a la propriété de factorisation de Quillen.}
\end{enumerate}
\textit{Cela implique donc que $\Phi$ $(=\widetilde{C_0}=\widetilde{\Phi_0})$ est formé des flèches qui sont facteurs directs de flèches de $\Phi_0$. Donc si $\Phi_0$ est stable par facteurs directs, on a $\Phi_0=\Phi$,} i.e. \textit{$\Phi_0$ est clos à gauche.}

\printedpage{10}
% head: A. Grothendieck : Dérivateurs, Ch. XV, Édition des Universités de Montpellier II et Paris VII
% foot: Version 5 novembre 2010

\section*{2\quad Variation des factorisations $f=p_\alpha i_\alpha$ avec la partie $C_0$ génératrice}

Je dis que le système inductif ordinal de foncteurs $p_\alpha=p_\alpha^{C_0} : \uFl(\cM)\longrightarrow\uFl(\cM)$,
\msp{page 13}
muni des flèches
\[
i_\alpha^{C_0} \;:\; \mathrm{id}_{\uFl(\cM)} \longrightarrow p_\alpha^{C_0}\,,
\]
dépend fonctoriellement de $C_0$, quand on considère $C_0$ comme un élément de l'ensemble ordonné des \textit{petites} parties de $\Fl(\cM)$, formées de flèches coquarrables (donc de l'ensemble ordonné
\[
\mathfrak{P}_{\mathfrak{U}\text{-petites}}(\Fl_{\text{coquar.}}(\cM)) \quad ).
\]
Si on a une inclusion
\[
C_0 \;\hookrightarrow\; C_1\,,
\tag*{(2.1)}
\]
on en déduit pour tout ordinal $\alpha$ un homomorphisme fonctoriel
\[
u_\alpha \;=\; u_\alpha^{C_1,C_0} \;:\; X_\alpha^{C_0} \;\longrightarrow\; X_\alpha^{C_1}\,,
\tag*{(2.2)}
\]
rendant commutatifs les carrés déduits de $f : X\longrightarrow Y$ dans $\cM$,
\[
\begin{tikzcd}
X \arrow[r, "i_\alpha^{C_0}(f)"] \arrow[d, "\mathrm{id}"'] & X_\alpha^{C_0}(f) \arrow[r, "p_\alpha^{C_0}(f)"] \arrow[d, "u_\alpha(f)"] & Y \arrow[d, "\mathrm{id}"] \\
X \arrow[r, "i_\alpha^{C_1}(f)"] & X_\alpha^{C_1}(f) \arrow[r, "p_\alpha^{C_1}(f)"] & Y\,,
\end{tikzcd}
\tag*{(2.3)}
\]
et compatible avec les morphismes de transition pour $\alpha$ variable,
\[
\begin{tikzcd}
X_\alpha^{C_0} \arrow[r, "{\mathrm{trans}_{\beta,\alpha}^{C_0}}"] \arrow[d, "u_\alpha"'] & X_\beta^{C_0} \arrow[d, "u_\beta"] \\
X_\alpha^{C_1} \arrow[r, "{\mathrm{trans}_{\beta,\alpha}^{C_1}}"] & X_\beta^{C_1}\,.
\end{tikzcd}
\tag*{(2.4)}
\]

\msp{page 14}

Pour définir les $u_\alpha$, satisfaisant aux conditions précédentes, il suffit de les définir pour $\alpha=1$ ``et d'itérer'' :
\[
\begin{tikzcd}[column sep=small]
& X_1^{C_0}(f) \arrow[dd, "u_1(f)"] \arrow[dr, "p_1^{C_0}(f)"] & \\
X \arrow[ur, "i_1^{C_0}(f)"] \arrow[dr, "i_1^{C_1}(f)"'] & & Y \\
& X_1^{C_1}(f) \arrow[ur, "p_1^{C_1}(f)"'] & {}.
\end{tikzcd}
\tag*{(2.5)}
\]
Or $X_1^{C_0}(f)$ est défini comme une somme directe dans $X\backslash\cM/Y$, relative à un ensemble d'indices $\Delta(C_0,f)$ \textit{contenu} dans $\Delta(C_1,f)$ -- et la famille d'objets dont on prend la somme directe étant induite d'une famille (non petite) dont l'ensemble d'indices serait $\Delta(\Fl_{\text{coquar.}}(\cM),f)$, savoir formé de \textit{tous} les diagrammes carrés
\printedpage{11}
% head: A. Grothendieck : Dérivateurs, Ch. XV, Édition des Universités de Montpellier II et Paris VII
% foot: Version 5 novembre 2010
\[
D : \qquad
\begin{tikzcd}
A \arrow[r, "u"] \arrow[d, "i"'] & X \arrow[d, "f"] \\
B \arrow[r, "v"] & Y
\end{tikzcd}
\qquad\qquad \text{avec $i$ coquarrable sans plus ,}
\]
la définition des $X_D$, $i_D$, $p_D$ associés à $D$ ne dépendant que de $D$, et non de l'ensemble $C_0$ ou $C_1$ de flèches coquarrables qui contient $i$. Ceci dit, la flèche $u_1(f)$
\msp{page 15}
dans (2.5) n'est autre que l'inclusion canonique d'une somme amalgamée partielle, indexée par $\Delta(C_0,f)$, dans la somme amalgamée correspondante à l'ensemble d'indices plus grand. La fonctorialité de cette flèche $u_1(f)$, pour $f$ variable, ``est immédiate'', elle résulte de la commutativité du carré, associé à un morphisme $v=(v_X,v_Y)$ de flèches
\[
\begin{tikzcd}
X \arrow[r, "f"] \arrow[d, "v_X"'] & Y \arrow[d, "v_Y"] \\
X' \arrow[r, "f'"] & Y'\,,
\end{tikzcd}
\tag*{(2.6)}
\]
savoir le diagramme d'applications entre ensembles d'indices
\[
\begin{tikzcd}
\Delta(C_0,f) \arrow[r, hook, "u_1(f)"] \arrow[d, "{\begin{array}{r}\text{dépendance fonctorielle}\\ \text{en $f$ de $\Delta(C_0,f)$}\end{array}}"'] & \Delta(C_1,f) \arrow[d, "{\begin{array}{l}\text{dépendance fonctorielle}\\ \text{en $f$ de $\Delta(C_1,f)$}\end{array}}"] \\
\Delta(C_0,f') \arrow[r, hook, "u_1(f')"] & \Delta(C_1,f')
\end{tikzcd}
\tag*{(2.7)}
\]
(commutativité qui revient à dire que $\Delta(C_0,f)$ peut être considéré comme un foncteur en le couple $(C_0,f)$ comme objet de la catégorie produit $\underline{\mathfrak{P}_{\text{pet.}}(\Fl_{\text{coquar.}}(\cM))}\times\uFl(\cM)$ \ldots). Le carré commutatif (2.7) est la partie ``indices'' d'un carré commutatif entre familles d'objets $\mathscr{X}(C_0,f)$, $\mathscr{X}(C_0,f')$, $\mathscr{X}(C_1,f)$, $\mathscr{X}(C_1,f')$,
\msp{page 16}
indexés par ces ensembles, savoir les $X_D^{C_0}$ et les $X_D^{C_1}$ relatifs aux diagrammes $D$ impliquant, soit la flèche $f$, soit la flèche $f'$ \ldots

Or la construction canonique du système inductif ordinal $p_\alpha$, avec les $i_\alpha : \mathrm{id}_{\uFl(\cM)}\longrightarrow p_\alpha$, est fonctorielle par rapport à la donnée initiale $p_1$, $i_1$ -- je n'explicite pas le détail de la fonctorialité, bien que je devrais \ldots

\section*{3\quad Majoration des cardinaux}

Je suppose à présent que $\cM$ est muni d'une \textit{filtration cardinale}
\[
(\Filt^\pi)_{\pi\geq\pi_0} \qquad\quad (\pi_0,\ \pi \text{ cardinaux dans } \mathfrak{U})
\tag*{(3.1)}
\]
(SGA 4 I 9.12, p.~149), les $\Filt^\pi$ stables par sommes amalgamées simples (chaque fois qu'elles existent), donc aussi par sommes amalgamées infinies indexées par des ensembles
\printedpage{12}
% head: A. Grothendieck : Dérivateurs, Ch. XV, Édition des Universités de Montpellier II et Paris VII
% foot: Version 5 novembre 2010
d'indices de cardinal $\leq\pi$. On suppose de plus que tout objet de $\Filt^{\pi_0}$ soit accessible, et on désigne par \circled{$\pi_1$} le plus petit cardinal tel que $\pi_1\geq\pi_0$ et que les objets de $\Filt^{\pi_0}$ soient $\pi_1$-accessibles. Désignant, pour tout cardinal, par $\cM_\pi$ le sous-ensemble de $\mathrm{Ob}\,\cM$ formé des objets $\pi$-accessibles, on aura alors pour tout $\pi\geq\pi_1$ (\textit{loc. cit.} 9.16)
\[
\cM_\pi \;\subset\; \Filt^\pi \;\subset\; \cM_{(\pi^{\pi_0})} \qquad\quad (\text{si } \pi\geq\pi_1)\,,
\tag*{(3.2)}
\]

\msp{page 17}

donc si $\pi^{\pi_0}=\pi$, p.~ex. si $\pi=2^c$ avec $c\geq\pi_0$, on aura
\[
\Filt^\pi \;=\; \cM_\pi\,.
\tag*{(3.2 bis)}
\]
(Si on admet l'hypothèse du continu généralisée, on a $\pi^{\pi_0}=\pi$ pour tout $\pi>\pi_0$, donc $\Filt^\pi=\cM_\pi$ si $\pi\geq\pi_1$ et $\pi>\pi_0$.)

Je pose, pour tout cardinal $\pi\geq\pi_0$,
\[
\Fl^\pi(\cM) \;=\; \{\ f\in\Fl(\cM) \mid \text{source et but de $f$ sont dans } \Filt^\pi\ \}
\tag*{(3.3)}
\]
-- ce qui sans doute donne une filtration cardinale sur $\uFl(\cM)$ (mais ce n'est pas là notre souci \ldots). Mon but à présent, c'est de dégager des conditions sur $C_0\subset\Fl_{\text{coquar.}}(\cM)$ (petite partie de $\Fl(\cM)$) et sur le cardinal $\pi$, et l'ordinal $\alpha$, pour que l'on ait
\[
\begin{aligned}
f\in\Fl^\pi(\cM) \;\;\overset{?}{\Longrightarrow}\;\; & X_\alpha(f)\in\Filt^\pi \\
& (\textit{i.e.}\ p_\alpha(f)\in\Fl^\pi(\cM),\ \text{ou encore } i_\alpha(f)\in\Fl^\pi(\cM))\,.
\end{aligned}
\tag*{(3.4)}
\]
Plus généralement, je me propose de donner une estimation de la ``cardinalité'' de $X_\alpha(f)$, pour $C_0$ fixé, en fonction de $C_0$, de $f$, de $\alpha$, plus exactement, en fonction des trois cardinaux $\pi(C_0)$, $\pi(f)$, $\pi(\alpha)$ définis comme
\[
\left\{
\begin{array}{rcl}
\pi(C_0) & = & \text{plus petit cardinal tel que } C_0\subset\Fl^\pi(\cM) \\[1ex]
\pi(f) & = & \pi(\{f\}) \;=\; \text{plus petit cardinal tel que } f\in\Fl^\pi(\cM) \\[1ex]
\pi(\alpha) & = & \card\alpha\,.
\end{array}
\right.
\tag*{(3.5)}
\]
Donc je cherche une fonction
\msp{page 18}
cardinale de variables cardinales $\varphi$, telle que l'on ait
\[
X_\alpha(f) \;\in\; \Filt^\pi \qquad\quad \text{si} \qquad\quad \pi \;=\; \varphi(\underbrace{\pi(C_0)}_{r},\underbrace{\pi(f)}_{s},\underbrace{\pi(\alpha)}_{t})\,.
\tag*{(3.6)}
\]
L'estimation cruciale concernera sûrement le cas $\alpha=1$. Notons que $X_1^{C_0}(f)=X_1(f)$ est défini comme une somme amalgamée (en général infinie) d'objets $X_D$, pour un ensemble $\Delta(C_0,f)=\Delta(f)$ de diagrammes, savoir les diagrammes commutatifs
\[
\begin{tikzcd}
A \arrow[r] \arrow[d, "i"'] & X \arrow[d, "f"] \\
B \arrow[r] & Y
\end{tikzcd}
\]
\printedpage{13}
% head: A. Grothendieck : Dérivateurs, Ch. XV, Édition des Universités de Montpellier II et Paris VII
% foot: Version 5 novembre 2010
avec $i\in C_0$. Pour un $i$ fixé, et un $f$ fixé, le cardinal de l'ensemble des diagrammes de ce type est majoré par celui de $\mathrm{Hom}(A,X)\times\mathrm{Hom}(B,Y)$. On a $A,B\in\Filt^r$, $X,Y\in\Filt^s$, donc il nous faut une majoration de
\[
\card\mathrm{Hom}_{\cM}(A,X)\,, \qquad\quad A\;\in\;\Filt^r\,, \qquad\quad X\;\in\;\Filt^s\,.
\]

\msp{page 19}

Voici les majorations essentielles dans ce sens :

\textbf{Proposition 3.1.} \textit{Soit $\pi_0$ le cardinal initial de la filtration cardinale} $(\Filt^\pi)_{\pi\geq\pi_0}$ \textit{envisagé, $\pi_1$ le plus petit cardinal $\geq\pi_0$ tel que les objets de} $\Filt^{\pi_0}$ \textit{soient $\pi_1$-accessibles, $C_1=$} card $\underbrace{\widetilde{\uFl^{\pi_1}(\cM)}}_{=\uFl(\cM^{\pi_1})}$\pfn{4}{\textit{cf.} plus bas la notation $\widetilde{\mathcal{C}}$.}. \textit{Pour tout cardinal $\pi\geq\pi_0$, je désigne par $\cM^\pi$ la sous-catégorie strictement pleine de $\cM$ dont l'ensemble des objets est} $\Filt^\pi$. \textit{Enfin, si $\mathcal{C}$ est une catégorie, je désigne par $\widetilde{\mathcal{C}}$ l'ensemble des classes d'isomorphie d'objets de $\mathcal{C}$.}\pfn{5}{Signalons qu'on a toujours $\card\widetilde{\uFl(\mathcal{C})}\geq\card\widetilde{\mathcal{C}}$ car $\widetilde{\mathcal{C}}\hookrightarrow\widetilde{\uFl(\mathcal{C})}$ via $\mathrm{cl}(X)\longmapsto\mathrm{cl}(\mathrm{id}_X)$.}
% ?? uncertain: straight underlines under "le plus petit" and under "pleine" -- underlined / not underlined (a wavy tilde and an underbrace stand close by)

$1^\circ$) \textit{Si $r\geq\sup(\pi_1,C_1)$, on a}
\[
\card\,(\underbrace{\widetilde{\uFl^r(\cM)}}_{=\uFl(\cM^r)}) \;\leq\; 2^{(r^{\pi_0})} \quad\ (=2^r,\ \textit{si } r=2^c,\ \textit{avec } c\geq\pi_0)\,.
\tag*{(3.7)}
\]

$2^\circ$) \textit{Si $r,s\geq\sup(\pi_1,C_1)$, et si $A\in\Filt^r$, $X\in\Filt^s$, on a}
% (3.8) and (3.9) share one brace on the right
\[
\begin{array}{l}
\text{(3.8)}\\[3ex]
\text{(3.9)}
\end{array}
\qquad\qquad
\begin{array}{r@{\;\;\leq\;\;}l}
\card\mathrm{Hom}_{\cM}(A,X) & s^{(r^{\pi_0})}\\[3ex]
\card\widetilde{\uFl(\cM^r/X)} & s^{(r^{\pi_0})}
\end{array}
\qquad\quad
\left\{
\begin{array}{lll}
= & s^r & \textit{si } r=2^c,\ \textit{avec } c\geq\pi_0\,,\\[3ex]
= & s & \textit{si } s=2^{c'},\ \textit{avec } c'\geq r^{\pi_0}\,.
\end{array}
\right.
\]

\textbf{N.B.} Si on admet l'hypothèse du continu généralisée, qui implique que $\pi^{\pi'}=\pi$ si $\pi>\pi'$ ($\pi$ cardinal infini, $\pi'\geq1$), alors prenant $r$ tel que $r>\pi_0$, les formules précédentes donnent (avec de plus $r\geq\sup(\pi_1,C_1)$) :
\[
\card\widetilde{\uFl^r(\cM)} \;\;\leq\;\; 2^r\,,
\tag*{(3.7 bis)}
\]
\[
\card\mathrm{Hom}_{\cM}(A,X) \;\leq\; s^r \qquad\quad (\text{si } A\in\Filt^r,\ X\in\Filt^s)\,,
\tag*{(3.8 bis)}
\]
\[
\card\widetilde{\uFl(\cM^r/X)} \;\leq\; s^r \qquad\quad (\text{si } X\in\Filt^s)\,.
\tag*{(3.9 bis)}
\]
Sans hypothèse du continu, ces formules sont aussi valables pour $r$ de la forme $2^c$, avec $c\geq\pi_0$. Ce sont ces dernières formules qui on l'air les plus naturelles. Pour (3.8 bis) et % [sic] qui on l'air -> qui ont l'air, verb form
(3.9 bis), je signale que le deuxième membre $s^r$ est le cardinal de l'ensemble $\mathrm{Hom}_{\mathrm{Ens}}(A,X)$, où $A$, $X$ sont deux ensembles tels que $\card A=r$, $\card X=s$. On ne peut donc espérer améliorer (3.8 bis)\,! Pour la


\endgroup
% END UNIT Der-ChXV-05-11-10/Der-ChXV-05-11-10_p011-015.tex
% BEGIN UNIT Der-ChXV-05-11-10/Der-ChXV-05-11-10_p016-020.tex
\begingroup\makeatletter%
\newcommand{\printedpage}[1]{\marginpar{\footnotesize p.~#1}}%
\newcommand{\msp}[1]{\marginpar{\footnotesize [#1]}}%
\newcommand{\M}{\mathcal{M}}%
\newcommand{\Filt}{\mathrm{Filt}}%
\newcommand{\Fl}{\mathrm{Fl}}%
\newcommand{\uFl}{\underline{\mathrm{Fl}}}%
\newcommand{\uFilt}{\underline{\mathrm{Filt}}}%
\newcommand{\card}{\operatorname{card}}%
\newcommand{\Hom}{\operatorname{Hom}}%
\newcommand{\Ob}{\operatorname{Ob}}%
\newcommand{\wavy}[1]{\widetilde{#1}}%
\makeatother\let\printedpage\derpp\let\msp\dermsp%
\let\mspin\DERmspin%


\printedpage{14}
% head: A. Grothendieck : Dérivateurs, Ch. XV, Édition des Universités de Montpellier II et Paris VII
% foot: Version 5 novembre 2010
\msp{page 20}
première formule (3.7 bis), notons que l'on a égalité lorsque $\M$ est la catégorie des couples $(A,B)$, $A$ un ensemble et $B$ une partie de $A$, avec les flèches évidentes, et qu'on définit $\M^\pi$ (pour $\pi\geq\pi_0$) comme la sous-catégorie strictement pleine formée des $(A,B)$ tels que $\card A\geq\pi$.

\textsc{Démonstration de} 3.1. Prouvons d'abord la formule (3.8),
\[
\card\Hom(A,X)\ \leq\ s^{(r^{\pi_0})}\qquad\quad
\left\{\begin{array}{l}
\text{si } A\in\Filt^r,\ X\in\Filt^s,\\[1.5ex]
r\geq\pi_1,\ s\geq\sup(\pi_1,C_1).
\end{array}\right.
\]

(\textbf{N.B.} Je n'utiliserai pas $r\geq C_1$, postulé dans 3.1 $1^\circ$.) J'écris \texttt{[<<J'écris>> rayé]}
\[
\begin{aligned}
A\ &=\ \varinjlim_I A_i\ ,\quad \card I\ \leq\ r^{\pi_0}\ ,\quad A_i\ \in\ \Filt^{\pi_0}\ ,\\
X\ &=\ \varinjlim_J X_j\ ,\quad \card J\ \leq\ s^{\pi_1}\ ,\quad X_j\ \in\ \Filt^{\pi_1}\ ,\qquad\qquad J\text{ grand devant }\pi_1\ .
\end{aligned}
\]
On aura
\[
\Hom(A,X)\ \simeq\ \varprojlim_{i\in I}\ \underbrace{\Hom(A_i,\,\varinjlim_J X_j)}_{\overset{(\text{\footnotemark[6]})}{\simeq}\ \varinjlim\limits_{j\in J}\ \underbrace{\Hom(A_i,X_j)}_{\text{de cardinal }\leq C_1}}\ ,
\]
\footnotetext[6]{car les $A\in\Filt^{\pi_0}$ sont $\pi_1$-accessibles.}
donc pour tout $i\in I$
\[
\Hom(A_i,X)\ \leq\ \card(J)\times C_1\ \leq\ s^{\pi_1}\times C_1\ =\ s^{\pi_1}\ ,
\]
car $s\geq C_1$, donc $s^{\pi_1}\geq C_1$, d'où
\[
\card(\,\varprojlim_{j\in I}\Hom(A_i,X))\ \leq\ (s^{\pi_1})^{\card I}\ \leq\ \underbrace{(s^{\pi_1})^{(r^{\pi_0})}}_{s^{\pi_1 r^{\pi_0}}}\ . % [sic] j\in I -> i\in I, the limit is over the index i of A_i
\]
Comme $r\geq\pi_1$, on a $r^{\pi_0}\geq\pi_1$, donc $r^{\pi_0}\pi_1=r^{\pi_0}$, d'où
\[
\card\Hom(A,X)\ \leq\ s^{(r^{\pi_0})}\ ,
\]
c'est bien la formule (3.8). Pour en tirer une estimation pour $\M^r/X$ (avec $X\in\Filt^s$), \msp{page 21} je vais d'abord majorer le cardinal des classes d'isomorphie d'objets de $\M^r/X$, \textit{i.e.} des flèches à isomorphisme près (dans $\M^r$)
\[
A\ \longrightarrow\ X\ ,\qquad\quad X\text{ fixé dans }\Filt^s,\ A\text{ variable dans }\M^r\ .
\]
On aura
\[
A\ =\ \varinjlim_{i\in I} A_i\qquad\quad\text{avec } A_i\in\Filt^{\pi_0},\ I\text{ ensemble ordonné de cardinal}\leq r^{\pi_0}\ .
\]

\printedpage{15}
% head: A. Grothendieck : Dérivateurs, Ch. XV, Édition des Universités de Montpellier II et Paris VII
% foot: Version 5 novembre 2010
Pour $I$ fixé, le cardinal de l'ensemble des systèmes inductifs indexés par $I$, dans $\M^{\pi_0}/X$, à isomorphisme près, peut se majorer ainsi : pour tout $i\in I$, il faut choisir $A_i$ dans $\M^{\pi_0}$ (mais seule sa classe dans $\wavy{\M^{\pi_0}}$ importe), donc c'est majoré par card $\wavy{\M^{\pi_0}}\ \leq C_0$~\textsuperscript{(}\footnote[7]{On a posé ici $C_0=\card\wavy{\uFl^{\pi_0}}$ , donc $C_0\leq C_1$.}\textsuperscript{)}, et pour tous les $i$, par $\card I\times C_0\leq r^{\pi_0}\times C_0$, et si $r\geq C_0$, donc $r^{\pi_0}\geq C_0$, on trouve une majoration des choix possibles du système des \textit{objets} $A_i$ dans $\M^\pi$, à isomorphisme près, par $r^{\pi_0}$. Si on se donne de plus les morphismes de transition $A_i\longrightarrow A_j$ pour $i<j$, cela fait pour chaque paire comparable, élément de $I\times I$, un choix majoré par $C_0$, donc encore une majoration par $\ \underbrace{\card(I\times I)}_{=(\card I)^2=\card I}\cdot C_0\leq r^{\pi_0}\cdot C_0=r^{\pi_0}$. Donc l'ensemble des systèmes inductifs en question, à isomorphisme près, pour $I$ fixé, est de cardinal $\leq r^{\pi_0}\cdot r^{\pi_0}=r^{\pi_0}$. Enfin, les \textit{homomorphismes} d'un système inductif donné dans $X$ sont majorés par $\card I\times s^{(r^{\pi_0})}$ (par (3.8) déjà prouvé), donc par
\[
r^{\pi_0}\cdot s^{(r^{\pi_0})}\ =\ s^{(r^{\pi_0})}\qquad\quad(\text{puisque } r^{\pi_0}\leq s^{(r^{\pi_0})})\ .
\]

Au total, pour $I$ fixé, l'ensemble des systèmes inductifs à isomorphisme près, indexés par $I$, à valeurs dans $\M^r/X$, est \msp{page 22} de cardinal majoré par $\quad\underbrace{r^{\pi_0}}_{\text{\scriptsize\begin{tabular}{@{}c@{}}nombre des\\ systèmes inductifs\\ dans $\M^\pi$, à\\ isomorphisme près\end{tabular}}}\quad\cdot\quad\underbrace{s^{(r^{\pi_0})}}_{\text{\scriptsize\begin{tabular}{@{}c@{}}nombre des\\ flèches d'un\\ tel système inductif\\ dans $X$\end{tabular}}}\quad=s^{(r^{\pi_0})}$.

Mais il faut maintenant faire varier $I$, à isomorphisme près, dans la catégorie des ensembles ordonnés (filtrants \textit{etc.}) de cardinal $\leq\lambda=r^{\pi_0}$. On voit de suite que cet ensemble de classes d'isomorphie (qui peut s'interpréter comme un ensemble quotient d'un certain ensemble de parties d'un ensemble fixe $I_0$ de cardinal $\lambda$) est de cardinal $\leq 2^\lambda$. On trouve donc finalement
\[
\card\wavy{\M^r/X}\ \ \leq\ s^{(r^{\pi_0})}\cdot 2^{(r^{\pi_0})}\ =\ s^{(r^{\pi_0})}\ ,
\]
sous condition que
\[
r\ \geq\ \sup(\pi_1,C_0)\ ,\qquad\quad s\ \geq\ \sup(\pi_1,C_1)\ .
\]
Il faut prouver l'inégalité similaire pour card $\wavy{\uFl^r\,\M}$ \ (3.9). Or, comme pour $A$, $B$ dans $\M^r$, on a, si $r\geq\sup(\pi_1,C_1)$,
\[
\card\Hom_{\M}(A,B)\ \leq\ r^{(r^{\pi_0})}
\]
en vertu de (3.8) où on fait $s=r$, on trouve que
\[
\begin{aligned}
\card\wavy{\uFl(\M^r/X)}\ &\leq\ \Bigl(\card\wavy{\M^r/X}\,\Bigr)^2\cdot r^{(r^{\pi_0})}\\
&\leq\ s^{(r^{\pi_0})}\cdot r^{(r^{\pi_0})}\\
&=\ s^{(r^{\pi_0})}\qquad\text{si } s\geq r\ .
\end{aligned}
\]
(On va regarder plus bas les cas $s\leq r$~\textsuperscript{(}\footnote[8]{\textit{cf.} cor. 2, page 23.}\textsuperscript{)}). C'est bien la formule (3.9).

\printedpage{16}
% head: A. Grothendieck : Dérivateurs, Ch. XV, Édition des Universités de Montpellier II et Paris VII
% foot: Version 5 novembre 2010
\msp{page 23}
Venons-en à la catégorie $\M^r$ elle-même. L'argument de la page 21 montre que l'ensemble des classes d'isomorphie d'objets de $\M^r$ indexés par un $I$ fixé, de cardinal $\leq r^{\pi_0}$, comme $A=\varinjlim\limits_i A_i$ avec $A_i\in\Filt^{\pi_0}$, est de cardinal majoré par $r^{\pi_0}$ (si $r\geq C_0$). Si on fait varier $I$, on trouve une majoration par $r^{\pi_0}\cdot 2^{(r^{\pi_0})}=2^{(r^{\pi_0})}$. On en conclut comme tantôt que l'ensemble $\ \wavy{\Fl\,\M^r}$ est de cardinal majoré par
\[
(2^{(r^{\pi_0})})^2\times r^{(r^{\pi_0})}\qquad\quad(\text{si } r\geq\sup(\pi_1,C_1))\ ,
\]
donc par $r^{(r^{\pi_0})}$. Je dis qu'on a
\[
r^{(r^{\pi_0})}\ =\ 2^{(r^{\pi_0})}\ ,
\]
en effet, on a $\geq$, et pour prouver $\leq$, on note que $r<2^r$, donc $r^{(r^{\pi_0})}\leq(2^r)^{(r^{\pi_0})}=2^{r\cdot r^{\pi_0}}=2^{(r^{\pi_0})}$, O.K.! Cela établit donc 3.7. Cet argument prouve en fait le

\textbf{Lemme 3.2.} \textit{Soient $r$ un cardinal, $r'$ un cardinal $\geq r$. Alors on a}
\[
r^{r'}\ =\ 2^{r'}\qquad\quad(=s^{r'}\ \textit{pour tout}\ 2\leq s\leq r)\ .
\]

\textbf{Corollaire 1.} \textit{Soient $r\geq\sup(\pi_1,C_1)$, et $A,X\in\Filt^\pi$. Alors}
\[
\card\Hom_{\M}(A,X)\ \leq\ 2^{(r^{\pi_0})}\qquad\quad(=2^r\ \textit{si}\ r=2^c,\ \textit{avec}\ c\geq\pi_0)\ . \tag*{(3.10)}
\]

On applique (3.8), qui donne la majoration $r^{(r^{\pi_0})}$, et le lemme 3.2.

On notera qu'on trouve la même majoration pour le cardinal de $\wavy{\uFl(\M^r)}$ tout entier (3.7)! De même, (3.9) donne :

\textbf{Corollaire 2.} \textit{Soit $r\geq\sup(\pi_1,C_1)$, et $X\in\Ob\M^r$. Alors}
\[
\card\wavy{\uFl(\M^r/X)}\ \leq\ 2^{(r^{\pi_0})}\qquad\quad(=2^r\ \textit{si}\ r=2^c,\ \textit{avec}\ c\geq\pi_0)\ . \tag*{(3.11)}
\]
c'est donc la formule (3.9) pour $s\leq r$, laissée en suspens.

\msp{page 24}
Revenons à la majoration cardinale des $X_\alpha(f)$, où $f : X\longrightarrow Y$ est dans $\Fl^s(\M)$, \textit{i.e.} $X,Y\in\Filt^s$. On suppose de plus $C_0\subset\Fl^r(\M)$,
\[
\begin{aligned}
&C_0\ \subset\ \Fl^r(\M)\ =\ \Fl(\M^r)\ ,\qquad\quad(f:X\longrightarrow Y)\in\Fl^s(\M)\ ,\ \ \textit{i.e.}\ X,Y\in\Filt^s\ ,\\
&\text{avec } s\geq r\ .
\end{aligned} \tag*{(3.10)}
\]

\texttt{[Deux fois (3.10).]} Il y a une difficulté idiote, c'est que $\M^r$ est équivalente à une petite catégorie, mais n'est pas elle-même petite. Même si $C_0$ ne contenait qu'une seule classe d'isomorphie de flèches de $\M^r$, son cardinal pourrait être (pour $r$ fixé) arbitrairement élevé. Aussi, dans ce qui suit, je dois supposer la condition de “non pléthoricité cardinale” pour $C_0$ :
\[
\parbox{0.72\linewidth}{$\forall\,i\in C_0$, désignant par $C_0(i)$ l'ensemble des $i'\in C_0$ telles que $i\simeq i'$ dans $\Fl(\M)$, j'impose que $\card C_0(i)\leq 2^{(r^{\pi_0})}$ (ce qui est aussi la valeur de la majoration pour le cardinal de $\wavy{\uFl^r(\M)}=\wavy{\uFl(\M^r)}$\,).} \tag*{(3.11)}
\]

\printedpage{17}
% head: A. Grothendieck : Dérivateurs, Ch. XV, Édition des Universités de Montpellier II et Paris VII
% foot: Version 5 novembre 2010
\texttt{[Deux fois (3.11).]} Cela implique donc que l'on a $\card C_0\leq(2^{(r^{\pi_0})})^2=2^{(r^{\pi_0})}$,
\[
\card C_0\ \leq\ 2^{(r^{\pi_0})}\ , \tag*{(3.12)}
\]
qui est équivalente à (3.11) (moyennant $C_0\subset\Fl^r(\M)$, postulé dans (3.10)).

Il s'ensuit par (3.8) que l'on a
\[
\card\Delta(C_0,f)\ \leq\ \underbrace{\card(C_0)}_{\leq 2^{(r^{\pi_0})}}\cdot s^{(r^{\pi_0})}\cdot s^{(r^{\pi_0})}\ \leq\ s^{(r^{\pi_0})}\ .
\]

D'autre part, les $X_D$ dont on prend la somme amalgamée sous $X$, pour obtenir $X_1(f)$, sont dans $\M^s$, si on fait l'hypothèse suivante. \msp{page 25}
\[
\left|\ \parbox{0.72\linewidth}{Les $\M^\pi=\uFilt^\pi$ sont stables par sommes amalgamées simples, et par $\varinjlim$ filtrantes de cardinal $\leq\pi$ (pas nécessairement grandes devant $\pi_0$) ,}\right. \tag*{(3.13)}
\]
compte tenu que $r\leq s$, donc $\M^r\subset\M^s$, donc $A$, $B$, $X$ dans $X_D=B\amalg_A X$ sont dans $\M^s$. Les sommes amalgamées partielles \textit{finies} des $X_D$ sont donc aussi dans $\M^s$, par récurrence sur le cardinal de l'ensemble d'indices. La somme amalgamée infinie apparaît ainsi comme $\varinjlim$ filtrante d'objets de $\M^s$, l'ensemble d'indices étant l'ensemble $\mathfrak{P}^*_{\mathfrak{f}}\Delta(C_0,f)$ des parties finies non vides de $\Delta(C_0,f)$. Cet ensemble est de cardinal majoré par $s'$, si $s'$ est un cardinal \textit{infini} tel que $\card\Delta(C_0,f)\leq s'$. On peut donc prendre $s'=s^{(r^{\pi_0})}$. On a donc prouvé, compte tenu de (3.13) appliqué à $\M^{s'}$, le % ?? uncertain: symbol in X_D = B \amalg_A X -- \amalg or a doubled \sqcup; subscript of \mathfrak{P}^* -- fraktur f or roman f

\textbf{Lemme 3.3.} \textit{On suppose} \textup{(3.10), i.e.} $C_0\subset\Fl^r(\M)$, $f\in\Fl^s(\M)$, \textit{avec} $s\geq r\geq\sup(\pi_1,C_1)$, \textit{et de plus que $C_0$ satisfait la condition de non-pléthoricité} \textup{(3.11)} \textup{(}\textit{ce qu'on peut toujours supposer, quitte à réduire $C_0$ sans changer $\widetilde{C_0}=\Phi$ et $C_{0*}=\Psi$}\textup{)}\textit{, et que la filtration cardinale satisfait la condition de stabilité (anodine!)} \textup{(3.13)}\textit{. Posons}
\[
s'\ =\ s^{(r^{\pi_0})}\ . \tag*{(3.14)}
\]
\textit{Alors $X_1(C_0,f)$ est dans $\M^{s'}$. Plus généralement et plus joliment, pour $C_0$ et $r$ fixés} \textup{(}\textit{avec $C_0\subset\Fl^r(\M)$ et satisfaisant la condition de non-pléthoricité} \textup{(3.11)} \textit{relativement à $r$}\textup{)}\textit{, pour tout cardinal} \msp{page 26} \textit{de la forme}
\[
s'\ =\ s^t\ ,\qquad\quad\textit{avec}\ t\geq r^{\pi_0},\ s\geq 2 \tag*{(3.15)}
\]
\textit{(d'où $s'\geq 2^{(r^{\pi_0})}>r^{\pi_0}\geq r$), $\uFl^{s'}(\M)$ est stable sous le foncteur $p_1^{C_0}:\uFl(\M)\longrightarrow\uFl(\M)$.}

En effet, si on applique la première partie du lemme au couple $(r,s')$, on trouve que $X_1(f)$ est dans $\M^{s''}$, avec $s''=s'^{(r^{\pi_0})}$. Mais
\[
s''\ =\ s'^{(r^{\pi_0})}\ =\ (s^t)^{(r^{\pi_0})}\ =\ s^{tr^{\pi_0}}\ \underset{\substack{\vert\\ \text{\scriptsize car $t\geq r^{\pi_0}$,}\\ \text{\scriptsize donc $tr^{\pi_0}=t$}}}{=}\ s^t\ =\ s'\ ,
\]
\printedpage{18}
% head: A. Grothendieck : Dérivateurs, Ch. XV, Édition des Universités de Montpellier II et Paris VII
% foot: Version 5 novembre 2010
donc on voit que $f\in\Fl^{s'}$ implique $X_1(f)\in\M^{s'}$, ou ce qui revient au même, $p_1(f)\in\Fl^{s'}$.

\textbf{Corollaire.} \textit{Soient $C_0$, $r$, $s'=s^t$ $(t\geq r^{\pi_0},\ s\geq 2)$ comme dessus, et soit $\alpha$ un nombre ordinal tel que}
\[
\card\alpha\ \leq\ s'\ =\ s^t\ .
\]
\textit{Alors $\uFl^{s'}(\M)$ est stable par le foncteur $p_\alpha=p_\alpha^{C_0}$.}

On le voit aussitôt par récurrence transfinie, compte tenu du lemme 3.3, et de l'hypothèse de stabilité (3.13).

Pour résumer :

\textbf{Théorème 2.} \textit{Soit $\M$ une catégorie stable par petites $\varinjlim$ filtrantes, et munie d'une filtration cardinale} $(\Filt^\pi)_{\pi\geq\pi_0}$\textit{, telle que les} $\Filt^\pi$ \textit{soient stables par $\varinjlim$ filtrantes de cardinal $\leq\pi$ (pas forcément grandes devant $\pi_0$), et par les sommes amalgamées simples qui existent dans $\M$. On suppose que les éléments de} $\Filt^{\pi_0}$ \textit{soient accessibles, et on désigne par $\pi_1$ le plus petit cardinal $\geq\pi_0$, tel que les objets de} $\Filt^{\pi_0}$ \textit{soient $\pi_1$-accessibles. On pose} \msp{page 27}
\[
c_1\ =\ \card\Fl(\M^{\pi_1})\ =\ \card\Fl^{\pi_1}(\M)
\]
\textit{(où pour tout $\pi\geq\pi_0$, $\M^\pi$ désigne la sous-catégorie strictement pleine de $\M$ définie par} $\Filt^\pi$\textit{, et} $\Fl^\pi(\M)=\Fl(\M^\pi)$\textit{).}

\textit{Soit $C_0$ une petite partie de} $\Fl(\M)$\textit{, formée de flèches coquarrables, et posons}
\[
\Phi\ =\ \widetilde{C_0}\ =\ \Psi^*\ ,\qquad\quad\textit{où}\ \ \Psi=C_{0*}=\Phi_*\ .
\]
\textit{Soit $r$ un cardinal} $\geq\sup(\pi_1,c_1)$\textit{, et tel que l'on ait $C_0\subset\Fl^r(\M)$. On suppose que $C_0$ est irrédondant (}i.e.\textit{ deux objets de $C_0$ isomorphes sont égaux), ou plus généralement, que $C_0$ satisfait la très anodine condition de non-pléthoricité} \textup{(3.11)} \textit{relativement à $r$. On considère la suite transfinie des foncteurs $p_\alpha=p_\alpha^{C_0}$ sur $\uFl(\M)$} \textup{(}\textit{cf.} \textup{th. 1)}\textit{, donnant lieu aux factorisations fonctorielles d'une $f\in\uFl(\M)$, $f:X\longrightarrow Y$, en}
\[
X\ \xrightarrow{i_\alpha(f)}\ X_\alpha(f)\ \xrightarrow{p_\alpha(f)}\ Y\ . \tag*{(3.16)}
\]

\textit{On a alors ce qui suit :}

$1^\circ$) \textit{On a (pour mémoire) $i_\alpha(f)\in\Phi$ pour tout $\alpha$, et de plus,} si $\alpha$ est grand devant $r^{\pi_0}$, \textit{on a aussi $p_\alpha(f)\in\Psi$. Si $r'$ est un cardinal $>r^{\pi_0}$ (p. ex. le cardinal successeur de $r^{\pi_0}$), on peut prendre pour $\alpha$ le premier ordinal dont le cardinal soit égal à $r'$.}

$2^\circ$) \textit{Pour tout cardinal $s$ de la forme}
\[
s\ =\ c^t\qquad\quad\textit{avec}\ \ t\geq r^{\pi_0},\ c\geq 2
\]
\textit{(donc $s=2^t$ si on prend $c\leq t$,} cf. \textit{lemme 3.2), $\uFl^s(\M)$ est stable sous les $p_\alpha$ tels que} $\card\alpha\leq s$\textit{. Si on prend p. ex. pour $\alpha$ le premier ordinal dont le cardinal soit $>r^{\pi_0}$, alors $\uFl^s(\M)$ est stable par $p_\alpha$, qui par $1^\circ$) réalise donc une factorisation fonctorielle} \textup{(3.16)} \textit{avec $i_\alpha(f)\in\Phi$, $p_\alpha(f)\in\Psi$.}


\endgroup
% END UNIT Der-ChXV-05-11-10/Der-ChXV-05-11-10_p016-020.tex
% BEGIN UNIT Der-ChXV-05-11-10/Der-ChXV-05-11-10_p021-025.tex
\begingroup\makeatletter%
\newcommand{\printedpage}[1]{\marginpar{\footnotesize p.~#1}}%
\newcommand{\msp}[1]{\marginpar{\footnotesize [#1]}}%
\newcommand{\Fl}{\mathrm{Fl}}%
\newcommand{\id}{\mathrm{id}}%
\newcommand{\qpre}{\mathrel{\leftarrow\!\!\!\!\mapsto}}%
\makeatother\let\printedpage\derpp\let\msp\dermsp%
\let\mspin\DERmspin%


\printedpage{19}
% head: A. Grothendieck : Dérivateurs, Ch. XV, Édition des Universités de Montpellier II et Paris VII
% foot: Version 5 novembre 2010

\msp{page 28}
\textsc{Démonstration}. Tout a été prouvé, si on se rappelle que les objets de $\mathcal{M}^r$ sont $r^{\pi_0}$-accessibles (\textit{cf.} (3.2)).

\section*{4\quad Compléments sur les structures de Quillen}

\textbf{Définition 4.1.} Soit $\mathcal{M}$ une catégorie. On appelle \textit{précouple de Quillen} dans $\mathcal{M}$ un couple $(\Phi,\Psi)$ de parties de $\Fl(\mathcal{M})$, satisfaisant les deux conditions suivantes.
\begin{itemize}[leftmargin=2.5em]
\item[\textit{a})] On a $\Phi\longleftrightarrow\Psi$, \textit{i.e.} $\Psi\subset\Phi_*$, ou ce qui revient au même, $\Phi\subset\Psi^*$ -- soit encore : pour tout carré commutatif dans $\mathcal{M}$
\[
\begin{tikzcd}[row sep=large, column sep=large]
A \arrow[r, "u"] \arrow[d, "i"'] & X \arrow[d, "p"] \\
B \arrow[ur, dotted, "w" description] \arrow[r, "v"] & Y \;,
\end{tikzcd}
\tag*{(4.1.1)}
\]
avec $i\in\Phi$, $p\in\Psi$, il existe $w$ telle que $wi=u$, $pw=v$.
\item[\textit{b})] Toute $f\in\Fl(\mathcal{M})$ se factorise en
\[
f \;=\; p\,i\;, \qquad \text{avec } i\in\Phi,\, p\in\Psi\ .
\]
\end{itemize}
On écrit
\[
\Phi \;\qpre\; \Psi
\tag*{(4.1.2)}
\] % ?? uncertain: the printed relation sign -- a double arrow with a small vertical bar in the middle (rendered by \qpre)
pour indiquer que $(\Phi,\Psi)$ est un précouple de Quillen.

On dit que le précouple de Quillen est \textit{clos}, si $\Phi$ est clos à gauche, $\Psi$ clos à droit. On va voir plus bas (cor. 2, page 31) que cela implique (donc équivaut à) % [sic] clos à droit -> clos à droite, gender agreement with « à droite »
\[
\Phi \;=\; \Psi^*\;, \qquad \Psi \;=\; \Phi_*\ .
\tag*{(4.1.3)}
\]

\msp{page 29}
Un précouple de Quillen clos est aussi appelé un couple de Quillen.

\textbf{Proposition 4.2.} \textit{Soit} $(\Phi,\Psi)$ \textit{un précouple de Quillen dans} $\mathcal{M}$\textit{, et soit} $j\in\Fl(\mathcal{M})$\footnote[9]{\textbf{N.B.} On utilise seulement que $\Phi\longleftrightarrow\Psi$, et que $j$ peut se mettre sous la forme $pi$, avec $i\in\Phi$, $p\in\Psi$.}. % footnote mark printed as "(9)"
\textit{On a alors les conditions équivalentes sur} $j$ \textit{:}
\begin{itemize}[leftmargin=2.5em]
\item[a)] $j\in\widetilde{\Phi}$.
\item[b)] $j\in\Psi^*$.
\item[c)] $j$ \textit{est facteur direct d'un} $i\in\Phi$.
\item[c$'$)] $j$ \textit{est facteur direct d'un} $i\in\Phi$ \textit{ayant même source que} $j$\textit{, plus précisément, il existe un diagramme}
\[
\begin{tikzcd}[row sep=large, column sep=large]
A \arrow[r, "j"] \arrow[d, shift right=1ex, "\id_A"'] & B \arrow[d, shift right=1ex, "q"'] \\
A \arrow[u, dotted, shift right=1ex, "\id_A"'] \arrow[r, "i"] & B' \;, \arrow[u, dotted, shift right=1ex, "p"']
\end{tikzcd}
\tag*{(4.2.1)}
\]
\printedpage{20}
% head: A. Grothendieck : Dérivateurs, Ch. XV, Édition des Universités de Montpellier II et Paris VII
% foot: Version 5 novembre 2010
\textit{avec} $i\in\Phi$\textit{, et}
\[
qj \;=\; i\;, \qquad pi \;=\; j\;, \qquad pq \;=\; \id_B\;,
\tag*{(4.2.2)}
\]
\textit{de sorte que} $j$ \textit{est un facteur direct de} $i\in\Phi$\textit{, par} $(\id_A,q)$\textit{, admettant la rétraction} $(\id_A,p)$.
\item[d)] \textit{Pour tout diagramme commutatif traits pleins}
\[
\begin{tikzcd}[column sep=small]
A \arrow[rr, "j"] \arrow[dr, "i"'] & & B \arrow[dl, dotted, bend left=40, "q"] \\
& B' \arrow[ur, "p"'] &
\end{tikzcd}
\;,
\tag*{(4.2.3)}
\]
i.e. \textit{toute factorisation de} $j$ \textit{en} $pi$\textit{, si on a} $p\in\Psi$\textit{, il existe une section} $q$ \textit{de} $B'$ \textit{sur} $B$ \textit{dans} $A\backslash\mathcal{M}$\textit{,} i.e. \textit{une flèche} $q:B\longrightarrow B'$ \textit{telle que}
\[
pq \;=\; \id_B \qquad \textit{et} \qquad qj \;=\; i\ .
\]
\end{itemize}

\textsc{Démonstration} de la proposition. Elle se fait par la suite d'implications suivante, toutes $(\pm)$ tautologiques :
\[
(c') \;\Longrightarrow\; (c) \;\Longrightarrow\; (a) \;\Longrightarrow\; (b) \;\Longrightarrow\; (d) \;\Longrightarrow\; (c')\;,
\]
\msp{page 30}
dont les quatre premières utilisent seulement l'hypothèse $\Phi\longleftrightarrow\Psi$ (de $(a)$ de la définition 4.1), et la dernière seulement, $(d)\Longrightarrow(c')$, fait appel à l'hypothèse de factorisation de Quillen ($(b)$ dans 4.1) ; et encore il suffit de la supposer pour $j$ seulement !

$(c')\Longrightarrow(c)$ tautologie pure, $(c)\Longrightarrow(a)$ car on sait que $\widetilde{\Phi}$ est stable par facteurs directs (et contient $\Phi$), $(a)\Longrightarrow(b)$ car $\widetilde{\Phi}\subset\Psi^*$ ($\widetilde{\Phi}$ étant par définition la plus petite partie close à gauche de $\Fl(\mathcal{M})$ contenant $\Phi$, et $\Psi^*$ étant close à gauche), $(b)\Longrightarrow(d)$ car on écrit le diagramme donné (4.2.2) (traits pleins) comme % [sic] (4.2.2) -> (4.2.3), the diagram given in condition d) is (4.2.3)
\[
\begin{tikzcd}[row sep=large, column sep=large]
A \arrow[r, "u"] \arrow[d, "j\in\Psi^*"'] & X \arrow[d, "p\in\Psi"] \\
B \arrow[ur, dotted, "q" description] \arrow[r, "v"] & Y \;,
\end{tikzcd}
\tag*{$(*)$}
\]
en y faisant $Y=B$, $v=\id_B$, et en rebaptisant $X$ par $B'$, et $u$ par $i$ : l'hypothèse $(b)$, $j\in\Psi^*$, signifie que dans un diagramme traits pleins $(*)$, il existe toujours $q$ telle que $qj=u$, $pq=v$, et en l'exprimant dans le cas particulier où $v=\id_B$, on trouve $(d)$.

Reste à prouver $(d)\Longrightarrow(c')$, en utilisant cette fois l'hypothèse de factorisation sur $j$, \textit{i.e.} qu'on peut écrire $j$ sous la forme $pi$ :
\[
\begin{tikzcd}[column sep=small]
A \arrow[rr, "j"] \arrow[dr, "i"'] & & B \arrow[dl, dotted, bend left=40, "q"] \\
& B' \arrow[ur, "p"'] &
\end{tikzcd}
\;,
\]
avec $i\in\Phi$, $p\in\Psi$. En vertu de $(d)$, comme $p\in\Psi$, il existe donc $q:B\longrightarrow B'$ satisfaisant
\printedpage{21}
% head: A. Grothendieck : Dérivateurs, Ch. XV, Édition des Universités de Montpellier II et Paris VII
% foot: Version 5 novembre 2010
\msp{page 31}
\[
pq \;=\; \id_B\;, \qquad qj \;=\; i\;,
\]
ce qui prouve $(c')$, puisque $i\in\Phi$.

\textbf{Corollaire 1.} \textit{Soit} $(\Phi,\Psi)$ \textit{un précouple de Quillen. Alors on a}
\[
\begin{aligned}
\widetilde{\Phi} \;\;=\;\; \Psi^* \;\;=\;\; &\textit{ensemble des objets de } \underline{\Fl}(\mathcal{M}) \textit{ qui sont fac-}\\
&\textit{teurs directs d'éléments de } \Phi\ ,
\end{aligned}
\tag*{(4.3)}
\] % printed line break of the display falls inside "fac-teurs"; kept as printed
\textit{et dualement}
\[
\begin{aligned}
\overline{\Psi} \;\;=\;\; \Phi_* \;\;=\;\; &\textit{ensemble des objets de } \underline{\Fl}(\mathcal{M}) \textit{ qui sont fac-}\\
&\textit{teurs directs d'éléments de } \Psi\ .
\end{aligned}
\tag*{(4.4)}
\] % printed line break of the display falls inside "fac-teurs"; kept as printed

\textbf{Corollaire 2.} \textit{Les conditions suivantes sur le précouple de Quillen} $(\Phi,\Psi)$ \textit{sont équivalentes.}
\begin{itemize}[leftmargin=2.5em]
\item[a)] \textit{Le précouple est} clos, i.e. $\Phi=\widetilde{\Phi}$\textit{,} $\Psi=\overline{\Psi}$\textit{.}
\item[b)] $\Phi=\Psi^*$\textit{,} $\Psi=\Phi_*$\textit{.}
\item[c)] $\Phi$ \textit{et} $\Psi$ \textit{sont stables par facteurs directs.}
\end{itemize}
Soit maintenant
\[
W \;\subset\; \Fl(\mathcal{M})
\tag*{(4.5)}
\]
un localiseur, satisfaisant les conditions habituelles
\[
\left\{
\begin{aligned}
&\textit{a})\quad W \text{ contient les isomorphismes.}\\[1ex]
&\textit{b})\quad \text{Si } X\xrightarrow{\;f\;}Y\xrightarrow{\;g\;}Z \text{ dans } \mathcal{M}\text{, si deux parmi } f,\ g,\ gf \text{ sont}\\
&\phantom{\textit{b})}\quad \text{dans } W\text{, alors le troisième aussi.}
\end{aligned}
\right.
\tag*{(4.6)}
\]
($W$ modérément saturé). Supposons de plus données quatre autres parties
\[
C\ ,\ TC\ , \qquad F\ ,\ TF
\tag*{(4.7)}
\]
de $\Fl(\mathcal{M})$, satisfaisant les conditions suivantes.
\[
\left\{
\begin{aligned}
&\textit{a})\quad (C,TF) \text{ et } (TC,F) \text{ sont des précouples de Quillen.}\\[1ex]
&\textit{b})\quad TC\subset C\cap W,\ TF\subset F\cap W.
\end{aligned}
\right.
\tag*{(4.8)}
\]

\msp{page 32}
Ces propriétés s'explicitent ainsi :
\begin{itemize}[leftmargin=2.5em]
\item[\textbf{M\,1}] Propriétés de relèvement-prolongement, $C\longleftrightarrow TF$ et $TC\longleftrightarrow F$.
\item[\textbf{M\,2}] Propriété de factorisation de toute flèche $f$ de $\mathcal{M}$ en $f=pi$, avec, au choix, $i\in C$, $p\in TF$, ou $i\in TC$, $p\in F$.
\item[\textbf{M$'$}] $TC\subset C\cap W$, $TF\subset F\cap W$.
\end{itemize}

\printedpage{22}
% head: A. Grothendieck : Dérivateurs, Ch. XV, Édition des Universités de Montpellier II et Paris VII
% foot: Version 5 novembre 2010
C'est donc une partie des axiomes d'une structure de Quillen faible pour $W$ (XIII, p. 27, déf. 1) -- il manque les axiomes de stabilité M\,3, M\,4, et les axiomes de “propreté” (propriety) M$''$, savoir $F\subset\mathrm{Fib}_W$, $C\subset\mathrm{Cof}_W$.

\textbf{Proposition 4.3.} \textit{Soient} $\mathcal{M}$ \textit{une catégorie,} $(C,TF)$ \textit{et} $(TC,F)$ \textit{deux précouples de Quillen dans} $\mathcal{M}$\textit{, satisfaisant}
\[
TC \;\subset\; C\;, \qquad TF \;\subset\; F\;,
\tag*{(4.3.1)}
\]
\textit{et soit} $W$ \textit{une partie de} $\Fl(\mathcal{M})$ \textit{dont on ne suppose pour l'instant que} 4.6, (b)\footnote[10]{On utilise seulement que si $X\xrightarrow{\;f\;}Y\xrightarrow{\;g\;}Z$, et si $gf\in W$, alors $f\in W\Longleftrightarrow g\in W$. On n'a pas besoin que $W$ contienne les isomorphismes, ni qu'il soit stable par composition.}\textit{, et} % footnote mark printed as "(10)"
\[
W \;\supset\; TC\ ,\ TF\ ,
\tag*{(4.3.2)}
\]
\textit{donc on a}
\[
TC \;\subset\; C\cap W\ , \qquad TF \;\subset\; F\cap W\ .
\tag*{(4.3.3)}
\]
\textit{Considérons les conditions de prolongement-relèvement suivantes} (a), (a$'$)\textit{, et les conditions de factorisation} (b), (b$'$), (c), (d) (\textit{en se rappelant que} $\widetilde{TC}=F^*$\textit{,} $\overline{TF}=C_*$\textit{, en vertu de la proposition} 4.2 (\textit{cor.} 1) \textit{appliquée aux précouples de Quillen} $(TC,F)$\textit{,} $(C,TF)$) :
\begin{itemize}[leftmargin=2.5em]
\item[a)] $C\cap W\subset\widetilde{TC}$\quad (\textit{donc} $TC\subset C\cap W\subset\widetilde{TC}$).
\item[a$'$)] $F\cap W\subset\overline{TF}$\quad (\textit{donc} $TF\subset F\cap W\subset\overline{TF}$).
\item[b)] $W\subset TF\circ\widetilde{TC}$.
\item[b$'$)] $W\subset\overline{TF}\circ TC$.
\item[c)] $W\subset\overline{TF}\circ\widetilde{TC}$.
\item[d)] $W\subset TF{\circ}TC$ (\textit{donc, par} (4.3.2)\textit{,} $W=TF{\circ}TC$ \textit{si} $W$ \textit{est stable par composition}).\footnote[11]{\textbf{N.B.} La condition $(a)$ signifie aussi que $(C\cap W,F)$ est encore un précouple de Quillen, resp. $(b)$ signifie que $(C,F\cap W)$ en est un -- et alors cette paire de précouples de Quillen satisfait cette fois à $(d)$, \textit{i.e.} $W=TF'\circ TC'$.} % footnote mark printed as "(11)"
\end{itemize}

\msp{page 33}
\textit{Alors toutes les conditions sauf la dernière sont équivalentes} (\textit{et elles sont donc impliquées par cette dernière}). \textit{Donc lorsque} $TC$ \textit{ou} $TF$ \textit{sont stables par facteurs directs} (\textit{p. ex. si} $TC$ \textit{est clos à gauche, ou} $TF$ \textit{clos à droite}), \textit{toutes les conditions envisagées sont équivalentes} (\textit{puisque par le corollaire} 2 \textit{à la proposition} 4.2\textit{, on a alors} $TC=\widetilde{TC}$ \textit{resp.} $TF=\overline{TF}$\textit{, donc} (b$'$) \textit{resp.} (b) \textit{équivaut alors à} (d)).

\textsc{Démonstration}. Elle se fait par le diagramme d'implications
\[
\begin{tikzcd}[column sep=small, row sep=small]
& (d) \arrow[dl, Rightarrow] \arrow[dr, Rightarrow] & \\
(b) \arrow[dr, Rightarrow] & \mathrm{Q} & (b') \arrow[dl, Rightarrow] \\
& (c) \arrow[dl, Rightarrow] \arrow[dr, Rightarrow] & \\
(a) \arrow[uu, Rightarrow] & & (a') \arrow[uu, Rightarrow]
\end{tikzcd}
\;,
\]
\printedpage{23}
% head: A. Grothendieck : Dérivateurs, Ch. XV, Édition des Universités de Montpellier II et Paris VII
% foot: Version 5 novembre 2010
où les quatre implications du carré supérieur Q sont tautologiques. Il suffit donc d'établir les implications
\[
(c) \;\Longrightarrow\; (a) \;\Longrightarrow\; (b)\;,
\]
les implications $(c)\Longrightarrow(a')\Longrightarrow(b')$ en résultent par dualité.

Supposons donc $(c)$, \textit{i.e.}
\[
W \;\subset\; \overline{TF}\circ\widetilde{TC}\;,
\tag*{$(*)$}
\]
et prouvons $(a)$, \textit{i.e.} que
\[
C\cap W \;\subset\; \widetilde{TC}\ .
\]
Mais si $j\in C\cap W$, pour prouver que $j\in\widetilde{TC}$, écrivons $j=pi$, $i\in\widetilde{TC}$, $p\in\overline{TF}$, ce qui est possible par l'hypothèse $(*)$ et $j\in W$.

\msp{page 34}
\[
\begin{tikzcd}[column sep=small]
A \arrow[rr, "j\in C"] \arrow[dr, "i\in\widetilde{TC}"'] & & B \arrow[dl, dotted, bend left=60, "q"] \\
& B' \arrow[ur, "p\in\overline{TF}"'] &
\end{tikzcd}
\;.
\]
Utilisons maintenant l'hypothèse $j\in C$ et $C\longleftrightarrow\overline{TF}$, on trouve donc (prop. 3.2, $(a)\Longrightarrow(d)$) une flèche $q:B\longrightarrow B'$ telle que $pq=\id_B$, $qj=i$, donc $j$ apparaît comme facteur direct de $i$. Comme $i\in\widetilde{TC}$ et que $\widetilde{TC}$ est stable par facteurs directs, on en conclut $j\in\widetilde{TC}$, \quad \textit{q.e.d.}

$(a)\Longrightarrow(b)$. On suppose donc
\[
C\cap W \;\subset\; \widetilde{TC}\;,
\tag*{$(*)$}
\]
et on veut prouver $(b)$, \textit{i.e.}
\[
W \;\subset\; TF\circ\widetilde{TC}\ .
\]
Soit donc $f\in W$, et factorisons $f$ en $pi$, avec $i\in C$, $p\in TF$. (O.K. à cause du précouple de Quillen $(C,TF)$). Comme $TF\subset W$, on a $p\in W$, et comme $pi\in W$, on en conclut $i\in W$ (saturation modérée de $W$ -- on doit quand même l'utiliser !). On a donc $i\in C{\cap}W$, donc $i\in\widetilde{TC}$ par $(*)$, d'où $f\in TF\circ\widetilde{TC}$, \quad \textit{q.e.d.}

\textbf{Corollaire 1.} \textit{Supposons} $TC$ \textit{stable par facteurs directs. Alors toutes les conditions de la proposition} 4.3 \textit{sont équivalentes, et équivalent à la relation}
\[
C\cap W \;=\; TC\ .
\tag*{(4.3.4)}
\]
\textit{Dualement, si} $TF$ \textit{est stable par facteurs directs, les conditions de la proposition} 4.3 \textit{sont toutes équivalentes et équivalent à}
\[
F\cap W \;=\; TF\ .
\tag*{(4.3.4$'$)}
\]


\endgroup
% END UNIT Der-ChXV-05-11-10/Der-ChXV-05-11-10_p021-025.tex
% BEGIN UNIT Der-ChXV-05-11-10/Der-ChXV-05-11-10_p026-030.tex
\begingroup\makeatletter%
\newcommand{\printedpage}[1]{\marginpar{\footnotesize p.~#1}}%
\newcommand{\msp}[1]{\marginpar{\footnotesize [#1]}}%
\newcommand{\Fl}{\mathrm{Fl}}%
\newcommand{\Cof}{\mathrm{Cof}}%
\newcommand{\deq}{\overset{\text{déf}}{=}}%
\newcommand{\univ}{\mathrm{univ}}%
\makeatother\let\printedpage\derpp\let\msp\dermsp%
\let\mspin\DERmspin%


\printedpage{24}
% head: A. Grothendieck : Dérivateurs, Ch. XV, Édition des Universités de Montpellier II et Paris VII
% foot: Version 5 novembre 2010

\msp{page 35}
En effet, par le corollaire 2 de la proposition 4.2, l'hypothèse sur $TC$ implique que $TC = \widetilde{TC}$, donc comme $TC \subset C{\cap}W$, l'hypothèse $C{\cap}W \subset \widetilde{TC}$ $(= TC)$ équivaut à $C{\cap}W = TC$.

\textbf{Corollaire 2.} \textit{Supposons $C$ et $W$ stables par facteurs directs. Alors les conditions équivalentes de la proposition} 4.3 \textit{équivalent à}
\[
C\cap W \;=\; \widetilde{TC}\ . \tag*{(4.3.5)}
\]
\textit{Dualement, si $F$, $W$ sont stables par facteurs directs, les conditions équivalentes de la proposition} 4.3 \textit{équivalent à}
\[
F\cap W \;=\; \overline{TF}\ . \tag*{(4.3.5$'$)}
\]

En effet, si $C$ et $W$ sont stables par facteurs directs, il en est de même de $C\cap W$, et dès lors si $C\cap W\subset\widetilde{TC}$, comme $C\cap W$ contient $TC$, on en conclut $\widetilde{C\cap W}=\widetilde{TC}$. Mais comme $C\cap W\longleftrightarrow TC$ et $C\cap W$ est stable par facteurs directs, $C\cap W$ est clos à gauche par corollaire 2 de la proposition 4.2, \textit{i.e.} $C\cap W=\widetilde{C\cap W}$, d'où $C\cap W=\widetilde{TC}$, \quad \textit{q.e.d.}

\textbf{Corollaire 3.} \textit{Supposons les conditions équivalentes} (a) \textit{à} (c) \textit{de} 4.3 \textit{satisfaites, et de plus} $W$ stable par facteurs directs\footnote[12]{En fait, il suffit de supposer que l'on a $W\supset\widetilde{TC}$, $\overline{TF}$, \textit{i.e.} que tout facteur direct d'un élément de $W$ qui est soit dans $C$, soit dans $F$, est encore dans $W$. Alors $(\widetilde{C},\widetilde{TC},\bar{F},\overline{TF})$ définissent une structure de Quillen close, et on sait que cela implique qu'en fait, $W$ lui aussi est stable par facteurs directs.}. \textit{Alors les conditions préliminaires} (4.3.1) \textit{et} (4.3.2) \textit{restent satisfaites quand on remplace $(C,TF)$ et $(TC,F)$ par leurs clôtures $(C',TF')=(\widetilde{C},\overline{TF})$ et $(TC',F')=(\widetilde{TC},\bar{F})$, et les conditions équivalentes de la proposition} 4.3 \textit{sont encore satisfaites pour ces couples, sous les formes fortes}
\[
\left\{
\begin{aligned}
C'\cap W &\;=\; TC' && (=\widetilde{TC'}\,)\\
F'\cap W &\;=\; TF' && (=\overline{TF'}\,)\\
W &\;=\; TF'\circ TC'\ .
\end{aligned}
\right.
\tag*{(4.3.6)}
\]

\msp{page 36}
\textit{Le triple $(W,C',F')$ est un triple de Quillen clos $($correspondant à une “closed model category” de Quillen, du moins si on suppose $\mathcal{M}$ stable par $\varinjlim$ finies et $\varprojlim$ finies, et que $W$ est modérément saturé$)$.}\footnote[13]{Il en résulte le\par
\textbf{Corollaire 4.} \textit{Soient $W$ un localiseur modérément saturé, et $(C,TF)$, $(TC,F)$ deux précouples de Quillen tels que $TC\subset C\cap W$, $TF\subset F\cap W$. Alors les conditions suivantes sont équivalentes.}
\begin{enumerate}[leftmargin=*]
\item[a)] $TC=C\cap W$.
\item[b)] $TF=F\cap W$.
\item[c)] $W=TF\circ TC$.
\item[d)] $(W,C,F)$ \textit{est un triple de Quillen clos.}
\end{enumerate}}

\textsc{Démonstration.} Les relations (4.3.1) impliquent
\[
\widetilde{TC}\;\subset\;\widetilde{C}\qquad\text{et}\qquad\overline{TF}\;\subset\;\bar{F}\ ,
\]
\printedpage{25}
% head: A. Grothendieck : Dérivateurs, Ch. XV, Édition des Universités de Montpellier II et Paris VII
% foot: Version 5 novembre 2010
\textit{i.e.}
\[
TC'\;\subset\;C'\ ,\qquad TF'\;\subset\;F'\ ,
\]
\textit{i.e.} (4.3.1) pour les nouveaux couples de Quillen. La relation 4.3.2, $W\supset TC,TF$, implique
\[
W\;\supset\;\widetilde{TC}\ ,\ \overline{TF}\ ,\qquad \textit{i.e.}\qquad W\;\supset\;TC'\ ,\ TF'\ ,
\]
en vertu du corollaire 2 à la proposition 4.2 et du fait que $W$ est stable par facteurs directs ; donc on a (4.3.2) pour les nouveaux couples. Comme par hypothèse on a
\[
W\;=\;\overline{TF}\circ\widetilde{TC}\ ,
\]
ceci s'écrit
\[
W\;=\;TF'\circ TC'\ ,
\]
donc les nouveaux couples satisfont aux mêmes conditions, qui cette fois, les couples étant clos, prennent la forme forte (4.3.6).

Prouvons que le triple $(W,C',F')$ satisfait aux axiomes de Quillen d'un triple de Quillen clos. On a bien, par (4.3.6),
\[
\left\{
\begin{aligned}
C' &\;=\;(TF')^*\;=\;(F'\cap W)^*\ , &\qquad F' &\;=\;(TC')_*\;=\;(C'\cap W)_*\ ,\\
W &\;=\;TF'\circ TC'\;=\;C'_*\circ F'^*\ ,
\end{aligned}
\right.
\]
qui correspond à la définition des catégories de modèles de Quillen \textit{clos} : les axiomes M\,1, M\,2, M\,3, M\,4 de Quillen sont réalisés de façon triviale (M\,3 et M\,4 parce que
\msp{page 37}
$C$ et $F$, $TC$ et $TF$ sont clos à gauche resp. à droite). L'axiome M\,5 est la saturation modérée de $W$.

\textbf{4.4.} Voici l'idée pour construire plus ou moins “toutes” les structures de Quillen closes $(C,TC,F,TF)$ associées à un localiseur donné $W\subset\mathcal{M}$. On suppose que $\mathcal{M}$ satisfait aux conditions du paragraphe précédent, concernant l'existence de $\varinjlim$ filtrantes, d'une filtration cardinale, l'accessibilité dans $\mathcal{M}$ (celle des éléments de $\mathrm{Filt}^{\pi_0}$). On part d'une partie petite
\[
C_0\;\subset\;\Fl(\mathcal{M})
\]
telle que
\[
C_{0*}\;\subset\;W\ , \tag*{(4.4.1)}
\]
et on pose
\[
\left\{
\begin{aligned}
C &\;=\;\widetilde{C_0}\ , &\qquad TC &\;=\;C\cap W\ ,\\
TF &\;=\;C_*\;=\;C_{0*}\ , &\qquad F &\;=\;(TC)_*\ .
\end{aligned}
\right.
\tag*{(4.4.2)}
\]
Il y a un problème seulement, vu que l'on sait déjà que $(C,TF)$ est un couple de Quillen par le théorème de factorisation dans §1 (p. 10) : il faut la factorisation pour le couple $(TC,F\;=\;(TC)_*)$. À ce moment là, les conditions préliminaires (4.3.1), (4.3.2) de la proposition 4.3 sont satisfaites :
\[
TC\;\subset C\ ,\qquad TF\;\subset F\ ,\qquad W\;\supset\;TC\ ,\ TF\ ,
\]
ainsi que la condition $(a)$, $C\cap W\subset\widetilde{TC}$, puisque $C\cap W=TC$. Donc on trouve, par l'équivalence des conditions $(a)$, $(a')$, $(c)$, et vu que $TC$ et $TF$ sont clos,
\[
TF\;=\;F\cap W\ ,\qquad W\;=\;TF\circ TC\ ,
\]
\printedpage{26}
% head: A. Grothendieck : Dérivateurs, Ch. XV, Édition des Universités de Montpellier II et Paris VII
% foot: Version 5 novembre 2010
\msp{page 38}
d'autre part l'hypothèse sur $W$ implique que $W$ est stable par facteurs directs, $C$ l'est aussi, étant close à gauche, donc $TC=C{\cap}W$ aussi, donc par proposition 4.1, corollaire 2, on a $TC\;=\;\widetilde{TC}$, $TC$ est close à gauche. On gagne donc, par la dernière assertion du corollaire 3, proposition 4.3 (p. 35).

En fait, indépendamment des conditions particulières sur $\mathcal{M}$, $W$, on vient de prouver ceci :

\textbf{Lemme 4.4.1} (corollaire de la proposition 4.3.). \textit{Soient $\mathcal{M}$ une catégorie, $W$ un localiseur dans $\mathcal{M}$ modérément saturé et stable par facteurs directs. Alors les structures de Quillen closes $(C,F)$ associées à $W$ s'obtiennent ainsi : on prend une partie $C_0$ de} $\Fl(\mathcal{M})$ \textit{telle que}
\[
TF\;\deq\;C_{0*}\;\subset\;W\ ,
\]
\textit{on pose}
\[
C\;=\;\widetilde{C_0}\;=\;TF^*\ ,\qquad TC\;=\;C\cap W\ ,\qquad F\;=\;(TC)_*\ ,
\]
\textit{et on doit se débrouiller pour que $(C,TF)$ et $(TC,F)$ satisfassent à la condition de factorisation de Quillen,} i.e. \textit{soient des précouples de Quillen.} (En fait, $C$, $TF$ et $F$ étant clos par construction, et $TC=C\cap W$ étant stable par facteurs directs ($C$ et $W$ l'étant), ce seront même des \textit{couples} de Quillen.)

\msp{page 39}
Les hypothèses supplémentaires faites tantôt sur $\mathcal{M}$ servaient à assurer que si je prends $C_0$ petite, alors la factorisation marche pour $(C,TF)$ automatiquement. Il reste le problème pour $TC$, $F$.

Notons que si elle marche, alors $TC=\widetilde{TC}=F^*$, donc $TC$ est clos à gauche, et \textit{a fortiori} satisfait les conditions de stabilité correspondantes :
\[
\left\{
\begin{tabular}{rp{0.6\textwidth}}
\textit{a}) & Stable par cochangement de base, par composition, par facteurs directs (ces deux dernières sont assurées automatiquement, car $C$, $W$ sont stables).\\[1ex]
\textit{b}) & Stable par $\varinjlim$ ordinales strictes.
\end{tabular}
\right.
\tag*{(4.4.3)}
\]

L'hypothèse de stabilité sur $W$ sert (entre autres) à garantir la stabilité $(b)$ pour $TC$ (en fait, il suffirait pour cela seulement que $W$ soit stable par $\varinjlim$ ordinales strictes). Il reste la stabilité de $TC=C{\cap}W$ par cochangement de base, qui est problématique car $W$ n'est en général \textit{pas} stable par cochangement de base. C'est pour assurer cette stabilité que je me vois pratiquement contraint ici à imposer
\[
C_0\;\subset\;\Cof_W\ , \tag*{(4.4.4)}
\]
je dis qu'avec l'hypothèse de stabilité faite sur $W$, cela implique
\[
C\;\deq\;\widetilde{C_0}\subset\Cof_W\ . \tag*{(4.4.5)}
\]
\printedpage{27}
% head: A. Grothendieck : Dérivateurs, Ch. XV, Édition des Universités de Montpellier II et Paris VII
% foot: Version 5 novembre 2010
\msp{page 40}
En effet, on a :

\textbf{Proposition 4.4.2.} \textit{Soit $\mathcal{M}$ stable par petites $\varinjlim$ inductives filtrantes, et $W\subset$} $\Fl(\mathcal{M})$ \textit{un localiseur $($modérément saturé$)$} stable par lesdites limites. \textit{Alors} $\Cof_W$ \textit{satisfait les conditions de stabilité} (4.4.3) (“\textit{stabilité de Quillen gauche}”), \textit{et est stable aussi par $\varinjlim$ filtrantes quelconques.}

En effet, la condition $(a)$ est satisfaite sans condition supplémentaire sur $W$. D'autre part, il est immédiat que cette condition par rapport aux $\varinjlim$ filtrantes implique que $\Cof_W$ \textit{est stable par limites inductives filtrantes} dans $\underline{\Fl}(\mathcal{M})$.

\textbf{Corollaire 4.4.3.} \textit{Sous les conditions de la proposition} 4.4.2, $\Cof_W\cap W=W^{\univ}$ $($\textit{ensemble des $W$-équivalences co-universelles dans $\mathcal{M}$}$)$ \textit{est stable par limites inductives filtrantes, et donc “stable à gauche” en sens de} (4.4.3).

\textbf{Corollaire 4.4.4.} \textit{Sous les conditions de la proposition} 4.4.2, \textit{supposons de plus que les sommes amalgamées simples existent dans $\mathcal{M}$, et soit $C_0$ une petite partie de} $\Cof_W$ \textit{dont les éléments soient accessibles, et considérons}
\[
C\;=\;\widetilde{C_0}\ ,\qquad TC\;\deq\;C\cap W\ .
\]
\textit{Alors on a}
\[
C\;=\;\Cof_W\ ,\qquad TC\;=\;C\cap W^{\univ}\;\subset\;W^{\univ}\ ,
\]
\msp{page 41}
\textit{et $C$ et $TC$ sont $Q$-stables à gauche.}

En effet, par le théorème 1 du paragraphe 1 (p. 10) et son corollaire, $\widetilde{C_0}$ est la plus petite partie de $\Fl(\mathcal{M})$ contenant $C_0$ et Q-stable à gauche. Comme $\Cof_W$ est Q-stable à gauche et contient $C_0$, on a donc $C=\widetilde{C_0}\subset\Cof_W$. Donc $TC\deq C{\cap}W\subset\Cof_W{\cap}W\deq W^{\univ}$, donc $TC=C\cap W^{\univ}$. Comme $W^{\univ}$ est lui aussi $Q$-stable à gauche (corollaire précédent), il en est de même de $C\cap W$, \quad \textit{q.e.d.}

\textbf{Corollaire 4.4.5.} \textit{Supposons que tout objet de $\mathcal{M}$ soit accessible. Sous les conditions du corollaire précédent, posant}
\[
TF\;=\;C_{0*}\;=\;C_*\ ,\qquad F\;=\;(TC)_*\ ,
\]
\textit{le couple $(C,TF)$ est un couple de Quillen. Pour qu'il en soit de même du couple $(TC,F)$, il suffit qu'il existe une} petite \textit{partie $TC_0$ de $C\cap W$, telle que $\widetilde{TC_0}=TC$ $(\deq C\cap W)$.}

La première assertion est, on l'a vu tantôt, le théorème 1 du paragraphe 1. La deuxième résulte de ce même théorème. On peut préciser dans ce sens :

\textbf{Proposition 4.4.6.} \textit{Soit $\mathcal{M}$ une catégorie stable par petites $\varinjlim$ filtrantes, par sommes amalgamées, et dont tous les objets sont accessibles, $W$ un localiseur}
\msp{page 42}
\textit{modérément saturé dans $\mathcal{M}$, stable par $\varinjlim$ filtrantes, $C_0$ une} petite \textit{partie de} $\Fl(\mathcal{M})$, \textit{telle que}
\[
C_0\;\subset\;\Cof_W\ ,\qquad \underbrace{C_{0*}}_{(=TF)}\;\subset\;W\ ,
\]
\printedpage{28}
% head: A. Grothendieck : Dérivateurs, Ch. XV, Édition des Universités de Montpellier II et Paris VII
% foot: Version 5 novembre 2010
\textit{posons}
\[
C\;=\;\widetilde{C_0}\qquad TC\;=\;C\cap W\ ,\qquad TF\;=\;C_*\ ,\qquad F\;=\;(TC)_*\ ,
\]
\textit{et soit $TC_0$ une petite partie de $TC=C\cap W$. On a alors}
\[
\left\{
\begin{aligned}
C &\;\subset\;\Cof_W\ , &\qquad TC &\;=\;C\cap W^{\univ}\quad(\textit{pour mémoire})\\
\widetilde{TC_0} &\;\subset\;TC\ , & F_0 &\supset F\quad(\textit{ce dernier trivial})\ ,
\end{aligned}
\right.
\tag*{(4.4.6)}
\]
\textit{et les conditions suivantes sont équivalentes :}
\begin{enumerate}[leftmargin=*]
\item[a)] $\widetilde{TC_0}=TC\quad(\textit{on a}\subset)$\,;
\item[a$'$)] $F_0=F\quad(\textit{on a}\supset)$\,;
\item[b)] $TF=F_0\cap W\quad(\textit{on a}\subset)$\footnote[14]{\textbf{N.B.} On voit facilement que $\boxed{TF=F\cap W,}$ mais ça ne suffit pas pour avoir $TF=F_0\cap W$ pour $TC_0$ convenable.}\,;
\item[c)] $W=TF\circ\widetilde{TC_0}\quad(\textit{on a}\supset)$\footnote[15]{\textbf{N.B.} On sait que $\boxed{W=TF\circ TC,}$ ce qui n'est nullement suffisant pour impliquer $(c)$\,!}\,;
\end{enumerate}
\textit{et impliquent que $(W,C,F)$ est un triple de Quillen clos.}

\textsc{Démonstration.} L'inclusion $\widetilde{TC_0}\;\subset\;TC$ résulte de la Q-stabilité à gauche de $TC$ (cor. 4.4.4) et du fait que $TC_0\subset TC$ est petite (on applique alors le théorème 1 tout comme pour $C\subset\Cof_W\ldots$). L'équivalence de $(a)$, $(b)$, $(c)$ n'est autre que la proposition 4.3 appliquée à $W$ et aux deux couples de Quillen $(C,TF)$, $(\widetilde{TC_0},F_0)$ (compte tenu du théorème 1 du paragraphe 1). Qu'on a alors un triple de Quillen clos a été vu dans le corollaire 3 de la proposition 4.3 (p. 35). Il reste à prouver l'équivalence de $(a)$ et de $(a')$. On a $(a)\Longrightarrow(a')$, puisque $(TC_0)_*=(\widetilde{TC_0})_*=F_0$. Inversement, $(a')$ implique $\widetilde{TC_0}=\widetilde{TC}$, donc le théorème de factorisation pour $(TC,F)$ en vertu du théorème de factorisation, corollaire 1 (p. 11) appliqué à $TC_0\subset TC\subset\widetilde{TC_0}$ et à $F_0=F=(TC_0)_*$.

\msp{page 43}
Ainsi, le seul problème, une fois trouvé $C_0$ (petite) satisfaisant $C_0\subset\Cof_W$, $C_{0*}\subset W$, c'est de trouver une \textit{petite} partie $TC_0\subset TC$ qui “engendre $TC\;=\;C\cap W$ ”, \textit{i.e.} qui satisfasse aux conditions équivalentes $(a)$, $(a')$, $(b)$, $(c)$ de la proposition précédente. J'ai passé des jours à essayer de prouver l'existence de $TC_0$, en supposant $\mathcal{M}$ munie d'une filtration cardinale, et $W$ une partie accessible de $\mathcal{M}$, sans y être arrivé. Bien sûr, tout ce qui nous importe vraiment, c'est que $(TC,F)$ satisfasse à la condition de factorisation. Mais je ne vois aucun moyen de le prouver, si ce n'est en “engendrant” à gauche $TC$ par un \textit{petit} sous-ensemble $TC_0$.

\section*{5\quad Résultats techniques préliminaires sur les parties accessibles de $\mathrm{Ob}\,\mathcal{M}$, stables par limites inductives filtrantes}

Je crois pourtant que je peux arriver à quelque chose d'utilisable, en travaillant systématiquement avec des $C$ ($\subset\Cof_W$, Q-closes à gauche, Q-engendrées à gauche par une \textit{petite} partie $C_0$ de $C$) qui sont stables par toutes limites inductives filtrantes.


\endgroup
% END UNIT Der-ChXV-05-11-10/Der-ChXV-05-11-10_p026-030.tex
% BEGIN UNIT Der-ChXV-05-11-10/Der-ChXV-05-11-10_p031-035.tex
\begingroup\makeatletter%
\newcommand{\printedpage}[1]{\marginpar{\footnotesize p.~#1}}%
\newcommand{\msp}[1]{\marginpar{\footnotesize [#1]}}%
\newcommand{\cM}{\mathcal{M}}%
\newcommand{\cC}{\mathcal{C}}%
\newcommand{\Filt}{\mathrm{Filt}}%
\newcommand{\uFilt}{\underline{\mathrm{Filt}}}%
\newcommand{\Ss}{\mathrm{Ss}}%
\newcommand{\Ob}{\mathrm{Ob}\,}%
\newcommand{\Img}{\mathrm{Im}\,}%
\newcommand{\Fl}{\mathrm{Fl}}%
\newcommand{\uFl}{\underline{\mathrm{Fl}}}%
\newcommand{\card}{\mathrm{card}\,}%
\newcommand{\isoto}{\xrightarrow{\ \sim\ }}%
\newcommand{\hookto}{\lhook\joinrel\longrightarrow}%
\newcommand{\circled}[1]{\tikz[baseline=(c.base)]\node[draw,circle,inner sep=1pt](c){#1};}%
\makeatletter%
\renewcommand\@makefnmark{\hbox{\textsuperscript{\normalfont(\@thefnmark)}}}%
\renewcommand\@makefntext[1]{\parindent 1em\noindent\hbox to 1.8em{\hss\textsuperscript{\normalfont\@thefnmark}}#1}%
\makeatother%
\makeatother\let\printedpage\derpp\let\msp\dermsp%
\let\mspin\DERmspin%


\printedpage{29}
% head: A. Grothendieck : Dérivateurs, Ch. XV, Édition des Universités de Montpellier II et Paris VII
% foot: Version 5 novembre 2010

Je suppose à nouveau sur $\cM$ les conditions suivantes :
\msp{page 44}

\begin{enumerate}[nosep]
\item[\textit{a})] $\cM$ stable par petites limites inductives (et pas seulement filtrantes).
\item[\textit{b})] $\cM$ muni d'une filtration cardinale $(\Filt^{\pi})_{\pi\geq\pi_0}$.
\item[\textit{c})] Les objets de $\cM$ sont accessibles (\textbf{N.B.} Il suffit de le supposer pour les objets de $\Filt^{\pi_0}$).
\end{enumerate}

La condition ($b$) implique que $\cM$ contient une petite famille génératrice par épimorphismes stricts (p. ex. une sous-catégorie $\cC$ petite de $\uFilt^{\pi_0} = \cM^{\pi_0}$ telle que l'inclusion $\cC \longrightarrow \uFilt^{\pi_0}$ soit une équivalence de catégories). Par SGA 4 I 8.12.7, il en résulte que $\cM$ est stable par petites limites projectives. Il s'ensuit le

\textbf{Lemme 5.1} (sous les conditions ($a$), ($b$)). \textit{Pour tout objet $X$ de $\cM$, l'ensemble} $\Ss(X)$ \textit{des sous-objets de $X$, muni de la relation d'ordre naturelle ``d'inclusion'', admet des} inf \textit{et des} sup \textit{quelconques.}

Des inf, car ceux-ci s'expriment comme des limites projectives dans $\cM$. Or l'existence des inf implique celle des sup, puisque $\sup_J = \inf_{J'}$, où $J'$ est l'ensemble des majorants de $J$ dans $\Ss(X)$.

\textbf{Corollaire 5.2.} \textit{Pour toute flèche $f : Z \longrightarrow X$ dans $\cM$, il existe, parmi les sous-objets $X'$ de $X$ par lesquelles $f$ se factorise, un plus petit} % [sic] lesquelles -> lesquels, "sous-objets" is masculine
\msp{page 45}
\textit{(qu'on notera} $\Img f$\textit{). Le morphisme $Z \longrightarrow \Img f$ est un épimorphisme.}

Cette dernière assertion provient de celle-ci :

\textbf{Lemme 5.3.} \textit{Soient $\cM$ une catégorie où les noyaux des doubles flèches existent. Soit $f : Z \longrightarrow Y$ une flèche de $\cM$. Alors $g$ est un épimorphisme si et seulement si $g$ ne se factorise pas par un sous-objet} strict $Y_0$ \textit{de $Y$ différent de $Y$.} % [sic] g -> f, the arrow has been named f

On va faire une hypothèse supplémentaire (pratiquement anodine) sur les $\Filt^{\pi}$ :

\begin{enumerate}[nosep]
\item[\textit{d})] $\Filt^{\pi}$ est stable par limites inductives indexées par des catégories d'indices $I$ telles que $\card \Fl(I) \leq \pi$\footnote[16]{Cette condition sera nommée ($d_0$) par la suite.}. (Il suffit de supposer la stabilité par objet initial, sommes amalgamées simples, et $\varinjlim$ filtrantes indexées par un ensemble ordonné $I$ de cardinal $\leq \pi$.) De plus, si $f : Z \longrightarrow X$ est une flèche de $\cM$ telle que $Z$ ou $X$ est dans $\Filt^{\pi}$, alors $\Img f \in \Filt^{\pi}$.
\end{enumerate}

Cette dernière condition sur $\Img f$ se décompose en deux :
\msp{page 46}

\begin{enumerate}[nosep]
\item[$d_1$)] Tout sous-objet d'un objet de $\Filt^{\pi}$ est dans $\Filt^{\pi}$.
\item[$d_2$)] Si $Z$ est dans $\Filt^{\pi}$, alors tout quotient $Z'$ de $Z$ tel que $\Img p = Z'$ (où $p : Z \longrightarrow Z'$ est le morphisme canonique) est dans $\Filt^{\pi}$.\footnote[17]{Il semblerait qu'on n'aura pas besoin de ($d_1$), seulement de ($d_2$). Mais en fait, ($d_1$) doit être conséquence de ($d_2$), et ($d_2$) est conséquence de la première partie de ($d$) (stabilité par $\varinjlim$ finies) si tout épimorphisme dans $\cM$ est strict, donc effectif (si les $\Filt^{\pi}$ sont stables par $\varprojlim$ finies aussi \ldots).}
\end{enumerate}

\printedpage{30}
% head: A. Grothendieck : Dérivateurs, Ch. XV, Édition des Universités de Montpellier II et Paris VII
% foot: Version 5 novembre 2010

Bien sûr, pratiquement il sera toujours vrai que $\Filt^{\pi}$ est stable par passage aux sous-objets et aux quotients d'un objet (pas seulement les quotients particuliers dont il est question dans ($d_2$). % [sic] closing parenthesis of "(pas seulement" missing, parentheses unbalanced as printed
D'ailleurs, lorsque tout bimorphisme (\textit{i.e.} monomorphisme \textit{et} épimorphisme) de $\cM$ est un isomorphisme, alors ($d_2$) équivaut bien à la stabilité de $\Filt^{\pi}$ par quotients d'objets.

\textbf{Proposition 5.3.} \texttt{[Deux fois 5.3.]} \textit{Supposons les conditions} (a), (b), (d) \textit{vérifiées pour $\cM$. Soient $\pi$ un cardinal $\geq \pi_0$, et $X$ un objet de $\cM$. Soit} $\Ss_{\pi}(X)$ \textit{l'ensemble des sous-objets de $X$ qui sont dans} $\Filt^{\pi}$\textit{. Cet ensemble, ordonné par la relation d'ordre induite par} $\Ss(X)$\textit{, est grand devant $\pi$ (et en particulier filtrant). Si de plus, les limites inductives filtrantes de monomorphismes dans $\cM$ sont des monomorphismes}\footnote[18]{C'est une \textit{condition d'exactitude} de $\cM$, que je vais noter comme :\par
\noindent\circled{e}\quad Les monomorphismes dans $\cM$ sont stables par $\varinjlim$ filtrantes dans $\uFl(\cM)$.} \textit{(p. ex. si les limites inductives filtrantes sont exactes à gauche), ou si les objets de} $\Filt^{\pi}$ \textit{sont $\pi$-accessibles, on a $\varinjlim\limits_{\alpha\in\Ss_{\pi}(X)} Z_{\alpha} \isoto X$}\footnote[19]{Dire aussi que $\Ss_{\pi}(X)$ est $\pi$-adapté, et le système inductif $(Z_{\alpha})_{\alpha}$ aussi.}.

\msp{page 47}
\textsc{Démonstration}. Soit $(Z_{\alpha})_{\alpha\in A}$ une famille de sous-objets de $X$ tels que $Z_{\alpha} \in \Filt^{\pi}$, et $\card A \leq \pi$. Soit $Z = \coprod_{\alpha\in A} Z_{\alpha}$, on a $Z \in \Filt^{\pi}$ par ($c$). Le plus petit sous-objet de $X$ majorant les $Z_{\alpha}$ n'est autre que $\Img (Z \longrightarrow X)$, qui est dans $\Filt^{\pi}$ par ($d_1$). Prouvons que $\varinjlim\limits_{\alpha\in\Ss_{\pi}(X)} Z_{\alpha} \isoto X$ est un isomorphisme, quand on sait que c'est un monomorphisme. Or
\[
X \;=\; \varinjlim_{I} X_i \ ,
\]
où $(X_i)_{i\in I}$ est un système inductif filtrant, $I$ ensemble ordonné grand devant $\pi$, les $X_i$ dans $\Filt^{\pi}$. Soit, pour tout $i \in I$, $X'_i$ l'image de $X_i$ dans $X$. On a donc $X'_i \in \Filt^{\pi}$ par ($d_2$). Utilisant le fait que les $X_i \longrightarrow X'_i$ sont des épimorphismes, on voit de suite que $X = \varinjlim\limits_{I} X'_i$. Considérons l'application
\[
\begin{aligned}
I \quad &\longrightarrow \quad \Ss_{\pi}(X)\\
i \quad &\longmapsto \quad X'_i \ .
\end{aligned}
\]
C'est une application croissante, soit $J$ son image, partie filtrante grande devant $\pi$ de $\Ss_{\pi}(X)$, telle que
\[
\varinjlim_{\alpha\in J} Z_{\alpha} \ \isoto \ X \ .
\]
\msp{page 48}
On a donc
\[
\underbrace{\varinjlim_{\alpha\in J} Z_{\alpha}}_{\simeq X} \ \longrightarrow \ \varinjlim_{\alpha\in\Ss_{\pi}X} Z_{\alpha} \ \xrightarrow{\ \varphi\ } \ X \ ,
\]
le composé étant l'identité. Si on sait que $\varphi$ est un monomorphisme, il en résulte que c'est un isomorphisme (monomorphisme ayant une section). Or $\varphi$ sera un monomorphisme si toute $\varinjlim$ filtrante de monomorphismes est un monomorphisme, et alors on gagne.

\printedpage{31}
% head: A. Grothendieck : Dérivateurs, Ch. XV, Édition des Universités de Montpellier II et Paris VII
% foot: Version 5 novembre 2010

Si on suppose que les éléments de $\Filt^{\pi}$ sont $\pi$-accessibles, on gagne aussi, car on trouve :

\textbf{Corollaire 5.4.} \textit{Si les objets de} $\Filt^{\pi}$ \textit{sont $\pi$-accessibles (p. ex. si $\pi \geq \pi_1$ et $\pi = \pi^{\pi_0}$, p. ex. $\pi \geq \pi_1$ et $\pi = 2^{c}$, avec $c \geq \pi_0$,} cf. (3.2) \textit{p. 16), alors pour toute partie filtrante $J$ de} $\Ss(X)$ \textit{grande devant $\pi$, telle que $\varinjlim\limits_{\alpha\in J} Z_{\alpha} \isoto X$, $J$ est cofinale dans} $\Ss_{\pi}(X)$ \textit{(donc $\varinjlim\limits_{\alpha\in J} Z_{\alpha} \isoto \varinjlim\limits_{\alpha\in\Ss(X)} Z_{\alpha}$).}

C'est immédiat.

\msp{page 49}
\textbf{Scholie 5.5.} Quand on fait sur $\cM$ l'hypothèse d'exactitude assez anodine que les $\varinjlim$ filtrantes de monomorphismes sont des monomorphismes, \textit{ou} qu'on se borne aux cardinaux $\pi$ tels que les objets de $\Filt^{\pi}$ soient $\pi$-accessibles, on peut \textit{canoniser} la représentation d'un objet $X$ comme $\varinjlim\limits_{i\in I} X_i$, avec $X_i \in \Filt^{\pi}$ et $I$ grand devant $\pi$, en prenant $I = \Ss_{\pi}(X)$, et le système inductif correspondant. On a de plus l'agréable propriété supplémentaire que voici : toute partie filtrante $J$ de $I$, de cardinal $\leq \pi$, admet une borne supérieure $i_J$, et (si on suppose que $\cM$ satisfait la condition ($e$), page 46, note 18), on aura
\[
X_{i_J} \;=\; \varinjlim_{i\in J} X_i \ ;
\]
en d'autres termes, on a une représentation de $X$ par un système inductif (parfaitement) \textit{$\pi$-adapté} (\textit{cf.} XIII, p. 85, 86).

\msp{page 50}
Soit maintenant
\[
M \;\subset\; \Ob\cM
\]
une partie de $\Ob\cM$, stable par la relation d'isomorphie. Si $M$ est $\pi$-accessible (suivant la condition $\mathrm{b}_{\pi}$ de XIII, p. 88), alors en posant
\[
M^{\pi} \;=\; M \cap \Filt^{\pi} \ ,
\]
$M^{\pi}$ détermine $M$ par la condition suivante sur $X \in \Ob\cM$.
\[
\begin{aligned}
&X \in M \iff \text{le sous-ensemble } \Ss^{M}_{\pi}(X) \text{ de } \Ss_{\pi}(X) \text{ formé des } X_{\alpha} \in \Ss_{\pi}(X)\\
&\text{tels que } X_{\alpha} \in M \ (\iff X_{\alpha} \in M^{\pi}) \text{ est cofinal dans } \Ss_{\pi}(X).
\end{aligned}
\tag*{(5.1)}
\]
Notons que lorsque $X \in \Filt^{\pi}$, ce critère est trivialement valable, puisque alors $\Ss_{\pi}(X)$ admet $X$ comme plus grand élément, et les sous-ensembles cofinaux sont ceux qui contiennent cet objet final. Donc (5.1) est intéressant seulement pour les $X \in \Ob\cM \backslash \Filt^{\pi}$. Notons que pour \textit{toute} partie $M^{\pi}$ de $\Filt^{\pi}$, on peut définir une partie $M$ de $\Ob\cM$ par la condition (5.1), et on aura trivialement
\[
M^{\pi} \;=\; M \cap \Filt^{\pi} \ . \tag*{(5.2)}
\]
Mais sans doute, bien que $M$ satisfasse (5.1), il n'est pas pour autant nécessairement $\pi$-accessible. En effet, pour qu'il le soit, il faut qu'il soit stable par $\varinjlim$
\printedpage{32}
% head: A. Grothendieck : Dérivateurs, Ch. XV, Édition des Universités de Montpellier II et Paris VII
% foot: Version 5 novembre 2010
\msp{page 51}
grandes devant $\pi$ d'éléments de $M^{\pi}$, ce qui implique que tout facteur direct d'un élément de $M^{\pi}$ est dans $M$, donc aussi dans $M^{\pi}$, d'après la condition de stabilité ($d$) sur $\Filt^{\pi}$. C'est là une condition particulière sur la partie $M^{\pi}$ de $\Filt^{\pi}$. Aussi, on devra avoir une condition sur $\cM$ plus forte que la validité de (5.1), savoir :
\[
\parbox{0.72\textwidth}{Pour tout $X \in \Ob\cM$ et pour toute partie $J$ de $\Ss_{\pi}(X)$, telle que
\begin{enumerate}[nosep]
\item[\textit{a})] $J$ est cofinale dans $\Ss_{\pi}(X)$ et
\item[\textit{b})] Si $J' \subset J$ est une partie filtrante de $J$, de cardinal $\leq \pi$, alors sa borne supérieure $\alpha_{J'}$ dans $\Ss_{\pi}(X)$ appartient à $J$,
\end{enumerate}
on a $X \in M$ si et seulement si $J^{M} \overset{\text{déf}}{=} J \cap \Ss^{M}_{\pi}(X)$ est cofinal dans $J$ (ou ce qui revient au même, dans $\Ss_{\pi}(X)$)\,\footnotemark[20].}
\tag*{(5.3)}
\]
\footnotetext[20]{Au lieu de $J$ cofinale dans $\Ss_{\pi}(X)$, il suffit même que l'on ait $\varinjlim\limits_{\alpha\in J} X_{\alpha} \isoto X$ -- la condition correspondante sur $M$ devient donc plus exigeante. Mais si les objets de $\Filt^{\pi}$ sont $\pi$-accessibles, cette condition est en fait équivalente à la condition plus faible.}

Si on admet déjà le ``il suffit'' de (5.1), alors il suffit dans (5.3) de supposer $X \in M$ et exiger qu'alors $J^{M}$ soit cofinal dans $J$.

Mais notons la
\msp{page 52}

\textbf{Proposition 5.6.}\footnote[21]{Faux tel quel, vrai si on suppose que $M^{\pi}$ est stable par $\varinjlim$ filtrantes.} \textit{Supposons $\pi$ choisi de telle façon que les objets de} $\Filt^{\pi}$ \textit{soient $\pi$-accessibles (p. ex. $\pi = 2^{c}$, avec $c \geq \pi_1$,} cf. (3.2), \textit{p. 16). Alors, pour une partie $M$ de} $\Ob\cM$\textit{, $M$ est $\pi$-accessible si et seulement si elle satisfait au critère} (5.3).

On a vu que c'est nécessaire (tautologiquement, compte tenu de la proposition 5.3 et du scholie 5.5, bas de la page 49). Prouvons que c'est suffisant, en établissant d'abord le

\textbf{Lemme 5.7.} \textit{Sous les conditions de} 5.6 \textit{sur $\pi$, et la condition} (e) \textit{sur $\cM$} (p. 46), \textit{pour toute représentation $\pi$-adaptée $X = \varinjlim\limits_{i\in I} X_i$ d'un objet $X$ de $\cM$, le sous-ensemble $I'$ des $i \in I$ tels que $X_i \longrightarrow X$ soit un monomorphisme, est cofinal dans $I$. De plus, $I'$ est stable par bornes supérieures dans $I$ de parties filtrantes $J$ de $I'$ de cardinal $\leq \pi$.}

Montrons comment ce lemme 5.7 implique le ``il suffit'' de la proposition 5.6. Soit donc $X = \varinjlim\limits_{i\in I} X_i$ une représentation $\pi$-adaptée de $X$ ($X_i \in \Filt^{\pi}$, $I$ grand devant $\pi$, \textit{etc.}). Il faut prouver que, si 5.3 est valable, alors $X \in M$ si et seulement si l'ensemble $I_M$ est cofinal dans $I$.

\msp{page 53}
Supposons $X \in M$. Considérons le sous-ensemble $I' \subset I$ décrit dans le lemme. On a une application $I' \longrightarrow \Ss_{\pi}(X)$, soit $i \longmapsto X'_i = \Img (X_i \hookto X)$, et on a
\[
X \simeq \varinjlim_{i\in I} X_i \ \simeq \ \varinjlim_{i\in I'} X_i \ \simeq \ \varinjlim_{i\in I'} X'_i \ \simeq \ \varinjlim_{\alpha\in J} Z_{\alpha} \ ,
\]
où $J$ est l'image de $I'$ dans $\Ss_{\pi}(X)$. Comme les objets de $\Filt^{\pi}$ sont $\pi$-accessibles, on en conclut que $J$ est cofinal dans $\Ss_{\pi}(X)$. Par 5.3, il s'ensuit que l'ensemble $J_M$ des $\alpha \in J$
\printedpage{33}
% head: A. Grothendieck : Dérivateurs, Ch. XV, Édition des Universités de Montpellier II et Paris VII
% foot: Version 5 novembre 2010
tels que $Z_{\alpha} \in M$ est cofinal dans $J$, donc $\neq \emptyset$. Donc $I'_M$, qui est l'image inverse de $J_M$ dans $I'$, est en tous cas non vide. Appliquant le même résultat en remplaçant $I'$ par le sous-ensemble $I'_{\geq i'_0}$, on voit que $I'_M$ est cofinal dans $I'$, donc aussi dans $I$, puisque $I'$ est cofinal dans $I$ ; donc $I_M \supset I'_M$ est cofinal dans $I$, \quad \textit{q.e.d.}

Inversement, supposons que $I_M$ soit cofinal dans $I$, et prouvons $X \in M$. La difficulté ici provient du fait que la représentation $X = \varinjlim\limits_{i\in I_M} X_i$ n'est plus adaptée à $\pi$ -- $I_M$ n'est pas nécessairement stable dans $I$ par bornes supérieures (dans $I$) de parties filtrantes dans $I_M$, $I'$ de cardinal $\leq \pi$. En somme, il faut
\msp{page 54}
prouver que $M$ est stable par $\varinjlim$ grandes devant $\pi$ sans plus, et je doute maintenant que ça résulte des seules hypothèses faites ici, qui ne concernent que les systèmes inductifs où les morphismes de transition sont des monomorphismes.

Voici, je crois, un contre-exemple. Prenons
\[
M \quad = \quad \parbox[t]{0.6\textwidth}{ensemble des objets de $\cM$ isomorphes à un objet fixe $X_0 \in \Filt^{\pi}$, tel que $X_0$ n'ait pas d'endomorphisme qui soit un monomorphisme autre que l'identité.}
\tag*{(5.4)}
\]
Alors la condition 5.3 est satisfaite. Car :
\begin{enumerate}[nosep]
\item[\textit{a})] Si $X \in M$, \textit{i.e.} $X \simeq X_0$, alors $X \in \Filt^{\pi}$, donc $\Ss_{\pi}(X)$ admet $X$ comme objet final, et $J$, étant cofinal, contient cet élément, donc $J_M$ (qui le contient aussi) est cofinal.
\end{enumerate}
Et :
\begin{enumerate}[nosep]
\item[\textit{b})] Si $J_M$ est cofinal, l'hypothèse (5.4) implique que l'ensemble des morphismes de transition entre les $Z_{\alpha}$ ($\alpha \in J_M$) est formé d'isomorphismes, donc la limite inductive $X$ des $Z_{\alpha}$ en question est isomorphe à chacun d'eux, donc $X \simeq X_0$, \textit{i.e.} $X \in M$.
\end{enumerate}
Pourtant, $M$ n'est $\pi$-accessible que s'il est stable par $\varinjlim$ grandes devant $\pi$, ce qui implique qu'il est stable par \textit{facteurs directs.} Mais l'hypothèse
\msp{page 55}
(5.4) n'exclut nullement que $X_0$ admette un projecteur $p$ ($p^{2} = p$) non trivial. Celui-ci aura une image $X'$ dans $\cM$, qui sera facteur direct de $X_0$ non isomorphe à $X_0$, \textit{i.e.} non dans $M$.

Ces difficultés disparaissent si on suppose que $M^{\pi}$ est stable par $\varinjlim$ filtrantes indexées par des ensembles ordonnés $I$ de cardinal $\leq \pi$. En effet, dans ce cas, dans la démonstration précédente, le sous-ensemble $I_M$ est stable par bornes supérieures (dans $I$) de parties $J$ de $I_M$ filtrantes et de cardinal $\leq \pi$. Donc le système inductif $(X_i)_{i\in I_M}$ est encore $\pi$-adapté, et on peut lui appliquer le lemme 5.7, pour trouver un sous-ensemble cofinal $I'_M \subset I_M$, tel que les $X_i \longrightarrow X$ pour $i \in I'_M$ soient des monomorphismes. L'image $J$ de $I'_M$ dans $\Ss_{\pi}(X)$ est alors une partie de $\Ss_{\pi}(X)$ telle que $\varinjlim\limits_{\alpha\in J} Z_{\alpha} = X$, donc (par l'hypothèse faite sur $\pi$) cofinale dans $\Ss_{\pi}(X)$ (5.4), avec $Z_{\alpha} \in M^{\pi}$ pour tout $\alpha \in J$. Donc par 5.3 on a $X \in M$, \textit{q.e.d.}

\msp{page 56}
Il reste à prouver le lemme 5.7. Avec les hypothèses et notations du lemme, soit donc $i_0 \in I$, il faut trouver $i' \in I'$ tel que $i' \geq i_0$. Considérons l'image $X'_i$ de $X_i$ dans $X$, et


\endgroup
% END UNIT Der-ChXV-05-11-10/Der-ChXV-05-11-10_p031-035.tex
% BEGIN UNIT Der-ChXV-05-11-10/Der-ChXV-05-11-10_p036-040.tex
\begingroup\makeatletter%
\newcommand{\printedpage}[1]{\marginpar{\footnotesize p.~#1}}%
\newcommand{\msp}[1]{\marginpar{\footnotesize [#1]}}%
\newcommand{\Filt}{\mathrm{Filt}}%
\newcommand{\Ob}{\mathrm{Ob}}%
\newcommand{\Ss}{\mathrm{Ss}}%
\newcommand{\M}{\mathcal{M}}%
\newcommand{\Ci}{\mathcal{C}}%
\makeatother\let\printedpage\derpp\let\msp\dermsp%
\let\mspin\DERmspin%


\printedpage{34}
% head: A. Grothendieck : Dérivateurs, Ch. XV, Édition des Universités de Montpellier II et Paris VII
% foot: Version 5 novembre 2010
l'inclusion
\[
\underbrace{X'_{i_0}}_{\in\, \Filt^\pi} \;\longrightarrow\; X \;\simeq\; \varinjlim X_i \;.
\]
Comme $X'_{i_0} \in \Filt^\pi$, donc est $\pi$-accessible, il existe un $i_1 > i_0$ tel que $X'_{i_0} \longrightarrow X$ se factorise par
\[
X'_{i_0} \;\longrightarrow\; X_{i_1} \;.
\]
Appliquant ce même résultat à $X_{i_1}$, $X'_{i_1}$ \textit{etc.}, on trouve une suite infinie de morphismes
\[
\begin{tikzcd}[column sep=small]
X_{i_0} \arrow[dr] & & X_{i_1} \arrow[dr] & & X_{i_2} \arrow[dr] & \\
& X'_{i_0} \arrow[ur] & & X'_{i_1} \arrow[ur] & & X'_{i_2} \quad \cdots
\end{tikzcd}
\]
telle que les morphismes composés de deux en deux soient, sur la ligne supérieure, les transitions
\[
X_{i_m} \longrightarrow X_{i_n} \qquad (\; i_m \;<\; i_n \;) \;,
\]
et sur la ligne inférieure, les inclusions
\[
X'_{i_m} \hookrightarrow X'_{i_n} \qquad (\; \hookrightarrow X \;) \;.
\]
Soit $i_\omega = \sup_n i_n$ dans $I$, donc
\[
X_{i_\omega} \;=\; \varinjlim_{n} X_{i_n} \;\simeq\; \varinjlim_{n} X'_{i_n} \;\simeq\; X'_{i_\omega} \;,
\]
donc $X_{i_\omega} \longrightarrow X$ est un monomorphisme.

\msp{page 57}

La dernière assertion de 5.7 résulte aussitôt du fait qu'une limite inductive filtrante de monomorphismes dans $\M$ est un monomorphisme (condition $(e)$, p.~46).

On peut préciser la proposition 5.6 obtenue ainsi :

\textbf{Proposition 5.7}. \texttt{[Deux fois 5.7.]} \textit{Soient $\M$ une catégorie satisfaisant les conditions} (a), (b), (c), (d) (p.~44, 45) et (e) (p.~46). \textit{Soit $\M$ une partie de} $\Ob\M$\textit{, stable par isomorphie, $\pi$ un cardinal tel que les objets de} $\Filt^\pi$ \textit{soient $\pi$-accessibles. Les conditions suivantes sont équivalentes :} % [sic] Soit $\M$ (calligraphique) -> Soit $M$, la partie de Ob M est notée M partout ailleurs

\begin{enumerate}[label={}]
\item[a)] \textit{$M$ est $\pi$-accessible, et stable par $\varinjlim$ filtrantes.}
\item[b)] \textit{$M$ est $\pi$-accessible, et $M \cap \Filt^\pi$ est stable par $\varinjlim$ filtrante de cardinal $\leq \pi$.}
\item[c)] \textit{Il existe une partie $M^\pi$ de} $\Filt^\pi$\textit{, stable par $\varinjlim$ filtrantes de cardinal $\leq \pi$ et telle que pour tout $X \in \Ob\M$, on ait $X \in M$ si et seulement si l'ensemble $J_M$ des $Z_\alpha \in J$ tels que $Z_\alpha \in M^\pi$ est cofinal dans $J$.}
\end{enumerate}

(\textit{De plus, sous la condition} (c), \textit{on a}
\printedpage{35}
% head: A. Grothendieck : Dérivateurs, Ch. XV, Édition des Universités de Montpellier II et Paris VII
% foot: Version 5 novembre 2010
\msp{page 58}
$M^\pi = M \cap \Filt^\pi$.)
\begin{enumerate}[label={}]
\item[c$'$)] \textit{Il existe une partie $M_0$ de} $\Filt^\pi$ \textit{telle que $M$ soit la plus petite partie de $\M$ stable par limites inductives filtrantes.} (\textit{De plus, on peut alors prendre $M_0$} petite.)\footnote[22]{\textbf{N.B.} On a
\[
\begin{tikzcd}[ampersand replacement=\&]
(a) \arrow[r, Leftrightarrow] \arrow[d, Rightarrow] \& (b) \arrow[d, Rightarrow] \\
(c') \arrow[r, Rightarrow, "?"] \& (c)
\end{tikzcd}
\quad ,
\]
sans conditions sur $\pi$ (sauf $\pi \geq \pi_0$).}.
\end{enumerate}

\textbf{Corollaire 5.8.} \textit{Les applications}
\[
M \;\longmapsto\; M^\pi = M \cap \Filt^\pi \;,
\]
\[
M^\pi \quad\longmapsto\quad \parbox[t]{0.5\textwidth}{\textit{l'ensemble des $X \in \Ob\M$ tels que l'ensemble $\Ss_\pi^{M^\pi}(X)$ des $Z_\alpha \in \Ss_\pi(X)$ qui sont dans $M^\pi$ soit cofinal} (= \textit{plus petit sous-ensemble de $\Ob\M$ stable par $\varinjlim$ filtrantes})}
\]
\textit{établissent des bijections, inverses l'une de l'autre, entre l'ensemble des $M \subset \Ob\M$ satisfaisant les conditions équivalentes de la proposition} 5.7\textit{, et l'ensemble des parties $M^\pi$ de} $\Filt^\pi$ \textit{stables par $\varinjlim$ filtrantes de cardinal $\leq \pi$.}

\textsc{Démonstration de} 5.7. L'implication $(b) \Longrightarrow (c)$ est tautologique, compte tenu de la définition de la $\pi$-accessibilité de $M$, et de la proposition 5.3. Et $(a) \Longrightarrow (b)$ provient de la stabilité de $\Filt^\pi$ par $\varinjlim$ filtrantes de cardinal $\leq \pi$ (condition $(d)$).

Il reste donc à prouver $(c) \Longrightarrow (a)$. (C'est ici qu'on doit utiliser l'hypothèse particulière sur $\pi$.) Une fois ceci vu, on saura l'équivalence de $(a)$, $(b)$, $(c)$, et d'autre part, il est immédiat que $(a) \Longrightarrow (c')$, en
\msp{page 59}
prenant $M_0 = M \cap \Filt^\pi$. D'autre part, $(c') \Longrightarrow (c)$ (sans hypothèses sur $\pi\,$?), car on peut supposer que la partie $M_0$ de $\Filt^\pi$ est stable par $\varinjlim$ filtrantes de cardinal $\leq \pi$. Mais alors, connaissant l'équivalence de $(a)$, $(b)$ $(c)$, d'où aussitôt le corollaire, il s'ensuit que l'enveloppe stable par $\varinjlim$ filtrantes de $M_0$ dans $\M$ n'est autre que la partie accessible déduite de $M_0$ dont il est question dans $(c)$, et on gagne. % [sic] $(b)$ $(c)$ -> $(b)$, $(c)$, virgule manquante

Donc tout revient à prouver que, si tous les objets de $\Filt^\pi$ sont accessibles, alors $(c)$ implique $(a)$. Montrons d'abord que $(c)$ implique le critère d'appartenance (5.3) (p.~51), donc par la proposition 5.6 “revue et corrigée” dans la note 21, il s'ensuit déjà que $(c) \Longrightarrow (b)$, et il restera à établir $(b) \Longrightarrow (a)$, \textit{i.e.} que $(b)$ implique la stabilité de $M$ par $\varinjlim$ filtrantes quelconques (pas seulement grandes devant $\pi$).

Que $(c)$ implique (5.3) va résulter du lemme suivant, appliqué à $I = \Ss_\pi(X)$ :

\textbf{Lemme 5.8.} \texttt{[Deux fois 5.8.]} \textit{Soit $I$ un ensemble ordonné} $\pi$-adapté, i.e. \textit{filtrant et tel que toute partie filtrante $K$ de $I$, de cardinal $\leq \pi$, ait une borne supérieure $i_K$ dans $I$. Je dis qu'une partie $J$ de $I$ est} $\pi$-cofinale, \textit{si elle est cofinale, et si pour toute partie filtrante $K$ de $J$ de cardinal $\leq \pi$, sa borne supérieure $i_K$ dans $I$ est}
\printedpage{36}
% head: A. Grothendieck : Dérivateurs, Ch. XV, Édition des Universités de Montpellier II et Paris VII
% foot: Version 5 novembre 2010
\msp{page 60}
\textit{dans $J$.} (\textit{Il s'ensuit donc que $J$ est lui aussi $\pi$-adapté, pour la structure d'ordre induite.}) \textit{Ceci dit, l'intersection de deux parties $\pi$-cofinales de $I$ est encore $\pi$-cofinale.}

Il suffit de voir qu'elle est cofinale, car sa stabilité par bornes supérieures dans $I$ de parties filtrantes de cardinal $\leq \pi$ est tautologique (si elle est vraie pour des $J_\alpha$, elle est vraie pour $\bigcap_\alpha J_\alpha$). Soit $i_0 \in I$, et construisons un élément de $J \cap J'$ qui majore $i_0$. On construit de proche en proche une double suite
\[
\begin{tikzcd}
i_1 \arrow[dr] & & i_2 \arrow[dr] & & \cdots \\
& i'_1 \arrow[ur] & & i'_2 \arrow[ur] &
\end{tikzcd}
\]
avec les $i_\alpha$ dans $J$, les $i'_\alpha$ dans $J'$, et $i_1 \geq i_0$. On a donc
\[
\sup i_\alpha \;=\; \sup i'_\alpha \;,
\]
et si $i$ est la valeur commune des deux membres, on a $i \in J \cap J'$.

Ceci vu, pour prouver la validité de (5.3), on note que la partie $J' = \Ss_\pi^{(M^\pi)}(X)$ de $I = \Ss_\pi(X)$, formée des $Z_\alpha \subset X$ tels que $Z_\alpha \in M^\pi$, est stable par bornes supérieures de parties filtrantes de cardinal $\leq \pi$, vu que $M^\pi$ est stable
\msp{page 61}
par $\varinjlim$ de cardinal $\leq \pi$. Donc si $J$ est une partie $\pi$-cofinale de $\Ss_\pi(X)$, $J \cap J'$ est cofinale dans $J$, \textit{i.e.} dans $I$, si et seulement si $J'$ est cofinale dans $I$ (\textit{i.e.} si et seulement si $X \in M$, par la condition $(c)$). Car $J \cap J'$ cofinale $\Longrightarrow$ $J'$ cofinale tautologiquement, et $J'$ cofinale implique $J'$ $\pi$-cofinale, donc $J \cap J'$ $\pi$-cofinale par le lemme 5.8, donc cofinale.

Reste à prouver seulement que $(b) \Longrightarrow (a)$, \textit{i.e.} la stabilité de $M$ par limites inductives filtrantes quelconques, moyennant $(b)$. Ici à nouveau, cela ne doit pas utiliser l'hypothèse particulière sur $\pi$, ni même l'hypothèse d'exactitude $(e)$ sur $\M$, mais être une histoire générale d'accessibilité :

\textbf{Proposition 5.9.} \textit{Soit $\M$ une catégorie satisfaisant} (a), (b), (c), (d$_0$) (pages 44, 45), \textit{$\pi \geq \pi_0$ un cardinal, $M$ une partie $\pi$-accessible de} $\Ob\M$, i.e. \textit{satisfaisant la condition $b_\pi$ de} XIII, p.~88. \textit{Alors $M$ est stable par $\varinjlim$ filtrantes si et seulement si $M^\pi \overset{\text{déf}}{=} M \cap \Filt^\pi$ est stable par $\varinjlim$ filtrantes de cardinal $\leq \pi$.}

Je n'ai pas envie d'expliciter ici la démonstration, ayant l'impression pénible de refaire à perpette la même démonstration (ici et dans XIII). Je devrai revenir là-dessus dans une situation plus décantée.

\msp{page 62}

Je vais quand même esquisser la démonstration. Soit $X = \varinjlim\limits_{I} X_i$, $I$ filtrant, les $X_i$ dans $M$, il faut prouver que $X \in M$ (moyennant l'hypothèse de stabilité sur $M^\pi = M \cap \Filt^\pi$).
\printedpage{37}
% head: A. Grothendieck : Dérivateurs, Ch. XV, Édition des Universités de Montpellier II et Paris VII
% foot: Version 5 novembre 2010
\[
\begin{tikzcd}
& \Ci \arrow[r] \arrow[d] \arrow[dl, dotted, "\mathfrak{z}"'] \arrow[dr, phantom, "\text{\footnotesize cart.}"] & I \arrow[d, "\varphi"] \\
\underline{M}^\pi & \underline{M}^\pi/\underline{M} \arrow[l, dotted, "\text{\footnotesize “source”}"] \arrow[r, "\text{\footnotesize “but”}"'] & \underline{M}
\end{tikzcd}
\] % resolved: the PDF glyph is fraktur z; see corrections.csv, leaf 4.

$\Ci$ est une catégorie cofibrée sur $I$, dont les fibres sont
\[
\Ci_i \;\simeq M^\pi/X_i \;,
\]
or comme $X_i \in M$, on voit, à cause de la condition de $\pi$-accessibilité, que $M^\pi/X_i$ est \textit{filtrante.} On va le dégager en

\textbf{Lemme 5.10.} \textit{Soit $M$ une partie $\pi$-accessible de $\M$. Alors pour tout $Y \in M$, $M^\pi/Y$ est une catégorie} filtrante, \textit{et $\Ss_\pi(Y)$ y est cofinale.}

\msp{page 63}

Considérons le foncteur canonique
\[
I \;=\; \Ss_\pi(Y) \xrightarrow{\;\varphi\;} \underline{M}^\pi/Y \;=\; I' \;,
\]
où la catégorie source $I$ est filtrante. Par SGA 4 I 8.1.3, ce foncteur est cofinal si et seulement s'il satisfait deux conditions F\,1, F\,2, et alors $I'$ est également filtrante. La condition F\,1 est la condition habituelle que tout $i' \in \Ob I'$ est “majoré” par un $\varphi(i)$, \textit{i.e.} tout $Z \longrightarrow Y$, avec $Z \in M^\pi$, est “majoré” par un \textit{monomorphisme} $Z_\alpha \hookrightarrow Y$, avec $Z_\alpha \in M^\pi$, ce qui résulte du fait que c'est majoré par $Z_\alpha \hookrightarrow Y$ avec $Z_\alpha \in \Filt^\pi$, \textit{i.e.} $Z_\alpha \in \Ss_\pi(Y)$, et que parmi ces $Z_\alpha$ ceux qui sont dans $M^\pi$ sont cofinaux (vu que $Y \in M$). L'autre condition\footnote[23]{Comme $I$ est filtrante.} est qu'une double flèche $i' \overset{f'}{\underset{g'}{\rightrightarrows}} \varphi(i)$ peut
\msp{page 64}
s'égaliser par une flèche $\varphi(h)$, avec $h : i \longrightarrow j$. Mais en fait, pour un $i' \in \Ob I'$, \textit{i.e.} un $Z \longrightarrow Y$ donné ($Z \in M^\pi$), et un $i \in \Ob I$, \textit{i.e.} un $Z_\alpha \hookrightarrow Y$ donné ($Z_\alpha \in M^\pi$), il existe au plus \textit{une} flèche dans $I' = \underline{M}^\pi/Y$ de $Z \longrightarrow Y$ dans $Z_\alpha \hookrightarrow Y$. Donc la condition F\,2 est trivialement satisfaite (on prend $h = \mathrm{id}_i$).

\textbf{Corollaire 5.11.} \textit{Si $(X_i)_{i \in I}$ est un système inductif filtrant, avec les $X_i$ dans $M^\pi$, alors la catégorie $\Ci = \underline{M}^\pi/\underline{M} \times_{\underline{M}} (I,\varphi)$, construite page} 62, \textit{est filtrante, et le foncteur canonique} (\textit{fibrant}) \textit{$\Ci \longrightarrow I$ est cofinal.}

Que $\Ci$ soit filtrante et $\Ci \longrightarrow I$ cofinal résulte formellement du fait que $I$ l'est, que les fibres de $\Ci$ le sont, et que $\Ci$ est cofibrée sur $I$. (Mérite un lemme, à toutes fins utiles.)

\textbf{N.B.} On peut remplacer $\Ci$ par $\Ci_0$ dont les fibres sont les
\[
\Ss_\pi^{M^\pi}(X_i) \;\subset\; \underline{M}^\pi \;; \qquad \text{cofinale dans } \Ci_i = \underline{M}^\pi/X_i \;,
\]
\msp{page 65}
on trouve une sous-catégorie de $\Ci$ cofinale dans $\Ci$, et qui est cette fois ordonnée tout comme $I$.

\printedpage{38}
% head: A. Grothendieck : Dérivateurs, Ch. XV, Édition des Universités de Montpellier II et Paris VII
% foot: Version 5 novembre 2010
On a
\[
X_i \;=\; \varinjlim_{Z_\alpha \in \Ci_i} Z_\alpha \;,
\]
donc
\[
X \;=\; \varinjlim_{\alpha \in \Ci} Z_\alpha \;,
\]
au total on est arrivé à représenter $X$ comme $\varinjlim$ filtrante d'éléments $Z_\alpha$ qui sont (non seulement dans $M$, mais même) dans $M^\pi$.

Changeons la notation, en prenant
\[
X \;=\; \varinjlim_{i \in I} X_i \;, \qquad I \text{ ensemble ordonné filtrant, } X_i \in M^\pi \;\; \forall\, i \in I \;,
\]
soit $I'$ l'ensemble des parties \textit{filtrantes} $J \subset I$ telles que $\mathrm{card}\, J \leq \pi$, $I'$ ordonné par inclusion. Alors $I'$ est grand devant $\pi$, et on a un système inductif
\[
(X_J)_{J \in I'} \;, \qquad X_J \;=\; \varinjlim_{i \in J} X_i \;.
\]
Par l'hypothèse de stabilité de $M^\pi$ (qui n'a pas encore servi), les $X_J$ sont dans $M^\pi$. On a donc
\[
X \;=\; \varinjlim_{J \in I'} X_J \;, \qquad \text{les } X_J \text{ dans } M^\pi\text{, } I' \text{ grande devant } \pi \;,
\]
donc $X \in M$ puisque $M$ est $\pi$-admissible, \quad \textit{q.e.d.}

\msp{page 66}

\textbf{Définition 5.12.}
\begin{enumerate}[label={}]
\item[1)] Soit $\M$ une catégorie satisfaisant les conditions $(a)$, $(b)$, $(c)$, $(d)$ (p.~44-46). Soit $\pi \geq \pi_0$ un cardinal. Une partie $M$ de $\Ob\M$ est dite \textit{L-$\pi$-accessible,} si elle est $\pi$-accessible et stable par petites $\varinjlim$ filtrantes\footnote[24]{On dit alors que $M$ est \textit{L-stable.}}, ou ce qui revient au même (5.9), si elle est $\pi$-accessible et si $M^\pi = M \cap \Filt^\pi$ est stable par $\varinjlim$ filtrantes de cardinal $\leq \pi$. (\textbf{N.B.} Si les objets de $\Filt^\pi$ sont $\pi$-accessibles, 5.8 nous dit que cela signifie aussi qu'il suffit pour cela qu'il existe une partie $M_0$ de $\Filt^\pi$ telle que $M$ soit la L-enveloppe (enveloppe stable par $\varinjlim$ filtrantes) de $M_0$ -- condition qui est de toutes façons toujours nécessaire, en prenant $M_0 = M^\pi$.)
\item[2)] Soit $M$ une partie de $\Ob\M$. On dit qu'elle est \textit{L-accessible,} si elle satisfait à l'une des quatre conditions équivalentes suivantes :
\begin{enumerate}[label={}]
\item[\textit{a})] Il existe un cardinal $\pi$ tel que $M$ soit L-$\pi$-accessible.
\item[\textit{a}$'$)] $M$ est accessible et L-stable.
\item[\textit{b})] Il existe une petite partie $M_0$ de $\Ob\M$, telle que $M$ soit la L-enveloppe de $M_0$\footnote[25]{Cette condition ne dépend pas de la filtration cardinale choisie de $\M$.}.
\item[\textit{b}$'$)] Il existe une partie $M_0$ de $\M$ contenue dans une partie $\Filt^\pi$ (ou ce qui revient au même, telle que la sous-catégorie pleine $\underline{M}_0$ de $\M$ soit essentiellement petite, \textit{i.e.} équivalente à une petite catégorie, telle que $M$ soit la \textit{L}-enveloppe de $M_0$).
\end{enumerate}
\end{enumerate}


\endgroup
% END UNIT Der-ChXV-05-11-10/Der-ChXV-05-11-10_p036-040.tex
% BEGIN UNIT Der-ChXV-05-11-10/Der-ChXV-05-11-10_p041-045.tex
\begingroup\makeatletter%
\newcommand{\printedpage}[1]{\marginpar{\footnotesize p.~#1}}%
\newcommand{\msp}[1]{\marginpar{\footnotesize [#1]}}%
\newcommand{\Filt}{\mathrm{Filt}}%
\newcommand{\Fl}{\mathrm{Fl}}%
\newcommand{\Ob}{\mathrm{Ob}}%
\newcommand{\card}{\operatorname{card}}%
\newcommand{\uHom}{\underline{\mathrm{Hom}}}%
\newcommand{\uFl}{\underline{\mathrm{Fl}}}%
\newcommand{\into}{\lhook\joinrel\longrightarrow}%
\makeatother\let\printedpage\derpp\let\msp\dermsp%
\let\mspin\DERmspin%


\printedpage{39}
% head: A. Grothendieck : Dérivateurs, Ch. XV, Édition des Universités de Montpellier II et Paris VII
% foot: Version 5 novembre 2010

\msp{page 67}
L'équivalence de $(a)$ et $(a')$, ou de $(b)$ et $(b')$ est plus ou moins tautologique. L'équivalence de $(a)$ et $(b)$ résulte aussitôt du N.B. inclus dans $1^{\circ}$, compte tenu que les cardinaux $\pi$ tels que les objets de $\Filt^{\pi}$ soient $\pi$-accessibles sont cofinaux dans l'ensemble des cardinaux dans $\mathfrak{U}$. (\textit{Cf.} 3.2, page 16).

\textbf{Proposition 5.13.} \textit{Toute réunion finie de parties L-$\pi$-accessibles de} $\Ob\mathcal{M}$ \textit{est L-$\pi$-accessible. L'intersection d'une famille $(M_{\alpha})_{\alpha\in A}$ de parties L-$\pi$-accessibles, avec} $\card A \leq \pi$\textit{, est L-$\pi$-accessible.}

\textsc{Démonstration}. La première assertion provient du fait qu'une réunion \textit{finie} de parties L-stables de $M$ est L-stable, et qu'on a itou pour la notion de $\pi$-accessibilité. Les deux choses se prouvent d'ailleurs pratiquement de la même façon, en utilisant le

\textbf{Lemme 5.14.} \textit{Soit $I$ un ensemble ordonné filtrant, $(I_{\alpha})_{\alpha\in A}$ une famille} finie \textit{non vide de parties de $I$, telle que $\bigcup I_{\alpha} = I$. Alors il existe $\alpha \in A$ tel que $I_{\alpha}$ soit cofinal.}

\msp{page 68}
\textsc{Démonstration} par l'absurde. Si pour aucun $\alpha$, $I_{\alpha}$ n'était cofinal, on en conclurait que pour tout $\alpha$ il existe un $i_{\alpha}$ tel que $I/i_{\alpha} \cap I_{\alpha} = \emptyset$. Soit $i$ majorant les $i_{\alpha}$, on aura donc $I/i \cap I_{\alpha} = \emptyset \ \ \forall\, \alpha \in A$, contrairement à l'hypothèse que $\bigcup I_{\alpha} = I$.

Le cas d'une \textit{intersection} de parties L-$\pi$-accessibles : elle est bien sûr L-stable, donc il reste à prouver qu'elle est $\pi$-accessible, \textit{i.e.} satisfait la condition $\mathrm{b}_{\pi}$ de \textit{loc. cit.}

\msp{page 69}
Soit donc
\[
X \;=\; \varinjlim_{i\in I} X_{i} \; , \qquad\quad X_{i} \in \Filt^{\pi} \; ,
\]
une représentation $\pi$-adaptée de $X$, avec $I$ un ensemble ordonné $\pi$-adapté. Donc on a
\[
\begin{aligned}
X \in M_{\alpha} \quad &\Longleftrightarrow \quad \text{l'ensemble } I_{M_{\alpha}} \text{ des } i \in I \text{ tels que } X_{i} \in M_{\alpha}\\
&\phantom{\Longleftrightarrow} \quad \text{est cofinal dans } I.
\end{aligned}
\tag*{$(*)$}
\]
D'ailleurs, comme $M_{\alpha}^{\pi} = M_{\alpha} \cap \Filt^{\pi}$ est stable par $\varinjlim$ filtrante de cardinal $\leq \pi$, on voit tout de suite que la partie $I_{M_{\alpha}}$ de $I$ est stable par bornes supérieures dans $I$ de parties $J$ de $I_{M_{\alpha}}$ filtrantes et de cardinal $\leq \pi$. Donc $I_{M_{\alpha}}$ cofinal équivaut à $I_{M_{\alpha}}$ $\pi$-\textit{cofinal} (\textit{cf.} 5.8, page 59). Soit $M = \bigcap_{\alpha\in A} M_{\alpha}$, je veux prouver que
\[
X \in M \quad \Longleftrightarrow \quad \text{l'ensemble } I_{M} \ (= \textstyle\bigcap I_{M_{\alpha}}) \text{ est cofinal dans } I.
\tag*{$(**)$}
\]
Or $(*)$ implique que $X \in M$ $(= \bigcap M_{\alpha})$ équivaut à $I_{M_{\alpha}}$ $\pi$-cofinal pour tout $\alpha$. Donc au total on veut prouver
\[
I_{M_{\alpha}} \text{ cofinal dans } I \text{ pour tout } \alpha \quad\ \Longleftrightarrow \quad\ \textstyle\bigcap I_{M_{\alpha}} \text{ est cofinal dans } I.
\tag*{$(***)$}
\]

\msp{page 70}
On a bien sûr $\Longleftarrow$, donc il reste à prouver $\Longrightarrow$. Si l'ensemble d'indices $A$ est fini, cela résulte du lemme 5.8, compte tenu que la cofinalité des $I_{M_{\alpha}}$ implique leur $\pi$-cofinalité. On doit donc simplement améliorer 5.8, par le

\printedpage{40}
% head: A. Grothendieck : Dérivateurs, Ch. XV, Édition des Universités de Montpellier II et Paris VII
% foot: Version 5 novembre 2010

\textbf{Lemme 5.15.} \textit{Soient $\pi$ un cardinal, $I$ un ensemble ordonné filtrant $\pi$-adapté, et soit $(I_{\alpha})_{\alpha\in A}$ une famille de parties $I_{\alpha}$ $\pi$-cofinales} (cf. 5.8), \textit{telle que} $\card A \leq \pi$. \textit{Alors $\bigcap I_{\alpha}$ est $\pi$-cofinale.}

Il suffit de prouver qu'elle est cofinale, car la propriété de stabilité par bornes supérieures de parties filtrantes de cardinal $\leq \pi$ est évidente. Quitte à remplacer $I$ par $I/i_{0}$, et les $I_{\alpha}$ par leurs traces sur $I/i_{0}$, la question est simplement de prouver que $\bigcap_{\alpha} I_{\alpha} \neq \emptyset$. On va construire un élément $s$ de $\bigcap I_{\alpha}$ par induction transfinie. Prenons un bon ordre sur $A$, et prouvons par récurrence transfinie
\msp{page 71}
que pour tout $\alpha \in A$, l'ensemble
\[
J_{\alpha} \;=\; \bigcap_{\beta<\alpha} I_{\beta}
\]
est cofinal. Si on sait que $J_{\alpha}$ est cofinal, il en est de même de $J_{\alpha+1} = J_{\alpha} \cap I_{\alpha}$, par 5.8. Donc on est ramené à prouver que $J_{\alpha}$ est cofinal si $\alpha$ est un ordinal limite, donc à prouver le lemme 5.15 pour une suite transfinie \textit{décroissante} de parties $\pi$-cofinales $J_{\beta}$ $(\beta < \alpha)$ quand $\card \alpha \leq \pi$, telle que $J_{\beta} = \bigcap_{\beta'<\beta} J_{\beta'}$ quand $\beta$ est limite. À prouver que l'intersection est $\pi$-cofinale -- et on est réduit à prouver qu'elle est non vide. On va construire par récurrence transfinie une suite \textit{croissante} d'éléments
\[
i_{\beta} \;\in\; J_{\beta} \qquad\quad (\, \beta \;\leq\; \alpha \,) \; .
\]
Si $i_{\beta}$ est construit, comme $J_{\beta+1}$ est cofinal, il existe $i_{\beta+1} \geq i_{\beta}$ dans $J_{\beta+1}$, O.K. Si $\beta$ est un ordinal limite, on prend
\[
i_{\beta} \;=\; \sup_{\beta'<\beta} i_{\beta'} \; ,
\]
$i_{\beta}$ est dans chaque $I_{\beta_{0}}$ pour $\beta_{0} < \beta$,
\msp{page 72}
puisque les $i_{\beta'}$ pour $\beta_{0} \leq \beta' < \beta$ sont tous dans $J_{\beta_{0}} \supset J_{\beta'}$, donc aussi leur limite, donc
\[
i_{\beta} \;\in\; \bigcap_{\beta_{0}<\beta} J_{\beta_{0}} \;=\; J_{\beta} \; .
\]
Donc la récurrence transfinie marche, et nous donne donc un élément $s_{\alpha}$ de $J_{\alpha}$\,!

Je vais prouver aussi en forme l'assertion pour $\bigcup_{\alpha\in A} M_{\alpha}$, quand $A$ est fini. Posant cette fois $M = \bigcup_{\alpha} M_{\alpha}$, il faut prouver, moyennant les relations $(*)$ (p. 68), qu'on a
\[
\begin{aligned}
&X \in M &&\Longleftrightarrow \quad &&\text{l'ensemble } \ I_{M} \ \ (= \ \textstyle\bigcup_{\alpha} I_{M_{\alpha}})\\
&(\textit{i.e.} \ \exists \ \alpha \text{ avec } X \in M_{\alpha}) && &&\text{est cofinal dans } I \ ,
\end{aligned}
\tag*{$(**)$}
\]
et moyennant $(*)$ le premier terme de cette équivalence à établir signifie qu'il existe $\alpha$ avec $I_{M_{\alpha}}$ cofinal dans $I$. Bien sûr, si un des $I_{M_{\alpha}}$ est cofinal dans $I$, a fortiori leur réunion. Mais inversement, si la réunion est cofinale, elle est filtrante et réunion des parties (filtrantes) $I_{M_{\alpha}}$, dont l'une doit donc être cofinale dans elle par le lemme 5.14, donc aussi dans $I$, \textit{q.e.d.}

\msp{page 73}
\textbf{Corollaire 5.16} (de la proposition 5.13). \textit{Toute réunion finie de parties L-accessibles de $\mathcal{M}$ est L-accessible. Toute intersection $\bigcap M_{\alpha}$ d'une famille petite de parties L-accessibles de $\mathcal{M}$, est L-accessible.}

\printedpage{41}
% head: A. Grothendieck : Dérivateurs, Ch. XV, Édition des Universités de Montpellier II et Paris VII
% foot: Version 5 novembre 2010

\section*{6\quad Essai avorté pour un critère d'engendrement des parties L-accessibles de $\Fl(\mathcal{M})$.}

Soit $\mathcal{M}$ une catégorie munie d'une filtration cardinale $(\Filt^{\pi})_{\pi\geq\pi_{0}}$, et soit $S$ un objet de Cat, considérons
\[
\mathcal{M}(S) \;=\; \uHom(S,\mathcal{M}) \; ,
\]
et munissons $\mathcal{M}(S)$ de la filtration par les
\[
\mathcal{M}^{\pi}(S) \;=\; \uHom(S,\mathcal{M}^{\pi} = \underline{\Filt}^{\pi}(\mathcal{M})) \; .
\]
Je dis que c'est là une filtration cardinale de $\mathcal{M}(S)$. Il faut vérifier les conditions $(a)$, $(b)$, $(c)$ de SGA 4 I 9.12. Les conditions $(a)$ ($\Filt^{\pi}$ essentiellement petite) et $(b)$ (condition $\mathrm{L}_{\pi_{0}}$, et stabilité de $\Filt^{\pi}$ par $\varinjlim\limits_{I}$ , $I$ filtrante grande devant $\pi_{0}$ et de cardinal $\leq \pi$) passent aussitôt de $\mathcal{M}$ à $\mathcal{M}(S)$. La condition $(c)$ est plus délicate (représentation
\msp{page 74}
d'un objet comme $\varinjlim$ filtrante grande devant $\pi$ d'objets de $\Filt^{\pi}$). Je ne vais pas refaire le travail, fait dans un cas beaucoup plus général et plus délicat dans SGA 4 I §9, savoir dans le corollaire 9.24 du théorème 9.22 (p. 177). La seule restriction, c'est qu'on y suppose les objets de $\mathcal{M}$ accessibles, et désignant par $\pi_{1}$ le plus petit cardinal tel que $\pi_{1} \geq \pi_{0}$ et que les objets de $\Filt^{\pi_{0}}$ soient $\pi_{1}$-accessibles, on se borne pour la filtration cardinale de $\uHom(S,\mathcal{M})$ aux cardinaux \fbox{$\pi \geq c = 2^{c_{0}}$,} où
\[
c_{0} \;=\; \sup(\pi_{1}, \card \Fl(S)) \; .
\]
C'est une restriction raisonnable et anodine.

Pour l'instant, je suis intéressé surtout au cas où $S = \boldsymbol{\Delta}_{1}$, donc à
\[
\uHom(\boldsymbol{\Delta}_{1},\mathcal{M}) \;=\; \uFl(\mathcal{M}) \; ,
\]
auquel cas on a $c_{0} = \pi_{1}$, donc il suffit de se borner aux cardinaux $\pi \geq 2^{\pi_{1}}$ pour avoir une bonne filtration cardinale sur $\uFl(\mathcal{M})$.

C'est donc l'ensemble des deux conditions $(b)$, $(c)$ sur $\mathcal{M}$ (p. 44) qui est stable par passage aux $\uHom(S,\mathcal{M}) = \mathcal{M}(S)$. Il en est de même de la condition $(a)$ (stabilité par petites limites inductives), et aussi de la condition\msp{page 75}
$(d)$ (p. 45). De même la condition d'exactitude $(e)$ (p. 46) sur les $\varinjlim$ filtrantes (elles doivent transformer les monomorphismes en monomorphismes).

Pour être tranquille, je vais faire sur $\mathcal{M}$ les hypothèses $(a)$ à $(e)$ inclus (bien que $(e)$ me semble un peu artificielle, on devrait pouvoir l'éliminer, dans un second souffle\,!). Ces conditions ont donc la propriété de stabilité souhaitable par passage aux $\mathcal{M}(S)$, et en particulier par passage à $\Fl(\mathcal{M})$. Quitte à augmenter $\pi_{0}$, on peut donc supposer que les $\Fl^{\pi}(\mathcal{M}) = \mathrm{Hom}(\boldsymbol{\Delta}_{1},\mathcal{M}^{\pi})$ $(\pi \geq \pi_{0})$ forment aussi une filtration cardinale de $\uFl(\mathcal{M})$.

On dispose donc des notions de parties $\pi$-accessibles, accessibles, L-$\pi$-accessibles, L-accessibles de $\Fl(\mathcal{M}) \ = \ \Ob\uFl(\mathcal{M})$. Ces notions sont stables par réunions et intersections finies, et les notions absolues (sans $\pi$) stables par intersections ``quelconques'' (\textit{i.e.} pour des \textit{petites} familles de parties). Ce sera là une chose particulièrement\printedpage{42}%
% head: A. Grothendieck : Dérivateurs, Ch. XV, Édition des Universités de Montpellier II et Paris VII
% foot: Version 5 novembre 2010
\msp{page 76}
cruciale, le cas d'une intersection $C \cap W$\,!

Le résultat principal, pour lequel j'ai fait les gammes du paragraphe précédent, est le suivant :

\setcounter{footnote}{25}
\textbf{Théorème 6.1}\footnote{Démonstration canulée, énoncé peut-être \textit{faux}\,!!}. \textit{Soit $\mathcal{M}$ une catégorie satisfaisant les conditions dites :}
\begin{enumerate}[label={}]
\item[a)] \textit{Stabilité par petites} $\varinjlim$ \textit{(pas nécessairement filtrantes).}
\item[b)] $\mathcal{M}$ \textit{munie d'une filtration cardinale} $(\Filt^{\pi})_{\pi\geq\pi_{0}}$.
\item[c)] \textit{Les objets de $\mathcal{M}$ sont accessibles.}
\item[d)] \textit{Les} $\Filt^{\pi}$ \textit{sont stables par} $\varinjlim$ \textit{indexées par des catégories $I$ telles que} $\card \Fl(I) \leq \pi$\textit{, et si $Z \in \Filt^{\pi}$, tout quotient $Z'$ de $Z$ (tel que $Z \longrightarrow Z'$ ne se factorise pas par un sous-objet $Z'_{0}$ de $Z'$, différent de $Z'$) est dans} $\Filt^{\pi}$.
\item[e)] \textit{Les limites inductives filtrantes dans $\mathcal{M}$ transforment les monomorphismes en monomorphismes.}
\end{enumerate}
\textit{Soit $C$ une partie L-accessible de} $\Fl(\mathcal{M}) = \Ob\uFl(\mathcal{M})$, i.e. (5.12) \textit{une partie stable par} $\varinjlim$ \textit{filtrantes, et accessible. Supposons de plus $C$ formée de monomorphismes}\footnote{Finalement, on se débarrasse de cette hypothèse.}\textit{, et que $C$ contient les} isomorphismes. \textit{Alors il existe une partie petite $C_{0} \subset C$, telle que $C \subset \widetilde{C_{0}}$. Plus précisément, si $\pi$ est un cardinal tel que $C$ soit $\pi$-accessible, et si on pose $C^{\pi} = C \cap \Fl^{\pi}(\mathcal{M})$, alors on a}
\msp{page 77}
\[
C \;\subset\; \widetilde{C^{\pi}} \;=\; \widetilde{C_{0}} \; ,
\tag*{(6.1.1)}
\]
\textit{où $C_{0}$ est n'importe quelle petite partie de $C^{\pi}$ telle que tout objet de $C^{\pi}$ soit isomorphe à un objet de $C_{0}$.}

\textbf{Corollaire 6.2.} \textit{On a $C_{*} = C_{0\,*}$, $\widetilde{C} = \widetilde{C_{0}}$.}

\textbf{Corollaire 6.3.} \textit{Conditions équivalentes sur la partie L-accessible $C$ de} $\Fl(\mathcal{M})_{\mathrm{mono}}$\textit{, contenant les isomorphismes :}
\begin{enumerate}[label={}]
\item[a)] \textit{$C = \widetilde{C}$.}
\item[b)] \textit{$C$ est stable par cochangement de base et par composition.}
\item[b$'$)] \textit{$C$ est stable par cochangement de base, et par sommes amalgamées simples,} i.e. \textit{si $A \longrightarrow B_{1}$, $A \longrightarrow B_{2}$ sont dans $C$, alors $A \longrightarrow B_{1} \amalg_{A} B_{2}$ aussi.}
\item[c)] \textit{$C^{\pi}$ est stable dans} $\Fl^{\pi}(\mathcal{M}) = \Fl(\mathcal{M}^{\pi})$ \textit{par cochangement de base et par composition.}
\item[c$'$)] \textit{$C^{\pi}$ est stable dans} $\Fl(\mathcal{M}^{\pi})$ \textit{par cochangement de base et par sommes amalgamées} (cf. b$'$).
\end{enumerate}
\textit{Quand ces conditions sont satisfaites, alors $(C,C_{*})$ est un couple de Quillen, justiciable du théorème} 1 \textit{du paragraphe} 1 (p. 10).

La dernière assertion résulte du théorème 6.1 et du corollaire au théorème 1 du paragraphe 1. On a
\printedpage{43}%
% head: A. Grothendieck : Dérivateurs, Ch. XV, Édition des Universités de Montpellier II et Paris VII
% foot: Version 5 novembre 2010
\msp{page 78}
d'autre part tautologiquement
\[
\begin{tikzcd}
(a) \arrow[r, Rightarrow] & (b) \arrow[r, Rightarrow] \arrow[d, Rightarrow] & (b') \arrow[d, Rightarrow] \\
 & (c) \arrow[r, Rightarrow] & (c') \ ,
\end{tikzcd}
\]
et il reste à prouver $(c') \Longrightarrow (a)$, ce qu'on fera plus bas.

Prouvons donc 6.1. Dans (6.1.1) l'égalité $\widetilde{C_{0}} = \widetilde{C^{\pi}}$ est tautologique. Donc tout revient à prouver
\[
C \;\subset\; \widetilde{C^{\pi}} \; .
\]
Une difficulté ici, c'est qu'on ne sait \textit{pas} que $\widetilde{C^{\pi}}$ soit stable par $\varinjlim$ filtrantes -- on le sait seulement pour $C$. Soit donc
\[
A \;\into\; B
\]
une flèche dans $C$, prouvons qu'elle est dans $\widetilde{C^{\pi}}$. Pour ceci, on va construire, par induction transfinie, une suite transfinie \textit{strictement} croissante de sous-objets $A_{\alpha}$ de $B$ contenant $A = A_{0}$,
\[
\underbrace{A}_{=\,A_{0}} \;\; \overset{i_{\alpha}}{\into} \;\; A_{\alpha} \;\; \overset{i_{\beta,\alpha}}{\into} \;\; A_{\beta} \xrightarrow{\ j_{\beta}\ } B \; ,
\]
telle que l'on ait
\[
i_{\alpha} \;\in\; \widetilde{C^{\pi}} \; , \qquad\quad \text{plus généralement les } i_{\beta,\alpha} \;\in\; \widetilde{C^{\pi}} \; , \qquad\quad j_{\alpha} \;\in\; C \; ,
\]
et que l'on ait
\[
A_{\alpha} \;=\; \varinjlim_{\beta<\alpha} A_{\beta} \qquad\quad \text{si } \alpha \text{ ordinal limite} \; .
\]

\msp{page 79}
Cela précise déjà la définition récurrente pour $\alpha = \alpha_{0}$ ordinal limite. Notons que le $A_{\alpha_{0}}$ construit en termes des $A_{\beta}$ $(\beta < \alpha_{0})$ satisfait bien aux conditions que
\begin{enumerate}[label={}]
\item[\textit{a})] les $i_{\alpha_{0},\beta} : A_{\beta} \longrightarrow A_{\alpha_{0}}$ sont dans $\widetilde{C^{\pi}}$, car $\widetilde{C^{\pi}}$ est stable par ``limites inductives ordinales strictes'' (\textit{cf.} 1.19, p. 8),
\item[\textit{b})] $j_{\alpha_{0}} \, : \, A_{\alpha_{0}} \longrightarrow B$ est dans $C$, car c'est la limite inductive filtrante du système des $A_{\beta} \longrightarrow B$ $(\beta < \alpha_{0})$ qui sont dans $C$, et $C$ est stable par ces limites.
\end{enumerate}
Il reste donc à préciser comment on passe de $A_{\alpha}$ à $A_{\alpha+1}$. Cela revient à prouver le

\textbf{Lemme 6.4.} \textit{Soit $f : A \into B$ dans $C$, telle que $f$ ne soit pas un isomorphisme. Alors il existe un sous-objet $A_{1}$ de $B$,} distinct de $A$, \textit{tel que $A \longrightarrow A_{1}$ soit dans $\widetilde{C^{\pi}}$, et $A_{1} \longrightarrow B$ dans $C$.}

Si ce lemme est admis, cela signifie que la construction transfinie se continue aussi longtemps que l'on trouve des $A_{\alpha}$ tels que $A_{\alpha} \neq B$, \textit{i.e.} $A_{\alpha} \into B$ ne soit pas un isomorphisme. Mais comme les $A_{\alpha}$ sont tous différents, si $\pi$ est l'ensemble des sous-objets de $B$ (on sait que c'est un


\endgroup
% END UNIT Der-ChXV-05-11-10/Der-ChXV-05-11-10_p041-045.tex
% BEGIN UNIT Der-ChXV-05-11-10/Der-ChXV-05-11-10_p046-050.tex
\begingroup\makeatletter%
\newcommand{\printedpage}[1]{\marginpar{\footnotesize p.~#1}}%
\newcommand{\msp}[1]{\marginpar{\footnotesize [#1]}}%
\newcommand{\M}{\mathcal{M}}%
\newcommand{\cercle}[1]{\tikz[baseline=(X.base)]\node[draw,circle,inner sep=0.6pt](X){#1};}%
\newcommand{\Filt}{\mathrm{Filt}}%
\newcommand{\Mon}{\mathrm{Mon}}%
\newcommand{\Ss}{\mathrm{Ss}}%
\newcommand{\Fl}{\mathrm{Fl}}%
\newcommand{\Ob}{\mathrm{Ob}}%
\newcommand{\Cof}{\mathrm{Cof}}%
\makeatother\let\printedpage\derpp\let\msp\dermsp%
\let\mspin\DERmspin%


\printedpage{44}
% head: A. Grothendieck : Dérivateurs, Ch. XV, Édition des Universités de Montpellier II et Paris VII

\msp{page 80}
ensemble petit, à cause des hypothèses faites sur $\M$), les ordinaux $\alpha$ pour lesquels $A_\alpha$ existe satisfont $\mathrm{card}\,A_\alpha \leq \pi$. Donc il doit exister un $\alpha$ tel que $A_\alpha = B$. Mais alors $i_\alpha : A \longrightarrow A_\alpha$ est égal à $f$, donc $f \in \widetilde{C^{\pi}}$, \quad \textit{q.e.d.}

Il reste donc à prouver 6.4. Pour ceci, soit $I$ l'ensemble des sous-objets de $B$ qui sont dans $\Filt^{\pi}$, ordonné par inclusion. (Si $B$ est dans $\Filt^{\pi}$, $A$ aussi, donc $f$ est déjà dans $C^{\pi} \subset \widetilde{C^{\pi}}$, et il n'y a rien à prouver.) C'est un ensemble filtrant grand devant $\pi$ (et même $\pi$-\textit{adapté}), dont la $\varinjlim$ est $B$ (5.3, p. 46). Considérons les
\[
A_i \;=\; B_i \cap A \;,
\]
ils sont dans $\Filt^{\pi}$, comme sous-objets d'objets de $\Filt^{\pi}$ (condition $(d)$ sur la filtration cardinale). Ainsi, la flèche $f$ apparaît comme limite inductive filtrante des flèches
\[
f_i \;:\; A_i \longrightarrow B_i \;.
\]
Il est trivial que le système des $A_i$ est cofinal dans $\Ss_{\pi}(A)$, donc $A = \varinjlim A_i$, d'où $f = \varinjlim f_i$.

\msp{page 81}
D'autre part, je dis que le système inductif des $f_i$ dans $\underline{\Fl}(\M)$ est \textit{bien adapté à $\pi$.} On sait déjà que $I$ est bien adapté à $\pi$, le sup dans $I$ d'une partie filtrante $J$ de cardinal $\leq \pi$ étant le sup dans l'ensemble ordonné $\Ss(B)$ (lequel sup se trouve dans $\Ss_{\pi}(B)$). Il est essentiel, pour notre argument, de savoir que si $i_J = \sup J$, alors non seulement $B_{i_J} = \varinjlim\limits_{i\in J} B_i$, mais \marginpar{$\|$}% marginal double vertical bar printed in the left margin at this line
aussi $A_{i_J} = \varinjlim\limits_{i\in J} A_i$, de sorte que le système inductif des $f_i$ est $\pi$-adapté dans $\underline{\Fl}(\M)$. Il semble qu'il faille donc une condition d'exactitude sur $\M$ voisine de $(e)$, savoir :

\begin{itemize}
\item[\cercle{e$'$}] Les foncteurs $\varinjlim$ filtrantes commutent au changement de base.
\end{itemize}

(Il le suffit ici dans le cas très particulier d'un changement de base par un monomorphisme $A \longrightarrow B$, et d'une limite inductive filtrante de sous-objets de $B$ pour des morphismes d'inclusion comme transitions.) Notons que \cercle{e}, \cercle{e$'$} se trouvent généralement comme conséquence de la condition % [sic] Il le suffit -> Il suffit, ungrammatical as printed

\begin{itemize}
\item[(e$_0$)] Les $\varinjlim$ filtrantes dans $I$ sont exactes à gauche.
\end{itemize}

\msp{page 82}
Ceci posé, on voit donc, $C$ étant $\pi$-accessible et $I$ grand devant $\pi$, que l'ensemble $I_C$ des $i \in I$ tels que $f_i \in C$ est \textit{cofinal} dans $I$. Mais comme $f_i \in \Filt^{\pi}$, $f_i \in C$ équivaut à $f_i \in C^{\pi}$. Comme $A \longrightarrow B$ n'est pas un isomorphisme, il existe un $i_0$ tel que $B_{i_0} \not\subset A$, puis un $i \geq i_0$ tel que $f_i \in C^{\pi}$, et aussi $B_i \not\subset A$ (puisque $B_i \supset B_{i_0}$). Ainsi on a prouvé ceci :

\begin{itemize}
\item[($*$)] Sous les hypothèses de (6.4) et moyennant l'hypothèse \cercle{e$'$} sur $\M$, il existe un sous-objet $B'$ de $B$ satisfaisant les conditions suivantes.
\begin{itemize}
\item[\textit{a})] $B' \in \Filt^{\pi}$.
\item[\textit{b})] $B' \not\subset A$.
\item[\textit{c})] Posant $A' = B' \cap A$, l'inclusion
\[
i' \;:\; A' \;\longrightarrow\; B'
\]
est dans $C^{\pi}$.
\end{itemize}

\printedpage{45}
% head: A. Grothendieck : Dérivateurs, Ch. XV, Édition des Universités de Montpellier II et Paris VII

De plus, l'ensemble des $J \subset \Ss_{\pi}(B)$ formé des $B'$ en question est cofinal dans $\Ss_{\pi}(B)$, et même $\pi$-cofinal.
\end{itemize}

Il faut prouver encore ce dernier point -- que $J$ (dont on sait qu'il est cofinal) est même $\pi$-cofinal, \textit{i.e.} stable par sup (dans $\Ss_{\pi}(B)$) des parties filtrantes $J' \subset J$ telles que $\mathrm{card}\, J' \leq \pi$. Mais cela résulte aussitôt du fait que $C$ étant stable par

\msp{page 83}
limites inductives filtrantes, $C^{\pi} = C \cap \Filt^{\pi}$ est stable par $\varinjlim$ filtrantes de cardinal $\leq \pi$, en vertu de la condition $(d)$ sur $\Filt^{\pi}$. Considérons maintenant le carré cocartésien
\[
\begin{tikzcd}
A' \arrow[r, hook, "i'"] \arrow[d, hook] & B' \arrow[d] & & \\
A \arrow[r, "i"] & C' = A \amalg_{A'} B' \arrow[r, "j"] & B \;.
\end{tikzcd}
\]
Comme $i' \in C^{\pi}$, on a $i \in \widetilde{C^{\pi}}$ (car $\widetilde{C^{\pi}}$ stable par cochangement de base). D'autre part, j'ai besoin du fait que $j$ soit dans $C$, et je ne sais pas même qu'il est un monomorphisme ! Aussi il faut repasser aux notations indicielles
\[
\begin{tikzcd}
A_i \arrow[r, hook, "\alpha_i"] \arrow[d, hook] & B_i \arrow[d] & \\
A \arrow[r, hook, "\beta_i"] & C_i \arrow[r, "\gamma_i"] & B \;.
\end{tikzcd}
\]
Comme la formation des sommes amalgamées commute aux limites inductives filtrantes, on voit que $\varinjlim C_i \simeq A \amalg_{\varinjlim A_i} \varinjlim B_i = A \amalg_A B = B$, donc la limite inductive des $\gamma_i$ est un isomorphisme (essentiellement $\mathrm{id}_B$). D'autre part, le même argument montre que le système inductif des $C_i$ $(i \in J)$ dans $\M$ est $\pi$-adapté, donc aussi celui des $\gamma_i$\footnote[28]{\textit{Sauf} qu'on ne sait pas que les $C_i$ soient dans $\Filt^{\pi}$ -- en fait, ils ne seront pas en général, puisqu'on ne suppose pas que $A, B \in \Filt^{\pi}$.}
\msp{page 84}
dans $\underline{\Fl}(\M)$. Comme $\varinjlim \gamma_i$ est un isomorphisme, il est dans $C$. Comme $C$ est $\pi$-accessible, il s'ensuit que l'ensemble $J'$ formé des $i \in J$ tels que $\gamma_i \in C$ est lui aussi cofinal\footnote[29]{\textit{Canulé, cf.} annotation précédente.}. Donc on trouve $i$ avec $\alpha_i \in C^{\pi}$, $\gamma_i \in C$, ce qui établit le lemme 6.4, et achève la démonstration du théorème 6.1.

\section*{7\quad Construction de triples de Quillen clos avec $\boldsymbol{C} = \mathbf{Mon} \subset \mathbf{Cof}_{\boldsymbol{W}}$}

Finalement, tous les beaux développements des paragraphes 5, 6 ne me donnent pas le “sweeping result” souhaité. Peut-être est-il faux -- peut-être la propriété d'engendrement
\printedpage{46}
% head: A. Grothendieck : Dérivateurs, Ch. XV, Édition des Universités de Montpellier II et Paris VII
“$^{\sim}$” par une petite partie (pour les ensembles accessibles, ou L-accessibles) est-elle beaucoup plus rarement satisfaite, que je ne l'aurais imaginé. Pour bien faire, il faudrait que je traite jusqu'au % resolved: the quoted raised wave denotes Q-closure; see corrections.csv, leaf 13.
\msp{page 85}
bout quelques exemples, p. ex. celui de $\Cof_W$, et de $W^{\mathrm{univ}}$ dans Cat, pour un localiseur fondamental plus ou moins général, et surtout dans le cas $W = \mathbf{W}_{\infty}$. Sûrement $\Cof_W$ et $W^{\mathrm{univ}}$ sont accessibles (il faudrait que je le vérifie pourtant), et si le “théorème” hypothétique du paragraphe 6 est faux pour eux, du moins en aurai-je le cœur net. Mais il reste pourtant la question : si $C = \widetilde{C}$ est stable par $\varinjlim$ filtrantes, et de la forme $\widetilde{C_0}$, $C_0$ petit, et si $W$ est L-accessible également, alors $C \cap W$ est-il aussi de la forme $\widetilde{TC_0}$, avec $TC_0$ petit ? Je viens encore d'essayer de le prouver, en supposant de plus $C \subset$ mono, et ne vois aucun moyen d'y parvenir. Il faudrait donc un contre-exemple pour dissiper ces perplexités ! Et le mieux serait dans Cat, tant qu'à faire.

Il surnage pourtant un résultat tangible positif, que je vais maintenant essayer de bien cerner et prouver :

\msp{page 86}
\textbf{Théorème 7.1.} \textit{Soit $\M$ une catégorie stable par petites $\varinjlim$ filtrantes, $M_0 \subset \Ob\,\M$ un sous-ensemble de} $\Ob\,\M$ \textit{générateur par monomorphismes stricts, plus précisément on suppose ceci :}

\begin{quote}
\itshape
Pour tout monomorphisme $i : A \hookrightarrow B$ dans $\M$ qui n'est pas un isomorphisme, il existe un monomorphisme $B' \overset{j}{\hookrightarrow} B$ ayant les propriétés suivantes :
\begin{itemize}
\item[\textup{a)}] $B' \in M_0$.
\item[\textup{b)}] $j$ ne se factorise pas par $i$.
\item[\textup{c)}] $A' = A \times_B B'$ existe dans $\M$, ainsi que $C = A \amalg_{A'} B'$, $A'$ est dans $M_0$, et dans le diagramme canonique
\[
\begin{tikzcd}
A' \arrow[r, hook, "i'"] \arrow[d, hook] & B' \arrow[d, hook] \arrow[dr, hook, dotted, "j"] & \\
A \arrow[r, hook, "\alpha"] \arrow[rr, hook, dotted, bend right=20, "i"'] & C \arrow[r, "\beta"] & B
\end{tikzcd}
\tag*{(7.1.1)}
\]
la flèche $\beta$ est un monomorphisme.
\end{itemize}
\end{quote}

(\textbf{N.B.} Ces conditions sont satisfaites si on suppose, en plus de la condition que $M_0$ soit générateur pour les monomorphismes, que $M_0$ contient, avec tout objet $B'$, et pour tout objet $A'$ qui est soit un quotient, soit un sous-objet de $B'$, un objet $A'_0$ isomorphe à $A'$ ; et si de plus $\M$ satisfait aux conditions d'exactitude suivantes :
\begin{itemize}
\item[1$^{\circ}$)] Toute flèche $f$ de $\M$ se factorise en $ip$, avec $i$ monomorphisme, $p$ épimorphisme, et :
\item[2$^{\circ}$)] Si $A_1$, $A_2$ sont deux sous-objets d'un objet $B$ de $\M$, alors l'intersection $A_0 = A_1 \times_B A_2$ existe, ainsi que la
\printedpage{47}
% head: A. Grothendieck : Dérivateurs, Ch. XV, Édition des Universités de Montpellier II et Paris VII
\msp{page 87}
somme amalgamée $C = A_1 \amalg_{A_0} A_2$, et le morphisme canonique $C \longrightarrow B$ est un monomorphisme.)
\end{itemize}

\textit{Supposons de plus que $\M$ satisfait la condition} \cercle{e}, p. 76, i.e. \textit{que toute $\varinjlim$ filtrante de monomorphismes est un monomorphisme.}

1$^{\circ}$) \textit{Soit} $\Mon \subset \Fl(\M)$ \textit{l'ensemble des monomorphismes, et} $\Mon_0$ \textit{le sous-ensemble de} $\Mon$ \textit{formé des $i \in$} $\Mon$ \textit{tels que source et but appartiennent à $M_0$. On a alors}
\[
\Mon \;\subset\; \widetilde{\Mon_0} \quad , \qquad \text{\textit{d'où}} \quad
\left\{
\begin{aligned}
\widetilde{\Mon} \quad &= \quad \widetilde{\Mon_0} \\
(\Mon)_{*} \quad &= \quad (\Mon_0)_{*} \;.\text{\footnotemark[30]}
\end{aligned}
\right.
\tag*{(7.1.2)}
\]
\footnotetext[30]{\textbf{N.B.} On a là un critère d'injectivité d'objets dans les catégories $\M/Y$ -- il suffit de tester la condition d'injectivité par des monomorphismes $i$ qui sont dans $\Mon_0$.}

2$^{\circ}$) \textit{Si $M_0$ est} petit \textit{et si les objets de $M_0$ sont accessibles} (p. ex. $\M$ satisfait \cercle{c}, p. 76), \textit{et si} $\Mon$ \textit{est stable par cochangement de base, alors il existe une factorisation fonctorielle de toute flèche $f$ de $\M$ en}
\[
f \;=\; p\, i \;,
\]
\textit{avec $i \in$} $\Mon$, \textit{$p \in$} $(\Mon)_{*}$. (\textbf{N.B.} Une flèche $p : X \longrightarrow Y$ de $\M$ est dans $(\Mon)_{*}$ si et seulement si elle fait de $X$ un objet \textit{injectif} de $\M/Y$ -- donc l'assertion de factorisation assure que dans les catégories induites $\M/Y$ “il y a assez d'injectifs”.) \textit{De plus, on a}
\[
\Mon \;=\; \widetilde{\Mon_0} \quad .
\tag*{(7.1.3)}
\]

3$^{\circ}$) \textit{Soit $W$ une partie de} $\Fl(\M)$ \textit{qui est} L-accessible (i.e. \textit{accessible, et stable par limites inductives filtrantes}). (\textit{On suppose à présent que $\M$ satisfait les conditions} (a), (b), (c), (d), (e) \textit{de la page 76, d'où une bonne théorie des parties L-accessibles de} $\Ob\,(\M)$ \textit{ou de} $\Fl(\M)$.) \textit{On suppose, plus précisément, que $W$ satisfait les conditions :}

\msp{page 88}
\[
\left\{
\begin{array}{@{\quad}l@{\quad}p{0.6\textwidth}}
\text{a)} & \textit{$W$ est L-accessible} (\textit{i.e.} stable par $\varinjlim$ filtrantes, et accessible). \\[1.5ex]
\text{b)} & \textit{Soient $\alpha$, $\beta$ composables, avec $\beta\alpha$ et $\alpha$ dans $W$, alors $\beta \in W$.} (\textbf{N.B.} Cette condition sera satisfaite si $W$ est un localiseur modérément saturé, et c'est à cette situation que nous allons appliquer notre énoncé.) \\[1.5ex]
\text{c)} & $\Mon \cap W$ \textit{est stable par cochangement de base.}
\end{array}
\right.
\tag*{(7.1.4)}
\]
\textit{et que $\M$ satisfait les conditions supplémentaires}
\[
\left\{
\begin{array}{@{\quad}l@{\quad}p{0.6\textwidth}}
\text{e$'$)} & \textit{Les $\varinjlim$ filtrantes commutent au changement de base.} \\[1.5ex]
\text{f)} & \textit{La condition d'exactitude} (7.2.1) \textit{ci-dessous.}
\end{array}
\right.
\tag*{(7.1.5)}
\]
\textit{Soit $\pi$ un cardinal tel que $W$ soit $\pi$-accessible, et $M_0 \subset$} $\Filt^{\pi}$ (\textit{donc $M_0$ essentiellement petite}). \textit{Alors on a}
\[
W \cap \Mon \;\subset\; \widetilde{(W \cap \Mon)^{\pi}}
\tag*{(7.1.6)}
\]
\printedpage{48}
% head: A. Grothendieck : Dérivateurs, Ch. XV, Édition des Universités de Montpellier II et Paris VII
(\textit{où pour toute partie $C$ de} $\Fl(\M)$, \textit{on pose} $C^{\pi} = C \cap \Fl^{\pi}(\M)$), \textit{donc}
\[
\widetilde{W \cap \Mon} \;=\; \widetilde{(W \cap \Mon)^{\pi}} \;,
\tag*{(7.1.7)}
\]
\textit{donc $\widetilde{W \cap \Mon}$ est Q-engendré par une} petite \textit{partie de} $W \cap \Mon$.

\textbf{Corollaire 7.2.} \textit{Soit $\M$ satisfaisant} (a), (b), (c), (d), (e) \textit{et aux conditions supplémentaires :}
\[
\left\{
\begin{array}{@{\quad}l@{\quad}p{0.6\textwidth}}
\text{e$'$)} & \textit{Les $\varinjlim$ filtrantes commutent au changement de base.} \\[1.5ex]
\text{f)} & \textit{Si $A_1$, $A_2$ sont deux sous-objets d'un objet $B$ de $\M$, et $A_0 = A_1 \cap A_2$, alors $A_1 \amalg_{A_0} A_2 \hookrightarrow B$ est un monomorphisme}\,\footnotemark[31].
\end{array}
\right.
\tag*{(7.2.1)}
\]
\footnotetext[31]{Une autre façon de le dire : si $C = \sup_{\Ss(B)}(A_1, A_2)$, le carré cartésien
\[
\begin{tikzcd}[ampersand replacement=\&]
A_0 \arrow[r] \arrow[d] \& A_1 \arrow[d] \\
A_2 \arrow[r] \& C
\end{tikzcd}
\]
est aussi cocartésien.}

\textit{Soit $W$ une partie de} $\Fl(\M)$ \textit{satisfaisant les conditions suivantes.}
\[
\left\{
\begin{array}{@{\quad}l@{\quad}p{0.6\textwidth}}
\text{a)} & \textit{$W$ est un localiseur modérément saturé.} \\[1.5ex]
\text{b)} & \textit{$W$ est L-accessible,} i.e. \textit{il est stable par $\varinjlim$ filtrantes, et est accessible.} \\[1.5ex]
\text{c)} & $\Mon \cap W$ \textit{est stable par cochangement de base,} i.e. \textit{si on a un diagramme} cocartésien
\begin{center}
\begin{tikzcd}[ampersand replacement=\&]
A \arrow[r, hook, "i"] \arrow[d, "\alpha"'] \arrow[dr, phantom, "\scriptstyle\text{cocart.}"] \& B \arrow[d, "\beta"] \\
A' \arrow[r, "i'"] \& B'
\end{tikzcd}
\end{center}
\textit{avec $i$ monomorphisme appartenant à $W$, alors $i'$ est un monomorphisme appartenant à $W$.}\,\footnotemark[32]
\end{array}
\right.
\tag*{(7.2.3)}
\]
\footnotetext[32]{\textit{Cf.} plus bas pour une quatrième condition $(d)$ sur $W$, qui assure que $W$ s'insère dans un triple de Quillen.}

\msp{page 89}
\textit{Sous ces conditions, on a}
\[
\underbrace{\Mon}_{C} \cap W \;=\; \widetilde{W \cap \Mon} \quad , \qquad \text{i.e. } \Mon \cap W \text{ \textit{est Q-clos à gauche}} \;,
\tag*{(7.2.4)}
\]


\endgroup
% END UNIT Der-ChXV-05-11-10/Der-ChXV-05-11-10_p046-050.tex
% BEGIN UNIT Der-ChXV-05-11-10/Der-ChXV-05-11-10_p051-055.tex
\begingroup\makeatletter%
\newcommand{\printedpage}[1]{\marginpar{\footnotesize p.~#1}}%
\newcommand{\msp}[1]{\marginpar{\footnotesize [#1]}}%
\newcommand{\Mon}{\mathrm{Mon}}%
\newcommand{\Ob}{\mathrm{Ob}}%
\newcommand{\Fl}{\mathrm{Fl}}%
\newcommand{\Ss}{\mathrm{Ss}}%
\newcommand{\card}{\mathrm{card}}%
\newcommand{\cercle}[1]{\tikz[baseline=(c.base)]\node[draw,circle,inner sep=1pt](c){#1};}%
\makeatother\let\printedpage\derpp\let\msp\dermsp%
\let\mspin\DERmspin%


\printedpage{49}
% head: A. Grothendieck : Dérivateurs, Ch. XV, Édition des Universités de Montpellier II et Paris VII
% foot: Version 5 novembre 2010

\textit{et}
\[
\text{\textit{il existe une petite partie} } TC_0 \subset \Mon, \text{ \textit{telle que} } \Mon \cap W = \widetilde{TC_0}\ . \tag*{(7.2.5)} % [sic] (7.2.5) -> numéro probablement autre (peut-être (7.2.3)), le même numéro (7.2.5) est repris plus bas pour une autre formule et précède (7.2.4)
\]

\textit{(De façon plus précise, si} $\pi \geq \pi_0$ \textit{est un cardinal tel que} $W$ \textit{soit} $\pi$\textit{-accessible, on peut prendre} $TC_0 = (\Mon \cap W)^{\pi}$ \textit{(ou une partie} petite \textit{de} $(\Mon \cap W)^{\pi}$ \textit{ayant les mêmes classes d'isomorphie).) Donc} (th.~1 de §1) \textit{le couple} $(TC = \Mon \cap W,\ F = (TC)_*)$ \textit{est un couple de Quillen clos, et il y a même une factorisation} fonctorielle \textit{de toute flèche} $f$ \textit{en}
\[
f \ = \ p\, i \ ; \qquad i \ \in \ TC \ , \qquad p \ \in \ F \ . \tag*{(7.2.4)}
\]
\textit{Si on pose}
\[
C \ = \ \Mon \ , \qquad TC \ = \ \Mon \cap W \ , \qquad F \ = \ (TC)_* \ , \qquad TF \ = \ C_* \ , \tag*{(7.2.5)}
\]
\textit{le triple} $(W, C, F)$ \textit{est un triple de Quillen clos, si et seulement si on a les conditions supplémentaires suivantes.}
\[
\left\{
\begin{array}{lll}
\text{(7.2.6)} & \cercle{g}\ \text{\textit{(sur $\mathcal{M}$ seulement)}} & \parbox[t]{0.36\textwidth}{Mon \textit{stable par cochangement de base.}} \\[4ex]
\text{(7.2.7)} & \underbrace{TF}_{\overset{\text{déf}}{=}\,(\Mon)_*} \ \subset \ W & \parbox[t]{0.36\textwidth}{i.e. \textit{si} d) \textit{toute} $f : X \longrightarrow Y$ \textit{qui fait de} $X$ \textit{un objet injectif de} $\mathcal{M}/Y$\textit{, est dans} $W$ \textit{(donc aussi dans} $W_{\mathrm{univ}}$\textit{).}}
\end{array}
\right.
\]

\msp{page 90}

\textsc{Démonstration de} 7.1. On peut donner un énoncé un peu plus délicat, en supposant d'emblée donnée, non une partie $M_0$ de $\Ob \mathcal{M}$, mais une partie
\[
\Mon_0 \ \subset \ \Mon \ , \tag*{(7.1)}
\]
satisfaisant la condition suivante

\begin{quote}
Pour tout monomorphisme $i : A \longrightarrow B$ dans $\mathcal{M}$, il existe $i' \in \Mon_0$, $i' : A' \longrightarrow B'$, et une flèche $i' \longrightarrow i$ dans $\underline{\Fl}(\mathcal{M})$,
\[
\begin{tikzcd}[row sep=large, column sep=large]
A' \arrow[r, "{i'}"] \arrow[d] & B' \arrow[d] \\
A \arrow[r, "i"] & B \ ,
\end{tikzcd}
\tag*{(7.2)}
\]
telle que la flèche correspondante
\[
C \ = \ A \amalg_{A'} B' \ \longrightarrow \ B
\]
soit un monomorphisme, et que $B' \longrightarrow B$ ne se factorise pas par $A$.
\end{quote}

C'est cela qui était visé par la condition (7.1.1) sur $M_0$, qui implique (7.2) sur $\Mon_0$ défini en termes de $M_0$ comme dit dans 1°.

\printedpage{50}
% head: A. Grothendieck : Dérivateurs, Ch. XV, Édition des Universités de Montpellier II et Paris VII
% foot: Version 5 novembre 2010

Dans ce qui suit, on suppose donnée une partie $\Mon_0 \subset \Fl(\mathcal{M})$ satisfaisant (7.2). En pratique, elle sera souvent déduite d'une partie $M_0$ de $\Ob \mathcal{M}$, génératrice pour les monomorphismes, comme explicité dans 1°, mais on n'a pas besoin du fait que $i' \longrightarrow i$ soit un monomorphisme (\textit{i.e.} (7.1.1) est plus forte que ce qu'il nous faut). P. ex. si $B' \in M_0$ et $B' \longrightarrow B$ est donné, ne se factorisant pas par $A$, on pourra prendre $A' = B' \times_B A$ et $i' : A' \longrightarrow B'$ (pour assurer $i' \in \Mon_0$, il faudra donc supposer $M_0$ stable par passage aux sous-objets), et il y a des chances

\msp{page 91}

que $C \longrightarrow B$ sera alors bel et bien un monomorphisme \ldots

1°) Il faut prouver
\[
\Mon \ \subset \ \widetilde{\Mon_0}
\]
moyennant (7.2) (et la stabilité de $\mathcal{M}$ par $\varinjlim$ filtrantes). Cela va résulter directement du

\textbf{Lemme 7.3.} \textit{Soient $\mathcal{M}$ une catégorie stable par limites inductives filtrantes, et $C_0$, $C$ des parties de} $\Fl(\mathcal{M})$\textit{. On supposera}
\begin{enumerate}[label={}]
\item[\cercle{a}] \textit{$C$ formée de monomorphismes, et}
\item[\cercle{b}] \textit{que la condition suivante (généralisant} (7.2) \textit{ci-dessus) soit vérifiée.}
\end{enumerate}

\begin{quote}
\textit{Pour toute flèche $i : A \hookrightarrow B$ dans $C$ telle que $i$ ne soit pas un isomorphisme, il existe $i' : A' \longrightarrow B'$ dans $\widetilde{C_0}$, et une flèche $i' \longrightarrow i$ dans} $\underline{\Fl}(\mathcal{M})$\textit{, telle que dans le diagramme}
\[
\begin{tikzcd}[row sep=large, column sep=large]
A' \arrow[r, "{i'}"] \arrow[d, hook] & B' \arrow[d] \arrow[rd, dotted] & \\
A \arrow[r, "\alpha"] \arrow[rr, dotted, bend right, "i"'] & D \ = \ A \amalg_{A'} B' \arrow[r, "\beta"] & B \ ,
\end{tikzcd}
\tag*{(7.3.1)}
\]
\textit{la flèche $\beta$ soit dans $C$, et que $B' \longrightarrow B$ ne se factorise pas par $A$ (ou, ce qui revient au même, que le sous-objet de $B$ défini par $D$ et $\beta$ ne soit pas égal à celui défini par $A$ et $i$).}
\end{quote}

\textit{On suppose de plus}
\begin{enumerate}[label={}]
\item[\cercle{c}] \textit{$C$ stable par} $\varinjlim$ \textit{filtrantes, et}
\item[\cercle{d}] \textit{que pour tout $B \in \Ob \mathcal{M}$, l'ensemble des sous-objets de $B$ est petit.} (Condition des plus anodines! C'est une condition \textit{sur} $\mathcal{M}$, n'impliquant pas $C$).
\end{enumerate}
\textit{Alors on a}
\[
C \ \subset \ \widetilde{C_0} \ , \qquad \text{\textit{donc} } \widetilde{C} \ \subset \ \widetilde{C_0} \ , \tag*{(7.3.2)}
\]
\textit{donc, si $C_0 \subset C$, on conclut}

\msp{page 92}
\[
\widetilde{C} \ = \ \widetilde{C_0} \ . \tag*{(7.3.3)}
\]

\printedpage{51}
% head: A. Grothendieck : Dérivateurs, Ch. XV, Édition des Universités de Montpellier II et Paris VII
% foot: Version 5 novembre 2010

(Ce lemme implique aussitôt la première assertion de 7.1, 1°.) Prouvons ce lemme, \textit{i.e.} que toute $i \in C$ est dans $\widetilde{C_0}$, en construisant par récurrence transfinie une suite strictement croissante de sous-objets de $B$,
\[
A_\alpha \ \overset{j_\alpha}{\hookrightarrow} \ B \ , \qquad \text{avec les } j_\alpha \in C \ ,
\]
et avec :

\cercle{1°} $A_0 = A$. On prouvera que $A \xrightarrow{\ i_\alpha \,=\, i_{\alpha,0}\ } A_\alpha$ dans $\widetilde{C_0}$, et il y aura un ordinal final $\alpha$ tel que $A_\alpha = B$\,; on en conclura donc que $i = i_\alpha$ est dans $\widetilde{C_0}$.

Si $i$ est un isomorphisme, on sait que $i \in \widetilde{C_0}$, il n'y a rien à prouver (la construction transfinie s'arrête au cran 0). Les deux pas de la construction, après le pas initial 1°, sont :

\cercle{2°} Soit $A_\alpha$ construit. Si $A_\alpha = B$, la construction est arrêtée. Si $A_\alpha \xrightarrow{j_\alpha} B$ n'est pas un isomorphisme, on applique (7.3.1) à $j_\alpha \in C$, pour trouver un $D \overset{j}{\hookrightarrow} B$, donc $j \in C$ et $D \supsetneqq A_\alpha$. De plus,
\[
i_{\alpha+1,\alpha} \ : \ A_\alpha \ \longrightarrow \ A_{\alpha+1} \ \overset{\text{déf}}{=} \ D \qquad \text{est dans } \widetilde{C_0} \ , \tag*{(7.3.4)}
\]
car déduit de $i' \in \widetilde{C_0}$ par cochangement de base.

\cercle{3°} Si $\alpha$ est ordinal limite, on pose
\[
A_\alpha \ = \ \varinjlim_{\beta < \alpha} A_\beta \ , \tag*{(7.3.5)}
\]

\msp{page 93}

et la flèche canonique
\[
A_\alpha \ \xrightarrow{j_\alpha} \ B
\]
est dans $C$, car les $j_\beta$ le sont et $C$ est stable par limites inductives filtrantes.

On a ainsi achevé la description des systèmes transfinis. Pour tout $\alpha$, on a
\[
\card \alpha \ \leq \ \card \Ss(B) \ ,
\]
puisque l'application $\beta \longmapsto A_\beta$ de l'ensemble $I_\alpha = \{\beta \mid \beta < \alpha\}$ dans $\Ss(B)$ est \textit{strictement} croissante. Donc il existe un cardinal $\alpha$ tel que $\card \alpha \leq \card \Ss(B)$, tel que $A_\alpha = B$. Mais alors $i \in \widetilde{C_0}$, comme limite inductive (ou composé) transfinie de flèches appartenant à $\widetilde{C_0}$. (Il suffit de la condition $\forall\ \beta$, $i_{\beta+1,\beta} \in \widetilde{C_0}$ pour impliquer que tous les morphismes de transition $i_{\beta',\beta}$ sont dans $\widetilde{C_0}$.)

Revenons à 7.1, 2°. Il reste à prouver que si Mon est stable par cochangement de base, alors on a
\[
\Mon \ = \ \widetilde{\Mon_0} \quad .
\]

Mais comme Mon est déjà stable par composition, et par limites inductives filtrantes quelconques (donc aussi par facteurs directs), il s'ensuit par le corollaire 2 du théorème 1, §1 (p. 11), que $\Mon = \widetilde{\Mon_0}$\ \ . L'énoncé de factorisation est le théorème 1 du §1.

\printedpage{52}
% head: A. Grothendieck : Dérivateurs, Ch. XV, Édition des Universités de Montpellier II et Paris VII
% foot: Version 5 novembre 2010

\msp{page 94}

Prouvons 7.1, 3°, \textit{i.e.} la formule (7.1.6),
\[
W \cap \Mon \ \subset \ \widetilde{(W \cap \Mon)^{\pi}} \ .
\]

Pour ceci, on applique encore le lemme 7.3 à
\[
C_0 \ = \ (W \cap \Mon)^{\pi} \ , \qquad C \ = \ W \cap \Mon \ .
\]

Les conditions (\textit{a}), (\textit{b}), (\textit{c}), (\textit{d}) du lemme sont vérifiées. C'est clair pour toutes, sauf pour (\textit{b}), qui demande une vérification non tautologique. Soit donc
\[
f \ : \ A \hookrightarrow B \ , \qquad f \in \Mon \cap W,\ f \text{ non isomorphisme} \ .
\]

On considère $B$ comme limite inductive\footnote[33]{\textit{D'un système inductif $\pi$-adapté.}}
\[
B \ = \ \varinjlim_{I} B_i \ , \qquad I \ = \ \Ss_{\pi}(B) \ .
\]
Pour tout $i$, on pose
\[
A_i \ = \ A \cap B_i \ = \ A \times_B B_i \ ,
\]
et
\[
f_i \ : \ A_i \ \hookrightarrow \ B_i \ ,
\]
donc
\[
f \ = \ \varinjlim_{i \in I} f_i \ .
\]
On a besoin de savoir que le système inductif des flèches $f_i$ est $\pi$-adapté, ce qui nous amène à postuler la condition d'exactitude supplémentaire suivante sur $\mathcal{M}$ :

\msp{page 95}
\[
e')\quad \text{les } \varinjlim \text{ filtrantes dans } \mathcal{M} \text{ commutent au changement de base}\ , \tag*{(7.4)}
\]

il suffit même de le supposer dans le cas très particulier qu'on a envisagé ici.

Comme $W$ est $\pi$-accessible, et $(f_i)_{i \in I}$ est $\pi$-adapté, comme $\varinjlim f_i \in W$, il s'ensuit que l'ensemble $I_W$ des $i \in I$ tels que $f_i \in W$ est cofinal. Comme $A \overset{f}{\hookrightarrow} B$ n'est pas un isomorphisme, il existe donc un $i \in I_W$ tel que $B_i \hookrightarrow B$ ne se factorise pas par $A$. On a alors le diagramme
\[
\begin{tikzcd}[row sep=large, column sep=large]
A_i \arrow[r, hook, "f_i"] \arrow[d, hook] \arrow[rd, phantom, "\text{\footnotesize cocart.}"] & B_i \arrow[d, hook] \arrow[rd, dotted, hook, "\text{\footnotesize incl.}"] & \\
A \arrow[r, hook, "\alpha_i"] \arrow[rr, dotted, hook, bend right, "f"'] & C_i \arrow[r, "\beta_i"] & B \ ,
\end{tikzcd}
\]

\printedpage{53}
% head: A. Grothendieck : Dérivateurs, Ch. XV, Édition des Universités de Montpellier II et Paris VII
% foot: Version 5 novembre 2010
avec $f_i \in (C \cap W)^{\pi} = C_0$, donc $\alpha_i \in \widetilde{C_0}$. On a besoin de savoir que $\alpha_i \in W$, donc une condition supplémentaire sur $W$ :

\begin{quote}
Si on a un carré cocartésien
\[
\begin{tikzcd}[row sep=large, column sep=large]
A \arrow[r, hook, "i"] \arrow[d] & B \arrow[d] \\
A' \arrow[r, "{i'}"] & B' \ ,
\end{tikzcd}
\tag*{(7.5)}
\]
avec $i \in \Mon \cap W$, alors $i' \in W$\footnote[34]{\textbf{N.B.} Cette condition est vérifiée.}.
\end{quote}

(j'ai renforcé (7.1.4) en conséquence).

\msp{page 96}

Comme par hypothèse $f = \beta_i \alpha_i$ est dans $W$, et $\alpha_i \in W$, il s'ensuit par (7.1.4), (\textit{a}) que $\beta_i \in W$. Il faudra que nous sachions de plus que $\beta_i$ est un monomorphisme, donc $\beta_i \in \Mon \cap W = C$, ce qui achèvera d'établir la condition manquante (\textit{b}) de 7.3. C'est la condition d'exactitude (7.2.1), qu'il me faut visiblement rajouter aux hypothèses dans 7.1, 3°. Ayant fait cela, le lemme 7.3 s'applique et établit 7.1, 3°.

\textsc{Démonstration de} 7.2. Posons
\[
C \ = \ \Mon \ , \qquad TC \ = \ C \cap W \ , \qquad TC_0 \ = \ TC^{\pi} \ ,
\]
donc par 7.1, 3°, on a
\[
\underbrace{TC_0}_{\substack{\text{essentiellement}\\ \text{petite}}} \ \subset \ TC \ \subset \ \widetilde{TC_0} \ = \ \widetilde{TC} \ ,
\]
et par le théorème 1, §1 (p. 10-12), cor. 2, pour vérifier que $TC = \widetilde{TC_0}$ (\textit{i.e.} $TC$ est Q-clos à gauche), il revient au même de voir qu'il est Q-stable à gauche, \textit{i.e.} stable par composition, par cochangement de base, par $\varinjlim$ ordinales strictes, par facteurs directs. Or les deux derniers sont impliqués par le fait que $TC = C \cap W$ est stable par $\varinjlim$ filtrantes, vu que $W$ l'est, et $C$ aussi par la condition \cercle{e} (p. 76) sur $\mathcal{M}$. D'autre part, $C = \Mon$ et $W$ étant stables par composition, itou pour $C \cap W$. Donc il reste le cochangement de

\msp{page 97}

base, et c'est assuré par 7.2.3, (\textit{c}) (dont on avait déjà besoin dans 7.1, 3°, pour établir $C \cap W \subset \widetilde{TC_0}$). (\textbf{N.B.} les hypothèses de 7.2 sont pratiquement celles de 7.1, 3°, sauf qu'on suppose que $W$ est un localiseur (sous-entendu : modérément saturé), de façon précise on n'a utilisé, jusqu'à présent, comme hypothèse supplémentaire, que la stabilité de $W$ par composition.)

Il reste à établir que $(W, C = \Mon, F = (TC)_*)$ est un triple de Quillen clos, si et seulement si les conditions (7.2.6) et (7.2.7), \textit{i.e.}
\[
\left\{
\begin{aligned}
& C = \Mon \ \ \text{stable par cochangement de base,} \\
& TF = C_* = (\Mon)_* \subset W,
\end{aligned}
\right.
\]


\endgroup
% END UNIT Der-ChXV-05-11-10/Der-ChXV-05-11-10_p051-055.tex
% BEGIN UNIT Der-ChXV-05-11-10/Der-ChXV-05-11-10_p056-060.tex
\begingroup\makeatletter%
\newcommand{\printedpage}[1]{\marginpar{\footnotesize p.~#1}}%
\newcommand{\msp}[1]{\marginpar{\footnotesize [#1]}}%
\newcommand{\Mon}{\mathrm{Mon}}%
\newcommand{\Fl}{\mathrm{Fl}}%
\newcommand{\Cof}{\mathrm{Cof}}%
\newcommand{\Fib}{\mathrm{Fib}}%
\newcommand{\Cat}{\mathrm{Cat}}%
\newcommand{\univ}{\mathrm{univ}}%
\newcommand{\chap}{^{\wedge}}%
\makeatother\let\printedpage\derpp\let\msp\dermsp%
% Shared placement decision BUILD-MARGIN-FR-XV-58
\def\DERmargintarget{58}%
\def\printedpage#1{\def\DERmarginarg{#1}\ifx\DERmarginarg\DERmargintarget\marginpar{\raggedright\footnotesize\color{black!55}\hspace*{6pt}p.~#1}\else\derpp{#1}\fi}%
\let\mspin\DERmspin%


\printedpage{54}
% head: A. Grothendieck : Dérivateurs, Ch. XV, Édition des Universités de Montpellier II et Paris VII
% foot: Version 5 novembre 2010
sont satisfaites. La nécessité est évidente, par définition des triples de Quillen clos, et pour la suffisance, on note qu'il résulte du théorème 7.1, 2$^\circ$, et de la partie déjà prouvée de 7.2, que $(C,TF)$, $(TC,F)$ sont des couples de Quillen clos, et on a de plus
\[
TC \;=\; C\cap W\ , \qquad TF \;\subset\; F\cap W
\]
(puisque $TF\subset W$ par (7.2.7), et $TF\subset F$ puisque $TC\subset C$). On sait
\msp{page 98}
(paragraphe 4.3, corollaire 4, page 36) qu'il en résulte que $(W,C,F)$ est un triple de Quillen clos. Cela achève la démonstration de 7.1 et 7.2.

Résumons l'essentiel :

\textbf{Corollaire 7.4.}\footnote[35]{\textbf{N.B.} C'est finalement ce corollaire qui est le résultat principal de la présente section.} \textit{Soit $\mathcal{M}$ une catégorie satisfaisant les conditions} (a), (b), (c), (d), (e) (p.~76) \textit{et} (e$'$), (f) ((7.2.1), p.~88), \textit{enfin} (g) (7.2.6), i.e. \textit{stabilité de l'ensemble} Mon \textit{des monomorphismes de $\mathcal{M}$ par cochangement de base. Soit $W\subset \Fl(\mathcal{M})$ un localiseur dans $\mathcal{M}$, qui est} stable par $\varinjlim$ filtrantes \textit{et de plus} (\textit{condition plus subtile, mais pratiquement anodine}) accessible. \textit{On pose}
\[
C \;=\; \Mon\ , \qquad TC \;=\; C\cap W\ , \qquad TF \;=\; C_*\ , \qquad F \;=\; (TC)_*\ .
\]
\textit{Alors :}
\begin{itemize}[leftmargin=*]
\item[1$^\circ$)] \textit{$(C,TF)$ est un couple de Quillen} (\textit{énoncé indépendant de la donnée de $W$}), \textit{et même un couple de Quillen à factorisation fonctorielle.}
\item[2$^\circ$)] \textit{Pour que $(TC,F)$ soit un couple de Quillen, il faut et il suffit que $TC$ soit stable par cochangement de base.} (\textit{Et alors il y a factorisation fonctorielle.})\footnote[36]{\textbf{N.B.} Dans les cas 1$^\circ$ et 2$^\circ$, $C$, $TC=C\cap W$ sont Q-engendrés à gauche par des \textit{petites} parties.}.
\item[3$^\circ$)] \textit{Pour que $(W,C,F)$ soit un triple de Quillen clos, il faut et il suffit que les conditions suivantes soient satisfaites.}
\[
\left\{
\begin{aligned}
&\text{a)}\quad TC \overset{\text{déf}}{=} W\cap \Mon\ \textit{stable par cochangement de base.}\\
&\text{b)}\quad TF \overset{\text{déf}}{=} (\Mon)_* = \textit{applications de type injectif} \subset W.
\end{aligned}
\right.
\tag*{(7.4.1)}
\]
\item[4$^\circ$)] \msp{page 99}\textit{Pour que $(W,C,F)$ soit un triple de Quillen clos et propre à gauche} (i.e. $C\subset \Cof_W$), \textit{il faut et il suffit qu'on ait la condition} (7.4.1), (b) (i.e. \textit{$W$ contient les morphismes de type injectif}), \textit{et de plus}
\[
\begin{aligned}
&C\subset \Cof_W,\ \textit{ou encore : pour tout carré cocartésien}\\
&\qquad\qquad\begin{tikzcd}[ampersand replacement=\&, row sep=large, column sep=large]
A \arrow[r, "i"] \arrow[d, "\alpha"'] \& B \arrow[d, "\beta"] \\
A' \arrow[r, "i'"] \& B'\ ,
\end{tikzcd}\\
&\textit{avec $i$ un monomorphisme, alors } \alpha\in W \Longrightarrow \beta\in W.
\end{aligned}
\tag*{(7.4.2)}
\]

\item[5$^\circ$)] \printedpage{55}
% head: A. Grothendieck : Dérivateurs, Ch. XV, Édition des Universités de Montpellier II et Paris VII
% foot: Version 5 novembre 2010
\textit{Pour que $(W,C,F)$ soit un triple de Quillen clos et} propre, \textit{il faut et il suffit que les deux conditions de 4$^\circ$ soient satisfaites, ainsi que l'on ait :}
\[
\begin{aligned}
&F\subset \Fib_W,\ \textit{ou encore : pour tout carré cartésien}\\
&\qquad\qquad\begin{tikzcd}[ampersand replacement=\&, row sep=large, column sep=large]
X \arrow[r, "q"] \arrow[d, "f"'] \& X' \arrow[d, "f'"] \\
Y \arrow[r, "p"'] \& Y'\ ,
\end{tikzcd}\\
&\textit{si } f\in F \ (\text{i.e. } f\in (C\cap W)_*)\ \textit{et } p\in W,\ \textit{on a } q\in W.
\end{aligned}
\tag*{(7.4.3)}
\]
\end{itemize}

\textsc{Démonstration de} 7.4. Les parties 1$^\circ$, 2$^\circ$, 3$^\circ$ ne sont qu'un résumé de ce qui a été établi. Pour 4$^\circ$, il faut voir que les conditions dites impliquent (7.4.1), $(a)$, or on a $C\subset \Cof_W$, donc $C\cap W\subset \Cof_W\cap W = W^{\univ}$
\msp{page 100}
donc $C\cap W = C\cap W^{\univ}$, qui est stable par cochangement de base, puisque $C$ et $W^{\univ}$ le sont. Quant à 5$^\circ$, c'est une tautologie, compte tenu de 4$^\circ$.

\section*{8\quad Application aux catégories de modèles $(A\chap,W_A)$}

Soit $A$ dans Cat, et considérons dans $A\chap$ le localiseur $W_A$, image inverse par
\[
\left\{
\begin{aligned}
i_A : A\chap \quad &\longrightarrow \quad \Cat\\
F \quad &\longmapsto \quad A/F
\end{aligned}
\right.
\tag*{(8.1)}
\]
d'un localiseur fondamental donné $W$, satisfaisant les conditions
\[
\mathrm{W}\,(\underbrace{1}_{\text{sat.}},\underbrace{2}_{\text{loc.}},\underbrace{3}_{\text{objet fin.}}\,,\underbrace{5}_{\text{conn. }\pi_0}\,,\underbrace{6}_{\varinjlim},\underbrace{7\text{ bis}}_{\substack{\text{carrés coc.}\\ \text{ou ``M.V.''}}}\,,\underbrace{8}_{\text{acc.}}),\quad
\text{\parbox[t]{4.3cm}{où W\,(8) est l'accessibilité de $W$ (donc $W$ est L-accessible s'il satisfait aussi W\,(6)).}}
\tag*{(8.2)}
\]

Pour pouvoir cependant appliquer les développements sur les parties accessibles d'une grosse catégorie du §5, il faut vérifier les conditions $(a)$, $(b)$, $(c)$, $(d)$, $(e)$ du §5, tant dans Cat, que (surtout) dans $A\chap$.

a) Stabilité par petites $\varinjlim$ filtrantes, O.K.

b) Filtration cardinale. Bien connu pour $A\chap = \underline{\mathrm{Hom}}(A,\mathrm{Ens})$, à partir de celle de Ens par les cardinaux (\textit{cf.} SGA 4 I, 9.24, p.~177). Pour Cat, on définit une filtration cardinale
\[
(\Cat^{\pi})_{\pi\geq\pi_0}\ , \qquad \text{où } \pi_0 = \text{cardinal des ensembles dénombrables}\ ,
\tag*{(8.3)}
\]
\printedpage{56}
% head: A. Grothendieck : Dérivateurs, Ch. XV, Édition des Universités de Montpellier II et Paris VII
% foot: Version 5 novembre 2010
\msp{page 101}
par la condition, pour $A$ dans Cat,
\[
A\in \mathrm{Filt}^{\pi}(\Cat) \;=\; \Cat^{\pi} \qquad\Longleftrightarrow\qquad \operatorname{card}\Fl(A) \;\leq\; \pi\ .
\tag*{(8.4)}
\]

Il faut vérifier les conditions $(a)$, $(b)$, $(c)$ de \textit{loc. cit.} 9.12. La condition $(a)$ est évidente, $(b)$ est contenue dans la condition plus forte de stabilité de $\Cat^{\pi}$ par $\varinjlim\limits_{I}$, pour $I$ une catégorie telle que $\operatorname{card}\Fl(I)\leq\pi$ (c'est la condition $(d_0)$ de ces notes, p.~45-46). Notons que $\Cat^{\pi}$ est également stable par $\varprojlim$ finies (et si $\pi$ est de la forme $2^c$, il est stable même par limites projectives $\varprojlim\limits_{I}$, avec $\operatorname{card}\Fl(I)\leq c\ \ldots$). La condition $(c)$ se vérifie aisément par la méthode du §5, en notant que pour $X$ dans Cat, l'ensemble
\[
\mathrm{Ss}_{\pi}(X)
\]
des sous-catégories de $X$ qui sont dans $\Cat^{\pi}$ est filtrant croissant et grand devant $\pi$, et que $X$ en est la limite inductive. D'autre part, il est clair aussi que si $X\in\Cat^{\pi'}$ avec $\pi'>\pi$, alors
\[
\operatorname{card}\mathrm{Ss}_{\pi}(X) \;\leq\; \pi'^{\pi}\ ,
\]
ce qui achève d'établir la condition $(c)$.

De plus, la condition $(d)$ de p.~45 se vérifie immédiatement pour les $\Cat^{\pi}$, et aussi pour les $A^{\wedge\pi}$ (pour $\pi\geq\operatorname{card}\Fl(A)$).

Aussi que les monomorphismes sont stables par $\varinjlim$ filtrantes (condition $(e)$, page 46).

\msp{page 102}
Passons aux conditions supplémentaires $(e')$, $(f)$, $(g)$ exigées dans le corollaire 7.4. La condition $(e')$ (commutativité des $\varinjlim$ filtrantes au changement de base) est satisfaite dans $A\chap$ et dans Cat. La condition $(f)$ dit que pour deux sous-objets $A_1$, $A_2$ d'un objet $B$, si $A_0=A_1\cap A_2$, alors
\[
A_1\amalg_{A_0}A_2 \longrightarrow B
\]
est un monomorphisme. Je doute que c'est vrai dans Cat, et ce n'est pas grave, car c'est à $A\chap$, non à Cat, qu'on voudra appliquer 7.4 -- et c'est évidemment vrai dans $A\chap$, et dans tout topos, l'étant dans (Ens). Enfin, la stabilité des monomorphismes par cochangement de base est sûrement fausse dans Cat, mais vraie dans $A\chap$ -- donc c'est O.K.

Comme $W$ est stable par $\varinjlim$ filtrantes (W\,(6)), et que $i_A$ y commute, $W_A$ est également stable. De même, $W_A$ est L-accessible, en vertu du

\textbf{Lemme 8.1.} \textit{Soient $\mathcal{M}$, $\mathcal{M}'$ deux catégories satisfaisant} (a), (b), (c), (d), (e), \textit{et $\varphi:\mathcal{M}\longrightarrow\mathcal{M}'$ un foncteur commutant aux $\varinjlim$ filtrantes. Alors pour toute partie L-accessible $M'$ de $\mathcal{M}'$, $\varphi^{-1}(M')$ est L-accessible.}

\msp{page 103}
\textit{De façon précise, soit $\pi$ un cardinal tel que $M'$ soit $\pi$-accessible et que $\varphi$ applique $\mathcal{M}^{\pi}$ dans $\mathcal{M}'^{\pi}$ (par exemple, pour cette dernière condition, il suffit $\pi=2^c$ pour $c$ assez grand,} cf. loc. cit. 9.14, p.~153). \textit{Alors $M=\varphi^{-1}(M')$ est $\pi$-accessible.}

Vérification immédiate.

On a donc le

\printedpage{57}
% head: A. Grothendieck : Dérivateurs, Ch. XV, Édition des Universités de Montpellier II et Paris VII
% foot: Version 5 novembre 2010
\textbf{Lemme 8.2.} \textit{Soit $A$ dans} Cat. \textit{Alors $(A\chap,W_A)$ satisfait aux conditions préliminaires du corollaire} 7.4, \textit{ainsi qu'aux conditions suivantes.}
\[
\underbrace{\Mon(A\chap)}_{\substack{\textit{ensemble des}\\ \textit{monomorphismes}\\ \textit{dans } A\chap}} \quad\subset\ \Cof_{W_A}\ ,
\tag*{(8.2.1)}
\]
\[
\Mon\cap W_A \qquad \textit{stable par cochangement de base}\ ,
\tag*{(8.2.2)}
\]
\textit{en d'autres termes} (\textit{compte tenu que} Mon \textit{est stable par cochangement de base}), \textit{pour un carré cocartésien}
\[
\begin{tikzcd}[row sep=large, column sep=large]
A \arrow[r, hook, "i"] \arrow[dr, "\alpha"'] & B \arrow[dr, "\beta"] & \\
& A' \arrow[r, hook, "i'"] & B'\ ,
\end{tikzcd}
\tag*{(8.2.3)}
\]
\textit{avec $i$ un monomorphisme, si $\alpha\in W_A$ (resp. $i\in W_A$), on a $\beta\in W_A$ (resp. $i'\in W_A$). En fait, on a mieux :}
\[
i\ \in\ W_A \ \Longleftrightarrow\ i'\ \in\ W_A\ , \qquad \alpha\ \in\ W_A \ \Longleftrightarrow\ \beta\ \in\ W_A\ .
\tag*{(8.2.4)}
\]

\msp{page 104}
Il reste à vérifier (8.2.4), qui a été vu dans XII \ldots

On conclut donc de 7.4 1$^\circ$, 2$^\circ$, 3$^\circ$ :

\textbf{Corollaire 8.3.} \textit{Posons}
\[
\left\{
\begin{aligned}
C \;&=\; \Mon(A\chap)\ , & TC \;&=\; C\cap W\ ,\\
F \;&=\; (TC)_*\ , & TF \;&=\; C_*\ .
\end{aligned}
\right.
\tag*{(8.3.1)}
\]
\textit{Alors $(C,TF)$ et $(TC,F)$ sont des couples de Quillen à factorisation fonctorielle.} (\textbf{N.B.} \textit{$C$ et $TC$ étant de plus Q-engendrés à gauche par des petites parties de} $\Fl(A\chap)$.) \textit{Pour que $(W_A,C,F)$ soit un triple de Quillen clos, il faut et il suffit qu'on ait}
\[
\underbrace{TF}_{\overset{\text{déf}}{=}\,(\Mon)_*} \quad\subset\ W_A\ ,
\tag*{(8.3.2)}
\]
i.e. \textit{que toute flèche $f:X\longrightarrow Y$ de type injectif dans $A\chap$} (i.e. \textit{qui fait de $X$ un objet injectif de $A\chap/Y$}) \textit{soit dans $W_A$.}

\textit{Quand cette condition est satisfaite, le triple de Quillen $(W_A,C,F)$ est propre à gauche, par} 8.2.1.

Il n'est pas clair qu'il soit forcément propre-tout-court, \textit{i.e.} qu'on ait
\[
F\ \subset \Fib_W \qquad (?)\ .
\tag*{(8.3.3)}
\]

Nous y reviendrons par la suite. Pour l'instant, la chose essentielle, c'est d'examiner la condition (8.3.2) sur $A$.

\printedpage{58}
% head: A. Grothendieck : Dérivateurs, Ch. XV, Édition des Universités de Montpellier II et Paris VII
% foot: Version 5 novembre 2010
\msp{page 105}
\textbf{Théorème 8.4.} \textit{Soit $W$ un localiseur fondamental dans} Cat \textit{satisfaisant} $\mathrm{W}\,(1,2,3,6,7,8)$, \textit{notations de} XII, (i.e. \textit{tous les axiomes envisagés jusqu'à présent, sauf peut-être celui} W\,(4) \textit{des fibrations} \texttt{[et celui W\,(5) de connexité]}), \textit{supposons de plus $W$ stable par facteurs directs} (\textit{renforcement} W\,(1\,bis) \textit{de l'axiome de saturation modérée} W\,(1)).\footnote[37]{\textbf{N.B.} Pour la validité de 1$^\circ$, 3$^\circ$ on n'a pas besoin de W\,(6) ( $\varinjlim$ ) ni de W\,(8) (accessibilité), ni même de W\,(7) (axiome du carré cocartésien). Ce dernier est nécessaire déjà pour assurer la partie $(a)$ du corollaire 8.5. Mais pour la partie 2$^\circ$ du théorème 8.4, on a besoin de \textit{tous} les axiomes du localiseur fondamental $W$, sauf seulement W\,(4) \texttt{[et W\,(5)]}.} \textit{Soit $A$ un objet de} Cat, \textit{et $W_A=i_A^{-1}(W)$ le localiseur de $A\chap$ déduit de $W$. Avec les notations de} 8.3 :
\begin{itemize}[leftmargin=*]
\item[1$^\circ$)] \textit{Les conditions suivantes sont équivalentes.}
\begin{itemize}[leftmargin=*]
\item[a)] \textit{On a} (8.3.2),
\[
TF\ \subset\ W_A\ ,
\]
i.e. \textit{toute flèche de type injectif dans $A\chap$ est dans $W_A$.}
\item[b)] \textit{Soit}
\[
L \;=\; L_A
\tag*{(8.4.1)}
\]
\textit{l'élément de Lawvere dans $A\chap$, donc le préfaisceau}
\[
a \quad\longmapsto\quad
\begin{aligned}
&\textit{ensemble des sous-préfaisceaux de $a$}\\
&(\text{i.e. }\textit{ensemble des ouverts de $A/a$},\\
&\text{i.e. }\textit{des cribles de $A/a$})\ ,
\end{aligned}
\]
\textit{et considérons la flèche $L\longrightarrow e_A$ (objet final de $A\chap$). On veut que cette flèche soit $W_A$-asphérique,} i.e. \textit{soit universellement dans $W_A$ -- ou ce qui revient au même, que pour tout $a\in A$, $L_A\times a$ $(=L_{A/a})$ soit un objet $W_A$-asphérique}
\msp{page 106}
\textit{de $A\chap$ (ou, ce qui revient au même, de $(A/a)\chap$).}
\item[c)] \textit{Il existe un objet $I$ de $A\chap$ tel que $I\longrightarrow e_A$ soit $W_A$-asphérique, avec deux sections $\delta_0$, $\delta_1$ de $I$ sur $e_A$, avec} $\mathrm{Ker}(\delta_0,\delta_1)=\emptyset_A$ \textit{(préfaisceau vide).}
\item[d)] \textit{Les catégories localisées $A/a=B$ sont des} catégories-test \textit{pour $W$} (cf. \textit{Pursuing Stacks}, \textit{§65, notamment pages 172-173), ce qui revient à dire que l'on a ceci, pour chacune des catégories $B=A/a$ :}
\[
\begin{aligned}
&\textit{Les foncteurs}\\
&\left\{
\begin{aligned}
&i_B:\ \ B\chap \ \longrightarrow\ \Cat\ \ (\ F\ \longmapsto\ B/F\ )\\
&j_B:\ \ \Cat\ \longrightarrow\ B\chap\ \ (\ C\ \longmapsto\ (b\longmapsto \mathrm{Hom}(B/b,C))\ )
\end{aligned}
\right.
\end{aligned}
\tag*{(8.4.2)}
\]
(\textit{où $j_B$ est l'adjoint à droite de $i_B$}) \textit{sont tous les deux ``model preserving''}, i.e.
\[
\left\{
\begin{aligned}
W_B \;&=\; i_B^{-1}(W) \qquad \textit{(par définition de $W_B$, en fait)}\\
W \;&=\; j_B^{-1}(W_B)\ ,
\end{aligned}
\right.
\tag*{(8.4.3)}
\]
\end{itemize}
\end{itemize}


\endgroup
% END UNIT Der-ChXV-05-11-10/Der-ChXV-05-11-10_p056-060.tex
% BEGIN UNIT Der-ChXV-05-11-10/Der-ChXV-05-11-10_p061-065.tex
\begingroup\makeatletter%
\newcommand{\printedpage}[1]{\marginpar{\footnotesize p.~#1}}%
\newcommand{\msp}[1]{\marginpar{\footnotesize [#1]}}%
\newcommand{\Simp}{\boldsymbol{\Delta}}%
\newcommand{\Cub}{\square}%
\makeatother\let\printedpage\derpp\let\msp\dermsp%
\let\mspin\DERmspin%


\printedpage{59}
% head: A. Grothendieck : Dérivateurs, Ch. XV, Édition des Universités de Montpellier II et Paris VII

\begin{itemize}[label={},leftmargin=4.5em]
\item \textit{et les foncteurs localisés}
\[
\begin{tikzcd}
\mathrm{Hot}_{B,W} \arrow[r, shift left, "\overline{i_B}"] & \mathrm{Hot}_W \arrow[l, shift left, "\overline{j_B}"]
\end{tikzcd}
\]
\textit{sont des équivalences de catégories. De plus, ces équivalences de catégories sont quasi-inverses l'une de l'autre, les morphismes d'adjonction étant déduits des morphismes d'adjonction pour le couple de foncteurs adjoints $(i_B, j_B)$, ces morphismes d'adjonction étant de plus dans $W$ resp. dans $W_B$.}\footnote[38]{\textbf{N.B.} La condition $(d)$, d'apparence très forte, est équivalente à la condition d'apparence plus anodine :\par
$d'$) le foncteur $j_A : \mathrm{Cat} \longrightarrow A^\wedge$ transforme éléments de $W$ en flèches $W_A$-asphériques ;\par\noindent
et il suffit même de l'exiger pour $C \longrightarrow e$ dans Cat, où $C$ est un objet de Cat ayant objet final.\par\noindent
\texttt{[La condition ($d'$) ci-dessus est trop forte et n'est pas équivalente à la condition ($d$) ; en revanche la condition plus faible envisagée ensuite est bien équivalente à ($d$)].}}
\end{itemize}

\msp{page 107}

\begin{enumerate}[leftmargin=2.5em]
\item[$2^\circ$)] \textit{Si les conditions équivalentes de $1^\circ$ sont satisfaites, alors le triple}
\[
(W_A\,,\ C = \mathrm{Mon}_{A^\wedge}\,,\ F = (C \cap W_A)_*)
\]
\textit{est un triple de Quillen (clos) propre à gauche.}
\item[$3^\circ$)] \textit{Si les conditions équivalentes de $1^\circ$ sont satisfaites, et si de plus $A$ est $W$-asphérique, alors (et alors seulement) $A$ elle-même est une} catégorie-test \textit{pour $W$, ce qui revient à dire que les conditions sur la catégorie $B$ dans} (d) \textit{sont satisfaites quand $B$ est} soit \textit{$A$, soit une catégorie localisée $A/a$.}
\end{enumerate}

\textbf{Corollaire 8.5.} \textit{Les conditions équivalentes de $1^\circ$ et de $3^\circ$ sont satisfaites dans chacun des cas suivants (et ceci quel que soit $W$, satisfaisant aux hypothèses dites).}
\begin{enumerate}[leftmargin=2.5em]
\item[a)] $A = \Simp$, $\Simp^{\sim}$, $\Cub$, $\Cub^{\sim}$ \textit{(catégories des simplexes combinatoires standard, ordonnés ou non, et des cubes combinatoires standard, ordonnés ou non \ldots)}\footnote[39]{$\Cub^{\sim}$ est essentiellement la catégorie $A$ “universelle”, stable par produits finis, et contenant un “intervalle séparant” $(I,\delta_0,\delta_1)$ comme dans $(b)$.}.
\item[b)] \textit{$A$ est stable par produits finis}\footnote[40]{En fait, il suffit que $I \longrightarrow e$ soit quarrable, \textit{i.e.} que $A$ ait un objet final, et que les produits $I \times a$ ($a \in \mathrm{Ob}\,A$) existent dans $A$.}\textit{, et contient un objet $I$ tel que $I$ ait deux sections $\delta_0$, $\delta_1$ sur l'objet final, avec $\delta_0$, $\delta_1$ disjointes,} i.e. $\mathrm{Ker}(\delta_0,\delta_1) = \emptyset_{A^\wedge}$ \textit{(sous-préfaisceau vide du préfaisceau final).}
\end{enumerate}

\msp{page 108}

Le cas $(b)$ du corollaire est justiciable de la condition $(c)$ dans $1^\circ$. (De plus, $A$ ayant un objet final, est $W$-asphérique, donc on est dans les conditions de $3^\circ$.) Les cas $\Simp$, $\Simp^{\sim}$, $\Cub$ (et $\Cub^{\sim}$, qui tombe en fait dans le cas $(b)$) sont traités dans Pursuing Stacks \S34 (p. 47-50) \texttt{[voir aussi \S36]}.

\textsc{Démonstration du théorème} 8.4. La partie essentielle, techniquement, est bien sûr $1^\circ$\footnote[41]{La partie $2^\circ$ en résulte, grâce au corollaire 8.3.}. On a les implications $\pm$ tautologiques
\[
(a) \ \Longrightarrow\ (b) \ \Longrightarrow\ (c)\ ,
\]
où $(a) \Longrightarrow (b)$ provient du fait que $L$ est évidemment un objet injectif de $A^\wedge$. (L'objet de Lawvere d'un topos quelconque est un objet injectif.) Il reste à établir
\[
(c) \ \Longrightarrow\ (d) \ \Longrightarrow\ (a)\ .
\]
% foot: Version 5 novembre 2010

\printedpage{60}
% head: A. Grothendieck : Dérivateurs, Ch. XV, Édition des Universités de Montpellier II et Paris VII

Que l'on ait $(c) \Longrightarrow (d)$ (et $(c) \Longrightarrow (b)$) est établi dans Pursuing Stacks, voir le résumé des notions $W$-test dans Pursuing Stacks, \S65, p. 168-179; l'énoncé pertinent ici est la proposition 3, page 173.

Il reste à établir $(d) \Longrightarrow (a)$ – mais à présent on joue sur du velours! Je note ici le lemme suivant, dont la place serait plutôt à la fin du \S4 des généralités sur les triples de Quillen :

\msp{page 109}

\textbf{Lemme 8.6.} \textit{Soit $\mathcal{M}$ une catégorie, $(C,TF)$ un précouple de Quillen dans $\mathcal{M}$, et $W \subset \mathrm{Fl}(\mathcal{M})$ une partie de} $\mathrm{Fl}(\mathcal{M})$ \textit{stable par facteurs directs. Conditions équivalentes :}
\begin{enumerate}[leftmargin=2.5em]
\item[a)] \textit{$TF \subset W$.}
\item[b)] \textit{Toute flèche $f$ dans $\mathcal{M}$ se factorise en $f = pi$, avec $i \in C$, $p \in W$.}
\item[b$'$)] \textit{Comme} (b), \textit{avec $f$ dans $TF$.}
\end{enumerate}

\textsc{Démonstration}. Les implications $(a) \Longrightarrow (b) \Longrightarrow (b')$ sont tautologiques et ne font pas appel à l'hypothèse de stabilité sur $W$. Mais pour un $f \in TF$ donné, pour voir qu'il est dans $W$, si $W$ est stable par facteurs directs, il suffit que cet-$f$-là se factorise en $pi$, avec $i \in C$, $p \in W$. En effet, on voit aussitôt qu'alors $f$ est facteur direct de $p$ (car $i \in C$), et comme $p \in W$, on a $i \in W$, \quad \textit{q.e.d.} % [sic] $i \in W$ -> $f \in W$, the conclusion to be proved concerns f

Revenons à 8.4, démonstration de $(d) \Longrightarrow (a)$.

\texttt{[Suit un passage rayé, en bas de la page 109 et en haut de la page 110.]}

\msp{page 110}

Finalement, mon projet de démonstration se modifie, car bien que $(d)$ ait l'air très, très fort comme condition, telle quelle je ne vois pas, finalement, comment en déduire $(a)$. Mais en fait on peut considérer que l'équivalence de $(b)$, $(c)$, $(d)$ est mise ici pour mémoire, elle est établie dans \textit{loc. cit.}, donc je n'ai à me préoccuper que de l'implication $(c) \Longrightarrow (a)$, qui va résulter
\msp{page 111}
du lemme 8.6, qui nous ramène à prouver que toute flèche $f : X \longrightarrow Y$ dans $A^\wedge$ se factorise en $pi$, avec $i$ monomorphisme et $p \in W_A$. On fait l'argument du mapping-cone \texttt{[plutôt mapping-cylinder]}, en partant de
\[
\begin{tikzcd}[column sep=small]
e_A \amalg e_A \arrow[rr, hook, "{(\delta_0,\delta_1)}", "\text{mono}"'] \arrow[dr] & & I \arrow[dl, "(W_A)_{\mathrm{univ}}"] \\
& e_A\ , &
\end{tikzcd}
\]
d'où pour tout $X$ dans $A^\wedge$, en prenant le produit par $X$,
\[
\begin{tikzcd}[column sep=small]
X \amalg X \arrow[rr, hook, "{(\delta_0^X,\delta_1^X)}", "\text{mono}"'] \arrow[dr, "{(\mathrm{id},\mathrm{id})}"] & & X \times I \arrow[dl, "\mathrm{pr}_1^X \in W_A"] \\
& X\ , &
\end{tikzcd}
\]
% foot: Version 5 novembre 2010
\printedpage{61}
% head: A. Grothendieck : Dérivateurs, Ch. XV, Édition des Universités de Montpellier II et Paris VII
et on construit le mapping-cone \texttt{[plutôt mapping-cylinder]} de $f$ par le diagramme cocartésien
\[
\begin{tikzcd}[column sep=4em,row sep=4em]
 & X \arrow[r, "f"] \arrow[d, hook, shift right, "\delta_0^X"'] \arrow[dr, phantom, "\text{cocart.}"] & Y \arrow[d, hook, shift right, "\beta"'] \\
X \arrow[r, dotted, "\delta_1^X"] \arrow[rr, dotted, bend right, "i_f"'] & X \times I \arrow[r, "f'"] \arrow[u, dotted, shift right, "\mathrm{pr}_1^X"'] & Z(f)\ , \arrow[u, dotted, shift right, "p_f"']
\end{tikzcd}
\qquad
\begin{array}{c}
p_f \text{ déduit de} \\
\mathrm{pr}_1^X : X \times I \to X
\end{array}
\tag*{(8.6.1)}
\]
la factorisation cherchée de $f$ est
\[
f \ = \ p_f\, i_f\ , \qquad \text{où}\quad i_f = f'\delta_1^X\ .
\]
Ici, comme $\mathrm{pr}_1^X \in W_A$, on en conclut que
\msp{page 112}
$\delta_0^X$, $\delta_1^X$ (qui en sont des sections) sont dans $W_A$, comme $\delta_0^X$ est un monomorphisme, on a donc $\delta_0^X \in TC$, donc ($TC$ étant stable par cochangement de base) $\beta \in W_A$, donc enfin $p_f$ (qui est une rétraction de $Z(f)$ sur $Y$) est aussi dans $W_A$. Il reste à voir que $i_f$ est un monomorphisme, ce qui se voit sur le diagramme plus complet (\textit{cf.} XII, p. 230) que j'ai sous les yeux : $i_f$ apparaît comme composé de deux monomorphismes $X \longrightarrow Y \amalg X$ et $Y \amalg X \longrightarrow Z(f)$ \ldots{} Le théorème 8.4 est démontré.

Je vais quand même isoler l'argument formel sempiternel (que j'ai utilisé surtout sous forme duale) dans un

\textbf{Lemme 8.7.} \textit{Soit $\mathcal{M}$ une catégorie, $W$ un localiseur modérément saturé dans $\mathcal{M}$, $C$ une partie stable par composition et par cochangement de base. On se propose, pour $X$ fixé dans $\mathcal{M}$, de trouver une factorisation fonctorielle de toute flèche $f : X \longrightarrow Y$} (i.e. \textit{du morphisme costructural d'un $Y$ dans $X\backslash\mathcal{M}$) en $p_f i_f$, avec $i_f \in C$, $p_f \in W$. On suppose pour ceci qu'on a}
\msp{page 113}
\textit{une factorisation de ce type pour $X \amalg X \xrightarrow{(\,\mathrm{id}\,,\mathrm{id}\,)} X$,} i.e. \textit{un objet $I_X$ et une factorisation}
\[
\begin{tikzcd}[column sep=6em]
X \amalg X \arrow[r, hook, "{i_X = (\delta_0^X,\delta_1^X)}"] \arrow[rr, dotted, bend right, "{(\mathrm{id}_X,\mathrm{id}_X)}"'] & I_X \arrow[r, "p_X"] & X\ ,
\end{tikzcd}
\]
\textit{avec $i_X \in C$, $p_X \in W$. On associe alors à toute $f$ le mapping cone} \texttt{[plutôt cylinder]} \textit{$Z(f)$ comme dans le diagramme} (8.6.1) \textit{(où il faut lire $I_X$ au lieu de $X \times I$), et on a alors $p_f \in W$, $i_f \in C$ comme voulu. Mais il faut supposer, pour savoir que $i_f \in C$, que $X \in \mathcal{M}_C$ (objets $C$-cofibrants),} i.e. \textit{que pour tout $Y$, $X \longrightarrow Y \amalg X$ est dans $C$. (De plus, on a supposé tacitement $\mathcal{M}$ stable par $\varinjlim$ finies.)}
% foot: Version 5 novembre 2010

\printedpage{62}
% head: A. Grothendieck : Dérivateurs, Ch. XV, Édition des Universités de Montpellier II et Paris VII

\textit{Le diagramme complet est}
\[
\left\{
\begin{array}{c}
\begin{tikzcd}[ampersand replacement=\&,column sep=4em,row sep=3.5em]
 \& X \arrow[r, "f"] \arrow[d, "\mathrm{inj}_0^X"'] \arrow[dr, phantom, "\text{cocart.}"] \& Y \arrow[d] \\
X \arrow[r, "\mathrm{inj}_1^X"] \arrow[dr, dotted, "\delta_1^X"'] \& X \amalg X \arrow[r, "f''"] \arrow[d, "i_X"'] \arrow[dr, phantom, "\text{cocart.}"] \& Y \amalg X \arrow[d, "j"] \\
 \& I_X \arrow[r, "f'"'] \arrow[uu, dotted, bend right=45, "p_X"{pos=.12,swap}] \& Z(f) \arrow[uu, dotted, bend right=45, "p_f"{pos=.25,swap}]
\end{tikzcd}
\\[1ex]
f = p_f i_f, \quad \textit{où} \quad i_f = f' \circ \delta_1^X = j \circ (\underbrace{f'' \circ \mathrm{inj}_1^X}_{\lambda})\ .
\end{array}
\right.
\tag*{(8.7.1)}
\]

\msp{page 114}

On a $i_X \in C$ par hypothèse, donc $j \in C$ (stabilité de $C$ par cochangement de base), on a $\lambda : X \longrightarrow Y \amalg X$ dans $C$ car $X$ dans $\mathcal{M}_C$, donc le composé $i_f = j \circ \lambda$ est dans $C$ (stabilité de $C$ par composition) – on gagne!

\textbf{Commentaires.} Pour l'équivalence de $(a)$, $(b)$, $(c)$ dans le théorème 8.4, $1^\circ$, on n'a donc pas eu à faire appel aux développements de Pursuing Stacks, qui servent seulement à établir l'équivalence des trois conditions en question avec la condition “modélisante” $(d)$. Pour résumer :
\begin{enumerate}[leftmargin=2.5em]
\item[$1^\circ$)] Les implications
\[
(a) \ \Longrightarrow\ (\,(b) \ \Longleftrightarrow\ (c) \ \Longleftrightarrow\ (d)\,)
\]
sont valables sous les seules conditions
\[
\mathrm{W}\,(1,\ 2\,\mathrm{bis},\ 3)\ ,
\]
et pour l'équivalence de $(a)$ aux autres conditions, il suffit de
\[
\mathrm{W}\,(1\,\mathrm{bis},\ 2\,\mathrm{bis},\ 3)
\]
(où W\,(1\,bis) renforce W\,(1) par la condition des facteurs directs, tandis que W\,(2\,bis) affaiblit W\,(2), en exigeant seulement que les morphismes $W$-asphériques soient dans $W$). Les conclusions de 8.4, $3^\circ$, relatives au cas où $A$ est $W$-asphérique, sont valables également sous les conditions W\,(1, 2\,bis, 3) sans plus.
\end{enumerate}

\msp{page 115}

\begin{enumerate}[leftmargin=2.5em]
\item[$2^\circ$)] Le corollaire 8.5, se référant aux conditions $(b)$, $(c)$, $(d)$ de $1^\circ$ (à l'exclusion de $(a)$, qui suppose $W$ stable par facteurs directs), est valable pour un $W$ satisfaisant
\[
\mathrm{W}\,(1,\ 2\,\mathrm{bis},\ 3,\ \underbrace{7\,\mathrm{bis}}_{\substack{\text{carré}\\ \text{cocartésien}}}\,)\ .
\]
% foot: Version 5 novembre 2010
\printedpage{63}
% head: A. Grothendieck : Dérivateurs, Ch. XV, Édition des Universités de Montpellier II et Paris VII
\item[$3^\circ$)] Pour la validité de 8.4, $2^\circ$, \textit{i.e.} pour avoir un triple de Quillen $(W_A,C,F)$, il nous faut (presque) tout l'éventail des conditions sur $W$,
\[
\mathrm{W}\,(1\,\mathrm{bis},\ 2\,\mathrm{bis},\ 3,\ \underbrace{\ 6\ }_{\text{limites}},\ \underbrace{7\,\mathrm{bis}}_{\substack{\text{carré}\\ \text{cocartésien}}},\ \underbrace{\ 8\ }_{\text{accessibilité}})\ ^{\footnotemark[42]}\ ,
\]
\footnotetext[42]{\textbf{N.B.} On devrait renuméroter les axiomes W\,(\ldots).}
à l'exclusion de W\,(4), et en se contentant pour W\,(2) de la forme affaiblie W\,(2\,bis) sans plus. (Même l'axiome de connexité W\,(5) n'a pas servi, mais je l'avais d'abord inclus, car il y a peu de chances qu'on aura jamais à travailler avec des $W$ qui n'y satisfont pas! Il se pourrait d'ailleurs, en présence des autres axiomes, ou seulement de W\,(1\,bis, 2\,bis, 3, 7\,bis), que si W\,(5) n'est \textit{pas} satisfait, on ait $W = \mathrm{Fl}\,\mathrm{Cat}$ tout entier, donc $W_A = \mathrm{Fl}\,A^\wedge$!)
\end{enumerate}

\msp{page 116}

\section*{9\quad Révision des notations, pour les axiomes d'un localiseur fondamental\footnote[43]{\textit{Cf.} XVII \texttt{[plutôt XVI]} pour une version définitive (?).}}

\begin{description}[style=sameline,leftmargin=6.5em]
\item[Loc\,(1)] (Saturation).
\begin{enumerate}[leftmargin=2.5em]
\item[$a$)] Les isomorphismes sont dans $W$.
\item[$b$)] Si $f$, $g$ sont composables, et deux parmi $f$, $g$, $gf$ sont dans $W$, le troisième aussi.
\item[$c$)] Si $X \underset{g}{\overset{f}{\rightleftarrows}} Y$, si $gf$ et $fg$ sont dans $W$, alors $f$ et $g$ aussi.
\end{enumerate}
\end{description}
\noindent Variante renforcée :
\begin{description}[style=sameline,leftmargin=6.5em]
\item[Loc\,(1\,bis)] Loc\,(1) plus
\begin{enumerate}[leftmargin=2.5em]
\item[$d$)] $W$ est stable par facteurs directs.
\end{enumerate}
\item[Loc\,(2)] (Axiome de l'objet final). Si $X$ dans Cat a un objet final, alors $X$ est $W$-asphérique, \textit{i.e.} $X \longrightarrow e$ est dans $W$.
\item[Loc\,(3)] (Axiome de localisation, forme faible). Soit $f : X \longrightarrow Y$ tel que pour tout $y \in Y$, $X/y$ soit $W$-asphérique (\textit{i.e.} le morphisme $f/y : X/y \longrightarrow Y/y$ déduit de $f$ par changement de base $Y/y \longrightarrow Y$ est dans $W$). Alors $f \in W$.
\end{description}
\noindent Variante renforcée de Loc\,(3) :
\begin{description}[style=sameline,leftmargin=6.5em]
\item[Loc\,(3\,bis)] (Axiome de localisation, forme forte). Soit $f : X \longrightarrow Y$ un morphisme dans $\mathrm{Cat}/S$, tel que $\forall\ s \in S$, $f/s : X/s \longrightarrow Y/s$ (déduit de $f$ par changement de base $S/s \longrightarrow S$) soit dans $W$. Alors $f \in W$. (\textbf{N.B.} Sauf erreur, il suffit le l'exiger si $X$, $Y$ sont Cat-fibrés scindés sur $S$, et $f$ $S$-cartésien compatible avec le scindage – en supposant Loc\,(1, 2)). % [sic] il suffit le l'exiger -> il suffit de l'exiger, "le" for "de"
\end{description}
% foot: Version 5 novembre 2010


\endgroup
% END UNIT Der-ChXV-05-11-10/Der-ChXV-05-11-10_p061-065.tex
% BEGIN UNIT Der-ChXV-05-11-10/Der-ChXV-05-11-10_p066-070.tex
\begingroup\makeatletter%
\newcommand{\printedpage}[1]{\marginpar{\footnotesize p.~#1}}%
\newcommand{\msp}[1]{\marginpar{\footnotesize [#1]}}%
\newcommand{\cercle}[1]{\tikz[baseline=(c.base)]\node[draw,circle,inner sep=0.8pt](c){#1};}%
\newcommand{\Cat}{\mathrm{Cat}}%
\newcommand{\id}{\mathrm{id}}%
\newcommand{\inj}{\mathrm{inj}}%
\newcommand{\Fib}{\mathrm{Fib}}%
\newcommand{\Fl}{\mathrm{Fl}}%
\makeatother\let\printedpage\derpp\let\msp\dermsp%
\let\mspin\DERmspin%


\printedpage{64}
% head: A. Grothendieck : Dérivateurs, Ch. XV, Édition des Universités de Montpellier II et Paris VII
% foot: Version 5 novembre 2010

\msp{page 117}

\begin{description}[labelwidth=5.5em,leftmargin=6.5em,style=sameline]
\item[Loc\,(4)] (Axiome du carré cocartésien). Soit $X$ dans Cat, $X_1$, $X_2$ deux ouverts de réunion $X$, $X_0 = X_1\cap X_2$, d'où un diagramme cocartésien dans Cat,
\[
\begin{tikzcd}
X_0 \arrow[r, "i_1"] \arrow[d, "i_2"'] & X_1 \arrow[d, "j_2"] \\
X_2 \arrow[r, "j_1"] & X\ .
\end{tikzcd}
\]
On a alors $i_1\in W \iff j_1\in W$ (et symétriquement $i_2\in W \iff j_2\in W$).

\item[Loc\,(6)]\footnote[44]{Il y aura un Loc\,(5) plus “basique” plus bas.} (Axiome des $W$-fibrations). Il y a des variantes \cercle{a}, (a$'$), \cercle{b}, (b$'$), (b$''$), (c), (c$'$), \textit{cf.} XII, p. 53-59. La forme la plus forte \cercle{a} est celle-ci :

Soit $f : X\longrightarrow Y$ dans Cat une $W$-fibration au sens ancien, \textit{i.e.} telle que pour toute $Y''\longrightarrow Y'$ dans $\Cat/Y$, avec $Y'$, $Y''$ contractiles, la flèche $X''\longrightarrow X'$ déduite par le changement de base $X\longrightarrow Y$ soit dans $W$. Alors $f$ est dans $\Fib_W$, \textit{i.e.} le foncteur de changement de base par $f$, $\Cat/Y\longrightarrow\Cat/X$, transforme flèches dans $W$ en flèches dans $W$.

\item[Loc\,(7)] (Axiome des limites inductives). $W$ est stable par limites inductives filtrantes.
\end{description}

Forme renforcée de Loc\,(7) :

\begin{description}[labelwidth=5.5em,leftmargin=6.5em,style=sameline]
\item[Loc\,(7\,bis)] (L-accessibilité). $W$ est stable par $\varinjlim$ filtrantes \textit{et} accessible, \textit{i.e.} il existe une petite partie $W_\pi$ de $\Fl(\Cat)$, telle que $W$ soit la plus petite partie de $\Fl(\Cat)$ stable par $\varinjlim$ filtrantes, contenant $W_\pi$.
\end{description}

\msp{page 118}

\begin{description}[labelwidth=5.5em,leftmargin=6.5em,style=sameline]
\item[Loc\,(8)]\footnote[45]{J'ai changé la numérotation, pour tenir compte d'un Loc\,(5) plus bas, qui est encore un axiome de stabilité \texttt{[initialement les axiomes Loc\,(6) à Loc\,(8) étaient notés Loc\,(5) à Loc\,(7)]}.} (Axiome de connexité). $f\in W\Longrightarrow \pi_0(f)$ bijectif.
\end{description}

\textbf{N.B.} Tous les axiomes, sauf le dernier (et un chouïa de l'avant dernier) sont des axiomes de stabilité – ils disent tous que $W$ est “grand” : il contient certaines flèches, et s'il contient telles flèches, il contient telles autres. Loc\,(8), et dans une certaine mesure Loc\,(7\,bis), sont des axiomes “limitatifs” sur la taille de $W$ : Loc\,(7) nous dit que $W\subset\mathbf{W}_0$ (= ensemble % [sic] Loc (7) -> Loc (8), apparently: the inclusion in W_0 is the axiome de connexité Loc (8)
des $f\in\Fl\,\Cat$ telles que $\pi_0(f)$ soit bijectif), Loc\,(7\,bis) nous dit l'existence d'un cardinal tel que $W$ soit égal à la L-enveloppe de $W^\pi$ (donc $W\subset\text{L-env}(W^\pi)$).

J'ai essayé de ranger les axiomes par ordre d'importance pratique. J'ai fini par rajouter l'axiome de connexité à la fin, car dans les développements théorique sur les localiseurs % [sic] théorique -> théoriques, agreement with "développements"
généraux, il n'intervient pratiquement jamais. Mais je rappelle (XII, p. 20, th. 2) que si $W$ satisfait Loc\,(1, 2, 3, 6), et si Loc\,(8) n'est \textit{pas} satisfait, alors $W=\Fl(\Cat)$\,; donc en présence de Loc\,(1, 2, 3, 6), Loc\,(8) \textit{équivaut} à $W\neq\Fl(\Cat)$ tout entier.

\msp{page 119}

Les axiomes qui ont servi pour 8.4, 1°, 3° et 8.5, $(b)$ sont
\[
\mathrm{Loc}\,(1\,\mathrm{bis},\ 2,\ 3)\ .
\]

\printedpage{65}
% head: A. Grothendieck : Dérivateurs, Ch. XV, Édition des Universités de Montpellier II et Paris VII
% foot: Version 5 novembre 2010

Ceux qui sont nécessaires pour 8.5, $(a)$ sont
\[
\mathrm{Loc}\,(1\,\mathrm{bis},\ 2,\ 3,\ 4)\ .
\]
Enfin, ceux nécessaires pour 8.4, 2° sont
\[
\mathrm{Loc}\,(1\,\mathrm{bis},\ 2,\ 3,\ 4,\ 7\,\mathrm{bis})\ ,
\]
et ils impliquent la \textit{saturation forte de $W$} (\textit{cf.} paragraphe suivant).

Il me faut réparer un oubli, ayant traité un peu par le mépris (!) les questions de composantes connexes, jusqu'à présent.

\textbf{Proposition 9.1.} \textit{Supposons que $W$ satisfait} Loc\,(1, 2, 3, 4) \textit{sans plus. Soit $f_i : X_i\longrightarrow Y_i$ $(i\in I)$ une famille} finie \textit{de flèches de} Cat. \textit{Pour que $f=\coprod f_i : \coprod X_i\longrightarrow\coprod Y_i$ soit dans $W$, il suffit que $\forall\ i\in I$, on ait $f_i\in W$. Si $W$ satisfait de plus à} Loc\,(7) (\textit{axiome des limites inductives}), \textit{alors cette conclusion reste valable sans condition de finitude sur $I$.}

\textsc{Démonstration}. Le cas $I$ fini se traite par récurrence sur $n=\operatorname{card}I$, on est ramené aussitôt au cas où $n=2$\footnote[46]{Le cas infini s'y ramène sous l'hypothèse Loc\,(7), par passage à la limite.}. On voit d'abord le

\msp{page 120}

\textbf{Corollaire 9.2.} \textit{Sous l'hypothèse} Loc\,(1, 4), \textit{soit $f : X\longrightarrow Y$ une flèche de} Cat, \textit{et $Z$ un objet de} Cat. \textit{Alors $f\in W$ si et seulement si $f\amalg\id_Z : X\amalg Z\longrightarrow Y\amalg Z$ est dans $W$.}

Immédiat sur le carré cocartésien, où $\inj_0^X$ est une immersion ouverte de Dwyer,
\[
\begin{tikzcd}[column sep=large, row sep=large]
X \arrow[r, "f"] \arrow[d, "\inj_0^X"'] & Y \arrow[d, "\inj_0^Y"] \\
X\amalg Z \arrow[r, "f'=f\amalg\id_Z"'] & Y\amalg Z\ ,
\end{tikzcd}
\]
qui montre que $f\in W\iff f'\in W$.

En itérant ce résultat, on trouve, si
\[
f_1\ :\ X_1\longrightarrow Y_1\ ,\qquad f_2\ :\ X_2\longrightarrow Y_2
\]
sont donnés, que
\[
f_1\ \in\ W\ ,\ f_2\ \in\ W\quad\Longrightarrow\quad f_1\amalg f_2\ \in\ W\ ,
\]
en regardant le diagramme de deux carrés cocartésiens :
\[
\begin{tikzcd}[column sep=large, row sep=large]
X_1 \arrow[r, "f"] \arrow[d] & Y_1 \arrow[d] \\ % [sic] f -> f_1, the arrow X_1 -> Y_1 is f_1 (cf. f'_1 = f_1 ∐ id_{X_2})
X_1\amalg X_2 \arrow[r, "f'_1=f_1\amalg\id_{X_2}"'] \arrow[d, "f'_2=\id_{X_1}\amalg f_2"'] & Y_1\amalg X_2 \arrow[d, "g_2=\id_{Y_1}\amalg f_2"] \\
X_1\amalg Y_2 \arrow[r, "g_1=f_1\amalg\id_{Y_2}"'] & Y_1\amalg Y_2\ .
\end{tikzcd}
\]

\printedpage{66}
% head: A. Grothendieck : Dérivateurs, Ch. XV, Édition des Universités de Montpellier II et Paris VII
% foot: Version 5 novembre 2010

Si $f_1$ et $f_2$ sont dans $W$, il en est de même de $f'_1$, $f'_2$, $g_1$, $g_2$, par 9.2, donc aussi de $f_1\amalg f_2=g_2f'_1$.

\msp{page 121}

Mais je ne vois aucun moyen de prouver la réciproque, que $f_1\amalg f_2\in W\Longrightarrow f_1,f_2\in W$. Grâce à ce qui a été prouvé, on est ramené au cas où $f_1$, $f_2$ sont des immersions ouvertes de Dwyer, mais il ne semble pas que cela nous avance d'un cheveu\,! En somme, il faut savoir que les applications canoniques $Y_i\hookrightarrow\coprod Y_i$ sont dans $\Fib_W$. Comme $Y_i$ est un sommand direct, on est ramené au cas d'une immersion canonique $Y_1\hookrightarrow Y_1\amalg Y_2=Y$ : si $f : X\longrightarrow Y$ est dans $W$, en est-il de même de $X_1=X\times_Y Y_1\longrightarrow Y_1$\,? C'est une propriété qui, bien sûr, serait conséquence de l'ancien W\,(4\,a) (XII, p. 53), ou seulement W\,(4\,a$'$), donc de l'actuel Loc\,(6\,a$'$). Il suffit même de la forme W\,(4\,c$'$) $=$ Loc\,(6\,c$'$) la plus faible de toutes, comme on voit en décomposant $X_1\longrightarrow Y$ en $X\xrightarrow{\ i\ }X'\xrightarrow{\ p\ }Y$, avec $i\in W$ (plus précisément, $i_1\ :\ X_1\longrightarrow X'_1$ et $i_2\ :\ X_2\longrightarrow X'_2$ dans $W$\,!), et $p\ \in\ \mathrm{Parf}_W$. Par Loc\,(6\,c$'$), $f\in W$, donc $p\in W$ implique que les fibres de $p$ sont $W$-asphériques, ce qui implique que $p$ est dans $W_{\mathrm{univ}}$, donc $p_1 : X'_1\longrightarrow Y_1$ est dans $W$.

\msp{page 122}

Donc notons par acquis de conscience :

\textbf{Corollaire 9.3.} \textit{Si $W$ satisfait} Loc\,(1, 2, 3, 6\,c$'$) (\textit{où} Loc\,(6\,c$'$) \textit{est la variante la plus faible de toutes de l'axiome des fibrations}), \textit{alors pour toute famille $(f_i : X_i\longrightarrow Y_i)_{i\in I}$ de flèches dans} Cat, \textit{$f=\coprod f_i\in W\Longrightarrow$ les $f_i$ sont dans $W$ $\forall\ i\in I$}\footnote[47]{En fait, je ne sais prouver l'implication en sens inverse, moyennant seulement Loc\,(1, 2, 3, 6\,c$'$).}. \textit{Et si $i : Y_0\hookrightarrow Y$ est une flèche qui est à la fois immersion ouverte et immersion fermée} (i.e. \textit{un isomorphisme avec un sommand direct de $Y$}), \textit{alors $i$ est dans} $\Fib_W$, i.e. \textit{si $f : X\longrightarrow X'$ est une flèche dans} $\Cat/Y$, \textit{alors $f\in W$ implique $(f_0 : X_0\longrightarrow Y_0)\in W$, où $X_0=X|Y_0$, $X'_0=X'|Y_0$.} % [sic] f_0 : X_0 -> Y_0 -> f_0 : X_0 -> X'_0, apparently: f is a flèche X -> X' and X'_0 is defined but otherwise unused

Mais il est assez scandaleux d'avoir à utiliser un axiome aussi fort, pour en conclure un résultat aussi anodin\,! Je note ceci :

\textbf{Proposition 9.4.} \textit{Soit $W\subset\Fl(\Cat)$, satisfaisant} Loc\,(1) \textit{sans plus. Les conditions suivantes sont équivalentes} (\textit{et impliquées par} Loc\,(1, 2, 3, 4, 6\,c$'$)).
\begin{enumerate}[label={}, leftmargin=2.5em]
\item[a)] \textit{Si $f_1 : X_1\longrightarrow Y_1$, $f_2 : X_2\longrightarrow Y_2$ sont deux flèches de} Cat, \textit{alors}
\[
f_1\amalg f_2\in W\iff f_1\in W,\ f_2\in W\,.
\]
\item[b)] \textit{Les deux conditions suivantes sur $W$ sont satisfaites :}

\msp{page 123}

\begin{enumerate}[label={}, leftmargin=2.5em]
\item[1°)] \textit{Si $f\ :\ X\longrightarrow Y$ est dans $W$, alors il l'est coüniversellement sous $\emptyset$,} i.e. \textit{pour tout $Z$ dans} Cat, \textit{$f\amalg\id_Z : X\amalg Z\longrightarrow Y\amalg Z$ est dans $W$.} (\textbf{N.B.} \textit{Ceci équivaut à : toute somme finie de flèches $f_i\in W$ est dans $W$ – et itou même pour sommes infinies, si} Loc\,(7), \textit{axiome des limites, est satisfait.})
\item[2°)] \textit{Soit $Y_0\hookrightarrow Y$ une inclusion d'un sommand direct d'un objet de} Cat. \textit{Alors si $X\xrightarrow{\ f\ }Y$ est dans $W$, de même $X_0=X|Y_0\longrightarrow Y_0$.} (\textbf{N.B.} \textit{Ceci équivaut à la condition que l'inclusion d'un sommand direct est toujours une $W$-fibration.})
\end{enumerate}

\textit{De plus, si $\pi_0(Y)$ est fini, ou si $W$ satisfait l'axiome des limites} Loc\,(7), \textit{il suffit d'exiger 2° pour $Y_0$ une composante connexe de $Y$.}

\printedpage{67}
% head: A. Grothendieck : Dérivateurs, Ch. XV, Édition des Universités de Montpellier II et Paris VII
% foot: Version 5 novembre 2010

\item[c)] \textit{Pour toute famille} finie \textit{$(f_i)_{i\in I}$ de flèches $f_i : X_i\longrightarrow Y_i$ dans} Cat,
\[
\coprod f_i\in W\iff f_i\ \in\ W\ \ \textit{pour tout}\ i\in I.
\]
\end{enumerate}
\textit{De plus, si $W$ satisfait} Loc\,(7), \textit{alors il est inutile dans} (c) \textit{de se borner à $I$ fini.} (\textit{Et en tous cas, les conditions équivalentes} (a), (b), (c) \textit{impliquent pour} toute \textit{famille de flèches $(f_i)_{i\in I}$, $\coprod f_i\in W\Longrightarrow f_i\in W$ $\forall\ i\in I$.})

\msp{page 124}

D'après 9.3, ces conditions sont satisfaites moyennant Loc\,(1, 2, 3, 4, 6\,c$'$), mais Loc\,(6\,c$'$) est un marteau-pilon pour une mouche\,! On aurait envie plutôt d'ajouter un axiome aux axiomes Loc\,(1) à Loc\,(4) et Loc\,(6) à Loc\,(8), qui avec Loc\,(8) serait l'axiome le plus anodin de tous.

\begin{description}[labelwidth=5.5em,leftmargin=6.5em,style=sameline]
\item[Loc\,(5)] (Axiome de décomposition, forme faible). Si $(f_i)_{i\in I}$ est une famille finie de flèches de Cat, alors $f=\coprod f_i$ est dans $W$ si et seulement si les $f_i$ le sont. (Ce sont les conditions équivalentes de 9.4.)\footnote[48]{Il peut être commode aussi d'introduire la condition encore plus faible :
\begin{description}[nosep,labelwidth=4.5em,leftmargin=5em,style=sameline]
\item[Loc\,(5\,a)] Si $f_1$, $f_2$ sont dans $W$, alors $f_1\amalg f_2\in W$.\\
C'est aussi la condition $(b)$, 1° de la proposition 9.4.
\item[Loc\,(5\,b)] Si $f_1\amalg f_2$ est dans $W$, alors $f_1$ et $f_2$ aussi.\\
$\iff$ Si $Y_0\subset Y$ est un sommand direct, alors $Y_0\hookrightarrow Y$ est dans $\Fib_W$.
\end{description}}

\item[Loc\,(5\,bis)] (Axiome de décomposition, forme forte). Comme Loc\,(5), mais sans supposer $I$ finie.\footnote[49]{\texttt{[Réécrivons :]}
\begin{description}[nosep,labelwidth=4.5em,leftmargin=5em,style=sameline]
\item[Loc\,(5)] \textit{a}) Si $f_1,f_2\in W$, alors $f_1\amalg f_2\in W$.\\
\textit{b}) Si $f_1\amalg f_2\in W$, alors $f_1,f_2\in W$\\
($\iff b'$) tout $Y_0\hookrightarrow Y$ sommand direct est dans $\Fib_W$\,).
\item[Loc\,(5\,bis)] Comme Loc\,(5), mais $(a)$ remplacé par ($a$\,bis) une famille (pas nécessairement finie).
\end{description}}
\end{description}

Ainsi, moyennant l'axiome des limites Loc\,(7), les deux axiomes précédents sont équivalents. De plus, notons ceci :

\textbf{Corollaire 9.5} (de la proposition 9.4). \textit{L'axiome} Loc\,(5\,bis) \textit{équivaut à chacune des propriétés suivantes :}
\begin{enumerate}[label={}, leftmargin=3em]
\item[\cercle{a}] \textit{Pour tout $Y$ de} Cat, \textit{décomposé en ses composantes connexes $Y_i$ $(i\in I=\pi_0(Y))$, et pour toute flèche $f\ :\ X\longrightarrow Y$ dans} Cat, \textit{on a $f\ \in\ W$ si et seulement si les $f_i : X_i=X|Y_i\longrightarrow Y_i$ sont dans $W$.}
\item[\cercle{b}] \textit{Pour $Y$, $Y_i$ comme dans} \cercle{a}, \textit{pour toute}

\msp{page 125}

\textit{flèche $f\ :\ X\longrightarrow X'$ dans $\Cat/Y$, on a $f\ \in\ W$ si et seulement si $\forall\ i\ \in\ I$, le morphisme $f_i : X_i\longrightarrow X'_i$ au-dessus de $Y_i$ est dans $W$.}
\end{enumerate}

\textbf{Corollaire 9.6.} \textit{Si l'axiome} Loc\,(5) \textit{est satisfait, alors l'axiome de connexité} Loc\,(8) \textit{équivaut à ceci : pour toute flèche $f\ :\ X\longrightarrow Y$ dans $W$, avec $Y$ $0$-connexe, $X$ est $0$-connexe.}

\printedpage{68}
% head: A. Grothendieck : Dérivateurs, Ch. XV, Édition des Universités de Montpellier II et Paris VII
% foot: Version 5 novembre 2010

N.B. Les axiomes Loc\,(8) (connexité) et Loc\,(5\,bis) (de décomposition) ont l'air de nature voisine, étant liés l'un et l'autre à la décomposition des objets de Cat en composantes connexes. (Néanmoins, je pense qu'ils sont indépendants.) Je présume qu'on n'aura jamais à travailler vraiment avec des $W$ qui ne satisfont à ces deux axiomes\,! (En plus de Loc\,(1, 2, 3), qui eux constituent le “minimum vital”.) Mais apparemment il n'y en a pas besoin dans les développements théoriques généraux jusqu'ici. (Mais l'axiome de connexité Loc\,(8) pourrait jouer un rôle, quand on parle de la détermination réciproque des structures de contractibilité, et de $W$-asphéricité, \textit{cf.} Pursuing Stacks. % [sic] "(Mais l'axiome ... Pursuing Stacks." -> closing parenthesis missing, the parenthesis opened at "(Mais" is never closed

\msp{page 126}

Notons ceci :

\textbf{Proposition 9.7.}\footnote[50]{\textit{Cf.} version plus forte dans 9.9 plus bas (débarrassée de l'hypothèse Loc\,(5).)} \textit{Supposons que $W$ satisfait} Loc\,(1, 2, 3, 5), i.e. \textit{c'est un localiseur fondamental satisfaisant la condition de décomposition} (\textit{forme faible}). \textit{Si l'axiome} Loc\,(8) \textit{de connexité n'est} pas \textit{satisfait, alors la sous-catégorie} $\mathrm{Hot}^*_W$ \textit{des $X\ \in$} Cat \textit{non vides est équivalente à la catégorie ponctuelle} (\textit{donc} $\mathrm{Hot}_W$ \textit{est équivalente soit à la catégorie ponctuelle, si $\emptyset\longrightarrow e$ est un isomorphisme dans} $\mathrm{Hot}_W$, \textit{soit à $\boldsymbol{\Delta}_1$ dans le cas contraire}\footnote[51]{Le premier cas est aussi celui où $\exists\ f\in W$ avec $\pi_0(f)$ \textit{non surjectif,} ou encore $\emptyset\longrightarrow X$ dans $W$, avec $X$ non vide.}.)

Ceci est une variante du th. 2 de XII (p. 20), où on faisait l'hypothèse Loc\,(6\,c$'$), plus forte que Loc\,(5) et même que Loc\,(5\,bis), mais avec la conclusion plus forte que $W=\Fl(\Cat)$.

\textsc{Démonstration}. Soit
\[
f\ :\ X\longrightarrow Y\ ,\qquad f\ \in\ W\ ,
\]
telle que $\pi_0(f)$ ne soit pas \textit{surjective.} Si $Y_0$ est une composante connexe qui n'est pas dans $\mathrm{Im}\pi_0(f)$, alors $X_0=X|Y_0$ est vide, et par Loc\,(7), $\emptyset\longrightarrow Y_0$ est dans $W$. Donc $\forall\ Z$ dans Cat, % [sic] Loc (7) -> Loc (5), apparently: the step X_0 -> Y_0 in W uses the decomposition hypothesis Loc (5) of 9.7, not Loc (7)
$\underbrace{\emptyset\times Z}_{=\,\emptyset}\longrightarrow Y_0\times Z$ est dans $W$, et pour toute flèche $g : Z\longrightarrow Z'$, $\id_{Y_0}\times g : Y_0\times Z\longrightarrow Y_0\times Z'$ est dans $W$. Prenant $g$ telle que $\pi_0(g)$ ne soit pas injective, p. ex. $e\amalg e\longrightarrow e$, on trouve une flèche $\varphi=\id_{Y_0}\times g$ dans $W$ telle que $\pi_0(\varphi)$ ne soit pas injective.

\msp{page 127}

Soit donc $\varphi\ :\ I\longrightarrow J$ dans $W$, avec $\pi_0(\varphi)$ non injective. Procédant comme tantôt en utilisant Loc\,(5), on est ramené au cas où de plus $J$ est $0$-connexe. (\textbf{N.B.} Dans les deux cas, on aurait pu aussi invoquer 9.6.) Choisissons deux éléments $i_0$, $i_1$ dans $I$ dans deux composantes connexes distinctes, soient $j_0$, $j_1$ leurs images dans $J$. On va traiter $(I,i_0,i_1)$ et $(J,j_0,j_1)$ comme des “intervalles homotopiques”. Le fait que $i_0$, $i_1$ appartiennent à des composantes distinctes de $I$ implique que deux flèches $f_0,f_1 : X\rightrightarrows Y$ sont “homotopes par rapport à $(I,i_0,i_1)$”, \textit{i.e.} il existe $h\ :\ X\times I\longrightarrow Y$ dans Cat telle que $f_0\ =\ h\partial_0^X$, $f_1=h\partial_1^X$, où
\[
\partial_0^X\ ,\ \partial_1^X\ :\ X\ \longrightarrow\ I\times X
\]
sont données par $x\longmapsto(i_\alpha,x)$ $(\alpha\in\{0,1\})$. Pour trouver $h$, on décompose $I$ en
\[
I\ =\ I_0\amalg I_1\ ,\qquad\quad \text{avec}\quad i_0\in I_0,\ i_1\in I_1\ ,
\]
d'où $I\times X\ \simeq\ (I_0\times X)\amalg(I_1\times X)$, et on définit $h$ par ses restrictions $h_0$, $h_1$ aux deux morceaux $I_0\times X$, $I_1\times X$, savoir
\[
h_0\ =\ f_0\circ\mathrm{pr}_2^{I_0}\ ,\qquad\quad h_1\ =\ f_1\circ\mathrm{pr}_2^{I_1}\ ,
\]


\endgroup
% END UNIT Der-ChXV-05-11-10/Der-ChXV-05-11-10_p066-070.tex
% BEGIN UNIT Der-ChXV-05-11-10/Der-ChXV-05-11-10_p071-075.tex
\begingroup\makeatletter%
\newcommand{\printedpage}[1]{\marginpar{\footnotesize p.~#1}}%
\newcommand{\msp}[1]{\marginpar{\footnotesize [#1]}}%
\newcommand{\marginbar}{\marginpar{$\Vert$}}%
\newcommand{\Hot}{\mathrm{Hot}}%
\newcommand{\Hom}{\mathrm{Hom}}%
\newcommand{\Cat}{\mathrm{Cat}}%
\newcommand{\hot}{\mathrm{hot}}%
\newcommand{\Ob}{\mathrm{Ob}}%
\newcommand{\card}{\operatorname{card}}%
\newcommand{\pr}{\mathrm{pr}}%
\newcommand{\id}{\mathrm{id}}%
\makeatletter%
\renewcommand\@makefnmark{\hbox{\@textsuperscript{\normalfont(\@thefnmark)}}}%
\renewcommand\@makefntext[1]{\parindent 1em\noindent\hbox{\@textsuperscript{\normalfont\@thefnmark}}\,#1}%
\makeatother%
\makeatother\let\printedpage\derpp\let\msp\dermsp%
\let\mspin\DERmspin%


\printedpage{69}
% head: A. Grothendieck : Dérivateurs, Ch. XV, Édition des Universités de Montpellier II et Paris VII
% foot: Version 5 novembre 2010
où
\[
\pr_2^{I_\alpha}\ :\ I_\alpha\times X \longrightarrow X\ .
\]

\msp{page 128}Considérons maintenant le diagramme (traits pleins)
\[
\begin{tikzcd}[column sep=small]
& X\amalg X \arrow[dl, "{(\partial_0^X,\partial_1^X)}"'] \arrow[dr, "{(\partial_0^{\prime X},\partial_1^{\prime X})}"] & \\
I\times X \arrow[rr, "\varphi\times\id_X\,=\,\varphi_X"] \arrow[dr, "h"'] & & J\times X \arrow[dl, dotted, "h'"] \\
& Y\ , &
\end{tikzcd}
\]
où $\partial_0^{\prime X}$, $\partial_1^{\prime X}$ sont définis comme $\partial_0^X$, $\partial_1^X$, mais utilisant $j_0,j_1\in\Ob J$ au lieu de $i_0,i_1\in\Ob I$. S'il existait dans Cat une flèche $h'$ (pointillée) telle que $h'\circ\varphi_X=h$, il s'ensuivrait que $h'$ est une homotopie de $f_0$ et $f_1$ suivant $(J,j_0,j_1)$. Comme $J$ est 0-connexe, il s'ensuivrait aussitôt que $f_0$, $f_1$ sont homotopes au sens habituel dans Cat, donc que $f_0$ et $f_1$ ont même image dans $\Hot_W$. Mais si on ne peut trouver une flèche $h'$ dans Cat, du moins peut-on la trouver dans $\Hot_W$, puisque $\varphi_X\in W$, donc un isomorphisme dans $\Hot_W$ (et $h'$ dans $\Hom_W(J\times X,Y)$ sera même uniquement déterminé). Donc dans $\Hot_W$, on a
\[
f_0\ =\ h'\partial_0^{\prime X}\ ,\qquad\quad f_1\ =\ h'\partial_1^{\prime X}\ .
\]

\msp{page 129}Pour prouver que
\[
f_0\ =\ f_1\qquad\text{dans}\quad\Hom_W(X,Y)\ ,
\]
il suffit de voir que
\[
\partial_0^{\prime X}\ =\ \partial_1^{\prime X}\qquad\text{dans}\quad\Hom_W(X,J\times X)\ .
\]
Or en fait $\partial_0^{\prime X}$ et $\partial_1^{\prime X}$ sont homotopes, car $e\underset{j_1}{\overset{j_0}{\rightrightarrows}}J$ sont homotopes ($j_0$ et $j_1$ appartenant à la même composante connexe de $J$), donc aussi leur produit par $\id_X$.

On voit donc que pour deux objets quelconques $X$, $Y$ de Cat, toutes les flèches de $X$ dans $Y$ sont égales dans Hot. L'argument fait prouve même que toutes les flèches \textit{dans} $\Hot_W$ de $X$ dans $Y$ sont égales. Cela signifie que la catégorie $\Hot_W^*$ (où pour deux objets quelconques $X$, $Y$, $\Hom_{\Cat}(X,Y)$, et a fortiori $\Hom_W(X,Y)$ est $\neq\emptyset$) est équivalente à la catégorie ponctuelle. Quant à $\Hot_W$ elle-même, dont l'ensemble des objets est $\{\emptyset\}\amalg\Ob\Hot_W^*$, avec $\emptyset$ objet initial, elle sera donc équivalente à la catégorie ponctuelle, ou bien à $\boldsymbol{\Delta}_1$, selon que $\emptyset\longrightarrow e$ est un isomorphisme dans $\Hot_W$, ou non.

\msp{page 130}\textbf{Commentaire 9.8.} Si les axiomes Loc\,(1, 2, 3, 5)~\footnote[52]{En fait, ce commentaire vaut même en l'absence de l'hypothèse Loc\,(5), \textit{cf.} 9.9 ci-dessous.} sont satisfaits, alors l'axiome de connexité Loc\,(8) \textit{équivaut} à chacun des suivants :

\printedpage{70}
% head: A. Grothendieck : Dérivateurs, Ch. XV, Édition des Universités de Montpellier II et Paris VII
% foot: Version 5 novembre 2010
\begin{enumerate}[leftmargin=*]
\item[\textit{a})] $\Hot_W$ n'est équivalente ni à $e$, ni à $\boldsymbol{\Delta}_1$.
\item[\textit{b})] $\Hot_W$ a une infinité de classes d'isomorphie.
\item[\textit{c})] $f_0:X_0\longrightarrow Y_0$ (flèche donnée d'avance dans Cat, telle que $\pi_0(f_0)$ n'est pas un isomorphisme, et $X_0\neq\emptyset$) n'est \textit{pas} un isomorphisme de $\Hot_W$.
\item[\textit{d})] $X_0$ et $Y_0$ ne sont pas isomorphes dans $\Hot_W$ (où $X_0$, $Y_0$ sont deux objets donnés d'avance dans Cat, avec $X_0$, $Y_0$ non vides et $\card\pi_0(X_0)\neq\card\pi_0(Y_0)$).
\end{enumerate}

% marginal mark: a double vertical rule in the left margin beside the first two lines of this paragraph
\marginbar On peut donc dire que l'axiome de connexité joue le rôle d'un \textit{axiome de non trivialité} pour la catégorie $\Hot_W$ localisée de Cat. Il sera satisfait automatiquement pour tous les localiseurs fondamentaux $W$ qu'on sera amené à construire ou à introduire. Mais on n'aura pratiquement jamais à l'invoquer comme une hypothèse, en vue de certaines conclusions.

Je m'aperçois finalement que dans la proposition 9.7 l'hypothèse Loc\,(5) est inutile :

\textbf{Proposition 9.9.}~\footnote[53]{\textit{Cf.} meilleur énoncé 9.12 (p. 133 bis).} \textit{Soit $W$ un localiseur fondamental dans} Cat (i.e. \textit{on a} Loc\,(1, 2, 3) \textit{sans plus}), \textit{plus la condition que la} somme \textit{de $f_1,f_2\in W$ est dans $W$}~\footnote[54]{Cette condition Loc\,(5\,a) est vérifiée notamment si $W$ satisfait de plus l'un des deux axiomes Loc\,(4) (\textit{cf.} 9.1), ou Loc\,(6\,c$'$) (\textit{cf.} 9.3).}. \textit{Si l'axiome de connexité n'est pas satisfait, alors} $\Hot_W$ \textit{est équivalente à la catégorie ponctuelle} (\textit{si $\emptyset\longrightarrow e$ est un isomorphisme dans} $\Hot_W$~\footnote[55]{Cela signifie aussi qu'il existe $Z\neq\emptyset$ dans Cat, tel que $\emptyset\longrightarrow Z$ soit un isomorphisme dans $\Hot_W$.}) \textit{ou à $\boldsymbol{\Delta}_1$} (\textit{dans le cas contraire}), \textit{et} $\Hot_W^*$ ($\overset{\text{déf}}{=}\Hot_W\backslash\{\emptyset\}$) \textit{est équivalente à la catégorie ponctuelle.}

\msp{page 131}La condition sur les sommes assure que dans $\Hot_W$ la somme de deux objets existe, et que $\Cat\longrightarrow\Hot_W$ commute aux sommes de deux objets. Ceci dit, soit
\[
f\ :\ X\longrightarrow Y \tag*{$(*)$} % the label (*) is printed at the left margin
\]
une flèche telle que $f\in W$, mais $\pi_0(f)$ non surjectif. Supposons d'abord $X\neq\emptyset$ (cas le moins évident). Si $Y_0$ est un sommand non vide de $Y$ tel que $f$ se factorise par $Y_1=Y\backslash Y_0$, alors $f$ est somme de
\[
f_1\ :\ X\ =\ X_1\longrightarrow Y_1\qquad\quad\text{et}\qquad\quad f_0\ :\ X_0\ =\ \emptyset\longrightarrow Y_0
\]
dans Cat, donc aussi dans $\Hot_W$. Donc on a une somme $f\ =\ f_0\amalg f_1$ dans $\Hot_W$ qui est un isomorphisme, donc pour tout $Z$ dans $\Hot_W$,
\[
\Hom_W(Y_0,Z)\times\Hom_W(Y_1,Z)\ \xrightarrow{\ \sim\ }\ \underbrace{\Hom_W(X_0,Z)}_{=\,e}\times\Hom_W(X_1,Z)
\]
(\textbf{N.B.} $\emptyset$ est objet initial de $\Hot_W$, l'étant de Cat), \textit{i.e.} le composé
\[
Z(Y_0)\times Z(Y_1)\ \xrightarrow{\pr_2}\ Z(Y_1)\ \xrightarrow{Z(f_1)}\ Z(X_1)
\]
est bijectif. Comme $X\ =\ X_1\neq\emptyset$, donc aussi $Y_1\neq\emptyset$, et $Y_0\neq\emptyset$, les ensembles $Z(Y_0)$, $Z(Y_1)$, $Z(X_1)$ sont $\neq\emptyset$ si $Z\neq\emptyset_{\Cat}$. Le composé étant bijectif, la première flèche $\pr_2$ est injective, ce qui implique
\printedpage{71}
% head: A. Grothendieck : Dérivateurs, Ch. XV, Édition des Universités de Montpellier II et Paris VII
% foot: Version 5 novembre 2010
\msp{page 132}que $Z(Y_0)$ est réduit à un point (les deux facteurs $Z(Y_0)$, $Z(Y_1)$ étant $\neq\emptyset$). Cela signifie que $Y_0$ est un objet initial dans $\Hot_W^*$. Il s'ensuit notamment que $Y_0\amalg Y_0$ est aussi un objet initial, et l'application canonique $Y_0\amalg Y_0\longrightarrow Y_0$ est donc un isomorphisme dans $\Hot_W$, alors que l'application induite sur les $\pi_0$ n'est pas injective. (C'est $\pi_0(Y_0)\amalg\pi_0(Y_0)\longrightarrow\ \pi_0(Y_0)$, avec $\pi_0(Y_0)$ non vide.) Ce cas va être traité ci-dessous.

Si dans $(*)$ $X$ est vide, $Y$ non vide (puisque $\pi_0(f)$ non surjectif), alors $Y$ est objet initial de $\Hot_W$, et on voit encore que $Y\amalg Y\longrightarrow Y$ est un isomorphisme dans $\Hot_W$, bien que n'induisant pas un isomorphisme des $\pi_0$. La proposition 9.9 va donc résulter du

\textbf{Lemme 9.10.} \textit{Soit $W$ localiseur fondamental sans plus} (\textit{axiomes} Loc\,(1, 2, 3)). \textit{Supposons qu'il existe $f:I\longrightarrow J$ dans} Cat \textit{tel que} $\hot_W(f)$ \textit{soit un isomorphisme, mais que $\pi_0(f)$} ne soit pas injectif~\footnote[56]{Ou bien tel que $I$ soit 0-connexe, et $J$ non 0-connexe (\textit{i.e.} $\card\pi_0(I)\ =\ 1$, $\card\pi_0(J)\ \geq\ 2$). Ce deuxième cas se traite par la même méthode essentiellement.}. \textit{Alors} $\Hot_W^*$ \textit{est équivalente à la catégorie ponctuelle $e$, et} $\Hot_W$ \textit{est équivalente à $e$} (\textit{si $\emptyset\longrightarrow e$ est un isomorphisme dans} $\Hot_W$) \textit{ou à $\boldsymbol{\Delta}_1$} (\textit{dans le cas contraire}).

\msp{page 133}La démonstration est celle donnée pour 9.7 : on prouve que si $X$, $Y$ sont deux objets de Cat, \textit{non vides}, alors deux flèches $f_0,f_1:X\rightrightarrows Y$ sont égales dans $\Hot_W$, par un argument d'homotopie, en choisissant deux points $i_0$, $i_1$ de $I$ qui sont dans deux composantes connexes distinctes, et tels que leurs images $j_0$, $j_1$ dans $J$ soient dans la même composante.

\textbf{Proposition 9.11.} \textit{Soit $W$ un localiseur fondamental} (i.e. \textit{on a} Loc\,(1, 2, 3)). \textit{Alors les conditions suivantes sont équivalentes.}
\begin{enumerate}[leftmargin=*]
\item[a)] \textit{On a} Loc\,(5\,bis) (\textit{axiome de décomposition, forme forte}) \textit{et} Loc\,(8) (\textit{axiome de connexité}).
\item[b)] \textit{On a} Loc\,(5\,bis), \textit{et} $\Hot_W$ \textit{n'est pas équivalente à $e$ ou à $\boldsymbol{\Delta}_1$} (cf. \textit{commentaire} 9.8).
\item[c)] \textit{Pour toute $f:X\longrightarrow Y$, on a $f\in W$ si et seulement si $\pi_0(f)$ est un isomorphisme, et pour toute composante connexe $X_i$ de $X$, si $Y_i$ désigne la composante connexe correspondante de $Y$, la flèche induite $f_i:X_i\longrightarrow Y_i$ est dans $W$.}
\end{enumerate}

La raison pour laquelle je n'isole pas cette propriété en un des axiomes Loc\,(\ldots), c'est que ce n'est \textit{pas} un axiome de stabilité, mais un axiome de nature mixte.

\msp{page 133 bis}\textbf{Proposition 9.12.} \textit{Soit $W$ un localiseur fondamental} (\textit{sans plus}), i.e. \textit{satisfaisant} Loc\,(1, 2, 3). \textit{Les conditions suivantes sont équivalentes.}
\begin{enumerate}[leftmargin=*]
\item[a)] \textit{$W$ satisfait l'axiome de connexité} Loc\,(8), i.e. $f\in W\Longrightarrow\pi_0(f)$ \textit{bijectif.}
\item[b)] \textit{Si $f,g:X\rightrightarrows Y$ dans} Cat \textit{sont telles que} $\hot_W(f)=\hot_W(g)$, \textit{alors} $\pi_0(f)=\pi_0(g)$.
\item[c)] \textit{Il existe $X$, $Y$ dans} Cat \textit{tel que l'application} $\hot_W^{X,Y}:\Hom_{\Cat}(X,Y)\longrightarrow\Hom_W(X,Y)$ \textit{ne soit pas constante.}
\item[c$'$)] \textit{La sous-catégorie} $\Hot_W^*$ \textit{de} $\Hot_W$, \textit{formée des objets non vides, n'est pas équivalente à la catégorie ponctuelle.}
\item[c$''$)] \textit{La catégorie} $\Hot_W$ \textit{n'est pas équivalente à la catégorie $e$ ni à la catégorie $\boldsymbol{\Delta}_1$.}
\item[d)] \textit{L'ensemble des classes d'isomorphie d'objets de} $\Hot_W$ \textit{est infini.}
\end{enumerate}

\printedpage{72}
% head: A. Grothendieck : Dérivateurs, Ch. XV, Édition des Universités de Montpellier II et Paris VII
% foot: Version 5 novembre 2010
\textsc{Démonstration}. On a $(a)\Longrightarrow(b)$, puisque $(a)$ équivaut à dire que le foncteur $\pi_0:\Cat\longrightarrow\mathrm{Ens}$ se factorise par $\Hot_W$. D'ailleurs le foncteur correspondant
\[
\Hot_W\ \longrightarrow\ (\mathrm{Ens})
\]
est essentiellement surjectif -- il y a un inverse à isomorphisme près, $I\longmapsto\hot_W(I)$. Cela prouve que $(a)$ implique $(d)$, et on a $(d)\Longrightarrow(c')$, $(d)\Longrightarrow(c'')$ tautologiquement. Un argument plus ou moins tautologique montre d'ailleurs que si $\Hot_W^*$ est équivalente à $e$, alors $\Hot_W$ est équivalente à $e$ ou à $\boldsymbol{\Delta}_1$ (car $\emptyset$ est objet initial de $\Hot_W$) -- donc $(c'')\Longrightarrow(c')$. Donc on a plus ou moins tautologiquement
\[
\begin{tikzcd}
(\mathrm{a}) \arrow[r, Rightarrow] \arrow[d, Rightarrow] & (\mathrm{d}) \arrow[r, Rightarrow] & (\mathrm{c}'') \arrow[r, Rightarrow] & (\mathrm{c}') \\
(\mathrm{b}) \arrow[rr, Rightarrow] & & (\mathrm{c}) \arrow[ru, Rightarrow] &
\end{tikzcd}
\]
où les implications $(b)\Longrightarrow(c)\Longrightarrow(c')$ sont elles aussi triviales \texttt{[c'est plutôt l'implication (c$'$) $\Longrightarrow$ (c) qui est triviale~; l'implication (c) $\Longrightarrow$ (c$'$) est démontrée après l'énoncé du corollaire 9.13, avant la preuve de ce dernier]}. Donc
\msp{page 133 ter}tout revient à prouver que $(c)\Longrightarrow(a)$, ou encore :

\textbf{Corollaire 9.13.} \textit{Si l'axiome de connexité n'est pas satisfait, alors si $X$, $Y$ sont deux objets de} Cat, \textit{toutes les flèches $X\longrightarrow Y$ dans} Cat \textit{sont égales dans} $\Hot_W$, i.e. \textit{l'application} $\Hom_{\Cat}(X,Y)\longrightarrow\Hom_W(X,Y)$ \textit{est constante.}

Cela implique donc que tous les objets de Cat$^*$, \textit{i.e.} tous les objets non vides de Cat, sont isomorphes dans $\Hot_W$, et on trouve un isomorphisme canonique entre elles, avec transitivité. Soit alors $\mathcal{H}$ la catégorie quotient de Cat, ayant mêmes objets que Cat, et ayant exactement une flèche entre deux objets non vides de Cat et exactement une flèche de $\emptyset$ dans tout autre objet de Cat (et pas de flèche d'un objet non vide de Cat dans $\emptyset$). C'est donc une catégorie équivalente à $\boldsymbol{\Delta}_1$. On a vu que $\Cat\longrightarrow\Hot_W$ se factorise en
\[
\Cat\ \xrightarrow{\text{can.}}\ \mathcal{H}\ \longrightarrow\ \Hot_W\ .
\]
Si $\Cat\longrightarrow\mathcal{H}$ transforme toute flèche appartenant à $W$ en isomorphisme, alors il est immédiat que $\mathcal{H}$ est solution au problème universel définissant $\Hot_W$, donc que $\mathcal{H}\longrightarrow\Hot_W$ est un isomorphisme. Cela signifie aussi que $W$ ne contient pas de flèche $\emptyset\longrightarrow X$ avec $X$ non vide (auquel cas $W$ contient toutes les flèches). Donc dans ce cas $\Hot_W\simeq\boldsymbol{\Delta}_1$. Dans le cas contraire, soit $\mathcal{H}'$ la catégorie déduite de $\mathcal{H}$ en inversant les flèches $\emptyset\longrightarrow X$, \textit{i.e.} le groupoïde 1-connexe sur l'ensemble d'objets $\Ob\Cat$. Ce n'est plus une catégorie quotient de Cat, mais on a encore une factorisation
\[
\Cat\ \xrightarrow{\text{épi}}\ \mathcal{H}\ \ \xrightarrow{\text{loc.}}\ \mathcal{H}'\ \longrightarrow\ \Hot_W\ ,
\]
\printedpage{73}
% head: A. Grothendieck : Dérivateurs, Ch. XV, Édition des Universités de Montpellier II et Paris VII
% foot: Version 5 novembre 2010
\msp{page 133 quater}et cette fois c'est $\mathcal{H}'$ qui est solution du problème universel ($\Cat\longrightarrow\mathcal{H}'$ transformant flèches de $W$ en isomorphismes). Donc dans ce cas, $\mathcal{H}'\longrightarrow\Hot_W$ est un isomorphisme, et $\Hot_W$ est équivalente à $e$. Dans les deux cas, $\Hot_W^*$ est équivalente à $e$.

Donc, la proposition 9.12 résulte du corollaire 9.13, qu'il s'agit à présent de démontrer.

Soit $f:I\longrightarrow J$ dans Cat telle que $\pi_0(f)$ ne soit pas injective. On a vu dans la démonstration de 9.7 comment on en déduit (en choisissant deux éléments $i_0$, $i_1$ dans deux composantes connexes distinctes de $I$, telles que $j_0=f(i_0)$ et $j_1=f(i_1)$ soient dans une même composante connexe de $J$) que pour $X$, $Y$ dans Cat, deux flèches $f_0,f_1:X\rightrightarrows Y$ sont toujours égales dans $\Hot_W$ (\textit{cf.} lemme 9.10).

Soit maintenant
\[
f\ :\ I\longrightarrow J
\]
dans Cat telle que $\pi_0(f)$ ne soit pas surjective, donc $J=J_0\amalg J_1$, avec $f$ se factorisant par $J_0$, et $J_1\neq\emptyset$. Considérons alors la catégorie discrète à deux objets,
\[
\varepsilon\ =\ \{0,1\}\ ,
\]
et deux flèches
\[
g_0,\ g_1\ :\ J\ \rightrightarrows\ \varepsilon\ ,
\]
$g_0$ étant constante de valeur 0, et $g_1$ étant égale à 0 sur $J_0$, à 1 sur $J_1$. On a
\[
g_0f\ =\ g_1f\ ,
\]
donc, comme $f\in W$,
\[
g_0\ =\ g_1\qquad\quad\text{dans}\quad\Hot_W\ .
\]
On en conclut, en prenant un point $j_1$ dans $J_1$, et l'inclusion $\alpha_{j_1}:e\longrightarrow J$ correspondante, que
\msp{page 133 quinquies}
\[
g_0\alpha_{j_1}\ =\ g_1\alpha_{j_1}\ ,
\]
\textit{i.e.} les deux sections de $\varepsilon$ sur $e$ sont égales dans $\Hot_W$. Comme le foncteur $\Cat\longrightarrow\Hot_W$ commute aux produits finis \texttt{[cette propriété n'est pas établie sous les seules conditions Loc\,(1,2,3)~; elle est démontrée dans Pursuing Stacks, section 64, en remplaçant Loc\,(3) par l'axiome plus fort Loc\,(3\,bis)]}, il s'ensuit que pour $X$ dans Cat, les deux flèches
\[
X\ \underset{\partial_1^X}{\overset{\partial_0^X}{\rightrightarrows}}\ X\times\varepsilon\ =\ X\amalg X
\]
sont égales dans $\Hot_W$. Si alors
\[
f_0,\ f_1\ :\ X\ \rightrightarrows\ Y
\]
sont deux flèches dans Cat, elles sont de la forme
\[
f_0\ =\ h\partial_0^X\ ,\qquad\quad f_1\ =\ h\partial_1^X\ ,
\]
avec
\[
h\ =\ (f_0,f_1)\ :\ X\amalg X\ \longrightarrow\ X\ , % [sic] X\amalg X \longrightarrow X -> X\amalg X \longrightarrow Y, h\partial_i^X = f_i : X \to Y requires target Y
\]
et comme $\partial_0^X=\partial_1^X$ dans $\Hot_W$, on en conclut $f_0=f_1$ dans $\Hot_W$, \quad \textit{q.e.d.}


\endgroup
% END UNIT Der-ChXV-05-11-10/Der-ChXV-05-11-10_p071-075.tex
% BEGIN UNIT Der-ChXV-05-11-10/Der-ChXV-05-11-10_p076-077.tex
\begingroup\makeatletter%
\newcommand{\printedpage}[1]{\marginpar{\footnotesize p.~#1}}%
\newcommand{\msp}[1]{\marginpar{\footnotesize [#1]}}%
\newcommand{\HOT}{\mathrm{HOT}}%
\newcommand{\hot}{\mathrm{hot}}%
\makeatother\let\printedpage\derpp\let\msp\dermsp%
\let\mspin\DERmspin%


\printedpage{74}
% head: A. Grothendieck : Dérivateurs, Ch. XV, Édition des Universités de Montpellier II et Paris VII
% foot: Version 5 novembre 2010

\msp{page 134}

\section*{10\quad Application au dérivateur $\HOT_W$}

\textbf{Théorème 10.1.} \textit{Soit $W$ un localiseur dans} Cat \textit{satisfaisant les conditions} Loc\,(1, 2, 3, 4, 7\,bis) (\textit{axiomes de saturation forme forte} (\textit{facteurs directs}), \textit{de l'objet final, de localisation forme faible, du carré cocartésien, de la L-accessibilité}~\textsuperscript{(57)}\footnotetext[57]{\textbf{N.B.} Loc\,(1\,bis), \textit{i.e.} $W = W^{\natural}$, est déjà impliqué par la condition que $W$ soit stable par $\varinjlim$ filtrantes.}). \textit{Soit} $\HOT_W$ \textit{le prédérivateur sur} (Cat) \textit{associé à} $(\mathrm{Cat}, W)$ \textit{de la façon habituelle,}
\[
I \longmapsto \mathrm{Cat}(I)\,W(I)^{-1} \;=\; \underline{\mathrm{Hom}}(I^{\circ}, \mathrm{Cat})\,W(I)^{-1}\,. \tag*{(10.1.1)}
\]
\textit{On a alors ce qui suit.}

1\,$^{\circ}$) \textit{Pour toute flèche $u : I \longrightarrow J$ dans} Cat, \textit{le foncteur $u_!$} (\textit{adjoint à gauche de $u^*$}) \textit{existe} (\textbf{N.B.} On n'a pas besoin pour ceci de Loc\,(4) ni de Loc\,(7\,bis), il suffit que $W$ soit un localiseur fondamental \texttt{[on a néanmoins besoin de l'axiome de localisation forme forte Loc\,(3\,bis)]}), \textit{et il se calcule de la façon standard} (\textit{axiomes} Der~4\,bis \textit{et} Der~5\,bis). \textit{Le foncteur $u_*$, adjoint à droite de $u^*$, existe}~\textsuperscript{(58)}\footnotetext[58]{\textit{Toujours\,! Cf.} XIX, p. 29.} [\textit{si $u$ est de la forme $v \times \mathrm{id}_S$, où $v : I_0 \longrightarrow J_0$ est une flèche entre catégories ordonnées finies, et $S$ un objet quelconque de} Cat], \textit{et ce foncteur $u_*$ se calcule de façon standard} (\textit{axiomes} Der~4 \textit{et} Der~5). \textit{Donc on a une théorie des carrés $W$-cartésiens et des carrés $W$-cocartésiens dans les catégories} $\HOT_W(S)$, \textit{et des suites}
\msp{page 135}
\textit{longues de suspension et de cosuspension dans les catégories} $\HOT^{\bullet}_W(S)$. (\textit{Plus exactement, ce sont des théories relatives aux dérivateurs induits} $\HOT_{W,S}$, $\HOT^{\bullet}_{W,S}$ \textit{sur $S$,} $I \longmapsto \HOT_W(S \times I)$ \textit{resp.} $I \longmapsto \HOT^{\bullet}_W(S \times I)$ \ldots). \textit{De plus, les carrés cocartésiens et les carrés cartésiens relatifs aux dérivateurs induits} $\HOT_{W,S}$ \textit{sont} exacts, i.e. \textit{les suites longues de suspension et de cosuspension sont exactes} (\textit{axiomes d'exactitude à droite et à gauche} Der~6\,bis \textit{et} Der~6).

2\,$^{\circ}$) \textit{Le prédérivateur} $\HOT_W$ \textit{satisfait aussi l'axiome} Der~1 \textit{de décomposition finie}~\textsuperscript{(59)}\footnotetext[59]{Sans aucun axiome sur $W$\,! \texttt{[On a besoin simplement que les identités soient dans} $W$\texttt{.]}} (\textit{pour} $I = \coprod_{\alpha} I_{\alpha}$), \textit{et l'axiome} Der~2 \textit{de localisation. De plus, $W$ est fortement saturé,} i.e. $f \in W \iff \hot_W(f)$ \textit{isomorphisme.}

3\,$^{\circ}$) \textit{Concernant l'axiome} Der~3 \textit{et ses variantes} (\textit{cf.} I, pages 111-112, 125), \textit{on a ceci :}

a)~\textsuperscript{(60)}\footnotetext[60]{On n'a besoin ici que de Loc\,(1\,ter, 2, 3) (Loc\,(1\,ter) est la saturation forte, \textit{i.e.} $f \in W \iff \hot_W(f)$ isomorphisme).}. \textit{Soit $f : S' \longrightarrow S$ une flèche dans} Cat. \textit{Pour que ce soit une} $\HOT_W$\textit{-équivalence,} i.e. \textit{pour que le foncteur} $f^* : \HOT_W(S) \longrightarrow \HOT_W(S')$ \textit{induise une} équivalence \textit{entre les sous-catégories pleines formées des objets} constants \textit{sur $S$, $S'$, il faut et il suffit que $f$ soit dans $W$. L'énoncé correspondant dans les dérivateurs induits} $\HOT_{W,I}$ \textit{est également valable, pour $I \neq \emptyset$. Donc a fortiori,} $\HOT_W$ \textit{et} $\HOT_{W,I}$ \textit{satisfont}
\msp{page 136}
\printedpage{75}
% head: A. Grothendieck : Dérivateurs, Ch. XV, Édition des Universités de Montpellier II et Paris VII
% foot: Version 5 novembre 2010
\textit{l'axiome de l'objet final,} i.e. \textit{si $S$ dans} Cat \textit{a un objet final, alors $S \longrightarrow e$ est une} $\HOT_W$\textit{-équivalence, et même une} $\HOT_{W,I}$\textit{-équivalence pour toute $I$ dans} Cat.

b) \textit{Supposons que $W$ satisfasse} Loc\,(1\,ter, 2, 3, 5\,b) (\textit{donc l'axiome} Loc\,(5\,b) \textit{des $W$-fibrations, variante} b). \textit{Alors pour toute $f : S' \longrightarrow S$ dans $W$,} $\HOT^{\mathrm{lc}}_W(S) \longrightarrow \HOT^{\mathrm{lc}}_W(S')$ \textit{est une} équivalence de catégories (\textit{la réciproque étant clair par} (a)). \textit{Itou pour} % [sic] clair -> claire, accord avec « réciproque »
\[
\HOT^{\mathrm{lc}}_{W,I}(S) \longrightarrow \HOT^{\mathrm{lc}}_{W,I}(S') \qquad (\text{i.e. \textit{pour}}\ \ \HOT^{\mathrm{lc}/I}_W(S \times I) \longrightarrow \HOT^{\mathrm{lc}/I}_W(S' \times I)\ ).
\]
\textit{En d'autres termes, pour le dérivateur} $\HOT_W$ \textit{et pour tous les} $\HOT_{W,I}$ ($\mathbf{D}$ \textit{étant un des ces dérivateurs})~\textsuperscript{(61)}\footnotetext[61]{À vérifier avec soin dans le cas $\mathbf{D} = \HOT_{W,I}$.}, \textit{toute $\mathbf{D}$-équivalence est une $\mathbf{D}$-équivalence} 2\textit{-complète,} i.e. \textit{pour $f : S' \longrightarrow S$, pour que $f^* : \mathbf{D}(S) \longrightarrow \mathbf{D}(S')$ induise une équivalence entre les sous-catégories $\mathbf{D}^{\mathrm{lc}}(S) \longrightarrow \mathbf{D}^{\mathrm{lc}}(S')$, il suffit déjà qu'elle induise un foncteur pleinement fidèle sur les sous-catégories pleines formées des objets constants, et ceci équivaut} (\textit{si $I \neq \emptyset$}) \textit{à la condition que $f \in W$.} % [sic] des ces dérivateurs -> de ces dérivateurs, faute de frappe

\msp{page 137}

\textbf{Commentaire.} Dans cet énoncé récapitulatif de ce que je sais prouver à présent pour $\HOT_W$, j'ai un peu mélangé les provenances diverses – ce qui est mis “pour mémoire” de VI (et qui ne dépend pas de Loc\,(4, 7 ou 7\,bis)), \textit{i.e.} ce qui dépend pour l'essentiel de la construction de $u_!$ donnée dans VI (donc ce qui concerne $u_!$ dans 1$^{\circ}$, \textit{moins} l'exactitude à droite, et tout ce qui concerne 3$^{\circ}$, $(a)$, $(b)$).

La démonstration de 10.1 résultera des développements du paragraphe suivant \texttt{[qui n'existe pas]}, compte tenu du résultat principal du paragraphe 8 (th. 8.4), et du paragraphe 7 (cor. 7.4), qui sont les ingrédients “durs” du théorème 10.1.


\endgroup
% END UNIT Der-ChXV-05-11-10/Der-ChXV-05-11-10_p076-077.tex
\clearpage\setcounter{footnote}{0}\phantomsection\addcontentsline{toc}{section}{Chapitre XVI}\markboth{Chapitre XVI}{Chapitre XVI}
% BEGIN UNIT der-ch16/der-ch16_p001-005.tex
\begingroup\makeatletter%
\newcommand{\printedpage}[1]{\marginpar{\footnotesize p.~#1}}%
\newcommand{\msp}[1]{\marginpar{\footnotesize [#1]}}%
\newcommand{\circled}[1]{\tikz[baseline=(c.base)]\node[draw,circle,inner sep=1pt](c){#1};}%
\newcommand{\Cat}{\mathrm{Cat}}%
\newcommand{\Ob}{\mathrm{Ob}\,}%
\newcommand{\hot}{\mathrm{hot}}%
\makeatother\let\printedpage\derpp\let\msp\dermsp%
\let\mspin\DERmspin%


% ===================== page 1 (PDF) =====================
\printedpage{[non numérotée]}
% head: (none printed)

\begin{center}
{\sffamily\LARGE LES DÉRIVATEURS}

\bigskip

{\large\textbf{ALEXANDRE GROTHENDIECK}}

\vspace{3em}

{\large\textbf{Chapitre XVI}}

\medskip

{\large\textbf{Les localiseurs fondamentaux}}
\end{center}

\par

% ===================== page 2 (PDF) =====================
\printedpage{[non numérotée]}
% head: (none printed)

\par\noindent\textbf{Note des éditeurs de l’édition transcrite.}\par
Ce texte a été déchiffré et transcrit en \LaTeXe{} par M. Künzer. Il a été édité par M. Künzer, J. Malgoire, et G. Maltsiniotis. La transcription est aussi fidèle que possible au manuscrit. Seules quelques corrections mineures évidentes ont été effectuées. Pour les rares ajouts ou commentaires des éditeurs, ainsi que pour la numérotation originale des pages du manuscrit, les caractères de machine à écrire \texttt{[typewriter]} entre crochets sont utilisés.

Cette édition est provisoire. Les remarques, commentaires et corrections sont bienvenus. Envoyer un message à :

\textbf{[e-mail address removed]}

\hfill G. Maltsiniotis\hspace*{2cm}

\par

% ===================== page 3 (PDF) =====================
\printedpage{1}
% head: A. Grothendieck : Dérivateurs, Ch. XVI, Édition des Universités de Montpellier II et Paris VII
% foot: Version 18 mars 2019

\msp{page 1}

\section*{1\quad Axiomes Loc(1) à Loc(3) (et variantes) : localiseurs fondamentaux}

J’ai envie de reprendre l’exposé des axiomes utiles pour un localiseur fondamental
\[
W \subset \mathrm{Fl}(\Cat),
\]
en me guidant cette fois plus sur la signification géométrique des axiomes, que sur les alias $\pm$ techniques qu’on rencontre (ou ne rencontre pas) pour les vérifier. Je révise donc dans cet esprit la présentation de XV §9.

\begin{description}
\item[Loc (1)] (\textit{saturation}).
\begin{itemize}
\item[\textit{a})] Iso $\subset W$.
\item[\textit{b})] Si deux parmi $f$, $g$, $gf$ sont dans $W$, aussi le troisième.
\item[\textit{c})] Si $gf$ et $fg$ sont dans $W$, $f$ et $g$ aussi.
\end{itemize}
\end{description}

Les conditions (\textit{a}) et (\textit{b}) toutes seules, je vais les appeler (pour un $W \subset \mathrm{Fl}\,\mathcal{M}$ général) \textit{saturation faible} du localiseur. Les conditions (\textit{a}), (\textit{b}), (\textit{c}), \textit{saturation modérée} (“mild saturation”). Autres variantes :

\begin{description}
\item[Loc (1 bis)] Les conditions (\textit{a}), (\textit{b}), (\textit{c}) du haut, plus
\begin{itemize}
\item[\textit{d})] $W$ stable par facteurs directs.
\end{itemize}
\end{description}

C’est bien sûr impliqué par :

\begin{description}
\item[Loc (1 ter)] $W$ est \textit{fortement saturé, i.e.} si $f \in \Cat$ est telle que $\hot_W(f)$ est un isomorphisme, alors $f \in W$.
\item[Loc (2)] (\textit{Objet final}). Si $X$ dans Cat a un objet final, alors $X$ est $W$-asphérique, \textit{i.e.} $X \longrightarrow e$ est dans $W$.
\end{description}

Parfois, cet axiome est conséquence de la variante affaiblie suivante.

\begin{description}
\item[Loc (2$_0$)] (Axiome du segment). $\boldsymbol{\Delta}^1$ est $W$-asphérique.
\end{description}

\msp{page 2}

\begin{description}
\item[Loc (3)] (\textit{localisation,} forme faible). Soit $f : X \longrightarrow S$ une flèche dans Cat $W$\textit{-asphérique, i.e.} telle que $\forall\ s \in \Ob S$, la catégorie localisée $X/s$ sur $S$ en $s$, $X/s = X \times_S S/s$, soit $W$-asphérique. Alors $f \in W$.\footnote[1]{Dans une version finale, il faudrait intervertir les notations pour Loc (3) et pour Loc (3 bis), car ce dernier m’apparaît comme le plus important des deux - il est satisfait pour tous les localiseurs qui ont été utilisées (à ma connaissance), et est valable pour tout localiseur de la forme $W_{\mathbf{D}}$, où $\mathbf{D}$ est un dérivateur (ou seulement un prédérivateur satisfaisant Der (4,5) \texttt{[il faut aussi Der (2) ; la numérotation des axiomes est celle du chapitre I, p. 68]}).} % mark printed as superscript (1); [sic] utilisées -> utilisés, agreement with "localiseurs"
\end{description}

\textbf{Rappel 1.} Loc (1, 2, 3) impliquent que l’on a
\[
W = W^\circ \qquad \textit{i.e. } f \in W \iff f^\circ \in W,
\]
\printedpage{2}
% head: A. Grothendieck : Dérivateurs, Ch. XVI, Édition des Universités de Montpellier II et Paris VII
% foot: Version 18 mars 2019
où on dénote par $X \longmapsto X^\circ$, $f \longmapsto f^\circ$ le foncteur de passage à la catégorie opposée, foncteur
\[
X \longmapsto X^\circ : \Cat \xrightarrow{\ \sim\ } \Cat
\]
qui est un \textit{automorphisme} de Cat (donc commute à toutes $\varinjlim$ et $\varprojlim$ ). Il s’ensuit que l’ensemble d’axiomes Loc (1, 2, 3) équivaut à Loc (1, 2$^{\mathrm{opp}}$, 3$^{\mathrm{opp}}$), où Loc (2$^{\mathrm{opp}}$) est l’axiome “dual” de Loc (2), savoir que tout objet de Cat ayant un objet initial est $W$-asphérique, et de même Loc (3$^{\mathrm{opp}}$) l’axiome que toute $f : X \longrightarrow S$ $W$\textit{-coasphérique, i.e.} telle que les catégories colocalisées $s\backslash X$ de $X$ sur $S$ soient $W$-asphériques, est dans $W$.

\textbf{Définition 1.1.} On appelle \textit{localiseur fondamental} une partie $W$ de Fl(Cat) satisfaisant les axiomes Loc (1, 2, 3) (\textit{i.e.} saturation modérée, axiome de l’objet final, axiome de localisation).

\msp{page 3}

\textbf{Rappel 1.2.} Si $W$ est un localiseur fondamental, alors tout \textit{homotopisme} $f : X \longrightarrow Y$ de Cat est dans $W$ (c’est pour cela essentiellement que l’on a imposé (\textit{c}) dans la saturation modérée Loc (1)). Si $f, g : X \rightrightarrows Y$ sont homotopes (\textit{i.e.} appartiennent à une même composante connexe de $\underline{\mathrm{Hom}}(X, Y)$), alors $\hot_W(f) = \hot_W(g)$ et $f \in W \iff g \in W$.

Voici la forme forte de Loc (3).

\begin{description}
\item[Loc (3 bis)] (\textit{Localisation,} forme forte). Soit $S$ dans Cat et $f : X \longrightarrow Y$ un morphisme dans $\Cat/S$ tel que “$f$ soit une $W$-équivalence localement sur $S$”, \textit{i.e.} pour tout $s \in \Ob S$, le morphisme induit
\[
f/s : X/s \longrightarrow Y/s
\]
pour les catégories localisées de $X$, $Y$ sur $S$ en $s$, soit dans $W$. Alors $f \in W$.
\end{description}

L’axiome Loc (3) n’est autre que le cas particulier de Loc (3 bis) relatif au cas où $Y = S$.

\textbf{Proposition 1.3.} \textit{Soit $W \subset \mathrm{Fl}(\mathcal{M})$. Alors on a les équivalences suivantes}
\[
\begin{aligned}
\text{Loc (1, 2, 3)} &\iff \text{Loc (1, $2_0$, $3'$)}\\
\text{Loc (1, 2, 3 bis)} &\iff \text{Loc (1, $2_0$, $3'$ bis),}
\end{aligned}
\]
\textit{où les axiomes} Loc ($3'$) \textit{et} Loc ($3'$ bis) \textit{se formulent ainsi :}

\msp{page 4}

\begin{description}
\item[Loc (3$'$)] \circled{\textit{a}} (Axiome des Cat-fibrations, forme directe faible). Soient $S$ dans Cat, $X$ une catégorie fibrée scindée sur $S$, telle que les fibres $X_s$ ($s \in \Ob S$) soient $W$-asphériques. Alors $X \longrightarrow S$ est dans $W$.

\circled{\textit{b}} Soit $X$ dans $\Cat/S$, alors $i_X : X \longrightarrow \Phi_S(X)$ est dans $W$\footnote[2]{(\textit{b}) concerne la \textit{comparaison} d’une $X$ quelconque sur $S$, avec la catégorie fibrée scindée $\Phi_S(X)$ sur $S$.}. % mark printed as superscript (2)
\end{description}

\par

% ===================== page 5 (PDF) =====================
\printedpage{3}
% head: A. Grothendieck : Dérivateurs, Ch. XVI, Édition des Universités de Montpellier II et Paris VII
% foot: Version 18 mars 2019

\begin{description}
\item[Loc (3$'$ bis)] \circled{\textit{a}} (Axiome des Cat-fibrations, forme directe forte). Soient $S$ dans Cat, $X$, $Y$ deux catégories fibrées scindées sur $S$, $f : X \longrightarrow Y$ un $S$-foncteur cartésien compatible aux scindages, tel que $\forall\ s \in \Ob S$, le foncteur induit sur les fibres
\[
f_s : X_s \longrightarrow Y_s
\]
soit dans $W$. Alors $f \in W$.

\circled{\textit{b}} Comme Loc ($3'$) (\textit{b}).
\end{description}

\textbf{N.B.} En fait, moyennant Loc (1, 2), les énoncés Loc ($3'$), Loc ($3'$ bis) impliquent les énoncés similaires sans scindages, et sans même supposer $f$ cartésien dans Loc ($3'$ bis), \textit{cf.} démonstration ci-dessous.

\textsc{Démonstration.} \textit{a}) Loc (1, 2, 3) $\Longrightarrow$ Loc (1, $2_0$, $3'$). Il suffit de prouver Loc ($3'$). Prouvons Loc ($3'$ $a$). Pour ceci, sous les conditions de Loc ($3'$ $a$), prouvons que $X \longrightarrow S$ est $W$-coasphérique, \textit{i.e.} que les catégories
\[
s\backslash X = X \times_S s\backslash S
\]
sont $W$-asphériques. Il en résultera bien que $f \in W$ en vertu de Loc (3) appliqué à $f^\circ$ et de $W = W^\circ$. Mais
\msp{page 5}
considérons l’inclusion canonique
\[
i_s : X_s \lhook\joinrel\longrightarrow s\backslash X,
\]
il résulte du fait que $p : X \longrightarrow S$ est Cat-fibrant que ce foncteur est $W$-asphérique (il a même un adjoint à droite, \textit{i.e.} les catégories $X_s/u$, où $u \in \Ob(s\backslash X)$, ont un objet final). Donc par Loc (3) on a $i_s \in W$, et comme $X_s$ est $W$-asphérique par hypothèse, de même $s\backslash X$, q.e.d.

Prouvons Loc ($3'$ $b$) : On note que $i_X : X \longrightarrow \Phi_S(X)$ est $W$\textit{-coasphérique} (il a même un adjoint à gauche), donc $i_X^\circ$ est $W$-asphérique, donc dans $W$, donc $i_X^\circ \in W^\circ$ car $W = W^\circ$. % [sic] i_X^\circ \in W^\circ -> i_X \in W (or i_X \in W^\circ), apparently intended: from i_X^\circ \in W one gets i_X \in W^\circ = W

$a$ bis) Loc (1, 2, 3 bis) $\Longrightarrow$ Loc (1, $2_0$, $3'$ bis). Cette fois il faut prouver seulement Loc ($3'$ bis). Pour (\textit{b}), c’est déjà fait. Reste à prouver (\textit{a}). Si $f$ est comme dans l’énoncé Loc ($3'$ bis $a$), prouvons que pour tout $s$ dans $S$, le morphisme induit
\[
s\backslash f : s\backslash X \longrightarrow s\backslash Y
\]
est dans $W$. Comme $W = W^\circ$, on aura donc $(s\backslash f)^\circ \in W$, et appliquant Loc (3 bis) à $f^\circ$, on trouve $f^\circ \in W$, donc $f \in W$, et on aura fini. Mais $s\backslash f$ s’insère dans le carré commutatif
\msp{page 6}
\[
\begin{tikzcd}
X_s \arrow[r, "f_s"] \arrow[d, "i_s^X"'] & Y_s \arrow[d, "i_s^Y"] \\
s\backslash X \arrow[r, "s\backslash f"] & s\backslash Y,
\end{tikzcd}
\]


\endgroup
% END UNIT der-ch16/der-ch16_p001-005.tex
% BEGIN UNIT der-ch16/der-ch16_p006-010.tex
\begingroup\makeatletter%
\newcommand{\printedpage}[1]{\marginpar{\footnotesize p.~#1}}%
\newcommand{\msp}[1]{\marginpar{\footnotesize [#1]}}%
\newcommand{\Loc}{\mathrm{Loc}}%
\makeatletter%
\renewcommand\@makefnmark{\hbox{\textsuperscript{(\@thefnmark)}}}%
\renewcommand\@makefntext[1]{\noindent\textsuperscript{\@thefnmark}\,#1}%
\makeatother%
\makeatother\let\printedpage\derpp\let\msp\dermsp%
\let\mspin\DERmspin%

\setcounter{footnote}{2}

\printedpage{4}
% head: A. Grothendieck : Dérivateurs, Ch. XVI, Édition des Universités de Montpellier II et Paris VII
où, on l'a vu dans $(a)$, $i_s^X$ et $i_s^Y$ sont dans $W$. Par hypothèse on a $f_s \in W$, donc $s\backslash f \in W$, q.e.d.

\textit{b}) Loc (1, $2_0$, $3'$) $\Longrightarrow$ Loc (1, 2, 3). Prouvons d'abord Loc (2). Appliquant Loc ($3'$ \textit{a}) à la projection
\[
\boldsymbol{\Delta}^1 \times X \longrightarrow X \qquad \text{(où $X \in \operatorname{Ob}\mathrm{Cat}$),}
\]
qui est fibrée scindée à fibres isomorphes à $\boldsymbol{\Delta}^1$, donc $W$-asphériques par Loc ($2_0$), on voit que cette projection est dans $W$. Cela permet d'appliquer les arguments standard d'homotopie (\textit{cf.} XII, p. 1 \ldots) pour conclure qu'un homotopisme est dans $W$, donc une catégorie contractile est $W$-asphérique, a fortiori une catégorie avec objet final (ou avec objet initial) est $W$-asphérique, \textit{i.e.} on a Loc (2).

Prouvons Loc (3). Comme Loc (1, 2, 3) $\Longleftrightarrow$ Loc (1, $2^\circ$, $3^\circ$), il suffit de prouver Loc ($3^\circ$), vu qu'on a prouvé Loc ($2^\circ$). Soit
\[
f : X \longrightarrow S \qquad \text{$W$-\textit{co}asphérique,}
\]
prouvons $f \in W$. Introduisons la catégorie \textit{fibrée} sur $S$ dont les fibres sont les $s\backslash X$, soit $\Phi_S(X)$, et considérons l'inclusion canonique
\[
i_X : X \longrightarrow \Phi_S(X).
\]
\msp{page 7}Par Loc ($3'$ \textit{b}) elle est dans $W$. D'autre part, les fibres $s\backslash X$ de $\Phi_S(X)$ sont $W$-coasphériques par hypothèse, donc par Loc ($3'$ \textit{a}) $\Phi_S(X) \longrightarrow S$ est dans $W$. Comme $i_X \in W$, il s'ensuit $(f : X \longrightarrow S) \in W$, q.e.d.

\textit{b} bis) Loc (1, $2_0$, $3'$ bis) $\Longrightarrow$ Loc (1, 2, 3 bis). On a déjà prouvé Loc (1, 2, 3) (vu que Loc ($3'$ bis) contient Loc ($3'$) comme cas particulier). Prouvons donc Loc (3 bis) en admettant déjà tout ça. Il suffira encore de le prouver sous forme duale, puisque
\[
\text{Loc (1, 2, 3 bis)} \iff \text{Loc(1, $2^\circ$, 3 bis$^\circ$),}
\]
\msp{page 8}et que Loc ($2^\circ$) est connu par Loc (1, 2, 3) $\Longleftrightarrow$ Loc (1, $2^\circ$, $3^\circ$). Il reste donc à prouver que si
\[
f : X \longrightarrow Y \qquad \text{dans } \mathrm{Cat}/S
\]
est telle que
\[
s\backslash f : s\backslash X \longrightarrow s\backslash Y
\]
est dans $W$ pour tout $s \in \operatorname{Ob} S$, alors $f \in W$. Considérons alors le carré commutatif
\[
\begin{tikzcd}
X \arrow[r, "i_X"] \arrow[d, "f"'] & \Phi_S(X) \arrow[d, "\Phi_S(f)"] \\
Y \arrow[r, "i_Y"] & \Phi_S(Y),
\end{tikzcd}
\]
où $i_X, i_Y \in W$ par Loc ($3'$ bis \textit{b}). L'hypothèse sur $f$ signifie que $\Phi_S(f)$ induit des $W$-équivalences sur les fibres, donc par Loc ($3'$ bis) $\Phi_S(f) \in W$, donc $f \in W$ par Loc (1), q.e.d.
% foot: Version 18 mars 2019

\printedpage{5}
% head: A. Grothendieck : Dérivateurs, Ch. XVI, Édition des Universités de Montpellier II et Paris VII
\section*{2\quad Axiomes Loc (4) et Loc (5) (des sommes directes, de connexité)\protect\footnote{Il s'agit ici d'axiomes qu'on peut considérer comme étant de nature plus ou moins “triviale”. Dans tous les cas à ma connaissance, ils sont bel et bien trivialement vérifiés.}}

\begin{description}[leftmargin=7em, labelwidth=6em]
\item[Loc (4)] (Axiome des \textit{sommes directes}). Pour une famille $(f_i)_{i\in I}$ de flèches dans Cat, si $f = \coprod f_i$, on a $f \in W$ si et seulement si $f_i \in W$ pour tout $i \in I$.
\end{description}

Cet axiome se décompose en deux, en le “si” et le “seulement si”, que je vais formuler sous la forme la plus économique :

\textbf{Proposition 2.1.} \textit{L'axiome} Loc (4) \textit{des sommes équivaut à la conjonction des deux axiomes suivants :}

\begin{description}[leftmargin=7em, labelwidth=6em]
\item[Loc (4 a)] \msp{page 9}Soit $(f_i)_{i\in I}$ une famille de flèches dans Cat, si les $f_i \in W$, alors $f = \coprod f_i \in W$.
\item[Loc (4 b)] Soit $S$ dans Cat, et $S_0 \hookrightarrow S$ un sommand direct (sous-catégorie à la fois ouverte et fermée). Soit $X$ dans Cat$/S$. Si $X \longrightarrow S$ est dans $W$, alors $X_0 = X|S_0 \longrightarrow S_0$ aussi. (En d'autres termes, l'inclusion $S_0 \hookrightarrow S$ est une “$W$-fibration”, \textit{i.e.} est dans $\mathrm{Fib}_W$.)
\end{description}

\textbf{Corollaire 2.2.} \textit{Considérons la variante de} Loc (4 a) \textit{obtenue en se bornant aux familles finies.} (\textit{Elle équivaut à} Loc (4) \textit{lorsque $W$ est stable par} $\varinjlim$ \textit{filtrantes,} cf. \textit{axiome} Loc (8) \textit{plus bas}). \textit{Cette variante équivaut à}

\begin{description}[leftmargin=7.5em, labelwidth=6.5em]
\item[Loc (4 a$_0$)] Tout objet $X$ de Cat est “$W$-cofibrant”, \textit{i.e.} $\emptyset \longrightarrow X$ est dans $\mathrm{Cof}_W$, \textit{i.e.} pour toute $f : Y \longrightarrow Y'$ qui est dans $W$, le morphisme correspondant $\mathrm{id}_X \amalg f : X \amalg Y \longrightarrow X \amalg Y'$ est dans $W$.
\end{description}

\textbf{Rappel 2.3.} De toutes façons, si $W$ est un localiseur fondamental (ce que nous supposerons désormais, sauf mention du contraire), l'analogue \textit{fini} de Loc (4 a) pour les produits est vrai, \textit{i.e.} si on a une famille \textit{finie} de flèches $(f_i)_{i\in I}$ dans Cat, si les $f_i$ sont dans $W$, alors $f = \prod f_i \in W$. \texttt{[Pour cela, il faut supposer que $W$ satisfait à l'axiome de localisation, forme \textit{forte}, Loc 3 bis.]} L'analogue pour une famille infinie n'est vérifié pratiquement pour aucun localiseur fondamental.

\msp{page 10}\textbf{Proposition 2.4.} \textit{Soit $W$ un localiseur fondamental. Alors les produits finis existent dans} $\mathrm{Hot}_W$\textit{, et le foncteur} $\mathrm{hot}_W : \mathrm{Cat} \longrightarrow \mathrm{Hot}_W \overset{\textit{déf}}{=} (\mathrm{Cat})W^{-1}$ \textit{y commute. Si $W$ satisfait à} Loc (4 a) (resp. Loc (4 a$_0$)), \textit{alors les sommes quelconques} (resp. \textit{finies}) \textit{existent dans} $\mathrm{Hot}_W$\textit{, et le foncteur} $\mathrm{hot}_W$ \textit{y commute. De plus, dans ce cas les sommes sont distributives par rapport aux produits (vu qu'il en est ainsi dans} Cat).

La démonstration résulte simplement du résultat général :

\textbf{Lemme 2.5.} \textit{Soit $(\mathcal{M}, W)$ une catégorie de modèles ($W$ partie quelconque de} $\mathrm{Fl}\,\mathcal{M}$). \textit{Supposons que les produits} (resp. \textit{les produits finis,} resp. \textit{les produits binaires}) \textit{existent dans $\mathcal{M}$, et que de plus pour toute famille $(f_i)_{i\in I}$ de flèches de $\mathcal{M}$} (resp. \textit{toute famille finie,}
% foot: Version 18 mars 2019
\printedpage{6}
% head: A. Grothendieck : Dérivateurs, Ch. XVI, Édition des Universités de Montpellier II et Paris VII
resp. \textit{famille binaire}) \textit{$f_i \in W$ pour tout $i \in I$ implique $f = \prod f_i \in W$. Alors les produits du type envisagé existent dans} $\mathrm{Ho}_W = \mathcal{M}W^{-1}$\textit{, et le foncteur} $\mathrm{ho}_W : \mathcal{M} \longrightarrow \mathcal{M}W^{-1}$ \textit{y commute. Enfin, si $\mathcal{M}$ admet un objet final $e$, alors $e$ est aussi objet final de} $\mathrm{Ho}_W$.

\texttt{[Il n'est pas clair que l'assertion soit vraie dans cette généralité pour les produits \textit{infinis} ; elle est vraie si $W$ est la classe des équivalences faibles d'une \textit{catégorie de modèles de Quillen.}]}

\msp{page 11}On en conclut les énoncés duaux concernant les sommes directes et les objets initiaux.

La démonstration est un cas particulier du résultat bien utile suivant, appliqué au couple de foncteurs adjoints
\[
\begin{tikzcd}
\mathcal{M} \arrow[r, shift left, "\mathrm{diag}"] & \mathcal{M}^I \arrow[l, shift left, "\text{\scriptsize\shortstack{produits\\ dans $\mathcal{M}$\\ indexés par $I$}}"]
\end{tikzcd}
\qquad\qquad :
\]

\textbf{Lemme 2.6}\footnote{\textit{cf.} XII, lemme page 248.}. \textit{Soient $(\mathcal{M}, W)$, $(\mathcal{M}', W')$ deux catégories de modèles, et}
\[
\begin{tikzcd}
\mathcal{M} \arrow[r, shift left, "\varphi"] & \mathcal{M}' \arrow[l, shift left, "\psi"]
\end{tikzcd}
\]
\textit{un couple de foncteurs adjoints compatibles aux localiseurs,} i.e. \textit{tels que}
\[
\varphi(W) \subset W', \qquad \psi(W') \subset W.
\]
\textit{Alors les foncteurs localisés}
\[
\begin{tikzcd}
\mathcal{M}W^{-1} \arrow[r, shift left, "\bar{\varphi}"] & \mathcal{M}'W'^{-1} \arrow[l, shift left, "\bar{\psi}"]
\end{tikzcd}
\]
\textit{sont également adjoints, les flèches d'adjonction pour $\bar{\varphi}\bar{\psi}$ et $\bar{\psi}\bar{\varphi}$ étant déduites de celles pour $\varphi\psi$ et $\psi\varphi$ par composition avec les foncteurs de localisation $\mathcal{M} \longrightarrow \mathcal{M}W^{-1}$, $\mathcal{M}' \longrightarrow \mathcal{M}'W'^{-1}$.}

\begin{description}[leftmargin=7em, labelwidth=6em]
\item[Loc (5)] \msp{page 12}(Axiome de \textit{connexité} ou de \textit{non-trivialité}).

Soit $f \in W$, alors $\pi_0(f)$ est bijectif.
\end{description}

Cette propriété de $W$ est d'une utilisation constante en pratique, quand on travaille avec $W$ au lieu de se borner à construire des $W$ et faire joujou avec les axiomes et leur stabilité par passage à des $\underline{\mathrm{Hom}}(I, M)$ (où ici $\mathcal{M} = \mathrm{Cat}$). D'autre part, on peut aussi le regarder comme un axiome de \textit{non-trivialité} pour $\mathrm{Hot}_W$. Je rappelle à ce sujet :

\textbf{Proposition 2.7.} \textit{Soit $W$ un localiseur fondamental. Conditions équivalentes :}
\begin{itemize}[leftmargin=3em]
\item[a)] $W$ \textit{satisfait l'axiome de connexité} Loc (5).
\item[b)] \textit{Le foncteur} $\pi_0 : \mathrm{Cat} \longrightarrow \mathrm{Ens}$ \textit{se factorise par un foncteur}
% foot: Version 18 mars 2019
\printedpage{7}
% head: A. Grothendieck : Dérivateurs, Ch. XVI, Édition des Universités de Montpellier II et Paris VII
\[
\begin{tikzcd}
\mathrm{Cat} \arrow[r, dotted] \arrow[rr, dotted, bend right, "\pi_0"'] & \mathrm{Hot}_W \arrow[r, "\pi_0"] & \mathrm{Ens}.
\end{tikzcd}
\]
\item[c)] \textit{La catégorie} $\mathrm{Hot}^*_W = \mathrm{Hot}_W\backslash\{\emptyset\}$ \textit{n'est pas équivalente à la catégorie ponctuelle.}
\item[c$'$)] \textit{La catégorie} $\mathrm{Hot}_W$ \textit{n'est pas équivalente à $e$ ni à $\boldsymbol{\Delta}^1$.}\footnote{\textbf{NB} Si $W$ est fortement saturée (condition Loc (1 ter), alors les conditions ci-contre (\textit{i.e.} Loc (5)) équivalent à ceci : % [sic] "(condition Loc (1 ter)," -> "(condition Loc (1 ter)),", parenthèse ouvrante non refermée
\begin{itemize}[leftmargin=2.5em]
\item[d)] \textit{$W$ n'est égal ni à Fl}(Cat) \textit{tout entier, ni à Fl}(Cat$^*$) $\cup$ $\{\mathrm{id}_\emptyset\}$ (i.e. \textit{Fl}(Cat)$\backslash W$ \textit{n'est ni vide, ni égal à l'ensemble des flèches $\emptyset \longrightarrow X$ avec $X \neq \emptyset$}).
\end{itemize}
Si $W$ satisfait Loc (7 D$_{W_\omega}$), alors $W$ satisfait Loc (5) si et seulement si $W \neq \mathrm{Fl}(\mathrm{Cat})$ (\textit{cf.} XII, th. 2, p. 20).}
\end{itemize}

Pour la démonstration, et pour d'autres formes équivalentes de l'axiome de connexité, \msp{page 13}je renvoie à XV 9.12 (p. 133 bis). Pour le nom “axiome de connexité”, je peux offrir l'explication suivante 2.8 et 2.9.

\textbf{Corollaire 2.8.} \textit{Si l'axiome de connexité} Loc (5) \textit{est satisfait, alors les objets $W$-asphériques sont} 0\textit{-connexes} (\textit{tautologie}), \textit{et la réciproque est vraie si $W$ est fortement saturé.} (\textit{Ou aussi, si $W$ satisfait l'axiome des fibrations} Loc (7) \textit{ci-dessous, puisque par} XII, th. 2, p. 20, \textit{dans ce cas,} Loc (5) \textit{équivaut à $W \neq \mathrm{Fl}(\mathcal{M})$ \ldots}).

Car si l'axiome de connexité n'est pas satisfait, alors pour toute $X$ non vide dans Cat, $X \longrightarrow e$ est un isomorphisme, donc ($W$ étant fortement saturé) $X \longrightarrow e$ est dans $W$, \textit{i.e.} $X$ est $W$-asphérique, donc il y a des objets $W$-asphériques non 0-connexe. Notons aussi : % [sic] "non 0-connexe" -> "non 0-connexes", accord avec "des objets"

\textbf{Corollaire 2.9.} \textit{Supposons que $W$ satisfait à l'axiome} Loc (4) \textit{des sommes directes} (\textit{ou seulement à} Loc (4 b)). \textit{Pour que $W$ satisfasse l'axiome de connexité, il faut et il suffit que pour toute flèche $f : X \longrightarrow Y$ avec $Y$} 0\textit{-connexe, si $f$ est dans $W$, alors $X$ est} 0\textit{-connexe.} (\textit{C'est nécessaire bien sûr sans supposer} Loc (4 b) \textit{ni rien du tout.})

\msp{page 14}\texttt{[Les numéros 2.7, 2.8 et 2.9 apparaissent deux fois.]}

\textbf{2.7.} Considérons maintenant le foncteur composé
\[
\mathrm{Ens} \xrightarrow{I \longmapsto I_{\mathrm{Cat}}} \mathrm{Cat} \xrightarrow{\quad \mathrm{hot}_W \quad} \mathrm{Hot}_W, \qquad I \longmapsto I_{\mathrm{Hot}_W} \text{ ou } I_W, \tag*{(2.7.1)}
\]
où le premier foncteur est celui qui associe à l'ensemble $I$ la catégorie discrète $I_{\mathrm{Cat}}$ ayant $I$ comme ensemble d'objets. (Le plus souvent, on identifiera $I_{\mathrm{Cat}}$ avec $I$ dans les notations, tout comme pour un ensemble ordonné $J$, on identifie $J$ et la “catégorie ordonné” % [sic] "catégorie ordonné" -> "catégorie ordonnée", accord
définie par $J$. Ici $I$ peut être considéré comme muni de l'ordre discret.) On a d'ailleurs
\[
I_{\mathrm{Cat}} \simeq I \times e \ \ (\overset{\text{déf}}{=} \text{somme de $I$ copies de $e$}), \tag*{(2.7.2)}
\]
donc quand Loc (4 a) est satisfait, donc $\mathrm{hot}_W$ commute aux sommes directes, on aura un isomorphisme canonique
\[
\begin{aligned}
I_W \simeq I \times e_W \quad &\text{(où $e_W$ est l'objet final de $\mathrm{Hot}_W$,} \\
&\text{\textit{i.e.} $e$ considéré comme objet de $\mathrm{Hot}_W$).}
\end{aligned} \tag*{(2.7.3)}
\]
% foot: Version 18 mars 2019

\printedpage{8}
% head: A. Grothendieck : Dérivateurs, Ch. XVI, Édition des Universités de Montpellier II et Paris VII
\textbf{Proposition 2.8.} \textit{Supposons l'axiome de connexité} Loc (5) \textit{satisfait, d'où un foncteur}
\[
\pi_0 : \mathrm{Hot}_W \longrightarrow \mathrm{Ens}. \tag*{(2.8.1)}
\]
\textit{Ce foncteur est} adjoint à gauche \textit{du foncteur $I \longmapsto I_W$} (2.7.1), i.e. \textit{on a une bijection fonctorielle en $(X, I)$ dans} $\mathrm{Hot}_W \times \mathrm{Ens}$ :
\msp{page 15}
\[
\mathrm{Hom}_{\mathrm{Ens}}(\pi_0(X), I) \simeq \mathrm{Hom}_W(X, I_W). \tag*{(2.8.2)}
\]
\textit{Le foncteur $\pi_0$ est également} adjoint à droite \textit{du} même \textit{foncteur $I \longmapsto I_W$,} i.e. \textit{on a un bijection fonctorielle} % [sic] "un bijection" -> "une bijection", accord
\[
\mathrm{Hom}_{\mathrm{Ens}}(I, \pi_0(X)) \simeq \mathrm{Hom}_W(I_W, X). \tag*{(2.8.3)}
\]

Prenant $I = e$ (ensemble ponctuel), on trouve le :

\textbf{Corollaire 2.9.} \textit{On a l'isomorphisme fonctoriel en $X$ dans} $\mathrm{Hot}_W$
\[
\boxed{\mathrm{Hom}_W(e_W, X) \simeq \pi_0(X).} \tag*{(2.8.4)}
\]

\textbf{Commentaire 2.10.} La formule d'adjonction (2.8.2) est d'ailleurs familière, chaque fois qu'on a une catégorie $\mathcal{H}$ (telle Cat, où $\mathrm{Hot}_W$ \textit{etc.}) % [sic] "où" -> "ou", conjonction de coordination
munie d'un foncteur $\pi_0$, $\pi_0(X)$ correspondant à l'intuition d'une décomposition d'un objet en “composantes connexes”. (Je ne vais pas expliciter ici les conditions générales sous lesquelles on définit un tel foncteur.) Par contre, l'adjonction en sens inverse de ces deux mêmes foncteurs est un phénomène tout à fait spécial à des catégories \msp{page 16}du type $\mathrm{Hot}_W$. Si on tient compte de la description de $I_W$ comme $I \times e_W$, la formule d'adjonction inhabituelle (2.8.3) équivaut, bien sûr, à la formule (2.8.4) - elle s'en déduit en “élevant à la puissance $I$”. On peut renverser la vapeur, et dans une catégorie générale $\mathcal{H}$, \textit{définir} un foncteur $\pi_0 : \mathcal{H} \longrightarrow \mathrm{Ens}$ par le formule (2.8.4), % [sic] "par le formule" -> "par la formule", accord
\[
\pi_0(X) \overset{\text{déf}}{=} \mathrm{Hom}(e_{\mathcal{H}}, X) \qquad (= \Gamma_{\mathcal{H}}(X), \text{ avec une notation familière}),
\]
où $e_{\mathcal{H}}$ est un objet final de $\mathcal{H}$. On aura donc, en posant $I_{\mathcal{H}} = I \times e_{\mathcal{H}}$, la formule
\[
\mathrm{Hom}(I_{\mathcal{H}}, X) \xrightarrow{\ \sim\ } \mathrm{Hom}(I, \pi_0(X))
\]
(\textit{i.e.} la formule (2.8.3)), disant que l'on a un couple $(I_{\mathcal{H}}, \pi_0)$ de foncteurs adjoints. Supposons d'autre part que pour tout ensemble $I$, l'application canonique d'adjonction
\[
I \longrightarrow \pi_0(I_{\mathcal{H}}) = \mathrm{Hom}_{\mathcal{H}}(e_{\mathcal{H}}, e_{\mathcal{H}} \times I)
\]
soit \textit{bijective}\footnote{ça a tendance à être ultrafaux si $\mathcal{H}$ est additive. Par contre, vrai si $\mathcal{H}$ est une catégorie à sommes distributives par rapport aux produits (voire universelles), dont l'objet final est \textit{connexe} \ldots}, \textit{i.e.} que $I \longmapsto I_{\mathcal{H}}$ soit \textit{pleinement fidèle.} On définit alors, pour tout $I$, une application canonique
% foot: Version 18 mars 2019


\endgroup
% END UNIT der-ch16/der-ch16_p006-010.tex
% BEGIN UNIT der-ch16/der-ch16_p011-015.tex
\begingroup\makeatletter%
\newcommand{\printedpage}[1]{\marginpar{\footnotesize p.~#1}}%
\newcommand{\msp}[1]{\marginpar{\footnotesize [#1]}}%
\newcommand{\Hom}{\operatorname{Hom}}%
\newcommand{\Ob}{\operatorname{Ob}}%
\newcommand{\Fl}{\operatorname{Fl}}%
\newcommand{\Hot}{\mathrm{Hot}}%
\newcommand{\Cat}{\mathrm{Cat}}%
\newcommand{\Ens}{\mathrm{Ens}}%
\newcommand{\id}{\mathrm{id}}%
\makeatletter%
\renewcommand\@makefnmark{\hbox{\@textsuperscript{\normalfont(\@thefnmark)}}}%
\renewcommand\@makefntext[1]{\parindent 1em\noindent\hb@xt@1.8em{\hss\@textsuperscript{\normalfont\@thefnmark}}#1}%
\makeatother%
\makeatother\let\printedpage\derpp\let\msp\dermsp%
\let\mspin\DERmspin%


\printedpage{9}
% head: A. Grothendieck : Dérivateurs, Ch. XVI, Édition des Universités de Montpellier II et Paris VII
% foot: Version 18 mars 2019

\[
\begin{tikzcd}[column sep=small]
\Hom(X,I_{\mathcal{H}}) \arrow[rr] \arrow[dr,dotted,"\pi_0^{X,I_{\mathcal{H}}}"'] & & \Hom(\pi_0(X),I) \arrow[dl,dotted,"\sim"'] \\ % ?? uncertain: label of the dotted right arrow, a tilde printed along the arrow -- \sim
& \Hom(\pi_0(X),\pi_0(I_{\mathcal{H}})). &
\end{tikzcd}
\]

La “bi-adjonction” des foncteurs $\pi_0$ et $I\longmapsto I_{\mathcal{H}}$ revient à dire que cette flèche est bijective,
\msp{17}
ou encore, que la flèche $\pi_0^{X,I_{\mathcal{H}}}$ :
\[
\Hom(X,I_{\mathcal{H}})\longrightarrow\Hom(\Gamma_{\mathcal{H}}(X),\Gamma_{\mathcal{H}}(I_{\mathcal{H}})\simeq I)
\]
soit bijective.

\textbf{Corollaire 2.11.} \textit{Le foncteur} $I\longmapsto I_W$ \textit{de} Ens \textit{dans} $\Hot_W$ \textit{est pleinement fidèle,} i.e. \textit{l'homomorphisme d'adjonction}
\[
\pi_0(I_W)\longrightarrow I
\]
\textit{est bijectif.} (\textit{Où encore,} % [sic] Où encore -> Ou encore, the accent makes no sense here
\textit{$\pi_0$ est un foncteur le localisation.}) % [sic] le localisation -> de localisation, wrong word

En effet, cette bijectivité est évidente quand on travaille dans Cat, et s'en déduit aussitôt dans $\Hot_W$.

\textsc{Démonstration de} 2.8. La formule d'adjonction (2.8.2) résulte de la formule correspondante pour Cat, évidente, et du lemme 2.6 sur les foncteurs adjoints entre catégories de modèles. L'adjonction dans l'autre sens se réduit à la formule (2.8.4). On a des flèches canoniques
\[
\begin{tikzcd}
\Hom_W(e_W,X) \arrow[r,shift left,"\alpha"] & \pi_0(X), \arrow[l,shift left,"\beta"]
\end{tikzcd}
\]
où $\alpha$ est $\pi_0^{e_W,X}$, compte tenu que $\pi_0(e_W)=\pi_0(e_{\Cat})=e_{\Ens}$, et $\Hom(e_{\Ens},\pi_0(X))\simeq\pi_0(X)$. La flèche $\beta$ est la flèche composée
\[
\Ob X\xrightarrow{\ \sim\ }\Hom_{\Cat}(e_{\Cat},X)\xrightarrow{\mathrm{hot}_W}\Hom_W(e_W,X)
\]
\msp{18}
en notant que cette flèche passe au quotient en $\pi_0(X)\longrightarrow\Hom_W(e_W,X)$, vu que deux applications homotopes de $e_{\Cat}$ dans $X$ sont égales dans $\Hot_W$. Je dis que
\[
\beta\alpha=\id\ ,\qquad\qquad \alpha\beta=\id\ .
\]
La relation $\alpha\beta=\id$ revient à dire que la flèche $i_x : e_{\Cat}\longrightarrow X$ définie par $x\in\Ob X$, donne $\pi_0(i_x) : e_{\Ens}\longrightarrow\pi_0(X)$ dont la valeur est la classe de $x$ dans $\pi_0(X)$, c'est tautologique. Pour prouver que $\beta\alpha=\id$, il suffit alors de prouver que $\beta$ est surjective, \textit{i.e.} le

\textbf{Corollaire 2.12.} \textit{Pour tout $X$ dans} Cat, \textit{l'application}
\[
\underbrace{\Hom_{\Cat}(e_{\Cat},X)}_{\simeq\Ob X}\longrightarrow\Hom_W(e_W,X)
\]
\printedpage{10}
% head: A. Grothendieck : Dérivateurs, Ch. XVI, Édition des Universités de Montpellier II et Paris VII
% foot: Version 18 mars 2019
\textit{est} surjective (\textit{moyennant l'axiome de connexité} Loc (5)).

\textsc{Démonstration de} 2.12. Si $X=\emptyset$, l'assertion revient à dire que $\Hom_W(e_W,\emptyset)=\emptyset$, mais plus généralement, on a le

\textbf{Corollaire 2.13.} (Moyennant Loc (5).) \textit{Soit $X$ non vide dans} Cat, \textit{alors} $\Hom_W(X,\emptyset)=\emptyset$. (\textit{Donc $\emptyset$ est un objet initial} strict \textit{de} $\Hot_W$.)

En effet, si on avait $X\longrightarrow\emptyset$ dans $\Hot_W$,
\msp{page 13'}\texttt{[il y a un décalage dans la numérotation]}
on en déduirait $\pi_0(X)\longrightarrow\emptyset$, ou $\pi_0(X)=\emptyset$, absurde\,!

Reste le cas où $X$ est non vide. Cela résulte du fait général :

\textbf{Lemme 2.14.} \textit{Soit $\mathcal{M}$ une catégorie ayant un objet initial} strict (\textit{si} $\Hom_{\mathcal{M}}(X,\emptyset_{\mathcal{M}})\neq\emptyset$, \textit{on a} $X=\emptyset_{\mathcal{M}}$), \textit{un objet final $e$, et telle que pour tout $X\neq\emptyset_{\mathcal{M}}$, $X\longrightarrow e$ ait une section,} i.e. $\Hom(e,X)\neq\emptyset$. \textit{Soit} $W\subset\Fl\mathcal{M}$
\begin{itemize}[leftmargin=3em,labelwidth=2.5em,align=right]
\item[1$^{\circ}$)] \textit{ne contenant pas de flèches} $\emptyset_{\mathcal{M}}\longrightarrow X$ \textit{avec} $X\neq\emptyset_{\mathcal{M}}$, \textit{et}
\item[2\,$^{\circ}$)] \textit{tel que pour tout diagramme}
\begin{tikzcd}[baseline=(\tikzcdmatrixname-1-1.base)]
X_1 & X_2 \arrow[l,"u_2"'] \\
e \arrow[u,"u_1"] \arrow[ur,dotted,"s"'] &
\end{tikzcd}
\textit{avec} $u_2\in W$ $\exists$ $s$ dans $\mathcal{M}$ \textit{tel que} $u_2s=u_1$

\textit{dans} $\mathcal{M}W^{-1}$.
\end{itemize}
\textit{Alors}
\begin{itemize}[leftmargin=3em,labelwidth=2.5em,align=right]
\item[a)] $\emptyset_{\mathcal{M}}$ \textit{est un objet initial} strict \textit{de} $\mathcal{M}W^{-1}$.
\item[b)] \textit{Pour tout} $X\neq\emptyset_{\mathcal{M}}$, $\Hom_{\mathcal{M}}(e,X)\longrightarrow\Hom_W(e,X)$ \textit{est surjectif.}
\end{itemize}

\textsc{Démonstration}. Que $\emptyset_{\mathcal{M}}$ soit objet initial de $\mathcal{M}W^{-1}$ a été vu (dernière assertion du lemme 2.5, forme duale), prouvons qu'il est strict, \textit{i.e.} que si $\Hom_W(X,\emptyset_{\mathcal{M}})\neq\emptyset$, on a $X=\emptyset_{\mathcal{M}}$. En effet, si on avait $u : X\longrightarrow\emptyset_{\mathcal{M}}$ dans $\mathcal{M}W^{-1}$, $u$ serait représenté par
\[
X=X_0\xleftarrow{\ W\ }X_1\longrightarrow X_2\xleftarrow{\ W\ }\cdots X_{2n-2}\longrightarrow X_{2n-1}\xleftarrow{\ W\ }X_{2n}=\emptyset_{\mathcal{M}}.
\]
On procède par récurrence sur la demi-longueur de cette chaîne pour voir que $X=\emptyset_{\mathcal{M}}$. En effet, l'hypothèse sur $W$ implique que $X_{2n-1}=\emptyset_{\mathcal{M}}$, puis celle que $\emptyset_{\mathcal{M}}$ est strict implique
$X_{2n-1}=\emptyset_{\mathcal{M}}$, d'où la conclusion. % [sic] X_{2n-1} -> X_{2n-2}, the strictness of the initial object gives X_{2n-2}

\msp{page 14'}
Prouvons la surjectivité de $\Hom_{\mathcal{M}}(e,X)\longrightarrow\Hom_W(e,X)$ pour $X\neq\emptyset_{\mathcal{M}}$, en montrant que tout chemin réduit de longueur 2n+1 est “homotope” à un chemin réduit de longueur $2n-1$ :
\[
e=X_0\longrightarrow X_1\xleftarrow{\ W\ }X_2\longrightarrow X_3\longleftarrow\ \cdots\ \longleftarrow X_{2n}\longrightarrow X_{2n+1}=X.
\]
En effet, on peut remplacer la section
\[
\begin{tikzcd}
e=X_0 \arrow[r,"u_1"'] \arrow[rr,dotted,bend right,"s"'] & X_1 & X_2 \arrow[l,"W"',"u_2"] \arrow[r,"u_3"'] & X_3 & X_4 \arrow[l,"W"']
\end{tikzcd}
\]
\printedpage{11}
% head: A. Grothendieck : Dérivateurs, Ch. XVI, Édition des Universités de Montpellier II et Paris VII
% foot: Version 18 mars 2019
par
\[
\begin{tikzcd}
e \arrow[r,"u_3s"'] & X_3 & \cdots, \arrow[l,"W"']
\end{tikzcd}
\]
où $s$ est une section de $X_2$ sur $e$ telle que
\[
u_2s=u_1\qquad\quad\textit{dans}\qquad\quad\mathcal{M}W^{-1}\ .
\]
Cela prouve donc par récurrence sur $n$ que le chemin est homotope à un chemin de longueur 1, $e\longrightarrow X$, q.e.d.

Il faut donc, dans le cas actuel, vérifier 1$^{\circ}$) (c'est essentiellement le corollaire 2.12.), et 2$^{\circ}$)
\[
\begin{tikzcd}
X_1 & X_2 \arrow[l,"u_2\in W"'] & \\
e \arrow[u,"u_1"] \arrow[ur,dotted,"s"'] & & .
\end{tikzcd}
\]
Mais comme $u_2\in W$, on a $\pi_0(u_2)$ isomorphisme, donc
\msp{page 15'}
surjectif, donc $\exists$ $s$ tel que $u_2s$ soit homotope à $u_1$ dans Cat, ce qui implique $u_2s=u_1$ dans $\Hot_W$, q.e.d.

\textbf{Commentaire final.} L'axiome de connexité Loc (5) diffère de tous les autres axiomes Loc (1) à Loc (8) (mais à l'exception de Loc (8 bis) d'accessibilité) par le fait que ce n'est pas un axiome de stabilité. Les autres axiomes Loc (1) à Loc (8) (et aussi Loc (8 bis) d'accessibilité et même les axiomes ultérieures % [sic] ultérieures -> ultérieurs, agreement with axiomes
d'exactitude~\footnote[7]{il n'y en aura en fait pas besoin, grâce à la providentielle théorie de K. S. \textsc{Brown}\,!}) sont satisfaits si $W=\Fl\mathcal{M}$ - l'axiome de connexité est donc le seul qui s'oppose à la \textit{trivialité} de la théorie homotopique définie par $W$. C'est dire qu'il ne peut être conséquence de l'un des autres. Par contre, l'axiome Loc (4) des sommes directes n'est pas indépendant des autres. Signalons à ce propos (rappel) :
\[
\left\{
\begin{array}{rl}
\text{Loc (4 a$_0$)} & \text{est conséquence de l'axiome Loc (6 a) du carré cocartésien (forme directe).}\\[1ex]
\text{Loc (4 a)} & \text{est conséquence de Loc (6 a) précédent}\\
 & \text{+ l'axiome des limites inductives filtrantes Loc (8).}\\[3ex]
\text{Loc (4 b)} & \text{est conséquence de l'axiome Loc (7) des fibrations. Donc}\\[1ex]
\text{Loc (4)} & \text{est conséquence de Loc (6 a, 7, 8).}
\end{array}
\right.
\]

\printedpage{12}
% head: A. Grothendieck : Dérivateurs, Ch. XVI, Édition des Universités de Montpellier II et Paris VII
% foot: Version 18 mars 2019
\msp{page 16'}

\section*{3\quad Axiome Loc (6) du carré cocartésien (ou de “Mayer-Vietoris”, ou des $W$-cofibrations)~\footnote[8]{C'est le nom “\textit{axiome des $W$-cofibrations}” qui m'apparaît à présent comme le meilleur, -- c'est ce nom qui me paraît le mieux saisir la signification conceptuelle et non technique de l'axiome.}}

Je vois plusieurs variantes, et n'ai pas encore l'esprit totalement net sur les implications mutuelles. Je commence par la variante plus compliquée, mais dont on aura besoin en tous cas :

\begin{itemize}[leftmargin=5em,labelwidth=4em,align=left]
\item[\textbf{Loc 6}] (Axiome des $W$-cofibrations, ou lemme des 5, ou axiome des homomorphismes de carrés cocartésiens). Considérons dans Cat un diagramme commutatif cubique
\[
\begin{tikzcd}[row sep=small,column sep=small]
X_0 \arrow[rr] \arrow[dr] \arrow[dd,"f_0"'] & & X_2 \arrow[dr] \arrow[dd,"f_2" near start] & \\
& X_1 \arrow[rr,crossing over] & & X_3 \arrow[dd,"f_3"] \\
Y_0 \arrow[rr] \arrow[dr] & & Y_2 \arrow[dr] & \\
& Y_1 \arrow[rr] \arrow[from=uu,crossing over,"f_1"' near start] & & Y_3,
\end{tikzcd}
\]
\end{itemize}
homomorphisme $f=(f_0,f_1,f_2,f_3)$ d'un diagramme carré $X$ dans un autre $Y$. On suppose que $X$ provient d'un objet $X_3$ de Cat de deux ouverts $X_1$, $X_2$ de $X_3$ \textit{recouvrant} $X_3$, et de leur intersection $X_0=X_1\cap X_2$, et de même pour $Y$. (Appelons ces diagrammes dans Cat des diagrammes MV (“\textit{ouverts}”) - “\textit{Mayer-Vietoris}”.) On a alors ce qui suit~\footnote[9]{Il est bon d'introduire des propriétés du diagramme
\begin{itemize}[leftmargin=3em,labelwidth=2.5em,align=right]
\item[1$^{\circ}$)] $f_0\in W$.
\item[2$^{\circ}$)] $f_1$ et $f_2$ dans $W$.
\item[3$^{\circ}$)] $f_3$ dans $W$.
\end{itemize}
et de formuler les parties a) b) c) de Loc (6) ainsi
\begin{itemize}[leftmargin=3em,labelwidth=2.5em,align=right]
\item[a)] 1$^{\circ}$ + 2$^{\circ}$ $\Longrightarrow$ 3$^{\circ}$.
\item[b)] 2$^{\circ}$ + 3$^{\circ}$ $\Longrightarrow$ 1$^{\circ}$.
\item[c)] 1$^{\circ}$ + 3$^{\circ}$ $\Longrightarrow$ 2$^{\circ}$.
\end{itemize}} :
\begin{itemize}[leftmargin=3em,labelwidth=2.5em,align=right]
\item[a)] $f_0,f_1,f_2\in W\Longrightarrow f_3\in W$~\footnote[10]{N.B. il semble que dans Loc (6 a), il suffit de l'énoncé en supposant $f_0$, $f_2$ des \textit{isomorphismes,} \textit{cf.} p. 22.}.
\item[b)] $f_1,f_2,f_3\in W\Longrightarrow f_0\in W$.
\item[c)] $f_0,f_3\in W\Longrightarrow f_1,f_2\in W$~\footnote[11]{Il ne faudrait pas inclure c dans Loc 6 \textit{cf.} p. 26.}.
\end{itemize}
\noindent\footnote[12]{N.B. On verra plus bas (3.3) que Loc 6 c implique déjà Loc (6 a, b), \textit{i.e.} Loc (6) tout entier. Et si $W$ satisfait Loc (1 bis), \textit{i.e.} stable par facteurs directs, alors Loc (6 a, b) $\Longrightarrow$ Loc (6 c).}.

\printedpage{13}
% head: A. Grothendieck : Dérivateurs, Ch. XVI, Édition des Universités de Montpellier II et Paris VII
% foot: Version 18 mars 2019
\msp{page 17'}
\noindent\footnote[13]{Je vais désigner par Loc (6) l'axiome Loc (6 a, b), sans y inclure c.}.

\textbf{N.B.} Je distingue ces conditions comme Loc (6 a), Loc (6 b), Loc (6 c). Je présume qu'on n'aura jamais à faire usage à Loc (6 b), sans supposer également Loc (6 a), mais l'inverse n'est sûrement pas le cas. Ainsi, $W_0=\{f\in\Fl\Cat\mid\pi_0(f)\text{ bijectif}\}$ satisfait a), mais non b). Je présume (sans avoir vérifié) que c'est pareil pour les $W_n$, $n\in\mathbf{N}$, de la “\textit{$n$-équivalence homotopique}”.

Il se pourrait que Loc (6 a) et Loc (6 a, b) soient équivalentes à la forme plus faible, et plus simple, suivante. (Si oui, c'est sous la forme la plus simple, bien sûr, qu'il convient de les énoncer.)

\begin{itemize}[leftmargin=5em,labelwidth=4em,align=left]
\item[\textbf{Loc 6$'$}] (Axiome des carrés cocartésiens, ou de Mayer-Vietoris). Soit
\[
\begin{tikzcd}[row sep=small,column sep=small]
X_0 \arrow[rr,"i_2"] \arrow[dr,"i_1"] & & X_2 \arrow[dr,"j_1"] & \\
& X_1 \arrow[rr,"j_2"] & & X_3
\end{tikzcd}
\]
un diagramme MV dans Cat. Alors
\begin{itemize}[leftmargin=3em,labelwidth=2.5em,align=right]
\item[a)] Si $i_1\in W$, $j_1\in W$.
\item[b)] Si $j_1\in W$, $i_1\in W$.
\end{itemize}
\end{itemize}

\msp{page 18'}
\textbf{Proposition 3.1.}
\[
\begin{aligned}
\text{Loc (6 a)}\quad &\Longrightarrow\quad \text{Loc (6$'$ a),}\\
\text{Loc (6 b)}\quad &\Longrightarrow\quad \text{Loc (6$'$ b),}\\
\text{Loc (6 c)}\quad &\Longrightarrow\quad \text{Loc (6$'$ a, b)}\ \footnotemark[14].
\end{aligned}
\]
\footnotetext[14]{N.B. on a même Loc (6 c) $\Longrightarrow$ $\underbrace{\text{Loc 6 (a, b)}}_{\text{Loc (6)}}$.}

Ça se voit sur le diagramme cubique
\[
\text{(Q)}\qquad\qquad
\begin{tikzcd}[row sep=small,column sep=small]
X_0 \arrow[rr,"i_2"] \arrow[dr,"i_1"] & & X_2 \arrow[dr,"j_1"] \arrow[from=dd,"f_2=\id"' near start] & \\
& X_1 \arrow[rr,"j_2",crossing over] & & X_3 \\
X_0 \arrow[uu,"f_0=\id"] \arrow[rr,"i_2"] \arrow[dr,"\id"] & & X_2 \arrow[dr,"\id"] & \\
& X_0 \arrow[rr,"i_2"] \arrow[uu,crossing over,"f_1=i_1" near end] & & X_2, \arrow[uu,"f_3=j_1"']
\end{tikzcd}
\]
où les $f_0$, $f_1$, $f_2$, $f_3$ de l'énoncé de Loc (6) deviennent $\id_{X_0}$, $i_1$, $\id_{X_2}$, $j_1$. Loc (6 a) donne donc $i_1\in W\Longrightarrow j_1\in W$, \textit{i.e.} Loc (6$'$ a), et Loc (6 c) $j_1\in W\Longrightarrow i_1\in W$, \textit{i.e.} Loc (6$'$ b). Pour Loc (6 b) $\Longrightarrow$ Loc (6$'$ b), on utilise (\textit{cf.} XII, p. 78) le cube


\endgroup
% END UNIT der-ch16/der-ch16_p011-015.tex
% BEGIN UNIT der-ch16/der-ch16_p016-020.tex
\begingroup\makeatletter%
\newcommand{\printedpage}[1]{\marginpar{\footnotesize p.~#1}}%
\newcommand{\msp}[1]{\marginpar{\footnotesize [#1]}}%
\newcommand{\Cof}{\mathrm{Cof}}%
\newcommand{\Fl}{\mathrm{Fl}}%
\newcommand{\Mono}{\mathrm{Mono}}%
\newcommand{\univ}{\mathrm{univ}}%
\newcommand{\id}{\mathrm{id}}%
\makeatletter%
\renewcommand\@makefnmark{\hbox{\textsuperscript{\normalfont(\@thefnmark)}}}%
\renewcommand\@makefntext[1]{\parindent 1em\noindent\hbox{\textsuperscript{\normalfont\@thefnmark}}\,#1}%
\makeatother%
\makeatother\let\printedpage\derpp\let\msp\dermsp%
\let\mspin\DERmspin%


\noindent\printedpage{14}% head: A. Grothendieck : Dérivateurs, Ch. XVI, Édition des Universités de Montpellier II et Paris VII
% foot: Version 18 mars 2019
\[
\begin{tikzcd}[ampersand replacement=\&, row sep=small, column sep=small]
X_0 \arrow[rr, "i_2"] \arrow[dr, "i_1"] \arrow[dd, "i_1"'] \& \& X_2 \arrow[dr, "j_1"] \arrow[dd, "j_1" near start] \& \\
\& X_1 \arrow[rr, "j_2" near start, crossing over] \& \& X_3 \arrow[dd, "\id"] \\
X_1 \arrow[rr, "j_2" near end] \arrow[dr, "\id"'] \& \& X_3 \arrow[dr, "\id"'] \& \\
\& X_1 \arrow[rr, "j_2"] \arrow[from=2-2, "\id"' near start, crossing over] \& \& X_3,
\end{tikzcd}
\tag*{(Q$'$)} % checked against PDF leaf 16: id labels the front vertical X_1 to X_1, not the horizontal X_1 to X_3.
\]
où $f_0 = i_1$, $f_1 = \id_{X_1}$, $f_1 = j_1$, $f_3 = \id_{X_3}$, donc % [sic] $f_1 = j_1$ -> $f_2 = j_1$, $f_1$ is already given as $\id_{X_1}$ and $f_2$ is missing
\[
j_1 \in W \Longrightarrow i_1 \in W, \ \ \textit{i.e.}\ \text{Loc (6 b)} \Longrightarrow \text{Loc (6$'$ b)},
\]
et si on a Loc (6 c), on voit que $i_1 \in W \Longrightarrow j_1 \in W$, \textit{i.e.} Loc (6$'$ a).

\textbf{Commentaire 3.2.} Loc (6 a) m'apparaît comme le “minimum vital” pour établir l'existence d'une théorie des cofibres $W$-homotopique dans Cat. Pour cela, il faut en effet disposer d'un critère serviable assurant qu'une flèche $i : A \longrightarrow B$ dans Cat est dans $\Cof_W$.

\msp{page 19}On trouve, quand on dispose de Loc (6 a) (sous forme duale, \textit{i.e.} pour les \textit{fermés,} c'est kif-kif), que les \textit{immersions ouvertes de Dwyer} sont dans $\Cof_W$, ce qui suffit pour avoir \textit{une bonne théorie des cofibres $W$-homotopiques, dans laquelle les carrés MV sont $W$-cocartésiens} (chose a laquelle on tient absolument). % [sic] a laquelle -> à laquelle, missing accent

Quant à Loc (6 b), si je dispose à la place de Loc (6$'$ b) + Loc (6 a), je dois m'en tirer pour prouver Loc (6). Dans XII Appendice (pages 265 ff.) je prouve, paraît-il, que (avec les notations actuelles) Loc (6$'$) implique Loc (6), et même Loc (6 c) si $W$ est stable par facteurs directs (condition Loc (1 bis)). Mais la démonstration est incomplète, puisque j'admets tacitement que les immersions de Dwyer sont dans $\Cof_W$, chose que je n'avais établie que moyennant l'axiome W (7), en fait moyennant Loc (6 a) (ce qui est une partie de W (7)).

\msp{page 20}Mais ici me vient l'idée du lemme de K. S. \textsc{Brown}, cité par \textsc{Baues}, que je vais redonner ici :

\textbf{Lemme 3.2.1.} (K. S. \textsc{Brown}) \textit{Soit $(\mathcal{M}, W)$ un localiseur} ($W$ \textit{modérément saturé}), \textit{muni d'un ensemble $C \subset \Fl\,\mathcal{M}$ satisfaisant les conditions :}
\begin{enumerate}
\item[1)] \textit{$C$ stable par composition et par cochangement de base} (N.B. \textit{on n'a à supposer l'existence des sommes amalgamées}
$\begin{tikzcd}[ampersand replacement=\&, baseline=(\tikzcdmatrixname-1-1.base), row sep=small, column sep=small]
X_0 \arrow[rr, "i_1"] \arrow[dr, "i_2"'] \& \& X_1 \arrow[dr] \& \\
\& X_2 \arrow[rr] \& \& X_3
\end{tikzcd}$
\textit{que si $i_1 \in C$}), \textit{contient les isomorphismes.}
\item[2)] \textit{$\mathcal{M}$ contient un objet initial, $\emptyset_{\mathcal{M}}$, et pour tout objet $X$ de $\mathcal{M}$, $\emptyset_{\mathcal{M}} \longrightarrow X$ est dans $C$.} (N.B. (1) \textit{et} (2) \textit{impliquent que $\mathcal{M}$ est stable par sommes finies.})
\item[3)] \textit{Toute flèche $f$ de $\mathcal{M}$ se factorise en $pi$, avec $i \in C$, $p \in W$.}
\item[4)] \textit{$C \cap W$ est stable par cochangement de base,} i.e. \textit{si}
\[
\begin{tikzcd}[ampersand replacement=\&, row sep=small, column sep=small]
X_0 \arrow[rr, "i_1"] \arrow[dr, "i_2"'] \& \& X_1 \arrow[dr] \& \\
\& X_2 \arrow[rr, "j_1"'] \& \& X_3
\end{tikzcd}
\tag*{$(*)$}
\]
\printedpage{15}% head: A. Grothendieck : Dérivateurs, Ch. XVI, Édition des Universités de Montpellier II et Paris VII
% foot: Version 18 mars 2019
\textit{est cocartésien et $i_1 \in C \cap W$, alors $j_1 \in C \cap W$.} (N.B. \textit{$j_1 \in C$ est déjà contenu dans} (1), \textit{il reste à exiger $j_1 \in W$.})
\end{enumerate}
\textit{Sous ces conditions, on a $C \subset \Cof_W$.}

On peut aussi énoncer le lemme ainsi

\textbf{Corollaire 3.2.2.} \textit{Supposons que $C$ satisfait} (1), (2), (3). \textit{Alors}
\msp{page 21}\textit{les conditions suivantes sur $C$ sont équivalentes.}
\begin{enumerate}
\item[a)] \textit{Condition} (4) \textit{ci-dessus,} i.e. \textit{stabilité de $C \cap W$ par cochangement de base.}
\item[a$'$)] \textit{$C \cap W \subset W^{\univ}$.}
\item[b)] \textit{$C \subset \Cof_W$.}
\item[b$'$)] \textit{Pour tout carré cocartésien}
\[
\begin{tikzcd}[ampersand replacement=\&, row sep=small, column sep=small]
X_0 \arrow[rr, "i"] \arrow[dr, "p"'] \arrow[to=2-4, phantom, "\textit{coc.}" description] \& \& X_1 \arrow[dr, "q"] \& \\
\& X'_0 \arrow[rr, "j"] \& \& X'_1
\end{tikzcd}
\]
\textit{avec $i \in C$, $p \in W$, on a $q \in W$.}
\end{enumerate}
\textsc{Démonstration}. On a
\[
\begin{aligned}
&\text{b} \iff \text{b}'\\
&\text{a} \iff \text{a}'
\end{aligned}
\]
car $C$ est stable par cochangement de base. On a b $\Longrightarrow$ a$'$ car $\Cof_W \cap W = W^{\univ}$ grâce à l'hypothèse 3°). On a a $\Longrightarrow$ b$'$ par le lemme de Brown.

On va appliquer ceci au cas $\mathcal{M} = \mathrm{Cat}$, $C =$ immersions ouvertes de Dwyer. Les conditions 1), 2) ne dépendent pas de $W$, elles sont vérifiées, et 3) aussi dès que $W$ est un localiseur fondamental. Quant à 4), c'est une conséquence de l'ancien W (7 bis), \textit{i.e.} Loc (6$'$) (on a
\msp{page 22}besoin apparemment de Loc (6$'$ a \textit{et} b).) On en conclut $C \subset \Cof_W$ - c'est merveilleux! (Mais je ne sais pas prouver le lemme \dots)

Cela légitime la démonstration faite dans XII \textit{loc. cit.}, et on trouve la

\textbf{Proposition 3.3.} \textit{La condition} Loc (6$'$) (\textit{axiome du carré cocartésien, forme forte} (a) \textit{et} (b)) équivaut \textit{à} Loc (6 a, b) (\textit{axiome des homomorphismes des carrés cocartésiens}) \textit{et est} impliqué \textit{par} Loc (6 c). \textit{Si $W$ satisfait à} Loc (1 bis), i.e. \textit{est stable par facteurs directs, alors} Loc (6$'$) équivaut \textit{à} Loc (6 c) \textit{et à} Loc (6 a b c).

\textbf{Commentaire 3.4.} Ainsi, on trouve une théorie des cofibres $W$-homotopiques \textit{soit} sous l'hypothèse Loc (6$'$) (équivalente à Loc (6) plus compliquée), soit sous l'hypothèse plus faible (et plus compliquée) Loc (6 a).

% left margin: a vertical line is printed alongside the following paragraph
Dans le premier cas, c'est une théorie de nature très particulière, puisque les carrés $W$-cocartésiens sont “réversibles” (\textit{i.e.} on a non seulement $i_1 \in W \Longrightarrow j_1 \in W$, mais aussi l'implication inverse).

Il reste donc la question si Loc (6$'$ a) suffit pour une théorie des cofibres $W$-homotopiques, ou ce qui revient sûrement au même, si ça implique déjà Loc (6 a) (donc lui est équivalent). C'est entièrement suspendu à la \textit{question} :

\printedpage{16}% head: A. Grothendieck : Dérivateurs, Ch. XVI, Édition des Universités de Montpellier II et Paris VII
% foot: Version 18 mars 2019
\msp{page 23}% left margin: a question mark followed by a double vertical line alongside the first three lines of the following paragraph
\marginpar{\footnotesize ?\,$\Vert$}l'axiome Loc (6$'$ a) (axiome du carré cocartésien, forme directe) implique-t-il que l'ensemble des immersions ouvertes de Dwyer qui sont dans $W$ (immersions de Dwyer $W$-triviales) est stable par cochangement de base ? En fait, on voit qu'il suffirait pour cela de vérifier l'axiome du cube (ou des homomorphismes de carrés cocartésiens), forme directe Loc (6 a), dans le cas où $f_0$ et $f_2$ sont des identités, donc à prouver qu'alors $f_1 \in W \Longrightarrow f_3 \in W$ (ce qui donnerait directement $\mathrm{Dw} \subset \Cof_W$). Mais je ne vois aucun moyen de le prouver \dots

Supposons maintenant que la condition Loc (6 a) est vérifiée, donc qu'on dispose d'une “bonne théorie des cofibres $W$-homotopiques”, donc aussi d'une théorie des carrés $W$-cocartésiens (\textit{cf.} XII, p. 258-262). Il y a deux questions qui se posent, concernant la validité, avec ces hypothèses plus faibles, des développements de XII :
\begin{enumerate}
\item[1°)] Soit un carré commutatif
\[
\begin{tikzcd}[ampersand replacement=\&, row sep=small, column sep=small]
X_0 \arrow[rr] \arrow[dr] \& \& X_2 \arrow[dr] \& \\
\& X_1 \arrow[rr] \& \& X_3
\end{tikzcd}
\]
dans Cat, est-il vrai que ce carré est $W$-cocartésien si et seulement si la flèche
\[
\text{canonique}\ \displaystyle\int
\begin{tikzcd}[ampersand replacement=\&, row sep=small, column sep=small]
\& X_1 \\
X_0 \arrow[ur] \arrow[dr] \& \\
\& X_2
\end{tikzcd}
\quad \longrightarrow \ X_3 \text{ est dans } W \text{ ? }\text{\footnotemark[15]}
\]
\footnotetext[15]{Mieux, dans tous les cas, quand on a un diagramme
$\begin{tikzcd}[ampersand replacement=\&, baseline=(\tikzcdmatrixname-2-1.base), row sep=tiny, column sep=small]
\& X_1 \\
X_0 \arrow[ur] \arrow[dr] \& \\
\& X_2
\end{tikzcd}$, son intégrale $\displaystyle\int$ est une somme amalgamée $W$-homotopique dudit diagramme (moyennant Loc (6 a) sans plus, et même seulement Loc (6$'$ a)).}
\end{enumerate}
\msp{page 24}
\begin{enumerate}
\item[2°)] Est-il encore vrai que si $A$ est dans Cat, et $u : F \lhook\joinrel\longrightarrow G$ un homomorphisme dans $A^{\wedge}$, alors $u \in \Cof_{W_A}$, et si $u \in W$, $u \in W_A^{\univ}$, donc a-t-on
\[
\Mono_{A^{\wedge}} \subset \Cof_{W_A}, \qquad \Mono_{A^{\wedge}} \cap W_A \subset W_A^{\univ} \text{ ? }\text{\footnotemark[16]}
\]
\footnotetext[16]{N.B. Loc (6$'$ a) suffit au lieu de ( Loc (6 a)).}
\end{enumerate}
Mais 1°) est prouvé dans XII, prop. 20, p. 262 - je viens de vérifier que la démonstration est valable dans les hypothèses actuelles \footnote[17]{Mais elle utilise de plus Loc (3 bis). Mais avec une autre démonstration plus simple que je viens de rajouter (p 264), on peut s'en passer (et on obtient une vison plus complète de la situation).}, pourvu que je vérifie que le corollaire 3, p. 149 dans XII est valable, \textit{i.e.} % [sic] (footnote 17) vison -> vision, missing letter
\begin{enumerate}
\item[1° bis)] Si
$\begin{tikzcd}[ampersand replacement=\&, baseline=(\tikzcdmatrixname-1-1.base)]
X_0 \arrow[r] \arrow[d] \& X_1 \arrow[d] \\
X_2 \arrow[r] \& X
\end{tikzcd}$
est un carré MV dans Cat, alors
$\displaystyle\int
\begin{tikzcd}[ampersand replacement=\&, row sep=small, column sep=small]
\& X_1 \\
X_0 \arrow[ur] \arrow[dr] \& \\
\& X_2
\end{tikzcd}
\quad \longrightarrow \quad X$ est dans $W$ (donc, moyennant 1°), le carré est $W$-cocartésien).
\end{enumerate}

\printedpage{17}% head: A. Grothendieck : Dérivateurs, Ch. XVI, Édition des Universités de Montpellier II et Paris VII
% foot: Version 18 mars 2019
Pour le vérifier, je suis amené à vérifier que le théorème 4 dans loc. cit. (p. 146) est encore valable.

Je viens de le faire, et ai annoté en conséquence le théorème 4 et ses corollaires (certains résultats (cor. 4 et cor. 5) utilisant W (7 bis) dans toute sa force, \textit{i.e.} Loc (6$'$), \textit{i.e.} Loc (6 a \textit{et} b). Mais le corollaire 3 du théorème 4 est O.K.

Je viens aussi de vérifier la validité de 2°), en vérifiant la démonstration du corollaire du théorème 5 (XII, p. 168) et en l'annotant en conséquence.

\msp{page 25}3.5. Résumons : on a
\[
\text{Loc (6 c)} \ \Longrightarrow \ \text{Loc (6)} \ \iff \ \text{Loc (6$'$)}
\tag*{(3.5.1)}
\]
et
\[
\text{Si } \ W = W^{\natural}\,, \ \text{ on a } \ \text{Loc (6 c)} \ \iff \ \text{Loc (6)} \ \iff \ \text{Loc (6$'$)}.
\tag*{(3.5.2)}
\]
D'autre part
\[
\text{Loc (6 a)} \ \Longrightarrow \ \text{Loc (6$'$ a)}\,,
\tag*{(3.5.3)}
\]
et j'ignore si l'implication inverse
\[
\text{Loc (6$'$ a)} \ \overset{?}{\Longrightarrow} \ \text{Loc (6 a)}
\]
est valable \footnote[18]{C'est le cas en présence de Loc (8, 9).}. Pour établir Loc (6 a), il suffit de le faire en tous cas quand $f_0$ et $f_2$ sont des isomorphismes.

L'axiome Loc (6 a) semble le “minimum vital” pour que $W$ donne lieu à une “bonne théorie des cofibres $W$-homotopiques”, donc aussi des carrés $W$-cocartésiens. De façon plus précise, j'ai envie de prouver ceci :

\textbf{Proposition 3.6.} \textit{L'axiome} Loc (6 a) \textit{équivaut à la conjonction des deux propriétés suivantes.}

\printedpage{18}% head: A. Grothendieck : Dérivateurs, Ch. XVI, Édition des Universités de Montpellier II et Paris VII
% foot: Version 18 mars 2019
\begin{itemize}[leftmargin=7em, labelwidth=6em]
\item[\textbf{Loc (6 A)}] Cat a “assez de $W$-cofibrations”,
\msp{page 26}\textit{i.e.} toute flèche $f$ de Cat se factorise en $f = pi$, avec $i \in \Cof_W$ et $p \in W$. (Cela signifie aussi que Cat muni de $(W, \Cof_W)$ est une catégorie à cofibrations de K. S. \textsc{Brown} - \textsc{Anderson}, et tous les résultats de la théorie de \textsc{Brown-Quillen} sont applicables. En particulier, on a une bonne théorie des carrés $W$-cocartésiens.)
\item[\textbf{Loc (6 B)}] Si $X$ dans Cat est réunion de deux ouverts $U$, $V$, d'intersection $S$, alors le carré
\[
\begin{tikzcd}[ampersand replacement=\&, row sep=large, column sep=large]
S \arrow[r] \arrow[d] \& U \arrow[d] \\
V \arrow[r] \& X
\end{tikzcd}
\]
est $W$-cocartésien \footnote[19]{En fait, Loc (6 B) est conséquence de Loc (6 A), laquelle donc est \textit{équivalente} à Loc (6 a). (\textit{Cf.} esquisse de démonstration dans 4.4.15, pages 66-68.)}.
\end{itemize}
D'autre part, la condition Loc (6 A) peut être remplacée par la condition (en apparence plus faible)
\begin{itemize}[leftmargin=7em, labelwidth=6em]
\item[\textbf{Loc (6 A$'$)}] Tout morphisme de Dwyer est dans $\Cof_W$.
\end{itemize}
\textbf{Corollaire 3.7.} \textit{La condition} Loc (6), i.e. Loc (6 a, b) (\textit{ou encore la condition équivalente} Loc (6$'$)) \textit{équivaut à la conjonction des conditions} Loc (6 A) (\textit{ou} Loc (6 A$'$)) \textit{et} Loc (6 B) \textit{ci-dessus}), \textit{et de la condition}
\begin{itemize}[leftmargin=7em, labelwidth=6em]
\item[\textbf{Loc (6 C)}] Pour tout carré $W$-cocartésien
\[
\begin{tikzcd}[ampersand replacement=\&, row sep=large]
X_0 \arrow[r, "i"] \arrow[d] \& X_1 \arrow[d] \\
X_2 \arrow[r, "j"] \& X
\end{tikzcd}
\]
dans Cat, $j \in W \Longrightarrow i \in W$.
\end{itemize}
\msp{page 27}D'autre part, je soupçonne que Loc (6 c) n'est guère vérifiée que si $W = W^{\natural}$ (cas qui est le plus commun bien sûr), auquel cas Loc (6 c) \textit{équivaut} à Loc (6). Il n'y a sans doute pas lieu de l'inclure dans Loc (6), pas plus que $W = W^{\natural}$ dans Loc (1).


\endgroup
% END UNIT der-ch16/der-ch16_p016-020.tex
% BEGIN UNIT der-ch16/der-ch16_p021-025.tex
\begingroup\makeatletter%
\newcommand{\printedpage}[1]{\marginpar{\footnotesize p.~#1}}%
\newcommand{\msp}[1]{\marginpar{\footnotesize [#1]}}%
\newcommand{\Fib}{\mathrm{Fib}}%
\newcommand{\Fibf}{\mathrm{Fibf}}%
\newcommand{\Parf}{\mathrm{Parf}}%
\newcommand{\Catfib}{\mathrm{Catfib}}%
\newcommand{\Fibsc}{\mathrm{Fibsc}}%
\newcommand{\Cat}{\mathrm{Cat}}%
\newcommand{\HOT}{\mathrm{HOT}}%
\makeatother\let\printedpage\derpp\let\msp\dermsp%
\let\mspin\DERmspin%


\printedpage{19}
% head: A. Grothendieck : Dérivateurs, Ch. XVI, Édition des Universités de Montpellier II et Paris VII

\section*{4\quad Axiome Loc (7) des $W$-fibrations}

4.1. Rappelons la notion de \textit{$W$-fibration faible} (XII p. 51, 52) : la version la plus forte des 8 conditions équivalentes pour que $f : X \longrightarrow S$ soit une fibration faible est celle-ci.

\begin{itemize}[leftmargin=2cm]
\item[\textbf{(Ff)}] $\forall\ S'' \longrightarrow S'$ dans $\Cat/S$, avec $S'$, $S''$ \textit{contractiles,} le morphisme correspondant $X'' = X \times_S S'' \longrightarrow X' = X \times_S S'$ est dans $W$.
\end{itemize}

Quand $W$ est fortement saturé, alors ceci implique la condition, en apparence plus forte :

\begin{itemize}[leftmargin=2cm]
\item[\textbf{(Ff$'$)}] $\forall\ S'' \longrightarrow S'$ dans $\Cat/S$, telle que $S'' \longrightarrow S'$ soit un homotopisme, $X'' \longrightarrow X'$ est dans $W$.
\end{itemize}

Je désigne par $\Fibf_W$ l'ensemble des fibrations faibles. Je résume dans un diagramme d'inclusions les principales classes de flèches dans Cat que nous avons rencontrées, s'apparentant à des classes de “fibrations” (pour $W$, dans un sens convenable) :

\msp{page 28}

\[
\begin{tikzcd}
 & & \Fib_W \arrow[d] & & \\
\Parf_\omega \arrow[r] & \Parf_W \arrow[r] & \Fibf_W & \underbrace{\Catfib^W}_{\mathcal{F}_W} \arrow[l] & \underbrace{\Catfib^\omega}_{\mathcal{F}_\omega} \arrow[l] \ ,
\end{tikzcd}
\]
où $\Fib_W$ désigne les $W$-fibrations ($f : S' \longrightarrow S$ est dans $\Fib_W$ si et seulement si le foncteur changement de base $X \longmapsto X' = X \times_S S'$ de $\Cat/S$ dans $\Cat/S'$ est compatible avec les localiseurs $W_S$ et $W_{S'}$), Catfib les Cat-fibrations, $\Catfib^W$ (resp. $\Catfib^\omega$) les Cat-fibrations telles que les foncteurs changement de base soient dans $W$ (resp. soient dans $W_\omega$, \textit{cf.} déf. plus bas) \footnote[20]{N.B. $\Catfib_0^W = \Catfib \cap \Fibf_W$, $\Catfib_0^\omega = \Catfib \cap \Fibf_{W_\omega}$}, $\Parf_W$ les morphismes $W$-parfaits, \textit{i.e.} $W$-lisses et $W$-propres, enfin $\Parf_\omega$ les morphismes $W_\omega$-parfaits, où $W_\omega$ est le plus petit localiseur satisfaisant Loc (1, 2, 3). Parmi les espèces de “fibrations” envisagées, c'est cette dernière seulement, avec $\Catfib^\omega$, qui est indépendante de $W$.

À ce titre, elle me paraît jouer un rôle similaire, dual, à celui des morphismes de Dwyer, et je présume qu'on peut leur faire jouer un rôle similaire. En tous \msp{page 29} cas, on sait que toute flèche $f$ de Cat se factorise en $pi$, avec $i \in W_\omega$ et $p \in \Parf_\omega$, donc $i \in W$ pour tout $W$ satisfaisant Loc (1, 2, 3). On sait aussi que $\Catfib^\omega$ satisfait l'énoncé de factorisation forte, $f$ se factorise en $f = pi$, $i \in W_\omega$, $p \in \Catfib_0^\omega$. \footnote[21]{Donc la propriété de factorisation est valable pour toutes ces classes, sauf peut-être $\Fib_W$.}

Notons aussi que $\Parf_\omega$ et $\Parf_W$, ainsi que $\Fib_W$, sont stables non seulement par changement de base (c'est le cas des 5 classes envisagées, qui de plus contiennent tous les isomorphismes), mais aussi par composition. (De plus, $\Fib_W$ est stable par facteurs directs, et je présume qu'il en est de même de $\mathrm{Liss}_W$ et de $\mathrm{Prop}_W$, donc de $\Parf_W$, en particulier de $\Parf_{W_\omega}$, mais je n'ai pas pris la peine de le vérifier.) Par contre, a priori $\Fibf_W$ n'est pas en général stable par composition, et $\Catfib^W$, $\Catfib^\omega$ n'ont pas de raison
% foot: Version 18 mars 2019
\printedpage{20}
% head: A. Grothendieck : Dérivateurs, Ch. XVI, Édition des Universités de Montpellier II et Paris VII
non plus d'être stables par composition \footnote[22]{il faudrait expliciter par des exemples.}.

4.2. Finalement, il paraît plus “éclairant” d'introduire, au lieu du localiseur fondamental $W_\omega$ bien déterminée, un localiseur fondamental % [sic] déterminée -> déterminé, accord avec « localiseur » (masculin)
\[
W_0 \subset \mathrm{Fl}(\Cat),
\]
qui jouera le rôle d'une sorte de “localiseur de référence” pour formuler certaines propriétés de $W$. On supposera toujours que
\[
W \supset W_0.
\]

\msp{page 30}

Suivant les cas, $W_0$ pourra être $W_\omega$, ou quelque autre localiseur fondamental défini comme le “plus petit” satisfaisant à des axiomes de stabilité donnés. Ou ça pourra être $\mathbf{W}_\infty$ (les équivalences faibles habituelles), quand on s'intéresse aux localiseurs qui contiennent $\mathbf{W}_\infty$ (les seuls, apparemment, avec lesquels on ait travaillé jusqu'à présent en homotopie). Ça pourra être aussi $W$, et les énoncés obtenus en spécialisant les résultats obtenus au cas où $W = W_0$ ne sont nullement tautologiques. Le diagramme de la page 28 se remplace alors par le diagramme d'inclusions suivant, plus éclairant pour notre propos :
\[
\begin{tikzcd}
 & \Fib_W \arrow[d] & \\
\Parf_W \arrow[r] & \Fibf_W & \mathcal{F}_W \arrow[l] \\
\Parf_{W_0} \arrow[r] \arrow[u] & \Fibf_{W_0} \arrow[u] & \mathcal{F}_{W_0} \arrow[l] \arrow[u] \\
 & \Fib_{W_0}, \arrow[u] &
\end{tikzcd}
\tag*{(4.2.1)}
\]
où pour tout localiseur fondamental $W$, je pose pour abréger
\[
\mathcal{F}_W = \Catfib^W,
\]
ce qui définit donc également $\mathcal{F}_{W_0}$. La classe $\Fib_W$ \msp{page 31} doit être regardée comme le point de mire de notre attention ; la plupart des propriétés de $W$ que nous allons passer en revue ici, s'exprimeront par une inclusion de la forme
\[
\mathcal{F} \subset \Fib_W,
\]
où $\mathcal{F}$ est une des sept autres classes du diagramme (4.2.1) ; à l'exception tout au plus du cas $\mathcal{F} = \Fib_{W_0}$, \textit{i.e.} de l'inclusion $\Fib_{W_0} \subset \Fib_W$. Comme propriété sur $W$, cette inclusion paraît très faible et pour cela pas utile à grand chose, sauf bien sûr dans le cas où $\Fib_{W_0} \longrightarrow \Fibf_{W_0}$ est une égalité, de sorte que l'inclusion envisagée équivaut à $\Fibf_{W_0} \subset \Fib_W$ (et implique que $\Parf_{W_0}$ et $\mathcal{F}_{W_0}$ sont eux aussi contenus dans $\Fib_W$). C'est là par contre, comme nous allons voir, une propriété assez forte et fort utile de $W$, impliquant notamment le fait que $(\Cat, W, \Fib_W)$ est une catégorie à fibrations.
% foot: Version 18 mars 2019

\printedpage{21}
% head: A. Grothendieck : Dérivateurs, Ch. XVI, Édition des Universités de Montpellier II et Paris VII

\msp{page 32}

Aux fins de référence ultérieure, je vais formuler maintenant les divers aspects principaux de \textit{l'axiome des $W$-fibrations} Loc (7). La formulation à laquelle j'étais parvenu (\textit{cf.} I, VI, XII) utilisait la notion que j'appelle à présent “$W$-fibrations \textit{faible}” (l'ex $W$-fibration, désormais détrônée !), et sous sa forme la plus forte s'écrivait comme l'inclusion % [sic] W-fibrations faible -> W-fibration faible, accord du nombre
\[
\Fibf_W \subset \Fib_W.
\]
Cette notion de $W$-fibration faible, tout comme celle de morphisme $W$-parfait, si utiles soient-elles, m'apparaît cependant un peu technique, et je préfère donner une formulation qui ne fasse appel qu'à des notions plus familières et d'apparence plus élémentaires. C'est pourquoi je vais remplacer $\Fibf_W$ par la sous-classe $\mathcal{F}_W = \Catfib^W$, ce qui (on le verra) revient en fait au même.

\msp{page 33}

Cette classe $\mathcal{F}_W$ est extrêmement naturelle, et s'impose même dans le yoga des dérivateurs. En effet, dans la construction de $\HOT_W(S)$ pour $S$ dans Cat, on passe à une catégorie de fractions de la catégorie
\[
\Cat(S) = \underline{\mathrm{Hom}}(S, \Cat),
\]
par un localiseur $W(S)$ (que nous avions généralement noté $W_S^{\mathrm{g}}$ ou $W_S^{\mathrm{u}}$ \ldots). Cette catégorie $\Cat(S)$ est d'autre part isomorphe à la catégorie
\[
\Fibsc\, S
\]
des catégories fibrées et scindées sur $S$, avec comme flèches les $S$-foncteurs qui respectent le scindage (et en particulier cartésiens) (\textit{cf.} VI, et notamment le diagramme p. 175). On a vu (\textit{loc. cit.}) que la catégorie localisée $\HOT_W(S)$ est équivalente à la catégorie localisée de la catégorie très voisine
\[
\Fib\, S \qquad \text{(que je préférerais maintenant noter $\Catfib\, S$ ou $\Catfib_S$),}
\]
formé des catégories \textit{fibrées} sur $S$, avec comme \msp{page 34} flèches les seuls foncteurs cartésiens. (On fera attention cependant que le foncteur canonique $\Fibsc\, S \longrightarrow \Fib\, S$ est fidèle mais non pleinement, néanmoins il induit une équivalence sur les catégories localisées, en définissant de façon évidente, par la $W$-équivalence fibre par fibre, \textit{i.e.} comme $W_S^{\mathrm{u}}$, le localiseur pertinent sur $\Fib\, S$.) Donc pour cette raison de principe, et d'ailleurs tout au long du développement de la théorie de $\HOT_W$, la catégorie $\Fib\, S$ et par là, les foncteurs fibrants, ont joué un rôle essentiel. La classe de flèches que nous notons ici $\mathcal{F}_W$ ou $\Catfib^W$ (on pourrait aussi la noter $\Fib^W$, si ce n'était pour le risque de confusion avec $\Fib_W$) s'interprète alors très naturellement à partir de là. Une flèche $f : X \longrightarrow S$ est dans $\mathcal{F}_W$ si et seulement si elle est a) fibrante, et par là définit de façon particulièrement directe un objet de $\HOT_W(S)$, puisque à $S$-équivalence près, $X/S$ provient d'un objet de $\Fibsc\, S$ ou encore de $\Cat(S)$, \msp{page 35} et b) l'objet dans $\HOT_W(S)$ défini par $f$ est dans $\HOT_W^{\mathrm{lc}}(S)$, \textit{i.e.} dans l'image essentielle de la sous-catégorie strictement pleine $\Fibsc_0^W S$ de $\Fibsc\, S$ (par le foncteur canonique
% foot: Version 18 mars 2019
\printedpage{22}
% head: A. Grothendieck : Dérivateurs, Ch. XVI, Édition des Universités de Montpellier II et Paris VII
$\Fibsc \longrightarrow \HOT_W(S)$), ou ce qui revient au même (du moins si $W$ est fortement saturée) que le foncteur % [sic] saturée -> saturé, accord avec W (« W fortement saturé » ailleurs)
\[
\begin{aligned}
S^{\circ} \quad &\longrightarrow \quad \mathrm{Hot}_W \\
s \quad &\longmapsto \quad \mathrm{hot}_W(X_s)
\end{aligned}
\]
soit un foncteur “localement constant”. Ainsi, on considérera les $f \in \mathcal{F}_W$ comme les flèches les plus économiques et les plus naturelles dans Cat, qui définissent des $W$-types d'homotopie relative localement constants.

Notons aussi que toutes les classes de flèches (4.2.1) sont stables par changement de base (et elles contiennent les isomorphismes). Les catégories $\Fib_W$, $\Parf_W$ et leurs variantes $W_0$ sont aussi stables par composition, par contre, ce n'est pas le cas pour $\Fibf_W$ et $\mathcal{F}_W$ et leurs variantes $W_0$, sauf axiomes particuliers sur $W$ ou sur $W_0$). Il en résulte aussitôt \msp{page 36} que si $\mathcal{F}$ est une de ces classes, la condition % [sic] « sur W_0). » -> parenthèse fermante sans parenthèse ouvrante correspondante
\[
\mathcal{F} \subset \Fib_W
\]
équivaut à celle ci :

\begin{itemize}[label={},leftmargin=3.5cm]
\item Pour tout diagramme cartésien
\[
\begin{tikzcd}
X \arrow[d, "f"'] & X' \arrow[l, "q"'] \arrow[d, "f'"] \\
S & S', \arrow[l, "p"]
\end{tikzcd}
\tag*{(4.2.2)}
\]
avec $f \in \mathcal{F}$ et $p \in W$, on a $q \in W$.

(N.B. On a aussi $f' \in \mathcal{F}$, par la stabilité rappelée plus haut.)
\end{itemize}

4.3. Voici donc (\textit{cf.} plus bas) l'axiome pertinent, Loc (7) avec ses trois aspects a, b, c \footnote[23]{J'y ai rajouté finalement un quatrième F (comme “fibration”).}. On a vu (et on reverra) que l'on a les implications
\[
\text{a} \qquad \Longrightarrow \qquad \text{b} \qquad \Longrightarrow \qquad \text{c} \qquad (\Longrightarrow \qquad \text{F}).
\]
Quand on référera à Loc (7) sans autre précision, il sera entendu qu'il s'agit de la forme la plus forte Loc (7 a), qui implique donc les deux autres. Il n'en faut pas moins pour être totalement à l'aise. Mais je rappelle que la variante affaiblie Loc (7b) est suffisante pour faire la théorie des foncteurs $f_!$ pour les catégories $\HOT_W(S)$ \footnote[24]{Ici ma mémoire me joue un tour - on n'a nullement besoin pour cela de l'axiome déjà très fort Loc (7b), mais seulement de Loc (3 bis), dont il semble bien qu'il soit vérifiée pour tous les localiseurs fondamentaux “raisonnables”. % [sic] qu'il soit vérifiée -> qu'il soit vérifié, accord avec « il »
} (\textit{cf.} VI).

\msp{page 37}

C'est d'ailleurs une chose surprenante, puisque via le yoga général des dérivateurs associés aux catégories de modèles, on s'attendrait plutôt que pour avoir des adjoints à gauche $f_!$ des foncteurs $f^*$ d'image inverse, il faut faire sur $W$ des hypothèses du type \textit{$W$-cofibration}
% foot: Version 18 mars 2019
\printedpage{23}
% head: A. Grothendieck : Dérivateurs, Ch. XVI, Édition des Universités de Montpellier II et Paris VII
(essentiellement, Loc (6 a)). Il me semble peu probable que Loc (7 b), ni même Loc (7), implique Loc (6 a), ni même sa variante plus faible Loc (6$'$ a) (axiome du carré cocartésien, ou de Mayer-Vietoris, forme faible). Si d'autre part ces axiomes ne sont pas \textit{vérifiés,} alors que Loc (7) l'est, il y a des chances que cela donnera un exemple où il y a un formalisme des $f_!$, sans que l'axiome d'exactitude à gauche soit vérifié.

Quant à la forme Loc (7 c), plus faible encore, elle suffit pourtant (comme nous le verrons plus bas) à assurer que Cat admet suffisamment de $W$-fibrations \msp{page 38} (résultat qui n'a rien d'évident, et ne pouvait m'apparaître qu'après avoir pris connaissance du substantiel lemme de K. S. \textsc{Brown}, dans son beau travail où il introduit les catégories à fibrations et à cofibrations).

Nous donnerons aussi plus bas des cas fort simples, où la condition Loc (7 b) entraîne déjà Loc (7) ($W$ fortement saturé, \textit{i.e.} Loc (1 bis)), voire même où Loc (7 c) entraîne Loc (7) (cas où $W$ satisfait à l'axiome Loc (8) des limites inductives). Je signale en passant que je ne connais pas de localiseur fondamental satisfaisant l'axiome Loc (8) (qui implique déjà $W = W^{\natural}$, \textit{i.e.} Loc (1 bis)), sans que $W$ soit fortement saturé (\textit{i.e.} en satisfasse même Loc (1 ter)).

\begin{itemize}[label={},leftmargin=0.6cm]
\item \textbf{Loc (7)}
\begin{itemize}[leftmargin=1.2cm]
\item[a)] On a $\mathcal{F}_W \subset \Fib_W$, \textit{i.e.} pour tout diagramme cartésien (4.2.2) dans Cat, avec $f \in \mathcal{F}_W$ (\textit{i.e.} $f$ Cat-fibrant et faiblement $W$-fibrant) et $p \in W$, on a $q \in W$.
\item[b)] Pour tout diagramme commutatif
\[
\begin{tikzcd}[column sep=small]
X \arrow[rr, "f"] \arrow[dr, "p"'] & & Y \arrow[dl, "q"] \\
 & S &
\end{tikzcd}
\]
\msp{page 39}
avec $p, q \in \mathcal{F}_W$ et $f \in W$, on a $f \in W_S^{\mathrm{u}}$ (ou, ce qui revient au même comme $X$, $Y$ sont $W$-lisses sur $S$, $f_s \in W$ pour tout $s \in \mathrm{Ob}\, S$) \footnote[25]{\textbf{N.B.} La démonstration de $B(\mathcal{F}_W) \Longrightarrow B(\Fibf_W)$ montre en fait qu'il suffit de postuler b) lorsque $f$ est cartésien, et même lorsque $X$, $Y$ sont scindées sur $S$ et $f$ compatible avec le scindage.}.
\item[c)] On a
\[
\mathcal{F}_W \cap W \subset W^{\mathrm{u}},
\]
en d'autres termes, pour toute flèche $f : X \longrightarrow S$ dans Cat telle que $f \in \mathcal{F}_W$ (on pourrait appeler ces flèches des Cat-$W$-fibrations, étant entendu qu'elles sont aussi des Cat-fibrations, mais non nécessairement des $W$-fibrations), si $f \in W$, alors $f \in W_S^{\mathrm{u}}$ (ou, ce qui revient au même, les fibres $X_s$ de $f$ ($s \in \mathrm{Ob}\, S$) sont $W$-asphériques).
\end{itemize}
\end{itemize}
% foot: Version 18 mars 2019


\endgroup
% END UNIT der-ch16/der-ch16_p021-025.tex
% BEGIN UNIT der-ch16/der-ch16_p026-030.tex
\begingroup\makeatletter%
\newcommand{\printedpage}[1]{\marginpar{\footnotesize p.~#1}}%
\newcommand{\msp}[1]{\marginpar{\footnotesize [#1]}}%
\newcommand{\cF}{\mathcal{F}}%
\newcommand{\cM}{\mathcal{M}}%
\newcommand{\Cat}{\mathrm{Cat}}%
\newcommand{\Fib}{\mathrm{Fib}}%
\newcommand{\Fibf}{\mathrm{Fibf}}%
\newcommand{\Parf}{\mathrm{Parf}}%
\newcommand{\Fl}{\mathrm{Fl}}%
\newcommand{\Comp}{\mathrm{Comp}}%
\newcommand{\Catfib}{\mathrm{Catfib}}%
\newcommand{\dottedimplies}{\mathrel{\tikz[baseline=-0.6ex]{\draw[dotted, double, double distance=1.5pt, -{Implies}] (0,0) -- (0.9,0);}}}%
\makeatother\let\printedpage\derpp\let\msp\dermsp%
\let\mspin\DERmspin%


\printedpage{24}
% head: A. Grothendieck : Dérivateurs, Ch. XVI, Édition des Universités de Montpellier II et Paris VII
% foot: Version 18 mars 2019

\begin{itemize}[leftmargin=4em]
\item[F)] \footnote[26]{rajouté après coup.} Cat a suffisamment de $W$-fibrations, \textit{i.e.} Cat muni de $(W, \Fib_W)$ est une catégorie à cofibrations (propre, et même de Brown, \textit{i.e.} $X \longrightarrow e$ est dans $\Fib_W$ pour tout $X$ dans Cat. % [sic] (propre, et même de Brown, -> parenthesis never closed, closing ")" apparently missing
\end{itemize}

Il est clair que c est le cas particulier de b obtenu en prenant $Y = S$, $q = \mathrm{id}_S$, donc b $\Longrightarrow$ c. L'implication a $\Longrightarrow$ b est moins évidente, elle est établie d'abord dans XII, et nous y reviendrons ci-dessous (4.4.2 ci-dessous). Le reste du présent paragraphe est consacré à étudier diverses variantes des conditions a) b) c) précédentes et les implications entre celles-ci.

\msp{page 40}

4.4. Soit, plus généralement,
\[
\cF \subset \Fl(\Cat)
\]
un ensemble de flèches dans Cat, en pensant plus particulièrement au cas où $\cF$ est un des sept ensembles du diagramme (4.2.1) (p. 30) autres que $\Fib_W$. Nous pouvons formuler trois propriétés similaires pour $\cF$ à celles formulées dans Loc (7 a b c) dans le cas de $\cF_W$, notons ces propriétés
\[
\mathrm{A}(\cF)\,,\ \mathrm{B}(\cF)\,,\ \mathrm{C}(\cF)\,,\ \text{\footnotemark[27]}
\]
\footnotetext[27]{\textbf{N.B.} Ces conditions ont un sens si $\cF \subset \Fl(\cM)$, si $(\cM, W)$ est une catégorie de modèles quelconque.}
donc
\[
\left\{
\begin{array}{lcl}
\mathrm{A}(\cF) & : & \cF \subset \Fib_W. \\[1ex]
\mathrm{B}(\cF) & : & \Fl(\cF_S) \cap W_S \subset W_S^{\mathrm{u}} \\
 & & \text{(où $\cF_S$ désigne la sous-catégorie pleine de $\Cat/S$} \\
 & & \text{formée des $f : X \longrightarrow S$ avec $f \in \cF$).} \\[1ex]
\mathrm{C}(\cF) & : & \cF \cap W \subset W^{\mathrm{u}}.
\end{array}
\right.
\tag*{(4.4.1)}
\]
Il est clair qu'on a $\mathrm{B}(\cF) \Longrightarrow \mathrm{C}(\cF)$, et je vais donner une condition sur $\cF$ qui implique qu'on a aussi $\mathrm{A}(\cF) \Longrightarrow \mathrm{B}(\cF)$.

\msp{page 41}

(Nous l'appliquerons dans les cas suivants : % [sic] (Nous -> opening parenthesis never closed in this paragraph, closing ")" apparently missing
\[
\cF \;=\; \cF_W,\ \Fibf_W,\ \Parf_W,\ \cF_{W_0},\ \Fibf_{W_0}.
\tag*{(4.4.2)}
\]
Je signale d'ailleurs que l'on verra plus bas (cor. 4.4.10) que
\[
\mathrm{A}(\cF_W) \iff \mathrm{A}(\Fibf_W), \qquad \mathrm{A}(\cF_{W_0}) \iff \mathrm{A}(\Fibf_{W_0}),
\tag*{(4.4.3)}
\]
et de même pour les variantes B, C, de sorte que pour l'implication que nous voulons établir, il n'y a pas lieu de distinguer les cas de $\cF_W$ et $\Fibf_W$, de $\cF_{W_0}$ et $\Fibf_{W_0}$. D'ailleurs il suffira de regarder les ensembles $\cF_{W_0}$, $\Fibf_{W_0}$, $\Parf_{W_0}$, puisque le cas des ensembles relatifs à $W$ en sont un cas particulier (en prenant $W_0 = W$). \footnote[28]{Cela fera donc a priori $3 \times 6 = 18$ conditions différentes, mais à cause de (4.4.3) il n'y en a (à équivalence près) que $4 \times 3 = 12$, et en fait deux parmi ces quatre triples de conditions se réduisent à une seule (pour $\cF = \cF_W$ et $\cF = \cF_{W_0}$), de sorte qu'il ne reste au total que \textit{huit} conditions différentes, que nous résumons dans le diagramme récapitulatif (4.4.7.1) de la page 48. Il y en a même seulement \textit{quatre,} $\mathrm{A}_{W_0}$, $\mathrm{B}_{W_0}$, $\mathrm{C}_{W_0}$, $\mathrm{D}_{W_0}$, comme fonctions de $W_0$, les quatre autres A, B, C, D s'obtiennent en spécialisant au cas où $W_0 = W$.}

\printedpage{25}
% head: A. Grothendieck : Dérivateurs, Ch. XVI, Édition des Universités de Montpellier II et Paris VII
% foot: Version 18 mars 2019

\textbf{Lemme 4.4.1.} \footnote[29]{Vaut dans une catégorie de modèles quelconque.} \textit{Supposons que $\cF$ satisfasse la condition suivante : toute flèche $f$ de} Cat \textit{se factorise en $f = pi$, avec $i \in W$, $f \in \cF$. Alors la condition} $\mathrm{A}(\cF)$ \textit{sur $\cF$ implique que dans} Cat \textit{il y a “assez de $W$-fibrations”, et par suite} (\textit{par un argument formel bien} % [sic] f \in \cF -> p \in \cF, the factorisation is f = pi and lemme 4.4.3 3) restates this condition with p \in \cF
\msp{page 42}
\textit{connu}), \textit{que l'on a} $\mathrm{B}(\Fib_W)$, \textit{et a fortiori} (\textit{comme par hypothèse} $\cF \subset \Fib_W$) $\mathrm{B}(\cF)$.

\textbf{N.B.} Il est clair en effet que si on a $\cF \subset \cF'$, alors
\[
\mathrm{A}(\cF') \Longrightarrow \mathrm{A}(\cF) \qquad \text{et de même pour B, C.}
\tag*{$(*)$}
\]
Or la condition préliminaire sur $\cF$ est bien satisfaite dans le cas d'un ensemble de l'une des formes $\cF \;=\; \cF_W$, $\Fibf_W$, $\Parf_W$. Il suffit de le voir dans le cas des sous-ensembles $\cF_W$, $\Parf_W$ de $\Fibf_W$, \textit{i.e.} que toute flèche $f$ de Cat se factorise en
\[
f = pi,
\]
avec $i \in W$, et $p$ étant, au choix, dans $\Parf_W$, ou dans $\cF_W$. Or dans le cas de $\Parf_W$, \marginpar{$\|$}c'est ce qui est établie dans VII (Catégories de chemins (2)) grâce au formalisme des $\underline{\mathrm{Ch}}(S) \longrightarrow S \times S$. Pour trouver une factorisation avec $f = qj$, $j \in W$, $q \in \cF_W$, il suffit de voir que $p$ se factorise ainsi en $p = q_0 j_0$, on aura % [sic] établie -> établi, agreement with "ce qui"
\[
f = \underbrace{q_0}_{\in \cF_{W_0}} \underbrace{(j_0 i)}_{\in W},
\]
donc on peut supposer $f \in \Parf_W$. Mais si $X$ est $W$-parfait sur $S$, donc un
\msp{page 43}
$W$-fibré faible sur $S$, alors dans la factorisation
\begin{center}
\begin{tikzcd}[column sep=small]
X \arrow[rr, "i"] \arrow[dr, "f"'] & & \Phi_S(X) \arrow[dl, "p"] \\
 & S &
\end{tikzcd}
\end{center}
avec $i \in W$ et $p \in \Catfib$ (\textit{cf.} VI), on sait que $p$ est aussi une $W$-fibration faible, \textit{i.e.} $p \in \cF_W$, d'où la conclusion. Ainsi on trouve le

\textbf{Corollaire 4.4.2.} \textit{Supposons que $\cF$ soit l'un des six ensembles} (4.4.2) (i.e. \textit{l'un des ensembles du diagramme} (4.2.1), \textit{à l'exclusion de $\mathit{Fib}_W$ et $\mathit{Fib}_{W_0}$}). \textit{Alors on a}
\[
\mathrm{A}(\cF) \;\Longrightarrow\; \mathrm{B}(\cF) \;\Longrightarrow\; \mathrm{C}(\cF)\,.
\]

\textbf{Lemme 4.4.3.} \textit{Soit $(\cM, W)$ une catégorie de modèles, et soit $\cF \subset \Fl(\cM)$ satisfaisant les conditions suivantes}
\begin{itemize}[leftmargin=3em]
\item[1)] \textit{$\cF$ stable par changement de base, composition, contient les isomorphismes.}
\item[2)] \textit{$\cM$ a un objet final, et pour tout $X \in \mathrm{Ob}\,\cM$, $X \longrightarrow e$ est dans $\cF$.}
\item[3)] \textit{Toute flèche $f$ de $\cM$ se factorise en $f = pi$, avec $i \in W$, $p \in \cF$} (\textit{c'est la condition du lemme} 4.4.1).
\end{itemize}

\printedpage{26}
% head: A. Grothendieck : Dérivateurs, Ch. XVI, Édition des Universités de Montpellier II et Paris VII
% foot: Version 18 mars 2019
\noindent\textit{Alors les conditions} $\mathrm{A}(\cF)$, $\mathrm{B}(\cF)$, $\mathrm{C}(\cF)$ \textit{sur $\cF$ sont équivalentes.}

En effet, en vertu de 4.4.1, il reste à prouver l'implication
\[
\mathrm{C}(\cF) \;\Longrightarrow\; \mathrm{A}(\cF),
\]

\msp{page 44}

\textit{i.e.} ceci : si $\cF$ satisfait aussi la condition
\begin{itemize}[leftmargin=3em]
\item[4)] $\cF \cap W \subset W^{\mathrm{u}}$ (ou ce qui revient au même, comme $\cF$ est stable par changement de base, que $\cF \cap W$ est stable par changement de base),
\end{itemize}
alors on en conclut
\[
\cF \subset \Fib_W.
\]
Or ceci n'est autre que le lemme de K. S. \textsc{Brown}, sous la forme duale de celle donnée plus haut (3.2.1, page 20).

Parmi les six ensembles de flèches envisagés dans (4.4.2), les hypothèses 1) 2) 3) sont satisfaites pour $\Parf_W$ et $\Parf_{W_0}$, à l'exclusion des quatre autres. (Les autres satisfont à toutes les conditions, à la seule exception de la stabilité par composition \footnote[30]{\textbf{N.B.} Pour $\Fib_{W_0}$, ce sont les conditions 1) 2) qui sont vérifiées, et par contre, la condition 3) qui fait problème. Notons là $\mathrm{F}(W_0)$. Si elle est satisfaite (\textit{i.e.} si Cat a assez de $W_0$-fibrations), alors on a aussi équivalence des conditions $\mathrm{A}(\Fib_{W_0})$, $\mathrm{B}(\Fib_{W_0})$, $\mathrm{C}(\Fib_{W_0})$.}). On trouve donc :

\textbf{Corollaire 4.4.4.} \textit{Les conditions} $\mathrm{A}(\Parf_W)$, $\mathrm{B}(\Parf_W)$, $\mathrm{C}(\Parf_W)$ \textit{sont équivalentes} (nous désignons cette condition sur $W$ par D). \textit{De même les conditions} $\mathrm{A}(\Parf_{W_0})$, $\mathrm{B}(\Parf_{W_0})$, $\mathrm{C}(\Parf_{W_0})$ \textit{sont équivalentes.}

\msp{page 45}

(On désignera la condition obtenue sur le couple $W \supset W_0$ de localiseurs fondamentaux par $D_{W_0}$.)

Ainsi on a (comme $\Parf_{W_0} \subset \Parf_W$)
\[
\mathrm{D} = \mathrm{D}_W \;\Longrightarrow\; \mathrm{D}_{W_0}\,,
\tag*{(4.4.4)}
\]
et si $W'_0 \subset W_0$, on aura
\[
\mathrm{D}_{W_0} \;\Longrightarrow\; \mathrm{D}_{W'_0}\,.
\tag*{(4.4.5)}
\]

\textbf{Corollaire 4.4.5.} \textit{Désignons par} A, B, C \textit{respectivement les conditions} Loc (7 a), Loc (7 b), Loc (7 c). \textit{On a alors les implications}
\[
\mathrm{A} \;\Longrightarrow\; \mathrm{B} \;\Longrightarrow\; \mathrm{C}\,.
\tag*{(4.4.5.1)}
\]
\textit{De plus, si on désigne} (\textit{pour un ensemble $\cF$ de flèches de} Cat) \textit{par} $\Comp(\cF)$ \textit{la condition sur $\cF$ que $\cF$ soit stable par composition, on a les équivalences}
\[
\mathrm{A} \iff (\mathrm{C} + \Comp(\cF_W)) \iff (\mathrm{C} + \Comp(\Fibf_W)).
\tag*{(4.4.5.2)}
\]

\printedpage{27}
% head: A. Grothendieck : Dérivateurs, Ch. XVI, Édition des Universités de Montpellier II et Paris VII
% foot: Version 18 mars 2019

\textsc{Démonstration}. La première assertion est contenue dans le corollaire 4.4.2, appliqué au cas $\cF = \cF_W$. D'autre part, si $\cF_W$ satisfait à $\Comp(\cF_W)$, il satisfait aux conditions 1) 2) 3) du lemme de Brown 4.4.3, donc joint à C, il implique A. Il reste, pour la première équivalence, à prouver que
\[
A \;\Longrightarrow\; \Comp(\cF_W).
\]

\msp{page 46}

Mais on a par définition
\[
\cF_W = \Catfib \cap \Fibf_W
\tag*{$(*)$}
\]
et on sait que Catfib est stable par composition. Donc $\cF_W$ l'est aussi, pourvu que le composé $gf$ de deux flèches dans $\cF_W$ soit dans $\Fibf_W$. Or par hypothèse A, qui signifie $\cF_W \subset \Fib_W$, on a $f, g \in \Fib_W$, donc $gf \in \Fib_W$, puisque $\Fib_W$ est stable par composition, a fortiori $gf \in \Fibf_W$.

D'autre part, la formule précédente $(*)$ montre que $\Comp(\Fibf_W) \Longrightarrow \Comp(\cF_W)$, donc on a gratis l'implication $\Longleftarrow$ dans la deuxième équivalence (4.4.5.2). L'implication en sens inverse revient alors à
\[
\mathrm{A} \;\Longrightarrow\; \Comp(\Fibf_W),
\]

\msp{page 47}

ce qui sera évident grâce au résultat (donné plus bas)
\[
\mathrm{A} = \mathrm{A}(\cF_W) \iff \mathrm{A}(\Fibf_W) \qquad \text{et de même pour B, C,}
\]
qui signifie donc que
\[
\cF_W \subset \Fib_W \qquad \text{implique} \qquad \Fibf_W \subset \Fib_W \qquad (\text{donc } \Fibf_W = \Fib_W),
\]
or $\Fib_W$ étant stable par composition, il en est de même pour $\Fibf_W$, q.e.d.

\textbf{Corollaire 4.4.6.} \textit{Supposons que $\cF_{W_0}$ soit stable par composition, ce qui est le cas notamment} (\textit{par} 4.4.5.2 \textit{appliqué à $W_0$ au lieu de $W$}) \textit{si $W_0$ satisfait} Loc (7). \textit{Alors les conditions} $\mathrm{A}(\cF_{W_0})$, $\mathrm{B}(\cF_{W_0})$, $\mathrm{C}(\cF_{W_0})$ (\textit{que nous notons aussi simplement} $\mathrm{A}_{W_0}$, $\mathrm{B}_{W_0}$, $\mathrm{C}_{W_0}$) \textit{sont équivalentes.} (\textbf{\textit{N.B.}} \textit{Si on ne fait pas l'hypothèse précédente sur $W_0$, on aura seulement} $\mathrm{A}_{W_0} \Longrightarrow \mathrm{B}_{W_0} \Longrightarrow \mathrm{C}_{W_0}$, \textit{en vertu de} 4.4.2.)

4.4.7. On a fait le tour à présent de toutes les conditions de type “fibration” sur $W$, ou sur le couple $(W, W_0)$, qui m'ont parues utiles. Je vais les déployer ici dans un diagramme (plus ou moins) récapitulatif : % [sic] parues -> paru, agreement of the past participle

\msp{page 48}

\printedpage{28}
% head: A. Grothendieck : Dérivateurs, Ch. XVI, Édition des Universités de Montpellier II et Paris VII
% foot: Version 18 mars 2019

\begin{center}
\begin{tikzcd}[column sep=2.3em, row sep=small]
 & \overbrace{\mathrm{A}_{W_0}}^{\cF_{W_0},\,\Fibf_{W_0}} \arrow[dd, Rightarrow] & & \\
\overbrace{\mathrm{A}}^{\cF_W,\,\Fibf_W} \arrow[ur, Rightarrow] \arrow[dd, Rightarrow] & & & \\
 & \mathrm{B}_{W_0} \arrow[dd, Rightarrow] & & \\
\mathrm{B} \arrow[ur, Rightarrow] \arrow[dd, Rightarrow] & & & \\
 & \mathrm{C}_{W_0} \arrow[dr, Rightarrow, "\bullet"] \arrow[uuuu, Rightarrow, dotted, bend right=40, "\text{\scriptsize si } W_0\ \text{\scriptsize (7)}"'] & & \\
\mathrm{C} \arrow[ur, Rightarrow] \arrow[dr, Rightarrow, "\bullet"] & & \overbrace{\mathrm{D}_{W_0}}^{\Parf_{W_0}} \arrow[r, Rightarrow, "\bullet"] & \mathrm{F} : (\Cat, W, \Fib_W) \text{ une catégorie à fibrations.} \\
 & \overbrace{\mathrm{D}}^{\Parf_W} \arrow[ur, Rightarrow] & &
\end{tikzcd}
\end{center}

\begin{center}
\begin{tikzcd}[row sep=large]
 & \mathrm{A} \arrow[r, Rightarrow] \arrow[d, Rightarrow] & \mathrm{A}_{W_0} \arrow[d, Rightarrow] & \\
\cF_W, \Fibf_W \smash{\left\{\vrule height 2.6cm depth 2.6cm width 0pt\right.} & \mathrm{B} \arrow[r, Rightarrow] \arrow[d, Rightarrow] & \mathrm{B}_{W_0} \arrow[d, Rightarrow] & \smash{\left.\vrule height 2.6cm depth 2.6cm width 0pt\right\}} \quad \text{\scriptsize\begin{tabular}{@{}l@{}}version revue\\ du diagramme\end{tabular}} \\
 & \mathrm{C} \arrow[r, Rightarrow] \arrow[d, Rightarrow] & \mathrm{C}_{W_0} \arrow[d, Rightarrow] & \\
 & \underbrace{\mathrm{D}}_{\Parf_W} \arrow[r, Rightarrow] & \underbrace{\mathrm{D}_{W_0}}_{\Parf_{W_0}} \arrow[r, Rightarrow] & \mathrm{F},
\end{tikzcd}
\end{center}

diagramme que je vais commenter. Toutes les implications en traits pleins du diagramme ont été déjà obtenues ou sont tautologiques, à l'exception des trois $\overset{\bullet}{\Longrightarrow}$ sur lesquels je vais revenir. L'implication en pointillés $\mathrm{C}_{W_0} \dottedimplies \mathrm{A}_{W_0}$ est une implication conditionnelle, soumise à la condition marquée sur le diagramme comme $W_0$ (7), ce qui signifie : $W_0$ satisfait à Loc (7). Cette implication, \textit{i.e.} l'équivalence de $\mathrm{A}_{W_0}$, $\mathrm{B}_{W_0}$, $\mathrm{C}_{W_0}$ sous la condition $W_0$ (7), a été vue dans 4.4.6. J'ai marqué, au dessus des deux groupes (A, B, C) et ($\mathrm{A}_{W_0}$, $\mathrm{B}_{W_0}$, $\mathrm{C}_{W_0}$) ainsi qu'au dessus de D et de D$'$ (qui correspondent chacun à un tel groupe de conditions $\mathrm{A}(\cF)$, $\mathrm{B}(\cF)$, $\mathrm{C}(\cF)$ toutes équivalentes, avec $\cF \;=\; \Parf_W$ ou $\cF = \Parf_{W_0}$) les
\msp{page 49}
ensembles de flèches qui leur donnent naissance respectivement : on notera qu'il y en a deux pour les pour les deux triplets nommés en premier, A, B, C pouvant se définir via $\cF_W$ ou $\Fib_W$ indifféremment, et de même $\mathrm{A}_{W_0}$, $\mathrm{B}_{W_0}$, $\mathrm{C}_{W_0}$ via $\cF_{W_0}$ ou $\Fib_{W_0}$, comme déjà énoncé (4.4.3) et comme nous le verrons plus bas (cor. 4.4.10, pages 54-55). Par contre, les triplets réduits D et $\mathrm{D}_{W_0}$ sont définis chacun par un $\cF$ seulement, savoir par $\Parf_W$ et % [sic] pour les pour les -> pour les, words repeated; [sic] Fib_W / Fib_{W_0} -> Fibf_W / Fibf_{W_0}, the diagram and (4.4.3) give Fibf


\endgroup
% END UNIT der-ch16/der-ch16_p026-030.tex
% BEGIN UNIT der-ch16/der-ch16_p031-035.tex
\begingroup\makeatletter%
\newcommand{\printedpage}[1]{\marginpar{\footnotesize p.~#1}}%
\newcommand{\msp}[1]{\marginpar{\footnotesize [#1]}}%
\newcommand{\Fib}{\mathrm{Fib}}%
\newcommand{\Fibf}{\mathrm{Fibf}}%
\newcommand{\Parf}{\mathrm{Parf}}%
\newcommand{\Fl}{\mathrm{Fl}}%
\newcommand{\inc}{\mathrm{inc}}%
\newcommand{\Ob}{\mathrm{Ob}}%
\newcommand{\dotimp}{\overset{\bullet}{\Longrightarrow}}%
\newcommand{\dottedimpl}{\mathrel{\begin{tikzcd}[ampersand replacement=\&, column sep=small, baseline=-0.5ex]{}\arrow[r,Rightarrow,dotted]\&{}\end{tikzcd}}}%
\newcommand{\circled}[1]{\tikz[baseline=(c.base)]\node[draw,circle,inner sep=1pt](c){#1};}%
\makeatother\let\printedpage\derpp\let\msp\dermsp%
\let\mspin\DERmspin%

\setcounter{footnote}{30}

\printedpage{29}
% head: A. Grothendieck : Dérivateurs, Ch. XVI, Édition des Universités de Montpellier II et Paris VII
% foot: Version 18 mars 2019
$\Parf_{W_0}$ respectivement. (En tout, cela fait intervenir six ensembles $\mathcal{F}$ distincts, à savoir les six ensembles du diagramme (4.2.1) autres que $\Fib_W$ et $\Fib_{W_0}$.)

L'implication non tautologique $\mathrm{C} \dotimp \mathrm{D}$ est un cas particulier de l'implication $\mathrm{C}_{W_0} \dotimp \mathrm{D}_{W_0}$, laquelle est contenue dans ce qui a été annoncé à l'instant, à savoir $\mathrm{C}_{W_0}(\mathcal{F}_{W_0}) \Longrightarrow \mathrm{C}_{W_0}(\Fib_{W_0})$, compte tenu que $\Parf_{W_0} \subset \Fibf_{W_0}$ (\textit{cf.} $(*)$, p. 42). Faisant
\msp{page 50}
$W_0 = W$, on trouve donc l'implication $\mathrm{A} \Longrightarrow \mathrm{D}$, qui est cependant contenue dans les autres implications du diagramme, comme composée dans la chaîne d'implications
\[
\mathrm{A} \Longrightarrow \mathrm{B} \Longrightarrow \mathrm{C} \Longrightarrow \mathrm{D} \Longrightarrow \mathrm{D}_{W_0} \dotimp \mathrm{F}.
\]

Il reste, dans ces commentaires généraux, à dire un mot de cette dernière implication
\[
\mathrm{D}_{W_0} \dotimp \mathrm{F},
\]
\textit{i.e.} du fait que la plus faible de toutes les conditions A à $\mathrm{D}_{W_0}$ envisagées ici implique déjà que Cat a assez de $W$-fibrations, \textit{i.e.} est une catégorie à cofibrations pour $(W, \Fib_W)$. C'est clair quand on prend $\mathrm{D}_{W_0}$ sous la forme $\mathrm{A}(\Parf_{W_0})$, \textit{i.e.}
\[
\Parf_{W_0} \subset W.
\]
Comme toute flèche $f$ de Cat se factorise en $f = pi$, avec
\[
i \in W_0 \subset W \quad\text{et}\quad f \in \Parf_{W_0} \subset \Fib_W\, , % [sic] f \in \Parf_{W_0} -> p \in \Parf_{W_0}, the factor of f = pi lying in Parf is p
\]
on trouve la factorisation avec $i \in W$, $f \in \Fib_W$. % [sic] f \in \Fib_W -> p \in \Fib_W, the factorisation f = pi requires p \in \Fib_W

\msp{page 51}
4.4.8. Pour terminer d'établir les implications du diagramme p. 48, il reste seulement à établir les équivalences
\[
\left\{
\begin{aligned}
&\mathrm{A}(\mathcal{F}_{W_0}) &&\Longleftrightarrow &&\mathrm{A}(\Fibf_{W_0})\\
&\mathrm{B}(\mathcal{F}_{W_0}) &&\Longleftrightarrow &&\mathrm{B}(\Fibf_{W_0})\\
&\mathrm{C}(\mathcal{F}_{W_0}) &&\Longleftrightarrow &&\mathrm{C}(\Fibf_{W_0}),
\end{aligned}
\right.
\tag*{(4.4.8.1)}
\]
qui donneront les équivalences similaires pour $W$, en prenant $W_0 = W$.

La première équivalence (4.4.8.1) sera conséquence du lemme ($\pm$ tautologique) suivant :

\textbf{Lemme 4.4.9.} \textit{Soient} $(\mathcal{M}, W)$ \textit{une catégorie de modèles,} $\mathcal{F}, \mathcal{G} \subset \Fl(\mathcal{M})$, \textit{tels que tout} $f \in \mathcal{F}$, $f : X \longrightarrow S$, \textit{s'écrive} $f = pi$, \textit{avec} $i \in W^u_S$, $p \in \mathcal{G}$. \textit{Alors} $\mathrm{A}(\mathcal{G}) \Longrightarrow \mathrm{A}(\mathcal{F})$, i.e.
\[
\mathcal{G} \subset \Fib_W \;\Longrightarrow\; \mathcal{F} \subset \Fib_W,
\tag*{(4.4.9.1)}
\]
\textit{et de même} $\mathrm{C}(\mathcal{G}) \Longrightarrow \mathrm{C}(\mathcal{F})$, i.e.
\[
\mathcal{G} \cap W \subset W^u \;\Longrightarrow\; \mathcal{F} \cap W \subset W^u.
\tag*{(4.4.9.2)}
\]

\textsc{Démonstration}. A) En effet, si $\mathcal{G} \subset \Fib_W$, la factorisation $f = pi$ de $f$ est telle que $p \in \Fib_W$. Or un composé $pi$, avec $p \in \Fib_W$ et $i \in W^{\mathrm{u}}_S$, est dans $\Fib_W$, comme on vérifie aussitôt sur le diagramme de carrés cartésiens
\printedpage{30}
% head: A. Grothendieck : Dérivateurs, Ch. XVI, Édition des Universités de Montpellier II et Paris VII
% foot: Version 18 mars 2019
\[
\begin{tikzcd}
X \arrow[d, "W^{\mathrm{u}}_S"'] & X' \arrow[l] \arrow[d, "(W)"'] & X'' \arrow[l] \arrow[d, "(W)"'] \\
\overline{X} \arrow[d, "\Fib_W"'] & \overline{X'} \arrow[l] \arrow[d] & \overline{X''} \arrow[l, "(W)"'] \arrow[d] \\
S & S' \arrow[l] & S'' \arrow[l, "W"']
\end{tikzcd}
\]
sur lequel on lit que $X'' \longrightarrow X'$ est dans $W$, q.e.d.

\msp{page 52}
Ce lemme s'applique notamment au cas où $\mathcal{F} = \Fibf_{W_0}$, $\mathcal{G} = \mathcal{F}_{W_0}$, car $f$ se factorise en
\[
\begin{tikzcd}
X \arrow[dr, "i"] \arrow[dd, "f"'] & \\
& \Phi_S(X) \arrow[dl, "p"] \\
S, &
\end{tikzcd}
\]
et on sait que $i \in W^{\mathrm{u}}_{0\,S} \subset W^{\mathrm{u}}_S$, et $p \in \mathcal{F}_{W_0} \subset \mathcal{F}_W$, ce qui sont les conditions voulues.

C) Démonstration similaire : Soit $f \in \mathcal{F} \cap W$, factorisant en $f = pi$, avec $i \in W^{\mathrm{u}}_S$, $g \in \mathcal{G}$, on aura alors $g \in W$, donc $g \in \mathcal{G} \cap W \subset W^{\mathrm{u}}$ (par hypothèse $\mathrm{C}(\mathcal{G})$), a fortiori on a $g \in W^{\mathrm{u}}_S$, donc $f \in W^{\mathrm{u}}_S$ comme composé de deux flèches de $W^{\mathrm{u}}_S$, et comme le but de $f$ est $S$ lui-même, cela signifie $f \in W^{\mathrm{u}}$, q.e.d.
\[
\begin{tikzcd}
X \arrow[dr, "i\in W^{\mathrm{u}}_S"] \arrow[dd, "\mathcal{F}\cap W"'] & \\
& \bar{X} \arrow[dl, "g"] \\
S &
\end{tikzcd}
\]

Je ne vois pas de façon pour prouver également
\[
\mathrm{B}(\mathcal{G}) \qquad\quad \Longrightarrow \qquad\quad \mathrm{B}(\mathcal{F}),
\]
avec les seules hypothèses du lemme précédent sur le couple $\mathcal{F}$, $\mathcal{G}$. Il faudrait que pour tout $S$ dans $\mathcal{M}$, et toute flèche
\msp{page 53}
$u : X \longrightarrow Y$ dans $\mathcal{M}/S$, on puisse insérer le diagramme obtenu dans $\mathcal{M}$
\[
\begin{tikzcd}[column sep=small]
X \arrow[rr, "u"] \arrow[dr, "f"'] & & Y \arrow[dl, "g"] \\
& S &
\end{tikzcd}
\]
\printedpage{31}
% head: A. Grothendieck : Dérivateurs, Ch. XVI, Édition des Universités de Montpellier II et Paris VII
% foot: Version 18 mars 2019
dans un diagramme commutatif ci-dessous
\[
\begin{tikzcd}[column sep=small]
X \arrow[rr, "u"] \arrow[d, "i"'] & & Y \arrow[d, "j"] \\
\bar{X} \arrow[rr, "\bar{u}"] \arrow[dr, "p"'] & & \bar{Y} \arrow[dl, "q"] \\
& S &
\end{tikzcd}
\tag*{(4.4.9.3)}
\]
avec
\[
f = pi, \qquad\quad g = qj, \qquad\quad i, j \in W^{\mathrm{u}}_S, \qquad\quad p, q \in \mathcal{G}.
\]
Ce sera le cas en tous cas si pour tout $S$, et un objet variable $X$ de $\mathcal{M}/S$, on peut trouver une factorisation \textit{fonctorielle en} $X$ de son morphisme structural $f_X$ en $p_X i_X$, avec $i_X \in W^{\mathrm{u}}_S$ et $p_X \in \mathcal{G}$, \textit{i.e.} si $i$, $p$ se déduisent d'un foncteur
\[
\varphi : \mathcal{F}_S \longrightarrow \mathcal{M}/S \ \text{\footnotemark},
\]
\footnotetext{$\mathcal{F}_S$ est la sous-catégorie pleine de $\mathcal{M}/S$ formée des $X$ tels que $f_X \in \mathcal{F}$.}
et d'un morphisme de foncteurs
\[
i : \underbrace{\inc_{\mathcal{F}_S}}_{\substack{\text{foncteur}\\ \text{d'inclusion}}} \longrightarrow \varphi.
\]
Admettant l'existence d'un diagramme ayant les propriétés dites, pour toute flèche $u$ dans $\mathcal{F}_S$, ou seulement toute flèche $u$ qui soit aussi dans $W_S$, nous pouvons prouver maintenant que $u \in W_S \Longrightarrow u \in W^{\mathrm{u}}_S$, pourvu que la propriété similaire soit vrai pour $\mathcal{G}$ à la place de $\mathcal{F}$.

\msp{page 54}
En effet, l'hypothèse sur $u \in W$ implique $\bar{u} \in W$, et l'hypothèse sur $\mathcal{G}$ implique dès lors $\bar{u} \in W^{\mathrm{u}}_S$. Comme $W^{\mathrm{u}}_S$ est un localiseur dans $\mathcal{M}/S$, on en conclut bien $u \in W^{\mathrm{u}}_S$, q.e.d.

Ainsi on a obtenu

\textbf{Corollaire 4.4.10.} \textit{Supposons, sous les conditions du lemme} 4.4.9, \textit{que tout diagramme commutatif triangulaire} $\begin{tikzcd}[column sep=small, row sep=small, ampersand replacement=\&]
X \arrow[rr, "u"] \arrow[dr, "f"'] \& \& Y \arrow[dl, "g"] \\
\& S \&
\end{tikzcd}$ \textit{s'insère dans un diagramme du type} (4.4.9.3) \textit{comme ci-dessus} (\textit{avec les relations explicites sur les flèches qui y figurent}). \textit{Alors}
\[
\mathrm{B}(\mathcal{G}) \;\Longrightarrow\; \mathrm{B}(\mathcal{F}).
\]
\textit{Donc si de plus} $\mathcal{G} \subset \mathcal{F}$, \textit{de sorte que} $\mathrm{A}(\mathcal{F}) \Longrightarrow \mathrm{A}(\mathcal{G})$ \textit{tautologiquement, et de même pour} B \textit{et pour} C, \textit{on trouve}
\[
\begin{aligned}
\mathrm{A}(\mathcal{F}) &\;\Longleftrightarrow\; \mathrm{A}(\mathcal{G})\\
\mathrm{B}(\mathcal{F}) &\;\Longleftrightarrow\; \mathrm{B}(\mathcal{G})\\
\mathrm{C}(\mathcal{F}) &\;\Longleftrightarrow\; \mathrm{C}(\mathcal{G}).
\end{aligned}
\]

Nous pouvons appliquer ce corollaire dans le cas qui nous occupe ici, avec
\[
\mathcal{G} \overset{\text{déf}}{=} \mathcal{F}_{W_0} \subset \mathcal{F} \overset{\text{déf}}{=} \Fibf_{W_0},
\]
\printedpage{32}
% head: A. Grothendieck : Dérivateurs, Ch. XVI, Édition des Universités de Montpellier II et Paris VII
% foot: Version 18 mars 2019
\msp{page 55}
vu que dans ce cas on a bien une factorisation \textit{fonctorielle} d'un $f \in \Fibf_{W_0}$, $f : X \longrightarrow S$, en $f = pi$, $i \in W^{\mathrm{u}}_{0\,S} \subset W^{\mathrm{u}}_S$, $p \in \mathcal{F}_{W_0} = \mathcal{G}$, savoir celle rappelée page 43 :
\[
\begin{tikzcd}[column sep=small]
X \arrow[rr, "i_X"] \arrow[dr, "f=f_X"'] & & \bar{X} = \Phi_S(X) \arrow[dl, "p_X"] \\
& S. &
\end{tikzcd}
\]

Cela prouve donc les équivalences annoncées (4.4.8.1) et achève d'établir la validité des implications du diagramme de la page 48.

Ces implications, à la seule exception de l'implication en pointillés $\mathrm{C}_{W_0} \dottedimpl \mathrm{A}_{W_0}$ (qui suppose qu'on ait $W_0$ (7)), ne font appel à aucune hypothèse particulière sur les localiseurs fondamentaux $W$ et $W_0$ (à part l'inclusion $W \supset W_0$). Mais nous pouvons, en utilisant le lemme 4.4.10, donner des implications conditionnelles remarquables, sous la seule condition que $W$, ou $W_0$, satisfassent à l'anodin \textit{axiome}
\msp{page 56}
\textit{des limites} Loc (8), et même seulement sous la forme faible de stabilité des localiseurs par $\varinjlim$ de \textit{suites}. Donnons-les d'abord sous forme technique, avant de déployer les conséquences dans un théorème scholie :

\textbf{Corollaire 4.4.11.} \textit{Supposons que} $W_0$ \textit{satisfasse à l'axiome} Loc (8), \textit{ou seulement} Loc (8 bis) (\textit{stabilité par} $\varinjlim$ \textit{de suites}). \textit{Alors} $\mathrm{D}_{W_0} \Longrightarrow \mathrm{A}_{W_0}$ (\textit{donc les conditions} $\mathrm{A}_{W_0}$, $\mathrm{B}_{W_0}$, $\mathrm{C}_{W_0}$, $\mathrm{D}_{W_0}$ \textit{sont équivalentes}). \textit{En particulier} (\textit{faisant} $W_0 = W$), \textit{si} $W$ \textit{satisfait à} Loc (8 bis), \textit{alors on a} $\mathrm{D}_W \Longrightarrow \mathrm{A}_W$ (\textit{donc les conditions} $\mathrm{A}_W$, $\mathrm{B}_W$, $\mathrm{C}_W$, $\mathrm{D}_W$ \textit{sont équivalentes}).

En effet, l'implication à prouver s'écrit aussi
\[
\mathrm{A}(\Parf_{W_0}) \;\Longleftrightarrow\; \mathrm{A}(\Fibf_{W_0}).
\]

L'implication non tautologique $\Longrightarrow$ est justiciable du lemme 4.4.9. Il faut vérifier que moyennant Loc (8 bis) pour $W_0$, tout $f \in \Fib_{W_0}$, $f : X \longrightarrow S$, se factorise en $f = pi$, avec $i \in W^{\mathrm{u}}_0$, $p \in \Parf_{W_0}$. Or on sait\footnote{et c'est là une trouvaille tout à fait inattendue, un peu comme le lemme de K. S. \textsc{Brown} \ldots} (VI, p. 193 ff.) % [sic] f \in \Fib_{W_0} -> f \in \Fibf_{W_0}, the implication to be proved concerns A(\Fibf_{W_0}) and the next line has \Fibf_{W_0}
\msp{page 57}
que l'axiome Loc (8 bis) implique que $f \in \Fibf_{W_0}$ se factorise en
\[
\begin{tikzcd}[column sep=small]
X \arrow[rr, "i_\infty"] \arrow[dr, "f"'] & & \bar{X} = \Phi_\infty(X) \arrow[dl, "p_\infty"] \\
& S, &
\end{tikzcd}
\qquad
\text{\parbox[t]{0.33\textwidth}{\footnotesize (catégorie de chemins infinis d'un côté, relativisée sur $S$, \textit{i.e.} des couples $(x,c)$, avec $x \in \Ob X$ et $c : f(x) \longrightarrow s$ un chemin infini dans $S$, de source $f(x)$.)}}
\]
avec $p_\infty \in \Parf_{W_0}$ (comme on veut), et $i_\infty \in W^{\mathrm{u}}_{0\,S}$, et comme $W^{\mathrm{u}}_{0\,S} \subset W^{\mathrm{u}}_S$, on a bien $i_\infty \in W^{\mathrm{u}}_S$.

\textbf{Corollaire 4.4.12.} \textit{Supposons que} $W$ et $W_0$ \textit{satisfassent à} Loc (8 bis). \textit{Alors on a} $\mathrm{D}_{W_0} \Longrightarrow \mathrm{A}$ (\textit{donc les huit conditions} A, B, C, D, $\mathrm{A}_{W_0}$, $\mathrm{B}_{W_0}$, $\mathrm{C}_{W_0}$, $\mathrm{D}_{W_0}$ \textit{sont équivalentes}).

\printedpage{33}
% head: A. Grothendieck : Dérivateurs, Ch. XVI, Édition des Universités de Montpellier II et Paris VII
% foot: Version 18 mars 2019
Cela signifie en effet que
\[
\mathrm{A}(\Parf_{W_0}) \;\Longrightarrow\; \mathrm{A}(\Fibf_W),
\]
et c'est encore justiciable du lemme 4.4.9. En effet, la factorisation ci-dessus d'une flèche $f$ ne dépend pas du choix d'un localiseur, et $p_\infty$ est $W$-parfait pour \textit{tout} localiseur satisfaisant Loc (8 bis), donc en particulier pour $W_0$. Et la flèche $i_\infty$ est dans $W^{\mathrm{u}}_S$, pour tout localiseur satisfaisant Loc (8 bis), donc pour $W_0$.

\msp{page 58}
Je vais récapituler les résultats obtenus dans cette section dans un théorème, qui constituera le résultat principal du présent paragraphe sur les axiomes de $W$-fibrations.

\textbf{Théorème 4.4.13.} \textit{Soit} $W$ \textit{un localiseur fondamental dans} Cat. \textit{Pour tout localiseur} $W_0$ \textit{contenu dans} $W$, \textit{on a explicité un ensemble de quatre conditions sur la paire} $W_0$, $W$, \textit{qui pour} $W$ \textit{fixé, dépendent de} $W_0$ \textit{seul, notées}
\[
\mathrm{A}_{W_0},\ \mathrm{B}_{W_0},\ \mathrm{C}_{W_0},\ \mathrm{D}_{W_0},
\tag*{(4.4.13.1)}
\]
\textit{et définies respectivement comme}
\[
\mathrm{A}(\mathcal{F}_{W_0}),\ \mathrm{B}(\mathcal{F}_{W_0}),\ \mathrm{C}(\mathcal{F}_{W_0}), \qquad\quad \textit{et}
\]
\[
\mathrm{A}(\Parf_{W_0}) \Bigl(\Longleftrightarrow\; \mathrm{B}(\Parf_{W_0}) \;\Longleftrightarrow\; \mathrm{C}(\Parf_{W_0})\Bigr),
\]
\textit{avec les notations} $\mathrm{A}(\cdot)$, $\mathrm{B}(\cdot)$, $\mathrm{C}(\cdot)$ \textit{introduites dans} (4.4.1). \textit{Faisant} $W_0 = W$, \textit{on trouve les conditions sur} $W$
\[
\mathrm{A},\ \mathrm{B},\ \mathrm{C},\ \mathrm{D}
\tag*{(4.4.13.2)}
\]
\textit{avec} $\mathrm{A} = \mathrm{A}_W$, $\mathrm{B} = \mathrm{B}_W$, $\mathrm{C} = \mathrm{C}_W$, $\mathrm{D} = \mathrm{D}_W$, \textit{dont les trois premières ne sont autres que} Loc (7 a), Loc (7 b), Loc (7 c) (cf. 4.3, \textit{page} 38).

\circled{1} \textit{On a pour ces conditions les implications suivantes.}

\begin{enumerate}[leftmargin=2.5em]
\item[a)] \textit{On a les implications}
\[
\mathrm{A}_{W_0} \;\Longrightarrow\; \mathrm{B}_{W_0} \;\Longrightarrow\; \mathrm{C}_{W_0} \;\Longrightarrow\; \mathrm{D}_{W_0},
\tag*{(4.4.13.3)}
\]
(\textit{d'où en particulier}

\msp{page 59}
\[
\mathrm{A}_W \;\Longrightarrow\; \mathrm{B}_W \;\Longrightarrow\; \mathrm{C}_W \;\Longrightarrow\; \mathrm{D}_W\ ).
\tag*{(4.4.13.4)}
\]
\item[b)] \textit{Si} $W'_0 \subset W_0$ \textit{est un localiseur fondamental contenu dans} $W_0$, \textit{on a} (\textit{tautologiquement})
\[
\mathrm{A}_{W_0} \;\Longrightarrow\; \mathrm{A}_{W'_0}, \qquad\quad \mathrm{B}_{W_0} \;\Longrightarrow\; \mathrm{B}_{W'_0}, \qquad\quad \mathrm{C}_{W_0} \;\Longrightarrow\; \mathrm{C}_{W'_0}.
\]
\textit{On trouve donc, par} (a) \textit{et} (b) (\textit{y faisant pour le couple} $(W, W_0)$ \textit{comme} $(W_0, W'_0)$),
\end{enumerate}


\endgroup
% END UNIT der-ch16/der-ch16_p031-035.tex
% BEGIN UNIT der-ch16/der-ch16_p036-040.tex
\begingroup\makeatletter%
\newcommand{\printedpage}[1]{\marginpar{\footnotesize p.~#1}}%
\newcommand{\msp}[1]{\marginpar{\footnotesize [#1]}}%
\newcommand{\cercle}[1]{\tikz[baseline=(c.base)]\node[draw,circle,inner sep=1pt](c){#1};}%
\newcommand{\fourche}{\begin{tikzcd}[ampersand replacement=\&, row sep=0.5em, column sep=1em, baseline=(\tikzcdmatrixname-2-1.base)]%
 \& X_1 \\%
 X_0 \arrow[ur] \arrow[dr] \& \\%
 \& X_2%
\end{tikzcd}}%
\newcommand{\carre}{\begin{tikzcd}[ampersand replacement=\&, row sep=0.5em, column sep=1em, baseline=(\tikzcdmatrixname-2-1.base)]%
 \& X_1 \arrow[dr] \& \\%
 X_0 \arrow[ur] \arrow[dr] \& \& X \\%
 \& X_2 \arrow[ur] \&%
\end{tikzcd}}%
\makeatletter%
\renewcommand{\@makefnmark}{\hbox{\textsuperscript{\normalfont(\@thefnmark)}}}%
\renewcommand{\@makefntext}[1]{\parindent 1em\noindent\hbox to 1.8em{\hss\textsuperscript{\normalfont\@thefnmark}}\,#1}%
\makeatother%
\makeatother\let\printedpage\derpp\let\msp\dermsp%
\let\mspin\DERmspin%


\begin{enumerate}[leftmargin=2.5em]
\item[] \printedpage{34}%
% head: A. Grothendieck : Dérivateurs, Ch. XVI, Édition des Universités de Montpellier II et Paris VII
\textit{le rectangle d'implications}
\[
\begin{tikzcd}[row sep=large, column sep=large, execute at end picture={
\draw[decorate,decoration={brace,amplitude=6pt}] ([xshift=-14pt]A.west |- C.south) -- ([xshift=-14pt]A.north west) node[midway,left=8pt]{$\mathcal{F}_W,\mathrm{Fibf}_W$};
\draw[decorate,decoration={brace,amplitude=6pt}] ([xshift=14pt]AW.north east) -- ([xshift=14pt]AW.east |- CW.south) node[midway,right=8pt]{$\mathcal{F}_{W_0},\mathrm{Fibf}_{W_0}$};
}]
|[alias=A]| \mathrm{A} \arrow[r,Rightarrow] \arrow[d,Rightarrow] & |[alias=AW]| \mathrm{A}_{W_0} \arrow[d,Rightarrow] \\
\mathrm{B} \arrow[r,Rightarrow] \arrow[d,Rightarrow] & \mathrm{B}_{W_0} \arrow[d,Rightarrow] \\
|[alias=C]| \mathrm{C} \arrow[r,Rightarrow] \arrow[d,Rightarrow] & |[alias=CW]| \mathrm{C}_{W_0} \arrow[d,Rightarrow] \\
\underbrace{\mathrm{D}}_{\mathrm{Parf}_W} \arrow[r,Rightarrow] & \underbrace{\mathrm{D}_{W_0}}_{\mathrm{Parf}_{W_0}}\,,
\end{tikzcd}
\tag*{(4.4.13.5)}
\]
\textit{où on a rappelé, à côté des accolades, quels sont l'ensemble, ou les ensembles $\mathcal{F}\subset\mathrm{Fl}(\mathrm{Cat})$ qui définissent les conditions envisagées (comme l'une des conditions $\mathrm{A}(\mathcal{F})$, $\mathrm{B}(\mathcal{F})$, $\mathrm{C}(\mathcal{F})$ ou $\mathrm{D}(\mathcal{F})=\mathrm{A}(\mathcal{F})$ dans le cas où les trois conditions sont équivalentes). La plus faible des huit conditions envisagées est $\mathrm{D}_{W_0}$, la plus forte} A. (\textit{Et pour $W_0$ variable, on trouve la condition la plus faible de toutes en prenant $D_{W_\omega}$, où $W_\omega$ est le plus petit des localiseurs fondamentaux dans} Cat.)

\item[c)] \textit{On a de plus, pour la condition la plus faible $\mathrm{D}_{W_0}$, l'implication}
\[
\mathrm{D}_{W_0} \implies \mathrm{F},
\]
\msp{page 60}%
\textit{où} $\mathrm{F}=\text{Loc (7 F)}$ \textit{est la condition que} Cat \textit{possède suffisamment de $W$-fibrations,} i.e. \textit{que c'est une catégorie de fibrations pour} $(W,\mathrm{Fib}_W)$.
\end{enumerate}

\noindent\cercle{2}\enspace{\itshape Si $W$ et $W_0$ satisfont à \textup{Loc (8 bis)} (axiome des limites dénombrables), alors \textup{$\mathrm{D}_{W_0}\Longrightarrow\mathrm{A}$, i.e.} les 8 conditions envisagées sont équivalentes. Donc si $W$ satisfait \textup{Loc (8 bis)}, considérant le plus petit localiseur fondamental $W_L$ satisfaisant à \textup{Loc (8 bis)} (de sorte que $W\supset W_L$), on voit que les huit conditions
\[
\mathrm{A},\ \mathrm{B},\ \mathrm{C},\ \mathrm{D},\ \mathrm{A}_{W_L},\ \mathrm{B}_{W_L},\ \mathrm{C}_{W_L},\ \mathrm{D}_{W_L}
\]
sont équivalentes~\footnote[33]{N.B. La forme la plus faible, $\mathrm{D}_{W_{\mathrm{L}}}$, de ces conditions, est la relation \fbox{$\mathrm{Parf}_{W_{\mathrm{L}}}\cap W\subset W^{\mathrm{u}}$.}}. En particulier, sous la même hypothèse (fort anodine) sur $W$, les conditions \textup{Loc (7) $=$ Loc (7 a), Loc (7 b)} et \textup{Loc (7 c)} sont équivalentes, et elles équivalent aussi à la condition \textup{$\mathrm{D}=\mathrm{A}(\mathrm{Parf}_W)$, i.e.} à
\[
\mathrm{Parf}_W\subset\mathrm{Fib}_W \qquad (\textit{cas particulier de } \text{Loc (7 a)}),
\]
et à ses deux variantes d'apparence plus faible \textup{$\mathrm{B}(\mathrm{Parf}_W)$, $\mathrm{C}(\mathrm{Parf}_W)$} (cas particulier de \textup{Loc (7 b) et de Loc (7 c)} respectivement).}

% foot: Version 18 mars 2019

\noindent\cercle{3}\enspace{\itshape Supposons que $W_0$ satisfasse à \textup{Loc (8 bis)} (sans hypothèse particulière sur $W$). Alors
\printedpage{35}%
% head: A. Grothendieck : Dérivateurs, Ch. XVI, Édition des Universités de Montpellier II et Paris VII
\msp{page 61}%
les quatre conditions $\mathrm{A}_{W_0}$, $\mathrm{B}_{W_0}$, $\mathrm{C}_{W_0}$, $\mathrm{D}_{W_0}$ sont équivalentes, et le diagramme d'implications de \textup{1°)} (p. \textup{59}) se réduit donc au diagramme linéaire
\[
\mathrm{A} \implies \mathrm{B} \implies \mathrm{C} \implies \mathrm{D} \implies \mathrm{D}_{W_0}\ (\Longleftrightarrow \mathrm{A}_{W_0}\ldots),
\]
dont le premier terme est \textup{Loc (7 a) $=$ Loc (7)}, et le dernier équivaut à l'inclusion
\[
\mathcal{F}_{W_0}\subset\mathrm{Fib}_W,
\]
cas particulier de \textup{$\mathrm{A}=$ Loc (7 a)}, ou encore à l'inclusion (en apparence plus forte)
\[
\mathrm{Fibf}_{W_0}\subset\mathrm{Fib}_W,
\]
(cas particulier de la forme équivalente
\[
\mathrm{Fibf}_W\subset\mathrm{Fib}_W,
\]
de l'axiome} Loc (6 a)).

\noindent\cercle{4}\enspace\textit{Enfin, on a les implications plus techniques conditionnelles suivantes.}
\begin{enumerate}[leftmargin=2.5em]
\item[a)] \textit{Si $W_0$ satisfait à Loc (7 F), alors les conditions $\mathrm{A}_{W_0}$, $\mathrm{B}_{W_0}$, $\mathrm{C}_{W_0}$ sont équivalentes, donc le diagramme (4.4.13.5) se réduit au diagramme}
\[
\begin{tikzcd}[row sep=large, column sep=large]
\mathrm{A} \arrow[d,Rightarrow] & \\
\mathrm{B} \arrow[d,Rightarrow] & \\
\mathrm{C} \arrow[r,Rightarrow] \arrow[d,Rightarrow] & \mathrm{C}_{W_0} \arrow[d,Rightarrow] \\
\mathrm{D} \arrow[r,Rightarrow] & \mathrm{D}_{W_0}\ .
\end{tikzcd}
\]

\item[b)] \textit{Si $W$ est fortement saturé} (\textit{axiome} Loc (1 ter), \textit{qui}
\msp{page 62}%
\textit{est satisfait dans la grande majorité des cas, et est même impératif pour construire un prédérivateur satisfaisant à l'axiome de localisation} Der 2) \textit{et satisfait à} Loc (3 bis) (\textit{axiome de localisation fort,} cf. §1), \textit{alors on a} $\mathrm{B}\Longrightarrow\mathrm{A}$, \textit{donc}
\[
\mathrm{A} \iff \mathrm{B},
\]
i.e. Loc (7 a) $=$ Loc (7) \textit{équivaut à} Loc (7 b).
\end{enumerate}

\noindent\textsc{Démonstration}. Tout ce qui est énoncé a été vu dans les pages précédentes, à la seule exception de 4°, b. Mais cette implication $\mathrm{B}\Longrightarrow\mathrm{A}$ pour $W$ fortement saturé est prouvé dans XII (prop. 1, p. 54). Je rappelle que la démonstration utilise de façon essentielle le % [sic] prouvé -> prouvée, accord avec « implication »
% foot: Version 18 mars 2019
\printedpage{36}%
% head: A. Grothendieck : Dérivateurs, Ch. XVI, Édition des Universités de Montpellier II et Paris VII
formalisme du prédérivateur $\mathrm{HOT}_W$ (construit et étudié surtout dans VI), et la théorie des foncteurs $f_!$, et notamment le résultat (valable si $W$ satisfait à Loc (3 bis, 7 b)) que pour une flèche $f:S'\longrightarrow S$ dans $W$, la flèche correspondante
\[
f^*_{\mathrm{lc}}:\mathrm{HOT}^{\mathrm{lc}}_W(S)\longrightarrow\mathrm{HOT}^{\mathrm{lc}}_W(S')
\]
est une équivalence de catégories. Donc l'implication $\mathrm{B}\Longrightarrow\mathrm{A}$, elle aussi, est un résultat qui n'a rien d'évident.

\noindent\msp{page 63}%
4.4.14. Je résume les implications obtenues jusqu'à présent, y compris les implications conditionnelles, dans le diagramme récapitulatif suivant :
\[
\begin{tikzpicture}[baseline=(C.base),font=\small]
\node[draw,circle,inner sep=1pt] (A) at (0,0) {$\mathrm A$};
\node (AW) at (2,0) {$\mathrm A_{W_0}$};
\node (B) at (0,-1.2) {$\mathrm B$};
\node (BW) at (2,-1.2) {$\mathrm B_{W_0}$};
\node (C) at (0,-2.4) {$\mathrm C$};
\node (CW) at (2,-2.4) {$\mathrm C_{W_0}$};
\node (D) at (0,-3.6) {$\mathrm D$};
\node (DW) at (2,-3.6) {$\mathrm D_{W_0}$};
\node (F) at (7.8,-3.6) {$\mathrm F$};
\node[draw,circle,inner sep=1pt] (Ab) at (2,-4.8) {$\mathrm A$};
\node[anchor=west] at (2.25,-4.8) {$\,,$};
\foreach \p/\q in {A/AW,B/BW,C/CW,D/DW,A/B,B/C,C/D,AW/BW,BW/CW,CW/DW,DW/F}
  \draw[double,double distance=1pt,-{Implies}] (\p) -- (\q);
\draw[dotted,->] (B.west) to[out=145,in=215] node[left,font=\scriptsize] {$W$ (1 ter, 3 bis)} (A.west);
\draw[dotted,->] (CW.east) to[out=70,in=-70] node[right,font=\scriptsize] {$W_0$ (7)} (AW.east);
\draw[dotted,->] (DW.east) to[out=60,in=-40,looseness=1.45] node[right,font=\scriptsize] {$W_0$ (8 bis)} (AW.east);
\draw[dotted,->] (D.south east) -- node[below left,font=\scriptsize] {$W$ (8 bis)} (Ab.north west);
\draw[dotted,->] (DW) -- node[right,font=\scriptsize] {$W$ (8 bis) $+$ $W_0$ (8 bis)} (Ab);
\node[anchor=south west,align=left,font=\scriptsize] at (3.4,-3.45) {N.B. on verra sur l'exemple de $\mathbf{W}_0$\\
que cette implication est \textit{stricte,}\\
même sous l'hypothèse\\
Loc (1 ter, 2 bis, 4, 5, 6 a, 8, 9).};
\end{tikzpicture}
\]
où j'ai cerclé la condition la plus forte et la plus importante de toutes, $\mathrm{A}=$ Loc (7), qui apparaît en deux extrémités opposées du diagramme.

Je voudrais encore faire un commentaire sur les ingrédients des résultats prouvés ici, et notamment de l'équivalence des huit conditions A à $\mathrm{D}_{W_0}$, quand $W$ et $W_0$ satisfont Loc (8 bis) :

\msp{page 64}%
\begin{enumerate}[leftmargin=2.5em]
\item[1)] La théorie des \textit{catégories de chemins} \underline{Ch}$(X)$ (VII), s'appuyant sur celle des \textit{morphismes lisses, propres, parfaits.}
\item[2)] Le lemme de K. S. \textsc{Brown} (3.2.1, page 20).
\item[3)] Pour l'implication $\mathrm{B}\implies\mathrm{A}$, \textit{i.e.} Loc (7 b) $\implies$ Loc (7) (sous l'hypothèse $W$ (1 ter, 3 bis)), j'utilise la théorie du prédérivateur $\mathrm{Hot}_W$ (VI), et plus particulièrement celle des foncteurs $f_!$ (définis sous la seule hypothèse $W$ (3 bis), renforçant la condition automatique $W$ (3) (faisant partie de la définition d'un localiseur fondamental)).
\end{enumerate}
% foot: Version 18 mars 2019

\leavevmode\printedpage{37}%
% head: A. Grothendieck : Dérivateurs, Ch. XVI, Édition des Universités de Montpellier II et Paris VII
Il est clair que je ne serais pas arrivé à y voir clair, sans le travail fondamental de K. S. \textsc{Brown}. Il reste pourtant encore des questions ouvertes. Ainsi, si on ne suppose pas que $W$ satisfasse à Loc (8 bis), mais au besoin à Loc (3 ter) (et je n'ai pas eu l'occasion encore de travailler avec un localiseur qui n'y satisfasse \ldots), est-il vrai néanmoins que l'on a encore l'implication
\[
\mathrm{D} \qquad \Longrightarrow \qquad \mathrm{A},
\]
de façon que les conditions A, B, C, D soient équivalentes (et en particulier Loc (7 c) implique déjà Loc (7) dans toute sa force) ? Il est clair dans tout cela que
\msp{page 65}%
je ne me suis pas suffisamment attaché à développer des exemples et des contre-exemples. Sans compter qu'il me semble que tous les localiseurs fondamentaux sur Cat qu'on rencontre “en pratique” satisfont déjà à Loc (8) et a fortiori à Loc (8 bis), donc pour eux on a bien l'implication escomptée. La même remarque s'applique à l'implication hypothétique
\[
\mathrm{D}_{W_0} \overset{?}{\implies} \mathrm{A},
\]
sous hypothèse Loc (8 bis) sur $W_0$ et sur $W$. Un candidat évident pour un contre-exemple déjà à la première question serait $W_\omega=$ plus petit localiseur fondamental sur Cat, un qu'il est presque évident qu'il ne satisfait \textit{pas} à Loc (8) ni même Loc (8 bis). Mais il n'y a aucune raison qu'il satisfasse non plus à F, de telle façon que (s'il n'y satisfait pas, donc non plus à D) l'implication $\mathrm{D}\Longrightarrow\mathrm{A}$ est bel et bien valable. Cela nous amène à la question, plus substantielle, si on n'aurait pas l'implication % [sic] un qu'il est -> dont il est (?), construction défectueuse
\msp{page 66}%
\[
\mathrm{F} \overset{?}{\implies} \mathrm{A}\ \text{\footnotemark[34]},
\]
\footnotetext[34]{C'est faux, \textit{cf.} plus bas, l'implication Loc (7 $\mathrm{D}_{W_\omega}$) $\Longrightarrow$ Loc (7 F) est \textit{stricte.}}%
du moins sous des conditions convenables assez anodines sur $W$, genre Loc (1 ter, 3 bis, 8 bis).

C'est sans doute là la question la plus substantielle et la plus intéressante qui est soulevée par le diagramme de la page 63, et qui ne fait plus appel à l'introduction d'un localiseur fondamental $W_0$ auxiliaire. Vu le lemme de Brown, il ne me semble pas absolument exclu qu'un tel miracle ait bel et bien lieu !

4.4.15 (Digression sur l'axiome Loc (6) des cofibrations.) Cela attire mon attention sur la question en quelque sorte duale : Si Cat a assez de $W$-\textit{co}fibrations, \textit{i.e.} si muni de $(W,\mathrm{Cof}_W)$ c'est une catégorie à cofibrations, cela implique-t-il déjà l'axiome Loc (6 a) (qui avait été dégagé, essentiellement, pour assurer qu'il en soit bien ainsi). La proposition 3.6. (p. 25) nous dit qu'il revient au même de prouver que, si $W$ satisfait Cof, alors les carrés de M. V. ouverts sont
\msp{page 67}%
$W$-cocartésien. Mais il me semble à présent que nous pouvons le prouver, sous la seule condition supplémentaire Loc (3 bis)~\footnote[35]{et peut-être, de plus, Loc (1 ter), \textit{i.e.} $W$ fortement saturée.}, laquelle assure l'existence des foncteurs $f_!$, et % [sic] cocartésien -> cocartésiens, accord avec « les carrés » ; [sic] (note 35) saturée -> saturé, accord avec « W » (localiseur)
% foot: Version 18 mars 2019
\printedpage{38}%
% head: A. Grothendieck : Dérivateurs, Ch. XVI, Édition des Universités de Montpellier II et Paris VII
en particulier la construction de Holim d'un diagramme $\fourche$ \quad comme $\int\fourche$ \quad .
Donc dire qu'un diagramme de $\mathcal{M}$ ouvert $\carre$ est $W$-cocartésien, doit revenir à dire que la flèche canonique
\[
\int\fourche \qquad \longrightarrow \qquad X
\]
est dans $W$ ; du moins si $W$ est supposé fortement saturé. Or cela a été bel et bien vu dans XII 3.4 (théorème 4 et son \textit{corollaire} 3, p. 149). Dans loc. cit. j'énonce les conclusions sous l'hypothèse que j'appelais W (7 bis) (l'actuel Loc (6)), mais je viens de vérifier que pour sa partie \textit{directe,} donc aussi pour ses corollaires 1 à 3, on n'a pas besoin
\msp{page 68}%
de cette hypothèse (comme je l'indique en annotations marginales).

4.16. Le résultat que je viens de prouver (sans vérifier avec soin la démonstration) me paraît aussi délicat, qu'il est satisfaisant. Cela rend la “conjecture” hésitante de 4.4.14, savoir que
\[
\mathrm{F} \overset{?}{\Longrightarrow} \mathrm{A} \qquad \text{donc} \qquad \mathrm{F}\Longleftrightarrow\mathrm{A}
\]
en supposant que $W$ satisfait Loc (1 ter, 3 bis, 8 bis), nettement plus engageante. Car on commence à s'habituer aux miracles les plus époustouflants \ldots

Pourtant, je viens de vérifier que cette conjecture est fausse, sur l'exemple de
\[
\mathbf{W}_0=\{f\in\mathrm{Fl}(\mathrm{Cat})\mid\pi_0(f)\ \mathrm{iso}\}\,,
\]
qui satisfait à F, mais pas à A, ni même à $\mathrm{D}_{W_\omega}$. Ainsi l'implication
\[
\mathrm{D}_{W_\omega} \implies \mathrm{F}
\]
et a fortiori les implications
\[
\left\{
\begin{aligned}
&\text{Loc (7 a)} \implies \mathrm{F},\ \text{Loc (7 b)} \implies \mathrm{F},\ \text{Loc (7 c)} \implies \mathrm{F}\\
&\mathrm{D}_{W_0} \implies \mathrm{F}
\end{aligned}
\right.
\]
sont \textit{strictes,} même en se limitant à des localiseurs fondamentaux $W$ satisfaisant (comme $\mathbf{W}_0$) à
\[
\text{Loc (1 ter, 2, 3 bis, 4, 5, 6 a, 8, 9)}\ \text{\footnotemark[36]}.
\]
\footnotetext[36]{\textbf{N.B. $\mathbf{W}_0$} ne satisfait \textit{pas} Loc (6 b), ni Loc (7 $\mathrm{D}_{W_\omega}$) (\textit{cf.} plus bas), par contre Loc (7 F).}
% foot: Version 18 mars 2019


\endgroup
% END UNIT der-ch16/der-ch16_p036-040.tex
% BEGIN UNIT der-ch16/der-ch16_p041-045.tex
\begingroup\makeatletter%
\newcommand{\printedpage}[1]{\marginpar{\footnotesize p.~#1}}%
\newcommand{\msp}[1]{\marginpar{\footnotesize [#1]}}%
\newcommand{\uFl}{\underline{\mathrm{Fl}}}%
\newcommand{\uCat}{\underline{\mathrm{Cat}}}%
\newcommand{\Parf}{\mathrm{Parf}}%
\makeatother\let\printedpage\derpp\let\msp\dermsp%
\let\mspin\DERmspin%


\printedpage{39}
% head: A. Grothendieck : Dérivateurs, Ch. XVI, Édition des Universités de Montpellier II et Paris VII

\msp{page 69}
Que $\mathbf{W}_0$ satisfasse à cet ensemble d'axiomes se vérifie par AQT \texttt{[âne qui trotte]}. (\textbf{N.B.} Qu'il satisfasse à Loc (1 ter, 2, 3 bis) est déjà conséquence du fait que $\mathbf{W}_0$ est de la forme $W_{\mathbf{D}}$, où $\mathbf{D}$ ($=\mathbf{D}_{\mathbf{W}_0}$ en l'occurrence) est un prédérivateur satisfaisant Der (3, 4, 5), ce qui est bien le cas pour $\mathbf{D}_{\mathbf{W}_0}$, qui est un vrai \textit{dérivateur.}) Il faut vérifier que tout $f$ dans Cat se factorise en $f = pi$, avec $i \in \mathbf{W}_0$, $p$ une $\mathbf{W}_0$-fibration. Or si $f : X \longrightarrow Y$, $I = \pi_0(X)$, on prend $\bar{X} = \coprod_{\alpha\in I} Y_{f(\alpha)}$, d'où une factorisation $f = pi$, où $\pi_0(i)$ est un isomorphisme, \textit{i.e.} $i \in \mathbf{W}_0$, et où $p : \bar{X} \longrightarrow Y$ est, au dessus de la composante connexe $Y_j$ de $Y$, donné par le produit $\bar{X}|Y_j \simeq Y_j \times \underbrace{\pi_0(f)^{-1}(j)}_{\text{catégorie discrète}}$ . Le morphisme $p$ est bien une $\mathbf{W}_0$-fibration, car grâce à Loc (4), il suffit de le vérifier sur chaque $Y_j$, et là c'est déduit d'une flèche $I_j \longrightarrow e$ (une $\mathbf{W}_0$-fibration) par changement de base $Y_j \longrightarrow e$.

D'autre part, que $\mathbf{W}_0$ ne satisfait \textit{pas} même à $\mathrm{D}_{W_\omega}$ a déjà été vu (I, p. 87),
\msp{page 70}
en réalisant l'application du $n$-revêtement $f : S^1 \longrightarrow S^1$ dans Cat par un morphisme $W_\omega$-parfait (en fait, par une Cat-fibration localement triviale), chose très facile\footnote[37]{Autre façon de le voir : on \textit{sait} (XII) que $\mathrm{W}$ (5, 7 $\mathrm{D}_W$) $\Longrightarrow W \subset \mathbf{W}_\infty$, or on a $\mathbf{W}_\infty \subset \mathbf{W}_0$, et l'inclusion est \textit{stricte\,!}}.

Je pense qu'à présent l'ensemble des différentes variantes de l'axiome des $W$-fibrations est à peu près compris, à cela près qu'il n'est toujours pas clair, en dehors de l'axiome Loc (8 bis), si les implications du diagramme
\[
\begin{tikzcd}[row sep=large, column sep=large]
\mathrm{A} \arrow[r, Rightarrow] \arrow[d, Rightarrow] & \mathrm{A}_{W_0} \arrow[d, Rightarrow] \\
\mathrm{B} \arrow[r, Rightarrow] \arrow[d, Rightarrow] & \mathrm{B}_{W_0} \arrow[d, Rightarrow] \\
\mathrm{C} \arrow[r, Rightarrow] \arrow[d, Rightarrow] & \mathrm{C}_{W_0} \arrow[d, Rightarrow] \\
\mathrm{D} \arrow[r, Rightarrow] & \mathrm{D}_{W_0}
\end{tikzcd}
\]
sont strictes ou non. Mais il me semble que tous les localiseurs fondamentaux intéressants rencontrés jusqu'à présent satisfont Loc (8), et même Loc (1 ter, 2, 3 bis, 4, 5, 6 a, 8, 9) \marginpar{?}(peut-être même Loc (7 F) - il faudrait que je teste sur les $\mathbf{W}_n$), donc la perplexité non résolu me paraît sans grande conséquence. % [sic] résolu -> résolue, accord avec « la perplexité » ; « ? » imprimé dans la marge gauche à la hauteur de cette ligne

\msp{page 71}

\printedpage{40}
% head: A. Grothendieck : Dérivateurs, Ch. XVI, Édition des Universités de Montpellier II et Paris VII

\section*{5\quad Axiome des limites Loc (8), et l'axiome Loc (9) d'accessibilité}

C'est l'axiome :

\textbf{Loc (8)}\quad $W$ est stable par $\varinjlim$ filtrantes.

Une variante affaiblie utile est la suivante :

\textbf{Loc (8 bis)}\quad $W$ est stable par $\varinjlim$ dénombrables.

Explicitement, si on a deux systèmes inductifs filtrants dans Cat
\[
\begin{aligned}
\underline{X} &= (X_i)_{i\in I}\\
\underline{Y} &= (Y_i)_{i\in I},
\end{aligned}
\]
et une flèche
\[
\underline{f} = (f_i)_{i\in I} : \underline{X} \longrightarrow \underline{Y},
\]
donnée par des $f_i : X_i \longrightarrow Y_i$, posant $X = \varinjlim X_i$, $Y = \varinjlim Y_i$, $f = \varinjlim f_i : X \longrightarrow Y$, Loc (8) nous dit que si les $f_i \in W$, alors $f \in W$. Ici $I$ est
\msp{page 72}
une (petite) catégorie filtrante quelconque, et quitte à la remplacer par une catégorie cofinale, on peut supposer $I$ un ensemble ordonné. L'axiome Loc (8 bis) correspond au cas où $I = \mathbf{N}$.

Notons que grâce à la construction de la catégorie des chemins, nous avons une factorisation, fonctorielle en $f$, de toute flèche $f : X \longrightarrow Y$ de Cat en $f = pi$ :
\[
X \xrightarrow{\quad i = i_f\quad} \bar{X}(f) = \bar{X} \xrightarrow{\quad p = p_f\quad} Y,
\]
avec $i \in W_\omega$, $p \in \Parf_{W_\omega}$, a fortiori $i \in W$, $p \in \Parf_W$. Ce foncteur $f \longmapsto (\bar{X}(f), p_f)$ commute aux limites inductives. Ainsi, si on applique ceci à chacun des $f_\alpha : X_\alpha \longrightarrow Y_\alpha$, le factorisant en
\[
X_\alpha \xrightarrow{i_\alpha} \bar{X}_\alpha \xrightarrow{p_\alpha} Y_\alpha,
\]
on trouve un système inductif de diagrammes de type $\boldsymbol{\Delta}^2$ (factorisant les $f_i$), et en passant à la $\varinjlim$ on trouve
\[
X \xrightarrow{i} \bar{X} \xrightarrow{p} Y,
\]
qui n'est autre que la factorisation canonique
\msp{page 73}
de $f = \varinjlim f_i$ en $pi$, avec $i \in W_\omega$, $p \in \Parf_{W_\omega}$ (N.B. Alors que ni $W_\omega$, ni $\Parf_{W_\omega}$ ne sont stables par $\varinjlim$ filtrantes, à coup sûr\,!). On a donc $i \in W$, et la condition $f \in W$, \textit{i.e.} $pi \in W$, équivaut donc à $p \in W$. Cela nous ramène, pour vérifier l'axiome Loc (8) (resp. Loc (8 bis)) au cas où les $f_i$ sont dans $\Parf_{W_\omega}$.

Supposons maintenant que $W$ satisfasse à $\mathrm{D}_{W_\omega}$ (la plus faible de toutes les variantes envisagées ici de l'axiome des $W$-fibrations Loc (7), à la seule exception de F). Alors on a
\[
p \in W \qquad\quad \Longleftrightarrow \qquad\quad \text{les fibres de $p$ sont $W$-asphériques,}
\]
\printedpage{41}
% head: A. Grothendieck : Dérivateurs, Ch. XVI, Édition des Universités de Montpellier II et Paris VII
l'implication $\Longrightarrow$ résultant de l'hypothèse $\mathrm{D}_{W_\omega}$ $\Parf_{W_\omega} \cap W \subset W^{\mathrm{u}}$, et $\Longleftarrow$ résultant du sorite des morphismes $W$-lisses \textit{etc.}, via Loc (3). De même, l'hypothèse que les $f_i \in W$ équivaut à celle que les fibres des $p_i$ sont $W$-asphériques. D'autre part, il est immédiat que
\msp{page 74}
toute fibre de $p$ est isomorphe à la $\varinjlim\limits_{I}$ des fibres correspondantes des $p_\alpha$. De cette façon, on trouve :

\textbf{Proposition 5.1.} \textit{Supposons que $W$ satisfasse l'axiome des $W$-fibrations} Loc (7 $\mathrm{D}_{W_\omega}$) (\textit{ce qui est donc le cas s'il satisfait un des axiomes plus forts} Loc (7 a \textit{ou} b \textit{ou} c)). \textit{Pour que $W$ satisfasse à} Loc (8) (resp. \textit{à} Loc (8 bis)), \textit{il faut et il suffit que l'ensemble des objets dans} Cat \textit{qui sont $W$-asphériques soit stable par $\varinjlim$ filtrante} (resp. \textit{par $\varinjlim$ de suites dénombrables}).

5.2. Notons que la catégorie Cat est un paratopos, a fortiori elle est \textit{accessible} (\textit{cf.} XVIII, §8), et il en est donc de même de $\uFl(\mathrm{Cat})$. Ainsi on a dans $\mathrm{Fl}(\mathrm{Cat}) = \mathrm{Ob}\,\uFl(\mathrm{Cat})$ la notion de \textit{partie accessible.} On dira donc qu'une partie $W$ de $\mathrm{Fl}\,\mathrm{Cat}$ (plus généralement d'une catégorie $\mathrm{Fl}(\mathcal{M})$, où $\mathcal{M}$ est une catégorie accessible) est \textit{accessible,} si elle l'est par rapport à la catégorie ambiante
\msp{page 75}
$\uFl\,\mathrm{Cat}$. Ceci posé, l'axiome Loc (9) est le suivant.

\textbf{Loc (9)}\quad $W$ est une partie \textit{accessible} de $\mathrm{Fl}(\mathrm{Cat}) = \mathrm{Ob}\,\uFl(\mathrm{Cat})$.

5.3. Rappelons que cette notion un peu technique se simplifie considérablement, si on suppose que $W$ satisfait l'axiome de limites Loc (8). Alors $W$ est accessible si et seulement si il existe une \textit{petite} partie $W_0$ de Fl(Cat), telle que $W$ soit la ``L$_1$-enveloppe'' de $W_0$, \textit{i.e.} la plus petite partie de Fl(Cat) contenant $W_0$ et L$_1$-stable (\textit{i.e.} stable par $\varinjlim$ filtrantes). S'il en est ainsi, il est clair qu'on peut prendre $W_0$ de la forme
\[
W_\pi = W \cap \mathrm{Fl}_\pi(\mathrm{Cat}),
\]
où $\pi$ est un cardinal infini convenable, et où $\mathrm{Fl}_\pi(\mathrm{Cat}) \overset{\text{déf}}{=} \mathrm{Fl}(\uCat_\pi)$, $\uCat_\pi$ désignant la sous-catégorie strictement pleine de Cat formée des objets $X$ tels que $\mathrm{card}\,X \leq \pi$ (où par définition, $\mathrm{card}\,X = \mathrm{card}\,\mathrm{Fl}\,X$).

\msp{page 76}
Notons que la filtration qu'on vient de décrire sur Cat (indexée par les petits cardinaux infinis) n'est autre que sa filtration canonique, \textit{i.e.} $\uCat_\pi$ est la sous-catégorie pleine de Cat formée des objets $\pi$-accessibles de Cat, et de même dans $\uFl(\mathrm{Cat})$. On trouve ainsi, grâce aux développements de XVIII, §13, la caractérisation suivante, que je rappelle ici pour mémoire (en me bornant pour simplifier au cas où $W$ satisfait à Loc (8)) :

\textbf{Proposition 5.3.1.}\footnote[38]{\textit{cf. loc. cit.} 13.36.} \textit{Soit $W$ une partie de} Fl(Cat) \textit{stable par $\varinjlim$ filtrantes, et soient pour tout cardinal infini $\pi$, $W_\pi = W \cap \mathrm{Fl}_\pi(\mathrm{Cat})$, et $\underline{W}$, $\underline{W}_\pi$ les sous-catégories pleines correspondantes de} $\uFl(\mathrm{Cat})$. \textit{On dit que $W$ est} $\pi$-accessible \textit{si $W$ est la L-enveloppe de $W_\pi$ (donc $W$ est accessible si et seulement s'il existe un cardinal infini $\pi$ tel que $W$ soit $\pi$-accessible). Ceci posé, les conditions suivantes sont équivalentes :}

\printedpage{42}
% head: A. Grothendieck : Dérivateurs, Ch. XVI, Édition des Universités de Montpellier II et Paris VII

\begin{enumerate}[label={}]
\item[a)] \textit{$W$ est $\pi$-accessible.}
\item[b)] \textit{Pour tout système inductif $(f_i)_{i\in I}$,}
\msp{page 77}
\textit{$L_1$-$\pi$-adapté dans} $\uFl\,\mathrm{Cat}$ (cf. \textit{ci-dessous}), \textit{de limite inductive $f$, on a $f \in W$ si et seulement si la sous-catégorie pleine $I_W$ des $i \in \mathit{Ob}I$ tels que $f_i \in W$ est cofinale dans $I$.}
\item[b$'$)] (\textit{Quand on a choisi, pour tout $f \in$} Fl(Cat)\textit{, une représentation $L_1$-$\pi$-adaptée de $f$ comme $\varinjlim$}) : \textit{Si $f \in W$, alors la sous-catégorie pleine $I_W$ définie dans} b) \textit{est cofinale dans $I$.}
\end{enumerate}

Rappelons ici la notion de système inductif L$_1$-$\pi$-adapté. Il s'agit des conditions :
\begin{enumerate}[label={}]
\item[1°)] $I$ est grand devant $\pi$, et stable par $\varinjlim$ filtrantes de cardinal $\leq \pi$.
\item[2°)] Le foncteur $i \longmapsto x_i : I \longrightarrow \uFl(\mathrm{Cat})$ commute aux $\varinjlim$ en question.
\item[3°)] Les $x_i$ sont dans $\uFl(\mathrm{Cat})_\pi \overset{\text{déf}}{=} \uFl_\pi(\mathrm{Cat})$.
\end{enumerate}
Notons que dans le cas d'un ``paratopos parabélien'' tel que Cat ou $\uFl\,\mathrm{Cat}$, il y a une façon canonique, $\pi$ étant donné, de représenter un objet $f$ comme $\varinjlim$ d'un système inductif L$_1$-$\pi$-adapté. Ici,
\[
f : X \longrightarrow Y
\]
étant donné, on considère l'ensemble $I$ des
\msp{page 78}
couples $(X_i \subset X, Y_i \subset Y)$ de sous-catégories de $X$ et de $Y$, telles que l'on ait $X_i, Y_i \in \mathrm{Ob}\,\mathrm{Cat}_\pi$ (\textit{i.e.} $\mathrm{card}\,(X_i), \mathrm{card}\,(Y_i) \leq \pi$) et que $f(X_i) \subset Y_i$. Il est immédiat que $I$ est filtrant, pour la relation d'ordre évidente
\[
(X_i, Y_i) \subset (X_j, Y_j) \qquad\quad \Longleftrightarrow \qquad\quad X_i \subset Y_i \ \text{ et } \ X_j \subset Y_j, % [sic] X_i \subset Y_i et X_j \subset Y_j -> X_i \subset X_j et Y_i \subset Y_j, la relation d'ordre compare les couples (X_i,Y_i) et (X_j,Y_j)
\]
et même grand devant $\pi$, que $I$ est stable par $\varinjlim$ filtrantes de cardinal $\leq \pi$, et que le foncteur
\[
i \longmapsto (f_i : X_i \longrightarrow Y_i) : I \longrightarrow \uFl\,\mathrm{Cat}
\]
commute auxdites $\varinjlim$ , enfin les $f_i$ sont dans $\uFl(\mathrm{Cat})_\pi$ par construction. Ainsi on a un système inductif L$_1$-$\pi$-adapté, et il est clair que sa limite inductive est canoniquement isomorphe à $f$. Ceci vu :

\textbf{Corollaire 5.3.2.} \textit{La condition que $W$ soit $\pi$-accessible signifie donc que pour tout $f : X \longrightarrow Y \in W$, et pour toute sous-catégorie $X_i$ de $X$ et $Y_i$ de $Y$, de cardinal $\leq \pi$ (avec $f(X_i) \subset Y_i$), mais en fait c'est inutile, quitte à agrandir $Y_i$), il existe des sous-catégories} % [sic] (avec f(X_i) \subset Y_i), mais ... agrandir Y_i) -> (avec f(X_i) \subset Y_i, mais ... agrandir Y_i), parenthèse fermante surnuméraire
\msp{page 79}
\textit{$X_j \supset X_i$ et $Y_j \supset Y_i$ de cardinal $\leq \pi$, telles que $f$ induise $f_j : X_j \longrightarrow Y_j$ et que l'on ait $f_j \in W$.}

La notion d'accessibilité de $W$ me semble aussi bien comprise que celle de limites, je n'y vois plus aucun mystère. D'autre part, il me semble que tous les localiseurs fondamentaux utilisés à ce jour y satisfont (tout comme à l'axiome Loc (8)). Ici, je vais me borner à établir :

\printedpage{43}
% head: A. Grothendieck : Dérivateurs, Ch. XVI, Édition des Universités de Montpellier II et Paris VII

\textbf{Théorème 5.4.}\footnote[39]{Mérite bien le nom de \textit{théorème,} surtout si on y inclut toute la liste des axiomes Loc mentionnés ci-contre\,! Ce théorème n'a rien d'évident, à partir de la définition cohomologique de $\mathbf{W}_\infty$ - c'est un résultat (``bien connu'') qui me paraît profond.} \textit{Le localiseur fondamental $\mathbf{W}_\infty$ (équivalences faibles ``habituelles'') satisfait à l'axiome} Loc (9) (\textit{en fait, il satisfait à} tous \textit{les axiomes vu jusqu'à présent, savoir} Loc (1 ter, 2, 3 bis, 4, 5, 6, 7, 8, 9)). \textit{Plus précisément, $\mathbf{W}_\infty$ est $\aleph_0$-accessible (où $\aleph_0$, ou $\pi_0$, est le premier cardinal infini).} % [sic] les axiomes vu -> les axiomes vus, accord du participe

Je ne reviens pas ici sur la vérification des autres axiomes, et me borne à prouver Loc (9), l'accessibilité, en utilisant ce qui peut m'être utile. Je me
\msp{page 80}
borne à indiquer la démonstration (assez technique) dans ses grandes lignes. Pour le moment, $W$ est à nouveau un localiseur fondamental quelconque, sauf qu'il satisfait à Loc (8).

\textbf{Lemme 5.5.}\footnote[40]{On suppose que $W$ satisfasse à Loc (8).}
\begin{enumerate}[label={}]
\item[a)] \textit{Pour que le critère} 5.3.1 b) \textit{de $\pi$-accessibilité de $W$ soit satisfait, il faut et il suffit qu'il le soit lorsque les $f_i$ sont $W_\omega$-parfaits.}
\item[b)] \textit{Supposons que $\mathit{Parf}_{W_\omega} \cap W \subset W^u$,} i.e. \textit{que $W$ satisfait à} Loc (7 $\mathrm{D}_{W_\omega}$) (\textit{la plus faible des variantes de l'axiome des $W$-fibrations, à l'exception seulement de} Loc (7 F)). \textit{Supposons de plus (pour être tranquille) que $W$ satisfait} \texttt{[phrase non terminée]}. \texttt{[?Alors]} \textit{$W$ est $\pi$-accessible si et seulement si l'ensemble} $\mathrm{Asph}_W \subset \mathrm{Ob}\,\mathrm{Cat}$ \textit{des objets $W$-asphériques de} Cat \textit{est $\pi$-accessible. Ou encore, que la partie} $\mathrm{Asph}_W \subset \mathrm{Ob}\,\mathrm{Cat}_{\textit{0-connexe}}$ \textit{soit $\pi$-accessible, où} $\mathrm{Cat}_{\textit{0-connexe}}$ \textit{désigne la sous-catégorie pleine de} Cat \textit{formée des objets} 0-\textit{connexes.}
\end{enumerate}
Pour l'essentiel, la démonstration est la même que celle de 5.1, en un peu plus technique. L'outil essentiel est toujours la factorisation fonctorielle $f = pi$ d'une flèche de Cat, avec $i \in W_\omega$, $p \in \Parf_{W_\omega}$, et la propriété de
\msp{page 81}
commutation aux $\varinjlim$ filtrantes de ladite factorisation.

Revenons maintenant au cas de $W = \mathbf{W}_\infty$. On doit prouver que $\mathrm{Asph}_{\mathbf{W}_\infty} \subset \mathrm{Ob}\,\mathrm{Cat}_{\text{0-connexe}}$ est $\pi_0$-accessible. Donc à prouver le critère correspondant à 5.3.1 b), en se plaçant maintenant dans $\mathrm{Cat}_{\text{0-connexe}}$ au lieu de $\mathrm{Fl}\,\mathrm{Cat}$. On a donc
\[
X = \varinjlim_{i\in I} X_i,
\]
représentation L$_1$-$\pi$-admissible de $X$ $\mathbf{W}_\infty$-asphérique, à prouver que pour tout $i$ dans $I$, $\exists j$ dans $I$ et $i \longrightarrow j$, avec $X_j$ $\mathbf{W}_\infty$-asphérique.

J'admets ici le lemme AQT \texttt{[âne qui trotte]} :

\textbf{Lemme 5.6.} \textit{Soit $\mathbf{W}_0 = \{f \in \mathrm{Fl}\,\mathrm{Cat} \mid \pi_0(f) \ \ \textit{isomorphisme}\}$. Alors $\mathbf{W}_0$ est $\pi_0$-accessible (et $L_1$-stable). A fortiori, l'ensemble des objets} 0-\textit{connexes de} Cat \textit{est $\pi_0$-accessible.}

Ceci nous ramène, quitte à remplacer $I$ par une catégorie cofinale, à supposer que les $X_i$ sont 0-connexes. (Mais j'ai oublié que cette réduction est déjà faite\,! Donc il faut voir 5.6 comme ingrédient


\endgroup
% END UNIT der-ch16/der-ch16_p041-045.tex
% BEGIN UNIT der-ch16/der-ch16_p046-050.tex
\begingroup\makeatletter%
\newcommand{\printedpage}[1]{\marginpar{\footnotesize p.~#1}}%
\newcommand{\msp}[1]{\marginpar{\footnotesize [#1]}}%
\makeatletter%
\renewcommand\@makefnmark{\hbox{\textsuperscript{(\@thefnmark)}}}%
\renewcommand\@makefntext[1]{\parindent 1em\noindent\hbox to 1.8em{\hss\textsuperscript{\@thefnmark}\,}#1}%
\makeatother%
\makeatother\let\printedpage\derpp\let\msp\dermsp%
\let\mspin\DERmspin%


\printedpage{44}
% head: A. Grothendieck : Dérivateurs, Ch. XVI, Édition des Universités de Montpellier II et Paris VII
% foot: Version 18 mars 2019

\msp{page 82}
de la démonstration de 5.5 b).)

Soit $x$ un objet (point) de $X$, $x_i$ son image dans $X_i$. Considérons (via ma construction standard des catégories de chemins) les $\Omega^1(X,x)$, $\Omega^1(X_i,x_i)$. On trouve
\[ \Omega^1(X,x) \simeq \varinjlim_{I} \Omega^1(X_i,x_i), \]
et on vérifie que c'est là encore un système inductif $\mathrm{L}_1$-$\pi_0$-adapté. En itérant, on trouve pour $n \in \mathbf{N}^*$
\[ \Omega^n(X,x) \simeq \varinjlim_{I} \Omega^n(X_i,x_i), \]
limite inductive $\mathrm{L}_1$-$\pi_0$-adaptée. Utilisant 5.6, on en conclut que pour tout $n$, le sous-ensemble (supposons $I$ un ensemble ordonné, que diable ça suffit !)
\[ I_n = \{ i \in I \mid \Omega^n(X_i,x_i) \ \text{ est 0-connexe} \} \]
est cofinal dans $I$, de plus il est stable par $\varinjlim$ filtrantes de cardinal $\leq \pi_0$. Donc % checked against PDF leaf 46: correct as printed; Theorem 5.4 explicitly uses pi_0 for the first infinite cardinal.
\[ I_\infty \overset{\text{déf}}{=} \bigcap I_n \]
est aussi cofinal dans $I$ (\textit{cf.} XVIII, 13.10.2, p. 302).

\msp{page 83}
Or $i \in I_\infty$ signifie que tous les $\Omega^n(X_i,x_i)$ sont 0-connexes, ce qui signifie aussi que $X_i$ est $\mathbf{W}_\infty$-asphérique (puisqu'il est déjà 0-connexe) (\textit{cf.} XII, cor. p. 190). Cela achève la démonstration du théorème 5.4.

\section*{6\quad La condition Loc (5, 7 D$_{W_\omega}$) implique $W \subset \mathbf{W}_\infty$}

Je vais redonner ici, pour mémoire, la démonstration du théorème 1 de XII :

\textbf{Théorème 6.1.} \textit{Soit $W$ un localiseur fondamental satisfaisant à} D$_{W_\omega}$ \textit{(où $W_\omega$ est le plus petit de tous les localiseurs fondamentaux). Alors on a $W = \mathrm{Fl}\,\mathrm{Cat}$,} ou \textit{$W \subset \mathbf{W}_\infty$ (suivant que $W$ ne satisfait pas à l'axiome de connexité} Loc (5), \textit{ou qu'il y satisfait.})

Je vais faire la démonstration en supposant que $W$ satisfait à Loc (5). Pour la démonstration du fait que dans le cas contraire, on a $W = \mathrm{Fl}(\mathrm{Cat})$, je renvoie à XII, th. 2, p. 20. (Voir aussi le résultat très voisin, plus général et plus faible, rappelé dans prop. 2.7, page 12.)

\textbf{Lemme 6.2.} \textit{Soient $W$, $W'$ deux localiseurs fondamentaux, $W$ satisfaisant à} Loc (7 D$_{W_\omega}$). \textit{Pour qu'on ait $W \subset W'$, il faut et il suffit qu'on ait} $\mathrm{Asph}_W \subset \mathrm{Asph}_{W'}$.

\msp{page 84}
Il suffit bien sûr de prouver “il suffit”, donc supposons que tout objet $W$-asphérique est $W'$-asphérique, et prouvons que $W \subset W'$. Soit donc $f : X \longrightarrow Y$ dans $W$, prouvons $f \in W'$. Pour ceci, j'utilise la factorisation canonique $f = pi$,
\[ X \xrightarrow{\ i \in W_\omega\ } \bar{X} \xrightarrow{\ p \in \mathrm{Parf}_{W_\omega}\ } Y. \]
\printedpage{45}
% head: A. Grothendieck : Dérivateurs, Ch. XVI, Édition des Universités de Montpellier II et Paris VII
% foot: Version 18 mars 2019
On a donc aussi $i \in W$, et comme $f = pi \in W$, on a $p \in W$, donc $p \in \mathrm{Parf}_{W_\omega} \cap W$. Or D$_{W_\omega}$ signifie que cet ensemble est contenu dans $W^{\mathrm{u}}$, donc on a $p \in W^{\mathrm{u}}$, et en particulier les fibres de $p$ sont $W$-asphériques. Elles sont donc $W'$-asphériques, et comme $p$ est $W_\omega$-parfait et a fortiori $W'$-parfait, il s'ensuit que $p \in W'$. On a aussi $i \in W_\omega \subset W'$, donc $i \in W'$, donc aussi $f = pi \in W'$, q.e.d.

Revenons alors à la démonstration de 6.1. En vertu du lemme, il suffit de prouver que si $X$ dans Cat est $W$-asphérique, il est $\mathbf{W}_\infty$-asphérique. Mais par hypothèse Loc (5) on sait déjà que $X$ est 0-connexe, et choisissant un point $x$ dans $X$, qu'il soit $\mathbf{W}_\infty$-asphérique s'exprime par la condition
\msp{page 85}
que les $\Omega^n(X,x)$ (pour $n \geq 1$) sont 0-connexes (\textit{cf.} XII, cor. p. 190 déjà cité). Pour ceci, à cause de Loc (5) il suffit de prouver qu'ils sont $W$-asphériques, et par induction il suffit de prouver que $\Omega^1(X,x)$ l'est. Or par définition de $\Omega^1(X,x)$, c'est là la fibre en $x$ d'une flèche dans Cat
\[ E = \underline{\mathrm{Ch}}(X,x,-) \xrightarrow{\ p \in \mathrm{Parf}_{W_\omega}\ } X, \]
où la source $E$ est $W_\omega$-asphérique, a fortiori $W$-asphérique. Comme $X$ l'est également, on a donc $p \in W$, donc $p \in \mathrm{Parf}_{W_\omega} \cap W \subset W^{\mathrm{u}}$, donc $p$ est à fibres $W$-asphériques, et en particulier $\Omega^1(X,x)$ est $W$-asphérique, q.e.d.

\textbf{Corollaire 6.3.} (Caractérisation de $\mathbf{W}_\infty$.) \textit{$\mathbf{W}_\infty$ est le plus grand localiseur fondamental satisfaisant à} Loc (5, 7 D$_{W_\omega}$) \textit{- donc aussi le plus grand localiseur fondamental satisfaisant à} Loc (5, 7).

Je pense que nous trouverons plus bas une caractérisation “interne” de $\mathbf{W}_\infty$, via un “axiome de descente” (très fort) Loc (10),
\msp{page 86}
comme l'unique localiseur fondamental satisfaisant Loc (5, 7 D$_{W_\omega}$, 10) (ou encore Loc (5, 7, 10)).

6.4. \textit{Commentaire sur les axiomes des $W$-fibrations du types} Loc (7 \ldots), \textit{et notamment sur} Loc (7 D$_{W_\omega}$). \textit{Relation conjecturale avec les} “types d'homotopie simples” \textit{de} \textsc{Whitehead}. % [sic] du types -> du type (or des types), grammatical agreement

a) Il me semble que \textit{tous} les localiseurs fondamentaux avec lesquels on a travaillé tacitement jusqu'à présent (les $\mathbf{W}_n$ ($n \in \mathbf{N}^*$), équivalence homologique et cohomologique, équivalence rationnelle \ldots) correspondent à des localiseurs $W \supset \mathbf{W}_\infty$. Par le théorème 6.1, il s'ensuit donc qu'\textit{aucun} d'eux, à part $\mathbf{W}_\infty$ lui-même, ne satisfait à la plus faible version D$_{W_\omega}$ de toutes de l'axiome Loc (7 \ldots), à la seule exception de Loc (7 F). Ainsi, les localiseurs $W$ qui y satisfont sont à l'exception de D$_{W_\omega}$ lui-même, d'un type tout nouveau - des localiseurs qui sont strictement contenus dans $\mathbf{W}_\infty$.

b) Ceci soulève, bien sûr, la question de l'existence de localiseurs fondamentaux contenus dans $\mathbf{W}_\infty$ et
\msp{page 87}
distincts de $\mathbf{W}_\infty$. Je n'ai pas même prouvé que $W_\omega \neq \mathbf{W}_\infty$, alors que je n'ai pas le moindre doute que tel est bien le cas. \texttt{[Cisinski a démontré que} $W_\omega = \mathbf{W}_\infty$\texttt{.]} A fortiori, j'ignore si le plus petit localiseur fondamental satisfaisant Loc (7 D$_{W_\omega}$) est différent de $\mathbf{W}_\infty$. Je serais tout particulièrement intéressé par le plus petit localiseur $\tilde{\mathbf{W}}_\infty \subset \mathbf{W}_\infty$ satisfaisant tous les axiomes de stabilité Loc (1 ter, 2, 3 bis, 4, 6, 7, 8) (\textbf{N.B.} il satisfait
\printedpage{46}
% head: A. Grothendieck : Dérivateurs, Ch. XVI, Édition des Universités de Montpellier II et Paris VII
% foot: Version 18 mars 2019
nécessairement Loc (5) car $\tilde{\mathbf{W}}_\infty \subset \mathbf{W}_\infty \subset \mathbf{W}_0$, de plus Loc (7) équivaut, en présence de Loc (8), à la forme très faible Loc (7 D$_{W_\omega}$) : $\mathrm{Parf}_{W_\omega} \cap W \subset W^{\mathrm{u}}$. Enfin, on montre que Loc (4) est conséquence des autres axiomes, \textit{cf.} plus bas.) Je conjecture à présent que $\tilde{\mathbf{W}}_\infty$ (qui donne lieu à un formalisme homotopique tout aussi complet que la sempiternelle “équivalence faible”) est bel et bien différent de $\mathbf{W}_\infty$ \texttt{[non par Cisinski]}. Et je suspecte
\msp{page 88}
que la catégorie $\mathrm{Hot}_{\tilde{\mathbf{W}}_\infty}$ mérite d'être appelée catégorie des \textit{“types d'homotopie simples”} (à la \textsc{Whitehead})\footnote[41]{Canulé (?), \textit{cf.} commentaire page 92.}. Un premier test décisif, c'est de vérifier que l'invariant cohomologique principal de \textsc{Whitehead} d'un type d'homotopie simple, savoir un élément dans un $K(\mathbf{Z}[\pi_1])$ (si je me rappelle bien\footnote[42]{non, c'est un invariant plus fin \ldots}) - ou quelque chose d'approchant, peut se définir déjà en termes des objets dans $\mathrm{Hot}_{\bullet\,\tilde{\mathbf{W}}_\infty}$ (le Hot pointé). Et deuxio, si on associe à une “maquette simpliciale” finie (décrivant combinatoirement une triangulation finie) l'ensemble ordonné de ses simplexes, vu comme un objet de Cat, et l'objet correspondant dans $\mathrm{Hot}_{\tilde{\mathbf{W}}_\infty}$, celui-ci ne change pas\footnote[43]{C'est évident même pour $W_\omega$, puisque le type d'homotopie \textit{stricte ne change pas !} - Non, il y a ici confusion entre homotopie topologique et homotopie combinatoire, \textit{cf.} p. 92 \ldots}, à isomorphisme canonique près, quand on effectue sur ce “polyèdre” envisagé une opération élé\msp{page 89}mentaire à la \textsc{Whitehead}. Cela suggère d'ailleurs d'essayer de définir de telles opérations élémentaires dans Cat tout entier, ou tout au moins dans la sous-catégorie des (petits) ensembles ordonnés, dans l'espoir de parvenir (peut-être) à une construction plus directe du localiseur pertinent pour le type d'homotopie simple.

c) Je termine par un commentaire sur le sens de l'axiome Loc (7 D$_{W_\omega}$), par opposition avec l'axiome plus faible Loc (7 F) qui assure seulement l'existence d'une “structure à fibrations” $(W, \mathrm{Fib}_W)$ sur Cat. Je signale aussi en passant que Loc (7 F), tout comme Loc (5) et comme Loc (9), et à la différence de tous les autres axiomes passés en revue jusqu'à présent, n'est \textit{pas} un axiome de stabilité : il affirme, non que certaines flèches (sous des hypothèses convenables) sont dans $W$, mais l'existence d'une factorisation. Aussi il n'est
\msp{page 90}
nullement clair que l'intersection de deux localiseurs fondamentaux satisfaisant à Loc (7 F) y satisfasse.

Loc (7 D$_{W_\omega}$) est exactement ce qu'il faut pour assurer que la factorisation canonique d'une flèche $f : X \longrightarrow Y$ de Cat en $f = pi$, avec $i \in W_\omega$, $p \in \mathrm{Parf}_{W_\omega}$, est aussi une factorisation du type stipulé par l'axiome de factorisation Loc (7 F) : on a (bien sûr) $i \in W$, mais aussi $p \in \mathrm{Fib}_W$ (car D$_{W_\omega}$ équivaut à sa forme forte $\mathrm{Parf}_{W_\omega} \subset \mathrm{Fib}_W$). Cela signifie aussi que la factorisation canonique de $f$ peut être utilisée pour “calculer” ses fibres homotopiques, non seulement en les points de $Y$, mais aussi par rapport à un morphisme de changement de base $Y' \longrightarrow Y$.

En l'absence de cet axiome,
\msp{page 91}
\printedpage{47}
% head: A. Grothendieck : Dérivateurs, Ch. XVI, Édition des Universités de Montpellier II et Paris VII
% foot: Version 18 mars 2019
admettant seulement Loc (7 F), on sait qu'il existe une théorie des fibres $W$-homotopiques, mais on n'a a priori aucun moyen pratique pour “calculer” lesdites fibres, faute de disposer d'un critère maniable suffisant pour assurer $p \in \mathrm{Fib}_W$. J'ai donc l'impression que l'axiome Loc (7 F) va être d'un maigre usage. (Il est vrai que l'exemple de $\mathbf{W}_0$ devrait me rendre prudent ! Il me faudrait regarder le cas des $\mathbf{W}_n$, s'ils satisfont Loc (7 F) \ldots). À mon sens, l'axiome Loc (7 D$_{W_\omega}$) est l'exact “minimum vital” pour être à l'aise dans le maniement des fibres $W$-homotopiques. Il est remarquable que ceci implique qu'on doive avoir $W \subset \mathbf{W}_\infty$, et que cela exclut du même coup \textit{toutes} les variantes considérées à ce jour de la notion de type d'homotopie !

\msp{page 92}
6.5. Finalement, la relation entrevue avec les “types d'homotopie simples” de \textsc{Whitehead} s'évanouit\footnote[44]{Mais peut-être quand même pas - la question reste ouverte.} : en effet, ces derniers sont strictement plus fin que les types d'homotopie (\textit{stricts}) ordinaires, et l'invariant de Whitehead n'est \textit{pas} un invariant du type d'homotopie ordinaire. Il est vrai que lorsqu'on considère le squelette simplicial d'un polyèdre triangulé comme objet de Cat, l'homotopie (continue) ordinaire entre deux tels polyèdres s'interprète, dans Cat, non pas (me semble-t-il) par l'existence d'un homotopisme entre les objets correspondants de Cat, mais par leur isomorphie dans $\mathrm{Hot}_{\mathbf{W}_\infty} \overset{\text{déf}}{=} \mathrm{Cat}(\mathbf{W}_\infty)^{-1}$. Mais le procédé d'approximations simpliciales de \textsc{Brouwer} montre que si les polyèdres $X$ et $Y$ sont homotopes, alors il existe des subdivisions simpliciales $X'$ et $Y'$ de $X$ et de $Y$ qui sont “combinatoirement homotopes” (dans un sens convenable). Donc je présume que les maquettes simpliciales de $X'$ et de $Y'$ sont alors homotopes dans Cat. Il s'ensuivrait que % [sic] plus fin -> plus fins, agreement with "ces derniers"
\msp{page 93}
$X$ et $Y$, j'entends leurs maquettes, sont bien homotopes dans Cat, s'il est vrai qu'une subdivision simpliciale (il suffirait la subdivision barycentrique, en itérant \ldots) $X'$ de $X$ a une maquette dans Cat qui lui est homotope. La question soulevée ici me paraît intéressante, il faut visiblement en trouver la réponse. Elle se pose déjà pour tout ensemble ordonnée $I$, qu'on va supposer fini pour simplifier, en prenant l'ensemble (ordonné par inclusion) $I' = \mathrm{sd}(I)$ des parties totalement ordonnées de $I$. On sait que les réalisations topologiques de $I$ et de $\mathrm{sd}(I)$ sont canoniquement isomorphes, mais $I$ et $\mathrm{sd}(I)$ sont-ils homotopes dans Cat ? On a bien une application croissante canonique % [sic] ensemble ordonnée -> ensemble ordonné, agreement
\[ f : \mathrm{sd}(I) \longrightarrow I, \]
qui associe à toute partie totalement ordonnée de $I$ son plus grand élément. Mais existe-t-il une application croissante en sens inverse, qui soit un homotopisme inverse du précédent à homotopie près ? Je viens de vérifier dans le cas

\msp{page 94} manque

\msp{page 95}
du bord du 2-simplexe

\begin{center}
\begin{tikzpicture}[baseline=(b.base)]
\node (b) at (0.95,0.3) {};
\draw (0,0) -- (1.9,0) -- (0.95,0.95) -- cycle;
\fill (0,0) circle (2.2pt);
\fill (1.9,0) circle (2.2pt);
\fill (0.95,0.95) circle (2.2pt);
\end{tikzpicture}
\qquad\qquad\qquad ou \qquad\qquad\qquad
\begin{tikzpicture}[baseline=(c.base)]
\node (c) at (0,-0.1) {};
\draw (0,0) circle (0.48);
\fill (90:0.48) circle (2.2pt);
\fill (210:0.48) circle (2.2pt);
\fill (330:0.48) circle (2.2pt);
\end{tikzpicture}
\qquad\quad ,
\end{center}

subdivision barycentrique
\printedpage{48}
% head: A. Grothendieck : Dérivateurs, Ch. XVI, Édition des Universités de Montpellier II et Paris VII
% foot: Version 18 mars 2019

\begin{center}
\begin{tikzpicture}[baseline=(b.base)]
\node (b) at (0.95,0.3) {};
\draw (0,0) -- (0.95,0.95) -- (1.9,0);
\fill (0,0) circle (2.2pt);
\fill (0.95,0) circle (2.2pt);
\fill (1.9,0) circle (2.2pt);
\fill (0.475,0.475) circle (2.2pt);
\fill (1.425,0.475) circle (2.2pt);
\fill (0.95,0.95) circle (2.2pt);
\end{tikzpicture}
\qquad\qquad\qquad ou \qquad\qquad\qquad
\begin{tikzpicture}[baseline=(c.base)]
\node (c) at (0,-0.1) {};
\draw (0,0) circle (0.48);
\fill (30:0.48) circle (2.2pt);
\fill (90:0.48) circle (2.2pt);
\fill (150:0.48) circle (2.2pt);
\fill (210:0.48) circle (2.2pt);
\fill (270:0.48) circle (2.2pt);
\fill (330:0.48) circle (2.2pt);
\end{tikzpicture}
\end{center}

(hexagone), que pour toute application croissante $g : I \longrightarrow \mathrm{sd}(I)$, sa réalisation topologique est homotope à une application constante, donc $g$ ne peut être un homotopisme dans Cat (car cela impliquerait que sa réalisation topologique est un homotopisme).

Ceci montre donc que la notion d'homotopisme dans Cat est \textit{strictement plus fine} que celle pour les réalisations topologiques, ce qui redonne à nouveau substance à l'espoir que l'invariant de Whitehead pourrait se définir comme invariant d'homotopie sur Cat, ou du moins sur une sous-catégorie pleine convenable, telle celle des ensembles ordonnés (éventuellement restreints à être finis), ou des maquettes simpliciales. Et même, que l'ensemble des Wh-types d'homotopie (types d'homotopie simples) peut se
\msp{page 96}
plonger canoniquement dans un $\mathrm{Hot}_W$ convenable, où $W$ serait un localiseur fondamental strictement plus fin que $\mathbf{W}_\infty$, p.ex. (qui sait ?) $\tilde{\mathbf{W}}_\infty$. Le premier travail qu'on aurait envie de faire, c'est de construire un tel $W$ \textit{ad hoc}, pour que l'invariant de Whitehead puisse s'écrire comme un foncteur sur $\mathrm{Hot}_{\bullet W}$. Pour commencer, il faudrait déjà s'assurer que cet invariant peut être vu comme un invariant du type d'homotopie, \textit{i.e.} qu'il s'interprète comme un foncteur sur $\mathrm{Hot}_{\bullet}\mathrm{Cat} \simeq \mathrm{Cat}(\,\underbrace{\text{homotop}}_{\substack{\text{ensemble des}\\ \text{homotopismes}}}\,)^{-1}$. Ce travail suggérera peut-être aussi la définition des “invariants de Whitehead supérieurs” (impliquant les $\pi_i$, non seulement le $\pi_1$), lesquels pourraient peut-être servir (avec celui de dimension 1) à décrire le localiseur hypothétique $W$ (“de Whitehead”).

\msp{page 97}

\section*{7\quad Commentaires sur les axiomes}

7.1. La principale raison d'être d'un localiseur fondamental $W$, c'est de définir une théorie d'homotopie’, et plus précisément un prédérivateur $\mathrm{HOT}_W$, dont on espère que ce sera même un dérivateur, ou qu'il satisfasse du moins à certains des axiomes des dérivateurs. Les axiomes Loc (1) à Loc (9) sont en majeure partie orientés vers cet objectif, à cela près que Loc (1, 2, 3) constitue un “minimum vital” absolu, pour pouvoir commencer à travailler avec $W$. Les principales questions qui se posent me paraissent les suivantes : % [sic] d'homotopie’ -> d'homotopie (or “théorie d'homotopie”), closing quote without opening quote
\begin{enumerate}[label={}, leftmargin=2.5em]
\item[1)] Quand Cat munie de $(W, \mathrm{Cof}_W)$, resp. de $(W, \mathrm{Fib}_W)$, est-elle une catégorie à cofibrations resp. à fibrations, de façon que les résultats de K. S. \textsc{Brown} soient applicables ?
\item[2)] Sous quelles conditions, sur $W$ et sur $f$, $f_!$ et $f_*$ existent-ils ? Et quand a-t-on, de plus, exactitude des carrés cocartésiens resp. cartésiens ?

\msp{page 98}
\item[3)] Plus généralement, en termes des axiomes imposés à $W$, quelles sont les axiomes Der (\ldots) que sont vérifiés pour le prédérivateur $\mathrm{HOT}_W$ ? (Ce qui pourrait aussi dépendre du choix du domaine Diag $\subset$ Cat. Le prédérivateur est défini a priori sur Cat tout entier). En particulier, % [sic] quelles sont les axiomes -> quels sont les axiomes, agreement; [sic] que sont vérifiés -> qui sont vérifiés, relative pronoun
\end{enumerate}


\endgroup
% END UNIT der-ch16/der-ch16_p046-050.tex
% BEGIN UNIT der-ch16/der-ch16_p051-055.tex
\begingroup\makeatletter%
\newcommand{\printedpage}[1]{\marginpar{\footnotesize p.~#1}}%
\newcommand{\msp}[1]{\marginpar{\footnotesize [#1]}}%
\makeatletter%
\renewcommand{\@makefnmark}{\hbox{\normalfont\@textsuperscript{(\@thefnmark)}}}%
\renewcommand{\@makefntext}[1]{\noindent\hbox{\normalfont\@textsuperscript{\@thefnmark}}\,#1}%
\makeatother%
\makeatother\let\printedpage\derpp\let\msp\dermsp%
\let\mspin\DERmspin%


\printedpage{49}
% head: A. Grothendieck : Dérivateurs, Ch. XVI, Édition des Universités de Montpellier II et Paris VII
% foot: Version 18 mars 2019

\begin{enumerate}[label=4)]
\item quand $\mathrm{HOT}_W$ est-il bel et bien un dérivateur sur Cat tout entier, \textit{i.e.} quand satisfait-il à tous les axiomes ?
\end{enumerate}

Une question qui me paraît moins importante (notamment que 1°) est la suivante :

\begin{enumerate}[label=1 bis)]
\item Quand $W$ s'insère-t-il dans un triple de Quillen clos $(W,C,F)$ (propre de préférence).
\end{enumerate}

Je vais passer en revue ces questions une à une.

7.2. On a vu (mais sans que j'aie fait une vérification soigneuse pour la réciproque) que $(W,\mathrm{Cof}_W)$ définit une structure à cofibrations si l'axiome \textit{Loc (6 a)} est satisfait, et que cette condition est aussi nécessaire si Loc (3 bis) est satisfait\footnote[45]{Vérifier quand même si Loc (3 bis) n'est pas déjà utilisé dans la partie directe ! \textbf{O.K.}}.

\msp{page 99}
Alors les immersions ouvertes de Dwyer (ou au choix les immersions fermées de Dwyer) sont dans $\mathrm{Cof}_W$, et la construction standard du mapping-cone définit, pour toute flèche $f$ dans Cat, une factorisation \textit{fonctorielle}
\[
\begin{aligned}
f = pi \quad & i \text{ immersion ouverte de Dwyer (donc dans } \mathrm{Cof}_W\text{),}\\
& p \in W.
\end{aligned}
\]
Cela, plus la construction du foncteur $f_!$, et en particulier des $W$-sommes amalgamées
\begin{tikzcd}[row sep=tiny, column sep=small, baseline=(\tikzcdmatrixname-2-1.base)]
 & X_1 \\
\displaystyle\int X_0 \arrow[ur] \arrow[dr] & \\
 & X_2
\end{tikzcd}, données dans VI (théorie de $\mathrm{HOT}_W$) et XII, nous donne une maîtrise parfaite (il me semble) sur la “structure homotopique gauche” définie par $W$, sous le seul axiome Loc (6 a) (Axiome du cube, forme directe). Je rappelle que les foncteurs $f_!$ sont construits dans VI sous la seule hypothèse \textit{Loc (3 bis)} ($f$ flèche quelconque de Cat), mais sauf erreur, la construction
\msp{page 100}
et l'interprétation $W$-homotopique de
\begin{tikzcd}[row sep=tiny, column sep=small, baseline=(\tikzcdmatrixname-2-1.base)]
 & X_1 \\
\displaystyle\int X_0 \arrow[ur] \arrow[dr] & \\
 & X_2
\end{tikzcd}
donnée dans XII ne dépend pas même de cette hypothèse\footnote[46]{\textit{cf.} XII, page 264} - mais je ne peux le garantir. Il me semble que XII (et notamment le long §5 sur les sommes amalgamées) est écrit sous l'hypothèse que W (1, 2, 3), \textit{i.e.} Loc (1, 2, 3bis), sont satisfaits.

7.3. C'est une chose assez extraordinaire que les foncteurs $f_!$ pour $\mathrm{HOT}_W$ existent, sous la seule hypothèse \textit{Loc (3 bis)} (axiome de localisation, forme forte) - et ils se “calculent” d'une façon féeriquement simple ! Je suis à peu près sûr que sous cette seule hypothèse, \textit{l'axiome d'exactitude pour les carrés cocartésiens} n'est \textit{pas} satisfait\footnote[47]{Il faudrait un exemple !}. Je ne vois aucun moyen de le vérifier, sauf via la théorie de \textsc{Brown}, sous l'hypothèse \textit{Loc (6 a)} (\textit{cf.} 7.2).

Mais sous la seule condition Loc (3 bis) le calcul standard des fibres d'un $f_!(\xi)$
\printedpage{50}
% head: A. Grothendieck : Dérivateurs, Ch. XVI, Édition des Universités de Montpellier II et Paris VII
% foot: Version 18 mars 2019
\msp{page 101}
est valable, en d'autre termes le prédérivateur $\mathrm{HOT}_W$ satisfait Der (4$^*$, 5$^*$) - et de plus Der (6$^*$) (exactitude “à gauche”) si $W$ satisfait Loc (6 a). % [sic] d'autre termes -> d'autres termes, missing plural s

7.4. Venons-en à la question si Cat muni de $(W,\mathrm{Fib}_W)$ est une catégorie à fibrations. C'est l'axiome \textit{Loc (7 F),} mais je ne connais guère de moyen de le vérifier (en dehors du cas très, très particulier $W = \mathbf{W}_0$) que moyennant tout au moins l'axiome \textit{Loc (7 $D_{W_\omega}$),} ou l'une ou l'autre des variantes plus fortes de Loc (7). De Loc (7 $\mathrm{D}_{W_\omega}$) il a été abondamment question précédemment, \textit{cf.} notamment 6.4 (pages 89-91). En mon sens, cet axiome assure une maîtrise comparable sur la structure $W$-homotopique droite (fibres $W$-homotopiques, suites exactes de cosuspension) que Loc (6 a) pour la
\msp{page 102}
structure gauche. En particulier, on a la factorisation \textit{fonctorielle} de toute flèche $f$ de Cat en
\[
f = pi, \qquad i \in W_\omega \subset W,\ p \in \mathrm{Parf}_{W_\omega} \subset \mathrm{Fib}_W .
\]
Mais \textit{l'axiome Loc (7 $D_{W_\omega}$) me paraît de nature considérablement plus forte que Loc (6 a),} lequel (si je ne m'abuse) est vérifié pour tous les localiseurs qui ont été envisagés jusqu'à présent. (Mais j'avoue n'avoir pas regardé le cas des $\mathbf{W}_n$, $n \in \mathbf{N}$, sauf $\mathbf{W}_0$.) Par contre, l'axiome Loc (7 $\mathrm{D}_{W_\omega}$) implique $W \subset \mathbf{W}_\infty$ (théorème 6.1), relation vraiment draconienne (\textit{cf.} commentaire 6.4 a).

Par ailleurs, le seul axiome Loc (7 $\mathrm{D}_{W_\omega}$), pour autant que je sache, n'implique pas l'existence de $f_*$ pour \textit{toute} flèche dans Cat, mais seulement pour des
\msp{page 103}
flèches de la forme $f \times \mathrm{id}_X : S \times X \longrightarrow T \times X$, où $f : S \longrightarrow T$ est un morphisme entre ensembles ordonnés finis. Ou plus généralement, entre ensembles ordonnés $S$, $T$ tels que $S^\circ$ soit artinien, \textit{i.e.} $\mathrm{S}$ noethérien, et que $\mathrm{Esp}(S^\circ)$ soit noethérien (cette dernière hypothèse étant inutile sous l'axiome Loc (8) des limites, \textit{cf.} XVII, th. 3.1, p. 50). Donc ce n'est pas le résultat rêvé.

7.5. Le seul cas que je connaisse où $f_*$ existe pour \textit{toute} flèche dans Cat, est le cas où $W$ satisfait à
\[
\textit{Loc (6 a, 8, 9)}\ \ \text{\footnotemark[48]\,\footnotemark[49]}, \tag*{$(*)$} % [sic] marks printed (48) (49) -> the two footnotes printed below are numbered 49 and 50, numbering mismatch between marks and footnotes
\]
\footnotetext[49]{Vérifier s'il ne faut pas aussi la condition Loc (3 bis) en plus de Loc (6 a, 8, 9). Vérification faite, c'est O.K. (sans Loc (3 bis)).}
\footnotetext[50]{\textbf{N.B.} Après nouvelle réflexion, il apparaît qu'il suffit même Loc (6$'$ a, 8, 9).}
c'est à dire l'axiome des $W$-cofibrations faibles, déjà abondamment mentionné et enfin ceux des limites et d'accessibilité, \textit{i.e.} $W$ est $\mathrm{L}_1$-accessible dans \underline{Fl}(Cat). Le résultat-clef pertinent est donné dans XV 8.4 (p. 105-107), qu'il faut joindre à XV 7.4 (p. 98-99) pour en faire
\msp{page 104}
bon usage, et en déduire le théorème 10.1 (XV, p. 134-135). Mieux que l'existence des $f_*$, l'ensemble d'axiomes étonnamment faible que je viens d'expliciter suffit à assurer que $\mathrm{HOT}_W$ est un \textit{dérivateur} sur Cat tout entier, en plein sens du terme, \textit{i.e.} qu'il satisfait à Der (1, 2, 3, 4, 5, 6, 4$^*$, 5$^*$, 6$^*$). Ceci est obtenu en remplaçant (grâce à XV 8.4) Cat par
\printedpage{51}
% head: A. Grothendieck : Dérivateurs, Ch. XVI, Édition des Universités de Montpellier II et Paris VII
% foot: Version 18 mars 2019
$A^\wedge$, où $A$ est une catégorie-test convenable, $W$ par $W_A \subset \mathrm{Fl}\, A^\wedge$, et en prouvant que $W_A$ s'insère dans un triple de Quillen clos $(W,C,F)$, avec $C = \mathrm{Mono}(A^\wedge)$, ce qui détermine $F$ comme $(C \cap W_A)_*$. On applique alors XV 7.4 à cette situation (résultats résumés dans XIX, th. 2.3 (p. 19), 2.6 (p. 25), 2.9 (p. 28). On retrouve l'existence des $f_!$ (avec un mode de calcul bien différent de celui dans Cat !), et en plus on trouve l'existence de $f_*$ (plus l'exactitude dans cor. [?] 5
\msp{page 105}
des carrés $W$-cartésiens) - ce qui pour moi, pendant un bon moment, était \textit{le} gros problème pour la théorie de $\mathrm{HOT}_W$ !).

Je trouve époustouflant qu'on parvienne à un résultat aussi fort avec des axiomes aussi ridiculement faibles. Plus particulièrement, qu'on ait toutes les propriétés désirables de la structure homotopique droite, sans avoir à faire sur $W$ aucune hypothèse du type $W$-fibration ! Ces axiomes n'impliquent \textit{pas} l'axiome Loc (7) (équivalent à Loc (7 $\mathrm{D}_{W_\omega}$) en présence de Loc (8)), comme on a vu sur l'exemple de $\mathbf{W}_0$. Impliquent-ils tout au moins Loc (7 F), \textit{i.e.} l'existence d'assez de $W$-fibrations ? Je ne hasarderais aucune conjecture à ce sujet ! Mais il faut bien se dire qu'en l'absence de Loc (7) (et même
\msp{page 106}
en admettant Loc (7 F)), tout en sachant que la structure homotopique droite associée à $W$ est parfaite à tous égards, on n'a a priori aucune maîtrise sur celle-ci, notamment pour le calcul des fibres $W$-homotopiques et la construction des carrés $W$-cartésiens. Même en ramenant toutes ces constructions à $A^\wedge$ (prenant, si on y tient, $A = \boldsymbol{\Delta}$, \textit{i.e.} se ramenant au contexte semisimplicial), quoi qu'on sache qu'il y a une factorisation \textit{fonctorielle} $f = pi$ avec $i \in \mathrm{TC} \subset W_A$, $p \in F \subset \mathrm{Fib}_{W_A}$, je ne vois aucune façon concrète de l'expliciter. À moins que la considération des chemins infinis, qui (en l'absence de Loc (7)) s'avère inopérante dans Cat sauve la mise dans $\boldsymbol{\Delta}^\wedge$ ? ?

\msp{page 107}
Question qui devrait être facile à élucider, mais sur laquelle je ne veux pas m'attarder à présent.

Notons en passant que le résultat de saturation forte de \textsc{Quillen}, pour les triples de Quillen clos, montre que l'ensemble d'axiomes $(*)$ implique l'axiome Loc (1 ter), \textit{i.e.} que $W$ est \textit{fortement saturé.}

Notons aussi que contrairement à ce qui a lieu pour les $f_!$, je ne sais prouver l'existence de $f_*$, et notamment des suites de cosuspension, que sous des conditions (soit Loc (7 $\mathrm{D}_{W_\omega}$), soit Loc (6a, 8, 9)) qui assurent déjà que les carrés $W$-cartésiens sont exacts, et par suite les suites de cosuspension exactes.

Je pense que dans les pages précédentes, j'ai dit tout ce que je sais dire de pertinent, concernant les questions 1, 2, 4 soulevées dans 7.1. Ces réponses me parais\msp{page 108}sent satisfaisantes - plus satisfaisantes que je ne l'aurais osé espérer pendant les premiers deux mois de mes cogitations ! Il reste les questions 3 et 1 bis, que je vais à présent examiner.

\printedpage{52}
% head: A. Grothendieck : Dérivateurs, Ch. XVI, Édition des Universités de Montpellier II et Paris VII
% foot: Version 18 mars 2019
7.6. Sur les différents axiomes des dérivateurs : sous quelles conditions sur $W$ un de ces axiomes est-il satisfait pour $\mathrm{HOT}_W$\footnote[51]{Pour les axiomes Der (1) à Der (5$^*$), \textit{cf.} I, p. 67-68. Les axiomes Der (6, 6$^*$) sont les axiomes d'exactitude, dont la nécessité m'est apparue seulement ultérieurement.} ? Pour ce qui est des axiomes Der (4, 5, 6, 4$^*$, 5$^*$, 6$^*$), je crois avoir dit tout ce que je sais en dire, et avoir des critères ce validité parfaitement maniables, dans les pages précédentes. Résumons pourtant : % [sic] ce validité -> de validité, misprint
\[
\left\{
\begin{array}{lcl}
\text{structure à cofibrations} & : & \text{Loc (6) (+ Loc (3 bis) ?)} \\[1ex]
\text{structure à fibrations} & : & \text{Loc (7 }\mathrm{D}_{W_\omega}\text{)} \\[1ex]
\text{Der (4}^*\text{, 5}^*\text{)} & : & \text{Loc (3 bis)} \\[1ex]
\text{Der (4}^*\text{, 5}^*\text{, 6}^*\text{)} & : & \text{Loc (3 bis, 6a) ou Loc (6}'\text{ a, 8, 9)} \\[1ex]
\text{Der (4, 5, 6) (sur Cat) (et tout le reste !)} & : & \text{Loc (6}'\text{ a, 8, 9)} \\[1ex]
\text{Der (4, 5, 6) (sur Diag $=$ ens. ordonnés finis)} & : & \text{Loc (7 }\mathrm{D}_{W_\omega}\text{) (ou seulement Loc (7 F)).}
\end{array}
\right.
\]

Il reste les axiomes Der (1) (décomposition), Der (2) (localisation), Der (3) (objet unité), que je vais passer à présent en revue. On vérifie
\msp{page 109}
de façon quasi-tautologique que l'axiome Der 1 est satisfait sans condition sur $W$ (pas même Loc (2, 3) !) - en fait, il est satisfait pour tout prédérivateur $\mathbf{D}_{(\mathcal{M},W)}$ défini par une catégorie de modèles.

L'axiome de localisation Der (2) est par contre de nature délicate. Il est valable, plus généralement, pour une catégorie de modèles $(\mathcal{M},W)$ (définissant un prédérivateur $\mathbf{D} = \mathbf{D}_{(\mathcal{M},W)}$) et un objet $S$ de Cat, lorsque $(\mathcal{M}(S),W(S))$ est associé à une structure à fibrations ou à cofibrations, ce qui par K. S. \textsc{Brown} implique que toute flèche de $\mathrm{HO}_W(S) = \mathcal{M}(S)W(S)^{-1}$ se “réalise” par une flèche de $\mathcal{M}(S)$ - et alors le résultat est évident. La condition que je viens de dire est satisfaite lorsque $(\mathcal{M},W)$ lui-même est associé à une structure de cofibrations ou de fibrations,
\msp{page 110}
et si de plus on suppose que l'on est dans l'un de deux cas suivants : % [sic] l'un de deux cas -> l'un des deux cas, apparent misprint
\begin{enumerate}[label=\alph*)]
\item $S$ est un ensemble ordonné fini (ce qui nous amène à nous restreindre, à défaut de mieux, à Diag $=$ sous-catégorie de Cat formée des ensembles ordonnés finis).
\item La factorisation stipulée dans l'axiome de factorisation $f = pi$ (avec $i \in \mathrm{Cof}$, $p \in W$ (cas de cofibrations) ou $i \in W$, $p \in \mathrm{Fib}$ (cas de fibrations)) peut être choisie fonctorielle en $f$.
\end{enumerate}
Or, dans le cas de Cat, dans les cas où on sait prouver que $W$ est associé à une structure à cofibrations (axiome \textit{Loc (6 a)}) où à fibrations (axiome \textit{Loc (7 $D_{W_\omega}$)}), on trouve bel et bien une factorisation fonctorielle. Donc chacun de ces axiomes implique Der (2). % [sic] où à fibrations -> ou à fibrations, accent misprint

\msp{page 111}
Reste l'axiome de contractibilité : si $S$ dans Cat est contractile, est-il $\mathrm{HOT}_W$-asphérique, \textit{i.e.} le foncteur
\[
\mathrm{HOT}_W(e) = \mathrm{Hot}_W \longrightarrow \mathrm{HOT}_W(S)
\]
\printedpage{53}
% head: A. Grothendieck : Dérivateurs, Ch. XVI, Édition des Universités de Montpellier II et Paris VII
% foot: Version 18 mars 2019
est-il pleinement fidèle ? Quand on dispose de la théorie de $f_!$ “version Cat” (donc sous le seul axiome \textit{Loc (3 bis)}), on trouve ceci (noté dans VI, p. 202)\footnote[52]{moyennant Loc (3 bis).} :
\begin{quote}
Soit $f : S \longrightarrow T$ une flèche dans Cat. Pour que ce soit une $\mathrm{HOT}_W$-équivalence, \textit{i.e.} pour que $f^* : \mathrm{HOT}_W(T) \longrightarrow \mathrm{HOT}_W(S)$ induise un foncteur \textit{pleinement fidèle} sur la sous-catégorie des objets constants de $\mathrm{HOT}_W(T)$, il faut et il suffit que $\mathrm{hot}_W(f) : \mathrm{hot}_W(S) \longrightarrow \mathrm{hot}_W(T)$ soit un isomorphisme. Donc il \textit{suffit} que $f \in W$, et cette condition est aussi nécessaire si $W$ est fortement saturé (axiome Loc (1 ter)).
\end{quote}
Par acquit de conscience, j'explicite ici la démonstration. L'existence des foncteurs
\msp{page 112}
$f_!$, en particulier des objets d'homologie $H_\bullet^{\mathrm{HOT}_W}(S,\xi)$ (où simplement $H_\bullet^W(S,\xi)$) pour $\xi \in \mathrm{HOT}_W(e) = \mathrm{Hot}_W$, permet d'interpréter homologiquement la condition de $\mathrm{HOT}_W$-équivalence, comme signifiant que pour tout $\xi = \mathrm{hot}_W(X)$ dans $\mathrm{Hot}_W$, la flèche canonique dans $\mathrm{Hot}_W$ % [sic] (où simplement -> (ou simplement, accent misprint
\[
H_\bullet^W(S,\xi) \longrightarrow H_\bullet^W(T,\xi) \tag*{$(*)$}
\]
est un isomorphisme. Or le calcul des $f_!$ (dans le cas de $p : S \longrightarrow e$ et de $T \longrightarrow e$) nous montre que les deux membres de $(*)$ ne sont autres (à isomorphisme près) que $\mathrm{hot}_W(X \times S)$ et $\mathrm{hot}_W(X \times T)$, la flèche étant induite par
\[
\mathrm{id}_X \times f : X \times S \longrightarrow X \times T. \tag*{$(*)$}
\]
Donc la condition envisagée signifie que cette flèche est un isomorphisme dans $\mathrm{Hot}_W$ pour tout $X$ dans Cat. Prenant $X = e$, on voit que cela implique que $\mathrm{hot}_W(f)$ est un isomorphisme, et cela suffit, vu que le
\msp{page 113}
foncteur $\mathrm{hot}_W : \mathrm{Cat} \longrightarrow \mathrm{Hot}_W$ commute aux produits finis (il suffit de Loc (1, 2, 3)) \texttt{[plutôt Loc (1, 2, 3 bis)]} et que la flèche $(*)$ dans $\mathrm{Hot}_W$ s'identifie à
\[
\mathrm{id}_{\mathrm{hot}_W(X)} \times \mathrm{hot}_W(f)\,.
\]
Cela montre donc que si $S$ est un objet $W$-asphérique, \textit{i.e.} tel que $S \longrightarrow e$ soit dans $W$ (et a fortiori, si $S$ est contractile), alors $S \longrightarrow e$ est $\mathrm{HOT}_W$-asphérique.

Pour résumer, on trouve le tableau
\[
\left\{
\begin{array}{lcl}
\text{Der (1)} & : & \text{sans hypothèse sur } W \\[1ex]
\text{Der (2)} & : & \text{Loc (6 a) \textit{ou} Loc (7 }\mathrm{D}_{W_\omega}\text{)} \\[1ex]
\text{Der (3)} & : & \text{Loc (3 bis) \textit{ou} Loc (6}'\text{ a, 8, 9)\ \ \footnotemark[53].}
\end{array}
\right.
\]
\footnotetext[53]{\textbf{N.B.} On verra plus bas que Loc (6 a, 8, 9) implique déjà Loc (3 bis).}

7.7. Il faut seulement justifier encore que Loc (6$'$ a, 8, 9) implique Der (3). Mais alors, on l'a vu, $\mathrm{HOT}_W$ peut être défini par un triple de Quillen très spécial. Mais même pour un triple de Quillen spécial sans plus $(\mathcal{M},W)$, le prédérivateur $\mathbf{D}$ qu'il définit satisfait (en plus de Der (4$^*$, 5$^*$, 6$^*$)) aux conditions Der (1) (trivialement), Der (2)


\endgroup
% END UNIT der-ch16/der-ch16_p051-055.tex
% BEGIN UNIT der-ch16/der-ch16_p056-059.tex
\begingroup\makeatletter%
\newcommand{\printedpage}[1]{\marginpar{\footnotesize p.~#1}}%
\newcommand{\msp}[1]{\marginpar{\footnotesize [#1]}}%
\newcommand{\circled}[1]{\tikz[baseline=(c.base)]\node[draw,circle,inner sep=1pt](c){#1};}%
\makeatother\let\printedpage\derpp\let\msp\dermsp%
\let\mspin\DERmspin%


\printedpage{54}
% head: A. Grothendieck : Dérivateurs, Ch. XVI, Édition des Universités de Montpellier II et Paris VII
\msp{page 114}
(\textit{cf.} p. 110, cas b)) et \textit{de plus} Der (3). (Donc dans le cas d'un triple de Quillen \textit{très} spécial, le prédérivateur obtenu est bel et bien un dérivateur.) Je vais expliciter ici la démonstration.

Je vais désigner par $(\mathcal{M},\Sigma)$ la catégorie de modèles donnée ($\Sigma\subset\operatorname{Fl}\mathcal{M}$), par $\mathbf{D}=\mathbf{D}_{\mathcal{M},W}$ le prédérivateur correspondant, et par
\[ W=W_D \] % [sic] W_D -> W_{\mathbf{D}}, index printed as italic D, elsewhere bold D
le localiseur dans Cat associé, formé des ``$\mathbf{D}$-équivalences''. On voit tout de suite que $W$ est fortement saturé (vrai pour tout $W_{\mathbf{D}}$, $\mathbf{D}$ un prédérivateur), \textit{i.e.} satisfait Loc (1 ter), a fortiori Loc (1), d'autre part il satisfait à Loc (3 bis$^*$), comme conséquence du fait que $\mathbf{D}$ satisfait à Loc (4$^*$, 5$^*$). A fortiori il satisfait à Loc (3$^*$), et pour voir que tout $X$ contractile dans Cat est $W$-asphérique, \textit{i.e.} $X\longrightarrow e$ dans $W$, on sait (\textit{cf.} §1)
\msp{page 115}
qu'il suffit de prouver qu'il en est ainsi si $X$ a un objet initial (axiome Loc (2$^*$) dual de l'axiome Loc (2)). À ce moment, on saura que $W$ est un localiseur fondamental, et de plus il satisfait Loc (3$^*$ bis), donc aussi Loc (3 bis) (forme forte de l'axiome de localisation).

Donc \textit{on peut supposer} que si $S$ a un objet initial, alors $S$ est $W$-asphérique, \textit{i.e.} $S\longrightarrow e$ est dans $W$, \textit{i.e.} $p^*:\mathbf{D}(e)\longrightarrow\mathbf{D}(S)$ est pleinement fidèle. Il revient au même de dire que la flèche d'adjonction
\[ p_!p^*\xi\longrightarrow\xi \]
est un isomorphisme, pour tout $\xi$ dans $\mathbf{D}(e)$. Or le calcul de $p_!(\eta)$ (pour $\eta\in\operatorname{Ob}\mathbf{D}(S)$, \textit{i.e.} $\eta\in\operatorname{Ob}\mathcal{M}(S)$) se fait ainsi : on choisit une flèche
\[ \eta'\xrightarrow{\Sigma_S}\eta \]
avec $\eta'$ ``cofibrant'' (pour la structure à cofibrations pertinente sur $\mathcal{M}(S)$), et on
\msp{page 116}
pose \footnote[54]{opération $p_!$ au niveau de $\mathcal{M}(S)$.}
\[ p_!(\operatorname{ho}_{W,S}(\eta))=\operatorname{ho}_W(p_!(\eta')). \]
Quand $\eta=p^*(\xi)$, la flèche d'adjonction pertinente dans $\operatorname{Ho}_\Sigma=\mathbf{D}(e)$ s'obtient par composition
\[ p_!(\eta')\longrightarrow p_!(p^*(\xi))\xrightarrow{\text{adj.}}\xi, \tag*{$(*)$} \]
où ici $p_!$ est encore pris dans le sens ordinaire ($\mathcal{M}(S)\longrightarrow\mathcal{M}(e)=\mathcal{M}$). Or $S$ ayant un objet initial 0, on a
\[ \begin{aligned}
p_!(\eta') &\ \simeq\ \eta'(0)\\
p_!(p^*(\xi)) &\ \simeq\ p^*(\xi)(0)\ \simeq\ \xi,
\end{aligned} \]
et la flèche d'adjonction dans $(*)$ est l'identité, la première flèche est la flèche $\eta'(0)\longrightarrow\eta(0)$ induite par $\eta'\longrightarrow\eta$, elle est dans $\Sigma$, donc $(*)$ est dans $\Sigma$, donc c'est un isomorphisme dans $\operatorname{Ho}_\Sigma=\mathbf{D}(e)$, q.e.d.

Le même argument marche d'ailleurs si $\mathcal{M}$ est seulement une catégorie à cofibrations ou à fibrations, mais en se bornant alors au cas où $S$ est un ensemble ordonné fini.
% foot: Version 18 mars 2019

% Original hard page break suppressed for the reflowed reader.
\printedpage{55}
% head: A. Grothendieck : Dérivateurs, Ch. XVI, Édition des Universités de Montpellier II et Paris VII
\msp{page 117}
Je vais dégager à part ce qu'on vient d'obtenir :

\textbf{Proposition 7.7.1.} \textit{Soit} $(\mathcal{M},W)$ \textit{une catégorie de modèles, associée à un triple de Quillen clos} spécial $(W,C,F)$\textit{. Alors le prédérivateur} $\mathbf{D}=\mathbf{D}_{\mathcal{M},W}$ \textit{sur} Cat \textit{satisfait aux axiomes} Der (1, 2, 3, 4$^*$, 5$^*$, 6$^*$)\textit{. Si de plus le triple de Quillen est} très \textit{spécial} (\textit{donc en particulier} $C=\operatorname{Mono}(\mathcal{M})$), \textit{alors il satisfait de plus à} Der (4, 5, 6), i.e. \textit{c'est un dérivateur au plein sens du terme. D'autre part, le localiseur} $W=W_{\mathbf{D}}$ \textit{dans} Cat \textit{associé à} $\mathbf{D}$ (\textit{dans le cas d'un triple de Quillen spécial sans plus}) \textit{est un} localiseur fondamental, \textit{satisfaisant de plus} Loc (1 ter, 3 bis).

\textbf{N.B.} Ce qui est dit ici pour le cas d'un triple de Quillen très spécial (qui ne fait pas appel à la théorie des localiseurs
\msp{page 118}
fondamentaux dans Cat) montre à nouveau que si $W$ localiseur fondamental dans Cat satisfait à Loc (6$'$ a, 8, 9), alors $\mathrm{HOT}_W$ est un vrai dérivateur. On a vu de plus dans 7.6 (page 111) que l'on a alors $W=W_{\mathrm{HOT}_W}$ \footnote[55]{Mais c'était sous l'hypothèse Loc (3 bis) ! Gare au cercle vicieux.}, donc il s'ensuit :

\textbf{Corollaire 7.7.2 :} \textit{On a l'implication}
\[ \text{Loc } (6'\text{ a, 8, 9})\ \Longrightarrow\ \text{Loc(1 ter, 3 bis)}\ \text{\footnotemark[56]}. \]
\footnotetext[56]{Loc (3 bis) pas prouvé, probablement faux.}

Notons aussi, pour mémoire :

\textbf{Corollaire 7.7.3.} \textit{Soit} $(\mathcal{M},W)$ \textit{une catégorie de modèles, associée à une catégorie à cofibrations. Alors sur le domaine} Diag \textit{formé des ensembles ordonnés finis, le prédérivateur} $\mathbf{D}=\mathbf{D}_{\mathcal{M},W}$ \textit{satisfait aux axiomes} Der (1, 2, 3, 4$^*$, 5$^*$, 6$^*$), \textit{et dualement. Donc si} $W$ \textit{est associé à la fois à une structure de cofibrations et de fibrations,} $\mathbf{D}$ \textit{est un dérivateur sur} Diag.

\msp{page 119}
\texttt{[Les six lignes qui suivent semblent hors contexte.]}

\textit{alors la flèche de changement de base}
\[ p^*f_!\longrightarrow f'_!q^* \]
\textit{entre foncteurs} $\mathrm{HOT}_W(X)\longrightarrow\mathrm{HOT}_W(S')$ \textit{est un isomorphisme.}

C'est pratiquement trivial, compte tenu, dans le cas $f$ $W$-lisse, que si $Z$ sur $X$ est $W$-lisse sur $X$, il l'est sur $S$ ; et dans le cas $p$ $W$-propre, en utilisant le dernier énoncé (``oubli'') dans \circled{6}, p. 199.

7.8. Je termine en mettant en vedette ce qui, à présent, me paraît la substance principale de tous ces ``rappels'' et comme la quintessence de la théorie des localiseurs fondamentaux
% margin: three vertical rules printed in the left margin, beside the footnotes and down to the foot of the page
% foot: Version 18 mars 2019
\printedpage{56}
% head: A. Grothendieck : Dérivateurs, Ch. XVI, Édition des Universités de Montpellier II et Paris VII
développée jusqu'à présent :
\[ \left\{ \begin{array}{rcl}
\text{Loc (3 bis)} & \Longrightarrow & \text{Der (4$^*$, 5$^*$)}\\[1ex]
\text{Loc (3 bis, 6 a)} & \Longrightarrow & \text{Der (1, 2, 3, 4$^*$, 5$^*$, 6$^*$)}\\[1ex]
\text{Loc (6$'$ a, 8, 9)} & \Longrightarrow & \text{tous les axiomes des dérivateurs (+ Loc (1 ter)).}
\end{array} \right. \]

Je considère les axiomes Loc (3 bis), Loc (8), Loc (9) comme des axiomes ``anodins'' - tous les localiseurs qu'on définit par voie cohomologique ou homotopique ordinaire, y satisfont. Il reste l'axiome Loc (6 a) ou Loc (6$'$ a) des $W$-cofibrations (forme faible),
\msp{page 120}
qui visiblement joue un rôle crucial dans la théorie de $\mathrm{HOT}_W$. Par contre, les axiomes Loc (4) (décomposition), Loc (5) (connexité) et Loc (7) et ses variantes (et bien sûr aussi Loc (6 b)) n'interviennent absolument pas ! C'est particulièrement impressionnant, en ce qui concerne l'axiome des $W$-fibrations Loc (7) et ses variantes (toutes équivalentes à Loc (7) en présence de Loc (8), à la seule exception de Loc (7 F)). Comme on a vu, Loc (7) n'est nullement impliqué par Loc (6 a, 8, 9) ($\mathbf{W}_0$ n'y satisfait pas), tout au plus se pourrait-il que Loc (7 F) le soit. Il en est de même de l'axiome de connexité Loc (5), comme on voit sur l'exemple de $W=\mathrm{Fl}(\mathrm{Cat})$, qui satisfait à tous les axiomes vus jusqu'à présent, à la seule exception
\msp{page 121}
de Loc (5). Quant à la question de l'indépendance de l'axiome Loc (4) par rapport à Loc (6 a, 8, 9), nous allons examiner ce genre de questions au paragraphe suivant.

L'ensemble d'axiomes Loc (6$'$ a, 8, 9) me paraît à présent si important, que j'ai envie de lui donner un nom, et d'appeler \textit{localiseur} (fondamental) \textit{spécial} dans Cat \footnote[57]{J'hésite si je ne dois pas rajouter l'axiome Loc (3 bis), qui assure que $W=W_{\mathbf{D}}$, $\mathbf{D}=\mathrm{HOT}_W$.}, un localiseur fondamental qui y satisfait. Ce sont donc aussi ceux pour lesquels nous savons que $\mathrm{HOT}_W$ est un vrai \textit{dérivateur.} Si on ne peut écarter comme ``pathologique'' une structure aussi riche et délicate qu'un dérivateur, il s'ensuit aussitôt que l'on ne peut non plus écarter les localiseurs spéciaux dans Cat (même si on affecte de regarder de haut les
\msp{page 122}
localiseurs fondamentaux généraux !). En effet, $W$ se récupère de façon unique à l'aide de $\mathbf{D}=\mathrm{HOT}_W$, ou de tout dérivateur équivalent à celui-ci, comme l'ensemble des $\mathbf{D}$-équivalences dans Cat \footnote[58]{Mais ceci n'est prouvé que moyennant Loc (3 bis).} : c'est le localiseur fondamental (correspondant à une notion d'équivalence ``universelle'', \textit{i.e.} dans la catégorie Cat de tous les ``espaces'' élémentaires) canoniquement associé au dérivateur. Ce n'est donc nullement un localiseur de référence fixe assez fin, tel $\mathbf{W}_\infty$, qui est apte à jouer un rôle !

Cela nous amène à nouveau à une question du genre déjà soulevé : si $W\subset\mathbf{W}_\infty$ est le plus petit localiseur qui satisfait à Loc (6$'$ a, 8) (ou Loc (3 bis, 6$'$ a, 8)), peut-on dire a) que $W$ est même \textit{spécial}, \textit{i.e.}
\msp{page 123}
satisfait également Loc (9) (accessibilité) ?, et de plus b) a-t-on $W\neq\mathbf{W}_\infty$ ? Je suis persuadé que c'est le cas (il y a moins de mérite pour ce $W$-là d'être différent de $\mathbf{W}_\infty$,
% foot: Version 18 mars 2019
\printedpage{57}
% head: A. Grothendieck : Dérivateurs, Ch. XVI, Édition des Universités de Montpellier II et Paris VII
que pour le $\tilde{\mathbf{W}}_\infty$ envisagé dans 6.4 (p. 87)).

7.9. Oubli : la question 1 bis (p. 98), sur l'existence d'une structure de Quillen close sur Cat, correspondant à $W$. Le seul résultat que je connaisse dans ce sens est celui de \textsc{Thomason}, construisant explicitement une telle structure de Quillen (de plus propre), dans le cas où $W=\mathbf{W}_\infty$. Je n'ai pas vraiment regardé la démonstration de \textsc{Thomason}, et ignore si celle-ci peut s'étendre à des localiseurs fondamentaux plus généraux, par exemple à ceux satisfaisant Loc (6$'$ a, 8, 9). S'il en était ainsi, cela prouverait d'ailleurs que ces localiseurs
\msp{page 124}
ont assez de cofibrations et assez de fibrations, donc qu'ils satisfont Loc (6 a) et Loc (7 F). (Pour ce qui est de Loc (6 a), il semble qu'on puisse le prouver directement très simplement, \textit{cf.} le paragraphe suivant. Mais ici encore, il faudrait regarder quelques exemples (équivalences homologiques ou cohomologiques, $\mathbf{W}_n$, \textit{i.e.} les $n$-équivalences d'homotopie), et regarder déjà si Loc (7 F) est satisfait.)

Il ne me semble pas que la théorie de \textsc{Thomason} ait été jusqu'à présent \texttt{[mots illisibles]} être très utile, techniquement parlant. Psychologiquement, je crois qu'il a été important, pour donner confiance que Cat est vraiment une excellente catégorie de modèles pour Hot. J'ai été vraiment enchanté et stimulé d'apprendre ce théorème. Mais la notion d'application de Dwyer (que j'ai apprise en 1983 dans l'article de \textsc{Thomason}) a été, techniquement, beaucoup plus importante pour moi : c'était exactement ce qui manquait à ma panoplie de notions et d'intuitions ``homotopiques'' dans Cat \dots
% foot: Version 18 mars 2019


\endgroup
% END UNIT der-ch16/der-ch16_p056-059.tex
\clearpage\setcounter{footnote}{0}\phantomsection\addcontentsline{toc}{section}{Chapitre XVII}\markboth{Chapitre XVII}{Chapitre XVII}
% BEGIN UNIT der-ch17/der-ch17_p001-005.tex
\begingroup\makeatletter%
\newcommand{\printedpage}[1]{\marginpar{\footnotesize p.~#1}}%
\newcommand{\msp}[1]{\marginpar{\footnotesize [#1]}}%
\newcommand{\Fl}{\mathrm{Fl}}%
\makeatother\let\printedpage\derpp\let\msp\dermsp%
\let\mspin\DERmspin%


% ===================== PDF page 1 =====================
\printedpage{--}
% head: (none)
% no page number printed on this page

\begin{center}
\vspace*{5cm}

{\Large\textsf{LES DÉRIVATEURS}}

\bigskip

\textbf{ALEXANDRE GROTHENDIECK}

\vspace{2cm}

\textbf{Chapitre XVII}

\medskip

\textbf{Catégories à fibrations et à cofibrations}\\
\textbf{(Modèles (3))}
\end{center}

\par

% ===================== PDF page 2 =====================
\printedpage{--}
% head: (none)
% no page number printed on this page

\vspace*{5cm}

\par\noindent\textbf{Note des éditeurs de l’édition transcrite.}\par
Ce texte a été déchiffré et transcrit en \LaTeXe{} par M. Künzer. Il a été édité par M. Künzer, J. Malgoire, et G. Maltsiniotis. La transcription est aussi fidèle que possible au manuscrit. Seules quelques corrections mineures évidentes ont été effectuées. Pour les rares ajouts ou commentaires des éditeurs, ainsi que pour la numérotation originale des pages du manuscrit, les caractères de machine à écrire \texttt{[typewriter]} entre crochets sont utilisés.

Cette édition est provisoire. Les remarques, commentaires et corrections sont bienvenus. Envoyer un message à :

\textbf{[e-mail address removed]}

\hspace*{\fill}G. Maltsiniotis\hspace*{2cm}

\par

% ===================== PDF page 3 =====================
\printedpage{1}
% head: A. Grothendieck : Dérivateurs, Ch. XVII, Édition des Universités de Montpellier II et Paris VII
% foot: Version 9 février 2019

\msp{page 1}

\section*{1\quad Préliminaires\footnote{\textbf{N.B.} Ces préliminaires seront complètement repris et dépassés dans le paragraphe 2. À ne pas lire !}}

\textbf{Proposition 1.} \textit{Soit $\mathcal{M}$ une catégorie stable par $\varinjlim$ quelconques, $C \subset \Fl(\mathcal{M})$ contenant les isomorphismes, stable par cochangement de base et par composition} transfinie. \textit{Soit, pour tout petit ensemble ordonné $I$, $C_I \subset \Fl(\mathcal{M}(I^{\circ}))$ l'ensemble des flèches $u_\bullet : X_\bullet \longrightarrow Y_\bullet$ telles que $\forall\ i_0 \in I$, dans le diagramme} % ?? uncertain: superscript of I -- \circ or degree sign

\[
\begin{tikzcd}[column sep=large]
X'_{i_0} = \varinjlim\limits_{i<i_0} X_i \arrow[r] \arrow[d, "\alpha_{i_0} = \varinjlim\limits_{i<i_0} u_i"'] \arrow[dr, phantom, "\text{cocart.}"] & X_{i_0} \arrow[d, "\beta_{i_0}"'] \arrow[dd, dotted, shift left=3ex, "u_{i_0}"] & \\
Y'_{i_0} = \varinjlim\limits_{i<i_0} Y_i \arrow[r] & \tilde{Y}_{i_0} \arrow[d, "\tilde{u}_{i_0}"'] & \text{\footnotesize (défini comme somme amalgamée)} \\
& Y_{i_0}\ , &
\end{tikzcd}
\tag*{$(*)$}
\]

\textit{la flèche $\tilde{u}_{i_0}$ soit dans $C$\footnote{Condition stable par passage de $I$ à un sous-ensemble ordonné \textit{ouvert} $I'$ de $I$.}. Je dis que si $I$ est} bien ordonné, \textit{alors $u_\bullet \in C_I$ implique}
\begin{itemize}[leftmargin=2.5em]
\item[a)] \textit{que $\varinjlim u_i : \varinjlim X_i \longrightarrow \varinjlim Y_i$ est dans $C$, et}
\item[b)] $\forall\ i \in I$, $(u_i : X_i \longrightarrow Y_i) \in C$\footnote{Dire aussi que l'on a
\begin{itemize}[leftmargin=2em]
\item[b$'$)] \textit{Si $i' \leq i$, alors dans le diagramme similaire à $(*)$,}
\end{itemize}
\[
\begin{tikzcd}[ampersand replacement=\&]
X'_{i'} \arrow[r] \arrow[d] \arrow[dr, phantom, "\text{cocart.}"] \& X'_i \arrow[d] \\
Y'_{i'} \arrow[r] \& \tilde{Y}'_i \arrow[d, "u_{i',i}"] \\
\& Y'_i\ ,
\end{tikzcd}
\]
\textit{on a} $u_{i',i} \in C$ (on retrouve $(*)$ en faisant $i' = i_0$, $i = i_0 + 1$). Le cas b) s'obtient en faisant $i' = 0$, $i$ quelconque, d'où $(X'_i \longrightarrow Y'_i) \in C$. L'appliquant à $i+1$ au lieu de $i$, on trouve $(u_i : X_i \longrightarrow Y_i) \in C$.}.
\end{itemize}

. % a lone full stop printed on its own line

\textbf{Corollaire.} \textit{Soit $W \subset \mathcal{M}$ localiseur modérément saturé, tel que $TC \overset{\text{déf}}{=} C \cap W$ soit stable par cochangement de base, et par composition} transfinie. \textit{Soit $u_\bullet : X_\bullet \longrightarrow Y_\bullet$ dans $C_I$. Alors $u_\bullet \in W_{I^{\circ}}$ (i.e. $\forall\ i \in I$, on a $u_i \in W$) si et seulement si $u_\bullet \in (TC)_I$, i.e. si et seulement si pour tout $i \in I$, la flèche $\tilde{u}_i$ dans $(*)$ est dans $TC$, ou ce qui}
\msp{page 2 n'existe pas}
\printedpage{2}
% head: A. Grothendieck : Dérivateurs, Ch. XVII, Édition des Universités de Montpellier II et Paris VII
% foot: Version 9 février 2019
\msp{page 3}
\textit{revient au même, dans $W$ (puisque par hypothèse elle est déjà dans $C$). Si cette condition est satisfaite, alors $\varinjlim u_i : \varinjlim X_i \longrightarrow \varinjlim Y_i$ est dans $W$.}

% ===================== PDF page 4 (continued) =====================

\textsc{Démonstration} du corollaire. Supposons $u_\bullet \in (TC)_I$, alors par la proposition, $\varinjlim\limits_{i\in I} u_i$ est dans $TC$, ou ce qui revient au même, dans $W$. D'autre part, appliquant ce même résultat aux systèmes inductifs induits sur $I_{<i_0}$, on trouve que dans $(*)$, la flèche $\alpha_{i_0}$ est dans $TC$, donc aussi $\beta_{i_0}$ par cochangement de base, donc aussi $u_{i_0} = \tilde{u}_{i_0}\beta_{i_0}$. Donc il reste à prouver qu'inversement, si on suppose les $u_{i_0} \in W$, alors les $\tilde{u}_{i_0}$ sont eux aussi dans $W$ (i.e. dans $TC = C \cap W$), ou ce qui revient au même, que les $\beta_{i_0}$ sont dans $W$, ou encore dans $TC$. Mais comme $TC$ est stable par cochangement de base, il suffit pour cela de prouver que $\alpha_{i_0} \in TC$, ou encore $\alpha_{i_0} \in C$. Donc on est ramené à prouver que
\begin{itemize}[leftmargin=2.5em]
\item[c)] $u \in W_{I^{\circ}} \cap C_I \Longrightarrow (\,\varinjlim\limits_{i<i_0} u_i) \in W$ pour tout $i_0 \in I$,
\end{itemize}
donc, changeant de notation, à prouver \texttt{[phrase incomplète]}. Nous prouverons ensemble a) b) c) par induction transfinie sur le nombre ordinal $\alpha$ tel que $I \simeq I_\alpha = \{\text{ensemble des ordinaux } i < \alpha\}$.

\msp{page 4}

1$^\circ$) Cas $\alpha = 0$, i.e. $I = \emptyset$ : il est trivial.

2$^\circ$) Si c'est ok pour $I_\alpha$, aussi pour $I_{\alpha+1} = \{\beta \mid \beta \leq \alpha\}$. Ici, $I_{\alpha+1}$ ayant un plus grand élément $\alpha$, a) est contenu dans b) (pour $i = \alpha$, plus grand élément de $I_{\alpha+1}$.) Pour prouver b), il reste à prouver que $u_\alpha \in C$. On regarde le diagramme $(*)$ pour $i_0 = \alpha$. Par l'hypothèse de récurrence (forme a), appliquée à $I_\alpha$, on a $\alpha_{i_0} \in C$, donc par cochangement de base $\beta_{i_0} \in C$, donc par composition $u_{i_0} = \tilde{u}_{i_0}\beta_{i_0} \in C$. Pour prouver c), il reste à prouver que $\varinjlim\limits_{i<\alpha} u_i \in W$. Mais par l'hypothèse de récurrence appliquée à $I_\alpha = \{i \mid i < \alpha\}$, on sait déjà que $(u_i)_{i<\alpha} \in (TC)_{I_\alpha}$, donc par la partie prouvée du corollaire ou de la proposition, $\varinjlim\limits_{i<\alpha} u_i \in TC$.

3$^\circ$) Si $\alpha$ est un ordinal limite, et si a) b) c) sont vrais pour $I_\beta$ avec $\beta < \alpha$, ils le sont aussi pour $I_\alpha$. Cette fois, b) est contenu dans l'hypothèse de
\msp{page 5}
récurrence, il faut prouver a), i.e. que
\[
u_\alpha \overset{\text{déf}}{=} \varinjlim_{i<\alpha} u_i \ :\ \underbrace{X_\alpha}_{\overset{\text{déf}}{=} \varinjlim\limits_{i<\alpha} X_i} \longrightarrow \underbrace{Y_\alpha}_{\overset{\text{déf}}{=} \varinjlim\limits_{i<\alpha} Y_i}
\]
est dans $C$. Quant à c), il est déjà contenu dans la hypothèse de récurrence. Donc il reste finalement à prouver
\[
u_\alpha \ \in \ C
\]
quand on suppose que la proposition est prouvée pour les ordinaux $\beta < \alpha$.

On va construire un troisième système ordinal sur $I_\alpha$, en introduisant pour tout $0 \leq i < \alpha$, i.e. $i \in I_\alpha$, $Z_i$ par le carré cocartésien
\printedpage{3}
% head: A. Grothendieck : Dérivateurs, Ch. XVII, Édition des Universités de Montpellier II et Paris VII
% foot: Version 9 février 2019
\[
\begin{tikzcd}
X_i \arrow[r] \arrow[d] \arrow[dr, phantom, "\text{cocart.}"] & X_\alpha \arrow[d] \\
Y_i \arrow[r] & Z_i\ .
\end{tikzcd}
\]

Les flèches de transition entre les $Z_i$ sont évidentes, et on aura $\varinjlim\limits_{i} Z_i \xrightarrow{\ \sim\ } Y_\alpha$. Ainsi, $u_\alpha : X_\alpha (= Z_{-1}) \longrightarrow Y_\alpha$ est décomposé en composé transfini de morphismes
\msp{page 6}
\[
Z_i \xrightarrow{v_{i+1}} Z_{i+1} \qquad (-1 \leq i < \alpha)
\tag*{$(*)$}
\]
et
\[
\varinjlim_{i<i_0} Z_i \xrightarrow{v_{i_0}} Z_{i_0} \qquad i_0 \text{ ordinal limite } < \alpha.
\tag*{$(**)$}
\]

(\textbf{N.B.} Le cas où $i_0$ n'est \textit{pas} ordinal limite n'est autre que $(*)$ pour $i = \beta_0 - 1$.) Tout revient à prouver que ces flèches sont dans $C$ (donc leur composé transfini aussi). % [sic] $\beta_0 - 1$ -> $i_0 - 1$, the case discussed is that of $i_0$

Cas $(*)$.
\[
\begin{tikzcd}
X_i \arrow[r] \arrow[d, "u_i"'] \arrow[dr, phantom, "\text{cocart.}"] & X_{i+1} \arrow[r] \arrow[d, "\beta_{i+1}"'] \arrow[dr, phantom, "\text{cocart.}"] & X_\alpha \arrow[d] \\
Y_i \arrow[r] & \tilde{Y}_{i+1} \arrow[r] \arrow[d, "\tilde{u}_{i+1}"'] \arrow[dr, phantom, "\text{cocart.}"] & Z_i \arrow[d, "v_{i+1}"] \\
& Y_{i+1} \arrow[r] & Z_{i+1}\ .
\end{tikzcd}
\tag*{\footnotemark}
\]
\footnotetext{Pour se donner confiance -- on pourrait traiter toute de suite le cas général.} % [sic] toute de suite -> tout de suite, French idiom

Comme $\tilde{u}_{i+1} \in C$, $v_{i+1} \in C$ par cochangement de base, OK. Il faut cependant traiter à part le cas $i = 0$, i.e. la flèche $X_\alpha \longrightarrow Z_0$, qui est incluse dans
\[
\begin{tikzcd}
X_0 \arrow[r] \arrow[d, "u_0"'] \arrow[dr, phantom, "\text{cocart.}"] & X_\alpha = Z_{-1} \arrow[d, "v_0"] \\
Y_0 \arrow[r] & Z_0\ ,
\end{tikzcd}
\]
où $u_0 \in C$, donc $v_0 \in C$, OK.

Cas $(**)$. Si $i_0 < \alpha$ est ordinal limite, inspectons $(**)$. On a par passage à la limite
\[
Z'_{i_0} \overset{\text{déf}}{=} \varinjlim_{i<i_0} Z_i \ \simeq\ Y'_{i_0} \amalg_{X'_{i_0}} X_\alpha
\]
(cf. diagramme $(*)$ pour la définition de $X'_{i_0}$, $Y'_{i_0}$), s'insérant dans le diagramme suivant greffé sur $(*)$


\endgroup
% END UNIT der-ch17/der-ch17_p001-005.tex
% BEGIN UNIT der-ch17/der-ch17_p006-010.tex
\begingroup\makeatletter%
\newcommand{\printedpage}[1]{\marginpar{\footnotesize p.~#1}}%
\newcommand{\msp}[1]{\marginpar{\footnotesize [#1]}}%
\newcommand{\defeq}{\overset{\textup{déf}}{=}}%
\newcommand{\cocart}{\textup{cocart.}}%
\makeatother\let\printedpage\derpp\let\msp\dermsp%
\let\mspin\DERmspin%

\setcounter{footnote}{4}

\printedpage{4}
% head: A. Grothendieck : Dérivateurs, Ch. XVII, Édition des Universités de Montpellier II et Paris VII
% foot: Version 9 février 2019

\msp{page 7}

\[
\begin{tikzcd}
X'_{i_0} \arrow[r] \arrow[d, "\alpha_{i_0}"'] \arrow[dr, phantom, "\cocart"] & X_{i_0} \arrow[r] \arrow[d, "\beta_{i_0}"'] \arrow[dr, phantom, "\cocart"] & X_{\alpha} \arrow[d] \\
Y'_{i_0} \arrow[r] & \tilde{Y}_{i_0} \arrow[r] \arrow[d, "\tilde{u}_{i_0}"'] \arrow[dr, phantom, "\cocart"] & Z'_{i_0} \arrow[d, "v_{i_0}"] \\
& Y_{i_0} \arrow[r] & Z_{i_0}\ .
\end{tikzcd}
\]

On a $\tilde{u}_{i_0} \in C$ par hypothèse, donc $v_{i_0} \in C$, ce qui achève la démonstration.

\textbf{Proposition 2.} \textit{$(\mathcal{M}, C)$ et le cas échéant $W$, comme dans proposition 1 et son corollaire, $I$ ensemble ordonné localement fini, i.e. tel que $\forall\ i_0 \in I$, $I/i_0$ soit fini. Soit}
\[
u_{\bullet}\ : X_{\bullet} \longrightarrow Y_{\bullet}
\]
\textit{une flèche dans\footnote{en fait, l'hypothèse que $X_{\bullet}$, $Y_{\bullet}$ dans $\mathcal{M}(I^{\circ})_{\mathrm{c}}$ est inutile, tant dans la proposition 2 que dans ses deux corollaires.} $\mathcal{M}(I^{\circ})_{\mathrm{c}} =$ sous-catégorie pleine de systèmes inductifs sur $I$, à valeurs dans $\mathcal{M}$, tels que $\underbrace{\emptyset_{\mathcal{M}(I^{\circ})}}_{\substack{\textup{objet}\\ \textup{initial}\\ \textup{de } \mathcal{M}(I^{\circ})}} \longrightarrow Z_{\bullet}$ soit dans $C_I$ (définie dans prop. 1), i.e. tel que $\forall\ i_0 \in I$, la flèche}
\[
Z'_{i_0} \ = \ \varinjlim_{i < i_0} Z_i \ \longrightarrow \ Z_{i_0}
\]
\textit{soit dans $C$, et supposons $u_{\bullet} \in C_I$. Alors}
\begin{enumerate}
\item[a)] $\varinjlim\limits_{i \in I} u_{\bullet} : X_I \defeq \varinjlim\limits_{i \in I} X_i \longrightarrow Y_I \defeq \varinjlim\limits_{i \in I} Y_i$ \textit{est dans $C$,}
\msp{page 8}\\
\textit{et}
\item[b)] \textit{$\forall\ i_0 \in I$, on a $\alpha_{i_0}, \beta_{i_0} \in C$ dans le diagramme $(*)$ p. 1.}
\end{enumerate}

\textbf{Corollaire.} \textit{De plus (supposant donné $W$ comme dans prop. 1, cor. 1), les conditions suivantes sur une $u_{\bullet}$ entre $X_{\bullet}$, $Y_{\bullet}$ $C$-cofibrants sont équivalentes.}
\begin{enumerate}
\item[$1^{\circ}$)] \textit{$u_{\bullet} \in (TC)_I$, i.e. les $\tilde{u}_{i_0} \in TC$ ($\forall\ i_0 \in I$).}
\item[$2^{\circ}$)] \textit{$u_{\bullet} \in T_I \cap W_{I^{\circ}}$, i.e. $\forall\ i_0 \in I$, $\tilde{u}_{i_0} \in C$, $u_{i_0} \in W$.} % [sic] T_I -> C_I, the explanation that follows gives \tilde u_{i_0} \in C, i.e. the condition C_I
\item[$3^{\circ}$)] \textit{Pour tout $i_0 \in I$, on a}
\[
\alpha_{i_0},\ \beta_{i_0},\ \tilde{u}_{i_0},\ u_{i_0} \ \in \ TC
\]
\textit{(i.e. $\alpha_{i_0}, \tilde{u}_{i_0} \in TC$, ce qui implique $\beta_{i_0} \in TC$, puis $u_{i_0} \in TC$ par cochangement de base et par composition).}
\end{enumerate}
\textit{Et ces conditions impliquent}
\begin{enumerate}
\item[$4^{\circ}$)] \textit{que $u_I : X_I \longrightarrow Y_I$ est dans $TC$.}
\end{enumerate}

Voyons comment ou dans quelle mesure le corollaire se déduit de la proposition. On a bien sur % [sic] sur -> sûr, accent missing (printed "bien sûr" elsewhere)
\printedpage{5}
% head: A. Grothendieck : Dérivateurs, Ch. XVII, Édition des Universités de Montpellier II et Paris VII
% foot: Version 9 février 2019
\[
\begin{tikzcd}[column sep=small]
& 3^{\circ} \arrow[dl, Rightarrow] \arrow[dr, Rightarrow] & \\
1^{\circ} & & 2^{\circ}
\end{tikzcd}
\]
tautologiquement. D'ailleurs, la proposition appliquée à $TC$ au lieu de $X$ montre qu'on a $1^{\circ} \Longrightarrow 3^{\circ}$ \& $4^{\circ}$, car en l'appliquant aux ensembles $I_{<i_0}$ on trouve que $\alpha_{i_0} \in TC\ \forall\ i_0 \in I$. Donc la proposition 2 implique % [sic] X -> C, the proposition is applied to TC instead of C
\[
(1^{\circ} \Longleftrightarrow 3^{\circ}) \Longrightarrow 2^{\circ}\ (\&\ 4^{\circ}).
\]

\msp{page 9}

Donc pour prouver $2^{\circ} \Longrightarrow 3^{\circ}$, il resterait à prouver que $2^{\circ}$ suffit à impliquer $4^{\circ}$ (et ceci quel que soit $I$ ensemble ordonné localement fini) -- car en appliquant ceci aux $I_{<i_0}$, on trouve que $2^{\circ}$ implique $\alpha_{i_0} \in TC = C \cap W$ (non seulement $\alpha_{i_0} \in C$, selon la proposition 2), donc $\beta_{i_0} \in TC \subset W$, donc $\tilde{u}_{i_0} \in W$ par saturation, donc $\tilde{u}_{i_0} \in TC$, d'où $3^{\circ}$.

Notons d'autre part que le b) dans la proposition 2 est conséquence de a), appliqué à la restriction de $u_{\bullet}$ à $I_{<i_0}$. Au total, il y a donc à prouver seulement les deux choses
\[
\left\{
\begin{aligned}
&\text{a)} & u_{\bullet} \in C_I &\quad\Longrightarrow\quad u_I \in C \\
&\text{b)} & u_{\bullet} \in C_I \cap W &\quad\Longrightarrow\quad u_I \in C \cap W\ \ (\defeq TC) \\
& & &\qquad\qquad\ \ (\text{et } u_{\bullet} \in (TC)_I)
\end{aligned}
\right.
\tag*{$(*)$}
\]

Nous le prouverons par réduction au cas particulier de la proposition 1, déjà prouvée.

J'ai d'ailleurs aussi envie de prouver le

\textbf{Corollaire 2.} \textit{$\mathcal{M}$, $C$ (et éventuellement $W$). Les hypothèses sur $X_{\bullet}, Y_{\bullet} \in \mathcal{M}(I^{\circ})_{\mathrm{c}}$ et sur $u_{\bullet} \in C_I$, enfin sur $I$ localement fini, sont celles de la proposition 2. Pour toute partie $J$ de $I$, soit}
\[
u_J \defeq \varinjlim_{i \in J} u_i \ :\ X_J \defeq \varinjlim_{i \in J} X_i \ \longrightarrow\ Y_J \defeq \varinjlim_{i \in J} Y_i\ .
\]
\textit{Soient $I' \subset I''$ deux parties} ouvertes \textit{de $I$, et considérons le diagramme}

\msp{page 10}

\[
\begin{tikzcd}
X_{I'} \arrow[r, "\tau_X"] \arrow[d, "\alpha = u_{I'}"'] \arrow[dr, phantom, "\cocart"] & X_{I''} \arrow[d, "\beta"'] \arrow[dd, bend left=30, "u_{I''}"] & \\
Y_{I'} \arrow[r, "\tilde{\tau}_Y"] \arrow[dr, dotted, "\tau_Y"'] & \tilde{Y}_{I''} \arrow[d, "\tilde{u}"'] & \text{\footnotesize (défini comme somme amalgamée)} \\
& Y_{I''}\ . &
\end{tikzcd}
\tag*{$(*)$}
\]

\printedpage{6}
% head: A. Grothendieck : Dérivateurs, Ch. XVII, Édition des Universités de Montpellier II et Paris VII
% foot: Version 9 février 2019

\textit{Alors on a}
\[
\left\{
\begin{aligned}
&\textit{$u_{I'}, u_{I''} \in C$ (donc aussi $\beta \in C$) (i.e. pour tout $J$ ouvert dans $I$,}\\
&\textit{$u_J : X_J \longrightarrow Y_J$ est dans $C$)\,\footnotemark, et $\tau_X, \tau_Y \in C$ (donc aussi $\tilde{\tau}_Y \in C$)}\\
&\textit{(i.e. pour $I' \subset I''$ inclusion d'ouverts, et $X_{\bullet}$ $C$-cofibrant dans $\mathcal{M}(I^{\circ})$,}\\
&\textit{on a $X_{I'} \longrightarrow X_{I''}$ dans $C$)\,\footnotemark, et enfin, même}\\
&\qquad\qquad\qquad\qquad\boxed{\tilde{u} \in C.}
\end{aligned}
\right.
\]
\addtocounter{footnote}{-1}\footnotetext{C'est bien sûr contenu dans prop. 2, appliquée à $u_{\bullet}|_J$.}\stepcounter{footnote}\footnotetext{C'est pour cette assertion $\tau_X \in C$ (d'où $\tilde{\tau}_Y, \tau_Y \in C$, moyennant $\tilde{u} \in C$) qu'on a besoin de l'hypothèse $X_{\bullet} \in \mathcal{M}(I^{\circ})_{\mathrm{c}}$, $\tilde{u} \in C$ vrai sans cette hypothèse.}

\textbf{N.B.} Le cas de la proposition 2 sous la forme $(*)$ a), page 9, (i.e. aussi le fait que les $u_{I'}$, $u_{I''}$, \dots\ sont dans $C$) est obtenu en prenant $I'' = I$, $I' = \emptyset$, le diagramme devient
\[
\begin{tikzcd}
\emptyset \arrow[r] \arrow[d, "\mathrm{id}"'] \arrow[dr, phantom, "\cocart"] & X_I \arrow[d, "\mathrm{id}"] \\
\emptyset \arrow[r] & X_I \arrow[d, "\tilde{u} = u_I"] \\
& Y_I\ ,
\end{tikzcd}
\]
et la conclusion est donc $u_I \in C$\footnote{Le cas $(Y_{I'} \xrightarrow{\tau_Y} Y_{I''}) \in C$ est obtenu en prenant $X_{\bullet} = \emptyset_{\bullet}$, on aura
\[
\begin{tikzcd}[ampersand replacement=\&]
\emptyset \arrow[r, "\tau_X", "\sim"'] \arrow[d] \arrow[dr, phantom, "\cocart"] \& \emptyset \arrow[d] \\
Y_{I'} \arrow[r, "\mathrm{id}"] \arrow[dr, "\tau_Y"'] \& \tilde{Y}_{I''} = Y_{I'} \arrow[d, "\tilde{u} = \tau_Y"] \\
\& Y_{I''}\ .
\end{tikzcd}
\]}. La \textit{définition} de $C_I$, i.e. la condition $u \in C_I$, se retrouve en exigeant $\tilde{u} \in C$ dans $(*)$ ci-dessus, dans le cas $I'' = I/i_0$, $I' = I_{<i_0}$. Il est remarquable
\msp{page 11}
que cette définition de $C_I$ donne une conclusion aussi forte que $\tilde{u} \in C$ pour \textit{tout} couple $I' \subset I''$ de sous-ensembles ouverts de $I$.

C'est finalement sous cette forme qu'on va prouver la proposition 2 -- c'est \textit{ça} à présent qui me semble \textit{le} bon énoncé dans cette direction. Avec le complément, bien sûr, que si de plus $u_{\bullet} \in W_{I^{\circ}}$, i.e. les $u_i$ sont dans $W$, alors $u_{\bullet}$ est même dans $(TC)_I$, ou encore
\printedpage{7}
% head: A. Grothendieck : Dérivateurs, Ch. XVII, Édition des Universités de Montpellier II et Paris VII
% foot: Version 9 février 2019
simplement que les $u_{I'}$ \textit{sont dans} $TC$ (d'où il résultera donc que $\tilde{u}$ dans $(*)$, p. 10, est aussi dans $TC$).

On va procéder par construction transfinie. On peut supposer bien sûr que $I'' = I$. On va quand même d'abord prouver le cas particulier où $I' = \emptyset$, i.e. prouver au moins que $u_I \in C$. On construit par récurrence transfinie une suite strictement croissante d'ouverts
\[
I_{\alpha} \ \subset \ I,
\]
aboutissant à
\[
I_{\alpha_0} \ = \ I,
\]
en posant
\[
\left\{
\begin{aligned}
I_0 \ &= \ \emptyset\,, \\
I_{\alpha} \ &= \ \textstyle\bigcup_{\beta < \alpha} I_{\beta} && \text{si $\alpha$ ordinal limite}\,, \\
I_{\alpha+1} \ &= \ I_{\alpha} \cup (I/s_{\alpha}) && \text{si $I_{\alpha} \neq I$, et en choisissant $s_{\alpha} \in I \backslash I_{\alpha}$}\,.
\end{aligned}
\right.
\]

\msp{page 12}

(On finit par avoir $I_{\alpha_0} = I$, bien sûr.) On pose
\[
\begin{aligned}
X_{\alpha} \ &= \ X_{I_{\alpha}} \ = \ \varinjlim_{i \in I_{\alpha}} X_i \\
Y_{\alpha} \ &= \ Y_{I_{\alpha}} \ = \ \varinjlim_{i \in I_{\alpha}} Y_i \\
u_{\alpha} \ &= \ u_{I_{\alpha}} \ \ :\ \ X_{\alpha} \longrightarrow Y_{\alpha}\ ,
\end{aligned}
\]
homomorphisme de systèmes inductifs transfinis. On aura
\[
X_{\alpha_0} \ = \ X_I, \qquad Y_{\alpha_0} \ = \ Y_I, \qquad u_{\alpha_0} \ = \ u_I,
\]
et on doit donc prouver que $u_{\alpha_0} \in C$. Pour ceci, on applique la proposition 1, il suffit donc de prouver que $(u_{\alpha}) \in C_{I_{\alpha_0+1}}$ (où $I_{\alpha_0+1} = \{\text{ordinaux } i \mid i < \alpha_0 + 1, \text{ i.e. } i \leq \alpha_0\}$). Donc il faut prouver que les flèches $\tilde{u}_{i_0}$ dans $(*)$ p. 1 (où cette fois les $i$, $i_0$ sont des ordinaux, $i_0 \leq \alpha_0$) sont dans $C$. Je vais distinguer les deux cas où $\alpha = i_0$ est ordinal limite, qui est trivial car alors $X_{i_0} \xleftarrow{\ \sim\ } \varinjlim\limits_{i < i_0} X_i$, $Y_{i_0} \xleftarrow{\ \sim\ } \varinjlim\limits_{i < i_0} Y_i$, et $\tilde{u}_{i_0} = \mathrm{id}$, et le cas $\alpha$ non limite, i.e.
\[
\begin{tikzcd}
X_{I_{\alpha-1}} \arrow[r] \arrow[d, "u_{\alpha-1}"'] & X_{I_{\alpha}} \arrow[d, "\beta_{\alpha}"'] \arrow[dd, bend left=30, "u_{\alpha}"] \\
Y_{I_{\alpha-1}} \arrow[r] & \tilde{Y}_{I_{\alpha}} \arrow[d, "\tilde{u}_{\alpha}"'] \\
& Y_{I_{\alpha}}\ ,
\end{tikzcd}
\]
\printedpage{8}
% head: A. Grothendieck : Dérivateurs, Ch. XVII, Édition des Universités de Montpellier II et Paris VII
% foot: Version 9 février 2019
\msp{page 13}
à prouver que $\tilde{u}_{\alpha} \in C$. Lorsqu'on suppose de plus $u_0 \in W_{I_0}$, i.e. les $u_i \in W$, il faut de plus prouver que $\tilde{u}_{\alpha} \in TC = C \cap W$. La proposition 2 et son corollaire 1 résulteront de là. Pour prouver $\tilde{u}_{\alpha} \in C$ resp. $\tilde{u}_{\alpha} \in TC$, on peut supposer $I_{\alpha} = I$ (quitte à changer de notation). On pose donc
\[
\begin{aligned}
I_{\alpha} \ &= \ I, & I_{\alpha-1} \ &= \ I' \\
I/s_{\alpha-1} \ &= \ I_0, & I_0 \cap I' \ &= \ I'_0\ ,
\end{aligned}
\]
de sorte qu'on a un diagramme commutatif d'inclusions ouvertes
\[
\begin{tikzcd}
I_0 \cap I' = I'_0 \arrow[r, hook] \arrow[d, hook] & I_0 \arrow[d, hook] \\
I' \arrow[r, hook] & I = I_0 \cup I'.
\end{tikzcd}
\]
On utilisera seulement le fait que $I_0$, donc $I'_0$, est fini, i.e. que dans le couple $(I' \subset I)$, $I$ peut s'écrire sous la forme $I = I' \cup I_0$, avec $I_0$ un ouvert \textit{fini,} pour en conclure que dans le diagramme
\[
\begin{tikzcd}
X_{I'} \arrow[r, "\tau^X_{I',I}"] \arrow[d, "u_{I'}"'] \arrow[dr, phantom, "\cocart"] & X_I \arrow[d, "v_{I',I}"'] \arrow[dd, bend left=30, "u_I"] \\
Y_{I'} \arrow[r] \arrow[dr, "\tau^Y_{I',I}"'] & \tilde{Y}_I \arrow[d, "\substack{\tilde{u} =\\ u_{I',I}}"'] \\
& Y_I
\end{tikzcd}
\]
on a $\tilde{u} \in C$ resp. $\tilde{u} \in TC = C \cap W$.

\msp{page 14}

Il faut des notations générales, d'une part pour une inclusion $J \subset K$ d'ouverts de $I$, donnant lieu au diagramme
\[
\begin{tikzcd}
X_J \arrow[r, "\tau^X_{J,K}"] \arrow[d, "u_J"'] \arrow[dr, phantom, "\cocart"] & X_K \arrow[d, "\tilde{u}_{J,K}"'] \arrow[dd, dotted, bend left=30, "u_K"] & \\
Y_J \arrow[r] \arrow[dr, dotted, "\tau^Y_{J,K}"'] & Y_{J,K} \arrow[d, "u_{J,K}"'] & \text{\footnotesize (anciennement $\tilde{Y}_K$)} \\
& Y_K\ , &
\end{tikzcd}
\tag*{(1)}
\]
et pour des inclusions de trois ouverts emboîtés $J \subset K \subset L$ de $I$,


\endgroup
% END UNIT der-ch17/der-ch17_p006-010.tex
% BEGIN UNIT der-ch17/der-ch17_p011-015.tex
\begingroup\makeatletter%
\newcommand{\printedpage}[1]{\marginpar{\footnotesize p.~#1}}%
\newcommand{\msp}[1]{\marginpar{\footnotesize [#1]}}%
\newcommand{\cocart}{\text{\scriptsize cocart.}}%
\makeatletter%
  \renewcommand\@makefnmark{\hbox{\textsuperscript{\normalfont(\@thefnmark)}}}%
  \renewcommand\@makefntext[1]{\parindent 1em\noindent\hbox to 1.8em{\hss\textsuperscript{\normalfont\@thefnmark}}#1}%
\makeatother%
\makeatother\let\printedpage\derpp\let\msp\dermsp%
\let\mspin\DERmspin%


\printedpage{9}
% head: A. Grothendieck : Dérivateurs, Ch. XVII, Édition des Universités de Montpellier II et Paris VII
% foot: Version 9 février 2019
\[
\begin{tikzcd}
X_J \arrow[r] \arrow[d, "u_J"'] \arrow[dr, phantom, "\cocart"] & X_K \arrow[r] \arrow[d, "\tilde{u}_{J,K}"'] \arrow[dr, phantom, "\cocart"] & X_L \arrow[d, "\tilde{u}_{J,L}"'] \arrow[dd, bend left=40, "\tilde{u}_{K,L}"] \\
Y_J \arrow[r] & Y_{J,K} \arrow[r] \arrow[d, "u_{J,K}"'] \arrow[dr, phantom, "\cocart"] & Y_{J,L} \arrow[d, "u^L_{J,K}"'] \arrow[dd, bend left=40, "u_{J,L}"] \\
 & Y_K \arrow[r] & Y_{K,L} \arrow[d, "u_{K,L}"'] \\
 & & Y_L\ ,
\end{tikzcd}
\tag*{(2)}
\]
donnant lieu à une nouvelle flèche $u^L_{J,K}$, déduite de $u_{J,K}$ par le cochangement de base \msp{page 15} $Y_{J,K} \longrightarrow Y_{J,L}$. Nous utiliserons ce diagramme dans les deux cas où $(J \subset K \subset L)$ est l'un des deux triplets
\[
\left\{
\begin{aligned}
& I'_0 \hookrightarrow I_0 \hookrightarrow I \\
& I'_0 \hookrightarrow I' \hookrightarrow I,
\end{aligned}
\right.
\]
issus du carré cocartésien (dans Cat) (carré MV)
\[
\begin{tikzcd}
I'_0 \arrow[r, hook] \arrow[d, hook] & I_0 \arrow[d, hook] \\
I' \arrow[r, hook] & I\ .
\end{tikzcd}
\tag*{(3)}
\]

Il faut insérer les deux diagrammes correspondants du type $(*)$ \texttt{[de la page 10]} dans un même diagramme. On notera à ce propos que les carrés dans $\mathcal{M}$ déduits de (3) en lui appliquant $J \longmapsto X_J$ et $J \longmapsto Y_J$
\[
\begin{matrix} (4_X) \\[12ex] (4_Y) \end{matrix}
\qquad\qquad
\begin{tikzcd}[column sep=small, row sep=small]
X_{I'_0} \arrow[rr] \arrow[dr] \arrow[dd, dotted, "u_{I'_0}"] & & X_{I_0} \arrow[dr] \arrow[dd, dotted, "u_{I_0}"] & \\
 & X_{I'} \arrow[rr] \arrow[dd, dotted, "u_{I'}"'] & & X_I \arrow[dd, dotted, "u_I"'] \\
Y_{I'_0} \arrow[rr] \arrow[dr] & & Y_{I_0} \arrow[dr] & \\
 & Y_{I'} \arrow[rr] & & Y_I
\end{tikzcd}
\qquad\qquad
\left.\vrule height 9ex depth 9ex width 0pt\right\}\ (4)
\]
sont cocartésien dans $\mathcal{M}$. (Transitivité des $\varinjlim$ \ldots). % [sic] cocartésien -> cocartésiens, agreement with « les carrés »

\msp{page 16}
On a un homomorphisme $(u_{I'_0}, u_{I_0}, u_{I'}, u_I)$ déduit de $u$ du carré $(4_X)$ dans les carré $(4_Y)$, % [sic] les carré -> le carré, singular noun with plural article
et cet homomorphisme lui aussi s'insère dans un diagramme, incluant les deux avatars déjà signalé de (2). Voici le diagramme pertinent : % [sic] signalé -> signalés, agreement with « les deux avatars »
\printedpage{10}
% head: A. Grothendieck : Dérivateurs, Ch. XVII, Édition des Universités de Montpellier II et Paris VII
% foot: Version 9 février 2019
% the footnote mark (9) stands in the printed label of the display; \footnotemark/\footnotetext used because \footnote cannot be set inside a display tag
\[
\begin{tikzcd}[column sep=small, row sep=small]
 & X_{I_0} \arrow[rrr] \arrow[dd] & & & X_I \arrow[dd] \\
X_{I'_0} \arrow[ur] \arrow[rr] \arrow[dd] & & X_{I'} \arrow[urr] \arrow[dd] & & \\
 & Y_{I'_0,I_0} \arrow[rrr] \arrow[dd, "u_{I'_0,I_0}"] & & & Y_{I'_0,I} \arrow[dd, "u^I_{I'_0,I_0}"] \arrow[dddl] \\
Y_{I'_0} \arrow[ur] \arrow[rr] \arrow[dr, dotted] \arrow[dddrr, dotted] & & Y_{I'_0,I'} \arrow[urr] \arrow[ddd] & & \\
 & Y_{I_0} \arrow[rrr] \arrow[dddrrr, dotted] & & & Y_{I_0,I} \arrow[ddd, "u_{I_0,I}"] \\
 & & & Y_{I',I} \arrow[ddr, "u_{I',I}"] & \\
 & & Y_{I'} \arrow[ur] \arrow[drr, dotted] & & \\
 & & & & Y_I\ .
\end{tikzcd}
\tag*{(5~\footnotemark[9])}
\]
\footnotetext[9]{diagramme imbitable, à remplacer par diagramme (6) ci-dessus.}

Le cube supérieur est le cube dans $\mathcal{M}$ à 6 faces cocartésiens déduit des trois flèches de $X_{I'_0}$ dans $X_{I_0}$, $X_{I'}$, $Y_{I'_0}$. % [sic] faces cocartésiens -> faces cocartésiennes, gender agreement
On aura besoin de savoir que la \textit{base} du cube est un carré cocartésien. La partie antérieure du diagramme est l'avatar de (2) correspondant aux inclusions $I'_0 \subset I' \subset I$, la partie derrière l'avatar pour l'inclusion $I'_0 \subset I_0 \subset I$. Donc les deux autres carrés verticaux, en dessous du cube supérieur, sont également cocartésiens, ainsi que le carré oblique en pointillés, lequel n'est autre que la base $(4_Y)$ du diagramme (4). De ces caractères du diagramme (5) résulte aisément (quoique \msp{page 17} j'ai mis une bonne journée de gribouillis en bonheur-la-chance, avant de m'en rendre compte enfin \ldots) que le carré en losange à traits pleins
\[
\begin{tikzcd}[column sep=small, row sep=small]
 & Y_{I'_0,I} \arrow[dl] \arrow[dr] & \\
Y_{I',I} \arrow[dr] & & Y_{I_0,I} \arrow[dl] \\
 & Y_I &
\end{tikzcd}
\]
est également cocartésien. On le voit trivialement sur le diagramme déployé
\printedpage{11}
% head: A. Grothendieck : Dérivateurs, Ch. XVII, Édition des Universités de Montpellier II et Paris VII
% foot: Version 9 février 2019
\[
\begin{tikzcd}[column sep=tiny, row sep=small]
 & & Y_{I'_0} \arrow[dl] \arrow[dr] & & \\
 & Y_{I'_0,I'} \arrow[dl] \arrow[dr] & \cocart & Y_{I'_0,I_0} \arrow[dl] \arrow[dr, "u_{I'_0,I_0}"] & \\
Y_{I'} \arrow[dr] & \cocart & Y_{I'_0,I} \arrow[dl] \arrow[dr, "u^I_{I'_0,I_0}"] & \cocart & Y_{I_0} \arrow[dl] \\
 & Y_{I',I} \arrow[dr, "u_{I',I}"] & \cocart & Y_{I_0,I} \arrow[dl] & \\
 & & {?\simeq Y_I \simeq Y_{I'} \amalg_{Y_{I'_0}} Y_{I_0}\ ,} & &
\end{tikzcd} % ?? uncertain: coproduct sign -- \amalg or \sqcup
\tag*{(6)}
\]
qui est finalement le diagramme pertinent (le diagramme (5) ne servait à rien, c'est un mauvais bout par lequel j'ai pris les choses). On voit sur le diagramme (6) que pour prouver $u_{I',I} \in C$ resp. $TC$, il \msp{page 18} suffit de prouver $u_{I'_0,I_0} \in C$ resp. $TC$. Mais c'est là une assertion relative à l'inclusion d'ensembles ordonnés $I'_0 \subset I_0$ \textit{finis.} On est donc réduit à prouver le lemme suivant, où $I_0$ devient $I$ :

\textbf{Lemme.} \textit{Soit $I$ un ensemble ordonné fini, $u_\bullet : X_\bullet \longrightarrow Y_\bullet$ une flèche dans $\mathcal{M}(I^\circ)$, où $\mathcal{M}$ est une catégorie munie d'un ensemble $C \subset \mathrm{Fl}(\mathcal{M})$ de flèches tel que $C$ soit stable par cochangement de base et par composition, et contienne les isomorphismes. On suppose $u_\bullet \in C_I$, i.e. que pour tout $i \in I$, la flèche $u_{I_{<i},I/i}$ est dans $C$. Alors pour toute partie ouverte $I'$ de $I$, la flèche $u_{I',I}$ est dans $C$}~\footnote[10]{\textbf{N.B.} Bien sûr, $u_{J,K} \in C$ si $J \subset K$ sont deux ouverts de $I$, comme on voit en appliquant le lemme à $K$ au lieu de $I$.}\textit{.}

\textbf{Corollaire.} \textit{Si de plus $W \ \subset\ \mathrm{Fl}(\mathcal{M})$ est un localiseur modérément saturé, tel que $TC \overset{\textup{déf}}{=} W \cap C$ soit stable par cochangement de base (et l'est de toutes façons par composition, et contient les isomorphismes), alors on a}
\[
C_I \cap W_{I^\circ} \ = \ (TC)_I,
\]
\textit{donc si $u_\bullet \in C_I$ (cf. ci-dessus) et les $u_i \in W$, alors $u_\bullet \in (TC)_I$, donc en appliquant ce qui est dit pour $C$ à $TC$, on a $u_{I',I} \in TC$ pour toute} \msp{page 19} \textit{partie ouverte $I'$ de $I$.}

On prouvera en même temps le lemme et son corollaire, en procédant par récurrence sur $\operatorname{card} I = n$. Si $n = 0$, i.e. $I = \emptyset$, les énoncés sont triviaux. Par transitivité, en utilisant les diagrammes du type (2), on est ramené au cas où $I = I' \cup \{i_0\}$, où $i_0 \in I \backslash I'$. On pose encore $I_0 = I/i_0$, $I'_0 = I_0 \cap I'$, et le raisonnement qu'on vient de faire nous montre qu'il suffit de prouver que $u_{I'_0,I_0} \in C$ resp. $\in TC$. Si $I'_0 \neq I$, cela résulte de l'hypothèse de récurrence. % [sic] I'_0 \neq I -> I_0 \neq I, the alternative « Sinon, i_0 est un objet final de I » requires I_0 = I
Sinon, $i_0$ est un objet final de $I$, et $I' = I_{<i_0}$, et alors $u_{I',I} \in C$ est inclus dans la définition de $C_I$. Reste à prouver alors que $u_{I',I} \in TC$ si les $u_i \in W$. On a le diagramme
\printedpage{12}
% head: A. Grothendieck : Dérivateurs, Ch. XVII, Édition des Universités de Montpellier II et Paris VII
% foot: Version 9 février 2019
\[
\begin{tikzcd}
X_{I'} \arrow[r] \arrow[d, "u_{I'}"'] \arrow[dr, phantom, "\cocart"] & X_I \arrow[d, "\tilde{u}_{I',I}"'] \arrow[dd, dotted, bend left=40, "u_I"] \\
Y_{I'} \arrow[r] & Y_{I',I} \arrow[d, "u_{I',I}"'] \\
 & Y_I\ ,
\end{tikzcd}
\]
\msp{page 20} et l'hypothèse de récurrence implique $u_{I'} \in TC$, donc $\tilde{u}_{I',I} \in TC$ par cochangement de base. D'autre part, la flèche $u_I : X_I \longrightarrow Y_I$ s'identifie à $X_{i_0} \longrightarrow Y_{i_0}$ ($i_0$ étant élément final de $I$), elle est donc dans $W$. Donc par saturation, $u_{I',I} \in W$, donc $u_{I',I} \in TC$, q.e.d.

Ainsi, la proposition 2 est prouvée, ainsi que son corollaire 1 -- et on voit que si dans la proposition 2 et son corollaire, on suppose $I$ \textit{fini,} alors il est inutile dans les hypothèses pour $C$, resp. $TC = C \cap W$, de supposer la stabilité par composition \textit{transfinie,} la composition ordinaire suffit. Mais il reste à prouver le corollaire 2, que je viens de prouver dans le cas $I$ \textit{fini,} sans avoir à faire appel à une hypothèse ``transfinie''. Il faut maintenant prouver que c'est encore OK dans le cas localement fini, en supposant $C$ stable par composition transfinie. (Il n'y a plus à revenir sur le cas de $W$, car on sait déjà maintenant que $C_I \cap W_{I^\circ} = (TC)_I$, par le corollaire 1 déjà prouvé.)

\msp{page 21}
La démonstration est essentiellement la même que pour le cas particulier déjà traité $I' = \emptyset$, $I'' = I$, et on aurait pu aussi bien rédiger la démonstration dans ce cadre, à l'aide d'une suite strictement croissante transfinie d'ouverts dans $I$, $I_\alpha$, avec \fbox{$I_0 = I',$} $I_{\alpha_0} = I''$ $(= I)$, et
\[
\begin{aligned}
I_\alpha \ &= \ I_{\alpha-1} \cup I/s_{\alpha-1} && \text{si $\alpha$ non limite, $I_{\alpha-1} \neq I$, en choisissant $s_{\alpha-1} \in I \backslash I_{\alpha-1}$,} \\
I_\alpha \ &= \ \textstyle\bigcup_{\beta<\alpha} I_\beta && \text{si $\alpha$ limite,}
\end{aligned}
\]
et en posant encore
\[
X_\alpha \ = \ X_{I_\alpha}, \qquad Y_\alpha \ = \ Y_{I_\alpha}, \qquad u_\alpha = u_{I_\alpha} \ : \ X_\alpha \longrightarrow Y_\alpha.
\]
Soit $\alpha_0$ tel que $I_{\alpha_0} = I$ (indice terminal de la suite), il faut prouver que dans
\[
\begin{tikzcd}
X_0 \arrow[r] \arrow[d] & X_{\alpha_0} \arrow[d] \\
Y_0 \arrow[r] & Y_{0,\alpha_0} \arrow[d, "u_{0,\alpha_0}"] \\
 & Y_{\alpha_0}\ ,
\end{tikzcd}
\]
\printedpage{13}
% head: A. Grothendieck : Dérivateurs, Ch. XVII, Édition des Universités de Montpellier II et Paris VII
% foot: Version 9 février 2019
\msp{page 22} on a
\[
u_{0,\alpha_0} \ \in \ C.
\]

Compte tenu du complément marginal à la proposition 1 (qui se prouve de la même façon \ldots), on est ramené à prouver que l'homomorphisme des systèmes ordinaux envisagés est dans $C_{I_{\alpha_0+1}}$ (où $I_{\alpha_0+1}$ est l'ensemble des ordinaux $\leq \alpha_0$). Or qu'il en soit ainsi se prouve exactement comme tantôt, à cela près qu'on doit d'abord faire une vérification de $(*)$ p.~1 ($\tilde{u}_\bullet \in C$) pour $i_0 = 0$. Alors le diagramme $(*)$ p.~1 est
\[
\begin{tikzcd}
\emptyset \arrow[r] \arrow[d, "\mathrm{iso}"] & X_0 = X_{I'} \arrow[d, "\mathrm{iso}"', "\mathrm{id}"] \\
\emptyset \arrow[r] & X_0 = X_{I'} \arrow[d, "u_0 = u_{I'}"] \\
 & Y_0 = Y_{I'}\ ,
\end{tikzcd}
\]
et il faut donc vérifier $u_0 = u_{I'} \in C$. Or ceci est contenu dans la proposition 2 elle-même, déjà prouvé, qu'on doit donc utiliser. % [sic] prouvé -> prouvée, agreement with « la proposition 2 »
Le reste se prouve comme avant. Ainsi, le corollaire 2 de la proposition 2 est également prouvé.

\textbf{Commentaires :} Il y a un malaise, car les proposition 1 et proposition 2 (ainsi que leurs corollaires) ont essentiellement la même conclusion, avec des hypothèses sur $I$ qui sont \msp{page 23} différentes : bien ordonné dans un cas, localement fini dans un autre. L'idée naturelle, c'est de donner un énoncé commun, qui coiffe les deux à la fois, avec $I$ artinien, i.e. toute suite décroissante est stationnaire. Ce qu'il faut, c'est une hypothèse sur $I$ d'une part, sur $u_\bullet : X_\bullet \longrightarrow Y_\bullet$ dans $\mathcal{M}(I^\circ)$ d'autre part, qui assure ceci :

\begin{enumerate}
\item[$(*)$] $\forall\ i_0 \in I$, et toute partie ouverte $J'$ de $J = I/i_0$, la flèche $u_{J',J}$ dans le diagramme
\[
\begin{tikzcd}
X_{J'} \arrow[r] \arrow[d, "u_{J'}"'] \arrow[dr, phantom, "\cocart"] & X_J \arrow[d, "\tilde{u}_{J',J}"'] \arrow[dd, bend left=40, "u_J"] \\
Y_{J'} \arrow[r] & Y_{J',J} \arrow[d, "u_{J',J}"'] \\
 & Y_J
\end{tikzcd}
\]
est dans $C$.
\end{enumerate}

\textbf{Question :} est-il exacte qu'il suffit, pour la validité de $(*)$, de supposer % [sic] exacte -> exact, agreement with « il »
\begin{enumerate}
\item[a)] $u_\bullet \in C_I$, i.e. $u_{J',J} \in C$ quand $J = I/i_0$, $J' = J \backslash \{i_0\}$, et
\item[b)] \textit{$I$ localement fini ?}
\end{enumerate}


\endgroup
% END UNIT der-ch17/der-ch17_p011-015.tex
% BEGIN UNIT der-ch17/der-ch17_p016-020.tex
\begingroup\makeatletter%
\newcommand{\printedpage}[1]{\marginpar{\footnotesize p.~#1}}%
\newcommand{\msp}[1]{\marginpar{\footnotesize [#1]}}%
\newcommand{\Fl}{\mathrm{Fl}}%
\newcommand{\M}{\mathcal{M}}%
\makeatother\let\printedpage\derpp\let\msp\dermsp%
\let\mspin\DERmspin%


\printedpage{14}
% head: A. Grothendieck : Dérivateurs, Ch. XVII, Édition des Universités de Montpellier II et Paris VII

\msp{page 24}

\textbf{Lemme.} \textit{Soit $P \subset I$ l'ensemble des $i_0 \in I$ tels que pour tout ouvert $K$ de $J = I/i_0$, la flèche $u_{K,J}$ est dans $C$. Ceci dit, soit $i_0 \in I$ tel que tout $i < i_0$ soit dans $P$. Alors $i_0 \in P$.} (On suppose $u_\bullet \in C_I$, mais $I$ quelconque.)

En effet, soit $K \subset J = I/i_0$ un ouvert. Si $K = J$, il est trivial que $u_{K,J} \in C$ (c'est un isomorphisme). Si $K = J$, on a % [sic] Si K = J (second occurrence) -> Si K \neq J, the first case already treats K = J
\[
K \ \subset\ J' \overset{\text{déf}}{=} J\backslash\{i_0\} \ \subset\ J,
\]
et pour prouver que $u_{K,J} \in C$, il suffit de prouver que $u_{K,J'}, u_{J',J} \in C$. Or $u_{J',J} \in C$ par hypothèse $u_\bullet \in C_I$. Et $u_{K,J'} \in C$ par ce que a été prouvé pour le cas de $J'$, compte tenu que par hypothèse sur $i_0$, tout $i \in J'$ (i.e. tout $i < i_0$) est dans $P$. % [sic] par ce que a été prouvé -> par ce qui a été prouvé, grammar

\textbf{Corollaire.} \textit{Si $I$ est artinien, $P = I$.}

Car toute partie $P$ de $I$ ayant la propriété qu'on vient de dire ($I/i_0 \subset P \Longrightarrow i_0 \in P$) est égale à $I$. Si en effet il y avait $x_0 \notin P$, il existerait $x_1 < x_0$ avec $x_0 \notin P$, et par récurrence on construirait une suite strictement décroissante dans $I$, absurde. % [sic] avec x_0 \notin P -> avec x_1 \notin P, the new element must lie outside P

\msp{page 25}

\section*{2\quad Description de $C_I \subset \Fl(\M(I^\circ))$, pour $I$ ensemble ordonné}

Soit $\M$ une catégorie, stable par (petites) $\varinjlim$ , $C \subset \Fl(\M)$ (“cofibrations”). On suppose
\[
\left\{
\begin{array}{lll}
\text{a)} & C & \text{contient les isomorphismes et est stable}\\
 & & \text{par cochangement de base.}\\[1ex]
\text{b)} & C & \text{stable par composition transfinie.}
\end{array}
\right.
\tag*{(2.1)}
\]

Soit $I$ un (petit) ensemble ordonné, et considérons
\[
\M(I^\circ) \ =\ \underline{\mathrm{Hom}}(I,\M),
\]
on va définir une partie
\[
C_I \ \subset\ \Fl(\M(I^\circ))
\]
ainsi, pour $u_\bullet : X_\bullet \longrightarrow Y_\bullet$ flèche dans $\M(I^\circ)$ :
\[
u_\bullet \in C_I \quad\Longleftrightarrow\quad
\begin{array}{l}
\forall\ i_0 \in I \text{ et pour tout ouvert } J \subset I/i_0 = K,\\
\text{la flèche } u_{J,K} \text{ est dans } C.
\end{array}
\tag*{(2.2)}
\]

Voici la définition de $u_{J,K}$. Pour tout ouvert $J$ de $I$, on pose
\[
\left\{
\begin{array}{rclrcl}
X_J & = & \varinjlim\limits_{i\in J} X_i, & Y_J & = & \varinjlim\limits_{i\in J} Y_i,\\[2ex]
u_J : X_J & \longrightarrow & Y_J, & u_J & = & \varinjlim\limits_{i\in J} u_i,
\end{array}
\right.
\tag*{(2.3)}
\]
et on considère, pour $K \subset J$, le diagramme
% foot: Version 9 février 2019

\printedpage{15}
% head: A. Grothendieck : Dérivateurs, Ch. XVII, Édition des Universités de Montpellier II et Paris VII

\msp{page 26}

\[
\begin{tikzcd}
X_K \arrow[r, "\tau^X_{K,J}"] \arrow[d, "u_K"'] \arrow[dr, phantom, "\text{\scriptsize cocart.}"] & X_J \arrow[d, "\tilde{u}_{K,J}"'] \arrow[dd, bend left=40, "u_J"] \\
Y_K \arrow[r, "\tau^Y_{K,J}"] & Y_{K,J} \rlap{$\ = Y_K \amalg_{X_K} X_J$} \arrow[d, "u_{K,J}"'] \\
 & Y_J \ .
\end{tikzcd}
\]

\textbf{Théorème 2.1.} \textit{Soit $\M$ $\mathfrak{U}$-catégorie stable par petites $\varinjlim$ , $C \subset \Fl(\M)$, satisfaisant les conditions de stabilité} \textup{(2.1)} \textit{(par cochangement de base et par composition transfinie), $I$ un petit ensemble ordonné, $u_\bullet : X_\bullet \longrightarrow Y_\bullet$ une flèche dans $\M(I^\circ)$.}
\begin{itemize}[leftmargin=3em]
\item[$1^\circ$)] \textit{On suppose $u_\bullet \in C_I$ (cf.} \textup{(2.2)}\textit{). Soient $K \subset J$ deux ouverts de $I$. Alors $u_{K,J}$ (défini dans} \textup{(2.4)}\textit{) est dans $C$.}
\item[$2^\circ$)] \textit{Supposons $I$ artinien. Pour qu'on ait $u_\bullet \in C_I$, i.e. pour que $u_{K,J}$ soit dans $C$ pour tout couple $(K \subset J)$ d'ouverts tel que $J$ ait un plus grand élément, i.e. soit de la forme $I/i_0$ (i.e. aussi, par} $1^\circ$\textit{), pour} \textup{tout} \textit{couple d'ouverts $J \subset K$), (il faut et) il suffit qu'il en soit ainsi dans le cas particulier où $J = I/i_0$, $J' = I'/i_0 \overset{\text{déf}}{=} I/i_0 - \{i_0\} = \{i \in I \mid i < i_0\}$.}

\msp{page 27}
\item[$3^\circ$)] \textit{Supposons encore $I$ artinien}\footnote[11]{J'ignore si cette hypothèse est nécessaire, cf. commentaire p. 39.}\textit{. Soit $W \subset \Fl(\M)$ un localiseur modérément saturé, tel que $TC \overset{\text{déf}}{=} C \cap W$ satisfasse également aux conditions de stabilité} \textup{(2.1)} \textit{(par cochangement de base, par composition transfinie). Les conditions suivantes sur une flèche $u_\bullet : X_\bullet \longrightarrow Y_\bullet$ dans $\M(I^\circ)$ sont équivalentes.}
\begin{itemize}[leftmargin=2em]
\item[a)] \textit{$u_\bullet \in (TC)_I$, i.e. pour tout $i_0 \in I$ et $K$ ouvert de $J = I/i_0$, la flèche $u_{K,J}$ est dans $TC$.}
\item[b)] \textit{$u_\bullet \in C_I \cap W_{I^\circ}$, i.e. pour tout $i_0 \in I$, on a $u_{i_0} \in W$, et pour tout ouvert $K$ de $I/i_0 = J$, on a $u_{K,J} \in C$.}
\end{itemize}
\end{itemize}

Donnons tout de suite les corollaires principaux.

\textbf{Corollaire 2.2.} \textit{Supposons $(u_\bullet : X_\bullet \longrightarrow Y_\bullet) \in C_I$. Alors}
\begin{itemize}[leftmargin=3em]
\item[a)] \textit{$\forall\ i \in I$, $u_i : X_i \longrightarrow Y_i$ est dans $C$.}
\item[b)] \textit{La flèche}
\[
u_I \ :\ X_I \longrightarrow Y_I
\tag*{(2.2.1)}
\]
\textit{est dans $C$. Plus généralement, pour toute partie ouverte $J$ de $I$, la flèche}
\[
u_J \ :\ X_J \longrightarrow Y_J
\tag*{(2.2.2)}
\]
\textit{est dans $C$.}
\end{itemize}
% foot: Version 9 février 2019

\printedpage{16}
% head: A. Grothendieck : Dérivateurs, Ch. XVII, Édition des Universités de Montpellier II et Paris VII

\msp{page 28}

Pour a), on applique $1^\circ$) (ou simplement la \textit{définition} de $C_I$) dans le cas de l'inclusion d'ouverts $\emptyset \hookrightarrow I/i$, vu qu'on a
\[
u_{\emptyset,I/i} \ \simeq\ u_i \ .
\]
Pour b), on applique $1^\circ$) à l'inclusion $\emptyset \longrightarrow I$, vu que
\[
u_{\emptyset,I} \ \simeq\ u_I \ ,
\]
et plus généralement
\[
u_{\emptyset,J} \ \simeq\ u_J \ .
\]

\textbf{Corollaire 2.3.} \textit{Supposons $u_\bullet \in C_I \cap W_{I^\circ}$. Alors pour toute inclusion d'ouverts $K \subset J$ de $I$, on a $u_{K,J} \in TC$, et en particulier (prenant $K = \emptyset$), on a $u_J \in TC$. Enfin, pour tout $i \in I$, on a $u_i \in TC$.}

\textbf{Corollaire 2.4.} \textit{Soit $X_\bullet$ dans $\M(I^\circ)_{\mathrm{c}}$ un système inductif $C$-cofibrant, i.e. tel que $\emptyset_I \longrightarrow X_\bullet$ soit dans $C_I$, i.e. tel que pour tout $i_0 \in I$ et tout ouvert $K$ de $J = I/i_0$, la flèche de transition $\tau_{K,J} : X_K = \varinjlim\limits_{i\in K} X_i \longrightarrow X_J \simeq X_{i_0}$ soit dans $C$. Alors pour} \textup{tout} \textit{couple d'ouverts $K \subset J$, la flèche}
\[
\tau_{K,J} \ :\ \varinjlim_{i\in K} X_i \longrightarrow \varinjlim_{i\in J} X_i
\]
\textit{est dans $C$}\footnote[12]{et réciproquement, caractérise les objets $C$-cofibrants.}\textit{. En particulier (faisant $K = I/k$, $J = I/j$ pour $k \leq j$), si $k \leq j$ sont dans $I$, le morphisme de transition}
\[
\tau_{k,j} \ :\ X_k \longrightarrow X_j
\]
\msp{page 29}
\textit{est dans $C$. Enfin, si $W \subset \Fl(\M)$ est comme dans} \textup{(2.1.\ $3^\circ$)}\textit{, et si $\emptyset_I \longrightarrow X$ est dans $W_{I^\circ}$, i.e. les $X_i$ sont $W$-coasphériques, alors $X_\bullet$ est dans $\M(I^\circ)_{TC}$ (où $TC = C \cap W$), et pour tout ouvert $J$ de $I$, $X_J = \varinjlim\limits_{i\in J} X_i$ est dans $\M_{TC}$, donc $X_J$ est $W$-coasphérique, et pour tout couple d'ouverts $K \subset J$, la flèche de transition $\tau_{K,J}$ est dans $TC$.}

Signalons enfin le complément suivant au théorème 2.1, qui résultera de sa démonstration :

\textbf{Corollaire 2.5.} \textit{Supposons que l'espace topologique associé à $I$ (dont les ouverts sont les ouverts de $I$, et les $I/i$ forment une base d'ouverts) soit} \textup{noethérien}\textit{, i.e. toute suite croissante d'ouverts est stationnaire. Alors toutes les conclusions dans le théorème} \textup{2.1} \textit{et ses corollaires restent valables, en supposant seulement la stabilité de $C$ ou de $TC = C \cap W$ par composition simple (au lieu de composition transfinie), en plus des deux autres stabilités.}
% foot: Version 9 février 2019

\printedpage{17}
% head: A. Grothendieck : Dérivateurs, Ch. XVII, Édition des Universités de Montpellier II et Paris VII

\msp{page 30}

\textsc{Démonstration} du théorème 2.1 et du corollaire 2.5.

2.1. $1^\circ$) On peut supposer $J = I$, $K \subset I$. On construit une suite transfinie strictement croissante d'ouverts $I_\alpha$, avec
\[
\left\{
\setlength{\arraycolsep}{4pt}
\begin{array}{rcll}
I_0 & = & K, & \\[1ex]
I_{\alpha+1} & = & I_\alpha \cup I/s_\alpha & \text{où } s_\alpha \in I\backslash I_\alpha \text{ (si } I_\alpha \neq I\text{, sinon, la construction transfinie s'arrête),}\\[1ex]
I_\alpha & = & \bigcup_{\beta<\alpha} I_\beta & \text{si } \alpha \text{ ordinal limite.}
\end{array}
\right.
\]

Il existe un $\alpha_0$ tel qu'on ait $I_{\alpha_0} = I$. On voit alors comme au paragraphe 1
\msp{page 31}
que pour tout $\alpha$, on a
\[
u_{I_\alpha,I_{\alpha+1}} \ \in\ C,
\tag*{$(*)$}
\]
car il est déduit de $u_{(I/s_\alpha)\cap I_\alpha,I/s_\alpha}$ par cochangement de base. Cf. diagramme (6), p. 17, qu'il faudrait dégager en lemme, débarrassé des particularités techniques de la situation, i.e. un lemme sur les homomorphismes de carrés cartésiens :

\textbf{Lemme 2.6.} \textit{Soit un homomorphisme de carrés cocartésiens $X_\bullet$, $Y_\bullet$, dans une catégorie $\M$ avec sommes amalgamées (\textbf{N.B.} il suffit que les sommes amalgamées intervenant ci-dessus existent) :}
\[
\begin{tikzcd}[row sep=small, column sep=small]
X_0 \arrow[rr] \arrow[dr] \arrow[dd, "f_0"'] & & X_1 \arrow[dr] \arrow[dd, "f_1" near end] & \\
 & X'_0 \arrow[rr, crossing over] & & X'_1 \arrow[dd, "f'_1"] \\
Y_0 \arrow[rr] \arrow[dr] & & Y_1 \arrow[dr] & \\
 & Y'_0 \arrow[rr] \arrow[from=uu, crossing over, "f'_0" near start] & & Y'_1 \ .
\end{tikzcd}
\tag*{(2.6.1)}
\]

\textit{Soient $Z_0$, $Z_1$ les deux carrés latéraux et considérons, pour tout carré commutatif}
$Z = \begin{tikzcd}[ampersand replacement=\&, baseline=(\tikzcdmatrixname.center)]
X \arrow[r] \arrow[d] \& X' \arrow[d] \\
Y \arrow[r] \& Y'
\end{tikzcd}$ \textit{, la flèche $u_Z$ du diagramme}
\[
\begin{tikzcd}
X \arrow[r] \arrow[d] \arrow[dr, phantom, "\text{\scriptsize cocart.}"] & X' \arrow[d] & \\
Y \arrow[r] & \tilde{Z} \arrow[r, "u_Z"'] & Y' \ .
\end{tikzcd}
\]
% foot: Version 9 février 2019

\printedpage{18}
% head: A. Grothendieck : Dérivateurs, Ch. XVII, Édition des Universités de Montpellier II et Paris VII

\msp{page 32}

\textit{Alors la flèche $u_{Z_1}$ se déduit de $u_{Z_0}$ par cochangement de base, de façon plus précise, on a un carré cocartésien, où} $T = \begin{tikzcd}[ampersand replacement=\&, baseline=(\tikzcdmatrixname.center)]
X_0 \arrow[r] \arrow[d] \& X_1 \arrow[d] \\
Y_0 \arrow[r] \& Y_1
\end{tikzcd}$ \textit{et} $T' = \begin{tikzcd}[ampersand replacement=\&, baseline=(\tikzcdmatrixname.center)]
X'_0 \arrow[r] \arrow[d] \& X'_1 \arrow[d] \\
Y'_0 \arrow[r] \& Y'_1
\end{tikzcd}$ \textit{désignent la face postérieure et antérieure du cube} \textup{(2.6.1) :}
\[
\begin{tikzcd}
\tilde{Z}_0 \arrow[r, "u_{Z_0}"] \arrow[d] \arrow[dr, phantom, "\text{\scriptsize cocart.}"] & Y'_0 \arrow[d] \\
\tilde{Z}_1 \arrow[r, "u_{Z_1}"] & Y'_1 \ ,
\end{tikzcd}
\tag*{$(*)$}
\]
\textit{qu'on voit sur la base du diagramme}
\[
\begin{tikzcd}[row sep=small, column sep=small]
X_0 \arrow[rr] \arrow[dr] \arrow[dd, "f_0"'] & & X_1 \arrow[dr] \arrow[dd, "f_1" near end] & & \\
 & X'_0 \arrow[rr, crossing over] & & X'_1 \arrow[dd] \arrow[dddr, "f'_1"] & \\
Y_0 \arrow[rr] \arrow[dr] \arrow[dr, phantom, shift left=2ex, "\text{\tiny cocart.}"] & & Y_1 \arrow[dr] & & \\
 & \tilde{Z}'_0 \arrow[rr] \arrow[dr, "u_{Z_0}"'] \arrow[from=uu, crossing over] & & \tilde{Z}'_1 \arrow[dr, "u_{Z_1}"'] & \\
 & & Y'_0 \arrow[rr] \arrow[from=uuul, crossing over, "f'_0" near end] \arrow[ur, phantom, "\text{\tiny cocart.}"] & & Y'_1 \ ,
\end{tikzcd} % ?? uncertain: arrow to which the label f'_0 belongs and exact positions of the two "cocart." labels -- f'_0 printed beside the long diagonal X'_0 -> Y'_0 near Z̃'_0 -> Z̃'_1
\tag*{$(**)$}
\]
\textit{dans lequel le cube est un cube totalement cocartésien, déduit du carré cocartésien $X$ en l'“exponentiant” par la flèche $f_0 : X_0 \longrightarrow Y_0$} \texttt{[plutôt, par la flèche} $(f_0,f_1)$ \texttt{de} $\underline{\Fl}(\M)$\texttt{]}\textit{.}

La base du cube, et le carré
\msp{page 33}
composé de celui-ci et de $(*)$ (base du diagramme polyédral $(**)$) étant cartésien, il en est de même de $(*)$.

\textbf{Lemme 2.7.} \textit{Soit $(I_\alpha)_{\alpha\leq\alpha_0}$ une suite ordinale croissante d'ouverts de $I$. Alors $u_{I_0,I_{\alpha_0}} : Y_{I_0,I_{\alpha_0}} \longrightarrow Y_{I_{\alpha_0}}$ est composé transfini, indexé par l'ensemble des ordinaux $-1 \leq \alpha < \alpha_0$, d'une suite de flèches $v_\alpha$, où chaque $v_\alpha$ est déduit de $u_{I_\alpha,I_{\alpha+1}}$ par cochangement de base (et de $u_{I_0} : X_{I_0} \longrightarrow Y_{I_0}$ pour $\alpha = 0$).}

Pour le voir, on introduit la suite transfinie d'objets $Z_\alpha = Y_{I_\alpha,I_{\alpha_0}}$, définis par le diagramme
% foot: Version 9 février 2019


\endgroup
% END UNIT der-ch17/der-ch17_p016-020.tex
% BEGIN UNIT der-ch17/der-ch17_p021-025.tex
\begingroup\makeatletter%
\newcommand{\printedpage}[1]{\marginpar{\footnotesize p.~#1}}%
\newcommand{\msp}[1]{\marginpar{\footnotesize [#1]}}%
\newcommand{\M}{\mathcal{M}}%
\newcommand{\Mc}{\mathcal{M}_{\mathrm{c}}}%
\makeatother\let\printedpage\derpp\let\msp\dermsp%
\let\mspin\DERmspin%


\printedpage{19}
% head: A. Grothendieck : Dérivateurs, Ch. XVII, Édition des Universités de Montpellier II et Paris VII

\[
\begin{tikzcd}
X_{I_\alpha} \arrow[r] \arrow[d] & X_{I_{\alpha_0}} \arrow[d] \\
Y_{I_\alpha} \arrow[r] & Y_{I_\alpha,I_{\alpha_0}} \rlap{${}= Z_\alpha$} \arrow[d] \\
& Y_{I_{\alpha_0}} \rlap{\ ,}
\end{tikzcd}
\]
on a $Z_0 = Y_{I_\alpha,I_{\alpha_0}}$, $Z_{\alpha_0} = Y_{I_\alpha,I_{\alpha_0}}$, $Z_\alpha = \varinjlim\limits_{\beta<\alpha} Z_\beta$ si $\alpha$ \msp{34} ordinal limite. Donc il reste à déterminer les $v_{\alpha,\alpha+1} : Z_\alpha \longrightarrow Z_{\alpha+1}$ (pour $\alpha < \alpha_0$), ce qu'on voit sur le diagramme
\[
\begin{tikzcd}
X_{I_\alpha} \arrow[r] \arrow[d] & X_{I_{\alpha+1}} \arrow[r] \arrow[d] & X_{I_{\alpha_0}} \rlap{${}= Z_0$} \arrow[d] \\
Y_{I_\alpha} \arrow[r] & Y_{I_\alpha,I_{\alpha+1}} \arrow[r] \arrow[d, "u_{I_\alpha,I_{\alpha+1}}"] & Y_{I_\alpha,I_{\alpha_0}} \rlap{${}= Z_\alpha$} \arrow[d, "v_{\alpha,\alpha+1}"] \\
& Y_{I_{\alpha+1}} \arrow[r] & Y_{I_{\alpha+1},I_{\alpha_0}} \rlap{${}= Z_{\alpha+1}\ .$}
\end{tikzcd}
\]

Reprenons la démonstration de 2.1. On va prouver 1° par les deux lemmes précédents. Notons que lorsque l'espace sous-jacent à $I$ est noethérien, alors la construction transfinie est forcément finie, donc on n'a pas besoin de la stabilité de $C$ par composés transfinis infinis.

\msp{35}

Prouvons 2°. On suppose donc que pour tout $i_0 \in I$, on ait $u_{I_{<\alpha_0},I/i_0} \in C$, on veut en déduire que $u_\bullet \in C_I$, i.e. $u_{J,I/i_0} \in C$ pour \textit{tout} ouvert.

\noindent\textbf{Lemme 2.8.} (Sous hypothèse artinienne sur $I$.) \textit{Soit $P \subset I$ l'ensemble des $i_0 \in I$ tels que pour $u_\bullet|I/i_0$ soit dans $C_{I/i_0}$, i.e. que pour tout $i \leq i_0$, et tout ouvert $K$ de $I/i$, on ait $u_{K,I/i} \in C$. Soit $i_0 \in I$ tel que $I_{<i_0} \subset P$. Alors $i_0 \in P$.}

\noindent\textsc{Démonstration}. Il reste à prouver que si $K \subset I/i_0 = J$ est un ouvert, alors $u_{K,J} \in C$. On peut supposer $K \neq I/i_0$, i.e. $K \subset I_{<i_0} \overset{\text{déf}}{=} J'$. Mais alors par hypothèse initiale sur $u_\bullet$, on a $u_{J',J} \in C$, et d'autre part, par hypothèse sur $i_0$, $u_\bullet|J'$ est dans $C_{J'}$. Donc par le théorème 2.1. 1°, on a $u_{K,J'} \in C$. Donc par transitivité $u_{K,J} \in C$, q.e.d.

\noindent\textbf{Corollaire 2.9.} \textit{Si $I$ est artinien, alors $P = I$, i.e. $u_\bullet \in C_I$.}

Cela prouve donc 2.1. 2°. Notons que $\mathrm{Esp}(I)$ est noethérien, il en est de même de tout ouvert (tel $J'$ ci-dessus), donc \msp{36} pour la conclusion du lemme, inutile de supposer $C$ stable par composition infinie.

\printedpage{20}
% head: A. Grothendieck : Dérivateurs, Ch. XVII, Édition des Universités de Montpellier II et Paris VII

Prouvons 2.1. 3°. On a a) $\Longrightarrow$ b), car par 2.1. 1°, $u_\bullet \in (TC)_I$ implique que pour $i \in I$, $u_i : X_i \longrightarrow Y_i$ est dans $TC$ (cf. cor. 2.1 a). Donc il reste à prouver b) $\Longrightarrow$ a), i.e. si $u_\bullet \in C \cap W_{I^\circ}$, alors $u_\bullet \in (TC)_I$, i.e. les $u_{K,I/i_0}$ sont dans $TC$. On sait déjà que $u_{K,I/i_0} \in C$ par hypothèse $u_\bullet \in C_I$, donc il faut prouver que c'est dans $W$. On est ramené à prouver que $u_K \in W$, i.e. $u_K \in TC$, car par le diagramme (où $J = I/i_0$)
\[
\begin{tikzcd}
X_K \arrow[r] \arrow[d, "u_K"'] \arrow[dr, phantom, "\text{cocart.}"] & X_J \arrow[d, "\tilde{u}_{K,J}"'] \arrow[dd, bend left=50, "u_J"] \\
Y_K \arrow[r] & Y_{K,J} \arrow[d, "u_{K,J}"'] \\
& Y_J \rlap{\ ,}
\end{tikzcd}
\]
$u_K \in TC$ implique $\tilde{u}_{K,J} \in TC$, et comme $u_J \simeq u_{i_0} \in W$ par hypothèse, il en résulte bien $u_{K,J} \in W$ par saturation. Mais on va s'y prendre autrement :

\msp{37}

\noindent\textbf{Lemme 2.9.} \textit{$I$ ensemble ordonné quelconque, $C$, $W$ comme dans th.} 2.1, 1° \textit{et} 3°, \textit{$u_\bullet : X_\bullet \longrightarrow Y_\bullet$ dans $C_I \cap W_{I^\circ}$. Soit $P$ l'ensemble des $i_0 \in I$ tels que $u_\bullet|I/i_0$ soit dans $(TC)_{I/i_0}$. Si $i_0 \in I$ est tel que $I_{<i_0} \subset P$, alors $i_0 \in P$.}

En effet, il reste à prouver que pour tout ouvert $K$ de $J = I/i_0$, $K \neq J$, on a $u_{K,J} \in TC$. Mais pour ceci, il suffit de voir que $u_K, u_J \in W$ (comme on venait de le voir). Or $u_J \simeq u_{i_0} \in W$ par hypothèse. D'autre part $u_\bullet|K \in (TC_K)$ puisque $K \subset I_{<i_0} \subset P$, donc par le théorème 2.1, 1°, on conclut que $u_K \in TC$.

Si maintenant on suppose $I$ artinien, il résulte de 3.9 que $P = I$, donc $u_\bullet \in (TC)_I$, q.e.d. % [sic] 3.9 -> 2.9, the lemma just stated is Lemme 2.9

Cela achève la démonstration du théorème 2.1, et aussi du corollaire 2.5, vu que pour la démonstration du lemme 2.9 également, les hypothèses infinies sur $C$, $C \cap W$ sont inutiles si $\mathrm{Top}(I)$ est noethérien, donc ses ouverts sont aussi noethériens.

\noindent\textbf{Remarques.}
\begin{enumerate}[leftmargin=2.5em]
\item[1)] Dans 2.1. 2°, l'hypothèse $I$ artinien n'est par superflue, comme on voit en prenant $I = \mathbf{N}^{\mathrm{opp}} \simeq \mathbf{Z}^{\leq 0}$, \msp{38} et pour une flèche % [sic] n'est par superflue -> n'est pas superflue, negation "pas" intended
\[ \emptyset_{\M(I)} \longrightarrow Y_\bullet\ . \]

Dire que cette flèche est dans $C_I$ signifie que les morphismes de transition $Y_i \xrightarrow{u_i} Y_{i+1}$ sont dans $C$, \textit{et} que les $Y_i$ sont $C$-cofibrants, i.e. $\emptyset_{\M} \longrightarrow Y_i$ est dans $C$, pour tout $i$. Mais la condition envisagée dans 2.1. 1° est seulement que les $u_i$ sont dans $C$. Donc en prenant pour $Y_\bullet$ un système inductif constant de valeur $Y$, avec $Y$ non $C$-cofibrant, i.e. $\emptyset_{\M} \longrightarrow Y$ non dans $C$, on a un contre-exemple.

Avec le même $I$, on voit que $u_\bullet : X_\bullet \longrightarrow Y_\bullet$ est dans $C$ si et seulement si
\begin{enumerate}[leftmargin=2.5em]
\item[a)] les $u_i \in C$, et
\item[b)] les $u_{i,i+1} : Y_{i,i+1} \longrightarrow Y_{i+1}$ sont dans $C$.
\end{enumerate}

\printedpage{21}
% head: A. Grothendieck : Dérivateurs, Ch. XVII, Édition des Universités de Montpellier II et Paris VII

Si de plus les $u_i$ sont dans $W$, donc dans $TC$, on voit comme ci-dessus que $u_{i,i+1} \in TC$ pour tout $i$, donc $u_\bullet \in (TC)_I$. Donc ici on a encore
\[ (TC)_I \ = \ C_I \cap W_{I^\circ}\ , \]
quoique $I$ ne soit pas artinien. J'ignore si l'hypothèse artinienne dans 2.1. 3° est superflue ou non.
\item[2)] Soit $u_\bullet \in \mathrm{Fl}(\M(I^\circ))$\footnote[13]{$I^\circ = I^{\mathrm{opp}}$.}, donc $u_\bullet$ est un objet de $\underline{\mathrm{Fl}}(\M(I^\circ)) \simeq \M(I^\circ \times \boldsymbol{\Delta}^1) = \M((I \times (\boldsymbol{\Delta}^1)^{\mathrm{opp}})^{\mathrm{opp}})$. On vérifie que pour que cet objet soit $C$-cofibrant, il faut et il suffit que
\begin{enumerate}[leftmargin=2.5em]
\item[1°)] $X_\bullet$ soit $C$-cofibrant, et
\item[2°)] $u_\bullet \in C_I$ (ce qui implique que $Y_\bullet$ est aussi cofibrant).
\end{enumerate}
\end{enumerate}
\msp{39}

\section*{3\quad Catégories de modèles à cofibrations}

Soit $\M$ une catégorie munie de deux ensembles de flèches
\[ W,\ C\ \subset\ \mathrm{Fl}(\M)\ . \]
On suppose que $\M$ contient un objet initial, et :
\begin{description}[leftmargin=3em, labelwidth=2.5em]
\item[Cof$_0$] $W$ est modérément saturé, i.e. contient les isomorphismes et si $f$, $g$ sont composables, et si deux parmi $f$, $g$, $gf$ sont dans $W$, le troisième aussi.
\item[Cof$_1$] $C$ contient les isomorphismes, est stable par composition et par cochangement de base.
\item[Cof$_2$] $TC \overset{\text{déf}}{=} C \cap W$ est stable par cochangement de base.
\item[Cof$_3$] Toute flèche $f$ dans $\M$ s'écrit
\[ f \ = \ pi \]
avec $i \in C$, $p \in W$.
\end{description}
Ce sont les axiomes de K. S. \textsc{Brown} (forme duale), qui suppose de plus que les objets de $\M$ sont $C$-cofibrants, i.e. $\emptyset_{\M} \longrightarrow X$ est dans $C$. Je ne ferai pas cette hypothèse. La sous-catégorie $\Mc$ de $\M$ formée des \msp{40} objets $C$-cofibrants satisfait aux axiomes de \textsc{Brown}, de plus on prouve que
\[ \mathrm{Ho}_W(\Mc) \xrightarrow{\ \sim\ } \mathrm{Ho}_W(\M) \]
est une équivalence. Ainsi les résultats de \textsc{Brown} relatifs à $\Mc$, muni des ensembles de flèches induits par $W$, $C$, sont applicables à $\mathrm{Ho}_W(\M)$ et à ses avatars (tels $\mathrm{Ho}_W(S\backslash\M)$, $\mathrm{Ho}_W(\M/T)$ et $\mathrm{Ho}_W(S\backslash\underset{u_0}{\M}/T)$, \ où $u_0 : S \longrightarrow T$ est donné). Il se trouve que pour l'essentiel, les résultats homotopiques standard (notamment suite exacte de cofibrations relatifs aux opérations homotopiques “à droite” (ou faut-il dire à gauche ? il s'agit pourtant de questions d'exactitude à droite, non à gauche)) sont valables. Il me faudra revenir dessus de façon plus circonstanciée par la suite. Notamment, il y a une bonne théorie de carrés $W$-cocartésiens (et \textit{dans} $\Mc$\textit{,} il y a assez de $W$-cofibrations\footnote[14]{Plus précisément, par K. S. \textsc{Brown}, \textit{dans} $\Mc$\textit{,} les objets de $C$ sont des $W_{\mathrm{c}}$-cofibrations.}, \ldots),
\printedpage{22}
% head: A. Grothendieck : Dérivateurs, Ch. XVII, Édition des Universités de Montpellier II et Paris VII
\msp{41}
et de plus, ceux-ci sont \textit{exacts.} La catégorie $\mathrm{Ho}_W(\M) \simeq \mathrm{Ho}_W(\Mc)$ “se calcule” par un calcul de fractions à gauche dans $\pi\Mc$, où $\pi\Mc$ est la catégorie déduite de $\Mc$ en y remplaçant les $\mathrm{Hom}_{\Mc}(X,Y) = \mathrm{Hom}_{\M}(X,Y)$ par le quotient par une relation d'homotopie convenable :
\begin{description}[leftmargin=9em, labelwidth=8em, font=\normalfont]
\item[$f,g : X \rightrightarrows Y$] sont QB-homotopes (au sens de \textit{Quillen-Brown}) s'il existe un $X' \xrightarrow{\alpha} X$ dans $W$ (et on peut supposer $X' \in \mathrm{Ob}\,\Mc$ également) tel que $f\alpha \overset{1}{\sim} g\alpha$ au sens de \textsc{Quillen}, i.e. il existe un objet cylindre
\begin{tikzcd}[baseline=(current bounding box.center)]
I_{X'} \arrow[r, "\in W"] & X' \\
X' \vee X' \arrow[u, "\in C"] \arrow[ur, dotted] &
\end{tikzcd}
pour $X'$, tel que $(f\alpha, g\alpha)\ :\ X' \vee X' \longrightarrow Y$ se factorise par un $h' : I_{X'} \longrightarrow Y$.
\end{description}

C'est de ce critère que résulte l'exactitude (à gauche) des carrés $W$-cocartésiens.

Je veux voir comment une structure de cofibrations sur $\M$ s'étend aux catégories $\M(I^\circ) = \underline{\mathrm{Hom}}(I, M)$. On a dans $\M(I^\circ)$ \msp{42} l'ensemble $W_{I^\circ}$, et $C_I$ (il vaut mieux sans doute le noter aussi $C_{I^\circ}$) défini au paragraphe précédent. Il s'agit de vérifier les axiomes. L'axiome $\mathrm{Cof}_0$ est immédiat, et bien connu pour $I$ une catégorie quelconque. % [sic] \underline{Hom}(I, M) -> \underline{Hom}(I, \mathcal{M}), italic M printed where the script \mathcal{M} is used everywhere else

$\mathrm{Cof}_1$ : $C_{I^\circ}$ contient les isomorphismes, c'est clair, parce que $C$ les contient. Stable par composition $X_\bullet \xrightarrow{u} Y_\bullet \xrightarrow{v} Z_\bullet$, regardons $K \subset J$, posons $w = v \circ u$,
\[
\begin{tikzcd}[execute at end picture={\draw[decorate,decoration={brace,mirror,amplitude=6pt}] ([yshift=-2pt]B.south west) -- ([yshift=-2pt]E.south east) node[midway,below=7pt]{$\scriptstyle w_{K,J}$};}]
X_K \arrow[r] \arrow[d] & X_J \arrow[d] & & \\
Y_K \arrow[r] \arrow[d] & Y_{K,J} \arrow[r, "u_{K,J}"] \arrow[d] & Y_J \arrow[d] & \\
Z_K \arrow[r] & |[alias=B]| \bullet \arrow[r, "u'_{K,J}"] & Z_{K,J} \arrow[r, "v_{K,J}"] & |[alias=E]| Z_J \rlap{\ ,}
\end{tikzcd}
\]
\vspace{1ex}

$w_{K,J}$ est composé de $v_{K,J}$, qui est dans $C$, et $u'_{K,J}$, qui l'est car déduit de $u_{K,J}$ par cochangement de base.

(\textbf{N.B.} On a déjà fait le raisonnement avec un seul $u_\bullet$, et \textit{trois} ouverts $J \subset K \subset L$. Il \msp{43} y aurait lieu d'introduire la notion de $C$-carré
\begin{tikzcd}[baseline=(current bounding box.center)]
X_0 \arrow[r] \arrow[d] & X_1 \arrow[d] \\
X_2 \arrow[r] & X_3
\end{tikzcd}, comme un carré tel que
\printedpage{23}
% head: A. Grothendieck : Dérivateurs, Ch. XVII, Édition des Universités de Montpellier II et Paris VII
$X_1 \amalg_{X_0} X_2 \longrightarrow X_3$ soit dans $C$ (\textbf{N.B.} notion symétrique pour les foncteurs de symétrie des carrés), et dire que le composé de deux $C$-carrés est un $C$-carré.)

Cochangement de base
\[
\begin{tikzcd}
X_\bullet \arrow[r, "u"] \arrow[d] & Y_\bullet \arrow[d] \\
X'_\bullet \arrow[r, "u'"] & Y'_\bullet \rlap{\ .}
\end{tikzcd}
\]

Considérons $K \subset J$, on a
\[
\begin{tikzcd}[row sep=large, column sep=large]
& X_K \arrow[r] \arrow[d] \arrow[ddl] \arrow[dr, phantom, "\text{\scriptsize cocart.}"] & X_J \arrow[d] \arrow[dr] & \\
& Y_K \arrow[r] \arrow[ddl] \arrow[dl, phantom, "\text{\scriptsize cocart.}"] & Y_{K,J} \arrow[r, "u_{K,J}"] \arrow[ddl] \arrow[d, phantom, "\text{\scriptsize cocartésien}"] & Y_J \arrow[ddl] \\
X'_K \arrow[r, crossing over] \arrow[d] \arrow[dr, phantom, "\text{\scriptsize cocart.}"] & X'_J \arrow[d] \arrow[dr, crossing over] \arrow[from=1-3, crossing over] & {} & \\
Y'_K \arrow[r] & Y'_{K,J} \arrow[r, "u'_{K,J}"] & Y'_J \rlap{\ ,} &
\end{tikzcd}
\]
où les trois carrés latéraux marqués “cocart.” sont cocartésiens, ainsi que le carré en plan incliné. Il en résulte, d'abord que le quatrième carré latéral du cube est lui aussi cocartésien, \msp{44} puis que le carré ayant deux côtés $u_{K,J}$, $u'_{K,J}$ est cocartésien. (\textbf{N.B.} comparer le diagramme similaire p. 32 -- il y aurait un lemme commun à dégager, distinguant deux cas \ldots)

$\mathrm{Cof}_2$. Pour prouver que $C_{I^\circ} \cap W_{I^\circ}$ est stable par cochangement de base, je ne vois pas de façon de le faire, si je ne sais déjà qu'on a
\[ C_{I^\circ} \cap W_{I^\circ} \ = \ (TC)_{I^\circ}\ , \tag*{$(*)$} \]
et alors on applique le résultat précédent (où $W$ n'entrait pas) à $TC$ au lieu de $C$. Mais pour avoir $(*)$, il faut que je suppose qu'on est dans l'un des deux cas suivants (de façon à pouvoir appliquer le théorème 2.1 et son corollaire 2.5).
\begin{enumerate}[leftmargin=2.5em]
\item[a)] $I$ fini, ou plus généralement $I$ artinien \textit{et} $\mathrm{Top}(I)$ est noethérien.
\item[b)] $I$ artinien, et de plus $C$ et $TC$ sont stables par composition transfinie.
\end{enumerate}
$\mathrm{Cof}_3$. L'axiome de facorisation me paraît a priori le plus délicat\footnote[15]{L'argument qui suit s'applique à d'autres situations de factorisation, par exemple pour un précouple de Quillen donné $(\Phi, \Psi)$ dans $\M$.}. Mais cela sort finalement sans trop de mal. Soit donné % [sic] facorisation -> factorisation, missing letter
\[ f \ : \ X_\bullet \longrightarrow Y_\bullet \]


\endgroup
% END UNIT der-ch17/der-ch17_p021-025.tex
% BEGIN UNIT der-ch17/der-ch17_p026-030.tex
\begingroup\makeatletter%
\newcommand{\printedpage}[1]{\marginpar{\footnotesize p.~#1}}%
\newcommand{\msp}[1]{\marginpar{\footnotesize [#1]}}%
\makeatother\let\printedpage\derpp\let\msp\dermsp%
\let\mspin\DERmspin%


\printedpage{24}
% head: A. Grothendieck : Dérivateurs, Ch. XVII, Édition des Universités de Montpellier II et Paris VII

\msp{45}dans $\mathcal{M}(I^\circ)$, considérons les triples
\[
(J,\ Z,\ i,\ p)\ ,
\]
où $J\subset I$ est un ouvert, $Z$ un objet de $\mathcal{M}(J^\circ)$, et $i$, $p$ sont des flèches
\[
\left\{
\begin{aligned}
& X|J \xrightarrow{\ i\ } Z \xleftarrow{\ p\ } Y|J \\
& \text{avec } i\in C_{J^\circ},\ p\in W_{J^\circ}.
\end{aligned}
\right.
\]
On ordonne les quadruples par prolongement. Si on a un ensemble totalement ordonné, de $(J_\alpha,Z_\alpha,i_\alpha,p_\alpha)$, posant $J=\bigcup I_\alpha$, on trouve de $Z$, $i$, $p$ uniques sur $J$ qui prolongent % [sic] \bigcup I_\alpha -> \bigcup J_\alpha, the union is of the open sets J_\alpha
les $Z_\alpha$, $i_\alpha$, $p_\alpha$ donnés, de plus on aura $i\in C_{J^\circ}$, $p\in W_{J^\circ}$, car l'appartenance à $C_{J^\circ}$ resp. $W_{J^\circ}$ est question locale sur $J$. Donc par Zorn, on est ramené à prouver que si $(J,Z,i,p)$ est maximal, on a $J=I$ – d'ou la factorisation voulue $f=pi$. Donc à prouver que si % [sic] d'ou -> d'où, missing accent
$J\neq I$, le quadruple n'est pas maximal. Considérons le fermé non vide $I\backslash J$. $I$ est supposé artinien, il y a un \textit{élément minimal} $i_0$ de $I\backslash J$. Posons
\msp{46}
\[
J' \;=\; J\cup\{i_0\} \;=\; J\cup I/i_0,
\]
et montrons qu'on peut prolonger le triplet $Z$, $i$, $p$ en $Z'$, $j'$, $p'$ sur $J'$. Posant $K'=I/i_0$, % [sic] j' -> i', the triplet is Z, i, p
$K=J\cap K=I_{<i_0}$, on a donc % [sic] J\cap K -> J\cap K', as in the diagram below
\[
\begin{tikzcd}
K=J\cap K' \arrow[r, hook] \arrow[d, hook] & K'=I/i_0 \arrow[d, hook] \\
J \arrow[r, hook] & J'=J\cup K'.
\end{tikzcd}
\]

On voit alors que les prolongements de $(Z,i,p)$ de $J$ à $J'$ correspondent à ceux de $(Z|K',i|K',p|K')$ à $K$. Pour ceci, on note que pour une flèche $\varphi$ d'objets de $\mathcal{M}(J'^\circ)$, cette flèche est dans $C_{J'^\circ}$ resp. $W_{J'^\circ}$, si et seulement si ses restrictions à $J$ et à $K'$ le sont.

On est donc ramené au cas standard d'un ensemble ordonné avec objet final $e_0$, ensemble ordonné qu'on va noter $I$, et l'ensemble ordonné $I'$ des objets $i<e_0$, i.e. $I'=I\backslash\{i_0\}$. % [sic] I\backslash\{i_0\} -> I\backslash\{e_0\}, I' consists of the i < e_0
On a $f:X_\bullet\longrightarrow Y_\bullet$, flèche dans $\mathcal{M}(I^\circ)$, une factorisation de $f|I'$ en $p'\circ i'$ avec $i'\in C_{I'^\circ}$, $p'\in W_{I^\circ}$, % [sic] W_{I^\circ} -> W_{I'^\circ}, factorisation of f|I' on I'
\msp{47}
et il s'agit de la prolonger en une factorisation similaire de $f$. On note que la donnée d'un objet $X_\bullet$ de $\mathcal{M}(I^\circ)$ équivaut à celle d'un objet $X'_\bullet$ de $\mathcal{M}(I'^\circ)$ (la restriction de $X$ à $I'$) et d'une flèche
\[
X'_{I'} \xrightarrow{\ \alpha_X\ } X_{e_0}=X_0
\]
dans $\mathcal{M}$, où $X'_{I'}=\varinjlim\limits_{i'\in I'}X'_{i'}$. Donc la donnée de la flèche $f:X_\bullet\longrightarrow Y_\bullet$ revient à celle de $f'=f|I':X'_\bullet\longrightarrow Y'_\bullet$, et d'une flèche $f_0:X_0\longrightarrow Y_0$ rendant commutatif le carré
\printedpage{25}
% head: A. Grothendieck : Dérivateurs, Ch. XVII, Édition des Universités de Montpellier II et Paris VII
\[
\begin{tikzcd}
X'_{I'} \arrow[r, "\alpha_X"] \arrow[d, "f'_{I'}"'] & X_0 \arrow[d, "f_0"] \\
Y'_{I'} \arrow[r, "\alpha_Y"] & Y_0\ .
\end{tikzcd}
\]

De même, trouver une factorisation de $f$ en $pi$, qui prolonge une factorisation donnée $f'=p'i'$ sur $J'$, revient à la donnée d'un diagramme commutatif
\[
\begin{tikzcd}[column sep=small, row sep=small]
& X'_{I'} \arrow[rr, "\alpha_X"] \arrow[dl, "i'_{I'}"'] \arrow[dd, "f'_{I'}"] & & X_0 \arrow[dl, "i_0"'] \arrow[dd, "f_0"] \\
Z_{I'} \arrow[rr, "\alpha_Z"'] \arrow[dr, "p'_{I'}"'] & & Z_0 \arrow[dr, "p_0"'] & \\
& Y'_{I'} \arrow[rr, "\alpha_Y"'] & & Y_0\ ,
\end{tikzcd}
\tag*{$(*)$}
\]
\msp{48 vide}
\msp{49}
où l'objet $Z_0$, et les trois flèches $\alpha_Z$, $i_0$, $p_0$ sont à déterminer. Si on a $i'\in C_{I'^\circ}$, $p\in W_{I'^\circ}$, % [sic] p -> p', the given factorisation is f'=p'i'
alors on conclut par théorème 2.1. $1^\circ$ (cor. 2.2) que
\[
i'_{I'}\ \in\ C\ ,
\]
et si on veut que $i\in C_{I^\circ}$, $p\in W_{I^\circ}$, il reste seulement à exiger
\[
p_0\in W, \qquad (i'_{I',I}:\underbrace{Z_{I'}\amalg_{X_{I'}}X_0}_{Z'_0}\longrightarrow Z_0)\in C,
\]
d'après le théorème 2.1. $2^\circ$. On va alors construire le diagramme $(*)$, à partir du diagramme similaire ci-dessous
\[
\begin{tikzcd}
& X_{I'} \arrow[rr, "\alpha_X"] \arrow[dl, "i'_{I'}\in C"'] \arrow[dd, "\text{cocart.}" {pos=0.25}, "f'_{I'}" {pos=0.75}] & & X_0 \arrow[dl, "i'_0\in C"'] \arrow[dd, "f_0" {pos=0.25}] \arrow[dr, dotted, "i_0=j_0i'_0"] & \\
Z_{I'} \arrow[rr, "\alpha'_Z"] \arrow[dr, "p'_{I'}"'] & & Z'_0=Z_{I',I} \arrow[rr, dotted, "j_0\in C"] \arrow[dr, "p'_0"'] & & Z_0 \arrow[dl, dotted, "p_0\in W"] \\
& Y_{I'} \arrow[rr, "\alpha_Y"] & & Y_0 &
\end{tikzcd}
\]
(avec $Z_0$ de $(*)$ remplacé par $Z'_0=Z_{I',I}$) en factorisant la flèche canonique $p'_0:Z'_0\longrightarrow Y_0$ en
\[
p'_0 \;=\; p_0j_0\ , \qquad \text{avec } p_0\in W,\ j_0\in C\ ,
\]
et en posant
\[
i_0 \;=\; j_0i'_0\ ,
\]
ce qui achève la construction.

On a ainsi prouvé le

\printedpage{26}
% head: A. Grothendieck : Dérivateurs, Ch. XVII, Édition des Universités de Montpellier II et Paris VII

\msp{50}\textbf{Théorème 3.1.} \textit{Soit $(\mathcal{M},W,C)$ une catégorie à cofibrations (cf. page \textup{39}), $\mathcal{M}$ stable par $\varinjlim$ finies. Pour tout ensemble odonné $I$, considérons le tuple $(\mathcal{M}(I^\circ),W_{I^\circ},C_{I^\circ})$ correspondant, où $C_{I^\circ}$ est défini au paragraphe \textup{2} (cf. à ce sujet le théorème \textup{2.1}). Supposons que $I$ soit un ensemble ordonné \textup{artinien} (p.ex. $I$ localement fini, i.e. les $I/i$ finis). Supposons de plus qu'on soit dans l'un des deux cas suivants :} % [sic] odonné -> ordonné, missing letter
\begin{itemize}[leftmargin=2.5em]
\item[a)] \textit{L'espace \textup{Esp}$(I)$ est noethérien, i.e. toute suite croissante d'ouverts de $I$ est stationnaire (p.ex. $I$ fini).}
\item[b)] \textit{$C$ et $TC=C\cap W$ sont stables, non seulement par composition ordinaire, mais aussi par composition transfinie. Dans ce cas, on suppose aussi $\mathcal{M}$ stable par petites $\varinjlim$ filtrantes (donc par petites $\varinjlim$ quelconques).}
\end{itemize}
\textit{Sous ces conditions, $(\mathcal{M}(I^\circ),W_{I^\circ},C_{I^\circ})$ est également une catégorie à cofibrations.}

\textbf{Proposition 3.2.} \textit{Soit $u:I\longrightarrow I'$ une application croissante entre ensembles ordonnés, d'où un couple de foncteurs adjoints}
\[
\begin{tikzcd}
\mathcal{M}(I^\circ) \arrow[r, shift left, "u_{!}"] & \mathcal{M}(I'^\circ) \arrow[l, shift left, "u^*"]
\end{tikzcd}\ .
\tag*{(3.2.1)}
\]
\msp{51}
\begin{itemize}[leftmargin=2.5em]
\item[a)] \textit{Le foncteur $u^*$ est compatible avec les localiseurs $W_{I'^\circ}$, $W_{I^\circ}$. Si $I$ est discrète, ou $u$ étale, alors $u^*$ applique $C_{I'^\circ}$ dans $C_{I^\circ}$. Il en est encore de même quand $u$ est localement cofiltrant, i.e. si pour tout $i'\in\operatorname{Ob}I'$, la catégorie $i'\backslash I$ est localement filtrante, i.e. ses composantes connexes sont cofiltrantes (condition $\mathrm{PS}_{\mathrm{g}}$ \textup{1} sur $i'\backslash I$)} \textsuperscript{(}\footnote[16]{inutile de parler de ``faiblement'' localement cofiltrant, car il s'agit d'ensembles ordonnés où il n'y a pas à distinguer entre notions forte et faible.}\textsuperscript{)}.
\item[b)] \textit{Dans tous les cas, $u_!$ transforme $C_{I^\circ}$ en $C_{I'^\circ}$, $TC_{I^\circ}$ en $TC_{I'^\circ}$, enfin $\mathcal{M}(I^\circ)_{\mathrm{c}}$ en $\mathcal{M}(I'^\circ)_{\mathrm{c}}$.}
\end{itemize}

\textsc{Démonstration}. Prouvons a), dans le cas le plus gênant, où on suppose seulement que $u$ est localement faiblement cofiltrant. On applique II 7.4 (p. 92), qui nous dit que $\forall i\in\operatorname{Ob}I$,
\msp{52}
considérant le morphisme induit par $u$
\[
I/i \longrightarrow I'/i' \qquad (\text{où } i'=u(i)),
\]
celui-ci est \textit{(globalement)} faiblement cofiltrant, et que pour tout ouvert $U$ de $I/i$, désignant par $U'$ l'ouvert engendré par son image, le morphisme induit $v:U\longrightarrow U'$ est encore faiblement cofiltrant, et a fortiori \textit{cofinal.} Soit alors $f':X'\longrightarrow Y'$ une flèche dans $C_{I'^\circ}$, prouvons que
\[
\underbrace{u^*(f')}_{f}\ :\ \underbrace{u^*(X')}_{X}\ \longrightarrow\ \underbrace{u^*(Y')}_{Y}
\]
est dans $C_{I^\circ}$, i.e. que pour tout $i\in I$, et $U$ ouvert de $I/i$, dans le diagramme canonique
\printedpage{27}
% head: A. Grothendieck : Dérivateurs, Ch. XVII, Édition des Universités de Montpellier II et Paris VII
\[
\begin{tikzcd}
X_U \arrow[r] \arrow[d] \arrow[dr, phantom, "\text{cocart.}"] & X(i) \arrow[d] \\
Y_U \arrow[r] & Y_{U,i} \arrow[d, "v_{U,i}"] \\
& Y(i)\ ,
\end{tikzcd}
\tag*{$(*)$}
\]
\msp{53}
la flèche $v_{U,i}$ est dans $C$. Le diagramme $(*)$ s'envoie dans le diagramme similaire pour $f':X'\longrightarrow Y'$, $i'=u(i)$, et l'ouvert $U'\subset I'/i'$ engendré par $u(U)$ :
\[
\begin{tikzcd}[column sep=small, row sep=small]
X_U \arrow[rr] \arrow[dr, "f_U"] \arrow[dd, "c_U^{X'}"'] & & X(i) \arrow[dr] \arrow[dd, "\text{cocart.}"' {pos=0.2}, "\mathrm{id}"' {pos=0.75}] & & \\
& Y_U \arrow[rr] \arrow[dd, "c_U^{Y'}"' {pos=0.3}] & & Y_{U,i} \arrow[dr, "v_{U,i}"] \arrow[dd, "c^{f'}_{U,i}"] & \\
X'_{U'} \arrow[rr, "\text{cocart.}"' {pos=0.6}] \arrow[dr, "f'_{U'}"'] & & X'(i') \arrow[dr] & & Y(i) \arrow[dd, "\mathrm{id}"] \\
& Y'_{U'} \arrow[rr] & & Y'_{U',i'} \arrow[dr, "v_{U',i'}"'] & \\
& & & & Y'(i')\ .
\end{tikzcd}
\tag*{$(**)$}
\]

Par cofinalité, les deux flèches verticales $c_U^{X'}$ et $c_U^{Y'}$ sont des isomorphismes, en plus de $X(i)\xrightarrow{\ id\ }X'(i')$, donc la quatrième flèche du cube, $c^{f'}_{U,i}$, est un isomorphisme. Comme $v_{U',i'}$ est dans $C$, $v_{U,i}$ aussi, on gagne.

Prouvons b). Soit donc
\[
f\ :\ X\longrightarrow Y
\]
une flèche dans $C_I$, et soit
\[
f' \;=\; u_!(f)\ :\ u_!(X)=X' \longrightarrow u_!(Y)=Y',
\]
prouvons que c'est dans $C_{I'^\circ}$, donc que
\msp{54}
pour deux ouverts $U'\subset V'$ de $I'$ (où on peut supposer que $V'=I'/i'$), dans le diagramme-base de $(**)$, $v_{U',V'}$ soit dans $C$. Soit $U=u^{-1}(U')$, $V=u^{-1}(V')$, et considérons le % ?? uncertain: hyphen at line end in "diagramme-base" -- kept as a compound hyphen, or "diagrammebase"
diagramme similaire à $(**)$, défini par les deux inclusions d'ouverts $U'\subset V'$, $U\subset V$. Cette fois, comme la formation de $u_!$ commute aux restrictions aux ouverts, par transitivité des $\varinjlim$ , les trois flèches vertivales $c^X_{U'}$, $c^Y_{U'}$, $c^X_{V'}$ sont des isomorphismes, donc aussi $c^f_{U',V'}$. De % [sic] vertivales -> verticales, misprint
même, la flèche $Y_V\longrightarrow Y'_{V'}$ (à la place de $Y(i)\longrightarrow Y'(i')$) est un isomorphisme. Comme $v_{U,V}\in C$, il en est de même de $v_{U',V'}$, ce qui prouve que $f'\in C_{I'^\circ}$. On procède de même pour $TC_I$. Enfin, l'assertion que $u_!$ applique $\mathcal{M}(I^\circ)_{\mathrm{c}}$ dans $\mathcal{M}(I'^\circ)_{\mathrm{c}}$ résulte des définitions et de $u_!(C_{I^\circ})\subset C_{I'^\circ}$.

\textbf{Corollaire 3.3} (de prop. 3.2 et th. 3.1). \textit{Soit $u:I\longrightarrow I'$ comme dans \textup{3.2}. Supposons de plus $I$ artinien, et}
\printedpage{28}
% head: A. Grothendieck : Dérivateurs, Ch. XVII, Édition des Universités de Montpellier II et Paris VII
\msp{55}
\textit{qu'on soit pour $\mathcal{M}$, $W$, $C$, $I$ dans l'une des conditions \textup{a), b)} du théorème \textup{3.1}. Alors $u_!$ induit une flèche}
\[
u_{!\,\mathrm{cof}}\ :\ \mathcal{M}(I^\circ)_{\mathrm{c}} \longrightarrow \mathcal{M}(I'^\circ)_{\mathrm{c}}\ \subset\ \mathcal{M}(I'^\circ)
\]
\textit{qui transforme flèches de $W(I^\circ)$ en flèches de $W(I'^\circ)$ (i.e. compatible avec les localiseurs). Il induit donc un foncteur sur les localisés, d'où un}
\[
u_!^W\ :\ \mathrm{Ho}_W(I^\circ) \longrightarrow \mathrm{Ho}_W(I'^\circ) \qquad \mbox{\textsuperscript{(}\footnotemark[17]\textsuperscript{)}}.
\tag*{(3.3.1)}
\]
\footnotetext[17]{\textbf{N.B.} Cette flèche est définie, même sans supposer $I'$ ordonné, car $u_{!\,\mathrm{cof}}$ transforme encore flèches de $W(I^\circ)$ en flèches de $W(I'^\circ)$. Cf. commentaire p. 86.}

\textsc{Démonstration}. Par le théorème 3.1, on sait que $\mathcal{M}(I^\circ)$, $W(I^\circ)$, $C(I^\circ)$ est une structure de catégorie de modèles à cofibrations. On sait par ailleurs que $TC_{I^\circ}=C_{I^\circ}\cap W(I^\circ)$, par th. 2.1. $3^\circ$ (p. 27). Donc par 3.2, $u_!$ applique $C_{I^\circ}\cap W(I^\circ)$ dans $TC_{I'^\circ}=C_{I'^\circ}\cap W(I'^\circ)$, et a fortiori, en des flèches dans $W(I'^\circ)$. Par ailleurs, par la théorie des catégories à cofibrations, on trouve que toute flèche $f$ entre objets cofibrants est composée $f=pi$, où $i\in C_{I^\circ}$, et où $p$ est dans $W(I^\circ)$, et de plus \textit{inverse à gauche} d'une flèche dans $C_{I^\circ}\cap W(I^\circ)$, i.e. d'une flèche de $(TC)_{I^\circ}$. Quand
\msp{56}
$f\in W(I^\circ)$, alors $i\in W(I^\circ)$, donc $i\in W(I^\circ)\cap C_{I^\circ}=(TC)_{I^\circ}$. Au total, toute flèche $f\in W(I^\circ)_{\mathrm{c}}$ est composée $pi$ dans $\mathcal{M}(I^\circ)_{\mathrm{c}}$, où $i\in(TC)_{I^\circ}$, et $p$ est inverse à gauche d'une $j\in(TC)_{I^\circ}$. Comme $u_!$ transforme $(TC)_{I^\circ}$ en $(TC)_{I'^\circ}$, il transforme $f$ en $f'=p'i'$, avec $i'\in(TC)_{I'^\circ}\subset W_{I'^\circ}$, et $p'$ inverse à gauche d'une $j'=u_!(j)$ qui est également dans $W_{I'^\circ}$. Mais alors $i'\in W(I'^\circ)$, donc aussi $f'=p'i'\in W_{I'^\circ}$, ce qui prouve que $u_!(W(I^\circ)_{\mathrm{c}})\subset W(I'^\circ)$. La dernière assertion provient du lemme suivant, appliqué à $(\mathcal{M}(I^\circ),W(I^\circ),C_{I^\circ})$ au lieu de $(\mathcal{M},W,C)$ :

\textbf{Lemme 3.4.} \textit{Soit $(\mathcal{M},W,C)$ une catégorie à cofibrations. Alors l'inclusion $\mathcal{M}_{\mathrm{c}}\subset\mathcal{M}$ induit une équivalence}
\[
\mathcal{M}_{\mathrm{c}}W_{\mathrm{c}}^{-1} \xrightarrow{\ \sim\ } \mathcal{M}W^{-1}=\mathrm{Ho}_W\ .
\]

Nous donnerons la démonstration (par AQT \texttt{[âne qui trotte]}) ultérieurement.

Cela achève la démonstration de 3.3.

\textbf{3.5.} Il se pose la question si le foncteur $u_!^W$ de (3.3.1) est bien un adjoint à gauche
\msp{57}
du foncteur en sens inverse $u^*_W$, induit par $u^*$. Notons que c'est le cas si $u^*$ transforme $\mathcal{M}(I'^\circ)_{\mathrm{c}}$ en $\mathcal{M}(I^\circ)_{\mathrm{c}}$, et quand on suppose que $I'$ est également artinien et satisfait également à l'une des hypothèses a) b) du théorème 3.1. Car alors $(\mathcal{M}(I^\circ),W(I^\circ),C_{I^\circ})$ est une catégorie à cofibrations par le théorème 3.1, on peut lui appliquer le lemme précédent, donc $u^*_W$ et $u_!^W$ se réalisent à partir d'un couple de foncteurs
\[
\begin{tikzcd}
\mathcal{M}(I^\circ)_{\mathrm{c}} \arrow[r, shift left, "u_{!}^{\mathrm{c}}"] & \mathcal{M}(I'^\circ)_{\mathrm{c}} \arrow[l, shift left, "u_{\mathrm{c}}^*"]
\end{tikzcd}\ ,
\]


\endgroup
% END UNIT der-ch17/der-ch17_p026-030.tex
% BEGIN UNIT der-ch17/der-ch17_p031-035.tex
\begingroup\makeatletter%
\newcommand{\printedpage}[1]{\marginpar{\footnotesize p.~#1}}%
\newcommand{\msp}[1]{\marginpar{\footnotesize [#1]}}%
\newcommand{\Esp}{\operatorname{Esp}}%
\newcommand{\Fl}{\operatorname{Fl}}%
\newcommand{\Ob}{\operatorname{Ob}}%
\newcommand{\Hom}{\operatorname{Hom}}%
\newcommand{\Ho}{\operatorname{Ho}}%
\newcommand{\id}{\mathrm{id}}%
\newcommand{\pr}{\mathrm{pr}}%
\newcommand{\M}{\mathcal{M}}%
\makeatother\let\printedpage\derpp\let\msp\dermsp%
\let\mspin\DERmspin%


\printedpage{29}
% head: A. Grothendieck : Dérivateurs, Ch. XVII, Édition des Universités de Montpellier II et Paris VII
% foot: Version 9 février 2019
lesquels sont visiblement adjoint l'un de l'autre (puisque induits par un couple de foncteurs % [sic] adjoint -> adjoints, accord avec « lesquels »
adjoints $u_!$, $u^*$), et de plus compatibles aux localiseurs. Donc on peut appliquer le lemme des foncteurs adjoints pour les localisés. Pour résumer :

\textbf{Corollaire 3.5.} \textit{Supposons que $I$, $I'$ soient artiniens, et qu'on soit dans l'un des deux cas suivants :}
\begin{enumerate}
\item[a)] \textit{Les espaces} $\Esp(I)$, $\Esp(I')$ \textit{sont noethériens (p.ex. $I$, $I'$ finis).}
\msp{page 58}
\item[b)] \textit{$C$ et $TC = C\cap W$ sont stables par composition transfinie.}
\end{enumerate}
\textit{Supposons de plus que $u^*$ applique $\M(I'^\circ)_{\mathrm{c}}$ dans $\M(I^\circ)_{\mathrm{c}}$, ce qui est le cas notamment si $u$ est localement cofiltrant, i.e. les ensembles ordonnés localisés $i'\backslash I$ sont localement cofiltrants. Alors le foncteur $u_!^W$ construit dans (3.3.1) est adjoint à gauche de $u_W^*$.}

\textbf{Remarque 3.6.} Quand on a un couple de foncteurs adjoints
\[
\begin{tikzcd}
\mathcal{C} \arrow[r, shift left, "u_!"] & \mathcal{C}' \arrow[l, shift left, "u^*"]
\end{tikzcd}
\ ,
\]
avec les localiseurs $W$, $W'$ de $\mathcal{C}$, $\mathcal{C}'$, $u^*$ compatible aux localiseurs, et si de plus on a une sous-catégorie pleine $\mathcal{C}_0$ de $\mathcal{C}$ telle que
\begin{enumerate}
\item[a)] $u_!|\mathcal{C}_0$ transforme flèches de $W$ en flèches de $W'$, et
\item[b)] l'inclusion $\mathcal{C}_0 \hookrightarrow \mathcal{C}$ induit une \textit{équivalence} $\mathcal{C}_0W_0^{-1} \xrightarrow{\sim} \mathcal{C}W^{-1}$ (où $W_0 = W\cap\Fl(\mathcal{C}_0)$),
\end{enumerate}
on trouve deux foncteurs canoniques
\[
\begin{tikzcd}
\mathcal{C}W^{-1} \arrow[r, shift left, "u_!^W"] & \mathcal{C}'W'^{-1} \arrow[l, shift left, "u_W^*"]
\end{tikzcd}
\ ,
\]
où $u_!^W$ est défini seulement à isomorphisme canonique près.

\msp{page 59}
Contrairement à ce que prétend \textsc{Anderson}, il n'est par vrai en général que ces deux % [sic] par -> pas, faute de frappe
foncteurs soient adjoints, sans hypothèses supplémentaires. Voici un contrexemple : on prend $\mathcal{C} = \mathcal{C}'$, $u_! = u^* = \id_{\mathcal{C}}$, $W' = \operatorname{Is}(\mathcal{C})$, $W = \Fl(\mathcal{C})$, donc $u_W^*$ est le foncteur canonique
\[
\underbrace{\Pi\,\mathcal{C}}_{\text{groupoïde fondamental}} \qquad \longleftarrow \mathcal{C}\ .
\]
Supposons que $\mathcal{C}$ soit 1-connexe, i.e. $\Pi\,\mathcal{C}$ soit équivalente à $e$. Soit $e\in\Ob\mathcal{C}$ tel que $\operatorname{End}_{\mathcal{C}}(e)$ soit le groupe unité, et prenons $\mathcal{C}_0 = \{e\}$. Alors l'hypothèse sur $\mathcal{C}_0$ est satisfaite. Mais si $u_!^W$ était adjoint à gauche de $u_W^*$, $e$ serait objet initial de $\mathcal{C}$. Par exemple si $\mathcal{C}$ a un objet final $e$ (ce qui implique que $\mathcal{C}$ est 1-connexe et $e$ rigide), cet objet final devrait être initial si \textsc{Anderson} avait raison avec “it is clear that \ldots”.

\msp{page 60}
Néanmoins :

\textbf{Théorème 3.7.} \textit{Soit $(\M,W,C)$ une catégorie à cofibrations stable par $\varinjlim$ finies, et $u : I \longrightarrow I'$ une application croissante entre ensembles ordonnés artiniens. On suppose qu'on est sous l'une des hypothèses suivantes.}
\begin{enumerate}
\item[a)] $\Esp(I)$, $\Esp(I')$ \textit{sont noethériens, p.ex. $I$, $I'$ finis.}

\printedpage{30}
% head: A. Grothendieck : Dérivateurs, Ch. XVII, Édition des Universités de Montpellier II et Paris VII
% foot: Version 9 février 2019
\item[b)] \textit{$C$ et $TC \overset{\text{déf}}{=} C\cap W$ sont stables par composition transfinie, et $\M$ stable par petites $\varinjlim$ filtrantes.}
\end{enumerate}
\textit{Alors le foncteur} $u_!^W : \Ho_W(I^\circ) \longrightarrow \Ho_W(I'^\circ)$ \textit{construit dans} (3.3.1) \textit{est adjoint à gauche du foncteur $u_W^*$ en sens opposé.}

\textsc{Démonstration.} On utilise le diagramme bien connu
\[
\begin{tikzcd}[row sep=large, column sep=large]
 & I \arrow[d, "\partial_1"'] \arrow[dd, bend left=60, "w"] \\
I \arrow[r, shift left, "\partial_0\ (\text{ouv.})"] \arrow[d, "u"'] \arrow[dr, phantom, "\text{cocart.}"] & I\times\boldsymbol{\Delta}^1 \arrow[l, dotted, shift left, "\pr_1"] \arrow[d] \\
I' \arrow[r, shift left, "\text{imm. ouv.}"] & C(u) \arrow[l, dotted, shift left, "v"]
\end{tikzcd}
\]
pour factoriser $u$ en
\[
u \;=\; vw\ ,
\]
où $w$ est une \textit{immersion fermée,} et $v$ est déduit de $\pr_1 : I\times\boldsymbol{\Delta}^1 \longrightarrow I$ par push-out. C'est une rétraction du mapping-cone C(u) sur l'ouvert $I'$, dont le complémentaire fermé s'identifie à $I$ via $w$.

\msp{page 61}
Le foncteur de recollement entre $I'$, $I$
\[
H \;:\; I'^\circ\times I \;\longrightarrow\; \mathrm{Ens}
\]
est donné par
\[
\underbrace{H(i',i)}_{=\,\Hom_{C(u)}(i',i)} \;=\; \Hom_{I'}(i',u(i))\ . \tag*{$(*)$}
\]
La rétraction $v$ de $C(u)$ dans $I'$ applique $i$ dans $u(i)$, et une flèche $i'\longrightarrow i$ correspondant à $i'\longrightarrow u(i)$ dans $I'$ via $(*)$, est transformé par $v$ en cette même flèche $i'\longrightarrow u(i)$. Calculons, % [sic] transformé -> transformée, accord avec « une flèche »
pour $i'_0$ donné, la catégorie $i'_0\backslash\Bigl(C(u),\ v : C(u)\longrightarrow I'\Bigr)$, des couple d'un objet $\xi$ de $C(u)$ % [sic] des couple -> des couples, accord du pluriel
et d'une flèche $i'_0\longrightarrow v(\xi)$ dans $I'$. On voit que $i'_0\backslash C(u)$ admet l'ouvert $i'_0\backslash I'$, avec le complémentaire fermé $i'_0\backslash I$,
\[
\begin{tikzcd}[column sep=large, row sep=huge]
I' \arrow[r, hook, "\text{ouvert}"'] & C(u) \arrow[d] & I \arrow[l, hook', "\begin{array}{l}\text{fermé}\\ \text{complémentaire}\\ \text{de } I'\end{array}"'] \\
 & I'\ ,
\end{tikzcd}
\]
ces deux catégories étant recollées par le bifoncteur
\[
\begin{aligned}
(i'_0\backslash I')^\circ \times (i'_0\backslash I) &\quad\longrightarrow\quad \mathrm{Ens}\\
\Bigl((\alpha : i'_0\longrightarrow i'), (\beta : i'_0\longrightarrow u(i), i)\Bigr) &\quad\longmapsto\quad \Hom_{i'_0\backslash I'}(i',u(i))
\end{aligned}
\]
\printedpage{31}
% head: A. Grothendieck : Dérivateurs, Ch. XVII, Édition des Universités de Montpellier II et Paris VII
% foot: Version 9 février 2019
induit par $(*)$\,; i.e. les flèches de $(\alpha : i'_0\longrightarrow i') \in \Ob\,(i'_0\backslash I')$ dans $(\beta : i'_0\longrightarrow u(i)) \in i'_0\backslash I$ sont les flèches $i' \xrightarrow{\gamma} u(i)$ telles que le triangle
\[
\begin{tikzcd}
 & i'_0 \arrow[dl, "\alpha"'] \arrow[dr, "\beta"] & \\
i' \arrow[rr, "\gamma"] & & u(i)
\end{tikzcd}
\]
commute.

\msp{page 62}
Si $i' = i'_0$, $\alpha = \id_{i'_0}$, alors cette commutativité revient à $\gamma = \beta$, ce qui signifie qu'il y a une flèche et une seule de $\alpha$ dans $(i,\beta)$. Cela montre donc que $(i'_0,\id_{i'_0})$ est un objet initial de $i'_0\backslash C(u)$. On a donc prouvé le lemme :

\textbf{Lemme 3.8.} \textit{Soit $u : I\longrightarrow I'$ une flèche dans} (Cat)\textit{, factorisons-le en}
\[
I \xrightarrow{\;w\;} C(u) \xrightarrow{\;v\;} I'\ ,
\]
\textit{où $C(u)$ est le mapping-cone défini page} 60\textit{, $w$ une immersion fermée de complémentaire $I'$, $v$ la rétraction canonique de $C(u)$ sur $I'$. Ceci posé, pour tout $i'_0\in\Ob I'$, la catégorie $i'_0\backslash(C(u),v)$ admet un objet initial – donc $v : C(u)\longrightarrow I'$ est (tout ce qu'il y a de) cofiltrante.}

Nous pouvons maintenant prouver 3.7. Tout d'abord, les hypothèses faites sur $I$, $I'$ se transportent à $C(u)$ (hypothèse artinienne, et éventuellement hypothèse noethérienne sur $\Esp(C(u))$). Donc
\msp{page 63}
on peut appliquer les résultats précédents, notamment 3.5, aux flèches $v$, $w$. Le lemme précédent montre bien que $v$ satisfait aux conditions de 3.5, donc le foncteur $v_!^W$ construit dans 3.3.1 est adjoint à gauche de $v_W^*$. Considérons d'autre part le foncteur $w_!^W$ construit dans 3.3.1, je dis que la même conclusion s'applique. Cela provient du fait que $w$ est une immersion fermée. En effet, dans ce cas il est immédiat que
\[
w_! \;:\; \M(I^\circ) \;\longrightarrow\; \M(I'^\circ), \qquad\quad \text{où } I = C(u), % [sic] I = C(u) -> I' = C(u), w a pour but le mapping-cone C(u)
\]
est compatible avec les localiseurs, puisque l'on a
\[
\left\{
\begin{array}{lcll}
w_!(X)(j) & \simeq & X(j) & \text{si } j\in I,\\[1ex]
w_!(X)(j) & \simeq & \emptyset_{\M} & \text{si } j\in J\backslash I,
\end{array}
\right.
\]
par le calcul des fibres
\[
w_!(X)(j) \;\simeq\; \varinjlim_{i\in I/j} X(i)\ .
\]
Donc la définition de $w_!^W$ donnée dans 3.3.1 revient en fait à passer directement aux catégories localisées pour $w_!$, sans avoir à se restreindre à $\M(I^\circ)_{\mathrm{c}}$, et le lemme des foncteurs adjoints nous dit que
\printedpage{32}
% head: A. Grothendieck : Dérivateurs, Ch. XVII, Édition des Universités de Montpellier II et Paris VII
% foot: Version 9 février 2019
\msp{page 64}
ce foncteur est adjoint à gauche de $w_W^*$.

Comme on a évidemment
\[
u_!^W \;\simeq\; v_!^W\circ w_!^W\ , \qquad\quad u_W^* \;\simeq\; w_W^*\circ v_W^*\ ,
\]
on en conclut que $u_!^W$ est bien adjoint à gauche de $u_W^*$, q.e.d.

\textbf{Corollaire 3.9.} \textit{Soit $(\M,W,C)$ une catégorie à cofibrations. Alors le prédérivateur $\mathbf{D}_W$, restreint à la catégorie} $\mathrm{Diag}_0$ \textit{formée des catégories $I$} ordonnées \textit{artiniennes et telles que} $\Esp(I)$ \textit{soit noethérien, satisfait aux axiomes} Dér 1 (décomposition), Dér 2 (localisation), Dér 3 (objet final), Dér 4 bis (existence de $u_!$), Dér 5 bis (mode de calcul des $u_!$), Dér 6 (exactitude des carrés $W$-cocartésiens). \textit{Si $C$, $TC = C\cap W$ sont stables par composition transfinie, on a le même résultat sur la catégorie} Diag \textit{formée des catégories $I$ ordonnées artiniennes sans plus.}

\msp{page 65}
\textbf{3.10.} Ainsi, pour avoir un dérivateur au plein sens du terme, soit sur $\mathrm{Diag}_0$, soit sur Diag, il ne manque que Dér 4 (existence des $u_*$), Dér 5 (calcul des $u_*$) et Dér 6 (exactitude des carrés $W$-cartésiens). On va les obtenir en principe à l'aide d'une structure de \textit{catégorie à fibrations} $(\M,F,W)$ associé à $(\M,W)$, i.e. l'espèce de structure duale de celle introduite % [sic] associé -> associée, accord avec « structure »
page 39. Ces deux structures n'ont pas à être reliées entre elles d'une façon autre que par le fait qu'elles correspondent au même $W$. Elles sont considérablement plus faciles à construire que des structures de Quillen $(W,C,F)$, où $C$, $F$ doivent être reliés de façon extrêmement stricte, avec les propriétés $C \longleftrightarrow \underbrace{TF}_{=\,F\cap W}\ ,\ \underbrace{TC}_{=\,C\cap W} \longleftrightarrow F$ (dont il n'est pas question dans l'axiomatique de K. S. \textsc{Brown}), et factorisation pour $(C,TF)$ et $(TC,F)$.

\msp{page 66}
Mais revenant à la construction des $u_!^W$, il reste le problème de les construire pour une flèche $u : I'\longrightarrow I$ \textit{quelconque} dans Cat. % [sic] u : I' -> I -> u : I -> I', u_!^W va de Ho_W(I°) dans Ho_W(I'°)

À vrai dire, la construction du foncteur
\[
u_!^W \;:\; \Ho_W(I^\circ) \;\longrightarrow\; \Ho_W(I'^\circ) \tag*{(3.3.1)}
\]
est valable sans supposer que $I'$ soit une catégorie ordonnée – ça peut être une catégorie quelconque. Cela provient du fait que le foncteur induit par $u_!$,
\[
u_{!\,\mathrm{cof}} \;:\; \M(I^\circ)_{\mathrm{c}} \;\longrightarrow\; \M(I'^\circ)
\]
est encore compatible avec les localiseurs, comme on voit par le calcul des fibres
\[
u_!(X)(i) \;\simeq\; \varinjlim_{i\in I/i'} X(i)\ . % [sic] u_!(X)(i) -> u_!(X)(i'), la fibre est prise en i'
\]
Ainsi, on est ramené aux cas où $I' = e$, quitte à remplacer $I$ par $I/i'$, et en notant que % [sic] aux cas -> au cas, singulier
$(I/i'\longrightarrow I$ étant étale) l'image inverse de $X$ sur $I/i'$ est cofibrant (cf. 3.2.a). % [sic] cofibrant -> cofibrante, accord avec « l'image inverse »

\msp{page 67}
Mais il se présente à nouveau la perplexité si le foncteur (3.3.1) ainsi construit est bien adjoint à gauche de $u_W^*$. On peut songer à utiliser encore la même factorisation en $u = vw$, en
\printedpage{33}
% head: A. Grothendieck : Dérivateurs, Ch. XVII, Édition des Universités de Montpellier II et Paris VII
% foot: Version 9 février 2019
utilisant le mapping-cone. La difficulté, cependant, c'est qu'à présent $J = C(u)$ n'est plus une catégorie ordonnée, donc on n'a pas une définition toute prête de $C_{J^\circ}\subset\Fl(\M(J^\circ))$, donnant lieu à une sous-catégorie $\M(J^\circ)_{\mathrm{c}}$ etc. En fait, il semblerait qu'on soit pratiquement paralysé, tant qu'on n'arrive à construire sur $\M(J^\circ)$ également une structure de catégorie à cofibrations, se comportant bien par rapport au foncteurs $v_!$. La % [sic] au foncteurs -> aux foncteurs, accord du pluriel
“démonstration” d'\textsc{Anderson} de l'existence de $u_!$ pour \textit{toute} flèche dans Cat, foire totalement (même pour $I\longrightarrow e$, $I$ ensemble ordonné fini – mais là j'ai pu rétablir la situation), car il se ramène au cas où $I$ est ordonné localement fini, et il
\msp{page 68}
admet sans plus que \textit{dans ce cas} $u_!^W$ est adjoint à gauche de $u_W^*$, et \textit{de plus,} que si $u$ est $\mathbf{W}_0$-connexe, donc $u^* : \M(I'^\circ)\longrightarrow\M(I^\circ)$ pleinement fidèle, qu'il en soit de même de $u_W^*$ déduit par passage aux localisés. Or même si on admet que $u_!^W$ est adjoint à gauche de
$u_*^W$, faute de disposer d'un morphisme d'adjonction tout près “en bas” (dans $\Ho_W(I'^\circ)$), % [sic] u_*^W -> u_W^*, notation employée partout ailleurs
déduit du morphisme d'adjonction dans $\M(I^\circ)$, il n'est absolument pas clair que $u_W^*$ soit pleinement fidèle, i.e. $u_!^W$ un foncteur de localisation. Le contrexemple 3.6 montre bien que, même si on a adjonction (i.e. si $e$ est un objet initial de $\mathcal{C} = \mathcal{C}'$), le foncteur $u^* = \id_{\mathcal{C}}$ a beau être pleinement fidèle (et même une équivalence de catégories), le foncteur localisé n'a aucune raison de l'être : cela pouverait que tout objet de Cat ayant un objet initial % [sic] pouverait -> prouverait, faute de frappe
est équivalent à $e$\,!

\msp{page 69}
\textbf{3.11.} Disons qu'on a une \textit{catégorie à cofibrations stricte,} si dans l'axiome $\mathrm{Cof}_3$ de factorisation (cf. p. 39), la factorisation peut être prise fonctoriellement en $f$. Il revient $\pm$ au même de dire que l'axiome de factorisation reste valable, quand on pense à une catégorie $\M(I)$ pour $I$ dans Cat, en y définissant $W(I)$ et $C(I)$ (à ne pas confondre avec $C_I$) fibre par fibre. Cela revient au même tout au moins si $\M$ est équivalente à une petite catégorie. Sinon, il faut supposer fonctorialité pour toute \textit{petite} sous-catégorie pleine de $\underline{\mathrm{Fl}}(\M)$. Il vaut mieux définir la catégorie à cofibrations stricte par cette propriété de stabilité.

On voit alors que le prédérivateur $\mathbf{D} = \mathbf{D}_{\M}$ sur Cat est un dérivateur \textit{relativement à} $\mathcal{D} = \mathrm{Diag}_0$ défini plus haut (catégories ordonnées $I$ artiniennes, telles que $\Esp(I)$ soit noethérien), ou même relativement à $\mathcal{D} = \mathrm{Diag}$ (catégories ordonnées $I$ artiniennes sans plus) lorsque $C$ et $TC$ sont stables par composition transfinie. J'entends par là que $u_!$ existe quand $u$ est une flèche de $\mathcal{D}$,
\msp{page 70}
et que ça reste valable quand on passe à un dérivateur induit $\mathbf{D}_Z$ ($Z$ dans Cat quelconque)\,; en d'autre termes, que $v_!$ existe quand $v$ est de la forme $\id_Z\times u$, où $u : I\longrightarrow I'$ est une % [sic] d'autre termes -> d'autres termes, accord du pluriel
flèche de $\mathcal{D}$, et que le calcul se fait \textit{au dessus de} $I'$ selon la règle Dér 5 bis. Cela signifie donc que la flèche de changement de base pour les $f_!$ est un isomorphismes dans le cas % [sic] un isomorphismes -> un isomorphisme, accord du singulier
suivant :
\[
\begin{tikzcd}[row sep=huge, column sep=huge]
I\times Z \arrow[d, "u\times\id_Z"] & (i'\backslash I)\times Z \arrow[l] \arrow[d, "\pr_2"] \\
I'\times Z & \{i'\}\times Z \arrow[l]
\end{tikzcd}
\ ,
\]
où $i'$ est un objet quelconque de $I'$.


\endgroup
% END UNIT der-ch17/der-ch17_p031-035.tex
% BEGIN UNIT der-ch17/der-ch17_p036-040.tex
\begingroup\makeatletter%
\newcommand{\printedpage}[1]{\marginpar{\footnotesize p.~#1}}%
\newcommand{\msp}[1]{\marginpar{\footnotesize [#1]}}%
\newcommand{\circled}[1]{\tikz[baseline=(c.base)]\node[draw,circle,inner sep=1pt](c){#1};}%
\newcommand{\defeq}{\overset{\text{déf}}{=}}%
\newcommand{\postilde}{\tilde{}}%
\newcommand{\Fl}{\mathrm{Fl}}%
\newcommand{\Cof}{\mathrm{Cof}}%
\makeatother\let\printedpage\derpp\let\msp\dermsp%
\let\mspin\DERmspin%


\printedpage{34}
% head: A. Grothendieck : Dérivateurs, Ch. XVII, Édition des Universités de Montpellier II et Paris VII
% foot: Version 9 février 2019
L'hypothèse \textit{stricte} pour une catégorie à cofibrations me paraît assez anodine, dans la plupart des cas à ma connaissance, elle est satisfaite. Mais même moyennant cette hypothèse, je ne vois aucun moyen de prouver l'existence de $u_!$ pour \textit{toute} flèche de Cat, et c'est une situation bien sûr très gênante. On voudrait absolument dégager des conditions qui assurent la validité de Dér 4 bis et de Dér 5 bis (existence de $u_!$ pour
\msp{page 71}
\textit{toute} flèche de Cat, et mode de calcul des fibres \dots).

\section*{4\quad Catégories à cofibrations spéciales}

\textbf{4.1.} Je suppose à présent que
\begin{itemize}[leftmargin=3em]
\item[\circled{a}] $\mathcal{M}$ est stable par $\varinjlim$ filtrantes, et
\item[\circled{b}] que ses objets sont accessibles,
\end{itemize}
de sorte qu'on est sous les conditions d'application du théorème de factorisation (XV, \S1). On suppose :
\begin{itemize}[leftmargin=3em]
\item[\circled{c}] donné $W, C \subset \Fl(\mathcal{M})$, satisfaisant les conditions $\Cof_1$, $\Cof_2$ de page 39. Pour $\Cof_3$, cf. plus bas.
\item[\circled{d}] $C$, $TC \defeq C \cap W$ \textit{stables par composition transfinie} (c'est le cas si la structure provient d'une structure de Quillen close $(W, C, F)$).
\end{itemize}
\[
TF \defeq C_* \;\subset\; \Fl(\mathcal{M})\ , \tag*{(4.1)}
\]
il résulte du théorème de factorisations que toute flèche $f$ de $\mathcal{M}$ se factorise en
\[
f \;=\; pi\ , \qquad p \in TF,\ i \in C\ ,
\]
\textit{pourvu} qu'on fasse l'hypothèse
\begin{itemize}[leftmargin=3em]
\item[\circled{e}] \mbox{}
\[
\exists\, C_0 \subset C,\ \textit{petite} \text{ partie de } C, \text{ telle que } C \subset \tilde{C}_0, \text{ i.e. } TF = (C_0)_*. \tag*{(4.3)}
\]
\end{itemize}
Sous ces conditions, si on suppose
\begin{itemize}[leftmargin=3em]
\item[\circled{f}] \mbox{}
\[
TF \subset W\ , \tag*{(4.4)}
\]
\end{itemize}
\msp{page 72}
alors l'axiome de factorisations $\Cof_3$ est satisfait. Inversement, si $\Cof_3$ est satisfait, et si $W$ est stable par facteurs directs, alors on a 4.4. Donc si $W$ est stable par facteurs directs, $\Cof_3$ équivaut à (4.4).

\textbf{4.2.} Il me semble que ces conditions, en pratique, ne seront guère jamais remplies, que si $W$ (qui est donné en tout premier lieu, satisfaisant $\Cof_0$ \textsuperscript{(}\footnote[18]{i.e. saturation faible.}\textsuperscript{)}), n'est pas stable par $\varinjlim$
\printedpage{35}
% head: A. Grothendieck : Dérivateurs, Ch. XVII, Édition des Universités de Montpellier II et Paris VII
% foot: Version 9 février 2019
filtrantes. À ce moment, pour construire la structure à cofibrations, on part de $C_0$ petite partie de $\Fl(\mathcal{M})$, prise telle que l'on ait
\[
(C_0)_* \;\subset\; W\ , \tag*{(4.5)}
\]
on prend pour $C$ au choix le saturé $C'_0$ de $C_0$ par cochangement de base et composition transfinie, ou $C = \tilde{C}_0 =$ ensemble des facteurs directs de $C'_0$. Les conditions $\Cof_1$ (stabilités de $C$), $\Cof_3$ (factorisation) sont satisfaites automatiquement, et il reste à s'occuper de $\Cof_2$, qui revient à la stabilité de $C{\cap}W$ par cochangement de base (la stabilité par composition et que ça contienne les isomorphismes, étant immédiates). C'est une condition délicate, et la seule façon standard que je vois pour l'assurer, c'est en supposant
\[
C_0 \;\subset\; \Cof_W\ ; \tag*{(4.6)}
\]
\msp{page 73}
la stabilité de $W$ par $\varinjlim$ filtrantes (donc aussi par facteurs directs) assure alors qu'on a également
\[
C \;\subset\; \Cof_W\ , \tag*{(4.7)}
\]
et par suite
\[
\left\{
\begin{aligned}
& C \cap W \subset \Cof_W \cap W = W^{\mathrm{univ}}, \\
& \text{d'où }\ TC \defeq C \cap W = C \cap W^{\mathrm{univ}}\ ,
\end{aligned}
\right. \tag*{(4.8)}
\]
ce qui assure la stabilité de $C \cap W$ par cochangement de base. La stabilité de $TC$ par composition transfinie résulte de celle de $C$, et de la stabilité de $W$ par composition transfinie.

\textbf{4.3.} Je reviens aux conditions générales\circled{a} à\circled{f} de 4.1, sans supposer nécessairement $W$ stable par $\varinjlim$ filtrantes. Je veux à présent étendre la structure à cofibrations aux $\mathcal{M}(I)$, $I$ quelconque dans Cat. Nous n'avons pas pour ceci (je crois) à utiliser les développements des paragraphes 2, 3, mais nous ferons le lien avec ceux-ci par la suite.

Désignons par $I_0$ la catégorie discrète ayant mêmes objets que $I$, d'où un foncteur canonique
\[
v = v_I \ :\ I_0 \longrightarrow I\ . \tag*{(4.9)}
\]
L'idée est de définir la structure de $\mathcal{M}(I)$ via la structure évidente de $\mathcal{M}(I_0) \simeq \mathcal{M}^{I_0}$, en jouant sur le couple de foncteurs adjoints
\[
\begin{tikzcd}
\mathcal{M}(I_0) \arrow[r, shift left, "\varphi = v_!"] & \mathcal{M}(I) \arrow[l, shift left, "\psi = v^*"]
\end{tikzcd}
\ .
\]
\msp{page 74}
Nous allons présenter les choses de façon un peu plus générale.

\textbf{Définition 4.4.} J'appelle structure de cofibrations \textit{spéciale} une structure $(\mathcal{M}, W, C)$ satisfaisant les conditions a) à f) de 4.1, que je reformule :
\begin{itemize}[leftmargin=3em]
\item[\circled{a}] $\mathcal{M}$ stable par petites limites inductives, et les objets de $\mathcal{M}$ sont accessibles.
\printedpage{36}
% head: A. Grothendieck : Dérivateurs, Ch. XVII, Édition des Universités de Montpellier II et Paris VII
% foot: Version 9 février 2019
\item[\circled{b}] $W$ est faiblement saturé (i.e. $\Cof_0$ de p. 39), $C$ stable par cochangement de base et par composition, $TC \defeq C{\cap}W$ stable par cochangement de base (et par composition automatiquement).
\item[\circled{b$'$}] $C$ et $TC$ sont stables même par composition \textit{transfinie.}
\item[\circled{c}] $\exists$ petite partie $C_0 \subset C$ telle que $C \subset \tilde{C}_0$ ($\Longrightarrow \tilde{C} = \tilde{C}_0$).
\item[\circled{d}] $TF \defeq C_* \subset W$ (ce qui implique, joint à\circled{a} à\circled{c}, l'axiome de factorisation $\Cof_3$ de page 39).
\end{itemize}

\textbf{4.5.} Nous allons donner une recette générale d'extension d'une telle structure, quand on dispose d'un couple de foncteurs adjoints
\[
\begin{tikzcd}
\mathcal{M} \arrow[r, shift left, "v_!"] & \mathcal{M}' \arrow[l, shift left, "v^*"]
\end{tikzcd}
\ ,
\]
$\mathcal{M}$ étant muni d'une structure de cofibrations spéciale $(W, C)$, et $\mathcal{M}'$ satisfaisant la condition liminaire\circled{a}. On définira alors les parties $W'$, $C'$ de $\Fl(\mathcal{M}')$
\msp{page 75}
par
\[
\left\{
\begin{aligned}
W' \;&=\; (v^*)^{-1}(W) \\
C' \;&=\; \text{clôture stable par cochangement de base} \\
&\phantom{\;=\;}\ \ \text{et composition transfinie de } v_!(C),
\end{aligned}
\right. \tag*{(4.5.1)}
\]
donc
\[
v_!(C) \;\subset\; C' \;\subset\; (v_!(C))\postilde \;=\; \tilde{C}'\ . \tag*{(4.5.2)}
\]
Je dis que $(W', C')$, et au choix $(W', \tilde{C}')$, est une structure de cofibrations spéciale sur $\mathcal{M}'$. On va vérifier les différentes conditions de 4.4. La condition a) est déjà acquise.

b) \textsuperscript{(}\footnote[19]{\textbf{N.B.} Il vaut mieux commencer par prouver c).}\textsuperscript{)}. La saturation faible de $W'$ résulte de celle de $W$, via (4.5.1). La stabilité de $C'$ par cochangement de base et par composition (même transfinie) est acquise par définition de $C'$, et itou pour $\tilde{C}'$ (étant Q-close à gauche). Stabilité de $TC' \defeq C'{\cap}W'$ par cochangement de base ? Il suffit de voir que si on a un carré cocartésien
\[
\begin{tikzcd}
X'_0 \arrow[r, "i'"] \arrow[d] \arrow[dr, phantom, "\text{\scriptsize cocart.}"] & X'_1 \arrow[d] \\
X'_2 \arrow[r, "j'"] & X'_3
\end{tikzcd}
\]
dans $\mathcal{M}'$, avec $i'$ dans $TC' = C' \cap W'$, alors $j' \in W'$ (car on sait déjà $j' \in C'$), i.e.
\marginpar{$||$}% marginal sign || printed in the left margin of this line
$v^*(j') \in W'$. Nous allons supposer que $v^*$ \textit{commute aux sommes amalgamées,} et on est donc ramené à prouver que $v^*(i') \in TC$. Donc % [sic] $v^*(j') \in W'$ -> $v^*(j') \in W$, since $W' = (v^*)^{-1}(W)$
\printedpage{37}
% head: A. Grothendieck : Dérivateurs, Ch. XVII, Édition des Universités de Montpellier II et Paris VII
% foot: Version 9 février 2019
\msp{page 76}
on doit prouver le

\textbf{Lemme 4.5.2.} \textit{Sous les conditions} (4.5.3) \textit{et} (4.5.4) \textit{ci-dessous, on a}
\[
v^*(C') \;\subset\; C\ , \qquad v^*(TC') \;\subset\; TC\ ,
\]
\textit{et si $C = \tilde{C}$, on a $v^*(\tilde{C}') \subset C$, $v^*(T\tilde{C}') \subset TC$.}

Pour $v^*(C') \subset C$, il faut d'abord prouver
\[
v^*(v_!(C)) \;\subset\; C\ , \tag*{(4.5.3)}
\]
et je ne vois aucun moyen -- il faut le poser comme hypothèse. De plus, que $v^*$ commute aux $\varinjlim$ quelconques (cochangement de base déjà admis -- à présent on a besoin des limites inductives filtrantes)
\[
v^* \text{ commute aux petites } \varinjlim \text{ quelconques} \tag*{(4.5.4)}
\]
pour conclure de (4.5.3) que $v^*(C' =$ clôture de $v_!(C)$ pour cochangement de base et composition transfinie) est également dans $C$. Passant à $v^*(TC') \subset TC$, par ce qui précède, $v^*(C') \subset C$, il suffit de prouver $v^*(TC') \subset W$, mais $v^*(TC') \subset v^*(W') \subset W$ par définition de $W'$.

Supposons $C = \tilde{C}$, i.e. $C$ stable par facteurs directs, et prouvons $v^*(\tilde{C}') \subset C$, $v^*(T\tilde{C}') \subset TC$. Par c) plus bas, $\tilde{C}'$ est formé des facteurs directs
\msp{page 77}
d'éléments de $C'$, donc $v^*(C') \subset C$ implique $v^*(\tilde{C}') \subset C$ puisque $C$ stable par facteurs directs.

Prouvons enfin $v^*(T\tilde{C}') \subset TC \defeq C \cap W$. Comme $T\tilde{C}' = \tilde{C}' \cap W'$, cela résulte de $v^*(\tilde{C}') \subset C$, et de $v^*(W') \subset W$.

Prouvons b$'$). La stabilité de $C'$ et de $\tilde{C}'$ par composition transfinie est acquise, il reste à prouver la stabilité de $TC' = C' \cap W'$ et de $T\tilde{C}' = \tilde{C}' \cap W'$. Il suffit de voir que le composé transfinie $f$ en question est dans $W'$, donc que $v^*(f) \in W$, mais cela résulte du fait que $v^*$ commute à la composition transfinie par (4.5.4), et que $v^*$ transforme $TC'$ et $T\tilde{C}'$ dans $TC$.

\textit{Prouvons} c) pour $C'$ (donc aussi pour $\tilde{C}'$). Soit $C_0$ une petite partie de $C$ telle que $C \subset \tilde{C}_0$, i.e. $C$ est contenu dans la clôture de $C_0$ par les trois opérations de cochangement de base, de composition transfinie, des facteurs directs. On voit alors que $v_!(C)$ est dans la clôture de $v_!(C_0)$ pour les opérations, donc aussi $C'$, et on peut alors prendre $C'_0 = v_!(C_0)$.

\textit{Prouvons enfin} d), c'est la même assertion pour $C'$ et pour $\tilde{C}'$, puisque
\[
C'_* \;=\; \tilde{C}'_*\ (\ \defeq\ TF'\ )\ ,
\]
\msp{page 78}
il faut donc prouver
\[
TF' \defeq C'_* \;\subset\; W'\ .
\]
Cela va résulter du
\printedpage{38}
% head: A. Grothendieck : Dérivateurs, Ch. XVII, Édition des Universités de Montpellier II et Paris VII
% foot: Version 9 février 2019

\textbf{Lemme 4.5.3.} \textit{On a $TF' = (v^*)^{-1}(TF)$.}

C'est formel par définition de $TF'$ comme
\[
TF' \;=\; (v_!(C))_*\ ,
\]
de $TF$ comme
\[
TF \;=\; C_*\ ,
\]
et de l'adjonction de $v_!$, $v^*$. Comme $TF \subset W$, on en conclut
\[
TF' \;=\; (v^*)^{-1}(TF) \;\subset\; (v^*)^{-1}(W) = W'\ ,
\]
q.e.d.

Donc on a prouvé ceci :

\textbf{Proposition 4.6.} \textit{Considérons un couple de foncteurs adjoints}
\[
\begin{tikzcd}
\mathcal{M} \arrow[r, shift left, "v_!"] & \mathcal{M}' \arrow[l, shift left, "v^*"]
\end{tikzcd}
\ ,
\]
\textit{où $\mathcal{M}$ est muni d'une structure de cofibrations} spéciale $(W, C)$\textit{. On suppose de plus les conditions suivantes satisfaites.}
\[
\begin{aligned}
&1^\circ)\quad \mathcal{M}' \textit{ stable par petites } \varinjlim\ \textit{, ses éléments sont accessibles.} \\
&2^\circ)\quad v^* \textit{ commute aux petites } \varinjlim\ . \\
&3^\circ)\quad v^*(v_!(C)) \subset C.
\end{aligned} \tag*{(4.6.1)}
\]
\msp{page 79}
\textit{Considérons les parties $W'$, $C'$ de} $\Fl(\mathcal{M}')$ \textit{définies par}
\[
\left\{
\begin{aligned}
W' \;&=\; (v^*)^{-1}(W), \\
C' \;&=\; \textit{clôture de } v_!(C) \textit{ par cochangement de base et par composition} \\
&\phantom{\;=\;}\ \textit{transfinie.}
\end{aligned}
\right. \tag*{(4.6.2)}
\]
\textit{Alors $(\mathcal{M}', W', C')$ est une structure à cofibrations spéciale. Si de plus $C = \tilde{C}$, i.e. $C$ est stable par facteurs directs, alors $(\mathcal{M}', W', \tilde{C}')$ est également une structure à cofibrations spéciale} \textsuperscript{(}\footnote[20]{\textbf{N.B.} C'est là un fait général pour les catégories à cofibrations spéciales, et devrait être explicité plus bas (cf. 5.1). L'hypothèse sur $C$ est inutile ici (ou plutôt, il suffit au lieu de cela que $W$ soit stable par facteurs directs \textit{et} par composition transfinie, ce qui n'est pas une hypothèse plus réjouissante \dots).}\textsuperscript{)}.

\textbf{4.7.} J'aimerais, en vue d'énoncés du type 3.7 (établissant l'existence de foncteurs du type $u_!^W$), avoir une prise sur $TC' = C'{\cap}W'$, et sur $T\tilde{C}' = \tilde{C}'{\cap}W'$, de façon précise, j'aimerais avoir la relation
\[
v_!(TC)\ \subset TC' \subset v_!(TC)\postilde\ , \tag*{(4.7.1)}
\]
voire la relation légèrement plus forte
\[
\underbrace{v_!(TC)\ \subset TC'}_{} \ \overset{\text{tautol.}}{\subset}\ \underbrace{T\tilde{C}' \;\subset\; (v_!(TC))\postilde}_{}\ , \tag*{(4.7.2)}
\]


\endgroup
% END UNIT der-ch17/der-ch17_p036-040.tex
% BEGIN UNIT der-ch17/der-ch17_p041-045.tex
\begingroup\makeatletter%
\newcommand{\printedpage}[1]{\marginpar{\footnotesize p.~#1}}%
\newcommand{\msp}[1]{\marginpar{\footnotesize [#1]}}%
\newcommand{\defeq}{\overset{\text{déf}}{=}}%
\newcommand{\fnm}[1]{\textsuperscript{(}\footnotemark[#1]\textsuperscript{)}}%
\newcommand{\cl}{^{\tilde{}}}% tilde printed as exponent after a symbol or a parenthesis
\makeatother\let\printedpage\derpp\let\msp\dermsp%
\let\mspin\DERmspin%


\printedpage{39}
% head: A. Grothendieck : Dérivateurs, Ch. XVII, Édition des Universités de Montpellier II et Paris VII
qui se décompose en deux
\[
\left\{
\begin{array}{rcl}
v_{!}(TC) & \subset & TC' \\
T\tilde{C}' & \subset & (v_{!}(TC))\cl.
\end{array}
\right.
\tag*{(4.7.3)}
\]
La moins délicate des deux est la \msp{page 80} première, $v_{!}(TC) \subset TC'$, qui revient, en vertu des définitions, impliquant
\[
v_{!}(TC) \subset v_{!}(C) \subset C' \, ,
\]
à
\[
v_{!}(TC) \;\subset\; W' \; ,
\]
i.e. par définition de $W'$, à
\[
v^{*}v_{!}(TC) \;\subset\; TC \; ,
\tag*{(4.7.4)}
\]
ce qui est une hypothèse de type (4.6.1) $3^{\circ}$), à joindre aux hypothèses (4.6.1).

La relation hypothétique
\[
\underbrace{T\tilde{C}'}_{=\,\tilde{C}'\cap W'} \;\subset\; (v_{!}(TC))\cl, \quad \text{ou seulement} \quad TC' \subset (v_{!}(TC))\cl,
\tag*{$(*)$}
\]
m'apparaît de nature nettement plus délicate. Comme $T\tilde{C}' \supset TC' \supset v_{!}(TC)$, elle équivaut à l'\textit{égalité} de $(T\tilde{C}')_{*}$ resp. $(TC')_{*}$ et $(v_{!}(TC))_{*} = (v^{*})^{-1}(TC)_{*}$. Posons
\[
\left\{
\begin{array}{lcl}
(TC)_{*} & \defeq & F \\
(TC')_{*} & \defeq & F' \;\supset\; (T\tilde{C}')_{*} \;=\; F'_{\sim} \; .
\end{array}
\right.
\tag*{(4.7.5)}
\]
On a donc, en vertu de l'inclusion
\[
v_{!}(TC) \;\subset\; TC' \; ( \;\subset\; T\tilde{C}' \; ) \; ,
\]
la relation
\[
F'_{\sim} \;\subset\; F' \;\subset\; (v^{*})^{-1}(F) \; ,
\tag*{(4.7.6)}
\]
et la question est donc si on a égalité $F' = (v^{*})^{-1}(F)$, voire $F'_{\sim} = (v^{*})^{-1}(F)$ (d'où $F'_{\sim} = F'$).

\msp{page 81}

En fait, on montre assez facilement qu'on a $F'_{\sim} = F'$, i.e. $T\tilde{C}' \defeq \tilde{C}' \cap W' \subset (C' \cap W')\cl$ -- c'est une histoire générale de catégories à cofibrations spéciales \textsuperscript{(}\footnote[21]{cf. prop. 5.1 ci-dessous.}\textsuperscript{)}. Mais je ne vois guère comment prouver (sous des hypothèses convenables) que $F' = (v^{*})^{-1}(F)$. Il me faudra revenir là-dessus au besoin \textsuperscript{(}\footnote[22]{voir cependant prop. 5.8 ci-dessous.}\textsuperscript{)}. Pour l'instant, je me bornerai à récolter le fruit de 4.6 dans un cas particulier :

\textbf{Proposition 4.8.} \textit{Soit} $(\mathcal{M}, W, C)$ \textit{une catégorie à cofibrations spéciale} (déf. 4.4, page 94)\textit{. Soit $I$ objet de} Cat\textit{, $I_0$ la catégorie discrète associée, considérons $v : I_0 \longrightarrow I$ et} % [sic] page 94 -> page 74, apparently: the same déf. 4.4 is cited as "page 74" on printed p. 42
% foot: Version 9 février 2019
\printedpage{40}
% head: A. Grothendieck : Dérivateurs, Ch. XVII, Édition des Universités de Montpellier II et Paris VII
\[
\begin{tikzcd}
\mathcal{M}(I_0) \arrow[r, shift left, "v_{!}"] & \mathcal{M}(I) \arrow[l, shift left, "v^{*}"]
\end{tikzcd} \; ,
\]
\textit{enfin munissons $\mathcal{M}(I_0) = \mathcal{M}^{I_0}$ de la structure de catégorie à cofibrations,} produit \textit{de celles des facteurs, i.e. définissons $W(I_0)$, $C(I_0)$ “fibre à fibre”. (Cela coïncide avec la définition générale donnée au paragraphe} 2\textit{, pour $I_0$ muni de la structure d'ordre discret, et fait de $\mathcal{M}(I_0)$ une catégorie à cofibrations spéciale.) Alors le couple $(v_{!}, v^{*})$ satisfait} \msp{page 82} \textit{aux conditions de la proposition} 4.6\textit{, ainsi qu'à la condition supplémentaire à} (4.6.1)\textit{,}
\[
4^{\circ}) \qquad v^{*}(v_{!}(TC(I))) \;\subset\; TC(I) \; .
\tag*{(4.8.1)}
\]
\textit{Donc on trouve une structure de catégorie à cofibrations spéciale sur $\mathcal{M}(I)$, en posant}
\[
W(I) \;=\; (v^{*})^{-1}(W(I_0))
\tag*{(4.8.2)}
\]
\textit{(c'est en fait équivalent à la définition habituelle de $W(I)$) et}
\[
C(I) \quad = \quad
\begin{array}[t]{l}
\textit{clôture de } v_{!}(C(I_0)) \textit{ pour les opérations de cochangement de}\\
\textit{base et composition transfinie.}
\end{array}
\tag*{(4.8.3)}
\]
\textit{On a de plus sous ces conditions :}
\[
\underbrace{TF(I)}_{\defeq\,(C(I))_{*}} \;=\; (v^{*})^{-1}( \; \underbrace{TF(I_0)}_{\substack{\defeq\,(C_{I_0})_{*},\\ \textit{se définit aussi}\\ \textit{fibre par fibre}}} \; ) \; ,
\tag*{(4.8.8)}
\]
\[
\left\{
\begin{array}{rcl}
v_{!}( \, \overbrace{(TC)(I_0)}^{C(I_0)\cap W(I_0)} \, ) & \subset & TC(I) \qquad\qquad\qquad \text{\fnm{23}} \\
F(I) \defeq (TC(I))_{*} & \subset & (v^{*})^{-1}(F(I_0)).
\end{array}
\right.
\tag*{(4.8.9)}
\]
\footnotetext[23]{Ces inclusions devraient être des égalités, mais je ne sais pas le prouver sans faire des hypothèses plus fortes, du type \textit{catégorie de Quillen spéciale.}}

\textsc{Démonstration}. On sait déjà que $\mathcal{M}(I)$ est stable par petites $\varinjlim$ , et que tout élément est accessible, comme il résulte aussitôt des propriétés idoines de $\mathcal{M}$. Il est clair aussi que $v^{*}$ commute aux petites $\varinjlim$ . Il reste à prouver (4.6.1. $3^{\circ}$), et $4^{\circ}$ dans (4.8.1), i.e. \msp{page 83} que $v^{*}v_{!} : \mathcal{M}(I_0) \longrightarrow \mathcal{M}(I_0)$ applique $C$ dans $C$, $TC$ dans $TC$. Or on a
\[
v_{!}(F_0)(i) \;=\; \varinjlim_{(\alpha,j)\in i\backslash I_0} F_0(j) \qquad\quad \text{\fnm{24}}
\]
\footnotetext[24]{où $j \in \mathrm{Ob}\, I_0 = \mathrm{Ob}\, I$, et $\alpha : i \longrightarrow v(j) = j$.}
et $i\backslash I_0$ est la catégorie discrète au dessus de $I_0$, dont la fibre en $j \in \mathrm{Ob}\, I_0 \;=\; \mathrm{Ob}\, I$ est $\mathrm{Hom}_I(i,j)$ :
\[
i\backslash I_0 \;\;=\;\; \coprod_{j\in \mathrm{Ob}\, I} \text{(catégorie discrète associée à $\mathrm{Hom}_I(i,j)$)} \; ,
\]
% foot: Version 9 février 2019
\printedpage{41}
% head: A. Grothendieck : Dérivateurs, Ch. XVII, Édition des Universités de Montpellier II et Paris VII
donc
\[
v_{!}(F_0)(i) \;=\; \coprod_{j\in \mathrm{Ob}\, I = \mathrm{Ob}\, I_0} \underbrace{F_0(j)^{\mathrm{Hom}_I(i,j)}}_{\substack{\text{somme directe}\\ \text{de } \mathrm{Hom}_I(i,j)\\ \text{copies de } F_0(j)}} \, ,
\]
donc on trouve
\[
v^{*}v_{!}F_0( \, \underbrace{\; i \;}_{\in \mathrm{Ob}\, I_0 \,=\, \mathrm{Ob}\, I} \, ) \;=\; v_{!}(F_0)(i) \;=\; \text{\tikz[baseline=0pt]\draw[->] (0,0) .. controls (0.4,0) .. (0.5,0.45);} \quad .
\]
% the curved arrow points up to the underbraced expression F_0(j)^{Hom_I(i,j)} of the preceding display
Comme $C$ et $TC$ sont stables par sommes directes quelconques (l'étant par sommes directes finies, à cause des conditions de stabilité finies sur $C$ et $TC$, et aussi par somme infinies à cause de l'axiome de stabilité transfinie \dots\textsuperscript{(}\footnote[25]{mérite un lemme général \dots}\textsuperscript{)}), il s'ensuit bien que $v^{*}v_{!}$ applique $C(I_0)$ et $TC(I_0)$ en eux-mêmes. % [sic] somme infinies -> sommes infinies, agreement with the plural adjective

Ainsi, la première assertion de 4.8 résulte de 4.6. De plus, (4.8.8) et (4.8.9) ont été vus dans 4.7, dans les conditions générales de 4.6 et en supposant de plus le $4^{\circ}$ de (4.7.4).

\msp{page 84}

\section*{5\quad Relation des catégories à cofibrations spéciales avec les catégories de Quillen spéciales}

Soit $(\mathcal{M}, W, C)$ une catégorie à cofibrations spéciale. J'introduis
\[
TF \;=\; C_{*}, \qquad\quad TC \;=\; C \cap W, \qquad\quad F \;=\; (TC)_{*} \; ,
\tag*{(5.1)}
\]
donc on a
\[
\begin{tikzcd}[row sep=large, column sep=large]
TC \arrow[r, hook] \arrow[d, leftrightarrow] & C \arrow[d, leftrightarrow, "\shortmid"{marking, sloped}] \\
F & TF \arrow[l, hook']
\end{tikzcd} \; ,
\tag*{(5.2)}
\]
$(C, TF)$ étant un précouple de Quillen clos à droite. On sait alors (XV 4.2 cor. 1, page 31) que $\tilde{C}$ est formé des facteurs directs d'éléments de $C$. Posons
\[
T\tilde{C} \;=\; \tilde{C} \cap W \; ,
\tag*{(5.3)}
\]
d'où
\[
TC \;\subset\; T\tilde{C} \; .
\tag*{(5.4)}
\]

\textbf{Proposition 5.1} \textit{Soit $(\mathcal{M}, W, C)$ une catégorie à cofibrations spéciale. On a $T\tilde{C} \subset (TC)\cl$ (et même $T\tilde{C} \subset (TC)^{\natural}$ \textsuperscript{(}\footnote[26]{\textbf{N.B.} Pour une partie $A$ de $\mathrm{Fl}(\mathcal{M})$, $A^{\natural}$ désigne l'ensemble des facteurs directs d'éléments de $A$.}\textsuperscript{)}), ou ce qui revient au même en vertu de (5.4), $TC$ et $T\tilde{C}$ ont même $Q$-clôture à gauche :}
\[
T\tilde{C} = (T\tilde{C})\cl = (TC)\cl \qquad\quad \textit{(et même } T\tilde{C} \subset (TC)^{\natural}\textit{),}
\tag*{(5.5)}
\]
% foot: Version 9 février 2019
\printedpage{42}
% head: A. Grothendieck : Dérivateurs, Ch. XVII, Édition des Universités de Montpellier II et Paris VII
\textit{d'où}
\[
(\tilde{C})_{*} = C_{*} \, (= TF), \qquad\quad (T\tilde{C})_{*} = (TC)_{*} \, (= F) \qquad\quad \text{\fnm{27}}.
\tag*{(5.6)}
\]
\footnotetext[27]{Ces relations ne dépendent pas de l'hypothèse $W$ stable par facteurs directs.}
\textit{De plus, si $W$ stable par facteurs directs (ou si seulement tout facteur direct d'une flèche de $TC$ est dans $W$), et si de plus $W$ est stable par composition transfinie, alors $(\mathcal{M}, W, \tilde{C})$ est une catégorie à cofibrations spéciale. De plus, $T\tilde{C}$ est $Q$-saturé à gauche, donc} (par 5.5)
\[
T\tilde{C} \;=\; (TC)^{\natural} \; .
\tag*{(5.7)}
\]

\msp{page 85}

\textsc{Démonstration}. Prouvons
\[
T\tilde{C} \;\subset\; (TC)\cl \, ,
\]
et plus précisément, que tout élément de $T\tilde{C}$ est facteur direct d'un élément de $TC$ :
\[
T\tilde{C} = (TC)^{\natural} \; .
\]
Soit donc $f \in T\tilde{C} \defeq \tilde{C}\cap W$, considérons leur factorisation de $f$ en $pi$, avec $i \in C$, $p \in TF$. Comme $TF \subset W$, on a $p \in W$, donc $i \in W$, donc $i \in C \cap W \defeq TC$. Comme $f \in C$, $p \in TF$, il existe $q : Y \longrightarrow X'$ tel que $pq = \mathrm{id}_Y$, $qf = i$,
\[
\begin{tikzcd}[column sep=small, row sep=small]
X \arrow[rr, "f\in C\cap W"] \arrow[dr, "i\in C"'] & & Y \arrow[dl, dotted, shift right, "q"'] \\
& X' \arrow[ur, shift right, "p\in TF"'] &
\end{tikzcd}
\]
donc $f$ apparaît comme facteur direct de $i \in TC$, q.e.d.

Prouvons que $(\mathcal{M}, W, \tilde{C})$ est une catégorie à cofibrations spéciale (déf. 4.4, page 74). Condition a) sur $\mathcal{M}$ acquise, condition b), la partie sur $W$, aussi. La stabilité de $\tilde{C}$ par cochangement de base et par composition OK, aussi (pour b$'$) la stabilité par composition transfinie. Stabilité de $T\tilde{C} = \tilde{C} \cap W$ par cochangement de base : il reste à prouver que si on a un carré cocartésien
\[
\begin{tikzcd}[row sep=large, column sep=large]
X \arrow[r, "i"] \arrow[d] \arrow[dr, phantom, "\text{cocart.}"] & Y \arrow[d] \\
X' \arrow[r, "i'"] & Y'
\end{tikzcd} \; ,
\]
si $i \in \tilde{C}\cap W$, alors $i' \in W$ (car on sait déjà que $i' \in \tilde{C}$). Mais la démonstration précédente montre que c'est facteur direct \textit{dans} $X\backslash\mathcal{M}$ d'une flèche dans $TC$, donc $i'$ est facteur direct d'une flèche de $TC \subset W$. Donc si $W$ est stable par facteurs directs, $i' \in W$.
% foot: Version 9 février 2019

\printedpage{43}
% head: A. Grothendieck : Dérivateurs, Ch. XVII, Édition des Universités de Montpellier II et Paris VII
\msp{page 86}

Cela achève de prouver b). Pour b$'$), il reste à voir que $T\tilde{C}$ est stable par composition transfinie. Je ne vois pas d'autre façon de le voir, qu'en supposant que $W$ lui-même est stable par composition transfinie (ce qui paraît à peine moins exigeant que la stabilité de $W$ par $\varinjlim$ \texttt{[filtrantes]}). Cela achève de prouver b, b$'$). La condition c) sur $C$ implique la même condition sur $\tilde{C}$. Enfin, comme $\tilde{C}_{*} = C_{*}$, la condition d) est la même pour $C$ et $\tilde{C}$. Cela achève de prouver que $(\mathcal{M}, W, \tilde{C})$ est une catégorie à cofibrations spéciale. (\textbf{N.B.} on avait besoin de deux hypothèses supplémentaires sur $W$, stabilité par facteurs directs et par composition transfinie, pour établir que $T\tilde{C}$ est stable par cochangement de base et par composition transfinie). Le dernière assertion de 5.1 et la relation (5.7) qui en résulte, revient à dire que $T\tilde{C} = \tilde{C} \cap W$ est stable par facteurs directs, ce qui est évident par la même propriété de $\tilde{C}$, et l'hypothèse faite sur $W$ de stabilité par facteurs directs (ou seulement que les facteurs directs d'éléments de $T\tilde{C}$ sont dans $W$, donc dans $T\tilde{C}$). % [sic] Le dernière -> La dernière, gender agreement with "assertion"

\msp{page 87}

\textbf{Proposition 5.2.} \textit{Soit $(\mathcal{M}, W, C)$ une catégorie à cofibrations spéciale. Alors avec les notations de} (5.1)\textit{, on a}
\[
W \;=\; TF \circ TC,
\tag*{(5.8)}
\]
\[
TF \;=\; F \cap W.
\tag*{(5.9)}
\]

\textsc{Démonstration}. Soit $f \in W$, écrivons $f = pi$, avec $i \in C$, $p \in TF$, donc $p \in W$, donc $i \in W$, donc $i \in C \cap W = TC$, d'où (5.8).

Prouvons (5.9). On a $TF \subset F \cap W$, prouvons inversement qu'une $f \in F \cap W$ est dans $TF$. Écrivons $f = pi$, avec $p \in TF$, $i \in TC$.
\[
\begin{tikzcd}[column sep=small, row sep=small]
X \arrow[rr, "f\in F\cap W"] \arrow[dr, "i\in TC"'] & & Y \\
& X' \arrow[ur, "p\in TF"'] &
\end{tikzcd}
\]
Comme $f \in F$, $i \in TC$, et $TC \longleftrightarrow F$ par définition de $F$, $\exists\, r : X' \longrightarrow X$ tel que $ri = \mathrm{id}_X$, $fr = p$. Donc $f$ est facteur direct de $p \in TF$, et comme $TF = C_{*}$ est stable par facteurs directs, on conclut que $f \in TF$.

\textbf{Corollaire 5.3.} \textit{Soit $(\mathcal{M}, W, C)$ une catégorie à cofibrations spéciale. Pour que $(W, C, F)$, $(F \defeq (TC)_{*})$ soit un triple de Quillen dans $\mathcal{M}$, il faut et il suffit que le couple $(TC, F)$ satisfasse à l'axiome de factorisation (i.e. que ce soit un précouple de Quillen). Pour que cette structure de Quillen soit clos, il faut et il suffit que $C$ soit stable par facteurs directs.} % [sic] soit clos -> soit close, agreement with "cette structure"

\msp{page 88}

\textsc{Démonstration}. La nécessité de l'axiome de factorisation pour $(TC, F)$, pour avoir un triple de Quillen, est tautologique. Pour la suffisance, ayant déjà les deux précouples de Quillen $(C, TF)$, $(TC, F)$ avec
% foot: Version 9 février 2019


\endgroup
% END UNIT der-ch17/der-ch17_p041-045.tex
% BEGIN UNIT der-ch17/der-ch17_p046-050.tex
\begingroup\makeatletter%
\newcommand{\printedpage}[1]{\marginpar{\footnotesize p.~#1}}%
\newcommand{\msp}[1]{\marginpar{\footnotesize [#1]}}%
\newcommand{\cl}{^{\tilde{}}}% tilde printed as exponent after a symbol or a parenthesis
\newcommand{\defeq}{\overset{\text{déf}}{=}}%
\newcommand{\encircled}[1]{\tikz[baseline=(X.base)]\node[draw,circle,inner sep=1pt](X){#1};}%
\newcommand{\encoval}[1]{\tikz[baseline=(X.base)]\node[draw,rounded corners=4pt,inner sep=2pt](X){#1};}%
\makeatother\let\printedpage\derpp\let\msp\dermsp%
\let\mspin\DERmspin%


\printedpage{44}
% head: A. Grothendieck : Dérivateurs, Ch. XVII, Édition des Universités de Montpellier II et Paris VII
\[
\begin{tikzcd}
TC \arrow[r, hook] \arrow[d, leftrightarrow] & C \arrow[d, leftrightarrow] \\
F & TF \arrow[l, hook'] \ ,
\end{tikzcd}
\tag*{(5.2)}
\]
et $TC = C \cap W$, le fait que $(W, C, F)$ est un triple de Quillen est le corollaire 4 p. 36 de XV : Les conditions $TC = C \cap W$, $TF = F \cap W$, $W = TF \circ TW$ sont équivalentes, et équivalentes au fait que $(W, C, F)$ soit un triple de Quillen.

Par la définition de Quillen des triples \textit{clos,} par les propriétés
\[
\underbrace{W = TF \circ TC}_{\substack{\text{vraie pour tous}\\ \text{le triples}\\ \text{de Quillen}}}, \qquad C = (TF)^*, \qquad \underbrace{F = (TC)_*}_{\substack{\text{vraie dans}\\ \text{le cas actuel}}}\ , % [sic] le triples -> les triples, agreement with "tous"
\]
il reste à exprimer que $C = (TF)^*$, ce qui équivaut à $C$ stable par facteurs directs (cf. loc. cit. cor. 1, p. 31), q.e.d.

\textbf{Remarque 5.4.} Les énoncés 5.2 et 5.3, et la relation $T\tilde{C} \subset (TC)^\natural$ dans proposition 5.1, ne font pas appel à la totalité de l'hypothèse que $(\mathcal{M}, W, C)$ soit une catégorie à cofibrations spéciale, mais seulement que l'on ait une catégorie à cofibrations, avec $C \mathrel{\leftarrow\!\!\!+\!\!\!\rightarrow} \underbrace{TF}_{\defeq C_*} \subset W$, % ?? uncertain: arrow symbol between C and TF -- double-headed arrow crossed by a small bar/plus sign
\msp{page 89}
i.e. une catégorie avec $W$, $C$ satisfaisant Cof$_0$, Cof$_1$, Cof$_2$ de la page 39, plus $TF \defeq C_* \subset W$, plus l'axiome de factorisation pour $(C, TF)$ (lequel implique Cof$_3$).

\textbf{Corollaire 5.5.} \textit{Soit $(\mathcal{M}, W, C)$ une catégorie à cofibrations spéciale, telle qu'il existe une} petite \textit{partie $TC_0 \subset TC = C \cap W$, avec}
\[
TC \ \subset \ TC_0\cl\ ,
\tag*{(5.5.1)}
\]
\textit{ou ce qui revient au même, $(TC)\cl = (TC_0)\cl$, i.e.}
\[
\underbrace{F}_{\defeq (TC)_*} \ = \ (TC_0)_*\ .
\tag*{(5.5.2)}
\]
\textit{Alors $(W, C, F)$ est un triple de Quillen. Ce triple est clos si et seulement si $C$ est stable par facteurs directs.}

En effet, le théorème de factorisation dans XV §1 implique que moyennant l'hypothèse faite sur $TC_0 \subset TC$, le couple $(TC, F)$ satisfait à la condition de factorisation. En effet, par les conditions b, b$'$ de la définition 4.4, $TC$ a les propriétés de stabilité voulues (contient les isomorphismes, est stable par cochangement de base et par composition transfinie).

\msp{page 90}
\textbf{Définition 5.6.} Les triples de Quillen déduits comme dans 5.5 d'une catégorie à cofibrations spéciale sont dits \textit{spéciaux.} Ainsi un triple de Quillen $(W, C, F)$ dans $\mathcal{M}$ est spécial s'il satisfait les conditions a) à d) dans la définition (4.4), concernant $(W, C)$ exclusivement, plus les deux conditions suivantes.
% foot: Version 9 février 2019

\printedpage{45}
% head: A. Grothendieck : Dérivateurs, Ch. XVII, Édition des Universités de Montpellier II et Paris VII
\begin{itemize}[leftmargin=3em]
\item[\encircled{c$'$}] $\exists$ petite partie $TC_0$ de $TC$, telle que $TC \subset (TC_0)\cl$.
\item[\encircled{e}] $F = (TC)_*$.
\end{itemize}

\textbf{Corollaire 5.7.} \textit{Soit $(W, C, F)$ un triple de Quillen} clos \textit{dans $\mathcal{M}$. Pour que ce triple soit spécial, il faut et il suffit qu'il satisfasse les conditions suivantes.}
\begin{itemize}[leftmargin=3em]
\item[\encircled{a}] \textit{$\mathcal{M}$ stable par petites $\varinjlim$\,, et ses éléments sont accessibles.}
\item[\encoval{c, c$'$}] \textit{Il existe une petite partie $C_0$ de $C$, $TC_0$ de $TC$, telles que l'on ait}
\[
C \ = \ C_0\cl\ , \qquad TC \ = \ (TC_0)\cl\ .
\tag*{(5.7.1)}
\]
\end{itemize}
En effet, toutes les autres conditions a) à e), y inclus c$'$) de 5.6, sont automatiquement satisfaites.

\msp{page 91}
\textbf{Proposition 5.8.} \textit{Considérons un couple de foncteurs adjoints}
\[
\begin{tikzcd}
\mathcal{M} \arrow[r, shift left, "v_!"] & \mathcal{M}' \arrow[l, shift left, "v^*"] \ ,
\end{tikzcd}
\tag*{(5.8.1)}
\]
\textit{où $\mathcal{M}$ est munie d'un triple de Quillen spécial clos $(W, C, F = (TC)_*)$, avec $TC = C \cap W$} (cf. définition 5.6). \textit{On suppose satisfaites les conditions suivantes} (comparer prop. 4.6, relations (4.6.1), et (4.7.4)) :
\[
\left\{
\begin{array}{ll}
\text{a)} & \mathcal{M}' \ \textit{stable par petites}\ \varinjlim\ \textit{, ses éléments sont accessibles.}\\[1.5ex]
\text{b)} & v^* \ \textit{commute aux petites}\ \varinjlim\ .\\[1.5ex]
\text{c)} & v^*(v_!(C)) \subset C,\ v^*(v_!(TC)) \subset TC.
\end{array}
\right.
\tag*{(5.8.2)}
\]
\textit{Définissons $(W', C', F')$ dans $\mathcal{M}'$} (comparer prop. 4.6, notamment 4.6.2) \textit{par}
\[
\left\{
\begin{array}{lcl}
W' & = & (v^*)^{-1}(W)\\[1ex]
C' & = & v_!(C)\cl\\[1ex]
F' & = & (TC')_*\ , \qquad \textit{où}\ \ TC' \ \defeq \ C' \cap W'\ .
\end{array}
\right.
\tag*{(5.8.3)}
\]
\textit{Alors $(W', C', F')$ est un triple de Quillen clos spécial dans $\mathcal{M}'$, et on a de plus}
\[
F' \ = \ (v^*)^{-1}(F)\ , \qquad TF' \ = \ (v^*)^{-1}(TF)\ ,
\tag*{(5.8.4)}
\]
\[
TC' \ = \ (v_!(TC))\cl\ .
\tag*{(5.8.5)}
\]

\msp{page 92}
\textsc{Démonstration}. Par 4.6, on sait déjà que $(\mathcal{M}', W', C')$ est une catégorie à cofibrations spéciale. (\textbf{N.B.} Ce qui est noté $C'$ ici, était noté $\tilde{C}'$ dans la proposition 4.5. On doit utiliser la dernière assertion de 4.6, où on suppose $C$ stable par facteurs directs, OK car $(W, C, F)$ supposé clos.) Pour prouver que $(W', C', F' \defeq (TC')_*)$ est un triple de Quillen
% foot: Version 9 février 2019
\printedpage{46}
% head: A. Grothendieck : Dérivateurs, Ch. XVII, Édition des Universités de Montpellier II et Paris VII
spécial, il reste, par 5.3 et définition 5.6, à prouver qu'il existe une petite partie $TC'_0$ de $TC'$, telle que
\[
(\ TC'_0 \ \subset\ )\ TC' \ \subset \ (TC'_0)\cl\ .
\tag*{(5.8.6)}
\]
(Ce triple de Quillen sera de plus clos, puisque $C'$ est Q-clos à gauche par définition.) Soit donc $TC_0$ une petite partie de $TC$ telle que
\[
TC \ \subset \ (TC_0)\cl\ ,
\tag*{(5.8.7)}
\]
qui existe par hypothèse sur $(W, C, F)$. On prendra
\[
TC'_0 \ \defeq \ v_!(TC_0)\ .
\tag*{(5.8.8)}
\]
On a vu dans 4.7 qu'on a
\[
v_!(TC) \ \subset \ TC'
\tag*{(5.8.9)}
\]
(résulte de l'hypothèse $v^*(v_!(TC)) \subset TC$ dans (5.8.2. c)), a fortiori
\[
TC'_0 \ = \ v_!(TC_0) \ \subset \ TC'\ ,
\]
\msp{page 93}
il reste à prouver
\[
TC' \ \overset{?}{\subset} \ (TC'_0)\cl\ ,
\tag*{(5.8.10)}
\]
ce qui équivaut à $(TC')\cl \subset (TC'_0)\cl$, ou encore à
\[
(TC'_0)_* \ \overset{?}{\subset} \ (TC')_* \ \defeq \ F'\ ,
\tag*{(5.8.10 bis)}
\]
alors qu'on a a priori l'inclusion inverse (déduite de $TC'_0 \subset TC'$)
\[
F' \ \subset \ (TC'_0)_*\ .
\tag*{(5.8.11)}
\]
Or par adjonction, on a
\[
\begin{aligned}
(\ \underbrace{TC'_0}_{v_!(TC_0)}\ )_* \ &= \ (v_!(TC_0))_*\\
&= \ (v^*)^{-1}(\ \underbrace{(TC_0)_*}_{= (TC)_* \defeq F}\ )\\
&= \ (v^*)^{-1}(F)\ ,
\end{aligned}
\]
donc (5.8.11) s'écrit
\[
F' \ \subset \ (v^*)^{-1}(F)\ , \qquad \text{i.e.}\ \ v^*(F') \ \subset \ F
\tag*{(5.8.12)}
\]
(comparer (4.7.6)), et il faut montrer que cette inclusion est une égalité (sur quoi j'avais buté p. 80, 81 \ldots), ce qui n'est autre que la première des deux relations (5.8.4). Mais c'est sous la forme initiale (5.8.10) (transposée de (5.8.4)) que nous allons prouver la chose, finalement. Notons que l'on a
% foot: Version 9 février 2019
\printedpage{47}
% head: A. Grothendieck : Dérivateurs, Ch. XVII, Édition des Universités de Montpellier II et Paris VII
\msp{page 94}
\[
(\underbrace{v_!(TC_0)}_{TC'_0})\cl \ = \ (v_!(TC))\cl\ ,
\]
puisque $TC$ est la clôture de $TC_0$ pour les opérations de cochangement de base, de composition transfinie, et de facteurs directs, et que $v_!$ commute à ces opérations. Donc on doit montrer simplement
\[
TC' \ \subset \ (v_!(TC))\cl
\tag*{(5.8.12 bis)}
\]
(où il n'est plus question de $TC_0$). Mais le théorème de factorisation s'applique au couple
\[
\Big( (v_!(TC))\cl,\ (v_!(TC))_* = (v^*)^{-1}(F) \Big)\ ,
\]
puisque $(v_!(TC))\cl = (TC'_0)\cl$, $TC'_0$ petit. Si donc on a $f' \in TC' = C' \cap W'$, factorisons $f'$ en $p'i'$, avec
\[
i' \ \in \ (v_!(TC))\cl\ , \qquad p' \ \in \ (v^*)^{-1}(F)\ .
\tag*{(5.8.13)}
\]
\[
\begin{tikzcd}[column sep=small, row sep=small]
X' \arrow[rr, "f' \in C' \cap W'"] \arrow[dr, "i'"'] & & Y' \arrow[dl, dotted, shift right, "q'"'] \\
& X'_1 \arrow[ur, shift right, "p'"'] &
\end{tikzcd}
\tag*{$(*)$}
\]
Notons qu'on a
\msp{page 95}
\[
(v_!(TC))\cl \ \subset \ TC' \ (\ \defeq \ C' \cap W'\ )
\tag*{(5.8.14)}
\]
car $v_!(TC) \subset TC' \subset C'$, donc $(v_!(TC))\cl \subset \tilde{C}' = C'$, et il reste à prouver que $v_!(TC)\cl \subset W'$, i.e. $v^*((v_!(TC))\cl) \subset W$. Or je dis qu'on a même
\[
v^*((v_!(TC))\cl) \ \subset \ TC\ ,
\]
i.e.
\[
v_!(TC)\cl \ \subset \ (v^*)^{-1}(TC)\ .
\tag*{(5.8.15)}
\]
Cela résulte du fait que $TC$ est Q-clos, donc Q-saturé à gauche, ce qui implique, comme $v^*$ commute aux petites $\varinjlim$\,, que $(v^*)^{-1}(T)$ l'est également. D'autre part, il contient $v_!(TC)$, donc aussi le Q-saturé à gauche de $v_!(TC)$. Mais ce dernier est \textit{égal} à la Q-clôture $(v_!(TC))\cl$, par le théorème de factorisation, vu que $TC'_0 \subset v_!(TC) \subset (TC'_0)\cl$, avec $TC'_0$ petit. % [sic] (v^*)^{-1}(T) -> (v^*)^{-1}(TC), the set in question is TC

Ainsi, (5.8.13) et (5.8.14) impliquent $i' \in TC' = C' \cap W'$, donc $i' \in W'$, d'où $p' \in W'$, i.e. $v^*(p') \in W$, et comme par hypothèse
% foot: Version 9 février 2019
\printedpage{48}
% head: A. Grothendieck : Dérivateurs, Ch. XVII, Édition des Universités de Montpellier II et Paris VII
\msp{page 96}
(5.8.13) on a $v^*(p') \in F$, on a $v^*(p') \in TF$, i.e.
\[
p' \ \in \ (v^*)^{-1}(TF) \ = \ TF' \ (\ \defeq \ C'\ )\ , % [sic] C' -> C'_*, cf. (5.8.16) TF' ( = C'_* )
\]
l'égalité $(v^*)^{-1}(TF) = TF'$, i.e.
\[
(v^*)^{-1}(C_*) \ = \ (v_!(C))_*\ ,
\]
résultant de l'adjonction de $v_!$, $v^*$. Au total, on a donc
\[
p' \ \in \ TF' \ (\ = \ C'_*\ )\ ,
\tag*{(5.8.16)}
\]
et comme $f' \in C'$, le diagramme $(*)$, page 114, nous montre qu'il existe
\[
q' \ : \ Y' \longrightarrow X'_1\ ,
\]
avec $p'q' = \mathrm{id}$, $q'f' = i'$. Donc $f'$ est facteur direct de $i' \in v_!(TC)\cl$, donc on a $f' \in v_!(TC)\cl$, d'où la relation
\[
TC' \ \subset \ v_!(TC)\cl\ ,
\]
ce qui achève de prouver que $(W', C', F')$ est un triple de Quillen spécial et clos.

Les deux relations (5.8.4) ont été prouvées chemin faisant (la deuxième étant essentiellement triviale, et la première représentant la difficulté principale dans la démonstration). Reste à établir (5.8.5). Mais on sait déjà (5.8.9) que
\[
v_!(TC) \ \subset \ TC'\ ,
\]
\msp{page 97}
d'où
\[
v_!(TC)\cl \ \subset \ TC',
\]
puisque $TC'$ est Q-clos à gauche (ou parce que ça a été prouvé dans (5.8.14)), on a déjà prouvé l'inclusion en sens inverse (5.8.12 bis), assertion équivalente à (5.8.4) (et qui représentait la principale difficulté).

\textbf{Remarque 5.9.} Si on ne suppose pas que la structure sur $\mathcal{M}$ soit close, alors la démonstration donnée de 5.8 montre qu'on trouve pourtant sur $\mathcal{M}'$ une structure de Quillen spéciale, en définissant cette fois $C'$ comme dans 4.6. Mais je doute qu'on aura jamais à travailler avec des triples de Quillen spéciaux qui ne soient aussi clos. De toute façon, si $W$ est stable par petites $\varinjlim$ filtrantes\footnote[28]{ce qui implique la stabilité par facteurs directs, et implique aussi la même propriété pour $W' \defeq (v^*)^{-1}(W)$.} (et je doute qu'on trouve des triples de Quillen spéciaux en dehors de ce cas), alors si $(W, C, F)$ est un triple de Quillen spécial, il en est de même de $(W, \tilde{C} = C^\natural, F)$, et je présume qu'on a tout intérêt à remplacer le triple initial par celui-ci. % footnote mark printed in the text as superscript (28)
% foot: Version 9 février 2019


\endgroup
% END UNIT der-ch17/der-ch17_p046-050.tex
% BEGIN UNIT der-ch17/der-ch17_p051-055.tex
\begingroup\makeatletter%
\newcommand{\printedpage}[1]{\marginpar{\footnotesize p.~#1}}%
\newcommand{\msp}[1]{\marginpar{\footnotesize [#1]}}%
\newcommand{\Dl}{\boldsymbol{\Delta}}%
\newcommand{\tl}{\tilde{}}%
\newcommand{\Ob}{\mathrm{Ob}}%
\newcommand{\Fl}{\mathrm{Fl}}%
\makeatother\let\printedpage\derpp\let\msp\dermsp%
\let\mspin\DERmspin%


\printedpage{49}
% head: A. Grothendieck : Dérivateurs, Ch. XVII, Édition des Universités de Montpellier II et Paris VII
% foot: Version 9 février 2019

\msp{page 98}

\section*{6\quad Application à la construction des dérivateurs (associés aux triples de Quillen spéciaux)}

\textbf{Théorème 6.1.}\footnote[29]{C'est le résultat principal dans cette partie XVII des notes.} \textit{Soit $(W,C,F)$ un triple de Quillen clos spécial sur la catégorie $\mathcal{M}$} (cf. déf. 5.6, p. 90, et cor. 5.7). \textit{Pour tout $I$ dans} Cat\textit{, considérons la catégorie}
\[
\mathcal{M}(I) \;=\; \underline{\mathrm{Hom}}(I^{\circ},\mathcal{M}) \;, \tag*{(6.1.1)}
\]
\textit{et considérons}
\[
W(I),\ F(I) \;\subset\; \Fl(\mathcal{M}(I)) \;,
\]
\textit{définis par}
\[
\begin{aligned}
&\text{\textit{$u : X \longrightarrow Y$ dans $\mathcal{M}(I)$ est dans $W(I)$ (resp. $F(I)$) si et seulement}}\\
&\text{\textit{si $\forall\, i \in \Ob I$, $u(i) : X(i) \longrightarrow Y(i)$ est dans $W$ (resp. $F$).}}
\end{aligned}
\tag*{(6.1.2)}
\]

\begin{itemize}[leftmargin=2.5em]
\item[a)] \textit{Il existe une structure de Quillen close (nécessairement unique) dans $\mathcal{M}(I)$, ayant $W(I)$ et $F(I)$ comme ensembles des quasi-isomorphismes et de fibrations respectivement. On a donc aussi}
\[
\left\{
\begin{aligned}
TF(I) \;&=\; F(I)\cap W(I) \;=\; \{u\in \Fl(\mathcal{M}(I)) \mid u(i)\in TF \ \ \forall\, i\in \Ob I\},\\
C(I) \;&=\; TF(I)^{*}.
\end{aligned}
\right.
\tag*{(6.1.3)}
\]
\textit{Cette structure de Quillen est} spéciale.

\msp{page 99}
\item[b)] \textit{Soit $I_0$ la catégorie discrète ayant même ensemble d'objets que $I$, et}
\[
v_I \ \ \text{\textit{ou}}\ \ v \;:\; I_0 \longrightarrow I
\tag*{(6.1.4)}
\]
\textit{le foncteur canonique. Désignons par}
\[
C(I) \;=\; (TF(I))^{*}, \qquad TC(I) \;=\; C(I)\cap W(I) \;=\; F(I)^{*}
\tag*{(6.1.5)}
\]
\textit{les ensembles de cofibrations resp. de cofibrations triviales dans $\mathcal{M}(I)$, et utilisons les mêmes notations $C(I_0)$, $TC(I_0)$ pour $I_0$. On a alors}
\[
\begin{aligned}
&\text{\textit{$u : X_0 \longrightarrow Y_0$ dans $\mathcal{M}(I_0)$ est dans $C_0$ (resp. $TC_0$) si et seulement si}}\\
&\text{\textit{pour tout $i$ dans $\Ob I$, $u(i) : X_0(i) \longrightarrow Y_0(i)$ est dans $C$ (resp. $TC$).}}
\end{aligned}
\tag*{(6.1.6)}
\]
\textit{Enfin on a}
\[
\left\{
\begin{aligned}
C(I) \;&=\; (v_!(C(I_0)))\tl\\
TC(I) \;&=\; (v_!(TC(I_0)))\tl\,.
\end{aligned}
\right.
\tag*{(6.1.7)} % ?? uncertain: shape of the small superscript mark after the closing parenthesis -- tilde (\tilde{}) or breve (\breve{})
\]

\item[c)] \textit{Soit}
\[
u \;:\; I \longrightarrow I'
\]
\textit{une flèche dans} Cat\textit{, considérons le couple des foncteurs adjoints}
\[
\begin{tikzcd}
\mathcal{M}(I) \arrow[r, shift left, "u_!"] & \mathcal{M}(I') \arrow[l, shift left, "u^{*}"]
\end{tikzcd}
\;.
\tag*{(6.1.8)}
\]

\printedpage{50}
% head: A. Grothendieck : Dérivateurs, Ch. XVII, Édition des Universités de Montpellier II et Paris VII
% foot: Version 9 février 2019

\msp{page 100}
\textit{Ce couple satisfait aux conditions de Quillen pour les couples de foncteurs adjoints entre catégories de Quillen} (Hom. Alg. Chp. I, §4, th. 3, p. 4.5), \textit{de sorte qu'ils définissent par passage aux localisés un couple des foncteurs adjoints}
\[
\begin{tikzcd}
\mathrm{Ho}_W(I) \arrow[r, shift left, "u_!^{W}"] & \mathrm{Ho}_W(I') \arrow[l, shift left, "u^{*}_{W}"]
\end{tikzcd}
\;.
\tag*{(6.1.9)}
\]
\textit{On a, plus précisément, les relations}
\[
\left\{
\begin{aligned}
u_!(C(I)) \;&\subset\; C(I')\\
u_!(TC(I)) \;&\subset\; TC(I')\;,
\end{aligned}
\right.
\tag*{(6.1.10)}
\]
\textit{d'où}
\[
(u_!(\mathcal{M}(I)_{\mathrm{c}})) \;\subset\; \mathcal{M}(I')_{\mathrm{c}}\;,
\tag*{(6.1.11)}
\]
\textit{et de plus}
\[
\begin{aligned}
&\text{\textit{pour toute flèche $f$ dans $\mathcal{M}(I)_{\mathrm{c}}$, telle que $f \in W(I)$, on a}}\\
&\text{\textit{$u_!(f) \in W(I')$, i.e. $(u_!)_{\mathrm{c}} : \mathcal{M}(I)_{\mathrm{c}} \longrightarrow \mathcal{M}(I')_{\mathrm{c}}$ est compatible}}\\
&\text{\textit{avec les localiseurs.}}
\end{aligned}
\tag*{(6.1.12)}
\]
\textit{D'autre part on a}
\[
\left\{
\begin{aligned}
u^{*}(F(I')) \;&\subset\; F(I)\\
u^{*}(TF(I')) \;&\subset\; TF(I)\\
u^{*}(W(I')) \;&\subset\; W(I)\;,
\end{aligned}
\right.
\tag*{(6.1.13)}
\]
\textit{donc}
\[
\left\{
\begin{aligned}
&u^{*}(\mathcal{M}(I')_{\mathrm{f}}) \;\subset\; \mathcal{M}(I)_{\mathrm{f}}\\
&u^{*}_{\mathrm{f}} : \mathcal{M}(I')_{\mathrm{f}} \longrightarrow \mathcal{M}(I)_{\mathrm{f}} \text{ \textit{est compatible avec les localiseurs.}}
\end{aligned}
\right.
\tag*{(6.1.14)}
\]

\msp{page 101}
\item[d)] \textit{Supposons que $I$ soit une catégorie ordonnée, et considérons l'ensemble $C_I \subset \Fl(\mathcal{M}(I))$ défini dans §2}\footnote[30]{Au paragraphe 2, c'était noté $C_{I^{\circ}}$ au lieu de $C_I$.} (cf. page 25, cf. aussi th. 2.1, p. 26, et ses corollaires). \textit{On a alors}
\[
\left\{
\begin{aligned}
C_I \;&\supset\; C(I)\\
TC_I \;&\supset\; TC(I)\;,
\end{aligned}
\right.
\tag*{(6.1.15)}
\]
\textit{et on a égalité quand $I$ est un ensemble ordonné} artinien.
\end{itemize}

\textsc{Démonstration}. L'assertion a) est un cas particulier de 5.8, appliqué au couple de foncteurs adjoints
\[
\begin{tikzcd}
\mathcal{M}(I_0) \arrow[r, shift left, "v_!"] & \mathcal{M}(I) \arrow[l, shift left, "v^{*}"]
\end{tikzcd}
\;.
\]
Les conditions (5.8.2) sont satisfaites, cela a été vu dans (4.8). La proposition 5.8 nous dit alors l'existence d'une structure de triple de Quillen close sur $\mathcal{M}(I)$, dont l'ensemble des quasi-isomorphismes $W(I)$ et l'ensemble des fibrations $F(I)$, sont donnés par la première formule (5.8.3) et par la première formule (5.8.4), ce qui équivaut aux définitions (6.1.2). (\textbf{N.B.} Le fait que $\mathcal{M}(I_0)$, avec la structure $W(I_0)$, $C(I_0)$, $F(I_0)$ évidente (produit), est
\printedpage{51}
% head: A. Grothendieck : Dérivateurs, Ch. XVII, Édition des Universités de Montpellier II et Paris VII
% foot: Version 9 février 2019
\msp{page 102}
bien une catégorie de Quillen spéciale close, est $\pm$ tautologique.)

\textit{Prouvons} b). La relation (6.1.6) est tautologique. Les relations (6.1.7) sont contenues dans 5.8 (cf. deuxième relation dans (5.8.3) pour $C' = v_!(C)\tl$, i.e. $C(I) = (v_!(C(I_0)))\tl$, et (5.8.5) pour $TC(I) = (v_!(TC(I_0)))\tl$).

\textit{Prouvons} c). La première assertion est, par \textsc{Quillen}, un cas particulier des relations (6.1.10) à (6.1.14), c'est elles qu'il s'agit d'établir. Les relations (6.1.13) et par suite (6.1.14) sont tautologiques en vertu des définitions, donc il s'agit d'établir \textit{(6.1.10)} et (6.1.12).

Pour les formules (6.1.10), j'utilise le diagramme commutatif
\[
\begin{tikzcd}[row sep=large, column sep=large]
I \arrow[d, "u"'] & I_0 \arrow[l, "v"'] \arrow[d, "u_0"] \\
I' & I'_0\;, \arrow[l, "v'"']
\end{tikzcd}
\]
et je note que les foncteurs $u_!^{\mathcal{M}}$, $u_{0\,!}^{\mathcal{M}}$, commutant aux petites limites inductives, se
\msp{page 103}
comportent bien pour l'opération
\[
\Phi \longmapsto \Phi\tl\,,
\]
quand $\Phi$ est une partie de $\Fl(\mathcal{M}(I))$ resp. de $\Fl(\mathcal{M}(I_0))$ qui est Q-accessible, i.e. telle qu'il existe une \textit{petite} partie $\Phi_0 \subset \Phi$ avec $\Phi \subset \tilde{\Phi}_0$, plus précisément, on a
\[
u_!(\Phi\tl) \;\subset\; (u_!(\Phi))\tl\,.
\]
Ceci noté, on a par b)
\[
C(I) \;=\; (v_!(C(I_0)))\tl\,,
\]
d'où
\[
u_!(C(I)) \;\subset\; (\underbrace{u_!(v_!(C(I_0)))}_{(uv)_!(C(I_0))})\tl\,,
\]
et le deuxième membre s'écrit aussi
\[
(v'_!u_{0\,!}(C(I_0)))\tl \;\overset{?}{\subset}\; (v'_!(I'_0))\tl \;=\; C(I')\;. % [sic] (v'_!(I'_0)) -> (v'_!(C(I'_0))), the C( ) around I'_0 apparently omitted
\]
Donc tout revient à montrer que
\[
u_{0\,!}(C(I_0)) \;\subset\; C(I'_0)\;,
\tag*{(6.1.16)}
\]
ce qui résulte aussitôt du calcul de $u_{0\,!}$ à coups de sommes directes, et du fait que $C$ est stable par sommes directes quelconques. Cela prouve donc $u_!(C(I)) \subset C(I')$, et on prouve de même $u_!(TC(I)) \subset TC(I')$, qui se ramène de même au cas particulier
\[
u_{0\,!}(TC(I_0)) \;\subset\; TC(I'_0)\;.
\tag*{(6.1.17)}
\]

\printedpage{52}
% head: A. Grothendieck : Dérivateurs, Ch. XVII, Édition des Universités de Montpellier II et Paris VII
% foot: Version 9 février 2019

\msp{page 104}

Il reste à prouver (6.1.12), et on procède comme pour 3.3 (p. 54), en utilisant le fait qu'une flèche $f$ dans $\mathcal{M}(I)_{\mathrm{c}}$ qui est dans $W(I)$, s'écrit dans $\mathcal{M}(I)_{\mathrm{c}}$ sous la forme $f = pi$, où $i \in TC(I)$, et où $p$ est inverse à gauche (rétraction) d'une flèche $j \in TC(I)$. On aura donc $u_!(f) = u_!(p)u_!(i)$, et $u_!(i) \in TC(I')$ par (6.1.10), et $u_!(p)$ est inverse à gauche de $u_!(j) \in TC(I')$. On a alors $u_!(i), u_!(j), u_!(p) \in W(I')$, d'où $u_!(f) \in W(I')$, q.e.d.

Prouvons enfin d). Compte tenu de la définition de $C_I$, pour établir $C(I) \subset C_I$, il suffit de voir ceci :

\textbf{Corollaire 6.2.}\footnote[31]{\textbf{N.B.} Ce résultat se généralise, sous les conditions générales dans 6.2.5.} \textit{Soit $I$ quelconque dans} Cat\textit{, soient $U \subset V$ deux fermés dans $I$ (i.e. des ouverts dans $I^{\circ}$), et soit $f : X_{\bullet} \longrightarrow Y_{\bullet}$ une flèche dans $C(I)$ (resp. dans $TC(I)$). Posons encore}
\msp{page 105}
\[
X_U \;=\; \varinjlim_{i\in U} X(i), \qquad X_V \;=\; \varinjlim_{i\in U} X(i)\;, % [sic] second index i\in U -> i\in V, X_V should be the limit over V
\]
\textit{et de même pour $Y$, et considérons le diagramme}
\[
\begin{tikzcd}[row sep=large, column sep=large]
X_U \arrow[r, "\text{trans.}"] \arrow[d, "u_U"'] \arrow[dr, phantom, "\text{cocart.}"] & X_V \arrow[d] \arrow[dd, dotted, bend left=60, "u_V"] \\
Y_U \arrow[r] \arrow[dr, dotted, "\text{trans.}"'] & Y_{U,V} \arrow[d, "u_{U,V}"'] \\
 & Y_V\;.
\end{tikzcd}
\tag*{(6.2.1)}
\]
\textit{Alors la flèche $u_{U,V}$ est dans $C$ (resp. $TC$).}

\textsc{Démonstration} de 6.2. Admettons que la restriction d'une flèche de $C(I)$ resp. $TC(I)$ à un fermé $V$ de $I$ est dans $C(V)$ resp. $TC(V)$ (cf. \dots{} ci-dessous). Cela nous permet de supposer $V = I$. La donnée de $U$, qui est un ouvert de $I^{\circ}$, équivaut à la donnée d'une flèche $v : I^{\circ} \longrightarrow \Dl^1$ (en prenant la fibre en 0), ou encore à la donnée d'une flèche $u : I \longrightarrow (\Dl^1)^{\circ} \simeq \Dl^1$ (en prenant le fibre en 1). Or % [sic] le fibre -> la fibre, gender (cf. "la fibre en 0" just before)
\[
u \;:\; I \longrightarrow \Dl^1\;, \qquad U \;=\; I_1 \overset{\text{déf}}{=} u^{-1}(\{1\})\;.
\tag*{(6.2.2)}
\]

\msp{page 106}
On a alors
\[
X_U \;=\; u_!(X)(1), \quad X_I \;=\; u_!(X)(0)\;,
\]
et le morphisme de transition $X_U \longrightarrow X_I$ est déduit de la flèche $0 \xrightarrow{\ \partial\ } 1$ dans $\Dl^1$, et de même pour $Y$. Soit
\[
\xi \;=\; u_!(X), \qquad \eta \;=\; u_!(Y)\;,
\]
\printedpage{53}
% head: A. Grothendieck : Dérivateurs, Ch. XVII, Édition des Universités de Montpellier II et Paris VII
% foot: Version 9 février 2019
donc le diagramme (6.2.1) s'identifie au diagramme déduit de
\[
\varphi \;:\; \xi \longrightarrow \eta\;, \qquad \varphi \;=\; u_!(f)\;,
\]
comme
\[
\begin{tikzcd}[row sep=large, column sep=large]
\xi(1) \arrow[r, "\xi(\partial)"] \arrow[d, "\varphi(1)"'] \arrow[dr, phantom, "\text{cocart.}"] & \xi(0) \arrow[d] \arrow[dd, dotted, bend left=60, "\varphi(0)"] \\
\eta(1) \arrow[r] \arrow[dr, dotted, "\eta(0)"'] & \tilde{\eta} \arrow[d, "\tilde{\varphi}"'] \\
 & \eta(0) \qquad .
\end{tikzcd}
\tag*{(6.2.3)}
\]
Or par 6.1.10 appliqué à $u : I \longrightarrow \Dl^1$, la relation $f \in C(I)$ resp. $f \in TC(I)$ implique $\varphi \in C(\Dl^1)$ resp. $\varphi \in TC(\Dl^1)$. Donc pour établir 6.1, on est ramené au cas particulier où $I = \Dl^1$, $U = \{1\}$, $V = I$, donc à prouver le

\textbf{Lemme 6.2.4.} \textit{Soit $\varphi : \xi \longrightarrow \eta$ une flèche dans $\mathcal{M}(\Dl^1)$. Si $\varphi \in C(\Dl^1)$ (resp.}
\msp{page 107}
\textit{$\varphi \in TC(\Dl^1)$), alors la flèche $\tilde{\varphi}$ dans le diagramme} (6.2.3) \textit{est dans $C$ (resp. $TC$).}

Plus généralement, et plus précisément :

\textbf{Lemme 6.2.5.} \textit{Soit $\mathcal{M}$ une catégorie stable par sommes amalgamées, et $\Phi$ une partie Q-close à gauche} (on prendra $\Phi = C$, et $\Phi = TC$, dans le cas du lemme 6.2.4). \textit{Soit $\Psi = \Phi_{*}$, et considérons, pour toute $I$ dans} Cat\textit{, $\Phi(I)$, $\Psi\langle I\rangle$ dans $\mathcal{M}(I)$, définis par :}
\[
\begin{aligned}
&\text{\textit{$u : X \longrightarrow Y$ dans $\mathcal{M}(I)$ est dans $\Psi\langle I\rangle$ si et seulement si $\forall\, i \in \Ob I$,}}\\
&\text{\textit{$u(i) : X(i) \longrightarrow Y(i)$ est dans $\Psi$,}}
\end{aligned}
\tag*{(6.2.5.1)}
\]
\[
\Phi(I) \;=\; \Psi\langle I\rangle^{*}.
\tag*{(6.2.5.2)}
\]

Prenons le cas $I = \Dl^1 = \{0 \xrightarrow{\ \partial\ } 1\}$, et soit $u : X \longrightarrow Y$ une flèche dans $\mathcal{M}(\Dl^1)$, correspondant à un carré commutatif
\[
\begin{tikzcd}[row sep=large, column sep=large]
X(1) \arrow[r, "X(\partial)"] \arrow[d, "u(1)"'] & X(0) \arrow[d, "u(0)"] \\
Y(1) \arrow[r, "Y(\partial)"] & Y(0)\;.
\end{tikzcd}
\tag*{(6.2.5.3)}
\]

Considérons le diagramme correspondant


\endgroup
% END UNIT der-ch17/der-ch17_p051-055.tex
% BEGIN UNIT der-ch17/der-ch17_p056-060.tex
\begingroup\makeatletter%
\newcommand{\printedpage}[1]{\marginpar{\footnotesize p.~#1}}%
\newcommand{\msp}[1]{\marginpar{\footnotesize [#1]}}%
\makeatother\let\printedpage\derpp\let\msp\dermsp%
\let\mspin\DERmspin%


\noindent\printedpage{54}\msp{page 108}
% head: A. Grothendieck : Dérivateurs, Ch. XVII, Édition des Universités de Montpellier II et Paris VII
\[
\begin{tikzcd}
X(1) \arrow[r, "X(\partial)"] \arrow[d, "u(1)"'] \arrow[dr, phantom, "\text{cocart.}"] & X(0) \arrow[d] \arrow[dd, dotted, bend left=40, "u(0)"] \\
Y(1) \arrow[r] \arrow[dr, dotted, "Y(\partial)"'] & \tilde{Y} \arrow[d, "\tilde{u}"'] \\
 & Y(0)
\end{tikzcd}
\qquad .
\tag*{(6.2.5.4)}
\]

Pour qu'on ait $u \in \Phi(\boldsymbol{\Delta}^1)$, il faut et il suffit que l'on ait
\[
u(0),\ u(1),\ \tilde{u} \ \in\ \Phi\ .
\tag*{(6.2.5.5)}
\]

J'ai déjà vérifié la chose (sous forme duale), mais ne retrouve plus l'endroit (peut-être n'ai-je pas pris des notes). Je n'ai pas envie de rechercher la démonstration ici.

Pour achever de prouver 6.1. d), i.e. l'égalité $C(I) = C_I$ pour $I$ artinien, j'ai envie de prouver encore quelque chose de plus générale que voici : % [sic] plus générale -> plus général, agreement with « quelque chose »

\textbf{Corollaire 6.3.} \textit{Soient $(\mathcal{M}, \Phi, \Psi)$ et $I \longmapsto \Phi(I) \subset \mathrm{Fl}(\mathcal{M}(I))$ comme dans \textup{6.2.5.} Soit $I$ un ensemble ordonné, et considérons aussi $\Phi_I$ défini (dans la situation duale) page \textup{25 (2.2)} (mais avec des notations duales \dots). Je dis que l'on a}
\[
\Phi_I \ \supset\ \Phi(I)\ ,
\tag*{(6.3.1)}
\]
\msp{page 109}\textit{et on a égalité si $I$ est noethérien, i.e. $I^{\circ}$ est artinien.}

\textsc{Démonstration}. La démonstration de $\Phi(I) \subset \Phi_I$ est celle qu'on vient de donner. Il reste à prouver que si $I$ est noethérien, on a $\Phi_I \subset \Phi(I)$. Étant donnée $i : A \longrightarrow B$ dans $\Phi_I$, il faut prouver que $u \in \Phi(I)$, % [sic] u -> i, the arrow in question is i : A -> B
i.e. que pour un diagramme commutatif dans $\mathcal{M}(I)$,
\[
\begin{tikzcd}
A \arrow[r, "\alpha"] \arrow[d, "i"'] & X \arrow[d, "p"] \\
B \arrow[r, "\beta"'] \arrow[ur, dotted, "h"] & Y
\end{tikzcd}
\tag*{(6.3.2)}
\]
avec $p \in \Psi\langle I\rangle$, i.e. $p(i) : X(i) \longrightarrow Y(i)$ dans $\Psi$ pour tout $i$, $\exists\ h : B \longrightarrow X$ laissant le diagramme commutatif. On va revenir aux notations duales, en travaillant dans $\mathcal{M}(I^{\circ})$, avec $I$ artinien. On considère l'ensemble $\mathcal{H}$ des couple % [sic] des couple -> des couples, plural
$(U, h_U)$ d'un ouvert $U$ de $I$ et d'un $h_U : B|U \longrightarrow X|U$, rendant commutatif la restriction à $U$ de (6.3.2). % [sic] commutatif -> commutative, agreement with « la restriction »
On ordonne $\mathcal{H}$ par prolongement, on note que toute partie totalement ordonnée $(U_\alpha, h_{U_\alpha})$ dans $\mathcal{H}$ a une borne supérieure, savoir $U = \bigcup_\alpha U_\alpha$, $h_U$ défini par les $h_{U_\alpha}$ sur les $U_\alpha$.
% foot: Version 9 février 2019

\printedpage{55}\msp{page 110}
% head: A. Grothendieck : Dérivateurs, Ch. XVII, Édition des Universités de Montpellier II et Paris VII
(\textbf{N.B.} L'ensemble $\mathcal{H}$ n'est pas vide, car $(\emptyset, h_{\emptyset})$ est dans $\mathcal{H}$ -- mais c'est contenu dans l'énoncé précédent que toute partie totalement ordonnée de $\mathcal{H}$ y a une borne supérieure \dots). Par Zorn, il existe un élément maximal $(U, h_U)$. À prouver que $U = I$, et on aura terminé. Ou encore que si $(U, h_U)$ tel que $U \neq I$, ce n'est \textit{pas} maximal. Soit $F = I\backslash U \neq \emptyset$ le fermé complémentaire. Par l'hypothèse artinienne sur $I$, $F$ a un élément minimal $i_0$. Soit
\[
U' \ =\ U \cup \overbrace{I/i_0}^{V} \ =\ U \cup \{i_0\}\ ,
\]
on va prolonger $h_U$ en $h_{U'}$ sur $U'$, on aura terminé. Pour ceci, soit $V_0 = V \cap U = V\backslash\{i_0\}$, on est ramené à prolonger la restriction $h_{V_0}$ de $h_U$ à $V_0$, de $V_0$ à $V$. Au total, on peut supposer $I = I/i_0$ muni d'un élément final $i_0$, $U = I\backslash\{i_0\}$, et montrer qu'une $(U, h_U) \in \mathcal{H}$ se prolonge à $I$. Soient $A'$, $B'$, $X'$, $Y'$ les restrictions de $A$, $B$, $X$, $Y$ à $U$ (qu'on pourra aussi noter $I'$). La donné de $A$ % [sic] La donné -> La donnée, missing feminine ending
équivaut à celle de $A'$, celle de $A_0 = A(i_0) \in \mathrm{Ob}\,\mathcal{M}$, et d'une flèche de
\[
A'_{I'} \ \overset{\text{déf}}{=}\ \varinjlim\limits_{i' \in I'} A'(i') \ \xrightarrow{\ \lambda_A\ }\ A_0\ (\ \simeq A_I\ )\ ,
\]
\msp{page 111}et de même pour $B$, $X$, $Y$. Plus précisément, la catégorie $\mathcal{M}(I^{\circ})$ s'interprète comme celle des triples $X = (X', X_0, \lambda_X)$, où $X' \in \mathcal{M}(I'^{\circ})$, $X_0 \in \mathcal{M}$, et $\lambda_X : X'_{I'} \longrightarrow X_0$. À ce compte, la donnée du diagramme (6.3.2) équivaut à celle d'un diagramme carré similaire
\[
\begin{tikzcd}
A' \arrow[r, "\alpha'"] \arrow[d, "i'"'] & X' \arrow[d, "p'"] \\
B' \arrow[r, "\beta'"'] \arrow[ur, dotted, "h'"] & Y'
\end{tikzcd}
\tag*{(6.3.3)}
\]
\textit{dans} $\mathcal{M}(I'^{\circ})$, d'où un diagramme dans $\mathcal{M}$
\[
\begin{tikzcd}
A'_{I'} \arrow[r, "\alpha'_{I'}"] \arrow[d, "i'_{I'}"'] & X'_{I'} \arrow[d, "p'_{I'}"] \\
B'_{I'} \arrow[r, "\beta'_{I'}"'] \arrow[ur, dotted, "h'_{I'}"] & Y'_{I'}\ ,
\end{tikzcd}
\tag*{(6.3.4)}
\]
d'un autre diagramme carré \textit{dans} $\mathcal{M}$
\[
\begin{tikzcd}
A_0 \arrow[r, "\alpha_0"] \arrow[d, "i_0"'] & X_0 \arrow[d, "p_0"] \\
B_0 \arrow[r, "\beta_0"'] & Y_0\ ,
\end{tikzcd}
\tag*{(6.3.5)}
\]
et enfin d'un homomorphisme $\lambda = (\lambda_A, \lambda_B, \lambda_X, \lambda_Y)$ du carré (6.3.4) dans le carré (6.3.5) :
% foot: Version 9 février 2019
\printedpage{56}\msp{page 112}
% head: A. Grothendieck : Dérivateurs, Ch. XVII, Édition des Universités de Montpellier II et Paris VII
\[
\begin{tikzcd}
A'_{I'} \arrow[rr, "\alpha'_{I'}"] \arrow[dr, "i'_{I'}"'] \arrow[dd, "\lambda_A"'] & & X'_{I'} \arrow[dr, "p'_{I'}"] \arrow[dd, "\lambda_X" near start] & \\
 & B'_{I'} \arrow[ur, dotted, "h'_{I'}"] \arrow[rr, crossing over, "\beta'_{I'}" near end] & & Y'_{I'} \arrow[dd, "\lambda_Y"] \\
A_0 \arrow[rr, "\alpha_0" near start] \arrow[dr, "i_0"'] & & X_0 \arrow[dr, "p_0"] & \\
 & B_0 \arrow[from=2-2, crossing over, "\lambda_B"' near start] \arrow[ur, dotted, "h_0"'] \arrow[rr, "\beta_0"] & & Y_0
\end{tikzcd}
\qquad\qquad \text{\footnotemark[32]}
\tag*{(6.3.6)}
\]
\footnotetext[32]{\textbf{N.B.} La condition $p \in \Psi\langle I\rangle$ s'exprime par $p' \in \Psi\langle I'\rangle$ et \fbox{$p_0 \in \Psi.$}}

D'autre part, la question de prolonger $(I', h') \in \mathcal{H}$ en $(I, h)$, revient à celle de trouver $h_0 : B_0 \longrightarrow X_0$ dans $\mathcal{M}$, de façon qu'avec $h'_{I'}$ cela fasse encore un diagramme commutatif de (6.3.6), ce qui revient à dire que d'une part les deux triangles formés dans le carré de base sont commutatifs :
\[
p_0h_0 \ =\ \beta_0, \qquad\quad h_0i_0 \ =\ \alpha_0\ ,
\tag*{(6.3.7)}
\]
et d'autre part le carré diagonal défini par $h'_{I'}$, $h_0$ soit commutatif (ce qui exprime qu'on a une flèche $h$ de $B$ dans $X$) :
\[
\lambda_X h'_{X'} \ =\ h_0\lambda_B\ . % [sic] h'_{X'} -> h'_{I'}, the arrow is h'_{I'} everywhere else (cf. (6.3.16))
\tag*{(6.3.8)}
\]

\msp{page 113}Mais le diagramme (6.3.6), où la flèche en traits pointillés denses $h'_{I'}$ est déjà acquise, et $h_0$ en pointillés espacés reste à trouver, donne naissance au diagramme, dans lequel est mis en évidence la flèche $\tilde{i} \in \Phi$ qui exprime l'hypothèse $i \in C(i)$, % [sic] C(i) -> C(I), cf. i' \in C(I')
compte tenu de $i' \in C(I')$ :
\[
\begin{tikzcd}
A'_{I'} \arrow[rr, "\alpha'_{I'}"] \arrow[dr, "i'_{I'}"'] \arrow[dd, "\lambda_A"'] & & X'_{I'} \arrow[dr, "p'_{I'}"] \arrow[dd, "\lambda_X" near start] & \\
 & B'_{I'} \arrow[ur, dotted, "h'_{I'}"] \arrow[rr, crossing over, "\beta'_{I'}" near end] & & Y'_{I'} \arrow[dd, "\lambda_Y"] \\
A_0 \arrow[ur, phantom, "\text{cocart.}"] \arrow[rr, "\alpha_0" near start] \arrow[dr, "r"'] \arrow[ddr, dotted, "i_0"'] & & X_0 \arrow[dr, "p_0\in\Psi"] & \\
 & \tilde{B} \arrow[from=2-2, crossing over, "s"' near start] \arrow[ur, dotted, "\tilde{h}_0"] \arrow[rr, dotted, "\tilde{\beta}_0"] \arrow[d, "\tilde{i}\in\Phi"'] & & Y_0 \\
 & B_0\ , \arrow[uur, dotted, "h_0"'] \arrow[urr, "\beta_0"'] & &
\end{tikzcd}
\tag*{(6.3.9)}
\]
% foot: Version 9 février 2019
\printedpage{57}
% head: A. Grothendieck : Dérivateurs, Ch. XVII, Édition des Universités de Montpellier II et Paris VII
où la flèche $\tilde{h}_0$ en traits pointillés denses de $\tilde{B} \overset{\text{déf}}{=} A_0 \amalg_{A'_{I'}} B'_{I'}$ dans $X_0$, est définie par les conditions de commutativité
\[
\tilde{h}_0r \ =\ \alpha_0, \qquad\quad \tilde{h}_0s \ =\ \lambda_X h'_{I'}\ ,
\tag*{(6.3.10)}
\]
ce qui est licite, vu que
\[
\begin{aligned}
\alpha_0\lambda_A \ &=\ (\lambda_X h'_{I'})i'_{I'} \\
&=\ \lambda_X(h'_{I'}i'_{I'}) \\
&=\ \lambda_X\alpha'_{I'}.
\end{aligned}
\]

\msp{page 114}Quant à la flèche $\tilde{\beta}_0 : \tilde{B} \longrightarrow Y_0$, elle est donnée par la condition de commutativité
\[
\tilde{\beta}_0r \ =\ p_0\alpha_0, \qquad\quad \tilde{\beta}_0s \ =\ \lambda_Y\beta'_{I'}\ ,
\tag*{(6.3.11)}
\]
ce qui est encore licite, puisque
\[
(p_0\alpha_0)\lambda_A \ =\ (\lambda_Y\beta'_{I'})i'_{I'}
\]
en vertu de la commutativité du cube (6.3.6). Pour établir la commutativité totale du diagramme (6.3.9) (traits pleins, et traits pointillés denses pour $h'_{I'}$, $\tilde{h}_0$ et $\tilde{\beta}_0$ -- à l'exclusion de $h_0$ pas encore construit), il faut encore vérifier d'une part les relation % [sic] les relation -> les relations, plural
\[
\left\{
\begin{aligned}
\tilde{h}_0r \ &=\ \alpha_0, \quad p_0\tilde{h}_0 \ =\ \tilde{\beta}_0, \\
\lambda_X h'_{I'} \ &=\ \tilde{h}_0s
\end{aligned}
\right.
\qquad\qquad \text{\footnotemark[33]},
\tag*{(6.3.12)}
\]
\footnotetext[33]{\textbf{N.B.} la première et la troisième de ces relations sont déjà connues, cf. (6.3.10).}
établissant la commutativité du cube supérieur avec le carré diagonal (diagramme du type (6.3.6), la flèche pointillée $h_0$ y incluse), d'autre part les relations
\[
\left\{
\begin{aligned}
\tilde{i}r \ &=\ i_0 \qquad \text{(et $\tilde{i}s = \lambda_B$ pour mémoire)} \\
\beta_0\tilde{i} \ &=\ \tilde{\beta}_0\ ,
\end{aligned}
\right.
\tag*{(6.3.13)}
\]
où la première ligne est déjà connue par définition de $\tilde{B}$ et de $\tilde{i}$. La deuxième
\msp{page 115}se ramène aux deux relations
\[
\begin{aligned}
(\beta_0\tilde{i})r \ &=\ \beta_0(\tilde{i}r) \\
&=\ \beta_0i_0 \\
&=\ p_0\alpha_0 \\
&=\ \tilde{\beta}_0r\ ,
\end{aligned}
\]
OK par commutativité du carré de base de (6.3.5),
\[
\begin{aligned}
(\beta_0\tilde{i})s \ &=\ \beta_0(\tilde{i}s) \\
&=\ \beta_0\lambda_B \\
&=\ \lambda_Y\beta'_{I'} \\
&=\ \tilde{\beta}_0s\ ,
\end{aligned}
\]
% foot: Version 9 février 2019
\printedpage{58}
% head: A. Grothendieck : Dérivateurs, Ch. XVII, Édition des Universités de Montpellier II et Paris VII
OK par commutativité du carré antérieur dans (6.3.6).

Donc il reste à vérifier la deuxième relation (6.3.12),
\[
p_0\tilde{h}_0 \ =\ \tilde{\beta}_0\ ,
\]
ce qui se ramène à
\[
\begin{aligned}
(p_0\tilde{h}_0)r \ &=\ p_0(\tilde{h}_0r) \\
&=\ p_0\alpha_0 \\
&=\ \tilde{\beta}_0r\ ,
\end{aligned}
\]
OK par (6.3.11),
\[
\begin{aligned}
(p_0\tilde{h}_0)s \ &=\ p_0(\tilde{h}_0s) \\
&=\ p_0(\lambda_X h'_{I'}) \\
&=\ (p_0\lambda_X)h'_{I'} \\
&=\ (\lambda_Y p'_{I'})h'_{I'} \\
&=\ \lambda_Y(p'_{I'}h'_{I'}) \\
&=\ \lambda_Y\beta'_{I'} \\
&=\ \tilde{\beta}_0s\ ,
\end{aligned}
\]
OK.

Ceci posé, je dis que les contraintes (6.3.7), (6.3.8) sur $h_0$ reviennent à la commutativité des deux triangles dans les carré % [sic] les carré -> le carré, singular (one square (6.3.14))
\[
\begin{tikzcd}
\tilde{B} \arrow[r, "\tilde{h}_0"] \arrow[d, "\tilde{i}\in\Phi"'] & X_0 \arrow[d, "p_0\in\Psi"] \\
B_0 \arrow[r, "\beta_0"'] \arrow[ur, dotted, "h_0"] & Y_0\ ,
\end{tikzcd}
\tag*{(6.3.14)}
\]
\msp{page 116}i.e. aux deux relations
\[
h_0\tilde{i} \ =\ \tilde{h}_0\ , \qquad\quad p_0h_0 \ =\ \beta_0\ .
\tag*{(6.3.15)}
\]

La deuxième de ces relations n'est autre que la première relation (6.3.7), et il reste à voir que la première relation (6.3.15) équivaut à l'ensemble des deux relations restantes dans (6.3.7), (6.3.8) :
\[
h_0i_0 \ =\ \alpha_0\ , \qquad\quad \lambda_X h'_{I'} \ =\ h_0\underbrace{\lambda_B}_{\tilde{i}s}
\tag*{(6.3.16)}
\]
% foot: Version 9 février 2019


\endgroup
% END UNIT der-ch17/der-ch17_p056-060.tex
% BEGIN UNIT der-ch17/der-ch17_p061-062.tex
\begingroup\makeatletter%
\newcommand{\printedpage}[1]{\marginpar{\footnotesize p.~#1}}%
\newcommand{\msp}[1]{\marginpar{\footnotesize [#1]}}%
\newcommand{\cercle}[1]{\tikz[baseline=(c.base)]\node[draw,circle,inner sep=1pt](c){#1};}%
\makeatother\let\printedpage\derpp\let\msp\dermsp%
\let\mspin\DERmspin%


\printedpage{59}
% head: A. Grothendieck : Dérivateurs, Ch. XVII, Édition des Universités de Montpellier II et Paris VII
% foot: Version 9 février 2019

Or la relation en question $h_0\tilde{i} = \tilde{h}_0$ équivaut à l'ensemble des deux relations
\[
\begin{aligned}
(h_0\tilde{i})r &= h_0(\tilde{i}r)\\
&= h_0i_0\\
&\overset{?}{=} \alpha_0\\
&= \tilde{h}_0r,
\end{aligned}
\]
équivalent à $h_0i_0 = \alpha_0$,
\[
\begin{aligned}
(h_0\tilde{i})s &= h_0(\tilde{i}s)\\
&= h_0\lambda_B\\
&\overset{?}{=} \lambda_Xh'_{I'}\\
&= \tilde{h}_0s,
\end{aligned}
\]
équivalent à $h_0\lambda_B = \lambda_Xh'_{I'}$, c'est OK.

Ceci posé, dans le carré (6.3.14), les hypothèses $i \in \Phi(I)$, $p \in \Psi\langle I\rangle$ impliquent $\tilde{i}_0 \in \Phi$, $p_o \in \Psi$. Donc la flèche cherchée $h_0$ existe, q.e.d. % [sic] p_o -> p_0, index printed as letter o, elsewhere index 0 (h_0, i_0)

Cela achève la démonstration du corollaire 6.3, et par là aussi, celle de la dernière partie, d), de 6.1. On est heureux\,!

\msp{page 117}

\textbf{Proposition 6.4.}\footnote[34]{Sous réserve d'existence des $\varinjlim$ et des $\varprojlim$ pertinentes dans $\mathcal{M}$. On peut supposer pour simplifier $\mathcal{M}$ stable par $\varinjlim$ et $\varprojlim$ .} % mark printed as superscript (34)
\textit{Soit $\mathcal{M}$ une catégorie, et $\Phi \longleftrightarrow \Psi$ un couple Q-clos en Q-dualité ($\Phi = \Psi^*$, $\Psi = \Phi_*$). Soit $u : I \longrightarrow I'$ une flèche dans} Cat\textit{, d'où une suite de foncteurs adjoints}
\[
\begin{tikzcd}[column sep=large]
\mathcal{M}(I) \arrow[r, shift left=3, "u_!"] \arrow[r, shift right=3, "u_*"'] & \mathcal{M}(I') \arrow[l, "u^*"']
\end{tikzcd}\ .
\]
\textit{On définit $\Phi(I)$, $\Psi\langle I\rangle$, $\Phi(I')$, $\Psi\langle I'\rangle$ comme dans} 6.2.5. \textit{On a alors :}

\cercle{a}
\[
\left\{
\begin{aligned}
u_!(\Phi(I)) &\ \subset\ \Phi(I')\ ,\\
u^*(\Psi\langle I'\rangle) &\ \subset\ \Psi\langle I\rangle\ .
\end{aligned}
\right.
\tag*{(6.4.1)}
\]

\cercle{b} \textit{On a l'équivalence}
\[
u^*(\Phi(I')) \subset \Phi(I) \qquad\Longleftrightarrow\qquad u_*(\Psi\langle I\rangle) \subset \Psi\langle I'\rangle\ .
\tag*{(6.4.2)}
\]

\cercle{c} \textit{Les relations équivalentes} (6.4.2) \textit{ont lieu dans chacun des trois cas suivants :}
\begin{enumerate}[label={}, leftmargin=2.5em]
\item 1°) \textit{$u$ est coétale.}
\item 2°) \textit{$I$ est discrète.}
\item 3°) \textit{Pour tout $i' \in \operatorname{Ob} I'$, les composantes connexes de $I/i'$ admettent un objet final.}
\end{enumerate}

\par
\printedpage{60}
% head: A. Grothendieck : Dérivateurs, Ch. XVII, Édition des Universités de Montpellier II et Paris VII
% foot: Version 9 février 2019

\textbf{Corollaire 6.5.} \textit{Soit $(W,C,F)$ un triple de Quillen clos spécial dans $\mathcal{M}$, et définissons, pour toute $I$ dans} Cat\textit{, $C(I)$, $TC(I)$ et $F(I)$, $TF(I)$ comme dans le théorème} 6.1. a)\textit{. On a alors}

\msp{page 118}

\textit{les équivalences}
\[
\Bigl(u^*(C(I'))\ \subset\ C(I)\Bigr) \quad\Longleftrightarrow\quad \Bigl(u_*(TF)\ \subset\ TF'\Bigr)\ ,
\tag*{(6.5.1)}
\]
\[
\Bigl(u^*(TC(I'))\ \subset\ TC(I)\Bigr) \quad\Longleftrightarrow\quad \Bigl(u_*(F)\ \subset\ F'\Bigr)\ .
\tag*{(6.5.2)}
\]
\textit{Les conditions équivalentes} (6.5.1) \textit{ainsi que} (6.5.2) \textit{sont satisfaites quand on est dans l'un des cas} 1°) 2°) 3°) \textit{de} 6.4. c).

\textsc{Démonstration} de 6.4. Les assertions \cercle{a} , \cercle{b} sont $\pm$ tautologiques, par adjonction, comparer XIII, prop. 2, et cor., page 46. Pour prouver c), il suffit de traiter le cas 3°), car 1° $\Longrightarrow$ 3°, 2° $\Longrightarrow$ 3°. Mais le calcul de
\[
u_*(X)(i')\ \simeq\ \varprojlim_{i\in I/i'} X(i)
\]
montre que si les composantes connexes $(I/i')_\alpha$ de $I/i'$ ont chacune un objet final $e_\alpha$, alors
\[
u_*(X)(i')\ \simeq\ \prod_\alpha X(e_\alpha)\ .
\]
Or $\Psi$, étant clos à droite, est stable par produits quelconques, donc $X \in \Psi\langle I\rangle$ implique que les $u_*(X)(i')$ sont dans $\Psi$, donc $u_*$ est dans $\Psi\langle I'\rangle$, q.e.d.


\endgroup
% END UNIT der-ch17/der-ch17_p061-062.tex
\clearpage\setcounter{footnote}{0}\phantomsection\addcontentsline{toc}{section}{Chapitre XVIII}\markboth{Chapitre XVIII}{Chapitre XVIII}
% BEGIN UNIT der-ch18/der-ch18_p001-005.tex
\begingroup\makeatletter%
\newcommand{\printedpage}[1]{\marginpar{\footnotesize p.~#1}}%
\newcommand{\msp}[1]{\marginpar{\footnotesize [#1]}}%
\newcommand{\uHom}{\underline{\mathrm{Hom}}}%
\newcommand{\Ens}{\mathrm{Ens}}%
\newcommand{\Fais}{\mathrm{Fais}}%
\newcommand{\pstop}{\mathrm{pstop}}%
\makeatother\let\printedpage\derpp\let\msp\dermsp%
\let\mspin\DERmspin%


\printedpage{} % no page number printed
% head: (none)

\begin{center}
{\Large\textsf{LES DÉRIVATEURS}}

\bigskip

\textbf{ALEXANDRE GROTHENDIECK}

\bigskip\bigskip

\textbf{Chapitre XVIII}

\medskip

\textbf{Catégories et ensembles accessibles}
\end{center}

\par

\printedpage{} % no page number printed
% head: (none)

\par\noindent\textbf{Note des éditeurs de l’édition transcrite.}\par
Ce texte a été déchiffré et transcrit en \LaTeXe{} par M. Künzer. Il a été édité par M. Künzer, J. Malgoire, et G. Maltsiniotis. La transcription est aussi fidèle que possible au manuscrit. Seules quelques corrections mineures évidentes ont été effectuées. Pour les rares ajouts ou commentaires des éditeurs, ainsi que pour la numérotation originale des pages du manuscrit, les caractères de machine à écrire \texttt{[typewriter]} entre crochets sont utilisés. Quelques notes de bas de page ajoutées par les éditeurs sont signalées par «~N.~Éd.~». La numérotation des pages de ce chapitre commence à la page 50, car à l'origine ce chapitre était la suite du chapitre XIX qui se termine au milieu de la page 50.

Cette édition est provisoire. Les remarques, commentaires et corrections sont bienvenus. Envoyer un message à :

\textbf{[e-mail address removed]}

\hfill G. Maltsiniotis\hspace*{2cm}

\par

\printedpage{1}
% head: A. Grothendieck : Dérivateurs, Ch. XVIII, Édition des Universités de Montpellier II et Paris VII
% foot: Version 17 décembre 2025

\msp{page 50}

\section*{4\quad Paratopos et l'opération $\mathcal{M}\boxtimes\mathcal{N}$}

\subsection*{4.1\quad Pseudo-topos. Problème du produit.}

J'appelle \textit{catégorie pseudo-$\mathfrak{U}$-topos}, ou simplement \textit{pseudo-topos}, une $\mathfrak{U}$-catégorie $\mathcal{M}$ stable par $\mathfrak{U}$-limites inductives, et qui admet une \textit{petite} sous-catégorie $\mathcal{M}_0$ \textit{strictement génératrice} (i.e. génératrice par épimorphismes stricts). Alors $\mathcal{M}$ admet aussi des $\mathfrak{U}$-limites projectives (SGA 4 I 8.12.7). (Mais l'exemple de (Gr), cf. loc. cit. 8.12.9, montre que $\mathcal{M}^\circ$ n'est par forcément, pour autant, un pseudo-$\mathfrak{U}$-topos, faute à $\mathcal{M}$ d'admettre une petite sous-catégorie strictement \textit{co}génératrice. Mais le théorème 2.6 \texttt{[du chapitre XIX]} donne des conditions sur $\mathcal{M}$ (surtout du type exactitude) % [sic] par forcément -> pas forcément, misprint; also the parenthesis opened at "(Mais" is never closed
\msp{page 51}
pour que $\mathcal{M}$ admette assez d'injectifs -- ce qui impliquera alors l'existence d'une petite sous-catégorie strictement cogénératrice, donc que $\mathcal{M}^\circ$ est également un pseudo-topos (\footnote{C'est le cas notamment si $\mathcal{M}$ est un topos.}).

Je visualise les catégories pseudo-topos comme étant la catégorie des “faisceaux” sur un objet $\mathcal{X}$, visualisé comme une sorte d'“espace” que je distingue de la catégorie $\mathcal{M}=\Fais(\mathcal{X})$ des faisceaux sur $\mathcal{X}$. C'est $\mathcal{X}$ qui sera pour ceci le “pseudo-topos”. Les pseudo-topos forment une 2-catégorie, la variance dans $\mathcal{X}$ étant opposée de celle dans $\Fais(\mathcal{X})=\mathcal{M}$. De façon précise, si $\mathcal{X}$ correspond à $\mathcal{M}$, $\mathcal{Y}$ à $\mathcal{N}$, un morphisme $f$ de $\mathcal{X}$ dans $\mathcal{Y}$ sera donné par un foncteur
\[
f^* \;:\; \mathcal{N} \longrightarrow \mathcal{M}
\tag*{(4.1.1)}
\]
en sens inverse, commutant aux $\varinjlim$ inductives quelconques (\footnote{Dire ici que les objets de $\mathcal{M}$ sont visualisés comme des “ouverts” généralisés de $\mathcal{X}$. On a $\mathcal{M}\xrightarrow{\sim}\uHom^!(\mathcal{M}^\circ,\Ens)$, i.e. $\mathcal{M}$ s'interprète bel et bien comme les “faisceaux” sur cette catégorie $\mathcal{M}$ d'ouverts \ldots}). (\textbf{N.B.} si $\mathcal{X}$, $\mathcal{Y}$ sont même des topos, i.e. $\mathcal{M}$, $\mathcal{N}$ des catégories-topos, et si on veut que $f$ soit un morphisme de topos, on exigera de plus que $f^*$ soit de plus \textit{exact à gauche}. On dira qu'un \textit{morphisme de pseudo-topos est exact} si de plus $f^*$ est exact à gauche.)

\msp{page 52}
En vertu du dual de loc. cit. 8.12.7, cela équivaut au fait que $f^*$ admette un adjoint à droite $f_*$,
\[
f_* \;:\; \mathcal{M} \longrightarrow \mathcal{N}\;,
\tag*{(4.1.2)}
\]
lequel commute donc aux petites $\varprojlim$\,. La donnée de $f_*$ satisfaisant à cette propriété détermine donc $f^*$ à isomorphisme unique près. Mais en général, il ne suffit pas de se donner un tel $f_*$, pour pouvoir affirmer qu'il corresponde à un $f^*$, i.e. qu'il admette un adjoint à droite\,; pour ceci, il faudrait qu'on sache que les foncteurs $\mathcal{M}\longrightarrow\Ens$ qui commutent aux petites $\varprojlim$ sont représentables, ce qui sera le cas si $\mathcal{M}$ admet une petite sous-catégorie \textit{co}génératrice, i.e. si $\mathcal{M}^\circ$ est également une catégorie pseudo-topos. % [sic] qu'il admette un adjoint à droite -> adjoint à gauche, f^* is the left adjoint of f_*

\par

\printedpage{2}
% head: A. Grothendieck : Dérivateurs, Ch. XVIII, Édition des Universités de Montpellier II et Paris VII
% foot: Version 17 décembre 2025

Il est sans doute judicieux de procéder comme dans le cas des topos, et d'appeler \textit{morphisme de pseudo-topos} $f:\mathcal{X}\longrightarrow\mathcal{Y}$ un quadruple
\[
(f^*,f_*,\alpha,\beta)
\tag*{(4.1.3)}
\]
\msp{page 53}
de deux foncteurs adjoints $(f^*,f_*)$, et des morphismes d'adjonction. Quand $f^*$ est une équivalence, ce qui équivaut à $f_*$ une équivalence, ou $\alpha$ et $\beta$ des isomorphismes, on dira qu'on a une \textit{équivalence de pseudo-topos}, et à toutes fins pratiques, on pourra alors regarder $\mathcal{X}$ et $\mathcal{Y}$ comme des objets essentiellement identiques.

Pour deux pseudo-topos $\mathcal{X}$, $\mathcal{Y}$, j'ai envie d'introduire (comme dans le cas des topos) un pseudo-topos produit
\[
\mathcal{X}\times\mathcal{Y}\;,
\tag*{(4.1.4)}
\]
qui ait les propriétés d'un 2-produit dans la 2-catégorie des pseudo-topos (\footnote{Il apparaît finalement que ce qu'on construit n'est \textit{pas} un 2-produit dans la 2-catégorie envisagée, donc il vaut mieux le noter par un signe tel que $\mathcal{X}\boxtimes\mathcal{Y}$. Mais si $\mathcal{X}$, $\mathcal{Y}$ sont deux topos, $\mathcal{X}\boxtimes\mathcal{Y}$ est un 2-produit dans la 2-catégorie des topos.}). Je devrais préciser que par définition
\[
\begin{aligned}
\uHom_{\pstop}(\mathcal{X},\mathcal{Y}) \;&\simeq\; \uHom_!(\mathcal{F}(\mathcal{Y}),\mathcal{F}(\mathcal{X}))^\circ\\
&\simeq\; \uHom^{!'}(\mathcal{F}(\mathcal{X}),\mathcal{F}(\mathcal{Y})),
\end{aligned}
\tag*{(4.1.5)}
\]
où l'indice ! (resp. l'exposant !) indique la sous-catégorie pleine dans la catégorie $\uHom$ envisagée formée des foncteurs qui commutent aux petites $\varinjlim$ (resp. aux petites
\msp{page 54}
$\varprojlim$\,), et l'accent dans le $\uHom^{!'}$ dans le dernier membre indique qu'on exige de plus que le foncteur en question ait un adjoint à gauche.

Soient donc
\[
\mathcal{M} \;=\; \mathcal{F}(\mathcal{X})\;, \qquad \mathcal{N} \;=\; \mathcal{F}(\mathcal{Y})\;,
\tag*{(4.1.6)}
\]
et soit $\mathcal{Z}$ un troisième pseudo-topos, décrit par
\[
\mathcal{P} \;=\; \mathcal{F}(\mathcal{Z})\;.
\tag*{(4.1.7)}
\]
On veut expliciter, à équivalence près, la catégorie produit $\uHom_{\pstop}(\mathcal{Z},\mathcal{X})\times\uHom_{\pstop}(\mathcal{Z},\mathcal{Y})$, i.e.
\[
\uHom_!(\mathcal{M},\mathcal{P})^\circ \;\times\; \uHom_!(\mathcal{N},\mathcal{Q})^\circ.
\tag*{(4.1.8)}
\] % [sic] \mathcal{Q} -> \mathcal{P}, \mathcal{Q} is not defined; \mathcal{Z} is described by \mathcal{P} (4.1.7)
Je vais introduire pour cela une catégorie
\printedpage{3}
% head: A. Grothendieck : Dérivateurs, Ch. XVIII, Édition des Universités de Montpellier II et Paris VII
% foot: Version 17 décembre 2025
\msp{page 55}
\[
\mathcal{M}\boxtimes\mathcal{N}
\tag*{(4.1.9)}
\]
par la formule
\[
\mathcal{M}\boxtimes\mathcal{N} \;\simeq\; \uHom^{!!}(\mathcal{M}^\circ\times\mathcal{N}^\circ,\Ens)\;,
\tag*{(4.1.10)}
\]
où le signe $\uHom^{!!}$ désigne la sous-catégorie pleine de la catégorie $\uHom$, formée des bi-contrafoncteurs % doubt resolved: retain the diplomatic bi-contrafoncteurs reading; see corrections.csv, PDF leaf 5
\[
F \;:\; \mathcal{M}^\circ\times\mathcal{N}^\circ \longrightarrow \Ens\;,
\tag*{(4.1.10.1)}
\]
tels que pour tout $x$ dans $\mathcal{M}$, le foncteur
\[
y \longmapsto F(x,y) \;:\; \mathcal{N}^\circ \longrightarrow \Ens
\tag*{(4.1.10.2)}
\]
soit dans $\uHom^!(\mathcal{N}^\circ,\Ens)$, i.e. transforme petites $\varinjlim$ en $\varprojlim$\,, et que, symétriquement, pour tout $y$ dans $\mathcal{N}$, le foncteur
\[
x \longmapsto F(x,y) \;:\; \mathcal{M}^\circ \longrightarrow \Ens
\tag*{(4.1.10.3)}
\]
soit dans $\uHom^!(\mathcal{M}^\circ,\Ens)$. Mais cela signifie aussi que ces foncteurs sont \textit{représentables}\,; donc on trouve aussi
\[
\mathcal{M}\boxtimes\mathcal{N} \xleftarrow{\sim} \uHom^!(\mathcal{N}^\circ,\mathcal{M})
\tag*{(4.1.11)}
\]
et
\[
\mathcal{M}\boxtimes\mathcal{N} \xleftarrow{\sim} \uHom^!(\mathcal{M}^\circ,\mathcal{N})
\tag*{(4.1.12)}
\]
\msp{page 56}
en utilisant (dans chacune de ces formules) dans deux côtés les équivalences de catégories
\[
\left\{
\begin{aligned}
&\mathcal{M} &&\xrightarrow{\sim} &&\uHom^!(\mathcal{M}^\circ,\Ens)\\
&\mathcal{N} &&\xrightarrow{\sim} &&\uHom^!(\mathcal{N}^\circ,\Ens)\;.
\end{aligned}
\right.
\tag*{(4.1.13)}
\]
On peut dire aussi que la catégorie $\mathcal{M}\boxtimes\mathcal{N}$ est équivalente à celle des \textit{couples} de foncteurs
\[
\left\{
\begin{aligned}
&\varphi &&: &&\mathcal{M}^\circ &&\longrightarrow &&\mathcal{N}\\
&\psi &&: &&\mathcal{N}^\circ &&\longrightarrow &&\mathcal{M}\;,
\end{aligned}
\right.
\tag*{(4.1.14)}
\]
munis de deux morphismes d'adjonction
\[
\left\{
\begin{aligned}
&\psi\varphi^\circ &&\xleftarrow{\alpha} &&\mathrm{id}_{\mathcal{M}}\\
&\varphi\psi^\circ &&\xleftarrow{\beta} &&\mathrm{id}_{\mathcal{N}}
\end{aligned}
\right.
\tag*{(4.1.15)}
\]
satisfaisant les deux compatibilités habituelles exprimant l'adjonction de $\varphi^\circ$ et de $\psi$ (ou, de façon équivalente, de $\psi^\circ$ et de $\varphi$). Le foncteur
\[
(\varphi,\psi) \longmapsto \varphi \qquad (\text{resp.}\;\; (\varphi,\psi) \longmapsto \psi)
\]
donne une équivalence de la catégorie de ces couples avec $\uHom^!(\mathcal{M}^\circ,\mathcal{N})$ (resp. avec $\uHom^!(\mathcal{N}^\circ,\mathcal{M})$), de sorte qu'on a un carré commutatif d'équivalences canoniques de foncteurs :


\endgroup
% END UNIT der-ch18/der-ch18_p001-005.tex
% BEGIN UNIT der-ch18/der-ch18_p006-010.tex
\begingroup\makeatletter%
\newcommand{\printedpage}[1]{\marginpar{\footnotesize p.~#1}}%
\newcommand{\msp}[1]{\marginpar{\footnotesize [#1]}}%
\newcommand{\cM}{\mathcal{M}}%
\newcommand{\cN}{\mathcal{N}}%
\newcommand{\cX}{\mathcal{X}}%
\newcommand{\cY}{\mathcal{Y}}%
\newcommand{\cF}{\mathcal{F}}%
\newcommand{\uHom}{\underline{\mathrm{Hom}}}%
\newcommand{\Hom}{\mathrm{Hom}}%
\newcommand{\Ens}{\mathrm{Ens}}%
\newcommand{\Top}{\mathrm{Top}}%
\newcommand{\Ob}{\mathrm{Ob}\,}%
\newcommand{\defeq}{\overset{\text{déf}}{=}}%
\newcommand{\bt}{\boxtimes}%
\newcommand{\tbt}{\mathbin{\widetilde{\boxtimes}}}%
\newcommand{\linj}{\lhook\joinrel\longrightarrow}%
\makeatother\let\printedpage\derpp\let\msp\dermsp%
\let\mspin\DERmspin%


\printedpage{4}
% head: A. Grothendieck : Dérivateurs, Ch. XVIII, Édition des Universités de Montpellier II et Paris VII

\msp{page 57}

\[
\begin{tikzcd}[column sep=small]
 & \cM \tbt \cN \arrow[dl, "\sim" sloped] \arrow[dr, "\sim" sloped] & \\
\uHom^{!}(\cM^{\circ},\cN) \arrow[dr, "\sim" sloped] & & \uHom^{!}(\cN^{\circ},\cM) \arrow[dl, "\sim" sloped] \\
 & \underbrace{\cM \bt \cN}_{\defeq \uHom^{!!}(\cM^{\circ}\times\cN^{\circ},\Ens)} &
\end{tikzcd}
\tag*{(4.1.16)}
\]

\texttt{[où $\cM \tbt \cN$ est la catégorie des couples définie dans la page précédente.]}

Je dis que $\cM \bt \cN$ \textit{est une catégorie pseudo-topos}. Il est clair, \texttt{[en]} outre, qu'elle admet des petites $\varprojlim$ , qui se calculent argument par argument dans $\uHom(\cM^{\circ}\times\cN^{\circ},\Ens)$, ou au choix dans $\uHom^{!}(\cM^{\circ},\cN)$, ou $\uHom^{!}(\cN^{\circ},\cM)$. Comment construire des petites $\varinjlim$ ? Il n'est pas question d'invoquer un argument de cogénérateurs. Il faut d'ailleurs trouver aussi une petite sous-catégorie strictement génératrice. Visiblement, il y a quelque chose de non trivial à démontrer !

\section*{4.2\quad Cafouillages sur $\cM \bt \cN$}

Revenons à la situation générale d'un pseudo-topos $\cX$, décrit par la catégorie pseudo-topos $\cM$. Si $\cN$ est une catégorie quelconque, on définit la catégorie

\msp{page 58}
\[
\mathrm{Fais}(\cX,\cN) \qquad \text{ou} \qquad \cF(\cX,\cN)
\]
des faisceaux sur $\cX$, à valeurs dans $\cN$, par
\[
\cF(\cX,\cN) \ \defeq\ \uHom^{!}(\cM^{\circ},\cN) \ .
\tag*{(4.2.1)}
\]
On a donc un foncteur pleinement fidèle
\[
\cF(\cX,\cN) \ \defeq\ \uHom^{!}(\cM^{\circ},\cN) \ \linj\ \Hom^{!}(\cN^{\circ},\overline{\cM}) \ ,
\tag*{(4.2.2)}
\]
où $\overline{\cM}$ désigne la sous-catégorie de $\cM^{\wedge}$ (\textbf{N.B.} $\cM^{\wedge}$ n'est pas en général une $\mathcal{U}$-catégorie qu'à cela ne tienne) formée des foncteurs \textit{représentables} $\cM^{\circ}\longrightarrow\Ens$ (laquelle est une $\mathfrak{U}$-catégorie, équivalente à $\cM$ par l'inclusion canonique $\cM \linj \overline{\cM}$). Donc on a à isomorphisme unique près une factorisation de (4.2.2) en
\[
\cF(\cX,\cN) \ \overset{\text{pl. fid.}}{\lhook\joinrel\relbar\joinrel\relbar\joinrel\longrightarrow}\ \uHom^{!}(\cN^{\circ},\cM) \ ,
\tag*{(4.2.3)}
\]
l'image essentielle étant formée des foncteurs
\[
\psi \ :\ \cN^{\circ} \longrightarrow \cM
\]
tels que pour tout $x$ dans $\cM$,
\[
y \longmapsto \Hom(y,\psi(x)) \ :\ \cN^{\circ} \longrightarrow \Ens
\]
soit représentable, i.e. des foncteurs qui
% foot: Version 17 décembre 2025
\printedpage{5}
% head: A. Grothendieck : Dérivateurs, Ch. XVIII, Édition des Universités de Montpellier II et Paris VII
\msp{page 59}
admettent un \textit{adjoint à gauche}. Quand $\cN$ lui-même est une catégorie pseudo-topos, cette condition est automatique (elle résulte de l'hypothèse que $\psi$ commute aux petites $\varprojlim$ ). Donc dans ce cas, (4.2.3) est une équivalence de catégories.

Soit maintenant $S$ ($S$ comme “site”) une \textit{petite} sous-catégorie de $\cM$, génératrice par épimorphismes stricts. On sait que cela signifie que le foncteur composé
\[
\cM \xrightarrow{\ \sim\ } \uHom^{!}(\cM^{\circ},\Ens) \xrightarrow{\text{restriction à } S} S^{\wedge} \ (\ \defeq \uHom(S^{\circ},\Ens)\ )
\tag*{(4.2.4)}
\]
est pleinement fidèle, et ceci implique plus généralement, que pour toute catégorie $\cN$, le foncteur restriction
\[
\cF(\cX,\cN) \ \defeq\ \uHom^{!}(\cM^{\circ},\cN) \ \overset{\text{pl. fid.}}{\lhook\joinrel\relbar\joinrel\relbar\joinrel\longrightarrow}\ \uHom(S^{\circ},\cN)
\tag*{(4.2.5)}
\]
est pleinement fidèle. En effet, le deuxième membre, via le plongement pleinement fidèle
\[
\cN \ \overset{\text{pl. fid.}}{\lhook\joinrel\relbar\joinrel\relbar\joinrel\longrightarrow}\ \uHom^{!}(\cN^{\circ},\Ens) \ \linj\ \cN^{\wedge} \ ,
\]
se plonge dans
\[
\uHom(S^{\circ},\uHom^{!}(\cN^{\circ},\Ens)) \ \simeq\ \Hom^{!}(\cN^{\circ},S^{\wedge}) \ ,
\]
et la pleine fidélité de (4.2.5) équivaut donc à celle du foncteur composé
\[
\cF(\cX,\cN) \ \defeq\ \uHom^{!}(\cM^{\circ},\cN) \ \longrightarrow\ \uHom^{!}(\cN^{\circ},S^{\wedge}) \ \simeq\ \uHom(S^{\circ},\uHom^{!}(\cN^{\circ},\Ens)) \ .
\tag*{(4.2.6)}
\]

\msp{page 60}
Or on a commutativité dans le carré
\[
\begin{tikzcd}[column sep=large, row sep=large]
\cF(\cX,\cN) \defeq \uHom^{!}(\cM^{\circ},\cN) \arrow[r, "\text{\scriptsize(4.2.5)}"] \arrow[d, hook, "\text{\scriptsize(4.2.3)}"', "\text{\scriptsize pleinement fidèle}"] & \uHom(S^{\circ},\cN) \arrow[d, hook, "\text{\scriptsize pleinement fidèle}"] \\
\uHom^{!}(\cN^{\circ},\cM) \arrow[r, hook, "\text{\scriptsize(4.2.4)}", "\text{\scriptsize pleinement fidèle}"'] & \underbrace{\uHom^{!}(\cN^{\circ},S^{\wedge})}_{\simeq\, \uHom(S^{\circ},\uHom^{!}(\cN^{\circ},\Ens))}
\end{tikzcd}
\tag*{(4.2.7)}
\]
(où les flèches verticales sont des équivalences, si $\cN$ est lui-même une catégorie pseudo-topos). Notons d'autre part que l'on a une équivalence
\[
\uHom(S^{\circ},\cN) \ \xleftarrow{\ \sim\ }\ \uHom^{!}(S^{\wedge\circ},\cN) \ \defeq\ \cF(\Top(S),\cN) \ ,
\tag*{(4.2.8)}
\]
où $\Top(S)$ est le pseudo-topos (en fait, un topos) défini par la catégorie pseudo-topos $S^{\wedge}$. Pour ceci, notons qu'on a
\[
\uHom^{!}(S^{\wedge\circ},\cN) \ \xrightarrow{\ \sim\ }\ \uHom_{!}(S^{\wedge},\cN^{\circ})^{\circ} \ ,
\]
et d'autre part le foncteur restriction à $S \linj S^{\wedge}$ donne
\[
\uHom_{!}(S^{\wedge},\cN^{\circ}) \ \xrightarrow{\ \sim\ }\ \uHom(S,\cN^{\circ}) \ ,
\]
d'où
% foot: Version 17 décembre 2025
\printedpage{6}
% head: A. Grothendieck : Dérivateurs, Ch. XVIII, Édition des Universités de Montpellier II et Paris VII
\msp{page 61}
\[
\uHom(S^{\wedge\circ},\cN) \ \xrightarrow{\ \sim\ }\ (\uHom(S,\cN^{\circ}))^{\circ} \ \xrightarrow{\ \sim\ }\ \uHom(S^{\circ},\cN) \ ,
\]
ce qui n'est autre que (4.2.8). Ainsi :

\textbf{Proposition 4.2.1.} \textit{Soient $\cM=\cF(\cX)$ une catégorie pseudo-topos, munie d'une sous-catégorie pleine strictement génératrice $S$, $\cN$ une $\mathfrak{U}$-catégorie quelconque, alors le foncteur “restriction à $S$” suivant est pleinement fidèle.}
\[
\cF(\cX,\cN) \ \defeq\ \uHom^{!}(\cM^{\circ},\cN) \ \longrightarrow\ \uHom(S^{\circ},\cN) \ \xleftarrow{\ \sim\ }\ \underbrace{\uHom(S^{\wedge\circ},\cN)}_{\defeq\, \cF(\Top(S),\cN)} \ ,
\tag*{(4.2.1.1)}
\]
\textit{et il s'insère dans le carré commutatif (4.2.7) de foncteurs pleinement fidèles.}

\textbf{Corollaire 4.2.2.} \textit{Supposons que $\cN=\cF(\cY)$ soit également une catégorie pseudo-topos, munie d'une sous-catégorie pleine $T$ strictement génératrice. Alors}
\[
\cF(\cX,\cN) \ \simeq\ \cF(\cY,\cM) \ \simeq\ \underbrace{\uHom^{!!}(\cM^{\circ}\times\cN^{\circ},\Ens)}_{=\,\cM\bt\cN} \ ,
\]
\textit{et le foncteur restriction}
\[
\begin{aligned}
\cM\bt\cN \ \defeq\ \uHom^{!!}(\cM^{\circ}\times\cN^{\circ},\Ens) \quad &\longrightarrow \quad \uHom(S^{\circ}\times T^{\circ},\Ens)\\
&\simeq \quad (S\times T)^{\wedge}\\
&\defeq \quad \cF(\Top(S\times T))
\end{aligned}
\tag*{(4.2.2.1)}
\]
\textit{est pleinement fidèle.}

La première assertion a été vue dans 4.1. Pour la pleine fidélité de (4.2.2.1), on note que ce foncteur s'interprète, à équivalence de catégories près, comme le composé du foncteur pleinement fidèle (4.2.1.1), et du foncteur
\msp{page 62}
similaire relatif à $\cN$ (donc lui aussi pleinement fidèle)
\[
\cF(\cY,S^{\wedge}) \ \defeq\ \uHom^{!}(\cN^{\wedge},S^{\wedge}) \ \xrightarrow{\ \text{restr. à } T\ }\ \underbrace{\uHom(T^{\circ},S^{\wedge})}_{\uHom(S^{\circ}\times T^{\circ},\Ens)} \ ,
\tag*{(4.2.2.2)} % [sic] \uHom^{!}(\cN^{\wedge},S^{\wedge}) -> \uHom^{!}(\cN^{\circ},S^{\wedge}), cf. (4.2.2.3) where the same F(Y,S^) is Hom!(N°,S^)
\]
compte tenu de l'équivalence de symétrie
\[
\cF(\cY,S^{\wedge}) \ \defeq\ \uHom^{!}(\cN^{\circ},S^{\wedge}) \ \simeq\ \uHom^{!}(S^{\wedge\circ},\cN) \ .
\tag*{(4.2.2.3)}
\]

Au total, quand $\cM$, $\cN$ sont tous deux des catégories pseudo-topos, munies de sous-catégories strictement génératrices $S$ resp. $T$, on a un plongement pleinement fidèle canonique
\[
\cM\bt\cN \ \overset{\rho}{\linj}\ (S\times T)^{\wedge} \ ,
\tag*{(4.2.9)}
\]
% foot: Version 17 décembre 2025
\printedpage{7}
% head: A. Grothendieck : Dérivateurs, Ch. XVIII, Édition des Universités de Montpellier II et Paris VII
qu'on interprète par la suite comme $(f_{S}^{\cM}\times f_{T}^{\cN})_{*}$, où
\[
\left\{
\begin{aligned}
f_{S}^{\cM} \ &:\ \cX \ \longrightarrow\ \Top S\\
f_{T}^{\cN} \ &:\ \cY \ \longrightarrow\ \Top T
\end{aligned}
\right.
\tag*{(4.2.10)}
\]
sont les morphismes de pseudo-topos, donnés par les plongements pleinement fidèles (foncteurs image directe)
\[
\left\{
\begin{aligned}
(f_{S}^{\cM})_{*} \ &:\ \cM \ \longrightarrow\ S^{\wedge}\\
(f_{T}^{\cN})_{*} \ &:\ \cN \ \longrightarrow\ T^{\wedge} \ ,
\end{aligned}
\right.
\tag*{(4.2.11)}
\]
lesquels admettent bel et bien des adjoints à gauche
\msp{page 63}
\[
\left\{
\begin{aligned}
(f_{S}^{\cM})^{*} \ &:\ S^{\wedge} \ \longrightarrow\ \cM\\
(f_{T}^{\cN})^{*} \ &:\ T^{\wedge} \ \longrightarrow\ \cN \ ,
\end{aligned}
\right.
\tag*{(4.2.12)}
\]
à savoir les prolongements canoniques (par commutativité aux petites $\varinjlim$ ) des foncteurs d'inclusion $S \linj \cM$, $T \linj \cN$.

La question essentielle, à présent, c'est visiblement de construire un adjoint à gauche pour le foncteur (4.2.9). Il revient au même, pour $(s,t)$ dans $S\times T$, de définir un objet
\[
s\bt t \ \in\ \Ob\cM\bt\cN \ ,
\tag*{(4.2.13)}
\]
tel que l'on ait, pour $F$ variable dans $\cM\bt\cN$, un isomorphisme fonctoriel en $F$,
\[
\underbrace{\rho(F)(s,t)}_{\defeq\, F(s,t)} \ \simeq\ \Hom(s\bt t,F) \ .
\]
Donc il faut prouver ceci :

\textbf{Lemme 4.2.3} (\footnote[4]{Pas prouvé ; peut-être faux !}). \textit{Soient $\cM$, $\cN$ deux catégories pseudo-topos, et $s\in\Ob\cM$, $t\in\Ob\cN$, et considérons sur $\cM\bt\cN\defeq\uHom^{!!}(\cM^{\circ}\times\cN^{\circ},\Ens)$ le foncteur $F\longmapsto F(s,t)$. Ce foncteur est} représentable \textit{(par un objet que je note $s\bt t$).}

L'idée naturelle serait de définir $s\bt t$ par
\[
(s\bt t)(x,y) \ =\ s(x)\times t(y) \ =\ \Hom_{\cM}(x,s)\times\Hom_{\cN}(y,t) \ ,
\]

\msp{page 64}
ce qui définit bien un foncteur
\[
\cM^{\circ}\times\cN^{\circ} \longrightarrow \Ens \ .
\]
Malheureusement, pour un des arguments fixé, disons pour $y$ fixé, ce foncteur n'est \textit{pas} compatible aux limites projectives, p.ex. il ne transforme pas l'objet final de $\cM^{\circ}$ correspondant à l'objet initial de $\cM$) en objet final de Ens, mais en $t(y)=\Hom_{\cN}(y,t)$. % [sic] "de $\cM$)" -> closing parenthesis without opening one, unbalanced as printed
% foot: Version 17 décembre 2025
\printedpage{8}
% head: A. Grothendieck : Dérivateurs, Ch. XVIII, Édition des Universités de Montpellier II et Paris VII
Néanmoins, je pense que $s\bt t$ doit bel et bien se définir comme produit de deux objets de $\cM\bt\cN$, dépendant fonctoriellement de $s$ resp. de $t$ (\footnote[5]{également faux, cf. plus bas.}). On est donc amené dès à présent, qu'on le veuille ou non, à essayer de définir des foncteurs canoniques, commutant aux petites $\varinjlim$ ,
\[
\cM \longrightarrow \cM\bt\cN \longleftarrow \cN
\tag*{(4.2.13)} % [sic] (4.2.13) -> (4.2.14), the label (4.2.13) is already used for s⊠t ∈ Ob M⊠N
\]
(lesquels seront les foncteurs image inverse pour les projections canoniques
\begin{tikzcd}[row sep=small, column sep=small, baseline=(\tikzcdmatrixname-2-1.base)]
 & \cX \\
\cX\times\cY \arrow[ur] \arrow[dr] & \\
 & \cY
\end{tikzcd}
) (\footnote[6]{Le cas où $\cM$, $\cN$ sont des catégories-topos, nous montre que ces foncteurs devraient être décrits comme $x\longmapsto x\bt e_{\cN}$, $y\longmapsto e_{\cM}\bt y$. Mais il n'est pas clair qu'on ait en dehors du cas des topos, et même en admettant l'existence du $x\bt y$, la formule (valable pour un topos) $x\bt y\simeq(x\bt e_{\cN})\times(e_{\cM}\bt y)$.}).

\section*{4.3\quad L'opération $\cM\bt\cN$ et l'axiome d'accessibilité}

J'ai passé la journée hier, essayant d'établir l'existence des $x\bt y$ dans $\cM\bt\cN$, sans y parvenir – du moins pas sans faire une hypothèse supplémentaire d'accessibilité,
\msp{page 65}
dont il devient évident à présent qu'il faut la faire de toutes façons. Je vais quand même énoncer le positif de mes perplexités.

\textbf{Proposition-perplexité 4.3.1.} \textit{Soient $\cM$, $\cN$ deux catégories pseudo-topos, d'où $\cM\bt\cN$ comme dans 4.1. Soient $S\subset\cM$, $T\subset\cN$ deux petites sous-catégories strictement génératrices, et considérons l'inclusion pleinement fidèle correspondante :}
\[
\begin{tikzcd}[column sep=large]
\underbrace{\cM\bt\cN}_{\defeq\, \uHom^{!!}(\cM^{\circ}\times\cN^{\circ},\Ens)} \arrow[r, hook, "\alpha"'] & S^{\wedge}\bt T^{\wedge} \ \simeq\ (S\times T)^{\wedge}. \arrow[l, dotted, bend right, "\varphi\,?"']
\end{tikzcd}
\tag*{(4.3.1.1)}
\]
\textit{Considérons les conditions suivantes :}
\begin{enumerate}[label={}]
\item[\textit{a)}] \textit{$\cM\bt\cN$ est un pseudo-topos.}
\item[\textit{b)}] \textit{Pour $x\in\Ob\cM$, $y\in\Ob\cN$, $x\bt y\in\Ob\cM\bt\cN$ existe, i.e. le foncteur (commutant aux petites $\varprojlim$ )}
\[
F \ \longmapsto\ F(x,y) \ , \qquad \cM\bt\cN \longrightarrow \Ens
\tag*{(4.3.1.2)}
\]
\textit{est représentable.}
\item[\textit{b$'$)}] \textit{Comme b), avec $x\in\Ob S$, $y\in\Ob T$.}
\item[\textit{c)}] \textit{Le foncteur $\alpha$ (lequel commute aux petites $\varprojlim$ ) admet un adjoint à gauche}
\[
\varphi \ :\ (S\times T)^{\wedge} \longrightarrow \cM\bt\cN \ .
\]
\end{enumerate}
\textit{On a alors les implications}
\[
b \ \Longrightarrow\ b' \ \Longleftrightarrow\ c \ \Longrightarrow\ a \ .
\tag*{(4.3.1.3)}
\]
% foot: Version 17 décembre 2025


\endgroup
% END UNIT der-ch18/der-ch18_p006-010.tex
% BEGIN UNIT der-ch18/der-ch18_p011-015.tex
\begingroup\makeatletter%
\newcommand{\printedpage}[1]{\marginpar{\footnotesize p.~#1}}%
\newcommand{\msp}[1]{\marginpar{\footnotesize [#1]}}%
\newcommand{\Hom}{\operatorname{Hom}}%
\newcommand{\Ob}{\operatorname{Ob}}%
\newcommand{\pr}{\operatorname{pr}}%
\newcommand{\card}{\operatorname{card}}%
\newcommand{\Fl}{\operatorname{Fl}}%
\newcommand{\Filt}{\operatorname{Filt}}%
\newcommand{\Ind}{\operatorname{Ind}}%
\makeatother\let\printedpage\derpp\let\msp\dermsp%
\let\mspin\DERmspin%

\setcounter{footnote}{6}

\printedpage{9}
% head: A. Grothendieck : Dérivateurs, Ch. XVIII, Édition des Universités de Montpellier II et Paris VII

\noindent\textit{De plus, si c est satisfait, on a ceci : pour $s\in S$, $t\in T$, considérons}
\[
(s,t)\ \in\ (S\times T)\ \subset\ (S\times T)^\wedge\ \simeq\ S^\wedge\boxtimes T^\wedge\ ,
\]
\textit{qui s'explicite aussi comme}
\[
\left\{
\begin{aligned}
&\rlap{$\displaystyle (s,t)\ =\ \pr_1^*(s)\times\pr_2^*(t)\ \simeq\ s\overset{S^\wedge,T^\wedge}{\boxtimes}t\in\Ob\,(S\times T)^\wedge\ ,$}\\
&\textit{i.e.}\ (s,t)(s',t')\ =\ s(s')\times t(t') &&=\ \Hom_S(s',s)\times\Hom_T(t',t)\\
& &&=\ \Hom_{S\times T}((s',t'),(s,t))\ ,
\end{aligned}
\right.
\tag*{(4.3.1.4)}
\]
\[
\begin{tikzcd}
S\times T \arrow[r, "\mathrm{pr}_2"] \arrow[d, "\mathrm{pr}_1"] & T\\
S\ ,
\end{tikzcd}
\]
\msp{page 66}\textit{on a}
% note: the label (4.3.1.4) is printed for this display as well as for the preceding one
\[
s\overset{\mathcal M,\mathcal N}{\boxtimes}t\ \simeq\ \varphi(s,t)\ ,
\tag*{(4.3.1.4)}
\]
\textit{et $\varphi(S\times T)=\{s\boxtimes t\mid s\in S,\ t\in T\}$ est une partie strictement génératrice de $\mathcal M\boxtimes\mathcal N$ (}\footnote{\textbf{N.B.} On a bien sûr $(s,t)=s\overset{S^\wedge,T^\wedge}{\boxtimes}t$.}\textit{).}

\noindent\textbf{N.B.} Je ne connais pas d'exemple où les conditions envisagées ne soient satisfaites, mais doute pourtant à présent que même la plus faible le soit toujours, i.e. que $\mathcal M\boxtimes\mathcal N$ soit un pseudo-topos dès que $\mathcal M$ et $\mathcal N$ le sont.

\noindent\textsc{Démonstration de} 4.3.1. L'implication b $\Longrightarrow$ b$'$ est tautologique. L'implication b$'$ $\Longleftrightarrow$ c résulte du

% marginal vertical bar printed beside Lemme 4.3.2 (from "Lemme 4.3.2." to "soit représentable.")
\noindent\textbf{Lemme 4.3.2.} \textit{Soient $\mathcal E$, $X$ des $\mathfrak U$-catégories avec $X$ petite, $\mathcal E$ stable par petites $\varprojlim$ ,}
\[
\alpha\ :\ \mathcal E\longrightarrow X^\wedge
\]
\textit{un foncteur. Pour que $\alpha$ admette un adjoint à gauche $\varphi$, il faut et il suffit que pour tout $x\in\Ob X$, le foncteur}
\[
\left\{
\begin{array}{ccl}
\xi & \longmapsto & \Hom_{X^\wedge}(x,\alpha(\xi))\ =\ \alpha(\xi)(x)\\
\mathcal E^\circ & \longrightarrow & \mathrm{Ens}
\end{array}
\right.
\]
\textit{soit représentable.}

\noindent C'est évidemment nécessaire. La suffisance résulte du fait que $X\subset X^\wedge$ est stricte%
% foot: Version 17 décembre 2025
\printedpage{10}%
% head: A. Grothendieck : Dérivateurs, Ch. XVIII, Édition des Universités de Montpellier II et Paris VII
\msp{page 67}%
ment génératrice dans $X^\wedge$, i.e. que tout objet $F$ de $X^\wedge$ est limite inductive $F=\varinjlim\limits_I F_i$, les $F_i$ dans $X$ (on prend $I=X/F$), d'où
\[
\begin{aligned}
\Hom(F,\alpha(\xi))\ &\simeq\ \varprojlim_I\ \underbrace{\Hom(F_i,\alpha(\xi))}_{\Hom(\rho(F_i),\xi)}\\
&\simeq\ \Hom(\varinjlim_I\,(\rho(F_i),\xi))\ ,
\end{aligned}
\]
donc le foncteur $\xi\longmapsto\Hom(F,\alpha(\xi))$ est représenté par $\rho(F)\overset{\text{déf}}{=}\varinjlim\limits_I\rho(F_i)$.

\noindent Prouvons c $\Longrightarrow$ a. Pour l'existence des petites $\varinjlim$ dans $\mathcal M\boxtimes\mathcal N$, et l'existence d'une partie strictement génératrice, on invoque le lemme général :

% marginal vertical bar printed beside Lemme 4.3.3 (from "Lemme 4.3.3." to "comme")
\noindent\textbf{Lemme 4.3.3.} \textit{Soit $\alpha:\mathcal E\hookrightarrow\mathcal F$ un foncteur pleinement fidèle admettant un adjoint à gauche $\rho$.}
\begin{enumerate}[label={}]
\item[\textcircled{a}] \textit{Si les $\varinjlim$ d'un type donné $I$ existent dans $\mathcal F$, elles existent dans $\mathcal E$, et se calculent comme}
\[
\varinjlim_I{}_{\mathcal E}\,\xi_i\ =\ \rho(\varinjlim_I{}_{\mathcal F}\,\alpha(\xi_i))\ .
\tag*{(4.3.3.1)}
\]
\item[\textcircled{b}] \textit{Si $U$ est une petite partie strictement génératrice de $\mathcal F$, $\rho(U)$ est une petite partie strictement génératrice dans $\mathcal F$.}
\end{enumerate}

\noindent En effet, si $\xi$ désigne le deuxième membre \texttt{[de l'égalité 4.3.3.1]}, et $\eta$ un objet variable de $\mathcal E$, on aura
\[
\begin{aligned}
\Hom_{\mathcal E}(\xi,\eta)\ &\overset{\text{adj.}}{\simeq}\ \Hom_{\mathcal F}(\varinjlim_I\alpha(\xi_i),\alpha(\eta))\\
&\simeq\ \varprojlim_I\quad\underbrace{\Hom_{\mathcal F}(\alpha(\xi_i),\alpha(\eta))}_{\simeq\,\Hom_{\mathcal E}(\xi_i,\eta)\ \ \text{car $\alpha$ pleinement fidèle}}\quad .
\end{aligned}
\]
\msp{page 68}
\[
\simeq\ \ \varprojlim_I\Hom_{\mathcal E}(\xi_i,\eta)\ ,
\]
ce qui prouve que $\xi$ est une limite inductive des $\xi_i$ dans $\mathcal E$ (\footnote{Pour prouver b), on note que la pleine fidélité de $\alpha$ implique que $\rho$ est essentiellement surjectif ; comme $\rho$ commute aux petites $\varinjlim$ et que tout objet de $\mathcal F$ est de la forme $\varprojlim\limits_I x_i$, les $x_i$ dans $U$, on aura $\rho(x)=\varinjlim\limits_I\rho(x_i)$, les $\rho(x_i)$ dans $\rho(U)$, on gagne.}).

\noindent Ceci prouve que (moyennant c)) $\mathcal M\boxtimes\mathcal N$ est stable par petites $\varinjlim$ , donc il reste à prouver que $\mathcal M\boxtimes\mathcal N$ admet une petite famille strictement génératrice. Cela résultera du b) dans le lemme, puisque $S\times T$ engendre strictement $(S\times T)^\wedge$. Cela prouve en même temps la dernière assertion de la proposition, sous réserve de prouver la formule (4.1.3.4). Il faut donc prouver un isomorphisme fonctoriel en $F\in\Ob\mathcal M\boxtimes\mathcal N$, % [sic] (4.1.3.4) -> (4.3.1.4), the formula s ⊠ t ≃ φ(s,t) to be proved is labelled (4.3.1.4) on p. 9
\[
\Hom(\rho(s,t),F)\ \simeq\ F(s,t)\ ,
\]
mais par adjonction on a
\[
\Hom(\rho(s,t),F)\ \simeq\ \Hom((s,t),\alpha(F))\ =\ \alpha(F)(s,t)\ =\ F(s,t)\ ,
\]
OK.

% foot: Version 17 décembre 2025
\printedpage{11}
% head: A. Grothendieck : Dérivateurs, Ch. XVIII, Édition des Universités de Montpellier II et Paris VII

\noindent\msp{page 69}\textbf{Proposition 4.3.4.} \textit{Soit $\mathcal M$ une $\mathfrak U$-catégorie. Les conditions suivantes sont équivalentes.}
\begin{enumerate}[label={}]
\item[\textit{a)}] \textit{$\mathcal M$ est une catégorie $\mathfrak U$-pseudo-topos, i.e. elle est stable par $\mathfrak U$-limites inductives, et admet une $\mathfrak U$-petite sous-catégorie $S$ strictement génératrice.}
\item[\textit{b)}] \textit{Il existe une $\mathfrak U$-petite catégorie $S$, et un foncteur $i:\mathcal M\hookrightarrow S^\wedge$} pleinement fidèle, \textit{ayant un adjoint à gauche.}
\end{enumerate}

\noindent\textsc{Démonstration}. On a a) $\Longrightarrow$ b), car il suffit de prendre pour $S$ dans b) celui spécifié dans a), l'hypothèse sur $S$ implique que le foncteur canonique
\[
\mathcal M\ \longrightarrow\ S^\wedge\ ,\qquad\quad \text{composé de } \mathcal M\hookrightarrow\ \mathcal M^\wedge\ \xrightarrow{\ \text{restriction}\ }S^\wedge
\]
est pleinement fidèle. Qu'il admette un adjoint à gauche résulte du lemme 4.3.2. L'implication b) $\Longrightarrow$ a) résulte du lemme 4.3.3.

\noindent\textbf{Corollaire 4.3.5.} \textit{Soit $\mathcal M$ une $\mathfrak U$-catégorie. Pour que $\mathcal M$ soit un $\mathfrak U$-pseudo-topos} accessible \textit{(i.e. dont tous les objets sont accessibles), il faut et il suffit qu'il existe $S$ dans} Cat \textit{et un foncteur pleinement fidèle et} accessible $\mathcal M\overset{\alpha}{\hookrightarrow}S^\wedge$\textit{, admettant un adjoint à gauche $\rho$.}

\noindent\msp{page 70}\textsc{Démonstration}. Il suffit de noter que si $S$, $\alpha$ sont comme dans 4.3.4 b), alors $\alpha$ est accessible si et seulement si les $\rho(S)$ sont formées d'objets accessibles dans $\mathcal M$ ; comme cette partie est strictement génératrice, on sait que cela implique que \textit{tout} objet de $\mathcal M$ est accessible. (Cf. SGA 4 I 9.9.)

\noindent Je vais essayer de donner un énoncé plus complet :

\noindent\textbf{Théorème 4.3.6.} \textit{Soit $\mathcal M$ une $\mathfrak U$-catégorie. Conditions équivalentes :}
\begin{enumerate}[label={}]
\item[\textit{a)}] \textit{$\mathcal M$ stable par petites $\varinjlim$ , et elle admet une petite sous-catégorie pleine $S$, strictement génératrice, et formée d'objets accessibles de $\mathcal M$.}
\item[\textit{a$'$)}] \textit{$\mathcal M$ est une catégorie $\mathfrak U$-pseudo-topos faiblement accessible, i.e. satisfait a), avec en plus} tout \textit{élément de $\mathcal M$ (pas seulement ceux de $S$) accessible dans $\mathcal M$.}
\item[\textit{b)}] \textit{Il existe une petite catégorie $S$, et un foncteur $\alpha:\mathcal M\hookrightarrow S^\wedge$} pleinement fidèle et accessible, \textit{tel que $\alpha$ admette un adjoint à gauche.}
\item[\textit{c)}] \msp{page 71}\textit{$\mathcal M$ est stable par petites $\varinjlim$ , est accessible, et admet une filtration cardinale} $(\Filt^\pi\mathcal M)_{\pi\geqslant\pi_0}$.
\item[\textit{d)}] \textit{Il existe un cardinal $\pi_1$, une petite catégorie $\mathcal M_{\pi_1}$, stable par limites inductives de cardinal $\leqslant\pi_1$, tels qu'on ait une équivalence de catégories}
\[
\mathcal M\ \xleftarrow{\ \sim\ }\ \Ind_{\pi_1}\mathcal M_{\pi_1}\ ,
\tag*{(4.3.6.1)}
\]
\textit{où le deuxième membre désigne la sous-catégorie pleine de} $\Ind(\mathcal M_{\pi_1})$ \textit{formée des ind-objets $(X_i)_{i\in I}$ de $\mathcal M_{\pi_1}$ tels que $I$ soit grand devant $\pi_1$.}
\end{enumerate}

\noindent\textsc{Démonstration de} 4.3.6 (\footnote{Pour une démonstration complète (celle donnée ici canule à partir d'un certain moment), cf. §4.6 (p. 139-148).}). On a déjà vu a$'$ $\Longrightarrow$ a $\Longrightarrow$ b $\Longrightarrow$ a$'$, i.e. l'équivalence de a, a$'$, b, soit PTA cette condition. Nous allons prouver les implications
\[
\mathrm{PTA}\ \Longrightarrow\ \mathrm{c}\ \Longrightarrow\ \mathrm{d}\ \Longrightarrow\ \mathrm{PTA}\ ,
\]
% foot: Version 17 décembre 2025
\printedpage{12}%
% head: A. Grothendieck : Dérivateurs, Ch. XVIII, Édition des Universités de Montpellier II et Paris VII
ce qui établira le théorème.

\noindent\textit{Prouvons PTA $\Longrightarrow$ c.} Soit $S$ une petite sous-catégorie pleine strictement génératrice. Ainsi, \msp{page 72}pour tout $x$ dans $\mathcal M$ on a
\[
x\ =\ \varinjlim_{s\in S/x}s
\tag*{(4.3.6.2)}
\]
(cf. SGA 4 I 7.2 (i)). Soit $\pi_0=\card\Fl(S)$, et pour tout cardinal $\pi\geqslant\pi_0$, soit
\[
\begin{array}{lcl}
\Filt^\pi\mathcal M & = & \text{sous-catégorie pleine de $\mathcal M$ formée des objets $x$ tels que}\\
 & & \card\,(\Fl(S/x))\leqslant\pi\ .
\end{array}
\tag*{(4.3.6.3)}
\]
Je veux prouver que l'on trouve ainsi une filtration cardinale sur $\mathcal M$, i.e. que les conditions a) b) c) de SGA 4 I 9.12 sont satisfaites. C'est prouvé dans loc. cit. 9.13 sous des hypothèses légèrement différentes, la condition d'accessibilité de $\mathcal M$ étant remplacée par celle (beaucoup moins sympathique!) que les familles de flèches épimorphiques strictes dans $\mathcal M$ sont épimorphiques strictes universelles. Cette condition n'est utilisée que dans la démonstration de la condition b), à savoir la stabilité de $\Filt^\pi\mathcal M$ par $\varinjlim\limits_I$ de cardinal $\leqslant\pi$ (sans avoir d'ailleurs à supposer $I$ filtrante, ni a fortiori grand devant $\pi_0$). Il faut d'ailleurs dire que la définition de $\Filt^\pi\mathcal M$ donnée dans loc. cit. n'est pas tout à fait la même que (4.3.6.3), \msp{page 73}mais c'est
% the printed line break "car-/dinal" inside the display is kept as printed
\[
\begin{array}{lcl}
\Filt'^{\pi}\mathcal M & = & \text{sous-catégorie pleine de $\mathcal M$ formée des objets $x$ tels qu'il}\\
 & & \text{existe une famille épimorphique stricte $x_i\xrightarrow{\alpha_i}x$ de car-}\\
 & & \text{dinal $\leqslant\pi$, avec $x_i\in S\ \forall\ i\in I$.}
\end{array}
\tag*{(4.3.6.3$'$)}
\]
Ainsi on a
\[
\Filt^\pi\mathcal M\ \subset\ \Filt'^{\pi}\mathcal M\ ,
\]
mais dans loc. cit. on prouve que pour $x$ dans $\Filt'^{\pi}\mathcal M$, on a
\[
\card\Fl(S/x)\ \leqslant\ \pi^{\pi_0}\ ,
\]
d'où
\[
\Filt^\pi\mathcal M\ \subset\ \Filt'^{\pi}\mathcal M\ \subset\ \Filt^{\pi^{\pi_0}}\mathcal M\ ,
\]
donc si $\pi=\pi^{\pi_0}$, p.ex. $\pi=2^c$ avec $c\geqslant\pi_0$, alors on trouve
\[
\Filt^\pi\mathcal M\ =\ \Filt'^{\pi}\mathcal M\ .
\]

\noindent Donc les deux filtrations on l'air $\pm$ équivalentes. Mais la démonstration donnée de la stabilité de $\Filt'^{\pi}$ semble vraiment faire appel de façon essentielle à la condition d'universalité des familles épimorphiques strictes (donc épimorphiques effectives, avec les hypothèses faites impliquant l'existence de petites $\varprojlim$ ). Donc il me faut bon gré, mal gré, refaire le travail! % [sic] on l'air -> ont l'air, verb avoir (3rd person plural) intended

% foot: Version 17 décembre 2025
\printedpage{13}
% head: A. Grothendieck : Dérivateurs, Ch. XVIII, Édition des Universités de Montpellier II et Paris VII

\noindent\msp{page 74}Condition a) : $\Filt^\pi$ équivalente à une petite catégorie. C'est immédiat, puisque par (4.3.6.2), $\Filt^\pi$ est formée de $\varprojlim\limits_I$ dans $\mathcal M$ d'objets de $S$ (petite catégorie), avec $I$ telle que $\card\Fl I\leqslant\pi$.

\noindent Condition b) : $\Filt^\pi$ stable par $\varprojlim\limits_I$ , quand $\card\Fl I\leqslant\pi$ (\footnote{quand $I$ filtrante et grande devant $\pi_0$ -- je ne le prouve que dans ce cas, hélas!}). Soit donc $(x_i)_{i\in I}$, où
\[
\begin{array}{ccc}
I & \longrightarrow & \mathcal M\\[1ex]
i & \longmapsto & x_i
\end{array}
\]
un système inductif dans $\mathcal M$, avec des $x_i$ dans $\Filt^\pi$. Considérons le produit fibré
\[
\begin{tikzcd}
I \arrow[d, "\varphi"'] & S/I = J \arrow[l, "\text{cofibré}"'] \arrow[d] \arrow[dr, dotted, "\psi"] & & \\
\mathcal M & S/\mathcal M \arrow[l, "\text{but}"', "\text{Cat-cofibré}"] \arrow[r, "\text{source}"] & S \arrow[r, hook] & \mathcal M\ ,
\end{tikzcd}
\tag*{(4.3.6.4)}
\]
où $J$ est Cat-cofibré sur $I$, sa fibre en $i\in\Ob I$ étant isomorphe à $S/x_i$, donc $\card\Fl J_i\leqslant\pi$. Comme par hypothèse $\card\Fl I\leqslant\pi$, on en conclut que $\card\Fl J\leqslant\pi^3=\pi$. (\textbf{N.B.} L'ensemble des flèches $v:j_0\longrightarrow j_1$ de $J$ au dessus d'une flèche $u:i_0\longrightarrow i_1$ de $I$ est “contenu” dans l'ensemble produit $\Ob J_{i_0}\times\Fl J_{i_1}$, par l'application
\[
v\longmapsto(\operatorname{source}v,v':u_*(j_0)\longrightarrow j_1)\ ,
\]
donc de cardinal majoré par $\pi^2$.) On \msp{page 75}a alors
\[
x\ \overset{\text{déf}}{=}\ \varinjlim_I x_i\ =\ \varinjlim_J\ \underbrace{\xi_j}_{=\,\psi(j)}\qquad\quad\mbox{(\footnotemark)}\ ,
\]
\footnotetext{On utilise la propriété de distributivité des $\varinjlim$ , relative à un foncteur Cat-cofibrant
\[
\begin{tikzcd}[ampersand replacement=\&]
J \arrow[r] \arrow[dr] \& I\\
\& \mathcal M\ ,
\end{tikzcd}
\]
explicitée dessus.}%
où cette fois la limite est une limite d'objets de $S$, prise dans $\mathcal M$. Il s'ensuit que $x\in\Filt'^{\pi}$, ce n'est pas tout à fait $x\in\Filt^\pi$. (Ça implique $x\in\Filt^{\pi^{\pi_0}}$, donc $x\in\Filt^\pi$ si $\pi=\pi^{\pi_0}$, p.ex. $\pi=2^c$ avec $c\geqslant\pi_0$.) Mais c'est le moment d'utiliser le fait que les objets de $S$ sont accessibles, et on prendra $\pi_0$ tel que les objets de $S$ soient $\pi_0$-accessibles, et cette fois on prendra quand même $I$ filtrante et grande devant $\pi_0$ (pour nous borner à ce qui est exigé dans la condition b) de loc. cit.). On doit prouver que $\card\,(\Fl(S/x))\leqslant\pi$ (ce qui signifie $x\in\Filt^\pi$). Mais comme les objets de $S$ sont $\pi_0$-accessibles et $I$ grande devant $\pi_0$, on voit aussitôt que
\[
S/x\ \simeq\ \varinjlim_i S/x_i\ ,
\]
% foot: Version 17 décembre 2025


\endgroup
% END UNIT der-ch18/der-ch18_p011-015.tex
% BEGIN UNIT der-ch18/der-ch18_p016-020.tex
\begingroup\makeatletter%
\newcommand{\printedpage}[1]{\marginpar{\footnotesize p.~#1}}%
\newcommand{\msp}[1]{\marginpar{\footnotesize [#1]}}%
\newcommand{\uHom}{\underline{\mathrm{Hom}}}%
\newcommand{\Ind}{\mathrm{Ind}}%
\newcommand{\Hom}{\mathrm{Hom}}%
\newcommand{\card}{\mathrm{card}}%
\newcommand{\Fl}{\mathrm{Fl}}%
\newcommand{\Ob}{\mathrm{Ob}}%
\newcommand{\Filt}{\mathrm{Filt}}%
\newcommand{\hookxrightarrow}[2][]{\lhook\joinrel\xrightarrow[#1]{#2}}%
\makeatother\let\printedpage\derpp\let\msp\dermsp%
\let\mspin\DERmspin%

\setcounter{footnote}{11}

\printedpage{14}
% head: A. Grothendieck : Dérivateurs, Ch. XVIII, Édition des Universités de Montpellier II et Paris VII

d'où
\[
\Fl\,S/x \;\simeq\; \varinjlim_{i} \Fl\,S/x_i \;,
\]
donc
\[
\card\Fl(S/x) \;\leqslant\; (\underbrace{\card I}_{\leqslant\pi})\underbrace{(\sup_{i}\card\Fl(S/x_i))}_{\leqslant\pi} \;\leqslant \pi^2 = \pi \;,
\]

\hfill q.e.d.

\msp{page 76}
% margin: hand-drawn bracket mark in the left margin beside the next two lines
Prouvons c) (\footnote{\textit{Démonstration canulée}, ne marche que si $\pi^{\pi_0} = \pi$\,!!}), i.e. que pour tout $\pi' > \pi \geqslant \pi_0$ et tout $x \in \Ob\Filt^{\pi'}$, on peut écrire
\[
x \;=\; \varinjlim_{I} x_i \qquad I \text{ grande devant } \pi\text{, les } x_i \text{ dans } \Filt^{\pi}\text{, } \card I \leqslant \pi'^{\pi}. \tag*{(4.3.6.5)}
\]
% margin: hand-drawn bracket mark in the left margin beside the next three lines
La démonstration se fait comme dans loc. cit. (p. 151-152), en changeant légèrement la définition de $I$, comme désignant l'ensemble des sous-catégories pleines $i$ de $S/x$ qui sont filtrantes, de cardinal $\leqslant \pi$ \textit{et de plus grandes devant} $\pi_0$ (\footnote{Il n'est pas sûr que $I$ soit grande devant $\pi$, sauf si on suppose $\pi^{\pi_0} = \pi$\,!}). (On pourra supposer $S/x$ filtrante en augmentant $S$ de façon à ce que $S$ soit stable dans $\mathcal{M}$ par petites $\varinjlim$ finies, donc les catégories $S/x$ aussi, ce qui implique qu'elles sont filtrantes. Il est \texttt{[phrase incomplète]})

On a donc prouvé PTA $\Longrightarrow$ c).

\textit{Prouvons c) $\Longrightarrow$ d).} Soit
\[
\pi_1 \;=\; 2^c \qquad \text{avec} \qquad c \;\geqslant\; \pi_0 \;.
\]
On sait alors que
\[
\Filt^{\pi_1}\mathcal{M} \;=\; \mathcal{M}_{\pi_1} \;\overset{\text{déf}}{=}\; \text{sous-catégorie pleine de $\mathcal{M}$ des objets $\pi_1$-accessibles de $\mathcal{M}$}
\]
(loc. cit. 9.16), ce qui implique notamment que $\mathcal{M}_{\pi_1}$ est stable par $\varinjlim$ de cardinal $\leqslant \pi_1$, en vertu de loc. cit. 9.9 (page 145). De plus,
\msp{page 77}
par loc. cit. 9.18 b) (p. 156), on a
\[
\Ind_{\pi_1}(\mathcal{M}_{\pi_1}) \xrightarrow{\;\sim\;} \mathcal{M} \;,
\]
ce qui prouve d).

\textit{Prouvons d) $\Longrightarrow$ PTA.} Par loc. cit. 9.18 a) (p. 156), on voit que les objets de $\mathcal{M}_{\pi_1}$ sont $\pi_1$-accessibles dans $\Ind_{\pi_1}(\mathcal{M}_{\pi_1})$, d'autre part il est clair que $\mathcal{M}_{\pi_1}$ est une sous-catégorie pleine strictement génératrice de $\Ind_{\pi_1}(\mathcal{M}_{\pi_1})$, donc il reste à voir que $\Ind_{\pi_1}(\mathcal{M}_{\pi_1})$ est stable par petites $\varinjlim$ , i.e. le

\textbf{Lemme 4.3.7.} \textit{Soit $\mathcal{M}_\pi$ une $\mathfrak{U}$-catégorie stable par limites inductives de cardinal $\leqslant \pi$, où $\pi$ est un cardinal infini donné. Alors} $\Ind_\pi\mathcal{M}_\pi$ \textit{est stable par petites limites inductives. De plus, $\mathcal{M} = \Ind_\pi\mathcal{M}_\pi$, munie de $\mathcal{M}_\pi \longrightarrow \mathcal{M}$, est solution d'un problème $2$-universel, à savoir
\printedpage{15}
% head: A. Grothendieck : Dérivateurs, Ch. XVIII, Édition des Universités de Montpellier II et Paris VII
que pour toute catégorie $\mathcal{N}$ stable par petites limites inductives, le foncteur restriction induit une équivalence}
\[
\uHom_{!}(\mathcal{M},\mathcal{N}) \xrightarrow{\;\sim\;} \uHom_{!\,\pi}(\mathcal{M}_\pi,\mathcal{N}) \;, \tag*{(4.3.7.1)}
\]
\textit{où le signe ! (resp. $!\,\pi$) désigne la sous-catégorie pleine de la catégorie} $\uHom$\textit{, formée des foncteurs qui commutent aux petites} $\varinjlim$ \textit{(resp. aux limites inductives de cardinal $\leqslant \pi$).}

\msp{page 78}
Prouvons d'abord que les $\varinjlim$ de cardinal $\leqslant \pi$ \textit{dans} $\mathcal{M}_\pi$ sont aussi des $\varinjlim$ \textit{dans} $\mathcal{M} = \Ind_\pi\mathcal{M}_\pi$ (\footnote{Mais \textit{attention}, $\mathcal{M}_\pi \longrightarrow \Ind(\mathcal{M}_\pi)$ ne commute \textit{pas} à ces $\varinjlim$ !}). On a donc
\[
x \xleftarrow{\;\sim\;} \varinjlim_{I} x_i \qquad \textit{dans } \mathcal{M}_\pi \;,
\]
i.e.
\[
\Hom(x,y) \xrightarrow{\;\sim\;} \varprojlim_{I} \Hom(x_i,y) \qquad \text{si } y \in \Ob\mathcal{M}_\pi \;, \tag*{$(*)$}
\]
prouvons que cela reste vrai pour $y$ dans $\mathcal{M}$, i.e. $y$ de la forme
\[
y \;=\; \varinjlim_{J} y_j \;, \qquad J \text{ grande devant } \pi\text{, les } y_j \text{ dans } \mathcal{M}_\pi.
\]
On a en effet
\[
\begin{aligned}
\Hom(x,y) \quad &\simeq \quad \varinjlim_{J} \Hom(x,y_j)\\
&\overset{(*)}{\simeq} \quad \varinjlim_{J}\, \varprojlim_{I} \Hom(x_i,y_j)\\
&\overset{\text{SGA 4 I 9.8}}{\simeq} \quad \varprojlim_{I}\, \varinjlim_{J} \Hom(x_i,y_j)\\
&\simeq \quad \varprojlim_{I} \Hom(x_i,y) \;,
\end{aligned}
\]
au total
\[
\Hom(x,y) \;\simeq\; \varprojlim_{I} \Hom(x_i,y),
\]
q.e.d.

Pour prouver que les petites $\varinjlim$ sont représentables dans $\mathcal{M} = \Ind_\pi(\mathcal{M}_\pi)$, comme elles le sont dans $\Ind(\mathcal{M}_\pi)$ (SGA 4 I 8.9.5 b, on utilise l'hypothèse que les $\varinjlim$ \textit{finies} sont représentables dans $\mathcal{M}_\pi$), il reste à voir que $\Ind_\pi(\mathcal{M}_\pi)$ est stable dans $\mathcal{M}_\pi$ par ces limites.

\msp{page 79}

\section*{4.4\quad Étude de $\Ind_\pi(C)$ (pour $C$ petite catégorie stable par $\varinjlim$ de cardinal $\leqslant \pi$)}

\textbf{4.4.1.} Finalement, la démonstration du lemme 4.3.7 n'est pas tellement évidente que je l'anticipais, et il me faut lui consacrer des préliminaires techniques, que je vais réunir ici. Dans toute la suite, $C$ représente une $\mathfrak{U}$-catégorie équivalente à une petite catégorie.

\printedpage{16}
% head: A. Grothendieck : Dérivateurs, Ch. XVIII, Édition des Universités de Montpellier II et Paris VII

\textbf{Proposition 4.4.2.} \textit{Soit $F \in C^{\wedge}$ un foncteur ind-représentable, i.e. tel que $C/F$ soit filtrante. Soit d'autre part $i \longmapsto x_i$ un foncteur $I \longrightarrow C$, où $I$ est une petite catégorie filtrante, et soit}
\[
\alpha \;:\; x \;=\; \varinjlim_{I}{}_{C^{\wedge}} x_i \longrightarrow F \qquad \text{(\footnotemark)} \tag*{(4.4.2.1)}
\]
\footnotetext{La donnée de $\alpha$ équivaut aussi à celle d'un homomorphisme d'ind-objets $(x_i)_{i\in I} \longrightarrow F$, $F$ étant considéré comme ind-objet (généralisé) indexé par $C/F$.}
\textit{un système projectif de flèches $x_i \longrightarrow F$ (}\textbf{N.B.} \textit{la limite inductive $x$ est prise dans $C^{\wedge}$, et est calculée argument par argument), ce qui équivaut à une factorisation $\varphi_\alpha : I \longrightarrow C/F$ de $I \longrightarrow C$ en}
\[
I \xrightarrow{\;\varphi_\alpha\;} C/F \hookxrightarrow[\psi]{\textit{can.}} C \;.
\]
\textit{Ceci dit, pour que $\alpha$ (4.4.2.1) soit un isomorphisme dans} $\Ind'(C) \subset C^{\wedge}$ \textit{(sous-catégorie pleine des foncteurs ind-représentables sur $C$), il faut et il suffit que $\varphi_\alpha$ soit} cofinal\textit{, ou ce qui revient au même, $I$ étant filtrante, qu'il satisfasse les deux conditions habituelles :}
\begin{itemize}[leftmargin=2.5em]
\item[F\,1)] $\forall\, \xi \in \Ob C/F$, $\exists\, i \in \Ob I$ \textit{et une flèche} $\xi \longrightarrow \varphi_\alpha(i)$.
\item[F\,2)] $\forall\, \xi \in \Ob C/F$, \textit{et} $i \in \Ob I$, \textit{et toute double-flèche}

\msp{page 80}
\[
\xi \rightrightarrows \varphi_\alpha(i) \;,
\]
\textit{il existe une flèche $i \xrightarrow{\;u\;} j$ dans $I$ telle que $\varphi_\alpha(u)$ égalise la double-flèche.}
\end{itemize}
\textsc{Démonstration}. Si $\varphi_\alpha$ est cofinal, on en conclut par définition que $\varinjlim \varphi_\alpha \longrightarrow \varinjlim \psi\varphi_\alpha$ est un isomorphisme, or cette flèche n'est autre que $\alpha$. Inversement, supposons que cette flèche soit un isomorphisme, donc $C/F \xleftarrow{\;\sim\;} C/\varinjlim x_i$, et il reste à voir que le foncteur canonique
\[
I \;\longrightarrow\; C/\varinjlim{}_{C^{\wedge}} x_i
\]
est cofinal, i.e. satisfait F\,1) et F\,2). La vérification est $\pm$ tautologique, en utilisant le fait que pour tout $y$ dans $C$, on a
\[
\Hom(y,\, \varinjlim{}_{C^{\wedge}} x_i) \;\xleftarrow{\;\sim\;}\; \varinjlim{}_{\mathrm{Ens}}\Hom(y,x_i),
\]
et du “calcul” des limites inductives filtrantes dans Ens.

% margin: vertical bar in the left margin beside Corollaire 4.4.3
\textbf{Corollaire 4.4.3.} \textit{Soit $\pi$ un ordinal infini. Pour que $F$ soit dans} $\Ind'_\pi(C)$ \textit{(image essentielle de} $\Ind_\pi(C)$ \textit{dans $C^{\wedge}$), il faut et il suffit qu'il existe une catégorie ordonnée grande devant $\pi$ $I$, et un foncteur cofinal $I \longrightarrow C/F$.}

Il nous faut donc donner des conditions, sur une catégorie $J$ $(= C/F)$, pour qu'il existe une catégorie ordonnée $I$ grande devant $\pi$, et un foncteur cofinal $I \longrightarrow J$.

\msp{page 81}
Une telle catégorie s'appellera \textit{grande devant} $\pi$. (Il est immédiat que dans le cas où $J$ est elle-même ordonnée, cette terminologie est cohérente avec celle déjà adoptée pour les ensembles ordonnés, i.e. que s'il existe une application croissante cofinale $\varphi : I \longrightarrow J$, avec $I$ grand devant $\pi$, alors $J$ est lui-même grand devant $\pi$.) Une telle catégorie est
\printedpage{17}
% head: A. Grothendieck : Dérivateurs, Ch. XVIII, Édition des Universités de Montpellier II et Paris VII
automatiquement filtrante (SGA 4 I 8.1.3 b)). Notons aussi que si on a un foncteur cofinal $J \longrightarrow J'$, si $J$ est grande devant $\pi$, $J'$ l'est aussi (par transitivité de la cofinalité).

\textbf{Proposition 4.4.4.} \textit{Soient $J$ une $\mathfrak{U}$-catégorie, $\pi$ un cardinal infini. Conditions équivalentes :}
\begin{itemize}[leftmargin=4em]
\item[\textit{a)}] \textit{$J$ est grande devant $\pi$, i.e. il existe un ensemble ordonné grand devant $\pi$ et un foncteur cofinal $\varphi : I \longrightarrow J$.}
\item[\textit{b)}] \textit{$J$ satisfait les deux conditions suivantes :}
\begin{itemize}[leftmargin=4em]
\item[$\mathrm{PF}_{\pi}\,1$] $\mathrm{Ord}\,J$ \textit{est grande devant $\pi$, i.e. pour toute famille $(y_\alpha)_{\alpha\in A}$ d'objets de $J$, avec} $\card A \leqslant \pi$\textit{, il existe un majorant $y$ dans $J$ (i.e. un $y$ et des flèches $y_\alpha \xrightarrow{\;v_\alpha\;} y$).}
\item[$\mathrm{PF}_{\pi}\,2$] \textit{Pour $y_0$, $y_1$ dans $J$ et tout ensemble $(u_\alpha)_{\alpha\in A}$ de flèches $y_0 \xrightarrow{\;u_\alpha\;} y_1$ avec} $\card A \leqslant \pi$\textit{, il existe une flèche $y_1 \xrightarrow{\;u'\;} y_2$ qui les égalise, i.e. telle que $u'u_\alpha = u'u_\beta \;\forall\, \alpha,\beta \in A$. (Il suffit de l'exiger quand} $\card A \geqslant 2$.\textit{)}
\end{itemize}
\item[\textit{c)}] \textit{Pour toute partie $A$ de} $\Fl J$\textit{, avec} $\card A \leqslant \pi$\textit{, existe une sous-catégorie $J'$ de $J$, avec} $\card\Fl J' \leqslant \pi$\textit{, telle que $A \subset \Fl J'$ et que $J'$ admette un objet final.}
\item[\textit{d)}] \textit{Soit $I$ l'ensemble des sous-catégories $J'$ de $J$ telles que} $\card J' \leqslant \pi$ \textit{et ayant un objet final}

\msp{page 82}
\textit{et ordonnons $I$ par inclusion. Alors $I$ est grand devant $\pi$, de plus le foncteur canonique $I \longrightarrow J$ $(J' \longmapsto e_{J'} =$ objet final de $J')$ est cofinal} (\footnote{Introduire $I(J)$ avant l'énoncé de la proposition. \textbf{N.B.} La condition de cofinalité $I(J) \longrightarrow J$ n'est pas automatique, comme on voit en prenant $J = B_G$, $G = \{\pm 1\}$.},\footnote{\textbf{N.B.} L'hypothèse de cofinalité est superflue si $J$ est rigide, i.e. tout endomorphisme d'un objet de $J$ est identique.},\footnote{\ N. Éd.\ \ Dans cette condition (d), le “\textit{de plus le foncteur \dots\ est cofinal}” est ajouté par Grothendieck \textit{a posteriori}, ce qui explique sa vaine tentative de l'établir dans l'implication d $\Longrightarrow$ a ci-dessous, implication qui devient évidente après cet ajout. Par ailleurs, les conditions (c) et (d) ne sont \textit{pas} impliquées par les conditions équivalentes (a) et (b). Par exemple, si $J$ est la catégorie ayant un seul objet et un seul morphisme non identique $p$ tel que $p^2 = p$, alors $J$ satisfait aux conditions (a) et (b), mais \textit{pas} aux conditions (c) et (d).}).
\end{itemize}
(\footnote{Autres conditions :
\begin{itemize}[leftmargin=3em]
\item[e)] Pour tout ensemble $A$ avec $\card A \leqslant \pi$, $\mathrm{diag}_J : J \longrightarrow J^A$ est cofinal.
\item[e$'$)] (Si $A_0$ est un ensemble donné, de cardinal $= \pi$.) $J \longrightarrow J^{A_0}$ est cofinal.
\end{itemize}
Voir aussi 4.9.6.9.1 et 4.9.6.9.2, pages 237, 238.})

\textsc{Démonstration}.
\begin{itemize}[leftmargin=7em,labelwidth=6em,align=left]
\item[a $\;\Longrightarrow\;$ b] vérification immédiate, en utilisant le critère F 1, F 2 de la cofinalité pour $I$ filtrante.
\item[b $\;\Longrightarrow\;$ c] par AQT \texttt{[âne qui trotte]}.
\item[c $\;\Longrightarrow\;$ d] par AQT \texttt{[âne qui trotte]}.
\end{itemize}
% foot: Version 17 décembre 2025

\printedpage{18}
% head: A. Grothendieck : Dérivateurs, Ch. XVIII, Édition des Universités de Montpellier II et Paris VII
\begin{itemize}[leftmargin=7em,labelwidth=6em,align=left]
\item[d $\;\Longrightarrow\;$ a] Considérons l'ensemble ordonné grand devant $\pi$ construit dans l'énoncé de d), et le \textit{foncteur}
\[
I \;\longrightarrow\; J
\]
donné par $J' \longmapsto e_{J'}$ $=$ objet final de $J'$. (Vérification que c'est un foncteur immédiate.) Ce foncteur est surjectif, donc satisfait la condition F 1 de cofinalité. La condition F 2 résulte aussi de d) en regardant les trois objets $J'$, $j \xrightarrow{\;u_1\;} e_{J'}$, $j \xrightarrow{\;u_2\;} e_{J'}$ de $I$,
\[
\begin{tikzcd}[column sep=small, row sep=small]
J', & e_{J'} \\
j \arrow[ur, shift left, "u_1"] \arrow[ur, shift right, "u_2"'] &
\end{tikzcd}
\]
et en prenant un $J''$ qui les majore tous trois, alors $e_{J'} \longrightarrow e_{J''}$ égalise $u_1$, $u_2$ (\footnote{Un peu bref, le cas où $e_{J'} = j$ échappe à cet argument, et demande qu'on \texttt{[phrase incomplète]}}).
\end{itemize}

% margin: vertical bar in the left margin beside Corollaire 4.4.5
\textbf{Corollaire 4.4.5.} \textit{Soit $F \in \Ob C^{\wedge}$. Pour que $F$ appartienne à} $\Ind'_\pi(C)$\textit{, il faut et il suffit que $C/F$ satisfasse à l'ensemble des quatre conditions équivalentes de 4.4.4.}

\textbf{Corollaire 4.4.6.} \textit{Soit $J$ une catégorie stable par} $\varinjlim$ \textit{de cardinal $\leqslant \pi$. Alors $J$ est grande devant $\pi$.}

En effet, les conditions $\mathrm{PF}_{\pi}\,1$, $\mathrm{PF}_{\pi}\,2$ sont visiblement vérifiées.

\msp{page 83}
\textbf{4.4.6.} Je me donne un foncteur % [sic] 4.4.6. -> 4.4.7. (?), numéro déjà employé par le Corollaire 4.4.6 qui précède
\[
J \xrightarrow{\;\varphi\;} \Ind'_\pi C \;, \tag*{(4.4.6.1)}
\]
$J$ une petite catégorie, et on cherche une $\varinjlim$ dans $\Ind'_\pi C$, en postulant au besoin l'existence des $\varinjlim$ de cardinal $\leqslant \pi$ dans $C$, et l'hypothèse $\card\Fl J \leqslant \pi$. Le procédé le plus naturel pour ramener la question de l'existence des $\varinjlim$ dans $\Ind'_\pi C$, à la question similaire dans $C$, est explicité dans SGA 4 I 8.8.3. On considère l'inclusion pleinement fidèle
\[
C \;\hookxrightarrow{\;\alpha\;}\; \Ind'_\pi C \;, \tag*{(4.4.6.2)}
\]
laquelle commute aux $\varinjlim$ de cardinal $\leqslant \pi$ (celles qui sont représentables dans $C$, cf. p. 78 et référence à SGA 4 I 9.8), d'où
\[
\uHom(J,C) \;\hookxrightarrow{\;\alpha\;}\; \uHom(J,\Ind'_\pi C) \tag*{(4.4.6.3)}
\]
également pleinement fidèle, et commutant aux $\varinjlim$ de cardinal $\leqslant \pi$ (représentables dans le premier membre). Cette flèche s'insère dans le carré commutatif relatif à la donnée d'un $j \in \Ob J$, et des foncteurs d'inclusion
\[
e \xrightarrow{\;\varepsilon_j\;} J \;, \qquad e_j(e) \;=\; j \;; \tag*{(4.4.6.4)} % [sic] e_j(e) -> \varepsilon_j(e), la flèche est notée \varepsilon_j
\]
% foot: Version 17 décembre 2025


\endgroup
% END UNIT der-ch18/der-ch18_p016-020.tex
% BEGIN UNIT der-ch18/der-ch18_p021-025.tex
\begingroup\makeatletter%
\newcommand{\printedpage}[1]{\marginpar{\footnotesize p.~#1}}%
\newcommand{\msp}[1]{\marginpar{\footnotesize [#1]}}%
\newcommand{\uHom}{\underline{\mathrm{Hom}}}%
\newcommand{\Hom}{\mathrm{Hom}}%
\newcommand{\Ind}{\mathrm{Ind}}%
\newcommand{\Ob}{\operatorname{Ob}}%
\newcommand{\Fl}{\operatorname{Fl}}%
\newcommand{\card}{\operatorname{card}}%
\makeatother\let\printedpage\derpp\let\msp\dermsp%
\let\mspin\DERmspin%


\printedpage{19}
% head: A. Grothendieck : Dérivateurs, Ch. XVIII, Édition des Universités de Montpellier II et Paris VII
% foot: Version 17 décembre 2025

\msp{page 84}
\[
\begin{tikzcd}
\uHom(J,C) \arrow[r, "\alpha^J"] \arrow[d, "\varepsilon_j^*"] & \uHom(J,\Ind'_\pi C) \arrow[d, "{\varepsilon_j^{*\prime}}"] & \ni \varphi \arrow[d, mapsto] \\
C \arrow[r, "\alpha"] & \Ind'_\pi(C) & \ni \varphi(j)\ ,
\end{tikzcd}
\qquad (\text{\footnotemark[21]})
\tag*{(4.4.6.2)}
\]
\footnotetext[21]{Les foncteurs $\varepsilon_j^*$, $\varepsilon_j^{*\prime}$ sont les foncteurs évaluation en $j$.}
lequel induit un foncteur
\[
\underbrace{\uHom(J/C)/\varphi}_{\overset{\text{déf}}{=}\, I} \xrightarrow{\varepsilon_j^\varphi} C/\varphi(j)\ . \tag*{(4.4.6.3)} % [sic] \uHom(J/C) -> \uHom(J,C), everywhere else written Hom(J,C)
\]
Si $C$ est stable par $\varinjlim$ relatives à un ensemble $K$ du type de conditions $\leqslant \pi$ (p.ex. $\varinjlim$ finies), il en est de même de $\uHom(J,C)$, et aussi de $\uHom(J,C)/\varphi$, et de $C/\varphi(j)$, en vertu du fait que $\alpha^J$ et $\alpha$ commutent aux $\varinjlim$ du type envisagé, de plus (4.4.6.3) commute aux limites en question. Supposons donc $C$ stable tout au moins par $\varinjlim$ finies, donc les catégories source et but dans (4.4.6.3) sont filtrantes. (Le fait que le foncteur commute aux $\varinjlim$ finies ne semble pas utile dans le contexte présent.) Cela nous servira pour avoir le critère commode de cofinalité SGA 4 I 8.1.3 b) par F\,1, F\,2.

\msp{page 85}
Rappelons que la cofinalité des foncteurs (4.4.6.3), quand elle a lieu ( ?), implique le fait que $\varphi$ est une limite inductive du foncteur composé
\[
\begin{tikzcd}[row sep=small]
I \arrow[r, "\psi_J"] \arrow[d, equal] & \uHom(J,C) \arrow[r, "\alpha^J"] \arrow[dd, dotted, "\varepsilon_j^*"] & \uHom(J,\Ind'_\pi(C)) \arrow[dd, dotted, "{\varepsilon_j^{*\prime}}"] \\
\uHom(J,C)/\varphi \arrow[d, "\varepsilon_j^\varphi"'] & & \\
C/\varphi(j) \arrow[r, dotted, "\psi_j"] & C \arrow[r, dotted, "\alpha"] & \Ind'_\pi(C)\ ,
\end{tikzcd}
\tag*{(4.4.6.4)}
\]
suivant l'ensemble ordonné filtrant $I$, plus précisément que $\varphi$ est la $\varinjlim$ “argument par argument en $j$”, ce qui signifie que pour tout $j \in \Ob J$, $\varphi(j) = \varepsilon_j^*(\varphi)$ est limite inductive de $\varepsilon_j^*(\alpha^J\psi) = (\alpha\psi_j)\varepsilon_j^\varphi$. (Or on sait en effet que $\varphi(j)$ est la limite de $\alpha\psi$, puisque $\varphi(j) \in \Ob C^\wedge$ est ind-représentable, donc aussi celle de $(\alpha\psi_j)\varepsilon_j^\varphi$ si $\varepsilon_j^\varphi$ est cofinal.)

Ainsi, admettant la cofinalité de (4.4.6.3) et interprétant $\psi_J$ dans (4.4.6.4)

\msp{page 86}
\[
\psi_J \ :\ I \longrightarrow \uHom(J,C)
\]
comme un foncteur
\[
I \times J \xrightarrow{\Psi} C\ , % [sic] \Psi -> \psi, the functor I x J -> C is called \psi in the next formula and in (4.4.7.1)
\]
on trouve la “représentation indicielle” de $\varphi$
\[
\varphi(j) \ =\ \underbrace{\varinjlim_{i\in I} \psi(i,j)}_{\substack{\text{limite inductive}\\ \text{prise dans } \Ind'_\pi(C)\\ \text{d'objets de } C \hookrightarrow \Ind'_\pi(C)}} \qquad ,
\]

\printedpage{20}
% head: A. Grothendieck : Dérivateurs, Ch. XVIII, Édition des Universités de Montpellier II et Paris VII
% foot: Version 17 décembre 2025
d'où
\[
\begin{aligned}
\underbrace{\varinjlim_{j\in J}{}_{\,\Ind'_\pi C}\varphi(j)}_{\text{sous réserve d'existence}} \quad &= \quad \varinjlim_{j\in J}\, \varinjlim_{i\in I} \psi(i,j)\\
&= \quad \varinjlim_{i\in I}\, \varinjlim_{j\in J} \psi(i,j)\ ,
\end{aligned}
\]
toutes les limites inductives étant prises dans $\Ind'_\pi(C)$. Mais si les $\varinjlim$ de type $J$ sont représentables dans $C$, et $\card J \leqslant \pi$, on sait que la $\varinjlim\limits_{j\in J} \psi(i,j)$ \textit{dans} $C$ est aussi leur $\varinjlim$ dans $\Ind'_\pi C$, et on trouve donc que $\varinjlim\limits_{j\in J} \varphi(j)$ existe dans $\Ind'_\pi(C)$ si et seulement si $\varinjlim\limits_{i\in I} \Psi(i)$ existe dans $\Ind'_\pi(C)$, où pour $i \in I$, on pose
\[
\Psi(i) \ =\ \varinjlim_{j\in J} \psi(i,j) \qquad \textit{dans } C.
\]

\msp{page 87}
Résumons ce que nous avons obtenu :

\textbf{Lemme 4.4.7.} \textit{Soit $C$ une petite catégorie, $\pi$ un cardinal infini, $J$ une catégorie telle que $\card \Fl J \leqslant \pi$, $\varphi : J \longrightarrow \Ind'_\pi(C)$ un foncteur. Supposons les $\varinjlim$ de type $J$ existant dans $C$, et que les foncteurs $\varepsilon_j^\varphi$ ($j \in \Ob J$) de (4.4.6.3) soient \textup{cofinaux} (ce qui s'explicite de façon simple surtout si source et but sont filtrants, chose acquise si $C$ est stable par $\varinjlim$ finies). Alors $\varphi$ a une “représentation indicielle canonique”}
\[
\psi \ :\ I \times J \longrightarrow C \tag*{(4.4.7.1)}
\]
\textit{avec}
\[
I \ =\ \uHom(J,C)/\varphi\ , \tag*{(4.4.7.2)}
\]
\textit{i.e. on a l'isomorphisme fonctoriel en $j \in \Ob J$}
\[
\varphi(j) \ \simeq\ \varinjlim_{I} \psi(i,j) \qquad (\text{\footnotemark[22]})\ . \tag*{(4.4.7.3)}
\]
\footnotetext[22]{limite inductive dans $\Ind'_\pi$ d'objets de $C$.}
\textit{Posons pour tout $i \in \Ob I$}
\[
\Psi(i) \ =\ \varinjlim_{J} \psi(i,j) \qquad (\text{\footnotemark[23]})\ . \tag*{(4.4.7.4)}
\]
\footnotetext[23]{limite inductive prise dans $C$, mais valable aussi dans $\Ind'_\pi C$.}
\textit{Alors $\varinjlim\limits_{j\in J} \varphi(j)$ (dans $\Ind'_\pi C$) existe si et seulement si}
\msp{page 88}
\textit{la $\varinjlim\limits_{i\in I} \Psi(i)$ existe (dans $\Ind'_\pi C$ également), et les deux sont égales.}

\textbf{Corollaire 4.4.8.} \textit{Supposons que $C$ soit stable par $\varinjlim$ de cardinal $\leqslant \pi$ (pas seulement par celles de type $J$). Alors $\uHom(J,C)$ et $I = \uHom(J,C)/\varphi$ ont la même propriété, $I$ est grande devant $\pi$, et $\varinjlim\limits_{j\in J} \varphi(j) \ =\ \varinjlim\limits_{i\in I} \Psi(i)$ existe dans $\Ind'_\pi(C)$, sous réserve de cofinalité des foncteurs (4.4.6.3).}

\printedpage{21}
% head: A. Grothendieck : Dérivateurs, Ch. XVIII, Édition des Universités de Montpellier II et Paris VII
% foot: Version 17 décembre 2025
Il suffit en effet de vérifier ce

\textbf{Lemme 4.4.8.1.} \textit{Soit $I$ une catégorie grande devant $\pi$, et $\Psi : I \longrightarrow C$ un foncteur. Alors $\varinjlim\limits_{i} \Psi(i)$ existe dans $\Ind'_\pi(C)$, et c'est la limite inductive dans $C^\wedge$, ou encore “c'est” l'ind-objet grand devant $\pi$ représenté par $\Psi$.} (En d'autres termes (\footnote[24]{C'est abusif, cette stabilité est un énoncé nettement plus fort! Je le prouve page 99.}), $\Ind'_\pi(C) \subset C^\wedge$ est une sous-catégorie strictement pleine \textit{stable} par $\varinjlim$ grandes devant $\pi$.)

\msp{page 89}
Ce lemme est une tautologie : si $F$ désigne la limite inductive dans $C^\wedge$, il est clair par définition (en prenant un foncteur $I_0 \longrightarrow I$ cofinal, avec $I_0$ \textit{ordonnée} grande devant $\pi$, de sorte que la $\varinjlim$ peut être prise suivant $I_0$) que $F$ est dans $\Ind'_\pi(C)$, et alors c'est a fortiori une $\varinjlim$ dans $\Ind'_\pi C$.

Pour vérifier l'existence des $\varinjlim$ de cardinal $\leqslant \pi$ dans $\Ind'_\pi C$, dans le cas où $C$ lui-même est stable par lesdites limites, il suffit donc, grâce au corollaire 4.4.8, de prouver la cofinalité des foncteurs (4.4.6.3). Cela va résulter du

\textbf{Lemme 4.4.9.} \textit{Soient $C$ une petite catégorie, $\pi$ un cardinal infini, $J$ une petite catégorie telle que pour tous $j, j' \in \Ob J$, on ait $\card \Hom(j,j') \leqslant \pi$. Sous ces conditions, pour tout $j \in \Ob J$, le foncteur (4.4.6.3) \textup{admet un adjoint à gauche,} donc a fortiori il est cofinal (car les catégories colocalisées au dessus de $C/\varphi(j)$ ont un objet initial, a fortiori elles sont $0$-connexes).}

\msp{page 90}
\textsc{Démonstration de} 4.4.8. Revenons au diagramme (4.4.6.2), dont on déduit (4.4.6.3), nous allons le réécrire de façon simplifiée :
\[
\begin{tikzcd}
B \arrow[r, "\beta"] \arrow[d, shift left, "\varepsilon"] & B' \arrow[d, shift left, "\varepsilon'"] & \ni y'_0 \arrow[d, mapsto] \\
C \arrow[r, "\gamma"] \arrow[u, dotted, shift left, "\lambda"] & C' \arrow[u, dotted, shift left, "\lambda'"] & \ni z'_0\ .
\end{tikzcd}
\tag*{(4.4.9.1)}
\]
Je dis que $\varepsilon$, $\varepsilon'$ ont \textit{des adjoints à gauche} $\lambda$ resp. $\lambda'$, et que l'homomorphisme canonique de “changement de base”
\[
\lambda'\gamma \longrightarrow \beta\lambda \tag*{(4.4.9.2)}
\]
est un \textit{isomorphisme.} En effet, l'existence des adjoints de $\varepsilon$, $\varepsilon'$, à savoir des foncteurs $\varepsilon_{j\,!}$ et $\varepsilon'_{j\,!}$ relatifs à l'inclusion
\[
\varepsilon_j \ :\ e \longrightarrow J
\]
et aux catégories de foncteurs $\uHom(J,C)$ resp. $\uHom(J,\Ind'_\pi C)$, résultant de l'existence des sommes de type $\Hom(j,j')$ ($j' \in \Ob J$) dans $C$ (déjà acquise), \textit{et dans} $\Ind'_\pi C$\textit{ ;} et le fait que le foncteur $C \longrightarrow \Ind'_\pi C$ commute
\msp{page 91}
auxdites sommes (chose également déjà acquise, puisque $\card\, (\Hom(j,j')) \leqslant \pi$) implique alors que (4.4.4.2) est un isomorphisme. Moyennant la vérification en suspens (l'existence % [sic] (4.4.4.2) -> (4.4.9.2), the isomorphism in question is the change-of-base arrow (4.4.9.2)
\printedpage{22}
% head: A. Grothendieck : Dérivateurs, Ch. XVIII, Édition des Universités de Montpellier II et Paris VII
% foot: Version 17 décembre 2025
des sommes de cardinal $\leqslant \pi$ dans $\Ind'_\pi(C)$), le lemme 4.4.9 résultera du résultat plus général :

\textbf{Lemme 4.4.10.} \textit{Considérons un carré commutatif (4.4.9.1) (traits pleins) de catégories, où $\varepsilon$ admet un adjoint à gauche $\lambda$, $\varepsilon'$ un adjoint à gauche $\lambda'$, et la flèche de changement de base (4.4.9.2) est un isomorphisme. Alors pour tout $y'_0$ dans $B'$, posant $\varepsilon'(y'_0) = z'_0$, le foncteur induit par $(\varepsilon, \varepsilon')$,}
\[
\overline{B} \ =\ B/y'_0 \xrightarrow{\bar{\varepsilon}} \overline{C} = C/z'_0\ ,
\]
\textit{admet un adjoint à gauche $\overline{\lambda}$, donné par}
\[
\begin{tikzcd}[column sep=small]
\overline{\lambda}(\gamma(z) \xrightarrow{v} z'_0) \ \ = & \Bigl(\lambda'\gamma(z) \simeq \beta(\lambda(z)) \arrow[r, "{\lambda'(v)}"] \arrow[ddr, "{\substack{\overline{\lambda}(v) \,=\, \text{composé}\\ \text{adj} \circ \lambda'(v)}}"'] & \lambda'(z'_0) \arrow[d, equal] \\
 & & \lambda'\varepsilon'(y'_0) \arrow[d, "\text{adj.}"] \\
 & & y'_0\Bigr)\ \ .
\end{tikzcd}
\]

\msp{page 92}
La démonstration se fait par un AQT \texttt{[âne qui trotte]} fastidieux (j'avoue ne pas avoir poussé jusqu'au dernier détail \ldots).

Pour terminer la démonstration de 4.4.9, il reste à prouver un dernier ( ?)

\textbf{Lemme 4.4.11.} \textit{Soit $C$ une petite catégorie, stable par sommes de cardinal $\leqslant \pi$. Alors $\Ind'_\pi C$ est stable par sommes de cardinal $\leqslant \pi$ (et, pour mémoire, le foncteur $\alpha : C \longrightarrow \Ind'_\pi C$ y commute).}

Appliquant le corollaire 4.4.8 au lemme 4.4.7, on est ramené à prouver la cofinalité de (4.4.6.3) quand $J$ est discrète, plus le fait que $\uHom(J,C)/\varphi$ est grande devant $\pi$, dans ce cas. Mais la donnée de $\varphi$ revient à celle d'une famille $(\varphi_j)_{j\in J}$ d'objets de $\Ind'_\pi(C)$, et le foncteur (4.4.6.3) s'explicite comme la projection
\[
\prod_{j\in J} (C/\varphi_j) \longrightarrow C/\varphi_j
\]
du produit total sur un des facteurs. Donc on est ramené à ceci :

\msp{page 93}
\textbf{Lemme 4.4.12.} \textit{Soit $(C_j)_{j\in J}$ une famille de petites catégories grandes devant $\pi$, indexée par un petit ensemble $J$. Alors $\prod C_j$ est grande devant $\pi$, et les projections (pour $j_0 \in J$) $\prod C_j \longrightarrow C_{j_0}$ sont cofinales.}

Prouvons d'abord la cofinalité, qui est le plus important pour nous. C'est évident si $J = \{j_0\}$, sinon on a $J = J' \cup \{j_0\}$, et $\prod_{j\in J} C_j \simeq C' \times C_{j_0}$, où $C' = \prod_{j\in J'} C_j$, donc il faut voir que $C' \times C_{j_0} \longrightarrow C_{j_0}$ est cofinal, i.e. que toute catégorie $\xi\backslash(C' \times C_{j_0})$ (pour $\xi \in \Ob C_{j_0}$) est 0-connexe, or elle est isomorphe à $C' \times (\xi\backslash C_{j_0})$, et $\xi\backslash C_{j_0}$ étant 0-connexe, il reste à voir que $C'$ est 0-connexe. On fera attention à cet égard qu'un produit \textit{infini} de catégories 0-connexes n'est pas forcément 0-connexe, donc la cofinalité n'est pas totalement tautologique! Mais si on prouve que la catégorie $\prod_{j\in J} C_j$ est
\printedpage{23}
% head: A. Grothendieck : Dérivateurs, Ch. XVIII, Édition des Universités de Montpellier II et Paris VII
% foot: Version 17 décembre 2025
\msp{page 94}
grande devant $\pi$, donc filtrante, cela implique qu'elle est 0-connexe, et on peut appliquer ce même résultat à $C' = \prod_{j\in J'} C_j$, et on gagne.

Donc on est ramené à prouver que $\prod C_j$ est grande devant $\pi$. Ceci se voit soit sur le critère $\mathrm{PF}_\pi\,1$, $\mathrm{PF}_\pi\,2$ de 4.4.4 b, soit en se ramenant au cas où les $C_j$ sont ordonnées, mais cela a l'air plus long. Notons que l'on prouve de même qu'un (petit) produit de catégories filtrantes est encore une catégorie filtrante.

\textbf{Remarques.}
\begin{enumerate}[label={}, leftmargin=*]
\item[1.] Cet argument ne montre nullement (même si $C$ est stable par $\varinjlim$ de cardinal $\leqslant \pi$) que pour tout $\varphi \in \uHom(J,\Ind(C))$, le foncteur similaire à (4.4.6.3)
\[
\uHom(J,C)/\varphi \longrightarrow C/\varphi(j)
\]
soit cofinal, et je doute qu'il le soit toujours. La difficulté provient du fait que $C \longrightarrow \Ind\, C$ ne commute pas aux sommes de cardinal $\leqslant \pi$, seulement aux sommes finies,
\msp{page 95}
donc le lemme 4.4.10 ne peut être invoqué.
\item[2.] Nous venons de prouver, via 4.4.8 et 4.4.9 (dont la démonstration est achevée avec 4.4.12), que si $C$ est une petite catégorie stable par $\varinjlim$ de cardinal $\leqslant \pi$, alors $\Ind_\pi C$ est également stable par ces limites. De plus, les systèmes inductifs de ce type $J$ dans $\Ind_\pi C$ admettent une représentation indicielle
\[
\psi \ :\ I \times J \longrightarrow C\ ,
\]
avec $I$ grande devant $\pi$. Et on peut toujours supposer, quitte à prendre $I' \longrightarrow I$ cofinal avec $I'$ ensemble ordonné grand devant $\pi$, que $I$ est de plus un ensemble ordonné.
\end{enumerate}

C'est là le moyen donc, pour la démonstration du

\textbf{Théorème 4.4.13.} \textit{Soit $C$ une catégorie équivalente à une petite catégorie, et stable par $\varinjlim$ de cardinal $\leqslant \pi$, où $\pi$ est un cardinal infini donné. Alors :}
\msp{page 96}

\begin{enumerate}[label={}, leftmargin=*]
\item[\textit{a)}] \textit{(Pour mémoire.) Le foncteur canonique}
\[
C \longrightarrow \Ind_\pi(C) \xrightarrow{\ \sim\ } \Ind'_\pi(C)
\]
\textit{commute aux $\varinjlim$ de cardinal $\leqslant \pi$.}
\item[\textit{b)}] \textit{(Pour mémoire.) Les systèmes inductifs de cardinal $\leqslant \pi$ dans $\Ind_\pi(C)$ (i.e. les foncteurs $\varphi : J \longrightarrow \Ind_\pi(C)$, avec $\card \Fl(J) \leqslant \pi$) admettent une représentation indicielle}
\[
\psi \ :\ I \times J \longrightarrow C\ ,
\]
\textit{avec $I$ ensemble ordonné grand devant $\pi$, et admettent une limite inductive dans $\Ind_\pi C$ qui se calculent par}
\[
\underbrace{\varinjlim_{j\in J} \varphi(j)}_{\text{dans } \Ind_\pi(C)} \ \simeq\ \text{“}\,\varinjlim_{i\in I}\,\text{”}\Psi(i)
\]
\end{enumerate}


\endgroup
% END UNIT der-ch18/der-ch18_p021-025.tex
% BEGIN UNIT der-ch18/der-ch18_p026-030.tex
\begingroup\makeatletter%
\newcommand{\printedpage}[1]{\marginpar{\footnotesize p.~#1}}%
\newcommand{\msp}[1]{\marginpar{\footnotesize [#1]}}%
\newcommand{\Ind}{\operatorname{Ind}}%
\newcommand{\uHom}{\underline{\mathrm{Hom}}}%
\newcommand{\Ob}{\operatorname{Ob}}%
\newcommand{\card}{\operatorname{card}}%
\newcommand{\Fl}{\mathrm{Fl}}%
\newcommand{\PF}{\mathrm{PF}}%
\makeatother\let\printedpage\derpp\let\msp\dermsp%
\let\mspin\DERmspin%


\printedpage{24}
% head: A. Grothendieck : Dérivateurs, Ch. XVIII, Édition des Universités de Montpellier II et Paris VII
% foot: Version 17 décembre 2025

\textit{avec}
\[
\Psi(i) \;=\; \varinjlim_{j\in J}{}_{C}\,\psi(i,j)\;.
\]
\textit{c)} $\Ind_\pi C$ \textit{est stable par petites limites inductives} (\footnote{rajouter aussi, en d), le corollaire 4.4.16, page 101.}).

Il reste à établir ce dernier point. Comme on a déjà les $\varinjlim$ finies par b, il reste à prouver l'existence des petites $\varinjlim$ filtrantes, ou au choix des sommes infinies, ce qui sera plus joli en fait. Mais
\msp{page 97}
on pourra aussi bien traiter le cas général d'une $\varinjlim$ de type $J$ quelconque ($J$ petite). Pour ceci, soit
\[
I \;=\; I(J) \;=\; \begin{aligned}[t]
&\text{ensemble des sous-catégories } J' \text{ de } J \text{ telles que}\\
&\card\Fl(J')\leqslant\pi,
\end{aligned}
\]
ordonnons $I$ par inclusion. Il est évident que $I\neq\emptyset$ (car $\emptyset\in I$), et même que $I$ \textit{est grand devant} $\pi$, car si $(J'_\alpha)_{\alpha\in A}$ avec $\card(A)\leqslant\pi$ est une famille d'éléments de $I$, alors la sous-catégorie de $J$ engendrée par les $J'_\alpha$ est encore dans $I$ (vérification immédiate), et majore les $J'_\alpha$. En fait, la considération de $I(J)$ associée à $J$ va nous permettre de prouver ceci, qui généralise 4.4.13 c), compte tenu de b) :

\textbf{Corollaire 4.4.14} (\footnote{Comparer aussi 4.11.3, page 270.}). \textit{Soit $\mathcal{M}$ une catégorie, $\pi$ un cardinal infini tel que}
\begin{enumerate}
\item[a)] \textit{$\mathcal{M}$ est}
\msp{page 98}
\textit{stable par $\varinjlim$ de cardinal $\leqslant\pi$, et}
\item[b)] \textit{$\mathcal{M}$ est stable par $\varinjlim$ sur les ensembles ordonnés $I$ grands devant $\pi$.}
\end{enumerate}
\textit{Alors $\mathcal{M}$ est stable par petites $\varinjlim$ .}

(\footnote{\textbf{Corollaire 4.4.14.1.} $\mathcal{M}$ étant comme dans 4.4.14, et $f:\mathcal{M}\longrightarrow\mathcal{E}$ étant un foncteur qui commute aux $\varinjlim$ de cardinal $\leqslant\pi$, et aux petites $\varinjlim$ grandes devant $\pi$, alors $f$ commute aux petites $\varinjlim$ .}).

En effet, si $\varphi:J\longrightarrow\mathcal{M}$ est donné, on en construit
\[
\psi\;:\;I\;=\;I(J)\;\longrightarrow\;\mathcal{M}
\]
par
\[
\psi(J')\;=\;\varinjlim_{J'}\varphi(J')\;,
\]
qui existe par hypothèse sur $\mathcal{M}$. Il est alors immédiat que $\varinjlim\limits_{J}\varphi$ existe si et seulement si $\varinjlim\limits_{I}\psi$ existe, or cette existence est assurée du fait que $I$ est grand devant $\pi$, et l'hypothèse sur $\mathcal{M}$.

Pour achever de prouver le théorème 4.4.13, on est donc ramené au cas particulier des limites inductives $\varinjlim\limits_{J}\varphi(j)$, quand $\varphi:J\longrightarrow\Ind'_\pi C$ est donné, avec $J$ ensemble
\printedpage{25}
% head: A. Grothendieck : Dérivateurs, Ch. XVIII, Édition des Universités de Montpellier II et Paris VII
% foot: Version 17 décembre 2025
\msp{page 99}
ordonné grand devant $\pi$. Mais ceci va résulter du résultat plus précis :

\textbf{Corollaire 4.4.15} (\footnote{comparer à 4.7.3.1, qu'il y aurait lieu d'expliciter ici.}). \textit{Soit $C$ une petite catégorie. Alors la sous-catégorie strictement pleine} $\Ind'_\pi C$ \textit{de $C^\wedge$ est stable dans $C^\wedge$ par $\varinjlim$ grandes devant $\pi$, en d'autres termes, ces $\varinjlim$} existent dans $\Ind'_\pi(C)$, \textit{et le foncteur d'inclusion y commute.} % doubt resolved: existent is upright in the PDF; see corrections.csv, PDF leaf 27

\textsc{Démonstration}. Soit $I$ ensemble ordonné, grand devant $\pi$, et $i\longmapsto F_i$ un foncteur $I\longrightarrow\Ind'_\pi(C)$, $F\;=\;\varinjlim\limits_{I}F_i$ dans $C^\wedge$, prouvons que $F\in\Ob\Ind'_\pi(C)$, i.e. (4.4.3) que $C/F$ est grande devant $\pi$. J'utilise le critère b) de 4.4.4, il faut vérifier $\PF_\pi\,1$, $\PF_\pi\,2$. Soient
\[
x_\alpha\xrightarrow{\;u_\alpha\;}F\;=\;\varinjlim F_i
\]
des objets de $C/F$ indexés par $A$ avec $\card A\leqslant\pi$. Chaque $u_\alpha$ provient de $x_\alpha\longrightarrow F_{i_\alpha}$, et en prenant un majorant commun des $i_\alpha$ ($I$ grand devant $\pi$), on trouve un $i\in I$ et des $v_\alpha:x_\alpha\longrightarrow F_i$. Comme $F_i$ est dans $\Ind'_\pi(C)$, i.e. $C/F_i$ grande devant $\pi$, considérant les $v_\alpha$ comme des objets de $C/F_i$, ils sont majorés par un objet $x\longrightarrow F_i$ de $C/F_i$. On a donc
\[
\begin{tikzcd}
x_\alpha \arrow[r] \arrow[d] & F_{i_\alpha} \arrow[d] \arrow[dr, dotted] & \\
x \arrow[r] & F_i \arrow[r] & F\;,
\end{tikzcd}
\]
qui montre que les $x_\alpha\longrightarrow F$ sont majorés par $x\longrightarrow F$,
\msp{page 100}
ce qui prouve $\PF_\pi\,1$. Prouvons $\PF_\pi\,2$, i.e. soient $x\xrightarrow{\;u\;}F$, $y\xrightarrow{\;v\;}F$ deux objets de $C/F$, et $(f_\alpha)_{\alpha\in A}$ une famille de cardinal $\leqslant\pi$ de flèches de $u$ dans $v$. On sait que $u$, $v$ proviennent de $u_i:x\longrightarrow F_i$, $v_i:y\longrightarrow F_i$ pour $i$ assez grand.
\[
\begin{tikzcd}[column sep=small]
x \arrow[rr, "f_\alpha"] \arrow[dr, "u"'] & & y \arrow[dl, "v"] \\
& F &
\end{tikzcd}
\]
\[
\begin{tikzcd}[column sep=small]
x \arrow[rr, "f_\alpha"] \arrow[dr, "u_i"'] & & y \arrow[dl, "v_i"] \\
& F_i &
\end{tikzcd}
\]

Pour tout $\alpha$, on peut choisir $i_\alpha\geqslant i$ tel que $v_if_\alpha$ et $u_i$ deviennent égaux en composant avec $F_i\longrightarrow F_{i_\alpha}$, et comme $\card A\leqslant\pi$ et $I$ grand devant $\pi$, il existe un majorant $i'$ des $i_\alpha$. Au total, on peut supposer (remplaçant $i$ par $i'$) que $v_if_\alpha=u_i$ pour tout $\alpha$. Alors $u_i$, $v_i$ étant considérés comme objets de $C/F_i$, les $f_\alpha$ sont des flèches de $C/F_i$ de $u_i$ dans $v_i$. Comme $C/F_i$ est grand devant $\pi$ (par l'hypothèse $F_i\in\Ob\Ind'_\pi(C)$), on trouve qu'il existe un égalisateur $(y,v_i)\xrightarrow{\;g\;}(z,w_i)$.
\printedpage{26}
% head: A. Grothendieck : Dérivateurs, Ch. XVIII, Édition des Universités de Montpellier II et Paris VII
% foot: Version 17 décembre 2025
\[
\begin{tikzcd}
x \arrow[r, shift left=2, "(f_\alpha)"] \arrow[r] \arrow[r, shift right=2] \arrow[dr, "u_i"'] & y \arrow[r, "g"] \arrow[d, "v_i"'] & z \arrow[dl, "w_i"] \\
& F_i &
\end{tikzcd}
\]

Alors l'image de $g$ dans $C/F$ (via $F_i\longrightarrow F$) égalise les flèches $f_\alpha$, en tant que flèches de $C/F$, ce qui prouve $\PF_\pi\,2$ et achève la démonstration de 4.4.14, et par là aussi celle du théorème 4.4.13, et du lemme 4.3. Par là est aussi achevée la démonstration de la première assertion du lemme
\msp{page 101}
4.3.7 (p. 77), ce qui achève la démonstration du théorème 4.3.6 (p. 70-71). Il reste à établir la propriété universelle énoncée dans 4.3.7, que je vais rappeler :

\textbf{Corollaire 4.4.16} (\footnote{Pourrait être incorporé en partie d) au théorème 4.4.13, car c'est un complément important, et pas tellement évident finalement.}). \textit{Soient $C$ une catégorie équivalente à une petite catégorie, stable par $\varinjlim$ de cardinal $\leqslant\pi$, où $\pi$ est un cardinal infini donné, et considérons le foncteur d'inclusion canonique}
\[
C\xrightarrow{\;\alpha\;}\Ind_\pi(C)\;. \tag*{(4.4.16.1)}
\]
\textit{Ce foncteur est} 2\textit{-universel pour les foncteurs $f:C\longrightarrow\mathcal{E}$, où $\mathcal{E}$ est une $\mathcal{U}$-catégorie stable par petites $\varinjlim$ , et $f$ un foncteur commutant à ces limites inductives. En d'autres termes, on a}
\begin{enumerate}
\item[\textit{1$^\circ$)}] \textit{$\alpha$ possède les propriétés précédentes, i.e.} $\Ind_\pi(C)$ \textit{est stable par $\varinjlim$ , et le foncteur $\alpha$ commute aux $\varinjlim$ de cardinal $\leqslant\pi$.}
\item[\textit{2$^\circ$)}] \textit{Pour toute $\mathfrak{U}$-catégorie $\mathcal{E}$ stable par petites limites inductives, le foncteur}
\[
\begin{aligned}
g\;&\longmapsto\;g\circ\alpha\\
\uHom_{!}(\Ind_\pi C,\mathcal{E})\;&\longrightarrow\;\uHom_{!\,\pi}(C,\mathcal{E})
\end{aligned}
\tag*{(4.4.16.2)}
\]
\textit{est une équivalence de catégories, où l'indice $!$ (resp. $!\,\pi$) indique la sous-catégorie strictement pleine de la catégorie} $\uHom$
\msp{page 102}
\textit{envisagée, formée des foncteurs qui commutent aux petites $\varinjlim$ (resp. à celles de cardinal $\leqslant\pi$).}
\end{enumerate}

On peut donc dire aussi que $\Ind_\pi(C)$ joue le rôle d'une “enveloppe” stable par petites $\varinjlim$ de la catégorie $C$, quand les $\varinjlim$ de cardinal $\leqslant\pi$, lesquelles existent déjà dans $C$, sont respectées par les plongement de $C$ dans son enveloppe. C'est un résultat nullement tautologique qu'il suffit pour cela d'ajouter à $C$ des $\varinjlim$ \textit{grandes devant} $\pi$ d'objets de $C$, et que plus précisément, la catégorie enveloppe est réalisée par $\Ind_\pi C$ (ou, au choix et moins pléthoriquement, par $\Ind'_\pi C\subset C^\wedge$). % [sic] les plongement -> le plongement, number agreement

\textsc{Démonstration}. Le 1°) est contenu dans 4.4.13 a) et c). Prouvons donc 2°). Il est évident que $\alpha$ est un foncteur pleinement fidèle dont l'image essentielle engendre strictement
\printedpage{27}
% head: A. Grothendieck : Dérivateurs, Ch. XVIII, Édition des Universités de Montpellier II et Paris VII
% foot: Version 17 décembre 2025
$\mathcal{M}=\Ind_\pi C$, d'où résulte que le foncteur (4.4.16.2) est pleinement fidèle, étant induit par le foncteur pleinement fidèle
\[
\alpha^*_{\mathcal{E}}\;:\;\uHom(\mathcal{M},\mathcal{E})\;\longrightarrow\;\uHom(C,\mathcal{E})
\]
\msp{page 103}
sur les sous-catégorie pleines envisagées. Il reste donc à montrer que (4.4.16.2) est essentiellement surjectif. Soit donc $f:C\longrightarrow\mathcal{E}$ dans $\uHom_{!\,\pi}(C,\mathcal{E})$, considérons son extension canonique % [sic] les sous-catégorie pleines -> les sous-catégories pleines, number agreement
\[
\left\{
\begin{aligned}
&Lf\;:\;\Ind(C)\;\longrightarrow\;\mathcal{E}\\
&Lf((X_i)_{i\in I})\;=\;\varinjlim_{i\in I}f(X_i)
\end{aligned}
\right.
\tag*{(4.4.16.3)}
\]
(cf. SGA 4 I 8.7.3), et soit
\[
\psi(f)\;=\;Lf|\Ind_\pi(C)\;. \tag*{(4.4.16.4)}
\]
Rappelons que $Lf$ commute aux petites $\varinjlim$ filtrantes, mais comme le foncteur d'inclusion
\[
\Ind_\pi C\;\longrightarrow\;\Ind(C) \tag*{(4.4.16.5)}
\]
ne commute pas à ces $\varinjlim$ , il n'est pas clair a priori que $\psi(f)$ commute auxdites limites, ni a fortiori qu'il commute aux petites $\varinjlim$ sans plus. Mais on a vu que (4.4.16.5) commute aux $\varinjlim$ \textit{grandes devant} $\pi$, ce qui est contenu dans le corollaire 4.4.15 (vu que les limites inductives \textit{filtrantes} dans $\Ind(C)$ sont les limites \textit{dans} $C^\wedge$) : dire que $\Ind'_\pi(C)\longrightarrow\Ind'(C)$ commute à un certain type de $\varinjlim$ filtrantes, revient à dire qu'il en est ainsi du composé avec $\Ind'_\pi(C)\hookrightarrow C^\wedge$, ce qui est établi par 4.4.15. De ceci résulte que $\psi(f)$ \textit{commute aux} $\varinjlim$ \textit{grandes devant} $\pi$. Pour établir que $g=\psi(f)$ commute
\msp{page 104}
aux petites $\varinjlim$ , il reste donc à voir qu'il commute aux $\varinjlim$ de cardinal $\leqslant\pi$, en vertu de 4.4.14.1 (marge page 98). Pour l'établir, on utilise le théorème 4.4.13 b, en représentant un foncteur $\varphi:J\longrightarrow\Ind_\pi(C)$ (où $\card\Fl J\leqslant\pi$) sous forme indicielle par
\[
\Phi\;:\;I\times J\;\longrightarrow\;C\;, \tag*{(4.4.16.6)}
\]
$I$ étant un ensemble ordonné grand devant $\pi$. On a donc
\[
\varphi(j)\;=\;\varinjlim_{i\in I}{}_{\Ind_\pi C}\,\Phi(i,j)\;,
\]
et comme $g=\psi(f)$ commute aux $\varinjlim$ de type $I$ par ce qu'on a vu, on aura
\[
g(\varphi(j))\;=\;\varinjlim_{i\in I}{}_{(\mathcal{E})}\;g(\Phi(i,j))\;,
\]
d'où
\[
\varinjlim_{j\in J}{}_{(\mathcal{E})}\,g(\varphi(j))\;=\;\varinjlim_{j\in J}\;\varinjlim_{i\in I}\;\underbrace{g(\Phi(i,j))}_{=\,f(\Phi(i,j))}\qquad\text{(limites dans $\mathcal{E}$)}\;. \tag*{$(*)$}
\]
\printedpage{28}
% head: A. Grothendieck : Dérivateurs, Ch. XVIII, Édition des Universités de Montpellier II et Paris VII
% foot: Version 17 décembre 2025

D'autre part, on a
\[
\begin{aligned}
\varinjlim_{j\in J}{}_{\Ind_\pi}\,\varphi(j)\;&\simeq\;\varinjlim_{j\in J}\;\varinjlim_{i\in I}\;\Phi(i,j)\\
&\simeq\;\varinjlim_{i\in I}\;\underbrace{\varinjlim_{j\in J}\;\Phi(i,j)}_{\substack{\text{peut être calculé}\\ \text{dans $C$, soit}\\ \Psi(i)\,\in\,\Ob C}}\;,
\end{aligned}
\]
donc
\[
\varinjlim_{j\in J}{}_{\Ind_\pi}\,\varphi(j)\simeq\varinjlim_{i\in I}\;\Psi(i)\;, \tag*{$(**)$}
\]
\msp{page 105}
d'où en appliquant $g$, qui commute aux $\varinjlim$ de type $I$, et se rappelant maintenant que $g|C\simeq f$,
\[
g(\,\varinjlim_{J}{}_{\Ind_\pi C}\,\varphi(j))\;\simeq\;\varinjlim_{i\in I}\;f(\Psi(i))\;. \tag*{$(*{*}*)$}
\]
On a d'autre part
\[
\Psi(i)\;\overset{\text{déf}}{=}\;\varinjlim_{j\in J}\;\Phi(i,j)\;,\qquad\qquad \varinjlim\;\text{ dans } C\;,
\]
et comme $f$ commute aux $\varinjlim$ de type $J$ (car $\card\Fl J\leqslant\pi$), on en conclut
\[
f(\Psi(i))\;\simeq\;\varinjlim_{j\in J}\;f(\Phi(i,j))\;,
\]
donc par $(*{*}*)$
\[
g(\,\varinjlim_{j\in J}\;\varphi(j))\;\simeq\;\varinjlim_{i\in I}\;\varinjlim_{j\in J}\;f(\Phi(i,j))\;. \tag*{$(*{*}{*}*)$}
\]
En composant avec $(*)$ et utilisant à nouveau la commutativité de $\varinjlim\limits_{I}$ à $\varinjlim\limits_{J}$ , on trouve donc
\[
g(\,\varinjlim\;\varphi(j))\;\overset{\sim}{\longleftarrow}\;\varinjlim\;g(\varphi(j))\;,
\]
ce qui prouve que $g$ commute aux limites inductives de type $J$, et achève de prouver que $g$ est \texttt{[dans]} $\uHom_{!}(\Ind_\pi(C),\mathcal{E})$. Comme $g|C\simeq f$, on en conclut que le foncteur (4.4.16.2) est essentiellement surjectif, ce qui achève de prouver 4.4.16. De plus, on a construit un foncteur quasi-inverse explicite, $\psi$,
\msp{page 106}
s'insérant dans le diagramme commutatif
\[
\begin{tikzcd}[column sep=large, row sep=large]
\uHom_{!\,\pi}(C,\mathcal{E}) \arrow[r, "\psi", "(\sim)"'] \arrow[d, hook, "\text{incl. pl. fid.}"] & \uHom_{!}(\Ind_\pi(C),\mathcal{E}) \arrow[r] \arrow[d, hook, dotted, "\text{pl. fid.}"] & \uHom_{!\,\pi}(\Ind_\pi(C),\mathcal{E}) \\
\uHom(C,\mathcal{E}) \arrow[r, "L", "\substack{\sim\\ \text{(cf. SGA 4 I 8.7.3)}}"'] & \uHom_{!\,\mathrm{filt}}(\Ind(C),\mathcal{E}) \arrow[ur, "R\ \text{(restriction)}"'] &
\end{tikzcd}
\qquad.
\tag*{(4.4.16.6)} % [sic] (4.4.16.6) -> (4.4.16.7)?, the label (4.4.16.6) is already used for Phi on the preceding page
\]


\endgroup
% END UNIT der-ch18/der-ch18_p026-030.tex
% BEGIN UNIT der-ch18/der-ch18_p031-035.tex
\begingroup\makeatletter%
\newcommand{\printedpage}[1]{\marginpar{\footnotesize p.~#1}}%
\newcommand{\msp}[1]{\marginpar{\footnotesize [#1]}}%
\newcommand{\card}{\operatorname{card}}%
\newcommand{\Fl}{\operatorname{Fl}}%
\newcommand{\Ob}{\operatorname{Ob}}%
\newcommand{\Ind}{\operatorname{Ind}}%
\newcommand{\Filt}{\operatorname{Filt}}%
\newcommand{\Hom}{\operatorname{Hom}}%
\newcommand{\Aut}{\operatorname{Aut}}%
\newcommand{\uHom}{\underline{\mathrm{Hom}}}%
\newcommand{\uFilt}{\underline{\mathrm{Filt}}}%
\newcommand{\circled}[1]{\tikz[baseline=(c.base)]\node[draw,circle,inner sep=1pt](c){#1};}%
\makeatother\let\printedpage\derpp\let\msp\dermsp%
\let\mspin\DERmspin%

\setcounter{footnote}{29}

\printedpage{29}
% head: A. Grothendieck : Dérivateurs, Ch. XVIII, Édition des Universités de Montpellier II et Paris VII
% foot: Version 17 décembre 2025

\section*{4.5\quad Filtration cardinale canonique d'une catégorie $\mathfrak{U}$-paratopos}

\textbf{Définition 4.5.1.} On appelle \textit{catégorie $\mathfrak{U}$-paratopos} une catégorie $\mathcal{M}$ $\mathfrak{U}$-pseudo-topos et accessible, i.e. qui satisfait aux conditions équivalentes du théorème 4.3.6 (p. 70). On dit que c'est une \textit{catégorie $\pi$-$\mathfrak{U}$-paratopos}, si elle est $\pi$-accessible, i.e. si elle admet une petite sous-catégorie génératrice dont les objets soient $\pi$-accessible, ou ce qui revient au même, si % [sic] soient π-accessible -> soient π-accessibles, accord du pluriel
$\mathcal{M}$ est équivalente à une catégorie de la forme $\Ind'_{\pi}(\mathcal{M}_{\pi})$, où $\mathcal{M}_{\pi}$ est une petite catégorie (ou : une catégorie équivalente à une petite catégorie), et stable par $\varinjlim$ de cardinal $\leqslant \pi$. (Et on peut prendre pour $\mathcal{M}_{\pi}$ la sous-catégorie pleine de $\mathcal{M}$ formée des objets
\msp{page 107}
$\pi$-accessibles).

Dans la suite, nous allons, plus généralement, considérer une catégorie $C$ équivalente à une petite catégorie, stable par $\varinjlim$ de cardinal $\leqslant \pi_0$ (\footnote{Il n'y a pas lieu de le supposer tout de suite.}), et la catégorie
\[
\mathcal{M} \ \overset{\text{déf}}{=}\ \Ind'_{\pi_0}C \ \subset\ C^{\wedge}\ ,
\tag*{(4.5.1)}
\]
où $\pi_0$ est un cardinal infini donné. Je me propose d'introduire sur $\mathcal{M}$ une filtration cardinale canonique :

\textbf{Proposition 4.5.2} (\footnote{N'est prouvé (et n'est vrai) que si $\pi^{\pi_0} = \pi$, voir rectification 4.5.6, page 116.}). \textit{Soit $F \in \Ob\mathcal{M} \subset \Ob C^{\wedge}$, et soit $\pi$ un cardinal $\geqslant \pi_1 =$} $\sup(\pi_0, \card\Fl C_0)$\textit{, où $C_0$ est une petite sous-catégorie pleine de $C$ telle que $C_0 \overset{\sim}{\hookrightarrow} C$. Les conditions suivantes sont équivalentes :}
\begin{enumerate}[leftmargin=3em]
\item[\textit{a)}] $\card\Fl(C_0/F) \leqslant \pi$.
\item[\textit{a$'$)}] $\card\Ob\,(C_0/F) \leqslant \pi$.
\item[\textit{a$''$)}] \textit{Pour tout $x$ dans $C_0$,} $\card F(x) \leqslant \pi$.
\item[\textit{b)}] \textit{Il existe une petite catégorie filtrante $I$ (pas nécessairement ordonnée),} grande devant $\pi_0$\textit{, avec} $\card I \leqslant \pi$\textit{, et un système inductif $(x_i)_{i\in I}$ dans $C$, tel que l'on ait}
\[
F \ \simeq\ \varinjlim_{i\in I} x_i\ ,
\tag*{(4.5.2.1)}
\]
\textit{où la $\varinjlim$ est prise dans $\mathcal{M}$, ou ce qui revient au même (en vertu de 4.4.15, $I$ étant grande devant $\pi_0$) dans $C^{\wedge}$, i.e. $F$ est isomorphe au foncteur ind-représentable défini par l'ind-objet $\underline{x} = (x_i)_{i\in I}$ de $C$.}
\item[\textit{c)}] \textit{$F$ est $\pi$-accessible dans $\mathcal{M}$.}
\item[\textit{c$'$)}] \textit{$F$ est $\pi$-accessible dans $C^{\wedge}$.}
\end{enumerate}
\msp{page 108}

\textsc{Démonstration}. On a les implications $\pm$ tautologiques
\[
\begin{tikzcd}[row sep=small]
 & \mathrm{a}' \arrow[r, Rightarrow] & \mathrm{a}'' & \\
\mathrm{a} \arrow[ur, Rightarrow] \arrow[dr, Rightarrow] & & & \\
 & \mathrm{b} \arrow[r, Rightarrow] & \mathrm{c}' \arrow[r, Rightarrow] & \mathrm{c}\ ,
\end{tikzcd}
\tag*{$(*)$}
\]
\printedpage{30}
% head: A. Grothendieck : Dérivateurs, Ch. XVIII, Édition des Universités de Montpellier II et Paris VII
% foot: Version 17 décembre 2025
où a$'$ $\Longrightarrow$ a$''$ provient de l'inclusion canonique $F(x) = \Hom_{\mathcal{M}}(x,F) \hookrightarrow \Ob C/F$, a $\Longrightarrow$ b s'obtient en prenant $I = C_0/F$ et notant que $I$ est stable par limites inductives de cardinal $\leqslant \pi_0$ ($C$ l'étant), donc est grande devant $\pi_0$ (\footnote{On n'a pas besoin de cette stabilité, de toutes façons $C_0/F$ est grande devant $\pi_0$, en vertu de 4.4.3, p. 80.}). On a b $\Longrightarrow$ c$'$ en vertu du fait que la sous-catégorie strictement pleine $C^{\wedge}_{\pi}$ de $C^{\wedge}$ définie par
\[
\left\{
\begin{array}{lcl}
C^{\wedge}_{\pi} & = & \text{sous-catégorie strictement pleine de } C^{\wedge} \text{ formée des}\\
 & & F \in \mathcal{M} \text{ qui sont } \pi\text{-accessibles.}
\end{array}
\right.
\tag*{(4.5.2.2)}
\]
est stable par $\varinjlim$ de cardinal $\leqslant \pi$ (SGA 4 I 9.9). Enfin c$'$ $\Longrightarrow$ c provient du fait que le foncteur d'inclusion $\mathcal{M} \hookrightarrow C^{\wedge}$ est pleinement fidèle et $\pi_0$-accessible, i.e. commute aux limites inductives grandes devant $\pi_0$, et a fortiori à celles grandes devant $\pi$. Par les implications du diagramme ci-dessus, il suffit donc de prouver qu'on a
% marginal mark (handwritten squiggle) in the left margin of the edition beside the following display
\[
\mathrm{c} \overset{(\mbox{\footnotemark[33]})}{\Longrightarrow} a \qquad\quad \text{et} \qquad\quad \mathrm{a}'' \ \Longrightarrow\ \mathrm{a}\ .
\]
\footnotetext[33]{n'est prouvé (et n'est vrai) que si $\pi^{\pi_0} = \pi$.}
\setcounter{footnote}{33}

\msp{page 109}

\textit{Prouvons} a$''$ $\Longrightarrow$ a. Considérons l'application
\[
\Fl C/F \ \longrightarrow\ \Fl C
\tag*{$(*)$}
\]
induit par le foncteur étale $C' = C/F \longrightarrow C$, alors l'ensemble des flèches au-dessus d'une % [sic] induit -> induite, accord avec « l'application »
flèche $\alpha : x \longrightarrow y$ de $C$ est en correspondance biunivoque avec $\Ob C'_y = F(y)$, puisqu'une telle flèche $\alpha'$ est connue quand on connaît son but, qui est un objet quelconque de $F(y)$ ($C' \longrightarrow C$ étant étale, i.e. fibrant à fibres discrètes). Donc le fibre de $(*)$ en $\alpha$ est de % [sic] le fibre -> la fibre, genre
cardinal $= \card F(y) \leqslant \pi$, donc on a
\[
\card\Fl C/F \ \leqslant\ \pi \underbrace{\card\Fl C}_{\leqslant \pi_1} \ \leqslant\ \pi \underbrace{\ \pi_1\ }_{\leqslant \pi} \ \leqslant\ \pi^2 = \pi\ ,
\]
ce qui prouve a.

Prouvons enfin c $\Longrightarrow$ a. Pour ceci, nous allons utiliser la première inclusion $\mathcal{M}_{\pi} \subset \Filt^{\pi}\mathcal{M}$ de SGA 4 I 9.16, qui s'applique dès lors que nous prouvons le

% marginal mark (handwritten squiggle) in the left margin of the edition beside Corollaire 4.5.3
\textbf{Corollaire 4.5.3} (\footnote{Pas prouvé, la condition c) des filtrations cardinales peut-être pas satisfaite. Sans doute \textit{faux} en général.}). \textit{Sous les conditions de 4.5.2, soit, pour tout cardinal $\pi \geqslant \pi_1$,} $\Filt^{\pi}(\mathcal{M})$ \textit{la sous-catégorie strictement pleine de $\mathcal{M}$ formée des $F \in \Ob\mathcal{M}$ qui satisfont la condition a) de 4.5.2. Alors la famille des} $(\Filt^{\pi})_{\pi\geqslant\pi_1}$ \textit{est une}
\msp{page 110}
\textit{filtration cardinale de $\mathcal{M}$.}

La condition a) de loc. cit. 9.12 (les $\uFilt^{\pi}$ sont équivalentes à des petites catégories) est évidente. Prouvons b), i.e. la stabilité de $\uFilt^{\pi}$ par $\varinjlim$ filtrantes $F = \varinjlim F_i$ \textit{grandes devant} $\pi_0$, de cardinal $\leqslant \pi$. Comme $\mathcal{M} \hookrightarrow C^{\wedge}$ commute à ces limites, pour $x$ dans $C$
\[
F(x) \ =\ \varinjlim_{i} F_i(x)\ ,
\]
\printedpage{31}
% head: A. Grothendieck : Dérivateurs, Ch. XVIII, Édition des Universités de Montpellier II et Paris VII
% foot: Version 17 décembre 2025
d'où $\card F(x) \leqslant \card\Ob I \cdot \pi \leqslant \pi \cdot \pi = \pi$, ce qui (par l'équivalence a $\Longleftrightarrow$ a$''$) implique $F \in \Ob\Filt^{\pi}$. Prouvons la condition c) de loc. cit. Pour ceci, nous aurons besoin du fait, plus précis que le b), qu'on vient d'établir :

\textbf{Corollaire 4.5.4} (\footnote{Prouvé seulement si $\pi^{\pi_0} = \pi$\,! Sans doute \textit{faux} en général.}). \textit{Pour tout $\pi \geqslant \pi_1$,} $\Filt^{\pi}$ \textit{est stable par $\varinjlim$ (dans $\mathcal{M}$) de cardinal $\leqslant \pi$ (pas nécessairement filtrantes, a fortiori pas nécessairement grandes devant $\pi_0$).}

Il suffit, en vertu de ce qui est déjà prouvé de du lemme 4.4.14 (\footnote{Il faut une variante du lemme 4.4.14, pour les $\varinjlim$ de cardinal $\leqslant \pi'$, qui se prouve de la même façon.}), de montrer que $\Filt^{\pi}$ % [sic] de du lemme -> du lemme, mot répété
est stable par $\varinjlim$ \textit{de cardinal} $\leqslant \pi_0$. Si $\varphi : J \longrightarrow \Filt^{\pi}$ est un foncteur avec $\card J \leqslant \pi_0$, on le met sous forme indicielle
\[
\psi\ :\ I \times J \ \longrightarrow\ C\ ,
\]
grâce à 4.14.3.b, avec $I$ grande devant $\pi_0$.

\msp{page 111}

Je rappelle qu'on pouvait prendre
\[
I \ =\ \uHom(J,C)/\varphi\ ,
\]
où $\uHom(J,C) \hookrightarrow \uHom(J, \mathcal{M} = \Ind'_{\pi}C)$, quitte à ne pas exiger que $I$ soit ordonnée. On a envie de prouver que $\card\Fl I \leqslant \pi$. L'ennui, c'est qu'on arrive seulement à la majoration
\[
\card\Fl I \ \leqslant\ \pi^{\pi_0}\ .
\]
(Même si $J$ est discrète, on n'a pas mieux, puisqu'alors
\[
I \ \simeq\ \prod_{j\in J} C/\varphi(j)\ , \qquad \text{d'où} \quad \Fl I = \prod_{j\in J} \Fl C/\varphi(j)\ ,
\]
et si $\card\Fl C/\varphi(j) = \pi$ pour tout $j$, on trouve $\card\Fl I = \pi^{\pi_0}$ -- il n'y a pas moyen de faire mieux\,!)

Donc je n'arrive pas à prouver directement 4.5.4, ce qui m'aurait arrangé pour prouver la condition c) de loc. cit. Mais à vrai dire, on a fait ici le travail de la démonstration du théorème 4.3.6, partie b $\Longrightarrow$ c -- on y a déjà prouvé que la filtration qu'on vient de construire (par la condition a) de 4.5.2) est une filtration cardinale. La condition c) de SGA 4 I 9.12 est prouvée à la page 76. Mais je vais refaire
\msp{page 112}
la démonstration plus sérieusement.

Soit donc $F \in \Ob\mathcal{M}$, considérons
\[
J \ =\ C/F\ ,
\]
c'est une catégorie stable par $\varinjlim$ de cardinal $\leqslant \pi_0$, donc grande devant $\pi_0$, et on a
\[
F \ \simeq\ \varinjlim_{j\in J} x_j \qquad\quad (\,\varinjlim \ \text{dans } \mathcal{M})
\]
\printedpage{32}
% head: A. Grothendieck : Dérivateurs, Ch. XVIII, Édition des Universités de Montpellier II et Paris VII
% foot: Version 17 décembre 2025
en désignant par $j \longmapsto x_j$ le foncteur canonique $C/F \longrightarrow C$. Soit alors $I$ l'ensemble des sous-catégories $J'$ de $J$ qui sont \textit{de cardinal $\leqslant \pi$ et grandes devant $\pi_0$, $I$} ordonné par inclusion, et pour tout $J' \in I$, posons
\[
X_{J'} \ =\ \varinjlim_{j\in J'} x_j\ .
\]
Comme $J'$ est grande devant $\pi_0$, la limite inductive dans $\mathcal{M}$ est aussi une limite dans $C^{\wedge}$, d'où résulte aussitôt, comme $\card\Ob J' \leqslant \card\Fl J' \leqslant \pi$, que $X_{J'}$ est dans $\Filt^{\pi}$. Alors
\[
J' \ \longmapsto\ X_{J'} \ \in\ \Ob\Filt^{\pi}
\]
est un foncteur
\[
I \ \longrightarrow\ \Filt^{\pi}\ ,
\]
dont la limite inductive est celle des $x_j$, donc isomorphe à $F$. D'autre part, il est immédiat que $I$ est grand devant $\pi$, i.e. que si $(J'_{\alpha})_{\alpha\in A}$ est
\msp{page 113}
une famille d'éléments de $I$, celle-ci est majorée par un $J' \in I$. Ceci résulte du

% marginal mark (handwritten squiggle) in the left margin of the edition beside Lemme 4.5.5
\textbf{Lemme 4.5.5} (\footnote{N'est prouvé \textit{que} si $\pi^{\pi_0} = \pi$\,! Sans doute faux en général.}). \textit{Soient $\pi$, $\pi_0$, avec $\pi > \pi_0$, des cardinaux infinis, et soit $J$ une catégorie grande devant $\pi_0$. Alors pour toute partie $\Phi$ de} $\Fl J$\textit{, telle que} $\card\Phi \leqslant \pi$\textit{, il existe une sous-catégorie $J'$ de $J$} grande devant $\pi_0$\textit{, telle que} $\card J' \leqslant \pi$ (\footnote{C'est OK si on exige seulement $\card J' \leqslant \pi^{\pi_0}$.})\textit{, et} $\Fl J' \supset \Phi$.

Ce lemme étant admis, on a donc représenté un objet quelconque $F$ de $\mathcal{M}$ comme une $\varinjlim\limits_{i\in I} X_i$, avec $I$ ensemble ordonné grand devant $\pi$, les $X_i$ dans $\Filt^{\pi}$. Il reste à prouver que si $\pi' > \pi$ et $F \in \Filt^{\pi'}$, alors on peut prendre $I$ tel que $\card I \leqslant \pi'^{\pi}$. Mais $F \in \Filt^{\pi'}$ signifie que $J = C/F$ satisfait $\card\Fl J \leqslant \pi'$. Alors l'ensemble des parties de cardinal $\leqslant \pi$ de $\Fl J$ a un cardinal majoré par $\pi'^{\pi}$, a fortiori $\card I \leqslant \pi'^{\pi}$, ce qui achève de prouver que $(\Filt^{\pi})_{\pi\geqslant\pi_1}$ est une filtration cardinale sur $\mathcal{M}$.

\msp{page 114}

Pour prouver le lemme, on commence par saturer $\Phi$ par composition finie, on trouve une sous-catégorie $J'_0$ de $J$ telle que $\card\Fl J'_0 \leqslant \pi_0$ \ldots{} Mais je m'aperçois que pour prouver le lemme, j'ai encore besoin de $\pi^{\pi_0} = \pi$ -- une enveloppe de $J'$ grande devant $\pi_0$ va avoir un cardinal majoré (au mieux) par $\pi^{\pi_0}$.

Prenons p.ex.
\[
J \ =\ \mathfrak{P}(S) \qquad\quad \text{(ensemble ordonné)}
\]
avec $\card S = \pi > \pi_0$, donc $J$ est grand devant $\pi_0$. Considérons l'application
\[
S \ \longrightarrow\ J = \mathfrak{P}(S)\ , \qquad\quad s \longmapsto \{s\}\ ,
\]
et prenons pour $\Phi \subset \Fl J$ l'ensemble des $\mathrm{id}_{\{s\}}$. Donc on cherche une partie $J'$ de $J$, contenant $S$, et grande devant $\pi_0$ pour l'ordre induit par $J$. On voit tout de suite que la plus petite partie $J'$ faisant l'affaire est l'ensemble $\mathfrak{P}_{\pi_0}(S)$ des parties de $S$ de cardinal $\leqslant \pi_0$. Mais on trouve
\[
\card J' \cdot \pi_0! \ =\ \pi^{\pi_0}\ ,
\]
\printedpage{33}
% head: A. Grothendieck : Dérivateurs, Ch. XVIII, Édition des Universités de Montpellier II et Paris VII
% foot: Version 17 décembre 2025
\msp{page 115}
où
\[
\begin{aligned}
\pi_0! \ &=\ \card\Aut\,(E_0) \qquad\quad (E_0 \text{ ensemble de cardinal } \pi_0)\\
&=\ \pi_0^{\pi_0}\\
&=\ 2^{\pi_0}\ ,
\end{aligned}
\]
donc
\[
\card J' \cdot 2^{\pi_0} \ =\ \pi^{\pi_0}\ .
\tag*{$(*)$}
\]
D'autre part, il est immédiat que $\card J' \geqslant \card\mathfrak{P}(E_0) = 2^{\pi_0}$ (où $E_0$ est une partie de $S$ de cardinal $\pi_0$), donc le premier membre de $(*)$ est égal à $\card J'$, donc on trouve
\[
\card J' \ =\ \pi^{\pi_0}\ ,
\]
sorry\,!

Finalement, je me suis convaincu que la proposition 4.5.2, et les corollaires et lemmes précédents par lesquels j'ai essayé de l'établir, n'est vrai que pour ces cardinaux $\pi > \pi_0$ tels que
\[
\pi^{\pi_0} \ =\ \pi\ ,
\]
ce qui, moyennant l'hypothèse du continu généralisée, équivaut à
\[
\pi^{\pi_0} \ <\ 2^{\pi}\ (\ =\ \pi^{\pi}\ )
\]
(ou encore à $\pi^{\pi_0} \neq 2^{\pi}$ $(= \pi^{\pi})$). Je vais donner
\msp{page 116}
ici un énoncé rectifié.

\textbf{Proposition 4.5.6.}
\begin{enumerate}[leftmargin=3em]
\item[\circled{1\textsuperscript{o}}] \textit{Sous les hypothèses générales de 4.5.2 sur $C$, $\pi_0$, $\pi_1$, $\pi$, les conditions a, a$'$, a$''$, b, c$'$ de cette proposition sont équivalentes (moyennant $\pi \geqslant \pi_1$)} (\footnote{\textbf{N.B.} Pour cette équivalence, l'hypothèse que $C$ soit stable par $\varinjlim$ de cardinal $\leqslant \pi$ est inutile, et aussi pour que c$'$) implique c), et lui soit équivalente si $\pi = \pi^{\pi_0} \geqslant \pi_1$, cf. 4.7.1, p. 161.})\textit{, et impliquent la condition c ($F$ est $\pi$-accessible). Quand $\pi^{\pi_0} = \pi$ (notamment quand $\pi = 2^{c}$, avec $c \geqslant \pi_0$), alors toutes les conditions a à c$'$ de 4.5.2 sont équivalentes.}
\item[\circled{2\textsuperscript{o}}] \textit{Plus généralement et plus précisément, pour $\pi \geqslant \pi_1$ donné, considérons les conditions sur $\pi$ :}
\begin{enumerate}[leftmargin=2em]
\item[\textit{1)}] $\pi^{\pi_0} = \pi$.
\item[\textit{2)}] $\Filt^{\pi}(\mathcal{M})$ \textit{(définie dans 4.5.3) est stable par $\varinjlim$ de cardinal $\leqslant \pi$.}
\item[\textit{3)}] \textit{$\forall$ $F$ dans $\mathcal{M}$, on a $F \simeq \varinjlim\limits_{I} X_i$, avec $I$ ensemble ordonné grand devant $\pi$, les $X_i$ dans} $\Filt^{\pi}(\mathcal{M})$\textit{. De plus, si $\pi' > \pi$ et $F \in$} $\Filt^{\pi'}(\mathcal{M})$\textit{, on peut prendre $I$ de cardinal $\leqslant \pi'^{\pi}$.}
\item[\textit{4)}] \textit{Comme 3), sans la majoration de cardinaux, et en supposant seulement $I$ catégorie grande devant $\pi$ (pas nécessairement catégorie ordonné).} % [sic] catégorie ordonné -> catégorie ordonnée, accord
\item[\textit{5)}] \textit{La condition c) de 4.5.2 implique les autres conditions (elles sont équivalentes, par 1°), en d'autres termes}
\[
\mathcal{M}_{\pi} \ \subset\ \Filt^{\pi} \qquad\quad \textit{(donc } \mathcal{M}_{\pi} = \Filt^{\pi}\textit{, puisqu'on sait que } \Filt^{\pi} \subset \mathcal{M}_{\pi}\textit{)}\ ,
\]
\end{enumerate}
\end{enumerate}


\endgroup
% END UNIT der-ch18/der-ch18_p031-035.tex
% BEGIN UNIT der-ch18/der-ch18_p036-040.tex
\begingroup\makeatletter%
\newcommand{\printedpage}[1]{\marginpar{\footnotesize p.~#1}}%
\newcommand{\msp}[1]{\marginpar{\footnotesize [#1]}}%
\newcommand{\Filt}{\operatorname{Filt}}%
\newcommand{\Ob}{\operatorname{Ob}}%
\newcommand{\card}{\operatorname{card}}%
\newcommand{\Fl}{\operatorname{Fl}}%
\newcommand{\Hom}{\operatorname{Hom}}%
\newcommand{\uHom}{\underline{\operatorname{Hom}}}%
\newcommand{\Ens}{\mathrm{Ens}}%
\newcommand{\Ind}{\operatorname{Ind}}%
\newcommand{\circled}[1]{\tikz[baseline=(c.base)]\node[draw,circle,inner sep=1pt](c){#1};}%
\makeatother\let\printedpage\derpp\let\msp\dermsp%
\let\mspin\DERmspin%

\setcounter{footnote}{39}

\printedpage{34}
% head: A. Grothendieck : Dérivateurs, Ch. XVIII, Édition des Universités de Montpellier II et Paris VII
% foot: Version 17 décembre 2025

\msp{page 117}

\begin{quote}
\textit{où $\mathcal{M}_\pi$ désigne la sous-catégorie strictement pleine de $\mathcal{M}$ formée des objets $\pi$-accessibles.}
\end{quote}

\begin{quote}
\textit{Ceci posé, on a les implications}
\[
\mathrm{1)} \ \Longrightarrow\ \mathrm{2)} \ \Longrightarrow\ \mathrm{3)} \ \Longrightarrow\ \mathrm{4)} \ \Longrightarrow\ \mathrm{5)} \ ,
\]
\textit{et dans le cas particulier où $\mathcal{M} = $} (Ens), \textit{$C$ = sous-catégorie pleine formée des ensembles de cardinal $\leqslant \pi_0$, on a aussi} 5) $\Longrightarrow$ 1), \textit{i.e. toutes ces conditions sont équivalentes.}
\end{quote}

\textsc{Démonstration.}

1) $\Longrightarrow$ 2). Cf. démonstration de 4.5.4.

2) $\Longrightarrow$ 3). La catégorie $J = \Filt^\pi/F$ est stable par $\varinjlim$ de cardinal $\leqslant \pi$, $\Filt^\pi$ l'étant, donc $J$ est grande devant $\pi$, d'ailleurs $F = \varinjlim\limits_{j\in J} X_j$, où $j \longmapsto X_j$ est le foncteur source $\Filt^\pi/F \longrightarrow \Filt^\pi$. D'où la représentation voulue de $F$, quitte à remplacer $J$ par un ensemble ordonné cofinal $I$, p.ex. en prenant pour $I$ l'ensemble des sous-catégories $J$ de $I$ telles que $\card \Fl J \leqslant \pi$ \texttt{[cette deuxième partie de la phrase “p.ex. ...” demande correction]}. La majoration $\card I \leqslant \pi'^{\pi}$, pour $F \in \Ob \Filt^{\pi'}$, se fait par AQT \texttt{[âne qui trotte]}.

L'implication 3) $\Longrightarrow$ 4) est tautologique.

4) $\Longrightarrow$ 5) (\footnote{comparer le lemme 4.5.14, page 133.}), car si $F \in \Ob \mathcal{M}_\pi$, on trouve par 4) que $F$ est facteur direct d'un objet de $\Filt^\pi$, et $\Filt^\pi$ étant stable par $\varinjlim$
\msp{page 118}
grandes devant $\pi_0$, cela implique que $\Filt^\pi$ est stable par facteurs directs, d'où $F \in \Ob \Filt^\pi$.

Enfin, si $\mathcal{M} = \Ens$, $C = \Ens(\pi_0)$. On note que
\[
\Filt^\pi \ = \ \{ \ X \in \Ob \Ens \mid \card\,(X)^{\pi_0} \ \leqslant \ \pi \ \} \ .
\]
D'autre part, on trouve facilement
\[
\begin{aligned}
(\Ens)_\pi \quad &\overset{\text{déf}}{=} \quad \text{sous-catégorie pleine de (Ens) formée} \\
&\phantom{\overset{\text{déf}}{=}} \quad \text{des objets $\pi$-accessibles} \\
&= \quad \{ X \in \Ob \Ens \mid \card X \leqslant \pi \} \ .
\end{aligned}
\]
Donc on a $(\Ens)_\pi \subset \Filt^\pi$ si et seulement si $\pi^{\pi_0} \leqslant \pi$, i.e. $\pi^{\pi_0} = \pi$, q.e.d. (\footnote{Cf. démonstration de la partie 1°) de 4.5.6, plus bas (bas de page et page suivante).}).

\textbf{Corollaire 4.5.7.} \textit{Pour tout cardinal $\pi \geqslant \pi_1$, on a}
\[
\Filt^\pi \ \subset \ \mathcal{M}_\pi \ \subset \ \Filt^{\pi^{\pi_0}} \ . \tag*{(4.5.7.1)}
\]

L'inclusion $\Filt^\pi \subset \mathcal{M}_\pi$ est l'implication a $\Longrightarrow$ c de la partie 1°) de 4.5.6. Pour prouver $\mathcal{M}_\pi \subset \Filt^{\pi^{\pi_0}}$, on pose $\pi' = \pi^{\pi_0}$ et on note que $\pi'^{\pi_0} = \pi^{\pi_0 \cdot \pi_0} = \pi^{\pi_0} = \pi'$. Donc on peut
\printedpage{35}
% head: A. Grothendieck : Dérivateurs, Ch. XVIII, Édition des Universités de Montpellier II et Paris VII
% foot: Version 17 décembre 2025
appliquer 1° à $\pi'$ au lieu de $\pi$, et on trouve que $\mathcal{M}_{\pi'} = \Filt^{\pi'}$. Comme $\mathcal{M}_\pi \subset \mathcal{M}_{\pi'}$, on en conclut $\mathcal{M}_\pi \subset \Filt^{\pi'}$, q.e.d.

J'ai oublié de démontrer le 1°) dans 4.5.6. Vu le diagramme d'implications $(\ast)$, p. 108, et l'implication $\mathrm{a}'' \Longrightarrow \mathrm{a}$ déjà prouvé (p. 109),
\msp{page 119}
il reste à prouver seulement
\[
\mathrm{c}' \ \Longrightarrow \ \mathrm{a}'' \ ,
\]
i.e. que $F$ $\pi$-accessible \textit{dans} $C^\wedge$ implique $\card F(x) \leqslant \pi \ \forall \ x \in C$. Ceci va résulter du résultat un peu plus général :

\textbf{Lemme 4.5.8.}
\begin{itemize}[leftmargin=3em]
\item[\textit{1°)}] \textit{Dans la catégorie $\mathcal{E} = $} (Ens), \textit{les objets $\pi$-accessibles ($\pi$ un cardinal infini) sont les ensembles $X$ tels que} $\card X \leqslant \pi$.
\item[\textit{2°)}] (\footnote{Cf. généralisation 4.14.18, p. 391 b) et c).}). \textit{Soit $C$ une catégorie équivalente à une petite catégorie, telle que $\forall \ x, y \in \Ob C$, on ait} $\card \Hom(x,y) \leqslant \pi$. \textit{Soit $\mathcal{E}$ une $\mathfrak{U}$-catégorie stable par $\varinjlim$ grandes devant $\pi$, et faiblement $\pi$-accessible (i.e. admettant une sous-catégorie génératrice formée d'objets $\pi$-accessibles), et stable par produits de cardinal $\leqslant \pi$. Alors pour qu'un objet $F$ de $\uHom(C,\mathcal{E})$ soit $\pi$-accessible, il} faut \textit{que $\forall \ x \in \Ob C$, $F(x)$ soit un objet $\pi$-accessible de $\mathcal{E}$, et cette condition est aussi suffisante si} $\card \operatorname{\text{réd}} \Fl(C) \leqslant \pi$ (\footnote{même \textit{sans} condition de stabilité de $\mathcal{E}$ par produits de cardinal $\leqslant \pi$.}).
\end{itemize}

\textsc{Démonstration.} 1°) $\card X \leqslant \pi \Longrightarrow X$ $\pi$-accessible dans $\mathcal{E} = (\Ens)$, car $X = \coprod_{x \in X} \{x\}$, et on applique SGA 4 I 9.9. Inversement, supposons $X$ $\pi$-accessible, et écrivons $X \simeq \varinjlim\limits_{i \in I} X_i$, où $I$ est l'ensemble des parties de $X$ qui sont de cardinal $\leqslant \pi$. Comme $I$ est grande devant $\pi$, l'application identique $X \xrightarrow{\ \mathrm{id}\ } X$ se factorise par un ensemble $X_i$, donc $X_i = X$, donc $\card X = \card X_i \leqslant \pi$, q.e.d.

\msp{page 120}

2°) Considérons, pour $x \in \Ob C$, le foncteur
\[
e \xrightarrow{\ \varepsilon_x\ } C \ ,
\]
d'où un foncteur
\[
\varepsilon_x^* \ : \ \uHom(C,\mathcal{E}) \longrightarrow \mathcal{E} \ , \qquad F \longmapsto F(x) \ ,
\]
admettant un adjoint à droite
\[
\varepsilon_{x\,*} \ : \ \mathcal{E} \ \longrightarrow \ \uHom(C,\mathcal{E}) \ ,
\]
donné par
\[
\varepsilon_{x\,*}(E)(y) \ = \ E^{\Hom(x,y)}
\]
(sous réserve d'existence des produits de cardinal $\leqslant \pi$). Ce foncteur $\varepsilon_{x\,*}$ est $\pi$-accessible en vertu de loc. cit. 9.8.

Ceci posé, soit $F$ un objet $\pi$-accessible de $\uHom(C,\mathcal{E})$, on voit que $\varepsilon_x^*(F)$ est $\pi$-accessible, $\varepsilon_x^*$ ayant un adjoint à droite $\pi$-accessible :
\[
E \ \longmapsto \ \Hom_{\mathcal{E}}(\underbrace{F(x)}_{\varepsilon_x^*(F)}, E) \ \simeq \ \Hom_{\uHom(C,\mathcal{E})}(F, \varepsilon_{x\,*}(E))
\]
\printedpage{36}
% head: A. Grothendieck : Dérivateurs, Ch. XVIII, Édition des Universités de Montpellier II et Paris VII
% foot: Version 17 décembre 2025
est composé du foncteur $\pi$-accessible $\varepsilon_{x\,*}$, et du foncteur covariant $\pi$-accessible représenté par $F$ dans $\uHom(C,\mathcal{E})$.

Inversement, si les $F(x)$ sont $\pi$-accessibles, prouvons que $F$ l'est, moyennant $\card \Fl(C) \leqslant \pi$. Se fait par AQT \texttt{[âne qui trotte]}, je ne vais pas l'expliciter ici.

\msp{page 121}

\textbf{Proposition 4.5.9.} \textit{Soit $C$ une catégorie équivalente à une petite catégorie, et soit $\mathcal{M} = $} $\Ind'_{\pi_0}(C)$, \textit{$\pi_0$ un cardinal infini donné. On suppose $C$ stable par $\varinjlim$ de cardinal $\leqslant \pi_0$. Pour tout cardinal $\pi \geqslant \pi_0$, soit $\mathcal{M}_\pi \supset C$ la sous-catégorie pleine de $\mathcal{M}$ formée des objets $\pi$-accessibles de $\mathcal{M}$. Soit d'autre part} $\pi_1 = \sup(\pi_0, \card \Fl C)$. \textit{Alors $(\mathcal{M}_\pi)_{\pi \geqslant \pi_1}$ est une filtration cardinale de $\mathcal{M}$} (\footnote{\textbf{N.B.} Toutes les conditions des filtrations cardinales sont vérifiées à l'exclusion, quand $\pi \geqslant \pi_0$, de celle de la majoration de cardinaux $\card I \leqslant \pi'^{\pi}$ condition c).})\textit{, et de plus, pour $\pi \geqslant \pi_0$, $\mathcal{M}_\pi$ est stable dans $\mathcal{M}$ par $\varinjlim$ de cardinal $\leqslant \pi$.}

\textsc{Démonstration.} La stabilité de $\mathcal{M}_\pi$ est un rappel, cf. SGA 4 I 9.9. Comme on a
\[
\mathcal{M}_\pi \ \subset \ \Filt^{\pi^{\pi_0}}
\]
pour $\pi \geqslant \pi_1$ (cf. 4.5.7), on voit donc que les $\mathcal{M}_\pi$ sont équivalentes à des petites catégories (condition a) des filtrations cardinales). La condition b) des filtrations cardinales est contenue dans la propriété de stabilité des $\mathcal{M}_\pi$ qu'on vient de rappeler. Prouvons enfin la condition c). Soit donc $\pi \geqslant \pi_0$, et soit $x$ dans $\mathcal{M}$, considérons la catégorie
\[
I \ = \ \mathcal{M}_\pi/x \ ,
\]
elle est stable par $\varinjlim$ de cardinal $\leqslant \pi$, d'après la stabilité similaire de $\mathcal{M}_\pi$, donc $I$ est grande devant $\pi$ et équivalente à une petite catégorie, d'autre part, on a
\msp{page 122}
\[
x \ = \ \varinjlim_{i \in I} x_i \ ,
\]
les $x_i$ dans $\mathcal{M}_\pi$ ( $\varinjlim$ dans $\mathcal{M}$ des foncteurs $I = \mathcal{M}_\pi/x \longrightarrow \mathcal{M}_\pi$). Quitte à remplacer $I$ par une catégorie ordonnée cofinale, cela prouve la partie qualitative de c. Supposons enfin $x \in \Ob \mathcal{M}_{\pi'}$, où $\pi' > \pi$, montrons que l'on a alors
\[
\card I_0, \ \card I' \ \leqslant \ \pi'^{\pi} \qquad \text{(\footnotemark, \footnotemark)} \ , \tag*{(4.5.9.1)}
\]
\addtocounter{footnote}{-1}%
\footnotetext{où $I_0$, $I'_0$ sont des sous-catégories réduites de $I$, $I'$ telles que $I_0 \hookrightarrow I$, $I'_0 \hookrightarrow I'$ soient des équivalences \texttt{[Grothendieck a biffé l'indice 0 de $I'$ dans 4.5.9.1]}.}%
\stepcounter{footnote}%
\footnotetext{On ne le prouve que si $\pi \geqslant \pi_1$.}%
ce qui achèvera de prouver que l'on a une filtration cardinale. Il nous faut une majoration des cardinaux.

\textbf{Lemme 4.5.10} (\footnote{Donner énonce commun avec 4.5.11, et prouver d'abord 4.5.11.}). \textit{Soit $\pi$ un cardinal $\geqslant \pi_1$. Alors on a} % [sic] énonce -> énoncé, missing accent (in footnote 47)
\[
\card \Fl((\Filt^\pi)_0) \ \leqslant \ 2^\pi \ . \tag*{(4.5.10.1)}
\]

On a en effet, en désignant par $C^\wedge(\pi)$ la sous-catégorie strictement pleine de $C^\wedge$ formée des $F : C^\circ \longrightarrow \Ens$ avec $\card F(x) \leqslant \pi$ pour tout $x$ dans $C$ (i.e. $F$ se factorise par $\Ens(\pi)$), et par $C^\wedge(\pi)_0$ une
\printedpage{37}
% head: A. Grothendieck : Dérivateurs, Ch. XVIII, Édition des Universités de Montpellier II et Paris VII
% foot: Version 17 décembre 2025
\msp{page 123}
sous-catégorie réduite exhaustive,
\[
\card C^\wedge(\pi)_0 \ \leqslant \ 2^\pi \ , \tag*{(4.5.10.2)}
\]
d'où résulte (4.5.10.1), puisque pour $\pi \geqslant \pi'$, $\Filt^\pi$ est une sous-catégorie pleine de $C^\wedge(\pi')$. Pour établir (4.5.10.2), on note que l'ensemble des classes à isomorphie près de $\Ens(\pi')$ s'identifie à l'ensemble des cardinaux $\leqslant \pi'$, lequel ensemble est de cardinal $\leqslant \pi'$. Remplaçant $\Ens(\pi)$ par $\Ens(\pi')_0$, $C$ par $C_0$, on doit majorer $\card \Fl(\uHom(C_0, \Ens(\pi')_0))$. Or $\Ob \uHom(C_0, \Ens(\pi')_0)$ est contenu dans l'ensemble des applications de $\Fl C_0$ dans $\Fl(\Ens(\pi')_0)$, le premier est majoré par $\pi_1$, le deuxième par $\pi^\pi \times \pi \times \pi = \pi^{2\pi} = \pi^\pi = 2^\pi$ (\textbf{N.B.} $\pi < 2^\pi$, d'où $\pi^\pi \leqslant (2^\pi)^\pi = 2^{\pi^2} = 2^\pi$), donc $\Ob \uHom$ majoré par $(2^\pi)^{\pi_1} = 2^{\pi\pi_1} = 2^\pi$. D'autre part, si $F$, $G$ sont dans $\uHom(C_0, \Ens(\pi_0))$, $\Hom(F,G)$ est majoré par $(\pi^\pi)^{\pi_1} = \pi^\pi = 2^\pi$, donc $\Fl \uHom$ est majoré par $2^\pi \times 2^\pi \times 2^\pi$, q.e.d.

\textbf{Corollaire 4.5.11.} \textit{Soit $F \in \Ob \Filt^\pi$, $F' \in \Ob \Filt^{\pi'}$, avec $\pi' \geqslant \pi \geqslant \pi_1$. Alors}
\[
\card \Hom(F, F') \ \leqslant \ \pi'^{\pi} \ . \tag*{(4.5.11.1)}
\]

\msp{page 124}

C'est essentiellement la même démonstration, légèrement plus simple. (On aurait dû commencer par là.)

\textbf{Corollaire 4.5.12.} \textit{Si $\pi' \geqslant \pi \geqslant \pi_1$, on a pour $x \in \Ob \mathcal{M}_{\pi'}$}
\[
\card\,(\mathcal{M}_\pi/x)_0 \ \leqslant \ \pi'^{(\pi^{\pi_0})} \ . \tag*{(4.5.12)}
\]

En effet,
\[
\card \Ob\,(\mathcal{M}_\pi/x)_0 \ \overset{\text{par 4.5.11}}{\leqslant} \ \underbrace{\underbrace{\card \Ob\,(\mathcal{M}_\pi)_0}_{\leqslant 2^{(\pi^{\pi_0})} \text{par 4.5.10}} \times (\pi'^{\pi_0})^{\pi^{\pi_0}}}_{= \pi'^{(\pi^{\pi_0})}} \ ,
\]
i.e.
\[
\card \Ob\,(\mathcal{M}_\pi/x)_0 \ \leqslant \ \pi'^{(\pi^{\pi_0})} \ ,
\]
d'où
\[
\card \Fl(\mathcal{M}_\pi/x)_0 \ \leqslant \ \pi'^{(\pi^{\pi_0})} \times \pi'^{(\pi^{\pi_0})} \times (\pi'^{\pi_0})^{(\pi^{\pi_0})} \ \leqslant \ \pi'^{(\pi^{\pi_0})} \ .
\]
Cela donne donc (revenant à 4.5.9.1)
\[
\card I_0 \ \leqslant \ \pi'^{(\pi^{\pi_0})} \ .
\]

D'autre part, $I'$ peut se définir en termes de $I_0$ comme l'ensemble des sous-catégories de $I_0$ qui sont de cardinal $\leqslant \pi$, donc leur nombre est majoré par $(\card I_0)^\pi \ \leqslant \ (\pi'^{\pi^{\pi_0}})^\pi = \pi'^{\pi\pi^{\pi_0}} = \pi'^{(\pi^{\pi_0})}$.

Mais ce n'est pas tout à fait la majoration exigée, puisque
\[
\pi'^{\pi} \leqslant \pi'^{(\pi^{\pi_0})} \qquad \text{(\footnotemark)},
\]
\footnotetext{\textbf{N.B.} On a égalité si on a $\pi' = 2^{c'}$ avec $c' \geqslant \pi^{\pi_0}$, soit $\pi = 2^c$ avec $c \geqslant \pi_0$.}%
et qu'il est possible, a priori, que l'inégalité
\printedpage{38}
% head: A. Grothendieck : Dérivateurs, Ch. XVIII, Édition des Universités de Montpellier II et Paris VII
% foot: Version 17 décembre 2025
\msp{page 125}
soit stricte ! Apparemment, les ``majorations''
\[
\mathcal{M}_\pi \ \subset \ \Filt^{\pi^{\pi_0}} \ , \qquad \mathcal{M}_{\pi'} \ \subset \ \Filt^{\pi'^{\pi_0}}
\]
sont trop grossières. Aussi il s'agit de mieux appréhender les catégories $\mathcal{M}_\pi$, pour $\pi \geqslant \pi_0$ (\footnote{Mais on peut utiliser 4.5.10, 4.5.11 pour conclure que pour $x \in \Ob \mathcal{M}_{\pi'} \subset \Ob \Filt^{\pi'^{\pi_0}}$, on a $\card \Fl(C_0/x) \leqslant (\pi'^{\pi_0})^{\pi_0} = \pi'^{\pi_0}$, on aura donc $x = \varinjlim\limits_{J} x_j$ ($x_j \in \Ob C$, $\varinjlim$ filtrante, $\card \Fl J \leqslant \pi'^{\pi_0}$), et en introduisant l'ensemble $I$ des sous-catégories $J' \subset J$ qui sont de cardinal $\leqslant \pi$, et posant $x_{J'} = \varinjlim\limits_{j \in J'} x_j$, on trouve $x = \varinjlim\limits_{J' \in I} x_{J'}$, où cette fois $I$ est un \textit{ensemble ordonné} grand devant $\pi$, les $x_{J'} \in \mathcal{M}_\pi$, et $\card I \leqslant (\card J)^\pi \leqslant (\pi'^{\pi_0})^\pi = \pi'^{\pi\pi_0} = \pi'^{\pi}$, c'est la bonne majoration !}). Je vais hasarder une

\textbf{Proposition 4.5.13.} \textit{Les hypothèses et notations sont celles de 4.5.9. Soit $\pi$ un cardinal $\geqslant \pi_0$, et soit $x \in \Ob \mathcal{M}$. Considérons les conditions :}
\begin{itemize}[leftmargin=3em]
\item[\textit{a})] \textit{Il existe une catégorie $I$, avec} $\card \Fl I \leqslant \pi$, \textit{et un foncteur $i \longmapsto x_i : I \longrightarrow C$, tel que $x$ soit isomorphe à $\varinjlim x_i$.}
\item[\textit{a}$'$)] \textit{Comme a), avec $I$ filtrante.}
\item[\textit{a}$''$)] \textit{Comme a), avec $I$ ordonnée filtrante.}
\item[\textit{b})] \textit{Il existe un foncteur} cofinal \textit{$I \longrightarrow C/x$, avec} $\card \Fl I \leqslant \pi$.
\item[\textit{b}$'$)] \textit{La catégorie filtrante $C/x$ est ``essentiellement de cardinal $\leqslant \pi$'', i.e. il existe une catégorie ordonnée} \texttt{[non, cf. b$''$)]} filtrante $I$, $\card \Fl I \leqslant \pi$, \textit{et un foncteur cofinal $I \longrightarrow C/x$.}
\item[\textit{b}$''$)] \textit{Comme b$'$), avec de plus $I$ ordonnée.}
\item[\textit{c})] \textit{Il existe une famille de flèches $(u_\alpha)_{\alpha \in A}$, $x_\alpha \longrightarrow x$, épimorphique stricte, indexée par un ensemble $A$ tel que} $\card A \leqslant \pi$.
\item[\textit{d})] \textit{On a $x \in \Ob \mathcal{M}_\pi$, i.e. $x$ est $\pi$-accessible dans $\mathcal{M}$.}
\end{itemize}

\msp{page 126}

\textit{Ceci posé :}
\begin{itemize}[leftmargin=3em]
\item[\circled{1\textsuperscript{\textit{o}}}] \textit{Les six premières conditions $a$, $a'$, $a''$, $b$, $b'$, $b''$ sont équivalentes, soit $P(\pi)$, et impliquent les deux dernières,}
\[
\begin{tikzcd}[column sep=small]
 & P(\pi) \arrow[dl, Rightarrow] \arrow[dr, Rightarrow] & \\
c & & d \ .
\end{tikzcd}
\]
\textit{Si $\pi \geqslant \pi_1$, alors $P(\pi)$ est équivalente à $c$.}
\item[\circled{2\textsuperscript{\textit{o}}}] \textit{Soit $\Filt'^{\pi} \subset \mathcal{M}$ la sous-catégorie strictement pleine de $\mathcal{M}$ formée des objets qui satisfont la condition $P(\pi)$ ci-dessus, et $\Filt^\pi$ la sous-catégorie strictement pleine définie dans 4.5.3 (p. 109). On a alors}
\[
\Filt^\pi \ \subset \ \Filt'^{\pi} \ \subset \ \mathcal{M}_\pi \ , \tag*{(4.5.13.1)}
\]
\textit{et pour qu'on ait égalité des trois membres, il suffit qu'on ait $\pi^{\pi_0} = \pi$. (Cf. prop. 4.5.6, page 116). De plus, $\Filt'^{\pi}$ est stable dans $\mathcal{M}$ par} sommes \textit{de cardinal $\leqslant \pi$.}
\item[\circled{3\textsuperscript{\textit{o}}}] \textit{Les conditions suivantes sur $C$, $\pi_0$, $\pi$ sont équivalentes.}
\end{itemize}


\endgroup
% END UNIT der-ch18/der-ch18_p036-040.tex
% BEGIN UNIT der-ch18/der-ch18_p041-045.tex
\begingroup\makeatletter%
\newcommand{\printedpage}[1]{\marginpar{\footnotesize p.~#1}}%
\newcommand{\msp}[1]{\marginpar{\footnotesize [#1]}}%
\newcommand{\Filt}{\operatorname{Filt}}%
\newcommand{\card}{\operatorname{card}}%
\newcommand{\Fl}{\operatorname{Fl}}%
\newcommand{\Ob}{\operatorname{Ob}}%
\newcommand{\Ind}{\operatorname{Ind}}%
\makeatother\let\printedpage\derpp\let\msp\dermsp%
\let\mspin\DERmspin%

\setcounter{footnote}{49}

\printedpage{39}
% head: A. Grothendieck : Dérivateurs, Ch. XVIII, Édition des Universités de Montpellier II et Paris VII

\begin{itemize}
\item[\textit{1)}] $\Filt'^{\pi}$ \textit{est stable dans $\mathcal{M}$ par} $\varinjlim$ \textit{de cardinal $\leqslant \pi$.}
\item[\textit{1$'$)}] $\Filt'^{\pi}$ \textit{est stable par conoyaux de double-flèches.}
\item[\textit{2)}] \textit{Tout objet $x$ de $\mathcal{M}$ est de la forme $\varinjlim\limits_{I} x_i$, où $x_i$ dans} $\Filt'^{\pi}$\textit{, $I$ grande devant $\pi$.}
\item[\textit{3)}] \textit{La condition d) de 1°) implique $P(\pi)$ (donc lui est équivalente), i.e. on a}
\[ \mathcal{M}_{\pi} \ = \ \Filt'^{\pi}. \]
\end{itemize}

\msp{page 127}
\textit{Ces conditions sont satisfaites dans les trois cas suivants :}
\begin{itemize}
\item[\textit{A)}] \textit{$\pi^{\pi_0} = \pi$ (pour mémoire, c'est contenu dans 2°).}
\item[\textit{B)}] \textit{Les familles épimorphiques strictes dans $\mathcal{M}$ sont épimorphiques strictes universelles.}
\item[\textit{C)}] \textit{Les limites inductives filtrantes dans $\mathcal{M}$ commutent au changement de base (ce qui est le cas notamment si les} $\varinjlim$ \textit{filtrantes sont exactes à gauche)} (\footnote{Je suppose que cette condition revient à la même condition sur $C$, pour les $\varinjlim$ filtrantes de cardinal $\leqslant \pi$.}).
\end{itemize}

\textbf{Remarque.} La filtration de $\mathcal{M}$ pour les $\Filt'^{\pi}$ n'est autre que celle envisagée dans SGA 4 I 9.13, où on fait de plus l'hypothèse B) sur $\mathcal{M}$, qui après-coup n'était plus tellement de mon goût. Toujours est-il que d'après (4.5.13.1), cette filtration se rapproche plus de la ``bonne'' filtration par les $\mathcal{M}_{\pi}$, que celle que j'avais eu encore naguère tendance à lui préférer, celle par les $\Filt^{\pi}$. Elle a d'ailleurs sur cette dernière l'avantage, dans le cas envisagé B), d'être une filtration cardinale (ce qui n'est pas nécessairement le cas pour $\Filt^{\pi}$, du moins s'il existe des cardinaux infinis $\pi \geqslant \pi_0$ avec $\pi^{\pi_0} \neq \pi$). Je soupçonne maintenant qu'en général (et s'il existe des cardinaux infinis $\pi > \pi_0$ tels que $\pi^{\pi_0} \neq \pi$), en dehors des conditions A) B) C), la filtration par les $(\Filt'^{\pi})$ n'est \textit{pas} une
\msp{page 128}
filtration cardinale. En tous cas, 4.5.6 et la présente proposition, 3°, montrent que si soit $(\Filt^{\pi})_{\pi\geqslant\pi_2}$, soit $(\Filt'^{\pi})_{\pi\geqslant\pi_2}$ est une filtration cardinale ($\pi_2$ étant un cardinal donné $\geqslant \pi_1$), alors elle coïncide avec la filtration par les $\mathcal{M}_{\pi}$. Donc apparemment, dans tous les cas où je sais construire une filtration cardinale, sur un paratopos du moins (i.e. sur une catégorie $\mathcal{M}$ stable par petites $\varinjlim$ , et dont tous les objets sont accessibles -- l'existence d'une filtration cardinale impliquant de plus celle d'une petite sous-catégorie strictement génératrice \dots), celle-ci n'est autre que celle par les $\mathcal{M}_{\pi}$ -- chose dont je ne m'étais par aperçu en rédigeant SGA 4 I 9. % [sic] m'étais par aperçu -> m'étais pas aperçu, "par" for the negation "pas"

\textsc{Démonstration de 4.5.13.}

\textit{Prouvons 1°.} On a les implications tautologiques
\[
\begin{tikzcd}[row sep=large]
\mathrm{b}'' \arrow[r, Rightarrow] \arrow[d, Rightarrow] & \mathrm{b}' \arrow[r, Rightarrow] \arrow[d, Rightarrow] & \mathrm{b} \arrow[d, Rightarrow] \\
\mathrm{a}'' \arrow[r, Rightarrow] & \mathrm{a}' \arrow[r, Rightarrow] & \mathrm{a} \text{ ,}
\end{tikzcd}
\qquad\qquad \text{(\footnotemark) ,}
\]
\footnotetext{compte tenu que $C$ est strictement génératrice dans $\mathcal{M}$, donc $x \simeq \varinjlim\limits_{i\in C/x} x_i$.}
et pour prouver l'équivalence de ces six conditions, il suffit de prouver
\[ \mathrm{a} \ \Longrightarrow \ \mathrm{b}'' \ . \]
% foot: Version 17 décembre 2025

\printedpage{40}
% head: A. Grothendieck : Dérivateurs, Ch. XVIII, Édition des Universités de Montpellier II et Paris VII

\msp{page 129}
Si on a
\[ x \ \simeq \ \varinjlim_{I} x_i \ , \qquad\quad \card \Fl I \ \leqslant \ \pi \ , \]
soit $I'$ l'ensemble des \textit{sous-diagrammes} $J \subset I$ finis de $I$, ordonné par inclusion, donc $I'$ est un ensemble ordonné filtrant, et du fait que $C$ est stable par $\varinjlim$ finies, on trouve un foncteur
\[ J \longmapsto x_J \ = \ \varinjlim_{i\in J} x_i \ : \ I' \longrightarrow C \ , \]
et on aura
\[ x \ \simeq \ \varinjlim_{J\in I'} x_J \ . \]
D'autre part, on a
\[
\begin{aligned}
\card I' \ &= \ \card \mathfrak{P}_f( \ \underbrace{\Fl(I)}_{\text{de card $\leqslant \pi$}} \ ) \\
&\leqslant \ \pi + \pi^2 + \ \cdots \ + \pi^n + \ \dots \\
&\leqslant \ \pi\aleph_0 \ = \ \pi \ ,
\end{aligned}
\]
q.e.d. Pour établir 1°, il reste à établir
\[
\begin{tikzcd}[column sep=small]
& P(\pi) \arrow[dl, Rightarrow] \arrow[dr, Rightarrow] & \\
\mathrm{c} & & \mathrm{d} \text{ ,}
\end{tikzcd}
\]
et l'implication inverse $\mathrm{c} \Longrightarrow P(\pi)$ quand $\pi \geqslant \pi_1$. Mais $P(\pi) \Longrightarrow \mathrm{c}$ est une tautologie, et $P(\pi) \Longrightarrow \mathrm{d}$ puisque $\mathcal{M}_{\pi}$ est stable par $\varinjlim$ de cardinal $\leqslant \pi$. Reste à prouver que $\mathrm{c} \Longrightarrow P(\pi)$ si $\pi \geqslant \pi_1$.

\msp{page 130}
Supposons donc qu'on ait une famille épimorphique stricte
\[ x_{\alpha} \longrightarrow x \qquad\quad (\alpha \in A,\ \card A \leqslant \pi) \ , \]
cela signifie donc que $x$ est limite inductive d'un système inductif indexé par la catégorie $B$ ayant comme ensemble d'objets $A \amalg (A\times A)$, les flèches non identiques étant représentées par les flèches canoniques
\[
\begin{tikzcd}[column sep=small]
& x_{\alpha} \times_x x_{\beta} \arrow[dl, "p_{\alpha,\beta}"'] \arrow[dr, "q_{\alpha,\beta}"] & \\
x_{\alpha} & & x_{\beta} \text{ .}
\end{tikzcd}
\]

(\textbf{N.B.} $B$ n'est par une catégorie ordonnée, car pour l'objet $(\alpha,\alpha)$, représenté par $x_{\alpha}\times_x x_{\alpha}$, % [sic] n'est par une -> n'est pas une, "par" for the negation "pas"
il y a \textit{deux} flèches dans l'objet $\alpha$, représenté par $x_{\alpha}$. Mais on a évidemment $\card \Fl(B) \leqslant \pi$.) Si on savait que $C$ est stable par produits binaires et par passage à des sous-objets (ici les $x_{\alpha} \times_x x_{\beta} \ \subset \ x_{\alpha} \times x_{\beta}$), on aurait prouvé la condition a). À défaut, utilisons encore le fait que $C$ est génératrice par épimorphismes (même stricts) et choisissons pour tout couple $(\alpha,\beta)$ un ensemble épimorphique de flèches :
\[ \underbrace{x_{\alpha,\beta,\gamma}}_{\in \Ob C_0} \xrightarrow{u_{\alpha,\beta,\gamma}} x_{\alpha} \times_x x_{\beta} \ , \]
% foot: Version 17 décembre 2025

\printedpage{41}
% head: A. Grothendieck : Dérivateurs, Ch. XVIII, Édition des Universités de Montpellier II et Paris VII
où $C_0$ est une sous-catégorie réduite exhaustive de $C$. La flèche $u_{\alpha,\beta,\gamma}$ est connue quand on connaît les composés
\msp{page 131}
\[
\begin{tikzcd}
& x_{\alpha} \\
x_{\alpha,\beta,\gamma} \arrow[ur, "v_{\alpha,\beta,\gamma}"] \arrow[dr, "w_{\alpha,\beta,\gamma}"'] & \\
& x_{\beta} \text{ ,}
\end{tikzcd}
\]
donc pour $x_{\alpha,\beta,\gamma} \in \Ob C_0$ fixé, les choix possibles pour $u_{\alpha,\beta,\gamma}$ sont de cardinal majoré par $\pi_1\pi_1 = \pi_1$, de plus on a $\card \Ob C_0 \leqslant \pi_1$, donc le nombre total des flèches $v_{\alpha,\beta,\gamma}$, $w_{\alpha,\beta,\gamma}$ (pour $\alpha$, $\beta$ fixés) est majoré par $\pi_1 \cdot \pi_1 = \pi_1$. Donc le nombre total des flèches $u_{\alpha,\beta,\gamma}$, pour $\alpha$, $\beta$, $\gamma$ variables, est majoré par $\pi \cdot \pi_1 = \pi$ (si $\pi \geqslant \pi_1$). Donc $x$ est représenté par une catégorie $I$ ayant comme objets l'ensemble somme de $A$ et de l'ensemble des indices triples $\alpha,\beta,\gamma$, et pour chaque tel indice, exactement deux flèches $v_{\alpha,\beta,\gamma}$, $w_{\alpha,\beta,\gamma}$ non identiques de source cet objet, et allant respectivement dans $x_{\alpha}$, $x_{\beta}$. On a donc $\card \Fl I \leqslant \pi$, et on a un foncteur $I \longrightarrow C$, $\alpha \longmapsto x_{\alpha}$, $(\alpha,\beta,\gamma) \longmapsto x_{\alpha,\beta,\gamma}$, tel que $x \ = \ \varinjlim\limits_{i\in I} x_i$. Cela achève la démonstration de 1°).

\msp{page 132}
\textit{Prouvons 2°.} La première inclusion (4.5.13.1) est tautologique, la deuxième ne fait qu'exprimer $P(\pi) \Longrightarrow \mathrm{d}$. Le fait que $\pi^{\pi_0} = \pi$ implique l'égalité des membres extrêmes est un rappel de 4.5.6. Reste à prouver que $\Filt'^{\pi}$ est stable dans $\mathcal{M}$ par sommes de cardinal $\leqslant \pi$. Mais si $(x_{\alpha})_{\alpha\in A}$ est une famille d'objets de $\Filt'^{\pi}$, avec $\card A \leqslant \pi$, et si $x = \coprod_{\alpha\in A} x_{\alpha}$, prenant pour tout $x_{\alpha}$ un foncteur $\varphi_{\alpha} : I_{\alpha} \longrightarrow C$, avec $\card \Fl I_{\alpha} \leqslant \pi$, tel que $x_{\alpha}$ soit la $\varinjlim$ dans $\mathcal{M}$ de ce foncteur, on trouve que $x$ est $\varinjlim$ du foncteur $\varphi : I \longrightarrow C$ défini par les $\varphi_{\alpha}$, où $I = \coprod_{\alpha} I_{\alpha}$ (somme dans Cat). Il est évident que $\card \Fl(I) \leqslant \pi\cdot\pi = \pi$, ce qui achève de prouver 2°.

\textit{Prouvons 3°.} On a les implications tautologiques $1 \Longrightarrow 1'$, et $1' \Longrightarrow 1$, puisque $\Filt'^{\pi}$ est stable par sommes de cardinal $\leqslant \pi$. Donc $1 \Longleftrightarrow 1'$, et implique que pour tout $x$ dans $\mathcal{M}$,
\msp{page 133}
$I \ = \ \Filt'^{\pi}/x$ est stable par $\varinjlim$ de cardinal $\leqslant \pi$, donc est grand devant $\pi$. Comme $x \simeq \varinjlim\limits_{i\in I} x_i$, cela prouve que $1 \Longrightarrow 2$. L'implication $2 \Longrightarrow 3$ est standard, et résulte de $\Filt'^{\pi} \subset \mathcal{M}_{\pi}$, je vais expliciter dans un

\textbf{Lemme 4.5.14.} \textit{Soient $\mathcal{M}$ une catégorie stable par petites} $\varinjlim$ \textit{, $\pi$ un cardinal infini, $\mathcal{M}_0$ une sous-catégorie strictement pleine de $\mathcal{M}$. Considérons les conditions}
\begin{itemize}
\item[\textit{a)}] \textit{$\mathcal{M}_0$ est une sous-catégorie strictement génératrice, et stable dans $\mathcal{M}$ par} $\varinjlim$ \textit{de cardinal $\leqslant \pi$.}
\item[\textit{b)}] \textit{Pour tout objet de $\mathcal{M}$, on a un isomorphisme}
\[ x \ \simeq \ \varinjlim_{i\in I} x_i \ , \]
\textit{les $x_i$ dans $\mathcal{M}_0$, $I$ grande devant $\pi$.}
% foot: Version 17 décembre 2025

\printedpage{42}
% head: A. Grothendieck : Dérivateurs, Ch. XVIII, Édition des Universités de Montpellier II et Paris VII
\item[\textit{b$'$)}] \textit{Comme b), avec $I$ catégorie} ordonnée.
\item[\textit{c)}] \textit{$\mathcal{M}_{\pi} \subset \mathcal{M}_0$ (où $\mathcal{M}_{\pi}$ est la sous-catégorie strictement pleine de $\mathcal{M}$ formée des objets $\pi$-accessibles).}
\end{itemize}
\textit{On a donc les implications}
\[ a \ \Longrightarrow \ b \ \Longleftrightarrow \ b' \ \Longrightarrow \ c \ , \]
\textit{et si on suppose $\mathcal{M}_0 \subset \mathcal{M}_{\pi}$, et $\mathcal{M}_{\pi}$ sous-catégorie strictement génératrice de $\mathcal{M}$, alors les quatre conditions $a$, $b$, $b'$, $c$ sont équivalentes, et équivalent à $\mathcal{M}_0 = \mathcal{M}_{\pi}$} (\footnote{cf. démonstration de proposition 4.5.6, $4 \Longrightarrow 5$, page 117.}).

\msp{page 134}
Enfin (revenant à la démonstration de 4.5.13, 3°), l'implication $3) \Longrightarrow 1)$ provient du fait que $\mathcal{M}_{\pi}$ est stable par $\varinjlim$ de cardinal $\leqslant \pi$.

Reste à prouver que ces conditions, i.e. la stabilité de $\Filt'^{\pi}$ par conoyaux de double-flèches, est satisfaite dans les cas A) B) C), le cas A) étant d'ailleurs déjà connu. Le cas B) résulte immédiatement de la transitivité des familles épimorphiques strictes universelles, et c'est d'ailleurs aussi contenu explicitement (stabilité de $\Filt'^{\pi}$ par $\varinjlim$ de cardinal $\leqslant \pi$) dans SGA 4 I 9.13 (sinon dans l'énoncé, du moins c'est signalé dans la démonstration). En rédigeant SGA 4 I 9, je savais donc qu'on a $\mathcal{M}_{\pi} \subset \Filt'^{\pi}$ pour tout $\pi$, mais l'inclusion, plus évidente pourtant dans le cas d'espèce, $\Filt'^{\pi} \subset \mathcal{M}_{\pi}$ (quand on suppose les objets de $C$ $\pi_1$-accessibles) m'avait apparemment échappé\,?

Reste à traiter le cas C). $\Big($J'ai l'impression qu'il est un peu plus général que le cas B), i.e. que la condition B) sur $\mathcal{M}$ implique C), mais non inversement. Ainsi, la condition C) est vérifiée dès que $\mathcal{M}$ correspond à une espèce de structure algébrique définie en
\msp{page 135}
termes de $\varprojlim$ \textit{finies}, i.e. quand $\mathcal{M}$ est de la forme $\Ind(C)$, $C$ stable par $\varinjlim$ \textit{finies}, alors que je doute que la condition B) soit toujours satisfaite pour ces catégories. Par exemple dans la catégorie Cat (qui est un paratopos de cette espèce), les $\varinjlim$ filtrantes sont exactes à gauche, mais je ne crois pas que les familles épimorphiques stricts soient universelles. % [sic] épimorphiques stricts -> épimorphiques strictes, agreement with "familles"
Prenons p.ex. deux homomorphismes de groupes $G_0 \xrightarrow{u_0} G$, $G_1 \xrightarrow{u_1} G$ tels que $G_0 * G_1 \longrightarrow G$ soit surjectif, alors cette famille à deux éléments est épimorphique stricte dans Gr et dans Cat (\footnote{!, c'est \textit{faux} en général, mais vrai dans le cas particulier plus bas.}), mais n'est pas forcément épimorphique stricte universelle, même si p.ex. $G_0 = G_1 = H$, $G = H \times H$, et $u_0, u_1$ sont des injections $H \rightrightarrows H \times H$, données par $u_0(h) = (h,1)$, et $u_1(h) = (1,h)$. Si je prends l'application diagonale $H \longrightarrow H \times H$, les images inverses de $u_0$, $u_1$ sont le même objet de $\mathcal{M}/H$, savoir $e \longrightarrow H$, et cette famille n'est épimorphique stricte que si $H$ est le groupe unité.$\Big)$

Le cas C) va résulter du résultat plus précis :

\textbf{Corollaire 4.5.15.} \textit{Sous les conditions générales de 4.5.13, supposons que dans $\mathcal{M}$ les} $\varinjlim$
\msp{page 136}
filtrantes \textit{commutent au changement de base (p.ex. les} $\varinjlim$ \textit{filtrantes sont exactes). Alors pour tout $\pi \geqslant \pi_1$,} $\Filt'^{\pi}$ \textit{est formée des $x \in \Ob \mathcal{M}$ satisfaisant la condition suivante.}
% foot: Version 17 décembre 2025

\printedpage{43}
% head: A. Grothendieck : Dérivateurs, Ch. XVIII, Édition des Universités de Montpellier II et Paris VII
\begin{itemize}
\item[\textit{c$'$)}] \textit{Il existe une famille de flèches $(u_{\alpha})_{\alpha\in A}$, $u_{\alpha} : x_{\alpha} \longrightarrow x$, dans $\mathcal{M}$, qui est} épimorphique stricte universelle \textit{(i.e. couvrante pour la topologie canonique de $\mathcal{M}$), avec $x_{\alpha} \in \Ob C$ $\forall\ \alpha \in A$, et} $\card A \leqslant \pi$.
\end{itemize}
\textsc{Démonstration} : il faut prouver que cette condition équivaut à la condition c) de 4.5.13 1°). Comme on a $\mathrm{c}' \Longrightarrow \mathrm{c}$ trivialement, il reste à montrer que $\mathrm{c} \Longrightarrow \mathrm{c}'$. Mais on a vu dans 4.5.13 1°) que c implique a$''$, i.e. qu'on a un isomorphisme
\[ x \ = \ \varinjlim_{i\in I} x_i \ , \]
les $x_i$ dans $C$, et $I$ ensemble \textit{ordonné filtrant}, avec $\card I \leqslant \pi$. Il reste à voir que si $x$ dans $\mathcal{M}$ est une limite inductive \textit{filtrante} d'objets $x_i$, alors l'ensemble des flèches $x_i \longrightarrow x$ est épimorphique stricte \textit{universelle.} Or cela est immédiat, car pour $x' \longrightarrow x$, on a $x' = \varinjlim\limits_{i} x'_i$ (par hypothèse C), où $x'_i = x' \times_x x_i$, donc la famille $x'_i \longrightarrow x'$ est épimorphique stricte, q.e.d.

Prouvons maintenant que dans le cas C, la condition de 3°, 1 est satisfaite, i.e. que $\Filt'^{\pi}$ est stable par $\varinjlim$ de cardinal $\leqslant \pi$. Cela
\msp{page 137}
résulte aussitôt du

\textbf{Corollaire 4.5.16.} (de 4.5.15). \textit{Sous les conditions de 4.5.15 (i.e. la condition C) de 4.5.13, 3°), soit $x$ un objet de $\mathcal{M}$, et supposons qu'il existe une famille} épimorphique \underline{universelle} $(x_{\alpha} \longrightarrow x)_{\alpha\in A}$ \textit{de flèches de $\mathcal{M}$, avec $x_{\alpha} \in$} $\Ob \Filt'^{\pi}$ \textit{pour tout $\alpha \in A$, et} $\card A \leqslant \pi$. \textit{Alors $x \in$} $\Ob \Filt'^{\pi}$.

En effet, l'hypothèse sur $x_{\alpha}$ équivaut, en vertu de 4.5.15, à l'existence d'une famille épimorphique stricte universelle de flèches $(x_{\alpha\beta} \longrightarrow x_{\alpha})_{\beta\in B_{\alpha}}$, avec $\card B_{\alpha} \leqslant \pi$ et les $x_{\alpha\beta}$ dans $C$. Soit $B = \coprod_{\alpha\in A} B_{\alpha}$, on a alors $\card B \leqslant \pi^2 = \pi$, et d'autre part, la famille des flèches composées $x_{\alpha\beta} \longrightarrow x_{\alpha} \longrightarrow x$ indexée par $B$ est épimorphique stricte universelle (par transitivité de ce type de familles), donc $x \in \Ob \Filt'^{\pi}$, q.e.d.

Avec tout ça, on n'a toujours pas donné de caractérisation de $\mathcal{M}_{\pi}$ lui-même, qui est meilleure que $\Filt^{\pi}$ et $\Filt'^{\pi}$ du point de vue catégorique, car seuls les $\mathcal{M}_{\pi}$ définissent une filtration cardinale, et de plus ils ont les propriétés de stabilité désirables dans ce
\msp{page 138}
contexte, pour les $\varinjlim$ de cardinal $\leqslant \pi$. Le seul critère non tautologique que j'arrive à dégager dans le cas général (pour tout $\pi \geqslant \pi_0$, et sans hypothèse d'exactitude supplémentaire sur $\mathcal{M}$) est le suivant, beaucoup moins ``explicite'' ou ``calculatoire'' que les critères très concrets pour $\Filt^{\pi}$ et pour $\Filt'^{\pi}$ :

\textbf{Proposition 4.5.17.} \textit{Soit $\pi$ un cardinal $\geqslant \pi_0$. Alors $\mathcal{M}_{\pi}$ est égal à la plus petite sous-catégorie strictement pleine de $\mathcal{M}$, contenant $C$, et stable par} $\varinjlim$ \textit{de cardinal $\leqslant \pi$.} (Il est important de noter qu'il \textit{ne suffit pas} de dire : stable par $\varinjlim\limits_{I}$ , avec \textit{$I$ grande devant} $\pi_0$, et $\card I \leqslant \pi$. On trouve alors $\Filt^{\pi}$, qui risque d'être une catégorie strictement plus petite que $\mathcal{M}_{\pi}$. Ce n'est pas non plus la catégorie des $\varinjlim\limits_{I} x_i$, avec $x_i \in C$ $\forall\ i \in \Ob I$, et $\card \Fl I \leqslant \pi$ -- puisque cette catégorie-là est $\Filt'^{\pi}$.)

\textsc{Démonstration}. Soit $\mathcal{M}'$ la catégorie de l'énoncé, on a donc $\mathcal{M}' \subset \mathcal{M}_{\pi}$, puisque $\mathcal{M}_{\pi}$ contient $C$ et est stable par les $\varinjlim$ de cardinal $\leqslant \pi$. Inversement, le lemme 4.5.14, $\mathrm{a} \Longrightarrow \mathrm{c}$.
% foot: Version 17 décembre 2025


\endgroup
% END UNIT der-ch18/der-ch18_p041-045.tex
% BEGIN UNIT der-ch18/der-ch18_p046-050.tex
\begingroup\makeatletter%
\newcommand{\printedpage}[1]{\marginpar{\footnotesize p.~#1}}%
\newcommand{\msp}[1]{\marginpar{\footnotesize [#1]}}%
\newcommand{\cardred}{\mathop{\mathrm{card}\,\textup{réd}}\nolimits}%
\newcommand{\Ind}{\mathrm{Ind}}%
\newcommand{\circled}[1]{\tikz[baseline=(c.base)]\node[draw,circle,inner sep=1pt](c){#1};}%
\makeatother\let\printedpage\derpp\let\msp\dermsp%
\let\mspin\DERmspin%


\printedpage{44}
% head: A. Grothendieck : Dérivateurs, Ch. XVIII, Édition des Universités de Montpellier II et Paris VII
% foot: Version 17 décembre 2025

\msp{page 139}

\section*{4.6\quad Caractérisation des catégories paratopos (démonstration de 4.3.6)}

Il est temps de revenir à la démonstration du théorème 4.3.6 (p. 70-71), celle que j'avais donnée étant incomplète (et erronée – j'affirmais que $(\mathrm{Filt}^{\pi})_{\pi\geqslant\pi_0}$ est une filtration cardinale, ce qui est faux si on n'admet que $\pi^{\pi_0}=\pi$ pour tout $\pi\geqslant\pi_1$).

On avait vu l'équivalence des conditions a, a$'$, b, soit PTA. Je vais achever la démonstration par le schéma logique
\[
\text{PTA} \ \Longrightarrow\ \text{d)} \ \Longrightarrow\ \text{c)} \ \Longrightarrow\ \text{PTA}\ . \tag*{$(*)$}
\]

Prouvons PTA $\Longrightarrow$ d), i.e. que si $\mathcal{M}$ est un pseudo-topos accessible, alors existe un cardinal infini $\pi_0$, et une petite catégorie $C_0$ stable par $\varinjlim$ de cardinal $\leqslant\pi_0$, et une équivalence de catégories % [sic] alors existe -> alors il existe, missing subject "il"
\[
\mathcal{M} \xrightarrow{\ \sim\ } \Ind'_{\pi_0}(C_0)\ . \tag*{(4.6.1)}
\]

Plus précisément, on a ceci :

\textbf{Proposition 4.6.1.} \textit{Soit $\mathcal{M}$ un paratopos, i.e. une $\mathfrak{U}$-catégorie $\mathcal{M}$ stable par petites $\varinjlim$ , admettant une sous-catégorie pleine génératrice $S$ équivalent à une petite catégorie, formée d'objets accessibles de $\mathcal{M}$. Soit $\pi_0$ un cardinal tel que les objets de $S$ soient $\pi_0$-accessibles (i.e. $S\subset\mathcal{M}_{\pi_0}$), et soit $C$ la plus petite sous-catégorie strictement pleine de $\mathcal{M}$, contenant $S$ et stable par $\varinjlim$ de cardinal $\leqslant\pi_0$. Alors} % [sic] équivalent -> équivalente, agreement with the feminine "sous-catégorie"
\begin{itemize}[leftmargin=2.5em]
\item[\textit{1°)}] \textit{$C$ est équivalente à une petite catégorie.}

\msp{page 140}

\item[\textit{2°)}] \textit{$C=\mathcal{M}_{\pi_0}$, sous-catégorie strictement pleine de $\mathcal{M}$ formée des objets $\pi_0$-accessibles.}
\item[\textit{3°)}] \textit{Le foncteur canonique}
\[
\Ind'_{\pi_0}(C) \xrightarrow[\sim]{\texttt{[}\varphi\texttt{]}} \mathcal{M} % ?? uncertain: font of the brackets around \varphi above the arrow -- typewriter (editorial) brackets or ordinary math brackets
\]
\textit{(défini à isomorphisme unique près par la condition que son composé avec l'inclusion $C\hookrightarrow\Ind'_{\pi_0}(C)$ soit isomorphe à l'inclusion $C\longrightarrow\mathcal{M}$, et que $\varphi$ commute aux petites $\varinjlim$ , ou encore aux $\varinjlim$ filtrantes grandes devant $\pi_0$ – cf. 4.4.16, page 101) est une équivalence de catégories} (\footnote[54]{donc on a aussi $\Ind'_{\pi_0}(C_0)\xrightarrow{\ \sim\ }\mathcal{M}$, où $C_0$ est une \textit{petite} catégorie équivalente à $C$.}).
\end{itemize}

\textsc{Démonstration}. 1°) On va construire $C$ par induction transfinie, en posant
\[
\begin{aligned}
C_0 &\quad=\quad && \text{sous-catégorie strictement pleine de $\mathcal{M}$ engendrée par $S$.}\\
C_{\alpha+1} &\quad=\quad && \text{sous-catégorie strictement pleine de $\mathcal{M}$ engendrée par les $\varinjlim\limits_{I}$ dans}\\
& && \text{$\mathcal{M}$ d'objets de $C_\alpha$, avec $I$ catégorie telle que $\operatorname{card}\operatorname{Fl}I\leqslant\pi_0$.}\\
C_\alpha &\quad=\quad && \textstyle\bigcup_{\alpha'<\alpha}C_{\alpha'}\ \text{ si $\alpha$ est un ordinal limite.}
\end{aligned}
\]

Cela définit de façon unique les $C_\alpha$ (pour tout petit cardinal $\alpha$). Par la définition de $C$, on voit par induction transfinie que les $C_\alpha$ sont contenues dans $C$. Soit $\alpha_0$
\printedpage{45}
% head: A. Grothendieck : Dérivateurs, Ch. XVIII, Édition des Universités de Montpellier II et Paris VII
% foot: Version 17 décembre 2025
\msp{page 141}
% ?? uncertain: a hand-drawn mark (shaped roughly like a "2" or a bracket) stands in the left margin beside lines 3--5 of this page -- marginal annotation reproduced in the image, not typeset text
le premier ordinal tel que $\operatorname{card}\alpha_0>\pi_0$, donc l'ensemble $I_{\alpha_0}$ des cardinaux $\alpha<\alpha_0$ est grand devant $\pi_0$. Je dis qu'on a $C_{\alpha_0}=C$. Pour ceci, il suffit de voir que $C_{\alpha_0}$ est stable par $\varinjlim$ de cardinal $\leqslant\pi$. Mais si $I$ est une catégorie telle que $\operatorname{card}\operatorname{Fl}I\leqslant\pi_0$, et $\varphi : I\longrightarrow C_{\alpha_0}=\bigcup_{\alpha<\alpha_0}C_\alpha$ un foncteur, alors la définition de $\alpha_0$ implique que $\varphi$ se factorise en $\varphi_\alpha : I\longrightarrow C_\alpha$, avec $\alpha<\alpha_0$ (\footnote[55]{Un peu court. Il faut invoquer le fait que $I_{\alpha_0}$ est grande devant $\pi_0$.}). Donc $\varinjlim\varphi\in\operatorname{Ob}C_{\alpha+1}\subset\operatorname{Ob}C_{\alpha_0}$, ce qui montre la stabilité requise. % [sic] de cardinal \leqslant\pi -> de cardinal \leqslant\pi_0, the stability to be proved concerns cardinal \leqslant\pi_0

Pour voir que $C$ est équivalente à une petite catégorie, il suffit de le voir pour les $C_\alpha$, en procédant par induction transfinie. Cela résulte du lemme suivant, immédiat :

\textbf{Lemme 4.6.1.1.}
\begin{itemize}[leftmargin=2.5em]
\item[\textit{1°)}] \textit{La limite inductive $C$ d'un système inductive filtrant $(C_i)_{i\in I}$ de catégories $C_i$ équivalentes à des petites catégories est équivalente à une petite catégorie. Si $\cardred(C_i)$ $\leqslant\pi$ pour tout $i$, et $\operatorname{card}I\leqslant\pi$, alors $\cardred(C)\leqslant\pi$ ($\pi$ étant un} % [sic] système inductive -> système inductif, agreement with the masculine "système"

\msp{page 142}

\textit{cardinal infini). Ici, pour toute catégorie $\mathcal{A}$, je pose $\operatorname{card}\mathcal{A}=\operatorname{card}\operatorname{Fl}\mathcal{A}$, $\cardred(\mathcal{A})$ $=\operatorname{card}(\mathcal{A}_0)$, où $\mathcal{A}_0$ est une catégorie réduite équivalente à $\mathcal{A}$.}
\item[\textit{2°)}] \textit{Soit $C_0\subset\mathcal{M}$ une sous-catégorie pleine d'une catégorie $\mathcal{M}$, $\pi_0$ un cardinal infini, et soit $C'_0$ la sous-catégorie strictement pleine de $\mathcal{M}$ engendrée par les limites inductives dans $\mathcal{M}$ (représentables dans $\mathcal{M}$) de foncteurs $I\longrightarrow C_0$, avec $I$ catégorie telle que $\operatorname{card}I\leqslant\pi_0$. Si $C_0$ est équivalente à une petite catégorie, il en est de même de $C'_0$. Plus précisément, si $\cardred(C_0)\leqslant\pi_1$, $\pi_1$ un cardinal infini, alors $\cardred(C'_0)\leqslant$ $\pi_1^{\pi_0}$} (\footnote[56]{Cette majoration n'est par prouvée et probablement fausse si on ne fait pas des hypothèses un peu plus strictes.}). % [sic] n'est par prouvée -> n'est pas prouvée, misprint in footnote 56
\end{itemize}

\textbf{Corollaire 4.6.1.2.} \textit{$\mathcal{M}$ et $C_0$ étant comme dans 2°) ci-dessus, soit $C$ la plus petite sous-catégorie strictement pleine de $\mathcal{M}$ contenant $C_0$, et stable par $\varinjlim$ de cardinal $\leqslant\pi_0$ (celles qui sont représentables dans $\mathcal{M}$). Si $C_0$ est équivalente à une petite catégorie, il en est de même de $C$. Plus précisément, si $\cardred(C_0)\leqslant\pi_1$, où $\pi_1$ est un cardinal infini, alors}

\msp{page 143}

\textit{on a $\cardred(C)\leqslant\pi_1^{\pi_0}$} (\footnote[57]{majoration non prouvée, probablement fausse, cf. annotation ci-dessus.}).

Le 3°) se prouve par la démonstration transfinie précédente, on prouve par récurrence transfinie que
\[
\cardred C_\alpha \ \leqslant\ \pi_1^{\pi_0}
\]
pour tout $\alpha$. C'est OK pour $\alpha=0$ par hypothèse (puisque $\pi_1\leqslant\pi_1^{\pi_0}$). Si c'est vrai pour $C_\alpha$, alors 2°) assure que
\[
\cardred C_\alpha \ \leqslant\ (\pi_1^{\pi_0})^{\pi_0} \ =\ \pi_1^{\pi_0\cdot\pi_0} \ =\ \pi_1^{\pi_0}\ , % [sic] C_\alpha -> C_{\alpha+1}, this is the successor step of the induction
\]
enfin, si $\alpha$ est un ordinal limite $<\alpha_0$, on aura $C_\alpha=\bigcup_{\alpha'<\alpha}C_{\alpha'}$, donc
\[
\cardred C_\alpha \ \longrightarrow\ \pi_0\cdot\pi_1^{\pi_0} \ =\ \pi_1^{\pi_0} % [sic] \longrightarrow -> \leqslant, an inequality is intended
\]
(vu que $\pi_1^{\pi_0}\geqslant\pi_0$), en vertu du 1° ci-dessus, et compte tenu du fait que $\operatorname{card}I_\alpha\leqslant\pi_0$. Soit $\pi'_0$ le cardinal successeur de $\pi_0$, alors la même majoration nous donne
\[
\cardred C_{\alpha_0} \ \leqslant\ \pi'_0\pi_1^{\pi_0}\ .
\]
\printedpage{46}
% head: A. Grothendieck : Dérivateurs, Ch. XVIII, Édition des Universités de Montpellier II et Paris VII
% foot: Version 17 décembre 2025
Mais on a $\pi_1^{\pi_0}\geqslant2^{\pi_0}>\pi_0$, donc $\pi_1^{\pi_0}\geqslant\pi'_0$, donc $\pi_1^{\pi_0}\geqslant\pi'_0$, donc $\pi'_0\pi_1^{\pi_0}=\pi_1^{\pi_0}$, d'où
\[
\cardred C_{\alpha_0} \ \leqslant\ \pi_1^{\pi_0}\ ,
\]
ce qui est la majoration annoncée.

\msp{page 144}

Prouvons 4.6.1 2°, i.e. $C=\mathcal{M}_{\pi_0}$. On a $C\subset\mathcal{M}_{\pi_0}$, puisque par hypothèse $S\subset\mathcal{M}_{\pi_0}$, et on sait que $\mathcal{M}_{\pi_0}$ est stable par $\varinjlim$ de cardinal $\leqslant\pi_0$. D'autre part, on a $\mathcal{M}_{\pi_0}\subset C$, en vertu de 4.5.19 a $\Longrightarrow$ c appliqué à $C$, qui est strictement génératrice dans $\mathcal{M}$, $S$ l'étant, et $C\supset S$, et stable par $\varinjlim$ de cardinal $\leqslant\pi_0$ par construction.

\textit{Prouvons enfin 3°}. Cela équivaut à ceci :

\textbf{Lemme 4.6.1.3.} \textit{Soit $\mathcal{M}$ une $\mathfrak{U}$-catégorie stable par} petites \textit{$\varinjlim$ , $\pi_0$ un cardinal, $\mathcal{M}_0$ une sous-catégorie de $\mathcal{M}$} strictement génératrice, \textit{stable par $\varinjlim$ de cardinal $\leqslant\pi_0$, et formée d'objets $\pi_0$-accessibles. (En vertu de 4.6.1 2° déjà prouvée, appliqué à $S=\mathcal{M}_0$, si $\mathcal{M}_0$ est strictement pleine, on a donc $\mathcal{M}_0=\mathcal{M}_{\pi_0}$.) Alors le foncteur canonique}
\[
\Ind'_{\pi_0}(\mathcal{M}_0) \ \longrightarrow\ \mathcal{M}
\]
\textit{est une équivalence de catégories.}

Le fait que ce facteur est pleinement fidèle résulte aussitôt de l'hypothèse que les objets de $\mathcal{M}_0$ soient $\pi_0$-accessibles : % [sic] facteur -> foncteur, misprint
\[
\begin{aligned}
\operatorname{Hom}_{\mathcal{M}}(\varinjlim_{i\in I}\underbrace{X_i}_{\in\mathcal{M}_0},\ \varinjlim_{j\in J}\underbrace{Y_j}_{\in\mathcal{M}_0}\ ) &\ \simeq\ \varprojlim_{i\in I}\ \underbrace{\operatorname{Hom}_{\mathcal{M}}(X_i,\ \varinjlim_{j\in J}Y_j)}_{\varinjlim\limits_{j\in J}\operatorname{Hom}(X_i,Y_j)}\\
&\ \simeq\ \underline{\operatorname{Hom}}_{\Ind_{\pi_0}(\mathcal{M})}((X_i),(Y_j))\ .
\end{aligned}
\]

\msp{page 145}

Reste à voir que ce foncteur est essentiellement surjectif. Comme $\mathcal{M}_0$ est strictement génératrice, il s'ensuit que tout $x\in\mathcal{M}$ s'écrit
\[
x \ =\ \varinjlim_{i\in I}x_i\ ,
\]
où $I=\mathcal{M}_0/x$, et les $x_i\in\mathcal{M}_0$. L'hypothèse que $\mathcal{M}_0$ est stable par $\varinjlim$ de cardinal $\leqslant\pi_0$ implique la même propriété pour $I=\mathcal{M}_0/x$, donc $I$ est grande devant $\pi_0$. Donc $(x_i)_{i\in I}$ définit un objet de $\Ind_{\pi_0}\mathcal{M}_0$, dont l'image dans $\mathcal{M}$ est isomorphe à $x$, q.e.d. (Comparer SGA 4 I 9.18.)

\textbf{Corollaire 4.6.2.} \textit{Sous les conditions de 4.6.1, prenons pour $\pi_0$} le plus petit \textit{cardinal tel que les objets de $S$ soient $\pi_0$-accessibles. Pour tout $\pi\geqslant\pi_0$, soit $\mathcal{M}_\pi$ la sous-catégorie strictement pleine de $\mathcal{M}$ formée des objets $\pi$-accessibles, qui est aussi (par 4.6.1 2°) l'enveloppe de $C$ dans $\mathcal{M}$ stable par $\varinjlim$ de cardinal $\leqslant\pi$.}
\begin{itemize}[leftmargin=2.5em]
\item[\textit{1°)}] \textit{Soit $\pi_1=\cardred S$. On a alors}
\[
\cardred\mathcal{M}_\pi \ \leqslant\ \pi_1^{\pi} \qquad (=2^{\pi}\ \textit{si}\ \pi\geqslant\pi_0). \tag*{(4.6.2.1)}
\]
\printedpage{47}
% head: A. Grothendieck : Dérivateurs, Ch. XVIII, Édition des Universités de Montpellier II et Paris VII
% foot: Version 17 décembre 2025
\msp{page 146}
\item[\textit{2°)}] \textit{$(\mathcal{M}_\pi)_{\pi\geqslant\pi_1}$ est une filtration cardinale de $\mathcal{M}$, où $\pi_1=\sup(\pi_0,\operatorname{card}S)$.}
\end{itemize}

\textsc{Démonstration}. La majoration (4.6.2.1) est contenue dans 4.6.1.2, où $\pi_0$ est devenu $\pi$. Si $\pi\geqslant\pi_1$, on a
\[
\pi_1^{\pi} \ =\ 2^{\pi} \tag*{$(*)$}
\]
car on a $\pi_1\geqslant2$, d'où $\pi_1^{\pi}\geqslant2^{\pi}$, d'autre part, on a $\pi_1\leqslant2^{\pi_1}$, d'où $\pi_1^{\pi}\leqslant(2^{\pi_1})^{\pi}=2^{\pi_1\pi}=2^{\pi}$, puisque $\pi_1\pi=\pi$ pour $\pi\geqslant\pi_1$ ; d'où $(*)$.

L'assertion 2° est mise pour mémoire ; en vertu de $\mathcal{M}\simeq\Ind'_{\pi_0}\mathcal{M}_{\pi_0}$, on peut supposer que l'on a $\mathcal{M}=\Ind'_{\pi_0}(\mathcal{M}_{\pi_0})$, et on peut alors appliquer 4.5.9. On aurait d'ailleurs envie de dire que $(\mathcal{M}_\pi)_{\pi\geqslant\pi_0}$ est une filtration cardinale, sans se borner aux $\pi\geqslant\pi_1$. C'est vrai sauf pour la majoration de cardinaux dans la condition c) de la définition des filtrations cardinales (SGA 4 I 9.12) :
\[
\underbrace{x}_{\in\operatorname{Ob}\mathcal{M}_{\pi'}} \ =\ \varinjlim_{i\in I}\underbrace{x_i}_{\in\operatorname{Ob}\mathcal{M}_{\pi}}\ , \qquad\qquad \text{$I$ grande devant $\pi$ et $\operatorname{card}I\leqslant\pi'^{\pi}$}
\]
pour $\pi'\geqslant\pi$. Il faut supposer, pour cette majoration, que $\pi\geqslant\pi_1$.

\msp{page 147}

On peut à présent achever la démonstration de 4.3.6, en prouvant les deux dernières implications $(*)$ (p. 139). On a d $\Longrightarrow$ c, i.e. si $\mathcal{M}\simeq\Ind'_{\pi_0}(C)$, $C$ stable par limites inductives de cardinal $\leqslant\pi_0$, et $C$ équivalente à une petite catégorie, alors $\mathcal{M}$ est stable par petites $\varinjlim$ (4.4.13, p. 95-96), ses éléments sont accessibles, et elle admet une filtration cardinale (cf. 4.5.9). Mais il faut vérifier quand même que les éléments de $C$ sont $\pi_0$-accessibles. C'est contenu dans SGA 4 I 9.18, qui affirme que la pleine fidélité du foncteur $\Ind_{\pi_0}(C)\longrightarrow\mathcal{M}$ \textit{équivaut} au fait que les éléments de $C$ sont $\pi_0$-accessibles dans $\mathcal{M}$. La démonstration dans loc. cit. n'était pas claire, le “il suffit” a été vu bas de la page 144. Le “il faut” ne me paraît pas clair, avec la généralité donnée dans loc. cit. – et à présent je vois que c'est même faux, comme on voit en posant pour $C$ la catégorie unité $e$, donc $\Ind'_{\pi_0}C\simeq e$, et l'assertion affirmant que si $e$ est immergée pleinement dans $\mathcal{M}$ (donc par un objet \textit{rigide} de $\mathcal{M}$, i.e. un objet $a$ tel que $\operatorname{End}a=\{\mathrm{id}_a\}$), alors $a$ est $\aleph_0$-accessible. C'est faux dans la catégorie des anneaux, où $\mathbf{R}$ est un objet rigide non $\aleph_0$-accessible !

\msp{page 148}

Dans le cas qui nous occupe, $\mathcal{M}=\Ind'_{\pi_0}C$, il faut visiblement utiliser le calcul explicite des $\varinjlim$ grandes devant $\pi_0$ dans $\Ind'_{\pi_0}C\subset C^{\wedge}$, elles se calculent en effet “argument par argument”, cf. 4.4.15.

Cela achève la démonstration de c $\Longrightarrow$ d dans $(*)$, page 139. Enfin l'implication c $\Longrightarrow$ PTA est tautologique, l'existence d'une filtration cardinale impliquant celle d'une petite sous-catégorie strictement génératrice. Cela achève la démonstration de 4.3.6.

Je vais redonner une formulation légèrement plus précise de 4.3.6., en précisant les cardinaux, et en y rajoutant une condition supplémentaire, inspirée par SGA 4 I 9.11 (revue et corrigée, cf. une annotation à l'exemplaire \texttt{[de Grothendieck]} de SGA 4 I (\footnote[58]{N.~Éd.\ \ Cette annotation dit “cf. résultat plus simple ci-dessus, si $C$ est \textit{strictement} génératrice” et
% footnote 58 continues at the foot of printed page 48
le résultat en question, également en annotation marginale sur l'exemplaire de Grothendieck, est :\par
\textbf{Corollaire.} \textit{Soit $F$ admettant} \texttt{[une]} \textit{sous-catégorie génératrice par épimorphismes stricts, formée d'objets accessibles, alors} tout \textit{objet de $F$ est accessible (utiliser 7.2. (i) b)).}})).

\printedpage{48}
% head: A. Grothendieck : Dérivateurs, Ch. XVIII, Édition des Universités de Montpellier II et Paris VII
% foot: Version 17 décembre 2025

\textbf{Théorème 4.6.3.} \textit{Soit $\mathcal{M}$ une $\mathfrak{U}$-catégorie.}
\begin{itemize}[leftmargin=2.5em]
\item[\circled{1\textsuperscript{\textit{o}}}] \textit{Les conditions suivantes sont équivalentes pour un cardinal infini $\pi_0$ :}
\begin{itemize}[leftmargin=3em]
\item[$a_{\pi_0}$\textit{)}] \textit{$\mathcal{M}$ est une catégorie $\mathfrak{U}$-$\pi_0$-paratopos, i.e. $\mathcal{M}$ est stable par petites $\varinjlim$ , et contient une petite sous-catégorie pleine} \texttt{[strictement]} \textit{génératrice $S$, formée d'objets $\pi_0$-accessibles de $\mathcal{M}$.}

\msp{page 149}

\item[$b_{\pi_0}$\textit{)}] \textit{Il existe une petite catégorie $S$, et un foncteur pleinement fidèle et $\pi_0$-accessible $\mathcal{M}\xrightarrow{\ i\ }S^{\wedge}$, admettant un adjoint à gauche.}
\item[$c_{\pi_0}$\textit{)}] \textit{Il existe une petite catégorie $C$, stable par $\varinjlim$ de cardinal $\leqslant\pi_0$, telle que $\mathcal{M}\simeq\Ind'_{\pi_0}(C)$.}
\item[$c'_{\pi_0}$\textit{)}] \textit{Comme c), avec $C$ seulement équivalente à une petite catégorie.}
\end{itemize}
\item[\textit{2°)}] \textit{Ces conditions impliquent que $\mathcal{M}$ est stable par petites $\varinjlim$ et petites $\varprojlim$ , et la condition :}
\begin{itemize}[leftmargin=3em]
\item[$e_{\pi_0}$\textit{)}] \textit{Pour toute catégorie $I$ grande devant $\pi_0$, le foncteur}
\[
L \ =\ \varinjlim_{I}\ :\ \underline{\operatorname{Hom}}(I,\mathcal{M}) \ \longrightarrow\ \mathcal{M}
\]
\textit{est exact à gauche (donc exact).}
\end{itemize}
\textit{Cette condition implique notamment ceci :}
\begin{itemize}[leftmargin=3em]
\item[$\alpha(\pi_0)$] \textit{Sur la catégorie des double-flèches de $\mathcal{M}$, le foncteur} Ker \textit{(à valeurs dans $\mathcal{M}$) est $\pi_0$-accessible.}
\item[$\beta(\pi_0)$] \textit{Pour toute flèche $y\leqslant x$ dans $\mathcal{M}$, le foncteur changement de base $\mathcal{M}/x\longrightarrow\mathcal{M}/y$ est $\pi_0$-accessible.}
\end{itemize}
(\footnotemark[59]). % the text of footnote 59 is not printed on pages 46--50 of the PDF

\msp{page 150}

\item[\textit{3°)}] \textit{Les conditions suivantes sont équivalentes.}
\begin{itemize}[leftmargin=3em]
\item[\textit{a)}] \textit{$\mathcal{M}$ est une catégorie $\mathfrak{U}$-paratopos, i.e. est stable par petites $\varinjlim$ , et contient une petite sous-catégorie pleine $S$ génératrice par épimorphismes strictes, et formée d'objets accessibles.} % [sic] épimorphismes strictes -> épimorphismes stricts, agreement with the masculine "épimorphismes"
\item[\textit{a$'$)}] \textit{Comme a), mais en exigeant que $\mathcal{M}$ soit faiblement accessible, i.e. que} tout \textit{objet de $\mathcal{M}$ soit accessible.}
\item[\textit{b)}] \textit{Il existe une petite catégorie $S$ et un foncteur pleinement fidèle et accessible $\mathcal{M}\hookrightarrow S^{\wedge}$ admettant un adjoint à gauche.}
\item[\circled{c}] \textit{Il existe un cardinal infini $\pi_0$, une petite catégorie $S$, et une équivalence $\mathcal{M}\simeq$ $\Ind'_{\pi_0}(S)$.}
\item[\textit{c$'$)}] \textit{Comme c, en supposant seulement $S$ équivalente à une petite catégorie.}
\item[\textit{d)}] \textit{$\mathcal{M}$ est stable par petites $\varinjlim$ , est faiblement accessible, i.e. tous ses objets sont accessibles, et admet une filtration cardinale.}
\end{itemize}
\end{itemize}


\endgroup
% END UNIT der-ch18/der-ch18_p046-050.tex
% BEGIN UNIT der-ch18/der-ch18_p051-055.tex
\begingroup\makeatletter%
\newcommand{\printedpage}[1]{\marginpar{\footnotesize p.~#1}}%
\newcommand{\msp}[1]{\marginpar{\footnotesize [#1]}}%
\newcommand{\circled}[1]{\tikz[baseline=(c.base)]\node[draw,circle,inner sep=1pt](c){#1};}%
\newcommand{\Mc}{\mathcal{M}}%
\newcommand{\U}{\mathfrak{U}}%
\makeatother\let\printedpage\derpp\let\msp\dermsp%
\let\mspin\DERmspin%


\printedpage{49}
% head: A. Grothendieck : Dérivateurs, Ch. XVIII, Édition des Universités de Montpellier II et Paris VII
% foot: Version 17 décembre 2025
% The mark of footnote 59 is on the preceding page of the edition; its text is printed at the foot of this page.
\footnotetext[59]{\textbf{Question.} Condition $e_{\pi_0}$ (voire, le couple $\alpha(\pi_0) + \beta(\pi_0)$ plus l'existence d'une petite sous-catégorie strictement génératrice (et la stabilité de $\Mc$ par petites $\varinjlim$ ), implique-t-il que $\Mc$ est un $\pi_0$-paratopos ? Je présume que non.)}

\begin{itemize}[leftmargin=4em]
\item[\circled{e}] (\footnote[60]{C'est là la caractérisation intrinsèque la plus jolie des paratopos par la modalité “économie”, l'autre, l'énoncé structurel, est \circled{c} .}). \textit{$\Mc$ est stable par petites $\varinjlim$ , contient une petite sous-catégorie pleine strictement génératrice, et il existe un cardinal infini $\pi_0$ tel que la condition $\mathrm{e}_{\pi_0}$ de 2°) soit satisfaite.}
\item[\circled{e$'$}] \textit{Comme e, mais en exigeant, au lieu de $\mathrm{e}_{\pi_0}$, la condition plus faible conjonction de $\alpha(\pi_0)$ et $\beta(\pi_0)$ dans 2°.}
\end{itemize}

\msp{page 151}

\textsc{Démonstration}. La partie 1° a déjà été démontrée. La première assertions dans 2°), savoir la stabilité de $\Mc$ par petites $\varprojlim$ , est bien connue (résulte de stabilité par petites $\varinjlim$ , et de l'existence d'une petite famille génératrice, cf. SGA 4 I 8.12.7, version duale de l'annotation marginale). La validité de la condition $\mathrm{e}_{\pi_0}$ résulte du fait qu'elle est valable dans $C^{\wedge}$ (où $C$ est la catégorie mentionnée dans 1°, $\mathrm{c}_{\pi_0}$), et du fait que les $\varprojlim$ dans $\Mc$, et les $\varinjlim$ grandes devant $\pi_0$ dans $\Mc$, se calculent dans $C^{\wedge}$. On trouve d'ailleurs de même : % [sic] assertions -> assertion, singular subject "La première"

\textbf{Corollaire 4.6.4.} \textit{Si $\Mc$ est une catégorie $\U$-$\pi_0$-paratopos, alors pour tout cardinal $\pi \geqslant \pi_0$, et tout catégorie $J$ telle que} $\mathrm{card}\, J \leqslant \pi$\textit{, le foncteur} % [sic] tout catégorie -> toute catégorie, gender agreement
\[
\underline{\mathrm{Hom}}(J,\Mc) \xrightarrow{\ \varprojlim\limits_{J}\ } \Mc
\]
\textit{est $\pi$-accessible, ou ce qui revient au même, pour toute petite catégorie $I$ grande devant $\pi$, le foncteur}
\[
\underline{\mathrm{Hom}}(I,M) \xrightarrow{\ \varinjlim\limits_{I}\ } \Mc % [sic] M -> \mathcal{M}, the first display has \mathcal{M}
\]
\textit{commute aux $\varprojlim$ de cardinal $\leqslant \pi$.}

\msp{page 152}

D'autre part, le fait que $\mathrm{e}_{\pi_0}$ implique $\alpha(\pi_0)$ et $\beta(\pi_0)$ est immédiat.

La partie 3°), en ce qui concerne l'équivalence des conditions a) à d), est contenue dans 1°, et l'équivalence avec d) a été vue également (la seule chose non tautologique étant que la condition c), savoir $\Mc \simeq \mathrm{Ind}_{\pi_0}(C)$ avec $C$ petite et stable par $\varinjlim$ de cardinal $\leqslant \pi_0$, implique l'existence d'une filtration cardinale, ce qui a été vu dans 4.5.9, et réaffirmé dans 4.6.2 2° ci-dessus). D'autre part, que ces conditions a) à d) impliquent \circled{e} , et que \circled{e} implique e$'$, a été vu dans 2°. Donc il reste à voir que les conditions de e$'$ impliquent a), et pour ceci, il reste à établir que les éléments de $\Mc$ sont accessibles. Mais c'est ce qui est établi dans SGA 4 I 9.11, à cela près que l'énoncé de la condition c), deuxième cas, est canulé, et devrait être remplacé par la condition $\beta(\pi_0)$ (la condition $\alpha(\pi_0)$ figurant dans la condition a) de loc. cit.).

\msp{page 153}

\textbf{Remarque 4.6.5.} Les développements des sections 4.4, 4.5, 4.6 précisent et affirment considérablement, à bien des égards, la théorie des filtrations cardinales de SGA 4 I 9. Il en est ainsi, tout au moins, dans le cas des $\U$-catégories $\Mc$ stables par petites $\varinjlim$ . Dans loc. cit. on travaille avec des $\U$-catégories légèrement plus générales, celles satisfaisant
\printedpage{50}
% head: A. Grothendieck : Dérivateurs, Ch. XVIII, Édition des Universités de Montpellier II et Paris VII
% foot: Version 17 décembre 2025
la \textit{condition} $L$ : il existe un cardinal infini $\pi_0$ tel que $\Mc$ soit stable par petites $\varinjlim$ \textit{grandes devant} $\pi_0$. Je présume cependant qu'on doit pouvoir améliorer de façon similaire la théorie des filtrations cardinales \textit{accessibles}, en s'inspirant du corollaire 9.18 de loc. cit., établissant une équivalence de catégories
\[
\Mc \;\simeq\; \mathrm{Ind}_{\pi_0}(C)\;,
\]
pour un cardinal $\pi_0$ et une petite catégorie $C$ convenable. Ici, on ne peut plus supposer $C$ stable par tel ou tel type de $\varinjlim$ . Néanmoins, 4.4.15 nous dit que sans autre hypothèse sur $C$, $\Mc \overset{\text{déf}}{=} \mathrm{Ind}_{\pi_0}(C)$ satisfait bien à $\mathrm{L}_{\pi_0}$, i.e. est stable (et même stable dans $C^{\wedge}$)

\msp{page 154}

par petites $\varinjlim$ grandes devant $\pi_0$.

Ceci vu, j'ai envie d'appeler \textit{accessible} une catégorie $\Mc$ qui admet une filtration cardinal, et qui est de plus faiblement accessible \texttt{[i.e.]} dont tout objet est accessible. Et de hasarder une % [sic] filtration cardinal -> filtration cardinale, gender agreement

\textbf{Conjecture 4.6.6} (\footnote[61]{Tout ceci sera considérablement précisé dans la théorie qui sera développé dans les sections 4.7, 4.16.}). \textit{Soit $\Mc$ une $\U$-catégorie. Conditions équivalentes :} % [sic] (in footnote 61) développé -> développée, agreement with "la théorie"
\begin{itemize}[leftmargin=4em]
\item[\textit{a)}] \textit{$\Mc$ satisfait la condition $L$, et admet une petite catégorie strictement génératrice formée d'objets accessibles.}
% printed: a vertical double arrow $\Uparrow$ at the left, between a) and b)
\item[\llap{$\Uparrow$\ }\textit{b)}] \textit{Il existe une petite catégorie $C$, et un cardinal infini $\pi_0$, tel que $\Mc$ soit équivalente à la catégorie $\mathrm{Ind}'_{\pi_0}(C)$.}
% printed: a vertical double arrow $\Uparrow$ at the left, between b) and c)
\item[\llap{$\Uparrow$\ }\textit{c)}] \textit{$\Mc$ est “accessible”, i.e. est faiblement accessible (tout objet accessible) et admet une filtration cardinale.}
\end{itemize}
On a b $\Longrightarrow$ a $\Longleftarrow$ c de façon tautologique, et c $\Longrightarrow$ b a été vu dans loc. cit. 9.18, donc on a la suite
\[
\mathrm{c} \;\Longrightarrow\; \mathrm{b} \;\Longrightarrow\; \mathrm{a}\;,
\]
donc il resterait à établir a $\Longrightarrow$ c.

\msp{page 155}

Le travail accompli (et la sueur versée) dans 4.5 nous préservera d'ailleurs, le moment venu, d'erreurs fort naturelles dans la description d'une filtration cardinale, à partir d'une sous-catégorie $S$ strictement génératrice : ce n'est ni $\mathrm{Filt}^{\pi}$, ni $\mathrm{Filt}'^{\pi}$, mais quelque chose de plus subtil -- peut-être faut-il prendre encore la catégorie $\Mc_{\pi}$, ou alors saturer $S$ par les limites inductives de cardinal $\leqslant \pi$ qui veulent bien être représentables dans $\Mc$ : on trouve alors une $\mathrm{Filt}''^{\pi} \subset \Mc_{\pi}$, et si ces $\mathrm{Filt}''^{\pi}$ donnent bien une filtration cardinale, on doit avoir forcément $\Mc_{\pi} \subset \mathrm{Filt}''^{\pi}$, donc finalement $\mathrm{Filt}''^{\pi} = \Mc_{\pi}$. Donc je m'attends à ce qu'en tout état de cause, la filtration cardinale (si elle existe) se fasse par les $\Mc_{\pi}$. La condition a) de loc. cit. 9.12 ($\Mc_{\pi}$ équivalente à une petite catégorie), et la majoration de cardinaux dans loc. cit. c), ne doit pas faire de difficultés, et la condition b) est OK. La difficulté

\msp{page 156}

provient du fait que $\Mc_{\pi}$ n'étant pas même stable par sommes et conoyaux de double-flèches (quand ceux-ci ne sont pas représentés dans $\Mc$), la catégorie $\Mc_{\pi}/x$ (pour $x$ dans $\Mc$) n'a aucune raison d'être filtrante, pour autant que je vois, ni a fortiori
\printedpage{51}
% head: A. Grothendieck : Dérivateurs, Ch. XVIII, Édition des Universités de Montpellier II et Paris VII
% foot: Version 17 décembre 2025
d'être grande devant $\pi$. Donc il n'est pas trop clair comment dès lors représenter $x$ comme $\varinjlim\limits_{I} x_i$, les $x_i$ dans $\Mc_{\pi}$ et $I$ grande devant $\pi$.

Il faut bien dire que dans SGA 4 I 9, je ne donne aucun \textit{exemple} d'une filtration cardinale sur une catégorie où les petites $\varinjlim$ quelconques ne seraient représentables (cf. loc. cit. 9.13). Peut-être après tout est-ce contre-nature de vouloir elles au delà ? ? Il faut peut-être attendre d'y être forcé, pour essayer de développer une théorie en dehors du cadre des paratopos.

\msp{page 157}

\textbf{4.6.7.} Nous sommes à présent en mesure de reformuler les kyrielles de conditions un peu techniques a) à f) du théorème 2.6 (p. 25 dans XVIII) concernant le caractère “spécial à gauche” de l'ensemble $\mathrm{Mon}(\Mc) \subset \mathrm{Fl}(\Mc)$ :

\textbf{Proposition 4.6.7} (\footnote[62]{Peut-être on pourrait appeler ces catégories des catégories “\textit{paratopos parabéliens}” ? ?}). \textit{Les conditions préliminaires sur la $\U$-catégorie $\Mc$, assurant la validité du théorème XVIII 2.6, sous une forme très légèrement modifié (savoir en modifiant dans 2.6 d) la condition de stabilité de} $\mathrm{Filt}^{\pi}$ \textit{par quotients, en une propriété d'exactitude d) ci-dessous), équivaut à l'ensemble de six conditions suivantes a) à f).} % [sic] modifié -> modifiée, agreement with "forme"; équivaut -> équivalent, plural subject "Les conditions"
\[
\left\{
\begin{array}{ll}
\textit{a)} & \Mc\ \textit{est stable par petites}\ \varinjlim\ . \\[1.5ex]
\textit{b)} & \Mc\ \textit{admet une petite sous-catégorie strictement génératrice.} \\[1.5ex]
\textit{c)} & \textit{Les limites inductives filtrantes dans}\ \Mc\ \textit{sont exactes.}
\end{array}
\right.
\]
\textit{Ces conditions impliquent que $\Mc$ est une} catégorie $\U$-paratopos \textit{(cf. 4.6.3, 3°), et en particulier que les objets de $\Mc$ sont} accessibles, \textit{et que $\Mc$ admet une filtration cardinale $(\Mc_{\pi})_{\pi \geqslant \pi_0}$ telle que pour tout $\pi \geqslant \pi_0$, $\Mc_{\pi}$ soit stable par $\varinjlim$ de cardinal $\leqslant \pi$ (en fait, on aura $\Mc_{\pi} =$ sous-catégorie pleine des objets $\pi$-accessibles de $\Mc$). Aussi, elles impliquent la stabilité de $\Mc$ par petites $\varprojlim$ . D'autre part, les trois conditions suivantes sont des conditions de stabilité pour les monomorphismes.}
\begin{itemize}[leftmargin=4em]
\item[\textit{d)}] \textit{Soit $f : X \longrightarrow Y$ une flèche dans $\Mc$ et $R = X \times_Y X \rightrightarrows X$, alors $X/R \longrightarrow Y$ est un monomorphisme -- en d'autres termes, $f$ se factorise en un composé $ip$, avec $i$ monomorphe et $p$ un épimorphisme effectif.}
\end{itemize}

\msp{page 158}

\begin{itemize}[leftmargin=4em]
\item[\textit{e)}] \textit{Si $A_1$, $A_2$ sont deux sous-objets d'un objet $B$ de $\Mc$, et $A_0 = A_1 \cap A_2 \overset{\text{déf}}{=} A_1 \times_B A_2$, alors $A_1 \amalg_{A_0} A_2 \longrightarrow B$ est un monomorphisme. (Noter la proche parenté des conditions d) et e).)}
\item[\textit{f)}] \textit{L'ensemble} $\mathrm{Mon}(\Mc)$ \textit{des monomorphismes dans $\Mc$ est stable par cochangement de base.}
\end{itemize}
En fait, il n'est pas clair que ces conditions, et notamment d), suffisent à entraîner la stabilité des $\Mc_{\pi}$ par quotients (seulement par quotients effectifs), comme postulé dans 2.6 d). Mais dans l'énoncé et la démonstration donnée dans XV cor. 7.4 (p. 98-99), on se contentait d'une forme un peu plus faible de cette propriété de stabilité, en se bornant aux quotients $Z'$ de $Z \in \mathrm{Ob}\,\Mc_{\pi}$ \textit{tels} qu'il n'existe pas de sous-objet $Z'_0$ de $Z'$, distinct de $Z'$, et tel que $Z \longrightarrow Z'$ se factorise par $Z'_0$ ; on exigeait donc que $Z' \in \mathrm{Ob}\,\Mc_{\pi}$. Mais si la condition d) est satisfaite, il s'ensuit aussitôt qu'un quotient $Z'$ de $Z$ ayant cette propriété supplémentaire est nécessairement effectif, d'où résulte bien que $Z' \in \mathrm{Ob}\,\Mc_{\pi}$.

\printedpage{52}
% head: A. Grothendieck : Dérivateurs, Ch. XVIII, Édition des Universités de Montpellier II et Paris VII
% foot: Version 17 décembre 2025
\msp{page 159}

\textbf{Corollaire 4.6.8.} \textit{Les conditions d), e), f) de 4.6.7 sont satisfaites si $\Mc$ est une catégorie abélienne. Donc l'ensemble des conditions a) à f) est satisfait dans chacun des deux cas suivants :}
\begin{itemize}[leftmargin=4em]
\item[\textit{1)}] \textit{$\Mc$ est une catégorie abélienne où les limites inductives filtrantes sont représentables et exactes, et admettant une petite catégorie strictement génératrice} (\footnote[63]{C'est le cas étudié dans “Tohoku”.}).
\item[\textit{2)}] \textit{$\Mc$ est un topos.}
\end{itemize}
\textbf{Remarque 4.6.9.} Je ne suis par sûr de connaître d'autres exemples où les conditions de 4.6.7 sont vérifiées, que ceux signalés dans 4.6.8 ci-dessus (\footnote[64]{ou qui s'en déduisent trivialement (produits etc.)}). On notera que le cas 1°) qui y est mentionné est celui de Tohoku, où j'établis l'existence de “suffisamment d'objets injectifs” dans une telle catégorie abélienne. Dans Tohoku, je me borne à supposer l'existence d'une petite sous-catégorie, \textit{génératrice pour les monomorphismes}, condition en apparence plus faible que celle énoncée dans 4.6.8 1°). Mais dans une catégorie abélienne, les deux notions (et leurs diverses variantes dans SGA 4 I 7.1) sont équivalentes en vertu de SGA 4 I 7.3 (i). % [sic] par sûr -> pas sûr, negation "ne ... pas"

\msp{page 160}

Par contre, il semble que les conditions e) f) ne sont pratiquement jamais vérifiées pour les catégories définies par les espèces de structure algébrique “définissables par produits finis” sans plus, tel les groupes, les anneaux (de toutes sortes), etc. Mais la catégorie des groupes \textit{commutatifs}, étant abélienne, est une exception. (On peut d'ailleurs l'interpréter comme $\Mc \boxtimes \Mc$, où $\Mc = (\mathrm{Gr})$, et on peut se demander dès lors s'il y a quelque généralisations aux catégories paratopos de cette forme, $\Mc$ étant une catégorie paratopos \ldots) % [sic] tel les groupes -> tels les groupes, agreement; quelque généralisations -> quelques généralisations, plural

\msp{page 161}

\section*{4.7\quad $\mathrm{Ind}_{\pi}C$ ($C$ quelconque) et filtrations cardinales généralisées.}

Je sens le besoin de me lancer dans des développements à ce sujet, pour parvenir à une meilleure compréhension de la notion de \textit{partie accessible} d'un paratopos, et du théorème de stabilité pour les paratopos par certaines opérations-clef sur ceux-ci.

Dans la suite, $C$ désigne une catégorie équivalente à une petite catégorie, $C_0$ une sous-catégorie réduite exhaustive, $\pi_0$ un cardinal infini. Nous posons
\[
\Mc \;=\; \mathrm{Ind}_{\pi_0}C \;\subset\; C^{\wedge} \qquad \text{(\footnotemark[65])}\;. \tag*{(4.7.1)}
\]
\footnotetext[65]{\textbf{N.B.} Je désigne maintenant par $\mathrm{Ind}C$, $\mathrm{Ind}_{\pi}C$ ce qui précédemment était noté $\mathrm{Ind}'C$, $\mathrm{Ind}'_{\pi}C$.}
On se propose d'introduire sur $\Mc$ diverses filtrations $\pm$ “cardinales”, par des propriétés correspondantes sur un objet $F$ de $\Mc \subset C^{\wedge}$, que je vais à présent introduire, en m'inspirant de 4.5.2 (p. 107) et 4.5.6 (p. 116).

\textbf{Proposition 4.7.1.} \textit{Posons} $\pi_1 = \sup(\pi_0, \mathrm{card}\, C_0)$. \textit{Soit $\pi \geqslant \pi_0$ un cardinal, et considérons les conditions suivantes (dépendant de $\pi$) sur un objet $F$ de $\Mc$ :}
\begin{itemize}[leftmargin=4em]
\item[\textit{a)}] $\mathrm{card}\, C_0/F \leqslant \pi$ \textit{(i.e.} $\mathrm{card}\,\mathrm{Fl}\, C_0/F \leqslant \pi$\textit{).}
\item[\textit{a$'$)}] $\mathrm{card}\,\mathrm{Ob}\, C_0/F \leqslant \pi$.
\item[\textit{a$''$)}] $\mathrm{card}\, F(x) \leqslant \pi$ \textit{pour tout $x \in$} $\mathrm{Ob}\, C$ \textit{(}\textbf{N.B.} \textit{il suffit de l'exiger pour $x$ dans $C_0$).}
\end{itemize}

\printedpage{53}
% head: A. Grothendieck : Dérivateurs, Ch. XVIII, Édition des Universités de Montpellier II et Paris VII
% foot: Version 17 décembre 2025
\msp{page 162}

\begin{itemize}[leftmargin=4em]
\item[\textit{b)}] \textit{On a $F = \varinjlim\limits_{i \in I} x_i$, avec les $x_i \in$} $\mathrm{Ob}\, C$\textit{, $I$ une catégorie (pas nécessairement ordonnée) grande devant $\pi_0$, et telle que} $\mathrm{card}\, I \leqslant \pi$. \textit{(}\textbf{N.B.} \textit{en vertu de 4.4.2 (p.79), pour $I$ donnée, ceci équivaut à l'existence d'un foncteur cofinal $I \longrightarrow C/F$, ou encore $I \longrightarrow C_0/F$.)} (\footnote[66]{\textbf{N.B.} On désignera par $\mathrm{Filt}^{\pi} \subset \Mc$ la sous-catégorie strictement pleine de $\Mc$ formée des $F$ satisfaisant cette condition.}).
\item[\textit{b$'$)}] \textit{On a $F \in$} $\mathrm{Ob}\,\mathrm{Filt}'^{\pi}$\textit{, où} $\mathrm{Filt}'^{\pi}$ \textit{est la plus petite sous-catégorie strictement pleine de $\Mc$ contenant $C$, et stable par $\varinjlim$ grandes devant $\pi_0$ et de cardinal $\leqslant \pi$.}
\item[\textit{c)}] \textit{On a $F \in$} $\mathrm{Ob}\,\Mc_{\pi}$\textit{, i.e. $F$ est un objet $\pi$-accessible de $\Mc$.}
\item[\textit{c$'$)}] \textit{On a $F \in$} $\mathrm{Ob}\, C^{\wedge}_{\pi}$\textit{, i.e. $F$ est un objet $\pi$-accessible de $C^{\wedge}$.}
\end{itemize}
\textit{Ceci posé, on a le diagramme d'implications suivant.}
\[
\begin{tikzcd}[column sep=large, row sep=large]
\mathrm{a} \arrow[r, Rightarrow] \arrow[dr, Rightarrow] & \mathrm{a}' \arrow[r, Rightarrow] & \mathrm{a}'' \arrow[dr, Rightarrow] & & \\
& \mathrm{b} \arrow[r, Rightarrow] & \mathrm{b}' \arrow[r, Rightarrow] & \mathrm{c}' \arrow[dr, Rightarrow] \arrow[ulll, dotted, bend right=50, "\text{si } \pi \geqslant \pi_1"'] & \\
& & & & \mathrm{c}\;, \arrow[ull, dotted, bend left=30, "\text{si } \pi = \pi^{\pi_0}"']
\end{tikzcd}
\]
\textit{où les deux implications en pointillés, c$'$ $\Longrightarrow$ a et c $\Longrightarrow$ b, sont conditionnelles soumises à la condition $\pi \geqslant \pi_1$, resp. $\pi = \pi^{\pi_0}$.}

\msp{page 163}

\textbf{Corollaire 4.7.2.}
\begin{itemize}[leftmargin=4em]
\item[\textit{1°)}] \textit{Si $\pi \geqslant \pi_1$, les six conditions a, a$'$, a$''$, b, b$'$, c$'$ sont équivalentes, et la condition en question (qu'on peut écrire $F \in$} $\mathrm{Ob}\,\mathrm{Filt}^{\pi}$ \textit{où $F \in$} $\mathrm{Ob}\,\mathrm{Filt}'^{\pi}$ \textit{indifféremment) implique c :} % [sic] où -> ou, alternative "ou ... indifféremment"
\[
\mathrm{Filt}^{\pi} \;=\; \mathrm{Filt}'^{\pi} \;\subset\; \Mc_{\pi} \qquad \textit{si } \pi \geqslant \pi_1 \tag*{(4.7.2.1)}
\]
\textit{(alors qu'en général, pour $\pi \geqslant \pi_0$, on a}
\[
\mathrm{Filt}^{\pi} \;\subset\; \mathrm{Filt}'^{\pi}\ (\,\subset C^{\wedge}_{\pi} \cap \Mc\,) \;\subset\; \Mc_{\pi} \qquad \textit{si } \pi \geqslant \pi_0\textit{).} \tag*{(4.7.2.2)}
\]
\item[\textit{2°)}] \textit{Quand $\pi = \pi^{\pi_0}$, p.ex. quand $\pi$ est de la forme $2^{c}$, avec $c \geqslant \pi_0$, alors les conditions b, b$'$, c, c$'$ sont équivalentes, i.e. on a}
\[
\mathrm{Filt}^{\pi} \;=\; \mathrm{Filt}'^{\pi} \;=\; (C^{\wedge})_{\pi} \cap \Mc \;=\; \Mc_{\pi}\;. \tag*{(4.7.3.2)} % [sic] (4.7.3.2) -> (4.7.2.3), the neighbouring labels are 4.7.2.1, 4.7.2.2, 4.7.2.4
\]
\item[\textit{3°)}] \textit{Quand $\pi \geqslant \pi_1$} et \textit{$\pi = \pi^{\pi_0}$, alors les sept conditions a), a$'$), a$''$), b) , b$'$), c) , c$'$) sont équivalentes, donc en particulier}
\[
\mathrm{Filt}^{\pi} \;=\; \mathrm{Filt}'^{\pi} \;=\; C^{\wedge}_{\pi} \cap \Mc \;=\; \Mc_{\pi} \qquad (\pi \geqslant \pi_1 \textit{ et } \pi = \pi^{\pi_0}\textit{).} \tag*{(4.7.2.4)}
\]
\end{itemize}


\endgroup
% END UNIT der-ch18/der-ch18_p051-055.tex
% BEGIN UNIT der-ch18/der-ch18_p056-060.tex
\begingroup\makeatletter%
\newcommand{\printedpage}[1]{\marginpar{\footnotesize p.~#1}}%
\newcommand{\msp}[1]{\marginpar{\footnotesize [#1]}}%
\newcommand{\Filt}{\mathrm{Filt}}%
\newcommand{\card}{\operatorname{card}}%
\newcommand{\Ob}{\operatorname{Ob}}%
\newcommand{\Fl}{\operatorname{Fl}}%
\newcommand{\Ind}{\mathrm{Ind}}%
\newcommand{\cardred}{\text{card réd}\,}%
\makeatother\let\printedpage\derpp\let\msp\dermsp%
\let\mspin\DERmspin%


\printedpage{54}
% head: A. Grothendieck : Dérivateurs, Ch. XVIII, Édition des Universités de Montpellier II et Paris VII

\textsc{Démonstration} de 4.7.1. Les implications a $\Longrightarrow$ a$'$ $\Longrightarrow$ a$''$ sont tautologiques, ainsi que b $\Longrightarrow$ b$'$. L’implication b$'$ $\Longrightarrow$ c$'$ provient de la stabilité de $C^{\wedge}_{\pi}$ par $\varinjlim$ de cardinal $\leqslant \pi$, et du fait que les $\varinjlim$ du type envisagé (en particulier, grandes devant $\pi_0$) \textit{dans} $\mathcal{M}$ sont les mêmes que dans $C^{\wedge}$, i.e. se calculent “argument par argument” (cf. 4.4.15, p. 99). Ceci, plus le fait que l’inclusion $\mathcal{M} \hookrightarrow C^{\wedge}$ est pleinement fidèle, implique aussi c$'$ $\Longrightarrow$ c.

\msp{page 164}
L’implication a$''$ $\Longrightarrow$ c$'$ se voit comme dans 4.5.6 (p. 119) en appliquant le lemme 4.5.8 1° et 2°, au cas de $C^{\wedge} = \underline{\mathrm{Hom}}(C^{\circ}, \mathcal{E})$, où $\mathcal{E} =$ (Ens). L’implication conditionnelle inverse c$'$ $\Longrightarrow$ a$''$, pour $\pi \geqslant \pi_1$, se voit comme dans loc. cit., en appliquant ce lemme 4.5.8. Il reste donc à établir encore les trois implications
\[
\left\{
\begin{array}{ccll}
\text{a} & \Longrightarrow & \text{b} & \\[4pt]
\text{a}'' & \Longrightarrow & \text{a} & \text{si } \pi \geqslant \pi_1 \\[4pt]
\text{c} & \Longrightarrow & \text{b} & \text{si } \pi = \pi^{\pi_0}.
\end{array}
\right. \tag*{$(*)$}
\]
Mais a $\Longrightarrow$ b est tautologique, en prenant $I = C_0/F$, puisqu’on sait que $C_0/F$ est grande devant $\pi_0$ et que $F = \varinjlim\limits_{i \in I} x_i$, pratiquement par définition de $\mathcal{M} = \underbrace{\Ind_{\pi}(C)}_{\simeq \Ind_{\pi} C_0} \subset C^{\wedge}$. Donc il reste à établir les deux implications conditionnelles de $(*)$. L’implication a$''$ $\Longrightarrow$ a pour $\pi \geqslant \pi_1$ s’établit comme dans 4.5.2 (cf. page 109), c’est essentiellement sorital. Reste à établir
\[
\text{c} \;\Longrightarrow\; \text{b} \qquad \text{si} \;\; \pi \;=\; \pi^{\pi_0} \;. \tag*{$(*)$}
\]
Nous aurons besoin du résultat auxiliaire :

\msp{page 165}
\textbf{Lemme 4.7.3.} \textit{Soit $\pi \geqslant \pi_0$ tel que $\pi^{\pi_0} = \pi$ (donc en fait $\pi > \pi_0$). Pour tout objet $F$ de $\mathcal{M}$, on a}
\[
F \;=\; \varinjlim_{i \in I} x_i \qquad \text{(\footnotemark[67])}\;. \tag*{(4.7.3.1)}
\]
\footnotetext[67]{les $x_i$ dans $\Filt^{\pi}$, $I$ ensemble ordonné grand devant $\pi$.}
\textit{Si de plus $\pi' > \pi$ et $F \in \Ob \Filt^{\pi'}$, alors on peut prendre $I$ telle que}
\[
\card I \;\leqslant\; \pi'^{\pi} \;. \tag*{(4.7.3.2)}
\]

En effet, on a
\[
F \;=\; \varinjlim_{j \in J} x_j \qquad \text{avec } J = C_0/F \qquad \text{(\footnotemark[68])}\;,
\]
\footnotetext[68]{où $J$ peut aussi se remplacer par une catégorie cofinale dans [?] $C_0/F$ grande devant $\pi_0$.}
avec $J$ grande devant $\pi_0$. Soit alors
\[
I \;\;=\;\; \parbox[t]{0.62\textwidth}{ensemble des sous-catégories $J'$ de $J$ \textit{grandes devant} $\pi_0$, et telles que $\card J' \leqslant \pi$.} \tag*{(4.7.3.3)}
\]

\printedpage{55}
% head: A. Grothendieck : Dérivateurs, Ch. XVIII, Édition des Universités de Montpellier II et Paris VII

Je dis que $I$ (ordonné par inclusion) \textit{est grand devant} $\pi$ (\footnote[69]{cf. démonstration 4.7.5, page 168-171.}). Cela résulte aussitôt du lemme 4.5.5, p. 113, (énoncé d’abord, de façon erronée, en supposant seulement $\pi > \pi_0$ et non $\pi^{\pi_0} = \pi$). Posons alors, pour tout $J' \in I$,
\[
x_{J'} \;=\; \varinjlim_{j \in J'} x_j \;,
\]
on a donc
\[
x_{J'} \;\in\; \Ob \Filt^{\pi}
\]
par définition de $I$ et de $\Filt^{\pi}$. D’autre part,
\[
J' \;\longmapsto\; x_{J'} \;:\; I \;\longrightarrow\; \Filt^{\pi}
\]
\msp{page 166}
est un foncteur, donc
\[
\varinjlim_{J' \in I} x_{J'} \;\simeq\; \varinjlim_{j \in J} x_j \;\simeq\; x \;, % [sic] x -> F, the colimit is the object F of (4.7.3.1)
\]
d’où la première assertion de 4.7.3. Reste la majoration de cardinaux pour $F \in \Ob \Filt^{\pi}$. Dans ce cas, quitte à remplacer $C_0/F$ par une catégorie cofinale grande devant $\pi_0$, $\card J \leqslant \pi'$. On voit alors tout de suite que $\card I \leqslant \pi'^{\pi_0}$, ce qui est la majoration annoncée.

Nous pouvons maintenant prouver l’implication $(*)$, page 164, i.e.
\[
\mathcal{M}_{\pi} \;\subset\; \Filt^{\pi} \qquad (\text{pour } \pi = \pi^{\pi_0}).
\]
Soit donc $F$ dans $\mathcal{M}_{\pi}$, i.e. $\pi$-accessible, écrivons $F$ sous la forme (4.7.3.2), avec $I$ grande devant $\pi$. Par définition de la $\pi$-accessibilité, appliqué à $F \xrightarrow{\mathrm{id}} F \simeq \varinjlim F_i$, on voit que $F$ est un facteur direct d’un $F_i$, donc d’un objet de $\Filt^{\pi}$. Faute de savoir que $\Filt^{\pi}$ est stable par facteurs directs, on ne peut en conclure que $F$ est dans $\Filt^{\pi}$ (aussi je ne peux prouver c $\Longrightarrow$ b), mais seulement que $F$ est dans $\Filt'^{\pi}$, qui, lui, est stable par $\varinjlim$ grandes devant $\pi_0$, de cardinal $\pi$ donné $> \pi_0$, donc aussi \msp{page 167} stable par facteurs directs. Cela achève la démonstration de 4.7.1, et en même temps on obtient :

% marginal vertical bar beside Corollaire 4.7.4 and its statement
\textbf{Corollaire 4.7.4.} \textit{Si $\pi \geqslant \pi_0$ est tel que $\pi^{\pi_0} = \pi$, alors on a}
\[
\Filt'^{\pi} \;=\; (\Filt^{\pi})^{\natural} \;, \tag*{(4.7.4.1)}
\]
\textit{i.e.} $\Filt'^{\pi}$ \textit{est formée des objets $F$ de $\mathcal{M}$ qui sont facteurs directs d’un objet de} $\Filt^{\pi}$.

En effet, comme $\Filt^{\pi} \subset \Filt'^{\pi}$, et que $(\Filt^{\pi})^{\natural} \subset \Filt'^{\pi}$, et que $\Filt'^{\pi}$ est stable par facteurs directs, on a $(\Filt^{\pi})^{\natural} \subset \Filt'^{\pi}$ pour tout $\pi \geqslant \pi_0$. Mais inversement, comme $\Filt'^{\pi} \subset \mathcal{M}_{\pi}$, l’argument précédent montre que pour $\pi = \pi^{\pi_0}$, on a $\Filt'^{\pi} \subset \mathcal{M}_{\pi} \subset (\Filt^{\pi})^{\natural}$, ce qui prouve 4.7.3.

\textbf{Remarque 4.7.4.1} \texttt{[deux fois 4.7.4.1]}. Bien sûr, par 4.7.2, 1°, dès que $\pi \geqslant \pi_1$, on a même $\Filt^{\pi} = \Filt'^{\pi}$ (sans avoir à supposer $\pi = \pi^{\pi_0}$), donc \textit{si} $\pi > \pi_0$ (ce qui sera la cas si $\pi_1 > \pi_0$, i.e. $\pi_1 \neq \pi_0$). $\Filt'^{\pi}$ étant stable par facteurs directs, $\Filt^{\pi}$ l’est aussi, (4.7.4.1) est satisfait trivialement. Mais si $\pi = \pi_1 = \pi_0$, on a $\Filt^{\pi} = \Filt^{\pi_0} = \Filt'^{\pi_0} =$ sous-catégorie des foncteurs \textit{représentables} dans $C$ (équivalent à $C$), et celle-ci n’est égale à $C^{\natural}$ dans $\mathcal{M}$ que si $C$ est karoubienne, cf. \ldots{} ci-dessous. % [sic] la cas -> le cas, gender of "cas"

\printedpage{56}
% head: A. Grothendieck : Dérivateurs, Ch. XVIII, Édition des Universités de Montpellier II et Paris VII

\msp{page 168}
\textbf{4.7.5.} Pour que la démonstration de 4.7.1 soit complète, il me faudrait encore donner une démonstration de 4.5.5, p. 113, qui est le lemme technique clef pour établir le

\textbf{Lemme 4.7.5.1} (utilisé dans la démonstration de 4.7.3). \textit{Soient $\pi_0$, $\pi$ deux cardinaux infinis, avec $\pi \geqslant \pi_0$ et $\pi^{\pi_0} = \pi$ (donc $\pi > \pi_0$). Soit $J$ une catégorie grande devant $\pi_0$. Alors l’ensemble $I$, ordonné par inclusion, des sous-catégories $J'$ de $J$ qui sont grandes devant $\pi_0$, et telles que} $\card J' \leqslant \pi$, \textit{est grand devant $\pi$. (De plus, si} $\card J \leqslant \pi'$, \textit{on a} $\card I \leqslant \pi'^{\pi}$\textit{.)}

Il faut donc montrer le

\textbf{Corollaire 4.7.5.2} (\footnote[70]{Comparer avec la démonstration de 4.6.1.2.}) (c’est le lemme 4.5.5 non encore démontré). \textit{Sous les conditions de 4.7.5, soit $K \subset \mathrm{Fl} J$ tel que} $\card K \leqslant \pi$. \textit{Alors il existe $J' \in I$ tel que $K \subset J'$.}

On sait, $\pi$ étant infini, que la catégorie $K'$ engendrée par $K$ est aussi de cardinal $\leqslant \pi$, donc on peut supposer que $K$ “est” une sous-catégorie de $J$. On va construire une $J' \in I$ contenant $K$, en construisant par récurrence transfinie une suite croissante $K_{\alpha}$ de sous-catégories de $J$, avec $\card K_{\alpha} \leqslant \pi$, indexée par les ordinaux $\alpha$ tels que \fbox{$\card \alpha \leqslant \pi$} :

\msp{page 169}
\begin{enumerate}
\item[1°)] $K_0 = K$.
\item[2°)] $K_{\alpha} = \bigcup_{\beta < \alpha} K_{\beta}$ ($= \varinjlim K_{\beta}$) si $\alpha$ est un ordinal limite (donc $\card K_{\alpha} \leqslant \pi$, puisque les $\card K_{\beta} \leqslant \pi$).
\item[3°)] $K_{\alpha+1}$ déduit de $K_{\alpha}$ de la façon suivante. Si $K_{\alpha}$ est grande devant $\pi_0$, on prend $K_{\alpha+1} = K_{\alpha}$ (donc, compte tenu de 2°, on aura $K_{\alpha'} = K_{\alpha}$ pour $\alpha' \geqslant \alpha$ -- la construction transfinie s’arrête \ldots). Sinon, $K_{\alpha+1}$ s’obtient ainsi :
\begin{enumerate}
\item[a)] pour toute partie $A$ de cardinal $\leqslant \pi_0$ de $\Ob K_{\alpha}$, on choisit un majorant $x_A$ dans $J$ (on a envie de le prendre dans $K_{\alpha}$, s’il y a un majorant dans $K_{\alpha}$)\,;
\item[b)] pour tout couple $x$, $y$ d’objets de $K_{\alpha}$, et toute partie de cardinal $\leqslant \pi_0$ de $\mathrm{Hom}_{K_{\alpha}}(x,y)$, on choisit une flèche $u_B : y \longrightarrow z_B$ dans $J$, qui égalise les $v_{\beta} \in B$ (et on prend cette flèche dans $K_{\alpha}$, quand c’est possible)\,;
\item[c)] on prend pour $K_{\alpha+1}$ la plus petite sous-catégorie contenant les $\mathrm{id}_{x_A}$ et les $u_B$.
\end{enumerate}
\end{enumerate}

Il faut vérifier que
\[
\card K_{\alpha+1} \;\leqslant\; \pi \;.
\]
Le nombre des choix possibles pour $A \in \mathfrak{P}_{\leqslant \pi_0}(\Ob J)$ \msp{page 170} est $\leqslant \pi^{\pi_0} = \pi$, le nombre des choix possibles pour $B \in \mathfrak{P}_{\pi_0}(\mathrm{Fl}\, J)$ est aussi $\leqslant \pi^{\pi_0} = \pi$, donc le nombre des flèches de la forme $\mathrm{id}_{x_A}$ ou $u_B$ est $\leqslant \pi + \pi = \pi$. Mais alors la sous-catégorie $K_{\alpha+1}$ de $J$ engendrée par cet ensemble de flèches est de cardinal $\leqslant 1$. Soit alors % [sic] \leqslant 1 -> \leqslant \pi, the statement to be verified is card K_{\alpha+1} \leqslant \pi
\printedpage{57}
% head: A. Grothendieck : Dérivateurs, Ch. XVIII, Édition des Universités de Montpellier II et Paris VII
$\alpha_0$ le plus petit ordinal tel que
\[
\card \alpha_0 \;>\; \pi_0 \;,
\]
on a donc
\[
\card \alpha_0 \;=\; \text{cardinal successeur de } \pi_0 \leqslant \pi \;,
\]
puisque $\pi^{\pi_0} = \pi$, d’où $\pi > \pi_0$ (car $\pi^{\pi_0} = \pi_0^{\pi_0} \geqslant 2^{\pi_0} > \pi_0$). Ainsi, $\alpha_0$ fait partie des ordinaux $\alpha$ admissibles (i.e. tels que $\card \alpha \leqslant \pi$). Je dis que $K_{\alpha_0}$ \textit{est grande devant} $\pi_0$ (ce qui prouvera 4.7.5.2, en prenant $J' = K_{\alpha_0}$). Cela résulte immédiatement du fait que
\[
K_{\alpha_0} \;=\; \bigcup_{\alpha < \alpha_0} K_{\alpha} \qquad (\text{réunion \textit{croissante}}),
\]
% double marginal vertical bar beside the following two lines
et que l’ensemble des $\alpha < \alpha_0$ est \textit{grande devant} $\pi_0$ (\footnote[71]{expliciter ceci en un

\textbf{Lemme.} \textit{Soit $\pi_0$ un cardinal, $\alpha_0 = \alpha_{\pi'_0}$ ($\pi'_0$ le successeur de $\pi_0$) le plus petit ordinal tel que} $\card \alpha_0 > \pi_0$, \textit{alors $I_{\alpha_0} = \{\alpha \mid \alpha < \alpha_0\}$ est grand devant $\pi_0$.}}). Donc tout ensemble d’objets ou de flèches de $K_{\alpha_0}$, de cardinal $\leqslant \pi_0$, est contenu dans un $K_{\alpha}$ avec $\alpha < \alpha_0$.

\msp{page 171}
Cela implique, s’il s’agit d’un ensemble d’objets, qu’il y a un majorant commun dans $K_{\alpha+1}$, a fortiori dans $K_{\alpha_0}$. S’il s’agit d’une partie d’un $\mathrm{Hom}_{K_{\alpha}}(x,y)$, qu’il y a un égalisateur $y \longrightarrow z$ dans $K_{\alpha+1}$, a fortiori dans $K_{\alpha_0}$. Donc le critère $\mathrm{PF}_{\pi}$ 1,2 de 4.4.4 (page 89) est satisfait par $K_{\alpha_0}$, ce qui achève la démonstration de 4.7.5.2.

\textbf{Corollaire 4.7.6} (de 4.7.1). \textit{Pour tout cardinal $\pi \geqslant \pi_0$, on a}
\[
\underbrace{\Filt'^{\pi}}_{\overset{\textit{si } \pi \geqslant \pi_1}{=}\ \Filt^{\pi}} \;\subset\; \mathcal{M}_{\pi} \;\subset\; \underbrace{\Filt^{\pi^{\pi_0}}}_{\overset{\textit{si } \pi^{\pi_0} \geqslant \pi_1}{=}\ \Filt^{\pi^{\pi_0}}} \qquad (= (\Filt^{\pi^{\pi_0}})^{\natural} \textit{ par cor. 4.7.4})\,, \tag*{(4.7.6.1)} % ?? uncertain: tiny condition under the second brace -- "si \pi^{\pi_0} \geqslant \pi_1" or "si \pi^{\pi}_0 \geqslant \pi_1"
\]
\textit{donc si $\pi \geqslant \pi_1$, on a aussi}
\[
\Filt^{\pi} \;\subset\; \mathcal{M}_{\pi} \;\subset\; \Filt^{\pi^{\pi_0}} \qquad \textit{(si $\pi^{\pi_0} \geqslant \pi_1$).} \tag*{(4.7.6.2)}
\]

En effet, $\Filt'^{\pi} \subset \mathcal{M}_{\pi}$ n’est autre que l’implication b$'$ $\Longrightarrow$ c de 4.7.1, d’autre part, on a $\pi \leqslant \pi^{\pi_0}$, d’où
\[
\mathcal{M}_{\pi} \;\subset\; \mathcal{M}_{\pi^{\pi_0}} \;=\; \Filt'^{\pi^{\pi_0}} \;,
\]
où la dernière égalité provient de $\mathcal{M}_{\pi'} = \Filt'^{\pi'}$ si $\pi'^{\pi_0} = \pi'$, i.e. de (4.7.2.3) (page 163), en faisant $\pi' = \pi^{\pi_0}$, compte tenu que $\pi'^{\pi_0} = (\pi^{\pi_0})^{\pi_0} = \pi^{\pi_0 \cdot \pi_0} = \pi^{\pi_0} = \pi$.

\msp{page 172}
\textbf{Proposition 4.7.7.} \textit{Considérons $C$, $C_0$, $\mathcal{M} = \Ind_{\pi_0}(C)$, $\pi_1 = \sup(\pi_0, \card \mathrm{Fl}\, C_0)$ comme précédemment. Pour tout cardinal $\pi \geqslant \pi_0$, soit $\mathcal{F}^{\pi}$ la sous-catégorie strictement pleine} $\Filt'^{\pi}$ \textit{(resp. $(C^{\wedge})_{\pi} \cap \mathcal{M}$, resp. $\mathcal{M}_{\pi}$) dans $\mathcal{M}$. Alors la filtration croissante $(\mathcal{F}^{\pi})_{\pi \geqslant \pi_0}$ de $\mathcal{M}$ par les sous-catégories $\mathcal{F}^{\pi}$ satisfait les conditions suivantes.}
\begin{enumerate}
\item[\textit{a)}] \textit{Pour tout $\pi \geqslant \pi_0$, $\mathcal{F}^{\pi}$ est équivalente à une petite catégorie.}
\item[\textit{b)}] \textit{$\mathcal{M}$ satisfait à $\mathrm{L}_{\pi_0}$ (i.e. est stable par $\varinjlim$ filtrantes grandes devant $\pi_0$), et pour tout $\pi \geqslant \pi_0$, $\mathcal{F}^{\pi}$ est stable par les limites filtrantes grandes devant $\pi_0$ et de cardinal $\leqslant \pi_0$ [?].}
\printedpage{58}
% head: A. Grothendieck : Dérivateurs, Ch. XVIII, Édition des Universités de Montpellier II et Paris VII
\item[\textit{c)}] \textit{Si $\pi \geqslant \pi_0$ est} tel que \fbox{$\pi^{\pi_0} = \pi$,} \textit{alors pour tout $F$ dans $\mathcal{M}$, on a un isomorphisme}
\[
F \;=\; \varinjlim_{i \in I} x_i \;,
\]
\textit{les $x_i$ dans $\mathcal{F}^{\pi}$, et $I$ catégorie ordonnée grande devant $\pi_0$. Si de plus $\pi' > \pi$ et $F \in \Ob F^{\pi'}$, on peut prendre $I$ telle que} % [sic] F^{\pi'} (italic F) -> \mathcal{F}^{\pi'}, the filtration is denoted \mathcal{F}^{\pi}
\[
\card I \;\leqslant\; \pi'^{\pi} \qquad \underline{\textit{si}}\ \pi'^{\pi} \geqslant \pi_1. \tag*{(4.7.7.1)}
\]
\item[\textit{d)}] \textit{Si $\pi \geqslant \pi_0$, on a $\mathcal{F}^{\pi} \subset \mathcal{M}_{\pi}$, i.e. les objets de $\mathcal{F}^{\pi}$ sont $\pi$-accessibles.}
\end{enumerate}
\textbf{N.B.} Ce sont là exactement, à part d), les conditions a), b), c) des filtrations cardinales de SGA 4 I 9.7.12, \msp{page 173} à cela près seulement que dans c), on se restreint ici aux ordinaux $\pi$ tels que $\pi^{\pi_0} = \pi$ (\footnote[72]{et \textit{de plus}, dans la majoration (4.7.7.1), aux cardinaux $\pi' > \pi$ tels que l’on ait de plus $\pi'^{\pi} \geqslant \pi_1$ (ce qui est automatique si on a $\pi \geqslant \pi_1$, ce qui est le cas, vu que $\pi = \pi^{\pi_0} \geqslant 2^{\pi_0}$, si on a $2^{\pi_0} \geqslant \pi_1$.)}). La condition supplémentaire d) est une “condition d’accessibilité”, cf. plus bas.

\textsc{Démonstration de} 4.7.7. a) Comme on a
\[
\Filt'^{\pi} \;\subset\; (C^{\wedge})_{\pi} \cap \mathcal{M} \;\subset\; \mathcal{M}_{\pi} \;,
\]
on est réduit à prouver que $\mathcal{M}_{\pi}$ est équivalente à une petite catégorie. On est ramené au cas $\pi \geqslant \pi_1$, $\pi = \pi^{\pi_0}$, auquel cas on a $\mathcal{M}_{\pi} = \Filt^{\pi}$. En fait, on trouve :

\textbf{Corollaire 4.7.7.1} \texttt{[deux fois 4.7.7.1]} (\footnote[73]{Comparer avec 4.7.19.}). \textit{Si $\pi \geqslant \pi_1$, on a}
\[
\left\{
\begin{array}{rcll}
\cardred \Filt^{\pi} & \leqslant & 2^{\pi} & \\[4pt]
\cardred \Filt'^{\pi} & \leqslant & \cardred \mathcal{M}_{\pi} \;\leqslant\; 2^{(\pi^{\pi_0})} & \qquad \textit{($= 2^{\pi}$ si $\pi^{\pi_0} = \pi$).}
\end{array}
\right.
\]

La deuxième majoration résulte de la première et de
\[
\mathcal{F}^{\pi} \;\subset\; \mathcal{M}_{\pi} \;\subset\; \Filt^{\pi^{\pi_0}} \qquad (\text{pour } \pi^{\pi_0} \geqslant \pi_1) \qquad \text{(\footnotemark[74]).}
\]
\footnotetext[74]{Cf. (4.7.6.2).}
Pour la première, on utilise la condition a$''$) de 4.7.1 (qui caractérise $\Filt^{\pi}$ pour $\pi \geqslant \pi_1$). Je passe sur les détails de la vérification.

\msp{page 174}
Prouvons 4.7.2 b). La condition $\mathrm{L}_{\pi_0}$ est connue, reste la stabilité de $\mathcal{F}^{\pi}$. Dans le cas de $\Filt'^{\pi}$, elle est acquise par définition. Dans les cas $C^{\wedge}_{\pi} \cap \mathcal{M}$ et $\mathcal{M}_{\pi}$, elle résulte de la stabilité de la notion de $\pi$-accessibilité par $\varinjlim$ de cardinal $\leqslant \pi$, soit dans $C^{\wedge}$, soit dans $\mathcal{M}$ (compte tenu que $C^{\wedge}$ et $\mathcal{M}$ ont une famille génératrice $C$ formée d’objets $\pi_0$-accessibles, donc $\pi$-accessibles).

Prouvons c) : c’est essentiellement le lemme 4.7.3, compte tenu que l’on a
\[
\Filt^{\pi} \;\subset\; \mathcal{F}^{\pi} \;,
\]


\endgroup
% END UNIT der-ch18/der-ch18_p056-060.tex
% BEGIN UNIT der-ch18/der-ch18_p061-065.tex
\begingroup\makeatletter%
\newcommand{\printedpage}[1]{\marginpar{\footnotesize p.~#1}}%
\newcommand{\msp}[1]{\marginpar{\footnotesize [#1]}}%
\newcommand{\Ob}{\operatorname{Ob}}%
\newcommand{\Filt}{\mathrm{Filt}}%
\newcommand{\Ind}{\operatorname{Ind}}%
\newcommand{\card}{\operatorname{card}}%
\newcommand{\cardred}{\operatorname{card\,r\acute{e}d}}%
\makeatother\let\printedpage\derpp\let\msp\dermsp%
\let\mspin\DERmspin%


\noindent\printedpage{59}% head: A. Grothendieck : Dérivateurs, Ch. XVIII, Édition des Universités de Montpellier II et Paris VII
% foot: Version 17 décembre 2025
sauf qu'il faut faire attention à la majoration de cardinaux, $\card I \leqslant \pi'^{\pi}$. Mais si $F \in \Ob\mathcal{F}^{\pi'}$, comme par (4.7.6.1)
\[
\mathcal{F}^{\pi'} \;\subset\; \underbrace{\mathcal{M}_{\pi'}}_{\subset\mathcal{M}_{\pi'^{\pi}}\subset\Filt'^{\pi'^{\pi}}} \;\subset\; \Filt'^{\pi'^{\pi_0}} ,
\]
donc a fortiori
\[
\mathcal{F}^{\pi'} \;\subset\; \Filt^{\pi'^{\pi}} \underbrace{\;=\;}_{\text{si } \pi'^{\pi}\geqslant\pi_1} \Filt^{\pi'^{\pi}} , % [sic] first \Filt^{\pi'^{\pi}} -> \Filt'^{\pi'^{\pi}}, otherwise the equality under the brace would be trivial (cf. the preceding display)
\]
d'où, par 4.7.3, avec $\pi'$ remplacé par $\pi'^{\pi}=\pi''$, la possibilité de prendre $I$ avec
\[
\card I \;\leqslant\; \pi''^{\pi} \;=\; (\pi'^{\pi})^{\pi} \;=\; \pi'^{\pi\cdot\pi} \;=\; \pi'^{\pi} ,
\]
ce qui achève la démonstration.

\noindent\msp{page 175}\textbf{Remarque 4.7.7.2.} J'ai été obligé, dans l'énoncé de la condition c), de faire intervenir un deuxième cardinal $\pi_1$, et ceci faut de savoir si, pour $\pi''$ tel que $\pi''=\pi''^{\pi_0}$, $\Filt^{\pi''}$ est stable par facteurs directs. (Il faudrait, pour bien faire, soit prouver qu'il est stable, donc qu'on a dans 4.7.1 c $\Longrightarrow$ b, i.e. $\mathcal{M}_\pi\subset\Filt^{\pi}$ (et non seulement c $\Longrightarrow$ b$'$, i.e. $\mathcal{M}_\pi\subset\Filt'^{\pi}$) si $\pi=\pi^{\pi_0}$, ou bien faire un contre-exemple!) Cet ingrédient un peu étrange, on sera hélas obligé de le traîner dans la définition des filtrations cardinales généralisées, si on ne veut s'interdire de regarder dans le cas présent $(\mathcal{F}^{\pi})_{\pi\geqslant\pi_0}$ comme une telle filtration, et ne s'autoriser que de regarder $(\mathcal{F}^{\pi})_{\pi\geqslant\pi_1}$ comme telle, ce qui me paraîtrait ridicule. Pratiquement (même si théoriquement l'introduction de $\pi_1$ dans la description d'une filtration cardinale généralisée était bel et bien nécessaire) on n'aura guère des cas où il faille faire mention de $\pi_1$ dans la condition c. En effet, $\pi\geqslant\pi_0$ et $\pi^{\pi_0}=\pi$ impliquent $\pi\geqslant 2^{\pi_0}$, et dès lors $\pi'^{\pi}\geqslant 2^{(2^{\pi_0})}$. Si donc $\pi_1\leqslant 2^{(2^{\pi_0})}$, i.e. $\card C_0\leqslant 2^{(2^{\pi_0})}$, alors la condition $\pi^{\pi'}\geqslant\pi_1$ dans c) est automatiquement satisfaite, et il n'y a donc pas à la mentionner. Or il est bien rare qu'on ait à faire à des petites catégories $C_0$, dont le cardinal soit $>$ puissance % [sic] faut -> faute, "faute de savoir" intended; [sic] \pi^{\pi'}\geqslant\pi_1 -> \pi'^{\pi}\geqslant\pi_1, the condition is written \pi'^{\pi}\geqslant\pi_1 elsewhere (cf. 4.7.8)
\msp{page 176}
du continu $c$, et a fortiori, telle que $\card C_0 > 2^{c}$. Si par contre $\card C_0\leqslant 2^{c}$, comme $2^{c}\leqslant 2^{(2^{\pi_0})}$ (puisque $\pi_0\geqslant\aleph_0$), on aura bien $\pi_1\leqslant 2^{(2^{\pi_0})}$, et on est tranquille.

\textbf{Définition 4.7.8.} Soit $\mathcal{M}$ une $\mathfrak{U}$-catégorie. On appelle $\mathfrak{U}$-\textit{filtration cardinale généralisée} (ou simplement filtration cardinale généralisée, quand $\mathfrak{U}$ est sous-entendu) la donnée d'un cardinal infini $\pi_0$, et d'une filtration $(\mathcal{F}^{\pi})_{\pi\geqslant\pi_0}$ de $\mathcal{M}$ par des sous-catégories strictement pleines $\mathcal{F}^{\pi}$, indexées par les $\mathfrak{U}$-cardinaux $\geqslant\pi_0$, et satisfaisant les conditions a) b) c) énoncées dans 4.7.7, où dans c) on pose
\[
\pi_1 \;=\; \sup(\pi_0,\cardred\mathcal{F}^{\pi_0}) \, . \tag*{(4.7.8.1)}
\]
(Donc si $\cardred\mathcal{F}^{\pi_0}\leqslant 2^{(2^{\pi_0})}$, il n'y a pas à tenir compte de la condition $\pi'^{\pi}\geqslant\pi_1$, cf. 4.7.7.1.) On dit que la filtration cardinale généralisée (fcg) est \textit{accessible}, si elle satisfait de plus la condition d). On dit qu'elle est \textit{canonique} si on a même
\[
\mathcal{F}^{\pi} \;=\; \mathcal{M}_{\pi} \tag*{(4.7.8.2)}
\]
pour tout $\pi\geqslant\pi_0$.

\textbf{Proposition 4.7.9.} \textit{Soient $(\mathcal{F}^{\pi})_{\pi\geqslant\pi_0}$ une fcg sur la $\mathfrak{U}$-catégorie $\mathcal{M}$, et $\pi>\pi_0$ un cardinal tel que $\pi^{\pi_0}=\pi$. Alors on a}
\printedpage{60}% head: A. Grothendieck : Dérivateurs, Ch. XVIII, Édition des Universités de Montpellier II et Paris VII
% foot: Version 17 décembre 2025
\msp{page 177}
\[
\left\{
\begin{aligned}
&\mathcal{M}_\pi\subset\mathcal{F}^{\pi}\quad\text{\textit{(si $\pi=\pi^{\pi_0}$),}}\\
&\text{\textit{donc} }\mathcal{M}_\pi=\mathcal{F}^{\pi}\text{ \textit{si de plus la fcg est} accessible.}
\end{aligned}
\right.
\tag*{(4.7.9.1)}
\]

\textsc{Démonstration} (\footnote[75]{\texttt{[Dans ce qui suit, $\mathcal{F}^{\pi}$ a parfois été écrit $\mathcal{F}_{\pi}$, souvent corrigé en $\mathcal{F}^{\pi}$ par l'auteur. Nous écrivons $\mathcal{F}^{\pi}$ partout.]}}). Appliquant la condition c) des fcg à un $F\in\Ob\mathcal{M}_\pi$, on trouve que $F$ est facteur direct d'un $x_i\in\Ob\mathcal{F}^{\pi}$. Comme $\mathcal{F}^{\pi}$ est stable par $\varinjlim$ grandes devant $\pi_0$ et de cardinal $\leqslant\pi$, et que \textit{de plus} $\pi>\pi_0$, on conclut que :

\textbf{Lemme 4.7.9.1} \texttt{[deux fois 4.7.9.1]}. \textit{Si $\pi>\pi_0$ (\footnote[76]{Mais $\mathcal{F}^{\pi_0}$ n'est pas forcément stable par facteurs directs, comme on voit en prenant $\mathcal{M}=\Ind_{\pi_0}(C)$, $C$ une petite catégorie non karoubienne, et la filtration cardinale pour les $\Filt'^{\pi}$ -- on note que $\Filt'^{\pi_0}=\Filt^{\pi_0}\simeq e$, vu qu'une catégorie grande devant $\pi_0$ et de cardinal $\leqslant\pi_0$ est cofinale à $e$.}), $\mathcal{F}^{\pi}$ est stable par facteurs directs.}

D'où résulte que $F\in\mathcal{F}^{\pi}$, d'où (4.7.9.1).

\textbf{Corollaire 4.7.10.} \textit{Supposons la fcg accessible. Pour tout cardinal $\pi\geqslant\pi_0$, on a}
\[
\mathcal{F}^{\pi} \;\subset\; \mathcal{M}_\pi \;\subset\; \mathcal{F}^{\pi^{\pi_0}} \; ( \;=\; \mathcal{M}_{\pi^{\pi_0}}) \, . \tag*{(4.7.10.1)}
\]

En effet, $\mathcal{F}^{\pi}\subset\mathcal{M}_\pi$ par définition de fcg accessibles, et $\mathcal{M}_\pi=\mathcal{M}_{\pi^{\pi_0}}$ (puisque $\pi\leqslant\pi^{\pi_0}$), or $\mathcal{M}_{\pi^{\pi_0}}=\mathcal{F}^{(\pi^{\pi_0})}$ en vertu de (4.7.9.1) appliqué à $\pi'=\pi^{\pi_0}$ au lieu de $\pi$ (vu que $\pi'^{\pi_0}=\pi'$), d'où (4.7.10.1). On notera que cette inclusion généralise (4.7.9.1), elle donne $\mathcal{F}^{\pi}=\mathcal{M}_\pi$ si $\pi=\pi^{\pi_0}$.

\textbf{4.7.11.} Il me semble que les seules filtrations cardinales généralisées utiles sont celles qui sont \textit{accessibles.} Par la suite, je vais me borner à ce cas. Comme alors $\mathcal{F}^{\pi}\subset\mathcal{M}_\pi$, on en conclut
\msp{page 178}
que le foncteur canonique
\[
\Ind_\pi(\mathcal{F}^{\pi}) \;\hookrightarrow\; \mathcal{M} \tag*{(4.7.11.1)}
\]
prolongeant l'inclusion $\mathcal{F}^{\pi}\subset\mathcal{M}$, de façon à commuter aux $\varinjlim$ grandes devant $\pi$, est \textit{pleinement fidèle.} De plus, lorsque $\pi^{\pi_0}=\pi$, la condition c) de fcg nous dit entre autres que ce foncteur est essentiellement surjectif, donc que c'est une \textit{équivalence de catégories} :
\[
\Ind_\pi(\mathcal{F}^{\pi}) \xrightarrow{\ \sim\ } \mathcal{M} \qquad (\text{si } \pi=\pi^{\pi_0}). \tag*{(4.7.11.2)}
\]

\textbf{Proposition 4.7.12.} \textit{Soit $(\mathcal{F}^{\pi})_{\pi\geqslant\pi_0}$ une fcg accessible (fcga) sur la $\mathfrak{U}$-catégorie $\mathcal{M}$, et soit $\pi$ un cardinal $\geqslant\pi_0$. Les conditions suivantes sont équivalentes.}
\begin{enumerate}[label=\textit{\alph*)}]
\item \textit{Le foncteur pleinement fidèle (4.7.11.1) est essentiellement surjectif, i.e. est une équivalence de catégories.}
\item \textit{Pour tout $F\in\Ob\mathcal{M}$, $\mathcal{F}^{\pi}/F$ est une catégorie filtrante et grande devant $\pi$.}
\end{enumerate}

\noindent\printedpage{61}% head: A. Grothendieck : Dérivateurs, Ch. XVIII, Édition des Universités de Montpellier II et Paris VII
% foot: Version 17 décembre 2025
\textit{Quand ces conditions sont satisfaites, et si $\pi>\pi_0$, on a}
\[
\mathcal{F}^{\pi} \;=\; \mathcal{M}_\pi \, . \tag*{(4.7.12.1)}
\]

\textsc{Démonstration}. L'équivalence de a) et b) résulte aussitôt de 4.4.2 et de la définition des catégories grandes devant $\pi$ (p.~81). Que ces
\msp{page 179}
conditions impliquent $\mathcal{F}^{\pi}=\mathcal{M}_\pi$, i.e. $\mathcal{M}_\pi\subset\mathcal{F}^{\pi}$ dans le cas où $\pi>\pi_0$, se voit comme (4.7.9.1) en utilisant le lemme 4.7.9.1.

\textbf{Définition 4.7.13.}
\begin{enumerate}[label={$\arabic*^{\circ}$)}]
\item Un cardinal $\pi\geqslant\pi_0$ est dit \textit{cardinal utile} pour la fcga $(\mathcal{F}^{\pi})_{\pi\geqslant\pi_0}$, s'il satisfait les conditions équivalentes a) et b) de 4.7.12, et si de plus dans le cas où $\pi=\pi_0$, on a $\mathcal{F}^{\pi}=\mathcal{M}_\pi$. (Donc si $\pi$ est un cardinal utile, on a toujours $\mathcal{F}^{\pi}=\mathcal{M}_\pi$.) Dans le cas contraire, on dit que $\pi$ est un \textit{cardinal superflu} (pour le fcga donnée). % [sic] le fcga donnée -> la fcga donnée, gender agreement (feminine)
\item Un cardinal $\pi$ est dit \textit{régulier} pour la fcg donnée, s'il est utile, et si de plus la majoration de cardinaux énoncée dans (4.7.7.1) est valable.
\end{enumerate}

\textbf{4.7.14} Ainsi, tous les cardinaux $\pi\geqslant\pi_0$ tels que $\pi^{\pi_0}=\pi$ sont utiles, et même réguliers pour la filtration cardinale \texttt{[généralisée]} donnée. Notons que si tous les cardinaux $\pi\geqslant\pi_0$ sont utiles, alors on a $\mathcal{F}^{\pi}=\mathcal{M}_\pi$ pour tout $\pi\geqslant\pi_0$, donc la fcga est canonique. Notons aussi que dans tous les cas, si $(\mathcal{F}^{\pi})_{\pi\geqslant\pi_0}$ est une fcga, alors $(\mathcal{M}_\pi)_{\pi\geqslant\pi_0}$ en est une aussi.

\noindent\msp{page 180}En effet, la condition a) des fcg (cf. 4.7.7) résulte de la même condition pour $(\mathcal{F}^{\pi})$, et de (4.7.10.1), $\mathcal{M}_\pi\subset\mathcal{F}^{\pi^{\pi_0}}$. La condition b) est une propriété générale de la $\pi$-accessibilité, compte tenu que $(\mathcal{F}^{\pi_0})_0$ est une petite sous-catégorie génératrice formée d'objets $\pi_0$-accessibles de $\mathcal{M}$. La condition d) est satisfaite tautologiquement, et il reste à vérifier la condition c). Mais on a vu que pour $\pi=\pi^{\pi_0}$, on a $\mathcal{M}_\pi=\mathcal{F}^{\pi}$ (4.7.9.1), donc la condition c) pour $(\mathcal{M}_\pi)$ résulte de la même condition pour $(\mathcal{F}^{\pi})$. Ainsi, on trouve :

\textbf{Proposition 4.7.15.} \textit{Si $(\mathcal{F}^{\pi})_{\pi\geqslant\pi_0}$ est une fcga sur la $\mathfrak{U}$-catégorie $\mathcal{M}$, il en est de même de $(\mathcal{M}_\pi)_{\pi\geqslant\pi_0}$. De plus, tout cardinal utile (resp. régulier) pour la première filtration, l'est pour la seconde (appelée la fcg} canonique associée \textit{à la filtration cardinale} \texttt{[généralisée]} \textit{donnée).}

La dernière assertion résulte du fait que si $\pi$ est utile pour $(\mathcal{F}^{\pi})$, alors on a $\mathcal{F}^{\pi}=\mathcal{M}_\pi$.

\textbf{4.7.16.} Il me semble qu'on peut dire que la fcg canonique associée à une fcga donnée
\msp{page 181}
est toujours ``meilleure'' que celle-ci (si celle-ci n'est canonique, i.e. si les deux ne sont égales), à cause des deux faits suivants :
\begin{enumerate}[label=\alph*)]
\item justement parce qu'elle est \textit{canonique}, i.e. ne dépend d'aucun autre choix que celui du seul cardinal $\pi_0$, et
\item parce qu'il y a a priori plus des cardinaux utiles et de cardinaux réguliers pour la fcg canonique.
\end{enumerate}
Donc au total, il me semble que les filtrations cardinales \texttt{[généralisées]} les seules utiles sont celles qui sont canoniques.

\noindent\printedpage{62}% head: A. Grothendieck : Dérivateurs, Ch. XVIII, Édition des Universités de Montpellier II et Paris VII
% foot: Version 17 décembre 2025
Notons qu'une filtration cardinale (au sens de SGA 4 I 9, que je conserve ici) qui est de plus \textit{accessible} (i.e. $\mathcal{F}^{\pi}\subset\mathcal{M}_\pi$ pour tout $\pi\geqslant\pi_0$), n'est autre qu'une fcga canonique, telle que \textit{tous} les cardinaux sont réguliers (a fortiori utiles) pour ladite filtration. Je présume (sans avoir cependant construit d'exemple) que les filtrations cardinales \texttt{[généralisées]} canoniques des catégories de la forme $\mathcal{M}=\Ind_{\pi_0}(C)$, $C$ une petite catégorie \textit{karoubienne} (pour que $\pi_0$ soit un cardinal utile) ne sont pas forcément des fc (non généralisées, i.e. au sens SGA), i.e. je présume qu'il arrive qu'il y ait des cardinaux $\pi\geqslant\pi_0$ non réguliers, et même superflus.

\noindent\msp{page 182}Il me faut enfin étudier ce qui se passe quand on change le cardinal initial $\pi_0$ d'une fcga.

\textbf{Proposition 4.7.17.} \textit{Soit $(\mathcal{F}^{\pi})_{\pi\geqslant\pi_0}$ une fcga sur la $\mathfrak{U}$-catégorie $\mathcal{M}$, et soit $\pi'_0$ un cardinal $\geqslant\pi_0$. Alors $(\mathcal{F}^{\pi})_{\pi\geqslant\pi'_0}$ est une fcga sur $\mathcal{M}$, qui est canonique si celle de départ est canonique. Pour un cardinal $\pi\geqslant\pi'_0$, celui-ci est utile pour $(\mathcal{F}^{\pi})_{\pi\geqslant\pi'_0}$ si et seulement s'il l'est pour $(\mathcal{F}^{\pi})_{\pi\geqslant\pi_0}$. Pour qu'il soit régulier pour $(\mathcal{F}^{\pi})_{\pi\geqslant\pi'_0}$, il} suffit \textit{qu'il le soit pour $(\mathcal{F}^{\pi})_{\pi\geqslant\pi_0}$, et c'est aussi nécessaire quand on suppose de plus que $2^{\pi}\geqslant\pi'_1\overset{\text{déf}}{=}\sup(\pi'_0,\cardred\mathcal{F}^{\pi'_0})$, a fortiori quand $\pi\geqslant\pi'_1$.}

\textsc{Démonstration}. Les conditions a) b) sont trivialement satisfaites par la deuxième filtration, étant satisfaites par la première, ainsi que d) si elle l'est par la première. Il en est de même de la condition c) si $\pi=\pi^{\pi'_0}$, qui implique en effet $\pi^{\pi_0}=\pi$, puisqu'en général
\[
\pi \;\leqslant\; \pi^{\pi_0} \;\leqslant\; \pi^{\pi'_0} \, .
\]
Reste à examiner la majoration de cardinaux
\[
\card I \;\longrightarrow\; \pi'^{\pi} \qquad \text{si } F\in\Ob\mathcal{F}^{\pi'} \ \underline{\text{et}}\ \pi'^{\pi}\geqslant\pi'_1 \, , % [sic] \longrightarrow -> \leqslant, a majoration (inequality) of cardinals is meant
\]
où cette fois $\pi'_1=\sup(\pi'_0,\cardred\mathcal{F}^{\pi'_0})$. En tout état de cause, on a $\pi'_1\geqslant\pi_1$, donc
\msp{page 183}
$\pi'^{\pi}\geqslant\pi'_1$ implique $\pi'^{\pi}\geqslant\pi_1$, donc on peut appliquer la condition c) des fcg pour $(\mathcal{F}^{\pi})_{\pi\geqslant\pi_0}$, pour conclure qu'elle reste encore vérifiée. Donc en somme, \textit{toutes} les conditions des fcg (resp. fcga) sont trivialement vérifiées pour la fcg \textit{restreinte} de la première. De même, le fait qu'un cardinal $\pi'\geqslant\pi'_0$ est utile est trivialement indépendant de quelle des deux filtrations on considère, elle ne dépend que de la donnée du cardinal $\pi'$ et de la sous-catégorie $\mathcal{F}^{\pi'}$ de $\mathcal{M}$. Reste la majoration de cardinaux, dans le cas où on s'intéresse à la \textit{régularité} de $\pi'$ pour les fcga envisagées. La majoration exigée est la même, mais la condition préalable de validité de cette majoration, savoir
\[
\pi'^{\pi} \;\geqslant\; \pi'_1 \qquad \text{resp.} \qquad \pi'^{\pi} \;\geqslant\; \pi_1 \, ,
\]
n'est pas la même. Comme $\pi'_1\geqslant\pi_1$, on voit que la condition préalable est moins restrictive dans le cas de la première filtration, que de la seconde, d'où le ``il suffit'' dans la dernière assertion. Mais lorsque $2^{\pi}\geqslant\pi'_1$, même la plus restrictive est automatique\msp{page 184}ment vérifiée, puisque $\pi'^{\pi}\geqslant 2^{\pi}\geqslant\pi'_1$, et alors la régularité de $\pi$ pour la deuxième filtration implique la régularité pour le première, q.e.d. % [sic] le première -> la première, gender agreement (feminine)

\noindent\printedpage{63}% head: A. Grothendieck : Dérivateurs, Ch. XVIII, Édition des Universités de Montpellier II et Paris VII
% foot: Version 17 décembre 2025
\textbf{Corollaire 4.7.18.} \textit{Supposons $\pi_0$ régulier pour la fcga (\footnote[77]{Il suffit pour cela que $\pi_0$ soit utile, i.e. $\mathcal{M}\xleftarrow{\ \sim\ }\Ind_{\pi_0}(\mathcal{F}^{\pi_0}=\mathcal{M}_{\pi_0})$, cela implique régulier, cf. 4.8.5.}). Si on restreint la fcga $(\mathcal{F}^{\pi})_{\pi\geqslant\pi_0}$ aux cardinaux $\pi\geqslant\pi'_0$, où $\pi'_0\geqslant\pi_1$, alors on trouve une fcga où la majoration de cardinaux (pour $\pi=\pi^{\pi_0}$) est valable sans condition préalable $\pi'^{\pi}\geqslant\pi'_1$ ; plus précisément, on a alors}
\[
\pi'_1 \;\leqslant\; 2^{({\pi'_0}^{\pi_0})} \qquad \text{\textit{(a fortiori $\pi'_1\leqslant 2^{(2^{\pi'_0})}$)},} \tag*{(4.7.18.1)}
\]
\textit{donc la condition préalable $\pi'^{\pi}\geqslant\pi'_1$ est automatiquement vérifiée si $\pi\geqslant\pi_0$, $\pi=\pi^{\pi'_0}$, $\pi'>\pi$ (cf. remarque 4.7.7.2).}

Ceci équivaut donc à la majoration de cardinaux suivante (appliquée au cas de $\pi=\pi'_0$).

\textbf{Lemme 4.7.19.} \textit{Supposons $\pi_0$ régulier pour la fcga $(\mathcal{F}^{\pi})_{\pi\geqslant\pi_0}$. Alors pour tout cardinal $\pi\geqslant\pi_1$ ($\overset{\text{déf}}{=}\sup(\pi_0,\cardred\mathcal{F}^{\pi_0})$), on a}
\[
\cardred\mathcal{F}^{\pi} \;\leqslant\; 2^{(\pi^{\pi_0})} \, , \tag*{(4.7.19.1)}
\]
\textit{et si $\pi'\geqslant\pi_1$ et $F\in\Ob\mathcal{F}^{\pi'}$, on a}
\[
\cardred\mathcal{F}^{\pi}/F \;\leqslant\; \pi'^{(\pi^{\pi_0})} \qquad \text{(\footnotemark[78])} \, . \tag*{(4.7.19.2)}
\]
\footnotetext[78]{\textbf{N.B.} J'ignore si on peut remplacer dans ces majorations $\pi^{\pi_0}$ par $\pi_0$.}

\noindent\msp{page 185}\textsc{Démonstration du lemme}. Comme $\pi_0$ est régulier, posant $C=\mathcal{F}^{\pi_0}=\mathcal{M}_{\pi_0}$, on trouve $\mathcal{M}\simeq\Ind_{\pi_0}(C)$, donc on peut supposer $\mathcal{M}=\Ind_{\pi_0}(C)$, et appliquer les résultats prouvés pour les catégories de cette forme. Comme $\mathcal{F}^{\pi}\subset\mathcal{M}_\pi$, $\mathcal{F}^{\pi'}\subset\mathcal{M}_{\pi'}$, les majorations (4.7.9.1) et (4.7.9.2) sont contenues dans % [sic] (4.7.9.1) et (4.7.9.2) -> (4.7.19.1) et (4.7.19.2), the majorations of Lemme 4.7.19 are meant
\[
\left\{
\begin{array}{rcll}
\cardred\mathcal{M}_\pi & \leqslant & 2^{(\pi^{\pi_0})} & \\
\cardred\mathcal{M}_\pi/F & \leqslant & \pi'^{(\pi^{\pi_0})} & \text{si } F\in\Ob\mathcal{M}_{\pi'},\ \text{avec}\\
 & & & \pi'\geqslant\pi\geqslant\pi_1.
\end{array}
\right\}
\ \text{pour } \pi\geqslant\pi_1 \qquad \text{(\footnotemark[79])}. \tag*{$(*)$}
\]
\footnotetext[79]{Comparer avec majorations $\cardred\Filt^{\pi}\leqslant 2^{\pi}$, $\cardred\Filt^{\pi}/F\leqslant\pi'^{\pi}$, si $F\in\Ob\Filt^{\pi'}$, avec $\pi'\geqslant\pi\geqslant\pi_1$, majorations plus agréables que celles de (4.7.19). Il faudrait quand même s'assurer sur un exemple si les majorations plus sympathiques ne seraient pas valables dans le cas général.}
La première de ces majorations n'est autre que la deuxième majoration dans 4.7.7.1. Notons qu'on a bien
\[
2^{(\pi^{\pi_0})} \;\leqslant\; 2^{(2^{\pi})} \, ,
\]
ce qui résulte en effet de
\[
\pi^{\pi_0} \;\leqslant\; 2^{\pi} \, ,
\]
qu'on prouve en notant que $\pi<2^{\pi}$, d'où $\pi^{\pi_0}\leqslant(2^{\pi})^{\pi_0}=2^{\pi\cdot\pi_0}=2^{\pi}$ puisque $\pi\geqslant\pi_0$. Il reste à établir la deuxième majoration
\[
\cardred\mathcal{M}_\pi/F \;\leqslant\; \pi'^{(\pi^{\pi_0})} \qquad \text{si } F\in\Ob\mathcal{M}_{\pi'},\ \text{avec } \pi'\geqslant\pi\geqslant\pi_1 \, . \tag*{$(*)$}
\]
(On n'en a pas besoin pour 4.7.18, mais ce sera sans doute utile par ailleurs.) On utilise $\mathcal{M}_{\pi'}\subset\Filt^{\pi'^{\pi_0}}$ si $\pi'\geqslant\pi_1$ (4.7.6.2), donc si $F\in\Ob\mathcal{M}_{\pi'}$, on a


\endgroup
% END UNIT der-ch18/der-ch18_p061-065.tex
% BEGIN UNIT der-ch18/der-ch18_p066-070.tex
\begingroup\makeatletter%
\newcommand{\printedpage}[1]{\marginpar{\footnotesize p.~#1}}%
\newcommand{\msp}[1]{\marginpar{\footnotesize [#1]}}%
\newcommand{\Ob}{\operatorname{Ob}}%
\newcommand{\Fl}{\operatorname{Fl}}%
\newcommand{\card}{\operatorname{card}}%
\newcommand{\Hom}{\operatorname{Hom}}%
\newcommand{\Filt}{\operatorname{Filt}}%
\newcommand{\Ind}{\operatorname{Ind}}%
\newcommand{\circled}[1]{\tikz[baseline=(c.base)]\node[draw,circle,inner sep=0.8pt](c){#1};}%
\makeatother\let\printedpage\derpp\let\msp\dermsp%
\let\mspin\DERmspin%


\printedpage{64}
% head: A. Grothendieck : Dérivateurs, Ch. XVIII, Édition des Universités de Montpellier II et Paris VII

\msp{page 186}

\[ \card F(x) \longrightarrow \pi'^{\pi_0} \qquad \forall\, x \in \Ob C , \] % [sic] \longrightarrow -> \leqslant, a cardinal inequality is expected (cf. the next display)

donc si $G \in \Ob \mathcal{M}_\pi \subset \Ob \Filt^{\pi^{\pi_0}}$, d'où
\[ \card G(x) \leqslant \pi^{\pi_0} \qquad \forall\, x \in \Ob C , \]
on trouve
\[ \card \Hom(G,F) \leqslant \underbrace{\underbrace{\card(\Ob C)}_{\leqslant \pi_1}\cdot(\pi'^{\pi_0})^{(\pi^{\pi_0})}}_{\pi'^{\pi_0\cdot\pi^{\pi_0}} = \pi'^{(\pi^{\pi_0})}} , \]
d'où
\[ \card \Ob \mathcal{M}_\pi/F \leqslant \underbrace{\underbrace{\card(\Ob \mathcal{M}_\pi)}_{\leqslant 2^{(\pi^{\pi_0})}}\cdot\pi'^{(\pi^{\pi_0})}}_{\gamma = \pi'^{\pi^{\pi_0}}} , \]
et enfin
\[ \card \Fl \mathcal{M}_\pi/F \leqslant \underbrace{\underbrace{\gamma\cdot\gamma}_{\gamma = \pi'^{(\pi^{\pi_0})}}\cdot\underbrace{\card(\Fl \mathcal{M}_\pi)}_{\leqslant 2^{(\pi^{\pi_0})}}}_{\pi'^{(\pi^{\pi_0})}} , \]
ce qui est la majoration annoncée (la même que celle prouvée dans 4.5.12, la démonstration étant essentiellement la même).

\msp{page 187}

\textbf{4.7.20.} Une fcg canonique est dite \textit{maximale} (où \textit{strictement canonique}) si elle n'est par déduite d'une fcg canonique strictement plus étendue $(\mathcal{F}^\pi)_{\pi\geqslant\pi_{00}}$ (donc avec $\pi_{00} < \pi_0$). Donc si $\mathcal{M}$ admet une fcga, alors elle admet même une fcg strictement canonique, et cette dernière est de plus unique. On va revenir là-dessus dans la section suivante. % [sic] où strictement -> ou strictement, alternative is meant; [sic] n'est par -> n'est pas, negation

\section*{4.8\quad Catégories $\mathfrak{U}$-accessibles}

\textbf{Proposition 4.8.1.} \textit{Soit $\mathcal{M}$ une $\mathfrak{U}$-catégorie, $\pi_0$ un cardinal infini. Considérons les conditions suivantes :}
\begin{itemize}[label={}]
\item[\textit{a)}] \textit{$\mathcal{M}$ satisfait à $L_{\pi_0}$ (i.e. est stable par petites $\varinjlim$ grandes devant $\pi_0$), et contient une petite sous-catégorie pleine $S$} strictement génératrice, \textit{formée d'objets $\pi_0$-accessibles.}
\item[\textit{b)}] \textit{Il existe une petite catégorie $C$, et une équivalence de catégories}
\[ \mathcal{M} \simeq \Ind_{\pi_0}(C) . \]
\item[\textit{c)}] \textit{$\mathcal{M}$ admet une filtration cardinale généralisée accessible $(\mathcal{F}^\pi)_{\pi\geqslant\pi_0}$, ayant $\pi_0$ comme cardinal initial, et telle que $\pi_0$ soit un cardinal utile pour ladite filtration.}
\end{itemize}
% foot: Version 17 décembre 2025

\printedpage{65}
% head: A. Grothendieck : Dérivateurs, Ch. XVIII, Édition des Universités de Montpellier II et Paris VII

\msp{page 188}
\begin{itemize}[label={}]
\item[\textit{c$'$)}] \textit{Comme c), avec $\pi_0$ cardinal régulier (et non seulement utile) pour la filtration cardinale généralisée accessible en question, et cette filtration étant supposée de plus} canonique.
\end{itemize}
\textit{On a alors le diagramme d'implications}
\[ c' \Longleftrightarrow c \Longleftrightarrow b \Longrightarrow a . \]

\textbf{4.8.1.1.} Il me faut ici préciser la terminologie, sur ce que j'entends par sous-catégorie (pleine) $C$ \textit{strictement génératrice} dans $\mathcal{M}$. Pour ceci, je me réfère à SGA 4 I 7.2, (i), aux conditions équivalentes b) et c) de loc. cit. :
\begin{itemize}[label={}]
\item[\circled{b}] $\forall\, F \in \Ob \mathcal{M}$, la flèche canonique
\[ \varinjlim_{i\in C/F} x_i \longrightarrow F \]
est un isomorphisme.
\item[\circled{c}] Le foncteur canonique
\[ \begin{array}{ccc}
\mathcal{M} & \longrightarrow & C^{\wedge} \\[1ex]
F & \longmapsto & (\underbrace{x}_{\in \Ob C} \longmapsto \underbrace{\Hom_{\mathcal{M}}(x,F)}_{\in \Ob (\mathrm{Ens})})
\end{array} \]
est pleinement fidèle.
\end{itemize}
Ces conditions impliquent :
\begin{itemize}[label={}]
\item[a)] $C$ est génératrice dans $\mathcal{M}$ par épimorphismes stricts, i.e. pour tout $F$ dans $\mathcal{M}$, la famille des flèches $x \longrightarrow F$ de source dans $C$, de but $F$, est épimorphique stricte.
\end{itemize}
Et lorsque $\mathcal{M}$ est stable par produits fibrés, cette condition a) est équivalente aux précédentes. Dans loc. cit. j'affirme

\msp{page 189}

cette équivalence sans hypothèse restrictive sur $\mathcal{M}$ – mais en annotation marginale j'affirme que cette assertion est \textit{fausse}, en dehors de l'hypothèse d'existence des produits fibrés dans $\mathcal{M}$. Pourtant, je viens de vérifier la démonstration donnée dans loc. cit. de l'implication a $\Longrightarrow$ b, et elle est correcte. Mais peut-être l'annotation résulte-t-elle d'une confusion entre la condition a), et la condition suivante, plus faible a priori, et qui lui ressemble comme un frère :
\begin{itemize}[label={}]
\item[a$'$)] Il existe une famille de flèches $u_i : x_i \longrightarrow F$, de source dans $C$, et de but $F$, qui est épimorphique stricte (cf. SGA 4 I 10.3).
\end{itemize}
On a envie de dire que quand on \textit{agrandisse} une famille de flèches épimorphique stricte, en rajoutant d'autres flèches de même but, elle doit rester épimorphique stricte (“a fortiori”) – ce qui donne l'impression que a$'$) doit impliquer a), donc lui être équivalent. Mais si on essaye de prouver cette propriété hypothétique % [sic] agrandisse -> agrandit, verb form
% foot: Version 17 décembre 2025
\printedpage{66}
% head: A. Grothendieck : Dérivateurs, Ch. XVIII, Édition des Universités de Montpellier II et Paris VII
\msp{page 190}
de stabilité pour les familles épimorphiques strictes, on voit qu'on n'y arrive que si on suppose que la famille initiale est de plus \textit{épimorphique universelle} (si les produits fibrés existent dans $\mathcal{M}$), ou (sous la réserve précédente) si pour tout $x' \longrightarrow F$, le crible de $\mathcal{M}/F'$ image inverse du crible de $\mathcal{M}/F$ engendré par les $x_i \longrightarrow F$ est un “crible épimorphique”. Cela fait entrevoir en effet qu'il doit y avoir des cas où la condition a$'$) est satisfaite, mais non la condition a) (équivalente à b et à c).

Mais nous sommes alors confrontés à la situation très déplaisante : si $C$ est une sous-catégorie strictement génératrice de $\mathcal{M}$, il ne s'ensuit pas pour autant qu'une sous-catégorie pleine $C'$ de $\mathcal{M}$ telle que $C' \supset C$ soit également strictement génératrice\,! Or j'ai toujours utilisé précédemment, sans y réfléchir à deux fois, qui si $C$ est strictement génératrice, il en est de même de $C' \supset C$. Donc il faudra que je sois à présent très attentif, % [sic] qui si -> que si, conjunction
\msp{page 191}
et que je vérifie si cette erreur $\pm$ tacite rend inopérantes certaines démonstrations, voire même, rend faux certains énoncés. Aussi, il me faudrait donner un exemple où $C$ est strictement génératrice, mais $C' \supset C$ ne l'est pas.

Finalement, réflexion faite, ces ennuis craints disparaissent. Je vais énoncer un résultat en forme :

\textbf{Proposition 4.8.2.} \textit{Soient $C$, $\mathcal{M}$ deux $\mathfrak{U}$-catégories, $i : C \longrightarrow \mathcal{M}$ un foncteur, d'où un foncteur}
\[ \left\{\begin{aligned}
&\varphi : \mathcal{M} \longrightarrow C^{\wedge}\\
&\varphi(X) = (x \longmapsto \Hom_{\mathcal{M}}(i(x),X))
\end{aligned}\right. \tag*{(4.8.2.1)} \]
\textit{(foncteur composé $\mathcal{M} \xrightarrow{\textit{can.}} \mathcal{M}^{\wedge} \underbrace{\xrightarrow{i^*}}_{\substack{\textit{i.e. composition}\\ \textit{avec } i}} C^{\wedge}$). Pour tout $X$ dans $\mathcal{M}$, on considère la catégorie $C/X \simeq C/\varphi(X)$, et l'injection} [\,?] \textit{canonique (dans $\mathcal{M}^{\vee}$)}
\[ \varinjlim_{x\in C/X} i(x) \longrightarrow X . \tag*{(4.8.2.2)} \]
\begin{itemize}[label={}]
\item[\textit{1°)}] \textit{Considérons les conditions suivantes.}
\begin{itemize}[label={}]
\item[\circled{a}] \textit{Pour tout $X \in \Ob \mathcal{M}$, la famille des flèches $i(x) \longrightarrow X$, indexée par $\Ob C/X$, est épimorphique stricte.}

\msp{page 192}
\item[\circled{b}] \textit{Pour tout $X \in \Ob \mathcal{M}$, la flèche canonique 4.8.2.2. est un isomorphisme, i.e. fait de $X$ une $\varinjlim$ des foncteurs canoniques $C/X \longrightarrow \mathcal{M}$.}
\item[\circled{c}] \textit{Le foncteur $\varphi$ (4.8.2.1) est pleinement fidèle.}
\end{itemize}
\textit{On a alors les implications}
\[ \mathrm{b} \Longleftrightarrow \mathrm{c} \Longrightarrow \mathrm{a} , \]
\textit{et si on suppose $i$ pleinement fidèle, alors les trois conditions} a, b, c \textit{sont équivalentes.} (\textbf{N.B.} On dit que le foncteur $i : C \longrightarrow \mathcal{M}$ est \textit{strictement générateur}, si les conditions équivalentes b, c sont satisfaites. Si $C$ est une sous-catégorie de $\mathcal{M}$ (pas
% foot: Version 17 décembre 2025
\printedpage{67}
% head: A. Grothendieck : Dérivateurs, Ch. XVIII, Édition des Universités de Montpellier II et Paris VII
nécessairement pleine), on dit que $C$ est \textit{strictement génératrice} si le foncteur d'inclusion est strictement générateur, i.e. satisfait aux conditions équivalentes b), c). Pour ceci, il faut donc que la condition a) soit satisfaite, i.e. que pour tout $X$ dans $\mathcal{M}$, la famille des flèches $x \longrightarrow X$ de source dans $C$, de but $X$, soit épimorphique stricte, et cette condition est aussi suffisante quand $C$ est sous-catégorie pleine de $\mathcal{M}$.)
\item[\textit{2°)}] (\footnote[80]{\texttt{[2° se trouve sur la page 193.]}}). \textit{Supposons que $i$ soit strictement générateur} (\footnote[81]{Il suffit même qu'il satisfasse la condition plus faible a) de 1°).}). \textit{Alors pour tout objet $X$ de $\mathcal{M}$, la famille des flèches $i(x) \longrightarrow X$, indexée par $\Ob C/X$, est non seulement épimorphique stricte (condition a) dans 1°), mais de plus, toute famille de flèches de but $X$, incluant la famille précédente, est également épimorphique stricte.}
\item[\textit{3°)}] \textit{Considérons un diagramme commutatif de foncteurs de $\mathfrak{U}$-catégories :}

\msp{page 193}
\[ \begin{tikzcd}[column sep=small]
C \arrow[rr, "\alpha"] \arrow[dr, "i"'] & & C' \arrow[dl, "i'"] \\
& E &
\end{tikzcd} \]
\textit{Si le foncteur $i$ est strictement générateur} (\footnote[82]{En fait, il suffit qu'il satisfasse la condition plus faible c) de 1°).}), \textit{et si de plus $i'$ est pleinement fidèle, alors $i'$ est strictement générateur.}
\end{itemize}

\textbf{Corollaire 4.8.3.} \textit{Soit $\mathcal{M}$ une $\mathfrak{U}$-catégorie, $C$ une sous-catégorie strictement génératrice, $C'$ une sous-catégorie} pleine \textit{de $\mathcal{M}$ qui contient $C$, alors $C'$ est strictement génératrice.}

C'est en effet un cas particulier de 4.8.2 3°.

\msp{page 194}

\textsc{Démonstration de 4.8.2}. La démonstration de 1°) est celle de SGA 4 I 7.2 (i). (Où on suppose que $C$ est une sous-catégorie pleine de $\mathcal{M}$, mais la démonstration donnée prouve en fait l'énoncé donné ici.) Prouvons 2°). Rajoutons donc à la famille des
\[ i(x) \longrightarrow X \qquad (x \in \Ob C/X) \]
une famille
\[ X_\alpha \xrightarrow{p_\alpha} X \qquad (\alpha \in A) \]
de nouvelles flèches, et prouvons que la famille totale ainsi obtenue est encore épimorphique stricte. On doit donc prouver que si $Y$ est un objet de $\mathcal{M}$, et si on a une famille de flèches $u_x : i(x) \longrightarrow Y$, $u_\alpha : X_\alpha \longrightarrow Y$, satisfaisant la condition de compatibilité habituelle pour des couples de flèches :
\[ (v_x,v_y) \qquad\qquad \begin{tikzcd}
& i(x) \arrow[dr, dotted, "p_x"] & \\
Z \arrow[ur, "v_x"] \arrow[dr, "v_y"'] & & X\ , \\
& i(y) \arrow[ur, dotted, "p_y"'] &
\end{tikzcd} \]
ou
% foot: Version 17 décembre 2025
\printedpage{68}
% head: A. Grothendieck : Dérivateurs, Ch. XVIII, Édition des Universités de Montpellier II et Paris VII
\[ (v_x,v_\alpha) \qquad\qquad \begin{tikzcd}
& i(x) \arrow[dr, dotted, "p_x"] & \\
Z \arrow[ur, "v_x"] \arrow[dr, "v_\alpha"'] & & X\ , \\
& X_\alpha \arrow[ur, dotted, "p_\alpha"'] &
\end{tikzcd} \]
\msp{page 195}
ou
\[ (v_\alpha,v_\beta) \qquad\qquad \begin{tikzcd}
& X_\alpha \arrow[dr, dotted] & \\
Z \arrow[ur, "v_\alpha"] \arrow[dr, "v_\beta"'] & & X\ , \\
& X_\beta \arrow[ur, dotted] &
\end{tikzcd} \]
savoir les relations
\[ \left\{\begin{aligned}
u_xv_x &= u_yv_y\\
u_xv_x &= u_\alpha v_\alpha\\
u_\alpha v_\alpha &= u_\beta v_\beta ,
\end{aligned}\right. \tag*{($*$)} \]
alors il existe une flèche
\[ u : X \longrightarrow Y \]
(nécessairement unique, d'après la propriété c) d'épimorphicité stricte de la famille initiale indexée par $C/X$) telle qu'on ait
\[ \left\{\begin{aligned}
up_x &= v_x \quad \text{pour tout } x \in \Ob C/X\\ % [sic] v_x -> u_x, the arrow u p_x goes to Y (cf. up_\alpha = u_\alpha below)
up_\alpha &= v_\alpha \quad \text{pour tout } \alpha \in A. % [sic] v_\alpha -> u_\alpha, cf. the relation up_\alpha = u_\alpha below
\end{aligned}\right. \tag*{($*$)} \]
Or d'après la propriété a), il existe un unique $u : X \longrightarrow Y$, de telle façon que la première des relations ($*$) soit satisfaite, pour tout $x \in \Ob C/X$. Donc il s'agit seulement de prouver que pour cet $u$, on a pour tout $\alpha$ la relation

\msp{page 196}
\[ up_\alpha = u_\alpha . \]
Or, en appliquant la propriété a) à $X_\alpha$ au lieu de $X$, on veut que la famille des flèches % [sic] veut -> voit, sense of the sentence
\[ i(y) \xrightarrow{q_y} X_\alpha \]
(indexée par $C/X_\alpha$) est épimorphique stricte, a fortiori épimorphique. Donc pour prouver $up_\alpha = u_\alpha$, il suffit de prouver qu'on a
\[ (up_\alpha)q_y = u_\alpha q_y \tag*{($**$)} \]
pour toute $y$ dans $C/X_\alpha$. Mais c'est là une conséquence de la compatibilité ($*$) (deuxième formule), appliquée au cas suivant,
% foot: Version 17 décembre 2025


\endgroup
% END UNIT der-ch18/der-ch18_p066-070.tex
% BEGIN UNIT der-ch18/der-ch18_p071-075.tex
\begingroup\makeatletter%
\newcommand{\printedpage}[1]{\marginpar{\footnotesize p.~#1}}%
\newcommand{\msp}[1]{\marginpar{\footnotesize [#1]}}%
\newcommand{\Ob}{\operatorname{Ob}}%
\newcommand{\Ind}{\operatorname{Ind}}%
\newcommand{\Filt}{\operatorname{Filt}}%
\newcommand{\card}{\operatorname{card}}%
\newcommand{\End}{\operatorname{End}}%
\newcommand{\Fl}{\operatorname{Fl}}%
\newcommand{\Kar}{\operatorname{Kar}}%
\makeatother\let\printedpage\derpp\let\msp\dermsp%
\let\mspin\DERmspin%


\printedpage{69}
% head: A. Grothendieck : Dérivateurs, Ch. XVIII, Édition des Universités de Montpellier II et Paris VII

\[
\begin{tikzcd}[column sep=small]
 & i(y) \arrow[dr, dotted, "p_\alpha q_y=p_x"] & \\
Z=i(y) \arrow[ur, "\mathrm{id}"] \arrow[dr, "q_y"'] & & X \ , \\
 & X_\alpha \arrow[ur, dotted, "p_\alpha"'] &
\end{tikzcd}
\]
on a donc, en posant $x=(y,p_\alpha q_y)\in\Ob C/X$,
\[
\underbrace{u_x\circ\mathrm{id}_{i(y)}}_{u_x} \ = \ u_\alpha q_y \ ,
\]
donc pour prouver $(**)$, il reste à vérifier
\[
u_x \ = \ \underbrace{(up_\alpha)q_y}_{u(p_\alpha q_y)\,=\,up_x} \ ,
\]
\msp{page 197}
qui est en effet vérifiée par définition de $u_x$. Cela achève de prouver $2^\circ$.

Prouvons $3^\circ$. Comme $i'$ est pleinement fidèle, en vertu de $1^\circ$, pour prouver que $i'$ est strictement génératrice, il suffit d'établir que pour tout objet $X$ de $\mathcal{M}$, la famille des flèches $i'(x')\longrightarrow X$, indexée par $\Ob C'/X$, est épimorphique stricte. Or cette famille contient la famille des $\underbrace{i'(\alpha(x))}_{=\,i(x)}\longrightarrow X$, indexée par $\Ob C/X$. Donc il résulte de l'hypothèse sur $i$, et de $2^\circ$, que cette famille est épimorphique stricte. On a gagné ! Ouf \dots

\textsc{Démonstration de 4.8.1.} On a
\[
\mathrm{c}' \ \Longrightarrow \ \mathrm{c} \ \Longrightarrow \ \mathrm{b} \ \Longrightarrow \ \mathrm{a} \, ,
\]
la première implication tautologique, la seconde résultant du fait que $\pi_0$ cardinal régulier pour la fcga envisagée, implique que
\[
\mathcal{M} \ \xleftarrow{\ \sim\ } \ \Ind_{\pi_0}(\mathcal{F}^{\pi_0}) \ ,
\]
$\mathcal{F}^{\pi_0}$ équivalente à une petite catégorie, d'où b).

\msp{page 198}
Enfin b) $\Longrightarrow$ a) résulte du fait que $\Ind_{\pi_0}(C)$ est stable par $\varinjlim$ grandes devant $\pi_0$ (4.4.15), et que $C$ est une sous-catégorie strictement génératrice de $\Ind_{\pi_0}(C)$, formée d'objets $\pi_0$-accessibles. Il nous faut expliciter le

\textbf{Lemme 4.8.4.} \textit{Soit $C$ une $\mathfrak{U}$-catégorie, $\pi_0$ un cardinal infini. Alors $C$ est strictement génératrice dans $\mathcal{M}=\Ind C$ et dans $\mathcal{M}=\Ind_{\pi_0}C$.}

\msp{page 199}
En effet, le foncteur canonique $\mathcal{M}\longrightarrow C^\wedge$ est pleinement fidèle, et on applique le critère c) de 4.8.2.

Pour achever de prouver 4.8.1, il reste donc à établir l'implication
\[
\mathrm{a} \ \Longrightarrow \ \mathrm{c}' \, .
\]

\printedpage{70}
% head: A. Grothendieck : Dérivateurs, Ch. XVIII, Édition des Universités de Montpellier II et Paris VII

Supposons donc que $\mathcal{M}$ satisfasse à $\mathrm{L}_{\pi_0}$, et admette une petite sous-catégorie strictement génératrice $S$, formée d'objets $\pi_0$-accessibles, donc un foncteur canonique
\[
\Ind_{\pi_0}S \ \longrightarrow \ \mathcal{M} \ ,
\]
prolongeant l'inclusion $S\hookrightarrow\mathcal{M}$, de façon à commuter aux $\varinjlim$ grandes devant $\pi_0$, et du fait que $S\subset\mathcal{M}_{\pi_0}$, ce foncteur est \textit{pleinement fidèle.} Mais comment prouver qu'il est essentiellement surjectif ? Ce ne sera d'ailleurs pas vrai sans autre hypothèse sur $S$. L'idée naturelle est de saturer $S$ pour les $\varinjlim$
\msp{page 200}
de cardinal $\leqslant\pi_0$ représentables dans $\mathcal{M}$, et notamment par facteurs directs, mais même alors (et même en admettant que $\mathcal{M}_{\pi_0}$ lui-même est équivalent à une petite catégorie, et prenant $C=\mathcal{M}_{\pi_0}$), il n'est pourtant pas clair que le foncteur pleinement fidèle $\Ind_{\pi_0}C\longrightarrow\mathcal{M}$ soit essentiellement surjectif, ce qui équivaudrait aussi à dire que pour tout $F$ dans $\mathcal{M}$, $C/F$ est grande devant $\pi_0$. Il n'est en effet pas même clair que $C/F$ soit filtrante.

Donc finalement, je doute qu'on ait a) $\Longrightarrow$ c$'$) ou seulement a) $\Longrightarrow$ b), la situation n'est pas bien comprise visiblement.

Mais prouvons quand même
\[
\mathrm{b} \ \Longrightarrow \ \mathrm{c}' \, .
\]
Pour ceci, on peut supposer $\mathcal{M}=\Ind_{\pi_0}C$, $C$ petite. On sait alors que la filtration de $\mathcal{M}$ par les $(\mathcal{M}_\pi)_{\pi\geqslant\pi_0}$ est une fcga (prop. 4.7.7 et définition 4.7.8), et même une fcg \textit{canonique.} De plus, je dis que le cardinal $\pi_0$ est \textit{utile} et même régulier pour cette filtration (ce qui achèvera de prouver c$'$). On a un résultat un peu plus précis :

\msp{page 201}
\textbf{Corollaire 4.8.5.} \textit{Soient $C$ une catégorie équivalente à une petite catégorie, $\pi_0$ un cardinal infini, $\mathcal{M}=\Ind_{\pi_0}(C)\subset C^\wedge$. Alors :}
\begin{enumerate}
\item[\textit{1$^\circ$)}] $\pi_0$ \textit{est un cardinal} utile \textit{et même régulier pour la fcg canonique} $(\mathcal{M}_\pi)_{\pi\geqslant\pi_0}$ \textit{de} $\mathcal{M}$.
\item[\textit{2$^\circ$)}] $\mathcal{M}_{\pi_0}=C^\natural$ \textit{(sous-catégorie pleine de $\mathcal{M}=\Ind_{\pi_0}C$ formée des facteurs directs dans $\mathcal{M}$ d'objets de $C$), et $\mathcal{M}_{\pi_0}$ est une enveloppe de Karoubi de $C$.}
\end{enumerate}
(\footnote{\textbf{Corollaire 4.8.5.1.} \textit{Si $C$ est karoubienne, alors $\mathcal{M}_{\pi_0}\subset C^\wedge$ est la catégorie formée des foncteurs} représentables \textit{dans} $C\overset{\sim}{\hookrightarrow}\mathcal{M}_{\pi_0}$.}, \footnote{\textbf{4.8.5.2.} Expliciter aussi ceci : Soit $u : C\longrightarrow C'$ un foncteur, avec $C$, $C'$ équivalentes à des petites catégories. Conditions équivalentes :
\begin{enumerate}
\item[a)] $u$ induit $\Kar C\xrightarrow{\ \sim\ }\Kar C'$.
\item[a$'$)] $u$ pleinement fidèle, et tout objet de $C'$ est isomorphe à un facteur direct d'un objet de l'image essentielle.
\item[b)] $\Ind(u) : \Ind(C)\longrightarrow\Ind(C')$ est une équivalence.
\item[c)] $\Ind_{\pi_0}(u) : \Ind_{\pi_0}(C)\longrightarrow\Ind_{\pi_0}(C')$ est une équivalence ($\pi_0$ un cardinal infini donné).
\item[d)] $u_! : C^\wedge\longrightarrow C'^\wedge$ est une équivalence.
\item[d$'$)] $u_* : C'^\wedge\longrightarrow C^\wedge$ est une équivalence \texttt{[plutôt} $u_* : C^\wedge\longrightarrow C'^\wedge$\texttt{]}.
\item[d$''$)] $u^* : C'^\wedge\longrightarrow C^\wedge$ est une équivalence.
\end{enumerate}}).

\printedpage{71}
% head: A. Grothendieck : Dérivateurs, Ch. XVIII, Édition des Universités de Montpellier II et Paris VII

\textsc{Démonstration.}

$1^\circ$) On sait déjà que
\[
\Ind_{\pi_0}\mathcal{M}_{\pi_0} \ \longrightarrow \ \mathcal{M} \tag*{$(*)$}
\]
est pleinement fidèle, et pour voir que c'est essentiellement surjectif, on note que l'on a
\[
\mathcal{M} \ \simeq \ \Ind_{\pi_0}C \ \hookrightarrow \ \Ind_{\pi_0}\mathcal{M}_{\pi_0} \ \longrightarrow \ \mathcal{M} \ ,
\]
le composé étant isomorphe au foncteur identique, est essentiellement surjectif, donc aussi $(*)$. Cela prouve donc que $\pi_0$ est un cardinal utile pour la filtration envisagée, prouvons que ce cardinal est même régulier, donc que pour
\[
F \ \in \ \Ob\mathcal{M}_\pi \ , \qquad \text{avec } \pi>\pi_0,\ \pi^{\pi_0}\geqslant\pi_1
\]
on a
\[
F \ = \ \varinjlim_I x_i \ , \qquad x_i\in\Ob\mathcal{M}_{\pi_0},\ I \text{ grande devant } \pi_0,\ \card I\leqslant\pi^{\pi_0}.
\]
\msp{page 202}
Pour ceci, on utilise (4.7.6.2), page 171,
\[
\mathcal{M}_\pi \ \subset \ \Filt^{\pi^{\pi_0}} \qquad (\text{si } \pi\geqslant\pi_0,\ \pi^{\pi_0}\geqslant\pi_1),
\]
et la définition de $\Filt^{\pi'}$ (pour $\pi'\geqslant\pi_0$) nous montre que c'est là la conclusion cherchée.

Prouvons donc $2^\circ$). Comme $\mathcal{M}_{\pi_0}$ est stable par les $\varinjlim$ inductives finies qui sont représentables dans $\mathcal{M}$ (SGA 4 I 9.9), il est stable par facteurs directs, et comme il contient $C$, on a donc $C^\natural\subset\mathcal{M}_{\pi_0}$. L'inclusion inverse $\mathcal{M}_{\pi_0}\subset C^\natural$ se voit de la façon standard, en écrivant un objet $F$ de $\mathcal{M}_{\pi_0}$ comme $\varinjlim\limits_I x_i$, les $x_i$ dans $C$, $I$ grande devant $\pi_0$, alors $F$ est facteur direct d'un des $x_i$.

Le fait que $C^\natural$ est une \textit{enveloppe de Karoubi} de $C$ dans $\mathcal{M}\subset C^\wedge$ provient du fait que $\mathcal{M}$ est stable par facteurs directs, et que $C$ est une sous-catégorie pleine de $\mathcal{M}$, cf. SGA 4 I (\,?\,\dots).

\textbf{4.8.6. Questions ouvertes.} On a ainsi prouvé 4.8.1 revu et corrigé. Il reste la question :
\begin{enumerate}
\item[1)] Donner un exemple où 4.8.1 a) est vérifié, mais non b) etc. (\footnote{Il suffit de trouver un cardinal non régulier, pour la fcg canonique d'une $\mathcal{M}=\Ind_{\pi_0}C$.}).
\item[2)] Dégager des conditions supplémentaires serviables qui, jointes à a), impliquent b), donc c, c$'$.
\end{enumerate}

\msp{page 203}
Signalons aussi pour mémoire les questions suivantes.
\begin{enumerate}
\item[3)] Soit $\mathcal{M}=\Ind_{\pi_0}(C)$, et $\pi>\pi_0$. Alors $\Filt^\pi$ est-il stable par facteurs directs ? Cela est-il vrai du moins si $\pi=\pi^{\pi_0}$, ou ce qui revient au même (cf. 4.7.3), a-t-on dans ce cas $\Filt^\pi=\Filt'^{\pi}$ ? Si oui, il serait inutile dans la description 4.7.7 des filtrations cardinales, dans la majoration de cardinaux (4.7.7.1), $\card I\leqslant\pi'^{\pi}$, de mettre l'hypothèse restrictive $\pi'^{\pi}\geqslant\pi_1$.
\end{enumerate}

Pour avoir une réponse affirmative à cette question, pour $\pi$ donné, il suffirait (et probablement faudrait) d'avoir une réponse affirmative à celle-ci : Soient $\mathcal{D}$, $\mathcal{D}_0$ deux catégories grandes devant $\pi_0$, $\mathcal{D}_0$ facteur direct de $\mathcal{D}$, $\mathcal{D}$ admettant une catégorie cofinale $I$ grande devant $\pi_0$, de cardinal $\leqslant\pi$. Alors en est-il de même de $\mathcal{D}_0$ ? Si $p$ est une rétraction de $\mathcal{D}$ sur $\mathcal{D}_0$, la difficulté est que je ne sais pas si $p$ est foncteur cofinal. En fait, je vois que
\printedpage{72}
% head: A. Grothendieck : Dérivateurs, Ch. XVIII, Édition des Universités de Montpellier II et Paris VII
\msp{page 204}
ce n'est pas toujours le cas, même si $\mathcal{D}$ est finie. Prenons p.ex. pour $\mathcal{D}_0$ une catégorie ayant objet final $e$, et une flèche $x\longrightarrow e$ qui n'est pas un isomorphisme, et admettant une section (p.ex. $\mathcal{D}_0$ la sous-catégorie pleine de $\boldsymbol{\Delta}$ formée par $\boldsymbol{\Delta}_0$, $\boldsymbol{\Delta}_1$ : \begin{tikzcd}[cramped] \boldsymbol{\Delta}_0 \arrow[r, shift left, "\partial_0"] \arrow[r, shift right, "\partial_1"'] & \boldsymbol{\Delta}_1 \arrow[l, bend left=60, "\kappa"] \end{tikzcd} ). Soit $\mathcal{D}$ déduit de $\mathcal{D}_0$ en ajoutant un objet final $\tilde{e}$. Alors $\mathcal{D}_0$, $\mathcal{D}$ sont filtrantes, et on obtient une rétraction de $\mathcal{D}$ sur $\mathcal{D}_0$, envoyant $\tilde{e}$ sur $x$. Ce foncteur n'est pas cofinal, car un foncteur cofinal doit envoyer $\tilde{e}$ en $e$.

Plus généralement, si $\mathcal{D}_0$ est une catégorie telle qu'il existe une flèche
\[
\alpha : e_{\mathcal{D}_0^\wedge} \ \longrightarrow \ x_0 \tag*{$(*)$}
\]
dans $\mathcal{D}_0^\wedge$, avec $x_0\in\Ob\mathcal{D}_0$, alors on voit que $\mathcal{D}_0$ est grande devant $\pi_0$ quel que soit le cardinal infini $\pi_0$, et que d'autre part, désignant encore par $\mathcal{D}$ la catégorie déduite de $\mathcal{D}_0$ en lui ajoutant un objet final, il y a une rétraction de $\mathcal{D}$ sur $\mathcal{D}_0$ qui envoie $e_{\mathcal{D}}$ dans $x_0$. Cette rétraction n'est cofinale que si $x_0$ est un objet final de $\mathcal{D}_0$, i.e. si et seulement si $(*)$ est un isomorphisme. Supposons qu'il n'en soit pas ainsi. Mais il y a un foncteur cofinal de la sous-catégorie $I_0$ de $\mathcal{D}_0$, telle que $\Ob I_0=\{x_0\}$, $\Fl I_0=\End_{I_0}(x_0)=\{p,\mathrm{id}_{x_0}\}$, où $p$ est le composé $x_0\xrightarrow{\text{can.}}e_{\mathcal{D}_0^\wedge}\xrightarrow{\alpha}x_0$. Donc pas de contre-exemple de cette façon !

\msp{page 205}
\begin{enumerate}
\item[4)] Considérons $\mathcal{M}=\Ind_{\pi_0}(C)$, $C$ petite, et considérons la filtration canonique $(\mathcal{M}_\pi)_{\pi\geqslant\pi_0}$. Donner des exemples où il y a des cardinaux non utiles, ou mieux, ou pour tout cardinal $\pi$, existe un cardinal $\pi'>\pi$ qui soit non utile. (Cela donnera l'exemple demandé dans $1^\circ$.) D'autre part, donner des exemples de cardinaux utiles qui ne soient réguliers -- si possible, comme tantôt, de tels cardinaux aussi grands qu'on veut.
\end{enumerate}

\textbf{Définition 4.8.7.} Soit $\mathcal{M}$ une $\mathfrak{U}$-catégorie, $\pi_0$ un cardinal infini. On dit que $\mathcal{M}$ est $\pi_0$-\textit{accessible} si $\mathcal{M}$ est équivalente à une catégorie $\Ind_{\pi_0}(C)$, $C$ une petite catégorie, i.e. si elle satisfait aux conditions équivalentes b, c, c$'$ de 4.8.1 (impliquant la condition a), sans doute strictement plus faible). On dit que $\mathcal{M}$ est accessible s'il existe un cardinal tel que $\mathcal{M}$ soit $\pi_0$-accessible (\textbf{N.B.} si $\pi_0$ est un tel cardinal, alors tout cardinal $\pi>\pi_0$ tel que $\pi^{\pi_0}=\pi$, notamment les cardinaux $\pi'^{\pi_0}$, avec $\pi'\geqslant 2$, ont la même propriété), ou ce
\msp{page 206}
qui revient au même d'après 4.8.1, si $\mathcal{M}$ admet une fcga, ou encore une fcg canonique.

\textbf{Commentaire 4.8.8.} Cette définition est en contradiction avec la terminologie introduite dans SGA 4 I 9.4 (p. 114) ; on désignera dorénavant sous les noms de ``catégorie faiblement $\pi_0$-accessible'', ``catégorie faiblement accessible'' les notions introduites dans 9.4 -- i.e. les catégories admettant une sous-catégorie pleine génératrice (pas même strictement génératrice) dont les objets sont $\pi_0$-accessibles (resp. accessibles).

Notons que si $\mathcal{M}$ est $\pi_0$-accessible (resp. accessible), la condition de 4.8.1 nous dit que $\mathcal{M}$ est faiblement $\pi_0$-accessible (resp. faiblement accessible), et même qu'il existe une petite sous-catégorie pleine \textit{strictement} génératrice, qui soit formée d'objets $\pi_0$-accessibles (resp. accessibles). Je doute fort que cette dernière condition nécessaire soit aussi suffisante pour
\printedpage{73}
% head: A. Grothendieck : Dérivateurs, Ch. XVIII, Édition des Universités de Montpellier II et Paris VII
que $\mathcal{M}$ soit $\pi_0$-accessible (resp. accessible). Mais en vertu de 4.6.3.3 (p. 150), il en est ainsi si on suppose $C$ stable par petites $\varinjlim$ , ou ce qui revient au même, moyennant $\mathrm{L}_{\pi_0}$, stable par limites inductives de cardinal $\leqslant\pi_0$.

\msp{page 207}
\textbf{Proposition 4.8.9.} \textit{Soit $\mathcal{M}$ une $\mathfrak{U}$-catégorie accessible (4.8.7), $\pi_0$ un cardinal tel que $\mathcal{M}$ satisfait à $L_{\pi_0}$. Conditions équivalentes :}
\begin{enumerate}
\item[\textit{a)}] \textit{$\mathcal{M}$ est une catégorie paratopos (cf. 4.6.3, $3^\circ$, p. 150).}
\item[\textit{b)}] \textit{$\mathcal{M}$ est stable par petites limites inductives.}
\item[\textit{b$'$)}] \textit{$\mathcal{M}$ est stable par limites inductives de cardinal $\leqslant\pi_0$.}
\item[\textit{b$''$)}] \textit{(Si $\mathcal{M}$ est équivalente à une catégories} $\Ind_{\pi_0}C$\textit{, avec $C$ une petite catégorie} karoubienne.\textit{) La catégorie $C$ est stable (en elle-même) par limites inductives de cardinal $\leqslant\pi_0$.} % [sic] une catégories -> une catégorie, accord du nombre
\item[\textit{c)}] \textit{$\mathcal{M}$ est stable par petites limites projectives.}
\item[\textit{c$'$)}] \textit{(Sous l'hypothèse $\mathcal{M}\simeq\Ind_{\pi_0}(C)$, et quand on se donne un cardinal $\pi$ tel que $\pi=\pi^{\pi_0}\geqslant\pi_1=\sup(\pi_0,\card\textup{réd}\,C)$.) $\mathcal{M}_\pi$ est stable par $\varprojlim$ finies, et $\mathcal{M}$ est stable par petits produits infinis (ou encore, par $\varprojlim$ filtrantes).}
\item[\textit{c$''$)}] \textit{Il existe un cardinal infini $\pi$, et une petite catégorie $\mathcal{D}$ stable par $\varprojlim$ finies, telle que $\mathcal{M}\simeq\Ind_\pi(\mathcal{D})$, de plus $\mathcal{M}$ stable par produits infinis (ou encore par $\varprojlim$ filtrantes).}
\end{enumerate}
(\footnote{\textbf{N.B.} L'équivalence de c$'$, c$''$ avec les autres conditions n'est pas établie (ici, mais seulement dans 4.11 plus bas.)}, \footnote{\textbf{N.B.} Les critères c$'$, c$''$ me paraissent plus techniques et moins utiles \texttt{[c$'$, c$''$ ci-dessus sont mis en parenthèses]}.}).

\textsc{Démonstration.} Les implications suivantes sont tautologiques, ou déjà $\pm$ connues :
\[
\begin{tikzcd}[column sep=large, row sep=large]
\mathrm{a} \arrow[r, Rightarrow] & \mathrm{b} \arrow[r, Rightarrow, "\text{cf. 4.4.14}"] \arrow[dl, Rightarrow] & \mathrm{b}' \arrow[r, Rightarrow, "\bullet"] & \mathrm{b}'' \\
\mathrm{c}\ ( \arrow[r, Rightarrow, "\bullet"] & \mathrm{c}' \arrow[r, Rightarrow] & \mathrm{c}'')\ , &
\end{tikzcd}
\]
\msp{page 208}
où seules les implications pointées b$'\Longrightarrow$ b$''$ et c $\Longrightarrow$ c$'$ demandent une démonstration. Pour b$'\Longrightarrow$ b$''$, on note que l'hypothèse $C$ karoubienne implique que $C\simeq\mathcal{M}_{\pi_0}$ (4.8.5, $2^\circ$), or on sait que $\mathcal{M}_{\pi_0}$ est stable dans $\mathcal{M}$ par limites inductives de cardinal $\leqslant\pi_0$ (et celles-ci sont représentables dans $\mathcal{M}$ en vertu de b$'$), donc $\mathcal{M}_{\pi_0}$ admet des $\varinjlim$ du type envisagé, donc $C$ aussi.

c $\Longrightarrow$ c$'$. Notons que si $\mathcal{M}=\Ind_{\pi_0}(C)$, alors les $\varprojlim$ représentables dans $\mathcal{M}$ se calculent dans $C^\wedge$. Il s'ensuit que si $\pi\geqslant\pi_1$, alors $\Filt^\pi$ est stable par $\varprojlim$ finies ($\Filt^\pi$ étant donné alors par la condition a$''$ de 4.7.1), d'autre part, si de plus $\pi^{\pi_0}=\pi$, $\pi$ est un cardinal utile, donc $\mathcal{M}\simeq\Ind_\pi(\mathcal{D})$, où $\mathcal{D}=\Filt^\pi=\mathcal{M}_\pi$. Il suffit maintenant de voir que $\Ind_\pi(\mathcal{D})$ est alors stable par $\varprojlim$ finies, pour conclure qu'il en est de même de $\mathcal{M}$, donc que $\mathcal{M}$ est stable par limites projectives quelconques, s'il l'est par petites produites ou par petites $\varprojlim$ filtrantes, d'où c$'$. Il reste donc à prouver le lemme suivant. % [sic] petites produites -> petits produits, accord


\endgroup
% END UNIT der-ch18/der-ch18_p071-075.tex
% BEGIN UNIT der-ch18/der-ch18_p076-080.tex
\begingroup\makeatletter%
\newcommand{\printedpage}[1]{\marginpar{\footnotesize p.~#1}}%
\newcommand{\msp}[1]{\marginpar{\footnotesize [#1]}}%
\newcommand{\Ind}{\operatorname{Ind}}%
\newcommand{\Filt}{\operatorname{Filt}}%
\newcommand{\card}{\operatorname{card}}%
\newcommand{\red}{\operatorname{\text{réd}}}%
\newcommand{\Ob}{\operatorname{Ob}}%
\newcommand{\Hom}{\operatorname{Hom}}%
\newcommand{\uHom}{\underline{\operatorname{Hom}}}%
\newcommand{\End}{\operatorname{End}}%
\newcommand{\Ens}{\mathrm{Ens}}%
\newcommand{\cM}{\mathcal{M}}%
\newcommand{\cD}{\mathcal{D}}%
\newcommand{\defeq}{\overset{\text{déf}}{=}}%
\makeatother\let\printedpage\derpp\let\msp\dermsp%
\let\mspin\DERmspin%


\printedpage{74}
% head: A. Grothendieck : Dérivateurs, Ch. XVIII, Édition des Universités de Montpellier II et Paris VII

\textbf{Lemme 4.8.10.} \textit{Soient $C$ une petite catégorie, $\pi_0$ un cardinal infini, $\pi_1 = \sup(\pi_0, \card\red C)$. Supposons $\Ind_{\pi_0}$ stable par $\varprojlim$ finies (resp. et d'un certain type). Alors pour tout $\pi \geqslant \pi_1$, $\Filt^{\pi}$ est stable par lesdites $\varprojlim$ . (Donc aussi $\cM_{\pi} = \Filt^{\pi}$, si $\pi = \pi^{\pi_0}$.) Plus} % [sic] \Ind_{\pi_0} stable -> \Ind_{\pi_0}C stable, the argument C appears to be omitted (cf. the following sentence)
\msp{page 209}
\textit{généralement, soit $J$ une petite catégorie telle que $\Ind_{\pi_0}C$ soit stable par $\varprojlim$ de type $J$, et soit $c = \card\red J$, supposons que $\pi^{c} = \pi$ (p.ex. $c$ fini). Alors $\Filt^{\pi}$ est stable dans $\cM$ par $\varprojlim$ de type $J$. (Donc il en est de même de $\cM_{\pi} = \Filt^{\pi}$, si de plus $\pi^{\pi_0} = \pi$.)}

On utilise le critère a$''$ de 4.7.1, savoir $\card F(x) \leqslant \pi$ pour tout $x \in C$, pour l'appartenance à $\Filt^{\pi}$. On peut supposer $J$ réduite, donc $\card\Ob J \leqslant c$. Comme le foncteur d'inclusion $\Ind_{\pi_0}(C) \longrightarrow C^{\wedge}$ commute aux $\varprojlim$ quelconques, le calcul de $\varprojlim\limits_{j\in J} x_j$ dans $C^{\wedge}$ fibre par fibre montre que dans $C^{\wedge}$, on a
\[
X \ \defeq\ \varinjlim_{j\in J} x_j \ = \ \underbrace{\prod_{j\in\Ob J} x_j}_{\text{produit dans } C^{\wedge}} \quad ,
\] % [sic] \varinjlim_{j\in J} x_j -> \varprojlim_{j\in J} x_j, the text concerns the projective limit (cf. the preceding sentence and (4.8.10.1))
d'où
\[
\card X(x) \ \leqslant\ \prod_{j\in\Ob J} \underbrace{\card x_j(x)}_{\leqslant\pi} \ \leqslant\ \pi^{c}\ ,
\]
d'où (sans autre hypothèse sur $\pi$ que $\pi \geqslant \pi_1$)
\[
\varprojlim_{j\in J} x_j \ \in\ \Ob\Filt^{\pi^{c}}\ , \tag*{(4.8.10.1)}
\]
d'où la conclusion voulue si $\pi^{c} = \pi$, q.e.d.

\msp{page 210}

Nous allons prouver les implications supplémentaires
\[
\text{b}'' \ \Longrightarrow\ \text{b}\ , \qquad \text{c} \ \Longrightarrow\ \text{b}\ , \qquad \text{c}'' \ \overset{?}{\Longrightarrow}\ \text{c}\ ,
\]
ce qui achève la démonstration.

L'implication b$''$ $\Longrightarrow$ a est contenue dans 4.6.3.

Pour l'implication c $\Longrightarrow$ b, on note que la stabilité de $\cM$ par petites $\varinjlim$ équivaut à l'assertion que toute petite limite projective de foncteurs $\cM \xrightarrow{\ \xi_i\ } \Ens$ représentables, est représentable. Or les objets de $\cM$ étant accessibles ($\cM$ étant accessible), les foncteurs représentables $\xi_i$ sont accessibles, donc aussi leur $\varprojlim$ (SGA 4 I 9.7). D'autre part, le foncteur $\xi$ est exact à gauche, comme $\varprojlim$ de foncteurs représentables donc exactes à gauche. Donc le fait que $\xi$ est représentable résulte du % [sic] exactes -> exacts, agreement with « foncteurs »

\textbf{Lemme 4.8.11.} \textit{Soit $\cM$ une catégorie accessible, stable par $\varprojlim$ finies, $\cM \longrightarrow$} Ens \textit{un foncteur. Pour que $\xi$ soit représentable, il faut et il suffit qu'il soit exact à gauche et accessible. Donc (SGA 4 I 8.12.8.1) si $\cM$ est stable par petites $\varprojlim$ , $\xi$ est représentable si et seulement s'il commute auxdites $\varprojlim$ , et qu'il soit accessible.}

% foot: Version 17 décembre 2025

\printedpage{75}
% head: A. Grothendieck : Dérivateurs, Ch. XVIII, Édition des Universités de Montpellier II et Paris VII

\msp{page 211}

Le ``il faut'' est remis pour mémoire, on vient de le voir comme conséquence du fait que tout objet de $\cM$ est accessible. Le ``il suffit'' est essentiellement SGA 4 I 8.13.2.b), sauf que dans loc. cit. on suppose $\cM$ munie d'une fc \textit{stricte.} Mais la démonstration donnée dans loc. cit. montre aussi pour une fcga, en choisissant dans loc. cit. (p. 137) $\pi$ tel que $\xi$ soit $\pi$-accessible, et de plus $\pi$ soit utile pour la filtration cardinale généralisée envisagée.

Reste à prouver c$''$ $\overset{?}{\Longrightarrow}$ c, sous réserve que ce soit bien vrai. Ce serait une conséquence du lemme suivant (où $\cD$, $\pi$ sont remplacés par $C$, $\pi_0$).

\textbf{Lemme 4.8.12} (?) (\footnote[88]{Lemme non prouvé ici -- cf. démonstration 4.11.7 (non achevée).}). \textit{Soient $C_0$ une petite catégorie, $\pi_0$ un cardinal infini. Si $C$ est stable par $\varprojlim$ finies, alors $\Ind_{\pi_0}C$ également.} % [sic] C_0 -> C, the statement then speaks only of C

L'assertion similaire pour $\Ind(C)$, au lieu de $\Ind_{\pi}C$, est établie dans SGA 4 I 8.9.1 c). Mais il n'est pas clair sans plus que $\Ind_{\pi}C$ soit stable dans $\Ind(C)$ par $\varprojlim$ finies, dans le cas envisagé. De façon précise, si $F = \varprojlim\limits_{j\in J} F_j$, les $F_j$ dans $\Ind_{\pi}C$,
\msp{page 212}
$F$ dans $\Ind(C) \subset C^{\wedge}$, on sait que $C/F$ est filtrant et il faut voir qu'il est filtrante grande devant $\pi_0$. On est ramené aux trois cas où $J = \emptyset$, $J = e \amalg e$, et $J = \vcenter{\hbox{\begin{tikzcd}[ampersand replacement=\&, column sep=small, row sep=small]
\bullet \arrow[dr] \& \& \bullet \arrow[dl] \\
\& \bullet \&
\end{tikzcd}}}$. Les deux premiers cas sont essentiellement triviaux (car $e$, et un produit de deux catégories grandes devant $\pi_0$, sont grandes devant $\pi_0$), le troisième cas par contre n'a pas l'air trivial a priori. C'est lié visiblement à la question des représentations indicielles de foncteurs $J \longrightarrow \Ind_{\pi_0}(C)$, pour $\card\red J \leqslant \pi_0$ -- i.e. la validité de 4.4.13 b) en dehors de l'hypothèse que $C$ soit stable par $\varinjlim$ de cardinal $\leqslant \pi_0$. Il nous faudra y revenir très prochainement, car cette question m'apparaît à présent une question-clef, pour les propriétés de stabilité de la notion de catégorie accessible (en dehors d'aucune hypothèse d'existence de $\varinjlim$ ou $\varprojlim$ ). % [sic] qu'il est filtrante -> qu'il est filtrant, agreement with « il » (C/F)

% marginal handwritten mark (curved stroke) in the left margin beside Proposition 4.8.13
\textbf{Proposition 4.8.13} (\footnote[89]{\textit{Pas prouvé.} Mais OK si $\pi_0 > \card\red C$, cf. 4.15.10, p. 462.}). \textit{Soit $C$ une catégorie équivalente à une petite catégorie, et soit $\pi_0 \geqslant \card\red C$ un cardinal infini. Alors $\cM = \Ind_{\pi_0}C = \cM_{\pi_0}$ (i.e. tous les objets de $\Ind_{\pi_0}C$ sont $\pi_0$-accessibles), et $\cM$ est une enveloppe de Karoubi de $C$.}

\msp{page 213}

Signalons tout de suite les corollaires suivants.

% marginal handwritten mark (curved stroke) in the left margin beside Corollaire 4.8.14
\textbf{Corollaire 4.8.14} (\footnote[90]{\textit{Pas prouvé.} Si finalement, cf. 4.15.10, et 4.15.11, pages 462-463.}). \textit{Soit $\cM$ une $\mathfrak{U}$-catégorie. Conditions équivalentes :}
\begin{enumerate}[leftmargin=*]
\item[\textit{a)}] \textit{$\cM$ est accessible, et est équivalent à une petite catégorie.}
\item[\textit{a$'$)}] \textit{$\cM$ est accessible, et il existe un cardinal infini $\pi_0$ tel que tous les objets de $\cM$ soient $\pi_0$-accessibles (i.e. $\cM = \cM_{\pi_0}$).}
\item[\textit{b)}] \textit{$\cM$ est équivalente à une petite catégorie, et est karoubienne, i.e. stable par facteurs directs} (\footnote[91]{\textbf{N.B.} Ceci implique que $\cM$ et $\cM^{\circ}$ sont accessible, et je conjecture que la réciproque est vraie également.}). % [sic] sont accessible -> sont accessibles, plural agreement (footnote 91)
\end{enumerate}

L'équivalence de a) et a$'$) est claire, puisque $\cM$ admet une filtration cardinale généralisée canonique $(\cM_{\pi})_{\pi\geqslant\pi_0}$, formée des catégories équivalentes à des petites catégories (condition
% foot: Version 17 décembre 2025
\printedpage{76}
% head: A. Grothendieck : Dérivateurs, Ch. XVIII, Édition des Universités de Montpellier II et Paris VII
a) des fcg), et telle que toute telle sous-catégorie soit contenue dans un des $\cM_{\pi}$ (résulte de la condition b) des fcg). Comme $\cM_{\pi_0}$ est équivalente à une petite catégorie, et est karoubienne (étant stable par $\varinjlim$ filtrantes grandes devant $\pi_0$), il est clair aussi que a$'$) $\Longrightarrow$ b). Reste à prouver b) $\Longrightarrow$ a$'$), or cela résulte aussitôt de 4.8.13 appliqué à $\cM$, qui montre de plus qu'on peut prendre pour $\pi_0$ n'importe quel cardinal infini tel que $\pi_0 \geqslant \card\red\cM$.

\msp{page 214}

\textbf{4.8.15.} \textsc{Démonstration de} 4.8.13. Nous allons prouver :

\textbf{Lemme 4.8.15.1.} \textit{Sous les conditions de 4.8.13, soit $F \in \Ob\cM$, $x \in \Ob C$, alors}
\[
\card F(x) \ \leqslant\ \pi_0\ . \tag*{$(*)$}
\]

Ceci signifie donc, en vertu de 4.7.7, a$''$ $\Longrightarrow$ b pour $\pi \geqslant \pi_1$, compte tenu qu'ici $\pi_J = \pi_0$, $\pi = \pi_0$, qu'on a $F \in \Filt^{\pi_0}$, d'où a fortiori $F \in \cM_{\pi_0}$, ce qui implique $\cM = \cM_{\pi_0}$, i.e. la proposition 4.8.13, compte tenu de 4.8.5. Donc tout revient à prouver $(*)$.

Écrivons
\[
F \ = \ \varinjlim_{i\in I} x_i\ , \qquad\quad \text{les } x_i \in \Ob C,\ I \text{ ensemble ordonné grand devant } \pi_0,
\]
d'où
\[
F(x) \ = \ \varinjlim_{i\in I} \Hom_C(x, x_i)\ .
\]
On peut supposer
\[
\card C \ = \ \pi_0\ ,
\]
quitte à remplacer $C$ par une catégorie réduite équivalente.

Pour tout $\xi \in \Ob C$, soit
\[
I_{\xi} \ = \ \{ i \in I \mid x_i = \xi \}\ ,
\]
alors
\[
I \ = \ \bigcup_{\xi\in\Ob C} I_{\xi}\ ,
\]
\msp{page 215}
réunion d'une famille de parties de $I$ de cardinal $\card\Ob C \leqslant \pi_0$. Comme $I$ est grand devant $\pi_0$, on voit que cela implique que l'ensemble des $I_{\xi}$ est cofinal dans $I$. Quitte à remplacer $I$ par $I_{\xi}$, on peut supposer $I_{\xi} = I$, i.e. que tous les $x_i$ sont égaux à un même élément $\xi$. Quitte à remplacer $C$ par la sous-catégorie pleine $C_0 = \{\xi\}$, on peut supposer $C$ réduite à un seul élément. (Car si on prouve que $F$, en tant qu'objet de $\Ind_{\pi_0}(C_0)$, est facteur direct d'un objet de $C_0$ (i.e. de $\xi$), il en sera de même de $F$ en tant qu'objet de $\Ind C$). Donc on peut supposer
\[
C \ = \ \{\xi\}\ ,
\]
donc $C$ donné par le monoïde
\[
M \ = \ \End_C(\xi)\ , \qquad\quad \card M \ \leqslant\ \pi\ ,
\]
et la donnée du système inductif $(x_i)_{i\in I}$ revient à la donnée d'un système d'éléments
\[
\alpha_{ji} \ \in\ M \qquad\quad (i, j \in I,\ i \leqslant j)
\]
% foot: Version 17 décembre 2025
\printedpage{77}
% head: A. Grothendieck : Dérivateurs, Ch. XVIII, Édition des Universités de Montpellier II et Paris VII
satisfaisant
\[
\alpha_{ii} \ = \ 1\ , \qquad\quad \alpha_{kj}\alpha_{ji} \ = \ \alpha_{ki} \ \ \text{si } i \leqslant j \leqslant k.
\]
On a
\[
F(\xi) \ = \ \varinjlim_{i\in I} M_i \qquad\quad
\left\{
\begin{array}{l}
\text{où } M_i = M \ \forall\, i \in I, \text{ et } M_i \longrightarrow M_j \\
\text{est donné par} \\[2ex]
\qquad\qquad u \longmapsto \alpha_{ji}u\ . \\[4ex]
\end{array}
\right.
\]

\msp{page 216}

J'ai pu me ramener, par des arguments similaires (\footnote[92]{Pour ces réductions, cf. 4.15.8 (p. 453-463). L'ultime réduction se trouve dans 4.15.8.9 (p. 461), et son corollaire 4.15.10.}), au cas où le monoïde $M = \End(\xi)$ est un monoïde de monomorphismes, i.e.
\[
uv \ = \ uv' \ \ \Longrightarrow\ \ v \ = \ v' \qquad\quad \text{si } u, v, v' \in \cM\ ,
\] % [sic] u, v, v' \in \cM -> u, v, v' \in M, the elements belong to the monoid M
ce qui implique que les flèches de transition dans le système inductif des $M_i = M$ sont des monomorphismes. Il s'ensuit que l'on ne peut avoir $\card F(\xi) \leqslant \pi_0$ ($I$ étant grand devant $\pi_0$) que si $\alpha_{ji_0}$ est un isomorphisme pour $j \geqslant i_0$, et $i_0$ assez grand, i.e. si le système inductif induit pour $j \geqslant i_0$ est constant, a fortiori que le système inductif envisagé est essentiellement constant. Mais je viens de passer près de quatre heures pour essayer de prouver ce point, sans y être arrivé ! Cela semble plus délicat à prouver que je ne le prévoyais, mais je crois toujours que c'est bel et bien vrai.

\msp{page 217}

\subsection*{4.9\quad Le théorème de représentation indicielle pour $\Ind_{\pi}(C)$}

\textbf{4.9.1} (\footnote[93]{\texttt{[Les numéros 4.9.1 à 4.9.6 sont utilisés deux fois, 4.9.3 trois fois.]}}). Soient $C$ une catégorie équivalente à une petite catégorie, $\pi_0$ un cardinal infini,
\[
\cM \ = \ \Ind_{\pi_0}(C) \ \subset\ C^{\wedge}\ , \tag*{(4.9.1)}
\]
$J$ une catégorie équivalente à une petite catégorie. Considérons l'inclusion pleinement fidèle $C \hookrightarrow \cM$, et l'inclusion correspondante de $\mathfrak{U}$-catégories
\[
\uHom(J, C) \ \hookrightarrow\ \uHom(J, \cM) \qquad\quad \text{(pleinement fidèle).} \tag*{(4.9.2)}
\]
On sait, comme $C \subset \cM_{\pi_0}$, que les objets de $\uHom(J, C)$ donnent des objets $\pi_0$-\textit{accessibles} de $\uHom(J, \cM)$, du moins si on suppose
\[
\card\red J \ \leqslant\ \pi_0\ , \tag*{(4.9.3)}
\]
ce que nous ferons désormais. (Cf. 4.5.8, 2°, p. 119.) Comme le deuxième membre de (4.9.2) satisfait $\mathrm{L}_{\pi_0}$ ($\cM$ y satisfait), il s'ensuit que ce foncteur (4.9.2) se prolonge en
\[
\Ind_{\pi_0}(\uHom(J, C)) \ \hookrightarrow\ \uHom(J,\ \underbrace{\cM}_{=\,\Ind_{\pi_0}(C)}\ )\ , \tag*{(4.9.3)}
\]
% foot: Version 17 décembre 2025
\printedpage{78}
% head: A. Grothendieck : Dérivateurs, Ch. XVIII, Édition des Universités de Montpellier II et Paris VII
et ce foncteur est \textit{pleinement fidèle} grâce au résultat d'accessibilité qu'on vient d'énoncer. Je m'intéresse à établir des cas où ce foncteur est essentiellement surjectif, i.e. est une \textit{équivalence} de catégories. Notons que pour que
\[
\varphi \ :\ J \longrightarrow \cM
\]
appartienne à l'image essentielle, il faut et il suffit qu'il existe un foncteur
\[
\Phi \ :\ I \ \longrightarrow\ \uHom(J, C)\ ,\ I \text{ grande devant } \pi_0 \text{ (et on peut supposer } I \text{ ordonnée)}\ , \tag*{(4.9.4)}
\]

\msp{page 218}

tel qu'on ait
\[
\varphi \ \xleftarrow{\ \sim\ }\ \varinjlim_{i\in I} \Phi(i)\ , \tag*{(4.9.5)}
\]
ou ce qui revient au même, tel qu'on ait
\[
\varphi(j) \ \xleftarrow{\ \sim\ }\ \varinjlim_{i\in I} \Phi(i, j)\ , \qquad\quad \text{pour } j \in J\ , \tag*{(4.9.6)}
\]
où on a posé
\[
\Phi(i, j) \ = \ \Phi(i)(j)\ ,
\]
interprétant donc $\Phi$ comme un foncteur
\[
\Phi \ :\ I \times J \ \longrightarrow\ C\ ; \tag*{(4.9.7)}
\]
la loi fonctorielle de $\varphi(j)$ donnée par (4.9.6) étant évidente. En d'autres termes, les $\varphi$ dans l'image essentielle de (4.4.3) sont ceux qui admettent une \textit{représentation indicielle} par un foncteur $\Phi$ sous forme (4.9.4), ou sous forme (4.9.7), avec \textit{de plus} $I$ grand devant $\pi_0$. Je rappelle a ce sujet la % [sic] a ce sujet -> à ce sujet, missing accent

\textbf{Proposition 4.9.2.} \textit{Sous les conditions précédentes sur $C$, $\pi_0$, $J$ (notamment} $\card\red J \leqslant \pi_0$\textit{), si $\varphi : J \longrightarrow \cM \defeq \Ind_{\pi_0}(C)$, pour que $\varphi$ appartienne à l'image essentielle du foncteur pleinement fidèle (4.9.3), il faut et il suffit qu'elle satisfasse les deux conditions suivantes :}

\msp{page 219}

\begin{enumerate}[leftmargin=*]
\item[\textit{a)}] $\uHom(J, C)/\varphi$ \textit{est grande devant $\pi_0$.}
\item[\textit{b)}] \textit{Pour tout $j \in J$, le foncteur}
\[
\uHom(J, C)/\varphi \ \longrightarrow\ C/\varphi(j) \qquad\quad
\left\{
\parbox{0.35\textwidth}{\textit{foncteur entre catégories filtrantes grandes devant $\pi_0$, en vertu de a),}}
\right. \tag*{(4.9.2.1)}
\]
\textit{induit par le foncteur}
\[
\left\{
\begin{array}{rcl}
\uHom(J, C) & \longrightarrow & C \\
F & \longmapsto & F(j)\ ,
\end{array}
\right. \tag*{(4.9.2.2)}
\]
\textit{est} cofinal.
\end{enumerate}
% foot: Version 17 décembre 2025


\endgroup
% END UNIT der-ch18/der-ch18_p076-080.tex
% BEGIN UNIT der-ch18/der-ch18_p081-085.tex
\begingroup\makeatletter%
\newcommand{\printedpage}[1]{\marginpar{\footnotesize p.~#1}}%
\newcommand{\msp}[1]{\marginpar{\footnotesize [#1]}}%
\newcommand{\uHom}{\underline{\mathrm{Hom}}}%
\newcommand{\Ind}{\mathrm{Ind}}%
\newcommand{\Ob}{\operatorname{Ob}}%
\newcommand{\Fl}{\operatorname{Fl}}%
\newcommand{\card}{\operatorname{card}}%
\newcommand{\cardred}{\mathop{\textup{card\,réd}}\nolimits}%
\newcommand{\cM}{\mathcal{M}}%
\newcommand{\rbox}[1]{\tikz[baseline=(X.base)]\node[draw,rounded corners=4pt,inner sep=2pt](X){\textup{#1}};}%
\makeatother\let\printedpage\derpp\let\msp\dermsp%
\let\mspin\DERmspin%


\printedpage{79}
% head: A. Grothendieck : Dérivateurs, Ch. XVIII, Édition des Universités de Montpellier II et Paris VII

Le résultat principal dans la présente section est le suivant :

\textbf{Théorème 4.9.3.} \textit{Sous les conditions générales de 4.9.1 pour $C$, $\pi_0$, $J$ (notamment} $\cardred J \leqslant \pi_0$\textit{), le foncteur pleinement fidèle (4.9.3) est une équivalence de catégories dans chacun des deux cas suivants :}
\begin{itemize}[label={}]
{\itshape
\item[a)] $C$ est stable par $\varinjlim$ filtrantes dénombrables.
\item[b)] Il existe un cardinal $\pi_0 < \pi$, et un cardinal $\pi'_0 > \pi_0$ \textup{(}\textbf{\upshape N.B.} On peut supposer que $\pi'_0$ est le successeur de $\pi_0$, donc $\pi'_0 \leqslant \pi$), tels que $C$ soit stable par limites inductives grandes devant $\pi_0$, et de cardinal $\pi'_0$.
}
\end{itemize}
(\footnote[94]{Cas commun : il existe un ordinal limite $\alpha$, tel que 1) $\card \alpha \leqslant \pi_0$. 2) Les limites inductives \texttt{[de]} type $I_\alpha = \{\beta \mid \beta < \alpha\}$ sont représentables dans $C$.}).

\textbf{Remarque 4.9.4.}
\begin{itemize}[label={}]
\item[a)] Ce théorème, cas a), contient le théorème 4.4.13, partie b) (p. 95-96), où on suppose que $C$ est stable par limites inductives de cardinal $\leqslant \pi_0$, l'hypothèse a) de 4.9.3 étant considérablement plus faible.
\item[b)] \msp{page 220}Je présume qu'en dehors d'une hypothèse restrictive du type a) ou b) ci-dessus, le foncteur pleinement fidèle (4.9.3) n'est pas toujours essentiellement surjectif, même dans le cas particulier où $J$ est une catégorie finie, p.ex. $J = B_G$, $G = \{\pm 1\}$. Pour un contre-exemple dans le cas très voisin où on prend $\cM = \Ind(C)$ (au lieu de $\Ind_{\pi_0}(C)$), et où on considère le foncteur (pleinement fidèle si $J$ est une catégorie équivalent à une catégorie \textit{finie}) % [sic] équivalent -> équivalente, agreement with « une catégorie »
\[
\Ind(\uHom(J,C)) \longrightarrow \uHom(J,\Ind(C)),
\]
je renvoie à l'exercice SGA 4 I 8.8.8.
\end{itemize}
L'application principale du théorème 4.9.3 consiste en le

\textbf{Corollaire 4.9.5.} \textit{Soient $\cM$ une catégorie} accessible \textit{(4.8.7), $J$ une catégorie équivalente à une petite catégorie. Alors la catégorie $\cM(J^\circ) = \uHom(J,\cM)$ est accessible. Plus précisément, soient $\pi_0$ un cardinal} utile \textit{pour $\cM$ (i.e. tel que $\Ind_{\pi_0}(\cM_{\pi_0}) \xrightarrow{\sim} \cM$, i.e. un cardinal tel que $\cM$ soit $\pi_0$-accessible), soit $\pi$ un cardinal infini tel que $\pi^{\pi_0} = \pi \geqslant \cardred J$} (\footnote[95]{On a donc $\pi = \pi^{\pi_0} \geqslant 2^{\pi_0} > \pi_0$.}), \textit{i.e. $\pi = c^{\pi_0}$, où $c$ est n'importe quel cardinal $c \geqslant \sup(\pi_0, \cardred J)$, alors $\uHom(J,\cM)$ est $\pi$-accessible, plus précisément on a une équivalence de catégories} % doubt resolved against PDF; see corrections.csv
\[
\uHom(J,\cM) \xleftarrow{\ \sim\ } \Ind_\pi(\uHom(J,\cM_\pi)) \ . \tag*{(4.9.5.1)}
\]

\msp{page 221}On sait en effet, avec les hypothèses faites sur $\pi_0$, $\pi$, que $\pi$ est utile pour $\cM$, i.e. que
\[
\cM \xleftarrow{\ \sim\ } \Ind_\pi(\cM_\pi) \ ,
\]
d'autre part, $\cM_\pi$ est stable par $\varinjlim$ grandes devant $\pi_0$ et de cardinal $\leqslant \pi$. On peut donc appliquer 4.9.3 b), en prenant $C = \cM_\pi$, $\pi$ et $\pi_0$ comme ici (donc $\pi_0 < 2^{\pi_0} \leqslant \pi$), et $\pi'_0 = \pi$.

\printedpage{80}
% head: A. Grothendieck : Dérivateurs, Ch. XVIII, Édition des Universités de Montpellier II et Paris VII

\textbf{4.9.6} (\footnote[96]{Résultats préliminaires (à la démonstration de 4.9.3).}). Le reste de la présente section 4.9 est consacré à la démonstration de 4.9.3, i.e. à la vérification, pour $J \longrightarrow \cM$, des conditions a), b) de 4.9.2. Dans toute la suite, on se fixe donc une
\[
\varphi \ : \ J \longrightarrow \cM \ = \ \Ind_\pi(C) \ . \tag*{(4.9.6.1)}
\]
On peut supposer $J$ réduite, donc petite et $\card J \leqslant \pi_0$. On va considérer le foncteur déduit du foncteur (4.9.2.1)
\[
\uHom(J,C)/\varphi \longrightarrow \prod_{j \in \Ob J} C/\varphi(j) \ , \tag*{(4.9.6.2)}
\]
dont les foncteurs en question se déduisent en composant avec les foncteurs projection
\[
\prod C/\varphi(j) \longrightarrow C/\varphi(j_0) \ . \tag*{(4.9.6.3)}
\]
Comme les catégories facteurs $C/\varphi(j)$ sont filtrantes (et même grandes devant $\pi$), leur produit est également filtrant, et même grand devant $\pi$ (immédiat), et de même
\msp{page 222}
pour les produits partiels. Il s'ensuit aussitôt que les projections (4.9.6.3) sont cofinales, en vertu du

\textbf{Lemme 4.9.6.1.} \textit{Soient $I$, $I'$ deux catégories avec $I'$ \textup{0}-connexe (p.ex. filtrante), alors le foncteur projection $pr_1 : I \times I' \longrightarrow I$ est cofinal.}

En effet, pour $i' \in \Ob I'$ \texttt{[plutôt} $i \in \Ob I$\texttt{]}, on a
\[
i\backslash(I \times I') \ \simeq \ (i\backslash I) \times I' \ ,
\]
or $i\backslash I$ et $I'$ étant 0-connexe, il en est de même de leur produit.

\textbf{Lemme 4.9.6.2} (\footnote[97]{Attention, il n'est pas vrai qu'un produit infini de catégories 0-connexes est 0-connexe !}). \textit{Soit $(I_\alpha)_{\alpha \in A}$ une famille de catégories} filtrantes \textit{(resp. grandes devant $\pi$). Alors $I_A \stackrel{\text{déf}}{=} \prod_{\alpha \in A} I_\alpha$ est filtrante (resp. grande devant $\pi$), et pour toute partie $B \subset A$, la projection canonique $I_A \longrightarrow I_B$ est cofinal.} % [sic] cofinal -> cofinale, agreement with « la projection »

(\footnote[98]{Cf. démonstration élégante ci-dessous, 4.9.6.6.}). La stabilité de la propriété “filtrante” ou “grande devant $\pi$” par AQT \texttt{[âne qui trotte]}. On a, si $B' = A\backslash B$, $I_A \simeq I_B \times I_{B'} \longrightarrow I_B$, et on applique le lemme 4.9.6.1 pour conclure que ce foncteur est cofinal, $I_{B'}$ étant filtrante, donc 0-connexe.

\msp{page 223}Le foncteur (4.9.6.3) étant cofinal, pour établir que (4.9.2.1) l'est, il suffit de prouver qu'il en est ainsi du foncteur (4.9.6.2), ce que nous ferons. Notons de plus que ledit foncteur est \textit{fidèle}, car il s'insère dans le carré commutatif
\[
\begin{tikzcd}[column sep=large, row sep=large]
\uHom(J,C)/\varphi \arrow[r] \arrow[d] & \prod C/\varphi(j) \arrow[d] \\
\uHom(J,C) \arrow[r] & C^{\Ob J}\ ,
\end{tikzcd}
\]
\printedpage{81}
% head: A. Grothendieck : Dérivateurs, Ch. XVIII, Édition des Universités de Montpellier II et Paris VII
où les flèches verticales et la deuxième flèche horizontale sont fidèles. En définitive, pour établir le théorème 4.9.3, nous allons étudier le foncteur \textit{fidèle} (4.9.6.2)
\[
\begin{aligned}
&f \ : \ \uHom(J,C)/\Phi \ \longrightarrow \ \textstyle\prod_{j \in J_0} C/\varphi(j) \\ % [sic] \Phi -> \varphi, the functor is written \varphi everywhere else
&\text{(où on a posé } J_0 = \Ob J,\ J_1 = \Fl J) \ ,
\end{aligned}
\tag*{$(*)$}
\]
et nous allons prouver :
\begin{itemize}[label={}]
\item[\textcircled{a}] Ce foncteur fidèle $(*)$ est \textit{cofinal.}
\item[\textcircled{b}] La catégorie source $\uHom(J,C)/\varphi$ est grande devant $\pi_0$.
\end{itemize}
Nous allons pour ceci dégager un lemme auxiliaire sur la cofinalité des foncteurs \textit{fidèles} :

\msp{page 224}\textbf{Lemme 4.9.6.3} (\footnote[99]{Ce lemme fait partie du sorite général sur la cofinalité et les catégories filtrantes.}). \textit{Soit $f : I \longrightarrow I'$ un foncteur, et considérons les conditions :}
{\itshape
\begin{itemize}[label={}]
\item[a)] $I$ filtrante, $f$ cofinal, i.e. (en présence de $I$ filtrante) $f$ satisfait les conditions familières \rbox{F 1}, \rbox{F 2}.
\item[b)] $I'$ est filtrante, et les $i'\backslash I$ (pour $i' \in \Ob I'$) sont filtrantes.
\item[b$'$)] $I'$ \textup{0}-connexe, satisfait PF (2) \textup{(}\footnote[100]{p.ex. PF 2 est satisfaite si $I'$ est \textit{ordonnée}, sans avoir à supposer d'avance $I'$ filtrante, seulement 0-connexe.}\textup{)} (égalisation des double-flèches).
\item[b$''$)] $I'$ est filtrante, les $i'\backslash I$ pseudo-filtrantes, i.e. satisfont à la condition PF $1'$.
\item[c)] $I'$ filtrante, et $f$ satisfait à F 1, ainsi que la condition suivante.
\begin{itemize}[label={}]
\item[\rbox{F 3}] Pour tout couple $(i_0,i_1)$ d'objets de $I$ et toute flèche $f(i_0) \xrightarrow{u'} f(i_1)$, il existe un objet $i_2$ de $I$, et une flèche $v' : f(i_1) \longrightarrow f(i_2)$, tels que $v'$ et $v'u'$ soient relevables en des flèches $i_1 \xrightarrow{v} i_2$ et $i_0 \xrightarrow{w} i_2$ (i.e. il existe $i_2$ et $v : i_1 \longrightarrow i_2$, $w : i_0 \longrightarrow i_2$ tels qu'on ait $f(v)u' = f(w)$).
\end{itemize}
\end{itemize}
\[
\begin{tikzcd}
i_0 \arrow[rr, bend left=20, "w"] & i_1 \arrow[r, "v"] & i_2 % doubt resolved against PDF; see corrections.csv
\end{tikzcd}
\]
\[
\begin{tikzcd}
f(i_0) \arrow[r, "u'"] \arrow[rr, dotted, bend right=20, "f(w)"'] & f(i_1) \arrow[r, "f(v) = v'"] & f(i_2)
\end{tikzcd}
\]
\begin{itemize}[label={}]
\item[1)] On a les implications
\[
a \iff b \iff b' \implies b'' \iff c \ ,
\]
et si $f$ est \textup{fidèle} (p.ex. $I$ ordonnée), les cinq conditions a), b), b$'$), b$''$), c) sont équivalentes.
\end{itemize}
}

\printedpage{82}
% head: A. Grothendieck : Dérivateurs, Ch. XVIII, Édition des Universités de Montpellier II et Paris VII
{\itshape
\begin{itemize}[label={}]
\item[2)] \msp{page 225}Soit $\pi$ un cardinal infini, et considérons les conditions :
\begin{itemize}[label={}]
\item[$a_\pi$)] $I$ grande devant $\pi$, $f$ cofinal, i.e. F 1, F 2.
\item[$b_\pi$)] $I'$ grande devant $\pi$, les $i'\backslash I$ ($i' \in \Ob I'$) aussi.
\item[$b_\pi{}''$)] $I'$ grande devant $\pi$, les $i'\backslash I$ ($i' \in \Ob I'$) sont $\pi$-pseudo-filtrantes, i.e. satisfont PS $1_\pi{}'$. % [sic] PS 1_π' -> PF 1_π', the conditions are elsewhere called PF
\item[$c_\pi$)] $I'$ grande devant $\pi$, et $f$ satisfait F 1 et
\begin{itemize}[label={}]
\item[\rbox{F $3_\pi$}] \textup{(}\footnote[101]{\textbf{N.B.} La condition F $3_\pi$ est tautologiquement satisfaite si $f$ est \textit{pleinement fidèle.}}\textup{)}. Pour toute famille $(i_\alpha)_{\alpha \in A}$ d'objets de $I$, avec \textup{card} $A \leqslant \pi$, et $i \in \Ob I$, et une famille de flèches $u'_\alpha : f(i_\alpha) \longrightarrow f(i)$, existe $j \in \Ob I$ et une flèche $v' : f(i) \longrightarrow f(j)$, telle que $v'$ \textup{et} les $v'u'_\alpha$ proviennent de flèches de $I$, i.e. il existe des flèches $v : i \longrightarrow j$, $w_\alpha : i_\alpha \longrightarrow j$ telles que $f(v)u'_\alpha = f(w_\alpha)$ $\forall\ \alpha \in A$.
\end{itemize}
\end{itemize}
\[
\begin{tikzcd}[row sep=huge]
i_\alpha \arrow[d, dash] \arrow[rr, dotted, bend left=20, "w_\alpha"] & i \arrow[r, dotted, "v"] \arrow[d, dash] & j \arrow[d, dash] \\
f(i_\alpha) \arrow[r, "u'_\alpha"] & f(i) \arrow[r, "v' = f(v)"] & f(j)
\end{tikzcd}
\]
Ceci posé, on a
\[
a_\pi \iff b_\pi \implies c_\pi \implies b_\pi{}'' \ ,
\]
et si $f$ est fidèle, ou $I$ ordonné, les quatre conditions $a_\pi$, $b_\pi$, $b_\pi{}''$, $c_\pi$ sont équivalentes.
\end{itemize}
}
(\footnote[102]{Démonstration de la proposition donnée dans 4.9.6.7, page 227.}).

\textbf{4.9.6.4. N.B.} Dans le cas qui nous intéresse, où $f$ est le foncteur $(*)$ (p. 223), on sait que $I'$ est grande devant $\pi$. On voudrait établir que $I$ l'est et que $f$ est cofinal, i.e. la condition $\mathrm{a}_\pi$. On le prouvera en prenant $\mathrm{b}_\pi$ (qui implique $\mathrm{a}_\pi$ puisque $f$ est fidèle). Il reste donc à
\msp{page 226}
prouver, sur le foncteur $f$ dans $(*)$, p. 223, les deux conditions suivantes :
\begin{itemize}[label={}]
\item[F 1)] Pour tout système de flèches dans $\cM$
\[
X(j) \xrightarrow{p(j)} \varphi(j) \ , \qquad j \in J_0 = \Ob J, \text{ les } X(j) \text{ dans } \Ob C \ ,
\]
il existe un foncteur $F \ : \ J \longrightarrow C$ et $F \xrightarrow{q} \varphi$, enfin un système de flèches $X(j) \xrightarrow{u(j)} F(j)$ dans les $C/\varphi(j)$, i.e. rendant commutatif les diagrammes % [sic] commutatif -> commutatifs, agreement with « les diagrammes »
\[
\begin{tikzcd}[column sep=small]
X(j) \arrow[rr, "u(j)"] \arrow[dr, "p(j)"'] & & F(j) \arrow[dl, "q(j)"] \\
& \varphi(j)\ . &
\end{tikzcd}
\]
\item[F $3_\pi$)] Pour tout système $(F_\alpha)_{\alpha \in A}$ ($\card A \leqslant \pi$) d'éléments dans $\uHom(C,J)$, et des flèches % [sic] \uHom(C,J) -> \uHom(J,C), the F_α are functors J -> C
\[
F_\alpha \xrightarrow{p_\alpha} \varphi \qquad \text{dans } \uHom(J,\cM),
\]
et pour $G$ dans $\uHom(J,C)$ et $G \xrightarrow{q} \varphi$ dans $\uHom(J,\cM)$, étant donné un système de flèches $u_\alpha(j)$,
\printedpage{83}
% head: A. Grothendieck : Dérivateurs, Ch. XVIII, Édition des Universités de Montpellier II et Paris VII
\[
\begin{tikzcd}[column sep=small]
F_\alpha(j) \arrow[rr, "u_\alpha(j)", "\text{comm.}"'] \arrow[dr, "p_\alpha(j)"'] & & G(j) \arrow[dl, "q(j)"] \\
& \varphi(j)\ , &
\end{tikzcd}
\qquad u_\alpha(j) \text{ flèche dans } C/\varphi(j) \ ,
\]
il existe $H \xrightarrow{q} \varphi$ dans $\uHom(J,C)$ et une flèche $G \xrightarrow{v} H$ dans $\uHom(J,C)/\varphi$, telle que pour tout $\alpha \in A$, le système de flèches composés
\msp{page 227}
\[
\begin{tikzcd}[row sep=huge]
F_\alpha(j) \arrow[r, "u_\alpha(j)", "\text{comm.}"'] \arrow[dr, dotted, "p_\alpha(j)"'] & G(j) \arrow[r, "v(j)", "\text{comm.}"'] \arrow[d, dotted, "q(j)"] & H(j) \arrow[dl, dotted, "q(j)"] \\
& \varphi(j) &
\end{tikzcd}
\]
(pour $j \in J_0 \in \Ob J$) provienne d'un \textit{foncteur} $w_\alpha : F_\alpha \longrightarrow H$. % [sic] J_0 \in \Ob J -> J_0 = \Ob J, J_0 was defined as Ob J
\end{itemize}
Pour prouver 4.9.3, nous sommes donc ramené à prouver les deux propositions F 1), F $3_\pi$) précédentes, sous l'une des deux conditions a) ou b) de 4.9.3.

\textbf{Corollaire 4.9.6.5} (de 4.9.6.3). \textit{Soit $(f_\alpha : I_\alpha \longrightarrow I'_\alpha)$ une famille de foncteurs cofinaux, les catégories source étant filtrantes. Alors le foncteur produit $f = \prod f_\alpha : I = \prod_\alpha I_\alpha \longrightarrow I' = \prod_\alpha I'_\alpha$ est cofinal.}

En effet, il faut vérifier que les catégories $i'\backslash I$, pour $i' \in \Ob I'$, sont 0-connexes. Or si $i' = (i'_\alpha)_{\alpha \in A}$, on a
\[
i'\backslash I \ \simeq \ \prod_\alpha (i'_\alpha\backslash I_\alpha) \ ,
\]
et par 4.9.6.3, 1°) a $\Longrightarrow$ c, les $i'_\alpha\backslash I_\alpha$ sont filtrantes, donc leur produit est filtrant
\msp{page 228}
(cf. 4.4.6.2), donc 0-connexe, q.e.d. % [sic] 4.4.6.2 -> 4.9.6.2, the lemma on products of filtrantes is 4.9.6.2

\textbf{Corollaire 4.4.6.6.} \textit{Soit $\pi$ un cardinal infini. Tout produit de catégories grandes devant $\pi$ est grand devant $\pi$.} % [sic] 4.4.6.6 -> 4.9.6.6, numbering of the section and footnote 98

Soit $I = \prod_{\alpha \in A} I_\alpha$, soit $B$ un ensemble de cardinal $\leqslant \pi$, il faut prouver que $I \xrightarrow{\mathrm{diag}_I} I^B$ est cofinal. Or ce foncteur s'identifie au produit des foncteurs similaires $I_\alpha \longrightarrow I_\alpha^B$, lesquels sont cofinaux, les catégories source étant de plus filtrantes. Donc on peut conclure par le corollaire 4.9.6.5.

\textbf{4.9.6.7.} \textsc{Démonstration de} 4.9.6.3. (Partie 1°.)

a $\Longrightarrow$ b. On sait déjà que a) implique que $I'$ est filtrante (SGA 4 I 8.1.3 b). Reste à prouver que pour tout $i' \in \Ob I'$, $i'\backslash I$ est filtrante (\footnote[103]{Court-circuit : on sait déjà que $i'\backslash I$ est 0-connexe, d'autre part elle est \textit{localement} filtrante, car coétale sur $I$ qui est filtrante, donc elle est filtrante.}). Deux conditions PF $1'$ et PF 2 à vérifier (en plus du fait que $i'\backslash I \neq \emptyset$, qui est connu par cofinalité de $f$, impliquant que $i'\backslash I$ est 0-connexe).
\begin{itemize}[label={}]
\item[1)] Soient $(i,u')$, $(j,v')$ deux éléments de $i'\backslash I$, il faut voir qu'ils sont majorés par un troisième élément. Or $I$ étant filtrante, $i$, $j$ sont majorés par $k$. On aurait fini donc si on savait que $f(u_1)u' = f(v_1)v'$, soit $w'$, car alors $(k,w')$ ferait l'affaire.
\end{itemize}


\endgroup
% END UNIT der-ch18/der-ch18_p081-085.tex
% BEGIN UNIT der-ch18/der-ch18_p086-090.tex
\begingroup\makeatletter%
\newcommand{\printedpage}[1]{\marginpar{\footnotesize p.~#1}}%
\newcommand{\msp}[1]{\marginpar{\footnotesize [#1]}}%
\newcommand{\bs}{\backslash}%
\newcommand{\encadre}[1]{\tikz[baseline=(X.base)]\node[draw,rounded corners,inner sep=2pt](X){#1};}%
\makeatother\let\printedpage\derpp\let\msp\dermsp%
\let\mspin\DERmspin%


\printedpage{84}
% head: A. Grothendieck : Dérivateurs, Ch. XVIII, Édition des Universités de Montpellier II et Paris VII

\[
\begin{tikzcd}
{} & i \arrow[drr, "u_1"] \arrow[dd, dash] & {} & {} \\
{} & {} & j \arrow[r, "v_1"] & k \arrow[dd, dash] \\
i' \arrow[r, "u'"] \arrow[drr, "v'"'] & f(i) \arrow[drr] & {} & {} \\
{} & {} & f(j) \arrow[r] \arrow[uu, dash, crossing over] & f(k)
\end{tikzcd}
\]

\begin{itemize}[label={}]
\item \msp{page 229}Mais on applique F 2 sur $f$ à la double-flèche en question $i' \rightrightarrows f(k)$, on trouve $w : k \longrightarrow \bar{k}$,
\[
\begin{tikzcd}
{} & k \arrow[r, "w"] \arrow[dd, dash] & \bar{k} \arrow[dd, dash] \\
{} & {} & {} \\
i' \arrow[r, shift left] \arrow[r, shift right] & f(k) \arrow[r] & f(\bar{k})
\end{tikzcd}
\]
telle que
\[
\underbrace{f(w)(f(u_1)u')}_{= f(wu_1)u'} = \underbrace{f(w)(f(v_1)v')}_{f(wv_1)v'} \ ,
\]
et remplaçant $k$ par $\bar{k}$, et $u_1$, $v_1$ par $wu_1$, $wv_1$, on trouve ce qu'on veut.
\end{itemize}

\begin{enumerate}[label=\arabic*), start=2]
\item PF (2) pour $i'\bs I$, i.e. on se donne un diagramme ci-contre \texttt{[ci-dessous]},
\[
\begin{tikzcd}
{} & i \arrow[r, shift left, "u_1"] \arrow[r, shift right, "u_2"'] \arrow[dd, dash] & j \arrow[r, dotted, "w"] \arrow[dd, dash] & j' \\
{} & {} & {} & {} \\
i' \arrow[r, "u'"] & f(i) \arrow[r, shift left, "f(u_1)"] \arrow[r, shift right, "f(u_2)"'] & f(j) \arrow[r, dotted, "f(w)"] & f(j') \ ,
\end{tikzcd}
\]
avec $f(u_1)u' = f(u_2)u'$, soit $v'$ (double-flèche dans $i'\bs I$ de $(i,u')$ dans $(j,v')$), on veut l'égaliser à droite. Or $I$ étant filtrant, il existe $w : j \longrightarrow j'$ telle que $wu_1 = wu_2$, or cette flèche définit une flèche $(j,v') \longrightarrow (j', f(w')v')$, qui égalise la double-flèche donnée. % [sic] f(w')v' -> f(w)v', the arrow just defined is w : j --> j'
\end{enumerate}

b $\Longrightarrow$ b$'$ est clair.

b) $\Longrightarrow$ a). Comme les $i'\bs I$ sont filtrantes, donc 0-connexes, $f$ est cofinal, il reste à prouver que $I$ est filtrant, donc à vérifier PF 0 ($I$ 0-connexe), PF 1, PF 2 pour $I$. $I$ est 0-connexe, car $I'$ l'est, et $f$ cofinal, i.e. $\mathbf{W}_0$-coasphérique, d'où PF 0.

\msp{page 230}Prouvons PF 1, i.e. que si
\begin{tikzcd}[row sep=tiny, column sep=small, baseline=(\tikzcdmatrixname-2-1.base)]
{} & i_1 \\
i_0 \arrow[ur] \arrow[dr] & {} \\
{} & i_2
\end{tikzcd}
est donné, on peut compléter en un carré commutatif.
\[
\begin{tikzcd}[row sep=small, column sep=small]
{} & i_1 \arrow[dr, dotted, "v_1"] & {} \\
i_0 \arrow[ur, "u_1"] \arrow[dr, "u_2"'] & {} & i_3 \\
{} & i_2 \arrow[ur, dotted, "v_2"'] & {}
\end{tikzcd}
\]
% foot: Version 17 décembre 2025

\printedpage{85}
% head: A. Grothendieck : Dérivateurs, Ch. XVIII, Édition des Universités de Montpellier II et Paris VII

\[
\begin{tikzcd}[row sep=small, column sep=small]
{} & f(i_1) \arrow[dr, "f(v_1)"] & {} \\
f(i_0) \arrow[ur, "f(u_1)"] \arrow[dr, "f(u_2)"'] & {} & f(i_3) \\
{} & f(i_2) \arrow[ur, "f(v_2)"'] & {}
\end{tikzcd}
\]

Mais considérons les objets $\xi_1 = (i_1, f(u_1))$, $\xi_2 = (i_2, f(u_2))$ de $f(i_0)\bs I$, et utilisant PF 1$'$ pour cette catégorie, on trouve $\xi_3 = (i_3, w')$ dans ladite catégorie, et des flèches de $\xi_1$, $\xi_2$ dans $\xi_3$, i.e. $v_1 : i_1 \longrightarrow i_3$, \texttt{[}$v_2 : i_2 \longrightarrow i_3$\texttt{,]} telles que $\underbrace{f(v_1)f(u_1)}_{f(v_1u_1)} = w'$, $\underbrace{f(v_2)f(u_2)}_{f(v_2u_2)} = w'$, donc
\[
f(v_1u_1) \ = \ f(v_2u_2) \ .
\]
Ainsi, on a une double-flèche $i_0 \overset{v_1u_1}{\underset{v_2u_2}{\rightrightarrows}} i_3$, égalisée par $f$. Si on prouve PF 2 sur $I$, on trouve $i_3 \xrightarrow{\alpha} \bar{i}_3$ telle que $\alpha(v_1u_1) = \alpha(v_2u_2)$, i.e. $(\alpha v_1)u_1 = (\alpha v_2)u_2$, et quitte à remplacer $i_3$ par $\bar{i}_3$, $v_1$, $v_2$ par $\alpha v_1$, $\alpha v_2$, on a gagné.

Reste à prouver PF 2 sur $I$, i.e. égaliser une double-flèche
\[
i_0 \ \overset{v_1}{\underset{v_2}{\rightrightarrows}} \ i_1 \ .
\]

\msp{page 231}Or $I'$ satisfaisant à PF 2, on peut égaliser la double-flèche $f(u_1)$, $f(u_2)$ par une flèche $f(i_1) \xrightarrow{v'} i'$ dans $I'$, et comme $i'\bs I$ non vide, on peut supposer que $i'$ est de la forme $f(i_2)$.
\[
\begin{tikzcd}
i_0 \arrow[r, shift left, "u_1"] \arrow[r, shift right, "u_2"'] \arrow[dd, dash] & i_1 \arrow[r] & i_2 \\
{} & {} & {} \\
f(i_0) \arrow[r, shift left, "f(u_1)"] \arrow[r, shift right, "f(u_2)"'] & f(i_1) \arrow[r, "v'"] & f(i_2) % doubt resolved against PDF; see corrections.csv
\end{tikzcd}
\]
Admettons pour un instant la condition F 3 de c), que nous allons appliquer à $(i_1, i_2, v')$. Il existe donc $i_2 \xrightarrow{v} \bar{i}_2$, $w : i_1 \longrightarrow \bar{i}_2$ tels que l'on ait $f(w)v' = f(v)$. Il s'ensuit que $f(w)$ égalise $f(u_1)$ et $f(u_2)$, i.e. qu'on a $f(wu_1) = f(wu_2)$, i.e. la double-flèche initiale est remplacée par $(\bar{u}_1 = wu_1, \bar{u}_2 = wu_2)$, telle que $f(\bar{u}_1) = f(\bar{u}_2)$. En somme, on est ramené au cas d'une double-flèche $(u_1, u_2)$ telle que $f(u_1) = f(u_2)$, soit $v'$. Mais alors soit $\xi_1 = (i_0, v')$, $\xi_0 = (i_0, \mathrm{id}_{f(i_0)})$, on a donc une double-flèche $(u_1, u_2) : \xi_0 \rightrightarrows \xi_1$ dans $f(i_0)\bs I$. Cette catégorie % [sic] \xi_1 = (i_0, v') -> \xi_1 = (i_1, v'), u_1, u_2 go from i_0 to i_1
\msp{page 232}est filtrante, donc cette double-flèche est égalisable à droite, ce qui implique aussitôt que la double-flèche $(v_1, v_2)$ \textit{dans $I$} est aussi égalisable à droite.

Donc il reste à prouver que b$'$) implique la condition F 3. Mais considérons les deux objets $\xi_0 = (i_0, \mathrm{id}_{f(i_0)})$, $\xi_1 = (i_1, u')$ de $f(i_0)\bs I$, et utilisons la condition de filtration PF 1 sur cette catégorie, donc on trouve un objet $\xi_2 = (i_2, v')$, et deux flèches
\begin{tikzcd}[row sep=tiny, column sep=small, baseline=(\tikzcdmatrixname-2-2.base)]
\xi_0 \arrow[dr, "w"] & {} \\
{} & \xi_2 \\
\xi_1 \arrow[ur, "v"'] & {}
\end{tikzcd}
, i.e. $i_0 \xrightarrow{w} i_2$
% foot: Version 17 décembre 2025
\printedpage{86}
% head: A. Grothendieck : Dérivateurs, Ch. XVIII, Édition des Universités de Montpellier II et Paris VII
et $i_1 \xrightarrow{v} i_2$ telles qu'on ait $f(w)\mathrm{id}_{i_0} = f(v)u' = v'$, d'où $f(w) = f(v)u'$, ce qui est la condition F 3.
\[
\begin{tikzcd}
i_0 \arrow[rr, dotted, bend left, "w"] \arrow[dd, dash] & i_1 \arrow[r, dotted, "v"] \arrow[dd, dash] & i_2 \arrow[dd, dash] \\
{} & {} & {} \\
f(i_0) \arrow[r, "u'"] \arrow[rr, dotted, bend right, "v'"'] & f(i_1) & f(i_2)
\end{tikzcd}
\]

Ainsi on a établi
\[
\text{a} \ \Longleftrightarrow \ \text{b} \ \Longleftrightarrow \ \text{b}' \ ,
\]
prouvons maintenant

a $\Longrightarrow$ c. On sait déjà que a) implique $I'$ filtrante, et il reste à établir F 3. Or $I$ étant filtrante, il existe des flèches
\begin{tikzcd}[row sep=tiny, column sep=small, baseline=(\tikzcdmatrixname-2-2.base)]
i_0 \arrow[dr, "w"] & {} \\
{} & i_2 \\
i_1 \arrow[ur, "v"'] & {}
\end{tikzcd}
. On
\msp{page 233}n'a pas forcément $f(w) = f(v)u'$, on a là une double-flèche $f(i_0) = f(i_1)$ dans $I'$. Mais par la condition F 2 sur $f$, il existe $i_2 \xrightarrow{\alpha} \bar{i}_2$ dans $I$ telle que $f(\alpha)$ égalise ladite double-flèche, donc $f(\underbrace{\alpha w}_{\bar{w}}) = f(\underbrace{\alpha v}_{\bar{v}})u'$, et quitte à remplacer $i_2$ par $\bar{i}_2$, et $w$, $v$ par $\bar{w} = \alpha w$, $\bar{v} = \alpha v$, on gagne. % [sic] f(i_0) = f(i_1) -> f(i_0) \rightrightarrows f(i_2), the double arrow f(w), f(v)u' goes from f(i_0) to f(i_2)

Reste à prouver (pour la partie $1^\circ$) que si $f$ est \textit{fidèle}, alors c) $\Longrightarrow$ b). Il reste donc à prouver que c) implique que pour $i'$ dans $I'$, $i'\bs I$ est filtrante. On sait déjà par F 1 que $i'\bs I$ est non vide, il reste à prouver les conditions PF 1$'$, PF 2 pour $i'\bs I$. Comme $I'$ filtrante, l'on peut compléter
\[
\begin{tikzcd}[row sep=small, column sep=small]
{} & f(i_1) \arrow[dr, dotted] & {} \\
i' \arrow[ur] \arrow[dr] & {} & j' \ , \\
{} & f(i_2) \arrow[ur, dotted] & {}
\end{tikzcd}
\]
et par F 1 on peut supposer $j'$ de la forme $f(i_3)$.
% foot: Version 17 décembre 2025

\printedpage{87}
% head: A. Grothendieck : Dérivateurs, Ch. XVIII, Édition des Universités de Montpellier II et Paris VII

\[
\begin{tikzcd}[column sep=small]
{} & i_1 \arrow[drr, dotted, "v_1\,?"] \arrow[dd, dash] & {} & {} & {} \\
{} & {} & {} & i_3 \arrow[r] \arrow[dd, dash] & \bar{i}_3 \arrow[dd, dash] \\
{} & f(i_1) \arrow[drr, "v_1'"] & i_2 \arrow[ur, dotted, "v_2\,?"'] \arrow[dd, dash, crossing over] & {} & {} \\
i' \arrow[ur, "u_1'"] \arrow[dr, "u_2'"'] & {} & {} & f(i_3) \arrow[r] & f(\bar{i}_3) \\
{} & {} & f(i_2) \arrow[ur, "v_2'"'] & {} & {}
\end{tikzcd}
\]
Appliquant F 3 à $i_1$, $i_3$, $v_1'$, on peut supposer (quitte à remplacer $i_3$ par $\bar{i}_3$)
\msp{page 234}que $v_1'$ provient d'un $v_1 : i_1 \longrightarrow i_3$. Appliquant à nouveau F 3 à $i_2$, $i_3$, $v_2'$, on peut supposer que $v_2'$ également provient d'un $v_2 : i_2 \longrightarrow i_3$. Alors $\xi_3 = (i_3, v_1'u_1' = v_2'u_2')$ est un majorant commun de $\xi_1 = (i_1, u_1')$, $\xi_2 = (i_2, u_2')$.

Prouvons enfin PF 2 pour $i'\bs I$, donc on a un diagramme ci-contre \texttt{[ci-dessous]} avec $f(u_1)u' = f(u_2)u'$ (double-flèche dans $i'\bs I$), on veut égaliser dans $i'\bs I$.
\[
\begin{tikzcd}
{} & i_0 \arrow[r, shift left, "u_1"] \arrow[r, shift right, "u_2"'] & i_1 \arrow[rr, bend left, "w"] & i_2 \arrow[r] & \bar{i}_2 \\
{} & {} & {} & {} & {} \\
i' \arrow[r, "u'"] & f(i_0) \arrow[r, shift left, "f(u_1)"] \arrow[r, shift right, "f(u_2)"'] & f(i_1) \arrow[r, dotted, "v'"] & f(i_2) \arrow[r] & f(\bar{i}_2)
\end{tikzcd}
\]
Mais $I'$ étant filtrante, on peut égaliser $(f(u_1), f(u_2))$ par $f(i_1) \longrightarrow j'$, et à cause de F 1 on peut supposer $j' = f(i_2)$. Appliquant encore F 3 à $(i_1, i_2, v')$, on trouve $i_1 \longrightarrow \bar{i}_2$ tel que $f(w)$ égalise $f(u_1)$, $f(u_2)$ :
\[
f(wu_1) \ = \ f(wu_2) \ .
\]
Mais $f$ étant \textit{fidèle}, cela implique $wu_1 = wu_2$, et on a donc égalisé la double-flèche de $i'\bs I$ par $w : (i_1, -) \longrightarrow (\bar{i}_2, -)$.

\subsection*{4.9.6.8. N.B.}

\begin{enumerate}[label=\arabic*$^\circ$)]
\item Si on ne suppose pas que $f$ soit fidèle, on voit que c) implique pourtant que les $i'\bs I$ satisfont PF 1$'$ (elles sont donc pseudo-filtrantes et 0-connexes). Donc $f$ est cofinal.

\msp{page 235}On voit, par le même argument, que $I$ satisfait PF 1$'$, i.e. est pseudo-filtrante. Ainsi, si on introduit la condition :
\begin{itemize}[label=\encadre{b$''$}]
\item $I'$ est filtrante, les $i'\bs I$ ($i \in \mathrm{Ob}\, I'$) sont pseudo-filtrantes (i.e. satisfont PF 1$'$) (donc $f$ cofinal).
\end{itemize}
On trouve les implications
\[
(\text{a} \ \Longleftrightarrow \ \text{b} \ \Longleftrightarrow \ \text{b}') \ \ \Longrightarrow \ \ (\text{b}'' \ \Longleftrightarrow \ \text{c}) \ ,
\]
et les conditions équivalentes c), b$''$) impliquent que $I$ est également pseudo-filtrante, et que pour toute double-flèche $\xi_0 \overset{u_1}{\underset{u_2}{\rightrightarrows}} \xi_1$ dans un $i'\bs I$ ou dans $I$, il existe une flèche $\xi_1 \xrightarrow{v} \xi_2$ dans la même catégorie, telle que $vu_1$ et $vu_2 \ : \ \xi_0 \rightrightarrows \xi_2$ aient \textit{même image}
% foot: Version 17 décembre 2025
\printedpage{88}
% head: A. Grothendieck : Dérivateurs, Ch. XVIII, Édition des Universités de Montpellier II et Paris VII
\textit{dans $I'$}. Donc si $f$ est fidèle, alors ces conditions c, c$'$ \texttt{[plutôt b$''$]} impliquent que $I$ et les $i'\bs I$ satisfont aussi la condition PF 2 de la double-flèche, donc sont filtrantes, donc dans ce cas les cinq conditions a) b) b$'$) c) c$'$) \texttt{[plutôt b$''$)]} sont équivalentes. Mais aussi lorsque $I$ est ordonnée, on sait [\,?] que la condition de la double-flèche dans $I$ est satisfaite automatiquement.

\item \msp{page 236}Notons aussi que dans la condition b$'$), on ne peut supprimer l'hypothèse que $I'$ satisfait PF 2 (en plus de $I'$ 0-connexe et les $i'\bs I$ filtrantes). Pour le voir, prenons $I' = I$, $f = \mathrm{id}$, donc l'hypothèse a ($\Longleftrightarrow$ b $\Longleftrightarrow$ b$'$) signifie que $I$ est filtrante, $f$ étant bien sûr cofinal. Mais b$'$ affaiblie signifie que $I$ est 0-connexe et les $i\bs I$ sont filtrantes, et on sait que ceci n'implique \textit{pas} que $I$ soit filtrante, même si on suppose $I$ pseudo-filtrante. Un exemple simple est donné par $I$ définie par
\[
\mathrm{Ob}\, I \ = \ \mathbf{N} \ , \qquad \mathrm{Hom}_I(i,j) \ = \ \left\{
\begin{array}{rcccc}
\emptyset & \text{si} & i & > & j \\[1ex]
\{\mathrm{id}_i\} & \text{si} & i & = & j \\[1ex]
\{\pm 1\} & \text{si} & i & < & j \ ,
\end{array}
\right.
\]
la composition des flèches se faisant par la multiplication des signes $\varepsilon_{k,j}\varepsilon'_{j,i} = (\varepsilon\varepsilon')_{k,i}$ (en posant $\mathrm{id}_i = (+1)_{i,i}$). Les catégories $i\bs I$ sont isomorphes entre elles et sont \textit{ordonnées} et pseudo-filtrantes, donc filtrantes, d'autre part, $I$ est pseudo-filtrante, pourtant elle n'est pas filtrante, car ne satisfait pas PF 2.

\item D'autre part, si $f$ n'est pas fidèle, alors c) (équivalente à $b''$) n'implique pas a), comme on voit en prenant $I' = e$ ; alors b$''$) équivaut à $I$ pseudo-filtrante, qui n'implique pas que $I$ soit filtrante.
\end{enumerate}

\subsection*{\msp{page 237}\textbf{4.9.6.9.} \textsc{Démonstration de 4.9.6.3, partie} $2^\circ$.}

C'est essentiellement “la même” démonstration que pour la partie $1^\circ$).

a$_\pi$) $\Longrightarrow$ b$_\pi$). On peut court-circuiter la démonstration directe, en regardant, pour un ensemble $A$ de cardinal $\leqslant \pi$, le diagramme
\[
\begin{tikzcd}
I \arrow[r, "\delta_A"] \arrow[d, "f"'] & I^A \arrow[d, "f^A"] \\
I' \arrow[r, "\delta_{A'}"] & I'^A
\end{tikzcd}
\]
\texttt{[plutôt $\delta'_A$]}. Comme $f$ est cofinal, $I$ filtrante, $f^A$ est aussi cofinal (4.9.8.5), d'autre part, $I$ grande devant $\pi$ signifie que $\delta_A$ est cofinal ($\forall$ $A$ avec $\mathrm{card}\, A \leqslant \pi$). Donc $f^A\delta_A = \delta_{A'}f$ est cofinal, donc $\delta_{A'}f$ et $f$ le sont, donc $\delta_{A'}$ aussi, ce qui signifie que $I'$ est grande devant $\pi$.

Reste à prouver que pour $i'$ dans $I'$, $i'\bs I$ est grande devant $\pi$. Cela résulte du fait que $i'\bs I \xrightarrow{p} I$ est coétale, que $i\bs I$ est 0-connexe, et du

\textbf{Lemme 4.9.6.9.1.} \textit{Soit $I$ une catégorie filtrante (resp. grande devant $\pi$). Alors pour toute $J$ sur $I$ coétale sur $I$ et $0$-connexe, $J$ est filtrante (resp. grande devant $\pi$).}
% foot: Version 17 décembre 2025


\endgroup
% END UNIT der-ch18/der-ch18_p086-090.tex
% BEGIN UNIT der-ch18/der-ch18_p091-095.tex
\begingroup\makeatletter%
\newcommand{\printedpage}[1]{\marginpar{\footnotesize p.~#1}}%
\newcommand{\msp}[1]{\marginpar{\footnotesize [#1]}}%
\newcommand{\uGamma}{\underline{\Gamma}}%
\newcommand{\Ind}{\operatorname{Ind}}%
\newcommand{\Ob}{\operatorname{Ob}}%
\newcommand{\card}{\operatorname{card}}%
\newcommand{\trois}[1]{\mathrel{\overset{#1}{\vcenter{\offinterlineskip\hbox{$\longrightarrow$}\vskip0.5pt\hbox{$\longrightarrow$}\vskip0.5pt\hbox{$\longrightarrow$}}}}}%
\makeatother\let\printedpage\derpp\let\msp\dermsp%
\let\mspin\DERmspin%


\printedpage{89}
% head: A. Grothendieck : Dérivateurs, Ch. XVIII, Édition des Universités de Montpellier II et Paris VII
% foot: Version 17 décembre 2025

Le cas filtrant est déjà connu (cf. II \dots). Pour le cas “grande devant $\pi$”, on va utiliser le

\textbf{Corollaire 4.9.6.9.2.} \textit{Soit $J$ une catégorie. Pour que $J$ soit grande devant $\pi$, il}
\msp{page 238}
\textit{faut et il suffit qu'elle soit $0$-connexe, et satisfasse PF$_\pi$ 2 (cf. page 81), ainsi que la variante suivante de PF$_\pi$ 1$'$.}
\begin{itemize}[leftmargin=4em]
\item[\textit{PF$_\pi$ 1)}] \textit{Pour tout objet $x_0$ de $J$, $x_0\backslash J$ satisfait PF$_\pi$ 1$'$, i.e. pour toute famille de flèches $x_0 \xrightarrow{u_\alpha} y_\alpha$ ($\alpha \in A$, $\card A \leqslant \pi$), existe une flèche $x_0 \xrightarrow{v} z$ et des flèches $v_\alpha : y_\alpha \longrightarrow z$, avec $v_\alpha u_\alpha = v \ \forall\, \alpha \in A$.}
\end{itemize}

\textsc{Démonstration}. Si $J$ est grande devant $\pi$, alors PF$_\pi$ 1) résulte aussitôt de la conjonction de PF$_\pi$ 1$'$ et PF$_\pi$ 2, d'où le “il faut”. Inversement, supposons que $J$ satisfasse PF$_\pi$ 1, PF$_\pi$ 2, et soit 0-connexe, il faut prouver qu'elle satisfait PF$_\pi$ 1$'$. Mais on sait déjà que les conditions en question, qui impliquent PF 1, PF 2, 0-connexe, impliquent que $J$ est filtrante. Donc si on a une famille $(x_\alpha)_{\alpha\in A}$, $\card A \leqslant \pi$, choisissant un $x_0 \in \Ob J$, on peut pour tout $\alpha$ trouver $y_\alpha$ qui majore $x_0$ et $x_\alpha$, d'où des flèches
\begin{tikzcd}[row sep=0.4em, column sep=1em, baseline=(\tikzcdmatrixname-2-2.base)]
x_\alpha \arrow[dr] & \\
& y_\alpha \\
x_0 \arrow[ur] &
\end{tikzcd}, on peut leurs appliquer PF$_\pi$ 1, d'où des flèches $y_\alpha \longrightarrow z$, et en composant des $x_\alpha \longrightarrow z$, q.e.d. % [sic] leurs appliquer -> leur appliquer, grammar

Prouvons maintenant le lemme. Il suffit (par le corollaire précédent) de prouver que $J$ satisfait à PF$_\pi$ 1 et PF$_\pi$ 2. Or PF$_\pi$ 1 est immédiat, puisque $\forall\, x_0$ dans $J$, $x_0\backslash J \xrightarrow{\ \sim\ } y_0\backslash I$, où $y_0$ est
\msp{page 239}
l'image de $x_0$ dans $I$. Quant à PF$_\pi$ 2, si on a des $x_0 \trois{u_\alpha} x_1$ dans $J$, d'où des $g(x_0) \trois{g(u_\alpha)} g(x_1)$, on égalise dans $I$ par $g(x_1) \xrightarrow{w} y$, mais comme $g$ est coétale, $\exists\, x_1 \xrightarrow{v} x_2$ dans $J$ telle que $g(v) = w$, donc on a $g(vu_\alpha) = g(vu_\beta)$ pour $\alpha, \beta \in A$, et comme $g$ est fidèle, on en conclut $vu_\alpha = vu_\beta$, donc $v$ égalise $(u_\alpha)_{\alpha\in A}$, q.e.d.

Prouvons b$_\pi$) $\Longrightarrow$ a$_\pi$) : Si $I'$ et les $i'\backslash I$ sont grandes devant $\pi$, alors $I$ l'est aussi (et $f$ est cofinal, les $i'\backslash I$ étant 0-connexes, car filtrantes). Comme on sait déjà que $I$ est 0-connexe ($I'$ l'étant et $f$ cofinal), il reste, par 4.9.6.9.2, à vérifier PF$_\pi$ 1 et PF$_\pi$ 2.

\msp{page 240}
La démonstration donnée dans 1° pour le cas similaire filtrant (implication b$'$ $\Longrightarrow$ a, pages 231-232), s'applique mutatis mutandis. (La propriété F 3 a déjà été établie dans la partie 1°), comme conséquence de b, et peut donc être utilisée librement.)

Donc on a prouvé a$_\pi$ $\Longleftrightarrow$ b$_\pi$. Prouvons a$_\pi$ $\Longrightarrow$ c$_\pi$. On sait déjà que a$_\pi$ implique que $I'$ est grande devant $\pi$, et bien sûr la condition F 1, reste à établir la validité de la condition F 3$_\pi$. Par F 3, $\forall\, \alpha \in A$ on peut trouver $j_\alpha$ dans $I$ et $w^0_\alpha : i_\alpha \longrightarrow j_\alpha$, $v^0_\alpha : i \longrightarrow j$ tels que % [sic] v^0_\alpha : i \longrightarrow j -> v^0_\alpha : i \longrightarrow j_\alpha, the following diagram has v^0_\alpha : i \to j_\alpha
\[
f(v_\alpha)u'_\alpha \ = \ f(w_\alpha) \ .
\]

\printedpage{90}
% head: A. Grothendieck : Dérivateurs, Ch. XVIII, Édition des Universités de Montpellier II et Paris VII
% foot: Version 17 décembre 2025

\[
\begin{tikzcd}[column sep=large, row sep=huge]
i_\alpha \arrow[rr, bend left=20, "w^0_\alpha"] \arrow[d, dash] & i \arrow[r, "v^0_\alpha"] \arrow[d, dash] & j_\alpha \arrow[r, "\lambda_\alpha"] \arrow[d, dash] & j \arrow[d, dash] \\
f(i_\alpha) \arrow[r, "u'_\alpha"] & f(i) \arrow[r, "f(v_\alpha)"] & f(j_\alpha) \arrow[r] & f(j)
\end{tikzcd}
\]

Comme $I$ satisfait PF 1$_\pi$ (cf. 4.9.6.9.2), il existe $h$ dans $I$ et des $\lambda_\alpha : j_\alpha \longrightarrow j$, tels que $\lambda_\alpha v^0_\alpha = \lambda$ soit indépendant de $\alpha$, alors on a pour tout $\alpha \in A$ % [sic] h -> j, the arrows \lambda_\alpha go to j
\[
f(\lambda)u'_\alpha \ \Bigl( \ = \ f(\lambda_\alpha)\underbrace{f(v_\alpha)u'_\alpha}_{f(w_\alpha)} \ \Bigr) \ = \ f(\underbrace{\lambda_\alpha w^0_\alpha}_{w_\alpha}) \ ,
\]
donc $\lambda : i \longrightarrow j$ et les $w_\alpha = \lambda_\alpha w^0_\alpha : i_\alpha \longrightarrow j$ satisfont aux conditions voulues.

c$_\pi$ $\Longrightarrow$ b$''_\pi$. Il faut prouver que si c$_\pi$ est satisfait, alors les $i'\backslash I$ satisfont PF 1$'_\pi$.

\msp{page 241}
Soit donc $\xi_\alpha = (i_\alpha, u'_\alpha)$ ($\alpha \in A$, $\card A \leqslant \pi$) une famille d'éléments de $i'\backslash I$, comme $I'$ satisfait à PF 1$_\pi$ (cf. 4.9.6.9.2), il existe $j'$ dans $I'$ et des flèches $u''_\alpha : f(i_\alpha) \longrightarrow j'$, telles que $u''_\alpha u'_\alpha = u'$ indépendant de $\alpha$. Comme $f$ satisfait F 1, on peut supposer que $j'$ est de la forme $f(j)$. Appliquant F 3$_\pi$ aux $i_\alpha$, $j$, $u''_\alpha$, on trouve $v : j \longrightarrow k$ et des $w_\alpha : i_\alpha \longrightarrow k$, tels que
\[
f(v)u''_\alpha \ = \ f(w_\alpha) \qquad \forall\, \alpha \in A \ ,
\]
d'où
\[
f(w_\alpha)u'_\alpha \ = \ f(v)\,\underbrace{u''_\alpha u'_\alpha}_{u'} \ = \ f(v)u' \qquad \text{indépendant de $\alpha$,}
\]
soit $w'$, donc les $w_\alpha$ définissent une flèche des $\xi_\alpha$ dans $\xi = (k, w')$, q.e.d.

b$''_\pi$ $\Longrightarrow$ b$_\pi$ si $f$ est fidèle. On est ramené à prouver que si b$''_\pi$ est satisfait, $f$ fidèle, alors les catégories $i'\backslash I$, dont on sait déjà qu'elles satisfont PF 1$_\pi$, satisfont également PF 2$_\pi$. Soient donc $\xi = (i, u')$, $\eta = (j, v')$ deux objets de $i'\backslash I$, et $(u_\alpha)_{\alpha\in A}$ une famille de flèches $\xi \longrightarrow \eta$ (avec $\card A \leqslant \pi$), donc des flèches dans $I$ telles que $f(u_\alpha)u' = v' \ \forall\, \alpha \in A$.
\[
\begin{tikzcd}[column sep=large]
i \arrow[r, shift left=2, "u_\alpha"] \arrow[r] \arrow[r, shift right=2] & j \arrow[r, dotted, "v"] & k
\end{tikzcd}
\]
\[
\begin{tikzcd}[column sep=large]
i' \arrow[r, "u'"] \arrow[rr, bend right=15, "v'"'] & f(i) \arrow[r, shift left=2] \arrow[r] \arrow[r, shift right=2] & f(j) \arrow[r] & f(k)
\end{tikzcd}
\]

Comme $I'$ est grande devant $\pi$, on peut égaliser par $f(j) \longrightarrow k'$, et comme $k'\backslash I \neq \emptyset$, on peut supposer $k'$ de la forme $f(k)$. Or
\msp{page 242}
d'après l'implication b$''$ $\Longrightarrow$ a $\Longrightarrow$ c de la partie 1°) (cf. page 235) et comme b$''_\pi$ $\Longrightarrow$ b$''$, on sait déjà que $f$ satisfait à F 3. En l'appliquant à $j$, $k$ et $f(j) \longrightarrow f(k)$, on peut
\printedpage{91}
% head: A. Grothendieck : Dérivateurs, Ch. XVIII, Édition des Universités de Montpellier II et Paris VII
% foot: Version 17 décembre 2025
supposer que $f(j) \longrightarrow f(k)$ provient de $v : j \longrightarrow k$. On aura donc $f(v)f(u_\alpha) = f(v)f(u_\beta)$ $\forall\, \alpha, \beta \in A$, donc $f(vu_\alpha) = f(vu_\beta)$. Comme $f$ est fidèle, cela implique que $vu_\alpha = vu_\beta$, i.e. $v$ égalise les $u_\alpha$. D'ailleurs, $v$ peut être considérée comme une flèche \textit{dans} $i'\backslash I$, de $(j, v')$ dans $(k, f(v)v')$. Q.e.d.

\textbf{4.9.6.10.} Quand $I$ est \textit{ordonné}, alors les quatre conditions a$_\pi$, b$_\pi$, c$_\pi$, $b''_\pi$ sont équivalentes, car les $i'\backslash I$ sont ordonnés, et il est claire qu'alors b$''_\pi$ $\Longleftrightarrow$ b$_\pi$. % [sic] claire -> clair, grammar

\msp{page 243}

\section*{4.10\quad Le théorème de représentation indicielle : démonstration}

On va d'abord reformuler le théorème 4.9.3 sous forme plus générale, en termes de catégories \textit{cofibrées}. Soit
\[
\mathcal{D} \longrightarrow J \tag*{(4.10.1)}
\]
un foncteur cofibrant, où $J$ est équivalent à une petite catégorie, et de même pour les catégories fibres $\mathcal{D}_j$ ($j \in \Ob J$). On se donne un cardinal infini $\pi$. Soit
\[
\mathcal{M}_j \ = \ \Ind_\pi(\mathcal{D}_j) \qquad (j \in \Ob J) \ , \tag*{(4.10.2)}
\]
alors les $\mathcal{M}_j$ sont les fibres d'une catégorie cofibrée $\mathcal{N}$ sur $J$, que je désigne aussi par $\Ind_\pi(\mathcal{D}/J)$, donc on pose
\[
\mathcal{N} \ = \ \Ind_\pi(\mathcal{D}/J) \ . \tag*{(4.10.3)}
\]
On a donc une inclusion pleinement fidèle de catégories cofibrées sur $J$
\[
\mathcal{D} \hookrightarrow \mathcal{N} \ ,
\]
d'où une inclusion
\[
\uGamma(\mathcal{D}/J) \longrightarrow \uGamma(\mathcal{N}/J) \ . \tag*{(4.10.4)}
\]

\msp{page 244}
Dans le cas où $\mathcal{D}$ est de la forme $C \times J$, alors $\mathcal{N}$ est de la forme $\mathcal{M} \times J$, avec $\mathcal{M} = \Ind_\pi(C)$, et l'inclusion (4.10.3) n'est autre (à isomorphisme près) que
\[
\underline{\operatorname{Hom}}(J, C) \longrightarrow \underline{\operatorname{Hom}}(J, \mathcal{M}) \ .
\]
Revenant au cas général, on voit encore que :
\begin{enumerate}[label={}, leftmargin=3em]
\item[1°)] Le deuxième membre de (4.10.3) est stable par $\varinjlim$ grandes devant $\pi$, lesquelles se calculent fibre à fibre.
\item[2°)] Le premier membre de (4.10.3) est formé d'objets $\pi$-accessibles du second.
\end{enumerate}
D'où un foncteur \textit{pleinement} fidèle
\[
\Ind_\pi(\uGamma(\mathcal{D}/J)) \hookrightarrow \uGamma(\mathcal{N}/J) \ . \tag*{(4.10.4)}
\]
\setcounter{footnote}{103}%
(\footnote{\texttt{[Les numéros 4.10.1 à 4.10.4 sont utilisés deux fois.]}}).

\printedpage{92}
% head: A. Grothendieck : Dérivateurs, Ch. XVIII, Édition des Universités de Montpellier II et Paris VII
% foot: Version 17 décembre 2025

\textbf{Théorème 4.10.1.} \textit{Supposons qu'on ait}
\[
\mathrm{card}\,\textup{réd}\, J \ \leqslant \ \pi \ , \tag*{(4.10.1.1)}
\]
\textit{et qu'il existe un ordinal limite $\alpha$, tel que :}
\begin{enumerate}[label={}, leftmargin=3em]
\item[\textit{1)}] $\card \alpha \leqslant \pi$.
\item[\textit{2)}] \textit{Les limites inductives de}
\msp{page 245}
\textit{type $I_\alpha = \{\beta \mid \beta < \alpha\}$ sont représentables dans les catégories $\mathcal{D}_j$ ($j \in \Ob J$), et pour toute flèche $\alpha : j \longrightarrow j'$ dans $J$ \texttt{[trop de $\alpha$]}, $\alpha^* : \mathcal{D}_j \longrightarrow \mathcal{D}_{j'}$ commute auxdites limites.}
\end{enumerate}
\textit{Alors le foncteur canonique pleinement fidèle (4.10.4) est une équivalence de catégories.}

Ceci va être contenu dans le résultat un peu plus précis :

\textbf{Corollaire 4.10.2.} \textit{Les hypothèses étant celles de 4.10.1, avec en plus}
\[
\card J \ \leqslant \ \pi \tag*{(4.10.2.1)}
\]
\textit{(qui équivaut à (4.10.1.1) si $J$ est réduite, cas auquel on peut évidemment se réduire dans 4.10.1). Soit $\varphi \in \Ob \uGamma(\mathcal{N}/J)$, et considérons le foncteur canonique} fidèle \textit{:}
\[
\uGamma(\mathcal{D}/J)/\varphi \longrightarrow \prod_{j\in J} C_j/\varphi(j) \ . \tag*{(4.10.2.2)}
\]
\textit{Alors la catégorie source $\uGamma(\mathcal{D}/J)/\varphi$ est grande devant $\pi$, et le foncteur envisagé est cofinal.}

\msp{page 246}
De ceci, le théorème résulte de façon essentiellement soritale, tout comme dans le cas particulier 4.9.3 précédent (cf. page 223) (\footnote{Un peu court\,! Il faudrait expliciter l'argument avec soin – il y en a pour deux ou trois pages\,!}).

Pour établir le corollaire, comme la catégorie but de (4.10.2.2) est grande devant $\pi$, comme produit de telles catégories (cf. 4.9.6.6, page 228), et le foncteur fidèle, il suffit en vertu de 4.9.6.3, 2°, de vérifier les conditions F 1 et F 3$_\pi$ (page 225) pour le foncteur envisagé.

On va interpréter le foncteur (4.10.2.2) en termes d'une catégorie cofibrée
\[
\mathcal{D}_\varphi \longrightarrow J \ , \tag*{(4.10.5)}
\]
dont les fibres sont les $\mathcal{D}_j/\varphi(j)$. On note que si $\alpha : j \longrightarrow j'$ est une flèche de $J$, alors on a
\[
\left\{
\begin{aligned}
&\alpha_* : C_j \ \longrightarrow \ C_{j'} \\
&\text{et} \ \ \alpha_*(\varphi(j)) \ \longrightarrow \ \varphi(j') \ ,
\end{aligned}
\right.
\]

\printedpage{93}
% head: A. Grothendieck : Dérivateurs, Ch. XVIII, Édition des Universités de Montpellier II et Paris VII
% foot: Version 17 décembre 2025

\msp{page 247}
d'où un foncteur composé
\[
C_j/\varphi(j) \ \longrightarrow \ C_{j'}/\alpha^*(\varphi(j)) \ \longrightarrow \ C_{j'}/\varphi(j') \ ,
\]
avec transitivité à isomorphisme canonique près etc. On note que l'on a un isomorphisme canonique
\[
\uGamma(\mathcal{D}_\varphi/J) \ \simeq \ \uGamma(\mathcal{D}/J)/\varphi \ ,
\]
de sorte que le foncteur fidèle canonique (4.10.2.2) s'identifie au foncteur
\[
\uGamma(\mathcal{D}_\varphi/J) \longrightarrow \prod_{j\in\Ob J} (\mathcal{D}_\varphi)_j \ . \tag*{(4.10.6)}
\]

Le corollaire 4.10.2 sera alors une conséquence du résultat plus général suivant :

\textbf{Corollaire 4.10.3.} \textit{Soit $\mathcal{E} \longrightarrow J$ un foncteur cofibrant, et $\pi$ un cardinal infini. On suppose :}
\begin{enumerate}[label={}, leftmargin=3em]
\item[\textit{a)}] $\card J \leqslant \pi$.
\item[\textit{b)}] \textit{Les catégories fibres $\mathcal{E}_j$ sont équivalentes à des petites catégories, et grandes devant $\pi$.}
\item[\textit{c)}] \textit{Il existe un ordinal limite tel que :}
\begin{enumerate}[label={}, leftmargin=3em]
\item[\textit{1°)}] $\card \alpha \leqslant \pi$.
\item[\textit{2°)}] \msp{page 248}\textit{Les limites inductives de type $I_\alpha$ sont représentables dans les catégories $\mathcal{E}_j$, et les foncteurs cochangement de base $\lambda_* : \mathcal{E}_j \longrightarrow \mathcal{E}_{j'}$ y commutent. (C'est le cas p.ex. si les limites inductives filtrantes} dénombrables \textit{sont représentables dans les $\mathcal{E}_j$, et si les $\alpha_*$ y commutent\,; ou s'il existe un cardinal $\pi_0 \leqslant \pi$ et un cardinal $\pi'_0 > \pi_0$, tel que les limites inductives grandes devant $\pi_0$ et de cardinal $\leqslant \pi'_0$ soient représentables dans les $\mathcal{E}_j$, et que les $\alpha_*$ y commutent.)}
\end{enumerate}
\end{enumerate}
\textit{Sous ces conditions, le foncteur canonique fidèle}
\[
\uGamma(\mathcal{E}/J) \longrightarrow \prod_{j\in\Ob J} \mathcal{E}_j \tag*{(4.10.3.1)}
\]
\textit{est cofinal, et la catégorie source est grande devant $\pi$.}

Il est $\pm$ immédiat que les hypothèses de 4.10.1 impliquent que celles de 4.10.3 sont satisfaites pour $\mathcal{E} = \mathcal{D}_\varphi \longrightarrow J$ (4.10.5). Pour
\msp{page 249}
ceci, on utilise le

% marginal vertical bar printed in the left margin beside the following statement
\textbf{Lemme 4.10.4.} \textit{Soient $C$ une catégorie, et $\pi$ un cardinal infini. Alors le foncteur d'inclusion $C \longrightarrow \mathcal{M} \overset{\textup{déf}}{=} \Ind_\pi C$ commute aux $\varinjlim$ de cardinal $\leqslant \pi$.}

Soit en effet $I$ telle que $\card I \leqslant \pi$, et $(x_i)_{i\in I}$ un système inductif sur $C$, y ayant une limite $x$. Soit d'autre part $X = \varinjlim\limits_{I'} \xi_{i'}$ un objet de $\Ind_\pi C$, avec $I'$ grande devant $\pi$. On a


\endgroup
% END UNIT der-ch18/der-ch18_p091-095.tex
% BEGIN UNIT der-ch18/der-ch18_p096-100.tex
\begingroup\makeatletter%
\newcommand{\printedpage}[1]{\marginpar{\footnotesize p.~#1}}%
\newcommand{\msp}[1]{\marginpar{\footnotesize [#1]}}%
\newcommand{\uGamma}{\underline{\Gamma}}%
\newcommand{\Hom}{\operatorname{Hom}}%
\newcommand{\Ob}{\mathrm{Ob}\,}%
\newcommand{\card}{\mathrm{card}\,}%
\makeatother\let\printedpage\derpp\let\msp\dermsp%
\let\mspin\DERmspin%


\printedpage{94}
% head: A. Grothendieck : Dérivateurs, Ch. XVIII, Édition des Universités de Montpellier II et Paris VII
donc
\[
\begin{aligned}
\Hom_{\mathcal{M}}(x,X) &\simeq \varinjlim_{I'} \Hom(x,\xi_{i'})\\
&\simeq \varinjlim_{I'} \varprojlim_{I} \Hom(x_i,\xi_{i'})\\
&\overset{\text{SGA 4 I 9.8}}{\simeq} \varprojlim_{I} \underbrace{\varinjlim_{I'} \Hom(x_i,\xi_{i'})}_{\simeq \Hom(x_i,X)}\\
&\simeq \varprojlim_{I} \Hom(x_i,X)\ ,
\end{aligned}
\]
ce qui prouve que $x$ est aussi une $\varinjlim$ des $x_i$ dans $\mathcal{M}$ tout entier.

\msp{250}
Ainsi, on est réduit à prouver 4.10.3. Appliquant 4.9.6 2°, on est ramené à prouver, pour le foncteur envisagé, les conditions F 1 et F $3_\pi$ (page 225).

Dans les arguments qui suivent, on va utiliser les conditions a), b) de 4.10.3, mais non, pour l'instant, la condition c.

\textbf{Lemme 4.10.5.} \textit{Sous les conditions de 4.10.3 a) b) (à l'exclusion de c), pour que (4.10.3.1) soit cofinal et $\Gamma(\mathcal{E}/J)$ grande devant $\pi$, il faut et il suffit que le foncteur envisagé satisfasse à la condition F 1, rappelée ci-dessous, et à la variante F$'$ 1 de celle-ci énoncée ci-dessous :}
\begin{itemize}[leftmargin=3em]
\item[F\,1)] {\itshape Pour tout système $(X(j))_{j\in\Ob J}$ d'éléments $X(j)\in\Ob\mathcal{E}_j$, existe une \textup{section} $G$ de $\mathcal{E}$ sur $J$, et un système de flèches
\[
X(j) \longrightarrow G(j) \qquad (j\in ObJ)\ . \tag*{\textup{(4.10.5.1)}}
\]}
\item[F$'$1)] {\itshape Sous les conditions de \textup{F\,1} précédent, supposons de plus une section $F$ de $\mathcal{E}$ sur $J$ \msp{251} et un système de flèches
\[
F(j) \longrightarrow X(j) \qquad (j\in\Ob J)\ ,
\]
alors $G$ et le système de flèches (4.10.5.1) peut être choisi de telle façon que les flèches composées
\[
F(j) \longrightarrow X(j) \longrightarrow G(j)
\]
définissent un homomorphisme de sections
\[
F \longrightarrow G
\]
(dans la catégorie $\uGamma(\mathcal{E}/J)$).}
\end{itemize}
(\footnote[106]{\textbf{N.B.} Les conditions F 1 et F$\,'$\,1 équivalent à dire que le foncteur envisagé satisfait aux conditions F1 et F3, vu encore (en vertu de 4.9.6.3, 1°) que $\uGamma(\mathcal{E}/J)$ est \textit{filtrante} sans plus, et le foncteur fidèle (4.10.3.1) cofinal. Cela implique déjà que $\uGamma(\mathcal{E}/J)$ est grande devant $\pi$.}).

\textsc{Démonstration}. Il faut voir que F 1 \& F$'$ 1 équivaut à F 1 \& F $3_\pi$. L'implication $\Longleftarrow$ est $\pm$ immédiate, c'est de $\Longrightarrow$ que nous aurons besoin, donc
\[
\text{F 1 \& F}'\text{ 1} \quad\Longrightarrow\quad \text{F } 3_\pi\ .
\]
% foot: Version 17 décembre 2025

\printedpage{95}
% head: A. Grothendieck : Dérivateurs, Ch. XVIII, Édition des Universités de Montpellier II et Paris VII
Pour ceci, réénonçons F $3_\pi$ dans le contexte présent. On se donne une famille $(F_\alpha)_{\alpha\in A}$, $\card A\leqslant\pi$, d'objets de $\uGamma(\mathcal{E}/J)$, un objet $F$ de la même catégorie, et pour tout $j$ dans $\Ob J$, un système de flèches
\[
F_\alpha(j) \xrightarrow{u_\alpha(j)} F(j)\ ,
\]
mais sans supposer que pour $\alpha$ fixé, $j$ variable

\msp{252}
cela définisse un homomorphisme de sections $F_\alpha\longrightarrow F$. On cherche une section $G$ de $\mathcal{E}$ sur $J$, et un homomorphisme de sections $F\longrightarrow G$, tels que pour tout $\alpha$, le système des composés
\[
F_\alpha(j) \longrightarrow F(j) \longrightarrow G(j)
\]
définisse un homomorphisme de sections $F_\alpha\longrightarrow G_\alpha$. Soit donc $\lambda : j\longrightarrow j'$ une flèche de $J$, considérons les carrés (\textit{non commutatifs}) % [sic] F_\alpha -> G_\alpha -> F_\alpha -> G, the section sought is G (cf. p. 95 below: F_\alpha -> G)
\[
\begin{tikzcd}
F_\alpha(j) \arrow[r,"u_\alpha(j)"] \arrow[d,"F_\alpha(\lambda)"'] \arrow[dr,phantom,"!"] & F(j) \arrow[d,"F(\lambda)"] \\
F_\alpha(j') \arrow[r,"u_\alpha(j')"] & F(j')
\end{tikzcd}
\tag*{$(*)$}
\]
(indexés par $\alpha\in A$). Pour $j'$ fixé dans $J$, le nombre de ces carrés est majoré par $\card\Ob(J/j')\times\card A\leqslant\pi\cdot\pi=\pi$. Il s'ensuit, comme $\mathcal{E}_{j'}$ est grande devant $\pi$, qu'il existe un objet $X(j')$ de $\mathcal{E}_{j'}$, et une flèche $v(j') : F(j')\longrightarrow X(j')$ de telle façon que pour tout $\alpha\in A$, et tout $\lambda\in\Ob J/j'$, le carré déduit de $(*)$

\msp{253}
en composant avec $F(j')\longrightarrow X(j')$ soit commutatif. On fait ce choix pour tout objet $j'$ de $J$. Appliquons alors F$'$ 1 à $F$ et au système de flèches $F(j)\xrightarrow{v(j)}X(j)$, on trouve $G$ et un système de flèches
\[
w(j)\ :\ X(j) \longrightarrow G(j)
\]
de telle façon que les flèches composées
\[
\begin{tikzcd}[column sep=small]
w(j)v(j) : & F(j) \arrow[rr] \arrow[dr,dotted] & & G(j) \\
& & X(j) \arrow[ur,dotted] &
\end{tikzcd}
\]
définissent un homomorphisme de sections. Je dis que cet homomorphisme a la propriété voulue pour F $3_\pi$, i.e. que pour tout $\alpha\in A$, le système des composés
\[
F_\alpha(j) \longrightarrow F(j) \longrightarrow G(j)
\]
(pour $j\in\Ob J$ variable) définit un homomorphisme de sections $F_\alpha\longrightarrow G$. En effet, il suffit pour cela de considérer, pour une flèche $\lambda$ de $J$, le
% foot: Version 17 décembre 2025
\printedpage{96}
% head: A. Grothendieck : Dérivateurs, Ch. XVIII, Édition des Universités de Montpellier II et Paris VII
\msp{254}
diagramme
\[
\begin{tikzcd}
F_\alpha(j) \arrow[r,"u_\alpha(j)"] \arrow[d,"F_\alpha(\lambda)"'] \arrow[dr,phantom,"!"] & F(j) \arrow[r,"v(j)"] \arrow[d,"F(\lambda)"'] & X(j) \arrow[r,"w(j)"] \arrow[d,phantom,"\text{comm.}"] & G(j) \arrow[d,"G(\lambda)"] \\
F_\alpha(j') \arrow[r,"u_\alpha(j')"] & F(j') \arrow[r,"v(j')"] & X(j') \arrow[r,"w(j')"] & G(j')\ .
\end{tikzcd}
\]
Le lemme 4.10.5 étant prouvé, tout revient donc à prouver la condition F 1, et son complément F$'$ 1, qui se prouvera dans la même haleine.

\textbf{4.10.6.} On se donne donc le système d'objets
\[
X(j)\ \in\ \Ob\mathcal{E}_j \qquad (j\in\Ob J)\ , \tag*{(4.10.6.1)}
\]
et le cas échéant, un objet $F$ de $\uGamma(\mathcal{E}/J)$, et un système de flèches
\[
F(j) \xrightarrow{u_j} X(j)\ , \tag*{(4.6.10.2)} % [sic] (4.6.10.2) -> (4.10.6.2), numbering of 4.10.6; cited as (4.10.6.2) below
\]
et on se propose de construire $G$ de $\uGamma(\mathcal{E}/J)$, et un système de flèches
\[
X(j) \xrightarrow{w(j)} G(j)\ , \tag*{(4.10.6.3)}
\]
de telle façon que (si on a la donnée supplémentaire des $u(j)$) les flèches composées
\[
F(j) \xrightarrow{u(j)} X(j) \xrightarrow{v(j)} G(j) % [sic] v(j) -> w(j), the arrow X(j) -> G(j) is w(j) in (4.10.6.3)
\]

\msp{255}
soient fonctorielles en $j$.

On va faire une construction transfinie sur les ordinaux $\alpha$ tels que
\[
\card\alpha\ \leqslant\ \pi
\]
d'un système inductif
\[
(X_\alpha(j))_\alpha
\]
d'objets de $\prod_{j\in\Ob J}\mathcal{E}_j$. On va d'abord ignorer la donnée éventuelle (4.10.6.2). On posera
\[
X_0(j)\ =\ X(j)\ , \tag*{(4.10.6.4)}
\]
et on va donner la prescription pour passer de $X_\alpha$ à $X_{\alpha+1}$. La construction inclura aussi, pour toute flèche $\lambda : j'\longrightarrow j$ de $J$ et tout $\alpha$, la donnée de
\[
X_\alpha(j) \xrightarrow{t_\alpha(\lambda)} X_{\alpha+1}(j)\ , \qquad \text{flèche (dans $\mathcal{E}$) au dessus de $\lambda$,} \tag*{(4.10.6.5)} % [sic] X_\alpha(j) -> X_\alpha(j'), the source of an arrow over \lambda : j' -> j (cf. p. 97: X_\alpha(j') -> X_{\alpha+1}(j))
\]
et ces flèches, et les flèches de transition $X_\alpha(j)\xrightarrow{\tau_\alpha(j)}X_{\alpha+1}(j)$, doivent
% foot: Version 17 décembre 2025
\printedpage{97}
% head: A. Grothendieck : Dérivateurs, Ch. XVIII, Édition des Universités de Montpellier II et Paris VII
\msp{256}
satisfaire aux trois conditions suivantes :
\begin{itemize}[leftmargin=3em]
\item[(i)] Si $\lambda$ est une flèche identique $\mathrm{id}_J : J\longrightarrow J$, alors \fbox{$t_\alpha(\lambda)=\tau_\alpha(j)$.} % [sic] id_J : J -> J -> id_j : j -> j, the identity arrow of an object j of J
\item[(ii)] Dans le diagramme
\[
\begin{tikzcd}
X_\alpha(j') \arrow[r,"\tau_\alpha(j')"] \arrow[dr,"t_\alpha(\lambda)"] & X_{\alpha+1}(j') \arrow[r] \arrow[dr,"t_{\alpha+1}(\lambda)"] \arrow[d,phantom,"\bullet"] & X_{\alpha+2}(j') \arrow[r] \arrow[dr] & X_{\alpha+3}(j') \arrow[r] \arrow[dr,dotted] & {} \\
X_\alpha(j') \arrow[r] & X_{\alpha+1}(j') \arrow[r,"\tau_{\alpha+1}(j')"] & X_{\alpha+2}(j') \arrow[r,"\tau_{\alpha+2}(j')"] & X_{\alpha+3}(j') \arrow[r] & {,}
\end{tikzcd}
\] % [sic] bottom row printed with j' (X_\alpha(j'), X_{\alpha+1}(j'), \tau_{\alpha+1}(j'), X_{\alpha+2}(j'), \tau_{\alpha+2}(j'), X_{\alpha+3}(j')) -> j, the arrows t(\lambda) go to the objects at j (cf. the text below: \tau_{\alpha+1}(j), \tau_{\alpha+2}(j))
$\tau_{\alpha+2}(j)$ “rend le carré $\bullet$ commutatif”, i.e. égalise $t_{\alpha+1}(\lambda)\tau_\alpha(j')$ et $\tau_{\alpha+1}(j)t_\alpha(\lambda)$ :
\[
\boxed{\tau_{\alpha+2}(j)t_{\alpha+1}(\lambda)\tau_\alpha(j')\ =\ \tau_{\alpha+2}(j)\tau_{\alpha+1}(j)t_\alpha(\lambda)\ .}
\]
\item[(iii)] Si on a deux flèches consécutives $\lambda : j''\longrightarrow j'$, $\mu : j'\longrightarrow j$, on a donc le diagramme
\[
\begin{tikzcd}
X_\alpha(j'') \arrow[r] \arrow[dr,"t_\alpha(\lambda)"] \arrow[ddr,"t_\alpha(\mu\lambda)"] & X_{\alpha+1}(j'') \arrow[r] \arrow[dr] & X_{\alpha+2}(j'') \arrow[r] \arrow[dr] & X_{\alpha+3}(j'') \\
X_\alpha(j') \arrow[r] & X_{\alpha+1}(j') \arrow[r] \arrow[dr,"t_{\alpha+1}(\mu)"] & X_{\alpha+2}(j') \arrow[r] \arrow[dr] & X_{\alpha+3}(j') \\
X_\alpha(j) \arrow[r] & X_{\alpha+1}(j) \arrow[r,"\tau_{\alpha+1}(j)"] & X_{\alpha+2}(j) \arrow[r,"\tau_{\alpha+2}(j)"] & X_{\alpha+3}(j)\ ,
\end{tikzcd}
\]
\[
\boxed{\tau_{\alpha+2}(j)t_{\alpha+1}(\mu)t_\alpha(\lambda)\ =\ \tau_{\alpha+2}(j)\tau_{\alpha+1}(j)t_\alpha(\mu\lambda)}
\]
\end{itemize}

\msp{257}
Ainsi, pour définir les $X_{\alpha+1}(j)$, il faut tenir compte de toutes ces contraintes. Donc :
\begin{itemize}[leftmargin=3em]
\item[1°)] Pour $j$ fixé, $X_{\alpha+1}(j)$ doit être but des flèches
\[
X_\alpha(j') \xrightarrow{t_\alpha(\lambda)} X_{\alpha+1}(j)\ ,
\]
pour les $\lambda\ :\ j'\longrightarrow j$ dans $\Ob J/j$ -- y compris pour $\lambda=\mathrm{id}_j$, auquel cas $t_\alpha(\lambda)\ :\ X_\alpha(j)\longrightarrow X_{\alpha+1}(j)$ jouera le rôle du morphisme de transition $\tau_\alpha(j)$.
\item[2°)] Si $\alpha$ est de la forme $\alpha'+2$, i.e. s'il n'est ni ordinal limite, ni successeur d'un ordinal limite, alors on exige de plus, pour les morphismes de transition $\tau_\alpha(j)$, d'égaliser les couples de flèches prévus dans (ii) et (iii) ci-dessus.
\end{itemize}
Comme l'ensemble de ces données dans 1°), ou de ces contraintes dans 2°),
\msp{258}
est de cardinal $\leqslant\pi$, et que la catégorie $\mathcal{E}_j$ est grande devant $\pi$, on voit qu'on peut bel et bien trouver $X_{\alpha+1}(j)$ et des $t_\alpha(\lambda)\ :\ X_\alpha(j')\longrightarrow X_{\alpha+1}(j)$, de façon à satisfaire (si $\alpha=\alpha'{+}2$) aux contraintes rappelées dans 2°. Cela achève la description de la construction de $X_{\alpha+1}(\bullet)$, en termes de $X_\beta(\bullet)$ $(\beta\leqslant\alpha)$ et des $t_\beta(\lambda)$ $(\beta<\alpha)$.

Il faut enfin définir $X_\alpha(\bullet)$ pour $\alpha$ un ordinal limite. Cette fois, on lui demande seulement de recevoir les $X_\beta(\bullet)$ pour $\beta<\alpha$. Comme l'ensemble $I_\alpha$ de ces $\beta$ est de cardinal $\leqslant\pi$, et que les $\mathcal{E}_j$ sont grandes devant $\pi$, on peut bel et bien trouver les $X_\alpha(j)$.
% foot: Version 17 décembre 2025

\printedpage{98}
% head: A. Grothendieck : Dérivateurs, Ch. XVIII, Édition des Universités de Montpellier II et Paris VII
\msp{259}
On a donc prouvé l'essentiel du

\textbf{Lemme 4.10.7.} \textit{Sous les conditions de 4.10.3 a) b), et pour un système $X(\bullet)\in\Ob\prod_{j\in\Ob J}\mathcal{E}_j$ donné, on peut trouver un système inductif $(X_\alpha(i))_\alpha$ dans $\prod\mathcal{E}_j$, indexé par les ordinaux $\alpha$ telle que} $\card\alpha\leqslant\pi$, \textit{et pour toute flèche $\lambda : j'\longrightarrow j$ dans $J$ des flèches $t_\alpha(\lambda) : X_\alpha(j')\longrightarrow X_\alpha(j)$, de telle façon que l'on ait les conditions suivantes.} % [sic] (X_\alpha(i))_\alpha -> (X_\alpha(j))_\alpha, the index is j \in Ob J; [sic] X_\alpha(j') -> X_\alpha(j) -> X_\alpha(j') -> X_{\alpha+1}(j), cf. (4.10.6.5) and p. 97
\begin{itemize}[leftmargin=3em]
\item[\textit{1°)}] \textit{$X_0(\bullet)=X(\bullet)$.}
\item[\textit{2°)}] \textit{Les conditions (i) (ii) (iii) (p. 256).}
\item[\textit{3°)}] \textit{Si on se donne $F$ dans \textup{Ob}$\,\uGamma(\mathcal{E}/J)$, et des flèches $u(j) : F(j)\longrightarrow X(j)=X_0(j)$, on peut supposer que pour toute flèche $\lambda : j'\longrightarrow j$ dans $J$, on a}
\[
\begin{tikzcd}
F(j') \arrow[r,"u(j')"] \arrow[d,"F(\lambda)"'] & X_0(j') \arrow[dr,"t_0(\lambda)"] & \\
F(j) \arrow[r,"u(j)"] & X_0(j) \arrow[r,"\tau_0(j)"] & X_1(j)
\end{tikzcd}
\qquad\qquad \text{commutatif.}
\tag*{$(*)$}
\]
\end{itemize}

\msp{260}
Il reste seulement à prouver ce dernier complément 3°. C'est une contrainte uniquement sur la construction de $X_1(\bullet)$ et des $t_0(\lambda)$. Supposons donc que l'on ait construit $X_1(\bullet)$ sans tenir compte de cette contrainte. Pour $j$ fixé, il faut donc “rendre commutatif” tous ces carrés, i.e. toutes ces double-flèches $(t_0(\lambda)u(j'),\tau_0(j)u(j)F(\lambda))$ de but $X_1(j)$. Comme l'ensemble de ces double-flèches est de cardinal $\card\Ob\mathcal{E}/j\leqslant\pi$, on peut trouver $X_1(j)\longrightarrow X'_1(j)$ dans $\mathcal{E}_j$, égalisant toutes ces double-flèches, donc quitte à remplacer $X_1(j)$ par $X'_1(j)$, on gagne.

\msp{261}
Il est temps enfin de faire appel à l'hypothèse c) de 4.10.3. Soit donc $\alpha$ un ordinal limite, comme stipulé dans cette condition. Comme $\card\alpha\leqslant\pi$, a fortiori $\card\beta\leqslant\pi$ si $\beta<\alpha$, on a donc un système inductif
\[
(X_\beta(\bullet))_{\beta\in I_\alpha}
\]
dans $\prod_{j\in J}\mathcal{E}_j$. Comme les $\varinjlim$ de type $I_\alpha$ existent dans $\mathcal{E}_j$, on peut donc considérer
\[
G(j)\ =\ \varinjlim_{\beta\in I_\alpha} X_\beta(j)\ \in\ \Ob\mathcal{E}_j\ .
\]
Si $\lambda : j'\longrightarrow j$ est une flèche de $J$, on en déduit une $\lambda$-flèche
\[
G(\lambda)\ :\ G(j)=\varinjlim X_\beta(j) \longrightarrow G(j')\simeq\varinjlim X_{\beta+1}(j')
\]
par $G(\lambda)=\varinjlim t_\beta(\lambda)$. Il faut
\msp{262}
faire attention que les $t_\beta(\lambda)$ ne définissent pas un homomorphisme de systèmes projectifs, mais par la condition (ii), p. 256, ils définissent un homomorphisme d'ind-objets, donc
% foot: Version 17 décembre 2025


\endgroup
% END UNIT der-ch18/der-ch18_p096-100.tex
% BEGIN UNIT der-ch18/der-ch18_p101-105.tex
\begingroup\makeatletter%
\newcommand{\printedpage}[1]{\marginpar{\footnotesize p.~#1}}%
\newcommand{\msp}[1]{\marginpar{\footnotesize [#1]}}%
\newcommand{\cercle}[1]{\tikz[baseline=(X.base)]\node[draw,circle,inner sep=1pt](X){\upshape #1};}%
\newcommand{\defeq}{\overset{\text{\upshape déf}}{=}}%
\newcommand{\cardred}{\operatorname{card}\,\text{\upshape réd}\,}%
\newcommand{\Hom}{\underline{\operatorname{Hom}}}%
\newcommand{\uGamma}{\underline{\Gamma}}%
\newcommand{\Ind}{\operatorname{Ind}}%
\newcommand{\into}{\lhook\joinrel\longrightarrow}%
\makeatother\let\printedpage\derpp\let\msp\dermsp%
\let\mspin\DERmspin%


\printedpage{99}
% head: A. Grothendieck : Dérivateurs, Ch. XVIII, Édition des Universités de Montpellier II et Paris VII
% foot: Version 17 décembre 2025

un homomorphisme sur les $\varinjlim$ (\footnote{Attention, c'est ici qu'on utilise tacitement le fait que le foncteur $\lambda_* : \mathcal{E}_{j'} \longrightarrow \mathcal{E}_j$ commute aux $\varinjlim$ de type $I_\alpha$, ce qui implique que les inclusions $\mathcal{E}_j \into \mathcal{E}$ commutent à ce type de $\varinjlim$ .}). La condition (i) assure que $\lambda \longmapsto G(\lambda)$ transforme identités en identités, et la condition (iii) que l'on a $G(\mu\lambda) = G(\mu)G(\lambda)$. Donc $G$ apparaît comme une section de $\mathcal{E}$ sur $J$. On a d'autre part des flèches
\[
X(j) \defeq X_0(j) \longrightarrow G(j) = \varinjlim_{\beta\in I_\alpha} X_\beta(j) \ ,
\]
ce qui prouve F (1). D'autre part, la condition 3°) du lemme prouve la validité de F$'$ (1). Cela achève la démonstration de 4.10.3, et par là de 4.10.1 et de 4.9.3.

\msp{page 263}

\textbf{Corollaire 4.10.8} (du théorème 4.10.1). \textit{Soit $\mathcal{N} \longrightarrow J$ un foncteur cofibrant, tel que les fibres soient des catégories accessibles (4.8), les foncteurs cochangement de base $\lambda_*$ : $\mathcal{N}_{j'} \longrightarrow \mathcal{N}_j$ accessibles, enfin $J$ équivalente à une petite catégorie. Alors la catégorie $\mathcal{M} = \uGamma(\mathcal{N}/J)$ est accessible. Plus précisément, soit $\pi_0$ un cardinal infini tel que les catégories $\mathcal{N}_j$ soient $\pi_0$-accessibles, et que les foncteurs $\lambda_*$ soient $\pi_0$-accessibles, et soit $\pi$ un cardinal infini tel que l'on ait :}
\begin{enumerate}[leftmargin=*]
\item[\cercle{a}] $\pi^{\pi_0} = \pi$.
\item[\cercle{b}] $\forall\ \lambda : j' \longrightarrow j$ \textit{dans} $J$, $\lambda_* : \mathcal{E}_j \longrightarrow \mathcal{E}_{j'}$ \textit{applique} $(\mathcal{E}_j)_{\pi_0}$ \textit{dans} $(\mathcal{E}_{j'})_{\pi_0}$.
\item[\cercle{c}] $\cardred J \leqslant \pi$.
\end{enumerate}
\textit{Alors $\mathcal{M}$ est $\pi$-accessible, et plus précisément, on a}
\[
\mathcal{M} \ \simeq \ \Ind_\pi(\mathcal{M}_\pi) \ ,
\]
\textit{où $\mathcal{M}_\pi$ est la sous-catégorie pleine de $\mathcal{M}$ formée des sections $F$ de $\mathcal{M}$ sur $J$ telles que l'on ait $F(j) \in (\mathcal{E}_j)_\pi \ \forall\ j \in \operatorname{Ob} J$,} \msp{page 264} \textit{$(\mathcal{E}_j)_\pi$ désignant la sous-catégorie pleine de $\mathcal{E}_j$ formée des objets $\pi$-accessibles de $\mathcal{E}_j$.}

En effet, l'hypothèse au début implique l'existence de $\pi_0$ et de $\pi$ satisfaisant les conditions dites. Les conditions a) et b) sur $\pi$, plus le fait que les foncteurs $\lambda_*$ soient $\pi_0$-accessibles, impliquent que la sous-catégorie strictement pleine $\mathcal{N}_\pi$ de $\mathcal{N}$ dont les fibres sont les $(\mathcal{N}_j)_\pi$, est une sous-catégorie cofibrée, et que de plus $\mathcal{N}$ est $J$-équivalent à $\Ind_\pi(\mathcal{N}_\pi/J)$. D'autre part, les fibres de $\mathcal{N}_\pi/J$ sont stables par $\varinjlim$ grandes devant $\pi_0$ et de cardinal $\leqslant \pi$, donc la condition “ordinale” de 4.10.1 est vérifiée (en y faisant $\mathcal{D} = \mathcal{N}_\pi$). Comme de plus $\operatorname{card} J \leqslant \pi$, on en conclut par 4.10.1 que $\mathcal{M} \simeq \Ind_\pi(\mathcal{M}_\pi)$.

\msp{page 265}

\textbf{Corollaire 4.10.9.} \textit{Soit $\mathcal{N} \longrightarrow J$ une “cofibration en paratopos” (\footnote{terme ambigu, cf. la double explication.}), i.e. un foncteur cofibrant tel que les fibres soient des catégories paratopos, et les foncteurs cochangement de base soient des foncteurs images inverses (resp. images directes) pour des morphismes de paratopos, i.e. qu'ils commutent aux (petites) $\varinjlim$ (resp. aux petites $\varprojlim$ ). Supposons de plus $J$ équivalente à une petite catégorie. Alors $\uGamma(\mathcal{N}/J)$ est une catégorie paratopos.}

En effet, c'est une catégorie accessible en vertu de 4.10.8, d'autre part l'hypothèse sur les $\lambda_*$ implique aussitôt que $\uGamma(\mathcal{N}/J)$ est stable par petites $\varprojlim$ (resp. par petites $\varprojlim$
\printedpage{100}
% head: A. Grothendieck : Dérivateurs, Ch. XVIII, Édition des Universités de Montpellier II et Paris VII
% foot: Version 17 décembre 2025
\texttt{[plutôt} $\varinjlim$ \texttt{?]}). Donc $\uGamma(\mathcal{N}/J)$ est une catégorie paratopos en vertu de 4.8.9 b $\Longrightarrow$ a (resp. c $\Longrightarrow$ a).

\textbf{Corollaire 4.10.10.} \textit{Soit $\mathcal{N} \longrightarrow J$ une catégorie} fibrée \textit{(non cofibrée!), telle que les fibres soient des catégories paratopos et les foncteurs} \msp{page 266} \textit{changement de base soient des foncteurs images directes pour des morphismes de paratopos (i.e. des foncteurs accessibles, commutant aux petites $\varprojlim$ ). Supposons de plus $J$ équivalente à une petite catégorie. Alors $\uGamma(\mathcal{N}/J)$ est une catégorie paratopos.}

En effet, l'hypothèse implique que pour toute flèche $\lambda : j' \longrightarrow j$ dans $J$, le foncteur $\lambda_*$ : $\mathcal{N}_{j'} \longrightarrow \mathcal{N}_j$ admet un adjoint à gauche $\lambda^* : \mathcal{N}_j \longrightarrow \mathcal{N}_{j'}$, d'où s'ensuit que $\mathcal{N} \longrightarrow J$ est aussi cofibrant, les foncteurs de cochangement de base commutant aux petites $\varinjlim$ . Donc on peut appliquer 4.8.9, qui donne le résultat voulu.

\textbf{Question 4.10.11.} \textit{Soit $\mathcal{N} \longrightarrow J$ un foncteur} fibrant, \textit{tel que les $\mathcal{N}_j$ soient des catégories accessibles, et les foncteurs $\lambda^*$ de changement de base accessibles, $J$ étant équivalente à une petite catégorie. Je} \msp{page 267} \textit{conjecture que sous ces conditions, $\uGamma(\mathcal{N}/J) \defeq \Hom(J,\mathcal{N})$ est accessible, et qu'il en est de même de $\Hom\,\underline{\operatorname{cart}}_J(J,\mathcal{N})$.}

C'est ce qui est prouvé dans SGA 4 I 9.24 sous des hypothèses légèrement plus fortes (à l'aide de l'ancienne notion trop forte à mon gré, de filtrations cardinales). Je serais étonné que la démonstration (un peu longuette) de loc. cit. ne s'applique pas au cas général. Mais on aimerait avoir un énoncé qui coiffe celui-ci et 4.10.8, dans le cas d'un foncteur $\mathcal{N} \longrightarrow J$ qui ne serait supposé ni fibrant ni cofibrant, mais \ldots(???).

\section*{4.11\quad Application à l'existence des limites dans une catégorie de la forme $\Ind_{\pi_0}(C)$.}

\textbf{Proposition 4.11.1.} \textit{Soit} \msp{page 268} \textit{$C$ une catégorie équivalente à une petite catégorie, $\pi$ un cardinal infini, posons $\mathcal{M} = \Ind_\pi(C)$. Soit $\Phi$ un ensemble de catégories de cardinal réduit $\leqslant \pi$. On se propose d'examiner si $\mathcal{M}$ est “stable par $\varinjlim$ de type $\Phi$” (i.e. stable par $\varinjlim$ de type $I$, quel que soit $I \in \Phi$). Il suffit pour cela que $C$ est stable par $\varinjlim$ de type $\Phi$, et que de plus il existe un ordinal limite $\alpha$ avec $\operatorname{card}\alpha \leqslant \pi$ et tel que $C$ soit stable par $\varinjlim$ de type $I_\alpha$ (p.ex. tel que $I_\alpha \in \Phi$). Inversement, si $\mathcal{M}$ est stable par $\varinjlim$ de type $\Phi$, et si on suppose de plus $C$ karoubienne (i.e. stable par facteurs directs), alors $C$ est stable par $\varinjlim$ de type $\Phi$. Dans l'un et l'autre cas, le foncteur d'inclusion $C \into \mathcal{M}$ commute aux $\varinjlim$ envisagées.}

(\footnote{Expliciter ici les corollaires des types 4.13.20, 4.13.21 (p. 341).}).

\textsc{Démonstration}. a) Condition suffisante. On applique 4.9.3 pour conclure que
\printedpage{101}
% head: A. Grothendieck : Dérivateurs, Ch. XVIII, Édition des Universités de Montpellier II et Paris VII
% foot: Version 17 décembre 2025
\msp{page 269}
pour toute catégorie $J$ de cardinal réduit $\leqslant \pi$ (en particulier pour $J \in \Phi$), tout foncteur $\varphi : J \longrightarrow \mathcal{M} = \Ind_\pi(C)$ admet une représentation indicielle $\Phi : J \times I \longrightarrow C$, avec $I$ grande devant $\pi$. Ceci, plus la stabilité de $C$ par $\varinjlim$ de type $J$, suffit à établir l'existence de la $\varinjlim$ de $\varphi$.

La réciproque provient du fait que si $C$ est karoubienne, alors son image essentielle dans $\mathcal{M}$ n'est autre que $\mathcal{M}_\pi$, laquelle est stable dans $\mathcal{M}$ par toute $\varinjlim$ de cardinal $\leqslant \pi$ représentable dans $\mathcal{M}$. Si $\mathcal{M}$ est stable par $\varinjlim$ de type $\Phi$, il en est de même de $\mathcal{M}_\pi$, donc de $C$, q.e.d.

\textbf{Corollaire 4.11.2.} \textit{Pour que $\mathcal{M}$ soit stable par $\varinjlim$ filtrantes (resp. par petites $\varinjlim$ ), il suffit que $C$ le soit par $\varinjlim$ filtrantes (resp. par $\varinjlim$ )} de cardinal $\leqslant \pi$, \textit{et c'est aussi nécessaire si $C$ est karoubienne.}

Supposons $C$ stable par $\varinjlim$ filtrantes \msp{page 270} (resp. par $\varinjlim$ de cardinal $\leqslant \pi$, alors par 4.11.1 $\mathcal{M}$ est stable par ce même type de limites inductives. D'autre part, il est stable par $\varinjlim$ grandes devant $\pi$. Il suffit donc de prouver le

\textbf{Lemme 4.11.3.} \textit{Soit $\mathcal{M}$ une catégorie, $\pi$ un cardinal. Supposons que $\mathcal{M}$ soit stable par $\varinjlim$ filtrantes de cardinal $\leqslant \pi$ (resp. par $\varinjlim$ de cardinal $\leqslant \pi$, resp. par sommes des cardinal $\leqslant \pi$),} et \textit{par $\varinjlim$ grandes devant $\pi$. Alors $\mathcal{M}$ est stable par petites $\varinjlim$ filtrantes (resp. par petites $\varinjlim$ , resp. par petites sommes).} % [sic] des cardinal -> de cardinal, agreement (cf. "de cardinal" elsewhere in the statement)

Ce lemme se prouve comme 4.4.14, page 97.

Reste la réciproque dans 4.11.2, qui résulte de l'assertion similaire dans 4.11.1.

\textbf{Corollaire 4.11.4.} \textit{Quand $\mathcal{M} = \Ind_\pi C$ est-il stable par sommes ? S'il existe un ordinal $\alpha$ tel que $\operatorname{card}\alpha \leqslant \pi$, et que $C$ soit stable par $\varinjlim$ de type $I_\alpha$, alors il suffit que $C$ soit stable par sommes de cardinal $\leqslant \pi$. Cette condition est aussi nécessaire si $C$ est karoubienne.}

(\footnote{Pour autres variantes de 4.11.1, cf. 4.13.19 (p. 340) et corollaires.}).

Se prouve comme 4.11.2, en utilisant 4.11.3.

\msp{page 271}

\textbf{Proposition 4.11.5} (\footnote{?, démonstration pas achevée.}). \textit{Soit $\Phi$ un ensemble de catégories $I$ équivalentes à des catégories finies, et “rigides”, i.e. pour tout $i \in \operatorname{Ob} I$, $\Hom_I(i,i) = \{\mathrm{id}_i\}$. Soient $C$ une catégorie équivalente à une petite catégorie, $\pi$ un cardinal infini, $\mathcal{M} = \Ind_\pi C$. Pour que $\mathcal{M}$ soit stable par $\varinjlim$ (resp. par $\varprojlim$ ) de type $\Phi$, il} suffit \textit{que $C$ le soit. Dans le cas non respé, si $C$ est karoubienne, cette condition est aussi nécessaire.}

Cette dernière assertion est contenue dans l'assertion similaire dans 4.11.1. Pour la première assertion, cela résulte de la façon habituelle du

\textbf{Lemme 4.11.6} (\footnote{?, démonstration non explicitée est peut-être faux!}). \textit{Soit $J$ une catégorie équivalente à une catégorie finie, $J$ rigide, $C$ une catégorie équivalente à une petite catégorie, $\pi$ un cardinal infini. Alors tout foncteur}
\printedpage{102}
% head: A. Grothendieck : Dérivateurs, Ch. XVIII, Édition des Universités de Montpellier II et Paris VII
% foot: Version 17 décembre 2025
\textit{$J \longrightarrow \mathcal{M} \defeq \Ind_\pi C$ admet une représentation indicielle $J \times I \longrightarrow C$, avec $I$ grande devant $\pi$, en d'autres termes, le} \msp{page 272} \textit{foncteur canonique pleinement fidèle}
\[
\Ind_\pi(\Hom(J,C)) \ \longrightarrow \ \Hom(J,\mathcal{M})
\]
\textit{déduit de l'inclusion $C \into \mathcal{M}$, d'où $\Hom(J,C) \into \Hom(J,\mathcal{M})$ (cf. 4.9), est une équivalence de catégories.} (\textbf{N.B.} J'ignore s'il en est ainsi sans hypothèse de rigidité sur $J$.)

\textsc{Démonstration}. Cela se prouve comme le résultat similaire pour $\Ind(C)$ au lieu de $\Ind_\pi C$, cf. SGA 4 I 8.8.5.

\msp{page 273}

% a vertical rule is printed in the left margin beside Corollaire 4.11.7
\textbf{Corollaire 4.11.7} (\footnote{sous réserve de validité de 4.11.5.}). \textit{Pour que $\mathcal{M} = \Ind_\pi C$ soit stable par $\varinjlim$ (resp. par $\varprojlim$ )} finies, \textit{il} suffit \textit{que $C$ le soit. Et c'est aussi nécessaire dans le cas non respé, lorsqu'on suppose $C$ karoubienne.}

\section*{4.12\quad Foncteurs accessibles entre catégories accessibles}

\textbf{Proposition 4.12.1} (\footnote{Remonterait à 4.8.}). \textit{Soient $\mathcal{M}$, $\mathcal{N}$ deux catégories accessibles (4.8), et $f : \mathcal{M} \longrightarrow \mathcal{N}$ un foncteur accessible. Alors il existe des cardinaux infinis $\pi_0$, $\pi_1$ tels que $\pi_1 \geqslant \pi_0$ et que pour tout cardinal $\pi \geqslant \pi_1$, on ait}
\[
f(\mathcal{M}_\pi) \ \subset \ \mathcal{N}_{\pi^{\pi_0}} \ , \tag*{(4.12.1)} % [sic] (4.12.1) -> (4.12.1.1), the formula is cited below as (4.12.1.1) and the next one is numbered (4.12.1.2)
\]
\textit{et en particulier}
\[
f(\mathcal{M}_\pi) \ \subset \ \mathcal{N}_\pi \qquad \textit{si } \pi = \pi^{\pi_0}. \tag*{(4.12.1.2)}
\]
\textit{De façon plus précise, on peut prendre pour $\pi_0$ un cardinal utile pour $\mathcal{M}$} \msp{page 274} \textit{et tel que $f$ soit $\pi_0$-accessible (p.ex. le plus petit cardinal de ce type). Soit alors $\pi_1$ un cardinal tel qu'on ait $f(\mathcal{M}_{\pi_0}) \subset \mathcal{N}_{\pi_1}$, $\pi_1 \geqslant \cardred \mathcal{M}_{\pi_0}$, et enfin tel que $\mathcal{N}_{\pi_1}$ soit} \texttt{[une]} \textit{sous-catégorie strictement génératrice de $\mathcal{N}$. Le couple $(\pi_0,\pi_1)$ satisfait aux conditions énoncées.}

\textsc{Démonstration}. On aura en effet
\[
\mathcal{N}_1 \ \xleftarrow{\ \sim\ } \Ind_{\pi_0}(\mathcal{M}_{\pi_0})
\]
et $f$ commutant aux $\varinjlim$ grandes devant $\pi_0$, sera connu (à isomorphisme unique près) par sa restriction à $\mathcal{M}_{\pi_0}$, laquelle se factorise par $\mathcal{N}_{\pi_1}$,
\[
f|\mathcal{M}_{\pi_0} : \mathcal{M}_{\pi_0} \ \longrightarrow \ \mathcal{N}_{\pi_1} \ .
\]
Pour prouver (4.12.1.1), posant $\pi' = \pi^{\pi_0} \geqslant \pi$, et notant que $\pi'^{\pi_0} = \pi_0$ \texttt{[plutôt} $\pi'^{\pi_0} = \pi'$\texttt{]},
\[
f(\mathcal{M}_\pi) \ \subset \ f(\mathcal{M}_{\pi'}) \ \overset{?}{\subset} \ \mathcal{N}_{\pi'} \ ,
\]
\printedpage{103}
% head: A. Grothendieck : Dérivateurs, Ch. XVIII, Édition des Universités de Montpellier II et Paris VII
% foot: Version 17 décembre 2025
on voit qu'il suffit de prouver (4.12.1.2) (pour $\pi'$). Donc on peut supposer $\pi = \pi^{\pi_0}$ et prouver (4.12.1.2). Mais comme $\pi_1 \geqslant \cardred \mathcal{M}_{\pi_0}$, on sait que si $\pi$ est tel que $\pi \geqslant \pi_1$, $\pi^{\pi_0} = \pi$, alors $\mathcal{M}_\pi$ est formé des $\varinjlim\limits_{I} x_i$, les $x_i \in \mathcal{N}_{\pi_0}$, et $I$ grande devant $\pi_0$, et $\operatorname{card} I \leqslant \pi$.

\msp{page 275}

On a alors ($f$ étant $\pi_0$-accessible)
\[
f(x) \ = \ \varinjlim_{I} f(x_i) \ ,
\]
les $f(x_i)$ sont dans $\mathcal{N}_{\pi_1}$, donc dans $\mathcal{N}_\pi$, or $\mathcal{N}_\pi$ stable par $\varinjlim$ de cardinal $\leqslant \pi$ (puisque $\mathcal{N}$ est engendrée par $\mathcal{N}_{\pi_1}$, formée d'objets $\pi_1$-accessibles et a fortiori $\pi$-accessibles, cf. SGA 4 I 9.9). D'où $f(x) \in \mathcal{N}_\pi$, q.e.d.

\textbf{Corollaire 4.12.2.} \textit{Supposons que $\pi_0$ soit aussi utile pour $\mathcal{N}$, et que $\pi_1 \geqslant \cardred \mathcal{N}_{\pi_0}$. Soit $\pi$ un cardinal $\geqslant \pi_1$, et tel que $\pi^{\pi_0} = \pi$. Alors $f$ s'identifie, à équivalence de catégories et à isomorphisme de foncteurs près, au foncteur}
\[
\Ind(f_\pi) \ : \ \Ind_\pi(\mathcal{M}_\pi) \ \longrightarrow \ \Ind_\pi(\mathcal{N}_\pi) \ ,
\]
\textit{où $f_\pi : \mathcal{M}_\pi \longrightarrow \mathcal{N}_\pi$ est le foncteur induit par $f$.}

\textbf{Proposition 4.12.3} (\footnote{remonte à 4.7.}). \textit{Soient $C$ une catégorie équivalente à une petite catégorie, $\mathcal{N}$ une $\mathfrak{U}$-catégorie, $\pi_0$ un cardinal infini.} \msp{page 276} \textit{On suppose que $\mathcal{N}$ satisfait $L_{\pi_0}$ (stabilité par $\varinjlim$ grandes devant $\pi_0$). Soit $\mathcal{M} = \Ind_{\pi_0} C$, $\Hom_{!\,\pi_0}(\mathcal{M},\mathcal{N})$ la sous-catégorie strictement pleine de $\Hom(\mathcal{M},\mathcal{N})$ formée des $f$ : $\mathcal{M} \longrightarrow \mathcal{N}$ $\pi_0$-accessibles. Alors le foncteur restriction}
\[
\Hom_{!\,\pi_0}(\mathcal{M},\mathcal{N}) \ \longrightarrow \ \Hom(C,\mathcal{N}) \tag*{(4.12.3.1)}
\]
\textit{est une équivalence de catégories. On trouve un foncteur quasi-inverse comme composé}
\[
\Hom(C,\mathcal{N}) \ \longrightarrow \ \Hom(\underbrace{\Ind_{\pi_0}(C)}_{\mathcal{M}},\Ind_{\pi_0}(\mathcal{N})) \ \longrightarrow \ \Hom(\mathcal{M},\mathcal{N}) \ , \tag*{(4.12.3.2)}
\]
\textit{où la deuxième flèche est le foncteur composition $F \longmapsto L \circ F$ avec le foncteur limite inductive $L : \Ind_{\pi_0}\mathcal{N} \longrightarrow \mathcal{N}$. Le foncteur composé (4.11.3.2) a une image contenue dans $\Hom_{!\,\pi_0}(\mathcal{M},\mathcal{N})$ et induit une équivalence entre cette dernière catégorie et $\Hom(C,\mathcal{N})$.} % [sic] (4.11.3.2) -> (4.12.3.2), the composed functor is the display just labelled (4.12.3.2)

\msp{page 277}

\textit{En d'autres termes, on peut dire que le foncteur d'inclusion canonique}
\[
C \ \longrightarrow \ \Ind_{\pi_0} C \ = \ \mathcal{M}
\]
\textit{est} 2-universel \textit{pour les foncteurs de $C$ dans des catégories $\mathcal{N}$ satisfaisant $L_{\pi_0}$, celles-ci étant considérées comme $0$-objets d'une $2$-catégorie où les $1$-objets sont les foncteurs $\mathcal{N} \longrightarrow \mathcal{N}'$ $\pi_0$-accessibles, i.e. commutent aux $\varinjlim$ grandes devant $\pi_0$.}

Cet énoncé est l'analogue, pour $C$ petite quelconque, de 4.4.16 (page 101), valable pour $C$ stable par $\varinjlim$ de cardinal $\leqslant \pi_0$. Il y a bien sûr d'autres variantes, en relation avec la commutation à d'autres types de $\varinjlim$ ( $\varinjlim$ filtrantes, sommes, limites inductives finies etc.) envisagées dans la section 4.11 précédente.

La démonstration de 4.12.3 est soritale.


\endgroup
% END UNIT der-ch18/der-ch18_p101-105.tex
% BEGIN UNIT der-ch18/der-ch18_p106-110.tex
\begingroup\makeatletter%
\newcommand{\printedpage}[1]{\marginpar{\footnotesize p.~#1}}%
\newcommand{\msp}[1]{\marginpar{\footnotesize [#1]}}%
\newcommand{\M}{\mathcal{M}}%
\newcommand{\N}{\mathcal{N}}%
\newcommand{\cH}{\mathcal{H}}%
\newcommand{\Hom}{\underline{\mathrm{Hom}}}%
\newcommand{\Ind}{\operatorname{Ind}}%
\newcommand{\Ob}{\operatorname{Ob}}%
\newcommand{\Fl}{\operatorname{Fl}}%
\newcommand{\card}{\operatorname{card}}%
\newcommand{\red}{\mathop{\textup{réd}}\nolimits}%
\newcommand{\uS}{\underline{S}}%
\makeatother\let\printedpage\derpp\let\msp\dermsp%
\let\mspin\DERmspin%


\printedpage{104}
% head: A. Grothendieck : Dérivateurs, Ch. XVIII, Édition des Universités de Montpellier II et Paris VII
% foot: Version 17 décembre 2025
\msp{page 278}
\textbf{Corollaire 4.12.4.} \textit{Soient $\M$ et $\N$ deux catégories accessibles, et soit $\pi_0$ un cardinal infini tel que $\M$ soit $\pi_0$-accessible, i.e. que $\pi_0$ soit un cardinal utile pour $\M$. Soit $\pi_1 = \sup(\pi_0,\allowbreak \card\red\M_{\pi_0})$. Alors pour tout cardinal $\pi \geqslant \pi_1$ utile pour $\M$ (et en particulier, pour tout $\pi \geqslant \pi_1$ tel que $\pi^{\pi_0} = \pi$), la catégorie des foncteurs $\M \longrightarrow \N$ qui sont $\pi$-accessibles est elle-même accessible.}

En effet, on aura $\M \simeq \Ind_\pi\M_\pi$, d'où par 4.12.3
\[
\Hom_{!\,\pi}(\M,\N) \xrightarrow{\sim} \Hom(\M_\pi,\N) ,
\]
or $\Hom(\M_\pi,\N)$ est accessible par 4.9.5 (p. 220). On peut préciser d'ailleurs ce point. Supposons $\pi_0$ choisi tel que $\pi_0$ soit utile non seulement pour $\M$, mais pour $\N$ également. Rappelons que l'on a $\card\Fl\M_\pi \leqslant 2^\pi$ (cf. 4.7.7, page 173).

\msp{page 279}
Posons donc
\[
\pi' = 2^\pi ,
\]
on a donc $\pi'^{\pi_0} = 2^{\pi\pi_0} = 2^\pi = \pi'$, et $\card\red\M_\pi \leqslant \pi'$. Donc 4.9.5 nous donne que
\[
\Hom(\M_\pi,\N) \simeq \Ind_{\pi'}(\Hom(\M_\pi,\M_{\pi'})) , % [sic] \M_{\pi'} -> \N_{\pi'}, the second argument should lie in N, as in Corollaire 4.12.5
\]
d'où :

\textbf{Corollaire 4.12.5.} \textit{Sous les conditions de 4.12.4, si $\pi_0$ est utile pour $\M$} et \textit{pour $\N$, alors pour $\pi \geqslant \pi_1$, $\Hom_{!\,\pi}(\M,\N)$ est $\pi'$-accessible, où $\pi' = 2^\pi$, de façon plus précise, on a une équivalence de catégories}
\[
\Hom_{!\,\pi}(\M,\N) \simeq \Ind_{\pi'}(\Hom(\M_\pi,\N_{\pi'})) .
\]

\textbf{Remarque 4.12.6.} Si $\M$, $\N$ sont accessibles, la catégorie $\Hom_{\mathrm{acc}}(\M,\N)$ des foncteurs accessibles n'est par toujours accessible. Si $\N$ satisfait $\mathrm{L}_{\pi_0}$, cette catégorie satisfait $\mathrm{L}_{\pi_0}$, % [sic] par toujours -> pas toujours, negation "n'est pas" intended
et les $\varinjlim$ grandes devant $\pi_0$ se calculent argument par argument.

\msp{page 280}
Si on pouvait trouver une sous-catégorie (équivalente à une petite catégorie) $\cH_\pi$ de $\cH = \Hom_{\mathrm{acc}}(\M,\N)$, pour un cardinal $\pi \geqslant \pi_0$ telle que tout objet $F$ de $\cH$ s'écrive $\varinjlim\limits_{I} F_i$, les $F_i$ dans $\cH_\pi$, $I$ grande devant $\pi$, alors si les objets de $\cH_\pi$ sont des foncteurs $\pi'$-accessibles de $\M$ dans $\N$, il en serait de même de tout $F$ dans $\cH$. Donc il existerait un cardinal infini $\pi'$ tel que tout foncteur accessible de $\M$ dans $\N$ soit $\pi'$-accessible. Mais sauf erreur, si p.ex. $\M = \N = \mathrm{Ens}$, il n'existe pas de tel cardinal. Je ne veux pas m'arrêter ici sur cette question, et faire un contre-exemple.

\msp{page 281}

\section*{4.13\quad Ensembles accessibles d'objets d'une catégorie accessible}

\textbf{Définition 4.13.1.} Soient $\M$ une catégorie accessible, $S$ une partie de $\Ob\M$, $\uS$ la sous-catégorie pleine de $\M$ définie par $S$. On dit que $S$ est une \textit{partie accessible} de $\Ob\M$, ou encore que $\uS$ est une \textit{sous-catégorie “accessible dans $\M$”}, si
\printedpage{105}
% head: A. Grothendieck : Dérivateurs, Ch. XVIII, Édition des Universités de Montpellier II et Paris VII
% foot: Version 17 décembre 2025
\begin{enumerate}
\item[a)] $\uS$ est strictement pleine, i.e. pour $x$, $y$ deux objets isomorphes de $\M$, $x \in S$ implique $y \in S$,
\item[b)] $\uS$ est accessible,
\item[c)] le foncteur d'inclusion $\uS \longrightarrow \M$ est accessible.
\end{enumerate}
Soient $\pi_0$ un cardinal utile pour $\M$, i.e. tel que $\M \simeq \Ind_{\pi_0}(\M_{\pi_0})$, soit $\pi_1 = \sup(\pi_0, \card\red\M_{\pi_0})$. Soit d'autre part $\pi'_0$ un cardinal tel que

\msp{page 282}
$\uS$ (supposé accessible dans $\M$) soit $\pi_0$-\textit{accessible}, i.e. $\pi'_0$-utile pour $\uS$, et de plus tel % [sic] \pi_0-accessible -> \pi'_0-accessible, the following "i.e. \pi'_0-utile" refers to \pi'_0
que $f : \uS \hookrightarrow \M$ soit $\pi'_0$-accessible. Soit $\pi'_1 \geqslant \pi'_0$ un cardinal tel que $\uS_{\pi'_0} \subset \M_{\pi'_1}$, $\pi'_1 \geqslant \card\red\uS_{\pi'_0}$. Il résulte de 4.12.1 que pour tout cardinal
\[
\pi \geqslant \pi'_1 , \qquad \text{tel que } \pi^{\pi'_0} = \pi, \tag*{$(*)$}
\]
on doit avoir
\[
\uS_\pi \subset \M_\pi , \qquad \text{donc } \uS_\pi \subset \M_\pi \cap \uS,
\]
d'ailleurs, comme $\uS \hookrightarrow \M$ est $\pi'_0$-accessible et a fortiori $\pi$-accessible, et pleinement fidèle, il s'ensuit aussitôt que $\M_\pi \cap S \subset S_\pi$, d'où
\[
S_\pi = \M \cap S . % [sic] \M \cap S -> \M_\pi \cap S, the preceding lines concern \M_\pi \cap S
\]

\msp{page 283}
D'ailleurs, la condition $(*)$ implique que $\pi$ est utile pour $\uS$, donc
\[
\uS \xleftarrow{\sim} \Ind_\pi(\uS_\pi) .
\]
$\Big($ D'ailleurs, \textit{si} $\pi'_0 \geqslant \pi_0$, la relation $\pi^{\pi'_0} = \pi$ implique a fortiori $\pi^{\pi_0} = \pi$, et si $\pi'_1 \geqslant \pi_1$, on voit que $(*)$ implique que $\pi$ est utile aussi pour $\M$, i.e.
\[
\M \xleftarrow{\sim} \Ind_\pi(\M_\pi) .
\]
D'autre part, le foncteur d'inclusion $\uS \hookrightarrow \M$ se déduit essentiellement de $\uS_\pi \hookrightarrow \M_\pi$ en passant aux catégories $\Ind_\pi$. (\footnote[116]{\texttt{[Le contenu de cette parenthèse est]} “facultatif”.}). $\Big)$

Ceci prouve entre autres le “il faut” dans la proposition qui suit :

\textbf{Proposition 4.13.2.} \textit{Soient $\M$ une catégorie accessible, $\pi_0$ un cardinal infini utile pour $\M$, $S$ une partie de} $\Ob\M$, \textit{$\uS$ la sous-catégorie pleine de $\M$ correspondante.}

\msp{page 284}
\begin{enumerate}
\item[\textit{a)}] \textit{Pour que $S$ soit accessible, il faut et il suffit qu'il existe un cardinal infini $\pi$ et une partie $S(\pi)$ de $\M_\pi$, tels que $S$ soit formé des objets de $\M$ qui sont isomorphes à un objet de la forme $\varinjlim\limits_{i\in I} x_i$, avec les $x_i$ dans $S(\pi)$ et $I$ grande devant $\pi$. En d'autres termes, il faut et il suffit qu'il existe un cardinal infini $\pi_1$ et une sous-catégorie pleine % [sic] \pi_1 -> \pi, the rest of the statement uses \pi
$\uS(\pi)$ de $\M_\pi$, tels que $\uS$ soit l'image essentielle du foncteur canonique pleinement fidèle}
\[
\Ind_\pi(\uS(\pi)) \longrightarrow \M \qquad \bigl(\xleftarrow{\sim} \Ind_\pi(\M)\ \text{(\footnotemark[117])}\bigr) , \tag*{(4.13.2.1)} % [sic] \Ind_\pi(\M) -> \Ind_\pi(\M_\pi), footnote 117 refers to \pi utile pour \M, i.e. \M \simeq \Ind_\pi(\M_\pi)
\]
\footnotetext[117]{si $\pi$ utile pour $\M$.}
\textit{prolongeant l'inclusion canonique $\uS(\pi) \hookrightarrow \M$ (cf. 4.12.3).}
\item[\textit{b)}] \printedpage{106}
% head: A. Grothendieck : Dérivateurs, Ch. XVIII, Édition des Universités de Montpellier II et Paris VII
% foot: Version 17 décembre 2025
\textit{Si $S$ est accessible, on peut trouver des cardinaux $\pi'_0$, $\pi'_1$ tels que $\pi'_1 \geqslant \pi'_0 \geqslant \pi_0$ (où $\pi_0$ est un cardinal donné utile pour $\M$)}
\msp{page 285}
\textit{et tels que tout cardinal $\pi \geqslant \pi'_1$ avec $\pi^{\pi'_0} = \pi$ satisfasse aux conditions de c), en prenant $S(\pi) = S \cap \Ob\M_\pi$, i.e. $\uS(\pi) = \uS \cap \M_\pi$. Sous ces conditions, $\pi$ est utile pour $\uS$} et \textit{pour $\M$, $S(\pi)$ est une sous-catégorie} strictement \textit{pleine de $\M_\pi$, stable dans $\M_\pi$ par facteurs directs (ou ce qui revient au même, $\M_\pi$ étant karoubienne, une sous-catégorie karoubienne de $\M_\pi$).}
\item[\textit{c)}] \textit{On peut prendre $\pi'_0$ le plus petit cardinal $\geqslant \pi_0$ utile pour $\uS$ et tel que $\uS \hookrightarrow \M$ soit $\pi'_0$-accessible, et $\pi'_1 = \sup(\pi'_0, \card\red\uS_{\pi'_0})$.}
\end{enumerate}

La partie b) ainsi que c) a déjà été prouvée, ce qui implique le “il faut” dans la partie a). Il reste à prouver le “il suffit” dans a). Mais l'hypothèse implique que
\[
\uS \xleftarrow{\sim} \Ind_\pi(S_\pi) ,
\]
ce qui implique que $\uS$ est accessible et même $\pi$-accessible, et elle implique aussi que $\uS$ est une sous-catégorie \textit{strictement} pleine de $\M$. Il reste à voir

\msp{page 286}
\texttt{[que]} le foncteur d'inclusion est accessible, en fait il est $\pi$-accessible, puisque (4.13.2.1) l'est (cf. 4.12.3).

\textbf{Définition 4.13.3.} Soient $\M$, $S \subset \Ob\M$, $\uS \subset \M$ comme dans 4.12.1, $\pi$ un cardinal infini. On dit que $S$ est une partie $\pi$-\textit{accessible} de $\Ob\M$, ou que $\uS$ est une sous-catégorie $\pi$-accessible dans $\M$, si $S$, $\pi$ satisfont les conditions de 4.13.3 a), i.e. $\uS$ est l'image essentielle % [sic] 4.13.3 a) -> 4.13.2 a), a self-reference of the definition, the conditions a) are those of Proposition 4.13.2
de $\Ind_\pi(S(\pi))$, où $S(\pi)$ est une sous-catégorie pleine de $\M_\pi$.

Notons que, quitte à renforcer $S(\pi)$ par $S'(\pi) = {\natural}$-clôture de $S(\pi)$ dans $\M_\pi$, on peut supposer que $S(\pi)$ est une sous-catégorie \textit{strictement pleine et karoubienne}

\msp{page 287}
de $\M_\pi$. On sait alors que, $S(\pi)$ étant karoubienne, $S(\pi) \longrightarrow (\underbrace{\Ind_\pi(S(\pi))}_{})_\pi$ est une équivalence de catégories. Comme $\Ind_\pi(S(\pi)) \simeq \uS$ et que $S(\pi)$ est strictement pleine dans $\M_\pi$, donc aussi dans $\M$ et a fortiori dans $\uS$, que l'on a
\[
\uS(\pi) = \uS_\pi , \qquad \text{sous-catégorie de } \uS \text{ des objets } \pi\text{-accessibles de } S. \tag*{(4.13.1)}
\]
Je dis qu'on aura aussi
\[
\uS(\pi) = \uS \cap \M_\pi . \tag*{(4.13.2)}
\]
Comme c est clair, il reste à voir que $\uS{\cap}\M_\pi \subset \uS(\pi)$, or il est immédiat que $\uS{\cap}\M_\pi \subset S_\pi$, % [sic] c -> an inclusion sign (\supset), a letter c printed where the reverse inclusion seems meant
et on conclut par (4.13.1). On trouve ainsi :

\textbf{Proposition 4.13.4.} \textit{Soit $\M$ une catégorie accessible, $\pi$ un cardinal infini. Soit $\Sigma_\pi$ l'ensemble des}

\msp{page 288}
\textit{sous-catégories strictement pleines $\uS(\pi)$ de $\M_\pi$ qui sont karoubiennes, ou ce qui revient au même, stables dans $\M_\pi$ par facteurs directs. Soit $\Sigma'_\pi$ l'ensemble des sous-catégories de}
\printedpage{107}
% head: A. Grothendieck : Dérivateurs, Ch. XVIII, Édition des Universités de Montpellier II et Paris VII
% foot: Version 17 décembre 2025
\textit{$\M$ qui sont $\pi$-accessibles dans $\M$. Considérons les applications}
\[
\begin{tikzcd}
\Sigma_\pi \arrow[r, shift left, "\alpha_\pi"] & \Sigma'_\pi \arrow[l, shift left, "\beta_\pi"]
\end{tikzcd}
\tag*{(4.13.4.1)}
\]
\textit{définies ainsi : Si $\uS(\pi) \in \Sigma_\pi$, $\alpha_\pi(\uS(\pi))$ est l'image essentielle du foncteur non pleinement fidèle} $\Ind_\pi(\uS(\pi)) \longrightarrow \M$. \textit{Si $\uS \in \Sigma'_\pi$, $\beta_\pi(\uS) = \uS \cap \M_\pi$. Les applications $\alpha_\pi$, $\beta_\pi$ sont des bijections, inverses l'une de l'autre.}

On n'a pas eu à supposer que $\pi$ soit utile pour $\M$, ni que $\pi$ soit minoré par un cardinal utile pour $\M$, i.e. qu'on ait $\pi \geqslant \pi_0$, où $\pi_0$ est le plus petit cardinal utile pour $\M$.

\msp{page 289}
\textbf{4.13.5.} Soit $\uS$ une sous-catégorie \textit{strictement pleine} de $\M$ (catégorie accessible), $\pi'_0$ un cardinal infini. Dire que $\uS$ est $\pi'_0$-accessible dans $\M$ signifie que les conditions suivantes sont satisfaites.
\begin{enumerate}
\item[1)] $\pi'_0$ est utile pour $\uS$, i.e. $\uS \xleftarrow{\sim} \Ind_{\pi'_0}(\uS_{\pi'_0})$.
\item[2)] $\uS \hookrightarrow \M$ est $\pi'_0$-accessible.
\item[3)] $\uS_{\pi'_0} \hookrightarrow \M_{\pi'_0}$ (qui implique que $S_{\pi'_0}$ est équivalente à une petite catégorie).
\end{enumerate}
Supposons ces conditions satisfaites, et soit $\pi$ un cardinal $> \pi'_0$. Est-ce que $\uS$ est $\pi$-accessible dans $\M$, i.e. les conditions similaires sont-elles satisfaites pour $\pi$ ? C'est évident pour 2°), mais risque d'être faux pour 1°, 2° \texttt{[plutôt 3°\,?]}. Mais faisons appel à 4.13.2 c). Posons
\[
\pi'_1 = \sup(\pi'_0, \card\red S_{\pi'_0}, \underbrace{\pi_0}_{\substack{\text{plus petit cardinal}\\ \text{infini utile}\\ \text{pour } \M}} ) .
\]

\msp{page 290}
Alors pour $\pi \geqslant \pi'_1$ tel que $\pi^{\pi'_0} = \pi$, $\uS$ est $\pi$-accessible dans $\M$ (4.13.2 b et c). Supposons p.ex.
\[
\pi'_0 \geqslant \pi_0
\]
(pratiquement ce sera toujours vérifié) et de plus
\[
\pi'_0 \geqslant \card\red\M_{\pi_0} ,
\]
i.e. supposons
\[
\pi'_0 \geqslant \sup(\pi_0, \card\red\M_{\pi_0}) = \pi_1 .
\]
Alors
\[
\pi'_1 = \sup(\pi'_0, \card\red S_{\pi'_0}) \leqslant \sup(\pi'_0, \card\red\M_{\pi'_0}) \leqslant 2^{({\pi'_0}^{\pi_0})}
\]
(4.7.7.1, p. 177). Donc on trouve ceci :

\textbf{Corollaire 4.13.6.} \textit{Soit $\pi'_0$ un cardinal $\geqslant \pi_1 = \sup(\pi_0, \card\red\M_{\pi_0})$ (\footnote[118]{\textbf{N.B.} $\pi_0$ est le plus petit cardinal utile pour $\M$.}), et soit $\uS$ une sous-catégorie $\pi'_0$-accessible dans $\M$. Alors pour tout cardinal $\pi$ tel que $\pi \geqslant 2^{({\pi'_0}^{\pi_0})}$, et tel que l'on ait $\pi^{\pi'_0} = \pi$ (i.e. $\pi$ de la forme $c^{\pi'_0}$, avec $c$ un cardinal $\geqslant 2^{({\pi'_0}^{\pi_0})}$), $\uS$ est $\pi$-accessible dans $\M$. (De plus, ces cardinaux sont réguliers pour $\uS$} et \textit{pour $\M$, donc $\uS \simeq \Ind_\pi(\uS_\pi)$, $\M \simeq \Ind_\pi(\M_\pi)$).}

(\footnote[119]{Enchaîner ici avec 4.13.23 etc.}, \footnote[120]{Mettre ici le lemme 4.16.7.11 signalé en marge p. 292.}).

\printedpage{108}
% head: A. Grothendieck : Dérivateurs, Ch. XVIII, Édition des Universités de Montpellier II et Paris VII
% foot: Version 17 décembre 2025
\msp{page 291}
\textbf{Proposition 4.13.7} (\footnote[121]{Le rejeter après 4.13.23.}). \textit{Soient $\M$ une catégorie accessible, $(S_\alpha)_{\alpha\in A}$ une famille de parties de} $\Ob\M$, \textit{indexée par un petit ensemble $A$, $(\uS_\alpha)_{\alpha\in A}$ la famille des sous-catégories accessibles dans $\M$ correspondante. Soit $S = \bigcap_{\alpha\in A} S_\alpha$, $\uS = \bigcap_{\alpha\in A} \uS_\alpha$, de sorte que $\uS$ est la sous-catégorie pleine de $\M$ définie par $S \subset$} $\Ob\M$. \textit{Ceci posé, si les $S_\alpha$ sont accessibles, il en est de même de $S$ – ou encore, si les $\uS_\alpha$ sont accessibles dans $\M$, il en est de même de $\uS$.}

Cela résulte du résultat plus précis :

\textbf{Corollaire 4.13.8} (\footnote[122]{Cf. reformulation plus raisonnable et plus correcte page 302, 302 bis.}). \textit{Avec les notations de 4.13.7 et 4.13.6, soit, pour tout $\alpha \in A$, $\pi'_\alpha$ un cardinal $\geqslant \pi_1$ tel que $S_\alpha$ soit $\pi'_\alpha$-accessible. Soit}
\[
\begin{aligned}
\pi'_0 &= \sup\nolimits_{\alpha\in A} \pi'_\alpha \\
\pi''_0 &= \sup(\card A, \sup\nolimits_\alpha 2^{({\pi'_\alpha}^{\pi_0})})
\end{aligned}
\]
\textit{et soit $\pi \geqslant \pi''_0$ un cardinal tel que $\pi^{\pi'_0} = \pi$, i.e. de la forme $c^{\pi'_0}$, avec $c \geqslant \pi'_0$. Alors $S$ est $\pi$-accessible (i.e. $\uS$ est $\pi$-accessible dans $\M$). En particulier, $S$ est $2^{\pi''_0}$-accessible.} (\footnote[123]{Ça a l'air un peu abracadabrant ! Il faut de plus supposer que $\pi$ n'est pas cardinal limite !}).

\msp{page 292}
\textsc{Démonstration}. Soit $\alpha \in A$. Comme $\pi'_\alpha \leqslant \pi'_0$ et $\pi^{\pi'_0} = \pi$, on a a fortiori
\[
\pi^{\pi'_\alpha} = \pi , \qquad \text{de plus } \pi \geqslant \pi''_0 \geqslant 2^{(\pi_\alpha^{\pi_0})} , % doubt resolved against PDF; see corrections.csv
\]
donc en vertu de 4.13.6, $\uS_\alpha$ est $\pi$-accessible dans $\M$. D'autre part, comme $\uS_\alpha$ est $\pi'_\alpha$-accessible, il est stable dans $\M$ par $\varinjlim$ grandes devant $\pi'_\alpha$, a fortiori par $\varinjlim$ grandes devant $\pi'_0$ ($\geqslant \pi'_\alpha$). D'ailleurs, $\pi^{\pi'_0} = \pi$ implique $\pi \geqslant 2^{\pi_0} > \pi'_0$. Donc 4.13.8 résulte du

\textbf{Corollaire 4.13.9.} \textit{Soit $\pi > \pi_0$ un cardinal tel que $\uS_\alpha$ soit $\pi$-accessible dans $\M$ pour tout $\alpha \in A$, et tel que} $\card A \leqslant \pi$. \textit{Supposons de plus qu'il existe un cardinal $\pi'_0$ tel que $\pi_0 \leqslant \pi'_0 < \pi$, tel que les $\uS_\alpha$ soient stables dans $\M$ par $\varinjlim$ grandes devant $\pi'_0$. Alors $\uS = \bigcap_{\alpha\in A} \uS_\alpha$ est $\pi$-accessible dans $\M$.}

(\footnote[124]{Ce serait le moment de formuler ici le Lemme 4.16.7.11 (p. 488), qui devrait même remonter avec les premières généralités sur les parties $\pi$-accessibles dès avant 4.13.7 (immédiat après 4.13.6).}).

Posons, pour tout $\alpha \in A$,
\[
\begin{aligned}
\uS_\alpha(\pi) &\overset{\text{déf}}{=} \uS_\alpha \cap \M_\pi \\
\uS(\pi) &\overset{\text{déf}}{=} \uS \cap \M_\pi = \bigcap \uS_\alpha(\pi) .
\end{aligned}
\]
Chaque $\uS_\alpha(\pi)$ est stable dans $\M$ par $\varinjlim$ grandes devant $\pi'_0$ et de cardinal $\leqslant \pi$ (puisqu'il en est ainsi de $\uS_\alpha$ et de $\M_\pi$) (donc il en est de même de leur intersection

\msp{page 293}
$\uS(\pi)$. A fortiori, $\uS(\pi)$ est une sous-catégorie strictement pleine et karoubienne de $\M_\pi$.)


\endgroup
% END UNIT der-ch18/der-ch18_p106-110.tex
% BEGIN UNIT der-ch18/der-ch18_p111-115.tex
\begingroup\makeatletter%
\newcommand{\printedpage}[1]{\marginpar{\footnotesize p.~#1}}%
\newcommand{\msp}[1]{\marginpar{\footnotesize [#1]}}%
\newcommand{\cercle}[1]{\tikz[baseline=(c.base)]\node[draw,circle,inner sep=1pt](c){#1};}%
\newcommand{\incl}{\lhook\joinrel\longrightarrow}%
\newcommand{\card}{\operatorname{card}}%
\newcommand{\Ob}{\operatorname{Ob}}%
\newcommand{\Ind}{\operatorname{Ind}}%
\makeatother\let\printedpage\derpp\let\msp\dermsp%
\let\mspin\DERmspin%


\printedpage{109}
% head: A. Grothendieck : Dérivateurs, Ch. XVIII, Édition des Universités de Montpellier II et Paris VII

D'autre part, chaque $\underline{S}_\alpha$ étant stable dans $\mathcal{M}$ par $\varinjlim$ grandes devant $\pi'_0$, il en est de même de leur intersection $\underline{S}$. A fortiori, celle-ci est stable par $\varinjlim$ grandes devant $\pi$, donc le foncteur pleinement fidèle canonique $\Ind_\pi(\underline{S}(\pi))\longrightarrow\mathcal{M}$ se factorise par $S$ en
\[ \Ind_\pi\underline{S}(\pi) \incl \underline{S} \; , \]
il faut prouver que ce foncteur est essentiellement surjectif. Soit donc $x\in S=\Ob\underline{S}$, tout revient à prouver que
\[ \varinjlim_{i\in\underline{S}(\pi)/x} x_i \xrightarrow{\ \sim\ } x \; , \tag*{$(*)$} \]
et que $\underline{S}(\pi)/x$ est grande devant $\pi$.

Notons que, comme $\pi'_0\geqslant\pi_0$, et $\pi^{\pi'_0}=\pi$, on aura $\pi^{\pi_0}=\pi$, donc $\pi$ est cardinal utile pour $\mathcal{M}$, donc $I=\mathcal{M}_\pi/x$ est grande devant $x$, et on a % [sic] grande devant $x$ -> grande devant $\pi$, apparently intended (one speaks of being large before a cardinal)
\[ \varinjlim_{i\in I=\mathcal{M}_\pi/x} x_i \xrightarrow{\ \sim\ } x \; . \tag*{$(**)$} \]
Donc, si on admet que $\underline{S}(\pi)/x$ est grande devant $\pi$,
\msp{page 294}
le fait que $(*)$ soit un isomorphisme équivaut à dire que le foncteur canonique
\[ J \overset{\text{déf}}{=} \underline{S}(\pi)/x \incl \mathcal{M}_\pi \; = \; I \] % [sic] $\mathcal{M}_\pi = I$ -> $\mathcal{M}_\pi/x = I$, since $I=\mathcal{M}_\pi/x$ above and in the display for $J_\alpha$ below
est \textit{cofinal} (cf. 4.4.2, p. 79). Mais on voit tout de suite que ce foncteur est un foncteur d'inclusion d'une sous-catégorie strictement pleine (stable par $\varinjlim$ grandes devant $\pi'_0$, de cardinal $\leqslant\pi$).

Notons maintenant que la flèche $(**)$ similaire à $(*)$, avec $\underline{S}(\pi)$ remplacée par $\underline{S}_\alpha(\pi)$, est bel et bien un isomorphisme, puisque $x\in\Ob\underline{S}\subset\Ob\underline{S}_\alpha$, et que $\underline{S}_\alpha$ est $\pi$-accessible. Cela s'interprète par le fait que le foncteur d'inclusion
\[ J_\alpha \overset{\text{déf}}{=} \underline{S}_\alpha(\pi)/x \incl \mathcal{M}_\pi/x \; = \; I \]
est cofinal. C'est encore une inclusion d'une sous-catégorie $J_\alpha$ strictement pleine de $I$, stable par $\varinjlim$ grandes devant $\pi'_0$ et de cardinal $\leqslant\pi$, et on a
\[ J \; = \; \bigcap J_\alpha \; . \]

Donc la condition voulue résultera du
\msp{page 295}

% marginal hand-drawn mark (manuscript) in the left margin beside the following lemma
\textbf{Lemme 4.13.10} (\footnote[125]{Il faut supposer que $\pi$ ne soit pas un cardinal limite, pour pouvoir conclure que $I$ est stable par $\varinjlim$ de type $I_{\alpha_\pi}$, et les $J_r$ stables dans $I$ par ces limites.},\footnote[126]{Cf. forme plus générale du lemme dans 4.13.10.2, p. 302.}). \textit{Soient $\pi'_0$ et $\pi>\pi'_0$ des cardinaux infinis, $I$ une catégorie équivalente à une petite catégorie, $(J_\alpha)_{\alpha\in A}$ une famille de sous-catégories $J_\alpha$ de $I$. On suppose :}
\begin{itemize}
\item[\cercle{a}] \textit{$I$ est grande devant $\pi$ (}\footnote[127]{Il suffit que $I$ soit grande devant tout $\pi'<\pi$.}\textit{).} % footnote 127 printed at the foot of p. 110
\item[\cercle{b}] \textit{$I$ est stable par $\varinjlim$ de cardinal $\leqslant\pi$, et grandes devant $\pi'_0$ (}\footnote[128]{Il suffit ici la stabilité par $\varinjlim$ de type $I_{\alpha_\pi}$, cf. page 302 [?].}\textit{).} % footnote 128 printed at the foot of p. 110
\end{itemize}
% foot: Version 17 décembre 2025

\printedpage{110}
% head: A. Grothendieck : Dérivateurs, Ch. XVIII, Édition des Universités de Montpellier II et Paris VII
\begin{itemize}
\item[\cercle{c}] \textit{$\forall\ r\in A$, $J_r$ est stable dans $I$ par les limites inductives envisagés dans b) (en particulier, $J_r$ est une sous-catégorie} strictement pleine \textit{de $I$).} % [sic] envisagés -> envisagées, agreement with « limites »
\item[\cercle{d}] \textit{$\forall\, r\in A$, $J_r$ est une sous-catégorie} cofinale \textit{de $I$ (}grande devant $\pi$\textit{).}
\item[\cercle{e}] $\card A\leqslant\pi$.
\end{itemize}
\textit{Soit $J=\bigcap_{r\in A}J_r$. Alors $J$ est encore cofinale dans $I$, et grande devant $\pi$ (et de plus stable dans $I$ par $\varinjlim$ de cardinal $\leqslant\pi$ et grandes devant $\pi'_0$).}

En d'autres termes, si $I$ est grande devant $\pi$, l'intersection $J$ d'une famille $(J_\alpha)_{\alpha\in A}$ de sous-catégories strictement pleines cofinales et grandes devant $\pi$, est encore
\msp{page 296}
cofinale et grande devant $\pi$, \textit{pourvu que} $\card A\leqslant\pi$, et que l'on postule une propriété de stabilité de $I$ et des $J_\alpha$ dans $I$, pour les limites inductives de cardinal $\leqslant\pi$ et grandes devant un cardinal convenable $\pi'_0<\pi$ donné.

\textbf{4.13.10.1.} Notons d'ailleurs que si $f : J\longrightarrow I$ est un foncteur cofinal \textit{pleinement fidèle}, avec $I$ grande devant $\pi$, alors $J$ est automatiquement grande devant $\pi$. Cela résulte p.ex. de 4.9.6.3 2° (p. 225) par le critère $\mathrm{c}_\pi$ (vu que pour $f$ pleinement fidèle, la condition F $3_\pi$ est tautologiquement satisfaite). On peut aussi le voir ainsi : Soit $B$ un ensemble avec $\card B\leqslant\pi$, il faut prouver que $J\longrightarrow J^B$ est cofinal, or on a
\[
\begin{tikzcd}
J \arrow[r, "\text{cof.\,?}"] \arrow[d, "\text{cof.}"'] & J^B \arrow[d, "\text{cof.}"] \\
I \arrow[r, "\text{cof.}"] & I^B \; ,
\end{tikzcd}
\]
où trois des quatre flèches du carré sont cofinales, $J^B\longrightarrow I^B$ l'étant comme produit de foncteurs cofinaux entre catégories \textit{filtrantes} (cf. 4.9.6.5, page 227) -- on sait en effet déjà que $J$ est filtrante.

\msp{page 297}
Donc le composé
\[ J \longrightarrow J^B \xrightarrow{\text{cof.}} I^B \]
est cofinal. Mais comme le deuxième foncteur $J^B\longrightarrow I^B$ est pleinement fidèle, il s'ensuit aussitôt que $J\longrightarrow J^B$ est cofinal (car $\xi\backslash J$ pour $\xi\in\Ob J^B$ est le même, à isomorphisme près, qu'on considère $\xi$ comme objet de $J^B$, ou de $I^B$).

Ainsi dans l'énoncé de la condition d) du lemme, il est inutile de spécifier que $J_\alpha$ est grande devant $\pi$, et de même dans la conclusion sur $J$.

Comme $J\incl I$ est pleinement fidèle, pour voir que ce foncteur est cofinal, il suffit de voir qu'il satisfait F 1, i.e. que pour tout $i\in\Ob I$, existe $j\in J$ avec $i\longrightarrow j$. On va construire $j$ comme limite d'un système inductif ordinal d'objets $i_\alpha$ de $I$, indexé par les ordinaux $\alpha$ tels que $\card\alpha<\pi$. Désignons par $\alpha_\pi$ le plus petite ordinal tel que % [sic] le plus petite ordinal -> le plus petit ordinal, agreement with « ordinal »
% foot: Version 17 décembre 2025
\printedpage{111}
% head: A. Grothendieck : Dérivateurs, Ch. XVIII, Édition des Universités de Montpellier II et Paris VII
\msp{page 298}
% marginal hand-drawn mark (manuscript) in the left margin beside the following two lines
$\card\alpha_\pi=\pi$. L'ensemble ordonné $I_{\alpha_\pi}$ des ordinaux $\alpha<\alpha_\pi$ étant grande devant $\pi'_0$ (\footnote[129]{faux si $\pi$ est un cardinal limite.}) et de cardinal $\pi$, on pourra donc prendre la limite $\varinjlim\limits_{\alpha\in I_{\alpha_0}} i_\alpha$ dans $I$. % [sic] $I_{\alpha_0}$ -> $I_{\alpha_\pi}$, the index set just defined is $I_{\alpha_\pi}$

Construction de $(i_\alpha)_{\alpha\in I_{\alpha_0}}$. Soit % [sic] $I_{\alpha_0}$ -> $I_{\alpha_\pi}$, the index set just defined is $I_{\alpha_\pi}$
\[ \alpha_\pi \; = \; \text{plus petit ordinal } \alpha \text{ tel que } \card\alpha=\pi, \tag*{(1)} \]
donc c'est l'ordinal caractérisé par
\[ \left\{
\begin{aligned}
&\card\alpha_\pi \; = \; \pi \\
&\alpha<\alpha_\pi \;\; \Longrightarrow \;\; \card\alpha \; < \; \pi \; ,
\end{aligned}
\right. \tag*{(2)} \]
% marginal hand-drawn mark (manuscript) in the left margin beside the following paragraph
ou encore par le fait que l'ensemble ordonné $I_{\alpha_\pi}$ des ordinaux $\alpha<\alpha_\pi$ est \textit{parfaitement bien ordonné}. Ainsi, $I_{\alpha_\pi}$ \textit{est grande devant $\pi'$ pour tout cardinal $\pi'<\pi$} (\footnote[130]{Faux si $\pi$ est cardinal limite.}). Dans les conditions b) c) de 4.3.10, il suffira de supposer $I$ stable par $\varinjlim$ de type $I_{\alpha_\pi}$, et les $J_r$ ($r\in A$) stables dans $I$ par ce type de limites. % [sic] 4.3.10 -> 4.13.10, the conditions b) c) are those of Lemme 4.13.10

\textbf{Lemme 4.13.10.1.} \textit{Soit $\pi$ un cardinal infini, $A$ un ensemble non vide tel que} $\card A\leqslant\pi$. \textit{Alors}
\msp{page 299}
\textit{il existe une application}
\[ \varphi \; : \; I_{\alpha_\pi} \longrightarrow A \]
\textit{telle que pour tout $r\in A$, $\varphi^{-1}(\{r\})$ soit une partie cofinale de $I_{\alpha_\pi}$, ou ce que est plus précis,} $\card\varphi^{-1}(\{r\})=\pi$. % [sic] ce que est -> ce qui est, relative pronoun subject

En effet, il suffit de considérer un ensemble $A'$ de cardinal $\pi$ (p.ex. $A'=I_{\alpha_\pi}$), l'application
\[ \psi \; = \; \mathrm{pr}_1 \; : \; A\times A'\longrightarrow A \; , \]
et de se donner une bijection
\[ A\times A' \xrightarrow{\ \sim\ } I_{\alpha_\pi} \]
laquelle existe puisque
\[ \card\,(A\times A') \; = \; \card A\cdot\card A' \; = \; \card A' \; = \; \pi \; = \; \card I_{\alpha_\pi} \; , \]
puisque $0\neq\card A\leqslant\card A'=\pi$ et $\pi$ infini. On suppose choisie $\varphi$ comme dans 4.13.10.1, et on construira le système inductif $(i_\alpha)_{\alpha\in I_{\alpha_\pi}}$ de telle façon qu'on ait
\[ i_\alpha \; \in \; J_{\varphi(\alpha)} \; . \]

Posant
\[ j \; = \; \varinjlim_{\alpha\in I_{\alpha_\pi}} i_\alpha \; , \]
% foot: Version 17 décembre 2025
\printedpage{112}
% head: A. Grothendieck : Dérivateurs, Ch. XVIII, Édition des Universités de Montpellier II et Paris VII
\msp{page 300}
pour tout $r$, comme
\[ I_r \overset{\text{déf}}{=} \varphi^{-1}(r) \; \subset \; I_{\alpha_\pi} \]
est cofinal dans $I_{\alpha_\pi}$, on aura donc
\[ j \; = \; \varinjlim_{\alpha\in I_r} i_\alpha \; , \]
et en vertu de $(*)$, les $i_\alpha$ pour $\alpha\in I_r$ sont dans $J_r$, donc par passage à la limite, utilisant l'hypothèse c), on trouve que $j\in J_r$, donc
\[ j \; \in \; \bigcap_{r\in A} J_r \; = \; J \; , \]
et comme on choisira
\[ i_0 \; = \; i \; , \tag*{$(**)$} \]
on aura bien trouvé $i\longrightarrow j$ avec $i\in J$. Donc il suffit de construire le système inductif $(i_\alpha)$ satisfaisant $(*)$ et $(**)$. % [sic] avec $i\in J$ -> avec $j\in J$, it is $j$ that was shown to lie in $J$

Si le système inductif est construit jusqu'au cran $i_\alpha\in\Ob J_{\varphi(\alpha)}$, on continue par
\[ i_\alpha \longrightarrow i_{\alpha+1} \; \in \; \Ob J_{\varphi(\alpha+1)} \; , \]
ce qui est possible puisque $J_{\varphi(\alpha+1)}$ est
\msp{page 301}
cofinale. Donc il reste à construire $i_\alpha$ pour $\alpha$ ordinal limite. Or l'ensemble des $i_\beta$ pour $\beta<\alpha$ est de cardinal $<\pi$, et comme par a) $I$ est grande devant $\pi$ (il suffit qu'il le soit devant tout cardinal $<\pi$), il existe un $i\in\Ob I$ qui majore tous les $i_\beta$. Choisissons des
\[ i_\beta \xrightarrow{v_\beta} i \; . \]

On n'aura pas a priori commutativité dans
\[
\begin{tikzcd}
i_{\beta'} \arrow[r, "\tau_{\beta,\beta'}"] \arrow[rr, bend right=15, "v_{\beta'}"'] & i_\beta \arrow[r, "v_\beta"] & i \; ,
\end{tikzcd}
\] % doubt resolved against PDF; see corrections.csv
mais comme l'ensemble des double-flèches $i_{\beta'}\mathrel{\overset{v_\beta\tau_{\beta,\beta'}}{\underset{v_{\beta'}}{\rightrightarrows}}} i$ est de cardinal $<\pi$, $I$ étant grande devant $\pi$, il existe $i\longrightarrow i'$ qui les égalise tous. Comme $J_{\varphi(\alpha)}$ est cofinal dans $I$, il existe $i''\in J_{\varphi(\alpha)}$, et $i\longrightarrow i''$. On prend $i_\alpha=i''$, avec les homomorphismes composés
\[
\begin{tikzcd}
i_\beta \arrow[r, "v_\beta"] \arrow[rrr, dotted, bend right=12, "\tau_{\alpha,\beta}"'] & i \arrow[r] & i' \arrow[r] & i''
\end{tikzcd}
\] % doubt resolved against PDF; see corrections.csv
% foot: Version 17 décembre 2025
\printedpage{113}
% head: A. Grothendieck : Dérivateurs, Ch. XVIII, Édition des Universités de Montpellier II et Paris VII
\msp{page 302}
comme morphismes de transition. Cela achève la démonstration de 4.13.10 et par là de 4.13.9, 4.13.8.

\textbf{Remarque 4.13.10.2.} Voici ce qu'on a eu à utiliser vraiment, pour établir la conclusion de 4.13.10, en termes de $I$, $\pi$, $(J_r)_{r\in A}$ (et sans introduire de $\pi'_0$).
\begin{itemize}
\item[\cercle{a$'$}] $I$ est grande devant tout cardinal $\pi'<\pi$.
\item[\cercle{b$'$}] $I$ est stable par $\varinjlim$ de type $I_{\alpha_\pi}$.
\item[\cercle{c$'$}] $\forall\, r\in A$, $J_r$ est stable dans $I$ par les $\varinjlim$ de type $I_{\alpha_\pi}$.
\item[\cercle{d$'$}] $\forall\, r\in A$, $J_r$ est cofinale dans $I$.
\item[\cercle{e$'$}] $\card A\leqslant\pi$.
\end{itemize}
(\textbf{N.B.} d$'$) et e$'$) ne sont autres que d) et e) de 4.3.10, tandis que a$'$), b$'$), c$'$) sont des versions affaiblies de a), b), c) respectivement.) % [sic] 4.3.10 -> 4.13.10, the conditions d) e) are those of Lemme 4.13.10

\textbf{4.13.8 (reformulation).} \textit{Soit pour $r\in A$, $\pi''_r\geqslant\pi'_r\geqslant\pi_0$ (= cardinal utile pour $\mathcal{M}$ donné) comme spécifié dans 4.13.2 b) (et c), soit $\pi'=\sup_{r\in A}\pi'_r$, $\pi''=\sup_{r\in A}\pi''_r$, de sorte que pour tout cardinal $\pi$ avec $\pi^{\pi'_0}=\pi\geqslant\pi''$, $\underline{S}_r$ soit $\pi$-accessible dans $\mathcal{M}$ (et $\pi$ utile pour $\underline{S}_r$} et \textit{pour $\mathcal{M}$). Soit $c$ un cardinal} non limite \textit{tel que $c\geqslant$} $\card A$ \textit{et $c>\pi'$ (p.ex. le successeur de} $\sup(\card A,\pi')$\textit{). Prenons $\pi$ tel que}
\[ \pi^{\pi'} \; = \; \pi \; \geqslant \; \sup(\pi'',c) \; . \]
\textit{Alors $\underline{S}=\bigcap_{r\in A}\underline{S}_r$ est $\pi$-accessible.}

En effet, si $\alpha_c$ est le premier ordinal de cardinal $c$, comme $c$ n'est pas limite, $I_{\alpha_c}$ est grande devant tout cardinal $c'<c$, en particulier devant $\pi'$. Or pour tout $x\in\Ob\mathcal{M}$, $I=\mathcal{M}_\pi/x$ et les $J_r=\underline{S}_r/x$ sont stables par $\varinjlim$ grandes devant $\pi'$ et de cardinal $\leqslant\pi$,
\msp{page 302 bis}
et les $J_r$ stables dans $I$ pour ces limites, donc il en est en particulier ainsi pour les $\varinjlim$ de type $I_{\alpha_c}$. On peut donc appliquer le résultat 4.13.10.2 pour conclure que si les $J_r$ sont cofinaux dans $I$ (qui est grande devant $\pi$, $\pi$ étant utile pour $\mathcal{M}$), alors $J=\underline{S}/x$ l'est également. Cela implique alors, comme on a vu, que $\underline{S}$ est $\pi$-accessible dans $\mathcal{M}$.

\textbf{4.13.10.3.} On peut reformuler la reformulation précédente, en utilisant la terminologie introduite plus bas (4.13.23) :

\textit{Soit $c$ un cardinal, $\alpha_c$ le plus petit ordinal $\alpha$ tel que} $\card\alpha=c$, \textit{$\pi$ un cardinal $\geqslant c$. On suppose}
\begin{itemize}
\item[$(*)$] \textit{$\pi$ est utile pour $\mathcal{M}$, et $\mathcal{M}$ est stable par $\varinjlim$ de type $I_{\alpha_c}$. (Cette deuxième condition est le cas p.ex. si $c$ est} non limite \textit{et $c>\pi_0$, $\pi_0$ étant un cardinal tel que $\mathcal{M}$ satisfait \textit{2}$_{\pi_0}$, p.ex. un cardinal} utile \textit{pour $\mathcal{M}$ (}\footnote[131]{Car alors $I_{\alpha_c}$ est grande devant $\pi_0$.}\textit{)).}
\end{itemize}
\textit{Alors l'ensemble des sous-catégories $\underline{S}$ de $\mathcal{M}$ qui sont $\alpha_c$-$\pi$-accessible dans $\mathcal{M}$ (i.e. qui sont $\pi$-accessibles et stables dans $\mathcal{M}$ par $\varinjlim$ de type $I_{\alpha_c}$) est stable par intersection de cardinal $\leqslant c$. Plus généralement, si $\Phi$ est un ensemble de catégories telles que} % [sic] accessible -> accessibles, agreement with « sous-catégories »
\begin{itemize}
\item[\textit{1°)}] \textit{$I\in\Phi\Longrightarrow$} $\card I\leqslant\pi$ \textit{et}
\item[\textit{2°)}] \textit{$I_{\alpha_c}\in\Phi$,}
\end{itemize}
\textit{alors l'ensemble des sous-catégories de $\mathcal{M}$ qui sont $\Phi$-$\pi$-accessibles dans $\mathcal{M}$, est stable par intersection de cardinal $\leqslant c$.}
% foot: Version 17 décembre 2025


\endgroup
% END UNIT der-ch18/der-ch18_p111-115.tex
% BEGIN UNIT der-ch18/der-ch18_p116-120.tex
\begingroup\makeatletter%
\newcommand{\printedpage}[1]{\marginpar{\footnotesize p.~#1}}%
\newcommand{\msp}[1]{\marginpar{\footnotesize [#1]}}%
\newcommand{\ilim}{\varinjlim\limits}%
\newcommand{\uS}{\underline{S}}%
\newcommand{\Ob}{\operatorname{Ob}}%
\newcommand{\circlabel}[1]{\tikz[baseline=(c.base)]\node[draw,circle,inner sep=1pt](c){#1};}%
\makeatother\let\printedpage\derpp\let\msp\dermsp%
\let\mspin\DERmspin%


\printedpage{114}
% head: A. Grothendieck : Dérivateurs, Ch. XVIII, Édition des Universités de Montpellier II et Paris VII
% foot: Version 17 décembre 2025

\msp{page 303}\noindent\textbf{Définition 4.13.11} (\footnote[132]{Refondre avec 4.13.23 etc.}). Soient $\mathcal{M}$ une catégorie accessible, $\pi_0$ un cardinal utile pour $\mathcal{M}$, $\pi > \pi_0$ un cardinal. Une catégorie $I$ est dite $\pi_0$-$\pi$-\textit{adaptée} si
\begin{enumerate}
\item[a)] $I$ est équivalente à une petite catégorie,
\item[b)] $I$ est grande devant $\pi$, et
\item[c)] $I$ est stable par $\varinjlim$ de cardinal $\leqslant \pi$, et grandes devant $\pi_0$ (\footnote[133]{\textbf{N.B.} Si $\pi = \pi_0$, la condition c) équivaut à dire que $I$ est stable par facteurs directs.}).
\end{enumerate}
\noindent Un système inductif $(x_i)_{i\in I}$ de type $I$ dans $\mathcal{M}$ est dit $\pi_0$-$\pi$-\textit{adapté}, si $I$ est $\pi_0$-$\pi$-adaptée, i.e. satisfait les conditions a) b) c), et si de plus on a
\begin{enumerate}
\item[d)] le foncteur $i \longmapsto x_i : I \longrightarrow \mathcal{M}$ commute aux $\varinjlim$ (grandes devant $\pi_0$ et de cardinal $\leqslant \pi$) envisagées dans c), et
\item[e)] les $x_i$ sont dans $\mathcal{M}_\pi$.
\end{enumerate}

\noindent\textbf{Proposition 4.13.12.} \textit{Soient $\mathcal{M}$ une catégorie $\pi_0$-accessible, et $\pi > \pi_0$ un cardinal utile pour $\mathcal{M}$. Alors tout $x \in \Ob\mathcal{M}$ admet une représentation comme $x \xleftarrow{\ \sim\ } \ilim_{j\in J} x_j$, où le système inductif $(x_j)_{j\in J}$ dans $\mathcal{M}$ est $\pi_0$-$\pi$-adapté. De façon précise, on peut prendre $J = \mathcal{M}_\pi/x$.}

\msp{page 304}\noindent En effet, la vérification des propriétés a) à d) pour le système inductif canonique, de $\varinjlim$ isomorphe à $x$, est immédiate.

La proposition précédente ne dit pas si on peut prendre pour $J$ un ensemble ordonné. Dans ce cas, la condition c) signifie que pour toute partie $J'$ de $J$ grande devant $\pi_0$ et de cardinal $\leqslant \pi$, $J$ admet une borne supérieure $i_J$ dans $J$, et la condition d) que l'on a alors $x_{i_J} \xleftarrow{\ \sim\ } \ilim_{i\in J'} x_i$.

L'idée naturelle, pour construire un système inductif $\pi_0$-$\pi$-adapté $(x_j)_{j\in J}$, avec $J$ ordonnée, en partant de celui qu'on a canoniquement, avec $J = \mathcal{M}_\pi/x$, est de prendre pour $I$ l'ensemble, ordonné par inclusion, des sous-catégories $J'$ de $J_0$ (sous-catégorie pleine petite de $J$ équivalente à $J$) qui sont de cardinal $\leqslant \pi$ et grandes devant $\pi_0$. (Comparer (4.7.3.3), page 165). Cet ensemble ordonné est bel et bien (petit et) \textit{grand devant} $\pi$, si on suppose que $\pi^{\pi_0} = \pi$ (cf. démonstration 4.7.5, p. 168-171).

\msp{page 305}\noindent\textbf{Lemme 4.13.13.} \textit{Soient $\pi_0$, $\pi > \pi_0$ des cardinaux infinis tels que $\pi^{\pi_0} = \pi$, et $J$ une petite catégorie, grande devant $\pi$. Soit $I$ l'ensemble des sous-catégories $J'$ de $J$ qui sont grandes devant $\pi_0$, et de cardinal $\leqslant \pi$. Alors $I$ est $\pi_0$-$\pi$-adaptée.}

Il reste à voir que pour tout ensemble ordonné $A$ grand devant $\pi_0$ et de cardinal $\leqslant \pi$, si $r \longmapsto J_r$ est une famille croissante d'éléments de $I$ indexée par $A$, alors
\[
J \ \overset{\text{déf}}{=} \ \bigcup_{r\in A} J_r \ = \ \varinjlim_{r\in A} J_r
\]
est dans $I$, i.e. est une sous-catégorie de $I$ de cardinal $\leqslant \pi$ (ce qui est clair) et \textit{grande devant} $\pi_0$. Cela va résulter du

\printedpage{115}
% head: A. Grothendieck : Dérivateurs, Ch. XVIII, Édition des Universités de Montpellier II et Paris VII
% foot: Version 17 décembre 2025

\noindent\textbf{Lemme 4.13.13.1.} \textit{Soit $(J_r)_{r\in A}$ un système inductif de catégories grandes devant $\pi_0$, indexé par un ensemble ordonné $A$ grand devant $\pi_0$. Alors $J = \ilim_{r\in A} J_r$ est grande devant $\pi_0$.}

\msp{page 306}\noindent Il faut voir que si $B$ est un ensemble de cardinal $\leqslant \pi_0$, alors le foncteur diagonal
\[
\delta \ : \ J \longrightarrow J^B \tag*{$(*)$}
\]
est cofinal. Or comme $A$ est grande devant $\pi_0$, $\operatorname{card} B \leqslant \pi_0$, les produits indexés par $B$ commutent aux $\varinjlim$ de type $A$ (SGA 4 I 9.8 appliqué à Cat et à $\pi_0$). Donc la flèche $(*)$ s'interprète comme limite inductive des flèches similaires
\[
\delta_r \ : \ J_r \longrightarrow J_r^B \ , \tag*{$(**)$}
\]
lesquelles sont cofinales, donc $\delta$ est cofinal comme limite inductive de foncteurs cofinaux.

\noindent\textbf{Lemme 4.13.13.2.} \textit{Sous les hypothèses de 4.13.13, soient $\mathcal{M}$ une catégorie satisfaisant $\mathrm{L}_{\pi_0}$, et $(x_j)_{j\in J}$ (où $J \longrightarrow \mathcal{M}$) un système inductif dans $\mathcal{M}$, indexé par $J$. Pour toute $J' \in I$, posons}
\[
x_{J'} \ = \ \varinjlim_{j\in J'} x_j
\]
\msp{page 307}\textit{(qui existe dans $\mathcal{M}$, puisque $\mathcal{M}$ satisfait à $\mathrm{L}_{\pi_0}$). Alors les $(x_{J'})_{J'\in I}$ forment un système inductif dans $\mathcal{M}$, et $I \longrightarrow \mathcal{M}$ commute aux limites inductives grandes devant $\pi_0$ et de cardinal $\leqslant \pi$. Si $\mathcal{M}$ est $\pi_0$-accessible, et les $x_j$ ($j \in \Ob J$) sont dans $\mathcal{M}_\pi$, alors $(x_{J'})_{J'\in I}$ est un système inductif $\pi_0$-$\pi$-adapté dans $\mathcal{M}$.}

Tout ceci est essentiellement trivial, et donne :

\noindent\textbf{Corollaire 4.13.14} (de 4.13.12). \textit{Sous les conditions de 4.13.12, si $\pi^{\pi_0} = \pi$, on peut pour tout $x \in \Ob\mathcal{M}$ trouver une représentation $x \simeq \ilim_{i\in I} x_i$, où $(x_i)_{i\in I}$ est un système inductif $\pi_0$-$\pi$-adapté avec $I$} ensemble ordonné. \textit{De plus, si $\pi' > \pi$ et si $x \in \mathcal{M}_{\pi'}$, on peut prendre $I$ tel qu'on ait}
\[
\operatorname{card} I \ \leqslant \ {\pi'}^{\pi} \ .
\]

Ce dernier complément est immédiat,\msp{page 308} via les majorations habituelles.

\noindent\textbf{Théorème 4.13.15} (\footnote[134]{Cf. version plus générale et plus symétrique du théorème 4.13.15 dans 4.13.22.}). \textit{Soient $\mathcal{M}$ une catégorie $\pi_0$-accessible, et $\pi > \pi'_0 \geqslant \pi_0$ tels que $\pi^{\pi_0} = \pi$} (\footnote[135]{On peut se borner à supposer $\pi$ utile pour $\mathcal{M}$ (au lieu de $\pi^{\pi_0} = \pi$), alors on trouve que a, a$'$, b, c sont équivalentes et impliquent c$'$, mais ne sont peut-être pas équivalentes à c$'$.})\textit{. Soit $\uS$ une sous-catégorie strictement pleine de $\mathcal{M}$, et soit $\uS(\pi) = \uS \cap \mathcal{M}_\pi$. Les conditions suivantes sont équivalentes.}
\begin{enumerate}
\item[\circlabel{a}] \textit{$\uS$ est $\pi$-accessible dans $\mathcal{M}$ (4.13.3) et $\uS$ est stable dans $\mathcal{M}$ par limites inductives grandes devant $\pi'_0$.}
\end{enumerate}

\printedpage{116}
% head: A. Grothendieck : Dérivateurs, Ch. XVIII, Édition des Universités de Montpellier II et Paris VII
% foot: Version 17 décembre 2025

\begin{enumerate}
\item[a$'$)] \textit{$\uS$ est $\pi$-accessible dans $\mathcal{M}$, et $\uS(\pi)$ est stable dans $\mathcal{M}_\pi$ par limites inductives grandes devant $\pi'_0$ et de cardinal $\leqslant \pi$} (\footnote[136]{donner un nom à cette dernière condition, soit $\mathrm{L}(\pi'_0, \pi)$ (pour $\uS(\pi)$ dans $\mathcal{M}_\pi$).})\textit{.}
\item[b)] \textit{Pour tout objet $x$ de $\mathcal{M}$, on a ceci : la sous-catégorie strictement pleine $\uS(\pi)/x$ de $J = \mathcal{M}_\pi/x$ est cofinale dans $\mathcal{M}_\pi/x$ si et seulement si $x \in \Ob\uS$. De plus, $\uS(\pi)$ est stable dans $\mathcal{M}_\pi$ pour $\varinjlim$ grandes devant $\pi'_0$ et de cardinal $\leqslant \pi$.}
\item[\circlabel{c}] \textit{Pour tout objet $x$ de $\mathcal{M}$ et pour toute représentation $x \simeq \ilim_{i\in I} x_i$, où $(x_i)_{i\in I}$ est un système inductif $\pi'_0$-$\pi$-adapté, (4.13.11), on a ceci : la sous-catégorie strictement pleine $I_{\uS}$ de $I$ formée des $i \in \Ob I$ tels que $x_i \in \Ob\uS$, est cofinale dans $I$ si et seulement si $x \in \Ob\uS$. De plus, $S(\pi)$ est stable dans $\mathcal{M}_\pi$ par $\varinjlim$ grandes devant $\pi'_0$ et de cardinal $\leqslant \pi$.} % [sic] S(\pi) -> \underline{S}(\pi), printed without the underline used everywhere else for this subcategory
\end{enumerate}
\msp{page 309}
\begin{enumerate}
\item[c$'$)] \textit{Comme d)} \texttt{[plutôt c)]}\textit{, avec $I$ ensemble ordonné.}
\end{enumerate}
\noindent\textbf{Corollaire 4.13.16.} \textit{Les hypothèses sur $\mathcal{M}$, $\pi_0$, $\uS$ étant celles de 4.13.15, pour qu'il existe des cardinaux $\pi > \pi'_0 \geqslant \pi_0$ satisfaisant les conditions $\pi^{\pi'_0} = \pi$ (a fortiori $\pi^{\pi_0} = \pi$) et les conditions équivalentes a) à d$'$)} \texttt{[plutôt c$'$)]} \textit{de 4.13.15, il (faut et il) suffit que $\uS$ soit accessible dans $\mathcal{M}$. De façon plus précise, si $\uS$ est accessible dans $\mathcal{M}$, soit $\pi'_0 \geqslant \pi_0$ un cardinal tel que $\uS$ soit $\pi'_0$-accessible dans $\mathcal{M}$, et soit $\pi$ un cardinal $> \pi'_0$ tel que $\pi^{\pi'_0} = \pi$. Alors $(\uS, \pi'_0, \pi)$ satisfait aux conditions équivalentes de 4.13.15.}

En effet, le “il faut” est évident sur la condition a) de 4.13.15. Pour le “il suffit” on note que $\uS$ $\pi'_0$-accessible dans $\mathcal{M}$ implique que $\uS$ est stable dans $\mathcal{M}$ par $\varinjlim$ grandes devant $\pi'_0$. Donc pour vérifier a), il suffit de vérifier que $\uS$ est $\pi$-accessible dans $\mathcal{M}$. Or comme $\uS$ est\msp{page 310} $\pi'_0$-accessible et que $\pi > \pi'_0$ satisfait $\pi^{\pi'_0} = \pi$, on en conclut que $\uS$ est $\uS_\pi$-accessible, i.e. $\uS \xleftarrow{\ \sim\ } \mathrm{Ind}_\pi(\uS_\pi)$. Donc en vertu de 4.13.3, il reste à voir que $\uS_\pi \subset \mathcal{M}_\pi$, ou ce qui revient au même, que $\uS_\pi \subset \uS(\pi) \overset{\text{déf}}{=} \uS \cap \mathcal{M}_\pi$. (En fait, on aura même égalité, car il est clair que $\uS(\pi) \subset \uS_\pi \ldots$). Or $\uS(\pi)$ est une sous-catégorie de $\uS$ stable par $\varinjlim$ grandes devant $\pi'_0$ et de cardinal $\leqslant \pi$ (puisque $\uS$ et $\mathcal{M}_\pi$ sont stables par ces limites \textit{dans} $\mathcal{M}$, et que ces limites dans $\uS$ et dans $\mathcal{M}$ sont itou), et $\uS(\pi)$ contient $S_{\pi'_0}$. Donc $\uS(\pi)$ contient $\uS_\pi$, en vertu de 4.7.1 (page 161), implication c) $\Longrightarrow$ b$'$), appliquée à $\uS$, $\pi'_0$, $\pi$ au lieu de $\mathcal{M}$, $\pi_0$, $\pi$. Cela prouve donc le corollaire 4.13.16. % [sic] S_{\pi'_0} -> \underline{S}_{\pi'_0}, printed without the underline

\textsc{Démonstration de} 4.13.15. Voici les implications $\pm$ tautologiques,
\[
\begin{tikzcd}[row sep=small]
 & & & \mathrm{c}' \\
\mathrm{a} \arrow[r, Rightarrow] & \mathrm{a}' \arrow[r, Rightarrow, "\bullet"] & \mathrm{c} \arrow[ur, Rightarrow] \arrow[dr, Rightarrow] & \\
 & & & \mathrm{b}\ ,
\end{tikzcd}
\tag*{$(*)$}
\]
\msp{page 311}que je vais expliciter un peu. On a a $\Longrightarrow$ a$'$ puisque si a est satisfait, $\uS$ \textit{et} $\mathcal{M}_\pi$ sont stables \textit{dans} $\mathcal{M}$ par $\varinjlim$ grandes devant $\pi'_0$ et de cardinal $\leqslant \pi$, donc aussi leur intersection $\uS(\pi)$,
\printedpage{117}
% head: A. Grothendieck : Dérivateurs, Ch. XVIII, Édition des Universités de Montpellier II et Paris VII
% foot: Version 17 décembre 2025
qui est donc aussi stable dans $\mathcal{M}_\pi$ par lesdites limites. L'implication c $\Longrightarrow$ c$'$ est une tautologie, et c $\Longrightarrow$ b aussi, compte tenu de 4.13.12, qui nous dit que le système inductif canonique $(x_i)_{i\in I}$, avec $I = \mathcal{M}_\pi/x$, de limite $x$ est $\pi_0$-$\pi$-adapté. (\textbf{N.B.} Comme $\pi^{\pi'_0} = \pi$ et $\pi'_0 \geqslant \pi_0$, on a aussi $\pi^{\pi_0} = \pi$, donc $\pi$ est utile pour $\mathcal{M}$, de sorte que 4.13.12 s'applique.)

Reste à prouver la “$\pm$ tautologie” a$'$ $\Longrightarrow$ c. Soit donc une représentation $\pi_0$-$\pi$-adaptée
\[
x \ \simeq \ \varinjlim_{i\in I} x_i \ ,
\]
prouvons que
\[
x \ \in \ \Ob\uS \qquad\quad \Longleftrightarrow \qquad\quad I_{\uS} \text{ cofinal dans } I.
\]
L'implication $\Longleftarrow$ provient du fait que, $\uS$ étant $\pi$-accessible dans $\mathcal{M}$, est stable dans $\mathcal{M}$ par $\varinjlim$ grandes devant $\pi$ (\footnote[137]{On n'utilise ici, pour $\Longleftarrow$, de l'hypothèse $(x_i)_{i\in I}$ $\pi_0$-$\pi$-adapté, que le seul fait que $I$ soit grande devant $\pi$.}). Reste à prouver\msp{page 312} que si $x \in \Ob\uS$, alors $I_{\uS}$ est cofinal dans $I$. Considérons donc le diagramme de catégories
\[
\begin{tikzcd}
I \arrow[r, "\varphi"] & J = \mathcal{M}_\pi/x \\
I' = I_{\uS} \arrow[u, hook] \arrow[r, "\varphi'"] & J' = \uS(\pi)/x \ , \arrow[u, hook, "\alpha"']
\end{tikzcd}
\]
c'est un carré cartésien ($I' = I_{\uS}$ est l'image inverse par $\varphi : I \longrightarrow J$ de la sous-catégorie strictement pleine $J'$ de $J$). On a les propriétés suivantes
\begin{enumerate}
\item[1)] $I$ est grande devant $\pi$, et $\varphi$ cofinal ($\Longrightarrow$ $J$ grande devant $\pi$).
\item[2)] $I$ stable par $\varinjlim$ grandes devant $\pi'_0$ et de cardinal $\leqslant \pi$, et $\varphi$ commute à ces $\varinjlim$ .
\item[3)] $J'$ est stable par $\varinjlim$ du type précédent et $\alpha$ y commute (\footnote[138]{C'est ici qu'intervient l'hypothèse de stabilité faite dans b), car d'après b), $\uS$ est $\pi$-accessible dans $\mathcal{M}$.}).
\item[4)] $J'$ cofinal dans $J$.
\end{enumerate}
\noindent On vient de conclure que $I' = \varphi^{-1}(J')$\msp{page 313} est cofinale dans $I$. Comme c'est une sous-catégorie pleine (et même strictement pleine), il suffit de prouver la propriété F 1, i.e. que pour tout $i \in \Ob I$, existe $i' \in \Ob I'$ et une flèche $i \longrightarrow i'$.

On va construire encore $i'$ par construction transfinie, sur l'intervalle ordinal $I_{\alpha_\pi}$, en construisant un système inductif
\[
(i_\alpha)_{\alpha\in I_{\alpha_\pi}} \ , \qquad\quad i_\alpha \in \Ob I \ ,
\]
d'objets de $I$, et un système inductif $(j_\alpha)_{\alpha\in I_{\alpha_\pi}}$, $j_\alpha \in \Ob J'$, flèches de transition notées $\tau_{\beta,\alpha} : i_\alpha \longrightarrow i_\beta$, $\tau'_{\beta,\alpha} : j_\alpha \longrightarrow j_\beta$, enfin pour tout $\alpha$ dans $I_{\alpha_\pi}$ des flèches
\[
\begin{tikzcd}
\varphi(i_\alpha) \arrow[r, "u_\alpha"] \arrow[rr, dotted, bend right, "\varphi(\tau_{\alpha+1,\alpha})"'] & j_\alpha \arrow[r, "v_\alpha"] \arrow[rr, dotted, bend left, "\tau'_{\alpha+1,\alpha}"] & \varphi(i_{\alpha+1})\ ( \arrow[r, dotted, "u_{\alpha+1}"] & j_{\alpha+1}\ )
\end{tikzcd}
\tag*{$(*)$}
\]
\printedpage{118}
% head: A. Grothendieck : Dérivateurs, Ch. XVIII, Édition des Universités de Montpellier II et Paris VII
% foot: Version 17 décembre 2025
dans $J$, telles que
\[
\left\{
\begin{aligned}
v_\alpha u_\alpha \ &= \ \varphi(\tau_{\alpha+1,\alpha}) \\
u_{\alpha+1} v_\alpha \ &= \ \tau'_{\alpha+1,\alpha}
\end{aligned}
\right.
\qquad\qquad \text{(\footnotemark[139])} \ .
\tag*{$(**)$}
\]
\footnotetext[139]{cf. aussi $(****)$, page suivante.}

\msp{page 314}\noindent Montrons comment la construction d'un tel système d'objets et de flèches avec
\[
i_0 \ = \ i \tag*{$(***)$}
\]
va impliquer l'existence de $i' \in \Ob I'$ et de $i \longrightarrow i'$. En effet, comme $I_{\alpha_\pi}$ est grande devant $\pi'_0$ (il est grand devant tout cardinal $< \pi$), il existe
\[
i' \ \overset{\text{déf}}{=} \ \varinjlim_{i\in I_{\alpha_\pi}} i_\alpha % [sic] i\in I_{\alpha_\pi} -> \alpha\in I_{\alpha_\pi}, the index of the limit is \alpha
\]
dans $I$, et on a une flèche de transition
\[
i \ = \ i_0 \longrightarrow i' \ ,
\]
donc il reste à prouver que $i' \in \Ob I'$, i.e. $\varphi(i') \in \Ob J'$. Or $\varphi$ commutant aux $\varinjlim$ de type $I_{\alpha_\pi}$, on aura
\[
\varphi(i) \ = \ \varinjlim \varphi(i_\alpha) \ .
\]
Il me faut encore des compatibilités
\[
\begin{tikzcd}
\varphi(i_\alpha) \arrow[r, "u_\alpha"] \arrow[d, "\varphi(\tau_{\beta,\alpha})"'] & j_\alpha \arrow[d, "\tau'_{\beta,\alpha}"] & {} \arrow[d, phantom, "\text{et}" description] & j_\alpha \arrow[r, "v_\alpha"] \arrow[d, "\tau'_{\beta,\alpha}"'] & \varphi(i_{\alpha+1}) \arrow[d, "\varphi(\tau_{\beta+1,\alpha+1})"] \\
\varphi(i_\beta) \arrow[r, "u_\beta"] & j_\alpha & {} & j_\beta \arrow[r, "v_\beta"] & \varphi(i_{\beta+1})
\end{tikzcd}
\tag*{$(****)$} % [sic] j_\alpha (lower right corner of the first square) -> j_\beta, target of \tau'_{\beta,\alpha} and of u_\beta
\]
oubliées tantôt.

\msp{page 315}\noindent Les compatibilités $(****)$ impliquent qu'on a des homomorphismes
\[
\begin{tikzcd}
\varphi(i) \arrow[r, shift left, "u"] & j \arrow[l, shift left, "v"]
\end{tikzcd}
\qquad\qquad \text{où } j \overset{\text{déf}}{=} \varinjlim_{\alpha\in I_{\alpha_\pi}} j_\alpha \quad \text{(limite \textit{dans} $J'$, donc dans $J$)} \ ,
\]
et les compatibilités $(**)$ impliquent que
\[
vu \ = \ \mathrm{id}_{\varphi(i)} \ , \qquad\quad uv = \mathrm{id}_J \ , % [sic] \mathrm{id}_J -> \mathrm{id}_j, uv is an endomorphism of j
\]
donc que $u$, $v$ sont des isomorphismes inverses l'un de l'autre. Comme $J'$ est une sous-catégorie \textit{strictement} pleine de $J$, il s'ensuit bien que $\varphi(i) \in J'$, i.e. $i \in I'$, q.e.d.

Donc tout revient à construire les
\[
\underbrace{i_\alpha}_{\in \Ob I} \ , \ \tau_{\beta,\alpha} \ , \qquad\quad \underbrace{j_\alpha}_{\in \Ob J'} \ , \ \tau'_{\beta,\alpha} \ , \qquad\quad u_\alpha \ , \ v_\alpha \ ,
\]


\endgroup
% END UNIT der-ch18/der-ch18_p116-120.tex
% BEGIN UNIT der-ch18/der-ch18_p121-125.tex
\begingroup\makeatletter%
\newcommand{\printedpage}[1]{\marginpar{\footnotesize p.~#1}}%
\newcommand{\msp}[1]{\marginpar{\footnotesize [#1]}}%
\newcommand{\circled}[1]{\tikz[baseline=(c.base)]\node[draw,circle,inner sep=1pt](c){#1};}%
\newcommand{\ovalbox}[1]{\tikz[baseline=(c.base)]\node[draw,rounded corners=6pt,inner sep=2pt](c){#1};}%
\newcommand{\Ob}{\operatorname{Ob}}%
\newcommand{\card}{\operatorname{card}}%
\newsavebox{\dctauAB}%
\newsavebox{\dctauABone}%
\newsavebox{\dctauAAzero}%
\makeatother\let\printedpage\derpp\let\msp\dermsp%
\let\mspin\DERmspin%


\sbox{\dctauAB}{\circled{$\scriptstyle\tau_{\alpha,\beta}$}}% boxes for the circled labels of ms. 318 (p. 120)
\sbox{\dctauABone}{\circled{$\scriptstyle\tau_{\alpha,\beta+1}$}}
\sbox{\dctauAAzero}{\circled{$\scriptstyle\tau_{\alpha,\alpha_0}$}}
\printedpage{119}
% head: A. Grothendieck : Dérivateurs, Ch. XVIII, Édition des Universités de Montpellier II et Paris VII
satisfaisant aux contraintes $(**)$, $(***)$, $(****)$. On va faire la construction en trois étapes.

\circled{1} Départ :
\[
i_0 \;=\; i\ .
\]
Comme $J'$ est cofinal dans $J$, existe $j_0\in\Ob J'$ et
\[
\varphi(i_0)\xrightarrow{u_0} j_0\ .
\]

\msp{page 316}

\circled{2} Supposons que les $i_\beta$, $j_\beta$ soient construits pour $\beta<\alpha$, ainsi que les $\tau_{\beta',\beta}:i_\beta\longrightarrow i_{\beta'}$, $\tau'_{\beta',\beta}:j_\beta\longrightarrow j_{\beta'}$ pour $\beta,\beta'<\alpha$, enfin les $u_\beta:\varphi(i_\beta)\longrightarrow j_\beta$, $v_\beta:j_\beta\longrightarrow\varphi(i_{\beta+1})$ pour $\beta<\alpha$ resp. pour $\beta+1<\alpha$, de façon à satisfaire aux contraintes $(**)$, $(****)$, i.e. commutativités dans
\[
\left\{
\begin{array}{l}
\left.
\begin{array}{ll}
\circled{1} &
\begin{tikzcd}[ampersand replacement=\&]
\varphi(i_\beta) \arrow[r, "u_\beta"] \arrow[rr, bend right=15, "\varphi(\tau_{\beta+1,\beta})"'] \& j_\beta \arrow[r, "v_\beta"] \& \varphi(i_{\beta+1})
\end{tikzcd}
\\[5ex]
\circled{2} &
\begin{tikzcd}[ampersand replacement=\&]
j_\beta \arrow[r, "v_\beta"] \arrow[rr, dotted, bend right=15, "\tau'_{\beta+1,\beta}"'] \& \varphi(i_{\beta+1}) \arrow[r, "u_{\beta+1}"] \& j_{\beta+1}
\end{tikzcd}
\end{array}
\right\}\quad \textit{si}\ \beta+1<\alpha
\\[8ex]
\begin{array}{ll}
\circled{3} &
\begin{tikzcd}[ampersand replacement=\&]
\varphi(i_\beta) \arrow[r, "u_\beta"] \arrow[d, "\varphi(\tau_{\beta',\beta})"'] \& j_\beta \arrow[d, "\tau'_{\beta',\beta}"] \\
\varphi(i_{\beta'}) \arrow[r, "u_{\beta'}"] \& j_{\beta'}
\end{tikzcd}
\qquad\qquad \text{si } \beta<\beta'<\alpha
\end{array}
\\[10ex]
\begin{array}{ll}
\circled{4} &
\begin{tikzcd}[ampersand replacement=\&]
j_\beta \arrow[r, "v_\beta"] \arrow[d, "\tau'_{\beta',\beta}"'] \& \varphi(i_{\beta+1}) \arrow[d, "\varphi(\tau_{\beta'+1,\beta+1})"] \\
j_{\beta'} \arrow[r, "v_{\beta'}"] \& \varphi(i_{\beta'+1})
\end{tikzcd}
\qquad\qquad \text{si } \beta<\beta' \text{ et } \beta'+1<\alpha\ ,
\end{array}
\end{array}
\right.
\]

\msp{page 317}

il faut construire
% foot: Version 17 décembre 2025

\printedpage{120}
% head: A. Grothendieck : Dérivateurs, Ch. XVIII, Édition des Universités de Montpellier II et Paris VII
\begin{itemize}[leftmargin=*]
\item[1)] $i_\alpha\in\Ob I$,\quad et un homomorphisme $(i_\beta)_{\beta\in I_\alpha}\longrightarrow i_\alpha$, i.e. des $\tau_{\alpha,\beta}$ : $i_\beta\longrightarrow i_\alpha$ $(\beta\in I_\alpha)$, compatibles aux translations $\tau_{\beta',\beta}$ pour $\beta<\beta'<\alpha$.
\item[2)] $j_\alpha\in\Ob J'$,\quad et un homomorphisme $(j_\beta)_{\beta\in I_\alpha}\longrightarrow j_\alpha$, i.e. des $\tau'_{\alpha,\beta}$ : $j_\beta\longrightarrow j_\alpha$ $(\beta\in I_\alpha)$, compatibles aux translations $\tau'_{\beta',\beta}$ pour $\beta<\beta'<\alpha$.
\item[3)] Une flèche
\[
u_\alpha:\varphi(i_\alpha)\longrightarrow j_\alpha\ .
\]
\item[4)] Si $\alpha$ n'est pas ordinal limite, i.e. $\alpha=\alpha_0+1$, une flèche
\[
v_{\alpha_0}:j_{\alpha_0}\longrightarrow\varphi(i_{\alpha_0+1})=\varphi(i_\alpha)\ ;
\]
\end{itemize}

et ceci de façon à satisfaire aux contraintes \circled{1}\,\circled{2}\,\circled{3}\,\circled{4} quand on y suppose $\beta+1\leqslant\alpha$ (au lieu de $\beta'<\alpha$) dans \circled{1} et \circled{2}, et $\beta'\leqslant\alpha$ (au lieu de $\beta'<\alpha$) dans \circled{3}, et $\beta'+1\leqslant\alpha$ (au lieu de $\beta'+1<\alpha$) dans \circled{4}. Donc pour les contraintes $1^\circ$, $2^\circ$, $4^\circ$, il n'y a en tenir compte \textit{que} si $\alpha$ n'est pas ordinal limite, donc $\alpha=\alpha_0+1$, et en y faisant % [sic] il n'y a en tenir -> il n'y a à en tenir, the preposition « à » appears to be missing
\[
\beta' \;=\; \alpha_0\ .
\]
Quant à la contrainte $3^\circ$, il faut en
\msp{page 318}
tenir compte pour $\beta'=\alpha$. Donc au total, il faut exiger
\[
(3_\alpha)(\beta)\qquad\qquad
\begin{tikzcd}
\varphi(i_\beta) \arrow[r, "u_\beta"] \arrow[d, "\varphi(\usebox{\dctauAB})"'] & j_\beta \arrow[d, "\tau'_{\alpha,\beta}" {draw, circle, inner sep=1pt}] \\
\varphi(i_\alpha) \arrow[r, "u_\alpha" {draw, circle, inner sep=1pt}] & j_\beta % [sic] j_\beta -> j_\alpha, the target of \tau'_{\alpha,\beta} and u_\alpha is j_\alpha
\end{tikzcd}
\qquad \text{commutatif pour } \beta<\alpha\ (\footnotemark[140]),
\]
\footnotetext[140]{\textbf{N.B.} Il suffit de l'exiger pour $\beta=\alpha_0$, dans le cas $\alpha=\alpha_0+1$, cf. ci-dessous, p. 323.}
(par restriction sur $\alpha$), et si $\alpha=\alpha_0+1$ est ordinal non limite, il faut exiger de plus
\[
(4_\alpha)(\beta)\qquad\qquad
\begin{tikzcd}
j_\beta \arrow[r, "v_\beta"] \arrow[d, "\tau'_{\alpha_0,\beta}"'] & \varphi(i_{\beta+1}) \arrow[d, "\varphi(\usebox{\dctauABone})"] \\
j_{\alpha_0} \arrow[r, "v_{\alpha_0}" {draw, circle, inner sep=1pt}] & \varphi(i_{\alpha_0+1})=\varphi(i_\alpha)
\end{tikzcd}
\qquad \beta<\alpha_0
\]
\[
(1_\alpha)\qquad\qquad
\begin{tikzcd}
\varphi(i_{\alpha_0}) \arrow[r, "u_{\alpha_0}"] \arrow[rr, bend right=15, "\varphi(\usebox{\dctauAAzero})"'] & j_{\alpha_0} \arrow[r, "v_{\alpha_0}" {draw, circle, inner sep=1pt}] & \varphi(i_{\alpha_0+1})=\varphi(i_\alpha)
\end{tikzcd}
\]
% foot: Version 17 décembre 2025

\printedpage{121}
% head: A. Grothendieck : Dérivateurs, Ch. XVIII, Édition des Universités de Montpellier II et Paris VII
\[
(2_\alpha)\qquad\qquad
\begin{tikzcd}
j_{\alpha_0} \arrow[r, "v_{\alpha_0}" {draw, circle, inner sep=1pt}] \arrow[rr, bend right=15, "\tau'_{\alpha,\alpha_0}"' {draw, circle, inner sep=1pt}] & \varphi(i_{\alpha_0+1})=\varphi(i_\alpha) \arrow[r, "u_\alpha" {draw, circle, inner sep=1pt}] & j_\alpha\ .
\end{tikzcd}
\]

Dans ces compatibilités, j'ai cerclé les flèches $\tau_{\alpha,\beta}$, $\tau'_{\alpha,\beta}$ $(\beta<\alpha)$, $\Bigl(\ \tau_{\alpha,\beta+1}\ (\beta+1<\alpha)\ \Bigr)$, $v_{\alpha_0}$, $\Bigl(\ \tau_{\alpha,\alpha_0}\ \Bigr)$, $u_\alpha$, $\Bigl(\ \tau'_{\alpha,\alpha_0}\ \Bigr)$, qui ont été introduites dans les donnés 1) 2) 3) 4), page 317. % [sic] donnés -> données, the noun « données » is meant

\msp{page 319}

On va distinguer donc les deux cas.

\textit{2 A)} $\alpha=\alpha_0+1$.

Pour construire $i_\alpha$ et les $\tau_{\alpha,\beta}:i_\beta\longrightarrow i_\alpha$, il suffit de construire $i_\alpha$ et
\[
\circled{$\tau_{\alpha,\alpha_0}$}\ :\ i_{\alpha_0}\ \longrightarrow\ i_\alpha=i_{\alpha_0+1}\ .
\]
De même pour $j_\alpha$ et les $\tau'_{\alpha,\beta}$, il suffit de construire
\[
\circled{$\tau'_{\alpha,\alpha_0}$}\ :\ j_{\alpha_0}\ \longrightarrow\ j_\alpha\ .
\]
Il faut de plus se donner
\[
j_{\alpha_0}\ \xrightarrow{\circled{$v_{\alpha_0}$}}\ \varphi(\underbrace{\,i_\alpha\,}_{=i_{\alpha_0+1}})\ \xrightarrow{\circled{$u_\alpha$}} j_\alpha\ ,
\]
donc en tout deux objets $i_\alpha$, $j_\alpha$ (dans $I$ et dans $J'$ respectivement) et quatre flèches $\tau_{\alpha,\alpha_0}$, $\tau'_{\alpha,\alpha_0}$, $v_{\alpha_0}$, $u_\alpha$. Il faut quatre compatibilités (page 318), où les deux dernières $3_\alpha(\beta)$, $4_\alpha(\beta)$ dépendent d'un paramètre $\beta<\alpha$, i.e. $\beta\leqslant\alpha_0$, ou de $\beta<\alpha_0$.

\msp{page 320 biffé}

\msp{page 321}

On va d'abord s'occuper de deux premières compatibilités $1_\alpha$, $2_\alpha$, en termes de la donnée de
\[
i_{\alpha_0}\ \in\ \Ob I\ ,\qquad j_{\alpha_0}\ \in\ \Ob J'\ ,\qquad u_{\alpha_0}\ :\ \varphi(i_{\alpha_0})\ \longrightarrow\ j_{\alpha_0}\ .
\]
Les deux compatibilités voulues pour les deux objets $i_\alpha$, $j_\alpha$ et les quatre flèches $\tau_{\alpha,\alpha_0}$, $\tau'_{\alpha,\alpha_0}$, $v_{\alpha_0}$, $u_\alpha$ sont indiquées dans le diagramme
\[
(*)\qquad\qquad
\begin{tikzcd}
\varphi(i_{\alpha_0}) \arrow[r, "u_{\alpha_0}"] \arrow[rr, bend right=15, "\varphi(\tau_{\alpha,\alpha_0})"'] & j_{\alpha_0} \arrow[r, "v_{\alpha_0}\ ?"] \arrow[rr, bend left=20, "\tau'_{\alpha,\alpha_0}"] & \varphi(\underbrace{\,i_\alpha\,}_{?}) \arrow[r, "u_\alpha\ ?"] & j_\alpha\ .
\end{tikzcd}
\]

On note que comme $\varphi$ est cofinal, $\exists\ i_\alpha=i\in\Ob I$ et $j_{\alpha_0}\xrightarrow{v_{\alpha_0}}\varphi(i)$. Mais on veut que le composé
\[
w\ :\ \varphi(i_{\alpha_0})\ \xrightarrow{u_{\alpha_0}}\ j_{\alpha_0}\ \xrightarrow{v_{\alpha_0}}\ \varphi(i)
\]
soit de la forme $\varphi(\tau)$, où $\tau=\tau_{\alpha,\alpha_0}:i_{\alpha_0}\longrightarrow i_\alpha$ est une flèche dans $I$. Comme $\varphi$ est un foncteur cofinal de catégories filtrantes, on sait
% foot: Version 17 décembre 2025
\printedpage{122}
% head: A. Grothendieck : Dérivateurs, Ch. XVIII, Édition des Universités de Montpellier II et Paris VII
\msp{page 322}
qu'on peut trouver $i\xrightarrow{\rho}i'$ dans $I$, tel que la flèche composée $\varphi(i_{\alpha_0})\xrightarrow{w}\varphi(i)\xrightarrow{\varphi(\rho)}\varphi(i')$ soit de la forme $\varphi(\tau)$. Donc quitte à remplacer $i$ par $i'$, et $v_{\alpha_0}$ par le composé $v'_{\alpha_0}$ : $j_{\alpha_0}\xrightarrow{v_{\alpha_0}}\varphi(i)\xrightarrow{\varphi(\rho)}\varphi(i')$, on peut supposer que $v_{\alpha_0}u_{\alpha_0}$ est de la forme $\varphi(\tau)$. Cela nous donne donc en principe l'objet $i=i_\alpha$, et les flèches $v_{\alpha_0}:j_{\alpha_0}\longrightarrow\varphi(i)$ et $\tau_{\alpha,\alpha_0}:i_{\alpha_0}\longrightarrow i_\alpha$, de façon à satisfaire à la première compatibilité dans $(*)$. Il reste à construire $j_\alpha$, $\tau'_{\alpha,\alpha_0}$, $u_\alpha$ de façon à satisfaire à la deuxième compatibilité $(*)$. Or $J'$ étant cofinal dans $J$, il existe $j_\alpha$ dans $J'$ et une flèche $u_\alpha:\varphi(i_\alpha)\longrightarrow j_\alpha$. D'autre part, comme $J'\longrightarrow J$ est pleinement fidèle, la flèche composée $j_{\alpha_0}\xrightarrow{v_{\alpha_0}}\varphi(i_\alpha)\xrightarrow{u_\alpha}j_\alpha$ dans $J$, entre objets de $J'$, provient d'une unique flèche $\tau'_{\alpha,\alpha_0}:j_{\alpha_0}\longrightarrow j_\alpha$ dans $J'$.

\msp{page 323}

Ainsi, on a construit les deux objets et les quatre flèches voulues, de façon à satisfaire aux contraintes $1_\alpha$, $2_\alpha$ de la page 318. Les contraintes $3_\alpha(\beta)$, $4_\alpha(\beta)$ sont-elles alors satisfaites automatiquement ?
\[
\begin{tikzcd}
\varphi(i_\beta) \arrow[r, "u_\beta"] \arrow[d] & j_\beta \arrow[d] \\
\varphi(i_{\alpha_0}) \arrow[r, "u_{\alpha_0}"] \arrow[d, "\varphi(\tau_{\alpha,\alpha_0})"'] & j_{\alpha_0} \arrow[d, "\tau'_{\alpha,\alpha_0}"] \\
\varphi(i_\alpha) \arrow[r, "u_\alpha"] & j_\alpha
\end{tikzcd}
\qquad\qquad \beta<\alpha, \text{ i.e. } \beta\leqslant\alpha_0\ .
\]
Le carré supérieur est commutatif par hypothèse de récurrence, pour que le carré total (composé) soit commutatif, il suffit que le carré inférieur le soit. On a donc une double-flèche $(\tau'_{\alpha,\alpha_0}u_{\alpha_0},u_\alpha\varphi(\tau_{\alpha,\alpha_0})):\varphi(i_\alpha)\rightrightarrows j_\alpha$, elle n'est peut-être pas triviale, mais on sait, $J'$ étant filtrante, qu'on peut l'égaliser par une flèche $j_\alpha\longrightarrow\bar{j}_\alpha$ dans $J'$. Donc quitte à remplacer $j_\alpha$ par % [sic] \varphi(i_\alpha) -> \varphi(i_{\alpha_0}), both composites of the double arrow start at \varphi(i_{\alpha_0})
\msp{page 324}
$\bar{j}_\alpha$, $u_\alpha$ par le composé $\bar{u}_\alpha:\varphi(i_0)\longrightarrow j_\alpha\longrightarrow\bar{j}_\alpha$ et $\tau'_{\alpha,\alpha_0}$ par le composé $\bar{\tau}'_{\alpha,\alpha_0}:j_{\alpha_0}\longrightarrow j_\alpha\longrightarrow\bar{j}_\alpha$ (sans toucher à $i_\alpha$, $v_{\alpha_0}$, $\tau_{\alpha_0,\alpha}$), on peut obtenir la commutativité, d'où $3_\alpha(\beta)$. % [sic] \varphi(i_0) -> \varphi(i_\alpha), u_\alpha has source \varphi(i_\alpha); [sic] \tau_{\alpha_0,\alpha} -> \tau_{\alpha,\alpha_0}, the arrow is written \tau_{\alpha,\alpha_0} everywhere else

Reste la contrainte $4_\alpha(\beta)$. Pour ceci, considérons le diagramme composite
\[
\begin{tikzcd}
j_\beta \arrow[r, "v_\beta"] \arrow[d, "\tau'_{\alpha_0,\beta}"'] & \varphi(i_{\beta+1}) \arrow[d, "\varphi(\tau_{\alpha,\beta+1})"] \\
j_{\alpha_0} \arrow[r, "v_{\alpha_0}"] & \varphi(i_{\alpha_0+1})=\varphi(i_\alpha)\ .
\end{tikzcd}
\]

Si $\alpha_0$ n'est pas un ordinal limite, mais de la forme $\alpha_1+1$, on peut encore se ramener au cas où $\beta=\alpha_1$, en procédant comme ci-dessus. Mais on peut procéder dans tous les cas, sans un tel artifice (\,!), en considérant les double-flèches
\[
(*)\qquad\qquad
\begin{tikzcd}
j_\beta \arrow[r, shift left, "\varphi(\tau_{\alpha,\beta+1})v_\beta"] \arrow[r, shift right, "v_{\alpha_0}\tau'_{\alpha_0,\beta}"'] & \varphi(i_\alpha)
\end{tikzcd}
\]
% foot: Version 17 décembre 2025
\printedpage{123}
% head: A. Grothendieck : Dérivateurs, Ch. XVIII, Édition des Universités de Montpellier II et Paris VII
déduites du carré précédent. Comme
\msp{page 325}
l'ensemble de ces double-flèches est de cardinal $\card\alpha<\pi$, et que $J$ est grande devant $\pi$ (il suffirait qu'il le soit pour les cardinaux $\pi'<\pi$), on peut trouver une flèche $\varphi(i_\alpha)\xrightarrow{u}j$ dans $J$ qui égalise toutes ces double-flèches. Comme $\varphi$ est cofinal, on peut supposer que $j$ est de la forme $\varphi(i)$. Comme $\varphi$ est un foncteur cofinal de catégories filtrantes, il existe alors $i\xrightarrow{\rho}i'$ dans $I$, tel que le composé $\varphi(i_\alpha)\xrightarrow{u}\varphi(i)\xrightarrow{\varphi(\rho)}\varphi(i')$ soit de la forme $\varphi(\sigma)$, pour $\sigma:i_\alpha\longrightarrow i'$. Au total, on a trouvé $\sigma:i_\alpha\longrightarrow i'$, tel que $\varphi(\sigma)$ $(=\varphi(\rho)u):\varphi(i_\alpha)\longrightarrow\varphi(i')$ égalise toutes les double-flèches envisagées dans $(*)$, p. 324. Donc il s'impose de remplacer $i_\alpha$ par $\bar{i}_\alpha=i'$, $v_{\alpha_0}$ par le composé $j_{\alpha_0}\longrightarrow\varphi(i_\alpha)\xrightarrow{\varphi(\sigma)}\varphi(\bar{i}_\alpha)$, $\tau_{\alpha,\alpha_0}$
\msp{page 326}
par le composé $\bar{\tau}_{\alpha,\alpha_0}:i_{\alpha_0}\xrightarrow{\tau_{\alpha,\alpha_0}}i_\alpha\xrightarrow{\sigma}\bar{i}_\alpha$.

Pour résumer, voici la recette pour construire, d'abord \ovalbox{$i_\alpha$, $v_{\alpha_0}$, $\tau_{\alpha,\alpha_0}$,}et \textit{ensuite} \ovalbox{$j_\alpha$, $u_\alpha$, $\tau'_{\alpha,\alpha_0}$}.

On procède d'abord comme page 321, 322 pour construire le premier triple, de façon à satisfaire à la première compatibilité $(*)$ (p. 321), i.e. à $1_\alpha$. On “corrige” ensuite ce choix de $i_\alpha$, $v_{\alpha_0}$, $\tau_{\alpha,\alpha_0}$, en remplaçant $i_\alpha$ par $\bar{i}_\alpha$ pour $\sigma:i_\alpha\longrightarrow\bar{i}_\alpha$, $v_{\alpha_0}$ par $\bar{v}_{\alpha_0}=\varphi(\sigma)v_{\alpha_0}$, et $\tau_{\alpha,\alpha_0}$ par $\bar{\tau}_{\alpha,\alpha_0}=\sigma\tau_{\alpha,\alpha_0}$, de sorte à obtenir les compatibilités $4_\alpha(\beta)$ pour $\beta<\alpha_0$, comme on vient de l'indiquer. On ne touchera plus à ces choix de $i_\alpha$, $v_{\alpha_0}$, $\tau_{\alpha,\alpha_0}$, dans la suite de la construction. On construit $(j_\alpha,u_\alpha,\tau'_{\alpha,\alpha_0})$ de façon à satisfaire à la deuxième compatibilité $(*)$ (p. 321), comme
\msp{page 327}
indiqué page 322. Ensuite, on corrige ce choix, en remplaçant $j_\alpha$ par $\bar{j}_\alpha$, avec $\rho:j_\alpha\longrightarrow\bar{j}_\alpha$, et $u_\alpha$ par $\bar{u}_\alpha=\rho u_\alpha$, $\tau'_{\alpha,\alpha_0}$ par $\bar{\tau}'_{\alpha,\alpha_0}=\rho\tau'_{\alpha,\alpha_0}$, de façon à satisfaire aux contraintes $3_\alpha(\beta)$ de la page 318.

Cela achève la construction dans le cas où $\alpha=\alpha_0+1$.

\textit{Cas 2 B : $\alpha$ est ordinal limite.}

Il faut donc construire les données 1) 2) 3) de la page 317 (à l'exclusion de la donnée 4), de façon à satisfaire à la seule contrainte $3_\alpha(\beta)$.

On commence à construire $i_\alpha$ et les $\tau_{\beta,\alpha}$ dans $I$. C'est possible, car $I_\alpha=\{\beta\mid\beta<\alpha\}$ est de cardinal $\pi_\alpha<\pi$ – il suffit donc de supposer $I$ grande devant les cardinaux $\pi'<\pi$.

On construit ensuite de même $J_\alpha$ dans $J'$ et les $\tau'_{\beta,\alpha}$ $(\beta<\alpha)$, en utilisant le % [sic] J_\alpha -> j_\alpha, the object of J' is meant
\msp{page 328}
fait que $J'$ est grande devant $\card\alpha=\card I_\alpha$. On doit de plus construire une flèche
\[
u_\alpha\ :\ \varphi(i_\alpha)\ \longrightarrow\ j_\alpha\ .
\]
C'est possible, car $J'$ étant cofinale dans $J$, il existe $j'$ dans $J'$ et $\varphi(i_\alpha)\longrightarrow J'$, et $J'$ étant filtrante, il existe $\bar{j}_\alpha$ et des flèches % [sic] \varphi(i_\alpha)\longrightarrow J' -> \varphi(i_\alpha)\longrightarrow j', the arrow goes to the object j'
$\begin{tikzcd}[ampersand replacement=\&, baseline=(\tikzcdmatrixname-1-1.base)]
j_\alpha \arrow[rr] \& \& j_\alpha \\
\& j' \arrow[ur] \&
\end{tikzcd}$, donc quitte à remplacer $j_\alpha$ par $\bar{j}_\alpha$, les $\tau'_{\alpha,\beta}$ par des $\bar{\tau}'_{\alpha,\beta}$, on trouve bien une flèche $u_\alpha:\varphi(i_\alpha)\longrightarrow j_\alpha$. Il faut ensuite s'occuper des % [sic] target j_\alpha in the small diagram -> \bar{j}_\alpha, the two arrows go to the object \bar{j}_\alpha just introduced
% foot: Version 17 décembre 2025


\endgroup
% END UNIT der-ch18/der-ch18_p121-125.tex
% BEGIN UNIT der-ch18/der-ch18_p126-130.tex
\begingroup\makeatletter%
\newcommand{\printedpage}[1]{\marginpar{\footnotesize p.~#1}}%
\newcommand{\msp}[1]{\marginpar{\footnotesize [#1]}}%
\DeclareRobustCommand{\card}{\operatorname{card}}%
\DeclareRobustCommand{\Ob}{\operatorname{Ob}}%
\DeclareRobustCommand{\Ind}{\operatorname{Ind}}%
\makeatother\let\printedpage\derpp\let\msp\dermsp%
\let\mspin\DERmspin%


\printedpage{124}
% head: A. Grothendieck : Dérivateurs, Ch. XVIII, Édition des Universités de Montpellier II et Paris VII

compatibilités $3_\alpha(\beta)$, donc des double-flèches
\[
\begin{tikzcd}
\varphi(i_\beta) \arrow[r, shift left, "\tau'_{\alpha,\beta}u_\beta"] \arrow[r, shift right, "u_\alpha\varphi(\tau_{\alpha,\beta})"'] & j_\alpha \; .
\end{tikzcd}
\]
Comme l'ensemble de ces double-flèches est de cardinal $\card\alpha < \pi$, et que $J'$ est grande devant ledit cardinal, on
\msp{page 329}
trouve donc une flèche
\[
\rho : j_\alpha \longrightarrow j
\]
dans $J$ qui les égalise toutes. Comme $J'$ est cofinal, on peut supposer $j$ dans $J'$, soit $\bar{j}_\alpha$. Donc quitte à remplacer $j_\alpha$ par $\bar{j}_\alpha$, et les $\tau'_{\alpha,\beta}$ par $\bar{\tau}'_{\alpha,\beta} = \rho\tau'_{\alpha,\beta}$, enfin $u_\alpha : \varphi(i_\alpha) \longrightarrow j_\alpha$ par $\bar{u}_\alpha = \rho u_\alpha : \varphi(i_\alpha) \longrightarrow \bar{j}_\alpha$, on satisfait à toutes les contraintes voulues.

Je résume le résultat (``$\pm$ tautologique''!) prouvé dans les pages 313-329.

% a vertical change bar is printed in the left margin alongside the statement from "Lemme-clef" to "suivantes :"
\textbf{Lemme-clef 4.13.17.} \textit{Considérons un carré cartésien dans} Cat
\[
\begin{tikzcd}
I \arrow[r, "\varphi"] & J \\
I' \arrow[u, hook] \arrow[r] & J' \arrow[u, hook, "\psi"'] \; .
\end{tikzcd}
\]

\textit{On fait les hypothèses suivantes :}

\msp{page 330}

\begin{enumerate}[label={}]
\item[\textit{1)}] \textit{$\varphi$ et $\psi$ sont cofinaux, $\psi$ est l'inclusion d'une sous-catégorie strictement pleine $J'$ de $J$, $I$ est filtrant (donc $J$ et $J'$ aussi).}
\end{enumerate}

\textit{On cherche des conditions pour que l'inclusion de la sous-catégorie strictement pleine $I' \longrightarrow I$ soit cofinale (auquel cas on voit tout de suite que $I'$ est filtrante, et $I' \longrightarrow J'$ est aussi cofinal). Pour ceci, on suppose qu'il existe un \textup{ordinal limite} $\alpha_0$ ayant les propriétés supplémentaires suivantes :}

\begin{enumerate}[label={}]
\item[\textit{2)}] \textit{$I$ et $J'$ sont stables par $\varinjlim$ de type $I_{\alpha_0}$ $(= \{\beta \mid \beta \ \textit{ordinal}\ < \alpha_0\})$, et les foncteurs $\varphi$, $\psi$ commutent à ces limites.}
\item[\textit{3)}] \textit{Pour tout $\alpha \in I_{\alpha_0}$, $I$ est grande devant $\card\alpha$. (Donc, $\varphi$ étant cofinal, $J$ est également grande devant $\card\alpha$, et de même, $\psi$ étant pleinement fidèle, et cofinal, $J'$ est grande devant $\card\alpha$.)}
\end{enumerate}

\textit{Sous ces conditions, $I'$ est cofinal dans $I$ (donc aussi grande devant $\card\alpha$, pour tout $\alpha \in I_{\alpha_0}$).}

La démonstration est celle donnée ci-dessus. On n'a pas eu à utiliser que l'ordinal $\alpha_0 = \alpha_\pi$ était de cette espèce particulière, seulement que c'était un cardinal limite.

\printedpage{125}
% head: A. Grothendieck : Dérivateurs, Ch. XVIII, Édition des Universités de Montpellier II et Paris VII

\msp{page 331}

Ce lemme implique, comme on l'a vu à la page 312, l'implication b $\Longrightarrow$ d dans le diagramme d'implications $(*)$ page 310, lequel est donc à présent prouvé.

Pour achever de prouver le théorème 4.13.15, il reste à prouver les implications
\[
\begin{tikzcd}[row sep=small]
\mathrm{c}' \arrow[dr, Rightarrow] & & \\
& \mathrm{a}' \arrow[r, Rightarrow] & \mathrm{a} \; . \\
\mathrm{b} \arrow[ur, Rightarrow] & &
\end{tikzcd}
\]

Démontrons donc (b ou c$'$) $\Longrightarrow$ a$'$. Il reste à prouver que $\underline{S}$ est $\pi$-accessible dans $\mathcal{M}$, i.e. que le foncteur pleinement fidèle
\[
\Ind_\pi(\underline{S}(\pi)) \longrightarrow \mathcal{M} \tag*{$(*)$}
\]
a pour image essentielle $\underline{S}$. Donc il faut prouver que si $x \in \Ob\mathcal{M}$, alors
\[
x \in \Ob\underline{S} \quad \overset{(?)}{\Longleftrightarrow} \quad
\begin{minipage}[c]{0.55\textwidth}
$x$ est dans l'image essentielle de $\Ind_\pi(\underline{S}(\pi))$, i.e. $x$ est isomorphe à $\varinjlim\limits_{i\in I} x_i$, les $x_i$ dans $\Ob\underline{S}(\pi)$, $I$ grande devant $\pi$.
\end{minipage}
\tag*{$(*)$}
\]

\msp{page 332}

$\Bigg($ Supposons d'abord que $\underline{S}$ satisfasse b) (\footnote[141]{Maladroit, on peut traiter les deux cas b) et c) en même temps, cf. bas de page 333.}). Donc on sait que
\[
x \in \Ob\underline{S} \qquad \Longleftrightarrow \qquad \underline{S}/x \text{ est cofinale dans } \mathcal{M}_\pi/x \; . \tag*{$(**)$}
\]

Nous voulons déduire l'équivalence $(*)$ de l'équivalence $(**)$.

Prouvons $\Longrightarrow$ dans $(*)$, i.e. supposons $x \in \Ob\underline{S}$, et prouvons que $x$ est dans l'image essentielle de $\Ind_\pi\underline{S}(\pi)$. En vertu de $(**)$, $\underline{S}(\pi)/x$ est une sous-catégorie (strictement pleine) cofinale dans $\mathcal{M}_\pi/x$, donc comme $\mathcal{M}_\pi/x$ est grande devant $\pi$, $\underline{S}(\pi)/x$ l'est aussi. Comme on sait que
\[
x \ = \ \varinjlim_{i\in\mathcal{M}_\pi/x} x_i
\]
($\pi$ étant cardinal utile pour $\mathcal{M}$), il s'ensuit par cofinalité qu'on a aussi
\[
x \ = \ \varinjlim_{i\in\underline{S}(\pi)/x} x_i \; ,
\]
et comme $\underline{S}(\pi)/x$ est grande devant $\pi$, cela implique que $x$ appartient à l'image essentielle de $\Ind_\pi(\underline{S}(\pi))$.

Prouvons $\Longleftarrow$ dans $(*)$, i.e. supposons que $x$ soit dans l'image essentielle de $\Ind_\pi(\underline{S}(\pi))$,
\printedpage{126}
% head: A. Grothendieck : Dérivateurs, Ch. XVIII, Édition des Universités de Montpellier II et Paris VII
\msp{page 333}
\[
x \ = \ \varinjlim_{i\in I} x_i \; , \qquad\quad \text{les } x_i \text{ dans } \underline{S}(\pi),\ I \text{ grande devant } \pi.
\]

Alors le foncteur composé $i \longmapsto x_i$,
\[
I \ \longrightarrow \ \underline{S}(\pi)/x \ \hookrightarrow \ \mathcal{M}_\pi/x
\]
est cofinal (cf. 4.4.2, page 79), donc $\underline{S}(\pi)/x \longrightarrow \mathcal{M}_\pi/x$ satisfait la condition F 1 (puisque le composé le satisfait), et comme $\underline{S}(\pi)/x \longrightarrow \mathcal{M}_\pi/x$ est pleinement fidèle, cela implique déjà que ce foncteur est cofinal. Donc par $(**)$ on en conclut $x \in \Ob\underline{S}$, q.e.d.

Supposons maintenant que $S$ satisfasse c$'$ (au lieu de b), et prouvons encore l'équivalence $(*)$. Pour ceci, écrivons
\[
\square \qquad x \ = \ \varinjlim_{i\in I} x_i \qquad\quad
\left\{
\begin{minipage}[c]{0.5\textwidth}
les $x_i \in \Ob\mathcal{M}_\pi$, $I$ \textit{ensemble ordonné}, $(x_i)_{i\in I}$ est $\pi_0$-$\pi$-adaptée (\footnotemark[142]).
\end{minipage}
\right.
\]
\footnotetext[142]{\textbf{N.B.} Ça existe, grâce à 4.13.14.}

ce qui est possible puisque $\pi$ est utile pour $\mathcal{M}$ $\Bigg)$. Définissons $I_{\underline{S}}$ comme dans l'énoncé de c), donc comme image inverse de $\underline{S}(\pi)/x \hookrightarrow \mathcal{M}_\pi/x$ par $I \longrightarrow \mathcal{M}_\pi$. Par hypothèse d$'$) \texttt{[plutôt c$'$)]}, on a
\[
x \in \Ob\underline{S} \qquad \Longleftrightarrow \qquad I_{\underline{S}} \text{ est cofinal dans } I \; . \tag*{$(**')$}
\]

En utilisant $(**')$, prouvons encore $(*)$ (page 331). En fait, on utilisera seulement qu'on a
\msp{page 334}
l'isomorphisme $\square$, avec $I$ grande devant $\pi$, et qu'on a l'équivalence $(**')$, sans avoir à utiliser que $I$ est ordonnée. Donc la démonstration vaut aussi dans le cas c).

Prouvons $\Longrightarrow$ dans $(*)$, i.e. supposons $x \in \Ob\underline{S}$, prouvons que $x$ est dans l'image essentielle de $\Ind_\pi(\underline{S}(\pi))$. En vertu de $(**')$, $I_{\underline{S}}$ est cofinale dans $I$, et comme $I_{\underline{S}} \longrightarrow I$ est une inclusion pleinement fidèle, il s'ensuit que $I_{\underline{S}}$ est également grande devant $\pi$, d'autre part on aura par cofinalité
\[
x \ = \ \varinjlim_{i\in I_{\underline{S}}} x_i \; ,
\]
où cette fois les $x_i$ sont dans $\underline{S}(\pi) = \underline{S}\cap\mathcal{M}_\pi$, d'où résulte que $x$ est dans l'image essentielle de $\Ind_\pi(\underline{S}(\pi))$.

Prouvons $\Longleftarrow$ dans $(*)$, i.e. supposons que l'on ait
\[
x \ = \ \varinjlim_{j\in J} x_j \; , \qquad\quad \text{les } x_j \text{ dans } \underline{S}(\pi),\ J \text{ grande devant } \pi \; .
\]

\printedpage{127}
% head: A. Grothendieck : Dérivateurs, Ch. XVIII, Édition des Universités de Montpellier II et Paris VII

\msp{page 335}

En invoquant à nouveau 4.4.2, on trouve que $J \longrightarrow \mathcal{M}_\pi/x$ est cofinal, et comme ce foncteur se factorise par le foncteur d'inclusion pleinement fidèle $\underline{S}(\pi)/x \hookrightarrow \mathcal{M}_\pi/x$, il s'ensuit comme tantôt que ce dernier est cofinal. Nous allons admettre pour un moment que ceci implique que l'image inverse $I_{\underline{S}}$ de cette sous-catégorie strictement pleine, par le foncteur également cofinal $I \longrightarrow \mathcal{M}_\pi/x$, est cofinal dans $I$. Ceci étant, l'équivalence $(**')$ (p. 333) implique $x \in \Ob\underline{S}$, ce qu'on voulait.

Donc il reste à prouver que si
\[
J' \ = \ \underline{S}(\pi)/x \hookrightarrow \mathcal{M}_\pi/x \ = \ J
\]
est cofinal, alors $I' = I_{\underline{S}} \longrightarrow I$ l'est aussi. Il semble qu'il faille bel et bien à nouveau invoquer les conditions de stabilité $\mathrm{L}(\pi_0,\pi)$ faites dans d$'$) \texttt{[plutôt c$'$)]} ce qui n'avait pas été nécessaire pour le cas c) et
\msp{page 336}
le lemme-clef (un peu) vache 4.13.17, où on prendra (comme dans la démonstration écrite
avant l'énoncé de cet lemme) % [sic] cet lemme -> ce lemme, demonstrative adjective "cet" before a consonant
$\alpha = \alpha_\pi \overset{\text{déf}}{=}$ plus petit ordinal $\alpha$ tel que $\card\alpha = \pi$. Donc $I_\alpha$ est de cardinal $\pi$, et grande devant tout cardinal $\pi' < \pi$, et en particulier devant $\pi'_0$. L'hypothèse $2^\circ$) de 4.13.17 est vérifiée pour $I$, par définition d'un système inductif $\pi_0$-$\pi$-adapté, et pour $J'$, grâce à l'hypothèse $\mathrm{L}(\pi_0,\pi)$ sur $\underline{S}(\pi)$. De même, l'hypothèse $3^\circ$) ($I$ grande devant tout cardinal $< \pi$) est satisfaite par définition des systèmes $\pi_0$-$\pi$-adaptés -- $I$ est même grande devant $\pi$. L'hypothèse $1^\circ$) a déjà été vue.

On a donc prouvé (par voie circulaire)
\[
\begin{tikzcd}[row sep=small]
& & & \mathrm{c}' & \\
\mathrm{a} \arrow[r, Rightarrow] & \Bigg(\mathrm{a}' \arrow[r, Leftrightarrow] & \mathrm{c} \arrow[ur, Leftrightarrow] \arrow[dr, Rightarrow] & & \Bigg)\, , \\ % Verified against PDF leaf 9: the oblique equivalence has both arrowheads.
& & & \mathrm{b} &
\end{tikzcd}
\]

\texttt{[et aussi b $\Longrightarrow$ c,]} en ayant à invoquer le lemme-clef vache pour les \textit{deux} implications
\[
\mathrm{a}' \ \Longrightarrow \ \mathrm{c} \qquad\quad \text{et} \qquad\quad \mathrm{c}' \ \Longrightarrow \ \mathrm{a}'
\]
(donc pour deux implications \textit{en sens inverses}, chose étrange \dots).

\msp{page 337}

Il reste seulement encore à prouver l'implication
\[
\mathrm{a}' \ \Longrightarrow \ \mathrm{a} \; ,
\]
ce qui est essentiellement sorital.

Cela achève la démonstration du critère d'accessibilité 4.13.15 d'une sous-catégorie pleine $\underline{S}$ de $\mathcal{M}$.

On va donner une variante un peu décantée de 4.13.15, qui (avec des hypothèses un peu plus faibles) capte l'essentiel de sa substance. L'énoncé 4.13.15 en est une conséquence quasi immédiate :

\printedpage{128}
% head: A. Grothendieck : Dérivateurs, Ch. XVIII, Édition des Universités de Montpellier II et Paris VII

\msp{page 338}

\textbf{Corollaire 4.13.18} (de 4.13.15). \textit{Soient $\mathcal{M}$ une catégorie accessible, $\pi$ un cardinal utile pour $\mathcal{M}$, $\alpha_0$ un \textup{ordinal limite} tel que l'on ait $\card\alpha_0 \leqslant \pi$. Soit $I_{\alpha_0} = \{\alpha \mid \alpha < \alpha_0\}$, et supposons que $\mathcal{M}$ soit stable par $\varinjlim$ de type $I_{\alpha_0}$.}

\textit{Soit alors $\underline{S}$ une sous-catégorie strictement pleine de $\mathcal{M}$, $\underline{S}(\pi) = \underline{S}\cap\mathcal{M}_\pi$. Désignons par $L(\alpha_0)$ la condition : $\underline{S}(\pi)$ est stable dans $\mathcal{M}_\pi$ (ou ce qui revient au même, dans $\mathcal{M}$) par $\varinjlim$ de type $I_{\alpha_0}$. Alors les conditions suivantes sont équivalentes.}

\begin{enumerate}[label={}]
\item[\textit{a)}] \textit{$\underline{S}$ est $\pi$-accessible dans $\mathcal{M}$, et est stable dans $\mathcal{M}$ par $\varinjlim$ de type $I_{\alpha_0}$.}
\item[\textit{a$'$)}] \textit{$\underline{S}$ est $\pi$-accessible et $\underline{S}(\pi)$ satisfait $L(\alpha_0)$.}
\item[\textit{b)}] \textit{$\underline{S}(\pi)$ satisfait $L(\alpha_0)$, et pour tout $x \in \Ob\mathcal{M}$, on a ceci : $x \in \Ob\underline{S}$ si et seulement si la sous-catégorie strictement pleine $\underline{S}(\pi)/x$ de $\mathcal{M}(\pi)/x$ est cofinale dans $\mathcal{M}(\pi)/x$} \texttt{[lire $\mathcal{M}_\pi/x$]}.
\item[\textit{c)}] \textit{$\underline{S}(\pi)$ satisfait $L(\alpha_0)$, de plus pour tout $x \in \Ob\mathcal{M}$, et toute représentation}
\[
x \ \simeq \ \varinjlim_{i\in I} x_i
\]
\msp{page 339}
\textit{qui est ``$\alpha_0$-$\pi$-adaptée'', i.e. telle que}
\begin{enumerate}[label={}]
\item[\textit{$1^\circ$)}] \textit{$I$ équivalente à une petite catégorie, et grande devant $\pi$,}
\item[\textit{$2^\circ$)}] \textit{les $\varinjlim$ de type $I_{\alpha_0}$ existent dans $I$,}
\item[\textit{$3^\circ$)}] \textit{le foncteur $I \longrightarrow \mathcal{M}$, $i \longmapsto x_i$, commute aux $\varinjlim$ de type $I_{\alpha_0}$, et enfin}
\item[\textit{$4^\circ$)}] \textit{les $x_i \in \Ob\mathcal{M}_\pi$}
\end{enumerate}
\textit{(comparer 4.13.11), on a ceci (où on désigne par $I_{\underline{S}}$ la sous-catégorie strictement pleine de $I$ formée des $i$ tels que $x_i \in \Ob\underline{S}$) : on a $x \in \Ob\underline{S}$ si et seulement si $I_{\underline{S}}$ est cofinale dans $I$.}
\item[\textit{c$'$)}] \textit{(\textup{Si} $I_{\alpha_0}$ est grande devant $\pi_0$, et $\pi^{\pi_0} = \pi$, pour pouvoir invoquer 4.13.14, cf. p. 335.) Comme c), mais en supposant $I$ ordonnée.}
\end{enumerate}
(\footnote[143]{\textbf{N.B.} Si on ne suppose par $I_{\alpha_0}$ grande devant $\pi_0$ et % [sic] suppose par -> suppose pas, negation "ne ... pas"
$\pi^{\pi_0} = \pi$, alors toutes les conditions de a) à c) (a, a$'$, b, c) sont équivalentes et \textit{impliquent} c, mais ne lui sont peut-être pas équivalentes.}).

\textsc{Démonstration}. Elle se fait suivant le même diagramme logique que celle de 4.13.15,
\[
\begin{tikzcd}[row sep=small]
& & & \mathrm{c}' \arrow[dr, Rightarrow, "\textbf{!}"] & \nearrow\ \begin{tabular}[b]{@{}r@{}}\footnotesize ici, il faut supposer\\ \footnotesize $I_{\alpha_0}$ grande devant $\pi_0$\\ \footnotesize et $\pi^{\pi_0} = \pi$.\end{tabular} & \\
\mathrm{a} \arrow[r, Rightarrow] & \mathrm{a}' \arrow[r, Rightarrow, "\textbf{!}"] & \mathrm{c} \arrow[ur, Rightarrow] \arrow[dr, Rightarrow] & & \mathrm{a}' \arrow[r, Rightarrow] & \mathrm{a} \; , \\
& & & \mathrm{b} \arrow[ur, Rightarrow] & &
\end{tikzcd}
\]
où on a marqué du signe\,! les deux implications qui sont ``vaches'', en ce sens qu'elles font appel au lemme-clef 4.13.17 (pages 329-330), qui se prouve par une construction transfinie un peu vissée. La démonstration des sept implications du diagramme est essentiellement
celle donné dans les pages précédentes 311-337, % [sic] celle donné -> celle donnée, agreement with feminine "celle"


\endgroup
% END UNIT der-ch18/der-ch18_p126-130.tex
% BEGIN UNIT der-ch18/der-ch18_p131-135.tex
\begingroup\makeatletter%
\newcommand{\printedpage}[1]{\marginpar{\footnotesize p.~#1}}%
\newcommand{\msp}[1]{\marginpar{\footnotesize [#1]}}%
\newcommand{\Ind}{\operatorname{Ind}}%
\newcommand{\Ob}{\operatorname{Ob}}%
\newcommand{\Ord}{\operatorname{Ord}}%
\newcommand{\card}{\operatorname{card}}%
\newcommand{\cardred}{\mathop{\textup{card\,réd}}\nolimits}%
\newcommand{\cercle}[1]{\tikz[baseline=(c.base)]\node[draw,circle,inner sep=1pt](c){\textup{#1}};}%
\makeatother\let\printedpage\derpp\let\msp\dermsp%
\let\mspin\DERmspin%

\setcounter{footnote}{143}

\printedpage{129}
% head: A. Grothendieck : Dérivateurs, Ch. XVIII, Édition des Universités de Montpellier II et Paris VII
% foot: Version 17 décembre 2025

\msp{page 340}
à l'exception de a$'$ $\Longrightarrow$ a, qui est un peu plus simple dans le cas actuel. On va la dégager sous forme de lemme un peu plus général. (\textbf{N.B.} La démonstration de a$'$ $\Longrightarrow$ a dans 4.13.15 avait été omise comme “soritale”\,!, cf. page 337\,\dots)

\textbf{Lemme 4.13.19.} \textit{Soient $\mathcal{M}$ une catégorie accessible dans $\mathcal{M}$, $\pi$ un cardinal utile pour $\mathcal{M}$ et tel que $\underline{S}$ soit $\pi$-accessible, i.e. $\underline{S}$ est l'image essentielle de $\Ind_\pi(\underline{S}(\pi))$ par le foncteur pleinement fidèle canonique $\Ind_\pi(\underline{S}(\pi))\longrightarrow\mathcal{M}$, où $\underline{S}(\pi)=\underline{S}\cap\mathcal{M}_\pi$. On suppose qu'il existe un ordinal limite $\alpha$, tel que $\card\alpha\leqslant\pi$ et tel que $\mathcal{M}_\pi$ et $\underline{S}(\pi)$ $(=\underline{S}_\pi)$ soient stables par $\varinjlim$ de type $I_\alpha$ soient stables par $\varinjlim$ de type $I_\alpha$, et que le foncteur d'inclusion $\underline{S}(\pi)\longrightarrow\mathcal{M}_\pi$ y commute. Soit $\Phi$ un ensemble de catégories telles que $\forall\ I\in\Phi$, on ait $\cardred I\leqslant\pi$, et supposons que $\mathcal{M}_\pi$ soit stable par $\varinjlim$ de type $\Phi$, et que $\underline{S}(\pi)$ soit stable dans $\mathcal{M}_\pi$ par lesdites limites (ou ce qui revient au même, qu'il soit stable en lui-même par ces $\varinjlim$\,, et que le foncteur d'inclusion y commute). Alors $\mathcal{M}$ et $\underline{S}$ sont stables par $\varinjlim$ de type $\Phi$, et le foncteur d'inclusion $\underline{S}\longrightarrow\mathcal{M}$ y commute, i.e. $\mathcal{M}$ est stable par ces limites, et $\underline{S}$ stable dans $\mathcal{M}$ par ces mêmes limites.} % [sic] « soient stables par lim de type I_alpha » printed twice -> once, repetition

\msp{page 341}
\textsc{Démonstration}. On applique 4.11.1 (p. 267-268) séparément à $\mathcal{M}\simeq\Ind_\pi\mathcal{M}_\pi$ et à $\underline{S}\simeq\Ind_\pi\underline{S}(\pi)$, pour conclure que $\mathcal{M}$ et $\underline{S}$ sont stables séparément par $\varinjlim$ de type $\Phi$. Il reste à voir que le foncteur $\underline{S}\longrightarrow\mathcal{M}$ y commute. Cela se voit comme dans la démonstration de 4.11.1, on obtient plus généralement ceci :

\textbf{Lemme 4.13.20} (cor. à 4.11.1). \textit{Soient $f_0 : C\longrightarrow C'$ un foncteur entre catégories équivalentes à des petites catégories, $\pi$ un cardinal infini, $f:\mathcal{M}\stackrel{\textup{déf}}{=}\Ind_\pi C\longrightarrow\mathcal{M}'\stackrel{\textup{déf}}{=}\Ind_\pi C'$ l'extension canonique $\Ind_\pi(f_0)$ de $f_0$. Soient $\Phi$ et $\alpha$ comme dans 4.11.1, et supposons que $I_\alpha\in\Phi$ et que $f_0$ commute aux $\varinjlim$ de type $\Phi$. Alors $f$ commute à ces mêmes limites inductives.} (\footnote{Il vaut mieux prendre des foncteurs $f_0 : C\longrightarrow\mathcal{N}$, et $\mathcal{N}$ [?] que $f_0$ commute aux $\varinjlim$ de type $\Phi$, donc il en est de même de son extension $\Ind_\pi(C,\mathcal{N})\longrightarrow\mathcal{N}$. (\textbf{N.B.} On suppose que $\mathcal{N}$ satisfait $\mathrm{L}_\pi$.)}).

\textbf{Corollaire 4.13.21.} \textit{Sous les conditions de 4.13.20, supposons que $\Phi$ contienne les catégories de cardinal $\leqslant\pi$ (resp. les catégories filtrantes de cardinal $\leqslant\pi$, resp. les catégories filtrantes grandes devant $\pi_0$ de cardinal $\leqslant\pi$, resp. les catégories discrètes de cardinal $\leqslant\pi$). Alors $f$ commute aux petites $\varinjlim$ (resp. aux petites $\varinjlim$ filtrantes, resp. aux petites $\varinjlim$ filtrantes grandes devant $\pi_0$, resp. aux petites sommes). (Réciproque claire et “laissée au lecteur”.)}

Se voit comme 4.4.14 et 4.4.14.1 (p. 97, 98).

\msp{page 342}
Voici enfin, il me semble, la forme définitive (\,?) de 4.13.15, 4.13.18.

\textbf{Théorème 4.13.22} (\footnote{cf. variante 4.13.36, page 358.}). \textit{Soient $\mathcal{M}$ une catégorie accessible, $\pi$ un cardinal infini, utile pour $\mathcal{M}$, et $\Phi$ un ensemble de catégories satisfaisant les deux conditions :}
\begin{itemize}[leftmargin=2.5em]
\item[\textit{1°)}] $J\in\Phi$ \textit{implique} $\cardred J\leqslant\pi$.
\item[\textit{2°)}] \textit{Il existe un} ordinal limite $\alpha_0$ \textit{avec} $I_{\alpha_0}\in\Phi$.
\end{itemize}
\textit{On suppose $\mathcal{M}$ stable par $\varinjlim$ de type $\Phi$. Soient $\underline{S}$ une sous-catégorie} strictement pleine \textit{de $\mathcal{M}$, et $\underline{S}(\pi)=\underline{S}\cap\mathcal{M}_\pi$. Désignons par $L(\varphi)$ la condition sur $\underline{S}(\pi)$ : $\underline{S}(\pi)$ est stable dans} % [sic] L(\varphi) -> L(\Phi), the condition is later written L(\Phi)
\printedpage{130}
% head: A. Grothendieck : Dérivateurs, Ch. XVIII, Édition des Universités de Montpellier II et Paris VII
% foot: Version 17 décembre 2025
\textit{$\mathcal{M}$ (où ce qui revient au même, dans $\mathcal{M}_\pi$) par $\varinjlim$ de type $\Phi$. Ceci posé, les conditions suivantes sur $\underline{S}$ sont équivalentes :} % [sic] où ce qui revient -> ou ce qui revient, conjunction "ou"
\begin{itemize}[leftmargin=2.5em]
\item[\cercle{a}] \textit{$\underline{S}$ est $\pi$-accessible dans $\mathcal{M}$, et stable dans $\mathcal{M}$ par $\varinjlim$ de type $\Phi$.}
\item[a$'$)] \textit{$\underline{S}$ est $\pi$-accessible dans $\mathcal{M}$, et $\underline{S}(\pi)$ satisfait à $L(\Phi)$.}
\end{itemize}
\msp{page 343}
\begin{itemize}[leftmargin=2.5em]
\item[b)] \textit{$\underline{S}(\pi)$ satisfait à $L(\Phi)$, et pour tout $x\in\Ob\mathcal{M}$, on a ceci : $x\in\Ob\underline{S}$ si et seulement si la sous-catégorie strictement pleine $\underline{S}(\pi)/x$ de $\mathcal{M}(\pi)/x$ est cofinale dans $\mathcal{M}(\pi)/x$} \texttt{[lire} $\mathcal{M}_\pi/x$\texttt{]}.
\item[\cercle{c}] \textit{$\underline{S}(\pi)$ satisfait à $L(\Phi)$, et pour tout $x\in\Ob\mathcal{M}$ et toute représentation $x\simeq\varinjlim\limits_{i\in I}x_i$ qui est $\Phi$-$\pi$-adaptée (cf. définition plus bas), on a ceci : $x\in\Ob\underline{S}$ si et seulement si $I_{\underline{S}}$ est cofinale dans $I$, où $I_{\underline{S}}$ est la sous-catégorie strictement pleine de $I$ formée des $i\in\Ob I$ tels que $x_i\in\Ob\underline{S}$ (ou encore, $x_i\in\Ob\underline{S}(\pi)$).} \textbf{N.B.} \textit{On dit qu'un système inductif $(x_i)_{i\in I}$ dans $\mathcal{M}$ est $\Phi$-$\pi$-adapté si}
\begin{itemize}[leftmargin=2em]
\item[\textit{1)}] \textit{$I$ équivalente à une petite catégorie, et $I$ grande devant $\pi$,}
\item[\textit{2)}] \textit{les $\varinjlim$ de type $\Phi$ existent dans $I$,}
\item[\textit{3)}] \textit{le foncteur $i\longmapsto x_i$ de $I$ dans $\mathcal{M}$ commute aux $\varinjlim$ de type $\Phi$, et enfin}
\item[\textit{4)}] \textit{les $x_i\in\Ob\mathcal{M}_\pi$.}
\end{itemize}
\item[c$'$)] \textit{(Si $\pi^{\pi_0}=\pi$, où $\pi_0$ est un cardinal $\leqslant\pi$ utile pour $\mathcal{M}$, et si $J\in\Phi$ implique que $J$ est grande devant $\pi_0$.) Comme c), avec $I$} ordonnée.
\end{itemize}
(\footnote{\textbf{N.B.} Quand les conditions restrictives ci-contre \texttt{[parenthèse dans c$'$)]} ne sont pas satisfaites, alors a, a$'$, b, c sont encore équivalentes, et impliquent c$'$, mais ne sont peut-être pas équivalentes à c$'$).},\footnote{\textbf{N.B.} Il faudrait réénoncer les conditions annexes [?] (Si \dots), cf. p. 359-360.}).

\textbf{Définition 4.13.23} (\footnote{\textbf{N.B.} Il y aurait intérêt à remonter cette définition (\textit{sans} b), dès avant 4.13.15.}). Soient $\mathcal{M}$ une catégorie accessible, $\pi$ un cardinal, $\Phi$ un ensemble de catégories, satisfaisant les deux conditions :
\begin{itemize}[leftmargin=2.5em]
\item[a)] $J\in\Phi\Longrightarrow\cardred J\leqslant\pi$.
\item[b)] Il existe un ordinal limite $\alpha_0$, tel que $I_{\alpha_0}\in\Phi$ (\footnote{Si c'est pour poser une définition, on peut se passer de la condition b).}).
\end{itemize}
\msp{page 344}
Soit $S$ une partie de $\Ob\mathcal{M}$, $\underline{S}$ la sous-catégorie pleine qu'elle définit. On dit que $S$ est une \textit{partie $\Phi$-$\pi$-accessible de} $\Ob\mathcal{M}$, ou que $\underline{S}$ est une \textit{sous-catégorie $\Phi$-$\pi$-accessible dans} $\mathcal{M}$, si les conditions suivantes sont satisfaites.
\begin{itemize}[leftmargin=2.5em]
\item[1°)] $\pi$ est utile pour $\mathcal{M}$, $\mathcal{M}$ est stable par $\varinjlim$ de type $\Phi$ (\footnote{Condition préliminaire sur $\mathcal{M}$.}).
\item[2°)] $\underline{S}$ est une sous-catégorie $\pi$-accessible dans $\mathcal{M}$.
\item[3°)] $\underline{S}$ est stable dans $\mathcal{M}$ par $\varinjlim$ de type $\Phi$.
\end{itemize}
(\footnote{\textbf{N.B.} Quand $\Phi=\{I_{\alpha_0}\}$, on dit aussi que $\underline{S}$ est $\alpha_0$-$\pi$-accessible dans $\mathcal{M}$. Si $\Phi$ est formé des catégories de cardinal $\leqslant\pi$ et grandes devant $\pi_0$, on dit que $\underline{S}$ est $\pi_0$-$\pi$-accessible dans $\mathcal{M}$.}).

\textbf{N.B.} Quand la condition préliminaire 1°) \textit{sur} $\mathcal{M}$ est satisfaite, la validité des conditions 2°) 3°) équivaut donc à la condition que $\underline{S}$ soit \textit{strictement} pleine dans $\mathcal{M}$, et que de
\printedpage{131}
% head: A. Grothendieck : Dérivateurs, Ch. XVIII, Édition des Universités de Montpellier II et Paris VII
% foot: Version 17 décembre 2025
plus $\underline{S}$ satisfasse la condition a) du théorème 4.13.22. D'autre part, cette condition est équivalente à chacune des conditions a$'$, b, c de ce même théorème, et aussi à la dernière condition c$'$, sous réserve qu'on ait $\pi^{\pi_0}=\pi$ (pour un cardinal $\pi_0<\pi$ utile pour $\mathcal{M}$) et que $J\in\Phi$ implique que $J$ est grande devant $\pi_0$.

Conceptuellement, c'est la définition donnée dans 4.13.23, i.e. la condition \cercle{a} de 4.13.22, qui paraît la plus claire, ou sa variante immédiate a$'$. Mais du point de vue pratique, du travail avec les parties accessibles, c'est par propriété \cercle{c} qui paraît la plus commode. Le poids principal du théorème 4.13.22 est bien
\msp{page 345}
l'implication a $\Longrightarrow$ c (via a $\Longrightarrow$ a$'$ $\overset{!}{\Longrightarrow}$ c). Donc il faut surtout voir ce théorème comme donnant un critère remarquable d'appartenance à $\Ob\underline{S}$ d'un objet $x$ de $\mathcal{M}$ qui est présenté par un système inductif $\Phi$-$\pi$-adapté :
\[ x\ \in\ \Ob\underline{S}\qquad\quad\Longleftrightarrow\qquad\quad I_{\underline{S}}\ \text{cofinal}\ , \tag*{$(*)$} \]
et dans cette équivalence, c'est l'implication $\Longrightarrow$ qui est la partie non soritale, et remarquable.

\textbf{Proposition 4.13.24.} \textit{Soient $\mathcal{M}$, $\pi$, $\Phi$ comme dans 4.13.23} (\footnote{\textbf{N.B.} La condition $I_{\alpha_0}\in\Phi$ n'est pas utile.})\textit{, on suppose que la condition 1° de 4.13.23 est satisfaite ($\pi$ utile pour $\mathcal{M}$, $\mathcal{M}$ stable par $\varinjlim$ de type $\Phi$). Soit ${\Sigma'}^{\Phi}_{\pi}$ l'ensemble des sous-catégories $\underline{S}$ de $\mathcal{M}$ qui sont $\Phi$-$\pi$-accessibles dans $\mathcal{M}$, et $\Sigma^{\Phi}_{\pi}$ l'ensemble des sous-catégories $\underline{S}(\pi)$ de $\mathcal{M}_\pi$ qui sont strictement pleines, et stables dans $\mathcal{M}_\pi$ par $\varinjlim$ de type $\Phi$. (Ces sous-catégories sont forcément karoubiennes, car stables par $\varinjlim$ de type $I_{\alpha_0}$, donc aussi par facteurs directs.) Alors les bijections (4.13.4.1) (page 288) induisent des bijections inverses l'une à l'autre}
\[
\begin{tikzcd}
\Sigma^{\Phi}_{\pi} \arrow[r, shift left] & {\Sigma'}^{\Phi}_{\pi} \arrow[l, shift left]
\end{tikzcd}\ .
\]

\msp{page 346}
% marge : trait manuscrit (reproduit en image) en regard de la Proposition 4.13.25
\textbf{Proposition 4.13.25.} \textit{Soient $\mathcal{M}$, $\pi$, $\Phi$ comme dans 4.13.23, avec $I_{\alpha_\pi}\in\Phi$ ($\alpha_\pi$ étant le plus petit ordinal de cardinal $\Phi$} \texttt{[plutôt} $\pi$\texttt{]}\textit{). Alors l'ensemble ${\Sigma'}^{\Phi}_{\pi}$ des sous-catégories de $\mathcal{M}$ $\Phi$-$\pi$-accessibles dans $\mathcal{M}$ est stable par intersection de familles de cardinal $\leqslant\pi$} (\footnote{Si $\pi$ n'est pas cardinal limite. Sinon, il faut dire : familles de cardinal $<\pi$. Cf. 4.13.10.3, où je donne un énoncé rectifié légèrement plus général.}).

(\textbf{N.B.} La bijection $\Sigma^{\Phi}_{\pi}\xrightarrow{\ \sim\ }{\Sigma'}^{\Phi}_{\pi}$ de 4.13.24 respecte les relations d'ordre naturelles d'inclusion, donc transforme intersections en intersections.)

Se voit comme 4.13.7 (qui est contenu dans l'énoncé actuel), en utilisant le lemme 4.13.10.

\textbf{4.13.26.} Soient $\mathcal{M}$, $\pi$, $\Phi$ comme dans 4.13.23, et soit $\underline{S}$ une sous-catégorie \textit{strictement pleine} de $\mathcal{M}$, $f:\underline{S}\longrightarrow\mathcal{M}$ le foncteur d'inclusion. La condition que $\underline{S}$ soit $\Phi$-$\pi$-accessible dans $\mathcal{M}$ se décompose en :
\begin{itemize}[leftmargin=2.5em]
\item[1)] $\pi$ est utile pour $\mathcal{M}$ \textit{et} pour $\underline{S}$.
\item[2)] $\mathcal{M}$ et $\underline{S}$ sont stables par $\varinjlim$ de type $\Phi$.
\item[3)] $f$ commute auxdites limites.
\item[4)] $f$ est $\pi$-accessible.
\item[5)] $f(\underline{S}_\pi)\subset\mathcal{M}_\pi$.
\end{itemize}

\printedpage{132}
% head: A. Grothendieck : Dérivateurs, Ch. XVIII, Édition des Universités de Montpellier II et Paris VII
% foot: Version 17 décembre 2025
\msp{page 347}
Soit alors $\pi'>\pi$ un autre cardinal, quand peut-on assurer que si $\underline{S}$ est $\Phi$-$\pi$-accessible dans $\mathcal{M}$, il soit aussi $\Phi$-$\pi'$-accessible\,? Les conditions 2) 3) 4) sont visiblement stables par passage de $\pi$ à $\pi'$. Restent les conditions 1) et 5). On sait que si $\pi'^{\pi}=\pi'$, alors la condition 1) reste valable pour $\pi'$. Plus généralement, si $\pi_0$ est un cardinal $\leqslant\pi$ tel que $\pi_0$ soit utile pour $\mathcal{M}$ et pour $\underline{S}$, alors si $\pi'>\pi_0$ est tel que $\pi'^{\pi_0}=\pi'$, alors $\pi'$ est aussi utile pour $\mathcal{M}$ et pour $\underline{S}$. Si on admet l'HCG (hypothèse du continu généralisée), il en est ainsi pour tout $\pi'>\pi_0$, donc pour tout $\pi'>\pi$.

Reste la condition 5), pour laquelle une réponse est donnée par 4.12.1 (p. 273). Si $\pi_0\leqslant\pi$ est choisi de façon que $f$ soit $\pi_0$-accessible (p.ex. $\pi_0=\pi$), et si
\[
\begin{aligned}
\pi'_1 &= \cardred\underline{S}_{\pi_0} \qquad\quad (\,\leqslant\cardred\underline{S}(\pi)\,)\\
\pi''_1 &= \sup(\pi'_1,\pi)\ ,
\end{aligned}
\]
on voit que si $\pi'\geqslant\pi''_1$ (i.e. $\pi'\geqslant\pi'_1,\pi$) est tel que $\pi'^{\pi_0}=\pi'$, alors on a bien
\[ f(\underline{S}_{\pi'})\ \subset\ \mathcal{M}_{\pi'}\ . \]

\msp{page 348}
Si donc de plus $\pi'>\pi\geqslant\pi_0$ (comme supposé au début), et si on suppose l'HCG, alors la condition $\pi'^{\pi_0}=\pi'$ est automatique, et il reste la seule condition $\pi'\geqslant\pi''_1$, i.e. en définitive $\pi'\geqslant\cardred\underline{S}_{\pi_0}$. Pour résumer :

\textbf{Proposition 4.13.27.} \textit{Sous les conditions de 4.13.23 (sans condition b) sur $\Phi$), soit $\underline{S}$ une sous-catégorie de $\mathcal{M}$ qui est $\Phi$-$\pi$-adaptée. Soit $\pi_0\leqslant\pi$ un cardinal qui est utile pour $\underline{S}$ et pour $\mathcal{M}$, et tel que $f:\underline{S}\longrightarrow\mathcal{M}$ soit $\pi_0$-accessible, i.e. tel que $\underline{S}$ soit stable dans $\mathcal{M}$ par $\varinjlim$ grandes devant $\pi_0$ (on peut p.ex. prendre $\pi_0=\pi$). Alors pour tout cardinal $\pi'>\pi$ tel que}
\begin{itemize}[leftmargin=2.5em]
\item[\textit{a)}] \textit{$\pi'^{\pi_0}=\pi'$ et}
\item[\textit{b)}] \textit{$\pi'\geqslant\cardred\underline{S}_{\pi_0}$,}
\end{itemize}
\textit{$\underline{S}$ est $\Phi$-$\pi'$-accessible dans $\mathcal{M}$. Donc si $\pi'>\pi$ est tel que $\pi'^{\pi}=\pi'$, $\pi'>\cardred\underline{S}_\pi$, alors $\underline{S}$ est $\Phi$-$\pi'$-accessible. (Il suffit p.ex. de prendre $\pi'=2^c$, avec $c\geqslant\sup(\pi,\cardred\underline{S}(\pi))$.) Si on admet l'HCG, il suffit donc qu'on ait $\pi'>\sup(\pi,\cardred\underline{S}(\pi))$.}

\msp{page 349}
\textbf{Proposition 4.13.28} (\footnote{À rejeter plus bas, après 4.13.31, qui est la première application importante de 4.13.22, alors que 4.13.28 se prouve sans avoir à utiliser 4.13.22 ou ses variantes.}). \textit{Soient $\mathcal{M}$, $\pi$, $\Phi$ comme dans 4.13.23. Supposons de plus que $J\in\Phi$ implique que $J$ est} filtrante. \textit{Alors l'ensemble $\tilde{\Sigma}^{\Phi}_{\pi}$ des sous-ensembles $\Phi$-$\pi$-accessibles de $\mathcal{M}$ est stable par réunions finies. Plus généralement, soit $\lambda$ un cardinal $<\pi$ tel que $\forall\ J\in\Phi$, $J$ soit grande devant $\lambda$. Alors l'ensemble $\tilde{\Sigma}^{\Phi}_{\pi}$ est stable par réunions $(S_r)_{r\in A}$, indexées par des ensembles $A$ de cardinal $\leqslant\lambda$.}

\textbf{Corollaire 4.13.29.} \textit{L'ensemble des parties accessibles de} $\Ob\mathcal{M}$ \textit{est stable par “petites” réunions -- i.e. si $(S_r)_{r\in A}$ est une petite famille de parties accessibles de} $\Ob\mathcal{M}$, \textit{alors $\bigcup_{r\in A}S_r$ est aussi accessible.}

La démonstration de 4.13.28 repose sur le lemme suivant (je n'explicite par la démonstration) : % [sic] par la démonstration -> pas la démonstration, sense of the sentence

\printedpage{133}
% head: A. Grothendieck : Dérivateurs, Ch. XVIII, Édition des Universités de Montpellier II et Paris VII
% foot: Version 17 décembre 2025
\msp{page 350}
\textbf{Lemme 4.13.30.} \textit{Soit $I$ une catégorie grande devant $\pi$, et $(J_r)_{r\in A}$ une famille des sous-catégories pleines telles que $\card A\leqslant\pi$, telle que $\bigcup_{r\in A}\Ob J_r$ soit cofinal dans $\Ord(I)$, i.e. pour tout $i\in\Ob I$ $\exists\ j\in\bigcup_{r\in A}\Ob J_r$, et une flèche $i\longrightarrow j$. Alors il existe $r\in A$ tel que $J_r$ soit cofinal dans $I$.}

\textbf{Théorème 4.13.31.} \textit{Soient $\mathcal{M}$, $\mathcal{N}$ des catégories accessibles, $f:\mathcal{M}\longrightarrow\mathcal{N}$ un foncteur accessible, $T\subset\Ob\mathcal{N}$ une partie de $\mathcal{N}$, $\underline{T}\subset\mathcal{N}$ la sous-catégorie strictement pleine correspondante de $\mathcal{N}$, $S=f^{-1}(T)$, $\underline{S}=f^{-1}(\underline{T})$ la sous-catégorie pleine de $\mathcal{M}$ associée. Si $T$ est accessible, alors $S$ est accessible\,; en d'autres termes, si $\underline{T}$ est accessible dans $\mathcal{N}$, $\underline{S}$ est accessible dans $\mathcal{M}$.}

Cela résulte de l'énoncé plus précis :

\textbf{Corollaire 4.13.32.} \textit{Soient $\pi_0$ et $\pi>\pi_0$ des cardinaux ayant les propriétés suivantes :}

\msp{page 351}
\begin{itemize}[leftmargin=2.5em]
\item[\textit{a)}] \textit{$\underline{T}$ est $\pi_0$-$\pi$-accessible dans $\mathcal{N}$ (4.13.23).}
\item[\textit{b)}] \textit{$\pi_0$ est utile pour $\mathcal{M}$ (il l'est pour $\mathcal{N}$, à cause de a)).}
\item[\textit{c)}] \textit{$f(\mathcal{M}_\pi)\subset\mathcal{N}_\pi$.}
\item[\textit{d)}] \textit{$\mathcal{M}$ satisfait à $L_{\pi_0}$, et $f$ est $\pi_0$-accessible.}
\end{itemize}
\textit{(L'existence d'un tel couple $(\pi_0,\pi)$ résulte aussitôt de 4.13.16 et de 4.12.1.) Alors $\underline{S}$ est $\pi_0$-$\pi$-accessible dans $\mathcal{N}$.}

Comme $f$ commute aux $\varinjlim$ grandes devant $\pi_0$ en vertu de d), et que $\underline{T}$ est stable par ces limites par a) (cf. 4.13.23, 3°), il en est donc de même de $\underline{S}$, i.e. le 3° de la définition 4.13.23 est satisfait. D'autre part, la condition 1° ($\pi$ utile pour $\mathcal{M}$, $\mathcal{M}$ stable par $\varinjlim$ grandes devant $\pi_0$ et de cardinal $\leqslant\pi$) est vérifiée par b) et d). Donc il reste à vérifier la condition 2°, i.e. que $\underline{S}$ est $\pi$-accessible. Posons
\[
\begin{aligned}
\underline{T}(\pi) &= \underline{T}\cap\mathcal{N}_\pi\\
\underline{S}(\pi) &= \underline{S}\cap\mathcal{M}_\pi\ \ \underbrace{\ \subset\ }_{\text{grâce à c)}}\ \ f^{-1}(\underline{T}(\pi))\ .
\end{aligned}
\]
Comme $\underline{S}$ est stable dans $\mathcal{M}$ par $\varinjlim$ grandes devant $\pi_0$ (lesquelles existent dans $\mathcal{M}$ par d), et a fortiori par $\varinjlim$ grandes devant $\pi$, on a un foncteur canonique pleinement fidèle
\[ \Ind_\pi\underline{S}(\pi)\ \longrightarrow\ \underline{S}\ . \]

\msp{page 352}
D'ailleurs, $\underline{T}$ étant strictement pleine dans $\mathcal{N}$, $\underline{S}$ est strictement pleine dans $\mathcal{M}$, et il reste à vérifier que le foncteur précédent est essentiellement surjectif. Soit donc $x\in\Ob\underline{S}$. On aura, grâce à b),
\[ x\ =\ \varinjlim_{I}x_i\ ,\qquad\quad x_i\in\mathcal{M}_\pi,\ I=\mathcal{M}_\pi/x\ . \]
D'ailleurs, le système inductif précédent est $\pi_0$-$\pi$-adapté, i.e. (p. 339) :
\begin{itemize}[leftmargin=2.5em]
\item[1)] $I$ équivalente à une petite catégorie, $I$ grande devant $\pi$ (résulte du fait que $\pi$ est utile pour $\mathcal{M}$, cf.\,\cercle{b}\,).
\end{itemize}


\endgroup
% END UNIT der-ch18/der-ch18_p131-135.tex
% BEGIN UNIT der-ch18/der-ch18_p136-140.tex
\begingroup\makeatletter%
\newcommand{\printedpage}[1]{\marginpar{\footnotesize p.~#1}}%
\newcommand{\msp}[1]{\marginpar{\footnotesize [#1]}}%
\newcommand{\Ob}{\operatorname{Ob}}%
\newcommand{\Fl}{\operatorname{Fl}}%
\newcommand{\card}{\operatorname{card}}%
\newcommand{\Ind}{\operatorname{Ind}}%
\newcommand{\Hom}{\operatorname{Hom}}%
\newcommand{\id}{\mathrm{id}}%
\newcommand{\limind}{\varinjlim}%
\newcommand{\cercle}[1]{\tikz[baseline=(c.base)]\node[draw,circle,inner sep=1pt](c){#1};}%
\newcommand{\defeq}{\overset{\text{déf}}{=}}%
\makeatother\let\printedpage\derpp\let\msp\dermsp%
\let\mspin\DERmspin%


\printedpage{134}
% head: A. Grothendieck : Dérivateurs, Ch. XVIII, Édition des Universités de Montpellier II et Paris VII
% foot: Version 17 décembre 2025
\begin{itemize}[leftmargin=2.5em]
\item[2)] Les limites inductives grandes devant $\pi_0$ est de cardinal $\leqslant \pi$ existent dans $I$ (résulte de $\mathrm{L}_{\pi_0}$ sur $\mathcal{M}$ (cf.\,\cercle{d})). % [sic] est de cardinal -> et de cardinal, coordination ("grandes devant π0 et de cardinal ≤ π")
\item[3)] Le foncteur $\xi : I \longrightarrow \mathcal{M}$ y commute.
\item[4)] Les $x_i$ sont dans $\mathcal{M}_\pi$.
\end{itemize}
Comme $f$ commute aux $\limind$ grandes devant $\pi_0$ (cf.\,\cercle{1}), a fortiori à celles grandes devant $\pi$, on a donc
\[ f(x) \;\simeq\; \varinjlim_{I} f(x_i) \ , \]
je dis que le système inductif $(f(x_i))_{i\in I}$ dans $\mathcal{N}$ est $\pi_0$-$\pi$-adapté. Les conditions 1), 2) sont \textit{ne varietur} (ne concernent que $I$), le foncteur $I \longrightarrow \mathcal{N}$ est le composé $I \xrightarrow{\xi} \mathcal{M} \xrightarrow{f} \mathcal{N}$, donc commute aux $\limind$ grandes devant $\pi_0$ et de cardinal $\leqslant \pi$ (\footnote[155]{cette commutation devrait avoir un nom, condition $\mathrm{L}(\pi_0,\pi)$ sur un foncteur.}), puisqu'il en est ainsi de $\xi$ et de $f$, enfin les $f(x_i)$ sont dans $\mathcal{N}_\pi$ puisque par \cercle{c} $f(\mathcal{M}_\pi) \subset \mathcal{N}_\pi$.

\msp{page 353}

Par ailleurs $f(x) \in \Ob \underline{T}$ (puisque $x \in \Ob \underline{S}$, où $\underline{S} \defeq f^{-1}(\underline{T})$). Donc on peut appliquer 4.13.22 a $\Longrightarrow$ c à la situation $(\mathcal{N}, \pi_0, \pi, \underline{S})$ et au système inductif $(f(x_i))_{i\in I}$, pour conclure que la sous-catégorie $I'$ de $I$ formée des $i$ tels que $f(x_i) \in \Ob \underline{S}$ \texttt{[plutôt Ob}~$\underline{T}$\texttt{]}, i.e. tels que $x_i \in \underline{S}$, donc $x_i \in \underline{S}(\pi)$, est \textit{cofinale} dans $I$. Donc on a
\[ x \;\simeq\; \varinjlim_{i\in I'} x_i \ , \]
avec cette fois les $x_i$ dans $\underline{S}(\pi)$, et $I'$ grande devant $\pi$ (car cofinale et pleine dans $I$, qui est grande devant $\pi$), ce qui prouve que $x$ est dans l'image essentielle de $\Ind_\pi(\underline{S}(\pi))$, q.e.d.

\textbf{Question 4.13.33.} On vient de voir que si $f : \mathcal{M} \longrightarrow \mathcal{N}$ est un foncteur accessible entre catégories accessibles, alors l'image inverse d'une partie accessible de $\Ob \mathcal{N}$ est une partie accessible de $\Ob \mathcal{M}$. Je conjecture que l'énoncé similaire, non pas pour les images directes, mais pour les \textit{images essentielles}, est
\msp{page 354}
vraie, i.e. que pour toute sous-catégorie $\underline{S}$ de $\mathcal{M}$ accessible dans $\mathcal{M}$, son image essentielle dans $\mathcal{N}$ est accessible dans $\mathcal{N}$ (\footnote[156]{Faux, même si on prend la $\natural$-clôture de l'image essentielle (“image $\natural$-essentielle”).}). Dans cette question, on peut supposer que $\underline{S} = \mathcal{M}$, et la question revient à celle-ci : tout foncteur accessible $f : \mathcal{M} \longrightarrow \mathcal{N}$ entre catégories accessibles se factorise-t-il en un foncteur (accessible) \textit{essentiellement surjectif}, et un foncteur accessible pleinement fidèle $\mathcal{M} \longrightarrow \underline{T} \longrightarrow \mathcal{N}$ (où $\underline{T}$ est elle-même accessible) (\footnote[157]{Il faut se borner tout au moins à essentiellement $\natural$-surjectif – mais même cela est trop optimiste.})~? Si oui, cela me paraîtrait peut-être le résultat le plus profond (peut-être le seul véritablement profond) en théorie des catégories. % [sic] vraie -> vrai, agreement with « l'énoncé similaire »

Voici une question qui me paraît très voisine. Soient $\mathcal{M}$ une catégorie accessible, $A$ une sous-catégories pleine petite. Supposons que $\mathcal{M}$ satisfasse à $\mathrm{L}_{\pi_0}$ (p.ex. $\pi_0$ utile pour $\mathcal{M}$), et soit $\underline{S}$ la plus petite sous-catégorie strictement pleine de $\mathcal{M}$ contenant $A$, et stable par limites grandes devant $\pi_0$. Je conjecture que $\underline{S}$ est accessible dans $\mathcal{M}$. % [sic] sous-catégories -> sous-catégorie, singular after « une »

\msp{page 355}

De même, si $\Phi$ est un petit ensemble de petites catégories, tel que $\mathcal{M}$ soit stable par $\limind$ de type $\Phi$, on peut prendre l'enveloppe $\underline{S}$ de $A$ pour les $\limind$ grandes devant $\pi_0$ \textit{et} les
\printedpage{135}
% head: A. Grothendieck : Dérivateurs, Ch. XVIII, Édition des Universités de Montpellier II et Paris VII
% foot: Version 17 décembre 2025
$\limind$ de type $\Phi$, et je conjecture encore que $\underline{S}$ est accessible. Dans cette direction, on a du moins le résultat très partiel suivant.

\textbf{Proposition 4.13.34.} \textit{Soit $\mathcal{M}$ une catégorie accessible stable par petites $\limind$ filtrantes, et soit $\underline{S}$ une sous-catégorie strictement pleine de $\mathcal{M}$. Conditions équivalentes :}
\begin{itemize}[leftmargin=2.5em]
\item[\textit{a)}] \textit{$\underline{S}$ est accessible dans $\mathcal{M}$, et stable dans $\mathcal{M}$ par petites $\limind$ filtrantes.}
\item[\textit{b)}] \textit{Il existe une petite sous-catégorie pleine $A$ de $\mathcal{M}$, telle que $\underline{S}$ soit la plus petite sous-catégorie strictement pleine de $\mathcal{M}$ contenant $A$ et stable par petites $\limind$ filtrantes.}
\end{itemize}

\msp{page 356}

On a a) $\Longrightarrow$ b), car si $\underline{S}$ est $\pi$-accessible dans $\mathcal{M}$, il est immédiat (sans autre condition de stabilité sur $\underline{S}$) que $\underline{S}$ est la plus petite sous-catégorie strictement pleine de $\mathcal{M}$ contenant $\underline{S}(\pi) \defeq \underline{S}\cap\mathcal{M}_\pi$ et stable par $\limind$ grandes devant $\pi$. Si donc on sait déjà que $\underline{S}$ est même stable par \textit{toutes} $\limind$ filtrantes petites, il s'ensuit que $\underline{S}$ est la plus petite sous-catégorie de $\mathcal{M}$ ayant cette propriété de stabilité plus forte, et contenant $A$.

Prouvons b $\Longrightarrow$ a. Soit $\pi$ un cardinal utile pour $\mathcal{M}$, et tel que $A \subset \mathcal{M}_\pi$. Soit $B$ la plus petite sous-catégorie strictement pleine de $\mathcal{M}$ contenant $A$ et stable par $\limind$ filtrantes de cardinal $\leqslant \pi$. On a donc $B \hookrightarrow \mathcal{M}_\pi$ et $B \subset \underline{S}$, et $\underline{S}$ est aussi la L-enveloppe de $B$. D'ailleurs $B$ est stable par
\msp{page 357}
facteurs directs. Considérons alors l'image essentielle $\underline{S}'$ de $\Ind_\pi(B)$ dans $\mathcal{M}$. On a $\underline{S}' \subset \underline{S}$, et tout revient à prouver que $\underline{S}' = \underline{S}$, d'où résultera en effet que $\underline{S}$ est $\pi$-accessible dans $\mathcal{M}$, a fortiori accessible dans $\mathcal{M}$, d'où a). Pour prouver que $\underline{S}' = \underline{S}$, par définition de $\underline{S}$ comme la L-enveloppe de $A$, il suffit de voir que $\underline{S}'$ est stable par petites $\limind$ filtrantes. Or cela résulte du fait que $B \subset \mathcal{M}_\pi$ est stable dans $\mathcal{M}_\pi$ par $\limind$ filtrantes de cardinal $\leqslant \pi$, et de 4.13.19, qui implique tout au moins, comme $B = \underline{S}' \cap \mathcal{M}_\pi$ (en vertu de 4.13.4) que $\underline{S}'$ est stable par $\limind$ filtrantes \textit{de cardinal} $\leqslant \pi$. Qu'il soit même stable par petites $\limind$ filtrante quelconques, résulte de ceci, et du fait qu'il l'est par $\limind$ grandes devant $\pi$. (Cf. 4.13.21 \dots) % [sic] filtrante quelconques -> filtrantes quelconques, plural agreement

\msp{page 358}

\textbf{Définition 4.13.35.} Soient $\mathcal{M}$ une catégorie accessible, stable par petites $\limind$ filtrantes, $S$ une partie de $\Ob \mathcal{M}$, $\underline{S}$ la sous-catégorie pleine correspondante dans $\mathcal{M}$. On dit que $S$ est \textit{L-accessible}, ou que $\underline{S}$ est L-accessible dans $\mathcal{M}$, si $\underline{S}$ est accessible dans $\mathcal{M}$, et de plus stable par petites $\limind$ filtrantes (\footnote[158]{\textbf{N.B.} L'ensemble des parties L-accessibles de $\Ob \mathcal{M}$ est stable par petites intersections, et par réunions \textit{finies}. L'image inverse d'une partie L-accessible par un foncteur (\footnotemark[159]) $f : \mathcal{M} \longrightarrow \mathcal{N}$ de catégories accessibles est L-accessible.}).\footnotetext[159]{un \textit{$L_0$-foncteur} : qui commute aux petites $\limind$ filtrantes.}

\textbf{Proposition 4.13.36.} \textit{Soient $\mathcal{M}$, $S$, $\underline{S}$ comme dans 4.13.35, et soit $\pi$ un cardinal utile pour $\mathcal{M}$. Supposons $\underline{S}$ stable dans $\mathcal{M}$ par petites $\limind$ filtrantes, ou seulement $\underline{S}(\pi) \defeq \underline{S} \cap \mathcal{M}_\pi$ stable par petites $\limind$ filtrantes de cardinal $\leqslant \pi$. Conditions équivalentes :}
\begin{itemize}[leftmargin=2.5em]
\item[\textit{a)}] \textit{$\underline{S}$ est $\pi$-accessible. (Il sera automatiquement stable par petites $\limind$ filtrantes, donc L-accessible.)}
\item[\textit{b)}] \textit{Pour tout $x \in \Ob \mathcal{M}$, on a $x \in \Ob \underline{S}$ si et seulement si $\underline{S}(\pi)/x$ est cofinale dans $\mathcal{M}_\pi/x$.}
\printedpage{136}
% head: A. Grothendieck : Dérivateurs, Ch. XVIII, Édition des Universités de Montpellier II et Paris VII
% foot: Version 17 décembre 2025
\item[\textit{c)}] \textit{Pour tout $x \in \Ob \mathcal{M}$, et toute représentation $x \simeq \varinjlim\limits_{i\in I} x_i$}
\msp{page 359}
\textit{qui est L-$\pi$-adaptée (cf. page 343, c), on a ceci : $x \in \Ob S$ si et seulement si $I_{\underline{S}}$ est cofinal dans $I$.}
\item[\textit{c}$'$\textit{)}] \textit{Comme c), avec $I$ ordonnée.}
\item[\textit{c}$''$\textit{)}] \textit{(Si on a choisi, pour tout $x \in \Ob \mathcal{M}$,} une \textit{représentation $\pi$-adaptée $x \simeq \varinjlim\limits_{I} x_i$) : $x \in \Ob \underline{S} \Longleftrightarrow$ l'ensemble $I_{\underline{S}}$ est cofinal dans $I$.}
\end{itemize}
\textsc{Démonstration}. C'est essentiellement l'énoncé de 4.13.22 dans le cas particulier où $\Phi =$ ensemble des catégories \textit{filtrantes} de cardinal $\leqslant \pi$, mais où, pour simplifier, on a supposé d'avance que $\underline{S}(\pi)$ satisfait à $\mathrm{L}(\Phi)$. L'équivalence de a) b) c) est contenue dans 4.13.22, ainsi que c $\Longrightarrow$ c$'$. Mais pour c$'$ $\Longrightarrow$ a, l'hypothèse accessoire $\pi^{\pi_0} = \pi$ et que les $J \in \Phi$ seront grandes devant $\pi_0$, n'est pas satisfaite ici, et il faut un démonstration indépendante, calquée sur celle des pages 333-336. Il suffit pour cela, à l'endroit \texttt{[de]} $\square$, page 333, d'invoquer le % [sic] un démonstration -> une démonstration, gender agreement

\textbf{Lemme 4.13.37.} \textit{Soient $\mathcal{M}$ une catégorie}
\msp{page 360}
\textit{accessible stable par $\limind$ filtrantes, $\pi$ un cardinal utile pour $\mathcal{M}$. Alors pour tout $x \in \Ob \mathcal{M}$, on a une représentation}
\[ x \;=\; \varinjlim_{i\in I} x_i \]
\textit{L-$\pi$-adaptée (i.e. avec}
\begin{itemize}[leftmargin=2.5em]
\item[\textit{1)}] \textit{$I$ équivalente à une petite catégorie, et $I$ grande devant $\pi$,}
\item[\textit{2)}] \textit{$I$ stable par $\limind$ filtrantes de cardinal $\leqslant \pi$,}
\item[\textit{3)}] \textit{$I \longrightarrow \mathcal{M}$ commute à ces $\limind$ ,}
\item[\textit{d)}] \textit{$x_i \in \Ob \mathcal{M}_\pi \ \forall\, i \in I$),} % [sic] d) -> 4), the items are numbered 1), 2), 3)
\end{itemize}
\textit{avec de plus $I$} ordonné.

En effet, on part de
\[ x \;=\; \varinjlim_{j\in J} x_j \ , \qquad\quad J \;\simeq\; \underbrace{(\mathcal{M}_\pi/x)_0}_{\substack{\text{sous-catégorie réduite}\\ \text{équivalente à } \mathcal{M}_\pi/x}} \qquad . \]
On définit $I$ comme l'ensemble des sous-catégories \textit{filtrantes} $J'$ de $J$, de cardinal \texttt{[“cardinal” biffé]} $\leqslant \pi$, $I$ ordonné par inclusion, et on pose, pour $J' \in I$,
\[ x_{J'} \;=\; \varinjlim_{j\in J'} x_j \ . \]
Il est immédiat qu'on obtient ainsi un système inductif dans $\mathcal{M}$, indexé par $I$.

\msp{page 361}

D'autre part, il est clair que $I$ est petite, je dis que $I$ est stable par $\limind$ filtrantes de cardinal $\leqslant \pi$, et que le foncteur $J' \longmapsto x_{J'} : I \longrightarrow \mathcal{M}$ y commute, enfin que les $x_{J'}$ sont dans $\mathcal{M}_\pi$. Il reste à voir que $I$ est grande devant $\pi$, i.e. que si on a $(J'_r)_{r\in A}$, avec $\card A \leqslant \pi$, les $J'_r$ dans $I$, alors $\exists\, J' \in I$ qui majore les $J'_r$. Soit $K$ la sous-catégorie de $J$ engendrée par les $J'_r$, on voit facilement que $\card K \leqslant \pi$. Tout \texttt{[“tout” biffé]} revient à prouver :

\printedpage{137}
% head: A. Grothendieck : Dérivateurs, Ch. XVIII, Édition des Universités de Montpellier II et Paris VII
% foot: Version 17 décembre 2025
\textbf{Lemme 4.13.38} (\footnote[160]{\textbf{N.B.} Pour la démonstration de 4.13.37, il suffit de prouver 4.13.38 dans le cas où $J$ est ordonnée, auquel cas les contorsions de la fin sont inutiles.}). \textit{Soient $J$ une catégorie grande devant $\pi$, $K$ une sous-catégorie de $J$ de cardinal $\leqslant \pi$. Alors il existe une sous-catégorie $K'$} filtrante \textit{de $J$, contenant $K$, de cardinal $\leqslant \pi$. Si on suppose que l'on n'est pas, dans le cas où $J$ contient un objet maximal, et que $K$ contienne tous les objets maximaux de $J$,}
\msp{page 362}
\textit{on peut même prendre $K'$ telle que $\Ob K' = \Ob K \cup \{e\}$, où $e$ est un} objet final \textit{de $K'$, et $K$ est une sous-catégorie pleine de $K'$. (Donc si $K' \neq K$, $K'$ est isomorphe au cône à droite construit sur $K$, “déduit de $K$ en lui ajoutant un objet final”.)}

\textsc{Démonstration}. Si $K$ admet un objet final, on prend $K' = K$. Sinon, comme $J$ est grande devant $\card \Ob K$, il existe un $j \in \Ob J$ qui majore tous les objets de $K$. Choisissons, pour tout $k \in \Ob K$, une flèche
\[ \alpha_k \;:\; k \longrightarrow j \ . \]
Pour toute flèche
\[ u \;:\; k' \longrightarrow k \]
dans $K$, considérons la double-flèche
\[ (\alpha_k \circ u, \alpha_{k'}) \;:\; k' \rightrightarrows j \ . \]
L'ensemble de ces double-flèches, de but $j$, est de cardinal $\leqslant \card \Fl K \leqslant \pi$, donc $J$ étant grande devant $\pi$, il existe une flèche $\rho : j \longrightarrow e$ dans $J$ qui égalise toutes ces double-flèches. Soit $K' \subset J$, ayant comme objets $\Ob K \cup \{e\}$, et ayant comme flèches, en plus de
\msp{page 363}
celles de $K$, les flèches $\rho\alpha_k : k \longrightarrow e$ (pour $k \in \Ob K$), et $\id_e$. \textit{Si} $e \notin \Ob K$, on voit que c'est bien une sous-catégorie, admettant $e$ comme objet final et $K$ comme sous-catégorie pleine. Si on avait $e \in \Ob K$, et s'il existe un objet $e'$ de $J$ qui majore $e$ et qui n'est pas dans $K$, on se ramène aussitôt au cas précédent, en remplaçant $e$ par $e'$.

$\Biggl($ (\footnote[161]{La construction qui suit est contenue dans l'énoncé plus général et plus précis 4.13.39.}). Donc il reste le cas où il existe un objet $e$ dans $K$ qui est un plus grand élément dans $K$, et tel que tout objet de $J$ majorant $e$ soit dans $K$. Cela implique, $J$ étant filtrante, que $e$ est un plus grand objet de $J$. (Car si $x \in J$, $\exists\, x'$ majorant commun de $x$ et $e$, donc $x'$ est dans $K$, donc majoré par $e$, qui majore donc $x$.) On voit comme tantôt qu'il existe une flèche $p : e \longrightarrow j$ dans $J$, qui égalise toutes les double-flèches $k \rightrightarrows e$ de but $e$ \textit{dans}~$K$. Quitte à composer avec $j \longrightarrow e$, on peut supposer $j = e$. On a donc un endomorphisme $p$ de $e$ dans $J$, tel que $pu = pv$ pour toute double-flèche $k \overset{u}{\underset{v}{\rightrightarrows}} e$ dans $K$, de but $e$.

\msp{page 364}

Pour $k \in \Ob K$, soit
\[ f_k \;:\; k \longrightarrow e \]
la valeur commune des $pu$, pour $u \in \Hom_K(k,e)$. On voit tout de suite que si on a $k' \xrightarrow{v} k$ dans $K$, alors
\[ f_{k'} \;=\; f_k \circ v \ . \]
\printedpage{138}
% head: A. Grothendieck : Dérivateurs, Ch. XVIII, Édition des Universités de Montpellier II et Paris VII
% foot: Version 17 décembre 2025
En particulier, prenant $k' = e$, on trouve
\[ f_k \circ v \;=\; f_e \,(=\ p\ ) \qquad\quad \text{si } v : e \longrightarrow k \ . \]
Soit donc $K'$ la sous-catégorie de $J$ ayant mêmes objets que $K$, et comme flèches, en plus de celles de $K$, les flèches composées $k \xrightarrow{u} e \xrightarrow{p^r} e \xrightarrow{v} k'$, avec $u$, $v$ dans $K$, et $r \geqslant 1$, i.e. les flèches composées $k \xrightarrow{f_k} e \xrightarrow{p^{r'}} e \xrightarrow{v} k'$, avec $k, k' \in \Ob K$, $v \in \Fl K$, $r' \geqslant 0$. Il faut voir que c'est stable par composition (pour avoir bien une sous-catégorie). C'est clair pour $vu$ si $u$ ou $v$ sont dans $\Fl K$. Si \texttt{[ni]} $u$ ni $v$ ne le sont, on doit prendre un composé
\[
\begin{tikzcd}[column sep=large]
k \arrow[r, "u_1"] & e \arrow[r, "p^r", "r\geqslant 1"'] & e \arrow[r, "u_2"] \arrow[rr, bend right=20] & k' \arrow[r, "v_1"] & e \arrow[r, "p^s", "s\geqslant 1"'] & e \arrow[r, "v_2"] & k'' \ ,
\end{tikzcd}
\]
avec $u_1$, $u_2$, $v_1$, $v_2$ dans $\Fl K$, $r, s \geqslant 1$. Or on a
\[ p(v_1u_2) \;=\; p \ , \qquad\quad \text{d'où} \qquad\quad p^s(v_1u_2) \;=\; p^{s-1}(pv_1u_2) \;=\; p^{s-1}p \;=\; p^s \ , \]
donc le composé est
\[ k \xrightarrow{\;u_1\;} e \xrightarrow{\;p^{r+s}\;} e \xrightarrow{\;v_2\;} k'' \ , \]
qui
\msp{page 365}
est donc bien dans $\Fl K'$.

Ainsi, $K'$ est une sous-catégorie de $J$, mais est-elle filtrante ? La condition PS 1$'$ est évidente, puisque $e$ est un plus grand élément de $J$, reste la condition PS 2 de la double-flèche. On voit aisément que cette condition est satisfaite dans $K'$, si et seulement si elle est satisfaite pour toute double-flèche
\[ e \overset{p^r}{\underset{p^s}{\rightrightarrows}} e \qquad\quad \text{avec } r, s \geqslant 0 \ . \]
(En effet, on se ramène aussitôt aux double-flèches de but $e$, et on note que toute flèche de $K'$ de but $e$ s'écrit comme un composé $k \xrightarrow{u} e \xrightarrow{p^r} e \xrightarrow{v} e$, avec $r \geqslant 0$ et $u$, $v$ des flèches dans $K'$. En composant avec $e \xrightarrow{p} e$, on trouve alors, comme $pv = pu$, $k \xrightarrow{u} e \xrightarrow{p^{r+1}} e$, soit $k \xrightarrow{f_k} e \xrightarrow{p^r} e$.) Mais on voit aussitôt que la possibilité d'égaliser $(p^r, p^s)$ dans $K'$ équivaut à l'existence de $t \geqslant 0$ tel que $p^{r+t} = p^{s+t}$. Donc $K'$ est filtrante si et seulement si $r \longmapsto p^r$ est une suite stationnaire (faire $r = 0$, $s = 1$, d'où $p^r = p^{r+1}$ pour $r$ grand). Ce n'est
\msp{page 366}
pas forcément le cas.

Posons donc $K = K_0$, $K' = K_1$, et construisons par induction une suite croissante de sous-catégories $K_i$ ($i \in \mathbf{N}$), avec $\Ob K_i = \Ob J$, et (pour $i \geqslant 1$) $K_i$ étant déduites de $K_{i-1}$ en lui rajoutant une flèche $p_i : e \longrightarrow e$, comme ci-dessus. Ainsi toute double-flèche dans un $K_i$ peut s'égaliser dans $K_{i+1}$. Prenons $K_\infty = \limind K_i$, $K_\infty$ est une sous-catégorie \textit{filtrante} de $J$, contenant $K$ et de cardinal $\leqslant \pi$, q.e.d. $\Biggr)$

On peut préciser cette construction ainsi :


\endgroup
% END UNIT der-ch18/der-ch18_p136-140.tex
% BEGIN UNIT der-ch18/der-ch18_p141-145.tex
\begingroup\makeatletter%
\newcommand{\printedpage}[1]{\marginpar{\footnotesize p.~#1}}%
\newcommand{\msp}[1]{\marginpar{\footnotesize [#1]}}%
\newcommand{\uHom}{\underline{\mathrm{Hom}}}%
\newcommand{\Fl}{\operatorname{Fl}}%
\newcommand{\Ob}{\operatorname{Ob}}%
\newcommand{\End}{\operatorname{End}}%
\newcommand{\Ind}{\operatorname{Ind}}%
\newcommand{\id}{\mathrm{id}}%
\newcommand{\cardred}{\mathop{\textup{card\,réd}}}%
\newcommand{\card}{\operatorname{card}}%
\newcommand{\edbr}[1]{\textup{\texttt{#1}}}%
\makeatother\let\printedpage\derpp\let\msp\dermsp%
\let\mspin\DERmspin%


\printedpage{139}
% head: A. Grothendieck : Dérivateurs, Ch. XVIII, Édition des Universités de Montpellier II et Paris VII
% foot: Version 17 décembre 2025

\noindent\textbf{Corollaire 4.13.39.} {\itshape Sous les conditions de 4.13.38, supposons que $K$ contienne un objet $e$ maximal dans $J$. Soit $\Sigma$ l'ensemble des endomorphismes $p$ de $e$ dans $J$, tels que pour toute double-flèche $k \overset{u}{\underset{v}{\rightrightarrows}} e$ \textup{dans} $K$, on ait $pu=pv$. Alors $\Sigma$ est stable par composition, et il existe une partie non vide $\Sigma_\infty$ \edbr{[}$=\Sigma_0$\edbr{]} de $\Sigma$, telle que pour $p,p'\in\Sigma_0$, existe
\msp{page 367}
$q\in\Sigma_0$ tel que $qp=qp'$ (i.e. telle que la sous-catégorie $J_0$ de $J$ dont l'ensemble des objets est $\{e\}$, et l'ensemble des flèches $\{\id_e\}\cup\Sigma_\infty$, soit filtrante). Si $\Sigma_0$ est choisie ainsi, soit $\mathcal{F}$ l'ensemble réunion de $\Fl K$ et de l'ensemble des flèches composées
\[
k \xrightarrow{\ u\ } e \xrightarrow{\ p\ } e \xrightarrow{\ v\ } k'\ , \tag*{$(*)$}
\]
avec $u,v\in\Fl J$, $p\in\Sigma_\infty$. Alors $\mathcal{F}$ est l'ensemble des flèches d'une sous-catégorie $K'$ de $J$, ayant mêmes objets que $J$, et admettant $e$ pour plus grand élément. De plus, $K'$ contient $K$ et est filtrante.}

\noindent\textsc{Démonstration}. Que $\Sigma$ soit stable par composition est clair. Pour construire $\Sigma_\infty$ ayant les propriétés énoncées, on note que, $J$ étant grande devant $\pi$, $\Sigma$ est non vide (\footnote[162]{cf. argument plus bas pour construction de $p_i$.}). On construit par récurrence une suite d'éléments
\[
p_0\ ,\ p_1\ ,\ \ldots\ ,\ p_n\ ,\ \ldots
\]
de $\Sigma$, en prenant pour $p_0$ un élément quelconque de $\Sigma$, et pour $i\geqslant 1$, prenant $p_i\in\Sigma$ tel que pour toute double-flèche $(p,p')$ du sous-monoïde $M_{i-1}$ de $\End_J(e)$ engendré par $(p_j)_{0\leqslant j\leqslant i-1}$, on ait $p_ip=p_ip'$. Un tel $p_i$ existe, car prenant l'ensemble des double-flèches de la forme précédente $(p,p')$, avec $p,p'\in M_{i-1}$, et de celles de la forme $(u,v) : k \rightrightarrows e$, avec $u,v\in\Fl J$, l'ensemble de
\msp{page 368}
toutes ces double-flèches est de cardinal $\leqslant\pi$ ($M_{i-1}$ étant dénombrable, donc de cardinal $\leqslant\pi$), donc il existe une flèche $e\longrightarrow j$ dans $J$ qui les égalise toutes ($J$ étant grande devant $\pi$). Quitte à composer avec $j\longrightarrow e$ ($e$ étant objet maximal de $J$), on trouve le $p_i : e\longrightarrow e$ cherché.

Soit alors $\Sigma_0$ le sous-ensemble multiplicatif de $\Sigma$ engendré par les $p_i$ ($i\in\mathbf{N}$), limite inductive des ensembles similaires $\Sigma_i$ engendrés par les $p_j$ pour $j\in[0,i]$. Si $p$, $p'$ sont dans $\Sigma_\infty$, il sont dans un même $\Sigma_i$, donc on aura $p_{i+1}p=p_{i+1}p'$, ce qui prouve que $\Sigma_\infty$ a la propriété filtrante exigée. % [sic] il sont -> ils sont, agreement of the verb with the plural subject

Prouvons maintenant que $\mathcal{F}$ est l'ensemble de flèches d'une sous-catégorie $K'$ de $J$, ayant $\Ob K$ comme ensemble d'objets. Il reste à établir la stabilité par composition $vu$. Mais si $u$ et $v$ sont dans $\Fl K$, ou si l'un des deux est dans $\Fl K$, l'autre non (donc de la forme $(*)$), cette stabilité est évidente. Reste à voir le cas où ni \edbr{[}$u$\edbr{]} ni $v$ sont dans $\Fl K$, donc sont tous deux de la forme $(*)$ :
\[
k \xrightarrow{\ u\ } e \xrightarrow{\ p\ } e \xrightarrow{\ v\ } k' \xrightarrow{\ u'\ } e \xrightarrow{\ p'\ } e \xrightarrow{\ v'\ } k''\ ,
\]
avec $u,v,u',v'\in\Fl K$, $p,p'\in\Sigma_\infty$. Comme $(u'v,\id_e)$ est une double-flèche de
\printedpage{140}
% head: A. Grothendieck : Dérivateurs, Ch. XVIII, Édition des Universités de Montpellier II et Paris VII
% foot: Version 17 décembre 2025
\msp{page 369}
$K$ de but $e$, on a
\[
p'u'v \;=\; p'\id_e \;=\; p'\ ,
\]
donc le composé en question s'écrit aussi
\[
k \xrightarrow{\ u\ } e \xrightarrow{\ p'p\ } e \xrightarrow{\ v'\ } k''\ ,
\]
où $p'p\in\Sigma_\infty$ puisque $\Sigma_\infty$ est stable par composition, donc il est dans $\mathcal{F}$, q.e.d.

Il est clair que $K'$ contient $K$, et ayant mêmes objets que $K$ qui admet $e$ comme plus grand élément, que $e$ est un plus grand élément de $K'$. Reste à prouver que $K'$ est filtrante. Comme $K'$ a un plus grand élément, il reste à voir que toute double-flèche $k\rightrightarrows l$ dans $K'$ peut être égalisée dans $K'$. Quitte à composer avec $l\longrightarrow e$, on peut supposer $l=e$. Or on voit que toute flèche $k\longrightarrow e$ dans $K'$ est de la forme
\[
k \xrightarrow{\ u\ } e \xrightarrow{\ p\ } e \xrightarrow{\ v\ } e\ ,
\]
avec $u,v\in\Fl K$ et $p\in\Sigma_\infty\cup\{\id_e\}$. Quitte à composer avec $p_0\in\Sigma_\infty$ fixé, cette flèche devient (puisque $p_0v=p_0\id_e=p_0$)
\[
k \xrightarrow{\ u\ } e \xrightarrow{\ p_0p\ } e\ , \qquad\quad \text{i.e. } p_0pu\ .
\]
Pour une flèche $v'p'u' : k\longrightarrow e$, on trouve $p_0p'u'$. Soit alors $q\in\Sigma_0$ tel que $q(p_0p)=q(p_0p')$, soit $r$ ($q$ existe puisque $p_0p,p_0p'\in\Sigma_\infty$). Donc $(q(p_0pu),q(p_0p'u'))=(ru,ru')$, et on a $ru=ru'$ puisque $r\in\Sigma_\infty$. Au total, $qp_0$ égalise les deux flèches $vpu$ et $v'p'u'$, q.e.d.

\msp{page 370}

\section*{4.14\quad Critères d'accessibilité pour grosses catégories}

\noindent\textbf{Théorème 4.14.1.} {\itshape Soient $\mathcal{C}$ une catégorie équivalente à une petite catégorie, $\Phi$, $\Psi$ deux petits ensembles de catégories, équivalentes à des petites catégories. Pour toute catégorie $I$, on désigne par $I_{\bullet}$ (resp. ${}_{\bullet}I$) le cône droit (resp. gauche) sur $I$, déduit de $I$ en ajoutant un objet final (resp. initial) $e=e_I$. Soit $\Lambda_1$ (resp. $\Lambda_2$) un ensemble de couples $(I,\varphi)$ (resp. $(I,\psi)$), où $I\in\Phi$ (resp. $I\in\Psi$), et où
\[
\varphi\ :\ I \longrightarrow \mathcal{C} \qquad\quad \text{(resp. }\psi\ :\ I \longrightarrow \mathcal{C}\ \ \text{)} \tag*{\textup{(4.14.1.1)}}
\]
est un foncteur. Pour toute catégorie $\mathcal{N}$ stable (\footnote[163]{Cette propriété de stabilité n'est pas nécessaire pour la définition de $\uHom^{\Lambda_2}_{\Lambda_1}(\mathcal{C},\mathcal{N})$, et il est presque sûr que le théorème est valable sans faire cette hypothèse sur $\mathcal{N}$. Cf. 4.14.16, où on voit que c'est en tous cas inutile pour $\Psi$.}), par $\varinjlim$ de type $\Phi$ et par $\varprojlim$ de type $\Psi$ on désigne par
\[
\uHom^{\Lambda_2}_{\Lambda_1}(\mathcal{C},\mathcal{N})\ \subset\ \uHom(\mathcal{C},\mathcal{N}) \tag*{\textup{(4.14.1.2)}}
\]
la sous-catégorie strictement pleine de $\uHom(\mathcal{C},\mathcal{N})$ formée des foncteurs $f:\mathcal{C}\longrightarrow\mathcal{N}$ ayant la propriété suivante :
\begin{enumerate}[leftmargin=2.5em]
\item[\textit{a)}] $\forall\ (I,\varphi)\in\Lambda_1$, le foncteur composé
\[
I_{\bullet} \xrightarrow{\ \varphi\ } \mathcal{C} \xrightarrow{\ f\ } \mathcal{N}
\]
est $I$-exact à gauche, i.e. fait de $(f\varphi)(e_I)$ une $\varinjlim\limits_{I}$ de $f\varphi|I$.
\printedpage{141}
% head: A. Grothendieck : Dérivateurs, Ch. XVIII, Édition des Universités de Montpellier II et Paris VII
% foot: Version 17 décembre 2025
\msp{page 371}
\item[\textit{b)}] Pour tout $(I,\psi)\in\Lambda_2$, le foncteur composé
\[
{}_{\bullet}I \xrightarrow{\ \psi\ } \mathcal{C} \xrightarrow{\ f\ } \mathcal{N}
\]
est $I$-exact à droite, i.e. fait de $(f\psi)(e_I)$ une limite projective de $(f\psi)|I$.
\end{enumerate}
\noindent Avec ces notations : si $\mathcal{N}$ est une catégorie accessible, il en est de même de $\uHom^{\Lambda_2}_{\Lambda_1}(\mathcal{C},\mathcal{N})$.}

\noindent\textbf{Commentaire.} On utilisera ce théorème surtout quand on a $\Lambda_2=\emptyset$ ou $\Lambda_1=\emptyset$, auquel cas on note $\uHom_{\Lambda_1}(\mathcal{C},\mathcal{N})$, resp. $\uHom^{\Lambda_2}(\mathcal{C},\mathcal{N})$, la catégorie $\uHom^{\Lambda_2}_{\Lambda_1}(\mathcal{C},\mathcal{N})$. On appliquera le théorème surtout quand $\mathcal{C}$ elle-même est stable par $\varinjlim$ de type $\Phi$ et par $\varprojlim$ de type $\Psi$, et quand $\Lambda_1$ (resp. $\Lambda_2$) est déjà formé de foncteurs $\varphi$ (resp. $\psi$) qui sont $I$-exacts à gauche (resp. à droite). Donc $\uHom^{\Lambda_2}_{\Lambda_1}(\mathcal{C},\mathcal{N})$ est alors formée des $f:\mathcal{C}\longrightarrow\mathcal{N}$ qui commutent à certaines $\varinjlim$ de type $\Phi$, et à certaines $\varprojlim$ de type $\Psi$. En particulier, on trouve :

\msp{page 372}

\noindent\textbf{Corollaire 4.14.2.} {\itshape Soient $\mathcal{C}$ une catégorie équivalente à une petite catégorie, $\mathcal{N}$ une $\mathfrak{U}$-catégorie, $\Phi$ et $\Psi$ deux petits ensembles de catégories $I$, équivalentes à des petites catégories. On suppose $\mathcal{C}$ et $\mathcal{N}$ stables par $\varinjlim$ de type $\Phi$, par $\varprojlim$ de type $\Psi$. Soit $\uHom^{\Psi}_{\Phi}(\mathcal{C},\mathcal{N})$ la sous-catégorie strictement pleine de $\uHom(\mathcal{C},\mathcal{N})$ formée des $f:\mathcal{C}\longrightarrow\mathcal{N}$ qui commutent aux $\varinjlim$ de type $\Phi$ et aux $\varprojlim$ de type $\Psi$. Alors, si $\mathcal{N}$ est accessible, il en est de même de $\uHom^{\Psi}_{\Phi}(\mathcal{C},\mathcal{N})$.}

Voici des cas particuliers remarquables de ce corollaire :

\noindent\textbf{Corollaire 4.14.3.} {\itshape Soient $\mathcal{M}$, $\mathcal{N}$ deux catégories accessibles, et $\pi_0$ un cardinal infini, tel que $\mathcal{M}$ satisfasse à $L_{\pi_0}$. Soit $\uHom_{\pi_0}(\mathcal{M},\mathcal{N})$ la sous-catégorie strictement pleine de $\uHom(\mathcal{M},\mathcal{N})$ formée des foncteurs $\pi_0$-accessibles. Alors $\uHom_{\pi_0}(\mathcal{M},\mathcal{N})$ est accessible.}

\msp{page 373}

(Le cas où $\pi_0$ est utile pour $\mathcal{M}$ a déjà été traité.)

\noindent\textsc{Démonstration}. Dans le cas général, soit $\pi>\pi_0$ un cardinal utile pour $\mathcal{M}$, de sorte qu'on a
\[
\mathcal{M} \xleftarrow{\ \sim\ } \Ind_\pi(\mathcal{C}) \qquad\quad \text{avec } \mathcal{C}=\mathcal{M}_\pi\ .
\]
On peut supposer $\mathcal{M}=\Ind_\pi(\mathcal{C})$. Le fait que $\mathcal{M}$ satisfasse à $\mathrm{L}_{\pi_0}$ implique que $\mathcal{C}$ est stable par $\varinjlim$ grandes devant $\pi_0$, et de cardinal $\leqslant\pi$. Soit $\Phi$ l'ensemble des catégories de cette nature. $\Phi$ n'est pas petit, mais il existe $\Phi_0\subset\Phi$ petit, tel que toute $I\in\Phi$ soit isomorphe à une $I_0\in\Phi_0$, ce qui suffira pour la validité de 4.14.1 et 4.14.2 (où on fera ici $\Psi=\emptyset$). On peut considérer $\uHom_{!\,\pi_0}(\mathcal{M},\mathcal{N})$ comme une sous-catégorie strictement pleine de $\uHom_{!\,\pi}(\mathcal{M},\mathcal{N})\xrightarrow{\ \sim\ }\uHom(\mathcal{C},\mathcal{N})$ (cf. 4.12.3, p. 275). Pour un foncteur $f:\mathcal{M}\longrightarrow\mathcal{N}$ qui est $\pi$-accessible, dire qu'il est même $\pi_0$-accessible revient à dire que sa restriction à $\mathcal{C}$, i.e. son image dans $\uHom(\mathcal{C},\mathcal{N})$, commute aux $\varinjlim$ de type $\Phi$. (Cf. 4.13.20, page 341 (\footnote[164]{La référence n'est pas tout à fait adéquate. Il faudrait réunir dans un paragraphe à part tous les énoncés de ce type, dont certains utilisent le théorème de représentation indicielle.}), où on prend $\alpha=\alpha_\pi$, notant que $J_{\alpha_\pi}$ grande devant $\pi_0$ puisque $\pi>\pi_0$).

\msp{page 374}

Donc on trouve
\[
\uHom_{!\,\pi_0}(\mathcal{M},\mathcal{N}) \xrightarrow{\ \sim\ } \uHom_{\Phi}(\mathcal{C},\mathcal{N})\ ,
\]
\printedpage{142}
% head: A. Grothendieck : Dérivateurs, Ch. XVIII, Édition des Universités de Montpellier II et Paris VII
% foot: Version 17 décembre 2025
et le corollaire 4.14.2 donne la conclusion voulue.

\noindent\textbf{Corollaire 4.14.4.} {\itshape Soient $\mathcal{M}$, $\mathcal{N}$ deux catégories accessibles, $\Phi$ et $\Psi$ deux ensembles des catégories comme dans 4.14.2, on suppose que $\mathcal{M}$ et $\mathcal{N}$ satisfassent $L_{\pi_0}$, et \edbr{[soient]} stables par $\varinjlim$ de type $\Phi$ et par $\varprojlim$ de type $\Psi$. Soit $\uHom^{\Psi}_{\Phi,\pi_0}(\mathcal{M},\mathcal{N})$ la sous-catégorie strictement pleine de $\uHom(\mathcal{M},\mathcal{N})$ formée des foncteurs $\pi_0$-accessibles qui commutent aux $\varinjlim$ de type $\Phi$, et aux $\varprojlim$ de type $\Psi$. Alors $\uHom^{\Psi}_{\Phi,\pi_0}(\mathcal{M},\mathcal{N})$ est accessible.}

\noindent\textsc{Démonstration}. Cette fois, on choisit $\pi>\pi_0$ tel qu'on ait $\cardred I\leqslant\pi$ pour tout $I$ dans $\Phi$ ou dans $\Psi$. La catégorie envisagée s'interprète alors (posant encore $\mathcal{C}=\mathcal{M}_\pi$) comme
\msp{page 375}
\[
\uHom^{\Psi}_{\Phi,\pi_0}(\mathcal{M},\mathcal{N}) \xrightarrow{\ \sim\ } \uHom^{\Psi}_{\Phi'}(\mathcal{C},\mathcal{N})\ ,
\]
où $\Phi'$ est formée des catégories qui sont soit dans $\Phi$, ou qui sont grandes devant $\pi_0$ et de cardinal $\leqslant\pi$. Il suffit donc d'appliquer encore 4.14.2.

\noindent\textbf{Corollaire 4.14.5.} {\itshape Considérons les trois ensembles suivants :
\begin{enumerate}[leftmargin=2.5em]
\item[\textit{a)}] $\Phi_1$ est formée de toutes les catégories équivalentes à des petites catégories.
\item[\textit{b)}] $\Phi_2$ comme a), mais en restreignant aux catégories filtrantes.
\item[\textit{c)}] $\Phi_3$ formé des catégories discrètes, et de celles qui sont grandes devant $\pi_0$ ($\pi_0$ un cardinal infini donné).
\end{enumerate}
\noindent Soit $\Phi$ l'un des ensembles de catégories $\Phi_1$, $\Phi_2$, $\Phi_3$, $\Phi_2\cup\Phi_3$ ou la réunion d'un de ceux-ci et d'un petit ensemble $\Phi_0$ de catégories essentiellement petites (\footnote[165]{On pourrait faire ici un énoncé qui englobe 4.14.4 en rajoutant un
\begin{enumerate}[leftmargin=2.5em]
\item[d)] $\Phi_4$ formé des catégories grandes devant $\pi_0$,
\end{enumerate}
\noindent et prendre $\Phi$ de la forme de $\Phi_1$, $\Phi_2$, $\Phi_4$, $\Phi_3$, $\Phi_3\cup\Phi_2$, $\Phi_3\cup\Phi_4$, ou la réunion de ceux-ci et d'un $\Phi_0$.}). Soit $\Psi$ un petit ensemble de catégories. Soient $\mathcal{M}$, $\mathcal{N}$ des catégories accessibles, stables par $\varinjlim$ de type $\Phi$, et par $\varprojlim$ de type $\Psi$. Alors $\uHom^{\Psi}_{\Phi}(\mathcal{M},\mathcal{N})$ est accessible.}

\msp{page 376}

\noindent\textsc{Démonstration}. Soit $\pi'_0\geqslant\pi_0$ un cardinal utile pour $\mathcal{M}$, $\pi$ un cardinal ayant les propriétés suivantes.
\begin{enumerate}[leftmargin=2.5em]
\item[a)] $\pi\geqslant\sup(\pi'_0,\cardred\mathcal{M}_{\pi'_0})$, $\pi>\pi_0$.
\item[b)] $\pi^{\pi'_0}=\pi$.
\item[c)] $I\in\Phi_0\cup\Psi\Longrightarrow\cardred I\leqslant\pi$.
\end{enumerate}
\noindent Soit $\mathcal{C}=\mathcal{M}_\pi$, d'où $\mathcal{M}\simeq\Ind_\pi(\mathcal{C})$. Dans tous les cas envisagés, $\Phi$ contient l'ensemble des catégories essentiellement petites qui sont grandes devant $\pi_0$, a fortiori celles qui sont grandes devant $\pi$, donc la catégorie $\uHom^{\Psi}_{\Phi}$ s'interprète comme une sous-catégorie strictement pleine de
\[
\uHom(\mathcal{C},\mathcal{N})\ .
\]
Si d'autre part $\Phi_\pi$ désigne l'ensemble des $I\in\Phi$ telles que $\card I\leqslant\pi$, on trouve
\[
\uHom^{\Psi}_{\Phi}(\mathcal{M},\mathcal{N})\ \simeq\ \uHom^{\Psi}_{\Phi_\pi}(\mathcal{C},\mathcal{N})\ ,
\]
et on applique 4.14.2.

\printedpage{143}
% head: A. Grothendieck : Dérivateurs, Ch. XVIII, Édition des Universités de Montpellier II et Paris VII
% foot: Version 17 décembre 2025
\msp{page 377}

En particulier, on trouve :

\noindent\textbf{Corollaire 4.14.6.} {\itshape Soient $\mathcal{M}$, $\mathcal{N}$ deux catégories accessibles, supposons $\mathcal{M}$ stable par petites $\varinjlim$ (i.e. $\mathcal{M}$ est une catégorie paratopos). Alors la catégorie $\uHom_{!}(\mathcal{M},\mathcal{N})$ des foncteurs $f:\mathcal{M}\longrightarrow\mathcal{N}$ qui commutent aux petites $\varinjlim$ est accessible. Si $\mathcal{N}$ est elle aussi un paratopos, $\uHom_{!}(\mathcal{M},\mathcal{N})$ est un paratopos.}

\noindent\textsc{Démonstration}. La première assertion est le cas particulier de 4.14.5 relatif au cas où $\Phi=\Phi_0$, $\Psi=\emptyset$. La deuxième assertion en résulte, puisque $\uHom_{!}(\mathcal{M},\mathcal{N})$ est visiblement stable par petites $\varinjlim$ , donc c'est un paratopos étant accessible (4.8.9, page 207).

\noindent\textbf{Corollaire 4.14.7.} {\itshape Soient $\Phi$, $\Psi$, $\mathcal{M}$ comme dans 4.14.5. Alors pour toute catégorie accessible $\mathcal{N}$, la catégorie
\msp{page 378}
\[
\uHom^{\Phi}_{\Psi}(\mathcal{M}^{\circ},\mathcal{N})\ \simeq\ (\uHom^{\Psi}_{\Phi}(\mathcal{M},\mathcal{N}^{\circ}))^{\circ}
\]
est accessible.}

\noindent\textsc{Démonstration}. On procède comme pour 4.14.5, pour expliciter
\[
\uHom^{\Psi}_{\Phi}(\mathcal{M},\mathcal{N}^{\circ}) \xrightarrow{\ \sim\ } \uHom^{\Psi}_{\Phi_\pi}(\mathcal{C},\mathcal{N}^{\circ})\ ,
\]
d'où
\[
\uHom^{\Phi}_{\Psi}(\mathcal{M}^{\circ},\mathcal{N}) \xrightarrow{\ \sim\ } \uHom^{\Phi_\pi}_{\Psi}(\mathcal{C}^{\circ},\mathcal{N})\ ,
\]
et on applique 4.14.2 à $\mathcal{C}^{\circ}$, $\Psi$, $\Phi_\pi$.

\noindent\textbf{Corollaire 4.14.8.} {\itshape Soient $\mathcal{M}$ et $\mathcal{N}$ deux catégories paratopos. Alors
\[
\mathcal{M}\boxtimes\mathcal{N}\ \simeq\ \uHom^{!}(\mathcal{M}^{\circ},\mathcal{N}) \qquad\quad \text{(cf. 4.2 et 4.3)}
\]
est un paratopos.}

En effet, en vertu de 4.14.7 (avec $\Phi=\Phi_\pi$ [?], $\Psi=\emptyset$), c'est une catégorie accessible.

\msp{page 379}

D'autre part, il est immédiat qu'elle est stable par petites $\varprojlim$ , $\mathcal{N}$ l'étant. Donc on conclut encore par 4.8.9 (page 207) que c'est une catégorie paratopos.

\noindent\textbf{Corollaire 4.14.9.} {\itshape Soient $\mathcal{M}$, $\mathcal{N}$ deux catégories paratopos, et $\uHom^{!}(\mathcal{M},\mathcal{N})$ la catégorie des foncteurs $\mathcal{M}\longrightarrow\mathcal{N}$ qui commutent aux petites $\varprojlim$ . On a alors une équivalence de catégories
\[
\begin{aligned}
\uHom^{!}(\mathcal{M},\mathcal{N}) &\ \xrightarrow{\ \sim\ }\ \uHom_{!}(\mathcal{N},\mathcal{M})^{\circ}\\
f &\ \longmapsto\ \text{adjoint à gauche de } f\ ,
\end{aligned}
\]
\noindent\textup{\textbf{!}}\ \textup{\textbf{!}} et par suite la catégorie \textup{\textbf{opposée}} à $\uHom^{!}(\mathcal{M},\mathcal{N})$ est accessible.} % Verified against PDF leaf 25: a nonverbal marginal curved mark; see the correction register.

Ceci résulte du

\noindent\textbf{Lemme 4.14.10} (pour mémoire). {\itshape Soient $\mathcal{M}$ \edbr{[une]} catégorie pseudo-topos, $f:\mathcal{M}\longrightarrow\mathcal{N}$ un foncteur. Pour que $f$ admette un adjoint à gauche (resp. à droite), il faut et il suffit que $f$ commute aux petites $\varprojlim$ (resp. aux petites $\varinjlim$ ).}


\endgroup
% END UNIT der-ch18/der-ch18_p141-145.tex
% BEGIN UNIT der-ch18/der-ch18_p146-150.tex
\begingroup\makeatletter%
\newcommand{\printedpage}[1]{\marginpar{\footnotesize p.~#1}}%
\newcommand{\msp}[1]{\marginpar{\footnotesize [#1]}}%
\newcommand{\uHom}{\underline{\mathrm{Hom}}}%
\newcommand{\uFl}{\underline{\mathrm{Fl}}}%
\newcommand{\uIs}{\underline{\mathrm{Is}}}%
\newcommand{\uMon}{\underline{\mathrm{Mon}}}%
\newcommand{\uEp}{\underline{\textup{Ép}}}%
\newcommand{\uGamma}{\underline{\Gamma}}%
\newcommand{\defeq}{\overset{\textup{déf}}{=}}%
\newcommand{\cardred}{\mathop{\textup{card\,réd}}}%
\newcommand{\Ob}{\operatorname{Ob}}%
\newcommand{\Ind}{\operatorname{Ind}}%
\newcommand{\Ens}{\mathrm{Ens}}%
\newcommand{\circled}[1]{\tikz[baseline=(c.base)]\node[circle,draw,inner sep=1pt](c){\upshape #1};}%
\makeatother\let\printedpage\derpp\let\msp\dermsp%
\let\mspin\DERmspin%


\printedpage{144}
% head: A. Grothendieck : Dérivateurs, Ch. XVIII, Édition des Universités de Montpellier II et Paris VII
% foot: Version 17 décembre 2025

\msp{page 380}
Ceci équivaut au cas particulier :

\textbf{Corollaire 4.14.11.} \textit{Soit $\mathcal{M}$ un paratopos, et $f$ un foncteur covariant (resp. contravariant) de $\mathcal{M}$ dans \textup{Ens}. Pour que $f$ soit représentable, il faut et il suffit que $f$ commute aux $\varprojlim$ .}

Ça a été vu -- le cas des contrafoncteurs étant le plus évident (cf. 4.1 -- il suffit que $\mathcal{M}$ soit un paratopos, et même que $\mathcal{M}$ admette une petite sous-catégorie pleine génératrice, cf. SGA 4 I 8.12.7 (pour l'énoncé dual). Dans les cas d'un foncteur covariant, la difficulté provient du fait que $\mathcal{M}$ n'admet pas forcément de petite sous-catégorie cogénératrice. Il faut utiliser l'existence d'une filtration cardinale (généralisée), cf. loc. cit. 8.3.12.) % [sic] Dans les cas -> Dans le cas, singular "un foncteur covariant"

Après tous ces beaux corollaires, on est mûr pour s'attaquer à la démonstration de 4.14.1 !

\msp{page 381}
\textsc{Démonstration de} 4.14.1. On sait que $\mathcal{M} = \uHom(\mathcal{C},\mathcal{N})$ est accessible (4.9.5, p. 220), et $\uHom^{\Lambda_2}_{\Lambda_1}(\mathcal{C},\mathcal{N})$ est une sous-catégorie strictement pleine de cette catégorie. Il suffit de prouver que c'est une sous-catégorie accessible dans $\mathcal{M}$ (4.13.1, p. 281). Pour tout $\xi = (I,\varphi) \in \Lambda_1$, on a un foncteur
\[
L_\xi \ : \ \mathcal{M} \ = \ \uHom(\mathcal{C},\mathcal{N}) \longrightarrow \uFl(\mathcal{N}) \ ,
\]
donné par
\[
L_\xi(f) \ = \ \text{flèche canonique } \varinjlim_{i\in I} f\varphi(i) \longrightarrow f\varphi(e_I) \ .
\]
Je dis que le foncteur $L_\xi$ est \textit{accessible.} En effet, soit $\pi_0$ un cardinal tel que $\mathcal{N}$ satisfasse à $\mathrm{L}_{\pi_0}$, il en est donc de même de $\mathcal{M} = \uHom(\mathcal{C},\mathcal{N})$ et de $\uFl(\mathcal{N}) \defeq \uHom(\boldsymbol{\Delta}_1,\mathcal{N})$, les limites en question se calculent argument par argument. Je dis que $L_\xi$ est $\pi_0$-accessible, i.e. commute aux $\varinjlim$ grandes devant $\pi_0$. Il suffit de voir qu'il en est ainsi des foncteurs composés avec les foncteurs
\[
\begin{tikzcd}
\uFl(\mathcal{N}) \arrow[r, shift left, "\sigma"] \arrow[r, shift right, "\rho"'] & \mathcal{N}
\end{tikzcd}
\]
``source'' et ``but'', i.e. des foncteurs
\[
f \longmapsto \varinjlim_{i\in I} f\varphi(i) \qquad \text{et} \qquad f \longmapsto f\varphi(e_I) \ .
\]

\msp{page 382}
Or le premier foncteur, $f \longmapsto \varinjlim\limits_{i\in I} f\varphi(i)$, est le composé de
\[
f \longmapsto f\circ\varphi_I \ : \ \uHom(\mathcal{C},\mathcal{N}) \longrightarrow \uHom(I,\mathcal{N})
\]
(lequel commute aux $\varinjlim$ qui sont représentables dans $\mathcal{N}$, et en particulier aux $\varinjlim$ grandes devant $\pi_0$) et du foncteur
\[
\varinjlim_{I} \ : \ \uHom(I,N) \longrightarrow \mathcal{N} \ , % [sic] \uHom(I,N) -> \uHom(I,\mathcal{N}), italic N printed instead of the script N used everywhere else
\]
qui lui aussi commute aux limites inductives représentables dans $\mathcal{N}$, et en particulier à celles grandes devant $\pi_0$. Quant au foncteur $f \longmapsto (f\circ\varphi)(e_I) = f(\varphi(e_I))$, il est clair qu'il commute également aux $\varinjlim$ représentables dans $\mathcal{N}$, donc il est $\pi_0$-accessible.

\printedpage{145}
% head: A. Grothendieck : Dérivateurs, Ch. XVIII, Édition des Universités de Montpellier II et Paris VII
% foot: Version 17 décembre 2025

Ceci posé, notons que
\[
\uHom_{\Lambda_1}(\mathcal{C},\mathcal{N}) \ = \ \bigcap_{\xi\in\Lambda_1} \uHom_{\{\xi\}}(\mathcal{C},\mathcal{N}) \ ,
\]
et que $\uHom_\xi(\mathcal{C},\mathcal{N})$ est l'image inverse, par le foncteur accessible $L_\xi$ précédent, de la sous-catégorie strictement pleine $\uIs(\mathcal{N})$ de $\uFl(\mathcal{N})$ formée des \textit{isomorphismes} dans $\mathcal{N}$. Si on prouve que cette dernière est accessible dans $\uFl(\mathcal{N})$, il s'ensuit que
\msp{page 383}
$\uHom_\xi(\mathcal{C},\mathcal{N})$ est accessible dans $\mathcal{M}$ (4.3.31, page 350). Donc, $\Phi$ étant petite, il s'ensuit que $\uHom_{\Lambda_1}(\mathcal{C},\mathcal{N})$ est accessible dans $\mathcal{M}$ (4.13.7, page 291).

On prouve de la même façon que $\uHom^{\Lambda_2}(\mathcal{C},\mathcal{N})$ est accessible dans $\mathcal{M} = \mathrm{Hom}(\mathcal{C},\mathcal{N})$. Ici, on doit considérer les foncteurs % [sic] \mathrm{Hom}(\mathcal{C},\mathcal{N}) -> \uHom(\mathcal{C},\mathcal{N}), printed without underline, elsewhere underlined
\[
L'_\xi \ : \ \mathcal{M} = \uHom(\mathcal{C},\mathcal{N}) \longrightarrow \uFl(\mathcal{N})
\]
définis par
\[
L'_\xi \ = \ \text{flèche canonique } f\varphi(e'_I) \longrightarrow \varprojlim_{i\in I} f\varphi(i) \ .
\]
Pour prouver que ce foncteur est accessible, on procède comme tantôt, en prenant cependant $\pi_0$ assez grand pour qu'on ait
\[
\left\{
\begin{array}{l}
I \in \Psi \ \Longrightarrow \ \cardred I \leqslant \pi_0. \\[1ex]
\mathcal{N}_{\pi_0} \text{ est génératrice dans } \mathcal{N}.
\end{array}
\right.
\]
Je dis qu'alors $L'_\xi$ est $\pi_0$-accessible : en effet, en procédant comme tantôt pour $L_\xi$, on est ramené à voir que
\[
\varinjlim_{I} \uHom(I,\mathcal{N}) \longrightarrow \mathcal{N}
\]
est $\pi$-accessible. Or cela est contenu dans SGA 4 I 9.8.

\msp{page 384}
Ainsi, on a prouvé que les deux sous-catégories strictement pleines
\[
\uHom_{\Lambda_1}(\mathcal{C},\mathcal{N}) \ , \qquad \uHom^{\Lambda_2}(\mathcal{C},\mathcal{N})
\]
de $\mathcal{M} \defeq \uHom(\mathcal{C},\mathcal{N})$ sont accessibles dans $\mathcal{M}$. Il en est donc de même de leur intersection (4.13.7)
\[
\uHom^{\Lambda_2}_{\Lambda_1}(\mathcal{C},\mathcal{N}) \ = \ \uHom_{\Lambda_1}(\mathcal{C},\mathcal{N}) \cap \uHom^{\Lambda_2}(\mathcal{C},\mathcal{N}) \ ,
\]
q.e.d.

Il reste à prouver le

\textbf{Lemme 4.14.12.} \textit{Soit $\mathcal{N}$ une catégorie accessible. Alors la catégorie $\uFl(\mathcal{N}) \defeq \uHom(\boldsymbol{\Delta}_1,\mathcal{N})$ est accessible (4.9.5), et la sous-catégorie strictement pleine $\uIs(\mathcal{N})$ formée des flèches inversibles dans $\mathcal{N}$ est accessible dans $\uFl(\mathcal{N})$. (Il en est de même de la sous-catégorie strictement pleine $\uMon(\mathcal{N})$ formée des monomorphismes dans $\mathcal{N}$.)}

\textsc{Démonstration}. Soit $\mathcal{C}$ une petite sous-catégorie strictement génératrice de $\mathcal{N}$, et considérons la petite famille des foncteurs
\[
h_x \ = \ (y \longmapsto \mathrm{Hom}(x,y)) \ : \ \mathcal{N} \longrightarrow \Ens \ ,
\]
\printedpage{146}
% head: A. Grothendieck : Dérivateurs, Ch. XVIII, Édition des Universités de Montpellier II et Paris VII
% foot: Version 17 décembre 2025
\msp{page 385}
d'où des foncteurs
\[
\uFl(h_\xi) \ : \ \uFl(\mathcal{N}) \longrightarrow \uFl(\Ens) \ . % [sic] h_\xi -> h_x, the functors defined just before are h_x
\]
Pour qu'une flèche $u : y \longrightarrow z$ dans $\mathcal{N}$ soit isomorphe (resp. monomorphe), il faut et il suffit que pour tout $x \in \Ob\mathcal{C}$, $h_x(u)$ le soit. Ainsi, $\uIs(\mathcal{N})$ (resp. $\uMon(\mathcal{N})$) apparaît comme l'intersection des images inverses, par les foncteurs $\uFl(h_\xi)$ ci-dessus, de la sous-catégorie $\uIs(\Ens)$ (resp. $\uMon(\Ens)$) de $\uFl(\Ens)$. Cela nous ramène au cas où $\mathcal{N} = (\Ens)$. Mais si $\mathcal{C}$ désigne la sous-catégorie strictement pleine de Ens formée des ensembles \textit{finis}, ceux-ci forment une famille strictement génératrice formée d'objets de \textit{présentation finie}, et comme Ens est stable par petites $\varinjlim$ , donc par $\varinjlim$ filtrantes, il s'ensuit qu'on a % [sic] h_\xi -> h_x, the functors defined just before are h_x
\[
\Ind(\mathcal{C}) \longrightarrow (\Ens) \ , \tag*{$(*)$}
\]
d'où aussitôt (SGA 4 I 8.8.5, p. 104)

\msp{page 386}
\[
\Ind(\uFl(\mathcal{C})) \xrightarrow{\ \sim\ } \uFl(\Ens) \ . \tag*{$(*)$}
\]
Donc le résultat d'accessibilité voulu résulte des deux lemmes suivants :

\textbf{Lemme 4.14.13.} \textit{L'équivalence $(*)$ précédente induit des équivalences}
\[
\begin{aligned}
\Ind(\uIs(\mathcal{C})) \quad &\xrightarrow{\ \sim\ } \quad \uIs(\Ens) \\
\Ind(\uMon(\mathcal{C})) \quad &\xrightarrow{\ \sim\ } \quad \uMon(\Ens)
\end{aligned}
\]
\textit{(où $\mathcal{C} = (\mathrm{Ensf})$ est la sous-catégorie pleine de \textup{(Ens)} formée des ensembles \textup{finis}.)}

\textbf{Lemme 4.14.14} (\footnote[166]{\textbf{N.B.} À équivalence près, on trouve les catégories accessibles $\mathcal{M}$ stables par petites $\varinjlim$ filtrantes, et telles que $\mathcal{M}_{\mathrm{pf}}$ (sous-catégorie des objets de présentation \textit{finie}) soit strictement génératrice.}). \textit{Soit $\mathcal{C}$ une catégorie équivalente à une petite catégorie. Alors :}
\begin{itemize}[nosep]
\item[\textit{a)}] \textit{$\mathcal{M} \defeq \Ind(\mathcal{C})$ est une catégorie accessible.}
\item[\textit{b)}] \textit{Pour toute sous-catégorie pleine $\mathcal{D}$ dans $\mathcal{C}$, l'image essentielle de $\Ind\mathcal{D}$ dans $\mathcal{M} = \Ind(\mathcal{C})$ est une sous-catégorie L-accessible dans $\mathcal{M}$.}
\end{itemize}
\textsc{Démonstration} des deux lemmes immédiate et ``laissée au lecteur''. Le b) dans 4.14.14 provient du fait que le foncteur $\Ind(\mathcal{D}) \longrightarrow \Ind(\mathcal{C})$ provenant d'un foncteur $\mathcal{D} \longrightarrow \mathcal{C}$ (ici un
\msp{page 387}
foncteur d'inclusion) est un L-foncteur (i.e. commute aux petites $\varinjlim$ filtrantes), et a fortiori est accessible (et même $\aleph_0$-accessible).

\textbf{4.14.15.} On ne peut s'empêcher de se demander si pour $\mathcal{N}$ accessible, la sous-catégorie strictement pleine $\uEp(\mathcal{N})$ dans $\uFl(\mathcal{N})$, formée des épimorphismes dans $\mathcal{N}$, est accessible dans $\mathcal{N}$. Il est immédiat que si $\mathcal{N}$ satisfait à $\mathrm{L}_{\pi_0}$, alors $\uEp(\mathcal{N})$ est stable dans $\uFl(\mathcal{N})$ par les $\varinjlim$ de type $\mathrm{L}_{\pi_0}$, a fortiori par facteurs directs. Elle aurait donc dû mal à ne pas être accessible dans $\uFl(\mathcal{N})$. Notons quand même : % [sic] dû mal -> du mal, "avoir du mal à"

\textbf{Corollaire 4.14.15.1} (de 4.14.13). \textit{Soit $\mathcal{N}$ une catégorie accessible stable par sommes amalgamées $y \amalg_x z$} (\footnote[167]{Je suis presque sûr que cette hypothèse est superflue. Comparer page 390.}). \textit{Alors $\uEp(\mathcal{N})$ est une sous-catégorie accessible dans $\uFl(\mathcal{N})$.}

\printedpage{147}
% head: A. Grothendieck : Dérivateurs, Ch. XVIII, Édition des Universités de Montpellier II et Paris VII
% foot: Version 17 décembre 2025

\msp{page 388}
Considérons en effet le foncteur
\[
\begin{aligned}
\gamma \ : \ \uFl(\mathcal{N}) \quad &\longrightarrow \quad \uFl(\mathcal{N}) \\
(x \longrightarrow y) \quad &\longmapsto \quad (y \amalg_x y \xrightarrow{\ \ \text{codiag}\ \ } y) \ .
\end{aligned}
\]
Ce foncteur commute aux types de $\varinjlim$ qui sont représentables dans $\mathcal{N}$, en particulier, il est accessible. D'autre part, $\uEp(\mathcal{N})$ est l'image inverse par $\gamma$ de la sous-catégorie strictement pleine $\uIs(\mathcal{N})$ de $\uFl(\mathcal{N})$. Donc le corollaire résulte du fait que par 4.14.13 cette dernière est accessible dans $\uFl(\mathcal{N})$.

\textbf{4.14.16.} Je présume que 4.14.1 reste valable sans supposer $\mathcal{N}$ stable par les $\varinjlim$ de type $\Phi$ et par les $\varprojlim$ de type $\Psi$. Si on a une essentiellement petite catégorie $I$, la question est de savoir si la sous-catégorie strictement pleine $\uHom_{\square}(I_\bullet,\mathcal{N})$ de $\uHom(I_\bullet,\mathcal{N})$, formée des foncteurs satisfaisant
\msp{page 389}
la condition a) de 4.14.1, est accessible dans $\uHom(I_\bullet,\mathcal{N})$, et de même pour $\uHom^{\square}({}_\bullet I,\mathcal{N})$. Pour ce deuxième cas, ça a l'air OK, car si $\mathcal{C}$ est une petite sous-catégorie strictement génératrice de $\mathcal{N}$, et si pour tout $x \in \Ob$ \texttt{[}$\mathcal{C}$\texttt{]} on considère encore le foncteur
\[
h_x \ : \ \mathcal{N} \longrightarrow \Ens \ ,
\]
d'où
\[
\bar{h}_x \ : \ \uHom({}_\bullet I,\mathcal{N}) \longrightarrow \uHom({}_\bullet I,\Ens) \ ,
\]
on voit que $\uHom^{\square}({}_\bullet I,\mathcal{N})$ est l'intersection des images inverses des sous-catégories $\uHom^{\square}({}_\bullet I,\Ens)$ de $\uHom({}_\bullet I,\Ens)$. Donc ces dernières étant accessibles dans $\uHom({}_\bullet I,\Ens)$, par 4.14.1 ((Ens) étant stable par petites $\varprojlim$ ), il s'ensuit bien que $\uHom^{\square}({}_\bullet I,\mathcal{N})$ l'est dans $\uHom({}_\bullet I,\mathcal{N})$, q.e.d.

Le cas de $\uHom_{\square}(I_\bullet,\mathcal{N})$ dans $\uHom(I_\bullet,\mathcal{N})$ semble plus délicat, dans le cas où
\msp{page 390}
\[
I = 
\begin{tikzcd}[sep=small]
 & {} \\
\bullet \arrow[ur] \arrow[dr] & \\
 & {}
\end{tikzcd}
\quad , \qquad \text{donc} \qquad
I = 
\begin{tikzcd}[sep=small]
 & {} \arrow[dr] & \\
\bullet \arrow[ur] \arrow[dr] & & {} \\
 & {} \arrow[ur] &
\end{tikzcd}
\quad ,
\]
si l'énoncé envisagé est vrai, on en conclura que $\uEp(\mathcal{N})$ est accessible dans $\uFl(\mathcal{N})$, pour toute catégorie accessible, car c'est l'image inverse de $\uHom_{\square}(I_\bullet,\mathcal{N})$ par le foncteur canonique (accessible)
\[
\begin{aligned}
\uFl(\mathcal{N}) \quad &\longrightarrow \quad \uHom(I_\bullet,\mathcal{N}) \\
(x \xrightarrow{\ u\ } y) \quad &\longmapsto \quad
\left(
\begin{tikzcd}[sep=small, ampersand replacement=\&]
 \& y \arrow[dr, "\mathrm{id}_y"] \& \\
x \arrow[ur, "u"] \arrow[dr, "u"'] \& \& y \\
 \& y \arrow[ur, "\mathrm{id}_y"'] \&
\end{tikzcd}
\right) \ .
\end{aligned}
\]

\printedpage{148}
% head: A. Grothendieck : Dérivateurs, Ch. XVIII, Édition des Universités de Montpellier II et Paris VII
% foot: Version 17 décembre 2025

C'est dire que la généralisation envisagée de 4.14.1 n'est pas oiseuse, mais est sans doute utile.

\msp{page 391}
\textbf{Théorème 4.14.17} (\footnote[168]{Cf. généralisation dans \texttt{[le]} théorème page 422, th. 4.14.26.}). \textit{Soit $p : \mathcal{E} \longrightarrow \mathcal{C}$ un foncteur cofibrant ou fibrant, avec $\mathcal{C}$ essentiellement petite (i.e. équivalente à une petite catégorie), et tel que les catégories fibres soient accessibles, et les foncteurs cochangement de base (cas cofibrant) ou de changement de base (cas fibrant) associés aux flèches $u$ de $\mathcal{C}$ soient accessibles. Alors la catégorie $\uGamma(\mathcal{E}/\mathcal{C})$ des sections de $\mathcal{E}$ sur $\mathcal{C}$ est accessible, et pour tout $s \in \Ob\mathcal{C}$, le foncteur ``évaluation en $s$'' $\varepsilon_s : \uGamma(\mathcal{E}/\mathcal{C}) \longrightarrow \mathcal{E}_s$ est accessible.}

\textbf{Corollaire 4.14.18.} \textit{Soit $\pi_0$ un cardinal infini tel que les catégories fibres satisfassent la condition $L_{\pi_0}$, et que les foncteurs cochangement de base (resp. changement de base) entre icelles satisfassent également $L_{\pi_0}$. (Un tel $\pi_0$ existe.)}
\begin{itemize}[nosep]
\item[\circled{a}] \textit{Alors $\mathcal{M} = \uGamma(\mathcal{E}/\mathcal{C})$ satisfait la condition $\underline{\mathrm{L}}_{\pi_0}$, et les foncteurs $\varepsilon_s$ aussi, i.e. les $\varinjlim$ grandes devant $\pi_0$ dans $\uGamma(\mathcal{E}/\mathcal{C})$ se calculent argument par argument.}
\item[\circled{b}] \textit{Soit un cardinal $\pi \geqslant \sup(\pi_0,\cardred\mathcal{C})$. Pour qu'un objet $f$ de $\mathcal{M}$ appartienne à $\mathcal{M}_\pi$, il \textup{suffit} que $\forall\ s \in \Ob C$, $f(s) \defeq \varepsilon_s(x)$ \textup{\texttt{[plutôt} $\varepsilon_s(f)$\texttt{]}} soit dans $\mathcal{E}_{s\,\pi}$. (Comparer 4.5.8, page 119.)} % [sic] \Ob C -> \Ob\mathcal{C}, italic C printed instead of script C
\item[\circled{c}] \textit{Supposons $pi_0$ choisi de façon} % [sic] pi_0 -> \pi_0, the letters "pi" printed instead of the Greek letter
\msp{page 392}

\textit{à satisfaire les conditions supplémentaires suivantes : $\pi_0$ est utile pour $\mathcal{M} = \uGamma(\mathcal{E}/\mathcal{C})$, et pour chacune des catégories fibres $\mathcal{E}_s$ ($s \in \Ob\mathcal{C}$) (un tel $\pi_0$ existe). Soit}
\[
\pi_1 \ = \ \sup(\pi_0, \cardred\mathcal{C}, \cardred\mathcal{M}_{\pi_0}) \ .
\]
\textit{Alors pour tout cardinal $\pi \geqslant \pi_1$ et tel que $\pi^{\pi_0} = \pi$, on a}
\[
\mathcal{M}_\pi \ = \ \left\{
\begin{array}{l}
\textit{sous-catégorie strictement pleine de $\mathcal{M}$ formée des $f \in$} \\
\textit{$\Ob\mathcal{M}$ telles que $f(s) \in \Ob\mathcal{E}_{s\,\pi}$ $\forall\ s \in S$ \textup{\texttt{[plutôt}} $s \in$} \\
\textit{$\Ob C$\textup{\texttt{]}}.} % [sic] \Ob C -> \Ob\mathcal{C}, italic C printed instead of script C
\end{array}
\right.
\]
\textit{Si de plus $\pi$ est choisi tel qu'on ait}
\[
\pi \ \geqslant \ \sup_{s\in\Ob\mathcal{C}} \cardred\mathcal{E}_{s\,\pi_0} \ ,
\]
\textit{alors la sous-catégorie strictement pleine $\mathcal{E}(\pi)$ de $\mathcal{E}$ telle que $\Ob\mathcal{E}(\pi) = \bigcup_{s\in\Ob\mathcal{C}} \Ob\mathcal{E}_{s\,\pi}$ est une sous-catégorie $\mathcal{C}$-cofibrée (resp. $\mathcal{C}$-fibrée) de $\mathcal{E}$, et on a}
\[
\uGamma(\mathcal{E}/\mathcal{C})_\pi \xleftarrow{\ \sim\ } \uGamma(\mathcal{E}(\pi)/C), % [sic] C -> \mathcal{C}, italic C printed instead of script C
\]
\textit{donc ($\pi$ étant utile pour $\mathcal{E}$) on a une équivalence de catégories}
\[
\uGamma(\mathcal{E}/\mathcal{C}) \xleftarrow{\ \sim\ } \Ind_\pi(\uGamma(\mathcal{E}(\pi)/\mathcal{C})) \ .
\]
\end{itemize}


\endgroup
% END UNIT der-ch18/der-ch18_p146-150.tex
% BEGIN UNIT der-ch18/der-ch18_p151-155.tex
\begingroup\makeatletter%
\newcommand{\printedpage}[1]{\marginpar{\footnotesize p.~#1}}%
\newcommand{\msp}[1]{\marginpar{\footnotesize [#1]}}%
\newcommand{\uGamma}{\underline{\Gamma}}%
\newcommand{\uFl}{\underline{\mathrm{Fl}}}%
\newcommand{\uIs}{\underline{\mathrm{Is}}}%
\newcommand{\uHom}{\underline{\mathrm{Hom}}}%
\newcommand{\Ob}{\mathrm{Ob}\,}%
\newcommand{\Fl}{\mathrm{Fl}\,}%
\newcommand{\id}{\mathrm{id}}%
\newcommand{\pr}{\mathrm{pr}}%
\newcommand{\diag}{\mathrm{diag}\,}%
\newcommand{\cardred}{\operatorname{card\,r\acute{e}d}}%
\newcommand{\cocart}{\mathrm{cocart}}%
\newcommand{\cart}{\mathrm{cart}}%
\newcommand{\cc}{\mathrm{c}}%
\newcommand{\circled}[1]{\tikz[baseline=(c.base)]\node[draw,circle,inner sep=0.8pt](c){\upshape #1};}%
\makeatother\let\printedpage\derpp\let\msp\dermsp%
\let\mspin\DERmspin%


\printedpage{149}
% head: A. Grothendieck : Dérivateurs, Ch. XVIII, Édition des Universités de Montpellier II et Paris VII

\msp{page 393}
Toutes ces assertions se prouvent facilement dans l’ordre où elles sont données, en utilisant 4.14.17 et les résultats généraux déjà prouvés. Bien sûr, on peut commencer par remplacer $\mathcal{C}$ par une sous-catégorie réduite $\mathcal{C}_0$ équivalente et $\mathcal{E}$ par $\mathcal{E}_0 = \mathcal{E}|\mathcal{C}$ \texttt{[plutôt} $\mathcal{E}|\mathcal{C}_0$\texttt{]}, de façon à se réduire (dans b) et c)) au cas où $\operatorname{card}\mathcal{C} \leqslant \pi$.

\textbf{Corollaire 4.14.19.} \textit{Sous les conditions de 4.14.17, la sous-catégorie $\uGamma_{\cocart}(\mathcal{E}/\mathcal{C})$ de $\mathcal{M} = \uGamma(\mathcal{E}/\mathcal{C})$ (cas $p$ cofibrant), resp. la sous-catégorie $\uGamma_{\cart}(\mathcal{E}/\mathcal{C})$ (cas $p$ fibrant) est accessible dans $\mathcal{M}$, donc accessible et le foncteur d’inclusion} \texttt{[est]} \textit{accessible.}

(\textbf{N.B.} Ces deux catégories sont strictement pleines dans $\uGamma(\mathcal{E}/\mathcal{C})$, donc l’assertion qu’elles soient accessibles dans celle-ci \textit{équivaut} à la dernière assertion.)

\textsc{Démonstration de} 4.14.19. On peut supposer $\mathcal{C}$ réduite, donc petite. Considérons le cas où $p$ est cofibrant. Pour toute flèche $\lambda : s \longrightarrow t$ dans $\mathcal{C}$, considérons le foncteur
\[
\varphi_\lambda \ : \ \mathcal{M} \longrightarrow \uFl(\mathcal{E}_t)
\]
défini par
\msp{page 394}
\[
\varphi_\lambda(f) \ = \ (\lambda_*(f(s)) \xrightarrow{f(\lambda)^\natural} f(t)) \ .
\]
C’est un foncteur accessible, car son composé avec les foncteurs source et but sont les foncteurs $\lambda_*\varepsilon_s$ et $\varepsilon_t$, donc accessibles, $\varepsilon_s$, $\varepsilon_t$ et $\lambda_*$ étant accessibles. Considérons la sous-catégorie strictement pleine
\[
\uIs(\mathcal{E}_t) \ \subset \ \uFl(\mathcal{E}_t) \ ,
\]
on sait qu’elle est accessible dans $\uFl(\mathcal{E}_t)$ (4.14.12). Donc son image inverse par $\varphi_\lambda$ est accessible dans $\mathcal{M}$, et par suite aussi
\[
\uGamma_{\cart}(\mathcal{E}/\mathcal{C}) \ \overset{\text{déf}}{=} \ \bigcap_{\lambda\in\Ob\mathcal{C}} \varphi_\lambda^{-1}(\uIs(\mathcal{E}_t))
\]
% [sic] \uGamma_{\cart} -> \uGamma_{\cocart}, the case treated is p cofibrant
% [sic] \lambda\in\Ob\mathcal{C} -> \lambda\in\Fl\mathcal{C}, \lambda runs over arrows of \mathcal{C}
comme intersection d’une \textit{petite} famille de sous-catégories accessibles dans $\mathcal{M}$, q.e.d. Le cas où $p$ est fibrant se traite de façon similaire, en introduisant les foncteurs
\[
\psi_\lambda \ : \ \mathcal{M} = \uGamma(\mathcal{E}/\mathcal{C}) \longrightarrow \uFl(\mathcal{E}_s)
\]
donnés par
\[
\psi_\lambda(f) \ = \ (f(s) \xrightarrow{f(\lambda)^\flat} \lambda^* f(t)) \ .
\]

\msp{page 395}
\textbf{Corollaire 4.14.20.} \textit{Soit $\pi_0$ un cardinal infini choisi comme dans 4.14.18 a).}
\begin{itemize}[leftmargin=2.5em]
\item[\circled{a}] \textit{Alors $\mathcal{M}_{\cc} = \uGamma_{\cocart}(\mathcal{E}/\mathcal{C})$ (resp. $\mathcal{M}_{\cc} = \uGamma_{\cart}$) est $L_{\pi_0}$-stable dans $\mathcal{M} = \uGamma(\mathcal{E}/\mathcal{C})$, ou ce qui revient au même (comme elle est strictement pleine), elle satisfait $L_{\pi_0}$, et le foncteur d’inclusion satisfait à $L_{\pi_0}$. Les $\varinjlim$ grandes devant $\pi_0$ dans cette catégorie se calculent argument par argument.}
\item[\circled{b}] \textit{Pour tout cardinal $\pi \geqslant \pi_0$, tout objet de $\mathcal{M}_{\cc}$ qui est $\pi$-accessible dans $\mathcal{M}$ l’est (tautologiquement) dans $\mathcal{M}_{\cc}$. Si $\pi \geqslant \sup(\pi_0, \cardred\mathcal{C})$, alors (par 4.14.18 b) tout objet $f$ de $\mathcal{M}_{\cc}$ tel que $f(s) \in \Ob\mathcal{E}_{s\,\pi}$ $\forall\ s \in S$} \texttt{[plutôt} $s \in \Ob C$\texttt{]} \textit{est dans $\mathcal{M}_\pi$} \textup{et} \textit{dans $\mathcal{M}_{\cc\,\pi}$.}
% foot: Version 17 décembre 2025

\printedpage{150}
% head: A. Grothendieck : Dérivateurs, Ch. XVIII, Édition des Universités de Montpellier II et Paris VII

\item[\circled{c}] \textit{Supposons $\pi_0$ choisi comme dans 4.14.18 c, et que de plus $\mathcal{M}_{\cc}$ soit $\pi_0$-accessible dans $\mathcal{M}$. (Un tel $\pi_0$ existe.) On a donc, pour $\pi \geqslant \pi_1 = \sup(\pi_0, \cardred\mathcal{C}, \cardred\mathcal{M}_\pi)$ tel que $\pi^{\pi_0} = \pi$,}
\[
(\mathcal{M}_{\cc})_\pi \ = \ \mathcal{M}_\pi{\cap}\mathcal{M}_{\cc} \ = \quad
\begin{array}{l}
\textit{Sous-catégorie strictement pleine de $\mathcal{M}_{\cc}$ formée des $f$}\\
\textit{telles que $f(s) \in \Ob\mathcal{E}_{s\,\pi}$ $\forall\ s \in \Ob\mathcal{C}$.}
\end{array}
\]

\msp{page 396}
\textit{Si de plus $\pi$ est choisi tel qu’on ait $\pi \geqslant \sup_{s\in\Ob\mathcal{C}} \cardred\mathcal{E}_{s\,\pi_0}$, alors}
\[
\left\{
\begin{array}{rrcl}
 & \mathcal{M}_{\cc\,\pi} & \xleftarrow{\ \sim\ } & \uGamma_{\cocart}(\mathcal{E}(\pi)/\mathcal{C})\\[1ex]
\textit{resp.} & \mathcal{M}_{\cc\,\pi} & \xleftarrow{\ \sim\ } & \uGamma_{\cart}(\mathcal{E}(\pi)/\mathcal{C}) \ ,
\end{array}
\right.
\]
\textit{donc ($\pi$ étant utile pour $\mathcal{M}_{\cc}$) une équivalence de catégories}
\[
\begin{array}{rrcl}
 & \uGamma_{\cocart}(\mathcal{E}/\mathcal{C}) & \xleftarrow{\ \sim\ } & \mathrm{Ind}_\pi(\uGamma_{\cocart}(\mathcal{E}(\pi)/\mathcal{C}))\\
\textit{resp.}\quad & \uGamma_{\cart}(\mathcal{E}/\mathcal{C}) & \xleftarrow{\ \sim\ } & \mathrm{Ind}_\pi(\uGamma_{\cart}(\mathcal{E}(\pi)/\mathcal{C})).
\end{array}
\]
\end{itemize}

\textsc{Démonstration}. Tout est clair dans l’ordre où c’est donné en utilisant les deux corollaires précédents.

\textbf{Corollaire 4.14.21.} \textit{Considérons un diagramme de foncteurs}
\[
\begin{tikzcd}[column sep=small, row sep=small]
\mathcal{E}_1 \arrow[dr, "p_1"] & & \mathcal{E}_2 \arrow[dl, "p_2"'] \\
& \mathcal{E}_0 \ , &
\end{tikzcd}
\tag*{$(*)$}
\]
\textit{les catégories et les foncteurs étant accessibles.}

\msp{page 397}
\begin{itemize}[leftmargin=2.5em]
\item[\textit{a)}] \textit{Le produit 2-fibré}
\[
\mathcal{M}_0 \ = \ \mathcal{E}_1 \overset{2}{\times}_{\mathcal{E}_0} \mathcal{E}_2
\]
\textit{est accessible, et les foncteurs canoniques $\mathcal{M} \longrightarrow \mathcal{E}_1$, $\mathcal{M} \longrightarrow \mathcal{E}_2$ sont accessibles.} % [sic] \mathcal{M} \longrightarrow \mathcal{E}_1, \mathcal{M} \longrightarrow \mathcal{E}_2 -> \mathcal{M}_0 \longrightarrow \mathcal{E}_1, \mathcal{M}_0 \longrightarrow \mathcal{E}_2, assertion a) concerns \mathcal{M}_0
\item[\textit{b)}] (\footnote[169]{Ceci est contenu dans 4.14.22.1 ci-dessus (p. 406), appliqué à $\Pi = \mathcal{E}_1{\times}\mathcal{E}_2$, $\Pi' = \mathcal{E}_0$, et $F, G : \Pi \rightrightarrows \Pi'$ donnés par $F = p_1\circ\pr_1$ et $G = p_2\circ\pr_2$. D’autre part, a) et c) résultent facilement de b).}). \textit{La catégorie $\mathcal{M}$ des triples $(x_1,x_2,u)$ avec $x_1 \in \Ob\mathcal{E}_1$, $x_2 \in \Ob\mathcal{E}_2$, et}
\[
u \ : \ p_1(x_1) \ \longrightarrow \ p_2(x_2) \ ,
\]
\textit{est accessible, et les foncteurs canoniques $\mathcal{M} \longrightarrow \mathcal{E}_1$, $\mathcal{M} \longrightarrow \mathcal{E}_2$ sont accessibles.}
\item[\textit{c)}] \textit{La sous-catégorie strictement pleine $\mathcal{M}_0$ de $\mathcal{M}$ est accessible dans $\mathcal{M}$.}
\end{itemize}

\noindent\textsc{Démonstration}. Soit $\mathcal{C}$ la catégorie
\[
\begin{tikzcd}[column sep=small]
1 \arrow[dr] & & 2 \arrow[dl] \\
& 0 \ , &
\end{tikzcd}
\]
et $\mathcal{E}$ la catégorie cofibrée sur $\mathcal{C}$ définie par le diagramme $(*)$. On est sous les conditions de 4.14.17. La catégorie $\Gamma = \uGamma(\mathcal{E}/\mathcal{C})$ s’interprète
% foot: Version 17 décembre 2025
\printedpage{151}
% head: A. Grothendieck : Dérivateurs, Ch. XVIII, Édition des Universités de Montpellier II et Paris VII
\msp{page 398}
comme la catégorie des quintuplets $(x_0,x_1,x_2,u_1,u_2)$, avec $x_i \in \Ob\mathcal{E}$ pour $i \in \{0,1,2\}$, et $u_1, u_2 \in \Fl\mathcal{E}_0$,
\[
\begin{tikzcd}[column sep=small, row sep=small]
p_1(x_1) \arrow[dr, "u_1"] & & p_2(x_2) \arrow[dl, "u_2"'] \\
& x_0 \ . &
\end{tikzcd}
\tag*{$(**)$}
\]
En vertu de 4.14.18, cette catégorie $\uGamma$ est accessible. La catégorie $\mathcal{M}$ est équivalente à la sous-catégorie strictement pleine $\Gamma_0$ de $\Gamma$, formée des quintuplets pour lesquels on a $u_2 \in \uIs\mathcal{E}_0$. On a en effet des foncteurs quasi-inverses l’un de l’autre % [sic] \Gamma_0 -> \Gamma_1, the subcategory is called \Gamma_1 in what follows
\[
\begin{tikzcd}
\mathcal{M} \arrow[r, shift left, "\alpha"] & \Gamma_1 \arrow[l, shift left, "\beta"]
\end{tikzcd}
\]
donnés par
\[
\begin{aligned}
\alpha(x_1,x_2,u) \ &= \ (p_2(x_2),x_1,x_2,u,\id_{p_2(x_2)})\\
\beta(x_0,x_1,x_2,u_1,u_2) \ &= \ (x_1,x_2,u_2^{-1}u_1) \ ,
\end{aligned}
\]
on a alors
\[
\begin{aligned}
\beta\alpha \quad &= \quad \id_{\mathcal{M}}\\
\alpha\beta \quad &\xrightarrow{\ \sim\ } \quad \id_{\Gamma_1} \ .
\end{aligned}
\]
D’autre part, $\Gamma_1$ est évidemment accessible dans $\mathcal{M}$, comme image inverse de $\uIs(\mathcal{E}_0)$ (qui % [sic] accessible dans \mathcal{M} -> accessible dans \Gamma, \Gamma_1 is a subcategory of \Gamma
\msp{page 399}
est accessible dans $\uFl(\mathcal{E}_0)$ par le foncteur accessible
\[
\uGamma \ \longrightarrow \ \uFl(\mathcal{E}_0)
\]
donné par $(x_1, x_1, x_2, u_1, u_2) \longmapsto u_2$. Cela prouve que $\mathcal{M}$, qui est équivalente à $\Gamma_1$, est également accessible, et l’accessibilité du foncteur canonique $\mathcal{M} \longrightarrow \mathcal{E}_i$ ($i \in \{1,2\}$) résulte de celle de $\Gamma \longrightarrow \mathcal{E}_i$, impliquant celle du composé $\Gamma_1 \hookrightarrow \Gamma \longrightarrow \mathcal{E}_i$. Cela prouve donc b). Pour prouver a), il suffit de voir que $\mathcal{M}_0$ est accessible dans $\mathcal{M}$, or c’est encore l’image inverse de $\uIs(\mathcal{E}_0) \subset \uFl(\mathcal{E}_0)$ par le foncteur accessible canonique % [sic] (x_1, x_1, x_2, u_1, u_2) -> (x_0, x_1, x_2, u_1, u_2), the quintuplets begin with x_0
\[
\mathcal{M} \ \longrightarrow \ \uFl(\mathcal{E}_0) \ ,
\]
donc $\mathcal{M}$ est bien accessible dans $\mathcal{M}$, comme image inverse d’une sous-catégorie accessible dans $\uFl(\mathcal{E})$. On peut d’ailleurs aussi interpréter $\mathcal{M}$, à équivalence de catégories près, comme % [sic] donc \mathcal{M} est bien accessible -> donc \mathcal{M}_0 est bien accessible, \mathcal{M}_0 is the subcategory in question
% [sic] \uFl(\mathcal{E}) -> \uFl(\mathcal{E}_0), as in the preceding display
% [sic] interpréter \mathcal{M} -> interpréter \mathcal{M}_0, as in the following display
\msp{page 400}
\[
\mathcal{M}_0 \ \simeq \ \uGamma_{\cocart}(\mathcal{E}/\mathcal{C}) \ ,
\]
les foncteurs canoniques $\mathcal{M}_0 \longrightarrow \mathcal{E}_i$ ($i \in \{0,1\}$) s’interprétant comme des foncteurs évaluation, ce qui donne a) directement, comme cas particulier de 4.14.19. % [sic] i \in \{0,1\} -> i \in \{1,2\}, the canonical functors of a) go to \mathcal{E}_1, \mathcal{E}_2
% foot: Version 17 décembre 2025

\printedpage{152}
% head: A. Grothendieck : Dérivateurs, Ch. XVIII, Édition des Universités de Montpellier II et Paris VII

\textbf{4.14.22.} \textsc{Démonstration de} 4.14.17. Le cas où $p$ est cofibrant a déjà été démontré (th. 4.10.8, p. 263). Nous allons voir que le cas où $p$ est fibrant se ramène au théorème précédent (dans un cas très particulier), en utilisant à fond les résultats généraux sur les parties accessibles d’une catégorie accessible. Notons que les assertions a) et b) du corollaire 4.14.18 se prouvent de façon AQT \texttt{[âne qui trotte]} sans avoir à faire appel à 4.14.17. Le a) implique déjà que les foncteurs évaluation $\varepsilon_s : \mathcal{M} \longrightarrow \mathcal{E}_s$ sont accessibles. Il reste donc à prouver
\msp{page 401}
que $\mathcal{M} = \uGamma(\mathcal{E}/\mathcal{C})$ est accessible, dans le cas où $p$ est \textit{fibrant}. On supposera de plus que $\mathcal{C}$ est petite (quitte à remplacer $\mathcal{C}$ par une petite sous-catégorie pleine équivalente).

Soient
\[
\begin{aligned}
\Pi \ &= \ \prod\nolimits_{s\in\Ob\mathcal{C}} \mathcal{E}_s\\
\Pi' \ &= \ \prod\nolimits_{\lambda\in\Fl\mathcal{C}} \mathcal{E}_{s(\lambda)} \ ,
\end{aligned}
\]
où $s(\lambda)$ désigne la \textit{source} $s$ de la flèche $\lambda : s \longrightarrow t$ dans $\mathcal{C}$. $\Pi$, $\Pi'$ sont des catégories accessibles, comme (petits) produits de catégories accessibles. (Cas particulier trivial de 4.10.8, pour une catégorie d’indices discrète.) On va définir deux foncteurs
\[
\Pi \ \overset{F}{\underset{G}{\rightrightarrows}} \ \Pi'
\]
en définissant pour tout $\lambda : s_0 \longrightarrow t_0$ dans $\Fl\mathcal{C}$ deux foncteurs
\msp{page 402}
\[
\left\{
\begin{array}{lcll}
\Pi & \overset{F_\lambda}{\underset{G_\lambda}{\rightrightarrows}} & \mathcal{E}_{s(\lambda)=s_0} & \\
F_\lambda & = & \pr_{s_0} & \text{projection de } \Pi = \prod\nolimits_s \mathcal{E}_s \text{ sur le facteur } \mathcal{E}_{s_0}\\
G_\lambda & = & \lambda^*\pr_{s_1} & \text{composé } \Pi \xrightarrow{\pr_{s_1}} \mathcal{E}_{s_1} \xrightarrow{\lambda^*} \mathcal{E}_{s_0}\\
 & & & \text{\texttt{[plutôt} } \lambda^*\pr_{t_0} \text{ \texttt{resp.} } \Pi \xrightarrow{\pr_{t_0}} \mathcal{E}_{t_0} \xrightarrow{\lambda^*} \mathcal{E}_{s_0}\text{\texttt{]}}
\end{array}
\right.
\qquad \bigl(\mbox{\footnotemark[170]}, \mbox{\footnotemark[171]}\bigr) \ .
\]
\footnotetext[170]{Il est immédiat que les $F_\lambda$, $G_\lambda$, donc $F$ et $G$, sont des foncteurs \textit{accessibles}.}
\footnotetext[171]{\textbf{N.B.} On suppose choisi pour tout $\lambda \in \Fl\mathcal{C}$ un foncteur image inverse $\lambda^* : \mathcal{E}_t \longrightarrow \mathcal{E}_s$. On suppose que si $\lambda$ est une flèche identique, $\id_s$, \texttt{[on]} choisit $\lambda^* = \id_{\mathcal{E}_s}$.}

Soit $f \in \mathcal{M} = \uGamma(\mathcal{E}/\mathcal{C})$, alors pour tout $\lambda \in \Fl\mathcal{C}$, $f$ définit une flèche
\[
\alpha_\lambda(f) \ : \quad \underbrace{f(s)}_{= F_\lambda(\varepsilon(f))} \quad \xrightarrow{f(\lambda)^\flat} \quad \underbrace{\lambda^*(f(t))}_{= G_\lambda(\varepsilon(f))} \ ,
\]
\textit{fonctorielle} en $f$. On a un foncteur
\[
\varepsilon \ = \ (\varepsilon_s)_{s\in\Ob\mathcal{C}} \ : \ \mathcal{M} \ \longrightarrow \ \Pi \ ,
\]
et pour tout $\lambda \in \Fl\mathcal{C}$ un homomorphisme fonctoriel
\[
\alpha_\lambda \ : \ F_\lambda \circ \varepsilon \ \longrightarrow \ G_\lambda \circ \varepsilon \ ,
\]
d’où un homomorphisme fonctoriel
\[
\alpha \ = \ (\alpha_\lambda)_{\lambda\in\Fl\mathcal{C}} \ : \ F \circ \varepsilon \ \longrightarrow \ G \circ \varepsilon \ .
\]
Soit alors $\mathcal{M}'$ la catégorie suivante :
% foot: Version 17 décembre 2025
\printedpage{153}
% head: A. Grothendieck : Dérivateurs, Ch. XVIII, Édition des Universités de Montpellier II et Paris VII
\msp{page 403}
\[
\mathcal{M}' \quad = \quad \parbox[t]{0.5\textwidth}{catégorie des couples $(X,u)$, avec $X$ dans $\Ob\Pi$, et $u : F(X) \longrightarrow G(X)$ une flèche dans $\Pi'$.}
\]
L’homomorphisme fonctoriel $\alpha : F \circ \varepsilon \longrightarrow G \circ \varepsilon$ définit un foncteur
\[
\mathcal{M} \ \longrightarrow \ \mathcal{M}' \ ,
\]
et il est immédiat que ce foncteur est pleinement fidèle. En fait, la catégorie $\mathcal{M}'$ est isomorphe à la catégorie formée des couples $(f_0,f_1)$ de deux applications
\[
\begin{array}{lcccccc}
f_0 : & \Ob\mathcal{C} & \longrightarrow & \Ob\mathcal{E} & s & \longmapsto & f_0(s)\\[1ex]
f_1 : & \Fl\mathcal{C} & \longrightarrow & \Fl\mathcal{E} & \lambda & \longmapsto & f_1(\lambda)
\end{array}
\]
de telle façon que $f_1$ et $f_0$ soient compatibles avec les applications source et but, i.e. qu’ils définissent un homomorphisme de \textit{diagrammes} $f : \diag\mathcal{C} \longrightarrow \diag\mathcal{E}$, et que cet homomorphisme soit inverse à droite de $\mathrm{diag}(p) : \diag\mathcal{E} \longrightarrow \diag\mathcal{C}$.

\msp{page 404}
Le foncteur précédent $\mathcal{M} \hookrightarrow \mathcal{M}'$ s’interprète comme le foncteur qui associe à toute sections $f$ de $\mathcal{E}$ sur $\mathcal{C}$, le morphisme de diagrammes correspondant. C’est donc l’inclusion d’une sous-catégorie \textit{strictement pleine} de $\mathcal{M}'$. Donc pour prouver que $\mathcal{M}$ est accessible, il suffit de prouver : % [sic] toute sections -> toute section, agreement of toute
\begin{itemize}[leftmargin=2.5em]
\item[a)] $\mathcal{M}'$ est accessible.
\item[b)] $\mathcal{M}$ est accessible dans $\mathcal{M}'$.
\end{itemize}
Prouvons d’abord b) (en admettant a). La sous-catégorie $\mathcal{M}$ de $\mathcal{M}'$ est définie par deux types des conditions, sur un objet $f = (f_0,f_1)$ de $\mathcal{M}'$ :
\begin{itemize}[leftmargin=2.5em]
\item[1°)] $\forall\ s \in \Ob\mathcal{C}$, $f_1(\id_s) = \id_{f_0(s)}$.
\item[2°)] Si $s \xrightarrow{\lambda} t \xrightarrow{\mu} u$ sont deux flèches composables dans $\mathcal{C}$, alors
\[
f_1(\mu\lambda) \ = \ f_1(\mu)f_1(\lambda) \ .
\]
\end{itemize}
Ces conditions s’explicitent l’une et \texttt{[l’autre par]} l’égalité d’une double-flèche dans $\mathcal{E}_s$,
\msp{page 405}
savoir
\[
f_0(s) \overset{f_1(\id_s)}{\underset{\id_{f_0(s)}}{\rightrightarrows}} f_0(s)
\]
dans le cas 1°), et
\[
f_0(s) \overset{f_1(\mu\lambda)^\flat}{\underset{(f_1(\mu)f_1(\lambda))^\flat}{\rightrightarrows}} (\mu\lambda)^*(f_0(u))
\]
dans le cas 2°). Dans l’un et l’autre cas, on a un foncteur
\[
\mathcal{M}' \ \xrightarrow{\gamma_s \text{ resp. } \gamma_{\lambda,\mu}} \ \uHom(I,\mathcal{E}_s)
\]
% foot: Version 17 décembre 2025


\endgroup
% END UNIT der-ch18/der-ch18_p151-155.tex
% BEGIN UNIT der-ch18/der-ch18_p156-160.tex
\begingroup\makeatletter%
\newcommand{\printedpage}[1]{\marginpar{\footnotesize p.~#1}}%
\newcommand{\msp}[1]{\marginpar{\footnotesize [#1]}}%
\newcommand{\uHom}{\underline{\mathrm{Hom}}}%
\newcommand{\uFl}{\underline{\mathrm{Fl}}}%
\newcommand{\uIs}{\underline{\mathrm{Is}}}%
\newcommand{\uGamma}{\underline{\Gamma}}%
\newcommand{\bDelta}{\boldsymbol{\Delta}}%
\newcommand{\Ens}{\mathrm{Ens}}%
\newcommand{\Ob}{\mathrm{Ob}\,}%
\newcommand{\Fl}{\mathrm{Fl}}%
\newcommand{\Hom}{\mathrm{Hom}}%
\makeatother\let\printedpage\derpp\let\msp\dermsp%
\let\mspin\DERmspin%


\printedpage{154}
% head: A. Grothendieck : Dérivateurs, Ch. XVIII, Édition des Universités de Montpellier II et Paris VII
dans la catégorie des double-flèches de $\mathcal{E}_s$ ($I$ désignant la catégorie double-flèche type $0\rightrightarrows 1$), foncteur qui est (trivialement) accessible. Soit alors $I\xrightarrow{\;i\;}\bDelta_1$ le foncteur canonique, et
\[
\underbrace{\uHom(\bDelta_1,\mathcal{E}_s)}_{\simeq\,\uFl(\mathcal{E}_s)}\ \hookrightarrow\ \uHom(I,\mathcal{E}_s)
\]
le foncteur $u\longmapsto u\circ i$, qui est un foncteur strictement pleine fidèle, ayant comme image la sous-catégorie des double-flèches $(u,v)$ telles que $u=v$. Comme le foncteur d'inclusion est accessible et $\uHom(\bDelta_1,\mathcal{E}_s)$ et $\uHom(I,\mathcal{E}_s)$ accessibles
\msp{page 406}
(cas particulier de 4.14.1, ou aussi de 4.10.8!), on voit que la sous-catégorie strictement pleine de $\uHom(I,\mathcal{E}_s)$ des double-flèches égalisées dans $\mathcal{E}_s$, est accessible dans $\uHom(I,\mathcal{E}_s)$. Donc, son image inverse dans $\mathcal{M}'$ est accessible. Donc l'intersection de ces images inverses pour $s\in\Ob\mathcal{C}$ et $(\lambda,\mu)\in\Fl_2(\mathcal{C})$ variables, l'est également ($\mathcal{C}$ étant petite).

Donc tout revient à prouver que $\mathcal{M}'$ est accessible, ce qui est contenu dans le

% vertical bar printed in the left margin alongside Lemme 4.14.22.1
\textbf{Lemme 4.14.22.1.} \textit{Soient $\Pi$, $\Pi'$ deux catégories accessibles, $F,G:\Pi\rightrightarrows\Pi'$ deux foncteurs accessibles, et $\mathcal{M}'$ la catégorie des couples $(X,u)$, avec $X\in\Ob\Pi$ et $u:F(x)\longrightarrow G(x)$ une flèche de $\Pi'$. Alors $\mathcal{M}'$ est accessible, et le foncteur canonique $\mathcal{M}'\longrightarrow\Pi$ est accessible.} % [sic] (X,u), X \in Ob Pi ... F(x) -> the same letter (X or x) throughout, capital X and lowercase x used for the same object

\textsc{Démonstration du Lemme.} La donnée de $\Pi$, $\Pi'$, $F$, $G$ définit une catégorie cofibrée $\mathcal{F}$ sur la catégorie $I=(0\rightrightarrows 1)$
\msp{page 407}
précédente. La catégorie $\uGamma(\mathcal{F}/I)$ s'interprète, à isomorphisme près, comme la catégorie des triples $(x,u,v)$ avec $x\in\Ob\Pi$, $x'\in\Ob\Pi'$, et $u$ et $v$ des flèches
\[
\begin{tikzcd}[column sep=small]
F(x) \arrow[dr, "u"'] & & G(x) \arrow[dl, "v"] \\
& x' &
\end{tikzcd}
\]
dans $\Pi'$. La catégorie $\mathcal{M}'$ est équivalente à la sous-catégorie strictement pleine de $\uGamma(\mathcal{F}/I)$, formée des triples $(x,u,v)$ tels que $v$ soit un isomorphisme, i.e. à la sous-catégorie de $\uGamma(\mathcal{F}/I)$, image inverse de $\uIs(\Pi')$ par le foncteur $(x,u,v)\longmapsto v$,
\[
\uGamma(\mathcal{F}/I)\ \longrightarrow\ \uFl(\Pi')\ .
\]
Or par 4.10.8, $\uGamma(\mathcal{F}/I)$ est accessible, et le foncteur précédent est bien sûr accessible. Enfin, par 4.14.12,
\msp{page 408}
$\uIs(\Pi')$ est accessible dans la catégorie accessible $\uFl(\Pi')$. Donc son image inverse $\mathcal{M}'$ dans $\uGamma(\mathcal{F}/I)$ est accessible dans $\uGamma(\mathcal{F}/I)$, donc accessible. Donc il en est de même de la catégorie équivalente, q.e.d.

Ainsi, grâce aux résultats généraux sur les parties accessibles des catégories accessibles, le théorème 4.14.17, dans le cas où $p$ est fibrant, est ramené aux cas cofibrant, dans le cas très particulier où la catégorie d'indices est $I=(0\rightrightarrows 1)$. On pourrait d'ailleurs faire une réduction similaire dans le cas où $p$ est cofibrant, par une démonstration pratiquement % [sic] aux cas cofibrant -> au cas cofibrant, singular noun
% foot: Version 17 décembre 2025
\printedpage{155}
% head: A. Grothendieck : Dérivateurs, Ch. XVIII, Édition des Universités de Montpellier II et Paris VII
identique. En d'autres termes, le théorème 4.10.8 se ramène au cas où la catégorie d'indices est $I=(0\rightrightarrows 1)$. Mais je suspecte que la démonstration du théorème dans ce cas particulier, n'est pas essentiellement plus simple que dans le cas général.

\msp{page 409}

\textbf{4.14.23.} On peut se demander de trouver un énoncé assurant l'accessibilité d'une catégorie $\uGamma(\mathcal{E}/\mathcal{C})$, qui coiffe les deux cas envisagés dans 4.14.17. Voici ce que je peux proposer. Je vais faire l'hypothèse suivante sur les foncteur $p:\mathcal{E}\longrightarrow\mathcal{C}$ (en plus de celle que $\mathcal{C}$ soit essentiellement petite) : % [sic] les foncteur -> le foncteur, singular functor p

\begin{itemize}[leftmargin=*]
\item[H)] Pour tout $\lambda\in\Fl\mathcal{C}$, la catégorie $\mathcal{E}_\lambda=\mathcal{E}\times_{\mathcal{C}}(\bDelta_1,i_\lambda)$, où $i_\lambda:\bDelta_1\longrightarrow\mathcal{C}$ est le foncteur défini par $\lambda$, est accessible, et le foncteur canonique $\mathcal{E}_\lambda\xrightarrow{\;p_\lambda\;}\bDelta_1$ est accessible.
\end{itemize}

Notons que si $\lambda:s\longrightarrow t$ dans $\mathcal{C}$ est tel que $\mathcal{E}_\lambda$ et $p_\lambda$ soient accessibles, alors $\mathcal{E}_s$ et $\mathcal{E}_t$ sont accessible, car isomorphes aux catégories image inverse par $p_\lambda:\mathcal{E}_\lambda\longrightarrow\bDelta_1$ des sous-catégories $\{0\}$ resp. $\{1\}$ de $\Delta_1$, lesquelles sont accessibles dans la catégorie accessible $\bDelta_1$, donc % [sic] sont accessible -> sont accessibles, plural agreement
\msp{page 410}
$\mathcal{E}_{\lambda_0}$ et $\mathcal{E}_{\lambda_1}$ sont accessibles dans $\mathcal{E}_\lambda$, donc sont accessibles, donc aussi $\mathcal{E}_s$ et $\mathcal{E}_t$. Ainsi, l'hypothèse H) implique que les catégories fibres sont accessibles. Notons de plus :

\textbf{Lemme 4.14.23.1.} \textit{Supposons que $p:\mathcal{E}\longrightarrow\mathcal{C}$ soit cofibrant ou fibrant, et soit $(\lambda:s\longrightarrow t)\in\Fl(\mathcal{C})$. Pour que $\mathcal{E}_\lambda$ soit accessible et $\mathcal{E}_\lambda\xrightarrow{\;p_\lambda\;}\bDelta_1$ accessible, il faut et il suffit que $\mathcal{E}_s$, $\mathcal{E}_t$ soient accessibles et $\lambda_*:\mathcal{E}_s\longrightarrow\mathcal{E}_t$ resp. $\lambda^*:\mathcal{E}_t\longrightarrow\mathcal{E}_s$ soit un foncteur accessible.}

\textsc{Démonstration.} On peut supposer que $\mathcal{C}=\bDelta_1$, et que $\mathcal{E}$ est la catégorie cofibrée (resp. fibrée) sur $\bDelta_1$ définie par un foncteur $\varphi:\mathcal{E}_0\longrightarrow\mathcal{E}_1$ (resp. $\psi:\mathcal{E}_1\longrightarrow\mathcal{E}_0$), et à montrer que $\mathcal{E}$ et $p:\mathcal{E}\longrightarrow\bDelta_1$ sont accessibles si et seulement si $\mathcal{E}_0$ et $\mathcal{E}_1$ et $\varphi$ (resp. $\psi$) le sont.

\msp{page 411}

Prouvons le “il faut”. On a déjà vu que si $\mathcal{E}$, $p$ sont accessibles, $\mathcal{E}_0$ et $\mathcal{E}_1$ sont accessibles dans $\mathcal{E}_s$, donc accessibles, et les foncteurs d'inclusion % [sic] dans E_s -> dans E, here C = Delta_1 and the subcategories lie in E
\[
\mathcal{E}_0\ \hookrightarrow\ \mathcal{E}\ \hookleftarrow\ \mathcal{E}_1
\]
accessibles. Dire que $p$ est cofibrant (resp. fibrant) signifie que le foncteur d'inclusion $\mathcal{E}_1\hookrightarrow\mathcal{E}$ admet un adjoint à droite $\bar{\varphi}$ (resp. que le foncteur $\mathcal{E}_0\hookrightarrow\mathcal{E}$ admet un adjoint à gauche $\tilde{\psi}$). Alors $\varphi$ (resp. $\psi$) est le composé $\mathcal{E}_0\longrightarrow\mathcal{E}\xrightarrow{\;\bar{\varphi}\;}\mathcal{E}_1$ (resp. le composé $\mathcal{E}_1\hookrightarrow\mathcal{E}\xrightarrow{\;\tilde{\psi}\;}\mathcal{E}_1$). Donc pour prouver que $\varphi$ (resp. $\psi$) est accessible, il suffit de prouver que $\bar{\varphi}$ (resp. $\tilde{\psi}$) l'est. Cela est un cas particulier du % [sic] E -> E_1 (target of psi-tilde composite) -> E_0, psi goes from E_1 to E_0

\textbf{Lemme 4.14.23.2.} \textit{Soient $\mathcal{E}$, $\mathcal{E}'$ deux catégories accessibles} (\footnote[172]{\textbf{N.B.} Il suffit que $\mathcal{E}$ admette une petite sous-catégorie génératrice, et que les objets de $\mathcal{E}'$ soient accessibles.}), \textit{et $\mathcal{E}\underset{G}{\overset{F}{\rightleftarrows}}\mathcal{E}'$ un couple de foncteurs adjoints entre $\mathcal{E}$ et $\mathcal{E}'$. Alors $F$ et $G$ sont accessibles.}

\msp{page 412}

En effet, $F$ est accessible, car il commute aux petites $\varinjlim$ quelconques, a fortiori il satisfait $\mathrm{L}_{\pi_0}$ (où $\pi_0$ est tel que $\mathcal{E}$, $\mathcal{E}'$ satisfassent $\mathrm{L}_{\pi_0}$). Prouvons que $G$ est accessible. Si $S$ est un
% foot: Version 17 décembre 2025
\printedpage{156}
% head: A. Grothendieck : Dérivateurs, Ch. XVIII, Édition des Universités de Montpellier II et Paris VII
petit ensemble d'éléments générateur dans $\mathcal{E}$, de sorte que la petite famille des foncteurs
\[
h_\xi\ =\ (x\longmapsto\Hom_{\mathcal{E}}(\xi,x))\ :\ \mathcal{E}\ \longrightarrow\ \Ens
\]
est conservative, pour prouver que $G$ est accessible, il suffit de prouver que les $h_\xi\circ G$ ($\xi\in S\subset\Ob\mathcal{E}$) le sont. Or
\[
L_\xi\circ G(x')\ =\ \Hom(\xi,G(x'))\ \simeq\ \Hom_{\mathcal{E}'}(F(\xi),x')\ , % [sic] L_xi -> h_xi, the functors were named h_xi
\]
et comme $F(\xi)$ est un objet accessible de $\mathcal{E}'$, le foncteur
\[
x'\ \longmapsto\ (L_\xi\circ G)(x')\ \simeq\ \Hom_{\mathcal{E}'}(F(\xi),x') % [sic] L_xi -> h_xi, the functors were named h_xi
\]
est bien accessible, q.e.d.

Cela prouve le “il faut” dans 4.14.23.1. Prouvons le “il suffit”, i.e. supposons $\mathcal{E}_s$, $\mathcal{E}_t$ accessibles, et $\varphi$ (resp. $\psi$) accessible.

\msp{page 413}

Dans les deux cas, $\mathcal{E}_0$ est un ouvert de $\mathcal{E}$, $\mathcal{E}_1$ est le fermé complémentaire, et $\mathcal{E}$ est la catégorie “recollée” à l'aide du bifoncteur
\[
H\ :\ \mathcal{E}_0^{\circ}\times\mathcal{E}_1\ \longrightarrow\ \Ens
\]
donné par
\[
\left\{
\begin{aligned}
H(x_0,x_1)\ &=\ \Hom(\varphi(x_0),x_1)\\
&\hspace{-2em}\text{resp.}\\
H(x_0,x_1)\ &=\ \Hom(x_0,\psi(x_1))\ ,
\end{aligned}
\right.
\]
suivant qu'on est dans le cas $p$ cofibrant ou $p$ fibrant. On va utiliser la

\textbf{Proposition 4.14.24} (\footnote[173]{cf. énoncés plus généraux et démonstration sympathique, 4.15.1 et 4.15.4 (p. 425, 432).}). \textit{Soit $\mathcal{E}$ une catégorie, $\mathcal{E}_0$ une sous-catégorie ouverte, $\mathcal{E}_1$ la sous-catégorie fermée complémentaire. Conditions équivalentes :}
\begin{itemize}[leftmargin=*]
\item[\textit{a)}] \textit{$\mathcal{E}$ est accessible, et $\mathcal{E}\longrightarrow\bDelta_1$ (de fibres $\mathcal{E}_0$, $\mathcal{E}_1$) est accessible.}
\item[\textit{b)}] \textit{$\mathcal{E}$ est accessible, et $\mathcal{E}_0$, $\mathcal{E}_1$ sont accessibles dans $\mathcal{E}$.}
\item[\textit{c)}] \textit{$\mathcal{E}_0$, $\mathcal{E}_1$ sont accessibles, et pour tout $x_0\in\Ob\mathcal{E}_0$, le foncteur}
\[
x_1\ \longmapsto\ \uHom_{\mathcal{E}}(x_0,x_1)\ :\ \mathcal{E}_1\ \longrightarrow\ \Ens
\]
\msp{page 414}
\textit{est accessible, et il existe $\pi$ tel que pour tout $x_1\in\Ob\mathcal{E}_1$, le foncteur}
\[
x_0\ \longmapsto\ \uHom_{\mathcal{E}}(x_0,x_1)^{\circ}\ :\ \mathcal{E}_0\ \longrightarrow\ (\Ens)^{\circ}
\]
\textit{soit $\pi$-accessible.}
\end{itemize}
% foot: Version 17 décembre 2025

\printedpage{157}
% head: A. Grothendieck : Dérivateurs, Ch. XVIII, Édition des Universités de Montpellier II et Paris VII
Admettons ce lemme \texttt{[plutôt cette proposition]}, achevons la démonstration de 4.14.23.1. On suppose $\mathcal{E}_0$, $\mathcal{E}_1$ et $\varphi$ (resp. $\psi$) accessibles, on veut prouver la condition a) de 4.14.23.3. Pour ceci, il suffit de prouver la condition c). Or on a, dans le cas $p$ cofibrant, % [sic] 4.14.23.3 -> 4.14.24, the conditions a), c) are those of Proposition 4.14.24
\[
\Hom_{\mathcal{E}}(x_0,x_1)\ \simeq\ \Hom_{\mathcal{E}_1}(\varphi(x_0),x_1)\ ,
\]
c'est bien accessible en $x_1$, puisque $\mathcal{E}_1$ est accessible, et $\pi$-accessible en $x_0$ si $\varphi$ l'est, comme composé de $\varphi$ avec $\mathcal{E}_1\longrightarrow\Ens$, $\Hom(-,x_1)$, qui l'est. Dans le cas $p$ fibrant, on a
\[
\Hom_{\mathcal{E}}(x_0,x_1)\ \simeq\ \Hom_{\mathcal{E}_0}(x_0,\psi(x_1))\ ,
\]
c'est accessible en $x_1$, comme le composé du foncteur accessible $\psi:\mathcal{E}_1\longrightarrow\mathcal{E}_0$, et du foncteur également accessible $\Hom(x_0,-):\mathcal{E}_0\longrightarrow\Ens$, qui est $\pi$-accessible en $x_0$ tautologiquement.

Pour achever la démonstration de 4.14.23.1, il reste donc à prouver la proposition 4.14.24. On a déjà noté qu'on a a $\Longrightarrow$ b, et b $\Longrightarrow$ c clair, puisque le foncteur envisagé $\mathcal{E}\longrightarrow\Ens$ est le composé $\mathcal{E}_1\overset{\text{incl.}}{\hookrightarrow}\mathcal{E}\xrightarrow{\;\Hom_{\mathcal{E}}(x_0,-)\;}\Ens$, est accessible comme composé de deux
\msp{page 415}
foncteurs accessibles. Reste donc à prouver c $\Longrightarrow$ a. Soit donc $\pi_0$ un cardinal utile pour $\mathcal{E}_0$, et tel que les foncteurs
\[
x_0\ \longmapsto\ H(x_0,x_1)\ :\ \mathcal{E}_0\longrightarrow\Ens^{\circ}
\]
soient $\pi_0$-accessibles, soit $\pi$ un cardinal $\geqslant\sup(\pi_0,\mathrm{card}\,\text{réd}\,\mathcal{E}_{0\,\pi_0})$ tel que $\pi^{\pi_0}=\pi$ (donc utile pour $\mathcal{E}_0$), utile pour $\mathcal{E}_1$, et tel que les foncteurs
\[
x_1\ \longmapsto\ H(x_0,x_1)\ :\ \mathcal{E}_1\longrightarrow\Ens
\]
soient $\pi$-accessibles, pour $x_0\in\Ob\mathcal{E}_{0\,\pi_0}$. Soit $\mathcal{E}(\pi)$ la sous-catégorie strictement pleine de $\mathcal{E}$, telle que
\[
\Ob\mathcal{E}(\pi)\ =\ \Ob\mathcal{E}_{0\,\pi}\cup\Ob\mathcal{E}_{1\,\pi}\ .
\]
Je dis que
\begin{itemize}[leftmargin=*]
\item[1°)] $\mathcal{E}$ satisfait à $\mathrm{L}_\pi$, et $\mathcal{E}_0$ et $\mathcal{E}_1$ sont stables dans $\mathcal{E}$ sous $\mathrm{L}_\pi$, les foncteurs $\mathcal{E}_0\longrightarrow\mathcal{E}$, $\mathcal{E}_1\longrightarrow\mathcal{E}$, $\mathcal{E}\longrightarrow\bDelta_1$ satisfont $\mathrm{L}_\pi$ (\footnote[174]{il suffit que $\mathcal{E}_0$ et $\mathcal{E}_1$ satisfassent $\mathrm{L}_\pi$, et que
\[
x_0\ \longmapsto\ H(x_0,x_1)\ :\ \mathcal{E}_0\longrightarrow\Ens^{\circ}
\]
satisfasse $\mathrm{L}_\pi$ $\forall\ x\in\Ob\mathcal{E}$.}).
\item[2°)] $\mathcal{E}(\pi)=\mathcal{E}_\pi$ (d'où un foncteur pleinement fidèle
\[
\mathrm{Ind}_\pi(\mathcal{E}(\pi))\ \longrightarrow\ \mathcal{E}\ \ ).
\]
\item[3°)] Ledit foncteur est essentiellement surjectif,
\msp{page 416}
donc une équivalence de catégories.
\end{itemize}
Notons que 2°) et 3°) impliquent que $\mathcal{E}$ est accessible, et 1°) que $\mathcal{E}\longrightarrow\bDelta_1$ est accessible, d'où la condition a) du lemme qu'il s'agit d'établir. Donc tout revient à prouver les 1°, 2°, 3° ci-dessus.
% foot: Version 17 décembre 2025

\printedpage{158}
% head: A. Grothendieck : Dérivateurs, Ch. XVIII, Édition des Universités de Montpellier II et Paris VII
1°) Comme $\mathcal{E}_1$ est fermée dans $\mathcal{E}$, il est clair que $\mathcal{E}$ est stable par toutes $\varinjlim$ représentables dans $\mathcal{E}$, et plus précisément, que pour un foncteur $I\longrightarrow\mathcal{E}_1$, sa limite inductive dans $\mathcal{E}_1$ existe si et seulement si elle existe dans $\mathcal{E}$, et les deux sont isomorphes. Supposons $I$ grande devant $\pi$, et soit $\varphi:I\longrightarrow\mathcal{E}$, $i\longmapsto x_i$. Distinguons deux cas.

a) $\varphi$ prend ses valeurs dans $\mathcal{E}_0$. Soit $x_0$ sa limite inductive \textit{dans} $\mathcal{E}_0$. Je dis que c'est aussi une limite inductive dans $\mathcal{E}$. Il faut prouver que pour $x_1\in\Ob\mathcal{E}_1$, la flèche
\[
\Hom_{\mathcal{E}}(x_0,x_1)\ \longrightarrow\ \varprojlim_{i\in I}\Hom_{\mathcal{E}}(x_i,x_1)
\]
est isomorphe. Or c'est le cas
\msp{page 417}
par hypothèse sur $\pi$,
\[
\xi\ \longmapsto H(\xi,x_1)\ :\ \mathcal{E}_0\ \longrightarrow\ (\Ens)^{\circ}
\]
satisfait à $\mathrm{L}_\pi$ pour tout $x_1$ dans $\Ob\mathcal{E}_1$.

b) $\varphi$ ne se factorise pas par $\mathcal{E}_0$, i.e. $\exists\ i_0\in\Ob I$ avec $\varphi(i_0)\in\Ob\mathcal{E}_1$. Alors on a $\varphi(i)\in\Ob\mathcal{E}_1$ pour tout $i\in\Ob I[i_0]$, où $I[i_0]$ désigne la sous-catégorie strictement pleine de $I$ formée des $i$ tels que $i\geqslant i_0$. Cette catégorie est cofinale dans $I$, donc grande devant $\pi$ comme $I$, et pour que $\varinjlim\limits_{I}x_i$ existe, il faut et il suffit que $\varinjlim\limits_{I[i_0]}x_i$ existe, ce qui est le cas d'après les cas a) déjà traité. % [sic] les cas a) déjà traité -> le cas a) déjà traité, singular agreement

Cela prouve donc en même temps que $\mathcal{E}$ satisfait à $\mathrm{L}_\pi$, et que $\mathcal{E}_0$ et $\mathcal{E}_1$ sont $\mathrm{L}_\pi$-stables. La démonstration montre en même temps que le foncteur $p:\mathcal{E}\longrightarrow\bDelta_1$ satisfait à $\mathrm{L}_\pi$.

2°) De ceci résulte aussitôt que $\mathcal{E}_\pi\subset\mathcal{E}(\pi)$. Prouvons qu'on a aussi
\[
\mathcal{E}(\pi)\ \subset\ \mathcal{E}_\pi\ ,
\]
et distinguons encore deux cas, pour
\msp{page 418}
$x\in\Ob E(\pi)$, suivant que $x$ est dans $\mathcal{E}_{0\,\pi}$ ou $\mathcal{E}_{1\,\pi}$. S'il est dans $\mathcal{E}_{1\,\pi}$, pour prouver que % [sic] E(pi) -> \mathcal{E}(pi), script letter elsewhere
\[
\Hom_{\mathcal{E}}(x,\underbrace{\varinjlim y_i}_{y})\ \xleftarrow{\ \sim\ }\ \varinjlim_{i}\Hom_{\mathcal{E}}(x,y_i)
\tag*{$(*)$}
\]
(pour un système inductif $(y_i)_{i\in I}$ dans $\mathcal{E}$ grand devant $\pi$), si les $y_i$ sont dans $\mathcal{E}_0$, donc $y\in\mathcal{E}_0$, alors les deux membres sont vides, et c'est OK. Sinon, on peut supposer (quitte à remplacer $I$ par un $I[i_0]$) que les $x_i$ sont dans $\mathcal{E}_1$, et alors c'est OK aussi, puisque $y$ est aussi la $\varinjlim$ dans $\mathcal{E}_1$, d'après 1°). % [sic] les x_i -> les y_i, the inductive system is (y_i)

Si $x$ est dans $\mathcal{E}_{0\,\pi}$, alors $(*)$ est clair si les $y_i$ sont dans $\mathcal{E}_0$. Sinon, on peut supposer encore que les $y_i$ sont dans $\mathcal{E}_0$, et $(*)$ s'écrit % [sic] y_i sont dans E_0 (second occurrence) -> E_1, the case "sinon" and the following "les y_i dans E_1"
\[
H(x,\varinjlim y_i)\ \xleftarrow{\ \sim\ }\ \varinjlim H(x,y_i)\ ,
\]
pour $x\in\Ob\mathcal{E}_{0\,\pi}$, les $y_i$ dans $\mathcal{E}_1$, $I$ grande devant $\pi$. En d'autres termes, il faut prouver que le foncteur
\[
x_1\ \longmapsto\ H(x_0,x_1)\ :\ \mathcal{E}_1\ \longrightarrow\ \Ens
\]
% foot: Version 17 décembre 2025


\endgroup
% END UNIT der-ch18/der-ch18_p156-160.tex
% BEGIN UNIT der-ch18/der-ch18_p161-165.tex
\begingroup\makeatletter%
\derleqno%
\newcommand{\printedpage}[1]{\marginpar{\footnotesize p.~#1}}%
\newcommand{\msp}[1]{\marginpar{\footnotesize [#1]}}%
\newcommand{\cE}{\mathcal{E}}%
\newcommand{\cC}{\mathcal{C}}%
\newcommand{\cM}{\mathcal{M}}%
\newcommand{\Ob}{\operatorname{Ob}}%
\newcommand{\Fl}{\operatorname{Fl}}%
\newcommand{\Hom}{\operatorname{Hom}}%
\newcommand{\End}{\operatorname{End}}%
\newcommand{\Ind}{\operatorname{Ind}}%
\newcommand{\Ens}{\mathrm{Ens}}%
\newcommand{\pr}{\mathrm{pr}}%
\newcommand{\id}{\mathrm{id}}%
\newcommand{\card}{\operatorname{card}}%
\newcommand{\cardred}{\operatorname{card}\,\text{\upshape réd}}%
\newcommand{\defeq}{\overset{\text{\upshape déf}}{=}}%
\makeatother\let\printedpage\derpp\let\msp\dermsp%
\let\mspin\DERmspin%

\setcounter{footnote}{174}

%=====================================================================
\printedpage{159}
% head: A. Grothendieck : Dérivateurs, Ch. XVIII, Édition des Universités de Montpellier II et Paris VII

\msp{page 419}
est $\pi$-accessible, quand $x_0$ est dans $\cE_{0\,\pi}$. Or par hypothèse sur $\pi$, il est $\pi$-accessible si $x_0$ est dans $\cE_{0\,\pi_0}$. Mais si $x$ est dans $\cE_{0\,\pi}$, on aura % [sic] x -> x_0, the formula that follows concerns x_0
\[
x_0 \;=\; \varinjlim_{i\in I} x_{0\,i} \qquad
\left\{
\begin{array}{l}
I\ \ \text{grande}\ \ \text{devant}\ \ \pi_0,\\
\text{les}\ \ x_{0\,i}\ \ \text{dans}\ \ \cE_{0\,\pi_0},\\
\card I \leqslant \pi
\end{array}
\right.
\]
(car $\pi_0$ est utile pour $\cE_0$ par hypothèse, et $\pi$ satisfait $\pi^{\pi_0} = \pi \geqslant \cardred \cE_{\pi_0}$). Comme $x_0 \longrightarrow H(x_0,x_1) : \cE_0 \longrightarrow \Ens$ satisfait à $\mathrm{L}_{\pi_0}$, pour tout $x_1$ dans $\cE_1$, on aura donc % [sic] x_0 \longrightarrow H(x_0,x_1) -> x_0 \longmapsto H(x_0,x_1), elsewhere such functors are written with \longmapsto
\[
H(x_0) \;\simeq\; \varprojlim_{i\in I} H(x_{0\,i})
\]
(où pour $x_0$ dans $\cE_0$, $H(x)$ désigne le foncteur $x_1 \longmapsto H(x_0,x_1) : \cE_1 \longrightarrow (\Ens)$). Comme % [sic] H(x) -> H(x_0), the functor depends on x_0
les foncteurs $H(x_{0\,i})$ sont $\pi$-accessibles, et $\card I \leqslant \pi$, il s'ensuit que leur $\varprojlim\limits_{I}$ $H(x_0)$ est $\pi$-accessible (SGA 4 I 9.6 (ii)). Cela achève de prouver donc 2°).

3°). Pour prouver que $\Ind_\pi(\cE(\pi)) \longrightarrow \cE$ est essentiellement surjectif, il suffit de voir que tout $x \in \Ob\cE_0$ et tout $x \in \Ob\cE_1$
\msp{page 420}
est dans l'image essentielle. Or cela résulte de l'hypothèse que $\pi$ est utile pour $\cE_0$ et $\cE_1$, et de 1°) qui implique que $\cE_0 \longrightarrow \cE$ et $\cE_1 \longrightarrow \cE$ satisfont $\mathrm{L}_\pi$.

Cela achève de prouver la proposition 4.14.24, et de même le résultat suivant, plus précis que l'implication c) $\Longrightarrow$ a).

\textbf{Corollaire 4.14.25.} \textit{Soient $\cE$ une catégorie, $\cE_0$ une sous-catégorie ouverte, $\cE_1$ la sous-catégorie fermée complémentaire,}
\[
H \;=\; ((x_0,x_1) \longmapsto \Hom_{\cE}(x_0,x_1)) \;:\; \cE_0^\circ \times \cE_1 \;\longrightarrow\; \Ens
\]
\textit{le foncteur recollement.}
\begin{enumerate}
\item[\textit{1°)}] \textit{Soit $\pi$ un cardinal infini. Pour ce que $\cE$ et $\cE_0$ satisfassent à $L_\pi$, et le foncteur d'inclusion $\cE_0 \hookrightarrow \cE$ itou, il faut et il suffit que $\cE_0$ satisfasse à $L_\pi$, et que pour tout $x_1$ dans $\cE_1$, le foncteur $x_0 \longmapsto H(x_0,x_1) : \cE_0 \longrightarrow \Ens^\circ$ satisfasse à $L_\pi$. Alors $\cE_1$ et $\cE_1 \hookrightarrow \cE$ satisfont aussi à $L_\pi$,}
\msp{page 421}
\textit{et on a}
\[
\cE_\pi \;\subset\; \cE(\pi) \;\defeq\;
\begin{array}{l}
\text{\itshape sous-catégorie\ \ strictement\ \ pleine\ \ de\ \ $\cE$\ \ telle\ \ que}\\
\Ob\cE(\pi) = \Ob\cE_{0\,\pi} \cup \Ob\cE_{1\,\pi}.
\end{array}
\]
\item[\textit{2°)}] \textit{Supposons que la condition c) de 4.14.24 soit satisfaite, et soit $\pi_0$ un cardinal infini tel que $\pi_0$ soit utile pour $\cE_0$, et que les foncteurs $x_0 \longmapsto H(x_0,x_1) : \cE_0 \longrightarrow \Ens^\circ$ définis par les $x_1 \in \Ob\cE_1$ soient $\pi_0$-accessibles. Soit $\pi$ un cardinal tel que $\pi^{\pi_0} = \pi \geqslant \sup(\pi_0, \cardred\cE_{0\,\pi_0})$, et tel que pour tout $x_0 \in \Ob\cE_{\pi_0}$, le foncteur $x_1 \longmapsto H(x_0,x_1) : \cE_1 \longrightarrow \Ens$ soit $\pi$-accessible. Alors on a}
\[
\cE_\pi \;=\; \cE(\pi)\ .
\]
% foot: Version 17 décembre 2025

%=====================================================================
\printedpage{160}
% head: A. Grothendieck : Dérivateurs, Ch. XVIII, Édition des Universités de Montpellier II et Paris VII

\item[\textit{3°)}] \textit{Si $\pi_0$ et $\pi$ sont comme dans 2°), et si de plus $\pi$ est utile pour $\cE_1$, alors $\pi$ est utile pour $\cE$, donc (en vertu de 2°)}
\[
\cE \;\xleftarrow{\ \sim\ }\; \Ind_\pi(\cE(\pi))\ .
\]
\end{enumerate}
Nous pouvons maintenant énoncer la généralisation cherchée de 4.14.17 :

\msp{page 422}
\textbf{Théorème 4.14.26} (généralisation de 4.14.13). \textit{Soit $p : \cE \longrightarrow \cC$ un foncteur, avec $\cC$ petite. On suppose que $p$ satisfait à l'hypothèse H) de la page 409, qui équivaut (moyennant 4.14.24) à la suivante :}
\begin{enumerate}
\item[\textit{H$'$)}] \textit{$\forall\ s \in \Ob\cC$, la catégorie fibre $\cE_s$ est accessible, et pour tout $\lambda : s \longrightarrow t$ dans $\Fl\cC$, le foncteur}
\[
\left\{
\begin{array}{lcl}
H_\lambda : \cE_s^\circ \times \cE_t & \longrightarrow & \Ens\\[1ex]
\qquad (x_0,x_1) & \longmapsto & \Hom_\lambda(x_0,x_1)
\end{array}
\right.
\]
\textit{satisfait les deux conditions suivantes :}
\begin{enumerate}
\item[\textit{a)}] \textit{$\forall\ x_0 \in \Ob\cE_s^\circ$, $x_1 \longmapsto H_\lambda(x_0,x_1) : \cE_t \longrightarrow \Ens$ est accessible.}
\item[\textit{b)}] \textit{Il existe un cardinal infini $\pi_\lambda$, tel que pour tout $x_1 \in \Ob\cE_t$, le foncteur $x_0 \longmapsto H_\lambda(x_0,x_1) : \cE_s \longrightarrow (\Ens)^\circ$ soit $\pi_\lambda$-accessible.}
\end{enumerate}
\end{enumerate}
\textit{Ceci posé, la catégorie $\underline{\Gamma}(\cE/\cC)$ des sections de $\cE$ sur $\cC$ est accessible, et les foncteurs $\varepsilon_s : \underline{\Gamma}(\cE/\cC) \longrightarrow \cE_s$ sont accessibles. On peut préciser de façon similaire à 4.14.18 (énoncé en forme laissé au lecteur).}

\msp{page 423}
\textsc{Démonstration} de 4.14.26. Elle se fait exactement comme celle de 4.14.17 (p. 400 ff.), mais est conceptuellement plus simple. On suppose $\cC$ petite, et on introduit les catégories
\[
\Pi \;=\; \prod_{s\,\in\,\Ob\cC} \cE_s\ , \qquad \Pi' \;=\; \prod_{\lambda\,\in\,\Fl\cC} \cE_\lambda\ ,
\]
où
\[
\cE_\lambda \;=\; \cE \times_{\cC} (\boldsymbol{\Delta}_1, i_\lambda)\ ,
\]
et les deux foncteurs canoniques
\[
\Pi \;\underset{G}{\overset{F}{\rightrightarrows}}\; \Pi'
\]
de composantes
\[
\begin{aligned}
F_\lambda \;&=\; i_\lambda \circ \pr_{s(\lambda)}\\
G_\lambda \;&=\; j_\lambda \circ \pr_{b(\lambda)}\ ,
\end{aligned}
\]
où pour $\lambda : s \longrightarrow t$, on pose
\[
s(\lambda) \;=\; s\ , \qquad b(\lambda) \;=\; t\ ,
\]
et $i_\lambda$, $j_\lambda$ sont les inclusions canoniques
% In the print, the vertical arrow p_lambda and Delta_1 are enclosed together in a pair of large parentheses.
\[
\begin{tikzcd}
\cE_s \arrow[r,"i_\lambda"] \arrow[d,dash,"\wr"] & \cE_\lambda \arrow[dd,"p_\lambda"] & \cE_t\ . \arrow[l,"j_\lambda"'] \arrow[d,dash,"\wr"] \\
(\cE_\lambda)_0 & & (\cE_\lambda)_1 \\
& \boldsymbol{\Delta}_1 &
\end{tikzcd}
\]
% foot: Version 17 décembre 2025

%=====================================================================
\printedpage{161}
% head: A. Grothendieck : Dérivateurs, Ch. XVIII, Édition des Universités de Montpellier II et Paris VII

\msp{page 424}
On désigne par $\cM'$ la catégorie
\[
\cM' \;=\;
\begin{array}{l}
\text{catégorie des couples } (X,u)\text{, avec}\\
X \in \Ob\Pi,\ u : F(x) \longrightarrow G(x),
\end{array}
\] % [sic] F(x), G(x) -> F(X), G(X), the object is named X
qui est constructible par 4.14.22.1 est accessible (car $\Pi$ et $\pi'$ et les foncteurs $F$ et $G$ le sont, par l'hypothèse H) faite sur $p$). On peut l'interpréter comme la catégorie des sections % [sic] \pi' -> \Pi', the category introduced is \Pi'
de $\mathrm{Diag}(\cE)$ sur $\mathrm{Diag}(\cC)$. On a un foncteur strictement pleinement fidèle
\[
\cM \;\defeq\; \underline{\Gamma}(\cE/\cC) \;\hookrightarrow\; \cM'\ ,
\]
qui permet d'interpréter $\cM$ comme une sous-catégorie strictement pleine de $\cM'$. On voit alors comme dans le cas de 4.14.17 que cette sous-catégorie est accessible dans $\cM'$ -- c'est essentiellement sorital. Il s'ensuit que $\cM$ est accessible et que les foncteurs $\cM \longrightarrow \cE_s$ le sont aussi, comme composés de
\[
\cM \;\overset{\text{\upshape acc.}}{\hookrightarrow}\; \cM' \;\overset{\text{\upshape acc.}}{\hookrightarrow}\; \Pi \;\xrightarrow{\ \ \pr_s\ \ \text{\upshape acc.}\ \ }\; \cE_s\ ,
\]

\hfill q.e.d.

\msp{page 425}

\subsection*{4.15\quad Catégorie accessible sur une autre}

\textbf{Proposition 4.15.1.} \textit{Soit $f : \cE \longrightarrow \cC$ un foncteur avec $\cC$ petite. On suppose de plus satisfaite la condition H (p. 409), qu'on explicite ici sous la forme équivalente H$'$ :}
\begin{enumerate}
\item[\textit{H 1)}] \textit{Pour tout $s \in \Ob\cC$, $\cE_s$ est accessible.}
\item[\textit{H 2)}] \textit{Pour toute $\lambda \in \Fl\cC$, $\lambda : s \longrightarrow t$, existe un cardinal infini $\pi_\lambda$ tel que pour tout $x_1 \in \Ob\cE_t$, le foncteur}
\[
x_0 \;\longmapsto\; \Hom_\lambda(x_0,x_1) \;:\; \cE_s \;\longrightarrow\; (\Ens)^\circ
\]
\textit{est $\pi_\lambda$-accessible ($\cE_s$ satisfaisant de plus à $L_{\pi_\lambda}$).}
\item[\textit{H 3)}] \textit{Pour toute $\lambda \in \Fl\cC$, $\lambda : s \longrightarrow t$, et pour tout $x_0 \in \Ob\cE_s$, le foncteur}
\[
x_1 \;\longmapsto\; \Hom_\lambda(x_0,x_1) \;:\; \cE_t \;\longrightarrow\; (\Ens)
\]
\textit{est accessible.}
\end{enumerate}
\textit{Alors les foncteurs d'inclusion $i_s : \cE_s \longrightarrow \cE$ sont accessibles. De plus, si $\cC$ ``est sans projecteurs'' (i.e. si tout endomorphisme $p$ d'un objet $s$ de $\cC$ tel que $p^2 = p$, est l'identité), alors $\cE$ est accessible, et $f : \cE \longrightarrow \cC$ est accessible.}

Cela résulte du résultat plus précis :
% foot: Version 17 décembre 2025

%=====================================================================
\printedpage{162}
% head: A. Grothendieck : Dérivateurs, Ch. XVIII, Édition des Universités de Montpellier II et Paris VII

\msp{page 426}
\textbf{Corollaire 4.15.2} (\footnote{Pour a) et b) on n'utilise que H 1 et H 2, à l'exclusion de H 3.}). \textit{Les hypothèses sur $f : \cE \longrightarrow \cC$ sont celles de 4.15.1 ($\cC$ petite, et $f$ satisfait $H' =$ H 1) 2) 3)).}
\begin{enumerate}
\item[\textit{a)}] \textit{Soit $s \in \Ob\cC$ et soit $\pi_1$ un cardinal infini tel que pour toute $\lambda : s \longrightarrow t$ (flèche de $\cC$ de source $t$), on ait $\pi_1 \geqslant \pi_\lambda$ (où $\pi_\lambda$ est défini dans l'énoncé de H 2). Alors le foncteur $i_s : \cE_s \longrightarrow \cE$ est $\pi_1$-accessible.} % [sic] source t -> source s, lambda : s -> t has source s (cf. p. 163 "de source s")
\item[\textit{b)}] \textit{On suppose $\cC$ sans projecteurs. Soit $\pi_1$ un cardinal $\geqslant \sup(\pi_\lambda, \card\cC)$. Alors $\cE$ satisfait $L_{\pi_1}$, et $f : \cE \longrightarrow \cC$ est $\pi$-accessible.}
\item[\textit{c)}] \textit{Supposons encore que $\cC$ soit sans projecteurs (p.ex. une catégorie ``rigide'', p.ex. une catégorie préordonnée). Soient $\pi_0$, $\pi_1$ des cardinaux tels que :}
\begin{enumerate}
\item[\textcircled{\footnotesize $1^o$}] $\forall\ s \in \Ob\cC$, $\pi_0$ est utile pour $\cE_s$ \textit{(un tel $\pi_0$ existe grâce à H 1 et au fait que $\cC$, donc $\Ob\cC$ petits).}
\item[\textcircled{\footnotesize $2^o$}] $\pi_1 \geqslant \sup(\pi_0, \card\cC, \sup_{\lambda\in\Fl\cC}\pi_\lambda, \sup_{s\in\Ob\cC} \cardred \cE_{s\,\pi_0})$.
\item[\textcircled{\footnotesize $3^o$}] \textit{$\forall\lambda \in \Fl\cC$, et $x_0 \in \Ob\cE_{s,\pi_0}$, le foncteur $H_\lambda(x_0)$ de H 3, $x_1 \longmapsto \Hom_\lambda(x_0,x_1) : \cE_t \longrightarrow \Ens$ est $\pi_1$-accessible.}
\end{enumerate}
\textit{Soit $\pi$ un cardinal tel que $\pi^{\pi_0} = \pi \geqslant \pi_1$. Alors par b), $\cE$ satisfait $L_{\pi_1}$, donc $L_\pi$, et $f : \cE \longrightarrow \cC$ est $\pi_1$-accessible. De plus,}
\[
i_s(\cE_{s\,\pi}) \;\subset\; \cE_\pi\ ,
\]
\textit{ou encore}
\[
\cE(\pi) \;\subset\; \cE_\pi\ , \tag*{$(*)$}
\]
\textit{où $\cE(\pi)$ est la sous-catégorie pleine de $\cE$ telle que}
\[
\Ob\cE(\pi) \;=\; \bigcup_{s\in\Ob\cC} \Ob\,(\cE_{s\,\pi})\ .
\]
\textit{Comme $\cE_\pi$ satisfait $L_{\pi_1}$ par b), donc a fortiori $L_\pi$, on a donc par $(*)$ un foncteur} pleinement fidèle
\[
\Ind_\pi(\cE_{s\,\pi}) \;\xrightarrow{\ \sim\ }\; \cE\ ,
\]
\textit{et ce foncteur est une équivalence de catégories (donc $\cE$ est accessible).}
\item[\textit{d)}] \textit{En fait, sous les hypothèses de c) sur $\cC$, $\pi_0$, $\pi_1$, $\pi$, on a $\cE(\pi) = \cE_\pi$ (non seulement $\cE(\pi) \subset \cE_\pi$).}
\end{enumerate}

\msp{page 427}
\textsc{Démonstration de 4.15.1 et 4.15.2.} Il est clair que 4.15.1 est contenu dans 4.15.2 a) et c). Nous allons donc prouver 4.15.2.

a) Soit $(x_i)_{i\in I}$ un système inductif dans $\cE_s$, avec $I$ grand devant $\pi_0$, et $x$ sa limite inductive dans $\cE_s$ (qui existe, car $\cE_s$ satisfait à $\mathrm{L}_{\pi_\lambda}$, donc à $\mathrm{L}_{\pi_0}$). Il faut prouver que c'est une limite dans $\cE$, donc que pour tout $y \in \Ob\cE$, la flèche canonique
\[
\Hom_{\cE}(x,y) \;\longrightarrow\; \varprojlim_{I} \Hom_{\cE}(x_i,y) \tag*{$(*)$}
\]
% foot: Version 17 décembre 2025

%=====================================================================
\printedpage{163}
% head: A. Grothendieck : Dérivateurs, Ch. XVIII, Édition des Universités de Montpellier II et Paris VII
est un isomorphisme. Or $\Hom_{\cE}(x,y)$ et les $\Hom_{\cE}(x_i,y)$ sont des objets de $\Ens/\Hom_{\cC}(s,t)$, où $s = f(x)$, $t = f(y)$, et les flèches de transition sont des flèches dans $\Ens/\Hom_{\cC}(s,t)$. Cela montre que la bijectivité de $(*)$ revient à celle de
\[
\Hom_\lambda(x,y) \;\longrightarrow\; \varprojlim_{I} \Hom_\lambda(x_i,y) \tag*{$(*)$}
\]
pour tout $\lambda \in \Hom_{\cC}(x,y)$ \texttt{[plutôt} $\Hom_{\cC}(f(x),f(y))$\texttt{]}. Pour $\lambda$ fixé, et $y$ fixé dans $\Ob\cE_t$, la bijectivité dans $(*)$, pour tout système inductif grand devant $\pi_1$ dans $\cE_s$, de limite $x$, signifie que le foncteur
\[
x \;\longmapsto\; \Hom_\lambda(x,y) \;:\; \cE_s \;\longrightarrow\; (\Ens)^\circ
\]
est $\pi_0$-accessible. Or c'est bien le cas, par hypothèse sur $\pi_1$, quel que soit la flèche $\lambda : s \longrightarrow t$ de source $s$, et quel que soit $y \in \Ob\cE_t$.

\msp{page 428}
Soit $(x_i)_{i\in I}$ un système inductif grand devant $\pi_1$, et $(s_i)_{i\in I}$ l'image de ce système dans $\cC$. Nous prouverons plus bas le

% In the left margin, beside the statement of Lemme 4.15.3 and beside the following paragraph "Mais attention ...", the edition prints a wavy marginal sign (one for each).
\textbf{Lemme 4.15.3} (\footnote{Cf. démonstration page 452.},\footnote{sous réserve de la validité de 4.8.13 (Mais c'est à présent établi, sous une forme suffisante, dans 4.15.10).}). \textit{Soit $\cC$ une catégorie, $\pi$ un cardinal, tels que} $\card\cC \leqslant \pi$\textit{, et que $\cC$ soit sans projecteurs. Alors pour tout système inductif $(s_i)_{i\in I}$ grand devant $\pi$ dans $\cC$, existe une partie cofinale $J$ dans $I$ telle que le système induit $(s_i)_{i\in J}$ soit} constant, \textit{i.e. tel que les $s_i$ soient égaux à un $s \in \cC$, et que les morphismes de transition soient tous égaux à} $\id_s$.

Mais attention, on s'appuie pour prouver 4.15.3 sur le théorème 4.8.13 (page 212-216), que je n'ai toujours pas prouvé. Dans le cas où $\cC$ est rigide (i.e. $\End_{\cC}(s) = \{\id_s\}\ \forall\ s \in \Ob\cC$), la démonstration est cependant immédiate, cf. plus bas.

\msp{page 429}
Pour prouver que $x = \varinjlim\limits_{I} x_i$ existe dans $\cE$, et commute à $f$, on peut remplacer $I$ par une sous-catégorie cofinale. Cela nous ramène, grâce au lemme, au cas où $(s_i)_{i\in I}$ est constant, donc au cas où le système inductif provient d'un système inductif dans une catégorie fibre $\cE_s$. Mais par a), on sait que sa limite $x$ dans $\cE_s$ est une limite dans $\cE$, d'où l'existence de cette dernière, et la commutation à $f$, puisque $\varinjlim s_i = s$ ( $\varinjlim$ d'un système inductif filtrant constant).

c) Il faut prouver que si $x \in \Ob\cE_s$ est $\pi$-accessible dans $\cE_s$, il l'est dans $\cE$, i.e. que pour tout système inductif $(y_i)_{i\in I}$ grande devant $\pi$, la flèche canonique
\[
\varinjlim_{i} \Hom(x,y_i) \;\longrightarrow\; \Hom(x,y) \tag*{$(**)$}
\]
est bijective. Or comme par 2°) $\pi \geqslant \pi_1 \geqslant \sup(\card S, \sup_\lambda \pi_\lambda)$, il résulte encore du lemme % [sic] card S -> card \mathcal{C}, no S is defined; the condition 2° bounds card \mathcal{C}
que, quitte à remplacer $I$ par une sous-catégorie cofinale, on peut supposer que le système inductif $(y_i)_{i\in I}$ provient d'un système dans une catégorie fibre $\cE_t$, et que $y$ est sa $\varinjlim$ dans $\cE_t$. Mais alors les ensembles $\Hom(x,y)$ et $\Hom(x,y_i)$, et les morphismes de transition entre celles-ci, sont encore dans $\Ens/\Hom(s,t)$, et on est ramené
% foot: Version 17 décembre 2025


\endgroup
% END UNIT der-ch18/der-ch18_p161-165.tex
% BEGIN UNIT der-ch18/der-ch18_p166-170.tex
\begingroup\makeatletter%
\newcommand{\printedpage}[1]{\marginpar{\footnotesize p.~#1}}%
\newcommand{\msp}[1]{\marginpar{\footnotesize [#1]}}%
\newcommand{\cE}{\mathcal{E}}%
\newcommand{\cC}{\mathcal{C}}%
\newcommand{\Ob}{\operatorname{Ob}}%
\newcommand{\Fl}{\operatorname{Fl}}%
\newcommand{\Hom}{\operatorname{Hom}}%
\newcommand{\Ind}{\operatorname{Ind}}%
\newcommand{\card}{\operatorname{card}}%
\newcommand{\red}{\operatorname{\text{réd}}}%
\newcommand{\id}{\mathrm{id}}%
\newcommand{\Ens}{\mathrm{Ens}}%
\newcommand{\circled}[1]{\tikz[baseline=(c.base)]\node[draw,circle,inner sep=1pt](c){#1};}%
\makeatother\let\printedpage\derpp\let\msp\dermsp%
\let\mspin\DERmspin%


\printedpage{164}
% head: A. Grothendieck : Dérivateurs, Ch. XVIII, Édition des Universités de Montpellier II et Paris VII

\msp{page 430}à prouver que pour tout $\lambda : s \longrightarrow t$, la flèche
\[
\varinjlim_{i} \Hom_{\lambda}(x, y_i) \ \longrightarrow \ \Hom_{\lambda}(x, y)
\]
est bijective. Pour $x$ et $\lambda$ fixés, et des systèmes inductifs variables dans $\cE_t$, grands devant $\pi$, cela signifie que le foncteur
\[
H_{\lambda}(x) \ = \ (y \longmapsto \Hom_{\lambda}(x, y)) \ : \ \cE_t \ \longrightarrow \ \Ens
\]
est $\pi$-accessible, et ceci quel que soit
\[
x \ \in \ \Ob \cE_{s\,\pi} \ .
\]
Or c'est le cas si $s \in \Ob \cE_{s\,\pi_0}$, en vertu de l'hypothèse $3^{\circ}$) (alors $H_{\lambda}(x_0)$ est $\pi_1$-accessible, % [sic] $s \in \Ob \cE_{s\,\pi_0}$ -> $x_0 \in \Ob \cE_{s\,\pi_0}$, the condition concerns an object of $\cE_{s\,\pi_0}$
et a fortiori $\pi$-accessible). Il faut passer de là au cas où $x \in \Ob \cE_{s\,\pi}$. Or comme $\pi_0$ est utile pour $\cE_s$ par $1^{\circ}$, on a
\[
x \ = \ \varinjlim_{j \in J} x_j \ , \qquad x_j \in \Ob \cE_{s\,\pi_0},\ J \text{ grande devant } \pi_0 \ ,
\]
et comme par $2^{\circ}$, $\pi^{\pi_0} = \pi \geqslant \pi_1 \geqslant \sup(\pi_0, \card \red \cE_{s\,\pi_0})$, on en conclut qu'on peut prendre $J$ de sorte que
\[
\card J \ \leqslant \ \pi \ .
\]
D'autre part, en vertu de $\pi \geqslant \pi_1 \ \geqslant \ \pi_{\lambda}$, le foncteur
\[
\left\{
\begin{array}{lcl}
x_0 & \longmapsto & H_{\lambda}(x_0) \\[1ex]
\cE_s & \longrightarrow & (\underline{\Hom}(\cE_s, \Ens))^{\circ}
\end{array}
\right.
\]
est $\pi_0$-accessible, donc on a
\[
H_{\lambda}(x_0) \ \simeq \ \varprojlim_{j \in J} H_{\lambda}(x_j) \ .
\]

\msp{page 431}Comme $x_j \ \in \ \Ob \cE_{s\,\pi_0}$, $H(x_j)$ est un foncteur $\pi_1$-accessible, et comme $\card J \leqslant \pi$, la $\varprojlim\limits_{j \in J} H_{\lambda}(x_j)$ est aussi $\pi$-accessible (cf. SGA 4 I 9.8).

Ainsi on a prouvé $\cE_{s\,\pi} \subset \cE(\pi)$ pour tout $s \in \Ob \cC$, i.e. l'inclusion
\[
\cE(\pi) \ \subset \ \cE_{\pi} \ ,
\]
d'où le foncteur pleinement fidèle
\[
\Ind_{\pi}(\cE(\pi)) \longrightarrow \cE \ , \tag*{$(*)$}
\]
comme $\cE$ satisfait à $\mathrm{L}_{\pi}$. Il reste à prouver que c'est essentiellement surjectif. Mais si $x \in \Ob \cE$, $s = f(x)$, on a $x \in \Ob \cE_s$, et comme $\pi_0$ est utile pour $\cE_s$ par $1^{\circ}$), et $\pi^{\pi_0} = \pi > \pi_0$,
$\pi$ est également utile pour $\pi$, donc on a \textit{dans} $\cE_s$ % [sic] utile pour $\pi$ -> utile pour $\cE_s$, usefulness is a property relative to the category $\cE_s$
\[
x \ = \ \varinjlim_{I} x_i \ , \qquad \text{les } x_i \in \cE_{s\,\pi} \subset \cE(\pi),\ I \text{ grande devant } \pi \ ,
\]
\printedpage{165}
% head: A. Grothendieck : Dérivateurs, Ch. XVIII, Édition des Universités de Montpellier II et Paris VII
et d'après a) (comme $\pi \geqslant \pi_1 \geqslant \sup_{\lambda \in \Fl \cC} \pi_{\lambda}$) $x$ est aussi limite des $x_i$ dans $\cE$, donc il est dans l'image essentielle de $(*)$.

d) Reste à prouver que l'on a $\cE_{\pi} \ \subset \ \cE(\pi)$. Mais on sait (par 4.8.~\dots), comme $\cE \ \simeq \ \Ind_{\pi}(\cE(\pi))$, que $\cE_{\pi} \ = \ \cE(\pi)^{\natural}$, i.e. que pour $x \ \in \ \Ob \cE_{\pi}$, $\exists \ y \ \in \ \cE(\pi)$ et un projecteur $p$ dans $x$, tel que $x$ soit isomorphe à l'image de $p$ dans $\cE$. Or $f(p)$ est un projecteur dans $f(y) = s$, donc par hypothèse sur $\cC$, $f(p) = \id_s$, i.e. $p$ est un projecteur dans $\cE_s$. Comme donc $x$ est dans $\cE_s$, et

\msp{page 432}même dans $\cE_{s\,\pi}$, puisque $\cE_{s\,\pi}$ est stable dans $\cE_s$ par facteurs directs. Cela prouve que $x \in \Ob \cE(\pi)$, q.e.d.

On peut présenter 4.15.1 sous une forme plus sympathique.

\textbf{Corollaire 4.15.4.} \textit{Soit $f : \cE \longrightarrow \cC$ un foncteur de $\mathfrak{U}$-catégories. Les conditions suivantes sont équivalentes :}
\begin{enumerate}
\item[\textit{a)}] \textit{Pour tout foncteur $\cC' \longrightarrow \cC$, avec $\cC'$ petite et sans projecteurs, $\cE' \overset{\text{déf}}{=} \cE \times_{\cC} \cC'$ est accessible, et $f' : \cE' \longrightarrow \cC'$ est accessible.}
\item[\textit{b)}] \textit{Comme a), mais avec $\cC' = \boldsymbol{\Delta}_1$, en d'autres termes, c'est la condition $H$ p.~409 : pour toute $\lambda \in \Fl \cC$, $\cE_{\lambda} = \cE \times_{\cC} (\boldsymbol{\Delta}_1, i_{\lambda})$ est accessible, et $\cE_{\lambda} \longrightarrow \boldsymbol{\Delta}_1$ est accessible.}
\item[\textit{c)}] \textit{La condition $H' = H\ \mathit{1} + H\ \mathit{2} + H\ \mathit{3}$, explicitée dans l'énoncé de 4.15.1.}
\end{enumerate}
\textsc{Démonstration}. Elle se fait circulaire,
\[
\mathrm{a} \ \Longrightarrow \ \mathrm{b} \ \Longrightarrow \ \mathrm{c} \ \Longrightarrow \ \mathrm{a} \ ,
\]
l'implication a $\Longrightarrow$ b étant tautologique, b $\Longrightarrow$ c l'étant presque (cf. 4.14.24, a $\Longrightarrow$ b $\Longrightarrow$ c). En fait, les conditions b et c sont des conditions sur les catégories $\cE_s \xrightarrow{f_{\lambda}} \boldsymbol{\Delta}_1$ pour $\lambda \in \Fl \cC$, donc stables par changement de base. Pour prouver

\msp{page 433}c $\Longrightarrow$ a, on peut donc supposer que $\cC' = \cC$, $\cC' \longrightarrow \cC$ le foncteur identique, et prouver que si c) est satisfaite et si $\cC$ est petite et sans projecteurs, alors $\cE$ est accessible et $f$ accessible. Or c'est ce qu'affirme 4.15.1.

\textbf{Remarque 4.15.5.} Si $f : \cE \longrightarrow \cC$ est un foncteur pleinement fidèle, il est immédiat que la condition H est satisfaite (car $\cE_{\lambda} \longrightarrow \boldsymbol{\Delta}_1$ sera pleinement fidèle, donc $\cE_{\lambda}$ est équivalente
à $e$ où à $\boldsymbol{\Delta}_1$, donc accessible, et tout foncteur de $\cE_{\lambda}$ dans une catégorie quelconque est % [sic] où -> ou, disjunction "e ou Δ1"
accessible \dots). Mais si on suppose $\cC$ petite \textit{sans} supposer $\cC$ sans projecteurs, il ne s'ensuit pas que $\cE$ soit accessible, ou ce qui revient au même (étant essentiellement petite), que $\cE$ soit karoubienne. (L'exemple le plus trivial est celui où $\cE = \cC$, $f = \id_{\cC}$, $\cE$ n'est karoubienne que si $\cC$ l'est !) Mais même en supposant $\cC$ une petite catégorie karoubienne, il ne s'ensuit pas forcément que $\cE$ l'est. En effet, \textit{toute} petite catégorie $\cE$ (karoubienne ou non) se plonge comme sous-catégorie pleine dans son enveloppe karoubienne $\cC$ -- il suffit donc de prendre $\cE$ non karoubienne. L'exemple

\msp{page 434}le plus économique est
\[
\Ob \cE \ = \ \{0\} \ , \qquad \Fl(\cE) \ = \ \{\id_0, p\} \ , \qquad \text{avec } p^2 = p \neq \id_0 \ ,
\]
\printedpage{166}
% head: A. Grothendieck : Dérivateurs, Ch. XVIII, Édition des Universités de Montpellier II et Paris VII
l'enveloppe de Karoubi de $\cE$ est donnée par
\[
P \ = \ \left\{
\begin{tikzcd}
1 \arrow[r, shift right, "\sigma"'] & 0 \arrow[l, shift right, "\rho"'] \arrow[loop right, "p"]
\end{tikzcd}
\qquad
\begin{aligned}
\rho\sigma \ &= \ \id_1 \\
\sigma\rho \ &= \ p \ \neq \ \id_0
\end{aligned}
\quad \Longrightarrow \ p^2 = p
\right\} .
\]
Cela soulève la question de trouver un énoncé similaire à 4.15.1, donnant des conditions serviables simples, plus fortes que $\mathrm{H}' = \mathrm{H}$ (1+2+3), sur $f : \cE \longrightarrow \cC$, pour que $\cE$ soit accessible et $f$ accessible, en supposant $\cC$ \textit{karoubienne}. Ou encore, des conditions nécessaires et
suffisants serviables (plus fortes que 4.15.4) sur $f : \cE \longrightarrow \cC$, pour que toute \textit{petite} catégorie % [sic] suffisants -> suffisantes, agreement with "conditions"
\textit{karoubienne} $\cC'$, et tout foncteur $\cC' \longrightarrow \cC$, la catégorie $\cE' = \cE \times_{\cC} \cC'$ soit accessible, et le foncteur $\cE' \longrightarrow \cC'$ accessible. Une idée naturelle, c'est qu'il suffirait que cette conclusion soit vraie quand $\cC'$ est une des deux catégories $\boldsymbol{\Delta}_1$, $P$ (où $P$ est la catégorie à deux objets, à trois flèches non identiques, définie plus haut). Mais je doute qu'il en soit ainsi (\footnote[178]{Ça devrait marcher par contre si on prend les catégories $0 \longrightarrow 1$,
\[
\begin{array}{c}
2 \longrightarrow \underbrace{\begin{tikzcd}[ampersand replacement=\&]
1 \arrow[r, shift right, "\sigma"'] \& 0 \arrow[l, shift right, "\rho"'] \arrow[loop right, "p"]
\end{tikzcd}}_{P}
\quad \text{et} \quad
2 \longleftarrow \underbrace{\begin{tikzcd}[ampersand replacement=\&]
1 \arrow[r, shift right, "\sigma"'] \& 0 \arrow[l, shift right, "\rho"'] \arrow[loop right, "p"]
\end{tikzcd}}_{P}
\qquad (\text{\footnotemark[179]}) \ , \\
\nwarrow \ \begin{array}{c} \text{3 objets et} \\ \text{5 flèches non} \\ \text{identiques} \end{array} \ \nearrow
\end{array}
\]
grâce à 4.15.7, page 440.}\footnotetext[179]{ces catégories représentent respectivement les foncteurs $\mathrm{Cat} \longrightarrow \Ens$ associant à une catégorie $\cC$ les diagrammes dans $\cC$ formés de $(x, y, p, u)$, avec $p$ projecteur dans $p$ \texttt{[plutôt dans $\cC$]}, et $u : y \longrightarrow x^{\natural}$ (i.e. $u : y \longrightarrow x$ avec $u = pu$), resp. $u : x^{\natural} \longrightarrow y$ (i.e. $u : x \longrightarrow y$ avec $u = up$).}).

On va quand même essayer un énoncé un peu très particulier dans ce sens :

% Verified on PDF leaf8: a non-textual wavy marginal mark beside items a), b); no mathematical token omitted.
\msp{page 435}\textbf{Proposition 4.15.6} (\footnote[180]{\textit{Faux ?}}). \textit{Soit $f : \cE \longrightarrow \cC$ un foncteur de $\mathfrak{U}$-catégories. On suppose :}
\begin{enumerate}
\item[\textit{a)}] \textit{$f$ est fibrant ou cofibrant} (\footnote[181]{Le cas fibrant n'est pas prouvé, peut-être \textit{faux} -- et de même pour le cas cofibrant !}).
\item[\textit{b)}] \textit{Les catégories fibres $\cE_s$ ($s \in \Ob \cC$) sont accessibles, et pour toute flèche $\lambda : s \longrightarrow t$ dans $\cC$, le foncteur changement de base $\lambda^* : \cE_t \longrightarrow \cE_s$ (resp. cochangement de base $\lambda_* : \cE_s \longrightarrow \cE_t$) est accessible.}
\item[\textit{c)}] \textit{$\cC$ est essentiellement petite et karoubienne.}
\end{enumerate}
\textit{Alors $\cE$ est accessible, et $f$ est accessible.}

\textsc{Démonstration}. On a déjà vu que les conditions a) et b) impliquent $\mathrm{H}'$, i.e. H (1+2+3). D'ailleurs, si $\cC'$ est une catégorie et $\cC' \xrightarrow{\ \sim\ } \cC$ une équivalence, donc alors $\cE' \overset{\text{déf}}{=} \cE \times_{\cC} \cC' \longrightarrow \cE$ est une équivalence, donc la conclusion voulue pour $\cE$, $f$ équivaut à la même conclusion pour $\cE'$, $f' : \cE' \longrightarrow \cC'$. Donc on peut supposer que $\cC$ est petite (et même réduite, si on y tient). En vertu de 4.15.2 a), on voit donc que les foncteurs d'inclusion $\cE_s \longrightarrow \cE$ sont accessibles (et même $\pi_1$-accessibles) dès que $\pi_1 \geqslant \sup_{\lambda \in \Fl \cC} \pi_{\lambda}$. On voit d'ailleurs

\msp{page 436}par AQT \texttt{[âne qui trotte]} qu'on peut trouver des cardinaux $\pi_0$, $\pi_1$ ayant les propriétés suivantes :

\printedpage{167}
% head: A. Grothendieck : Dérivateurs, Ch. XVIII, Édition des Universités de Montpellier II et Paris VII
\begin{enumerate}
\item[\circled{1\textsuperscript{o}}] $\pi_0$ est utile pour $\cE_s$, $\forall \ s \in \Ob \cC$.
\item[$2^{\circ}$)] $\pi_1 \geqslant \sup(\pi_0, \card \cC, \sup_{\lambda \in \Fl \cC} \pi_{\lambda}, \sup_{s \in \Ob \cC} \card \red \cE_s)$.
\item[\circled{3\textsuperscript{o}}] $\forall \ (\lambda : s \longrightarrow t) \in \Fl \cC$ et $x_0 \in \Ob \cE_{s\,\pi_0}$, le foncteur
\[
H_{\lambda}(x_0) \ : \ (x_1 \longmapsto \Hom_{\lambda}(x_0, x_1)) \ : \ \cE_t \ \longrightarrow \ (\Ens)
\]
est $\pi_1$-accessible. (\textbf{N.B.} Si $f$ est cofibrant, cela signifie que $\forall \ \lambda \in \Fl \cC$, $\lambda : s \longrightarrow t$,
\fbox{$\lambda_*(\cE_{s\,\pi_0}) \ \subset \ \cE_{s\,\pi_1}$}. Si $f$ est fibrant, il suffit d'exiger que $\forall \ \lambda \in \Fl \cC$, $\lambda^* : \cE_t \longrightarrow \cE_s$ % [sic] $\cE_{s\,\pi_1}$ -> $\cE_{t\,\pi_1}$, $\lambda_*$ goes to $\cE_t$, cf. the later display "(contenu dans 3°)"
soit $\pi_1$-accessible.) Mais dans ce cas, on veut aussi \fbox{$\lambda^*(\cE_{s\,\pi_0}) \ \subset \ \cE_{s\,\pi_1}$}. % [sic] $\lambda^*(\cE_{s\,\pi_0})$ -> $\lambda^*(\cE_{t\,\pi_0})$, $\lambda^*$ is defined on $\cE_t$, cf. the later display "(contenu dans 3°)"
\item[\circled{4\textsuperscript{o}}] Les foncteurs $\lambda^*$ (cas $f$ fibrant) resp. $\lambda_*$ (cas $f$ cofibrant) sont $\pi_1$-accessibles.
\end{enumerate}
Soit alors $\pi$ un cardinal tel que
\[
\pi \ = \ \pi^{\pi_0} \ \geqslant \ \pi_1 \ , \tag*{$(*)$}
\]
et soit $\cE(\pi)$ la sous-catégorie (strictement) pleine de $\cE$ définie par
\[
\Ob \cE(\pi) \ = \ \bigcup_{s \, \in \, \Ob \cC} \Ob \cE_{s\,\pi} \ .
\]
Grâce à $(*)$, et à $1^{\circ}$, $4^{\circ}$ et à
\[
\left\{
\begin{array}{l}
\pi_1 \ \geqslant \ \sup_s(\card \red \cE_{s\,\pi_0}) \qquad \text{(contenu dans $2^{\circ}$)} \\[1ex]
\lambda^*(\cE_{t\,\pi_0}) \ \subset \ \cE_{s\,\pi_1} \qquad \text{resp.} \qquad \lambda_*(\cE_{s\,\pi_0}) \ \subset \ \cE_{t\,\pi_1} \qquad \text{(contenu dans $3^{\circ}$)} \ ,
\end{array}
\right.
\]
on voit que $\cE(\pi)$ est une sous-catégorie

\msp{page 437}fibrée, resp. une sous-catégorie cofibrée, sur $\cC$ (le foncteur $\cE(\pi) \longrightarrow \cE$ étant donc cartésien resp. cocartésien). À $\cC$-équivalence près, $\cE$ se reconstitue alors à partir de $\cE(\pi)$, en passant au $\Ind_{\pi}(\cE_{s\,\pi})$, qui forment un système fibré resp. cofibré sur $\cC$ par les foncteurs $\Ind_{\pi}(\lambda^*_{\pi})$ resp. $\Ind_{\pi}(\lambda^{\pi}_*)$ \dots

Prouvons que $\cE$ satisfait à $\mathrm{L}_{\pi}$ et que $f : \cE \longrightarrow \cC$ est $\pi$-accessible (\footnote[182]{Ici on utilise seulement $\pi \geqslant \card \red \cC$.}). Comme $\cC$ est karoubienne, et $\pi \geqslant \pi_1 \geqslant \card \cC$, $\cC$ satisfait à $\mathrm{L}_{\pi}$ et les ind-objets grands devant $\pi$ de $\cC$ sont essentiellement constants (4.8.13). Soit alors $I$ grande devant $\pi$, et $I \longrightarrow \cE$, $i \longmapsto x_i$, un foncteur, il faut construire une $\varinjlim$ dans $\cE$, en utilisant la $\varinjlim$ déjà connue $s$ des $s_i = f(x_i)$. On va prendre un $x$ dans $\cE_s$. Prenons d'abord le cas où $f$ est cofibrant. On voit que les $\beta_i(x_i)$ forment un système inductif dans $\cE_s$, puisque pour $i = j$, on a
\[
\beta_{i\,*}(x_j) \ = \ \beta_{j\,*} \underbrace{\beta_{ji\,*}(x_i)}_{\longrightarrow \, x_j} \ \longrightarrow \ \beta_{j\,*}(x_j) \ . % [sic] $\beta_{i\,*}(x_j)$ -> $\beta_{i\,*}(x_i)$, the object transported from $s_i$ is $x_i$
\]
\[
\begin{tikzcd}[row sep=small]
x_i \arrow[rr, "\alpha_i"] & & x \, ? \\
s_i \arrow[rr, "\beta_i"] \arrow[dr, "\beta_{ji}"'] & & s \\
& s_j \arrow[ur, "\beta_j"'] &
\end{tikzcd}
\]
On pose

\printedpage{168}
% head: A. Grothendieck : Dérivateurs, Ch. XVIII, Édition des Universités de Montpellier II et Paris VII
\msp{page 438}
\[
x \ = \ \varinjlim_{I} \beta_{i\,*}(x_i) \ , \qquad \varinjlim \ \textit{dans } \cE_s \ .
\]
On sait que c'est une $\varinjlim$ dans $\cE$, donc que pour tout $y$ dans $\cE$, on a
\[
\Hom(x, y) \xrightarrow{\ \sim\ } \varprojlim_{I} \Hom(\beta_{i\,*}(x_i), y) \ ,
\]
ou ce qui revient au même, pour toute flèche $\lambda : s \longrightarrow t$, et $y$ dans $\Ob \cE_t$, on a
\[
\underbrace{\Hom_{\lambda}(x, y)}_{= \, \Hom_{\cE_t}(\lambda_* x, y)} \ \simeq \ \varprojlim \Hom_{\cE_t}(\underbrace{\lambda_* \beta_{i\,*}(x_i)}_{(\lambda_i \beta_i)_*(x_i)}, y) \ .
\]
Cela résulte aussi du fait que $\lambda_*$ est $\pi$-accessible (étant $\pi_1$-accessible) \dots{} On vérifie facilement que c'est une $\varinjlim$ des $x_i$ dans $\cE$. Et bien sûr, cette limite commute avec $f$ par construction.

Mais dans le cas où $f$ est fibrant, je ne vois comment faire pour construire une limite dans $\cE_s$. Il est vrai qu'on peut supposer que tous les $s_i$ sont égaux à $s$, mais les transitions $\beta_{ji}$ ne sont pas l'identité, ni même des isomorphismes. Il faudrait prendre $\cC = P$ pour essayer de comprendre -- je ne vais pas insister (\footnote[183]{J'ai vérifié que si $f$ est fibrant, alors $\cE$ est \textit{karoubienne}. Il y a donc de l'espoir \dots}).

\msp{page 439}Dans le cas où $f$ est fibrant, admettons (comme hypothèse supplémentaire au besoin) que $\cE$ satisfasse $L$ et que $f$ soit accessible, et prenons $\pi$ assez grand pour que $\cE$, $f$ satisfassent $\mathrm{L}_{\pi}$. Nous pouvons alors terminer la démonstration de 4.15.6. Il faut prouver
\[
\cE(\pi) \ \subset \ \cE_{\pi} \ ,
\]
donc que si $x \in \Ob \cE_{s\,\pi}$, et $(y_i)_{i \in I}$ est un système inductif dans $\cE$ grand devant $\pi$, de limite $y$, on a
\[
\varinjlim \Hom_{\cE}(x, y_i) \xrightarrow{\ \sim\ } \Hom_{\cE}(x, y) \ .
\]
Donc il faut prouver que toute flèche $x \xrightarrow{u} y$ se factorise, de façon essentiellement unique $u_i$, par un $y_i$. Mais la flèche $s \xrightarrow{\ \lambda \, = \, f(u)\ } t$ dans $\cC$ se factorise, de façon essentiellement unique, par des $\lambda_i \, : \, s \longrightarrow t_i$. Donc on est ramené à trouver un \textit{$\lambda_i$-morphisme} $x \longrightarrow y_i$ pour $i$ grand, qui factorise. Dans le cas fibrant, cela signifie qu'on cherche $x \longrightarrow \lambda^*_i(y_i)$ qui factorise $x \longrightarrow \lambda^*(y)$, or les $\lambda^*_i(y_i)$ forment un système inductif dans $\cE_s$, et sa limite \textit{dans} $\cE_s$ est $\lambda^*(y)$ (\textit{à vérifier}). Comme $x$ est dans $\cE_{s\,\pi}$, l'hom\texttt{[omorphisme]} donné $x \longrightarrow \lambda^*(y)$ se factorise de

\msp{page 440}façon essentiellement unique par un $\lambda^*_i(y_i)$.

Cette fois par contre, c'est le cas $f$ cofibrant qui pose problème. En fait, je n'arrive pas à m'en tirer.

Ainsi, on trouve bien un foncteur
\[
\Ind_{\pi}(\cE(\pi)) \ \longrightarrow \ \cE \ ,
\]


\endgroup
% END UNIT der-ch18/der-ch18_p166-170.tex
% BEGIN UNIT der-ch18/der-ch18_p171-175.tex
\begingroup\makeatletter%
\newcommand{\printedpage}[1]{\marginpar{\footnotesize p.~#1}}%
\newcommand{\msp}[1]{\marginpar{\footnotesize [#1]}}%
\newcommand{\ux}{\underline{x}}%
\newcommand{\cC}{\mathcal{C}}%
\newcommand{\Ob}{\operatorname{Ob}}%
\newcommand{\card}{\operatorname{card}}%
\newcommand{\id}{\mathrm{id}}%
\newcommand{\Hom}{\operatorname{Hom}}%
\newcommand{\Ind}{\operatorname{Ind}}%
\newcommand{\Img}{\operatorname{Im}}%
\newcommand{\Ker}{\operatorname{Ker}}%
\newcommand{\defeq}{\overset{\text{\upshape déf}}{=}}%
\makeatother\let\printedpage\derpp\let\msp\dermsp%
\let\mspin\DERmspin%

\setcounter{footnote}{183}

\printedpage{169}
% head: A. Grothendieck : Dérivateurs, Ch. XVIII, Édition des Universités de Montpellier II et Paris VII
qui est essentiellement surjectif, mais faute de savoir que $\mathcal{E}(\pi) \subset \mathcal{E}_\pi$ dans le cas cofibrant, il n'est pas clair que ce foncteur soit une équivalence de catégories.

\textbf{4.15.7.} Il me faut reporter la démonstration du lemme 4.15.3 (p. 428), en admettant 4.8.13. (Non encore démontré -- je vais revenir plus bas sur la démonstration de 4.8.13.) On va utiliser le

\textbf{Lemme 4.15.7.1.} \textit{Soient $\cC$ une catégorie, $\ux = (x_i)_{i\in I}$ un ind-objet de $\cC$. Pour que $\ux$ soit essentiellement constant, i.e. isomorphe, dans} $\Ind(\cC)$\textit{, au ind-objet constant défini par un objet}

\msp{page 441}
\textit{de $\cC$, il faut et il suffit qu'il existe un foncteur cofinal $\varphi : J \longrightarrow I$, avec $J$ filtrant, de telle façon que le système réindexé $(x_j)_{j\in J}$, avec $x_j = x_{\varphi(j)}$, ait la propriété suivante : pour tout $j \in J$, on peut trouver un projecteur $p_j$ dans $x_j$, de telle façon que :}
\begin{itemize}
\item[\textit{1°)}] \textit{Pour toute flèche $\lambda : j \longrightarrow j'$ dans $J$, on ait commutativité dans}
\[
\begin{tikzcd}
x_j \arrow[r, "x_\lambda"] \arrow[d, "p_j"] & x_{j'} \arrow[d, "p_{j'}"] \\
x_j \arrow[r, "x_\lambda"] & x_{j'}
\end{tikzcd}
\]
\textit{(où les $x_\lambda$ sont les flèches de transition), i.e. $p_{j'}x_\lambda = x_\lambda p_j$.}
\item[\textit{2°)}] \textit{Pour tout $j \in \Ob J$, il existe une flèche $\lambda : j \longrightarrow j'$ dans $J$ telle que $p_{j'}x_\lambda = x_\lambda p_j = x_\lambda$.}
\item[\textit{3°)}] \textit{Pour tout $j \in \Ob J$, le projecteur $p_j$ dans $x_j$ admet un image $x_j^\natural$.} % [sic] un image -> une image, gender agreement
\item[\textit{4°)}] \textit{Pour toute flèche $\lambda : j \longrightarrow j'$ dans $J$, la flèche $x_\lambda^\natural : x_j^\natural \longrightarrow x_{j'}^\natural$ induite par $x_\lambda$ (grâce à 1°) est un isomorphisme.}
\end{itemize}
\textit{De plus, si on peut supposer $J$ de la forme $i_0\backslash I$, le foncteur cofinal $J \longrightarrow I$ étant le foncteur canonique $i_0\backslash I \longrightarrow I$ ($i_0$ étant un objet fixé convenable de $I$)} \texttt{[phrase incomplète]}

\msp{page 442}
Notons tout de suite deux corollaires.

\textbf{Corollaire 4.15.7.2.} \textit{Supposons que $\cC$ soit sans projecteurs. Alors $\ux = (x_i)_{i\in I}$ est un ind-objet essentiellement constant si et seulement s'il existe un foncteur cofinal, $J \longrightarrow I$, avec $J$ filtrant (\footnote{et on peut supposer $J = i_0\backslash I \xrightarrow{\text{can.}} I$}), de telle façon que le système inductif réindexé $(x_j)_{j\in J}$ soit isomorphe} (dans la catégorie $\underline{\mathrm{Hom}}(J,\cC)$) \textit{à un système inductif} constant \textit{(i.e. défini par un foncteur constant $J \longrightarrow \cC$ : tous les $x_j$ égaux à un même $x$, et les morphismes de transition des identités).}

La suffisance est claire, pour la nécessité on applique 4.15.7.1, qui nous dit qu'on peut réindexer de façon à satisfaire aux conditions 1°) à 4°) (pour des projecteurs $p_j$ convenables dans les $x_j$). Mais l'hypothèse sur $\cC$ implique que $p_j = \id_{x_j}\ \forall\, j \in \Ob J$, et les conditions 1°, 2°, 3° sont satisfaites tautologiquement pour un tel choix, la condition 4° signifie que les morphismes de transition sont des isomorphismes. Mais il est immédiat que pour un système inductif $(x_j)_{j\in J}$ donné, avec $J$ 1-\textit{connexe} (p.ex. $J$ filtrant), la condition que les $x_j$ sont des isomorphismes % [sic] x_j -> x_\lambda, the transition morphisms are meant
est nécessaire et suffisante pour que ce système soit isomorphe à un système inductif constant, d'où la conclusion.
% foot: Version 17 décembre 2025

\printedpage{170}
% head: A. Grothendieck : Dérivateurs, Ch. XVIII, Édition des Universités de Montpellier II et Paris VII
\msp{page 443}
\textbf{Corollaire 4.15.7.3.} \textit{Supposons que $I$ soit grande devant $\pi_0$, où $\pi_0$ est un cardinal infini tel que $\pi_0 \geqslant \card \cC$. Alors $(x_i)_{i\in I}$ est essentiellement constant si et seulement si on peut trouver un foncteur cofinal $J \longrightarrow I$ avec $J$ grande devant $\pi_0$ (et ordonnée si on veut) (\footnote{La démonstration (légèrement réarrangée) montre qu'on peut prendre pour $J$ une sous-catégorie (pas pleine) d'une catégorie $i_1\backslash I$, pour un objet convenable $i_1$ de $I$. On a donc $\card J \leqslant \pi$ si $\card I \leqslant \pi$, $\pi$ étant un cardinal infini donné.}), de telle façon que le système réindexé $(x_j)_{j\in J}$ satisfasse les conditions suivantes :}
\begin{itemize}
\item[\textit{1°)}] \textit{L'application $j \longmapsto x_j$ de $\Ob J$ dans $\Ob\cC$ est constante, i.e. il existe un $x_0 \in \Ob\cC$ tel que $x_j = x_0\ \forall\, j \in \Ob J$.}
\item[\textit{2°)}] \textit{Il existe un projecteur $p$ dans $x_0$ tel que pour $j < j'$ dans $J$, on ait $\tau_{j',j} = p$.}
\item[\textit{3°)}] \textit{Le projecteur $p$ admet une image dans $\cC$.}
\end{itemize}
\textsc{Démonstration}. Ces conditions impliquent évidemment celles de 4.15.7.1, en prenant $p_j = p\ \forall\, j \in \Ob J$, du moins si $J$ n'admet de plus grand élément (ce qui fait une difficulté pour 2°). Mais si $J$ admet un objet final, il en est de même de $I$, et il est clair alors que tout système inductif indexé par $I$ est essentiellement constant. Donc, dans ces cas, les conditions 1°) 2°) 3°) ci-dessus impliquent que $\ux$ est essentiellement constant. Inverse\msp{page 444}ment, supposons que $\ux$ soit essentiellement constant, et prouvons que l'on peut réindexer comme indiqué. Le cas où $I$ admet un objet final est trivial, on prendra $J = e$, $p = \id$. Supposons donc que $I$ n'ait pas d'objet final. Comme $I$ grande devant $\pi_0$ et $\pi_0 > \card\Ob\cC$, on voit qu'il existe un $x_0 \in \Ob\cC$ tel que la sous-catégorie pleine $I' = I_{x_0}$ de $I$, formée des $i$ tels que $x_i = x_0$, soit cofinale dans $I$. Comme $(x_i)_{i\in I}$ est essentiellement constant, il en est de même de $(x_{i'})_{i'\in I'}$, et on peut lui appliquer le critère 4.15.7.1, avec \marginpar{$\Vert$}% double vertical bar printed in the left margin beside the two following lines
$J \longrightarrow I'$ cofinal. La démonstration montrera de plus que si $I'$ est grand devant $\pi_0$ (cardinal infini sans autres conditions), on peut supposer $J$ grande devant $\pi_0$ (\footnote{en tenir compte dans la démonstration de 4.15.7.1.}). De plus, il est clair qu'on peut supposer $J$ ordonnée. Comme $\pi_0 \geqslant \card\Hom(x_0,x_0)$, on voit encore qu'il existe un projecteur $p$ dans $x_0$ tel que la partie $J_1 = J'$ de $J$, formée des $j$ tels que $p_j = p$, soit cofinale. Étant cofinale, elle est également grande devant $\pi_0$. Donc on peut supposer que tous les $p_j$ sont égaux à un même $p$. Les conditions 1°-4° deviennent :
\begin{itemize}
\item[1°)] $p\tau_{j'j} = \tau_{j'j}p$ pour $j \leqslant j'$.
\item[2°)] Pour tout $j$, existe $j' \geqslant j$ avec
\[
p\tau_{j'j} \;=\; \tau_{j'j}p \;=\; \tau_{jj'}\ . % [sic] \tau_{jj'} -> \tau_{j'j}, index order used everywhere else
\]
\item[3°)] L'image $x_0^\natural$ du projecteur $p$ dans $x_0$ existe.
\item[4°)] L'endomorphisme $\tau^\natural_{j'j}$ de $x_0^\natural$ est un isomorphisme, pour $j \leqslant j'$.
\end{itemize}
\msp{page 445}
Or on a le

\textbf{Lemme 4.15.7.4.} \textit{Soient $\cC$ une catégorie, $x_0$ un objet de $\cC$, $p$ un projecteur dans $x_0$, $\tau$ un endomorphisme de $x_0$ satisfaisant les relations}
\[
\tau p \;=\; p\tau \;=\; \tau\ , \tag*{$(*)$}
\]
% foot: Version 17 décembre 2025
\printedpage{171}
% head: A. Grothendieck : Dérivateurs, Ch. XVIII, Édition des Universités de Montpellier II et Paris VII
\textit{et tel que l'automorphisme de $x_0^\natural \defeq$} $\Img p$ \textit{(défini p.ex. comme objet de $\cC^\wedge$, $x_0^\natural =$} $\Ker(\id_{x_0},p)$ % [sic] "(défini ... Ker(id_{x_0},p)" -> a closing parenthesis is missing, the parenthesis opened before "défini" is never closed
\textit{soit l'identité. Alors on a $\tau = p$.}

\textsc{Démonstration de} 4.15.7.4. On est ramené aussitôt, par le procédé standard, au cas où $\cC = \mathrm{Ens}$, où la vérification est immédiate : les automorphismes $\tau$ de l'ensemble $E$ $(= x)$ qui satisfait $(*)$ correspondent biunivoquement aux endomorphismes $u$ de $E^\natural = pE$, via
\[
\begin{cases}
u \;\longmapsto\; \sigma u \rho \\
\tau \;\longmapsto\; \tau|E^\natural
\end{cases}
\qquad
\text{(où }
\begin{tikzcd}[column sep=small]
E^\natural \arrow[r, hook, shift left, "\sigma"] & E \arrow[l, shift left, "\rho"]
\end{tikzcd}
\text{ sont les flèches canoniques)}\ ,
\]
donc $\tau|E^\natural = \id \Longrightarrow u = \sigma\rho = p$.

Choisissons un $j_0 \in J$, et considérons les homomorphismes
\[
\tau^\natural_{j,j_0} \;:\; x_0^\natural \longrightarrow x_0^\natural
\]
induits par les $\tau_{j,j_0}$ avec $j \geqslant j_0$. Ici encore, comme l'ensemble des isomorphismes possibles est de

\msp{page 446}
cardinal $\leqslant \pi_1$, et $J/j_0$ est grande devant $\pi_0$, on voit qu'il existe $J' \subset j_0\backslash J$, cofinale dans $j_0\backslash J$, donc dans $J$, telle que pour $j \in J'$, les $\tau^\natural_{j,j_0}$ soient égaux à un même $\tau : x_0^\natural \xrightarrow{\sim} x_0^\natural$. Donc si $j', j'' \in J'$, avec $j' \leqslant j''$, on a
\[
\underbrace{\tau_{j''j_0}}_{\tau} \;=\; \tau_{j''j'}\,\underbrace{\tau_{j'j_0}}_{\tau}\ ,
\]
d'où
\[
\tau_{j''j'} \;=\; \id_{x_0^\natural}\ ,
\]
puisque la flèche $\tau$ est un isomorphisme. Donc le système inductif réindexé par $J'$, la condition 4° est remplacée par la condition plus stricte \fbox{4° bis) $\tau^\natural_{j',j} = \id_{x_0^\natural}$.} Mais la condition 2° devient alors, grâce à 4.15.7.4,
\[
\boxed{\text{2° bis) Pour tout } j\text{, existe } j' \geqslant j \text{ avec } \tau_{j',j} = p}\ .
\]
Donc, on peut supposer 1°, 2° bis, 3°, 4° bis (on désigne à nouveau par $J$, au lieu de $J'$, l'ensemble ordonné grand devant $\pi_0$ indexant).

Soit $K$ l'ensemble ordonné défini par

\msp{page 447}
\[
\begin{cases}
\Ob K \;=\; \Ob J \\
\text{si } j,j' \in \Ob J\text{, on a } j \overset{K}{\leqslant} j' \text{ si et seulement si } j \overset{\mathrm{J}}{=} j'\text{, ou } j \overset{J}{<} j' \text{ et } \tau_{j',j} = p\,.
\end{cases}
\]
Il est immédiat que c'est une \textit{relation d'ordre} sur $J$ (car $p^2 = p$), plus fin que la relation d'ordre initiale. Le système inductif induit sur $K$ satisfait aux conditions 1°, 2°, 3° de 4.15.7.3. D'autre part, $K \longrightarrow J$ est cofinal, donc aussi le composé $K \longrightarrow J \longrightarrow I$. Il reste à prouver que $K$ est grande devant $\pi_0$, donc que toute partie $K' \subset K$ telle que $\card K' \leqslant \pi_0$ admet un majorant dans $K$. Elle admet un majorant \textit{dans} $J$, soit $j$. Par hypothèse 2 bis, il existe $j' \geqslant j$ tel que $\tau_{j'j} = p$, et on peut supposer $j' > j$. (En effet, si $\tau_{j'j} = p$, on voit que
% foot: Version 17 décembre 2025
\printedpage{172}
% head: A. Grothendieck : Dérivateurs, Ch. XVIII, Édition des Universités de Montpellier II et Paris VII
\msp{page 448}
la forme initiale 2°) de la condition 2° bis que l'on a $\tau_{j''j} = p$ pour tout $j'' \geqslant j'$, donc si $j$ n'est par objet final de $J$, % [sic] par -> pas, misprint
on peut supposer $j' > j$. Mais $J$ n'a pas d'objet final, car sinon $I$ en aurait un, cas que l'on avait éliminé. Donc il existe $j'$ tel que $j' \overset{K}{>} j$.) Or on a ceci (qu'on applique au cas $i \in K'$) :
\[
i \leqslant j \ \text{ et } \ j \overset{K}{<} j' \qquad \Longrightarrow \qquad i \overset{K}{<} j'\ .
\]
\[
\underbrace{x_i}_{= x_0} \xrightarrow{\tau_{j,i}} x_j \xrightarrow{\tau_{j',j}=p} x_{j'}\ ,
\]
\[
\tau_{j'i} \;=\; \tau_{j'j}\tau_{ji} = \underbrace{p\tau_{ji}}_{\overset{\text{1°}}{=}\,\tau_{ji}p} \;\overset{?}{=}\; p\ ,
\]
on a en effet
\[
\left.
\begin{aligned}
\tau_{j',i} &\;=\; p\tau_{j',i} \\
\tau_{j',i} &\;-\; \tau_{j',i}p % [sic] - -> =, the second relation of (*) is meant
\end{aligned}
\right\}
\quad \overset{\text{par 4.15.7.4}}{\Longrightarrow} \quad % Doubt resolved against PDF leaf 14: no underline; the two strokes belong to the implication arrow.
\tau_{j',i} \;=\; p\ .
\]
Cela achève la démonstration du corollaire 4.15.7.3.

\textbf{Corollaire 4.15.7.5} (de 4.15.7.3). \textit{Soit $(x_i)_{i\in I}$ un système inductif essentiellement constant grand devant $\pi_0 \geqslant$} $\card\cC$\textit{, dans $\cC$} sans projecteurs. \textit{Alors il existe $J \longrightarrow I$ cofinal, avec $J$ ordonné grand devant $\pi_0$, et tel que le système inductif induit sur $J$ soit} constant.

\textsc{Démontrons le lemme} 4.15.7.1. Il faut expliciter l'existence de deux flèches dans $\Ind(\cC)$,
\[
z \underset{\beta}{\overset{\alpha}{\rightleftarrows}} \text{"}\varinjlim_{i\in I}\text{"} x_i \;=\; \ux\ ,
\]
avec $\beta\alpha = \id_z$, $\alpha\beta = \id_{\ux}$. La donnée

\msp{page 449}
de $x$ équivaut à celle d'un élément de $\varinjlim\limits_{i\in I} \Hom(z,x_i)$, et est donc donné par un % [sic] x -> \alpha, parallel with "la donnée de \beta" below
\[
\alpha_{i_0} \;:\; z \longrightarrow x_{i_0}\ ,
\]
la donnée de $\beta$ équivaut à celle d'un élément de $\varprojlim\limits_{i\in I} \Hom(x_i,z)$, i.e. de flèches
\[
\beta_i \;:\; x_i \longrightarrow z\ ,
\]
avec la condition de compatibilité habituelle. Quitte à remplacer $I$ par $I'$, avec $I' \longrightarrow I$ cofinal, on peut supposer $I$ ordonnée (et grande devant $\pi_0$, si $I$ initiale l'était), et la condition de compatibilité sur les $\beta_i$ s'écrit
\[
\beta_i\tau_{ji} \;=\; \beta_j \qquad \text{si } i \leqslant j \text{ dans } I,
\]
où $\tau_{ji} : x_i \longrightarrow x_j$ désigne la flèche de transition. Il faut expliciter, en termes de ces données, les conditions
\[
\alpha\beta \;=\; \id_z\ , \qquad \beta\alpha \;=\; \id_{\ux}\ .
\]
% foot: Version 17 décembre 2025

\printedpage{173}
% head: A. Grothendieck : Dérivateurs, Ch. XVIII, Édition des Universités de Montpellier II et Paris VII
\[
\begin{tikzcd}
z \arrow[r, "\alpha_{i_0}"] \arrow[dr] & x_{i_0} \arrow[r, "\tau_{i,i_0}"] \arrow[d, "\beta_{i_0}"] & x_i \arrow[dl, "\beta_i"] \\
 & z &
\end{tikzcd}
\]
La relation $\alpha\beta = \id_z$ équivaut à
\[
\beta_{i_0}\alpha_{i_0} \;=\; \id_z\ , \tag*{$(*_{i_0})$}
\]
elle implique d'ailleurs, en posant pour $i \geqslant i_0$,
\[
\alpha_i \;\defeq\; \tau_{i,i_0}\alpha_{i_0} \;:\; z \longrightarrow x_i\ ,
\]
que l'on a pour $i \geqslant i_0$
\[
\beta_i\alpha_i \;=\; \id_z \qquad \forall\, i \geqslant i_0\ . \tag*{$(*_i)$}
\]

\msp{page 450}
Posons alors
\[
p_i \;=\; \alpha_i\beta_i \;:\; x_i \longrightarrow x_i\ ,
\]
il est clair qu'on a
\[
\begin{cases}
p_i^2 \;=\; p_i \qquad \text{($p_i$ un projecteur)} \\
\alpha_i \;:\; z \longrightarrow x_i \qquad \text{induit un isomorphisme} \\
\qquad z \xrightarrow{\sim} x_i^\natural \;\defeq\; \Img p_i \;=\; \Ker(\id_{x_i},p_i)\ .
\end{cases}
\]
L'ensemble ordonné cofinal cherché $J$ est $i_0\backslash I$, il est filtrant, et est grand devant $\pi_0$ si $I$ l'est. À vrai dire, la réduction au cas $I$ ordonné est inutile, on peut prendre directement $J = i_0\backslash I$, il est immédiat que $J$ est filtrante ($I$ l'étant), et $i_0\backslash I \longrightarrow I$ cofinal, enfin que si $I$ est grande devant un cardinal infini $\pi_0$, il en est de même de $i_0\backslash I$.

Exprimons que $\beta\alpha = \id_{\ux}$ ; or on voit que $\beta\alpha$ est défini par l'endomorphisme $(p_j)_{j\in J}$ du système inductif $(x_j)_{j\in J}$ -- le fait que ce soit bien un endomorphisme de ce système, i.e. qu'on ait la relation 1°) de 4.15.7.1,
\[
p_{j'}x_\lambda \;=\; x_\lambda p_\lambda \qquad \text{si } \lambda : j \longrightarrow j' \text{ dans } J, % [sic] p_\lambda -> p_j, cf. relation 1°) of 4.15.7.1
\]
est immédiat, ça s'écrit
\[
\alpha_{j'}\underbrace{\beta_{j'}x_\lambda}_{\beta_j} \;=\; \underbrace{x_\lambda\alpha_j}_{\alpha_{j'}}\beta_j\ , \qquad \text{or les deux membres sont égaux à } \alpha_{j'}\beta_j.
\]
% foot: Version 17 décembre 2025


\endgroup
% END UNIT der-ch18/der-ch18_p171-175.tex
% BEGIN UNIT der-ch18/der-ch18_p176-180.tex
\begingroup\makeatletter%
\newcommand{\printedpage}[1]{\marginpar{\footnotesize p.~#1}}%
\newcommand{\msp}[1]{\marginpar{\footnotesize [#1]}}%
\newcommand{\card}{\operatorname{card}}%
\makeatother\let\printedpage\derpp\let\msp\dermsp%
\let\mspin\DERmspin%


\printedpage{174}
% head: A. Grothendieck : Dérivateurs, Ch. XVIII, Édition des Universités de Montpellier II et Paris VII

\msp{page 451}
Dire que c'est l'identité, est exprimé par la condition 2°) sur $(x_j,p_j)_{j\in J}$. La condition 3° est clair par construction, $p_j$ admettant une image isomorphe à $z$ par $\alpha_j$. D'autre part, on voit \texttt{[tout]} de suite que pour une flèche $\lambda : j \longrightarrow j'$ de $J$, on a commutativité dans % [sic] est clair -> est claire, agreement with "condition"
\[
\begin{tikzcd}
x_j \arrow[r, "x_\lambda"] & x_{j'} \\
x_j^\natural \arrow[u, hook] \arrow[r, "x_\lambda"] & x_{j'}^\natural \arrow[u, hook] \\
z \arrow[u, "\wr", "\overline{\alpha_j}"'] \arrow[r, "\mathrm{id}_z", "\sim"'] & z \ , \arrow[u, "\wr", "\overline{\alpha_{j'}}"']
\end{tikzcd}
\]
où $\overline{\alpha_j}$, $\overline{\alpha_{j'}}$ sont les isomorphismes définis par $\alpha_j$, $\alpha_{j'}$ respectivement. Cela prouve que $x_\lambda^\natural$ est un isomorphisme.

Cela prouve que si $(x_i)_{i\in I}$ est essentiellement constant, il existe $J\longrightarrow I$ comme indiqué. L'inverse est également vrai, car on peut supposer $J = I$, puisque $(x_i)_{i\in I}$ et $(x_j)_{j\in J}$ définissent des ind-objets isomorphes. Alors le système inductif des $x_j^\natural$ a des morphismes de transition des isomorphismes, et comme $J$ est filtrante donc 1-connexe, on en déduit facilement qu'il est isomorphe
\msp{page 452}
à un système inductif constant, de valeur $z$. En choisissant un tel isomorphisme
\[
z \xrightarrow[\sim]{\overline{\alpha_i}} x_i^\natural \ ,
\]
on trouve $z \xrightarrow{\alpha} \underline{x}$, et le système des $p_i$ définit des
\[
\begin{tikzcd}
x_i \arrow[r, "\overline{p_i}"] \arrow[rr, bend right=12, "\beta_i \overset{\text{déf}}{=} \overline{\alpha_j}^{-1}\overline{p_i}"'] & x_i^\natural \arrow[r, "\overline{\alpha_j}^{-1}"] & z \ ,
\end{tikzcd}
\] % [sic] \overline{\alpha_j} (index j) -> \overline{\alpha_i}, the arrow goes from x_i^\natural, cf. z -> x_i^\natural above
et on voit facilement que l'on aura
\[
\beta\alpha \ = \ id_z \ , \qquad \alpha\beta \ = \ \mathrm{id}_{\underline{x}} \ ,
\]

\hfill q.e.d.

\textsc{Je peux maintenant prouver} 4.15.3. C'est en effet essentiellement la conjonction des 4.15.7.5 (p. 448), et de 4.8.13, qui nous dit essentiellement que si $\mathcal{C}$ est \textit{karoubienne} (condition vérifiée si $\mathcal{C}$ est sans projecteurs), alors tout système inductif dans $\mathcal{C}$, grand devant $\pi_0$, où $\pi_0 \geqslant \card\,\text{réd}\,\mathcal{C}$, est essentiellement constant,

\hfill q.e.d.

\printedpage{175}
% head: A. Grothendieck : Dérivateurs, Ch. XVIII, Édition des Universités de Montpellier II et Paris VII

\msp{page 453}
\textbf{4.15.8.} Je veux revenir ici sur la démonstration inachevée de 4.8.13 (p. 212-216). Je reprends la démonstration où je l'avais laissée (p. 215). Il faut prouver que
\[
F(\xi) \ = \ \varinjlim_{i\in I} M_i
\]
est de cardinal $\leqslant \pi_0$. J'ai annoncé qu'on peut se ramener au cas où $M$ est un monoïde de monomorphismes, i.e. que $uv = uv' \Longrightarrow v = v'$ (pour $u,v,v' \in M$). Je vais expliciter cet argument.

Pour tout $i\in I$, soit
\[
N_i \ = \ \left\{ u \in M \ \middle| \ \begin{array}{l}
\text{l'ensemble } I_{i,u} \text{ des } j\in I \text{ tels}\\
j \geqslant i,\ \alpha_{ji} = u \text{ est cofinal}\\
\text{dans } I\text{, i.e. pour tout } j_0\in I,\\
\text{existe } j \in I_{i,u}\text{, avec } j \geqslant j_0
\end{array} \right\} .
\tag*{(4.15.8.1)}
\]
Comme $I$, donc $i\backslash I$ est grande devant $\pi_0$, et $\card M \leqslant \pi_0$, on voit aussitôt que
\[
N_i \ \neq \ \emptyset \ .
\tag*{(4.15.8.2)}
\]
$\Big($ Si d'autre part $i' \geqslant i$, on trouve
\[
N_i \ = \ N_{i'}\alpha_{i'i} \ .
\tag*{(4.15.8.3)}
\]
L'inclusion $N_{i'}\alpha_{i'i} \subset N_i$ est $\pm$ tautologique, prouvons $N_i \subset N_{i'}\alpha_{i'i}$. Soit $u \in N_i$, l'ensemble $J = I_{i,u} \cap (i'\backslash I)$ des $j \geqslant i'$ tels que $\alpha_{ji} = u$ est cofinal dans $I$, donc grand devant $\pi_0$.
\[
\begin{tikzcd}
i \arrow[r] \arrow[rr, bend right=12, "u = \alpha_{ji}"'] & i' \arrow[r, "\alpha_{ji'}"] & j
\end{tikzcd}
\]
Donc il existe un $v \in M$ tel que l'ensemble $J'$ de $j\in J$ tels que $\alpha_{ji'} = v$
\msp{page 454}
soit cofinal dans $J$, donc dans $I$. Comme cet ensemble $J'$ est contenu dans $I_{i',v}$, cela implique que $I_{i',v}$ est cofinal dans $I$, i.e. $v \in N_{i'}$. On a d'autre part, pour $j \in J'$,
\[
\underbrace{\alpha_{ji'}}_{v}\alpha_{i'i} \ = \ \underbrace{\alpha_{ji}}_{u} \ ,
\]
d'où $u = v\alpha_{i'i}$ avec $v\in N_{i'}$, d'où $u \in N_i\alpha_{i'i}$, q.e.d. $\Big)$ % [sic] N_i\alpha_{i'i} -> N_{i'}\alpha_{i'i}, the inclusion to be proved is N_i \subset N_{i'}\alpha_{i'i}

Pour $i \in I$, soit
\[
\alpha_i \ : \ \underbrace{\mathrm{Hom}(\xi, \underbrace{x_i}_{=\,\xi})}_{=\,M} \ \longrightarrow \ \underbrace{\mathrm{Hom}(\xi,F) \ = \ F(\xi)}_{\simeq\ \varinjlim\limits_{i} \mathrm{Hom}(\xi,x_i)}
\]
\printedpage{176}
% head: A. Grothendieck : Dérivateurs, Ch. XVIII, Édition des Universités de Montpellier II et Paris VII
le morphisme canonique. On trouve ceci :
\begin{center}
\makebox[0.20\linewidth][l]{(4.15.8.4)}%
\begin{minipage}[c]{0.78\linewidth}
Si $u \in N_i$, alors la relation d'équivalence $R_i$ dans $M$ définie par $\alpha_i$ ($v \overset{R_i}{\sim} v' \Longleftrightarrow \alpha_i(v) = \alpha_i(v')$) est égale à celle définie par l'application $v \longmapsto uv$ de $M$ dans $M$, i.e.
\[
\alpha_i(v) \ = \ \alpha_i(v') \qquad \Longleftrightarrow \qquad uv \ = \ uv' \ .
\]
\end{minipage}
\end{center}

C'est immédiate, puisque $I_{i,u}$ étant cofinal dans $I$, on a $\alpha_i(v) = \alpha_i(v')$ si et seulement si $\exists\ j \in I_{i,u}$ avec $\alpha_{ji}v = \alpha_{ji}v'$, or $\alpha_{ji} = u$. % [sic] C'est immédiate -> C'est immédiat, agreement

Pour $u \in M$, notons $R_u$ la relation d'équivalence dans $M$ définie par $v \longmapsto uv$. L'ensemble des relations d'équivalence dans \texttt{[}$M$\texttt{]}
\msp{page 455}
obtenu ainsi est de cardinal $\leqslant \card M \leqslant \pi_0$, donc l'ensemble des relations d'équivalences $R_i$ ($i \in I$) est lui aussi de cardinal $\leqslant \pi_0$. Cela montre, $\pi_0$ étant grand devant $\pi_0$, qu'il existe une relation d'équivalence de la forme $R_u$ ($u \in M$), et un ensemble cofinal $I'$ dans $I$, tel que $\forall\ i \in I'$ on ait $R_i = R_u$. Cela montre que, quitte à remplacer $I$ par l'ensemble ordonné cofinal $I'$, qu'on peut supposer que l'on a
\[
R_i \ = \ R_u \qquad \forall\, i \in I \ .
\tag*{(4.15.8.5)}
\]
Soit alors
\[
E \ = \ M/R_u \ = \ M/R_i \qquad (\forall\ i \in I) \ ,
\]
on voit donc que les $\alpha_{ji} \in M$ (pour $i \leqslant j$) sont tels que l'application
\[
\tau_{\alpha_{ji}} \ : \ v \ \longmapsto \ \alpha_{ji}v \ : \ M \ \longrightarrow \ M
\]
est compatible avec la relation d'équivalence $R_u$, donc passe au quotient en une application
\[
\overline{\alpha_{ji}} \ : \ E \ \longmapsto \ E \ ,
\]
laquelle est d'ailleurs \textit{injective.}

Soit $M'$ le sous-monoïde de $M$ formé des $u \in M$ qui ont la propriété de stabilité précédente, ainsi que la propriété d'injectivité qu'on vient de dire, et $\overline{M'}$ le monoïde image de $M'$ dans le monoïde
\[
\mathrm{Mon}(E) \ \subset \ \mathrm{End}_{\mathrm{Ens}}(E)
\]
\msp{page 456}
des applications injectives de $E$ dans lui-même. On voit donc que le système inductif donné, défini par les $\alpha_{ji}$, prend ses valeurs dans le sous-monoïde $M'$ de $M$ (i.e. de la sous-catégorie correspondante de la catégorie $B_M$), et que la limite inductive $F(\xi)$ qui nous intéresse (pour prouver $\card F(\xi) \leqslant \pi_0$) est isomorphe à la limite inductive du système inductif d'un système d'ensembles tous égaux à $E$, les morphismes de transition étant les $\overline{\alpha_{ji}} \in \overline{M'} \subset \mathrm{Mon}(E)$. De plus, on aura $\card \overline{M'} \leqslant \card M \leqslant \pi_0$. Le résultat $\card F(\xi) \leqslant \pi_0$ voulu résultera alors du

\textbf{Lemme-clef 4.15.8.1.} \textit{Soit $M$ un monoïde de monomorphismes (i.e. tel que $uv = uv'$ $\Longrightarrow v = v'$, pour $u,v,v' \in M$), et soit $(\alpha_{ji})_{i\leqslant j}$ un système inductif dans $B_M$, indexé par un ensemble ordonné $I$ grand devant $\pi_0$, où $\pi_0$ est un cardinal}
\printedpage{177}
% head: A. Grothendieck : Dérivateurs, Ch. XVIII, Édition des Universités de Montpellier II et Paris VII
\msp{page 457}
\textit{infini tel que} $\card M \leqslant \pi_0$. \textit{Alors il existe un $i_0 \in I$ tel que pour tout $i \geqslant i_0$, $\alpha_{i_0,i}$ soit un isomorphisme.}

Si on applique ce lemme au cas de $\overline{M'}$ et du $\overline{\alpha_{ji}}$, on trouve donc qu'il existe un $i_0 \in I$ tel que les $\overline{\alpha_{ji_0}}$ soient des isomorphismes. Cela implique que $F(\xi) \simeq \varinjlim\limits_{i} \underbrace{E_i}_{E}$ est isomorphe à $E_{i_0}$, donc est de cardinal égal à $\card E_{i_0} = \card E \leqslant \pi_0$, ce qui achèvera la démonstration.

D'ailleurs, l'énoncé 4.8.13 implique 4.15.8.1, car dans le système inductif des $(M_i = M, \alpha_{ji})$, de $\varinjlim$ égale à $F(\xi)$, s'il n'existait \textit{pas} de $i_0$ comme stipulé dans 4.15.8.1, alors on en conclurait aisément que l'on aurait $\card F(\xi) > \pi_0$, contrairement à 4.8.13. Nous allons dégager ceci dans un lemme :

\msp{page 458}
\textbf{Lemme 4.15.8.2.} \textit{Soit $(E_i)_{i\in I}$ un système inductif d'ensembles, indexé par un ensemble ordonné $I$ grand devant $\pi_0$. On suppose qu'il existe $i_0 \in I$ tel que}
\begin{enumerate}[nosep]
\item[\textit{a)}] \textit{$i \geqslant i_0$ implique} $\card E_i = \pi_0$, \textit{et}
\item[\textit{b)}] \textit{$j \geqslant i \geqslant i_0$ implique que $u_{ji} : E_i \longrightarrow E_j$ est un monomorphisme. Supposons de plus}
\item[\textit{c)}] \textit{$\forall\ i \geqslant i_0$ $\exists\ j > i$ tel que $u_{ji}$ ne soit} pas \textit{un isomorphisme (i.e. ne soit pas surjectif). Soit $E_\infty = \varinjlim\limits_{i\in I} E_i$. Alors} $\card E_\infty > \pi_0$.
\end{enumerate}
Ce lemme implique donc l'\textit{équivalence} de 4.8.13 et du “lemme-clef” 4.15.8.1. Pour prouver 4.15.8.2, nous allons prouver quelque chose d'un peu plus précis :

\textbf{Corollaire 4.15.8.3.} \textit{Sous les conditions de 4.15.8.2, soit $\alpha_0$ le premier ordinal tel que} $\card \alpha_0 > \pi_0$, \textit{et soit $I_{\alpha_0}$ l'ensemble des ordinaux $\beta < \alpha$ (ensemble ordonné qui est grand devant $\pi_0$). Il existe alors une application strictement croissante}
\[
I_{\alpha_0} \ \longrightarrow \ I \ , \qquad \alpha \ \longmapsto \ i_\alpha \ , \qquad \textit{appliquant } 0 \textit{ dans } i_0 \ ,
\]
\msp{page 459}
\textit{tel que pour tout $\alpha \in I_{\alpha_0}$, on ait}
\[
\underbrace{u_{i_{\alpha+1},i_\alpha}}_{\substack{\textit{c'est un } \text{monomorphisme,}\\ \textit{d'après b, vu que } i_{\alpha+1} > i_\alpha \geqslant i_0}} \qquad\qquad \text{\textit{n'est pas un isomorphisme} (\footnotemark[187]).}
\]
\footnotetext[187]{\textit{mieux} : $\forall\ \alpha \in I_{\alpha_0}$ $\varinjlim\limits_{\beta\in I_\alpha} E_{i_\beta} \longrightarrow E_{i_\alpha}$ n'est pas surjectif.}
\textit{On a alors}
\[
E'_\infty \ \overset{\text{déf}}{=} \ \underbrace{\varinjlim_{\alpha\in I_{\alpha_0}} E_{i_\alpha}}_{} \ \hookrightarrow \ \varinjlim_{i\in I} E_i \ \overset{\text{déf}}{=} E_\infty
\]
\textit{et}
\[
\card E'_\alpha \ \geqslant \ \underbrace{\card I_{\alpha_0}}_{=\ \textit{successeur de } \pi_0} \ > \ \pi_0 \ ,
\] % [sic] E'_\alpha -> E'_\infty, the following line has \card E'_\infty
\textit{d'où} $\card E_\infty \geqslant \card E'_\infty > \pi_0$.

\textsc{Démonstration}. La construction des $i_\alpha$ par récurrence transfinie est immédiate, en utilisant c) et le fait que $I$ est grande devant $\pi_0$, et que pour tout $\alpha \in I_{\alpha_0}$, l'ensemble $I_\alpha$
\printedpage{178}
% head: A. Grothendieck : Dérivateurs, Ch. XVIII, Édition des Universités de Montpellier II et Paris VII
des $\beta < \alpha$ est de cardinal $\leqslant \pi_0$. Le fait que $E'_\infty \longrightarrow E_\infty$ soit injectif résulte de b). Le fait que
\[
\card E'_\infty \ \geqslant \ \card I_{\alpha_0}
\]
résulte du fait que $E'_\infty$ est réunion croissante de sous-ensembles identifiés à $E_{i_\alpha}$, et pour tout
\msp{page 460}
$\alpha$ on peut trouver un
\[
\xi_\alpha \ \in \ E_{i_\alpha}\backslash \bigcup_{\beta<\alpha} E_{i_\beta} \ .
\]
L'application
\[
\begin{aligned}
\alpha \ &\longmapsto \ \xi_\alpha\\
I_{\alpha_0} \ &\longrightarrow \ E'_\infty
\end{aligned}
\]
est injective, donc $\card E'_\infty \geqslant \card I_{\alpha_0}$.

% vertical bar printed in the left margin beside the whole of Corollaire 4.15.8.4
\textbf{Corollaire 4.15.8.4.} \textit{Pour un $M$ et un $\pi_0$ donnés, pour que le lemme-clef 4.15.8.1 soit valable, il faut et il suffit qu'il le soit dans le cas où $I = I_{\alpha_0}$, $\alpha_0$ étant le premier ordinal tel que} $\card \alpha_0 > \alpha_0$. \textit{Il revient au même qu'on ne puisse trouver de système inductif $(\alpha_{ji})_{j\geqslant i}$ indexé par $I_{\alpha_0}$, à valeurs dans $M$, tel que pour $i < j$, on ait $\alpha_{ji}$ non inversible dans $M$.} % [sic] \card \alpha_0 > \alpha_0 -> \card \alpha_0 > \pi_0, cf. Corollaire 4.15.8.3

\msp{page 461}
Je vais résumer ce qui a été obtenu dans un énoncé récapitulatif :

\textbf{Théorème 4.15.9.} \textit{Soient $\mathcal{C}$ une petite catégorie, $\pi$ un cardinal tel que $\pi \geqslant$} $\card\mathcal{C}$. \textit{Les conditions suivantes sont équivalentes (et je conjecture qu'elles sont toujours satisfaites) :}
\begin{enumerate}[nosep]
\item[\textit{a)}] \textit{$\mathcal{M} \overset{\text{déf}}{=} \mathrm{Ind}_\pi\mathcal{C}$ est réduit à $\mathcal{M}_\pi$.}
\item[\textit{a$'$)}] \textit{(Si $\mathcal{C}$ est karoubienne.) Tout système inductif $\underline{x} = (x_i)_{i\in I}$ dans $\mathcal{C}$, grand devant $\pi$, est essentiellement constant.}
\item[\textit{a$''$)}] \textit{Pour tout système inductif $(x_i)_{i\in I}$ dans $\mathcal{C}$, grand devant $\pi$, il existe un foncteur cofinal $J \longrightarrow I$, avec $J$ ensemble ordonné grand devant $\pi$, et tel que le système réindexé $(x_j)_{j\in J}$ satisfasse les conditions suivantes :}
\begin{enumerate}[nosep]
\item[\textit{1)}] \textit{Tous les $x_j$ sont égaux à un même élément $x$ de $\mathcal{C}$.}
\item[\textit{2)}] \textit{Il existe un projecteur $p$ dans $x$, tel que l'on ait $\alpha_{j'j} = p$ pour $j < j'$ ($j,j' \in J$), où $\alpha_{j'j} : x_j \longrightarrow x_{j'}$ est la flèche de transition.}
\end{enumerate}
\item[\textit{b)}] \textit{Pour tout $x \in \mathrm{Ob}\,\mathcal{C}$, et pour tout monoïde quotient $M'$ de $M = \mathrm{End}_{\mathcal{C}}(x)$, tel que $M'$ soit un monoïde de monomorphismes (i.e. $uv = uv' \Longrightarrow v = v'$ si $u,v,v' \in M'$), et pour tout système inductif grand devant $\pi$ à valeurs dans $B_{M'}$ (définie par un système d'éléments $(\alpha_{ji})_{i\leqslant j,\ i,j\in I}$ d'un ensemble ordonné $I$ grand devant $\pi$, avec $\alpha_{kj}\alpha_{ji} = \alpha_{ki}$ pour $i \leqslant j \leqslant k$), il existe $i_0 \in I$ tel}
\msp{page 462}
\textit{que $\alpha_{ii_0} = 1$ pour $i \geqslant i_0$, i.e. $\alpha_{ji} = 1$ pour $j \geqslant i \geqslant i_0$.}
\item[\textit{b$'$)}] \textit{$M'$ étant un monoïde de monomorphismes comme dans b) (donc} $\card M' \leqslant \card\mathcal{C} \leqslant$ \textit{$\pi$), soit de plus $\alpha_0$ le plus petit ordinal tel que} $\card\alpha_0 > \pi$, \textit{et soit $I = I_{\alpha_0}$ l'ensemble ordonné des ordinaux $\alpha < \alpha_0$ (lequel ensemble est grand devant $\pi$). Il n'existe} pas \textit{de système inductif dans $B_{M'}$, tel que pour $i < j$ dans $I$, $\alpha_{ji}$ soit un élément non inversible dans $M'$ (donc $u \longmapsto \alpha_{ji}u : M' \longrightarrow M'$ est injectif mais non surjectif).}
\end{enumerate}


\endgroup
% END UNIT der-ch18/der-ch18_p176-180.tex
% BEGIN UNIT der-ch18/der-ch18_p181-185.tex
\begingroup\makeatletter%
\newcommand{\printedpage}[1]{\marginpar{\footnotesize p.~#1}}%
\newcommand{\msp}[1]{\marginpar{\footnotesize [#1]}}%
\newcommand{\card}{\operatorname{card}}%
\newcommand{\redop}{\mathop{\textnormal{réd}}\nolimits}%
\newcommand{\Hom}{\operatorname{Hom}}%
\newcommand{\Ob}{\operatorname{Ob}}%
\newcommand{\Ind}{\operatorname{Ind}}%
\newcommand{\Ens}{\mathrm{Ens}}%
\newcommand{\defeq}{\overset{\text{déf}}{=}}%
\newcommand{\into}{\lhook\joinrel\longrightarrow}%
\makeatother\let\printedpage\derpp\let\msp\dermsp%
\let\mspin\DERmspin%


\printedpage{179}
% head: A. Grothendieck : Dérivateurs, Ch. XVIII, Édition des Universités de Montpellier II et Paris VII
% foot: Version 17 décembre 2025

Signalons tout de suite le

\textbf{Corollaire 4.15.10.} \textit{Les conditions équivalentes précédentes sont satisfaites, si $\pi > \card\mathcal{C}$.}

En effet, soit $\alpha_1$ le plus petit ordinal tel que $\card\alpha_1 > \card\mathcal{C}$. Donc $\pi_1 \defeq \card\alpha_1$ est le successeur de $\pi_0 \defeq \card\mathcal{C}$, et on a $\pi \geqslant \pi_1$, donc (si $\alpha_0$ est comme explicité dans b$'$ ci-dessus) $\card\alpha_0 > \pi_1$, donc $\alpha_1 < \alpha_0$. Soit (avec les notations de b$'$), pour $\alpha \in I$ \texttt{[plutôt} $\alpha \in I_{\alpha_0}$\texttt{]}, i.e.\ $\alpha < \alpha_0$,
\[
M_\alpha \;=\; \varinjlim_{i<\alpha} \underbrace{M_i}_{=\,M}\ ,
\]
on voit donc, comme les $\alpha_{ji}$ sont injectifs et non surjectifs, que
\[
\card M_\alpha \;\geqslant\; \card\alpha\ ,
\]
\msp{463}
et que pour $\alpha < \beta$, on a
\[
M_\alpha \;\into\; M_\beta \qquad \text{injectif}\ .
\]
Appliquons ceci au cas où $\alpha = \alpha_1$, donc on trouve
\[
\card M_{\alpha_1} \;\geqslant\; \card\alpha_1 \;=\; \pi_1\ ,
\]
d'autre part, en prenant $\beta = \alpha_1 + 1$, d'où $M_\beta \simeq M$, on trouve une injection
\[
M_{\alpha_1} \;\into\; M\ ,
\]
d'où
\[
\pi_1 \;=\; \card M_{\alpha_1} \;\leqslant\; \card M \;=\; \pi_0\ ,
\]
une contradiction, ce qui établit b$'$.

\textbf{Corollaire 4.15.11.} \textit{Soit $\mathcal{C}$ une catégorie essentiellement petite, $\pi$ un cardinal tel que $\pi > \card\redop\mathcal{C}$. Soit $\mathcal{M} \defeq \Ind_\pi\mathcal{C}$, alors $\mathcal{M} = \mathcal{M}_\pi$, et c'est une enveloppe de Karoubi de $\mathcal{C}$.}

On peut supposer $\mathcal{C}$ réduite, donc $\card\mathcal{C} < \pi$, et on applique le corollaire précédent.

Ainsi, le résultat litigieux 4.8.13 (p.~212) est finalement prouvé sous une forme légèrement plus faible, avec $\pi > \card\redop\mathcal{C}$ au lieu de $\pi \geqslant \card\redop\mathcal{C}$. Du même coup, les applications de 4.8.13 qu'on a fait jusqu'à présent (notamment le lemme 4.15.3).

\msp{464}

\subsection*{4.16\quad Exemples et questions}

\textbf{4.16.0.} J'avais cru que tout foncteur “raisonnable” (ou plus précisément, tout foncteur “logiquement accessible”) entre catégories accessibles devrait être accessible. Ce n'est malheureusement pas vrai, puisqu'un composé de deux foncteurs Hom
\[
\Ens \longrightarrow \Ens\ , \qquad\quad E \longmapsto \underbrace{\Hom(\Hom(E,I),J)}_{J^{(I^E)}}
\]
\printedpage{180}
% head: A. Grothendieck : Dérivateurs, Ch. XVIII, Édition des Universités de Montpellier II et Paris VII
% foot: Version 17 décembre 2025
est non accessible, si $I$, $J$ sont tous deux de cardinal $\geqslant 2$. Prenant p.ex.\ $I = J = \{0,1\}$, on trouve le foncteur
\[
E \longmapsto \mathfrak{P}(\mathfrak{P}(E))\ ,
\]
qui n'est pas accessible.

\textbf{4.16.1.} Deuxième désillusion : j'avais cru que pour toute catégories accessible $\mathcal{M}$, toute % [sic] toute catégories -> toute catégorie, singular required after "toute"
sous-catégorie $\mathcal{N}$ strictement pleine et stable par facteurs directs “raisonnable” (p.ex. “logiquement accessible”) devrait être accessible dans $\mathcal{M}$. Il n'en est rien, même dans le cas particulier de l'image inverse de la sous-catégorie accessible dans Ens
\msp{465}
\[
\underbrace{\Ens^*}_{=\,\{E\,\in\Ens\,\mid\, E\neq\emptyset\}} \quad\subset\ \Ens
\]
par un contrafoncteur accessible \textit{représentable}
\[
\begin{aligned}
\mathcal{M} \quad &\longrightarrow \quad (\Ens)^\circ \\
x \quad &\longmapsto \quad \Hom_{\mathcal{M}}(x,z)\ ,
\end{aligned}
\]
i.e.\ la sous-catégorie ouverte (donc strictement pleine) $\mathcal{M}_z$ de $\mathcal{M}$ formée des objets $x$ tels que $\Hom_{\mathcal{M}}(x,z) \neq \emptyset$, n'est par toujours accessible dans $\mathcal{M}$. Ici, je ne suis pas parvenu à % [sic] n'est par toujours -> n'est pas toujours, misprint "par" for "pas"
un exemple aussi élémentaire (cf.\ appendice) -- j'ai pris $\mathcal{M} = (k\text{-alg.\ comm.})$, $k$ un corps infini et $z$ l'objet initial $k$ de $\mathcal{M}$ -- j'ai trouvé que $\mathcal{M}_z$ n'est pas $\mathrm{L}_\pi$-stable dans $\mathcal{M}$, quelque soit le cardinal $\pi$ % [sic] quelque soit -> quel que soit, standard spelling of the locution
-- mais cet exemple utilisait que $\mathfrak{U} = \mathfrak{U}_1$, plus petit univers contenant $\mathbf{N}$. J'ignore si pour tout autre univers, on peut trouver un corps $k$ de cardinal $\pi \in \mathfrak{U}$ assez grand, de façon que tout $\pi$-ultrafiltre sur un ensemble $E \in \mathfrak{U}$ soit trivial. Peut-être
\msp{466}
un tel exemple, où non la construction du couple $(\mathcal{M},\mathcal{N})$, mais la démonstration du fait que $\mathcal{N}$ n'est pas accessible dans $\mathcal{M}$, fait appel à des hypothèses particulières sur l'univers ambiant, ne doit pas être considéré comme (tout à fait !) “logiquement accessible”. Pour être vraiment probant, il faudrait un exemple de nature plus élémentaire.

% margin: a vertical bar is printed in the left margin beside the three lines from "L_pi-stable dans M pour pi assez grand" to "logiquement accessible) serait accessible dans M"
Il reste la possibilité qu'il y ait un théorème métamathématique qui assure que tout couple $(\mathcal{M},\mathcal{N})$ logiquement accessible, avec $\mathcal{M}$ accessible, et $\mathcal{N}$ sous-catégorie strictement pleine $\mathrm{L}_\pi$-stable dans $\mathcal{M}$ pour $\pi$ assez grand, soit accessible dans $\mathcal{M}$ (\footnote[188]{C'est faux, même pour une sous-catégorie fermée, cf.\ 4.16.7.10 ci-dessous.}). Ceci s'appliquerait en particulier au cas d'une sous-catégorie \textit{fermée} : toute sous-catégorie fermée (qui soit logiquement accessible) serait accessible dans $\mathcal{M}$. L'exemple-test le plus simple serait celui du complémentaire fermé $\mathcal{M}'_z$ de la sous-catégorie ouverte $\mathcal{M}_z$ ci-dessus (\footnote[189]{\textbf{N.B.} elles est accessible dans $\mathcal{M}$ si et seulement si $\mathcal{M}_z$ l'est ! % [sic] elles est -> elle est, agreement with singular verb
}), i.e.\ la catégorie des $x \in \Ob\mathcal{M}$ tels que $\Hom_{\mathcal{M}}(x,z) = \emptyset$, ou encore
\msp{467}
l'image inverse de la sous-catégorie fermée accessible
\[
e \;=\; \{\emptyset\} \;\subset\; (\Ens)
\]
réduite à $\{\emptyset\}$, par le foncteur représentable $x \longmapsto \Hom(x,z) : \mathcal{M} \longrightarrow \Ens$. Il se pourrait donc qu'une telle catégorie soit toujours accessible dans $\mathcal{M}$. Mais j'ai été incapable de le
\printedpage{181}
% head: A. Grothendieck : Dérivateurs, Ch. XVIII, Édition des Universités de Montpellier II et Paris VII
% foot: Version 17 décembre 2025
prouver, même dans le cas particulier ci-dessus, où $\mathcal{M} = (k\text{-alg.\ comm.})$, $k$ étant un corps commutatif infini. La difficulté est (dans tous les cas) la suivante : trouver, pour $z$ fixé, un cardinal $\pi > \pi_0$ tel que pour tout $x \in \Ob\mathcal{M}$, et tout système (ou \textit{un} système) $\pi_0$-$\pi$-adapté $(x_i)_{i\in I}$ de limite inductive $x_0$ ($\pi_0$ un cardinal donné, utile pour $\mathcal{M}$), \textit{si} $\Hom_{\mathcal{M}}(x,z)$ ($\simeq \varprojlim \Hom_{\mathcal{M}}(x_i,z)$) est vide, alors $\exists\ i_0 \in I$ tel que $\Hom_{\mathcal{M}}(x_{i_0},z)$ soit déjà vide. Dans l'exemple $\mathcal{M} = (k\text{-alg.\ comm.})$, il faut donc prouver que l'on peut trouver un cardinal $\pi$
tel que pour tout $k$-algèbre $A$ avec $\Hom_{\mathcal{M}}(A,k) = \emptyset$, il existe une sous-$k$-algèbre $A_i$ de % [sic] pour tout k-algèbre -> pour toute k-algèbre, "algèbre" is feminine
cardinal $\leqslant \pi$, telle que $\Hom_k(A_i,k) = \emptyset$ (\footnote[190]{C'est faux, cf.\ annotation marginale p.~466.}). J'ai tendance à croire que ça doit être vrai (et même sans
\msp{468}
supposer que $k$ soit un corps, et aussi en prenant la condition $\Hom_{\mathcal{M}}(A,Z) = \emptyset$, où $Z$ est un objet quelconque de $\mathcal{M}$). Ce genre de questions semble étonnement délicat \dots % [sic] étonnement -> étonnamment, adverb required

\textbf{4.16.2.} Il est possible, plus généralement, que la question suivante ait une réponse affirmative : Soit
\[
f \;:\; \mathcal{M} \longrightarrow \mathcal{N}^\circ
\]
un foncteur accessible, avec $\mathcal{M}$, $\mathcal{N}$ catégories accessibles, et soit $\underline{T}$ une sous-catégorie de $\mathcal{N}$ accessible dans $\mathcal{N}$, et telle qu'il existe un cardinal $\pi$ ayant les propriétés suivantes : $\mathcal{N}^\circ$ satisfait $\mathrm{L}_\pi$ (i.e.\ $\mathcal{N}$ stable par $\varprojlim$ grandes devant $\pi_1$). Alors $\underline{S} \defeq f^{-1}(\underline{T})$ est-elle accessible dans $\mathcal{M}$ (\footnote[191]{Faux, avec $\mathcal{N} = \Ens$ et un foncteur représentable $\mathcal{M} \longrightarrow \Ens^\circ$, prenant $I = \{\emptyset\}$, cf.\ annotations marginales p.~466, 467.}) ? Ceci généralise le cas $\mathcal{N} = (\Ens)$, $\underline{T} = \{\emptyset\}$, $f$ foncteur représenté par $z \in \Ob\mathcal{M}$, envisagé tout à l'heure. Avec le même type d'exemple, mais en prenant $\mathcal{N} = (\Ens)^*$
\msp{469}
au lieu de $\mathcal{N} = \{\emptyset\}$, on voit que la condition supplémentaire de stabilité sur $\underline{T}$ n'est pas
superflue. (\textbf{N.B.} Elle sont à assurer que $\underline{S}$ est $\mathrm{L}_\pi$-stable dans $\mathcal{M}$ pour $\pi$ grand.) % [sic] Elle sont -> Elles sont, agreement with plural verb

\textbf{4.16.3.} Questions apparentée : j'avais espéré que pour un foncteur accessible entre catégories accessibles % [sic] Questions apparentée -> Question apparentée, agreement
\[
f \;:\; \mathcal{N} \longrightarrow \mathcal{M}\ ,
\]
la $\natural$-image essentielle, formée des $y \in \Ob\mathcal{M}$ qui sont isomorphes à un facteur direct d'un objet de la forme $f(x)$, avec $x \in \Ob\mathcal{N}$, serait accessible dans $\mathcal{N}$. L'exemple ci-dessus (nullement élémentaire !) de la sous-catégorie $\mathcal{M}_z$ de $\mathcal{M}$ (formée des $x \in \Ob\mathcal{M}$ tels que $\Hom_{\mathcal{M}}(x,z) \neq \emptyset$) montre qu'il n'en est rien. En effet, $\mathcal{M}_z$ est l'image essentielle, et aussi l'image $\natural$-essentielle, du foncteur
\[
\mathcal{M}/z \longrightarrow \mathcal{M}\ ,
\]
or $\mathcal{M}/z$ est accessible, et le foncteur précédent accessible, si $\mathcal{M}$ est accessible (vérification par AQT \texttt{[âne qui trotte]}, cf.\ plus bas, énoncé plus général).

\msp{470}

On a l'impression que la raison pour laquelle on a de tels contre-exemples, c'est que l'on a des $\natural$-images essentielles qui ne sont $\mathrm{L}_\pi$-stables dans $\mathcal{M}$ pour aucun cardinal $\pi \in \mathfrak{U}$. Cela amène à modifier la conjecture ainsi : soit $\pi$ un cardinal tel que $\mathcal{M}$ satisfasse à $\mathrm{L}_\pi$,
\printedpage{182}
% head: A. Grothendieck : Dérivateurs, Ch. XVIII, Édition des Universités de Montpellier II et Paris VII
% foot: Version 17 décembre 2025
et soit $\underline{S}(\pi)$ la plus petite sous-catégorie strictement pleine de $\mathcal{M}$, $\mathrm{L}_\pi$-stable dans $\mathcal{M}$, et qui contienne l'image essentielle de $f$. (On pourrait appeler $\underline{S}(\pi)$ la $\mathrm{L}_\pi$-image essentielle.) Alors $\underline{S}(\pi)$ est accessible dans $\mathcal{M}$. De plus, la famille (décroissante) des $\underline{S}(\pi)$, pour $\pi$ variable, est stationnaire, i.e.\ il existe un cardinal $\pi$ tel que $\pi' \geqslant \pi$ implique $\underline{S}(\pi') = \underline{S}(\pi)$. Il est clair alors que $\underline{S}(\pi)$ est la plus petite sous-catégorie de $\mathcal{M}$ qui soit accessible dans $\mathcal{M}$ et qui contienne l'image essentielle de $f$. La même question se pose pour un foncteur $g^\circ : \mathcal{N}^\circ \longrightarrow \mathcal{M}$, où $g : \mathcal{N} \longrightarrow \mathcal{M}^\circ$ est un foncteur accessible, $\mathcal{N}$ et $\mathcal{M}$ accessibles. Je conjecture que la réponse est encore affirmative.

\msp{471}

% margin: a hand-drawn curved stroke is printed in the left margin beside the last three lines of the following paragraph ("pleine M de C^ qui soit ..." to "sans pourtant être accessible.")
Soit $\mathcal{M}$ une catégorie satisfaisant la condition L, et admettant une petite sous-catégorie pleine strictement génératrice $\mathcal{C}$, formée d'objets accessibles. On a donc un foncteur pleinement fidèle et accessible
\[
\mathcal{M} \;\into\; \mathcal{C}^\wedge\ ,
\]
$\mathcal{M}$ $\mathrm{L}_\pi$-stable dans $\mathcal{C}^\wedge$ pour $\pi$ convenable. Si $\overline{\mathcal{M}}$ est l'image essentielle, de sorte que $\mathcal{M} \simeq \overline{\mathcal{M}}$, c'est donc une sous-catégorie strictement pleine $\mathrm{L}_\pi$-stable dans $\mathcal{C}^\wedge$. Donc on s'attend à ce qu'elle soit accessible dans $\mathcal{C}^\wedge$ (si la situation décrite est “logiquement accessible”), ou ce qui revient au même, que $\mathcal{M}$ soit accessible. (Comparer 4.8.1, p.~187.) Mais il ne doit pas être difficile (\footnote[192]{OK en prenant un \textit{fermé} $\underline{S}$ non accessible dans $\mathcal{M}$ d'une catégorie accessible $\mathcal{M}$, qu'on plonge dans un $\mathcal{C}^\wedge$. Mais ces $\underline{S}$ n'a \textit{pas} une petite famille strictement génératrice ! % [sic] ces S n'a -> ce S n'a, plural determiner with singular verb
}), en procédant à un nombre $\pi_{\mathfrak{U}}$ de choix (où $\pi_{\mathfrak{U}} = \card\mathfrak{U}$), de faire quand même des exemples (“logiquement inaccessibles”) d'une sous-catégorie strictement pleine $\mathcal{M}$ de $\mathcal{C}^\wedge$ qui soit $\mathrm{L}_\pi$-stable, mais non accessible dans $\mathcal{C}^\wedge$. Cette $\mathcal{M}$ admet donc une petite sous-catégorie strictement génératrice formée d'objets accessibles (\footnote[193]{Ce n'est pas clair, il faudrait pour cela que le foncteur d'inclusion soit un adjoint à gauche !}) et satisfait à $\mathrm{L}_\pi$, sans pourtant être accessible.

\msp{472}

\textbf{4.16.4} (\footnote[194]{descendre après 4.16.5.}). Soit $\mathcal{M}$ une catégorie telle que $\mathcal{M}$ et $\mathcal{M}^\circ$ soient accessibles. Je conjecture qu'alors $\mathcal{M}$ est essentiellement petite. Plus généralement, supposons $\mathcal{M}$ accessible, et que $\mathcal{M}$ admette une petite sous-catégorie qui soit \textit{strictement co}génératrice. Alors $\mathcal{M}$ est-elle essentiellement petite ? Plus généralement, si $\mathcal{M}$ admet une petite sous-catégorie strictement génératrice, et aussi une petite sous-catégorie strictement cogénératrice, i.e.\ si $\mathcal{M}$ et $\mathcal{M}^\circ$ sont toutes deux équivalentes à des sous-catégories pleines d'une $\mathcal{C}^\wedge$, avec $\mathcal{C}$ petite, alors $\mathcal{M}$ est-elle essentiellement petite ?

\textbf{4.16.5.} Soit $I$ une petite catégorie, $I_\bullet$ le cône droit sur $I$ (obtenu en rajoutant un objet final $e$), $\mathcal{M}$ une catégorie accessible,
\[
\underline{\Hom}_{\square}(I_\bullet,\mathcal{M}) \;\into\; \underline{\Hom}(I_\bullet,\mathcal{M})
\]
la sous-catégorie strictement pleine de la catégorie (accessible) $\underline{\Hom}(I_\bullet,\mathcal{M})$, formée des foncteurs $f$ qui font de $f(e)$ la limite inductive dans $\mathcal{M}$ des $f(i)$ ($i \in \Ob I$). Cette sous-catégorie est $\mathrm{L}_\pi$-stable pour $\pi$ grand, donc il y a lieu
\msp{473}
de conjecturer qu'elle est accessible dans $\underline{\Hom}(I_\bullet,\mathcal{M})$. $\Big($ Je rappelle qu'on a vu dans 4.14 que c'est OK si les $\varinjlim$ de type $I$ sont représentables dans $\mathcal{M}$. Et qu'une réponse affirmative, dans le cas où
\printedpage{183}
% head: A. Grothendieck : Dérivateurs, Ch. XVIII, Édition des Universités de Montpellier II et Paris VII
% foot: Version 17 décembre 2025
\[
\begin{tikzcd}[row sep=small, column sep=small]
 & 1 \\
I = 0 \arrow[ur] \arrow[dr] & \\
 & 2
\end{tikzcd}
\quad ,\qquad\quad \text{donc} \qquad\quad
\begin{tikzcd}[row sep=small, column sep=small]
 & 1 \arrow[dr] & \\
I = 0 \arrow[ur] \arrow[dr] & & 3 \\
 & 2 \arrow[ur] &
\end{tikzcd}
\quad ,
\]
implique que la sous-catégorie $\underline{\text{Ép}}(\mathcal{M})$ dans $\underline{\mathrm{Fl}}(\mathcal{M})$ est accessible dans $\underline{\mathrm{Fl}}(\mathcal{M})$.$\Big)$

\textbf{N.B.} Si $\mathcal{M}$ admet une petite sous-catégorie \textit{co}génératrice, alors la question est justiciable de la conjecture 4.16.2 relative à un foncteur accessible $\mathcal{M} \longrightarrow \mathcal{N}^\circ$, ici $\underline{\Hom}(I_\bullet,\mathcal{M}) \longrightarrow \underline{\Hom}(I^\circ \times \mathcal{C},\Ens)^\circ \dots$

\textbf{4.16.6.} Il reste la question de l'accessibilité relative d'une catégorie $\mathcal{E}$ sur une petite catégorie $\mathcal{C}$. Je pense à présent, grâce à 4.15.9, et son corollaire 4.15.10, que l'énoncé suivant doit “tomber” par AQT \texttt{[âne qui trotte]} :

\textbf{Proposition 4.16.6.1.} \textit{Soit $f : \mathcal{E} \longrightarrow \mathcal{C}$ un foncteur, avec $\mathcal{C}$ petite catégorie karoubienne. On suppose que $f$ satisfait la condition $H' = H$ (1+2+3) de 4.15.1 (page 415), où il suffit H 2, H 3, pour toute flèche $\lambda$} non identique \textit{de $\mathcal{C}$. On suppose de plus que pour une catégorie $\mathcal{C}'$ d'un des deux types qu'on va décrire, et un foncteur $\mathcal{C}' \longrightarrow \mathcal{C}$, $\mathcal{E}' = \mathcal{E} \times_{\mathcal{C}} \mathcal{C}'$ soit accessible, et $\mathcal{E}' \longrightarrow \mathcal{C}'$ accessible.}

\msp{474}

$\Big($ \textit{On prend pour $\mathcal{C}'$ la catégorie}
\[
\mathcal{Q} = \left\{ \quad
\begin{tikzcd}[column sep=2.2em]
\xi_0 = \xi^\natural \arrow[r, shift left, "\sigma"] \arrow[dr, "u_0"'] & \xi \arrow[l, shift left, "\rho"] \arrow[d, "u = u_0\rho"] \arrow[loop right, "p"] \\
 & \eta
\end{tikzcd}
\qquad
\begin{array}{rclll}
\rho\sigma & = & \mathrm{id}_{\xi_0} & & \\[2ex]
\sigma\rho & = & p\ , & \text{\textit{d'où}} & \left\{\begin{array}{rcl} p^2 & = & p \\ p\sigma & = & \sigma \\ \rho p & = & \rho \end{array}\right. \\[4ex]
u & = & u_0\rho\ , & \text{\textit{d'où}} & \left\{\begin{array}{rcl} up & = & u \\ u\sigma & = & u_0 \end{array}\right.
\end{array}
\right.
\]
\textit{(trois objets $\xi$, $\xi_0$, $\eta$ et cinq flèches non identiques $\sigma$, $\rho$, $p$, $u_0$, $u$, liés par les relations ci-dessus), ou la catégorie (isomorphe à $Q^\circ$)} % Doubt resolved against PDF leaf 25: Q is ordinary mathematical italic, retained.
\[
\mathcal{Q}' = \left\{ \quad
\begin{tikzcd}
\xi_0 \arrow[r, shift left, "\sigma"] & \xi \arrow[l, shift left, "\rho"] \arrow[loop right, "p"] \\
 & \eta \arrow[ul, "u_0"] \arrow[u, "u"]
\end{tikzcd}
\qquad
\begin{array}{lcl}
\left.\begin{array}{rcl} \rho\sigma & = & \mathrm{id}_{\xi_0} \\ \sigma\rho & = & p \end{array}\right\} & \Longrightarrow & \left\{\begin{array}{rcl} p^2 & = & p \\ p\sigma & = & \sigma \\ \rho p & = & \rho \end{array}\right. \\[4ex]
\begin{array}{rcl} u & = & \sigma u_0 \end{array} & \Longrightarrow & \left\{\begin{array}{rcl} pu & = & u \\ \rho u & = & u_0 \end{array}\right.
\end{array}
\right.
\]
\textit{On se limite aux foncteurs $\mathcal{Q} \longrightarrow \mathcal{C}$, $\mathcal{Q}' \longrightarrow \mathcal{C}$ tels que $p$ soit transformé en une flèche} non identique \textit{de $\mathcal{C}$ (\footnote[195]{Je ne suis pas sûr si c'est raisonnable -- à vérifier avec soin !}) (donc la condition relative à $\mathcal{Q}$, $\mathcal{Q}'$ devient vide, si $\mathcal{C}$ est sans projecteurs).} $\Big)$

\textit{Sous ces conditions, $\mathcal{E}$ est accessible, et $f : \mathcal{E} \longrightarrow \mathcal{C}$ est accessible.}


\endgroup
% END UNIT der-ch18/der-ch18_p181-185.tex
% BEGIN UNIT der-ch18/der-ch18_p186-190.tex
\begingroup\makeatletter%
\newcommand{\printedpage}[1]{\marginpar{\footnotesize p.~#1}}%
\newcommand{\msp}[1]{\marginpar{\footnotesize [#1]}}%
\newcommand{\Ob}{\operatorname{Ob}}%
\newcommand{\Hom}{\operatorname{Hom}}%
\newcommand{\Ind}{\operatorname{Ind}}%
\newcommand{\card}{\operatorname{card}}%
\newcommand{\cM}{\mathcal{M}}%
\newcommand{\cE}{\mathcal{E}}%
\newcommand{\cC}{\mathcal{C}}%
\newcommand{\cQ}{\mathcal{Q}}%
\newcommand{\uS}{\underline{S}}%
\newcommand{\uT}{\underline{T}}%
\newcommand{\U}{\mathfrak{U}}%
\makeatother\let\printedpage\derpp\let\msp\dermsp%
\let\mspin\DERmspin%

\setcounter{footnote}{195}

\printedpage{184}
% head: A. Grothendieck : Dérivateurs, Ch. XVIII, Édition des Universités de Montpellier II et Paris VII
% foot: Version 17 décembre 2025

\msp{page 475}
On trouve alors le

\begingroup\itshape
\noindent\textbf{\upshape Corollaire 4.16.6.2.} Soit $f : \cE \longrightarrow \cC$ un foncteur. Conditions équivalentes :
\begin{enumerate}[label={}]
\item[a)] Pour toute catégorie $\cC'$ \textup{petite} et \textup{karoubienne}, et tout foncteur $\cC' \longrightarrow \cC$, $\cE' = \cE \times_{\cC} \cC'$ est accessible, et le foncteur $\cE' \longrightarrow \cC'$ est accessible.
\item[b)] Comme a), mais $\cC'$ étant une des trois catégories $\Delta_1$, $\cQ$, $\cQ'$ (où $\cQ$, $\cQ'$ sont définies dans 4.15.6.1 ci-dessus).
\item[b$'$)] \textup{(}\footnote{?, à vérifier avec soin !}\textup{)}. (Si $\cC$ est sans projecteurs.) Comme b), mais en se limitant au cas $\cC' = \boldsymbol{\Delta}_1$ (à l'exclusion de $\cQ$, $\cQ'$).
\end{enumerate}
\endgroup

\begingroup\itshape
\noindent\textbf{\upshape Corollaire 4.16.6.3.} Supposons $f : \cE \longrightarrow \cC$ fibrant ou cofibrant. Alors les conditions équivalentes précédentes impliquent que :
\begin{enumerate}[label={}]
\item[1$^\circ$)] Les $\cE_s$ sont accessibles ($s \in \Ob \cC$).
\item[2$^\circ$)] Pour toute flèche $\lambda : s \longrightarrow t$ dans $\cC$, le foncteur $\lambda^* : \cE_t \longrightarrow \cE_s$ (resp. $\cE_s \longrightarrow \cE_t$) est accessible.
\end{enumerate}
\noindent D'autre part, l'ensemble de ces conditions 1$^\circ$, 2$^\circ$) équivaut à la validité de b) ci-dessus, dans le cas $\cC' = \boldsymbol{\Delta}_1$, ou encore à la validité
\msp{page 476}
de a) ci-dessus, quand $\cC'$ est \textup{sans projecteurs}.
\endgroup

\textbf{Question 4.16.6.4.} \textit{Les conditions 1$^\circ$ et 2$^\circ$ précédentes impliquent-elles celles de 4.16.6.2 ? En d'autres termes, si la catégorie base $\cC$ est égale à $\cQ$ ou à $\cQ'$, alors ces conditions impliquent-elles que $\cE$ est accessible, et $f : \cE \longrightarrow \cC$ accessible ?}

Ça semble idiot de ne pas connaître la réponse !

\textbf{4.16.6.5.} J'ai envie d'appeler catégorie \textit{accessible sur $\cC$}, une catégorie $\cE$ sur $\cC$ satisfaisant les conditions équivalentes de 4.16.6.1 -- du moins dans le cas où on suppose $\cC$ petite. Pour que cette terminologie soit vraiment satisfaisant, il faudrait que l'on sache prouver que ces conditions impliquent ceci : pour toute catégorie accessible $\cM$ et tout foncteur accessible $\cM \longrightarrow \cC$, la catégorie $\cE \times_{\cC} \cM$ est accessible, et le foncteur $\cE \times_{\cC} \cM \longrightarrow \cM$ est accessible. % [sic] satisfaisant -> satisfaisante, accord avec « terminologie »

\msp{page 476bis}
L'idée, c'est qu'on appellerait (sous hypothèse de petitesse) \textit{catégorie accessible $\cE$} sur une autre $\cC$, une catégorie ayant la propriété qu'on vient de dire -- mais en supposant maintenant de plus que $\cC$ satisfasse à $\mathrm{L}_\pi$, j'imagine, pour tout foncteur accessible $\cM \longrightarrow \cC$, avec $\cM$ accessible, $\cE \times_{\cC} \cM$ est accessible et le foncteur canonique $\cE \times_{\cC} \cM \longrightarrow \cE$ est accessible.

\msp{page 477}
\subsection*{4.16.7. Retour sur les sous-catégories fermées d'une catégorie accessible (4.16.1).}

\textbf{Proposition 4.16.7.1.} \textit{Soient $\cM$ une catégorie accessible, $(\uS_i)_{i\in I}$ une famille de sous-catégories fermées accessibles dans $\cM$. On ne suppose que \textup{card}\,$I$ soit petit. Soit $\uS$ la sous-catégorie strictement pleine de $\cM$ telle que $\Ob \uS = \bigcup_{i\in I} \Ob \uS_i$. Alors pour que $\uS$ % [sic] On ne suppose que -> On ne suppose pas que, « pas » apparemment omis
\printedpage{185}
% head: A. Grothendieck : Dérivateurs, Ch. XVIII, Édition des Universités de Montpellier II et Paris VII
% foot: Version 17 décembre 2025
soit accessible dans $\cM$, il faut et il suffit qu'il existe une \textup{petite} partie $J \subset I$, telle que $\Ob \uS = \bigcup_{i\in J} \uS_i$. (En particulier, si $I$ est ordonnée, et $i \longmapsto \uS_i$ croissante et grande devant tout petit cardinal, il faut et il suffit qu'il existe un $i \in I$ tel que $\uS = \uS_i$, i.e. tel que $j \geqslant i$ implique $\uS_j = \uS_i$.)}

C'est évidemment suffisant, en vertu de 4.13\dots{} Prouvons inversement que si $\uS$ est accessible dans $\cM$, alors la condition précédente est satisfaite. Soit $\pi$ un cardinal $\geqslant \pi_0$ tel que $\uS$ soit $\pi$-accessible dans $\cM$, donc
\[
\uS \ \xleftarrow{\ \sim\ }\ \Ind_\pi \uS(\pi) \ , \qquad\qquad \text{où } \uS(\pi) = \uS \cap \cM_\pi.
\]

\msp{page 478}
Soit $\uS(\pi)_0$ une sous-catégorie réduite de $\uS(\pi)$, telle que $\uS(\pi)_0 \hookrightarrow \uS(\pi)$ soit une équivalence, de sorte que $\uS(\pi)_0$ est petite. Pour tout $x \in \Ob \uS(\pi)_0$, soit $i_x \in I$ tel que $x \in \Ob \uS_{i_x}$. Soit $J$ l'ensemble de tous les $i_x$, que est une petite partie de $I$, et soit $\uS_J$ la sous-catégorie strictement pleine de $\cM$ ayant pour ensemble d'objets $\bigcup_{j\in J} \Ob \uS_j$. Donc $\uS(\pi)_0$ est contenue dans $\uS_J$, donc aussi $\uS(\pi)$, puisque $\uS_J$ est strictement pleine : % [sic] que est -> qui est, pronom relatif sujet
\[
\uS(\pi) \ \hookrightarrow\ \uS_J \ .
\]
D'ailleurs $\uS_J$ est accessible dans $\cM$ (loc. cit.), et est fermée dans $\cM$ comme réunion de sous-catégories fermées, \textit{a fortiori elle satisfait à $L_\pi$}. Donc, comme $\uS(\pi) \subset \uS_J$, le foncteur essentiellement surjectif $\Ind_\pi(\uS(\pi)) \longrightarrow \uS$ se factorise par $\uS_J$. Donc $\uS_J \hookrightarrow \uS$ est essentiellement surjectif, et comme $\uS_J$ est strictement pleine, on en conclut $\uS_J = \uS$, q.e.d.

\msp{page 479}
\textbf{Proposition 4.16.7.2.} \textit{Soient $\cM$ une catégorie accessible, $x$ un objet de $\cM$, et ${}_x\cM$ la sous-catégorie fermée de $\cM$ engendrée par $x$, formée des $y \in \Ob \cM$ tels que $\Hom(x,y) \neq \emptyset$, i.e. l'image essentielle de $x\backslash\cM \longrightarrow \cM$. Alors ${}_x\cM$ est une sous-catégorie de $\cM$ accessible dans $\cM$.}

En effet, c'est l'image inverse de la sous-catégorie $\mathrm{Ens}^*$ accessible dans $\mathrm{Ens}$ par le foncteur représentable (donc accessible)
\[
h_x \ :\ y \longmapsto \Hom(x,y) \ :\ \cM \ \longrightarrow\ \mathrm{Ens} \ ,
\]
donc est accessible dans $\cM$ en vertu de 4.13\dots

\textbf{Corollaire 4.16.7.3.} \textit{Soit $\cM$ une catégorie accessible, $\uS$ une sous-catégorie fermée. Pour que $\uS$ soit accessible, il faut et il suffit qu'il existe une \textup{petite} partie $J$ de $\uS$, telle que $\Ob \uS = \bigcup_{x\in J} \Ob\,(\,{}_x\cM)$.}

En effet, on a
\[
\Ob \uS \ =\ \bigcup_{x\in\Ob\uS} \Ob\,(\,{}_x\cM) \ ,
\]
i.e. $\uS$ apparaît comme réunion de
\msp{page 480}
la famille (pas nécessairement petite) de sous-catégories fermées ${}_x\cM$ de $\cM$ (indexées par $x \in \Ob \uS$), accessibles dans $\cM$ en vertu de 4.16.7.2. Donc notre assertion est un cas particulier de 4.16.7.1.

\printedpage{186}
% head: A. Grothendieck : Dérivateurs, Ch. XVIII, Édition des Universités de Montpellier II et Paris VII
% foot: Version 17 décembre 2025

On peut aussi, pour toute partie petite $J$ de $\Ob \uS$, considérer la sous-catégorie fermée $\uS_J \subset \uS$ de $\cM$ engendrée par $J$. L'ensemble $I$ des parties $J$ en question, ordonné par inclusion, est grand devant tout cardinal $\pi \in \U$, et $\uS$ est réunion croissante des $\uS_J$. D'autre part, la proposition nous dit que $\uS$ est accessible dans $\cM$ si et seulement s'il existe un $J \in I$, tel que $\uS_J = \uS$, i.e. tel que $\uS_J = \uS_{J'}$ pour $J \subset J'$.

\begingroup\itshape
\noindent\textbf{\upshape Proposition 4.16.7.4.} Soit $\cM$ une catégorie accessible. Les conditions suivantes sont équivalentes :
\begin{enumerate}[label={}]
\item[a)] Il existe une sous-catégorie $\uS$ de $\cM$ fermée dans $\cM$, et non accessible dans $\cM$.
\item[b)] Soit $I$ l'ensemble des ordinaux $\alpha$ tels que $\card \alpha \in \U$, i.e. des
\msp{page 481}
ordinaux $\alpha$ qui sont dans $\U$. Donc $I = I_{\alpha_0}$, où $\alpha_0$ est le plus petit ordinal qui n'appartienne pas à $\U$, et $I$ est grand devant tout petit cardinal $\pi$, i.e. devant tout cardinal dans $\U$ \textup{(}\footnote{introduire $I = I_{\U}$ avant l'énoncé de la proposition.}\textup{)}. La condition envisagée est la suivante :

Il existe une suite transfini % [sic] transfini -> transfinie, accord avec « suite »
\[
\underline{x} \ =\ (x_\alpha)_{\alpha\in I}
\]
d'éléments de $\Ob \cM$, telle que l'on ait
\[
\Hom(x_\beta, x_\alpha) \ =\ \emptyset \qquad\qquad \textit{si } \beta < \alpha.
\]
De façon précise, pour une telle suite transfinie, soit $\uS$ ou ${}_{\underline{x}}\cM$ la sous-catégorie fermée de $\cM$ engendrée par la famille des $x_\alpha$, i.e. telle que $\Ob \uS = \bigcup_{\alpha\in I} \Ob\ {}_{x_\alpha}\cM = \bigcup_{\alpha\in I} x_\alpha\backslash\Ob \cM$. Alors ${}_{\underline{x}}\cM$ est une sous-catégorie fermée non accessible dans $\cM$, et toute sous-catégorie fermée non accessible dans $\cM$ s'obtient de cette façon.
\end{enumerate}
\endgroup

\textsc{Démonstration}. Considérons une suite $\underline{x} = (x_\alpha)_{\alpha\in I}$ comme dessus. Si $J$ est une petite partie de $I$, il existe un $\alpha \in I$ qui majore les $\beta \in J$, donc $\Hom(x_\beta, x_\alpha) = \emptyset$, pour $\beta \in J$,
\msp{page 482}
ce qui signifie que $x_\alpha$ n'appartient à aucune des $\uS_\beta = {}_{x_\beta}\cM$, pour $\beta \in J$. Donc $\uS$ n'est par réunion d'une petite sous-famille de la famille des ${}_{x_\alpha}\cM$, donc en vertu de 4.16.7.1, $\uS$ n'est par accessible dans $\cM$. % [sic] n'est par réunion -> n'est pas réunion ; n'est par accessible -> n'est pas accessible, « par » pour « pas » (2 fois)

Inversement, soit $\uS$ une telle sous-catégorie dans $\cM$. Considérons un bon ordre sur $\Ob \uS$ (objets $s_\alpha$), et construisons la suite des $x_\alpha$ ($\geqslant s_\alpha$) par récurrence transfinie (\footnote{\textbf{N.B.} On a $\card \uS \notin \U$, sinon, $\uS$ serait petite, donc accessible dans $\cM$, contrairement à l'hypothèse.}).

1) On choisit $x_0 =$ premier élément de $\Ob \uS$ (existe, car $\Ob \uS \neq \emptyset$, i.e. $\uS \neq \emptyset$, car sinon, $\uS$ serait accessible dans $\cM$).

2) Si les $x_\beta$ ($\beta < \alpha$) sont construits, on prend pour $x_\alpha$ le premier élément de $\uS$ majorant les $x_\beta$ et qui n'appartient pas à la sous-catégorie fermée de $\cM$ engendrée par les $x_\beta$ ($\beta < \alpha$). (Cette sous-catégorie fermée est $\subset \uS$, et $\neq \uS$, car $\uS$ n'est pas accessible dans $\cM$.) On a donc bien
\[
\Hom(x_\beta, x_\alpha) \ =\ \emptyset \qquad\qquad \text{si } \beta < \alpha.
\]

\printedpage{187}
% head: A. Grothendieck : Dérivateurs, Ch. XVIII, Édition des Universités de Montpellier II et Paris VII
% foot: Version 17 décembre 2025

\msp{page 483}
Ainsi, la récurrence transfinie se poursuit sur $I$ tout entier. On voit par récurrence transfinie qu'on a, pour tout $\alpha \in I$,
\[
s_\alpha \ \leqslant\ x_\alpha
\]
(card $x_\alpha > x_\beta \geqslant s_\beta$ pour $\beta < \alpha$, donc $x_\alpha > s_\beta$ pour tout $\beta < \alpha$, ce qui équivaut à $x_\alpha \geqslant s_\alpha$). De plus, on trouve aussi, par récurrence transfinie, que
\[
\beta \ <\ \alpha \qquad\quad \Longrightarrow \qquad\quad
\text{\parbox[c]{5.5cm}{$s_\beta$ appartient à la sous-catégorie fermée $\uS_\alpha$ de $\cM$ engendrée par les $x_{\beta'}$ avec $\beta' < \alpha$.}}
\quad ,
\tag*{$(*)$}
\]
puisque par construction $s_\alpha$ est le \textit{premier} des éléments de $\uS$ qui n'est \textit{pas} dans ladite sous-catégorie $\uS_\alpha$ de $\cM$. Cela montre que tout $s_\beta$ appartient à la sous-catégorie fermée ${}_{\underline{x}}\cM$ de $\cM$ engendrée par la famille de \textit{tous} les $x_\alpha$, donc on trouve que $\uS \subset {}_{\underline{x}}\cM \subset \uS$, donc $\uS = {}_{\underline{x}}\cM$, q.e.d.

\begingroup\itshape
\noindent\textbf{\upshape Proposition 4.16.7.5.} Pour un univers $\U$ donné (contenant $\mathcal{N}$ comme élément), les conditions suivantes sont équivalentes :
\begin{enumerate}[label={}]
\item[1)] Il existe une catégorie $\U$-accessible $\cM$ ayant une sous-catégorie fermée non accessible % [sic] 1) -> a), la démonstration renvoie à a) et b)
\msp{page 484}
dans $\cM$.
\item[b)] Il existe une catégorie $\cC \in \U$, et une suite transfinie d'éléments de
\[
\cC^\wedge \ \overset{\text{\upshape déf}}{=}\ \underline{\mathrm{Hom}}(\cC, \mathrm{Ens}_{\U}) \ ,
\]
$\underline{x} = (x_\alpha)_{\alpha\in I}$, indexée par l'ensemble de tous les ordinaux $\alpha \in \U$, et telle que
\[
\Hom(x_\alpha, x_\beta) \ =\ \emptyset \qquad\qquad \textit{si } \alpha < \beta.
\]
\end{enumerate}
\endgroup

\textsc{Démonstration}. On a b) $\Longrightarrow$ a) en prenant $\cM = \cC^\wedge$, en vertu de la proposition précédente. On voit de même a) $\Longrightarrow$ b), en plongeant $\cM$ dans une catégorie $\cC^\wedge$ ($\cC$ étant une petite sous-catégorie strictement génératrice de $\cM$).

Je n'ai pas été capable (\footnote{Mais si, finalement, cf. 4.16.7.10 plus bas.}) de construire une suite $(x_\alpha)_{\alpha\in I}$ comme exigée dans cette proposition, ni de prouver qu'une telle suite n'existe pas, même dans un cas particulier tel que $\cM = (k\text{-alg. comm.})$, $k$ un anneau commutatif. Je n'ai pas
\msp{page 485}
même prouvé qu'on ne puisse trouver une suite transfinie de \textit{corps} $k_\alpha$, ayant les propriétés dites.

A fortiori, je n'ai pu tirer au clair la question si une sous-catégorie $\underline{\mathrm{L}}$-stable d'une catégorie $\cM$ accessible est toujours accessible dans $\cM$ (\footnote{C'est faux, cf. plus bas.}). Il se pourrait qu'il n'y ait pas de tel cas. Il se pourrait aussi qu'il y en ait, même avec une sous-catégorie fermée (laquelle est stable par tout type de $\varinjlim$ représentable dans $\cM$, sauf la $\varinjlim$ sur une catégorie d'indices vide $\ldots$).

\printedpage{188}
% head: A. Grothendieck : Dérivateurs, Ch. XVIII, Édition des Universités de Montpellier II et Paris VII
% foot: Version 17 décembre 2025

Pour les sous-catégories \textit{ouvertes}, je rappelle que j'ai construit (dans l'appendice) un exemple d'une telle sous-catégorie, non accessible dans $\cM$. Au sujet des sous-catégories ouvertes, je signale ici l'énoncé (quasi tautologique) :

\msp{page 486}
\textbf{Proposition 4.16.7.6} (\footnote{Cette proposition et son corollaire sont contenus dans 4.16.7.8 ci-dessous, il n'y a pas lieu d'en faire un énoncé séparé.}). \textit{Soient $\cM$ une catégorie accessible, $\uS$ une sous-catégorie \textup{ouverte} de $\cM$. Pour que $\uS$ soit accessible dans $\cM$, il faut et il suffit qu'elle soit $\underline{L}$-stable dans $\cM$, i.e. qu'il existe un cardinal $\pi_0$ tel que $\cM$ satisfasse à $L_{\pi_0}$, et que $\uS$ soit stable dans $\cM$ par $\varinjlim$ de type $L_{\pi_0}$. De façon plus précise, si cette condition est satisfaite, alors pour tout cardinal $\pi$ tel que $\pi \geqslant \pi_0$ et que $\pi$ soit utile pour $\cM$, $\uS$ est $\pi$-accessible dans $\cM$. Ou encore : si $\pi$ est un cardinal utile pour $\cM$, alors $\uS$ est $\pi$-accessible dans $\cM$ si et seulement s'il est $L_\pi$-stable dans $\cM$.}

Cette dernière assertion résulte aussitôt du corollaire suivant, lui-même tautologique.

\textbf{Corollaire 4.16.7.7.} \textit{Soit $\pi$ un cardinal comme indiqué ci-dessus, soit $\uS(\pi) = \uS \cap \cM_\pi$, et soit $x \in \Ob \cM$ et
\[
x \ \simeq\ \varinjlim_{i\in I} x_i \qquad\qquad
\begin{cases}
x_i \in \Ob \cM_\pi \\[1ex]
I \ \textit{grande devant}\ \pi
\end{cases}
\]
une représentation de $x$ comme image essentielle d'un objet de $\Ind_\pi \cM_\pi$. Alors $x \in \Ob \uS$ $\Longleftrightarrow$ $x_i \in \Ob \uS(\pi)$ $\forall\, i \in I$.}

\msp{page 487}
J'ai finalement tiré au clair la question de l'accessibilité des sous-catégories ouvertes et des sous-catégories fermées d'une catégorie accessible :

\begingroup\itshape
\noindent\textbf{\upshape Proposition 4.16.7.8.} Soient $\cM$ une catégorie accessible, $\uS$ une sous-catégorie fermée de $\cM$, $\uT$ la sous-catégorie ouverte complémentaire de $\uS$, $\pi$ un cardinal infini utile pour $\cM$. Les conditions suivantes sont équivalentes :
\begin{enumerate}[label={}]
\item[a)] $\uS$ est $\pi$-accessible dans $\cM$.
\item[a$'$)] $\uS$ est le fermé engendré par $\uS(\pi) \overset{\text{\upshape déf}}{=} \uS \cap \cM_\pi$, i.e. pour $x \in \Ob \cM$, on a $x \in \Ob \uS$ si et seulement si $\exists\, \xi \in \Ob \uS(\pi)$ tel que $\Hom_{\cM}(\xi, x) \neq \emptyset$.
\item[b)] $\uT$ est $\pi$-accessible dans $\cM$.
\item[b$'$)] $\uT$ est $\underline{L}_\pi$-stable dans $\cM$.
\end{enumerate}
\endgroup

\noindent Notons tout de suite le

\begingroup\itshape
\noindent\textbf{\upshape Corollaire 4.16.7.9.} Pour $\cM$, $\uS$, $\uT$ comme dessus, conditions équivalentes :
\begin{enumerate}[label={}]
\item[a)] $\uS$ est accessible dans $\cM$.
\item[a$'$)] Il existe une \textup{petite} partie $A$ de $\Ob \cM$, telle que $\uS$ soit la sous-catégorie fermée ${}_A\cM$ engendrée par $A$, i.e. $\uS = \bigcup_{\xi\in A} {}_\xi\cM$, où ${}_\xi\cM$ est la sous-catégorie fermée de $\cM$ formée des $x \in \Ob \cM$ tel que $\Hom(\xi, x) \neq \emptyset$.
\item[b)] \msp{page 488}$\uT$ est accessible dans $\cM$.
\end{enumerate}
\endgroup


\endgroup
% END UNIT der-ch18/der-ch18_p186-190.tex
% BEGIN UNIT der-ch18/der-ch18_p191-195.tex
\begingroup\makeatletter%
\newcommand{\printedpage}[1]{\marginpar{\footnotesize p.~#1}}%
\newcommand{\msp}[1]{\marginpar{\footnotesize [#1]}}%
\newcommand{\cM}{\mathcal{M}}%
\newcommand{\cN}{\mathcal{N}}%
\newcommand{\cF}{\mathcal{F}}%
\newcommand{\cA}{\mathcal{A}}%
\newcommand{\Ob}{\operatorname{Ob}}%
\newcommand{\Hom}{\operatorname{Hom}}%
\newcommand{\Ind}{\operatorname{Ind}}%
\newcommand{\Ker}{\operatorname{Ker}}%
\newcommand{\Max}{\operatorname{Max}}%
\newcommand{\Ann}{\operatorname{Ann}}%
\newcommand{\defeq}{\overset{\text{déf}}{=}}%
\newcommand{\incl}{\lhook\joinrel\longrightarrow}%
\makeatother\let\printedpage\derpp\let\msp\dermsp%
\let\mspin\DERmspin%


\printedpage{189}
% head: A. Grothendieck : Dérivateurs, Ch. XVIII, Édition des Universités de Montpellier II et Paris VII

\begin{enumerate}
\item[\textit{b$'$)}] \textit{$\underline{T}$ est $\underline{L}$-stable dans $\cM$.}
\end{enumerate}

\noindent\textbf{Corollaire 4.16.7.10.} \textit{Soit $\cM$ une catégorie accessible. Les conditions suivantes sont équivalentes} (\footnote[202]{Exemple où ces conditions sont satisfaites : $\cM = (k\text{-alg. comm.})$, $k$ un corps commutatif infini, $\mathfrak{U}$ étant le plus petit univers tel que $\mathbf{N}\in\mathfrak{U}$.}).
\begin{enumerate}
\item[\textit{a)}] \textit{Il existe une sous-catégorie fermée de $\cM$, non accessible dans $\cM$.}
\item[\textit{b)}] \textit{Il existe une sous-catégorie ouverte dans $\cM$, non accessible dans $\cM$.}
\item[\textit{c)}] \textit{Il existe une suite transfinie}
\[
\underline{x} \ = \ (x_\alpha)_{\alpha\in I}
\]
\textit{d'objets de $\cM$, indexée par l'ensemble $I$ des ordinaux de l'univers de référence $\mathfrak{U}$, telle que l'on ait}
\[
\Hom(x_\alpha,x_\beta) \ = \ \emptyset \qquad \text{\textit{pour} } \alpha<\beta .
\]
\end{enumerate}

\noindent En effet, l'équivalence de a) et b) est claire par le corollaire précédent, et celle de a) et c) n'est autre que 4.16.7.4.

\noindent\textsc{Démonstration de} 4.16.7.8. Nous aurons à utiliser le

\noindent\textbf{Lemme 4.16.7.11.} \textit{Soient $\cM$ une catégorie accessible, $\pi$ un cardinal utile pour $\cM$, $\underline{S}$ une sous-catégorie pleine de $\cM$, $\underline{S}(\pi)\defeq\underline{S}\cap\cM_\pi$, d'où un foncteur pleinement fidèle}
\[
\varphi \ : \ \Ind_\pi(\underline{S}(\pi)) \longrightarrow \cM \ .
\]

\msp{page 489}
\begin{enumerate}
\item[\textit{a)}] \textit{Soit $x\in\Ob\cM$. Pour que $x$ appartienne à l'image essentielle du foncteur $\psi$, il faut et il suffit que la sous-catégorie pleine $\underline{S}(\pi)/x$ de la catégorie $\cM_\pi/x$ (qui est filtrante et même grande devant $\pi$) soit cofinale (donc aussi grande devant $\pi$).} % [sic] \psi -> \varphi, the functor of the lemma is called \varphi in the display above
\item[\textit{b)}] \textit{Pour que $\underline{S}$ soit $\pi$-accessible dans $\cM$, il faut et il suffit que pour tout $x\in\Ob\cM$, on ait}
\[
x \ \in \ \Ob\underline{S} \qquad\iff\qquad \underline{S}(\pi)/x \text{ \textit{est cofinale dans} } \cM_\pi/x
\]
\textit{(ou encore, que $\underline{S}$ satisfasse à $L_\pi$, et que pour tout $x$ dans $\cM$, on ait}
\[
x \ \in\Ob\underline{S} \qquad\Longrightarrow\qquad \underline{S}(\pi)/x \text{ \textit{cofinale dans} } \cM_\pi/x\text{\textit{.)}}
\]
\end{enumerate}

\noindent\textsc{Démonstration}. La partie b) est une conséquence tautologique de a), que je vais prouver. Dire que $x$ est dans l'image essentielle de $\underline{S}$, signifie que
\[
x \ \simeq \ \varinjlim_{i\in I} x_i \ , \tag*{$(*)$}
\]
avec $I$ grande devant $\pi$ et les $x_i$ dans $\underline{S}(\pi)$. On sait que $(*)$ implique ($I$ étant filtrante, et $\cM$ étant équivalente à $\Ind_\pi(\cM)$ puisque $\pi$ utile pour $\cM$) que $I\longrightarrow\cM_\pi/x$ est cofinale (4.4.2, page 79). Or ce foncteur se factorise par le foncteur d'inclusion pleinement fidèle $\underline{S}(\pi)/x\incl\cM_\pi/x$, lequel est cofinal si et seulement s'il satisfait la condition F 1. Il est immédiat que cette condition est satisfaite quand elle est satisfaite pour le foncteur composé en question, ce qui est bien le cas
\printedpage{190}
% head: A. Grothendieck : Dérivateurs, Ch. XVIII, Édition des Universités de Montpellier II et Paris VII
\msp{page 490}
par hypothèse.

\noindent Inversement, supposons que $\underline{S}(\pi)/x$ soit cofinale dans $\cM_\pi/x$. Comme on a
\[
x \xleftarrow{\ \sim\ } \varinjlim_{\cM_\pi/x} x_i \ ,
\]
$\pi$ étant utile pour $\cM$, on aura aussi par cofinalité
\[
x \xleftarrow{\ \sim\ } \varinjlim_{\underline{S}(\pi)/\pi} x_i \ , % [sic] \underline{S}(\pi)/\pi -> \underline{S}(\pi)/x, the cofinal subcategory is \underline{S}(\pi)/x
\]
or $\cM_\pi$ étant grande devant $\pi$, il en est de même de la sous-catégorie pleine cofinale $\underline{S}(\pi)/x$. Donc $x$ est dans l'image essentielle de $\Ind_\pi\underline{S}(\pi)$, q.e.d.

\noindent\textbf{Lemme 4.16.7.12.} \textit{Soient $I$ une catégorie filtrante, $J$ une sous-catégorie fermée, $K$ la sous-catégorie ouverte complémentaire. On a alors les équivalences}
\[
J \ \neq \ \emptyset \iff J \text{ \textit{cofinale} } \iff K \text{ \textit{non cofinale} } \iff K \ \neq \ I \ .
\]

\noindent\textsc{Démonstration} immédiate.

\noindent Je reviens à la démonstration de 4.16.7.8. On applique le lemme 4.16.7.11 à $\cM$, $\pi$, et aux deux
\msp{page 491}
sous-catégories $\underline{S}$, $\underline{T}$ de $\cM$, en notant que pour $x\in\Ob\cM$, $\underline{S}(\pi)/x$ est une sous-catégorie fermée de $\underline{S}(\pi)$, dont le complémentaire ouvert est $\underline{T}(\pi)/x$, de sorte qu'on peut appliquer le lemme 4.16.7.12. On trouve :
\[
\begin{aligned}
\underline{S}\ \pi\text{-accessible} \quad &\overset{4.16.7.11}{\iff} \quad \forall\ x\ \in\ \Ob\cM,\ \text{on a } (x\ \in\ \Ob\underline{S}\iff\underline{S}(\pi)/x \text{ cofinale dans}\\
&\qquad\quad \cM_\pi/x)\\
&\iff \quad \forall\, x\in\Ob\cM,\ \text{on a } \Bigl( \qquad \underbrace{x\notin\Ob\underline{S}}_{\text{(i.e. $x\in\Ob\underline{T}$, par 4.16.7.12)}}\\
&\qquad\qquad\qquad \iff \qquad \underbrace{\underline{S}(\pi)/x \text{ non cofinale dans } \cM_\pi/x}_{\text{(i.e. $\underline{T}(\pi)/x$ cofinale dans $\cM_\pi$, par 4.16.7.12)}} \ \Bigr)\\
&\overset{4.16.7.11}{\iff} \quad \underline{T} \text{ est } \pi\text{-accessible.}
\end{aligned}
\]
On a donc l'équivalence de a) et b). D'autre part, $\underline{S}(\pi)/x$ cofinale dans $\cM_\pi/x$ signifie simplement $\underline{S}(\pi)/x\neq\emptyset$ (4.16.7.12), donc $\underline{S}$ $\pi$-accessible signifie que $\underline{S}$ est formée des $x\in\Ob\cM$ tels que $\underline{S}(\pi)\neq\emptyset$, ce qui signifie aussi que $\underline{S}$ est le fermé dans $\cM$ engendré par $\underline{S}(\pi)$, d'où a $\iff$ a$'$. D'autre part, l'équivalence b $\iff$ b$'$ n'est autre que 4.6.7.6 (dernière formulation). % [sic] 4.6.7.6 -> 4.16.7.6, the surrounding references are all to 4.16.7.x

\msp{page 492}

\section*{4.17\quad Retour sur les paratopos, et sur les opérations $\cM\boxtimes\cN$, $\underline{\operatorname{Hom}}_{!}(\cM,\cN)$}

\noindent\texttt{[Ce paragraphe n'existe pas.]}

\printedpage{191}
% head: A. Grothendieck : Dérivateurs, Ch. XVIII, Édition des Universités de Montpellier II et Paris VII

\msp{page 1}

\section*{Appendice}

\subsection*{Ultraproduits de corps, et catégories accessibles}

\noindent Soit $(k_\xi)_{\xi\in E}$ une famille de corps, indexé par un ensemble $E$. $\Bigl($ Soit $\Phi$ un filtre non trivial dans $E$, i.e. une partie $\Phi\subset\mathfrak{P}(E)$ telle que :
\begin{enumerate}
\item[a)] $\emptyset\notin\Phi$, i.e. $A\in\Phi\Longrightarrow A\neq\emptyset$.
\item[b)] $A,B\in\Phi\Longrightarrow A\cap B\in\Phi$.
\item[c)] Si $A\in\Phi$, $A\subset B$, alors $B\in\Phi$.
\item[d)] $\Phi\neq\emptyset$ ($\iff E\in\Phi$).
\end{enumerate}
(Cela revient à dire que $\Phi$ est une partie fermée de l'ensemble ordonné $\mathfrak{P}(E)$, \textit{filtrante à gauche}, et ne contient pas le plus petit élément de $\mathfrak{P}(E)$, i.e. non égale à $\mathfrak{P}(E)$ tout entier.) $\Bigr)$ Pour toute partie $J$ de $E$, posons
\[
k_J \ = \ \prod_{\xi\in J} k_\xi \ ,
\]
on a donc un système projectif
\[
(k_J)_{J\in\mathfrak{P}(E)} \ , \qquad \text{i.e. } J\subset J' \text{ donne } k_{J'}\longrightarrow k_J \ ,
\]
et le système
\[
(k_J)_{J\in\Phi}
\]
est un système inductif d'annaux, indexé par l'ensemble ordonné filtrant $\Phi^o$. Posons % [sic] d'annaux -> d'anneaux, misspelling
\[
k_\Phi \ = \ \varinjlim_{J\in\Phi} k_J \ .
\]
Si $\Psi$ est un filtre plus fin que $\Phi$, i.e.
\msp{page 2}
tel que $\Psi\supset\Phi$, on trouve un homomorphisme canonique
\[
k_\Phi \longrightarrow k_\Psi \ ,
\]
donc si $\cF(E)=\cF$ est l'ensemble des filtres sur $E$, ordonnés par inclusion, on trouve un système inductif
\[
\Phi \longmapsto k_\Phi \ : \ \cF \longrightarrow \Ann \ .
\]
On n'a pas utilisé encore que les $k_\xi$ sont des corps, en fait il suffirait que ce soient des objets d'une catégorie $\cA$ stable par petits produis et par petits $\varinjlim$ filtrantes. Prenons maintenant $\cA=\Ann$ (catégorie des anneaux unitaires, non nécessairement commutatifs), et supposons que les $k_\xi$ soient des corps. % [sic] produis -> produits, misspelling

\printedpage{192}
% head: A. Grothendieck : Dérivateurs, Ch. XVIII, Édition des Universités de Montpellier II et Paris VII

\noindent\textbf{Proposition 1.} \textit{Supposons que les $k_\xi$ ($\xi\in E$) soient des} corps (\footnote[203]{NB Cette condition est nécessaire si on veut que $k_\Phi$ soit un corps pour tout ultrafiltre.}). \textit{Pour que $k_\Phi$ soit un corps, il faut et il suffit que $\Phi$ soit un ultrafiltre, i.e. que pour toute décomposition $E=A\amalg B$, on ait $A\in\Phi$ ou $B\in\Phi$.} (NB Si $A\in\Phi$, on a $B\notin\Phi$, puisque sinon, on aurait $A\cap B\in\Phi$, donc $\emptyset\in\Phi$, absurde.)

\noindent Nécessité : Soit, pour toute partie $A$ de $E$,
\[
f_A\in k_E \qquad \text{défini par} \qquad f_A(\xi) \ = \ \left\{
\begin{array}{lll}
1_{k_\xi} & \text{si} & \xi\in A\\[1ex]
0_{k_\xi} & \text{si} & \xi\in E\backslash A \ .
\end{array}\right.
\]

\msp{page 3}
\noindent Soit
\[
f_{A,\Phi} \ \in k_\Phi
\]
l'image de $f_A$ par l'homomorphisme canonique $k_E\longrightarrow k_\Phi$. Si $E=A\amalg B$, on aura
\[
f_A \ + \ f_B \ = \ 1 \ , \qquad f_A\, f_B \ = \ 0 \ ,
\]
d'où
\[
f_{A,\Phi} \ + \ f_{B,\Phi} \ = \ 1 \ , \qquad f_{A,\Phi}\, f_{B,\Phi} \ = \ 0 \ .
\]
Si $k_\Phi$ est un corps, cela implique qu'on a $f_{A,\Phi}=0$ ou $f_{B,\Phi}=0$. Mais si p.ex. $f_{A,\Phi}=0$, on voit que cela signifie qu'il existe un $J\in\Phi$ tel que $f_A|J=0$, ce qui signifie aussi que $J\subset B$. De même, si $f_{B,\Phi}=0$, on aura $J\subset A$. Donc $\Phi$ est un ultrafiltre.

\noindent Inversement, supposons que $\Phi$ soit un ultrafiltre. Prouvons que $k_\Phi$ est un corps, i.e. que pour $f\in k_\Phi$, $f\neq\emptyset\Longrightarrow f$ inversible. Notons que $f$ provient d'un élément d'un $k_J$ ($J\in\Phi$), donc est l'image canonique d'un $g\in k_E$. Dire que $f=0$ signifie que $\exists\ J\in\Phi$ tel que $g|J=0$. L'inverse $f\neq0$ signifie donc que pour tout $J\in\Phi$, on a $g|J\neq\emptyset$.

\msp{page 4}
\noindent Soit alors $A$ l'ensemble des $\xi\in E$ tel que $g(\xi)\neq\emptyset$, $B=E\backslash A$ l'ensemble des $\xi\in E$ tels que $g(\xi)\neq\emptyset$. On a donc $B\notin\Phi$ puisque $f\neq0$, et comme $\Phi$ est un ultrafiltre, on a donc $A\in\Phi$. Mais soit $g'\in k_E$ défini par $g'(\xi)=f(\xi)^{-1}$ si $\xi\in A$, $g'(\xi)=1$ si $\xi\in B$, et soit $f'$ l'image de $g'$ dans $k_\Phi$. On voit alors qu'on a % [sic] B = E\A ... tels que g(\xi) \neq \emptyset -> g(\xi) = 0, the complement of A is the set where g vanishes
\[
f\ f' \ = \ f' \ = \ f \ = \ 1 \ ,
\]
ce qui montre bien que $k_\Phi$ est un corps.

\noindent Supposons maintenant que tous les $k_\xi$ ($\xi\in E$) soient égaux à un même corps $k$. Alors les $k_J$ forment un système projectif de $k$-algèbres, les $k_\Phi$ un système inductif de $k$-algèbres. Si $k$ est commutatif, il en est de même des $k_J$, $k_\Phi$, et en particulier pour $\Phi$ un ultrafiltre, $k_\Phi$ est une \textit{extension} du corps $k$.

\noindent Revenons au cas général, mais les $k_\xi$ commutatifs. Soit $\cF_0\subset\cF$ l'ensemble des ultrafiltres sur $E$. Notons que pour tout $\Phi\in\cF$, $k_E\longrightarrow k_\Phi$ est surjectif. Donc si $\Phi\in\cF$, $k_\Phi$ s'identifie à un corps quotient de l'anneau $k_E$.

\printedpage{193}
% head: A. Grothendieck : Dérivateurs, Ch. XVIII, Édition des Universités de Montpellier II et Paris VII

\msp{page 5}
\noindent Pour tout anneau $\cA$, soit $\Max(\cA)$ le spectre maximal, i.e. l'ensemble des idéaux maximaux de $\cA$. On trouve donc une application canonique
\[
\left\{
\begin{array}{lll}
\cF_0 & \longrightarrow & \Max(k_E)\\[1ex]
\ \Phi & \longmapsto & \Ker(k_E\longrightarrow k_\Phi) \ \defeq \ \mathfrak{m}_\Phi \ .
\end{array}\right. \tag*{$(*)$}
\]

\noindent\textbf{Proposition 2.} \textit{L'application précédente est} bijective.

\noindent Prouvons d'abord qu'elle est injective, cela résulte aussitôt du

\noindent\textbf{Corollaire 1} (\footnote[204]{\textbf{NB} Pour $\Phi\in\cF$ quelconque, $\mathfrak{m}_\Phi\defeq\Ker(\chi_\Phi : k_E\longrightarrow k_\Phi)$ détermine $\Phi$, puisque
\[
A \ \in \ \Phi \qquad\iff\qquad \chi_\Phi(f_A) \ = \ 1
\]
(mais il faut supposer les $K_\xi\neq0$, i.e. $0_\xi\neq1_\xi\ \ldots$).}). \textit{Soit, pour $\Phi\in\cF_0$,}
\[
\chi_\Phi \ : \ k_E \longrightarrow k_\Phi
\]
\textit{l'application canonique. Alors pour $A\in\mathfrak{P}(E)$, on a}
\[
A \ \notin \ \Phi \qquad\iff\qquad \chi_\Phi(f_A) \ = \ 0 \ , \quad \text{\textit{i.e.} } f_A\in\mathfrak{m}_\Phi
\]
\textit{(i.e. $A\in\Phi\iff\chi_\Phi(f_A)\neq0$, i.e. $f_A\notin\mathfrak{m}_\Phi$)} (\footnote[205]{\textbf{NB} On a donc
\[
(\chi_\Phi(f_A) \ \neq \ 0) \qquad\iff\qquad (\chi_\Phi(f_A) \ = \ 1) \ .
\]}).

\noindent En effet, $\chi_\Phi(f_A)=0$ signifie qu'il existe $B\in\Phi$ tel que $f_A|B=0$, i.e. tel que $B\subset E\backslash A$ ; mais cela équivaut à $E\backslash A\in\Phi$, donc à $A\notin\Phi$ puisque $\Phi$ est un ultrafiltre, q.e.d.

\msp{page 6}
\noindent Reste à prouver que l'application $(*)$ est surjective. Soit donc $\mathfrak{m}$ un idéal maximal de $k_E$, et soit
\[
\Phi \ \defeq \ \{A\in\mathfrak{P}(E)\mid\chi(f_A) \ = \ 1\} \ ,
\]
où $\chi : k_E\longrightarrow k_E/\mathfrak{m}=\kappa(\mathfrak{m})$ est l'application canonique. Je dis que $\Phi$ est un ultrafiltre. En effet :
\begin{enumerate}
\item[a)] $\emptyset\notin\Phi$ car $f_\emptyset=0$, d'où $\chi(f_\emptyset)=0\neq1$ dans $k_\Phi$.
\item[b)] Si $A,B\in\Phi$, on a $A\cap B\in\Phi$, car $f_{A\cap B}=f_A\cdot f_B$, d'où $\chi(f_{A\cap B})=\chi(f_A)\chi(f_B)=1$.
\item[c)] Si $A\subset B$ et $A\in\Phi$, alors $B\in\Phi$. En effet, on a $f_A=f_Af_B$, d'où $\chi(f_A)=\chi(f_A)\chi(f_B)$, et comme $\chi(f_A)=1$, on trouve $\chi(f_B)=1$.
\item[d)] On a $E\in\Phi$, i.e. $\chi(f_E)=1$, or $f_E=1$, donc $\chi(f_E)=1$.
\end{enumerate}
Ainsi, $\Phi$ est un filtre -- on n'a utilisé que le fait que $A/\mathfrak{m}$ est un anneau $\neq0$ (pour conclure que $\emptyset\notin\Phi$). Prouvons que si $A/\mathfrak{m}$ est intègre, $\Phi$ est un \textit{ultrafiltre}, i.e. :
\begin{enumerate}
\item[e)] Si $A\notin\Phi$, alors $B\defeq E\backslash A\in\Phi$.
\end{enumerate}
En effet, on a $f_A+f_B=1$, $f_Af_B=0$, d'où


\endgroup
% END UNIT der-ch18/der-ch18_p191-195.tex
% BEGIN UNIT der-ch18/der-ch18_p196-200.tex
\begingroup\makeatletter%
\newcommand{\printedpage}[1]{\marginpar{\footnotesize p.~#1}}%
\newcommand{\msp}[1]{\marginpar{\footnotesize [#1]}}%
\newcommand{\limind}{\mathop{\underrightarrow{\mathrm{lim}}}}%
\newcommand{\limproj}{\mathop{\underleftarrow{\mathrm{lim}}}}%
\makeatother\let\printedpage\derpp\let\msp\dermsp%
\let\mspin\DERmspin%


\printedpage{194}
% head: A. Grothendieck : Dérivateurs, Ch. XVIII, Édition des Universités de Montpellier II et Paris VII

\noindent\hspace*{1em}\textcircled{a}\quad $\chi(f_A)+\chi(f_B)=1$, et\\
\hspace*{1em}\textcircled{b}\quad $\chi(f_A)\chi(f_B)=0$.\\
$\chi(f_A)\neq 1$, on conclut de \textcircled{a} $\chi(f_B)\neq 0$, d'où par \textcircled{b} $\chi(f_A)=0$, d'où de \textcircled{a} $\chi(f_B)=1$, d'où $B\in\Phi$.

\msp{page 7}
Prouvons enfin que $\mathfrak{m}=\mathfrak{m}_\Phi$. Comme $\mathfrak{m}_\Phi$ est un idéal maximal, il suffit de prouver
\[
\mathfrak{m}_\Phi\ \subset\ \mathfrak{m}\ .
\]
Or si $A\in\mathfrak{P}(E)$ est telle que $A\notin\Phi$, ou $f_A\in\mathfrak{m}_\Phi$ (cor.~1), on a aussi $f_A\in\mathfrak{m}$. Donc il suffit de prouver le

\textbf{Lemme.} \textit{Soit $\Phi\in\mathcal{F}_0$. Alors $\mathfrak{m}_\Phi$ est l'idéal engendré par les $f_A$ pour $A\in\mathfrak{P}(E)\backslash\Phi$. Plus précisément, les $f_A$ sont dans $\mathfrak{m}_\Phi$, et tout $f\in\mathfrak{m}_\Phi$ est de la forme $ff_A$, avec $A$ comme dessus.}

En effet, $f\in\mathfrak{m}_\Phi$ signifie qu'il existe $B\in\Phi$ tel que $f|B=0$. Mais cela signifie aussi que $f=ff_A$, où $A=E\backslash A$ \texttt{[plutôt} $A=E\backslash B$\texttt{]}, d'où $A\notin\Phi$.

Cela prouve la proposition 2, et en même temps :

\textbf{Corollaire 2.} \textit{Tout anneau quotient intègre de $k_E=\prod_{\xi\in E}k_\xi$ est un corps.}

Notons maintenant que si $\Phi\in\mathcal{F}$ est un filtre quelconque sur $E$, alors $k_\Phi$ est un anneau quotient de $k_E$, donc
\[
\operatorname{Max}(k_\Phi)\ \hookrightarrow\ \underbrace{\operatorname{Max}(k_E)}_{\simeq\,\mathcal{F}_0}\ . \tag*{$(*)$}
\]
Ceci posé, on trouve :

% a vertical bar is printed in the left margin beside the statement of Corollaire 3
\textbf{Corollaire 3} (\footnote[206]{Marche pour tout système d'anneaux $(k_\xi)_{\xi\in E}$ \textit{non nuls}, et $\Phi,\Psi\in\mathcal{F}$, $\Psi$ plus fin que $\Phi$, i.e. $\Phi\subset\Psi$ si et seulement si $\mathfrak{m}_\Phi\subset\mathfrak{m}_\Psi$ (pas la peine que $\Psi$ soit un ultrafiltre).}). \textit{L'inclusion précédente identifie} $\operatorname{Max}(k_\Phi)$ \textit{à l'ensemble des ultrafiltres $\Psi$ sur $E$ qui sont} plus fins \textit{que $\Phi$.}

En effet, si $\Psi$ est un tel ultrafiltre, on
\msp{page 8}
a le diagramme commutatif
\[
\begin{tikzcd}
k_E \arrow[r, "\chi_\Phi"] \arrow[rr, bend right=15, "\chi_\Psi"'] & k_\Phi \arrow[r] & k_\Psi
\end{tikzcd}\ , \tag*{$(**)$}
\]
ce qui montre que $\Psi$ \texttt{[plutôt} $\mathfrak{m}_\Psi$\texttt{]} est dans l'image de $(*)$. Inversement, supposons que $\Psi\in\mathcal{F}_0$ provienne de $\operatorname{Max}(k_\Phi)$, i.e. que l'on ait $\mathfrak{m}_\Phi\subset\mathfrak{m}_\Psi$ (où $\mathfrak{m}_\Phi$ est le noyau de $\chi_\Phi$), ou encore qu'on ait une factorisation $(**)$ de $\chi_\Psi$. Prouvons que $\Psi$ est plus fin que $\Phi$, i.e. que $\Psi\supset\Phi$. Soit donc $A\in\Phi$, prouvons que l'on a $A\in\Psi$. Mais $A\in\Phi$ implique \texttt{[le]} fait que $\chi_\Phi(f_A)=1$, donc son image dans $k_\Psi$ est égale à 1, donc on a $A\in\Psi$ par le corollaire 1.

\printedpage{195}
% head: A. Grothendieck : Dérivateurs, Ch. XVIII, Édition des Universités de Montpellier II et Paris VII

\msp{page 9}
Supposons à présent que tous les $k_\xi$ sont égaux \texttt{[à]} un même anneau \fbox{$\ k\neq 0\ $}\,. Donc $k_J$, $k_\Phi$ sont des `anneaux sous $k$' $k\longrightarrow k_J\longrightarrow k_\Phi$, d'ailleurs $k\longrightarrow k_J$ (si $J\neq\emptyset$) et $k\longrightarrow k_\Phi$ sont injectifs. On se pose la question quand $k\longrightarrow k_\Phi$ est \textit{bijectif}, i.e. surjectif. Comme $k_\Phi\simeq k_E/\mathfrak{m}_\Phi$, cela signifie donc que $k_E=k+\mathfrak{m}_\Phi$, i.e. que tout $f\in k_E$ s'écrit
\[
f\ =\ c\ +\ g\qquad
\begin{cases}
c\ \in\ k\ , & \text{i.e. } c(\xi)=\underbrace{c_0}_{\in\,k}\ \forall\,\xi\in E\ ,\\[2ex]
g\ \in\ \mathfrak{m}_\Phi & \text{i.e. } \exists\,A\in\Phi \text{ avec } g|A=0\ .
\end{cases}
\]
Mais pour $f\in k_E$, $c\in k$ et $A\in\Phi$ fixés, l'existence de $g\in k_E$ telle que $g|A=0$ et $f=c+g$ équivaut à $f|A=c_A$ (fonction constante dans $A$, de valeur $c$), i.e. $A\subset f^{-1}(\{c\})$. Pour $f$, $c$ fixés, l'existence de $A\in\Phi$ satisfaisant la relation équivaut à $f^{-1}(\{c\})\in\Phi$. Et pour $f$ fixé, l'existence de $c$ équivaut à celle d'une fibre $E_c$ de $f:E\longrightarrow k$, qui soit dans $\Phi$. Ainsi :

\textbf{Proposition 3.} \textit{Soit $\Phi\in\mathcal{F}$. Pour que $k\longrightarrow k_\Phi$ soit surjectif, il faut et il suffit que pour toute famille de parties} essentiellement disjointes \textit{$(E_\xi)_{\xi\in k}$ de $E$ indexée par $k$, de réunion $E$, il existe un $\xi\in k$ tel que $E_\xi\in\Phi$. A fortiori, il} faut \textit{que $\Phi$ soit un ultrafiltre (puisque} $\operatorname{card}k\geqslant 2$ \textit{par hypothèse).}

\msp{page 10}
En effet, la donnée d'une telle famille $(E_\xi)_{\xi\in k}$ équivaut à la donnée d'une application $E\longrightarrow k$, i.e. d'un élément de $f_E$. % [sic] $f_E$ -> $k_E$, an application $E\to k$ is an element of $k_E$

\textbf{Scholie.} La propriété qu'on vient d'expliciter pour un filtre $\Phi$ sur $E$, ne dépend de $k$ que par l'intermédiaire de $\pi=\operatorname{card}k$. On peut aussi l'énoncer en disant que pour toute \textit{partition} $(E_i)_{i\in I}$ de $E$ en un ensemble $I$ des parties tel que $\operatorname{card}I\leqslant\pi$, il existe un (et un seul) des $E_i$ qui soit dans $\Phi$. On dit alors que $\Phi$ est un $\pi$-\textit{ultrafiltre.}

Si $\pi\geqslant 2$ est fini, cette propriété équivaut à dire que $\Phi$ est un ultrafiltre. Par ailleurs, quel que soit $\pi$, les ultrafiltres $\Phi$ \textit{triviaux} (i.e. tels qu'il existe $\xi\in E$ avec $\{\xi\}\in\Phi$, ou encore tel que $\Phi=\{A\in\mathfrak{P}(E)\mid\xi\in A\}$) sont des $\pi$-ultrafiltres. Si $\pi$ est infini, il n'est pas clair qu'il existe des $\pi$-ultrafiltres non triviaux. On sait qu'il n'en existe sur $E$ \textit{que} si $\operatorname{card}E\notin\mathfrak{U}_1$, où $\mathfrak{U}_1$ est le plus petit univers tel que $\mathbf{N}\in\mathfrak{U}_1$, i.e. tel
\msp{page 11}
que $\mathfrak{U}$ contienne comme élément un ensemble infini.

Si $\mathfrak{U}=\mathfrak{U}_1$, et si $k$ est un corps infini, p.ex. $k=\mathbf{Q}$, alors pour tout ensemble $E\in\mathfrak{U}$ et pour tout ultrafiltre \textit{non trivial} $\Phi$ sur $E$, $k_E$ est une extension \textit{non triviale} de $k$. Mais si $\mathfrak{U}=\mathfrak{U}_\alpha$ est un univers quelconque, et $\pi$ le premier cardinal dans $\mathfrak{U}$ qui n'appartienne pas à un univers $\mathfrak{U}'\subset\mathfrak{U}$ strictement plus petit, je présume que pour tout $E\in\mathfrak{U}$, tout $\pi$-ultrafiltre sur $E$ est trivial. % [sic] $k_E$ est une extension -> $k_\Phi$ est une extension, the extension of $k$ attached to the ultrafilter $\Phi$ is $k_\Phi$

\printedpage{196}
% head: A. Grothendieck : Dérivateurs, Ch. XVIII, Édition des Universités de Montpellier II et Paris VII

\subsection*{Application à un exemple d'une sous-catégorie non accessible.}

Soit $\mathcal{M}$ une catégorie $\mathfrak{U}$-accessible, et $a\in\operatorname{Ob}\mathcal{M}$, d'où un foncteur
\[
\begin{aligned}
f^o\ :\ \mathcal{M}^o\ &\longrightarrow\ \mathrm{Ens}\\
x\ &\longmapsto\ \operatorname{Hom}_{\mathcal{M}}(x,a)\ \overset{\text{déf}}{=}\ a(x)\ ,
\end{aligned}
\]
foncteur commutant aux $\limproj$ quelconques, donc tel que $f:\mathcal{M}\longrightarrow(\mathrm{Ens})^o$ commute aux $\limind$ quelconques, et en particulier $f$ est accessible. Considérons la sous-catégorie strictement pleine
\msp{page 12}
\[
\underbrace{\mathrm{Ens}^*}_{=\,\mathrm{Ens}\backslash\{\emptyset\}}\ \subset\ \mathrm{Ens}
\]
des ensembles non vides, on sait qu'elle est accessible, mais elle n'est par stable par $\limproj$ de type $\mathrm{L}_\pi$, quelque grand qu'on prenne $\pi$. % [sic] n'est par stable -> n'est pas stable, typo

Considérons
\[
\underbrace{\underline{S}\ =\ \underline{S}_a}_{\overset{\text{déf}}{=}\,f^{-1}(\mathrm{Ens}^*)}\ \subset\ \mathcal{M}\ ,
\]
sous-catégorie strictement pleine de $\mathcal{M}$ formée des $x\in\operatorname{Ob}\mathcal{M}$ tels que
\[
a(x)\ \neq\ \emptyset\ ,\qquad\quad \text{i.e. } x\leqslant a \text{ pour le préordre canonique de } \mathcal{M}\ .
\]
C'est donc une sous-catégorie \textit{ouverte} de $\mathcal{M}$. Son complémentaire fermé $\underline{S}'=\underline{S}'_a$ est la sous-catégorie fermée image inverse, par $f$, de la sous-catégorie strictement pleine de Ens réduite au seul objet $\{\emptyset\}$. Cette sous-catégorie de Ens est la catégorie finale type. Elle est stable dans Ens par $\limind$ quelconques, et par $\limproj$ filtrantes quelconques. Ainsi, $\underline{S}'$ est stable dans $\mathcal{M}$ par $\limind$ filtrantes quelconques. Ainsi, $\underline{S}'$ est stable dans $\mathcal{M}$ par $\limind$ filtrantes quelconques, en particulier, elle est $\mathrm{L}_\pi$-stable pour $\pi$ assez grand. Explicitons :

\msp{page 13}
\[
\begin{array}{cl}
\underbrace{\underline{S}'\ =\ \underline{S}'_a}_{}\ \subset\ \mathcal{M} & \\
\Big\|\ \text{\scriptsize déf} & \text{sous-catégorie fermée de } \mathcal{M}\\
f^{-1}(\{\emptyset\})\ = & \text{formée des } x\ \in\ \operatorname{Ob}\mathcal{M} \text{ tels}\\
 & \text{que } a(x)=\emptyset.
\end{array}
\]

Comme $\underline{S}'$, étant fermée dans $\mathcal{M}$, est stable par $\limind$ non vides et en particulier $\mathrm{L}_\pi$-stable pour $\pi$ grand, il y a plus de raison de s'attendre à ce que $\underline{S}'$ soit accessible dans $\mathcal{M}$ que pour son complémentaire $\underline{S}$, dont il n'est pas clair qu'il ait cette propriété de stabilité, mais seulement stabilité par $\limproj$ non vides, et en particulier par facteurs directs.

En fait, nous allons voir qu'il peut arriver que $\underline{S}$ et $\underline{S}'$ soient tous deux non accessibles dans $\mathcal{M}$ :

\textbf{Proposition 4.} \textit{Soient $\mathfrak{U}$ un univers, contenant un cardinal infini $\pi_0$ tel que pour tout $E\in\mathfrak{U}$, tout $\pi_0$-ultrafiltre sur $E$ soit trivial. Soit alors $k\in\mathfrak{U}$ un corps de cardinal $\pi'_0\geqslant\pi_0$, et soit $\mathcal{M}$ la catégorie}
\printedpage{197}
% head: A. Grothendieck : Dérivateurs, Ch. XVIII, Édition des Universités de Montpellier II et Paris VII
\msp{page 14}
\textit{des $k$-algèbres commutatives. Donc $\mathcal{M}$ est un $\mathfrak{U}$-paratopos, admettant $k$ comme objet initial. Soit $\underline{S}'$ la sous-catégorie strictement pleine, stable par petites $\limind$ non vides, formée des objets $A$ de $\mathcal{M}$ tels que $\operatorname{Hom}_{\mathcal{M}}(A,k)=\emptyset$, et $\underline{S}$ la sous-catégorie strictement pleine complémentaire, formée des objets $A$ tels que $\operatorname{Hom}_{\mathcal{M}}(A,k)=\emptyset$. Alors $\underline{S}'$ est une sous-catégorie} fermée \textit{dans $\mathcal{M}$, et $\underline{S}$ la sous-catégorie ouverte complémentaire, donc $\underline{S}$ (tout comme $\underline{S}'$) est stable par facteurs directs. Ceci posé :} % [sic] second $\operatorname{Hom}_{\mathcal{M}}(A,k)=\emptyset$ -> $\operatorname{Hom}_{\mathcal{M}}(A,k)\neq\emptyset$, $\underline{S}$ is the complement of $\underline{S}'$
\begin{enumerate}[label={}]
\item[\textit{a)}] \textit{Il n'existe pas de cardinal $\pi$ tel que $\underline{S}$ soit $L_\pi$-stable dans $\mathcal{M}$, a fortiori $\underline{S}$ n'est pas accessible dans $\mathcal{M}$.}
\item[\textit{b)}] (\footnote[207]{?, pas prouvé.}). \textit{Bien que $\underline{S}'$ soit $L_0$-stable dans $\mathcal{M}$ (i.e. stable par $\limind$ filtrantes, et même, on l'a dit, par toute $\limind$ non vide), $\underline{S}'$ n'est pourtant pas accessible dans $\mathcal{M}$.}
\end{enumerate}
\noindent\textsc{Démonstration}. L'assertion a) va résulter du

\textbf{Corollaire 1.} \textit{Sous les conditions de la proposition 4, soient $E\in\mathfrak{U}$ un ensemble infini, $\Phi_0$ le filtre des complémentaires des parties finies de $E$. Alors $\operatorname{Hom}_{\mathcal{M}}(k_{\Phi_0},k)=\emptyset$, a fortiori, pour tout filtre $\Phi$ sur $E$ plus fin que $\Phi_0$, on a $\operatorname{Hom}_{\mathcal{M}}(k_\Phi,k)=\emptyset$.}

\msp{page 15}
On a posé ici, bien sûr, pour tout filtre $\Phi$ sur $E$,
\[
k_\Phi\ =\ \limind_{J\in\Phi}k_J\ .
\]
On sait que pour toute $k$-algèbre commutative $\mathcal{A}$, $\operatorname{Hom}(\mathcal{A},k)$ est en correspondance 1-1 avec la partie de $\operatorname{Max}(\mathcal{A})$ formée des idéaux maximaux $\mathfrak{m}$ de $\mathcal{A}$ tels que $k\longrightarrow\mathcal{A}/\mathfrak{m}$ soit un isomorphisme. Si $\mathcal{A}$ est de la forme $k_\Phi$, $\Phi$ un filtre sur $E$, cet ensemble est en correspondance biunivoque avec l'ensemble des $\pi'_0$-ultrafiltres $\Psi$ sur $E$ plus fins que $\Phi$, où $\pi'_0=\operatorname{card}(k)$ (proposition 3, et cor.~3 de prop.~2, page 7). Comme $\pi'_0\geqslant\pi_0$, un tel ultrafiltre est trivial par hypothèse. Donc on trouve
\[
\operatorname{Hom}_{\mathcal{M}}(k_\Phi,k)\ \simeq\ \text{ensemble des } \xi\in E \text{ tels que } \{\xi\}=\textstyle\bigcap_{A\in\Phi}A\ . \tag*{$(*)$}
\]
Mais dire que $\Phi$ est plus fin que $\Phi_0$ signifie aussi que pour tout $\xi\in E$, on a $E\backslash\{\xi\}\in\Phi$, i.e. $\exists\,A\in\Phi$ avec $A\subset E\backslash\{\xi\}$, i.e. $\{\xi\}\notin\bigcap_{A\in\Phi}A$. Donc on a alors $\operatorname{Hom}_{\mathcal{M}}(k_\Phi,k)=\emptyset$.

La conclusion du corollaire 1 s'exprime par la relation
\[
k_\Phi\ \notin\ \operatorname{Ob}\underline{S}\ ,\qquad\quad \text{i.e. } k_\Phi\in\operatorname{Ob}\underline{S}'\ ,
\]
d'autre part, on a par définition
\[
k_\Phi\ =\ \limind_{J\in\Phi}k_J\ ,
\]
\msp{page 16}
et il est clair (comme ces $J\in\Phi$ sont $\neq\emptyset$), que les $k_J$ en question sont dans $\underline{S}$. Ainsi, pour un filtre $\Phi$ comme dans le corollaire 1, on trouve un système inductif d'objets de $\underline{S}$, indexé par $\Phi^o$, dont la limite inductive n'est pas dans $\underline{S}$. Donc le a) de la proposition 4 se déduit en notant que pour tout cardinal $\pi\in\mathfrak{U}$, on peut trouver un ensemble $E\in\mathfrak{U}$, et un filtre $\Phi$ sur $E$ plus fin que $\Phi_0$, tel que l'ensemble ordonné $\Phi^o$ soit grand devant $\pi$. Pour trouver un tel $\Phi$, il suffit de prendre $E$ de cardinal $\pi'>\pi$, et pour $\Phi$ l'ensemble
\printedpage{198}
% head: A. Grothendieck : Dérivateurs, Ch. XVIII, Édition des Universités de Montpellier II et Paris VII
des complémentaires des parties de $E$ de cardinal $\leqslant\pi$, de sorte que $\Phi^o$ est isomorphe à l'ensemble ordonné des parties de $E$ de cardinal $\leqslant\pi$, lequel est bien grand devant $\pi$. Cela prouve le a) de la proposition 4.

Prouvons maintenant le b). Notons le

\textbf{Lemme.} \textit{Soit $\mathcal{M}$ une catégorie accessible stable par $\limind$ de type $I_\alpha$, $\alpha$ un ordinal limite donné, $\underline{S}'$ une sous-catégorie fermée, $\pi$ un cardinal utile}
\msp{page 17}
\textit{pour $\mathcal{M}$, tel que} $\operatorname{card}\alpha\leqslant\pi$\textit{. Pour tout $x\in\operatorname{Ob}\underline{S}'$, choisissons une représentation $\omega$-$\pi$-adaptée ($\omega$ étant le premier ordinal infini)}
\[
x\ \simeq\ \limind_{i\in I_x}x_i
\]
\textit{(donc $I_x$ grande devant $\pi$, stable par $\limind$ de type $I_\alpha$, $i\longmapsto x_i:I\longrightarrow\mathcal{M}$ commute auxdites limites, enfin les $x_i\in\operatorname{Ob}\mathcal{M}_\pi$). Pour que $\underline{S}'$ soit $\pi$-accessible dans $\mathcal{M}$, il faut et il suffit que pour tout $x\in\operatorname{Ob}\underline{S}'$, il existe un $i\in I_x$ tel que $x_i\in\operatorname{Ob}\underline{S}'$.}

En effet, notons que $\underline{S}'$ étant fermée, est stable dans $\mathcal{M}$ par toute $\limind$ non vide représentable dans $\mathcal{M}$, et en particulier par $\limind$ de type $I_\alpha$. Donc on peut appliquer 4.13.22 (p.~342) à $\mathcal{M}$, $\underline{S}'$, $\pi$ et à $\Phi=\{I_\omega\}$. Alors a) $\Longrightarrow$ c) implique que si $\underline{S}'$ est $\pi$-accessible dans $\mathcal{M}$, alors pour tout $x\in\operatorname{Ob}\underline{S}'$, l'ensemble $I_{x,S'}$ des $i\in I_x$ tels que $x_i\in\operatorname{Ob}\underline{S}$ est cofinal dans $I_x$. Cela implique qu'il est non vide. (Et réciproquement d'ailleurs, $\underline{S}'$ étant un fermé dans $\mathcal{M}$, puisque si $x_i\in\underline{S}'$, alors pour $j\in\operatorname{Ob}I_x$ avec $i\leqslant j$, on a $x_j\in\underline{S}'$, i.e. $I_{x,\underline{S}'}$ contient l'image essentielle de $i\backslash I_x$, laquelle est cofinale dans $I$ \dots) Cela prouve le `il faut'. Prouvons le `il suffit'. Donc prouvons que $\underline{S}'$ est égale à l'image essentielle de $\operatorname{Ind}_\pi(\underline{S}'(\pi))$. Elle la contient, puisque $\underline{S}'$ est stable par $\limind$ de type $\mathrm{L}_\pi$. % Doubt resolved against PDF leaf 40: the first S-prime subscript is not underlined, retained. % [sic] $x_i\in\operatorname{Ob}\underline{S}$ -> $x_i\in\operatorname{Ob}\underline{S}'$, the set $I_{x,\underline{S}'}$ concerns the $x_i$ in $\underline{S}'$

\msp{page 18}
Donc il reste à prouver que tout $x\in\operatorname{Ob}\underline{S}'$ est dans l'image essentielle. Or l'hypothèse nous assure que $x\simeq\limind\limits_{i\in I_x}x_i$ est isomorphe à $\limind\limits_{i\in I_{x,\underline{S}'}}x_i$ (car $I_{x,\underline{S}'}$ est cofinale dans $I_x$), où cette fois les $x_i$ sont dans $\underline{S}'$. Comme $I_x$ grande devant $\pi$, de même la sous-catégorie pleine cofinale $I_{x,\underline{S}'}$, donc on gagne.

\textbf{Corollaire.} \textit{Sous les conditions précédentes, si $\underline{S}'$ est $\pi$-accessible, il en est de même de l'ouvert complémentaire $\underline{S}$,} pourvu \textit{que $\underline{S}$ soit stable dans $\mathcal{M}$ par $\limind$ de type $L_\pi$.}

Il faut prouver que $\underline{S}$ est égale à l'image essentielle de $\operatorname{Ind}_\pi\underline{S}(\pi)$. L'hypothèse `pourvu' assure que cette image essentielle est $\subset\underline{S}$. Inversement, tout [\,?] $x\in\operatorname{Ob}\underline{S}$ est dans cette image essentielle. On aura
\[
x\ \simeq\ \limind_{i\in I}x_i\ ,\qquad\quad x_i\in\operatorname{Ob}\mathcal{M}_\pi,\ I=\mathcal{M}_\pi/x,\ \text{grande devant } \pi\ ,
\]
et comme $x\in\operatorname{Ob}\underline{S}$ et $\underline{S}$ ouverte, les $x_i$ sont dans $\underline{S}$, donc dans $\underline{S}(\pi)$, OK.

Dans le cas qui nous occupe, où $\mathcal{M}=(k\text{-alg.})$, comme $\underline{S}$ n'est justement \textit{pas} stable dans $\mathcal{M}$ par $\mathrm{L}_\pi$ (quel que soit $\pi$) (c'est cela qui nous faisait conclure que
\msp{page 19}
$\underline{S}$ n'est pas accessible dans $\mathcal{M}$), on ne peut donc conclure de cette inaccessibilité de $\underline{S}$, celle de $\underline{S}'$. Notons le


\endgroup
% END UNIT der-ch18/der-ch18_p196-200.tex
% BEGIN UNIT der-ch18/der-ch18_p201-201.tex
\begingroup\makeatletter%
\newcommand{\printedpage}[1]{\marginpar{\footnotesize p.~#1}}%
\newcommand{\msp}[1]{\marginpar{\footnotesize [#1]}}%
\newcommand{\Hom}{\operatorname{Hom}}%
\newcommand{\card}{\operatorname{card}}%
\makeatother\let\printedpage\derpp\let\msp\dermsp%
\let\mspin\DERmspin%


\printedpage{199}
% head: A. Grothendieck : Dérivateurs, Ch. XVIII, Édition des Universités de Montpellier II et Paris VII
% foot: Version 17 décembre 2025

\noindent\textbf{Corollaire.} \textit{Soit $\pi$ un cardinal infini. Pour que la sous-catégorie fermée $\underline{S}'$ de $\mathcal{M}$ (formée de $k$-algèbres commutatives $A$ telles que $\Hom_{\mathcal{M}}(A,k)=\emptyset$) soit $\pi$-accessible, il faut et il suffit que pour tout $k$-algèbre $A$ (de cardinal ${}>\pi$), telle que $\Hom_{\mathcal{M}}(A,k)=\emptyset$, existe une sous-algèbre $A_i$} de cardinal $\leqslant\pi$, \textit{telle que $\Hom_{\mathcal{M}}(A_i,k)=\emptyset$.} % [sic] tout k-algèbre -> toute k-algèbre, « algèbre » est féminin

\noindent(On applique le lemme, en prenant pour $\alpha$ le premier ordinal infini.)

\noindent\textbf{Question.} \textit{Cela est-il vérifié, en prenant simplement $\pi\geqslant\card k$~?} La question de l'existence d'un cardinal $\pi$, ayant la propriété indiquée plus haut, me paraît délicate, et j'ignore la réponse. Si on veut construire un contre-exemple pour tout $\pi$, les développements précédents suggèrent de prendre pour $A$ un produit infini d'extensions non triviales de $k$, ou un anneau $k_{\Phi}$ comme dans la démonstration de la proposition 4 a).


\endgroup
% END UNIT der-ch18/der-ch18_p201-201.tex
\clearpage\setcounter{footnote}{0}\phantomsection\addcontentsline{toc}{section}{Chapitre XIX}\markboth{Chapitre XIX}{Chapitre XIX}
% BEGIN UNIT der-ch19/der-ch19_p001-005.tex
\begingroup\makeatletter%
\newcommand{\printedpage}[1]{\marginpar{\footnotesize p.~#1}}%
\newcommand{\msp}[1]{\marginpar{\footnotesize [#1]}}%
\newcommand{\Qo}{\mathrel{\leftharpoondown\joinrel\relbar\joinrel\rightharpoonup}}   % lower half-head left, upper half-head right
\newcommand{\Qorev}{\mathrel{\leftharpoonup\joinrel\relbar\joinrel\rightharpoondown}} % mirrored form of (1.2)
\newcommand{\Qog}{\mathrel{\leftarrow\joinrel\relbar\joinrel\rightharpoonup}}         % full head left
\newcommand{\Qod}{\mathrel{\leftharpoondown\joinrel\relbar\joinrel\rightarrow}}       % full head right
\newcommand{\Qogd}{\mathrel{\leftarrow\joinrel\relbar\joinrel\rightarrow}}            % full heads both sides
\newcommand{\Qx}[2]{\mathrel{#1\joinrel\mapstochar#2}}%
\newcommand{\Qxx}[2]{\mathrel{#1\joinrel\mapstochar\mkern2mu\mapstochar#2}}%
\newcommand{\tld}{\tilde{}}   % left Q-closure mark (small tilde in exponent)
\newcommand{\rcl}{\bar{}}       % right Q-closure mark (horizontal bar in exponent)
\makeatother\let\printedpage\derpp\let\msp\dermsp%
\let\mspin\DERmspin%


\printedpage{(non numérotée)}
% head: (none; title page)

\begin{center}
{\Large\textsf{LES DÉRIVATEURS}}

\bigskip

\textbf{ALEXANDRE GROTHENDIECK}

\bigskip\bigskip

\textbf{Chapitre XIX}

\medskip

\textbf{Modèles (4)}
\end{center}

\par

\printedpage{(non numérotée)}
% head: (none; editorial note)

\par\noindent\textbf{Note des éditeurs de l’édition transcrite.}\par
Ce texte a été déchiffré et transcrit en \LaTeXe\ par M. Künzer. Il a été édité par M. Künzer, J. Malgoire, et G. Maltsiniotis. La transcription est aussi fidèle que possible au manuscrit. Seules quelques corrections mineures évidentes ont été effectuées. Pour les rares ajouts ou commentaires des éditeurs, ainsi que pour la numérotation originale des pages du manuscrit, les caractères de machine à écrire \texttt{[typewriter]} entre crochets sont utilisés.

Cette édition est provisoire. Les remarques, commentaires et corrections sont bienvenus. Envoyer un message à :

\textbf{[e-mail address removed]}

\hfill G. Maltsiniotis\hspace*{2cm}

\par

\printedpage{1}
% head: A. Grothendieck : Dérivateurs, Ch. XIX, Édition des Universités de Montpellier II et Paris VII
% foot: Version 11 février 2019

\msp{page 1}

\section*{1\quad Rappels sur les couples de Quillen (couples de Quillen spéciaux)}

Je vais reprendre les notations et la terminologie pour les relations de LLP \texttt{[left lifting property]} et RLP \texttt{[right lifting property]} de Quillen. J'écris (pour une catégorie $\mathcal{M}$, et $\Phi, \Psi \subset \mathrm{Fl}(\mathcal{M})$)
\[
\Phi \Qo \Psi \qquad \iff \qquad (\Phi \subset \Psi^{*}) \qquad \iff \qquad (\Psi \subset \Phi_{*})\ ,
\tag*{(1.1)}
\]
i.e. $\Phi$ a la LLP par rapport à $\Psi$, ou encore $\Psi$ la RLP par rapport à $\Phi$. Je dis aussi que $(\Phi, \Psi)$ est un \textit{couple} Q-\textit{orthogonal.} Cette notion n'est pas symétrique, mais je me permets parfois (exceptionnellement) d'écrire
\[
\Psi \Qorev \Phi \qquad \text{comme synonyme de} \qquad \Phi \Qo \Psi\ .
\tag*{(1.2)}
\]
$\Phi$ est Q-\textit{orthogonal à gauche} de $\Psi$, $\Psi$ est Q-\textit{orthogonal à droite} à $\Phi$, et la Q-orthogonalité à gauche (resp. à droite) est notée par le demi-sommet de flèche pointant vers le bas (resp. vers le haut). J'écris
\[
\left\{
\begin{aligned}
&\Phi \quad \Qog \quad \Psi \quad \iff \quad (\Phi = \Psi^{*})\\
&\Phi \quad \Qod \quad \Psi \quad \iff \quad (\Psi = \Phi_{*})\\
&\Phi \quad \Qogd \quad \Psi \quad \iff \quad (\Phi = \Psi^{*} \ \text{ et } \ \Psi = \Phi_{*})\ .
\end{aligned}
\right.
\tag*{(1.3)}
\]
Ces trois cas de Q-orthogonalité renforcée sont donc ceux où on a de plus $\Phi = \Phi^{\tld}$
resp. $\Psi = \Psi^{\rcl}$ resp. $\Phi = \Phi^{\tld}$ et $\Psi = \Psi^{\rcl}$. On dit que $\Phi$ est Q-\textit{clos à gauche} si $\Phi = \Phi^{\tld}$, que % reading resolved: DER_XIX_corrections.csv
$\Psi$ est Q-\textit{clos à droite} si $\Psi = \Psi^{\rcl}$, que le couple Q-orthogonal $(\Phi, \Psi)$ est Q-\textit{clos à gauche} (resp. \textit{à droite}) si $\Phi$ (resp. $\Psi$) est Q-clos à gauche (resp. à droite). On dit qu'il est Q-\textit{clos} si les deux conditions % reading resolved: DER_XIX_corrections.csv

\msp{page 2}

sont satisfaites simultanément. Quand aucune confusion n'est à craindre, il m'arrivera de dire simplement que $\Phi$ est Q-clos, ou que $\Psi$ est Q-clos, sans spécifier si c'est à gauche où % [sic] où -> ou, conjunction "or" intended
à droite. Mais pour un couple Q-orthogonal $(\Phi, \Psi)$ (ou je préfère l'écrire $\Phi \Qo \Psi$), Q-clos % [sic] ou -> où, relative adverb intended
signifie toujours Q-clos à gauche et à droite.

Quand un couple Q-orthogonal $\Phi \Qo \Psi$ a la propriété de factorisation de Quillen, je l'indique par la barre verticale croisant l'horizontale, ainsi
\[
\Phi \ \Qx{\leftharpoondown}{\rightharpoonup} \ \Psi\ ,
\tag*{(1.4)}
\]
et on définit de même le sens des graphismes
\[
\left\{
\begin{aligned}
&\Phi \quad \Qx{\leftharpoondown}{\rightharpoonup} \quad \Psi\\
&\Phi \quad \Qx{\leftarrow}{\rightharpoonup} \quad \Psi\\
&\Phi \quad \Qx{\leftharpoondown}{\rightarrow} \quad \Psi\\
&\Phi \quad \Qx{\leftarrow}{\rightarrow} \quad \Psi\ .
\end{aligned}
\right.
\tag*{(1.3)} % [sic] (1.3) -> (1.5), number already used above and (1.5) is missing
\]

\par

\printedpage{2}
% head: A. Grothendieck : Dérivateurs, Ch. XIX, Édition des Universités de Montpellier II et Paris VII
% foot: Version 11 février 2019

Un couple Q-\textit{orthogonal} Q-\textit{factorisant} $\Phi \Qx{\leftharpoondown}{\rightharpoonup} \Psi$ est appelé un \textit{précouple de Quillen.} S'il est de plus Q-clos, i.e. si on a $\Phi \Qx{\leftarrow}{\rightarrow} \Psi$, on dit que c'est un \textit{couple de Quillen.}

$\Phi$ est dit Q-\textit{saturé à gauche} s'il contient les isomorphismes et est stable par cochangement de base, par (petite) composition transfinie, et par facteurs directs. Dualement $\Psi$ est dit Q-\textit{saturé à droite} s'il

\msp{page 3}

contient les isomorphismes, est stable par changement de base, par composition transfinie, et par facteurs directs. Si $\Phi$ est Q-saturé à gauche, il est stable par (petites) sommes directes, i.e. si les $f_i : A_i \longrightarrow B_i$ ($i \in I$, $I$ petit) sont dans $\Phi$, $f = \coprod f_i : \coprod A_i \longrightarrow \coprod B_i$ aussi, et de même pour (petites) sommes amalgamées finies ou infinies (quand on se place dans un $A\backslash\mathcal{M}$ au lieu de $\mathcal{M}$). On a l'énoncé dual pour $\Psi$, quand $\Psi$ est Q-saturé à droite.

En fait, il y a des notions plus faibles de Q-saturation que celle que je viens de dire : Q-\textit{saturation simple} (à gauche : $\Phi$ contient les isomorphismes, stable par composition binaire, par cochangement de base), Q-\textit{saturation transfinie} (on exige de plus stabilité par composition transfinie), enfin Q-\textit{saturation totale,} ou Q-saturation tout court (quand on exige de plus la stabilité par facteurs directs). On a les implications
\[
\left\{
\begin{aligned}
&\text{Q-clos à gauche}\\
&\Downarrow\\
&\text{Q-saturé (total)}\\
&\Downarrow\\
&\text{Q-saturé transfini}\\
&\Downarrow\\
&\text{Q-saturé simple.}
\end{aligned}
\right.
\tag*{(1.6)}
\]
Pour $\Phi \subset \mathrm{Fl}(\mathcal{M})$, on a sa Q-clôture (à gauche) $\Phi^{\tld}$ (qui est la plus petite partie Q-close à gauche de $\mathrm{Fl}(\mathcal{M})$ contenant $\Phi$) et de même son Q-saturé (total) à gauche, son Q-saturé transfini, son Q-saturé simple

\msp{page 4}

avec les inclusions tautologiques
\[
\Phi \ \subset \
\underbrace{\Phi'}_{\substack{\text{Q-saturé}\\ \text{à gauche simple}}}
\ \subset \
\underbrace{\Phi''}_{\substack{\text{Q-saturé}\\ \text{à gauche transfini}}}
\ \subset \
\underbrace{\Phi'''}_{\substack{\text{Q-saturé}\\ \text{à gauche total}}}
\ \subset \
\underbrace{\Phi^{\tld}}_{\substack{\text{Q-clôture}\\ \text{à gauche}}}
\ ,
\tag*{(1.7)}
\]
et dualement pour $\Psi$. Entre $\Phi$ et $\Phi'$ il y a encore un $\Phi'_0$,
\[
\Phi \ \subset \ \Phi'_0 \ \subset \ \Phi'\ ,
\tag*{(1.8)}
\]
qui est le Q-\textit{saturé pour cochangement de base} (contenant aussi, par définition, tous les isomorphismes). Pour construire tous ces saturés, on procède ainsi :

\par

\printedpage{3}
% head: A. Grothendieck : Dérivateurs, Ch. XIX, Édition des Universités de Montpellier II et Paris VII
% foot: Version 11 février 2019

\begin{enumerate}[leftmargin=*]
\item[1)] $\Phi'_0 \quad = \quad$ ensemble des flèches $i : A \longrightarrow B$ s'insérant dans un carré cocartésien
\[
\begin{tikzcd}
A_0 \arrow[r, "i_0"] \arrow[d] & B_0 \arrow[d]\\
A \arrow[r, "i"] & B\ ,
\end{tikzcd}
\]
avec $i_0 \in \Phi$ ou $i_0$ une identité (auquel cas on trouve pour $i$ les isomorphismes).

\item[2)] $\Phi' \quad = \quad$ ensemble des flèches composées multiples
\[
i_n \, i_{n-1} \ \cdots \ i_1\ ,
\]
avec $i_\alpha \in \Phi'_0$.

\item[3)] $\Phi'' \quad = \quad$ ensemble des composés transfinis de flèches $\in \Phi'_0$.
\end{enumerate}

On a envie de poser
\[
\Phi''' \ = \ \Phi^{\prime\prime\natural} \ \overset{\text{déf}}{=} \ \text{ensemble des facteurs directs d'éléments de } \Phi'',
\]
et c'est un fait que $\Phi^{\prime\prime\natural} \subset \Phi'''$ trivialement, mais il n'est par clair en général qu'on ait % [sic] par -> pas, typo
\msp{page 5}
égalité, faute de pouvoir montrer que $\Phi^{\prime\prime\natural}$ est stable, non seulement par facteurs directs, mais aussi par cochangement de base et par composition transfinie. Cependant, rappelons :

\textbf{Théorème 1.1.} \textit{Supposons $\mathcal{M}$ stable par petites $\varinjlim$ , et que ses éléments soient accessibles. Supposons de plus que $\Phi$ soit petit, plus généralement, qu'il existe une partie $\Phi_0 \subset \Phi$ petite telle que $\Phi \subset \Phi_0{}^{\tld}$, donc}
\[
\Psi \ \overset{\text{déf}}{=} \ \Phi_{*} \ = \ (\Phi_0)_{*} \qquad \text{\footnotemark}.
\tag*{(1.1.1)} % footnote mark printed as superscript (1)
\]
\footnotetext{\textbf{N.B.} On a donc $\Phi \Qod \Psi$, a fortiori $\Phi_0 \Qod \Psi$.}
\textit{Sous ces conditions, si $\Phi''_0$ désigne le \textup{Q}-saturé à gauche transfini de $\Phi_0$, alors on a $\Phi''_0 \Qx{\leftharpoondown}{\rightarrow} \Psi$, i.e. toute flèche $f$ de $\mathcal{M}$ se factorise en $f = pi$, avec $i \in \Phi''_0$, $p \in \Psi$\,\footnote{Il y a même factorisation \textit{fonctorielle} en $f$, propriété qui mérite une notation, p.ex. $\Phi''_0 \Qxx{\leftharpoondown}{\rightarrow} \Psi$.}. A fortiori, si $\Phi''$ désigne le \textup{Q}-saturé à gauche transfini de $\Phi$, on a} % footnote mark printed as superscript (2)
\[
\Phi'' \ \Qx{\leftharpoondown}{\rightarrow} \ \Psi\ .
\]

\bigskip

\textbf{Corollaire 1.2.} \textit{Sous les conditions du théorème, on a}
\[
\Phi''' \ \text{(Q-saturé total à gauche de $\Phi$)} \ = \
\underbrace{\Phi^{\prime\prime\natural}}_{\substack{\text{ensemble des}\\ \text{facteurs directs d'éléments}\\ \text{de $\Phi''$}}}
\ = \ \Phi^{\tld}\ .
\tag*{(1.2.1)}
\]
\textit{Plus généralement (en apparence seulement!), si $\Phi_1$ est une partie de \textup{$\mathrm{Fl}(\mathcal{M})$} qui est \textup{Q}-saturée à gauche transfinie, et telle que}
\[
\Phi_0 \ \subset \ \Phi_1 \ \subset \ \Phi^{\tld} \ = \ \Phi_0{}^{\tld}\ ,
\]


\endgroup
% END UNIT der-ch19/der-ch19_p001-005.tex
% BEGIN UNIT der-ch19/der-ch19_p006-010.tex
\begingroup\makeatletter%
\derleqno%
\newcommand{\printedpage}[1]{\marginpar{\footnotesize p.~#1}}%
\newcommand{\msp}[1]{\marginpar{\footnotesize [#1]}}%
\newcommand{\Qorth}{\mathrel{\leftharpoondown\mkern-6mu\relbar\mkern-6mu\rightharpoonup}}%
\newcommand{\Qorthx}{\mathrel{\leftharpoondown\mkern-9mu{+}\mkern-9mu\rightharpoonup}}%
\newcommand{\Qorthxfull}{\mathrel{\leftharpoondown\mkern-9mu{+}\mkern-9mu\rightarrow}}%
\newcommand{\Qcouple}{\mathrel{\leftarrow\mkern-9mu{+}\mkern-9mu\rightarrow}}%
\newcommand{\clos}{\longleftrightarrow}%
\newcommand{\fl}{\longrightarrow}%
\makeatother\let\printedpage\derpp\let\msp\dermsp%
\let\mspin\DERmspin%


\printedpage{4}
% head: A. Grothendieck : Dérivateurs, Ch. XIX, Édition des Universités de Montpellier II et Paris VII
% foot: Version 11 février 2019

\msp{page 6}
\textit{alors}
\begin{enumerate}[label={}]
\item[a)] \textit{on a $\Phi_1 \Qorthxfull \Psi$ et} % reading resolved: DER_XIX_corrections.csv
\item[b)] $\Phi_1^{\sim} = \Phi^{\sim} = \Phi_1^{\natural}$.
\end{enumerate}
Ce corollaire résulte du théorème en vertu de XV prop. 4.2, dont je rappelle ici l'essentiel :

\textbf{Proposition 1.3.} \textit{Soit $\Phi \Qorthx \Psi$. Alors $\Phi^{\sim} = \Phi^{\natural} = \Psi^{*}$, et dualement on a $\Psi^{-} = \Psi^{\natural} = \Phi_{*}$.}

Notons aussi, dans cet ordre d'idées :

\textbf{Corollaire 1.4.} \textit{Soit $\Phi \Qorth \Psi$, et soit $W$ une partie quelconque de $\mathrm{Fl}(\mathcal{M})$. Pour qu'on ait $\Phi \subset W^{\natural}$, il suffit que toute flèche de $\mathcal{M}$ se factorise en $f = pi$, avec $i \in W$, $p \in \Psi$. (Et cette condition est évidemment aussi nécessaire si $\Phi \Qorthx \Psi$, et $W$ stable par facteurs directs.)}

On dit qu'une partie $\Phi_0$ de $\Phi$ est Q-\textit{génératrice à gauche} de $\Phi$ si elle satisfait la condition
\[
\Phi \ \subset \Phi_0^{\sim}\ , \tag*{(1.9)}
\]
et on dit que $\Phi$ est une partie \textit{spéciale} si $\Phi$ admet une \textit{petite} partie Q-génératrice à gauche. Si $\Phi_0 \subset \Phi$ est petit, et si $\mathcal{M}$ satisfait les conditions du théorème 1.1, alors $\Phi_0$ est une partie Q-génératrice (à gauche) si et seulement si son Q-saturé total (à gauche) $\Phi_0'''$ contient $\Phi$ (condition en apparence plus forte que $\Phi \subset \Phi_0^{\sim}$, puisque en général on a $\Phi_0''' \subset \Phi_0^{\sim}$). Cela signifie donc aussi, via 1.2, que tout \msp{page 7} élément de $\Phi$ est facteur direct d'un composé transfini de flèches déduites de
\[
\Phi_0 \cup (\text{Identités de } \mathcal{M})
\]
par cochangement de base.

Je rappelle encore le sorite $\pm$ tautologique des foncteurs adjoints :

\textbf{Proposition 1.5.} \textit{Soit}
\[
\begin{tikzcd}
\mathcal{M}' \arrow[r, shift right=1, "v^{*}"'] & \mathcal{M} \arrow[l, shift right=1, "v_{!}"']
\end{tikzcd}
\tag*{(1.5.1)}
\]
\textit{un couple de foncteurs adjoints, soient}
\[
\Phi \ \subset\ \mathrm{Fl}(\mathcal{M})\ , \qquad \Psi' \ \subset\ \mathrm{Fl}(\mathcal{M}')\ ,
\]
\textit{alors on a}
\[
(v_{!}(\Phi) \Qorth \Psi') \qquad \Longleftrightarrow \qquad (\Phi \Qorth v^{*}(\Psi'))\ . \tag*{(1.5.2)}
\]

\textbf{Corollaire 1.6.} \textit{Pour toute partie $\Phi$ de} $\mathrm{Fl}(\mathcal{M})$,
\[
(v_{!}(\Phi))_{*} \ =\ (v^{*})^{-1}(\Phi_{*})\ . \tag*{(1.6.1)}
\]

\printedpage{5}
% head: A. Grothendieck : Dérivateurs, Ch. XIX, Édition des Universités de Montpellier II et Paris VII
% foot: Version 11 février 2019

En effet, appliquant 1.5 à $\Phi$ et à $\{f'\}$, où $f' \in \mathrm{Fl}(\mathcal{M}')$, on trouve
\[
f' \in (v_{!}(\Phi))_{*} \qquad \Longleftrightarrow \qquad v^{*}(f') \in \Phi_{*}\ ,
\]
ce qui équivaut à (1.6.1). Dualement :

\textbf{Corollaire 1.7.} \textit{Pour toute partie $\Psi'$ de $\mathcal{M}'$, on a}
\[
(v^{*}(\Psi'))^{*} \ =\ (v_{!})^{-1}(\Psi'^{*}). \tag*{(1.7.1)}
\]

En effet, si $f \in \mathrm{Fl}(\mathcal{M})$, on a par (1.5.2)
\[
f \in (v^{*}(\Psi'))^{*} \qquad \Longleftrightarrow \qquad v_{!}(f) \in \Psi^{*}\ . % [sic] \Psi^{*} -> \Psi'^{*}, cf. (1.7.1)
\]
\textsuperscript{(}\footnote[3]{donner ici 1.9 page 9.}\textsuperscript{)}.

\msp{page 8}
Un couple Q-orthogonal $\Phi \Qorth \Psi$ est dit \textit{spécial} s'il satisfait les conditions suivantes.
\[
\left\{
\begin{array}{ll}
\text{a)} & \parbox[t]{0.7\textwidth}{$\mathcal{M}$ est stable par petites $\varinjlim$ , les objets de $\mathcal{M}$ sont accessibles. (On dit alors que $\mathcal{M}$ est “spéciale”.)} \\[3ex]
\text{b)} & \text{Le couple est clos : } \Phi \clos \Psi. \\[1ex]
\text{c)} & \Phi \text{ est \textit{spéciale.}}
\end{array}
\right.
\tag*{(1.10)}
\]
Cela implique donc, par 1.1, que $\Phi \Qcouple \Psi$ est un couple de Quillen. Si $\mathcal{M}$ satisfait a), les couples Q-orthogonaux spéciaux dans $\mathcal{M}$, ou encore les couples de Quillen spéciaux, correspondent donc biunivoquement aux parties $\Phi$ de $\mathrm{Fl}(\mathcal{M})$ qui sont \textit{spéciales} et closes à gauche (ou ce qui revient au même, Q-saturées à gauche), en associant à une telle $\Phi$ le couple $\Phi \Qcouple \Phi_{*}$.

Ceci posé :

\marginpar{$\Vert$}% double vertical bar printed in the left margin beside Corollaire 1.8
\textbf{Corollaire 1.8.} \textit{Supposons que $\mathcal{M}$, $\mathcal{M}'$ satisfassent la condition} a) \textit{de} (1.10), \textit{et soit $(\Phi, \Psi)$ un couple de Quillen spécial dans $\mathcal{M}$. Alors $\Bigl(\Phi' = v_{!}(\Phi)^{\sim}, \Psi' = (v^{*})^{-1}(\Psi)\Bigr)$ est un couple de Quillen spécial dans $\mathcal{M}'$.}

En effet, par le corollaire 1.6 on a
\[
\Phi'_{*} \ =\ (v_{!}(\Phi))_{*} \ =\ (v^{*})^{-1}(\Psi) \ =\ \Psi'\ ,
\]
donc $\Phi' \clos \Psi'$. Il reste à prouver que $\Phi'$ est spécial. Soit $\Phi_0$ une partie petite de $\Phi$ telle que $\Phi \subset \Phi_0^{\sim}$, posons $\Phi'_0 \ =\ v_{!}(\Phi_0) \subset \Phi'$, \msp{page 9} il suffit de voir qu'on a
\[
v_{!}(\Phi) \ \subset\ v_{!}(\Phi_0)^{\sim}\ , \tag*{($*$)}
\]
d'où résultera qu'on a $\Phi' \overset{\text{déf}}{=} (v_{!}(\Phi))^{\sim} \subset v_{!}(\Phi_0)^{\sim}$. Mais on a
\[
v_{!}(\Phi_0)^{\sim} \ =\ (v_{!}(\Phi_0)_{*})^{*}\ ,
\]
\printedpage{6}
% head: A. Grothendieck : Dérivateurs, Ch. XIX, Édition des Universités de Montpellier II et Paris VII
% foot: Version 11 février 2019
et
\[
v_{!}(\Phi_0)_{*} \ =\ (v^{*})^{-1}(\underbrace{\Phi_{0\,*}}_{\Psi}) \ =\ (v^{*})^{-1}(\Psi) \ \overset{\text{déf}}{=}\ \Psi'\ ,
\]
donc
\[
v_{!}(\Phi_0)^{\sim} \ =\ \Psi'^{*} \ =\ \Phi' \ =\ v_{!}(\Phi)^{\sim}\ ,
\]
d'où ($*$). Cet argument ne fait pas appel à l'hypothèse “spéciale” sur $\mathcal{M}$, et peut s'expliciter en le

\textbf{Corollaire 1.9.} \textsuperscript{(}\footnote[4]{devrait passer avant 1.8.}\textsuperscript{)} \textit{Sous les conditions générales de} 1.5 \textit{($v_{!} : \mathcal{M} \fl \mathcal{M}'$ étant un foncteur admettant un adjoint à droite), soient $\Phi_0$, $\Phi$ deux parties de} $\mathrm{Fl}(\mathcal{M})$ \textit{telles que $\Phi \subset \Phi_0^{\sim}$ (ou ce qui revient au même, $\Phi^{\sim} \subset \Phi_0^{\sim}$). Alors $v_{!}(\Phi) \subset v_{!}(\Phi_0)^{\sim}$, i.e. $v_{!}(\Phi)^{\sim} \subset v_{!}(\Phi_0)^{\sim}$.}

En effet, cette relation équivaut aussi à $v_{!}(\Phi_0)_{*} \ \subset\ v_{!}(\Phi)_{*}$, c'est-à-dire (par 1.6) $(v^{*})^{-1}(\Phi_{0\,*}) \ \subset\ (v^{*})^{-1}(\Phi_{*})$, et est impliqué par $\Phi_{0\,*} \ \subset\ \Phi_{*}$. Cette dernière d'autre part équivaut à l'hypothèse $\Phi^{\sim} \subset \Phi_0^{\sim}$, q.e.d.

\msp{page 10}
Si $\Psi$ est une partie de $\mathrm{Fl}(\mathcal{M})$, on en déduit pour toute $I$ dans Cat une partie
\[
\left\{
\begin{array}{l}
\Psi(I) \subset \mathrm{Fl}(\mathcal{M}(I)), \quad \text{avec}\ \ \mathcal{M}(I) = \underline{\mathrm{Hom}}(I^{\circ}, \mathcal{M})\ , \\[1ex]
f \in \Psi(I) \iff f(i) \in \Psi \ \ \forall\, i \in \mathrm{Ob}\, I\ .
\end{array}
\right.
\tag*{(1.11)}
\]
Si on introduit la catégorie discrète $I_0$ sous-jacente à $\mathrm{Ob}\, I$, et le foncteur canonique
\[
v_I = v \ :\ I_0 \fl I\ , \tag*{(1.12)}
\]
on trouve un couple de foncteurs adjoints
\[
\begin{tikzcd}[row sep=small]
\mathcal{M}(I) \arrow[r, shift right=1, "v^{*}"'] & \mathcal{M}(I_0) \arrow[l, shift right=1, "v_{!}"'] \\
( & ) \arrow[l, dotted, "v_{*}"]
\end{tikzcd}
\ \simeq\ \mathcal{M}^{I_0}\ ,
\tag*{(1.13)}
\]
et on a
\[
\Psi(I) \ =\ (v^{*})^{-1}(\Psi(I_0))\ .
\]
Supposons maintenant qu'on ait
\[
\left\{
\begin{array}{rcl}
\Psi & = & \Phi_{*} \\[1ex]
\text{d'où}\ \ \Psi(I_0) & = & \Phi(I_0)_{*}
\end{array}
\right.
\tag*{(1.14)}
\]
(vérification immédiate), on conclut alors par 1.6
\[
\Psi(I) \ =\ (v_{!}(\Phi(I_0)))_{*} \ =\ (\Phi\langle I\rangle)_{*}\ , \tag*{(1.15)}
\]
avec
\[
\Phi\langle I\rangle \ \overset{\text{déf}}{=}\ v_{!}(\Phi(I_0))^{\sim} \ =\ \Psi(I)^{*}\ . \tag*{(1.16)}
\]
Si on a un foncteur
\[
u \ :\ I \fl I'\ , \tag*{(1.17)}
\]
d'où un triple de foncteurs adjoints
\printedpage{7}
% head: A. Grothendieck : Dérivateurs, Ch. XIX, Édition des Universités de Montpellier II et Paris VII
% foot: Version 11 février 2019
\msp{page 11}
\[
\left\{
\begin{tikzcd}[column sep=large]
\mathcal{M}(I) \arrow[r, shift left=3, "u_{!}"] \arrow[r, shift right=3, "u_{*}"] & \mathcal{M}(I') \arrow[l, "u^{*}"']
\end{tikzcd}
\right.
,
\tag*{(1.18)}
\]
on aura de façon $\pm$ tautologique
\[
\left\{
\begin{array}{rcl}
u_{!}(\Phi\langle I\rangle) & \subset & \Phi\langle I'\rangle \\
u^{*}(\Psi(I')) & \subset & \Psi(I)\ .
\end{array}
\right.
\tag*{(1.19)}
\]
Le corollaire 1.8 donne :

\textbf{Proposition 1.10.} \textit{Supposons que $\Phi \Qcouple \Psi$ soit un couple de Quillen} spécial \textit{dans $\mathcal{M}$, i.e. de la forme $(\Phi, \Phi_{*})$, avec $\Phi$ partie Q-close à gauche spéciale dans $\mathcal{M}$ ($\mathcal{M}$ satisfaisant la condition} (1.10. a))\textit{. Alors pour toute $I$ dans} Cat, \textit{$(\Phi\langle I\rangle, \Psi(I))$ est un couple spécial dans $\mathcal{M}(I)$, en particulier il satisfait au théorème de factorisation.}

Considérons la construction duale, qui nous fait introduire, à partir d'un couple $\Phi \clos \Psi$ dans $\mathcal{M}$, le couple
\[
\Phi(I) \ \clos\ \Psi[I]\ , \tag*{(1.20)}
\]
où $\Phi(I)$ est défini par la formule générale (1.11) (où $\Psi$ est remplacé par $\Phi$), et où $\Psi[I]$ est défini soit comme $\Phi(I)_{*}$, soit comme
\[
\Psi[I] \ =\ v_{*}(\Psi(I_0))^{-} \ =\ \Phi(I)_{*}\ . \tag*{(1.21)}
\]

\msp{page 12}
On a alors, pour $u : I \fl I'$
\[
\left\{
\begin{array}{rcl}
u^{*}(\Phi(I')) & \subset & \Phi(I) \\
u_{*}(\Psi[I]) & \subset & \Psi[I']\ .
\end{array}
\right.
\tag*{(1.22)}
\]
Il se pose encore la question si le couple $(\Phi(I), \Psi[I])$ satisfait au théorème de factorisation, sous des hypothèses convenables sur $(\Phi, \Psi)$. Je peux seulement hasarder la

\textbf{Question-conjecture 1.11.} \textit{Supposons que $(\Phi, \Psi)$ soit un couple de Quillen} spécial, \textit{en est-il de même de $(\Phi(I), \Psi[I])$ pour toute $I$ dans} (Cat) \textit{? En d'autres termes, si $\mathcal{M}$ est spéciale, condition} (1.10. a))\textit{, et si $\Phi$ est une partie Q-close à gauche} spéciale \textit{de $\mathcal{M}$, i.e. admettant une petite partie Q-génératrice, est-il vrai que pour toute $I$ dans} Cat, \textit{$\Phi(I)$ est également spéciale ? Je conjecture que oui, en faisant au besoin des hypothèses supplémentaires sur $\mathcal{M}$ (existence d'une filtration cardinale ?). P.ex. si $\Phi_0$ est Q-génératrice dans $\Phi$, et petite, alors $\Phi_0(I)$ est-elle Q-génératrice dans $\Phi(I)$ ?}

\textbf{Proposition 1.12} \textsuperscript{(}\footnote[5]{cf. XVII 6.4 (p. 117), mais où les notations pour $\Phi\langle I\rangle$, $\Psi(I)$ sont inverses.}\textsuperscript{)}. \textit{Soit $\Phi \clos \Psi$ dans $\mathcal{M}$, soit $u : I \fl I'$ une flèche dans} Cat, \textit{d'où}
\[
\left\{
\begin{tikzcd}[column sep=large]
\mathcal{M}(I) \arrow[r, shift left=3, "u_{!}"] \arrow[r, shift right=3, "u_{*}"] & \mathcal{M}(I') \arrow[l, "u^{*}"']
\end{tikzcd}
\right.
.
\tag*{(1.12.1)}
\]
\printedpage{8}
% head: A. Grothendieck : Dérivateurs, Ch. XIX, Édition des Universités de Montpellier II et Paris VII
% foot: Version 11 février 2019
\textit{On a alors} (1.19), \textit{i.e.}
\[
\left\{
\begin{array}{rcl}
u_{!}(\Phi\langle I\rangle) & \subset & \Phi\langle I'\rangle \\
u^{*}(\Psi(I')) & \subset & \Psi(I)\ ,
\end{array}
\right.
\tag*{(1.13.1)}
\]
\textit{et de plus l'équivalence}
\[
u^{*}(\Phi(I')) \subset \Phi\langle I'\rangle \qquad \Longleftrightarrow \qquad u_{*}(\Psi(I)) \subset \Psi(I')\ . \tag*{(1.13.2)}
\]
\textit{Ces relations équivalentes sont satisfaites dans chacun des trois cas}~\textsuperscript{(}\footnote[6]{\textbf{N.B.} On a $1^{\circ} \Longrightarrow 3^{\circ}$, $2^{\circ} \Longrightarrow 3^{\circ}$.}\textsuperscript{)} \textit{:}
\begin{enumerate}[label={}]
\item[1°)] \textit{$u$ est coétale.}
\item[2°)] \textit{$I$ est discrète.}
\item[3°)] $\forall\, i' \in \mathrm{Ob}\, I'$\textit{, toute composante connexe de $I/i'$ admet un objet final.}
\end{enumerate}
\msp{page 13}

\section*{2\quad Rappels sur les triples de Quillen clos et les triples de Quillen spéciaux}

Soit $(\mathcal{M}, W)$ une \textit{catégorie de modèles,} i.e. $\mathcal{M}$ munie d'un localiseur $W \subset \mathrm{Fl}(\mathcal{M})$, i.e. une partie satisfaisant
\[
\parbox{0.75\textwidth}{$W$ contient les isomorphismes, et si $f$, $g$ sont des flèches composables, si deux parmi $f$, $g$, $gf$ sont dans $W$, la troisième aussi \textit{(saturation faible).}}
\tag*{(2.1)}
\]
Je rappelle que la donnée d'un \textit{triple de Quillen clos} $(W, C, F)$\textit{,} ou d'une structure de Quillen close, sur $\mathcal{M}$, ayant $W$ comme ensemble de quasi-isomorphismes, revient à la donnée de deux couples de Quillen
\[
C \Qcouple TF\ , \qquad TC \ \Qcouple\ F\ , \tag*{(2.2)}
\]
satisfaisant
\[
TC \ \subset\ C \cap W\ , \qquad TF \ \subset\ F \cap W\ , \tag*{(2.3)}
\]
plus une des trois conditions équivalentes :
\[
\left\{
\begin{array}{lrcll}
\text{a)} & TC & = & C \cap W & \qquad (\text{i.e. } C \cap W \subset TC) \\[1ex]
\text{b)} & TF & = & F \cap W & \qquad (\text{i.e. } F \cap W \subset TF) \\[1ex]
\text{c)} & W & = & TF \circ TC\ . &
\end{array}
\right.
\tag*{(2.4)}
\]
Rappelons aussi que ceci implique que $W$ \textit{est fortement saturée,} et a fortiori que $W = W^{\natural}$, i.e. $W$ stable par facteurs directs.

Rappelons aussi qu'en vertu de (1.4), on a :


\endgroup
% END UNIT der-ch19/der-ch19_p006-010.tex
% BEGIN UNIT der-ch19/der-ch19_p011-015.tex
\begingroup\makeatletter%
\newcommand{\printedpage}[1]{\marginpar{\footnotesize p.~#1}}%
\newcommand{\msp}[1]{\marginpar{\footnotesize [#1]}}%
\newcommand{\M}{\mathcal{M}}%
\newcommand{\Fl}{\mathrm{Fl}}%
\newcommand{\Ho}{\mathrm{Ho}}%
\newcommand{\cc}{\mathrm{c}}%
\newcommand{\ddarr}{\mathrel{\twoheadleftarrow\joinrel\twoheadrightarrow}}%
\newcommand{\harp}{\mathrel{\leftharpoondown\joinrel\relbar\joinrel\rightharpoonup}}%
\makeatletter%
\renewcommand\@makefnmark{\hbox{\@textsuperscript{\normalfont(\@thefnmark)}}}%
\renewcommand\@makefntext[1]{\noindent\hbox{\@textsuperscript{\normalfont\@thefnmark}}#1}%
\newcommand{\Qcouple}{\mathrel{\leftarrow\mkern-9mu{+}\mkern-9mu\rightarrow}}%
\makeatother%
\makeatother\let\printedpage\derpp\let\msp\dermsp%
\let\mspin\DERmspin%


\printedpage{9}
% head: A. Grothendieck : Dérivateurs, Ch. XIX, Édition des Universités de Montpellier II et Paris VII
% foot: Version 11 février 2019

\msp{page 14}

\begin{itemize}[leftmargin=4em]
\item[(2.5)] Étant donnés les deux couples (2.2), avec $TC \subset C$ (équivalent à $TF \subset F$), pour qu'on ait la première relation (2.3), i.e. $TC \subset W$, il faut, et il suffit si $W = W^{\natural}$, que toute flèche $f$ de $\M$ (ou seulement toute flèche de $TC$) se factorise en $f = pi$, avec $i \in W$, $p \in F$. Dualement, pour qu'on ait $TF \subset W$, il faut et il suffit si $W = W^{\natural}$, que toute flèche $f$ de $\M$ (ou seulement toute flèche de $TF$) se factorise en $f = pi$, avec $i \in W$, $f \in C$. % [sic] i \in W, f \in C -> i \in C, p \in W (?), dual of the preceding factorisation with p \in F
\end{itemize}

On dit que la structure de Quillen $(W,C,F)$ est \textit{spéciale} si les deux couples de Quillen (2.3) sont spéciaux, i.e. si les conditions suivantes sont satisfaites.

\begin{itemize}[leftmargin=4em]
\item[(2.6)]
\begin{itemize}
\item[a)] $\M$ est stable par petites $\varinjlim$, et ses éléments sont accessibles ($\M$ “spéciale”).
\item[b)] Il existe des \textit{petites} parties $C_0 \subset C$, $TC_0 \subset TC$ qui soient Q-génératrices à gauche, i.e. telles que l'on ait
\[
\left\{
\begin{array}{rclcrclcrcl}
C & \subset & C_0^{\sim}\ , & \text{i.e.} & C & = & C_0^{\sim}\ , & \text{i.e.} & TF & = & C_{0*}\ ;\\
TC & \subset & (TC_0)^{\sim}\ , & \text{i.e.} & TC & = & (TC_0)^{\sim}\ , & \text{i.e.} & F & = & (TC_0)_*\ .
\end{array}\right. % reading resolved: DER_XIX_corrections.csv
\]
\end{itemize}
\end{itemize}

Ainsi, la donnée d'un tel triple de Quillen \textit{spécial} dans $\M$, quand $\M$ satisfait (2.6. a), revient à la donnée des deux parties
\msp{page 15}
\[
TC \ \subset\ C
\tag*{(2.7)}
\]
de $\Fl(\M)$, satisfaisant les conditions :
\[
\left\{
\begin{array}{ll}
C,\ TC\ \text{“spéciales”} &
\left\{\begin{array}{ll}
\text{a)} & \text{$C$, $TC$ sont Q-closes à gauche.}\\[1ex]
\text{b)} & \parbox[t]{0.45\textwidth}{$C$, $TC$ sont Q-engendrées à gauche par des \textit{petites} parties $C_0$, $TC_0$ de $C$, $TC$ respectivement.}
\end{array}\right.\\[5ex]
 & \quad\ \begin{array}{ll}
\text{c)} & TC \subset W,\ TF \overset{\text{déf}}{=} C_* \subset W.\\[1ex]
\text{d)} & \parbox[t]{0.45\textwidth}{On a l'une des trois conditions équivalentes supplémentaires de (2.4), i.e. :}\\[5ex]
 & \begin{array}{lrcllrcl}
\alpha) & TC & = & C \cap W\ , & \text{i.e.} & C \cap W & \subset & TC\ .\\
\beta) & TF & = & F \cap W\ , & \text{i.e.} & F \cap W & \subset & TF\ .\\
\gamma) & W & = & TF \circ TC\ . & & & &
\end{array}
\end{array}
\end{array}\right.
\tag*{(2.8)}
\]

\printedpage{10}
% head: A. Grothendieck : Dérivateurs, Ch. XIX, Édition des Universités de Montpellier II et Paris VII
% foot: Version 11 février 2019

\textbf{Proposition 2.1.} \textit{Soit}
\[
\begin{tikzcd}
\M' \arrow[r, shift right, "v^*"'] & \M \arrow[l, shift right, "v_!"']
\end{tikzcd}
\tag*{(2.11)} % [sic] (2.11) -> (2.1.1), numbering sequence (2.1.2), (2.1.3), ...
\]
\textit{un couple de foncteurs adjoints, $\M$ étant muni d'une structure de Quillen close} $(W,C,F)$. \textit{Posons}
\[
\left\{
\begin{array}{rcl}
W' & = & (v^*)^{-1}(W)\\
F' & = & (v^*)^{-1}(F)\\
TF' & = & F' \cap W' \ = \ (v^*)^{-1}(TF)\\
C' & = & v_!(C)^{\sim}\\
TC' & = & v_!(TC)^{\sim}\ .
\end{array}\right.
\tag*{(2.1.2)}
\]

\begin{itemize}
\item[1°)] \textit{On a alors}
\[
\left\{
\begin{array}{rclcrcl}
C' & \longleftrightarrow & TF'\ , & & TC' & \longleftrightarrow & F'\ ,\\
TC' & \subset & C'\ , & & TF' & \subset & F'\ .
\end{array}\right.
\tag*{(2.1.3)}
\]

\textit{Pour que les données} (2.1.2) \textit{dans $\M'$ y correspondent à un triple de Quillen $(W',C',F')$, il faut et il suffit que les conditions suivantes soient satisfaites :}
\msp{page 16}
\[
\left\{
\begin{array}{ll}
\text{a)} & \parbox[t]{0.7\textwidth}{\textit{Factorisation pour $C' \longleftrightarrow TF'$ et $TC' \longleftrightarrow F'$, i.e. les relations $C' \Qcouple TF'$ et $TC' \Qcouple F'$,}}\\[4ex]
\text{b)} & TC' \subset W',
\end{array}\right. % reading resolved: DER_XIX_corrections.csv
\tag*{(2.1.4)}
\]
\textit{ou encore, au lieu de} b)\textit{, la condition équivalente} (en présence de a)) :
\[
\left\{
\begin{array}{ll}
\text{b$'$)} & \parbox[t]{0.7\textwidth}{\textit{Toute flèche $f'$ dans $\M'$ (non seulement $f' \in TC'$) se factorise en $f' = p'i'$, avec $i' \in W'$, $p' \in F'$.}}
\end{array}\right.
\tag*{(2.1.5)}
\]
\item[2°)] \textit{Supposons} (2.1.4. a) \textit{satisfaite (factorisation). Alors, pour que $(W',C',F')$ soit un triple de Quillen clos, il} suffit \textit{que l'on ait}
\[
\text{b$''$)}\qquad \boxed{\ v^*v_!(TC)\ \subset\ TC\ }\ ,
\tag*{(2.1.6)}
\]
\textit{et que de plus on soit dans l'un des trois cas suivants :}
\begin{itemize}
\item[A)] \textit{$v^*$ admet un adjoint à droite $v_*$.}
\item[B)] \textit{$v^*$ commute aux petites} $\varinjlim$\textit{, $(W,C,F)$ est spécial, et $\M'$ est spéciale, i.e. satisfait à la condition} (2.6. a).
\item[C)] \textit{$v^*$ est pleinement fidèle, et $v^*v_!(F) \subset F$.}
\end{itemize}
\end{itemize}
Notons tout de suite le

\textbf{Corollaire 2.2.} \textit{Supposons que $(W,C,F)$ soit un triple de Quillen spécial, que $\M'$ soit spéciale, que $v^*$ commute aux petites} $\varinjlim$\textit{, enfin que l'on ait} (2.1.6)\textit{, i.e. $v^*v_!(TC) \subset TC$. Alors $(W',C',F')$ est un triple de Quillen spécial.}

\printedpage{11}
% head: A. Grothendieck : Dérivateurs, Ch. XIX, Édition des Universités de Montpellier II et Paris VII
% foot: Version 11 février 2019

\msp{page 17}
En effet, en vertu du 2.1. 2°, cas B, pour voir que $(W',C',F')$ soit un triple de Quillen clos, il est à vérifier (2.1.6. b$''$). Mais cela résulte de 1.8, appliqué aux deux couples de Quillen spéciaux $(C,TF)$, $(TC,F)$ dans $\M$ : les couples $(C',TC')$ et $(TC',F')$ sont des couples de Quillen spéciaux, ce qui prouve en même temps que le triple $(W',C',F')$ est spécial, q.e.d. % [sic] (C',TC') -> (C',TF'), image of the couple (C,TF)

Démontrons 2.1. Les premières deux relations (2.1.3) sont contenues dans 1.6, les relations $TC' \subset C'$, $TF' \subset F'$ sont tautologiques. Le critère (2.1.4) pour que $(W',C',F')$ soit un triple de Quillen clos, est contenu dans le rappel de la page 13, compte tenu que l'on sait déjà que $TF' = F' \cap W'$ (relation (2.4. b)). Que la condition (2.1.4. b), en présence de (2.1.4. a), soit équivalente à la condition (2.1.5. b$'$), est contenu dans la première assertion (2.5), en notant que l'on a $W' = W'^{\natural}$ (ce qui résulte de la définition de $W'$ et de $W = W^{\natural}$).

Il reste à prouver 2°), i.e. que la condition b) est satisfaite moyennant a) et b$''$) (2.1.6) dans chacun des trois cas A), B), C).

\textit{Cas A.} Par définition de $W' = (v^*)^{-1}(W)$, pour prouver que $TC' \subset W'$, il suffit de prouver
\[
v^*(TC')\ \subset\ W'\ ,  % [sic] W' -> W, v^*(TC') lies in M and W' = (v^*)^{-1}(W)
\]
et a fortiori, il suffit de prouver
\[
v^*(TC')\ \subset\ TC\ ,\qquad\quad \text{i.e. } v^*(TC') \harp F\ .
\tag*{(2.1.7)} % reading resolved: DER_XIX_corrections.csv
\]
Or par 1.5, appliqué au couple de foncteurs
\msp{page 18}
adjoints $(v^*,v_*)$, cette relation équivaut à
\[
TC' \harp v_*(F)\ .
\]
Comme $TC' \overset{\text{déf}}{=} v_!(TC)^{\sim}$, ceci équivaut à
\[
v_!(TC) \harp v_*(F)\ ,
\]
ou encore (en réappliquant 1.5 à $(v_!(TC),F)$ au lieu de $(TC',F)$), à
\[
v^*v_!(TC) \harp F\ ,\qquad\quad \text{i.e. } v^*v_!(TC) \subset \underbrace{TC}_{=\,F^*}\ ,
\]
or c'est là l'hypothèse b$''$).

\textit{Cas B.} Il suffit encore de prouver 2.1.7. Or, du fait que $v_!(TC)$ admet un petit ensemble Q-générateur à gauche, $v_!(TC)^{\sim}$ est la clôture de $v_!(TC)$ pour les opérations de cochangement de base, de composition transfinie, de facteurs directs. Comme $v^*$ commute à toutes ces opérations (comme il commute aux petites $\varinjlim$), et comme $TC$ dans $\Fl(\M)$ est stable par ces opérations, on voit donc que pour montrer que $v^*$ applique $v_!(TC)^{\sim}$ dans $TC$, il suffit de voir qu'il applique $v_!(TC)$ dans $TC$, ce qui est encore l'hypothèse b$''$).

\textit{Cas C.} On va vérifier b) sous la forme équivalente b$'$). Soit donc $f' \in \Fl(\M')$, $f' : X' \longrightarrow Y'$,
\printedpage{12}
% head: A. Grothendieck : Dérivateurs, Ch. XIX, Édition des Universités de Montpellier II et Paris VII
% foot: Version 11 février 2019
\msp{page 19}
posons
\[
f\ =\ v^*(f'),
\]
et factorisons $f$ en $pi$, avec $i \in TC$, $p \in F$, d'où
\[
v_!v^*(f') = v_!(f) = v_!(p)v_!(i)\ .
\]
Mais $v^*$ étant pleinement fidèle, $v_!v^* \xrightarrow{\ \sim\ } \mathrm{id}$, donc on trouve une factorisation de $f'$ en $v_!(p)v_!(i)$. Il suffit donc de voir que
\[
v_!(p)\ \in\ F'\ ,\qquad\quad v_!(i)\ \in\ W'\ ,
\]
ce qui équivaut à
\[
v^*v_!(p) \in F\ ,\qquad\quad v^*v_!(i) \in W\ ,
\]
et résulte des hypothèses que $v^*v_!(F) \subset F$ et $v^*v_!(TC) \subset TC \subset W$, q.e.d.

\textbf{N.B.} J'ai noté les cas A) et C) par acquis de conscience, sans application précise en vue. Le cas le plus utile sera justiciable du cas B), et mérite d'être réexplicité en un

\textbf{Théorème 2.3}\footnote[7]{cf. scholie 2.10, qui devrait précéder le théorème 2.3.}. \textit{Soit $(W,C,F)$ un triple de Quillen spécial dans $\M$. Pour toute $I$ dans} Cat\textit{, considérons le triple $(W(I),C\langle I\rangle,F(I))$ dans $\M(I) = \underline{\mathrm{Hom}}(I^{\circ},\M)$, où les notations sont celles de page 10, donc $W(I)$, $F(I)$ et $TF(I)$ se définissent fibre par fibre, et $C\langle I\rangle = (F(I))^* = (v_{I\,!}(C(I_0)))^{\sim}$}\footnote[8]{où $v_I : I_0 \longrightarrow I$ ($I_0$ catégorie discrète ayant mêmes objets que I).}\textit{.}
\begin{itemize}
\item[1°)] \textit{Pour toute $I$ dans} Cat\textit{, $(W(I),C\langle I\rangle,F(I))$ est}
\msp{page 20}
\textit{un triple de Quillen spécial dans $\M(I)$.}
\item[2°)] \textit{Si $u : I \longrightarrow I'$ est une flèche dans} Cat\textit{, d'où}
\[
\begin{tikzcd}
\M(I) \arrow[r, shift left, "u_!"] & \M(I'), \arrow[l, shift left, "u^*"]
\end{tikzcd}
\tag*{(2.3.1)}
\]
\textit{alors}
\[
\left\{
\begin{array}{l}
\left.\begin{array}{rcl}
u^*(W(I')) & \subset & W(I)\\
u^*(F(I')) & \subset & F(I)\\
u^*(TF(I')) & \subset & TF(I)
\end{array}\right\}\textit{tautologie}\\
\begin{array}{rcl}
u_!(C\langle I\rangle) & \subset & C\langle I'\rangle\\
u_!(TC\langle I\rangle) & \subset & TC\langle I'\rangle\ ,
\end{array}
\end{array}\right.
\tag*{(2.3.2)}
\]
\textit{enfin, $u_!$ induit un foncteur}
\[
\M(I)_{\cc} \longrightarrow \M(I')_{\cc} \qquad\quad \text{\footnotemark[9]}
\tag*{(2.3.3)}
\]
\footnotetext[9]{Sous-catégories des objets cofibrants.}
compatible aux localiseurs $W(I)_{\cc} = W(I)|\M(I)_{\cc}$ \textit{et} $W(I')_{\cc} = W(I')|\M(I')_{\cc}$\textit{, donc induit un foncteur}
\[
u_!^W\ :\ \Ho_W(I) \longrightarrow \Ho_W(I')\ .
\tag*{(2.3.4)}
\]
\printedpage{13}
% head: A. Grothendieck : Dérivateurs, Ch. XIX, Édition des Universités de Montpellier II et Paris VII
% foot: Version 11 février 2019
\textit{Ce foncteur est un adjoint à gauche du foncteur $u^*_W$,}
\[
u^*_W\ :\ \Ho_W(I') \longrightarrow \Ho_W(I)\ ,
\tag*{(2.3.5)}
\]
\textit{induit par $u^*$} (3.3.1). % [sic] (3.3.1) -> (2.3.1), the adjunction displayed in (2.3.1)
\end{itemize}
\textsc{Démonstration de 2.3.}

\textit{Prouvons} 1°). On considère comme dans 1.10 le foncteur
\[
v = v_I\ :\ I_0 \longrightarrow I
\]
et le couple de foncteurs adjoints
\[
\begin{tikzcd}
\M(I) \arrow[r, shift right, "v^*"'] & \M(I_0)\ , \arrow[l, shift right, "v_!"']
\end{tikzcd}
\tag*{(2.3.6)}
\]
auquel on peut appliquer 2.2. Il reste
\msp{page 21}
à vérifier seulement
\[
v^*v_!(TC)\ \subset\ TC\ ,
\tag*{(2.3.7)}
\]
or on a, pour $X = (X_i)_{i \in I_0} \in \M(I_0)$,
\[
\begin{aligned}
v_!(X)(i) &\ \simeq\ \textstyle\coprod_{j \in I_0} v_{j\,!}(X_j)(i)\\
&\ \simeq\ \textstyle\coprod_{j \in I_0}\ \ \underbrace{X_j^{\mathrm{Hom}(i,j)}}_{\substack{\text{somme directe}\\ \text{de }\mathrm{Hom}(i,j)\text{ copies}}}\qquad ,
\end{aligned}
\]
et la relation (2.3.7) résulte du fait que $TC$, étant Q-clos à gauche, est stable par petites sommes directes.

\textit{Prouvons} 2°). Les formules (2.3.2) sont essentiellement tautologiques, les deux dernières résultent de la troisième et de la deuxième, en vertu de 1.5 appliqué au couple de foncteurs adjoints $(u_!,u^*)$. La relation $u_!(C\langle I\rangle) \subset C\langle I'\rangle$, et le fait que $u_!$ transforme objet initial en objet initial, impliquent (2.3.3). Que $u_!$ restreint à $\M(I)_{\cc}$ soit compatible aux localiseurs provient de $u_!(TC\langle I\rangle) \subset TC\langle I'\rangle$, et du fait que toute flèche $f$ dans $\M(I)_{\cc}$ qui est dans $W(I)$, se décompose en $f = pi$, avec $i \in TC$ et $p$ inverse à gauche d'une $j \in TC\langle I\rangle$. Ceci vu, joint aux deux premières formules (2.3.2) montre qu'on est sous les conditions
\msp{page 22}
du théorème d'adjonction de Quillen, d'où l'adjonction de $u_!^W$ et de $u^*_W$, q.e.d.

\textbf{Corollaire 2.4.} \textit{Soit $W$ un localiseur dans $\M$, tel qu'il existe un triple de Quillen spécial $(W,C,F)$ incluant $W$. Alors le prédérivateur} $\Ho_W$ \textit{sur} Cat \textit{satisfait aux axiomes Dér 1 (décomposition), Dér 2 (localisation), Dér 3 (objet final), Der 4 bis (existence des $u_!^W$), Dér 5 bis (calcul des fibres des $u_!^W$), Dér 6 bis (exactitude à gauche, i.e. exactitude à gauche des carrés $W$-cocartésiens dans les} $\Ho_W(I)$\textit{), Dér 6 (exactitude à droite des carrés $W$-cartésiens).} % [sic] Der 4 bis -> Dér 4 bis, accent missing as printed

On vient de prouver Dér 4 bis, et Dér 5 bis est pratiquement contenu dans le mode d'obtenir de $u_!$, compte tenu de proposition 1.12, qui implique que pour $X$ dans $\M(I)_{\cc}$, et $i'$ dans $I'$, l'image inverse de $X$ dans $\M(i'\backslash I)$ est dans $\M(i'\backslash I)_{\cc}$ \dots


\endgroup
% END UNIT der-ch19/der-ch19_p011-015.tex
% BEGIN UNIT der-ch19/der-ch19_p016-020.tex
\begingroup\makeatletter%
\derleqno%
\newcommand{\printedpage}[1]{\marginpar{\footnotesize p.~#1}}%
\newcommand{\msp}[1]{\marginpar{\footnotesize [#1]}}%
\newcommand{\deq}{\overset{\text{déf}}{=}}%
\newcommand{\Ho}{\mathrm{Ho}}%
\newcommand{\Mono}{\mathrm{Mono}}%
\newcommand{\Fl}{\mathrm{Fl}}%
\newcommand{\Filt}{\mathrm{Filt}}%
\newcommand{\Cof}{\mathrm{Cof}}%
\newcommand{\qpair}{\mathrel{\leftharpoondown\joinrel\relbar\joinrel\rightharpoonup}}%
\makeatletter%
\renewcommand\@makefnmark{\hbox{\@textsuperscript{\normalfont(\@thefnmark)}}}%
\renewcommand\@makefntext[1]{\parindent 1em\noindent\hbox{\@textsuperscript{\normalfont\@thefnmark}}#1}%
\newcommand{\Qcouple}{\mathrel{\leftarrow\mkern-9mu{+}\mkern-9mu\rightarrow}}%
\makeatother%
\makeatother\let\printedpage\derpp\let\msp\dermsp%
\let\mspin\DERmspin%


\printedpage{14}
% head: A. Grothendieck : Dérivateurs, Ch. XIX, Édition des Universités de Montpellier II et Paris VII
% foot: Version 11 février 2019

\msp{page 23}
L'axiome de décomposition est satisfait pour tout prédérivateur de la forme $\mathbf{D}_{\mathcal{M},W}$. L'axiome d'exactitude à gauche Dér 5 bis est valable, pourvu que les $(\mathcal{M}(I),W(I))$ s'insèrent dans une structure à cofibrations, ce qui est le cas. De même pour l'axiome Dér 2 de localisation, pourvu que de plus $W$ soit fortement saturé, ce qui est également le cas. Enfin, l'axiome de l'objet final dit que si $I$ admet un objet final, alors
\[ \Ho_W(e) \longrightarrow \Ho_W(I) \]
est pleinement fidèle, ou encore que le morphisme d'adjonction
\[ u_!u^*(\xi) \longrightarrow \xi \]
est un isomorphisme. Mais en vertu de 1.12, cas 3°, dans ce cas $u^*$ applique $\mathcal{M}(e)_{\mathrm{c}} = \mathcal{M}_{\mathrm{c}}$ dans $\mathcal{M}(I)_{\mathrm{c}}$, donc le calcul du couple de foncteurs adjoints $u_!^W$, $u_W^*$ peut se faire en utilisant les restrictions de $u_!$, $u^*$ à $\mathcal{M}(I)_{\mathrm{c}}$ et $\mathcal{M}_{\mathrm{c}}$ respectivement. Comme le morphisme d'adjonction $u_!^{\mathrm{c}}u_{\mathrm{c}}^* \longrightarrow \mathrm{id}$ est un isomorphisme, il en est donc de même de $u_!^W u_W^* \longrightarrow \mathrm{id}$, qui s'en déduit en localisant, q.e.d.

\textbf{2.5.} Pour achever de construire $\Ho_W$ comme un dérivateur sur Cat tout entier, il faut \msp{page 24} donc établir Dér 4, Dér 5 : existence des $u_*^W$, et calcul des fibres de $u_*^W$. Quant à Dér 6, il est acquis automatiquement, du fait que chaque $(\mathcal{M}(I),W(I))$ s'insère dans une structure de fibration $(W(I),F(I))$. (Oubli dans la démonstration de 2.4.) L'idée naturelle est d'utiliser à présent la structure de triple
\[ (W(I),C(I),\ \underbrace{F[I]}_{\deq TC(I)_*}\ )\ , \tag*{(2.5.1)} \]
mais la difficulté est de prouver la factorisation pour un couple
\[ C(I) \longleftrightarrow TF[I]\ , \qquad \underbrace{TC(I)}_{\deq C(I)\cap W(I)} \quad \longleftrightarrow F[I]\ , \tag*{(2.5.2)} \]
faute de pouvoir montrer, dans les conditions générales de 2.3 seulement (triple de Quillen spécial) que $C(I)$ et $TC(I)$ sont également spéciaux. (Cf. question 1.11.) Mais moyennant des conditions convenables sur $\mathcal{M}$, on arrivera à s'en tirer si $C = \mathit{mono}$. L'énoncé pertinent, un peu technique, se trouve dans XV cor. 7.4 (p. 98-99), que je vais rappeler ici.

\printedpage{15}
% head: A. Grothendieck : Dérivateurs, Ch. XIX, Édition des Universités de Montpellier II et Paris VII
% foot: Version 11 février 2019

\msp{page 25}
\textbf{Théorème 2.6.} \textit{Soit $\mathcal{M}$ une catégorie satisfaisant les conditions suivantes.}
\[ \left\{ \begin{minipage}{0.72\textwidth}\itshape
\begin{itemize}[leftmargin=2.5em]
\item[\upshape a)] $\mathcal{M}$ est “spéciale”, i.e. stable par petites $\varinjlim$ , et ses éléments sont accessibles.
\item[\upshape b)] Les limites inductives filtrantes dans $\mathcal{M}$ sont exactes à gauche (et dans $\mathcal{M}$ les $\varprojlim$ finies sont représentables – résulte de \textup{a} et \textup{c}).
\item[\upshape c)] $\mathcal{M}$ admet une filtration cardinale $(\Filt^\pi)_{\pi\geq\pi_0}$.
\item[\upshape d)] Les $\Filt^\pi$ sont stables par $\varprojlim$ indexées par des catégories $I$ telles que $\operatorname{card}\Fl(I) \leq \pi$, et par quotients ($Z \in \Filt^\pi$ et $Z \longrightarrow Z'$ épimorphisme implique $Z' \in \Filt^\pi$).
\item[\upshape e)] Si $A_1$, $A_2$ sont deux sous-objets d'un objet $B$ de $\mathcal{M}$, et $A_0 = A_1 \cap A_2$, alors $A_1 \amalg_{A_0} A_2 \longrightarrow B$ est un monomorphisme.
\item[\upshape f)] Si on a un carré cocartésien
\begin{center}
\begin{tikzcd}
A_0 \arrow[r, "i"] \arrow[d] & A_1 \arrow[d] \\
A_2 \arrow[r, "i'"] & A_3
\end{tikzcd}
\end{center}
avec $i$ monomorphisme, alors $i'$ est un monomorphisme.
\end{itemize}
\end{minipage}\right. \tag*{(2.6.1)} \]
$^{(10,11,12,13)}$\footnotetext[10]{\textbf{N.B.} a + c signifie aussi que $\mathcal{M}$ est une “catégorie $\mathfrak{U}$-paratopos”, et implique l'existence d'une filtration cardinale satisfaisant d) (sauf l'assertion concernant les quotients $Z'$ de $Z$). Donc il s'agit d'un $\mathfrak{U}$-paratopos satisfaisant de plus b) e) f). (Est-ce alors un topos\,? \textit{Non,} car e), f) sont O.K. dans les \textit{catégories abéliennes\,!})}\footnotetext[11]{\textbf{N.B.} La condition d) est sans doute inutile, du moins si les épimorphismes sont effectifs, car d) sera satisfaite en tous cas pour $\pi = 2^c$, avec $c \geq \pi_0$.}\footnotetext[12]{\textbf{N.B.} La condition e) n'est \textit{pas} satisfaite pour la catégorie (Gr), et le théorème est en défaut pour cette catégorie, laquelle n'admet \textit{pas} “assez d'injectifs”, cf. SGA 4 I 8.12.9. Par contre, e) et f) sont satisfaites par les catégories \textit{abéliennes.}}\footnotetext[13]{Cf. XVIII 4.6.7 (p. 157) pour une formulation plus sympathique de l'ensemble des conditions a) à f) pour $\mathcal{M}$.}. \textit{(On dira alors que $\mathcal{M}$ est} “très spéciale”\textit{). Soit de plus}
\[ W \ \subset \Fl(\mathcal{M}) \tag*{(2.6.2)} \]
\textit{un localiseur dans $\mathcal{M}$, qui satisfait les conditions}
\[ \left\{ \begin{minipage}{0.72\textwidth}\itshape
\begin{itemize}[leftmargin=2.5em]
\item[\upshape a)] $W$ est stable par $\varinjlim$ filtrantes $^{\textup{(14)}}$.
\item[\upshape b)] $W$ est de plus $\mathrm{L}_1$-accessible, i.e. il existe une \textup{petite} partie $W_0 \subset W$ telle que $W$ soit la $\mathrm{L}_1$-enveloppe de $W$, i.e. que $W$ soit la plus petite partie de $\Fl(\mathcal{M})$ contenant $W_0$, et stable par $\varinjlim$ filtrantes. % [sic] la L_1-enveloppe de W -> la L_1-enveloppe de W_0, W_0 being the small generating part just introduced
\end{itemize}
\end{minipage}\right. \tag*{(2.6.3)} \]
\footnotetext[14]{\textbf{N.B.} a) implique déjà que $W = W^\natural$.}

\printedpage{16}
% head: A. Grothendieck : Dérivateurs, Ch. XIX, Édition des Universités de Montpellier II et Paris VII
% foot: Version 11 février 2019

\msp{page 26}
\textit{Posons alors}
\[ C \ = \ \Mono(\mathcal{M}), \qquad TC \ = \ C\cap W\ , \qquad TF \ = \ C_*\ , \qquad F = (TC)_*\ . \tag*{(2.6.4)} \]
\textit{Sous ces conditions :}
\begin{itemize}[leftmargin=2.5em]
\item[1°)] \textit{$C \Qcouple TF$, plus précisément $C \qpair TF$ est un couple de Quillen spécial.} % reading resolved: DER_XIX_corrections.csv
\item[2°)] \textit{Pour que $TC \qpair F$ soit un couple de Quillen, il faut et il suffit que $TC$ soit stable par cochangement de base, et alors c'est un couple de Quillen spécial.} % reading resolved: DER_XIX_corrections.csv
\item[3°)] \textit{Pour que $(W,C,F)$ soit un triple de Quillen (lequel sera nécessairement spécial par \textup{1°, 2°}, et en particulier clos), il faut et il suffit que les conditions suivantes soient satisfaites.}
\[ \left\{ \begin{array}{ll}
\text{a)} & TC\ (\deq W\cap\Mono)\ \text{\itshape est stable par cochangement de base.}\\[1ex]
\text{b)} & TF\ (\deq (\Mono)_*) \ = \ \text{\itshape “applications de type injectif”}\ \subset W.
\end{array}\right. \tag*{(2.6.5)~$^{(15)}$} \]
\footnotetext[15]{\textbf{N.B.} Ces conditions (2.6.5) équivalent à dire que $(W,C)$ est une structure de catégorie à cofibrations sur $\mathcal{M}$.}
\textit{Et la condition} b) \textit{équivaut encore à :}
\[ \begin{minipage}{0.7\textwidth}
\begin{itemize}[leftmargin=2.5em]
\item[b$'$)] \textit{Toute flèche $f$ de $\mathcal{M}$ (ou seulement de $TF$) se factorise en $f = pi$, avec $i$ monomorphisme, $p \in W$.}
\end{itemize}
\end{minipage} \tag*{(2.6.6)} \]
\end{itemize}

\textbf{Corollaire 2.7.} \textit{$\mathcal{M}$ et $W$ étant comme dans} (2.6.1), (2.6.3), \textit{pour que $(W,C,F)$ soit un triple de Quillen} propre à gauche \textit{(i.e. $C \subset \Cof_W$), il faut et il suffit que $W$ satisfasse la condition} (2.6.5. b) \textit{(i.e. $W$ contient les flèches de type injectif), et de plus}
\msp{page 27}

\begin{itemize}[leftmargin=7em]
\item[] $C \subset \Cof_W$, \textit{i.e. la condition :}
\item[a$'$)] \textit{Pour tout diagramme cocartésien}
\[ \begin{tikzcd}
A_0 \arrow[r, "i"] \arrow[d, "f"'] & A_1 \arrow[d, "f'"'] \\
A_2 \arrow[r, "i'"] & A_3
\end{tikzcd} \tag*{(2.7.7)} \]
\textit{avec $i$ monomorphisme, si $f \in W$, alors $f' \in W$. (Stabilité de $W$ par cochangement de base par monomorphismes.)}
\end{itemize}

\textit{Sous ces conditions, ce triple de Quillen est spécial, et a fortiori clos.}

Ce corollaire résulte de 2.6. 3°), puisque $C \subset \Cof_W$ implique $TC \deq C\cap W = C\cap W^{\mathrm{univ}}$ (car $\Cof_W \cap W \ = \ W^{\mathrm{univ}}$), donc que $TC$ est stable par cochangement de base, d'où la condition (2.6.5. a).

\printedpage{17}
% head: A. Grothendieck : Dérivateurs, Ch. XIX, Édition des Universités de Montpellier II et Paris VII
% foot: Version 11 février 2019

\textbf{2.8.} On dit qu'une \textit{structure de Quillen} $(W,C,F)$ sur $\mathcal{M}$ est \textit{très spéciale,} si elle est obtenue comme au théorème 2.1., i.e. si on a les conditions :
\[ \left\{ \begin{minipage}{0.72\textwidth}
\begin{itemize}[leftmargin=2.5em]
\item[a)] $\mathcal{M}$ est très spéciale, i.e. satisfait aux conditions (2.6.1).
\item[b)] $W$ est stable par petites $\varinjlim$ filtrantes, et est accessible (conditions (2.6.3)).
\item[c)] $C = \Mono$.
\end{itemize}
\end{minipage}\right. \tag*{(2.8.1)} \]
Donc pour qu'un localiseur $W$ dans $\mathcal{M}$ corresponde à une structure de Quillen très spéciale \msp{page 28} dans $\mathcal{M}$ (laquelle sera forcément unique), il faut et il suffit que $\mathcal{M}$ et $W$ satisfassent les conditions (2.8.1. a)\,b)), et que de plus on ait :
\[ \left\{ \begin{minipage}{0.65\textwidth}
\begin{itemize}[leftmargin=2.5em]
\item[a)] $TC \deq \Mono\cap W$ est stable par cochangement de base.
\item[b)] $TF \deq (\Mono)_* \subset W$, i.e. $W$ contient les flèches de type injectif dans $\mathcal{M}$.
\end{itemize}
\end{minipage}\right. \tag*{(2.8.2)} \]
\footnote[16]{On dit aussi que le localiseur $W$ est très spécial, s'il définit une structure de Quillen très spéciale.}.

On trouve maintenant le

\textbf{Théorème 2.9.} \textit{Soit} $(W, C = \Mono, F = (TC)_* = (W\cap\Mono)_*)$ \textit{un triple de Quillen très spécial dans une catégorie $\mathcal{M}$} (2.8). \textit{Alors :}
\begin{itemize}[leftmargin=2.5em]
\item[1°)] \textit{Pour toute $I$ dans} Cat, $(W(I),C(I),F[I])$ \textit{est un triple de Quillen très spécial dans $\mathcal{M}(I)$, et on a en particulier}
\[ C(I) \ = \ \Mono(\mathcal{M}(I))\ , \tag*{(2.9.1)} \]
\textit{et de plus}
\[ F[I]\cap W(I) \ = \ (TF)[I]\ . \tag*{(2.9.2)} \]
\item[2°)] \textit{Pour toute flèche $u : I \longrightarrow I'$ dans} Cat, \textit{d'où le couple de foncteurs adjoints}
\[ \begin{tikzcd}
\mathcal{M}(I) \arrow[r, shift right, "u_*"'] & \mathcal{M}(I') \arrow[l, shift right, "u^*"']
\end{tikzcd}\ , \tag*{(2.9.3)} \]
\textit{on a}
\[ \left\{ \begin{array}{rcl}
u^*(W(I')) & \subset & W(I)\\
u^*(C(I')) & \subset & C(I)\\
u^*(TC(I')) & \subset & TC(I)\\
u_*(F[I]) & \subset & F[I']\\
u_*(TF[I]) & \subset & TF[I']\ ,
\end{array}\right. \tag*{(2.9.4)} \]
\textit{enfin $u_*$ induit un foncteur}
\printedpage{18}
% head: A. Grothendieck : Dérivateurs, Ch. XIX, Édition des Universités de Montpellier II et Paris VII
% foot: Version 11 février 2019
\msp{page 29}
\[ u_*^{\mathrm{f}} \ :\ \mathcal{M}(I)_{\mathrm{f}} \longrightarrow \mathcal{M}(I')_{\mathrm{f}} \qquad\quad ^{(17)} \tag*{(2.9.5)} \]
\footnotetext[17]{sous-catégories pleines des objets fibrants.}
compatible aux localiseurs $W(I)_{\mathrm{f}}$ \textit{et} $W(I')_{\mathrm{f}}$ \textit{induits par $W(I)$, $W(I')$ respectivement, donc induit un foncteur}
\[ u_*^W \ :\ \Ho_W(I) \longrightarrow \mathit{Ho}_W(I')\ . \tag*{(2.9.6)} \] % [sic] italic Ho_W(I') -> upright Ho_W(I'), inconsistent with the upright Ho_W elsewhere
\textit{Ce foncteur est un} adjoint à droite \textit{du foncteur}
\[ u_W^* \ :\ \mathit{Ho}_W(I') \longrightarrow \Ho_W(I) \tag*{(2.9.7)} \] % [sic] italic Ho_W(I') -> upright Ho_W(I'), inconsistent with the upright Ho_W elsewhere
\textit{induit par $u^*$} (2.9.3).
\end{itemize}
\textsc{Démonstration de 2.9.} Pour prouver 1°), il faut vérifier d'abord que les conditions préliminaires (2.8.1. a)\,b)) sur $\mathcal{M}$, $W$ sont stables par passage de $\mathcal{M}$, $W$ à $\mathcal{M}(I)$, $W(I)$, ce qui se fait par AQT \texttt{[âne qui trotte]}. D'ailleurs on vérifie tout de suite que
\[ C(I) \ = \ \Mono(\mathcal{M}(I))\ , \]
et le théorème 2.6 (1°, 2°) implique alors que
\[ \begin{array}{ccccc}
C(I) & \qpair & C(I)_* & \deq & TF[I]\\
\underbrace{TC(I)}_{\deq C(I)\cap W(I)} & \qpair & (TC(I))_* & \deq & F[I]
\end{array} \] % reading resolved: DER_XIX_corrections.csv
sont des couples de Quillen spéciaux. (En notant pour le deuxième, que $TC(I)$ est bien \msp{page 30} stable par cochangement de base, $TC$ l'étant.) Pour vérifier que l'on a un triple de Quillen très spécial, il reste en vertu de 2.8 à vérifier que
\[ TF[I] \ \subset\ W(I)\ , \]
ce qui s'écrit
\[ v^*(TF[I]) \ \subset\ W(I_0)\ , \]
où
\[ v \ :\ I_0 \longrightarrow I \]
est défini comme dans (1.12) (p. 10). En fait, on aura même
\[ v^*(TF[I]) \ = \ v^*v_*(TF(I_0)) \ \subset\ \underbrace{TF(I_0)}_{= F(I_0)\cap W(I_0)}\quad , \]
et la démonstration est duale de celle de
\[ v^*v_!(TC(I_0)) \ \subset\ TC(I_0) \]
(cf. (2.3.7), p. 21), elle ne dépend pas des questions de factorisation dans $\mathcal{M}(I)$, et provient ici de la stabilité de $TF$ par petits produits. Cela prouve 1°.

La démonstration de 2° est maintenant la duale de celle de (2.3. 2°). Comme cela nous fait répéter deux fois la même démonstration, cela montre qu'on n'a pas “monté le sorite” tout à fait avec assez de soin. Donc à titre de rattrapage, j'ajoute un


\endgroup
% END UNIT der-ch19/der-ch19_p016-020.tex
% BEGIN UNIT der-ch19/der-ch19_p021-025.tex
\begingroup\makeatletter%
\newcommand{\printedpage}[1]{\marginpar{\footnotesize p.~#1}}%
\newcommand{\msp}[1]{\marginpar{\footnotesize [#1]}}%
\newcommand{\M}{\mathcal{M}}%
\newcommand{\Cat}{\mathrm{Cat}}%
\newcommand{\Ho}{\mathrm{Ho}}%
\newcommand{\Fl}{\mathrm{Fl}}%
\newcommand{\Cof}{\mathrm{Cof}}%
\newcommand{\Diag}{\mathrm{Diag}}%
\newcommand{\id}{\mathrm{id}}%
\newcommand{\deq}{\overset{\text{déf}}{=}}%
\newcommand{\Qs}{\tilde{\ }}%
\newcommand{\fact}{\mathrel{\leftarrow\mkern-11mu\twoheadrightarrow}}%
\newcommand{\harp}{\mathrel{\leftharpoondown\joinrel\relbar\joinrel\rightharpoonup}}%
\makeatother\let\printedpage\derpp\let\msp\dermsp%
\let\mspin\DERmspin%


\printedpage{19}
% head: A. Grothendieck : Dérivateurs, Ch. XIX, Édition des Universités de Montpellier II et Paris VII
% foot: Version 11 février 2019

\msp{page 31}
\textbf{Scholie 2.10.} \textit{Soit $(W,C,F)$ un triple de Quillen clos dans $\M$. Pour $I$ dans \textup{Cat}, pour que $(W(I),C\langle I\rangle,F(I))$ soit un triple de Quillen clos pour $\M(I)$, il faut et il suffit que les couples $C\langle I\rangle\longleftrightarrow TF(I)$ et $TC\langle I\rangle\longleftrightarrow F(I)$ soient factorisants. Supposons que cette condition soit satisfaite pour $I$ et pour $I'$, et soit $u : I\longrightarrow I'$ une flèche de $I$ dans $I'$. Alors les assertions de \textup{2.3. 2°} sont valables, en particulier on trouve un foncteur $u^{W}_{!} : \Ho_W(I)\longrightarrow\Ho_W(I')$, adjoint à gauche de $u^{*}_{W}$\footnote[18]{Dire aussi que le calcul des fibres de $u^{W}_{!}(X)$, prévu dans Dér 5 bis, est OK.}. Ces résultats se dualisent aussitôt (par $\M\longmapsto\M^{\circ}$, $I\longmapsto I^{\circ}$) pour les triples $(W(I),C(I),F[I])$, et pour la construction de $u^{W}_{*}$, adjoint à droite de $u^{*}_{W}$. En particulier, pour que le triple précédent soit un triple de Quillen (forcément clos), il faut et il suffit que les deux couples $C(I)\longleftrightarrow TF[I]$ et $TC(I)\longleftrightarrow F[I]$ soient factorisants ; et quand cette condition est satisfaite pour $I$ et pour $I'$ (pour une flèche $u : I\longrightarrow I'$ dans \textup{Cat} donnée), alors $u^{W}_{*}$ existe et se construit dualement de la construction de $u^{W}_{!}$ explicitée dans \textup{2.3. 2°}.}

\msp{page 32}
% a vertical bar is printed in the left margin alongside Corollaire 2.11
\textbf{Corollaire 2.11.} \textit{Le prédérivateur $\mathbf{D}_{\M,W}$ défini pour un localiseur \textup{très spécial} $W$ dans $\M$ est un dérivateur en plein sens du terme, i.e. il satisfait à tous les axiomes Dér 1 à Dér 6 (y inclus Dér 4 bis et Dér 5 bis et Dér 6 bis) des dérivateurs.}

\textbf{N.B.} À ma propre surprise, je n'ai pas été capable de trouver un cas d'un localiseur $W$ qui définisse un dérivateur sur Cat tout entier (avec existence de $u^{W}_{!}$, $u^{W}_{*}$), \textit{que} pour $W$ un localiseur très spécial, donc s'insérant dans un triple de Quillen très spécial -- ou quand on peut s'y ramener de la façon standard. (Par exemple dans le cas où $\M=\Cat$, $W$ un localiseur fondamental satisfaisant Loc 6 a) et Loc 8, 8 bis)\footnote[19]{axiome des carrés cocartésiens (version directe), axiome des limites inductives, accessibilité.} -- les monomorphismes dans $W$ se refusent formellement à être stable par cochangement de base !) J'ignore totalement si le résultat mirobolant ``prouvé'' par \textsc{Anderson} est vrai ou faux. Ça a l'air presque trop beau pour être vrai \ldots{}D'ailleurs, je ne connais pas de structure de Quillen spéciale qui ne soit très spéciale -- désolé ! % [sic] stable -> stables, agreement with « les monomorphismes »

\msp{page 33}

\section*{3\quad Cas des catégories à cofibrations, des catégories à fibrations}

Je rappelle les axiomes d'une \textit{structure à cofibrations} $(W,C)$ sur une catégorie $\M$.

\begin{itemize}[leftmargin=4em,labelwidth=3.5em,align=left]
\item[Cof$_0$ :] $W$ est un localiseur.
\item[Cof$_1$ :] $C$ contient les isomorphismes, est stable par cochangement de base, par composition.
\item[Cof$_2$ :] $C\cap W$ est stable par cochangement de base (et aussi par composition, et contient les isomorphismes, en vertu de Cof$_0$, Cof$_1$).
\item[Cof$_3$ :] Toute flèche $f$ dans $\M$ se factorise en $f=pi$, avec $i\in C$, $p\in W$.
\end{itemize}

(Axiomes de K. S. \textsc{Brown}, légèrement généralisés par \textsc{Anderson}). Il y a un ``axiome
\printedpage{20}
% head: A. Grothendieck : Dérivateurs, Ch. XIX, Édition des Universités de Montpellier II et Paris VII
% foot: Version 11 février 2019
de continuité'' ou \textit{``axiome transfini''} facultatif, qui est le suivant :

\begin{itemize}[leftmargin=4em,labelwidth=3.5em,align=left]
\item[Cof$_4$ :] $\M$ stable par petites $\varinjlim$ filtrantes, et $C$ et $TC = C\cap W$ sont stables par composition transfinie.
\end{itemize}

Ainsi, Cof$_1$, Cof$_2$, Cof$_4$ peuvent se résumer en disant que $C$ et $TC$ sont Q-saturés transfiniment à gauche.

La structure de cofibration est dite \textit{spéciale} si elle satisfait Cof$_4$, si $\M$ est spéciale, et si $C$ admet une \textit{petite} partie $C_0$ qui soit Q-génératrice à gauche, i.e. telle que
\[
C \ \subset\ C_0\Qs\,, \qquad \text{i.e.}\quad \tilde{C}=C_0\Qs. \tag*{(3.1)}
\]

\msp{page 34}
Alors, par le théorème de factorisation 1.1, on a
\[
C \ \fact\ TF \ \deq\ C_* \ ,
\]
et l'axiome de factorisation Cof$_3$ implique (cf. 1.4) $TF\subset W^{\natural}$. Plus généralement, on a :

\textbf{Proposition 3.1.} \textit{Supposons que $\M$ soit spéciale (stable par petites $\varinjlim$ , et éléments accessibles), que $(W,C)$ satisfasse à Cof$_0$, Cof$_1$, et que $C$ soit stable par composition transfinie (ce qui, joint à Cof$_1$, s'énonce en disant que $C$ est Q-saturé transfini à gauche), enfin que $C$ admette une petite partie $C_0$ Q-génératrice à gauche, de sorte que $C\fact TF\deq C_*$. Alors pour que l'axiome de factorisation Cof$_3$ soit satisfait, il suffit qu'on ait}
\[
TF \ \subset\ W \ , \tag*{(3.1.1)}
\]
\textit{et il faut qu'on ait $TF\subset W^{\natural}$, donc si $W=W^{\natural}$, l'axiome Cof$_3$ équivaut à} (3.1.1).

La structure de cofibration est dite \textit{semi-close} si $C$ est Q-clos à gauche, i.e. si on a $C=C\Qs$ (ce qui implique que $C$ est stable par composition transfinie), \textit{close} si de plus $TC=(TC)\Qs$ (ce qui implique l'axiome transfini Cof$_4$).

\msp{page 35}
Je ne reviens pas ici sur la théorie de K. S. \textsc{Brown} des catégories à cofibrations. C'est à ma connaissance l'ensemble d'axiomes le plus économique connu, pour assurer d'une part une bonne prise sur $\Ho_W$ (par une notion correspondante d'homotopie d'applications, et un calcul de fractions dans $\Pi\M_{\mathrm{c}}$), et d'autre part l'existence de la suite \textit{exacte} de cofibrations, donc une théorie de cofibres $W$-homotopiques. (\textbf{N.B.} Quand la structure est \textit{propre,} i.e. $C\subset\Cof_W$, alors l'existence de cette structure revient à dire simplement que ``$W$ admet assez de $W$-cofibrations'', i.e. que toute $f\in\Fl(\M)$ se factorise en $pi$, avec $i\in\Cof_W$, $p\in W$ ; et on peut donc prendre $C=\Cof_W$. Dans le cas où $W$ est stable par petites $\varinjlim$ filtrantes, $C=\Cof_W$ et $TC=W^{\mathrm{univ}}$ satisfont de plus l'axiome transfini Cof$_4$.) Enfin, K. S. \textsc{Brown} montre que toute flèche $i : X\longrightarrow Y$ dans $C$, \textit{avec $X$, $Y$ cofibrants} (i.e. $\emptyset_{\M}\longrightarrow X$, $\emptyset_{\M}\longrightarrow Y$ dans $C$) est dans $\Cof_W$ (``Lemme de Brown'').

Je m'intéresse ici à la question d'étendre la structure à cofibrations à une catégorie de la forme $\M(I)$, et d'utiliser cette extension pour construire
\msp{page 36}
les foncteurs $u^{W}_{!}$. Notons déjà que la théorie de K. S. \textsc{Brown}, avec les compléments d'\textsc{Anderson}\footnote[20]{revus et arrangés par mes soins.}, donne essentiellement ceci :

\printedpage{21}
% head: A. Grothendieck : Dérivateurs, Ch. XIX, Édition des Universités de Montpellier II et Paris VII
% foot: Version 11 février 2019

\textbf{Proposition 3.2.} \textit{Soit $W\subset\Fl(\M)$ un localiseur associé à une structure à cofibrations $(C,W)$ dans $\M$.}
\begin{itemize}[leftmargin=3em,labelwidth=2.5em,align=right]
\item[1°)] \textit{Sur la catégorie \textup{$\Diag_0$} des ensembles ordonnés \textup{finis}, ce prédérivateur satisfait aux axiomes Dér 1 (décomposition), Dér 2 (localisation), Dér 3 (objet final), Dér 4 bis (existence de $u^{W}_{!}$), Dér 5 bis (calcul des fibres de $u^{W}_{!}$), Dér 6 bis (exactitude à gauche).}
\item[2°)] \textit{C'est encore vrai sur la sous-catégorie \textup{$\Diag_1$} dans \textup{Cat} plus grande, formée des ensembles ordonnés $I$ \textup{noethériens} tels que l'espace topologique associé à $I^{\circ}$ soit noethérien (toute suite croissante d'objets de $I$, ou d'ouverts de $I^{\circ}$, est stationnaire).}
\item[3°)] \textit{Enfin, si l'axiome transfini Cof$_4$ est satisfait, i.e. si $C$ et $TC$ sont stables par composition transfinie, alors la conclusion de \textup{1°} reste valable sur la sous-catégorie pleine \textup{$\Diag_2$} de \textup{Cat} formée des ensembles ordonnés qui sont noethériens sans plus.}
\end{itemize}

\textbf{Corollaire 3.3.} \textit{Supposons que l'on ait Cof$_3$ sous la forme plus forte d'une factorisation
\msp{page 37}
$f=pi$ \textup{fonctorielle} en $f$, ce qui est le cas en particulier si la structure à cofibrations est spéciale. Alors pour toute $I$ dans \textup{Cat}, $\M(I)$ muni de $(W(I),C(I))$ est une catégorie à cofibrations, laquelle satisfait l'axiome transfini Cof$_4$ si $(W,C)$ y satisfait. Donc on peut lui appliquer les résultats de \textup{3.2}, qui impliquent notamment que les axiomes Dér 1, 2, 3 ( ?) sont valables sur \textup{Cat} tout entier, et que $u^{W}_{!}$ existe quand $u$ est de la forme $\id_I\times v$, avec $v : J\longrightarrow K$ une flèche dans \textup{$\Diag_1$} (resp. dans \textup{$\Diag_2$}, si l'axiome transfini est satisfait).}

\textbf{3.4.} La clef de la théorie est dans l'extension de la structure à cofibrations $(W,C)$ de $\M$ à une catégorie $\M(I)$, où $I$ est un ensemble ordonné, en $(W(I),C_I)$. Les faits principaux sont énoncés dans XVII, dans les théorèmes p. 26-27, p. 50, p. 60. C'est quand $I$ est \textit{noethérien,} et qu'on suppose soit que $\mathrm{Esp}(I^{\circ})$ est également noethérien (p.ex. $I$ fini), soit que $(W,C)$ satisfait l'axiome transfini Cof$_4$, qu'on prouve que $(W(I),C_I)$ est également une structure à cofibrations sur $\M(I)$.

\msp{page 38}
Un fait important pour la construction des foncteurs $u^{W}_{!}$, pour $u : I\longrightarrow I'$ une flèche entre catégories ordonnées noethériennes, est la relation
\[
C_I\cap W(I) \ = \ (TC)_I \ , \tag*{(3.4.1)}
\]
d'où résulte que l'on a, pour $u$ comme dessus,
\[
u_!(C_I\cap W(I)) \ \subset\ C_{I'}\cap W(I') \ , \tag*{(3.4.2)}
\]
ingrédient essentiel de la construction du foncteur $u^{W}_{!}$. Pour étendre cette construction au cas d'une flèche $\pm$ quelconque dans Cat, il se pose donc le problème d'étendre la définition de $C_I$, et les formules (3.4.1) et (3.4.2), au cas de catégories $\pm$ quelconques dans Cat. Or on a prouvé :

\textbf{Proposition 3.5.} \textit{Si $\Phi=\Psi^*$ est une partie \textup{Q-close} à gauche de \textup{$\Fl(\M)$}, et $I$ un ensemble ordonné noethérien, $\Phi_I$ l'extension de $\Phi$ en une partie de \textup{$\Fl(\M(I))$} définie par \textup{\textsc{Anderson}}, alors on a}
\[
\Phi_I \ = \ \Psi(I)^* \ = \ v_!(\Phi(I_0))\Qs \ , \tag*{(3.5.1)}
\]
\textit{où $v : I_0\longrightarrow I$ est le foncteur canonique de la catégorie discrète définie par \textup{$\mathrm{Ob}\,I$}.}

\printedpage{22}
% head: A. Grothendieck : Dérivateurs, Ch. XIX, Édition des Universités de Montpellier II et Paris VII
% foot: Version 11 février 2019

C'est (avec des notations différentes) le lemme 6.2.5 (p. 107) de XVII, dont je ne reprends pas ici la démonstration.

Ceci donne la clef d'une définition de $\Phi_I$ comme $\Phi\langle I\rangle$ (avec les notations des paragraphes 1, 2),
\msp{page 39}
pour $I$ quelconque dans Cat,
\[
\Phi_I \ \deq\ \Phi\langle I\rangle \ \deq\ \Psi(I)^* \ = \ v_!(\Phi(I_0))\Qs \ , \tag*{(3.5.2)}
\]
où $\Psi=\Phi_*$, et où on suppose $\Phi$ Q-clos à gauche, i.e. $\Phi=\Psi^*$. Dans le cas où on suppose seulement $\Phi$ Q-saturé transfini sans plus, on aurait envie de poser
\[
\Phi_I \ \deq\ \text{Q-saturé transfini de } v_!(\Phi(I_0)) \ , \tag*{(3.5.3)}
\]
et je présume que pour $I$ ordonné noethérien, cela coïncide avec la définition d'\textsc{Anderson}. Mais je vais me borner pour la suite au cas où $C=\Phi$ est Q-clos à gauche.

Revenons alors au cas d'une structure à cofibrations $(W,C)$ dans $\M$, en supposant $C$ et $TC=C\cap W$ Q-clos à gauche. Définissons pour toute $I$ dans Cat $C_I=C\langle I\rangle$ comme dans (3.5.2). On a alors, pour $u : I\longrightarrow I'$ une flèche dans Cat,
\[
u_!(C\langle I\rangle) \ \subset\ C\langle I'\rangle \ , \tag*{(3.5.4)}
\]
mais on n'y a guère de raison d'espérer avoir aussi (3.4.2), sans disposer de la formule
\[
C\langle I\rangle\cap W(I) \ \overset{?}{=}\ (TC)\langle I\rangle \qquad\qquad \text{\footnote[21]{\textbf{N.B.} On a l'inclusion $\supset$, car $v^*v_!TC(I_0)\subset TC(I_0)$.}}. \tag*{(3.5.5)}
\]
On a prouvé que cette formule est valable si $I$ est ordonné noethérien (XVII th. 2.1, p. 26-27). Cela
\msp{page 40}
nous amène à la

\textbf{Question 3.6.} \textit{Soit $\M$ une catégorie stable par petites limites, et soit $(W,C)$ une structure à cofibrations dans $\M$ telle que $C$ et $TC=C\cap W$ soient Q-clos à gauche. Soit $I$ un objet de \textup{Cat}, a-t-on alors}
\[
C\langle I\rangle\cap W(I) \ = \ (TC)\langle I\rangle \text{ ?} \tag*{(3.6.1)}
\]

On a vu au paragraphe précédent qu'il en est ainsi si $(W,C)$ s'insère dans un triple de Quillen clos $(W,C,F=(TC)_*)$, et quand on suppose de plus que les couples
\[
C\langle I\rangle \ \longleftrightarrow\ TF(I) \ , \qquad TC\langle I\rangle \ \longleftrightarrow\ F(I)
\]
sont factorisant, auquel cas $(W(I),C\langle I\rangle,F(I))$ est une structure de Quillen close dans $\M(I)$. D'ailleurs, $(W,C,F)$ est une structure de Quillen (nécessairement close) si et seulement si les couples $C\longleftrightarrow \underbrace{TF}_{\deq C_*}$, $TC\longleftrightarrow \underbrace{\ F\ }_{\deq (TC)_*}$ eux-mêmes sont factorisants, et si % [sic] factorisant -> factorisants, agreement with « les couples »
\printedpage{23}
% head: A. Grothendieck : Dérivateurs, Ch. XIX, Édition des Universités de Montpellier II et Paris VII
% foot: Version 11 février 2019
de plus $TF\subset W$ (ce qui est automatique si $W=W^{\natural}$, car on aura en tous cas, par Cof$_4$, $TF\subset W^{\natural}$).

Mais peut-on s'en tirer avec moins, pour pouvoir conclure à (3.6.1) ? Ne suffirait-il pas de savoir que $TC\langle I\rangle\longleftrightarrow F(I)$ est factorisant (ce qui n'a guère de chances d'être vrai pour une $I$ $\pm$ générale, si on n'a déjà $TC\fact F$, factorisation dans $\M$ lui-même). C'est bien le cas :

\msp{page 41}
\textbf{Proposition 3.7.} \textit{Sous les conditions générales de \textup{(3.6)} sur $\M$, $W$, $C$, $I$, supposons que le couple $TC\langle I\rangle\longleftrightarrow F(I)$ (où $F\deq (TC)_*$) soit factorisant, i.e. $TC\langle I\rangle\fact F(I)$. Alors la formule \textup{(3.6.1)} est vraie.}

En effet, si $f\in C\langle I\rangle\cap W(I)$, on factorise $f$ en $pi$, avec $i\in TC\langle I\rangle$, $p\in F(I)$. Comme $f,i\in W(I)$, on aura $p\in W(I)$, i.e. $p\in TF(I)$, et à cause de $\underbrace{C\langle I\rangle}_{\ni\, f}\harp\underbrace{TF(I)}_{\ni\, p}$, il existe $p : Y\longrightarrow X'$ telle que $pq=\id_Y$, $qf=i$, % reading resolved: DER_XIX_corrections.csv; % [sic] p : Y -> q : Y, the arrow Y -> X' is q (see the diagram and pq = id_Y)
\[
\begin{tikzcd}[column sep=small, row sep=large]
X \arrow[rr, "f\in C\langle I\rangle\cap W(I)"] \arrow[dr, "i\in TC\langle I\rangle"'] & & Y \arrow[dl, dotted, shift right=1ex, "q"'] \\
& X' \arrow[ur, "p\in F(I)"'] \ , &
\end{tikzcd}
\]
donc $f$ est facteur direct de $i\in TC\langle I\rangle$. Comme $TC\langle I\rangle$ est stable par facteurs directs, on gagne.

\textbf{Corollaire 3.8.} \textit{Soit $(W,C,F)$ une structure de Quillen close dans $\M$, et supposons que $TC=C\cap W$ soit une partie Q-spéciale de $\M$, i.e. le couple $TC\fact F$ un couple spécial (i.e. $\M$ est \textup{spéciale}, et il existe une petite partie $TC_0$ Q-génératrice à gauche pour $TC$). Alors pour toute $I$ dans \textup{Cat}, on a}
\[
C\langle I\rangle\cap W(I) \ = \ TC\langle I\rangle \ . \tag*{(3.8.1)}
\]

\msp{page 42}
En apparence, on n'a pas eu à supposer que $C$ soit elle aussi une partie Q-spéciale de $\Fl(\M)$. Mais comme on a besoin de la factorisation pour $C\longleftrightarrow TF$ (pour avoir au moins une structure de cofibrations sur $\M$), on n'a guère le moyen de l'établir sans supposer $C$ spéciale. Sans compter que pour pouvoir expliciter (3.8.1) pour la construction de $u^{W}_{!}$, on a besoin de savoir que $(W(I),C\langle I\rangle)$ définit bien une structure à cofibrations sur $\M(I)$, ce qui ne peut guère se faire (il me semble) sans avoir la factorisation pour $C\langle I\rangle\longleftrightarrow TF(I)$ également ; chose que je n'arrive à établir pour un $I$ général dans Cat, que moyennant l'hypothèse de spécialité sur $C$ également.

\textbf{3.9.} Pour en avoir le cœur net, il faudrait quelques exemples concrets et illuminants, du type suivant :

a) \texttt{[Il n'y a pas de b).]} Structure de Quillen close $(W,C,F)$ sur $\M$, telle que pour le prédérivateur associé, $u^{W}_{!}$ n'existe pas toujours, ou $u^{W}_{*}$ n'existe par toujours. Je pense notamment au cas d'une flèche $u : I\longrightarrow I'$ avec $I'=e$ (présumant que si dans ce cas, $u^{W}_{!}$ resp. $u^{W}_{*}$ existe, alors il existe dans tous les cas). Un tel exemple serait encore plus % [sic] n'existe par toujours -> n'existe pas toujours, typing error


\endgroup
% END UNIT der-ch19/der-ch19_p021-025.tex
% BEGIN UNIT der-ch19/der-ch19_p026-028.tex
\begingroup\makeatletter%
\newcommand{\printedpage}[1]{\marginpar{\footnotesize p.~#1}}%
\newcommand{\msp}[1]{\marginpar{\footnotesize [#1]}}%
\newcommand{\limind}{\varinjlim}%
\newcommand{\limproj}{\varprojlim}%
\newcommand{\Fl}{\mathrm{Fl}}%
\newcommand{\Is}{\mathrm{Is}}%
\newcommand{\Ho}{\mathrm{Ho}}%
\newcommand{\qpair}{\mathrel{\leftarrow\mkern-4mu\mapstochar\mkern-2mu\rightarrow}}%
\newcommand{\Qcouple}{\mathrel{\leftarrow\mkern-9mu{+}\mkern-9mu\rightarrow}}%
\makeatother\let\printedpage\derpp\let\msp\dermsp%
\let\mspin\DERmspin%


\printedpage{24}
% head: A. Grothendieck : Dérivateurs, Ch. XIX, Édition des Universités de Montpellier II et Paris VII
% foot: Version 11 février 2019
\msp{page 43}
instructif, bien sûr, si $\mathcal{M}$ est stable par petites $\limind$ , que ses objets soient accessibles etc. et que $W$ soit L-accessible (i.e. stable par petites $\limind$ filtrantes, et de plus accessible \dots). Le premier candidat pour $I$ auquel je pense serait un $B_G$ ($G$ un groupe, p.ex. $G=\pm1$), c'est en effet l'exemple aux antipodes des ensembles ordonnés, lesquels, eux, ne semblent pas faire trop de problème (moyennant l'hypothèse noethérienne tout au moins).

Les exemples suivants ne sont pas aussi convaincants que ça, hélas!

\textbf{Exemple 3.10.} Soit $\mathcal{M}$ une catégorie quelconque, alors les couples $(\Fl(\mathcal{M}),\Is(\mathcal{M}))$ et $(\Is(\mathcal{M}),\Fl(\mathcal{M}))$ sont tous deux des couples de Quillen,
\[
\Fl(\mathcal{M}) \Qcouple \Is(\mathcal{M}) \ , \qquad \Is(\mathcal{M}) \Qcouple \Fl(\mathcal{M}) \ . \tag*{(3.10.1)} % reading resolved: DER_XIX_corrections.csv
\]
Il s'ensuit que si on pose
\[
W \ = \ \Is(\mathcal{M}) \ , \qquad C \ = \ F \ = \ \Fl(\mathcal{M}) \ , \tag*{(3.10.2)}
\]
$(W,C,F)$ est un triple de Quillen clos \textsuperscript{(}\footnote[22]{au moins si on admet que $\mathcal{M}$ est stable par $\limind$ et $\limproj$ \textit{finies.}}\textsuperscript{)}. D'ailleurs pour toute $I$ dans Cat, on a $W(I)=\Is(\mathcal{M}(I))$, donc
\[
\Ho_W(I) \ = \ \mathcal{M}(I) \ . \tag*{(3.10.3)}
\]
Donc la question de l'existence des foncteurs $u_!^W$, $u_*^W$ se réduit à celle des foncteurs $u_!$, $u_*$ habituels, adjoints à gauche et à droite de $u^* : \mathcal{M}(I') \longrightarrow \mathcal{M}(I)$.

\msp{page 44}
On sait bien que l'existence de ces foncteurs pour $u$ donné, fait appel à l'existence de certaines $\limind$ resp. $\limproj$ dans $\mathcal{M}$; p.ex. si $I'=e$, il est nécessaire et suffisant que les $\limind$ (resp. $\limproj$ ) de type $I^\circ$ soient représentables dans $\mathcal{M}$. Mais avec les hypothèses faites sur $\mathcal{M}$, il n'y a pas de problème si $u : I \longrightarrow I'$ est un foncteur entre catégories \textit{finies.} On va donner un exemple où il n'en est plus ainsi.

\textbf{Exemple 3.11.} Soit $k$ un anneau commutatif, et $\mathcal{M}=K^{\mathrm{b}}(k\text{-Ab})$ la catégorie des complexes de cochaînes bornés de $k$-modules, $W$ l'ensemble des quasi-isomorphismes, donc
\[
\Ho_W \ = \ D^{\mathrm{b}}(k\text{-Ab}) \ ,
\]
catégorie dérivée bornée de la catégorie abélienne des $k$-modules. On sait, en utilisant l'existence d'assez de projectifs dans $k$-Ab, que cette structure est associée à une structure de Quillen close $(W,C,F)$, décrite par \textsc{Quillen}
\msp{page 45}
dans Homotopical Algebra -- à cela près qu'il prend $D^-(k)$ au lieu de $D^{\mathrm{b}}(k)$. Peut-être, pour vérifier les axiomes, doit-on supposer que $k$ est de dimension cohomologique finie.

Prenons $I=B_G$, où $G$ est un groupe (p.ex. $G=\{\pm1\}$), et $u : I \longrightarrow e$. Alors $\mathcal{M}(I)$ est la catégorie des complexes bornés de $k[G]$-modules. On voit que $u_!^W$ resp. $u_*^W$ existe si et seulement si pour tout tel complexe borné, ou simplement pour un $k[G]$-module sans plus $M$, le complexe $L\Gamma^G(M)$ resp. $R\Gamma_G(M)$ appartient à $D^{\mathrm{b}}(k)$, ou ce qui revient au même,
\printedpage{25}
% head: A. Grothendieck : Dérivateurs, Ch. XIX, Édition des Universités de Montpellier II et Paris VII
% foot: Version 11 février 2019
si et seulement si on a
\[
\begin{aligned}
H_i(G,M) \ &= \ 0 \quad \text{pour } i \gg 0 \ ,\\
\text{resp.}\quad H^i(G,M) \ &= \ 0 \quad \text{pour } i \gg 0 \ ,
\end{aligned}
\]
pour tout $k[G]$-module $M$. Ce n'est pratiquement jamais le cas; quand on prend p.ex. $k=\mathbf{Z}$ et $G$ groupe fini, cela n'est le cas que si $G=e$. L'exemple le plus économique me paraît $G=\pm1$, $k=\mathbf{F}_2$.

\msp{page 46}
Donc ici $u_!^W$ et $u_*^W$ n'existent pas, pour $I \longrightarrow e$ avec $I$ catégorie finie. Mais aussi $\mathcal{M}$ n'est pas stable par petites $\limind$ .

\textbf{Exemple 3.12.} Soit $\mathcal{M}$ une catégorie muni d'un localiseur $W$ fortement saturé, tel que $\mathcal{M} \longrightarrow \underbrace{\mathcal{M}W^{-1}}_{=\mathcal{H}}$ admette un adjoint -- à gauche disons. On a donc une situation de deux foncteurs adjoints % [sic] muni -> munie, agreement with « une catégorie »
\[
\begin{tikzcd}[column sep=huge]
\mathcal{H} \arrow[r, shift left, "v_!\ \text{pleinement fidèle}"] & \mathcal{M} \arrow[l, shift left, "v^*\ \text{localisation}"]
\end{tikzcd}
\ , \tag*{(3.12.1)}
\]
et $W=(v^*)^{-1}(\mathrm{Iso}(\mathcal{H}))$. On trouve que pour toute $I$ dans Cat, on a
\[
\mathrm{HO}_W(I) \ \simeq \ \mathcal{H}(I) \ \overset{\text{déf}}{=} \ \mathrm{Hom}(I^\circ,\mathcal{H}). \tag*{(3.12.2)} % [sic] HO_W -> Ho_W, notation Ho_W used in (3.10.3) and Exemple 3.11
\]
On peut donc songer à faire un exemple où dans $\mathcal{M}$ les petites $\limind$ et $\limproj$ existent, mais non dans $\mathcal{H}$, de sorte que les $u_!^W$ ou $u_*^W$ n'existent pas toujours. Mais ce semble difficile à construire, vu que dans la situation envisagé, on voit facilement que les petites $\limproj$ existent dans $\mathcal{H}$. Donc si dans $\mathcal{H}$ % [sic] envisagé -> envisagée, agreement with « la situation »
\msp{page 47}
il existe une petite famille cogénératrice (et pour ceci, il suffit qu'une telle famille existe dans $\mathcal{M}$), il y aura également des petites $\limind$ , en vertu de SGA 4 I 8.12.7. De même, elles existent si $\mathcal{H}$ est une sous-catégorie pleine de $\mathcal{M}$ \textit{stable par petites} $\limind$ , et il paraît pratiquement impossible de faire un exemple (3.12.1) sans que $\mathcal{H}$ ait cette propriété de stabilité. Sans compter qu'il y a même la perplexité suivante :

\marginpar{! \textbar}\textbf{Question 3.13.} \textit{Existe-t-il une $\mathfrak{U}$-catégorie (\,$\mathfrak{U}$ un univers donné) où les petites $\limind$ soient représentables, sans qu'il en soit de même des petites $\limproj$ ?} \textbf{N.B.} Une telle catégorie n'admet pas de petite famille génératrice, a fortiori elle n'est pas équivalente à une petite catégorie.

\msp{page 48}
\textbf{Question 3.14.} \textsc{Anderson} prétend que pour le prédérivateur défini par une catégorie à cofibrations satisfaisant l'axiome Cof$_4$ transfini, les foncteurs $u_!^W$ existent. Sa démonstration est totalement canulée, néanmoins je n'arrive pas à faire de contre-exemple. Le résultat qu'il annonce serait-il vrai? Je ne connais pas même d'exemple d'un localiseur $W$, dans une catégorie $\mathcal{M}$ \textit{avec petites $\limind$ et petites $\limproj$ ,} tel que $u_!^W$ ou $u_*^W$ n'existent pas toujours -- mais je suis persuadé qu'il y a des exemples (même avec $\mathcal{M}$ une petite catégorie?). J'en suis moins convaincu si on suppose $W$ stable par $\limind$ filtrantes, et au besoin, de plus, que les éléments de $\mathcal{M}$ sont accessibles (plus éventuellement d'autres propriétés catégoriques anodines sur $\mathcal{M}$, telles l'existence d'une filtration
\printedpage{26}
% head: A. Grothendieck : Dérivateurs, Ch. XIX, Édition des Universités de Montpellier II et Paris VII
% foot: Version 11 février 2019
\msp{page 49}
cardinale etc.)

\textbf{Question 3.15.} Cela me ramène à la question particulièrement élusive et irritante, de l'existence des petits ensembles Q-générateurs, dans une partie $\Phi$ de $\Fl(\mathcal{M})$, quand on suppose (disons) que $\mathcal{M}$ est stable par petites $\limproj$ , que ses éléments sont accessibles et qu'il admet une filtration cardinale. Ainsi on a la question : une partie L-accessible de $\Fl(\mathcal{M})$, qu'on suppose au besoin Q-saturé à gauche (ou ce qui revient au même, contenant les isomorphismes, et stable par cochangement de base, par composition et par facteurs directs) admet-elle une \textit{petite partie} $\Phi_0$ Q-génératrice à gauche? % [sic] Q-saturé -> Q-saturée, agreement with « une partie »

\marginpar{! \textbar}D'autre part, je n'ai pas à présent construit un seul exemple d'un couple $\Phi \longleftrightarrow \Psi$, dans une catégorie $\mathcal{M}$, qui ne soit factorisant! Le cas où $\mathcal{M}$ est spéciale, était bien sûr le plus intéressant. Visiblement, dans tout
\msp{page 50}
ça ce sont les exemples qui manquent. Le clef d'une compréhension semble bien être dans les questions de factorisation entre ensembles $\Phi \longleftrightarrow \Psi$. Il faudrait en faire une étude systématique, mais je recule devant la perspective d'une digression qui me prendrait peut-être encore des semaines! % [sic] Le clef -> La clef, « clef » is feminine


\endgroup
% END UNIT der-ch19/der-ch19_p026-028.tex
\end{document}
